% Release engine: XeLaTeX (the LuaTeX compatibility layer breaks Xy-pic output).
\documentclass[12pt]{../fichiers-configuration-latex/amsbook}
\usepackage[a4paper,left=3cm,right=3cm,marginpar=0.5cm,marginparsep=0.8cm,vmargin=2.8cm]{geometry}
% BEGIN INCLUDED SGA4_English_current_buildable_source_20260906/source/fichiers-configuration-latex/paquets.tex
\usepackage{polyglossia}
%\usepackage[francais]{babel}
\usepackage{fontspec}
\usepackage[math-style=ISO,bold-style=upright,sans-style=upright]{unicode-math}
%\usepackage[mathletters]{ucs}
%\usepackage[utf8]{inputenc}
%\usepackage[T1]{fontenc}
\usepackage[unicode,hypertexnames=false,colorlinks=true,linkcolor=blue,citecolor=Sepia,urlcolor=magenta]{hyperref}
\setmainlanguage{french}
\setmainfont[Ligatures=TeX]{Linux Libertine O}
% variante possible : TeXGyrePagella ?
\setmathfont{XITS Math}

% maths
\usepackage{amsmath}
%\usepackage{amssymb}
%\usepackage{mathrsfs}
%\usepackage{stmaryrd}
%\usepackage{euscript}

% couleurs
\usepackage[usenames,dvipsnames]{xcolor}

% graphiques
\usepackage{graphicx}
\usepackage[all]{xy}
%\usepackage{tikz}
%\usetikzlibrary{matrix,arrows,calc}
%\tikzset{column sep=3em,row sep=2em,text height=1.5ex, text depth=0.25ex,>=angle 90}
%\tikzset{description/.style={fill=white,inner sep=1pt}}

% hyperliens, 'etiquettes
\usepackage{xr-hyper}
\usepackage{../fichiers-configuration-latex/smfthm}
%\ifx\cacheretiquettes\undefined
%\usepackage[inline]{showlabels} \else \fi % pour rendre visibles les 'etiquettes
%\usepackage{../fichiers-configuration-latex/makeidx} % pour faire un index
%\usepackage[pagebackref,unicode]{hyperref}  % pour voir o`u sont cit'ees les r'ef'erences
% hyperref is loaded once above with the cumulative-reader destination policy.
%\hypersetup{breaklinks=true} %permet le retour `a la ligne dans les liens trop longs. Incompatible avec pdflatex
%\renewcommand*{\backref}[1]{{\small [cit\'e page(s)~#1.]}}

\usepackage{xspace}
\usepackage{pifont} % pour expose XVII

% END INCLUDED SGA4_English_current_buildable_source_20260906/source/fichiers-configuration-latex/paquets.tex

% BEGIN INCLUDED SGA4_English_current_buildable_source_20260906/source/fichiers-configuration-latex/numerotation.tex
\renewcommand{\thefootnote}{\roman{footnote}} % pour mettre les notes en bas de page sous la forme (ii) et non (2), qui ressemble trop (?) a un exposant.

\newcounter{pageoriginal}
\newcommand{\pageoriginale}{\refstepcounter{pageoriginal}\marginpar{\footnotesize\textbf{\thepageoriginal}}\xspace}
\let\pageoriginaled\pageoriginale

\usepackage{ifthen}
\newboolean{original}
\def\sisi{\ifthenelse{\boolean{original}}}
\setboolean{original}{false}

\newcounter{cnote}
\def\nde#1{\ifthenelse{\boolean{original}}{}{\stepcounter{cnote}\def\thefootnote{\arabic{cnote}}\protect\footnote[\value{cnote}]{N.D.E.: #1}}}
\def\ndemark{\ifthenelse{\boolean{original}}{}{\stepcounter{cnote}\def\thefootnote{\arabic{cnote}}\protect\footnotemark[\value{cnote}]}}
\def\ndetext#1{\ifthenelse{\boolean{original}}{}{\def\thefootnote{\arabic{cnote}}\protect\footnotetext[\value{cnote}]{N.D.E.: #1}}}

% END INCLUDED SGA4_English_current_buildable_source_20260906/source/fichiers-configuration-latex/numerotation.tex

% BEGIN INCLUDED SGA4_English_current_buildable_source_20260906/source/fichiers-configuration-latex/macros.tex
\setmainlanguage{french}
\setotherlanguage{english}


%% Fixed, engine-independent release identifier.  The inherited dynamic Git
%% probe used Lua io.popen, exposed build-state-dependent output, and printed
%% a visible question mark outside a Git checkout.
\def\showgitstatus{English release r2, 2026-07-23}


% macros pour LaTeX
\renewcommand{\mathsf}{}

% Police différente pour U+2113 SCRIPT SMALL L
\setmathfont[range="2113]{latinmodern-math.otf}

% Give every inherited semantic label its own stable PDF destination instead
% of inheriting whichever automatic equation/page anchor happened to be live.
% The legacy label remains as a compatibility alias used by existing edges.
\makeatletter
\def\etiquette[#1]#2{%
  \Hy@raisedlink{\hyper@anchorstart{sga4.semantic.#2}\hyper@anchorend}%
  \def\@currentHref{sga4.semantic.#2}%
  \def\@currentlabel{#1}%
  \label{#2}%
}
\makeatother
\def\eqhyperref[#1]#2{(\hyperref[#1]{#2})}
\def\pagehyperref[#1]#2{(page \hyperref[#1]{#2})}
\def\hyperrefext[#1][#2]#3{\hyperref[#2]{#1 #3}}
\def\eqhyperrefext[#1][#2]#3{\hyperref[#2]{#1 #3}}     %cf. exposes 17
\def\neweqhyperrefext[#1][#2]#3{\hyperref[#2]{#1 #3}}  %cf. exposes 17

\def\shat#1{\protect\widehat{#1}}
\def\stilde#1{\protect\widetilde{#1}}
%%%%%%%%%%%%%%%%%%


\def\etcn{etc.\xspace}
\def\etco{etc...\xspace}
\def\etc{\sisi{\etco}{\etcn}}
\def\ldotsn{\sisi{...}{\etcn}}

\def\cf{{cf.}\kern.3em}
\def\Cf{{Cf.}\kern.3em}
\def\ie{{i.e.}\kern.3em}
\def\iev{{i.e.},\xspace}
\def\loccito{loc.\kern2pt cit.\xspace}
\def\loccitn{\textit{loc.\kern2pt cit.}\xspace}
\def\loccit{\sisi{\loccito}{\loccitn}}
\def\EGA{\textup{EGA}\kern.2em}
\def\SGA{\textup{SGA}\kern.3em}

\def\quoi{}

\def\quav{\quad,\quad}
\def\quaet{\quad \text{et}\quad}
\def\fibreproduct#1#2#3{#1 ×_{#3} #2}
\let\fprod\fibreproduct
\def\ccheck#1#2{\mbox{${\check{\check{\rm #1}}}^{\lower2pt\hbox{$\scriptstyle #2$}}$}}
\def\sheaf#1{\protect\mathcal{#1}}
\def\vieuxsheaf#1{{#1}} % cf. problème avec sheaf{a\Phi_i} par exemple...
\DeclareMathOperator{\Hom}{Hom}
\DeclareMathOperator{\can}{can}
\DeclareMathOperator{\Ind}{Ind}
\DeclareMathOperator{\ind}{ind}
%\DeclareMathOperator{\Hhomtop}{\mathcal{H}\!\mathit{omtop}}
\DeclareMathOperator{\Hhomtop}{\mathbf{Homtop}}
\DeclareMathOperator{\Point}{Point}
\DeclareMathOperator{\Pro}{Pro}
\DeclareMathOperator{\Top}{Top}
\DeclareMathOperator{\Mod}{Mod}
\DeclareMathOperator{\Cart}{Cart}
%\DeclareMathOperator{\Mmorsite}{\mathcal{M}\!\mathit{orsite}}
\DeclareMathOperator{\Mmorsite}{\mathbf{Morsite}}
\DeclareMathOperator{\loc}{loc}
\DeclareMathOperator{\Loc}{Loc}
\DeclareMathOperator{\Pic}{Pic}
\DeclareMathOperator{\Et}{Et}
\DeclareMathOperator{\card}{card}
\DeclareMathOperator{\Ens}{Ens}
\DeclareMathOperator{\h}{\underline{H}}
\DeclareMathOperator{\Fonct}{Fonct}
\DeclareMathOperator{\Ker}{Ker}
\DeclareMathOperator{\but}{but}
\DeclareMathOperator{\Ob}{Ob}
\DeclareMathOperator{\ob}{ob}
\DeclareMathOperator{\Lib}{Lib}
\DeclareMathOperator{\esj}{esj}
\DeclareMathOperator{\id}{id}
\DeclareMathOperator{\Id}{Id}
\DeclareMathOperator{\ad}{ad}
\DeclareMathOperator{\Iim}{Im}
\DeclareMathOperator{\pr}{pr}
\DeclareMathOperator{\Cat}{Cat}
\DeclareMathOperator{\car}{car}
\DeclareMathOperator{\SysInd}{SysInd}
\DeclareMathOperator{\Sup}{Sup}
\DeclareMathOperator{\Coker}{Coker}
\DeclareMathOperator{\Filt}{Filt}
\DeclareMathOperator{\acc}{acc}
\DeclareMathOperator{\sur}{sur}
\DeclareMathOperator{\cat}{cat}
\DeclareMathOperator{\Ord}{Ord}
\DeclareMathOperator{\Gr}{Gr}
\DeclareMathOperator{\Cov}{Cov}
\DeclareMathOperator{\Caf}{Caf}
\DeclareMathOperator{\coker}{coker}
\DeclareMathOperator{\coim}{coim}
\DeclareMathOperator{\im}{im}
\DeclareMathOperator{\Ouv}{Ouv}
\DeclareMathOperator{\diag}{diag}
\DeclareMathOperator{\Esp}{Esp}
\DeclareMathOperator{\TOP}{TOP}
\DeclareMathOperator{\Unif}{Unif}
\DeclareMathOperator{\Iind}{\mathbf{Ind}}
\DeclareMathOperator{\Hhomcart}{\mathbf{Homcart}}
\DeclareMathOperator{\Hhomtopess}{\mathbf{Homtopess}}
\DeclareMathOperator{\Ppoints}{\mathbf{Points}}
\DeclareMathOperator{\Ppoint}{\mathbf{Point}}
\DeclareMathOperator{\Ttors}{\mathbf{Tors}}
\DeclareMathOperator{\Ssex}{\mathbf{Sex}}
\DeclareMathOperator{\Ffib}{\mathbf{Fib}}
\DeclareMathOperator{\Vvois}{\mathbf{Vois}}
\DeclareMathOperator{\Ppointess}{\mathbf{Pointess}}
\DeclareMathOperator{\scrHom}{\mathbf{Hom}}
\DeclareMathOperator{\sheafPic}{\mathbf{Pic}}
\DeclareMathOperator{\sheafAut}{\mathbf{Aut}}
\DeclareMathOperator{\Hhomtoplocan}{\mathbf{Homtoplocan}}
\DeclareMathOperator{\Hhomtopan}{\mathbf{Homtopan}}
\DeclareMathOperator{\Pptgeom}{\mathbf{Ptgéom}}
\DeclareMathOperator{\Morsite}{Morsite}
\DeclareMathOperator{\Fib}{Fib}
\DeclareMathOperator{\Fais}{Fais}
\DeclareMathOperator{\kar}{kar}
\DeclareMathOperator{\Groupes}{Groupes}
\DeclareMathOperator{\Restr}{Restr}
\DeclareMathOperator{\Isom}{Isom}
\DeclareMathOperator{\Prol}{Prol}
\DeclareMathOperator{\Res}{Res}
\DeclareMathOperator{\Tors}{Tors}
\DeclareMathOperator{\esp}{esp}
\DeclareMathOperator{\es}{es}
\DeclareMathOperator{\od}{od}
\DeclareMathOperator{\gr}{gr}
\DeclareMathOperator{\sob}{sob}
\DeclareMathOperator{\dfn}{dfn}
\DeclareMathOperator{\Sch}{Sch}
\DeclareMathOperator{\ouv}{ouv}
\DeclareMathOperator{\cons}{cons}
\DeclareMathOperator{\Nerf}{Nerf}
\DeclareMathOperator{\Points}{Points}
\DeclareMathOperator{\Espsob}{Espsob}
\DeclareMathOperator{\Homtop}{Homtop}
\DeclareMathOperator{\Pt}{Pt}
\DeclareMathOperator{\Quot}{Quot}
\DeclareMathOperator{\glob}{glob}
\DeclareMathOperator{\Homfib}{Homfib}
\DeclareMathOperator{\Qouv}{Qouv}
\DeclareMathOperator{\Quotouv}{Quotouv}
\DeclareMathOperator{\ens}{ens}
\DeclareMathOperator{\ab}{ab}
\DeclareMathOperator{\an}{an}
\DeclareMathOperator{\fs}{fs}
\DeclareMathOperator{\red}{red}
\DeclareMathOperator{\Ab}{Ab}
\DeclareMathOperator{\cosq}{cosq}
\DeclareMathOperator{\Ann}{Ann}
\DeclareMathOperator{\Zar}{Zar}


%%%%%%%%%%%%
%%%%% SGA4-2
\DeclareMathOperator{\Sex}{Sex}
\DeclareMathOperator{\vois}{vois}
\DeclareMathOperator{\Vois}{Vois}
\DeclareMathOperator{\Ext}{Ext}
\DeclareMathOperator{\ext}{ext}
\DeclareMathOperator{\Fer}{Fer}
\DeclareMathOperator{\Add}{Add}
\DeclareMathOperator{\Cocart}{Cocart}
\DeclareMathOperator{\cosk}{cosk}
\DeclareMathOperator{\plat}{plat}

\DeclareMathOperator{\faisc}{faisc}
\DeclareMathOperator{\sk}{sk}
\DeclareMathOperator{\ch}{ch}

\def\HH{\mathbf{H}}
\def\qc{\mathrm{qc}}
\def\et{\mathrm{et}}
\def\coh{\mathrm{coh}}

\def\Hvr#1{%
\begin{array}[b]{l}
{\kern2pt\scriptstyle r}\\[-8pt]
{\kern1pt\scriptstyle \vee}\\[-5pt]
H^{#1}
\end{array}}

\def\Hvv#1{%
\begin{array}[b]{l}
\kern2pt\rotatebox{90}{${\scriptscriptstyle\ll}$}\\[-6pt]
H^{#1}
\end{array}}

\def\Hvi#1{%
\begin{array}[b]{l}
{\kern0.2pt\scriptscriptstyle \infty}\\[-8pt]
{\kern1pt\scriptstyle \vee}\\[-5pt]
H^{#1}
\end{array}}

\def\Hv#1{%
\begin{array}[b]{l}
{\kern1pt\scriptstyle \vee}\\[-5pt]
H^{#1}
\end{array}}

\DeclareMathOperator{\cart}{cart}
\DeclareMathOperator{\Ffaisc}{\mathbf{Faisc}}


\def\Ev#1{%
\begin{array}[b]{l}
{\kern1pt\scriptstyle \vee}\\[-5pt]
E'_{#1}
\end{array}}

\DeclareMathOperator{\Cont}{Cont}
\DeclareMathOperator{\Ccont}{\mathscr{C}\!\mathit{ont}}
\DeclareMathOperator{\Site}{Site}
\newcommand\MM{\mathbf{M}}
\def\limtop#1{{\displaystyle{\varprojlim_{#1}\!\mbox{\rm top }}}}
\def\Limtop{\mathop{\varprojlim\!\mbox{\rm top}}}
\DeclareMathOperator{\spec}{spec}
\def\ss{\scriptstyle{\mathscr{S}}}
\DeclareMathOperator{\Gal}{Gal}
\def\smallm{\scriptstyle{\mathscr{M}}}
\DeclareMathOperator{\End}{End}
\DeclareMathOperator{\ass}{ass}
\DeclareMathOperator{\zar}{zar}
\DeclareMathOperator{\Ogg}{Ogg}
\DeclareMathOperator{\res}{res}
\DeclareMathOperator{\tr}{tr}
\DeclareMathOperator{\Supp}{Supp}
\DeclareMathOperator{\codim}{codim}
\DeclareMathOperator{\rk}{rk}
\DeclareMathOperator{\cur}{cur}
\DeclareMathOperator{\cdqc}{cdqc}
\newcommand\DD{\mathbf{D}}
\DeclareMathOperator{\cl}{cl}
\def\JJ{\mathbf{J}}
\def\FF{\mathbf{F}}
\def\BB{\mathbf{B}}
\DeclareMathOperator{\Alget}{Alget}
\DeclareMathOperator{\Ass}{Ass}
\DeclareMathOperator{\supp}{supp}
\DeclareMathOperator{\rad}{rad}
\DeclareMathOperator{\prof}{prof}
\DeclareMathOperator{\cint}{int}
\DeclareMathOperator{\sic}{sic}
\DeclareMathOperator{\Tor}{Tor}
\DeclareMathOperator{\tors}{tors}
\DeclareMathOperator{\Sym}{Sym}
\DeclareMathOperator{\parf}{parf}
\DeclareMathOperator{\tor}{tor}
\DeclareMathOperator{\defn}{defn}
\DeclareMathOperator{\Div}{Div}
\DeclareMathOperator{\Pol}{Pol}
\DeclareMathOperator{\Fex}{Fex}
\DeclareMathOperator{\Lie}{Lie}
\DeclareMathOperator{\tot}{tot}
\DeclareMathOperator{\Aut}{Aut}
\DeclareMathOperator{\HOM}{HOM}
\DeclareMathOperator{\pch}{pch}
\DeclareMathOperator{\Ch}{Ch}
\DeclareMathOperator{\EXT}{EXT}
\DeclareMathOperator{\PIC}{PIC}
\DeclareMathOperator{\Rf}{Rf}
\DeclareMathOperator{\inj}{inj}
\newcommand\VV{\mathbf{V}}

\def\Xv{%
\begin{array}[b]{l}
{\kern1pt\scriptstyle \vee}\\[-5pt]
X
\end{array}}

%%%%%%%
\newcommand\TT{\mathbf{T}}
\newcommand\XX{\mathbf{X}}
\newcommand\EE{\mathbf{E}}
%%%%%%

\def\qedsymbol{\textup{C.Q.F.D.}}
\let\CQFD\qed
\let\cqfd\qed

%%%%%\bigdot
\font\bigsymb=cmsy10 at 4pt
\def\bigdot{{\kern1.2pt\raise 1.5pt\hbox{\bigsymb\char15}}}
%%%%%

\def\Lright{\mathop{\vcenter{\hbox{\hbox{$\longrightarrow$}\lower8pt \hbox{\kern-16pt\hbox{$\longrightarrow$}}}}}}

\def\Aa{{\mathcal A}}
\def\Bb{{\mathcal B}}
\def\Cc{{\mathcal C}}
\def\Dd{{\mathcal D}}
\def\Ee{{\mathcal E}}
\def\Ff{{\mathcal F}}
\def\Gg{{\mathcal G}}
\def\Hh{{\mathcal H}}
\def\Ii{{\mathcal I}}
\def\Jj{{\mathcal J}}
\def\Kk{{\mathcal K}}
\def\Ll{{\mathcal L}}
\def\Mm{{\mathcal M}}
\def\Nn{{\mathcal N}}
\def\Oo{{\mathcal O}}
\def\Pp{{\mathcal P}}
\def\Qq{{\mathcal Q}}
\def\Rr{{\mathcal R}}
\def\Ss{{\mathcal S}}
\def\Tt{{\mathcal T}}
\def\Uu{{\mathcal U}}
\def\Vv{{\mathcal V}}
\def\Ww{{\mathcal W}}
\def\Xx{{\mathcal X}}
\def\Yy{{\mathcal Y}}
\def\Zz{{\mathcal Z}}

\newcommand{\CC}{\mathbf{C}}
\newcommand{\GG}{\mathbf{G}}
\newcommand{\KK}{\mathbf{K}}
\newcommand{\LL}{\mathbf{L}}
\newcommand{\NN}{\mathbf{N}}
\newcommand{\PP}{\mathbf{P}}
\newcommand\QQ{\mathbf{Q}}
\newcommand\RR{\mathbf{R}}
\newcommand\ZZ{\mathbf{Z}}

\def\to{\mathchoice{\longrightarrow}{\rightarrow}{\rightarrow}{\rightarrow}}
\def\from{\mathchoice{\longleftarrow}{\leftarrow}{\leftarrow}{\leftarrow}}
\def\mto{\mathchoice{\longmapsto}{\mapsto}{\mapsto}{\mapsto}}
\def\mfrom{\mathrel{\from\joinrel\mapstochar\,}}
\def\hto{\mathrel{\lhook\joinrel\to}}
\def\hfrom{\mathrel{\from\joinrel\rhook}}

\def\implies{\mathchoice{\Longrightarrow}{\Rightarrow}{\Rightarrow}{\Rightarrow}}
\def\implied{\mathchoice{\Longleftarrow}{\Leftarrow}{\Leftarrow}{\Leftarrow}}
\def\iff{\mathchoice{\Longleftrightarrow}{\Leftrightarrow}{\Leftrightarrow}{\Leftrightarrow}}

\def\To#1{\mathchoice{\xrightarrow{\labelstyle~~#1~}}{\stackrel{#1}{\longrightarrow}}{}{}}
\def\Mto#1{\mathrel{\mapstochar\To{#1}}}
\def\From#1{\mathchoice{\xleftarrow{\labelstyle~~#1~}}{\stackrel{#1}{\longleftarrow}}{}{}}

\newcommand{\isomto}{\overset{\sim}{\longrightarrow}}
\newcommand{\isomfrom}{\overset{\sim}{\longleftarrow}}

\newcommand{\IFF}{\unskip~$\iff$~\ignorespaces}
\newcommand{\THEN}{\unskip~$\implies$~\ignorespaces}

\let\leq\leqslant
\let\geq\geqslant
\let\le\leqslant
\let\ge\geqslant

\def\cf{{cf.}\kern.3em}
\def\Cf{{Cf.}\kern.3em}
\def\ie{{i.e.}\kern.3em}
\def\iev{{i.e.},\xspace}
%\def\loccito{loc.\kern2pt cit.\xspace}
%\def\loccitn{\textit{loc.\kern2pt cit.}\xspace}
%\def\loccit{\sisi{\loccito}{\loccitn}}
\def\EGA{\textup{EGA}\kern.2em}
\def\SGA{\textup{SGA}\kern.3em}

\let\oldS\S
\def\S{\oldS\kern .15em}
\def\ptbl{.\kern .2em}
\def\resp{resp.\kern.3em}
\def\Exp{Exp.\kern.3em}

\let\bar\overline
\let\hat\widehat
\let\tilde\widetilde
\let\dpl\displaystyle
\let\ldots\dots
\let\emptyset\varnothing
\let\skpt\skippointrait
\let\epsilon\varepsilon

\def\sfrac#1#2{{#1}/{#2}}

\DeclareMathOperator{\cd}{cd}
%\DeclareMathOperator{\Hhom}{\mathcal{H}\!\mathit{om}}
\DeclareMathOperator{\Hhom}{\mathbf{Hom}}
%\DeclareMathOperator{\Eext}{\mathcal{E}\!\mathit{xt}}
\DeclareMathOperator{\Eext}{\mathbf{Ext}}
\DeclareMathOperator{\Spec}{Spec}
\DeclareMathOperator{\Tr}{Tr}

\let\emptyset=\varnothing

%%%%%%%%%%%%%%%
%%% expose XVII
%%%%%%%%%%%%%%%

\makeatletter
\let\old@optexpose\optexpose
\newcommand{\empty@optexpose}[1]{}
\newcommand{\paren@optexpose}[1]{\old@optexpose{#1}(}
\newcommand{\refnoexpose}[1]{\let\optexpose\empty@optexpose\ref{#1}\let\optexpose\old@optexpose}%\refnoexpose{X.cor5.2} engendre "5.2" au lieu de "X 5.2"
\newcommand{\neweqref}[1]{\let\optexpose\paren@optexpose\textup{\ref{#1})}\let\optexpose\old@optexpose}%\neweqref{XII.eq4.1} engendre "XII (4.1)" au lieu de "(XII 4.1)"
\makeatother
\newcommand{\refvb}[2][VI${}_{\text{B}}$ ]{\sisi{#1\refnoexpose{#2}}{\ref{#2}}}

\newcommand{\leftexp}[2]{{\vphantom{#2}}^{#1}{#2}}
\newcommand{\AAA}{\mathbf{A}}
\newcommand{\Sf}{\mathfrak{S}}
\DeclareMathOperator*{\quotprojlim}{\text{``$\varprojlim$''}}
\DeclareMathOperator*{\quotinjlim}{\text{``$\varinjlim$''}}
\DeclareMathOperator{\Fl}{F\sisi{\ell}{l}}
\DeclareMathOperator{\Pl}{P\sisi{\ell}{l}}
%\DeclareMathOperator{\cl}{c\sisi{\ell}{l}}
\DeclareMathOperator{\tdf}{tdf}
\DeclareMathOperator{\rg}{rg}
\newcommand{\iInd}{\sisi{I}{i}nd}
\newcommand{\iind}{ind}
\newcommand{\ppro}{pro}
\newcommand{\tTor}{\sisi{T}{t}or}
\newcommand{\ttor}{tor}
\newcommand{\chap}{Chap.}
\newcommand{\Ffex}{\Ff\mathit{ex}}
\newcommand{\sheafk}{\underline{k}}
\newcommand{\sheafn}{\underline{n}}
\newcommand{\sheafGamma}{\underline{\Gamma}}

% Repeated references to an already typeset footnote must not use
% \footnotemark without a matching \footnotetext: Hyperref would then create
% a live annotation to a destination that cannot exist.  This helper keeps the
% current printed footnote mark and links it to an explicit stable note target.
\makeatletter
\newcommand{\sgarepeatedfootnotemark}[1]{%
  \begingroup
  \edef\@thefnmark{\thefootnote}%
  \hyperlink{#1}{\@makefnmark}%
  \endgroup}
\makeatother

\sisi{
\DeclareMathOperator{\simarrow}{\rule{.15mm}{1.8mm}\kern -3.4pt\rightsquigarrow}
\newcommand{\sheafell}{\underline{\ell}}
\newcommand{\sheafZZ}{\underline{\ZZ}}
}{
\newcommand{\nderef}[1]{la N.D.E.~(\ref{#1}) page~\pageref{#1}}
}

%%%%%%%%%%%%%%%%%
%%%%% expose XVII
%%%%%%%%%%%%%%%%%

% END INCLUDED SGA4_English_current_buildable_source_20260906/source/fichiers-configuration-latex/macros.tex

\providecommand{\SGAArchiveOnly}[1]{}

\hypersetup{
  pdftitle={SGA 4: Theory of Topoi and Étale Cohomology of Schemes -- English Translation, Exposés I--XIX including V bis},
  pdfauthor={M. Artin; A. Grothendieck; J.-L. Verdier},
  pdfsubject={English translation of SGA 4 proper, Exposés I--XIX including V bis},
  pdfkeywords={SGA 4, topoi, étale cohomology, English translation, exposés I--XIX},
  pdfcreator={XeLaTeX}
}

\begin{document}
\frontmatter
\begin{center}
{\Large Theory of Topoi and Étale Cohomology of Schemes.\\
Volumes 1 to 3}\\
\end{center}

\vskip1cm
\begin{center}
M.~Artin, A.~Grothendieck, J.-L. Verdier
\end{center}

\vskip1cm
\begin{center}
With the participation of N.~Bourbaki, P.~Deligne, B.~Saint-Donat
\end{center}

\selectlanguage{english}
\SGAArchiveOnly{%
\begin{center}
\small Edition: \showgitstatus
\end{center}
}

\clearpage
\SGAArchiveOnly{%
\thispagestyle{empty}
\vspace*{\fill}
\begin{center}
\begin{minipage}{0.88\textwidth}
\small
\textbf{Source and rights notice.} This reader is based on the frozen
Orgogozo SGA~4 TeX snapshot at commit \texttt{71766d9}; the frozen source
archive has SHA-256
\texttt{DA2D939D3BD66B03E7BEFF4353550F7A27982847AE4F798D0872F1DB5D64C7DC}.
No license grant is asserted here for the underlying French text, this
English translation, or third-party TeX support files. Rights remain with
their respective holders; inclusion and provenance metadata do not by
themselves confer reuse permission.
\end{minipage}
\end{center}
\vspace*{\fill}
\clearpage
}
\tableofcontents

\mainmatter

\def\danslelivre{1}

\def\thechapter{\Roman{chapter}}
\def\chaptername{Exposé}

\makeatletter
\def\input@path{{../01/}}
\makeatother
% BEGIN INCLUDED SGA4_English_current_buildable_source_20260906/source/01/01_English.tex
\ifx\danslelivre\undefined
% Release engine: XeLaTeX (the LuaTeX compatibility layer breaks Xy-pic output).
\documentclass[10pt]{../fichiers-configuration-latex/smfart}
% BEGIN INCLUDED SGA4_English_current_buildable_source_20260906/source/fichiers-configuration-latex/paquets.tex
\usepackage{polyglossia}
%\usepackage[francais]{babel}
\usepackage{fontspec}
\usepackage[math-style=ISO,bold-style=upright,sans-style=upright]{unicode-math}
%\usepackage[mathletters]{ucs}
%\usepackage[utf8]{inputenc}
%\usepackage[T1]{fontenc}
\usepackage[unicode,hypertexnames=false,colorlinks=true,linkcolor=blue,citecolor=Sepia,urlcolor=magenta]{hyperref}
\setmainlanguage{french}
\setmainfont[Ligatures=TeX]{Linux Libertine O}
% variante possible : TeXGyrePagella ?
\setmathfont{XITS Math}

% maths
\usepackage{amsmath}
%\usepackage{amssymb}
%\usepackage{mathrsfs}
%\usepackage{stmaryrd}
%\usepackage{euscript}

% couleurs
\usepackage[usenames,dvipsnames]{xcolor}

% graphiques
\usepackage{graphicx}
\usepackage[all]{xy}
%\usepackage{tikz}
%\usetikzlibrary{matrix,arrows,calc}
%\tikzset{column sep=3em,row sep=2em,text height=1.5ex, text depth=0.25ex,>=angle 90}
%\tikzset{description/.style={fill=white,inner sep=1pt}}

% hyperliens, 'etiquettes
\usepackage{xr-hyper}
\usepackage{../fichiers-configuration-latex/smfthm}
%\ifx\cacheretiquettes\undefined
%\usepackage[inline]{showlabels} \else \fi % pour rendre visibles les 'etiquettes
%\usepackage{../fichiers-configuration-latex/makeidx} % pour faire un index
%\usepackage[pagebackref,unicode]{hyperref}  % pour voir o`u sont cit'ees les r'ef'erences
% hyperref is loaded once above with the cumulative-reader destination policy.
%\hypersetup{breaklinks=true} %permet le retour `a la ligne dans les liens trop longs. Incompatible avec pdflatex
%\renewcommand*{\backref}[1]{{\small [cit\'e page(s)~#1.]}}

\usepackage{xspace}
\usepackage{pifont} % pour expose XVII

% END INCLUDED SGA4_English_current_buildable_source_20260906/source/fichiers-configuration-latex/paquets.tex

% BEGIN INCLUDED SGA4_English_current_buildable_source_20260906/source/fichiers-configuration-latex/numerotation.tex
\renewcommand{\thefootnote}{\roman{footnote}} % pour mettre les notes en bas de page sous la forme (ii) et non (2), qui ressemble trop (?) a un exposant.

\newcounter{pageoriginal}
\newcommand{\pageoriginale}{\refstepcounter{pageoriginal}\marginpar{\footnotesize\textbf{\thepageoriginal}}\xspace}
\let\pageoriginaled\pageoriginale

\usepackage{ifthen}
\newboolean{original}
\def\sisi{\ifthenelse{\boolean{original}}}
\setboolean{original}{false}

\newcounter{cnote}
\def\nde#1{\ifthenelse{\boolean{original}}{}{\stepcounter{cnote}\def\thefootnote{\arabic{cnote}}\protect\footnote[\value{cnote}]{N.D.E.: #1}}}
\def\ndemark{\ifthenelse{\boolean{original}}{}{\stepcounter{cnote}\def\thefootnote{\arabic{cnote}}\protect\footnotemark[\value{cnote}]}}
\def\ndetext#1{\ifthenelse{\boolean{original}}{}{\def\thefootnote{\arabic{cnote}}\protect\footnotetext[\value{cnote}]{N.D.E.: #1}}}

% END INCLUDED SGA4_English_current_buildable_source_20260906/source/fichiers-configuration-latex/numerotation.tex

% BEGIN INCLUDED SGA4_English_current_buildable_source_20260906/source/fichiers-configuration-latex/macros.tex
\setmainlanguage{french}
\setotherlanguage{english}


%% Fixed, engine-independent release identifier.  The inherited dynamic Git
%% probe used Lua io.popen, exposed build-state-dependent output, and printed
%% a visible question mark outside a Git checkout.
\def\showgitstatus{English release r2, 2026-07-23}


% macros pour LaTeX
\renewcommand{\mathsf}{}

% Police différente pour U+2113 SCRIPT SMALL L
\setmathfont[range="2113]{latinmodern-math.otf}

% Give every inherited semantic label its own stable PDF destination instead
% of inheriting whichever automatic equation/page anchor happened to be live.
% The legacy label remains as a compatibility alias used by existing edges.
\makeatletter
\def\etiquette[#1]#2{%
  \Hy@raisedlink{\hyper@anchorstart{sga4.semantic.#2}\hyper@anchorend}%
  \def\@currentHref{sga4.semantic.#2}%
  \def\@currentlabel{#1}%
  \label{#2}%
}
\makeatother
\def\eqhyperref[#1]#2{(\hyperref[#1]{#2})}
\def\pagehyperref[#1]#2{(page \hyperref[#1]{#2})}
\def\hyperrefext[#1][#2]#3{\hyperref[#2]{#1 #3}}
\def\eqhyperrefext[#1][#2]#3{\hyperref[#2]{#1 #3}}     %cf. exposes 17
\def\neweqhyperrefext[#1][#2]#3{\hyperref[#2]{#1 #3}}  %cf. exposes 17

\def\shat#1{\protect\widehat{#1}}
\def\stilde#1{\protect\widetilde{#1}}
%%%%%%%%%%%%%%%%%%


\def\etcn{etc.\xspace}
\def\etco{etc...\xspace}
\def\etc{\sisi{\etco}{\etcn}}
\def\ldotsn{\sisi{...}{\etcn}}

\def\cf{{cf.}\kern.3em}
\def\Cf{{Cf.}\kern.3em}
\def\ie{{i.e.}\kern.3em}
\def\iev{{i.e.},\xspace}
\def\loccito{loc.\kern2pt cit.\xspace}
\def\loccitn{\textit{loc.\kern2pt cit.}\xspace}
\def\loccit{\sisi{\loccito}{\loccitn}}
\def\EGA{\textup{EGA}\kern.2em}
\def\SGA{\textup{SGA}\kern.3em}

\def\quoi{}

\def\quav{\quad,\quad}
\def\quaet{\quad \text{et}\quad}
\def\fibreproduct#1#2#3{#1 ×_{#3} #2}
\let\fprod\fibreproduct
\def\ccheck#1#2{\mbox{${\check{\check{\rm #1}}}^{\lower2pt\hbox{$\scriptstyle #2$}}$}}
\def\sheaf#1{\protect\mathcal{#1}}
\def\vieuxsheaf#1{{#1}} % cf. problème avec sheaf{a\Phi_i} par exemple...
\DeclareMathOperator{\Hom}{Hom}
\DeclareMathOperator{\can}{can}
\DeclareMathOperator{\Ind}{Ind}
\DeclareMathOperator{\ind}{ind}
%\DeclareMathOperator{\Hhomtop}{\mathcal{H}\!\mathit{omtop}}
\DeclareMathOperator{\Hhomtop}{\mathbf{Homtop}}
\DeclareMathOperator{\Point}{Point}
\DeclareMathOperator{\Pro}{Pro}
\DeclareMathOperator{\Top}{Top}
\DeclareMathOperator{\Mod}{Mod}
\DeclareMathOperator{\Cart}{Cart}
%\DeclareMathOperator{\Mmorsite}{\mathcal{M}\!\mathit{orsite}}
\DeclareMathOperator{\Mmorsite}{\mathbf{Morsite}}
\DeclareMathOperator{\loc}{loc}
\DeclareMathOperator{\Loc}{Loc}
\DeclareMathOperator{\Pic}{Pic}
\DeclareMathOperator{\Et}{Et}
\DeclareMathOperator{\card}{card}
\DeclareMathOperator{\Ens}{Ens}
\DeclareMathOperator{\h}{\underline{H}}
\DeclareMathOperator{\Fonct}{Fonct}
\DeclareMathOperator{\Ker}{Ker}
\DeclareMathOperator{\but}{but}
\DeclareMathOperator{\Ob}{Ob}
\DeclareMathOperator{\ob}{ob}
\DeclareMathOperator{\Lib}{Lib}
\DeclareMathOperator{\esj}{esj}
\DeclareMathOperator{\id}{id}
\DeclareMathOperator{\Id}{Id}
\DeclareMathOperator{\ad}{ad}
\DeclareMathOperator{\Iim}{Im}
\DeclareMathOperator{\pr}{pr}
\DeclareMathOperator{\Cat}{Cat}
\DeclareMathOperator{\car}{car}
\DeclareMathOperator{\SysInd}{SysInd}
\DeclareMathOperator{\Sup}{Sup}
\DeclareMathOperator{\Coker}{Coker}
\DeclareMathOperator{\Filt}{Filt}
\DeclareMathOperator{\acc}{acc}
\DeclareMathOperator{\sur}{sur}
\DeclareMathOperator{\cat}{cat}
\DeclareMathOperator{\Ord}{Ord}
\DeclareMathOperator{\Gr}{Gr}
\DeclareMathOperator{\Cov}{Cov}
\DeclareMathOperator{\Caf}{Caf}
\DeclareMathOperator{\coker}{coker}
\DeclareMathOperator{\coim}{coim}
\DeclareMathOperator{\im}{im}
\DeclareMathOperator{\Ouv}{Ouv}
\DeclareMathOperator{\diag}{diag}
\DeclareMathOperator{\Esp}{Esp}
\DeclareMathOperator{\TOP}{TOP}
\DeclareMathOperator{\Unif}{Unif}
\DeclareMathOperator{\Iind}{\mathbf{Ind}}
\DeclareMathOperator{\Hhomcart}{\mathbf{Homcart}}
\DeclareMathOperator{\Hhomtopess}{\mathbf{Homtopess}}
\DeclareMathOperator{\Ppoints}{\mathbf{Points}}
\DeclareMathOperator{\Ppoint}{\mathbf{Point}}
\DeclareMathOperator{\Ttors}{\mathbf{Tors}}
\DeclareMathOperator{\Ssex}{\mathbf{Sex}}
\DeclareMathOperator{\Ffib}{\mathbf{Fib}}
\DeclareMathOperator{\Vvois}{\mathbf{Vois}}
\DeclareMathOperator{\Ppointess}{\mathbf{Pointess}}
\DeclareMathOperator{\scrHom}{\mathbf{Hom}}
\DeclareMathOperator{\sheafPic}{\mathbf{Pic}}
\DeclareMathOperator{\sheafAut}{\mathbf{Aut}}
\DeclareMathOperator{\Hhomtoplocan}{\mathbf{Homtoplocan}}
\DeclareMathOperator{\Hhomtopan}{\mathbf{Homtopan}}
\DeclareMathOperator{\Pptgeom}{\mathbf{Ptgéom}}
\DeclareMathOperator{\Morsite}{Morsite}
\DeclareMathOperator{\Fib}{Fib}
\DeclareMathOperator{\Fais}{Fais}
\DeclareMathOperator{\kar}{kar}
\DeclareMathOperator{\Groupes}{Groupes}
\DeclareMathOperator{\Restr}{Restr}
\DeclareMathOperator{\Isom}{Isom}
\DeclareMathOperator{\Prol}{Prol}
\DeclareMathOperator{\Res}{Res}
\DeclareMathOperator{\Tors}{Tors}
\DeclareMathOperator{\esp}{esp}
\DeclareMathOperator{\es}{es}
\DeclareMathOperator{\od}{od}
\DeclareMathOperator{\gr}{gr}
\DeclareMathOperator{\sob}{sob}
\DeclareMathOperator{\dfn}{dfn}
\DeclareMathOperator{\Sch}{Sch}
\DeclareMathOperator{\ouv}{ouv}
\DeclareMathOperator{\cons}{cons}
\DeclareMathOperator{\Nerf}{Nerf}
\DeclareMathOperator{\Points}{Points}
\DeclareMathOperator{\Espsob}{Espsob}
\DeclareMathOperator{\Homtop}{Homtop}
\DeclareMathOperator{\Pt}{Pt}
\DeclareMathOperator{\Quot}{Quot}
\DeclareMathOperator{\glob}{glob}
\DeclareMathOperator{\Homfib}{Homfib}
\DeclareMathOperator{\Qouv}{Qouv}
\DeclareMathOperator{\Quotouv}{Quotouv}
\DeclareMathOperator{\ens}{ens}
\DeclareMathOperator{\ab}{ab}
\DeclareMathOperator{\an}{an}
\DeclareMathOperator{\fs}{fs}
\DeclareMathOperator{\red}{red}
\DeclareMathOperator{\Ab}{Ab}
\DeclareMathOperator{\cosq}{cosq}
\DeclareMathOperator{\Ann}{Ann}
\DeclareMathOperator{\Zar}{Zar}


%%%%%%%%%%%%
%%%%% SGA4-2
\DeclareMathOperator{\Sex}{Sex}
\DeclareMathOperator{\vois}{vois}
\DeclareMathOperator{\Vois}{Vois}
\DeclareMathOperator{\Ext}{Ext}
\DeclareMathOperator{\ext}{ext}
\DeclareMathOperator{\Fer}{Fer}
\DeclareMathOperator{\Add}{Add}
\DeclareMathOperator{\Cocart}{Cocart}
\DeclareMathOperator{\cosk}{cosk}
\DeclareMathOperator{\plat}{plat}

\DeclareMathOperator{\faisc}{faisc}
\DeclareMathOperator{\sk}{sk}
\DeclareMathOperator{\ch}{ch}

\def\HH{\mathbf{H}}
\def\qc{\mathrm{qc}}
\def\et{\mathrm{et}}
\def\coh{\mathrm{coh}}

\def\Hvr#1{%
\begin{array}[b]{l}
{\kern2pt\scriptstyle r}\\[-8pt]
{\kern1pt\scriptstyle \vee}\\[-5pt]
H^{#1}
\end{array}}

\def\Hvv#1{%
\begin{array}[b]{l}
\kern2pt\rotatebox{90}{${\scriptscriptstyle\ll}$}\\[-6pt]
H^{#1}
\end{array}}

\def\Hvi#1{%
\begin{array}[b]{l}
{\kern0.2pt\scriptscriptstyle \infty}\\[-8pt]
{\kern1pt\scriptstyle \vee}\\[-5pt]
H^{#1}
\end{array}}

\def\Hv#1{%
\begin{array}[b]{l}
{\kern1pt\scriptstyle \vee}\\[-5pt]
H^{#1}
\end{array}}

\DeclareMathOperator{\cart}{cart}
\DeclareMathOperator{\Ffaisc}{\mathbf{Faisc}}


\def\Ev#1{%
\begin{array}[b]{l}
{\kern1pt\scriptstyle \vee}\\[-5pt]
E'_{#1}
\end{array}}

\DeclareMathOperator{\Cont}{Cont}
\DeclareMathOperator{\Ccont}{\mathscr{C}\!\mathit{ont}}
\DeclareMathOperator{\Site}{Site}
\newcommand\MM{\mathbf{M}}
\def\limtop#1{{\displaystyle{\varprojlim_{#1}\!\mbox{\rm top }}}}
\def\Limtop{\mathop{\varprojlim\!\mbox{\rm top}}}
\DeclareMathOperator{\spec}{spec}
\def\ss{\scriptstyle{\mathscr{S}}}
\DeclareMathOperator{\Gal}{Gal}
\def\smallm{\scriptstyle{\mathscr{M}}}
\DeclareMathOperator{\End}{End}
\DeclareMathOperator{\ass}{ass}
\DeclareMathOperator{\zar}{zar}
\DeclareMathOperator{\Ogg}{Ogg}
\DeclareMathOperator{\res}{res}
\DeclareMathOperator{\tr}{tr}
\DeclareMathOperator{\Supp}{Supp}
\DeclareMathOperator{\codim}{codim}
\DeclareMathOperator{\rk}{rk}
\DeclareMathOperator{\cur}{cur}
\DeclareMathOperator{\cdqc}{cdqc}
\newcommand\DD{\mathbf{D}}
\DeclareMathOperator{\cl}{cl}
\def\JJ{\mathbf{J}}
\def\FF{\mathbf{F}}
\def\BB{\mathbf{B}}
\DeclareMathOperator{\Alget}{Alget}
\DeclareMathOperator{\Ass}{Ass}
\DeclareMathOperator{\supp}{supp}
\DeclareMathOperator{\rad}{rad}
\DeclareMathOperator{\prof}{prof}
\DeclareMathOperator{\cint}{int}
\DeclareMathOperator{\sic}{sic}
\DeclareMathOperator{\Tor}{Tor}
\DeclareMathOperator{\tors}{tors}
\DeclareMathOperator{\Sym}{Sym}
\DeclareMathOperator{\parf}{parf}
\DeclareMathOperator{\tor}{tor}
\DeclareMathOperator{\defn}{defn}
\DeclareMathOperator{\Div}{Div}
\DeclareMathOperator{\Pol}{Pol}
\DeclareMathOperator{\Fex}{Fex}
\DeclareMathOperator{\Lie}{Lie}
\DeclareMathOperator{\tot}{tot}
\DeclareMathOperator{\Aut}{Aut}
\DeclareMathOperator{\HOM}{HOM}
\DeclareMathOperator{\pch}{pch}
\DeclareMathOperator{\Ch}{Ch}
\DeclareMathOperator{\EXT}{EXT}
\DeclareMathOperator{\PIC}{PIC}
\DeclareMathOperator{\Rf}{Rf}
\DeclareMathOperator{\inj}{inj}
\newcommand\VV{\mathbf{V}}

\def\Xv{%
\begin{array}[b]{l}
{\kern1pt\scriptstyle \vee}\\[-5pt]
X
\end{array}}

%%%%%%%
\newcommand\TT{\mathbf{T}}
\newcommand\XX{\mathbf{X}}
\newcommand\EE{\mathbf{E}}
%%%%%%

\def\qedsymbol{\textup{C.Q.F.D.}}
\let\CQFD\qed
\let\cqfd\qed

%%%%%\bigdot
\font\bigsymb=cmsy10 at 4pt
\def\bigdot{{\kern1.2pt\raise 1.5pt\hbox{\bigsymb\char15}}}
%%%%%

\def\Lright{\mathop{\vcenter{\hbox{\hbox{$\longrightarrow$}\lower8pt \hbox{\kern-16pt\hbox{$\longrightarrow$}}}}}}

\def\Aa{{\mathcal A}}
\def\Bb{{\mathcal B}}
\def\Cc{{\mathcal C}}
\def\Dd{{\mathcal D}}
\def\Ee{{\mathcal E}}
\def\Ff{{\mathcal F}}
\def\Gg{{\mathcal G}}
\def\Hh{{\mathcal H}}
\def\Ii{{\mathcal I}}
\def\Jj{{\mathcal J}}
\def\Kk{{\mathcal K}}
\def\Ll{{\mathcal L}}
\def\Mm{{\mathcal M}}
\def\Nn{{\mathcal N}}
\def\Oo{{\mathcal O}}
\def\Pp{{\mathcal P}}
\def\Qq{{\mathcal Q}}
\def\Rr{{\mathcal R}}
\def\Ss{{\mathcal S}}
\def\Tt{{\mathcal T}}
\def\Uu{{\mathcal U}}
\def\Vv{{\mathcal V}}
\def\Ww{{\mathcal W}}
\def\Xx{{\mathcal X}}
\def\Yy{{\mathcal Y}}
\def\Zz{{\mathcal Z}}

\newcommand{\CC}{\mathbf{C}}
\newcommand{\GG}{\mathbf{G}}
\newcommand{\KK}{\mathbf{K}}
\newcommand{\LL}{\mathbf{L}}
\newcommand{\NN}{\mathbf{N}}
\newcommand{\PP}{\mathbf{P}}
\newcommand\QQ{\mathbf{Q}}
\newcommand\RR{\mathbf{R}}
\newcommand\ZZ{\mathbf{Z}}

\def\to{\mathchoice{\longrightarrow}{\rightarrow}{\rightarrow}{\rightarrow}}
\def\from{\mathchoice{\longleftarrow}{\leftarrow}{\leftarrow}{\leftarrow}}
\def\mto{\mathchoice{\longmapsto}{\mapsto}{\mapsto}{\mapsto}}
\def\mfrom{\mathrel{\from\joinrel\mapstochar\,}}
\def\hto{\mathrel{\lhook\joinrel\to}}
\def\hfrom{\mathrel{\from\joinrel\rhook}}

\def\implies{\mathchoice{\Longrightarrow}{\Rightarrow}{\Rightarrow}{\Rightarrow}}
\def\implied{\mathchoice{\Longleftarrow}{\Leftarrow}{\Leftarrow}{\Leftarrow}}
\def\iff{\mathchoice{\Longleftrightarrow}{\Leftrightarrow}{\Leftrightarrow}{\Leftrightarrow}}

\def\To#1{\mathchoice{\xrightarrow{\labelstyle~~#1~}}{\stackrel{#1}{\longrightarrow}}{}{}}
\def\Mto#1{\mathrel{\mapstochar\To{#1}}}
\def\From#1{\mathchoice{\xleftarrow{\labelstyle~~#1~}}{\stackrel{#1}{\longleftarrow}}{}{}}

\newcommand{\isomto}{\overset{\sim}{\longrightarrow}}
\newcommand{\isomfrom}{\overset{\sim}{\longleftarrow}}

\newcommand{\IFF}{\unskip~$\iff$~\ignorespaces}
\newcommand{\THEN}{\unskip~$\implies$~\ignorespaces}

\let\leq\leqslant
\let\geq\geqslant
\let\le\leqslant
\let\ge\geqslant

\def\cf{{cf.}\kern.3em}
\def\Cf{{Cf.}\kern.3em}
\def\ie{{i.e.}\kern.3em}
\def\iev{{i.e.},\xspace}
%\def\loccito{loc.\kern2pt cit.\xspace}
%\def\loccitn{\textit{loc.\kern2pt cit.}\xspace}
%\def\loccit{\sisi{\loccito}{\loccitn}}
\def\EGA{\textup{EGA}\kern.2em}
\def\SGA{\textup{SGA}\kern.3em}

\let\oldS\S
\def\S{\oldS\kern .15em}
\def\ptbl{.\kern .2em}
\def\resp{resp.\kern.3em}
\def\Exp{Exp.\kern.3em}

\let\bar\overline
\let\hat\widehat
\let\tilde\widetilde
\let\dpl\displaystyle
\let\ldots\dots
\let\emptyset\varnothing
\let\skpt\skippointrait
\let\epsilon\varepsilon

\def\sfrac#1#2{{#1}/{#2}}

\DeclareMathOperator{\cd}{cd}
%\DeclareMathOperator{\Hhom}{\mathcal{H}\!\mathit{om}}
\DeclareMathOperator{\Hhom}{\mathbf{Hom}}
%\DeclareMathOperator{\Eext}{\mathcal{E}\!\mathit{xt}}
\DeclareMathOperator{\Eext}{\mathbf{Ext}}
\DeclareMathOperator{\Spec}{Spec}
\DeclareMathOperator{\Tr}{Tr}

\let\emptyset=\varnothing

%%%%%%%%%%%%%%%
%%% expose XVII
%%%%%%%%%%%%%%%

\makeatletter
\let\old@optexpose\optexpose
\newcommand{\empty@optexpose}[1]{}
\newcommand{\paren@optexpose}[1]{\old@optexpose{#1}(}
\newcommand{\refnoexpose}[1]{\let\optexpose\empty@optexpose\ref{#1}\let\optexpose\old@optexpose}%\refnoexpose{X.cor5.2} engendre "5.2" au lieu de "X 5.2"
\newcommand{\neweqref}[1]{\let\optexpose\paren@optexpose\textup{\ref{#1})}\let\optexpose\old@optexpose}%\neweqref{XII.eq4.1} engendre "XII (4.1)" au lieu de "(XII 4.1)"
\makeatother
\newcommand{\refvb}[2][VI${}_{\text{B}}$ ]{\sisi{#1\refnoexpose{#2}}{\ref{#2}}}

\newcommand{\leftexp}[2]{{\vphantom{#2}}^{#1}{#2}}
\newcommand{\AAA}{\mathbf{A}}
\newcommand{\Sf}{\mathfrak{S}}
\DeclareMathOperator*{\quotprojlim}{\text{``$\varprojlim$''}}
\DeclareMathOperator*{\quotinjlim}{\text{``$\varinjlim$''}}
\DeclareMathOperator{\Fl}{F\sisi{\ell}{l}}
\DeclareMathOperator{\Pl}{P\sisi{\ell}{l}}
%\DeclareMathOperator{\cl}{c\sisi{\ell}{l}}
\DeclareMathOperator{\tdf}{tdf}
\DeclareMathOperator{\rg}{rg}
\newcommand{\iInd}{\sisi{I}{i}nd}
\newcommand{\iind}{ind}
\newcommand{\ppro}{pro}
\newcommand{\tTor}{\sisi{T}{t}or}
\newcommand{\ttor}{tor}
\newcommand{\chap}{Chap.}
\newcommand{\Ffex}{\Ff\mathit{ex}}
\newcommand{\sheafk}{\underline{k}}
\newcommand{\sheafn}{\underline{n}}
\newcommand{\sheafGamma}{\underline{\Gamma}}

% Repeated references to an already typeset footnote must not use
% \footnotemark without a matching \footnotetext: Hyperref would then create
% a live annotation to a destination that cannot exist.  This helper keeps the
% current printed footnote mark and links it to an explicit stable note target.
\makeatletter
\newcommand{\sgarepeatedfootnotemark}[1]{%
  \begingroup
  \edef\@thefnmark{\thefootnote}%
  \hyperlink{#1}{\@makefnmark}%
  \endgroup}
\makeatother

\sisi{
\DeclareMathOperator{\simarrow}{\rule{.15mm}{1.8mm}\kern -3.4pt\rightsquigarrow}
\newcommand{\sheafell}{\underline{\ell}}
\newcommand{\sheafZZ}{\underline{\ZZ}}
}{
\newcommand{\nderef}[1]{la N.D.E.~(\ref{#1}) page~\pageref{#1}}
}

%%%%%%%%%%%%%%%%%
%%%%% expose XVII
%%%%%%%%%%%%%%%%%

% END INCLUDED SGA4_English_current_buildable_source_20260906/source/fichiers-configuration-latex/macros.tex


%\externaldocument{../01/01.tex}
\externaldocument{../02/01.tex}
\externaldocument{../03/03.tex}
\externaldocument{../04/04.tex}

\setboolean{original}{false}

\def\labelenumi{(\theenumi)}
\def\labelenumii{(\theenumii)}
\def\labelenumiii{(\theenumiii)}
\def\labelenumiv{(\theenumiv)}

\begin{document}

\begin{center}
I. Presheaves\\
A. Grothendieck and J.-L. Verdier
\end{center}

\selectlanguage{english}version: \showgitstatus

\tableofcontents

\else
\chapter{Presheaves}
\begin{center}
A. Grothendieck and J.-L. Verdier
\end{center}
\fi
\phantomsection\MakeLinkTarget*{sga4.exp.I}\label{sga4.exp.I}

% BEGIN INCLUDED SGA4_English_current_buildable_source_20260906/source/01/drafts/SGA4_I_LINES_46_100_English.texfrag
In\pageoriginale numbers 0 through 5 of this exposé, we present the
elementary, and most often well-known, properties of categories of
presheaves\footnote[1]{The reader may also consult SGA 3 I \S\S 1 to 3.}.
These properties are used constantly throughout the rest of the seminar,
and knowledge of them is essential for understanding the subsequent
exposés. The proofs are immediate; most often, they are omitted. In
numbers 6 through 9, we take up a few themes used on various occasions
later. The reader in a hurry %%
may omit them on a first reading. Number 10 fixes the %%
terminology used. Appendix 11 is %%
due to N. Bourbaki.


\setcounter{section}{-1}
\section{Universes}\etiquette[0]{I.sec0}

A \emph{universe} is a nonempty set\nde{This definition contradicts
Bourbaki's appendix \textit{below}. However, the usefulness of empty
universes seems debatable...} $𝒰$ having the following properties:
\begin{itemize}
\item[($𝒰$ 1)] If $x\in 𝒰$ and $y\in x$, then $y\in 𝒰$.

\item[($𝒰$ 2)] If $x$, $y\in 𝒰$, then $\{x,y\}\in 𝒰$.

\item[($𝒰$ 3)] If $x\in 𝒰$, then $\Pp(x)\in 𝒰$.

\item[($𝒰$ 4)] If $(x_i, i\in I\sisi{\in 𝒰}{})$ is a
  family of elements of $𝒰$\sisi{}{ and $I\in𝒰$}, then $\cup_{i\in I}x_i \in
  𝒰$.
\end{itemize}

From\pageoriginale the preceding axioms one readily deduces the following
properties:
\begin{itemize}
\item If $x\in 𝒰$, the set $\{x\}$ belongs to $𝒰$.

\item If $x$ is a subset of $y\in 𝒰$, then $x\in 𝒰$.

\item If $x$, $y\in 𝒰$, the ordered pair $(x,y)=\{\{x\},\{x,y\}\}$ (Kuratowski's
  definition) is an element of $𝒰$.

\item If $x$, $y\in 𝒰$, the union $x\cup y$ and %%
the product $x\times
   y$ are elements of $𝒰$.

\item If $(x_i, i\in I\in 𝒰)$ is a family
  of elements of $𝒰$,
   the product $\prod_{i\in I}x_i$ is an element of $𝒰$.

\item If $x\in 𝒰$, then $\card(x)< \card(𝒰)$ (strictly). In
   particular, the relation $𝒰\in 𝒰$ does not hold.
\end{itemize}

% END INCLUDED SGA4_English_current_buildable_source_20260906/source/01/drafts/SGA4_I_LINES_46_100_English.texfrag


% BEGIN INCLUDED SGA4_English_current_buildable_source_20260906/source/01/drafts/SGA4_I_LINES_102_196_English.texfrag
We can therefore carry out all the usual operations of set theory
starting from the elements of a universe without thereby causing the
final result to cease to be an element of the universe.

The first advantage of the notion of a universe is that it provides a
definition of the usual categories: the category of sets belonging to
the universe $𝒰$ ($𝒰$-$\Ens$), the category of
topological spaces belonging to the universe $𝒰$, the category of
commutative groups belonging to the universe $𝒰$
($𝒰$-$\Ab$), the category of categories belonging to
the universe $𝒰$\ldots.

However, the only known universe is the set of symbols of the form
$\{\{∅\}, \{\{∅\}, ∅\}\}$ \etc. (all
the elements of this universe are finite sets, and this universe
is countable). In particular, no universe is known that
contains an element of infinite cardinality. We are therefore led
to add
to the axioms of set theory the axiom: $(𝒰 A)$
\emph{For every set $x$ there exists a universe $𝒰$ such
that $x\in
  𝒰$.}

Since the intersection\pageoriginale of a family of universes is a
universe, it follows immediately that every set is
an element
of a smallest universe. One can show that axiom $(𝒰
A)$ is independent of the axioms of set theory.

We shall also add the axiom:


\noindent \emph{$(𝒰 B)$ Let $R\{x\}$ be a relation and let
$𝒰$ be a universe. If there
  exists an element $y\in 𝒰$ such that $R\{y\}$, then
  $\tau_xR\{x\}\in 𝒰$.}

The consistency of axioms $(𝒰 A)$ and $(𝒰 B)$
relative to the other axioms of set theory has not been
proved and, it seems, cannot be proved.

Let $𝒰$ be a universe and let $c(𝒰)$ be the supremum
of the cardinalities of the elements of $𝒰(c(𝒰)\leq
\card(𝒰))$. The cardinal $c(𝒰)$ has the following
properties:

$(FI)$ If $a<c(𝒰)$, then $2^a<c(𝒰)$.

$(FII)$ If $(a_i, i\in I)$ is a family of cardinals strictly
less than $c(𝒰)$ and if $\card(I)$ is %%
strictly less than $c(𝒰)$, then $\Sigma_{i\in I} a_i <
c(𝒰)$.

The cardinals that have properties $(FI)$ and $(FII)$ are
called \emph{strongly inaccessible cardinals.}

The cardinal 0 and the countably infinite cardinal are strongly
inaccessible.

Axiom $(𝒰 A)$ implies:

\emph{$(𝒰 A')$ Every cardinal is strictly less than a
  strongly inaccessible cardinal.}

Conversely, one can show (11) that the consistency of
$(UA')$ implies the consistency of $(UA)$, and that the
consistency of these axioms entails that of axiom
$(𝒰 B)$.

Call\pageoriginale an \emph{Artinian set} any set $E$
for which there is no infinite family $(x_n, n\in N)$ such
that $x_0 \in E$, $x_{n+1}\in x_n$. One can then show
[\loccit] that there is a one-to-one correspondence between the
strongly inaccessible cardinals and the Artinian universes, defined
as follows: to every strongly inaccessible cardinal $c$ we
associate the unique Artinian universe $U_c$ such that
$$
\card(𝒰_c)=c{\quoi}.
$$
Let $c<c'$ be two strongly inaccessible cardinals. We have:
$$
𝒰_c\in 𝒰_{c'}{\quoi}.
$$
In particular, the Artinian universes of cardinality less than a
given cardinal form a well-ordered set (with respect to the
membership relation). Axiom $(𝒰 A')$ is equivalent to
the axiom:

$(𝒰 A'')$. \emph{Every Artinian set is an element of an
Artinian
  universe.}

Observe that all the usual sets $(\ZZ, \QQ,\ldots)$ are
Artinian sets.

% END INCLUDED SGA4_English_current_buildable_source_20260906/source/01/drafts/SGA4_I_LINES_102_196_English.texfrag


% BEGIN INCLUDED SGA4_English_current_buildable_source_20260906/source/01/drafts/SGA4_I_LINES_197_410_English.texfrag
\section{$𝒰$-categories. Presheaves of sets}\etiquette[1]{I.sec1}

\setcounter{subsection}{-1}
\subsection{}\etiquette[1.0]{I.sec1.0}
\emph{Throughout the remainder of the seminar and unless expressly stated
  otherwise, the universes under consideration will contain an element
  of infinite cardinality.} Let $𝒰$ be a universe. A set is said to be
  $𝒰$-\emph{small} (or, when no confusion can result,
  \emph{small}) if it is \emph{isomorphic} to an element of
  $𝒰$. We shall also use the terminology: small group, small ring,
  small category\nde{A category $C$ is regarded as a set of arrows
  (equipped with the subset of objects, regarded as the set of identity
  arrows, and with the ``source, target'' maps). Thus, the expressions ``$C$ is an element of $𝒰$'' or
  ``$C$ is $𝒰$-small'' make sense.}
  \ldots We shall often assume, without
  explicit mention, that the schemes, topological spaces, indexing
  sets\ldots with which we work are $\in 𝒰$,
or\pageoriginale at least have cardinality $\in 𝒰$;
  nevertheless, many of the categories with which we shall work
  will not be $\in 𝒰$.

\begin{enonce*}{Definition 1.1}\etiquette[1.1]{I.def1.1}
Let $𝒰$ be a universe and $C$ a category. We say that $C$
is a $𝒰$-category if, for every pair $(x,y)$ of objects
of $C$, the set $\Hom_C(x,y)$ is $𝒰$-small.
\end{enonce*}

\setcounter{subsection}{1}
\subsubsection{}\etiquette[1.1.1]{I.sec1.1.1}
Let $C$ and $D$ be two categories and let $\text{Fonct}(C,D)$ be the
category of functors from $C$ to $D$. The following assertions
are immediate:
\begin{enumerate}
\item[a)] If $C$ and $D$ are elements of a universe $𝒰$
  (\resp  are $𝒰$-small), the category $\text{Fonct}(C,D)$ is an element
  of $𝒰$ (\resp  is $𝒰$-small).

\item[b)] If $C$ is $𝒰$-small and $D$ is a $𝒰$-category,
  $\text{Fonct}(C,D)$ is a $𝒰$-category.
\end{enumerate}

\begin{enonce*}[remark]{Remark 1.1.2}\etiquette[1.1.2]{I.rem1.1.2}
Let $D$ be a category having the following properties:
\begin{itemize}
\item[(C1)] The set $\ob(D)$ is contained in the universe $𝒰$.

\item[(C2)] For every pair $(x,y)$ of objects of $D$, the set
  $\Hom_D(x,y)$ is an element of $𝒰$.
\end{itemize}
(The usual categories constructed from a universe
$𝒰$ have these two properties: $𝒰$-$\Ens$,
$𝒰$-$\Ab$,\ldots). Let $C$ be a category belonging to
$𝒰$. Then the category $\text{Fonct}(C,D)$ does not
in general have properties (C1) and (C2). For example, the category
$\text{Fonct}(C,𝒰$-$\Ens$) has neither property
(C1) nor property (C2). This justifies the adopted definition of
$𝒰$-category, in preference to the more restrictive notion given
by conditions (C1) and (C2) above.
\end{enonce*}

\begin{enonce*}{Definition 1.2}\etiquette[1.2]{I.def1.2}
Let\pageoriginale $C$ be a category. The category of %%
presheaves of sets on $C$ relative to the universe $𝒰$
(or, when no confusion can result, the category of
presheaves on $C$) is the category of contravariant functors on
$C$ with values in the category of $𝒰$-sets.
\end{enonce*}

We denote by $\widehat{C}_{𝒰}$ (or more
simply, when no confusion can result, by
$\widehat{C}$\,) the category of presheaves
of sets on $C$ relative to the universe $𝒰$. The objects of
$\widehat{C}_{𝒰}$ are called
$𝒰$-presheaves (or simply presheaves) on $C$.
When $C$ is $𝒰$-small, the category
$\widehat{C}_{𝒰}$ is a
$𝒰$-category. When $C$ is a $𝒰$-category,
$\widehat{C}_{𝒰}$ is not necessarily a
$𝒰$-category.


\begin{enonce*}[remark]{Construction-Definition
    1.3}\etiquette[1.3]{I.condef1.3}
Let $x$ be an object of a $𝒰$-category $C$. The
$𝒰$-\emph{functor represented by} $x$ is the functor
$h_{𝒰}(x):C^\circ\rightarrow 𝒰$-$\Ens$ whose
construction is as follows \footnote[1]{$C^\circ$ will always denote the
category opposite
  to $C$.}. Let $y$ be an object of $C$.
\begin{enumerate}
\item[a)] If $\Hom_C(y,x)$ is an element of $𝒰$, then:
      $$
      h_{𝒰}(x)(y)=\Hom_C(y,x)
      $$

\item[b)] Suppose that $\Hom_C(y,x)$ is not an element of $𝒰$
  and let $R(Z,x,y)$ be the relation: ``The set $Z$ is the codomain of an
  isomorphism $\Hom_C(y,x)\xrightarrow{\sim}Z$''. We then set:
  $$
  h_{𝒰}(x)(y)=\tau_ZR(Z){\quoi}.
  $$
\end{enumerate}
(By axiom $(𝒰 B)$, $h_{𝒰}(x)(y)$ is an
element of $𝒰$). Let $R'(u,x,y)$ be the relation: « $u$ is
a bijection
$$
u:\Hom_C(y,x)\xrightarrow{\sim} h_{𝒰}(x)(y)\mbox{
»}{\quoi}.
$$
We\pageoriginale then set:
$$
\varphi(y,x)=\tau_uR'(u){\quoi}.
$$

It should be noted that, in cases (a) and (b), we have $a$ canonical
isomorphism:
$$
\varphi(y,x):\Hom_C(y,x)\xrightarrow{\sim}
h_{𝒰}(x)(y){\quoi}.
$$
(In case a) $\varphi(y,x)$ is the identity). Now let
$u:y\rightarrow y'$ be an arrow of $C$. By composition of morphisms, the
morphism $u$ defines a map:
$$
\Hom_C(u,x):\Hom_C(y',x)\longrightarrow \Hom_C(y,x){\quoi}.
$$
We then set
$$
h_{𝒰}(x)(u)=\varphi(y,x)\Hom_C(u,x)\varphi(y',x)^{{}-1}{\quoi}.
$$

It is immediate that $h_{𝒰}(x)$, so defined, is
a functor $C^\circ\rightarrow 𝒰$-$\Ens$.
\end{enonce*}

\setcounter{subsection}{3} \setcounter{subsubsection}{0}
\subsubsection{}\etiquette[1.3.1]{I.sec1.3.1}
Similarly, using the isomorphisms $\varphi$, for every
morphism $y:x\rightarrow x'$ we define a morphism of functors:
$$
h_{𝒰}(y):h_{𝒰}(x)\longrightarrow h_{𝒰}(x')
$$
and it is immediate that this defines a functor
$$
h_{𝒰}:C\longrightarrow
\widehat{C}_{𝒰}{\quoi}.
$$

\subsubsection{}\etiquette[1.3.2]{I.sec1.3.2}
Now let $𝒱$ be a universe such that $𝒰\subset
𝒱$. There is then a canonical fully faithful functor:
$$
𝒰\text{-}\Ens\hookrightarrow 𝒱\text{-}\Ens
$$
whence a fully faithful functor:
$$
\widehat{C}_{𝒰}\hookrightarrow
\widehat{C}_{𝒱}
$$
and\pageoriginale the diagram:

%\def\objectstyle{\mathrm}--here
\[
\xymatrix@R=1.75cm@C=3.5cm@M=.2cm{
C\ar[r]^{h_{\underline{U}}} \ar[dr]_{h_{\underline{V}}}& C{\widehat{\phantom{i}}}_{\underline{U}}\ar@{_{(}->}[d]\\
& C\widehat{\phantom{i}}_{\underline{V}} }
\]
commutes up to canonical isomorphism. When $C$ is an
element of $𝒰$, this canonical isomorphism is the identity.

\subsubsection{}\etiquette[1.3.3]{I.sec1.3.3}
\emph{In practice, the universe $𝒰$ is fixed
  once and for all and is not mentioned.} We then use the
  notations $\widehat{C}$ (for the category of
  $𝒰$-presheaves of sets) and
$$
h:C\longrightarrow \widehat{C}{\quoi}.
$$
For every object $x$ of $C$, the presheaf $h(x)$ is called the
\emph{presheaf represented by} $x$, and we shall always identify
the value at $y\in \ob(C)$ of the presheaf $h(x)$ with $\Hom_C(y,x)$.

\begin{enonce*}{Proposition 1.4}\etiquette[1.4]{I.prop1.4}
Let $C$ be a $𝒰$-category, $F$ a presheaf on $C$, and
$X$ an object of $C$. There is an isomorphism functorial in $X$ and
in $F$:
$$
i:\Hom_{\widehat{C}}(h(X),F)\xleftarrow{\sim}
F(X){\quoi},
$$
where the map $i$, when $F$ is of the form $h(Y)$, is none other
than the map:
$$
h:\Hom(h(X), h(Y))\longleftarrow \Hom(X,Y){\quoi}.
$$
In particular, the functor $h$ is fully faithful.
\end{enonce*}

\setcounter{subsection}{4} \setcounter{subsubsection}{0}
\subsubsection{}\etiquette[1.4.1]{I.sec1.4.1}
This proposition justifies the customary abuse of language identifying
an object of $C$ with the corresponding contravariant functor. A
presheaf\pageoriginale isomorphic to an object in the image of $h$ (or,
using the abuse of language noted above, isomorphic to an object
of $C$) is called a \emph{representable presheaf.}

\subsubsection{}\etiquette[1.4.2]{I.sec1.4.2}
Let $𝒱$ be a universe containing the universe $𝒰$. The
category of $𝒰$-presheaves is a \emph{full} subcategory
of the category of $𝒱$-presheaves, and consequently
a $𝒰$-presheaf is representable if and only
if its image in the category of $𝒱$-presheaves is a
representable $𝒱$-presheaf.

% END INCLUDED SGA4_English_current_buildable_source_20260906/source/01/drafts/SGA4_I_LINES_197_410_English.texfrag


% BEGIN INCLUDED SGA4_English_current_buildable_source_20260906/source/01/drafts/SGA4_I_LINES_411_824_English.texfrag
\section{Projective and inductive limits}\etiquette[2]{I.sec2}

Let $C$ be a $𝒰$-category and $I$ a small category.
Denote by $\Fonct(I,C)$ the category of functors from $I$ to $C$. The
category $\Fonct(I,C)$ is a $𝒰$-category. To every object
$X$ of $C$, associate the subcategory $𝒳$ of $C$ whose
only object is $X$ and whose only arrow is the identity of $X$.
Let $i_X:𝒳\rightarrow C$ denote the inclusion functor.
There is a unique functor $e_X:I\rightarrow 𝒳$, and
we denote by $k_X:I\rightarrow C$ the functor $i_X\circ
e_X$. (We say that $k_X$ is the constant functor with value $X$.) The
correspondence $X\mapsto k_X$ is clearly functorial in $X$,
which allows us to define, for every functor $G:I\rightarrow
C$, the presheaf $X\mapsto \Hom_{\Fonct(I,C)}(k_X,G)$.

\begin{enonce*}{Definition 2.1}\etiquette[2.1]{I.def2.1}
The projective limit of $G$, denoted by $\varprojlim_I G$, is the
presheaf:
$$
X\longmapsto \Hom_{\Fonct(I,C)}(k_X,G){\quoi}.
$$
When the presheaf $\varprojlim_I G$ is representable, we
also denote by $\varprojlim_I G$ an\pageoriginale object of $C$
representing it. Thus the object $\varprojlim_I G$ is defined only up to
isomorphism. When no confusion can result, we use
the abbreviated notation $\varprojlim G$.
\end{enonce*}

As $G$ varies, $\varprojlim G$ is a functor from the category
$\Fonct(I,C)$ to $\widehat{C}$.

\setcounter{subsection}{1} \setcounter{subsubsection}{0}
\subsubsection{}\etiquette[2.1.1]{I.sec2.1.1}
By symmetry (reversing the direction of the arrows
in $C$), one similarly defines the inductive limit of a functor: it is a
covariant functor on $C$ with values in the category of
$𝒰$-sets. We shall use the notations $\varinjlim_I
G$ or $\varinjlim G$.

Observe that products, fiber products, and kernels are projective
limits. Similarly, sums, amalgamated sums, and cokernels are
%%
inductive limits.

\begin{enonce*}{Definition 2.2}\etiquette[2.2]{I.def2.2}
Let $I$ and $C$ be two categories and $G:I\rightarrow C$ a functor.
The projective limit of $G$ is said to be representable if there exists
a universe $𝒰$ such that:
\begin{enumerate}
\renewcommand{\labelenumi}{\theenumi)}
\item The category $I$ is $𝒰$-small.

\item The category $C$ is a $𝒰$-category.

\item The presheaf $\varprojlim G$, with values in the
  category of $𝒰$-sets, is representable.
\end{enumerate}
\end{enonce*}

It follows from ${\rm n^\circ}$ 1 that the object $\varprojlim G$
representing the presheaf $\varprojlim G$ is independent, up to
isomorphism, of the universe $𝒰$. Note also that there
exists a smallest universe $𝒰'$ having properties
(1) and (2), and that the presheaf $\varprojlim G$ necessarily takes
values in the category of $𝒰'$-sets.

\setcounter{subsection}{2} \setcounter{subsubsection}{0}
\subsubsection{}\etiquette[2.2.1]{I.sec2.2.1}
Let\pageoriginale $I$ and $C$ be two categories. We say that
\emph{the $I$-projective limits in $C$ are
  representable} if, for every functor $G:I\rightarrow C$, the
  projective limit of $G$ is representable. Finally, let $C$ be a
  category and $𝒰$ a universe. We say that the
  $𝒰$-\emph{projective
  limits in $C$ are representable} if, for every
  category $I$ that is $𝒰$-small and every functor $G:I\rightarrow
  C$, the projective limit of $G$ is representable.

\begin{enonce*}{Proposition 2.3}\etiquette[2.3]{I.prop2.3}
Let $C$ be a category and $𝒰$ a universe\sisi{ that is
nonempty}{}. The following assertions are equivalent:
\begin{enumerate}
\renewcommand{\theenumi}{\roman{enumi}}
\renewcommand{\labelenumi}{\theenumi)}
\item The $𝒰$-projective limits in $C$ are representable.

\item Products indexed by a small set are
  representable, and kernels of pairs of arrows are
  representable.
\item Products indexed by a small set are
  representable, and fiber products are representable.
\end{enumerate}
\end{enonce*}

\begin{enonce*}[remark]{Proof}
It suffices to observe that there is an isomorphism, functorial in
$G$,
$$
\varprojlim G=\Ker (\prod_{i\in \ob(I)} %%
G(i)\rightrightarrows \prod_{u\in \rm{F\ell}(I)} G(\but(u))){\quoi},
$$
where the pair of arrows is defined by the morphisms
\begin{align*}
\prod_{i\in \Ob(I)}  G(i) %%
 & \xrightarrow{\pr_{\but(u)}}{} G(\but(u))\\
\prod_{i\in \Ob(I)} G(i) %%
 &\xrightarrow{\pr_{\text{source}(u)}}{} G(\text{source
}(u))\xrightarrow{G(u)} G(\but(u)){\quoi}.
\end{align*}
Moreover, it is clear that $\Ker (X  \displaystyle \Lright^{u}_{v}Y)=
X\prod_{X\times Y} X$, where the two morphisms from $X$ to $X\times Y$
are $\id_X \times u$ and $\id_X\times v$.
\end{enonce*}

\setcounter{subsection}{3} \setcounter{subsubsection}{0}
\subsubsection{}\etiquette[2.3.1]{I.sec2.3.1}
There are of course analogous definitions and assertions for
inductive limits, which we shall not spell out. The same applies to
finite projective\pageoriginale and inductive limits (\ie
those relative to a %%
finite category $I$).


\begin{enonce*}{Corollary 2.3.2}\etiquette[2.3.2]{I.cor2.3.2}
Denote by $𝒰$-$\Ens$ the category of
$𝒰$-sets and by $𝒰$-$\Ab$ the category of
abelian group objects in $𝒰$-$\Ens$. The
$𝒰$-projective and inductive limits in
$𝒰$-$\Ens$ and in $𝒰$-$\Ab$ are representable.
\end{enonce*}

\begin{enonce*}{Proposition 2.3.3}\etiquette[2.3.3]{I.prop2.3.3}
Let $I$ be a small category, $F:I\rightarrow 𝒰$-$\Ens$
a functor, and \linebreak %%
$G\subset\ob(I)$ a small set of objects of $I$ such that, for every
object $X$ of $I$, there exists an object $Y\in G$ and a morphism
$Y\rightarrow X$ (\resp $X\rightarrow Y$). Then $\varprojlim_I F$
(\resp $\varinjlim_I F$) is representable, and
$\card(\varprojlim_I F)\leq \Pi_{Y\in G}\card(F(Y))$ (\resp
$\card(\varinjlim_I F)\leq \Sigma_{Y\in G}\card(F(Y))$).
\end{enonce*}

\begin{enonce*}[remark]{Proof}
One immediately verifies that $\varprojlim_I F$ (\resp $\varinjlim_I
F$) is isomorphic to a subobject (\resp a quotient) of $\Pi_{Y\in
G}F(Y)$ (\resp $\amalg_{Y\in G}F(Y)$); hence the proposition.
\end{enonce*}

\begin{enonce*}{Definition 2.4.1}\etiquette[2.4.1]{I.def2.4.1}
Let $C$ be a category in which finite projective (\resp
inductive) limits are representable, and let $u:C\rightarrow C'$ be a
functor. We say that $u$ is left exact (\resp right exact) if it
``commutes'' with finite projective (\resp inductive) limits.
A functor that is both left and right exact is called an exact functor.
\end{enonce*}

\setcounter{subsection}{4} \setcounter{subsubsection}{1}
\subsubsection{}\etiquette[2.4.2]{I.sec2.4.2}
It follows from \hyperref[I.prop2.3]{2.3} that, for a functor to be exact\sisi{}{ on the left}, it is necessary and %% étiquette !
sufficient that it take the terminal object (= empty product) to the
terminal object, the product of two objects to the product of their
images, and the kernel of a pair of arrows to the kernel of the
image pair; equivalently, that it take the terminal object to the
terminal object and fiber products to fiber products\sisi{. (W}{ (w}e
assume that in $C$ finite $\varprojlim$ are
representable\sisi{.)}{).}

\setcounter{subsection}{5} \setcounter{subsubsection}{-1}
\subsubsection{}\etiquette[2.5.0]{I.sec2.5.0}
Let\pageoriginale $I$, $J$, and $C$ be three categories, and let $G:I\times
J\rightarrow C$ be a functor (\ie a functor from $I$ with values in the
category of functors from $J$ to $C$). Suppose that the projective
limits of the functors
\begin{align*}
G_i:&J\longrightarrow C\qquad i\in \ob(I)\\
    & j\longmapsto G(i\times j)
\end{align*}
are representable, and that the functor
$$
i\longmapsto \varprojlim G_i
$$
admits a representable projective limit. It is then clear that
the functor $G$ admits a representable projective limit and that there is
a canonical isomorphism
$$
\varprojlim_{I\times J} G\\
\xrightarrow{\sim}\varprojlim_I \varprojlim_J G_i{\quoi},
$$
and consequently a canonical isomorphism
$$
\varprojlim_I \varprojlim_J G_i\xrightarrow{\sim}\varprojlim_J
\varprojlim_I G_j{\quoi}.
$$
Hereafter we shall say more briefly that projective limits
commute with projective limits. Similarly, inductive limits
commute with inductive limits. But in general it is not true
that inductive limits commute with projective limits.

\setcounter{subsection}{5}
\begin{enonce*}{Definition 2.5}\etiquette[2.5]{I.def2.5}
Let $C$ be a category with representable fiber products, $I$ a
category, $G:I\rightarrow C$ a functor, $g:G\rightarrow X$ a
morphism from $G$ to an object of $C$ (\ie a morphism from $G$ to
the constant functor associated with $X$), and $m:Y\rightarrow X$ a morphism
in $C$. Let\pageoriginale $\fprod{G}{Y}{X}:I\rightarrow C$ be the
functor
$$
i\longmapsto \fprod{G(i)}{Y}{X} {\quoi}.
$$
We say that the inductive limit of $G$ is universal if, for every
object $X$, every morphism $g:G\rightarrow X$, and every morphism
$m:Y\rightarrow X$,
\begin{enumerate}
\item[a)] the inductive limit of the functor $\fprod{G}{Y}{X}$ is
  representable,

\item[b)] the canonical morphism $\varinjlim (\fprod{G}{Y}{X}
  )\rightarrow \fprod{(\varinjlim G)}{Y}{X}$ is an isomorphism.
\end{enumerate}
\end{enonce*}

\begin{enonce*}{Proposition 2.6}\etiquette[2.6]{I.prop2.6}
Let $𝒰$ be a universe. The inductive limits in
$𝒰$-$\Ens$ that are representable (in particular, the
$𝒰$-inductive limits $\hyperref[I.sec2.2.1]{(2.2.1)}$ in $𝒰$-$\Ens$) %% étiquette !
are universal.
\end{enonce*}
We shall also use another commutation result for
projective and inductive limits, which we now present.

\begin{enonce*}{Definition 2.7}\etiquette[2.7]{I.def2.7}
A category $I$ is pseudo-filtered if it has the following
properties:
\begin{itemize}
\item[PS 1)] Every diagram of the form
\[
\xymatrix@R=.5cm@C=1cm{
       & j\\
i \ar[ru]\ar[dr]     &  \\
       & j'
}
\]
can be embedded in a commutative diagram
\[
\xymatrix@R=.5cm@C=1cm{
  & j \ar[rd]&\\
i\ar[ur]\ar[rd] &   & k\\
  & j'\ar[ru]&
}
\]

\item[PS 2)] Every diagram of the form
\[
\xymatrix@R=.75cm@C=1.5cm@M=.25cm{
i\ar@<2.5pt>[r]^{u}\ar@<-2.5pt>[r]_{v} & j }
\]
%%
  can \pageoriginale be embedded in a diagram
\[
\xymatrix@C=1.75cm@M=.25cm{ i\ar@<2.5pt>[r]^{u}\ar@<-2.5pt>[r]_{v} &
j\ar[r]%%
^{w} & k\\
}
\]
  such that
  $$
  w\circ u=w\circ v{\quoi}.
  $$
\end{itemize}
A category $I$ is said to be filtered if it is pseudo-filtered,
nonempty, and connected, \ie if any two objects of $I$ can
be joined by a sequence of arrows (with no condition imposed
on the directions of the arrows); in the presence of {\rm
PS 2}, this also means that $I\neq ∅$ and that, for any two objects $a$, $b$ of
$I$, there is always an object $c$ of $I$ and arrows $a\sisi{$
$}{\rightarrow}c$ and $b\sisi{$ $}{\rightarrow}c$. We also say
that a category $I$ is cofiltered if $I^\circ$ is filtered.
\end{enonce*}

\begin{enonce*}{Example 2.7.1}\etiquette[2.7.1]{I.exem2.7.1}
If, in $I$, amalgamated sums (\resp sums of two objects)
and cokernels of pairs of parallel arrows are representable, then $I$
is pseudo-filtered (\resp filtered).
\end{enonce*}

\begin{enonce*}{Proposition 2.8}\etiquette[2.8]{I.prop2.8}
Let $𝒰$ be a universe. Filtered $𝒰$-inductive limits
in $𝒰$-$\Ens$ commute with finite projective
limits.
\end{enonce*}

We are immediately reduced to proving that filtered $𝒰$-inductive
limits commute with fiber products. The proof is
left to the reader. One may use the description of the limit
given by the following

\begin{enonce*}{Lemma 2.8.1}\etiquette[2.8.1]{I.lem2.8.1}
Let $I$ be a small filtered category and $i\mapsto X_i$ a
functor from $I$ to $𝒰$-$\Ens$. On the disjoint union
$\amalg_{i\in
 \ob I} X_i$, let $R$ be the relation:
\begin{itemize}
\item[(R)] Two elements $x_i\in X_i$ and $x_j\in X_j$ are
  related if there exists an object $k\in \ob I$ and two morphisms
  $u:i\rightarrow k$ and $v:j\rightarrow k$ such that the images in
  $X_k$ of $x_i$ and $x_j$ under the transition maps $u$ and
  $v$, respectively, are equal.
\end{itemize}
Then:\pageoriginale
\begin{enumerate}
\renewcommand{\labelenumi}{\theenumi)}
\item $R$ is an equivalence relation.

\item The quotient $\amalg_{i\in \ob I}X_i/R$ is canonically isomorphic
  to $\varinjlim_I X_i$.

\item Two elements $x_i\in X_i$ and $x_j\in X_j$ are
  respectively equivalent under $R$ to two elements of the
  same $X_k$.

\item For every $i\in \ob I$, two elements $\alpha$ and $\beta$ of $X_i$ are
  equivalent under $R$ if and only if there is a morphism
  $u:i\rightarrow j$ such that $u(\alpha)=u(\beta)$.
\end{enumerate}
\end{enonce*}

Assertion 1) follows from (PS 1) \hyperref[I.def2.7]{(2.7)}. To prove 2), one %% étiquette !
checks that $\amalg_{i\in \ob I}X_i/R$ has the universal
property of the inductive limit (only (PS 1) is used).
Assertion 3) follows from the fact that $I$ is connected. Assertion 4)
follows from (PS 2).

\begin{enonce*}{Corollary 2.9}\etiquette[2.9]{I.cor2.9}
Let $%%
\gamma$ be a kind of algebraic structure ``defined by finite
projective limits.'' (The reader is asked to give a mathematical meaning
to the preceding phrase. We merely note that the structures of groups,
abelian groups, rings, modules, etc. %%
are structures of this kind.) Denote by $𝒰$-$\gamma$ the
category of $\gamma$-objects in $𝒰$-$\Ens$. The functor that
associates with each object of $𝒰$-$\gamma$ its underlying set
commutes with filtered $𝒰$-inductive limits. Consequently,
filtered $𝒰$-inductive limits in $𝒰$-$\gamma$
commute with finite projective limits.
\end{enonce*}

\begin{enonce*}{Corollary 2.10}\etiquette[2.10]{I.cor2.10}
Pseudo-filtered $𝒰$-inductive limits in $𝒰$-$\Ab$
commute with finite projective limits.
\end{enonce*}

(One reduces to filtered limits by decomposing the indexing category
into connected components.)

We now record\pageoriginale a result that we shall use
constantly: let $C$ and $C'$ be two $𝒰$-categories, and let $u$ and
$v$ be two functors $\xymatrix@C=1.25cm{ C  \ar@<-.6ex>[r]_u & C'
\ar@<-.6ex>[l]_v }$, with $u$ \emph{left adjoint} to $v$. Recall
that this means that there is an isomorphism between the
bifunctors with values in $𝒰$-$\Ens$:
$$
\Hom_{C'}(u(X),X')\xrightarrow{\sim}\Hom_{C}(X,v(X')){\quoi}.
$$

\begin{enonce*}{Proposition 2.11}\etiquette[2.11]{I.prop2.11}
The functor $u$ commutes with representable inductive limits; the
functor $v$ commutes with representable projective limits.
\end{enonce*}

This assertion means that, for every category $I$ and every
functor $G:I\rightarrow C'$ such that the projective limit of $G$ is
representable, the functor $v\circ G$ has a representable
projective limit and there is a canonical isomorphism
$$
v(\varprojlim G)\longrightarrow \varprojlim(v\circ G){\quoi}.
$$
The reader may spell out the assertion concerning the
functor $u$.

Finally, we record a computation of projective limits in the case where the
indexing category has products.

\begin{enonce*}{Proposition 2.12}\etiquette[2.12]{I.prop2.12}
Let $I$ and $C$ be two categories, and $𝒱$ a universe. Assume
that the $𝒱$-projective limits in $C$ are
representable and that there is in $I$ a family of objects
$(i_\alpha)_{\alpha\in A}$, where $A$ is an element of $𝒱$,
such that:
\begin{enumerate}
\renewcommand{\labelenumi}{\theenumi)}
\item the products $i_\alpha\times i_\beta$ are representable
  for every pair $(\alpha,\beta)\in A\times A$,

\item every object of $I$ admits a morphism to at least one of the $i_\alpha$.
\end{enumerate}


Then, for every contravariant functor $F$ from $I$ to $C$,
$$
F:I^\circ\longrightarrow C{\quoi},
$$
the projective limit of $F$ exists, and there is an isomorphism
functorial in $F$:
$$
\varprojlim F\xrightarrow{\sim} \Ker \left(\prod_{\alpha\in
  A} F(i_\alpha)\rightrightarrows \prod_{(\alpha,\beta)\in A\times A}
F(i_\alpha \times i_\beta)\right){\quoi},
$$
where the\pageoriginale two arrows are defined by the projections
of the products $i_\alpha\times i_\beta$ onto their factors.
\end{enonce*}

% END INCLUDED SGA4_English_current_buildable_source_20260906/source/01/drafts/SGA4_I_LINES_411_824_English.texfrag


% BEGIN INCLUDED SGA4_English_current_buildable_source_20260906/source/01/drafts/SGA4_I_LINES_825_1024_English.texfrag
\section{Exactness properties of the category of presheaves}\etiquette[3]{I.sec3}

Let $C$ be a $𝒰$-category and let $\widehat{C}$ be the
category of presheaves on $C$. All the exactness properties of
$\widehat{C}$ follow from the following proposition:

\begin{enonce*}{Proposition 3.1}\etiquette[3.1]{I.prop3.1}
The $𝒰$-projective and inductive limits in
$\widehat{C}$ are representable. For every object $X$
of $C$, the functor on $\widehat{C}$
$$
F\longmapsto F(X)\qquad F\in \widehat{C}
$$
commutes with inductive and projective limits.

In other words, ``inductive and projective limits in
$\widehat{C}$ are computed argument by argument.''
We mention a few corollaries.
\end{enonce*}

\begin{enonce*}{Corollary 3.2}\etiquette[3.2]{I.cor3.2}
Let $\gamma$ be an algebraic structure defined by finite projective
limits. The category of contravariant functors with values in
$𝒰$-$\gamma$ is equivalent to the category of $\gamma$-objects of
$\widehat{C}$.
\end{enonce*}

\begin{enonce*}{Corollary 3.3}\etiquette[3.3]{I.cor3.3}
A morphism of $\widehat{C}$ that is both a monomorphism
and an epimorphism is an isomorphism. A morphism of
$\widehat{C}$ factors uniquely as an epimorphism followed
by a monomorphism. The representable inductive limits in
$\widehat{C}$ are universal. The filtered
$𝒰$-inductive limits in $\widehat{C}$ commute with
finite projective limits. The canonical functor $C\rightarrow
\widehat{C}$ commutes with representable projective
limits.
\end{enonce*}

More\pageoriginale generally, one may say that the category
$\widehat{C}$ inherits all the properties of the category of
$𝒰$-sets that involve inductive and projective limits.

\setcounter{subsection}{4} \setcounter{subsubsection}{-1}
\subsubsection{}\etiquette[3.4.0]{I.sec3.4.0}
Let $F$ be an object of $\widehat{C}$. We denote by
$C/F$ the following category. The objects of $C/F$ are pairs consisting
of an object $X$ of $C$ and a morphism $u$ from $X$ to $F$.
Let $(X,u)$ and $(Y,v)$ be two objects. A morphism from $(X,u)$ to
$(Y,v)$ is a morphism $g$ from $X$ to $Y$ such that the following
diagram is commutative:
\[
\xymatrix@=1.25cm{
X \ar[rr]^{g}\ar[rd]_{u}&   & Y\ar[ld]^{v}\\
            & F &.
}
\]

\begin{enonce*}{Proposition 3.4}\etiquette[3.4]{I.prop3.4}
With the notation of \hyperref[I.sec3.4.0]{(3.4.0)}, the source functor $C/F\rightarrow
\widehat{C}$ admits a representable inductive limit in
$\widehat{C}$. The canonical morphism
$$
\mathop{\underrightarrow{\text{\rm lim source}}}_{C/F}\;
(.)\longrightarrow F
$$
is an isomorphism.
\end{enonce*}


\begin{enonce*}{Corollary 3.5}\etiquette[3.5]{I.cor3.5}
Let $F$ and $H$ be two objects of $\widehat{C}$. There is
a canonical isomorphism
$$
\Hom(F,H)\xrightarrow{\sim}\varprojlim_{(X,u)\in \ob(C/F)}
H(X){\quoi}.
$$
\end{enonce*}

\begin{enonce*}[remark]{Proof}
The corollary follows immediately from \hyperref[I.prop3.4]{3.4}
and \hyperref[I.def1.2]{1.2}.
\end{enonce*}

\begin{enonce*}{Proposition 3.6}\etiquette[3.6]{I.prop3.6}
Let\pageoriginale $C$ be a $𝒰$-category, let $𝒱$ be a
universe containing $𝒰$, and let
$\nobreak{i:\widehat{C}_{𝒰}\rightarrow
\widehat{C}_{𝒱}}$ be the natural inclusion functor of the
corresponding categories of presheaves \hyperref[I.condef1.3]{(1.3)}. The functor $i$
commutes with inductive and projective limits. Moreover, for every
object $F$ of $\widehat{C}_{𝒰}$, every subobject of $i(F)$
is isomorphic to the image under $i$ of a unique subobject of $F$
(thus giving a bijection from the set of subobjects of $F$ to the
set of subobjects of $i(F)$).
\end{enonce*}

\section{Sieves}\etiquette[4]{I.sec4}

\begin{enonce*}{Definition 4.1}\etiquette[4.1]{I.def4.1}
Let $C$ be a category. A \emph{sieve} in the category $C$ is a full
subcategory $D$ of $C$ with the following property: every object of
$C$ from which there exists a morphism to an object of $D$ belongs to
$D$. If $X$ is an object of $C$, the sieves of $X$ are (by an abuse of
language) the sieves in the category $C/X$.
\end{enonce*}

Let $𝒰$ be a universe such that $C$ is a $𝒰$-category.
Let $\widehat{C}$ be the corresponding category of presheaves.
To every sieve of $X$ one associates a subobject of $X$ in
$\widehat{C}$ as follows: to every object $Y$ of $C$, one
assigns the set of morphisms $f:Y\rightarrow X$ such that the object
$(Y,f)$ belongs to the sieve.

\begin{enonce*}{Proposition 4.2}\etiquette[4.2]{I.prop4.2}
The map defined above establishes a %%
bijection between the set of sieves of $X$ and the set of subobjects
of $X$ in $\widehat{C}$.
\end{enonce*}

\begin{enonce*}[remark]{Proof}
We shall only \sisi{show that it}{describe what} is the inverse map. To every
subfunctor $R$ of $X$ we associate the category $C/R$ of objects of
$C$ over $R$ \sisi{(3.4)}{(\hyperref[I.prop3.4]{3.4})}. One immediately
checks that $C/R$ is a sieve of $X$.
\end{enonce*}

\begin{enonce*}[remark]{Remark 4.2.1}\etiquette[4.2.1]{I.rem4.2.1}
One\pageoriginale sees likewise that the sieves of $C$ are in canonical
bijective correspondence with the set of subfunctors of the ``terminal
functor'' on $C$ (the ``terminal'' object of $\widehat{C}$).
\end{enonce*}

\setcounter{subsection}{2}
\subsection{}\etiquette[4.3]{I.sec4.3}
Let $C$ be a $𝒰$-category. By an abuse of language, we shall
also call the subobjects of $X$ in the category $\widehat{C}$
the sieves of $X$. This abuse of language allows us, for every presheaf
$F$ and every sieve $R$ of $X$, to define
$\Hom_{\widehat{C}}(R,F)$ as the set of %%
morphisms from the functor $R$ to $F$. There is moreover a canonical
isomorphism, functorial in $F$ \hyperref[I.cor3.5]{(3.5)}:
$$
\Hom_{\widehat{C}}(R,F)\xrightarrow{\sim}\varprojlim_{C/R}
F(.){\quoi},
$$
which also gives a direct definition\nde{More simply, $C/R$ is identified
with $R$, so that one has the formula $
\Hom_{\widehat{C}}(R,F)\xrightarrow{\sim}\varprojlim_{R}
F{\quoi}, $ where, by an abuse of notation, $F$ denotes the composite of
$F$ and the tautological ``source'' functor $R\rightarrow C/X\rightarrow
C$.}. Likewise, \hyperref[I.prop4.2]{Proposition 4.2} allows us to transfer the usual operations
on functors to sieves. We mention the following:

\subsubsection{}\etiquette[4.3.1]{I.4.3.1}
\emph{Base change.} Let $R$ be a sieve of $X$ and let
$f:Y\rightarrow X$ be a morphism of objects of $C$. The fiber product
$\fprod{R}{Y}{X}$ is a sieve of $Y$, called the \emph{sieve
deduced
  from $R$ by base change.} The corresponding subcategory of $C/Y$
is the inverse image of the subcategory of $C/X$ defined by $R$ under
the canonical functor $C/Y\rightarrow C/X$ defined by $f$.

\subsubsection{}\etiquette[4.3.2]{I.sec4.3.2}
\emph{Order relation, intersection, union.} The inclusion relation on
the subfunctors of $X$ is an order relation. The union and intersection
of a family of sieves indexed by an arbitrary set may be defined as the
least upper bound and the greatest lower bound of the family of
corresponding subpresheaves.

\subsubsection{}\etiquette[4.3.3]{I.sec4.3.3}
\emph{Image, generated sieve.}\pageoriginale Let
$(F_\alpha)_{\alpha\in A}$ be a family of presheaves and, for every
$\alpha\in A$, let $f_\alpha:F_\alpha\rightarrow X$ be a morphism,
where $X$ is an object of $C$. The \emph{image} of this family of
morphisms is the union of the images of the $f_\alpha$. The image of
this family is therefore a sieve of $X$. In particular, if all the
$F_\alpha$ are objects of $C$, the image sieve will be called the
sieve generated by the morphisms $f_\alpha$. The category $C/R$ is
the full subcategory of $C/X$ formed by the objects $X'\rightarrow X$
over $X$ for which there exists an $X$-morphism from $X'$ to one of the
$F_{\alpha}$.

The reader may, as an exercise, translate the relations and operations
defined here on subfunctors into the language of the categories $C/R$.
One will then observe that these relations and operations do not depend
on the universe $𝒰$ for which $C$ is a $𝒰$-category, and that
they are therefore defined for every category, without having to specify
the universe to which the morphism sets belong; this could moreover have
been foreseen \textit{a priori} from \hyperref[I.prop3.6]{3.6}.

% END INCLUDED SGA4_English_current_buildable_source_20260906/source/01/drafts/SGA4_I_LINES_825_1024_English.texfrag


% BEGIN INCLUDED SGA4_English_current_buildable_source_20260906/source/01/drafts/SGA4_I_LINES_1025_1161_English.texfrag
\section{Functoriality of categories of presheaves}\etiquette[5]{I.sec5}

\setcounter{subsection}{-1}
\subsection{}\etiquette[5.0]{I.sec5.0}
Let $C$, $C'$, and $D$ be three categories, and let $u:C\rightarrow C'$ be a
functor. We denote by $u^\ast$ the functor \sisi{$$
u^\ast:\Hhom({C'}^\circ,D)\longrightarrow \Hhom(C^\circ, D),
G\longmapsto G\circ u{\quoi},
$$}{$$u^\ast:\left\{\begin{array}{ccc}
                      \Hhom({C'}^\circ,D) & \longrightarrow & \Hhom(C^\circ, D) \\
                      G & \longmapsto & G\circ u{\quoi}
                    \end{array}
\right.$$} obtained by composition with the functor $u$. The functor
$u^\ast$ commutes with inductive and projective limits.

\begin{enonce*}{Proposition 5.1}\etiquette[5.1]{I.prop5.1}
Suppose\pageoriginale that $C$ is small and that the
$𝒰$-inductive (
\resp projective) limits in $D$ are representable. The functor $u^\ast$
has a left adjoint $u_{!}$ (\resp a right adjoint $u_\ast$). Thus there
is an isomorphism
\begin{gather*}
\Hom_{\Hhom(C^\circ,D)}(F,u^\ast G)\xrightarrow{\sim}
\Hom_{\Hhom({C'}^\circ,D)}(u_{!}F,G)\\
(\mbox{\resp  } \Hom_{\Hhom(C^\circ,D)}(u^\ast G,F)\simeq
\Hom_{\Hhom({C'}^\circ,D)}(G,u_\ast F){\quoi}).
\end{gather*}
\end{enonce*}

\begin{enonce*}[remark]{Proof}
We shall give only the proof of the existence of the left adjoint functor.
The \resp part of the proposition will then follow formally from the
isomorphisms
$$
\Hhom(C^\circ,
D)\xrightarrow{\sim}\Hhom(C,D^\circ)^\circ\xrightarrow{\sim}\Hhom((C^\circ)^\circ,D^\circ)^\circ{\quoi}.
$$
\end{enonce*}

Let $Y$ be an object of $C'$. Denote by $I^Y_u$ the following category.
The objects of $I^Y_u$ are pairs $(X,m)$, where $X$ is an object of $C$
and $m$ is a morphism $Y\rightarrow u(X)$. Let $(X,m)$ and $(X',m')$
be two objects of $I^Y_u$. A morphism from $(X,m)$ to $(X',m')$ is a
morphism $\xi:X\rightarrow X'$ such that $m'=u(\xi)m$. Composition of
morphisms is defined in the evident way.

Let $f:Y\rightarrow Y'$ be a morphism of $C'$. By composition, the
morphism $f$ defines a functor $I^f_u:I^{Y'}_u\rightarrow I^Y_u$.

There is moreover a functor $\pr_Y:I^Y_u\rightarrow C$ that assigns the
object $X$ to the object $(X,m)$. Notice that the diagram
\[\vcenter{
\xymatrix@=1.5cm{
       &I^{Y'}_u \ar[rr]^{I^f_u} \ar[rd]_{\pr_{Y'}}& & I^Y_u \ar[ld]^{\pr_Y}\\
       & &   C   & \\
}}\leqno{(*)}
\]
is then commutative.

Now\pageoriginale let $F$ be a presheaf on $C$, and set
\hypertarget{sga4.exp.I.equation.5.1.1}{}$$
\displaylines{\rlap{(5.1.1)}\hfill u_{!}
  F(Y)=\varinjlim_{I^Y_u} F\sisi{}{\circ}\pr_Y(.){\quoi}.\hfill}
$$
The commutativity of diagram ($^\ast$) and the functoriality of the
inductive limit make $u_{!}F$ a presheaf on $C'$. Let us show that the %%
functor $u_{!}$ is a left adjoint of the functor $u^\ast$. For this, we
show that for every presheaf $G$ on $C'$, there is a functorial
isomorphism
$$
\Hom(u_{!}F,G)\xrightarrow{\sim}\Hom(F,u^\ast G){\quoi}.
$$
Let $\xi\in \Hom(u_{!}F,G)$. Thus, for every object $X$ of $C$, there
is a morphism
$$
\xi_X:u_{!}F(u(X))\longrightarrow G(u(X)){\quoi}.
$$
But $(u(X),\id_{u(X)})$ is an object of $I^{u(X)}_u$, and by the
definition of the inductive limit there is a canonical morphism
$$
F(X)\longrightarrow u_{!}F(u(X)){\quoi}.
$$
For every object $X$ of $C$, we therefore obtain a morphism
$$
\eta_X:F(X)\longrightarrow G(u(X))
$$
which is visibly functorial in $X$. Hence a morphism
$$
\eta:F\longrightarrow u^\ast G{\quoi}.
$$

Conversely, let $\eta\in \Hom(F,u^\ast G)$. For every object $Y$ of
$C'$, this gives a morphism of functors
$$
\eta_Y:F\sisi{}{\circ}\pr_Y\longrightarrow (u^\ast G)\circ\pr_Y{\quoi},
$$
and hence, by composition with the evident morphism from the functor
$(u^\ast G)\circ\pr_Y$ to the constant functor $G(Y)$, a morphism
$$
\xi_Y:u_{!}F(Y)\longrightarrow G(Y)
$$
which is functorial in $Y$. Hence a morphism
$$
\xi:u_{!}F\longrightarrow G{\quoi}.
$$

The\pageoriginale reader will check that the two maps thus defined are
inverse to one another, thereby completing the proof.

\begin{enonce*}{Proposition 5.2}\etiquette[5.2]{I.prop5.2}
Suppose that in $D$ the $𝒰$-inductive limits are representable,
finite projective limits are representable, and filtered
$𝒰$-inductive limits commute with finite projective limits. Suppose
also that finite projective limits are representable in $C$ and that
the functor $u$ is left exact $(\hyperref[I.cor2.3.2]{2.3.2})$. Then
finite projective limits are representable in $\Hhom(C^\circ,D)$ and
in $\Hhom({C'}^\circ,D)$, and the functor $u_{!}$ is left exact.
\end{enonce*}

\begin{enonce*}[remark]{Proof}
The first assertion is trivial. Let us prove the second. By the proof of
\hyperref[I.prop5.1]{5.1}, for every presheaf $F$ on $C$ and every object $Y$ of $C'$, one has
$$
u_{!}F(Y)\xrightarrow{\sim} \varinjlim_{I^Y_u}
F\sisi{}{\circ}\pr_Y{\quoi}.
$$
 It therefore suffices to show that the category $(I^Y_u)^\circ$
 satisfies axioms (PS 1) and (PS 2) \hyperref[I.def2.7]{(2.7)}, and that this category is
 connected. The verification is left to the reader.
\end{enonce*}

% END INCLUDED SGA4_English_current_buildable_source_20260906/source/01/drafts/SGA4_I_LINES_1025_1161_English.texfrag


% BEGIN INCLUDED SGA4_English_current_buildable_source_20260906/source/01/drafts/SGA4_I_LINES_1162_1346_English.texfrag
\setcounter{subsection}{2}
\subsection{}\etiquette[5.3]{I.sec5.3}
Specializing these results to the case where $D$ is the category of
$𝒰$-sets, one obtains a sequence of three functors:
$$
u_{!}, u^\ast, u_\ast {\quoi},
$$
which is a ``sequence of adjoint functors'' in the sense that
for any two consecutive functors in the sequence, the one on the right is
right adjoint to the other. Their\pageoriginale essential
properties are summarized in the following proposition:

\begin{enonce*}{Proposition 5.4}\etiquette[5.4]{I.prop5.4}
Let $C$ be a small category, $C'$ a $𝒰$-category, and
$u:C\rightarrow C'$ a functor.
\begin{enumerate}
\renewcommand{\labelenumi}{\theenumi)}
\item The functor $u^\ast:\widehat{C'}\rightarrow
  \widehat{C}$ commutes with inductive and
  projective limits.

\item The functor $u_\ast:\widehat{C}\rightarrow
  \widehat{C'}$ commutes with projective limits. For
  every presheaf $F$ on $C$ and every object $Y$ of $C'$, one has:
  $$
  u_\ast F(Y)\xrightarrow{\sim}
  \Hom_{\widehat{C}}(u^\ast(Y),F){\quoi}.
  $$

\hypertarget{sga4.exp.I.proposition.5.4.item.3}{}\item The functor $u_{!}:\widehat{C}\rightarrow
  \widehat{C'}$ commutes with inductive limits. The
  functor $u_{!}$ is defined only up to isomorphism, but it
  can always be chosen so that the diagram
\[
\xymatrix@C=2.75cm@R=1.5cm{
  C \ar[r]^{u}\ar[d]_{h}& C' \ar[d]^{h^\prime}\\
C{\widehat{\phantom{\;i}}}\ar[r]^{u_{!}}&C'{\widehat{\phantom{\;i}}}}
\]
($h$ and $h'$ are the canonical inclusion functors) is
commutative.

For every presheaf $F$ on $C$, one has:
$$
u_{!}F\xrightarrow{\sim} \varinjlim_{C/F} h'\sisi{}{\circ}u{\quoi}.
$$

\item If finite projective limits are representable in
  $C$ and if $u$ is left exact $\hyperref[I.cor2.3.2]{(2.3.2)}$, the functor $u_{!}$
  is left exact.
\end{enumerate}
\end{enonce*}

\begin{enonce*}[remark]{Proof}
Assertion (1) is trivial. Assertion (2) follows from the fact that
$u_\ast$ is a right adjoint functor, by (\hyperref[I.prop2.11]{2.11}) and
(\hyperref[I.prop1.4]{1.4}). The same is true of assertion (3), but one
also applies (\hyperref[I.prop3.4]{3.4}). Finally, assertion (4) is
simply \sisi{5.2}({\hyperref[I.prop5.2]{5.2}}), which can be applied thanks to
(\hyperref[I.def2.7]{2.7}).
\end{enonce*}

\begin{enonce*}{Proposition 5.5}\etiquette[5.5]{I.prop5.5}
Let\pageoriginale $C$ and $C'$ be two small categories and let
$\xymatrix@C=1cm{
 C \ar@<-2.5pt>[r]_{u} & C'\ar@<-2.5pt>[l]_{v}
}$ be a pair of functors, where $v$ is left adjoint to $u$. There
are then isomorphisms, compatible with the adjunction
isomorphisms:
\[
\xymatrix@R=.75cm{
v^{\ast} \ar[rr]^{\sim}&    & u_!      &\\
v_{\ast} \ar[rr]^{\sim}&    & u^{\ast} & . }
\]
\end{enonce*}

\begin{enonce*}[remark]{Proof}
It suffices to exhibit an isomorphism $v^\ast\xrightarrow{\sim}
u_{!}$; the other isomorphism follows from it by adjunction. Let $F$
be a presheaf on $C$ and let $Y$ be an object of $C'$. Then:
$$
v^\ast F(Y)\xrightarrow{\sim} \Hom(v(Y),F){\quoi}.
$$
Next, using (\hyperref[I.prop3.4]{3.4}):
$$
\Hom(v(Y),F)\xrightarrow{\sim} \varinjlim_{C/F} \Hom(v(Y),.){\quoi}.
$$
But $v$ is left adjoint to $u$, and hence:
$$
\varinjlim_{C/F}\Hom(v(Y),.)\xrightarrow{\sim}\varinjlim_{C/F}\Hom(Y,u(.)){\quoi}.
$$
Now using (\hyperref[I.prop5.4]{5.4}.3)), one obtains:
$$
\varinjlim_{C/F}\Hom(Y,u(.))\xrightarrow{\sim}
\Hom(Y,u_{!}F)\xrightarrow{\sim} u_{!}F(Y){\quoi}.
$$
Thus, for every object $Y$ of $C'$, we have determined an isomorphism
$v^\ast F(Y)\xrightarrow{\sim} u_{!}F(Y)$ that is clearly
functorial in $Y$ and in $F$, \cqfd
\end{enonce*}

\begin{enonce*}{Corollary 5.5.1}\etiquette[5.5.1]{I.cor5.5.1}
Let $u:C\rightarrow C'$ be a functor that admits a left adjoint.
The functor $u_{!}:\widehat{C}\rightarrow
\widehat{C'}$ commutes with projective limits
(recall that it commutes with inductive limits by $\hyperlink{sga4.exp.I.proposition.5.4.item.3}{(1.4.3)}$).
\end{enonce*}

\begin{enonce*}[remark]{Remark 5.5.2}\etiquette[5.5.2]{I.rem5.5.2}
One\pageoriginale thus obtains a ``sequence of \emph{four}
adjoint functors'' (\cf \hyperref[I.sec5.3]{5.3}):
$$
v_{!}, v^\ast=u_{!}, v_\ast = u^\ast, u_\ast{\quoi},
$$
of which the first three (\resp  the last three) therefore commute with
$\varinjlim$ (\resp  $\varprojlim$).
\end{enonce*}

\begin{enonce*}{Proposition 5.6}\etiquette[5.6]{I.prop5.6}
The hypotheses are those of $\hyperref[I.prop5.4]{5.4}$. The following
conditions are equivalent:
\begin{enumerate}
\renewcommand{\theenumi}{\roman{enumi}}
\renewcommand{\labelenumi}{\theenumi)}
\item The functor $u$ is fully faithful.

\item The functor $u_{!}$ is fully faithful.

\item The adjunction morphism
  $\id_{\widehat{C}}\rightarrow u^\ast u_{!}$ is an
  isomorphism.

\item The functor $u_\ast$ is fully faithful.

\item The adjunction morphism $u^\ast u_\ast\rightarrow
  \id_{\widehat{C}}$ is an isomorphism.
\end{enumerate}
\end{enonce*}

\begin{enonce*}[remark]{Proof}
It is clear that ii) $\Leftrightarrow$ iii) and iv) $\Leftrightarrow$
v) (general properties of adjoint functors), and that ii)
$\Rightarrow$ i) (\hyperref[I.prop5.4]{5.4}.3)). Let us show that i) $\Rightarrow$
iii). The functors $\id_{\widehat{C}}$, $u^\ast$, and
$u_{!}$ commute with inductive limits. By \hyperref[I.prop3.4]{(3.4)}, it therefore suffices
to prove that $H\rightarrow u^\ast u_{!} H$ is an
isomorphism when $H$ is representable, which is evident.
Let us show that %%
iii) is equivalent to %%
v). For every object $H$ (\resp  $K$) of $\widehat{C}$,
denote by $\Phi(H):H\rightarrow u^\ast u_{!} H$ (\resp  by
$\Psi(K):u^\ast u_\ast\rightarrow K$) the adjunction morphism. We
then have a commutative diagram:
\[
\xymatrix@C=3.5cm@L=6pt{
\Hom_{\widehat{C}}(H,u^{\ast}u_{\ast}
K)\ar@{-}[d]_{\wr}\ar[r]^{\quad \Hom(H,\Psi(K))}   & \Hom_{\widehat{C}}(H,K)\\
\Hom_{\widehat{C}}(u_{!}H,u_{\ast}K)\ar@{-}[d]_{\wr} &  \\
\Hom_{\widehat{C}}(u^{\ast}u_{!}H,K)\ar[uru]_{\quad
   \Hom(\Phi(H),K)} & .
}
\]

It\pageoriginale follows that $\Phi(H)$ is an isomorphism for every $H$
if and only if $\Psi(𝒦)$ is an isomorphism for every
$K$, \cqfd
\end{enonce*}

\begin{enonce*}[remark]{Remark 5.7}\etiquette[5.7]{I.rem5.7}
\begin{enumerate}
\renewcommand{\theenumi}{\alph{enumi}}
\renewcommand{\labelenumi}{\theenumi)}
\item The equivalences ii) $\Leftrightarrow$ iii)
  $\Leftrightarrow$ iv) $\Leftrightarrow$ v) are general results
  on adjoint functors.

\item The explicit form of $u_{!}$ \resp  $u_\ast$ given in the
  proof of \hyperref[I.prop1.4]{1.4} immediately shows that, under the general
  hypotheses of \hyperref[I.def1.1]{1.1}, if $u$ is fully faithful, then the
  adjunction morphism $\id\rightarrow u^\ast u_{!}$ (\resp  $u^\ast
  u_\ast\rightarrow \id$) is an isomorphism, \ie that $u_{!}$
  (\resp  $u_\ast$) is fully faithful.
\end{enumerate}
\end{enonce*}

\setcounter{subsection}{8} \setcounter{subsubsection}{-1}

% END INCLUDED SGA4_English_current_buildable_source_20260906/source/01/drafts/SGA4_I_LINES_1162_1346_English.texfrag


% BEGIN INCLUDED SGA4_English_current_buildable_source_20260906/source/01/drafts/SGA4_I_LINES_1347_1496_English.texfrag
\subsubsection{}\etiquette[5.8.0]{I.sec5.8.0}
Let $\gamma$ be a type of algebraic structure defined by finite
projective limits, let $𝒰$-$\gamma$-$\Ens$ be the category of
$\gamma$-objects of $𝒰$-$\Ens$, and let esj:
$𝒰$-$\gamma$-$\Ens\rightarrow 𝒰$-$\Ens$ be the underlying-set functor
(for simplicity, we suppose that the type of structure under consideration
has only one underlying set). Let $C$ be a category.
Composition with $\esj$ gives a functor denoted by
$$
\mbox{esj}^{\widehat{\phantom{,}}}:\Hhom(C^\circ,
𝒰\text{-}\gamma\text{-}\Ens)\longrightarrow
\widehat{C}{\quoi}.
$$
Since projective limits in $\widehat{C}$ are computed argument by
argument, the functor $\mbox{esj}^{\widehat{\phantom{,}}}$ factors as
an equivalence
\hypertarget{sga4.exp.I.equation.5.8.1}{}$$
\displaylines{\rlap{(5.8.1)}\hfill \Hhom(C^\circ,
  𝒰\text{-}\gamma\text{-}\Ens)\xrightarrow{\approx}
  \widehat{C}_\gamma{\quoi},\hfill}
$$
where $\widehat{C}_\gamma$ is the category of $\gamma$-objects of
$\widehat{C}$, followed by a functor, still denoted by
$$
\mbox{esj}^{\widehat{\phantom{,}}}:\widehat{C}_\gamma\longrightarrow
\widehat{C}{\quoi},
$$
and called the ``\emph{underlying presheaf of sets}'' functor.

\setcounter{subsubsection}{1}
\subsubsection{}\etiquette[5.8.2]{I.sec5.8.2}
Suppose\pageoriginale that the functor
$\esj:𝒰$-$\gamma$-$\Ens\rightarrow 𝒰$-$\Ens$ has a left adjoint
$\Lib:𝒰$-$\Ens\rightarrow 𝒰$-$\gamma$-$\Ens$
\footnote[1]{The functor $\Lib$ is the ``free generated
$\gamma$-object'' functor. Examples: free group, free abelian group,
free $A$-module, \etc} (in fact, one can show that this condition is
always satisfied). Composition with $\Lib$ gives a functor
$$
\Lib^{\widehat{\phantom{,}}}:\widehat{C}\longrightarrow
\Hhom(C^\circ, 𝒰\text{-}\gamma\text{-}\Ens)
$$
and, after composition with the equivalence \hyperlink{sga4.exp.I.equation.5.8.1}{(5.8.1)}, a functor still
denoted by
$$
\Lib^{\widehat{\phantom{,}}}:\widehat{C}\longrightarrow
\widehat{C}_\gamma
$$
and called the ``freely generated presheaf of $\gamma$-objects'' functor.
The functor $\Lib^{\widehat{\phantom{,}}}$ is left adjoint to the functor
$\esj^{\widehat{\phantom{,}}}$.

\begin{enonce*}{Proposition 5.8.3}\etiquette[5.8.3]{I.prop5.8.3}
Let $\gamma$ be a type of algebraic structure defined by finite
projective limits such that, in the category of $\gamma$-objects of
$𝒰$-$\Ens$, the $𝒰$-inductive limits are representable
\footnote[2]{One can show that this condition is always satisfied.};
let $C$ be a category belonging to $𝒰$, let $C'$ be a
$𝒰$-category, and let $u:C\rightarrow C'$ be a functor. Denote by
$\widehat{C}_\gamma$ (\resp $\widehat{C'}_\gamma$) the category of
$\gamma$-objects of $\widehat{C}$ (\resp $\widehat{C'}$), and by
${u^\ast}^\gamma$ the functor on $\gamma$-objects induced by the functor
$u^\ast$. It follows from \hyperref[I.prop5.1]{5.1} and the equivalence
\hyperlink{sga4.exp.I.equation.5.8.1}{5.8.1} that the functor ${u^\ast}^\gamma$ has a left (\resp right)
adjoint. This functor is denoted by $u_{!_\gamma}$ (\resp
$u_{\ast_\gamma}$).
\begin{enumerate}
\item The functor ${u^\ast}^\gamma$ commutes with inductive and
  projective limits. The diagram
\[\vcenter{
\xymatrix@C=3cm@R=1.5cm{
 \widehat{C'}_\gamma \ar[r]^{u^{\ast^{\gamma}}} \ar[d]_{\esj'\widehat{\phantom{a}}}&
\widehat{C}_\gamma  \ar[d]_{\esj^{\widehat{\phantom{a}}}}\\
 \widehat{C'} \ar[r]^{u_{\ast}}& \widehat{C}
}}\leqno{(*)}
\]
is\pageoriginale commutative ($\esj'^{\widehat{\phantom{,}}}$ and
$\esj^{\widehat{\phantom{,}}}$ denote the ``underlying-set'' functors).

\item The functor $u_{\ast_\gamma}$ commutes with projective limits. The
  diagram
\[\vcenter{
\xymatrix@C=3cm@R=1.5cm{
 \widehat{C}_\gamma \ar[r]^{u_{\ast_\gamma}} \ar[d]_{\esj\widehat{\phantom{a}}}&
  \widehat{C'}_\gamma\ar[d]_{\esj'\widehat{\phantom{a}}}\\
\widehat{C} \ar[r]^{u_{\ast}}&
\widehat{C'} }}\leqno{({*}{*})}
\]
commutes up to isomorphism.

\item The functor $u_{!_\gamma}$ commutes with inductive limits. Suppose
  that $\esj^{\widehat{\phantom{,}}}$
  (\resp ${\esj'}^{\widehat{\phantom{,}}}$) has a left adjoint
  $\Lib^{\widehat{\phantom{,}}}$
  (\resp ${\Lib'}^{\widehat{\phantom{,}}}$) $\hyperref[I.sec5.8.2]{(5.8.2)}$. The diagram
\[\vcenter{
\xymatrix@C=3.5cm@R=1.75cm{
 \widehat{C} \ar[r]^{u_{!}} \ar[d]_{\Lib\widehat{\phantom{a}}}&
  \widehat{C'} \ar[d]_{\Lib'\widehat{\phantom{a}}}\\
\widehat{C}_{\gamma} \ar[r]^{u_{!\gamma}} &
\widehat{C'}_{\gamma} }}\leqno{({*}{*}{*})}
\]
commutes up to isomorphism.
\end{enumerate}

Suppose that $u_{!}$ commutes with finite projective limits
(\hyperref[I.prop5.4]{5.4}) and (\hyperref[I.prop5.6]{5.6}). Then the diagram
\[\vcenter{
\xymatrix@C=3.5cm@R=1.5cm{
 \widehat{C}_\gamma \ar[r]^{u_{!\gamma}} \ar[d]_{\esj\widehat{\phantom{a}}}&
  \widehat{C'}_\gamma\ar[d]_{\esj'\widehat{\phantom{a}}}\\
\widehat{C} \ar[r]^{u_{!}}& \widehat{C'}
}}\leqno{({*}{*}{*}{*})}
\]
commutes up to isomorphism, and $u_{!\gamma}$ commutes with finite
projective limits.
\end{enonce*}

\begin{enonce*}[remark]{Proof}
Assertion (1) is evident. So is assertion (2), since $u_\ast$ is a right
adjoint and therefore (\hyperref[I.prop2.11]{2.11}) commutes with projective
limits and, in particular, with finite projective limits; this gives the
commutativity of diagram $({\ast}{\ast})$. The commutativity of diagram
$({\ast}{\ast}{\ast})$ follows from the uniqueness, up\pageoriginale to
isomorphism, of the left adjoint functor, and finally the commutativity of
diagram $({\ast}{\ast}{\ast}{\ast})$ follows immediately from the fact that
$u_{!}$ commutes with finite projective limits.
\end{enonce*}

\begin{enonce*}[remark]{Notation 5.9}\etiquette[5.9]{I.not5.9}
By an abuse of notation, the functors ${u^\ast}^\gamma$ and
$u_{\ast\gamma}$ will often be denoted by $u^\ast$ and $u_\ast$. In view
of the commutativity of diagrams $(\ast)$ and $({\ast}{\ast})$, this
cannot cause confusion. On the other hand, when $u_{!}$ does not commute
with finite projective limits, diagram $({\ast}{\ast}{\ast}{\ast})$ does
not commute up to isomorphism, and the notations $u_{!}$ and
$u_{!\gamma}$ must be used to avoid confusion.
\end{enonce*}

\setcounter{subsection}{9}

% END INCLUDED SGA4_English_current_buildable_source_20260906/source/01/drafts/SGA4_I_LINES_1347_1496_English.texfrag


% BEGIN INCLUDED SGA4_English_current_buildable_source_20260906/source/01/drafts/SGA4_I_LINES_1497_1738_English.texfrag
\subsection{}\etiquette[5.10]{I.sec5.10}
Let $C$ be a small category. For every object $X$ of
$\widehat{C}$, let $C/X$ denote the category of arrows with target
$X$ and source an object of $C$. The source functor defines a functor
$j_X:C/X\rightarrow C$. Let $m:Y\rightarrow X$ be a morphism of
$\widehat{C}$. Composition of morphisms defines a functor
$j_m:C/Y\rightarrow C/X$. The diagram
\[
\xymatrix@=1.5cm{
C/Y \ar[rr]^{j_m} \ar[rd]_{j_{Y}}&   & C/X\ar[ld]^{j_{X}}\\
    & C &
}
\]
is commutative. It follows from \hyperref[I.prop5.1]{5.1} that, for every
object $F$ of $(C/X)^{\widehat{\phantom{,}}}$ and every object $Y$ of
$C$, one has
\hypertarget{sga4.exp.I.equation.5.10.1}{}$$
\displaylines{\rlap{(5.10.1)}\hfill j_{X!}F(Y)=\coprod_{u\in
    \Hom_{\widehat{C}}(Y,X)} F(u){\quoi}.\hfill}
$$
Formula \hyperlink{sga4.exp.I.equation.5.10.1}{(5.10.1)} defines $j_{X!}$ when $C$ is a $𝒰$-category, and
one checks that the functor $j_{X!}$ thus defined is always right
adjoint\pageoriginale to the functor
$j^\ast_X:\widehat{C}\rightarrow
(C/X)^{\widehat{\phantom{,}}}$ (\hyperref[I.sec5.0]{5.0}).

\begin{enonce*}{Proposition 5.11}\etiquette[5.11]{I.prop5.11}
Let $C$ be a $𝒰$-category and let $X$ be a presheaf on $C$.
\begin{enumerate}
\renewcommand{\labelenumi}{\theenumi)}
\item The functor
$$
j_{X!}:(C/X)^{\widehat{\phantom{,}}}\longrightarrow
\widehat{C}
$$
factors through the category $\widehat{C}/X$:
$$
(C/X)^{\widehat{\phantom{,}}}\xrightarrow{e_X}
\widehat{C}/X\longrightarrow
\widehat{C}{\quoi}.
$$
The functor $e_X$ is an equivalence of categories.

\item The functor $e_{X}\circ j^\ast_X:\widehat{C}
\to  \widehat{C}/X$ is canonically isomorphic to the functor
  $H\mapsto (H\times X \xrightarrow{\pr_2} X)$.
\end{enumerate}
\end{enonce*}

\begin{enonce*}[remark]{Proof}
\begin{enumerate}
\renewcommand{\labelenumi}{\theenumi)}
\item Let $f$ be the terminal object of $(C/X)^{\widehat{\phantom{,}}}$.
There is a canonical isomorphism $f\xrightarrow{\sim}
\varinjlim_{Y\in \ob C/X}Y$, and hence
$j_{X!}(f)=\varinjlim_{Y\in \ob C/X}j_X(Y)$ \hyperref[I.prop5.4]{(5.4)}. But
$j_{X!}(f)\simeq X$, which gives the factorization. To show that $e_X$
is an equivalence, it is enough to exhibit a quasi-inverse functor. To
every object $H\rightarrow X$ of $\widehat{C}/X$, associate the
presheaf on $C/X$
$$
(Y\rightarrow X)\longmapsto
  \Hom_{\widehat{C}/X}((Y\rightarrow X), (H\rightarrow
  X)){\quoi}.
$$

\item The functor $e_X\circ j^\ast_X$ is right adjoint to the forgetful
functor, and therefore the functor $e_X\circ j^\ast$ is canonically
isomorphic to the functor $H\mapsto (H\times
  H\xrightarrow{\pr_2}X)$.
\end{enumerate}
\end{enonce*}

\setcounter{subsection}{11}
\subsection{}\etiquette[5.12]{I.sec5.12}
Let $m:Y\rightarrow X$ be a morphism of $\widehat{C}$.
By \hyperref[I.prop5.11]{(5.11)}, the morphism $m$ is canonically isomorphic to the image under
$e_X$ of an object of $(C/X)^{\widehat{\phantom{,}}}$, which we denote
by $[m]$. By restriction to the subcategories, the functor $e_X$ defines
an equivalence
$$
(C/X)/[m]\xrightarrow{e_m}C/Y{\quoi}.
$$
The\pageoriginale diagram
\[
\xymatrix{
  (C/X)/[m]\ar[rr]^{e_m} \ar[rd]_{j_{[m]}} &       & C/Y \ar[ld]^{j_m}\\
          & C/X   &
}
\]
commutes up to canonical isomorphism.

\subsection{}\etiquette[5.13]{I.sec5.13}

We record a result that will be useful in Exp.~VI. Let
$u:C\rightarrow C'$ be a functor between small categories. For every
object $H$ of $\widehat{C}$, denote by $u/H:C/H\rightarrow C'/u_{!}H$
the functor that assigns to every morphism $m:X\rightarrow H$ the
morphism $u_{!}X\xrightarrow{u_{!}m}u_{!}H$ (we know
(\hyperref[I.prop5.4]{5.4}) that we may always put $u_{!}X=uX$). The
following diagram is commutative:
\hypertarget{sga4.exp.I.equation.5.13.1}{}\[\vcenter{
\xymatrix@C=3cm@R=1.75cm{
C/H \ar[r]^{u/H}\ar[d]_{J_H} & C'/u_{!}H \ar[d]_{j_{u_{!}H}} \\
C\ar[r]^{u}   & C'\quad. }}\leqno{(5.13.1)}
\]
We therefore have a diagram commutative up to isomorphism:
%\[\vcenter{
%\xymatrix@C=2.5cm@R=1.5cm{ (C/H){\widehat{\phantom{a}}}
%\ar[r]^{(u/H)_{!}} \ar[d]_{(j_H)_!}&
%  (C'/u^·H){\widehat{\phantom{a}}} \ar[d]_{(j_{u!H})_{!}}\\
%  C{\widehat{\phantom{\;a}}} \ar[r]^{u_!}& C'{\widehat{\phantom{\;a}}}\quad,
%}}\leqno{(5.13.2)}
%\]
and since $(u/H)_{!}$ takes the terminal object of
$(C/H)^{\widehat{\phantom{,}}}$ to the terminal object of
$(C'/u_{!}H)^{\widehat{\phantom{,}}}$, the diagram
\hypertarget{sga4.exp.I.equation.5.13.3}{}\[\vcenter{
\xymatrix@C=3.5cm@R=1.5cm{ (C/H)\widehat{\phantom{a}}
\ar[r]^{(u/H)_!} \ar[d]_{e_H}&
  (C'/u_!H)\widehat{\phantom{a}} \ar[d]_{e_{u_{!}H}}\\
C\widehat{\phantom{a}}/H  \ar[r]^{u_{!}/H} &
C'\widehat{\phantom{a}}/u_{!}H }}\leqno{(5.13.3)}
\]
commutes\pageoriginale up to isomorphism \hyperref[I.def6.1]{(6.1)}.

\begin{enonce*}{Proposition 5.14}\etiquette[5.14]{I.prop5.14}
Let $u:C\rightarrow C'$ be a functor between small categories. Suppose
that $u$ has the following property:

{\rm(PPF)} For every object $X$ of $C$, the functor
$u/X:C/X\rightarrow C'/uX$ is fully faithful.

Then:
\begin{enumerate}
\renewcommand{\labelenumi}{\theenumi)}
\item Let $f$ be the terminal object of $\widehat{C}$. The functor $u$
factors as
$$
C=C/f\xrightarrow{u/f}C'/u_{!}f\xrightarrow{j_{u_{!}f}} C'{\quoi}.
$$
The functor $u/f$ is fully faithful.

\item The functor $u_{!}:\widehat{C}\rightarrow
  \widehat{C'}$ factors as
$$
\widehat{C}=(C/f)^{\widehat{\phantom{,}}}\xrightarrow{(u/f)_{!}}
(C'/u_{!}f)^{\widehat{\phantom{,}}}\xrightarrow{e_{u_{!}f}}\widehat{C'}/u_{!}f\longrightarrow
\widehat{C'}{\quoi},
$$
where the functor $(u/f)_{!}$ is fully faithful, the functor
$e_{u_{!}f}$ is an equivalence, and the functor
$\widehat{C'}/u_{!}f\rightarrow\widehat{C'}$ is the forgetful functor.

\item In particular, the functor $u_{!}$ is faithful, and therefore the
unit morphism $\id\xrightarrow{\Phi}u^\ast u_{!}$ is a monomorphism.
Moreover, for every morphism $\alpha:H\rightarrow K$ of
$\widehat{C}$, the diagram
\[
\xymatrix@C=3.5cm@R=1.5cm@M=.2cm{
H \ar@{^{(}->}[r]^{\Phi(H)}\ar[d]_{\alpha} & u^{*}u_{!}H \ar[d]^{u.u^{\bigdot}(\alpha)}\\
K \ar@{^{(}->}[r]^{\Phi(K)} &  u^{*}u_!K }
\]
is cartesian.
\end{enumerate}
\end{enonce*}

\begin{enonce*}[remark]{Proof}
\begin{enumerate}
\renewcommand{\labelenumi}{\theenumi)}
\item The\pageoriginale factorization comes from diagram \hyperlink{sga4.exp.I.equation.5.13.1}{(5.13.1)}. The
functor $u$ is faithful. Hence $u/f$ is faithful. Let us show that it is
fully faithful. Let $X$ and $Y$ be two objects of $C$, let
$\can_X:uX\rightarrow u_{!}f$ (\resp $\can_Y:uY\rightarrow u_{!}f$)
be the canonical morphisms, and let
\[
\xymatrix@=1.25cm{
uX \ar[rr]^{m} \ar[rd]_{\can_X}&               & uY \ar[ld]^{\can_Y}\\
   & u^{\bigdot}f &
}
\]
be a morphism of $C'/u_{!}f$. We have
$u_{!}f=\varinjlim_{Z\in \ob C}uZ$, and therefore \hyperref[I.prop3.1]{(3.1)}
$$
\Hom_{\widehat{C'}}(uX,u_{!}f)=\varinjlim_{Z\in
  \ob C}\Hom_{C'}(uX, uZ){\quoi}.
$$
By the definition of the inductive limit, to say that
$\can_Ym=\can_X$ is equivalent to saying that there exist
\begin{enumerate}
\renewcommand{\theenumii}{\alph{enumii}}
\renewcommand{\labelenumii}{\theenumii)}
\item a finite sequence of objects $X_i$ of $C$, $i\in [0,n]$,
  with $X_\circ=X$ and $X_n=Y$,

\item for every $i$, a morphism $m_i:uX\rightarrow
  uX_i$ $(m_\circ=\id_{X}, m_n=m)$,

\item for every pair $(i,i+1)$, a morphism $f_i:X_i\rightarrow
  X_{i+1}$ or else $f_i:X_{i+1}\rightarrow X_i$,
\end{enumerate}
such that the diagrams
\begin{displaymath}
\xymatrix@R=1.25cm@C=1.25cm{
   &  uX \ar[rd]^{\;\;m_{i+1} \qquad\quad\text{ or else}} \ar[ld]_{m_i}&\\
  uX_i\ar[rr]^{u(f_i)} &          &  uX_{i+1}
}\qquad \xymatrix@R=1.25cm@C=1.25cm{
      &  uX \ar[rd]^{\;\;m_{i+1}} \ar[ld]_{m_i}&\\
 uX_i &                                        &
   uX_{i+1}\ar[ll]_{u(f_i)}\quad ,
}
\end{displaymath}
are commutative.

It then follows immediately, by induction on $i$ and using property
(PPF), that $m_i$ is of the form $u(p_i)$. In particular $m=u(p)$, and
hence $u/f$ is fully faithful.

\item The\pageoriginale factorization is immediate. The functor
$(u/f)_{!}$ is fully faithful by \hyperref[I.prop5.6]{5.6}. The other
assertions follow from \hyperref[I.prop5.11]{5.11}.

\item The functor $u_{!}$ is the composite of the faithful forgetful
functor and fully faithful functors. It is therefore faithful. It follows
from the general properties of adjoint functors that the unit morphism
$\id\rightarrow u^\ast u_{!}$ is a monomorphism. By 2), the functor
$u_{!}$ appears as the composite of a fully faithful functor
$v:\widehat{C}\to\widehat{C'}/u_{!}f$ and the forgetful functor. The
functor $u^\ast$, right adjoint to $u_{!}$, is therefore the composite of
the ``product with $u_{!}f$'' functor, right adjoint to the forgetful
functor, and a functor $w$ right adjoint to $v$. Moreover, since $v$ is
fully faithful, the unit morphism $\id\rightarrow wv$ is an isomorphism.
The last assertion follows readily.
\end{enumerate}
\end{enonce*}


% END INCLUDED SGA4_English_current_buildable_source_20260906/source/01/drafts/SGA4_I_LINES_1497_1738_English.texfrag


% BEGIN INCLUDED SGA4_English_current_buildable_source_20260906/source/01/drafts/SGA4_I_LINES_1739_1946_English.texfrag
\section{Faithful functors and conservative functors}\etiquette[6]{I.sec6}\pageoriginale

\begin{enonce*}{Definition 6.1}\etiquette[6.1]{I.def6.1}
Let $E$ be a category and let $(\varphi_i:E\rightarrow F_i)_{i\in I}$
be a family of functors
$$
\varphi_i:E\longrightarrow F_i{\quoi}.
$$
The family of functors $(\varphi_i)$ is said to be faithful if, for every
pair of objects $X$, $Y$ of $E$ and every pair of arrows
$u$, $v:X\rightrightarrows Y$, the relation
$\varphi_i(u)=\varphi_i(v)$ for every $i\in I$ implies $u=v$ (in other
words, if the map $\Hom(X,Y)\rightarrow
\Pi_i\Hom(\varphi_i(X),\varphi_i(Y))$ defined by $(\varphi_i)$ is
injective). The family of functors $(\varphi_i)$ is said to be
conservative if every arrow $u$ of $E$ for which $\varphi_i(u)$ is an
isomorphism for every $i\in I$ is an isomorphism. We say that
$(\varphi_i)$ is conservative for monomorphisms (\resp for epimorphisms,
\resp \ldots) if the preceding condition holds whenever $u$ is a
monomorphism (\resp an epimorphism, \resp \ldots).
\end{enonce*}

\setcounter{subsection}{1} \setcounter{subsubsection}{0}
\subsubsection{}\etiquette[6.1.1]{I.sec6.1.1}
If we introduce the unique functor
$$
\varphi:E\longrightarrow F=\prod_{i\in I} F_i
$$
defined by the family of functors $(\varphi_i)$, it is clear that this
family is faithful (\resp conservative, \resp conservative for
monomorphisms, \resp\ldots) if and only if the functor $\varphi$ is
faithful (\resp conservative, \resp\ldots) (by which we mean that the
family consisting only of $\varphi$ is faithful, \resp conservative,
\resp\ldots). Thus, without major inconvenience, we could restrict
ourselves below to a family consisting of a single functor. Nevertheless,
for ease of future reference, we shall state the following results for
families.

The\pageoriginale notions of \hyperref[I.def6.1]{6.1} are especially useful when the
$\varphi_i$ satisfy suitable exactness properties, and in that case they
tend to coincide:

\begin{enonce*}{Proposition 6.2}\etiquette[6.2]{I.prop6.2}
The notation is that of \hyperref[I.prop5.1]{5.1}.
\begin{enumerate}
\renewcommand{\theenumi}{\roman{enumi}}
\renewcommand{\labelenumi}{(\theenumi)}
\item If the kernels of double arrows, or the cokernels of double arrows,
  are representable in $E$, and if the $\varphi_i$ commute with them,
  then one has the implication
$$
(\varphi_i)\mbox{ conservative } \Longrightarrow (\varphi_i)\mbox{
    faithful}{\quoi}.
$$

\item Suppose that fiber products (\resp amalgamated sums) are
  representable in $E$, and that the $\varphi_i$ commute with them.
  Suppose that $(\varphi_i)$ is faithful or conservative; then for every
  arrow $u$ of $E$, $u$ is a monomorphism (\resp an epimorphism) if and
  only if the same holds for $\varphi_i(u)$ for every $i\in I$.

\item Suppose that fiber products and amalgamated sums are representable
  in $E$ and that the $\varphi_i$ commute with them, and that every arrow
  in $E$ which is a bimorphism (\ie both a monomorphism and an
  epimorphism) is an isomorphism. Then one has the implication
$$
(\varphi_i)\mbox{ faithful } \Longrightarrow (\varphi_i)\mbox{
    conservative}{\quoi}.
$$

\item Suppose that fiber products (\resp amalgamated sums) are
  representable in $E$, and that the $\varphi_i$ commute with them. Then,
  if $(\varphi_i)$ is conservative for monomorphisms (\resp for
  epimorphisms), it is in fact conservative.

\item Let $D$ be a type of diagram, let $F:d\mapsto F(d)$ be a diagram of
  type $D$ in $E$, let $X$ be an object of $E$, and let
  $u=(u_d)_{d\in D}$ be a family of arrows
$$
X\longrightarrow F(d)\quad (\mbox{\resp  } F(d)\longrightarrow
X){\quoi}.
$$
Suppose\pageoriginale that $(\varphi_i)$ is conservative, that projective
(\resp inductive) limits of type $D$ are representable in $E$, and that
the $\varphi_i$ commute with them. Then $u$ makes $X$ a projective
(\resp inductive) limit of $F$ in $E$ if and only if, for every $i\in I$,
$\varphi_i(u)$ makes $\varphi_i(X)$ a projective (\resp inductive) limit
of $\varphi_i(D)$ in $E_i$.
\end{enumerate}
\end{enonce*}

\noindent \emph{Proof.}
\begin{enumerate}
\renewcommand{\theenumi}{\roman{enumi}}
\renewcommand{\labelenumi}{(\theenumi)}
\item For the non-\resp statement, given a double arrow
$u$, $v:X\rightrightarrows Y$, it suffices to express the equality $u=v$
by the condition that the inclusion $\Ker(u,v)\rightarrow X$ is an
isomorphism. Here and below, we omit repetition of the dual argument for
the dual statement.

\item If $(\varphi_i)$ is faithful, the condition that
$u:X\rightarrow Y$ be a monomorphism is expressed by the equality
$\pr_1=\pr_2$ for the fiber product $\fprod{X}{X}{Y}$. If $(\varphi_i)$
is conservative, it is expressed by the condition that the diagonal
morphism $\delta:X\rightarrow \fprod{X}{X}{Y}$ be an isomorphism.

\setcounter{enumi}{3}
\item Since in the latter argument the morphism $\delta$ is a
monomorphism, we see that it would in fact have been enough to suppose
that $(\varphi_i)$ is conservative for monomorphisms. But this then
implies that $(\varphi_i)$ is conservative without qualification.
Indeed, if $u\in %%
F\ell\,E$ is such that the $\varphi_i(u)$ are isomorphisms, the preceding
argument shows that these are monomorphisms, and hence they are
isomorphisms by the hypothesis on $(\varphi_i)$.

\setcounter{enumi}{2}
\item This is a trivial consequence of (ii).

\setcounter{enumi}{4}
\item This is a trivial consequence of the definitions.
\end{enumerate}

We record the following consequence of (i), (ii), and (iv):

\begin{enonce*}{Corollary 6.3}\etiquette[6.3]{I.cor6.3}
Suppose\pageoriginale that fiber products and amalgamated sums are
representable in $E$ and that the $\varphi_i$ commute with them, and that
kernels of double arrows or cokernels of double arrows are representable
and that the $\varphi_i$ commute with them. (It suffices, for example,
that finite projective limits and finite inductive limits be representable
in $E$, and that the $\varphi_i$ be exact functors.) Then the following
conditions are equivalent:
\begin{itemize}
\item[a)] $(\varphi_i)$ is faithful.

\item[b)] $(\varphi_i)$ is conservative.

\item[c)] $(\varphi_i)$ is conservative for monomorphisms.

\item[${\rm c')}$] $(\varphi_i)$ is conservative for epimorphisms.
\end{itemize}
\end{enonce*}
We also record, for reference:

\begin{enonce*}{Proposition 6.4}\etiquette[6.4]{I.prop6.4}
Let $\varphi:E\rightarrow F$ be a functor with a right adjoint $\psi$
(thus $\Hom(\varphi(X),Y)\simeq \Hom(X,\psi(Y))$). The functor $\varphi$
(\resp $\psi$) is faithful if and only if, for every element $X$ of $E$
(\resp every element $Y$ of $F$), the unit morphism
$X\rightarrow \psi\varphi(X)$ is a monomorphism. The functor $\varphi$
(\resp $\psi$) is fully faithful if and only if the preceding unit
morphism is an isomorphism.
\end{enonce*}
Indeed, if $X'$, $X$ are two objects of $E$, the map
\begin{equation*}
\Hom(X',X)\longrightarrow
  \Hom(\varphi(X'),\varphi(X))\tag{$\ast$}
\end{equation*}
is identified with the map obtained from
\begin{equation*}
X\longrightarrow \psi
  \varphi (X)\tag{${\ast}{\ast}$}
\end{equation*}
by applying the functor $\Hom(X',-)$. Thus, for $(\ast)$ to be a
monomorphism\pageoriginale (\resp an isomorphism) for every $X'$, with
$X$ fixed, it is necessary and sufficient that the unit morphism
$({\ast}{\ast})$ be a monomorphism (\resp an isomorphism).

\begin{enonce*}{Proposition 6.5}\etiquette[6.5]{I.prop6.5}
Let $p:E'\rightarrow E$ be a functor. The following conditions are %%
equivalent:
\begin{enumerate}
\renewcommand{\theenumi}{\roman{enumi}}
\renewcommand{\labelenumi}{(\theenumi)}
\item $p$ is faithful, conservative, and fibered {\rm(SGA 1 VI 6.1).}

\item $p$ is a fibered functor whose fibers are discrete categories.

\item For every $X'\in \ob E'$, the functor $E'_{/X'}\rightarrow
  E_{/p(X)}$ induced by $p$ is an equivalence of categories that is
  surjective on objects.

\item (When $E$ is a $𝒰$-category.) There exists an element
  $F\in \ob\widehat{C}$ and an equivalence of categories over $E$
  {\rm(SGA 1 VI 4.3)} $E'\xrightarrow{\approx}E_{/F}$ (where $E_{/F}$
  is the full subcategory of $\widehat{E}_{/F}$ formed by the arrows
  $X\rightarrow F$ whose source belongs to $E$).
\end{enumerate}
\end{enonce*}

\setcounter{subsection}{5} \setcounter{subsubsection}{0}

% END INCLUDED SGA4_English_current_buildable_source_20260906/source/01/drafts/SGA4_I_LINES_1739_1946_English.texfrag


% BEGIN INCLUDED SGA4_English_current_buildable_source_20260906/source/01/drafts/SGA4_I_LINES_1947_2034_English.texfrag
\subsubsection{}\etiquette[6.5.1]{I.sec6.5.1}
Recall that a category $C$ is said to be \emph{discrete} if it is a
\emph{groupoid} (\ie every arrow in it is invertible) and if it
is \emph{rigid} (\ie the automorphism group of every object
is reduced to the trivial group); equivalently, the
category is equivalent to the category $C'$ defined by a
set $I$ (with $\text{ob}C'=I$, and with the identity
arrows as its only arrows). When $C$ is already assumed to be a
groupoid, saying that $C$ is discrete amounts to saying that for
two objects $X$, $Y$ of $C$, there exists at most one arrow from $X$
to $Y$, \ie that $C$ is isomorphic to the category defined by a
preordered set.

The equivalence\pageoriginale of conditions (i) and (ii) of
\hyperref[I.prop6.5]{6.5} is an immediate consequence of the preceding observations
and of the following

\begin{enonce*}{Lemma 6.5.2}\etiquette[6.5.2]{I.lem6.5.2}
Let $p:E'\rightarrow E$ be a fibered functor. Then:
\begin{enumerate}
\renewcommand{\theenumi}{\roman{enumi}}
\renewcommand{\labelenumi}{(\theenumi)}
\item For $p$ to be conservative, it is necessary and sufficient that its
  fiber categories be groupoids.

\item For $p$ to be faithful, it is necessary and sufficient that its
  fiber categories be ordered categories.
\end{enumerate}
\end{enonce*}

\noindent \emph{Proof of} \hyperref[I.lem6.5.2]{6.5.2}.
\begin{enumerate}
\renewcommand{\theenumi}{\roman{enumi}}
\renewcommand{\labelenumi}{(\theenumi)}
\item Suppose $p$ is conservative. For every arrow $u$ in a fiber
  $E'_X$, $p(u)=\id_X$ is an isomorphism, so $u$ is an
  isomorphism in $E'$, hence also in $E'_X$ (since an inverse of
  $u$ in $E'$ will clearly be an inverse in $E'_X$). Thus $E'_X$
  is a groupoid. Conversely, suppose that the $E'_X$ are
  groupoids, and let $u'$ be an arrow of $E'$ such that $p(u')$
  is an isomorphism; let us prove that $u$ is an isomorphism. To
  this end, note that, since $p$ is fibered, we can factor
  $u':X'\rightarrow Y'$ as a composite $X'\rightarrow
  u^\ast(Y')\rightarrow Y'$, where the first arrow is an
  $X$-morphism (N.B. $X=p(X')$, $u=p(u')$), and the second is a
  cartesian morphism over $u$. The first arrow is
  an isomorphism since $E'_X$ is a groupoid, and the second
  is as well, since a \emph{cartesian} %%
morphism of a fibered category
  is clearly an isomorphism whenever its projection is an isomorphism.

\item Suppose $p$ is faithful, and let $X'$, $Y'$ be two objects of a
  fiber category $E_X$. Then two arrows from $X'$ to $Y'$ lie
  over the same arrow $\id_X$ of $E$, and hence are identical,
  so $E_X$ is ordered. Conversely,\pageoriginale suppose that the
  fiber categories are ordered, and let us prove that $p$ is
  faithful. Thus let $u'$, $v':X'\rightrightarrows Y'$ be
  arrows of $E'$ lying over the same arrow $u:X\rightarrow Y$
  of $E$. They then factor as $X'\rightrightarrows
  u^\ast(Y)\rightarrow Y'$, where the two arrows
  $X'\rightrightarrows u^\ast(Y')$ are arrows of $E'_X$ with
  the same source and target; they are therefore equal, and hence
  $u'=v'$, \cqfd
\end{enumerate}

Let us return to the proof of \hyperref[I.prop6.5]{6.5}. We have proved (i)
$\Leftrightarrow$ (ii). On the other hand, (ii) %%
$\Leftrightarrow$ (iv) is fairly clear: indeed, on the one hand, the
category $E_{/F}$ is fibered over $E$, with fiber categories the
discrete categories defined by the sets $F(X)$, as
follows immediately from the definitions; on the other hand, if $p$ is as
in (ii), then by the sorites of SGA 1 VI 8 the fibered category
$E'$ over $E$ is $E$-equivalent to the split category over $E$
defined by the functor $E\rightarrow (\Cat)$ determined by the
functor $F:E\rightarrow (\Ens)$, which assigns to every $X\in \ob E$
the set of isomorphism classes of objects of $E'_X$. Now this
split category is $E$-isomorphic to the category $E_{/F}$. Since
(iv) $\Rightarrow$ (iii) is clear, it remains to prove (iii)
$\Rightarrow$ (i). Now it is clear that for $p$ to be faithful
(\resp conservative), it is necessary and sufficient that the induced functors
$E'_{/X'}\rightarrow E_{/p(X)}$ be so; \textit{a fortiori}, it
suffices that these be fully faithful; therefore it remains only to
prove that (iii) implies that $p$ is fibered. But one
again sees that a functor $p$ is fibered if and only if the
induced functors $E'_{/X'}\rightarrow E_{/p(X')}$ are
fibered. This is in particular the case if they are equivalences of categories
that are surjective on objects.


% END INCLUDED SGA4_English_current_buildable_source_20260906/source/01/drafts/SGA4_I_LINES_1947_2034_English.texfrag


% BEGIN INCLUDED SGA4_English_current_buildable_source_20260906/source/01/drafts/SGA4_I_LINES_2035_2353_English.texfrag
\section{Generating and cogenerating
  subcategories}\etiquette[7]{I.sec7}\pageoriginale

\begin{enonce*}{Definition 7.1}\etiquette[7.1]{I.def7.1}
Let $E$ be a category and $C$ a full subcategory of $E$. We
say that $C$ is a subcategory of $E$ \emph{generating by strict
epimorphisms (\resp  by epimorphisms)} if, for every object
$X$ of $E$, the family of arrows of $E$ with target $X$ and source
$X'\in \ob C$ is epimorphic (\resp  strictly epimorphic)
(\hyperref[I.condef1.3]{1.3}). We say that $C$ is a subcategory of $E$
\emph{generating (\resp  generating for monomorphisms, \resp
generating for strict monomorphisms)} if, for every arrow
$u:Y\rightarrow X$ (\resp  every monomorphism of $E$, \resp  every
strict monomorphism of $E$ $\hyperref[I.sec10.5]{(10.5)}$) such that, for every $X'\in
\ob C$, the corresponding map $\Hom(X',Y)\rightarrow
\Hom(X',X)$ is bijective, $u$ is an isomorphism. Finally, we say
that a family $(X_i)$ of objects of $E$ is generating by
strict epimorphisms (\resp \ldots) if the full subcategory $C$
of $E$ generated by this family is generating by strict epimorphisms
(\resp \ldots).
\end{enonce*}

\setcounter{subsection}{1} \setcounter{subsubsection}{0}
\subsubsection{}\etiquette[7.1.1]{I.sec7.1.1}
Observe that, in terms of the family $(h_{X'})_{X'\in \ob C}$ of
functors
$$
h_{X'}:E\longrightarrow (\Ens)\; (X'\in \ob C)
$$
represented by the $X'\in \ob C$, the condition that
$C$ be generating (\resp  generating for
monomorphisms,
\resp  generating for strict monomorphisms) can be expressed by
saying that the family $(h_{X'})$ is \emph{conservative} (\resp  conservative
for monomorphisms, \resp  conservative for strict monomorphisms)
(\hyperref[I.def6.1]{6.1}). It also follows immediately from the
definitions that $C$ is generating by epimorphisms
if\pageoriginale and only if the family $(h_{X'})_{X'\in \ob C}$
is \emph{faithful} (\hyperref[I.def6.1]{6.1}). We shall also give below
(\hyperref[I.prop7.2]{7.2} (i)) an analogous interpretation of the condition
that $C$ be generating by strict epimorphisms.

\subsubsection{}\etiquette[7.1.2]{I.sec7.1.2}
As with the notions introduced in \hyperref[I.def6.1]{6.1}, the notions
of \hyperref[I.def7.1]{7.1} are especially useful when $E$ has suitable
exactness properties, in which case the various notions
introduced have a marked tendency to be all equivalent
(\hyperref[I.cor7.3]{7.3}). Thus the question of which of
these notions in \hyperref[I.def7.1]{7.1} should be regarded as the most
important scarcely arises; in the most important cases, these
notions coincide, and the term ``generating
subcategory''
may therefore be interpreted indifferently as
referring to
any one of the properties considered in \hyperref[I.def7.1]{7.1}
(for example the first, which is the strongest of all, as we
shall see (\hyperref[I.prop7.2]{7.2} (ii))).

\subsubsection{}\etiquette[7.1.3]{I.sec7.1.3}
Suppose that $E$ is a $𝒰$-category, and consider the
canonical functor
\begin{equation*}
\varphi:E\longrightarrow
  \widehat{C}=\Hhom(C, 𝒰\text{-}\Ens)\tag{7.1.3.1}
\end{equation*}\etiquette[7.1.3.1]{I.sec7.1.3.1}
which is the composite of the functors $E\rightarrow \widehat{E}\rightarrow
\widehat{C}$, where the first functor is
the canonical functor (\hyperref[I.sec1.3.3]{1.3.3}), and
the second is
the restriction functor to $C$. Observe that it is plainly
equivalent to say that the preceding functor $\varphi$ is
conservative (\resp  faithful), or to say that the family of functors
$h_{X'}:X\rightarrow \Hom(X',X)=\varphi(X)(X')$, for variable $X'\in\ob C$,
is a conservative (\resp  faithful) family, that is,
also (\hyperref[I.sec7.1.2]{7.1.2}), that $C$ is generating (\resp
generating by
epimorphisms). Similarly, $\varphi$ is\pageoriginale conservative for
monomorphisms (\resp  for strict monomorphisms) if and
only if the family of the $h_{X'}(X'\in \ob C)$ is conservative for
monomorphisms (\resp  for strict monomorphisms), \ie if
and only if the subcategory $C$ of $E$ is generating for
monomorphisms (\resp  for strict monomorphisms).

\begin{enonce*}{Proposition 7.2}\etiquette[7.2]{I.prop7.2}
Let $E$ be a $𝒰$-category and $C$ a full subcategory.
\begin{enumerate}
\renewcommand{\theenumi}{\roman{enumi}}
\renewcommand{\labelenumi}{(\theenumi)}
\item The following conditions are equivalent:
\begin{enumerate}
\renewcommand{\theenumii}{\alph{enumii}}
\renewcommand{\labelenumii}{\theenumii)}
\item $C$ is a subcategory generating by strict
  epimorphisms.

\item For every $X\in \ob E$, let $C_{/X}$ denote the
  full subcategory of $E_{/X}$ consisting of arrows
  $X'\rightarrow X$ with source $X'\in \ob C$. The natural arrow from the
  inclusion functor $j:C_{/X}\rightarrow E$ to the constant functor
  on $C_{/X}$ with value $X$ makes $X$ an inductive limit
  of $j$:
$$
X\xleftarrow{\sim} \lim_{C_{/X}}X'{\quoi}.
$$

\item The canonical functor $\varphi$ of (\hyperref[I.sec7.1.3.1]{7.1.3.1}) is fully
  faithful.
\end{enumerate}

\item Among the notions of \hyperref[I.def7.1]{7.1} there are the following
implications~:
\[
\xymatrix@C=-.7cm@R=.45cm{
  &  \txt{{\rm 1)} $C$ generating by\\ strict epimorphisms\\
  \emph{($\varphi$ fully faithful)}} \ar@{=>}[rd]&\\
 \txt{{\rm 2)} $C$ generating by\\ epimorphisms\\
 \emph{($\varphi$ faithful)}} &    &
 \txt{{\rm 3)} $C$ generating \phantom{aaa}\\
 \emph{($\varphi$ conservative)}} \ar@{=>}[d]\\
    &      &  \txt{{\rm 4)} $C$ generating for\\ monomorphisms\\
  \emph{($\varphi$ conservative for mon.)}}\ar@{=>}[d]\\
    &      &  \txt{{\rm 5)} $C$ generating for\\ strict monomorphisms\\
  \emph{($\varphi$ conservative for strict mon.).}}\\
}
\]

\item There\pageoriginale are the following conditional implications:
\begin{enumerate}
\renewcommand{\theenumii}{\alph{enumii}}
\renewcommand{\labelenumii}{\theenumii)}
\item If in $E$ epimorphic families of arrows are
  strictly epimorphic, then $2)\Rightarrow 1)$. If in $E$
  monomorphisms are strict, then $5)\Rightarrow 4)$.

\item If in $E$ the kernels of pairs of arrows (\resp  fiber
  products) are representable, then $3)\Rightarrow
  2)$ (\resp  $4)\Rightarrow 3)$).

\item If in $E$ every family of morphisms $X_i\rightarrow X$ with
  common target $X$ factors as a strictly epimorphic
  (\resp  epimorphic) family $X_i\rightarrow Y$ followed by a monomorphism
  (\resp  by a strict monomorphism) $Y\rightarrow X$, then
  $4)\Rightarrow 1)$ (\resp  $5)\Rightarrow 2)$).
\end{enumerate}
\end{enumerate}
\end{enonce*}
Let us immediately record the following.

\begin{enonce*}{Corollary 7.3}\etiquette[7.3]{I.cor7.3}
All the notions considered in $\hyperref[I.def6.1]{6.1}$ (and repeated in the
implication diagram in {\rm(ii)} above) are equivalent
in each of the following two cases:
\begin{enumerate}
\renewcommand{\theenumi}{\roman{enumi}}
\renewcommand{\labelenumi}{(\theenumi)}
\item In $E$, the kernels of pairs of arrows and the fiber
  products are representable, the monomorphisms are strict, and
  the epimorphic families of arrows are strictly
  epimorphic.

\item In $E$, every family $(X_i\rightarrow X)_{i\in I}$ of
  arrows with common target $X$ factors as an
  epimorphic family $(X_i\rightarrow Y)$ followed by a monomorphism
  $Y\rightarrow X$, every monomorphism of $E$ is strict, and every
  epimorphic family of arrows of $E$ is strict.
\end{enumerate}
\end{enonce*}

Indeed,\pageoriginale in case (i) we have 4) $\Rightarrow$ 3) and
3) $\Rightarrow$ 2) by b), and 5) $\Rightarrow$ 4) and 2)
$\Rightarrow$ 1) by a). In case (ii), by c), we have the
implications 4) $\Rightarrow$ 1) and 5) $\Rightarrow$ 2). We therefore
conclude by the implication diagram \hyperref[I.prop6.2]{6.2} (ii).

\medskip
\noindent \emph{Proof of} \hyperref[I.prop7.2]{7.2}.
\begin{enumerate}
\renewcommand{\theenumi}{\roman{enumi}}
\renewcommand{\labelenumi}{(\theenumi)}
\item The implication b) $\Rightarrow$ a) follows at once from the
  definitions. Let us prove a) $\Rightarrow$ b). Thus, under
  hypothesis a), we must prove that, for all $X$, $Y\in \ob E$,
  every system of arrows
$$
u_{X'}:X'\longrightarrow Y
$$
indexed by the $X'\in \ob C_{/X}$, such that
$u_{X'}f=u_{X''}$ for every arrow $f:X''\rightarrow X'$ in
$C_{/X}$, factors through an arrow (necessarily unique by
hypothesis a)) $X\rightarrow Y$. By hypothesis a), it suffices
to check that, for every object $Z$ of $E_{/X}$ and every pair of
morphisms $v':Z\rightarrow X'$, $v'':Z\rightarrow X''$ in
$E_{/X}$, with $X'$ and $X''$ in $C_{/X}$, we have
$u_{X'}v'=u_{X''}v''$. But, by hypothesis a), the family of
arrows $w:X'''\rightarrow Z$, with $X'''\in \ob C$, is
epimorphic, and it therefore suffices to check that, for every such
$w$, we have $(u_{X'}v')w=(u_{X''}v'')w$, which can also be written
$u_{X'}(v'w)=u_{X''}(v''w)$ and follows at once from the hypothesis made
on the family of the $u$.

Let us now prove the equivalence of conditions b) and c. For this,
observe that, for every $X\in \ob E$, the object $\varphi(X)$ of
$\widehat{C}$ is the inductive limit in $\widehat{C}$ of the canonical
functor $\widehat{C}_{/\varphi(X)}\rightarrow \widehat{C}$ (\hyperref[I.prop3.4]{3.4});
but $\widehat{C}_{/\varphi(X)}$ is canonically isomorphic to
$C_{/X}$, so that in $\widehat{C}$ we have
$$
\varphi(X)=\varinjlim_{C_{/X}}X'{\quoi},
$$
where\pageoriginale the object $X'$ of $C$ is identified with the functor
$\in \widehat{C}$ that it represents. Consequently, for a second
object $Y$ of $E$, we have a canonical isomorphism
\begin{align*}
\Hom(\varphi(X),\varphi(Y))\simeq\varprojlim_{C_{/X}}
\Hom(X',\varphi(Y))&\simeq
\varprojlim_{C_{/X}}\varphi(Y)(X')\hyperref[I.prop1.4]{(1.4)}\\
&\simeq \varprojlim_{C_{/C}} \Hom(X'%%
,Y){\quoi}.\tag{$\ast$}
\end{align*}
This being so, the map
$$
\Hom(X,Y)\longrightarrow \Hom(\varphi(X),\varphi(Y))
$$
defined by $\varphi$ is none other, via the isomorphism between the
extreme terms of $(\ast)$, than the map
$$
\Hom(X,Y)\longrightarrow \varprojlim_{C_{/X}} \Hom(X',Y)
$$
deduced from the inductive system of arrows $X'\rightarrow X$ indexed by
$C_{/X}$ considered in \hyperref[I.prop7.2]{7.2} (i) b). Thus the first
map is bijective for every $Y$ (with $X$ fixed) if and
only if $X$ is an inductive limit of the inclusion functor
$j:C_{/X}\rightarrow E$, which proves the equivalence of b) and c).

\item The implications 1) $\Rightarrow$ 2) and 3) $\Rightarrow$ 4)
  $\Rightarrow$ 5) are trivial by virtue of the
  definitions. The implication 1) $\Rightarrow$ 3) is obtained by
  interpreting property 1) as the full faithfulness of
  $\varphi$ by (i), and by observing that a fully
  faithful functor is conservative. But we have already observed
  (\hyperref[I.sec7.1.1]{7.1.1}) that 3)
  means that $\varphi$ is conservative.

\item Assertion\pageoriginale a) is a tautology. Assertion
  b) follows from \hyperref[I.prop6.2]{6.2} (i) (\resp  \hyperref[I.prop6.2]{6.2} (iv)) applied to the functor
  $\varphi$, taking into account that the latter is left exact. Finally, let us prove
  c). For $X\in \ob E$, consider the family of all
  the morphisms $X_i\rightarrow X$, with $X_i\in \ob C$; by
  hypothesis on $E$ it factors as a family $X_i\rightarrow
  Y$ that is effective epimorphic (\resp  epimorphic), followed by a
  monomorphism (\resp  by a strict monomorphism) $Y\rightarrow
  X$. Then hypothesis 4) (\resp  5)) implies that $Y\rightarrow
  X$ is an isomorphism, so the family considered is
  strictly epimorphic (\resp  epimorphic), \cqfd
\end{enumerate}

\begin{enonce*}{Proposition 7.4}\etiquette[7.4]{I.prop7.4}
Let $E$ be a category, $C$ a full subcategory of $E$, and $X$
an object of $E$. Suppose that $C$ is generating in $E$ (\resp  $C$
is generating for monomorphisms, and that the fiber product of two
subobjects of $X$ over $X$ is representable in $E$). Then a
strict subobject $\hyperref[I.sec10.11]{(10.11)}$ (\resp  a subobject) $X'$ of $X$ is
determined %%
when, for every $T\in \ob C$, one knows the subset of $\Hom(T,X)$
which is the image of $\Hom(T,X')$. Consequently, the cardinality of the set of
strict subobjects (\resp  of the set of subobjects) of $X$
is bounded above by $\prod_{T\in \ob C} 2^{\card
  \Hom(T,X)}$, and if $X\in \ob C$, it is bounded above by $2^{\card%%
F\ell\,C}$.
\end{enonce*}

Let us first prove the parenthetical assertion. Let $X'$ and $X''$ be two
subobjects of $X$ such that, for every $T\in \ob C$, the images of
$\Hom(T,X')$ and $\Hom(T,X'')$ in $\Hom(T,X)$ are equal. They
are therefore also equal to the image of $\Hom(T,X''')$, where $X'''$ is the
fiber product of $X'$ and $X''$ over $X$. Since $C$ is generating
for monomorphisms, it follows that the
monomorphisms\pageoriginale $X'''\rightarrow X'$ and
$X'''\rightarrow X''$ are isomorphisms, so $X'$ and $X''$
are equal, each being equal to the subobject $X'''$ of $X$.

Let us prove the nonparenthetical assertion. By definition of the notion of
strict subobject, it suffices to check that knowledge of the
subset $\Hom(T,X')$ of $\Hom(T,X)$ for every $T\in\ob C$ implies
knowledge of those pairs of arrows $\xymatrix{
X\ar@<2pt>[r]^{u,v}\ar@<-2pt>[r]& T }$ such that $ui=vi$, where
$i:X'\rightarrow X$ is the canonical injection. But since $C$ is
generating, the relation $ui=vi$ is equivalent to the relation
$(ui)f=(vi)f$ for every $f\in \Hom(T,X')$, \ie to $ug=vg$ for every
$g\in \Hom(T,X)$ coming from $\Hom(T,X')$ (\ie of the form if, with
$f\in \Hom(T,X')$), which proves our assertion.

\begin{enonce*}{Corollary 7.5}\etiquette[7.5]{I.cor7.5}
Let $E$ be a category, $C$ a full generating subcategory,
and
$X$ an object of $C$. Then a strict quotient (\hyperref[I.sec10.8]{10.8}) $X'$ of $X$
is determined when, for every $T\in \ob C$, one knows the subset of
$\Hom(T,X)^2$ consisting of pairs $(u,v)$ such that $qu=qv$, where
$q:X\rightarrow X'$ is the canonical morphism. Thus the cardinality of
the set of strict quotients of $X$ is bounded above by $\prod_{T\in
\ob C} 2^{\card\Hom(T,X)^2}$.
\end{enonce*}

Indeed, by %%
definition, a strict quotient $X'$ of $X$ is determined when one
knows, for every object $Y$ of $E$, the subset of $\Hom(Y,X)\times
\Hom(Y,X)$ consisting of pairs $(u,v)$ such that $qu=qv$, where
$q:X\rightarrow X'$ is the canonical morphism. But the relation
$qu=qv$ is equivalent to the relation $(qu)f=(qv)f$ for every morphism
$f:T\rightarrow Y$ with source $T$ in $C$, since $C$ is
generating. This relation may also be written $q(uf)=q(vf)$, which
proves the first assertion of \hyperref[I.cor7.5]{7.5}. The second follows
at once.

\setcounter{subsection}{5} \setcounter{subsubsection}{0}

% END INCLUDED SGA4_English_current_buildable_source_20260906/source/01/drafts/SGA4_I_LINES_2035_2353_English.texfrag


% BEGIN INCLUDED SGA4_English_current_buildable_source_20260906/source/01/drafts/SGA4_I_LINES_2354_2668_English.texfrag
\subsubsection{}\etiquette[7.5.1]{I.sec7.5.1}
We\pageoriginale can generalize \hyperref[I.cor7.5]{7.5} by
introducing, for a
family of objects $(X_i)_{i\in I}$ of $E$, the notion of a
\emph{strict quotient
  in $E$ of the family}, by which we mean a strictly epimorphic family
$(p_i:X_i\rightarrow X')_{i\in I}$ of
morphisms of $E$, --- with the understanding, as for ordinary
quotients, that we identify two such families
$(p_i:X_i\rightarrow X')$ and $(q_i:X_i\rightarrow X'')$ if one can
find an isomorphism $v:X'\rightarrow X''$ (necessarily
unique) such that $fp_i=q_i$ for every $i\in I$. With this
terminology, the proof of \hyperref[I.cor7.5]{7.5} applies {\it ne
varietur} to give the

\begin{enonce*}[remark]{Variant 7.5.2}\etiquette[7.5.2]{I.var7.5.2}
Let $E$, $C$ be as in \hyperref[I.cor7.5]{7.5}, and let $(X_i)_{i\in I}$ be a family
of objects of $E$. Then a strict quotient $X'$ of $(X_i)_{i\in I}$
in $E$ (\hyperref[I.sec7.5.1]{7.5.1}) is determined by the following data: for every
$T\in \ob C$,
and every pair $(i,j)\in I\times I$, the subset of $\Hom(T,X_i)\times
\Hom(T,X_j)$ consisting of the pairs $(u,v)$ such that $p_iu=p_jv$, where
for every $i\in I$, $p_i:X_i\rightarrow X'$ denotes the canonical
morphism. Consequently, the cardinality of the set of strict
quotients of $(X_i)_{i\in I}$ in $E$ is bounded above by $ \prod_{T\in
\ob C} \prod_{i, j\in I}2^{\card \Hom(T,X_i)\times
  \card\Hom(T,X_j)}{\quoi}.
$
\end{enonce*}

\setcounter{subsubsection}{2}
\subsubsection{}\etiquette[7.5.3]{I.sec7.5.3}
It is immediately clear that the analogous conclusion holds if one
assumes only that $C$ is generating for strict monomorphisms,
provided that the products $X_i\times X_j$
are representable and that one restricts to the \emph{effective} quotients
of the family $(X_{i})_{i\in I}$, \ie to the strict quotients
$X'$ such that the fiber products $\fprod{X_i}{X_{j}}{X'}$
are representable in $E$ (which is no restriction if
$E$ is stable under fiber products).

\begin{enonce*}{Proposition 7.6}\etiquette[7.6]{I.prop7.6}
Let\pageoriginale $E$ be a category, $C$ a full subcategory
of $E$ generating by epimorphisms
(\hyperref[I.def7.1]{7.1}), $\Pi_\circ=$ an
infinite cardinal $\geq \card  F\ell  C$, $\Pi\geq \Pi_\circ$ a
cardinal, and $(u_i:X_i\rightarrow X)_{i\in I}$ a universally epimorphic
family (\hyperref[I.sec10.3]{10.3}) in $E$ consisting of pullbackable morphisms
(\hyperref[I.sec10.7]{10.7}) with sources $X_i\in\ob C$, and such that $\card I\leq \Pi$.
Then $\card F\ell  C_{/X}\leq \Pi^{\Pi_\circ}$.
\end{enonce*}

Let $T\in \ob C$; it will suffice to prove that
\begin{equation*}
\card\Hom(T,X)\leq
  \Pi^{\Pi_\circ}{\quoi},\tag{7.6.1}
\end{equation*}\MakeLinkTarget*{sga4.exp.I.equation.7.6.1}\label{I.eq7.6.1}
for one then obtains successively
$$
\card\ob C_{/X}\leq(\Pi^{\Pi_\circ})\times\card \ob C_{/X}\leq
(\Pi^{\Pi_\circ})\times \Pi_\circ=\Pi^{\Pi_\circ}{\quoi},
$$
and finally $\card F\ell C_{/X}\leq ((\Pi^{\Pi_\circ})^2)\times
\Pi_\circ=\Pi^{\Pi_\circ}$, since for two objects $S$, $T$ of
$C_{/X}$, one has $\card \Hom_{C_{/X}}(S,T)\leq
\car\Hom_C(S,T)\leq\Pi_\circ$. To prove \ref{I.eq7.6.1}, first note the

\begin{enonce*}{Lemma 7.6.2}\etiquette[7.6.2]{I.lem7.6.2}
Let $J$ be a set such that $\card J=\Pi_\circ$. For every object
$T$ of $C$, and every morphism $f:T\rightarrow X$, there exists an
epimorphic family $(v_j:S_j\rightarrow T)_{j\in J}$, with sources
$S_j\in\ob C$, and, for every $j\in J$, an $i(j)\in I$ and a
$g_j:S_j\rightarrow X_{i(j)}$ such that $u_{i(j)}g_j=fv_j$.
\end{enonce*}

Indeed, the family of $\fprod{X_i}{T}{X}\rightarrow$ is
epimorphic by hypothesis; on the other hand, since $C$ is
generating
by strict epimorphisms, for every $i$ the family of arrows
$S\rightarrow \fprod{X_i}{T}{X}$ with source $S\in \ob C$ is
epimorphic, and hence, by transitivity, the family of arrows
$v:S\rightarrow T$ with source in $\ob C$ which factor through one
of the $\fprod{X_i}{T}{X}$ is epimorphic. These are precisely the arrows
$v:S\rightarrow T$ with source in $C$ for which there exist an
$i\in I$ and a $g:S\rightarrow X_i$ such that $u_ig=fv$. Since
the set of\pageoriginale these arrows $v:S\rightarrow T$ is
contained in $F\ell$ $C$, and therefore has cardinality $\leq \Pi_\circ$, it can
be indexed by the index set $J$, whence the lemma.

Now note that a morphism $f:T\rightarrow X$ is determined once
the $fv_j$ are known, and these are determined once the $g_j$ are known.
Thus $\card\Hom(T,X)$ is bounded above by the cardinality of the set of
families $(i(j),S_j,v_j,g_j)_{j\in J}$; since the cardinality of
the set of maps $j\rightarrow I$ is $\leq
\Pi^{\Pi_\circ}$, and since, for a fixed such map $j\mapsto
i(j)$, the cardinality of the set of corresponding families
$S_j$, $f_j$, $v_j$ is bounded above by that of the set of
maps from $J$ to $F\ell  C\times F\ell  C$, which is $\leq
\Pi^{\Pi_\circ}_\circ$ because $\card(F\ell C\times F\ell
C)=\Pi_\circ^2=\Pi_\circ$, we find that the left-hand side of
\hyperref[I.eq7.6.1]{(7.6.1)} is bounded above by $\Pi^{\Pi_\circ}\times
\Pi_\circ^{\Pi_\circ}=\Pi^{\Pi_\circ}$, which completes the
proof of \hyperref[I.prop7.6]{7.6}.

\begin{enonce*}{Proposition 7.7}\etiquette[7.7]{I.prop7.7}
Let $E$ be a category in which fiber products are representable,
$(X_\alpha)_{\alpha\in A}$ a generating family of objects of $E$,
and $(\varphi_i)_{i\in I}$ a family of functors
$\varphi_i:E\rightarrow E_i$ commuting with fiber products. The family
$(\varphi_i)$ is conservative (\hyperref[I.prop5.1]{5.1}) if and only if,
for every $\alpha \in A$ and every subobject $X'$ of
$X_\alpha$ distinct from $X_\alpha$, there exists an $i\in I$ such that
$\varphi_i(X')\rightarrow \varphi_i(X_\alpha)$ is not an
isomorphism. In this case, if $E$ is a $𝒰$-category and
if $A$ is $𝒰$-small, there exists a
$𝒰$-small subset $J$ of $I$ such that $(\varphi_j)_{j\in J}$ is
already a conservative family of functors.
\end{enonce*}

The necessity of the stated
conservativity condition being
obvious, let us prove its sufficiency. By \hyperref[I.def6.1]{6.1} (iv),
it suffices to
prove that $(\varphi_i)$ is conservative for monomorphisms.
Let $u:Y'\rightarrow Y$\pageoriginale be a monomorphism in $E$ which
is not an isomorphism; we must prove that there exists $i\in I$ such
that $\varphi_i(u)$ is not an isomorphism. By the hypothesis on
$(X_\alpha)$, there exist an $\alpha$ and a morphism
$X_\alpha\rightarrow Y$ which does not factor through $Y'$, in
other words, such that the inverse image $X'$ of $Y'$ is a
subobject of $X_\alpha$ distinct from $X_\alpha$. By hypothesis, there
exists an $i\in I$ such that $\varphi_i(X')\rightarrow \varphi(X)$ is not
an isomorphism. Since $\varphi_i(X')\simeq
\fprod{\varphi(X_\alpha)}{\varphi_i(Y')}{\varphi_{i(Y)}}$, it
indeed follows that $\varphi_i(Y')\rightarrow \varphi_i(Y)$ is not
an isomorphism.

The second assertion of \hyperref[I.prop7.7]{7.7} follows immediately from the
first, in view of \hyperref[I.prop7.4]{7.4}.

\begin{enonce*}{Corollary 7.7.1}\etiquette[7.7.1]{I.cor7.7.1}
Let $E$ be a $𝒰$-category in which fiber products are
representable, and which admits a $𝒰$-small generating family
of objects. Then every generating family
$(Y_i)_{i\in I}$ of $E$ contains a $𝒰$-small generating subfamily
$(Y_j)_{j\in J}(\card J\in 𝒰)$.
\end{enonce*}

Indeed, it suffices to apply \hyperref[I.prop7.7]{7.7} to a
$𝒰$-small generating family $(X_\alpha)_{\alpha\in A}$ of
$E$ and to the family of functors $\varphi_i$ $(i\in I)$ represented
by the $Y_i$.

\begin{enonce*}{Proposition 7.8}\etiquette[7.8]{I.prop7.8}
Let $E$ be a category, $C$ a full subcategory generating
by strict epimorphisms, $D$ a category, and $\Hhom'(E,D)$ the
full subcategory of $\Hhom(E,D)$ consisting of the functors which
commute with inductive limits of type $C_{/X}$, where $X$ is an
arbitrary object of $E$. Then the functor $F\rightarrow F| C$
induces a fully faithful functor
$$
\Hhom'(E,D)\longrightarrow \Hhom(C,D){\quoi}.
$$
Consequently,\pageoriginale if $\Hhom(C,D)$ is a
$𝒰$-category $(\hyperref[I.def1.1]{1.1}.)$, for example {\rm(\hyperref[I.sec1.1.1]{1.1.1}. b))} if $C$
is $𝒰$-small and $D$ is a $𝒰$-category, then
$\Hhom'(E,D)$ is also a $𝒰$-category.
\end{enonce*}

Let $F$, $G:E\rightarrow D$ be two functors in the first
member, and let $u:F\rightarrow G$ be a natural transformation. If $X=\varinjlim
X_i$ in $E$ and if $F$ and $G$ commute with the inductive limit
under consideration, then $u(X):F(X)\rightarrow G(X)$ is identified with the inductive limit
of the morphisms $u(X_i):F(X_i)\rightarrow G(X_i)$, and is therefore determined
once the $u(X_i)$ are known. This shows that the functor considered
in \hyperref[I.prop7.8]{7.8} is faithful, in view of the implication a)
$\Rightarrow$ b) in \hyperref[I.prop7.2]{7.2} (i). Conversely, let $v:F|
\sisi{D}{C}\rightarrow G|\sisi{D}{C}$ be a natural transformation; let us prove that it comes from a
natural transformation $u:F\rightarrow G$. For every $X\in \ob
E$, define
$$
u(X):F(X)=\varinjlim_{C_{/X}}F(X_i)\longrightarrow
G(X)=\varinjlim_{C_{/X}}G(X_i)
$$
as the inductive limit of the $v(X_i)$. It is immediate that one
thus obtains maps functorial in $X$, hence a
natural transformation $u:F\rightarrow G$, \sisi{in}{and} finally that the morphism induced by $u$ from
$F|C\sisi{'}{}$ to $G|C\sisi{'}{}$ is $v$, which completes the proof.


\setcounter{subsection}{8}
\subsection{Cogenerating families and subcategories}\etiquette[7.9]{I.sec7.9}

Let $E$ be a category and $C$ a full subcategory of $E$. We
say that $C$ is cogenerating by strict monomorphisms (\resp
cogenerating by monomorphisms, \resp cogenerating, \resp
cogenerating for epimorphisms, \resp cogenerating for
strict epimorphisms) if the full subcategory $C^\circ$ of
$E^\circ$ is generating by strict epimorphisms (\resp\ldotsn).
The terminology for families is analogous. All the results of this
section concerning\pageoriginale the notion of a generating subcategory
and its variants (\hyperref[I.def7.1]{7.1}) therefore yield
corresponding results
for the dual notions, which we leave to the reader
to formulate for their personal satisfaction.

\begin{enonce*}{Proposition 7.10}\etiquette[7.10]{I.prop7.10}
Let $E$ be a category, $C$ a full
generating subcategory
$\hyperref[I.def7.1]{(7.1)}$, and $D$ a full subcategory of $E$. The subcategory $D$ is
cogenerating $\hyperref[I.sec7.9]{(7.9)}$ if and only if, for every pair of parallel
arrows $\xymatrix{ T\ar@<2pt>[r]^{u,v}\ar@<-2pt>[r]& X }$ in $E$
with source $T\in \ob C$ and $u\neq v$, there exists an arrow
$w:X\rightarrow I$ with target $I\in \ob D$ such that $wu\neq wv$.
\end{enonce*}

By definition, to say that $D$ is cogenerating means that for
every pair of parallel arrows $\xymatrix{ Y\ar@<2pt>[r]^{u,v}\ar@<-2pt>[r]& X
}$ in $E$ such that $u\neq v$, there exists an arrow
$w:X\rightarrow I$, with $I\in\ob D$, such that $wu\neq wv$. Thus
\hyperref[I.prop7.10]{7.10} simply means that it suffices to test this
property when $Y\in \ob C$. Now, since $C$ is generating,
the hypothesis $u\neq v$ implies that there exists $f:T\rightarrow Y$ such
that $uf\neq vf$, whence, by hypothesis, there exists a
$w:X\rightarrow I$ with target $I\in\ob D$ such that $w(uf)\neq w(vf)$,
and hence $wu=wv$, \cqfd

\begin{enonce*}{Corollary 7.11}\etiquette[7.11]{I.cor7.11}
With the notation of \hyperref[I.prop7.10]{7.10}, suppose that the
objects $I$ of $D$ are injective objects of $E$, \ie such that
for every monomorphism $X\rightarrow Y$ in $E$, the corresponding map
$\Hom(Y,I)\rightarrow \Hom(X,I)$ is surjective.
Suppose moreover that every pair of parallel arrows $\xymatrix{
T\ar@<2pt>[r]^{u,v}\ar@<-2pt>[r]& X }$ in $E$ factors as an
epimorphic (\resp effective epimorphic) pair of parallel arrows
$\xymatrix{ Y\ar@<2pt>[r]^{u',v'}\ar@<-2pt>[r]& X }$ followed by a
monomorphism $i:X'\rightarrow X$. Then, in the criterion of
\hyperref[I.prop7.10]{7.10} for $D$ to be cogenerating, one may restrict
to the pairs $(u,v)$ which are epimorphic (\resp
effective epimorphic).
\end{enonce*}

We\pageoriginale conclude:

\begin{enonce*}{Corollary 7.12}\etiquette[7.12]{I.cor7.12}
Let $E$ be a $𝒰$-category admitting a small
generating subcategory $C$, and such that every object of $E$ is the
source of a monomorphism into an injective object of $E$. Suppose
moreover that, for every $T\in\ob C$, the coproduct $T\amalg T$ in $E$
is representable, and that every morphism with source $T\amalg T$
factors as an effective epimorphism followed by a monomorphism.
Then $E$ admits a small full cogenerating subcategory $D$.
More precisely, one may take $D$ such that $\card \ob D\leq
\prod_{T,T'\in \ob C}2^{\card(\Hom(T',T\coprod T))^2}$.
\end{enonce*}

Indeed, by \hyperref[I.cor7.11]{7.11}, it suffices, for every $T\in \ob C$ and every
effective quotient $X$ of $T\amalg T$, to choose an embedding
of $X$ into an injective object $I$ of $E$, and to take for $D$ the
full subcategory of $E$ spanned by these $I$. The conclusion
then follows from \hyperref[I.cor7.5]{7.5}.

\begin{enonce*}[remark]{Examples 7.13}\etiquette[7.13]{I.exem7.13}
To construct small cogenerating subcategories in terms
of small generating subcategories, we are thus led to
seek conditions under which a $𝒰$-category $E$
has ``enough injectives'', \ie every object
embeds into an injective object (by a monomorphism). It is well
known [Tohoku] that this condition holds in a
abelian $𝒰$-category with exact (small) filtered inductive
limits (axiom $AB$ 5 of \loccit) which admits a small
generating family. Moreover, the construction of \loccit is not
essentially tied to abelian categories, and also works
in the category of sheaves of $𝒰$-sets
on a topological space $X\in 𝒰$. We shall not state
here the exactness properties which make the construction in
question work, and shall merely point out that, in\pageoriginale the
special case of the category of sheaves of sets on $X$, one
immediately reduces to the case of \loccit as follows. Note
that if $\mathcal{O}_X$ is a ring on $X$, then every injective
object $I$ of the category of $\mathcal{O}_X$-modules is also
injective as an object of the category of sheaves of sets.
Indeed, if $F$ is a %%
sheaf of sets, and $\mathcal{O}_X[F]$ the ``free $\mathcal{O}_X$-module generated by $F$'', there is by
definition a homomorphism of sheaves of sets
\begin{equation*}
F\longrightarrow \mathcal{O}_X[F]\tag{$\ast$}
\end{equation*}
giving rise to an isomorphism, functorial in $F$,
$$
\Hom_{\mathcal{O}_X}(\mathcal{O}_X[F], I)\xrightarrow{\sim} \Hom
(F,I){\quoi}.
$$
Since the functor $F\mapsto \mathcal{O}_X[F]$ manifestly takes
monomorphisms to monomorphisms, it follows immediately
that if $I$ is an injective $\mathcal{O}_X$-module, it is also an
injective sheaf of sets. It follows that if the morphism
$(\ast)$ is a monomorphism (as is the case if one takes for
$\mathcal{O}_X$ a constant ring with value a ring $A\neq 0$),
then an embedding of $\mathcal{O}_X[F^\ast]$ into an
injective $\mathcal{O}_X$-module $I$ gives an embedding of $F$ into
the injective sheaf of sets underlying $I$.
\end{enonce*}

The same argument applies, without change, to the category
of sheaves of $𝒰$-sets on a $𝒰$-site, which
will be introduced in the following exposé.

The chief interest of the existence of a small
cogenerating subcategory
lies in the representability criterion
\hyperref[I.cor8.12.7]{8.12.7}
below.

\setcounter{section}{7}

% END INCLUDED SGA4_English_current_buildable_source_20260906/source/01/drafts/SGA4_I_LINES_2354_2668_English.texfrag


% BEGIN INCLUDED SGA4_English_current_buildable_source_20260906/source/01/drafts/SGA4_I_LINES_2669_2936_English.texfrag
\section{Ind-objects and pro-objects}\etiquette[8]{I.sec8}\pageoriginale

\subsection{Cofinal functors and cofinal subcategories}\etiquette[8.1]{I.sec8.1}

\begin{enonce*}{Definition 8.1.1}\etiquette[8.1.1]{I.def8.1.1}
Let
\hypertarget{sga4.exp.I.equation.8.1.1.1}{}$$
\displaylines{\rlap{$(8.1.1.1)$}\hfill \varphi:I\longrightarrow
I'\hfill}
$$
be a functor. The functor $\varphi$ is said to be cofinal if, for every
category $C$ and every functor $u:I'\rightarrow C$, regarding
$\varinjlim u$ and $\varinjlim u\circ\varphi$ as objects of
$\Hhom(G(𝒱\text{-}\Ens))^\circ$ (\hyperref[I.def2.1]{2.1},
\hyperref[I.sec2.3.1]{2.3.1}) (where $𝒱$ is a universe such that
$I$, $I'\in 𝒱$ and $C$ is a $𝒱$-category), the canonical morphism
\hypertarget{sga4.exp.I.equation.8.1.1.2}{}$$
\displaylines{\rlap{$(8.1.1.2)$}\hfill \varinjlim
  u\circ\varphi\longrightarrow \varinjlim u\hfill}
$$
is an isomorphism. When $\varphi$ is the inclusion functor of a
subcategory of $I'$, the subcategory $I$ is said to be cofinal if
$\varphi$ is cofinal.
\end{enonce*}

\setcounter{subsubsection}{1}
\subsubsection{}\etiquette[8.1.2]{I.sec8.1.2}
Notice that for given $\varphi$ and $u$, the bijectivity of \hyperlink{sga4.exp.I.equation.8.1.1.2}{(8.1.1.2)}
does not depend on the choice of the universe $𝒱$.

Returning to the definition of the terms occurring in \hyperlink{sga4.exp.I.equation.8.1.1.2}{(8.1.1.2)}
(\hyperref[I.def2.1]{2.1}, \hyperref[I.sec2.3.1]{2.3.1}), we see that to
say that $\varphi$ is cofinal also means that, for every universe $𝒱$
such that $I$ and $I'$ are $𝒱$-small, and every functor
$$
v:{I'}^\circ\longrightarrow (𝒱\text{-}\Ens){\quoi},
$$
the canonical homomorphism
\hypertarget{sga4.exp.I.equation.8.1.2.1}{}$$
\displaylines{\rlap{$(8.1.2.1)$}\hfill \varprojlim v\longrightarrow
  \varprojlim v\circ\varphi\hfill}
$$
is an isomorphism, \ie a bijection. To see that this condition is indeed
necessary, observe, on putting $v=v^\circ$, that it also says that the
condition of \hyperref[I.def8.1.1]{Definition 8.1.1} is satisfied when
$C=(𝒱\text{-}\Ens)^\circ$ (which\pageoriginale implies that the
inductive limits considered in \hyperlink{sga4.exp.I.equation.8.1.1.2}{(8.1.1.2)} are representable in $C$).

\setcounter{paragraph}{1}
\paragraph{}\etiquette[8.1.2.2]{I.sec8.1.2.2}
It follows immediately from the definition that \emph{the composite of
two cofinal functors is cofinal.}

\begin{enonce*}{Proposition 8.1.3}\etiquette[8.1.3]{I.prop8.1.3}
Let $\varphi:I\rightarrow I'$ be a functor.
\begin{enumerate}
\renewcommand{\theenumi}{\alph{enumi}}
\renewcommand{\labelenumi}{\theenumi)}
\item For $\varphi$ to be cofinal, it is necessary that $\varphi$
  satisfy the condition:

\begin{itemize}
\item[F 1)] For every object $i'$ of $I'$, there exists an object $i$ of
  $I$ such that $\varphi(i)$ is an upper bound of $i'$ (\ie such that
  $\Hom(i',\varphi(i))\neq \phi$).
\end{itemize}

\item Suppose that $I$ is filtered. For $\varphi$ to be cofinal, it is
  necessary and sufficient that $\varphi$ satisfy condition {\rm (F 1)}
  and the following condition:
\begin{itemize}
\item[F 2)] For every object $i$ of $I$ and every double arrow
$\xymatrix{ i'\ar@<2pt>[r]^{f',g'}\ar@<-2pt>[r]& \varphi(i) }$ in $I'$
with target $\varphi(i)$, there exists an arrow $h:i\rightarrow j$ in
$I$ such that $\varphi(h)f'=\varphi(h)g'$.
\end{itemize}
Moreover, if $\varphi$ is cofinal, $I'$ is filtered.

\item Suppose that $I'$ is filtered and $\varphi$ is fully faithful. For
  $\varphi$ to be cofinal, it is necessary and sufficient that it satisfy
  condition {\rm F 1)} of {\rm a)}. This implies that $I$ is filtered.
\end{enumerate}
\end{enonce*}

\noindent \emph{Proof.} For necessity in a) and b), use Definition
\hyperref[I.def8.1.1]{8.1.1} only in the case where $u$ is the canonical
inclusion functor (\hyperref[I.sec1.3.3]{1.3.3})
$$
u:I'\hookrightarrow \widehat{I}'
$$
(having chosen a universe $𝒰$ such that $I$ and $I'$ are
$𝒰$-small). We may then interpret \hyperlink{sga4.exp.I.equation.8.1.1.2}{(8.1.1.2)} as an arrow of
$\widehat{I}'$ (\hyperref[I.prop3.1]{3.1}) whose target is the
\emph{terminal} functor\pageoriginale on $I'$ (\hyperref[I.prop3.4]{3.4}).
It is immediately clear that condition F 1) says that this arrow is an
epimorphism of $\widehat{I}'$, \ie that it is surjective at each
argument (taking into account that inductive limits in $\widehat{I}'$
are computed ``argument by argument''). This proves a). Now suppose that
$I$ is filtered and $\varphi$ is cofinal. Then $I'$ is filtered: indeed,
the condition that two objects of $I'$ be bounded above by a third follows
at once from the same condition on $I$ and from F 1), and it remains only
to prove condition PS 2) of \hyperref[I.def2.7]{2.7}, \ie that for every
double arrow
$$
f',g':i'\rightrightarrows j'
$$
in $I'$, there exists an arrow $h':j'\rightarrow k'$ of $I'$ such that
$h'f'=h'g'$. By F~1), we may suppose that $k'$ is of the form
$\varphi(i)$. But since $\varphi$ is cofinal, for every object $i'$ of
$I'$ the set $\varinjlim_i\Hom(i',\varphi(i))$ consists of a single
element. The standard description of \emph{filtered} inductive limits in
(Ens) (\hyperref[I.lem2.8.1]{2.8.1}) therefore implies that there exists
an arrow $h:i\rightarrow j$ in $I$ such that
$\varphi(h)f'=\varphi(h)g'$. This completes the proof that $I'$ is
filtered and at the same time proves condition F~2).

To prove that the conditions in b) are sufficient for $\varphi$ to be
cofinal, use the form \hyperref[I.sec8.1.2]{8.1.2} of the definition. One
sees at once that condition F 1) implies that \hyperlink{sga4.exp.I.equation.8.1.2.1}{(8.1.2.1)} is a
monomorphism, \ie is injective at each argument, while condition F 2)
(together with F 1)) ensures that it is bijective (taking into account
the argument-by-argument computation of projective limits in
$\widehat{{I'}^\circ}$). This proves b).

Finally,\pageoriginale if $I'$ is filtered and $\varphi$ is fully
faithful, it is immediate that condition F 1) implies that $I$ is
filtered and implies condition F 2). Thus c) follows from a) and b).

\setcounter{subsubsection}{3}
\subsubsection{}\etiquette[8.1.4]{I.sec8.1.4}
When $I'$ is a filtered category and $I$ is a full subcategory of $I'$,
we see from \hyperref[I.prop8.1.3]{8.1.3} c) that the condition that $I$ be cofinal in $I'$
depends only on the subset $\Ob I$ of the preordered set $OB\,I'$ (for
the preorder relation ``$x\leq y$'' $\Leftrightarrow$
``$\Hom(x,y)\neq ∅$''). If $J'$ is a preordered set and $J$ a subset of
$J'$, we shall sometimes say that $J$ is a \emph{cofinal subset of} $J'$
when every element of $J'$ is bounded above by an element of $J$. When
$J'$ is filtered, this therefore means that the inclusion functor
$𝒥\hookrightarrow 𝒥'$ for the associated categories is cofinal.

\subsubsection{}\etiquette[8.1.5]{I.sec8.1.5}
In what follows, we shall use the notion of cofinal functor only when the
categories $I$ and $I'$ are filtered. Classically, one even restricted
oneself to categories associated with preordered sets (\ie categories in
which there is at most one arrow with given source and target). It turns
out, however, that this restriction is inconvenient in applications,
since the ``natural'' filtered categories arising in many applications
are \emph{not} ordered categories. The following result, %%
due to P. \textsc{Deligne}, shows nevertheless that there is no essential
difference between the two points of view:

\begin{enonce*}{Proposition 8.1.6}\etiquette[8.1.6]{I.prop8.1.6}
Let\pageoriginale $I$ be a small filtered category. Then there exist a
small ordered set $E$ and a cofinal functor $\varphi:ℰ\rightarrow I$,
where $ℰ$ denotes the category associated with $E$.
\end{enonce*}

First suppose that the preordered set $\Ob I$ has no greatest element.
A subdiagram of $I$ is a pair $D=(\mathcal{O},F)$ consisting of a subset
$F$ of $F\ell(I)$ and a subset $\mathcal{O}$ of $\Ob I$, such that for
every $f\in F$, the source and target of $f$ belong to $\mathcal{O}$. An
element $e$ of $\mathcal{O}$ is called a \emph{terminal object} of the
subdiagram $D$ if, for every $x\in D'$, the set
$\Hom(x,e)\cap F$ of arrows of $D$ with source $x$ and target $y$ has
exactly one element $f_x$, if $f_x=f_y\circ g$ for every arrow
$g:x\rightarrow y$ of $D$, and if $f_e=\id_e$. Let $E$ be the set of
\emph{finite} subdiagrams $D$ of $I$ having a \emph{unique} terminal
element $\varphi(D)$, and order $E$ by inclusion. If $D$, $D'$ are two
elements of $E$ and $D\subset D'$, there is a unique arrow
$\varphi(D',D)$ of $D'$ with source $\varphi(D)$ and target
$\varphi(D')$. If $D\subset D'\subset D''$ are inclusions in $E$, one
clearly has $\varphi(D'',D)=\varphi(D'',D')\varphi(D',D)$; one also has
$\varphi(D,D)=\id_{\varphi(D)}$. Thus we obtain a functor
$\varphi:ℰ\rightarrow I$, and it remains to prove that $E$ is filtered
and $\varphi$ is cofinal. In view of \hyperref[I.prop8.1.3]{8.1.3} b),
three verifications are required:
\begin{enumerate}
\renewcommand{\labelenumi}{\theenumi)}
\item Condition F 1) of \hyperref[I.prop8.1.3]{8.1.3}: for every $i\in I$, take for $D$ the
  subdiagram of $I$ consisting only of the object $i$ and its identity
  arrow; then $i\leq \varphi(D)=i$. This condition F 1) also shows that
  $E\neq ∅$.

\item Condition F 2): For every $i\in \Ob I$, every $D\in E$, and every
  double arrow $i\rightrightarrows \varphi(D)$, find a diagram $D'\in E$
  containing $D$ such that
  $\varphi(D',D):\varphi(D)\rightarrow \varphi(D')$ \emph{equalizes}
  the double arrow. Since $I$ is filtered, there\pageoriginale is an
object $j$ of $I$ and an arrow $f:\varphi(D)\rightarrow j$ that equalizes
the double arrow. Since $I$ has no greatest element, we may suppose that
$j$ is a \emph{strict} upper bound of $\varphi(D)$, which implies that it
is distinct from every other object of $D$. Let $D'$ be the subdiagram
of $I$ whose set of objects is $\mathcal{O}\cup\{j\}$ and whose set of
arrows consists of the arrows of $D$, together with the composites
$x\xrightarrow{f_x}\varphi(D)\xrightarrow{f}j$, where $x$ is an object
of $D$ and $f_x$ is the unique arrow of $D$ with source $x$ and target
$\varphi(D)$, and $\id_j$. It is then clear that $D'$ is a finite
subdiagram of $I$, that it has $j$ as its unique terminal object, hence
that $D'\in E$, and that $D'$ satisfies the required condition.

\renewcommand{\labelenumi}{$3^\circ$}
\item Two subdiagrams $D$, $D'\subset E$ are contained in the same
  $D''\in E$. Indeed, one may find a strict upper bound $j$ of
  $\varphi(D)$ and $\varphi(D')$,
\[
\xymatrix@C=1.5cm@R=1cm{
\varphi(D)\ar[rd]^f && \\
           &j &\\
\varphi(D')\ar[ru]_{f'}&\quad , }
\]
and take for $D''$ the subdiagram whose set of objects is the union of
the sets of objects of $D$ and $D'$ with $\{j\}$, and whose set of arrows
is the union of the sets of arrows of $D$ and $D'$, the set of composites
$f\circ f_x$ ($x$ an object of $D$) and $f'\circ f'_{x'}$ ($x'$ an
object of $D'$), and $\{\id_j\}$. This indeed gives a finite subdiagram
of $E$. Let us show that, after replacing $j$, $f$, $f'$ by $j'$, $gf$,
$gf'$ for a suitable $g:j\rightarrow j'$, we may arrange that $j$ is a
terminal object of $D'$. Equivalently, for every $x$ that is an object
of both $D$ and $D'$, one must have $ff_x=f'f'_{x'}$. This amounts to a
\emph{finite} set of double arrows\pageoriginale to be equalized by a
$g:j\rightarrow j'$, which is possible since $I$ is filtered.
\end{enumerate}

This completes the proof in the case considered. The general case is
reduced to it by introducing the filtered category $\NN$ associated with
the ordered set $\NN$ of natural integers, and observing that the category
$\NN\times I$ is filtered, has no greatest element, and that the
projection $\NN\times I\rightarrow I$ is a cofinal functor.

\begin{enonce*}{Corollary 8.1.7}\etiquette[8.1.7]{I.cor8.1.7}
Let $I$ be a $𝒰$-category. There exist a small filtered ordered set $E$
and a cofinal functor $ℰ\rightarrow I$ if and only if $I$ is filtered
and $\Ob I$ has a small cofinal subset $I'$ (\hyperref[I.sec8.1.4]{8.1.4}).
\end{enonce*}

Necessity follows from \hyperref[I.prop8.1.3]{8.1.3} a), and sufficiency follows from \hyperref[I.prop8.1.6]{8.1.6}
applied to the full subcategory of $I$ defined by $I'$.

\begin{enonce*}{Definition 8.1.8}\etiquette[8.1.8]{I.def8.1.8}
Let $I$ be a filtered category. We say that $I$ is essentially small if
$I$ is a $𝒰$-category and satisfies the equivalent conditions of
\hyperref[I.cor8.1.7]{8.1.7}.
\end{enonce*}

\setcounter{subsection}{1}

% END INCLUDED SGA4_English_current_buildable_source_20260906/source/01/drafts/SGA4_I_LINES_2669_2936_English.texfrag


% BEGIN INCLUDED SGA4_English_current_buildable_source_20260906/source/01/drafts/SGA4_I_LINES_2937_3118_English.texfrag
\subsection{Ind-objects and ind-representable functors}\etiquette[8.2]{I.sec8.2}

\subsubsection{}\etiquette[8.2.1]{I.sec8.2.1}
In what follows, we fix a $𝒰$-category $C$, which we always regard as
embedded in the category $\widehat{C}$ of $𝒰$-presheaves by means of
the canonical functor \hyperref[I.sec1.3.1]{(1.3.1)}
\hypertarget{sga4.exp.I.equation.8.2.1.1}{}$$
\displaylines{\rlap{$(8.2.1.1)$}\hfill h:C\hookrightarrow
  \widehat{C}{\quoi}.\hfill}
$$
An \emph{ind-object of} $C$ (or an \emph{inductive system} of $C$) is
any functor
\hypertarget{sga4.exp.I.equation.8.2.1.2}{}$$
\displaylines{\rlap{$(8.2.1.2)$}\hfill \varphi:I\longrightarrow
  C{\quoi},\hfill}
$$
where\pageoriginale $I$ is a \emph{filtered} category
(\hyperref[I.def2.7]{2.7}), called the \emph{index category} of the
ind-object. (Apart from this, $I$ is arbitrary; in particular, we do not
require $I$ to be a $𝒰$-category.) It will often be convenient to denote
an ind-object $\varphi$ by the indexed notation $(X_i)_{i\in \ob I}$, or
even (by another abuse of notation) $(X_i)_{i\in I}$, or $(X_i)_i$, or
$(X_i)$, where $X_i=\varphi(i)$ for $i\in \ob I$. This notation is no
more and no less abusive than the notation classically used for inductive
systems indexed by filtered ordered sets (where the transition morphisms
are likewise not mentioned). The $X_i$ are called the \emph{component
objects} of the ind-object $𝒳$. Note that, in the general case considered
here, the ``\emph{transition morphisms}'' $\varphi(f)$ of the inductive
system are indexed by the set $F\ell I$ of arrows of $I$, which in
general is not identified with a subset of $\Ob I\times\Ob I$; in other
words, for given $i$ and $j$, with $j$ an upper bound of $i$, there may
exist more than one transition morphism from $X_i$ to $X_j$.

\subsubsection{}\etiquette[8.2.2]{I.sec8.2.2}
The most useful ind-objects $𝒳=(X_i)_{i\in I}$ of $C$ are those for
which the index category $I$ is essentially small \hyperref[I.def8.1.8]{(8.1.8)}; such an
ind-object is called \emph{essentially small}. If $(X_i)_{i\in I}$ is
such an object, then, using the fact that small inductive limits are
representable in $\widehat{C}$ \hyperref[I.prop3.1]{(3.1)}, we may consider
\hypertarget{sga4.exp.I.equation.8.2.2.1}{}$$
\displaylines{\rlap{$(8.2.2.1)$}\hfill
L(𝒳)\mathrel{\mathop=^{\rm def}}
  \varinjlim h\cdot 𝒳 = \varinjlim_i X_i \in \ob
  \widehat{C}{\quoi},\hfill}
$$
where the last limit is taken \emph{in} $\widehat{C}$. Thus $L(𝒳)$ is
the presheaf
\hypertarget{sga4.exp.I.equation.8.2.2.2}{}$$
\displaylines{\rlap{$(8.2.2.2)$}\hfill L(𝒳):Y\longmapsto
  \varinjlim_i \Hom(Y,X_i)\hfill}
$$
on\pageoriginale $C$. This presheaf is called the \emph{presheaf
ind-represented by the ind-object $𝒳=(X_i)_{i\in I}$ of $C$.} A presheaf
$F$ on $C$ is called an \emph{ind-representable presheaf} if it is
isomorphic to a presheaf $L(𝒳)$ ind-represented by an essentially small
ind-object. It follows moreover from \hyperref[I.prop8.1.6]{8.1.6} that,
in this definition, we may suppose that $I$ is a small category, and even
that $I$ is associated with an ordered set $\in 𝒰$.

\subsubsection{}\etiquette[8.2.3]{I.sec8.2.3}
Consider an ind-object $𝒳=(X_i)_{i\in I}$ given by a functor
$\varphi:I\rightarrow C$, and let
$$
u:I'\longrightarrow I
$$
be a functor, where $I'$ is a second filtered category. Then the functor
$\varphi\circ u$ is an ind-object $𝒳'$ of $C$ with index category $I'$;
in indexed notation it is $(X_{u(i')})_{i'\in I'}$. It is called the
ind-object of $C$ \emph{deduced from $(X_i)_{i\in I}$ by a change of
index categories} using the functor $u$. If $I$ and $I'$ (\ie $𝒳$ and
$𝒳'$) are essentially small, one obtains a canonical morphism
\hypertarget{sga4.exp.I.equation.8.2.3.1}{}$$
\displaylines{\rlap{$(8.2.3.1)$}\hfill
  L(𝒳')=\varinjlim_{i'}X_{u(i')}\longrightarrow L(𝒳)=\varinjlim_i
  X_i\hfill}
$$
between the presheaves ind-represented by $𝒳'$ and $𝒳$, respectively.
When $u$ is a cofinal functor (\hyperref[I.def8.1.1]{8.1.1}), $𝒳$ is
essentially small if and only if $𝒳'$ is \hyperref[I.cor8.1.7]{(8.1.7)}, and the preceding
homomorphism \hyperlink{sga4.exp.I.equation.8.2.3.1}{(8.2.3.1)} is an isomorphism
\hypertarget{sga4.exp.I.equation.8.2.3.2}{}$$
\displaylines{\rlap{$(8.2.3.2)$}\hfill L(𝒳')\simeq
  L(𝒳){\quoi}.\hfill}
$$

\subsubsection{}\etiquette[8.2.4]{I.sec8.2.4}
Let\pageoriginale
\hypertarget{sga4.exp.I.equation.8.2.4.1}{}$$
\displaylines{\rlap{$(8.2.4.1)$}\hfill 𝒳=(X_i)_{i\in
    I},𝒴=(Y_j)_{j\in I}\hfill}
$$
be two essentially small ind-objects of $C$, indexed by the (filtered and
essentially small) categories $I$ and $J$. A \emph{morphism} from the
ind-object $𝒳$ to the ind-object $𝒴$ is any morphism
\hypertarget{sga4.exp.I.equation.8.2.4.2}{}$$
\displaylines{\rlap{$(8.2.4.2)$}\hfill L(𝒳)\longrightarrow
  L(𝒴)\hfill}
$$
between the presheaves that they ind-represent. We put
\hypertarget{sga4.exp.I.equation.8.2.4.3}{}$$
\displaylines{\rlap{$(8.2.4.3)$}\hfill
\Hom_{\text{ind}-\ob}(𝒳,𝒴)=\Hom_{\widehat{C}}(L(𝒳),L(𝒴)){\quoi}.
   \hfill}
$$
(The subscript ind-ob on $\Hom$ is omitted when there is no risk of
confusion.) Composition of morphisms of ind-objects of $𝒞$ is defined as
composition of the morphisms of the presheaves that they ind-represent.
If $𝒱\supset 𝒰$ is a universe containing $𝒰$, and if we restrict to
the essentially small ind-objects that belong to $𝒱$, these form the set
of objects of a category whose set of arrows consists of triples
$(𝒳,𝒴,u)$, where $𝒳$ and $𝒴$ are essentially small ind-objects
$\in 𝒱$ and $u:𝒳\rightarrow 𝒴$ is a morphism of ind-objects, \ie a
morphism $L(𝒳)\rightarrow L(𝒴)$. The category thus obtained is denoted by
\hypertarget{sga4.exp.I.equation.8.2.4.4}{}$$
\displaylines{\rlap{$(8.2.4.4)$}\hfill
  \Ind_{𝒱}(C,𝒰){\quoi}.\hfill}
$$
When $𝒱=𝒰$, it is denoted simply by
\hypertarget{sga4.exp.I.equation.8.2.4.5}{}$$
\displaylines{\rlap{$(8.2.4.5)$}\hfill \Ind_{𝒰}(C)\text{ or
}
  \Ind(C){\quoi},\hfill}
$$
and is called the \emph{category of ind-objects of} $C$ relative to
$𝒰$, with the mention of $𝒰$ omitted when no confusion is likely.

Clearly,\pageoriginale if $𝒱\subset 𝒱'$ are two universes containing
$𝒰$, then $\Ind_{𝒱}(C,𝒰)$ is a full subcategory of
$\Ind_{𝒱'}(,𝒰)$; it follows from \hyperref[I.sec8.2.3]{8.2.3} that the inclusion functor is an
\emph{equivalence of categories}:
\hypertarget{sga4.exp.I.equation.8.2.4.6}{}$$
\displaylines{\rlap{$(8.2.4.6)$}\hfill
\Ind_{𝒱}(C,𝒰)\xrightarrow{\approx}\Ind_{𝒱'}(C,𝒰){\quoi}.\hfill}
$$
This shows in particular that \emph{every essentially small ind-object of
$C$ defines an element of $\Ind(C)$, up to unique isomorphism.} This
observation justifies the fairly common abuse of language of identifying
every essentially small ind-object of $C$ with an object of $\Ind(C)$.

It follows at once from the definitions that, for every universe $𝒱$,
there is a canonical functor
\hypertarget{sga4.exp.I.equation.8.2.4.7}{}$$
\displaylines{\rlap{$(8.2.4.7)$}\hfill
  L:\Ind_{𝒱}(C,𝒰)\longrightarrow \widehat{C}\hfill}
$$
which is fully faithful. By \hyperlink{sga4.exp.I.equation.8.2.4.6}{(8.2.4.6)}, these functors are clearly
determined up to unique isomorphism once one of them is known, and in
particular once the canonical functor
\hypertarget{sga4.exp.I.equation.8.2.4.8}{}$$
\displaylines{\rlap{$(8.2.4.8)$}\hfill L:\Ind(C)\longrightarrow
  \widehat{C}{\quoi}.\hfill}
$$
is known. Since this functor is fully faithful, it is often used to
identify an object of the first member with the functor that it
ind-represents, or even to identify $\Ind(C)$ with its essential image in
$\widehat{C}$, consisting of the ind-representable presheaves. Note,
however, that this identification has distinctly more disadvantages than
the analogous identification of $C$ with a full subcategory of
$\widehat{C}$ by the functor $h$ \hyperlink{sga4.exp.I.equation.8.2.1.1}{(8.2.1.1)}, since, unlike the latter, the
functor \hyperlink{sga4.exp.I.equation.8.2.4.8}{(8.2.4.8)} is not in general injective on objects.


% END INCLUDED SGA4_English_current_buildable_source_20260906/source/01/drafts/SGA4_I_LINES_2937_3118_English.texfrag


% BEGIN INCLUDED SGA4_English_current_buildable_source_20260906/source/01/drafts/SGA4_I_LINES_3119_3227_English.texfrag
\subsubsection{}\etiquette[8.2.5]{I.sec8.2.5}
Let us express\pageoriginale Definition \hyperlink{sga4.exp.I.equation.8.2.4.3}{(8.2.4.3)} explicitly in terms
of the indexed expressions \hyperlink{sga4.exp.I.equation.8.2.4.1}{(8.2.4.1)} of the ind-objects under
consideration. Taking Definition \hyperlink{sga4.exp.I.equation.8.2.2.1}{(8.2.2.1)} into account, one obtains
$$
\Hom(𝒳,𝒴)\simeq \varprojlim_i
\Hom(X_i,𝒴)\mathrel{\mathop=^{\rm dfn}} \varprojlim_i
\Hom(X_i, L(𝒴))=\varprojlim_i\varinjlim_j \Hom_C(X_i,
Y_j){\quoi},
$$
where the last equality follows from \hyperlink{sga4.exp.I.equation.8.2.2.2}{(8.2.2.2)} applied to $𝒴$.
Thus one obtains a canonical bijection
\hypertarget{sga4.exp.I.equation.8.2.5.1}{}$$
\displaylines{\rlap{$(8.2.5.1)$}\hfill \Hom((X_i)_{i\in
I},(Y_j)_{j\in
    J})\simeq \varprojlim_{i\in I} \varinjlim_{j\in
    J}\Hom_C(X_i,Y_j){\quoi}.\hfill}
$$
We leave to the reader the task of expressing composition of morphisms of
ind-objects explicitly in terms of this formula. Notice that it follows
at once from this formula that the set $\Hom(𝒳,𝒴)$ of homomorphisms of
ind-objects is $𝒰$-\emph{small}, \ie that the categories \hyperlink{sga4.exp.I.equation.8.2.4.4}{(8.2.4.4)} are
$𝒰$-categories. Equivalently (by \hyperlink{sga4.exp.I.equation.8.2.4.6}{(8.2.4.6)}), the \emph{category
$\Ind(C)$ is a $𝒰$-category}, \ie the second member of \hyperlink{sga4.exp.I.equation.8.2.5.1}{(8.2.5.1)} is
small when $I$, $J\in 𝒰$. This follows at once from the fact that $C$
is a $𝒰$-category, so the sets $\Hom_C(X_i,Y_j)$ are small.

\subsubsection{}\etiquette[8.2.6]{I.sec8.2.6}
Let $I$ be an essentially small filtered category. It follows from
\hyperlink{sga4.exp.I.equation.8.2.5.1}{(8.2.5.1)}, or from the functoriality of the inductive limit of a functor
$I\rightarrow \widehat{C}$ with respect to that functor, that there is a
canonical functor
\hypertarget{sga4.exp.I.equation.8.2.6.1}{}$$
\displaylines{\rlap{$(8.2.6.1)$}\hfill \Hhom(I,C)\longrightarrow
  \Ind(C){\quoi}.\hfill}
$$
Note that this functor is in general \emph{not} fully faithful, nor even
faithful.

\begin{enonce*}[remark]{Remark 8.2.7}\etiquette[8.2.7]{I.sec8.2.7}
The\pageoriginale definitions of \hyperref[I.sec8.2.4]{8.2.4} make clear
the notion of \emph{isomorphism} of two (essentially small) ind-objects,
which also means the isomorphism of the presheaves that they
ind-represent. Note that this is a very weak notion of isomorphism when
compared with the notion of isomorphism in categories of the form
$\Hhom(I,C)$. Thus, many fairly natural conditions that may be imposed on
an ind-object (such as having its components in a given strictly full
subcategory, or being strict (\hyperref[I.sec8.12]{8.12} below))
\emph{are not} invariant under isomorphism. Hence, if one wishes to work
with notions invariant under isomorphism of ind-objects, the notions in
question must be ``saturated'' by passing to the strictly full
subcategory of $\Ind(C)$ generated by the subcategory of $\Ind(C)$ formed
by the objects satisfying the condition under consideration. See, for
example, the notion of essentially constant ind-object introduced below
(\hyperref[I.sec8.4]{8.4}).
\end{enonce*}

\begin{enonce*}[remark]{Exercise 8.2.8}\etiquette[8.2.8]{I.exer8.2.8}
Let $𝒰\subset 𝒱$ be two universes and let $C$ be a
$𝒰$-category belonging to $𝒱$. Denote by $\SysInd(C)$ the following
category:
\begin{enumerate}
\renewcommand{\theenumi}{\alph{enumi}}
\renewcommand{\labelenumi}{\theenumi)}
\item The objects of $\SysInd(C)$ are the functors $\varphi:I\rightarrow C$
where $I$ is an essentially $𝒰$-small filtered category belonging to
$𝒱$.

\item Let $(I,\varphi)$ and $(J,\psi)$ be two objects of $\SysInd(C)$.
A morphism in $\SysInd(C)$ from $(I,\varphi)$ to $(J,\psi)$ is a pair
$(m,u)$, where $m:I\rightarrow J$ is a functor and
$u:\varphi\rightarrow \psi\circ m$ is a morphism of functors. Composition
of morphisms in $\SysInd(C)$ is defined in the evident way.
\end{enumerate}
\end{enonce*}

Let\pageoriginale $S$ be the set of morphisms
$(m,u):(I,\varphi)\rightarrow (J,\Psi)$ of $\SysInd(C)$ such that $m$ is
cofinal and $u$ is an isomorphism. Let $(I,\varphi)$ be an object of
$\SysInd(C)$. The category $F\ell(I)$ of arrows of $I$ maps to $I$ by two
natural functors: source and target. These two functors are moreover
linked by the canonical morphism of functors $v:$ source $\rightarrow$
target. Hence there are two morphisms in $\SysInd(C)$ from
$(F\ell(I),\varphi\circ\text{ source})$ to $(I,\varphi)$:
$p_1(I,\varphi)=(\text{source, } \id)$ and
$p_2(I,\varphi)=(\mbox{target, } \varphi^\ast v:\varphi\circ
\text{source }\rightarrow \varphi\circ{\text{ target}})$. Let
$p:\SysInd(C)\rightarrow\Ind(C)$ be the evident functor. Let $B$ be a
category. Show that the functor $F\mapsto F\circ p$
$$
\Hhom(\Ind(C),B)\longrightarrow \Hhom(\SysInd(C),B)
$$
is fully faithful, and that a functor $G:\SysInd(C)\rightarrow B$ belongs
to its essential image if and only if it has the following two
properties:
\begin{enumerate}
\renewcommand{\labelenumi}{\theenumi)}
\item For every $s\in S$, $F(s)$ is an isomorphism of $B$.

\item For every object $(I,\varphi)$ of $\SysInd(C)$,
$F(p_1(I,\varphi))=F(p_2(I,\varphi))$.
\end{enumerate}


% END INCLUDED SGA4_English_current_buildable_source_20260906/source/01/drafts/SGA4_I_LINES_3119_3227_English.texfrag


% BEGIN INCLUDED SGA4_English_current_buildable_source_20260906/source/01/drafts/SGA4_I_LINES_3228_3408_English.texfrag
\subsection{Characterization of ind-representable functors}\etiquette[8.3]{I.sec8.3}

\begin{enonce*}{Proposition 8.3.1}\etiquette[8.3.1]{I.prop8.3.1}
Let $F$ be an ind-representable presheaf on $C$. Then $F$ is
left exact, \ie (\hyperref[I.cor2.3.2]{2.3.2}) for every finite category
$J$ and every
functor $\varphi:J\rightarrow C$ such that $\varinjlim \varphi$ is
representable (\ie $\varprojlim \varphi^\circ$ is representable),
the canonical map
$$
F(\varinjlim \varphi)\longrightarrow \varprojlim F\circ\varphi
$$
is bijective.
\end{enonce*}

Indeed,\pageoriginale since representable presheaves are
clearly left exact, it follows at once
from \hyperref[I.prop2.8]{2.8} that the same is true
of every filtered inductive limit of such functors, \cqfd

\setcounter{subsubsection}{1}
\subsubsection{}\etiquette[8.3.2]{I.sec8.3.2}
Given a presheaf $F$ on $C$, we shall often have occasion to work
with the category
\hypertarget{sga4.exp.I.equation.8.3.2.1}{}$$
\displaylines{\rlap{$(8.3.2.1)$}\hfill C_{/F}\hookrightarrow
\widehat{C}_{/F}{\quoi},\hfill}
$$
which denotes the full subcategory of the latter category formed by
arrows $X\rightarrow F$ \emph{whose source lies in} $C$. It follows from
\hyperref[I.prop1.4]{1.4} that the objects of this category may be identified with pairs
$(X,u)$, where $X$ is an object of $C$ and $u\in F(X)$. A morphism from
$(X,u)$ to $(X',u')$ is then interpreted as an arrow
$f:X\rightarrow X'$ in $C$ such that $F(f)(u')=u$. The
functor from $\widehat{C}_{/F}$ to $\widehat{C}$ that ``forgets the
target'' induces a functor (called \emph{canonical}).
\hypertarget{sga4.exp.I.equation.8.3.2.2}{}$$
\displaylines{\rlap{$(8.3.2.2)$}\hfill C_{/F}\longrightarrow
C{\quoi},\hfill}
$$
already considered in \hyperref[I.prop3.4]{3.4}, where it is proved that
the inductive limit of this
functor exists in $\widehat{C}$ (without any smallness condition on
$C_{/F}$!) and is canonically isomorphic to $F$.


\setcounter{paragraph}{2}
\paragraph{}\etiquette[8.3.2.3]{I.sec8.3.2.3}
Note that if finite inductive limits are
representable in $C$, and if $F$ is left exact (\ie takes them to
finite projective limits in $(\Ens)$), then the same is true
in $C_{/F}$, and a fortiori (\hyperref[I.exem2.7.1]{2.7.1}) $C_{/F}$ is filtered. If,
more generally, sums of two objects and
cokernels of pairs of parallel arrows are representable in $C$, and if $F$
takes them to products \resp kernels, then\pageoriginale $C_{/F}$
is closed under the same types of finite inductive limits, hence
it is filtered if and only if it is nonempty, \ie if and
only if the functor $F$ is not the constant functor with
value $∅$.

\begin{enonce*}{Theorem 8.3.3}\etiquette[8.3.3]{I.them8.3.3}
Let $C$ be a $𝒰$-category, and let $F\in \ob \widehat{C}$,
$F:C^\circ\rightarrow (𝒰\text{-}\Ens)$ be a
$𝒰$-presheaf on $C$. The following conditions are
equivalent:
\begin{itemize}
\item[(i)] $F$ is ind-representable (\hyperref[I.sec8.2.2]{8.2.2}).

\item[(ii)] The category $C_{/F}$ (\hyperref[I.sec8.3.2]{8.3.2}) is filtered and
essentially small.

\item[(iii)] (if finite inductive limits are
representable in $C$.) The functor $F$ is left exact, and $\ob C_{/F}$
has a small cofinal subset (\hyperref[I.sec8.1.4]{8.1.4}).

\item[(iii bis)] (If sums of two objects and
cokernels of pairs of parallel arrows are representable in $C$.) The functor $F$
takes sums of two objects of $C$ to products, and cokernels of
pairs of parallel arrows of $C$ to kernels; the category $C_{/F}$ is nonempty,
\ie the functor $F$ is not the constant functor with value
$∅$; finally, there exists a small family of objects of
$C_{/F}$ such that every object $X$ of $C_{/F}$ is dominated by an
object $X_i$, \ie admits an $F$-morphism $X_i\rightarrow X$.

\item[(iv)] (if the category $C$ is equivalent to a small
category.) The category $C_{/F}$ is filtered.

\item[(v)] (If the category $C$ is equivalent to a small
category, and if filtered inductive limits are representable in
it.) The functor $F$ is left exact.
\end{itemize}
\end{enonce*}

\noindent \emph{Proof.} (i)\pageoriginale $\Rightarrow$
(ii). Suppose that $F$ is ind-represented by $(X_i)_{i\in I}$, with $I$
small. Let us prove that $C_{/F}$ is filtered. Let $X$, $X'$ be two
objects of $C_{/F}$, \ie objects of $C$ equipped with morphisms
$X\rightarrow F$, $X'\rightarrow F$; let us prove that they have
a common upper bound in $C_{/F}$. The morphisms $X\rightarrow
F$, $X'\rightarrow F$ come from morphisms $X\rightarrow X_i$,
$X'\rightarrow X'_i$, and, after replacing $i$, $i'$ by a
common upper bound in $\Ob I$, we may assume that $i=i'$, and take
as a common upper bound of $X$, $X'$ the object $X_i$ equipped with the
canonical morphism $X_i\rightarrow F$. Now let $\xymatrix{
X\ar@<2pt>[r]^{f,g}\ar@<-2pt>[r]& X' \ar@<2pt>[r]^{h} & F }$
 be a pair of parallel arrows in $C_{/F}$; let us prove that it is
equalized by an arrow $X'\rightarrow X''$ in $C_{/F}$. Now
$h:X'\rightarrow F$ is given by a morphism $h_i:X'\rightarrow
X_i$, and the condition $hf=hg$ means that there exists
$\alpha:i\rightarrow j$ in $I$ such that
$\varphi(\alpha)(h_if)=\varphi(\alpha)(h_ig)$, \ie after
replacing $h_i$ by $\varphi(\alpha)h_i$, this morphism equalizes $f$ and $g$.
This proves that $C_{/F}$ is filtered. It is then immediate
that it is essentially small, since the objects
$(X_i\rightarrow F)_{\in \Ob I}$ form a small family in
$\Ob C_{/F}$ that is cofinal.

(ii) $\Rightarrow$ (i). Put $I=C_{/F}$, and consider the
canonical functor \hyperlink{sga4.exp.I.equation.8.3.2.1}{(8.3.2.1)}
$$
\varphi:I=C_{/F}\longrightarrow F{\quoi}.
$$
We know that the presheaf ind-represented by this ind-object of $C$ is
$F$ (\hyperref[I.prop3.4]{3.4}), so $F$ is ind-representable by
definition (\hyperref[I.sec8.2.2]{8.2.2}).

(ii) $\Leftrightarrow$ (iii). Indeed, (ii) $\Rightarrow$ (iii) because
an ind-representable functor is left exact \hyperref[I.prop8.3.1]{(8.3.1)}, and (iii)
$\Rightarrow$ (ii) because we noted (\hyperref[I.prop8.3.1]{8.3.1}) that if $F$
is left exact, then $C_{/F}$ is filtered, and\pageoriginale
we apply Definition \hyperref[I.def8.1.8]{8.1.8} to conclude that this
category is essentially small.

(ii) $\Leftrightarrow$ (iii bis) is proved in the same way as (ii)
$\Leftrightarrow$ (iii).

The equivalences (ii) $\Leftrightarrow$ (iv) and (iii)
$\Leftrightarrow$ (v) are immediate, since if $C$ is equivalent to
a small category, then $C_{/F}$ is clearly likewise equivalent
to a small category and \textit{a fortiori} it is
automatically essentially small as soon as it is filtered.
This completes the proof of \hyperref[I.them8.3.3]{8.3.3}.

\begin{enonce*}[remark]{Remark 8.3.4}\etiquette[8.3.4]{I.rem8.3.4}
Let $F:C'\rightarrow C''$ be a functor between $𝒰$-categories.
It appears that the notion of left exactness
(\hyperref[I.cor2.3.2]{2.3.2}) is of
little practical interest unless finite projective
limits are representable in $C'$. In the case where
$C''=(𝒰\text{-}\Ens)$, and where $C'$ is $𝒰$-small,
writing $C'$ in the form $C^\circ$ (thus $C={C'}^\circ)$, one
should regard the ``right notion'' that replaces,
in the general case, that of left exactness as the
ind-representability of $F$ (regarded as a
presheaf on
$C^\circ$; \ie the pro-representability of $F$,
regarded as a
functor on $C'$, \cf \hyperref[I.sec8.10]{8.10}). This notion does indeed coincide
with that
of left exactness when finite projective
limits are representable in $C'$ %%
 \textit{i.e.} finite inductive limits in $C$ are
                   %%
representable (\hyperref[I.them8.3.3]{8.3.3}). When $C'$ is no longer assumed
$𝒰$-small, or at least equivalent to a
$𝒰$-small category, the two notions for a functor $F$,
left exactness and pro-representability, no longer necessarily coincide,
even if finite projective limits in $C'$ are
representable (\cf \hyperref[I.exer8.12.9]{8.12.9}). In fact, it seems that
in\pageoriginale this case, left exact functors that are not
ind-representable should be regarded as pathological, while
the ``good'' objects remain the ind-representable functors;
compare \hyperref[I.cor8.13.3]{8.13.3}. For the case of functors
$F:C'\rightarrow C''$, with $C'$ and $C''$ once again arbitrary
$𝒰$-categories, we shall develop below
\hyperref[I.sec8.11.5]{(8.11.5)}, more generally, a notion that ``improves'' upon
left exactness (namely, that of a functor admitting a
pro-adjoint).
\end{enonce*}

% END INCLUDED SGA4_English_current_buildable_source_20260906/source/01/drafts/SGA4_I_LINES_3228_3408_English.texfrag


% BEGIN INCLUDED SGA4_English_current_buildable_source_20260906/source/01/drafts/SGA4_I_LINES_3409_3598_English.texfrag
\subsection{Constant and essentially constant ind-objects}\etiquette[8.4]{I.sec8.4}

Choose a terminal category $\sheaf{e}$, \ie one such that $\Ob
% 2013-5-24 compilation problem with \sheaf \XXX
\sheaf{e}$ and $F\ell$ $\sheaf{e}$ each consist of a single element (for
example, one may take for $\sheaf{e}$ the full subcategory
of $(\Ens)$ on the empty set, if a canonical choice is
desired). This is obviously a small filtered category. If
$X$ is an object of $C$, we associate with it the constant ind-object indexed
by $\sheaf{e}$, with value $X$, denoted by $c(X)$. It is clear that, as
$X$ varies, this gives a fully faithful functor
\hypertarget{sga4.exp.I.equation.8.4.1}{}$$
\displaylines{\rlap{$(8.4.1)$}\hfill c:C\longrightarrow
\Ind(C){\quoi},\hfill}
$$
which is moreover injective on objects, and by which we shall identify
$C$ with a full subcategory of $\Ind(C)$. Note that the
composite functor
\hypertarget{sga4.exp.I.equation.8.4.2}{}$$
\displaylines{\rlap{$(8.4.2)$}\hfill
C\xrightarrow{c}\Ind(C)\xrightarrow{L}\widehat{C}\hfill}
$$
is the canonical functor $h$ \hyperlink{sga4.exp.I.equation.8.2.1.1}{(8.2.1.1)}.

\setcounter{subsubsection}{2}
\subsubsection{}\etiquette[8.4.3]{I.sec8.4.3}
More\pageoriginale generally, let $I$ be a filtered category.
The constant functor on $I$ whose value is an object $X$ of $C$ is
also called \emph{the constant ind-object with value $X$ indexed by
$I$.} It is clear that, if $I$ is essentially small, this
ind-object ind-represents the functor $h(X)$ represented by $X$, hence
this ind-object is isomorphic to $c(X)$.

\subsubsection{}\etiquette[8.4.4]{I.sec8.4.4}
An ind-object $𝒳$ of $C$ is said to be \emph{essentially
constant} if it is isomorphic (in a category
$\Ind_{𝒱}(C,𝒱')$, for suitable $𝒱$, $𝒱'$
) to an ind-object of the form $c(X)$, $X\in \ob C$.
The object $X$, which is then determined up to canonical isomorphism,
is called the \emph{value} of the essentially constant ind-object
under consideration. Thus, by definition, the functor $c$ of \hyperlink{sga4.exp.I.equation.8.4.1}{(8.4.1)} induces an
equivalence of $C$ with the full subcategory of $\Ind(C)$ consisting
of the ind-objects that are essentially constant, this
subcategory being the essential image of the functor $c$.

Clearly a constant ind-object is essentially constant,
whereas the converse holds only in the trivial case in which $C$ is
empty or a one-point category.

\subsection{Filtered inductive limits in $\Ind(C)$}\etiquette[8.5]{I.sec8.5}

\begin{enonce*}{Proposition 8.5.1}\etiquette[8.5.1]{I.prop8.5.1}
Let $C$ be a $𝒰$-category. In $\Ind(C)$, small
filtered inductive limits are representable, and the
canonical functor $\hyperlink{sga4.exp.I.equation.8.2.4.8}{(8.2.4.8)}$
$$
L:\Ind(C)\longrightarrow \widehat{C}
$$
commutes with them.
\end{enonce*}

Since\pageoriginale $L$ is fully faithful, the
assertions of the proposition are equivalent to the following:

\begin{enonce*}{Corollary 8.5.2}\etiquette[8.5.2]{I.cor8.5.2}
Every small filtered inductive limit (in $\widehat{C}$) of
ind-representable presheaves is ind-representable.
\end{enonce*}

This follows easily from criterion \hyperref[I.them8.3.3]{8.3.3} (ii). The
details of the verification are left to the reader.

\setcounter{subsubsection}{2}
\subsubsection{}\etiquette[8.5.3]{I.sec8.5.3}
The ``tautological'' case of filtered inductive limits in
$\Ind(C)$ is that in which one starts with an essentially small
ind-object $𝒳=(X_i)_{i\in I}$ of $C$. Then, in $\widehat{C}$
and hence also in $\Ind(C)$ (or in $\widehat{C}$), one has:
\hypertarget{sga4.exp.I.equation.8.5.3.1}{}$$
\displaylines{\rlap{$(8.5.3.1)$}\hfill 𝒳=\varinjlim_I
X_i{\quoi}.\hfill}
$$
When writing this formula, care must be taken that it is not
an inductive limit in $C$, and that even when the latter
exists, it is not isomorphic in $\Ind(C)$ to $𝒳$: indeed,
the canonical functor \hyperlink{sga4.exp.I.equation.8.4.1}{(8.4.1)} $c:C\rightarrow \Ind(C)$
\emph{does not in general commute with filtered inductive limits.}
(\Cf \hyperref[I.exer8.5.5]{8.5.5} below.)

To avoid this possible confusion in writing
\hyperlink{sga4.exp.I.equation.8.5.3.1}{(8.5.3.1)}, ``some authors'' (= {P. Deligne}) prefer
to write it as
\hypertarget{sga4.exp.I.equation.8.5.3.2}{}$$
\displaylines{\rlap{$(8.5.3.2)$}\hfill
%𝒳=\lq\lq{\varinjlim_I}\textrm{ »\,}
𝒳=\textrm{« }{\varinjlim_I}\textrm{ »\,}X_i{\quoi},\hfill}
$$
the quotation marks indicating that the inductive limit is
taken in a category of ind-objects $\Ind(C)$. By extension, one
should then denote by $\textrm{« }{\lim}\textrm{ »}$ every
inductive-limit operation in $\Ind(C)$ (without the
components of the inductive system under consideration in $\Ind(C)$
necessarily lying in the image, or the essential image, of
$c:C\rightarrow \Ind(C)$).

\subsubsection{}\etiquette[8.5.4]{I.sec8.5.4}
To\pageoriginale compute the inductive or projective limit of a
functor
\hypertarget{sga4.exp.I.equation.8.5.4.1}{}$$
\displaylines{\rlap{$(8.5.4.1)$}\hfill \varphi:J\longrightarrow
\Ind(C){\quoi},\hfill}
$$
where $J$ is now an arbitrary category, it is often
convenient to make $\varphi$ explicit by means of a functor
\hypertarget{sga4.exp.I.equation.8.5.4.2}{}$$
\displaylines{\rlap{$(8.5.4.2)$}\hfill \Phi:J\times I\longrightarrow
C{\quoi},\hfill}
$$
where $I$ is an essentially small filtered category, or
preferably a small one. Such a $\Phi$ indeed defines a functor
$$
J\longmapsto (i\longmapsto \Phi (j,i)):J\longmapsto
\Hhom(I,C){\quoi},
$$
whence, by composing with the canonical arrow \hyperlink{sga4.exp.I.equation.8.2.6.1}{(8.2.6.1)}
$\Hhom(I,C)\rightarrow \Ind(C)$, a functor
\hypertarget{sga4.exp.I.equation.8.5.4.3}{}$$
\displaylines{\rlap{$(8.5.4.3)$}\hfill J\longmapsto (i\longmapsto
\Phi(j,i)):J\longrightarrow \Ind(C){\quoi}.\hfill}
$$
We shall say that $\Phi$ is an \emph{indexed expression of
$\varphi$, indexed by the (filtered) category $I$}, if the corresponding
functor \hyperlink{sga4.exp.I.equation.8.5.4.3}{(8.5.4.3)} is isomorphic to $\varphi$. We shall study
below (\hyperref[I.sec8.8]{8.8}) general conditions for the existence
of an indexed expression
for a given functor $\varphi$, and here we shall restrict ourselves to
starting from a functor $\varphi$ given in indexed form
\hyperlink{sga4.exp.I.equation.8.5.4.3}{(8.5.4.3)}, in the case where $J$ is a small
filtered category,
in order to indicate in this case how to compute
$\textrm{«}{\varinjlim}\textrm{»\,}\varphi=\textrm{«}{\varinjlim_j}\textrm{»\,}\varphi(j)$.
Formula \hyperlink{sga4.exp.I.equation.8.5.3.2}{(8.5.3.2)}, together with the associativity formula for
inductive limits (\hyperref[I.sec2.5.0]{2.5.0}), then immediately gives the ``\emph{tautological
computation}'' of
$\textrm{«}{\varinjlim}\textrm{»\,}\varphi$:
\hypertarget{sga4.exp.I.equation.8.5.4.4}{}$$
\displaylines{\rlap{$(8.5.4.4)$}\hfill
\textrm{«}{\varinjlim_J}\textrm{»\,}\varphi=\textrm{«}{\varinjlim_{J\times
I}}\textrm{»\,}\Phi{\quoi},\hfill}
$$
\ie\pageoriginale
\hypertarget{sga4.exp.I.equation.8.5.4.5}{}$$
\displaylines{\rlap{$(8.5.4.5)$}\hfill
\textrm{«}{\varinjlim_j}\textrm{»\,}(i\longmapsto \Phi
(j,i))=\textrm{«}{\varinjlim_{j,i}}\textrm{»\,}\Phi(j,i){\quoi},\hfill}
$$
Thus, the sought-for inductive system is none other than $\Phi$ itself
(the indexing category being $J\times I$, which is indeed filtered
since $J$ and $I$ are).

\begin{enonce*}[remark]{Exercise 8.5.5}\etiquette[8.5.5]{I.exer8.5.5}
Let $C$ be a $𝒰$-category
\begin{enumerate}
\renewcommand{\theenumi}{\alph{enumi}}
\renewcommand{\labelenumi}{\theenumi)}
\item Prove that the following conditions are equivalent:

\begin{itemize}
\item[(i)] In $C$, small inductive limits are representable,
and the functor $c:C\rightarrow \Ind(C)$ commutes with them.

\item[(iii)] The functor $c$ is an equivalence of categories.

\item[(ii)] In $C$, small inductive limits are
representable, and for every $X\in \Ob C$, the functor $Y\mapsto
\Hom(X,Y)$ commutes with these limits.

\item[(iv)] Small filtered inductive limits are
representable in $C$, and the functor $\lim c:\Ind C\rightarrow C$
(\cf \hyperlink{sga4.exp.I.equation.8.7.1.5}{8.7.1.5}) is fully faithful (or equivalently, an
equivalence).
\end{itemize}

\item If $C$ is a finite category, prove that the preceding
conditions are equivalent to the following: for every
projector in $C$
(\hyperref[I.sec10.6]{10.6}), its image is representable in $C$.
\end{enumerate}
\end{enonce*}

% END INCLUDED SGA4_English_current_buildable_source_20260906/source/01/drafts/SGA4_I_LINES_3409_3598_English.texfrag


% BEGIN INCLUDED SGA4_English_current_buildable_source_20260906/source/01/drafts/SGA4_I_LINES_3599_3796_English.texfrag
\subsection{Extension of a functor to ind-objects}\etiquette[8.6]{I.sec8.6}

\subsubsection{}\etiquette[8.6.1]{I.sec8.6.1}
Let
\hypertarget{sga4.exp.I.equation.8.6.1.1}{}$$
\displaylines{\rlap{$(8.6.1.1)$}\hfill f:C\longrightarrow C'\hfill}
$$
be a functor between $𝒰$-categories. For every ind-object
$$
\varphi:I\longrightarrow C
$$
of\pageoriginale $C$ ($I$ a filtered category), the composite functor
$f\circ\varphi$ is an ind-object of $C'$, also denoted by
$\Ind(f)(\varphi)$. In indexed notation, if $\varphi$ is written as
$𝒳=(X_i)_{i\in I}$, one obtains
\hypertarget{sga4.exp.I.equation.8.6.1.2}{}$$
\displaylines{\rlap{$(8.6.1.2)$}\hfill \Ind(f)(X_i)_{i\in
I}=(f(X_i))_{i\in I}{\quoi}.\hfill}
$$
Let $𝒴=(Y_j)_{j\in J}$ be a second inductive system of $C$. One then
obtains an evident map, also denoted by $u\rightarrow\Ind(f)(u)$:
\hypertarget{sga4.exp.I.equation.8.6.1.3}{}$$
\displaylines{\rlap{$(8.6.1.3)$}\hfill
\Hom(𝒳,𝒴)=\varprojlim_i \varinjlim_j \Hom
(X_i,Y_j)\rightarrow\hfill\cr \hfill
\Hom(\Ind(f)(𝒳),\Ind(f)(𝒴))\simeq \varprojlim_i
\varinjlim_j \Hom(f(X_i),f(Y_j)){\quoi}.\hfill}
$$
Referring to the unwritten formula for composition of morphisms of
ind-objects (\hyperref[I.sec8.2.5]{8.2.5}), one immediately sees that
these maps are compatible with composition of morphisms of ind-objects.
Thus, for every universe $𝒱$, we have obtained a functor
\hypertarget{sga4.exp.I.equation.8.6.1.4}{}$$
\displaylines{\rlap{$(8.6.1.4)$}\hfill
\Ind(f):\Ind_{𝒱}(C,𝒰)\longrightarrow
\Ind_{𝒱}(C',𝒰)\hfill}
$$
(notation of \hyperlink{sga4.exp.I.equation.8.2.4.5}{(8.2.4.5)}), and in particular a functor
\hypertarget{sga4.exp.I.equation.8.6.1.5}{}$$
\displaylines{\rlap{$(8.6.1.5)$}\hfill
\Ind(f):\Ind(C)\longrightarrow \Ind(C'){\quoi}.\hfill}
$$
If there is a second functor
$$
g:C'\longrightarrow C''{\quoi},
$$
then one clearly has an identity of functors between $\Ind$-categories
(or $\Ind_{𝒱}$-categories, as desired):
\hypertarget{sga4.exp.I.equation.8.6.1.6}{}$$
\displaylines{\rlap{$(8.6.1.6)$}\hfill \Ind(gf)=\Ind(g)\circ
\Ind(f){\quoi},\hfill}
$$
and,\pageoriginale when $C'=C$, the perennial formula
\hypertarget{sga4.exp.I.equation.8.6.1.7}{}$$
\displaylines{\rlap{$(8.6.1.7)$}\hfill
\Ind(\id_C)=\id_{\Ind(C)}{\quoi}.\hfill}
$$
One may therefore say, heuristically, that the category $\Ind(C)$
\emph{depends functorially on $C$} (covariantly)
\footnote[1]{For a contravariant dependence, under certain conditions,
\cf \hyperref[I.sec8.11]{8.11}.}. We leave to the reader the task of making explicit even a
2-\emph{functorial} dependence by defining, for every pair of
$𝒰$-categories $C$, $C'$, a canonical functor
\hypertarget{sga4.exp.I.equation.8.6.1.8}{}$$
\displaylines{\rlap{$(8.6.1.8)$}\hfill
\Iind(f):\Hhom(C,C')\longrightarrow \Hhom(\Ind(C),
\Ind(C')){\quoi},\hfill}
$$
and specifying the functorial nature of the identical isomorphisms
\hyperlink{sga4.exp.I.equation.8.6.1.6}{(8.6.1.6)} and \hyperlink{sga4.exp.I.equation.8.6.1.7}{(8.6.1.7)}.

\subsubsection{}\etiquette[8.6.2]{I.sec8.6.2}
Return to the functor \hyperlink{sga4.exp.I.equation.8.6.1.1}{(8.6.1.1)}, and consider the corresponding functor
\hyperref[I.prop5.1]{(5.1)}
\hypertarget{sga4.exp.I.equation.8.6.2.1}{}$$
\displaylines{\rlap{$(8.6.2.1)$}\hfill
f_{!}:\widehat{C}\longrightarrow \widehat{C'}{\quoi},\hfill}
$$
and the diagram of functors
\hypertarget{sga4.exp.I.equation.8.6.2.2}{}\[\vcenter{
\xymatrix@C=3.25cm@R=1.25cm{
\Ind(C) \ar[r]^{\ind(f)} \ar[d]_{L_C}& \Ind(C')\ar[d]^{L_{C'}}\\
\hat{C} \ar[r]^{f_!}& \hat{C}' }}\leqno{(8.6.2.2)}
\]
where the vertical arrows denote the canonical functors $L$ of
\hyperlink{sga4.exp.I.equation.8.2.4.8}{(8.2.4.8)}. Since the functor $f_{!}$ ``extends $f$'' and commutes with
inductive limits \hyperlink{sga4.exp.I.proposition.5.4.item.3}{(5.4.3)}, and the functors $L_C$ and $L_{C'}$ commute
with small filtered inductive limits (\hyperref[I.prop8.5.1]{8.5.1}), it
follows from \hyperlink{sga4.exp.I.equation.8.5.3.2}{(8.5.3.2)} that, for every ind-object $(X_i)_{i\in I}$ of
$C$, there is a canonical isomorphism in $\widehat{C'}$
$$
f_{!}(L_C((X_i)_{i\in I}))\simeq \varinjlim_i f_{!}L_C(X_i)\simeq
\varinjlim_i L_{C'}f(X_i)=L_{C'}(f(X_i)_{i\in I})
$$
\ie\pageoriginale a canonical isomorphism
$$
f_{!}L_C(𝒳)\simeq L_{C'}(\Ind(f)(𝒳)){\quoi}.
$$
We leave to the reader the verification that the latter is functorial,
\ie that it corresponds to a canonical isomorphism
\hypertarget{sga4.exp.I.equation.8.6.2.3}{}$$
\displaylines{\rlap{$(8.6.2.3)$}\hfill f_{!}L_C\simeq L_{C'}\circ
\Ind(f){\quoi}.\hfill}
$$
In other words, the \emph{square \hyperlink{sga4.exp.I.equation.8.6.2.2}{(8.6.2.2)} commutes up to canonical
isomorphism.} Taking into account the fact that $f_{!}$ commutes with
small inductive limits, and using \hyperref[I.prop8.5.1]{8.5.1}, we conclude:

\begin{enonce*}{Proposition 8.6.3}\etiquette[8.6.3]{I.prop8.6.3}
Let $f:C\rightarrow C'$ be a functor between $𝒰$-categories. Then the
functor
$$
\Ind(f):\Ind(C)\longrightarrow \Ind(C')
$$
commutes with filtered inductive limits. Moreover, it makes the following
diagram commutative:
\[
\xymatrix@C=2.5cm@R=1.5cm{
C       \ar[r]^{f}       \ar[d]_{c_{C}}& C' \ar[d]^{c_{C'}}\\
\Ind(C) \ar[r]^{\Ind(f)} & \Ind(C')\quad , }
\]
where the vertical arrows $c_C$, $c_{C'}$ are the canonical functors
$\hyperlink{sga4.exp.I.equation.8.4.1}{(8.4.1)}$.
\end{enonce*}
The last assertion is trivial from the definitions and has been included
for ease of reference. Notice moreover that the properties stated in
\hyperref[I.prop8.6.3]{8.6.3} \emph{characterize the functor} $\Ind(f)$ up to unique
isomorphism,\pageoriginale as the functor induced by $f_{!}$ \hyperlink{sga4.exp.I.equation.8.6.2.1}{(8.6.2.1)},
as follows from the proof of \hyperlink{sga4.exp.I.equation.8.6.2.3}{(8.6.2.3)} just given.

\begin{enonce*}{Proposition 8.6.4}\etiquette[8.6.4]{I.prop8.6.4}
The notation is that of $\hyperref[I.prop8.6.3]{8.6.3}$.
\begin{enumerate}
\renewcommand{\theenumi}{\alph{enumi}}
\renewcommand{\labelenumi}{\theenumi)}
\item The functor $\Ind(f)$ is faithful (\resp fully faithful) if and
only if $f$ is.

\item The functor $\Ind(f)$ is an equivalence of categories if and only
if $f$ is fully faithful and every object of $C'$ is isomorphic to a
direct factor \hyperref[I.sec10.6]{(10.6)} of an object in the image of $f$.
\end{enumerate}
\end{enonce*}

\noindent \emph{Proof.}
\begin{enumerate}
\renewcommand{\theenumi}{\alph{enumi}}
\renewcommand{\labelenumi}{\theenumi)}
\item Necessity clearly follows from the fact that $f$ is induced by
$\Ind(f)$. For sufficiency, it is enough to use the form \hyperlink{sga4.exp.I.equation.8.6.1.3}{(8.6.1.3)} of
$\Ind(f)$ on the sets $\Hom(𝒳,𝒴)$, recalling that filtered inductive
limits of sets and arbitrary projective limits take monomorphisms to
monomorphisms and isomorphisms to isomorphisms.

\item We may already suppose that $f$, hence $\Ind(f)$, is fully faithful.
Since every object of $\Ind(C')$ is a small filtered inductive limit of
objects of $C'$, the full faithfulness of $\Ind(f)$ implies that, for this
functor to be essentially surjective, it is equivalent that every object
$X'$ of $C'$ belong to the essential image. But if there is an
isomorphism
$$
X'\xrightarrow{\sim} \textrm{« }{\varinjlim}\textrm{ »\,}f(X_i){\quoi},
$$
this isomorphism factors through one of the $f(X_i)$, showing that $X'$
is a direct factor of $f(X_i)$ and proving necessity in b). For
sufficiency, using the full faithfulness of $\Ind(f)$, the assertion
follows immediately.
\end{enumerate}

\begin{enonce*}{Corollary 8.6.5}\etiquette[8.6.5]{I.cor8.6.5}
In\pageoriginale $\Ind(E)$ the images of projectors
(\hyperref[I.sec10.6]{10.6}) are representable, and the functor
$\Ind(f):\Ind(C)\rightarrow \Ind(C')$ commutes with the formation of
these images.
\end{enonce*}

Indeed, this follows from the fact that small filtered inductive limits
are representable in $\Ind(C)$ (\hyperref[I.prop8.5.1]{8.5.1}) and that
$\Ind(f)$ commutes with them (\hyperref[I.prop8.6.3]{8.6.3}), taking into
account that the image of a projector $p:X\rightarrow X$ may be
interpreted as the inductive limit of a filtered inductive system indexed
by $\NN$
$$
X\xrightarrow{P}X\xrightarrow{P}X\longrightarrow\ldots
$$
or, equivalently, as the inductive limit of the evident functor on the
filtered category $P$ having one object and one nonidentity arrow $\Pi$
such that $\Pi^2=\Pi$.


% END INCLUDED SGA4_English_current_buildable_source_20260906/source/01/drafts/SGA4_I_LINES_3599_3796_English.texfrag


% BEGIN INCLUDED SGA4_English_current_buildable_source_20260906/source/01/drafts/SGA4_I_LINES_3797_3955_English.texfrag
\subsection{}\etiquette[8.7]{I.sec8.7}
\emph{The functor $\varinjlim_C:\Ind(C)\rightarrow C$.
Universal characterizations of the category $\Ind(C)$.}

\subsubsection{}\etiquette[8.7.1]{I.sec8.7.1}
Let us return to a $𝒰$-category $C$ and the canonical functor
\hypertarget{sga4.exp.I.equation.8.7.1.1}{}$$
\displaylines{\rlap{$(8.7.1.1)$}\hfill c:C\longrightarrow
\Ind(C){\quoi}.\hfill}
$$
Let $𝒳=(X_i)_{i\in I}$ be an object of $\Ind(C)$. It follows
immediately from the definition of representable inductive limits
(\hyperref[I.def2.1]{2.1}, \hyperref[I.sec2.1.1]{2.1.1}) that
$\varinjlim_iX_i$ \emph{is representable in $C$ if and only if the left
adjoint of $c$ is defined at $𝒳$}, and that in this case
$\varinjlim_iX_i$ (computed in $C$) is precisely the value at $𝒳$ of
that left adjoint:
\hypertarget{sga4.exp.I.equation.8.7.1.2}{}$$
\displaylines{\rlap{$(8.7.1.2)$}\hfill \Hom_C(\varinjlim_i X_i,
Y)\simeq \Hom_{\Ind(C)}((X_i)_{i\in I}, c(Y)){\quoi}.\hfill}
$$
If\pageoriginale we denote by
\hypertarget{sga4.exp.I.equation.8.7.1.3}{}$$
\displaylines{\rlap{$(8.7.1.3)$}\hfill \Ind(C)'\subset
\Ind(C)\hfill}
$$
the full subcategory of $\Ind(C)$ formed by the ind-objects of $C$ that
admit an inductive limit \emph{in} $C$, the preceding observation shows
that this subcategory is \emph{strictly full}, and that there is a
canonical functor
\hypertarget{sga4.exp.I.equation.8.7.1.4}{}$$
\displaylines{\rlap{$(8.7.1.4)$}\hfill
\textstyle{\varinjlim}_C:\Ind(C)'\longrightarrow C{\quoi},\hfill}
$$
whose value at each object $(X_i)_{i\in I}$ of $\Ind(C)'$ is its
inductive limit in $C$. In the special case where small filtered
inductive limits are representable in $C$, one therefore obtains a
natural functor
\hypertarget{sga4.exp.I.equation.8.7.1.5}{}$$
\displaylines{\rlap{$(8.7.1.5)$}\hfill
{\varinjlim}_C:\Ind(C)\longrightarrow C{\quoi}.\hfill}
$$

\setcounter{paragraph}{5}
\paragraph{}\etiquette[8.7.1.6]{I.sec8.7.1.6}
Of course, the preceding functors may also be defined on categories of
the form $\Ind_{𝒱}(C,𝒰)'$ and $\Ind_{𝒱}(C,𝒰)$, but in view of
$\hyperlink{sga4.exp.I.equation.8.2.4.6}{(8.2.4.6)}$ they are already determined (up to unique isomorphism) by
the preceding functors corresponding to the case $𝒱=𝒰$.

\paragraph{}\etiquette[8.7.1.7]{I.sec8.7.1.7}
The preceding construction of $\varinjlim_C$ as a left adjoint functor
immediately implies that this functor \emph{commutes with arbitrary
inductive limits}, and in particular with small filtered inductive limits
(the latter being representable in $\Ind(C)$
(\hyperref[I.prop8.5.1]{8.5.1})). Notice also that $\Ind(C)'$ always
contains the essential image of $c$ (formed by the essentially constant
ind-objects), and that there is a canonical isomorphism, functorial in
$X\in\ob\Ind(C)'$:
\hypertarget{sga4.exp.I.equation.8.7.1.8}{}$$
\displaylines{\rlap{$(8.7.1.8)$}\hfill {\varinjlim}_C c(X)\simeq
X{\quoi},\hfill}
$$
\ie,\pageoriginale in the favorable case where $C$ is stable under small
inductive limits, there is a canonical isomorphism
\hypertarget{sga4.exp.I.equation.8.7.1.9}{}$$
\displaylines{\rlap{$(8.7.1.9)$}\hfill {\varinjlim}_C \circ c\simeq
\id_C{\quoi}.\hfill}
$$

\subsubsection{}\etiquette[8.7.2]{I.sec8.7.2}

Now consider a functor
\hypertarget{sga4.exp.I.equation.8.7.2.1}{}$$
\displaylines{\rlap{$(8.7.2.1)$}\hfill f:C\longrightarrow
E{\quoi},\hfill}
$$
where $E$ is a $𝒰$-category in which small inductive limits are
representable, so that there is a functor
$$
{\varinjlim}_E:\Ind(E)\longrightarrow E{\quoi}.
$$
Since we have also defined \hyperref[I.sec8.6]{(8.6)}
$$
\Ind(f):\Ind(C)\longrightarrow \Ind(E){\quoi},
$$
we may consider the composite
\hypertarget{sga4.exp.I.equation.8.7.2.2}{}$$
\displaylines{\rlap{$(8.7.2.2)$}\hfill \bar{f}={\varinjlim}_E\circ
\Ind(f):\Ind(C)\longrightarrow E{\quoi},\hfill}
$$
sometimes called the \emph{canonical extension of $f$ to ind-objects}
(but not to be confused with $\Ind(f)$!). As a composite of two functors
that commute with small filtered inductive limits
(\hyperref[I.prop8.6.3]{8.6.3}, \hyperref[I.sec8.7.1.7]{8.7.1.7}), this
functor itself commutes with small filtered inductive limits. Moreover,
the isomorphism \hyperlink{sga4.exp.I.equation.8.7.1.9}{(8.7.1.9)} applied to $E$, together with \hyperlink{sga4.exp.I.equation.8.6.2.3}{(8.6.2.3)}, shows
that $\bar f$ \emph{extends} $f$ up to isomorphism, \ie there is a
canonical isomorphism
\hypertarget{sga4.exp.I.equation.8.7.2.3}{}$$
\displaylines{\rlap{$(8.7.2.3)$}\hfill \bar{f}\circ c_C\simeq
f{\quoi}.\hfill}
$$
It\pageoriginale is moreover fairly clear, in view of the fact that every
object of $\Ind(C)$ is a small filtered inductive limit of objects of $C$,
that the two preceding properties still \emph{characterize} $\bar f$ up
to unique isomorphism. This can be sharpened to obtain at the same time a
universal characterization, up to equivalence, of $\Ind(C)$ among the
$𝒰$-categories $E$ in which small filtered inductive limits are
representable. If $E$ and $F$ are two such $𝒰$-categories, denote by
\hypertarget{sga4.exp.I.equation.8.7.2.4}{}$$
\displaylines{\rlap{$(8.7.2.4)$}\hfill \Hhom(E,F)'\subset
\Hhom(E,F)\hfill}
$$
the full subcategory of $\Hhom(E,F)$ formed by the functors that commute
with small filtered inductive limits. Then:

\begin{enonce*}{Proposition 8.7.3}\etiquette[8.7.3]{I.prop8.7.3}
Let $C$ be a $𝒰$-category, and use the notation \hyperlink{sga4.exp.I.equation.8.7.2.4}{(8.7.2.4)} above. Then
the canonical functor
$$
c_C:C\longrightarrow \Ind(C)
$$
is $2$-universal among the functors with source $C$ and target a
$𝒰$-category having representable small filtered inductive limits. In
other words, $\Ind(C)$ is such a category
(\hyperref[I.sec8.2.5]{8.2.5}, \hyperref[I.prop8.5.1]{8.5.1}), and if $E$
is such a category, the functor
\hypertarget{sga4.exp.I.equation.8.7.3.1}{}$$
\displaylines{\rlap{$(8.7.3.1)$}\hfill g\longmapsto g\circ
c_C:\Hhom(\Ind(C),E)'\longrightarrow \Hhom(C,E)\hfill}
$$
is an equivalence of categories.
\end{enonce*}

To prove the latter point, it is enough to verify that the map
$f\mapsto\bar f$ defined by \hyperlink{sga4.exp.I.equation.8.7.2.2}{(8.7.2.2)} can be refined to a \emph{functor}
\hypertarget{sga4.exp.I.equation.8.7.3.2}{}$$
\displaylines{\rlap{$(8.7.3.2)$}\hfill f\longmapsto
\bar{f}={\varinjlim}_E\cdot \Ind(f):\Hhom(C,E)\longrightarrow
\Hhom(\Ind(C),E)'{\quoi},\hfill}
$$
and\pageoriginale that the latter is a quasi-inverse of \hyperlink{sga4.exp.I.equation.8.7.3.1}{(8.7.3.1)}. The
details of the essentially trivial verification are left to the reader.
Alternatively, one may invoke \hyperref[I.prop7.8]{7.8} to conclude first that \hyperlink{sga4.exp.I.equation.8.7.3.1}{(8.7.3.1)} is
fully faithful, and then \hyperlink{sga4.exp.I.equation.8.7.2.3}{(8.7.2.3)} to conclude that it is essentially
surjective.

\setcounter{subsubsection}{3}

% END INCLUDED SGA4_English_current_buildable_source_20260906/source/01/drafts/SGA4_I_LINES_3797_3955_English.texfrag


% BEGIN INCLUDED SGA4_English_current_buildable_source_20260906/source/01/drafts/SGA4_I_LINES_3956_4068_English.texfrag
\subsubsection{}\etiquette[8.7.4]{I.sec8.7.4}
Return to a functor between $𝒰$-categories
\hypertarget{sga4.exp.I.equation.8.7.4.1}{}$$
\displaylines{\rlap{$(8.7.4.1)$}\hfill f:C\longrightarrow
E{\quoi},\hfill}
$$
where $E$ is a $𝒰$-category in which small filtered inductive limits
are representable, giving a functor
\hypertarget{sga4.exp.I.equation.8.7.4.2}{}$$
\displaylines{\rlap{$(8.7.4.2)$}\hfill \bar{f}:(X_i)_{i\in I}
\longmapsto \mathop{\lim}_i f(X_i):\Ind(C)\longrightarrow
E{\quoi}.\hfill}
$$
We propose to study conditions on $f$ ensuring that $\bar f$ is fully
faithful, \resp an equivalence of categories. Since $f$ is isomorphic to
the composite $\bar f\circ c_C$, and $c_C:C\rightarrow\Ind(C)$ is fully
faithful, we see that for $\bar f$ to be fully faithful, it is
\emph{necessary} that $f$ be so. Replacing $C$ by its essential image in
$E$, we therefore lose essentially no generality by supposing that $f$ is
the inclusion functor of a subcategory $C$ of $E$, as we shall do below
to simplify the notation.

\begin{enonce*}{Proposition 8.7.5}\etiquette[8.7.5]{I.prop8.7.5}
The notation is that of \hyperref[I.sec8.7.4]{8.7.4}.
\begin{enumerate}
\renewcommand{\theenumi}{\alph{enumi}}
\renewcommand{\labelenumi}{\theenumi)}
\item The functor $\bar f$ is fully faithful if and only if:
\begin{itemize}
\item[(i)] Every object $X$ of $C$ satisfies the following condition:

{\rm(PF)} For every small filtered inductive system $(Y_i)_{i\in I}$ in
$C$ whose inductive limit $\varinjlim_iY_i$ exists in\pageoriginale $E$,
the canonical map
\hypertarget{sga4.exp.I.equation.8.7.5.1}{}$$
\displaylines{\rlap{$(8.7.5.1)$}\hfill \varinjlim_i \Hom
(X,Y_i)\longrightarrow \Hom(X,\varinjlim_i Y_i)\hfill}
$$
is bijective.
\end{itemize}

\item The functor $\bar f$ is an equivalence of categories if and only if
$C$ satisfies condition {\rm(i)} and the following two conditions:
\begin{itemize}
\item[(ii)] $C$ is a generating subcategory of $E$ by strict
epimorphisms (\hyperref[I.def7.1]{7.1}).

\item[\rm(iii)] For every object $X$ of $E$, the full subcategory
$C_{/X}$ of $E_{/X}$ formed by the objects $X'$ over $X$ whose source is
in $\Ob C$ is filtered and essentially small.
\end{itemize}
\end{enumerate}

Condition {\rm(iii)} is satisfied in particular if $C$ is equivalent to
a small category, the finite inductive limits in $C$ are representable,
and the inclusion functor $f:C\rightarrow E$ commutes with them (\ie for
every finite category $J$ and every functor $\varphi:J\rightarrow C$, the
inductive limit of $f\varphi$ is representable and is isomorphic in $E$
to an object of $C$).
\end{enonce*}


\noindent \emph{Proof}.
\begin{enumerate}
\renewcommand{\theenumi}{\alph{enumi}}
\renewcommand{\labelenumi}{\theenumi)}
\item With the notation of condition (i), let $𝒴$ denote the ind-object
$(Y_i)$ and let $𝒳$ denote the constant ind-object defined by $X$. Then
\hyperlink{sga4.exp.I.equation.8.7.5.1}{(8.7.5.1)} is precisely the canonical map
$$
\Hom(𝒳, 𝒴)\longrightarrow \Hom
(\bar{f}(𝒳),\bar{f}(𝒴)){\quoi},
$$
so if $\bar f$ is fully faithful, this map is indeed bijective.
Conversely, if $𝒳=(X_j)_{j\in J}$ and $𝒴=(Y_i)_{i\in I}$ are arbitrary
ind-objects of $E$, then the map
$u:\Hom(𝒳,𝒴)\rightarrow\Hom(\bar f(𝒳),\bar f(𝒴))$ is\pageoriginale the
projective limit over $j$ of the maps
$u_j:\Hom(𝒳_j,𝒴)\rightarrow\Hom(\bar f(𝒳_j),𝒴)$, where $𝒳_j$ is the
constant ind-object with value $X_j$. Thus $u$ is bijective if the $u_j$
are, and (i) therefore implies that $\bar f$ is fully faithful.

\item Suppose that (i), (ii), and (iii) hold, and let us prove that
$\bar f$ is an equivalence. It remains to prove that it is essentially
surjective, hence that every $X\in\ob E$ belongs to the essential image.
But hypothesis (ii) means that $\varinjlim_{C_{/X}}Z\rightarrow X$ is an
isomorphism, while (iii) says that the category $C_{/X}$ is filtered and
essentially small. Thus $X$ is the image of the ind-object of $E$ defined
by the natural functor $C_{/X}\rightarrow E$. Conversely, suppose that
$\bar f$ is an equivalence. In view of a), it remains to verify (ii) and
(iii). We may suppose that $E=\Ind(C)$, with $C$ identified with the
subcategory $c(C)$ of $E$. But we know (\hyperref[I.them8.3.3]{8.3.3} (ii)) that for every
$X\in\Ind(C)$, identified if desired with the functor $F$ on $C^\circ$
that it ind-represents, $C_{/X}=C_{/F}$ is an essentially small filtered
category (which proves (iii)), and that the inductive limit in
$\widehat{C}$ of the functor $C_{/F}\rightarrow\widehat{C}$ is $F$. It
is therefore also so in the full subcategory $E$ of $\widehat C$, since
$F=X$ belongs to $E$, which proves (ii). It remains to prove the last
assertion of b). But the hypothesis on $C$ clearly implies that finite
inductive limits are representable in the category $C_{/X}$; \textit{a
fortiori}, the category $C_{/X}$ is filtered.
\end{enumerate}


% END INCLUDED SGA4_English_current_buildable_source_20260906/source/01/drafts/SGA4_I_LINES_3956_4068_English.texfrag


% BEGIN INCLUDED SGA4_English_current_buildable_source_20260906/source/01/drafts/SGA4_I_LINES_4069_4270_English.texfrag
\begin{enonce*}[remark]{Remark 8.7.5.2}\etiquette[8.7.5.2]{I.rem8.7.5.2}
The\pageoriginale preceding argument shows more generally that, when
condition (i) is satisfied, the essential image of the fully faithful
functor $\bar f$ \hyperlink{sga4.exp.I.equation.8.7.4.2}{(8.7.4.2)} consists of the $X\in\ob E$ such that the
category $C_{/X}$ is filtered and essentially small (a condition
automatically satisfied if $C$ satisfies the conditions stated at the end
of b)), \emph{and} such that the inductive limit of the canonical functor
$C_{/X}\rightarrow E$ is $X$.
\end{enonce*}

\begin{enonce*}{Corollary 8.7.6}\etiquette[8.7.6]{I.cor8.7.6}
Let $E_{PF}$ be the full subcategory of $E$ formed by the objects
satisfying condition $(PF)$ of \hyperref[I.prop8.7.5]{8.7.5} {\rm a)}.
Then, for every functor $J\rightarrow E_{PF}$, where $J$ is a finite
category, that admits an inductive limit $X$ in $E$, one has
$X\in\ob E_{PF}$. The canonical functor
$(X_i)_{i\in I}\rightarrow\displaystyle{\mathop{\varinjlim}_i}{}_EX_i$
deduced from the inclusion $g:E_{PF}\hookrightarrow E$
\hypertarget{sga4.exp.I.equation.8.7.6.1}{}$$
\displaylines{\rlap{$(8.7.6.1)$}\hfill
\bar{g}:\Ind(E_{PF})\longrightarrow E\hfill}
$$
is fully faithful. If finite inductive limits are representable in $E$
and $E_{PF}$ is equivalent to a small category, then the essential image
of the functor $\bar f$ consists of the objects $X$ of $E$ such that the
canonical morphism $\varinjlim_{(E_{PF})_{/X}}Y\rightarrow X$ is an
isomorphism. If, moreover, the subcategory $E_{PF}$ of $E$ is generating
by strict epimorphisms (\hyperref[I.def7.1]{7.1}), then the functor
\hyperlink{sga4.exp.I.equation.8.7.6.1}{(8.7.6.1)} is an equivalence of categories.
\end{enonce*}

All the assertions are evident, in the order in which they are stated,
in view of \hyperref[I.prop8.7.5]{8.7.5}, using
\hyperref[I.prop2.8]{2.8} for the first assertion.

\begin{enonce*}{Corollary 8.7.7}\etiquette[8.7.7]{I.cor8.7.7}
Suppose\pageoriginale that finite inductive limits are representable in
$E$. Let $C$ be a full subcategory of $E$, equivalent to a small
category, and consider the functor
$(X_i)_{i\in I}\rightarrow
{\displaystyle{\mathop{\varinjlim}_i}{}_E}X_i$:
\hypertarget{sga4.exp.I.equation.8.7.7.1}{}$$
\displaylines{\rlap{$(8.7.7.1)$}\hfill \bar{f}:\Ind(C)\rightarrow
E{\quoi}.\hfill}
$$
Let $E_{PF}$ be as in $\hyperref[I.cor8.7.6]{8.7.6}$. The following conditions are equivalent:
\begin{itemize}
\item[(i)] The functor $\bar f:\Ind(C)\rightarrow E$ is an equivalence.

\item[(ii)] The category $C$ is generating in $E$ by strict
epimorphisms, is contained in $E_{PF}$, and every object of $E_{PF}$ is
isomorphic in $E$ (or in $E_{PF}$, which is equivalent
(\hyperref[I.cor8.7.6]{8.7.6})) to a direct factor of an object of $C$.

When the subcategory $C$ of $E$ is stable under direct factors
(\hyperref[I.sec10.6]{10.6}), the preceding conditions are also
equivalent to the following:

\item[(ii bis)] $C=E_{PF}$, and $C$ is generating in $E$ by strict
epimorphisms.

\item[(iii)] The subcategory $C$ of $E$ is generating by strict
epimorphisms, is contained in $E_{PF}$, and finite inductive limits are
representable in it.

When, moreover, finite inductive limits are representable in $E$, these
conditions are also equivalent to

\item[(iii bis)] The subcategory $C$ of $E$ is generating by strict
epimorphisms, is contained in $E_{PF}$, and is stable in $E$ under finite
inductive limits (or, equivalently, under finite sums and cokernels of
double arrows).
\end{itemize}
\end{enonce*}

We\pageoriginale already know (\hyperref[I.prop8.7.5]{8.7.5}) that
condition (i) implies $C\subset E_{PF}$ and that $C$ is generating by
strict epimorphisms. Let us also prove that every object $X$ of $E_{PF}$
is then isomorphic to a direct factor of an object of $C$. Indeed, we know
that $X$ is a filtered inductive limit $\varinjlim_iX_i$ in $E$ of
objects of $C$, and by the definition of $E_{PF}$, for a suitable $i$ the
canonical homomorphism $X_i\rightarrow X$ has a left inverse, so $X$ is
indeed a direct factor of the object $X_i$ of $C$. This proves
(i) $\Rightarrow$ (ii) (without any hypothesis on $C$ or $E$, moreover).
Let us prove (ii) $\Rightarrow$ (i). Since $C$ is equivalent to a small
category and every object of $E_{PF}$ is a direct factor of an object of
$C$, we see that $E_{PF}$ is also equivalent to a small category. Thus,
by \hyperref[I.cor8.7.6]{8.7.6}, the functor \hyperlink{sga4.exp.I.equation.8.7.6.1}{(8.7.6.1)} is an equivalence
of categories, and we conclude by \hyperref[I.prop8.6.4]{8.6.4} b) applied
to the inclusion $E\hookrightarrow E_{PF}$.

When every direct factor in $E$ of an object of $C$ belongs to $C$, it is
clear that (ii) $\Leftrightarrow$ (ii bis). On the other hand, (ii bis)
implies (iii) \resp (iii bis) by \hyperref[I.cor8.7.6]{8.7.6}. Finally, (iii) (and
\textit{a fortiori} (iii bis)) implies (i) by
\hyperref[I.prop8.7.5]{8.7.5}. This completes the proof.

\begin{enonce*}[remark]{Exercise 8.7.8}\etiquette[8.7.8]{I.exer8.7.8}
(\emph{Karoubi envelopes}). Let $C$ be a $𝒰$-category, let
$E=\Ind(C)$, and let $E_{PF}\supset C$ be the full subcategory of $E$
defined in \hyperref[I.cor8.7.6]{8.7.6}.
\begin{enumerate}
\renewcommand{\theenumi}{\alph{enumi}}
\renewcommand{\labelenumi}{\theenumi)}
\item Prove that images of projectors are representable in $E_{PF}$ and
that every object of $E_{PF}$ is a direct factor of an object of $C$.

\item Call a category $F$ \emph{Karoubian} if images of projectors are
representable in it. Let
$$
\varphi:C\longrightarrow K
$$
be\pageoriginale a fully faithful functor such that $K$ is Karoubian and
every object of $K$ is a direct factor of an object in the image of $C$.
Prove that $\varphi$ is 2-\emph{universal} among functors from $C$ with
values in Karoubian categories, more precisely: for every Karoubian
category $F$, the functor
$$
g\longmapsto g\circ \varphi :\Hhom(K,F)\longrightarrow \Hhom(K,F)
$$
is an equivalence of categories. (Use the fact that every functor commutes
with images of projectors.) The category $K$ equipped with $\varphi$,
determined up to equivalence (itself determined up to unique isomorphism)
by the preceding properties, is called the \emph{Karoubi envelope}
$\widetilde C$ of $C$. (Compare also IV \hyperref[I.cor7.5]{7.5}.)

\item Deduce from a) and b) that $\Ind(C)_{PF}$, equipped with the functor
$C\rightarrow\Ind(C)_{PF}$ induced by $c:C\rightarrow\Ind(C)$, is a
Karoubi envelope of $C$.

\item Show that every functor $f:C\rightarrow C'$ extends in an
essentially unique way to a functor
$\widetilde f:\widetilde C\rightarrow\widetilde{C'}$ of Karoubi
envelopes, and that if these Karoubi envelopes are taken as in c),
$\widetilde f$ is the functor induced by
$\Ind(f):\Ind(C)\rightarrow\Ind(C')$. Show that the conditions of
\hyperref[I.sec8.5.4]{8.5.4} %%
b) %%
are also equivalent to the following: $\widetilde f$ is an equivalence of
categories.
\end{enumerate}
\end{enonce*}

\begin{enonce*}[remark]{Exercise 8.7.9}\etiquette[8.7.9]{I.exer8.7.9}
Let $E$ be a $𝒰$-category.
\begin{enumerate}
\renewcommand{\theenumi}{\alph{enumi}}
\renewcommand{\labelenumi}{\theenumi)}
\item Show that the following conditions are equivalent:

\begin{itemize}
\item[(i)] $E$ is equivalent to a category of the form $\Ind(C)$, where
$C$ is a $𝒰$-category.

\item[(i bis)] As\pageoriginale in (i), with $C$ moreover Karoubian
(\hyperref[I.exer8.7.8]{8.7.8} b)).

\item[(ii)] The subcategory $E_{PF}$ of $E$
(\hyperref[I.cor8.7.6]{8.7.6}) is generating by strict epimorphisms, and
for every object $X$ of $E$, $(E_{PF})_{/X}$ is an essentially small
filtered category.
\end{itemize}

Show that if $C$ is a \emph{Karoubian} $𝒰$-category such that $E$ is
equivalent to $\Ind(C)$, then $C$ is equivalent to $E_{PF}$ (by an
equivalence determined up to unique isomorphism). Establish a
2-equivalence between the 2-category formed by the categories $E$
satisfying the preceding conditions and the functors between them that
\emph{commute with small filtered inductive limits}, and the 2-category
formed by \emph{Karoubian} $𝒰$-categories and arbitrary functors between
them.

\item Show that the following conditions are equivalent:

\begin{itemize}
\item[(i)] $E$ is equivalent to a category of the form $\Ind(C)$, where
$C$ is a $𝒰$-category in which finite inductive limits are
representable.

\item[(i bis)] As in (i), but with $C$ moreover Karoubian.

\item[(ii)] Finite inductive limits in $E$ are representable, the
subcategory $E_{PF}$ of $E$ is generating by strict epimorphisms, and for
every object $X$ of $E$, there exists a small subset $I$ of
$\Ob E_{PF/X}$ such that every object of $E_{PF/X}$ is bounded above by
an object of $I$. (Use \hyperref[I.prop8.9.5]{8.9.5} b).)
\end{itemize}
\end{enumerate}
\end{enonce*}


% END INCLUDED SGA4_English_current_buildable_source_20260906/source/01/drafts/SGA4_I_LINES_4069_4270_English.texfrag


% BEGIN INCLUDED SGA4_English_current_buildable_source_20260906/source/01/drafts/SGA4_I_LINES_4271_4625_English.texfrag
\subsection{Indexed representation of a functor $J\rightarrow
\Ind(C)$}\etiquette[8.8]{I.sec8.8}\pageoriginale

\subsubsection{}\etiquette[8.8.1]{I.sec8.8.1}
Let
$$
\varphi:J\longrightarrow \Ind(C)\text{ \ie } \varphi\in \Ob
\Hhom(J,\Ind(C))
$$
be a functor, where $C$ is a $𝒰$-category. Recall \hyperref[I.sec8.5.4]{(8.5.4)}
that an \emph{indexed representation of} $\varphi$ is, by
definition, an ind-object of $\Hhom(J,C)$
$$
\psi:I\longrightarrow \Hhom(J,C){\quoi},
$$
with I an essentially small filtered category, together with a specified
isomorphism from its inductive limit in $\Hhom(J,\Ind(C))$ to $\varphi$.
Thus $\varphi$ admits an indexed
representation if and only if it is isomorphic to an essentially
small filtered inductive limit in $\Hhom(J,\Ind(C))$
of objects of the full subcategory $\Hhom(J,C)$. We shall say that
the category $J$ is \emph{admissible for} $C$ if every functor
$\varphi:J\rightarrow \Ind(C)$ admits an indexed representation.

Since in $\Ind(C)$ small filtered inductive limits are representable
(\hyperref[I.prop8.5.1]{8.5.1}), the same is true in
$\Hhom(J,\Ind(C))$, and they are computed “argument by argument”.
Consequently, the inclusion functor
\hypertarget{sga4.exp.I.equation.8.8.1.1}{}$$
\displaylines{\rlap{$(8.8.1.1)$}\hfill \Hhom(J,C)\hookrightarrow
\Hhom(J,\Ind(C))\hfill}
$$
extends canonically to a functor
\hypertarget{sga4.exp.I.equation.8.8.1.2}{}$$
\displaylines{\rlap{$(8.8.1.2)$}\hfill
\Ind(\Hhom(J,C))\longrightarrow \Hhom(J,\Ind(C))\hfill}
$$
\hyperref[I.sec8.7.2]{(8.7.2)}. The objects of the essential image of this functor are
       then precisely the $\varphi$ that admit an
       indexed representation; thus, to say that $J$ is admissible
       for $C$ also means that the functor \hyperlink{sga4.exp.I.equation.8.8.1.2}{(8.8.1.2)} is
       essentially surjective. In this connection, we note:

\begin{enonce*}{Proposition 8.8.2}\etiquette[8.8.2]{I.prop8.8.2}
If\pageoriginale the category $J$ is equivalent to a
finite category,
the canonical functor $\hyperlink{sga4.exp.I.equation.8.8.1.2}{(8.8.1.2)}$ is fully faithful; thus
$J$ is admissible for $C$ if and only if this functor is an equivalence
of categories.
\end{enonce*}

By \hyperref[I.prop8.7.5]{8.7.5} a), it suffices to prove that for
every small filtered inductive system $(\varphi_i)_{i\in I}$ in
$\Hhom(J,C)$ and every object $\varphi$ of $\Hhom(C,J)$, the
canonical map
\hypertarget{sga4.exp.I.equation.8.8.2.1}{}$$
\displaylines{\rlap{$(8.8.2.1)$}\hfill \varinjlim_i \Hom(\varphi,
\varphi_i)\longrightarrow \Hom(\varphi, \varinjlim_i
\varphi_i)\hfill}
$$
is bijective, where $\varinjlim_i \varphi_i$ denotes the inductive
limit taken in $\Hom(J,\Ind(C))$. Now, if $\varphi$, $\psi$
are two functors $J\rightarrow C$, there is an exact diagram
of sets, functorial in $\varphi$ and $\psi$:
$$
\Hom(\varphi,\psi)\longrightarrow \prod_{X\in \Ob J}
\Hom(\varphi(X),\psi(X))\rightrightarrows \prod_{\substack{f\in
F\ell J\\ f:X\rightarrow Y}} \Hom(\varphi(X), \psi (Y)){\quoi}.
$$
Since filtered inductive limits commute with kernels of
double arrows and with finite products, the bijectivity of
\hyperlink{sga4.exp.I.equation.8.8.2.1}{(8.8.2.1)}
follows when $J$ is finite. The case in which $J$ is equivalent to a
finite category reduces immediately to the preceding case.

\begin{enonce*}{Proposition 8.8.3}\etiquette[8.8.3]{I.prop8.8.3}
Let $\varphi:J\rightarrow \Ind(C)$ be a functor, with $J$
equivalent to a small category. For $\varphi$ to admit an
indexed representation, the following two conditions are sufficient,
and they are also necessary when $J$
is equivalent to a finite category:
\begin{enumerate}
\renewcommand{\theenumi}{\alph{enumi}}
\renewcommand{\labelenumi}{\theenumi)}
\item The category $\Hhom(J,C)_{/\varphi}$ (formed by the
arrows in $\Hhom(J,\Ind(C))$ with target $\varphi$ and source in
$\Hom(J,C)$) is filtered.

\item For\pageoriginale every object $j$ of $J$, the functor
\hypertarget{sga4.exp.I.equation.8.8.3.1}{}$$
\displaylines{\rlap{$(8.8.3.1)$}\hfill \psi\longmapsto
\psi(j):\Hhom(J,C)_{/\varphi}\longrightarrow C_{/\varphi(j)}\hfill}
$$
\end{enumerate}
is cofinal, \ie satisfies conditions {\rm %%
F 1), F 2)} of $\hyperref[I.prop8.1.3]{8.1.3}$.
\end{enonce*}

\noindent \emph{Proof}. Suppose a) and b) hold;
let us prove that $\varphi$ admits an indexed representation. We may
clearly suppose that $J$ is small. Let
\hypertarget{sga4.exp.I.equation.8.8.3.2}{}$$
\displaylines{\rlap{$(8.8.3.2)$}\hfill
I=\Hhom(J,C)_{/\varphi}{\quoi},\hfill}
$$
which is a filtered category by hypothesis, and first suppose that
$I$ is essentially small, so that the “inclusion” of $I$
in $\Hhom(J,C)$ defines an ind-object of $\Hhom(J,C)$, $i\mapsto
\psi_i$. Moreover, there is a canonical homomorphism
\hypertarget{sga4.exp.I.equation.8.8.3.3}{}$$
\displaylines{\rlap{$(8.8.3.3)$}\hfill \varinjlim_i
\psi_i\longrightarrow \varphi{\quoi},\hfill}
$$
given argument by argument by the homomorphism
$$
\varinjlim_I \psi_i (j)\longrightarrow
\varphi(j)=\varinjlim_{C_{/\varphi(j)}}X
$$
induced by the functor \hyperlink{sga4.exp.I.equation.8.8.3.1}{(8.8.3.1)}. Since %%
the latter is cofinal, we conclude that \hyperlink{sga4.exp.I.equation.8.8.3.3}{(8.8.3.3)} is an
isomorphism, which proves the assertion. In the case where we do not
suppose $I$ essentially small, it is enough to construct an
essentially small full subcategory $I'$ such that
\hyperlink{sga4.exp.I.equation.8.8.3.3}{(8.8.3.3)} remains an isomorphism when $\varinjlim_{I'}$ is used in
place of $\varinjlim_I$. For this purpose, it is enough that the composite functors
$$
I'\longrightarrow I\longrightarrow C_{/\varphi(j)}
$$
induced\pageoriginale by the functors \hyperlink{sga4.exp.I.equation.8.8.3.1}{(8.8.3.1)} all be
cofinal. We then conclude by the following lemma, whose
proof is immediate from criterion \hyperref[I.prop8.1.3]{(8.1.3 b)}) and is
left to the reader:

\begin{enonce*}{Lemma 8.8.3.4}\etiquette[8.8.3.4]{I.lem8.8.3.4}
Let $I$ be a filtered $𝒰$-category, and let
$$
f_j:I\longrightarrow I_j\quad (j\in J)
$$
be a small family of cofinal functors from $I$ to essentially small
filtered categories. Then there exists a
full subcategory $I'$ of $I$ which is filtered and
essentially small, and such that the functors induced by the
$f_j$ are cofinal.
\end{enonce*}

Finally, let us prove necessity in \hyperref[I.prop8.8.3]{8.8.3} when
$J$ is assumed equivalent to a finite category. An indexed representation
of $\varphi$ by means of an essentially small filtered indexing
category $I$ defines a functor $\psi$
\hypertarget{sga4.exp.I.equation.8.8.3.5}{}\[
\xymatrix@C=1.25cm@R=2cm{
 I \ar[rr]^(.4){\psi} \ar@{-->}[rrd]^{\psi_j}&&  \Hhom(J,C)_{/\varphi} \ar@{-->}[d] &\\
                        &     &   C_{/\varphi(j)}\quad ,\\
}\tag{8.8.3.5}
\]
and it is enough to prove that $\psi$ and each of the functors $\psi_j$
induced by it are cofinal: by \hyperref[I.prop8.1.3]{8.1.3} b), it
will then indeed follow that $\Hhom(J,C)_{/\varphi}$ is filtered, and by
definition \hyperref[I.def8.1.1]{8.1.1} that the vertical arrow in \hyperlink{sga4.exp.I.equation.8.8.3.5}{(8.8.3.5)} is
also cofinal. By \hyperref[I.prop8.8.2]{8.8.2}, we may identify $\varphi$ with
an object of $\Ind(E)$, where $E=\Hhom(J,C)$, and the fact that the functor
$$
\psi:I\longrightarrow E_{/\varphi}{\quoi},
$$
arising\pageoriginale from the ind-object $\varphi$ of $E$ indexed by $I$
is cofinal is a general fact, which follows immediately
from criteria (F 1) and (F 2) of \hyperref[I.prop8.1.3]{8.1.3}. The same
argument (with $E$, $\varphi$ replaced by $C$, $\varphi(j)$)
shows that $\psi_j$ is cofinal, which completes the proof.

\begin{enonce*}[remark]{Remark 8.8.4}\etiquette[8.8.4]{I.rem8.8.4}
\begin{enumerate}
\renewcommand{\theenumi}{\alph{enumi}}
\renewcommand{\labelenumi}{\theenumi)}
\item In the case where $J$ is equivalent to a finite
category, if $\varphi$ admits an indexed representation, we have seen that
\hyperlink{sga4.exp.I.equation.8.8.3.5}{(8.8.3.5)} is a cofinal functor, which implies that the category
$\Hhom(J,C)_{/\varphi}$ is itself essentially small.

\item Suppose that in $C$ finite inductive limits are representable.
Then the same is true in $E=\Hhom(J,C)$, and
these limits are computed argument by argument and are also
inductive limits in $\Hhom(J,\widehat{C})$ and \textit{a
fortiori} in $\Hhom(J,\Ind(C))$. It follows immediately that
condition a) of \hyperref[I.prop8.8.3]{8.8.3} is then automatically
satisfied, and it remains only to consider condition b). I do not know whether
it is automatically satisfied when $J$ is in addition assumed
small.
\end{enumerate}
\end{enonce*}
We now come to the main result of this section:

\begin{enonce*}{Proposition 8.8.5}\etiquette[8.8.5]{I.prop8.8.5}
Let $J$ be a category equivalent to a finite category, and
suppose in addition that $J$ is rigid, \ie that for every $j\in \Ob J$, every
endomorphism of $j$ is the identity. Then $J$ is admissible for
$C$ (for any $𝒰$-category $C$), \ie every
functor $J\rightarrow \Ind(C)$ admits an indexed representation.
\end{enonce*}

Upon \pageoriginale %%
replacing $J$ by an equivalent category if necessary, we may suppose
that $J$ is \emph{skeletal}, \ie that two isomorphic objects of $J$
are equal. Then $J$ is in fact finite. We proceed by
induction on $\text{card}\Ob J$; the case where this number is $\leq 0$
is trivial, so we shall suppose it is $\geq 1$. Let $j_\circ$ be a
\emph{maximal} object of $J$, \ie such that for every arrow
$j_\circ\rightarrow j$, there exists an arrow $j\rightarrow j_\circ$;
in view of the fact that $J$ is rigid, this implies that
$j_\circ\rightarrow j$ is an isomorphism, and since $J$ is
skeletal, it implies that $j_\circ=j$. Let $J'$ therefore be the full
subcategory obtained from $J$ by removing the object $j_\circ$, and let $J''$ be the
full subcategory consisting of the object $j_\circ$. The proposition
then follows from the

\begin{enonce*}{Lemma 8.8.5.1}\etiquette[8.8.5.1]{I.lem8.8.5.1}
Let $J$ be a finite category, and let $J'$ and $J''$ be two full
subcategories such that $\Ob J=\Ob J'\cup \Ob J''$, and such that for every
$j'\in \Ob J'$ and $J''\in \Ob J''$, one has
$\Hom(j'',j')=∅$. If $J'$ and $J''$ are admissible for
$C$, then so is $J$.
\end{enonce*}

It is enough to prove criteria a) and b) of \hyperref[I.prop8.8.3]{8.8.3},
which amounts to carrying out six elementary verifications, namely the
last two filteredness conditions for
$\Hhom(J,C)_{/\varphi}$, and conditions
F 1) and
F 2) for the functors \hyperlink{sga4.exp.I.equation.8.8.3.1}{(8.8.3.1)}, successively in the cases $j\in \Ob J'$ and
$j\in \Ob J''$; the latter conditions moreover imply
that $\Hhom(J,C)_{/\varphi}$ is nonempty. The rather
tedious verification presents no difficulty and is left to the
reader. (The editor sought in vain for an elegant proof
that would circumvent these unpleasant
verifications.)

\begin{enonce*}[remark]{Exercise 8.8.6}\etiquette[8.8.6]{I.exer8.8.6}
\begin{enumerate}
\renewcommand{\theenumi}{\alph{enumi}}
\renewcommand{\labelenumi}{\theenumi)}
\item Let\pageoriginale $J$, $J'$ be two categories admissible
for $C$, one of them being finite. Prove that $J\times J'$ is
admissible for $C$. (Use \hyperref[I.prop8.8.2]{8.8.2}.)

\item Prove that a small discrete category is admissible for
every category $C$.

\item Let $J$ be a category satisfying the following conditions:
(i) $J$ is rigid, (ii) $\Ob J$ is countable, (iii) For any two
objects $j$, $j'$ of $J$, $\Hom(j,j')$ is finite, (iv)
For every object $j$ of $J$, the set of objects of $J$ that are
bounded above by $j$, \ie that are sources of an arrow with target $j$, is
finite. Show that $J$ is admissible for every
$𝒰$-category $C$. (First show that $J$ can be written
as a filtered union of finite full subcategories $J_n$,
such that every object of $J$ bounded above by an object of $J_n$ belongs to
$J_n$. Then verify criteria a), b) of \hyperref[I.prop8.8.3]{8.8.3}.)
\end{enumerate}
\end{enonce*}

\begin{enonce*}[remark]{Exercise 8.8.7}\etiquette[8.8.7]{I.exer8.8.7}
In our notation we identify an ordered set with the
category it defines.
\begin{enumerate}
\renewcommand{\theenumi}{\alph{enumi}}
\renewcommand{\labelenumi}{\theenumi)}
\item Let $C$ be an ordered set. Let $C^{\widetilde{\phantom{,}}}$
be the set of subsets $A$ of $C$ that are filtered and such that
$x\in A$ and $y\leq x$ imply $y\in A$. Let $L_\circ:C\rightarrow
C^{\widetilde{\phantom{,}}}$ be the map that associates to every $x\in
C$ the set $L_\circ(x)$ consisting of those $y\in A$ such that $y\leq x$. Show
that $L$ is injective and that the order on $C$ is
induced by that on $C^{\widetilde{\phantom{,}}}$. For every $A\in
C^{\widetilde{\phantom{,}}}$, consider the inclusion
$f(A):A\rightarrow C$, %%
which is an ind-object of $C$; for every ind-object $\varphi:I\rightarrow
C$ of $C$, let $g(\varphi)\in C^{\widetilde{\phantom{,}}}$ be the
subset of $C$ consisting of the elements of $C$ bounded above by an element of
the form $\varphi(i)$. Show\pageoriginale that one thus obtains
two mutually quasi-inverse equivalences
\[
\xymatrix@C=1.5cm{ \Ind(C)
\ar@{<->}[r]^(.5){\displaystyle\mathop{f,g}_{\approx}}&
C^{\widetilde{\phantom{a}}}, }
\]
transforming the functor $L_\circ$ into the canonical functor
$L:C\rightarrow \Ind(C)$.

\item Suppose that for every $x\in C$, the set of elements
bounding $x$ from above (\resp from below) is finite. Show that every ind-object
(\resp pro-object) of $C$ is essentially constant, and indeed that
for every such $\varphi:I\rightarrow C$ (\resp
$\varphi:I^\circ\rightarrow C$) there exists an $i_\circ\in \Ob I$ such
that the functor induced on $I_{/i_\circ}$ is constant (\ie
takes every arrow to an isomorphism).

\item Let $C\subset\NN^\circ\times
\NN$ be formed by the pairs of natural numbers $(j,i)$ such that $i\geq
j$. Show that $\NN^{\widetilde{\phantom{,}}}$ is isomorphic to
the ordered set obtained from $\NN$ (identified with an ordered subset
of $\NN^{\widetilde{\phantom{,}}}$ as in a)) by
adjoining a greatest element $\infty$. Show that
$C^{\widetilde{\phantom{,}}}$ is isomorphic to $\NN^\circ\times
\NN^{\widetilde{\phantom{,}}}$. Consider the functor
$$
\varphi:\NN^\circ\longrightarrow C^{\widetilde{\phantom{,}}}\quav
\varphi(j)=(j,\infty){\quoi}.
$$
\end{enumerate}
\end{enonce*}

\noindent Show that:
\begin{enumerate}
\renewcommand{\labelenumi}{\theenumi)}
\item the projective limit of $\varphi$ is not representable in
$C^{\widetilde{\phantom{,}}}$, although filtered projective limits
in $C$ are representable (and are essentially
constant, by b)).

\item For every functor $\psi:\NN^\circ\rightarrow
C$, one has $\Hom(\psi,\varphi)=∅$, and \textit{a fortiori} the
functor $\varphi$ admits no indexed
representation.
\end{enumerate}

\begin{enonce*}[remark]{Exercise 8.8.8}\etiquette[8.8.8]{I.exer8.8.8}
Let\pageoriginale $X=\NN\amalg \NN$ be the disjoint union of two
copies of $\NN$, let $s$ be the symmetry of $X$, and for every $i\in \NN$,
let $X_i$ be the subset $\Delta_{i+1}\amalg\Delta_i$ of $\NN$ (where
$\Delta_i$ denotes the subset of $\NN$ formed by the $j$ such that $j\leq
i$). Let $C$ be the subcategory of the category of subsets
of $X$ whose objects are the $X_i$ and whose arrows are the
maps between the $X_i$ induced by the identity of $X$
or by its symmetry $s$, and let $C^{\widetilde{\phantom{,}}}$ be the
category defined analogously, but with the additional
object $X$.

\begin{enumerate}
\renewcommand{\theenumi}{\alph{enumi}}
\renewcommand{\labelenumi}{\theenumi)}
\item Show that $\Ind(C)$ is equivalent to
$C^{\widetilde{\phantom{,}}}$, with the canonical functor $C\rightarrow
\Ind(C)$ corresponding to the inclusion $C\rightarrow
C^{\widetilde{\phantom{,}}}$. Show that there exists no pair
$(Y,f)$ consisting of an object $Y$ of $C$ and an endomorphism $f$ of $Y$, together
with a morphism of objects with endomorphism in $\Ind(C)$ from $(Y,f)$ to
$(X,s)$.

\item Deduce that if $J$ is the category having a single object
and, in addition to the identity arrow, a single arrow of order 2, then
the functor $J\rightarrow \Ind(C)$ defined by $(X,s)$ admits no
indexed representation.
\end{enumerate}
\end{enonce*}
 

% END INCLUDED SGA4_English_current_buildable_source_20260906/source/01/drafts/SGA4_I_LINES_4271_4625_English.texfrag


% BEGIN INCLUDED SGA4_English_current_buildable_source_20260906/source/01/drafts/SGA4_I_LINES_4626_4853_English.texfrag
\subsection{Exactness properties of $\Ind(C)$}\etiquette[8.9]{I.sec8.9}

\begin{enonce*}{Proposition 8.9.1}\etiquette[8.9.1]{I.prop8.9.1}
Let $C$ be a $𝒰$-category.
\begin{enumerate}
\renewcommand{\theenumi}{\alph{enumi}}
\renewcommand{\labelenumi}{\theenumi)}
\item The canonical functors $L$ \hyperlink{sga4.exp.I.equation.8.2.4.8}{(8.2.4.8)} and $c$ \hyperlink{sga4.exp.I.equation.8.4.1}{(8.4.1)}
$$
C\xrightarrow{c}\Ind(C)\xrightarrow{L}\widehat{C}
$$
commute with projective limits.

\item Suppose that finite inductive limits are representable in $C$
and that $C$ is equivalent to a small
category.
Then small projective limits are representable in
$\Ind(C)$.

\item If\pageoriginale small projective limits
(\resp finite projective limits) are representable in $C$, then the
same is true in $\Ind(C)$. Let $J$ be a category that is finite
and rigid, or discrete; if projective limits of type $J$ are
representable in $C$, projective limits of the same type are
representable in $\Ind(C)$.

\item Small filtered inductive limits in $\Ind(C)$ (which
are representable by virtue of \hyperref[I.prop8.5.1]{8.5.1}) are left
exact, \ie commute with finite projective limits.
\end{enumerate}
\end{enonce*}

\noindent \emph{Proof}.
\begin{enumerate}
\renewcommand{\theenumi}{\alph{enumi}}
\renewcommand{\labelenumi}{\theenumi)}
\item The fact that $L$ commutes with projective limits follows from
the fact that it is fully faithful and from the computation of projective
limits in $\widehat{C}$ ``argument by argument,'' which
implies that, to verify that a projective system of morphisms
$F\rightarrow F_i$ $(i\in I)$ in $\widehat{C}$ exhibits $F$ as a
projective limit of the $F_i$, it suffices to verify that the analogous
assertion holds for the %%
projective systems of maps of sets
$\Hom(X,F)\rightarrow \Hom(X,F_i)$, for every $X\in \Ob C$. But
$C\subset \Ind(C)$.

The fact that $c$ commutes with projective limits follows formally
from the fact that $L$ does so, as does the composite $L\circ
c:C\rightarrow \widehat{C}$.

\item In view of a), the assertion amounts to saying that every
projective limit of ind-representable presheaves is
ind-representable, which follows at once from the criterion
\hyperref[I.them8.3.3]{8.3.3} (v),
since a projective limit of left exact functors
is left exact (which follows from the fact that ``projective
limits commute with projective limits'' (\hyperref[I.sec2.5.0]{2.5.0})).

\item This\pageoriginale follows formally from the analogous
  property of
$\widehat{C}$ (\hyperref[I.cor3.3]{3.3}), and from the fact that $L$ is conservative (being
fully faithful) and commutes with limits of the type under consideration (a) and
\hyperref[I.prop8.5.1]{8.5.1}).

\item It is well known (\cf \hyperref[I.prop2.3]{2.3}) that the
  representability of
finite projective limits is equivalent to that of projective limits
of the types

\begin{center}
empty $, \ldots ,\quad \vcenter{ \xymatrix@C=.25cm@R=.5cm{
                     \bigdot \ar[rd]&       &\bigdot\ar[ld]\\
                             &\bigdot&
}}$
\end{center}

\noindent (corresponding to certain finite ordered sets),
and the representability of small projective limits amounts to
that of products and finite projective limits. Thus the
second assertion made in (c) implies the first.
\end{enumerate}

\emph{Suppose first that $I$ is finite.} Using the result
\hyperref[I.prop8.8.5]{8.8.5} on the representability of %%
functors $\varphi:J\rightarrow \Ind(C)$ in indexed form
(\hyperref[I.sec8.5.4]{8.5.4})
\hypertarget{sga4.exp.I.equation.8.9.1.1}{}$$
\displaylines{\rlap{$(8.9.1.1)$}\hfill \phi:J\times I\longrightarrow
C{\quoi},\hfill}
$$
the existence of $\varprojlim \varphi$ is therefore a special case of the
following more precise and more general result:

\begin{enonce*}{Corollary \sisi{8.8.2}{8.9.2}}\etiquette[8.9.2]{I.cor8.9.2}
Consider a functor $\varphi:J\rightarrow \Ind(C)$ given in
indexed form $\Phi$ \hyperlink{sga4.exp.I.equation.8.9.1.1}{(8.9.1.1)}, where $J$ is a finite
category.
If for every $i\in \Ob I$, the partial functor $j\rightarrow
\Phi(j,i)$ has a projective (\resp inductive) limit representable
in $C$, then $\varphi$ has a projective (\resp inductive) limit
representable in $\Ind(C)$, and there is a canonical isomorphism
\hypertarget{sga4.exp.I.equation.8.9.2.1}{}$$
\displaylines{\rlap{$(8.9.2.1)$}\hfill
\varprojlim\varphi\simeq\textrm{« }{\varinjlim_i}\textrm{ »} \varprojlim_j
\Phi(j,i)\hfill}
$$
(\resp\pageoriginale
\hypertarget{sga4.exp.I.equation.8.9.2.2}{}$$
\displaylines{\rlap{$(8.9.2.2)$} \hfill \varinjlim \varphi \simeq
\textrm{« }{\varinjlim_i}\textrm{ »} \varprojlim_j\Phi(j,i)){\quoi},\hfill}
$$
where $\varprojlim_j$ (\resp. $\varinjlim_j$) is computed in $C$.
\end{enonce*}

For the first formula, one simply uses the fact that in $\Ind(C)$
filtered inductive limits commute with $\varprojlim_j$, and
that $c$ commutes with $\varprojlim_J$ \hyperref[I.prop8.9.1]{(8.9.1 d)} and a)):
$$
\varprojlim \varphi=\varprojlim_j \textrm{« }{\varinjlim_i}\textrm{ »\,} \Phi
(j,i) \xleftarrow{\sim} \textrm{« }{\varinjlim_i}\textrm{ »\,}
{\varprojlim_j}^{\Ind(C)} \Phi(j,i)\simeq \textrm{« }{\varinjlim_i}\textrm{ »}%%
\varprojlim_j{^C \ \Phi(j,i){\quoi}}.
$$
The second is proved in the same way, using the commutation of
inductive limits of type $I$ with inductive limits of type $J$
(\hyperref[I.sec2.5.0]{2.5.0}) and the fact that $c$ commutes with inductive limits of type
$J$, proved in \hyperref[I.rem8.9.4]{8.9.4} a) below.

\emph{The case where $J$ is discrete}. This case is contained in the following more
precise assertion:

\begin{enonce*}{Corollary 8.9.3}\etiquette[8.9.3]{I.cor8.9.3}
Let the $I_\alpha$ be essentially small filtered categories
indexed by a small set $A$, and for every $\alpha \in A$,
let $X(\alpha)=(X(\alpha)_{i_\alpha})_{i_\alpha \in I_\alpha}$ be an
ind-object of $C$ indexed by $I_\alpha$. Let $I$ be the product category
of the $I_\alpha$, and suppose that for every
$\sheaf{i}=(i_\alpha)_{\alpha\in A}\in I$, the product
$\prod_{\alpha\in A}%%
X (\alpha)_{i_\alpha}$ is representable in $C$. Then the product
$\prod_{\alpha \in A} %%
X (\alpha)$ is representable in $\Ind(C)$, and it is
canonically isomorphic to the ind-object indexed by $I$ given by the
formula
\hypertarget{sga4.exp.I.equation.8.9.3.1}{}$$
\displaylines{\rlap{$(8.9.3.1)$}\hfill \prod_{\alpha \in A}\;
\mathop{\textrm{« }\varinjlim\textrm{ »}}_{i_\alpha\in I_\alpha}\;
X(\alpha)_{i_\alpha}\simeq
\mathop{\textrm{« }\varinjlim\textrm{ »}}_{\sheaf{i}=(i_\alpha)_{\alpha\in
A}\in I}\; \left(\prod_{\alpha\in
A}A(\alpha)_{i_\alpha}\right){\quoi},\hfill}
$$
where\pageoriginale the product on the right-hand side is the product
computed in $C$. (Note that $I$ is filtered and essentially
small, since the $I_\alpha$ are.)
\end{enonce*}

Indeed, it is well known (and immediate by reduction to the case in which one
works in the category of sets) that the formula in question
holds when the limits are computed in $\widehat{C}$. The
conclusion then follows from the fact that $L$ commutes with the limits
in question (\hyperref[I.prop8.9.1]{8.9.1} a) and \hyperref[I.prop8.5.1]{8.5.1}).

\begin{enonce*}[remark]{Remarks 8.9.4}\etiquette[8.9.4]{I.rem8.9.4}
The proof given for \hyperref[I.cor8.9.2]{8.9.2} and \hyperref[I.cor8.9.3]{8.9.3} shows, more
generally,
that if for every $i\in \Ob I$, $\varprojlim_j \Phi(j,i)$ (\resp
$\varinjlim_j \Phi(j,i)$) \emph{computed in} $\Ind(C)$ is
representable, then the same is true of $\varprojlim_i\varphi$
(\resp of $\varinjlim_j\varphi$). This, together with the argument of c), shows
that for a given category %%
$J$ arising from a finite or discrete ordered set (or more
generally, one that is $C$-admissible (\hyperref[I.prop8.3.1]{8.3.1})), projective
limits (\resp inductive limits) of type $J$ are
representable in $\Ind(C)$ if (and only if) for every
functor $\varphi:J\rightarrow C$, the projective (\resp
inductive) limit of $\varphi$ \emph{computed in} $\Ind(C)$ is
representable. In the nonparenthetical case, this also means, by virtue
of (a), that every projective limit of type $J$ of representable
presheaves on $C$ is ind-representable.
\end{enonce*}

\begin{enonce*}{Proposition 8.9.5}\etiquette[8.9.5]{I.prop8.9.5}
Let $C$ be a $𝒰$-category.
\begin{enumerate}
\renewcommand{\theenumi}{\alph{enumi}}
\renewcommand{\labelenumi}{\theenumi)}
\item The canonical functor
$$
c:C\longrightarrow \Ind(C)
$$
is right exact (hence exact, in view of $\hyperref[I.prop8.9.1]{8.9.1}$ {\rm a)}).

\item If finite inductive limits (\resp finite sums) are
representable in $C$, then small inductive limits (\resp
small sums) are representable in $\Ind(C)$. Let $J$ be a
preordered\pageoriginale set that is finite or discrete; if inductive limits
of type $J$ are representable in $C$, the same is
true in $\Ind(C)$.
\end{enumerate}
\end{enonce*}

\noindent \emph{Proof.}
\begin{enumerate}
\renewcommand{\theenumi}{\alph{enumi}}
\renewcommand{\labelenumi}{\theenumi)}
\item Suppose that in $C$ one has $X\simeq \varinjlim_J X_j$, where
$J$ is a finite category and $X$ and the $X_j$ belong to $C$. Then,
for every ind-object $𝒴=(Y_i)_{i\in I}$ of $C$, one has:
\begin{gather*}
\Hom(X,𝒴)\simeq \varinjlim_i \Hom(X,Y_i)\simeq \varinjlim_i
\varprojlim_j \Hom(X_j,Y_i)\\
\mathrel{\mathop\simeq^{\hyperref[I.prop2.8]{2.8}}}\varprojlim_j \varinjlim_i \Hom(X_j,
Y_i) \simeq \varprojlim_j \Hom (X_j,𝒴){\quoi},
\end{gather*}
and the desired conclusion follows by comparing the outer
terms.

\item We already noted in \hyperref[I.rem8.8.4]{8.8.4} that the
  assertions made
follow from a) and \hyperref[I.cor8.9.2]{8.9.2}, at least for finite inductive
limits. To prove the stated conclusions in the infinite
cases, we are reduced to the case of sums, which can be interpreted
as a filtered inductive limit of finite sums, and we therefore conclude
by \hyperref[I.prop8.5.1]{8.5.1}.
\end{enumerate}

% END INCLUDED SGA4_English_current_buildable_source_20260906/source/01/drafts/SGA4_I_LINES_4626_4853_English.texfrag


% BEGIN INCLUDED SGA4_English_current_buildable_source_20260906/source/01/drafts/SGA4_I_LINES_4854_5078_English.texfrag
\begin{enonce*}[remark]{Remark 8.9.6}\etiquette[8.9.6]{I.rem8.9.6}
We have already observed that the functor $c$ does not in general commute with
filtered inductive limits, and hence does not commute with infinite sums
either, in contrast to what happens for $L:\Ind(C)\rightarrow
\widehat{C}$. On the other hand, apart from its commutation with
\emph{filtered} inductive limits (\hyperref[I.prop8.5.1]{8.5.1}), $L$ has
practically no commutation properties with any other type of inductive
limit (initial object, sum of two objects, cokernels of double arrows).
Thus, if $C$ admits an initial object $∅_C$, which is therefore an initial
object of $\Ind(C)$ by a), $∅_C$ is never an initial object of
$\widehat{C}$ (\ie identical to the constant presheaf with value $∅$),
since $\Hom(∅_C,∅_C)\neq ∅$!
\end{enonce*}

\begin{enonce*}{Proposition 8.9.7}\etiquette[8.9.7]{I.prop8.9.7}
Let\pageoriginale
$$
f:C\longrightarrow C'
$$
be a functor between $𝒰$-categories, and let $J$ be a finite rigid
category $\hyperref[I.prop8.8.5]{(8.8.5)}$. If inductive (\resp projective) limits of type $J$
are representable in $C$ and $f$ commutes with them, then
$\Ind(f):\Ind(C)\rightarrow \Ind(C')$ also commutes with limits of this
type.
\end{enonce*}

This follows at once from \hyperref[I.prop8.8.5]{8.8.5} and from the
calculation \hyperref[I.cor8.9.2]{8.9.2} of finite limits in a category of
ind-objects, for a functor represented in indexed form.

\begin{enonce*}{Corollary 8.9.8}\etiquette[8.9.8]{I.cor8.9.8}
If finite inductive (\resp projective) limits are representable in $C$,
and if $f$ is right exact (\resp left exact), then the same is true of
$\Ind(f)$; in the unparenthesized case, $\Ind(f)$ even commutes with all small
inductive limits.
\end{enonce*}

The first assertion follows from \hyperref[I.prop8.9.7]{8.9.7}. The second
follows from the first, taking into account that $\Ind(f)$ commutes with
filtered inductive limits (\hyperref[I.prop8.6.3]{8.6.3}).

\begin{enonce*}[remark]{Exercise 8.9.9}\etiquette[8.9.9]{I.exer8.9.9}
Let $C$ be a $𝒰$-category.
\begin{enumerate}
\renewcommand{\theenumi}{\alph{enumi}}
\renewcommand{\labelenumi}{\theenumi)}
\item Suppose that finite sums are representable in $C$. Show that if
\emph{finite} sums in $C$ are disjoint (\resp universal) (\cf II \hyperref[II.def4.5]{4.5}),
then in $\Ind(C)$ \emph{small} sums are disjoint (\resp universal).

\item Suppose that every morphism in $C$ factors as an epimorphism
followed by a monomorphism (\resp as an effective epimorphism followed\pageoriginale by
a monomorphism, \resp as an epimorphism followed by an effective
monomorphism, \resp as an effective epimorphism followed by an effective
monomorphism). Show that $\Ind(C)$ satisfies the same property. In the
last case, it follows that every epimorphism of $\Ind(C)$ is effective,
every monomorphism of $\Ind(C)$ is effective, and every bimorphism of
$\Ind(C)$ is an isomorphism; show that in this case every epimorphism
(\resp monomorphism) of $\Ind(C)$ can be represented by an inductive
system of epimorphisms (\resp monomorphisms) of $C$. If, in addition,
every epimorphism (\resp every monomorphism) in $C$ is assumed to be
\emph{universal}, the same property holds in $\Ind(C)$.

\item Suppose that $C$ is additive (\resp abelian); then so is $\Ind(C)$.

\item Let $X=\textrm{« }\varinjlim_I\textrm{ »\,}X_i$ be an ind-object of $C$,
and let $R\rightrightarrows X$ be an equivalence relation on $X$. For
every $i\in \Ob I$, let $R_i\rightrightarrows X_i$ be the equivalence
relation on $X_i$ (regarded as an object of $\Ind(C)$) induced by $R$
via $X_i\rightarrow X$. (We assume that the fiber product
$R_i=R\times_{X_XX}X_i \times X_i$ in $(\Ind(C))$ is representable by an
object of $\Ind(C)$, as is the case if finite projective limits are
representable in $C$.) Show that, for $R$ to be effective, it suffices
that the $R_i$ be effective.

\item Let $X$ be an object of $C$, and let $R$ be an equivalence relation
on $c(X)$, \ie on $X$ regarded as an object of $\Ind(C)$. Suppose that
fiber products are representable in $C$. For $R$ to be effective, it is
necessary that $R$ be of the form
$\textrm{« }\varinjlim_i\textrm{ »\,} R_i$, where the $R_i$ are equivalence
relations on $X$ (regarded as an object of $C$), and this condition
is\pageoriginale sufficient if equivalence relations in $C$ are assumed
to be effective.

\item Suppose that every morphism in $C$ factors as an effective
epimorphism followed by an effective monomorphism, and that finite
projective limits are representable. Let $X$ be an object of $C$, and
let $R$ be a subobject of $X\times X$, regarded as an object of
$\Ind(C)$. Write $R$ in the form $\textrm{« }\varinjlim\mbox{''
}Z_i$, where $(Z_i)_{i\in I}$ is a filtered increasing family of
subobjects of $X\times X$ \emph{in} $C$ (this is possible by b)). Show
that, for $R$ to be an equivalence relation on $X$, it is necessary and
sufficient that, for every $i\in \Ob I$, there exists an index
$j\in \Ob I$ whose corresponding subobject contains $s(Z_i)$ and
$Z_i\circ Z_i$, where $s$ is the symmetry of $X\times X$.

\item Suppose that in $C$ finite projective limits and finite sums are
representable, that every equivalence relation there is effective, and
that every morphism there factors as an effective
epimorphism followed by an effective monomorphism. Show that equivalence
relations in $\Ind(C)$ are universal if and only if $C$ satisfies the
following condition:
\begin{itemize}
\item[%%
ST)] For every object $X$ of $C$ and every subobject $Z$ of $X\times
X$
in $C$, if the sequence of subobjects $Z_n$ $(n\geq 0)$ of $X\times X$
in $C$ is defined recursively by $Z_\circ=Z$,
$Z_n=\Sup(Z_n,s(Z_n),Z_n\circ Z_n)$ (where $\Sup$ is taken in the set of
subobjects of $X\times X$, which exists by the hypotheses made on C),
then the sequence $(Z_n)_{n\geq 0}$ is \emph{stationary.}
\end{itemize}

\setcounter{enumi}{10}
\item Show that equivalence relations in $\Ind((\Ens))$ are not
necessarily effective.

$^\ast$ NB. For criteria ensuring that $\Ind(C)$ is a topos, \cf VI
8.9.9.${}_\ast$
\end{enumerate}
\end{enonce*}

\begin{enonce*}[remark]{Exercise 8.9.10}\etiquette[8.9.10]{I.exer8.9.10}
Let\pageoriginale $C$ be a $𝒰$-category. Put
$\Pro(C)=(\Ind(C^\circ))^\circ$ (\cf \hyperref[I.sec8.11]{8.11}).
\begin{enumerate}
\renewcommand{\theenumi}{\alph{enumi}}
\renewcommand{\labelenumi}{\theenumi)}
\item Suppose that finite sums (\resp small sums) are representable in
$C$. Show that if they are disjoint (\cf II \hyperref[II.def4.5]{4.5}), then $\Pro(C)$
satisfies the same %%
condition.

\item Let $J$ be a small set such that sums indexed by $J$ are
representable in $C$, and hence also in $\Pro(C)$. Let
$(X(\alpha))_{\alpha\in J}$ be a family of objects of $\Pro(C)$, with
$X(\alpha)=\textrm{« }\varprojlim_{I_\alpha}\textrm{ »\,}X_{i_\alpha}$.
Show that, for the sum of the $X(\alpha)$ in $\Ind(C)$ to be universal,
it suffices that the same hold for each of the families
$\varprojlim_J X_{i_\alpha}$, where $(i_\alpha)_{a\in I}\in
\prod_{\alpha\in J}I_\alpha$. Deduce that, if in $C$ sums of type $J$
are universal, then for the same to hold in $\Pro(C)$ it is necessary
and sufficient that, for every family $(X(\alpha))_{\alpha\in J}$ as
above, with $I_\alpha=I$ for every $\alpha\in J$, the canonical
homomorphism in $\Pro(C)$
$$
\textrm{« }\varprojlim_I\textrm{ »\,}\amalg_\alpha X(\alpha)_i\longleftarrow
\textrm{« }\varprojlim_{I^J}\textrm{ »\,} \amalg_{\alpha\in
J}X(\alpha)_{i_\alpha}
$$
be an isomorphism. Deduce that in $\Pro((\Ens))$, sums of type $J$ are
universal if and only if $J$ is finite.
\end{enumerate}
\end{enonce*}

\begin{enonce*}[remark]{Exercise 8.9.11}\etiquette[8.9.11]{I.exer8.9.11}
Let $C$ be a $𝒰$-category.
\begin{enumerate}
\renewcommand{\theenumi}{\alph{enumi}}
\renewcommand{\labelenumi}{\theenumi)}
\item Let $(X_i\rightarrow X)_{i\in I}$ be a family of morphisms in
$C$. Show that, for it to be epimorphic in $\Ind(C)$, it suffices that
there exists a \emph{finite} subfamily that is epimorphic in $C$. Prove
that this condition is also necessary when it is assumed that in $C$
every finite family of morphisms with target $X_i$ factors as an epimorphic
family,\pageoriginale followed by an effective monomorphism \hyperref[I.sec10.5]{(10.5)}. (For
this last assertion, let $P$ be the set of finite subsets of $I$, ordered
by inclusion, and let $X'=(X'_J)_{J\in P}$ be the ind-object formed by the
images of the finite subfamilies of the given family; finally, let
$X''=X\amalg_{X'}X=\textrm{« }\lim\textrm{ »\,}X\amalg_{X'_J}X$. Show
that the two canonical morphisms $X\rightrightarrows X''$ are distinct,
but coincide on the $X'_J$.)

\item Suppose that every finite family of morphisms with the same target
in $C$ factors as an epimorphic family followed by an effective
monomorphism. Prove that $\Ind(C)$ admits a small subcategory generating
by epimorphisms (\hyperref[I.def7.1]{7.1}) if and only if $C$ admits a
small subcategory $C'$ such that every object of $C$ is the target of a
\emph{finite} epimorphic family with source in $C'$. If $C$ is assumed
to admit a small generating subcategory (\hyperref[I.def7.1]{7.1}), then
$\Ind(C)$ admits a small subcategory generating by epimorphisms if and
only if $C$ is equivalent to a small category. (For this last statement,
use \hyperref[I.var7.5.2]{7.5.2}).

\item $\Ind((\Ens))$ does not admit a small generating subcategory.
\end{enumerate}
\end{enonce*}

\begin{enonce*}[remark]{Exercise 8.9.12}\etiquette[8.9.12]{I.8.9.12}
Let $C$ be a $𝒰$-category, and consider
$\Pro(C)=\Ind(C^\circ)^\circ$.
\begin{enumerate}
\renewcommand{\theenumi}{\alph{enumi}}
\renewcommand{\labelenumi}{\theenumi)}
\item Show that if $\Pro(C)$ admits a small subcategory $P'$ generating
by epimorphisms \hyperref[I.def7.1]{(7.1)}, then there exists a small subcategory $C'$ of $C$
that is generating by epimorphisms in $\Pro(C)$. (Take the category of
components of the objects of $P'$.)

\item Show\pageoriginale that the category $\Pro((\Ens))$ has no small
subcategory generating by epimorphisms. (If $C'$ is as in a), choose a
set $X$ whose cardinality strictly exceeds the cardinalities of the
objects of $C'$, and consider the pro-set $𝒳$ formed by the complements
in $X$ of subsets of cardinality $<\card(X)$. Show that, for every
nonempty set $Y$ of cardinality $<\card(X)$, one has $\Hom(Y,𝒳)=∅$, but
that $𝒳$ is not isomorphic to the constant pro-set $∅$.)
\end{enumerate}
\end{enonce*}

% END INCLUDED SGA4_English_current_buildable_source_20260906/source/01/drafts/SGA4_I_LINES_4854_5078_English.texfrag


% BEGIN INCLUDED SGA4_English_current_buildable_source_20260906/source/01/drafts/SGA4_I_LINES_5079_5229_English.texfrag
\subsection{Dual notions: pro-objects, pro-representable functors}\etiquette[8.10]{I.sec8.10}

Let $C$ be a $𝒰$-category. A \emph{pro-object} of $C$ is a functor
\hypertarget{sga4.exp.I.equation.8.10.1}{}$$
\displaylines{\rlap{$(8.10.1)$}\hfill \varphi:I^\circ\longrightarrow
C{\quoi},\hfill}
$$
where $I$ is an essentially small filtered category (called the index
category), and, as usual, $I^\circ$ denotes the opposite category. Let
$𝒱$ be a universe such that $𝒱\supset 𝒰$. One should note that the
pro-objects of $C$ indexed by $I\in 𝒱$ are in one-to-one
correspondence with the ind-objects of $C^\circ$ indexed by $I\in
𝒱$: to every such ind-object $\psi:I\rightarrow C^\circ$ one
associates the pro-object $\varphi=\psi^\circ:I^\circ\rightarrow C$.
Following this correspondence, one defines ``by transport of structure
and reversal of arrows'' the notion of a morphism between pro-objects
of $C$ from the analogous notion \hyperlink{sga4.exp.I.equation.8.2.4.2}{(8.2.4.2)}, \hyperlink{sga4.exp.I.equation.8.2.5.1}{(8.2.5.1)} for ind-objects,
and consequently defines the category of pro-objects of $C$ indexed by
$I\in 𝒱$, together with a canonical isomorphism
\hypertarget{sga4.exp.I.equation.8.10.2}{}$$
\displaylines{\rlap{$(8.10.2)$}\hfill
\Pro_{𝒱}(C,𝒰)\simeq \Ind_{𝒱}(C^\circ,
𝒰)^\circ{\quoi};\hfill}
$$
for\pageoriginale $𝒱=𝒰$ one simply speaks of the \emph{category of
pro-objects of} $C$, denoted $\Pro_{𝒰}(C)$ or simply $\Pro(C)$,
which is thus defined by
\hypertarget{sga4.exp.I.equation.8.10.3}{}$$
\displaylines{\rlap{$(8.10.3)$}\hfill
\Pro(C)=\Pro_{𝒰}(C,𝒰)\simeq
\Ind(C^\circ)^\circ{\quoi}.\hfill}
$$
If two pro-objects are given in \emph{indexed form}
\hypertarget{sga4.exp.I.equation.8.10.4}{}\begin{equation}
𝒳=(X_i)_{i\in I}\quav 𝒴=(Y_j)_{j\in
J}{\quoi},\tag{8.10.4}
\end{equation}
formula \hyperlink{sga4.exp.I.equation.8.2.5.1}{(8.2.5.1)} takes here the form
\hypertarget{sga4.exp.I.equation.8.10.5}{}\begin{equation}
\Hom(𝒳,𝒴)\simeq \varprojlim_j \varinjlim_i
\Hom(X_i,Y_j){\quoi}.\tag{8.10.5}
\end{equation}
The functor \hyperlink{sga4.exp.I.equation.8.4.1}{(8.4.1)}, applied to $C^\circ$, gives, upon passing to
opposite categories, a canonical functor (denoted by the same letter
$c$ when there is no risk of confusion) that identifies $C$ with a
full subcategory of $\Pro(C)$:
\hypertarget{sga4.exp.I.equation.8.10.6}{}\begin{equation}
c:C\longrightarrow\Pro(C){\quoi};\tag{8.10.6}
\end{equation}
it is in order to have such a functor, rather than
$C\rightarrow\Pro(C)^\circ$, that we ``reversed the arrows'' in the
definition of morphisms of pro-objects starting from the analogous
definition for ind-objects. Pro-objects in the essential image of
\hyperlink{sga4.exp.I.equation.8.10.6}{(8.10.6)} are again called \emph{essentially constant pro-objects.}

Put
\hypertarget{sga4.exp.I.equation.8.10.7}{}\begin{equation}
\check{C}=(C^\circ)^{\widehat{\phantom{,}}}=\Hhom(C,(\Ens)){\quoi},\tag{8.10.7}
\end{equation}
then the functor \hyperlink{sga4.exp.I.equation.8.2.4.7}{(8.2.4.7)} may be regarded as a canonical fully faithful
functor
\hypertarget{sga4.exp.I.equation.8.10.8}{}\begin{equation}
L:\Pro_{𝒱}(C,𝒰)\longrightarrow
\check{C}^\circ{\quoi},\tag{8.10.8}
\end{equation}\pageoriginale
and in particular one obtains
\hypertarget{sga4.exp.I.equation.8.10.9}{}\begin{equation}
L:\Pro(C)\hookrightarrow \check{C}^\circ{\quoi}.\tag{8.10.9}
\end{equation}

We leave it to the reader to translate, as the need arises, the results
of the preceding sections and of those that follow from the language of
ind-objects into that of pro-objects. Depending on the mathematical
context, one language or the other is more useful, not to mention the
cases in which both notions arise simultaneously, for example when one
must consider complex categories such as $\Pro(\Ind(C))$ or
$\Ind(\Pro(C))$. We confine ourselves to giving a few additional
indications, in order to fix the terminology and notation and to make
the ``yoga'' precise.


\setcounter{subsubsection}{9}
\subsubsection{}\etiquette[8.10.10]{I.sec8.10.10}
A \emph{covariant functor $F:C\rightarrow (\Ens)$} is called a
\emph{pro-representable functor} if it belongs to the essential image
of \hyperlink{sga4.exp.I.equation.8.10.9}{(8.10.9)} (or of \hyperlink{sga4.exp.I.equation.8.10.8}{8.10.8}, which amounts to the same thing); the full
subcategory of $\check{C}$ formed by these functors is therefore
equivalent to $\Pro(C)^\circ$. In general one will prefer to work in
the opposite category, the essential image \emph{as a subcategory of}
\hyperlink{sga4.exp.I.equation.8.10.9}{(8.10.9)}, which is therefore equivalent to $\Pro(C)$ itself. In
accordance with this usage, it is often preferable to regard covariant
functors
$$
F:C\longrightarrow (\Ens)
$$
as objects of $\Hhom(C,(\Ens))^\circ=\check{C}^\circ$, and consequently
to write the category of morphisms
$$
X\longrightarrow F\mbox{ in } \check{C}^\circ{\quoi},
$$
\ie\pageoriginale of pairs $(X,u)$ consisting of an object $X$ of $C$
and an element $u\in F(X)$, as
\hypertarget{sga4.exp.I.equation.8.10.11}{}\begin{equation}
{}_{\displaystyle F}\backslash C{\quoi}  %%
\ .\tag{8.10.11}
\end{equation}
With this convention, one thus obtains a \emph{covariant} functor
(source functor)
\hypertarget{sga4.exp.I.equation.8.10.12}{}\begin{equation}
{}_{\displaystyle F}\backslash C\rightarrow C{\quoi},\tag{8.10.12}
\end{equation}
and \hyperref[I.prop3.4]{3.4} is rewritten in the form
\hypertarget{sga4.exp.I.equation.8.10.13}{}\begin{equation}
F\xrightarrow{\sim}%%
\textrm{``}\varprojlim_{({}_{\displaystyle F}\backslash C)^\circ}\textrm{''\,}
X\tag{8.10.13}
\end{equation}

\setcounter{subsubsection}{13}
\subsubsection{}\etiquette[8.10.14]{I.sec8.10.14}
Criterion \hyperref[I.them8.3.3]{8.3.3}, applied to $C^\circ$, becomes
a pro-representability criterion for $F$: it is necessary and
sufficient that $({}_F\backslash C)^\circ$ be filtered and essentially
small, the latter condition being superfluous if $C$ is equivalent to a
small category; if, moreover, finite projective limits are representable
in $C$, this is equivalent to saying that $F$ commutes with them, \ie is
\emph{left exact.}

\subsubsection{}\etiquette[8.10.15]{I.sec8.10.15}
In $\Pro(C)$ small filtered projective limits are representable and the
functor \hyperlink{sga4.exp.I.equation.8.10.9}{(8.10.9)} commutes with them; in other words, this functor carries
projective limits in $\Pro(C)$ to \emph{inductive} limits in
$\check{C}=\Hhom(C,(\Ens))$, just like its composite
$C\rightarrow\check{C}^\circ$ with \hyperlink{sga4.exp.I.equation.8.10.6}{(8.10.6)}, which shares with it the
unfortunate property of being contravariant when regarded as taking
values in $\check{C}$. One should note that pro-representable\pageoriginale
functors from $C$ to $(\Ens)$ are the functors that are small filtered
\emph{inductive} limits of representable functors (and not projective
limits, as the terminology might perhaps suggest).


% END INCLUDED SGA4_English_current_buildable_source_20260906/source/01/drafts/SGA4_I_LINES_5079_5229_English.texfrag


% BEGIN INCLUDED SGA4_English_current_buildable_source_20260906/source/01/drafts/SGA4_I_LINES_5230_5389_English.texfrag
\subsection{Ind-adjoints and pro-adjoints}\etiquette[8.11]{I.sec8.11}

\subsubsection{}\etiquette[8.11.1]{I.sec8.11.1}
Let
\hypertarget{sga4.exp.I.equation.8.11.1.1}{}\begin{equation}
f:C\longrightarrow C'\tag{8.11.1.1}
\end{equation}
be a functor between $𝒰$-categories, whence a functor
$F \sisi{→}{↦} F\circ f$
\hypertarget{sga4.exp.I.equation.8.11.1.2}{}\begin{equation}
f^\ast:\widehat{C}'\longrightarrow\widehat{C}{\quoi}.\tag{8.11.1.2}
\end{equation}
We say that $f$ \emph{admits an ind-adjoint} if the preceding functor
carries ind-representable functors to ind-representable functors. Since
$f^\ast$ commutes with inductive limits, and since the full subcategory
of $\widehat{C}'$ formed by the ind-representable functors is stable
under small filtered inductive limits, it is equivalent to say that
$f$ admits an ind-adjoint, or that $f^\ast$ carries objects of $C'$
(identified with a full subcategory of $\widehat{C}'$) to
ind-representable functors on $C$, \ie
\hypertarget{sga4.exp.I.equation.8.11.1.3}{}\begin{equation}
X\longmapsto \Hom(f(X),Y')\mbox{ is ind-representable for every }
Y'\in \Ob C'{\quoi}.\tag{8.11.1.3}
\end{equation}
Another way to express the condition that $f$ admit an ind-adjoint is
to say that there exists a functor
\hypertarget{sga4.exp.I.equation.8.11.1.4}{}\begin{equation}
g:\Ind(C')\longrightarrow \Ind(C)\tag{8.11.1.4}
\end{equation}
essentially induced by $f^\ast$, \ie such that there is an isomorphism
of bifunctors
\hypertarget{sga4.exp.I.equation.8.11.1.5}{}\begin{equation}
\Hom_{\Ind(C)}(X,g(𝒴'))\simeq
\Hom_{\Ind(C')}(f(X),𝒴'){\quoi},\tag{8.11.1.5}
\end{equation}
where\pageoriginale $X\in \Ob C$ and $𝒴'\in \Ob \Ind(C')$. In fact,
it even suffices to have a functor
\hypertarget{sga4.exp.I.equation.8.11.1.6}{}\begin{equation}
g_\circ:C'\longrightarrow \Ind(C)\tag{8.11.1.6}
\end{equation}
and an isomorphism of bifunctors
\hypertarget{sga4.exp.I.equation.8.11.1.7}{}\begin{equation}
\Hom_{\Ind(C)}(X,g_\circ(Y'))\simeq \Hom_{C'}(f(X),Y')\tag{8.11.1.7}
\end{equation}
in $X\in \Ob C$, $Y'\in \Ob C'$. The functor $g$ (respectively
$g_\circ$) is evidently determined up to unique isomorphism by $f$,
and conversely $f$ can be reconstructed up to canonical isomorphism
from the knowledge of this $g$ (respectively $g_\circ$). It is clear
from \hyperlink{sga4.exp.I.equation.8.11.1.5}{(8.11.1.5)} that the functor $g$ commutes with inductive limits and
that it ``extends'' the functor $g_\circ$; this therefore determines it
up to unique isomorphism in terms of $g_\circ$
(\hyperref[I.sec8.7.2]{8.7.2}). The functor $g$, and sometimes also the
functor $g_\circ$ that it extends, is called the \emph{ind-adjoint
functor of} $f$ (there is no need here to specify: on the right, since
the other one, if it exists, will be called the pro-adjoint; cf.
\hyperref[I.sec8.11.5]{8.11.5} below).

Of course, when $f$ admits a right adjoint
$$
f^{\ad}:C'\longrightarrow C{\quoi},
$$
it admits an ind-adjoint $g$, and the latter is canonically isomorphic
to the canonical extension of $f^{\ad}$ to ind-objects:
\hypertarget{sga4.exp.I.equation.8.11.1.8}{}\begin{equation}
g\simeq\Ind(f^{\ad}){\quoi}.\tag{8.11.1.8}
\end{equation}
The notion of an ind-adjoint is thus a natural generalization of the
notion of a right adjoint.

Now consider\pageoriginale the canonical extension
\hypertarget{sga4.exp.I.equation.8.11.1.9}{}\begin{equation}
\Ind(f):\Ind(C)\longrightarrow\Ind(C'){\quoi}.\tag{8.11.1.9}
\end{equation}
We then have:

\begin{enonce*}{Proposition 8.11.2}\etiquette[8.11.2]{I.prop8.11.2}
For the functor $f:C\rightarrow C'$ to admit an ind-adjoint, it is
necessary and sufficient that the functor $\Ind(f)$ $\hyperlink{sga4.exp.I.equation.8.11.1.9}{(8.11.1.9)}$ admit
a right adjoint. The latter is canonically isomorphic to the
ind-adjoint $\hyperlink{sga4.exp.I.equation.8.11.1.4}{(8.11.1.4)}$.
\end{enonce*}

Sufficiency and the last assertion are immediate from the
ind-adjunction formula \hyperlink{sga4.exp.I.equation.8.11.1.5}{(8.11.1.5)}. For necessity, note that by passing
to the projective limit of this formula over the $X_i$, one obtains an
adjunction isomorphism for the ind-objects $𝒳=(X_i)_{i\in I}$ and
$𝒴'$:
\hypertarget{sga4.exp.I.equation.8.11.2.1}{}\begin{equation}
\Hom_{\Ind(C)}(𝒳, g(𝒴'))\simeq
\Hom_{\Ind(C')}(\Ind(f)(𝒳),𝒴'){\quoi},\tag{8.11.2.1}
\end{equation}
\cqfd

\begin{enonce*}{Corollary 8.11.3}\etiquette[8.11.3]{I.cor8.11.3}
If $f$ admits an ind-adjoint, then $\Ind(f)$ commutes with inductive
limits, and the ind-adjoint $g$ commutes with projective limits.
\end{enonce*}

\begin{enonce*}{Proposition 8.11.4}\etiquette[8.11.4]{I.prop8.11.4}
For the functor $f:C\rightarrow C'$ to admit an ind-adjoint, it is
necessary that $f$ be \sisi{left}{right} exact, and this condition is
sufficient when $C$ is equivalent to a small category.
\end{enonce*}

This follows from criterion \hyperlink{sga4.exp.I.equation.8.11.1.3}{(8.11.1.3)}, and from
\hyperref[I.prop8.3.1]{8.3.1} and \hyperref[I.them8.3.3]{8.3.3} (iv).

For another criterion in terms of the notion of an accessible functor,
cf. \hyperref[I.cor8.13.3]{8.13.3}.

\setcounter{subsubsection}{4}\etiquette[8.11.5]{I.sec8.11.5}
\subsubsection{}
Now\pageoriginale consider the functor $F\mapsto F\circ f$
\hypertarget{sga4.exp.I.equation.8.11.5.1}{}\begin{equation}
f^{\circ\ast}:\check{C}'\longrightarrow \check{C}\tag{8.11.5.1}
\end{equation}
induced by \sisi{$F$}{$f$}. We say that $f$ \emph{admits a pro-adjoint}
if the preceding functor carries pro-representable functors to
pro-representable functors, \ie if there exists a functor (called the
\emph{pro-adjoint functor of $f$})
\hypertarget{sga4.exp.I.equation.8.11.5.2}{}\begin{equation}
g:\Pro(\sisi{C}{C'})\longrightarrow \Pro(\sisi{C'}{C}){\quoi},\tag{8.11.5.2}
\end{equation}
and an isomorphism of bifunctors
\hypertarget{sga4.exp.I.equation.8.11.5.3}{}\begin{equation}
\Hom_{\Pro(C)}(g(𝒴'),𝒳)\simeq
\Hom_{\Pro(C')}(𝒴'\sisi{)}{},f(𝒳)){\quoi}.\tag{8.11.5.3}
\end{equation}
Of course, to say that $f$ admits a pro-adjoint means that $f^\circ$
admits an ind-adjoint, so that the notions and results for ind-adjoints
translate immediately into terms of pro-adjoints. We merely point out
that $f$ admits a pro-adjoint if and only if
$\Pro(f):\Pro(C)\rightarrow\Pro(C')$ admits a \emph{\sisi{right}{left}}
adjoint, and that in this case the preceding functor $g$ is such a
\sisi{right}{left} adjoint of $\Pro(\sisi{C}{f})$; and that for this it
is necessary that $f$ be \emph{left} exact, this condition also being
sufficient when $C$ is equivalent to a small category. In this case,
$f$ is therefore exact if and only if it admits both an ind-adjoint and
a pro-adjoint.

\begin{enonce*}[remark]{Example 8.11.6}\etiquette[8.11.6]{I.exem8.11.6}
Consider the case of a functor
$$
f:C\longrightarrow (\Ens){\quoi}.
$$
If this functor admits a pro-adjoint, it is pro-representable, and the
converse is true if and only if the full subcategory of $\check{C}$
formed by the pro-representable functors is stable under small products;
this\pageoriginale is the case in particular if $C$ is equivalent to a
small category (\hyperref[I.sec8.10.14]{8.10.14}) or if small products
are representable in $C$ (\hyperref[I.prop8.9.5]{8.9.5} b), applied to
$C^\circ$). Up to minor qualifications, one may thus say that for a
functor $f:C\rightarrow C'$, the existence of a pro-adjoint is the
\emph{natural generalization of the notion of pro-representability of
$f$}, which is defined when $C'=(\Ens)$.
\end{enonce*}


% END INCLUDED SGA4_English_current_buildable_source_20260906/source/01/drafts/SGA4_I_LINES_5230_5389_English.texfrag


% BEGIN INCLUDED SGA4_English_current_buildable_source_20260906/source/01/drafts/SGA4_I_LINES_5390_5467_English.texfrag
\subsection{Strict ind-objects and pro-objects. Application to a
representability criterion}\etiquette[8.12]{I.sec8.12}

\subsubsection{}\etiquette[8.12.1]{I.sec8.12.1}
Let
$$
𝒳=(X_i)_{i\in I}
$$
be an ind-object of the $𝒰$-category $C$, and let $F$ be the presheaf
that it ind-represents. It is then immediately clear that it is
equivalent to say that the canonical morphisms
$$
X_i\longrightarrow F=\textrm{``}\varinjlim_i\textrm{''\,}X_i
$$
are monomorphisms in $\Ind(C)$ (or, equivalently, in $\widehat{C}$,
\ie monomorphisms of functors ``argument by argument''), or to say that
for every arrow $i\rightarrow j$ of $I$, the corresponding transition
arrow
$$
X_i\longrightarrow X_j
$$
is a monomorphism. When these conditions are satisfied, and if moreover
$I$ is an ordered category, we say that $𝒳$ is a \emph{strict
ind-object}. One\pageoriginale should note that this condition is not
invariant under isomorphism of ind-objects; an ind-object will be called
\emph{essentially strict} if it is isomorphic to a strict ind-object. A
presheaf will be called \emph{strictly ind-representable} if it is
ind-representable by a strict ind-object; thus
$F=\textrm{``}\varinjlim_i\textrm{''\,}X_i$ is strictly
ind-representable if and only if $𝒳=(X_i)_{i\in I}$ is essentially
strict.

\paragraph{}\etiquette[8.12.1.1]{I.8.12.1.1}
Let $F$ be a presheaf on $C$, and consider the full subcategory of
$C_{/F}$ formed by the arrows $X\rightarrow F$ with source in $C$ that
are monomorphisms. We shall call it the \emph{category of representable
subfunctors of $F$}; it is the category associated with the ordered set
of representable subfunctors of $F$, ordered by the order induced from
that of the set of subobjects of $F$. The following then follows
immediately from the definitions:

\begin{enonce*}{Proposition 8.12.2}\etiquette[8.12.2]{I.prop8.12.2}
For the presheaf $F$ on $C$ to be strictly ind-representable, it is
necessary and sufficient that the (ordered) category $I$ of
representable subfunctors of $F$ be filtered and essentially small, and
that one have
\hypertarget{sga4.exp.I.equation.8.12.2.1}{}\begin{equation*}
\varinjlim_I X_i \xrightarrow{\sim}F{\quoi}.\tag{8.12.2.1}
\end{equation*}
When, for every object $X$ of $C$, the set of subobjects of $X$ in $C$
is small (for example, if $C$ admits a small generating subcategory
$\hyperref[I.prop7.4]{(7.4)}$), this implies that the category of representable subfunctors is
itself small.
\end{enonce*}

\setcounter{subsubsection}{2} \setcounter{paragraph}{1}
\paragraph{}\etiquette[8.12.2.2]{I.sec8.12.2.2}
If\pageoriginale $F$ is strictly ind-representable, there is thus a
preferred way of ind-representing it by an ind-object, and the latter is
even a strict ind-object: one takes the representation \hyperlink{sga4.exp.I.equation.8.12.2.1}{(8.12.2.1)}. A
strict ind-object $𝒴$ is called \emph{saturated} if it is isomorphic to
an ind-object of the form $(X_i)$ occurring in \hyperlink{sga4.exp.I.equation.8.12.2.1}{(8.12.2.1)}, for a
suitable $F$; thus, up to unique isomorphism, there exists a single
saturated strict ind-object isomorphic to the given strict ind-object,
namely the one considered in \hyperref[I.prop8.12.2]{8.12.2}, taking
$F=\varinjlim 𝒴$ (the limit in $\widehat{C}$).

\paragraph{}\etiquette[8.12.2.3]{I.sec8.12.2.3}
Suppose that it is already known that one can find a small cofinal subset
in the set $\Ob C_{/F}$, as is the case in particular if $F$ is
ind-representable; it then follows that the same condition is satisfied
in the full subcategory $I$ considered in
\hyperref[I.prop8.12.2]{8.12.2}, so in criterion
\hyperref[I.prop8.12.2]{8.12.2} one may omit the condition that $I$ be
essentially small.


% END INCLUDED SGA4_English_current_buildable_source_20260906/source/01/drafts/SGA4_I_LINES_5390_5467_English.texfrag


% BEGIN INCLUDED SGA4_English_current_buildable_source_20260906/source/01/drafts/SGA4_I_LINES_5468_5594_English.texfrag
\subsubsection{}\etiquette[8.12.3]{I.sec8.12.3}
Let $F$ be a presheaf on $C$, and consider an object
$$
u:X\longrightarrow F
$$
of $C_{/F}$. We say that $u$ (or the pair $(X,u)$) is \emph{minimal}
if, for every factorization of $u$ as
\hypertarget{sga4.exp.I.equation.8.12.3.1}{}\begin{equation}
X\xrightarrow{p}X'\xrightarrow{u'}F{\quoi},\tag{8.12.3.1}
\end{equation}
with $p$ a strict epimorphism (\hyperref[I.sec10.2]{10.2}), $p$ is an
isomorphism. Regarding $u$ as an element of $F(X)$ \hyperref[I.prop1.4]{(1.4)}, to say that
$u$ is minimal thus means that\pageoriginale every strict epimorphism
$p:X\rightarrow X'$ such that $u\in
\Iim(F(p):F(X')\rightarrow F(X))$ is an isomorphism. This notion is
clarified by part a) of the following lemma:

\begin{enonce*}{Lemma 8.12.4}\etiquette[8.12.4]{I.lem8.12.4}
Let $F$ be a presheaf on $C$.
\begin{enumerate}
\renewcommand{\theenumi}{\alph{enumi}}
\renewcommand{\labelenumi}{\theenumi)}
\item Suppose that $F$ carries cokernels to kernels, and consider a
morphism
$$
u:X\longrightarrow F{\quoi},
$$
with $X\in \Ob C$. For $u$ to be a monomorphism, it is sufficient,
when cokernels of pairs of arrows are representable in $C$, that $u$ be
minimal; this condition is also necessary if fiber products are
representable in $C$.

\item Suppose that finite $\varinjlim$ are representable in $C$. For
the subcategory of representable subfunctors of $F$
(\hyperref[I.8.12.1.1]{8.12.1.1}) to be filtered (not necessarily
small) and to have $F$ as its inductive limit, it is sufficient that
$F$ be left exact and that every morphism $u:X\rightarrow F$, with
$X\in \Ob C$, factor as
$$
X\rightarrow X'\xrightarrow{u'}F{\quoi},
$$
with $u'$ minimal; this condition is also necessary if fiber products
are representable in $C$.

\item Suppose that $C$ admits a small generating subcategory
(\hyperref[I.def7.1]{7.1}), and that $I$, the category of representable
subfunctors of $F$, is filtered; then $I$ is small.
\end{enumerate}
\end{enonce*}

\noindent \emph{Proof}.\pageoriginale
\begin{enumerate}
\renewcommand{\theenumi}{\alph{enumi}}
\renewcommand{\labelenumi}{\theenumi)}
\item Suppose $u$ is minimal; let us prove that $u$ is a monomorphism,
\ie that for every pair of arrows $v,v':Y\rightrightarrows X$ such that
$uv=uv'$, one has $v=v'$. Indeed, if $X'=\Coker(v,v')$, then $u$
factors as $X\xrightarrow{p}X'\xrightarrow{u'}F$ (since $F$ carries
cokernels to kernels), and since $p$ is a strict epimorphism by
construction, it follows that $p$ is an isomorphism, \ie $v=v'$.
Suppose that $u$ is a monomorphism; let us prove that it is minimal.
Consider a factorization \hyperlink{sga4.exp.I.equation.8.12.3.1}{(8.12.3.1)}; since $p$ is a strict epimorphism,
it is the cokernel of the canonical pair of arrows
$v,v':\fprod{X}{X}{X'}\rightrightarrows X$, and since
$(u'p)v=(u'p)v'$ and $u'p=u$ is a monomorphism, one has $v=v'$, hence
$p$ is an isomorphism.

\item Sufficiency: Since $F$ is left exact, $C_{/F}$ is filtered. Since
we know that $F=\varinjlim_{C_{/F}}X$, by
\hyperref[I.prop8.1.3]{8.1.3} c) it remains to prove that the full
subcategory $I$ of $C_{/F}$ consisting of the subobjects of $F$ is
cofinal in $C_{/F}$; by the ``sufficient'' part of a), this is precisely
what is ensured by the hypothesis that every object of $C_{/F}$ is
dominated by a ``minimal'' object. Necessity: Since every filtered
inductive limit of right exact functors is again right exact, the first
condition is trivially necessary. The second then follows from the
``necessary'' part of a).

\item A representable subfunctor $X\hookrightarrow F$ of $F$ is known
once one knows the full subcategory $C'_{/X}$ of $C'_{/F}$, where $C'$
is a fixed small generating subcategory of $C$. (Use the hypothesis
that $I$ is filtered.) Since $C'_{/F}$ is small, the set of its full
subcategories is small, whence the conclusion.
\end{enumerate}

\begin{enonce*}{Proposition 8.12.5}\etiquette[8.12.5]{I.prop8.12.5}
Let\pageoriginale $C$ be a $𝒰$-category in which finite inductive
limits are representable, and which admits a small generating
subcategory (\hyperref[I.def7.1]{7.1}). Let $F$ be a presheaf on $C$.
For $F$ to be strictly ind-representable, it is sufficient that it
satisfy the following two conditions, and these conditions are also
necessary if fiber products are representable in $C$:
\begin{enumerate}
\renewcommand{\theenumi}{\alph{enumi}}
\renewcommand{\labelenumi}{\theenumi)}
\item $F$ is left exact.

\item Every pair $(X,u)$, with $X\in \Ob C$ and $u\in F(X)$, is
dominated in $C_{/F}$ by a minimal pair $\hyperref[I.sec8.12.3]{(8.12.3)}$, \ie there exists a
minimal pair $(X',u')$ $(X'\in \Ob C, u'\in F(X'))$ and a morphism
$f:X\rightarrow X'$ such that $u=F(f)(u')$.
\end{enumerate}
\end{enonce*}

Sufficiency follows from \hyperref[I.prop8.12.2]{8.12.2} and
\hyperref[I.lem8.12.4]{8.12.4} b), c), and necessity from
\hyperref[I.prop8.3.1]{8.3.1} and \hyperref[I.lem8.12.4]{8.12.4} a).

\begin{enonce*}[remark]{Remark 8.12.6}\etiquette[8.12.6]{I.rem8.12.6}
When $F$ carries amalgamated sums in $C$ to fiber products, it is
immediately clear that for a given pair $(X,u)$ as in b), the set of
strict quotients $X'$ of $X$ such that
$u\in\Iim(F(X')\rightarrow F(X))$ is filtered decreasingly; this
implies that if the set of strict quotients of $X$ is Artinian, then
condition b) of \hyperref[I.prop8.12.5]{8.12.5} is automatically
satisfied. Thus, if the preceding condition on $X$ is satisfied for
every object $X$ of $C$ (one then also says, sometimes, that the objects
of $C^\circ$ are \emph{Artinian}), it follows from \hyperref[I.prop8.12.5]{8.12.5} that $F$ is
strictly ind-representable if and only if $F$ is left exact; in this
case, $F$ is therefore strictly ind-representable as soon as it is
ind-representable (\hyperref[I.prop8.3.1]{8.3.1}).
\end{enonce*}


% END INCLUDED SGA4_English_current_buildable_source_20260906/source/01/drafts/SGA4_I_LINES_5468_5594_English.texfrag


% BEGIN INCLUDED SGA4_English_current_buildable_source_20260906/source/01/drafts/SGA4_I_LINES_5595_5699_English.texfrag
\begin{enonce*}{Corollary 8.12.7}\etiquette[8.12.7]{I.cor8.12.7}
Let\pageoriginale $C$ be a $𝒰$-category in which small projective
limits are representable, and which admits a small cogenerating
subcategory (\hyperref[I.sec7.9]{7.9},
\hyperref[I.exem7.13]{7.13}). Then a functor
$F:C\rightarrow (\Ens)$ is representable if and only if $F$ commutes
with small projective limits.
\end{enonce*}
Necessity is clear. For sufficiency, apply
\hyperref[I.prop8.12.5]{8.12.5} to the opposite category $C^\circ$. To
prove first that $F$ is (strictly) pro-representable, it remains to show
that every pair $(X,u)$, $X\in \Ob C$ and $u\in F(X)$, is dominated by
a ``minimal'' pair. Now the family $(X_i)_{i\in I}$ of \emph{strict}
subobjects of $X$ such that
$u\in\Iim(F(X_i)\rightarrow F(X))$ is small
(\hyperref[I.cor7.5]{7.5} in dual form). Since $F$ is left exact, this
family is cofiltered (\hyperref[I.rem8.12.6]{8.12.6}), and since $F$
commutes with small projective limits, one sees that
$X'=\varprojlim_i X_i$ is a least object of this family. (Use the fact
that, being left exact, $F$ carries monomorphisms to monomorphisms.) If
$u'\in F(X')$ is the unique element whose image in $F(X)$ is $u$, one
then sees that $(X',u')$ is a minimal pair dominating $(X,u)$. This
proves that $F$ is pro-representable, and the conclusion then follows
from the

\begin{enonce*}{Lemma 8.12.8.1}\etiquette[8.12.8.1]{I.lem8.12.8.1}
Let $F:C\rightarrow(\Ens)$ be a functor, where $C$ is a $𝒰$-category
in which small projective limits are representable. For $F$ to be
representable, it is necessary and sufficient that it be
pro-representable and commute with small projective limits.
\end{enonce*}

Necessity\pageoriginale is clear; let us prove sufficiency. To say that
$F$ is representable evidently means that ${}_F\backslash C$ admits an
initial object (in fact, the initial objects of ${}_F\backslash C$ are
precisely the isomorphisms $F\rightarrow X$, \ie the representation
data for $F$). Now, since $F$ is pro-representable,
$({}_F\backslash C)^\circ$ is filtered and equivalent to a small
category, so $X=\varprojlim_{({}_F\backslash C)^\circ} X$ is
representable in $C$, and since $F$ commutes with the limit under
consideration, it follows that $X$ is an initial object of
${}_F\backslash C$, \cqfd

\begin{enonce*}{Corollary 8.12.8}\etiquette[8.12.8]{I.cor8.12.8}
Under the hypotheses on $C$ of $\hyperref[I.cor8.12.7]{8.12.7}$, let $f:C\rightarrow C'$ be a
functor from $C$ to a $𝒰$-category $C'$. For $f$ to admit a left
adjoint, it is necessary and sufficient that $f$ commute with small
projective limits.
\end{enonce*}

Indeed, this reduces immediately to \hyperref[I.cor8.12.7]{8.12.7}, by
applying that statement to composite functors of the form
$X\mapsto \Hom(Y',f(X))$.

\begin{enonce*}[remark]{Examples 8.12.9}\etiquette[8.12.9]{I.exer8.12.9}
As was pointed out in \hyperref[I.exem7.13]{7.13}, the hypotheses on
$C$ in \hyperref[I.cor8.12.7]{8.12.7} and
\hyperref[I.cor8.12.8]{8.12.8} are satisfied if $C$ is the category of
$𝒰$-sheaves of sets on a topological space $X\in 𝒰$ (or, more
generally, on a $𝒰$-site (II \hyperref[II.def3.0.2]{(3.0.2)}). We give an instructive
example\footnote[1]{(due to H. BASS).} showing that the hypothesis of
the existence of a small cogenerating subcategory $D$ of $C$ is not
superfluous in \hyperref[I.cor8.12.7]{8.12.7}. Take for $C$ the
category of groups that are elements of $𝒰$. Let $J$ be the set of
isomorphism classes of simple groups in $C$; for every $j\in J$ choose
a simple group $G_j$ in the class $j$, and let $I$ be the filtered
ordered set of $𝒰$-small subsets of $J$, and, for $i\in I$, let
$X_i=\amalg_{j\in i}G_j$. The $X_i$ then form a\pageoriginale
projective system $(X_i)_{i\in I}$ in $C$, whose transition morphisms
are epimorphisms, and the corresponding functor
\begin{equation}
F(X)=\varinjlim_i \Hom(X_i,X)\tag{$\ast$}
\end{equation}
takes its values in $(𝒰\text{-}\Ens)$ (although the index set $I$
clearly does not have cardinality in $𝒰$); more precisely, let us show
that for every small full subcategory $C_\circ$ of $C$, there exists an
$i_\circ\in I$ such that the restriction of $F$ to $C_\circ$ is
represented by $X_{i_\circ}$, which will prove both that $F$ takes
values in $𝒰\text{-}\Ens$ and that it commutes with small projective
limits. To prove our assertion, it suffices to note that for every
$X\in \Ob C_\circ$, the cardinality of the set $J(X)$ of $j\in J$ for
which there exists a nontrivial homomorphism from $G_j$ to $X$ is
necessarily small, since such a morphism is necessarily a monomorphism
($G_j$ being simple); consequently, if $i_\circ$ is the subset of $J$
given by the union of the $J(X)$ for $X\in\ob C_\circ$, then $i_\circ$
is small, \ie $i_\circ\in I$, and it has the required property. On the
other hand, it is clear that $F$ is not representable, since
$\card I\not\in 𝒰$. From this and \hyperref[I.cor8.12.7]{8.12.7} one therefore concludes that
\emph{the category $C$ of groups in $𝒰$ admits no small full
cogenerating subcategory.} Since, on the other hand, the object $\ZZ$
of $C$ is a generator, it then follows from the proof of
\hyperref[I.cor7.12]{7.12} that there exists a two-generator group $G$
that cannot be embedded in an injective object of the category $C$ of
groups in $𝒰$. It moreover seems plausible that $C$ has no injective
objects other than the trivial groups.
\end{enonce*}


% END INCLUDED SGA4_English_current_buildable_source_20260906/source/01/drafts/SGA4_I_LINES_5595_5699_English.texfrag


% BEGIN INCLUDED SGA4_English_current_buildable_source_20260906/source/01/drafts/SGA4_I_LINES_5700_5804_English.texfrag
\subsection{Pro-representable functors and accessible
  functors}\etiquette[8.13]{I.sec8.13}\pageoriginale

\subsubsection{}\etiquette[8.13.1]{I.sec8.13.1}
In this section, we use some notions and results from the following
section, in particular \hyperref[I.prop9.11]{9.11} and
\hyperref[I.prop9.13]{9.13}, to obtain a pro-representability criterion that
we shall use (incidentally) in IV \hyperref[IV.theo9.5.4]{9.16}. In what follows, $C$ denotes a
$𝒰$-category satisfying condition $L$ of
\hyperref[I.def9.1]{9.1} b), this condition being satisfied, for example, if
small filtered inductive limits are representable in $C$.

\begin{enonce*}{Proposition 8.13.2}\etiquette[8.13.2]{I.prop8.13.2}
Let $C$ be as above, and let
$$
f:C\longrightarrow(\Ens)
$$
be a functor.
\begin{enumerate}
\renewcommand{\theenumi}{\alph{enumi}}
\renewcommand{\labelenumi}{\theenumi)}
\item Suppose that every object of $C$ is accessible (\hyperref[I.def9.3]{9.3}). If $f$
is pro-representable, then $f$ is accessible (\hyperref[I.def9.2]{9.2}) and
left exact.

\item Suppose that finite projective limits are representable in
$C$, and that $C$ admits a cardinal filtration
(\hyperref[I.def9.12]{9.12}). If $f$ is accessible and left exact, then $f$ is
pro-representable.
\end{enumerate}
\end{enonce*}

\noindent \emph{Proof.}
\begin{enumerate}
\renewcommand{\theenumi}{\alph{enumi}}
\renewcommand{\labelenumi}{\theenumi)}
\item The hypothesis on $C$ means that the covariant representable
functors from $C$ to $(\Ens)$ are accessible. The same is therefore
true of every small inductive limit of such functors
(\hyperref[I.prop9.6]{9.6} (i)), and hence also of every
pro-representable functor.

\item By \hyperref[I.them8.3.3]{8.3.3} (iii), it remains to prove that
$\Ob({}_F\backslash C)^\circ$ contains a small cofinal subcategory.
By hypothesis, there exists a cardinal $\Pi$ such that $f$ is
$\Pi$-accessible. Let $(X,u)$,\pageoriginale $u\in
F(X)$, be an object of ${}_F\backslash C$. With the notation of \hyperref[I.def9.12]{9.12} c),
we then have $X=\varinjlim_i X_i$, with $I$ filtered and large relative to
$\Pi$ and the $X_i$ in $C'=\Filt^{\Pi}(C)$, whence
$F(X)\xleftarrow{\sim}\varinjlim_i F(X_i)$. This shows that the
small subcategory $({}_F\backslash C')^\circ$ is cofinal in
$({}_F\backslash C)^\circ$, and completes the proof.
\end{enumerate}

\begin{enonce*}{Corollary 8.13.3}\etiquette[8.13.3]{I.cor8.13.3}
Let $C$ be a $𝒰$-category satisfying the following
conditions:
\begin{enumerate}
\renewcommand{\theenumi}{\alph{enumi}}
\renewcommand{\labelenumi}{\theenumi)}
\item Finite projective limits are representable in $C$.

\item Small filtered inductive limits are representable in
$C$.

\item The functor $\Ker$ on the category of double arrows of
$C$ is accessible (for example, it commutes with small filtered
$\varinjlim$).

\item Every strict epimorphism of $C$ is a universal strict epimorphism $\hyperref[I.sec10.2]{(10.2)}$.

\item There exists a small subcategory of $C$ generating by
strict epimorphisms $\hyperref[I.def7.1]{(7.1)}$.
\end{enumerate}
Under these conditions, a functor $f:C\rightarrow (\Ens)$ is
pro-representable if and only if it is left exact and
accessible (\hyperref[I.def9.2]{9.2}).
\end{enonce*}

Indeed, the conditions of \hyperref[I.prop8.13.2]{8.13.2} a) and b) on $C$ are
satisfied, by \hyperref[I.prop9.11]{9.11} and \hyperref[I.prop9.13]{9.13}
respectively.


\begin{enonce*}{Corollary 8.13.4}\etiquette[8.13.4]{I.cor8.13.4}
Let $f:C\rightarrow C'$ be a functor between $𝒰$-categories
satisfying conditions {\rm a)} through {\rm e)} of \hyperref[I.cor8.13.3]{8.13.3}. For
$f$ to admit a pro-adjoint, it is necessary and sufficient that $f$ be
left exact and accessible.
\end{enonce*}

Since\pageoriginale the representable functors
$h_{Y'}:X'\rightarrow \Hom(Y',X'):C'\rightarrow(\Ens)$ are left exact
and accessible (\hyperref[I.prop9.11]{9.11}), if $f$ has the same
properties, so do its composites with the preceding functors;
these are therefore pro-representable by \hyperref[I.cor8.13.3]{8.13.3},
\ie $f$ admits a pro-adjoint. Conversely, suppose that $f$ admits a
pro-adjoint, and let $C'_1$ be a small generating subcategory
of $C$, and $\Pi$ a cardinal such that the $Y'\in \Ob C'$ are
$\Pi$-accessible; to prove that $f$ is $\Pi$-accessible, it
therefore suffices to prove the same of its composites with the
$h_{Y'}$, $Y'\in \Ob C'$; by hypothesis, these composites are
pro-representable, hence accessible (\hyperref[I.cor8.13.3]{8.13.3}),
and hence $\Pi$-accessible provided $\Pi$ is taken sufficiently large, \cqfd

% END INCLUDED SGA4_English_current_buildable_source_20260906/source/01/drafts/SGA4_I_LINES_5700_5804_English.texfrag


% BEGIN INCLUDED SGA4_English_current_buildable_source_20260906/source/01/drafts/SGA4_I_LINES_5805_5940_English.texfrag
\section[Accessible functors and small generating subcategories]{Accessible functors, cardinal filtrations, and the construction
of small generating subcategories}\etiquette[9]{I.sec9}

The present section, which is more technical in nature than the other
sections of this exposé, will be used in this seminar only in IV 9 and
VI 4, which are not used elsewhere in the Seminar. The reader should
therefore omit the present section, at least on a first reading!

\setcounter{subsection}{-1}
\subsection{}\etiquette[9.0]{I.sec9.0}
All categories considered in this section are assumed to be
$𝒰$-categories. Except for the small index categories $I,J\ldots$ that
we shall have occasion to use, the developments that follow will apply
primarily to ``large'' categories $E,F\ldots$ that are stable under
small filtered inductive limits. Most often, however, it will suffice
that a slightly weaker condition be satisfied\pageoriginale (condition
$L$ in \hyperref[I.def9.1]{9.1} below). All cardinals considered in the
present section are assumed to belong to $𝒰$.

Following a suggestion of P. \textsc{Deligne}, we shall study, for a
functor $f:E\rightarrow F$ between large categories, a condition that
$f$ commute with certain types of filtered inductive limits. This
condition is remarkably stable, and is satisfied by the most important
functors occurring in nature. The applications we have in mind for our
seminar are 9.13.3, 9.13.4 (used in VI 4), and especially
\hyperref[I.cor9.25]{9.25}, which gives, in a nontrivial case, the
existence of a small generating family in a category of sections of a
fibered category; this result will be used in IV \hyperref[IV.theo9.5.4]{9.16}.

\begin{enonce*}{Definition 9.1}\etiquette[9.1]{I.def9.1}
\begin{enumerate}
\renewcommand{\theenumi}{\alph{enumi}}
\renewcommand{\labelenumi}{\theenumi)}
\item Let $I$ be a preordered set and $pi$ a cardinal. We say that $I$
is large relative to $pi$ if $I$ is filtered and every subset of $I$
of cardinality $\leq \pi$ has an upper bound in $I$.

\item Let $E$ be a category. If $pi$ is a cardinal, we say that $E$
satisfies condition $L_{\pi}$ if, for every small ordered set $I$ large
relative to $\pi$, $E$ is stable under inductive limits of type $I$. We
say that $E$ satisfies condition $L$ if there exists a cardinal
$\pi\in 𝒰$ such that $E$ satisfies condition $L_{\pi}$.
\end{enumerate}
\end{enonce*}

\setcounter{subsection}{1}
\subsubsection{}\etiquette[9.1.1]{I.sec9.1.1}
When, in \hyperref[I.def9.1]{9.1} a), one has $\pi\geq 2$, the second
stated condition already implies that $I$ is filtered, and if $\pi$ is
finite, for $I$ to be large relative to $\pi$ simply means that $I$ is
filtered. In what follows we shall be concerned almost exclusively with
the case in which $\pi$ is infinite. Note that if $\pi,\pi'$ are two
cardinals such that $\pi'\geq\pi$, then $I$ being large relative to
$\pi'$ evidently implies that it is large relative to $\pi$.

\subsubsection{}\etiquette[9.1.2]{I.sec9.1.2}
As\pageoriginale announced in \hyperref[I.sec9.0]{9.0}, conditions
$L_{\pi}$ and $L$ should be regarded as technical variants of the
stronger condition of stability under small filtered inductive limits.
It is clear that if $\pi,\pi'$ are cardinals such that
$\pi'\geq\pi$, then condition $L_{\pi}$ implies condition $L_{\pi'}$.

\begin{enonce*}{Definition 9.2}\etiquette[9.2]{I.def9.2}
Let $f:E\rightarrow F$ be a functor. If $\Pi$ is a cardinal, we say
that $f$ is $\Pi$-preaccessible (respectively $\Pi$-accessible) if $E$
satisfies $L_{\Pi}$ (\hyperref[I.def9.1]{9.1}) and if, for every ordered
set $I\in 𝒰$ large relative to $\Pi$, and every inductive system
$(X_i)_{i\in I}$ in $E$ of type $I$, the canonical morphism
$$
\varinjlim f(X_i)\longrightarrow f(\varinjlim X_i)
$$
is a monomorphism (respectively an isomorphism). We say that $f$ is
preaccessible (respectively accessible) (relative to the universe
$𝒰$) if there exists a cardinal $\Pi\in 𝒰$ such that $f$ is
$\Pi$-preaccessible (respectively $\Pi$-accessible).
\end{enonce*}

The category of $\Pi$-accessible (respectively accessible) functors
from $E$ to $F$ will be denoted $\Hhom(E,F)_{\Pi}$, respectively
$\Hhom(E,F)_{\acc}$.

\setcounter{subsection}{2} \setcounter{subsubsection}{0}
\subsubsection{}\etiquette[9.2.1]{I.sec9.2.1}
Evidently, a functor commuting with small filtered $\varinjlim$ (for
example, a functor admitting a right adjoint) is $\Pi$-accessible for
every cardinal $\Pi\geq 2$.

\begin{enonce*}{Definition 9.3}\etiquette[9.3]{I.def9.3}
Let $E$ be a category, $X$ an object of $E$,
$$
h^\circ_X:E\longrightarrow (𝒰\text{-}\Ens)
$$
the covariant functor it represents, and $\Pi\in 𝒰$ a cardinal. We say
that $X$ is a \emph{$\Pi$-preaccessible (respectively
$\Pi$-accessible) object of} $E$ if the functor
$h^\circ_X$\pageoriginale is $\Pi$-preaccessible (respectively
$\Pi$-accessible); we say that $X$ is \emph{preaccessible (respectively
accessible)} if there exists a cardinal $\Pi\in 𝒰$ such that $X$ is a
$\Pi$-preaccessible (respectively $\Pi$-accessible) object of $E$.
\end{enonce*}

\setcounter{subsection}{3} \setcounter{subsubsection}{0}
\subsubsection{}\etiquette[9.3.1]{I.sec9.3.1}
We denote by $E_{\Pi}$ the full subcategory of $E$ formed by the
$\Pi$-accessible objects of $E$.

When \hyperref[I.def9.3]{Definition 9.3} is applied to a category of the form
$\Hhom(C,F)$, the terminology just introduced presents \textit{a
priori} an ambiguity with the analogous terminology introduced in \hyperref[I.def9.2]{9.2},
when the objects of $\Hhom(C,F)$ are interpreted as functors; there
does not, however, seem to be any serious risk of confusion.

\begin{enonce*}{Definition 9.4}\etiquette[9.4]{I.def9.4}
Let $E$ be a category and let $\pi\in 𝒰$ be a cardinal. We say that
$E$ is a $\pi$-preaccessible (respectively $\pi$-accessible) category
if there exists in $E$ a small full subcategory $C$ that is generating
(\hyperref[I.def7.1]{7.1}) and whose objects are $\pi$-preaccessible
(respectively $\pi$-accessible) (\hyperref[I.def9.3]{9.3}). We say that
$E$ is a preaccessible (respectively accessible) category if there
exists a cardinal $\pi\in 𝒰$ such that $E$ is $\pi$-preaccessible
(respectively $\pi$-accessible).
\end{enonce*}

For important examples, cf. \hyperref[I.rem9.11.3]{9.11.3} below.

% END INCLUDED SGA4_English_current_buildable_source_20260906/source/01/drafts/SGA4_I_LINES_5805_5940_English.texfrag


% BEGIN INCLUDED SGA4_English_current_buildable_source_20260906/source/01/drafts/SGA4_I_LINES_5941_6010_English.texfrag
\begin{enonce*}{Proposition 9.5}\etiquette[9.5]{I.prop9.5}
Let $f:E\rightarrow F$ be a functor between categories such that $F$
is accessible and every object of $E$ is accessible. Then, if
$f$ admits a left adjoint, $f$ is accessible.
\end{enonce*}

Indeed,\pageoriginale by hypothesis, $F$ admits a small conservative
family of representable functors $F\rightarrow
(𝒰\text{-}\Ens)$ that are accessible, so we are immediately reduced
to showing that the composite of $f$ with each of the preceding functors
is accessible. But since $f$ admits a left adjoint,
these composites are functors $E\rightarrow
(𝒰\text{-}\Ens)$ that are \emph{representable,} and hence accessible
by hypothesis on $E$.

\begin{enonce*}{Proposition 9.6}\etiquette[9.6]{I.prop9.6}
Let $E$ and $F$ be two categories, let $\pi\in 𝒰$ be a cardinal, and
consider the full subcategory $\Hhom(E,F)_{\Pi}$ of
$\Hhom(E,F)$ formed by the $\Pi$-accessible functors $\hyperref[I.def9.2]{(9.2)}$ from $E$
to $F$.
\begin{enumerate}
\renewcommand{\theenumi}{\roman{enumi}}
\renewcommand{\labelenumi}{(\theenumi)}
\item This subcategory is stable under every type of $\varinjlim$
that is representable in~$F$.

\item Suppose that $F$ is $\pi$-accessible (\hyperref[I.def9.4]{9.4}), and let $J$ be a
category such that $\card  F\ell \  J\leq \pi$ and that the
$\varprojlim$ of type $J$ are representable in $F$, so the
$\varprojlim$ of type $J$ are representable in $\Hhom(E,F)$.
Then the subcategory $\Hhom(E,F)_{\Pi}$ is stable under the
$\varprojlim$ of type $J$.
\end{enumerate}
\end{enonce*}

\begin{enonce*}{Corollary 9.7}\etiquette[9.7]{I.cor9.7}
Let $E$ and $F$ be two categories, with $F$ accessible (\hyperref[I.def9.4]{9.4}).
Then the full subcategory $\Hhom(E,F)_{\acc}$ of $\Hhom(E,F)$
formed by the accessible functors is stable under every type of inductive
or projective limit, indexed by a small category
$J$, that is representable in $F$ (and hence in $\Hhom(E,F)$).
\end{enonce*}

Proof of \hyperref[I.prop9.6]{9.6}. Assertion (i) follows trivially from
the commutation of the functor $\varinjlim$ with arbitrary inductive
limits. For (ii), let $(f_j)_{j\in J}$ be a projective system
of $\Pi$-accessible functors $E\rightarrow F$, and let $f=\varprojlim
f_j$\pageoriginale be its projective limit, computed ``argument
by argument.'' We prove that $f$ is $\Pi$-accessible, \ie that
for every ordered set $I\in 𝒰$ large relative to $\Pi$, and
every inductive system $(X_i)_{i\in I}$ in $E$, the
canonical morphism
$$
\varinjlim_i ((\varprojlim_j f_j)(X_i))\longrightarrow
(\varprojlim_j f_j) (\varinjlim_i X_i)
$$
is an isomorphism. Now the ``argument-by-argument'' computation of the
functor $\varprojlim_j f_j$ allows us to identify the preceding canonical
morphism with the canonical morphism
$$
\varinjlim_i \varprojlim_j f_j (X_i)\longrightarrow \varprojlim_j
\varinjlim_i f_j (X_i)
$$
associated with the bifunctor
$$
(i,j)\longmapsto f_j(X_i):I\times J\longrightarrow F{\quoi}.
$$
Thus \hyperref[I.prop9.6]{9.6} is a consequence of the following more
general assertion:


% END INCLUDED SGA4_English_current_buildable_source_20260906/source/01/drafts/SGA4_I_LINES_5941_6010_English.texfrag


% BEGIN INCLUDED SGA4_English_current_buildable_source_20260906/source/01/drafts/SGA4_I_LINES_6011_6089_English.texfrag
\begin{enonce*}{Corollary 9.8}\etiquette[9.8]{I.cor9.8}
Let $F$ be a $\Pi$-accessible category, $I$ an ordered set large
relative to $\Pi$, $J$ a small category such that $\card F\ell\leq
\Pi$ and such that projective limits of type $J$ are representable in
$F$, and $h:I\times J\rightarrow F$ a functor; then the canonical
morphism
\hypertarget{sga4.exp.I.equation.9.8.1}{}\begin{equation}
\varinjlim_i \varprojlim_j h(i,j)\longrightarrow \varprojlim_j
\varinjlim_i h(i,j)\tag{9.8.1}
\end{equation}
is an isomorphism. In other words, the functor
\hypertarget{sga4.exp.I.equation.9.8.2}{}\begin{equation}
\varinjlim_i:\Hhom(I,F)\longrightarrow F\tag{9.8.2}
\end{equation}
commutes with projective limits of type $J$ (for every small category
$J$\pageoriginale such that $\card F\ell\,(J)\leq\Pi$ and such that
projective limits of type $J$ are representable in $F$). Equivalently,
for every $J$ as above, the functor
\hypertarget{sga4.exp.I.equation.9.8.3}{}\begin{equation}
\varinjlim_j:\Hhom(J,F)\longrightarrow F\tag{9.8.3}
\end{equation}
is $\Pi$-accessible.
\end{enonce*}

By hypothesis, $F$ admits a conservative family of covariant
representable $\Pi$-accessible functors
$g:F\rightarrow(𝒰\text{-}\Ens)$, so it is immediately clear that it
suffices to prove \hyperref[I.cor9.8]{9.8} in the case $F=𝒰\text{-}\Ens$. We shall then
prove the assertion in the form that \hyperlink{sga4.exp.I.equation.9.8.2}{(9.8.2)} commutes with projective
limits of type $J$, where $\card F\ell\,J\leq\Pi$. Since $I$ is
filtered, we know that \hyperlink{sga4.exp.I.equation.9.8.2}{(9.8.2)} commutes with finite projective limits
\hyperref[I.prop2.8]{(2.8)}. By a standard argument (cf. \hyperref[I.prop2.3]{2.3}), it therefore remains to prove
that it commutes with products indexed by a set $J$ such that
$\card J\leq\Pi$. This reduces us to proving the bijectivity of
\hyperlink{sga4.exp.I.equation.9.8.1}{(9.8.1)} when $J$ is discrete. Let us prove injectivity: consider two
elements $a,b$ of the first member; they therefore come from
$\prod_jh(i_\circ,j)$ %%
for a suitable $i_\circ\in I$, say $\prod_j a(i_\circ,j)$ and
$\prod_j b(i_\circ,j)$; suppose that the elements
$\prod_jh(\infty,j)$ of the second member that they define are equal,
\ie that for every $j$, there exists $i(j)\geq i_\circ$ such that
$a(i_\circ,j)$ and $b(i_\circ,j)$ have the same image in $h(i,j)$.
Since $I$ is large relative to $\card J$, it follows that there exists
a common upper bound $i_1\in I$ of all the $i(j)$, which implies that
the two elements under consideration of $\prod_jh(i_\circ,j)$ have the
same image in $\prod_jh(i_1,j)$ and therefore define %%
the same element of the first member of \hyperlink{sga4.exp.I.equation.9.8.1}{(9.8.1)}, establishing
injectivity.\pageoriginale For surjectivity, let
$\prod_j a(\infty,j)$ be an element of the second member; thus, for
every $j\in J$, $a(\infty,j)$ comes from an element of $h(i(j),j)$ for
a suitable $i(j)\in I$. As before, one can find a common upper bound
$i\in I$ of the $i(j)$, so $\prod_j a(\infty,j)$ comes from an element
$\prod_j a(i,j)$ of $\prod_jh(i,j)$, and hence lies in the image of
\hyperlink{sga4.exp.I.equation.9.8.1}{(9.8.1)}. This completes the proof.

\begin{enonce*}{Corollary 9.9}\etiquette[9.9]{I.cor9.9}
Let $F$ be a $\Pi$-accessible (respectively accessible) category. Then,
for every category $J$ such that $\card F\ell\,J\leq\Pi$ (respectively
every small category $J$), and for every functor
$(X_j)_{j\in J}:J\rightarrow F$ whose projective limit in $F$ is
representable, if the $X_j$ are $\Pi$-accessible (respectively
accessible), then so is $\varprojlim X_j$.
\end{enonce*}

Indeed, the functor $F\rightarrow(𝒰\text{-}\Ens)$ represented by
$\varprojlim X_j$ is the projective limit of the functors represented
by the $X_j$, and one applies \hyperref[I.prop9.6]{9.6} (ii) to the
projective system formed by these functors.

\begin{enonce*}[remark]{Remark 9.10}\etiquette[9.10]{I.rem9.10}
In the statements \hyperref[I.prop9.6]{9.6},
\hyperref[I.cor9.7]{9.7}, \hyperref[I.cor9.8]{9.8}, and
\hyperref[I.cor9.9]{9.9}, one may everywhere replace the words
``$\Pi$-accessible'', ``accessible'' by ``$\Pi$-preaccessible'',
``preaccessible''. Indeed, the proof given also proves this variant of
the preceding statements.
\end{enonce*}


% END INCLUDED SGA4_English_current_buildable_source_20260906/source/01/drafts/SGA4_I_LINES_6011_6089_English.texfrag


% BEGIN INCLUDED SGA4_English_current_buildable_source_20260906/source/01/drafts/SGA4_I_LINES_6090_6222_English.texfrag
\begin{enonce*}{Proposition 9.11}\etiquette[9.11]{I.prop9.11}
Let $E$ be a category satisfying condition $L$
(\hyperref[I.def9.1]{9.1}), and let $C$ be a small full generating
subcategory (\hyperref[I.def7.1]{7.1}). Suppose that the following
condition is satisfied:
\begin{enumerate}
\renewcommand{\theenumi}{\alph{enumi}}
\renewcommand{\labelenumi}{\theenumi)}
\item $E$ is stable under kernels of pairs of arrows, and the functor
$\Ker$ on the category of pairs of arrows in $E$ is accessible (for
example, it\pageoriginale commutes with small filtered inductive
limits).

Then every object of $E$ is preaccessible (\hyperref[I.def9.3]{9.3}),
\textit{a fortiori} $E$ is preaccessible. Suppose that $C$ is even
generating by strict epimorphisms (\hyperref[I.def7.1]{7.1}), and
suppose that $E$ satisfies the following conditions:

\item $E$ is stable under fiber products.

\item Every small strictly epimorphic family in $E$ is universally
strictly epimorphic (\hyperref[I.sec10.3]{10.3}), or small direct sums
are representable in $E$ and every strict epimorphism of $E$ is a
universal strict epimorphism.
\end{enumerate}

Then every object of $E$ is accessible, \textit{a fortiori} $E$ is
accessible.
\end{enonce*}

The fact that, under (a), every object of $E$ is preaccessible follows
from the fact that for every object $X$ of $E$, the set of strict
subobjects of $X$ is small (\hyperref[I.prop7.4]{7.4}), and from the

\begin{enonce*}{Lemma 9.11.1}\etiquette[9.11.1]{I.lem9.11.1}
Under the conditions of (\hyperref[I.prop9.11]{9.11}a), let
$X\in\ob E$ and let $\Pi$ be a cardinal such that $E$ satisfies
$L_{\Pi}$ (\hyperref[I.def9.1]{9.1}), the functor $\Ker$ in
(\hyperref[I.prop9.11]{9.11}a) is $\Pi$-accessible, and the set of
strict subobjects of $X$ has cardinality bounded by $\Pi$. Then $X$ is
$\Pi$-preaccessible.
\end{enonce*}

Indeed, let $I$ be an ordered set large relative to $\Pi$, let
$(Y_i)_{i\in I}$ be an inductive system in $E$ with inductive limit
$Y$, and let $(u_i,v_i:X\rightrightarrows Y_i)$ be a pair of arrows
such that the composite pair
$ \xymatrix{ X \ar@<2pt>[r]^{u,v} \ar@<-2pt>[r]& Y }$ satisfies
$u=v$; let us prove that there exists $j\geq i$ in $I$ such that
$u_j=v_j$. To this end,\pageoriginale consider the inductive system of
pairs of arrows $(u_j,v_j:X\rightrightarrows Y_j)_{j\geq i}$, whose
inductive limit is $(u,v)$. By hypothesis one has
$X=\Ker(u,v)=\varinjlim_j\Ker(u_j,v_j)=\varinjlim_jX_j$, where
$X_j=\Ker(u_j,v_j)$. Now the $X_j$ are strict subobjects of $X$, so
there exists a subset $I'$ of $I$, formed by indices $j\geq i$, such
that $\card I'\leq\Pi$ and every $X_j$ is equal to one of the $X_{i'}$
$(i'\in I')$. Since $I$ is large relative to $\Pi$, there exists an
upper bound $j$ of $I'$ in $I$. Then $X_j$ contains all the $X_{j'}$
for $j'\geq i$, so $X_j\rightarrow X$ is an epimorphism (since the
family of the $X_{j'}\rightarrow X$ is epimorphic), and hence an
isomorphism because it is a strict monomorphism. Thus $u_j=v_j$, which
proves \hyperref[I.lem9.11.1]{9.11.1}.

The second assertion of \hyperref[I.prop9.11]{9.11} follows from the
first and from the

\begin{enonce*}{Lemma 9.11.2}\etiquette[9.11.2]{I.lem9.11.2}
Under the conditions of (\hyperref[I.prop9.11]{9.11}a), b), c), let
$X$ be an object of $E$, and let $\Pi$ be an infinite cardinal such
that $E$ satisfies $L_{\Pi}$ $\hyperref[I.def9.1]{(9.1)}$, that
$\card\ob C_{/X}\leq\Pi$, that for two objects $X'\rightarrow X$ and
$X''\rightarrow X$ of $C_{/X}$, $\fprod{X'}{X''}{X}$ is
$\Pi$-preaccessible, and finally that $X$ is $\Pi$-preaccessible.
Then $X$ is $\Pi$-accessible.
\end{enonce*}

With the notation of the proof of \hyperref[I.lem9.11.1]{9.11.1}, it
suffices to prove that every morphism $u:X\rightarrow Y$ comes from a
morphism $u_i:X\rightarrow Y_i$. Consider the family of morphisms
$Y_i\rightarrow Y$, which is strictly epimorphic; by b) and c), the
family of morphisms $\fprod{Y_i}{X}{Y}\rightarrow X$ is also strictly
epimorphic; on the other hand, for every $i$, since $C$ is generating
by strict epimorphisms, the family of arrows
$$
T\longrightarrow \fprod{Y_i}{X}{Y}
$$
with\pageoriginale source in $C$ is strictly epimorphic. In the second
alternative considered in c), one can find such a strictly epimorphic
arrow whose source is a (small) sum of objects of $C$. By the
transitivity of the notion of a universally strictly epimorphic family
(II \hyperref[II.def2.5]{2.5}), it then follows that the family of arrows
$T_\alpha\xrightarrow{f_\alpha}X$ with source in $C$ that factor
through one of the $\fprod{Y_i}{X}{Y}$ is strictly epimorphic. Let $J$
be the index set of this family; its cardinality is bounded by
$\card\ob C_{/X}\leq\Pi$. For every $\alpha\in J$, choose an
$i=i(\alpha)\in I$ and an $X$-morphism
$T_\alpha\rightarrow\fprod{Y_i}{X}{Y}$, or, equivalently, a
$v_\alpha:T_\alpha\rightarrow Y_i$ such that
$p_iv_\alpha=uf_\alpha$, where $p_i:Y_i\rightarrow Y$ is the
canonical morphism. Since $I$ is large relative to $\Pi$, one can
choose $i(\alpha)$ independently of $\alpha$, say equal to $i$. For
every pair of indices $\alpha,\beta\in J$, consider the composites
$$
\pr_1v_\alpha,\pr_2v_\beta:
\fprod{T_\alpha}{T_\beta}{X}\rightrightarrows Y_i{\quoi}.
$$
Their composites with $Y_i\rightarrow Y$ are equal, so, since
$\fprod{T_\alpha}{T_\beta}{X}$ is $\Pi$-preaccessible by hypothesis,
there exists an index $i'=i(\alpha,\beta)\geq i$ such that the
composites of the arrows under consideration with
$Y_i\rightarrow Y_{i'}$ are equal. Since the set of pairs
$\alpha,\beta$ has cardinality $\leq\Pi^2=\Pi$ ($\Pi$ being
infinite), it follows once again that $i'$ can be chosen independently
of $\alpha,\beta$. We may evidently suppose that $i'=i$. But then,
since the family $(f_\alpha:T_\alpha\rightarrow X)$ is strictly
epimorphic, one can find a morphism $u_i:X\rightarrow Y_i$ such that
$u_if_\alpha=v_\alpha$. One then has $p_iu_i=u$, since for every
$\alpha$ one has
$(p_iu_i)f_\alpha=p_i(u_if_\alpha)=p_iv_\alpha=uf_\alpha$, and the
family of the $f_\alpha$ is epimorphic. This completes the proof of
\hyperref[I.lem9.11.2]{9.11.2}.

\begin{enonce*}[remark]{Remark 9.11.3}\etiquette[9.11.3]{I.rem9.11.3}
As\pageoriginale a special case of \hyperref[I.prop9.11]{9.11}, one
sees that every object in the category $E$ is accessible in each of the
following two cases: 1) $E$ is a $𝒰$-abelian category with exact
filtered inductive limits and admitting a small generating subcategory
(here it is the \emph{second} alternative of c) that applies). 2) $E$
is the category of sheaves of sets on a topological space $X\in 𝒰$.
More generally, it suffices that $E$ be the category of sheaves of sets
on a $𝒰$-site (II \hyperref[II.def2.1]{2.1}), or that $E$ be a $𝒰$-topos (IV \hyperref[IV.def1.1]{1.1}).
\end{enonce*}


% END INCLUDED SGA4_English_current_buildable_source_20260906/source/01/drafts/SGA4_I_LINES_6090_6222_English.texfrag


% BEGIN INCLUDED SGA4_English_current_buildable_source_20260906/source/01/drafts/SGA4_I_LINES_6223_6376_English.texfrag
\begin{enonce*}{Definition 9.12}\etiquette[9.12]{I.def9.12}
Let $E$ be a category. A cardinal filtration of $E$ is an
increasing filtration $(\Filt^{\Pi}(E))_{\Pi\geq \Pi_\circ}$ of $E$
by strictly full subcategories $\Filt^{\Pi}(E)$,
indexed by the cardinals $\Pi\in 𝒰$ such that $\Pi\geq
\Pi_\circ$ (where $\Pi_\circ$ is a fixed infinite cardinal, depending on
the cardinal filtration under consideration), and satisfying the following
conditions:
\begin{enumerate}
\renewcommand{\theenumi}{\alph{enumi}}
\renewcommand{\labelenumi}{\theenumi)}
\item For every $\Pi\geq\Pi_\circ$, $\Filt^{\Pi}(E)$ is
equivalent to a small category.

\item $E$ satisfies $L_{\Pi_\circ}$ (\hyperref[I.def9.3]{9.3}), and for every
$\Pi\geq \Pi_\circ$, $\Filt^{\Pi}(E)$ is stable in $E$ under
filtered inductive limits indexed by ordered sets
$I$ large relative to $\Pi_\circ$ such that $\card I\leq \Pi$.

\item For every $\Pi\geq \Pi_\circ$, and every $X\in \ob E$, one can
find an isomorphism
$$
X\simeq \varinjlim_I X_i{\quoi},
$$
where $(X_i)_{i\in I}$ is an inductive system in $E$ indexed by an
ordered set $I$\pageoriginale large relative to $\Pi$
(\hyperref[I.def9.1]{9.1}), and where
the $X_i$ belong to $\Filt^{\Pi}(E)$. Moreover, if $X\in \ob
\Filt^{\Pi}(E)$, $\Pi'\geq \Pi$, one may choose $I$ such that $\card
I\leq {\Pi'}^{\Pi}$.
\end{enumerate}
\end{enonce*}

\setcounter{subsection}{12} \setcounter{subsubsection}{0}
\subsubsection{}\etiquette[9.12.1]{I.9.12.1}
Note that it follows from b) and c) that every $X\in \ob E$ belongs
to some $\Filt^{\Pi}(E)$, for $\Pi$ sufficiently large.

\begin{enonce*}[remark]{Example 9.12.2}\etiquette[9.12.2]{I.exem9.12.2}
Take $E=(𝒰\text{-}\Ens)$, $\Phi_\circ=\aleph_\circ$,
$\Filt^{\Pi}(E)=$ the full subcategory of $E$ consisting of the sets
such that $\card X\leq \Pi$. More generally:
\end{enonce*}

\begin{enonce*}{Proposition 9.13}\etiquette[9.13]{I.prop9.13}
Let $E$ be a $𝒰$-category. Suppose that $E$ is stable
under small filtered $\varinjlim$, under sums of two objects, and under
cokernels of pairs of arrows, that $E$ is stable under fiber products,
and that the epimorphisms in $E$ are universal strict
epimorphisms. Let $C$ be a small full subcategory of $E$
that is generating by strict epimorphisms. Let $\Pi_\circ$ be an
infinite cardinal $\geq \card F\ell\,  C$. For every cardinal $\Pi\geq
\Pi_\circ$, let $\Filt^{\Pi}(E)$ be the strictly full
subcategory of $E$ consisting of the objects $X$ of $E$ for which there exists a
strictly epimorphic family $(X_i\rightarrow X)_{i \in I}$ with target
$X$, such that $\card I\leq \Pi$ and $X_i\in \ob C$ for every
$i\in I$. Then $(\Filt^{\Pi}(E))_{\Pi\geq \Pi_\circ}$ is a
cardinal filtration of $E$. Moreover, for every $X\in \ob
\Filt^{\Pi}(E)$, the cardinality of the set of arrows of $C_{/X}$
(and \textit{a fortiori}, of the set of objects of $C_{/X}$) is
bounded by $\Pi^{\Pi_\circ}$.
\end{enonce*}

This last assertion is precisely \hyperref[I.prop7.6]{7.6}.

To show that this is a cardinal filtration, one must prove
conditions a), b), and c) of \hyperref[I.def9.12]{9.12}. Condition a)
follows at once from
\hyperref[I.var7.5.2]{7.5.2}. For b), suppose that $X=\varinjlim_I X_i$, with
$\card I\leq \Pi$ and the\pageoriginale $X_i\in \ob \Filt^{\Pi}(E)$.
Thus the family of canonical morphisms $X_i\rightarrow X$ is
strictly epimorphic, and by the hypothesis on the $X_i$, for
every $i\in I$ there exists a strictly epimorphic family $(X_{ij}\rightarrow
X_i)_{j\in J_i}$, with $\card J_i\leq \Pi$. Taking the sums
$\amalg_i X_i$ and $\amalg_{i,j}X_{ij}$, one sees by the transitivity of
universal strict epimorphisms (II \hyperref[II.def2.5]{2.5}) that the family of
composites $X_{ij}\rightarrow X_i\rightarrow X$ (indexed by
the disjoint union of the $J_i$, for $i\in I$) is strictly epimorphic.
But $\card J\leq \Pi^2=\Pi$, whence the conclusion. (NB. we have
used only the fact that $\varinjlim_I X_i$ exists, and not the
fact that $I$ is large relative to $\Pi_\circ$, or even that it is filtered.)
Finally, let us prove c). Note that for every $X\in \ob E$, the family of
arrows $X_i\rightarrow X$ with source in $\ob C$ is strictly
epimorphic, since $C$ is generating by strict
epimorphisms,
and is plainly small, so there exists a cardinal $\Pi\geq \Pi_\circ$
such that $X\in \ob \Filt^{\Pi}(E)$. It remains to prove that if
$\Pi'\geq \Pi\geq \Pi_\circ$ are cardinals, and if $X\in \ob
\Filt^{\Pi'}(E)$, then there is an isomorphism
$$
X\simeq \varinjlim_I X_i{\quoi},
$$
with $I$ large relative to $\Pi$, $\card I\leq {\Pi'}^{\Pi}$, and the $X_i$
in $\ob \Filt^{\Pi}(E)$. Now, since $C$ is generating by
strict epimorphisms, one has
\begin{equation*}
X\simeq \varinjlim_{C_{/X}} T\tag{$\ast$}{\quoi}.
\end{equation*}
Let us also note %%
that by successively adjoining to $C$ sums and cokernels of
pairs of arrows in $C$, and doing the same for the category thus obtained,
and so on, one reduces to the case where $C$ is stable under sums of two objects
in $E$ and under cokernels of pairs of arrows in $E$, without
losing the hypothesis $\Pi_\circ \geq \card F\ell\, C$. We may
suppose moreover that, if $E$ contains a\pageoriginale chosen initial
object $∅_E$, then $∅_E\in \ob C$. This implies
that the category $C_{/X}$ is stable under sums of two objects and
cokernels of pairs of arrows. It is then filtered, since it is
nonempty: if it were empty, the relation $(\ast)$
would show that $X$ is an initial object, and hence $C_{/X}$ would contain the
arrow $∅_E\rightarrow X$, a contradiction. Let
$I$ be the set, ordered by inclusion, of full subcategories
$i$ of $C_{/X}$ that are filtered and such that $\card \ob i\leq
\Pi$. For every $i\in I$, let $X_i\in \ob E$ be the inductive limit of the
composite functor $i\rightarrow C_{/X}\rightarrow E$, which exists by
hypothesis. Then plainly
$$
\varinjlim_I X_i \simeq \varinjlim_{C/X} T\simeq X{\quoi}.
$$
On the other hand, we have already noted that $\card \ob C_{/X}\leq
{\Pi'}^{\Pi_\circ}$. It follows that $I$ is large relative to $\Pi$, and
that $\card I\leq
({\Pi'}^{\Pi_\circ})^{\Pi}={\Pi'}^{\Pi\Pi_\circ}={\Pi'}^{\Pi}$ by
the following lemma, whose proof is left to the reader
(where we take $I=C_{/X}$, $c={\Pi'}^{\Pi_\circ}$):

\begin{enonce*}{Lemma 9.13.1}\etiquette[9.13.1]{I.lem9.13.1}
Let $J$ be a filtered category, and let $c$ and $\Pi$ be two cardinals such
that $c\geq \Sup(\card  F\ell\,  J,\Pi)$, and suppose that $\card
\Hom(j,J')\leq \Pi_\circ$ for every pair of objects $j$, $j'$ of
$J$. Let $I$ be the set of full filtered subcategories $i$
of $J$ such that $\card \ob i\leq \Pi$. Then, ordered by
inclusion, $I$ is large relative to $\Pi$, and one has $\card I\leq
c^{\Pi}$.
\end{enonce*}

This completes the proof of \hyperref[I.prop9.13]{9.13}.

\begin{enonce*}[remark]{Remark 9.13.2}\etiquette[9.13.2]{I.rem9.13.2}
A\pageoriginale slight additional effort should make it possible to
replace in \hyperref[I.prop9.13]{9.13} the hypothesis that $E$ is stable under
small filtered inductive limits by the hypothesis that $E$
satisfies $L$ (\hyperref[I.def9.1]{9.1}), provided $E$ is stable under
small sums, or that in $E$ every epimorphic \emph{family}
is universally epimorphic. In the proof
of c), one must then choose $\Pi_\circ$ such that $E$ satisfies $L_{\Pi_\circ}$,
and restrict oneself to those full subcategories $i$ of $C_{/X}$ that are
not only filtered, but such that the preordered set $\ob
i$ is large relative to $\Pi_\circ$. Using \hyperref[I.sec8]{Section 8} and the hypothesis that
$E$ satisfies $L_{\Pi_\circ}$, one then finds that the $X_i$
exist, and one should be able to conclude by a suitable variant of
\hyperref[I.lem9.13.1]{9.13.1}, which the editor has not verified.
\end{enonce*}


% END INCLUDED SGA4_English_current_buildable_source_20260906/source/01/drafts/SGA4_I_LINES_6223_6376_English.texfrag


% BEGIN INCLUDED SGA4_English_current_buildable_source_20260906/source/01/drafts/SGA4_I_LINES_6377_6439_English.texfrag
\begin{enonce*}{Proposition 9.14}\etiquette[9.14]{I.prop9.14}
Let $f:E\rightarrow F$ be an accessible functor
(\hyperref[I.def9.2]{9.2}) between two categories equipped with cardinal
filtrations $(\Filt^{\Pi}(E))_{\Pi\geq\Pi_\circ}$ and
$(\Filt^{\Pi}(F))_{\Pi\geq\Pi'_\circ}$. Then there exists a cardinal
$\Pi_1\geq\Sup(\Pi_\circ,\Pi'_\circ)$ such that for every cardinal
$\Pi\geq\Pi_1$, one has
\hypertarget{sga4.exp.I.equation.9.14.1}{}\begin{equation}
f(\Filt^{\Pi}(E))\subset \Filt^{\Pi^{\Pi_1}}(F){\quoi}.\tag{9.14.1}
\end{equation}
In particular, if $\Pi=2^c$ with $c\geq\Pi_1$, whence
$\Pi^{\Pi_1}=2^{c\Pi_1}=2^c=\Pi$, one has
\hypertarget{sga4.exp.I.equation.9.14.2}{}\begin{equation}
f(\Filt^{\Pi}(E))\subset \Filt^{\Pi}(F){\quoi}.\tag{9.14.2}
\end{equation}
\end{enonce*}

We immediately point out the

\begin{enonce*}{Corollary 9.15}\etiquette[9.15]{I.cor9.15}
Let\pageoriginale $E$ be a category equipped with two cardinal
filtrations\linebreak $(\Filt^{\Pi}(E))_{\Pi\geq\Pi_\circ}$ and
$({\Filt'}^{\Pi}(E))_{\Pi\geq %%
\Pi'_\circ}$. Then there exists a cardinal
$\Pi_1\geq\Sup(\Pi_\circ,\Pi'_\circ)$ such that, for every cardinal
$c\geq\Pi_1$, putting $\Pi=2^c$, one has
$\Filt^{\Pi}(E)={\Filt'}^{\Pi}(E)$.
\end{enonce*}

\noindent \emph{Proof of} \hyperref[I.prop9.14]{9.14}. Let $c$ be a cardinal such that
$c\geq\Sup(\Pi_\circ,\Pi'_\circ)$ and such that $f$ is
$c$-accessible. Let, moreover, $\Pi_1\geq c$ be such that
\begin{equation}
f(\Filt^c(E))\subset \Filt^{\Pi_1}(F){\quoi}.\tag{$\ast$}
\end{equation}
Such a $\Pi_1$ exists because $\Filt^c(E)$ is equivalent to a small
category \hyperref[I.def9.12]{(9.12 a)}), and hence so is $f(\Filt^c(E))$, so that one can
apply \hyperref[I.9.12.1]{9.12.1} to the objects of the latter to find a
$\Filt^{\Pi}(F)$ containing them all %%
(taking into account that the $\Filt^{\Pi}(F)$ are full
subcategories). Let then $\Pi\geq\Pi_1$ and
$X\in\ob\Filt^{\Pi}(E)$; let us prove that
$f(X)\in\Filt^{\Pi^{\Pi_1}}(F)$. Indeed, write
$$
X=\varinjlim_I X_i{\quoi},
$$
where the $X_i\in\ob\Filt^c(E)$, $I$ is large relative to $c$, and
$\card I\leq\Pi^c\leq\Pi^{\Pi_1}$
(\hyperref[I.def9.12]{9.12}c)). Since $f$ is $c$-accessible and $I$ is
large relative to $c$, one has
$$
f(X)\simeq \varinjlim_I f(X_i){\quoi},
$$
and since $\card I\leq\Pi^{\Pi_1}$ and
$f(X_i)\in\ob\Filt^{\Pi_1}(F)\subset
\ob\Filt^{\Pi^{\Pi_1}}(F)$ by $(\ast)$, one has
$f(X)\in\Filt^{\Pi^{\Pi_1}}(F)$ by
\hyperref[I.def9.12]{9.12} b), \cqfd

\setcounter{subsection}{15} \setcounter{subsubsection}{0}
\subsubsection{}\etiquette[9.15.1]{I.sec9.15.1}
The notion of cardinal filtration \hyperref[I.def9.12]{9.12} is of
interest mainly when the objects of $E$ are accessible. We point out
that it follows from \hyperref[I.prop9.11]{9.11} that this condition is
satisfied when, in addition to the hypotheses\pageoriginale of
\hyperref[I.prop9.13]{9.13}, one supposes that the functor $\Ker$ on
the category of pairs of arrows in $E$ is accessible. We also point
out:


% END INCLUDED SGA4_English_current_buildable_source_20260906/source/01/drafts/SGA4_I_LINES_6377_6439_English.texfrag


% BEGIN INCLUDED SGA4_English_current_buildable_source_20260906/source/01/drafts/SGA4_I_LINES_6440_6522_English.texfrag
\begin{enonce*}{Proposition 9.16}\etiquette[9.16]{I.prop9.16}
Let $E$ be a category equipped with a cardinal filtration
$(\Filt^{\Pi}(E))_{\Pi\geq\Pi_\circ}$. Suppose that the objects of
$\Filt^{\Pi_\circ}(E)$ are accessible objects of $E$; then
there exists a cardinal $\Pi_1\geq\Pi_\circ$ in $𝒰$ such that
the objects of $\Filt^{\Pi_\circ}(E)$ are $\Pi_1$-accessible; if
$\Pi_1$ is chosen in this way, then for every cardinal $\Pi\geq \Pi_1$,
we have (with the notation of \hyperref[I.sec9.3.1]{9.3.1}):
\hypertarget{sga4.exp.I.equation.9.16.1}{}\begin{equation}
E_{\Pi}\subset \Filt^{\Pi}(E)\subset
E_{(\Pi^{\Pi_\circ})}{\quoi}.\tag{9.16.1}
\end{equation}
\end{enonce*}

The existence of $\Pi_1$ follows immediately from the fact that
$\Filt^{\Pi_\circ}(E)$ is equivalent to a small
category (\hyperref[I.def9.12]{9.12}a)). Now let $\Pi\geq \Pi_1$. If $X$
  belongs to $\Filt^{\Pi}(E)$,
write $X\simeq \varinjlim_I X_i$ with $\card I\leq
\Pi^{\Pi_\circ}$, $X_i\in \ob \Filt^{\Pi_\circ}(E)$ for every $i$
(\hyperref[I.def9.12]{9.12}c)). It then follows from \hyperref[I.cor9.9]{9.9} that $X$ is
$\Pi^{\Pi_\circ}$-accessible, whence the second inclusion \hyperlink{sga4.exp.I.equation.9.16.1}{(9.16.1)}.
Suppose that $X$ is $\Pi$-accessible, and write $X\simeq
\varinjlim_I X_i$, with $I$ large relative to $\Pi$ and the $X_i\in \ob
\Filt^{\Pi}(E)$ (\hyperref[I.def9.12]{9.12}c)). By the hypothesis on $X$, the given
isomorphism $X\xrightarrow{\sim} \varinjlim_I X_i$ factors through one of the
$X_i$, so $X$ is isomorphic to a direct factor of that $X_i$. We
conclude that $X\in \ob \Filt^{\Pi}(E)$, and hence obtain the first inclusion
\hyperlink{sga4.exp.I.equation.9.16.1}{(9.16.1)}, by the following lemma.

\begin{enonce*}{Lemma 9.16.2}\etiquette[9.16.2]{I.lem9.16.2}
Every object of $E$ that is a direct factor of an object of
$\Filt^{\Pi}(E)$ belongs to $\Filt^{\Pi}(E)$.
\end{enonce*}

Indeed, if $X$ is the image of a projector $p$ on an object $Y$ of
$E$ (\ie an endomorphism $p$ such that $p^2=p$), and if $I$ is an
ordered\pageoriginale  filtered set, $X$ is the inductive limit
of the filtered inductive system $(Y_i)_{i\in I}$ defined by $Y_i=Y$
for every $i\in I$, $p:Y_i\rightarrow Y_j$ if $i<j$. Taking $I$
large relative to $\Pi_\circ$ and $\card I\leq \Pi$, we thus see that if
$Y\in \Filt^{\Pi}(E)$, then the same holds for $X$, by
\hyperref[I.def9.12]{9.12} b).

\begin{enonce*}{Corollary 9.17}\etiquette[9.17]{I.cor9.17}
Under the conditions of \hyperref[I.prop9.16]{9.16}, for every cardinal $c\geq
\Pi_1$, setting $\Pi=2^c$, $\Filt^{\Pi}(E)$ coincides with the
strictly full subcategory $E_{\Pi}$ of $E$ consisting of the
$\Pi$-accessible objects.
\end{enonce*}

\begin{enonce*}{Corollary 9.18}\etiquette[9.18]{I.cor9.18}
Let $E$ be a category satisfying condition $L_{\Pi}$
(\hyperref[I.def9.1]{9.1}), where $\Pi$ is a cardinal, and let $C$ be a
full subcategory
of $E$. Denote by $\Ind(C)_{\Pi}$ the full subcategory of
the category $\Ind(C)$ of ind-objects of $C$ (\hyperref[I.sec8.2]{8.2})
consisting of
ind-objects of the form $(X_i)_{i\in I}$, where $I$ is an ordered set
large relative to $\Pi$. Consider the canonical functor
\hypertarget{sga4.exp.I.equation.9.18.1}{}\begin{equation}
(X_i)_{i\in I}\longmapsto \varinjlim_I
X_i:\Ind(C)_{\Pi}\longrightarrow E{\quoi}.\tag{9.18.1}
\end{equation}
\begin{enumerate}
\renewcommand{\theenumi}{\alph{enumi}}
\renewcommand{\labelenumi}{\theenumi)}
\item This functor is fully faithful if and only if every object $X$
of $C$ is a $\Pi$-accessible object
(\hyperref[I.def9.3]{9.3}) of $E$.

\item Assume the conditions of \hyperref[I.cor9.17]{9.17}, in particular
$\Pi=2^c$, and take $C=\Filt^{\Pi}(E)$. Then the functor \hyperlink{sga4.exp.I.equation.9.18.1}{(9.18.1)}
is an equivalence of categories.
\end{enumerate}
\end{enonce*}

Assertion a) is an immediate generalization of \hyperref[I.prop8.7.5]{8.7.5}
a), and is proved in the same way. Part b) then follows from
\hyperref[I.cor9.17]{9.17} and condition \hyperref[I.def9.12]{9.12} c) on
cardinal filtrations, which implies that the functor under consideration is
essentially surjective.


% END INCLUDED SGA4_English_current_buildable_source_20260906/source/01/drafts/SGA4_I_LINES_6440_6522_English.texfrag


% BEGIN INCLUDED SGA4_English_current_buildable_source_20260906/source/01/drafts/SGA4_I_LINES_6523_6659_English.texfrag
\begin{enonce*}{Corollary 9.19}\etiquette[9.19]{I.cor9.19}
Under\pageoriginale the conditions of \hyperref[I.cor9.17]{9.17}, let
$C=\Filt^{\Pi}(E)$, let $F$ be a $𝒰$-category, and consider the
functor
\hypertarget{sga4.exp.I.equation.9.19.1}{}\begin{equation}
\Hhom(E,F)_{\Pi}\longrightarrow \Hhom(C,F)\tag{9.19.1}
\end{equation}
induced by the “restriction to $C$” functor $f\mapsto f| C$,
where the source of \hyperlink{sga4.exp.I.equation.9.19.1}{(9.19.1)} is the category of
$\Pi$-accessible functors from $E$ to $F$ (\hyperref[I.def9.2]{9.2}). The
preceding functor is
fully faithful. If $F$ satisfies condition $L_{\Pi}$
(\hyperref[I.def9.1]{9.1}), then the functor \hyperlink{sga4.exp.I.equation.9.19.1}{(9.19.1)} is an equivalence of
categories.
\end{enonce*}

The first assertion is proved as in \hyperref[I.prop7.8]{7.8}, using
\hyperref[I.cor9.18]{9.18} b). The second is obtained by constructing a
quasi-inverse functor to \hyperlink{sga4.exp.I.equation.9.19.1}{(9.19.1)}, associating to every
functor $f:C\rightarrow F$ the functor $(X_i)_{i\in I}\rightarrow
\varinjlim_i f(X_i)$ from $\Ind(E_{\Pi})$ to $F$, and using
\hyperref[I.cor9.18]{9.18} b) to deduce from it a functor $\bar{f}:E\rightarrow
F$. It remains only to %%
show that the%%
latter is $\Pi$-accessible. But this is proved in the same way as the
analogous assertion \hyperref[I.prop8.7.3]{8.7.3}.

\begin{enonce*}{Corollary 9.20}\etiquette[9.20]{I.cor9.20}
Let $E$ be a category admitting a cardinal filtration and such
that every object of $E$ is accessible (\cf \hyperref[I.prop9.16]{9.16}), and let $F$
be a category. Then:
\begin{enumerate}
\renewcommand{\theenumi}{\alph{enumi}}
\renewcommand{\labelenumi}{\theenumi)}
\item The category $\Hhom(E,F)_{acc}$ of accessible functors
from $E$ to $F$ (\hyperref[I.def9.2]{9.2}) is a
$𝒰$-category. ({\rm NB} we
recall (\hyperref[I.sec9.0]{9.0}) that the given categories $E$, $F$
are assumed
to be $𝒰$-categories.)

\item Suppose that $F$ is stable under small filtered inductive limits.
For every full subcategory $C$ of $E$ equivalent
to a small category, with inclusion functor $i:C\rightarrow E$,
consider the corresponding functor
$$
i_{!}:\Hhom(C,F)\longrightarrow \Hhom(E,F)
$$
$\hyperref[I.prop5.1]{(5.1)}$.\pageoriginale A functor $f:E\rightarrow F$ is
accessible if and only if there exists a small
subcategory $C$ of $E$ such that $f$ belongs to the essential image
of the preceding functor $i_{!}$.
\end{enumerate}
\end{enonce*}

\noindent \emph{Proof}.
\begin{enumerate}
\renewcommand{\theenumi}{\alph{enumi}}
\renewcommand{\labelenumi}{\theenumi)}
\item It suffices to prove that for every cardinal $\Pi$ such that $E$
satisfies $L_{\Pi}$, the full subcategory $\Hhom(E,F)_{\Pi}$ of
$\Hhom(E,F)_{acc}$ is a $𝒰$-category. It is
clearly enough to verify this for cardinals of the form $2^c$, with
$c$ sufficiently large. But then this follows from \hyperref[I.cor9.19]{9.19}, since
($C=\Filt^{\Pi}(E)$ being essentially small) $\Hhom(C,F)$ is
obviously a $𝒰$-category.

\item By transitivity of the formation of the functors $i_{!}$, one
may restrict in the statement to subcategories $C$ of the form
$\Filt^{\Pi}(E)$, where $\Pi$ is as in \hyperref[I.cor9.17]{9.17}. One then sees
readily that the composite functor $\Hhom(\Filt^{\Pi}(E),F)\rightarrow
\Hhom(E,F)_{\Pi}\rightarrow \Hhom(E,F)$, where the first arrow is
quasi-inverse to \hyperlink{sga4.exp.I.equation.9.19.1}{(9.19.1)}, and the second is the inclusion, is
none other than the functor $i_{!}$, up to isomorphism. Thus assertion
b) follows from \hyperref[I.cor9.19]{9.19}.
\end{enumerate}

\begin{enonce*}[remark]{*Exercise 9.20.1}\etiquette[9.20.1]{I.exer9.20.1}
(The present exercise uses the notions of site and topos,
developed in Exposés II and IV.) Let $E$ be a
$𝒰$-topos, let $𝒱$ be a universe such that $𝒰\in
𝒱$, let $C$ be a small full generating subcategory of the
$𝒰$-topos $E$, and let $\Pi_\circ\in 𝒰$ be an infinite cardinal
such that $\Pi_\circ\in \card F\ell\,  C$. For every cardinal $\Pi\geq
\Pi_\circ$, $\Pi\in 𝒰$, let $\Filt^{\Pi}(E)$ be the
strictly full subcategory of $E$ consisting of the objects $X$ such
that there exists a strictly epimorphic family\pageoriginale
$(X_i\rightarrow X)_{i\in I}$ with target $X$, such that $\card I\leq
\Pi$ and $X_i\in \Ob C$ for every $i\in I$.
\begin{enumerate}
\renewcommand{\theenumi}{\alph{enumi}}
\renewcommand{\labelenumi}{\theenumi)}
\item Show that $(\Filt^{\Pi}(E))_{\Pi\geq \Pi_\circ}$ is a
cardinal filtration of the $𝒰$-category $E$, and that one can
choose $\Pi_\circ$ such that for every $\Pi\geq \Pi_\circ$, one has
the inclusions
$$
E_{\Pi}\subset \Filt^{\Pi}(E)\subset E_{\Pi^{\Pi_\circ}}{\quoi},
$$
where, for every cardinal $c$, $E_c$ denotes the strictly full
subcategory of $E$ consisting of the $c$-accessible objects.

\item Choosing a cardinal $\Pi\geq \Pi_\circ$ such that
$\Pi=\Pi^{\Pi_\circ}$ (for example, $\Pi$ of the form $2^c$, with
$c>\Pi_\circ$), and putting $C=\Filt^{\Pi}(E)$, show that there is an
equivalence of categories
$$
\Ind(C)_{\Pi}\xrightarrow{\approx} E{\quoi},
$$
(notation $\Ind(C)_{\Pi}$ of \hyperref[I.cor9.18]{9.18}).

\item Retain the notation of b), and let $𝒱$ be a universe such
that $𝒰\subset 𝒱$, let $%%
E^{\sim{\phantom{,}}}_{\underline{V}}$ be the $𝒱$-topos of
$𝒱$-sheaves on the site $E$ (endowed with its canonical
topology). Denote by $\Ind(C,𝒱)_{\Pi}$ the category of
$𝒱\text{-}\Ind$-objects of $C$ indexed by%%
 preordered index sets \emph{large relative to} $\Pi$. Show that there is
an equivalence of categories
$$
\Ind(C,𝒱)_{\Pi}\xrightarrow{\approx}%%
E^{\tilde{\phantom{,}}}_{𝒱}{\quoi}.
$$

\item Let $c\in 𝒱$ be a cardinal, and let $\Ind(E,𝒱)'_c$ be the
full subcategory of $\Ind(E,𝒱)$ consisting of the
$𝒱$-ind-objects of $E$ indexed by a preordered set that
is large relative to every cardinal $<c$. Taking $c=\card 𝒰$,
show\pageoriginale that there is an equivalence of categories
$$
\displaylines{\hfill \Ind(E,𝒱)'_{c}\xrightarrow{\approx}%%
E^{\tilde{\phantom{,}}}_{𝒱}{\quoi}.\hfill{\llap{\textbf{*}}}}
$$
\end{enumerate}
\end{enonce*}

\setcounter{subsection}{20}

% END INCLUDED SGA4_English_current_buildable_source_20260906/source/01/drafts/SGA4_I_LINES_6523_6659_English.texfrag


% BEGIN INCLUDED SGA4_English_current_buildable_source_20260906/source/01/drafts/SGA4_I_LINES_6660_6750_English.texfrag
\subsection{}\etiquette[9.21]{I.sec9.21}
The present section \hyperref[I.sec9.21]{9.21} develops technical
preliminaries for the proof of Theorem \hyperref[I.them9.22]{9.22}
below, which is the main result of the present section
\hyperref[I.sec9]{9}. Let
\hypertarget{sga4.exp.I.equation.9.21.1}{}\begin{equation}
p:E\rightarrow B\tag{9.21.1}
\end{equation}
be a fibered functor, where $B$ is a small category, and where the
inverse-image functors
$$
f^\ast:E_\beta\longrightarrow E_\alpha
$$
associated with the arrows $f:\alpha\rightarrow\beta$ of $B$ are
accessible (\hyperref[I.def9.2]{9.2}). In particular, the fiber
categories $E_\alpha$ $(\alpha\in\ob B)$ satisfy condition $L$
(\hyperref[I.def9.1]{9.1}), and therefore, since $B$ is small, there
exists a cardinal $\Pi'_\circ\in 𝒰$ such that all the categories
$E_\alpha$ satisfy condition $L_{\Pi'_\circ}$. Let
$\Pi_\circ\in 𝒰$ be an infinite cardinal $>\Pi'_\circ$, so that one
has
\hypertarget{sga4.exp.I.equation.9.21.2}{}\begin{equation}
\Pi_{\circ}>\Pi'_{\circ}, E_\alpha \mbox{ satisfies }
L_{\Pi'_{\circ}} \mbox{ for every } \alpha\in \ob B.\tag{9.21.2}
\end{equation}
Moreover, the hypotheses made immediately imply the existence of a
cardinal $c\in 𝒰$ satisfying the following conditions:
\hypertarget{sga4.exp.I.equation.9.21.3}{}\begin{equation}
\left\{
\begin{tabular}{rp{6cm}}
a) & $c$ is infinite,\\
b) & $c\geq \card F\ell\,B$,\\
c) & for every arrow $f:\alpha\rightarrow\beta$ in $B$, the functor
$f^\ast:E_\beta\rightarrow E_\alpha$ is $c$-accessible.
\end{tabular}\right.
\tag{9.21.3}
\end{equation}
Suppose\pageoriginale moreover that, for every $\alpha\in\ob B$, one
can find a full subcategory $C_\alpha$ of $E_\alpha$ satisfying the
following conditions:
\hypertarget{sga4.exp.I.equation.9.21.4}{}\begin{equation}
\left\{
\begin{tabular}{rp{6cm}}
a) & For every $X\in\ob C_\alpha$, $X$ is $c$-accessible in
$E_\alpha$ \hyperref[I.def9.3]{(9.3)}.\\
b) & For every $X\in\ob E_\alpha$, one can find an isomorphism
$X\simeq\varinjlim_I X_i$ in $E_\alpha$, where the $X_i$ belong to
$C_\alpha$ and $I$ is an ordered set large relative to $c$.
\end{tabular}\right.
\tag{9.21.4}
\end{equation}

Let $\Pi$ be a cardinal such that
\hypertarget{sga4.exp.I.equation.9.21.5}{}\begin{equation}
\Pi\geq \Pi_\circ\quav\mbox{ hence }
\Pi>\Pi'_\circ{\quoi},\tag{9.21.5}
\end{equation}
and, for every $\alpha\in\ob B$, let
\hypertarget{sga4.exp.I.equation.9.21.6}{}\begin{equation}
C^{\Pi}_{\alpha}\subset E_{\alpha}\tag{9.21.6}
\end{equation}
be formed by the objects that can be represented in the form
$\varinjlim_I X_i$, where $I$ is an ordered set large relative to
$\Pi'_\circ$, such that $\card I\leq\Pi$, and where the $X_i$ belong
to $C_\alpha$. It is then clear, by \hyperref[I.var7.5.2]{7.5.2}, that
$C^{\Pi}_{\alpha}$ is essentially small, \ie is equivalent to a small
category.

Let
\hypertarget{sga4.exp.I.equation.9.21.7}{}\begin{equation}
F=\Hhom_B(B,E)\tag{9.21.7}
\end{equation}
be the category of sections of $E$ over $B$, and let
\hypertarget{sga4.exp.I.equation.9.21.8}{}\begin{equation}
F^{\Pi}\subset F\tag{9.21.8}
\end{equation}
be the strictly full subcategory of $F$ formed by the sections
$X:\alpha\mapsto X(\alpha)$ such that for every $\alpha\in\ob B$ one
has
$$
X(\alpha)\in C^{\Pi}_{\alpha}{\quoi}.
$$
Since\pageoriginale the $C^{\Pi}_{\alpha}$ are essentially small, it
is clear that the same is true of the category $F^{\Pi}$. We shall
show that this category is generating, and more precisely:


% END INCLUDED SGA4_English_current_buildable_source_20260906/source/01/drafts/SGA4_I_LINES_6660_6750_English.texfrag


% BEGIN INCLUDED SGA4_English_current_buildable_source_20260906/source/01/drafts/SGA4_I_LINES_6751_6838_English.texfrag
\begin{enonce*}{Lemma 9.21.9}\etiquette[9.21.9]{I.lem9.21.9}
Under the preceding conditions and with the preceding notation, the
following statements hold:
\begin{enumerate}
\renewcommand{\theenumi}{\roman{enumi}}
\renewcommand{\labelenumi}{(\theenumi)}
\item Every object of $F$ is accessible (\hyperref[I.def9.3]{9.3}). If
$d$ is a cardinal $\geq\Pi'_\circ$, and if
$\alpha\mapsto X(\alpha)$ is an element of $F$ such that, for every
$\alpha$, $X(\alpha)$ is $d$-accessible in $E_\alpha$, then $X$ is
$d$-accessible in $F$.

\item Suppose that $\Pi\leq c$, or that $\Pi'_\circ$ is finite (\ie
the $E_\alpha$ are stable under small filtered $\varinjlim$). Then
every object $X$ of $F$ is isomorphic to an object of the form
$\varinjlim_I X_i$, where the $X_i$ belong to $F^{\Pi}$ and $I$ is
large relative to $\Pi$.
\end{enumerate}
\end{enonce*}

To prove (i), note that by \hyperlink{sga4.exp.I.equation.9.21.4}{(9.21.4)} and
\hyperref[I.cor9.9]{9.9}, every object of $E_\alpha$ is accessible. On
the other hand, $F$ satisfies condition $L_{\Pi'_\circ}$ by the


\begin{enonce*}{Lemma 9.21.10}\etiquette[9.21.10]{I.lem9.21.10}
Let $p:E\rightarrow B$ be a fibered functor, let
$F=\Hhom_B(B,E)$, let $I$ be a category, and let $i\mapsto X_i$ be a
functor from $I$ to $F$. For $\varinjlim_I X_i$ to be representable in
$F$, it suffices that, for every $\alpha\in\ob B$,
$\varinjlim_I X_i(\alpha)$ be representable in the fiber category
$E_\alpha$; when this is so, $\varinjlim_I X_i$ is ``computed
argument by argument''. In particular, if the fiber categories satisfy
condition $L_{\Pi'_\circ}$ (where $\Pi'_\circ$ is a given cardinal),
then so does $F$.
\end{enonce*}

We\pageoriginale leave the details of the (easy) proof of
\hyperref[I.lem9.21.10]{9.21.10} to the reader, merely noting that it
is convenient to use the following result, whose proof is immediate:

\begin{enonce*}{Corollary 9.21.10.1}\etiquette[9.21.10.1]{I.cor9.21.10.1}
With the hypotheses and notation of \hyperref[I.lem9.21.10]{9.21.10}
for $p:E\rightarrow B$ and for $I$, for every $\alpha\in\ob B$, the
inclusion functor $E_\alpha\rightarrow E$ commutes with inductive
limits of type $I$.
\end{enonce*}

Returning now to the general conditions of
\hyperref[I.lem9.21.9]{9.21.9}, let $X$ be an object of $F$. Since, for
every $\alpha\in\ob B$, $X(\alpha)$ is accessible in $E_\alpha$,
there exists a cardinal $d\in 𝒰$ such that, for every $\alpha$,
$X(\alpha)$ is $d$-accessible in $E_\alpha$. We may choose
$d\geq\Pi'_\circ$, so that $F$ satisfies $L_d$ by
\hyperref[I.lem9.21.10]{9.21.10}. Using
\hyperref[I.lem9.21.10]{9.21.10} once more to compute the inductive
limits $\varinjlim_IY_i$ in $F$, with $I$ large relative to $d$, one
immediately sees that $X$ is $d$-accessible. This proves (i).

We shall now prove \hyperref[I.lem9.21.9]{9.21.9} (ii) in several
steps (\hyperlink{sga4.exp.I.equation.9.21.11}{(9.21.11)} to (\hyperref[I.lem9.21.16]{9.21.16})).

Let then
$$
X:\alpha\longmapsto X(\alpha)
$$
be an object of $F$. By \hyperlink{sga4.exp.I.equation.9.21.4}{(9.21.4 b)}), for every $\alpha\in\ob B$ one
can find an ordered set $I_\alpha$ large relative to $c$, and an
isomorphism
$X(\alpha)=\varinjlim_{i\in I_\alpha}X(\alpha)_i$, where
$i\mapsto X(\alpha)_i$ is an inductive system of type $I_\alpha$ in
$E_\alpha$, with the $X(\alpha)_i$ in $C_\alpha$. Consider the product
ordered set $I=\Pi_\alpha I_\alpha$; it is clear that it is large
relative to $c$, and that the preceding inductive
systems\pageoriginale give rise to inductive systems in the $E_\alpha$,
\hypertarget{sga4.exp.I.equation.9.21.11}{}\begin{equation}
i\longmapsto X(\alpha)_i\in \ob C_{\alpha}, i\in \ob
I{\quoi},\tag{9.21.11}
\end{equation}
indexed by the same ordered set $I$, and to isomorphisms
\hypertarget{sga4.exp.I.equation.9.21.12}{}\begin{equation}
X(\alpha)\xrightarrow{\sim}\varinjlim_I X(\alpha)_i\mbox{ in }
E_{\alpha}{\quoi},\tag{9.21.12}
\end{equation}
for every $\alpha\in\ob B$. In what follows we suppose that data
\hyperlink{sga4.exp.I.equation.9.21.11}{(9.21.11)} and \hyperlink{sga4.exp.I.equation.9.21.12}{(9.21.12)} have been fixed.


% END INCLUDED SGA4_English_current_buildable_source_20260906/source/01/drafts/SGA4_I_LINES_6751_6838_English.texfrag


% BEGIN INCLUDED SGA4_English_current_buildable_source_20260906/source/01/drafts/SGA4_I_LINES_6839_6963_English.texfrag
\begin{enonce*}{Lemma 9.21.13}\etiquette[9.21.13]{I.lem9.21.13}
Under the preceding conditions, one can find a map
$\varphi:I\rightarrow I$ such that $\varphi(i)\geq i$ for every
$i\in I$, and a map $(i,f)\mapsto\lambda(i,f)$ from $I\times %%
\mathrm{F}{}l\,B$ to $%%
\mathrm{F}{}l\,E$, satisfying the following conditions:
\begin{enumerate}
\renewcommand{\theenumi}{\alph{enumi}}
\renewcommand{\labelenumi}{\theenumi)}
\item For $i\in I$ and $(f:\alpha\rightarrow\beta)\in F\ell B$,
$\lambda(i,f)$ is an $f$-morphism from $X(\alpha)_i$ to
$X(\beta)_{\varphi(i)}$, making the diagram
\[
\xymatrix@C=3cm@R=.8cm{
 X(\alpha)\ar[r]^{X(f)} & X(\beta)\\
           &  X(\beta)_{\varphi(i)}\ar[u]\\
X(\alpha)_{i} \ar[uu]\ar[ru]^{\lambda(i,f)}&
                  \\
\alpha \ar[r]^{f}& \beta\quad , }
\]
commutative, where the vertical arrows are the canonical morphisms
deduced from \hyperlink{sga4.exp.I.equation.9.21.12}{(9.21.12)}.

\item For\pageoriginale every $i\in I$, and every
$(f:\alpha\rightarrow\beta)\in %%
\mathrm{F}{}l\,B$, the diagram
\[
\xymatrix@C=3cm@R=.9cm@L=6pt{
                      & X(\beta)_{\varphi^{2}(i)}\\
  X(\alpha)_\varphi(i) \ar[ru]^{\lambda(\varphi(i),f)}  & X(\beta)_{\varphi(i)} \ar[u]\\
X(\alpha)_i  \ar[ru]_{\lambda(i,f)} \ar[u]         \\
\alpha \ar[r]^{f} & \beta\quad . }
\]
is commutative.

\item For every pair of consecutive arrows
$\alpha\xrightarrow{f}\beta\xrightarrow{g}\gamma$ of $B$, the
following diagram is commutative:
\[
\xymatrix@C=2cm@L=6pt@R=.5cm{
X(\alpha) \ar[r]^{X(f)}& X(\beta) \ar[r]^{X(g)}            & X(\gamma)\\
          &                      & X(\gamma)_{\varphi^{2}(i)}\ar[u]\\
          & X(\beta)_{\varphi(i)} \ar[uu]\ar[ru]^{\lambda(\varphi(i),g)} & X(\gamma)_{\varphi(i)}\ar[u]\\
X(\alpha)_i \ar[uuu]\ar[ru]^{\lambda(i,f)} \ar[rru]_{\lambda(i,gf)}&
 \\
\alpha \ar[r]^{f}& \beta \ar[r]^{g}& \gamma\quad , }
\]
where the vertical arrows are those deduced from \hyperlink{sga4.exp.I.equation.9.21.12}{(9.21.12)}.
\end{enumerate}
\end{enonce*}

One\pageoriginale should moreover note that the commutativity of the
two upper trapezoids in c) is already contained in a), so c) in fact
asserts the commutativity of the lower triangle of the diagram under
consideration.

Let us prove \hyperref[I.lem9.21.13]{9.21.13}. Let $i\in I$, and let
$f:\alpha\rightarrow\beta$ be an arrow in $B$. Consider the composite
$X(\alpha)_i\rightarrow X(\alpha)\xrightarrow{X(f)}X(\beta)\simeq
\varinjlim_jX(\beta)_j$; it defines a morphism in $E$:
$$
X(\alpha)_i\rightarrow f^\ast(X(\beta))\simeq
f^\ast(\varinjlim_jX(\beta)_j)\simeq
\varinjlim_jf^\ast(X(\beta)_j){\quoi},
$$
where the last equality follows from the fact that $f^\ast$ is
$c$-accessible \hyperlink{sga4.exp.I.equation.9.21.3}{(9.21.3 c)}%%
 and that $I$ is large relative to $c$. By \hyperlink{sga4.exp.I.equation.9.21.4}{(9.21.4 a)}), since
$X(\alpha)_i\in\ob C_\alpha$, the morphism under consideration
factors through some $f^\ast(X(\beta)_j)$, where $j$ \textit{a priori}
depends on $i$ and on $f\in %%
\mathrm{F}{}l\,(B)$. But by \hyperlink{sga4.exp.I.equation.9.21.3}{(9.21.3 b)}), and since $I$ is large
relative to $c$, one may choose $j$ independently of $f$, say
$j=\varphi(i)$. Since $I$ is filtered, one may suppose
$\varphi(i)\geq i$. One thus obtains, for $i\in I$ and
$f\in %%
\mathrm{F}{}l\,B$, a morphism
$$
X(\alpha)_i\longrightarrow f^\ast(X(\beta)_{\varphi(i)}){\quoi},
$$
or, equivalently, an $f$-morphism
$$
\lambda(i,f):X(\alpha)_i\longrightarrow
X(\beta)_{\varphi(i)}{\quoi},
$$
which by construction makes the diagram
\[\vcenter{
\xymatrix@C=3cm{
X(\alpha)\ar[r]^{X(f)} & X(\beta)\\
          & X(\beta)_{\varphi(i)} \ar[u]\\
X(\alpha)_i \ar[uu] \ar[ru]^{\lambda(i,f)}&\\
\alpha \ar[r]^{f} & \beta \quad . }}\leqno{D(f)}
\]
commutative. Now\pageoriginale, for given $i,f,g$, consider the lower
triangle of the diagram in \hyperref[I.lem9.21.13]{9.21.13} c); it is not clear that it is
commutative, but the two composites
$X(\alpha)_i\rightrightarrows X(\gamma)_{\varphi^2(i)}$ become equal
after composition with
$X(\gamma)_{\varphi^2(i)}\rightarrow X(\gamma)$, as follows from the
commutativity of the two upper trapezoids and of the outer trapezoid,
which are respectively the diagrams $D(f)$, $D(g)$, and $D(gf)$.
Thus, since $X(\alpha)_i$ is $c$-accessible \hyperlink{sga4.exp.I.equation.9.21.4}{(9.21.4 a)}) and
$X(\gamma)\simeq\varinjlim_{j\in I}X(\gamma)_j$, with $I$ large
relative to $c$, it follows that one can find an element
$j\geq\varphi^2(i)$ of $I$ such that the two arrows under consideration
become equal after composition with
$X(\gamma)_{\varphi^2(i)}\rightarrow X(\gamma)_j$. A priori, $j$
depends on $i,f,g$. But, for fixed $i$, the set of possible pairs
$f,g$ has cardinality $\leq c^2=c$, by \hyperlink{sga4.exp.I.equation.9.21.3}{(9.21.3 a)} and b)), and hence,
since $I$ is large relative to $c$, one may choose $j$ independently of
$f$ and $g$, say
$j=\varphi'(i)\geq\varphi^2(i)\geq\varphi(i)$. For every $i\in I$
and $f:\alpha\rightarrow\beta$, let then
$$
\lambda'(i,f):X(\alpha)_i\longrightarrow X(\beta)_{\varphi'(i)}
$$
be the composite $f$-morphism
$X(\alpha)_i\rightarrow X(\beta)_{\varphi(i)}\rightarrow
X(\beta)_{\varphi'(i)}$, where the second arrow is the transition
morphism. It is then immediate from the construction that
$(\varphi',\lambda')$ satisfies conditions \hyperref[I.lem9.21.13]{9.21.13} a) and c) for
$(\varphi,\lambda)$. Proceeding in the same way for condition b), one
sees that $\varphi'$ may be chosen so that this condition is also
satisfied for $(\varphi',\lambda')$. This completes the proof of
\hyperref[I.lem9.21.13]{(9.21.13)}.

\setcounter{subsection}{21} \setcounter{subsubsection}{13}

% END INCLUDED SGA4_English_current_buildable_source_20260906/source/01/drafts/SGA4_I_LINES_6839_6963_English.texfrag


% BEGIN INCLUDED SGA4_English_current_buildable_source_20260906/source/01/drafts/SGA4_I_LINES_6964_7051_English.texfrag
\subsubsection{}\etiquette[9.21.14]{I.sec9.21.14}
Now\pageoriginale let
$$
J\subset I
$$
be a subset of $I$ satisfying the following conditions:
\begin{enumerate}
\renewcommand{\theenumi}{\alph{enumi}}
\renewcommand{\labelenumi}{\theenumi)}
\item $J$ is filtered,

\item for every $j\in J$, one has $\varphi(j)\in J$,

\item For every $\alpha\in\ob B$,
$X_J(\alpha)=\varinjlim_{i\in J}X(\alpha)_i$ is representable in
$E_\alpha$.
\end{enumerate}
Conditions a) and c) are satisfied in particular if $J$ is large
relative to $\Pi'_\circ$ \hyperlink{sga4.exp.I.equation.9.21.2}{(9.21.2)}. It then follows from
\hyperref[I.cor9.21.10.1]{9.21.10.1} that, for every
$\alpha\in\ob B$, $X_J(\alpha)$ is the inductive limit
$\varinjlim_{i\in J}X(\alpha)_i$ in $E$. Now, for an arrow
$f:\alpha\rightarrow\beta$ of $B$, the arrows $\lambda(i,f)$, with
$i$ varying in $J$, define, by (\hyperref[I.lem9.21.13]{9.21.13} b)),
a morphism of ind-objects from $(X(\alpha)_i)_{i\in J}$ to
$(X(\beta)_i)_{i\in J}$, whence a morphism
$$
X_J(f):X_J(\alpha)\rightarrow X_J(\beta)
$$
on the inductive limits, which is manifestly an $f$-morphism. Using
(\hyperref[I.lem9.21.13]{9.21.13} (c)), one obtains the transitivity
relations
$$
X_J(gf)=X_J(g)X_J(f){\quoi},
$$
so that a section $X_J\in\ob F$ of $E$ over $B$ has been defined.
Finally, (\hyperref[I.lem9.21.13]{9.21.13} a)) shows us that the
canonical morphisms
$$
X_J(\alpha)\longrightarrow X(\alpha)\quav \alpha\in \ob B{\quoi},
$$
are functorial in $\alpha$, so there is a canonical morphism
$$
u_J:X_J\longrightarrow X{\quoi}.
$$
On the\pageoriginale other hand, if
$$
J'\supset J
$$
is another subset of $I$ satisfying conditions a), b), c) above, one
obtains a canonical morphism
$$
u_{J',J}:X_J\longrightarrow X_{J'}{\quoi},
$$
so that the $X_J$ form an inductive system in $F$, parametrized by the
set (ordered by inclusion) of subsets $J$ of $I$ satisfying the
conditions under consideration. Finally, the morphisms $u_J$ above
define a morphism from this inductive system to (the constant inductive
system defined by) $X$.


\setcounter{subsection}{21} \setcounter{subsubsection}{14}
\subsubsection{}\etiquette[9.21.15]{I.sec9.21.15}
Now let $K$ be a set of subsets $J$ of $I$ satisfying conditions a),
b), c) considered in \hyperref[I.sec9.21.14]{9.21.14}, and suppose
that $K$ is filtered and has union $I$. It is then clear that one has
$$
\varinjlim_{J\in K}X_J\xrightarrow{\sim}X{\quoi},
$$
using \hyperref[I.cor9.21.10.1]{9.21.10.1}, which reduces us to
checking that there is an isomorphism argument by argument.

Take, for example, for $K$ the set of all subsets $J$ of $I$ that are
large relative to $\Pi'_\circ$, stable under $\varphi$, and such that
$\card J\leq\Pi$. Then, by the definition \hyperlink{sga4.exp.I.equation.9.21.6}{(9.21.6)} of
$C^{\Pi}_\alpha$, for every $\alpha\in\ob B$ one has
$X_J(\alpha)\in\ob C^{\Pi}_\alpha$, hence $X_J\in\ob F^{\Pi}$.
Consequently, \hyperref[I.lem9.21.9]{9.21.9} (ii) will be proved if we
establish that $K$ is large relative to $\Pi$ (and hence filtered) and
has union $I$. For this, it will evidently suffice to prove that every
subset $S$ of $I$ such that $\card S\leq\Pi$ is contained in a
$J\in K$. This will indeed follow\pageoriginale from the preliminary
hypothesis stated in \hyperref[I.lem9.21.9]{9.21.9} (ii), and from the


% END INCLUDED SGA4_English_current_buildable_source_20260906/source/01/drafts/SGA4_I_LINES_6964_7051_English.texfrag


% BEGIN INCLUDED SGA4_English_current_buildable_source_20260906/source/01/drafts/SGA4_I_LINES_7052_7130_English.texfrag
\begin{enonce*}{Lemma 9.21.16}\etiquette[9.21.16]{I.lem9.21.16}
Let $I$ be an ordered set, let $\Pi$ be an infinite cardinal, and let
$\varphi:I\rightarrow I$ be a map from $I$ to itself.
\begin{enumerate}
\renewcommand{\theenumi}{\alph{enumi}}
\renewcommand{\labelenumi}{\theenumi)}
\item Suppose $I$ is filtered. Then, for every subset $S$ of $I$ such
that $\card S\leq\Pi$, there exists a filtered subset $J\supset S$ of
$I$ such that $\varphi(J)\subset J$ and $\card(J)\leq\Pi$.

\item Let $\Pi'_\circ$ be a cardinal $<\Pi$, and suppose that $I$ is
large relative to $\Pi$. Then, for every subset $S$ of $I$ such that
$\card S\leq\Pi$, there exists a subset $J\supset S$ of $I$ large
relative to $\Pi'_\circ$, such that $\varphi(J)\subset J$ and
$\card J\leq\Pi$.
\end{enumerate}
\end{enonce*}

Let us prove b), for example (the proof of a) being similar and
simpler). Let $P$ be the set of subsets of $I$ of cardinality
$\leq\Pi$, and let $T\mapsto i_T:P\rightarrow I$ be a map such that,
for $T\in P$, $i_T$ is an upper bound of $T$ in $I$; the existence of
this map merely expresses the hypothesis that $I$ is large relative to
$\Pi$. Let $A$ be a well-ordered set such that $\card A=\Pi$, and such
that, for every $a\in A$, the set of $a'\leq a$ has cardinality
$<\Pi$. It follows in particular, since $\Pi'_\circ<\Pi$, that $A$
\emph{is large relative to} $\Pi'_\circ$. Define, by transfinite
recursion, maps
$$
S\longmapsto S_a:P\longrightarrow P\quad(a\in A){\quoi},
$$
by the following formulas:

$S_\circ=S$, $S_{a+1}=S_a\cup\varphi(S_a)\cup\{i_{S_a}\}$,
$S_a=\bigcup_{a'<a}S_{a'}$ if $a$ is a limit ordinal.

Then put,\pageoriginale for $S\in P$,
$$
J(S)=\bigcup_{a\in A}S_a{\quoi}.
$$
It is then clear that $J(S)\supset S$, that $J(S)$ is stable under
$\varphi$, that $\card J(S)\leq\Pi$ (since $\card A=\Pi$), and finally
that $J(S)$ is large relative to $\Pi'_\circ$ (using the fact that $A$
is large relative to $\Pi'_\circ$).

This proves \hyperref[I.lem9.21.16]{9.21.16} and completes the proof of
\hyperref[I.lem9.21.9]{9.21.9}.

We also point out a variant of \hyperref[I.lem9.21.9]{9.21.9}:

\begin{enonce*}{Lemma 9.21.17}\etiquette[9.21.17]{I.lem9.21.17}
The notation is that of \hyperref[I.lem9.21.9]{9.21.9}. Let $F'$ denote
the full subcategory $\Hhomcart_B(B,E)$ of $F=\Hhom_B(B,E)$ formed by
the cartesian sections of $E$ over $B$. Then:
\begin{enumerate}
\renewcommand{\theenumi}{\roman{enumi}}
\renewcommand{\labelenumi}{(\theenumi)}
\item Every object of $F'$ is accessible. More precisely, let
$\alpha\mapsto X(\alpha)$ be an element of $F'$, and let $d$ be a
cardinal such that $d\geq\Pi'_\circ$, $d\geq c$, and such that, for
every $\alpha\in\ob B$, $X(\alpha)$ is a $d$-accessible object of
$E_\alpha$. Then $X$ is a $d$-accessible object of $F'$.

\item Suppose $\Pi\leq c$ and that the functors
$f^\ast:E_\beta\rightarrow E_\alpha$ are
$\Pi'_\circ$-accessible, or that the categories $E_\alpha$ are stable
under small filtered inductive limits and the functors
$f^\ast:E_\beta\rightarrow E_\alpha$ commute with these inductive
limits. Then every object $X$ of $F'$ is isomorphic to an object of the
form $\varinjlim_I X_i$, where, for every $i\in I$, $X_i$ is an object
of the full subcategory ${F'}^{\Pi}=F'\cap F^{\Pi}$ of $F'$, and where
$I$ is large relative to $\Pi$.
\end{enumerate}
\end{enonce*}

The\pageoriginale proof being entirely analogous to that of
\hyperref[I.lem9.21.9]{9.21.9}, we shall merely indicate the points at
which the latter must be modified. First note:


% END INCLUDED SGA4_English_current_buildable_source_20260906/source/01/drafts/SGA4_I_LINES_7052_7130_English.texfrag


% BEGIN INCLUDED SGA4_English_current_buildable_source_20260906/source/01/drafts/SGA4_I_LINES_7131_7193_English.texfrag
\begin{enonce*}{Lemma 9.21.18}\etiquette[9.21.18]{I.lem9.21.18}
The statement of \hyperref[I.lem9.21.10]{9.21.10} remains valid when
$F=\Hhom_B(B,E)$ is replaced there by $F'=\Hhomcart_B(B,E)$, provided
that the inverse image functors $f^\ast:E_{\beta}\rightarrow
E_{\alpha}$ commute with inductive limits of type $I$.
\end{enonce*}

Thus, under the conditions of \hyperref[I.lem9.21.17]{9.21.17}, if $d$ is a cardinal such that
$d\geq c$ and $d\geq \Pi'_{\circ}$, it follows from the above and
from hypothesis \hyperlink{sga4.exp.I.equation.9.21.3}{(9.21.3 c)}) that, for every
ordered set $I$ large
relative to $d$, inductive limits of type $I$ are representable
in $F'$ and are computed argument by argument, whence the conclusion
\hyperref[I.lem9.21.17]{9.21.17} (i) follows by an immediate argument.

To prove (ii), one must supplement \hyperref[I.lem9.21.13]{9.21.13}:

\begin{enonce*}{Lemma 9.21.19}\etiquette[9.21.19]{I.lem9.21.19}
Under the conditions of \hyperref[I.lem9.21.13]{9.21.13}, suppose that $X\in \ob F'$ \ie
that $X$ is a cartesian section of $E$ over $B$. Then the conclusion can be
strengthened by asserting that there exists a function
$μ$ which assigns to every $i\in I$ and every arrow $f:\alpha\rightarrow
\beta$ of $B$ a morphism $μ(i,f):f^\ast
X(\beta)_i\rightarrow X(\alpha)_{\varphi(i)}$ in $E_{\alpha}$, such
that:

\begin{enumerate}
\renewcommand{\theenumi}{\alph{enumi}}
\renewcommand{\labelenumi}{\theenumi)}
\setcounter{enumi}{3}
\item For every $i\in I$ and every arrow $f:\alpha\rightarrow \beta$
of $B$, the following two diagrams are commutative:

\begin{displaymath}
\xymatrix@L=5pt@R=1.25cm@C=2.2cm{
                       & f^{\ast}X(\beta)_{\varphi^{2}(i)}\\
X(\alpha)_{\varphi(i)}\ar[ru]^{\lambda(\varphi (i),f)} &\\
                       & f^{\ast}X(\beta)_i \ar[lu]^{\mu(i,f)}\ar[uu]
} \qquad\qquad \xymatrix@L=5pt@R=1.35cm@C=2cm{
X(\alpha)_{\varphi^2(i)} & \\
                         & f^{\ast}X(\alpha)_{\varphi(i)} \ar[lu]_{\mu(\varphi(i),f)}\\
x(\alpha)_i \ar[uu]\ar[ru]_{\lambda(i,j)}             &. }
\end{displaymath}
\end{enumerate}
\end{enonce*}

To\pageoriginale prove this, one proceeds as in \hyperref[I.lem9.21.13]{9.21.13}, by
considering the composite morphism $f^\ast(X(\beta)_i)\rightarrow
f^\ast(X(\beta))\xrightarrow{X(f)^{-1}}X(\alpha)$, where (as already
in the formulation of condition d)) we identify notationally
an $f$-morphism $R\rightarrow S$ of $E$ with the corresponding morphism
$R\rightarrow f^\ast(S)$ of $E_{\alpha}$. Since
$X(\alpha)=\varinjlim_j X(\alpha)_j$, and $I$ is large relative to $c$, one
can factor the preceding morphism through one of the $X(\alpha)_j$, where
$j$ \textit{a priori} depends on $i$ \emph{and} on $f$, but can be
chosen independently of $f$, say $\psi(i)$. After enlarging the
function $\varphi$ of \hyperref[I.lem9.21.13]{9.21.13}, one may suppose that $\psi=\varphi$.
Proceeding as for conditions b) and c) of
\hyperref[I.lem9.21.13]{9.21.13}, one sees
that, after further enlarging the map $\varphi$, one may
suppose that the two diagrams of \hyperref[I.lem9.21.19]{9.21.19} d) are commutative.

\setcounter{subsection}{21} \setcounter{subsubsection}{19}

% END INCLUDED SGA4_English_current_buildable_source_20260906/source/01/drafts/SGA4_I_LINES_7131_7193_English.texfrag


% BEGIN INCLUDED SGA4_English_current_buildable_source_20260906/source/01/drafts/SGA4_I_LINES_7194_7240_English.texfrag
\subsubsection{}\etiquette[9.21.20]{I.sec9.21.20}
With the map $\varphi$ chosen as in
\hyperref[I.lem9.21.19]{9.21.19}, let us repeat
the argument of \hyperref[I.sec9.21.14]{9.21.14}, except that we must assume that $J$
satisfies, in addition to the stated conditions a), b), c), the following condition:
\begin{enumerate}
\renewcommand{\theenumi}{\alph{enumi}}
\renewcommand{\labelenumi}{\theenumi)}
\setcounter{enumi}{3}
\item For every arrow $f:\alpha\rightarrow \beta$ in $B$, the
functor $f^\ast:E_{\beta}\rightarrow E_{\alpha}$ commutes with
inductive limits of type $J$.
\end{enumerate}

We\pageoriginale claim that, under this additional condition,
the section $X_J$ of $E$ over $B$ is cartesian, \ie for every
arrow $f:\alpha\rightarrow \beta$ of $B$, the morphism
\begin{equation}
X_J(\alpha)\longrightarrow f^\ast X_J(\beta)\tag{$\ast$}
\end{equation}
is an isomorphism. Indeed, by condition d)
above, the target of the morphism under consideration is identified with
$\varinjlim_{i\in J}f^\ast(X(\beta)_i)$, and the morphism is obtained
by taking $\varinjlim_J$ of the morphism of ind-objects
in $E_{\alpha}$
$$
(X(\alpha)_i)_{i\in J}\longrightarrow (f^\ast)(X(\beta)_i)_{i\in
J}{\quoi},
$$
induced by the morphisms $\lambda(i,f):X(\alpha)_i\rightarrow
f^\ast(X(\beta)_{\varphi(i)})$ for $i\in J$. Now the conditions
set out in \hyperref[I.lem9.21.19]{9.21.19} ensure that
the preceding homomorphism
of ind-objects is in fact an isomorphism, an inverse being obtained
from the homomorphism induced by the system of $\mu(i,f)$, for $i\in J$.
This proves our assertion that $(\ast)$ is an %%
isomorphism, and hence that $X\in \ob F'$.

\subsubsection{}\etiquette[9.21.21]{I.sec9.21.21}
We shall now assume that the functors
$f^\ast:E_{\beta}\rightarrow E_{\alpha}$ are
$\Pi'_{\circ}$-\emph{accessible.} Then the conditions on $J$
considered in \hyperref[I.sec9.21.20]{9.21.20} are satisfied if $J$ is large relative to
$\Pi'_{\circ}$, stable under $\varphi$, and such that $\card J\leq \Pi$,
and we complete the proof of \hyperref[I.lem9.21.17]{9.21.17} just as in
\hyperref[I.sec9.21.15]{9.21.15}.

% END INCLUDED SGA4_English_current_buildable_source_20260906/source/01/drafts/SGA4_I_LINES_7194_7240_English.texfrag


% BEGIN INCLUDED SGA4_English_current_buildable_source_20260906/source/01/drafts/SGA4_I_LINES_7241_7301_English.texfrag
\begin{enonce*}{Theorem 9.22}\etiquette[9.22]{I.them9.22}
Let\pageoriginale $p:E\rightarrow B$ be a fibered functor, where $B$
is a category essentially equivalent to a small
category, and where for every arrow $f:\alpha\rightarrow \beta$ in
$B$ the inverse image functor $f^\ast:E_{\beta}\rightarrow
E_{\alpha}$ is accessible (\hyperref[I.def9.2]{9.2}). Suppose moreover that for every
$\alpha\in \ob E$, the fiber category $E_{\alpha}$ admits a
cardinal filtration (\hyperref[I.def9.12]{9.12}) and that every object of $E_{\alpha}$ is
accessible (\hyperref[I.def9.3]{9.3}). Then each of the categories
$F=\Hhom_B(B,E)$ and
$F'=\Hhomcart_B(B,E)$ has a small full subcategory
generating by strict epimorphisms (\hyperref[I.def7.1]{7.1}), and
every object of either category is accessible.
\end{enonce*}

It is immediate that the statement is not essentially changed when
$B$ is replaced by a full subcategory $B_{\circ}$ such that the
inclusion functor $B_{\circ}\rightarrow B$ is an equivalence,
and $E$ by $E_{\circ}=E\times_B B_{\circ}$. This allows us to
assume that $B$ is a small category.

For every $\alpha\in \ob B$, let $(\Filt^{\Pi}(E))_{\Pi\geq
\Pi_{\alpha}}$ be a cardinal filtration of $E_{\alpha}$. Let
$c_{\circ}\in 𝒰$ be a cardinal satisfying the following
conditions:
\hypertarget{sga4.exp.I.equation.9.22.1}{}\begin{equation}
\begin{cases}
& c_{\circ}\geq {\displaystyle{\mathop{\Sup}_{\alpha\in ob
B}}}\Pi_{\alpha}{\quav} c_{\circ}\geq \card %%
\mathrm{F}{}l\, B;\\
& \mbox{for every arrow $f:\alpha\rightarrow \beta$ of $B$,
$f^\ast:E_{\beta}\rightarrow E_{\alpha}$ is
$c_{\circ}$-accessible;}\\
& \mbox{for every $\alpha\in \ob B$, the objects of
$\Filt^{\Pi_{\alpha}}(E_{\alpha})$ are $c_\circ$-accessible.}
\end{cases}\tag{9.22.1}
\end{equation}
Set
$$
c=2^{c_{\circ}}{\quoi},
$$
so that $c>c_{\circ}$, and, for every $\alpha\in \ob B$, let
$$
C_{\alpha}=\Filt^c(E_{\alpha}){\quoi}.
$$
We\pageoriginale claim that the preliminary conditions of
\hyperref[I.lem9.21.9]{9.21.9} and
\hyperref[I.lem9.21.17]{9.21.17} are satisfied on taking
$\Pi=\Pi_{\circ}=c$. This is clear
for \hyperlink{sga4.exp.I.equation.9.21.2}{(9.21.2)} and \hyperlink{sga4.exp.I.equation.9.21.3}{(9.21.3)}. For \hyperlink{sga4.exp.I.equation.9.21.4}{(9.21.4 a)}),
this follows from \hyperref[I.cor9.17]{9.17},
and \hyperlink{sga4.exp.I.equation.9.21.4}{(9.21.4 b)}) follows from condition \hyperref[I.def9.12]{9.12} c)
for a cardinal
filtration. Let us also note that condition \hyperref[I.def9.12]{9.12} b) on
cardinal filtrations implies that for every $\alpha\in \ob B$, we
have $C^{\Pi}_{\alpha}=C_{\alpha}=\Filt^C(E_{\alpha})$, where
$C^{\Pi}_{\alpha}$ is defined in \hyperlink{sga4.exp.I.equation.9.21.6}{(9.21.6)}. Consequently,
\hyperref[I.them9.22]{9.22} follows
from \hyperref[I.lem9.21.9]{9.21.9} and \hyperref[I.lem9.21.17]{9.21.17}, and more
precisely, we have proved the


% END INCLUDED SGA4_English_current_buildable_source_20260906/source/01/drafts/SGA4_I_LINES_7241_7301_English.texfrag


% BEGIN INCLUDED SGA4_English_current_buildable_source_20260906/source/01/drafts/SGA4_I_LINES_7302_7379_English.texfrag
\begin{enonce*}{Corollary 9.23}\etiquette[9.23]{I.cor9.23}
Under the conditions of \hyperref[I.them9.22]{9.22}, if
$c_\circ\in 𝒰$ is a cardinal satisfying conditions \hyperlink{sga4.exp.I.equation.9.22.1}{(9.22.1)}, and if
$c=2^{c_\circ}$, then the subcategory $F^c$ of $F$ (respectively
${F'}^c$ of $F'$) formed by the $X$ such that
$X(\alpha)\in\ob\Filt^c(E_\alpha)$ for every $\alpha\in\ob B$ is
essentially small and generating by strict epimorphisms; more
precisely, every object $X$ of $F$ (respectively $F'$) is isomorphic to
an object of the form $\varinjlim_I X_i$, where the $X_i$ belong to
$F^c$ (respectively ${F'}^c$), and where $I$ is an ordered set large
relative to $c$.
\end{enonce*}

Under slightly stronger hypotheses in \hyperref[I.them9.22]{9.22}, one
can moreover make \hyperref[I.them9.22]{9.22} and
\hyperref[I.cor9.23]{9.23} considerably more precise:

\begin{enonce*}{Corollary 9.24}\etiquette[9.24]{I.cor9.24}
Under the conditions of \hyperref[I.them9.22]{9.22}, suppose that each
of the categories $E_\alpha$ $(\alpha\in\ob B)$ is stable under small
filtered inductive limits, and that, for every $\Pi\geq\Pi_\alpha$,
$\Filt^{\Pi}(E_\alpha)$ is stable under filtered inductive limits
indexed by filtered ordered sets $I$ such that $\card I\leq\Pi$ (which
slightly strengthens condition \hyperref[I.def9.12]{9.12} {\rm b)} for
cardinal\pageoriginale filtrations). In the case where $F'$ is being
considered, suppose moreover that, for every arrow
$f:\alpha\rightarrow\beta$ of $B$, the functor
$f^\ast:E_\beta\rightarrow E_\alpha$ commutes with small filtered
inductive limits. Finally, let $c_\circ\in 𝒰$ be a cardinal satisfying
conditions \hyperlink{sga4.exp.I.equation.9.22.1}{(9.22.1)}, and, for every cardinal
$\Pi\geq c=2^{c_\circ}$, consider the strictly full subcategory
$F^{\Pi}$ (respectively ${F'}^{\Pi}$) of $F$ (respectively $F'$)
formed by the $X$ such that
$X(\alpha)\in\Filt^{\Pi}(E_\alpha)$ for every $\alpha\in\ob B$. Then
the subcategories under consideration define a cardinal filtration
(\hyperref[I.def9.12]{9.12}) of $F$ (respectively $F'$).
\end{enonce*}

Conditions a), b), c) of \hyperref[I.def9.12]{9.12} must be checked.
Conditions a) and b) follow from the analogous conditions for the given
cardinal filtrations of the $E_\alpha$, and from
\hyperref[I.lem9.21.10]{9.21.10} (respectively
\hyperref[I.lem9.21.18]{9.21.18}, taking account of the second of
conditions \hyperlink{sga4.exp.I.equation.9.22.1}{(9.22.1)}). It remains to prove c), whose first part follows
immediately from \hyperref[I.lem9.21.9]{9.21.9} (respectively
\hyperref[I.lem9.21.17]{9.21.17}), by taking $\Pi'_\circ=0$ and
$\Pi_\circ=c$. It remains to prove that if
$\Pi'\geq\Pi\geq c$ are two cardinals, then for every $X$ in
$F^{\Pi'}$ (respectively ${F'}^{\Pi'}$) one has
$X\simeq\varinjlim_{i\in K}X_i$, with the $X_i$ in $F^{\Pi}$
(respectively ${F'}^{\Pi}$), $K$ large relative to $\Pi$, and
$\card K\leq{\Pi'}^{\Pi}$. For this, return to the proofs of
\hyperref[I.lem9.21.9]{9.21.9} (ii) and
\hyperref[I.lem9.21.17]{9.21.17} (ii). By condition
\hyperref[I.def9.12]{9.12} c) on the cardinal filtration of $E_\alpha$,
one may suppose that each $I_\alpha$ has cardinality $\leq\Pi'$, and
therefore one will have
$$
\card I\leq ({\Pi'}^{\Pi})^{\card \ob B}\leq
({\Pi'}^{\Pi})^c={\Pi'}^{\Pi c}={\Pi'}^{\Pi}{\quoi}.
$$
The index set $K$ used in \hyperref[I.sec9.21.15]{9.21.15} (respectively \hyperref[I.sec9.21.21]{9.21.21}) is the set of
subsets $J$ of $I$ that are filtered (\ie large relative to
$\Pi'_\circ=0$), stable under $\varphi$, and such that
$\card J\leq\Pi$. It is then immediate that
$$
\card K\leq (\card I)^{\Pi}\leq
({\Pi'}^{\Pi})^{\Pi}={\Pi'}^{\Pi^2}={\Pi'}^{\Pi}{\quoi},
$$
which\pageoriginale completes the proof of \hyperref[I.cor9.24]{9.24}.

To conclude, it is appropriate to give a statement freed from the
somewhat technical hypotheses of \hyperref[I.them9.22]{9.22}, replacing
them by the conjunction, for the $E_\alpha$, of the hypotheses that
occur in \hyperref[I.prop9.11]{9.11} (which ensure the accessibility of
the objects of the $E_\alpha$) and in \hyperref[I.prop9.13]{9.13}
(which ensure the existence of a cardinal filtration in $E_\alpha$):


% END INCLUDED SGA4_English_current_buildable_source_20260906/source/01/drafts/SGA4_I_LINES_7302_7379_English.texfrag


% BEGIN INCLUDED SGA4_English_current_buildable_source_20260906/source/01/drafts/SGA4_I_LINES_7380_7423_English.texfrag
\begin{enonce*}{Corollary 9.25}\etiquette[9.25]{I.cor9.25}
Let $p:E\rightarrow B$ be a fibered functor, where $B$ is a category
equivalent to a small category, and where, for every arrow
$f:\alpha\rightarrow\beta$ of $B$, the functor
$f^\ast:E_\beta\rightarrow E_\alpha$ is accessible
(\hyperref[I.def9.2]{9.2}). Suppose moreover that, for every
$\alpha\in\ob B$, the fiber category $E_\alpha$ satisfies the
following conditions:
\begin{enumerate}
\renewcommand{\theenumi}{\alph{enumi}}
\renewcommand{\labelenumi}{\theenumi)}
\item $E_\alpha$ admits a small full subcategory generating by strict
epimorphisms (\hyperref[I.def7.1]{7.1}).

\item $E_\alpha$ is stable under fiber products, kernels of pairs of
arrows, sums of two objects, and cokernels of pairs of arrows, and
finally under small filtered inductive limits
\footnote[1]{This last condition is probably unnecessary; cf. \hyperref[I.rem9.13.2]{9.13.2}.}.

\item The functor $\Ker(u,v)$ on the category of pairs of arrows of
$E_\alpha$ is accessible (for example, it commutes with small filtered
inductive limits).

\item Every strict epimorphism of $E_\alpha$
(\hyperref[I.sec10.2]{10.2}) is a universal strict epimorphism.
\end{enumerate}

Under these conditions, each of the categories $F=\Hhom_B(B,E)$ and
$F'=\Hhomcart_B(B,E)$ admits a small subcategory generating by strict
epimorphisms, and all its objects are accessible
(\hyperref[I.def9.3]{9.3}).
\end{enonce*}
Moreover,\pageoriginale \hyperref[I.cor9.24]{9.24} gives us:

\begin{enonce*}{Corollary 9.26}\etiquette[9.26]{I.cor9.26}
Under the conditions of $\hyperref[I.cor9.25]{9.25}$, and in the case where one considers
$F'=\Hhomcart_B(B,E)$, suppose that the functors
$f^\ast:E_\beta\rightarrow E_\alpha$ commute with small filtered
inductive limits. Then $F$ (respectively $F'$) admits a cardinal
filtration, which can be made explicit by the procedure of
\hyperref[I.cor9.24]{9.24}.
\end{enonce*}


% END INCLUDED SGA4_English_current_buildable_source_20260906/source/01/drafts/SGA4_I_LINES_7380_7423_English.texfrag


% BEGIN INCLUDED SGA4_English_current_buildable_source_20260906/source/01/drafts/SGA4_I_LINES_7424_7512_English.texfrag
\section{Glossary}\etiquette[10]{I.sec10}

For the reader's convenience, we collect here the definitions of some
terms used in the preceding sections. We denote a category by $C$.

\subsection{Cartesian, cocartesian}\etiquette[10.1]{I.sec10.1}

A diagram in $C$
\[
\xymatrix@C=3.25cm@R=1.5cm{
A\ar[r] \ar[d]& B \ar[d]\\
C\ar[r] & D }
\]
is called \emph{cartesian} if it is commutative and the canonical
morphism from $A$ to the fiber product $C\times_D B$ is an isomorphism.
It is called \emph{cocartesian} if the corresponding diagram in the
category opposite to \sisi{$𝒜$}{$C$} is cartesian.

\subsection{Epimorphism, strict epimorphism}\etiquette[10.2]{I.sec10.2}

\etc\pageoriginale \Cf \hyperref[I.sec10.3]{10.3}.

\subsection{Epimorphic family, strictly epimorphic
family}\etiquette[10.3]{I.sec10.3}

\etc: A family
$$
f_i:A_i\longrightarrow B, i\in I{\quoi},
$$
of arrows with the same target in $C$ is called an \emph{epimorphic
family} if, for every object $C$ of $C$, the map induced by the $f_i$,
\hypertarget{sga4.exp.I.equation.10.3.1}{}\begin{equation}
\Hom_{C}(B,C)\longrightarrow \prod_i
\Hom_{C}(A_i,C),\tag{10.3.1}
\end{equation}
is injective.

The family under consideration is called \emph{strictly epimorphic} if
the image of the map \hyperlink{sga4.exp.I.equation.10.3.1}{(10.3.1)} consists of the families $(g_i)$ such
that, for every object $R$ of $C$, every pair of indices $i,j\in I$,
and every pair of arrows $u:R\rightarrow A_i$, $v:R\rightarrow A_j$
with $f_iu=f_jv$, one also has $g_iu=g_jv$. When the fiber products
$A_i\times_BA_j$ are representable for $i,j\in I$, this is equivalent
to saying that the image of \hyperlink{sga4.exp.I.equation.10.3.1}{(10.3.1)} consists of the $(g_i)$ such that,
for every pair of indices $i,j\in I$, one has
$g_i\pr_1=g_j\pr_2$, where $\pr_1,\pr_2$ are the two projections of
$A_i\times_BA_j$. The family $(f_i)_{i\in I}$ is called an
\emph{effective epimorphic family} if the preceding condition is
satisfied, \ie if it is strictly epimorphic and the fiber products
$A_i\times_BA_j$ are representable.

The family $(f_i)_{i\in I}$ is called a \emph{universally epimorphic}
(respectively \emph{universally effective epimorphic}) family if the
morphisms $f_i$ are pullbackable (\hyperref[I.sec10.3]{10.3}), and if,
for every arrow $B'\rightarrow B$, the family of arrows
$f'_i:A'_i=A_i\times_BB'\rightarrow B'$, $i\in I$, deduced from the
family $(f_i)_{i\in I}$ by the base change $B'\rightarrow B$, is
epimorphic (respectively effective epimorphic).

The\pageoriginale notion of a \emph{universally strictly epimorphic
family} will be defined in (II \hyperref[II.def2.5]{2.5}, \hyperref[II.prop2.6]{2.6}). Note that when the morphisms
$f_i$ are pullbackable (for example, if fiber products are representable
in $C$), this notion coincides with that of a universally effective
epimorphic family (use II \hyperref[II.cor2.4]{2.4}). In the general case, the six variants of
the notion of epimorphicity under consideration are related by the
logical implications:
\[
\xymatrix@C=-1cm@R=.5cm{
  &   & \txt<3cm>{universally effective epimorphic} \ar@{=>}[lld]\ar@{=>}[rrd] &     & &\\
\txt<3cm>{effective epimorphic} \ar@{=>}[rdd]       &   &    &     &  \txt<3cm>{universally strictly epimorphic} \ar@{=>}[rd] \ar@{=>}[ld]&\\
 &   &    &   \txt{strictly\\ epimorphic}
         \ar@{=>}[lld]&     & \txt{universally \\epimorphic} \ar@{=>}[llld]\\
         & \txt{epimorphic}&    &     &      &
}
\]
A morphism $f:A\rightarrow B$ of $C$ is called an \emph{epimorphism
(respectively a strict epimorphism, respectively an effective
epimorphism, respectively a universal epimorphism, respectively a
universal effective epimorphism, respectively a universal strict
epimorphism)} if the family of morphisms consisting of the single
element $f$ is epimorphic (respectively~\ldots).


% END INCLUDED SGA4_English_current_buildable_source_20260906/source/01/drafts/SGA4_I_LINES_7424_7512_English.texfrag


% BEGIN INCLUDED SGA4_English_current_buildable_source_20260906/source/01/drafts/SGA4_I_LINES_7513_7636_English.texfrag
\subsection{Monomorphic family, strictly monomorphic
  family}\etiquette[10.4]{I.sec10.4}

\etc A family of arrows $f_i:A_i\rightarrow B$ of $C$ is called
\emph{monomorphic (respectively strictly monomorphic,\ldots)} if,
regarded as a family of arrows of the opposite category $C^{\circ}$,
it is epimorphic (respectively strictly epimorphic,\ldots). An arrow
$f:A\rightarrow B$ of $C$ is called a \emph{monomorphism (respectively
a strict monomorphism,\ldots)} if the family consisting only of $f$
is monomorphic (respectively\pageoriginale strictly
monomorphic,\ldots), \ie if $f$, regarded as an arrow of the opposite
category $C^{\circ}$, is an epimorphism (respectively a strict
epimorphism,\ldots).

\subsection{Monomorphism, strict
  monomorphism}\etiquette[10.5]{I.sec10.5}

\etc \Cf \hyperref[I.sec10.4]{10.4}.

\subsection{Projector, image of a projector, direct factor of an
object}\etiquette[10.6]{I.sec10.6}

Let $f:X\rightarrow X$ be an endomorphism of an object of $C$. The
endomorphism $f$ is called a \emph{projector} if $f^2=f$. The pair
$(f,\id_X)$ then has a representable cokernel $X'$ if and only if it
has a representable kernel $X''$; when this is the case, there exists a
unique isomorphism $u:X'\rightarrow X''$ in $C$ such that $iup=f$,
where $p:X\rightarrow X'$ and $i:X''\rightarrow X$ are the canonical
morphisms.\nde{Representability of the kernel, or of the cokernel, is
also equivalent to the existence of morphisms $i$, $p$, and an
isomorphism~$u$ such that $iup=f$ and $pi=u^{-1}$. In particular, this
property is preserved by every functor.}
One then generally identifies $X'$ with
$X''$ and calls it the \emph{image of the projector} $f$; the
projector $f$ is said to \emph{admit an image} if
$\Ker(f,\id_X)$ is representable, \ie if $\Coker(f,\id_X)$ is
representable. A subobject (respectively a quotient) $Y$ of $X$ is
called a \emph{direct factor} of $X$ if there exists a projector $f$
on $X$ that factors through $Y$ and has $Y$ as its image as a
subobject (respectively as a quotient) of $X$. Sometimes, by abuse of
language, an object of $C$ is called a \emph{direct factor of} $X$ if
it is isomorphic to the image of a projector on $X$.

\subsection{Pullbackable}\etiquette[10.7]{I.sec10.7}

An arrow $f:A\rightarrow B$ of $C$ is called \emph{pullbackable} if,
for every arrow $B'\rightarrow B$ of $C$, the fiber product
$A\times_B B'$ is representable in $C$.

\subsection{Quotient, strict quotient}\etiquette[10.8]{I.sec10.8}

\etc\pageoriginale

Let $A$ be an object of $C$. Two epimorphisms $f:A\rightarrow B$ and
$f':A\rightarrow B'$ with source $A$ are called \emph{equivalent} if
there exists an isomorphism $u:B\xrightarrow{\sim}B'$ such that
$uf=f'$. This defines an equivalence relation on the set of
epimorphisms with source $A$, whose classes are called the
\emph{quotients, or quotient objects,} of $A$. For each quotient of
$A$, one generally chooses an element of this class, say
$f:A\rightarrow B$, and one often speaks (by abuse of language) of
the quotient $B$ of $A$ (or of the quotient $f:A\rightarrow B$ of
$A$). The quotient $B$ is called a \emph{strict quotient}
(respectively an \emph{effective quotient, respectively a universal
quotient,}\ldots) if the morphism $f:A\rightarrow B$ is a strict
epimorphism (respectively an effective epimorphism, respectively a
universal epimorphism,\ldots) (\hyperref[I.sec10.3]{10.3}).

\subsection{Equivalence relations}\etiquette[10.9]{I.sec10.9}

Let $F$ and $G$ be two presheaves of sets on $C$. A diagram
$p_1$, $p_2:F\rightrightarrows G$ is called an \emph{equivalence
relation on} $G$ if, for every object $A$ of $C$, the map
$$
(p_1(A),p_2(A)):F(A)\longrightarrow G(A)\times G(A)
$$
induces a bijection from $F(A)$ onto the graph of an equivalence
relation on the set $G(A)$.

A diagram $p_1$, $p_2:B\rightrightarrows C$ in $C$ is called an
equivalence relation on $C$ if the corresponding diagram of
presheaves is an equivalence relation. When finite products and fiber
products are representable in $C$, a functor $\Phi:C\rightarrow C'$
commuting with finite products and fiber products carries equivalence
relations on $C$ to equivalence relations on $\Phi(C)$.

Let\pageoriginale $\pi:C\rightarrow D$ be a morphism of $C$ such that
the fiber product $C\times_D C$ is representable. Then the diagram
$\pr_1$, $\pr_2:C\times_D C\rightrightarrows C$ is an equivalence
relation on~$C$.

\subsection{Effective equivalence relation, universally effective
equivalence relation}\etiquette[10.10]{I.sec10.10}

An equivalence relation $p_1$, $p_2:R\rightrightarrows A$ on an object
$A$ of $C$ is called an \emph{effective equivalence relation} if
there exists a morphism $\pi:A\rightarrow B$ such that the square
\[
\xymatrix@C=3.75cm@R=1.5cm{
R\ar[r]^{p_2} \ar[d]_{p_1}& A \ar[d]^{\pi}\\
A\ar[r] & B }
\]
is cartesian and cocartesian. (The morphism $\pi$ is the cokernel of
the pair $(p_1,p_2)$ and is therefore determined up to unique
isomorphism.) The morphism $\pi$ is then an effective epimorphism; if
it is a universal effective epimorphism (\hyperref[I.sec10.3]{10.3}),
the \emph{equivalence relation} is said to be \emph{universally
effective.}

\subsection{Subobject, strict subobject}\etiquette[10.11]{I.sec10.11}

\etc These are the notions dual to those of quotient and strict
quotient \etc (\hyperref[I.sec10.8]{10.8}).

\pdfbookmark[1]{Bibliography}{sga4.exp.I.bibliography.main}
\begingroup
\makeatletter
\def\Hy@writebookmark#1#2#3#4#5{}
\makeatother
\begin{thebibliography}{1}
\bibitem[1]{sga4.exp.I.bib.I.1} \sisi{B. Mitchell, Theory of Categories, Academic Press
(1965).}{\textsc{B. Mitchell}, ``Theory of Categories,'' Academic
Press (1965).}
\end{thebibliography}
\endgroup


\newpage

\footnotetext[1]{We reproduce here, with his permission, secret papers
of N. \textsc{Bourbaki}. The references in this text are to his
learned work.}

% END INCLUDED SGA4_English_current_buildable_source_20260906/source/01/drafts/SGA4_I_LINES_7513_7636_English.texfrag


% BEGIN INCLUDED SGA4_English_current_buildable_source_20260906/source/01/drafts/SGA4_I_LINES_7637_7820_English.texfrag
\pdfbookmark[1]{II. Appendix: Universes (by N. Bourbaki)}{sga4.exp.I.appendix.II}
\begingroup
\makeatletter
\def\Hy@writebookmark#1#2#3#4#5{}
\makeatother
\section*{II. Appendix: Universes (by
  N. \textsc{Bourbaki}($^*$))}\pageoriginale
\endgroup
\markboth{II. APPENDIX: UNIVERSES}{II. APPENDIX: UNIVERSES}

\hypersetup{bookmarksdepth=2}
\makeatletter
\let\sgaFourSavedToclevelSection\toclevel@section
\def\toclevel@section{2}
\makeatother

\setcounter{section}{0}
\section{Definition and first properties of universes}%
\hypertarget{sga4.exp.I.appendix.II.sec.1}{}

\hypertarget{sga4.exp.I.appendix.II.sec.1.definition.1}{}
\begin{enonce*}{Definition 1}
A set $U$ is called a universe if it satisfies the following conditions:
\begin{itemize}
\hypertarget{sga4.exp.I.appendix.II.sec.1.definition.1.condition.U.I}{}
\item[(U.I)] if $x\in U$ and $y\in x$, then $y\in U$;

\hypertarget{sga4.exp.I.appendix.II.sec.1.definition.1.condition.U.II}{}
\item[(U.II)] if $x$, $y\in U$, then $\{x,y\}\in U$;

\hypertarget{sga4.exp.I.appendix.II.sec.1.definition.1.condition.U.III}{}
\item[(U.III)] if $x\in U$, then $\mathcal{P}(x)\in U$;

\hypertarget{sga4.exp.I.appendix.II.sec.1.definition.1.condition.U.IV}{}
\item[(U.IV)] if $(x_{\alpha})_{\alpha\in I}$ is a family
of elements of $U$, and if $I\in U$, then the union
$\bigcup_{\alpha\in I}x_{\alpha}$ belongs to $U$;

\hypertarget{sga4.exp.I.appendix.II.sec.1.definition.1.condition.U.question}{}
\item[(U.?)] if $x$, $y\in U$, then the ordered pair $(x,y)\in U$.
\end{itemize}
\end{enonce*}
\noindent N.B. Since it has, I believe, been decided that in future
editions the ordered pair will be defined à la \sisi{Kuratowaki}{Kuratowski}
by $(x,y)=\{\{x,y\}, \{x\}\}$, condition (U.?) is unnecessary, since
it follows from \hyperlink{sga4.exp.I.appendix.II.sec.1.definition.1.condition.U.II}{(U.II)}.

\hypertarget{sga4.exp.I.appendix.II.sec.1.examples}{}
\begin{enonce*}[remark]{Examples}
\begin{enumerate}
\renewcommand{\labelenumi}{\theenumi)}
\hypertarget{sga4.exp.I.appendix.II.sec.1.example.1}{}
\item The \emph{empty} set is a universe, denoted by $U_\circ$.

\hypertarget{sga4.exp.I.appendix.II.sec.1.example.2}{}
\item Consider the nonempty finite words formed with the four
symbols $''\{'',''\}'','',''$ and $''∅''$ (\cf Alg. I).
Define, by induction\pageoriginale on the length $n$ of such a word,
the notion of a \emph{meaningful} word:
\begin{enumerate}
\renewcommand{\theenumii}{\alph{enumii}}
\renewcommand{\labelenumii}{(\theenumii)}
\item For $n=1$, only the word $∅$ is meaningful;

\item for a word $A$ of length $n$ to be meaningful, it is necessary and
sufficient that there exist $p$ distinct meaningful words
$A_1,\ldots,A_p$ $(p\geq 1)$ of lengths $<n$ such that
$$
A=\{A_1,A_2,\ldots,A_p\}{\quoi}.
$$
\end{enumerate}
\end{enumerate}
\end{enonce*}
\noindent For example, $∅$, $\{∅\}$,
$\{\{∅\},\{\{∅\}\},\{∅,\{∅\}\}\}$
are meaningful words. It is clear, by induction on the
length, that every meaningful word denotes a term of Set
Theory; for example, if $A=\{A_1,A_2,\ldots,A_p\}$, A
denotes the set whose elements are $A_1$, $A_2,\ldots,$ and
$A_p$; these terms are plainly \emph{finite sets.} Let
$U_1$ be the set of sets thus obtained. It is readily checked
that $U_1$ satisfies conditions \hyperlink{sga4.exp.I.appendix.II.sec.1.definition.1.condition.U.I}{(U.I)}, \hyperlink{sga4.exp.I.appendix.II.sec.1.definition.1.condition.U.II}{(U.II)}, \hyperlink{sga4.exp.I.appendix.II.sec.1.definition.1.condition.U.III}{(U.III)}, and \hyperlink{sga4.exp.I.appendix.II.sec.1.definition.1.condition.U.IV}{(U.IV)}
of \hyperlink{sga4.exp.I.appendix.II.sec.1.definition.1}{Definition 1} (but not that idiotic condition (U.?), which is no
problem if one is willing to decanulate the ordered pair). Thus $U_1$ is a
\emph{universe}. Note that the elements of $U_1$ are finite, and
that $U_1$ is countable.

In the statements that follow, $U$ \emph{denotes a universe.}

\hypertarget{sga4.exp.I.appendix.II.sec.1.proposition.1}{}
\begin{enonce*}{Proposition 1}
If $X\in U$ and if $y\subset x$, then $y\in U$.
\end{enonce*}
Indeed, $\mathcal{P}(x)\in U$ by \hyperlink{sga4.exp.I.appendix.II.sec.1.definition.1.condition.U.III}{(U.III)}, whence $y\in
\mathcal{P}(x)$ and $y\in U$ by \hyperlink{sga4.exp.I.appendix.II.sec.1.definition.1.condition.U.I}{(U.I)}.

\hypertarget{sga4.exp.I.appendix.II.sec.1.corollary.to.proposition.1}{}
\begin{enonce*}{Corollary}
If $x\in U$, every quotient set $y$ of $x$ is an element of $U$.
\end{enonce*}
Indeed, $y$ is a subset of $\mathcal{P}(x)$. Hence $y\in U$ by
\hyperlink{sga4.exp.I.appendix.II.sec.1.definition.1.condition.U.III}{(U.III)} and \hyperlink{sga4.exp.I.appendix.II.sec.1.proposition.1}{Proposition 1}.

\hypertarget{sga4.exp.I.appendix.II.sec.1.proposition.2}{}
\begin{enonce*}{Proposition 2}
If\pageoriginale $x\in U$, then $\{x\}\in U$.
\end{enonce*}
This follows from \hyperlink{sga4.exp.I.appendix.II.sec.1.definition.1.condition.U.II}{(U.II)}, applied with $y=x$.

\hypertarget{sga4.exp.I.appendix.II.sec.1.proposition.3}{}
\begin{enonce*}{Proposition 3}
Every ordered pair, ordered triple, and ordered quadruple of elements of $U$ is an
element of $U$.
\end{enonce*}
This is true for ordered pairs by the N.B. or \hyperlink{sga4.exp.I.appendix.II.sec.1.definition.1.condition.U.question}{(U.?)}. The case of ordered
triples follows, since $(x,y,z)=((x,y),z)$, and then that of ordered
quadruples, since $(x,y,z,t)=((x,y,z),t)$.

\hypertarget{sga4.exp.I.appendix.II.sec.1.proposition.4}{}
\begin{enonce*}{Proposition 4}
If $X$, $Y\in U$, then $X\times Y\in U$.
\end{enonce*}
Indeed, for $x\in X$, $\{x\}\times Y$ is the union of the
family $\{(x,y)\}_{y\in Y}$; since $\{(x,y)\}\in U$ by
\hyperlink{sga4.exp.I.appendix.II.sec.1.definition.1.condition.U.I}{(U.I)}, \hyperlink{sga4.exp.I.appendix.II.sec.1.definition.1.condition.U.II}{(U.II)}, and \hyperlink{sga4.exp.I.appendix.II.sec.1.proposition.3}{Proposition 3}, we have $\{x\}\times Y\in U$ by
(U.IV). Finally, $X\times Y$ is the union of the family $(\{x\}\times
Y)_{x\in X}$; it is therefore an element of $U$, again by \hyperlink{sga4.exp.I.appendix.II.sec.1.definition.1.condition.U.IV}{(U.IV)}.

\hypertarget{sga4.exp.I.appendix.II.sec.1.corollary.to.proposition.4.1}{}
\begin{enonce*}{Corollary 1}
If $X$, $Y$, $Z,\ldots$ are elements of $U$, all the sets
of the scale constructed on $X$, $Y$, $Z,\ldots$ are elements
of $U$ {\rm(\cf Chapter IV, \S)}.
\end{enonce*}

This indeed follows by successive applications of \hyperlink{sga4.exp.I.appendix.II.sec.1.proposition.4}{Proposition 4} and
of (U.III).

\hypertarget{sga4.exp.I.appendix.II.sec.1.corollary.to.proposition.4.2}{}
\begin{enonce*}{Corollary 2}
If $(x_{\alpha})_{\alpha \in I}$ is a family of elements of $U$
and if $I\in U$, the sum set $\Sigma_{\alpha \in I} x_{\alpha}$
is an element of $U$.
\end{enonce*}
Indeed, this sum set is a subset of the product
$(\bigcup_{\alpha\in I} x_{\alpha})\times I$ (Chapter II), a product
both of whose factors are elements of $U$ (by (U.IV) and
the hypothesis). We then apply \hyperlink{sga4.exp.I.appendix.II.sec.1.proposition.4}{Propositions 4} and \hyperlink{sga4.exp.I.appendix.II.sec.1.proposition.1}{1}.

\hypertarget{sga4.exp.I.appendix.II.sec.1.proposition.5}{}
\begin{enonce*}{Proposition 5}
If\pageoriginale $X$ and $Y$ are elements of $U$, every
correspondence between $X$ and $Y$ (in particular, every map from
$X$ to $Y$) is an element of $U$.
\end{enonce*}
Indeed, such a correspondence $C$ is an ordered triple $(X,Y,\Gamma)$,
where $\Gamma$ is a subset of $X\times Y$ (the graph of $C$) (Chapter
II, \S). We have $\Gamma\in U$ by \hyperlink{sga4.exp.I.appendix.II.sec.1.proposition.4}{Propositions 4} and 1. Hence $C\in U$
by \hyperlink{sga4.exp.I.appendix.II.sec.1.proposition.3}{Proposition 3}.

\hypertarget{sga4.exp.I.appendix.II.sec.1.proposition.6}{}
\begin{enonce*}{Proposition 6}
If $X$, $Y\in U$, every set $Z$ of correspondences between $X$ and
$Y$ (in particular, of maps from $X$ to $Y$) is an element of
$U$.
\end{enonce*}
Indeed, $Z$ is a subset of $\{X\}\times \{Y\}\times
\mathcal{P}(X\times Y)$. Now this product is an element of $U$
by
\hyperlink{sga4.exp.I.appendix.II.sec.1.proposition.2}{Proposition 2} and the \hyperlink{sga4.exp.I.appendix.II.sec.1.corollary.to.proposition.4.1}{corollary to Proposition 4}. Therefore $Z\in
U$ by \hyperlink{sga4.exp.I.appendix.II.sec.1.proposition.1}{Proposition 1}.

\hypertarget{sga4.exp.I.appendix.II.sec.1.corollary.to.proposition.6}{}
\begin{enonce*}{Corollary}
If $(x_{\alpha})_{\alpha \in I}$ is a family of elements of $U$,
and if $I\in U$, then $\Pi_{\alpha\in I}x_{\alpha} \in U$.
\end{enonce*}
Indeed, this product is a set of maps from $I$ to
$\bigcup_{\alpha\in I}x_{\alpha}$, and this union is an element of
$U$ by \hyperlink{sga4.exp.I.appendix.II.sec.1.definition.1.condition.U.IV}{(U.IV)}.

\hypertarget{sga4.exp.I.appendix.II.sec.1.proposition.7}{}
\begin{enonce*}{Proposition 7}
If $X$ is a subset of $U$ whose cardinality is at most that of an
element of $U$, then $X$ is an element of $U$.
\end{enonce*}
Let $I$ be an element of $U$ such that $\card(X)\leq \card(I)$. There is a
surjection $i\mapsto x_i$ from a subset $I'$ of $I$ onto $X$. Then
$X$ is the union of the family $(\{x_i\})_{i\in I'}$; now this
union is an element of $U$ by \hyperlink{sga4.exp.I.appendix.II.sec.1.proposition.1}{Proposition 1}, \hyperlink{sga4.exp.I.appendix.II.sec.1.proposition.2}{Proposition 2}, and \hyperlink{sga4.exp.I.appendix.II.sec.1.definition.1.condition.U.IV}{(U.IV)}.

\hypertarget{sga4.exp.I.appendix.II.sec.1.corollary.to.proposition.7}{}
\begin{enonce*}{Corollary}
If\pageoriginale $U$ is nonempty, every finite subset of $U$ is an
element of $U$, and $U$ has elements of arbitrarily large finite cardinality.
\end{enonce*}
On the other hand, if $U=∅$, then $∅$ is a finite subset
of $∅$ but not an element of $∅$.

Indeed, if $U$ is nonempty, \hyperlink{sga4.exp.I.appendix.II.sec.1.proposition.2}{Proposition 2} shows that a one-element
set $x_{\circ}$ belongs to $U$. By induction on $n$,
put $x_{n+1}=\mathcal{P}(x_n)$. We %%
have $x_n\in U$ by \hyperlink{sga4.exp.I.appendix.II.sec.1.definition.1.condition.U.III}{(U.III)}, and $\card(x_{n+1})=2^{\card(x_n)}$, so
that $U$ has elements of arbitrarily large finite cardinality.

\begin{enonce*}[remark]{Remark}
It follows from \hyperlink{sga4.exp.I.appendix.II.sec.1.proposition.2}{Proposition 2} and the \hyperlink{sga4.exp.I.appendix.II.sec.1.corollary.to.proposition.7}{corollary to Proposition 7} that every
nonempty universe $U$ contains the universe $U_1$ of \hyperlink{sga4.exp.I.appendix.II.sec.1.example.2}{Example 2}; indeed,
$∅\in U$ by \hyperlink{sga4.exp.I.appendix.II.sec.1.proposition.1}{Proposition 1}. Thus $U_1$ is
the intersection of all nonempty universes. More generally:
\end{enonce*}

\hypertarget{sga4.exp.I.appendix.II.sec.1.proposition.8}{}
\begin{enonce*}{Proposition 8}
If $(U_{\lambda})_{\lambda\in L}$ is a nonempty family
of universes, then $U=\bigcap_{\lambda\in L}U_{\lambda}$ is a
universe.
\end{enonce*}
This follows at once from \hyperlink{sga4.exp.I.appendix.II.sec.1.definition.1}{Definition 1}.

% END INCLUDED SGA4_English_current_buildable_source_20260906/source/01/drafts/SGA4_I_LINES_7637_7820_English.texfrag


% BEGIN INCLUDED SGA4_English_current_buildable_source_20260906/source/01/drafts/SGA4_I_LINES_7821_7923_English.texfrag
\section{Universes and species of structures}%
\hypertarget{sga4.exp.I.appendix.II.sec.2}{}

Let $\mathcal{E}$ be a species of structure; suppose, to fix ideas
and simplify the exposition, that every structure of species
$\mathcal{E}$ is defined on an underlying set. Let $(X,S)$ be a
structure of species $\mathcal{E}$ (where $X$ is the underlying set
and $S$ the structure), and let $U$ be a \emph{universe.} If $X\in U$,
then all the constituent objects of the structure $S$ on $X$ are
elements of $U$ (by \hyperlink{sga4.exp.I.appendix.II.sec.1.corollary.to.proposition.4.1}{Cor. 1 of Prop. 4}, \hyperlink{sga4.exp.I.appendix.II.sec.1}{${\rm n^{\circ}}$ 1}, and by
\hyperlink{sga4.exp.I.appendix.II.sec.1.definition.1.condition.U.I}{(U.I)}), so the structure $(X,S)$ is an element of $U$.

Suppose\pageoriginale now that $\mathcal{E}$ is a species of structure
with \emph{morphisms}. If $X$ and $X'$ are elements of $U$ endowed
with structures of species $\mathcal{E}$, then the set of morphisms
from $X$ to $X'$ is again an element of $U$ (\hyperlink{sga4.exp.I.appendix.II.sec.1}{${\rm n^{\circ}}$ 1},
\hyperlink{sga4.exp.I.appendix.II.sec.1.proposition.6}{Prop. 6}).

Now consider the \emph{category} $(U\text{-}\mathcal{E})$ defined in
\hyperlink{sga4.exp.I.appendix.II.sec.1}{\S 1}, ${\rm n^{\circ}}$ 2, Ex. c). Since the objects and arrows of
$(U\text{-}\mathcal{E})$ are elements of $U$,
$(U\text{-}\mathcal{E})$ is a \emph{pair of subsets} of $U$, endowed
with a quadruple of maps; in general, it is not an element of $U$.
The notation $X\in (U\text{-}\mathcal{E})$ will mean that $X$ is an
element of $U$ endowed with a structure of species $\mathcal{E}$.

The category $(U\text{-}\mathcal{E})$ is \emph{stable} under many
operations:

\begin{enumerate}
\renewcommand{\theenumi}{\alph{enumi}}
\renewcommand{\labelenumi}{\theenumi)}
\item If $X\in(U\text{-}\mathcal{E})$, and if $X'$ is a subset of $X$ that
admits an \emph{induced structure}, then $X'\in
(U\text{-}\mathcal{E})$. This follows from \hyperlink{sga4.exp.I.appendix.II.sec.1.proposition.1}{Prop. 1}, \hyperlink{sga4.exp.I.appendix.II.sec.1}{${\rm
n^{\circ}}$ 1}.

\item If $X\in (U\text{-}\mathcal{E})$, and if $X''$ is a quotient set of
$X$ that admits a \emph{quotient structure}, then $X''\in
(U\text{-}\mathcal{E})$ (\hyperlink{sga4.exp.I.appendix.II.sec.1.corollary.to.proposition.1}{Cor. of Prop. 1}, \hyperlink{sga4.exp.I.appendix.II.sec.1}{${\rm n^{\circ}}$ 1}).

\item If $X$, $Y\in (U\text{-}\mathcal{E})$, and if $X\times Y$ admits a
\emph{product structure}, then $X\times Y\in
(U\text{-}\mathcal{E})$ (\hyperlink{sga4.exp.I.appendix.II.sec.1.proposition.4}{Prop. 4}, \hyperlink{sga4.exp.I.appendix.II.sec.1}{${\rm n^{\circ}}$ 1}). More
generally, if $(X_i)_{i\in I}$ is a family of elements of
$(U\text{-}\mathcal{E})$, \emph{if} $I\in U$, and if $\Pi_{i\in
I}X_i$ admits a \emph{product structure}, then $\Pi_{i\in I}X_i
\in (U\text{-}\mathcal{E})$ (\hyperlink{sga4.exp.I.appendix.II.sec.1.corollary.to.proposition.6}{Cor. of Prop. 6}). The analogous
assertions hold for \emph{sum structures} (\hyperlink{sga4.exp.I.appendix.II.sec.1.corollary.to.proposition.4.2}{Cor. 2 of Prop. 4}).

\item Let $I$ be a preordered set, and let $(X_i, f_{ij})_{i,j\in
I}$ be a \emph{projective system} (\resp \emph{inductive system}) of
sets endowed with structures of species $\mathcal{E}$ and morphisms;
let $L$ be the limit of this system, in the sense of Chap. III. If
$X_i\in U$ for every $i$, if $I\in U$, and if $L$ admits a
projective-limit structure (\resp an inductive-limit structure), then
$L\in (U\text{-}\mathcal{E})$: indeed, $L$ is a\pageoriginale subset
of $\Pi_{i\in I}X_i$ (\resp a quotient set of $\Sigma_{i\in I}X_i$),
and is therefore an element of $U$ by \hyperlink{sga4.exp.I.appendix.II.sec.1}{${\rm n^{\circ}}$ 1}.
\end{enumerate}

Here are a few more particular examples:
\begin{itemize}
\hypertarget{sga4.exp.I.appendix.II.sec.2.example.star.1}{}
\item[* 1)] Let $X$, $Y\in (U\text{-}\Top)$ be two \emph{locally
compact} spaces. Then the set $\mathscr{C}(X,Y)$ of continuous maps
from $X$ to $Y$, endowed with the topology of compact convergence
(General Topology $X$), is an element of $(U\text{-}\Top)$. Indeed,
$\mathscr{C}(X,Y)\in U$ by \hyperlink{sga4.exp.I.appendix.II.sec.1.proposition.6}{Prop. 6} of \hyperlink{sga4.exp.I.appendix.II.sec.1}{${\rm
n^{\circ}}$ 1}.${}_{\mbox{*}}$

\hypertarget{sga4.exp.I.appendix.II.sec.2.example.star.2}{}
\item[* 2)] Let $X$, $Y\in (U\text{-}\Ab)$. Then the group
$\Hom_{\ZZ}(X,Y)$ is an element of $(U\text{-}\Ab)$ by \hyperlink{sga4.exp.I.appendix.II.sec.1.proposition.6}{Prop. 6} of
\hyperlink{sga4.exp.I.appendix.II.sec.1}{${\rm n^{\circ}}$ 1}. If, in addition, $\ZZ\in U$, the tensor product
$X\otimes_{\ZZ}Y$ is an element of $(U\text{-}\Ab)$; indeed, this
tensor product is a quotient of $\ZZ^{X\times Y}$, which is a subset
of $\ZZ^{X\times Y}$, which in turn is an element of $U$ by \hyperlink{sga4.exp.I.appendix.II.sec.1.proposition.4}{Props. 4}
and \hyperlink{sga4.exp.I.appendix.II.sec.1.proposition.6}{6} of \hyperlink{sga4.exp.I.appendix.II.sec.1}{${\rm n^{\circ}}$ 1}.${}_{\mbox{*}}$

\hypertarget{sga4.exp.I.appendix.II.sec.2.example.star.3}{}
\item[* 3)] Let $X\in (U\text{-}\Unif)$ be a uniform space. Then its
\emph{completion} $\widehat{X}$ is an element of $(U\text{-}\Unif)$.
Indeed, $\widehat{X}$ is a set of equivalence classes of Cauchy
filters on $X$; now a filter on $X$ is an element of
$\mathcal{P}(\mathcal{P}(X))$, so an equivalence class of filters is
an element of $\mathcal{P}(\mathcal{P}(\mathcal{P}(X)))$, whence
$\widehat{X}$ is an element of
$\mathcal{P}(\mathcal{P}(\mathcal{P}(\mathcal{P}(X))))$, and hence an
element of $U$ by \hyperlink{sga4.exp.I.appendix.II.sec.1.definition.1.condition.U.III}{(U.III)} and \hyperlink{sga4.exp.I.appendix.II.sec.1.definition.1.condition.U.I}{(U.I)}.${}_{\mbox{*}}$
\end{itemize}

N.B. Moral: Bourbaki will have to take care to ``canonicalize'' its
constructions properly. For example, in constructions in which an
element $\infty$ is adjoined (Alexandroff compactification, projective
field), it will be advantageous to take $\infty =∅$\pageoriginale
(since $∅$ is an element of every nonempty universe) and to form the
disjoint union of $\{∅\}$ and the given set. Of course, one should
also ``canonicalize'' the two-element set used to construct disjoint
unions; thus $\{∅, \{∅\}\}$ seems to me a good candidate, since it
belongs to every nonempty universe.

% END INCLUDED SGA4_English_current_buildable_source_20260906/source/01/drafts/SGA4_I_LINES_7821_7923_English.texfrag


% BEGIN INCLUDED SGA4_English_current_buildable_source_20260906/source/01/drafts/SGA4_I_LINES_7924_7993_English.texfrag
\section{Universes and categories}%
\hypertarget{sga4.exp.I.appendix.II.sec.3}{}

Let $U$ be a universe and $C$ a category. We write $C\in U$ if
$\Ob(C)\in U$ and $%%
\mathrm{F}{}l\,(C)\in U$. This notation is justified by the fact that
$C$ is the sextuple consisting of $\Ob(C)$, $\mathrm{F}{}l\,(C)$, and
the four structure maps; thus, if $\Ob(C)\in U$ and
$\mathrm{F}{}l\,(C)\in U$, these four maps are elements of $U$
(\hyperlink{sga4.exp.I.appendix.II.sec.1}{${\rm n^{\circ}}$ 1}, \hyperlink{sga4.exp.I.appendix.II.sec.1.corollary.to.proposition.4.1}{Corollary 1} of \hyperlink{sga4.exp.I.appendix.II.sec.1.proposition.4}{Proposition 4}), and hence so is
the sextuple $C$ (\hyperlink{sga4.exp.I.appendix.II.sec.1}{${\rm n^{\circ}}$ 1}, \hyperlink{sga4.exp.I.appendix.II.sec.1.proposition.3}{Proposition 3}, \textit{cum
grano salis}). This is moreover a special case of \hyperlink{sga4.exp.I.appendix.II.sec.2}{${\rm n^{\circ}}$ 2},
if one considers the species of structure ``$\cat$'' of a category.
With the notation of \hyperlink{sga4.exp.I.appendix.II.sec.2}{${\rm n^{\circ}}$ 2}, the relation $C\in U$ is
also written $C\in(U\text{-}\cat)$.

One should note that a category such as $(U\text{-}\Ens)$,
$(U\text{-}\Ord)$, $(U\text{-}\Top)$, or $(U\text{-}\Gr)$ is
\emph{not}, in general, an element of $U$.

Let $C,D$ be two categories and $U$ a universe such that $C,D\in U$.
If $C'$ is a \emph{subcategory} of $C$, then $C'\in U$; the
\emph{product category} $C\times D$, the sum category of $C$ and $D$,
and the opposite categories $C^\circ$ and $D^\circ$ are also elements
of $U$ (cf. \hyperlink{sga4.exp.I.appendix.II.sec.2}{${\rm n^{\circ}}$ 2}). The \emph{functor category}
$E=\Hhom(C,D)$ is likewise an element of $U$: indeed, $\Ob(E)$ is a set
of pairs of maps $\Ob(C)\rightarrow\Ob(D)$,
$F\ell\,(C)\rightarrow F\ell\,(D)$, and hence $\Ob(E)\in U$ by
\hyperlink{sga4.exp.I.appendix.II.sec.1}{${\rm n^{\circ}}$ 1};\pageoriginale as for $F\ell\,(E)$, it is a set
of natural transformations, that is, of maps
$\Ob(C)\rightarrow F\ell\,(D)$, whence $F\ell\,(E)\in U$ by
\hyperlink{sga4.exp.I.appendix.II.sec.1.proposition.6}{Proposition 6} of \hyperlink{sga4.exp.I.appendix.II.sec.1}{${\rm n^{\circ}}$ 1}.

\hypertarget{sga4.exp.I.appendix.II.sec.3.proposition.9}{}
\begin{enonce*}{PROPOSITION 9}
Let $\mathcal{E}$ be a species of structure with morphisms, and let
$U,U'$ be two universes such that $U'\subset U$. Then
$(U'\text{-}\mathcal{E})$ is a full subcategory of
$(U\text{-}\mathcal{E})$.
\end{enonce*}

Indeed, if $x$ and $y$ are objects of $(U'\text{-}\mathcal{E})$, they
are by definition objects of $(U\text{-}\mathcal{E})$. As for
$\Hom(x,y)$, it is the same set in both categories (cf. \hyperlink{sga4.exp.I.appendix.II.sec.1}{\S 1},
${\rm n^{\circ}}$ 2, Example c)).

N.B. One should note that the hypothesis that $U$ and $U'$ are
universes %%
is unnecessary, so that \hyperlink{sga4.exp.I.appendix.II.sec.3.proposition.9}{Proposition 9} would be more
appropriately placed in \hyperlink{sga4.exp.I.appendix.II.sec.1}{\S 1}, ${\rm n^{\circ}}$ 4. But the Tribe asked
that it be put here.

\begin{enonce*}[remark]{Remark}
It may happen that the axioms of the species of structure $\mathcal{E}$
imply that the \emph{cardinals} of sets endowed with structures of
species $\mathcal{E}$ are \emph{bounded} by a fixed cardinal, say
$\sheaf{\sheaf{c}}$ (* for example, the species of structure of finite
group, finitely generated group, finitely generated module over a fixed
ring $A$, or finitely generated algebra over $A$ ${}_{\mbox{*}}$).
Suppose then that there exists an element $z$ of $U'$ such that
$\sheaf{\sheaf{c}}\leq\card(z)$. Then, for every set $x$ endowed with
a structure of species $\mathcal{E}$, there exists an element $x'$ of
$U'$ equipotent to $x$, for example a subset of $z$; endow $x'$ with
the structure deduced from that of $x$ by transport of structure; this
gives an element of $(U'\text{-}\mathcal{E})$. It follows that
\emph{the inclusion functor from $(U'\text{-}\mathcal{E})$ to
$(U\text{-}\mathcal{E})$ is then essentially surjective {\rm
(\hyperlink{sga4.exp.I.appendix.II.sec.4}{\S 4}, ${\rm n^{\circ}}$ 2, Definition 2)} and is therefore an
equivalence of categories} (\hyperlink{sga4.exp.I.appendix.II.sec.4}{\S 4}, ${\rm n^{\circ}}$ 3, Theorem 1).
\end{enonce*}

% END INCLUDED SGA4_English_current_buildable_source_20260906/source/01/drafts/SGA4_I_LINES_7924_7993_English.texfrag


% BEGIN INCLUDED SGA4_English_current_buildable_source_20260906/source/01/drafts/SGA4_I_LINES_7994_8025_English.texfrag
\section{The universe axiom}%
\hypertarget{sga4.exp.I.appendix.II.sec.4}{}\pageoriginale

The very pleasant stability results of \hyperlink{sga4.exp.I.appendix.II.sec.2}{${\rm n^{\circ}}$ 2} are of
interest only if they can be applied to something other than the two
small universes $U_\circ$ and $U_1$ described in \hyperlink{sga4.exp.I.appendix.II.sec.1}{${\rm n^{\circ}}$ 1}.
We shall therefore add the following axiom to the axioms of Set Theory:

\hypertarget{sga4.exp.I.appendix.II.sec.4.axiom.A.6}{}
(A.6) \emph{For every set $x$, there exists a universe $U$ such that
$x\in U$.}

This axiom implies that, if $(x_\alpha)_{\alpha\in I}$ is a
\emph{family} of sets, there exists a universe $U$ such that
$x_\alpha\in U$ for every $\alpha\in I$: indeed, it suffices to apply
\hyperlink{sga4.exp.I.appendix.II.sec.4.axiom.A.6}{(A.6)} to $x=\bigcup_{\alpha\in I}x_\alpha$ and then apply \hyperlink{sga4.exp.I.appendix.II.sec.1.proposition.1}{Proposition
1} of \hyperlink{sga4.exp.I.appendix.II.sec.1}{${\rm n^{\circ}}$ 1}.

In particular, given a \emph{category} $C$, there exists a universe
$U$ such that $C\in U$ in the sense of \hyperlink{sga4.exp.I.appendix.II.sec.3}{${\rm n^{\circ}}$ 3}: apply the
preceding result to the family $(\Ob(C),F\ell\,(C))$. This applies to
categories of the form $(V\text{-}\mathcal{E})$, where $\mathcal{E}$
is a species of structure with morphisms and $V$ is a set, for example
a universe; in general one has $V\neq U$.

For example, if $V$ is a universe, then
$\Ob(V\text{-}\Ens)=V$ and $F\ell\,(V\text{-}\Ens)\subset V$
(\hyperlink{sga4.exp.I.appendix.II.sec.1}{${\rm n^{\circ}}$ 1}, \hyperlink{sga4.exp.I.appendix.II.sec.1.proposition.5}{Proposition 5}). The universes $U$ such that
$(V\text{-}\Ens)\in U$ are therefore those such that $V\in U$. But
the relation $V\in V$ is impossible for a universe: indeed, for every
subset $A$ of $V$, one would have $A\in V$ (\hyperlink{sga4.exp.I.appendix.II.sec.1}{${\rm n^{\circ}}$ 1},
\hyperlink{sga4.exp.I.appendix.II.sec.1.proposition.1}{Proposition 1}), whence $\card(\mathcal{P}(V))\leq\card(V)$, which is
impossible.

% END INCLUDED SGA4_English_current_buildable_source_20260906/source/01/drafts/SGA4_I_LINES_7994_8025_English.texfrag


% BEGIN INCLUDED SGA4_English_current_buildable_source_20260906/source/01/drafts/SGA4_I_LINES_8026_8095_English.texfrag
\section{Universes and strongly inaccessible cardinals}%
\hypertarget{sga4.exp.I.appendix.II.sec.5}{}\pageoriginale

Let $U$ be a universe. By condition (U.I) of \hyperlink{sga4.exp.I.appendix.II.sec.1.definition.1}{Definition 1}, every
element $x$ of $U$ is a subset of $U$; hence
$\card(x)\leq\card(U)$. Since the cardinals of the subsets of $U$ form
a well-ordered set, the cardinal
\hypertarget{sga4.exp.I.appendix.II.sec.5.equation.1.cardinal.c.U}{}
\begin{equation}
\sheaf{\sheaf{c}}(U)={\sup}_{x\in U}\card(x)\tag{1}
\end{equation}
exists. Note that, for every cardinal
$\sheaf{\sheaf{c}}<\sheaf{\sheaf{c}}(U)$, there exists an element $x$
of $U$ such that $\card(x)=\sheaf{\sheaf{c}}$; indeed, by definition
there exists $y\in U$ such that
$\sheaf{\sheaf{c}}\leq\card(y)\leq\sheaf{\sheaf{c}}(U)$, and one takes
for $x$ a suitable subset of $y$. Conversely, if $x\in U$, then
$\card(x)<\sheaf{\sheaf{c}}(U)$; indeed,
$\mathcal{P}(x)\in U$ by \hyperlink{sga4.exp.I.appendix.II.sec.1.definition.1.condition.U.III}{(U.III)}, whence
$2^{\card(x)}\leq\sheaf{\sheaf{c}}(U)$. Thus the cardinal
$\sheaf{\sheaf{c}}(U)$ has the following properties:

\begin{enumerate}
\renewcommand{\labelenumi}{\theenumi)}
\item If $\sheaf{\sheaf{c}}$ is a cardinal
$<\sheaf{\sheaf{c}}(U)$, then
$2^{\sheaf{\sheaf{c}}}<\sheaf{\sheaf{c}}(U)$; indeed, if $x$ is an
element of $U$ of cardinality $\sheaf{\sheaf{c}}$, then
$\mathcal{P}(x)\in U$ by \hyperlink{sga4.exp.I.appendix.II.sec.1.definition.1.condition.U.III}{(U.III)}, whence
$2^{\sheaf{\sheaf{c}}}<\sheaf{\sheaf{c}}(U)$.

\item If $(\sheaf{\sheaf{c}}_\lambda)_{\lambda\in I}$ is a family of
cardinals such that
$\sheaf{\sheaf{c}}_\lambda<\sheaf{\sheaf{c}}(U)$ for every
$\lambda\in I$ and $\card(I)<\sheaf{\sheaf{c}}(U)$, then the sum
cardinal $\Sigma_{\lambda\in I}\sheaf{\sheaf{c}}_\lambda$ is
$<\sheaf{\sheaf{c}}(U)$. Indeed, let $x_\lambda\in U$ be such that
$\card(x_\lambda)=\sheaf{\sheaf{c}}_\lambda$; after replacing $I$ by
an equipotent index set if necessary, one may suppose that $I\in U$;
then the sum set of the $x_\lambda$ is an element of $U$ (\hyperlink{sga4.exp.I.appendix.II.sec.1.corollary.to.proposition.4.2}{Corollary 2}
of \hyperlink{sga4.exp.I.appendix.II.sec.1.proposition.4}{Proposition 4} of \hyperlink{sga4.exp.I.appendix.II.sec.1}{${\rm n^{\circ}}$ 1}), which proves our assertion.
\end{enumerate}

We make the following definition:

\hypertarget{sga4.exp.I.appendix.II.sec.5.definition.2}{}
\begin{enonce*}{DEFINITION 2}
A\pageoriginale cardinal $\sheaf{\sheaf{d}}$ is called strongly
inaccessible if:
\begin{itemize}
\hypertarget{sga4.exp.I.appendix.II.sec.5.definition.2.condition.FI.1}{}
\item[(FI.1)] If $\sheaf{\sheaf{c}}$ is a cardinal such that
$\sheaf{\sheaf{c}}<\sheaf{\sheaf{d}}$, then
$2^{\sheaf{\sheaf{c}}}<\sheaf{\sheaf{d}}$.

\hypertarget{sga4.exp.I.appendix.II.sec.5.definition.2.condition.FI.2}{}
\item[(FI.2)] If $(\sheaf{\sheaf{c}}_\lambda)_{\lambda\in I}$ is a
family of cardinals such that
$\sheaf{\sheaf{c}}_\lambda<\sheaf{\sheaf{d}}$ for every
$\lambda\in I$, and if $\card(I)<\sheaf{\sheaf{d}}$, then
$\Sigma_{\lambda\in I}\sheaf{\sheaf{c}}_\lambda<
\sheaf{\sheaf{d}}$.
\end{itemize}
\end{enonce*}

\begin{enonce*}[remark]{Examples}
The cardinal $0$ and the countably infinite cardinal are strongly
inaccessible. No nonzero finite cardinal is strongly inaccessible.
Thus we have just proved that the universe axiom \hyperlink{sga4.exp.I.appendix.II.sec.4.axiom.A.6}{(A.6)} implies the
statement:
\begin{itemize}
\hypertarget{sga4.exp.I.appendix.II.sec.5.axiom.A-prime.6}{}
\item[${\rm (A'.6)}$] \emph{Every cardinal is strictly bounded above
by a strongly inaccessible cardinal.}
\end{itemize}
Conversely:
\end{enonce*}

% END INCLUDED SGA4_English_current_buildable_source_20260906/source/01/drafts/SGA4_I_LINES_8026_8095_English.texfrag


% BEGIN INCLUDED SGA4_English_current_buildable_source_20260906/source/01/drafts/SGA4_I_LINES_8096_8244_English.texfrag
\hypertarget{sga4.exp.I.appendix.II.sec.5.theorem.1}{}
\begin{enonce*}{THEOREM %%
1} Statement \hyperlink{sga4.exp.I.appendix.II.sec.5.axiom.A-prime.6}{$(A'.6)$} implies the universe axiom (A.6).
\end{enonce*}
Indeed, let $A$ be a set. We must construct a universe $U$ of which
$A$ is an element. Define recursively a sequence $(A_n)_{n\geq0}$ of
sets by:
\begin{itemize}
\hypertarget{sga4.exp.I.appendix.II.sec.5.theorem.1.proof.equation.2}{}
\item[(2)] $A_0=A$, $A_{n+1}=$ the union of the elements of $A_n=$
the set of elements of the elements of $A_n$.
\end{itemize}
\noindent Put $B=\bigcup^\infty_{n=0}A_n$. By \hyperlink{sga4.exp.I.appendix.II.sec.5.axiom.A-prime.6}{$(A'.6)$}, let
$\sheaf{\sheaf{c}}$ be a strongly inaccessible cardinal such that
$\card(B)<\sheaf{\sheaf{c}}$.

There exists a \emph{well-ordered} set $I$ such that
$\card(I)=\sheaf{\sheaf{c}}$. After replacing $I$ by its smallest
segment of cardinality $\sheaf{\sheaf{c}}$ if necessary, we may suppose
that \emph{every segment of $I$ distinct from $I$ has cardinality
$<\sheaf{\sheaf{c}}$.} We\pageoriginale denote by $\mathcal{E}$ the
\emph{least element of $I$}. It follows from the hypothesis on segments
that $I$ has no greatest element (otherwise we would remove it), and
hence that every element of $I$ has a successor; for $\alpha\in I$, we
denote the \emph{successor} of $\alpha$ by $s(\alpha)$.

Now define, by \emph{transfinite recursion}, a family
$(B_\alpha)_{\alpha\in I}$ of sets by:
\hypertarget{sga4.exp.I.appendix.II.sec.5.theorem.1.proof.equation.3}{}
\begin{equation}
\begin{cases}
& B_{\mathcal{E}}=B\\
& B_{s(\alpha)}=B_{\alpha}\cup \mathcal{P}(B_{\alpha})\\
& B_{\alpha}={\displaystyle{\bigcup_{\beta<\alpha}}}B_{\beta}\mbox{
if $\alpha$ has no predecessor.}
\end{cases}\tag{3}
\end{equation}
Put $U=\bigcup_{\alpha\in I}B_\alpha$. We shall show that $U$ is the
desired universe.

First, $A\in U$, since $A=A_\circ$ is a subset of
$B=B_{\mathcal{E}}$, and therefore an element of
$\mathcal{P}(B_{\mathcal{E}})\subset B_{s(\mathcal{E})}$.

Next note that, for every $\alpha\in I$, one has:
\hypertarget{sga4.exp.I.appendix.II.sec.5.theorem.1.proof.equation.4}{}
\begin{equation}
\card(B_{\alpha})<\sheaf{\sheaf{c}}{\quoi}.\tag{4}
\end{equation}
Indeed, proceed by transfinite recursion on $\alpha$. This is true for
$\alpha=\mathcal{E}$ by hypothesis. From \hyperlink{sga4.exp.I.appendix.II.sec.5.theorem.1.proof.equation.4}{(4)} one obtains
$\card(B_{s(\alpha)})\leq
\card(B_\alpha)+2^{\card(B_\alpha)}<
\sheaf{\sheaf{c}}+\sheaf{\sheaf{c}}$ (by \hyperlink{sga4.exp.I.appendix.II.sec.5.definition.2.condition.FI.1}{(FI.1)})
$=\sheaf{\sheaf{c}}$ (since $\sheaf{\sheaf{c}}$ is infinite).
Finally, if $\alpha$ has no predecessor, the set $I'$ of
$\beta<\alpha$ has cardinality $<\sheaf{\sheaf{c}}$ by construction;
hence (if $\card(B_\beta)<\sheaf{\sheaf{c}}$ for every
$\beta<\alpha$),
$\card(B_\alpha)=\card(\bigcup_{\beta\in I'}B_\beta)
\leq\Sigma_{\beta\in I'}\card(B_\beta)<\sheaf{\sheaf{c}}$ (by
(FI.2)). This proves \hyperlink{sga4.exp.I.appendix.II.sec.5.theorem.1.proof.equation.4}{(4)}.

Let us now\pageoriginale show that
\hypertarget{sga4.exp.I.appendix.II.sec.5.theorem.1.proof.equation.5}{}
\begin{equation}
\card(x)<\sheaf{\sheaf{c}}\mbox{ for every }x\in U{\quoi}.\tag{5}
\end{equation}
Indeed, it suffices to show that, for every $\alpha\in I$, one has
``$\card(x)<\sheaf{\sheaf{c}}$ for every $x\in B_\alpha$''. Proceed
again by transfinite recursion on $\alpha$. This is true for
$\alpha=\mathcal{E}$, since, if $x\in B_{\mathcal{E}}=B$, there exists
$n\geq0$ such that $x\in A_n$; hence $x\subset A_{n+1}$ by \hyperlink{sga4.exp.I.appendix.II.sec.5.theorem.1.proof.equation.2}{(2)},
$x\subset B$, and
$\card(x)\leq\card(B)<\sheaf{\sheaf{c}}$. The case where $\alpha$ has
no predecessor is evident. Finally, if $x\in B_\alpha$ and
$\alpha=s(\beta)$, either $x\in B_\beta$, in which case the assertion
``$\card(x)<\sheaf{\sheaf{c}}$'' is true by recursion, or
$x\in\mathcal{P}(B_\beta)$, whence $x\subset B_\beta$ and the
assertion ``$\card(x)<\sheaf{\sheaf{c}}$'' is true by \hyperlink{sga4.exp.I.appendix.II.sec.5.theorem.1.proof.equation.4}{(4)}.

We are now in a position to prove that $U$ is indeed a \emph{universe}:

\begin{itemize}
\hypertarget{sga4.exp.I.appendix.II.sec.5.theorem.1.proof.clause.U.I}{}
\item[(U.I)] Let $x\in U$ and $y\in x$. We must show %%
that $y\in U$. In other words, we must show that, for every
$\alpha\in I$, one has:
$$
\textrm{``}x\in B_{\alpha}\mbox{ and }y\in x\Longrightarrow y\in
B_{\alpha}\textrm{''\,}{\quoi}.
$$
Proceed again by transfinite recursion on $\alpha$. This is true for
$\alpha=\mathcal{E}$, since $x\in B_{\mathcal{E}}=B$ implies
$x\in A_n$ for some $n$; hence, if $y\in x$, then $y\in A_{n+1}$,
whence $y\in B=B_{\mathcal{E}}$. Passage to an element $\alpha$ with no
predecessor is evident. Finally, pass from $\alpha$ to $s(\alpha)$: if
$x\in B_{s(\alpha)}$ and $y\in x$, either $x\in B_\alpha$, whence
$y\in B_\alpha\subset B_{s(\alpha)}$ by recursion, or
$x\in\mathcal{P}(B_\alpha)$%%
, whence again $y\in B_\alpha\subset B_{s(\alpha)}$.

\hypertarget{sga4.exp.I.appendix.II.sec.5.theorem.1.proof.clause.U.II}{}
\item[(U.II)] Let $x,y\in U$. Since the family
$(B_\alpha)_{\alpha\in I}$ is increasing by \hyperlink{sga4.exp.I.appendix.II.sec.5.theorem.1.proof.equation.3}{(3)}, there exists
$\alpha\in I$ such that $x,y\in B_\alpha$. Then
$\{x,y\}\in\mathcal{P}(B_\alpha)\subset B_{s(\alpha)}\subset U$.

\hypertarget{sga4.exp.I.appendix.II.sec.5.theorem.1.proof.clause.U.III}{}
\item[(U.III)] Let\pageoriginale $x\in U$. Then there exists
$\alpha\in I$ such that $x\in B_\alpha$. It therefore suffices to
show that, for every $\alpha\in I$, one has:
$$
\textrm{``}x\in B_{\alpha}\Longrightarrow \mathcal{P}(x)\in
B_{s(s(\alpha))}\textrm{''\,}{\quoi}.
$$
Proceed again by transfinite recursion on $\alpha$. If
$\alpha=\mathcal{E}$, then $x\in A_n$ for some $n$, whence
$y\subset A_{n+1}\subset B=B_{\mathcal{E}}$ for every subset $y$ of
$x$; thus $y\in\mathcal{P}(B_{\mathcal{E}})\subset
B_{s(\mathcal{E})}$ for every $y\in\mathcal{P}(x)$, whence
$\mathcal{P}(x)\subset B_{s(\mathcal{E})}$ and
$\mathcal{P}(x)\in\mathcal{P}(B_{s(\mathcal{E})})\subset
B_{s(s(\mathcal{E}))}$, so our assertion is true for
$\alpha=\mathcal{E}$. %%
Passage to an element $\alpha$ with no predecessor is immediate,
since $\beta<\alpha$ then implies $s(s(\beta))<\alpha$. Finally, pass
from $\alpha$ to $s(\alpha)$; let $x\in B_{s(\alpha)}$; if
$x\in B_\alpha$, then
$\mathcal{P}(x)\in B_{s(s(\alpha))}\subset B_{s(s(s(\alpha)))}$ by
recursion; if $x\in\mathcal{P}(B_\alpha)$, then $x\subset B_\alpha$,
whence $\mathcal{P}(x)\subset\mathcal{P}(B_\alpha)$ and
$\mathcal{P}(x)\in\mathcal{P}(\mathcal{P}(B_\alpha))\in
B_{s(s(s(\alpha)))}$.

\hypertarget{sga4.exp.I.appendix.II.sec.5.theorem.1.proof.clause.U.IV}{}
\item[(U.IV)] Let $(x_\lambda)_{\lambda\in K}$ be a family of
elements of $U$ such that $K\in U$. We must show that the union
$x=\bigcup_{\lambda\in K}x_\lambda$ is an element of $U$. For every
$\lambda\in K$, choose an $\alpha(\lambda)\in I$ such that
$x_\lambda\in B_{\alpha(\lambda)}$. Let us show that the set
$\alpha(K)$ of the $\alpha(\lambda)$ is \emph{bounded above} in $I$:
indeed, if it were not, one would have
$$
I=\bigcup_{\lambda%%
\in K}[\mathcal{E},\alpha(\lambda)]{\quoi},
$$
which would contradict \hyperlink{sga4.exp.I.appendix.II.sec.5.definition.2.condition.FI.2}{(FI.2)}, since
$\card(K)<\sheaf{\sheaf{c}}$ and
$\card([\mathcal{E},\alpha(\lambda)])<\sheaf{\sheaf{c}}$ for every
$\lambda\in K$. Let then $\beta$ be an upper bound of $\alpha(K)$;
one has $x_\lambda\in B_\beta$ for every $\lambda\in K$, since the
family $(B_\alpha)_{\alpha\in I}$ is increasing. Every element of
$x_\lambda$ then belongs to $B_\beta$ by what was proved in \hyperlink{sga4.exp.I.appendix.II.sec.5.theorem.1.proof.clause.U.I}{(U.I)},
whence $x=\bigcup_{\lambda\in K}x_\lambda\subset B_\beta$. By \hyperlink{sga4.exp.I.appendix.II.sec.5.theorem.1.proof.equation.3}{(3)}, it
follows that $x\in B_{s(\beta)}$, whence $x\in U$. \cqfd
\end{itemize}

% END INCLUDED SGA4_English_current_buildable_source_20260906/source/01/drafts/SGA4_I_LINES_8096_8244_English.texfrag


% BEGIN INCLUDED SGA4_English_current_buildable_source_20260906/source/01/drafts/SGA4_I_LINES_8245_8358_English.texfrag
\hypertarget{sga4.exp.I.appendix.II.sec.5.definition.3}{}
\begin{enonce*}{DEFINITION 3}
Let\pageoriginale $x$, $y$ be two sets and let $n$ be an integer
$\geq 0$. We say that $y$ is a component of order $n$ of $x$ if there
exists a sequence $(x_j)_{j=0,\ldots,n}$ such that $x_0=x$, $x_n=y$,
and $x_{j+1}\in x_j$ for $j=0,\ldots,n-1$.
\end{enonce*}

Thus $x$ is the only component of order 0 of $x$. The components of
order 1 (\resp 2) of $x$ are the elements of $x$ (\resp the elements
of the elements of $x$). We say that $y$ is a \emph{component} of $x$
if there exists $n\geq 0$ such that $y$ is a component of order $n$ of
$x$. The relation ``$y$ is a component of $x$'' is a \emph{preorder
relation}. By the selection--union scheme (Chapter II, \ldots), the
relation ``$y$ is a component of $x$'' is collectivizing with respect
to $y$, so that the components of $x$ form a \emph{set}.

\hypertarget{sga4.exp.I.appendix.II.sec.5.definition.4}{}
\begin{enonce*}{DEFINITION 4}
Let $\sheaf{\sheaf{c}}$ be a cardinal. A set $x$ is said to be of
type $\sheaf{\sheaf{c}}$ (\resp of strict type
$\sheaf{\sheaf{c}}$, of finite type) if all the components of $x$ have
cardinalities $\leq \sheaf{\sheaf{c}}$ (\resp
$<\sheaf{\sheaf{c}}$, finite cardinalities).
\end{enonce*}

\begin{enonce*}[remark]{Examples}
The elements of the universe $U_1$ (\hyperlink{sga4.exp.I.appendix.II.sec.1}{${\rm n^{\circ}}$ 1}, \hyperlink{sga4.exp.I.appendix.II.sec.1.example.2}{Example 2})
are all of finite type. If $\sheaf{\sheaf{c}}$ is a cardinal and $x$
is of type $\sheaf{\sheaf{c}}$ (\resp of strict type
$\sheaf{\sheaf{c}}$, of finite type), then every component of $x$ and
every subset of $x$ is of type $\sheaf{\sheaf{c}}$ (\resp of strict
type $\sheaf{\sheaf{c}}$, of finite type); likewise
$\mathcal{P}(x)$ is of type $2^{\sheaf{\sheaf{c}}}$ (\resp of strict
type $2^{\sheaf{\sheaf{c}}}$, of finite type). If $x$ is of type
$\sheaf{\sheaf{c}}$ and if $\sheaf{\sheaf{c}}\leq
\sheaf{\sheaf{c}}'$, then $x$ is of type $\sheaf{\sheaf{c}}'$.
\end{enonce*}

\hypertarget{sga4.exp.I.appendix.II.sec.5.lemma.of.section.5}{}
\begin{enonce*}{LEMMA}
Let\pageoriginale $\sheaf{\sheaf{c}}$ be an uncountable strongly
inaccessible cardinal, and let $X$ be a set of strict type
$\sheaf{\sheaf{c}}$. Then there exists a cardinal
$\sheaf{\sheaf{d}}<\sheaf{\sheaf{c}}$ such that $X$ is of type
$\sheaf{\sheaf{d}}$.
\end{enonce*}

Indeed, for every integer $n$, let $\sheaf{\sheaf{d}}_n$ denote the
cardinality of the set $X_n$ of components of order $n$ of $X$. We
have $\sheaf{\sheaf{d}}_0=1$,
$\sheaf{\sheaf{d}}_1=\card(X)<\sheaf{\sheaf{c}}$, and, since
$X_n=\bigcup_{y\in X_{n-1}} y$, induction on $n$ using \hyperlink{sga4.exp.I.appendix.II.sec.5.definition.2.condition.FI.2}{(FI.2)} gives
$\sheaf{\sheaf{d}}_n<\sheaf{\sheaf{c}}$ for every $n$. The
$\sheaf{\sheaf{d}}_n$ are therefore bounded above by the cardinal
$\sheaf{\sheaf{d}}=\Sigma_{j\geq 0}\sheaf{\sheaf{d}}_j$, which is
$<\sheaf{\sheaf{c}}$ by (FI.2) and the hypothesis that
$\sheaf{\sheaf{c}}$ is uncountable. If $Y$ is a component of order
$n$ of $X$, then $\card(Y)\leq \card(X_{n+1})\leq
\sheaf{\sheaf{d}}$. \cqfd

\hypertarget{sga4.exp.I.appendix.II.sec.5.proposition.10}{}
\begin{enonce*}{PROPOSITION 10}
Let $U$ be a universe and let $\sheaf{\sheaf{c}}$ be a strongly
inaccessible cardinal. Then the set $U'$ of those $x\in U$ that are
of strict type $\sheaf{\sheaf{c}}$ is a universe.
\end{enonce*}
Indeed, if $x$, $y\in U'$ and if $z\in x$, then plainly $z\in U'$
and $\{x,y\}\in U'$. If $x\in U'$, then
$\card(x)<\sheaf{\sheaf{c}}$, whence
$\card(\mathcal{P}(x))<\sheaf{\sheaf{c}}$ by \hyperlink{sga4.exp.I.appendix.II.sec.5.definition.2.condition.FI.1}{(FI.1)}; since a
component of $\mathcal{P}(x)$ is either $\mathcal{P}(x)$ itself, a
subset of $x$, or a component of $x$, we have
$\mathcal{P}(x)\in U'$. Finally, if
$(x_{\alpha})_{\alpha\in I}$ is a family of elements of $U'$ such that
$I\in U'$, then
$\card(\bigcup_{\alpha}x_{\alpha})<\sheaf{\sheaf{c}}$ by \hyperlink{sga4.exp.I.appendix.II.sec.5.definition.2.condition.FI.2}{(FI.2)};
since every component of order $>0$ of
$\bigcup_{\alpha}x_{\alpha}$ is a component of some $x_{\alpha}$, we
indeed have $\bigcup_{\alpha}x_{\alpha}\in U'$. \cqfd

\medskip
\noindent \emph{Remark on the cardinal} $\sheaf{\sheaf{c}}(U)$.
Let $U$ be a universe and let $\sheaf{\sheaf{c}}(U)$ be the cardinal
defined by \hyperlink{sga4.exp.I.appendix.II.sec.5.equation.1.cardinal.c.U}{(1)}, \ie
$$
\sheaf{\sheaf{c}}(U)={\sup}_{x\in U}\card(x){\quoi}.
$$
We\pageoriginale plainly have $\sheaf{\sheaf{c}}(U)\leq \card(U)$,
since, as we recall, $x\in U$ implies $x\subset U$. But the equality
$\sheaf{\sheaf{c}}(U)=\card(U)$ does not always hold. For example, let
$(x_{\alpha})_{\alpha\in I}$ be an \emph{uncountable} family of
symbols subject to the axioms:
$$
\text{«} x_{\alpha}=\{x_{\alpha}\}\mbox{ for every } \alpha\in I
\text{»} \quav \text{«} x_{\alpha}\neq x_{\beta}\mbox{ if }
\alpha\neq \beta \text{»}
$$
(in other words, $x_{\alpha}$ is a one-element set, namely itself).
As in \hyperlink{sga4.exp.I.appendix.II.sec.1.example.2}{Example 2} of \hyperlink{sga4.exp.I.appendix.II.sec.1}{${\rm n^{\circ}}$ 1}, form the ``significant words''
constructed from the symbols $∅$, $x_{\alpha}$, $\{,\}$ %%
and so forth. Each denotes a \emph{finite} set. Let $U$ be the set of
these words. One readily verifies that the conditions of \hyperlink{sga4.exp.I.appendix.II.sec.1.definition.1}{Definition 1}
are satisfied, so that $U$ is a \emph{universe}. Since the significant
words are finite sequences of elements of a set of cardinality
$\card(I)+4=\card(I)$, we have $\card(U)=\card(I)$ (Chapter III; in
fact, $\card(U)\geq \card(I)$ suffices, and this is obvious). On the
other hand, $\sheaf{\sheaf{c}}(U)$ is the countably infinite cardinal,
and hence $\sheaf{\sheaf{c}}(U)<\card(U)$.

The strict inequality $\sheaf{\sheaf{c}}(U)<\card(U)$ is due to our
having introduced sets $x_{\alpha}$ such that $x_{\alpha}\in
x_{\alpha}$. If one rules out horrors of this kind, one obtains
$\card(U)=\sheaf{\sheaf{c}}(U)$ for every universe $U$, as well as
other very pretty results ``which cannot possibly be useful.''
This is what we shall do in the next section.

% END INCLUDED SGA4_English_current_buildable_source_20260906/source/01/drafts/SGA4_I_LINES_8245_8358_English.texfrag


% BEGIN INCLUDED SGA4_English_current_buildable_source_20260906/source/01/drafts/SGA4_I_LINES_8359_8513_English.texfrag
\section{Artinian sets and universes}%
\hypertarget{sga4.exp.I.appendix.II.sec.6}{}

\hypertarget{sga4.exp.I.appendix.II.sec.6.definition.5}{}
\begin{enonce*}{DEFINITION 5}
A set $x$ is said to be Artinian if there is no infinite sequence
$(x_n)_{n\geq 0}$ such that $x_{\circ}=x$ and $x_{n+1}\in
x_n$ for every $n\geq 0$.
\end{enonce*}

\begin{enonce*}[remark]{Examples}
The sets $∅$, $\{∅\}$,
$\{∅,\{∅\}\}$, and, more generally, the elements of
the universe $U_1$ of \hyperlink{sga4.exp.I.appendix.II.sec.1}{${\rm n^{\circ}}$ 1} are Artinian.
\end{enonce*}

In\pageoriginale pictorial terms, if $x$ is Artinian, and if one takes
an element $x_1$ of $x$, then an element $x_2$ of $x_1$, \etc, the
process must stop, and one arrives at a component $x_n$ of $x$
that is \emph{empty}. In other words, Artinian sets are
``built from $∅$''; this will be made more precise
later (\cf \hyperlink{sga4.exp.I.appendix.II.sec.6.theorem.2.proof.claim.star}{(*)} in the proof of \hyperlink{sga4.exp.I.appendix.II.sec.6.theorem.2}{Theorem 2}).

Artinian sets plainly have the following properties:
\begin{itemize}
\hypertarget{sga4.exp.I.appendix.II.sec.6.property.AR.I}{}
\item[(AR.I)] If $x$ is Artinian, every subset of $x$ and every
component of $x$ is Artinian. For $x$ to be Artinian, it is necessary
and sufficient that every element of $x$ be Artinian.

\hypertarget{sga4.exp.I.appendix.II.sec.6.property.AR.II}{}
\item[(AR.II)] If $x$ and $y$ are Artinian, then $\{x,y\}$ is
Artinian.

\hypertarget{sga4.exp.I.appendix.II.sec.6.property.AR.III}{}
\item[(AR.III)] If $x$ is Artinian, $\mathcal{P}(x)$ is Artinian.

\hypertarget{sga4.exp.I.appendix.II.sec.6.property.AR.IV}{}
\item[(AR.IV)] Every union of Artinian sets is an Artinian set.
\end{itemize}

These properties immediately yield the following result:

\hypertarget{sga4.exp.I.appendix.II.sec.6.proposition.11}{}
\begin{enonce*}{PROPOSITION 11}
If $U$ is a universe, the set of Artinian elements of $U$ is
a universe, necessarily Artinian.
\end{enonce*}

\hypertarget{sga4.exp.I.appendix.II.sec.6.corollary.to.proposition.11}{}
\begin{enonce*}{COROLLARY}
If $x$ is an Artinian set, there exists an Artinian universe $U$
such that $x\in U$.
\end{enonce*}

Indeed, by axiom \hyperlink{sga4.exp.I.appendix.II.sec.4.axiom.A.6}{(A.6)}, $x$ is an element of a universe $V$; take
for $U$ the set of Artinian elements of $V$.

The following proposition is even less useful than the rest of ${\rm
n^{\circ}}$ 6:

\hypertarget{sga4.exp.I.appendix.II.sec.6.proposition.12}{}
\begin{enonce*}{PROPOSITION 12}
Let\pageoriginale $A$ be an Artinian set. Then:
\begin{enumerate}
\renewcommand{\theenumi}{\alph{enumi}}
\renewcommand{\labelenumi}{\theenumi)}
\hypertarget{sga4.exp.I.appendix.II.sec.6.proposition.12.part.a}{}
\item For every $x\in A$, one has $x\not\in x$;

\hypertarget{sga4.exp.I.appendix.II.sec.6.proposition.12.part.b}{}
\item If $x$, $y\in A$, one cannot have both $x\in y$ and $y\in
x$;

\hypertarget{sga4.exp.I.appendix.II.sec.6.proposition.12.part.c}{}
\item For every $x\in A$, the relation ``$y$ is a component of $z$'' on the components $y$, $z$ of $x$ is an order relation;

\hypertarget{sga4.exp.I.appendix.II.sec.6.proposition.12.part.d}{}
\item For every nonempty element $x$ of $A$, there exists $y\in x$
such that $x\cap y=∅$.
\end{enumerate}
\end{enonce*}
Indeed, the negation of a) (\resp of b)) entails the existence of an
infinite sequence $(x_n)_{n\geq 0}$ contradicting \hyperlink{sga4.exp.I.appendix.II.sec.6.definition.5}{Definition 5}, namely
$(x,x,x,\ldots)$ (\resp $(x,y,x,y,x,y,\ldots)$). The negation of c)
means that there exist $x\in A$, distinct components $y$, $z$ of $x$,
and membership chains:
$$
y\in y_1\in\ldots \in y_q\in z\quav z\in z_1\in\ldots\in z_r\in
y{\quoi};
$$
whence, as in \hyperlink{sga4.exp.I.appendix.II.sec.6.proposition.12.part.b}{b)}, an infinite sequence $(x_n)_{n\geq 0}$
contradicting \hyperlink{sga4.exp.I.appendix.II.sec.6.definition.5}{Definition 5}. Finally, if d) is false, there exists a
nonempty element $x$ of $A$ such that $x\cap y\neq ∅$ for
every $y\in x$; set $x_\circ=x$, and take for $x_1$ an
element of $x$; since $x\cap x_1\neq ∅$, take for
$x_2$ an element of $x\cap x_1$, and so on; more formally, define
by induction an infinite sequence $(x_n)_{n\geq 1}$
of elements of $x$ by means of $x_1\in x$, $x_{n+1}\in x_n\cap x$
for $n\geq 1$; then the sequence $(x_n)_{n\geq 0}$ contradicts
\hyperlink{sga4.exp.I.appendix.II.sec.6.definition.5}{Definition 5}.

\begin{enonce*}[remark]{Remark}
Let $B$ be a set. Then $B$ is Artinian if and only if every set $A$
whose elements are sets of components of $B$ satisfies \hyperlink{sga4.exp.I.appendix.II.sec.6.proposition.12.part.d}{condition d)}
of \hyperlink{sga4.exp.I.appendix.II.sec.6.proposition.12}{Proposition 12}. Indeed, necessity
follows from \hyperlink{sga4.exp.I.appendix.II.sec.6.property.AR.I}{(AR.I)} and\pageoriginale \hyperlink{sga4.exp.I.appendix.II.sec.6.proposition.12}{Proposition 12}. Conversely,
if $B$ is not Artinian, there exists an infinite sequence $(x_n)_{n\geq
0}$ with $x_{\circ}=B$ and $x_{n+1}\in x_n$ for every $n\geq 0$; then
take for $A$ the set whose sole element is the set $X$ of the $x_n$;
thus $A$ contains a nonempty element $X$ such that, for every
element $y$ of $X$, one has $y\cap X\neq ∅$ (indeed, $y$
is of the form $x_n$, and one has $x_{n+1}\in y\cap X$).
\end{enonce*}

\hypertarget{sga4.exp.I.appendix.II.sec.6.theorem.2}{}
\begin{enonce*}{THEOREM 2}
Let $\sheaf{\sheaf{c}}$ be an infinite cardinal. Then:
\begin{itemize}
\hypertarget{sga4.exp.I.appendix.II.sec.6.theorem.2.part.a}{}
\item[a)] The relation ``$x$ is an Artinian set of type $\sheaf{\sheaf{c}}$
(\resp of strict type $\sheaf{\sheaf{c}}$)'' %%
is collectivizing with respect to $x$; the set
$A_{\sheaf{\sheaf{c}}}$ of Artinian sets of type
$\sheaf{\sheaf{c}}$ has cardinality $2^{\sheaf{\sheaf{c}}}$.

\hypertarget{sga4.exp.I.appendix.II.sec.6.theorem.2.part.b}{}
\item[b)] If $\sheaf{\sheaf{c}}$ is strongly inaccessible, the set
$U_{\sheaf{\sheaf{c}}}$ of Artinian sets of strict type
$\sheaf{\sheaf{c}}$ is a universe of cardinality $\sheaf{\sheaf{c}}$;
the cardinal $\sheaf{\sheaf{c}}(U_{\sheaf{\sheaf{c}}})=\sup_{x\in
U_{\sheaf{\sheaf{c}}}}\card(x)$ is $\sheaf{\sheaf{c}}$.

\hypertarget{sga4.exp.I.appendix.II.sec.6.theorem.2.part.c}{}
\item[c)] If a universe $U$ contains an element of cardinality $\sheaf{\sheaf{c}}$,
every Artinian set of type $\sheaf{\sheaf{c}}$ belongs to $U$
(in other words, $A_{\sheaf{\sheaf{c}}}\subset U$).

\hypertarget{sga4.exp.I.appendix.II.sec.6.theorem.2.part.c-prime}{}
\item[${\rm c')}$] If a universe $U$ is nonempty, every Artinian
set of finite type is an element of $U$.
\end{itemize}
\end{enonce*}
Before proving \hyperlink{sga4.exp.I.appendix.II.sec.6.theorem.2}{Theorem 2}, let us deduce from it a few illuminating
corollaries:

\hypertarget{sga4.exp.I.appendix.II.sec.6.corollary.to.theorem.2.1}{}
\begin{enonce*}{COROLLARY 1}
If a universe $U$ is Artinian, then $\card(U)$ is strongly
inaccessible, and $U$ is the set of Artinian sets of strict type
$\card(U)$.
\end{enonce*}

Indeed,\pageoriginale set $\sheaf{\sheaf{c}}(U)=\sup_{x\in
U}\card(x)$. This is a strongly inaccessible cardinal (beginning of
\hyperlink{sga4.exp.I.appendix.II.sec.5}{${\rm n^{\circ}}$ 5}). First suppose that it is uncountable; then,
for every infinite cardinal $\sheaf{\sheaf{c}}<\sheaf{\sheaf{c}}(U)$,
every Artinian set of type $\sheaf{\sheaf{c}}$ is an element of
$U$ by \hyperlink{sga4.exp.I.appendix.II.sec.6.theorem.2.part.c}{c)}; hence every Artinian set of strict type
$\sheaf{\sheaf{c}}(U)$ is an element of $U$ by the \hyperlink{sga4.exp.I.appendix.II.sec.5.lemma.of.section.5}{lemma of ${\rm
n^{\circ}}$ 5}. This last assertion remains valid if
$\sheaf{\sheaf{c}}(U)$ is countable, by \hyperlink{sga4.exp.I.appendix.II.sec.6.theorem.2.part.c-prime}{${\rm c'}$)}. Conversely,
by (U.I), every set in $U$ is of strict type
$\sheaf{\sheaf{c}}(U)$. Thus $U$ is the universe
$U_{\sheaf{\sheaf{c}}(U)}$ of \hyperlink{sga4.exp.I.appendix.II.sec.6.theorem.2.part.b}{b)}, whence
$\sheaf{\sheaf{c}}(U)=\card(U)$ by b).

It follows from \hyperlink{sga4.exp.I.appendix.II.sec.6.corollary.to.theorem.2.1}{Corollary 1} that an Artinian universe is uniquely
\emph{determined} by its cardinality (which is moreover strongly
inaccessible). Thus there is a ``one-to-one correspondence'' between
Artinian universes and strongly inaccessible cardinals. In particular:

% END INCLUDED SGA4_English_current_buildable_source_20260906/source/01/drafts/SGA4_I_LINES_8359_8513_English.texfrag


% BEGIN INCLUDED SGA4_English_current_buildable_source_20260906/source/01/drafts/SGA4_I_LINES_8514_8642_English.texfrag
\hypertarget{sga4.exp.I.appendix.II.sec.6.corollary.to.theorem.2.2}{}
\begin{enonce*}{COROLLARY 2}
The inclusion relation $U\subset U'$ on Artinian universes is a
well-ordering relation.
\end{enonce*}

Indeed, the relation $\sheaf{\sheaf{c}}\leq \sheaf{\sheaf{c}}'$ on
cardinals is a well-ordering relation (Chapter III), and, with the
notation of \hyperlink{sga4.exp.I.appendix.II.sec.6.theorem.2.part.b}{b)} of \hyperlink{sga4.exp.I.appendix.II.sec.6.theorem.2}{Theorem 2}, the relations $\sheaf{\sheaf{c}}\leq
\sheaf{\sheaf{c}}'$ and $U_{\sheaf{\sheaf{c}}}\subset
U_{\sheaf{\sheaf{c}}'}$ are equivalent.

Observe that \hyperlink{sga4.exp.I.appendix.II.sec.6.theorem.2.part.b}{b)} of \hyperlink{sga4.exp.I.appendix.II.sec.6.theorem.2}{Theorem 2} gives a second proof of \hyperlink{sga4.exp.I.appendix.II.sec.5.theorem.1}{Theorem 1}
(\hyperlink{sga4.exp.I.appendix.II.sec.5}{${\rm n^{\circ}}$ 5}).

We now turn to the proof of \hyperlink{sga4.exp.I.appendix.II.sec.6.theorem.2}{Theorem 2}. Given a set $A$, we shall call
a \emph{chain} of $A$ any finite sequence
$(x_j)_{j=0,\ldots,n}$ such that $x_0=A$ and $x_{j+1}\in x_j$
for $j=0,\ldots,n-1$. The chains of $A$ form a\pageoriginale
\emph{set}, by the selection--union scheme; denote it by
$G(A)$. Given a chain $X=(x_j)_{j=0,\ldots,n}$, the chains of the
form $(x_i)_{i=0,\ldots,q}$ with $q\leq n$ will be called less than
$X$; this endows $G(A)$ with the structure of an \emph{ordered set}.
For $A$ to be Artinian, it is necessary and sufficient that $G(A)$
be a ``\emph{Noetherian}'' ordered set (i.e. one satisfying the
equivalent conditions of Chapter III, \S 6, ${\rm n^{\circ}}$ 5).

We shall prove that:

\hypertarget{sga4.exp.I.appendix.II.sec.6.theorem.2.proof.claim.star}{}
\noindent (*) \emph{If $A$ is Artinian, it is uniquely determined by
the isomorphism class of the ordered set $G(A)$.}

Indeed, given an ordered set $G$ and an element $g\in G$, let $S(g)$
denote the set of all $g'\in G$ such that $g<g'$ and such that
$g\leq h\leq g'$ implies $h=g$ or $h=g'$ (in other words, the set
of ``immediate successors'' of $g$). Consider the map $\theta$ that
assigns to every chain $X=(x_0,\ldots,x_n)$ in $G(A)$ the set $x_n$;
then, for every $X\in G(A)$, we have:
$$
\theta(X)=\{\theta(X')| X'\in S(X)\}{\quoi}.
$$
Since $A$ is the image under $\theta$ of the least element $X_0=(A)$
of $G(A)$, it will suffice to show that $\theta$ is uniquely
determined by the isomorphism class of the ordered set $G(A)$. This
follows from the following lemma:

\hypertarget{sga4.exp.I.appendix.II.sec.6.lemma.1}{}
\begin{enonce*}{LEMMA 1}
Let\pageoriginale $G$ be a Noetherian ordered set and let $\varphi$
be a function with domain $G$ such that, for every $g\in G$, one has:
\hypertarget{sga4.exp.I.appendix.II.sec.6.lemma.1.equation.1}{}
\begin{equation}
\varphi(g)=\{\varphi(g')| g'\in S(g)\}{\quoi}.\tag{1}
\end{equation}
Then $\varphi$ is uniquely determined. Moreover, if $U$ is a universe
containing an element equipotent to $G$, then $\varphi$ takes its
values in $U$.
\end{enonce*}

The hypothesis implies that $\varphi(g)=∅$ if $g$ is a maximal
element of $G$.

Indeed, let $\varphi$ and $\varphi'$ be two functions satisfying
\hyperlink{sga4.exp.I.appendix.II.sec.6.lemma.1.equation.1}{(1)}. If $\varphi\neq \varphi'$, the set of all $g\in G$ such that
$\varphi(g)\neq \varphi'(g)$ is nonempty and hence, since $G$ is
Noetherian, has a \emph{maximal} element $h$. We then have
$\varphi(g)=\varphi'(g)$ for every $g>h$, in particular for every
immediate successor $g\in S(h)$; hence $\varphi(h)=\varphi'(h)$ by
\hyperlink{sga4.exp.I.appendix.II.sec.6.lemma.1.equation.1}{(1)}, which is a contradiction. Thus indeed
$\varphi=\varphi'$. Now let $U$ be a universe containing an element
$x$ equipotent to $G$; let us show that $\varphi$ takes its values in
$U$. Otherwise, let $h$ be a maximal element among those $g\in G$
such that $\varphi(g)\not\in U$. Then $\varphi(g')\in U$ for every
$g'\in S(h)$, so that
$$
\varphi(h)=\{\varphi(g')| g'\in S(h)\}
$$
is a subset of $U$. Since its cardinality does not exceed
$\card(G)$, and hence does not exceed the cardinality of an element
of $U$, we have $\varphi(h)\in U$ (\hyperlink{sga4.exp.I.appendix.II.sec.1}{${\rm n^{\circ}}$ 1},
\hyperlink{sga4.exp.I.appendix.II.sec.1.proposition.7}{Proposition 7}). This contradiction shows that $\varphi$ takes its
values in $U$.

With this \pageoriginale %%
established, let us prove a) of \hyperlink{sga4.exp.I.appendix.II.sec.6.theorem.2}{Theorem 2}. We may confine ourselves to
the assertion without the \resp clause, since the other follows at
once. Let $\sheaf{\sheaf{c}}$ be an \emph{infinite} cardinal. If $A$
is a set of type $\sheaf{\sheaf{c}}$, then:
\hypertarget{sga4.exp.I.appendix.II.sec.6.theorem.2.proof.equation.2}{}
\begin{equation}
\card(G(A))\leq \sheaf{\sheaf{c}}{\quoi}%%
\ .\tag{2}
\end{equation}
Indeed, if $A_n$ denotes the set of components of order $n$ of $A$,
then $\card(A_1)\leq \sheaf{\sheaf{c}}$ and $\card(A_{n+1})\leq
\sheaf{\sheaf{c}}\cdot \card(A_n)$, whence $\card(A_n)\leq
\sheaf{\sheaf{c}}^n$ by induction. But
$\sheaf{\sheaf{c}}^n=\sheaf{\sheaf{c}}$, since
$\sheaf{\sheaf{c}}$ is infinite (Chapter III); hence
$\card(\bigcup^\infty_{n=0} A_n)\leq
\card(\NN)\cdot \sheaf{\sheaf{c}}=\sheaf{\sheaf{c}}$ (Chapter III).
Now $G(A)$ is a set of \emph{finite} sequences of elements of
$\bigcup_n A_n$, so inequality (2) indeed holds (Chapter III).
Now, if $E$ is a set of cardinality $\sheaf{\sheaf{c}}$, specifying
an order structure on a subset $E'$ of $E$ is equivalent to specifying
the subset of $E\times E$ consisting of the $(x,y)$ such that $x\in
E'$, $y\in E'$, and $x\leq y$. Consequently, by \hyperlink{sga4.exp.I.appendix.II.sec.6.theorem.2.proof.equation.2}{(2)}, the isomorphism
classes of the ordered sets $G(A)$ (where $A$ is of type
$\sheaf{\sheaf{c}}$) form a set
$\mathfrak{F}_{\sheaf{\sheaf{c}}}$, and
$\card(\mathfrak{F}_{\sheaf{\sheaf{c}}})\leq
\card(\mathcal{P}(E\times E))=2^{\sheaf{\sheaf{c}}}$. Let
$\mathfrak{F}'_{\sheaf{\sheaf{c}}}$ be the subset of
$\mathfrak{F}_{\sheaf{\sheaf{c}}}$ consisting of the classes of
Noetherian ordered sets having a least element. If, to every $G\in
\mathfrak{F}'_{\sheaf{\sheaf{c}}}$, one assigns the value of the
function $\varphi$ of \hyperlink{sga4.exp.I.appendix.II.sec.6.lemma.1}{Lemma 1} at the least element of $G$, one obtains
a map $\theta$ on $\mathfrak{F}'_{\sheaf{\sheaf{c}}}$ whose image
contains \emph{all} Artinian sets of type $\sheaf{\sheaf{c}}$. The
latter therefore do indeed form a set $A_{\sheaf{\sheaf{c}}}$, and
we have $\card(A_{\sheaf{\sheaf{c}}})\leq
\card(\mathfrak{F}'_{\sheaf{\sheaf{c}}})\leq
\card(\mathfrak{F}_{\sheaf{\sheaf{c}}})\leq
2^{\sheaf{\sheaf{c}}}$.

It remains %%
\pageoriginale to show that
$\card(A_{\sheaf{\sheaf{c}}})=2^{\sheaf{\sheaf{c}}}$. For this it
suffices to show that there exists an Artinian set $B$ of type
$\sheaf{\sheaf{c}}$ and cardinality $\sheaf{\sheaf{c}}$, since the
subsets of $B$ will then be elements of $A_{\sheaf{\sheaf{c}}}$.
This existence follows from the following lemma:

% END INCLUDED SGA4_English_current_buildable_source_20260906/source/01/drafts/SGA4_I_LINES_8514_8642_English.texfrag


% BEGIN INCLUDED SGA4_English_current_buildable_source_20260906/source/01/drafts/SGA4_I_LINES_8643_8705_English.texfrag
\hypertarget{sga4.exp.I.appendix.II.sec.6.lemma.2}{}
\begin{enonce*}{LEMMA 2}
For every cardinal $\sheaf{\sheaf{c}}$, there exists an Artinian set
$B$ of type $\sheaf{\sheaf{c}}$ and cardinality
$\sheaf{\sheaf{c}}$.
\end{enonce*}

Proceed by transfinite induction on $\sheaf{\sheaf{c}}$. For
$\sheaf{\sheaf{c}}=0$ there is no choice, and one takes $B=∅$. If
$\sheaf{\sheaf{c}}$ has a predecessor $\sheaf{\sheaf{c}}'$, let $B'$
be an Artinian set of type $\sheaf{\sheaf{c}}'$ and cardinality
$\sheaf{\sheaf{c}}'$; one then has
$\sheaf{\sheaf{c}}\leq2^{\sheaf{\sheaf{c}}'}$, so there exists a
subset $B$ of $\mathcal{P}(B')$ of cardinality
$\sheaf{\sheaf{c}}$; the elements of $B$ are subsets of $B'$ and
therefore have cardinalities
$\leq\sheaf{\sheaf{c}}'\leq\sheaf{\sheaf{c}}$; the higher-order
components of $B$ are components of $B'$, and hence also have
cardinalities $\leq\sheaf{\sheaf{c}}'\leq\sheaf{\sheaf{c}}$.
Finally, if $\sheaf{\sheaf{c}}$ has no predecessor, choose, for every
cardinal $\sheaf{\sheaf{c}}_\lambda<\sheaf{\sheaf{c}}$, an Artinian
set $B_\lambda$ of type $\sheaf{\sheaf{c}}_\lambda$ and cardinality
$\sheaf{\sheaf{c}}_\lambda$; then
$B=\bigcup_\lambda B_\lambda$ has the required property. This proves
\hyperlink{sga4.exp.I.appendix.II.sec.6.lemma.2}{Lemma 2} and completes the proof of part a).

We turn to b). Let $\sheaf{\sheaf{c}}$ be a strongly inaccessible
cardinal. By \hyperlink{sga4.exp.I.appendix.II.sec.6.theorem.2.part.a}{a)}, we already know that the Artinian sets of strict type
$\sheaf{\sheaf{c}}$ form a set $U_{\sheaf{\sheaf{c}}}$. The fact that
$U_{\sheaf{\sheaf{c}}}$ is a universe follows immediately from
properties \hyperlink{sga4.exp.I.appendix.II.sec.6.property.AR.I}{(AR.I)} to \hyperlink{sga4.exp.I.appendix.II.sec.6.property.AR.IV}{(AR.IV)} of Artinian sets (beginning of
\hyperlink{sga4.exp.I.appendix.II.sec.6}{${\rm n^{\circ}}$}), and\pageoriginale from cardinal estimates analogous
to those of \hyperlink{sga4.exp.I.appendix.II.sec.5.proposition.10}{Proposition 10} of \hyperlink{sga4.exp.I.appendix.II.sec.5}{${\rm n^{\circ}}$~5}. The relation
$\sup_{x\in U_{\sheaf{\sheaf{c}}}}\card(x)=\sheaf{\sheaf{c}}$ follows%%
 from \hyperlink{sga4.exp.I.appendix.II.sec.6.lemma.2}{Lemma 2}, applied to cardinals $<\sheaf{\sheaf{c}}$. Finally, to
show that
$\card(U_{\sheaf{\sheaf{c}}})=\sheaf{\sheaf{c}}$, first suppose that
$\sheaf{\sheaf{c}}$ is \emph{uncountable}; by%%
 the \hyperlink{sga4.exp.I.appendix.II.sec.5.lemma.of.section.5}{lemma of ${\rm n^{\circ}}$~5},
$U_{\sheaf{\sheaf{c}}}$ is the union
$\bigcup_{\sheaf{\sheaf{d}}<\sheaf{\sheaf{c}}}
A_{\sheaf{\sheaf{d}}}$, where $A_{\sheaf{\sheaf{d}}}$ denotes the
set of Artinian sets of type $\sheaf{\sheaf{d}}$; now
$\card(A_{\sheaf{\sheaf{d}}})=2^{\sheaf{\sheaf{d}}}
<\sheaf{\sheaf{c}}$ (by \hyperlink{sga4.exp.I.appendix.II.sec.6.theorem.2.part.a}{a)}); on the other hand, the set of cardinals
$\sheaf{\sheaf{d}}<\sheaf{\sheaf{c}}$ has cardinality
$\leq\sheaf{\sheaf{c}}$ (Chapter III); whence
$\card(U_{\sheaf{\sheaf{c}}})\leq
\sheaf{\sheaf{c}}\cdot\sheaf{\sheaf{c}}=\sheaf{\sheaf{c}}$, and also
$\card(U_{\sheaf{\sheaf{c}}})\geq\sheaf{\sheaf{c}}$, since
$\card(U_{\sheaf{\sheaf{c}}})\geq
\card(A_{\sheaf{\sheaf{d}}})=2^{\sheaf{\sheaf{d}}}$ for every
$\sheaf{\sheaf{d}}<\sheaf{\sheaf{c}}$.


The case $\sheaf{\sheaf{c}}=0$ being trivial, there remains the case
where $\sheaf{\sheaf{c}}$ is the \emph{countably infinite cardinal.}
In this case $U_{\sheaf{\sheaf{c}}}$ is the set of Artinian sets of
finite type (\ie finite together with all their components), and one
uses a pleasant result of a combinatorial nature.


% END INCLUDED SGA4_English_current_buildable_source_20260906/source/01/drafts/SGA4_I_LINES_8643_8705_English.texfrag


% BEGIN INCLUDED SGA4_English_current_buildable_source_20260906/source/01/drafts/SGA4_I_LINES_8706_8780_English.texfrag
\hypertarget{sga4.exp.I.appendix.II.sec.6.lemma.3}{}
\begin{enonce*}{LEMMA 3}
{\rm (D. K"onig?)}. Consider two infinite sequences
$(E_n)_{n\geq1}$, $(f_n)_{n\geq1}$ of finite sets $E_n$ and maps
$f_n:E_{n+1}\rightarrow E_n$. If there exists no infinite sequence
$(x_n)_{n%%
\geq1}$ such that $x_n\in E_n$ and $f_n(x_{n+1})=x_n$ for every
$n$, then $E_n$ is empty for all sufficiently large $n$.
\end{enonce*}

In other words, if all sequences $(x_n)$ such that
$f_n(x_{n+1})=x_n$ are \emph{finite}, their lengths are
\emph{bounded}. This can be expressed in terms of projective limits: a
projective limit of nonempty finite sets is nonempty (cf. General
Topology, Chapter I, 2nd ed., \S 9, ${\rm n^{\circ}}$ 6,
Proposition 8, $2^{\circ}$)).

Indeed,\pageoriginale argue by contradiction. If there exist nonempty
$E_n$ for arbitrarily large $n$, no $E_n$ is empty (since
$E_n=∅$ implies $E_{n+1}=∅$ in view of the existence of
$f_n:E_{n+1}\rightarrow E_n$). Call a finite sequence
$(x_j)_{1\leq j\leq n}$ \emph{coherent} if
$f_j(x_{j+1})=x_j$ for $j=1,\ldots,n-1$. We prove by induction on $n$
the existence of a coherent sequence $(a_1,\ldots,a_n)$
$(a_i\in E_i)$ that, for every $q\geq n$, can be extended to a
coherent sequence $(a_1,\ldots,a_n,x_{n+1},\ldots,x_q)$ of length
$q$. This is evident for $n=0$, since no $E_q$ is empty. Pass from $n$
to $n+1$. If, for every $x\in f_n^{-1}(\{a_n\})\subset E_{n+1}$,
all coherent sequences extending $(a_1,\ldots,a_n,x)$ had lengths
bounded by an integer $q(x)$, then all coherent sequences extending
$(a_1,\ldots,a_n)$ would have bounded lengths (by $\sup_xq(x)$), since
$E_{n+1}$ is finite; hence there exists
$a_{n+1}\in f_n^{-1}(\{a_n\})$ such that the coherent sequence
$(a_1,\ldots,a_n,a_{n+1})$ admits extensions of arbitrary length. We
then obtain an infinite sequence $(a_n)_{n%%
\geq1}$ contradicting the hypothesis.

It follows from \hyperlink{sga4.exp.I.appendix.II.sec.6.lemma.3}{Lemma 3} that if $A$ is an Artinian set of finite type,
then the ordered set $G(A)$ of its chains is \emph{finite}: indeed,
take for $E_n$ the set of chains
$(x_n\in x_{n-1}\in\ldots\in x_1\in A)$ with $n+1$ terms (which is
finite because the components of $A$ of order $\leq n+1$ are finite
in number and are all finite), and for $f_n$ the map
$(x_{n+1}\in x_n\in x_{n-1}\in\ldots)\mapsto
(x_n\in x_{n-1}\in\ldots)$. Now the set of isomorphism classes of
finite\pageoriginale ordered sets is \emph{countable}: indeed, giving
an order structure on a finite subset of $\NN$ is equivalent to giving
its graph, which is a finite subset of $\NN\times\NN$; on the other
hand, the set of finite subsets of a countable set is countable
(Chapter III). It follows from \hyperlink{sga4.exp.I.appendix.II.sec.6.theorem.2.proof.claim.star}{(*)} and \hyperlink{sga4.exp.I.appendix.II.sec.6.lemma.1}{Lemma 1} that the set
$\mathfrak{a}$ of Artinian sets of finite type is countable. It is
infinite by \hyperlink{sga4.exp.I.appendix.II.sec.6.lemma.2}{Lemma 2} (or, more simply, by using the sequence
$(z_n)_{n\geq0}$ defined by $z_0=∅$, $z_{n+1}=\{z_n\}$). This
completes the proof of b).

\begin{enonce*}[remark]{Remark}
We have just proved that if $A$ is an Artinian set of finite type, it
has only finitely many components. There therefore exists an integer
$n$ such that $A$ is of type $n$.
\end{enonce*}
We turn to the proof of c). Let $\sheaf{\sheaf{c}}$ be an infinite
cardinal, let $U$ be a universe admitting an element of cardinality
$\sheaf{\sheaf{c}}$, and let $A$ be an Artinian set of type
$\sheaf{\sheaf{c}}$. Then the ordered set $G(A)$ has cardinality
$\leq\sheaf{\sheaf{c}}$ (\hyperlink{sga4.exp.I.appendix.II.sec.6.theorem.2.proof.equation.2}{formula (2)} above). The second assertion of
\hyperlink{sga4.exp.I.appendix.II.sec.6.lemma.1}{Lemma 1}, applied to $G(A)$, shows that the map
$(x_0,\ldots,x_n)\rightsquigarrow x_n$ from $G(A)$ takes its values in
$U$. In other words, all the components of $A$ are elements of $U$.
This proves c). The proof of \hyperlink{sga4.exp.I.appendix.II.sec.6.theorem.2.part.c-prime}{${\rm c'}$)} is analogous: if $A$ is an
Artinian set of finite type, we have just seen that $G(A)$ is finite;
since $U$ is nonempty, it contains an element equipotent to $G(A)$ by
the \hyperlink{sga4.exp.I.appendix.II.sec.1.corollary.to.proposition.7}{corollary to Proposition 7} (\hyperlink{sga4.exp.I.appendix.II.sec.1}{${\rm n^{\circ}}$ 1}). \cqfd


% END INCLUDED SGA4_English_current_buildable_source_20260906/source/01/drafts/SGA4_I_LINES_8706_8780_English.texfrag


% BEGIN INCLUDED SGA4_English_current_buildable_source_20260906/source/01/drafts/SGA4_I_LINES_8781_8916_English.texfrag
\section{Vague metamathematical remarks}%
\hypertarget{sga4.exp.I.appendix.II.sec.7}{}\pageoriginale

\begin{enumerate}
\renewcommand{\theenumi}{\alph{enumi}}
\renewcommand{\labelenumi}{\theenumi)}
\item The axiom ``\emph{every set is Artinian}'' is
\emph{harmless}. Indeed, if one has a model $M$ of Set Theory, the set
$M'$ of the Artinian elements of $M$ is also a model (\cf \hyperlink{sga4.exp.I.appendix.II.sec.6.proposition.11}{Proposition
11}).

\item The universe axiom \hyperlink{sga4.exp.I.appendix.II.sec.4.axiom.A.6}{(A.6)} (and the equivalent axiom \hyperlink{sga4.exp.I.appendix.II.sec.5.axiom.A-prime.6}{$(A'.6)$} on
strongly inaccessible cardinals) is \emph{independent} of the rest of
Set Theory. Indeed, let $\sheaf{\sheaf{c}}$ be the first uncountable
strongly inaccessible cardinal. The universe
$U_{\sheaf{\sheaf{c}}}$ of Artinian sets of strict type
$\sheaf{\sheaf{c}}$ (\hyperlink{sga4.exp.I.appendix.II.sec.6.theorem.2}{Theorem 2}, \hyperlink{sga4.exp.I.appendix.II.sec.6.theorem.2.part.b}{b)}) is a model of Set Theory without
\hyperlink{sga4.exp.I.appendix.II.sec.4.axiom.A.6}{(A.6)}: the elements of $U_{\sheaf{\sheaf{c}}}$ are called ``sets,''
the ``membership relation'' is the restriction to
$U_{\sheaf{\sheaf{c}}}$ of the usual one, \etc The ``universes'' of
the model are therefore the ordinary universes that are elements of
$U_{\sheaf{\sheaf{c}}}$. But we have seen that the only universes
which are elements of $U_{\sheaf{\sheaf{c}}}$ are the two
peque${\rm\widetilde{n}}$os $U_0=∅$ and $U_1$. Thus $U_1$ is a
``set'' which is not an element of any ``universe.'' We therefore
have a model of Set Theory in which \hyperlink{sga4.exp.I.appendix.II.sec.4.axiom.A.6}{(A.6)} is false.

\item Bourbaki was too cautious in merely ``presuming'' that the
\emph{axiom of infinity} (A.5) is \emph{independent} of the preceding
axioms and schemes. It is indeed independent, since the countable
universe $U_1$ of Artinian sets of finite type is a model in which
(A.5) is false and the preceding axioms and schemes are true.

\item It would be very interesting to prove that the universe axiom
\hyperlink{sga4.exp.I.appendix.II.sec.4.axiom.A.6}{(A.6)} is harmless. That looks difficult---and in fact it is
unprovable, says Paul Cohen.
\end{enumerate}

The adjective\pageoriginale ``vague'' in the title means that, when
constructing models, we did not take the trouble to fiddle with the
symbol $\tau$ so that it could not take us outside the model. The
trick, if one is considering a model $M$, is to replace the ordinary
$\tau_x(ℛ(x))$ by:
$$
\tau_x(ℛ(x)\mbox{ and } x\in M){\quoi}.
$$
One still has to check that this really turns the ordinary
quantifiers into the quantifiers formerly called ``typical'':
$$
(\forall x \in M)\quav (\exists x \in M){\quoi}.
$$

\hypertarget{sga4.exp.I.appendix.II.sec.7.exercises}{}
\begin{enonce*}[remark]{Exercises}
\begin{enumerate}
\renewcommand{\labelenumi}{\theenumi)}
\hypertarget{sga4.exp.I.appendix.II.sec.7.exercise.1}{}
\item Let $n$ be an integer $\geq 1$. Show that the set of Artinian
sets of type $n$ is infinite (the sets $z_0=∅$, $z_1=\{∅\}$,
$z_{q+1}=\{z_q\}$ are of type 1).

\hypertarget{sga4.exp.I.appendix.II.sec.7.exercise.2}{}
\item Define the \emph{height} of a set $A$ to be the supremum
(finite or infinite) of the integers $n$ for which there exists a
sequence $x_n\in x_{n-1}\in \ldots \in x_0=A$. Show that the sets of
height $\leq n$ are finite and form a finite set, whose cardinality
$p_n$ can be calculated from $p_0=1$, $p_{n+1}=2^{p_n}$ (proceed by
induction on $n$, noting that the elements of a set $A$ of height
$\leq n$ are sets of height $\leq n-1$).

\hypertarget{sga4.exp.I.appendix.II.sec.7.exercise.3}{}
\item Let $G$ be a Noetherian ordered set having a least element
$g_0$, and suppose that, for every $h\in G$, the set of $g\in G$ such
that $g\leq h$ is finite and totally ordered.
\begin{enumerate}
\renewcommand{\theenumii}{\alph{enumii}}
\renewcommand{\labelenumii}{\theenumii)}
\hypertarget{sga4.exp.I.appendix.II.sec.7.exercise.3.part.a}{}
\item Show\pageoriginale that every element $h\neq g_0$ of $G$ has
a unique predecessor.

\hypertarget{sga4.exp.I.appendix.II.sec.7.exercise.3.part.b}{}
\item Let $\varphi$ be the map on $G$ defined in \hyperlink{sga4.exp.I.appendix.II.sec.6.lemma.1}{Lemma 1} (\ie
$\varphi(g)$ is the set of all $\varphi(g')$ as $g'$ ranges over the
successors of $g$). Put $A=\varphi(g_0)$. Show that $A$ is an Artinian
set. For $g\in G$, let $g_0<g_1<\ldots<g_n=g$ be the sequence of
elements $\leq g$; put
$f(g)=(\varphi(g_0),\varphi(g_1),\ldots,\varphi(g_n))$; show that $f$
is an increasing surjection from $G$ onto the ordered set $G(A)$ of
the text. Show that, if $f$ is \emph{injective}, it is an
\emph{isomorphism} from $G$ onto $G(A)$.

\hypertarget{sga4.exp.I.appendix.II.sec.7.exercise.3.part.c}{}
\item For $g\in G$, let $M_g$ be the set of upper bounds of $g$.
Suppose that, for every pair of distinct elements $g$, $g'$ having
the same predecessor $p$, the ordered sets $M_g$ and $M_{g'}$ are
not isomorphic. Show that the map $f$ of \hyperlink{sga4.exp.I.appendix.II.sec.7.exercise.3.part.b}{b)} is then an
\emph{isomorphism} from $G$ onto $G(A)$. [If $f(g)=f(g')$ with
$g\neq g'$, and if $g_0<g_1<\ldots<g_n=g$ and
$g'_0<g'_1<\ldots <g'_{n'}=g'$ are the sequences of elements
$\leq g$ and $\leq g'$, respectively, show that $n=n'$, and that one
may assume that $g$ and $g'$ have the same predecessor $p$; then
consider the set of $p\in G$ for which there exist two distinct
successors $g$, $g'$ of $p$ such that $f(g)=f(g')$, a maximal element
$q$ of this set, and two distinct successors $h$, $h'$ of $q$ such
that $f(h)=f(h')$; note that the restrictions of $f$ to $M_h$ and
$M_{h'}$ are injective and hence (by \hyperlink{sga4.exp.I.appendix.II.sec.7.exercise.3.part.b}{b)}) are isomorphisms from $M_h$
onto $G(\varphi(h))$ and from $M_{h'}$ onto $G(\varphi(h'))$; deduce
from the hypothesis that $M_h$ and $M_{h'}$ are not isomorphic that
$\varphi(h)\neq \varphi(h')$, which contradicts $f(h)=f(h')$].
\end{enumerate}
N.B.\pageoriginale \hyperlink{sga4.exp.I.appendix.II.sec.7.exercise.3}{Exercise 3}) gives very precise information on how
Artinian sets are ``made.'' We already knew that such a set $A$ is
determined by the isomorphism class of the ordered set $G(A)$ (\hyperlink{sga4.exp.I.appendix.II.sec.6.lemma.1}{Lemma
1}). We now know how to \emph{characterize} the ordered sets $G$
isomorphic to some $G(A)$:
\begin{itemize}
\item[$\alpha)$] $G$ is Noetherian and has a least element;

\item[$\beta)$] For every $g\in G$, the set of $h\in G$ such that
$h\leq g$ is totally ordered and finite (whence the existence and
uniqueness of the predecessor of every nonleast $g$);

\item[$\gamma)$] If $g$, $g'$ are distinct elements having the same
predecessor, the set $M_g$ of upper bounds of $g$ and the set
$M_{g'}$ of upper bounds of $g'$ are not isomorphic.
\end{itemize}

\hypertarget{sga4.exp.I.appendix.II.sec.7.exercise.4}{}
\item Let $(A_n)_{n\geq 0}$ be the sequence of finite \emph{ordinals}
defined by $A_0=∅$, $A_{n+1}=A_n\cup \{A_n\}$. Show that the
ordered set $G(A_n)$ has $2^n$ elements, and that the number of its
elements of height $q$ (in the sense of \hyperlink{sga4.exp.I.appendix.II.sec.7.exercise.2}{Exercise 2}) is
$\binom{n}{q}$.
\end{enumerate}
\end{enonce*}

% END INCLUDED SGA4_English_current_buildable_source_20260906/source/01/drafts/SGA4_I_LINES_8781_8916_English.texfrag


% BEGIN INCLUDED SGA4_English_current_buildable_source_20260906/source/01/drafts/SGA4_I_LINES_8917_8927_English.texfrag
\makeatletter
\let\toclevel@section\sgaFourSavedToclevelSection
\makeatother
\pdfbookmark[2]{Bibliography}{sga4.exp.I.appendix.II.bibliography}
\begingroup
\makeatletter
\def\Hy@writebookmark#1#2#3#4#5{}
\makeatother
\begin{thebibliography}{10}%%
\bibitem[1]{sga4.exp.I.bib.I.2} \sisi{B. \textsc{Mitchell}, Theory of categories, Academic Press
(1965).}{B. \textsc{Mitchell}, \textit{Theory of categories},
Academic Press, 1965.}
\end{thebibliography}
\endgroup
\hypersetup{bookmarksdepth=1}

\ifx\danslelivre\undefined
%\bibliography{../fichiers-configuration-latex/bibliographie}
%\bibliographystyle{../fichiers-configuration-latex/smfalpha}
\end{document}
\fi

% END INCLUDED SGA4_English_current_buildable_source_20260906/source/01/drafts/SGA4_I_LINES_8917_8927_English.texfrag


% END INCLUDED SGA4_English_current_buildable_source_20260906/source/01/01_English.tex


\makeatletter
\def\input@path{{../02/}}
\makeatother
% BEGIN INCLUDED SGA4_English_current_buildable_source_20260906/source/02/02_English.tex
\ifx\danslelivre\undefined
% Release engine: XeLaTeX (the LuaTeX compatibility layer breaks Xy-pic output).
\documentclass[10pt]{../fichiers-configuration-latex/smfart}
\usepackage{ifthen}
% BEGIN INCLUDED SGA4_English_current_buildable_source_20260906/source/fichiers-configuration-latex/paquets.tex
\usepackage{polyglossia}
%\usepackage[francais]{babel}
\usepackage{fontspec}
\usepackage[math-style=ISO,bold-style=upright,sans-style=upright]{unicode-math}
%\usepackage[mathletters]{ucs}
%\usepackage[utf8]{inputenc}
%\usepackage[T1]{fontenc}
\usepackage[unicode,hypertexnames=false,colorlinks=true,linkcolor=blue,citecolor=Sepia,urlcolor=magenta]{hyperref}
\setmainlanguage{french}
\setmainfont[Ligatures=TeX]{Linux Libertine O}
% variante possible : TeXGyrePagella ?
\setmathfont{XITS Math}

% maths
\usepackage{amsmath}
%\usepackage{amssymb}
%\usepackage{mathrsfs}
%\usepackage{stmaryrd}
%\usepackage{euscript}

% couleurs
\usepackage[usenames,dvipsnames]{xcolor}

% graphiques
\usepackage{graphicx}
\usepackage[all]{xy}
%\usepackage{tikz}
%\usetikzlibrary{matrix,arrows,calc}
%\tikzset{column sep=3em,row sep=2em,text height=1.5ex, text depth=0.25ex,>=angle 90}
%\tikzset{description/.style={fill=white,inner sep=1pt}}

% hyperliens, 'etiquettes
\usepackage{xr-hyper}
\usepackage{../fichiers-configuration-latex/smfthm}
%\ifx\cacheretiquettes\undefined
%\usepackage[inline]{showlabels} \else \fi % pour rendre visibles les 'etiquettes
%\usepackage{../fichiers-configuration-latex/makeidx} % pour faire un index
%\usepackage[pagebackref,unicode]{hyperref}  % pour voir o`u sont cit'ees les r'ef'erences
% hyperref is loaded once above with the cumulative-reader destination policy.
%\hypersetup{breaklinks=true} %permet le retour `a la ligne dans les liens trop longs. Incompatible avec pdflatex
%\renewcommand*{\backref}[1]{{\small [cit\'e page(s)~#1.]}}

\usepackage{xspace}
\usepackage{pifont} % pour expose XVII

% END INCLUDED SGA4_English_current_buildable_source_20260906/source/fichiers-configuration-latex/paquets.tex

% BEGIN INCLUDED SGA4_English_current_buildable_source_20260906/source/fichiers-configuration-latex/numerotation.tex
\renewcommand{\thefootnote}{\roman{footnote}} % pour mettre les notes en bas de page sous la forme (ii) et non (2), qui ressemble trop (?) a un exposant.

\newcounter{pageoriginal}
\newcommand{\pageoriginale}{\refstepcounter{pageoriginal}\marginpar{\footnotesize\textbf{\thepageoriginal}}\xspace}
\let\pageoriginaled\pageoriginale

\usepackage{ifthen}
\newboolean{original}
\def\sisi{\ifthenelse{\boolean{original}}}
\setboolean{original}{false}

\newcounter{cnote}
\def\nde#1{\ifthenelse{\boolean{original}}{}{\stepcounter{cnote}\def\thefootnote{\arabic{cnote}}\protect\footnote[\value{cnote}]{N.D.E.: #1}}}
\def\ndemark{\ifthenelse{\boolean{original}}{}{\stepcounter{cnote}\def\thefootnote{\arabic{cnote}}\protect\footnotemark[\value{cnote}]}}
\def\ndetext#1{\ifthenelse{\boolean{original}}{}{\def\thefootnote{\arabic{cnote}}\protect\footnotetext[\value{cnote}]{N.D.E.: #1}}}

% END INCLUDED SGA4_English_current_buildable_source_20260906/source/fichiers-configuration-latex/numerotation.tex

% BEGIN INCLUDED SGA4_English_current_buildable_source_20260906/source/fichiers-configuration-latex/macros.tex
\setmainlanguage{french}
\setotherlanguage{english}


%% Fixed, engine-independent release identifier.  The inherited dynamic Git
%% probe used Lua io.popen, exposed build-state-dependent output, and printed
%% a visible question mark outside a Git checkout.
\def\showgitstatus{English release r2, 2026-07-23}


% macros pour LaTeX
\renewcommand{\mathsf}{}

% Police différente pour U+2113 SCRIPT SMALL L
\setmathfont[range="2113]{latinmodern-math.otf}

% Give every inherited semantic label its own stable PDF destination instead
% of inheriting whichever automatic equation/page anchor happened to be live.
% The legacy label remains as a compatibility alias used by existing edges.
\makeatletter
\def\etiquette[#1]#2{%
  \Hy@raisedlink{\hyper@anchorstart{sga4.semantic.#2}\hyper@anchorend}%
  \def\@currentHref{sga4.semantic.#2}%
  \def\@currentlabel{#1}%
  \label{#2}%
}
\makeatother
\def\eqhyperref[#1]#2{(\hyperref[#1]{#2})}
\def\pagehyperref[#1]#2{(page \hyperref[#1]{#2})}
\def\hyperrefext[#1][#2]#3{\hyperref[#2]{#1 #3}}
\def\eqhyperrefext[#1][#2]#3{\hyperref[#2]{#1 #3}}     %cf. exposes 17
\def\neweqhyperrefext[#1][#2]#3{\hyperref[#2]{#1 #3}}  %cf. exposes 17

\def\shat#1{\protect\widehat{#1}}
\def\stilde#1{\protect\widetilde{#1}}
%%%%%%%%%%%%%%%%%%


\def\etcn{etc.\xspace}
\def\etco{etc...\xspace}
\def\etc{\sisi{\etco}{\etcn}}
\def\ldotsn{\sisi{...}{\etcn}}

\def\cf{{cf.}\kern.3em}
\def\Cf{{Cf.}\kern.3em}
\def\ie{{i.e.}\kern.3em}
\def\iev{{i.e.},\xspace}
\def\loccito{loc.\kern2pt cit.\xspace}
\def\loccitn{\textit{loc.\kern2pt cit.}\xspace}
\def\loccit{\sisi{\loccito}{\loccitn}}
\def\EGA{\textup{EGA}\kern.2em}
\def\SGA{\textup{SGA}\kern.3em}

\def\quoi{}

\def\quav{\quad,\quad}
\def\quaet{\quad \text{et}\quad}
\def\fibreproduct#1#2#3{#1 ×_{#3} #2}
\let\fprod\fibreproduct
\def\ccheck#1#2{\mbox{${\check{\check{\rm #1}}}^{\lower2pt\hbox{$\scriptstyle #2$}}$}}
\def\sheaf#1{\protect\mathcal{#1}}
\def\vieuxsheaf#1{{#1}} % cf. problème avec sheaf{a\Phi_i} par exemple...
\DeclareMathOperator{\Hom}{Hom}
\DeclareMathOperator{\can}{can}
\DeclareMathOperator{\Ind}{Ind}
\DeclareMathOperator{\ind}{ind}
%\DeclareMathOperator{\Hhomtop}{\mathcal{H}\!\mathit{omtop}}
\DeclareMathOperator{\Hhomtop}{\mathbf{Homtop}}
\DeclareMathOperator{\Point}{Point}
\DeclareMathOperator{\Pro}{Pro}
\DeclareMathOperator{\Top}{Top}
\DeclareMathOperator{\Mod}{Mod}
\DeclareMathOperator{\Cart}{Cart}
%\DeclareMathOperator{\Mmorsite}{\mathcal{M}\!\mathit{orsite}}
\DeclareMathOperator{\Mmorsite}{\mathbf{Morsite}}
\DeclareMathOperator{\loc}{loc}
\DeclareMathOperator{\Loc}{Loc}
\DeclareMathOperator{\Pic}{Pic}
\DeclareMathOperator{\Et}{Et}
\DeclareMathOperator{\card}{card}
\DeclareMathOperator{\Ens}{Ens}
\DeclareMathOperator{\h}{\underline{H}}
\DeclareMathOperator{\Fonct}{Fonct}
\DeclareMathOperator{\Ker}{Ker}
\DeclareMathOperator{\but}{but}
\DeclareMathOperator{\Ob}{Ob}
\DeclareMathOperator{\ob}{ob}
\DeclareMathOperator{\Lib}{Lib}
\DeclareMathOperator{\esj}{esj}
\DeclareMathOperator{\id}{id}
\DeclareMathOperator{\Id}{Id}
\DeclareMathOperator{\ad}{ad}
\DeclareMathOperator{\Iim}{Im}
\DeclareMathOperator{\pr}{pr}
\DeclareMathOperator{\Cat}{Cat}
\DeclareMathOperator{\car}{car}
\DeclareMathOperator{\SysInd}{SysInd}
\DeclareMathOperator{\Sup}{Sup}
\DeclareMathOperator{\Coker}{Coker}
\DeclareMathOperator{\Filt}{Filt}
\DeclareMathOperator{\acc}{acc}
\DeclareMathOperator{\sur}{sur}
\DeclareMathOperator{\cat}{cat}
\DeclareMathOperator{\Ord}{Ord}
\DeclareMathOperator{\Gr}{Gr}
\DeclareMathOperator{\Cov}{Cov}
\DeclareMathOperator{\Caf}{Caf}
\DeclareMathOperator{\coker}{coker}
\DeclareMathOperator{\coim}{coim}
\DeclareMathOperator{\im}{im}
\DeclareMathOperator{\Ouv}{Ouv}
\DeclareMathOperator{\diag}{diag}
\DeclareMathOperator{\Esp}{Esp}
\DeclareMathOperator{\TOP}{TOP}
\DeclareMathOperator{\Unif}{Unif}
\DeclareMathOperator{\Iind}{\mathbf{Ind}}
\DeclareMathOperator{\Hhomcart}{\mathbf{Homcart}}
\DeclareMathOperator{\Hhomtopess}{\mathbf{Homtopess}}
\DeclareMathOperator{\Ppoints}{\mathbf{Points}}
\DeclareMathOperator{\Ppoint}{\mathbf{Point}}
\DeclareMathOperator{\Ttors}{\mathbf{Tors}}
\DeclareMathOperator{\Ssex}{\mathbf{Sex}}
\DeclareMathOperator{\Ffib}{\mathbf{Fib}}
\DeclareMathOperator{\Vvois}{\mathbf{Vois}}
\DeclareMathOperator{\Ppointess}{\mathbf{Pointess}}
\DeclareMathOperator{\scrHom}{\mathbf{Hom}}
\DeclareMathOperator{\sheafPic}{\mathbf{Pic}}
\DeclareMathOperator{\sheafAut}{\mathbf{Aut}}
\DeclareMathOperator{\Hhomtoplocan}{\mathbf{Homtoplocan}}
\DeclareMathOperator{\Hhomtopan}{\mathbf{Homtopan}}
\DeclareMathOperator{\Pptgeom}{\mathbf{Ptgéom}}
\DeclareMathOperator{\Morsite}{Morsite}
\DeclareMathOperator{\Fib}{Fib}
\DeclareMathOperator{\Fais}{Fais}
\DeclareMathOperator{\kar}{kar}
\DeclareMathOperator{\Groupes}{Groupes}
\DeclareMathOperator{\Restr}{Restr}
\DeclareMathOperator{\Isom}{Isom}
\DeclareMathOperator{\Prol}{Prol}
\DeclareMathOperator{\Res}{Res}
\DeclareMathOperator{\Tors}{Tors}
\DeclareMathOperator{\esp}{esp}
\DeclareMathOperator{\es}{es}
\DeclareMathOperator{\od}{od}
\DeclareMathOperator{\gr}{gr}
\DeclareMathOperator{\sob}{sob}
\DeclareMathOperator{\dfn}{dfn}
\DeclareMathOperator{\Sch}{Sch}
\DeclareMathOperator{\ouv}{ouv}
\DeclareMathOperator{\cons}{cons}
\DeclareMathOperator{\Nerf}{Nerf}
\DeclareMathOperator{\Points}{Points}
\DeclareMathOperator{\Espsob}{Espsob}
\DeclareMathOperator{\Homtop}{Homtop}
\DeclareMathOperator{\Pt}{Pt}
\DeclareMathOperator{\Quot}{Quot}
\DeclareMathOperator{\glob}{glob}
\DeclareMathOperator{\Homfib}{Homfib}
\DeclareMathOperator{\Qouv}{Qouv}
\DeclareMathOperator{\Quotouv}{Quotouv}
\DeclareMathOperator{\ens}{ens}
\DeclareMathOperator{\ab}{ab}
\DeclareMathOperator{\an}{an}
\DeclareMathOperator{\fs}{fs}
\DeclareMathOperator{\red}{red}
\DeclareMathOperator{\Ab}{Ab}
\DeclareMathOperator{\cosq}{cosq}
\DeclareMathOperator{\Ann}{Ann}
\DeclareMathOperator{\Zar}{Zar}


%%%%%%%%%%%%
%%%%% SGA4-2
\DeclareMathOperator{\Sex}{Sex}
\DeclareMathOperator{\vois}{vois}
\DeclareMathOperator{\Vois}{Vois}
\DeclareMathOperator{\Ext}{Ext}
\DeclareMathOperator{\ext}{ext}
\DeclareMathOperator{\Fer}{Fer}
\DeclareMathOperator{\Add}{Add}
\DeclareMathOperator{\Cocart}{Cocart}
\DeclareMathOperator{\cosk}{cosk}
\DeclareMathOperator{\plat}{plat}

\DeclareMathOperator{\faisc}{faisc}
\DeclareMathOperator{\sk}{sk}
\DeclareMathOperator{\ch}{ch}

\def\HH{\mathbf{H}}
\def\qc{\mathrm{qc}}
\def\et{\mathrm{et}}
\def\coh{\mathrm{coh}}

\def\Hvr#1{%
\begin{array}[b]{l}
{\kern2pt\scriptstyle r}\\[-8pt]
{\kern1pt\scriptstyle \vee}\\[-5pt]
H^{#1}
\end{array}}

\def\Hvv#1{%
\begin{array}[b]{l}
\kern2pt\rotatebox{90}{${\scriptscriptstyle\ll}$}\\[-6pt]
H^{#1}
\end{array}}

\def\Hvi#1{%
\begin{array}[b]{l}
{\kern0.2pt\scriptscriptstyle \infty}\\[-8pt]
{\kern1pt\scriptstyle \vee}\\[-5pt]
H^{#1}
\end{array}}

\def\Hv#1{%
\begin{array}[b]{l}
{\kern1pt\scriptstyle \vee}\\[-5pt]
H^{#1}
\end{array}}

\DeclareMathOperator{\cart}{cart}
\DeclareMathOperator{\Ffaisc}{\mathbf{Faisc}}


\def\Ev#1{%
\begin{array}[b]{l}
{\kern1pt\scriptstyle \vee}\\[-5pt]
E'_{#1}
\end{array}}

\DeclareMathOperator{\Cont}{Cont}
\DeclareMathOperator{\Ccont}{\mathscr{C}\!\mathit{ont}}
\DeclareMathOperator{\Site}{Site}
\newcommand\MM{\mathbf{M}}
\def\limtop#1{{\displaystyle{\varprojlim_{#1}\!\mbox{\rm top }}}}
\def\Limtop{\mathop{\varprojlim\!\mbox{\rm top}}}
\DeclareMathOperator{\spec}{spec}
\def\ss{\scriptstyle{\mathscr{S}}}
\DeclareMathOperator{\Gal}{Gal}
\def\smallm{\scriptstyle{\mathscr{M}}}
\DeclareMathOperator{\End}{End}
\DeclareMathOperator{\ass}{ass}
\DeclareMathOperator{\zar}{zar}
\DeclareMathOperator{\Ogg}{Ogg}
\DeclareMathOperator{\res}{res}
\DeclareMathOperator{\tr}{tr}
\DeclareMathOperator{\Supp}{Supp}
\DeclareMathOperator{\codim}{codim}
\DeclareMathOperator{\rk}{rk}
\DeclareMathOperator{\cur}{cur}
\DeclareMathOperator{\cdqc}{cdqc}
\newcommand\DD{\mathbf{D}}
\DeclareMathOperator{\cl}{cl}
\def\JJ{\mathbf{J}}
\def\FF{\mathbf{F}}
\def\BB{\mathbf{B}}
\DeclareMathOperator{\Alget}{Alget}
\DeclareMathOperator{\Ass}{Ass}
\DeclareMathOperator{\supp}{supp}
\DeclareMathOperator{\rad}{rad}
\DeclareMathOperator{\prof}{prof}
\DeclareMathOperator{\cint}{int}
\DeclareMathOperator{\sic}{sic}
\DeclareMathOperator{\Tor}{Tor}
\DeclareMathOperator{\tors}{tors}
\DeclareMathOperator{\Sym}{Sym}
\DeclareMathOperator{\parf}{parf}
\DeclareMathOperator{\tor}{tor}
\DeclareMathOperator{\defn}{defn}
\DeclareMathOperator{\Div}{Div}
\DeclareMathOperator{\Pol}{Pol}
\DeclareMathOperator{\Fex}{Fex}
\DeclareMathOperator{\Lie}{Lie}
\DeclareMathOperator{\tot}{tot}
\DeclareMathOperator{\Aut}{Aut}
\DeclareMathOperator{\HOM}{HOM}
\DeclareMathOperator{\pch}{pch}
\DeclareMathOperator{\Ch}{Ch}
\DeclareMathOperator{\EXT}{EXT}
\DeclareMathOperator{\PIC}{PIC}
\DeclareMathOperator{\Rf}{Rf}
\DeclareMathOperator{\inj}{inj}
\newcommand\VV{\mathbf{V}}

\def\Xv{%
\begin{array}[b]{l}
{\kern1pt\scriptstyle \vee}\\[-5pt]
X
\end{array}}

%%%%%%%
\newcommand\TT{\mathbf{T}}
\newcommand\XX{\mathbf{X}}
\newcommand\EE{\mathbf{E}}
%%%%%%

\def\qedsymbol{\textup{C.Q.F.D.}}
\let\CQFD\qed
\let\cqfd\qed

%%%%%\bigdot
\font\bigsymb=cmsy10 at 4pt
\def\bigdot{{\kern1.2pt\raise 1.5pt\hbox{\bigsymb\char15}}}
%%%%%

\def\Lright{\mathop{\vcenter{\hbox{\hbox{$\longrightarrow$}\lower8pt \hbox{\kern-16pt\hbox{$\longrightarrow$}}}}}}

\def\Aa{{\mathcal A}}
\def\Bb{{\mathcal B}}
\def\Cc{{\mathcal C}}
\def\Dd{{\mathcal D}}
\def\Ee{{\mathcal E}}
\def\Ff{{\mathcal F}}
\def\Gg{{\mathcal G}}
\def\Hh{{\mathcal H}}
\def\Ii{{\mathcal I}}
\def\Jj{{\mathcal J}}
\def\Kk{{\mathcal K}}
\def\Ll{{\mathcal L}}
\def\Mm{{\mathcal M}}
\def\Nn{{\mathcal N}}
\def\Oo{{\mathcal O}}
\def\Pp{{\mathcal P}}
\def\Qq{{\mathcal Q}}
\def\Rr{{\mathcal R}}
\def\Ss{{\mathcal S}}
\def\Tt{{\mathcal T}}
\def\Uu{{\mathcal U}}
\def\Vv{{\mathcal V}}
\def\Ww{{\mathcal W}}
\def\Xx{{\mathcal X}}
\def\Yy{{\mathcal Y}}
\def\Zz{{\mathcal Z}}

\newcommand{\CC}{\mathbf{C}}
\newcommand{\GG}{\mathbf{G}}
\newcommand{\KK}{\mathbf{K}}
\newcommand{\LL}{\mathbf{L}}
\newcommand{\NN}{\mathbf{N}}
\newcommand{\PP}{\mathbf{P}}
\newcommand\QQ{\mathbf{Q}}
\newcommand\RR{\mathbf{R}}
\newcommand\ZZ{\mathbf{Z}}

\def\to{\mathchoice{\longrightarrow}{\rightarrow}{\rightarrow}{\rightarrow}}
\def\from{\mathchoice{\longleftarrow}{\leftarrow}{\leftarrow}{\leftarrow}}
\def\mto{\mathchoice{\longmapsto}{\mapsto}{\mapsto}{\mapsto}}
\def\mfrom{\mathrel{\from\joinrel\mapstochar\,}}
\def\hto{\mathrel{\lhook\joinrel\to}}
\def\hfrom{\mathrel{\from\joinrel\rhook}}

\def\implies{\mathchoice{\Longrightarrow}{\Rightarrow}{\Rightarrow}{\Rightarrow}}
\def\implied{\mathchoice{\Longleftarrow}{\Leftarrow}{\Leftarrow}{\Leftarrow}}
\def\iff{\mathchoice{\Longleftrightarrow}{\Leftrightarrow}{\Leftrightarrow}{\Leftrightarrow}}

\def\To#1{\mathchoice{\xrightarrow{\labelstyle~~#1~}}{\stackrel{#1}{\longrightarrow}}{}{}}
\def\Mto#1{\mathrel{\mapstochar\To{#1}}}
\def\From#1{\mathchoice{\xleftarrow{\labelstyle~~#1~}}{\stackrel{#1}{\longleftarrow}}{}{}}

\newcommand{\isomto}{\overset{\sim}{\longrightarrow}}
\newcommand{\isomfrom}{\overset{\sim}{\longleftarrow}}

\newcommand{\IFF}{\unskip~$\iff$~\ignorespaces}
\newcommand{\THEN}{\unskip~$\implies$~\ignorespaces}

\let\leq\leqslant
\let\geq\geqslant
\let\le\leqslant
\let\ge\geqslant

\def\cf{{cf.}\kern.3em}
\def\Cf{{Cf.}\kern.3em}
\def\ie{{i.e.}\kern.3em}
\def\iev{{i.e.},\xspace}
%\def\loccito{loc.\kern2pt cit.\xspace}
%\def\loccitn{\textit{loc.\kern2pt cit.}\xspace}
%\def\loccit{\sisi{\loccito}{\loccitn}}
\def\EGA{\textup{EGA}\kern.2em}
\def\SGA{\textup{SGA}\kern.3em}

\let\oldS\S
\def\S{\oldS\kern .15em}
\def\ptbl{.\kern .2em}
\def\resp{resp.\kern.3em}
\def\Exp{Exp.\kern.3em}

\let\bar\overline
\let\hat\widehat
\let\tilde\widetilde
\let\dpl\displaystyle
\let\ldots\dots
\let\emptyset\varnothing
\let\skpt\skippointrait
\let\epsilon\varepsilon

\def\sfrac#1#2{{#1}/{#2}}

\DeclareMathOperator{\cd}{cd}
%\DeclareMathOperator{\Hhom}{\mathcal{H}\!\mathit{om}}
\DeclareMathOperator{\Hhom}{\mathbf{Hom}}
%\DeclareMathOperator{\Eext}{\mathcal{E}\!\mathit{xt}}
\DeclareMathOperator{\Eext}{\mathbf{Ext}}
\DeclareMathOperator{\Spec}{Spec}
\DeclareMathOperator{\Tr}{Tr}

\let\emptyset=\varnothing

%%%%%%%%%%%%%%%
%%% expose XVII
%%%%%%%%%%%%%%%

\makeatletter
\let\old@optexpose\optexpose
\newcommand{\empty@optexpose}[1]{}
\newcommand{\paren@optexpose}[1]{\old@optexpose{#1}(}
\newcommand{\refnoexpose}[1]{\let\optexpose\empty@optexpose\ref{#1}\let\optexpose\old@optexpose}%\refnoexpose{X.cor5.2} engendre "5.2" au lieu de "X 5.2"
\newcommand{\neweqref}[1]{\let\optexpose\paren@optexpose\textup{\ref{#1})}\let\optexpose\old@optexpose}%\neweqref{XII.eq4.1} engendre "XII (4.1)" au lieu de "(XII 4.1)"
\makeatother
\newcommand{\refvb}[2][VI${}_{\text{B}}$ ]{\sisi{#1\refnoexpose{#2}}{\ref{#2}}}

\newcommand{\leftexp}[2]{{\vphantom{#2}}^{#1}{#2}}
\newcommand{\AAA}{\mathbf{A}}
\newcommand{\Sf}{\mathfrak{S}}
\DeclareMathOperator*{\quotprojlim}{\text{``$\varprojlim$''}}
\DeclareMathOperator*{\quotinjlim}{\text{``$\varinjlim$''}}
\DeclareMathOperator{\Fl}{F\sisi{\ell}{l}}
\DeclareMathOperator{\Pl}{P\sisi{\ell}{l}}
%\DeclareMathOperator{\cl}{c\sisi{\ell}{l}}
\DeclareMathOperator{\tdf}{tdf}
\DeclareMathOperator{\rg}{rg}
\newcommand{\iInd}{\sisi{I}{i}nd}
\newcommand{\iind}{ind}
\newcommand{\ppro}{pro}
\newcommand{\tTor}{\sisi{T}{t}or}
\newcommand{\ttor}{tor}
\newcommand{\chap}{Chap.}
\newcommand{\Ffex}{\Ff\mathit{ex}}
\newcommand{\sheafk}{\underline{k}}
\newcommand{\sheafn}{\underline{n}}
\newcommand{\sheafGamma}{\underline{\Gamma}}

% Repeated references to an already typeset footnote must not use
% \footnotemark without a matching \footnotetext: Hyperref would then create
% a live annotation to a destination that cannot exist.  This helper keeps the
% current printed footnote mark and links it to an explicit stable note target.
\makeatletter
\newcommand{\sgarepeatedfootnotemark}[1]{%
  \begingroup
  \edef\@thefnmark{\thefootnote}%
  \hyperlink{#1}{\@makefnmark}%
  \endgroup}
\makeatother

\sisi{
\DeclareMathOperator{\simarrow}{\rule{.15mm}{1.8mm}\kern -3.4pt\rightsquigarrow}
\newcommand{\sheafell}{\underline{\ell}}
\newcommand{\sheafZZ}{\underline{\ZZ}}
}{
\newcommand{\nderef}[1]{la N.D.E.~(\ref{#1}) page~\pageref{#1}}
}

%%%%%%%%%%%%%%%%%
%%%%% expose XVII
%%%%%%%%%%%%%%%%%

% END INCLUDED SGA4_English_current_buildable_source_20260906/source/fichiers-configuration-latex/macros.tex


\externaldocument{../01/01_English.tex}
%\externaldocument{../02/02_English.tex}
\externaldocument{../03/03.tex}
\externaldocument{../04/04.tex}

\begin{document}

\begin{center}
II. Topologies and sheaves\\
J.-L. Verdier
\end{center}

\selectlanguage{english}version: \showgitstatus
\tableofcontents

\else
\chapter{Topologies and sheaves}\MakeLinkTarget*{sga4.exp.II}\label{sga4.exp.II}
%\addtocontents{toc}{by J.-L. Verdier}
\begin{center}
J.-L. Verdier
\end{center}
\fi

% BEGIN INCLUDED SGA4_English_current_buildable_source_20260906/source/02/drafts/SGA4_II_LINES_31_212_English.texfrag
After\pageoriginale defining topologies and
pretopologies (\hyperref[II.sec1]{Section 1}) and sheaves of sets
(\hyperref[II.sec2]{Section 2}), in \hyperref[II.sec3]{Section 3} we turn to the
central theorem of this exposé (\hyperref[II.theo3.4]{3.4}): the existence
of the sheaf associated with a presheaf. In order to cover all the
cases encountered in practice, the existence of this functor is
proved for presheaves on a $𝒰$-site
(\hyperref[II.def3.0.2]{3.0.2}). The consequences of this theorem for
the behavior of inductive and projective limits in categories of
sheaves are drawn in \hyperref[II.sec4]{Sections 4} and 5. In \hyperref[II.sec6]{Section 6} we define and study
sheaves with values in arbitrary categories, immediately focusing
attention on sheaves of rings, groups,
modules, \etcn

\section{Topologies, covering families,
  pretopologies}\etiquette[1]{II.sec1}

\begin{enonce*}{Definition 1.1}\etiquette[1.1]{II.def1.1}
A topology on a category $C$ consists, for every
object $X$ of $C$, of a set $J(X)$ of sieves on $X$, subject to the
following axioms:
\begin{itemize}
\item[T 1)] {\rm{(Stability under base change).}} For every
  object $X$ of $C$, every sieve $R\in J(X)$, and every morphism
  $f:Y\rightarrow X$ $(Y\in \ob (C))$, the sieve $R\times_X Y$ on $Y$
  belongs to $J(Y)$.

\item[T 2)] {\rm(Local character).} If $R$ and $R'$ are two
  sieves on $X$, if $R\in J(X)$, and if, for every $Y\in \ob (C)$ and
  every morphism $Y\rightarrow R$, the sieve $R'\times_X Y$
  belongs\pageoriginale
  to $J(Y)$, then $R'$ belongs to $J(X)$.

\item[T 3)] For every object $X$ of $C$, $X$ belongs to $J(X)$.
\end{itemize}
\end{enonce*}

\setcounter{subsection}{1}
\subsubsection{}\etiquette[1.1.1]{II.subsubsec1.1.1}

The sieves belonging to $J(X)$ will be called the \emph{sieves
  covering $X$}, or the \emph{refinements of $X$}. It follows
  immediately from axioms (T1), (T2), and (T3) that the set of
  sieves covering $X$ is stable under finite intersections, and that
  every sieve containing a covering sieve is itself a covering
  sieve. Thus the set $J(X)$, ordered by inclusion, is
  cofiltered (\hyperrefext[I][I.sec8]{8}).

\subsubsection{}\etiquette[1.1.2]{II.subsubsec1.1.2}

Let $C$ be a category, and let $T$ and $T'$ be two topologies on $C$. The
topology $T$ is said to be \emph{finer than the topology} $T'$ if,
for every object $X$ of $C$, every refinement of $X$ for the
topology $T'$ is a refinement of $X$ for the topology $T$. This
defines an order structure on the set
of topologies.

\subsubsection{}\etiquette[1.1.3]{II.subsubsec1.1.3}

Let $(T_i)_{i\in I}$ be a family of topologies on $C$. For every
object $X$ of $C$, the set of sieves on $X$ that are
covering for all the topologies $T_i$ satisfies axioms
(T 1), (T 2), and (T 3), and hence defines a topology: the
\emph{intersection topology} of the $T_i$, \ie the infimum
of the $T_i$. It is the finest topology that is coarser than
all the topologies $T_i$. Consequently, the family $(T_i)_{i \in I}$
also admits a \emph{supremum}: the intersection topology of the
topologies finer than each of the~$T_i$.

\subsubsection{}\etiquette[1.1.4]{II.subsubsec1.1.4}
In\pageoriginale particular, taking $J(X)$ to be the set of all
sieves on $X$ gives a topology finer than every topology
on $C$; we call it the \emph{discrete topology}.

There is a topology coarser than every topology on $C$: the
topology for which $J(X)=\{X\}$ for every $X$ of $C$. This
topology is called the \emph{coarse} or
\emph{chaotic} topology.

\subsubsection{}\etiquette[1.1.5]{II.subsubsec1.1.5}
A category $C$ equipped with a topology is called a \emph{site.}
The category $C$ is called the \emph{category
  underlying the site.}

\begin{enonce*}{Definition 1.2}\etiquette[1.2]{II.def1.2}
Let $C$ be a site and $X$ an object of $C$. A family of morphisms
$(f_{\alpha}:X_{\alpha}\rightarrow X)$, $\alpha\in A$, is called a
covering family if the sieve generated by the family $f_{\alpha}$
{\rm(I~\hyperref[I.sec4.3.3]{4.3.3})} is a sieve covering $X$.
\end{enonce*}

\subsubsection{}\etiquette[1.1.6]{II.subsubsec1.1.6}
Let $C$ be a category. Suppose that, for each object $X$ of $C$,
we are given a set of families of morphisms with target $X$. There is
then a coarsest topology $T$ for which
the given families are covering, namely the
intersection (\hyperref[II.subsubsec1.1.3]{1.1.3}) of all the topologies in
question. This topology is called the \emph{topology generated
by the given sets of families
  of morphisms}. In general, it is difficult to describe all the
covering sieves of this topology.
The situation is more favorable, however, in the following case:

\begin{enonce*}{Definition 1.3}\etiquette[1.3]{II.def1.3}
Let $C$ be a category. A pretopology on $C$ consists, for
each object $X$ of $C$, of a set $\Cov(X)$ of families of
morphisms with target $X$, subject to the following
axioms:
\begin{itemize}
\item[PT 0)] For\pageoriginale every object $X$ of $C$, the morphisms
  in the families belonging to $\Cov(X)$ are pullbackable. (Recall
  that a morphism $Y\rightarrow X$ is called pullbackable if, for every
  morphism $Z\rightarrow X$, the fiber product $Z\times_X Y$ is
  representable.)

\item[PT 1)] For every object $X$ of $C$, every family
  $(X_{\alpha}\rightarrow X)_{\alpha \in A}$ belonging to
  $\Cov(X)$, and every morphism $Y\rightarrow X$ $(Y\in \ob (C))$, the
  family $(X_{\alpha}\times_X Y\rightarrow Y)_{\alpha \in A}$
  belongs to $\Cov(Y)$. (Stability under base change.)

\item[PT 2)] If $(X_{\alpha}\rightarrow X)_{\alpha \in A}$ belongs
  to $\Cov(X)$ and if, for every $\alpha\in A$,
  $(X_{\beta_{\alpha}}\rightarrow X_{\alpha})_{\beta_{\alpha}\in
  B_{\alpha}}$ belongs to $\Cov(X_{\alpha})$, then the family
  $(X_\gamma\rightarrow X)_{\gamma\in \amalg_{\alpha\in A}B_{\alpha}}$ (where,
  for every $\gamma=(\alpha,\beta_{\alpha})$, $\alpha\in A$,
  $\beta_{\alpha}\in B_{\alpha}$, the morphism $X_\gamma\rightarrow X$ is
  the composite morphism $X_\gamma=X_{\beta_{\alpha}}\rightarrow
  X_{\alpha}\rightarrow X$) belongs to $\Cov(X)$. (Stability under
  composition.)

\item[PT 3)] The family $X\xrightarrow{\id_X} X$ belongs to $\Cov(X)$.
\end{itemize}
\end{enonce*}

\setcounter{subsection}{3} \setcounter{subsubsection}{0}
\subsubsection{}\etiquette[1.3.1]{II.subsubsec1.3.1}

The definition in \hyperref[II.subsubsec1.1.6]{1.1.6} allows us to consider,
for a given \emph{pretopology} on $C$, the \emph{topology} on
$C$ generated by that pretopology. Note that if $C$ is a
category in which fiber products are representable, then every
topology $T$ on $C$ can be defined by a pretopology, namely
the one for which $\Cov(X)$ consists of \emph{all} the
families covering $X$ for the topology $T$.

\begin{enonce*}{Proposition 1.4}\etiquette[1.4]{II.prop1.4}
Let $C$ be a category, $E$ a pretopology on $C$, $T$ the
topology defined by the pretopology $E$
$(\hyperref[II.subsubsec1.1.6]{1.1.6})$, and $X$ an object of $C$. Denote by
$J_E(X)$ the set of sieves on $X$ generated by families
of the pretopology, and by $J_T(X)$ the set of refinements of
$X$ for the\pageoriginale topology $T$. Then $J_E(X)$ is cofinal
in $J_T(X)$. In other words, a sieve $R$ on $X$
belongs to $J_T(X)$ if and only if there exists a sieve
$R'\in J_E(X)$ such that $R'\subset R$.
\end{enonce*}

\begin{enonce*}[remark]{Proof}
For every object $X$ of $C$, let $J'(X)$ be the set of sieves on
$X$ that contain a sieve belonging to $J_E(X)$. Clearly,
$J'(X)\subset J_T(X)$. To prove that $J'(X)=J_T(X)$, it is enough to
show that the data $J'(X)$ define a topology on $C$. Now
the $J'(X)$ clearly satisfy axioms (T 1) and (T 3). It
therefore remains to verify (T 2). For this, it is enough to prove that if
$R'$ is a subsieve of $R\in J_E(X)$ such that, for every
$Y\rightarrow R$, the sieve $R'\times_X Y$ on $Y$ belongs to
$J'(Y)$, then the sieve $R'$ belongs to $J'(X)$. Now the sieve $R$ is
generated by a family $(X_{\alpha}\rightarrow X)$ belonging to
$\Cov(X)$, and for every $\alpha$ the sieve $R'\times_X X_{\alpha}$
contains a sieve generated by a family
$(X_{\beta_{\alpha}}\rightarrow X_{\alpha})$ belonging to
$\Cov(X_{\alpha})$. It follows that the sieve $R'$ contains a
sieve generated by the family $(X_{\beta_{\alpha}}\rightarrow X)$.
Thus, by axiom (PT 2), $R'$ contains a sieve belonging to $J_E(X)$
and consequently belongs to $J'(X)$, \cqfd
\end{enonce*}

% END INCLUDED SGA4_English_current_buildable_source_20260906/source/02/drafts/SGA4_II_LINES_31_212_English.texfrag


% BEGIN INCLUDED SGA4_English_current_buildable_source_20260906/source/02/drafts/SGA4_II_LINES_213_421_English.texfrag
\section{Sheaves of sets}\etiquette[2]{II.sec2}

\begin{enonce*}{Definition 2.1}\etiquette[2.1]{II.def2.1}
Let $C$ be a site whose underlying category is a
$𝒰$-category. A presheaf $F$ with values in
$𝒰-\Ens$ is said to be separated (\resp to be a sheaf) if, for
every sieve $R$ covering an object $X$ of $C$, the map
$$
\Hom_{C^{\widehat{\phantom{,}}}}(X,F)\longrightarrow
\Hom_{C^{\widehat{\phantom{,}}}}(R,F)
$$
is injective (\resp bijective). The full subcategory
of $C^{\widehat{\phantom{,}}}$ whose\pageoriginale objects are
the sheaves is called the category of sheaves
of sets on $C$, and is most often denoted by $\widehat{C}$
\footnote[1]{or
  $C^{\widetilde{\phantom{,}}}$, depending on the mood of the
  typewriter.}. When no ambiguity can result, we shall
simply say the category of sheaves on $C$; on the other hand,
we shall sometimes be more precise and say the category of sheaves with
values in $𝒰-\Ens$, or the category of
$𝒰$-sheaves.
\end{enonce*}

\begin{enonce*}{Proposition 2.2}\etiquette[2.2]{II.prop2.2}
Let $C$ be a $𝒰$-category, and let $\Ff=(F_i)_{i\in I}$ be a
family of $𝒰$-presheaves on $C$. For
each object $X$ of $C$, denote by $J_{\Ff}(X)$ the set of sieves
$R\rightarrow X$ such that, for every morphism $Y\rightarrow X$ in
$C$ with target $X$, the sieve $R\times_X Y$ has the following
property: the map
$$
\Hom_{C^{\widehat{\phantom{,}}}}(Y,F_i)\longrightarrow
\Hom_{C^{\widehat{\phantom{,}}}} (R\times_X Y, F_i)
$$
is bijective (\resp injective) for every $i\in I$. Then the
sets $J_{\Ff}(X)$ define a topology on $C$, which is
the finest topology for which each $F_i$ is a
sheaf (\resp separated presheaf).
\end{enonce*}

\begin{enonce*}[remark]{Proof}
The $J_{\Ff}(X)$ clearly satisfy axioms (T 1) and (T
3). It remains to show that they satisfy (T 2). For this it suffices
to show the following:
\begin{enumerate}
\renewcommand{\labelenumi}{\theenumi)}
\item If $R'\hookrightarrow R\hookrightarrow X$ are two sieves on
  $X$ such that $R\in J_{\Ff}(X)$ and such that, for every
  $Y\rightarrow R$ (where $Y$ is an object of $C$), the sieve
  $R'\times_X Y$ belongs to $J_{\Ff}(Y)$, then the sieve
  $R'$ belongs to $J_{\Ff}(X)$.

\item If $R'\hookrightarrow R\hookrightarrow X$ are two sieves on
  $X$ such that $R'\in J_{\Ff}(X)$, then the sieve $R$
  belongs to $J_{\Ff}(X)$.
\end{enumerate}
\end{enonce*}

Note that, in case 2), for every $Y\rightarrow R$ (where $Y$ is
an object of $C$), the sieve $R'\times_X Y$ belongs to $J_{\Ff}(Y)$.
Now, in cases 1) and 2), $R$ is the inductive limit of the objects $Y$ of
$C$ over $R$ (I \hyperref[I.prop3.4]{3.4}). The\pageoriginale inductive limits
in $C^{\widehat{\phantom{,}}}$ are universal
(\hyperrefext[I][I.cor3.3]{3.3}). Consequently, in cases 1) and 2), $R'$ is the inductive limit
of the $R'\times_X Y$. Now, in cases 1) and 2),
$R'\times_X Y$ belongs to $J_{\Ff}(Y)$. Passing to the
inductive limit over the objects $Y$ of $C$ over $R$, we deduce that
the map
$$
\Hom_{C^{\widehat{\phantom{,}}}}(R,F_i)\longrightarrow
\Hom_{C^{\widehat{\phantom{,}}}}(R',F_i)
$$
is a bijection (\resp an injection) for every $i\in I$. It
follows that the maps
$$
\Hom_{C^{\widehat{\phantom{,}}}}(X,F_i)\longrightarrow
\Hom_{C^{\widehat{\phantom{,}}}}(R',F_i)\quav
\Hom_{C^{\widehat{\phantom{,}}}}(X,F_i)\longrightarrow
\Hom_{C^{\widehat{\phantom{,}}}}(R,F_i)
$$
are, in cases 1) and 2), bijections (\resp injections).
Moreover, hypotheses 1) and 2) are plainly stable under
arbitrary base change $Y\rightarrow X$ ($Y$ an object of $C$).
It then follows that, for every object $Y$ of $C$ over $X$
and every $i\in I$, the maps
\begin{align*}
\Hom_{C^{\widehat{\phantom{,}}}}(Y,F_i)&\longrightarrow
\Hom_{C^{\widehat{\phantom{,}}}}(R\times_X Y, F_i)\\
\Hom_{C^{\widehat{\phantom{,}}}}(Y,F_i)&\longrightarrow
\Hom_{C^{\widehat{\phantom{,}}}}(R'\times_X Y, F_i)
\end{align*}
are in both cases bijections (\resp injections); this
shows that $R'$ and $R$ belong to $J_{\Ff}(X)$, \cqfd

\begin{enonce*}{Corollary 2.3}\etiquette[2.3]{II.cor2.3}
Let $C$ be a $𝒰$-category, and for every $X$ of $C$ let
$K(X)$ be a set of sieves on $X$. Suppose that the $K(X)$
satisfy axiom {\rm (T 1)} of {\rm(\hyperref[II.def1.1]{1.1})}. A
presheaf $F$ is a sheaf (\resp a separated presheaf)
for the topology generated (\hyperref[II.subsubsec1.1.6]{1.1.6})
by the $K(X)$ if and only if, for every object $X$ of $C$
and every sieve $R\in K(X)$, the map
$$
\Hom_{C^{\widehat{\phantom{,}}}}(X,F)\longrightarrow
\Hom_{C^{\widehat{\phantom{,}}}}(R,F)
$$
is a bijection (\resp an injection).
\end{enonce*}

\begin{enonce*}{Corollary 2.4}\etiquette[2.4]{II.cor2.4}
In\pageoriginale particular, let $C$ be a $𝒰$-category
equipped with a pretopology. A presheaf $F$ is a
sheaf (\resp a separated presheaf) if and only if,
for every object $X$ of $C$ and every covering family
$(X_{\alpha}\rightarrow X)$ belonging to $\Cov(X)$, the diagram
of sets
$$
F(X)\rightarrow \prod_{\alpha\in A}F(X_{\alpha})\rightrightarrows
\prod_{(\alpha,\beta)\in A\times A}F(X_{\alpha}\times_X X_{\beta})
$$
is exact (\resp the map
$$
F(X)\longrightarrow \prod_{\alpha\in A}F(X_{\alpha})
$$
is injective).
\end{enonce*}

\begin{enonce*}[remark]{Proof}
Apply \hyperref[II.cor2.3]{2.3}, then (\hyperrefext[I][I.cor3.5]{3.5}) and
(\hyperrefext[I][I.prop2.12]{2.12}).

Corollary \hyperref[II.cor2.4]{2.4} recovers the definition
given in [1].
\end{enonce*}

\begin{enonce*}{Definition 2.5}\etiquette[2.5]{II.def2.5}
Let $C$ be a $𝒰$-category. The canonical topology
of $C$ is the finest topology for which the representable functors
are sheaves (\hyperref[II.prop2.2]{2.2}). A sieve
covering $X$ for the canonical topology will be called a universally
strictly epimorphic sieve. A covering family for the
canonical topology will be called a universally strictly epimorphic
family. If, in addition, the morphisms of the covering family
are pullbackable, the family will be called universally effective
epimorphic $[2]$.
\end{enonce*}

\begin{enonce*}[remark]{Remark 2.5.1}\etiquette[2.5.1]{II.rem2.5.1}
For almost all the sites that have been used up to now,
the topology is \emph{coarser} than the canonical topology; in
other words, the representable functors on $C$ are
sheaves, \ie the covering families of $C$ are universally strictly
epimorphic. The\pageoriginale only exception to this
rule is the site $\widehat{C}$ studied in ${\rm n^{\circ}}$ 5
(whose topology is \emph{finer,} and most often strictly
finer, than the canonical topology).
\end{enonce*}

\begin{enonce*}{Proposition 2.6}\etiquette[2.6]{II.prop2.6}
A sieve $R\hookrightarrow X$ is universally strictly
epimorphic if and only if, for every object $Y\rightarrow X$
of $C$ over $X$, and for every object $Z$ of $C$,
the map
$$
\Hom_C(Y,Z)\longrightarrow \varprojlim_{C/(Y\times_X R)}\Hom_C (.,
Z)
$$
is a bijection.
\end{enonce*}

\begin{enonce*}[remark]{Proof}
Immediate upon applying \hyperref[II.prop2.2]{2.2} and then (\hyperrefext[I][I.sec5.3]{5.3}).
\end{enonce*}

\begin{enonce*}[remark]{Remarks 2.7}\etiquette[2.7]{II.rem2.7}
\begin{enumerate}
\renewcommand{\labelenumi}{\theenumi)}
\item Proposition \hyperref[II.prop2.6]{2.6} gives a characterization of the universally
  strictly epimorphic sieves of a category $C$
  that is independent of the universe in which the presheaves take
  their values, provided only that the morphism sets
  $\Hom(X,Y)$ of $C$ belong to that universe. It
  thus makes it possible to define the canonical topology for any
  category $C$, without having to specify the universes.

\item Let $C$ be a site whose underlying category is a
  $𝒰$-category, let $F$ be a $𝒰$-presheaf, and let
  $𝒱$ be a universe containing $𝒰$. The category
  underlying $C$ is a $𝒱$-category, and $F$ may
  be regarded as a $𝒱$-presheaf. The
  $𝒰$-presheaf $F$ is a sheaf (\resp a
  separated presheaf) if and only if the
  $𝒱$-presheaf $F$ is a sheaf (\resp a
  separated presheaf). In other words, the condition for a
  presheaf $F$ to be a sheaf (\resp a separated
  presheaf) does not depend (in the sense just\pageoriginale
  specified) on the universe in which the presheaf $F$ takes its
  values. In particular, let $C$ be a site and let $F$ be a
  $𝒰$-presheaf; we say that $F$ is a sheaf (\resp a
  separated presheaf) if there exists a universe $𝒱$
  containing $𝒰$ such that the category $C$ is a
  $𝒱$-category and such that $F$ is a $𝒱$-sheaf
  (\resp a separated $𝒱$-presheaf). This property
  does not depend on the universe $𝒱$.
\end{enumerate}
\end{enonce*}

% END INCLUDED SGA4_English_current_buildable_source_20260906/source/02/drafts/SGA4_II_LINES_213_421_English.texfrag


% BEGIN INCLUDED SGA4_English_current_buildable_source_20260906/source/02/drafts/SGA4_II_LINES_422_562_English.texfrag
\section{The sheaf associated with a presheaf}\etiquette[3]{II.sec3}

\begin{enonce*}{Definition 3.0.1}\etiquette[3.0.1]{II.def3.0.1}
Let $C$ be a site. A topologically generating family (or, when no
confusion can result, a generating family) of $C$ is a set $G$ of
objects of $C$ such that every object of $C$ is the target of a
covering family of morphisms of $C$ (\hyperref[II.def1.2]{1.2}) whose
sources are elements of $G$.
\end{enonce*}

\begin{enonce*}{Definition 3.0.2}\etiquette[3.0.2]{II.def3.0.2}
Let $𝒰$ be a universe. A $𝒰$-site is a site $C$ whose underlying
category is a $𝒰$-category
{\rm(\hyperrefext[I][I.def1.1]{1.1})}, and which possesses a small
topologically generating family. Let $C$ be a $𝒰$-category; a
$𝒰$-topology on $C$ is a topology on $C$ making $C$ a $𝒰$-site. A
site $C$ is called $𝒰$-small, or, by abuse of language, small, if its
underlying category is small {\rm(\hyperrefext[I][I.sec1.0]{1.0})}.
\end{enonce*}

\setcounter{subsubsection}{2}
\subsubsection{}\etiquette[3.0.3]{II.subsubsec3.0.3}
It follows immediately from the definitions that every topology finer
than a $𝒰$-topology is a $𝒰$-topology, and that a small site is a
$𝒰$-site.

\begin{enonce*}{Proposition 3.0.4}\etiquette[3.0.4]{II.prop3.0.4}
Let\pageoriginale $C$ be a $𝒰$-site and $G$ a small topologically
generating family of $C$. For every $X\in\ob C$, denote by $J_G(X)$
the set of sieves covering $X$ generated by a family of morphisms
$$
(Y_{\alpha}\xrightarrow{u_{\alpha}} X), \alpha\in A, \mbox{ where }
Y_{\alpha}\in G{\quoi}.
$$
\begin{enumerate}
\renewcommand{\labelenumi}{\theenumi)}
\item The set $J_G(X)$ is small.

\item The set $J_G(X)$ is cofinal in the set $J(X)$ of all sieves
covering $X$, ordered by inclusion.

\item For every $R\in J_G(X)$, there exists a small epimorphic family
{\rm(\hyperrefext[I][I.sec10.2]{10.2})}
$(u_\alpha:Y\rightarrow R)$, $\alpha\in A$, with $Y\in G$.
\end{enumerate}
\end{enonce*}

\begin{enonce*}[remark]{Proof}
\begin{enumerate}
\renewcommand{\labelenumi}{\theenumi)}
\item Put $A(X)=\amalg_{Y\in G}\hom(Y,X)$. The set $A(X)$ is small
(\hyperrefext[I][I.sec0]{0}) and
$\card(J_G(X))\leq2^{\card(A(X))}$. Consequently $J_G(X)$ is small.

\item Let $R\in J(X)$. Put $A(R)=\amalg_{Y\in G}\Hom(Y,R)$, and let
$R'$ be the sieve of $X$ generated by the family
$(Y\xrightarrow{u}R\hookrightarrow X)$, $u\in A(R)$. One has
$R'\subset R$, and it suffices to show that $R'$ is covering. By
axiom (T 2) of \hyperref[II.def1.1]{1.1}, it suffices to show that for
every morphism $Z\rightarrow R$, $Z\in\ob C$, the sieve
$R'\times_XZ$ of $Z$ is covering. Now the sieve $R'\times_XZ$
contains the sieve generated by the family of all morphisms
$Y\rightarrow Z$ with $Y\in G$, which is covering by hypothesis.
Therefore the sieve $R'\times_XZ$ of $Z$ is covering (axiom T2 of
\hyperref[II.def1.1]{1.1}).

\item Let $R\in J_G(X)$. The family $(Y\xrightarrow{u}R)$,
$u\in A(R)=\amalg_{Y\in G}\Hom(Y,R)$, is epimorphic by hypothesis.
Now, for every $Y\in\ob C$, $\Hom(Y,R)$ is contained in $\Hom(Y,X)$,
which is small. Consequently $A(R)$ is small.
\end{enumerate}
\end{enonce*}

\setcounter{subsubsection}{4}
\subsubsection{}\etiquette[3.0.5]{II.subsubsec3.0.5}
Let $C$ be a $𝒰$-site, $𝒱$ a universe such that $C\in𝒱$ and
$𝒰\subset𝒱$, and $G$ a $𝒰$-small topologically generating family of
$C$. The category $C^{\widehat{\phantom{,}}}$ of presheaves of
$𝒰$-sets on $C$ is a $𝒱$-category (I \hyperref[I.sec1.1.1]{1.1.1}). Let
$X$\pageoriginale be an object of $C$. The set $J(X)$ of sieves
covering $X$ is $𝒱$-small and, ordered by inclusion, is cofiltered
(\hyperref[II.subsubsec1.1.1]{1.1.1}). For every $𝒰$-presheaf $F$,
the inductive limit
$\varinjlim_{J(X)}\Hom_{C^{\widehat{\phantom{,}}}}(R,F)$ is therefore
representable by an element of $𝒱$
(\hyperrefext[I][I.def2.4.1]{2.4.1}). Moreover, it follows from
(\hyperref[II.prop3.0.4]{3.0.4} 3)) that, for every $R\in J_G(X)$,
$\Hom_{C^{\widehat{\phantom{,}}}}(R,F)$ is $𝒰$-small, and since
$J_G(X)$ is a $𝒰$-small set cofinal in $J(X)$ (\loccit), it follows
from (\hyperrefext[I][I.sec2.4.2]{2.4.2}) that
$\varinjlim_{J(X)}\Hom_{C^{\widehat{\phantom{,}}}}(R,F)$ is
$𝒰$-small. Choose, for every $F$ and every $X$, an element of $𝒰$
representing this inductive limit, and put
$$
LF(X)=\varinjlim_{J(X)}\Hom_{C^{\widehat{\phantom{,}}}}(.,F).
$$
Let $g:Y\rightarrow X$ be a morphism of $C$. The base-change functor
$g^\ast:J(X)\rightarrow J(Y)$ defines a map
$$
LF(g):LF(X)\longrightarrow LF(Y){\quoi},
$$
which makes $X\mapsto LF(X)$ a presheaf on $C$.

Since the morphism $\id_X:X\rightarrow X$ is an element of $J(X)$,
for every object $X$ of $C$ one has a map
$$
\ell(F)(X):F(X)\longrightarrow LF(X){\quoi},
$$
which visibly defines a morphism of functors
$$
\ell(F):F\longrightarrow LF{\quoi}.
$$
It is moreover clear that $F\mapsto LF$ is a functor in $F$ and that
the morphisms $\ell(F)$ define a morphism
$$
\ell:\Id\longrightarrow L{\quoi}.
$$
Finally, let $R\hookrightarrow X$ be a covering sieve of $X$. The
definition of $LF(X)$ and (\hyperrefext[I][I.prop1.4]{1.4}) provide a
map
$$
Z_R:\Hom_{C^{\widehat{\phantom{,}}}}(R,F)\longrightarrow
\Hom_{C^{\widehat{\phantom{,}}}}(X,LF){\quoi},
$$
and,\pageoriginale for every morphism $Y\xrightarrow{g}X$ of $C$, the
definition of the functor $LF$ shows that the following diagram is
commutative:
\[\vcenter{
\xymatrix{ \qquad Z_R :& \Hom_{C^{\widehat{\phantom{a}}}}(R,F)
\ar[r]\ar[d] &
\Hom_{C^{\widehat{\phantom{a}}}}(X,LF)\ar[d]\\
Z_{\fprod{R}{Y}{X}}:&  \Hom_{C^{\widehat{\phantom{a}}}}(\fprod{R}{Y}{X},F) \ar[r]&  \Hom_{C^{\widehat{\phantom{a}}}}(Y,LF)\\
}}\leqno{(*)}
\]
(the vertical arrows are the evident arrows).

\noindent Let us now copy [2].


% END INCLUDED SGA4_English_current_buildable_source_20260906/source/02/drafts/SGA4_II_LINES_422_562_English.texfrag


% BEGIN INCLUDED SGA4_English_current_buildable_source_20260906/source/02/drafts/SGA4_II_LINES_563_764_English.texfrag
\begin{enonce*}{Lemma 3.1}\etiquette[3.1]{II.lem3.1}
\begin{enumerate}
\renewcommand{\labelenumi}{\theenumi)}
\item For every covering sieve $i_R:R\hookrightarrow X$ and every
$u:R\rightarrow F$, the diagram
\[\vcenter{
\xymatrix@L=5pt@R=1.35cm@C=3.5cm@M=.2cm{
F\ar[r]^{\ell(F)} & LF\\
R\ar[u]^{u}\ar@{^{(}->}_{i_R}[r] & X \ar[u]_{Z_R(u)} }}\leqno{(**)}
\]
is commutative.

\item For every morphism $v:X\rightarrow LF$, there exists a covering
sieve $R$ of $X$ and a morphism $u:R\rightarrow F$ such that
$Z_R(u)=v$.

\item Let $Y$ be an object of $C$, and let
$u,v:Y\rightrightarrows F$ be two morphisms such that
$\ell(F)\circ u=\ell(F)\circ v$. Then the kernel of the pair $(u,v)$
is a covering sieve of $Y$.

\item Let $R,R'$ be two covering sieves of $X$, and let
$u:R\rightarrow F$, $u':R'\rightarrow F$ be two morphisms. For
$Z_R(u)=Z_{R'}(u')$, it is necessary and sufficient that $u$ and $u'$
coincide on a covering sieve $R''\hookrightarrow R\times_XR'$.
\end{enumerate}
\end{enonce*}

\begin{enonce*}[remark]{Proof}
The only nontrivial assertion is assertion 1). We must show that
$Z_R(u)\circ i_R=\ell(F)\circ u$. For this it suffices to show
(\hyperrefext[I][I.prop3.4]{3.4}) that the composites of these
morphisms with every morphism $Y\xrightarrow{g}R$ ($Y$ an object
of\pageoriginale $C$) are equal. Consider $f=i_R\circ g$ and the fiber
product $R\times_XY$:
\[\vcenter{
\xymatrix@L=5pt@R=1cm@C=3.5cm@M=.2cm{
F  \ar[r]^{\ell(F)}   & LF \\
R \ar[u]^{u}        \ar@{^{(}->}_{i_R}[r]  & X \ar[u]_{Z_R(u)} \\
\fprod{R}{Y}{X} \ar[u]^{g'} \ar[r]^(.35){\sim}_{i'} & Y
\ar[lu]^{g}\ar[u]_{f}. }}
\]

In the diagram above, $i'$ is a monomorphism admitting a section.
Consequently (\hyperref[II.def1.3]{1.3} 3) $i'$ is an isomorphism. By
the definition of the morphism $\ell(F)$, the morphism
$Z_{R\times_XY}(u\circ g')$ is equal to $\ell(F)\circ u\circ g$. On
the other hand, the commutativity of diagram (*) gives the equality
$Z_{R\times_XY}(u\circ g')=Z_R(u)\circ f$, and hence one indeed has
$\ell(F)\circ u\circ g=Z_R(u)\circ i_R\circ g$.
\end{enonce*}

\begin{enonce*}{Proposition 3.2}\etiquette[3.2]{II.prop3.2}
\begin{enumerate}
\renewcommand{\labelenumi}{\theenumi)}
\item The functor $L$ is left exact
{\rm(\hyperrefext[I][I.cor2.3.2]{2.3.2}).}

\item For every presheaf $F$, $LF$ is a separated presheaf.

\item The presheaf $F$ is separated if and only if the morphism
$\ell(F):F\rightarrow LF$ is a monomorphism. The presheaf $LF$ is then
a sheaf.

\item The following properties are equivalent:
\begin{enumerate}
\renewcommand{\theenumii}{\roman{enumii}}
\renewcommand{\labelenumii}{(\theenumii)}
\item $\ell(F):F\rightarrow LF$ is an isomorphism.

\item $F$ is a sheaf.
\end{enumerate}
\end{enumerate}
\end{enonce*}

\begin{enonce*}[remark]{Proof}
\begin{enumerate}
\renewcommand{\labelenumi}{\theenumi)}
\item It suffices to show (\hyperrefext[I][I.prop3.1]{3.1}) that, for
every object $X$ of $C$, the functor $F\mapsto LF(X)$ commutes with
finite projective limits. By the definition of a projective limit, the
functor
$$
F\longmapsto \Hom_{C^{\widehat{\phantom{,}}}}(R,F)\quad
R\hookrightarrow X\in \ob(J(X))
$$
commutes\pageoriginale with projective limits, and the inductive limit
$\varinjlim_{J(X)}$ commutes with finite projective limits because
$J(X)$ is a cofiltered set (\hyperrefext[I][I.def2.7]{2.7}).

\item Let $X$ be an object of $C$, and let
$f,g:X\rightrightarrows LF$ be two morphisms that coincide on a
covering sieve $R\hookrightarrow X$ of $X$. By
\hyperref[II.lem3.1]{3.1}\quad 2), there exists a covering sieve
$R'\hookrightarrow X$, which may always be supposed contained in $R$,
and two morphisms $u,v:R'\rightrightarrows F$ such that
$Z_{R'}(u)=f$ and $Z_{R'}(v)=g$. By
\hyperref[II.lem3.1]{3.1} 1), one then has
$\ell(F)\circ u=\ell(F)\circ v$. Consequently
(\hyperref[II.lem3.1]{3.1} 4)) $u$ and $v$ coincide on a covering
sieve $R''\hookrightarrow R'$. Let $w$ be the restriction of $u$ to
$R''$. One has $Z_{R''}(w)=Z_{R'}(u)=Z_{R'}(v)$, and consequently
$f=g$. Thus the presheaf $LF$ is separated.

\item If $F$ is separated, the morphism $\ell(F)$ is a monomorphism
because a filtered inductive limit of monomorphisms is a monomorphism.
If $\ell(F)$ is a monomorphism, the presheaf $F$ is a subpresheaf of a
separated presheaf, and therefore is separated. Let us show that $LF$
is then a sheaf. Let $i:R\hookrightarrow X$ be a sieve covering an
object $X$ of $C$, and let $u:R\rightarrow LF$ be a morphism. It
suffices to show that $u$ factors through $X$. Put
$R'=\fprod{F}{R}{LF}$ and $v=Z_{R'}(\pr_1)$:
\[
\xymatrix@L=5pt@R=1cm@C=1.5cm{
F \ar[rr]^{\ell(F)}               && LF                      &\\
R' \ar[u]^{\pr_1} \ar[rr]^{\pr_2}   && R \ar[r]^{i} \ar[u]_{u} & X \ar@{-}[lu]_{v}\\
R''\ar[u]^{p_1} \ar[rr]^{p_2}     && Y  \ar[u]_{m}               &
\quad . }
\]
To\pageoriginale show that $u=v\circ i$, it suffices to show (I \hyperref[I.prop3.4]{3.4})
that for every morphism $m:Y\rightarrow R$ ($Y$ an object of $C$),
$v\circ i\circ m=u\circ m$. Put $R''=Y\times_RR'$, and let $p_1,p_2$
be the projections. The projection
$\xymatrix@M=5pt{R''\ar@{^{(}->}[r]^{p_2}&Y}$ is a monomorphism and
makes $R''$ a covering sieve of $Y$ (\hyperref[II.lem3.1]{3.1} 2)).
One has $v\circ i\circ\pr_2=\ell(F)\circ\pr_1$
(\hyperref[II.lem3.1]{3.1} 1)), and hence
$v\circ i\circ\pr_2=u\circ m\circ p_2$. It follows that
$v\circ i\circ\pr_2\circ p_1=u\circ\pr_2\circ p_1$, \ie
$v\circ i\circ m\circ p_2=u\circ m\circ p_2$. Since the presheaf
$LF$ is separated, one has $v\circ i\circ m=u\circ m$, \cqfd

\item Clear.
\end{enumerate}
\end{enonce*}

\begin{enonce*}[remark]{Remark 3.3}\etiquette[3.3]{II.rem3.3}
Let $J'(X)$ be a cofinal subset of $J(X)$. One has
$$
LF(X)=\varinjlim_{J'(X)}\Hom_{C^{\widehat{\phantom{,}}}}(.,F){\quoi}.
$$

In particular, if the topology of $C$ is defined by a pretopology
$X\mapsto\Cov(X)$ (\hyperref[II.subsubsec1.1.3]{1.1.3}), the functor
$L$ can be described using the covering families of $\Cov(X)$
(\hyperrefext[I][I.prop2.12]{2.12} and
\hyperrefext[I][I.cor3.5]{3.5}). Writing the formulas explicitly, one
recovers the construction of [
\end{enonce*}

\begin{enonce*}{Theorem 3.4}\etiquette[3.4]{II.theo3.4}
Let $C$ be a $𝒰$-site. The inclusion functor
$i:C^{\widetilde{\phantom{,}}}\hookrightarrow
C^{\widehat{\phantom{,}}}$ from sheaves into presheaves admits a left
adjoint $\underline{a}$, which is left exact
(\hyperrefext[I][I.cor2.3.2]{2.3.2}):
\[
\xymatrix@=3cm@M=.2cm{
 \widetilde{C}  \ar @{^{(}->} @<-1ex>[r]_{i} & C^{\widehat{\phantom{,}}} \ar@<-1ex>[l]_{\underline{a}}
\qquad .}
\]

The functor $i\circ\underline{a}$ is canonically isomorphic to the
functor $L\circ L$ (cf. \hyperref[II.subsubsec3.0.5]{3.0.5}). For
every presheaf $F$, the adjunction morphism
$F\rightarrow i\circ\underline{a}(F)$ is deduced, via the preceding
isomorphism, from the morphism
$\ell(LF)\ell(F):F\rightarrow L\circ L(F)$.
\end{enonce*}

\begin{enonce*}{Definition 3.5}\etiquette[3.5]{II.def3.5}
The\pageoriginale sheaf $\underline{a}F$ is called the sheaf associated
with the presheaf $F$.
\end{enonce*}

Theorem \hyperref[II.theo3.4]{3.4} follows immediately from Proposition
\hyperref[II.prop3.2]{3.2}.

\begin{enonce*}{Proposition 3.6}\etiquette[3.6]{II.prop3.6}
Let $C$ be a $𝒰$-site and let $𝒱\supset𝒰$ be a universe. Denote by
$C^{\widehat{\phantom{,}}}_{𝒰}$ and
$C^{\widetilde{\phantom{,}}}_{𝒰}$ (respectively
$C^{\widehat{\phantom{,}}}_{𝒱}$ and
$C^{\widetilde{\phantom{,}}}_{𝒱}$) the categories of $𝒰$-presheaves
and $𝒰$-sheaves (respectively $𝒱$-presheaves and $𝒱$-sheaves), and
$\underline{a}_{𝒰}:C^{\widehat{\phantom{,}}}_{𝒰}\rightarrow
C^{\widetilde{\phantom{,}}}_{𝒰}$ (respectively
$\underline{a}_{𝒱}:C^{\widehat{\phantom{,}}}_{𝒱}\rightarrow
C^{\widetilde{\phantom{,}}}_{𝒱}$) the corresponding ``associated
sheaf'' functors. The diagram
\[
\xymatrix@C=3.5cm@R=1.25cm{ C^{\widehat{\phantom{a}}}_{U}
\ar[r]^{\underline{a}_{\underline{U}}}
\ar @{_{(}->}[d]& C^{\widetilde{\phantom{a}}}_{\underline{U}}\ar @{_{(}->}[d]\\
C^{\widehat{\phantom{a}}}_{\underline{V}}
\ar[r]^{\underline{a}_{\underline{V}}} &
C^{\widetilde{\phantom{a}}}_{\underline{V}} , }
\]
where the vertical functors are the canonical inclusions, is
commutative up to canonical isomorphism.
\end{enonce*}

\begin{enonce*}[remark]{Proof}
This follows from the construction of the functors $\underline{a}_U$
and $\underline{a}_V$ (\hyperref[II.theo3.4]{3.4} and
\hyperref[II.subsubsec3.0.5]{3.0.5}).
\end{enonce*}


% END INCLUDED SGA4_English_current_buildable_source_20260906/source/02/drafts/SGA4_II_LINES_563_764_English.texfrag


% BEGIN INCLUDED SGA4_English_current_buildable_source_20260906/source/02/drafts/SGA4_II_LINES_765_975_English.texfrag
\section{Exactness properties of the category of sheaves}\etiquette[4]{II.sec4}

The exactness properties of the category of sheaves follow from the
exactness properties of the category of presheaves via Theorem
\hyperref[II.theo3.4]{3.4}. This section makes this philosophy explicit
through some of the most useful standard statements.

\begin{enonce*}{Theorem 4.1}\etiquette[4.1]{II.theo4.1}
Let $C$ be a $𝒰$-site, let $\widetilde{C}$ be the category of
sheaves, let $\underline{a}:C^{\widehat{\phantom{,}}}\rightarrow
\widetilde{C}$ be the associated-sheaf functor, and let
$i:\widetilde{C}\rightarrow C^{\widehat{\phantom{,}}}$ be the
inclusion functor.
\begin{enumerate}
\renewcommand{\labelenumi}{\theenumi)}
\item The functor $\underline{a}$ commutes with inductive limits and is exact.

\item The\pageoriginale $𝒰$-inductive limits in
$\widetilde{C}$ are representable. For every small category $I$ and
every functor $E:I\rightarrow \widetilde{C}$, the canonical morphism
$$
\varinjlim_I E\longrightarrow \underline{a}(\varinjlim_I i\circ E)
$$
is an isomorphism.

\item The $𝒰$-projective limits in $\widetilde{C}$ are
representable. For every object $X$ of $C$, the functor on
$\widetilde{C}:F\mapsto F(X)$ commutes with projective limits,
\ie the inclusion functor $i:\widetilde{C}\rightarrow \widehat{C}$
commutes with projective limits.
\end{enumerate}
\end{enonce*}

\begin{enonce*}[remark]{Proof}
These properties follow essentially from Theorem
\hyperref[II.theo3.4]{3.4} and from
(\hyperrefext[I][I.prop2.11]{2.11}).
\end{enonce*}

Thus, in the category of sheaves, products indexed by an element of
$𝒰$, fiber products, sums indexed by an element of $𝒰$, amalgamated
sums, kernels, cokernels, images, and coimages are representable.

\begin{enonce*}{Corollary 4.1.1}\etiquette[4.1.1]{II.cor4.1.1}
Let $C$ be a $𝒰$-site and let $F$ be a sheaf of sets on $C$. The
canonical homomorphism
$$
\varinjlim_{C/F} \underline{a}X\longrightarrow F
$$
is an isomorphism.
\end{enonce*}

\begin{enonce*}[remark]{Proof}
This follows from (\hyperrefext[I][I.prop3.4]{3.4}) and the fact that
$\underline{a}$ commutes with inductive limits.
\end{enonce*}

\begin{enonce*}{Proposition 4.2}\etiquette[4.2]{II.prop4.2}
Every morphism of $\widetilde{C}$ that is both an epimorphism and a
monomorphism is an isomorphism.
\end{enonce*}

\begin{enonce*}[remark]{Proof}
Let $u:G\rightarrow H$ be a morphism of $\widetilde{C}$ that is both
an epimorphism and a monomorphism. First observe that the morphism $u$
is a monomorphism of\pageoriginale presheaves
(\hyperref[II.theo4.1]{4.1} 1)). Construct the amalgamated sum
$H\amalg_G H=K$ and the two canonical morphisms $i_1$,
$i_2:H\rightrightarrows K$ in the category of presheaves. Since $u$
is a monomorphism of presheaves, one immediately verifies that the
following diagram is cartesian and cocartesian
(\hyperrefext[I][I.sec10.1]{10.1}):
\[
\xymatrix@C=3cm@R=1.5cm{
G \ar[r]^{u} \ar[d]_{u}& H \ar[d]^{i_2}\\
H \ar[r]^{i_1}& K . }
\]
Applying the ``associated sheaf'' functor, one therefore obtains a
cartesian and cocartesian diagram in the category of sheaves
(\hyperref[II.theo4.1]{4.1} 1)):
\[\vcenter{
\xymatrix@C=3cm@R=1.5cm{
G \ar[r]^{u} \ar[d]_{u}& H \ar[d]^{\underline{a}i_2}\\
H \ar[r]^{\underline{a}i_1}& \underline{a}K . }}\leqno{(*)}
\]
Since $u$ is an epimorphism of sheaves, the morphism
$\underline{a}i_1$ is an isomorphism, and since diagram (*) is
cartesian, the morphism $u$ is an isomorphism.
\end{enonce*}

\begin{enonce*}{Proposition 4.3}\etiquette[4.3]{II.prop4.3}
\begin{enumerate}
\renewcommand{\labelenumi}{\theenumi)}
\item The inductive limits in $C^{\widetilde{\phantom{,}}}$ that are
representable are universal {\rm(\hyperrefext[I][I.def2.5]{2.5}).}

\item Every epimorphic family {\rm(\hyperrefext[I][I.sec10.2]{10.2})}
of morphisms is universally effective epimorphic
(\hyperref[II.prop2.6]{2.6}).

\item Every equivalence relation is universally effective {\rm(I
\hyperref[I.sec10.6]{10.6}).}

\item Filtered $𝒰$-inductive limits commute with finite projective
limits {\rm(\hyperrefext[I][I.prop2.6]{2.6})}.
\end{enumerate}
\end{enonce*}

\begin{enonce*}[remark]{Proof}
Assertions\pageoriginale 1) through 4) hold in the category of sets,
hence in the category of presheaves
(\hyperrefext[I][I.prop3.1]{3.1}). Assertions 1) and 4) then follow
immediately from the corresponding assertions for presheaves and from
\hyperref[II.theo4.1]{4.1}. Let us prove 3). Let $p_1$, $p_2:$
$\xymatrix{R\ar@<2pt>[r] \ar@<-2pt>[r]& X}$ be an equivalence
relation in $C^{\widetilde{\phantom{,}}}$. The diagram
$\xymatrix{R \ar@<2pt>[r]^{p_1} \ar@<-2pt>[r]_{p_2}&X}$ is then an
equivalence relation of presheaves
(\hyperref[II.cor4.1.1]{4.1.1}). There is therefore a morphism of
presheaves $u:X\rightarrow Y$ such that the diagram
\[
\xymatrix@C=3cm@R=1.25cm{
R \ar[r]^{p_2} \ar[d]_{p_1}& X\ar[d]^{u}\\
X \ar[r]^{u}& Y }
\]
is cartesian and cocartesian in the category of presheaves. Applying
the ``associated sheaf'' functor, one obtains a cartesian and
cocartesian diagram in the category of sheaves
(\hyperref[II.theo4.1]{4.1} 1)). It follows that
$\xymatrix{R \ar@<2pt>[r]^{p_1} \ar@<-2pt>[r]_{p_2}&X}$ is an
effective equivalence relation. Since every equivalence relation is
effective by what precedes, every equivalence relation is universally
effective. Let us prove 2). Let $(u_i:X_i\rightarrow X)$, $i\in I$,
be an epimorphic family in $C^{\widetilde{\phantom{,}}}$. By
(II\ \hyperref[II.prop2.6]{2.6}) and
(\hyperrefext[I][I.prop2.12]{2.12}), it suffices to prove the following
assertions:
\begin{enumerate}
\renewcommand{\theenumi}{\alph{enumi}}
\renewcommand{\labelenumi}{\theenumi)}
\item For every morphism of sheaves $v:Y\rightarrow X$, the family
of morphisms $\pr_{2,i}:X_i\times_X Y\rightarrow Y$, $i\in I$, is an
epimorphic family in $C^{\widetilde{\phantom{,}}}$.

\item For every sheaf $Z$, the following diagram of sets is exact:
$$
\Hom(X,Z)\rightarrow \prod_{i\in I} \Hom(X_i,Z)\rightrightarrows
\prod_{(i,k)\in I\times I}\Hom(X_i\times_X X_k,Z)
$$
\end{enumerate}
\end{enonce*}

Let\pageoriginale $X'$ be the union of the images of the morphisms
$u_i$, formed in the category of presheaves, and let
$j:X'\hookrightarrow X$ be the canonical injection. For every
$i\in I$, the morphism $u_i:X_i\rightarrow X$ factors as a morphism
$u'_i:X_i\rightarrow X'$ followed by $j$, and the family
$u'_i:X_i\rightarrow X'$, $i\in I$, is an epimorphic family in
$C^{\widehat{\phantom{,}}}$. Since $j$ is a monomorphism, the morphism
$\underline{a}j:\underline{a}X'\rightarrow X$ is a monomorphism of
sheaves (\hyperref[II.theo4.1]{4.1}). Since, for every $i\in I$, one
has $u_i=\underline{a}j\underline{a}u'_i$, $\underline{a}j$ is an
epimorphism of sheaves. It follows that $\underline{a}j$ is an
isomorphism (\hyperref[II.prop4.2]{4.2}). Let $v:Y\rightarrow X$ be a
morphism of sheaves. Base change then yields a morphism
$j_v:X'\times_X Y\rightarrow Y$ and morphisms
$u'_{i,v}:X_i\times_X Y\rightarrow X'\times_X Y$. Since
$\underline{a}$ commutes with fiber products, $\underline{a}j_v$ is
an isomorphism. Since epimorphic families in
$C^{\widehat{\phantom{,}}}$ remain epimorphic under base change, the
family $u'_{i,v}$, $i\in I$, is epimorphic in
$C^{\widehat{\phantom{,}}}$. Since $\underline{a}$ commutes with
inductive limits, the family
$\underline{a}u'_{i,v}:X_i\times_X Y\rightarrow
\underline{a}(X'\times_X Y)$, $i\in I$, is epimorphic in
$C^{\widetilde{\phantom{,}}}$, proving a). Let us prove b). Let $Z$
be a sheaf. Since $\underline{a}j$ is an isomorphism, the map
$$
\Hom(X,Z)\longrightarrow \Hom(X',Z)
$$
is a bijection. Since the $u'_i$, $i\in I$, form an epimorphic family
in $C^{\widehat{\phantom{,}}}$, the following diagram of sets is
exact:
$$
\Hom(X',Z)\rightarrow \prod_{i\in I} \Hom(X_i,Z)\rightrightarrows
\prod_{(i,k)\in I\times I}\Hom(X_i\times_{X'}X_k, Z)
$$
Finally, since $X'$ is a subpresheaf of $X$, the presheaf
$X_i\times_{X'}X_k$, $(i,k)\in I\times I$, is canonically isomorphic
to $X_i\times_X X_k$, proving b).

\begin{enonce*}[remark]{Remarks 4.3.1}\etiquette[4.3.1]{II.rem4.3.1}
\begin{enumerate}
\renewcommand{\labelenumi}{\theenumi)}
\item Proposition\pageoriginale \hyperref[II.prop4.3]{4.3} 2)
formally gives us a second proof of \hyperref[II.prop4.2]{4.2}: it is
clear that, in any category, a morphism that is both a universal
effective epimorphism and a monomorphism is in fact an isomorphism.

\item It follows from the preceding propositions that every morphism
in the category of sheaves factors uniquely as an effective
epimorphism followed by an effective monomorphism. This property
naturally generalizes, to nonadditive categories, axiom (AB2) for
abelian categories (coim $\simeq$ im) [Tohoku].
\end{enumerate}
\end{enonce*}

\setcounter{subsection}{4} \setcounter{subsubsection}{-1}

% END INCLUDED SGA4_English_current_buildable_source_20260906/source/02/drafts/SGA4_II_LINES_765_975_English.texfrag


% BEGIN INCLUDED SGA4_English_current_buildable_source_20260906/source/02/drafts/SGA4_II_LINES_976_1126_English.texfrag
\subsubsection{}\etiquette[4.4.0]{II.subsubsec4.4.0}
Let $C$ be a $𝒰$-site. Composing the functor $h:C\rightarrow
C^{\widehat{\phantom{,}}}$ (\hyperrefext[I][I.sec1.3.1]{1.3.1}) with
the associated sheaf functor yields a functor$$
\epsilon_C:C\longrightarrow C^{\widetilde{\phantom{,}}}{\quoi},
$$
called the \emph{canonical functor} from $C$ to
$C^{\widetilde{\phantom{,}}}$, which will be used constantly in what
follows. The functor $\epsilon_C$ commutes with finite projective
limits. When the topology on $C$ is coarser than the canonical
topology, $\epsilon_C$ is fully faithful and commutes with projective
limits; moreover, it is then defined even when $C$ is not assumed to
possess a small topologically generating family. We shall
not study in detail the behavior of $\epsilon_C$ with respect to
inductive limits (\cf [SGA 3 IV]). We shall, however, need the results
below.

\begin{enonce*}{Theorem 4.4}\etiquette[4.4]{II.theo4.4}
Let $C$ be a $𝒰$-site and let $(u_i:X_i\rightarrow X)$, $i\in
I$, be a family of morphisms of $C$ with target $X$. The following
conditions are equivalent \footnote[1]{When the site $C$ does not
necessarily possess a small topologically generating family and when
the topology on $C$ is coarser than the canonical topology, one has
ii) $\Rightarrow$ i) (\cf the proof of
\hyperref[II.theo4.4]{4.4}).}:
\begin{enumerate}
\renewcommand{\theenumi}{\roman{enumi}}
\renewcommand{\labelenumi}{\theenumi)}
\item The family $(\epsilon_C u_i:\epsilon_C X_i\rightarrow \epsilon_C
  X)$, $i\in I$, is an epimorphic family in
  $C^{\widetilde{\phantom{,}}}$ {\rm(\hyperrefext[I][I.sec10.2]{10.2}).}\pageoriginale

\item The family $(u_i:X_i\rightarrow X)$, $i\in I$, is a covering
  family of $C$ (\hyperref[II.def1.2]{1.2}).
\end{enumerate}
\end{enonce*}

\begin{enonce*}[remark]{Proof}
ii) $\Rightarrow$ i). Let $R\hookrightarrow X$ be the sieve generated
by the $u_i$, $i\in I$ (\hyperrefext[I][I.sec4.3.3]{4.3.3}), and let
$u'_i:X_i\rightarrow R$ be the morphisms induced by the $u_i$. The
morphisms $u'_i$, $i\in I$, form an epimorphic family in
$C^{\widehat{\phantom{,}}}$, and $R$ is a covering sieve of $X$.
Consequently, for every sheaf $F$, the map
$$
\Hom(X,F)\longrightarrow \Hom(R,F)
$$
is bijective, and the map
$$
\Hom(R,F)\longrightarrow \prod_i \Hom(X_i,F)
$$
is injective. Thus the map
$$
\Hom(X,F)\longrightarrow \prod_i\Hom(X_i,F)
$$
is injective, and consequently the map
$$
\Hom(\underline{a}X,F)\longrightarrow \prod_i \Hom(\underline{a}X_i,F)
$$
is injective, which proves i).

\noindent i) $\Rightarrow$ ii). With the notation introduced above,
let $J_R:R\hookrightarrow X$ be the canonical inclusion into $X$ of
the sieve generated by the $u_i$, $i\in I$. It follows from i) that
the morphism of sheaves
$\underline{a}J_R:\underline{a}R\rightarrow \underline{a}X$ is both an
epimorphism and a monomorphism in the category of sheaves. It is
therefore an isomorphism of sheaves (\hyperref[II.prop4.2]{4.2}). With the
notation of ${\rm n^{\circ}}$ 3, we have a commutative diagram:
\[
\xymatrix@C=2.5cm@R=1.5cm@M=.25cm{
X \ar[r]^{\ell(X)} & LX \ar[r]^{\ell(LX)} & \mathsf{\underline{a}}X\\
R \ar[r]^{\ell(R)} \ar@{^{(}->}[u]^{J_R}& LR \ar[r]^{\ell(LR)}
\ar@{^{(}->}[u]^{LJ_R}& \mathsf{\underline{a}}R
\ar[u]^{\wr}_{\mathsf{\underline{a}}J_R}. }
\]
First of\pageoriginale all, by \hyperref[II.lem3.1]{3.1} 2), there exist a covering sieve
$J_{S_1}:S_1\rightarrow X$ and a morphism
$u_1:S_1\rightarrow LR$ such that
\begin{equation}
\ell(LX)\circ \ell (X)\circ J_{S_1}=\underline{a}J_R\circ \ell (LR)\circ
u_1=\ell(LX)\circ LJ_R\circ u_1{\quoi}.\tag{4.6.1}
\end{equation}\etiquette[4.6.1]{II.eq4.6.1}
Put $m=\ell(X)\circ J_{S_1}$ and $n=LJ_R\circ u_1$. The morphisms
$m$ and $n$ are morphisms from $S_1$ to $LX$. Let
$S_2\hookrightarrow S_1$ be the kernel of the pair of arrows $(m,n)$.
For every object $Y$ of $C$ and every morphism $\alpha:Y\rightarrow S_1$,
one has, by \hyperref[II.eq4.6.1]{4.6.1}, $\ell(LX)\circ m\circ \alpha
=\ell(LX)\circ n\circ\alpha$. It then follows from
(\hyperref[II.lem3.1]{3.1} 3)) that the kernel $S_2\times_{S_1}Y$ of
the pair of arrows $(m\circ \alpha,n\circ \alpha)$ is a covering
sieve of $Y$. Consequently, by axiom (T 2) for topologies,
$J_{S_2}:S_2\hookrightarrow X$ is a covering sieve of $X$. Thus there
is a morphism $u_2:S_2\rightarrow LR$ and a commutative diagram:
\hypertarget{sga4.exp.II.equation.4.6.2}{}\[\vcenter{
\xymatrix@C=2.5cm@R=.8cm@M=.2cm{
S_2 \ar@{^{(}->}[rr]^{J_{S_{2}}} \ar[dd]_{u_2}&  &X \ar[dd]^{\ell(X)}\\
                  &R\ar@{^{(}->}[ru]_{J_R} \ar[ld]^{\ell(R)}&\\
LR \ar[rr]^{LJ_R}  &  &LX&. }}\leqno{(4.6.2)}
\]

Let $Y$ be an object of $C$ and let $\beta:Y\rightarrow S_2$ be a
morphism. By \hyperref[II.lem3.1]{3.1} 2) and 1), there exist a
covering sieve $J_{Q_1}:Q_1\hookrightarrow Y$ and a morphism
$v_1:Q_1\rightarrow R$ such that
$u_2\circ\beta\circ J_{Q_1}=\ell(R)\circ v_1$. Put
$f=J_{S_{2}}\circ \beta\circ J_{Q_1}$ and $g=J_R\circ v_1$. Let
$Q_2$ be the kernel of the pair $(f,g):Q_1\rightrightarrows X$, and
let $J_{Q_2}:Q_2\hookrightarrow Y$ be the canonical inclusion. For
every object $Z$ of $C$ and every morphism $\gamma:Z\rightarrow Q_1$,
the commutativity of (\hyperref[II.cor4.6.2]{4.6.2}) gives:
\begin{gather*}
\ell(X)\circ f\circ \gamma =\ell(X)\circ J_{S_2}\circ \beta\circ
J_{Q_1}\circ \gamma =LJ_R\circ u_2\circ \beta\circ J_{Q_1}\circ \gamma=\\
=LJ_R\circ \ell(R)\circ v_1\circ \gamma =\ell (X)\circ J_R\circ
v_1\circ \gamma = \ell(X)\circ g\circ \gamma{\quoi}.
\end{gather*}
It\pageoriginale then follows from \hyperref[II.lem3.1]{3.1} 3) that the
kernel $Q_2\times_{Q_1}Z$ of the pair $(f\circ\gamma,g\circ\gamma)$
is a covering sieve of $Z$. By axiom (T 2) for topologies,
$J_{Q_2}:Q_2\hookrightarrow Y$ is a covering sieve of $Y$. Thus, for
every morphism $\beta:Y\rightarrow S_2$, there exist a covering sieve
$J_{Q_2}:Q_2\hookrightarrow Y$ and a morphism $v_2:Q_2\rightarrow
R$ such that the following diagram is commutative:
\[
\xymatrix@C=2.5cm@R=1.5cm@M=.2cm{
Y\ar[r]^{\beta}   & S_2\ar[r]^{J_{S_2}} & X\\
Q_2 \ar[r]^{v_2} \ar@{^{(}->}[u]^{J_{Q_2}}& R
\ar@{^{(}->}[ru]_{J_{R}}  &}
\]

Denote by $S_2\cap R$ the fiber product of $S_2$ and $R$ over $X$.
The sieve (of $Y$) $(S_2\cap R)\times_{S_2}Y$ contains the sieve
$Q_2$, and hence $(S_2\cap R)\times_{S_2}Y$ is a covering sieve of
$Y$. It then follows from axiom (T 2) for topologies that
$S_2\cap R$ is a covering sieve of $X$, and consequently the sieve
$R$, which contains $S_2\cap R$, is a covering sieve of $X$, \cqfd
\end{enonce*}

\setcounter{subsubsection}{3}
\begin{enonce*}{Corollary 4.4.4}\etiquette[4.4.4]{II.cor4.4.4}
Let $T$ be the topology of the $\Uu$-site $C$, and let $T'$ (\resp
$T''$) be the finest topology on $C$ among those for which the
objects of $\widetilde{C}$ are sheaves \emph{(\resp
separated presheaves)} (\hyperref[II.prop2.2]{2.2}). Then $T=T'=T''$.
\end{enonce*}

Indeed, one trivially has $T\leq T'\leq T''$, and
\hyperref[II.theo4.4]{4.4} and \hyperref[II.prop2.2]{2.2} imply that
$T''\leq T$, \cqfd

% END INCLUDED SGA4_English_current_buildable_source_20260906/source/02/drafts/SGA4_II_LINES_976_1126_English.texfrag


% BEGIN INCLUDED SGA4_English_current_buildable_source_20260906/source/02/drafts/SGA4_II_LINES_1127_1298_English.texfrag
\begin{enonce*}{Definition 4.5}\etiquette[4.5]{II.def4.5}
An initial object of a category $A$ is an object $\phi_A$ that
represents the empty inductive limit, \ie such that, for every $X\in
\ob A$, there exists a unique arrow $\phi_A\rightarrow X$. An object
$\phi_A$ is called \emph{strict initial} if it is initial and every
morphism with target $\phi_A$ is an isomorphism. Let $(S_i)$, $i\in I$,
be a family of objects of a category $A$. Suppose that the sum
$S=\amalg_{i\in I} S_i$ is representable. The sum $S$ is called
disjoint if the structure morphisms $S_i\rightarrow S$ are
pullbackable, if they are monomorphisms, and if, for every pair $i$,
$j$ with $i\neq j$, the fiber product $S_i\times_S S_j$ is an initial
object of $A$. The sum $S$ is called disjoint and
universal\pageoriginale if it is disjoint and remains a disjoint sum
after every base change $T\rightarrow S$; it follows that, for every
pair $i$, $j$ with $i\neq j$, the objects $S_i\times_S S_j$ are strict
initial objects of $A$.
\end{enonce*}

\begin{enonce*}[remark]{Example 4.5.1}\etiquette[4.5.1]{II.exem4.5.1}
In the category of sets, sums are disjoint and universal. The same is
therefore true in every category of presheaves of sets
(\hyperrefext[I][I.prop3.1]{3.1}), and hence in every category of
sheaves of sets on a $𝒰$-site (\hyperref[II.theo4.1]{4.1}); in
particular, the initial object of $C^{\widetilde{\phantom{,}}}$ is
strict.
\end{enonce*}

\begin{enonce*}{Proposition 4.6}\etiquette[4.6]{II.prop4.6}
Let $C$ be a $𝒰$-category and let $(s_i:X_i\rightarrow X)$,
$i\in I$, be a family of morphisms of $C$. For every
$𝒰$-topology $\Tt$ on $C$ (\hyperref[II.def3.0.2]{3.0.2}), denote by
$C^{\widetilde{\phantom{,}}}_{\Tt}$ the corresponding category of
sheaves and by $\epsilon_{\Tt}:C\rightarrow
C^{\widetilde{\phantom{,}}}_{\Tt}$ the corresponding functor.
\begin{enumerate}
\renewcommand{\labelenumi}{\theenumi)}
\item Let $\Tt$ be a $𝒰$-topology such that
$$
\coprod_{i\in I}
\epsilon_{\Tt}(X_i)\xrightarrow{(\epsilon_{\Tt}(s_i))}\epsilon_{\Tt}(X)
$$
is an isomorphism. Then, for every topology $\Tt'$ finer than $\Tt$,
the morphism
$$
\coprod_{i\in
  I}\epsilon_{\Tt'}(X_i)\xrightarrow{(\epsilon_{\Tt'}(s_i))}\epsilon_{\Tt'}(X)
$$
is an isomorphism.

\item Let $\Tt$ be a $𝒰$-topology on $C$. The following
  properties {\rm(i)} and {\rm(ii)} are equivalent:
\begin{itemize}
\item[i)]
\begin{itemize}
\item[a)] The family $(s_i:X_i\rightarrow X)$, $i\in I$, is covering
  for $\Tt$.

\item[b)] For every $i\in I$, the diagonal morphism of presheaves
  $\Delta_i:X_i\hookrightarrow X_i\times_X X_i$\pageoriginale
  is sent by the associated sheaf functor (for $\Tt$) to an
  isomorphism (as is the case if the $s_i$ are monomorphisms).

\item[c)] For every pair $(i,j)$ of elements of $I$ with $i\neq j$,
  the presheaf $X_i\times_X X_j$ is sent by the associated sheaf
  functor (for $\Tt$) to the initial object of
  $C^{\widetilde{\phantom{,}}}_{\Tt}$.
\end{itemize}
\item[ii)] $\epsilon_{\Tt}(X)$ is the sum of the
  $\epsilon_{\Tt}(X_i)$.
\end{itemize}
\end{enumerate}
\end{enonce*}

\begin{enonce*}[remark]{Proof}
\begin{enumerate}
\renewcommand{\labelenumi}{\theenumi)}
\item A sheaf for $\Tt'$ is a sheaf for $\Tt$. Consequently, for
  every sheaf $F$ for $\Tt'$, the morphism
$$
\Hom_{C^{\widehat{\phantom{,}}}}(X,F)\longrightarrow \prod_{i\in I}
\Hom_{C^{\widehat{\phantom{,}}}}(X_i,F)
$$
is an isomorphism, which proves the assertion.

\item By definition, property ii) holds if and only if the morphism
  of presheaves
$$
\Phi=(h(s_i),i\in I):\coprod_i h(X_i)\longrightarrow h(X)
$$
is sent by the associated sheaf functor to an isomorphism, \ie
(\hyperref[II.prop4.2]{4.2}) if and only if $\underline{a}(\Phi)$ is
both an epimorphism and a monomorphism of sheaves. By
(\hyperref[II.theo4.4]{4.4}), the morphism $\underline{a}(\Phi)$ is an
epimorphism if and only if property a) holds. The diagonal morphism
$$
\Delta:\coprod_i h(X_i)\longrightarrow (\coprod_i
h(X_i))\times_{h(X)}(\coprod_i h(X_i))
$$
is the direct sum of a family $\Delta_{i,j}$, $(i,j)\in I\times I$,
of morphisms defined as follows:
\begin{itemize}
\item[$\beta)$] When $i=j$, $\Delta_{i,i}$ is the diagonal morphism
  $h(X_i)\rightarrow h(X_i)\times_{h(X)}h(X_i)$.

\item[$\gamma)$] When $i\neq j$, $\Delta_{i,j}$ is the morphism
  $\phi_{C^{\widehat{\phantom{,}}}}\rightarrow
  h(X_i)\times_{h(X)}h(X_j)$ ($\phi_{C^{\widehat{\phantom{,}}}}$
  denotes the initial object of $C^{\widehat{\phantom{,}}}$).
\end{itemize}
\end{enumerate}
\end{enonce*}

The\pageoriginale morphism $\underline{a}(\Phi)$ is a monomorphism if
and only if $\underline{a}(\Delta)$ is an isomorphism, and by
(\hyperref[II.theo4.1]{4.1}) $\underline{a}(\Delta)$ is an isomorphism
if and only if the $\underline{a}(\Delta_{i,j})$ are isomorphisms, i.e.
if and only if properties b) and c) hold.

\begin{enonce*}{Corollary 4.6.1}\etiquette[4.6.1]{II.cor4.6.1}
Let $C$ be a $𝒰$-category and $X$ an object of $C$.
\begin{enumerate}
\renewcommand{\labelenumi}{\theenumi)}
\item Let $\Tt$ be a $𝒰$-topology on $C$. The object $X$ of $C$ is
  sent by the associated sheaf functor (for the topology $\Tt$) to
  the initial object of $C^{\widetilde{\phantom{,}}}_{\Tt}$ if and
  only if the empty sieve covers $X$.

\item When $X$ is a strict initial object of $C$, the sheaf associated
  with $X$, for every $𝒰$-topology finer than the canonical topology,
  is an initial object.
\end{enumerate}
\end{enonce*}

\begin{enonce*}[remark]{Proof}
\begin{enumerate}
\renewcommand{\labelenumi}{\theenumi)}
\item Take the indexing set $I$ in \hyperref[II.prop4.6]{4.6} 2) to
be the empty set.

\item By 1) and \hyperref[II.prop4.6]{4.6} 1), it is enough to show
  that the empty sieve covers $X$ for the canonical topology; \ie by
  (\hyperref[II.prop2.6]{2.6}), it is enough to show that, for every
  object $Y\rightarrow X$ over $X$, $Y$ is an initial object of $C$,
  which follows from the definition (\hyperref[II.def4.5]{4.5}).
\end{enumerate}
\end{enonce*}

\begin{enonce*}{Corollary 4.6.2}\etiquette[4.6.2]{II.cor4.6.2}
Let $C$ be a $𝒰$-site and let $(s_i:X_i\rightarrow X)$, $i\in I$,
be a family of pullbackable morphisms of $C$ with the same target and
having the following properties:
\begin{itemize}
\item[$\alpha)$] The family of the $s_i$, $i\in I$, is covering.

\item[$\beta)$] For every $i\in I$, $s_i:X_i\rightarrow X$ is a
  monomorphism.

\item[$\gamma)$] For every pair $(i,j)$ of elements of $I$ with
  $i\neq j$, $X_i\times_X X_j$ is covered by the empty sieve, \ie for
  every sheaf $F$ on $C$, $F(X_i\times_X X_j)$ is a singleton set.
\end{itemize}
Then $\epsilon_C(X)$ is the sum of the $\epsilon_C(X_i)$, $i\in I$.
\end{enonce*}

\begin{enonce*}[remark]{Proof}
This follows\pageoriginale immediately from
\hyperref[II.prop4.6]{4.6} 2) and
\hyperref[II.cor4.6.1]{4.6.1} 1).
\end{enonce*}

% END INCLUDED SGA4_English_current_buildable_source_20260906/source/02/drafts/SGA4_II_LINES_1127_1298_English.texfrag


% BEGIN INCLUDED SGA4_English_current_buildable_source_20260906/source/02/drafts/SGA4_II_LINES_1299_1528_English.texfrag
\begin{enonce*}{Corollary 4.6.3}\etiquette[4.6.3]{II.cor4.6.3}
Let $C$ be a $𝒰$-category and let $(s_i:X_i\rightarrow X)$,
$i\in I$, be a family of pullbackable morphisms with the same target.
Suppose that the canonical topology of $C$ is a $𝒰$-topology. The
following conditions are equivalent:
\begin{enumerate}
\renewcommand{\theenumi}{\roman{enumi}}
\renewcommand{\labelenumi}{\theenumi)}
\item There exists a $𝒰$-topology $\Tt$ on $C$, coarser than the
canonical topology, such that $\epsilon_\Tt(X)$ is the sum of the
$\epsilon_\Tt(X_i)$, $i\in I$.

\item The object $X$ is the disjoint and universal sum
(\hyperref[II.def4.5]{4.5}) (in $C$) of the $X_i$.
\end{enumerate}
\end{enonce*}

\begin{enonce*}[remark]{Proof}
ii) $\Rightarrow$ i). Since $X$ is the universal sum of the $X_i$,
$i\in I$, the family of the $s_i$, $i\in I$, is covering for the
canonical topology of $C$ (\hyperref[II.prop2.6]{2.6}). Condition
$\alpha$) of \hyperref[II.cor4.6.2]{4.6.2} is therefore satisfied when
$C$ is endowed with the canonical topology. Condition $\beta$) is
evidently satisfied, and property $\gamma$) follows from
\hyperref[II.prop4.6]{4.6} 1) 2), whence i).

\noindent i) $\Rightarrow$ ii). Let $\Tt$ be a topology on $C$ coarser
than the canonical topology such that $\epsilon_\Tt(X)$ is the sum of
the $\epsilon_\Tt(X_i)$. For every sheaf $F$ for $\Tt$, the canonical
map $\Hom(X,F)\rightarrow\prod_i\Hom(X_i,F)$ is a bijection. Moreover,
every representable presheaf is a sheaf. It follows immediately that
$X$ is the sum in $C$ of the $X_i$, $i\in I$. Applying this to the
case where the set $I$ is empty, one sees that if an object of $C$ is
carried to the initial object of $C^{\widetilde{\phantom{,}}}$, that
object is an initial object of $C$. The functor
$\epsilon_\Tt:C\rightarrow C^{\widetilde{\phantom{,}}}$ is fully
faithful and commutes with finite projective limits. Consequently,
condition b) of \hyperref[II.cor4.6.2]{4.6.2} implies that the $s_i$,
$i\in I$, are monomorphisms of $C$, and condition c) implies, by what
precedes, that the $X_i\times_XX_j$, $i\neq j$, are initial objects of
$C$. Consequently $X$ is the disjoint sum of the\pageoriginale $X_i$.
Since $\epsilon_\Tt$ commutes with fiber products, this last property
is stable under every base change $Y\rightarrow X$, and therefore $X$
is the disjoint and universal sum of the $X_i$, $i\in I$.
\end{enonce*}

\begin{enonce*}{Corollary 4.6.4}\etiquette[4.6.4]{II.cor4.6.4}
Let $C$ be a $𝒰$-category and $X$ an object of $C$. Suppose that the
canonical topology on $C$ is a $𝒰$-topology. The following properties
are equivalent:
\begin{enumerate}
\renewcommand{\theenumi}{\roman{enumi}}
\renewcommand{\labelenumi}{\theenumi)}
\item There exists on $C$ a topology coarser than the canonical
topology such that $\epsilon_\Tt(X)$ is an initial object of
$C^{\widetilde{\phantom{,}}}_\Tt$.

\item $X$ is a strict initial object of $C$.
\end{enumerate}
\end{enonce*}

\begin{enonce*}[remark]{Proof}
Take the empty set for the set $I$ in \hyperref[II.cor4.6.3]{4.6.3}.
\end{enonce*}

\begin{enonce*}{Proposition 4.7}\etiquette[4.7]{II.prop4.7}
Let $C$ be a $𝒰$-category, let $C^{\widetilde{\phantom{,}}}$ be the
category of sheaves on $C$ for the canonical topology, let
$\epsilon_C:C\rightarrow C^{\widetilde{\phantom{,}}}$ be the canonical
functor (\hyperref[II.subsubsec4.4.0]{4.4.0}), and let
$R\rightrightarrows X$ be an equivalence relation in $C$ admitting a
cokernel $Y$ in $C$. Let $\Pi:X\rightarrow Y$ be the canonical
morphism and suppose it is pullbackable. Consider the following
properties:
\begin{enumerate}
\renewcommand{\theenumi}{\roman{enumi}}
\renewcommand{\labelenumi}{\theenumi)}
\item $\epsilon_C(Y)$ is the quotient of the equivalence relation
$\epsilon_C(R)\rightrightarrows\epsilon_C(X)$.

\item The equivalence relation $R\rightrightarrows X$ is universally
effective {\rm(cf. ${\rm n^{\circ}}$ 7).}
\end{enumerate}
One has {\rm{ii) $\Rightarrow$ i).}} When $C$ possesses a
$𝒰$-topology coarser than the canonical topology, one has
{\rm{i) $\Rightarrow$ ii).}}
\end{enonce*}

\begin{enonce*}[remark]{Proof}
ii) $\Rightarrow$ i). The canonical morphism
$R\rightarrow X\times_YX$ is an isomorphism. Let $S\rightarrow Y$ be
the sieve generated by $\pi:X\rightarrow Y$. For every presheaf $F$,
there is an exact diagram (\hyperrefext[I][I.prop2.12]{2.12}):
$$
\Hom(S,F)\rightarrow \Hom(X,F)\rightrightarrows \Hom(R,F){\quoi}.
$$
It\pageoriginale therefore suffices to show that, for every sheaf $F$
for the canonical topology, the map
$$
\Hom(X,F)\longrightarrow \Hom(S,F)
$$
is a bijection, \ie it suffices to show that $S$ is covering for the
canonical topology, or equivalently that the sieve generated by
$\pi:X\rightarrow Y$ is covering, which follows from \hyperref[II.def2.5]{Definition 2.5}.

\noindent i) $\Rightarrow$ ii). The functor
$\epsilon_C:C\rightarrow C^{\widetilde{\phantom{,}}}$ commutes with
finite projective limits and is fully faithful. The canonical morphism
$\epsilon_C(R)\rightarrow J_C(X)\times_{J_C(Y)}J_C(X)$ is an
isomorphism (\hyperref[II.prop4.3]{4.3}). Consequently the canonical
morphism $R\rightarrow X\times_YX$ is an isomorphism. The morphism
$\epsilon_C(X)\rightarrow\epsilon_C(Y)$ is an epimorphism, which
implies (\hyperref[II.theo4.4]{4.4}) that $\pi:X\rightarrow Y$ is a
covering morphism of $C$ for the canonical topology, \cqfd
\end{enonce*}

\begin{enonce*}{Proposition 4.8}\etiquette[4.8]{II.prop4.8}
Let $C$ be a $𝒰$-site. The category of sheaves on $C$ has the
following properties:
\begin{enumerate}
\renewcommand{\theenumi}{\alph{enumi}}
\renewcommand{\labelenumi}{\theenumi)}
\item Finite projective limits are representable.

\item Direct sums indexed by an element of $𝒰$ are representable.
They are disjoint and universal (\hyperref[II.def4.5]{4.5}).

\item Equivalence relations are universally effective.
\end{enumerate}
\end{enonce*}

This was proved in \hyperref[II.theo4.1]{4.1}, 2) and 3),
\hyperref[II.prop4.3]{4.3}, 3), and
\hyperref[II.exem4.5.1]{4.5.1}.

These properties have been singled out because they will later allow
us to characterize the $𝒰$-categories equivalent to categories of
sheaves on categories belonging to $𝒰$
(\hyperrefext[IV][IV.theo1.2]{1.2}).

\begin{enonce*}[remark]{Remark}
Recall that property a) is equivalent to the property:

There exists a terminal object, and fiber products are representable
(\hyperrefext[I][I.sec2.3.1]{2.3.1}).
\end{enonce*}

\setcounter{subsection}{8}
\subsection{}\etiquette[4.9]{II.subsec4.9}
Let\pageoriginale $C$ be a site whose underlying category is a
$𝒰$-category. Then the category $\widetilde{C}$ of $𝒰$-sheaves on
$C$ satisfies the conditions considered in
\hyperrefext[I][I.cor7.3]{7.3} ((i) or (ii), as desired), and therefore,
by \loccit, the various variants of the notion of a \emph{generating
family} in $\widetilde{C}$ considered in
\hyperrefext[I][I.def7.1]{7.1} coincide. This being understood:

\begin{enonce*}{Proposition 4.10}\etiquette[4.10]{II.prop4.10}
With the preceding notation, consider the canonical functor
$\epsilon:C\rightarrow\widetilde{C}$
(\hyperref[II.subsubsec4.4.0]{4.4.0}), and let $G\subset\ob C$ be a
topologically generating family in $C$
(\hyperref[II.def3.0.1]{3.0.1}). Then the family
$(\epsilon(X))_{X\in G}$ of objects of $\widetilde{C}$ is a generating
family. In particular, the family $(\epsilon(X))_{X\in\ob C}$ of
objects of $\widetilde{C}$ is generating.
\end{enonce*}

We must prove that every morphism $u:F\rightarrow F'$ in
$\widetilde{C}$ such that
\begin{equation}
\Hom(\epsilon(X),F)\longrightarrow \Hom(\epsilon(X),F')\tag{4.10.1}
\end{equation}\etiquette[4.10.1]{II.eq4.10.1}
is a bijection for every $X\in G$, is an isomorphism. The preceding
map identifies with the map
\begin{equation}
u(X):F(X)\longrightarrow F'(X){\quoi},\tag{4.10.2}
\end{equation}\etiquette[4.10.2]{II.eq4.10.2}
and we must show that if the latter is bijective for every $X\in G$,
then it is so for every $X\in\ob C$. Let
$(X_i\rightarrow X)_{i\in I}$ be a covering family of $X$ by objects
of $G$, and, for every pair of indices $(i,j)$ in $I$, consider the
set of all morphisms $X_{ijk}\rightarrow X_i\times_XX_j$ in
$\widehat{C}$, with $X_{ijk}\in G$ ($k$ varying in an index set
$I_{i,j}$). One then obtains a morphism of exact diagrams of sets
\[
\xymatrix{ F(X) \ar[r] \ar[d]& \displaystyle{{\prod_i}} F(X_i)
\ar[d]^{\wr} \ar@<3pt>[r] \ar@<-3pt>[r] & \displaystyle{\prod_{ijk}}
F(X_{ijk}) \ar[d]^{\wr}\\
F'(X) \ar[r] & \displaystyle{{\prod_i}} F'(X_i) \ar@<3pt>[r]
\ar@<-3pt>[r] & \displaystyle{\prod_{ijk}} F'(X_{ijk})&, }
\]
where\pageoriginale, by hypothesis, the last two vertical arrows are
bijections. The same is therefore true of the first vertical arrow,
\cqfd

\begin{enonce*}{Corollary 4.11}\etiquette[4.11]{II.cor4.11}
Suppose that $C$ is a $𝒰$-site (\hyperref[II.def3.0.2]{3.0.2}). Then
the category $\widetilde{C}$ is a $𝒰$-category and admits a
$𝒰$-small generating family.
\end{enonce*}

Indeed, by hypothesis one may take for $G$ in
\hyperref[II.prop4.10]{4.10} a small generating family, which proves
the existence of a small generating family. On the other hand, if
$F,F'\in\ob\widetilde{C}$, a morphism from $F$ to $F'$ is known once
one knows the morphism \eqhyperref[II.eq4.10.1]{4.10.1}, \ie
\eqhyperref[II.eq4.10.2]{4.10.2}, for every $X\in\ob G$ (I \hyperref[I.sec7.1.1]{7.1.1}),
whence the map
$$
\Hom(F,F')\longrightarrow \prod_{X\in G}\Hom(F(X),F'(X))
$$
is injective. Since the second member is $𝒰$-small, so is the first,
which proves that $\widetilde{C}$ is a $𝒰$-category.

\begin{enonce*}{Corollary 4.12}\etiquette[4.12]{II.cor4.12}
Let $C$ be a $𝒰$-site. Then, for every object $X$ of
$\widetilde{C}$, the set of subobjects of $X$ and the set of quotient
objects of $X$ are $𝒰$-small.
\end{enonce*}

This follows from \hyperrefext[I][I.prop7.4]{7.4}, respectively
\hyperrefext[I][I.cor7.5]{7.5}, which apply by
\hyperref[II.cor4.11]{4.11}.


% END INCLUDED SGA4_English_current_buildable_source_20260906/source/02/drafts/SGA4_II_LINES_1299_1528_English.texfrag


% BEGIN INCLUDED SGA4_English_current_buildable_source_20260906/source/02/drafts/SGA4_II_LINES_1529_1645_English.texfrag
\section{Extension of a topology from $C$ to
  $C^{\shat{\;\hbox{}}}$}\etiquette[5]{II.sec5}

\begin{enonce*}{Proposition 5.1}\etiquette[5.1]{II.prop5.1}
Let $C$ be a $𝒰$-site and let $f:H\rightarrow K$ be a morphism of
$C^{\widehat{\phantom{,}}}$. The following conditions are
equivalent:
\begin{itemize}
\item[i)] For every $X\rightarrow K$, with $X\in \ob C$, the image of
  the corresponding morphism $H\times_K X\rightarrow X$ is a
  covering sieve of $X$.

\item[ii)] The morphism $\underline{a}(f):\underline{a}H\rightarrow \underline{a}K$
  between the associated sheaves is an epimorphism\pageoriginale in
  $C^{\widetilde{\phantom{,}}}$.

\item[ii bis)] For every sheaf $F$ on $C$, the map
  $\Hom(K,F)\rightarrow \Hom(H,F)$ induced by $f$ is injective.
\end{itemize}
\end{enonce*}

\begin{enonce*}[remark]{Proof}
It is clear that ii) is equivalent to ii bis).
\end{enonce*}

Let us prove that (ii) $\Rightarrow$ (i). By
\hyperref[II.theo4.4]{4.4}, (i) means that
$u:\underline{a}(H\times_K X)\rightarrow \underline{a}(X)$ is an
epimorphism. Now, since $\underline{a}(H\times_K X)\simeq
\underline{a}(H)\times_{\underline{a}(K)}\underline{a}(X)$, because
$\underline{a}$ commutes with finite $\varprojlim$
(\hyperref[II.theo4.1]{4.1}), the morphism in question is obtained
from $\underline{a}(f):\underline{a}(H)\rightarrow
\underline{a}(K)$ by base change. Since the epimorphisms of
$\widetilde{C}$ are universal, this proves (ii) $\Rightarrow$ (i).
Conversely, since the family of all $X_i\rightarrow K$, with
$X_i\in\ob C$, is epimorphic, the same is true of the family of
$\underline{a}(X_i)\rightarrow \underline{a}(K)$ in $\widetilde{C}$
(\hyperref[II.theo4.1]{4.1}). Hence, to verify that
$\underline{a}(f):\underline{a}(H)\rightarrow \underline{a}(K)$ is
epimorphic, it suffices to do so after every base change of the
preceding kind $\underline{a}(X_i)\rightarrow \underline{a}(K)$,
which proves i) $\Rightarrow$ ii), \cqfd

\begin{enonce*}{Definition 5.2}\etiquette[5.2]{II.def5.2}
\begin{enumerate}
\renewcommand{\labelenumi}{\theenumi)}
\item A morphism $f:H\rightarrow K$ satisfying the three equivalent
  conditions of \hyperref[II.prop5.1]{5.1} is called a
  \emph{covering morphism}. A family of morphisms
  $f_i:H_i\rightarrow K$, $i\in I$, with the same target is called
  \emph{covering} if the corresponding morphism
  $f:\amalg_i H_i\rightarrow K$ is covering.

\item A morphism $f:H\rightarrow K$ is called \emph{bicovering} if it
  is covering and if the diagonal morphism
  $H\rightarrow H\times_K H$ is covering. A family
  $f_i:H_i\rightarrow K$, $i\in I$, with the same target is called
  bicovering if the corresponding morphism
  $f:\amalg_i H_i\rightarrow K$ is bicovering.
\end{enumerate}
\end{enonce*}

By condition i) of \hyperref[II.prop5.1]{5.1}, to say that a family
$(f_i:H_i\rightarrow K)$, $i\in I$, is covering means that, for every
morphism $X\rightarrow K$, with $X\in \ob C$, the image of the family
$H_i\times_K X\rightarrow X$, $i\in I$, is a covering sieve
of\pageoriginale $X$; or equivalently, by ii), that the family of
morphisms $\underline{a}(f_i):\underline{a}H_i\rightarrow
\underline{a}K$ in $C^{\widetilde{\phantom{,}}}$ is epimorphic
(taking into account that the functor $\underline{a}$ commutes with
direct sums (\hyperref[II.theo4.1]{4.1})).\footnote[1]{Note that the
property of a morphism or a family of morphisms of
$C^{\widehat{\phantom{,}}}$ being covering (\resp bicovering) is
stable under base change.}

\begin{enonce*}{Proposition 5.3}\etiquette[5.3]{II.prop5.3}
Let $C$ be a $𝒰$-site and let $f:H\rightarrow K$ be a morphism of
$C^{\widehat{\phantom{,}}}$. The following conditions are
equivalent:
\begin{itemize}
\item[i)] The morphism $f$ is bicovering
  (\hyperref[II.def5.2]{5.2}).

\item[i bis)] The morphism $f$ is covering and, for every object $X$
  of $C$ and every pair of morphisms $u$, $v:X\rightrightarrows H$
  such that $fu=fv$, the kernel of $(u,v)$ is a covering sieve of
  $X$.

\item[ii)] The morphism
  $\underline{a}(f):\underline{a}H\to \underline{a}K$ is an
  isomorphism in $C^{\widetilde{\phantom{,}}}$.

\item[ii bis)] For every sheaf $F$ on $C$, the map
  $\Hom(K,F)\rightarrow \Hom(H,F)$ is a bijection.
\end{itemize}
\end{enonce*}

\begin{enonce*}[remark]{Proof}
The equivalence i) $\Leftrightarrow$ i bis) follows from condition i)
of \hyperref[II.prop5.1]{5.1}, applied to the diagonal morphism
$H\rightarrow H\times_K H$; the equivalence ii) $\Leftrightarrow$ ii
bis) is trivial.

\noindent i) $\Rightarrow$ ii). The morphism
$\underline{a}(f):\underline{a}H\rightarrow \underline{a}K$ is an
epimorphism (\hyperref[II.prop5.1]{5.1}). Since the functor
$\underline{a}$ commutes with the formation of fiber products
(\hyperref[II.theo4.1]{4.1}), the diagonal morphism
$\underline{a}H\rightarrow
\underline{a}H\times_{\underline{a}K}\underline{a}H$ is an
epimorphism (\hyperref[II.prop5.1]{5.1}). Since a diagonal morphism is
always a monomorphism, it is an isomorphism
(\hyperref[II.prop4.2]{4.2}). It follows that the morphism
$\underline{a}(f)$ is a monomorphism. It is therefore an isomorphism
(\hyperref[II.prop4.2]{4.2}).

\noindent ii) $\Rightarrow$ i). Since the functor $\underline{a}$
commutes with the formation of fiber products, the morphism $f$ and
the diagonal morphism $H\rightarrow H\times_K H$ are carried by
$\underline{a}$ to isomorphisms. In particular, they are carried by
$\underline{a}$ to epimorphisms. They are therefore covering, \cqfd
\end{enonce*}

\setcounter{subsection}{3} \setcounter{subsubsection}{0}

% END INCLUDED SGA4_English_current_buildable_source_20260906/source/02/drafts/SGA4_II_LINES_1529_1645_English.texfrag


% BEGIN INCLUDED SGA4_English_current_buildable_source_20260906/source/02/drafts/SGA4_II_LINES_1646_1819_English.texfrag
\subsubsection{}\etiquette[5.3.1]{II.subsubsec5.3.1}
It\pageoriginale follows from \hyperref[II.prop5.3]{5.3}, and from
the fact that $\underline{a}$ commutes with direct sums, that a family
$(f_i:H_i\rightarrow K)$, $i\in I$, of morphisms of
$C^{\widehat{\phantom{,}}}$ is bicovering if and only if the $f_i$
induce an isomorphism from the direct sum
$\amalg_i \underline{a}H_i$ onto $\underline{a}K$ on the associated
sheaves, or equivalently if and only if, for every sheaf $F$, the map
$$
\Hom(K,F)\longrightarrow \prod_i \Hom(H_i,F)
$$
is bijective.

\begin{enonce*}{Proposition 5.4}\etiquette[5.4]{II.prop5.4}
Let $C$ be a $𝒰$-site. There exists a (necessarily unique) topology
on $C^{\widehat{\phantom{,}}}$ such that a family
$H_i\rightarrow K$ of arrows of $C^{\widehat{\phantom{,}}}$ with the
same target is covering for this topology if and only if it is
covering in the sense of \hyperref[II.def5.2]{5.2}. It is also the
coarsest topology on $C^{\widehat{\phantom{,}}}$ among the topologies
$T$ having the following properties:
\begin{enumerate}
\renewcommand{\theenumi}{\alph{enumi}}
\renewcommand{\labelenumi}{\theenumi)}
\item $T$ is finer than the canonical topology of
  $C^{\widehat{\phantom{,}}}$ (\ie every epimorphic family in
  $C^{\widehat{\phantom{,}}}$ is covering for $T$).

\item Every covering family in $C$ is covering in
  $C^{\widehat{\phantom{,}}}$.
\end{enumerate}
\end{enonce*}

\begin{enonce*}[remark]{Proof}
We shall confine ourselves to giving indications. We leave it to the
reader to show that the families which are covering in the sense of
\hyperref[II.def5.2]{5.2} are the covering families of a topology
$T_C$ on $C^{\widehat{\phantom{,}}}$ (one uses
\hyperref[II.theo4.1]{4.1}). The covering families of the canonical
topology on $C^{\widehat{\phantom{,}}}$ are the epimorphic families
in $C^{\widehat{\phantom{,}}}$ (\hyperref[II.prop2.6]{2.6} and I
\hyperref[I.prop3.1]{3.1}). Since $\underline{a}$ commutes with inductive limits
(\hyperref[II.theo4.1]{4.1}), the topology $T_C$ is finer than the
canonical topology of $C^{\widehat{\phantom{,}}}$. Moreover, the
covering families of $C$ are covering families for $T_C$
(\hyperref[II.prop5.1]{5.1}). Thus, if $T'$ denotes the coarsest
topology on $C^{\widehat{\phantom{,}}}$ having properties a) and b),
$T_C$ is finer than $T'$. Let $(f_i:H_i\rightarrow K)$, $i\in I$, be
a\pageoriginale covering family for $T_C$. We show that $(f_i,i\in
I)$ is a covering family for $T'$. Let
$s_i:H_i\rightarrow H=\amalg_i H_i$ be the canonical monomorphisms,
and let $f=(f_i,i\in I):H=\amalg_i H_i\rightarrow K$ be the morphism
defined by the $f_i$. The family of the $s_i$ is covering for $T'$.
Thus, by axiom (T~2) for topologies, to show that $(f_i,i\in I)$ is a
covering family for $T'$, it suffices to show that the morphism
$f:H\rightarrow K$ is covering for $T'$. There exists an epimorphic
family in $C^{\widehat{\phantom{,}}}$,
$u_{\lambda}:X_{\lambda}\rightarrow K$, $\lambda\in \Lambda$, with
$X_{\lambda}\in \ob C$ (\hyperrefext[I][I.prop3.4]{3.4}). Therefore,
again by axiom (T 2) for topologies, to show that $f:H\rightarrow K$
is covering for $T'$, it suffices to show that, for every
$\lambda\in \Lambda$, the morphism
$\pr_2:H\times_K X_{\lambda}\rightarrow X_{\lambda}$ is covering
for $T'$. Let $v_j:Y_j\rightarrow H\times_K X_{\lambda}$, $j\in J$,
with $Y_j\in \ob C$, be an epimorphic family in
$C^{\widehat{\phantom{,}}}$. The family $(\pr_2\circ v_j,j\in J)$ is
covering for $T_C$. It is therefore a covering family in $C$
(\hyperref[II.prop5.1]{5.1}), and hence a covering family for $T'$.
Consequently, the sieve generated by
$\pr_2:H\times_K X_{\lambda}\rightarrow X_{\lambda}$ contains a
covering sieve for $T'$. It is therefore covering for $T'$, and hence
$\pr_2:H\times_K X_{\lambda}\rightarrow X_{\lambda}$ is covering,
\cqfd
\end{enonce*}

\begin{enonce*}[remark]{Remark 5.4.1}\etiquette[5.4.1]{II.rem5.4.1}
The proof of \hyperref[II.prop5.4]{5.4} in fact shows that every
topology $T'$ on $C^{\widehat{\phantom{,}}}$ which is finer than the
canonical topology of $C^{\widehat{\phantom{,}}}$ is the coarsest of
the topologies $T$ on $C^{\widehat{\phantom{,}}}$ having the
following properties:
\begin{enumerate}
\renewcommand{\theenumi}{\alph{enumi}}
\renewcommand{\labelenumi}{\theenumi)}
\item $T$ is finer than the canonical topology of
  $C^{\widehat{\phantom{,}}}$.

\item Every family covering for $T'$ of the form
  $u_i:X_i\rightarrow X$, $i\in I$, where $X$ and the $X_i$ are
  objects of $C$, is covering for $T$.
\end{enumerate}
\end{enonce*}

In particular, every topology $T'$ on
$C^{\widehat{\phantom{,}}}$ which is finer than the canonical
topology is uniquely determined by the families
$u_i:X_i\rightarrow X$, $i\in I$, with $X_i$ and $X$ objects of $C$,
which are covering for $T'$.

\begin{enonce*}[remark]{Remark 5.4.2}\etiquette[5.4.2]{II.rem5.4.2}
It\pageoriginale is easy to show that, for every topology $T'$ on
$C^{\widehat{\phantom{,}}}$ which is finer than the canonical
topology, the set of families of morphisms $(X_i\rightarrow X)$,
$i\in I$, with the same target in $C$, which are covering for $T'$ is
the set of covering families of a topology on $C$. Thus
\hyperref[II.prop5.4]{5.4} and \hyperref[II.rem5.4.1]{5.4.1}
establish a one-to-one correspondence between the topologies on $C$
and the topologies on $C^{\widehat{\phantom{,}}}$ finer than the
canonical topology.
\end{enonce*}

\setcounter{subsection}{5} \setcounter{subsubsection}{-1}
\subsubsection{}\etiquette[5.5.0]{II.subsubsec5.5.0}
Let $C$ be a small category. Denote by $\Caf$ the set of strictly full
subcategories (every object isomorphic to an object of the
subcategory is an object of the subcategory) of
$C^{\widehat{\phantom{,}}}$ whose inclusion functor admits a left
adjoint which commutes with finite projective limits. Also denote by
$\Tt$ the set of topologies on $C$. Theorem
\hyperref[II.theo3.4]{3.4} defines a map
$$
\Phi:\Tt\longrightarrow \Caf{quoi}.
$$

We shall define a map in the opposite direction. For this purpose,
to every element
$\xymatrix{ e=i': C' \ar@<-4pt>[r]&
\ar@<-2pt>[l]_(.35){a'} C^{\widehat{\phantom{,}}} }$ of $\Caf$, with
$\underline{a}'$ left adjoint to $i'$, we must associate a topology
$T_e$ on $C$. For every object $X$ of $C$, define $J_e(X)$ to be the
set of subobjects of $X$ whose inclusion morphism is carried by
$\underline{a}'$ to an isomorphism. Using the hypotheses made on
$\underline{a}'$, one verifies immediately that this defines a
topology $T_e$ on $C$. We have therefore defined a map
$$
\Psi:\Caf\longrightarrow \Tt{quoi}.
$$

We then have the following result (due to J. GIRAUD):

\begin{enonce*}{Theorem 5.5}\etiquette[5.5]{II.theo5.5}
The\pageoriginale map $\Phi$ is a bijection, and $\Psi$ is its
inverse.
\end{enonce*}

\begin{enonce*}[remark]{Proof}
The map $\Psi\circ \Phi$ is the identity. Indeed, this follows
immediately from \hyperref[II.prop5.1]{5.1}. The map
$\Phi\circ \Psi$ is the identity. Indeed, let
$\xymatrix{ i': C' \ar@<-4pt>[r]&
\ar@<-2pt>[l]_(.35){{\underline{a}'}}
C^{\widehat{\phantom{,}}} }$ be an element of $\Caf$, let $T_e$ be
the topology corresponding to it under $\Psi$, and let $C_e$ be the
category of sheaves for $T_e$. Using the definition of the topology
$T_e$ and the definition of bicovering morphisms
(\hyperref[II.def5.2]{5.2}), one proves that the following assertions
are equivalent:
\begin{enumerate}
\renewcommand{\theenumi}{\roman{enumi}}
\renewcommand{\labelenumi}{\theenumi)}
\item The morphism $u$ of $C^{\widehat{\phantom{,}}}$ is bicovering
  for the topology $T_e$.

\item The morphism $u$ of $C^{\widehat{\phantom{,}}}$ is carried by
  $\underline{a}'$ to an isomorphism.
\end{enumerate}

It is clear that $C'$ is a full subcategory of $C_e$. It therefore
suffices to show that every sheaf $F$ for the topology $T_e$ is an
object of $C'$. But, by the equivalence above, the morphism
$F\rightarrow i'\circ a'(F)$ is bicovering, and since its source and
target are sheaves, it follows from \hyperref[II.prop5.3]{5.3} (ii)
that $F\rightarrow i'\circ a'(F)$ is an isomorphism, \cqfd
\end{enonce*}


% END INCLUDED SGA4_English_current_buildable_source_20260906/source/02/drafts/SGA4_II_LINES_1646_1819_English.texfrag


% BEGIN INCLUDED SGA4_English_current_buildable_source_20260906/source/02/drafts/SGA4_II_LINES_1820_2120_English.texfrag
\section{Sheaves with values in a category}\etiquette[6]{II.sec6}

\setcounter{subsection}{-1}
\subsection{}\etiquette[6.0]{II.subsec6.0}
Let $C$ and $D$ be two categories. A contravariant functor from
$C$ to $D$, $F:C^{\circ}\rightarrow D$, is called a
\emph{presheaf on $C$ with values in $D$.}

\begin{enonce*}{Definition 6.1}\etiquette[6.1]{II.def6.1}
Let $C$ be a site and $D$ a category. A presheaf
$F:C^{\circ}\rightarrow D$ is called a \emph{sheaf on $C$ with
values in} $D$ (or, more briefly, a sheaf with values in $D$) if, for
every object $S$ of $D$, the presheaf of sets
$$
X\longmapsto \Hom_D(S,F(X))\quav X\in \ob(C)
$$
is a sheaf. The full subcategory of $\Hhom(C^{\circ},D)$ formed by
the sheaves on $C$ with values in $D$ will be denoted by
$\Hhom^{\sim}(C^{\circ},D)$.
\end{enonce*}

\begin{enonce*}[remark]{Remarks 6.2}\etiquette[6.2]{II.rem6.2}
\begin{enumerate}
\renewcommand{\labelenumi}{\theenumi)}
\hypertarget{sga4.exp.II.remark.6.2.item.1}{}\item The\pageoriginale condition of \hyperref[II.def6.1]{6.1} means
  that for every object $X$ of $C$, every covering sieve $R$ of $X$,
  and every object $S$ of $D$, one has an isomorphism
  (\hyperrefext[I][I.cor3.5]{3.5} and \hyperref[II.def2.1]{2.1}):
$$
\Hom_D(S,F(X))\xrightarrow{\sim}\varprojlim_{C/R}F(.){\quoi}.
$$
Thus, when $D$ is the category of $𝒰$-sets, one recovers the
definition of sheaves of sets (\loccit).

\item Let $D'$ be a $𝒰$-category and let $G:D\rightarrow D'$ be a
  functor commuting with projective limits. By composition, the
  functor $G$ carries sheaves with values in $D$ to sheaves with
  values in $D'$.
\end{enumerate}
\end{enonce*}

\setcounter{subsection}{3} \setcounter{subsubsection}{-1}
\subsubsection{}\etiquette[6.3.0]{II.subsubsec6.3.0}
Let $\gamma$ be a type of algebraic structure defined by finite
projective limits (\hyperrefext[I][I.cor2.9]{2.9}), and let
$𝒰-\gamma$ be the category of $\gamma$-sets belonging to $𝒰$.
Denote by $\esj:𝒰-\gamma\rightarrow 𝒰-\Ens$ the ``underlying
set'' functor. The functor $\esj$ commutes with projective limits and
therefore, by the preceding remark, defines by composition the
functors
\begin{align*}
&\esj^{\widehat{\phantom{,}}} :
\Hhom(C^{\circ},𝒰-\gamma)\longrightarrow
C^{\widehat{\phantom{,}}}\\
&\esj^{\widetilde{\phantom{,}}} :
\Hhom^{\sim}(C^{\circ},𝒰-\gamma) \longrightarrow
C^{\widetilde{\phantom{,}}}{\quoi}.
\end{align*}

The functor $\esj^{\widetilde{\phantom{,}}}$ is called the
``underlying sheaf of sets'' functor. In fact, it factors canonically
through the category $\gamma-\widetilde{C}$ of $\gamma$-objects of
$C^{\widetilde{\phantom{,}}}$. Denoting again by
$\widetilde{\esj}$ the resulting functor
$$
\Hhom^{\sim}(C^{\circ}, 𝒰-\gamma)\longrightarrow
\gamma-\widetilde{C},
$$
we may state:

\begin{enonce*}{Proposition 6.3.1}\etiquette[6.3.1]{II.prop6.3.1}
The functor $\esj^{\widetilde{\phantom{,}}}$ establishes an
equivalence among the category of sheaves on $C$ with values in
$𝒰-\gamma$, the category of $\gamma$-objects in the category of
sheaves of sets, and the category of $\gamma$-objects in the category
of presheaves of sets whose underlying presheaf is a sheaf.
\end{enonce*}

\begin{enonce*}[remark]{Proof}
The\pageoriginale proof essentially uses
(\hyperrefext[I][I.cor3.2]{3.2}) and the fact that a finite projective
limit of sheaves in the category of presheaves is a sheaf
(\hyperref[II.theo4.1]{4.1}).
\end{enonce*}

\setcounter{subsection}{3} \setcounter{subsubsection}{1}
\subsubsection{}\etiquette[6.3.2]{II.subsubsec6.3.2}
The preceding proposition justifies the abuse of language whereby a
sheaf with values in $𝒰-\gamma$ is identified with the corresponding
$\gamma$-object of $C^{\widetilde{\phantom{,}}}$. We shall
henceforth make this abuse of language systematically.

We shall use the classical terminology: \emph{sheaves of groups,
sheaves of commutative groups} (which we shall most often call
\emph{abelian sheaves}), \emph{sheaves of rings, sheaves of
$A$-modules} \ldotsn

\subsubsection{}\etiquette[6.3.3]{II.subsubsec6.3.3}
We denote by
$$
C^{\widetilde{\phantom{,}}}_{\gamma}(\resp
C^{\widehat{\phantom{,}}}_{\gamma})
$$
the category of $\gamma$-objects of $C^{\widetilde{\phantom{,}}}$
(\resp of $C^{\widehat{\phantom{,}}}$). When $\gamma$ is the type of
structure ``$A$-module'', we also write
$C^{\widetilde{\phantom{,}}}_{A}$, or simply
$\widetilde{C}_{\ab}$ when $A=\ZZ$.

Suppose that $C$ is a $𝒰$-site. The ``associated sheaf'' functor is
left exact (\hyperref[II.theo4.1]{4.1}), and therefore:

\begin{enonce*}{Proposition 6.4}\etiquette[6.4]{II.prop6.4}
The inclusion functor
$i_{\gamma}:C^{\widetilde{\phantom{,}}}_{\gamma}\rightarrow
C^{\widehat{\phantom{,}}}_{\gamma}$ admits a left adjoint
$a_{\gamma}$ which is left exact; that is, one has an isomorphism
$$
\Hom_{C^{\widetilde{\phantom{,}}}_{\gamma}}(a_{\gamma},\;)\xrightarrow{\sim}\Hom_{C^{\widehat{\phantom{,}}}_{\gamma}}(\;,i_{\gamma}){\quoi}.
$$

Let $X$ be an object of $C^{\widehat{\phantom{,}}}_{\gamma}$. The
underlying sheaf of sets of $a_{\gamma}(X)$ is canonically isomorphic
to the sheaf of sets associated with the underlying presheaf of sets
of $X$. The morphism of underlying presheaves of sets of the adjunction
morphism $\id\rightarrow i_{\gamma}a_{\gamma}$ is identified with the
adjunction morphism applied to the underlying presheaf of sets.
\end{enonce*}

\begin{enonce*}{Proposition 6.5}\etiquette[6.5]{II.prop6.5}
Suppose\pageoriginale that the functor
$\esj:\gamma-𝒰\rightarrow 𝒰-\Ens$ admits a left adjoint
\footnote[1]{One proves that such an adjoint always exists [C.F.]:
  free group, free abelian group, free $A$-module \ldotsn}:
$$
\Hom_{𝒰-\Ens}(.,\esj(.))\xrightarrow{\sim}\Hom_{\gamma-𝒰}(\Lib(.),.){\quoi}.
$$

The functor
$\esj^{\widetilde{\phantom{,}}}:C^{\widetilde{\phantom{,}}}_{\gamma}\rightarrow
C^{\widetilde{\phantom{,}}}$ (\resp
$\esj^{\widehat{\phantom{,}}}:C^{\widehat{\phantom{,}}}_{\gamma}\rightarrow
C^{\widehat{\phantom{,}}}$) admits a left adjoint
$\Lib^{\widetilde{\phantom{,}}}$ (\resp
$\Lib^{\widehat{\phantom{,}}}$), and one has a canonical isomorphism
$$
\Lib^{\widetilde{\phantom{,}}}=a_{\gamma}\Lib^{\widehat{\phantom{,}}}
i{\quoi}.
$$
\end{enonce*}

\begin{enonce*}[remark]{Proof}
The proof is formal and is left to the reader.
\end{enonce*}

\begin{enonce*}{Corollary 6.6}\etiquette[6.6]{II.cor6.6}
Let $(X_i)$, $i\in I$, be a generating family of
$C^{\widetilde{\phantom{,}}}$ (\hyperref[II.subsec4.9]{4.9}). Then
the family $\Lib^{\widetilde{\phantom{,}}}(X_i)$ is a generating
family of $C^{\widetilde{\phantom{,}}}_{\gamma}$.
\end{enonce*}

\begin{enonce*}[remark]{Proof}
The proof is formal once one observes that the functor
$\esj^{\widetilde{\phantom{,}}}$ commutes with projective limits and
is conservative ($\esj^{\widetilde{\phantom{,}}}(u)$ is an
isomorphism $\Leftrightarrow$ $u$ is an isomorphism).
\end{enonce*}

\begin{enonce*}{Proposition 6.7}\etiquette[6.7]{II.prop6.7}
Let $A$ be a $𝒰$-sheaf of rings, or a small ring, and let
$C^{\widetilde{\phantom{,}}}_A$ (\resp
$C^{\widehat{\phantom{,}}}_A$) be the category of sheaves (\resp of
presheaves) of unital $A$-modules
$(\hyperref[II.subsubsec6.3.3]{6.3.3})$ on a $𝒰$-site $C$. Then
$C^{\widetilde{\phantom{,}}}_A$ is an abelian $𝒰$-category satisfying
the axioms {\rm (AB 5)} (``existence of left exact filtered inductive
limits'') and {\rm (AB 3)$^{\ast}$} (``existence of infinite
products'') of {\rm[\textsc{Tohoku}]}. It has a family of generators
indexed by an element of $𝒰$.
\end{enonce*}

\begin{enonce*}[remark]{Proof}
It is clear that the category $C^{\widehat{\phantom{,}}}_A$ is an
abelian category satisfying the axioms (AB 3), (AB 4), and (AB 5)
(\hyperrefext[I][I.prop2.8]{2.8} and
\hyperrefext[I][I.cor3.3]{3.3}). It then follows\pageoriginale from
\hyperref[II.prop6.4]{6.4} that the category
$C^{\widetilde{\phantom{,}}}_A$ is an additive category in which
kernels and cokernels are representable. More precisely, let $i_A$
and $a_A$ be the inclusion and associated-sheaf functors, respectively,
and let $u:X\rightarrow Y$ be a morphism of
$C^{\widetilde{\phantom{,}}}_A$. The functor $a_A$ is left exact and
commutes with inductive limits. The functor $i_A$ commutes with
projective limits. Consequently the functor
$i_Aa_A:C^{\widehat{\phantom{,}}}_A\rightarrow
C^{\widehat{\phantom{,}}}_A$ is an additive left exact functor. One
deduces canonical isomorphisms
\begin{gather*}
i_A(\Ker(u))\xrightarrow{\sim} \Ker(i_A(u)){\quoi},\tag{$\ast$}\\
\coker(u)\xrightarrow{\sim} a_A \coker
(i_A(u)){\quoi}.\tag{${\ast}{\ast}$}
\end{gather*}
Consequently there is a canonical isomorphism
\begin{equation}\coim(u)\xrightarrow{\sim} a_A \coim (i_A (u)){\quoi}.
\tag{${\ast}{\ast}{\ast}$}
\end{equation}
Let us show that the canonical morphism
$\coim(u)\rightarrow \im(u)$ is an isomorphism. For this, by (*), it
suffices to show that the sequence
$$
0\longrightarrow i_A(\coim(u))\longrightarrow i_A(Y)\longrightarrow
i_A(\coker(u))
$$
is exact. Using the isomorphisms $({\ast}{\ast})$ and
$({\ast}{\ast}{\ast})$, and observing that $a_Ai_A(Y)$ is isomorphic
to $Y$, one sees that this sequence is isomorphic to the sequence
$$
0\longrightarrow i_Aa_A(\coim(i_A(u)))\longrightarrow i_Aa_Ai_A
(Y)\longrightarrow i_Aa_A(\coker(i_A(u))){\quoi},
$$
which is precisely the image under the functor $i_Aa_A$ of the
sequence
$$
0\longrightarrow \coim(i_A(u))\longrightarrow i_A(Y)\longrightarrow
\coker(i_A(u)){\quoi}.
$$
This sequence is exact, and the functor $i_Aa_A$ is left exact. Thus
the category $C^{\widetilde{\phantom{,}}}_A$ is an abelian category.
The functor $a_A:C^{\widehat{\phantom{,}}}_A\rightarrow
C^{\widetilde{\phantom{,}}}_A$\pageoriginale is left exact and
commutes with arbitrary inductive limits. Consequently it is
additive, exact, and commutes with inductive limits. Since the category
$C^{\widehat{\phantom{,}}}_A$ satisfies axiom (AB 5), the same is true
of the category $C^{\widetilde{\phantom{,}}}_A$. Finally, it is clear
that in the category $C^{\widetilde{\phantom{,}}}_A$, products indexed
by an element of $𝒰$ are representable, and hence axiom
(AB 3)$^{\ast}$ is satisfied. This completes the proof of the first
assertion.
\end{enonce*}

For every ring $A$ which is an element of $𝒰$, the functor
$$
\esj:𝒰-A-\mbox{module}\longrightarrow 𝒰-\Ens
$$
admits a left adjoint $X\longmapsto X_A$ (the free $A$-module
generated by $X$). It
therefore suffices to apply \hyperref[II.prop4.10]{4.10} and
\hyperref[II.cor6.6]{6.6} to complete the proof.

\begin{enonce*}[remark]{Notation 6.8}\etiquette[6.8]{II.not6.8}
Let $H$ be a sheaf of sets. The sheaf of $A$-modules associated with
the presheaf (\cf notation \hyperref[II.prop6.5]{6.5})
$$
X\longmapsto \Lib_A H(X)\qquad X\in \ob C
$$
is denoted by $A_H$. When $H=\underline{a}(X)$ (the sheaf associated
with the presheaf represented by $X$), one sometimes simply writes
$A_X$ (by abuse of notation). The proof of \hyperref[II.prop6.7]{6.7}
shows that the family of sheaves of $A$-modules $A_X$,
$(X\in \ob(C))$, is a generating family, indexed by an element of
$𝒰$, of the category of sheaves of $A$-modules.
\end{enonce*}

\begin{enonce*}[remark]{Remark 6.9}\etiquette[6.9]{II.rem6.9}
According to [\textsc{Tôhoku}], \hyperref[II.prop6.7]{6.7} shows that
the category $C^{\widetilde{\phantom{,}}}_A$, when $A$ is an element
of $𝒰$, has enough injectives. It is also known that infinite products
are not necessarily exact in $C^{\widetilde{\phantom{,}}}_A$, and
therefore that, in general, the category
$C^{\widetilde{\phantom{,}}}_A$ does not have enough projectives
[\textsc{Roos}].
\end{enonce*}

\begin{thebibliography}{99}
\bibitem{sga4.exp.II.bib.1} M. \textsc{Artin};\pageoriginale Grothendieck's topologies.

\bibitem{sga4.exp.II.bib.2} M. \textsc{Demazure}: Séminaire de Géométrie Algébrique III,
  Exposé IV

Lecture Notes. Springer-Verlag.

\bibitem{sga4.exp.II.bib.3} A. \textsc{Grothendieck}: Sur quelques points d'Algèbre
Homologique.

 Tohoku Math. Journal.

\bibitem{sga4.exp.II.bib.4} J.E. \textsc{Roos}: CR
\end{thebibliography}

\ifx\danslelivre\undefined
%\bibliography{../fichiers-configuration-latex/bibliographie}
%\bibliographystyle{../fichiers-configuration-latex/smfalpha}
\end{document}
\fi

% END INCLUDED SGA4_English_current_buildable_source_20260906/source/02/drafts/SGA4_II_LINES_1820_2120_English.texfrag


% END INCLUDED SGA4_English_current_buildable_source_20260906/source/02/02_English.tex


\makeatletter
\def\input@path{{../03/}}
\makeatother
% BEGIN INCLUDED SGA4_English_current_buildable_source_20260906/source/03/03_English.tex
\ifx\danslelivre\undefined
% Release engine: XeLaTeX (the LuaTeX compatibility layer breaks Xy-pic output).
\documentclass[10pt]{../fichiers-configuration-latex/smfart}
% BEGIN INCLUDED SGA4_English_current_buildable_source_20260906/source/fichiers-configuration-latex/paquets.tex
\usepackage{polyglossia}
%\usepackage[francais]{babel}
\usepackage{fontspec}
\usepackage[math-style=ISO,bold-style=upright,sans-style=upright]{unicode-math}
%\usepackage[mathletters]{ucs}
%\usepackage[utf8]{inputenc}
%\usepackage[T1]{fontenc}
\usepackage[unicode,hypertexnames=false,colorlinks=true,linkcolor=blue,citecolor=Sepia,urlcolor=magenta]{hyperref}
\setmainlanguage{french}
\setmainfont[Ligatures=TeX]{Linux Libertine O}
% variante possible : TeXGyrePagella ?
\setmathfont{XITS Math}

% maths
\usepackage{amsmath}
%\usepackage{amssymb}
%\usepackage{mathrsfs}
%\usepackage{stmaryrd}
%\usepackage{euscript}

% couleurs
\usepackage[usenames,dvipsnames]{xcolor}

% graphiques
\usepackage{graphicx}
\usepackage[all]{xy}
%\usepackage{tikz}
%\usetikzlibrary{matrix,arrows,calc}
%\tikzset{column sep=3em,row sep=2em,text height=1.5ex, text depth=0.25ex,>=angle 90}
%\tikzset{description/.style={fill=white,inner sep=1pt}}

% hyperliens, 'etiquettes
\usepackage{xr-hyper}
\usepackage{../fichiers-configuration-latex/smfthm}
%\ifx\cacheretiquettes\undefined
%\usepackage[inline]{showlabels} \else \fi % pour rendre visibles les 'etiquettes
%\usepackage{../fichiers-configuration-latex/makeidx} % pour faire un index
%\usepackage[pagebackref,unicode]{hyperref}  % pour voir o`u sont cit'ees les r'ef'erences
% hyperref is loaded once above with the cumulative-reader destination policy.
%\hypersetup{breaklinks=true} %permet le retour `a la ligne dans les liens trop longs. Incompatible avec pdflatex
%\renewcommand*{\backref}[1]{{\small [cit\'e page(s)~#1.]}}

\usepackage{xspace}
\usepackage{pifont} % pour expose XVII

% END INCLUDED SGA4_English_current_buildable_source_20260906/source/fichiers-configuration-latex/paquets.tex

% BEGIN INCLUDED SGA4_English_current_buildable_source_20260906/source/fichiers-configuration-latex/numerotation.tex
\renewcommand{\thefootnote}{\roman{footnote}} % pour mettre les notes en bas de page sous la forme (ii) et non (2), qui ressemble trop (?) a un exposant.

\newcounter{pageoriginal}
\newcommand{\pageoriginale}{\refstepcounter{pageoriginal}\marginpar{\footnotesize\textbf{\thepageoriginal}}\xspace}
\let\pageoriginaled\pageoriginale

\usepackage{ifthen}
\newboolean{original}
\def\sisi{\ifthenelse{\boolean{original}}}
\setboolean{original}{false}

\newcounter{cnote}
\def\nde#1{\ifthenelse{\boolean{original}}{}{\stepcounter{cnote}\def\thefootnote{\arabic{cnote}}\protect\footnote[\value{cnote}]{N.D.E.: #1}}}
\def\ndemark{\ifthenelse{\boolean{original}}{}{\stepcounter{cnote}\def\thefootnote{\arabic{cnote}}\protect\footnotemark[\value{cnote}]}}
\def\ndetext#1{\ifthenelse{\boolean{original}}{}{\def\thefootnote{\arabic{cnote}}\protect\footnotetext[\value{cnote}]{N.D.E.: #1}}}

% END INCLUDED SGA4_English_current_buildable_source_20260906/source/fichiers-configuration-latex/numerotation.tex

% BEGIN INCLUDED SGA4_English_current_buildable_source_20260906/source/fichiers-configuration-latex/macros.tex
\setmainlanguage{french}
\setotherlanguage{english}


%% Fixed, engine-independent release identifier.  The inherited dynamic Git
%% probe used Lua io.popen, exposed build-state-dependent output, and printed
%% a visible question mark outside a Git checkout.
\def\showgitstatus{English release r2, 2026-07-23}


% macros pour LaTeX
\renewcommand{\mathsf}{}

% Police différente pour U+2113 SCRIPT SMALL L
\setmathfont[range="2113]{latinmodern-math.otf}

% Give every inherited semantic label its own stable PDF destination instead
% of inheriting whichever automatic equation/page anchor happened to be live.
% The legacy label remains as a compatibility alias used by existing edges.
\makeatletter
\def\etiquette[#1]#2{%
  \Hy@raisedlink{\hyper@anchorstart{sga4.semantic.#2}\hyper@anchorend}%
  \def\@currentHref{sga4.semantic.#2}%
  \def\@currentlabel{#1}%
  \label{#2}%
}
\makeatother
\def\eqhyperref[#1]#2{(\hyperref[#1]{#2})}
\def\pagehyperref[#1]#2{(page \hyperref[#1]{#2})}
\def\hyperrefext[#1][#2]#3{\hyperref[#2]{#1 #3}}
\def\eqhyperrefext[#1][#2]#3{\hyperref[#2]{#1 #3}}     %cf. exposes 17
\def\neweqhyperrefext[#1][#2]#3{\hyperref[#2]{#1 #3}}  %cf. exposes 17

\def\shat#1{\protect\widehat{#1}}
\def\stilde#1{\protect\widetilde{#1}}
%%%%%%%%%%%%%%%%%%


\def\etcn{etc.\xspace}
\def\etco{etc...\xspace}
\def\etc{\sisi{\etco}{\etcn}}
\def\ldotsn{\sisi{...}{\etcn}}

\def\cf{{cf.}\kern.3em}
\def\Cf{{Cf.}\kern.3em}
\def\ie{{i.e.}\kern.3em}
\def\iev{{i.e.},\xspace}
\def\loccito{loc.\kern2pt cit.\xspace}
\def\loccitn{\textit{loc.\kern2pt cit.}\xspace}
\def\loccit{\sisi{\loccito}{\loccitn}}
\def\EGA{\textup{EGA}\kern.2em}
\def\SGA{\textup{SGA}\kern.3em}

\def\quoi{}

\def\quav{\quad,\quad}
\def\quaet{\quad \text{et}\quad}
\def\fibreproduct#1#2#3{#1 ×_{#3} #2}
\let\fprod\fibreproduct
\def\ccheck#1#2{\mbox{${\check{\check{\rm #1}}}^{\lower2pt\hbox{$\scriptstyle #2$}}$}}
\def\sheaf#1{\protect\mathcal{#1}}
\def\vieuxsheaf#1{{#1}} % cf. problème avec sheaf{a\Phi_i} par exemple...
\DeclareMathOperator{\Hom}{Hom}
\DeclareMathOperator{\can}{can}
\DeclareMathOperator{\Ind}{Ind}
\DeclareMathOperator{\ind}{ind}
%\DeclareMathOperator{\Hhomtop}{\mathcal{H}\!\mathit{omtop}}
\DeclareMathOperator{\Hhomtop}{\mathbf{Homtop}}
\DeclareMathOperator{\Point}{Point}
\DeclareMathOperator{\Pro}{Pro}
\DeclareMathOperator{\Top}{Top}
\DeclareMathOperator{\Mod}{Mod}
\DeclareMathOperator{\Cart}{Cart}
%\DeclareMathOperator{\Mmorsite}{\mathcal{M}\!\mathit{orsite}}
\DeclareMathOperator{\Mmorsite}{\mathbf{Morsite}}
\DeclareMathOperator{\loc}{loc}
\DeclareMathOperator{\Loc}{Loc}
\DeclareMathOperator{\Pic}{Pic}
\DeclareMathOperator{\Et}{Et}
\DeclareMathOperator{\card}{card}
\DeclareMathOperator{\Ens}{Ens}
\DeclareMathOperator{\h}{\underline{H}}
\DeclareMathOperator{\Fonct}{Fonct}
\DeclareMathOperator{\Ker}{Ker}
\DeclareMathOperator{\but}{but}
\DeclareMathOperator{\Ob}{Ob}
\DeclareMathOperator{\ob}{ob}
\DeclareMathOperator{\Lib}{Lib}
\DeclareMathOperator{\esj}{esj}
\DeclareMathOperator{\id}{id}
\DeclareMathOperator{\Id}{Id}
\DeclareMathOperator{\ad}{ad}
\DeclareMathOperator{\Iim}{Im}
\DeclareMathOperator{\pr}{pr}
\DeclareMathOperator{\Cat}{Cat}
\DeclareMathOperator{\car}{car}
\DeclareMathOperator{\SysInd}{SysInd}
\DeclareMathOperator{\Sup}{Sup}
\DeclareMathOperator{\Coker}{Coker}
\DeclareMathOperator{\Filt}{Filt}
\DeclareMathOperator{\acc}{acc}
\DeclareMathOperator{\sur}{sur}
\DeclareMathOperator{\cat}{cat}
\DeclareMathOperator{\Ord}{Ord}
\DeclareMathOperator{\Gr}{Gr}
\DeclareMathOperator{\Cov}{Cov}
\DeclareMathOperator{\Caf}{Caf}
\DeclareMathOperator{\coker}{coker}
\DeclareMathOperator{\coim}{coim}
\DeclareMathOperator{\im}{im}
\DeclareMathOperator{\Ouv}{Ouv}
\DeclareMathOperator{\diag}{diag}
\DeclareMathOperator{\Esp}{Esp}
\DeclareMathOperator{\TOP}{TOP}
\DeclareMathOperator{\Unif}{Unif}
\DeclareMathOperator{\Iind}{\mathbf{Ind}}
\DeclareMathOperator{\Hhomcart}{\mathbf{Homcart}}
\DeclareMathOperator{\Hhomtopess}{\mathbf{Homtopess}}
\DeclareMathOperator{\Ppoints}{\mathbf{Points}}
\DeclareMathOperator{\Ppoint}{\mathbf{Point}}
\DeclareMathOperator{\Ttors}{\mathbf{Tors}}
\DeclareMathOperator{\Ssex}{\mathbf{Sex}}
\DeclareMathOperator{\Ffib}{\mathbf{Fib}}
\DeclareMathOperator{\Vvois}{\mathbf{Vois}}
\DeclareMathOperator{\Ppointess}{\mathbf{Pointess}}
\DeclareMathOperator{\scrHom}{\mathbf{Hom}}
\DeclareMathOperator{\sheafPic}{\mathbf{Pic}}
\DeclareMathOperator{\sheafAut}{\mathbf{Aut}}
\DeclareMathOperator{\Hhomtoplocan}{\mathbf{Homtoplocan}}
\DeclareMathOperator{\Hhomtopan}{\mathbf{Homtopan}}
\DeclareMathOperator{\Pptgeom}{\mathbf{Ptgéom}}
\DeclareMathOperator{\Morsite}{Morsite}
\DeclareMathOperator{\Fib}{Fib}
\DeclareMathOperator{\Fais}{Fais}
\DeclareMathOperator{\kar}{kar}
\DeclareMathOperator{\Groupes}{Groupes}
\DeclareMathOperator{\Restr}{Restr}
\DeclareMathOperator{\Isom}{Isom}
\DeclareMathOperator{\Prol}{Prol}
\DeclareMathOperator{\Res}{Res}
\DeclareMathOperator{\Tors}{Tors}
\DeclareMathOperator{\esp}{esp}
\DeclareMathOperator{\es}{es}
\DeclareMathOperator{\od}{od}
\DeclareMathOperator{\gr}{gr}
\DeclareMathOperator{\sob}{sob}
\DeclareMathOperator{\dfn}{dfn}
\DeclareMathOperator{\Sch}{Sch}
\DeclareMathOperator{\ouv}{ouv}
\DeclareMathOperator{\cons}{cons}
\DeclareMathOperator{\Nerf}{Nerf}
\DeclareMathOperator{\Points}{Points}
\DeclareMathOperator{\Espsob}{Espsob}
\DeclareMathOperator{\Homtop}{Homtop}
\DeclareMathOperator{\Pt}{Pt}
\DeclareMathOperator{\Quot}{Quot}
\DeclareMathOperator{\glob}{glob}
\DeclareMathOperator{\Homfib}{Homfib}
\DeclareMathOperator{\Qouv}{Qouv}
\DeclareMathOperator{\Quotouv}{Quotouv}
\DeclareMathOperator{\ens}{ens}
\DeclareMathOperator{\ab}{ab}
\DeclareMathOperator{\an}{an}
\DeclareMathOperator{\fs}{fs}
\DeclareMathOperator{\red}{red}
\DeclareMathOperator{\Ab}{Ab}
\DeclareMathOperator{\cosq}{cosq}
\DeclareMathOperator{\Ann}{Ann}
\DeclareMathOperator{\Zar}{Zar}


%%%%%%%%%%%%
%%%%% SGA4-2
\DeclareMathOperator{\Sex}{Sex}
\DeclareMathOperator{\vois}{vois}
\DeclareMathOperator{\Vois}{Vois}
\DeclareMathOperator{\Ext}{Ext}
\DeclareMathOperator{\ext}{ext}
\DeclareMathOperator{\Fer}{Fer}
\DeclareMathOperator{\Add}{Add}
\DeclareMathOperator{\Cocart}{Cocart}
\DeclareMathOperator{\cosk}{cosk}
\DeclareMathOperator{\plat}{plat}

\DeclareMathOperator{\faisc}{faisc}
\DeclareMathOperator{\sk}{sk}
\DeclareMathOperator{\ch}{ch}

\def\HH{\mathbf{H}}
\def\qc{\mathrm{qc}}
\def\et{\mathrm{et}}
\def\coh{\mathrm{coh}}

\def\Hvr#1{%
\begin{array}[b]{l}
{\kern2pt\scriptstyle r}\\[-8pt]
{\kern1pt\scriptstyle \vee}\\[-5pt]
H^{#1}
\end{array}}

\def\Hvv#1{%
\begin{array}[b]{l}
\kern2pt\rotatebox{90}{${\scriptscriptstyle\ll}$}\\[-6pt]
H^{#1}
\end{array}}

\def\Hvi#1{%
\begin{array}[b]{l}
{\kern0.2pt\scriptscriptstyle \infty}\\[-8pt]
{\kern1pt\scriptstyle \vee}\\[-5pt]
H^{#1}
\end{array}}

\def\Hv#1{%
\begin{array}[b]{l}
{\kern1pt\scriptstyle \vee}\\[-5pt]
H^{#1}
\end{array}}

\DeclareMathOperator{\cart}{cart}
\DeclareMathOperator{\Ffaisc}{\mathbf{Faisc}}


\def\Ev#1{%
\begin{array}[b]{l}
{\kern1pt\scriptstyle \vee}\\[-5pt]
E'_{#1}
\end{array}}

\DeclareMathOperator{\Cont}{Cont}
\DeclareMathOperator{\Ccont}{\mathscr{C}\!\mathit{ont}}
\DeclareMathOperator{\Site}{Site}
\newcommand\MM{\mathbf{M}}
\def\limtop#1{{\displaystyle{\varprojlim_{#1}\!\mbox{\rm top }}}}
\def\Limtop{\mathop{\varprojlim\!\mbox{\rm top}}}
\DeclareMathOperator{\spec}{spec}
\def\ss{\scriptstyle{\mathscr{S}}}
\DeclareMathOperator{\Gal}{Gal}
\def\smallm{\scriptstyle{\mathscr{M}}}
\DeclareMathOperator{\End}{End}
\DeclareMathOperator{\ass}{ass}
\DeclareMathOperator{\zar}{zar}
\DeclareMathOperator{\Ogg}{Ogg}
\DeclareMathOperator{\res}{res}
\DeclareMathOperator{\tr}{tr}
\DeclareMathOperator{\Supp}{Supp}
\DeclareMathOperator{\codim}{codim}
\DeclareMathOperator{\rk}{rk}
\DeclareMathOperator{\cur}{cur}
\DeclareMathOperator{\cdqc}{cdqc}
\newcommand\DD{\mathbf{D}}
\DeclareMathOperator{\cl}{cl}
\def\JJ{\mathbf{J}}
\def\FF{\mathbf{F}}
\def\BB{\mathbf{B}}
\DeclareMathOperator{\Alget}{Alget}
\DeclareMathOperator{\Ass}{Ass}
\DeclareMathOperator{\supp}{supp}
\DeclareMathOperator{\rad}{rad}
\DeclareMathOperator{\prof}{prof}
\DeclareMathOperator{\cint}{int}
\DeclareMathOperator{\sic}{sic}
\DeclareMathOperator{\Tor}{Tor}
\DeclareMathOperator{\tors}{tors}
\DeclareMathOperator{\Sym}{Sym}
\DeclareMathOperator{\parf}{parf}
\DeclareMathOperator{\tor}{tor}
\DeclareMathOperator{\defn}{defn}
\DeclareMathOperator{\Div}{Div}
\DeclareMathOperator{\Pol}{Pol}
\DeclareMathOperator{\Fex}{Fex}
\DeclareMathOperator{\Lie}{Lie}
\DeclareMathOperator{\tot}{tot}
\DeclareMathOperator{\Aut}{Aut}
\DeclareMathOperator{\HOM}{HOM}
\DeclareMathOperator{\pch}{pch}
\DeclareMathOperator{\Ch}{Ch}
\DeclareMathOperator{\EXT}{EXT}
\DeclareMathOperator{\PIC}{PIC}
\DeclareMathOperator{\Rf}{Rf}
\DeclareMathOperator{\inj}{inj}
\newcommand\VV{\mathbf{V}}

\def\Xv{%
\begin{array}[b]{l}
{\kern1pt\scriptstyle \vee}\\[-5pt]
X
\end{array}}

%%%%%%%
\newcommand\TT{\mathbf{T}}
\newcommand\XX{\mathbf{X}}
\newcommand\EE{\mathbf{E}}
%%%%%%

\def\qedsymbol{\textup{C.Q.F.D.}}
\let\CQFD\qed
\let\cqfd\qed

%%%%%\bigdot
\font\bigsymb=cmsy10 at 4pt
\def\bigdot{{\kern1.2pt\raise 1.5pt\hbox{\bigsymb\char15}}}
%%%%%

\def\Lright{\mathop{\vcenter{\hbox{\hbox{$\longrightarrow$}\lower8pt \hbox{\kern-16pt\hbox{$\longrightarrow$}}}}}}

\def\Aa{{\mathcal A}}
\def\Bb{{\mathcal B}}
\def\Cc{{\mathcal C}}
\def\Dd{{\mathcal D}}
\def\Ee{{\mathcal E}}
\def\Ff{{\mathcal F}}
\def\Gg{{\mathcal G}}
\def\Hh{{\mathcal H}}
\def\Ii{{\mathcal I}}
\def\Jj{{\mathcal J}}
\def\Kk{{\mathcal K}}
\def\Ll{{\mathcal L}}
\def\Mm{{\mathcal M}}
\def\Nn{{\mathcal N}}
\def\Oo{{\mathcal O}}
\def\Pp{{\mathcal P}}
\def\Qq{{\mathcal Q}}
\def\Rr{{\mathcal R}}
\def\Ss{{\mathcal S}}
\def\Tt{{\mathcal T}}
\def\Uu{{\mathcal U}}
\def\Vv{{\mathcal V}}
\def\Ww{{\mathcal W}}
\def\Xx{{\mathcal X}}
\def\Yy{{\mathcal Y}}
\def\Zz{{\mathcal Z}}

\newcommand{\CC}{\mathbf{C}}
\newcommand{\GG}{\mathbf{G}}
\newcommand{\KK}{\mathbf{K}}
\newcommand{\LL}{\mathbf{L}}
\newcommand{\NN}{\mathbf{N}}
\newcommand{\PP}{\mathbf{P}}
\newcommand\QQ{\mathbf{Q}}
\newcommand\RR{\mathbf{R}}
\newcommand\ZZ{\mathbf{Z}}

\def\to{\mathchoice{\longrightarrow}{\rightarrow}{\rightarrow}{\rightarrow}}
\def\from{\mathchoice{\longleftarrow}{\leftarrow}{\leftarrow}{\leftarrow}}
\def\mto{\mathchoice{\longmapsto}{\mapsto}{\mapsto}{\mapsto}}
\def\mfrom{\mathrel{\from\joinrel\mapstochar\,}}
\def\hto{\mathrel{\lhook\joinrel\to}}
\def\hfrom{\mathrel{\from\joinrel\rhook}}

\def\implies{\mathchoice{\Longrightarrow}{\Rightarrow}{\Rightarrow}{\Rightarrow}}
\def\implied{\mathchoice{\Longleftarrow}{\Leftarrow}{\Leftarrow}{\Leftarrow}}
\def\iff{\mathchoice{\Longleftrightarrow}{\Leftrightarrow}{\Leftrightarrow}{\Leftrightarrow}}

\def\To#1{\mathchoice{\xrightarrow{\labelstyle~~#1~}}{\stackrel{#1}{\longrightarrow}}{}{}}
\def\Mto#1{\mathrel{\mapstochar\To{#1}}}
\def\From#1{\mathchoice{\xleftarrow{\labelstyle~~#1~}}{\stackrel{#1}{\longleftarrow}}{}{}}

\newcommand{\isomto}{\overset{\sim}{\longrightarrow}}
\newcommand{\isomfrom}{\overset{\sim}{\longleftarrow}}

\newcommand{\IFF}{\unskip~$\iff$~\ignorespaces}
\newcommand{\THEN}{\unskip~$\implies$~\ignorespaces}

\let\leq\leqslant
\let\geq\geqslant
\let\le\leqslant
\let\ge\geqslant

\def\cf{{cf.}\kern.3em}
\def\Cf{{Cf.}\kern.3em}
\def\ie{{i.e.}\kern.3em}
\def\iev{{i.e.},\xspace}
%\def\loccito{loc.\kern2pt cit.\xspace}
%\def\loccitn{\textit{loc.\kern2pt cit.}\xspace}
%\def\loccit{\sisi{\loccito}{\loccitn}}
\def\EGA{\textup{EGA}\kern.2em}
\def\SGA{\textup{SGA}\kern.3em}

\let\oldS\S
\def\S{\oldS\kern .15em}
\def\ptbl{.\kern .2em}
\def\resp{resp.\kern.3em}
\def\Exp{Exp.\kern.3em}

\let\bar\overline
\let\hat\widehat
\let\tilde\widetilde
\let\dpl\displaystyle
\let\ldots\dots
\let\emptyset\varnothing
\let\skpt\skippointrait
\let\epsilon\varepsilon

\def\sfrac#1#2{{#1}/{#2}}

\DeclareMathOperator{\cd}{cd}
%\DeclareMathOperator{\Hhom}{\mathcal{H}\!\mathit{om}}
\DeclareMathOperator{\Hhom}{\mathbf{Hom}}
%\DeclareMathOperator{\Eext}{\mathcal{E}\!\mathit{xt}}
\DeclareMathOperator{\Eext}{\mathbf{Ext}}
\DeclareMathOperator{\Spec}{Spec}
\DeclareMathOperator{\Tr}{Tr}

\let\emptyset=\varnothing

%%%%%%%%%%%%%%%
%%% expose XVII
%%%%%%%%%%%%%%%

\makeatletter
\let\old@optexpose\optexpose
\newcommand{\empty@optexpose}[1]{}
\newcommand{\paren@optexpose}[1]{\old@optexpose{#1}(}
\newcommand{\refnoexpose}[1]{\let\optexpose\empty@optexpose\ref{#1}\let\optexpose\old@optexpose}%\refnoexpose{X.cor5.2} engendre "5.2" au lieu de "X 5.2"
\newcommand{\neweqref}[1]{\let\optexpose\paren@optexpose\textup{\ref{#1})}\let\optexpose\old@optexpose}%\neweqref{XII.eq4.1} engendre "XII (4.1)" au lieu de "(XII 4.1)"
\makeatother
\newcommand{\refvb}[2][VI${}_{\text{B}}$ ]{\sisi{#1\refnoexpose{#2}}{\ref{#2}}}

\newcommand{\leftexp}[2]{{\vphantom{#2}}^{#1}{#2}}
\newcommand{\AAA}{\mathbf{A}}
\newcommand{\Sf}{\mathfrak{S}}
\DeclareMathOperator*{\quotprojlim}{\text{``$\varprojlim$''}}
\DeclareMathOperator*{\quotinjlim}{\text{``$\varinjlim$''}}
\DeclareMathOperator{\Fl}{F\sisi{\ell}{l}}
\DeclareMathOperator{\Pl}{P\sisi{\ell}{l}}
%\DeclareMathOperator{\cl}{c\sisi{\ell}{l}}
\DeclareMathOperator{\tdf}{tdf}
\DeclareMathOperator{\rg}{rg}
\newcommand{\iInd}{\sisi{I}{i}nd}
\newcommand{\iind}{ind}
\newcommand{\ppro}{pro}
\newcommand{\tTor}{\sisi{T}{t}or}
\newcommand{\ttor}{tor}
\newcommand{\chap}{Chap.}
\newcommand{\Ffex}{\Ff\mathit{ex}}
\newcommand{\sheafk}{\underline{k}}
\newcommand{\sheafn}{\underline{n}}
\newcommand{\sheafGamma}{\underline{\Gamma}}

% Repeated references to an already typeset footnote must not use
% \footnotemark without a matching \footnotetext: Hyperref would then create
% a live annotation to a destination that cannot exist.  This helper keeps the
% current printed footnote mark and links it to an explicit stable note target.
\makeatletter
\newcommand{\sgarepeatedfootnotemark}[1]{%
  \begingroup
  \edef\@thefnmark{\thefootnote}%
  \hyperlink{#1}{\@makefnmark}%
  \endgroup}
\makeatother

\sisi{
\DeclareMathOperator{\simarrow}{\rule{.15mm}{1.8mm}\kern -3.4pt\rightsquigarrow}
\newcommand{\sheafell}{\underline{\ell}}
\newcommand{\sheafZZ}{\underline{\ZZ}}
}{
\newcommand{\nderef}[1]{la N.D.E.~(\ref{#1}) page~\pageref{#1}}
}

%%%%%%%%%%%%%%%%%
%%%%% expose XVII
%%%%%%%%%%%%%%%%%

% END INCLUDED SGA4_English_current_buildable_source_20260906/source/fichiers-configuration-latex/macros.tex


\externaldocument{../01/01_English.tex}
\externaldocument{../02/02_English.tex}
%\externaldocument{../03/03_English.tex}
\externaldocument{../04/04.tex}

\setboolean{original}{false}

\begin{document}

\begin{center}
III. Functoriality of categories of sheaves\\
J.-L. Verdier
\end{center}

\selectlanguage{english}version: \showgitstatus
\tableofcontents

\else
\chapter{Functoriality of categories of sheaves}
%\addtocontents{toc}{J.-L. Verdier}
\begin{center}
J.-L. Verdier
\end{center}
\fi
\phantomsection\MakeLinkTarget*{sga4.exp.III}\label{sga4.exp.III}

% BEGIN INCLUDED SGA4_English_current_buildable_source_20260906/source/03/drafts/SGA4_III_LINES_31_241_English.texfrag
In\pageoriginale \hyperrefext[I][I.sec5]{5}, we studied the behavior of categories of
presheaves with respect to functors between the categories of
arguments. In this exposé, we extend this study to the case of sites
and categories of sheaves (${\rm n^{\circ}}$ 1 and 2). After
introducing the induced topology (${\rm n^{\circ}}$ 3), we turn in
${\rm n^{\circ}}$ 4 to the comparison lemma, which will play an
important role in Exp. IV. In number 5, for the patient reader, we
study certain commutative diagrams related to localized categories
(of the form $C/X$). In Theorem \hyperrefext[I][I.prop5.1]{5.1}, we impose
smallness hypotheses on the categories under consideration. In
general, these hypotheses are not satisfied by the sites encountered
in practice ($𝒰$-sites). This necessitates a few contortions
(\hyperref[III.subsec4.2]{4.2}, \hyperref[III.subsec4.3]{4.3},
\hyperref[III.subsec4.4]{4.4}).

\section{Continuous functors}\etiquette[1]{III.sec1}

\begin{enonce*}{Definition 1.1}\etiquette[1.1]{III.def1.1}
Let $C$ and $C'$ be two $𝒰$-sites and let $u:C\rightarrow C'$ be a
functor between their underlying categories. We say that $u$ is
continuous if, for every sheaf of sets $F$ on $C'$, the presheaf
$X\mapsto F\circ u(X)$ on $C$ is a sheaf.
\end{enonce*}

This notion of continuity depends a priori on the universe $𝒰$
with respect to which the two sites are $𝒰$-sites. Proposition
\hyperref[III.prop1.5]{1.5} shows that it is independent of it.

%\setcounter{subsection}{1}
\setcounter{subsection}{10}
\subsection{}\etiquette[1.11]{III.subsec1.11}
%% original misprint: 1.11 should be 1.1.1!? \XXX
Denote\pageoriginale by $i:C^{\sim}\rightarrow
C^{\widehat{\phantom{,}}}$ (\resp by $i':{C'}^{\sim}\rightarrow
{C'}^{\widehat{\phantom{,}}}$) the canonical inclusion functor from
sheaves into presheaves. By Definition \hyperref[III.def1.1]{1.1}, the
functor $u$ is continuous if and only if there exists a functor
$u_s:{C'}^{\sim}\rightarrow C^{\sim}$ such that
$iu_s=u^\ast i'$ (where $u^\ast F=F\circ u$
(\hyperrefext[I][I.sec5.0]{5.0})).

\begin{enonce*}{Proposition 1.2}\etiquette[1.2]{III.prop1.2}
Let $C$ be a small site, $C'$ a $𝒰$-site, and
$u:C\rightarrow C'$ a functor between their underlying
categories. The following properties are equivalent:
\begin{enumerate}
\renewcommand{\theenumi}{\roman{enumi}}
\renewcommand{\labelenumi}{\theenumi)}
\item The functor $u$ is continuous.

\item For every object $X$ of $C$ and every sieve $R\hookrightarrow X$
  covering $X$, the morphism $u_{!}R\hookrightarrow u(X)$ is
  bicovering in ${C'}^{\widehat{\phantom{,}}}$ {\rm(\hyperrefext[II][II.def5.2]{5.2}
  and \hyperrefext[I][I.prop5.1]{5.1}).} \footnote[1]{``Covering'' in place of
  ``bicovering'' did not necessarily suffice; \cf Example
  \hyperref[III.exem1.9.3]{1.9.3} below.}

\item For every bicovering family $H_i\rightarrow K$, $i\in I$, in
  $C^{\widehat{\phantom{,}}}$, the family of morphisms
  $u_{!}H_i\rightarrow u_{!}K$, $i\in I$, is bicovering in
  ${C'}^{\widehat{\phantom{,}}}$.

\item There exists a functor
  $u^s:C^{\sim}\rightarrow
  {C'}^{\sim}$, commuting with inductive limits and
  ``extending $u$'', \ie such that the diagram of canonical functors
\[
\xymatrix@C=3cm@R=1.25cm{
  C \ar[r]^{u} \ar[d]_{\epsilon_C} & C' \ar[d]^{\epsilon_{C'}}\\
C{\widetilde{\phantom{A}}} \ar[r]^{u^s}&
  C'{\widetilde{\phantom{A}}},
}
\]
commutes up to isomorphism.
\end{enumerate}

Moreover, when $C$ is no longer necessarily a small site but is a
$𝒰$-site, one still has {\rm i) $\Leftrightarrow$ iv)}, and the
functor $u^s$ of {\rm iv)} is necessarily a left adjoint of the functor
$u_s:\widetilde{C}'\rightarrow \widetilde{C}$, and is consequently
determined up to a unique isomorphism.
\end{enonce*}

\begin{enonce*}[remark]{Proof}
The\pageoriginale proof of i) $\Leftrightarrow$ iv) when $C$ is a
$𝒰$-site will be given in \hyperref[III.subsec4.2]{4.2}. Assume that
$C$ is small. It is clear that iii) $\Rightarrow$ ii).
\begin{itemize}
\item[i) $\Rightarrow$ iii):] Let $H_i\rightarrow K$, $i\in I$, be a
  bicovering family in $C^{\widehat{\phantom{,}}}$. For every sheaf
  $F$ on $C'$, the presheaf $u^\ast F$
  (\hyperrefext[I][I.sec5.0]{5.0}) is a sheaf on $C$. Hence
  (\hyperrefext[II][II.prop5.3]{5.3}) the canonical map
$$
\Hom(K, u^\ast F)\longrightarrow \prod_i \Hom(H_i, u^\ast F){\quoi},
$$
is bijective. By adjunction (\hyperrefext[I][I.prop5.1]{5.1}), it follows
that the canonical map
$$
\Hom(u_{!}K,F)\longrightarrow \prod_i \Hom(u_{!}H_i, F)
$$
is bijective. Consequently (\hyperrefext[II][II.prop5.3]{5.3}) the family
$u_{!}H_i\rightarrow u_{!}K$, $i\in I$, is bicovering.

\item[ii) $\Rightarrow$ i):] Let $X$ be an object of $C$, let
  $R\hookrightarrow X$ be a covering sieve, and let $F$ be a sheaf on
  $C'$. By (\hyperrefext[II][II.prop5.3]{5.3}), the map
$$
\Hom(u_{!}X,F)\longrightarrow \Hom(u_{!}R,F)
$$
is bijective. By adjunction (\hyperrefext[I][I.prop5.1]{5.1}), it follows
that the map
$$
\Hom(X,u^\ast F)\longrightarrow \Hom(R, u^\ast F)
$$
is bijective. Consequently $u^\ast F$ is a sheaf.

\item[i) $\Rightarrow$ iv):] We use the customary notation
  $\underline{a}$ and $i$ (\resp $\underline{a}'$ and $i'$) for the
  ``associated sheaf'' and ``inclusion into presheaves'' functors.
  For every sheaf $G$ on $C$, set $u^sG=\underline{a}'u_{!}iG$.
  For every sheaf $F$ on $C'$, we have the following sequence of
  natural isomorphisms:
\begin{align*}
\Hom(G,u_sF) &\simeq \Hom(iG, iu_s F)\simeq \Hom(iG, u^\ast i'F)\\
             &\simeq \Hom(u_{!}iG, i'F)\simeq \Hom(\underline{a}'u_{!} iG,
             F){\quoi}.
\end{align*}
(The\pageoriginale first isomorphism follows from the full
faithfulness of $i$; the second follows from the definition of $u_s$;
the third and fourth are obtained by adjunction.) Consequently $u^s$
is left adjoint to $u_s$. In particular, $u^s$ commutes with inductive
limits. For every presheaf $K$ on $C$ and every sheaf $F$ on $C'$, we
have the following sequence of natural isomorphisms:
$\Hom(u^s\underline{a}K,F)\simeq
\Hom(\underline{a}K, u_s F)\simeq \Hom(K,iu_sF)\simeq \Hom(K,u^\ast
i'F)\simeq \Hom(u_{!}K, i'F)\simeq \Hom(\underline{a}'u_{!}K,F)$ (the
third isomorphism follows from the definition of $u_s$; the others are
obtained by adjunction); whence a functorial isomorphism
$u^s\underline{a}K\simeq \underline{a}'u_{!}K$. In particular, using
this isomorphism when $K$ is representable, and taking
\hyperrefext[I][I.prop5.4]{5.4} 3) into account, we obtain a diagram
commutative up to isomorphism:
\[
\xymatrix@C=2.75cm@R=1.25cm{
C \ar[d]_{\epsilon_C} \ar[r]^{u}& C'\ar[d]^{\epsilon_{C'}}\\
\widetilde{C} \ar[r]^{u^{S}}& \widetilde{C'} }
\]

\item[iv) $\Rightarrow$ ii):] Let $h:C\rightarrow
  C^{\widehat{\phantom{,}}}$ (\resp $h':C'\rightarrow
  {C'}^{\widehat{\phantom{,}}}$) be the functor that assigns to an
  object the presheaf it represents. Consider the diagram:
\[
\xymatrix@C=2.5cm@R=1cm{
C \ar[r]^{u}\ar[d]_{h}             &   C' \ar[d]^{h'}\\
C{\widehat{\phantom{i}}} \ar[d]_{\underline{a}} \ar[r]^{u_!}&
C'{\widehat{\phantom{i}}} \ar[d]^{\underline{a}'}\\
C{\widetilde{\phantom{A}}} \ar[r]^{u^{S}}&
C'{\widetilde{\phantom{A}}}. }
\]
The\pageoriginale upper square is commutative
(\hyperref[III.rem1.4]{1.4} 3)), and we have
$\underline{a}h=\epsilon_C$, $\underline{a}'h'=\epsilon_{C'}$. For
every object $K$ of $C^{\widehat{\phantom{,}}}$, we have a canonical
isomorphism (I \hyperref[I.prop3.4]{3.4}):
$$
\varinjlim_{(h(X)\rightarrow K)\in \ob (C/K)} h(X)\xrightarrow{\sim}
K{\quoi}.
$$
Since the functors $u_{!}$, $u^s$, $\underline{a}$, and
$\underline{a}'$ commute with inductive limits, we obtain
isomorphisms:
\begin{gather*}
\varinjlim_{(h(X)\rightarrow K)\in \ob (C/K)}
u^s\underline{a}h(X)\xrightarrow{\sim} u^s \underline{a}K{\quoi},\\
\varinjlim_{(h(X)\rightarrow K)\in \ob (C/K)}
\underline{a}'u_{!}h(X)\xrightarrow{\sim} \underline{a}'u_{!}K{\quoi}.
\end{gather*}
We have $u_{!}h=h'u$ and, by hypothesis, an isomorphism
$\underline{a}'h'u\simeq u^s \underline{a}h$, whence an isomorphism
$$
u^s\underline{a}K\simeq \underline{a}'u_{!}K{\quoi},
$$
which is immediately seen to be functorial in $K$. The diagram
\[
\xymatrix@C=2.5cm@R=1.35cm{
C\widehat{\phantom{a}} \ar[r]^{u_!} \ar[d]_{\underline{a}}&  C'\widehat{\phantom{a}}\ar[d]^{\underline{a}'}\\
C \widetilde{\phantom{A}} \ar[r]^{u^{S}}&  C',& }
\]
is therefore commutative up to isomorphism. Since $\underline{a}i$ is
isomorphic to the identity functor, we have an isomorphism
$u^s\simeq \underline{a}'u_{!}i$, whence the uniqueness of $u^s$. Let
$v:H\rightarrow K$ be a bicovering morphism in
$C^{\widehat{\phantom{,}}}$. Then $\underline{a}(v)$ is an isomorphism
(\hyperrefext[II][II.prop5.3]{5.3}), and consequently
$\underline{a}'u_{!}(v)$ (isomorphic to $u^s\underline{a}(v)$) is an
isomorphism. Thus $u_{!}(v)$ is carried by $\underline{a}'$ to an
isomorphism, and consequently $u_{!}(v)$ is bicovering
(\hyperrefext[II][II.prop5.3]{5.3}), whence ii).
\end{itemize}
\end{enonce*}

The additional assertion,\pageoriginale in the case where $C$ is
small, was proved in the course of the proof, \cqfd

\setcounter{subsection}{2} \setcounter{subsubsection}{0}

% END INCLUDED SGA4_English_current_buildable_source_20260906/source/03/drafts/SGA4_III_LINES_31_241_English.texfrag


% BEGIN INCLUDED SGA4_English_current_buildable_source_20260906/source/03/drafts/SGA4_III_LINES_242_461_English.texfrag
\subsubsection{}\etiquette[1.2.1]{III.subsubsec1.2.1}

The functor $u^s:C^{\sim}\rightarrow {C'}^{\sim}$ of
\hyperref[III.prop2.2]{2.2} iv) may often be interpreted as an
``inverse image'' functor for a ``morphism of topoi''
${C'}^{\sim}\rightarrow C^{\sim}$, \cf \hyperrefext[IV][IV.subsec4.9]{4.9}.
Its properties are summarized in the following proposition:

\begin{enonce*}{Proposition 1.3}\etiquette[1.3]{III.prop1.3}
Let $u:C\rightarrow C'$ be a continuous functor from a small site $C$
to a $U$-site $C'$.
\begin{enumerate}
\renewcommand{\labelenumi}{\theenumi)}
\item The functor $u^s$ is left adjoint to the functor $u_s$.

\item There is a canonical isomorphism $u^s\simeq \underline{a}'u_{!}i$.

\item There is a canonical isomorphism $u^s\underline{a}\simeq
  \underline{a}'u_{!}$.

\item The functor $u^s$ commutes with inductive limits.

\item When the functor $u_{!}$ is left exact (\cf \hyperrefext[I][I.prop5.4]{5.4} 4)),
  the functor $u^s$ is also left exact. More generally, the functor
  $u^s$ commutes with those types of finite projective limits with
  which the functor $u_{!}$ commutes.
\end{enumerate}
\end{enonce*}

\begin{enonce*}[remark]{Proof}
Assertions 1), 2), 3), and 4) were proved in the course of the proof
of \hyperref[III.prop1.2]{1.2}. Assertion 5) follows from the
isomorphism in 2), since $i$ commutes with projective limits and
$\underline{a}'$ commutes with finite projective limits
(\hyperrefext[II][II.theo4.1]{4.1}).
\end{enonce*}

\begin{enonce*}[remark]{Remark 1.4}\etiquette[1.4]{III.rem1.4}
Combining \hyperrefext[I][I.prop5.4]{5.4} 3) and
\hyperref[III.prop1.3]{1.3} 3), we obtain a diagram:
\[
\xymatrix@C=3.5cm@R=1cm{
C \ar[r]^{u} \ar[d]_{h}& C'\ar[d]_{h'}\\
C\widehat{\phantom{a}} \ar[r]^{u_!} \ar[d]_{\underline{a}}&
C'\widehat{\phantom{a}} \ar[d]_{\underline{a}'}\\
C\widetilde{\phantom{a}} \ar[r]& C'\widetilde{\phantom{a}}, }
\]
where\pageoriginale the upper square commutes and the lower square
commutes up to isomorphism. But the functors $\underline{a}$,
$\underline{a}'$, $u_{!}$, and $u^s$ are defined only up to
isomorphism. The reader may verify that the functors $\underline{a}$,
$\underline{a}'$, and $u^s$ can be chosen so that:
\begin{enumerate}
\renewcommand{\labelenumi}{\theenumi)}
\item The composite functors $\underline{a}h$ and
  $\underline{a}'h'$ are injective on the sets of objects;

\item the diagram
\[
\xymatrix@C=2cm@R=1.25cm{
C \ar[r]^{u} \ar[d]_{\underline{a}h}& C' \ar[d]_{\underline{a}'h'}\\
C{\widetilde{\phantom{a}}} \ar[r]^{u^{s}}&
C'{\widetilde{\phantom{a}}} }
\]
commutes.
\end{enumerate}

More precisely, the functors $\underline{a}$ and $\underline{a}'$ can
be chosen so as to satisfy condition 1). Once the functors
$\underline{a}$ and $\underline{a}'$ have been chosen, the functor
$u^s$ can be chosen so as to satisfy condition 2).
\end{enonce*}

\begin{enonce*}{Proposition 1.5}\etiquette[1.5]{III.prop1.5}
Let $𝒰\subset 𝒱$ be two universes, let $C$ and $C'$ be two
$𝒰$-sites, and let $u:C\rightarrow C'$ be a functor between the
underlying categories. Then $u$ is continuous relative to $𝒰$ if
and only if it is continuous relative to $𝒱$. Moreover, when $u$ is
continuous, denote by $u^s_{𝒰}$ (\resp $u^s_{𝒱}$) the functor between
the categories of $𝒰$-sheaves (\resp $𝒱$-sheaves) introduced in
\hyperref[III.prop1.2]{1.2} {\rm iv)}. Then the diagram
\[\vcenter{
\xymatrix@C=3cm@R=1.5cm{ \widetilde{C}_{𝒰}
\ar[r]^{u_{\underline{U}}^s} \ar[d]&
\widetilde{C'}_{\underline{U}} \ar[d]\\
 C\widetilde{\phantom{a}}_{\underline{V}} \ar[r]^{u_{\underline{V}}^s}
 &  C'\widetilde{\phantom{a}}_{\underline{V}} ,
}}\leqno{(1.5.1)}
\]\etiquette[1.5.1]{III.eq1.5.1}
where\pageoriginale the vertical functors are the canonical inclusion
functors, commutes up to canonical isomorphism.
\end{enonce*}

\begin{enonce*}[remark]{Proof}
We give the proof only in the case where $C$ is $𝒰$-small. The
general case will be proved in \hyperref[III.subsec4.3]{4.3}. It is
clear that if the functor $u^\ast$ carries every $𝒱$-sheaf on $C'$
to a $𝒱$-sheaf on $C$, then it carries every $𝒰$-sheaf on $C'$ to a
$𝒰$-sheaf on $C$. Suppose that $u^\ast$ carries every $𝒰$-sheaf on
$C'$ to a $𝒰$-sheaf on $C$. Denote by
$C^{\widehat{\phantom{,}}}_{𝒰}\hookrightarrow
C^{\widehat{\phantom{,}}}_{𝒱}$ (\resp
${C'}^{\widehat{\phantom{,}}}_{𝒰}\hookrightarrow
{C'}_{𝒱}^{\widehat{\phantom{,}}}$) the categories of
$𝒰$-presheaves and $𝒱$-presheaves, and by $u_{𝒰!}$ and $u_{𝒱!}$
the ``inverse image'' functors for $𝒰$-presheaves and
$𝒱$-presheaves, respectively. We have a commutative diagram
(\cf the explicit construction of $u_{!}$ in
\hyperrefext[I][I.prop5.1]{5.1}):
\[\vcenter{
  \xymatrix@C=3cm@R=1.25cm@M=.2cm{
C\widehat{\phantom{a}}_{\underline{U}} \ar[r]^{u_{{\underline{U}}!}}
\ar @{_{(}->}[d]&
C'\widehat{\phantom{a}}_{\underline{V}} \ar @{_{(}->}[d]\\
C\widehat{\phantom{a}}_{\underline{V}}
\ar[r]^{u_{{\underline{V}}!}}&
C'\widehat{\phantom{a}}_{\underline{V}}. }}\leqno{(*)}
\]
The inclusion functors from $𝒰$-presheaves into $𝒱$-presheaves are
fully faithful and commute with projective limits. It follows from
(\hyperrefext[II][II.prop5.1]{5.1} i)) that a bicovering morphism of
$𝒰$-presheaves is a bicovering morphism of $𝒱$-presheaves. Now let
$R\hookrightarrow X$ be a sieve covering an object $X$ of $C$. By
\hyperref[III.prop1.2]{1.2} ii), the morphism
$u_{𝒰!}R\rightarrow u_{𝒰!}X$ is bicovering. Using the commutativity of
$(\ast)$ and what precedes, we deduce that the morphism
$u_{𝒱!}R\rightarrow u_{𝒱!}X$ is bicovering. Consequently
(\hyperref[III.prop1.2]{1.2}), the functor $u^\ast$ carries
$𝒱$-sheaves on $C'$ to $𝒱$-sheaves on $C$. The commutativity of
\eqhyperref[III.eq1.5.1]{1.5.1} then follows from the commutativity of
$(\ast)$, from \hyperrefext[II][II.prop3.6]{3.6}, and from
\hyperref[III.prop1.3]{1.3} 2).
\end{enonce*}

\begin{enonce*}{Proposition 1.6}\etiquette[1.6]{III.prop1.6}
Let\pageoriginale $C$ and $C'$ be two $𝒰$-sites, and let
$u:C\rightarrow C'$ be a functor between the underlying categories.
Consider the following conditions:
\begin{enumerate}
\renewcommand{\theenumi}{\roman{enumi}}
\renewcommand{\labelenumi}{\theenumi)}
\item $u$ is continuous {\rm(\cf. \hyperref[III.def1.1]{1.1} and
  \hyperref[III.prop1.2]{1.2}).}

\item For every covering family $(X_i\rightarrow X)$, $i\in I$, in
  $C$, the family $(uX_i\rightarrow uX)$, $i\in I$, in $C'$ is
  covering.

We have the implication {\rm i) $\Rightarrow$ ii).}

Suppose that the topology of $C$ can be defined by a pretopology
{\rm $T$ (\hyperrefext[II][II.def1.3]{1.3})} and that the functor $u$
commutes with fiber products.\footnote[1]{In fact, it suffices that
$u$ commute with the fiber products occurring in base changes of
morphisms belonging to covering families for $T$.} Then conditions
{\rm i)} and {\rm ii)} are equivalent and are also equivalent to the
following condition:

\item The functor $u$ carries the covering families of $T$ to
  covering families in $C'$.
\end{enumerate}

In particular, suppose that fiber products in $C$ are representable
and that the functor $u$ commutes with fiber products. Then conditions
{\rm i)} and {\rm ii)} are equivalent.
\end{enonce*}

\begin{enonce*}[remark]{Proof}
Let us prove that i) $\Rightarrow$ ii). Let $X_i\rightarrow X$,
$i\in I$, be a covering family in $C$. For every sheaf $F$ on $C'$,
$u^\ast F$ is a sheaf. Consequently, the map
$$
\Hom(X,u^\ast F)\longrightarrow \prod_i\Hom(X_i,u^\ast F)
$$
is injective (\hyperrefext[II][II.prop5.1]{5.1}). Hence there is an
injective map
$$
\Hom(u_{!}X,F)\longrightarrow \prod_i \Hom(u_{!}X_i, F){\quoi}.
$$
Therefore the family $u_{!}X_i\rightarrow u_{!}X$, $i\in I$, is
covering in $C'$ (\hyperrefext[II][II.prop5.1]{5.1}).

Every covering family of the pretopology $T$ is a covering family in
$C$, whence ii) $\Rightarrow$ iii). Let us prove that iii)
$\Rightarrow$ i). Let $X$ be an object\pageoriginale of $C$ and let
$X_i\rightarrow X$, $i\in I$, be a covering family in $\Cov(X)$ (II
\hyperref[II.def1.3]{1.3}). The morphisms $X_i\rightarrow X$ are pullbackable, and the
functor $u$ commutes with fiber products. Consequently, the fiber
products $u(X_i)\times_{u(X)} u(X_j)$ are representable and
canonically isomorphic to $u(X_i\times_X X_j)$. Let $F$ be a sheaf on
$C'$. Then the diagram of sets
$$
F(u(X))\rightarrow \prod_i F(u(X_i))\rightrightarrows
\prod_{i,j}F(u(X_i\times_X X_j))
$$
is exact (\hyperrefext[II][II.def2.1]{2.1},
\hyperrefext[I][I.cor3.5]{3.5}, and \hyperrefext[I][I.prop2.12]{2.12}).
Consequently, the diagram of sets
$$
u^\ast F(X)\rightarrow \prod_i u^\ast F(X_i)\rightrightarrows
\prod_{i,j}u^\ast F(X_i\times_X X_j)
$$
is exact. Therefore $u^\ast F$ is a sheaf
(\hyperrefext[II][II.cor2.4]{2.4}), whence i). The final assertion of
\hyperref[III.prop1.6]{1.6} follows from what precedes and from the
fact that, when fiber products are representable in $C$, all covering
families in $C$\footnote[1]{indexed by sets belonging to a suitable
universe.} define a pretopology on $C$ whose associated topology is
the topology of $C$.
\end{enonce*}


% END INCLUDED SGA4_English_current_buildable_source_20260906/source/03/drafts/SGA4_III_LINES_242_461_English.texfrag


% BEGIN INCLUDED SGA4_English_current_buildable_source_20260906/source/03/drafts/SGA4_III_LINES_462_612_English.texfrag
\begin{enonce*}{Proposition 1.7}\etiquette[1.7]{III.prop1.7}
Let $C$ be a small site, $C'$ a $𝒰$-site, and
$u:C\rightarrow C'$ a continuous functor. Let $\gamma$ be an algebraic
structure defined by finite projective limits, such that in the category
of $\gamma$-objects of $𝒰-\Ens$ the $𝒰$-indexed inductive limits are
representable
\footnote[2]{In fact, one can show that this condition is always
  satisfied.}. We use the notation of Proposition {\rm II \hyperref[II.prop6.4]{6.4}.} The
functor $u_s$ commutes with projective limits and consequently defines a
functor $u^{\gamma}_s:C^{\sim}_{\gamma}\rightarrow
{C'}^{\sim}_{\gamma}$. The functor $u^{\gamma}_s$ has a left adjoint
$u^s_{\gamma}$ with the following properties:
\begin{enumerate}
\renewcommand{\labelenumi}{\theenumi)}
\item There\pageoriginale is a canonical isomorphism
  $u^s_{\gamma}\xrightarrow{\sim}\underline{a}'_{\gamma}u_{\gamma !}
  i_{\gamma}$, and $u^s_{\gamma}$ commutes with every finite projective
  limit with which $u_{\gamma !}$ commutes.

\item There is a canonical isomorphism
  $u^s_{\gamma}\underline{a}_{\gamma}\xrightarrow{\sim}
  \underline{a}'_{\gamma} u'_{\gamma !}$.

\item The functor $u^s_{\gamma}$ commutes with inductive limits.

\item If $u^s$ is left exact, the diagram {\rm(in the notation of II
  \hyperref[II.subsubsec6.3.0]{6.3.0})}
\[
\xymatrix@C=3.5cm@R=1.25cm{ C^{\widetilde{\phantom{a}}_\gamma}
\ar[r]^{u_{Y}^{s}}
\ar[d]_{\esj^{\widetilde{\phantom{a}}}}& C'\widetilde{\phantom{a}}_\gamma \ar[d]_{{\esj'}^{\widetilde{\phantom{a}}}}\\
C^{\widetilde{\phantom{a}}} \ar[r]^{u^{s}}  &
{C'}\widetilde{\phantom{a}} }
\]
is commutative up to isomorphism, and the functor $u^s_{\gamma}$ is
exact.

\item Suppose that the functor $\esj^{\sim}$
  (\resp $\esj^{\sim}$) has a left adjoint
  $\Lib^{\sim}$ (\resp
  ${\Lib'}^{\sim}$) {\rm (\hyperrefext[II][II.prop6.5]{6.5})}. There is a
  canonical isomorphism
  $u^s_{\gamma}\Lib^{\sim}\xrightarrow{\sim}
  {\Lib'}^{\sim} u^s$.
\end{enumerate}

Moreover, when $C$ is not necessarily a small site but is a $𝒰$-site,
the functor $u^{\gamma}_s$ has a left adjoint $u^s_{\gamma}$. Finally,
if $V$ is a universe containing $𝒰$ and if $u^s_{\gamma𝒰}$ (\resp
$u^s_{\gamma𝒱}$) denotes the left adjoint of $u_s$ relative to the
universe $𝒰$ (\resp $𝒱$), the following diagram, analogous to diagram
\eqhyperref[III.eq1.5.1]{1.5.1}:

\hypertarget{sga4.exp.III.equation.1.7.1}{}\[\vcenter{
\xymatrix@C=3.5cm@R=1.25cm{ C^{\widetilde{\phantom{a}}}_{\gamma
\underline{U}} \ar[r]^{u_{\gamma \underline{U}}^{s}}
\ar@{_{(}->}[d]&
\widetilde{C}'_{\gamma \underline{U}} \ar@{_{(}->}[d]\\
C^{\widetilde{\phantom{a}}}_{\gamma \underline{V}}
\ar[r]^{u^s_{r\underline{V}}}  &  \widetilde{C}'_{\gamma
\underline{V}} }}\leqno{(1.7.1)}
\]
is commutative up to canonical isomorphism.
\end{enonce*}

\begin{enonce*}[remark]{Proof}
The proof, in the case where $C$ is a small site, is left to the
reader. It will be completed in \hyperref[III.subsec4.3]{4.3} in the
case where $C$ is a $𝒰$-site.
\end{enonce*}

\begin{enonce*}[remark]{Notation 1.8}\etiquette[1.8]{III.not1.8}
Henceforth\pageoriginale we shall use the notation $u_s$ for the
functor $u^{\gamma}_s$. This notation is unlikely to cause confusion,
because the functor $u_s$ (defined on sheaves of sets) commutes with
finite projective limits.
\end{enonce*}

Likewise, when $u^s$ is left exact, we shall use the notation $u^s$ for
the functor $u^s_{\gamma}$. This abuse of notation is justified by
\hyperref[III.prop1.7]{1.7} 4).

\begin{enonce*}[remark]{Example 1.9.1}\etiquette[1.9.1]{III.exem1.9.1}
Let $X$ be a small topological space. Denote by $\Ouv(X)$ the category
of open subsets of $X$ (the objects of $\Ouv(X)$ are the open subsets
of $X$, and the morphisms of $\Ouv(X)$ are inclusions), equipped with
the canonical topology. A family $(U_i\subset U)$, $i\in I$, is
covering if and only if the open subsets $U_i$ cover $U$
(\hyperrefext[II][II.prop2.6]{2.6}). The sheaves of sets on $\Ouv(X)$
are therefore the sheaves of sets in the sense of [TF]. Let
$f:X\rightarrow Y$ be a continuous map. It induces a functor
$f^{-1}:\Ouv(Y)\rightarrow \Ouv(X)$: for every open subset $U$ of $Y$,
$f^{-1}(U)$ is the inverse image of $U$ under the map $f$. The functor
$f^{-1}$ is continuous (\hyperref[III.prop1.6]{1.6}). The functor
$f^{-1}_s:\Ouv(X)^{\sim}\rightarrow \Ouv(Y)^{\sim}$ is the direct-image
functor for sheaves, in the sense of [TF]. The functor
$(f^{-1})^s:\Ouv(Y)^{\sim}\rightarrow \Ouv(X)^{\sim}$ is the
inverse-image functor for sheaves in the sense of [TF]. The functor
$(f^{-1})^s$ commutes with finite projective limits, by
\hyperref[III.prop1.3]{1.3} 5).
\end{enonce*}

\begin{enonce*}[remark]{Example 1.9.2}\etiquette[1.9.2]{III.exem1.9.2}
Let $X$ be a topological space and let $Y\hookrightarrow X$ be an open
subset of $X$. Every open subset of $Y$ is an open subset of $X$, and
thus we obtain a functor $\Phi:\Ouv(Y)\rightarrow \Ouv(X)$. The functor
$\Phi$ is continuous (\hyperref[III.prop1.6]{1.6}). The functor
$\Phi_s:\Ouv(X)^{\sim}\rightarrow
\Ouv(Y)^{\sim}$ is the\pageoriginale ``restriction to $Y$'' functor.
The functor $\Phi^s_{\ab}:\Ouv(Y)^{\sim}_{\ab}\rightarrow
\Ouv(X)^{\sim}_{\ab}$ (\hyperrefext[II][II.subsubsec6.3.3]{6.3.3}, I \hyperref[III.prop1.7]{1.7})
is the ``extension by zero outside $Y$'' functor introduced in [TF].
The functor $\Phi^s:\Ouv(Y)^{\sim}\rightarrow
\Ouv(X)^{\sim}$ (\hyperref[III.prop1.3]{1.3}) is analogous to the
``extension by zero'' functor: for every sheaf $H$ on $Y$ and every
point $y\in Y$, the stalk of $\Phi^sH$ at $y$ is canonically
isomorphic to the stalk of $H$ at $y$; for every point $x\in X$,
$x\not\in Y$, the stalk of $\Phi^sH$ at $x$ is empty. The functor
$\Phi^s$ is called the ``extension by the empty set outside $Y$''
functor. The functor $\Phi^s$ commutes with fiber products and with
finite products over a nonempty index set, but it does not carry the
terminal object of $\Ouv(Y)^{\sim}$ to the terminal object of
$\Ouv(X)^{\sim}$.
\end{enonce*}

\begin{enonce*}[remark]{Example 1.9.3}\etiquette[1.9.3]{III.exem1.9.3}
Let $C$ and $C'$ be two small sites, and let $u:C\rightarrow C'$ be a
functor between the underlying categories that carries every covering
family of $C$ to a covering family of $C'$. The functor $u$ is not
continuous in general (\cf nevertheless \hyperref[III.prop1.6]{1.6}).
Here is a counterexample. Let $S$ be a set of infinite cardinality in
the universe $𝒰$, and let $C$ be the full subcategory of the category
of sets whose objects are the finite, nonempty subsets of $S$. Equip
$C$ with the canonical topology. Observe that the covering families of
$C$ are the surjective families of maps and that covering families are
nonempty. Observe also that, up to isomorphism, there are only two
constant sheaves on $C$: the sheaf with value the empty set and the
sheaf with value a one-element set. Set $C=C'$ and let
$u:C\rightarrow C'$ be a constant functor. The functor $u$ carries the
covering families of $C$ to covering families of $C'$. The functor $u$
is not continuous. Indeed, since the representable presheaves of $C'$
are sheaves, for every object $X$ of $C'$ there exists a
sheaf\pageoriginale $F$ on $C'$ such that the cardinality of $F(X)$ is
$\geq 2$, and consequently the presheaf $u^\ast F$ is not a sheaf.
\end{enonce*}


% END INCLUDED SGA4_English_current_buildable_source_20260906/source/03/drafts/SGA4_III_LINES_462_612_English.texfrag


% BEGIN INCLUDED SGA4_English_current_buildable_source_20260906/source/03/drafts/SGA4_III_LINES_613_861_English.texfrag
\section{Cocontinuous functors}\etiquette[2]{III.sec2}

\begin{enonce*}{Definition 2.1}\etiquette[2.1]{III.def2.1}
Let $C$ and $C'$ be two $𝒰$-sites, and let $u:C\rightarrow C'$ be a
functor between the underlying categories. We say that $u$ is
cocontinuous if it has the following property:

$(COC)$ For every object $Y$ of $C$ and every covering sieve
$R\hookrightarrow u(Y)$, the sieve on $Y$ generated by the arrows
$Z\rightarrow Y$ such that $u(Z)\rightarrow u(Y)$ factors through
$R$ covers~$Y$.
\end{enonce*}

\begin{enonce*}{Proposition 2.2}\etiquette[2.2]{III.prop2.2}
Let $C$ be a small site, let $C'$ be a $𝒰$-site, and let
$u:C\rightarrow C'$ be a functor between the underlying categories.
Denote by $\widehat{u}_{\ast}:C^{\widehat{\phantom{,}}}\rightarrow
{C'}^{\widehat{\phantom{,}}}$ the functor right adjoint to the
functor $F\mapsto F\circ u=\widehat{u}^{\ast}F$ {\rm
(\hyperrefext[I][I.prop5.1]{5.1}).} The functor $u$ is cocontinuous if and only if,
for every sheaf $G$ on $C$, $\widehat{u}_{\ast}G$ is a sheaf on $C'$.
\end{enonce*}

\begin{enonce*}[remark]{Proof}
\emph{The condition is sufficient.} Suppose that for every sheaf
$G$ on $C$, the presheaf $\widehat{u}_{\ast}G$ is a sheaf. Let $Y$
be an object of $C$, and let $R\hookrightarrow u(Y)$ be a covering
sieve. There is an isomorphism $\Hom(u(Y),
u_{\ast}G)\xrightarrow{\sim}\Hom(R,u_{\ast}G)$, hence, by adjunction
(\hyperrefext[I][I.prop5.1]{5.1}), an isomorphism
$\Hom(\widehat{u}^{\ast}\circ
u(Y),G)\xrightarrow{\sim}\Hom(\widehat{u}^{\ast}R,G)$. It follows
(\hyperrefext[II][II.prop5.3]{5.3}) that the morphism $\widehat{u}^{\ast}R\rightarrow
\widehat{u}^{\ast}u(Y)$ is bicovering. But the functor
$\widehat{u}^{\ast}$ commutes with projective limits, and therefore
the morphism $\widehat{u}^{\ast}R\rightarrow \widehat{u}^{\ast}u(Y)$
is a monomorphism. It is thus a covering monomorphism. On the other
hand, there is a canonical morphism
$Y\xrightarrow{p}\widehat{u}^{\ast}u(Y)$ deduced from the adjunction
morphism $\id_C\rightarrow
\widehat{u}^{\ast}u_{!}$\pageoriginale (\hyperrefext[I][I.prop5.4]{5.4} 3)). We then
make the base change $Y\xrightarrow{p}\widehat{u}^{\ast}u(Y)$. This
gives a covering sieve on $Y$:
$$
\widehat{u}^{\ast}R\times_{\widehat{u}^{\ast}u(Y)}Y\longrightarrow
Y{\quoi},
$$
and, on determining the arrows $Z\rightarrow Y$ that factor through
this sieve, one obtains the sieve described by property (COC).
\end{enonce*}

\emph{The condition is necessary.} Let $S$ be an object of $C'$
and let $R\hookrightarrow S$ be a covering sieve on $S$. We must
prove that, for every sheaf $G$ on $C$, there is an isomorphism
$$
\Hom(S,u_{\ast}G)\xrightarrow{\sim} \Hom(R,u_{\ast}G){\quoi}.
$$

Using Proposition \hyperrefext[II][II.prop5.3]{5.3} and the adjunction
isomorphisms, we see that it suffices to prove that the monomorphism
$\widehat{u}^{\ast}R\rightarrow \widehat{u}^{\ast}S$ is covering.
For this, it suffices to prove (\hyperrefext[II][II.prop5.1]{5.1}) that for every
base change $Y\rightarrow \widehat{u}^{\ast}S$, the sieve
$\widehat{u}^{\ast}R\times_{\widehat{u}^{\ast}S}Y$ is covering,
where $Y$ is an object of $C$. But, by the properties of adjoint
functors and \hyperrefext[I][I.prop5.4]{5.4} 3), the morphism
$Y\rightarrow \widehat{u}^{\ast}S$ factors through
$Y\xrightarrow{p}\widehat{u}^{\ast}u(Y)\rightarrow
\widehat{u}^{\ast}S$, and consequently the sieve
$\widehat{u}^{\ast}R\times_{\widehat{u}^{\ast}S}Y\hookrightarrow Y$
is obtained, by the base change
$Y\xrightarrow{p}\widehat{u}^{\ast}u(Y)$, from the monomorphism
$\xymatrix@1@M=.15cm{
\fprod{\hat{u}^{\ast}R}{\hat{u}^{\ast}u(Y)}{\hat{u}^{\ast}S}
\ar@{^{(}->}@<-1pt>[r]^(.6){q}&\hat{u}^{\ast}u(Y)}$.
Now the functor $\widehat{u}^{\ast}$ commutes with projective limits,
and consequently the monomorphism $q$ is the image under
$\widehat{u}^{\ast}$ of the sieve $R\times_S u(Y)\hookrightarrow u(Y)$,
which is covering. Hypothesis (COC) then allows us to conclude that
the sieve
$\widehat{u}^{\ast}R\times_{\widehat{u}^{\ast}S}Y\hookrightarrow Y$
is covering, \cqfd

\begin{enonce*}{Proposition 2.3}\etiquette[2.3]{III.prop2.3}
Let $C$ and $C'$ be two $𝒰$-sites, and let $u:C\rightarrow C'$ be a
cocontinuous functor. Denote by $u^{\ast}:\widehat{C}'\rightarrow
\widehat{C}$ the functor $u^{\ast}=\underline{a}\widehat{u}^{\ast}i'$,
where $\underline{a}$ denotes the ``associated sheaf'' functor for
$C^{\widehat{\phantom{,}}}$, $\widehat{u}^{\ast}$ the functor
$F\mapsto F\circ u$, and $i'$ the inclusion functor\pageoriginale
$\widetilde{C}'\rightarrow {C'}^{\widehat{\phantom{,}}}$.
\begin{enumerate}
\renewcommand{\labelenumi}{\theenumi)}
\item The functor $u^{\ast}$ commutes with inductive limits and is
  exact.

\item There is a canonical isomorphism $u^{\ast}\underline{a}'\simeq
  \underline{a}\widehat{u}^{\ast}$, where $\underline{a}'$ denotes
  the ``associated sheaf'' functor for
  ${C'}^{\widehat{\phantom{,}}}$.

\item The functor $u^\ast$ has a right adjoint
  $u_{\ast}:\widetilde{C}\rightarrow \widetilde{C}'$. When $C$ is
  small, the diagram
\[\vcenter{
\xymatrix@C=3.25cm@R=1.25cm@M=.25cm{
C\widetilde{\phantom{a}}\ar[r]^{u_\ast} \ar@{_{(}->}[d]_{i}&
\widetilde{C}' \ar@{_{(}->}[d]^{i'}\\
C\widehat{\phantom{a}} \ar[r]^{\widehat{u}_{\ast}}  &
C'\widehat{\phantom{a}}, }}\leqno{(2.3.1)}
\]\etiquette[2.3.1]{III.eq2.3.1}
in which the lower functor is right adjoint to the functor
$\widehat{u}^{\ast}$ (\hyperrefext[I][I.prop5.1]{5.1}), commutes up
to canonical isomorphism.

\item Let $𝒱\supset 𝒰$ be a universe. Denote by
  $u_{\ast_{𝒰}}:
  C^{\sim}_{𝒰}\rightarrow
  \widetilde{C}'_{𝒰}$ (\resp
  $u_{\ast_{𝒱}}:C^{\sim}_{𝒱}\rightarrow
  \widetilde{C}'_{𝒱}$) the right adjoint relative to the universe
  $𝒰$ (\resp $𝒱$) whose existence is asserted in $3$). The diagram
\[\vcenter{
\xymatrix@C=3.25cm@R=1.25cm@M=.2cm{
C\widetilde{\phantom{a}}_{\underline{U}}
\ar[r]^{u_{\ast_{\underline{U}}}} \ar@{_{(}->}[d]&
\widetilde{C}'_{\underline{U}} \ar@{_{(}->}[d]\\
C\widetilde{\phantom{a}}_{\underline{V}}
\ar[r]^{u_{\ast_{\underline{V}}}}& \widetilde{C}'_{\underline{V}}
}}\leqno{(2.3.2)}
\]\etiquette[2.3.2]{III.eq2.3.2}
commutes up to isomorphism.
\end{enumerate}
\end{enonce*}

\begin{enonce*}[remark]{Proof}
Assertions 1) and 2) follow immediately from
\hyperrefext[II][II.theo3.4]{3.4}. Assertion 3), when $C$ is small,
follows from \hyperref[III.prop2.2]{2.2}. When $C$ is a $𝒰$-site, it
will be proved in \hyperref[III.subsec4.4]{4.4}. Assertion 4), when
$C$ is small, follows from the analogous commutativity assertion when
the topologies on $C$ and $C'$ are chaotic and from
\hyperrefext[I][I.cor3.5]{3.5}. When the topologies on $C$ and $C'$
are chaotic,\pageoriginale the commutativity of
\eqhyperref[III.eq2.3.2]{2.3.2} follows immediately from the explicit
description of $u_{\ast}$ (\hyperrefext[I][I.prop5.1]{5.1}).
\end{enonce*}

\setcounter{subsection}{3}
\subsection{}\etiquette[2.4]{III.subsec2.4}
We shall not develop here the considerations concerning
$\gamma$-objects of categories of sheaves. It will suffice to note
that the functors $u^{\ast}$ and $u_{\ast}$ introduced in this
section always commute with finite projective limits (unlike what
happened in the preceding section for the functor $u^s$).
Consequently, they extend naturally to functors defined on
$\gamma$-objects; these functors are adjoint to one another and
commute with the ``underlying sheaf of sets'' functors. We shall make
the notational abuses mentioned in \hyperref[III.not1.8]{1.8}, using
$u^{\ast}$ and $u_{\ast}$ also for the extensions to $\gamma$-objects.

\begin{enonce*}{Proposition 2.5}\etiquette[2.5]{III.prop2.5}
Let $C$ and $C'$ be two $𝒰$-sites, and let $\xymatrix{ C
\ar@<2pt>[r]^{v}& C'\ar@<2pt>[l]^{u} }$ be a pair of functors, where
$v$ is left adjoint to $u$. The following properties are equivalent:
\begin{enumerate}
\renewcommand{\theenumi}{\roman{enumi}}
\renewcommand{\labelenumi}{\theenumi)}
\item The functor $u$ is continuous.

\item The functor $v$ is cocontinuous.
\end{enumerate}
Moreover, under these equivalent conditions, there are canonical
isomorphisms $v_{\ast}\simeq u_{s}$ and $v^{\ast}\simeq u^s$. In
particular, the functor $u^s$ commutes with finite projective limits.
\end{enonce*}

\begin{enonce*}[remark]{Proof}
The property of a functor being continuous or cocontinuous does not
depend on the universe $𝒰$ for which $C$ and $C'$ are $𝒰$-sites
(this is immediate for a cocontinuous functor, and follows from
\hyperref[III.prop1.5]{1.5} for a continuous functor). We may
therefore, after enlarging the universe, suppose that $C$ and $C'$
are small. The proposition then follows from
\hyperref[III.prop2.2]{2.2} and \hyperrefext[I][I.prop5.6]{5.6}.
\end{enonce*}

\begin{enonce*}{Proposition 2.6}\etiquette[2.6]{III.prop2.6}
Let\pageoriginale $u:C\rightarrow C'$ be a continuous and
cocontinuous functor between two $𝒰$-sites. Then the functor
$u^{\ast}:{C'}^{\sim}\rightarrow C^{\sim}$
$(\hyperref[III.prop2.3]{2.3})$ commutes with $𝒰$-inductive and
projective limits. Denote by $u_{!}:C^{\sim}\rightarrow {C'}^{\sim}$
and $u_{\ast}:C^{\sim}\rightarrow {C'}^{\sim}$ the functors left and
right adjoint to $u^{\ast}$, respectively. The functor $u_{!}$ is
fully faithful if and only if the functor $u_{\ast}$ is fully
faithful. When $u$ is fully faithful, $u_{!}$ is fully faithful, and
the converse holds when the topologies of $C$ and $C'$ are coarser
than the canonical topology.
\end{enonce*}

The functor $u^{\ast}$ has a right adjoint $u_{\ast}$
(\hyperref[III.prop2.3]{2.3}) and a left adjoint $u_{!}$
(\hyperref[III.prop1.2]{1.2}). The functor $u^{\ast}$ therefore
commutes with inductive and projective limits
(\hyperrefext[I][I.prop2.11]{2.11}). The fact that $u_{!}$ is fully
faithful if and only if $u_{\ast}$ is fully faithful is a general
property of adjoint functors (\hyperrefext[I][I.rem5.7]{5.7} 1).
When $u$ is fully faithful, the functor
$\widehat{u}_{\ast}:C^{\widehat{\phantom{,}}}_{𝒱}\rightarrow
{C'}^{\widehat{\phantom{,}}}_{𝒱}$ on $𝒱$-presheaves, for a
sufficiently large universe $𝒱$, is fully faithful
(\hyperrefext[I][I.rem5.7]{5.7}). Consequently, by the commutativity
of \eqhyperref[III.eq2.3.1]{2.3.1} and
\eqhyperref[III.eq2.3.2]{2.3.2}, the functor
$u_{\ast}:C^{\sim}\rightarrow \widetilde{C}'$ is fully faithful, and
hence so is $u_{!}$ by what precedes. The converse follows from the
existence of the commutative diagram in \hyperref[III.prop1.2]{1.2}
iv), taking into account the fact that the functors $\epsilon_C$ and
$\epsilon_{C'}$ are fully faithful when the topologies are coarser
than the canonical topology, \cqfd

\begin{enonce*}[remark]{Example 2.7}\etiquette[2.7]{III.exem2.7}
With the notation of \hyperref[III.exem1.9.2]{1.9.2}, the functor
$$
\Phi:\Ouv(Y)\longrightarrow \Ouv(X)
$$
is continuous (\hyperref[III.exem1.9.2]{1.9.2}) and cocontinuous
(\hyperref[III.prop2.2]{2.2}). Moreover, let $i:Y\rightarrow X$ be
the canonical inclusion\pageoriginale and let
$i^{-1}:\Ouv(X)\rightarrow \Ouv(Y)$ be the functor that it defines
(\hyperref[III.exem1.9.1]{1.9.1}). The functor $i^{-1}$ is right
adjoint to the functor $\Phi$. There is therefore a sequence of three
adjoint functors:
\begin{align*}
\Phi^s\quad ,\quad & \Phi_s\\
        & ‖\\
        &\Phi^{\ast}\quad,\quad \Phi_{\ast},
\end{align*}
and canonical isomorphisms $\Phi^{\ast}\simeq (i^{-1})^s$ and
$\Phi_{\ast}\simeq i^{-1}_s$.
\end{enonce*}

% END INCLUDED SGA4_English_current_buildable_source_20260906/source/03/drafts/SGA4_III_LINES_613_861_English.texfrag


% BEGIN INCLUDED SGA4_English_current_buildable_source_20260906/source/03/drafts/SGA4_III_LINES_862_1113_English.texfrag
\section{Induced topology}\etiquette[3]{III.sec3}

\subsection{}\etiquette[3.1]{III.subsec3.1}
Let $C'$ be a site, $C$ a category, and $u:C\rightarrow C'$ a
functor. For every universe $𝒰$ such that $C'$ is a
$𝒰$-site and $C$ is a $𝒰$-small category, denote by
$𝒯_{𝒰}$ the finest among the topologies $T$ on $C$
that make $u$ continuous (\hyperref[III.def1.1]{1.1}). (Such a topology
exists by \hyperrefext[II][II.prop2.2]{2.2}.) The topology $𝒯_{𝒰}$ does
not depend on the universe $𝒰$. Indeed, if $𝒱\supset
𝒰$ is a universe, then $𝒯_{𝒰}=𝒯_{𝒱}$
(\hyperref[III.def1.1]{1.1} and \hyperref[III.prop1.5]{1.5}). The topology
$𝒯_{𝒰}$ is called \emph{the topology induced on $C$ by the topology
  of $C'$ by means of the functor $u$}
  \footnote[1]{When no confusion can result, this topology is called
    the \emph{topology induced on $C$ by the topology of $C'$.}}.

\begin{enonce*}{Proposition 3.2}\etiquette[3.2]{III.prop3.2}
Let $C$ be a small category, $C'$ a $𝒰$-site,
$u:C\rightarrow C'$ a functor, and $𝒯$ the topology on $C$ induced
by $u$. Let $X$ be an object of $C$, and let $R\rightarrow X$ be a
sieve on $X$. The following properties are equivalent:
\begin{enumerate}
\renewcommand{\theenumi}{\roman{enumi}}
\renewcommand{\labelenumi}{\theenumi)}
\item The sieve $R\hookrightarrow X$ is covering for $𝒯$.

\item For every base change $Y\rightarrow X$, where $Y$ is an
  object of $C$, the morphism $u_{!}(R\times_X Y)\rightarrow u(Y)$ is
  bicovering in ${C'}^{\widehat{\phantom{,}}}$ {\rm (\hyperrefext[II][II.def5.2]{5.2}).}
\end{enumerate}
\end{enonce*}

\begin{enonce*}[remark]{Proof}
\begin{itemize} %%% \XXX display problem [2016-11-23]: crossed-out i in i) below
\item[i) $\Rightarrow$ ii):] This follows from axiom (T 1) for
  topologies and from \hyperref[III.prop1.2]{1.2}.

\item[ii) $\Rightarrow$ i):] For every sheaf $F$ on $C'$, the map
$$
\Hom(Y,u^{\ast}F)\longrightarrow \Hom(R\times_X Y, u^{\ast}F)
$$
is\pageoriginale isomorphic, by adjunction (\hyperrefext[I][I.prop5.1]{5.1}), to
the map
$$
\Hom(u(Y),F)\longrightarrow \Hom(u_{!}(R\times_X Y),F),
$$
which is bijective (\hyperrefext[II][II.prop5.3]{5.3}). Consequently
(\hyperrefext[II][II.prop2.2]{2.2}), the sieve $R\hookrightarrow X$ is covering
for $𝒯$.
\end{itemize}
\end{enonce*}

\begin{enonce*}{Corollary 3.3}\etiquette[3.3]{III.cor3.3}
Let $C'$ be a site, $C$ a category, $u:C\rightarrow C'$ a
functor, and $𝒯$ the topology on $C$ induced by the topology of
$C'$. Let $(X_i\rightarrow X)$, $i\in I$, be a family of pullbackable
morphisms in $C$, and suppose that $u$ commutes with fiber products
\footnote[2]{In fact, it suffices that $u$ commute with fiber products
  of the form $\fprod{X_i}{X'}{X}$.}. The following properties are
equivalent:
\begin{enumerate}
\renewcommand{\theenumi}{\roman{enumi}}
\renewcommand{\labelenumi}{\theenumi)}
\item The family $X_i\rightarrow X$, $i\in I$, is covering for
  $𝒯$.

\item The family $(u(X_i)\rightarrow u(X))$, $i\in I$, is covering.
\end{enumerate}
\end{enonce*}

\begin{enonce*}[remark]{Proof}
\begin{itemize}
\item[i) $\Rightarrow$ ii):] This follows from \hyperref[III.prop1.6]{1.6}.

\item[ii) $\Rightarrow$ i):] Let $𝒰$ be a universe such that $C$
  is $𝒰$-small and $C'$ is a $𝒰$-site. Let
  $R\hookrightarrow X$ be the sieve generated by the $X_i\rightarrow
  X$. The presheaf $R$ is the coequalizer of the pair of arrows
(\hyperrefext[I][I.prop2.12]{2.12} and \hyperrefext[I][I.cor3.5]{3.5})
$$
\coprod_{i,j} \fprod{X_i}{X_j}{X}\rightrightarrows \coprod_i X_i
$$
(the coproduct is taken here in $C^{\wedge}$).
Since the functor $u_{!}$ commutes with inductive limits
(\hyperrefext[I][I.prop5.4]{5.4}), the presheaf $u_{!}R$ is the coequalizer of the
pair of arrows
$$
\coprod_{i,j}u(\fprod{X_i}{X_j}{X})\rightrightarrows \coprod_i
u(X_i){\quoi}.
$$
Since the functor $u$ commutes with fiber products, the presheaf
$u_{!}R$ is the coequalizer of the pair of arrows
$$
\coprod_{i,j}\fprod{u(X_i)}{u(X_j)}{u(X)}\rightrightarrows \coprod_i
u(X_i){\quoi},
$$
and\pageoriginale consequently $u_{!}R\rightarrow u(X)$ is the sieve on
$u(X)$ generated by the $u(X_i)\rightarrow u(X)$, $i\in I$.
Using once again the fact that $u$ commutes with fiber products, one
shows that for every base change $Y\rightarrow X$, where $Y$ is an
object of $C$, $u_{!}(\fprod{R}{Y}{X})\rightarrow u(Y)$ is the sieve
on $u(Y)$ generated by the
$\fprod{u(X_i)}{u(Y)}{u(X)}\rightarrow u(Y)$, $i\in I$. Since the
family $u(X_i)\rightarrow u(X)$, $i\in I$, is covering, the family
$\fprod{u(X_i)}{u(Y)}{u(X)}\rightarrow u(Y)$, $i\in I$, is covering.
Thus $u_{!}(\fprod{R}{Y}{X})\rightarrow u(Y)$ is a covering sieve,
and consequently (\hyperref[III.prop3.2]{3.2}) $R\rightarrow X$ is
covering for $𝒯$, \cqfd
\end{itemize}
\end{enonce*}

\begin{enonce*}{Corollary 3.4}\etiquette[3.4]{III.cor3.4}
Let $C'$ be a $𝒰$-site, let $C$ be a full subcategory
of $C'$, and let $u:C\rightarrow C'$ be the inclusion functor. Suppose
that fiber products are representable in $C$ and that $u$
commutes with fiber products. The following conditions are
equivalent:
\begin{enumerate}
\renewcommand{\theenumi}{\roman{enumi}}
\renewcommand{\labelenumi}{\theenumi)}
\item
\begin{enumerate}
\renewcommand{\theenumii}{\alph{enumii}}
\renewcommand{\labelenumii}{\theenumii)}
\item For every object $X$ of $C$, every covering family
  $(Y_j\rightarrow X)$, $j\in J$, in $C'$ is refined by a covering
  family $X_i\rightarrow X$, $i\in I$, where the $X_i$ are objects
  of $C$.

\item There exists a small set $G$ of objects of $C$ such that every
  object of $C$ is the target of a covering family in $C'$ of
  morphisms whose sources belong to $G$.
\end{enumerate}

\item The topology induced by the topology of $C'$ is a
  $𝒰$-topology {\rm (I \hyperref[II.def3.0.2]{3.0.2})}, and the functor $u$ is continuous and
  cocontinuous for this topology.
\end{enumerate}
\end{enonce*}

\begin{enonce*}[remark]{Proof}
This follows immediately from \hyperref[III.cor3.3]{3.3} and
\hyperref[III.def2.1]{2.1}.
\end{enonce*}

\begin{enonce*}{Proposition 3.5}\etiquette[3.5]{III.prop3.5}
Let\pageoriginale $C$ be a $𝒰$-site and let
$\epsilon_C:C\rightarrow C^{\sim}$ be the canonical functor {\rm
(\hyperrefext[II][II.subsubsec4.4.0]{4.4.0})}. Equip $C^{\sim}$ with the canonical
topology. Then the topology of the site $C$ is the topology induced
by the topology of $C^{\sim}$.
\end{enonce*}

\begin{enonce*}[remark]{Proof}
The covering families of $C^{\sim}$ for the canonical topology are
the universally effective epimorphic families
(\hyperrefext[II][II.def2.5]{2.5}), that is (\hyperrefext[II][II.prop4.3]{4.3}), the
epimorphic families. Let $T$ be the topology of the site $C$, and let
$𝒯_{𝒰}$ be the finest among the topologies $T'$ on $C$ such that,
for every sheaf $F$ on $C^{\sim}$, $F\circ\epsilon_C$ is a sheaf for
$T'$. It suffices to show that $T=𝒯_{𝒰}$, for then
$𝒯_{𝒰}$ is a $𝒰$-topology and consequently $𝒯_{𝒰}$ is the induced
topology.
\begin{itemize}
\item[A)] \emph{The topology $T$ is finer than $𝒯_{𝒰}$}:
Let $(X_i\rightarrow X)$, $i\in I$, be a covering family for
$𝒯_{𝒰}$. Then, for every sheaf $F$ on $C^{\sim}$, the map
$F(\epsilon_C(X))\rightarrow\prod_i F(\epsilon_C(X_i))$ is injective,
and consequently (\hyperrefext[II][II.def5.2]{5.2})
$(\epsilon_C X_i\rightarrow \epsilon_CX)$, $i\in I$, is covering for
the canonical topology of $C^{\sim}$, hence
(\hyperrefext[II][II.def5.2]{5.2}) $(X_i\rightarrow X)$, $i\in I$, is covering
for $T$.

\item[B)] \emph{The topology $T$ is coarser than
$𝒯_{𝒰}$}: It suffices to show that
$\epsilon_C:C\rightarrow C^{\sim}$ is continuous
(\hyperref[III.def1.1]{1.1}). Now we shall prove in IV ${\rm n^{\circ}}$ 1
that every sheaf on $C^{\sim}$ is representable, and consequently,
for every sheaf $F$ on $C^{\sim}$, $F\circ\epsilon_C$ is a sheaf on
$C$. (We shall not use \hyperref[III.prop3.5]{3.5} before
\hyperrefext[IV][IV.sec1]{1}.)
\end{itemize}
We record two results that will be used in VI 7.
\end{enonce*}

\begin{enonce*}{Proposition 3.6}\etiquette[3.6]{III.prop3.6}
Let $(C_i)_{i\in I}$ be a family of sites, let $C$ be a category,
and, for every $i\in I$, let $u_i:C_i\rightarrow C$ be a functor.
Let $\sheaf{\Uu}$ be a universe such that the categories $C_i$ and
$C$ are $𝒰$-small. There exists a topology $𝒯_{𝒰}$ on $C$ that is
the coarsest among the topologies for which the $u_i$ are continuous.
The topology $𝒯_{𝒰}$ does not depend on the universe $𝒰$ relative
to which the categories under consideration are small.
\end{enonce*}

The last assertion of \hyperref[III.prop3.6]{3.6} follows from
\hyperref[III.prop1.5]{1.5}.

Let $T$ be a topology on $C$. The functors $u_i$ are continuous if
and only if, for every $i\in I$ and every covering sieve
$R\hookrightarrow X$ on an object $X$ of $C_i$,\pageoriginale the
morphism
$$
u_{i!}(R)\longrightarrow u_i(X)
$$
in $\widehat{C}$ is bicovering (\hyperref[III.prop1.2]{1.2} (ii)). Thus
\hyperref[III.prop3.6]{3.6} is a consequence of the following lemma.

\begin{enonce*}{Lemma 3.6.1}\etiquette[3.6.1]{III.lem3.6.1}
Let $C$ be a small category and let $(u_i:F_i\rightarrow G_i)_{i\in
I}$ be a family of arrows in $\widehat{C}$. There exists on $C$ a
topology that is coarsest among those making the morphisms $u_i$
covering (\resp bicovering) {\rm(\hyperrefext[II][II.def5.2]{5.2})}.
\end{enonce*}

To say that $u:F\rightarrow G$ is covering for a given topology $T$
means that, for every arrow $X\rightarrow G$, with $X\in\Ob C$, the
corresponding arrow $F\times_G X\rightarrow X$ is covering; equivalently,
the family of arrows $X'\rightarrow X$ in $C$ that factor through the
preceding arrow is covering. The existence of a coarsest topology
among those for which the $u_i:F_i\rightarrow G_i$ are covering thus
follows from I \hyperref[II.subsubsec1.1.6]{1.1.6}, proving the assertion of
\hyperref[III.lem3.6.1]{3.6.1} without the parenthetical alternative.
The parenthetical assertion follows by recalling that a morphism
$u:F\rightarrow G$ is bicovering if and only if the morphisms
$u:F\rightarrow G$ and $\diag_u:F\rightarrow F\times_G F$ are
covering.

\begin{enonce*}{Proposition 3.7}\etiquette[3.7]{III.prop3.7}
Let $(C_i)_{i\in I}$ be a family of sites, let $C$ be a category,
for every $i\in I$ let $u_i:C_i\rightarrow C$ be a functor, and let
$\sheaf{\Uu}$ be a universe such that the categories $C_i$ and $C$
are $𝒰$-small. There exists a topology $\mathfrak{T}_{𝒰}$ that is
the finest for which the $u_i$ are cocontinuous. The topology
$\mathfrak{T}_{𝒰}$ does not depend on the universe $𝒰$ relative to
which the categories under consideration are small.
\end{enonce*}

Let $𝒰$ be a universe relative to which the categories under
consideration are small. The functors $u_i$ are cocontinuous for a
topology $\mathfrak{T}$ on $C$ if and only if, for every $i\in I$ and
every sheaf $F$ on $C_i$, the presheaf
$\widehat{u}_{i^{\ast}}F$ is a sheaf for $\mathfrak{T}$
(\hyperref[III.prop2.2]{2.2}). It follows that the topology
$\mathfrak{T}_{𝒰}$ is the finest topology for which the presheaves
$\widehat{u}_{i^\ast}F$, $i\in I$, $F\in C_i^{\sim}$, are sheaves
(\hyperrefext[II][II.prop2.2]{2.2}). The last assertion follows from
\hyperref[III.prop2.2]{2.2}.


% END INCLUDED SGA4_English_current_buildable_source_20260906/source/03/drafts/SGA4_III_LINES_862_1113_English.texfrag


% BEGIN INCLUDED SGA4_English_current_buildable_source_20260906/source/03/drafts/SGA4_III_LINES_1114_1313_English.texfrag
\section{Comparison lemma}\etiquette[4]{III.sec4}\pageoriginale

\begin{enonce*}{Theorem 4.1}\etiquette[4.1]{III.theo4.1}
{\rm(comparison lemma)}. Let $C$ be a small category, let $C'$ be a
site whose underlying category is a $𝒰$-category, and let
$u:C\rightarrow C'$ be a fully faithful functor. Endow $C$ with the
topology induced by $u$ {\rm(\hyperref[III.subsec3.1]{3.1}).} Consider
the following properties:
\begin{enumerate}
\renewcommand{\theenumi}{\roman{enumi}}
\renewcommand{\labelenumi}{\theenumi)}
\item Every object of $C'$ can be covered by objects coming from $C$.

\item The functor $F\mapsto F\circ u$ induces an equivalence from the
category of sheaves on $C'$ to the category of sheaves on $C$.
\end{enumerate}

We always have {\rm i) $\Rightarrow$ ii).} When $C'$ is a $𝒰$-site
and the topology of $C'$ is coarser than the canonical topology, we
have {\rm ii) $\Rightarrow$ i).}
\end{enonce*}

\setcounter{subsection}{1}
\subsubsection{}\etiquette[4.1.1]{III.subsubsec4.1.1}
Let us first prove i) $\Rightarrow$ ii). The proof proceeds in two
steps.

%\emph{ %% causes compilation problems
First step. For every presheaf $H$ on $C'$, the adjunction morphism
$\varphi:u_{!}u^{\ast}H\rightarrow H$ {\rm(\hyperrefext[I][I.prop5.1]{5.1})} is
bicovering {\rm(\hyperrefext[II][II.prop5.3]{5.3})}, and the functor
$u_s:\widetilde{C'}\rightarrow \widetilde{C}$ $\hyperref[III.subsec1.11]{(1.1.1)}$ is fully
faithful.
%} %% cf. above

Let $C/H$ be the small category whose objects are the objects $Y$ of
$C$ equipped with a morphism $u(Y)\rightarrow H$, and whose
morphisms are the commutative diagrams
\[
\xymatrix@=1.25cm{
u(Y) \ar[rr]^{u(m)} \ar[rd]&     &u(Y')\ar[ld]\\
     &  H  & .
}
\]
We have $u^{\ast}H=\varinjlim_{Y\in \ob C/H} Y$
(\hyperrefext[I][I.prop3.4]{3.4}), and hence
(\hyperrefext[I][I.prop5.4]{5.4})
$u_{!}u^{\ast}H=\varinjlim_{Y\in \ob C/H} u(Y)$. The adjunction
morphism is the evident morphism arising from the description of
$u_{!}u^{\ast}H$ as an inductive limit. The morphism $\varphi$ has
the following property:\pageoriginale $(\ast)$ For every object $Y$
of $C$, every morphism $m:u(Y)\rightarrow H$ factors uniquely as the
canonical morphism $\alpha(m):u(Y)\rightarrow u_{!}u^{\ast}H$
followed by $\varphi$. It follows immediately from i) that
$\varphi$ is covering (\hyperrefext[II][II.prop5.1]{5.1}); let us show
that it is bicovering. Let
$p,q:Z\rightrightarrows u_{!}u^{\ast}H$ be two morphisms from an
object $Z$ of $C'$ such that $\varphi p=\varphi q$. For every object
$Y$ of $C$ and every morphism $n:u(Y)\rightarrow Z$, we have
$\varphi pn=\varphi qn$. Property $(\ast)$ then implies that
$pn=qn$, and since the $u(Y)$ cover $Z$, the kernel of $(p,q)$ is a
covering sieve of $Z$. Thus $\varphi$ is bicovering
(\hyperrefext[II][II.prop5.3]{5.3}). For every sheaf $H$ on $C'$,
$u^{\ast}H$ is a sheaf on $C$, denoted by $u_sH$ \hyperref[III.subsec1.11]{(1.1.1)}. Moreover,
$u^su_sH=\underline{a}'u_{!}u_sH$ (\hyperref[III.prop1.3]{1.3}), and
the adjunction morphism $u^su_sH\rightarrow H$ is obtained by
applying the ``associated sheaf'' functor to
$\varphi:u_{!}u^{\ast}H\rightarrow H$. Consequently
(\hyperrefext[II][II.prop5.3]{5.3}), the adjunction morphism
$u^su_sH\rightarrow H$ is an isomorphism. Hence
$u_s:\widetilde{C}'\rightarrow \widetilde{C}$ is fully faithful.

\noindent \emph{Second step. The functor $u$ is cocontinuous and the
functor $u^s:C^{\sim}\rightarrow {C'}^{\sim}$
$(\hyperref[III.prop1.2]{1.2})$ is fully faithful. Consequently
$u_s:{C'}^{\sim}\rightarrow C^{\sim}$ is an equivalence.} Let $Y$
be an object of $C$ and let $i:R\hookrightarrow u(Y)$ be a covering
sieve. Since the functor $u^{\ast}$ commutes with projective limits
(\hyperrefext[I][I.prop5.5]{5.5}), the morphism
$u^{\ast}(i):u^{\ast}(R)\hookrightarrow u^{\ast}u(Y)$ is a
monomorphism. Since $u$ is fully faithful, we have
$u^{\ast}u(Y)\simeq Y$, and thus obtain a sieve of $Y$,
$u^{\ast}(i):u^{\ast}(R)\rightarrow Y$. To show that $u$ is
cocontinuous, it suffices to show that
$u^{\ast}(i):u^{\ast}(R)\hookrightarrow Y$ is a covering sieve for
the induced topology on $C$ (\hyperref[III.def2.1]{2.1}). For this,
it suffices to show (\hyperref[III.prop3.2]{3.2}) that for every base
change $m:X\rightarrow Y$, the morphism
$u_{!}(u^{\ast}(R)\times_Y X)\rightarrow u(X)$ is bicovering. But
since $u^{\ast}$ commutes with\pageoriginale projective limits, we
have
$\fprod{u^{\ast}(R)}{X}{Y}=u^{\ast}(\fprod{R}{u(X)}{u(Y)})$, and
$\fprod{R}{u(X)}{u(Y)}$ is a covering sieve of $u(X)$. It therefore
suffices to show that for every object $Y$ of $C$ and every covering
sieve $i:R\hookrightarrow u(Y)$, the morphism in
${C'}^{\widehat{\phantom{,}}}$,
$u_{!}u^{\ast}(R)\xrightarrow{u_{!}u^{\ast}(i)}u(Y)$, is
bicovering. Now this morphism factors as the adjunction morphism
$u_{!}u^{\ast}(R)\rightarrow R$, which is bicovering (first step),
followed by the covering monomorphism $R\hookrightarrow u(Y)$. It is
therefore bicovering (\hyperrefext[II][II.prop5.3]{5.3}). This shows
that $u$ is cocontinuous. Since $u$ is continuous and fully
faithful, it follows from \hyperref[III.prop2.6]{2.6} (used in the
case where $C$ is small) that $u^s$ (denoted by $u_{!}$ in
\hyperref[III.prop2.6]{2.6}) is fully faithful. Since $u^s$ and
$u_s$ are adjoint to one another and are fully faithful, they are
quasi-inverse functors, and hence $u_s$ is an equivalence.

\setcounter{subsection}{1}
\subsubsection{}\etiquette[4.1.2]{III.subsubsec4.1.2}
Let us now prove ii) $\Rightarrow$ i). The objects $Y$ of $C$ form a
generating family of $C^{\sim}$ (\hyperrefext[II][II.prop4.10]{4.10}),
and consequently, for every object $X$ of $C'$, there exists an
epimorphic family (in $C^{\sim}$)
$v_i:Y_i\rightarrow u_s X$, $i\in I$. We have
$u_su(Y_i)=Y_i$, and, since $u_s$ is an equivalence of categories,
$v_i=u_s(w_i)$, where the $w_i:u(Y_i)\rightarrow X$ form an
epimorphic family in ${C'}^{\sim}$. It then follows from
\hyperrefext[II][II.prop5.1]{5.1} that the family
$w_i:u(Y_i)\rightarrow X$, $i\in I$, is covering.

\subsection{Completion of the proof of \hyperref[III.prop1.2]{1.2}}\etiquette[4.2]{III.subsec4.2}

Let $G$ be a full subcategory of $C$ whose objects form a small
topologically generating family of $C$ (II \hyperref[II.def3.0.1]{3.0.1}). Endow $G$ with
the topology induced by the inclusion functor $i:G\rightarrow C$
(\hyperref[III.subsec3.1]{3.1}). The functor
$(u\circ i)_s=i_s\circ u_s$ \hyperref[III.subsec1.11]{(1.1.1)} has a left adjoint by the first
part of the proof of \hyperref[III.prop1.2]{1.2}. Since $i_s$ is an
equivalence of\pageoriginale categories (\hyperref[III.theo4.1]{4.1}),
the functor $u_s$ has a left adjoint $u^s$, and there is a functorial
isomorphism $u^s\simeq (u\circ i)^s\circ i_s$. Let us show that the
diagram
\hypertarget{sga4.exp.III.equation.4.2.1}{}\[\vcenter{
\xymatrix@C=3cm@R=1.5cm{
C \ar[r]^{u} \ar[d]_{\epsilon_C}& C'\ar[d]^{\epsilon_{C'}}\\
C\widetilde{\phantom{a}} \ar[r]^{u^{s}}& \widetilde{C}'
}}\leqno{(4.2.1)}
\]
commutes up to isomorphism. For every sheaf $H$ on $C'$, we have
$\Hom_{\widetilde{C}'}(\epsilon_{C'}uX,H)\simeq u_sH(X)$ for every
$X\in \ob C$. Moreover,
$\Hom_{\widetilde{C}'}(u^s\epsilon_CX,H)\simeq
\Hom_{\widetilde{C}}(\epsilon_CX,u_sH)$ by adjunction, and then
$\Hom_{\widetilde{C}}(\epsilon_CX,u_sH)\simeq u_sH(X)$ by the
definition of $\epsilon_C$. Hence there is an isomorphism
$\epsilon_{C'}uX\simeq u^s\epsilon_CX$ for every $X\in \ob C$. This
proves i) $\Rightarrow$ iv). Let us prove iv) $\Rightarrow$ i). We
have a commutative diagram
\hypertarget{sga4.exp.III.equation.4.2.2}{}\[\vcenter{
\xymatrix@C=2.25cm@R=1.5cm{ G \ar[r]^{i} \ar[d]_{\epsilon_G} & C
\ar[r]^{u} \ar[d]_{\epsilon_C}&
C' \ar[d]_{\epsilon_{C'}}\\
\widetilde{G} \ar[r]^{i^{s}}& \widetilde{C} \ar[r]^{u^s}&
\widetilde{C'}, }}\leqno{(4.2.2)}
\]
and consequently, by the first part of the proof, the functor
$u\circ i:G\rightarrow C'$ is continuous. It follows at once from
\hyperref[III.theo4.1]{4.1} and \hyperref[III.def1.1]{1.1} that $u$
is continuous, whence iv) $\Rightarrow$ i). This also gives the
uniqueness of $u^s$, since the first part of the proof gives
uniqueness when $C$ is small.

\subsection{Completion of the proofs of \hyperref[III.prop1.5]{1.5} and \hyperref[III.prop1.7]{1.7}}\etiquette[4.3]{III.subsec4.3}
\pageoriginale

Let $u:C\rightarrow C'$ be a functor between two $𝒰$-sites, and let
$G$ be a full subcategory of $C$ whose set of objects is a small
topologically generating family of $C$ (II \hyperref[II.def3.0.1]{3.0.1}). Endow $G$ with
the topology induced by the inclusion functor $i:G\rightarrow C$.
It follows immediately from \hyperref[III.theo4.1]{4.1} and
\hyperref[III.def1.1]{1.1} that $u$ is continuous if and only if
$u\circ i$ is continuous. This observation and the first part of
the proof of \hyperref[III.prop1.5]{1.5} give the general case. The
same observation also reduces the proof of \hyperref[III.prop1.7]{1.7}
to the case where the category $C$ is small.

\subsection{Completion of the proof of \hyperref[III.prop2.3]{2.3}}\etiquette[4.4]{III.subsec4.4}

Let $u:C\rightarrow C'$ be a functor between two $𝒰$-sites, and let
$G$ be a full subcategory of $C$ whose set of objects is a small
topologically generating family of $C$
(\hyperrefext[II][II.def3.0.1]{3.0.1}), endowed with the topology
induced by the inclusion functor $i:G\rightarrow C$. It follows from
the proof of \hyperref[III.theo4.1]{4.1} (\hyperref[III.subsubsec4.1.1]{4.1.1}, second step) that
$i$ is cocontinuous. Consequently, when $u$ is cocontinuous, the
functor $u\circ i:G\rightarrow C'$ is cocontinuous. Moreover,
$i^{\ast}:\widetilde{C}\rightarrow \widetilde{G}$ is an equivalence
of categories (\hyperref[III.theo4.1]{4.1}). Therefore the functor
$u^{\ast}:\widetilde{C}'\rightarrow \widetilde{C}$ has a right
adjoint. This proves 3). To prove 4), it suffices to observe that
$(u\circ i)_{\ast}\simeq u_{\ast}\circ i_{\ast}$ and that
$i_{\ast}$ is an equivalence (\hyperref[III.theo4.1]{4.1}). The
commutativity of \eqhyperref[III.eq2.3.2]{2.3.2} then follows from
the commutativity of these diagrams when $C$ is small.


% END INCLUDED SGA4_English_current_buildable_source_20260906/source/03/drafts/SGA4_III_LINES_1114_1313_English.texfrag


% BEGIN INCLUDED SGA4_English_current_buildable_source_20260906/source/03/drafts/SGA4_III_LINES_1314_1551_English.texfrag
\section{Localization}\etiquette[5]{III.sec5}\pageoriginale

\subsection{}\etiquette[5.1]{III.subsec5.1}
Now let $C$ be a $𝒰$-site and $X$ an object of
$C^{\widehat{\phantom{,}}}$. Unless expressly stated otherwise, the
category $C/X$ will be endowed with the topology $𝒯$ induced by the
functor $j_X:C/X\rightarrow C$ (\hyperref[III.subsec3.1]{3.1}). The notation
$C/X$ will denote the category $C/X$ endowed with the topology $𝒯$.
Proposition \hyperrefext[I][I.prop5.11]{5.11} shows that the functor $j_{X!}$
commutes with fiber products and consequently sends every
monomorphism to a monomorphism. In particular, for every object
$(Y\rightarrow X)$ of $C/X$, the functor $j_{X!}$ establishes a
one-to-one correspondence between the sieves, in the category
$C/X$, on the object $(Y\rightarrow X)$ and the sieves, in $C$, on
the object $Y$.

\begin{enonce*}{Proposition 5.2}\etiquette[5.2]{III.prop5.2}
Let $C$ be a $𝒰$-site, $X$ an object of
$C^{\widehat{\phantom{,}}}$, and $j_X:C/X\rightarrow C$ the
corresponding continuous functor.
\begin{enumerate}
\renewcommand{\labelenumi}{\theenumi)}
\item A sieve $R$ on an object $(Z\rightarrow X)$ is covering in
$C/X$ if and only if the sieve $j_{X!}(R)\hookrightarrow Z$ is
covering in $C$.

\item The functor $j_X:C/X\rightarrow C$ is cocontinuous and
continuous.

\item Let $m:Y\rightarrow X$ be a morphism of
$C^{\widehat{\phantom{,}}}$. We then have the commutative diagram
\[
\xymatrix@C=1cm@R=1.25cm{
C/Y \ar[rr]^{j_m} \ar[rd]_{j_Y}&       &C/X \ar[ld]^{j_X}\\
    &  C    &.
}
\]
The topology induced by the functor $j_Y$ on $C/Y$ is equal to the
topology induced by $j_m$ on $C/Y$.

\item The topology induced by $j_X:C/X\rightarrow C$ is a
$𝒰$-topology (\hyperrefext[II][II.def3.0.2]{3.0.2}).
\end{enumerate}
\end{enonce*}

\begin{enonce*}[remark]{Proof}
\begin{enumerate}
\renewcommand{\labelenumi}{\theenumi)}
\item If\pageoriginale the sieve $R$ on $(Z\rightarrow X)$ is
covering, the sieve $j_{X!}(R)\hookrightarrow Z$ is covering
(\hyperref[III.prop1.6]{1.6}).

Conversely, if the sieve $j_{X!}(R)\rightarrow Z$ is covering in
$C$, one sees that the same is true for every sieve obtained by a
base change in $C/X$. The sieve $R$ is therefore covering
(\hyperref[III.prop3.2]{3.2}).

\item This follows immediately from 1) by applying
\hyperref[III.def2.1]{2.1}.

\item This follows immediately from the description of covering
sieves given in 1).

\item Let $(G_i)_{i\in I}$ be a small topologically generating
family of $C$. One verifies immediately that the small family
$(u:G_i\rightarrow X)$, $u\in \amalg_{i\in
I}\Hom_{C^{\widehat{\phantom{,}}}}(G_i, X)$, is topologically
generating in $C/X$.
\end{enumerate}
\end{enonce*}

\phantomsection\MakeLinkTarget*{sga4.exp.III.terminology.5.3}\label{sga4.exp.III.terminology.5.3}\label{III.subsec5.3}%
\noindent \emph{Terminology and notation} 5.3. By the preceding
proposition, the functor $j_X$ is both continuous and cocontinuous.
It therefore defines a sequence of three adjoint functors (4.3.2)
between the categories of sheaves of sets (\hyperref[III.prop1.3]{1.3}
and \hyperref[III.prop2.3]{2.3}):
\begin{gather*}
j^s_X:(C/X)^{\sim}\longrightarrow
C^{\sim}\\
j^{\ast}_X=j_{X,s}:C^{\sim}\longrightarrow
(C/X)^{\sim}\\
j_{X\ast}:(C/X)^{\sim}\longrightarrow C^{\sim}{\quoi}.
\end{gather*}

In the particular situation of Proposition \hyperref[III.prop5.2]{5.2},
we shall use the following terminology and notation:
\begin{enumerate}
\renewcommand{\labelenumi}{\theenumi)}
\item The functor $j_{X\ast}$ will be called the \emph{direct image
functor.}

\item The functor $j_{X,s}=j^{\ast}_X$ will be denoted by
$j^{\ast}_X$ and will be called the \emph{restriction functor to
$C/X$.}

\item The\pageoriginale functor $j^s_X$ on sheaves of sets will be
called the ``extension by the empty set to $C$''
functor and will be denoted by $j_{X!}$ (\cf 2.9.2).
\end{enumerate}

Thus we have a sequence of three adjoint functors between
$(C/X)^{\sim}$ and $C^{\sim}$:
$$
j_{X!},\; j^{\ast}_X,\; j_{X\ast}{\quoi}.
$$

\begin{enonce*}{Proposition 5.4}\etiquette[5.4]{III.prop5.4}
The functor $j_{X!}:(C/X)^{\sim}\rightarrow C^{\sim}$ factors
through the category $C^{\sim}/\underline{a}X$ ($\underline{a}$ is
the associated sheaf functor):

$(C/X)^{\sim}\xrightarrow{e^{\sim}_X} C^{\sim}/\underline{a}X\rightarrow
C^{\sim}$. The functor
$$
e^{\sim}_X:(C/X)^{\sim}\longrightarrow \widetilde{C}/\underline{a}X
$$
is an equivalence of categories. The restriction functor to $C/X$,
followed by the equivalence $e^{\sim}_X$, is isomorphic to the
``base change by $\underline{a}X\rightarrow \epsilon$'' functor
($\epsilon$ being the terminal object of $C^{\sim}$):
$F\mapsto (F\times
\underline{a}X\xrightarrow{\pr_2}\underline{a}X)$.
\end{enonce*}

\begin{enonce*}[remark]{Proof}
The image under $j_{X!}$ of the terminal object of $(C/X)^{\sim}$ is
the object $\underline{a}X$, which gives the factorization. To show
that the functor $e^{\sim}_X$ is an equivalence, we shall confine
ourselves to a few indications. By \hyperrefext[I][I.prop5.11]{5.11}, a
presheaf on $C/X$ is defined by a presheaf $F$ on $C$ equipped with
a morphism $F\rightarrow X$. One then proves that the following two
properties are equivalent:
\begin{enumerate}
\renewcommand{\theenumi}{\roman{enumi}}
\renewcommand{\labelenumi}{\theenumi)}
\item The presheaf on $C/X$ defined by $F\rightarrow X$ is a
sheaf.

\item The following diagram is cartesian (where $i$ and
$\underline{a}$ denote the inclusion in the category of presheaves
and the associated sheaf functor, respectively):
\[
\xymatrix@C=2.75cm@R=1cm{
F \ar[r]\ar[d]& i\mathsf{\underline{a}}F \ar[d]\\
X \ar[r]& i\mathsf{\underline{a}}X\quad . }
\]
\end{enumerate}

It\pageoriginale follows immediately that $e^{\sim}_X$ is an
equivalence. The last assertion is trivial.
\end{enonce*}

\begin{enonce*}{Proposition 5.5}\etiquette[5.5]{III.prop5.5}
\begin{enumerate}
\renewcommand{\labelenumi}{\theenumi)}
\item Let $C$ be a $𝒰$-site and $X$ an object of
$C^{\widehat{\phantom{,}}}$. The diagram below of categories and
functors commutes up to canonical isomorphisms:
\[
\xymatrix@C=1.5cm@R=1.5cm{ C/X \ar[r]^{h_X} \ar[d]&
(C/X)\widehat{\phantom{a}} \ar[d]^{e\widehat{\phantom{a}}_X}
\ar[r]^{\underline{a}_X}& (C/X)\widetilde{\phantom{a}} \ar[rr]^{i_X}
\ar[d]^{e\widetilde{\phantom{a}}_X}& &
(C/X)\widehat{\phantom{a}} \ar[d]^{e\widehat{\phantom{a}}/X}\\
C/X \ar[r]^{h/X}& C\widehat{\phantom{a}}/X \ar[r]^{\underline{a}/X}&
C\widetilde{\phantom{a}}/\underline{a}X \ar[r]^{i/\underline{a}X}&
C\widehat{\phantom{a}}/i\underline{a}X \ar[r]^{\pi}&
C\widehat{\phantom{a}}/X. }
\]

The notations $\underline{a}_X$ and $i_X$ denote the ``associated
sheaf'' and inclusion into presheaves functors for the site $C/X$.
The notation $\underline{a}/X$ denotes the natural extension of the
functor $\underline{a}$ (the associated sheaf functor for the site
$C$) to the category of arrows with target $X$; similarly for the
notation $i/\underline{a}X$. Finally, the notation $\pi$ denotes the
base-change functor along the canonical arrow
$X\rightarrow i\underline{a}X$.

\item Let, in addition, $m:Y\rightarrow X$ be a morphism of
$C^{\widehat{\phantom{,}}}$. The following diagram commutes up to
canonical isomorphism:
\[
\xymatrix@C=1.75cm@H=.5cm{ C\widetilde{\phantom{a}}
/\underline{a}X/\underline{a}Y \ar@{=}[dd]&
(C/X)\widetilde{\phantom{a}}/\underline{a}_X Y \ar[l]_{f}& \\
& & (C/X/Y)\widetilde{\phantom{a}} \ar[lu]_{g} \ar[ld]_{e\widetilde{\phantom{a}}_m}\\
C\widetilde{\phantom{a}}/\underline{a}Y & (C/Y)\widetilde{\phantom{a}}
\ar[l]_{e\widetilde{\phantom{a}}_Y}&. }
\]

The\pageoriginale arrow $f$ is the natural extension of the
equivalence $e^{\sim}/X$ to the category of arrows with target
$\underline{a}_XY$, and is consequently an equivalence of
categories. The arrow $g$ is precisely the equivalence of
\hyperref[III.prop5.4]{5.4} applied to the situation $(C/X)/[m]$.

\item Denote by
$j_{\underline{a}X!}:C^{\sim}/\underline{a}X\rightarrow C^{\sim}$ the
``forget $\underline{a}X$'' functor, and by
$j^{\ast}_{\underline{a}X}:C^{\sim}\rightarrow C^{\sim}/\underline{a}X$ the
``product with $\underline{a}X$'' functor. The following diagrams
commute up to canonical isomorphism:
\begin{displaymath}
\xymatrix@C=2.25cm@R=1.25cm{ (C/X)^{\widetilde{\phantom{a}}}
\ar[r]^{j_{X!}} \ar[d]   &
  C\widetilde{\phantom{a}} \ar@{=}[d]\\
C\widetilde{\phantom{a}}/\underline{a}X \ar[r]^{j_{\underline{a}X!}} &
C\widetilde{\phantom{a}} } \qquad {\raisebox{-1.75cm}{,}}\qquad
\xymatrix@C=2.5cm@R=1.25cm{
C{\widetilde{\phantom{a}}} \ar[r]^{j_X^{\ast}} \ar@{=}[d]&  (C/X)\widetilde{\phantom{a}}\ar[d]^{e\widetilde{\phantom{a}}/X}\\
C{\widetilde{\phantom{a}}} \ar[r]^{j_{\underline{a}X}^{\ast}} &
C{\widetilde{\phantom{a}}}/\underline{a}X. }
\end{displaymath}
\end{enumerate}
\end{enonce*}

\begin{enonce*}[remark]{Proof}
The proof of these assertions follows immediately from the
definition of the equivalences $e^{\sim}_X$
(\hyperref[III.prop5.4]{5.4}) and
$e^{\widehat{\phantom{\,}}}_X$ (\hyperrefext[I][I.prop5.11]{5.11}).
\end{enonce*}

\begin{thebibliography}{99}%
\bibitem[TF]{sga4.exp.III.bib.TF} R. Godement, théorie des faisceaux, Hermann,
1958. Act. Scient.

Ind ${\rm n^{\circ}}$ 1252, (Paris).
\end{thebibliography}


%I didn't get the label...
%---\widetilde{C}$ $(1.1.1)$ is fully faithful.}
%---\emph{Terminology and notation} 5.3. By the proposition
%---4.3.2
%---2.9.2

\ifx\danslelivre\undefined
\end{document}
\fi

% END INCLUDED SGA4_English_current_buildable_source_20260906/source/03/drafts/SGA4_III_LINES_1314_1551_English.texfrag


% END INCLUDED SGA4_English_current_buildable_source_20260906/source/03/03_English.tex


\makeatletter
\def\input@path{{../04/}}
\makeatother
% BEGIN INCLUDED SGA4_English_current_buildable_source_20260906/source/04/04_English.tex
\ifx\danslelivre\undefined
% Release engine: XeLaTeX (the LuaTeX compatibility layer breaks Xy-pic output).
\documentclass[10pt]{../fichiers-configuration-latex/smfart}
% BEGIN INCLUDED SGA4_English_current_buildable_source_20260906/source/fichiers-configuration-latex/paquets.tex
\usepackage{polyglossia}
%\usepackage[francais]{babel}
\usepackage{fontspec}
\usepackage[math-style=ISO,bold-style=upright,sans-style=upright]{unicode-math}
%\usepackage[mathletters]{ucs}
%\usepackage[utf8]{inputenc}
%\usepackage[T1]{fontenc}
\usepackage[unicode,hypertexnames=false,colorlinks=true,linkcolor=blue,citecolor=Sepia,urlcolor=magenta]{hyperref}
\setmainlanguage{french}
\setmainfont[Ligatures=TeX]{Linux Libertine O}
% variante possible : TeXGyrePagella ?
\setmathfont{XITS Math}

% maths
\usepackage{amsmath}
%\usepackage{amssymb}
%\usepackage{mathrsfs}
%\usepackage{stmaryrd}
%\usepackage{euscript}

% couleurs
\usepackage[usenames,dvipsnames]{xcolor}

% graphiques
\usepackage{graphicx}
\usepackage[all]{xy}
%\usepackage{tikz}
%\usetikzlibrary{matrix,arrows,calc}
%\tikzset{column sep=3em,row sep=2em,text height=1.5ex, text depth=0.25ex,>=angle 90}
%\tikzset{description/.style={fill=white,inner sep=1pt}}

% hyperliens, 'etiquettes
\usepackage{xr-hyper}
\usepackage{../fichiers-configuration-latex/smfthm}
%\ifx\cacheretiquettes\undefined
%\usepackage[inline]{showlabels} \else \fi % pour rendre visibles les 'etiquettes
%\usepackage{../fichiers-configuration-latex/makeidx} % pour faire un index
%\usepackage[pagebackref,unicode]{hyperref}  % pour voir o`u sont cit'ees les r'ef'erences
% hyperref is loaded once above with the cumulative-reader destination policy.
%\hypersetup{breaklinks=true} %permet le retour `a la ligne dans les liens trop longs. Incompatible avec pdflatex
%\renewcommand*{\backref}[1]{{\small [cit\'e page(s)~#1.]}}

\usepackage{xspace}
\usepackage{pifont} % pour expose XVII

% END INCLUDED SGA4_English_current_buildable_source_20260906/source/fichiers-configuration-latex/paquets.tex

% BEGIN INCLUDED SGA4_English_current_buildable_source_20260906/source/fichiers-configuration-latex/numerotation.tex
\renewcommand{\thefootnote}{\roman{footnote}} % pour mettre les notes en bas de page sous la forme (ii) et non (2), qui ressemble trop (?) a un exposant.

\newcounter{pageoriginal}
\newcommand{\pageoriginale}{\refstepcounter{pageoriginal}\marginpar{\footnotesize\textbf{\thepageoriginal}}\xspace}
\let\pageoriginaled\pageoriginale

\usepackage{ifthen}
\newboolean{original}
\def\sisi{\ifthenelse{\boolean{original}}}
\setboolean{original}{false}

\newcounter{cnote}
\def\nde#1{\ifthenelse{\boolean{original}}{}{\stepcounter{cnote}\def\thefootnote{\arabic{cnote}}\protect\footnote[\value{cnote}]{N.D.E.: #1}}}
\def\ndemark{\ifthenelse{\boolean{original}}{}{\stepcounter{cnote}\def\thefootnote{\arabic{cnote}}\protect\footnotemark[\value{cnote}]}}
\def\ndetext#1{\ifthenelse{\boolean{original}}{}{\def\thefootnote{\arabic{cnote}}\protect\footnotetext[\value{cnote}]{N.D.E.: #1}}}

% END INCLUDED SGA4_English_current_buildable_source_20260906/source/fichiers-configuration-latex/numerotation.tex

% BEGIN INCLUDED SGA4_English_current_buildable_source_20260906/source/fichiers-configuration-latex/macros.tex
\setmainlanguage{french}
\setotherlanguage{english}


%% Fixed, engine-independent release identifier.  The inherited dynamic Git
%% probe used Lua io.popen, exposed build-state-dependent output, and printed
%% a visible question mark outside a Git checkout.
\def\showgitstatus{English release r2, 2026-07-23}


% macros pour LaTeX
\renewcommand{\mathsf}{}

% Police différente pour U+2113 SCRIPT SMALL L
\setmathfont[range="2113]{latinmodern-math.otf}

% Give every inherited semantic label its own stable PDF destination instead
% of inheriting whichever automatic equation/page anchor happened to be live.
% The legacy label remains as a compatibility alias used by existing edges.
\makeatletter
\def\etiquette[#1]#2{%
  \Hy@raisedlink{\hyper@anchorstart{sga4.semantic.#2}\hyper@anchorend}%
  \def\@currentHref{sga4.semantic.#2}%
  \def\@currentlabel{#1}%
  \label{#2}%
}
\makeatother
\def\eqhyperref[#1]#2{(\hyperref[#1]{#2})}
\def\pagehyperref[#1]#2{(page \hyperref[#1]{#2})}
\def\hyperrefext[#1][#2]#3{\hyperref[#2]{#1 #3}}
\def\eqhyperrefext[#1][#2]#3{\hyperref[#2]{#1 #3}}     %cf. exposes 17
\def\neweqhyperrefext[#1][#2]#3{\hyperref[#2]{#1 #3}}  %cf. exposes 17

\def\shat#1{\protect\widehat{#1}}
\def\stilde#1{\protect\widetilde{#1}}
%%%%%%%%%%%%%%%%%%


\def\etcn{etc.\xspace}
\def\etco{etc...\xspace}
\def\etc{\sisi{\etco}{\etcn}}
\def\ldotsn{\sisi{...}{\etcn}}

\def\cf{{cf.}\kern.3em}
\def\Cf{{Cf.}\kern.3em}
\def\ie{{i.e.}\kern.3em}
\def\iev{{i.e.},\xspace}
\def\loccito{loc.\kern2pt cit.\xspace}
\def\loccitn{\textit{loc.\kern2pt cit.}\xspace}
\def\loccit{\sisi{\loccito}{\loccitn}}
\def\EGA{\textup{EGA}\kern.2em}
\def\SGA{\textup{SGA}\kern.3em}

\def\quoi{}

\def\quav{\quad,\quad}
\def\quaet{\quad \text{et}\quad}
\def\fibreproduct#1#2#3{#1 ×_{#3} #2}
\let\fprod\fibreproduct
\def\ccheck#1#2{\mbox{${\check{\check{\rm #1}}}^{\lower2pt\hbox{$\scriptstyle #2$}}$}}
\def\sheaf#1{\protect\mathcal{#1}}
\def\vieuxsheaf#1{{#1}} % cf. problème avec sheaf{a\Phi_i} par exemple...
\DeclareMathOperator{\Hom}{Hom}
\DeclareMathOperator{\can}{can}
\DeclareMathOperator{\Ind}{Ind}
\DeclareMathOperator{\ind}{ind}
%\DeclareMathOperator{\Hhomtop}{\mathcal{H}\!\mathit{omtop}}
\DeclareMathOperator{\Hhomtop}{\mathbf{Homtop}}
\DeclareMathOperator{\Point}{Point}
\DeclareMathOperator{\Pro}{Pro}
\DeclareMathOperator{\Top}{Top}
\DeclareMathOperator{\Mod}{Mod}
\DeclareMathOperator{\Cart}{Cart}
%\DeclareMathOperator{\Mmorsite}{\mathcal{M}\!\mathit{orsite}}
\DeclareMathOperator{\Mmorsite}{\mathbf{Morsite}}
\DeclareMathOperator{\loc}{loc}
\DeclareMathOperator{\Loc}{Loc}
\DeclareMathOperator{\Pic}{Pic}
\DeclareMathOperator{\Et}{Et}
\DeclareMathOperator{\card}{card}
\DeclareMathOperator{\Ens}{Ens}
\DeclareMathOperator{\h}{\underline{H}}
\DeclareMathOperator{\Fonct}{Fonct}
\DeclareMathOperator{\Ker}{Ker}
\DeclareMathOperator{\but}{but}
\DeclareMathOperator{\Ob}{Ob}
\DeclareMathOperator{\ob}{ob}
\DeclareMathOperator{\Lib}{Lib}
\DeclareMathOperator{\esj}{esj}
\DeclareMathOperator{\id}{id}
\DeclareMathOperator{\Id}{Id}
\DeclareMathOperator{\ad}{ad}
\DeclareMathOperator{\Iim}{Im}
\DeclareMathOperator{\pr}{pr}
\DeclareMathOperator{\Cat}{Cat}
\DeclareMathOperator{\car}{car}
\DeclareMathOperator{\SysInd}{SysInd}
\DeclareMathOperator{\Sup}{Sup}
\DeclareMathOperator{\Coker}{Coker}
\DeclareMathOperator{\Filt}{Filt}
\DeclareMathOperator{\acc}{acc}
\DeclareMathOperator{\sur}{sur}
\DeclareMathOperator{\cat}{cat}
\DeclareMathOperator{\Ord}{Ord}
\DeclareMathOperator{\Gr}{Gr}
\DeclareMathOperator{\Cov}{Cov}
\DeclareMathOperator{\Caf}{Caf}
\DeclareMathOperator{\coker}{coker}
\DeclareMathOperator{\coim}{coim}
\DeclareMathOperator{\im}{im}
\DeclareMathOperator{\Ouv}{Ouv}
\DeclareMathOperator{\diag}{diag}
\DeclareMathOperator{\Esp}{Esp}
\DeclareMathOperator{\TOP}{TOP}
\DeclareMathOperator{\Unif}{Unif}
\DeclareMathOperator{\Iind}{\mathbf{Ind}}
\DeclareMathOperator{\Hhomcart}{\mathbf{Homcart}}
\DeclareMathOperator{\Hhomtopess}{\mathbf{Homtopess}}
\DeclareMathOperator{\Ppoints}{\mathbf{Points}}
\DeclareMathOperator{\Ppoint}{\mathbf{Point}}
\DeclareMathOperator{\Ttors}{\mathbf{Tors}}
\DeclareMathOperator{\Ssex}{\mathbf{Sex}}
\DeclareMathOperator{\Ffib}{\mathbf{Fib}}
\DeclareMathOperator{\Vvois}{\mathbf{Vois}}
\DeclareMathOperator{\Ppointess}{\mathbf{Pointess}}
\DeclareMathOperator{\scrHom}{\mathbf{Hom}}
\DeclareMathOperator{\sheafPic}{\mathbf{Pic}}
\DeclareMathOperator{\sheafAut}{\mathbf{Aut}}
\DeclareMathOperator{\Hhomtoplocan}{\mathbf{Homtoplocan}}
\DeclareMathOperator{\Hhomtopan}{\mathbf{Homtopan}}
\DeclareMathOperator{\Pptgeom}{\mathbf{Ptgéom}}
\DeclareMathOperator{\Morsite}{Morsite}
\DeclareMathOperator{\Fib}{Fib}
\DeclareMathOperator{\Fais}{Fais}
\DeclareMathOperator{\kar}{kar}
\DeclareMathOperator{\Groupes}{Groupes}
\DeclareMathOperator{\Restr}{Restr}
\DeclareMathOperator{\Isom}{Isom}
\DeclareMathOperator{\Prol}{Prol}
\DeclareMathOperator{\Res}{Res}
\DeclareMathOperator{\Tors}{Tors}
\DeclareMathOperator{\esp}{esp}
\DeclareMathOperator{\es}{es}
\DeclareMathOperator{\od}{od}
\DeclareMathOperator{\gr}{gr}
\DeclareMathOperator{\sob}{sob}
\DeclareMathOperator{\dfn}{dfn}
\DeclareMathOperator{\Sch}{Sch}
\DeclareMathOperator{\ouv}{ouv}
\DeclareMathOperator{\cons}{cons}
\DeclareMathOperator{\Nerf}{Nerf}
\DeclareMathOperator{\Points}{Points}
\DeclareMathOperator{\Espsob}{Espsob}
\DeclareMathOperator{\Homtop}{Homtop}
\DeclareMathOperator{\Pt}{Pt}
\DeclareMathOperator{\Quot}{Quot}
\DeclareMathOperator{\glob}{glob}
\DeclareMathOperator{\Homfib}{Homfib}
\DeclareMathOperator{\Qouv}{Qouv}
\DeclareMathOperator{\Quotouv}{Quotouv}
\DeclareMathOperator{\ens}{ens}
\DeclareMathOperator{\ab}{ab}
\DeclareMathOperator{\an}{an}
\DeclareMathOperator{\fs}{fs}
\DeclareMathOperator{\red}{red}
\DeclareMathOperator{\Ab}{Ab}
\DeclareMathOperator{\cosq}{cosq}
\DeclareMathOperator{\Ann}{Ann}
\DeclareMathOperator{\Zar}{Zar}


%%%%%%%%%%%%
%%%%% SGA4-2
\DeclareMathOperator{\Sex}{Sex}
\DeclareMathOperator{\vois}{vois}
\DeclareMathOperator{\Vois}{Vois}
\DeclareMathOperator{\Ext}{Ext}
\DeclareMathOperator{\ext}{ext}
\DeclareMathOperator{\Fer}{Fer}
\DeclareMathOperator{\Add}{Add}
\DeclareMathOperator{\Cocart}{Cocart}
\DeclareMathOperator{\cosk}{cosk}
\DeclareMathOperator{\plat}{plat}

\DeclareMathOperator{\faisc}{faisc}
\DeclareMathOperator{\sk}{sk}
\DeclareMathOperator{\ch}{ch}

\def\HH{\mathbf{H}}
\def\qc{\mathrm{qc}}
\def\et{\mathrm{et}}
\def\coh{\mathrm{coh}}

\def\Hvr#1{%
\begin{array}[b]{l}
{\kern2pt\scriptstyle r}\\[-8pt]
{\kern1pt\scriptstyle \vee}\\[-5pt]
H^{#1}
\end{array}}

\def\Hvv#1{%
\begin{array}[b]{l}
\kern2pt\rotatebox{90}{${\scriptscriptstyle\ll}$}\\[-6pt]
H^{#1}
\end{array}}

\def\Hvi#1{%
\begin{array}[b]{l}
{\kern0.2pt\scriptscriptstyle \infty}\\[-8pt]
{\kern1pt\scriptstyle \vee}\\[-5pt]
H^{#1}
\end{array}}

\def\Hv#1{%
\begin{array}[b]{l}
{\kern1pt\scriptstyle \vee}\\[-5pt]
H^{#1}
\end{array}}

\DeclareMathOperator{\cart}{cart}
\DeclareMathOperator{\Ffaisc}{\mathbf{Faisc}}


\def\Ev#1{%
\begin{array}[b]{l}
{\kern1pt\scriptstyle \vee}\\[-5pt]
E'_{#1}
\end{array}}

\DeclareMathOperator{\Cont}{Cont}
\DeclareMathOperator{\Ccont}{\mathscr{C}\!\mathit{ont}}
\DeclareMathOperator{\Site}{Site}
\newcommand\MM{\mathbf{M}}
\def\limtop#1{{\displaystyle{\varprojlim_{#1}\!\mbox{\rm top }}}}
\def\Limtop{\mathop{\varprojlim\!\mbox{\rm top}}}
\DeclareMathOperator{\spec}{spec}
\def\ss{\scriptstyle{\mathscr{S}}}
\DeclareMathOperator{\Gal}{Gal}
\def\smallm{\scriptstyle{\mathscr{M}}}
\DeclareMathOperator{\End}{End}
\DeclareMathOperator{\ass}{ass}
\DeclareMathOperator{\zar}{zar}
\DeclareMathOperator{\Ogg}{Ogg}
\DeclareMathOperator{\res}{res}
\DeclareMathOperator{\tr}{tr}
\DeclareMathOperator{\Supp}{Supp}
\DeclareMathOperator{\codim}{codim}
\DeclareMathOperator{\rk}{rk}
\DeclareMathOperator{\cur}{cur}
\DeclareMathOperator{\cdqc}{cdqc}
\newcommand\DD{\mathbf{D}}
\DeclareMathOperator{\cl}{cl}
\def\JJ{\mathbf{J}}
\def\FF{\mathbf{F}}
\def\BB{\mathbf{B}}
\DeclareMathOperator{\Alget}{Alget}
\DeclareMathOperator{\Ass}{Ass}
\DeclareMathOperator{\supp}{supp}
\DeclareMathOperator{\rad}{rad}
\DeclareMathOperator{\prof}{prof}
\DeclareMathOperator{\cint}{int}
\DeclareMathOperator{\sic}{sic}
\DeclareMathOperator{\Tor}{Tor}
\DeclareMathOperator{\tors}{tors}
\DeclareMathOperator{\Sym}{Sym}
\DeclareMathOperator{\parf}{parf}
\DeclareMathOperator{\tor}{tor}
\DeclareMathOperator{\defn}{defn}
\DeclareMathOperator{\Div}{Div}
\DeclareMathOperator{\Pol}{Pol}
\DeclareMathOperator{\Fex}{Fex}
\DeclareMathOperator{\Lie}{Lie}
\DeclareMathOperator{\tot}{tot}
\DeclareMathOperator{\Aut}{Aut}
\DeclareMathOperator{\HOM}{HOM}
\DeclareMathOperator{\pch}{pch}
\DeclareMathOperator{\Ch}{Ch}
\DeclareMathOperator{\EXT}{EXT}
\DeclareMathOperator{\PIC}{PIC}
\DeclareMathOperator{\Rf}{Rf}
\DeclareMathOperator{\inj}{inj}
\newcommand\VV{\mathbf{V}}

\def\Xv{%
\begin{array}[b]{l}
{\kern1pt\scriptstyle \vee}\\[-5pt]
X
\end{array}}

%%%%%%%
\newcommand\TT{\mathbf{T}}
\newcommand\XX{\mathbf{X}}
\newcommand\EE{\mathbf{E}}
%%%%%%

\def\qedsymbol{\textup{C.Q.F.D.}}
\let\CQFD\qed
\let\cqfd\qed

%%%%%\bigdot
\font\bigsymb=cmsy10 at 4pt
\def\bigdot{{\kern1.2pt\raise 1.5pt\hbox{\bigsymb\char15}}}
%%%%%

\def\Lright{\mathop{\vcenter{\hbox{\hbox{$\longrightarrow$}\lower8pt \hbox{\kern-16pt\hbox{$\longrightarrow$}}}}}}

\def\Aa{{\mathcal A}}
\def\Bb{{\mathcal B}}
\def\Cc{{\mathcal C}}
\def\Dd{{\mathcal D}}
\def\Ee{{\mathcal E}}
\def\Ff{{\mathcal F}}
\def\Gg{{\mathcal G}}
\def\Hh{{\mathcal H}}
\def\Ii{{\mathcal I}}
\def\Jj{{\mathcal J}}
\def\Kk{{\mathcal K}}
\def\Ll{{\mathcal L}}
\def\Mm{{\mathcal M}}
\def\Nn{{\mathcal N}}
\def\Oo{{\mathcal O}}
\def\Pp{{\mathcal P}}
\def\Qq{{\mathcal Q}}
\def\Rr{{\mathcal R}}
\def\Ss{{\mathcal S}}
\def\Tt{{\mathcal T}}
\def\Uu{{\mathcal U}}
\def\Vv{{\mathcal V}}
\def\Ww{{\mathcal W}}
\def\Xx{{\mathcal X}}
\def\Yy{{\mathcal Y}}
\def\Zz{{\mathcal Z}}

\newcommand{\CC}{\mathbf{C}}
\newcommand{\GG}{\mathbf{G}}
\newcommand{\KK}{\mathbf{K}}
\newcommand{\LL}{\mathbf{L}}
\newcommand{\NN}{\mathbf{N}}
\newcommand{\PP}{\mathbf{P}}
\newcommand\QQ{\mathbf{Q}}
\newcommand\RR{\mathbf{R}}
\newcommand\ZZ{\mathbf{Z}}

\def\to{\mathchoice{\longrightarrow}{\rightarrow}{\rightarrow}{\rightarrow}}
\def\from{\mathchoice{\longleftarrow}{\leftarrow}{\leftarrow}{\leftarrow}}
\def\mto{\mathchoice{\longmapsto}{\mapsto}{\mapsto}{\mapsto}}
\def\mfrom{\mathrel{\from\joinrel\mapstochar\,}}
\def\hto{\mathrel{\lhook\joinrel\to}}
\def\hfrom{\mathrel{\from\joinrel\rhook}}

\def\implies{\mathchoice{\Longrightarrow}{\Rightarrow}{\Rightarrow}{\Rightarrow}}
\def\implied{\mathchoice{\Longleftarrow}{\Leftarrow}{\Leftarrow}{\Leftarrow}}
\def\iff{\mathchoice{\Longleftrightarrow}{\Leftrightarrow}{\Leftrightarrow}{\Leftrightarrow}}

\def\To#1{\mathchoice{\xrightarrow{\labelstyle~~#1~}}{\stackrel{#1}{\longrightarrow}}{}{}}
\def\Mto#1{\mathrel{\mapstochar\To{#1}}}
\def\From#1{\mathchoice{\xleftarrow{\labelstyle~~#1~}}{\stackrel{#1}{\longleftarrow}}{}{}}

\newcommand{\isomto}{\overset{\sim}{\longrightarrow}}
\newcommand{\isomfrom}{\overset{\sim}{\longleftarrow}}

\newcommand{\IFF}{\unskip~$\iff$~\ignorespaces}
\newcommand{\THEN}{\unskip~$\implies$~\ignorespaces}

\let\leq\leqslant
\let\geq\geqslant
\let\le\leqslant
\let\ge\geqslant

\def\cf{{cf.}\kern.3em}
\def\Cf{{Cf.}\kern.3em}
\def\ie{{i.e.}\kern.3em}
\def\iev{{i.e.},\xspace}
%\def\loccito{loc.\kern2pt cit.\xspace}
%\def\loccitn{\textit{loc.\kern2pt cit.}\xspace}
%\def\loccit{\sisi{\loccito}{\loccitn}}
\def\EGA{\textup{EGA}\kern.2em}
\def\SGA{\textup{SGA}\kern.3em}

\let\oldS\S
\def\S{\oldS\kern .15em}
\def\ptbl{.\kern .2em}
\def\resp{resp.\kern.3em}
\def\Exp{Exp.\kern.3em}

\let\bar\overline
\let\hat\widehat
\let\tilde\widetilde
\let\dpl\displaystyle
\let\ldots\dots
\let\emptyset\varnothing
\let\skpt\skippointrait
\let\epsilon\varepsilon

\def\sfrac#1#2{{#1}/{#2}}

\DeclareMathOperator{\cd}{cd}
%\DeclareMathOperator{\Hhom}{\mathcal{H}\!\mathit{om}}
\DeclareMathOperator{\Hhom}{\mathbf{Hom}}
%\DeclareMathOperator{\Eext}{\mathcal{E}\!\mathit{xt}}
\DeclareMathOperator{\Eext}{\mathbf{Ext}}
\DeclareMathOperator{\Spec}{Spec}
\DeclareMathOperator{\Tr}{Tr}

\let\emptyset=\varnothing

%%%%%%%%%%%%%%%
%%% expose XVII
%%%%%%%%%%%%%%%

\makeatletter
\let\old@optexpose\optexpose
\newcommand{\empty@optexpose}[1]{}
\newcommand{\paren@optexpose}[1]{\old@optexpose{#1}(}
\newcommand{\refnoexpose}[1]{\let\optexpose\empty@optexpose\ref{#1}\let\optexpose\old@optexpose}%\refnoexpose{X.cor5.2} engendre "5.2" au lieu de "X 5.2"
\newcommand{\neweqref}[1]{\let\optexpose\paren@optexpose\textup{\ref{#1})}\let\optexpose\old@optexpose}%\neweqref{XII.eq4.1} engendre "XII (4.1)" au lieu de "(XII 4.1)"
\makeatother
\newcommand{\refvb}[2][VI${}_{\text{B}}$ ]{\sisi{#1\refnoexpose{#2}}{\ref{#2}}}

\newcommand{\leftexp}[2]{{\vphantom{#2}}^{#1}{#2}}
\newcommand{\AAA}{\mathbf{A}}
\newcommand{\Sf}{\mathfrak{S}}
\DeclareMathOperator*{\quotprojlim}{\text{``$\varprojlim$''}}
\DeclareMathOperator*{\quotinjlim}{\text{``$\varinjlim$''}}
\DeclareMathOperator{\Fl}{F\sisi{\ell}{l}}
\DeclareMathOperator{\Pl}{P\sisi{\ell}{l}}
%\DeclareMathOperator{\cl}{c\sisi{\ell}{l}}
\DeclareMathOperator{\tdf}{tdf}
\DeclareMathOperator{\rg}{rg}
\newcommand{\iInd}{\sisi{I}{i}nd}
\newcommand{\iind}{ind}
\newcommand{\ppro}{pro}
\newcommand{\tTor}{\sisi{T}{t}or}
\newcommand{\ttor}{tor}
\newcommand{\chap}{Chap.}
\newcommand{\Ffex}{\Ff\mathit{ex}}
\newcommand{\sheafk}{\underline{k}}
\newcommand{\sheafn}{\underline{n}}
\newcommand{\sheafGamma}{\underline{\Gamma}}

% Repeated references to an already typeset footnote must not use
% \footnotemark without a matching \footnotetext: Hyperref would then create
% a live annotation to a destination that cannot exist.  This helper keeps the
% current printed footnote mark and links it to an explicit stable note target.
\makeatletter
\newcommand{\sgarepeatedfootnotemark}[1]{%
  \begingroup
  \edef\@thefnmark{\thefootnote}%
  \hyperlink{#1}{\@makefnmark}%
  \endgroup}
\makeatother

\sisi{
\DeclareMathOperator{\simarrow}{\rule{.15mm}{1.8mm}\kern -3.4pt\rightsquigarrow}
\newcommand{\sheafell}{\underline{\ell}}
\newcommand{\sheafZZ}{\underline{\ZZ}}
}{
\newcommand{\nderef}[1]{la N.D.E.~(\ref{#1}) page~\pageref{#1}}
}

%%%%%%%%%%%%%%%%%
%%%%% expose XVII
%%%%%%%%%%%%%%%%%

% END INCLUDED SGA4_English_current_buildable_source_20260906/source/fichiers-configuration-latex/macros.tex


\externaldocument{../01/01_English.tex}
\externaldocument{../02/02_English.tex}
\externaldocument{../03/03_English.tex}
%\externaldocument{../04/04_English.tex}

\setboolean{original}{false}

\begin{document}

\begin{center}
IV. Topoi\\
A. Grothendieck and J.-L. Verdier
\end{center}

\selectlanguage{english}version: \showgitstatus
\tableofcontents

\else
\chapter{Topoi}
%\addtocontents{toc}{A. Grothendieck and J.-L. Verdier}
\begin{center}
A. Grothendieck and J.-L. Verdier
\end{center}
\fi
\phantomsection\MakeLinkTarget*{sga4.exp.IV}\label{sga4.exp.IV}

% BEGIN INCLUDED SGA4_English_current_buildable_source_20260906/source/04/drafts/SGA4_IV_LINES_31_187_English.texfrag
\setcounter{pageoriginal}{298}
\setcounter{section}{-1}
\section{Introduction}\etiquette[0]{IV.sec0}\pageoriginale

\subsection{}\etiquette[0.1]{IV.subsec0.1}
We saw in \hyperref[sga4.exp.II]{II} various exactness properties of categories of the
form $\widetilde{𝒞}=$ the category of sheaves of sets on $𝒞$, where
$𝒞$ is a small site. These properties can be expressed by saying
that in many respects such categories (which we shall call
\emph{topoi}) inherit the familiar properties of the category
$(\Ens)$ of (small) sets. On the other hand, experience has shown
that various situations in Mathematics should be regarded
\emph{above all as a technical means of constructing the
  corresponding categories of sheaves} (of sets), \ie \emph{the
  corresponding ``topoi.''} It appears that all the genuinely
important notions attached to a site (for example its cohomological
invariants, studied in \hyperref[sga4.exp.V]{V}; various other ``topological''
invariants, such as its homotopy invariants recently studied by M.
\textsc{Artin} and B. \textsc{Mazur} [1]; and the notions studied in
J. \textsc{Giraud}'s book on noncommutative cohomology) are in fact
expressed directly in terms of the associated topos. From this point
of view, two sites should be regarded as essentially equivalent when
their associated topoi are equivalent categories. The datum of a
site (at least in the case, especially important in practice, where
its topology is coarser than its canonical topology) should be
regarded as amounting to the datum of a topos $ℰ$ (namely the
associated topos formed by the sheaves of sets on the site) together
with a generating family of objects of $ℰ$ (\cf\pageoriginale II
\hyperref[II.subsec4.9]{4.9}, and \hyperref[IV.cor1.2.1]{1.2.1} below). This viewpoint is
analogous to associating a group with a system of generators and a
system of relations among those generators, and attaching interest
rather to the structure of that group than to the system of
generators and relations used to generate it (which are regarded as
auxiliary data of the situation). Moreover, the ``comparison lemma''
III \hyperref[III.theo4.1]{4.1} provides numerous examples of pairs of sites $𝒞$, $𝒞'$ that
are not isomorphic, and are not even equivalent as categories, but
give rise to equivalent topoi; thus $𝒞$ and $𝒞'$ should be regarded
as essentially equivalent.

\subsection{}\etiquette[0.2]{IV.subsec0.2}
In the present exposé we give a characterization of topoi by simple
exactness properties (due to J. \textsc{Giraud}); we study the
natural notion of a morphism of topoi, inspired by the notion of a
continuous map from one topological space to another; and we develop
in the setting of topoi certain constructions familiar from ordinary
sheaf theory (sheaves $\Hhom$, tensor-product sheaves, supports).
Finally, following M. \textsc{Artin}, we show how a topos can be
reconstructed from an ``open'' part, its complementary ``closed''
part, and a certain left exact functor connecting them, which may in
fact be chosen almost arbitrarily.

\subsection{}\etiquette[0.3]{IV.subsec0.3}
This gives a procedure for gluing topoi which, when applied to topoi
arising from ordinary topological spaces, will in general produce a
topos no longer of the same type. This is a first indication of the
remarkable stability of the notion of a topos under various natural
constructions, a stability lacking for the notion of a topological
space (from which the notion of a\pageoriginale topos is inspired).
As a second remarkable example, let us also mention the relative
classifying topos of a group in a topos (\cf [1] or
\hyperref[IV.exer5.9]{5.9} below), inspired by the classical notion
of the classifying space of a topological group, and the notion of a
``moduli topos'' associated with various ``moduli problems'' in
Algebraic Geometry or Analytic Geometry [10] [13].

Other topoi, such as the \emph{étale topos} of a scheme
(systematically studied in the present Seminar beginning with \Exp
VII) or the \emph{crystalline topos} of a relative scheme [6], arise
naturally when one wishes to develop usable cohomology theories for
abstract algebraic varieties (and, more generally, schemes),
replacing the classical Betti cohomology of algebraic varieties over
the field of complex numbers.

\subsection{}\etiquette[0.4]{IV.subsec0.4}
Thus one may say that the notion of a topos, a natural outgrowth of
the \emph{sheaf-theoretic viewpoint} in Topology, in turn constitutes
a substantial enlargement of the notion of a topological
space\footnote[1]{\Cf [9], or \hyperref[IV.subsec4.1]{4.1} and
\hyperref[IV.subsec4.2]{4.2} below, for the precise relations between
  the notion of a topos and that of a topological space.}, embracing
a large number of situations which formerly were not regarded as
belonging to topological intuition. The characteristic feature of
such situations is the availability of a notion of ``localization,''
a notion formalized precisely by that of a site and, ultimately, by
that of a topos (through the topos associated with the site). As the
term ``topos'' itself is intended precisely to suggest, it seems
reasonable and legitimate to the authors of the present Seminar to
regard the subject of Topology as the study of \emph{topoi} (and not
merely of topological spaces).

\subsection{}\etiquette[0.5]{IV.subsec0.5}
It seemed useful to include in this general exposé on topoi a fairly
large number of examples, many of which are only remotely\pageoriginale related to
the original aim of this seminar (namely, the study of étale
cohomology). The reader in a hurry, interested exclusively in étale
cohomology, may of course omit these examples without inconvenience,
as well as most of the present exposé, consulting it only when the
need arises.

\section{Definition and characterization of topoi}\etiquette[1]{IV.sec1}

\begin{enonce*}{Definition 1.1}\etiquette[1.1]{IV.def1.1}
A category $E$ is called a $𝒰$-topos, or simply a topos when no
confusion is likely, if there exists a site $C\in 𝒰$ such that $E$ is
equivalent to the category $C^{\widetilde{\phantom{,}}}$ of
$𝒰$-sheaves of sets on $C$.
\end{enonce*}

\setcounter{subsection}{1}
\subsubsection{}\etiquette[1.1.1]{IV.subsubsec1.1.1}
Let $E$ be a $𝒰$-topos. We always regard $E$ as endowed with its
canonical topology (\hyperrefext[II][II.def2.5]{2.5}), which therefore
makes it a \emph{site}, and even, by \hyperref[IV.subsubsec1.1.2]{1.1.2}
d) below, a $𝒰$-site (\hyperrefext[II][II.def3.0.2]{3.0.2}). Unless
expressly stated otherwise, we shall consider no topology on $E$
other than the one just specified.

\subsubsection{}\etiquette[1.1.2]{IV.subsubsec1.1.2}
We saw in \hyperrefext[II][II.prop4.8]{4.8},
\hyperrefext[II][II.cor4.11]{4.11} that a $𝒰$-topos $E$ is a
$𝒰$-category (\hyperrefext[I][I.def1.1]{1.1}) satisfying the following
conditions:
\begin{enumerate}
\renewcommand{\theenumi}{\alph{enumi}}
\renewcommand{\labelenumi}{\theenumi)}
\item Finite projective limits are representable in $E$.

\item Direct sums indexed by an element of $𝒰$ are representable in
  $E$. They are disjoint and universal
  (\hyperrefext[II][II.def4.5]{4.5}).

\item Equivalence relations in $E$ are universally effective
  \sisi{(\hyperrefext[I][I.sec10.9]{10.9}).}{(\hyperrefext[I][I.sec10.10]{10.10}).}

\item $E$ admits a generating family
  (\hyperrefext[II][II.subsec4.9]{4.9}) indexed by an element of $𝒰$.
\end{enumerate}

In\pageoriginale fact, we shall see that these intrinsic properties
\emph{characterize} $𝒰$-topoi:

% END INCLUDED SGA4_English_current_buildable_source_20260906/source/04/drafts/SGA4_IV_LINES_31_187_English.texfrag


% BEGIN INCLUDED SGA4_English_current_buildable_source_20260906/source/04/drafts/SGA4_IV_LINES_188_382_English.texfrag
\begin{enonce*}{Theorem 1.2}\etiquette[1.2]{IV.theo1.2}
{\rm (J. Giraud).} Let $E$ be a $𝒰$-category. The following
properties are equivalent:
\begin{itemize}
\item[i)] $E$ is a $𝒰$-topos $(\hyperref[IV.def1.1]{1.1})$.

\item[ii)] $E$ satisfies conditions {\rm a), b), c),} and {\rm d)
  of \hyperref[IV.subsubsec1.1.2]{1.1.2}.}

\item[iii)] The $𝒰$-sheaves on $E$ for the canonical topology
  are representable, and $E$ has a small generating family
  (condition {\rm \hyperref[IV.subsubsec1.1.2]{1.1.2} d)}).

\item[iv)] There exist a category $C\in 𝒰$ and a fully faithful
  functor $i:E\rightarrow \widehat{C}$ (where $\widehat{C}$ denotes
  the category of $𝒰$-presheaves on $C$) having a left adjoint
  a which is left exact.

\item[${\rm i')}$] There exists a site $C\in 𝒰$ such that
  projective limits are representable in $C$ and the topology of
  $C$ is coarser than its canonical topology {\rm (II \hyperref[II.def2.5]{2.5}),} %% label!
  and such that $E$ is equivalent to the category
  $C^{\widetilde{\phantom{,}}}$ of $𝒰$-sheaves of sets on
  $C$.
\end{itemize}
\end{enonce*}

Moreover:

\begin{enonce*}{Corollary 1.2.1}\etiquette[1.2.1]{IV.cor1.2.1}
Let $E$ be a $𝒰$-topos, and let $C$ be a full subcategory of
$E$, endowed with the topology induced {\rm (\hyperrefext[III][III.subsec3.1]{3.1})}
by that of $E$ (\hyperref[IV.subsubsec1.1.1]{1.1.1}). Consider the
functor
$$
E\longrightarrow C^{\widetilde{\phantom{,}}}
$$
which associates to every $X\in \ob E$ the restriction to $C$ of
the functor represented by $X$. This functor is an equivalence of
categories if and only if $\ob C$ is a generating family of $E$.
\end{enonce*}

The equivalence\pageoriginale i) $\Leftrightarrow$ iv) follows
immediately from \hyperrefext[II][II.theo5.5]{5.5}, and we have already recalled
above that i) $\Rightarrow$ ii). Since ${\rm i'}$) $\Rightarrow$ i)
trivially, it remains to prove ii) $\Rightarrow$ iii) and iii)
$\Rightarrow$ ${\rm i'}$), which will be done in
\hyperref[IV.subsubsec1.2.4]{1.2.4} and \hyperref[IV.subsubsec1.2.3]{1.2.3} below. The
corollary is then obtained by observing that the functor in question
factors as $E\rightarrow E^{\widetilde{\phantom{,}}}\rightarrow
C^{\widetilde{\phantom{,}}}$; since $E\rightarrow
E^{\widetilde{\phantom{,}}}$ is an equivalence by virtue of
(iii), the question reduces to determining when
$E^{\widetilde{\phantom{,}}}\rightarrow C^{\widetilde{\phantom{,}}}$
is an equivalence. The conclusion then follows from the ``comparison
lemma'' \hyperrefext[III][III.theo4.1]{4.1}; \cf the proof in
\hyperref[IV.subsubsec1.2.3]{1.2.3} below.

\setcounter{subsection}{2} \setcounter{subsubsection}{2}
\subsubsection{}\etiquette[1.2.3]{IV.subsubsec1.2.3}
\emph{Proof of} iii) $\Rightarrow$ ${\rm i'}$). Let
$\mathfrak{X}=(X_i)_{i\in I}$, $I\in 𝒰$, be a small generating
family of $E$. Since $E$ is a $𝒰$-category, the set of
isomorphism classes of finite diagrams in $E$ whose objects are
members $X_i$ is $𝒰$-small. Consequently, the smallest set
$\mathfrak{X}'$ of objects of $E$ containing the finite projective
limits of objects of $\mathfrak{X}$ is a countable union of small
sets and is therefore small~\footnote[1]{Here we reason with classes
of objects of $E$ up to isomorphism.}. Setting
$\mathfrak{X}^{(n+1)}=\sisi{\mathfrak{X}^{(n)(1)}}{(\mathfrak{X}^{(n)})'}$, and
$\mathfrak{X}=\bigcup_n \mathfrak{X}^{(n)}$, we see that, after
enlarging the generating family, we may suppose that
$\mathfrak{X}$ is stable under finite projective limits. Let
$𝒱$ be a universe containing $𝒰$ such that $E$ is
$𝒱$-small. For every object $H$ of $E$, denote by $I(H)$
the set $\amalg_{i\in I}\Hom(X_i, H)$. Let $H$ be an object of $E$
and let $H'\hookrightarrow H$ be the subsheaf of $H$, for the
canonical topology, which is the ``union'' of the images of the
morphisms $u:X_i\rightarrow H$, $(u,i)\in I(H)$
(\hyperrefext[II][II.theo4.1]{4.1}). Since $H'$ is a subsheaf of a
$𝒰$-sheaf, $H'$ is a $𝒰$-sheaf. It is therefore
representable. Since the family $\mathfrak{X}$ is generating and,
for every $i\in I$, the map $\Hom(X_i, H')\rightarrow \Hom(X_i,H)$
is bijective, the morphism $H'\hookrightarrow H$ is an isomorphism
(\hyperrefext[II][II.subsec4.9]{4.9}). Consequently
(\sisi{\hyperrefext[II][II.def6.1]{6.1}}{\hyperrefext[II][II.def5.2]{5.2}}), the family
$(u:X_i\rightarrow H)$, $(u,i)\in I(H)$, is covering for the
canonical topology of $E$. Thus\pageoriginale every object of $E$
can be covered, for the canonical topology of $E$, by objects of
$\mathfrak{X}$. Let $C$ be the full subcategory of $E$ defined by
the objects of $\mathfrak{X}$, and let $u:C\hookrightarrow E$ be the
inclusion functor. The topology $\Cc$ induced on $C$ by the canonical
topology of $E$ is coarser than the canonical topology of $C$, and
when $𝒰$ contains an element of infinite cardinality, finite
projective limits are representable in $C$. It follows from
\hyperrefext[III][III.theo4.1]{4.1} that the functor $F\mapsto F\circ
u$ is an equivalence from $E$ to the category of $𝒰$-sheaves
on $C$ for the topology $\Cc$, \cqfd

\subsubsection{}\etiquette[1.2.4]{IV.subsubsec1.2.4}
\emph{Proof of} ii) $\Rightarrow$ iii). The proof consists of four
steps.

Let $𝒱$ be a universe containing $𝒰$, such that $E$ is
an element of $𝒱$. Let $\widetilde{E}$ be the category of
$𝒱$-sheaves on $E$ for the canonical topology, and let
$J_E:E\rightarrow \widetilde{E}$ be the canonical functor.

\paragraph{}\etiquette[1.2.4.1]{IV.par1.2.4.1}
Let $(g_i:G_i\rightarrow H)$, $i\in I\in 𝒰$, be an
epimorphic family in $\widetilde{E}$. If the $G_i$ and the fiber
products $\fprod{G_i}{G_j}{H}$ are representable, then the sheaf
$H$ is representable.

Indeed, the direct sum $\coprod_{i\in I}G_i$ is represented in
$\widetilde{E}$ (\hyperrefext[II][II.theo4.1]{4.1}) by a representable sheaf
$G$ (property b) and \hyperrefext[II][II.prop4.6]{4.6}). The same holds for the
direct sum $K=\coprod_{(i,j)\in I\times
I}\fprod{G_i}{G_j}{H}$. Moreover, the diagram $K\rightrightarrows
G\rightarrow H$ is exact and $K$ is the fiber square of $G$ over
$H$. (This latter property holds in the category of presheaves, and
hence in the category of sheaves (\hyperrefext[II][II.theo4.1]{4.1}).) It follows
from c) and \hyperrefext[II][II.prop4.7]{4.7} that the sheaf $H$ is
representable.

\paragraph{}\etiquette[1.2.4.2]{IV.par1.2.4.2}
Let $X_{\alpha}$, $\alpha\in A\in 𝒰$, be a generating
family of $E$. For every sheaf $H$, denote by $I(H)$ the set
$\coprod_{\alpha \in A}\Hom_E (X_{\alpha}, H)$. The family
$(u:X_{\alpha}\rightarrow H, (u,\alpha)\in
I(H))$\pageoriginale is epimorphic in $\widetilde{E}$. When $H$ is
a $𝒰$-sheaf, $I(H)$ is an element of $𝒰$.

Indeed, since every sheaf $H$ is the target of an epimorphic family
of morphisms whose sources are representable sheaves (
\hyperrefext[I][I.prop3.4]{3.4} and \hyperrefext[II][II.theo4.1]{4.1}), it is enough to prove the
first assertion when the sheaf $H$ is representable. Let $G$ then
be the image of the family $(u:X_{\alpha}\rightarrow H,
(u,\alpha)\in I(H))$. The morphism $G\rightarrow H$ is a
monomorphism. We are therefore in the situation of the first step,
since $\fprod{X_{\alpha}}{X_{\beta}}{G}$ is isomorphic to
$\fprod{X_{\alpha}}{X_{\beta}}{H}$, which is representable, and
moreover $I(H)$ is an element of $𝒰$. It follows that $G$ is
representable. But since $X_{\alpha}$ is a generating family,
$G\rightarrow H$ is an isomorphism. The last assertion is immediate
from the definition of $𝒰$-sheaves.

\paragraph{}\etiquette[1.2.4.3]{IV.par1.2.4.3}
Every subsheaf of a representable sheaf is representable.

Indeed, let $G\rightarrow H$ be a subsheaf of a representable
sheaf. Then $G$ is a $𝒰$-sheaf. The family
$(u:X_{\alpha}\rightarrow G, (u,\alpha)\in I(G))$ is therefore
epimorphic in $\widetilde{E}$ and indexed by an element of
$𝒰$. Moreover, the fiber products
$\fprod{X_{\alpha}}{X_{\beta}}{G}$ are isomorphic to the fiber
products $\fprod{X_{\alpha}}{X_{\beta}}{H}$, which are
representable. Thus we are in the situation of the first step and
$G$ is representable.

\paragraph{}\etiquette[1.2.4.4]{IV.par1.2.4.4}
Every $𝒰$-sheaf is representable.

Indeed, by the first and second steps, it is enough to show that the
fiber products $\fprod{X_{\alpha}}{X_{\beta}}{H}$ are representable.
But these fiber products are subobjects of the products
$\fprod{X_{\alpha}}{X_{\beta}}{}$. The conclusion therefore follows
from the third step. This completes the proof of Theorem
\hyperref[IV.theo1.2]{1.2}.

\begin{enonce*}[remark]{Remark 1.3}\etiquette[1.3]{IV.rem1.3}
Of\pageoriginale course, for a given $𝒰$-topos $E$, there is in
general no preferred way of representing it, up to equivalence, in
the form $C^{\widetilde{\phantom{,}}}$, with $C$ a \emph{small}
site; or, what amounts essentially to the same thing by virtue of
\hyperref[IV.cor1.2.1]{1.2.1} when one restricts to $C$ whose topology is
coarser than the canonical topology, there is no preferred
\emph{small} generating family in $E$. When one ceases to impose on
$C$ the condition that $C$ be small, there is, on the other hand (by virtue of
\hyperref[IV.theo1.2]{1.2} iii)), an entirely canonical choice of a
$𝒰$-site $C$ such that $E$ is equivalent to
$C^{\widetilde{\phantom{,}}}$, namely $E$ itself! This is one of the
technical reasons why in practice it is inconvenient to work only
with small sites: in fact, the most important sites of all, namely
$𝒰$-topoi, are not small! Moreover, the sites generating topoi
that arise in many questions in algebraic geometry (and even in
topology, \cf \hyperref[IV.subsec2.5]{2.5}) are not small either;
examples include the étale site of a scheme \sisi{(VIII.1).}{(VII.1).}
\end{enonce*}

% END INCLUDED SGA4_English_current_buildable_source_20260906/source/04/drafts/SGA4_IV_LINES_188_382_English.texfrag


% BEGIN INCLUDED SGA4_English_current_buildable_source_20260906/source/04/drafts/SGA4_IV_LINES_383_536_English.texfrag
\begin{enonce*}{Proposition 1.4}\etiquette[1.4]{IV.prop1.4}
Let $E$ be a $𝒰$-topos, let $E'$ be a category (not
necessarily a $𝒰$-category), and let $F$ be a presheaf
on $E$ with values in $E'$. In order that $F$ be a sheaf with
values in $E'$ {\rm(\hyperrefext[II][II.def6.1]{6.1}),} it is necessary and sufficient that
$F$ take $𝒰$-inductive limits in $E$ to projective
limits in $E'$.
\end{enonce*}

Let $𝒱$ be a universe containing $𝒰$ and such that $E'$ is
a $𝒱$-category. By composing $F$ with the functors
$\Hom(X',-):E'\rightarrow 𝒱\text{-}\Ens$ defined by the
objects $X'$ of $E'$, we are reduced to the case where
$E'=𝒱\text{-}\Ens$, with $𝒱$ a universe.
Suppose that $F$ takes $𝒰$-inductive limits to projective
limits. It then follows at once from the definitions that $F$ is
a sheaf, since a covering family\pageoriginale
$X_i\rightarrow X$ in $ℰ$, by the definition of the canonical
topology of $ℰ$, allows one to regard $X$ as an inductive limit
of the diagram $ \left(
\fprod{X_i}{X_j}{X}\raise6pt\hbox{$\mathop{\nearrow}^{\mbox{$X_i$}}_{\displaystyle{\kern-10pt\raise-8pt\hbox{$\searrow_{\mbox{$X_j$}}$}}}$}
\right) $. Suppose that $F$ is a sheaf, and let us prove that it
takes $𝒰$-inductive limits to projective limits. If
$𝒱\subset 𝒰$, we may suppose that $𝒱=𝒰$,
and it suffices to apply criterion \hyperref[IV.theo1.2]{1.2} iii). To
treat the general case, a little more work is needed. Let
$C$ be a full subcategory of $E$ generated by a generating family
indexed by an element of $𝒰$, and let
$u:C\rightarrow E$ be the inclusion functor. Endow $C$ with the
topology induced by the canonical topology of $E$ \sisi{(III 4.5).}{(III \hyperref[III.prop3.5]{3.5}).}
Denote by $C^{\widetilde{\phantom{,}}}_{𝒰}$ and
$C^{\widetilde{\phantom{,}}}_{𝒱}$ the categories of
sheaves on $C$, by $E^{\widetilde{\phantom{,}}}_{𝒱}$ the
category of $V$-sheaves on $E$ for the canonical topology,
by $J_E:E\rightarrow E_{𝒱}$ the canonical functor, and
by $i_{𝒰,𝒱}:C^{\widetilde{\phantom{,}}}_{𝒰}\rightarrow
C^{\widetilde{\phantom{,}}}_{𝒱}$ the inclusion functor. The
diagram
\[\vcenter{
\xymatrix@C=3.25cm@R=1.25cm{
C^{\widetilde{\phantom{a}}}_{\underline{U}} \ar[r]^{i_{\underline{U},\underline{V}}} & C^{\widetilde{\phantom{a}}}_{\underline{V}}\\
E  \ar[r]^{J_E} \ar[u]&
E^{\widetilde{\phantom{a}}}_{\underline{V}}\ar[u]\quad,
}}\leqno{(*)}
\]
where the vertical arrows are induced by the functor
$F\mapsto F\circ u$, is commutative. It follows from the explicit
construction of the associated sheaf functor (\hyperrefext[II][II.sec3]{3}) and from
\hyperrefext[II][II.theo4.1]{4.1} that the functor $i_{𝒰,𝒱}$ commutes
with $𝒰$-inductive limits. Moreover, it follows from
\hyperrefext[III][III.subsec5.1]{5.1} and \hyperref[IV.cor1.5]{1.5} that the vertical arrows
in $(\ast)$ are equivalences of categories. Consequently,
$J_E:E\rightarrow E^{\widetilde{\phantom{,}}}_{𝒱}$
commutes with $𝒰$-inductive limits. For every object $X$ of
$E$, we have
$$
F(X)\simeq
\Hom_{E^{\widetilde{\phantom{,}}}_{𝒱}}(J_E(X),F){\quoi}.
$$
Thus\pageoriginale $F$ takes $𝒰$-inductive limits
in $E$ to projective limits.

\begin{enonce*}{Corollary 1.5}\etiquette[1.5]{IV.cor1.5}
Let $E$ be a $𝒰$-topos, let $E'$ be a $𝒰$-category, and let
$f:E\rightarrow E'$ be a functor. In order that $f$ admit a right
adjoint, it is necessary and sufficient that $f$ commute with
$𝒰$-inductive limits.
\end{enonce*}

This is an immediate consequence of \hyperref[IV.prop1.4]{1.4} in the case
$E'=𝒰\text{-}\Ens$.

\begin{enonce*}{Corollary 1.6}\etiquette[1.6]{IV.cor1.6}
Let $E$ and $E'$ be two $𝒰$-topoi, and let $f:E\rightarrow E'$ be a
functor. The following conditions are equivalent:
\begin{enumerate}
\renewcommand{\theenumi}{\roman{enumi}}
\renewcommand{\labelenumi}{\theenumi)}
\item $f$ commutes with $𝒰$-inductive limits.

\item $f$ admits a right adjoint.

\item $f$ is continuous {\rm(\hyperrefext[III][III.def1.1]{1.1}).}
\end{enumerate}
\end{enonce*}

The equivalence i) $\Leftrightarrow$ ii) was proved in
\hyperref[IV.cor1.5]{1.5}. To prove i) $\Leftrightarrow$ iii), apply the
definition of continuous functors, choosing a universe
$𝒱$ such that $𝒰\in 𝒱$ (so $E$ and $E'$ are
$𝒱$-small sites). We must express that for every
$𝒱$-sheaf $F$ on $E'$, the composite $F\circ f$ is a
sheaf on $E$, that is (\hyperref[IV.prop1.4]{1.4}), that it takes
$𝒰$-inductive limits to projective limits. For this it is
sufficient that i) hold, since by \hyperref[IV.prop1.4]{1.4} $F$
itself takes $𝒰$-inductive limits to projective
limits; it is also necessary, as one sees by taking $F$ to be
a representable functor.

\begin{enonce*}{Corollary 1.7}\etiquette[1.7]{IV.cor1.7}
With the notation of \hyperref[IV.cor1.6]{1.6}, in order that $f$ be the
inverse-image functor $u^{\ast}$ for a morphism of topoi
$(\hyperref[IV.def3.1]{3.1})$ $u:E'\rightarrow E$, it is necessary and sufficient that
$f$ be left exact and take epimorphic families to epimorphic
families.
\end{enonce*}

The\pageoriginale necessity is trivial by definition (NB every
right exact functor takes epimorphisms to epimorphisms). The
sufficiency follows from \sisi{\hyperrefext[III][III.prop2.6]{2.6}}{\hyperrefext[III][III.prop1.6]{1.6}}, which implies
that $f$ is continuous under the stated conditions, and hence
commutes with $𝒰$-inductive limits by
\hyperref[IV.cor1.6]{1.6}.

\begin{enonce*}[remark]{Remark 1.8}\etiquette[1.8]{IV.rem1.8}
With the notation of \hyperref[IV.cor1.5]{1.5}, one can show that $f$ admits
a left adjoint if and only if it commutes with
$𝒰$-projective limits. In other words, a
\emph{covariant} functor $E\rightarrow (𝒰\text{-}\Ens)$ is
representable if and only if it commutes with $𝒰$-projective
limits \footnote[1]{A more general statement is found in
\hyperrefext[I][I.cor8.12.8]{8.12.8}, \hyperrefext[I][I.exer8.12.9]{8.12.9}; the outline of the proof
following a) through c) corresponds to the proof given in
\loccit}. We indicate only the principle of the proof, which
proceeds in three steps:
\begin{enumerate}
\renewcommand{\theenumi}{\alph{enumi}}
\renewcommand{\labelenumi}{\theenumi)}
\item Standard arguments [5, ${\rm n^{\circ}}$ 195, \S 3]
show that if $F$ commutes with $\varprojlim$, it is
prorepresentable by a strict projective system $(T_i)_{i\in
I}$, where $I$ is a filtered ordered set, \emph{not
necessarily small.} We may suppose that if $i>j$, then
$T_i\rightarrow T_j$ is not an isomorphism, and under this
hypothesis $F$ is representable if and only if $I$ is
small (which in fact implies that the projective system is
essentially constant).

\item To prove that $I$ is small, knowing that for every object $X$
of $E$ the set $F(X)=\varinjlim_i \Hom(T_i,X)$ is small, it suffices
to have a \emph{small cogenerating family}
$(X_j)_{j\in J}$ (\ie one that is generating for the opposite
category $E^{\circ}$). Now one shows that every $𝒰$-topos
$E$ has a small cogenerating family.

\item To prove this last point, one observes by standard
arguments [Toh] that every object $X$ of $E$ admits a monomorphism
into an ``injective'' object; and then that for every generating
family $(L_{\alpha})$ of $E$, if\pageoriginale each $L_{\alpha}$ is
embedded in this way into an injective object $I_{\alpha}$, the
family $(I_{\alpha})$ is cogenerating.
\end{enumerate}
\end{enonce*}

% END INCLUDED SGA4_English_current_buildable_source_20260906/source/04/drafts/SGA4_IV_LINES_383_536_English.texfrag


% BEGIN INCLUDED SGA4_English_current_buildable_source_20260906/source/04/drafts/SGA4_IV_LINES_537_947_English.texfrag
\section{Examples of topoi}\etiquette[2]{IV.sec2}

\setcounter{subsection}{-1}
\subsection{}\etiquette[2.0]{IV.subsec2.0}
We have brought together in this section a fairly large number of
typical examples of topoi, which the reader will already have had
occasion to encounter elsewhere, and which are intended to facilitate
his access to the ``yoga'' of topoi. For other examples (drawn from
algebraic geometry) of topologies on sites, giving rise to as many
topoi, he may consult \SGA 3 IV 6, and (for the étale topos) Exposé VII
of the present seminar. Since later on we shall hardly refer to the
present section except for questions of notation or terminology, we
leave it to the reader, as an exercise, to verify the statements with
which we have supplemented these examples for his general instruction.
All the examples of the present section will be made more precise in
4, where their functorial dependence on the data will be examined, and
in the following sections by way of illustration of the general notions
concerning topoi.

\subsection{Topos associated with a topological space}\etiquette[2.1]{IV.subsec2.1}

Let $X$ be a small topological space, and let $\Ouv(X)$ be the category
of open subsets of $X$, endowed with the canonical topology
\sisi{(III 2.9.1).}{(III \hyperref[III.exem1.9.1]{1.9.1}).} We shall denote by $\Top(X)$ the
topos of $𝒰$-sheaves on $\Ouv(X)$. This topos is equivalent to the
category of étale topological spaces over $X$, by associating with every
such space $X'$ over $X$ the sheaf
$U\mapsto \Gamma(X'/U)=\Hom_X(U,X')$ on $\Ouv(X)$
[TF]\nde{Here [TF] denotes reference [8].} By an abuse\pageoriginale
of language, we shall sometimes be unable to refrain from denoting by
the same letter $X$ the topos $\Top(X)$ defined by the topological space
$X$.

It is of course the preceding example that served chiefly as a guide
and as an intuitive basis for the development of the theory of topoi.
One should note, however, that the topoi obtained from topological
spaces are of a very special nature, owing in particular to the fact
that they have been described by sites $C=\Ouv(X)$ in which
\emph{all morphisms are monomorphisms} (and hence whose underlying
category reduces to a preordered set). It follows in particular that
the sheaves represented by the objects of $C$ are subsheaves of the
terminal sheaf, and consequently that \emph{the subsheaves of the
terminal sheaf form a generating family of the topos under
consideration}. This property is not shared by most topoi that arise
naturally in algebraic geometry or in algebra, \cf the examples below.
It is, up to minor qualifications, characteristic of topoi of the form
$\Top(X)$ (\cf \hyperref[IV.exer7.1.9]{7.1.9} below).

It is easily verified that the map $\Ouv(X)\rightarrow \Top(X)$,
which associates with every open subset of $X$ the sheaf it represents,
is a \emph{bijection} from $\Ouv(X)$ onto the set of subobjects of the
terminal object of $\Top(X)$; this bijection is even an isomorphism for
the natural order structures, \ie it induces an isomorphism of the
corresponding categories. This suggests that it should be possible to
reconstruct the topological space $X$, up to homeomorphism, from
$\Top(X)$ known up to equivalence. We shall see below
(\hyperref[IV.subsec4.2]{4.2}) that this is indeed the case, subject to
a slight restriction on $X$.

\subsection{Point topos or terminal topos, and empty topos or initial topos}\etiquette[2.2]{IV.subsec2.2}\pageoriginale

When $X$ is a topological space consisting of a single point, the
functor
$$
F\longrightarrow F(X):\Top(X)\longrightarrow 𝒰\mbox{-}\Ens
$$
is an equivalence of categories. This shows in particular that the
\emph{category} $𝒰$-$\Ens$ is a $𝒰$-\emph{topos.} We saw by
examples in \Exp II that this $𝒰$-topos is typical from the point
of view of exactness properties, in that the verification of many
properties (especially exactness properties) of general topoi reduces
to this particular topos. The interpretation given here in terms of
the one-point topological space justifies the abuse of language by
which a topos equivalent to the category $𝒰$-$\Ens$ is called a
\emph{point topos} (although, as a category, it is not at all
equivalent to the one-point category!). This is the terminology that
corresponds to the correct geometric intuition of the role played by
these topoi. By another abuse of language, a point topos is sometimes
called a \emph{terminal topos}, \cf \hyperref[IV.subsec4.3]{4.3}; we
shall say ``\emph{the} terminal topos'' for the topos
($𝒰$-$\Ens$).

When $X$ is the empty topological space, so that $\Ouv(X)$ is the
one-point category, a presheaf $F$ on $\Ouv(X)$ is a sheaf if and only
if its value at the unique object $\phi$ of $\Ouv(X)$ is a singleton
set. It follows that $\Top(∅)$ is isomorphic to the category of
singleton $𝒰$-sets, a category equivalent to the one-point category.
From this one concludes in particular that the one-point category (as
well as any $𝒰$-category equivalent to it) is a $𝒰$-topos. By an
abuse\pageoriginale of language, it is sometimes called the
\emph{empty topos} or the \emph{initial topos} (\cf
\hyperref[IV.subsec4.4]{4.4}); one should take care that it is not
equivalent to the empty category.

\subsection{Topos associated with a space with operators}\etiquette[2.3]{IV.subsec2.3}

Let $X$ be a topological space and $G$ a discrete group acting on $X$
by homeomorphisms. In \sisi{[Toh. \hyperref[IV.subsec5.1]{5.1}]}{[Toh]\nde{Here
[Toh] denotes reference [4].} \hyperref[IV.subsec5.1]{5.1}} the category
of $G$-\emph{sheaves on} $X$, or, as we shall also say, sheaves on
$(X,G)$, was defined; these are sheaves (of sets) on $X$
\emph{equipped} with actions of $G$ compatible with those of $G$ on
$X$. One sees at once, using Giraud's criterion
\hyperref[IV.theo1.2]{1.2} iii), that this %% \XXX ii) est probablement plus facile à vérifier
category is a topos (NB $𝒰$ is understood throughout), which we shall
simply denote by $\Top(X,G)$. When $G$ is the trivial group, we recover
the example of \hyperref[IV.subsec2.1]{2.1}; when $X$ is the one-point
space, we obtain the topos of sets on which $G$ acts on the left (or
$G$-\emph{sets}), also called the \emph{classifying topos of the
discrete group} $G$, and denoted by $B_G$. It is easily verified that
the only subobjects of the terminal object $e$ of the topos $B_G$ are
$e$ and the empty sheaf $\phi$; in particular, if $G$ is not the
trivial group, the subobjects of the terminal object of the classifying
topos $B_G$ do not form a generating family of $B_G$. Thus $B_G$ is
then not equivalent to a topos of the type $\Top(X)$ considered in
\hyperref[IV.subsec2.1]{2.1}.

The notion of a $G$-sheaf was introduced in \loccit in order to develop
the cohomological theory of abelian $G$-sheaves. Interpreting the latter
as the abelian sheaves of the topos $(X,G)$, that theory becomes part
of the theory of \Exp V, developed in the setting of general topoi.

More\pageoriginale generally, one may seek to associate an appropriate
topos with a topological space $X$ equipped with a topological group
$G$ (not necessarily discrete) of automorphisms, which would give rise
to an adequate cohomological theory. The same may be done in the
context of differentiable manifolds, real or complex analytic
manifolds, or schemes. This is indeed possible, \cf
\hyperref[IV.subsec2.5]{2.5} below.

\subsection{Classifying topos of a group object}\etiquette[2.4]{IV.subsec2.4}

Let $E$ be a topos, and let $G$ be a group object of $E$. Let $(E,G)$
be the category of objects of $E$ on which $G$ acts. It follows at
once from Giraud's criterion that this is a topos. It is called the
\emph{classifying topos} of the group object $G$, and is denoted by
$B_G$. When $E$ is the point topos (\hyperref[IV.subsec2.2]{2.2}),
\ie when $G$ is an ordinary group, we recover the classifying topos of
\hyperref[IV.subsec2.3]{2.3}.

The terminology adopted here is justified by the fact that the topos
$B_G$ plays a universal role in the classification of ``torsors'' (or
homogeneous principal bundles) under $G$, or more generally under
$G_{E'}=f^{\ast}(G)$, where $E'$ is a topos ``over $E$'', \ie endowed
with a morphism $f:E'\rightarrow E$ (\cf
\hyperref[IV.def3.1]{3.1} below). This role, made explicit in [3 Chap V]
or in \hyperref[IV.exer5.9]{5.9} below, shows that $B_G$ plays, in the
context of topoi, the same role as the classical classifying spaces of
topological groups in the homotopy theory of topological spaces. The
latter may be regarded (\cf \hyperref[IV.subsec2.5]{2.5}) as a weakened
version of the former, obtained by retaining from the classifying topos
only the ``homotopy type'' of that topos, in a suitable sense which
need not be specified here.

\subsection{``Big site'' and ``big topos'' of a topological space.
Classifying topos of a topological group}\footnote[1]{The introduction
of these sites and topoi is due to M. \textsc{Giraud}, who also brought
out their advantages over the traditional ``small'' site.}\etiquette[2.5]{IV.subsec2.5}\pageoriginale

Let $U-\Esp$, or simply $(\Esp)$, be the category of topological
spaces $\in 𝒰$. It is known that finite projective limits are
representable in $(\Esp)$. Consider on $(\Esp)$ the pretopology
(\hyperrefext[II][II.def1.3]{1.3}) for
which $\Cov(X)$ is the set of surjective families of open embeddings
$u_i:X_i\rightarrow X$. We regard $(\Esp)$ as a site by means of the
topology generated by the preceding pretopology. For every object $X$
of $(\Esp)$, regard the category
$$
(\Esp)_{/X}
$$
of objects of $(\Esp)$ over $X$, \ie topological spaces over $X$, as a
site, using the topology induced by that of $(\Esp)$ via the forgetful
functor $(\Esp)_{/X}\rightarrow (\Esp)$
(\hyperrefext[III][III.prop5.2]{5.2} 4)). This site is called the
\emph{big site} associated with $X$. One should note that it is not
$\in 𝒰$; nor is it a $𝒰$-site in the sense of
\hyperrefext[II][II.def3.0.2]{3.0.2}, so precautions are necessary when
applying the usual results to it. To remedy this inconvenience, one may
choose a universe $𝒱$ such that $𝒰\in 𝒱$, so that
$(\Esp)_{/X}$ becomes a $𝒱$-site, and work with the associated
$𝒱$-topos $(\Esp)_{/X}$, which may be denoted by $\TOP(X)$ and will
be called the \emph{big topos of} $X$. If one is reluctant to enlarge
$𝒰$, one may choose a cardinal $c$ bounding the cardinalities of $X$
and of all the topological spaces one intends to use in the arguments
(most often,
\sisi{$\Sup(\card X, \card ℛ)$}{$\Sup(\card X, \card \mathbb{R})$}
will suffice!), and replace $(\Esp)_{/X}$ by the subcategory
$(\Esp)'_{/X}$ consisting of those $X'$ over $X$ such that
$\card X'\leq c$, endowed with the induced topology; one then denotes
by $\TOP(X)$ the topos of sheaves on this\pageoriginale site. To fix
ideas, suppose that the first definition has been adopted.

The advantage of the big topos of $X$ over the small one is that the
site defining it contains $(\Esp)_{/X}$ as a full subcategory; since
the topology of this site is manifestly coarser than the canonical
topology, the canonical functor from $(\Esp)_{/X}$ to $\TOP(X)$,
associating with every space $X'$ over $X$ the sheaf it represents,
is seen to be \emph{fully faithful.} Consequently, \emph{a space $X'$
over $X$ is determined up to $X$-isomorphism by the sheaf
$(\in \Top(X))$ it defines}; thus the notion of a sheaf on (the big
site of) $X$ may be regarded as a \emph{generalization} of the notion
of a topological space over $X$, by means of which all constructions
of sheaf theory acquire a meaning for topological spaces over $X$.

Thus, when $G$ is a group object of the category $(\Esp)_{/X}$ of
topological spaces over $X$, one may associate with it the classifying
topos $B_G$ (\hyperref[IV.subsec2.4]{2.4}), hence classifying cohomology
groups, classifying homotopy groups, \etc (defined as the corresponding
invariants of the $𝒱$-topos $B_G$). In particular, when $X$ is a
one-point space, $G$ is identified with an ordinary topological group.
Subject to the usual local conditions ensuring that the singular
cohomology of the Cartesian products $G^n$ agrees with cohomology in
the sense of sheaves (for constant coefficients, say), for example if
$G$ is locally contractible, one may verify that the cohomology of the
classifying topos of $G$ is canonically isomorphic to that of the
classifying space of $G$ in the sense of topologists.

The introduction\pageoriginale of classifying topoi (via ``big
sites'') has over classifying spaces the advantage of providing a
richer theory, since in particular they provide useful cohomological
invariants for coefficients more general than constant or locally
constant coefficients. Moreover, the definition considered here
adapts in an evident way to the other customary contexts:
differentiable manifolds, real or complex analytic manifolds or
spaces, and schemes. This point of view makes it possible in particular
to connect the study of characteristic classes from the traditional
and the ``arithmetic'' points of view, by regarding the ``classical
groups'' as arising from schemes defined over the ring of integers;
\cf [\hyperref[IV.sec7]{7}] for indications in this direction.
Likewise, the general results of J. \textsc{Giraud}
[\hyperref[IV.sec3]{3}] on the classification of extensions of group
objects, developed in the very general and flexible setting of topoi,
may, by means of ``big topoi'', be specialized to results on the
classification of extensions of topological groups, or of real or
complex Lie groups, which seem to have been scarcely known to
topologists except in the case of extensions with abelian kernel
[\hyperref[IV.sec11]{11}].

\subsection{Topoi of the form $\widehat{C}$}\etiquette[2.6]{IV.subsec2.6}

Let $C$ be a small category. Then the category $\widehat{C}$ of
$U$-presheaves on $C$ is plainly a $𝒰$-topos, since it is of the
form $C^{\widetilde{\phantom{,}}}$, where $C$ is endowed with the
chaotic topology. We shall give below some details on the relations
between $C$ and $\widehat{C}$. Let us note here only that a topos $E$
equivalent to a topos of the form $\widehat{C}$ is of a rather special
nature, because \emph{it has a small generating family $(X_i)$
consisting of connected projective objects}, \ie objects $X$ such
that the functor $Y\mapsto \Hom(X,Y)$\pageoriginale takes epimorphisms
to epimorphisms and coproducts to coproducts: indeed, it suffices to
take in $\widehat{C}$ the generating family consisting of the functors
represented by the $X\in \ob C$. Moreover, note that if in a topos
$E$ one has a covering family $X_i\rightarrow X$, where the $X_i$ are
projective and connected, then every other covering family of $X$ is
refined by the former (\hyperrefext[I][I.sec4.3.2]{4.3.2},
\hyperrefext[I][I.sec4.3.3]{4.3.3}). Consequently, in a topos $E$ of
the form $\widehat{C}$ every object $X$ admits a covering family
refining all the others. A topos of the form $\Top(X)$
(\hyperref[IV.subsec2.1]{2.1}), where $X$ is a topological space whose
points are closed, has the preceding property only if $X$ is discrete.

When the category $C$ has a single object, $C$ is identified with a
monoid $G$. A presheaf on $C$ is then identified with a set on which
$G$ acts \emph{on the right} (since it is a functor
$G^{\circ}\rightarrow (\Ens)$), and $\widehat{C}$ is the topos of
sets with a monoid of right operators, which may also be denoted by
$B_{G^{\circ}}$ in view of \hyperref[IV.subsec2.3]{2.3}: it is the
\emph{topos of sets with monoid of operators $G^{\circ}$ (the monoid
opposite to $G$)}. When $G$ is a group, using the isomorphism
$g\mapsto g^{-1}$ from $G$ onto $G^{\circ}$, one recovers the
classifying topos $B_G$ of \hyperref[IV.subsec2.3]{2.3}.

\subsection{Classifying topos of a pro-group}\etiquette[2.7]{IV.subsec2.7}

\subsubsection{}\etiquette[2.7.1]{IV.subsubsec2.7.1}
Let $𝒢=(G_i)_{i\in I}$ be a projective system of groups, with
$G_i$, $I\in 𝒰$. Assume that the projective system is
\emph{strict}, \ie that the transition morphisms $G_j\rightarrow G_i$
are surjective. If $E$ is a set, by a (left, say) \emph{action of
$𝒢$ on $E$} we mean the following structure: a) a family
$(E_i)_{i\in I}$ of subsets of $E$, whose union is $E$; b) for every
$i\in I$, an action of the group $G_i$ on the set $E_i$; these data
are furthermore required to satisfy\pageoriginale the following
condition: for $j\geq i$, $E_i$ is the subset of $E_j$ consisting of
the elements fixed under the kernel of $G_j\rightarrow G_i$. We also
say that $𝒢$ \emph{acts on} $E$ (on the left) if a (left) action of
$𝒢$ on $E$ has been given. The sets $\in 𝒰$ endowed with an
action of $𝒢$ form a category in the evident way. By Giraud's
criterion, one sees at once that this category is a $𝒰$-topos. It
is denoted by $B_{𝒢}$ and called the \emph{classifying topos of}
$𝒢$. When $I$ has an initial object $i_{\circ}$, on putting
$G=G_{i_\circ}$ one recovers the classifying topos of
\hyperref[IV.subsec2.3]{2.3}.

\subsubsection{}\etiquette[2.7.2]{IV.subsubsec2.7.2}
Another important example is that in which the groups $G_i$ are
finite, so that
$$
G=\varprojlim G_i
$$
is a compact totally disconnected topological group, or
\emph{profinite group.} An action of $𝒢$ on $E$ then amounts to a
continuous action of $G$ on $E$, or, equivalently, to one for which
the stabilizer of every point of $E$ is an open subgroup of $G$.
The classifying topos $B_{𝒢}$ will also be denoted by $B_{G}$, where,
of course, $G$ is to be regarded as endowed with its profinite
topology.

\subsubsection{}\etiquette[2.7.3]{IV.subsubsec2.7.3}
Using the remarks of \hyperref[IV.subsec2.6]{2.6}, it is easy to verify
that the topos $B_{𝒢}$ defined by a strict projective system
$𝒢=(G_i)_{i\in I}$ of groups is equivalent to a topos of the form
$\widehat{C}$ only if this projective system is essentially constant;
in the case of a profinite group, this means that the group is in fact
finite.

\subsubsection{}\etiquette[2.7.4]{IV.subsubsec2.7.4}
The following geometric interpretation of the classifying topos $B_G$
of a discrete group $G$ is useful for giving a correct geometric
intuition for these topoi. (\Cf also, in the same vein,
\hyperref[IV.subsec4.5]{4.5}, \hyperref[IV.subsec5.8]{5.8},
\hyperref[IV.exer5.9]{5.9}, and \hyperref[IV.subsec7.2]{7.2} below.)
Let $X$ be a connected, locally connected, and\pageoriginale locally
simply connected topological space, let $x$ be a point of $X$, and let
$G$ be its fundamental group at $x$. (NB It is known that, up to
isomorphism, every discrete group $G$ may be obtained in this way.)
Then the Galois theory of coverings of $X$ gives an equivalence between
the category $B_G$ of $G$-sets and the category of \emph{étale
coverings} of $X$, \ie spaces $X'$ over $X$ that are locally
$X$-isomorphic to $X$-spaces of the form $X\times I$, where $I$ is a
discrete space. (Compare \SGA 1 V 4,5.) The latter can moreover be
interpreted as the category of locally constant sheaves on $X$, \ie
the category of locally constant objects (IX \hyperref[IX.subsec2.0]{2.0}) of the topos
$\Top(X)$.

When $G$ is a profinite group, there is an analogous geometric
interpretation of $B_G$ as the category of $X$-schemes that are
coproducts of finite étale coverings of a connected scheme $X$,
equipped with a geometric point $x$ and an isomorphism
$G\simeq \pi_1(X,x)$. It is again known that every profinite group may
be obtained as the fundamental group of a suitable connected scheme
(the spectrum of a field, if desired). Finally, pro-groups (not
necessarily profinite or essentially constant) also occur in the
classification of coverings of connected and locally connected spaces
that are not locally simply connected, or in the classification of
étale coverings, not necessarily finite or ind-finite, of connected
nonnormal schemes. For the latter case, \cf \SGA 3 X 6.

\begin{enonce*}[remark]{Exercise 2.7.5}\etiquette[2.7.5]{IV.exer2.7.5}
For a topos $E$, define the notions of connectedness, local
connectedness~\footnote[1]{\cf
\sisi{\hyperref[IV.exer8.7]{8.7} 1).}{\hyperref[IV.exer8.7]{8.7} l).}},
simple connectedness, and local simple connectedness. Define the notion
of a constant and locally constant object of $E$ (\cf IX \hyperref[IX.subsec2.0]{2.0}). Let
$f:E'\rightarrow E$\pageoriginale be a morphism of topoi
(\hyperref[IV.def3.1]{3.1}), with $E'$ simply connected (for example,
$E'$ the point topos (\hyperref[IV.subsec2.2]{2.2})), and $E$ connected
and locally connected. Define a strict pro-group
$\pi_1(E,f)=(G_i)_{i\in I}=𝒢$ (called the \emph{fundamental
pro-group of $E$ at $f$}) and an equivalence of categories between
$B_{𝒢}$ and the category of locally constant objects of
$E$.\footnote[2]{One may take inspiration from \SGA 3 X 6.} Show that,
when $E$ is locally simply connected, $\pi_1(E,f)$ is essentially
constant and hence is identified with an ordinary discrete group
$\pi_1(E,f)$, called the \emph{fundamental group of $E$ at $f$}.
Show that every strict pro-group $𝒢$ is isomorphic (as a pro-group)
to the fundamental pro-group of a suitable connected and locally
connected topos at a suitable $f$, with $E'$ the point topos (take
$E=B_{𝒢}$, and $f:E'\rightarrow E$ defined by the forgetful functor
$f^{\ast}:E\rightarrow E'=(\Ens)$). When $𝒢$ is essentially
constant, \ie isomorphic (as a pro-group) to an ordinary discrete
group, prove that one may take $E$ above to be locally simply connected
(again take $E=B_G$).
\end{enonce*}

\subsection{Example of a false topos}\etiquette[2.8]{IV.subsec2.8}

Let $𝒢=(G_i)_{i\in I}$ be a strict pro-group, where $I$ is a
filtered ordered set and where $i>j$ implies that
$G_i\rightarrow G_j$ is not an isomorphism. Suppose that
$\card(I)\not\in 𝒰$. Consider the category of sets $E\in 𝒰$
on which $𝒢$ acts on the left
(\hyperref[IV.subsubsec2.7.1]{2.7.1}). It is a $𝒰$-category, and one
sees as in \hyperref[IV.subsubsec2.7.1]{2.7.1} that this category
satisfies conditions a), b), and c) of
\hyperref[IV.subsubsec1.1.2]{1.1.2}. Nevertheless, it is not a
$𝒰$-topos, since it has no generating family that is
$𝒰$-small. Similarly, it is not a $𝒱$-topos for any universe
$𝒱$.

% END INCLUDED SGA4_English_current_buildable_source_20260906/source/04/drafts/SGA4_IV_LINES_537_947_English.texfrag


% BEGIN INCLUDED SGA4_English_current_buildable_source_20260906/source/04/drafts/SGA4_IV_LINES_948_1319_English.texfrag
\section{Morphisms of topoi}\etiquette[3]{IV.sec3}\pageoriginale

\begin{enonce*}{Definition 3.1}\etiquette[3.1]{IV.def3.1}
Let $E$ and $E'$ be two $𝒰$-topoi. A morphism\nde{Some authors
speak of a \emph{geometric morphism}.} from
$E$ to $E'$, or sometimes (by abuse of language) a continuous map
from $E$ to $E'$, is a triple $u=(u_{\ast}, u^{\ast},\varphi)$,
consisting of functors
$$
u_{\ast}:E\longrightarrow E'\quav u^{\ast}:E'\longrightarrow E
$$
and an ``adjunction'' isomorphism $\varphi$ of bifunctors
in $X'\in \Ob E'$, $Y\in \Ob E$:
$$
\varphi:\Hom_E (u^{\ast}(X'),Y)\xrightarrow{\sim} \Hom_{E'}(X',
u_{\ast}(Y)){\quoi},
$$
where the functor $u^{\ast}$ is moreover required to be left exact,
\ie to commute with finite projective limits. The functor $u_{\ast}$
is called the direct-image functor for the morphism of topoi $u$,
the functor $u^{\ast}$ is called the inverse-image functor for the
morphism of topoi $u$, and the isomorphism $\varphi$ is called the
adjunction isomorphism for $u$.
\end{enonce*}

\setcounter{subsection}{1}
\subsubsection{}\etiquette[3.1.1]{IV.subsubsec3.1.1}
Unless expressly stated otherwise, for a morphism of topoi
$u:E\rightarrow E'$, we shall henceforth denote the corresponding
direct- and inverse-image functors by $u_{\ast}$ and $u^{\ast}$
\footnote[1]{It is sometimes appropriate also to write $u^{-1}$
instead of $u^{\ast}$, \cf \hyperref[IV.subsubsec13.2.3]{13.2.3}.}.
Note that, since $u_{\ast}$ is right adjoint to $u^{\ast}$ and
$u^{\ast}$ is left adjoint to $u_{\ast}$ by the adjunction
isomorphism $\varphi$, each of the two functors $u_{\ast}$,
$u^{\ast}$ determines the other up to unique isomorphism, by the
well-known general facts about adjoint functors [\hyperref[IV.sec14]{14}].
In practice, depending on the circumstances, it may be more convenient
to define a morphism of topoi $u:E\rightarrow E'$ either by specifying
$u^{\ast}:E'\rightarrow E$, or by specifying
$u_{\ast}:E\rightarrow E'$; in the first case, one need only verify
that the given functor $u^{\ast}$ has\pageoriginale a right adjoint
and is left exact. In the second case, one must verify that the given functor
$u_{\ast}$ has a left adjoint which is left exact. In either case,
from the partial data, by \emph{choosing} an adjoint functor and an
adjunction isomorphism, one obtains a morphism of topoi
$u:E\rightarrow E'$, and the latter will be ``unique up to unique
isomorphism'' in terms of the datum $u^{\ast}$ \resp $u_{\ast}$, in
a sufficiently clear sense which will moreover be made completely
explicit below \hyperlink{sga4.exp.IV.subsubsection.3.2.1}{(3.2.1)}.

\subsubsection{}\etiquette[3.1.2]{IV.subsubsec3.1.2}
If $u:E\rightarrow E'$ is a morphism of topoi, it follows from the
properties of adjoint functors (\hyperrefext[I][I.prop2.11]{2.11}) that
the direct-image functor $u_{\ast}:E\rightarrow E'$ commutes with
projective limits, and the functor $u^{\ast}:E'\rightarrow E$ commutes
with inductive limits \footnote[2]{Moreover
(\hyperref[IV.cor1.5]{1.5} and \hyperref[IV.rem1.8]{1.8}), for a given functor
$u_{\ast}:E\rightarrow E'$ \resp $u^{\ast}:E'\rightarrow E$,
this functor has a left adjoint (\resp a right adjoint) if and only if
it commutes with $𝒰$-projective limits (\resp $𝒰$-inductive
limits).}\MakeLinkTarget*{sga4.exp.IV.footnote.3.1.1.2}\label{sga4[1]}. Since the latter is assumed in addition to
be left exact, \ie to commute with finite projective limits, we see in
particular that $u^{\ast}$ \emph{is exact}. We may therefore say that
\emph{it is the inverse-image functor $u^{\ast}$}, in the pair
$(u_{\ast},u^{\ast})$, \emph{which has the most remarkable exactness
properties.} These properties ensure that, for every type of
algebraic structure $\Sigma$ whose data can be described in terms of
arrows between the underlying sets and sets deduced from them by
repeated application of operations of finite projective limits and
arbitrary inductive limits, and for every ``object of $ℰ'$ endowed
with a structure of type $\Sigma$'' (a notion which makes sense by
virtue of the internal exactness properties of the topos $ℰ'$
(\hyperrefext[II][II.theo4.1]{4.1})), its image under $u^{\ast}$ is endowed
with the same structures. Rather than undertake the unrewarding task
of giving a precise meaning to this statement and formally justifying
it,\pageoriginale we advise the reader to spell it out and convince
himself of its validity for types of structure such as groups, rings,
modules over a ring, comodules over a coring, bialgebras over a ring,
and torsors under a group. (In these examples, the first three types
of structure are defined exclusively in terms of finite projective
limits, while the other notions implicitly involve constructions
which also use inductive limits.) Furthermore, the functor
$u^{\ast}$ ``\emph{commutes}'' with all the usual functorial
operations performed in terms of such structures, more precisely
with all operations which can be expressed in terms of
$\varinjlim$ and finite $\varprojlim$: constructions of free objects
(free groups or modules, for example) generated by an object, tensor
products (\cf \hyperref[IV.sec12]{\S 12} below), \etcn

As for the direct-image functor $u_{\ast}:E\rightarrow E'$, since it
commutes with projective limits, it consequently ``preserves'' every
algebraic structure on an object (or a family of objects) of $E$
which can be defined exclusively in terms of projective limits (such
as the structures of group, ring, or module over a ring among the
preceding examples). On the other hand, the functor $u_{\ast}$ is not
in general right exact, \ie it does not in general commute with finite
inductive limits, and it does not even in general take epimorphisms to
epimorphisms (it is, moreover, this failure of exactness of
$u_{\ast}$ which is the source of its cohomological properties; these
will be studied, from the viewpoint of commutative homological
algebra, in the following exposé). Consequently it does not, in
general, extend to a functor on objects of the type comodule,
\pageoriginale bialgebra, or torsor under a group, and does not in
general commute with operations such as ``free module generated by'',
tensor product of modules, \etc

\subsubsection{}\etiquette[3.1.3]{IV.subsubsec3.1.3}
In practice, when one is given a functor $f:E\rightarrow F$ from one
$𝒰$-topos to another, one should make all of its exactness properties
explicit, including the possible existence of left or right adjoint
functors (\cf the footnote on page \pageref{sga4[1]}), in order to
understand the ``geometric nature'' of $f$, an understanding which
will generally be an indispensable guide to a correct geometric
intuition of the situation. Thus, if it turns out that $f$ commutes
with arbitrary inductive limits and finite projective limits, one
should write $f$ in the form
$$
f=u^{\ast}{\quoi},
$$
where
$$
u:F\longrightarrow E
$$
is a morphism of topoi, \ie interpret $f$ as an ``inverse-image''
functor for a ``continuous map'' of topoi. When $f$ commutes with
arbitrary projective limits, so that it has a left adjoint, and if
the latter (which a priori commutes with arbitrary inductive limits)
\emph{also} commutes with finite projective limits, one should write
$f$ in the form
$$
f=v_{\ast}{\quoi},
$$
where
$$
v:E\longrightarrow F
$$
is a morphism of topoi. In certain cases it may happen that $f$
satisfies\pageoriginale both of the properties under consideration
(\cf \hyperref[IV.subsec4.10]{4.10} for an example). In that case one
should introduce simultaneously the morphisms of topoi
$$
u:F\longrightarrow E \quav v:E\longrightarrow F{\quoi},
$$
which give rise to a sequence of three adjoint functors
(\hyperrefext[I][I.sec5.3]{5.3}):
$$
v^{\ast} \quav v_{\ast}=u^{\ast}=f\quav u_{\ast}{\quoi}.
$$
One must then distinguish carefully between the morphisms of topoi
$u$ and $v$, lest one lose the geometric intuition of the situation.

In this connection, note that when between two topoi $E$, $F$ one has
a sequence of three adjoint functors
$$
e, f, g\; (e,g: F\rightrightarrows E, f:E\rightarrow F){\quoi},
$$
then $f$ commutes with arbitrary inductive and projective limits, and
therefore can always be written in the form $u^{\ast}$, where
$u:F\rightarrow E$ is a morphism of topoi. Thus $g$ can be written
$g=u_{\ast}$. On the other hand, of course, $e$ can be written in the
form $v^{\ast}$ (and then $f$ in the form $v_{\ast}$) only if it also
commutes with finite projective limits. This will evidently be the
case if $e$ is itself the right adjoint of a fourth functor $d$.
Without a condition of this kind, one will often write
$$
e=u_{!}
$$
for the left adjoint of an inverse-image functor $f=u^{\ast}$ when
such a left adjoint exists, this notation being suggested by the
example of an open immersion $u:X\rightarrow Y$ of topological
spaces. Similarly, if $e$ is of the form $v^{\ast}$, \ie if $f$ is
of the form $v_{\ast}$, one sometimes writes $g=v^{!}$ for the right
adjoint of a direct-image functor $v_{\ast}$, when this adjoint exists
\footnote[1]{\Cf below \hyperref[IV.subsec14.4]{14.4}, in the case of
abelian sheaves.}.

\subsection{}\etiquette[3.2]{IV.subsec3.2}
Let\pageoriginale $E$, $E'$ be two $𝒰$-topoi, and let
$$
u=(u_{\ast},u^{\ast},\varphi)\quav
v=(v_{\ast},v^{\ast},\Psi):E\rightrightarrows E'
$$
be two morphisms of topoi from $E$ to $E'$. A \emph{morphism from
$u$ to $v$} is any morphism from $u_{\ast}$ to $v_{\ast}$ (in the
sense of the category $\Hhom(E,E')$ of functors from $E$ to $E'$).
Morphisms of morphisms of topoi compose in the evident way, thereby
defining a category, which is in fact a
$\sheaf{\mathcal{U}}$-\emph{category} (\hyperrefext[I][I.prop7.8]{7.8}),
denoted
$$
\Hhomtop(E,E'){\quoi},
$$
and called the \emph{category of morphisms} (or \emph{continuous
maps}) \emph{from $E$ to $E'$}. One then has an evident functor
$$
u\longmapsto u_{\ast}:\Hhomtop(E,E')\longrightarrow
\Hhom(E,E'){\quoi},
$$
which is fully faithful by the definition of arrows in the first
member (but is not injective on objects).

\hypertarget{sga4.exp.IV.subsubsection.3.2.1}{}\subsubsection{}
Care must be taken that, if $u$ and $v$ are given as above, the theory
of adjoint functors supplies a canonical bijection
$$
\Hom(u_{\ast},v_{\ast})\xrightarrow{\sim}\Hom(v^{\ast},
u^{\ast}){\quoi};
$$
in particular, one obtains a \emph{contravariant functor} on
$\Hhomtop(E,E')$:
$$
u\longmapsto u^{\ast}:\Hhomtop(E,E')^{\circ}\longrightarrow
\Hhom(E,E'){\quoi}.
$$
In accordance with geometric intuition, according to which the
direct-image functor $u_{\ast}$ ``goes in the same direction'' as the
continuous map which gives rise to it, one should therefore define
the direction of arrows for morphisms between morphisms of topoi
\emph{in terms of the direct-image functors}, and not in terms of the
inverse-image functors (although the latter, as we have seen, possess
the exactness properties characteristic of the notion of morphism of
topoi).

\subsubsection{}\etiquette[3.2.2]{IV.subsubsec3.2.2}
Having\pageoriginale defined the category $\Hhomtop(E,E')$ of
morphisms from the topos $E$ to the topos $E'$, the notion of an
\emph{isomorphism} between two morphisms $u$, $v$ from $E$ to $E'$
is also defined. In practice there is no need to distinguish
essentially between two isomorphic morphisms of topoi, at least when
one has a \emph{canonical} isomorphism between the two (just as there
is often no need to distinguish between two objects of a category
when a canonical isomorphism between them is given). Let us point out
in this connection that, most often, when one deals with diagrams of
morphisms of topoi and questions of commutativity of such diagrams
(a notion which makes sense by \hyperlink{sga4.exp.IV.subsection.3.3}{(3.3)}), one means only commutativity up
to (``canonical'') isomorphism; by abuse of language, one then treats
these diagrams as actually commutative diagrams.

One might try to justify this abuse of language by introducing the
set \\ $\Homtop(E,E')/\Isom$ of morphisms from $E$ to $E'$ up to
isomorphism, and calling such an isomorphism class a morphism, instead
of following Definition \hyperref[IV.subsec3.2]{3.2}. But this runs
into the very serious difficulties which arise whenever one tries to
identify two isomorphic objects of a category without having a
canonical isomorphism between them. Experience shows that such a
viewpoint is impracticable, and that one must retain the ``fine''
notion \hyperref[IV.def3.1]{3.1} of a morphism between morphisms of
topoi, even if one is sometimes obliged to struggle with
compatibilities between canonical isomorphisms \footnote[1]{For
examples of such struggles (apparently successful), we refer the
reader to Mme M. \textsc{Hakim}'s book on relative schemes [9].}.

\hypertarget{sga4.exp.IV.subsection.3.3}{}\setcounter{subsection}{3} \setcounter{subsubsection}{0}
\subsubsection{}\etiquette[3.3.1]{IV.subsubsec3.3.1}
Let\pageoriginale $E$, $E'$, $E''$ be three $𝒰$-topoi, and consider
morphisms of topoi
$$
u:E\longrightarrow E'\quav v:E'\longrightarrow E''{\quoi}.
$$
The theory of adjoint functors then gives us an adjunction
isomorphism between the composite functors $v_{\ast}u_{\ast}$ and
$u^{\ast}v^{\ast}$, in terms of the adjunction isomorphisms for the
pairs $(u_{\ast},u^{\ast})$ and $(v_{\ast},v^{\ast})$. On the other
hand, the functor $u^{\ast}v^{\ast}$ is left exact, as the composite
of two left exact functors. Consequently one obtains a morphism from
$E$ to $E''$, called the \emph{composite of the morphisms $u$ and
$v$}, and denoted $vu$:
$$
vu:E\longrightarrow E''{\quoi}.
$$
One verifies at once that composition of morphisms is associative,
and that for every $𝒰$-topos $E$ there is a morphism from $E$ to
itself which is a two-sided identity for composition: it is the
morphism $(\id_E,\id_E,\varphi)$, where $\varphi$ is the evident
adjunction isomorphism of $\id_E$ with itself. Let $𝒱$ be a universe
such that $𝒰\in 𝒱$. One defines a category
$$
(𝒱\text{-}𝒰\text{-}\Top){\quoi},
$$
whose objects are the $𝒰$-topoi which belong to $𝒱$, whose arrows
are morphisms between such $𝒰$-topoi, and whose composition of arrows
is the one just described.

\subsubsection{}\etiquette[3.3.2]{IV.subsubsec3.3.2}
In fact, the composition map
$$
\Homtop(E, E')\times \Homtop(E',E'')\longrightarrow \Homtop(E,E'')
$$
is the map induced on objects by a ``functor of composition of
morphisms'':
$$
\Hhomtop(E,E')\times \Hhomtop(E',E'')\longrightarrow
\Hhomtop(E,E''){\quoi},
$$
whose\pageoriginale effect on arrows is the usual ``convolution
product'' operation for morphisms between functors (here,
direct-image functors). These composition functors satisfy a (strict)
associativity property, for four topoi $E$, $E'$, $E''$, $E'''$,
which refines the associativity of composition of morphisms of topoi.
One may also say, in language which is beginning to become familiar
[2], [9], that the $𝒰$-topoi are the objects (or 0-arrows) of a
2-\emph{category}, whose 1-arrows are the morphisms of topoi and
whose 2-arrows are morphisms of morphisms of topoi.

It is the fact that the $𝒰$-topoi (elements of a universe $𝒱$) form
a 2-category, rather than merely an ordinary category like ordinary
topological spaces, which from the technical viewpoint constitutes
the most important difference between the theory of topoi and that
of topological spaces. This fact is the source of certain technical
complications already alluded to, but also of facts essentially new
in comparison with traditional topology.

\subsection{}\etiquette[3.4]{IV.subsec3.4}
The fact that the $𝒰$-topoi (elements of a universe $𝒱$) form a
2-category (\hyperref[IV.subsubsec3.3.2]{3.3.2}) makes it possible in
particular to define the notion of \emph{equivalence of two
$𝒰$-topoi $E$, $E'$}: one says that $E$ and $E'$ are
\emph{equivalent} if there exist morphisms of topoi
$u:E\rightarrow E'$ and $v:E'\rightarrow E$ such that the composites
$uv$ and $vu$ are isomorphic respectively to the identity morphisms
of $E$ and $E'$; the morphisms $u$ and $v$ are then said to be
\pageoriginale \emph{quasi-inverse equivalences of one another.}

One sees at once that for the morphism of topoi $u:E\rightarrow E'$
to be an equivalence, it is necessary and sufficient that $u_{\ast}$
be an equivalence, or, what amounts to the same thing, that
$u^{\ast}$ be an equivalence. (Use the fact that a functor
$f:E\rightarrow E'$ between two topoi which is an equivalence is
both of the form $u_{\ast}$ and of the form $v^{\ast}$, with $u$ and
$v$ morphisms of topoi); and for $E$ and $E'$ to be equivalent in
the sense of the preceding paragraph, it is necessary and sufficient
that they be equivalent as categories (\ie as objects of the
2-category $(𝒱\text{-}\cat)$). As it should, this shows that the
notions of equivalence introduced do not depend on the choice of the
universe $𝒱$ in \hyperref[IV.subsubsec3.3.1]{3.3.1}.

\subsubsection{}\etiquette[3.4.1]{IV.subsubsec3.4.1}
In practice, there is most often no reason to distinguish essentially
between equivalent $𝒰$-topoi, just as there is often no reason to
distinguish essentially between two equivalent categories---provided,
however, that one has an explicit equivalence from one to the other,
or at least an equivalence defined up to unique isomorphism. Here it
is the notion of equivalence of topoi which replaces the traditional
notion of homeomorphism between two topological spaces. See the
example in \hyperref[IV.subsec4.2]{4.2} below for the precise relations
between these two notions.

% END INCLUDED SGA4_English_current_buildable_source_20260906/source/04/drafts/SGA4_IV_LINES_948_1319_English.texfrag


% BEGIN INCLUDED SGA4_English_current_buildable_source_20260906/source/04/drafts/SGA4_IV_LINES_1320_1659_English.texfrag
\section{Examples of morphisms of topoi}\etiquette[4]{IV.sec4}

Here we return to the examples of 2, using the notion of a
morphism of topoi. The comments of \hyperref[IV.subsec2.0]{2.0}
apply equally to the present section. The universe
$𝒰$ will generally be understood.

\subsection{The topos $\Top(X)$ for a varying topological space $X$}
\etiquette[4.1]{IV.subsec4.1}\pageoriginale

\subsubsection{}\etiquette[4.1.1]{IV.subsubsec4.1.1}
Let
$$
f:X\longrightarrow Y
$$
be a continuous map of topological spaces. We shall associate with it
canonically a morphism of topoi
$$
\Top(f)\mbox{ or } f:\Top(X)\rightarrow \Top(Y){\quoi},
$$
with the notation of \hyperref[IV.subsec2.1]{2.1}. When $\Top(X)$ is
defined as $\Ouv(X)^{\widetilde{\phantom{,}}}$, the most convenient
description of $\Top(f)$ is by means of the direct-image functor on
sheaves
$$
f_{\ast}:\Top(X)\longrightarrow \Top(Y){\quoi},
$$
defined by the formula
$$
f_{\ast}(F)=F\circ f^{-1}{\quoi},
$$
where
$$
f^{-1}:\Ouv(Y)\longrightarrow \Ouv(X)
$$
is the evident functor $U\rightarrow f^{-1}(U)$. We have already noted
that this functor is continuous and left exact, and therefore indeed
defines the functor $f_{\ast}$ above, which admits a left adjoint
$f^{\ast}$ that is left exact (\hyperrefext[III][III.exem1.9.1]{1.9.1}).
Of course, strictly speaking, the morphism of topoi $\Top(f)$ depends
on the choice of the left adjoint $f^{\ast}$ of $f_{\ast}$ and hence
is defined only up to canonical isomorphism. In what follows we shall
refrain from explicitly pointing out phenomena of this kind.

When one adopts the ``étale spaces'' point of view in defining
$\Top(X)$, it is the inverse-image functor
$$
f^{\ast}:\Top(Y)\longrightarrow \Top(X)
$$
that\pageoriginale is the most convenient for defining the morphism of
topoi $\Top(f)$, by simply setting
$$
f^{\ast}(Y')=\fprod{X}{Y'}{Y}
$$
for every étale space $Y'$ over $Y$; it is clear that the fiber product
is indeed an étale space over $X$, and that the functor $f^{\ast}$ so
obtained is left exact and commutes with arbitrary $\varinjlim$, and
therefore defines a morphism of topoi $\Top(f)$. For the compatibility
of the two definitions obtained, we refer to [TF].

For two composable continuous maps
$$
X\xrightarrow{f}Y\xrightarrow{g}Z{\quoi},
$$
there is a canonical isomorphism
$\Top(gf)\simeq\Top(g)\Top(f)$ of morphisms of topoi. These
\emph{transitivity isomorphisms}, for three composable continuous maps
$f$, $g$, $h$, satisfy a compatibility relation which we shall refrain
from writing here, and which is precisely the one considered in
\SGA 1 VI 7.4 B) (for $\epsilon = (\Esp)^{\circ}$). This may be
expressed by saying that, as $X$ varies in the category $(\Esp)$,
$$
X\longmapsto \Top(X)
$$
is a ``pseudofunctor''
\begin{equation*}
(𝒰\text{-}\esp)\longrightarrow
(𝒱\text{-}𝒰\text{-}\top){\quoi},\tag{4.1.1.1}
\end{equation*}\etiquette[4.1.1.1]{IV.eq4.1.1.1}
or, in the terminology of 2-categories, that one has a
\emph{non-strict} functor of 2-categories [9]. In practice, we shall
most often permit ourselves the abuse of language consisting in
identifying $\Top(gf)$ and $\Top(g)\Top(f)$, that is, reasoning as if
\eqhyperref[IV.eq4.1.1.1]{4.1.1.1} were an actual functor between
ordinary categories. We shall permit ourselves analogous abuses of
language in the other examples treated below.\pageoriginale

\subsubsection{}\etiquette[4.1.2]{IV.subsubsec4.1.2}
The preceding considerations extend immediately to topoi associated
with topological spaces equipped with operator groups
(\hyperref[IV.subsec2.3]{2.3}). If $f=(f^{\es},f^{\gr})$,
$$
f:(X,G)\longrightarrow (Y,H)
$$
is a morphism of spaces with operators (where
$$
f^{\es}:X\longrightarrow Y\quav f^{\gr}:G\longrightarrow H
$$
are respectively continuous maps and group homomorphisms, compatible
in the evident sense), one associates with it a morphism of topoi
$$
\Top(f)\mbox{ or } f:\Top(X,G)\longrightarrow \Top(Y,H){\quoi},
$$
whose definition is left to the reader. When the groups $G$ and $H$
are trivial, one recovers the definition of
\hyperref[IV.subsubsec4.1.1]{4.1.1}; when, on the other hand, the spaces
$X$ and $Y$ are each reduced to a point, the inverse-image functor
$f^{\ast}$ is the ``restriction of the operator group'' functor
$$
f^{\ast}:B_H\longrightarrow B_G{\quoi},
$$
which associates with every $H$-set the $G$-set that it defines by
means of $f:G\rightarrow H$. We shall encounter this example again in
other forms in \hyperref[IV.subsec4.5]{4.5} and
\hyperref[IV.subsubsec4.6.1]{4.6.1}.

\subsubsection{}\etiquette[4.1.3]{IV.subsubsec4.1.3}
Given a continuous map $f:X\rightarrow Y$ of topological spaces, one
also associates with it a morphism between the corresponding ``large
topoi'' (\hyperref[IV.subsec2.5]{2.5})
$$
\TOP(f)\mbox{ or } f:\TOP(X)\longrightarrow \TOP(Y){\quoi},
$$
most conveniently defined\pageoriginale by the inverse-image functor
$$
f^{\ast}:\TOP(Y)\longrightarrow \TOP(X){\quoi},
$$
which is simply the \emph{restriction functor}. This morphism
$\TOP(f)$ is a special case of the morphism called an ``inclusion''
for an induced topos, which will be studied in 5. Thus, subject to the
same qualifications as in \hyperref[IV.subsubsec4.1.1]{4.1.1}, the
topos $\TOP(X)$ may be regarded as a functor of $X$, as $X$ varies in
$(\Esp)$.

\subsection{Fidelity properties of $X\mapsto \Top(X)$}\etiquette[4.2]{IV.subsec4.2}

We propose to specify to what extent a topological space $X$ can be
reconstructed in terms of the topos $\Top(X)$. For this purpose it is
appropriate, for two spaces $X$ and $Y$, to describe the category of
morphisms from $\Top(X)$ to $\Top(Y)$
(\hyperref[IV.subsec3.2]{3.2}), so as to determine the fidelity
properties of the ``functor'' $X\mapsto \Top(X)$. We shall merely
state the results obtained, referring the reader to [9] for details.
The reader who wishes to carry out the verification himself may
consult Exercise \hyperref[IV.exer7.8]{7.8}.

\subsubsection{}\etiquette[4.2.1]{IV.subsubsec4.2.1}
Recall that a topological space $X$ is called \emph{sober} if every
irreducible closed subset of $X$ has exactly one generic point. Almost
all spaces used in practice are sober; this is true in particular of a
Hausdorff space, more generally of a space all of whose points are
closed, or of the underlying space of a scheme. With a topological
space $X$ one associates (\loccit or \EGA $0_I$, new edition) a sober
topological space $X_{\sob}$ and a continuous map\pageoriginale
\begin{equation*}
\varphi: X\longrightarrow X_{\sob}\tag{4.2.1.1}
\end{equation*}\etiquette[4.2.1.1]{IV.eq4.2.1.1}
which is \emph{universal} for continuous maps from $X$ to sober spaces;
in other words, one constructs a left adjoint $X\mapsto X_{\sob}$ to
the inclusion functor $(\Espsob)\rightarrow (\Esp)$ from the category
of sober spaces to that of ``arbitrary'' topological spaces (the
quotation marks serving as a reminder that a universe is involved!).
The explicit construction takes as the points of $X_{\sob}$ the
irreducible closed subsets of $X$, and as its open subsets the sets of
the form $U'$, where $U$ is an open subset of $X$ and
$U'\subset X_{\sob}$ denotes the set of irreducible closed subsets of
$X$ that meet $U$. The map
\eqhyperref[IV.eq4.2.1.1]{4.2.1.1} is obtained by associating with each
$x\in X$ the closure of $\{x\}$. The space $X$ is sober if and only if
the preceding map is bijective, and hence a homeomorphism.

One checks that the functor
$$
\varphi^{-1}:\Ouv(X_{\sob})\longrightarrow \Ouv(X)
$$
induced by $\varphi$ is an isomorphism, which implies that the morphism
of topoi
$$
\Top(\varphi):\Top(X)\longrightarrow \Top(X_{\sob})
$$
defined by $\varphi$ is also an isomorphism. This explains a priori
why $X_{\sob}$ must necessarily enter into the question of
reconstructing $X$ from $\Top(X)$: since the latter depends only on
$X_{\sob}$ up to isomorphism, the question can have an affirmative
answer only if $X$ is sober. We shall specify below
(\hyperref[IV.subsec7.1]{7.1}) how $X_{\sob}$ can in fact be
reconstructed\pageoriginale in terms of $\Top(X)$, by interpreting its
points as ``points'' of the topos $\Top(X)$ (or, equivalently, as fiber
functors).

\subsubsection{}\etiquette[4.2.2]{IV.subsubsec4.2.2}
On every topological space one should introduce the order relation
$\leq$ for which
$$
x\leq y\Longleftrightarrow \bar{\{x\}}\subset \bar{\{y\}}\mbox{ \ie
} x\in \bar{\{y\}}
$$
(one also expresses this by saying that $x$ is a
\emph{specialization} of $y$, or that $y$ is a \emph{generization} of
$x$). For a space of the form $X_{\sob}$, this is simply the inclusion
relation among the irreducible closed subsets of $X$.

With this understood, on the set of maps from a space $X$ to another
space $Y$ one should introduce the order relation called
``specialization'', deduced from that of $Y$, namely
$$
f\leq g\Longleftrightarrow f(x)\leq g(x)\mbox{ for every } x\in
X{\quoi}.
$$
With these conventions, one has the following result:

\subsubsection{}\etiquette[4.2.3]{IV.subsubsec4.2.3}
\emph{Let $X$, $Y$ be two topological spaces, with $Y$ sober, and let
$f$ and $g$ be two continuous maps from $X$ to $Y$. Then there is at
most one morphism from $\Top(f)$ to $\Top(g)$, and such a morphism
exists if and only if $g$ is a specialization of $f$. Finally, every
morphism of topoi $\Top(X)\rightarrow \Top(Y)$ is isomorphic to a
morphism of the form $\Top(f)$, where $f:X\rightarrow Y$ is a
continuous map (uniquely determined by the first assertion).}

If\pageoriginale, for an ordered set $I$, one defines the associated
category $\cat(I)$ by declaring that, for $i\geq j$, there is exactly
one arrow from $i$ to $j$, the preceding result may be summarized by
saying that there is a canonical equivalence of categories
$$
\cat(\Hom_{(\esp)}(X,Y))\xrightarrow{\approx}\Hhomtop(\Top(X),\Top(Y))
$$
(where the second member is defined in \hyperref[IV.subsec3.2]{3.2}).

The following is a formal consequence of these results:

\begin{enonce*}{Corollary 4.2.4}\etiquette[4.2.4]{IV.cor4.2.4}
\begin{enumerate}
\renewcommand{\theenumi}{\alph{enumi}}
\renewcommand{\labelenumi}{\theenumi)}
\item Let $f:X\rightarrow Y$ be a continuous map. Then
$\Top(f):\Top(X)\rightarrow \Top(Y)$ is an equivalence of topoi if
and only if $f_{\sob}:X_{\sob}\rightarrow Y_{\sob}$ is a
homeomorphism (and hence, when $X$ and $Y$ are sober, if and only if
$f$ is a homeomorphism).

\item Let $X$ and $Y$ be topological spaces. The topoi $\Top(X)$ and
$\Top(Y)$ are equivalent if and only if $X_{\sob}$ and $Y_{\sob}$ are
homeomorphic (and hence, if $X$ and $Y$ are sober, if and only if $X$
and $Y$ are homeomorphic).
\end{enumerate}
\end{enonce*}

\subsection{Morphisms into the terminal topos: constant objects of a topos,
section functors}\etiquette[4.3]{IV.subsec4.3}

Let $P$ (the initial letter of ``point'') denote the standard terminal
topos, \ie $P=(\Ens)$ (\hyperref[IV.subsec2.2]{2.2}). Let $E$ be any
topos. We shall see that \emph{up to unique isomorphism, there exists a
unique morphism of topoi}
\hypertarget{sga4.exp.IV.equation.4.3.1}{}$$
f:E\rightarrow P
$$
or, more precisely, that \emph{the category $\Hhomtop(E,P)$ is
equivalent to the one-point category}: between any two objects of this
category there is therefore a\pageoriginale unique arrow, and it is an
isomorphism. (This justifies to some extent the name ``terminal
topos''.) For this, recall \hyperlink{sga4.exp.IV.subsubsection.3.2.1}{(3.2.1)} that $\Hhomtop(E,P)$ is equivalent
to the full subcategory of $\Hhom(P,E)^{\circ}$ formed by the functors
$$
f^{\ast}:P=(\Ens)\longrightarrow E
$$
that commute with $\varinjlim$ and are left exact. Let $e$ be a
singleton set. Every set $X$ can then be written canonically as the
``coproduct of $X$ copies of $e$'', from which it follows that the
category of functors $g:(\Ens)\rightarrow E$ that commute with
$\varinjlim$ is equivalent to the category $E$, by associating with
each $g$ the object $g(e)$ of $E$. In terms of $T=g(e)$, the functor
$g$ is reconstructed, up to unique isomorphism, by
$g(I)=T\times I$, where one sets
$T\times I=\amalg_{i\in I}T_i$, with $T_i=T$ for every $i\in I$. The
fact that the functor of $I$ so defined by $T$ does indeed commute with
$\varinjlim$ follows from the fact that it is manifestly
\sisi{right}{left} adjoint to the functor
$X\mapsto \Hom(T,X)$ from $E$ to $(\Ens)$. This said, for $g$ to be
left exact, it is plainly necessary that $g(e)=T$ be a terminal object
of $E$ (since $e$ is a terminal object of $(\Ens)$), and it follows
easily from the fact that in $E$ ``coproducts are universal''
(\hyperref[IV.subsubsec1.1.2]{1.1.2} b)) that this condition is also
sufficient. Thus the category of inverse-image functors $f^{\ast}$ is
equivalent to the category of terminal objects of $E$, which is itself
plainly equivalent to the terminal category.

It follows from the foregoing that the choice of a morphism \hyperlink{sga4.exp.IV.equation.4.3.1}{(4.3.1)}
is essentially equivalent to the choice of a terminal object of $E$,
say $e_E$. In terms of this object, one then has canonical
isomorphisms of functors
\begin{equation*}
 f^{\ast}(I)\simeq e_E\times
 I=\mbox{ coproduct of } I \mbox{ copies of } e_E \quad \mbox{ for } I\in \ob
 (\Ens),\tag{4.3.2}
\end{equation*}\etiquette[4.3.2]{IV.eq4.3.2}
and\pageoriginale
\begin{equation*}
f_{\ast}(X)=\Hom(e_E,X) \quad \quad \mbox{ for } X\in \Ob E.\tag{4.3.3}
\end{equation*}\etiquette[4.3.3]{IV.eq4.3.3}


\setcounter{subsubsection}{3}
\subsubsection{}\etiquette[4.3.4]{IV.subsubsec4.3.4}
The preceding two functors will play an important role below. For a
set $I$, the object $f^{\ast}(I)=e_E\times I$ of
\eqhyperref[IV.eq4.3.2]{4.3.2} is called the \emph{constant object
with value} $I$ in $E$ (or, when $E$ is presented as a category
$C^{\widetilde{\phantom{,}}}$ in terms of a site $C$, the
\emph{constant sheaf with value $I$ on $C$}). It will also often be
denoted by $I_E$, or by $I_C$ when $E$ is defined by the site $C$. The
fact that $I\mapsto I_E$ is the inverse-image functor of a morphism of
topoi specifies its exactness properties, which imply in particular
that this functor respects all the usual kinds of algebraic structure,
transforming a group into a group object of $E$, and so on
(\hyperref[IV.subsubsec3.1.2]{3.1.2}). If, for example, one has a
group $G$ of $E$, one says that it is a \emph{constant group} (or,
where appropriate, a \emph{constant sheaf of groups}) if it is
isomorphic to a group of the form $G_{\circ E}$, where $G_{\circ}$ is
an ordinary group. The same terminology is used for every other kind
of ``algebraic'' structure, in the sense specified (more or less) in
\hyperref[IV.subsubsec3.1.2]{3.1.2}.

% END INCLUDED SGA4_English_current_buildable_source_20260906/source/04/drafts/SGA4_IV_LINES_1320_1659_English.texfrag


% BEGIN INCLUDED SGA4_English_current_buildable_source_20260906/source/04/drafts/SGA4_IV_LINES_1660_1806_English.texfrag
\subsubsection{}\etiquette[4.3.5]{IV.subsubsec4.3.5}
Notice that the functor $I\mapsto I_E$ is not necessarily fully
faithful (or even faithful: take $E$ to be the ``empty topos''
(\hyperref[IV.subsec2.2]{2.2})); thus a constant object of $E$ does
not in general determine, up to unique isomorphism, the set $I$ from
which it arises. We say that $E$ is $0$-\emph{acyclic}, or
\emph{connected and nonempty}, if the functor $I\mapsto I_E$ is fully
faithful. By the general properties of adjoint functors, this amounts
to saying that the adjunction morphism
$$
I\longrightarrow f_{\ast}(f^{\ast}(I))=f_{\ast}(I_E)=\Hom(e_E, I_E)
$$
is\pageoriginale an isomorphism (so that the functor
\eqhyperref[IV.eq4.3.3]{4.3.3} recovers the ``value'' of a constant
object of $E$). It is easily checked that this holds if and only if
$e_E$ is not the initial object $\phi_E$ of $E$, \ie if $E$ is not
an ``empty topos'' (which expresses the \emph{faithfulness} of the
functor $I\mapsto I_E$\footnote[1]{or equivalently the fact that this
functor is conservative (\hyperrefext[I][I.cor6.3]{6.3}).}), \emph{and}
if $e_E$ is a \emph{connected} object of $E$, \ie is not the sum of
two objects of $E$ which are not ``empty'' (\ie which are not initial
objects of $E$).

\subsubsection{}\etiquette[4.3.6]{IV.subsubsec4.3.6}
The functor \eqhyperref[IV.eq4.3.3]{4.3.3} is also often called the
\emph{sections functor} and is denoted by $\Gamma_E$, $\Gamma(E,-)$,
or simply $\Gamma$:
\hypertarget{sga4.exp.IV.equation.4.3.6.1}{}\begin{equation*}
\Hom(e_E,X)=\Gamma_E(X)=\Gamma(E,X)=\Gamma(X){\quoi}.\tag{4.3.6.1}
\end{equation*}
It commutes with arbitrary projective limits, but is not in general
right exact; its derived functors (on abelian group objects) will be
studied in the next exposé.

\subsection{Morphisms from the ``empty topos''}\etiquette[4.4]{IV.subsec4.4}

Let $\phi_{\text{top}}$ be an empty topos, corresponding to a category
of sheaves $\Phi$ equivalent to the terminal category
(\hyperref[IV.subsec2.2]{2.2}). Let $E$ be a topos. The category of
functors from $E$ to $\Phi$ is clearly equivalent to the one-point
category, and every such functor commutes with inductive and
projective limits (without any merit in this), and hence may be
regarded as an inverse-image functor for a morphism of topoi
$\phi_{\text{top}}\rightarrow E$. It follows that the \emph{category
$\Hhomtop(\phi_{\text{top}},E)$ is equivalent to the one-point
category}, and in particular that \emph{there exists, up to
unique\pageoriginale isomorphism, exactly one morphism of topoi}
\begin{equation*}
\phi_{\text{top}}\longrightarrow E{\quoi}.\tag{4.4.1}
\end{equation*}\etiquette[4.4.1]{IV.eq4.4.1}
This justifies to some extent the terminology ``initial topos''
introduced in \hyperref[IV.subsec2.2]{2.2}.

One may also determine the morphisms of topoi
\begin{equation*}
E\longrightarrow \phi_{\text{top}}{\quoi}\ ;\tag{4.4.2}
\end{equation*}\etiquette[4.4.2]{IV.eq4.4.2}
one checks at once that such a morphism exists if and only if the
initial object of $E$ is also a terminal object, \ie if and only if
$E$ itself is an ``empty topos,'' and that in this case the category
$\Hhomtop(E,\phi_{\text{top}})$ is again equivalent to the one-point
category. The unique morphism \eqhyperref[IV.eq4.4.2]{4.4.2} (up to
isomorphism) is then an equivalence of topoi.

\subsection{The classifying topos $B_G$ for a varying group $G$}\etiquette[4.5]{IV.subsec4.5}

\subsubsection{}\etiquette[4.5.1]{IV.subsubsec4.5.1}
Let $E$ be a topos, and let
$$
f:G\longrightarrow H
$$
be a morphism of group objects in $E$. It induces a ``restriction of
the group of operators'' functor
$$
f^{\ast}:B_H\longrightarrow B_G{\quoi},
$$
where the notation is that of \hyperref[IV.subsec2.4]{2.4}. It is
immediate that this functor commutes with inductive and projective
limits; a fortiori, it may be interpreted as an inverse-image
functor associated with a morphism of topoi
$$
B_f\mbox{ or } f:B_G\longrightarrow B_H{\quoi}.
$$
The\pageoriginale corresponding direct-image functor
$$
f_{\ast}:B_G\longrightarrow B_H
$$
is readily made explicit by the formula
$$
f_{\ast}(X)=\Hhom_G(H_s,X){\quoi},
$$
where $X$ is an object of \sisi{$X$}{$E$} with a left action of $G$,
where $H_s$ is $H$ regarded as endowed with the left actions by $G$
induced by \sisi{$F$}{$f$}, and where $\Hhom_G$ denotes the evident
subobject of the object $\Hhom$ defined below
(\hyperref[IV.subsec10.2]{10.2}); the left action of $H$ on
$\Hhom_G(H_s,X)$ is induced by the right actions of $H$ on $H_s$ by
\emph{right} translations.

Since the inverse-image functor $f^{\ast}$ commutes with arbitrary
\sisi{$\varinjlim$}{$\varprojlim$} (and not only with finite
\sisi{$\varinjlim$}{$\varprojlim$}), it is itself the right adjoint of
a functor
$$
f_{!}:B_G\longrightarrow B_H{\quoi},
$$
so that there is a sequence of three adjoint functors as in
\hyperref[IV.subsubsec3.1.3]{3.1.3}:
$$
f_{!},f^{\ast}, f_{\ast}{\quoi}.
$$
The functor $f_{!}$ is readily made explicit by the formula
$$
f_{!}(X)=H\times {}^GX{\quoi},
$$
where the right-hand side denotes the ``contracted product'' induced
by the actions of $G$ on $X$ (on the left) and on $H$ (on the right,
via right translations and $f$), defined as the quotient of
$H\times X$ by the action of $G$ given by
$$
g\circ (h,x)=(hg^{-1}, gx){\quoi}.
$$

The functor $f_{!}$, being a left adjoint, clearly commutes with
inductive limits, but is not in general left exact (\ie
it\pageoriginale cannot in turn be regarded as an inverse-image
functor for a morphism of topoi $B_H\rightarrow B_G$). In fact, one
checks easily that it can be left exact only if $f:G\rightarrow H$
is an isomorphism. Likewise, the functor $f_{\ast}$, which commutes
with projective limits, is not in general right exact and, a
fortiori, does not in general admit a right adjoint. At least when
$E$ is the punctual topos, \ie when $G$ and $H$ are ordinary groups,
$f_{\ast}$ is right exact only if $f$ is an isomorphism.

Given a second morphism of groups $g:H\rightarrow K$, one obtains,
as in \hyperref[IV.subsubsec4.1.1]{4.1.1}, transitivity up to
canonical isomorphism, so that, subject to the usual reservation, the
classifying topos $B_G$ may be regarded as depending
\emph{functorially} on the group $G$. We leave it to the reader to
generalize this functorial behavior to the case where both $G$
\emph{and} the topos $E$ vary.


% END INCLUDED SGA4_English_current_buildable_source_20260906/source/04/drafts/SGA4_IV_LINES_1660_1806_English.texfrag


% BEGIN INCLUDED SGA4_English_current_buildable_source_20260906/source/04/drafts/SGA4_IV_LINES_1807_1981_English.texfrag
\subsubsection{The topos $B_{𝒢}$ for a varying pro-group
  $𝒢$}\etiquette[4.5.2]{IV.subsubsec4.5.2}

We leave it to the reader to make precise the covariant character of
the topos $B_{𝒢}$ (\hyperref[IV.subsec2.7]{2.7}) with respect to $𝒢$,
by following the account given in \hyperref[IV.subsubsec4.5.1]{4.5.1}.
One should note, however, that in the case of a morphism
$f:𝒢\rightarrow ℋ$ of pro-groups which are not essentially
constant, the corresponding morphism of topoi
$B_{𝒢}\rightarrow B_{ℋ}$ does not in general permit the definition
of a functor $f_{!}$ (whose right adjoint would be the functor
$f^{\ast}$ restricting the pro-group of operators). Suppose, to
simplify the statement, that $𝒢$ and $ℋ$ are defined by
\emph{profinite} groups $G$ and $H$. It is then easily seen that
$f_{!}$ exists if and only if the image of the morphism in question
$f:G\rightarrow H$ has finite index in $H$, and in\pageoriginale
that case $f_{!}$ is given by the same formula as in
\hyperref[IV.subsec4.5]{4.5}.

\subsection{The topos $\widehat{C}$ for a varying category $C$}\etiquette[4.6]{IV.subsec4.6}

\subsubsection{}\etiquette[4.6.1]{IV.subsubsec4.6.1}
Let
$$
f:C\longrightarrow C'
$$
be a functor from a category $C\in 𝒰$ to another category $C'$, and
hence let
$$
f^{\ast}:\widehat{C}'\longrightarrow \widehat{C}\quav
f^{\ast}(F)=F\circ f{\quoi}.
$$
It is immediate that this functor commutes with projective and
inductive limits; a fortiori, it may be regarded as the inverse-image
functor for a morphism of topoi
$$
\widehat{f}\mbox{ or } f:\widehat{C}\rightarrow \widehat{C}'{\quoi}.
$$
The corresponding direct-image functor
$$
f_{\ast}:\widehat{C}\longrightarrow \widehat{C}'
$$
is none other than the functor likewise denoted $f_{\ast}$ in I \hyperref[I.prop5.1]{5.1}.
Moreover (as was to be expected from the fact that $f^{\ast}$ also
commutes with arbitrary $\varprojlim$), $f^{\ast}$ also admits a left
adjoint
$$
f_{!}:\widehat{C}\longrightarrow \widehat{C}'{\quoi},
$$
(which was likewise denoted $f_{!}$ in
\hyperrefext[I][I.prop5.1]{5.1}). Thus one obtains a sequence of three
adjoint functors
$$
f_{!}, f^{\ast}, f_{\ast}{\quoi},
$$
the first of which is moreover an extension of $f:C\rightarrow C'$
to the categories of presheaves (for the usual embeddings of $C$ and
$C'$ in the categories $\widehat{C}$ and $\widehat{C}'$). One should
note that $f_{!}$ is not in general left exact, nor is $f_{\ast}$ in
general right exact; this removes any ambiguity about the direction
of\pageoriginale variance of the topos $\widehat{C}$ associated with
the varying category $C$: this topos is a \emph{covariant} functor of
$C$, subject to the usual reservation arising from the transitivity
isomorphisms (\cf \hyperref[IV.subsubsec4.1.1]{4.1.1}). When the sets
of objects of $C$ and $C'$ each consist of a single element, so that
$C$ and $C'$ identify with monoids $G$ and $G'$, the corresponding
topoi are the classifying topoi $B_{G^{\circ}}$ and
$B_{{G'}^{\circ}}$, and one recovers their variance with respect to
the varying $G$, already encountered from other viewpoints in
\hyperref[IV.subsubsec4.1.2]{4.1.2} and
\hyperref[IV.subsec4.5]{4.5} (where the restriction to groups rather
than monoids was inessential).

\subsubsection{}\etiquette[4.6.2]{IV.subsubsec4.6.2}
The 2-functorial dependence of the topos $\widehat{C}$ on $C$ can be
made more precise by introducing, for two categories $C,C'\in 𝒰$, a
canonical functor
\begin{equation*}
\Hhom(C,C')\longrightarrow
\Hhomtop(\widehat{C},\widehat{C}'){\quoi}.\tag{4.6.2.1}
\end{equation*}\etiquette[4.6.2.1]{IV.eq4.6.2.1}
It remains to define this functor on arrows. For this, note that
$f\mapsto f^{\ast}$ identifies (up to equivalence of categories) the
right-hand side with a full subcategory of
$\Hhom(\widehat{C}',\widehat{C})^{\circ}$ \hyperlink{sga4.exp.IV.subsubsection.3.2.1}{(3.2.1)}. If
$f,g:C\rightarrow C'$ are two functors, every morphism of functors
$u:f\rightarrow g$ defines, for $F\in\ob\widehat{C}'$, a morphism of
functors $F\circ g\rightarrow F\circ f$ (since $F$ is a
\emph{contra}variant functor), hence a morphism
$g^{\ast}\rightarrow f^{\ast}$ and therefore a morphism
$\widehat{f}\rightarrow\widehat{g}$, as asserted.

When $C$ is the one-point category, $\widehat{C}$ is the punctual
topos (\hyperref[IV.subsec2.2]{2.2}), denoted $P$, and
\eqhyperref[IV.eq4.6.2.1]{4.6.2.1} is interpreted as a natural functor
\begin{equation*}
C'\longrightarrow \Hhomtop(P,\widehat{C}') \mathop{=}^{\dfn}
\Ppoints(\widehat{C}')\tag{4.6.2.2}
\end{equation*}\etiquette[4.6.2.2]{IV.eq4.6.2.2}
from $C'$ to the ``\emph{category of points}'' of $\widehat{C}'$,
which will be studied in 6. This\pageoriginale functor is not
necessarily an equivalence of categories \hyperlink{sga4.exp.IV.subsubsection.7.3.3}{(7.3.3)}; a fortiori,
\eqhyperref[IV.eq4.6.2.1]{4.6.2.1} is not necessarily an equivalence
of categories.

On the other hand, the functor \eqhyperref[IV.eq4.6.2.1]{4.6.2.1} is
always \emph{fully faithful}. To see this, note that if
$f,g:E\rightarrow E'$ are two morphisms of topoi for which $f_{!}$
and $g_{!}$ are defined (\hyperref[IV.subsubsec3.1.3]{3.1.3}), then
the theory of adjoint functors gives canonical isomorphisms
$$
\Hom(f_{!},g_{!})\simeq \Hom(g^{\ast}, f^{\ast})\simeq
\Hom(f_{\ast}, g_{\ast}) \mathop{=}^{\dfn} \Hom(f,g){\quoi}.
$$
Applying this to functors of the form $\widehat{f}$ and $\widehat{g}$
associated with $f,g:C\rightarrow C'$, one obtains the asserted
result, taking into account that the natural extension map
$\Hom(f,g)\rightarrow \Hom(f_{!},g_{!})$ is bijective (I \hyperref[I.prop7.8]{7.8}).

\subsubsection{}\etiquette[4.6.3]{IV.subsubsec4.6.3}
One may ask when the functor \eqhyperref[IV.eq4.6.2.1]{4.6.2.1} is an
equivalence of categories, \ie when it is essentially surjective.
This is a special case of the problem of determining all morphisms of
topoi $\widehat{C}\rightarrow\widehat{C}'$. More generally, if $C$ is
a category in $𝒰$ and $E$ is a topos, one may seek to determine the
morphisms of topoi
$$
f:\widehat{C}\longrightarrow E{\quoi}.
$$
The category of these morphisms is equivalent to the opposite of the
category of functors $f^{\ast}:E\rightarrow\widehat{C}$ commuting
with arbitrary inductive limits and finite projective limits. A
functor $E\rightarrow\widehat{C}$ may be interpreted as a functor
$E\times C^{\circ}\rightarrow(𝒰\text{-}\Ens)=(\Ens)$, or equivalently
as a functor $F:C^{\circ}\rightarrow\Hhom(E,(\Ens))$. The exactness
property in question is then expressed by requiring that, for every
object $X$ of $C$, the functor $F(X):E\rightarrow(\Ens)$ commute with
arbitrary inductive limits and finite projective limits (or, as we
shall say in 6, that $F(X)$ is a ``\emph{fiber\pageoriginale
functor}'' on $E$). This amounts to saying that $F(X)$ is the
inverse-image functor for a morphism of topoi $P\rightarrow E$, and
one therefore obtains finally a canonical equivalence of categories
\begin{equation*}
\Hhomtop(\widehat{C}, E) \xrightarrow{\approx} \Hhom(C,
\Ppoints(E)){\quoi},\tag{4.6.3.1}
\end{equation*}\etiquette[4.6.3.1]{IV.eq4.6.3.1}
where, for brevity, we write
$$
\Ppoints(E)=\Hhomtop(P,E)
$$
for the category of points of the topos $E$.

When $E$ is of the form $\widehat{C}'$, it follows immediately from
the definitions that the composite of
\eqhyperref[IV.eq4.6.2.1]{4.6.2.1} and the preceding equivalence
\eqhyperref[IV.eq4.6.3.1]{4.6.3.1} is the functor
\begin{equation*}
\Hhom(C,C')\longrightarrow
\Hhom(C,\Ppoints(\widehat{C}'))\tag{4.6.3.2}
\end{equation*}\etiquette[4.6.3.2]{IV.eq4.6.3.2}
defined by $F\mapsto i\circ F$, where
$i:C'\rightarrow\Ppoints(\widehat{C}')$ is the canonical embedding
\eqhyperref[IV.eq4.6.2.2]{4.6.2.2}. It follows at once that
\emph{for} \eqhyperref[IV.eq4.6.2.1]{4.6.2.1} \emph{to be an
equivalence, it is necessary and sufficient that $C$ be empty or
that} \eqhyperref[IV.eq4.6.2.2]{4.6.2.2} \emph{be essentially
surjective} (and hence an equivalence of categories). This last
condition on $C'$ is satisfied in certain interesting cases, notably
when $C'$ is the one-object category defined by a group $G$
(\hyperref[IV.subsubsec7.2.5]{7.2.5}).



% END INCLUDED SGA4_English_current_buildable_source_20260906/source/04/drafts/SGA4_IV_LINES_1807_1981_English.texfrag


% BEGIN INCLUDED SGA4_English_current_buildable_source_20260906/source/04/drafts/SGA4_IV_LINES_1982_2188_English.texfrag
\begin{enonce*}[remark]{Remark 4.6.4}\etiquette[4.6.4]{IV.rem4.6.4}\nde{For
a recent perspective on the notion of the homotopy type of a topos, see
[15] $\S$ \hyperref[IV.subsubsec7.1.6]{7.1.6}.}
The association of a topos $\widehat{C}$ with an arbitrary category
$C$ suggests that a category $C$ has invariants of a topological
nature (cohomology groups, homotopy groups, \etc), just as a topos
does. The cohomology groups of $\widehat{C}$ with coefficients in an
abelian group object $F$ (in the general sense studied in V) are
precisely the values of the derived functors
$\varprojlim^{(n)}$ of the functor $\varprojlim$, already familiar
to topologists.\pageoriginale J. L. Verdier and (independently)
D. G. Quillen verified that, when one restricts to constant, or more
generally locally constant, coefficients, these cohomology groups
identify with the cohomology groups of the semi-simplicial set
$\Nerf(C)$ canonically associated with $C$ [5, ${\rm n^{\circ}}$ 212,
prop. 4.1], and moreover that, up to isomorphism in the ``homotopy
category'' of [2], every semi-simplicial set can be obtained from a
category $C$. With a suitable notion of homotopy type for topoi,
which we do not specify here, one may say that the semi-simplicial
homotopy types of topologists are precisely the homotopy types of
topoi of the special form $\widehat{C}$, and more generally of topoi
$E$ in which every object admits a covering that refines every other
one (\hyperref[IV.subsec2.6]{2.6})
\footnote[1]{and, more generally still, of topoi that are ``locally
$\infty$-connected'' in an evident sense which we leave it to the
reader to make precise.}. In the absence of this condition on $E$,
its homotopy type can at best be expressed by a suitable
\emph{projective system} of semi-simplicial sets [1].
\end{enonce*}

\subsection{The topos $C^{\sim}$ for a variable site $C$
(cocontinuous functors)}\etiquette[4.7]{IV.subsec4.7}

Let
$$
f:C\longrightarrow C'
$$
be a cocontinuous functor (\hyperrefext[III][III.def2.1]{2.1}) between
sites in $𝒰$, \ie a functor such that the functor
$\widehat{f}_{\ast}$ from $\widehat{C}$ to $\widehat{C}'$
(\hyperref[IV.subsubsec4.6.1]{4.6.1}) maps $C^{\sim}$ into
$C'^{\sim}$, that is, induces a functor
\begin{equation*}
\widetilde{f}_{\ast}:\widetilde{C}\longrightarrow
{C'}^{\widetilde{\phantom{,}}}\tag{4.7.1}
\end{equation*}\etiquette[4.7.1]{IV.eq4.7.1}
making the following diagram of functors commutative:
\[\vcenter{
\xymatrix@C=2.5cm@R=1.25cm{ C^{\widetilde{\phantom{a}}}
\ar[r]^{\widetilde{f}_{\ast}} \ar[d]_{\sisi{i}{i_{\ast}}}&
C'^{\widetilde{\phantom{a}}} \ar[d]^{\sisi{i_{\ast}}{i_{\ast}'}}\\
\widehat{C}  \ar[r]^{\widehat{f}_{\ast}}       & \widehat{C}'\quad,
}}\leqno{(4.7.2)}
\]\etiquette[4.7.2]{IV.eq4.7.2}
where\pageoriginale \sisi{$i$}{$i_{\ast}$} and
\sisi{$i'$}{$i_{\ast}'$} are the inclusion functors. We saw
(\hyperrefext[III][III.prop2.3]{2.3}) that the functor
$\widetilde{f}_{\ast}$ then has a left adjoint
$\widetilde{f}^{\ast}$, and that the latter is left exact. In other
words, $\widetilde{f}_{\ast}$ is the direct-image functor associated
with a morphism of topoi
$$
\widetilde{f}\mbox{ or } f:C^{\sim}\longrightarrow
{C'}^{\widetilde{\phantom{,}}}{\quoi},
$$
the corresponding inverse-image functor being, of course,
$\widetilde{f}^{\ast}$. Taking left adjoints of the functors involved,
the commutative diagram \eqhyperref[IV.eq4.7.2]{4.7.2} also gives a
diagram commutative up to isomorphism:
\[\vcenter{
\xymatrix@C=2.5cm@R=1.5cm{
C^{\widetilde{\phantom{a}}} &  C'^{\widetilde{\phantom{a}}} \ar[l]_{\widetilde{f}^{\ast}}\\
\widehat{C}  \ar[u]^{\underline{a}}  & \widehat{C}'
\ar[l]_{\widehat{f}^{\ast}}\ar[u]_{\underline{a}'}\quad,
}}\leqno{(4.7.3)}
\]\etiquette[4.7.3]{IV.eq4.7.3}
where $\underline{a}$ and $\underline{a}'$ are the ``associated
sheaf'' functors. This diagram immediately recovers the formula
(\hyperrefext[III][III.prop2.3]{2.3})
$$
\widetilde{f}^{\ast}=\underline{a}\widehat{f}^{\ast} i'{\quoi}.
$$

The transitivity property for the morphisms of topoi
$\widehat{f}:\widehat{C}\rightarrow \widehat{C}'$ implies the same
property for the morphisms of topoi
$\widetilde{f}:C^{\widetilde{\phantom{,}}}\rightarrow
{C'}^{\widetilde{\phantom{,}}}$ associated with cocontinuous
functors. Thus one may say that \emph{the topos
$C^{\widetilde{\phantom{,}}}$ varies functorially with $C$} in a
\emph{covariant} manner, when cocontinuous functors are taken as the
``morphisms'' of sites.

In\pageoriginale all the cases encountered thus far, the
cocontinuous functor $f$ used is also continuous, that is
(\hyperrefext[III][III.def1.1]{1.1}), it ``extends'' to a functor
$$
\widetilde{f}_{!}:C^{\widetilde{\phantom{,}}}\longrightarrow
{C'}^{\widetilde{\phantom{,}}}
$$
commuting with inductive limits, which is left adjoint to
$\widetilde{f}^{\ast}$, so that there is a sequence of three adjoint
functors
$$
\widetilde{f}_{!}, \widetilde{f}^{\ast},
\widetilde{f}_{\ast}{\quoi}.
$$
One should note that the functor $\widetilde{f}_{!}$ is not in
general left exact, nor is $\widetilde{f}_{\ast}$ right exact. This
removes any ambiguity concerning the direction of the variance of
the topos $C^{\widetilde{\phantom{,}}}$, when $C$ varies through
functors assumed only to be cocontinuous, or even continuous and
cocontinuous.

\begin{enonce*}[remark]{Remark 4.7.4}\etiquette[4.7.4]{IV.rem4.7.4}
Given a morphism of topoi
$$
F:E\longrightarrow E'{\quoi},
$$
for there to exist a functor $F_{!}$ left adjoint to $F^{\ast}$, it
is necessary and sufficient that one be able to ``realize'' (up to
equivalence) $E$ and $E'$ in the form
$C^{\widetilde{\phantom{,}}}$ and ${C'}^{\widetilde{\phantom{,}}}$,
for two sites $C,C'\in 𝒰$, and that one be able to find a
\emph{continuous and cocontinuous} functor $f:C\rightarrow C'$ such
that $F$ identifies with $\widetilde{f}$. This is clearly sufficient;
for necessity, it is enough to take for $C$ and $C'$ small full
generating subcategories of $E$ and $E'$, respectively, such that
$$
F_{!}(\ob C)\subset \ob C'{\quoi},
$$
endowed with the topologies induced by those of $E$ and $E'$, and to
take for $f$ the functor induced by $F_{!}$ (\cf
\hyperref[IV.cor1.2.1]{1.2.1}). Recall
(\hyperref[IV.rem1.8]{1.8}) that the existence of $F_{!}$ also means
that $F^{\ast}$ commutes with projective limits, or,\pageoriginale
equivalently here since this functor is left exact, that it commutes
with products.
\end{enonce*}

\setcounter{subsubsection}{4}
\subsubsection{}\etiquette[4.7.5]{IV.subsubsec4.7.5}
If one asks which morphisms of topoi $F:E\rightarrow E'$ can be
realized by a cocontinuous (not necessarily continuous) functor of
sites, one likewise finds the following answer: the set of objects
$X$ of $E$ such that the functor
$X'\mapsto \Hom(X,F^{\ast}(X'))$ from $E'$ to $(\Ens)$ commutes with
projective limits (or, equivalently, with products) must be a
\emph{generating family of $E$.}\nde{This statement does not appear
to be correct: the condition expressed is equivalent to the assertion
that $F^{\ast}$ commutes with arbitrary products, or equivalently that
$F^{\ast}$ has a left adjoint $F_{!}$. For every site $C$, the
identity functor $F=id:C\longrightarrow C$ is cocontinuous when the
term on the right is endowed with the trivial topology. The functor
$F^{\ast}:\hat{C}\longrightarrow\widetilde{C}$ is then the associated
sheaf functor, and does not in general commute with arbitrary
products.}

\subsection{The morphism of topoi $\widetilde{C}\rightarrow
\widehat{C}$ for a site $C$}\etiquette[4.8]{IV.subsec4.8}

Let $C$ be a small site, with which the two topoi
$C^{\widetilde{\phantom{,}}}$ and $\widehat{C}$ are therefore
associated (the latter does not depend on the topology placed on
$C$). In II \hyperref[II.theo3.4]{3.4} we defined the ``associated sheaf'' functor
$$
\underline{a}:\widehat{C}\longrightarrow
C^{\widetilde{\phantom{,}}}{\quoi},
$$
and proved that it is left exact and commutes with inductive limits.
It is therefore the inverse-image functor associated with a morphism
of topoi
$$
p:\widetilde{C}\longrightarrow \widehat{C},\;
p^{\ast}=\underline{a}{\quoi},
$$
the corresponding direct-image functor being the canonical inclusion
$$
i=p_{\ast}:C^{\widetilde{\phantom{,}}}\longrightarrow
\widehat{C}{\quoi}.
$$
It is well known that the latter functor is not in general right
exact (its derived functors on abelian objects give rise to the
cohomology presheaves $\Xx^n(F)$ of V 2), and that $\underline{a}$
does not in general commute with arbitrary $\varprojlim$. This
removes any ambiguity concerning the direction of the ``natural''
morphism of topoi between $\widehat{C}$ and
$C^{\widetilde{\phantom{,}}}$.

Given a cocontinuous functor
$$
f:C\longrightarrow C'
$$
of\pageoriginale sites, one obtains a diagram of morphisms of topoi
\[
\xymatrix@C=2.25cm@R=1.25cm{ C^{\widetilde{\phantom{a}}} \ar[d]_{\sisi{p}{p_{\ast}}}
\ar[r]^{\widetilde{f}} &
C'^{\widetilde{\phantom{a}}} \ar[d]^{\sisi{p'}{p'_{\ast}}}\\
\widehat{C} \ar[r]^{\widehat{f}}&  \widehat{C}' }
\]
which is commutative up to canonical isomorphism: this is precisely
what the commutativity of the diagram of functors
\eqhyperref[IV.eq4.7.2]{4.7.2} expresses. Thus one may say that the
morphism of topoi
$p:C^{\widetilde{\phantom{,}}}\rightarrow \widehat{C}$ is functorial
in $C$, when $C$ is varied through \emph{cocontinuous} functors
between sites.


% END INCLUDED SGA4_English_current_buildable_source_20260906/source/04/drafts/SGA4_IV_LINES_1982_2188_English.texfrag


% BEGIN INCLUDED SGA4_English_current_buildable_source_20260906/source/04/drafts/SGA4_IV_LINES_2189_2354_English.texfrag
\subsection{Effect of a continuous functor of sites. Morphisms of sites}\etiquette[4.9]{IV.subsec4.9}

\subsubsection{}\etiquette[4.9.1]{IV.subsubsec4.9.1}
Let
$$
f:C\longrightarrow C'
$$
be a continuous functor from $C$ to $C'$, \ie such that the functor
$\widehat{f}^{\ast}:\widehat{C}'\rightarrow\widehat{C}$ maps
${C'}^{\widetilde{\phantom{,}}}$ into
$C^{\widetilde{\phantom{,}}}$, and hence induces a functor
$$
f_s:{C'}^{\widetilde{\phantom{,}}}\longrightarrow
C^{\widetilde{\phantom{,}}}{\quoi}.
$$
We saw (\hyperrefext[III][III.prop1.2]{1.2}) that the latter admits a
left adjoint
$$
f^s:C^{\widetilde{\phantom{,}}}\longrightarrow
{C'}^{\widetilde{\phantom{,}}}{\quoi},
$$
which moreover ``extends'' $f$ in the evident sense. One should note
that in general $f^s$ is not \sisi{right}{left} exact (even if $f$ is
also cocontinuous), nor does $f_s$ commute with inductive limits, so
that the datum of $f$ does not, without further hypotheses, describe
a morphism of topoi in either direction between
$C^{\widetilde{\phantom{,}}}$ and
${C'}^{\widetilde{\phantom{,}}}$. The case where $f$ is cocontinuous,
\ie where $f_s$ commutes with inductive limits and may\pageoriginale
therefore be regarded as an inverse-image functor for a morphism of
topoi $\widetilde{f}:C^{\widetilde{\phantom{,}}}\rightarrow
{C'}^{\widetilde{\phantom{,}}}$, was examined in
\hyperref[IV.subsec4.7]{4.7}. We shall examine the case where the
functor $f^s$ is left exact, and may therefore be regarded as an
inverse-image functor for a morphism of topoi in the \emph{opposite
direction}:
\begin{equation*}
\Top(f)=g:{C'}^{\sim}\longrightarrow C^{\sim}{\quoi}. \tag{4.9.1.1}
\end{equation*}\etiquette[4.9.1.1]{IV.eq4.9.1.1}
Notice that \emph{we have deliberately not denoted this morphism by
the letter $\widetilde{f}$ or $f$, so as to avoid confusion with the
situation of \hyperref[IV.subsec4.7]{4.7}, in accordance with the
general recommendations of \hyperref[IV.subsubsec3.1.3]{3.1.3}.} A
continuous functor $f:C\rightarrow C'$ is sometimes called a
\emph{morphism of sites from $C'$ to $C$} (note: \emph{not from $C$
to $C'$}) if the functor $f^s$ is left exact; in other words, if
there exists a morphism of topoi \eqhyperref[IV.eq4.9.1.1]{4.9.1.1}
such that the corresponding inverse-image functor
$$
g^{\ast}:C^{\widetilde{\phantom{,}}}\longrightarrow
{C'}^{\widetilde{\phantom{,}}}
$$
\emph{extends} the functor $f$, \ie makes the following diagram of
functors commute up to isomorphism:
\[
\xymatrix@C=2cm@R=1.5cm{
C \ar[d]_{\epsilon} \ar[r]^{f}& C' \ar[d]^{{\epsilon}'}\\
C^{\widetilde{\phantom{a}}} \ar[r]^{g^{\ast}} &
C'^{\widetilde{\phantom{a}}} &, }
\]
%% widetilde + \phantom: horrible hack \XXX
where $\epsilon$ and $\epsilon'$ are the canonical functors of II
\hyperref[II.subsubsec4.4.0]{4.4.0}.

\subsubsection{}\etiquette[4.9.2]{IV.subsubsec4.9.2}
In practice, one recognizes that a continuous functor
$f:C\rightarrow C'$ is a morphism of sites from $C'$ to $C$ from the
fact that finite projective limits are representable in $C$ and that
$f$ commutes with them (\hyperrefext[III][III.prop1.3]{1.3} 5)).
Subject\pageoriginale to the indicated condition on $C$ (which is
almost always satisfied in practice), and assuming in addition that
the topology of $C'$ is coarser than the canonical topology (also
almost always satisfied), the preceding sufficient condition ($f$
left exact) for $f$ to be a morphism of sites from $C'$ to $C$ is
also necessary.

\subsubsection{}\etiquette[4.9.3]{IV.subsubsec4.9.3}
In the spirit of the preceding discussion, if $C$ and $C'$ are two
$𝒰$-sites, the category of morphisms of sites
$\Mmorsite(C',C)$ from $C$ to $C'$ should be defined as the full
subcategory of the \emph{opposite category
$\Hhom(C,C')^{\circ}$} of the category of functors from $C$ to $C'$,
formed by those functors which are morphisms of sites (from $C'$ to
$C$). In this way one obtains a canonical functor (defined up to
unique isomorphism)
$$
\Mmorsite(C',C)\longrightarrow
\Hhomtop({C'}^{\widetilde{\phantom{,}}},
C^{\widetilde{\phantom{,}}}){\quoi}.
$$

\begin{enonce*}{Proposition 4.9.4}\etiquette[4.9.4]{IV.prop4.9.4}
Let $E$ be a $𝒰$-topos and $C$ a $𝒰$-site. The functor
$f\mapsto f^{\ast}| C=f^{\ast}\circ\epsilon_C$, associating with
every morphism of topoi
$f:E\rightarrow C^{\widetilde{\phantom{,}}}$ the ``restriction'' to
$C$ of the associated inverse-image functor
$f^{\ast}:C^{\widetilde{\phantom{,}}}\rightarrow E$, induces an
equivalence of categories
$$
\Hhomtop(E,C^{\widetilde{\phantom{,}}})\xrightarrow{\approx}\Mmorsite(E,C)(\hookrightarrow
\Hhom(C,E)^{\circ}){\quoi}.
$$
When finite $\varprojlim$ are representable in $C$, the corresponding
fully faithful functor
$$
\Hhomtop(E,C^{\widetilde{\phantom{,}}})\longrightarrow
\Hhom(C,E)^{\circ}
$$
has as its essential image the functors $g:C\rightarrow E$ which are
left exact and continuous, or equivalently those which are left exact
and carry\pageoriginale covering families to covering families.
\end{enonce*}

The last assertion follows from the first by
\hyperref[IV.subsubsec4.9.2]{4.9.2}. On the other hand, it follows
from \hyperrefext[III][III.prop1.2]{1.2} iv) and IV
\hyperref[IV.theo1.2]{1.2} iii) that the functor
$G\rightarrow G\circ\epsilon$ induces an equivalence between the
category of \emph{continuous} functors from
$C^{\widetilde{\phantom{,}}}$ to $E$ and the category of
\emph{continuous} functors from $C$ to $E$, a quasi-inverse being
given by $g\rightarrow g^s$ in the notation of \loccit On the other
hand, by definition, $g$ is a morphism of sites if and only if $g^s$
is the inverse-image functor associated with a morphism of topoi
$E\rightarrow C^{\widetilde{\phantom{,}}}$, whence the conclusion by
\hyperlink{sga4.exp.IV.subsubsection.3.2.1}{3.2.1}; the ``or equivalently'' follows from
\hyperrefext[III][III.prop1.6]{1.6}.

The main point to retain from \hyperref[IV.prop4.9.4]{4.9.4} is that
(when finite $\varprojlim$ are representable in $C$) ``it amounts to
the same thing'' to give a morphism of topoi
$f:E\rightarrow C^{\widetilde{\phantom{,}}}$ or to give a functor
$g:C\rightarrow E$ which is left exact and carries covering families
to covering families.

\setcounter{subsubsection}{4}
\subsubsection{}\etiquette[4.9.5]{IV.subsubsec4.9.5}
Using the developments of \hyperref[IV.subsubsec4.9.1]{4.9.1} and
\hyperref[IV.subsubsec4.9.3]{4.9.3}, one sees as usual that, for a
varying $𝒰$-site $C$, via the notions of a morphism of sites and a
morphism between morphisms of sites just made explicit, the topos
$C^{\widetilde{\phantom{,}}}$ depends functorially (or more exactly,
2-functorially) on the site $C$. Note that, with the terminology just
introduced, $C^{\widetilde{\phantom{,}}}$ depends
\emph{covariantly} on the site $C$.

\subsubsection{}\etiquette[4.9.6]{IV.subsubsec4.9.6}
It is immediate that (unlike what happens for the notion of a
cocontinuous functor, \cf \hyperref[IV.subsec4.7]{4.7} 4)
\emph{every} morphism of topoi $F:E\rightarrow E'$ can be realized
by means of a morphism of sites $f:C\rightarrow C'$ (\ie a continuous
functor $f:C'\rightarrow C$ such that \ldots): it suffices to choose
small full generating subcategories $C$\pageoriginale and $C'$ of
$E$ and $E'$, respectively, endowed with the topologies induced by
those of $E$ and $E'$, such that
$$
F^{\ast}(\ob C')\subset \ob C{\quoi},
$$
and to take for $f$ the functor induced by $F^{\ast}$. This explains
why most morphisms of topoi encountered in practice are indeed
described by morphisms of sites (rather than by cocontinuous functors
as in \hyperref[IV.subsec4.7]{4.7}).


% END INCLUDED SGA4_English_current_buildable_source_20260906/source/04/drafts/SGA4_IV_LINES_2189_2354_English.texfrag


% BEGIN INCLUDED SGA4_English_current_buildable_source_20260906/source/04/drafts/SGA4_IV_LINES_2355_2641_English.texfrag
\subsection{Relations between the small and large topoi associated with
a topological space $X$}\etiquette[4.10]{IV.subsec4.10}

We retain the notation of \hyperref[IV.subsec2.5]{2.5}. In particular,
$𝒱$ is a universe such that $𝒰\in 𝒱$. We depart
from the convention of \hyperref[IV.subsec2.1]{2.1} by using
$\Top(X)$ to denote the topos of $𝒱$-sheaves (rather than
$𝒰$-sheaves) on $X$. Thus, we shall work with
$𝒱$-topoi and not $𝒰$-topoi. (NB it would be
possible to retain $𝒰$-topoi by adopting the appropriate
convention for the definition of $\TOP(X)$, so that the latter is a
$𝒰$-topos, \cf \hyperref[IV.subsec2.5]{2.5}). With this understood,
we shall define TWO morphisms of topoi
\begin{equation*}
\begin{cases}
 f: \Top(X)\longrightarrow \TOP(X)\\
 g: \TOP(X)\longrightarrow \Top(X)\quav gf\simeq \id_{\Top(X)}{\quoi},
\end{cases}\tag{4.10.1}
\end{equation*}\etiquette[4.10.1]{IV.eq4.10.1}
giving rise to a sequence of three adjoint functors
\begin{equation*}
g^{\ast}=f_{!}\quav g_{\ast}=f^{\ast}\quav
g^{!}=f_{\ast}{\quoi},\tag{4.10.1.1}
\end{equation*}\etiquette[4.10.1.1]{IV.eq4.10.1.1}
with the notation of \hyperref[IV.subsubsec3.1.3]{3.1.3}. We shall
successively define the\pageoriginale first two functors in this
sequence, the third then being defined as the right adjoint of the
second.

\setcounter{subsubsection}{1}
\subsubsection{}\etiquette[4.10.2]{IV.subsubsec4.10.2}
The functor
$$
g_{\ast}=f^{\ast}:\TOP(X)\longrightarrow \Top(X)
$$
is defined simply as the \emph{restriction functor} from
$(\Esp)_{/X}$ to the site $\Ouv(X)$ of open subsets of $X$; it does
indeed take sheaves to sheaves, as follows immediately from the
definition. Using underlined letters for sheaves on the large site of
$X$, for such a sheaf $\underline{F}$ we denote by
$\underline{F}_X$ its restriction to the site $\Ouv(X)$, thereby
obtaining a functor
\begin{equation*}
\Restr:\underline{F}\longmapsto \underline{F}_X:\TOP(X)\longrightarrow
\Top(X){\quoi}.\tag{4.10.2.1}
\end{equation*}\etiquette[4.10.2.1]{IV.eq4.10.2.1}
It is clear that this functor commutes with $𝒱$-projective
limits, since these are computed pointwise (one could also invoke the
existence of the left adjoint constructed in
\hyperref[IV.subsubsec4.10.4]{4.10.4} below). I claim that it also
commutes with $𝒱$-inductive limits. To see this, we shall give a
very convenient interpretation of ``large'' sheaves on $X$, \ie of
objects of $\TOP(X)$, in terms of ordinary sheaves (NB these are
$𝒱$-sheaves) on topological spaces $X'$ over~$X$.

\subsubsection{}\etiquette[4.10.3]{IV.subsubsec4.10.3}
For a large sheaf $\underline{F}$ on $X$, and for every space $X'$
over $X$ (understood that $X'\in 𝒰$), one defines in the
obvious way, as in \hyperref[IV.subsubsec4.10.2]{4.10.2}, the
``small'' sheaf $\underline{F}_{X'}$, the restriction of
$\underline{F}$ to $X'$. If
$$
u:X''\longrightarrow X'
$$
is\pageoriginale a morphism in $(\Esp)_{/X}$, one then defines in
the obvious way a morphism
$u_{\ast}(\underline{F}_{X''})\rightarrow \underline{F}_{X'}$, or,
what amounts to the same thing, a ``\emph{transition morphism}''
\begin{equation*}
\varphi_u :u^{\ast}(\underline{F}_{X'})\longrightarrow
\underline{F}_{X''}{\quoi}. \tag{4.10.3.1}
\end{equation*}\etiquette[4.10.3.1]{IV.eq4.10.3.1}
As $u$ varies, these morphisms satisfy an evident transitivity
condition for a composite
$$
X'''\xrightarrow{u}X''\xrightarrow{v}X'
$$
of morphisms in $(\Esp)_{/X}$, whose explicit formulation is left to
the reader. In this way one obtains a natural functor from the
category $\TOP(X)$ of large sheaves on $X$ to the category of systems
$$
(\underline{F}_{X'})\quad(X'\in \ob (\Esp)_{/X}), (\varphi_u)\quad(u\in
Fl(\Esp)_{/X}){\quoi},
$$
satisfying the transitivity condition just described and, in
addition, such that for every morphism $u:X''\rightarrow X'$ which is
an \emph{open immersion} (or, more generally, a local homeomorphism),
the transition morphism $\varphi_u$ is an isomorphism. Since the
functors $u^{\ast}:\Top(X')\rightarrow \Top(X'')$ used to describe
the $\varphi_u$ commute with $𝒱$-inductive limits, it follows at
once that, in the preceding description of the objects of $\TOP(X)$
in terms of ``small'' sheaves $F_{X'}$, $𝒱$-inductive limits are
computed pointwise, \ie the functors of the form
$\underline{F}\rightarrow \underline{F}_{X'}$ commute with
$𝒱$-inductive limits.

In particular, this holds for the functor $\underline{F}\mapsto
\underline{F}_X$ considered in \hyperref[IV.subsubsec4.10.2]{4.10.2}.
It therefore does have a right adjoint (\hyperref[IV.cor1.5]{1.5}). It
remains to construct a left adjoint (whose existence also follows a
priori from \hyperref[IV.rem1.8]{1.8}) and to verify that the latter is
left exact. This\pageoriginale will complete the definition of the
morphisms of topoi announced in \eqhyperref[IV.eq4.10.1]{4.10.1} in
terms of the functors \eqhyperref[IV.eq4.10.1.1]{4.10.1.1}.

\subsubsection{}\etiquette[4.10.4]{IV.subsubsec4.10.4}
The functor
$$
g^{\ast}=f_{!}:\Top(X)\longrightarrow \TOP(X)
$$
is obtained by associating with every small sheaf $F$ on $X$ the
system of its inverse images $(F_{X'})$, $X'\in \ob (\Esp)_{/X}$,
under the structural morphisms $X'\rightarrow X$, the transition
morphism $\varphi_u$ for $u:X''\rightarrow X'$ being the transitivity
isomorphism for inverse images of sheaves
(\hyperref[IV.subsubsec4.1.1]{4.1.1}). One sees at once that this
gives a functor
$$
\Prol:\Top(X)\longrightarrow \TOP(X){\quoi},
$$
which is \emph{fully faithful}, and whose essential image consists of
the large sheaves $ℱ$ on $X$ for which all the transition
morphisms $\varphi_u$ of \sisi{4.11.3}{\eqhyperref[IV.eq4.10.3.1]{4.10.3.1}}
are isomorphisms. We leave to the reader the task of defining an
adjunction isomorphism between this canonical prolongation functor and
the restriction functor of \hyperref[IV.subsubsec4.10.2]{4.10.2},
proving that the latter is right adjoint to the former. This is
immediate from the description in \hyperref[IV.subsubsec4.10.3]{4.10.3}
of the category $\TOP(X)$.

\begin{enonce*}[remark]{Remarks 4.10.5}\etiquette[4.10.5]{IV.rem4.10.5}
\begin{enumerate}
\renewcommand{\theenumi}{\alph{enumi}}
\renewcommand{\labelenumi}{\theenumi)}
\item The construction in \hyperref[IV.subsubsec4.10.4]{4.10.4}
also shows that $g^{\ast}=f_{!}$ is fully faithful (or, what amounts
to the same thing, that $g_{\ast}=f^{\ast}$ is a localization
functor, or again that $g^{!}=f_{\ast}$ is fully faithful). The large
sheaves on $X$ belonging to the essential image of the functor
$g^{\ast}=f_{!}=\Prol$ deserve the name of large \emph{étale
sheaves} on $X$, since they form a category equivalent to the category
of ordinary sheaves on $X$, or again to that of étale spaces
over\pageoriginale $X$. It is moreover easy to see that if $X'$ is a
topological space over $X$, then the large sheaf on $X$ represented by
it is étale in the preceding sense if and only if $X'$ is an étale
space over $X$.

\item The full faithfulness of $f_{\ast}$ may also be expressed by
saying that the adjunction morphism
$f^{\ast}f_{\ast}\rightarrow \id_{\Top(X)}$ is an isomorphism, \ie
that $g_{\ast}f_{\ast}\simeq \id$, \ie that $gf\simeq \id$. Thus $g$
makes $\TOP(X)$ into a \emph{topos over} $\Top(X)$, admitting a
``section'' $f$ over $\Top(X)$.

\item The fact that the functor $g^{\ast}=f_{!}$ is fully faithful,
exact, and commutes with $𝒱$-inductive limits partially justifies
the very convenient point of view (due to J. \textsc{Giraud})
according to which, in practically all questions in the theory of
sheaves on $X$, it is harmless to replace the usual sheaves, or
``small'' sheaves, by the associated ``large'' sheaves. This applies
in particular to cohomological questions, since $g_{\ast}$ is exact,
the functors $R^i g_{\ast}$ (\sisi{V\quad}{\hyperrefext[V][V.sec5]{5}})
vanish for $i>0$, and hence (\sisi{V\quad}{\hyperrefext[V][V.prop5.4]{5.4}})
for every \emph{large} sheaf $ℱ$ on $X$ there are canonical
isomorphisms
$$
H^i(\TOP(X),ℱ)\simeq H^i(\Top(X),
ℱ_X)=H^i(X,F){\quoi}.
$$
Applying this to a sheaf of the form $\Prol(F)$, where $F$ is a small
sheaf on $X$, we obtain (since $\Prol(F)_X\simeq F$ canonically) a
canonical isomorphism
$$
H^i(X,F)=H^i(\TOP(X), \Prol(F)){\quoi}.
$$
Thus, \emph{the cohomological invariants of $X$, computed by means of
the small or the large topos of $X$, are essentially identical.} The
same result is moreover valid in noncommutative cohomology.
\end{enumerate}
\end{enonce*}

\begin{enonce*}[remark]{Exercise 4.10.6}\etiquette[4.10.6]{IV.exer4.10.6}
Let\pageoriginale $S$ be a $𝒰$-site whose topology is
coarser than the canonical topology, and let $\underline{M}$ be a
subset of $Fl$ $S$ satisfying the following conditions:
\begin{enumerate}
\renewcommand{\theenumi}{\alph{enumi}}
\renewcommand{\labelenumi}{\theenumi)}
\item The morphisms of $ℳ$ are pullbackable
(\sisi{\hyperrefext[I][I.sec10.3]{10.3}}{\hyperrefext[I][I.sec10.7]{10.7}}),
and $\underline{M}$ is stable under base change.

\item $\underline{M}$ contains the identity arrows and is stable under
composition.

\item An arrow $u:X\rightarrow Y$ such that there exists a covering
family $Y_i\rightarrow Y$, with the arrows
$\fprod{X}{Y_i}{Y}\rightarrow Y_i$ in $\underline{M}$, itself belongs
to $\underline{M}$.

\item For every $X\in \ob S$, every covering family of $X$ is refined
by a covering family $(f_i:X_i\rightarrow X)$, with
$f_i\in \underline{M}$.
\end{enumerate}

For every $X\in \ob S$, consider the site $S(X)$ (the ``small site of
$X$''), whose underlying category is the full subcategory of
$S_{/X}$ consisting of the objects whose structural morphism belongs
to $\underline{M}$, endowed with the topology induced
(\hyperrefext[III][III.subsec3.1]{3.1}) by that of $S$.
\begin{itemize}
\item[$1^{\circ})$] For every arrow $u:X\rightarrow Y$ of $S$, show
that base change by $u$ from $S(Y)$ to $S(X)$ is a continuous
functor, hence gives a functor commuting with inductive limits
(\sisi{III}{\hyperrefext[III][III.prop1.2]{1.2} iv)})
$$
S(u)^x:S(Y)^{\sim}\longrightarrow S(X)^{\sim}{\quoi}.
$$

\item[$2^{\circ})$] Define an equivalence between the topos
$S^{\sim}$ and the category of systems
$$
(F_X)\quad (X\in \ob S)\quav (\varphi_u)\quad (u\in Fl S)
$$
consisting\pageoriginale of objects $F_X\in \ob S(X)^{\sim}$ and,
for every arrow $u:X\rightarrow Y$ in $S$, of a morphism
$\varphi_u:S(u)^x(F_Y)\rightarrow F_X$, these systems being subject
to a transitivity condition for a composite $v\circ u$ of arrows of
$S$, and to the condition that $u\in \underline{M}$ imply that
$\varphi_u$ is an isomorphism.

\item[$3^{\circ})$] Define ``restriction'' and ``prolongation''
functors
$$
\Res_X:(S_{/X})^{\sim}\longrightarrow S(X)^{\sim},
\Prol_X:S(X)^{\sim}\longrightarrow (S_{/X})^{\sim}
$$

Show that $\Res_X$ commutes with small inductive and projective limits
and that $\Prol_X$ is \emph{fully faithful}, its essential image
consisting of the sheaves $F$ such that $\varphi_u$ is an isomorphism
for every arrow $u$ of $S_{/X}$.

\item[$4^{\circ})$] Define an adjunction morphism making $\Res_X$ the
right adjoint of $\Prol_X$. Deduce that there exists a morphism of
topoi
$$
f:S(X)^{\sim}\rightarrow (S_{/X})^{\sim}{\quoi}
$$
making $S(X)^{\sim}$ a subtopos of $(S_{/X})^{\sim}$ and such that
\begin{align*}
f_{!} &= \Prol_X\\
f^x &= \Res_X{\quoi}.
\end{align*}

Show that $\Prol_X$ takes abelian sheaves to abelian sheaves.

\item[$5^{\circ})$] Show that if $S(X)$ admits fiber products, then
$\Prol_X$ is exact. Deduce that there then exists a morphism of topoi
$g:(S_{/X})^{\sim}\rightarrow S(X)^{\sim}$ which is a left retraction
of $f$, \ie such that
\begin{align*}
g^x &= \Prol_X\\
g_x &= \Res_X{\quoi}.
\end{align*}
(For\pageoriginale an example in which $S(X)$ does not admit fiber
products, take for $S$ the category of schemes endowed with the étale
topology, for $M$ the smooth morphisms, and for $X$ a Noetherian
scheme of dimension $>0$.)

\item[$6^{\circ})$] Show that for every abelian sheaf $F$ of
$(S_{/X})^{\sim}$ there is an isomorphism
$$
H^q(X,F)\simeq H^q(X,\Res_X F)\qquad \forall q{\quoi}.
$$
Show, using hypercoverings for example, that for every abelian sheaf
$G$ of $S(X)^{\sim}$ there is an isomorphism
$$
H^q(X,G)\simeq H^q(X,\Prol_X G)\qquad \forall q{\quoi}.
$$

\item[$7^{\circ})$] Buy a chocolate medal for the editor.
\end{itemize}
\end{enonce*}

% END INCLUDED SGA4_English_current_buildable_source_20260906/source/04/drafts/SGA4_IV_LINES_2355_2641_English.texfrag


% BEGIN INCLUDED SGA4_English_current_buildable_source_20260906/source/04/drafts/SGA4_IV_LINES_2642_2987_English.texfrag
\section{Induced topos}\etiquette[5]{IV.sec5}

\subsection{}\etiquette[5.1]{IV.subsec5.1}
Let $E$ be a topos and $X$ an object of $E$. Then the category
$E_{/X}$ of objects of $E$ over $X$ is a topos, as follows, for
example, immediately from Giraud's criterion
\hyperref[IV.theo1.2]{1.2} ii). Alternatively, by \hyperref[IV.cor1.2.1]{1.2.1},
one may realize $E$ as a category of sheaves $C^{\sim}$, where
$C$ is a generating subcategory\pageoriginale of $E$ that may be
chosen so that $X\in \ob C$; one then knows that
$E_{/X}=C^{\sim}_{/X}$ is equivalent to $(C_{/X})^{\sim}$
(\hyperrefext[III][III.prop5.4]{5.4}), and hence is a topos.

The topos $E_{/X}$ is called the \emph{topos induced} on the object
$X$ of $E$.

\subsection{}\etiquette[5.2]{IV.subsec5.2}
We shall define a canonical morphism of topoi
\begin{equation*}
j_X:E_{/X}\longrightarrow E{\quoi},\tag{5.2.1}
\end{equation*}\etiquette[5.2.1]{IV.eq5.2.1}
which is called the \emph{inclusion morphism} of the induced topos
$E_{/X}$ into the ambient topos $E$, or, better (\cf
\hyperref[IV.subsec5.7]{5.7}), the \emph{localization morphism of $E$ at $X$.}
This morphism corresponds to a sequence of three adjoint functors
(\cf \hyperref[IV.subsubsec3.1.3]{3.1.3})
\begin{equation*}
j_{X!}\quav j^{\ast}_X\quav j_{X^{\ast}}{\quoi},\tag{5.2.2}
\end{equation*}\etiquette[5.2.2]{IV.eq5.2.2}
which may be described explicitly as follows:
\begin{enumerate}
\renewcommand{\theenumi}{\alph{enumi}}
\renewcommand{\labelenumi}{\theenumi)}
\item The functor
$$
j_{X!}:E_{/X}\longrightarrow E
$$
is the functor that \emph{forgets the structural arrow} from
$E_{/X}$ to $E$.

\item The functor
$$
j^{\ast}_X:E\longrightarrow E_{/X}
$$
is defined by
$$
j^{\ast}_X(Z)=(X\times Z, \pr_1){\quoi},
$$
where $\pr_1:X\times Z\rightarrow X$ is the first projection. It
may also be interpreted as the \emph{base-change functor} relative
to the morphism
$$
X\longrightarrow e_E{\quoi},
$$
where $e_E$ is the terminal object of $E$. By the definition of
the base-change functor, it is immediate that $j^{\ast}_X$ is
indeed right adjoint to $j_{X!}$. In particular, it commutes with
projective limits. It also commutes with inductive
limits,\pageoriginale by virtue of the special exactness properties
of topoi (\hyperrefext[II][II.sec4]{4}).

\item It follows from this last fact (\hyperref[IV.cor1.6]{1.6}) that $j^{\ast}_X$
has a right adjoint
$$
j_{X^{\ast}}:E_{/X}\longrightarrow E{\quoi}.
$$
For an object $X'$ over $X$, the object $j_{X^{\ast}}(X')$ is also
sometimes denoted by one of the symbols $\prod_{X/e_E}(X'/X)$,
$\Hhom_{X/e_E}(X,X')$, or $\Res_{X/e_E}(X'/X)$ (``Weil
restriction''), one or another of which is doubtless already
familiar to the learned reader of our modest work.
\end{enumerate}

\subsection{}\etiquette[5.3]{IV.subsec5.3}
One should note that the functor $j_{X!}$ commutes with fiber
products and takes monomorphisms to monomorphisms, but it is not in
general left exact for all that. It is left exact only if $X$ is a
terminal object of $E$ (since $X=j_{X!}(X,\id_X)$ and
$(X,\id_X)$ is a terminal object of $E_{/X}$), that is, if and only
if $j_X$ is in fact an equivalence of topoi (and hence all the
functors in \eqhyperref[IV.eq5.2.2]{5.2.2} are equivalences).
\footnote[1]{For a commutation property of $j_{X!}$ under change of
topos, \cf XVII \hyperref[XVII.lem5.1.2]{5.1.2}.}

Likewise, the functor $j_{X^{\ast}}$ is not in general right exact
and does not necessarily take epimorphisms to epimorphisms.

It follows from these observations that the direction of the
morphism of topoi relating $E_{/X}$ to $E$, for an object $X$ of the
topos $E$, is determined without any possible ambiguity.

\subsection{}\etiquette[5.4]{IV.subsec5.4}
The functor $j^{\ast}_X$ is often called the \emph{localization
functor} or the \emph{restriction functor}. The latter terminology
is justified by identifying the objects of $E$ \resp of $E_{/X}$
with the sheaves on $E$ \resp on $E_{/X}$ (\hyperref[IV.theo1.2]{1.2}
iii)), and observing that, under this identification, the adjunction
formula between $j_{X!}$ (the forgetful functor) and $j^{\ast}_X$
may be interpreted as saying that $j^{\ast}_X(F)$ \emph{is\pageoriginale
the restriction to $E_{/X}$ of the sheaf $F$ on $E$} (where
``restriction'' is understood in the generalized sense: composition
with the ``inclusion'' functor $j_{X!}:E_{/X}\rightarrow E$).
In accordance with notation familiar from other contexts, we shall
also often write, interchangeably:
\begin{equation*}
j^{\ast}_X(F)=F| X=F_X=\mbox{ restriction of } F\mbox{ to }
X{\quoi}.\tag{5.4.1}
\end{equation*}\etiquette[5.4.1]{IV.eq5.4.1}

Likewise, when $E$ is of the form $C^{\sim}$, where $C$ is a
(small) site, and $X$ is of the form $\epsilon(S)$ with $S\in
\ob C$, so that one has the equivalence of categories recalled in
\hyperref[IV.subsec5.1]{5.1}
$$
E_{/X}=C^{\sim}_{/\epsilon(S)}\xrightarrow{\approx}(C_{/S})^{\sim}
$$
(where $C_{/S}$ is endowed with the topology induced by that of
$C$), the functor $j^{\ast}_X$ is identified simply with the functor
that ``restricts'' a variable sheaf $F$ on $C$ to the category
$C_{/S}$.

These considerations also make it possible to foresee the important
role that induced topoi and the localization morphisms
\eqhyperref[IV.eq5.2.1]{5.2.1} will play in all questions in which one is led
to argue by ``localization on $E$'' (\cf 8), that is, in practically
all questions involving sites or topoi.

\subsection{}\etiquette[5.5]{IV.subsec5.5}
Let
$$
f:X\longrightarrow Y
$$
be an arrow of $E$, which thus allows us to regard $X$ (or, more
correctly, $(X,f)$) as an object of $E_{/Y}$. It is clear that there
is a canonical isomorphism
\begin{equation*}
(E_{/Y})_{/X}\xrightarrow{\sim}E_{/X}\tag{5.5.1}
\end{equation*}\etiquette[5.5.1]{IV.eq5.5.1}
(\emph{transitivity of induced topoi}). Applying the construction of
\hyperref[IV.subsec5.1]{5.1} to $E_{/Y}$ in place of $E$, one therefore
obtains a canonical morphism of topoi
\begin{equation*}
\loc(f)\mbox{ or } f:E_{/X}\longrightarrow E_{/Y}{\quoi},\tag{5.5.2}
\end{equation*}\etiquette[5.5.2]{IV.eq5.5.2}
also called the\pageoriginale \emph{localization morphism
(associated with $f$).} When $Y$ is the terminal object of $E$, one
essentially recovers \eqhyperref[IV.eq5.2.1]{5.2.1}. The localization morphism
for $f$ is associated with a sequence of three adjoint functors
\begin{equation*}
f_{!}\quav f^{\ast}\quav f_{\ast}{\quoi},\tag{5.5.3}
\end{equation*}\etiquette[5.5.3]{IV.eq5.5.3}
which may be interpreted respectively as a forgetful functor, a
restriction functor, and a functor denoted, according to preference
(with $X'$ as argument), by $\prod_{X/Y}(X'/X)$,
$\Hhom_{X/Y}(X,X')$, or $\Res_{X/Y}(X'/X)$.

\subsection{}\etiquette[5.6]{IV.subsec5.6}
The transitivity of the forgetful functors implies that, for two
composable morphisms
$$
X\xrightarrow{f} Y \xrightarrow{g} Z
$$
of $E$, there is a canonical isomorphism of morphisms of topoi
$$
\loc(gf)\simeq \loc(g)\loc(f):E_{/X}\longrightarrow
E_{/Y}\longrightarrow E_{/Z}{\quoi}.
$$
Moreover, for three composable morphisms, the preceding transitivity
isomorphisms satisfy the usual compatibility relation. Thus, subject
to the customary abuse of language, this allows one to regard
$$
X\longmapsto E_{/X}
$$
as a covariant functor from $E$ to the category
($𝒱$-$𝒰$-$\text{top}$) (\hyperref[IV.subsubsec3.3.1]{3.3.1}).
More precisely, one obtains a 2-functor (not strict in general) from
$E$ to the 2-category ($𝒱$-$𝒰$-$\text{top}$).

Before continuing with the general theory of induced topoi
(\hyperref[IV.subsec5.10]{5.10}), let us give a few instructive examples.

\subsection{}\etiquette[5.7]{IV.subsec5.7}
Let\pageoriginale $X$ be a topological space, and hence let
$\Top(X)$ be a topos (\hyperref[IV.subsec2.1]{2.1}). Let $X'$ be an object
of $\Top(X)$, which it will be convenient to interpret as an
\'etale space over $X$, $p:X'\rightarrow X$. The category
$\Top(X)_{/X'}$ is then identified with the category of \'etale
spaces $X''$ over $X$, equipped with an $X$-morphism
$X''\rightarrow X'$. It is known that such a morphism makes $X''$
an \'etale space over $X'$, and in this way one obtains a canonical
equivalence of topoi
$$
\Top(X)_{/X'}\xrightarrow{\approx} \Top(X'){\quoi}.
$$
The localization morphism is therefore identified with a morphism
of topoi $\Top(X')\rightarrow \Top(X)$, and one sees immediately
that the latter is none other than the morphism
$$
\Top(p):\Top(X')\longrightarrow \Top(X)
$$
associated with the structural continuous map $p:X'\rightarrow X$.
Intuitively, the localization morphism is thus simply the translation,
into the language of topoi, of the \'etale map $p:X'\rightarrow X$.
This both explains the appropriateness of the term ``localization
morphism'' and counsels caution in using the term ``inclusion
morphism'', which seems chiefly appropriate when $p$ is an open
immersion. It will consequently be prudent in
\hyperref[IV.subsec5.1]{5.1} to reserve the term ``inclusion morphism'' for
the case where the morphism $X\rightarrow e_E$ is a monomorphism,
\ie where $X$ is identified with a subobject of the terminal object
of $E$.

\subsection{}\etiquette[5.8]{IV.subsec5.8}
Let $E$ be a topos, $G$ a group in $E$, and $H$ a subgroup of $G$,
and let
$$
X=G/H
$$
be the quotient homogeneous space, regarded as an object of $E$
with a left action by $G$, \ie as an object of the classifying topos
$B_G$ (\hyperref[IV.subsec2.4]{2.4}). We\pageoriginale shall determine the
induced topos
$$
(B_G)_{/X}\quad (X=G/H){\quoi},
$$
by defining an equivalence of topoi
\begin{equation*}
c:B_H \xrightarrow{\approx} (B_G)_{/X}{\quoi},\tag{5.8.1}
\end{equation*}\etiquette[5.8.1]{IV.eq5.8.1}
such that the following diagram commutes up to canonical
isomorphism:
\[\vcenter{
\xymatrix@C=2cm{ B_H \ar[r]^{c} \ar[r] \ar
@{^{(}->}@<-9pt>[rr]_{B_i}& {(B_{G})_{/X}} \ar[r]^{j_X}& B_G
}}\leqno{(5.8.2)}
\]\etiquette[5.8.2]{IV.eq5.8.2}
where $j_X$ is the localization morphism in $B_G$, and where
$$
B_i:B_H\longrightarrow B_G
$$
is the morphism deduced from the inclusion $i:H\rightarrow G$
(\hyperref[IV.subsec4.5]{4.5}).

To define \eqhyperref[IV.eq5.8.1]{5.8.1}, we shall simply give the
descriptions of the inverse-image functor $c^{\ast}$ and the
direct-image functor $c_{\ast}$, leaving it to the reader to verify
that these are indeed quasi-inverse equivalences and give rise to
the commutative diagram \eqhyperref[IV.eq5.8.2]{5.8.2} of morphisms of
topoi. Let $e$ be the subobject of $E$ underlying $X$, the image of
the section of $X$ (over a chosen terminal object $e_E$ of $E$)
deduced from the unit section of $G$. If $X'$ is an object of $B_G$
over \sisi{$G$}{$X$}, then $c^{\ast}(X')$ denotes the inverse image
of $e$ in $X'$,
$$
c^{\ast}(X')=\fprod{X'}{e}{X}{\quoi},
$$
which is stable under the left action of $H$ on $X'$, obtained by
restricting the given action of $G$ on $X'$. This indeed defines a
functor
$$
c^{\ast}:(B_G)_{/X}\longrightarrow B_H{\quoi}.
$$
On the other\pageoriginale hand, if $X'$ is an object of $B_H$, then
the morphism from $X'$ to the terminal object $e_H$ of $B_H$ (\ie
the terminal object $e_E$ of $E$ with the trivial action of $H$)
gives, upon application of the functor $i_{!}$
(\hyperref[IV.subsec4.5]{4.5}), a morphism in $B_G$
$$
i_{!}(X')\longrightarrow i_{!}(e_H)\simeq X{\quoi},
$$
which allows $i_{!}(X')$ to be regarded as an object of
$(B_G)_{/X}$, and hence gives the desired functor
$$
c_{\ast}:X'\longmapsto (i_{!}(X')\longrightarrow
X):B_H\longrightarrow (B_G)_{/X}{\quoi}.
$$

\setcounter{subsubsection}{2}
\subsubsection{}\etiquette[5.8.3]{IV.subsubsec5.8.3}
Thus, by \eqhyperref[IV.eq5.8.2]{5.8.2}, one sees that if $G$ is a group
in a topos and $H$ a subgroup, \emph{the functor ``restriction of
the operators from $G$ to $H$'' may be interpreted as a localization
functor.} In particular, taking $H$ to be the unit subgroup, one sees
that the functor that \emph{forgets the actions of $G$} may be
interpreted as a localization functor. More generally, it follows
that if $(X,G)$ is an object of $E$ with an action of $G$ (\ie an
object of the classifying topos $B_G$), then the functor that
``forgets the actions of $G$'',
$$
E'=(B_G)_{/(X,G)}\longrightarrow E_{/X}
$$
may be interpreted as a localization functor with respect to a
suitable object $E_{G,(X,G)}$ of $E'$ covering the terminal object.
It suffices to take $E_{G,(X,G)}=E_G\times (X,G)$, giving the
cartesian diagram (where $E_G=G_S=G$ with the action of $G$ by left
translation)
\[\vcenter{
\xymatrix@C=2cm@R=1.5cm{
E_G \ar[d]& **[r]E_{G,(X,G)} = Z \ar[l] \ar[d]\\
e_G & (X,G) \ar[l]          \quad , }}\leqno{(5.8.3.1)}
\]\etiquette[5.8.3.1]{IV.eq5.8.3.1}
and\pageoriginale to observe that, by transitivity of induced topoi,
the induced topos $E'=((B_G)_{X/G)/Z}$ is equivalent to the topos
$((B_G)_{/E_G})_{/Z}$; since $(B_G)_{/E_G}$ is equivalent to $E$
by the functor $c^{\ast}$ of \eqhyperref[IV.eq5.8.1]{5.8.1} (with $H=e$),
it suffices to verify that this equivalence takes $Z$ to the object
$X$ of $E$ (which is immediate) in order to deduce the desired
equivalence of categories
$$
E'_{/Z}\xrightarrow{\approx}E_{/X}{\quoi}.
$$
We leave it to the reader to verify that the composite of this
equivalence with the localization functor
$$
E'=(B_G)_{/(X,G)}\longrightarrow E'_{/Z}
$$
is the functor that ``forgets the actions of $G$''. Thus one sees
that, upon passing to induced topoi, the cartesian diagram above gives
a diagram of morphisms of topoi, commutative up to canonical
isomorphism (\hyperref[IV.subsec5.6]{5.6}):
\[\vcenter{
  \xymatrix@C=1.25cm@R=1cm{
**[r](B_G)_{/E_G}\approx E \ar[d]_{j} & **[r](B_G)_{/Z} \approx
E'_{/Z} \ar[l]_(.3){i}
  \approx E_{/X} \ar[d]^{j'}\\
B_G & **[r](B_G)_{/(X,G)} =  E' \ar[l]_(.5){f}\quad,
}}\leqno{(5.8.3.2)}
\]\etiquette[5.8.3.2]{IV.eq5.8.3.2}
where the inverse-image functors associated with the vertical arrows
$j$, $j'$ are the functors that ``forget the actions of $G$'', and
where the inverse-image functor $i^{\ast}$ is identified with the
localization functor $E\rightarrow E_{/X}$. This diagram\pageoriginale is,
moreover, ``2-cartesian'' in the sense of \hyperref[IV.prop5.11]{5.11}
below.

\subsubsection{}\etiquette[5.8.4]{IV.subsubsec5.8.4}
We leave it to the reader to extend the preceding considerations to
the case of a pro-group $𝒢=(G_i)$, in the particular case where $H$
is defined as the ``kernel'' of $𝒢\rightarrow G_{i_{\circ}}$. (In
the case of profinite groups, this means restricting attention to
\emph{open} subgroups $H$ of $G$, \ie subgroups that are closed and
of \emph{finite index}.)

% END INCLUDED SGA4_English_current_buildable_source_20260906/source/04/drafts/SGA4_IV_LINES_2642_2987_English.texfrag


% BEGIN INCLUDED SGA4_English_current_buildable_source_20260906/source/04/drafts/SGA4_IV_LINES_2988_3412_English.texfrag
\begin{enonce*}[remark]{Exercise 5.9}\etiquette[5.9]{IV.exer5.9}
Let $E$ be a topos, $G$ a group of $E$, $B_G$ the classifying topos
of $G$ (\hyperref[IV.subsec2.4]{2.4}), and
$$
\pi:B_G\longrightarrow E
$$
the morphism of topoi deduced from the homomorphism from $G$ to the
unit group $e$ of $E$ (taking into account that $B_e\approx E$).
\begin{enumerate}
\renewcommand{\theenumi}{\alph{enumi}}
\renewcommand{\labelenumi}{\theenumi)}
\item Let
$$
E_G=G_s
$$
be the object of $B_G$ whose underlying $E$-object is $G$, with the
actions of $G$ defined by left translations. Regard $\pi^{\ast}(G)$
as a group of $B_G$ (it is the group $G$ of $E$ with the trivial
action of $G$ on it). Show that the morphism of $E$
$$
G_s\times G\longrightarrow G_s
$$
defined by the right translations of $G$ on $G_s$ is compatible
with the actions of $G$ on $G_s$ and on $G=\pi^{\ast}(G)$, and
defines a morphism of $B_G$
$$
E_G\times \pi^{\ast}(G)\longrightarrow E_G{\quoi}.
$$
Show that this morphism equips $E_G$ with a right action of the
$B_G$-group $\pi^{\ast}(G)$, and that this action makes $E_G$ a
\emph{torsor} under\pageoriginale $\pi^{\ast}(G)$ (that is, the
terminal object of $B_G$ can be covered by objects $X_i$ such that,
for every $i$, the restriction of $E_G$ to $X_i$ is isomorphic to
the trivial bundle with operators $(\pi^{\ast}(G)_{X_i})_d$).

\item Let
$$
q:E'\longrightarrow E
$$
be a topos over $E$. For every topos $B$ over $E$, define the
category $\Hhomtop_E(E',B)$ of morphisms of topoi from $E'$ to $B$
compatible, up to a specified isomorphism, with the structural
morphisms $E'\rightarrow E$ and $B\rightarrow E$. Taking $B=B_G$,
define, by the formula $f\mapsto f^{\ast}(E_G)$, an equivalence of
categories
$$
\Hhomtop_E(E',B_G)\longrightarrow \Ttors(E',q^{\ast}(G)){\quoi},
$$
where the right-hand side denotes the category of (right)
$q^{\ast}(G)$-torsors on $E'$.

\item Let $H$ be a subgroup of $G$, so that $\pi^{\ast}(H)$ is a
subgroup of $\pi^{\ast}(G)$ and hence acts on $E_G$ by restriction
of the operator group, giving a quotient object
$$
X=E_G/\pi^{\ast}(H){\quoi},
$$
which is simply the object $G/H$ endowed with the usual left action
of $G$. Let $e_G$ be the terminal object of $B_G$ (that is, the
terminal object $e_E$ of $E$ with the necessarily trivial action of
$G$), and consider in $B_G$ the diagram of morphisms
\[
\xymatrix@C=-1cm@R=.75cm{
E_G \ar[dd] \ar[rd]& \\
    & \qquad\qquad\quad X = E_{G}/\pi^{*}(H)\ar[ld]\\
e_G & }
\]
which,\pageoriginale upon passing to the topoi induced by $B_G$ on
these objects, gives a diagram of morphisms of topoi, commutative up
to canonical isomorphism (\hyperref[IV.subsec5.6]{5.6}). Show that
this diagram is equivalent to the diagram
\[
\xymatrix@C=.05cm@R=.85cm{
E & = & B_e \ar[dd] \ar[rrrd]& & &\\
  &   &                    &  & &B_H \ar[llld]\\
  &   & B_G                &  & & &,
}
\]
whose three arrows are the arrows associated by
\hyperref[IV.subsec4.5]{4.5} with the group morphisms
$e\rightarrow H\rightarrow G$ in $E$. (Thus, \emph{the classifying
topos $B_H$ may intuitively be interpreted as a homogeneous space
over $B_G$, with group $\pi^{\ast}(G)$, associated with the universal
torsor} (= homogeneous principal bundle) \emph{$E_G$ over $B_G$}.)

\item Reconsider the example of \hyperref[IV.subsec5.7]{5.7} in the
case where $X'$ is an \emph{\'etale covering} (locally trivial, \cf
\hyperref[IV.subsubsec2.7.4]{2.7.4}) of $X$, interpreting it in terms
of the considerations of the present exercise.
\end{enumerate}
\end{enonce*}

\setcounter{subsection}{9}
\subsection{Inverse images of induced topoi}\etiquette[5.10]{IV.subsec5.10}

Let
$$
f:E'\rightarrow E
$$
be a morphism of $𝒰$-topoi, let $X$ be an object of $E$, and let
$X'=f^{\ast}(X)$ be its inverse image in $E'$. We shall define a
diagram of morphisms of topoi
\[\vcenter{
\xymatrix@C=2.5cm@R=1.5cm{
E_{/X}  \ar[d]_{j_X} & E'_{/X'} \ar[l]_{f_{/X}} \ar[d]^{j_{X'}} \\
E      &  E'\ar[l]_{f}\quad , }}\leqno{(5.10.1)}
\]\etiquette[5.10.1]{IV.eq5.10.1}
where\pageoriginale $j_X$ and $j_{X'}$ are the localization
morphisms (\hyperref[IV.subsec5.2]{5.2}), and where $f_{/X}$ is
defined below; this diagram will be commutative up to canonical
isomorphism. We shall define $f_{/X}$ here by means of its inverse
image functor
\begin{equation*}
(f_{/X})^{\ast}:E_{/X}\longrightarrow E'_{/X'}{\quoi},\tag{5.10.2}
\end{equation*}\etiquette[5.10.2]{IV.eq5.10.2}
which we simply take to be the functor ``induced'' by $f^{\ast}$, in
the evident sense. Since the inclusion functors
$j_{X!}:E_{/X}\rightarrow E$ and
$j_{X'!}:E'_{/X'}\rightarrow E'$ commute with $𝒰$-inductive limits
and fiber products, and are conservative, it follows at once that
the functor $(f_{/X})^{\ast}$ commutes with $𝒰$-inductive limits and
fiber products, just as does the functor $f^{\ast}$ which induces
it. Moreover, since $(f_{/X})^{\ast}$ plainly takes the terminal
object $X$ to the terminal object $X'$, it is left exact and hence
is associated with a morphism of topoi
\begin{equation*}
f_{/X}:E'_{/X'}\longrightarrow E_{/X}{\quoi},\tag{5.10.3}
\end{equation*}\etiquette[5.10.3]{IV.eq5.10.3}
defined up to unique isomorphism
(\hyperref[IV.subsubsec3.1.1]{3.1.1}). On the other hand, the diagram
deduced from \eqhyperref[IV.eq5.10.1]{5.10.1} by passing to the
inverse-image functors is commutative (up to canonical isomorphism)
by construction and by the fact that $f^{\ast}$ commutes with the
products $X\times Y$. Thus \eqhyperref[IV.eq5.10.1]{5.10.1} itself
is commutative, up to a unique isomorphism inducing the canonical
isomorphism on the inverse-image functors of the two composites
$f\circ j_{X'}$ and $j_X\circ f_{/X}$ \hyperlink{sga4.exp.IV.subsubsection.3.2.1}{(3.2.1)}.

\setcounter{subsubsection}{3}
\subsubsection{}\etiquette[5.10.4]{IV.subsubsec5.10.4}
When $f:E'\rightarrow E$ is associated with a continuous map
$Y'\rightarrow Y$ of topological spaces
(\hyperref[IV.subsec4.1]{4.1}), so that $X$ is identified with an
\'etale space over $Y$, and $X'$ with the fiber product
$\fprod{X}{Y'}{Y}$, then, identifying the induced topoi
$\Top(Y)_{/X}$ and $\Top(Y')_{/X'}$ with $\Top(X)$ and $\Top(X')$
respectively (\hyperref[IV.subsec5.7]{5.7}), one sees that the
morphism $f_{/X}$ just constructed\pageoriginale is (up to unique
isomorphism) the morphism $\Top(X')\rightarrow\Top(X)$ deduced from
the canonical continuous map $X'\rightarrow X$. In the light of
this example, one may therefore say that the topos $E'_{/X'}$ plays
the role of a fiber product of $E_{/X}$ and $E'$ over $E$. This
intuition is made precise by the following result:

\begin{enonce*}{Proposition 5.11}\etiquette[5.11]{IV.prop5.11}
With the notation of \hyperref[IV.subsec5.10]{5.10}, the diagram
\eqhyperref[IV.eq5.10.1]{5.10.1}, equipped with the compatibility
isomorphism $\alpha:f\circ j_{X'}\rightarrow j_X\circ f_{/X}$, is
``$2$-cartesian.'' More precisely, for every $𝒰$-topos $F$, if to
each morphism of topoi $g:F\rightarrow E'_{/X'}$ one associates the
triple $(f_{/X}\circ g, j_{X'}\circ g,
\alpha\ast g)=(g_1,g_2,\beta)$, where
$g_1:F\rightarrow E_{/X}$ and $g_2:F\rightarrow E'$ are morphisms
of topoi and $\beta:j_X\circ g_1\rightarrow f\circ g_2$ is an
isomorphism of morphisms of topoi from $F$ to $E$, then one obtains
an equivalence between the category $\Hhomtop(F,E'_{/X'})$ and the
category
\sisi{$\fprod{\Hhomtop(F,E_{/X})^2}{\Hhomtop(F,E')}{\Hhomtop(F,E)}$}{$\Hhomtop(F,E_{/X})\times^{2}_{\Hhomtop(F,E)}\Hhomtop(F,E')$}
of all triples $(g_1,g_2,\beta)$ as above.
\end{enonce*}

(NB. A $2$ has been placed above the cartesian-product sign to
recall that this is not an ordinary fiber product of categories, but
a ``$2$-fiber product,'' in the sense made explicit in the
statement.)

To prove \hyperref[IV.prop5.11]{5.11}, we shall describe morphisms of
topoi explicitly by means of their associated inverse-image
functors. For this we shall need the auxiliary results given in
\hyperref[IV.subsubsec5.11.1]{5.11.1} through
\hyperref[IV.prop5.12]{5.12} below.

\setcounter{subsection}{11} \setcounter{subsubsection}{0}
\subsubsection{}\etiquette[5.11.1]{IV.subsubsec5.11.1}
Consider, in general, the situation in which two categories $E$ and
$E'$ are given, together with an object $X$ of $E$ which is assumed
to be pullbackable, that is, such that the functor
$$
j:A\longmapsto A\times X:E\longrightarrow E_{/X}
$$
is\pageoriginale defined, and finally a left exact functor
\begin{equation*}
\varphi:E_{/X}\longrightarrow E'{\quoi}.\tag{$\ast$}
\end{equation*}
Since $E_{/X}$ has a terminal object, namely $(X,\id_X)=X$, the same
is therefore true of $E'$, which has the terminal object
$$
e'\simeq \varphi(X){\quoi}.
$$
Having chosen the terminal object $e'$ of $E'$, we shall associate
with every $\varphi$ as above a pair
\begin{equation*}
(\Psi,u),\mbox{ with } \Psi=\varphi\circ j:E\rightarrow E',\mbox{ and
} u\in \Hom(e',\Psi(X)){\quoi}.\tag{${\ast}{\ast}$}
\end{equation*}
To define $u$, note that
$$
j(X)=X\times X
$$
has a canonical section $\delta_X$ over the terminal object $X$ of
$E_{/X}$, namely the diagonal section, and define
$$
u=\varphi(\delta_X):\varphi(X)=e'\longrightarrow \varphi(X\times
X)=\Psi(X){\quoi}.
$$
I claim that knowledge of the pair $({\ast}{\ast})$ makes it
possible to reconstruct the functor $\varphi$ of $(\ast)$ up to
unique isomorphism. For this, for every object
$$
p:X'\longrightarrow X
$$
of $E_{/X}$, consider the following cartesian diagram in $E_{/X}$:
\[
\xymatrix@C=1cm@R=1.5cm{
X' \ar[d] \ar[r] & **[r]j(X') = X'\times X \ar[d]^{j(p)}\\
X  \ar[r]^{\delta_X}& **[r]j(X) = X \times X   &, }
\]
where the first horizontal arrow is the graph morphism
$\Gamma_p=(\id_{X'},p)$. Applying the left exact functor $\varphi$ to
this diagram, and using the definition of $\Psi$ as $\varphi\circ
j$, gives a cartesian diagram
\[
\xymatrix@C=2.25cm@R=1.25cm{
\varphi(X',p) \ar[d] \ar[r]& \Psi (X') \ar[d]^{\Psi(p)}\\
e'  \ar[r]          & \Psi (X) \quad , }
\]
or,\pageoriginale in other words, a canonical isomorphism
\begin{equation*}
\varphi(X',p)\simeq
\fprod{\Psi(X')}{e'}{\Psi(X)}{\quoi}\tag{5.11.1.1}
\end{equation*}\etiquette[5.11.1.1]{IV.eq5.11.1.1}
It is immediately seen to be functorial in the object $(X',p)$ of
$E_{/X}$, which describes explicitly how the functor
$\varphi:E_{/X}\rightarrow E'$ is reconstructed, up to isomorphism,
from the pair $(\Psi,u)$ of $({\ast}{\ast})$. Notice also that the
functor $\Psi$ is left exact. More generally, $\Psi$ commutes with
every type of projective limit with which $\varphi$ commutes,
because $j$ commutes with projective limits.

Conversely, suppose now that finite projective limits are
representable in $E$ and $E'$, and start with a pair
$$
(\Psi,u),\Psi:E\longrightarrow E', u\in \Hom(e',\Psi(X)){\quoi},
$$
where the functor $\Psi$ is left exact. Formula
\eqhyperref[IV.eq5.11.1.1]{5.11.1.1} then defines a functor
$\varphi:E_{/X}\rightarrow E'$. One checks at once that this functor
is left exact and, more generally, that it commutes with every type
of projective limit representable in $E$ with which $\Psi$ commutes.
This follows immediately from the cartesian diagram
\[
\xymatrix@C=2.5cm@R=1cm{ \displaystyle{\varprojlim^{E}_{I} X'_i}
\ar[d] &
\displaystyle{\varprojlim^{E_{/X}}_{I} X'_i} \ar[d] \ar[l]\\
\displaystyle{\varprojlim^{E}_{I} X} & X \ar[l]_{\delta} }
\]
relating\pageoriginale the projective limits computed in $E$ and in
$E_{/X}$ (where $\delta$ denotes the diagonal morphism into the
projective limit of the constant functor with value $X$): apply the
functor $\Psi$ to it, obtaining a cartesian diagram in $E'$, and
compose the latter with the cartesian square deduced from
$\Psi(X)\leftarrow e'$. The resulting composite cartesian square
expresses the compatibility of $\varphi$ with the projective limit
under consideration.

The pair $(\Psi,u)$ is reconstructed up to isomorphism from
$\varphi$ by observing that there is a canonical isomorphism
\begin{equation*}
\Psi(X')=\varphi(j(X')){\quoi},\tag{5.11.1.2}
\end{equation*}\etiquette[5.11.1.2]{IV.eq5.11.1.2}
deduced from \eqhyperref[IV.eq5.11.1.1]{5.11.1.1} by replacing
$(X',p)$ there with $j(X')=(X'\times X,\pr_2)$ and observing that
$\Psi(X'\times X)\simeq\Psi(X')\times\Psi(X)$. The preceding
isomorphism is manifestly functorial in $X'$, that is, it gives an
isomorphism
$$
\Psi\simeq \varphi\circ j{\quoi},
$$
and one likewise checks that, under this isomorphism,
$u:e'\rightarrow\Psi(X)$ is identified with $\varphi(\delta_X)$.
We have thus obtained the essential part of the following result.

\begin{enonce*}{Lemma 5.11.2}\etiquette[5.11.2]{IV.lem5.11.2}
Let $E$ and $E'$ be two categories in which finite projective limits
are representable, let $X$ be an object of $E$, let $\Ssex(E,E')$
be the category of left exact functors $\Psi$ from $E$ to $E'$, and
let $\Ssex(E_{/X},E')$ be the analogous category of left exact
functors $\varphi$ from $E_{/X}$ to $E'$. There is then an
equivalence between the category $\Ssex(E_{/X},E')$ and the category
$\Ssex(E,E')_{/\Gamma}$ of pairs $(\Psi,u)$, where $\Psi\in\ob
\Ssex(E,E')$ and $u:e'\rightarrow\Psi(X)$ (with $e'$ the terminal
object of $E'$). This equivalence, as well as a quasi-inverse
functor, is described explicitly in
\hyperref[IV.subsubsec5.11.1]{5.11.1}. If $\varphi$ and $(\Psi,u)$
correspond, then $\varphi$ commutes with a given type of projective
limit if\pageoriginale and only if $\Psi$ does. If the base-change
functors in $E$ and $E'$ commute with a given type of inductive
limit, then $\varphi$ commutes with inductive limits of that type if
and only if $\Psi$ does.
\end{enonce*}

We leave it to the reader to make explicit the category structure on
$\Ssex(E,E')_{/\Gamma}$ corresponding to the pairs $(\Psi,u)$. If
the functor $\Gamma:\Ssex(E,E')\rightarrow(\Ens)$ is defined by
$\Gamma(\Psi)=\Hom(e',\Psi(X))$, then $u$ in the pair $(\Psi,u)$ may
be interpreted as an element of $\Gamma(\Psi)$, that is, as a
homomorphism $\Psi\rightarrow\Gamma$ in the category of presheaves
on $\Ssex(E,E')^{\circ}$. This justifies the notation
${}_{/\Gamma}$ used in the statement of
\hyperref[IV.subsubsec5.11.1]{5.11.1}. To finish the proof of
\hyperref[IV.lem5.11.2]{5.11.2}, it remains to make explicit in
\hyperref[IV.subsubsec5.11.1]{5.11.1} the functorial nature of the
constructions considered when $\varphi$, respectively $(\Psi,u)$,
varies. This is essentially trivial and is left to the reader. It
also remains to verify the assertion concerning commutation with
inductive limits; this too follows immediately from the explicit
formulas \eqhyperref[IV.eq5.11.1.1]{5.11.1.1} and
\eqhyperref[IV.eq5.11.1.2]{5.11.1.2} relating these functors.

In particular, when $E$ and $E'$ are $𝒰$-topoi, and morphisms of
topoi are interpreted by means of their associated inverse-image
functors, one obtains:

\begin{enonce*}{Proposition 5.12}\etiquette[5.12]{IV.prop5.12}
Let $E$ and $E'$ be two $𝒰$-topoi, let $X$ be an object of $E$, and
consider on $\Hhomtop(E',E)$ the contravariant functor
$$
\gamma:f\longmapsto
\Gamma(E',f^{\ast}(X))=\Hom(e',f^{\ast}(X)):\Hhomtop(E',E)^{\circ}\longrightarrow
(\Ens){\quoi}.
$$
There is then an equivalence of categories
\begin{equation*}
\Hhomtop(E',E_{/X})\xrightarrow{\approx}\Hhomtop(E',E)_{/\gamma}{\quoi},\tag{5.12.1}
\end{equation*}\etiquette[5.12.1]{IV.eq5.12.1}
where\pageoriginale the right-hand side is the full subcategory of
$\Hhomtop(E',E)^{\widehat{\phantom{,}}}_{/\gamma}$ consisting of
pairs $(f,u)$ with $f\in\Homtop(E',E)$ and $u\in\gamma(f)$, that is,
$u\in\Gamma(E',f^{\ast}(X))$. This functor is obtained by
associating with every morphism of topoi $h:E'\rightarrow E_{/X}$
the pair $(f,u)$ consisting of the composite
$$
f=j_X\circ h:E'\rightarrow E_{/X}\rightarrow E{\quoi},
$$
and the morphism
$$
u:e'=h^{\ast}(X,\id_X)\longrightarrow
f^{\ast}(X)=h^{\ast}(j^{\ast}_X(X))=h^{\ast}(X\times X)
$$
deduced from the diagonal morphism $\delta_X:X\rightarrow X\times
X$ by applying $h^{\ast}$ to it.
\end{enonce*}

\setcounter{subsection}{12}
\subsubsection{}\etiquette[5.12.2]{IV.subsubsec5.12.2}
Using \hyperref[IV.prop5.12]{5.12}, the proof of
\hyperref[IV.prop5.11]{5.11} becomes almost evident, and essentially
reduces to the following (which will serve in place of a formal
proof that would be no more instructive): to give a morphism of
topoi $g:F\rightarrow E'_{/X'}$ ``amounts to the same thing'' as to
give a morphism of topoi $g_2:F\rightarrow E'$ and a section $u$ of
$g^{\ast}_2(X')$. But
$g^{\ast}_2(X')=g^{\ast}_2(f^{\ast}(X))=(fg_2)^{\ast}(X)$, so giving
$u$ is also equivalent to giving a lifting of the morphism of topoi
$fg_2:F\rightarrow E$ to a morphism of topoi
$g_1:F\rightarrow E_{/X}$. This expresses precisely that giving a
morphism of topoi $g:F\rightarrow E'_{/X'}$ is essentially
equivalent to giving a triple $(g_1,g_2,\alpha)$ as in
\hyperref[IV.prop5.11]{5.11}.

\begin{enonce*}[remark]{Remark 5.13}\etiquette[5.13]{IV.rem5.13}
We shall see (\S 15\nde{This section does not exist. See, however,
[16], Proposition 4.47, for a general existence statement for fiber
products of topoi (noting that all geometric morphisms between
Grothendieck topoi are bounded).}) that for every diagram
\[
\xymatrix@C=2cm@R=1.25cm{
E_1 \ar[d]&\\
E   & E_2 \ar[l] }
\]
of\pageoriginale topoi, there exists a $2$-fiber product in the sense
of the $2$-category ($𝒱$-$𝒰$-$\Top$) of $𝒰$-topoi belonging to
$𝒱$ (independent, up to equivalence, of the choice of a universe
$𝒱$ such that $𝒰\in𝒱$). For natural examples of ``fiber products
of topoi'' other than \hyperref[IV.prop5.11]{5.11}, see the final
chapter of \textsc{Giraud}'s book [3].
\end{enonce*}

\begin{enonce*}[remark]{Exercise 5.14}\etiquette[5.14]{IV.exer5.14}
\begin{enumerate}
\renewcommand{\theenumi}{\alph{enumi}}
\renewcommand{\labelenumi}{\theenumi)}
\item Let $F$ and $G$ be two topoi over a topos $E$. Define the
category $\Hhomtop_E(F,G)$ of pairs $(f,\alpha)$, where
$f:F\rightarrow G$ is a morphism of topoi and $\alpha:p\rightarrow
q\circ f$ is an isomorphism of morphisms of topoi from $F$ to $E$
(where $p:F\rightarrow E$ and $q:G\rightarrow E$ are the structural
morphisms of topoi). Define pairing functors
$\Hhomtop_E(F,G)\times\Hhomtop_E(G,H)\rightarrow\Hhomtop_E(F,H)$,
and associativity isomorphisms.

\item Let $X$ and $Y$ be two objects of a topos $E$. Define a
natural functor from the discrete category defined by the set
$\Hom(X,Y)$ to the category
$\Hhomtop_E(E_{/X},E_{/Y})$. Establish compatibility with the
composition of morphisms in $E$ and with the pairings considered in
a).

\item Prove that the functor considered in b) is an equivalence of
categories. In particular, an object $X$ of a topos $E$ can be
reconstructed up to unique isomorphism when one knows, ``up to
$E$-equivalence,'' the induced topos $E_{/X}$ as a topos
\emph{over $E$.}~\footnote[1]{This result is due to P.
\textsc{Deligne}. The case where $E$ is the punctual topos had
previously been treated by Mme \textsc{Hakim}.}
\end{enumerate}
\end{enonce*}

% END INCLUDED SGA4_English_current_buildable_source_20260906/source/04/drafts/SGA4_IV_LINES_2988_3412_English.texfrag


% BEGIN INCLUDED SGA4_English_current_buildable_source_20260906/source/04/drafts/SGA4_IV_LINES_3413_3631_English.texfrag
\section{Points of a topos and fiber functors}\etiquette[6]{IV.sec6}

\setcounter{subsection}{-1}
\subsection{}\etiquette[6.0]{IV.subsec6.0}
Let $P$ be the standard $𝒰$-\emph{punctual topos}
(\hyperref[IV.subsec2.2]{2.2})
$$
P=(𝒰\text{-}\Ens){\quoi}.
$$
In view of the rather different intuitions attached on the one hand
to the symbol\pageoriginale $P$, representing a geometric object of
the nature of a point, and on the other hand to
$(𝒰\text{-}\Ens)$, representing a ``large category'' to be
interpreted as ``the category of sheaves on $P$,'' one should use one
or the other symbol exclusively according to context, although from
a strictly logical viewpoint they denote one and the same object.

\begin{enonce*}{Definition 6.1}\etiquette[6.1]{IV.def6.1}
Let $E$ be a $𝒰$-topos. A point of $E$ is a morphism of topoi
$P\rightarrow E$ from the punctual topos $P$
(\hyperref[IV.subsec6.0]{6.0}) to $E$. The category
$\Hhomtop(P,E)$ (\hyperref[IV.subsec3.2]{3.2}) is called the category
of points of $E$ and is denoted by $\Ppoints(E)$ or
$\sheaf{\Pt}(E)$. If $p:P\rightarrow E$ is a point of $E$, with
associated inverse-image functor
$p^{\ast}:E\rightarrow(𝒰\text{-}\Ens)$, and if $F$ is an object of
$E$, the set $p^{\ast}(F)$ is called the fiber of $F$ at $p$ and is
denoted by $F_p$.
\end{enonce*}

\setcounter{subsection}{1}
\subsubsection{}\etiquette[6.1.1]{IV.subsubsec6.1.1}
Of course, when $F$ is a group object (\resp ring object,
\resp\ldots) of $E$, its fiber $F_p$ at $p$ is a group (\resp a ring,
\resp\ldots) (\hyperref[IV.subsubsec3.1.2]{3.1.2}).

\setcounter{subsubsection}{1}

\subsubsection{}\etiquette[6.1.2]{IV.subsubsec6.1.2}
By \hyperlink{sga4.exp.IV.subsubsection.3.2.1}{3.2.1}, the category of points of $E$ is equivalent, via the
functor $p\mapsto p^{\ast}$, to the \emph{opposite} of the category
of functors
$$
\varphi:E\longrightarrow (𝒰\text{-}\Ens)
$$
which commute with $𝒰$-inductive limits and are left exact. Thus it
amounts essentially to the same thing to give a point of $E$ or to
give such a functor $\varphi$.

\begin{enonce*}{Definition 6.2}\etiquette[6.2]{IV.def6.2}
A fiber functor of the $𝒰$-topos $E$ is a functor
$\varphi:E\rightarrow(𝒰\text{-}\Ens)$ which commutes with
$𝒰$-inductive limits and is left exact. The category of fiber
functors on $E$, denoted by $\Ffib(E)$, is the\pageoriginale full
subcategory of $\Hhom(E,(𝒰\text{-}\Ens))$ formed by the fiber
functors on $E$.
\end{enonce*}

\setcounter{subsection}{2}
\subsubsection{}\etiquette[6.2.1]{IV.subsubsec6.2.1}
By \hyperref[IV.subsubsec6.1.2]{6.1.2}, the category $\Ffib(E)$ of
fiber functors on $E$ is therefore equivalent to the \emph{opposite}
of the category $\Ppoint(E)$ of points of $E$, via the functor
$p\mapsto p^{\ast}$ from the latter to the former:
\begin{equation*}
\Ppoint(E)\xrightarrow{\approx}\Ffib(E)^{\circ}{\quoi}.
\tag{6.2.1.1}
\end{equation*}\etiquette[6.2.1.1]{IV.eq6.2.1.1}

A functor
$$
\varphi:E\longrightarrow (𝒰\text{-}\Ens)
$$
is a fiber functor if and only if it is the fiber functor
$F\mapsto F_p$ associated with a point $p$ of $E$
(\hyperref[IV.def6.1]{6.1}). In other words, the functor
\eqhyperref[IV.eq6.2.1.1]{6.2.1.1} is surjective. Let us also note
that a functor $\varphi$ is a fiber functor if and only if it is left
exact and carries covering families to surjective families
(\hyperref[IV.cor1.7]{1.7}). This last condition may also be expressed
by saying that $\varphi$ commutes with $𝒰$-sums and carries
epimorphisms to epimorphisms.

\subsection{}\etiquette[6.3]{IV.subsec6.3}
Let $C\in 𝒰$ be a site. A \emph{point of the site} $C$ is a point of
the topos $C^{\sim}$, and the \emph{category of points of the site}
$C$ is the category
$$
\Ppoint(C)=\Ppoint(C^{\sim}){\quoi}.
$$
Consider, on the other hand, the functor
\begin{equation*}
\varphi\longrightarrow \varphi|
C=\varphi\circ\epsilon_C:\Ffib(C^{\sim})\longrightarrow
\Hhom(C,(𝒰\text{-}\Ens)){\quoi},\tag{6.3.1}
\end{equation*}\etiquette[6.3.1]{IV.eq6.3.1}
which is fully faithful and whose essential image is the full
subcategory $\Mmorsite((𝒰\text{-}\Ens),C)^{\circ}$
(\hyperref[IV.prop4.9.4]{4.9.4}), which we shall also denote by
$\Ffib(C)$. A functor
$$
\Psi:C\longrightarrow (𝒰\text{-}\Ens)
$$
is\pageoriginale called a \emph{fiber functor on the site} $C$ if it
belongs to the essential image of \eqhyperref[IV.eq6.3.1]{6.3.1},
just described. Such a functor is therefore \emph{continuous} (a
fortiori (III \hyperref[III.prop1.6]{1.6}), it carries covering families to covering
families, \ie to surjective families), and it is left exact. When
finite projective limits are representable in $C$ (a condition
satisfied in almost every case encountered in practice), the
preceding properties \emph{characterize} the fiber functors on $C$:
they are then precisely the functors
$\Psi:C\rightarrow(𝒰\text{-}\Ens)$ which are left exact and carry
covering families to surjective families
(\hyperref[IV.prop4.9.4]{4.9.4}).

One should retain that the functor $\varphi\rightarrow\varphi|C$ is
an equivalence of categories
$$
\Ffib(C^{\sim})\xrightarrow{\approx}\Ffib(C){\quoi},
$$
so that \emph{it amounts essentially to the same thing to give a
fiber functor on the topos $C^{\sim}$} (or equivalently a point of
that topos) \emph{or to give a fiber functor on the site $C$.} Given
a fiber functor $\Psi$ on $C$, arising up to isomorphism from a fiber
functor $\varphi$ on $C^{\sim}$, the latter is reconstructed from
$\Psi$, up to canonical isomorphism, by the formula
\begin{equation*}
\varphi(F)\simeq \varinjlim_{C_{/F}}\Psi (X){\quoi},\tag{6.3.2}
\end{equation*}\etiquette[6.3.2]{IV.eq6.3.2}
where $C_{/F}$ is the category of objects $X$ of $C$ endowed with a
morphism $X\rightarrow F$ (in $\widehat{C}$), \ie an element of
$F(X)$. Formula \eqhyperref[IV.eq6.3.2]{6.3.2} is an immediate
consequence of the fact that $\varphi$ commutes with inductive limits
and of the canonical isomorphism, functorial in $F$
(\hyperrefext[II][II.cor4.1.1]{4.1.1}):
$$
F\simeq \varinjlim_{C_{/F}}\epsilon_C(X){\quoi}.
$$
More\pageoriginale generally, when $G$ is a presheaf on $C$, there
is a canonical isomorphism, functorial in $G$:
\begin{equation*}
\varphi(\underline{a}G)\simeq
\varinjlim_{C_{/G}}\Psi(X){\quoi},\tag{6.3.3}
\end{equation*}\etiquette[6.3.3]{IV.eq6.3.3}
where $\underline{a}G$ is the sheaf associated with $G$
(\hyperrefext[II][II.cor4.1.1]{4.1.1}). Indeed, one has the formula
$$
\underline{a}G\simeq \varinjlim_{C_{/G}}\epsilon_C(X){\quoi}.
$$

\setcounter{subsection}{4} \setcounter{subsubsection}{-1}
\subsubsection{}\etiquette[6.4.0]{IV.subsubsec6.4.0}
We refer to \hyperrefext[I][I.def6.1]{6.1} for the notion of a
\emph{conservative family of functors}
$$
\varphi_i:E\longrightarrow E_i\qquad i\in I{\quoi}.
$$
Note that when $E$ and the $E_i$ are $𝒰$-topoi and the functors
$\varphi_i$ are inverse-image functors $f_i^{\ast}$ associated with
morphisms of topoi
$$
f_i:E_i\longrightarrow E{\quoi},
$$
all the exactness properties assumed in
\hyperrefext[I][I.prop6.2]{6.2}, \hyperrefext[I][I.cor6.3]{6.3}, and
\hyperrefext[I][I.prop6.4]{6.4} are satisfied, provided that in
\hyperrefext[I][I.prop6.2]{6.2} (v) $D$ is assumed finite. The
following conditions are then equivalent: $(f_i^{\ast})_{i\in I}$ is
conservative; it is conservative for \emph{monomorphisms}
(\hyperrefext[I][I.def6.1]{6.1}); it is conservative for
\emph{epimorphisms}; and it is \emph{faithful}.

When\pageoriginale the family of functors $f_i^{\ast}$ is
conservative, one also sometimes says that \emph{the family of
morphisms of topoi $(f_i)_{i\in I}$ is conservative}. In particular:

\begin{enonce*}{Definition 6.4.1}\etiquette[6.4.1]{IV.def6.4.1}
Let $E$ be a $𝒰$-topos and let $(p_i:P\rightarrow E)_{i\in I}$ be a
family of points of $E$. This family of points is called conservative
if the family of associated fiber functors $F\mapsto F_{p_i}$ from
$E$ to $(𝒰\text{-}\Ens)$ is conservative
(\hyperref[IV.subsubsec6.4.0]{6.4.0}). The topos $E$ is said to have
enough points if it admits a conservative family of points (\ie if
the family of all points of the topos is conservative).
\end{enonce*}

\setcounter{subsection}{4} \setcounter{subsubsection}{1}
\subsubsection{}\etiquette[6.4.2]{IV.subsubsec6.4.2}
Let us note that all $𝒰$-topoi used thus far have enough points. It
is nevertheless possible, ``on purpose,'' to construct topoi which
do not have enough points (\hyperref[IV.exer7.2.6]{7.2.6} e) and
\hyperref[IV.subsec7.4]{7.4}); such a topos is necessarily not
``empty'' in the geometric sense of \hyperref[IV.subsec2.2]{2.2}. A
topos having enough points admits a conservative family of points
indexed by some $I\in𝒰$ (\hyperref[IV.prop6.5]{6.5}). Note, however,
that if $E$ is a $𝒰$-topos, the set of isomorphism classes of points
of $E$ is not necessarily $𝒰$-small (\hyperref[IV.subsec7.3]{7.3}),
although it is so in many cases encountered in practice. Finally,
for an interesting existence theorem (due to P. Deligne) providing
enough points and covering all cases arising in algebraic
geometry\sisi{.}{, see \hyperrefext[VI][VI.sec9]{9}.}

\subsubsection{}\etiquette[6.4.3]{IV.subsubsec6.4.3}
When $(p_i)_{i\in I}$ is a conservative family of points of the
topos $E$, the remarks of \hyperrefext[I][I.prop6.2]{6.2} apply. In
particular, they imply that two arrows $u,v:F\rightrightarrows G$ in
$E$ are equal if and only if, for every $i\in I$, the induced arrows
on fibers $u_{p_i},v_{p_i}:F_{p_i}\rightrightarrows G_{p_i}$
are\pageoriginale equal; that an arrow $u:F\rightarrow G$ in $E$ is
a monomorphism (\resp an epimorphism) if and only if, for every
$i\in I$, the map $u_{p_i}:F_{p_i}\rightarrow G_{p_i}$ is so; that
an object $F$ of $E$ is initial (\resp terminal) if and only if, for
every $i\in I$, $F_{p_i}$ is empty (\resp consists of one point); and
that an object $F$ is a subobject of the terminal object $e_E$ if and
only if, for every $i\in I$, $F_{p_i}$ has at most one point.


% END INCLUDED SGA4_English_current_buildable_source_20260906/source/04/drafts/SGA4_IV_LINES_3413_3631_English.texfrag


% BEGIN INCLUDED SGA4_English_current_buildable_source_20260906/source/04/drafts/SGA4_IV_LINES_3632_3863_English.texfrag
\begin{enonce*}{Proposition 6.5}\etiquette[6.5]{IV.prop6.5}
\begin{enumerate}
\renewcommand{\theenumi}{\alph{enumi}}
\renewcommand{\labelenumi}{\theenumi)}
\item Let $C$ be a $𝒰$-site and let $(\varphi_i)_{i\in I}$ be a
family of fiber functors on $C$ (\hyperref[IV.subsec6.3]{6.3}). The
family $(p_i)_{i\in I}$ of corresponding points of the topos
$E=C^{\sim}$ is conservative if and only if every family
$(X_j\rightarrow X)_{j\in J}$ in $C$ for which, for every $i\in I$,
the corresponding family
$(\varphi_i(X_j)\rightarrow\varphi_i(X))_{j\in J}$ is surjective is
a covering family.

\item Let $E$ be a $𝒰$-topos. If $E$ has enough points
(\hyperref[IV.def6.4.1]{6.4.1}), then $E$ admits a $𝒰$-small
conservative family of points.
\end{enumerate}
\end{enonce*}



Assertion b) is a special case of \hyperrefext[I][I.prop7.7]{7.7}.
To prove a), apply \hyperrefext[I][I.prop7.7]{7.7} to the generating
family in $E$ formed by the $\epsilon(X)$, $X\in\ob C$. Recall
(\hyperrefext[II][II.theo4.4]{4.4}) that the family
$X_j\rightarrow X$ is covering if and only if the family
$\epsilon(X_j)\rightarrow\epsilon(X)$ is epimorphic. If the family of
points $(p_i)_{i\in I}$ is conservative, this amounts
(\hyperref[IV.subsubsec6.4.3]{6.4.3}) to saying that, for every
$i\in I$, the family $(X_j)_{p_i}\rightarrow(X)_{p_i}$ is
surjective, \ie that the family
$\varphi_i(X_j)\rightarrow\varphi_i(X)$ is surjective. This proves
the ``only if'' part of a). For the ``if'' part, apply the criterion
\hyperrefext[I][I.prop7.7]{7.7}, which reduces us to verifying that
every monomorphism $F\rightarrow\epsilon(X)$\pageoriginale for which
$F_{p_i}\rightarrow\epsilon(X)_{p_i}$ is an isomorphism for every
$i\in I$ is an isomorphism, or equivalently an epimorphism. Now one
can find a covering family of morphisms
$\epsilon(X_j)\rightarrow F$, and, after refining further, one may
suppose that the morphisms
$\epsilon(X_j)\rightarrow\epsilon(X)$ are induced by morphisms
$X_j\rightarrow X$. By hypothesis, for every $i\in I$, the composite
morphisms $(X_j)_{p_i}\rightarrow F_{p_i}\rightarrow(X)_{p_i}$ form
a surjective family (as the composite of two surjective families),
\ie the family
$\varphi_i(X_j)\rightarrow\varphi_i(X)$ is surjective. It follows
from the hypothesis that the family $X_j\to X$ is covering, hence
that the family $\epsilon(X_j)\rightarrow\epsilon(X)$ is
epimorphic, and a fortiori that $F\rightarrow\epsilon(X)$ is an
epimorphism.

\begin{enonce*}{Corollary 6.5.1}\etiquette[6.5.1]{IV.cor6.5.1}
Let $C$ be a category. A topology on $C$ making $C$ a $𝒰$-site with
enough fiber functors (\ie such that the associated topos has enough
points) is completely determined by the full subcategory $\Phi$ of
$\Hhom(C,\Ens)$ formed by the fiber functors on $C$.
\end{enonce*}

Indeed, by \hyperref[IV.prop6.5]{6.5} a), one can then describe the
covering families in terms of the family of fiber functors.

In fact, we shall see below
\eqhyperref[IV.subsubsec6.8.5]{6.8.5} that the
subcategory $\Phi$ is contained in the category of pro-representable
functors on $C$, so that
$\Phi^{\circ}\simeq\Ppoint(C^{\sim})$ identifies, up to equivalence,
with a full subcategory of $\Pro-C$.

Compare \hyperref[IV.cor6.5.1]{6.5.1} with \textsc{Giraud}'s theorem
\hyperrefext[II][II.theo5.5]{5.5}, which establishes a one-to-one
correspondence between the set of topologies on $C$ and a certain set
of full subcategories of $\Hhom(C^{\circ},\Ens)$.

\begin{enonce*}[remark]{Exercise 6.5.2}\etiquette[6.5.2]{IV.exer6.5.2}
Let\pageoriginale $C$ be a category equivalent to a category in $𝒰$,
and let $P$ be a subcategory of $\Pro(C)$. For every $p\in P$, denote
by $X\mapsto X_p$ the functor pro-represented by $p$. A family
$(X_i\rightarrow X)_{i\in I}$ in $C$ is called
$P$-\emph{covering} if, for every $p\in P$, the family
$(X_{ip}\rightarrow X_p)_{i\in I}$ is surjective. Show that there is
a topology $T_P$ on $C$ whose covering families are precisely the
$P$-covering families, and that $T_P$ is the finest topology on $C$
for which the functors $X\mapsto X_p$, $p\in P$, are fiber functors.
Show that the topology $T_P$ makes $C$ a site with \emph{enough}
fiber functors. Show that, for every topology $T$ on $C$, among the
topologies $T'$ finer than $T$ which have enough fiber functors there
is a coarsest one: it is the topology $T_P$, where $P$ is the
strictly full subcategory of $\Pro(C)$, equivalent to
$\Ppoint(C)$, described in
\eqhyperref[IV.subsubsec6.8.5]{6.8.5} below.
\end{enonce*}

\begin{enonce*}[remark]{Problem 6.5.3}\etiquette[6.5.3]{IV.prop6.5.3}
Characterize the subcategories (strictly full, stable under
$𝒰$-filtered projective limits, \ldots) $P$ of $\Pro(C)$ which can
arise from a topology on $C$ as the essential image of
$\Ppoint(C^{\sim})$ (\hyperref[IV.subsubsec6.8.5]{6.8.5}).
\footnote[1]{\cf \hyperref[IV.exer9.1.8]{9.1.8} a) and e) for
necessary conditions, and \hyperref[IV.exer9.1.12]{9.1.12} b) for
necessary and sufficient conditions in the cases of $\Top(X)$, with
$X$ a sober locally Noetherian space.}
\end{enonce*}

\setcounter{subsection}{5}
\subsection{}\etiquette[6.6]{IV.subsec6.6}
Let
$$
f:E\longrightarrow E'
$$
be a morphism of $𝒰$-topoi. It induces a canonical functor
$$
\Ppoint(f)=(p\longmapsto f\circ p):\Ppoint(E)\longrightarrow
\Point(E'){\quoi}.
$$
Given two composable morphisms of topoi
$$
E\xrightarrow{f} E'\xrightarrow{g} E''{\quoi},
$$
one\pageoriginale checks immediately that
$$
\Ppoint(gf)=\Ppoint(g)\circ\Ppoint(f):\Ppoint(E)\longrightarrow
\Ppoint(E')\longrightarrow \Ppoint(E''){\quoi}.
$$
This makes precise the functorial dependence (without any abuse of
language for once) of $\Ppoint(E)$ on the topos $E$: if $𝒱$ is a
universe such that $𝒰\in𝒱$, one obtains a (genuine!) functor
$$
\Ppoint:(𝒱\text{-}𝒰\text{-}\text{top})\longrightarrow
(𝒱\text{-}\cat)
$$
from the category of $𝒰$-topoi which are elements of $𝒱$ to the
category of categories which are elements of $𝒱$.

\subsection{}\etiquette[6.7]{IV.subsec6.7}
Let $E$ be a $𝒰$-topos. Every object $F$ of $E$ then defines a
contravariant functor
$$
p\longmapsto F_p=p^{*}(F):\Ppoint(E)^{\circ}\longrightarrow
(𝒰\text{-}\Ens){\quoi},
$$
and hence, for varying $F$, a canonical functor
\begin{equation*}
E\longrightarrow
\Ppoint(E)^{\widehat{\phantom{,}}}=\Hhom(\Ppoint(E)^{\circ},
(𝒰\text{-}\Ens)){\quoi},\tag{6.7.1}
\end{equation*}\etiquette[6.7.1]{IV.eq6.7.1}
which, by definition (\hyperref[IV.prop6.5]{6.5}), is faithful if and
only if $E$ has enough points. (The functor
\eqhyperref[IV.eq6.7.1]{6.7.1} bears a certain formal analogy to a
Fourier transform\ldots)

Let $X$ be an object of $E$, which therefore defines a presheaf
$\widehat{X}$ on
$$
C=\Ppoint(E){\quoi}.
$$
We may therefore define the category
$$
C_{/\widehat{X}}=\Ppoint(E)_{/\widehat{X}}
$$
of morphisms $p\rightarrow\widehat{X}$ in the category
$\widehat{C}$ of presheaves on $C$, \ie the category of pairs
$(p,\xi)$ where $p$ is a point of $E$ and $\xi$ is an element of
$\widehat{X}(p)=X_p$, the fiber of $X$ at $p$. With this understood,
I claim that \emph{there is a canonical equivalence of categories}
\begin{equation*}
C'=\Ppoint(E_{/X})\xrightarrow{\approx}C_{/\widehat{X}}\quav \mbox{
where }\quad C=\Ppoint(E){\quoi},\tag{6.7.2}
\end{equation*}\etiquette[6.7.2]{IV.eq6.7.2}
whose\pageoriginale composite with the inclusion functor
$C_{/\widehat{X}}\rightarrow C$ is none other than the functor
$$
\Ppoint(j_X):C'=\Ppoint(E_{/X})\longrightarrow C=\Ppoint(E)
$$
induced functorially (\hyperref[IV.subsec6.6]{6.6}) by the
localization morphism \eqhyperref[IV.eq5.2.1]{5.2.1}
$$
j_X:E_{/X}\longrightarrow E{\quoi}.
$$
This is simply the special case of \hyperref[IV.prop5.12]{5.12}
obtained by taking $E'=P=(\Ens)$.

\begin{enonce*}{Corollary 6.7.3}\etiquette[6.7.3]{IV.cor6.7.3}
Let $E$ be a topos. If $E$ has enough points, then so does every
induced topos $E_{/X}$, $X\in\ob E$. Conversely, if
$(X_i)_{i\in I}$ is a family of objects of $E$ covering the terminal
object and, for every $i\in I$, $E_{/X_i}$ has enough points, then
$E$ has enough points.
\end{enonce*}

For the first assertion, let $f:X'\rightarrow X''$ be a morphism in
$E_{/X}$ which is not an isomorphism. We prove that there exists a
point $q$ of $E_{/X}$ such that $f_q$ is not an isomorphism. Since
$E$ has enough points, there is a point $p$ of $E$ such that
$f_p:X'_p\rightarrow X''_p$ is not an isomorphism. This implies that
there exists $q\in X_p$ such that the map induced by $f_p$ on the
fibers $X'_q$ and $X''_q$ of $X'_p$ and $X''_p$ over $X_p$ at $q$
is not an isomorphism. But by the equivalence
\eqhyperref[IV.eq6.7.2]{6.7.2}, $q$ may be identified with a point of
$E_{/X}$, and the map $X'_q\rightarrow X''_q$ just considered is
none other than the map on fibers at $q$ induced by
$f:X'\rightarrow X''$, which proves the assertion.

Conversely, suppose that the $X_i$ cover the terminal object of $E$
and that the $E_{/X_i}$ have enough points. We prove the same for
$E$. Let\pageoriginale $f:X'\rightarrow X''$ be a morphism in $E$
which is not an isomorphism. There is then some $i\in I$ such that
$f_i:X'_i\rightarrow X''_i$ is not an isomorphism, where the
subscript $i$ denotes restriction to $X_i$. Hence there is a point
$p_i$ of $E_{/X_i}$ such that
$f_{ip_i}:X'_{ip_i}\rightarrow X''_{ip_i}$ is not an isomorphism.
If $p$ denotes the image of $p_i$ in $E$ under the localization
morphism $E_{/X_i}\rightarrow E$, this says that
$X'_p\rightarrow X''_p$ is not an isomorphism, which proves that $E$
has enough points.

\begin{enonce*}[remark]{Remark 6.7.4}\etiquette[6.7.4]{IV.rem6.7.4}
The preceding argument shows that if $(p_{\alpha})$ is a conservative
family of points of $E$, then the family of points of $E_{/X}$ lying
over one of the $p_{\alpha}$ (a family indexed by the disjoint union
of the $X_{p_{\alpha}}$) is conservative. Likewise, if the $X_i$
cover the terminal object of $E$ and, for every $i\in I$, $P_i$ is a
conservative set of points of $E_{/X_i}$, then the set of points of
$E$ which are the \emph{images} of the $p\in P_i$ under the
localization morphisms $E_{/X_i}\rightarrow E$ (with $i\in I$ and
$p$ varying) is conservative.
\end{enonce*}

\setcounter{subsection}{7}

% END INCLUDED SGA4_English_current_buildable_source_20260906/source/04/drafts/SGA4_IV_LINES_3632_3863_English.texfrag


% BEGIN INCLUDED SGA4_English_current_buildable_source_20260906/source/04/drafts/SGA4_IV_LINES_3864_4120_English.texfrag
\subsection{}\etiquette[6.8]{IV.subsec6.8}
Let $E$ be a topos and $p$ a point of $E$. A \emph{neighborhood
of the point $p$ of the topos $E$} is a pair $(X,u)$, where $X\in \ob E$ and
$u\in X_p$. By \eqhyperref[IV.eq6.7.2]{6.7.2}, one may also interpret $u$
as a lift of $p$ to a point of the induced topos $E_{/X}$.
Writing $\varphi_p$ for the fiber functor associated with the point $p$, one
may also interpret a neighborhood of the point $p$ as an object of
the category $E_{/\varphi_p}$ (the full subcategory of
$(E^\circ)^{\widehat{\phantom{,}}}_{/\varphi_p}$ formed by the objects
whose source lies in $E\subset \widehat{E}$). This latter
category, \ie the category of pairs $(X,u)$ as above,
will naturally be called the \emph{category of neighborhoods} of the
point $p$ of the topos $E$; it is also denoted by $\Vvois(p)$. Since
$\varphi_p$\pageoriginale is left exact and finite projective limits
are representable in $E$, finite projective limits are representable
in the category $\Vvois(p)$, and the canonical functor
$\Vvois(p)\rightarrow E$ commutes with them
(\hyperrefext[I][I.sec8.10.10]{8.10.10}). A fortiori, the opposite category
$\Vvois(p)^{\circ}$ is filtered. Moreover, it is clear
(\hyperrefext[I][I.prop3.4]{3.4}) that there is a canonical isomorphism, functorial in
$F\in \ob E$,
\begin{equation*}
\varphi_p(F)=F_p\simeq
\varinjlim_{\Vvois(p)^{\circ}}F(X){\quoi}.\tag{6.8.1}
\end{equation*}\etiquette[6.8.1]{IV.eq6.8.1}
In other words, \emph{the functor $\varphi_p$, ``fiber at the
point $p$,'' is isomorphic to the filtered inductive limit of the
functors represented in the topos $E$ by the neighborhoods of the point
$p$ of $E$.} In fact, \emph{the fiber functor is
pro-representable} (\hyperrefext[I][I.sec8]{8}), which also means that it is
represented by a pro-object $(X_i)_{i\in I}$ of $E$, where $I$ is
a filtered ordered set such that $I\in𝒰$. Indeed, by
\loccit, this is equivalent to saying that the category
$\Vvois(p)$ admits a small full cofinal subcategory. Now let
$C$ be a small full generating subcategory of $E$; I claim that
the full subcategory $C_{/\varphi_{p}}$ of
$\Vvois(p)=E_{/\varphi_p}$, formed by the neighborhoods $(X,u)$ of $p$
such that $X\in \ob C$ (a category that is obviously small), is
cofinal in $\Vvois(p)$. Indeed, let $(X,u)$ be a neighborhood of
$p$; one must find a neighborhood $(X',u')$ of $p$ lying over
$(X,u)$, with $X'\in \ob C$. There is a covering family
$X'_i\rightarrow X$, with the $X'_i$ in $C$, from which it follows that the
$X'_{ip}$ cover $X_p$; hence there exist an $X'=X'_{i_{\circ}}$ and
a $u'\in X'_p$ such that $(X',u')$ lies over $(X,u)$, as
claimed.

\setcounter{subsubsection}{1}
\subsubsection{}\etiquette[6.8.2]{IV.subsubsec6.8.2}
Let\pageoriginale $C$ be a $𝒰$-site and $p$ a point of
$C$, \ie a point of the topos $C^{\sim}$. By an abuse of
notation, we again denote by $\varphi_p$ the ``restriction'' of the fiber functor
$\varphi_p$ to $C$, \ie the composite $\varphi_p\circ \epsilon$; for
an object $X$ of $C$, we shall also often write
$$
X_p=\varphi_p(X)=\epsilon(X)_p{\quoi}.
$$
A \emph{neighborhood of the point $p$ in the site $C$} is a pair
$(X,u)$, where $X\in \ob C$ and $u\in X_p=\varphi_p(X)$. These neighborhoods
again form a category $C_{/\varphi_p}$, which may be denoted by
$\Vvois_C(p)$. Let us show that \emph{the category $\Vvois_C(p)^{\circ}$
is again filtered, that it admits a $𝒰$-small full cofinal
subcategory, and that there is an isomorphism, functorial
in the sheaf $F$ on $C$:}
\begin{equation*}
F_p\simeq \varinjlim_{\Vvois_C(p)^{\circ}} F(X){\quoi}.\tag{6.8.3}
\end{equation*}\etiquette[6.8.3]{IV.eq6.8.3}
First note, by repeating the argument of \eqhyperref[IV.eq6.8.1]{6.8.1}, that if
$C'$ is a full subcategory of $C$ that is topologically
generating (II \hyperref[II.def3.0.1]{3.0.1}), then $\Vvois_{C'}(p)^{\circ}$ is cofinal
in $\Vvois_C(p)^{\circ}$. Taking $C'\in𝒰$ to be small, we are
thus reduced, in proving the assertion, to the case where $C\in𝒰$.
In this case $\widehat{C}$ is a $𝒰$-topos, and there is an
associated-sheaf functor $\underline{a}:\widehat{C}\rightarrow
C^{\sim}$, which allows us to construct on $\widehat{C}$ the composite
fiber functor $\varphi_p\underline{a}$. Applying \eqhyperref[IV.eq6.8.1]{6.8.1} to the
topos $\widehat{C}$ and to its full generating
subcategory $C$, one finds that the category
$C_{/\varphi_p\underline{a}}^{\circ},$ which is manifestly isomorphic to
$\Vvois_C(p)^{\circ}$, is filtered, and that there is an isomorphism
functorial in the presheaf $P$ on $C$:
$$
\varphi_p\underline{a}(P)\simeq \varinjlim_X P(X){\quoi}.
$$
If $F$ is a sheaf on $C$, then $F=\underline{a}F$, and applying
the preceding isomorphism\pageoriginale to $F$ regarded as a
presheaf gives \eqhyperref[IV.eq6.8.3]{6.8.3}. At the same time, this argument
establishes the more general formula
\begin{equation*}
(\underline{a}P)_p\simeq \varinjlim_{\Vvois_C(p)^{\circ}}
P(X){\quoi},\tag{6.8.4}
\end{equation*}\etiquette[6.8.4]{IV.eq6.8.4}
an isomorphism functorial in an arbitrary presheaf $P$ on $C$, at
least when $C$ is small. The general case follows by
introducing a universe $𝒱$ such that $C\in𝒱$, by
regarding $\varphi_p$ on $C$ as a fiber functor with values
in $(𝒱\text{-}\Ens)$, and using the compatibility of the
formation of the associated sheaf with enlargement of universes
(\hyperrefext[II][II.prop3.6]{3.6}).

\setcounter{subsubsection}{4}
\subsubsection{}\etiquette[6.8.5]{IV.subsubsec6.8.5}
To every point $p$ of the topos $C^{\sim}$ we have associated a pro-object of
the $𝒰$-site $C$, defined by the canonical functor
$\Vvois_C(p)\rightarrow C$, taking into account the fact that the category of
neighborhoods of $p$ in $C$ is cofiltered and admits a small
full cofinal subcategory. One immediately verifies that, as $p$
varies, this gives a functor
\begin{equation*}
\Ppoint(C^{\sim})\longrightarrow \Pro(C){\quoi}.\tag{6.8.5.1}
\end{equation*}\etiquette[6.8.5.1]{IV.eq6.8.5.1}
This functor is \emph{fully faithful.} Via the associated
fiber functors, it may be interpreted as a functor
$\Fib(C)\rightarrow \Fib(C^{\sim})\rightarrow \Pro(C)^{\circ}$. The
first functor is an equivalence by \hyperref[IV.prop4.9.4]{4.9.4}
(recalled for fiber functors in \hyperref[IV.subsec6.3]{6.3}), and the
second is fully faithful by the general formalism
(\hyperrefext[I][I.sec8]{8}) of pro-objects.

\subsubsection{}\etiquette[6.8.6]{IV.subsubsec6.8.6}
Consider the special case where $C$ is $𝒰$-small and endowed with
the chaotic topology, so that $C^{\sim}=\widehat{C}$. I claim that
in this case the preceding functor is even an equivalence of
categories
\begin{equation*}
\Ppoint(\widehat{C})\xrightarrow{\approx}\Pro(C){\quoi}.\tag{6.8.6.1}
\end{equation*}\etiquette[6.8.6.1]{IV.eq6.8.6.1}
Indeed,\pageoriginale it remains to prove that this functor is
essentially surjective, which follows from the fact that the functors
of the form $P\mapsto P(X)$ on $\widehat{C}$ are fiber functors,
and from the evident fact that every filtered inductive limit of
fiber functors on a topos is again a fiber functor.

Identifying $C$ with a full subcategory of $\Pro(C)$
(\hyperrefext[I][I.sec8]{8}), the functor
$$
C\longrightarrow \Ppoint(\widehat{C})
$$
induced by a quasi-inverse to \sisi{(6.8.5.2)}{\eqhyperref[IV.eq6.8.6.1]{6.8.6.1}} is manifestly
isomorphic to the canonical functor already considered in
\eqhyperref[IV.eq4.6.2.2]{4.6.2.2};
the quasi-inverse in question may be regarded as its canonical
extension (\hyperrefext[I][I.sec8]{8}), taking into account the fact that
$\Ppoint(\widehat{C})$ is stable under small projective limits.

\subsubsection{}\etiquette[6.8.7]{IV.subsubsec6.8.7}
More generally, let $(X_i)_{i\in I}$ be a pro-object of the
$𝒰$-site $C$, and hence a functor on $C^{\sim}$
\begin{equation*}
\varphi(F)=\varinjlim_I F(X_i){\quoi},\tag{6.8.7.1}
\end{equation*}\etiquette[6.8.7.1]{IV.eq6.8.7.1}
and, in view of (\hyperref[IV.subsubsec6.8.5]{6.8.5}), it is natural to ask
when this functor is a fiber functor. The following
conditions are equivalent:
\begin{enumerate}
\renewcommand{\theenumi}{\roman{enumi}}
\renewcommand{\labelenumi}{(\theenumi)}
\item The preceding functor $\varphi$ is a fiber functor on
$C^{\sim}$.

\item The ``restriction'' of $\varphi$ to $C$ is a fiber functor
on the site $C$ (\hyperref[IV.subsec6.3]{6.3}).

\item For every covering family $Y_{\alpha}\rightarrow Y$ of an
object $Y$ of $C$, every $i_{\circ}\in I$, and every morphism
$X_{i_{\circ}}\rightarrow Y$, there exist an $i\geq i_{\circ}$, an
$\alpha$, and a morphism $X_i\rightarrow Y_{\alpha}$ making
the diagram
\[
\xymatrix@C=2cm@R=1.25cm{
X_i \ar[d] \ar[r]& Y_{\alpha} \ar[d]\\
X_{{i}_{o}} \ar[r]     &  Y \quad . }
\]
commutative.

Indeed,\pageoriginale condition (iii) simply expresses the fact
that $\varphi| C$ carries covering families to covering
families. The implications (i) $\Rightarrow$ (ii) $\Rightarrow$
(iii) are trivial, and it remains to prove (iii) $\Rightarrow$ (i).
Since the functor $\varphi$ is manifestly left exact, it
remains to prove (\hyperref[IV.prop4.9.4]{4.9.4}) that it carries covering
families $F_{\alpha}\rightarrow F$ in $C^{\sim}$ to covering
families, \ie that for every $i_{\circ}$ and every $x_{\circ}\in
F(X_{i_{\circ}})$, there exist an $i\geq i_{\circ}$, an $\alpha$, and a
$y\in F_{\alpha}(X_i)$ such that $x_{\circ}$ and $y$ have the same image
in $F(X_i)$. Since the family $F_{\alpha}\rightarrow F$ is covering, there
is a covering family $Z_{\gamma}\rightarrow X_{i_{\circ}}$
of $X_{i_{\circ}}$ such that, for every $\gamma$, the inverse image of
$x_{\circ}$ in $F(Z_{\gamma})$ lifts to an element of some
$F_{\alpha}(Z_{\gamma})$. By hypothesis (iii), there exist an $i\geq
i_{\circ}$ and an $X_{i_{\circ}}$-morphism $X_i\rightarrow
Z_{\gamma}$. The image of $x_{\circ}$ in $F(X_i)$ therefore lifts
to an element $y$ of some $F_{\alpha}(X_i)$, \cqfd
\end{enumerate}

\begin{enonce*}[remark]{Exercise 6.9}\etiquette[6.9]{IV.exer6.9}
Let $C\in𝒰$ be a site, and let $\varphi$ be a fiber functor on
$C$, pro-represented by an object $(X_i)_{i\in I}$ of $\Pro(C)$. For
every index $i\in I$, every object $Y$ of $C$, every morphism
$u:X_i\rightarrow Y$, and every covering sieve $R$ on $Y$, choose
an index $j\geq i$ such that the composite $X_j\rightarrow
X_i\xrightarrow{u}Y$ factors through $R$. For fixed $i$, let
$J(i)$ be the set formed by $i$ and by the $j$ obtained as $Y$, $u$,
and $R$ vary; for a subset $I'$ of $I$, likewise let $J(I')$ be the
union of the $J(i)$ for $i\in I'$. On the other hand, if $I'$ is a
finite subset of $I$, let $m(I')\in I$ be an upper bound of $I'$, and for
an arbitrary subset $I'$ of $I$, let $M(I')$ be the set of all $m(I'')$,
where $I''$ ranges over the set of finite subsets
of\pageoriginale $I'$. We assume that the function
$I'\mapsto m(I')$ is chosen so that, when $I'$ consists of a single element
$i$, one has $m(I')=i$; this implies that, for every subset $I'$ of
$I$, one has $I'\subset M(I')$. Consider the iterates $(MJ)^n$ of
the map $MJ$ from the set of subsets of $I$ to itself,
and, for every subset $I'$ of $I$, let $P(I')$ be the union of the
$(MJ)^n(I')$.
\begin{enumerate}
\renewcommand{\theenumi}{\alph{enumi}}
\renewcommand{\labelenumi}{\theenumi)}
\item Show that, for every subset $I'$ of $I$, $P(I')$ is a
filtered subset of $I$ for the induced order, and that $I$ is the
increasing filtered union of the $P(I')$ as $I'$ ranges over the set
of \emph{finite} subsets of $I$.

\item Show that there is a cardinal $c\in 𝒰$, depending
only on the cardinal $\card \text{fl}(C)=a$, such that, for every
finite subset $I'$ of $I$, one has $\card(P(I'))\leq c$. (If $a$ is infinite,
one may take $c=2^{a}$.)

\item Show, using \hyperref[IV.subsubsec6.8.7]{6.8.7}, that for every subset $I'$ of $I$
the functor $\varphi_{I'}$ on $C$, pro-represented by the pro-object
$(X_i)_{i\in P(I')}$, is a fiber functor on $C$, and that the fiber
functor $\varphi$ is the filtered inductive limit of the functors
$\varphi_{I'}$ as $I'$ ranges over the set of finite subsets
of $I$.

\item Deduce that there is a small full subcategory $\Phi$
of $\Ffib(C)\approxeq \Ffib(C^{\sim})$ such that every object of $\Ffib(C)$
is a filtered inductive limit of objects of $\Phi$. (If $a$ is
infinite, one may take $\Phi$ to have cardinal $\leq 2^{2^{a}}$; if $a$
is finite, one may of course take $\Phi=C$.)
\end{enumerate}
\end{enonce*}

\begin{enonce*}[remark]{Exercise 6.10}\etiquette[6.10]{IV.exer6.10}
Let $C$ be a $\mathcal{U}$-site, and let $D$ be a
$\sheaf{\mathcal{U}}$-category in which small filtered $\varinjlim$
are representable; let $\Fais(C,D)$ be the category of
sheaves on $C$ with values in $D$ (\hyperrefext[II][II.def6.1]{6.1}). By
extending \eqhyperref[IV.eq6.8.3]{6.8.3}, define a ``fiber functor''
\begin{equation*}
F\longrightarrow F_p:\Fais(C,D)\longrightarrow D{\quoi}.\tag{6.10.1}
\end{equation*}\etiquette[6.10.1]{IV.eq6.10.1}
If finite $\varprojlim$ are representable in $D$ and
commute with filtered $\varinjlim$, then the preceding functor
is left exact.
\end{enonce*}

% END INCLUDED SGA4_English_current_buildable_source_20260906/source/04/drafts/SGA4_IV_LINES_3864_4120_English.texfrag


% BEGIN INCLUDED SGA4_English_current_buildable_source_20260906/source/04/drafts/SGA4_IV_LINES_4121_4475_English.texfrag
\section{Examples of fiber functors and points of topoi}\etiquette[7]{IV.sec7}\pageoriginale

\subsection{The case of $\Top(X)$ for a topological space
  $X$}\etiquette[7.1]{IV.subsec7.1}

Let $x$ be a point of $X$. Regard the singleton $\{x\}$ as a
topological space, and consider the inclusion
$$
i_x:\{x\}\longrightarrow X{\quoi}.
$$
By \hyperref[IV.subsubsec4.1.1]{4.1.1}, it defines a morphism of topoi
(defined up to unique isomorphism)
\begin{equation*}
\Top(i_x):\Top(\{x\})\longrightarrow \Top(X){\quoi}.\tag{7.1.1}
\end{equation*}\etiquette[7.1.1]{IV.eq7.1.1}
On the other hand, one may identify $P$ with $\Top(\{x\})$
(\hyperref[IV.subsec2.2]{2.2}), and one thereby obtains a point $p_x$ of
$\Top(X)$:
\begin{equation*}
p_x:P\longrightarrow \Top(X){\quoi},\tag{7.1.2}
\end{equation*}\etiquette[7.1.2]{IV.eq7.1.2}
defined by $x$ up to unique isomorphism. The associated fiber functor
is given by the familiar formula $[TF]$, which moreover follows
immediately from \hyperrefext[I][I.prop5.1]{5.1}:
\begin{equation*}
p^{\ast}_x(F)=F_x=F_{p_x}=\varinjlim_{U\ni x}F(U){\quoi},\tag{7.1.3}
\end{equation*}\etiquette[7.1.3]{IV.eq7.1.3}
the limit being taken over the open neighborhoods $U$ of $x$ in
$X$.

Since $\Top(X)$ is known to be isomorphic to $\Top(X_{\sob})$
(\hyperref[IV.subsubsec4.2.1]{4.2.1}), one sees more generally that every point
of $X_{\sob}$, \ie every irreducible closed subset $Z$ of $X$,
defines a point of $\Top(X)$, or equivalently a fiber functor on
$\Top(X)$, which we shall still denote by $F\mapsto F_Z$. Returning to
its definition by formula \eqhyperref[IV.eq7.1.3]{7.1.3} on $X_{\sob}$, one obtains
\begin{equation*}
F_Z=\varinjlim_{U\cap Z\neq \phi} F(U){\quoi},\tag{7.1.4}
\end{equation*}\etiquette[7.1.4]{IV.eq7.1.4}
the\pageoriginale inductive limit being taken over the open subsets
$U$ of $X$ which meet $Z$.

\setcounter{subsubsection}{4}
\subsubsection{}\etiquette[7.1.5]{IV.subsubsec7.1.5}
It is well known that the fiber functors $F\mapsto F_x$ $(x\in X)$
already form a conservative family. This is especially clear from
the interpretation of sheaves on $X$ in terms of étale spaces, since
the fiber of $F$ at $x$ is then simply its fiber in the sense of
spaces fibered over $X$. Note that the index set of the conservative
family of fiber functors under consideration is $X$, and hence is
$𝒰$-small.

\subsubsection{}\etiquette[7.1.6]{IV.subsubsec7.1.6}
By \hyperref[IV.subsubsec4.2.3]{4.2.3}, the category $\Ppoint(\Top(X))$
is equivalent to the category associated with the set $X_{\sob}$
ordered by the specialization relation. In other words, every fiber
functor on $\Top(X)$ is isomorphic to a fiber functor
$F\mapsto F_Z$, where $Z$ is a uniquely determined element of
$X_{\sob}$, \ie a uniquely determined irreducible closed subset
of $X$; moreover, if $Z$, $Z'\in X_{\sob}$, then the set
$\Hom(F_Z,F_{Z'})$ is either empty or consists of a single point, the
latter case occurring if and only if $Z$ is a specialization of $Z'$
in $X_{\sob}$, \ie if and only if $Z\subset Z'$ as a subset of
$X$. When $X$ itself is sober, one may replace $X_{\sob}$ by $X$
itself in these statements.

Note that the automorphism group of a fiber functor of $\Top(X)$
(opposite to the automorphism group of the corresponding point of
$\Top(X)$) is always the trivial group (more generally,
\hyperref[IV.subsubsec4.2.3]{4.2.3} tells us the same thing about the
automorphism group of every morphism of topoi associated with
topological spaces). This is a phenomenon very special to the case
under consideration: see the example in \hyperref[IV.subsec7.2]{7.2}, as\pageoriginale
well as VIII 7. In the case of topoi associated with ``moduli
problems'' (\cf for example [10]), the automorphism groups of fiber
functors moreover have a remarkable interpretation as the
automorphism groups of the algebraic structures (over algebraically
closed fields) which one proposes to classify.

\subsubsection{}\etiquette[7.1.7]{IV.subsubsec7.1.7}
We have just seen how one can reconstruct a sober space $X$ (or the
sober space $X_{\sob}$ associated with an arbitrary topological
space), at least as a set ordered by the specialization relation,
from the topos $\Top(X)$ (which it even suffices to know up to
equivalence), as the set of isomorphism classes of ``points'' of
$\Top(X)$. From what was said in \hyperref[IV.subsec2.1]{2.1}, one also
reconstructs the topology of $X$, \ie the family of its open subsets,
as follows: for every subobject $U$ of the terminal object $e_E$ of
$E$, let $\Pt(U)$ be the set of $x\in X$ such that $U_x\neq\phi$
(which moreover identifies, by \eqhyperref[IV.eq6.7.2]{6.7.2}, with the set of
isomorphism classes of points of the induced topos $E_{/U}$). Then
$U\mapsto \Pt(U)$ is an isomorphism of ordered sets from the set of
subobjects of $e_E$ onto the set of open subsets of $X$.

\subsubsection{}\etiquette[7.1.8]{IV.subsubsec7.1.8}
The determination in \hyperref[IV.subsubsec7.1.6]{7.1.6} of the category of
points of the topos $\Top(X)$ defined by a topological space $X$
leads to the following terminology for points $p$, $p'$ of an
arbitrary topos $E$: one says that $p$ is a \emph{specialization}
of $p'$, or that $p'$ is a \emph{generization} of $p$, when there
exists a morphism from $p'$ to $p$ (in the sense of the category
$\Ppoint(E)$ of \hyperref[IV.def6.1]{6.1}). These relations are still
transitive, but, using \hyperref[IV.subsubsec6.8.6]{6.8.6}, one easily finds
examples in which $p$ and $p'$ are each a specialization of\pageoriginale the
other, without $p$ and $p'$ being isomorphic (much less equal). The
category of generizations of a point $p$ of the topos $E$, by which
we mean the category $\Ppoint(E)_{/p}$, plays in some respects a role
analogous to that of passing to the localization of a scheme at a
point. This role will be made precise in
\sisi{Chap. VI}{\hyperrefext[VI][VI.subsec8.4]{8.4}} by the construction of the topos obtained by
localizing $E$ at the point $p$, whose category of points is
canonically equivalent to the category of generizations of $p$.

\begin{enonce*}[remark]{Exercise 7.1.9}\etiquette[7.1.9]{IV.exer7.1.9}
Let $E$ be a $𝒰$-topos. Prove that $E$ is equivalent to a
topos of the form $\Top(X)$, where $X$ is a topological space $\in
𝒰$, if and only if it satisfies the following two conditions:
\begin{enumerate}
\renewcommand{\theenumi}{\alph{enumi}}
\renewcommand{\labelenumi}{\theenumi)}
\item The family of subobjects of the terminal object $e_E$ is
generating for the topos $E$.

\item $E$ has enough points (\hyperref[IV.prop6.5]{6.5}). (This condition is not
superfluous: \cf \hyperref[IV.subsec7.4]{7.4}.)
\end{enumerate}

Compare with \hyperref[IV.exer7.8]{7.8} c).
\end{enonce*}

\begin{enonce*}[remark]{Exercise 7.1.10}\etiquette[7.1.10]{IV.exer7.1.10}
Let $X$ be a topological space endowed with an operator group $G$
$(X,G\in𝒰)$, and let $E=\Top(X,G)$ (\hyperref[IV.subsec2.3]{2.3}).
\begin{enumerate}
\renewcommand{\theenumi}{\alph{enumi}}
\renewcommand{\labelenumi}{\theenumi)}
\item Show that, for every object $(X',G)$ of $E$, the induced topos
$E_{/(X',G)}$ is canonically equivalent to the topos $\Top(X',G)$
(compare \hyperref[IV.subsec5.7]{5.7}).

\item When $G$ acts properly and freely on $X'$ (Bourbaki,
Gen. Top. Chap. \hyperrefext[III][III.sec4]{4}), show that the morphism of
spaces with operators $(X',G)\rightarrow (X'/G,e)$ induces
(\hyperref[IV.subsubsec4.1.2]{4.1.2}) an equivalence of topoi
$$
\Top(X',G)\longrightarrow \Top(X'/G){\quoi}.
$$

\item Deduce\pageoriginale from b) that there exists an object $Z$ of
$E$ such that the induced topos $E_{/Z}$ is equivalent to $\Top(X)$,
the localization functor $E\rightarrow E_{/Z}$ being isomorphic to
the functor that ``forgets the actions of $G$''. (Imitate the
argument of \hyperref[IV.subsubsec5.8.3]{5.8.3}.)

\item Deduce from c) that every fiber functor on $E$ is induced, via
the functor that ``forgets the actions of $G$'', by a fiber functor
on $\Top(X)$, and can therefore be defined by a point of $X_{\sob}$.

\item Determine the structure of the category $\Ppoint(E)$ (up to
equivalence) in terms of the space with operators $(X_{\sob},G)$.
Conclude in particular that two points of $X_{\sob}$ define
isomorphic fiber functors on $E$ if and only if they are conjugate
under the action of $G$.
\end{enumerate}
\end{enonce*}

\begin{enonce*}[remark]{Exercise 7.1.11}\etiquette[7.1.11]{IV.exer7.1.11}
Let $X\in 𝒰$ be a topological space. Prove that the
$𝒱$-topos $\TOP(X)$ (\hyperref[IV.subsec2.5]{2.5}) has enough points.
More precisely, for every object $X'$ of $\TOP(X)$ and every $x'\in
X'$, define a fiber functor $F\mapsto F_{X',x'}$ on $\TOP(X)$ which
depends only on the ``germ'' of the space $(X',x')$ over $X$, and
prove that this family of fiber functors is conservative (and indexed
by a $𝒱$-small index set). Show, using a suitable filtered projective
system $(X'_i,x'_i)_{i\in I}$, with $I$ not $𝒰$-small, that one does
not obtain in this way all the fiber functors of the $𝒱$-topos
$\TOP(X)$. Show that the set of isomorphism classes of such fiber
functors does not have cardinality $\in 𝒱$.
\end{enonce*}

\subsection{Points of a classifying topos
  $B_G$}\etiquette[7.2]{IV.subsec7.2}\pageoriginale

Let $E$ be a $𝒰$-topos and $G$ a group of $E$, whence a classifying
topos $B_G$ (\hyperref[IV.subsec2.4]{2.4}). If $e$ is the unit group, the
morphisms $e\rightarrow G\rightarrow e$ define
(\hyperref[IV.subsec4.5]{4.5}) morphisms of topoi (taking into account that
$B_e\approx E$)
\begin{equation*}
E\longrightarrow B_G \longrightarrow E{\quoi},\tag{7.2.1}
\end{equation*}\etiquette[7.2.1]{IV.eq7.2.1}
whose composite is isomorphic to $\id_E$; passage to the categories
of points therefore gives functors
\begin{equation*}
\Ppoint(E)\longrightarrow \Ppoint(B_G)\longrightarrow
\Ppoint(E){\quoi},\tag{7.2.2}
\end{equation*}\etiquette[7.2.2]{IV.eq7.2.2}
whose composite is isomorphic to the identity. Use the fact that the
first morphism of topoi in \eqhyperref[IV.eq7.2.1]{7.2.1} identifies with a
localization morphism relative to the object $E_G=G$ of $B_G$
(\hyperref[IV.subsec5.8]{5.8}); by \eqhyperref[IV.eq6.7.2]{6.7.2}, it follows that the
first functor in \eqhyperref[IV.eq7.2.2]{7.2.2} identifies with the natural
``inclusion'' functor
\begin{equation*}
\Ppoint(B_G)_{/\widehat{E}_G}\longrightarrow
\Ppoint(B_G){\quoi},\tag{7.2.3}
\end{equation*}\etiquette[7.2.3]{IV.eq7.2.3}
where $\widehat{E}_G$ denotes the presheaf $p'\mapsto (E_G)_{p'}$ on
$\Ppoint(B_G)$. Since $E_G$ plainly covers the terminal object of
$B_G$, its fibers $(E_G)_p$ are nonempty, from which it follows that
\eqhyperref[IV.eq7.2.3]{7.2.3} is essentially surjective; this means that
\emph{every fiber functor on $B_G$ is induced (up to isomorphism) by
a fiber functor on $E$, via the functor that ``forgets the actions of
$G$'' $B_G\rightarrow E$.} One concludes in particular that
\emph{if $E$ has enough points (\resp a small conservative family of
points),} then the same is true of $B_G$.

\begin{enonce*}[remark]{Remark 7.2.4}\etiquette[7.2.4]{IV.rem7.2.4}
One can go further and determine (up to equivalence) the structure of
the category $C'=\Ppoint(B_G)$ in terms of the category
$C=\Ppoint(E)$ and the presheaf of groups
$$
\widehat{G}:p\longmapsto G_p: C^{\circ}\longrightarrow
(\text{Groups})
$$
on\pageoriginale it. Indeed, define a new category $C_1$, having the
same objects as $C$, and such that for two objects $p$, $q$ of $C$,
one has
$$
\Hom_{C_1}(p,q)=\Hom_C(p,q)\times \widehat{G}(p){\quoi},
$$
composition of arrows in $C_1$ being induced by that in $C$ and by
the group law of the presheaf $\widehat{G}$. The datum of a functor
$\varphi_1:C_1\rightarrow D$ into an arbitrary category $D$ is then
equivalent to the datum of a functor $\varphi:C\rightarrow D$
endowed with an action of the presheaf $\widehat{G}$ on the latter
(\ie the datum, for every $p\in \ob C$, of an action of
$\widehat{G}(p)$ on $\varphi(p)$, satisfying an evident functoriality
condition as $p$ varies). Using this observation, one defines a
canonical functor
\begin{equation*}
\varphi_1:C_1\longrightarrow C'{\quoi},\tag{7.2.4.1}
\end{equation*}\etiquette[7.2.4.1]{IV.eq7.2.4.1}
corresponding to the functor $\varphi:C\rightarrow C'$ of
\eqhyperref[IV.eq7.2.2]{7.2.2}, which associates with every point $p$ of $E$
the point $\varphi(p)$ induced on $B_G$, together with the natural
actions of $\widehat{G}(p)=G_p$ on the latter, arising from the
actions of $G_p$ on the $X_p$ as $X$ ranges over $B_G$. We leave to
the reader the proof that \eqhyperref[IV.eq7.2.4.1]{7.2.4.1} is an equivalence
of categories, using the structure \eqhyperref[IV.eq7.2.3]{7.2.3} of the first
functor in \eqhyperref[IV.eq7.2.2]{7.2.2}, and noting that the natural
inclusion functor $C\rightarrow C_1$ has an analogous structure,
relative to the inverse image $P_1$ of the presheaf
$P'=\widehat{E}_G$ on $C'$.
\end{enonce*}

\setcounter{subsubsection}{4}
\subsubsection{}\etiquette[7.2.5]{IV.subsubsec7.2.5}
One\pageoriginale concludes in particular that the \emph{isomorphism
classes of points of $B_G$ correspond exactly to the isomorphism
classes of points of $E$: every point of the classifying topos $B_G$
is induced, up to a} non-unique \emph{isomorphism, by a point of
$E$}. In particular, when $E$ is the punctual topos $P$, and hence
$G$ is an ordinary group, the category $\Ppoint(B_G)$ is a connected
groupoid with fundamental group $G$: every fiber functor on $B_G$ is
isomorphic (noncanonically) to the functor $\omega$ that ``forgets
the actions of $G$'', and the monoid of endomorphisms of the latter
functor is the group $G$ (hence every endomorphism of $\omega$ is an
automorphism).

\begin{enonce*}[remark]{Exercise 7.2.6}\etiquette[7.2.6]{IV.exer7.2.6}
\begin{enumerate}
\renewcommand{\theenumi}{\alph{enumi}}
\renewcommand{\labelenumi}{\theenumi)}
\item Let (Tors) be the category of pairs $(\Gamma,P)\in 𝒰$,
where $\Gamma$ is a group and $P$ is a right torsor under $\Gamma$.
The functor $(\Gamma,P)\mapsto \Gamma$
\begin{equation*}
(\Tors)\longrightarrow (\Groupes)\tag{7.2.6.1}
\end{equation*}\etiquette[7.2.6.1]{IV.eq7.2.6.1}
is a cofibered functor. With the notation of \hyperref[IV.rem7.2.4]{7.2.4},
consider the functor
$$
\widehat{G}:C^{\circ}\longrightarrow (\Groupes){\quoi},
$$
and let $D'$ be the cofibered category over $C^{\circ}$ obtained by
inverse image from the cofibered category \eqhyperref[IV.eq7.2.6.1]{7.2.6.1}, and
let $D$ be the corresponding category fibered over $C$, having the
same fiber categories as $D'$, so that for $p\in \ob C$ one has
$$
D_p=\mbox{ category of right torsors under $G_p$}{\quoi}.
$$
Construct canonical functors
\begin{equation*}
D\longrightarrow C'=\Ppoint(B_G)\quav C'\longrightarrow
D{\quoi},\tag{7.2.6.2}
\end{equation*}\etiquette[7.2.6.2]{IV.eq7.2.6.2}
\emph{quasi-inverse to one another.} In particular, conclude that
the category of points of $B_G$ above a given point $p$ of $E$ is
canonically equivalent to the category of right torsors under the
group $G_p$.

\item Let\pageoriginale $𝒢=(G_i)_{i\in I}$ be a strict
projective system of groups of the $𝒰$-topos $E$ ($I\in
𝒰$ a filtered preordered set), whence a classifying topos
$B_{𝒢}$, by imitating the definition of
\hyperref[IV.subsubsec2.7.1]{2.7.1}. Determine the category of points of
$B_{𝒢}$ by imitating a). (NB. Replace the categories
$(\Tors)$ and $(\Groupes)$ by the categories of projective systems
indexed by $I$ in these categories.) Conclude that if $I$ admits a
countable cofinal subset, then every fiber functor on $B_{𝒢}$ is
isomorphic to a functor induced by a fiber functor on $E$ (via the
functor that ``forgets the actions of $G$''), the latter being
determined up to non-unique isomorphism.

\item Taking $E$ to be the punctual topos, so that $𝒢$ is an
ordinary pro-group, show that the conclusion of b) is likewise valid
if $I$ is arbitrary but $𝒢$ is profinite (\ie the $G_i$ are
finite): every fiber functor is isomorphic to the functor that
``forgets the actions of $G$''.

\item Still taking the case where $E$ is the punctual topos, give an
example of a fiber functor on $B_{𝒢}$, for a suitable projective
system $𝒢=(G_i)_{i\in I}$, which is not isomorphic to the functor
that ``forgets the actions of $G$''.

\item Regard $𝒢=(G_i)$ as a group of the topos
$\widehat{I}$, and let $𝒯$ be a \emph{gerbe} on $I$ with
\emph{band} $G$ [\hyperref[IV.sec3]{3}]. Show how one can
``twist'' the classifying topos $B_G$ by means of the gerbe $T$, to
obtain a topos $B^T_{𝒢}$, an inductive limit of subcategories
$B^T(i)$ equivalent (noncanonically) to the classifying topoi
$B_{G_i}$. Show that the topos $B^T_{𝒢}$ admits a fiber functor if
and only if the gerbe $T$ is ``neutral'', by establishing an\pageoriginale
equivalence of categories between the category of fiber functors on
$B^T_{𝒢}$ and the category of sections of $T$ over $I$. Conclude,
if the $G_i$ are commutative, so that gerbes on $I$ with band $𝒢$
are classified by
$\displaystyle{\varprojlim_I}^{(2)}G_i=H^{2}(I,𝒢)$
(\loccit), that there is an example of a topos (not ``empty''
(\hyperref[IV.subsec2.2]{2.2})) of the form $B^T_{𝒢}$ which has no
points. (Take an example in which
$\displaystyle{\varprojlim_I}^{(2)}G_i\neq 0$.)
\end{enumerate}
\end{enonce*}

% END INCLUDED SGA4_English_current_buildable_source_20260906/source/04/drafts/SGA4_IV_LINES_4121_4475_English.texfrag


% BEGIN INCLUDED SGA4_English_current_buildable_source_20260906/source/04/drafts/SGA4_IV_LINES_4476_4830_English.texfrag
\subsection{Points of presheaf topoi $\widehat{C}$; examples of
$𝒰$-presheaf topoi $\widehat{C}$ whose category
of points is not equivalent to a small
category}\etiquette[7.3]{IV.subsec7.3}

Let $E$ be a $𝒰$-topos, and let $(\varphi_i)_{i\in I}$ be a filtered
family of fiber functors on $E$. It follows from the exactness
properties of the functors $\varinjlim$
(\sisi{I\quad and}{\hyperrefext[I][I.prop2.8]{2.8}}) that the functor
$\varphi=\varinjlim\varphi_i$ is also a fiber functor:
\emph{every filtered inductive limit of fiber functors is a fiber
functor.} Consequently, the category $\Ffib(E)$ admits $𝒰$-filtered
inductive limits (and the inclusion functor
$\Ffib(E)\rightarrow \Hhom(E,(𝒰\text{-}\Ens))$ commutes with them);
in other words, the category $\Ppoint(E)$ admits $𝒰$-filtered
projective limits. This fact was already used in Exercise
\sisi{\hyperref[IV.exer7.1.10]{7.1.10}}{\hyperref[IV.exer7.1.11]{7.1.11}}
to give an example of ``large'' categories of fiber functors.

\hypertarget{sga4.exp.IV.subsubsection.7.3.3}{}Considerably simpler examples can be constructed, with presheaf
topoi of the form $E=\widehat{C}$, $C\in𝒰$, by using
\eqhyperref[IV.eq6.8.6.1]{6.8.6.1}. This immediately gives examples
in which $\Ffib(\widehat{C})$ is not equivalent to a category
$\in𝒰$, \ie in which the cardinal of the set of isomorphism classes
of objects of this $𝒰$-category is not $\in𝒰$. This means
that the category of pro-objects $\Pro-(C)$ is not equivalent to a
category $\in 𝒰$. For example, it suffices to take
$D=C^{\circ}$ to be the category of finite sets of the form
$[0,n]$, where $n\geq 0$ is an\pageoriginale integer; in
this case $\Ind(D)\simeq \Pro(C)^{\circ}$ is equivalent to the
category of nonempty sets. The topos $E$ is then the familiar
category of \emph{cosimplicial sets}. One could also take $C$ to be
the category of nonempty finite sets; one then finds that the
category of points of the presheaf topos $\widehat{C}$ of simplicial
sets is equivalent to the category of profinite sets, or equivalently
to the category of totally disconnected compact spaces.

\subsection{Nonempty topoi without points}\etiquette[7.4]{IV.subsec7.4}

The following example is due to $P$. Deligne. (For another example,
\cf \hyperref[IV.exer7.2.6]{7.2.6} e).) Let $K$ be a compact space
equipped with a measure $\mu$, and let $U$ be the ordered set of
measurable subsets of $K$, modulo sets of measure zero. Make $U$
into a category $\underline{U}$ with $\ob \underline{U}=U$, the
morphisms of $\underline{U}$ being the ``inclusion morphisms''
between elements of $U$. Make $\underline{U}$ into a site by taking
the pretopology for which $\Cov(E)$ (for $E\in U$) consists of the
\emph{countable} families of elements $E_i$ of $U$ bounded above by
$E$, such that $E$ is the union of the $E_i$ modulo a set of measure
zero. This gives a topos $\Top(\mu)=\underline{U}^{\sim}$ whose set
of subobjects of the terminal object is a generating family (a topos
which appears to have escaped the attention of probabilists). This
topos is an ``empty topos'' (\hyperref[IV.subsec2.2]{2.2}) if and
only if $\mu=0$. On the other hand, the category of points of this
topos is equivalent to the discrete category defined by the set of
points $x\in K$ such that $\mu(\{x\})\neq 0$ (proof left to the
reader). It is therefore empty if $K$ has no such points, for
example if $K$ is the unit interval in the real line, with the
measure induced by Lebesgue measure.

\begin{enonce*}[remark]{Exercise 7.5}\etiquette[7.5]{IV.exer7.5}
(\emph{\sisi{Karoubian}{karoubian} categories\pageoriginale and
morphisms of topoi})\footnote[1]{Compare
\hyperrefext[I][I.exer8.7.8]{8.7.8}.}
\begin{enumerate}
\renewcommand{\theenumi}{\alph{enumi}}
\renewcommand{\labelenumi}{\theenumi)}
\item Let $C$ be a category, let $X$ be an object of $C$, and let
$p$ be an endomorphism of $X$. We say that $p$ is a
\emph{projector} if $p^2=p$. Prove that, if $p$ is a projector,
$\Ker(\id_X,p)$ is representable if and only if
$\Coker(\id_X,p)$ is representable, and that the two objects of $C$
so obtained are canonically isomorphic. We then say that the
projector $p$ \emph{admits an image,} and
$\Ker(\id_X,p)\simeq \Coker(\id_X,p)$ is called the \emph{image of
the projector $p$}; according to the context, it is identified with
either a subobject or a quotient object of $C$. An object isomorphic
to $\text{Im\,} p$, for a suitable projector $p$ in $X$, is called
a \emph{direct factor of the object $X$ of $C$.} We say that $C$ is
a \emph{category with direct factors}, or a
\emph{\sisi{Karoubian}{karoubian} category,} if every projector in
an object of $C$ admits an image. If $F:C\rightarrow C'$ is a
functor that commutes with kernels or cokernels, then $F$ sends a
projector admitting an image to a projector admitting an
image\nde{No condition on $F$ is needed here: every functor
preserves projectors admitting images.}.

\item Show that, for every category $C$, one can find a functor
$\varphi:C\rightarrow \kar (C)$ from $C$ to a Karoubian category
(determined up to equivalence) such that, for every Karoubian
category $C'$, the functor
$$
f\mapsto f\circ \varphi:\Hhom(\kar(C),C')\longrightarrow \Hhom(C,C')
$$
is an equivalence of categories; $\kar(C)$ is called the
\emph{Karoubi envelope} of the category $C$. (Hint: take
$\ob \kar (C)$ to be the set of pairs $(X,p)$, with $X\in \ob C$
and $p$ a projector in $X$, and take $\Hom((X,p),(Y,q))$ to be the
subset of $\Hom(X,Y)$ consisting of those $f:X\rightarrow Y$ such
that $f=qfp.$) Show that $\varphi$ is fully faithful.

\item Let $F:E\rightarrow E'$ be a functor from one category to
another. Show that if $F$ commutes with a given type of inductive or
projective limits,\pageoriginale then the same is true of every
direct factor of $F$. In particular, if $E$ and $E'$ are
$𝒰$-topoi and $F$ is an inverse-image functor for a morphism of
topoi $E'\rightarrow E$, then the same is true of every direct
factor of $F$; consequently, the category $\Hhomtop(E',E)$ is
Karoubian, and in particular the category $\Ppoint(E)$ is
Karoubian.

\item Show that, for every $𝒰$-category $C$, the category
$\Ind(C)$ is Karoubian (where $\Ind(C)$ is formed using inductive
systems in $C$ indexed by a filtered preordered set $I\in 𝒰$). (If
$C\in 𝒰$, use, for example, the fact
(\hyperref[IV.subsec7.3]{7.3}) that $\Ind(C)$ is equivalent to the
category
$\Ffib((C^{\circ})^{\widehat{\phantom{,}}})=\Ppoint((C^{\circ})^{\widehat{\phantom{,}}})^{\circ}.$)
Deduce another construction of $\kar(C)$ as the full subcategory of
$\Ind(C)$ formed by the images in $\Ind(C)$ of the projectors of
objects $X$ of $C$.
\end{enumerate}
\end{enonce*}

\begin{enonce*}[remark]{Exercise 7.6}\etiquette[7.6]{IV.exer7.6}
(\emph{Essential morphisms of topoi, essential points}).
\begin{enumerate}
\renewcommand{\theenumi}{\alph{enumi}}
\renewcommand{\labelenumi}{\theenumi)}
\item Let $f:E\rightarrow F$ be a morphism of topoi. Show that the
following conditions are equivalent:
\begin{itemize}
\item[i)] $f_{!}$ exists, \ie $f^{\ast}$ admits a left adjoint.

\item[ii)] $f^{\ast}$ commutes with $𝒰$-projective limits.

\item[${\rm ii')}$] $f^{\ast}$ commutes with $𝒰$-products.
\end{itemize}
We then say that $f$ is an \emph{essential morphism} from the topos
$E$ to the topos $F$.

Show that if $f$ satisfies these conditions, then the same is true
of every direct factor of $f$. (Use \hyperref[IV.rem1.8]{1.8} and
\hyperref[IV.exer7.5]{7.5} c).)

\item Let $E$ be a topos. An \emph{essential point of the topos}
$E$ is a point $p:P\rightarrow E$ of $E$ such that $p_{!}$ exists.
Show that if $E=\Top(X)$, where $X$ is a topological space, and if
$p$ is the point of $E$ defined by an $x\in X$, then $p$ is
essential\pageoriginale if and only if $x$ has a smallest open
neighborhood $U$ (which means, if the points of $X$ are closed, that
$x$ is an \emph{isolated point} of $X$). For a morphism of topoi
$f:E\rightarrow F$ to be essential (a)), it is necessary that $f$
send essential points to essential points, and this condition is
also sufficient when $E$ has enough essential points (\ie when the
family of these points is conservative) (\cf d)).

\item A point $p$ of the topos $E$ is essential if and only if the
associated fiber functor is representable by an object $X$ of $E$.
For the covariant functor $\varphi:E\rightarrow (\Ens)$ represented
by a given object $X$ of $E$ to be a fiber functor (\ie to define a
point, necessarily an essential point, of $E$), it is necessary and
sufficient that $X$ be nonempty and connected
(\hyperref[IV.subsubsec4.3.5]{4.3.5}), and projective (\ie that the
functor $\varphi$ send epimorphisms to epimorphisms). (Hint: first
show that $\varphi$ commutes with coproducts if and only if $X$ is
nonempty and connected, and then use the criterion
\hyperref[IV.prop4.9.4]{4.9.4}.) Deduce an equivalence between the
category $\Ppointess(E)$ and the full subcategory $P$ of $E$ formed
by the nonempty connected projective objects.

\item Show that the topology on $P$ induced by that of $E$ is the
chaotic topology. Deduce the equivalence of the following conditions
on $E$ (due to J. E. ROOS [12 c), prop. 1]): (i) The family of
essential points of $E$ is conservative; (ii) The full subcategory
$P$ of $E$ formed by the nonempty connected projective objects is
generating; (iii)~$E$ is equivalent to a presheaf topos of the form
$\widehat{C}$, where $C$ is a category equivalent to a category
$\in 𝒰$ (or, if one prefers, $C\in 𝒰$). (Use $c$) and the fact that,
if $P$ is generating, the natural functor $E\rightarrow P^{\sim}$
is an equivalence of\pageoriginale categories.)

\item The category of essential points of a topos $E$ is Karoubian
(\cf \hyperref[IV.exer7.5]{7.5} c)). Let $C$ be a category
equivalent to a category $\in 𝒰$. Prove that the canonical fully
faithful functor (\hyperref[IV.subsubsec4.6.2]{4.6.2})
$$
C\longrightarrow \Ppoint(\widehat{C})
$$
factors through a functor
$$
C\longrightarrow \Ppointess(\widehat{C}){\quoi},
$$
which makes $\Ppointess(\widehat{C})$ a \emph{Karoubi envelope}
(\hyperref[IV.exer7.5]{7.5} b)) \emph{of} $C$. In particular, this
functor is an equivalence of categories if and only if the category
$C$ is Karoubian. (Hint: using $c$), prove that every essential point
of $\widehat{C}$ is isomorphic to a direct factor of some
$X\in \ob C$.)

\item Let $C$ and $C'$ be two categories equivalent to categories
$\in 𝒰$. Prove that the canonical fully faithful functor
(\hyperref[IV.subsubsec4.6.2]{4.6.2})
$$
\Hhom(C,C')\longrightarrow \Hhomtop(\widehat{C},\widehat{C}')
$$
takes its values in the full subcategory
$\Hhomtopess(\widehat{C},\widehat{C}')$ of essential morphisms
\emph{of topoi} (a)), and that the induced functor
$$
\Hhom(C,C')\longrightarrow \Hhomtopess(\widehat{C},\widehat{C}')
$$
is an equivalence of categories if and only if $C$ is empty or $C'$
is Karoubian. (Use $g$) below.)

\item Let $C$ be a category equivalent to a category $\in 𝒰$, and
let $E$ be a topos. Define an equivalence between the category
$\Hhomtopess(\widehat{C},E)$ of essential morphisms of topoi
$\widehat{C}\rightarrow E$, and the category
$\Hhom(C, \Ppointess(E))$. (Use $b$)\pageoriginale and the fact that
$\widehat{C}$ has enough essential points.)

\item Let $C$ be a finite category. Prove that every point of the
presheaf topos $\widehat{C}$ is essential, and that
$\Ppoint(\widehat{C})$ is equivalent to a finite category. (Note
that the Karoubi envelope of $C$ is finite, and, using
\eqhyperref[IV.eq6.8.6.1]{6.8.6.1}, reduce to proving that if $C$ is
a finite Karoubian category, then the canonical functor
$C\rightarrow \Pro(C)$ is an equivalence of categories.) Using b),
deduce that \emph{every} morphism of topoi
$\widehat{C}'\rightarrow \widehat{C}$ is essential, and that
$\Hhom(C,C')\rightarrow \Hhomtop(\widehat{C},\widehat{C}')$ is an
equivalence if and only if $C$ is empty or $C'$ is Karoubian.

\item Let $X$ be a topological space, and let
$f:\Top(X)\rightarrow P$ be the morphism of topoi induced by the
continuous map from $X$ to the one-point space. Show that $f$ is
essential if and only if the space $X$ is \emph{locally connected},
\ie if it satisfies the following condition: for every $x\in X$ and
every open neighborhood $U$ of $x$, the connected component of $x$
in $U$ is a neighborhood of $x$; or equivalently: for every open
subset $U$ of $X$, the connected components of $U$ are open. (\cf
\hyperref[IV.exer8.7]{8.7} l) for a generalization).
\end{enumerate}
\end{enonce*}

\begin{enonce*}[remark]{Exercise 7.7}\etiquette[7.7]{IV.exer7.7}
(\emph{Unusual points of a classifying topos})\footnote[1]{\cf also
\hyperref[IV.exer7.2.6]{7.2.6} d).}.
\begin{enumerate}
\renewcommand{\theenumi}{\alph{enumi}}
\renewcommand{\labelenumi}{\theenumi)}
\item Let $G$ be a monoid. The monoid $G$ contains an element $g$
such that $g\neq e$ and $g^2=g$ if and only if the classifying topos
$B_G$ admits an essential point that is not isomorphic to the
trivial point (corresponding to the fiber functor ``forget the
$G$-action''). (Use \hyperref[IV.exer7.6]{7.6} e)).

\item Let $G$ be the additive monoid of the integers $\geq 0$.
Show that every essential point of the classifying topos $B_G$ is
isomorphic to the trivial point. Construct a point of $B_G$ that is
not isomorphic to the trivial point. Show that the category
$\Ppoint(B_G)$ is not equivalent to a category $\in 𝒰$.
\end{enumerate}
\end{enonce*}

\begin{enonce*}[remark]{Exercise 7.8}\etiquette[7.8]{IV.exer7.8}
(\emph{Topology on $\Point(E)$, and topoi associated with ordered
sets}).
\begin{enumerate}
\renewcommand{\theenumi}{\alph{enumi}}
\renewcommand{\labelenumi}{\theenumi)}
\item Let $E$ be a topos. Denote by $\Point(E)$ the set of
isomorphism classes of points of $E$. For every open $U$ of $E$,
let $U'$ be the set\pageoriginale of classes of points $p$ of $E$
such that $U_p\neq \phi$. Show that $U\mapsto U'$ is an
order-preserving map from the ordered set $\Ouv(E)$ of opens of $E$
to the set of subsets of $X=\Point(E)$, and that this map commutes
with finite Inf and arbitrary $\Sup$. Deduce that the subsets of
$X$ of the form $U'$ define a topology on $X$ (called the
\emph{canonical topology on the set of classes of points of the
topos $E$}). Show that if $E$ has enough points, the map
$U\mapsto U'$ is an isomorphism from the ordered set $\Ouv(E)$ onto
the ordered set $\Ouv(X)$.

\item Let $\Oo\in 𝒰$ be an ordered set admitting finite Inf and
arbitrary $\Sup$, and hence defining a category $\cat(\Oo)$ having
finite products and arbitrary coproducts. We say that a family of
morphisms $U_i\rightarrow U$ with common target $U$ is covering if
$U$ is the $\Sup$ of the $U_i$. Suppose that coproducts in
$\cat(\Oo)$ are universal, \ie that in $\Oo$ the arbitrary $\Sup$
commute with Inf. This defines a topology on $\cat(\Oo)$, making
it a site, and hence a topos $\cat(\Oo)^{\sim}$, also denoted simply
by $\Oo^{\sim}$. Show that $\Oo$ is canonically isomorphic to
$\Ouv(\Oo^{\sim})$.

Show that the category $\Hhomtop(E,\Oo^{\sim})$ is equivalent to the
category associated with the ordered set of morphisms of ordered
sets $\Oo\rightarrow \Ouv(E)$ that commute with finite Inf and
arbitrary $\Sup$. (Use the criterion
\hyperref[IV.prop4.9.4]{4.9.4}.) If $\Oo'$ is an ordered set
satisfying the same conditions as $\Oo$, deduce that
$\Hhomtop(\Oo^{\sim}, {\Oo'}^{\sim})$ is equivalent to the category
associated with the ordered set of all morphisms of ordered sets
${\Oo}'\rightarrow\Oo$ that commute with finite Inf and
arbitrary $\Sup$.

\item Show\pageoriginale that the topos $E$ is equivalent to a topos
of the form $\Oo^{\sim}$, with $\Oo$ as in b), if and only if the
family of opens of $E$ is generating in $E$. It is equivalent to
some $\Top(X)$ with $X\in 𝒰$ if and only if it satisfies the
preceding condition and has enough points. Assuming the first
condition, express the second in terms of the structure of the
ordered set $\Oo=\Ouv(E)$, using the last assertion of b).

\item Define a canonical morphism of topoi
$$
E\longrightarrow \Ouv(E)^{\sim}{\quoi},
$$
which is universal (up to equivalence) for morphisms from $E$ to
topoi generated by their opens (\cf c)).

\item Let $X'$ be a \emph{small} subset of $X=\Point(E)$, equipped
with the topology induced by that of $X$. Define a canonical
morphism of topoi
$$
\Ouv(E)^{\sim}\longrightarrow \Top(X'){\quoi},
$$
which is an equivalence when the family $X'$ of points of $E$ is
conservative. Then deduce a canonical morphism of topoi
$$
E\longrightarrow \Top(X'){\quoi}.
$$
Give a universal characterization (up to equivalence) of this
morphism of topoi with respect to morphisms from the topos $E$ to
topoi of the form $\Top(Y)$. (Note in this connection that, up to
equivalence, $\Top(X')$ does not depend on the chosen small
conservative family of fiber functors, or equivalently
(\hyperref[IV.subsec4.2]{4.2}), that $X'_{\sob}$ is independent, up
to homeomorphism, of the choice of such a family.)

\item Suppose\pageoriginale that the category $\Ppoint(E)$ is not
equivalent to a small category (\cf \hyperref[IV.subsec7.3]{7.3}),
and that $E$ admits a small conservative family of fiber functors.
Show that, if such a family is taken sufficiently large, the subset
$X'$ of $\Ppoint(E)$ that it defines is not a sober topological
space for the induced topology.

\item For the relations with Exercise \hyperref[IV.exer7.6]{7.6},
\cf \hyperref[IV.sec8]{8}.
\end{enumerate}
\end{enonce*}

% END INCLUDED SGA4_English_current_buildable_source_20260906/source/04/drafts/SGA4_IV_LINES_4476_4830_English.texfrag


% BEGIN INCLUDED SGA4_English_current_buildable_source_20260906/source/04/drafts/SGA4_IV_LINES_4831_5055_English.texfrag
\section{Localization. Opens of a topos}\etiquette[8]{IV.sec8}

\subsection{}\etiquette[8.1]{IV.subsec8.1}
Let $C$ be a category, and let $T(X)$ be a relation involving an
object $X$ of $C$. The relation $T(X)$ is said to be \emph{stable
under base change} (in $X$) if, for every morphism $X'\rightarrow X$
of $C$, the relation $T(X)$ implies $T(X')$. This amounts to saying
that the full subcategory $R$ of $C$ formed by the objects $X$ for
which $T(X)$ holds is a \emph{sieve}
(\hyperrefext[I][I.def4.1]{4.1}).

\begin{enonce*}[remark]{Remark 8.1.1}\etiquette[8.1.1]{IV.rem8.1.1}
Suppose that $C$ is a $𝒰$-category
(\hyperrefext[I][I.def1.1]{1.1}), so that $C$ may be identified with
a full subcategory of $\widehat{C}$
(\hyperrefext[I][I.prop1.4]{1.4}). It is sometimes convenient to
extend the definition of the relation $T$ on $\ob C$ to a relation
$\widehat{T}$ on an object $F$ of $\widehat{C}$, by defining
$\widehat{T}(F)$ to mean that $T(X)$ holds for every arrow
$X\rightarrow F$ in $\widehat{C}$ with $X\in\ob C$. The relation
$\widehat{T}(F)$ is clearly still stable under base change in $F$.
Again denoting by $R$ the subobject of the terminal object $e$ of
$\widehat{C}$ defined by the sieve $R$ of $C$
(\hyperrefext[I][I.rem4.2.1]{4.2.1}), the relation $\widehat{T}(F)$
may also be expressed as
$\Hom_C^{\widehat{\phantom{,}}}(F,R)\neq∅$, \ie it means that the
unique morphism $F\rightarrow e$ factors through the subobject $R$
of $e$.
\end{enonce*}

\subsection{}\etiquette[8.2]{IV.subsec8.2}
Now suppose that $C$ is a \emph{site}. A relation $T(X)$ in an
argument $X\in\ob C$ is said to be \emph{stable under descent} if,
for every covering\pageoriginale family
$(X_i:X_i\rightarrow X)_{i\in I}$ of an object $X$ of $C$, the
relation ``$T(X_i)$ for every $i\in I$'' implies $T(X)$. We say that
$T$ is a \emph{relation of local nature} if it is stable under base
change (\hyperref[IV.subsec8.1]{8.1}) and stable under descent. When
$C$ is a $𝒰$-category and $T$ is already assumed stable under base
change, \ie is defined by a sieve $R$ of $C$, which we identify with
a subpresheaf of the terminal presheaf on $C$, one sees immediately
that $T$ is of local nature if and only if $R$ is a \emph{sheaf}.
Thus \emph{relations of local nature in an argument $X\in\ob C$
correspond exactly to the subsheaves of the terminal sheaf of $C$,
\ie to the subobjects of the terminal object of the topos
$C^{\sim}$.} They therefore depend essentially only on the topos
$C^{\sim}$ defined by the site $C$ (like all important notions
associated with a site!). In terms of the subsheaf $R$ of the
terminal sheaf, the relation $T(X)$ is expressed as
$\Hom(\epsilon(X),R)\neq∅$. It extends canonically to the relation
$T^{\sim}(F):\Hom(F,R)\neq∅$ on the topos $C^{\sim}$.

\begin{enonce*}[remark]{Remark 8.2.1}\etiquette[8.2.1]{IV.rem8.2.1}
With the notation introduced in \hyperref[IV.rem8.1.1]{8.1.1}, for a
presheaf $F$ on $C$, since
$\Hom(F,R)\simeq\Hom(\underline{a}(F),R)$, the relation
$\widehat{T}(F)$ is equivalent to the relation
$\widehat{T}(\underline{a}(F))$.
\end{enonce*}

\begin{enonce*}{Definition 8.3}\etiquette[8.3]{IV.def8.3}
An open of a $𝒰$-topos $E$ is a subobject of the terminal object
$e_E$ of $E$; if $X$ is an object of $E$, an open of $X$ sometimes
means an open of the induced topos $E_{/X}$, \ie a subobject of $X$.

An open of a $𝒰$-site $C$ is an open of the associated $𝒰$-topos
$C^{\sim}$.
\end{enonce*}

Since the terminal objects of $E$ are canonically isomorphic, the
ambiguity introduced in \hyperref[IV.def8.3]{8.3} by the choice of
$e_E$ is harmless. One\pageoriginale may also, in a more intrinsic
but practically less convenient way, define the opens of $E$ as the
full subcategories $R$ of $E$ for which the relation
$F\in\ob R$, for an object $F$ of $E$, is of local nature. Likewise,
the opens of a $𝒰$-site $C$ correspond bijectively to the full
subcategories of $C$ for which the relation $F\in\ob R$, for an
object of $C$, is of local nature; this notion is therefore
essentially independent of the chosen universe $𝒰$. Finally, by what
was said in \hyperref[IV.subsec8.2]{8.2}, a relation of local nature
on a topos (or on a site) is essentially the same thing as an open of
that topos (or site).

\setcounter{subsection}{3}
\subsection{Examples of opens}\etiquette[8.4]{IV.subsec8.4}

\subsubsection{}\etiquette[8.4.1]{IV.subsubsec8.4.1}
Let $X$ be a topological space, and interpret $\Top(X)$
(\hyperref[IV.subsec2.1]{2.1}) as the category of étale spaces over
$X$. Since the monomorphisms in $\Top(X)$ are plainly the injective
maps, it follows that the opens of the topos $\Top(X)$ identify with
the opens of the space $X$. More generally, the opens of a topos
$\Top(X,G)$ (\hyperref[IV.subsec2.4]{2.4}) identify with the opens of
$X$ which are \emph{invariant} under the action of $G$.

\subsubsection{}\etiquette[8.4.2]{IV.subsubsec8.4.2}
The opens of a $𝒰$-topos $E$ form a $𝒰$-small set
(\hyperrefext[I][I.prop7.4]{7.4}), ordered by inclusion, admitting
arbitrary suprema and finite infima. Its least element is the initial
object (or ``empty object'') $\phi_E$, and its greatest element is
the terminal object $e_E$ of $E$. These two elements are identical
only if $E$ is an ``empty topos'' (\hyperref[IV.subsec2.2]{2.2}).

\subsubsection{}\etiquette[8.4.3]{IV.subsubsec8.4.3}
Let $G$ be a monoid object in $E$, giving a topos $B_G$
(\sisi{\hyperref[IV.subsec2.4]{2.4}}{\hyperref[IV.subsec2.6]{2.6}}),
the category of objects of $E$ on which $G$ acts on the left. The
terminal object $e_G$ of $B_G$ is the\pageoriginale terminal object
$e_E$ of $E$ with the trivial action of $G$. It follows that the
ordered set of opens of $B_G$ is canonically isomorphic to that of
the opens of $E$. In particular, if $E$ is the punctual topos, so
that $G$ is an ordinary monoid, the set of opens of $B_G$ consists
of exactly two elements, namely $\phi_{B_G}$ and $e_G$.

\subsubsection{}\etiquette[8.4.4]{IV.subsubsec8.4.4}
If $E$ is a topos of the form $\widehat{C}$, where $C$ is a
$𝒰$-category, then (\hyperref[IV.subsec8.1]{8.1}) the ordered set of
opens of $E$ is isomorphic to the ordered set of sieves of $C$.

\subsubsection{}\etiquette[8.4.5]{IV.subsubsec8.4.5}
Let $X$ be a sober nondiscrete topological space. Show that
$\TOP(X)$ (\hyperref[IV.subsec2.5]{2.5}) has opens which do not arise
from opens of $X$, \ie from opens of $\Top(X)$, by inverse image
under the canonical morphism \eqhyperref[IV.eq4.10.1]{4.10.1}
$\TOP(X)\rightarrow\Top(X)$.

\subsection{Examples of relations of local nature}\etiquette[8.5]{IV.subsec8.5}

For every object $X$ of the topos $E$, denote by
$F\rightarrow F_X$ the \emph{localization functor}
$j_X^{\ast}:E\rightarrow E_{/X}$, $Y\mapsto Y\times X$
(\hyperref[IV.subsec5.2]{5.2}).

\subsubsection{}\etiquette[8.5.1]{IV.subsubsec8.5.1}
Let $u:F\rightarrow G$ be a morphism of $E$. The following relation
in the argument $X\in\ob E$:

$u_X:F_X\rightarrow G_X$ is a monomorphism (\resp an epimorphism,
\resp an isomorphism)

\noindent is a relation of local nature. In particular, if $G$ is a
group object of $E$, the relation ``$G_X$ is the trivial group over
$X$'' is of local nature; the corresponding open of $E$ is called
the \emph{cosupport of the group $G$.}

\subsubsection{}\etiquette[8.5.2]{IV.subsubsec8.5.2}
Let\pageoriginale $u,v:F\rightrightarrows G$ be two morphisms of
$E$. The relation in the argument $X\in\ob E$
$$
u_X=v_X:F_X\longrightarrow G_X
$$
is a relation of local nature. In particular, if $G$ is a group
object of $E$, and $u$ is a \emph{section} of $G$
(\hyperref[IV.subsubsec4.3.6]{4.3.6}), the relation in $X$
$$
\mbox{the section $u_X$ of the group $G_X$ of $E_{/X}$ is zero}
$$
is of local nature; the corresponding open of $E$ is called the
\emph{cosupport of the section $u$ of the group $G$.}

\begin{enonce*}[remark]{Remark 8.5.3}\etiquette[8.5.3]{IV.rem8.5.3}
When $E=\Top(X)$, the cosupport of a sheaf of groups $G$, \resp of a
section of such a sheaf $G$, is simply the open of $X$ complementary
to the \emph{support of} $G$, \resp to the \emph{support of} $u$, in
the classical terminology [TF].
\end{enonce*}

\setcounter{subsubsection}{3}
\subsubsection{}\etiquette[8.5.4]{IV.subsubsec8.5.4}
Let $P(f)$ be a relation in the argument $f\in{\rm f\ell\,}C$, where
$C$ is a site. When fiber products are representable in $C$, the
relation $P(f)$ is said to be \emph{stable under base change (\resp
stable under descent, \resp of local nature) on the target} (or
\emph{on the base}, or ``\emph{downstairs}'') if, for every morphism
$f:X\rightarrow Y$ of $C$, the following relation in the argument
$Y'\in\ob C_{/Y}$:

``the morphism $f':X'=\fprod{X}{Y'}{Y}\rightarrow Y'$ obtained from
$f$ by base change along the structural morphism $Y'\rightarrow Y$
satisfies $P(f')$''

is stable under base change, \resp stable under descent, \resp is of
local nature. When the topology of $C$ is defined by means of a
pretopology, the relation $P(f)$ is said to be \emph{of local nature
on the source} (or \emph{``upstairs''}) if, for every morphism
$f:X\rightarrow Y$ of $C$ and every family
$(g_i:X_i\rightarrow X)_{i\in I}$\pageoriginale in $\Cov(X)$, the
relation $P(f)$ is equivalent to the relation ``$P(fg_i)$ for every
$i\in I$.'' When $\Cov(X)$ consists of all covering families of $C$,
this condition also means that the following relation in the object
$X'$ of the induced site $C_{/X}$
\begin{center}
$P(fg)$, where $g:X'\rightarrow X$ is the structural morphism
\end{center}
is of local nature. Note that if the relation $P(f)$ is of local
nature upstairs, it is a fortiori of local nature downstairs.

\subsubsection{}\etiquette[8.5.5]{IV.subsubsec8.5.5}
Take, for example, $C=(\Sch)$, the category of schemes in $𝒰$, with
the faithfully flat quasi-compact topology (\SGA 3 IV 6.3) or a
coarser topology. Each of the following properties of a morphism is
then of local nature downstairs:

surjective, radicial, universally open, universally closed, proper,
quasi-compact, quasi-compact and dominant, universal homeomorphism,
separated, quasi-separated, locally of finite type, locally of finite
presentation, of finite type, of finite presentation, an
isomorphism, a monomorphism, an open immersion, a closed immersion, a
quasi-compact immersion, affine, quasi-affine, finite, quasi-finite,
integral, flat, faithfully flat, neat, smooth, étale

\noindent (\EGA IV 2.6.4, 2.7.1 and \EGA IV 17.7.1). On the other
hand, endow $(\Sch)$ with the pretopology for which, for every scheme
$X\in𝒰$, $\Cov(X)$ consists of the families of morphisms
$(f_i:X_i\rightarrow X)_{i\in I}$ which are surjective and for which
the $f_i$ are flat and locally of finite presentation. Each of the
following properties is then \emph{of local nature upstairs}:

locally of finite type, locally of finite presentation, of finite
type, flat, neat, smooth, étale.

\noindent (\EGA IV 17.7.5, 17.7.7).

% END INCLUDED SGA4_English_current_buildable_source_20260906/source/04/drafts/SGA4_IV_LINES_4831_5055_English.texfrag


% BEGIN INCLUDED SGA4_English_current_buildable_source_20260906/source/04/drafts/SGA4_IV_LINES_5056_5253_English.texfrag
\begin{enonce*}[remark]{Exercise 8.6}\etiquette[8.6]{IV.exer8.6}
Let\pageoriginale $f:E\rightarrow F$ be a morphism of $𝒰$-topoi.
Consider the following relation in the argument $X\in\ob F$: the
induced morphism $E_{/f^{\ast}(X)}\rightarrow F_{/X}$ is an
equivalence of topoi. Prove that this relation is of local nature.
\end{enonce*}

\begin{enonce*}[remark]{Exercise 8.7}\etiquette[8.7]{IV.exer8.7}
(\emph{Partitions of a topos, sum of topoi}). Let $E$ be a
$𝒰$-topos. A family $(e_i)_{i\in I}$ of opens of $E$, \ie of
subobjects of the terminal object $e$ of $E$, is called a
\emph{partition of} $e$ (or of the topos $E$) if the canonical
morphism $\amalg_{i\in I}e_i\rightarrow e$ is an isomorphism, \ie
(\hyperrefext[II][II.cor4.6.2]{4.6.2}) if $e$ is the supremum of the
$e_i$ and $i\neq j$ implies $e_i\cap e_j=∅_E$.
\begin{enumerate}
\renewcommand{\theenumi}{\alph{enumi}}
\renewcommand{\labelenumi}{\theenumi)}
\item Show that a family $(e_i)_{i\in I}$ of objects of $E$ is a
partition of $E$ if and only if the functor
\begin{equation*}
E\longrightarrow \prod_i E_{/e_i}\tag{8.7.1}
\end{equation*}\etiquette[8.7.1]{IV.eq8.7.1}
defined by the localization functors $E\rightarrow E_{/e_i}$ is an
equivalence of categories.

\item Conversely, let $(E_i)_{i\in I}$ be a family of $𝒰$-topoi,
with $I\in𝒰$. Prove that $E=\Pi_{i\in I}E_i$ is a $𝒰$-topos (called
the \emph{sum topos} of the family $(E_i)_{i\in I}$). Define a
partition $(e_i)_{i\in I}$ of $E$ and equivalences of categories
$E_{/e_i}\approx E_i$ in such a way that the projection functors
$E\rightarrow E_i$ identify with the localization functors
$E\rightarrow E_{/e_i}$.

\item Let $E$ be the sum topos of the family $(E_i)_{i\in I}$. Show
that, for every $i\in I$, the projection functor $E\rightarrow E_i$
is of the form $u_i^{\ast}$, where
\begin{equation*}
u_i:E_i\longrightarrow E\tag{8.7.2}
\end{equation*}\etiquette[8.7.2]{IV.eq8.7.2}
is\pageoriginale a morphism of topoi. Prove that, for every
$𝒰$-topos $F$, the functor $v\mapsto(v\circ u_i)_{i\in I}$
\begin{equation*}
\Hhomtop(E,F)\longrightarrow \prod_i\Hhomtop(E_i,F)\tag{8.7.3}
\end{equation*}\etiquette[8.7.3]{IV.eq8.7.3}
is an equivalence of categories.

\item Let $(e_i)_{i\in I}$ be a partition of the topos $E$. For
every morphism of topoi $g:F\rightarrow E$, the family
$(f_i)_{i\in I}$, where $f_i=g^{\ast}(e_i)$, is a partition of $F$,
and the induced morphisms $g_i:F_{/f_i}\rightarrow E_{/e_i}$ make it
possible to reconstruct $g$ (up to unique isomorphism), the functor
$g^{\ast}$ identifying with the Cartesian product of the functors
$g_i^{\ast}$ (taking into account the equivalences of the type in
\hyperref[IV.exer8.7]{8.7} 1)). Deduce a complete description of the
category $\Hhomtop(F,E)$ in terms of categories of the form
$\Hhomtop(F',E_i)$, where $F'$ is a topos of the form $F_{/f}$ (with
$f$ a direct summand of the terminal object of $F$) and
$E_i=E_{/e_i}$.

\item In particular, if $F$ is connected and nonempty (\cf
\hyperref[IV.subsubsec4.3.5]{4.3.5}), \ie if, for every partition
$(f_i)_{i\in I}$ of $F$, there exists a unique $i\in I$ such that
$f_i\neq\phi_F$, prove that the canonical functor
\begin{equation*}
\coprod_{i\in I}\Hhomtop(F,E_i)\longrightarrow
\Hhomtop(F,E)\tag{8.7.4}
\end{equation*}\etiquette[8.7.4]{IV.eq8.7.4}
induced by the morphisms of topoi
\eqhyperref[IV.eq8.7.2]{8.7.2} is an equivalence of categories. More
particularly, the family of morphisms of topoi
\eqhyperref[IV.eq8.7.2]{8.7.2} induces an equivalence of categories
$$
\coprod_{i\in I}\Ppoint(E_i)\xrightarrow{\approx}\Ppoint(E){\quoi}.
$$

\item A\pageoriginale subset $I$ of the set of opens of $E$ is
called a \emph{reduced partition} of $E$ if the identity family
$(e_i)_{i\in I}$ indexed by $I$ is a partition and the $e_i$ are
different from $\phi_E$. Show that the refinement relation
(\hyperrefext[I][I.sec4.3.2]{4.3.2}) between reduced partitions makes
the set $P$ of reduced partitions into a $𝒰$-small ordered set in
which the supremum of any two elements exists. This gives a strict
projective system $(I_p)_{p\in P}$, regarded as a pro-set and denoted
by $\pi_{\circ}(E)$. It is called the \emph{pro-$\pi_{\circ}$ of the
topos} $E$.

\item Prove that $\pi_{\circ}(E)$ pro-represents the functor
$$
T\longmapsto f_{\ast}f^{\ast}(T)=\Gamma(E,
T_E):(𝒰\text{-}\Ens)\longrightarrow (𝒰\text{-}\Ens)
$$
associated with the canonical morphism $f:E\rightarrow P$ from $E$
to the punctual topos $P$ (\hyperref[IV.subsec4.3]{4.3}). In
particular, this functor is representable if and only if
$\pi_{\circ}(E)$ is essentially constant, \ie isomorphic (as a
pro-set) to an ``ordinary'' set, again denoted $\pi_{\circ}(E)$. The
pro-set $\pi_{\circ}(E)$ consists of one point if and only if $E$ is
``connected and nonempty.'' It is empty if and only if $E$ is the
``empty topos'' (\hyperref[IV.subsec2.2]{2.2}).

\item Suppose that $E$ is quasi-compact, \ie every covering of its
terminal object $e$ admits a finite subcovering. Show that
$\pi_{\circ}(E)$ is then a profinite set and hence may be identified
(using the familiar equivalence between pro-objects of the category
of finite sets and compact totally disconnected spaces) with a
compact totally disconnected space, also denoted $\pi_{\circ}(E)$.

\item Let $X$ be a topological space, and let $\pi_{\circ}(X)$ be the
set of connected components of $X$, regarded as a constant pro-set.
Define a canonical\pageoriginale morphism of pro-sets
\begin{equation*}
\pi_{\circ}(X)\longrightarrow
\pi_{\circ}(\Top(X)){\quoi},\tag{8.7.5}
\end{equation*}\etiquette[8.7.5]{IV.eq8.7.5}
or equivalently a canonical map
\begin{equation*}
\pi_{\circ}(X)\longrightarrow \varprojlim
\pi_{\circ}(\Top(X)){\quoi}.\tag{8.7.6}
\end{equation*}\etiquette[8.7.6]{IV.eq8.7.6}
Show that the first morphism is an isomorphism if and only if $X$ is
homeomorphic to a sum space of connected spaces (which is the case,
in particular, when $X$ is locally connected).

\item Suppose that $X$ is a quasi-compact scheme. Prove that
\eqhyperref[IV.eq8.7.6]{8.7.6} is bijective.

\item Establish the ``functorial character'' in $X$ of the morphisms
\eqhyperref[IV.eq8.7.5]{8.7.5} and
\eqhyperref[IV.eq8.7.6]{8.7.6}, first making precise the meaning of
this expression.


\item (Compare \hyperref[IV.exer7.6]{7.6} i)). Show that the
following two conditions on the $𝒰$-topos $E$ are equivalent: (i)
The canonical morphism from $E$ to the punctual topos is
``essential'' (Exercise \hyperref[IV.exer7.6]{7.6} a)), \ie the
functor $I\mapsto I_E$ from $(𝒰\text{-}\Ens)$ to $E$ commutes with
products indexed by sets in $𝒰$. (ii) For every object $X$ of $E$,
the induced topos $E_{/X}$ has a $\pi_{\circ}(E_{/X})$ which is an
ordinary set, \ie $X$ is isomorphic to the sum of a family of
connected objects of $E$. In this case $E$ is called a
\emph{locally connected topos} (compare
\sisi{IV~}{}\hyperref[IV.exer2.7.5]{2.7.5}).
\end{enumerate}
\end{enonce*}

\begin{enonce*}[remark]{Exercise 8.8}\etiquette[8.8]{IV.exer8.8}
(\emph{Dominant morphisms of topoi}).
\begin{enumerate}
\renewcommand{\theenumi}{\alph{enumi}}
\renewcommand{\labelenumi}{\theenumi)}
\item Let $f:E'\rightarrow E$ be a morphism of topoi. Show that the
following conditions are equivalent:

\begin{itemize}
\item[(i)] $f_{\ast}(\phi_{E'})=\phi_E$ (where $\phi_E$ and
$\phi_{E'}$ denote the initial objects of $E$ and $E'$, respectively).

\item[(ii)] For\pageoriginale every object $X$ of $E$ which is
``nonempty,'' \ie not isomorphic to $\phi_E$, the object
$f^{\ast}(X)$ is nonempty in $E'$.

\item[$\rm(ii')$] As in (ii), with $X$ a subobject of the terminal
object $e_E$ of $E$, \ie with $X$ an open of the topos $E$.
\end{itemize}

The morphism of topoi $f$ is then called \emph{dominant}.

\item Let $f:X'\rightarrow X$ be a continuous map of topological
spaces, giving a morphism of topoi
$\Top(f):\Top(X')\rightarrow\Top(X)$
(\hyperref[IV.subsec4.1]{4.1}). Show that $\Top(f)$ is dominant if
and only if $f$ is dominant, \ie if $f(X')$ is dense in $X$.

\item Show that if $f:E'\rightarrow E$ is a ``conservative''
morphism of topoi (\hyperref[IV.subsubsec6.4.0]{6.4.0}), \ie if
$f^{\ast}$ is faithful, then $f$ is dominant. Show that the converse
need not hold, even for a localization morphism
$E_{/U}\rightarrow E$, where $U$ is an open of the topos $E$.
(Show that in this case $f$ is conservative only if $U$ is the
terminal object of $E$ (and hence only if $f$ is an equivalence of
topoi), and use the example in b).)

\item Let $f:X'\rightarrow X$ be an arrow in a topos $E$. We call
$f$ a \emph{dominant morphism} in the topos $E$ if the corresponding
morphism of induced topoi $E_{/X'}\rightarrow E_{/X}$
\eqhyperref[IV.eq5.5.2]{5.5.2} is dominant. Show that this property
depends only on the image $f(X')$ of $X'$ in $X$, as an element of
the ordered set $\ouv(X)$ of subobjects of $X$. Show that the
morphism $E_{/X'}\rightarrow E_{/X}$ is conservative if and only if
$f$ is covering.
\end{enumerate}
\end{enonce*}

% END INCLUDED SGA4_English_current_buildable_source_20260906/source/04/drafts/SGA4_IV_LINES_5056_5253_English.texfrag


% BEGIN INCLUDED SGA4_English_current_buildable_source_20260906/source/04/drafts/SGA4_IV_LINES_5254_5490_English.texfrag
\section{Subtopoi and gluing of topoi}\etiquette[9]{IV.sec9}\pageoriginale

\subsection{Subtopoi and embeddings of topoi}\etiquette[9.1]{IV.subsec9.1}

\begin{enonce*}{Definition 9.1.1}\etiquette[9.1.1]{IV.def9.1.1}
\begin{enumerate}
\renewcommand{\theenumi}{\alph{enumi}}
\renewcommand{\labelenumi}{\theenumi)}
\item A morphism of topoi $f:F\rightarrow E$ is called an embedding
if the functor $f_{\ast}:F\rightarrow E$ is fully faithful (\ie if
the adjunction morphism $f^{\ast}f_{\ast}\rightarrow\id_F$ is an
isomorphism).

\item A subtopos of a topos $E$ is a strictly full subcategory $E'$
of $E$ which is a topos and for which the inclusion functor
$\alpha:E'\rightarrow E$ is of the form $i_{\ast}$, where
$i:E'\rightarrow E$ is a morphism of topoi (\ie
(\hyperref[IV.subsubsec3.1.1]{3.1.1}) $\alpha$ admits a left adjoint
$i^{\ast}$ which is left exact). The morphism $i$ (which is
determined up to unique isomorphism) is then called the inclusion
morphism of the subtopos $E'$ in $E$.
\end{enumerate}
\end{enonce*}

\begin{enonce*}[remark]{Remark 9.1.2}\etiquette[9.1.2]{IV.rem9.1.2}
\begin{enumerate}
\renewcommand{\theenumi}{\alph{enumi}}
\renewcommand{\labelenumi}{\theenumi)}
\item It is clear that the inclusion morphism of a subtopos is an
embedding, and that a subcategory $E'$ of a topos $E$ is a subtopos
if and only if it is the essential image of a direct-image functor
$f_{\ast}$ associated with a morphism of topoi $f:F\rightarrow E$
which is an embedding.

\item It follows immediately from the definitions that a morphism of
topoi $f:F\rightarrow E$ is an embedding if and only if it is
isomorphic to a composite morphism
$$
F\xrightarrow{g}E'\xrightarrow{i}E{\quoi},
$$
where $g$ is an equivalence of topoi and $i$ is the inclusion
morphism of a subtopos $E'$ of $E$. The latter is uniquely determined
by the preceding conditions as the essential image of $f_{\ast}$,
and the preceding factorization is likewise unique up to unique
isomorphism. In practice, of course, one should use $g$ to identify
$F$ with the subtopos $E'$ of $E$ (compare
\hyperref[IV.subsubsec3.4.1]{3.4.1}).

\item An equivalence of topoi (\hyperref[IV.subsec3.4]{3.4})
$u:E\xrightarrow{\approx}E_1$ defines in the evident way a bijection
between the set of subtopoi of $E$ and the set of subtopoi of $E_1$.
Indeed, since $u_{\ast}$ is an equivalence of categories, it
establishes\pageoriginale a bijection between the set of strictly
full subcategories of $E$ and the set of strictly full subcategories
of $E_1$ (to $E'\subset E$ corresponds the essential image $E'_1$ of
$u_{\ast}|E'$), and one sees immediately that $E'$ is a subtopos of
$E$ if and only if $E'_1$ is a subtopos of $E_1$. Thus, in studying
subtopoi, one may reduce to the case where $E$ is of the form
$\widetilde{C}$, with $C$ a small site. In this case,
\hyperrefext[II][II.theo5.5]{5.5} shows that a strictly full
subcategory $E'$ of $E$ is a subtopos if and only if its inclusion
functor $\alpha:E'\rightarrow E$ admits a left adjoint which is left
exact (which already implies that $E'$ is a topos), and that these
subtopoi are in one-to-one correspondence with the topologies $T'$
on $C$ finer than the given topology $T$ of $C$: to such a $T'$ one
associates the strictly full subcategory $(C,T')^{\sim}$ of
$C^{\sim}$ formed by the sheaves for the topology $T'$. This also
shows, for \emph{every} $𝒰$-topos $E$, that a strictly full
subcategory $E'$ of $E$ is a subtopos if and only if its inclusion
functor $\alpha:E'\rightarrow E$ admits a left adjoint which is left
exact; \emph{the ordered set $\mathfrak{E}(E)$ of subtopoi of $E$ is
$𝒰$-small, and every subset of $\mathfrak{E}(E)$ admits a supremum
and an infimum.}

\item If $E'$ and $E''$ are two subtopoi of a topos $E$, their
intersection is a subtopos of $E$. Indeed, since this intersection
is strictly full, one reduces to the case where $E$ is of the form
$C^{\sim}$. With the notation of \sisi{b)}{c)}, $E'$ and $E''$
then correspond to two topologies $T'$ and $T''$ on $C$ finer than
$T$, and it suffices to see that $E'\cap E''$ \emph{is the category
of sheaves for the topology $T'''$ which is the supremum of $T'$ and
$T''$}. We indicate the proof of this fact (which should have
appeared as a corollary after \hyperrefext[II][II.theo4.4]{4.4} or
\hyperrefext[II][II.theo5.5]{5.5}!). Let $E'''$ be the category of
sheaves for $T'''$, plainly contained in $E'$ and $E''$, and hence in
$E'\cap E''$. For every full subcategory $D$ of $E$, let $t(D)$ be
the finest topology among those for which the objects of $D$ are
sheaves\pageoriginale (\hyperrefext[II][II.prop2.2]{2.2}). The
inclusions $E'''\subset E'\cap E''\subset E',E''$ plainly imply
$t(E'),t(E'')\leq t(E'\cap E'')\leq t(E''')$. On the other hand,
(\sisi{\hyperrefext[II][II.theo4.4]{4.4} I}{\hyperrefext[II][II.cor4.4.4]{4.4.4}})
$t(E')=T'$, $t(E'')=T''$, and $t(E''')=T'''$, so the preceding
inequalities and the definition of $T'''$ as the supremum imply
$T'''=t(E'\cap E'')$, and hence $E'\cap E''\subset E'''$, \cqfd\
More generally, this argument shows that if $T'$ is the supremum of
an arbitrary family of topologies $(T_i)$, the category $E'$ of
sheaves for $T'$ is the intersection of the categories $E_i$ of
sheaves for the $T_i$. In particular, one obtains:
\end{enumerate}
\end{enonce*}

\begin{enonce*}[remark]{Proposition 9.1.3}\etiquette[9.1.3]{IV.prop9.1.3}
Let $E$ be a topos. For every family of subtopoi $(E_i)_{i\in I}$ of
$E$, the intersection $\bigcap_{i\in I}E_i$ is a subtopos of $E$.
\end{enonce*}

\begin{enonce*}{Proposition 9.1.4}\etiquette[9.1.4]{IV.prop9.1.4}
Let $i:E'\rightarrow E$ be an embedding of topoi, and let $F$ be a
topos. The functor
\begin{equation*}
f\longmapsto i\circ f:\Hhomtop(F,E')\longrightarrow
\Hhomtop(F,E)\tag{9.1.4.1}
\end{equation*}\etiquette[9.1.4.1]{IV.eq9.1.4.1}
is fully faithful, and its essential image consists of the morphisms
of topoi $g:F\rightarrow E$ for which the essential image of
$g_{\ast}$ is contained in that of $i_{\ast}$.
\end{enonce*}

Indeed, consider the evident commutative diagram
\[
\xymatrix@C=1.75cm@R=1cm@M=6pt{
\Hhomtop (F,E') \ar @{^{(}->}[d] \ar[r]& \Hhomtop (F,E) \ar @{^{(}->}[d]\\
\Hhom (F,E') \ar[r] &\Hhom (F,E) , }
\]
where the vertical arrows are the functors $h\mapsto h_{\ast}$,
which are fully faithful by the definition of $\Hhomtop$
(\sisi{IV}{}\hyperref[IV.subsec3.2]{3.2}). The second horizontal arrow
$G\mapsto i_{\ast}G$ is also fully faithful, since $i_{\ast}$ is, and
therefore so is the first horizontal arrow. It remains to prove the
characterization of the essential image of
\eqhyperref[IV.eq9.1.4.1]{9.1.4.1}. The necessity of the condition is
clear from the formula $(if)_{\ast}=i_{\ast}f_{\ast}$. Conversely,
suppose that the essential image of $g_{\ast}$ is contained in that
of $i_{\ast}$, \ie that one may write
$g_{\ast}=i_{\ast}f_{\ast}$, with $f_{\ast}:F\rightarrow E'$ a
functor. We must show that $g$ belongs to the essential image of
\eqhyperref[IV.eq9.1.4.1]{9.1.4.1}. This amounts to saying that
$f_{\ast}$ admits a left adjoint which is left exact.
Now\pageoriginale it is clear that $f_{\ast}$ admits
$g^{\ast}i_{\ast}$ as a left adjoint, and this functor is left exact,
being the composite of the left exact functors $g^{\ast}$ and
$i_{\ast}$, \cqfd

For a converse to \hyperref[IV.prop9.1.4]{9.1.4}, \cf Exercise
\hyperref[IV.exer9.1.6]{9.1.6}.

Taking $F$ to be the punctual topos
(\sisi{IV~}{}\hyperref[IV.subsec2.2]{2.2}), one obtains in particular
from \hyperref[IV.prop9.1.4]{9.1.4}:

\begin{enonce*}{Corollary 9.1.5}\etiquette[9.1.5]{IV.cor9.1.5}
If $f:E'\rightarrow E$ is an embedding morphism of topoi, the
corresponding functor
\begin{equation*}
\Point(f):\Ppoint(E')\longrightarrow \Ppoint(E)\tag{9.1.5.1}
\end{equation*}\etiquette[9.1.5.1]{IV.eq9.1.5.1}
is fully faithful.
\end{enonce*}

\begin{enonce*}[remark]{Exercise 9.1.6}\etiquette[9.1.6]{IV.exer9.1.6}
(\emph{Embeddings of topoi and 2-monomorphisms}).

Let $i:E'\rightarrow E$ be a morphism of topoi.

\begin{enumerate}
\renewcommand{\theenumi}{\alph{enumi}}
\renewcommand{\labelenumi}{\theenumi)}
\item If $F$ and $G$ are two topoi \emph{over} $E'$, define (with
the notation of \hyperref[IV.exer5.14]{5.14} a)) a canonical functor
\begin{equation*}
f\longmapsto f\ast i:\Hhomtop_{E'}(F,G)\longrightarrow
\Hhomtop_E(F,G){\quoi}.\tag{9.1.6.1}
\end{equation*}\etiquette[9.1.6.1]{IV.eq9.1.6.1}

\item Suppose that $i$, as a 1-arrow of the 2-category of
$𝒱\text{-}𝒰$-topoi (\hyperref[IV.subsubsec3.3.2]{3.3.2}), is a
$2$-monomorphism, \ie that for every $𝒰$-topos $F$ which is an
element of $𝒱$, the functor \hyperref[IV.eq9.1.4.1]{(9.1.4.1)} is fully % add a ref
faithful. Show that, for every pair of $𝒰$-topoi $F,G$ (not
necessarily in $𝒱$), the functor \hyperref[IV.eq9.1.6.1]{(9.1.6.1)} is %% add a ref
fully faithful.

\item Show that $i$ is an embedding if and only if $i$ is a
$2$-monomorphism. (For sufficiency, apply b) to induced topoi
$E'_{/X'}$ and $E'_{/Y'}$, and use
\hyperref[IV.exer5.14]{5.14} c)).

\item Let $g:F\rightarrow E$ be a morphism of topoi, and suppose that
the $2$-fiber product (\hyperref[IV.prop5.11]{5.11})
\sisi{$F'=\fprod{F^2}{E'}{E}$}{$F'=F\times_E^2E'$}
exists (a condition which is always satisfied, as P. DELIGNE has
shown\nde{See also [16], Proposition 4.47.}). Let
$j:F'\rightarrow F$ be the projection morphism. Prove that if $i$ is
an embedding, then so is $j$. (Use c), and note that the stability of
the notion of a 2-monomorphism under 2-base change is formal.) When
$E'$ is a subtopos of $E$ and $i:E'\rightarrow E$ is the canonical
inclusion morphism, the essential image of $F'$ under $j_{\ast}$
(which is a subtopos of $F$ independent\pageoriginale of the
indeterminacy in the construction of a 2-fiber product) is called the
\emph{subtopos of $F$ which is the inverse image under
$g:F\rightarrow E$ of the subtopos $E'$ of $E$.}

\item With the notation at the end of d), choose as 2-fiber product
the subtopos of $F$ which is the inverse image of the subtopos $E'$ of
$E$. Let $X$ be an object of $F$. Show that, for $X$ to belong to
$F'$, it is necessary that it satisfy the following condition: for
every object $U$ of $F$, writing
$g_U:F_{/U}\rightarrow E$ for the morphism of topoi induced by $g$,
one has $g_{U\ast}(X|U)\in\Ob E'$, where
$X|U=(X\times U\xrightarrow{\pr_2}U)$ is the ``restriction of $X$ to
$U$.'' (NB: Note that if $i_U:F_{/U}\rightarrow F$ is the canonical
inclusion, then $g_U=g\circ i_U$, whence
$g_{U\ast}(X|U)=g_{\ast}(i_{U\ast}(X|U))=g_{\ast}(\Hhom(U,X))$,
where $\Hhom(U,X)$ is the object defined in \hyperref[IV.prop10.1]{10.1} below.)

\noindent \emph{Problem.} Conversely, is the preceding condition on
$X$ sufficient for $X$ to belong to $F'$? This amounts to asking
whether the strictly full subcategory of $F$ formed by these objects
$X$ is a \emph{subtopos} of $F$.

\item Let $X$ be an object of $E$, let $X'=i^{\ast}(X)$, and let
$i_{/X}:E'_{/X'}\rightarrow E_{/X}$ be the morphism of topoi induced
by $i$ (\eqhyperref[IV.eq5.10.1]{5.10.1}). Show that if $i$ is an
embedding, then so is $i_{/X}$. (Use d) and
\hyperref[IV.prop5.11]{5.11}.) Conversely, if
$(X_{\alpha})_{\alpha}$ is a family of objects covering the terminal
object, and if $i_{/X_{\alpha}}$ is an embedding for every $\alpha$,
then so is $i$.
\end{enumerate}
\end{enonce*}

% END INCLUDED SGA4_English_current_buildable_source_20260906/source/04/drafts/SGA4_IV_LINES_5254_5490_English.texfrag


% BEGIN INCLUDED SGA4_English_current_buildable_source_20260906/source/04/drafts/SGA4_IV_LINES_5491_5778_English.texfrag
\begin{enonce*}[remark]{Exercise 9.1.7}\etiquette[9.1.7]{IV.exer9.1.7}
(\emph{Subtopoi and multiplicative systems of arrows}).

\begin{enumerate}
\renewcommand{\theenumi}{\alph{enumi}}
\renewcommand{\labelenumi}{\theenumi)}
\item Let $i:E'\rightarrow E$ be a morphism of topoi, let $S$ be the set of
arrows $u$ of $E$ such that $i^{\ast}(u)$ is an isomorphism, and
consider the canonical functor induced by $i^{\ast}$ (\cf [2,
Chap. I] or \sisi{VI}{the proof of \hyperrefext[VI][VI.prop6.2]{6.2}}):
\begin{equation*}
\rho:E[S^{-1}]\longrightarrow E'{\quoi}.\tag{9.1.7.1}
\end{equation*}\etiquette[9.1.7.1]{IV.eq9.1.7.1}

Show that $i$ is an embedding if and only if the preceding functor
$\rho$ is an equivalence of categories, and that in this case $S$
admits both a calculus of left fractions and a calculus of right
fractions. (Use [2, Chap. I, 1.3].) In this case, the essential image
of $i_{\ast}$ can be recovered from $S$ as consisting of those
$X\in \Ob \sisi{\not{E}}{(E)}$ such that for every\pageoriginale
$u:Y\rightarrow Y'$ in $S$, the map $\Hom(u,X):\Hom(Y',X)\rightarrow
\Hom(Y,X)$ is bijective.

\item We henceforth assume that i is an embedding. Show that, for
every topos $F$, the functor
$$
f^{\ast}\longmapsto f^{\ast}i^{\ast}:\Hhom(E',F)\longrightarrow
\Hhom(E,F)
$$
induces an equivalence between the category of the $f^{\ast}$ arising
from morphisms of topoi $f:F\rightarrow E'$ (\ie which are left exact
and admit a right adjoint) and the category of functors
$g^{\ast}:E\rightarrow F$ arising from morphisms of topoi
$g:F\rightarrow E$ and such that $g^{\ast}$ takes arrows of $S$ to
isomorphisms.

\item Let $S$ be a subset of $Fl$ $E$, where $E$ is a
$𝒰$-topos. Show that $S$ is associated with a subtopos $E'$
of $E$ by the procedure of a) if and only if $S$ satisfies the
following conditions:
\begin{itemize}
\item[ST 1)] $S$ contains the invertible arrows and is stable under
composition and base change.

\item[ST 2)] If a composite vu belongs to $S$, then $u\in S$ if
and only if $v\in S$.

\item[ST 3)] If $u:X\rightarrow Y$ is such that there exists a
covering family $Y_i\rightarrow Y$, $i\in I$, for which
$u_i:X_i=\fprod{X}{Y_i}{Y}\to Y_{i}$ belongs to $S$ for every $i\in
I$, then $u\in S$.
\end{itemize}

Show that this gives a one-to-one correspondence between the set of
subtopoi $E'$ of $E$ and the set of subsets $S$ of $Fl$ $E$
satisfying conditions ST 1), ST 2), and ST 3). (Given $S$, associate
with it a $𝒰$-topology $T'$, finer than the canonical topology
$T$ of $E$, whose covering families $(X_i\rightarrow X)_{i\in I}$
are those for which the inclusion $X'\hookrightarrow X$ of the image
of the $X_i$ is $\in S$. Show that the category of sheaves for
$T'$ is equivalent to a subtopos $E'$ of $E$, and use
\hyperrefext[II][II.theo4.4]{4.4} to prove that it recovers $S$.)

\item Show that $S$ is determined by the subset $S_0$ of $S$
consisting of the \emph{monomorphisms}. Show that the map
$S\mapsto S_0$ establishes a bijection between the set of subsets
$S$ of $Fl$ $E$ satisfying conditions ST 1) through ST 3) of c) and
the set of subsets of $Fl$ $E$ consisting of
monomorphisms\pageoriginale and satisfying the same conditions ST
1) through ST 3).

\item Let $(E'_i)_{i\in I}$ be a family of subtopoi of $E$, and let
$(S_i)_{i\in I}$ be the corresponding family of subsets of $Fl$ $E$.
Show that $\bigcap_{i\in I} S_i$ satisfies conditions ST 1) through
ST 3) of c), and hence corresponds to a subtopos $E'$ of $E$, and
that the latter is the \emph{subtopos generated by the} $E'_i$, \ie
the \emph{supremum} (\cf \hyperref[IV.rem9.1.2]{9.1.2} c)) of the
subtopoi $E'_i$.

\item Let $A$ be a subset of $Fl$ $E$. Show that there is a smallest
subset $S$ of $Fl$ $E$ among those containing $A$ and satisfying
conditions ST 1) through ST 3), and that $S$ corresponds to the
largest subtopos $E'$ of $E$ such that the corresponding
multiplicative system of $Fl$ $E$ contains $A$. Show that $S$ is
equal to the union of the subsets $A_n$ $(n\geq 0)$ of $Fl$ $E$,
constructed inductively as follows: $A_0=A$, and $A_{n+1}$ consists
of the arrows $u:X\rightarrow Y$ which are of one of the following
four types (i) through (iv):
\begin{enumerate}
\renewcommand{\theenumii}{\roman{enumii}}
\renewcommand{\labelenumii}{(\theenumii)}
\item $u$ is invertible, or is the composite of two arrows of $A_n$,
or is obtained from an arrow of $A_n$ by base change.

\item $u$ is one of the factors of a composite arrow whose other
factor and whose composite both belong to $A_n$.

\item $u$ is the coproduct of a small family of arrows $\in
A_n$.

\item There exists an epimorphism $Y'\rightarrow Y$ such that
$u'=X'=\fprod{X}{Y'}{Y}\rightarrow Y'$, obtained from $u$ by base
change, belongs to $A_n$.
\end{enumerate}

\item With the notation of e), let $A$ be the union of the $S_i$,
let $S$ be the subset of $Fl$ $E$ obtained from $A$ by the procedure
of f), and let $E'$ be the corresponding subtopos of $E$. Show that
$E'$ is the intersection (= infimum) of the subtopoi $E'_i$.
\end{enumerate}
\end{enonce*}



\begin{enonce*}[remark]{Exercise 9.1.7.2}\etiquette[9.1.7.2]{IV.exer9.1.7.2}
(\emph{Canonical factorization of a morphism of topoi}).
\begin{enumerate}
\renewcommand{\theenumi}{\alph{enumi}}
\renewcommand{\labelenumi}{\theenumi)}
\item Let $f:F\longrightarrow E$
be a morphism of topoi. Show that, up to isomorphism, one can find a
factorization of $f$ as
\begin{equation*}
F\xrightarrow{f'}E'\xrightarrow{i} E,\tag{9.1.7.3}
\end{equation*}\etiquette[9.1.7.3]{IV.eq9.1.7.3}
where $f'$ is \emph{conservative} (\ie (\hyperref[IV.subsubsec6.4.0]{6.4.0})
${f'}^{\ast}$ is conservative, or equivalently faithful)\pageoriginale
and $i$ is an embedding, and that this factorization is unique up to
equivalence (\cf \hyperref[IV.prop9.1.4]{9.1.4}). To do this, introduce
the set $S_f\subset Fl$ $E$ of arrows $u$ of $E$ such that
$f^{\ast}(u)$ is an isomorphism, prove that $S_f$ satisfies the
conditions of \hyperref[IV.exer9.1.7]{9.1.7} c), take for $E'$ the corresponding subtopos of
$E$, and define ${f'}^{\ast}$ by \hyperref[IV.eq9.1.7.1]{(9.1.7.1)}, which corresponds to a
morphism of topoi $f'$ by \hyperref[IV.exer9.1.7]{9.1.7} b).

\item The subtopos of $E$ defined by the embedding $i$ is called the
subtopos of $E$ which is the \emph{image of the morphism of topoi}
$f$. Show that it is the smallest of the subtopoi $E_1$ of $E$ such
that $f$ factors, up to isomorphism, through $E_1$, \ie
(\hyperref[IV.prop9.1.4]{9.1.4}) the smallest of the subtopoi of $E$
containing the image of $f_{\ast}$ (or equivalently the essential
image of $f_{\ast}$). For $f$ to be an embedding, it is necessary
and sufficient that $f$ induce an equivalence $f'$ of $F$ with its
image; for $f$ to be conservative, it is necessary and sufficient
that its image be all of $E$.

\item Suppose that $f=\Top(f_0)$ is associated with a continuous map
of sober topological spaces $f_0:Y\rightarrow X$
(\hyperref[IV.subsec4.1]{4.1}). Consider the canonical factorization of
$f_0$ as
$$
Y\xrightarrow{{f'}_0}X'=f_0(Y)\xrightarrow{i_0}X{\quoi},
$$
where $i_0$ is the inclusion. Show that the canonical factorization
of $f$ then identifies with the corresponding factorization
\[
\xymatrix@C=2cm@M=.2cm{ \Top(Y) \ar[r]^{\Top(f'_0)}& \Top(X')
\ar@{^{(}->}[r]^{\Top(i_0)}& \Top(X). }
\]
Let $Z$ be the smallest sober subspace of $X$ containing $X'=f_0(Y)$,
\ie the set of points $x$ of $X$ such that $\{\bar{x}\}\cap X'$
is dense in $\{\bar{x}\}$. (Note that $Z=X'$ if the points of $X$
are closed, or if $X'$ is a \emph{constructible} subset of the
locally Noetherian space $X$.) Then $f$ is conservative, \ie the
subtopos of $E=\Top(X)$ which is the image of $F=\Top(Y)$ under
$f=\Top(f_0)$ is equal to $E$ itself, if and only if $Z=X$, \ie if
and only if $f_0(Y)$ is \emph{very dense} (EGA IV 10.13) in $X$. This
makes precise the extent to which the notion of a conservative
morphism of topoi may legitimately be regarded as the natural
generalization of the notion of a \emph{surjective} continuous map
of topological spaces.

\item Let\pageoriginale
\[
\xymatrix{
Y \ar[d]& Y'\ar[l] \ar[d]\\
X & X'\ar[l] }
\]
be a Cartesian diagram of topological spaces. If $Y$ and $X'$ are
discrete spaces, and hence $Y'$ is discrete, then the corresponding
diagram of topoi is a 2-fiber product diagram
(\hyperref[IV.prop5.11]{5.11}). (NB. For a counterexample when $Y$
and $X'$ are no longer assumed discrete, with $Y$ and $X'$ subspaces
of the Hausdorff space $X$, \cf \hyperref[IV.exer9.1.10]{9.1.10} c)).
Deduce an example of a 2-fiber product diagram of topoi
\[
\xymatrix{
F \ar[d]_{f}& F'\ar[l] \ar[d]^{f'}\\
E & E'\ar[l] }
\]
such that $f$ is conservative but $f'$ is not conservative, more
precisely with $E'=$ the punctual topos and $F'=$ the empty topos
(\hyperref[IV.subsec2.2]{2.2}). (Take $Y$ such that the image of $Y$
in $X$ is very dense and $\neq X$, and take $X'$ to consist of a
point not belonging to the image.)

\item Let $(E_i)_{i\in I}$ be a family of subtopoi of a topos $E$.
Prove that $E$ is the supremum of the $E_i$ if and only if the family
of embedding morphisms $E_i\rightarrow E$ is conservative
(\hyperref[IV.subsubsec6.4.0]{6.4.0}).
\end{enumerate}
\end{enonce*}

\begin{enonce*}[remark]{Exercise 9.1.8}\etiquette[9.1.8]{IV.exer9.1.8}
(\emph{Subtopoi and points of $E$. The case of topological spaces}).

Let $E$ be a topos.
\begin{enumerate}
\renewcommand{\theenumi}{\alph{enumi}}
\renewcommand{\labelenumi}{\theenumi)}
\item Let $Z$ be a full subcategory of $\Ppoint(E)$ (or, what amounts
to the same thing, a subset of $\Ob \Ppoint(E)$). Show that the set
$S(Z)$ of arrows of $E$ which are taken to bijections by the fiber
functors corresponding to the $z\in \Ob Z$ satisfies the conditions
of \hyperref[IV.exer9.1.7]{9.1.7} c), and hence defines a subtopos $E_Z$ of $E$. Show that
this subtopos has enough points, and that $Z$ is equivalent to a full
subcategory of $\Ppoint(E_Z)$. Show that every subtopos $E'$ of $E$
which has enough points can be obtained by the preceding procedure.
(Take for $Z$ the essential image of the functor \hyperref[IV.eq9.1.5.1]{(9.1.5.1)}.) Show
that one obtains\pageoriginale a bijection between the set of
subtopoi $E'$ of $E$ having enough points and a subset of the set of
strictly full subcategories $Z$ of $\Ppoint(E)$ which are stable
under direct factors (\sisi{\hyperrefext[I][I.sec10]{10}}{\hyperrefext[I][I.sec10.6]{10.6}})
and under small filtered projective limits (\cf
\hyperref[IV.subsec7.3]{7.3}). (Problem: Which subcategories of
$\Ppoint(E)$ arise in this way? This is essentially, in another
form, the problem already raised in \hyperref[IV.prop6.5.3]{6.5.3}, in view of the dictionary
in \hyperref[IV.rem9.1.2]{9.1.2} c). Note that, besides the necessary
conditions just indicated, there are also subtler conditions on the
topological nature of $\Point(E)$, such as sobriety; \cf e)).

\item Let $(Z_i)_{i\in I}$ be a family of full subcategories of
$\Ppoint(E)$, and let $Z$ be the full subcategory which is the union
of the $Z_i$. Show that the corresponding subtopos $E_Z$ defined in
a) is the subtopos of $E$ generated by the $E_{Z_i}$. (Use \hyperref[IV.exer9.1.7]{9.1.7}
e).) Conclude that the subtopos of $E$ generated by a family of
subtopoi having enough points also has enough points.

\item Let $f:E'\rightarrow E$ be an embedding of topoi. Show that
the map $U\mapsto f^{\ast}(U)$
\begin{equation*}
\Ouv(E)\longrightarrow \Ouv(E')\tag{9.1.8.1}
\end{equation*}\etiquette[9.1.8.1]{IV.eq9.1.8.1}
is surjective, and conclude that the map
\begin{equation*}
\Point(E')\longrightarrow \Point(E)\tag{9.1.8.2}
\end{equation*}\etiquette[9.1.8.2]{IV.eq9.1.8.2}
induced by $f$ on the isomorphism classes of points is a
homeomorphism from $\Point(E')$ onto a subspace of $\Point(E)$, for
the topologies defined in \hyperref[IV.exer7.8]{7.8} a). (For
injectivity, use \hyperref[IV.cor9.1.5]{9.1.5}.)

\item Let $f:X'\rightarrow X$ be a continuous map, giving a morphism
of topoi (\hyperref[IV.subsec4.1]{4.1})
$$
\Top(f):\Top(X')\longrightarrow \Top(X){\quoi}.
$$
Prove that the latter is an embedding if and only if the topology of
$X'$ is the inverse image under $f$ of that of $X$ (\ie, when $X'$
is a Kolmogorov space, for example a sober space, if and only if $f$
induces a homeomorphism of $X'$ onto a subspace of $X$). (For the
necessity,\pageoriginale use c), and for the sufficiency use
\hyperref[IV.exer9.1.7.2]{9.1.7.2} c)).

\item Let $X$ be a sober topological space, and let $E=\Top(X)$. Using
a) and d), establish an isomorphism of ordered sets $X'\mapsto
E_{X'}$ between the set of sober subspaces $X'$ of $X$ (\ie, when
$X$ is Hausdorff, the set of \emph{all} subsets of $X$) and
the set of subtopoi of $E$ having enough points. Show that $E_{X'}$
is equivalent to $\Top(X')$. If $(X'_i)_{i\in I}$ is a family of
sober subspaces of $X$, show that there is a smallest sober subspace
$X'$ of $X$ containing the $X'_i$, equal to the union of the $X'_i$
if I is finite or $X$ is Hausdorff, and that $E_{X'}$ is the
subtopos of $E$ generated by the $E_{X'_i}$. (Use b)).

\item Give an example of a topos $E$ of the form $\Top(X)$ and a
subtopos $E'$ which does not have enough points, and which is even
nonempty (\hyperref[IV.subsec2.2]{2.2}) but has no points. (Take $E$
of the form $\widehat{U}$ and $E'=U^{\sim}$, where $U$ is the ordered
set defined in \hyperref[IV.subsec7.4]{7.4}. Or better, take the
example in \sisi{9.1. 10. b)}{\hyperref[IV.exer9.1.10]{9.1.10} b)},
which shows that one may take for $X$ any nonempty compact space
without isolated points.)
\end{enumerate}
\end{enonce*}

% END INCLUDED SGA4_English_current_buildable_source_20260906/source/04/drafts/SGA4_IV_LINES_5491_5778_English.texfrag


% BEGIN INCLUDED SGA4_English_current_buildable_source_20260906/source/04/drafts/SGA4_IV_LINES_5779_6107_English.texfrag
\begin{enonce*}[remark]{Exercise 9.1.9}\etiquette[9.1.9]{IV.exer9.1.9}
(\emph{The case of topoi defined by an ordered set}).

\begin{enumerate}
\renewcommand{\theenumi}{\alph{enumi}}
\renewcommand{\labelenumi}{\theenumi)}
\item We use the yoga of \hyperref[IV.exer7.8]{7.8} b) and c), which gives a dictionary
between, on the one hand, topoi generated by their opens and morphisms
of such topoi, and, on the other hand, ordered sets admitting
arbitrary $\Sup$, finite infima, and in which the infimum of two elements %% standardize notation sup/Sup/\Sup etc.
distributes over arbitrary $\Sup$; morphisms between such ordered
sets are increasing maps commuting with finite infima and arbitrary
$\Sup$.

\item Let $E$ be a topos and $f:E'\rightarrow E$ an embedding of
topoi. Show that, for every generating family $(X_i)_{i\in I}$ in
$E$, the family $(f^{\ast}(X_i))_{i \in I}$ is generating. Deduce
that if $E$ is generated by its opens, then the same is true of
$E'$. (Use \hyperref[IV.exer9.1.8]{9.1.8} c).)

\item Let $f:E'\rightarrow E$ be a morphism of topoi generated by
their opens, associated with a morphism $g:\Oo\rightarrow \Oo'$ of
the corresponding ordered sets. Show that $f$ is an embedding if
and only if $g$ is surjective, \emph{and} $g$\pageoriginale is
surjective on the graphs of the order relations, \ie the order on
$\Oo'$ is the quotient of that on $\Oo$.

\item Let $E$ be a topos of the form $\Oo^{\sim}$ as in a). Deduce
from b) and c) a bijective correspondence between the set of
subtopoi of $E$ and the set of equivalence relations $R$ on $\Oo$
having the following properties, analogous to conditions ST 1) to
ST 3) of \hyperref[IV.exer9.1.7]{9.1.7} c):
\begin{itemize}
\item[Q0 1)] The relation $R$ is compatible with finite infima; in
other words, if $U$ and $U'$ are equivalent under $R$, then so are
$U\cap V$ and $U'\cap V$ for every $V\in \Oo$.

\item[Q0 2)] The relation $R$ is compatible with arbitrary $\Sup$;
in other words, if $U_i$ is equivalent to $U'_i$ for every $i\in I$,
then $\Sup_i U_i$ is equivalent to $\Sup_i U'_i$.
\end{itemize}

Show that if subtopoi $E'_i$ of $E$ correspond to the equivalence
relations $R_i$, then the subtopos generated by the $E'_i$
corresponds to the intersection of the relations~$R_i$.
\end{enumerate}
\end{enonce*}

\begin{enonce*}[remark]{Exercise 9.1.10}\etiquette[9.1.10]{IV.exer9.1.10}
(\emph{Intersections of subspaces; new topoi without points; nonexistence
of a complement to a subtopos}).

Let $X$ be a topological space, $E=\Top(X)=\Ouv(X)^{\sim}$
(\hyperref[IV.subsec2.1]{2.1}).

\begin{enumerate}
\renewcommand{\theenumi}{\alph{enumi}}
\renewcommand{\labelenumi}{\theenumi)}
\item Let $R$ be the following relation on $\Ouv(X):(U,U')\in
R\Longleftrightarrow U\cap U'$ is dense in both $U$ and $U'$.
Show that $R$ is an equivalence relation satisfying conditions Q0
1) and Q0 2) of \hyperref[IV.exer9.1.9]{9.1.9} d), and therefore defines a
subtopos $E'$ of $E$. Show that, for a point $x\in X$, the
corresponding point belongs to the essential image of $\Ppoint(E')$
if and only if $x$ is \emph{thick}, \ie belongs to the closure of the
interior of $\bar{x}$. Show that $E'$ is the empty topos
(\hyperref[IV.subsec2.2]{2.2}) if and only if $X$ is empty.

\item If $X$ is sober, show that the category $\Ppoint(E')$ is
equivalent to the category defined by the ordered subset of $X$
(for the specialization relation) consisting of the thick points
of $X$. Show that if the points of $X$ are closed (\resp if $X$ is
Noetherian), then the thick points of $X$ are the isolated points of
$X$, \ie those for which $\{x\}$ is open (\resp are the maximal
points of $X$). Conclude that if $X$ is\pageoriginale a nonempty
Hausdorff topological space without isolated points, then the
subtopos $E'$ of $E=\Top(X)$ is a non-``empty'' topos having no
points.

\item Let $X'$ be a subset of $X$, and let $E_{X'}$ be the subtopos
of $E$ associated with $X'$ by the construction of
\hyperref[IV.exer9.1.8]{9.1.8} a). Show that if $X'$ is dense, then $E_{X'}$
contains $E'$. Show that if $X$ is sober and $X'$ and $X''$ are two
sober subsets of $X$, then their intersection $X'\cap X''$ is sober
and $E_{X'\cap X''}\subset E_{X'}\cap E_{X''}$, and that the
inclusion can be strict
\footnote[1]{See, however, \hyperref[IV.exer9.1.11]{9.1.11} b) below for the
case where $X$ is locally Noetherian.}. (Take a nonempty Hausdorff
space $X$ having two disjoint dense subsets $X'$, $X''$.) Show that
the preceding inclusion is an equality if and only if the subtopos
$E_{X'}\cap E_{X''}$ has enough points. Show that this is the case
if $X'$ or $X''$ is a \emph{locally closed} subset of~$X$.

\item Let $X'$ be a subset of $X$, and let $X''$ be the complement
of $X'$. A subtopos $F$ of $E$ contains the $E_{\{x\}}$ for $x\in
X''$ if and only if it contains $E_{X''}$ (\cf
\hyperref[IV.exer9.1.8]{9.1.8} b)). Conclude that if $X''$ is a union of
locally closed subsets of $X$ (for example, if the points of $X''$
are closed), then the set of subtopoi $F$ of $E$ whose intersection
with $E_{X'}$ is the empty subtopos $\Top(∅)$ has a greatest element
(which then deserves the name of \emph{weak complement of}
$E_{X'}$ in $E$; \cf \hyperref[IV.exer9.1.13]{9.1.13}) \emph{only if}
$E_{X'}\cap E_{X''}$ is the empty subtopos (a condition that is not
satisfied if $X\neq ∅$ and if $X'$ and $X''$ are both dense in $X$,
with $X$, $X'$, and $X''$ sober; \cf c)). Show that when this
condition fails, the infimum of two elements in the ordered set of
subtopoi of $E$ is not distributive with respect to arbitrary
$\Sup$. (The editor does not know whether it is distributive with
respect to finite $\Sup$; \cf, however, \hyperref[IV.exer9.1.11]{9.1.11} e)
and \hyperref[IV.exer9.1.12]{9.1.12} d)).
\end{enumerate}
\end{enonce*}

\begin{enonce*}[remark]{Exercise 9.1.11}\etiquette[9.1.11]{IV.exer9.1.11}
(\emph{Subtopoi of locally Noetherian topoi. The case of Noetherian
spaces}).
\begin{enumerate}
\renewcommand{\theenumi}{\alph{enumi}}
\renewcommand{\labelenumi}{\theenumi)}
\item Let $f:E'\rightarrow E$ be an embedding of topoi. Show that if
$X$ is a pre-Noetherian object of $E$ (\ie (\sisi{VI 1.}{\hyperrefext[VI][VI.def1.30]{1.30}}) every
increasing sequence of subobjects\pageoriginale of $X$ is stationary),
then $X'=f^{\ast}(X)$ is a pre-Noetherian object of $E'$.
(By introducing the induced topoi ${E'}_{/X'}$ and $E_{/X}$ and
using \hyperref[IV.exer9.1.6]{9.1.6} e), first reduce to the case where $X$
is the final object $e$ of $E$, so that $X'$ is the final object $e'$
of $E'$. Then observe that every increasing sequence of subobjects
of $e'$ defines, upon applying $f_{\ast}$, an increasing sequence of
subobjects of $f_{\ast}(e')=e$.) Conclude that if $E$ admits a
generating family consisting of pre-Noetherian objects (\resp if $E$
is a Noetherian topos, \resp locally Noetherian (\sisi{VI 2.)}{\hyperrefext[VI][VI.def2.11]{2.11}})), then
$E'$ satisfies the same condition. (Use \hyperref[IV.exer9.1.9]{9.1.9} b).)

\item Let $E$ be a locally Noetherian topos (\sisi{VI 2.}{\hyperrefext[VI][VI.def2.11]{2.11}}). Show that
every subtopos of $E$ is of the form $E_Z$ (\hyperref[IV.exer9.1.8]{9.1.8}
a)), for a suitable full subcategory $Z$ of $\Ppoint(E)$. (Use a)
and \textsc{Deligne}'s theorem that every locally Noetherian topos
has enough points.) Conclude that if $E$ is of the form $\Top(X)$,
where $X$ is a sober locally Noetherian topological space, then
$Z\mapsto E_Z$ is an isomorphism from the ordered set of sober
subspaces of $X$ to the ordered set of subtopoi of $E$. (Use
\hyperref[IV.exer9.1.8]{9.1.8} e).) The intersection of sober subspaces $Z_i$
is necessarily sober and defines the intersection of the subtopoi
$E_{Z_i}$ (in contrast with what can occur in the general case
(\hyperref[IV.exer9.1.10]{9.1.10} c)).

\item Let $X$ be a sober locally Noetherian space, let $X'$ be a
subset of $X$, and let $X''$ be its complement. Show that the
following conditions are equivalent:
\begin{enumerate}
\renewcommand{\theenumii}{\roman{enumii}}
\renewcommand{\labelenumii}{(\theenumii)}
\item $X'$ is constructible.

\item $X'$ and $X''$ are sober.

\item $X'$ is sober, and the subtopos $E_{X'}\approxeq \Top(X')$ of
$E=\Top(X)$ admits a ``complementary'' subtopos (\cf
\sisi{\hyperref[IV.exer9.1.13]{9.1.13} e)}{\hyperref[IV.exer9.1.13]{9.1.13}}).

\item The intersection of the subtopoi $E_{X'}\approx\Top(X')$ and
$E_{X''}\approx\Top(X'')$ of $E=\Top(X)$ is the ``empty'' subtopos
of $E$.
\end{enumerate}

Show that if $E'$ and $E''$ are two subtopoi of $E$, they are weak
complements if and only if they are complements
(\hyperref[IV.exer9.1.13]{9.1.13}).

\item Let\pageoriginale $X$ be a sober Noetherian space. Show that
the following conditions are equivalent:
\begin{itemize}
\item[(i)] Every sober subspace $X'$ of $X$ is constructible, \ie,
according to c), its complement is sober.

\item[(i bis)] For every maximal point $x$ of $X$, $X-\{x\}$ is
sober, \ie $\{x\}$ is constructible.

\item[(ii)] $X$ is finite.

\item[(iii)] Every subtopos of $\Top(X)$ admits a complementary
subtopos (\cf \hyperref[IV.exer9.1.13]{9.1.13}).

\item[(iv)] In the ordered set of subtopoi of $\Top(X)$, the
infimum of two elements is distributive with respect to arbitrary
$\Sup$.
\end{itemize}
(Use the argument of \hyperref[IV.exer9.1.10]{9.1.10} d), together with the
intersection formula proved in b).)

\item Let $X$ be a locally Noetherian space. Show that in the ordered
set of subtopoi of $\Top(X)$, the infimum of two elements is
distributive with respect to \emph{finite} $\Sup$. (Reduce to the
case where $X$ is sober, and use b).)

\item For every sober subspace $Z$ of the sober locally Noetherian
space $X$, let $T_Z$ be the topology on the category $\Ouv(X)$ of
open subsets of $X$ for which a family of inclusions
$U_i\rightarrow U$ is covering if and only if $U\cap
Z=\bigcup_i U_i \cap Z$. Show that $Z\mapsto T_Z$ establishes a
bijection (reversing inclusion relations) between the set of sober
subspaces $Z$ of $X$ and the set of topologies on $\Ouv(X)$ finer
than the canonical topology $T_X$. (Use b) and
\hyperref[IV.rem9.1.2]{9.1.2} c).)

\item Let $X$ be a sober topological space. A subset $X'$ of $X$ is
sober if and only if its complement is a union of locally closed
subsets; if $X$ is locally Noetherian, this also means that $X'$ is
``proconstructible'' (EGA IV 1.9.4). Show that there is a topology
$T_{\cons'}$ on $X$ such that the closed subsets of
$X_{\cons'}=(X,T_{\cons'})$ are the sober subsets of $X$. Assume
henceforth that $X$ is locally Noetherian, so that $X_{\cons'}$ is
none other than the space $X_{\cons}$\pageoriginale of EGA IV 9.1.13.
Deduce from b) and c) that the ordered set $\Sigma$ of subtopoi of
$\Top(X)$ is isomorphic to the ordered set of closed subsets of
$X_{\cons}$, and that, under this correspondence, the subtopoi
admitting a complement correspond to the subsets that are both open
and closed in $X_{\cons}$. Show that \sisi{$X$}{$X_{cons}$} is locally compact,
and that it is compact if and only if $X$ is Noetherian, \ie the
topos \sisi{$\Top(E)$}{$\Top(X)$} is Noetherian.

\item Let $E$ be a locally Noetherian topos. Suppose that it has the
property considered in \hyperref[IV.subsubsec9.1.14]{9.1.14} b) (a condition
that may perhaps always hold\ldots). Show that the essential images
of the canonical functors $\Ppoint(E')\rightarrow \Ppoint(E)$, where
$E'$ ranges over the set $\Sigma$ of subtopoi of $E$, define a
collection of subsets of the set $X=\Point(E)$ of isomorphism classes
of points of $E$ which is the collection of closed subsets for a
\emph{topology} $T_{\cons}$ \emph{on} $X$, and that in this way one
obtains an isomorphism of ordered sets between the set $\Sigma$ of
subtopoi of $E$ and the set of closed subsets of $X_{\cons}$. Conclude
that in $\Sigma$, the $\Sup$ of two elements is distributive with
respect to \emph{arbitrary} infima. Show that $X_{\cons}$, which is not
necessarily $𝒰$-small, has an associated sober space
(\hyperref[IV.subsubsec4.2.1]{4.2.1}) which is $𝒰$-small.
\end{enumerate}
\end{enonce*}

\begin{enonce*}[remark]{Exercise 9.1.12}\etiquette[9.1.12]{IV.exer9.1.12}
(\emph{Finite topoi}). A topos $E$ is said to be \emph{finite} if it
is equivalent to a topos of the form $\widehat{C}$, with $C$ a
finite category.
\begin{enumerate}
\renewcommand{\theenumi}{\alph{enumi}}
\renewcommand{\labelenumi}{\theenumi)}
\item (Dictionary). Let $𝒱$ be a universe such that $𝒰\in 𝒱$. Let
(Karfiness) be the 2-category defined as the full 2-subcategory of
the category $(\Cat)_{𝒱}$ consisting of categories that are elements
of $𝒱$, are \emph{Karoubian and equivalent to a finite category},
and let (Topfin) be the 2-category defined as the full 2-subcategory
of $(𝒱\text{-}𝒰\text{-}\Top)$
(\hyperref[IV.subsubsec3.3.1]{3.3.1}) consisting of the \emph{finite}
$𝒰$-topoi that are $\in 𝒱$. Show that there are mutually
quasi-inverse 2-equivalences:
\hypertarget{sga4.exp.IV.equation.9.1.12.1}{}\begin{align*}
C & \longmapsto \widehat{C}: (\text{Karfiness})\longrightarrow
(\text{Topfin})\tag{9.1.12.1}\\
E & \longmapsto \Ppoint(E): (\text{Topfin})\longrightarrow
(\text{Karfiness}).\tag{9.1.12.2.}
\end{align*}\etiquette[9.1.12.2]{IV.eq9.1.12.2}
(Use \hyperref[IV.exer7.6]{7.6} h), e).)

\item Show\pageoriginale that a finite topos is Noetherian
\sisi{(VI 2)}{\hyperrefext[VI][VI.def2.11]{2.11}}.

\item Show that every subtopos of a finite topos $E$ is a finite
topos. More precisely, if $E$ is equivalent to $\widehat{C}$, where
$C$ is a finite Karoubian category (or, more generally, is equivalent
to a finite category), show that one obtains an order isomorphism
between the set of subtopoi $E'$ of $E$ and the set of strictly full
subcategories $C'$ of $C$ that are Karoubian (or, equivalently,
stable under taking direct factors in $C$), by associating with every
$C'$ the essential image of
$f_{\ast}:\widehat{C}'\rightarrow \widehat{C}$, where
$f:C'\rightarrow C$ is the inclusion functor. (Use
\hyperref[IV.subsec5.6]{5.6} to verify that $f_{\ast}$ is fully faithful, and
b) and \hyperref[IV.exer9.1.11]{9.1.11} b) for the fact that every subtopos
$E'$ of $E$ is obtained in the indicated way.)

\item Let $E$ be a finite topos. Construct on the finite set
$X=\Point(E)$ of isomorphism classes of points of $E$ a topology
making it a sober topological space and such that the ordered set of
subtopoi of $E$ is canonically isomorphic to the set of closed subsets
of $X$. (Take the ``constructible topology'' $\Cc_{\cons}$ on $X$,
defined by the order relation $(x\leq y)\Leftrightarrow$ ($x$ is a
direct factor of $y$), for which the closure of a point $x$ consists
of the $y$ such that $y\leq x$.) In particular, the set $\Sigma$ of
subtopoi of $E$ is finite, and the binary infimum operation $\wedge$
is distributive with respect to arbitrary $\Sup$.

\item Let $C$ be a category equivalent to a finite category. Show
that, for every full subcategory $C'$ of $C$, there exists a topology
$T_{C'}$ on $C$ for which a family $(X_i\rightarrow X)_{i\in I}$ is
covering if and only if, for every $Y\in \Ob C'$, the family of maps
$\Hom(Y,X_i)\rightarrow \Hom(Y,X)$ $(i\in I)$ is surjective. Show
that the category of sheaves on $(C,T_{C'})$ is equivalent to
$\widehat{C}'$, and that $C'\mapsto T_{C'}$ is a bijection (reversing
the order relations) between the set of strictly full subcategories
$C'$ of $C$ that are Karoubian (or, equivalently, stable under taking
direct factors in $C$), and the set of topologies
(\hyperrefext[II][II.def1.1]{1.1}) on the category $C$. (Use $c$) and
\hyperref[IV.rem9.1.2]{9.1.2} c).)\pageoriginale In particular, if $C$ is a
site whose underlying category is equivalent to a finite category,
then the topos $C^{\sim}$ is a finite topos.

\item Let $f:C'\rightarrow C$ be a functor of $𝒰$-categories, with
$C$ equivalent to a finite category. Show that the morphism of topoi
$\widehat{f}:\widehat{C}'\rightarrow \widehat{C}$
(\hyperref[IV.subsec4.6]{4.6}) is conservative if and only if every object of
$C$ is isomorphic to a direct factor of an object in the image of
$f$. (Reduce to the case where $f$ is an inclusion
$f:C'\hookrightarrow C$, with $C'$ as in c).)

\item Let $C$ be the category having two objects $e$ (the final
object) and $a$, and, besides the identity arrows, three arrows
$f:a\rightarrow e$, $g:e\rightarrow a$, $p=gf$ (thus $p^2=p$).
Show that $E=\widehat{C}$ has, besides the subtopoi $E$ and
$\Top(∅)$, exactly one subtopos $E'$, namely
$E'=\widehat{\{e\}}$; consequently, for every subtopos $E''$ of $E$
with $E''\neq \Top(∅)$, one has $E'\cap E''\neq \Top(∅)$.
\end{enumerate}
\end{enonce*}

% END INCLUDED SGA4_English_current_buildable_source_20260906/source/04/drafts/SGA4_IV_LINES_5779_6107_English.texfrag


% BEGIN INCLUDED SGA4_English_current_buildable_source_20260906/source/04/drafts/SGA4_IV_LINES_6108_6401_English.texfrag
\begin{enonce*}[remark]{Exercise 9.1.13}\etiquette[9.1.13]{IV.exer9.1.13}
(\emph{Complementary subtopoi of a topos}).

Let $E$ be a topos and let $E'$ and $E''$ be two subtopoi. We say
that $E'$ and $E''$ are \emph{complementary} if
$E'\cap E''=\Top(∅)$ and $E'\vee E''=E$ (where $\vee$ denotes the
$\Sup$ in the ordered set of subtopoi of $E$).

Suppose that in the set $\Sigma$ of subtopoi of $E$, the infimum of
two elements is distributive over finite suprema.

\begin{enumerate}
\renewcommand{\theenumi}{\alph{enumi}}
\renewcommand{\labelenumi}{\theenumi)}
\item If $E'$ and $E''$ are complementary, then $E''$ is the largest
among the subtopoi $E''_1$ such that
$E'\cap E''_1=\Top(∅)$ (\ie $E''$ is a weak complement of $E'$, in
the terminology of \hyperref[IV.exer9.1.10]{9.1.10} d)).

\item Show that the following conditions on $E$ are equivalent:
\begin{itemize}
\item[(i)] If $E'$ and $E''$ are two subtopoi of $E$ such that $E''$
is a weak complement of $E'$, then $E''$ is a complement of $E'$.

\item[(ii)] For every subtopos $E'$ of $E$ such that $E'\neq E$,
there exists a subtopos $E''$ of $E$ such that
$E''\neq\Top(∅)$ and $E'\cap E''=\Top(∅)$.

\item[(i bis)] If $E'$ and $E''$ are two subtopoi of $E$ such that
$E''$ is maximal among the subtopoi $E''_1$ satisfying
$E'\cap E''_1=\Top(∅)$, then $E''$ is a complement of $E'$.
\end{itemize}

Show that even if $E$ is a finite topos
(\hyperref[IV.exer9.1.12]{9.1.12}), these\pageoriginale conditions
need not hold (\hyperref[IV.exer9.1.12]{9.1.12} g)). They do hold,
however, if $E$ is of the form $\Top(X)$, with $X$ a locally
Noetherian space. (Use \hyperref[IV.exer9.1.11]{9.1.11} b)).

\item A subtopos $E'$ of $E$ is called \emph{complemented} if it
admits a complement (compare \hyperref[IV.exer9.1.12]{9.1.12} c)).
Prove that the set of complemented subtopoi is stable under finite
infima and suprema and under complementation, and that
complementation transforms infima into suprema and suprema into
infima.

\item State the duals of observations a), b), and c) (which were in
fact trivial general statements about an abstract ordered set
$\Sigma$; note that the hypothesis made on $\Sigma$ is itself
self-dual).
\end{enumerate}
\end{enonce*}

\setcounter{subsubsection}{13}
\subsubsection{}\etiquette[9.1.14]{IV.subsubsec9.1.14}
\emph{Open questions}\nde{Positive answers to questions a) and d)
were given in [17].}. Let $E$ be a topos.

\begin{itemize}
\item[a)] In the set of subtopoi of $E$, is the infimum of two
elements distributive over finite suprema?

\item[b)] Let $E'$ and $E''$ be two subtopoi of $E$ such that
$E=E'\vee E''$. Is every point of $E$ then isomorphic to the image of
a point of $E'$ or a point of $E''$?

Note that if the answer to b) is affirmative for $E$ and all its
subtopoi, then the answer to a) is affirmative at least for subtopoi
having enough points (and hence for all subtopoi if $E$ is locally
Noetherian (\hyperref[IV.exer9.1.11]{9.1.11} b)). The answer to a)
(and a fortiori to b)) is nevertheless not known for all Noetherian
topoi. The answers to a) and b) are affirmative, however, if $E$ is
of the form $\Top(X)$, with $X$ a locally Noetherian space
(\hyperref[IV.exer9.1.11]{9.1.11}), or if $E$ is a finite topos
(\hyperref[IV.exer9.1.12]{9.1.12}).

Note that a) can be reformulated as follows (applied to all subtopoi
$F$ of $E$):

\item[$\rm a'$)] If $E'$ and $E''$ are two subtopoi of a topos $F$
such that the canonical morphism (\cf \hyperref[IV.exer8.7]{8.7} b))
\begin{equation*}
E'\amalg E''\longrightarrow F\tag{9.1.14.1}
\end{equation*}\etiquette[9.1.14.1]{IV.eq9.1.14.1}
is conservative (which also means that $F=E'\vee E''$
(\hyperref[IV.exer9.1.7.2]{9.1.7.2} e)), then the same is true after
2-base change by $F_1\hookrightarrow F$,\pageoriginale where $F_1$
is a subtopos of $F$.

There is an analogous reformulation $\rm b'$) of b), taking a base
change $F_1\rightarrow F$ by a \emph{punctual} topos $F_1$
(\hyperref[IV.subsec2.2]{2.2}). Thus affirmative answers to a) and b)
would follow from an affirmative answer to the following question
(where $F$ is any subtopos of $E$):

\item[c)] If $E'$ and $E''$ are two subtopoi of the topos $F$ whose
supremum is $F$, \ie such that the morphism of topoi
\eqhyperref[IV.eq9.1.14.1]{9.1.14.1} is conservative, is this
morphism even \emph{universally conservative}, \ie does it remain
conservative after every 2-base change $F_1\rightarrow F$?

This statement makes sense by a result of \textsc{Deligne} (see [16],
Proposition 4.47) asserting the existence of 2-fiber products of
topoi. See, however, the example in
\hyperref[IV.exer9.1.7.2]{9.1.7.2} d).

On the other hand, \hyperref[IV.exer9.1.11]{9.1.11} h) shows that an
affirmative answer to b) would, in the locally Noetherian case,
provide an affirmative answer to the following question:

\item[d)] In the ordered set $\Sigma$ of subtopoi of $E$, is the
infimum of two elements distributive over \emph{arbitrary} (not
necessarily finite) suprema?

This also means that $\Sigma$ may be interpreted as the ordered set
of ``closed subtopoi'' of a suitable topos $E_{\cons'}$ (namely
$\Sigma^{0^{\sim}}$, where the opposite ordered set $\Sigma^0$ of
$\Sigma$, regarded as a category, is endowed with its canonical
topology, \cf \hyperref[IV.exer7.8]{7.8} b)), which is generated by
the subobjects of the terminal object and is therefore very close to
an ordinary topological space. It would remain to study this topos
$E_{\cons'}$, taking inspiration from the examples
\hyperref[IV.exer9.1.11]{9.1.11} g) and
\hyperref[IV.exer9.1.12]{9.1.12} d), and in particular to determine
whether it is Noetherian when $E$ is. This amounts to the following
question for subtopoi $F$ of $E$, which makes sense independently of
d):

\item[e)] Let $F$ be a Noetherian topos, and let $(E_i)_{i\in I}$ be
a family of subtopoi of $F$ whose intersection is $\Top(∅)$. Is there
a finite subfamily with the same property?

In a rather different direction, recall also the
question\pageoriginale raised in \hyperref[IV.exer9.1.6]{9.1.6} e),
which deserves clarification.
\end{itemize}

\subsection{The case of open subtopoi}\etiquette[9.2]{IV.subsec9.2}

\subsubsection{}\etiquette[9.2.1]{IV.subsubsec9.2.1}

Let $U$ be an open of a topos $E$, \ie a subobject of the terminal
object $e$ of $E$ (\hyperref[IV.def8.3]{8.3}). Consider the
localization morphism (\hyperref[IV.subsec5.2]{5.2})
\begin{equation*}
j:E_{/U}\longrightarrow E{\quoi}.\tag{9.2.2}
\end{equation*}\etiquette[9.2.2]{IV.eq9.2.2}
Since $U\rightarrow e$ is a monomorphism, the functor $j_{!}$ (which
is interpreted as the forgetful functor) is fully faithful, and
hence so is its biadjoint $j_{\ast}$
(\sisi{1.5.7 a)}{\hyperrefext[I][I.rem5.7]{5.7}}). In other words,
\emph{the localization morphism $j$ in \hyperref[IV.eq9.2.2]{(9.2.2)} associated with an
open $U$ of the topos $E$ is an embedding of topoi}
(\hyperref[IV.def9.1.1]{9.1.1} a)).

An embedding of topoi $F\rightarrow E$ is called an \emph{open
embedding} if it is isomorphic to the embedding of topoi defined by
an open $U$ of $E$. A subtopos $E'$ (\hyperref[IV.def9.1.1]{9.1.1}
b)) of $E$ is called an \emph{open subtopos} if it is defined by an
open of $E$, \ie if the canonical inclusion morphism
$E'\rightarrow E$ is an open embedding. Note that the open $U$ of
$E$ associated with an open embedding $j:E'\rightarrow E$ is
uniquely determined as the subobject $j_{!}(e')$ of $e$, where $e'$
denotes the terminal object of $E'$. Thus one obtains a one-to-one
correspondence between opens $U$ of $E$ and open subtopoi of $E$.

\begin{enonce*}[remark]{Remark 9.2.3}\etiquette[9.2.3]{IV.rem9.2.3}
In accordance with the preceding definitions, the open subtopos of
$E$ associated with the open $U$ of $E$ is the strictly full
subcategory of $E$ which is the essential image of the direct-image
functor
$$
j_{\ast}:E_{/U}\longrightarrow E.
$$
One should not confuse this subcategory of $E$ with the strictly full
subcategory which is the essential image of $j_{!}$, formed by the
objects $X$ of $E$ for which the structural morphism $X\rightarrow e$
factors through $U$. These two subcategories plainly determine one
another and are canonically equivalent. This is why some authors may
have been tempted to call (and indeed have called) the strictly full
subcategory which is the essential image of $j_{!}$ the ``open
subtopos defined by $U$,'' since its description in terms of $U$ is
simpler than that\pageoriginale of the essential image of
$j_{\ast}$. This terminology causes no difficulty so long as one does
not intend to study kinds of subtopoi other than open subtopoi. Since
that is not our situation, we shall not follow the authors just
mentioned.

For the reader's convenience, we summarize the most important special
properties of open embeddings:
\end{enonce*}

\begin{enonce*}{Proposition 9.2.4}\etiquette[9.2.4]{IV.prop9.2.4}
Let
$$
j:\mathcal{U}\longrightarrow E
$$
be an open embedding of topoi
(\hyperref[IV.subsubsec9.2.1]{9.2.1}). Then the functor $j_{!}$ left
adjoint to $j^{\ast}$ exists, so that there is a sequence of three
adjoint functors
\[
\xymatrix@R=.35cm@C=.5cm{
           & j_!,\quad j^{\ast}, \quad j_{\ast} : &\\
&&\\
\mathcal{U} \ar @<-16pt>[rr]^{j_{\ast}} \ar
           @<14pt>[rr]^{j_{!}}&            &\ar @<1pt>[ll]_{j^{\ast}} E
}
\]

Moreover:
\begin{enumerate}
\renewcommand{\theenumi}{\alph{enumi}}
\renewcommand{\labelenumi}{\theenumi)}
\item The functors $j_{!}$ and $j_{\ast}$ are fully faithful.

\item The functor $j_{!}$ commutes with fiber products, with products
indexed by small sets $I\neq\phi$, and with projective limits indexed
by small cofiltered categories (\hyperrefext[I][I.def2.7]{2.7}).

\item For every object $X$ of $E$, the adjunction morphism
$$
j_{!}j^{\ast}(X)\longrightarrow X
$$
is a monomorphism.
\end{enumerate}
\end{enonce*}

\noindent \emph{Proof}. We may suppose that $j$ is the localization
morphism defined by an open $U$ of $E$. Then a) is recalled for
reference, b) follows from the interpretation of $j_{!}$ as a
forgetful functor and from the fact that $U\rightarrow e$ is a
monomorphism. Finally, c) is immediate on observing that the morphism
in question is obtained from the inclusion $U\rightarrow e$ by base
change along $X\rightarrow e$, since a monomorphism is carried to a
monomorphism by base change.

\begin{enonce*}[remark]{Remark 9.2.4.1}\etiquette[9.2.4.1]{IV.rem9.2.4.1}
There exist embeddings of topoi $j:E'\rightarrow E$ for which
$j_{!}$ exists but does not commute with products of two objects or
with cofiltered projective limits (Exercise
\hyperref[IV.exer9.5.9]{9.5.9} c)).
\end{enonce*}

\setcounter{subsubsection}{4}
\subsubsection{}\etiquette[9.2.5]{IV.subsubsec9.2.5}
Now\pageoriginale let
$$
g:E'\longrightarrow E
$$
be an arbitrary morphism of topoi, and let $U'=g^{\ast}(U)$ be the
inverse image of the open $U$ of $E$. We saw in
\hyperref[IV.prop5.11]{5.11} that the diagram of topoi
\[
\xymatrix@C=2cm@R=1.5cm{
E'/_{U'} \ar[d]_{g_U} \ar[r]^{j'}& E' \ar[d]^{g}\\
E/_{U}  \ar[r]^{j} & E }
\]
is 2-Cartesian. It follows in particular that the inverse image under
$g$ of an open subtopos of $E$ (\cf
\hyperref[IV.exer9.1.6]{9.1.6} d)) is an open subtopos of $E'$, or
equivalently that the notion of an open embedding of topoi
(\hyperref[IV.subsubsec9.2.1]{9.2.1}) is stable under 2-base changes
in the 2-category of $𝒰$-topoi (which are elements of a given
universe $𝒱$).

Also note that the morphism of topoi $g:E'\rightarrow E$ factors, up
to isomorphism, through $j:E_{/U}\rightarrow E$ if and only if
$j':E'_{/U'}\rightarrow E'$ is an equivalence of categories, which
also means that the inclusion $U'\hookrightarrow e'$ (where $e'$ is
the terminal object of $E'$) is an isomorphism. This makes
\hyperref[IV.prop9.1.4]{9.1.4} more precise by giving a simple
characterization of the essential image of the functor considered
in \loccit

\subsubsection{}\etiquette[9.2.6]{IV.subsubsec9.2.6}
Apply \hyperref[IV.subsubsec9.2.5]{9.2.5} when $E'$ is of the form
$E_{/V}$, where $V$ is a second open of $E$, and $g$ is the
localization morphism. Then $U'=U\cap V$, and the 2-Cartesian product
of $E_{/U}$ and $E_{/V}$ over $E$ identifies with
$E_{/U\cap V}$. It follows that \emph{a finite intersection of open
subtopoi of $E$ is an open subtopos of $E$}; more precisely, the map
\begin{equation*}
U\longmapsto \text{ the subtopos of $E$ defined by $U$}
\tag{9.2.6.1}
\end{equation*}\etiquette[9.2.6.1]{IV.eq9.2.6.1}
commutes with finite intersections.

\begin{enonce*}[remark]{Exercise 9.2.7}\etiquette[9.2.7]{IV.exer9.2.7}
Prove that the map \eqhyperref[IV.eq9.2.6.1]{9.2.6.1} also commutes with \emph{arbitrary}
$\Sup$. (Use \hyperref[IV.exer9.1.7.2]{9.1.7.2} e)).
\end{enonce*}

% END INCLUDED SGA4_English_current_buildable_source_20260906/source/04/drafts/SGA4_IV_LINES_6108_6401_English.texfrag


% BEGIN INCLUDED SGA4_English_current_buildable_source_20260906/source/04/drafts/SGA4_IV_LINES_6402_6597_English.texfrag
\subsection{Construction of the closed subtopos complementary to an
open subtopos}\etiquette[9.3]{IV.subsec9.3}\pageoriginale

\subsubsection{}\etiquette[9.3.1]{IV.subsubsec9.3.1}
Let $\mathsf{T}$ be a topological space and $U$ an open subset of
$\mathsf{T}$, so that $\Top(U)$ identifies with an open subtopos of
$\Top(\mathsf{T})$ (notation of \hyperref[IV.subsec2.1]{2.1}). Let
$Y$ be the closed topological subspace of $\mathsf{T}$ complementary
to $U$. Up to equivalence, $\Top(Y)$ may then be regarded as a
subtopos of $\Top(\mathsf{T})$, namely the subtopos formed by the
objects $F$ of $\Top(\mathsf{T})$ whose restriction to $U$ is
\sisi{the terminal topos of $U$}{the terminal object of $U$}. This
description of a subcategory of $E=\Top(\mathsf{T})$ makes sense
whenever one has a topos and an open $U$ of it, and we shall see that
it always gives a subtopos of $E$. Moreover, in the particular case
first considered, it is immediate (with the terminology introduced
in Exercise \hyperref[IV.exer9.1.13]{9.1.13}) that $\Top(Y)$ and
$\Top(U)$ are complementary subtopoi (use
\hyperref[IV.subsubsec9.2.5]{9.2.5} and
\hyperref[IV.exer9.1.7.2]{9.1.7.2} e)). The same will hold in the
general case, and we shall see that this property uniquely
characterizes the subtopos in question. It therefore deserves in
every respect the name of the closed subtopos complementary to the
given open $U$, or to the open subtopos defined by $U$. The details
of the construction of this topos $\mathcal{F}$ and of the
inverse-image functor $i^{\ast}:E\rightarrow\mathcal{F}$ will be
given in the present section, while Section
\hyperref[IV.subsec9.4]{9.4} develops its first properties.

\subsubsection{}\etiquette[9.3.2]{IV.subsubsec9.3.2}
Thus let $E$ be a topos and $U$ an open of $E$, as in
\hyperref[IV.subsubsec9.2.1]{9.2.1}, and consider the localization
morphism
\begin{equation*}
j:\mathcal{U}=E_{/U}\longrightarrow E{\quoi},\tag{9.3.2.1}
\end{equation*}\etiquette[9.3.2.1]{IV.eq9.3.2.1}
which is an open embedding.

For every object $X$ of $E$, put
\begin{equation*}
X_{\complement U}=U\coprod_{U\times X}
X{\quoi},\tag{9.3.2.2}
\end{equation*}\etiquette[9.3.2.2]{IV.eq9.3.2.2}
where the amalgamation is formed with respect to the canonical
projections
$$
\pr_1:U\times X\longrightarrow U,\quad \pr_2:U\times
X\longrightarrow X.
$$
Thus there is a \emph{cocartesian} diagram
(\hyperrefext[I][I.sec10]{10}) in $E$, depending functorially on
$X$:\pageoriginale
\[\vcenter{
\xymatrix@C=1.75cm@R=1.5cm{
U\times X \ar[r]^{\pr_2} \ar[d]_{\pr_1}& X \ar[d]\\
U \ar[r]        & X_{\complement U} }}\leqno{(9.3.2.3)}
\]\etiquette[9.3.2.3]{IV.eq9.3.2.3}
Note that in the example considered in
\hyperref[IV.subsubsec9.3.1]{9.3.1}, $X_{\complement U}$ is
canonically isomorphic to $i_{\ast}i^{\ast}(X)$ (where
$i:Y\rightarrow T$ is the inclusion), and
$X\rightarrow X_{\complement U}$ identifies with the adjunction
morphism.

\begin{enonce*}{Proposition 9.3.3}\etiquette[9.3.3]{IV.prop9.3.3}
With the preceding notation, the following statements hold:
\begin{enumerate}
\renewcommand{\theenumi}{\alph{enumi}}
\renewcommand{\labelenumi}{\theenumi)}
\item For every object $X$ of $E$, the following conditions are
equivalent:
\begin{itemize}
\item[(i)] There exists an object $Y$ of $E$ and an isomorphism
$X\simeq Y_{\complement U}$.

\item[(ii)] The canonical morphism
$X\xrightarrow{m}X_{\complement U}$ is an isomorphism.

\item[(iii)] The canonical morphism
$\pr_2:X\times U\rightarrow U$ is an isomorphism (\ie
$j^{\ast}(X)$ is a terminal object of $E_{/U}$).
\end{itemize}

\item For every morphism
$$
m:X\longrightarrow X'
$$
in $E$, the following conditions are equivalent:
\begin{itemize}
\item[${\rm (i')}$] The diagram
\[
\xymatrix@C=2cm@R=1.5cm{
X \ar[r]^{m} \ar[d]& X'\ar[d]\\
X_{\complement U} \ar[r]& X'_{\complement U} }
\]
is Cartesian.

\item[${\rm (ii')}$] The morphism
$m\times\id_U:X\times U\rightarrow X'\times U$ is an isomorphism,
\ie $j^{\ast}(m)$ is an isomorphism.
\end{itemize}

\item The functor
$$
X\longrightarrow X_{\complement U}
$$
from $E$ to $E$ is left exact, carries epimorphic families to
epimorphic families, and commutes with filtered inductive limits
and\pageoriginale amalgamated sums.
\end{enumerate}
\end{enonce*}

\noindent \emph{Proof}. First prove the proposition when $E$ is of
the form $\widehat{C}$, where $C$ is a small category. By \hyperref[I.prop3.1]{I.3.1} this
reduces to the case where $E$ is the ``punctual'' topos $(\Ens)$,
where the proposition is evident. One then passes to the general
case by the standard procedures of \hyperrefext[II][II.sec4]{4},
using the ``associated sheaf'' functor.

\begin{enonce*}{Proposition 9.3.4}\etiquette[9.3.4]{IV.prop9.3.4}
With the notation of \hyperref[IV.subsubsec9.3.2]{9.3.2}, the
strictly full subcategory $\mathcal{F}$ of $E$ formed by the objects
$X$ of $E$ such that $j^{\ast}(X)$ is a terminal object of
$\mathcal{U}$ \textup{(\cf \hyperref[IV.prop9.3.3]{9.3.3} a))} is a
subtopos of $E$: \ie (\hyperref[IV.def9.1.1]{9.1.1}) it is a topos,
and the inclusion functor $i_{\ast}:\mathcal{F}\rightarrow E$ is the
direct-image functor of a morphism of topoi
$i:\mathcal{F}\rightarrow E$, whose inverse-image functor is
$X\mapsto i^{\ast}X=X_{\complement U}$. The adjunction morphism
$X\rightarrow i_{\ast}i^{\ast}X$ is the canonical morphism
$X\rightarrow X_{\complement U}$ of
\eqhyperref[IV.eq9.3.2.3]{9.3.2.3}.
\end{enonce*}

\noindent \emph{Proof}: For every object $X$ of $E$, denote by
$u(X)$ the canonical morphism $X\rightarrow X_{\complement U}$. The
morphism $u(X)$ is functorial in $X$, and for every object $X$,
$u(X_{\complement U})$ is an isomorphism. It follows formally and
immediately from these two properties that the functor $i^{\ast}$ is
left adjoint to the functor $i_{\ast}$ and that the adjunction
morphism is $u(X)$. Since the functor
$X\rightarrow X_{\complement U}$ is left exact, the functor
$i^{\ast}$ is left exact. It then follows from
\hyperrefext[II][II.theo5.5]{5.5} that $\mathcal{F}$ is a topos, and
from Definition \hyperref[IV.def3.1]{3.1} that the pair
$(i^{\ast},i_{\ast})$ is a morphism of topoi.

\setcounter{subsubsection}{4}
\subsubsection{}\etiquette[9.3.5]{IV.subsubsec9.3.5}
The subtopos $\mathcal{F}$ of $E$ described in
\hyperref[IV.prop9.3.4]{9.3.4} is called the \emph{closed subtopos of
$E$ complementary to the open $U$ of $E$}, or, equivalently, to the
open subtopos of $E$ corresponding to $U$. In accordance with
\hyperref[IV.def9.1.1]{9.1.1} b), the morphism
$i:\mathcal{F}\rightarrow E$ constructed in
\hyperref[IV.prop9.3.4]{9.3.4} is called the \emph{inclusion
morphism}. A full subcategory $\mathcal{F}$ of $E$ is called a
\emph{closed subtopos} of $E$ if there exists an open $U$ of $E$ for
which $\mathcal{F}$ is the closed subtopos complementary to $U$.
Note that this $U$ is uniquely determined as
$i_{\ast}(\phi_{\mathcal{F}})$, where $\phi_{\mathcal{F}}$ is the
initial object of $\mathcal{F}$, equal to $i^{\ast}(\phi_E)$; it is
also called \emph{the open of $E$ complementary to the closed
subtopos $\mathcal{F}$}, and\pageoriginale the subtopos of $E$
defined by $U$ is also called the \emph{open subtopos of $E$
complementary to $\mathcal{F}$}. Finally, an embedding of topoi
(\hyperref[IV.def9.1.1]{9.1.1} a)) $i:F\rightarrow E$ is called a
\emph{closed embedding} if the corresponding subtopos of $E$ is
closed.

If $G$ is a group object on $E$ and $s\in\Hom(e_E,G)$, the
\emph{support} of $G$ (\resp of $s$) is the closed subtopos of $E$
complementary to the cosupport open of $G$ (\resp of $s$)
(\hyperref[IV.subsubsec8.5.1]{8.5.1},
\hyperref[IV.subsubsec8.5.2]{8.5.2}).

\subsubsection{}\etiquette[9.3.6]{IV.subsubsec9.3.6}
With the notation of \hyperref[IV.subsubsec9.3.2]{9.3.2} and
\hyperref[IV.prop9.3.4]{9.3.4}, one therefore has two morphisms of
topoi
\begin{equation*}
\mathcal{U}\xrightarrow{j}
E\xleftarrow{i}\mathcal{F}{\quoi},\tag{9.3.6.1}
\end{equation*}\etiquette[9.3.6.1]{IV.eq9.3.6.1}
defining five functors which form two sequences of adjoint functors,
$(j_{!},j^{\ast},j_{\ast})$ and $(i^{\ast},i_{\ast})$:
\[
\xymatrix@C=2cm@R=1.5cm@M=.2cm{ \Uu \ar@<6pt>[r]^{j_!}
\ar@<-13pt>[r]|{j_{\ast}} & E \ar@<3pt>[r]^{i^{\ast}}
\ar@<4pt>[l]|{j^{\ast}}& \Ff
\ar@<13pt>[l]|{i_{\ast}}.}\leqno{(9.3.6.2)}
\]\etiquette[9.3.6.2]{IV.eq9.3.6.2}

The functor
\begin{equation*}
\rho=i^{\ast}j_{\ast}:\mathcal{U}\longrightarrow
\mathcal{F}\tag{9.3.6.3}
\end{equation*}\etiquette[9.3.6.3]{IV.eq9.3.6.3}
is called the \emph{gluing functor} relative to the open $U$ of $E$.
More generally, if $\mathcal{U}$ and $\mathcal{F}$ are arbitrary
topoi, a functor $\rho:\mathcal{U}\rightarrow\mathcal{F}$ is called
a \emph{gluing functor} if it is left exact and accessible (\hyperref[I.def9.2]{I.9.2}),
or equivalently
(\sisi{\hyperref[IV.rem1.8]{1.8}\quad}{\hyperrefext[I][I.sec8]{8}})
if it admits a pro-adjoint. The reasons for this terminology will
appear in \hyperref[IV.prop9.4.1]{9.4.1} d) and
\hyperref[IV.theo9.5.4]{9.5.4} below.

% END INCLUDED SGA4_English_current_buildable_source_20260906/source/04/drafts/SGA4_IV_LINES_6402_6597_English.texfrag


% BEGIN INCLUDED SGA4_English_current_buildable_source_20260906/source/04/drafts/SGA4_IV_LINES_6598_6787_English.texfrag
\subsection{First properties of the closed subtopos
$\mathcal{F}$ and of $i:\mathcal{F}\rightarrow
E$}\etiquette[9.4]{IV.subsec9.4}

\begin{enonce*}{Proposition 9.4.1}\etiquette[9.4.1]{IV.prop9.4.1}
The notation is that of \hyperref[IV.subsubsec9.3.2]{9.3.2} and
\hyperref[IV.prop9.3.4]{9.3.4}.
\begin{enumerate}
\renewcommand{\theenumi}{\alph{enumi}}
\renewcommand{\labelenumi}{\theenumi)}
\item The functor $i_{\ast}:\mathcal{F}\rightarrow E$ carries
epimorphic families to epimorphic families. It commutes with the
formation of amalgamated sums and filtered inductive limits.

\item For every object $X$ of $E$ such that the projection morphism
$U\times X\rightarrow U$ is an epimorphism, the adjunction morphism
$X\rightarrow i_{\ast}i^{\ast}X$ is an epimorphism.

\hypertarget{sga4.exp.IV.proposition.9.4.1.item.c}{}\item The\pageoriginale pair of functors $(i^{\ast},j^{\ast})$ from
$E$ to $\mathcal{F}$ and $\mathcal{U}$, respectively, is
conservative {\rm(I \hyperref[I.def6.1]{6.1})}, \ie a morphism $m$ of $E$ is an
isomorphism if and only if $i^{\ast}(m)$ and $j^{\ast}(m)$ are
isomorphisms (equivalently, the pair $(i^{\ast},j^{\ast})$ is
faithful (\hyperref[IV.subsubsec6.4.0]{6.4.0})).

\item The gluing functor
$\rho=i^{\ast}j_{\ast}:\mathcal{U}\rightarrow\mathcal{F}$ is left
exact and accessible {\rm(\hyperrefext[I][I.def9.2]{9.2})}, or
equivalently {\rm(\hyperrefext[I][I.sec8]{8})} it admits a pro-adjoint
$\sigma:\Pro\mathcal{F}\rightarrow\Pro\mathcal{U}$.
\end{enumerate}
\end{enonce*}

\noindent \emph{Proof}: a) follows from \hyperref[IV.prop9.3.3]{9.3.3} c). For b), it suffices
to use the definition \eqhyperref[IV.eq9.3.2.2]{9.3.2.2} of
$X_{\complement U}$, taking into account that
$i_{\ast}i^{\ast}X=X_{\complement U}$
(\hyperref[IV.prop9.3.4]{9.3.4}). Assertion c) follows from
\hyperref[IV.prop9.3.3]{9.3.3} b). The functors $i^{\ast}$ and
$j_{\ast}$ are left exact, and hence so is $\rho$. The functor
$j_{\ast}$ is a direct-image functor for a morphism of topoi and is
therefore accessible (\hyperrefext[I][I.prop9.5]{9.5}); note that a
topos is an accessible category
(\hyperrefext[I][I.rem9.11.3]{9.11.3}). Since the functor $i^{\ast}$
commutes with inductive limits (it is an inverse-image functor), the
functor $\rho$, being a composite of two accessible functors, is
accessible.

\begin{enonce*}{Proposition 9.4.2}\etiquette[9.4.2]{IV.prop9.4.2}
Let $E'$ be a topos. Then the functor $h\mapsto i\circ h$,
$$
\Hhomtop(E',\mathcal{F})\longrightarrow \Hhomtop(E',E),
$$
is fully faithful, and its essential image consists of the morphisms
of topoi $g:E'\rightarrow E$ such that $g^{\ast}(U)$ is an initial
object $\phi_{E'}$ of $E'$.
\end{enonce*}

By \hyperref[IV.prop9.1.4]{9.1.4}, the functor in question is fully
faithful, and $g$ belongs to its essential image if and only if, for
every object $X'$ of $E'$, the object $g_{\ast}(X')$ of $E$ belongs
to $\mathcal{F}$, \ie if
\begin{equation*}
j^{\ast}(g_{\ast}(X'))\simeq e_{\mathcal{U}}\text{ for every }X'\in
\Ob E'{\quoi}.\tag{*}
\end{equation*}
Put $U'=g^{\ast}(U)$ and consider the diagram of topoi
\[
\xymatrix@C=1.5cm{
E'_{/U'}\ar[r]^{j'} \ar[d]_{g_U}& E'\ar[d]^{g}\\
E_{/U}   \ar[r]^{j}             & E\quad , }
\]
There\pageoriginale is plainly a canonical isomorphism
$$
j^{\ast}(g_{\ast}(X'))\simeq g_{U\ast}({j'}^{\ast}(X')){\quoi}.
$$
The functor ${j'}^{\ast}$ is also plainly essentially surjective, so
condition (*) is equivalent to
\begin{equation*}
g_{U\ast}(Y')\simeq e_{\mathcal{U}}\text{ for every }Y'\in\Ob
E'_{/U'}{\quoi}.\tag{**}
\end{equation*}
Using the adjunction between $g_{U\ast}$ and $g_U^{\ast}$, one sees
immediately that this means that $g_U^{\ast}(Y)$ is an initial object
of $E'_{/U'}$ for every object $Y$ of $E_{/U}$; or, using the fact
that the initial object of $E'_{/U'}$ is strict
(\hyperrefext[II][II.def4.5]{4.5}), that it suffices for this to hold
for $Y=e_{\mathcal{U}}$, the terminal object of $\mathcal{U}$. But
this also means that the terminal object $U'$ of $E'_{/U'}$ is
initial, which is manifestly equivalent to $U'$ being an initial
object $\phi_{E'}$ of $E'$, \cqfd

\begin{enonce*}{Corollary 9.4.3}\etiquette[9.4.3]{IV.cor9.4.3}
Let $g:E'\rightarrow E$ be a morphism of topoi, let
$U'=g^{\ast}(U)$, let $\mathcal{U}'=E'_{/U'}$, and let
$\mathcal{F}'$ be the closed subtopos of $E'$ complementary to the
open $U'$; finally, let $i':\mathcal{F}'\rightarrow E'$ be the
inclusion morphism. Then the morphism of topoi
$g\circ i':\mathcal{F}'\rightarrow E$ factors, up to isomorphism,
through a morphism of topoi
$g_{\mathcal{F}}:\mathcal{F}'\rightarrow\mathcal{F}$, and the
corresponding diagram of topoi
\[\vcenter{
\xymatrix@=.5cm{ \Ff' \ar[rr]^{i'}\ar[dd]_{g_{\Ff}}  &&
E'\ar[dd]^{g}
\\
                          &            &  \\
\Ff  \ar[rr]^(.4){i}  &  \ar@/^.5pc/[ru]  &   E }}\leqno{(9.4.3.1)}
\]\etiquette[9.4.3.1]{IV.eq9.4.3.1}
is $2$-Cartesian {\rm(\cf \hyperref[IV.prop5.11]{5.11}).}
\end{enonce*}

This follows formally from \hyperref[IV.prop9.4.2]{9.4.2} and the
definitions.

With the terminology introduced in \hyperref[IV.exer9.1.6]{9.1.6}
d), one may therefore say in particular that the inverse image under
a morphism of topoi $g:E'\rightarrow E$ of a \emph{closed} subtopos
$\mathcal{F}$ of $E$ is a \emph{closed} subtopos $\mathcal{F}'$ of
$E'$ (namely the closed subtopos complementary to the open
$U'=g^{\ast}(U)$, where $U$ is the open of $E$ complementary to
$\mathcal{F}$).

\begin{enonce*}{Corollary 9.4.4}\etiquette[9.4.4]{IV.cor9.4.4}
Let $E$ be a topos, $\mathcal{U}$ an open subtopos of $E$, and
$\mathcal{F}$ the complementary closed subtopos.
\begin{enumerate}
\renewcommand{\theenumi}{\alph{enumi}}
\renewcommand{\labelenumi}{\theenumi)}
\item One has\pageoriginale
$$
\mathcal{U}\cap\mathcal{F}\approxeq\Top(∅)\quav
\mathcal{U}\vee\mathcal{F}=E{\quoi},
$$
where $\vee$ denotes the $\Sup$ in the set of all subtopoi of $E$
{\rm(\hyperref[IV.rem9.1.2]{9.1.2} c))}. (In the terminology
introduced in \hyperref[IV.exer9.1.13]{9.1.13}, $\mathcal{U}$ and
$\mathcal{F}$ are complementary subtopoi of $E$.)

\item $\mathcal{U}$ (\resp $\mathcal{F}$) is the largest among the
subtopoi $E'$ of $E$ such that
$E'\cap\mathcal{F}\approxeq\Top(∅)$ (\resp
$\mathcal{U}\cap E'\approxeq\Top(∅)$).
\end{enumerate}
\end{enonce*}

\noindent \emph{Proof.}
\begin{enumerate}
\renewcommand{\theenumi}{\alph{enumi}}
\renewcommand{\labelenumi}{\theenumi)}
\item The first relation amounts to saying that if $E'$ is a topos
and $g:E'\rightarrow E$ is a morphism of topoi factoring, up to
isomorphism, through both $\mathcal{U}$ and $\mathcal{F}$, then
$E'\approxeq\Top(∅)$. By \hyperref[IV.subsubsec9.2.5]{9.2.5} and
\hyperref[IV.prop9.4.2]{9.4.2}, the hypothesis means that
$g^{\ast}(U)$ is both an initial and a terminal object of $E'$, which
proves the assertion. For the second relation, note that by
\hyperref[IV.exer9.1.7.2]{9.1.7.2} e) it is equivalent to
\hyperref[IV.prop9.4.1]{9.4.1} c).

\item Suppose that $E'\cap\mathcal{F}\approxeq\Top(∅)$. By \hyperref[IV.cor9.4.3]{9.4.3},
$E'\cap\mathcal{F}$ is equivalent to the subtopos of $E'$
complementary to $g^{\ast}(U)$, where $g:E'\rightarrow E$ is the
inclusion morphism. Thus the condition means that
$U'=g^{\ast}(U)$ is a terminal object of $E'$, \ie
(\hyperref[IV.subsubsec9.2.5]{9.2.5}) that
$E'\subset\mathcal{U}$. Suppose that
$\mathcal{U}\cap E'\approxeq\Top(∅)$. By
\hyperref[IV.prop5.11]{5.11}, $\mathcal{U}\cap E'$ is equivalent to
$E'_{/U'}$, so the condition means that $U'$ is an initial object of
$E'$, \ie (\hyperref[IV.prop9.4.2]{9.4.2}) that
$E'\subset\mathcal{F}$.
\end{enumerate}

\begin{enonce*}{Corollary 9.4.5}\etiquette[9.4.5]{IV.cor9.4.5}
Let $(\mathcal{F}_i)_{i\in I}$ be a family of closed subtopoi of the
topos $E$, and for every $i\in I$ let $U_i$ be the open of $E$
complementary to $\mathcal{F}_i$. Then the intersection of the
$\mathcal{F}_i$ as subtopoi of $E$
(\hyperref[IV.prop9.1.3]{9.1.3}) is a closed subtopos whose
complementary open is $U=\Sup_iU_i$.
\end{enonce*}

This follows formally from the universal property of the
intersection as a 2-product and from
\hyperref[IV.prop9.4.2]{9.4.2}, taking into account that for a
morphism of topoi $g:E'\rightarrow E$,
$g^{\ast}(U)=\Sup_i g^{\ast}(U_i)$ is an initial object if and only
if all the $g^{\ast}(U_i)$ are initial.

Similarly one proves:

\begin{enonce*}{Corollary 9.4.6}\etiquette[9.4.6]{IV.cor9.4.6}
With the notation of \hyperref[IV.cor9.4.5]{9.4.5}, suppose that $I$ is finite. Then the
supremum\pageoriginale of the subtopoi $\mathcal{F}_i$ of $E$
{\rm(\hyperref[IV.rem9.1.2]{9.1.2} c))} is a closed subtopos of $E$,
whose complementary open is $\bigcap_iU_i$.
\end{enonce*}


% END INCLUDED SGA4_English_current_buildable_source_20260906/source/04/drafts/SGA4_IV_LINES_6598_6787_English.texfrag


% BEGIN INCLUDED SGA4_English_current_buildable_source_20260906/source/04/drafts/SGA4_IV_LINES_6788_7070_English.texfrag
\begin{enonce*}[remark]{Exercise 9.4.7}\etiquette[9.4.7]{IV.exer9.4.7}
(\emph{Gluing of subtopoi}). Let $E$ be a topos.
\begin{enumerate}
\renewcommand{\theenumi}{\alph{enumi}}
\renewcommand{\labelenumi}{\theenumi)}
\item Let $E'$ be a subtopos of $E$, let $(S_i)_{i\in I}$ be a family
of objects of $E$ covering the terminal object, and, for every $i\in I$,
let $S'_i$ be the inverse image of $S_i$ in $E'$, and let $E'_i\approxeq
E'_{/S'_i}$ be the subtopos of $E_i=E_{/S_i}$ which is the inverse image,
under the localization morphism $E_i\rightarrow E$, of the subtopos $E'$ of
$E$ (\hyperref[IV.exer9.1.6]{9.1.6} d) and \hyperref[IV.prop5.11]{5.11}). Show that, for
an object $X$ of $E$, $X$ belongs to $E'$ if and only if, for
every $i\in I$, the inverse image $X| S_i$ of $X$ in $E_i$ belongs
to $E'_i$.

\item Suppose that $E$ is of the form $C^{\sim}$, where $C$ is a
$𝒰$-site. Show that the presheaf on $C$
\begin{center}
$S\mapsto$ set of subtopoi of the topos
$(C_{/S})^{\sim}\approxeq E_{/\varepsilon(S)}$
\end{center}
is a sheaf. For arbitrary $E$, deduce that there
exists an object of $E$ which represents the functor
$$
\text{$S\mapsto$ set of subtopoi of the topos $E_{/S}$;}
$$

\item With the notation of a), show that the subtopos
$E'$ of $E$ is an open subtopos (\resp a closed subtopos)
if and only if, for every $i\in I$, the same is true of the
induced subtopos $E'_i$ of $E_i$.
\end{enumerate}
\end{enonce*}

\begin{enonce*}[remark]{Exercise 9.4.8}\etiquette[9.4.8]{IV.exer9.4.8}
(\emph{Interior, exterior, closure, and boundary of a
subtopos.}) Let $E$ be a topos, $E'$ a subtopos, and
$i:E'\rightarrow E$ the inclusion morphism.
\begin{enumerate}
\renewcommand{\theenumi}{\alph{enumi}}
\renewcommand{\labelenumi}{\theenumi)}
\item Show that $i_{\ast}(\phi_{E'})$ is the largest open $U$
of $E$ such that $E'| \mathcal{U} \approxeq \Top(∅)$, where
$\mathcal{U}$ is the open subtopos of $E$ defined by $U$. This
$U$, or $\mathcal{U}$, is called the \emph{exterior} of the
subtopos $E'$ of $E$, and the closed subtopos of $E$ complementary
to $U$ is called the \emph{closure} of $E'$.

\item Show that there exists a largest open $V$ of $E$ such that the
corresponding open subtopos $\mathcal{V}$ of $E$ is contained
in $E'$. (Use \hyperref[IV.exer9.2.7]{9.2.7}.) This $V$, or
$\mathcal{V}$, is called the \emph{interior} of the subtopos $E'$ of $E$. The
opens $U$ and $V$ of $E$ are disjoint $(U\cap V=\phi_E)$; the
closed subtopos of $E$ complementary\pageoriginale to
$\Sup(U,V)=\mathcal{U}\amalg \mathcal{V}$ is called the
\emph{boundary} of the subtopos $E'$ of $E$.

\item The subtopos $E'$ is an open subtopos (\resp a closed subtopos) of $E$ if
and only if it is equal to its interior (\resp if it is
equal to its closure). The boundary of $E'$ is empty if and
only if $E'$ is both an open and a closed subtopos of
$E$, or equivalently, if it is the open subtopos of $E$ defined by an
open $V$ of $E$ which is a \emph{direct summand} of the terminal object
$e$, \ie such that there exists a subobject $U$ of $e$ with $e\simeq
U\amalg V$. The exterior of $E'$ is ``empty'' \ie the
closure is equal to $E$, if and only if
$i:E'\rightarrow E$ is dominant (\hyperref[IV.exer8.8]{8.8}).

\item Let $S$ be an object of $E$, let $F=E_{/S}$ be the induced topos, and let $F'\simeq
F_{/S'}$ be the subtopos of $F$ induced by the subtopos $E'$ of $E$.
Show that the exterior (\resp the closure, \resp
the interior, \resp
the boundary) of $F'$ is induced by the exterior (\resp\ldots) of
$E'$.
\end{enumerate}
\end{enonce*}

\begin{enonce*}[remark]{Exercise 9.4.9}\etiquette[9.4.9]{IV.exer9.4.9}
(\emph{Locally closed subtopoi of a topos.}) A subtopos $F$
of a topos $E$ is called a \emph{locally closed subtopos} of $E$
if it is the intersection of an open subtopos and a
closed subtopos.
\begin{enumerate}
\renewcommand{\theenumi}{\alph{enumi}}
\renewcommand{\labelenumi}{\theenumi)}
\item Prove that the intersection of a finite family of locally closed subtopoi
of $E$ is locally closed. If $g:E'\rightarrow
E$ is a morphism of topoi, prove that the inverse image
(\hyperref[IV.exer9.1.6]{9.1.6} d)) of a locally closed subtopos of $E$ is
a locally closed subtopos of $E'$.

\item Let $X$ be a topological space, and let $E=\Top(X)$. For every
subset $X'$ of $X$, consider the subtopos $T(X')$ of $E$ which is the essential
image of $\Top(X')$ in $\Top(X)$. Show that if $X'$ is
a locally closed subset, $T(X')$ is a locally closed subtopos
of $E$, the converse being true if, in addition, $X$ and
$X'$ are assumed sober. Prove that $X'\mapsto T(X')$ defines an order
isomorphism between the set of locally closed subsets
$X'$ of $X$ and the set of locally closed subtopoi of~$E$.

\item Let $F$ be a subtopos of $E$. Show that every open subtopos
of $F$\pageoriginale is induced by an open subtopos of $E$.
(Use \hyperref[IV.exer9.1.8]{9.1.8} c).) Show that $F$ is a locally closed
subtopos of $E$ if and only if it is an open subtopos
of its closure (\hyperref[IV.exer9.4.8]{9.4.8} a)). Deduce, using
\hyperref[IV.exer9.4.8]{9.4.8} d), that the property for $F$
of being a
locally closed subtopos of $E$ is a local property
on $E$.

\item Let $F$ be a locally closed subtopos of $E$. Show that,
among all the ways of writing $F$ as the intersection of an
open subtopos $\mathcal{U}$ and a closed subtopos
$\mathcal{F}$ of $E$, there is one for which $\mathcal{U}$ is the largest
possible, and $\mathcal{F}$ the smallest possible. (Take
$\mathcal{F} =$ the closure of $F$, and $\mathcal{U} =$ the largest
open subtopos inducing on $\mathcal{F}$ the subtopos $F$.)
Show that the formation of $\mathcal{U}$ and $\mathcal{F}$ is compatible
with localization on $E$.
\end{enumerate}
\end{enonce*}

\begin{enonce*}[remark]{Exercise 9.4.10}\etiquette[9.4.10]{IV.exer9.4.10}
Develop the notions of a \emph{constructible, locally
constructible, quasi-constructible, and locally quasi-constructible subtopos},
following the model of the familiar notions for topological
spaces (EGA ${\rm 0_{III}}$ \hyperref[IV.subsec9.1]{9.1}, EGA IV 10.1).
For a topos of the form $\Top(X)$, one recovers a
bijection between the set of constructible subsets (\resp\ldots)
of $X$, and the set of constructible subtopoi (\resp\ldots) of
$\Top(X)$.
\end{enonce*}

\subsection{The gluing theorem}\etiquette[9.5]{IV.subsec9.5}

\subsubsection{}\etiquette[9.5.1]{IV.subsubsec9.5.1}
Let $f:A\rightarrow B$ be a functor between two categories. Denote
by $(B,A,f)$ the following category: the objects of $(B,A,f)$ are
the triples
$$
(X,Y,u)
$$
where $X$ is an object of $B$, $Y$ an object of $A$, and $u$ a morphism
from $X$ to $f(Y)$; the morphisms between two objects $(X,Y,u)$ and
$(X',Y',u')$ are the pairs $(m,m')$, where $m$ is a morphism from
$X$ to $X'$ and $m'$ a morphism from $Y$ to $Y'$, such that the
following diagram is commutative:
\[
\xymatrix@C=1.5cm@R=1cm{
X \ar[r]^{u} \ar[d]_{m} & **[r]f(Y) \ar[d]^{f(m')}\\
X' \ar[r]^{u'}& **[r]f(Y') }
\]
the\pageoriginale composition of morphisms is defined in the
evident way.

\subsubsection{}\etiquette[9.5.2]{IV.subsubsec9.5.2}
Note that the category $(B,A,f)$ depends functorially
on the system of data $A$, $B$, $f$, in a sense whose precise
formulation we leave to the reader. In particular, if $f'$ is a
second functor from $A$ to $B$, then every morphism of functors
$$
f\xrightarrow{u} f'
$$
defines a functor $(B,A,f)\rightarrow (B,A,f')$, which is an
isomorphism if $u:f\rightarrow f'$ is one.

\subsubsection{}\etiquette[9.5.3]{IV.subsubsec9.5.3}
Return to the situation of \hyperref[IV.subsec9.3]{9.3}, and consider the
category $(\mathcal{F},\mathcal{U}, \rho)$ defined in
\hyperref[IV.subsubsec9.5.1]{9.5.1}. To each object $X$ of $E$ there corresponds an
object $(i^{\ast}X, j^{\ast}X, h(X):i^{\ast}X\rightarrow \rho
j^{\ast}X)$ of $(\mathcal{F}, \mathcal{U},\rho)$, where $h(X)$ is obtained
by applying the functor $i^{\ast}$ to the adjunction morphism
$X\rightarrow j_{\ast}j^{\ast}(X)$ (one has $\rho=i^{\ast}j_{\ast}$
\sisi{(9.10.)}{\hyperref[IV.eq9.3.6.3]{9.3.6.3}}). Thus this defines a functor
\begin{equation*}
\Phi:E\longrightarrow (\mathcal{F}, \mathcal{U},
\rho){\quoi}.\tag{9.5.3.1}
\end{equation*}\etiquette[9.5.3.1]{IV.eq9.5.3.1}

\begin{enonce*}{Theorem 9.5.4}\etiquette[9.5.4]{IV.theo9.5.4}
\begin{enumerate}
\renewcommand{\theenumi}{\alph{enumi}}
\renewcommand{\labelenumi}{\theenumi)}
\item Let $E$ be a topos, let $\mathcal{U}$ be an open subtopos of $E$, and let
$\mathcal{F}$ be the complementary closed subtopos. The functor
$$
\rho=i^{\ast}j_{\ast}:\mathcal{U} \longrightarrow \mathcal{F}
$$
is left exact and accessible, and the functor $\Phi$ $\hyperref[IV.eq9.5.3.1]{(9.5.3.1)}$
is an equivalence of categories.

\item Conversely, let $\mathcal{U}$ and $\mathcal{F}$ be two topoi, and let
$\rho:\mathcal{U}\rightarrow \mathcal{F}$ be a gluing functor
$(\hyperref[IV.subsubsec9.3.6]{9.3.6})$ (\ie accessible {\rm (\hyperrefext[I][I.def9.2]{9.2})}
and left exact). The category $E=(\mathcal{F},\mathcal{U},\rho)$
is a topos. Let $\phi_{\mathcal{F}}$ be an initial object of
$\mathcal{F}$, let $e_{\mathcal{U}}$ be a terminal object of $\mathcal{U}$, and
let $U=(\phi_{\mathcal{F}},e_{\mathcal{U}},
\phi_{\mathcal{F}}\rightarrow \rho(e_{\mathcal{U}}))\in \Ob E$. The
functor $X\mapsto (\phi_{\mathcal{F}},
X,\phi_{\mathcal{F}}\rightarrow \rho(X))$ induces an equivalence of
$\mathcal{U}$ with $E_{/U}$, hence with the open subtopos
$\mathcal{U}'$ of $E$ defined by $U$. The functor $Y\mapsto
(Y,e,Y\rightarrow \rho(e))$ is a closed embedding
$(\hyperref[IV.subsubsec9.3.5]{9.3.5})$, \ie it induces an equivalence of
$\mathcal{F}$ with a closed subtopos $\mathcal{F}'$ of $E$. The
subtopoi $\mathcal{U}'$ and $\mathcal{F}'$\pageoriginale are
complementary. Let $\rho':\mathcal{U}'\rightarrow \mathcal{F}'$
be the gluing functor. The functors $\rho$ and $\rho'$ are
compatible with the equivalences described above.
\end{enumerate}
\end{enonce*}

\begin{enonce*}{Definition 9.5.5}\etiquette[9.5.5]{IV.def9.5.5}
Let $\mathcal{U}$ and $\mathcal{F}$ be two topoi, and let
$\rho:\mathcal{U}\rightarrow \mathcal{F}$ be a gluing functor
$(\hyperref[IV.subsubsec9.3.6]{9.3.6})$. The topos $(\mathcal{F},
\mathcal{U},\rho)$ is called the topos constructed from
$\mathcal{U}$ and $\mathcal{F}$ by gluing by means of the
functor $\rho$.
\end{enonce*}

\setcounter{subsubsection}{5}
\subsubsection{}\etiquette[9.5.6]{IV.subsubsec9.5.6}
\emph{Proof of} \hyperref[IV.theo9.5.4]{9.5.4} a) \footnote[1]{For a
gluing statement more general than \hyperref[IV.def9.5.5]{9.5.5}, \cf
Exercise\stepcounter{cnote}\textsuperscript{\arabic{cnote}}%
\textup{ [N.D.E.: This exercise does not seem to exist.]}.} The first assertion merely recalls
\hyperref[IV.prop9.4.1]{9.4.1} d). The pair of functors $(i^{\ast},j^{\ast})$
is conservative, or equivalently, faithful (\hyperref[IV.prop9.4.1]{9.4.1} c)). A
fortiori, the functor $\Phi$ is faithful. Let us show that it is
fully faithful. It follows from \hyperref[IV.prop9.3.3]{9.3.3} b) that, for
every object $X$ of $E$, the following diagram is cartesian:
\[\vcenter{
\xymatrix@C=2.5cm{
X  \ar[r] \ar[d]               & j_{\ast}j^{\ast}X \ar[d]\\
i_{*}i^{*}X   \ar[r]     & i_{\ast}i^{\ast}j_{\ast}j^{\ast}X ;
}}\leqno{(*)}
\]
Let $X$ and $Y$ be two objects of $E$, and suppose that $(m,m')$:
\[\vcenter{
\xymatrix@C=2.75cm@R=1cm{
 i^{\ast}X  \ar[r] \ar[d]_{m}               & i^{\ast}j_{\ast}j^{\ast}X \ar[d]_{i^{\ast}j_{\ast}m'}\\
i^{*}Y   \ar[r]     & i^{\ast}j_{\ast}j^{\ast}Y ; }}
\]
is a morphism from $\Phi(X)$ to $\Phi(Y)$. Using the
functoriality of the fiber product and the diagrams $(\ast)$, one
obtains a morphism $w:X\rightarrow Y$ which makes the
following diagrams commutative:
\begin{displaymath}
\xymatrix@C=.5cm@R=1.25cm{
X \ar@{-}[rr]^{w} \ar[d]  &             &Y \ar[d]\\
 j_{\ast}j^{\ast}X \ar[rr]^{j_{\ast}m'}       &   &
j_{\ast}j^{\ast}Y , }\qquad\qquad \xymatrix@C=.5cm@R=1.25cm{
X \ar@{-}[rr]^{w} \ar[d]& &Y \ar[d]\\
 i_{\ast}i^{\ast}X \ar[rr]^{i_{\ast}m}&   & i_{\ast}i^{\ast}Y
}
\end{displaymath}
By\pageoriginale applying the functors $j^{\ast}$
and $i^{\ast}$ respectively, one obtains $j^{\ast}w=m'$ and $i^{\ast}w=m$. Let us
now show that the functor $\Phi$ is essentially surjective.
Let $(Y,X,u:Y\rightarrow i^{\ast}j_{\ast}(X))$ be an object of
$(\mathcal{F},\mathcal{U},i^{\ast} j_{\ast})$. From it one obtains a
diagram:
\[
\xymatrix@C=2cm{
                             & j_{\ast}X \ar[d]\\
i_{\ast}Y \ar[r]^{i_{\ast}u} & i_{\ast}i^{\ast}j_{\ast}X }
\]
Hence, on taking the fiber product, one obtains a cartesian diagram:
\[\vcenter{
  \xymatrix@C=2.75cm@R=1.25cm{
W \ar[r] \ar[d]& j_{\ast}X \ar[d]\\
i_{\ast}Y \ar[r]^{i_{\ast}u}& i_{\ast}i^{\ast}j_{\ast} X\\
}}\leqno{(**)}
\]

We leave it to the reader to show that $\Phi(W)$ is
isomorphic to $(Y,X,u:Y\rightarrow i^{\ast}j_{\ast}(X))$ and that the
diagram (${\ast}{\ast}$) is isomorphic to the diagram ($\ast$)
associated with $W$.

\hypertarget{sga4.exp.IV.subsubsection.9.5.7}{}\subsubsection{}

% END INCLUDED SGA4_English_current_buildable_source_20260906/source/04/drafts/SGA4_IV_LINES_6788_7070_English.texfrag


% BEGIN INCLUDED SGA4_English_current_buildable_source_20260906/source/04/drafts/SGA4_IV_LINES_7071_7243_English.texfrag
\emph{Proof} of \hyperref[IV.theo9.5.4]{9.5.4}
\sisi{1)}{b)}. We shall confine ourselves to showing that
$(\mathcal{F},\mathcal{U},\rho)$ is a topos. The proof of the other
assertions is left to the reader. We verify properties a), b), c),
d) of \hyperref[IV.subsubsec1.1.2]{1.1.2}.
\begin{enumerate}
\renewcommand{\theenumi}{\alph{enumi}}
\renewcommand{\labelenumi}{\theenumi)}
\item \emph{Finite projective limits are representable}: This is
clear because the categories $\mathcal{U}$ and $\mathcal{F}$ are
topoi and $\rho$ commutes with finite projective limits.

\item \emph{Direct sums are representable, disjoint, and universal}:
Let $(Y_i\xrightarrow{u_i}\rho(X_i))$, $i\in I$, be a family of
objects of $(\mathcal{F},\mathcal{U},\rho)$. By the definition of
inductive limits, there is a canonical morphism
$$
\coprod_{i\in I}\rho(X_i)\longrightarrow\rho(\coprod_{i\in I}X_i),
$$
and hence,\pageoriginale by composing with $\coprod u_i$, a morphism
$$
n:\coprod_{i\in I}Y_i\longrightarrow\rho(\coprod_{i\in I}X_i){\quoi}.
$$
One verifies immediately that this last object is the direct sum of
the given family and that this direct sum is disjoint and universal.

\item \emph{Equivalence relations are effective and universal}:

Let
\[
\xymatrix@C=2cm@R=1.5cm{
R_1 \ar@<3pt>[r] \ar@<-3pt>[r] \ar[d]& Y \ar[d]^{w}\\
\rho(R_2) \ar@<3pt>[r]^{f(u)} \ar@<-3pt>[r]_{f(v)} & \rho(X) }
\]
be an equivalence relation on the object
$Y\xrightarrow{w}\rho(X)$. Then
$\xymatrix@C=.6cm{R_1\ar@<-2pt>[r]\ar@<2pt>[r]&Y}$ is an equivalence
relation on $Y$ in $\mathcal{F}$, and
$\xymatrix{R_2\ar@<2pt>[r]^{u}\ar@<-2pt>[r]_{v}&X\\}$ is an
equivalence relation on $X$ in $\mathcal{U}$. The quotients
$Y/R_1$, $\rho(X)/\rho(R_2)$, and $X/R_2$ are representable in
$\mathcal{F}$ and $\mathcal{U}$ because these categories are topoi.
Moreover, by the definition of inductive limits, the canonical
morphism $\rho(X)\rightarrow\rho(X/R_2)$ factors through
$\rho(X)/\rho(R_2)$. By functoriality of inductive limits, one obtains
a canonical morphism
$Y/R_1\xrightarrow{t}\rho(X/R_2)$. One verifies that this last object
is the quotient of $Y\xrightarrow{w}\rho(X)$ by the equivalence
relation in question, that this quotient is effective (I~10), and
that all these properties are preserved under base change.

\item $(\mathcal{F},\mathcal{U},\rho)$ \emph{admits a small
generating family}. This property follows from
\sisi{(1.9.2 5)}{\hyperrefext[I][I.cor9.25]{9.25}} (taking into
account that $\rho$ is accessible), by taking
$B=\Delta_1=(0\xrightarrow{f}1)$, $E_0=\mathcal{F}$,
$E_1=\mathcal{U}$, and $f^{\ast}=\rho:E_1\rightarrow E_0$, whence
$\Hhom_B(B,E)=(\mathcal{F},\mathcal{U},\rho)$.

One may also avoid using
\sisi{(1.9.2 5)}{\hyperrefext[I][I.cor9.25]{9.25}} (whose proof as
given was rather arduous!) and use the hypothesis on $\rho$ in the
form that $\rho$ admits a pro-adjoint
$$
\sigma:\Pro(\mathcal{F})\longrightarrow\Pro(\mathcal{U}){\quoi},
$$
(this\pageoriginale interpretation (\hyperrefext[I][I.sec8]{8})
being based on \hyperrefext[I][I.sec9]{9}, whose proof is distinctly
more comprehensible). It suffices to note that if $C'$ (\resp $C''$)
is a generating subcategory of $\mathcal{U}$ (\resp $\mathcal{F}$),
and if, for every $X''\in\Ob C''$, one represents $\sigma(X'')$ as a
projective system $(\sigma(X'')_i)_{i\in I(X'')}$, where $I(X'')$ is
a small filtered ordered set, then the family of objects of $E$ which
are either of the form
$(\phi_{\mathcal{F}},X',\phi_{\mathcal{F}}\rightarrow\rho(X'))$ with
$X'\in\Ob C'$, or of the form
$(X'',\sigma(X'')_i,u_{X''}:X''\rightarrow
\rho(\sigma(X'')_i))$, where $X''\in\Ob C''$, $i\in I(X'')$, and
$u_{X''}$ corresponds under adjunction to the canonical morphism
$\sigma(X'')\rightarrow\sigma(X'')_i$, is a generating family
(plainly a small one) when $C'$ and $C''$ are chosen small. The
verification is indeed immediate and is left to the reader.
\end{enumerate}

\setcounter{subsubsection}{4}\setcounter{paragraph}{7}
\paragraph{}\etiquette[9.5.4.8]{IV.par9.5.4.8}
Use the notation of \hyperlink{sga4.exp.IV.subsubsection.9.5.7}{9.5.7} b), and let us explain how, up to canonical
isomorphisms of functors, the system of functors
\eqhyperref[IV.eq9.3.6.2]{9.3.6.2} can be made explicit when
$E=(\mathcal{F},\mathcal{U},\rho)$. The details verifying the
assertions below (which are essentially mechanical applications of
\hyperref[IV.theo9.5.4]{9.5.4}) are left to the reader.

\setcounter{subsubsection}{8}\setcounter{paragraph}{0}
\paragraph{}\etiquette[9.5.8.1]{IV.par9.5.8.1}
The functor
$$
i^{\ast}:(\mathcal{F},\mathcal{U},\rho)\longrightarrow\mathcal{F}
$$
is given by
$$
i^{\ast}(X,Y,u:X\longrightarrow\rho(Y))=X{\quoi}.
$$

\paragraph{}\etiquette[9.5.8.2]{IV.par9.5.8.2}
The functor
$$
i_{\ast}:\mathcal{F}\longrightarrow(\mathcal{F},\mathcal{U},\rho)
$$
is given by $i_{\ast}(X)=(X,e,X\rightarrow\rho(e))$, where $e$ is
the terminal object of $\mathcal{U}$.

\paragraph{}\etiquette[9.5.8.3]{IV.par9.5.8.3}
The adjunction morphism
$$
\id\longrightarrow i_{\ast}i^{\ast}
$$
is isomorphic to the functorial morphism
\[
\xymatrix@C=2cm{
X \ar[r]^{u} \ar[d]_{id}& \rho(Y) \ar[d]\\
X \ar[r]                & \rho(e) }
\]

\paragraph{}\etiquette[9.5.8.4]{IV.par9.5.8.4}
The\pageoriginale functor
$$
j^{\ast}:(\mathcal{F},\mathcal{U},\rho)\longrightarrow\mathcal{U}
$$
is given by\qquad $j^{\ast}(X,Y,u:X\longrightarrow\rho(Y))=Y$.

\paragraph{}\etiquette[9.5.8.5]{IV.par9.5.8.5}
The functor
$$
j_{!}:\mathcal{U}\longrightarrow(\mathcal{F},\mathcal{U},\rho)
$$
is given by\qquad
$j_{!}(Y)=(\phi_{\mathcal{F}},Y,
\phi_{\mathcal{F}}\longrightarrow\rho(Y))$

\noindent (where $\phi_{\mathcal{F}}$ is the initial object of
$\mathcal{F}$).

\paragraph{}\etiquette[9.5.8.6]{IV.par9.5.8.6}
The functor
$$
j_{\ast}:\mathcal{U}\longrightarrow(\mathcal{F},\mathcal{U},\rho)
$$
is given by\qquad
$j_{\ast}(Y)=(\rho(Y),Y,\id:\rho(Y)\longrightarrow\rho(Y))$.

\paragraph{}\etiquette[9.5.8.7]{IV.par9.5.8.7}
The adjunction morphism
$$
j_{!}j^{\ast}\longrightarrow\id
$$
is isomorphic to the functorial morphism
\[
\xymatrix{
\phi_{\Ff} \ar[r] \ar[d] & \rho(Y) \ar[d]^{\rho(\id)}\\
X \ar[r]^{u}                 & \rho(Y)     &. }
\]

\paragraph{}\etiquette[9.5.8.8]{IV.par9.5.8.8}
The adjunction morphism
$$
\id\longrightarrow j_{\ast}j^{\ast}
$$
is isomorphic to the functorial morphism
\[
\xymatrix{
X  \ar[r]^{u} \ar[d]_{u}& \rho(Y) \ar[d]^{\rho(\id)}\\
\rho(Y) \ar[r]^{\id}& \rho(Y) }
\]

% END INCLUDED SGA4_English_current_buildable_source_20260906/source/04/drafts/SGA4_IV_LINES_7071_7243_English.texfrag


% BEGIN INCLUDED SGA4_English_current_buildable_source_20260906/source/04/drafts/SGA4_IV_LINES_7244_7443_English.texfrag
\begin{enonce*}[remark]{Exercise 9.5.9}\etiquette[9.5.9]{IV.exer9.5.9}
(\emph{Exactness properties of a glued topos.}) Let
$\mathcal{U}$, $\mathcal{F}$ be two topoi, let
$\rho:\mathcal{U}\rightarrow \mathcal{F}$ be a gluing functor, let
$E=(\mathcal{F}, \mathcal{U},\rho)$ be the glued topos, and let
$j:\mathcal{U}\rightarrow E$ and $i:\mathcal{F}\rightarrow E$ be the
canonical morphisms of topoi, so that
$\rho=i^{\ast}j_{\ast}$.
\begin{enumerate}
\renewcommand{\theenumi}{\alph{enumi}}
\renewcommand{\labelenumi}{\theenumi)}
\item Prove that $i^{\ast}$ commutes with (small) products (\ie (\hyperref[IV.rem1.8]{1.8}
or \sisi{\hyperref[IV.rem1.8]{1.8}}{\hyperrefext[I][I.sec8]{8}}) the functor $i^{\ast}$ admits a left adjoint
$i_{!}$) if and only if the same holds for $\rho$, \ie
(\loccit) if and only if $\rho$ admits a left adjoint
$\sigma:\mathcal{F}\rightarrow \mathcal{U}$ (or equivalently if and
only if the pro-adjoint of $\rho$, restricted\pageoriginale to
$\mathcal{F}$ \eqhyperref[IV.eq9.6.3]{9.6.3}, factors, up to isomorphism, through
$\mathcal{U}$): one then has $\sigma=j^{\ast}i_{!}$. Thus the
gluing functors $\rho$ having this property correspond,
up to isomorphism, to the functors $\sigma:\mathcal{F}\rightarrow
\mathcal{U}$ which have a right adjoint, \ie (\hyperref[IV.cor1.6]{1.6}) which
commute with $𝒰$-small inductive limits, or equivalently which are
continuous. When $\mathcal{F}$ is equivalent to a topos
$C^{\sim}$, where $C$ is a $𝒰$-site, these functors
therefore correspond, up to isomorphism, to continuous functors
$C\rightarrow \mathcal{U}$ (\hyperrefext[III][III.prop1.2]{1.2} and \hyperrefext[III][III.prop1.7]{1.7}).

\item Suppose that the functor $i_{!}$ is defined. Show that
$i_{!}$ commutes with a specified type of small projective limits
if and only if $\sigma:\mathcal{F}\rightarrow
\mathcal{U}$ commutes with them. (For necessity, use the fact that $j^{\ast}$
commutes with small projective limits; for sufficiency,
use the equivalent description of $E$ in terms of $\sigma$
as consisting of the triples $(Y,X,u)$ with $Y\in \Ob \mathcal{F}$, $X\in
\Ob \mathcal{U}$, $u:\sigma (Y)\rightarrow X$.)

\item Deduce from a) and b) an example of a closed embedding of topoi
$i:\mathcal{F}\rightarrow E$ such that $i_{!}$ exists, but commutes
neither with fiber products nor with filtered projective limits. (It
suffices to find a continuous functor between two topoi which
has neither of these two exactness properties.)
\end{enumerate}
\end{enonce*}

\begin{enonce*}[remark]{Exercise 9.5.10}\etiquette[9.5.10]{IV.exer9.5.10}
(\emph{Gluing of topoi of the form $\widehat{C}$.}) Let $C$
be a small category, and let $E=\widehat{C}$, which is a
$𝒰$-topos.
\begin{enumerate}
\renewcommand{\theenumi}{\alph{enumi}}
\renewcommand{\labelenumi}{\theenumi)}
\item Recall that the opens $U$ of $E$ correspond to
\emph{sieves} $C'$ of $C$ \sisi{(1.4.1, \hyperref[IV.subsec4.2]{4.2})}{(\hyperrefext[I][I.def4.1]{4.1}, \hyperref[IV.subsubsec8.4.4]{8.4.4})}, with $E_{/U}$
then being equivalent to $\widehat{C'}$, and the functor
$j^{\ast}:E\rightarrow E_{/U}$ identifying with the restriction functor
\hyperref[IV.subsec5.4]{5.4}, \ie the inclusion morphism $j:E_{/U}\rightarrow E$ being
$\widehat{f}:\widehat{C'}\rightarrow \widehat{C}$
(\hyperref[IV.subsubsec4.6.1]{4.6.1}), where $f:C'\rightarrow C$ is the inclusion.
Let $C''$ be the full subcategory of $C$ complementary to $C'$
(\ie $\Ob C''$ is the complement of $\Ob C'$ in $\Ob C$),
and let $g:C''\rightarrow C$ be the inclusion. Then the closed subtopos
$\mathcal{F}$ of $E=\widehat{C}$, complementary to the open subtopos
defined by $C'$, is canonically equivalent to
$\widehat{C''}$, with the inclusion morphism $i:\mathcal{F}\rightarrow
E$ identifying with $\widehat{g}:\widehat{C''}\rightarrow
\widehat{C}$.

\item Consider\pageoriginale the functor
\begin{equation*}
h:{C'}^0\times C''\longrightarrow (\Ens),\quad
h(X',X'')=\Hom_C(X',X''){\quoi},\tag{9.5.10.1}
\end{equation*}\etiquette[9.5.10.1]{IV.eq9.5.10.1}
and note that $C$, together with its full subcategory $C'$, can be
reconstructed up to isomorphism from the categories $C'$,
$C''$ and the bifunctor $h$ (the latter may be chosen
arbitrarily). (Compare \hyperref[IV.subsec9.7.1]{9.7.1} below.) The datum
of $h$ is equivalent to that of a functor $C''\rightarrow
\widehat{C'}=\Hhom({C'}^0,(\Ens))$, or equivalently (\hyperrefext[III][III.prop1.2]{1.2} and
\hyperrefext[III][III.prop1.7]{1.7}) to that of a functor
\begin{equation*}
\sigma:\widehat{C''}\longrightarrow \widehat{C'}\tag{9.5.10.2}
\end{equation*}\etiquette[9.5.10.2]{IV.eq9.5.10.2}
which is continuous, \ie (\hyperref[IV.cor1.6]{1.6}) commutes with small inductive
limits, or equivalently, admits a right adjoint
\begin{equation*}
\rho:\widehat{C}'\longrightarrow \widehat{C''}{\quoi}.\tag{9.5.10.3}
\end{equation*}\etiquette[9.5.10.3]{IV.eq9.5.10.3}
Show that $\rho$ is precisely the gluing functor
associated with the open subtopos $\widehat{C}'$ of $E=\widehat{C}$.

\item Deduce from a) and b) that if one is given a
gluing functor \eqhyperref[IV.eq9.5.10.3]{9.5.10.3} between two presheaf categories,
and hence a glued topos $E$, the following conditions are
equivalent:
\begin{itemize}
\item[(i)] $E$ is equivalent to a topos of the form $\widehat{C}$.

\item[(ii)] The functor $\rho$ commutes with small products (or equivalently, it
admits a left adjoint \eqhyperref[IV.eq9.5.10.2]{9.5.10.2}).

\item[(iii)] The inclusion morphism $i:\widehat{C''}\rightarrow E$ is ``essential,'' \ie $i^{\ast}$ commutes with products, or equivalently
$i^{\ast}$ admits a left adjoint $i_{!}$.
\end{itemize}
When this is so, describe explicitly a category $C$ such that
$E\approx \widehat{C}$ using the bifunctor \eqhyperref[IV.eq9.5.10.1]{9.5.10.1}
defined by the restriction of \eqhyperref[IV.eq9.5.10.2]{9.5.10.2} to $C''$.

\item Show that a topos obtained by gluing two punctual topoi
need not be equivalent to a topos of the form
$\widehat{C}$.
\end{enumerate}
\end{enonce*}

\begin{enonce*}[remark]{Exercise 9.5.11}\etiquette[9.5.11]{IV.exer9.5.11}
(Morphism from a glued topos to a topos.)
\begin{enumerate}
\renewcommand{\theenumi}{\alph{enumi}}
\renewcommand{\labelenumi}{\theenumi)}
\item Let $\rho \sisi{}{: \mathcal{U}} \rightarrow \mathcal{F}$ be a functor between two
categories, let $E=(\mathcal{F},\mathcal{U},\rho)$\pageoriginale be the
``glued'' category of \hyperref[IV.subsubsec9.5.1]{9.5.1}, and let $F$ be an
arbitrary category. Consider the functor $f\mapsto \rho\circ f$:
$$
\bar{\rho}:\bar{\mathcal{U}}=\Hhom(F,\mathcal{U})\longrightarrow
\bar{\mathcal{F}}=\Hhom(F,\mathcal{F}){\quoi},
$$
and define an equivalence of categories
$$
\Hhom(F,E)\xrightarrow{\approx}(\bar{\mathcal{F}},\bar{\mathcal{U}},\bar{\rho}){\quoi}.
$$

\item Let $f:F\rightarrow E$ be a functor, and consider its composites
with the canonical functors $j^{\ast}:E\rightarrow \mathcal{U}$ and
$i^{\ast}:E\rightarrow \mathcal{F}$. Show that $f$ commutes with a
specified type of inductive limits (assumed to be representable in
$\mathcal{U}$ and $\mathcal{F}$, and hence in $E$) if and only if
$j^{\ast}f$ and $i^{\ast}f$ commute with them. Show that if $\rho$
commutes with a specified type of projective limits (which is then
representable in $E$), then $f$ commutes with them if and only if
the same holds for $j^{\ast}f$ and $i^{\ast}f$.

\item Now suppose that $\mathcal{U}$ and $\mathcal{F}$ are
topoi, and that
$\rho$ is a gluing functor. Let $F$ be a topos and
$f:F\rightarrow E$ a functor. Conclude from b) that $f$ is an
inverse-image functor associated with a morphism of topoi $E\rightarrow
F$ if and only if the same holds for the functors
$j^{\ast}f:F\rightarrow \mathcal{U}$ and $i^{\ast}f:F\rightarrow
\mathcal{F}$. Conclude from this and from a) that there is an equivalence
between the category $\Hhomtop(E,F)$ and the category of triples
$(f',f'',u)$, where $f'$ (\resp $f''$) is an object of
$\Hhomtop(\mathcal{U}, F)$ (\resp $\Hhomtop(\mathcal{F}, F)$), and where
$u:{f''}^{\ast}\rightarrow \rho\circ {f'}^{\ast}$ is a
homomorphism of functors $F\rightarrow \mathcal{F}$: if $f$ is an
object of $\Hhomtop(E,F)$ and $(f',f'',u)$ is the corresponding triple,
then $f'=f\circ j$, $f''=f\circ i$.

\item Specify in what sense c) may be said to provide a
characterization, up to equivalence (in a 2-category of topoi),
of the topos $E$ in terms of the triple $(\mathcal{U},\mathcal{F},\rho)$.
\end{enumerate}
\end{enonce*}

\begin{enonce*}[remark]{Exercise 9.5.12}\etiquette[9.5.12]{IV.exer9.5.12}
(\emph{Gluing of two topological spaces.})
\begin{enumerate}
\renewcommand{\theenumi}{\alph{enumi}}
\renewcommand{\labelenumi}{\theenumi)}
\item Let $X$ be a topological space, let $\rho:\Top(X)\rightarrow
(\Ens)=\Top(pt)$ be a gluing functor, \ie (\hyperrefext[I][I.sec8]{8}) a
pro-representable functor, let $A\in \Ob \Pro(\Top(X))$ be the object which
pro-represents it, and let $E$ be the topos obtained from $\rho$ by gluing.
Show that $E$ is equivalent to a topos of the form $\Top(Y)$ if
and only if $A$ is isomorphic to an object of $\Pro(\Ouv(X))$.
Conclude\pageoriginale that if $X$ is nonempty, one can find
$\rho$ such that $E$ is not equivalent to a topos of the form
$\Top(Y)$.

\item More generally, consider a gluing functor
$\rho:\Top(X)\rightarrow \Top(Y)$, where $X$, $Y$ are two topological
spaces. Deduce from it a map
$$
\phi:Y\longrightarrow \Ob \Pro(\Top(X))
$$
(the composite of $Y\rightarrow \Ob \Ppoint(\Top(Y))\rightarrow
\Ob\Pro(\Top(Y))$ (\hyperref[IV.subsubsec6.8.5]{6.8.5}) and the pro-adjoint
$\sigma$ of $\rho$). (N.B. for $y\in Y$, $\phi(y)$ identifies with the
\emph{fiber at $y$ of the gluing sheaf} \eqhyperref[IV.eq9.6.4]{9.6.4}
\emph{defined by} $\rho$.) Show that if $E$ is equivalent to a
topos of the form $\Top(Z)$, then $\phi$ takes its values in
the essential image of $\Pro(\Ouv(X))$.
\end{enumerate}
\end{enonce*}

\begin{enonce*}[remark]{Problem}
Is the converse true?
\end{enonce*}

% END INCLUDED SGA4_English_current_buildable_source_20260906/source/04/drafts/SGA4_IV_LINES_7244_7443_English.texfrag


% BEGIN INCLUDED SGA4_English_current_buildable_source_20260906/source/04/drafts/SGA4_IV_LINES_7444_7700_English.texfrag
\subsection{Gluing sheaf}\etiquette[9.6]{IV.subsec9.6}

We saw in \hyperref[IV.theo9.5.4]{9.5.4} that the datum of a topos $E$
equipped with an open $U$ amounts essentially to that of two topoi
$\mathcal{U}$ and $\mathcal{F}$ and a gluing functor
\begin{equation*}
\rho:\mathcal{U}\longrightarrow \mathcal{F}\tag{9.6.1}
\end{equation*}\etiquette[9.6.1]{IV.eq9.6.1}
\ie a functor admitting a pro-adjoint
\begin{equation*}
\sigma:\Pro(\mathcal{F})\longrightarrow
\Pro(\mathcal{U}){\quoi}.\tag{9.6.2}
\end{equation*}\etiquette[9.6.2]{IV.eq9.6.2}
By the sorites on pro-adjoint functors (\hyperrefext[I][I.sec8]{8}), the
functor $\rho$ is determined up to unique isomorphism (and consequently
the glued topos $E=(\mathcal{F},\mathcal{U},\rho)$ is determined up to
equivalence of topoi) once the functor $\sigma$ is known, or,
equivalently, once the functor induced by $\sigma$ is known:
\begin{equation*}
\sigma_{\circ}:\mathcal{F}\longrightarrow
\Pro(\mathcal{U}){\quoi}.\tag{9.6.3}
\end{equation*}\etiquette[9.6.3]{IV.eq9.6.3}
The functor $\sigma$ is determined, up to unique isomorphism, by the
property of commuting with small filtered projective limits and of
extending the functor $\sigma_\circ$. On the other hand, for a given
functor \eqhyperref[IV.eq9.6.3]{9.6.3} to be isomorphic to a pro-adjoint, it
is plainly necessary and sufficient that, for every object $X$ of
$\mathcal{U}$, the contravariant functor
$$
Y\longmapsto \Hom_{\Pro(\mathcal{U})}(\sigma_{\circ}(Y),X)
$$
on\pageoriginale $\mathcal{F}$ be representable, or equivalently be a
sheaf on $\mathcal{F}$ for the canonical topology
(\sisi{\hyperref[IV.subsubsec1.1.2]{1.1.2} (iii)}{\hyperref[IV.theo1.2]{1.2} iii)}). It follows at once that this
also means that the opposite functor
\begin{equation*}
\sigma_{\circ}^{\circ}:\mathcal{F}^{\circ}\longrightarrow
\Pro(\mathcal{U})^{\circ}\subset
\check{\mathcal{U}}=\Hhom(\mathcal{U},(\Ens))\tag{9.6.4}
\end{equation*}\etiquette[9.6.4]{IV.eq9.6.4}
is a \emph{sheaf} on $\mathcal{F}$ with values in
$\Pro(\mathcal{U})^{\circ}$ (\hyperrefext[II][II.def6.1]{6.1}). Thus, \emph{the
datum of a gluing functor} \eqhyperref[IV.eq9.6.1]{9.6.1} \emph{is also
equivalent} (up to isomorphism) \emph{to that of a sheaf on
$\mathcal{F}$ with values in} $\Pro(\mathcal{U})^{\circ}$. The sheaf
thus associated with a gluing functor (or with a situation
$(E,U,\mathcal{F})$ as in (\hyperref[IV.subsec9.3]{9.3})) is called the
\emph{gluing sheaf} associated with the gluing functor $\rho$ under
consideration (\resp with the open $U$ of the topos $E$).

\begin{enonce*}[remark]{Exercise 9.6.5}\etiquette[9.6.5]{IV.exer9.6.5}
Refine \hyperref[IV.theo9.5.4]{9.5.4} into a statement of 2-equivalence of
2-categories between the 2-category consisting of pairs $(E,U)$, where
$E$ is a topos and $U$ an open of the latter, and the 2-category consisting
of triples $(\mathcal{U},\mathcal{F},\rho)$, where $\mathcal{U}$ and
$\mathcal{F}$ are topoi and
$\rho:\mathcal{U}\rightarrow \mathcal{F}$ is a gluing functor (the
explicit description of the 2-categories in question being left to
the reader). Give a variant of this statement involving the
2-category of triples $(\mathcal{U},\mathcal{F}, \phi)$, where $\phi$
is now a sheaf on $\mathcal{F}$ with values in
$\Pro(\mathcal{U})^{\circ}$.
\end{enonce*}

\subsection{Points of a glued topos}\etiquette[9.7]{IV.subsec9.7}

\subsubsection{}\etiquette[9.7.1]{IV.subsec9.7.1}
Let $C$ be a category, and let
$$
u:D'\longrightarrow C\quav v:D''\longrightarrow C
$$
be two fully faithful functors such that the essential images $C'$ and
$C''$ of $u$ and $v$ satisfy
$$
\Ob C'\cap \Ob C''=\phi\quav \Ob C'\cup \Ob C''=\Ob C{\quoi}.
$$
Suppose, in addition, that
$$
X'\in \Ob C'\quav X''\in \Ob C''\Rightarrow
\Hom(X'',X')=\phi{\quoi},
$$
or, equivalently, that
$$
\Hom(v(A''), u(A'))=\phi\mbox{ for } A'\in \Ob D', A''\in \Ob
D''{\quoi}.
$$
It is then immediate that the category $C$ can be reconstructed, up
to canonical\pageoriginale isomorphism, from the categories $C'$ and
$C''$ and the functor
$$
(X',X'')\longmapsto \Hom(X',X''):{C'}^{\circ}\times
C''\longrightarrow (\Ens){\quoi};
$$
similarly, $C$ can be reconstructed up to equivalence (this
equivalence being defined up to unique isomorphism) from $D'$, $D''$,
and the functor
$$
(A',A'')\longmapsto \Hom(u(A'), v(A'')):{D'}^{\circ}\times
D''\longrightarrow (\Ens){\quoi}.
$$
This observation applies in particular to the situation described in
the following proposition, and shows that the category
$\Ppoint(E)$ is determined, up to equivalence, by the categories
$\Ppoint(\mathcal{U})$ and $\Ppoint(\mathcal{F})$ and by a certain
functor
$$
\Ppoint(\mathcal{U})^{\circ}\times
\Ppoint(\mathcal{F})\longrightarrow (\Ens)
$$
derived from the gluing functor $\rho:\mathcal{U}\rightarrow
\mathcal{F}$.

\begin{enonce*}{Proposition 9.7.2}\etiquette[9.7.2]{IV.prop9.7.2}
Let $E$ be a topos, let $U$ be an open of $E$, let
$\mathcal{U}=E_{/U}$, and let $\mathcal{F}$ be the closed subtopos of
$E$ complementary to $U$. Consider the functors $p\mapsto p\circ j$
and $q\mapsto q\circ i$ induced by the inclusion morphisms
$j:\mathcal{U}\rightarrow E$ and $i:\mathcal{F} \rightarrow E$:
\begin{equation*}
u:\Ppoint(\mathcal{U})\longrightarrow \Ppoint(E),\quad
v:\Ppoint(\mathcal{F})\longrightarrow
\Ppoint(E){\quoi}.\tag{9.7.2.1}
\end{equation*}\etiquette[9.7.2.1]{IV.eq9.7.2.1}
Then the following assertions hold:
\begin{enumerate}
\renewcommand{\theenumi}{\alph{enumi}}
\renewcommand{\labelenumi}{\theenumi)}
\item The preceding functors $u$ and $v$ are fully faithful, and every
point of $E$ belongs to the essential image of exactly one of the
functors $u$ and $v$.

\item Let $p$ and $q$ be two points of $E$ such that $p$ (\resp $q$)
belongs to the essential image of $u$ (\resp of $v$). Then
\begin{equation*}
\Hom(q,p)=\phi{\quoi}.\tag{9.7.2.2}
\end{equation*}\etiquette[9.7.2.2]{IV.eq9.7.2.2}

\item Let $p$ (\resp $q$) be a point of $\mathcal{U}$ (\resp of
$\mathcal{F}$). There is then an isomorphism, functorial in $(q,p)$:
\begin{equation*}
\Hom(u(p),v(q))\simeq \Hom(\Pro(\rho)(p),q)\simeq
\Hom(p,\sigma(q)){\quoi},\tag{9.7.2.3}
\end{equation*}\etiquette[9.7.2.3]{IV.eq9.7.2.3}
where
$$
\sigma:\Pro(\mathcal{F})\longrightarrow \Pro(\mathcal{U})
$$
is the functor $\eqhyperref[IV.eq9.6.2]{9.6.2}$, pro-adjoint to the gluing
functor $\rho:\mathcal{U}\rightarrow \mathcal{F}$, and where, in
writing
$$
\rho:\mathcal{U}\longrightarrow \mathcal{F}{\quoi},
$$
we have\pageoriginale identified (up to equivalence) the categories
$\Ppoint(\mathcal{U})$ and $\Ppoint(\mathcal{F})$ with full
subcategories of $\Pro(\mathcal{U})$ and $\Pro(\mathcal{F})$,
respectively (\hyperref[IV.subsubsec6.8.5]{6.8.5}).
\end{enumerate}
\end{enonce*}

\noindent \emph{Proof}
\begin{enumerate}
\renewcommand{\theenumi}{\alph{enumi}}
\renewcommand{\labelenumi}{\theenumi)}
\item The full faithfulness of $u$ and $v$ is known
(\hyperref[IV.prop9.1.4]{9.1.4}). The assertion concerning the essential
images of $u$ and $v$ follows from the criteria
\hyperref[IV.subsubsec9.2.5]{9.2.5} and \hyperref[IV.prop9.2.4]{9.2.4} for a
point $P\rightarrow E$ to factor, up to isomorphism, through
$\mathcal{U}$ \resp $\mathcal{F}$, taking into account the fact that
the punctual topos $P$ has only two opens, namely the initial object
and the terminal object of $P$.

\item In view of a), the assertion means that if the point
$p:P\rightarrow E$ factors through $\mathcal{U}$, then the same holds
for every point $p':P\rightarrow E$ such that
$\Hom(p',p)\neq \phi$. Now, if $\alpha :p'\rightarrow p$, it induces
$p^{\ast}(U)\rightarrow {p'}^{\ast}(U)$, and
then $p^{\ast}(U)=e_p$ implies
${p'}^{\ast}(U)=e_p$, \cqfd

\item By the lemma \hyperref[IV.lem9.7.2.4]{9.7.2.4} below, applied to the
inclusions $j:\mathcal{U}\rightarrow E$ and
$i:\mathcal{F}\rightarrow E$, there is a diagram of functors
commuting up to canonical isomorphisms:
\[
\xymatrix@C=2cm@R=.55cm@M=.18cm{ \Point(\Uu) \ar@{^{(}->}[d]
\ar[r]^{u}& \Point(E) \ar@{^{(}->}[d]&
\Point(\Gg) \ar@{^{(}->}[d] \ar[l]_{v}\\
\Pro(\Uu) \ar[r]^{\Pro(j_{\ast})}  & \Pro(E)   & \Pro(\Gg)
\ar[l]_{\Pro(i_{\ast})} , }
\]
where the vertical arrows denote the canonical fully faithful
functors (\hyperref[IV.subsubsec6.8.5]{6.8.5}). Moreover, since
$i_{\ast}$ admits a left adjoint $i^{\ast}$,
$\Pro(i_{\ast})$ admits the left adjoint $\Pro(i^{\ast})$
\hyperrefext[I][I.sec8.11.5]{8.11.5}, and one finally obtains the functorial isomorphisms
\begin{align*}
\Hom(u(p),v(q)) \simeq & \Hom_{\Pro(E)}(\Pro(j_{\ast})(p),
\Pro(i_{\ast})(q))\\
& \simeq \Hom_{\Pro(\mathcal{F})}(\Pro(i^{\ast})(\Pro(j_{\ast})(p)),q)\\
& \simeq \Hom_{\Pro(\mathcal{F})}(\Pro(i^{\ast}j_{\ast})(p),q),
\end{align*}
which is precisely the first formula in \hyperref[IV.eq9.7.2.3]{(9.7.2.3)}, the second then
being immediate from the definition of $\sigma$ as a pro-adjoint of
$\rho$. This proves \hyperref[IV.prop9.7.2]{9.7.2}, subject to the
following lemma.
\end{enumerate}

\begin{enonce*}{Lemma 9.7.2.4}\etiquette[9.7.2.4]{IV.lem9.7.2.4}
Let $f:F\rightarrow E$ be a morphism of topoi. Then the diagram of
functors
\[\vcenter{
\xymatrix@C=2cm@R=.5cm{
  \Point(F)\ar[r]^{p \mapsto f \circ p} \ar[d] & \Point(E)\ar[d]\\
\Pro(F) \ar[r]^{\Pro(f_{\ast})}     & \Pro(E) }}\leqno{(9.7.2.4.1)}
\]\etiquette[9.7.2.4.1]{IV.eq9.7.2.4.1}
commutes\pageoriginale up to canonical isomorphism (the vertical
arrows denoting the fully faithful functors
$\hyperref[IV.subsubsec6.8.5]{6.8.5}$).
\end{enonce*}

The verification is immediate from the definitions and is left to
the reader.

\begin{enonce*}{Corollary 9.7.3}\etiquette[9.7.3]{IV.cor9.7.3}
Let $E$ be a topos having enough points. Then every closed subtopos
$\mathcal{F}$ of $E$ also has enough points. More precisely, if
$(p_i)_{i\in I}$ is a conservative family of points of $E$, then the
family of points of $\mathcal{F}$ defined by the subfamily
$(p_i)_{i\in J}$ consisting of those $p_i$ which come from points of
$\mathcal{F}$ is conservative. (Compare
$\hyperref[IV.cor6.7.3]{6.7.3}$, $\hyperref[IV.rem6.7.4]{6.7.4}$.)
\end{enonce*}

Indeed, let $f:X\rightarrow Y$ be an arrow of $\mathcal{F}$ which is
carried to an isomorphism by the fiber functors associated with the
points under consideration of $\mathcal{F}$. Since
$i^{\ast}i_{\ast}\simeq \id_{\mathcal{F}}$, it follows that
$i_{\ast}(f):i_{\ast}(X)\rightarrow i_{\ast}(Y)$ is carried to an
isomorphism by the fiber functors associated with the $p_i$ for
$i\in J$; the same is true for the $p_i$ with $i\in I-J$, \ie those
coming from points of $U$, since $j^{\ast}i_{\ast}$ is the constant
functor with value the terminal object. Therefore $i_{\ast}(f)$ is an
isomorphism, and hence so is $f\simeq i^{\ast}i_{\ast}(f)$, \cqfd

\begin{enonce*}[remark]{Remark 9.7.4}\etiquette[9.7.4]{IV.rem9.7.4}
If $E$ is a topos having enough points, it is clear that an open $U$
of $E$ is determined once the full subcategory of $\Ppoint(E)$ which
is the essential image of $\Ppoint(E_{/U})$ is known; in fact, if
$C$ is a full subcategory of $\Ppoint(E)$ defining a conservative
family of points of $E$, it even suffices to know the subcategory
$C_{/U}$ consisting of those elements of $C$ which belong to the
essential image of $\Ppoint(U)$. It follows from this and from
\hyperref[IV.prop9.7.2]{9.7.2} a) that a closed subtopos
$\mathcal{F}$ of $E$ is determined once the essential image of
$\Ppoint(\mathcal{F})$ in $\Ppoint(E)$ is known, or even merely its
intersection with $C$.
\end{enonce*}

\begin{enonce*}[remark]{Exercise 9.7.5}\etiquette[9.7.5]{IV.exer9.7.5}

% END INCLUDED SGA4_English_current_buildable_source_20260906/source/04/drafts/SGA4_IV_LINES_7444_7700_English.texfrag


% BEGIN INCLUDED SGA4_English_current_buildable_source_20260906/source/04/drafts/SGA4_IV_LINES_7701_7780_English.texfrag
Let $E$ be a topos. For every subtopos $F$ of $E$, let $P(F)$ be the
subset of $\Point(E)$ formed by the isomorphism classes of points of
$E$ which factor through \sisi{$E$}{$F$}. Thus $P(F)$ is homeomorphic
to $\Point(F)$ (\hyperref[IV.exer9.1.8]{9.1.8} c)). Show that if $F$
is an open (\resp closed, \resp locally closed
(\hyperref[IV.exer9.4.9]{9.4.9})) subtopos of $E$, then $P(F)$ is an
open\pageoriginale (\resp closed, \resp locally closed) subset of
\sisi{$E$}{$\Point(E)$}. Show that when $E$ has enough points, the
same is true of every locally closed subtopos of $E$, and the map
$F\mapsto P(F)$ induces a bijection between the set of open (\resp
closed, \resp locally closed) subtopoi of $E$ and the set of open
(\resp closed, \resp locally closed) subsets of $\Point(E)$. Extend
the preceding considerations to the case where $\Point(E)$ is
replaced by any subspace $X$ of $\Point(E)$ corresponding to a
conservative family of points of $E$.
\end{enonce*}

\subsection{Complements concerning certain topoi associated with topological spaces}\etiquette[9.8]{IV.subsec9.8}

In this section we give a few complements in the form of exercises.
The reader is advised to look through them in order to become
accustomed to the viewpoint of topoi in various classical
situations.

\begin{enonce*}[remark]{Exercise 9.8.1}\etiquette[9.8.1]{IV.exer9.8.1}
(\emph{Descent of sheaves on a topological space.})

Let $f:X\rightarrow Y$ be a continuous map of topological spaces in
$𝒰$, giving a functor
$$
f^{\ast}:\Top(Y)\longrightarrow\Top(X){\quoi}.
$$
\begin{enumerate}
\renewcommand{\theenumi}{\alph{enumi}}
\renewcommand{\labelenumi}{\theenumi)}
\item The following conditions are equivalent:
\begin{enumerate}
\renewcommand{\theenumii}{\roman{enumii}}
\renewcommand{\labelenumii}{(\theenumii)}
\item $f^{\ast}$ is faithful.

\item $f^{\ast}$ is conservative.

\item $f(X)$ is a very dense subset (EGA IV 10.1.3) of
$Y$.
\end{enumerate}
It suffices for $f$ or $f_{\sob}$ to be surjective. Give an example
where $f^{\ast}$ is faithful and $f=f_{\sob}$ is not surjective
(take $X$ discrete, $f$ injective, $f(X)$ the set of closed points
of $Y$, and $Y$ a nondiscrete Noetherian topological space). Give an
example where $f_{\sob}$ is surjective but $f$ is not, and another
where $f$ is surjective but $f_{\sob}$ is not.

\item Suppose that $f(X)$ is very dense in $Y$ and that its
topology is the quotient topology induced from $X$. Prove that the
arrow $f$ of $(\Esp)$ is a descent morphism for the fibered category
of sheaves of sets on varying spaces, \ie that the functor
$f^{\ast}$ induces a \emph{fully faithful}\pageoriginale functor
from $\Top(Y)$ to the category of objects of $\Top(X)$ equipped with
descent data relative to $f:X\rightarrow Y$. Reduce to the case
where $f$ is surjective and follow the argument of VIII \hyperref[VIII.prop9.1]{9.1} given in
the setting of étale topologies of schemes. Give a converse.

\item Suppose that $f(X)$ is very dense in $Y$, that its
topology is the quotient topology induced from $X$, and that the
fibers of $f$ are connected. Prove that the functor $f^{\ast}$ is
fully faithful. (Follow the argument of \SGA 1 IX 3.4.) Give a
converse, at least if the points of $Y$ are closed.

\item Suppose that $Y$ is the union of open subsets $Y_i$ over which
$X$ admits sections. Prove that $f$ is then an \emph{effective
descent morphism} for the fibered category of sheaves of sets on
varying topological spaces, \ie that the functor considered in b) is
in fact an equivalence of categories.
\end{enumerate}
\end{enonce*}

% END INCLUDED SGA4_English_current_buildable_source_20260906/source/04/drafts/SGA4_IV_LINES_7701_7780_English.texfrag


% BEGIN INCLUDED SGA4_English_current_buildable_source_20260906/source/04/drafts/SGA4_IV_LINES_7781_8025_English.texfrag
\begin{enonce*}[remark]{Exercise 9.8.2}\etiquette[9.8.2]{IV.exer9.8.2}
(\emph{Quotient topoi of topological spaces. Topological
\'etendues.})

\begin{enumerate}
\renewcommand{\theenumi}{\alph{enumi}}
\renewcommand{\labelenumi}{\theenumi)}
\item Let
\hypertarget{sga4.exp.IV.equation.9.8.2.1}{}\[
\xymatrix@C=1.5cm@R=1.5cm@M=.2cm{
X_2 \ar@<6pt>[r] \ar@<-1.5pt>[r] \ar@<-8pt>[r]& X_1 \ar@<2pt>[r]^{p_1} \ar@<-4pt>[r]_{p_2}& X_0\\
}\leqno{(9.8.2.1)}
\]\etiquette[9.8.2.1]{IV.eq9.8.2.1}
be a semi-simplicial object truncated in degree $2$ in the category
$(\Esp)$ of topological spaces $\in 𝒰$; the inverse-image functors
then give a diagram of categories of sheaves
\[
\xymatrix@M=.2cm@C=1.25cm{
  \Top(X_0) \ar@<2pt>[r]^{p_1^{\ast}} \ar@<-2pt>[r]_{p_2^{\ast}}&
  \Top(X_1) \ar@<-5pt>[r] \ar[r] \ar@<5pt>[r]& \Top(X_2).
}
\]
(More precisely, one has a category cofibered over the category of
standard simplices $\Delta_n$, $0\leq n\leq 2$.) Show that the
category $\varprojlim$ of this diagram (\ie the category of sheaves
of sets on $X_0$ \emph{equipped} with a descent datum relative to
the diagram \sisi{(9.8.6.1)}{\hyperref[IV.eq9.8.2.1]{9.8.2.1}}, \ie with an
isomorphism $p^{\ast}_1(F)\simeq p^{\ast}_2(F)$ satisfying\ldots), is
a $𝒰$-topos $T$. (For a more general statement on the stability of
topoi under $\varprojlim$ operations, \cf \sisi{VI)}{\hyperrefext[VI][VI.sec8]{8}}.) Define a morphism
of topoi $q:\Top(X_0)\rightarrow T$, and prove that for every
$𝒰$-topos\pageoriginale $E$, $\Hhomtop(T,E)$ is equivalent, via
$q^{\ast}$, to the category $\varprojlim$ of the following diagram
of categories, deduced from \hyperref[IV.eq9.8.2.1]{(9.8.2.1)}:
\[
\xymatrix@M=.2cm{ \Hhomtop(\Top(X_0),E)\ar@<2pt>[r] \ar@<-4pt>[r]&
\Hhomtop(\Top(X_1),E)\ar[r]\ar@<6pt>[r]\ar@<-6pt>[r]&
\Hhomtop(\Top(X_2),E). }
\]
\item In particular, let $R$ be an equivalence relation in an object
$X$ of $(\Esp)$, \ie (\sisi{\hyperrefext[I][I.sec10]{10}}{\hyperrefext[I][I.sec10.9]{10.9}}) a subobject $R$ of
$X\times X$ (N.B. the inclusion $R\hookrightarrow X\times X$ is
continuous, but the topology of $R$ is not necessarily induced by
that of $X\times X$), such that, for every object $Y$ of $(\Esp)$,
$\Hom(Y,R)$ is the graph of an equivalence relation in $\Hom(Y,X)$.
This gives a diagram of the form \hyperref[IV.eq9.8.2.1]{(9.8.2.1)}, with $X_0=X$, $X_1=R$,
$X_2=\fprod{(R,p_2)}{R,p_1}{X}$, where $p_1$, $p_2$ are the two
projections from $R$ to $X$, and hence a topos, which we denote by
$\Top(X)/R$, or even $\Top(X/R)$, or indeed $X/R$, by abuse of
notation, and which plays the role of a quotient of $X$ by $R$ in
the sense made precise in a). Suppose that the topology of $R$ is
induced by that of $X\times X$ and that, denoting by $X/R$ the
ordinary quotient topological space, $X$ locally admits sections
over $X/R$. Prove that $T$ is then equivalent to $\Top(X/R)$.
(Use \hyperref[IV.exer9.8.1]{9.8.1} d).)

\item Let $H\rightarrow X$ be an injective morphism of topological
groups, \ie a monomorphism of group objects in $(\Esp)$, and let $R$
be the equivalence relation that it defines in $X$, \ie
$R=H\times X$, with $p_1=\pr_2$ and $p_2$ defined by the action of
$H$ on $X$ by left translations. Show that the topology of $R$ is
induced by that of $X\times X$ if and only if the topology of $H$ is
induced by that of $X$. Give examples in which this condition is not
satisfied and in which the ordinary quotient topology on $X/H$ is
the indiscrete topology, with a) $H=\mathbb{Z}\times\mathbb{Z}$,
$X=\mathbb{R}$, and b) $H=\mathbb{R}$,
$X=\mathbb{T}\times\mathbb{T}$, where
$\mathbb{T}=\mathbb{R}/\mathbb{Z}$ (``geodesics on the torus'').
Prove that the two topoi so obtained are equivalent and are not
equivalent to $\Top(X/H)$ (which is a final topos
(\hyperref[IV.subsec2.2]{2.2})), nor to any topos of the form $\Top(Y)$.

\item Let $G$ be a topological group acting on a topological space
$X$, whence, in the familiar way, a semi-simplicial object
\[
\xymatrix@M=.2cm{ ... G\times G \times X
\ar@<6pt>[r]\ar@<-1.5pt>[r]\ar@<-8pt>[r] &
  G\times X \ar@<3pt>[r] \ar@<-4pt>[r] & X.
}
\]
Show\pageoriginale that the topos $T$ deduced from it by a) identifies
with the category of spaces $X'$ over $X$ having $G$ as a group of
operators, such that $X'\rightarrow X$ is an \'{e}tale map
(compatible with the action of $G$). In particular, when $G$ is
discrete, one recovers the topos $\Top(X,G)$ of \hyperref[IV.subsec2.3]{2.3}.

The preceding considerations justify the notation $\Top(X)/G$, or
indeed $\Top(X/G)$ or even simply $X/G$, for the preceding topos $T$.
One should bear in mind, however, that when $G$ is a discrete group
(to fix ideas) acting \emph{properly} on $X$, so that the ordinary
quotient space $X/G$ has reasonably good properties, the natural
morphism of topoi $T\rightarrow\Top(X/G)$ deduced from the universal
characterization a) of $T$ is an equivalence of topoi only if $G$
acts freely, \ie without fixed points, \ie when
$G\times X\rightarrow X\times X$ is a monomorphism. Thus, in the
case of a ``pre-equivalence relation'' (or ``groupoid'' in the sense
of (\SGA 3 V 1)) which is not an equivalence relation, passage to
the quotient in the sense of topoi (or ``passage to the fine
quotient'') does not in general correspond (via the correspondence
$Y\mapsto\Top(Y)$) to the usual topological quotient.

\item Let $T$ be a $𝒰$-topos. Show that the following conditions are
equivalent:
\begin{itemize}
\item[(i)] There exists a family $(S_i)_{i\in I}$ of objects of $T$
covering the final object of $T$, such that for every $i\in I$ the
induced topos $T_{/S_i}$ is equivalent to a topos of the form
$\Top(X)$, with $X\in(\Esp)$.

\item[${\rm (i')}$] There exists an $S\in\Ob T$ covering the final
object such that the induced topos $T_{/S}$ is equivalent to a topos
of the form $\Top(X)$.

\item[(ii)] There exist a topological space $X$ and an
\emph{\'{e}tale} pre-equivalence relation $R$ in $X$ (in the sense of
the category $(\Esp)$) (\ie one for which $p_1:R\rightarrow X$ is an
\'{e}tale map), such that $T$ is equivalent to $\Top(X/R)$, with the
notation of b).
\end{itemize}

The\pageoriginale topos $T$ is then said to be \emph{locally a
topological space}, or alternatively a \emph{topological
\'{e}tendue} (or simply an \emph{\'{e}tendue}, when no confusion is
likely).

\item Show that the topos $T$ constructed in a) has enough points
and, more precisely, that it admits a conservative family of points
parametrized by the (small) set $X_0$. In particular, every
\'{e}tendue has a small conservative family of points, and hence has
enough points. When an \'{e}tendue $T$ is represented in the form
$\Top(X/R)$ as in e) ii), prove that every point of $T$ is isomorphic
to the image of a point of $\Top(X)$, and is therefore defined by an
ordinary point of $X_{\sob}$ (or of $X$ itself, when $X$ is sober).

\item Show that the topos $\Top(X/G)$ of \hyperref[IV.subsec2.3]{2.3} is an
\'{e}tendue. (Use its description in d) or \hyperref[IV.exer7.1.10]{7.1.10} c).)
Show that, for every $x\in X$, the monoid of endomorphisms of the
point of $\Top(X,G)=\Top(X)/G$ defined by $x$ is canonically
isomorphic to the stabilizer group $G_x$ of $x$. In particular, the
points of an \'{e}tendue may have nontrivial automorphism groups.
Determine $\Ppoint(T)$ for the \'{e}tendue $T$ defined by an \'{e}tale
pre-equivalence relation in a space $X$, and show that every
endomorphism of a point of $T$ is an automorphism.

\item Let $T$ be a topos. Show that the following conditions are
equivalent:
\begin{itemize}
\item[i)] For every point $p$ of $T$, every endomorphism of $p$
(resp. every automorphism of $p$) is the identity.

\item[ii)] For every topos $S$ having enough points and every
morphism of topoi $f:S\rightarrow T$, every endomorphism of $f$
(resp. every automorphism of $f$) is the identity.
\end{itemize}
Likewise, show that the following conditions are equivalent:
\begin{itemize}
\item[${\rm i')}$] For two points $p$, $q$ of $T$, there exists at
most one morphism from $p$ to $q$.

\item[${\rm ii')}$] For two morphisms $f$, $g:S\rightrightarrows T$
from a topos $S$ having enough\pageoriginale points to $T$, there
exists at most one morphism from $f$ to $g$.
\end{itemize}

When these latter conditions hold, the topos $T$ is said to be
\emph{pointwise rigid}, or simply \emph{rigid}. Show that if $X$ is
a topological space equipped with an \'{e}tale pre-equivalence
relation $R$, the quotient \'{e}tendue $T=\Top(X)/R$ is rigid if and
only if $R$ is an equivalence relation.

\item Let $(X,G)$ be a topological space with a discrete group of
operators, let $H$ be a normal subgroup of $G$ such that the induced
action of $H$ on $X$ is proper and free, and let $X'=X/H$,
$G'=G/H$, so that $G'$ acts on $X'$ in the evident way. Prove that
the morphism of topoi
$$
\Top(X,G)\longrightarrow \Top(X',G')
$$
induced by the canonical morphism $(X,G)\rightarrow(X',G')$ of
spaces with operators (\hyperref[IV.subsubsec4.1.2]{4.1.2}) is an equivalence of
topoi.

Show that the set of open subtopoi of $\Top(X/G)$ identifies with
the set of open subsets of $X$ stable under the action of $G$, or
equivalently with the set of open subsets of the quotient
topological space $X/G$. Deduce that $\Top(X/G)$ is connected if and
only if every subset of $X$ which is both open and closed and stable
under the action of $G$ is either empty or equal to $X$. When the
connected components of $X$ are open, $\Top(X/G)$ is connected if
and only if $G$ acts transitively on the set $\pi_0(X)$ of connected
components of $X$; more generally, $\pi_0(\Top(X/G))$
(\hyperref[IV.exer8.7]{8.7} g)) is an essentially constant pro-set
isomorphic to the quotient set $\pi_0(X)/G$.

Show that when $X$ is locally connected and locally simply
connected, and $T=\Top(X,G)$ is connected, \ie $G$ is transitive on
$\pi_0(X)$, then among all the ways of representing $T$ by means of
a space with operators $(X',G')$ (\cf above for some such
representations), there is one, unique up to nonunique isomorphism,
for which $X'$ is connected and simply connected. Show that choosing
such an $(X',G')$ amounts to choosing a ``universal covering'' of
the final object of $T$, and that $G'$ is isomorphic to the group
$\pi_1(T)$ relative to this universal covering (\cf
\hyperref[IV.exer2.7.5]{2.7.5}).

\noindent (Hint:\pageoriginale show that giving an equivalence of
$T$ with a $\Top(X,G)$ amounts to giving a $G_T$-torsor $P$ in $T$,
and that, if $(X,G)$ and $P$ correspond, there is a canonical
equivalence $\Top(X)\simeq T_{/P}$.)

Deduce from this a trivial proof of the last assertion of c).
Characterize the topoi $T$ equivalent to a $\Top(X,G)$, where $G$
is a discrete group acting on a locally connected and locally simply
connected topological space while transitively permuting its
connected components, as the connected, locally connected, and
locally simply connected \'{e}tendues whose universal covering $P$
of the final object $e_T$ of $T$ is such that the induced topos
$T_{/P}$ is equivalent to a topos of the form $\Top(X)$ (or, as we
shall simply say by abuse of language, $T_{/P}$ ``is'' a topological
space).

Give an example of an \'{e}tendue, locally isomorphic to $\mathbb{R}$,
which is connected and simply connected but is not a topological
space. (Take the universal covering of the line $\mathbb{R}$ with
doubled origin.)

\item A ringed topos $(T,{\Oo}_T)$ (\cf \hyperref[IV.sec11]{\S 11}) is called a
\emph{differentiable \'{e}tendue} (resp. a \emph{real-analytic
\'{e}tendue}, resp. a \emph{complex-analytic \'{e}tendue}) if there
exist objects $S_i$ of $T$ covering the final object such that the
induced ringed topos $(T_{/S_i},{\Oo}_T/S_i)$ is equivalent to the
ringed topos defined by a differentiable manifold with its sheaf of
real $C^{\infty}$ functions (resp. to the ringed topos defined by a
real-analytic space, resp. to the ringed topos defined by a
complex-analytic space). Give a constructive description of these
ringed topoi of the type in e) ii) (taking for $X$, respectively, a
differentiable manifold, a real-analytic space, or a
complex-analytic space). Give examples of such ringed topoi which are
not equivalent to ringed topoi associated with topological spaces,
by returning to the examples considered in c).
\end{enumerate}
\end{enonce*}

% END INCLUDED SGA4_English_current_buildable_source_20260906/source/04/drafts/SGA4_IV_LINES_7781_8025_English.texfrag


% BEGIN INCLUDED SGA4_English_current_buildable_source_20260906/source/04/drafts/SGA4_IV_LINES_8026_8245_English.texfrag

\begin{enonce*}[remark]{Exercise 9.8.3}\etiquette[9.8.3]{IV.exer9.8.3}
(\emph{Topoi\pageoriginale associated with local equivalence relations
and with foliations}) \footnote[1]{This exercise is given subject to
every reservation, since it was written hastily and insufficiently
checked. Mr. N. \textsc{Saavedra} checked parts a) through e). The
reader is requested to communicate any observations to us.}.

Let $X$ be a topological space.
\begin{enumerate}
\renewcommand{\theenumi}{\alph{enumi}}
\renewcommand{\labelenumi}{\theenumi)}
\item For every open subset $U$ of $X$, let $\Quot(U)$ be the set of
equivalence relations on $U$, and, for an open subset $V\subset U$,
consider the natural map $\Quot(U)\rightarrow \Quot(V)$. This gives a
presheaf $\Quot_X$ on $X$. Let $Q_X$, or simply $Q$, be the associated
sheaf, and let $r$ be a section of $Q$ (a ``\emph{local equivalence
relation on $X$}''). From the order relation on the set $\Quot(U)$ of
equivalence relations on an open subset $U$ of $X$, deduce on the
presheaf $\Quot$, and on the associated sheaf $Q$, an order structure,
\ie a structure of a presheaf \resp a sheaf with values in the category
of ordered sets.

\item For every equivalence relation $R$ on $X$, consider the section
$\loc(R)$ of $Q$ that it defines. Let $E(r)$ be the set of equivalence
relations on $X$ such that $\loc(R)$ is less fine than $r$, and let
$\glob(r)$ be the least upper bound of the equivalence relations in
$E(r)$ (its graph is the intersection of the graphs of the $R\in
E(r)$). Let $(U_x)_{x\in X}$ be a family of open neighborhoods of the
points $x\in X$, and, for every $x\in X$, let $R_x$ be an equivalence
relation on $U_x$ whose germ at $x$ is $r_x$. For every family $𝒱$
of open subsets $V_x$ $(x\in X, x\in V_x\subset U_x)$, let $R_{𝒱}$
be the equivalence relation on $X$ generated by the family of
relations $R_x| V_x$. Show that if $𝒱'\leq 𝒱$ (in the obvious
sense), then $R_{𝒱'}\geq R_{𝒱}$, and that $\glob(r)$ is the least
upper bound of the increasing filtered family of equivalence relations
$R_{𝒱}$. Give an example in which $\glob(r)$ does not belong to
$E(r)$, \ie in which there is an $a\in X$ such that, for\pageoriginale
every open neighborhood $W\subset U_a$ of $a$, there exist a
$𝒱=(V_x)$ and two points $y,z$ of $W$ which are equivalent under
$R_a$ but not under $R_{𝒱}$. Show that, for every $R\in\Quot(X)$,
one has $\glob\loc(R)\geq R$.

\item The local equivalence relation $r$ is said to be
\emph{precoherent (\resp coherent)} if $\glob(r)\in E(r)$, \ie if
$\loc(\glob(r))\leq r$ (\resp if, for every open subset $U$ of $X$,
the restriction of $r$ to $U$ is precoherent). It is said to be
globally coherent if it is coherent and if, in addition,
$r=\loc\glob(r)$. An equivalence relation $R$ on $X$ is said to be
\emph{locally coherent} if $r=\loc(R)$ is coherent, and
\emph{coherent} if, in addition, $R=\glob(r)$, \ie if
$R=\glob(\loc(R))$. Show that globally coherent local equivalence
relations $r$ on $X$ correspond bijectively to coherent equivalence
relations $R$ on $X$, by $r\mapsto\glob(r)$ and
$R\mapsto\loc(R)$.

\item Show that, in order for $r$ to be coherent, it is enough that,
for every $a\in X$ and every open neighborhood $W'\subset U_a$ of
$a$, there exist an open neighborhood $W\subset W'$ of $a$ such that
every equivalence class of $R_a|W$ is contained in a connected
component of an equivalence class of $R_a|W'$ (a fortiori, it is
enough that there exist a fundamental system of open neighborhoods
$W_a\subset U_a$ of $a$ such that the fibers of the induced
equivalence relation $R_a|W_a$ are connected). (Hint: reduce to
proving precoherence in the case where $r$ is defined by an
$R\in\Quot(X)$, and observe in that case that the fibers of the
equivalence relations $R_{𝒱}$ are relatively open---and hence also
relatively closed---subsets of the fibers of $R$.) Show that, in order
for a global equivalence relation $R$ to be coherent, it is enough
that it be locally coherent and that its fibers be connected; prove
that this condition is necessary if the fibers are closed. Give an
example of an equivalence relation with connected and locally
connected fibers which is not locally coherent (?).

\item Let\pageoriginale $f:X'\rightarrow X$ be a continuous map.
Define a natural homomorphism of presheaves
$f_{\ast}(\Quot_{X'})\leftarrow\Quot_X$, inducing a homomorphism of
sheaves $f_{\ast}(Q_{X'})\leftarrow Q_X$, and hence a homomorphism
$f^{\ast}(Q_X)\rightarrow Q_{X'}$. If $r$ is a local equivalence
relation on $X$, then $f$ is said to be a ``\emph{fiber map}'' for
the local equivalence relation $r$ on $X$ if the inverse image of $r$
is the coarse local equivalence relation on $X'$. For every
topological space $X'$, let $\Homfib_r(X',X)$ be the set of continuous
maps from $X'$ to $X$ which are fiber maps relative to $r$. Show that,
as $X'$ varies, this defines a contravariant functor in $X'$.

\item Suppose that $r$ is coherent. Show that the preceding functor is
representable by a topological space $X^r$ over $X$. Show that the
universal fiber map $\Phi:X^r\rightarrow X$ is bijective, and that
every $x\in X^r$ has an open neighborhood $U$ such that the map
$U\rightarrow X$ induced by $\Phi$ is a homeomorphism of $U$ onto its
image. Show that the connected components of $X^r$ are open and
correspond under the bijection $\Phi$ to the fibers of the equivalence
relation $R=\glob(r)$. Give an example in which the restriction of
$\Phi$ to the connected components of $X^r$ does not induce
homeomorphisms from those spaces onto their images in $X$.

\item The local equivalence relation $r$ is said to be \emph{open} if
it is a section of the subsheaf $\Qouv_X$ of $Q$ arising from the
subpresheaf $\Quotouv_X$ of $\Quot_X$ whose value on every open subset
$U$ of $X$ is the set of \emph{open} equivalence relations on $U$.
The local equivalence relation $r$ is said to be \emph{strictly open}
if it is coherent and open or, equivalently, if it is coherent and if,
for every open subset $U$ of $X$, $\glob(r|U)$ is an open equivalence
relation on $U$. An equivalence relation $R$ on $X$ is said to be
\emph{strictly open} if it is coherent and if $\loc(R)$ is a strictly
open local equivalence relation. Prove that, for the open relation $r$
to be strictly open, it is enough that it satisfy\pageoriginale the
condition considered in d); if $r$ is locally defined by an $R$ with
closed fibers, then that condition is also necessary.

\item Let $F$ be a sheaf of sets on $X$. For every open subset $U$ of
$X$, let $\Quot(U,F)$ be the set of pairs consisting of an equivalence
relation $R_U$ on $U$ and an equivalence relation $R_{F|U}$ on $F|U$
($F$ being regarded as an étale space over $X$), such that the
structure map $p:F|U\rightarrow U$ is compatible with the equivalence
relations $R_U$ and $R_{F|U}$, and such that the corresponding diagram
of topological spaces
\[
\xymatrix@C=2cm@R=1cm{
F|U \ar[r]\ar[d] \ar[d]_{p}& (F|U)/R_{F|U} \ar[d]_{q}\\
U   \ar[r]&  U/R_{U} }
\]
is Cartesian, with $q$ an étale map (so that $F|U$ is identified
with the inverse image of the sheaf $(F|U)/R_{F|U}$ on $U/R_U$).
Show that the $\Quot(U,F)$, for varying $U$, define a presheaf on $X$,
whose associated sheaf will be denoted by $Q(F/X)$. Define a
homomorphism of sheaves $Q(F/X)\rightarrow Q_X$. For a given section
$r$ of $Q_X$, an $r$-\emph{structure} on the sheaf $F$ is a section
of $Q(F/X)$ lying above the given section $r$ of $Q_X$. Define the
category $\Top(X/r)$ of $r$-sheaves on $X$, and define a conservative
and faithful functor, ``forgetting the $r$-structure'',
$\Top(X/r)\rightarrow\Top(X)$. Prove that finite projective limits and
finite inductive limits are representable in $\Top(X/r)$, and that the
preceding functor commutes with those limits. Deduce that finite sums
in $\Top(X/r)$ are disjoint and universal and that equivalence
relations are universally effective. Prove that $\Top(X/r)$ has a
small generating family. Give an example, with $r$ coherent, in which
$\Top(X/r)$ does not have infinite coproducts (?), and hence is not a
topos.

\item Let $f:X\rightarrow Y$ be a continuous map compatible with
$R=\glob(r)$. Define a functor
$$
f^{\ast}_r:\Top(Y)\longrightarrow\Top(X/r)
$$
whose\pageoriginale composite with the inclusion functor
$\Top(X/r)\rightarrow\Top(X)$ is $f^{\ast}$. Show that, if $r$ is
globally coherent and strictly open (\ie if $R$ is strictly open and
$r=\loc(R)$), and if $f$ is the canonical map
$X\rightarrow Y=X/R$, then $f^{\ast}_r$ is an equivalence of
categories. Consequently, $\Top(X/r)$ is a $𝒰$-topos, and
$\Top(X/r)\rightarrow\Top(X)$ is the inverse-image functor of a
morphism of topoi.

\item Deduce from i) that, if $r$ is strictly open, then $\Top(X/r)$
is a topos, and that the inclusion functor
$\Top(X/r)\rightarrow\Top(X)$ is the inverse-image functor associated
with a morphism of topoi
$$
p:\Top(X)\longrightarrow\Top(X/r){\quoi}.
$$

\item Define a functorial construction of the topos $\Top(X/r)$ and
of the preceding morphism of topoi $p$ with respect to the pair
$(X,r)$ consisting of a topological space $X$ equipped with a strictly
open local equivalence relation. Show that, if $X'\rightarrow X$ is an
étale map and $r'$ is the inverse image of $r$ on $X'$, then the
induced morphism $\Top(X'/r')\rightarrow\Top(X/r)$ is equivalent to a
localization morphism $\Top(X/r)_{/E}\rightarrow\Top(X/r)$, where $E$
is an object of $\Top(X/r)$ (determined up to a unique isomorphism).
The object $E$ covers the final object of $\Top(X/r)$ if and only if
the saturation under $R=\glob(r)$ of the image of $X'$ in $X$ is all
of $X$.

\item Deduce from k) that $\Top(X/r)$ is a \emph{rigid étendue}
(\hyperref[IV.exer9.8.2]{9.8.2} h)). Show that, if $X$ is sober, the
set of isomorphism classes of points of $\Top(X/r)$, endowed with its
canonical topology (\hyperref[IV.exer7.8]{7.8}), is homeomorphic to the
quotient topological space $X/\glob(r)$, and that, for any two points
of $\Top(X/r)$ arising from points $x$ and $x'$ of $X$, there is at
most one morphism from one to the other. Such a morphism exists if and
only if $x'$ and $x$ are equivalent modulo $\glob(r)$ to points $x'_1$
and $x_1$ such that $x_1$ is a generalization of $x'_1$.

\item Study the canonical morphism of topoi
$\Top(X)\rightarrow\Top(X/r)$, observing\pageoriginale that, for every
object $U$ of $\Top(X/r)$ such that the topos induced on $U$ ``is'' an
ordinary topological space, the induced morphism
$\Top(X)_{/p^{\ast}(U)}\rightarrow\Top(X/r)_{/U}$ ``is'' a continuous
map of ordinary topological spaces $X_U\rightarrow U$. Show that the
fibers of the map $X_U\rightarrow U$ are homeomorphic to connected
components of the space $X^r$ of f).

\item Let $X$ be a differentiable manifold equipped with a foliation,
\ie with a subsheaf $F$, locally a direct factor of the tangent sheaf
of $X$, which is stable under brackets. Define on $X$ a strictly open
local equivalence relation associated with the foliation, and define
on the quotient topos $T=\Top(X/r)$ a ring which makes it a
\emph{differentiable étendue} (\hyperref[IV.exer9.8.2]{9.8.2} j)).
Define an equivalence between the category of locally free modules on
the differentiable étendue $T$ (also called differentiable vector
bundles on $T$) and the category of locally free modules on $X$
equipped with a connection relative to the given subbundle $F$ of the
tangent bundle. Give variants in the real-analytic and complex-analytic
cases.
\end{enumerate}
\end{enonce*}

% END INCLUDED SGA4_English_current_buildable_source_20260906/source/04/drafts/SGA4_IV_LINES_8026_8245_English.texfrag


% BEGIN INCLUDED SGA4_English_current_buildable_source_20260906/source/04/drafts/SGA4_IV_LINES_8246_8573_English.texfrag
\section{Sheaves of morphisms}\etiquette[10]{IV.sec10}\pageoriginale

\begin{enonce*}{Proposition 10.1}\etiquette[10.1]{IV.prop10.1}
Let $E$ be a topos and let $X$ and $Y$ be two objects of $E$; the
functor
$Z\mapsto \Hom_E(Z\times X,Y)\simeq \Hom_{E_{/Z}}(X_Z,Y_Z)$ is
representable.
\end{enonce*}

Indeed, inductive limits are universal in $E$ (II \hyperref[II.prop4.3]{4.3}). Hence the
functor $Z\mapsto Z\times X$ commutes with inductive limits.
Consequently, the functor $Z\mapsto \Hom_E(Z\times X,Y)$ carries
inductive limits in the argument $Z$ to projective limits. It is
therefore representable \sisi{(IV~\hyperref[IV.prop1.4]{1.4} and
\hyperref[IV.theo1.2]{1.2})}{(\hyperref[IV.cor1.2.1]{1.2.1})}.

\setcounter{subsection}{1}
\subsection{}\etiquette[10.2]{IV.subsec10.2}
The object representing the functor $Z\mapsto \Hom_E(Z\times X,Y)$
is denoted by $\scrHom_E(X,Y)$ (or more simply by $\scrHom(X,Y)$) and
is called the \emph{sheaf of morphisms from $X$ to $Y$}. It is a
bifunctor in $X$ and $Y$. We thus have a trifunctorial isomorphism
\begin{equation*}
\Hom_E(Z,\scrHom_E(X,Y))\simeq \Hom_E(Z\times X,
Y){\quoi}.\tag{10.2.1}
\end{equation*}\etiquette[10.2.1]{IV.eq10.2.1}

It then follows from formula \eqhyperref[IV.eq10.2.1]{10.2.1} that the
bifunctor $(X,Y)\mapsto \scrHom(X,Y)$ carries inductive limits in the
argument $X$ (\resp projective limits in the argument $Y$) to
projective limits in $E$.

\begin{enonce*}{Proposition 10.3}\etiquette[10.3]{IV.prop10.3}
Let $v:E\rightarrow E'$ be a morphism of topoi. For a variable object
$X$ of $E$ and a variable object $Y$ of $E'$, there is a bifunctorial
isomorphism
\begin{equation*}
v_{\ast}\scrHom_E(v^{\ast}Y,X)\simeq
\scrHom_{E'}(Y,v_{\ast}X){\quoi}.\tag{10.3.1}
\end{equation*}\etiquette[10.3.1]{IV.eq10.3.1}
\end{enonce*}

\noindent \emph{Proof}. For every object $Z$ of $E'$, we have a
sequence of isomorphisms, functorial in all the arguments:
\begin{alignat*}{4}
&\Hom_{E'}(Z,v_{\ast}\scrHom_E(v^{\ast}Y,X)) &&\simeq
&& \Hom_E(v^{\ast}Z,\scrHom_E(v^{\ast}Y,X)) &&\quad \text{(adjunction)}\\
& \Hom_E(v^{\ast}Z,\scrHom_E(v^{\ast}Y,X)) && \simeq
&& \Hom_E(v^{\ast}Z\times v^{\ast}Y,X) &&\quad \eqhyperref[IV.eq10.2.1]{10.2.1}\\
&\Hom_E(v^{\ast}Z\times v^{\ast}Y,X) &&\simeq &&
\Hom_E(v^{\ast}(Z\times Y),X) &&\quad \text{($v^{\ast}$ is left
exact)}\\
&\Hom_E(v^{\ast}(Z\times Y),X) &&\simeq && \Hom_{E'}(Z\times Y,
v_{\ast}X) &&\quad \text{(adjunction)}\\
& \Hom_{E'}(Z\times Y,v_{\ast}X) &&\simeq &&
\Hom_{E'}(Z,\scrHom_{E'}(Y,v_{\ast}X)) &&\quad \eqhyperref[IV.eq10.2.1]{10.2.1}
\end{alignat*}
The two sides of \eqhyperref[IV.eq10.3.1]{10.3.1} therefore represent
isomorphic functors, \cqfd

\begin{enonce*}{Corollary 10.4}\etiquette[10.4]{IV.cor10.4}
Let\pageoriginale $C$ be a $𝒰$-site, $X$ a $𝒰$-presheaf on
$C$, $Y$ a $𝒰$-sheaf on $C$, and $𝒱$ a universe containing
$𝒰$ such that $C$ is $𝒱$-small. Let
$C^{\widehat{\phantom{,}}}_{𝒱}$ be the topos of $𝒱$-presheaves on
$C$, and let $C^{\sim}_{𝒰}$ be the topos of $𝒰$-sheaves on $C$.
The presheaf
$\scrHom_{C^{\widehat{\phantom{,}}}_{𝒱}}(X,Y)$ is a $𝒰$-sheaf. Let
$X\rightarrow \underline{a}X$ be the canonical morphism from $X$ to
its associated sheaf. The corresponding morphism
$\scrHom_{C^{\widehat{\phantom{,}}}_{𝒱}}(\underline{a}X,Y)\rightarrow
\scrHom_{C^{\widehat{\phantom{,}}}_{𝒱}}(X,Y)$ is an isomorphism. There
is a canonical isomorphism
$\scrHom_{C^{\widehat{\phantom{,}}}_{𝒱}}(\underline{a}X,Y)\simeq
\scrHom_{C^{\sim}_{𝒰}}(\underline{a}X,Y)$.
\end{enonce*}

\setcounter{subsection}{4} \setcounter{subsubsection}{0}
\subsubsection{}\etiquette[10.4.1]{IV.subsubsec10.4.1}
Suppose first that $C$ is a small site and that $𝒱=𝒰$. We then have
a morphism of topoi
$C^{\sim}_{𝒰}\rightarrow
C^{\widehat{\phantom{,}}}_{𝒰}$ (\sisi{IV~}{}\hyperref[IV.subsec4.8]{4.8}),
from which \eqhyperref[IV.eq10.3.1]{10.3.1} gives the assertions in
this case. To pass from this to the general case, note first that the
``associated sheaf'' functor is independent of the universe
(\hyperrefext[II][II.prop3.6]{3.6}), and that the topos of
$𝒱$-sheaves on $C$ is equivalent to the topos of $𝒱$-sheaves on
$C^{\sim}_{𝒰}$ (\hyperrefext[III][III.sec4]{4} and
\hyperref[IV.sec1]{1}). It therefore suffices to prove the following
lemma:

\begin{enonce*}{Lemma 10.4.2}\etiquette[10.4.2]{IV.lem10.4.2}
Let $E$ be a $𝒰$-topos, let $𝒱$ be a universe containing $𝒰$, let
$E^{\sim}_{𝒱}$ be the topos of $𝒱$-sheaves on $E$, and let $X$ and
$Y$ be two objects of $E$. There exists a canonical isomorphism
$\scrHom_E(X,Y)\simeq \scrHom_{E^{\sim}_{𝒱}}(X,Y)$.
\end{enonce*}

\setcounter{subsubsection}{2}
\subsubsection{}\etiquette[10.4.3]{IV.subsubsec10.4.3}
The functor $\epsilon_E:E\rightarrow E_{\widetilde{𝒱}}$ is fully
faithful and left exact. Consequently, for every object $Z$ of $E$,
there is an isomorphism
$$
\Hom_{E_{\widetilde{𝒱}}}(\epsilon_E Z,\epsilon_E
\scrHom_E(X,Y))\simeq \Hom_{E_{\widetilde{𝒱}}}(\epsilon_E
Z\times \epsilon_E X,\epsilon_E Y).
$$
But every object of $E_{\widetilde{𝒱}}$ is an inductive limit of
objects coming from $E$; the assertion therefore follows from
\eqhyperref[IV.eq10.2.1]{10.2.1}.

\subsection{}\etiquette[10.5]{IV.subsec10.5}
Let $v:E\rightarrow E'$ be a morphism of topoi. We shall define four
bifunctorial morphisms:
\begin{alignat*}{4}
\Phi_v & : v^{\ast}\scrHom(X,Y) &&\longrightarrow
\scrHom(v^{\ast}X,v^{\ast}Y) &\quav
& X \text{ and } Y \text{ objects of } E'{\quoi},\\
\Psi_v & : v_{\ast}\scrHom(X,Y) &&\longrightarrow \scrHom(v_{\ast}X,
v_{\ast}Y) &\quav & X \text{ and } Y \text{ objects of } E{\quoi},\\
\Xi_v & : \scrHom(v_{!}X,Y) &&\longrightarrow
v_{\ast}\scrHom(X,v^{\ast}Y) &\quav & X\in \ob E, Y\in \ob
E'{\quoi},\\
\Lambda_v & : v_{!}(X\times v^{\ast}Y) &&\longrightarrow
v_{!}(X)\times Y & & X\in \ob E, Y\in \ob E'{\quoi},
\end{alignat*}
where,\pageoriginale in order to define the morphisms $\Xi_v$ and
$\Lambda_v$, we suppose that the inverse-image functor $v^{\ast}$ of
$v$ admits a left adjoint $v_{!}$.

\subsubsection{}\etiquette[10.5.1]{IV.subsubsec10.5.1}
\emph{Definition of} $\Psi_v$. Let $X$ be an object of $E$. There is
an adjunction morphism $v^{\ast}v_{\ast}X\rightarrow X$, and hence,
for every object $Z$ of $E'$, a bifunctorial morphism
$v^{\ast}(Z\times v_{\ast}X)\rightarrow (v^{\ast}Z)\times X$. Thus,
for every object $Y$ of $E$, there is a trifunctorial morphism
$$
\Hom_E((v^{\ast}Z)\times X,Y)\longrightarrow
\Hom_E(v^{\ast}(Z\times v_{\ast}X),Y){\quoi}.
$$
By adjunction this gives a trifunctorial morphism
$$
\Hom_{E'}(Z,v_{\ast}\scrHom(X,Y))\longrightarrow \Hom_{E'}(Z,
\scrHom(v_{\ast}X,v_{\ast}Y)){\quoi};
$$
and hence the morphism $\Psi_v$.

\subsubsection{}\etiquette[10.5.2]{IV.subsubsec10.5.2}
\emph{Definition of} $\Lambda_v$. For every object $X$ of $E$, there
is an adjunction morphism $X\rightarrow v^{\ast}v_{!}X$, and hence,
for every object $Y$ of $E'$, a bifunctorial morphism
$X\times v^{\ast}Y\rightarrow v^{\ast}v_{!}(X)\times v^{\ast}Y
\simeq v^{\ast}(v_{!}(X)\times Y)$; by adjunction this gives the
morphism $\Lambda_v$.

\subsubsection{}\etiquette[10.5.3]{IV.subsubsec10.5.3}
\emph{Definition of} $\Xi_v$. Let $X$ be an object of $E$, and let
$Z$ and $Y$ be two objects of $E'$. The morphism
$\Lambda_v:v_{!}(X\times v^{\ast}Z)\rightarrow v_{!}(X)\times Z$
gives a trifunctorial morphism
$\Hom_{E'}(v_{!}(X)\times Z,Y)\rightarrow
\Hom_{E'}(v_{!}(X\times v^{\ast}Z),Y)$; by adjunction this gives a
trifunctorial morphism
$$
\Hom_{E'}(Z,\scrHom(v_{!}X,Y))\longrightarrow
\Hom_{E'}(Z,v_{\ast}\scrHom(X,v^{\ast}Y)){\quoi};
$$
and hence the morphism $\Xi_v$.

\subsubsection{}\etiquette[10.5.4]{IV.subsubsec10.5.4}
\emph{Definition of} $\Phi_v$. The trifunctorial morphism of
(\hyperref[IV.subsubsec10.5.3]{10.5.3})
$$
\Hom_{E'}(v_{!}(Z)\times X,Y)\longrightarrow
\Hom_{E'}(v_{!}(Z\times v^{\ast}X),Y){\quoi},
$$
where $Z$ is an object of $E$ and $X,Y$ are objects of $E'$, gives,
by adjunction, a trifunctorial morphism
$$
\Hom_E(Z,v^{\ast}\scrHom(X,Y))\longrightarrow
\Hom_E(Z,\scrHom(v^{\ast}X,v^{\ast}Y)){\quoi};
$$
and hence a definition of the morphism $\Phi_v$. There is a second
way to define $\Phi_v$. Let $X$, $Y$, and $Z$ be three objects of
$E'$. The functor $v^{\ast}:E'\rightarrow E$ gives a trifunctorial
morphism
$$
\Hom_{E'}(Z\times X,Y)\longrightarrow \Hom_E(v^{\ast}(Z)\times
v^{\ast}(X),v^{\ast}Y){\quoi};
$$
and hence, by adjunction, a trifunctorial morphism
$$
\Hom_{E'}(Z,\scrHom(X,Y))\longrightarrow
\Hom_{E'}(Z,v_{\ast}\scrHom(v^{\ast}X,v^{\ast}Y)){\quoi}.
$$
We\pageoriginale therefore have a bifunctorial morphism
$$
\scrHom(X,Y)\longrightarrow v_{\ast}\scrHom(v^{\ast}X,
v^{\ast}Y){\quoi};
$$
and hence, by adjunction, a morphism which the reader may verify is
indeed the morphism $\Phi_v$ defined above.

\begin{enonce*}{Proposition 10.6}\etiquette[10.6]{IV.prop10.6}
Let $v:E\rightarrow E'$ be a morphism of topoi.
\begin{enumerate}
\renewcommand{\labelenumi}{\theenumi)}
\item The morphism $\Psi_v$ (\hyperref[IV.subsubsec10.5.1]{10.5.1})
is an isomorphism for every object $Y$ of $E$ if and only if the
adjunction morphism $v^{\ast}v_{\ast}X\rightarrow X$ is an
isomorphism. In particular, $\Psi_v$ is an isomorphism for every pair
of objects $(X,Y)$ if and only if the functor
$v_{\ast}:E\rightarrow E'$ is fully faithful, \ie if $v$ is an
embedding of topoi.

\item The following conditions are equivalent:
\begin{enumerate}
\renewcommand{\theenumii}{\roman{enumii}}
\renewcommand{\labelenumii}{(\theenumii)}
\item For every pair $(X,Y)$ of objects of $E'$, the morphism
$\Phi_v$ (\hyperref[IV.subsubsec10.5.4]{10.5.4}) is an isomorphism.

\item The functor $v^{\ast}$ admits a left adjoint $v_{!}$, and for
every object $X$ of $E$ and every object $Y$ of $E'$, the morphism
$\Xi_v$ (\hyperref[IV.subsubsec10.5.3]{10.5.3}) is an isomorphism.

\item The functor $v^{\ast}$ admits a left adjoint $v_{!}$, and for
every object $X$ of $E$ and every object $Y$ of $E'$, the morphism
$\Lambda_v$ (\hyperref[IV.subsubsec10.5.2]{10.5.2}) is an isomorphism.
\end{enumerate}
\end{enumerate}
\end{enonce*}

\setcounter{subsection}{6} \setcounter{subsubsection}{0}
\subsubsection{}\etiquette[10.6.1]{IV.subsubsec10.6.1}
The first assertion follows immediately from
(\hyperref[IV.subsubsec10.5.1]{10.5.1}). It also follows immediately
from (\hyperref[IV.subsubsec10.5.2]{10.5.2}),
(\hyperref[IV.subsubsec10.5.3]{10.5.3}), and
(\hyperref[IV.subsubsec10.5.4]{10.5.4}) that
(iii) $\Leftrightarrow$ (ii) $\Leftrightarrow$ ((i) and the functor
$v^{\ast}$ admits a left adjoint $v_{!}$). It remains to show that
(i) implies that $v^{\ast}$ admits a left adjoint, \ie
(\hyperref[IV.rem1.8]{1.8}) that (i) implies that $v^{\ast}$ commutes
with small products. Let $e'$ be the final object of $E'$ and let $I$
be a small set. Since $v^{\ast}$ is left exact, $v^{\ast}(e')$ is a
final object of $E$, and since $v^{\ast}$ commutes with inductive
limits, $v^{\ast}(\amalg_I e')\simeq \amalg_I v^{\ast}(e')$. For
every object $Y$ of $E'$, we have
$\scrHom(\amalg_I e',Y)\simeq \Pi_I\scrHom(e',Y)\simeq \Pi_IY$ and
$\scrHom(\amalg_I v^{\ast}(e'),v^{\ast}Y)\simeq
\Pi_I\scrHom(v^{\ast}(e'),v^{\ast}Y)\simeq
\Pi_Iv^{\ast}Y$.\pageoriginale The functorial morphism $\Phi_v$
therefore induces an isomorphism
$v^{\ast}\Pi_IY\simeq \Pi_Iv^{\ast}Y$, which one checks is the
canonical isomorphism. \cqfd

\begin{enonce*}{Corollary 10.7}\etiquette[10.7]{IV.cor10.7}
Let $E$ be a topos, let $Z$ be an object of $E$, let
$j_Z:E_{/Z}\rightarrow E$ be the localization morphism
(\hyperref[IV.subsec5.2]{5.2}), and let $X,Y$ be two objects of $E$.
\begin{enumerate}
\renewcommand{\labelenumi}{\theenumi)}
\item There is a bifunctorial isomorphism
$$
j^{\ast}_Z\scrHom(X,Y)\simeq \scrHom(j^{\ast}_ZX,
j^{\ast}_ZY){\quoi}.
$$

\item Let $X'$ be an object of $E_{/Z}$. There is a bifunctorial
isomorphism
$$
\scrHom_E(j_{Z!}X',Y)\simeq
j_{Z\ast}\scrHom_{E/Z}(X',j^{\ast}_ZY){\quoi}.
$$

\item Let $Y'$ be an object of $E/Z$. When $Z$ is open in $E$, there
is a bifunctorial isomorphism
$$
j_{Z\ast}\scrHom(X',Y')\simeq \scrHom(j_{Z\ast}X',
j_{Z\ast}Y'){\quoi}.
$$
\end{enumerate}
\end{enonce*}

Assertions 1) and 2) are proved by noting that the morphism
$\Lambda_{j_Z}$ is an isomorphism. For assertion 3), it suffices to
note that $j_{Z\ast}$ is fully faithful, since $j_{Z!}$ is fully
faithful (\hyperrefext[I][I.rem5.7]{5.7}. a)).

\begin{enonce*}{Corollary 10.8}\etiquette[10.8]{IV.cor10.8}
Let $E$ be a topos, let $X$, $Y$, and $Z$ be three objects of $E$,
and let $j_Z:E_{/Z}\rightarrow E$ be the localization morphism.
\begin{enumerate}
\renewcommand{\labelenumi}{\theenumi)}
\item There is a canonical isomorphism
$$
j_{Z\ast}j^{\ast}_ZY\simeq \scrHom(Z,Y){\quoi}.
$$

\item There is a canonical isomorphism
$$
\Hom_E(Z,\scrHom(X,Y))\simeq
\Hom_{E_{/Z}}(j^{\ast}_ZX,j^{\ast}_ZY){\quoi}.
$$
\end{enumerate}
\end{enonce*}

Let $e_Z$ be the final object of $E_{/Z}$ (\ie the object
$\id_Z:Z\rightarrow Z$). It follows from
\eqhyperref[IV.eq10.2.1]{10.2.1} that there is a canonical
isomorphism
$$
\scrHom_{E/Z}(e_Z,j^{\ast}_ZY)\simeq j^{\ast}_ZY{\quoi};
$$
hence, by (\hyperref[IV.cor10.7]{10.7} 2), an isomorphism
$$
j_{Z\ast}j^{\ast}_ZY\simeq \scrHom(j_{Z!}e_Z,Y){\quoi.}
$$
Since $j_{Z!}e_Z=Z$, this proves 1). Let us prove 2). From
(\hyperref[IV.cor10.7]{10.7}) we obtain a canonical\pageoriginale
isomorphism
$$
\Hom_{E_{/Z}}(e_Z,j^{\ast}_Z\scrHom(X,Y))\simeq
\Hom_{E_{/Z}}(j^{\ast}_ZX,j^{\ast}_ZY){\quoi};
$$
hence, by adjunction on the first member,
$$
\Hom_E(Z,\scrHom(X,Y))\simeq
\Hom_{E_{/Z}}(j^{\ast}_ZX,j^{\ast}_ZY){\quoi}.
$$


% END INCLUDED SGA4_English_current_buildable_source_20260906/source/04/drafts/SGA4_IV_LINES_8246_8573_English.texfrag


% BEGIN INCLUDED SGA4_English_current_buildable_source_20260906/source/04/drafts/SGA4_IV_LINES_8574_8831_English.texfrag
\section{Ringed topoi, localization in ringed topoi}\etiquette[11]{IV.sec11}

\setcounter{subsection}{1} \setcounter{subsubsection}{0}
\subsubsection{}\etiquette[11.1.1]{IV.subsubsec11.1.1}
Let $𝒰$ be a universe. A $𝒰$-\emph{ringed topos} is a pair
$(E,A)$ where $E$ is a $𝒰$-topos and $A$ is an object endowed with
a ring structure. A $𝒰$-ringed site is a pair $(C,A)$ where $C$ is
a $𝒰$-site and $A$ is a $𝒰$-sheaf of rings on $C$. We omit mention
of the universe when the context leaves no room for confusion. To a
ringed site $(C,A)$ is associated the ringed topos $(C^{\sim},A)$.
To a ringed topos $(E,A)$ is associated the ringed site consisting
of the site $E$ and the sheaf of rings represented by $A$.

\subsubsection{}\etiquette[11.1.2]{IV.subsubsec11.1.2}
Let $(E,A)$ be a ringed topos. We denote by ${}_AE$ (\resp $E_A$)
the category of sheaves of unital left (\resp right) $A$-modules
(we will also write $A$-Modules). The category ${}_AE$ (\resp
$E_A$) is an abelian category
(\hyperrefext[II][II.prop6.7]{6.7}). Let $M$ and $N$ be two sheaves
of left (\resp right) $A$-modules. The commutative group of
$A$-Module morphisms from $M$ to $N$ is denoted by $\Hom_A(M,N)$.

\subsubsection{}\etiquette[11.1.3]{IV.subsubsec11.1.3}
Let $A$, $B$, and $C$ be three rings in a topos $E$, let $M$ be a
sheaf of $A-B$-bimodules, and let $N$ be a sheaf of
$A-C$-bimodules. Let $e$ be a final object of $E$, and put
$\Hom(e,B)=\Gamma(B)$ and $\Hom(e,C)=\Gamma(C)$. The
$A-B$-bimodule structure on $M$ gives a ring homomorphism from
$\Gamma(B)^{\circ}$ to the ring of endomorphisms of the $A$-module
$M$; similarly, the $A-C$-bimodule structure on $N$ gives a ring
homomorphism from $\Gamma(C)^{\circ}$ to the ring of endomorphisms
of the $A$-module $N$. By functoriality, we deduce a
$\Gamma(B)-\Gamma(C)$-bimodule structure on the commutative group
$\Hom_A(M,N)$. In particular, when $A$ is a sheaf of commutative
rings, the group $\Hom_A(M,N)$ is canonically endowed with a
$\Gamma(A)$-module structure.

\subsubsection{}\etiquette[11.1.4]{IV.subsubsec11.1.4}
Let\pageoriginale $E$ be a topos, let $A$ and $B$ be two rings in
$E$, let $u:A\rightarrow B$ be a morphism of sheaves of rings, and
let $M$ be a $B$-module (a left module, to fix ideas). The sheaf $M$
may be regarded as a sheaf of $A$-modules by means of $u$: for every
object $X$ of $E$, $M(X)$ is endowed with the $A(X)$-module
structure obtained from its $B(X)$-module structure and the homomorphism
$u(X):A(X)\rightarrow B(X)$ by restriction of scalars. We thus
obtain a functor ``\emph{restriction of scalars along} $u$''
$$
\Res(u):{}_BE\longrightarrow {}_AE{\quoi}.
$$

\subsubsection{}\etiquette[11.1.5]{IV.subsubsec11.1.5}
In particular, when $A=\ZZ$ (the constant sheaf
$\sisi{\underline{\underline{\ZZ}}}{\underline{\ZZ}}$, that is, the
sheaf associated to the constant presheaf $\ZZ$) and when
$u:\sisi{\underline{\underline{\ZZ}}}{\underline{\ZZ}}\rightarrow B$
is the unique canonical morphism, we obtain a functor from ${}_BE$
to the category
${}_{\sisi{\underline{\underline{\ZZ}}}{\underline{\ZZ}}}E$, which
is none other than the category of abelian sheaves, denoted by
$E_{Ab}$. This functor is called the \emph{underlying abelian-sheaf
functor}.

\begin{enonce*}{Proposition 11.1.6}\etiquette[11.1.6]{IV.prop11.1.6}
The restriction-of-scalars functor along $u$ commutes with inductive
and projective limits. It is conservative.
\end{enonce*}

The proposition is true for the restriction-of-scalars functor for
ordinary modules, that is, when $E$ is the punctual topos. It is
therefore true when $E$ is the topos of presheaves on a small site.
It then follows from the determination of inductive and projective
limits by means of the associated-sheaf functor
(\hyperrefext[II][II.prop6.4]{6.4}) that the proposition is true in
general.

\setcounter{subsection}{2} \setcounter{subsubsection}{0}
\subsubsection{}\etiquette[11.2.1]{IV.subsubsec11.2.1}
Let $(E,A)$ be a ringed topos, let $X$ be an object of $E$, and let
$j_X:E_{/X}\rightarrow E$ be the localization morphism
\sisi{(\hyperref[IV.sec8]{8})}{\hyperref[IV.subsec5.2]{5.2}}.
By \hyperrefext[III][III.prop1.7]{1.7} or
\hyperref[IV.subsubsec3.1.2]{3.1.2}, the sheaf $j^{\ast}_XA$ is
canonically endowed with a ring structure. The sheaf of rings
$j^{\ast}_XA$ is most often denoted by $A|X$, or, by abuse, simply
by $A$. Unless otherwise stated, the topos $E_{/X}$ will be ringed
by $A|X$. The sheaf $j_{X\ast}j^{\ast}_XA$ is canonically endowed
with a ring structure. The adjunction morphism
$A\rightarrow j_{X\ast}j^{\ast}_XA$ is a morphism of sheaves of
rings.

\subsubsection{}\etiquette[11.2.2]{IV.subsubsec11.2.2}
Let\pageoriginale, in addition, $M$ be an $A$-module (a left module,
to fix ideas). The sheaf $j^{\ast}_XM$ is endowed with an
$A|X$-module structure (\hyperrefext[III][III.prop1.7]{1.7}; or
\hyperref[IV.subsubsec3.1.2]{3.1.2}); hence there is a functor
$$
j^{\ast}_X:{}_AE\longrightarrow {}_{A|X}E_{/X}
$$
called the restriction functor to $E_{/X}$, or also the restriction
functor to $X$. The functor $j^{\ast}_X$ commutes with inductive and
projective limits (\loccit). In particular, it is exact. Now let $N$
be an $A|X$-Module. The sheaf $j_{X\ast}N$ is a sheaf of
$j_{X\ast}A|X$-Modules (\hyperrefext[III][III.prop1.7]{1.7} or
\hyperref[IV.subsubsec3.1.2]{3.1.2}); hence, by restriction of
scalars along the adjunction morphism
$A\rightarrow j_{X\ast}(A|X)$
(\hyperref[IV.subsubsec11.1.4]{11.1.4}), it is an $A$-Module, again
denoted by $j_{X\ast}N$. We have therefore defined a functor
$$
j_{X\ast}:{}_{A|X}E_{/X}\longrightarrow {}_AE
$$
which commutes with projective limits
(\hyperref[IV.prop11.1.6]{11.1.6} and
\hyperrefext[III][III.prop1.7]{1.7}). For every $A$-Module $M$, the
adjunction morphism $M\rightarrow j_{X\ast}j^{\ast}_XM$ is a
morphism of $A$-Modules. For every $A|X$-Module $N$ and every
$A$-Module $M$, the adjunction morphism
$M\rightarrow j_{X\ast}j^{\ast}_XM$ defines a morphism,
bifunctorial in $M$ and $N$,
\begin{equation*}
\Hom_{A|X}(j^{\ast}_XM,N)\longrightarrow
\Hom_A(M,j_{X\ast}N){\quoi}.\tag{11.2.2.1}
\end{equation*}\etiquette[11.2.2.1]{IV.eq11.2.2.1}
The latter morphism is an isomorphism; that is, the functors
$j^{\ast}_X$ and $j_{X\ast}$ for $A$-Modules are adjoint
(III \hyperref[III.prop1.7]{1.7}).

\subsubsection{}\etiquette[11.2.3]{IV.subsubsec11.2.3}
The functors $j^{\ast}_X$ and $j_{X\ast}$ for Modules commute with
the ``underlying set'' functor
(\hyperrefext[III][III.prop1.7]{1.7} 4)). They therefore commute with
the restriction-of-scalars functors and, in particular, with the
underlying abelian-sheaf functor.

\setcounter{subsection}{3} \setcounter{subsubsection}{0}
\begin{enonce*}{Proposition 11.3.1}\etiquette[11.3.1]{IV.prop11.3.1}
Let $(E,A)$ be a ringed topos and let $X$ be an object of $E$. The
functor
$$
j^{\ast}_X:{}_{A}E\longrightarrow {}_{A|X}E_{/X}
$$
admits a left adjoint, denoted by
$j_{X!}:{}_{A|X}E_{/X}\rightarrow {}_{A}E$ and called extension by
zero. The extension-by-zero functor is exact and faithful, and
commutes with inductive limits. The functors $j_{X!}$ for
Modules\pageoriginale commute with restriction-of-scalars functors
and, in particular, they commute with the underlying abelian-sheaf
functor.
\end{enonce*}

The existence of the functor $j_{X!}$ follows from
\hyperrefext[III][III.prop1.7]{1.7}. Since $j_{X!}$ is a left
adjoint, it commutes with inductive limits
(\hyperrefext[I][I.sec2]{2}). To prove the other assertions, first
suppose that $E$ is the topos of presheaves of sets on a small
category $C$ containing the object $X$. We then know that $E_{/X}$
is equivalent to the category of presheaves on $C_{/X}$ and that,
under this equivalence, the functor
$j^{\ast}_X:E\rightarrow E_{/X}$ is simply composition with the
forgetful functor $C_{/X}\rightarrow C$
(\hyperrefext[I][I.prop5.11]{5.11}). It then follows immediately
from the explicit construction of $j_{X!}$
(\hyperrefext[I][I.prop5.1]{5.1}) that, for every $A|X$-Module $N$
and every object $Y$ of $C$, one has
$$
j_{X!}N(Y)=\bigoplus_{u\in \Hom_C(Y,X)}N(u){\quoi};
$$
hence $j_{X!}$ is exact, and extension by zero commutes with
restriction-of-scalars functors in this case. In the general case,
we may suppose that $E$ is the topos of sheaves on a small site $C$
and that $X$ comes from an object of $C$
(\hyperref[IV.cor1.2.1]{1.2.1}). The topos $E_{/X}$ is then
equivalent to the topos of sheaves on $C_{/X}$ (endowed with the
induced topology) (\hyperrefext[III][III.prop5.4]{5.4}), and the
topos morphism $j_X:E_{/X}\rightarrow E$ comes from the forgetful
functor $C_{/X}\rightarrow C$, which is continuous and cocontinuous
(\hyperrefext[III][III.prop5.2]{5.2}). It then follows from
\hyperrefext[III][III.prop1.7]{1.7} that extension by zero for
sheaves is obtained by composing extension by zero for presheaves
with the associated-sheaf functor; hence the assertion of exactness
and the compatibility with restrictions of scalars.

To prove faithfulness, it is equivalent to show that, for every
$A|X$-Module $N$, the adjunction morphism
$$
\ad_N:N\longrightarrow j^{\ast}_X j_{X!}N
$$
is a monomorphism. Now, denoting by $j_{\widehat{X}!}$ the
extension-by-zero functor for presheaves and by $\underline{a}$ the
``associated sheaf'' functor, we have a commutative diagram
\[
\xymatrix@C=2.75cm@R=1.5cm{ N \ar[r]^{\ad_{N}}
\ar[rd]_{\ad_{\hat{N}}} & j^{\ast}_{X} \underline{a}
j_{\widehat{X}!} N \\
                   & \underline{a} j^{\ast}_X j_{\widehat X!} N,\ar[u]
}
\]
where\pageoriginale $\ad_{\widehat{N}}$ is the adjunction morphism
for presheaves. The morphism $\ad_{\widehat{N}}$ is a monomorphism
by what precedes, and therefore
$\underline{a}(\ad_{\widehat{N}})$ is a monomorphism. Moreover,
since the forgetful functor $C_{/X}\rightarrow C$ is continuous and
cocontinuous, the canonical morphism
$\underline{a}j^{\ast}_X\rightarrow j^{\ast}_X\underline{a}$ is an
isomorphism (\hyperrefext[III][III.prop2.3]{2.3}). Consequently,
$\ad_N$ is a monomorphism.

\begin{enonce*}[remark]{Remark 11.3.2}\etiquette[11.3.2]{IV.rem11.3.2}
Let $N$ be an $A|X$-Module. The sheaf of sets underlying $j_{X!}N$
is not, in general, isomorphic to the sheaf obtained by extending
the sheaf of sets underlying $N$ by the empty set
(\hyperref[sga4.exp.III.terminology.5.3]{III 5.3} and \hyperref[IV.subsec5.2]{5.2}). One must
therefore take care not to confuse the \emph{extension-by-zero}
functor denoted by
$j_{X!}:{}_{A|X}E_{/X}\rightarrow {}_{A}E$ in
\hyperref[IV.prop11.3.1]{11.3.1} with the extension-by-the-empty-set
functor, also denoted by $j_{X!}:E_{/X}\rightarrow E$, in III \hyperref[III.subsec5.3]{5.3}
and IV \hyperref[IV.subsec5.2]{5.2}. In most cases encountered in practice, the abuse of
notation just noted causes no confusion. When confusion is
nevertheless possible, we will use the notations $j^{\ab}_{X!}$ and
$j^{\ens}_{X!}$ for the extension-by-zero and
extension-by-the-empty-set functors, respectively.
\end{enonce*}

\begin{enonce*}{Proposition 11.3.3}\etiquette[11.3.3]{IV.prop11.3.3}
Let $(E,A)$ be a ringed topos and let $X$ be an object of $E$. The
$A$-Module $j_{X!}(A|X)$, most often denoted by $A_X$ or $A_{X,E}$,
is the free $A$-Module generated by $X$
{\rm (\hyperrefext[II][II.prop6.5]{6.5})}; that is, for every
$A$-Module $M$, there is a canonical isomorphism, functorial in $M$:
$$
\Hom_E(X,M)\xleftarrow{\sim} \Hom_A (A_X, M){\quoi}.
$$
Let $(X_i)_{i\in I}$ be a topologically generating family of $E$
{\rm (\hyperrefext[II][II.def3.0.1]{3.0.1}).} The family
$(A_{X_i})_{i\in I}$ is a generating family of the category of
$A$-Modules.
\end{enonce*}

Let $e_X$ be the final object of the topos $E_{/X}$
$(e_X=X\xrightarrow{\id}X)$. There is an isomorphism
$$
\Hom_{A|X}(A|X,j^{\ast}_XM)\xrightarrow{\sim}
\Hom_{E_{/X}}(e_X,j^{\ast}_XM){\quoi};
$$
deduced\pageoriginale from the unit section $e_X\rightarrow A|X$;
hence, by adjunction, there is an isomorphism
$$
\Hom_A(j^{\ab}_{X!}(A|X),
M)\xrightarrow{\sim}\Hom_E(j^{\ens}_{X!}(e_X),M){\quoi}.
$$
Since $j^{\ens}_{X!}(e_X)=X$, we obtain the asserted isomorphism.
The last assertion follows from \hyperrefext[II][II.cor6.6]{6.6}.

\begin{enonce*}[remark]{Remark 11.3.4}\etiquette[11.3.4]{IV.rem11.3.4}
Let $A$ and $B$ be two sheaves of rings on a topos $E$, let $X$ be
an object of $E$, and let $N$ be an $A|X$-$B| X$-bimodule. The
abelian sheaf obtained by extending $N$ by zero is canonically
endowed with an $A$-$B$-bimodule structure, as follows immediately
from its explicit description (\hyperref[IV.prop11.3.1]{11.3.1}).
In particular, the free $A$-module so obtained, $A_X$, is an
$A$-bimodule.
\end{enonce*}

\begin{enonce*}[remark]{Exercise 11.3.5}\etiquette[11.3.5]{IV.exer11.3.5}
Let $E$ be a topos, let $X$ be an object of $E$, let
$x:P\rightarrow E$ be a point of $E$ (\hyperref[IV.def6.1]{6.1}),
and let $(x_i)_{i\in I}$ be the family of points of $E_{/X}$ above
$x$ (in bijective correspondence with the fiber $X_x$
(\hyperref[IV.eq6.7.2]{6.7.2})). Show that, for every abelian sheaf
$M$ on $E_{/X}$, the fiber $(j_{X!}M)_x$ is canonically isomorphic
to $\oplus_i M_{x_i}$.
\end{enonce*}

% END INCLUDED SGA4_English_current_buildable_source_20260906/source/04/drafts/SGA4_IV_LINES_8574_8831_English.texfrag


% BEGIN INCLUDED SGA4_English_current_buildable_source_20260906/source/04/drafts/SGA4_IV_LINES_8832_9247_English.texfrag
\section{Operations on modules}\etiquette[12]{IV.sec12}

\begin{enonce*}{Proposition 12.1}\etiquette[12.1]{IV.prop12.1}
Let $(E,A)$ be a ringed topos, and let $M$ and $N$ be two left
$A$-modules (\resp right $A$-modules). The functor on $E$ which
assigns to every object $X$ of $E$ the commutative group
$\Hom_A(M,\scrHom_E(X,N))$ is representable by an abelian sheaf
denoted by $\scrHom_A(M,N)$ (or sometimes $\scrHom(M,N)$ when no
confusion can result). For every object $X$ of $E$, there is a
canonical isomorphism
\begin{equation*}
\scrHom_A(M,N)(X)\simeq \Hom_{A| X}(j^{\ast}_XM,
j^{\ast}_XN){\quoi}.\tag{12.1.1}
\end{equation*}\etiquette[12.1.1]{IV.eq12.1.1}
\end{enonce*}

There is a canonical isomorphism
$\scrHom_E(X,N)=j_{X\ast}j^{\ast}_XN$
(\hyperref[IV.cor10.8]{10.8}); consequently, $\scrHom_E(X,N)$ is
endowed, functorially in $X$, with a structure of left $A$-module
(\resp right $A$-module). The functor
$X\longmapsto\scrHom_E(X,N)$ transforms inductive limits in the
argument $X$ into projective limits of $A$-modules
(\hyperref[IV.subsec10.2]{10.2} and
\hyperref[IV.subsubsec11.2.3]{11.2.3}); consequently, the functor
$$
X\longmapsto \Hom_A(M,\scrHom_E(X,N))
$$
transforms inductive limits\pageoriginale in the argument $X$ into
projective limits of commutative groups, and hence into projective
limits of the underlying sets. It is therefore representable
(\hyperref[IV.prop1.4]{1.4} and \hyperref[IV.theo1.2]{1.2}), and the
object representing it is endowed with the structure of an abelian
sheaf. By definition, for every object $X$ of $E$, there is a
canonical isomorphism
\begin{equation*}
\scrHom_A(M,N)(X)=\Hom(X,\scrHom_A(M,N))\simeq
\Hom_A(M,\scrHom(X,N))\tag{12.1.2}
\end{equation*}\etiquette[12.1.2]{IV.eq12.1.2}
and the isomorphism (\hyperref[IV.eq12.1.1]{12.1.1}) follows from
\eqhyperref[IV.eq12.1.2]{12.1.2}, from the isomorphism
$\scrHom(X,N)\approx j_{X\ast}j^{\ast}_XN$
(\hyperref[IV.cor10.8]{10.8}), and from the adjunction formulas of 10.

\setcounter{subsection}{1}
\subsection{}\etiquette[12.2]{IV.subsec12.2}
The abelian sheaf $\scrHom_A(M,N)$ is called the \emph{sheaf of
$A$-module homomorphisms from $M$ to $N$}. It is a bifunctor in $M$
and $N$. It follows from its definition and from
(\hyperref[IV.subsec10.2]{10.2}) that it transforms projective limits
in the argument $N$ (\resp inductive limits in $M$) into projective
limits of abelian sheaves. In particular, it is left exact in both
arguments.

\begin{enonce*}{Proposition 12.3}\etiquette[12.3]{IV.prop12.3}
Let $(E,A)$ be a ringed topos, let $M$ and $N$ be two right
$A$-modules (\resp left $A$-modules), and let $X$ be an object of $E$:
\begin{enumerate}
\renewcommand{\theenumi}{\alph{enumi}}
\renewcommand{\labelenumi}{\theenumi)}
\item There are canonical isomorphisms:
$$
\scrHom_A(M,N)(X)\approx \Hom_A(M, j_{X\ast}j^{\ast}_XN)\simeq
\Hom_{A| X}(j^{\ast}_XM, j^{\ast}_XN)\simeq
\Hom_A(j_{X!}j^{\ast}_XM, N).
$$

\item There is a canonical isomorphism:
$$
\Phi_X:j^{\ast}_X\scrHom_A(M,N)\simeq
\scrHom_{A| X}(j^{\ast}_XM,j^{\ast}_XN).
$$
Let, in addition, $P$ be a right $A|X$-module (\resp a left
$A|X$-module):

\item There is a canonical isomorphism:
$$
j_{X\ast}\scrHom_{A| X}(j^{\ast}_XM,P)\simeq
\scrHom_A(M,j_{X\ast}P).
$$

\item There is a canonical isomorphism:
$$
\Xi_X:\scrHom_A(j_{X!}P,N)\simeq j_{X\ast}\scrHom_{A|X}
(P,j^{\ast}_XN).
$$
\end{enumerate}
\end{enonce*}

\noindent The isomorphisms in a) follow from
\eqhyperref[IV.eq12.1.2]{12.1.2}, from the isomorphism
$\scrHom_E(X,N)\simeq j_{X\ast}j^{\ast}_XN$
(\hyperref[IV.cor10.8]{10.8}), and from the adjunction formulas of 10.

\begin{itemize}
\item[Proof of b).] For every object $Y$ of $E_{/X}$, there is the
following sequence of isomorphisms:

$\Hom(Y,j^{\ast}_X\scrHom_A(M,N))\simeq
\Hom(j_{X!}Y,\scrHom_A(M,N))$ (adjunction for sheaves of sets)
\begin{alignat*}{2}
& \Hom(j_{X!}Y,\scrHom_A(M,N))\simeq
\Hom_A(M,\scrHom_E(j_{X!}Y,N))
&\quad& \hyperref[IV.eq12.1.2]{(12.1.2)}\\
& \Hom_A(M,\scrHom_E(j_{X!}Y,N))\simeq
\Hom_A(M,j_{X\ast}\scrHom_{E_{/X}}(Y,j^{\ast}_XN))
&& (\hyperref[IV.cor10.7]{10.7})\\
& \Hom_A(M,j_{X\ast}\scrHom_{E_{/X}}(Y,j^{\ast}_XN))\simeq
\Hom_{A|X}(j^{\ast}_XM,\scrHom_E(Y,j^{\ast}_XN))
&& \eqhyperref[IV.eq11.2.2.1]{11.2.2.1}\\
& \Hom_{A|X}(j^{\ast}_XM,\scrHom_E(Y,j^{\ast}_XN))\simeq
\Hom(Y,\scrHom_{A|X}(j^{\ast}_XM,j^{\ast}_XN))
&& \hyperref[IV.eq12.1.2]{(12.1.2)}
\end{alignat*}\pageoriginale

\item[Proof of c).] For every object $Y$ of $E$, there is the
following sequence of isomorphisms:
\begin{alignat*}{2}
& \Hom(Y,j_{X\ast}\scrHom_{A|X}(j^{\ast}_XM,P))\simeq
\Hom(j^{\ast}_XY,\scrHom_{A|X}(j^{\ast}_XM,P))
&\quad& \text{(adjunction)}\\
& \Hom(j^{\ast}_XY,\scrHom_{A|X}(j^{\ast}_XM,P))\simeq
\Hom_{A|X}(j^{\ast}_XM,\scrHom_{E_{/X}}(j^{\ast}_XY,P))
&& \hyperref[IV.eq12.1.2]{(12.1.2)}\\
& \Hom_{A|X}(j^{\ast}_XM,\scrHom_{E_{/X}}(j^{\ast}_XY,P))\simeq
\Hom_A(M,j_{X\ast}\scrHom_{E_{/X}}(j^{\ast}_XY,P))
&& \eqhyperref[IV.eq11.2.2.1]{11.2.2.1}\\
& \Hom_A(M,j_{X\ast}\scrHom_{E_{/X}}(j^{\ast}_XY,P))\simeq
\Hom_A(M,\scrHom_E(Y,j_{X\ast}P))
&& (\hyperref[IV.prop10.3]{10.3})\\
& \Hom_A(M,\scrHom_E(Y,j_{X\ast}P))\simeq
\Hom(Y,\scrHom_A(M,j_{X\ast}P))
&& \hyperref[IV.eq12.1.2]{(12.1.2)}
\end{alignat*}

\item[Proof of d).] For every object $Y$ of $E$, there is the
following sequence of isomorphisms:
\begin{alignat*}{2}
& \Hom(Y,\scrHom_A(j_{X!}P,N))\simeq
\Hom_A(j_{X!}P,\scrHom_E(Y,N))
&\quad& \hyperref[IV.eq12.1.2]{(12.1.2)}\\
& \Hom_A(j_{X!}P,\scrHom_E(Y,N))\simeq
\Hom_{A|X}(P,j^{\ast}_X\scrHom_E(Y,N))
&& (\hyperref[IV.prop11.3.1]{11.3.1})\\
& \Hom_{A|X}(P,j^{\ast}_X\scrHom_E(Y,N))\simeq
\Hom_{A|X}(P,\scrHom_{E_{/X}}(j^{\ast}_XY,j^{\ast}_XN))
&& (\hyperref[IV.cor10.7]{10.7})\\
& \Hom_{A|X}(P,\scrHom_{E_{/X}}(j^{\ast}_XY,j^{\ast}_XN))\simeq
\Hom(j^{\ast}_XY,\scrHom_{A|X}(P,j^{\ast}_XN))
&& \hyperref[IV.eq12.1.2]{(12.1.2)}\\
& \Hom(j^{\ast}_XY,\scrHom_{A|X}(P,j^{\ast}_XN))\simeq
\Hom(Y,j_{X\ast}\scrHom_{A|X}(P,j^{\ast}_XN))
&& \text{(adjunction)}
\end{alignat*}
\end{itemize}

\begin{enonce*}{Corollary 12.4}\etiquette[12.4]{IV.cor12.4}
Let $C$ be a $𝒰$-site, let $A'$ be a $𝒰$-presheaf of rings on
$C$, let $M$ be a $𝒰$-presheaf of $A'$-modules (left modules, to fix
ideas), and let $N$ be a $𝒰$-sheaf of $A'$-modules (left modules).
Denote by $A$ the sheaf associated with $A'$, so that the sheaves $N$
and $\underline{a}M$ (the sheaf associated with $M$) are $A$-modules.
The presheaf
$X\longmapsto\Hom_{A'|X}(j^{\ast}_XM,j^{\ast}_XN)$ is a
$𝒰$-abelian sheaf canonically isomorphic to
$\scrHom_A(\underline{a}M,N)$.
\end{enonce*}

Indeed, it follows from \hyperrefext[III][III.prop5.5]{5.5} and
\hyperrefext[III][III.prop2.3]{2.3} that the localization functor to
$C_{/X}$ commutes with the associated-sheaf functors (for $C$ and
$C_{/X}$). Consequently, there are isomorphisms, functorial in the
variable object $X$ of $C$:
\begin{align*}
\Hom_{A'|X}(j^{\ast}_XM,j^{\ast}_XN)
&\simeq \Hom_{A|X}(\underline{a}j^{\ast}_XM,j^{\ast}_XN)\\
&\simeq \Hom_{A|X}(j^{\ast}_X\underline{a}M,j^{\ast}_XN)
\simeq \scrHom_A(\underline{a}M,N)(X){\quoi}.
\end{align*}

\begin{enonce*}{Corollary 12.5}\etiquette[12.5]{IV.cor12.5}
Let\pageoriginale $E$ be a topos, let $A$, $B$, and $C$ be three
sheaves of rings on $E$, let $M$ be an $A$-$B$-bimodule, and let $N$
be an $A$-$C$-bimodule. The abelian sheaf $\scrHom_A(M,N)$ is
canonically endowed with a $B$-$C$-bimodule structure. In particular,
when $A$ is a sheaf of commutative rings and $M$ and $N$ are
$A$-modules, $\scrHom_A(M,N)$ is canonically endowed with an
$A$-module structure. The canonical isomorphisms b), c), and d) of
$\hyperref[IV.prop12.3]{12.3}$ are isomorphisms of bimodules.
\end{enonce*}

We shall confine ourselves to giving indications. For every object
$X$ of $E$, one has
$\scrHom_A(M,N)(X)\simeq\Hom_{A|X}(j^{\ast}_XM,j^{\ast}_XN)$
(\hyperref[IV.prop12.3]{12.3}). Consequently, the commutative group\\
$\scrHom_A(M,N)(X)$ is canonically endowed with a
$B(X)$-$C(X)$-bimodule structure
(\hyperref[IV.subsubsec11.1.3]{11.1.3}), which is readily checked to
be functorial in $X$. The other assertions are left to the reader.

\begin{enonce*}{Corollary 12.6}\etiquette[12.6]{IV.cor12.6}
Let $(E,A)$ be a ringed topos, let $X$ be an object of $E$, and let
$N$ be an $A$-module (a left module, to fix ideas). There are canonical
isomorphisms of $A$-modules:
$$
\scrHom_E(X,N)\simeq j_{X\ast}j^{\ast}_XN\simeq
\scrHom_A(A_X,N){\quoi};
$$
the $A$-module structure on $\scrHom_A(A_X,N)$ comes from the
bimodule structure on $A_X$ $(\hyperref[IV.rem11.3.4]{11.3.4})$.
\end{enonce*}

The first isomorphism follows from \hyperref[IV.cor10.7]{10.7}, and
the second from the isomorphism of $A$-modules
$N\simeq\scrHom_A(A,N)$ and from \hyperref[IV.prop12.3]{12.3}.

\begin{enonce*}{Proposition 12.7}\etiquette[12.7]{IV.prop12.7}
Let $(E,A)$ be a ringed topos, let $M$ be a right $A$-module, and let
$N$ be a left $A$-module. The functor which assigns to every abelian
sheaf $P$ the commutative group
$\Hom_A(M,\scrHom_{\sisi{\underline{\underline{\ZZ}}}{\underline{\ZZ}}}(N,P))$
is representable by an abelian sheaf denoted by $M\otimes_A N$ and
called the tensor product over $A$ of $M$ and $N$.
\end{enonce*}

There is a canonical injection of
$\Hom_A(M,\scrHom_{\sisi{\underline{\underline{\ZZ}}}{\underline{\ZZ}}}(N,P))$
into the set
$\Hom(M,\scrHom(N,P))\simeq\Hom(M\times N,P)$; returning to the
definitions, one sees immediately that a morphism $f$ of sheaves of
sets from $M\times N$ to $P$ comes from an element of
$\Hom_A(M,\scrHom_Z(N,P))$ if and only if, for every object $X$ of
$E$, the map $f(X):M(X)\times N(X)\rightarrow P(X)$ is an
$A(X)$-bilinear map\pageoriginale from $M(X)\times N(X)$ to $P(X)$.
We call morphisms of sheaves of sets from $M\times N$ to $P$ having
this property $A$-bilinear morphisms. Thus, the problem is to represent
the functor of $A$-bilinear morphisms from $M\times N$ to $P$. To do
this, one proceeds as in the ordinary case, \ie the case where $E$ is
the punctual topos. Consider the eight morphisms:
\begin{itemize}
\item[(i):] $M\times M\times N\longrightarrow M\times N$\qquad
$1\leq i\leq 3$

\item[(i):] $M\times N\times N\longrightarrow M\times N$\qquad
$4\leq i\leq 6$

\item[(i):] $M\times A\times N\longrightarrow M\times N$\qquad
$7\leq i\leq 8$
\end{itemize}
defined by the formulas:
\begin{tabbing}
\quad\= (1) \= $(m_1, m_2, n)$ \= $\longmapsto$ \= $(m_1+m_2, n)$\\
\> (2) \> $(m_1, m_2, n)$ \> $\longmapsto$ \> $(m_1, n)$\\
\> (3) \> $(m_1, m_2, n)$ \> $\longmapsto$ \> $(m_2, n)$\\
\> (4) \> $(m, n_1, n_2)$ \> $\longmapsto$ \> $(m, n_1+n_2)$\\
\> (5) \> $(m, n_1, n_2)$ \> $\longmapsto$ \> $(m, n_1)$\\
\> (6) \> $(m, n_1, n_2)$ \> $\longmapsto$ \> $(m, n_2)$\\
\> (7) \> $(m, a, n)$ \> $\longmapsto$ \> $(ma, n)$\\
\> (8) \> $(m, a, n)$ \> $\longmapsto$ \> $(m, an)$;
\end{tabbing}
here formula (i) describes the map (i)$(X)$ for the variable objects
$X$ of $E$. Let $L$ denote the ``free generated abelian sheaf''
functor. It is clear that the largest quotient (in the sense of
abelian sheaves) of $L(M\times N)$ which equalizes $L(1)$ with
$L(2)+L(3)$, $L(4)$ with $L(5)+L(6)$, and $L(7)$ with $L(8)$
represents the functor of $A$-bilinear morphisms from $M\times N$ to
$P$.

\setcounter{subsection}{7}





\subsection{}\etiquette[12.8]{IV.subsec12.8}
One observes that the functor
$P\longmapsto
\Hom_A(N,\scrHom_{\sisi{\underline{\underline{\ZZ}}}{\underline{\ZZ}}}(M,P))$
is also canonically isomorphic to the functor of $A$-bilinear maps
from $M\times N$ to $P$. Consequently, the tensor product
$M\otimes_A N$ also represents the functor
$\Hom_A(N,\scrHom_{\sisi{\underline{\underline{\ZZ}}}{\underline{\ZZ}}}(M,P))$.
Thus there are isomorphisms, functorial in all the arguments:
\hypertarget{sga4.exp.IV.equation.12.8.1}{}\begin{align*}%\etiquette[12.8.1]{IV.eq12.8.1}
\Hom_Z(M\otimes_A N,P)
&\simeq \Hom_A(M,\scrHom_{\sisi{\underline{\underline{\ZZ}}}{\underline{\ZZ}}}(N,P)){\quoi},\tag{12.8.1}\\
\Hom_Z(M\otimes_A N,P)
&\simeq \Hom_A(N,\scrHom_{\sisi{\underline{\underline{\ZZ}}}{\underline{\ZZ}}}(M,P)){\quoi}.\tag{12.8.2}
\end{align*}\hypertarget{sga4.exp.IV.equation.12.8.2}{}

\subsection{}\etiquette[12.9]{IV.subsec12.9}
It\pageoriginale follows from \hyperref[IV.subsec12.8]{(12.8)}, or alternatively from \hyperlink{sga4.exp.IV.equation.12.8.1}{(12.8.1)}
and (\hyperref[IV.subsec12.2]{12.2}), that the functor $\otimes_A$
commutes with inductive limits in both arguments.

\begin{enonce*}{Proposition 12.10}\etiquette[12.10]{IV.prop12.10}
Let $C$ be a $𝒰$-site, let $A'$ be a presheaf of rings on $C$, let
$M'$ (\resp $N'$) be a right (\resp left) $A'$-module, and let $A$,
$M$, and $N$ be the sheaves associated with $A'$, $M'$, and $N'$
respectively. The sheaf associated with the presheaf
$X\longmapsto M'(X)\otimes_{A'(X)}N'(X)$ $(X\in\ob C)$ is
canonically isomorphic to the sheaf $M\otimes_A N$.
\end{enonce*}

The presheaf $X\longmapsto M'(X)\otimes_{A'(X)}N'(X)$ can be
constructed from the presheaves $M'$, $N'$, and $A'$ by the operations
specified in the proof of \hyperref[IV.prop12.7]{12.7}. Since the
``associated sheaf'' functor commutes with finite projective and
inductive limits and with the ``free generated abelian object''
functors, formation of the tensor product commutes with the
``associated sheaf'' functor.

\begin{enonce*}{Proposition 12.11}\etiquette[12.11]{IV.prop12.11}
Let $(E,A)$ be a ringed topos, let $M$ be a right $A$-module, let $N$
be a left $A$-module, and let $X$ be an object of $E$.
\begin{enumerate}
\renewcommand{\theenumi}{\alph{enumi}}
\renewcommand{\labelenumi}{\theenumi)}
\item There is a canonical isomorphism
$$
j^{\ast}_X(M\otimes_A N)\simeq
(j^{\ast}_XM)\otimes_{A|X}(j^{\ast}_XN){\quoi}.
$$
Let, in addition, $P$ be a right $A|X$-module and $Q$ a left
$A|X$-module.

\item There are canonical isomorphisms (projection formulas):
\begin{align*}
j_{X!}(P\otimes_{A|X}j^{\ast}_XN)&\simeq(j_{X!}P)\otimes_A N\\
j_{X!}(j^{\ast}_XM\otimes_{A|X}Q)&\simeq
M\otimes_A(j_{X!}Q){\quoi}.
\end{align*}
\end{enumerate}
\end{enonce*}

\begin{itemize}
\item[Proof of a).] For every abelian sheaf $R$ on $E_{/X}$, there is
the following sequence of isomorphisms, functorial in $R$:
\begin{alignat*}{2}
&\Hom_{\sisi{\underline{\underline{\ZZ}}}{\underline{\ZZ}}}
(j^{\ast}_X(M\otimes_A N),R)\simeq
\Hom_{\sisi{\underline{\underline{\ZZ}}}{\underline{\ZZ}}}
(M\otimes_A N,j_{X\ast}R)
&\quad& \eqhyperref[IV.eq11.2.2.1]{11.2.2.1}\\
&\Hom_{\sisi{\underline{\underline{\ZZ}}}{\underline{\ZZ}}}
(M\otimes_A N,j_{X\ast}R)\simeq
\Hom_A(M,\scrHom_Z(N,j_{X\ast}R))
&& \hyperlink{sga4.exp.IV.equation.12.8.1}{(12.8.1)}\\
&\Hom_A(M,\scrHom_Z(N,j_{X\ast}R))\simeq
\Hom_A(M,j_{X\ast}\scrHom_{\sisi{\underline{\underline{\ZZ}}}{\underline{\ZZ}}}
(j^{\ast}_XN,R))
&& (\hyperref[IV.prop12.3]{12.3})\\
&\Hom_A(M,j_{X\ast}\scrHom_{\sisi{\underline{\underline{\ZZ}}}{\underline{\ZZ}}}
(j^{\ast}_XN,R))\simeq
\Hom_{A|X}(j^{\ast}_XM,
\scrHom_{\sisi{\underline{\underline{\ZZ}}}{\underline{\ZZ}}}
(j^{\ast}_XN,R))
&& \eqhyperref[IV.eq11.2.2.1]{11.2.2.1}\\
&\Hom_{A|X}(j^{\ast}_XM,
\scrHom_{\sisi{\underline{\underline{\ZZ}}}{\underline{\ZZ}}}
(j^{\ast}_XN,R))\simeq
\Hom_{\sisi{\underline{\underline{\ZZ}}}{\underline{\ZZ}}}
(j^{\ast}_XM\otimes_{A|X}j^{\ast}_XN,R)
&& \hyperlink{sga4.exp.IV.equation.12.8.1}{(12.8.1)}
\end{alignat*}

Let us exhibit\pageoriginale the first isomorphism in b). For every
abelian sheaf $R$, there is a sequence of isomorphisms, functorial in
$R$:
\begin{alignat*}{2}
&\Hom_Z(j_{X!}(P\otimes_{A|X}j^{\ast}_XN),R)\simeq
\Hom_Z(P\otimes_{A|X}j^{\ast}_XN,j^{\ast}_XR)
&\quad& (\hyperref[IV.prop11.3.1]{11.3.1})\\
&\Hom_Z(P\otimes_{A|X}j^{\ast}_XN,j^{\ast}_XR)\simeq
\Hom_{\sisi{A/X}{A|X}}(P,\scrHom_Z(j^{\ast}_XN,j^{\ast}_XR))
&& \hyperlink{sga4.exp.IV.equation.12.8.1}{(12.8.1)}\\
&\Hom_{A|X}(P,\scrHom_Z(j^{\ast}_XN,j^{\ast}_XR))\simeq
\Hom_{A|X}(P,j^{\ast}_X\scrHom_Z(N,R))
&& (\hyperref[IV.prop12.3]{12.3})\\
&\Hom_{A|X}(P,j^{\ast}_X\scrHom_Z(N,R))\simeq
\Hom_A(j_{X!}P,\scrHom_Z(N,R))
&& (\hyperref[IV.prop11.3.1]{11.3.1})\\
&\Hom_A(j_{X!}P,\scrHom_Z(N,R))\simeq
\Hom_Z(j_{X!}P\otimes_A N,R)
&& \hyperlink{sga4.exp.IV.equation.12.8.1}{(12.8.1)}
\end{alignat*}

The second isomorphism is obtained similarly using \hyperlink{sga4.exp.IV.equation.12.8.2}{(12.8.2)}.
\end{itemize}

\begin{enonce*}{Corollary 12.12}\etiquette[12.12]{IV.cor12.12}
Let $E$ be a topos, let $A$, $B$, and $C$ be three sheaves of rings
on $E$, let $M$ be a $B$-$A$-bimodule, and let $N$ be an
$A$-$C$-bimodule. The abelian sheaf $M\otimes_A N$ is canonically
endowed with a $B$-$C$-bimodule structure. In particular, when $A$ is
a sheaf of commutative rings and $M$ and $N$ are $A$-modules,
$M\otimes_AN$ is canonically endowed with an $A$-module structure.
The isomorphisms of \hyperref[IV.prop12.11]{12.11} are isomorphisms of
bimodules.
\end{enonce*}

For every object $X$ of $E$, $j^{\ast}_XM$ is a right $A|X$-module
on which the ring $B(X)$ acts on the left, and likewise $j^{\ast}_XN$
is a left $\sisi{A/X}{A|X}$-module on which $C(X)$ acts on the right. By
functoriality, the abelian sheaf
$j^{\ast}_XM\otimes_{\sisi{A/X}{A|X}}j^{\ast}_XN\simeq
j^{\ast}_X(M\otimes_A N)$ is endowed with the structure of a
$B(X)$-$C(X)$-bimodule, and this structure varies functorially in
$X$. Consequently, $M\otimes_A N$ is endowed with a $B$-$C$-bimodule
structure. This $B(X)$-$C(X)$-bimodule structure on
$j^{\ast}_X(M\otimes_A N)$ is reflected in the evident way on the
functor of ``$A$-bilinear morphisms''. One then observes that, when
$A$ is commutative and $M$ and $N$ are $A$-modules, the right and left
$A$-module structures so obtained are ``equal''; consequently,
$M\otimes_A N$ is an $A$-module in this case. The last assertion is
left to the reader.

\begin{enonce*}{Corollary 12.13}\etiquette[12.13]{IV.cor12.13}
Let $(E,A)$ be a ringed topos, let $M$ be a left\pageoriginale
$A$-module (\resp a right $A$-module), and let $X$ be an object of
$E$. With the notation $A_X$ of \hyperref[IV.prop11.3.3]{11.3.3},
there is a canonical isomorphism
$$
A_X\otimes_AM\simeq j_{X!}j^{\ast}_XM
\quad\text{(\resp }M\otimes_AA_X\simeq j_{X!}j^{\ast}_XM\text{)}{\quoi}.
$$
\end{enonce*}

\noindent This follows from the projection formulas
(\hyperref[IV.prop12.11]{12.11}).

\begin{enonce*}{Proposition 12.14}\etiquette[12.14]{IV.prop12.14}
Let $E$ be a topos, let $A$ and $B$ be two sheaves of rings on $E$,
let $M$ be a right $A$-module, let $N$ be an $A$-$B$-bimodule, and let
$P$ be a right $B$-module. There is a canonical isomorphism
$$
\Hom_B(M\otimes_A N,P)\simeq
\Hom_A(M,\scrHom_B(N,P)){\quoi}.
$$
\end{enonce*}

\noindent From \hyperlink{sga4.exp.IV.equation.12.8.1}{(12.8.1)} one obtains a canonical $A$-bilinear morphism
$\scrHom_Z(N,P)\times N\rightarrow P$; hence, by restricting it to
$\scrHom_B(N,P)\times N$ (which is a subsheaf), one obtains an
$A$-bilinear morphism $\scrHom_B(N,P)\times N\rightarrow P$, which is
immediately seen to be $B$-linear in the second factor. Thus there is
a canonical morphism of $B$-modules
$\scrHom_B(N,P)\otimes_A N\rightarrow P$; hence a canonical map,
functorial in $M$:
\begin{equation*}
\Hom_A(M,\scrHom_B(N,P))\longrightarrow
\Hom_B(M\otimes_A N,P){\quoi}.\tag{12.14.1}
\end{equation*}\etiquette[12.14.1]{IV.eq12.14.1}
Let us show that this morphism of functors is an isomorphism. Since
both sides transform inductive limits in $M$ into projective limits,
it is enough to show that \eqhyperref[IV.eq12.14.1]{12.14.1} is an
isomorphism when $M$ ranges over a generating family of the category
$E_A$. It is therefore enough to show that
\eqhyperref[IV.eq12.14.1]{12.14.1} is an isomorphism when $M=A_X$,
where $X$ is an object of $E$. The verification is then immediate.

% END INCLUDED SGA4_English_current_buildable_source_20260906/source/04/drafts/SGA4_IV_LINES_8832_9247_English.texfrag


% BEGIN INCLUDED SGA4_English_current_buildable_source_20260906/source/04/drafts/SGA4_IV_LINES_9248_9593_English.texfrag
\section{Morphisms of ringed topoi}\etiquette[13]{IV.sec13}

\begin{enonce*}{Definition 13.1}\etiquette[13.1]{IV.def13.1}
Let $(E,A)$ and $(E',A')$ be two ringed topoi. A morphism of ringed
topoi $u:(E,A)\rightarrow(E',A')$ is a pair $(m,\theta)$, where
$m:E\rightarrow E'$ is a morphism of topoi
\textup{(\hyperref[IV.def3.1]{3.1})} and $\theta:m^{\ast}A'\rightarrow A$ is
a morphism of Rings.
\end{enonce*}

\setcounter{subsection}{1}
\subsubsection{}\etiquette[13.1.1]{IV.subsubsec13.1.1}
Since the functor $m^{\ast}:E'\rightarrow E$ is left adjoint to the
functor $m_{\ast}:E\rightarrow E'$, giving a morphism of ringed
topoi $(m,\theta):(E,A)\rightarrow(E',A')$ amounts to giving a
morphism of topoi $m:E\rightarrow E'$ and a morphism of sheaves of
rings $\Theta:A'\rightarrow m_{\ast}A$.

\subsection{}\etiquette[13.2]{IV.subsec13.2}
To\pageoriginale a morphism of ringed topoi
$u=(m,\theta):(E,A)\rightarrow(E',A')$, we associate two notable
functors between the categories of Modules:

\subsubsection{The direct-image functor for
  Modules:}\etiquette[13.2.1]{IV.subsubsec13.2.1}
Let $M$ be a left (\resp right) $A$-Module. The object $m_{\ast}M$ is
canonically endowed with an $m_{\ast}A$-Module structure; restriction
of scalars along the canonical morphism
$\Theta:A'\rightarrow m_{\ast}A$ therefore gives an $A'$-Module,
denoted by $u_{\ast}(M)$ and called the \emph{direct image of $M$
under the morphism $u$}.

\subsubsection{The inverse-image functor for
  Modules:}\etiquette[13.2.2]{IV.subsubsec13.2.2}
Let $N$ be a left (\resp right) $A'$-Module. The object $m^{\ast}N$
of $E$ is canonically endowed with the structure of a left (\resp
right) $m^{\ast}A'$-Module. The left $A$-Module
$A\otimes_{m^{\ast}A'}m^{\ast}N$ (\resp the right $A$-Module
$m^{\ast}N\otimes_{m^{\ast}A'}A$), where $A$ is endowed with the
$m^{\ast}A'$-Module structure defined by
$\theta:m^{\ast}A'\rightarrow A$, is denoted by $u^{\ast}N$ and is
called the \emph{inverse image of the Module $N$ under the morphism
of ringed topoi} $u$.

\subsubsection{}\etiquette[13.2.3]{IV.subsubsec13.2.3}
Note that for every $A$-Module $M$, the sheaf of sets and the abelian
sheaf underlying $u_{\ast}M$ is the sheaf $m_{\ast}M$. Thus the
notation $u_{\ast}$ is most often also used for the functor
$m_{\ast}$, the direct image for sheaves of sets. On the other hand,
for an $A'$-Module $N$, the sheaf of sets or abelian sheaf underlying
$u^{\ast}N$ is in general neither equal nor isomorphic to
$m^{\ast}N$ (except when $\theta$ is an isomorphism). One must
therefore distinguish between the inverse image for Modules and the
inverse image for sheaves of sets or abelian sheaves. The notation
$u^{-1}:E'\rightarrow E$ is most often used for the functor
$m^{\ast}$, the inverse image for sheaves of sets. The functor
$u^{-1}$ is called the \emph{``set-theoretic'' inverse-image functor}
under the morphism of ringed topoi $u$, as opposed to the inverse
image ``in the sense of Modules'' $u^{\ast}$.

\begin{enonce*}{Definition 13.3}\etiquette[13.3]{IV.def13.3}
Let $(C,A)$ and $(C',A')$ be two $𝒰$-ringed sites. A morphism of
ringed sites $u:(C,A)\rightarrow(C',A')$ is a pair $(m,\theta)$,
where $m$ is a morphism from the site $C$ to the site $C'$
(\hyperref[IV.subsec4.9]{4.9}) and
$\theta:m^{\ast}A'\rightarrow A$ is\pageoriginale a morphism of
Rings.
\end{enonce*}

\setcounter{subsection}{3} \setcounter{subsubsection}{0}
\subsubsection{}\etiquette[13.3.1]{IV.subsubsec13.3.1}
A morphism of ringed topoi is a morphism of $𝒰$-ringed sites. A
morphism of ringed sites gives rise to a morphism between the
corresponding ringed topoi (\hyperref[IV.subsubsec11.1.1]{11.1.1}
and \hyperref[IV.subsubsec4.9.1]{4.9.1}).

\begin{enonce*}{Proposition 13.4}\etiquette[13.4]{IV.prop13.4}
Let $u:(E,A)\rightarrow(E',A')$ be a morphism of ringed topoi.
\begin{enumerate}
\renewcommand{\theenumi}{\alph{enumi}}
\renewcommand{\labelenumi}{\theenumi)}
\item For a variable left (\resp right) $A$-Module $M$ and a variable
left (\resp right) $A'$-Module $N$, there are canonical
bifunctorial isomorphisms (called adjunction isomorphisms):
\hypertarget{sga4.exp.IV.equation.13.4.1}{}\begin{align*}
\Hom_A(u^{\ast}N,M) &\simeq
\Hom_{A'}(N,u_{\ast}M){\quoi},\tag{13.4.1}\\
u_{\ast}\scrHom_A(u^{\ast}N,M) &\simeq
\scrHom_{A'}(N,u_{\ast}M){\quoi}.\tag{13.4.2}
\end{align*}\hypertarget{sga4.exp.IV.equation.13.4.2}{}

\item For a variable object $X$ of $E'$, there is a canonical
isomorphism, functorial in $X$:
\begin{equation*}
u^{\ast}A'_X\xrightarrow{\sim} A_{u^{-1}(X)}{\quoi},\tag{13.4.3}
\end{equation*}\etiquette[13.4.3]{IV.eq13.4.3}
where $A'_X$ and $A_{u^{-1}(X)}$ denote the generated free Modules
(\hyperref[IV.prop11.3.3]{11.3.3}).

\item When $A'$ is commutative and the canonical morphism
$u^{-1}A'\rightarrow A$ is central (\resp when the canonical
morphism $u^{-1}A'\rightarrow A$ is an isomorphism), there is, for a
variable right $A'$-Module $M$ and a variable left $A'$-Module $N$,
a canonical bifunctorial isomorphism:
\hypertarget{sga4.exp.IV.equation.13.4.4}{}\begin{align*}
u^{\ast}M\otimes_A u^{\ast}N &\simeq
u^{\ast}(M\otimes_{A'}N)\tag{13.4.4}\\
u^{\ast}M\otimes_A u^{\ast}N &\simeq
u^{-1}(M\otimes_{A'}N){\quoi}.\tag*{(\resp (13.4.5))}
\end{align*}\hypertarget{sga4.exp.IV.equation.13.4.5}{}
\end{enumerate}
\end{enonce*}

First, there is a canonical isomorphism
$\Hom_A(u^{\ast}N,M)\simeq \Hom_{u^{-1}A'}(u^{-1}N,M)$
(\sisi{\hyperref[IV.cor12.12]{12.12}}{\hyperref[IV.prop12.14]{12.14}}),
and then an isomorphism
$\Hom_{u^{-1}A'}(u^{-1}N,M)\simeq
\Hom_{A'}(N,u_{\ast}M)$
(\hyperrefext[III][III.prop1.7]{1.7}); this gives \hyperlink{sga4.exp.IV.equation.13.4.1}{(13.4.1)}. Let us
exhibit the isomorphism \hyperlink{sga4.exp.IV.equation.13.4.2}{(13.4.2)}. For every object $X$ of $E'$, there
is the following sequence of functorial isomorphisms:
\begin{alignat*}{3}
& \Hom_{E'}(X,u_{\ast}\scrHom_A(u^{\ast}N,M))\simeq && \Hom(
u^{-1}X,
 \scrHom_A(u^{\ast}N,M)) &\quad& \text{(adjunction){\quoi},}\\
&\Hom_A(u^{\ast}N,\scrHom_E(u^{-1}X,M)) \simeq &&\qquad ''\qquad ''\qquad
''\qquad '' && (\hyperref[IV.prop12.1]{12.1}){\quoi},\\
& \Hom_A(u^{\ast}N,\scrHom_E(u^{-1}X,M))\simeq &&\Hom_{A'}(N,
u_{\ast}\scrHom_E(u^{-1}X,M)) && \hyperlink{sga4.exp.IV.equation.13.4.1}{(13.4.1)}{\quoi},\\
& \Hom_{A'}(N,\scrHom_{E'}(X,u_{\ast}M))\simeq &&\qquad ''\qquad ''\qquad
 ''\qquad '' && \eqhyperref[IV.eq10.3.1]{10.3.1}{\quoi},\\
& \qquad ''\qquad ''\qquad ''\qquad '' \simeq &&
 \Hom_{E'}(X,\scrHom_{A'}(N,u_{\ast}M)) &&
 (\hyperref[IV.prop12.1]{12.1}){\quoi}.
\end{alignat*}
Formula \hyperlink{sga4.exp.IV.equation.13.4.5}{(13.4.5)}\pageoriginale then follows from formula \hyperlink{sga4.exp.IV.equation.13.4.2}{(13.4.2)}
by adjunction (the definition of the tensor product
(\hyperref[IV.prop12.7]{12.7})). Formula \hyperlink{sga4.exp.IV.equation.13.4.4}{(13.4.4)} follows from
\hyperlink{sga4.exp.IV.equation.13.4.5}{(13.4.5)}, using the commutativity of the tensor product when the
base ring is commutative (\hyperref[IV.subsec12.8]{12.8}). Finally,
to prove \eqhyperref[IV.eq13.4.3]{13.4.3}, consider the sequence of
isomorphisms
\begin{alignat*}{3}
& \Hom_A(u^{\ast}A'_X,M) && \simeq\Hom_{A'}(A'_X,u_{\ast}M)
&\quad& \hyperlink{sga4.exp.IV.equation.13.4.1}{(13.4.1)}{\quoi},\\
& \Hom_{E'}(X,u_{\ast}M) && \simeq\quad ''\qquad ''\qquad ''&&
(\hyperref[IV.prop11.3.3]{11.3.3}){\quoi},\\
&\quad ''\qquad ''\qquad ''&& \simeq\Hom_E(u^{-1}X,M) &&
\text{(adjunction)}{\quoi},\\
&\Hom_A(A_{u^{-1}X},M) &&\simeq\quad ''\qquad ''\qquad ''&&
(\hyperref[IV.prop11.3.3]{11.3.3}){\quoi}.
\end{alignat*}

\begin{enonce*}{Corollary 13.5}\etiquette[13.5]{IV.cor13.5}
Let $(E,A)$ be a ringed topos, let $x:P\rightarrow E$ be a point of
$E$, and let $F\mapsto F_x$ be the associated fiber functor
(\hyperref[IV.def6.1]{6.1}). The fiber functor at $x$ carries the
tensor product of $A$-Modules to the (ordinary) tensor product of
$A_x$-modules.
\end{enonce*}

The tensor product in $P$ is the ordinary tensor product of modules
($P$, the punctual topos, is the category of sets). The assertion
therefore follows from \hyperlink{sga4.exp.IV.equation.13.4.4}{(13.4.4)}.

\begin{enonce*}{Corollary 13.6}\etiquette[13.6]{IV.cor13.6}
Let $u:(E,A)\rightarrow(E',A')$ be a morphism of ringed topoi. The
direct-image functor $u_{\ast}$ for right or left Modules commutes
with projective limits and, in particular, is left exact. The
inverse-image functor $u^{\ast}$ for right or left Modules commutes
with inductive limits and, in particular, is right exact.
\end{enonce*}

This follows from formula \hyperlink{sga4.exp.IV.equation.13.4.1}{(13.4.1)}
(\hyperrefext[I][I.prop2.11]{2.11}).

\setcounter{subsection}{6}
\subsection{}\etiquette[13.7]{IV.subsec13.7}
Note that the inverse-image functor $u^{\ast}$ for Modules is not, in
general, exact, whereas the set-theoretic inverse-image functor
$u^{-1}$ is exact. The exactness of $u^{\ast}$ can nevertheless be
asserted when the canonical morphism $u^{-1}A'\rightarrow A$ is flat
(on the right\pageoriginale or on the left) (V \hyperref[V.def1.8]{1.8}). This is so, in
particular, when the canonical morphism $u^{-1}A'\rightarrow A$ is an
isomorphism.

\subsection{}\etiquette[13.8]{IV.subsec13.8}
Let $C$ be a $𝒰$-site and let $A'$ be a presheaf of rings on $C$;
denote its associated sheaf by $A$. The category of presheaves of
$A'$-modules which are sheaves is equivalent to the category of
$A$-Modules. The notation $\Hom_{A'}(M,N)$ is sometimes used for the
group $\Hom_A(M,N)$ (\hyperref[IV.subsubsec11.1.2]{11.1.2}).
Likewise, and by abuse, the notations $\scrHom_{A'}(M,N)$ and
$M\otimes_{A'}N$ are used for the sheaves $\scrHom_A(M,N)$ and
$M\otimes_A N$. When $A'=k_E$ is the constant presheaf defined by an
ordinary ring $k$, we also write $\scrHom_k$, $\otimes_k$ instead of
$\scrHom_{k_E}$, $\otimes_{k_E}$.

\begin{enonce*}[remark]{Exercise 13.9}\etiquette[13.9]{IV.exer13.9}
\emph{Locally ringed topoi} (\cf [\hyperref[IV.sec9]{9}] for further
information in the direction considered below).

Let $(E,{\Oo}_E)$ be a commutatively ringed topos. For $X\in \ob E$
and $f\in {\Oo}_E(X)$, let $X_f$ be the largest subobject of $X$ on
which $f$ is invertible.
\begin{enumerate}
\renewcommand{\theenumi}{\alph{enumi}}
\renewcommand{\labelenumi}{\theenumi)}
\item Show that, for $X'\in \ob E_{/X}$, $f_{X'}$ is invertible if
and only if the structural morphism $X'\rightarrow X$ factors through
$X_f$, and that, for $f,g\in {\Oo}_E(X)$, one has
$$
X_{fg}=X_f\cap X_g{\quoi}.
$$

\item Show that the following conditions (i) through (iii) are
equivalent:
\begin{enumerate}
\renewcommand{\theenumii}{\roman{enumii}}
\renewcommand{\labelenumii}{(\theenumii)}
\item For $X\in \ob E$ and $f,g\in {\Oo}_E(X)$, one has
$$
X_{f+g}\subset \Sup(X_f,X_g).
$$

\item For $X\in \ob E$ and $f,g\in {\Oo}_E(X)$ such that $f+g$ is
invertible, one has $X=\Sup(X_f,X_g)$.

\item For $X\in \ob E$ and $f\in {\Oo}_E(X)$, one has
$$
X=\Sup(X_f,X_{1-f}){\quoi}.
$$

Show that these conditions imply the following condition (iv),
and\pageoriginale are equivalent to it if $E$ has enough points:

\item For every point $p$ of $E$, the fiber ring ${\Oo}_{E,p}$ is a
\emph{local ring}.
\end{enumerate}

When the equivalent conditions (i) through (iii) hold, $(E,{\Oo}_E)$
is said to be a \emph{locally ringed topos}; likewise, a ringed site
is said to be a \emph{locally ringed site} if the corresponding
ringed topos is a locally ringed topos.

\item Let $(E,{\Oo}_E)$ and $(E',{\Oo}_{E'})$ be two locally ringed
topoi, and let $f$ be a morphism of ringed topoi from the first to the
second. Show that the following condition (i) implies condition (ii),
and is equivalent to it if $E$ has enough points:
\begin{enumerate}
\renewcommand{\theenumii}{\roman{enumii}}
\renewcommand{\labelenumii}{(\theenumii)}
\item For every $X'\in \ob E'$ and $s'\in {\Oo}_{E'}(X')$, setting
$X=f^{-1}(X')$ and $s=f^{\ast}(s')$, one has
$$
X_s=f^{-1}(X'_{s'}){\quoi}.
$$

\item For every point $p$ of $E$, setting $p'=f(p)$, the natural
homomorphism on the fibers
$$
{\Oo}_{E',p'}\longrightarrow {\Oo}_{E,p}
$$
is a \emph{local} homomorphism of local rings.
\end{enumerate}

When condition (i) holds, $f$ is said to be a \emph{morphism of
locally ringed topoi}, or also an \emph{admissible morphism of locally
ringed topoi} if there is a risk of confusion. We denote by
$\Hhomtoplocan(E,E')$ the full subcategory of the category
$\Hhomtopan(E,E')$ of all morphisms of ringed topoi from $E$ to $E'$
defined by the admissible morphisms from $E$ to $E'$.

\item Suppose that $(E,{\Oo}_E)$ is the ringed topos associated with
a ringed topological space $(X,{\Oo}_X)$
(\hyperref[IV.subsec2.1]{2.1}). Show that $(E,{\Oo}_E)$ is locally
ringed if and only if $(X,{\Oo}_X)$ is locally ringed, \ie if and
only if, for every $x\in X$, the fiber ${\Oo}_{X,x}$ is a local
ring. (Note that in criterion (iv) of b), it suffices to take $p$ in
a \emph{conservative} family of points of $E$.) Let\pageoriginale
$f:(X,{\Oo}_X)\rightarrow(X',{\Oo}_{X'})$ be a morphism of ringed
spaces, where ${\Oo}_X$ and ${\Oo}_{X'}$ are sheaves of local rings.
Show that the corresponding morphism of ringed topoi
$\Top(X,{\Oo}_X)\rightarrow \Top(X',{\Oo}_{X'})$ is admissible if and
only if the same is true of the morphism $f$, \ie if and only if,
for every $x\in X$,
$${\Oo}_{X',f(x)}\longrightarrow {\Oo}_{X,x}$$
is a \emph{local} homomorphism of local rings.

\item Let $(E,{\Oo}_E)$ and $(E',{\Oo}_{E'})$ be two locally ringed
topoi such that $E$ has enough points and $(E',{\Oo}_{E'})$ is
equivalent to the ringed topos defined by a \emph{scheme}
$(X',{\Oo}_{X'})$; prove that the category $\Hhomtoplocan(E,E')$ is
equivalent to a \emph{discrete} category, \ie that it is a rigid
groupoid (every morphism is an isomorphism, and the automorphism
groups of the objects are the trivial groups).

\noindent (Hint: reduce to the case where $(E,{\Oo}_E)$ is the
punctual topos ringed by a field.) Recall, on the other hand, that
$\Hhomtop(E,E')$ is not in general equivalent to a discrete
category, even if $E$ and $E'$ are topoi defined by schemes (and even
if $E$ is the punctual topos), \cf
\hyperref[IV.subsubsec4.2.3]{4.2.3}.

\item Let $E$ be a locally ringed topos, and let $P$ be the ringed
topos defined by the ringed space $\Spec k$, where $k$ is a field.
Show that the admissible morphisms of locally ringed topoi
$P\rightarrow E$ correspond to pairs $(p,u)$ consisting of a point
$p$ of $E$ and an injection of $k(p)$ into $k$, where $k(p)$ is the
residue field of the local ring ${\Oo}_{E,p}$. Generalize this to a
statement exhibiting the admissible morphisms from a punctual
locally ringed topos to $E$ (generalizing EGA I 2.4.4).

A \emph{geometric point of a locally ringed topos} $E$ is any
admissible morphism into $E$ from the locally ringed topos defined by
a ringed space of the form $\Spec(k)$, where $k$ is an
\emph{algebraically closed} field. Define the \emph{category of
geometric points of} $E$ (with the understanding that the
corresponding fields $k$ belong to $𝒰$), denoted by $\Pptgeom(E)$,
and a canonical functor $\Pptgeom(E)\rightarrow \Ppoint(E)$. Show
that this functor is faithful if\pageoriginale and only if $E$ has no
point, \ie $\Ppoint(E)$ is the empty category.

\item Let $E$ be a $𝒰$-topos and let $𝒱$ be a universe such that
$𝒰\in 𝒱$. Define a 2-category
$(𝒱\text{-}𝒰\text{-}\Top)_{/E}$ in terms of the object $E$ of the
2-category $(𝒱\text{-}𝒰\text{-}\Top)$
(\hyperref[IV.subsubsec3.3.1]{3.3.1}), following
\hyperref[IV.exer5.14]{5.14} a). When $E$ is ringed by a Ring $A$,
define similarly a 2-category
$(𝒱\text{-}𝒰\text{-}\text{Topan})_{/E}$, and when $E$ is locally
ringed, use the notion of admissible morphism (\cf c)) to define a
2-category
$(𝒱\text{-}𝒰\text{-}\text{Topanloc})_{/E}$ having $\Pptgeom(E)$
(\cf f)) as a full subcategory.
\end{enumerate}
\end{enonce*}

% END INCLUDED SGA4_English_current_buildable_source_20260906/source/04/drafts/SGA4_IV_LINES_9248_9593_English.texfrag


% BEGIN INCLUDED SGA4_English_current_buildable_source_20260906/source/04/drafts/SGA4_IV_LINES_9594_9852_English.texfrag
\section{Modules on a topos defined by gluing}\etiquette[14]{IV.sec14}

\subsection{}\etiquette[14.1]{IV.subsec14.1}
Let $(E,A)$ be a ringed topos, let $U$ be an open subtopos of $E$,
let $F$ be the complementary closed subtopos, and let
$j:U\rightarrow E$ and $i:F\rightarrow E$ be the canonical
morphisms (\hyperref[IV.subsec9.3]{9.3}). The topos $U$ is ringed by
$j^{\ast}A$ (\sisi{\hyperref[IV.subsubsec11.2.2]{11.2.2}}{\hyperref[IV.subsubsec11.2.1]{11.2.1}}),
and the morphism of topoi $j$, completed in the evident way, becomes
a morphism of ringed topoi. In this section, the topos $F$ will be
ringed by the sheaf $i^{\ast}A$, denoted by $A|F$, and the morphism
$i$, completed in the evident way, is then a morphism of ringed
topoi.


\subsection{}\etiquette[14.2]{IV.subsec14.2}
We therefore have two morphisms of ringed topoi:
$$
j:(U,A| U)\longrightarrow (E,A)\text{ and } i:(F,A| F)\longrightarrow
(E,A){\quoi};
$$
whence five functors between the corresponding categories of
modules (left modules, to fix ideas):
\[\vcenter{
\xymatrix@C=2cm@R=1.5cm@M=.25cm{ ({}_{A|U}U) \ar@<14pt>[r]^{j_!}
\ar@<-14pt>[r]^{j_{\ast}}&
  ({}_{A}E) \ar[l]_{j^{\ast}} \ar@<7pt>[r]^{i^{\ast}} &
  ({}_{A|F}F)\quad. \ar@<7pt>[l]_{i_{\ast}}
}}\leqno(14.2.1)
\]\etiquette[14.2.1]{IV.eq14.2.1}
Each functor in diagram \eqhyperref[IV.eq14.2.1]{14.2.1} is left
adjoint to the functor lying below it.

\subsection{}\etiquette[14.3]{IV.subsec14.3}
We\pageoriginale know that the functor $X\mapsto (j^{\ast}X,
i^{\ast}X;i^{\ast}X\rightarrow i^{\ast}j_{\ast}j^{\ast}X)$ is an
equivalence from $E$ to the topos $(U,F,i^{\ast}j_{\ast})$
(\hyperref[IV.theo9.5.4]{9.5.4}). This equivalence induces an
equivalence between the corresponding categories of modules, and
hence the functor $P\mapsto (j^{\ast}P,i^{\ast}P;
i^{\ast}P\rightarrow i^{\ast}j_{\ast}j^{\ast}P)$ is an equivalence,
denoted by $\Phi$, from the category ${}_AE$ to the category
$({}_{A| U}U,{}_{A| F}F,i^{\ast}j_{\ast})$. The functors in
\hyperref[IV.eq14.2.1]{(14.2.1)}, composed with $\Phi$ or $\Phi^{-1}$, are then the
functors:
\begin{align*}
\Phi\circ j_{\ast} &: M\mapsto (M,i^{\ast} j_{\ast}M;
i^{\ast}j_{\ast}M\xrightarrow{\id} i^{\ast}j_{\ast}M){\quoi},\\
j^{\ast}\circ \Phi^{-1} &: (M,N; N\rightarrow
i^{\ast}j_{\ast}M)\mapsto M{\quoi},\\
\Phi\circ j_{!} &: M\mapsto (M,0;0\rightarrow i^{\ast}
j_{\ast}M){\quoi},\\
\Phi\circ i_{\ast} &: N\mapsto (0,N;N\rightarrow 0){\quoi},\\
i^{\ast}\circ \Phi^{-1} &: (M,N;N\rightarrow
i^{\ast}j_{\ast}M)\mapsto N{\quoi}.
\end{align*}
The reader may, as an exercise, make the various adjunction
morphisms explicit.

\subsection{}\etiquette[14.4]{IV.subsec14.4}
Denote by
\begin{equation*}
i^{!}:{}_AE\longrightarrow {}_{A| F}F\tag{14.4.1}
\end{equation*}\etiquette[14.4.1]{IV.eq14.4.1}
the functor defined by the formula
\begin{equation*}
i^{!}\circ
\Phi^{-1}(M,N;N\xrightarrow{u}i^{\ast}j_{\ast}M)=\Ker(u){\quoi}.\tag{14.4.2}
\end{equation*}\etiquette[14.4.2]{IV.eq14.4.2}

\begin{enonce*}{Proposition 14.5}\etiquette[14.5]{IV.prop14.5}
The functor $i^{!}$ is right adjoint to the functor $i_{\ast}$. The
adjunction morphism $i_{\ast}i^{!}\rightarrow\id$ is a
monomorphism. The functors $i^{!}$ (for varying rings) commute with
restriction-of-scalars functors.
\end{enonce*}

It is clear that every morphism
$$
(0,N;0)\longrightarrow (M,N;N\xrightarrow{u}i^{\ast}j_{\ast}M)
$$
factors uniquely through $(0,\ker(u);0)$, which proves the
adjunction property. The other assertions are trivial.

\begin{enonce*}{Proposition 14.6}\etiquette[14.6]{IV.prop14.6}
For\pageoriginale every object $P$ of ${}_AE$, there are exact
sequences, functorial in $P$:
\hypertarget{sga4.exp.IV.equation.14.6.1}{}\begin{gather}
0\longrightarrow j_{!}j^{\ast}P\longrightarrow P\longrightarrow
i_{\ast}i^{\ast}P\longrightarrow 0{\quoi};\tag{14.6.1}\\
0\longrightarrow i_{\ast}i^{!}P\longrightarrow P\longrightarrow
j_{\ast}j^{\ast}P{\quoi},\tag{14.6.2}
\end{gather}\hypertarget{sga4.exp.IV.equation.14.6.2}{}
where the nontrivial arrows are the adjunction arrows.
\end{enonce*}

First note that a sequence
$$
(M',N';u')\longrightarrow(M,N;u)\longrightarrow(M'',N'';u'')
$$
in $({}_{A|U}U,{}_{A|F}F;i^{\ast}j_{\ast})$ is exact if and only
if the corresponding sequences
\begin{align*}
& M'\longrightarrow M\longrightarrow M''\\
& N'\longrightarrow N\longrightarrow N''
\end{align*}
are exact.

The sequences \hyperlink{sga4.exp.IV.equation.14.6.1}{(14.6.1)} and \hyperlink{sga4.exp.IV.equation.14.6.2}{(14.6.2)} are transformed by the
equivalence $\Phi$ into sequences of the type in
\eqhyperref[IV.subsec14.3]{14.3}:
\begin{align*}
& 0\longrightarrow (M,0;0)\longrightarrow(M,N;u)\longrightarrow
(0,N;0)\longrightarrow 0\\
& 0\longrightarrow(0,\ker(u);0)\longrightarrow
(M,N,u)\longrightarrow(M,i^{\ast}j_{\ast}M;\id){\quoi}.
\end{align*}
The exactness of these sequences is immediate.

\setcounter{subsection}{6}
\subsection{}\etiquette[14.7]{IV.subsec14.7}
The exact sequence \hyperlink{sga4.exp.IV.equation.14.6.2}{(14.6.2)} gives a new interpretation of the
functor $i_{\ast}i^{!}$. Indeed, let $X$ be an object of $E$.
Sequence \hyperlink{sga4.exp.IV.equation.14.6.2}{(14.6.2)} gives the exact sequence of abelian groups
$$
0\longrightarrow \Hom_E(X,i_{\ast}i^{!}P)\longrightarrow
\Hom_E(X,P)\longrightarrow \Hom_E(X,j_{\ast}j^{\ast}P){\quoi}.
$$
Also denote by $U$ the open subobject of $E$ which is the final
object of the open subtopos $U$. It follows from the adjunction
properties of the functors $j_{\ast}$, $j^{\ast}$, and
$j_{!}^{\ens}$ that the group $\Hom_E(X,j_{\ast}j^{\ast}P)$ is
canonically isomorphic to $\Hom_E(X\times U,P)$, and that the
morphism $\Hom_E(X,P)\rightarrow\Hom_E(X,j_{\ast}j^{\ast}P)$ is
precisely the morphism defined by the monomorphism
$X\times U\rightarrow X$ (\hyperref[IV.sec5]{5}). In other words
(\hyperref[IV.subsubsec8.5.2]{8.5.2}):

\begin{enonce*}{Proposition 14.8}\etiquette[14.8]{IV.prop14.8}
The\pageoriginale sections of $i_{\ast}i^{!}P$ on $E_{/X}$ are the
sections of $P$ on $E_{/X}$ whose support
$(\hyperref[IV.subsubsec9.3.5]{9.3.5})$ is contained in $F$
\sisi{(notation of $5.9.1$)}{}. In other words, $i_{\ast}i^{!}P$ is
the largest subsheaf of $P$ whose support is contained in $F$.
\end{enonce*}

\setcounter{subsection}{8}
\subsection{}\etiquette[14.9]{IV.subsec14.9}
By abuse of language, one sometimes says that $i_{\ast}i^{!}P$ is
the submodule of $P$ defined by the sections of $P$ supported in
$F$. It follows from \hyperref[IV.rem8.5.3]{8.5.3} that this
terminology merely extends to general topoi terminology used for
topoi of sheaves on topological spaces.

\begin{enonce*}{Proposition 14.10}\etiquette[14.10]{IV.prop14.10}
\begin{enumerate}
\renewcommand{\labelenumi}{\theenumi)}
\item The functor $P\mapsto i_{\ast}i^{\ast}P$ is isomorphic to the
functor $P\mapsto i_{\ast}i^{\ast}A\otimes_A P$.

\item The functor $P\mapsto i_{\ast}i^{!}P$ is isomorphic to the
functor $P\mapsto\scrHom_A(i_{\ast}i^{\ast}A,P)$.
\end{enumerate}
\end{enonce*}

The exact sequence \hyperlink{sga4.exp.IV.equation.14.6.1}{(14.6.1)}, in the particular case where $P$ is
the $A$-module $A$, reads
\begin{equation*}
0\longrightarrow A_U\longrightarrow A\longrightarrow
i_{\ast}i^{\ast}A\longrightarrow 0.\tag{14.10.1}
\end{equation*}\etiquette[14.10.1]{IV.eq14.10.1}
Here $A_U$ is the free $A$-module generated by the open subobject
$U$ corresponding to the open subtopos $U$ (9 and \hyperref[IV.prop11.3.3]{11.3.3}). For
every $A$-module $P$, the $A$-modules $j_{!}j^{\ast}P$ and
$j_{\ast}j^{\ast}P$ are canonically isomorphic respectively to
$A_U\otimes_A P$ and $\scrHom_A(A_U,P)$
(\hyperref[IV.prop12.3]{12.3} and \hyperref[IV.cor12.6]{12.6}).
Moreover, under these isomorphisms, the canonical morphisms
$j_{!}j^{\ast}P\rightarrow P$ and
$P\rightarrow j_{\ast}j^{\ast}P$ arise from the monomorphism
$A_U\rightarrow A$. The exact sequence
\hyperref[IV.eq14.10.1]{14.10.1} therefore gives two exact
sequences (\hyperref[IV.subsec12.2]{12.2} and
\hyperref[IV.prop12.11]{12.11}):
\begin{align*}
& j_{!}j^{\ast}P\longrightarrow P\longrightarrow
i_{\ast}i^{\ast}A\otimes_AP\longrightarrow 0{\quoi},\\
& 0\longrightarrow \scrHom_A(i_{\ast}i^{\ast}A,P)\longrightarrow
P\longrightarrow j_{\ast}j^{\ast}P{\quoi};
\end{align*}
and the stated isomorphisms follow by comparison with \hyperlink{sga4.exp.IV.equation.14.6.1}{(14.6.1)} and
\hyperlink{sga4.exp.IV.equation.14.6.2}{(14.6.2)}.

\begin{thebibliography}{99}\pageoriginale
\bibitem{sga4.exp.IV.bib.1} M. Artin and B. Mazur, Homotopy of varieties in the etale
topology, in: Proceedings of a Conference on Local Fields,
Driebergen 1966, Springer.
\bibitem{sga4.exp.IV.bib.2} P. Gabriel and G. Zisman, Calculus of fractions and
homotopy theory, Ergebnisse der Mathematik, Bd 35.

\bibitem{sga4.exp.IV.bib.3} J. Giraud, Algèbre homologique non commutative,
Grundlehren Springer 1971.

\bibitem{sga4.exp.IV.bib.4} A. Grothendieck, Sur quelques points d'Algèbre
homologique, Tohoku Math. Journal (cited as [Toh]).

\bibitem{sga4.exp.IV.bib.5} A. Grothendieck, Fondements de la Géométrie
Algébrique (collection of Bourbaki talks 1957/62), Secrétariat
mathématique, 11 rue P. Curie, Paris.

\bibitem{sga4.exp.IV.bib.6} A. Grothendieck, Crystals and the De Rham cohomology
of schemes (notes by I. Coates and O. Jussila), in: Dix exposés sur
la cohomologie des schémas, North Holland Pub. Cie., 1969.

\bibitem{sga4.exp.IV.bib.7} A. Grothendieck, Classes de Chern des
représentations linéaires des groupes discrets, in: Dix exposés sur
la cohomologie des schémas, North Holland Pub. Cie., 1969.

\bibitem{sga4.exp.IV.bib.8} R. Godement, Théorie des faisceaux,
Act. Scient. Ind. ${\rm n^{\circ}}$ 1 252 (1958), Hermann (Paris)
(cited as [TF]).

\bibitem{sga4.exp.IV.bib.9} M. Hakim, Topos annelés et schémas relatifs,
mimeographed thesis, Orsay 1967.

\bibitem{sga4.exp.IV.bib.10} D. Mumford, Picard groups of moduli problems, in
Arithmetic Algebraic Geometry, Harper's Series in Modern
Mathematics.

\bibitem{sga4.exp.IV.bib.11} Nguyen Dinh Ngoc, Thèse Sciences Mathématiques,
Paris 1963, ${\rm n^{\circ}}$ 4 995.

\bibitem{sga4.exp.IV.bib.12} J.E. Roos, Distributivité des $\varinjlim$ par
rapport aux $\varprojlim$ des topos

\begin{tabular}{@{}r@{\;}l@{\hspace{1cm}}r@{\;}l}
a) & CR t.259 p. 969-972 & b) & CR T. 259 p. 1605-1608\\
c) & CR t.259 p. 1801-1804 & & (August and September 1964).
\end{tabular}

\bibitem{sga4.exp.IV.bib.13} P. Deligne--D. Mumford, The irreducibility of
the space of curves of given genus, Pub. math. ${\rm n^{\circ}}$ 36
(1969).

\bibitem{sga4.exp.IV.bib.14} B. Mitchell, Theory of categories, Academic Press
(1965).

\bibitem{sga4.exp.IV.bib.15} J. Lurie, Higher Topos Theory, Annals of Mathematics Studies, 170. Princeton University Press,
Princeton, NJ, 2009. xviii+925 pp.

\bibitem{sga4.exp.IV.bib.16} P. Johnstone, Topos theory,
London Mathematical Society Monographs, Vol. 10.
Academic Press [Harcourt Brace Jovanovich, Publishers], London--New York, 1977. xxiii+367 pp.

\bibitem{sga4.exp.IV.bib.17} F. Borceux, G. Van den Bossche, Structure des topologies
d'un topos, Cahiers de topologie et géométrie différentielle catégoriques,
t. 25, n°1 (1984), p. 37--39.

\end{thebibliography}

\ifx\danslelivre\undefined
\end{document}
\fi

% END INCLUDED SGA4_English_current_buildable_source_20260906/source/04/drafts/SGA4_IV_LINES_9594_9852_English.texfrag


% END INCLUDED SGA4_English_current_buildable_source_20260906/source/04/04_English.tex


\makeatletter
\def\input@path{{../05/}}
\makeatother
% BEGIN INCLUDED SGA4_English_current_buildable_source_20260906/source/05/05_English.tex
\ifx\danslelivre\undefined
% Release engine: XeLaTeX (the LuaTeX compatibility layer breaks Xy-pic output).
\documentclass[10pt]{../fichiers-configuration-latex/smfart}
% BEGIN INCLUDED SGA4_English_current_buildable_source_20260906/source/fichiers-configuration-latex/paquets.tex
\usepackage{polyglossia}
%\usepackage[francais]{babel}
\usepackage{fontspec}
\usepackage[math-style=ISO,bold-style=upright,sans-style=upright]{unicode-math}
%\usepackage[mathletters]{ucs}
%\usepackage[utf8]{inputenc}
%\usepackage[T1]{fontenc}
\usepackage[unicode,hypertexnames=false,colorlinks=true,linkcolor=blue,citecolor=Sepia,urlcolor=magenta]{hyperref}
\setmainlanguage{french}
\setmainfont[Ligatures=TeX]{Linux Libertine O}
% variante possible : TeXGyrePagella ?
\setmathfont{XITS Math}

% maths
\usepackage{amsmath}
%\usepackage{amssymb}
%\usepackage{mathrsfs}
%\usepackage{stmaryrd}
%\usepackage{euscript}

% couleurs
\usepackage[usenames,dvipsnames]{xcolor}

% graphiques
\usepackage{graphicx}
\usepackage[all]{xy}
%\usepackage{tikz}
%\usetikzlibrary{matrix,arrows,calc}
%\tikzset{column sep=3em,row sep=2em,text height=1.5ex, text depth=0.25ex,>=angle 90}
%\tikzset{description/.style={fill=white,inner sep=1pt}}

% hyperliens, 'etiquettes
\usepackage{xr-hyper}
\usepackage{../fichiers-configuration-latex/smfthm}
%\ifx\cacheretiquettes\undefined
%\usepackage[inline]{showlabels} \else \fi % pour rendre visibles les 'etiquettes
%\usepackage{../fichiers-configuration-latex/makeidx} % pour faire un index
%\usepackage[pagebackref,unicode]{hyperref}  % pour voir o`u sont cit'ees les r'ef'erences
% hyperref is loaded once above with the cumulative-reader destination policy.
%\hypersetup{breaklinks=true} %permet le retour `a la ligne dans les liens trop longs. Incompatible avec pdflatex
%\renewcommand*{\backref}[1]{{\small [cit\'e page(s)~#1.]}}

\usepackage{xspace}
\usepackage{pifont} % pour expose XVII

% END INCLUDED SGA4_English_current_buildable_source_20260906/source/fichiers-configuration-latex/paquets.tex

% BEGIN INCLUDED SGA4_English_current_buildable_source_20260906/source/fichiers-configuration-latex/numerotation.tex
\renewcommand{\thefootnote}{\roman{footnote}} % pour mettre les notes en bas de page sous la forme (ii) et non (2), qui ressemble trop (?) a un exposant.

\newcounter{pageoriginal}
\newcommand{\pageoriginale}{\refstepcounter{pageoriginal}\marginpar{\footnotesize\textbf{\thepageoriginal}}\xspace}
\let\pageoriginaled\pageoriginale

\usepackage{ifthen}
\newboolean{original}
\def\sisi{\ifthenelse{\boolean{original}}}
\setboolean{original}{false}

\newcounter{cnote}
\def\nde#1{\ifthenelse{\boolean{original}}{}{\stepcounter{cnote}\def\thefootnote{\arabic{cnote}}\protect\footnote[\value{cnote}]{N.D.E.: #1}}}
\def\ndemark{\ifthenelse{\boolean{original}}{}{\stepcounter{cnote}\def\thefootnote{\arabic{cnote}}\protect\footnotemark[\value{cnote}]}}
\def\ndetext#1{\ifthenelse{\boolean{original}}{}{\def\thefootnote{\arabic{cnote}}\protect\footnotetext[\value{cnote}]{N.D.E.: #1}}}

% END INCLUDED SGA4_English_current_buildable_source_20260906/source/fichiers-configuration-latex/numerotation.tex

% BEGIN INCLUDED SGA4_English_current_buildable_source_20260906/source/fichiers-configuration-latex/macros.tex
\setmainlanguage{french}
\setotherlanguage{english}


%% Fixed, engine-independent release identifier.  The inherited dynamic Git
%% probe used Lua io.popen, exposed build-state-dependent output, and printed
%% a visible question mark outside a Git checkout.
\def\showgitstatus{English release r2, 2026-07-23}


% macros pour LaTeX
\renewcommand{\mathsf}{}

% Police différente pour U+2113 SCRIPT SMALL L
\setmathfont[range="2113]{latinmodern-math.otf}

% Give every inherited semantic label its own stable PDF destination instead
% of inheriting whichever automatic equation/page anchor happened to be live.
% The legacy label remains as a compatibility alias used by existing edges.
\makeatletter
\def\etiquette[#1]#2{%
  \Hy@raisedlink{\hyper@anchorstart{sga4.semantic.#2}\hyper@anchorend}%
  \def\@currentHref{sga4.semantic.#2}%
  \def\@currentlabel{#1}%
  \label{#2}%
}
\makeatother
\def\eqhyperref[#1]#2{(\hyperref[#1]{#2})}
\def\pagehyperref[#1]#2{(page \hyperref[#1]{#2})}
\def\hyperrefext[#1][#2]#3{\hyperref[#2]{#1 #3}}
\def\eqhyperrefext[#1][#2]#3{\hyperref[#2]{#1 #3}}     %cf. exposes 17
\def\neweqhyperrefext[#1][#2]#3{\hyperref[#2]{#1 #3}}  %cf. exposes 17

\def\shat#1{\protect\widehat{#1}}
\def\stilde#1{\protect\widetilde{#1}}
%%%%%%%%%%%%%%%%%%


\def\etcn{etc.\xspace}
\def\etco{etc...\xspace}
\def\etc{\sisi{\etco}{\etcn}}
\def\ldotsn{\sisi{...}{\etcn}}

\def\cf{{cf.}\kern.3em}
\def\Cf{{Cf.}\kern.3em}
\def\ie{{i.e.}\kern.3em}
\def\iev{{i.e.},\xspace}
\def\loccito{loc.\kern2pt cit.\xspace}
\def\loccitn{\textit{loc.\kern2pt cit.}\xspace}
\def\loccit{\sisi{\loccito}{\loccitn}}
\def\EGA{\textup{EGA}\kern.2em}
\def\SGA{\textup{SGA}\kern.3em}

\def\quoi{}

\def\quav{\quad,\quad}
\def\quaet{\quad \text{et}\quad}
\def\fibreproduct#1#2#3{#1 ×_{#3} #2}
\let\fprod\fibreproduct
\def\ccheck#1#2{\mbox{${\check{\check{\rm #1}}}^{\lower2pt\hbox{$\scriptstyle #2$}}$}}
\def\sheaf#1{\protect\mathcal{#1}}
\def\vieuxsheaf#1{{#1}} % cf. problème avec sheaf{a\Phi_i} par exemple...
\DeclareMathOperator{\Hom}{Hom}
\DeclareMathOperator{\can}{can}
\DeclareMathOperator{\Ind}{Ind}
\DeclareMathOperator{\ind}{ind}
%\DeclareMathOperator{\Hhomtop}{\mathcal{H}\!\mathit{omtop}}
\DeclareMathOperator{\Hhomtop}{\mathbf{Homtop}}
\DeclareMathOperator{\Point}{Point}
\DeclareMathOperator{\Pro}{Pro}
\DeclareMathOperator{\Top}{Top}
\DeclareMathOperator{\Mod}{Mod}
\DeclareMathOperator{\Cart}{Cart}
%\DeclareMathOperator{\Mmorsite}{\mathcal{M}\!\mathit{orsite}}
\DeclareMathOperator{\Mmorsite}{\mathbf{Morsite}}
\DeclareMathOperator{\loc}{loc}
\DeclareMathOperator{\Loc}{Loc}
\DeclareMathOperator{\Pic}{Pic}
\DeclareMathOperator{\Et}{Et}
\DeclareMathOperator{\card}{card}
\DeclareMathOperator{\Ens}{Ens}
\DeclareMathOperator{\h}{\underline{H}}
\DeclareMathOperator{\Fonct}{Fonct}
\DeclareMathOperator{\Ker}{Ker}
\DeclareMathOperator{\but}{but}
\DeclareMathOperator{\Ob}{Ob}
\DeclareMathOperator{\ob}{ob}
\DeclareMathOperator{\Lib}{Lib}
\DeclareMathOperator{\esj}{esj}
\DeclareMathOperator{\id}{id}
\DeclareMathOperator{\Id}{Id}
\DeclareMathOperator{\ad}{ad}
\DeclareMathOperator{\Iim}{Im}
\DeclareMathOperator{\pr}{pr}
\DeclareMathOperator{\Cat}{Cat}
\DeclareMathOperator{\car}{car}
\DeclareMathOperator{\SysInd}{SysInd}
\DeclareMathOperator{\Sup}{Sup}
\DeclareMathOperator{\Coker}{Coker}
\DeclareMathOperator{\Filt}{Filt}
\DeclareMathOperator{\acc}{acc}
\DeclareMathOperator{\sur}{sur}
\DeclareMathOperator{\cat}{cat}
\DeclareMathOperator{\Ord}{Ord}
\DeclareMathOperator{\Gr}{Gr}
\DeclareMathOperator{\Cov}{Cov}
\DeclareMathOperator{\Caf}{Caf}
\DeclareMathOperator{\coker}{coker}
\DeclareMathOperator{\coim}{coim}
\DeclareMathOperator{\im}{im}
\DeclareMathOperator{\Ouv}{Ouv}
\DeclareMathOperator{\diag}{diag}
\DeclareMathOperator{\Esp}{Esp}
\DeclareMathOperator{\TOP}{TOP}
\DeclareMathOperator{\Unif}{Unif}
\DeclareMathOperator{\Iind}{\mathbf{Ind}}
\DeclareMathOperator{\Hhomcart}{\mathbf{Homcart}}
\DeclareMathOperator{\Hhomtopess}{\mathbf{Homtopess}}
\DeclareMathOperator{\Ppoints}{\mathbf{Points}}
\DeclareMathOperator{\Ppoint}{\mathbf{Point}}
\DeclareMathOperator{\Ttors}{\mathbf{Tors}}
\DeclareMathOperator{\Ssex}{\mathbf{Sex}}
\DeclareMathOperator{\Ffib}{\mathbf{Fib}}
\DeclareMathOperator{\Vvois}{\mathbf{Vois}}
\DeclareMathOperator{\Ppointess}{\mathbf{Pointess}}
\DeclareMathOperator{\scrHom}{\mathbf{Hom}}
\DeclareMathOperator{\sheafPic}{\mathbf{Pic}}
\DeclareMathOperator{\sheafAut}{\mathbf{Aut}}
\DeclareMathOperator{\Hhomtoplocan}{\mathbf{Homtoplocan}}
\DeclareMathOperator{\Hhomtopan}{\mathbf{Homtopan}}
\DeclareMathOperator{\Pptgeom}{\mathbf{Ptgéom}}
\DeclareMathOperator{\Morsite}{Morsite}
\DeclareMathOperator{\Fib}{Fib}
\DeclareMathOperator{\Fais}{Fais}
\DeclareMathOperator{\kar}{kar}
\DeclareMathOperator{\Groupes}{Groupes}
\DeclareMathOperator{\Restr}{Restr}
\DeclareMathOperator{\Isom}{Isom}
\DeclareMathOperator{\Prol}{Prol}
\DeclareMathOperator{\Res}{Res}
\DeclareMathOperator{\Tors}{Tors}
\DeclareMathOperator{\esp}{esp}
\DeclareMathOperator{\es}{es}
\DeclareMathOperator{\od}{od}
\DeclareMathOperator{\gr}{gr}
\DeclareMathOperator{\sob}{sob}
\DeclareMathOperator{\dfn}{dfn}
\DeclareMathOperator{\Sch}{Sch}
\DeclareMathOperator{\ouv}{ouv}
\DeclareMathOperator{\cons}{cons}
\DeclareMathOperator{\Nerf}{Nerf}
\DeclareMathOperator{\Points}{Points}
\DeclareMathOperator{\Espsob}{Espsob}
\DeclareMathOperator{\Homtop}{Homtop}
\DeclareMathOperator{\Pt}{Pt}
\DeclareMathOperator{\Quot}{Quot}
\DeclareMathOperator{\glob}{glob}
\DeclareMathOperator{\Homfib}{Homfib}
\DeclareMathOperator{\Qouv}{Qouv}
\DeclareMathOperator{\Quotouv}{Quotouv}
\DeclareMathOperator{\ens}{ens}
\DeclareMathOperator{\ab}{ab}
\DeclareMathOperator{\an}{an}
\DeclareMathOperator{\fs}{fs}
\DeclareMathOperator{\red}{red}
\DeclareMathOperator{\Ab}{Ab}
\DeclareMathOperator{\cosq}{cosq}
\DeclareMathOperator{\Ann}{Ann}
\DeclareMathOperator{\Zar}{Zar}


%%%%%%%%%%%%
%%%%% SGA4-2
\DeclareMathOperator{\Sex}{Sex}
\DeclareMathOperator{\vois}{vois}
\DeclareMathOperator{\Vois}{Vois}
\DeclareMathOperator{\Ext}{Ext}
\DeclareMathOperator{\ext}{ext}
\DeclareMathOperator{\Fer}{Fer}
\DeclareMathOperator{\Add}{Add}
\DeclareMathOperator{\Cocart}{Cocart}
\DeclareMathOperator{\cosk}{cosk}
\DeclareMathOperator{\plat}{plat}

\DeclareMathOperator{\faisc}{faisc}
\DeclareMathOperator{\sk}{sk}
\DeclareMathOperator{\ch}{ch}

\def\HH{\mathbf{H}}
\def\qc{\mathrm{qc}}
\def\et{\mathrm{et}}
\def\coh{\mathrm{coh}}

\def\Hvr#1{%
\begin{array}[b]{l}
{\kern2pt\scriptstyle r}\\[-8pt]
{\kern1pt\scriptstyle \vee}\\[-5pt]
H^{#1}
\end{array}}

\def\Hvv#1{%
\begin{array}[b]{l}
\kern2pt\rotatebox{90}{${\scriptscriptstyle\ll}$}\\[-6pt]
H^{#1}
\end{array}}

\def\Hvi#1{%
\begin{array}[b]{l}
{\kern0.2pt\scriptscriptstyle \infty}\\[-8pt]
{\kern1pt\scriptstyle \vee}\\[-5pt]
H^{#1}
\end{array}}

\def\Hv#1{%
\begin{array}[b]{l}
{\kern1pt\scriptstyle \vee}\\[-5pt]
H^{#1}
\end{array}}

\DeclareMathOperator{\cart}{cart}
\DeclareMathOperator{\Ffaisc}{\mathbf{Faisc}}


\def\Ev#1{%
\begin{array}[b]{l}
{\kern1pt\scriptstyle \vee}\\[-5pt]
E'_{#1}
\end{array}}

\DeclareMathOperator{\Cont}{Cont}
\DeclareMathOperator{\Ccont}{\mathscr{C}\!\mathit{ont}}
\DeclareMathOperator{\Site}{Site}
\newcommand\MM{\mathbf{M}}
\def\limtop#1{{\displaystyle{\varprojlim_{#1}\!\mbox{\rm top }}}}
\def\Limtop{\mathop{\varprojlim\!\mbox{\rm top}}}
\DeclareMathOperator{\spec}{spec}
\def\ss{\scriptstyle{\mathscr{S}}}
\DeclareMathOperator{\Gal}{Gal}
\def\smallm{\scriptstyle{\mathscr{M}}}
\DeclareMathOperator{\End}{End}
\DeclareMathOperator{\ass}{ass}
\DeclareMathOperator{\zar}{zar}
\DeclareMathOperator{\Ogg}{Ogg}
\DeclareMathOperator{\res}{res}
\DeclareMathOperator{\tr}{tr}
\DeclareMathOperator{\Supp}{Supp}
\DeclareMathOperator{\codim}{codim}
\DeclareMathOperator{\rk}{rk}
\DeclareMathOperator{\cur}{cur}
\DeclareMathOperator{\cdqc}{cdqc}
\newcommand\DD{\mathbf{D}}
\DeclareMathOperator{\cl}{cl}
\def\JJ{\mathbf{J}}
\def\FF{\mathbf{F}}
\def\BB{\mathbf{B}}
\DeclareMathOperator{\Alget}{Alget}
\DeclareMathOperator{\Ass}{Ass}
\DeclareMathOperator{\supp}{supp}
\DeclareMathOperator{\rad}{rad}
\DeclareMathOperator{\prof}{prof}
\DeclareMathOperator{\cint}{int}
\DeclareMathOperator{\sic}{sic}
\DeclareMathOperator{\Tor}{Tor}
\DeclareMathOperator{\tors}{tors}
\DeclareMathOperator{\Sym}{Sym}
\DeclareMathOperator{\parf}{parf}
\DeclareMathOperator{\tor}{tor}
\DeclareMathOperator{\defn}{defn}
\DeclareMathOperator{\Div}{Div}
\DeclareMathOperator{\Pol}{Pol}
\DeclareMathOperator{\Fex}{Fex}
\DeclareMathOperator{\Lie}{Lie}
\DeclareMathOperator{\tot}{tot}
\DeclareMathOperator{\Aut}{Aut}
\DeclareMathOperator{\HOM}{HOM}
\DeclareMathOperator{\pch}{pch}
\DeclareMathOperator{\Ch}{Ch}
\DeclareMathOperator{\EXT}{EXT}
\DeclareMathOperator{\PIC}{PIC}
\DeclareMathOperator{\Rf}{Rf}
\DeclareMathOperator{\inj}{inj}
\newcommand\VV{\mathbf{V}}

\def\Xv{%
\begin{array}[b]{l}
{\kern1pt\scriptstyle \vee}\\[-5pt]
X
\end{array}}

%%%%%%%
\newcommand\TT{\mathbf{T}}
\newcommand\XX{\mathbf{X}}
\newcommand\EE{\mathbf{E}}
%%%%%%

\def\qedsymbol{\textup{C.Q.F.D.}}
\let\CQFD\qed
\let\cqfd\qed

%%%%%\bigdot
\font\bigsymb=cmsy10 at 4pt
\def\bigdot{{\kern1.2pt\raise 1.5pt\hbox{\bigsymb\char15}}}
%%%%%

\def\Lright{\mathop{\vcenter{\hbox{\hbox{$\longrightarrow$}\lower8pt \hbox{\kern-16pt\hbox{$\longrightarrow$}}}}}}

\def\Aa{{\mathcal A}}
\def\Bb{{\mathcal B}}
\def\Cc{{\mathcal C}}
\def\Dd{{\mathcal D}}
\def\Ee{{\mathcal E}}
\def\Ff{{\mathcal F}}
\def\Gg{{\mathcal G}}
\def\Hh{{\mathcal H}}
\def\Ii{{\mathcal I}}
\def\Jj{{\mathcal J}}
\def\Kk{{\mathcal K}}
\def\Ll{{\mathcal L}}
\def\Mm{{\mathcal M}}
\def\Nn{{\mathcal N}}
\def\Oo{{\mathcal O}}
\def\Pp{{\mathcal P}}
\def\Qq{{\mathcal Q}}
\def\Rr{{\mathcal R}}
\def\Ss{{\mathcal S}}
\def\Tt{{\mathcal T}}
\def\Uu{{\mathcal U}}
\def\Vv{{\mathcal V}}
\def\Ww{{\mathcal W}}
\def\Xx{{\mathcal X}}
\def\Yy{{\mathcal Y}}
\def\Zz{{\mathcal Z}}

\newcommand{\CC}{\mathbf{C}}
\newcommand{\GG}{\mathbf{G}}
\newcommand{\KK}{\mathbf{K}}
\newcommand{\LL}{\mathbf{L}}
\newcommand{\NN}{\mathbf{N}}
\newcommand{\PP}{\mathbf{P}}
\newcommand\QQ{\mathbf{Q}}
\newcommand\RR{\mathbf{R}}
\newcommand\ZZ{\mathbf{Z}}

\def\to{\mathchoice{\longrightarrow}{\rightarrow}{\rightarrow}{\rightarrow}}
\def\from{\mathchoice{\longleftarrow}{\leftarrow}{\leftarrow}{\leftarrow}}
\def\mto{\mathchoice{\longmapsto}{\mapsto}{\mapsto}{\mapsto}}
\def\mfrom{\mathrel{\from\joinrel\mapstochar\,}}
\def\hto{\mathrel{\lhook\joinrel\to}}
\def\hfrom{\mathrel{\from\joinrel\rhook}}

\def\implies{\mathchoice{\Longrightarrow}{\Rightarrow}{\Rightarrow}{\Rightarrow}}
\def\implied{\mathchoice{\Longleftarrow}{\Leftarrow}{\Leftarrow}{\Leftarrow}}
\def\iff{\mathchoice{\Longleftrightarrow}{\Leftrightarrow}{\Leftrightarrow}{\Leftrightarrow}}

\def\To#1{\mathchoice{\xrightarrow{\labelstyle~~#1~}}{\stackrel{#1}{\longrightarrow}}{}{}}
\def\Mto#1{\mathrel{\mapstochar\To{#1}}}
\def\From#1{\mathchoice{\xleftarrow{\labelstyle~~#1~}}{\stackrel{#1}{\longleftarrow}}{}{}}

\newcommand{\isomto}{\overset{\sim}{\longrightarrow}}
\newcommand{\isomfrom}{\overset{\sim}{\longleftarrow}}

\newcommand{\IFF}{\unskip~$\iff$~\ignorespaces}
\newcommand{\THEN}{\unskip~$\implies$~\ignorespaces}

\let\leq\leqslant
\let\geq\geqslant
\let\le\leqslant
\let\ge\geqslant

\def\cf{{cf.}\kern.3em}
\def\Cf{{Cf.}\kern.3em}
\def\ie{{i.e.}\kern.3em}
\def\iev{{i.e.},\xspace}
%\def\loccito{loc.\kern2pt cit.\xspace}
%\def\loccitn{\textit{loc.\kern2pt cit.}\xspace}
%\def\loccit{\sisi{\loccito}{\loccitn}}
\def\EGA{\textup{EGA}\kern.2em}
\def\SGA{\textup{SGA}\kern.3em}

\let\oldS\S
\def\S{\oldS\kern .15em}
\def\ptbl{.\kern .2em}
\def\resp{resp.\kern.3em}
\def\Exp{Exp.\kern.3em}

\let\bar\overline
\let\hat\widehat
\let\tilde\widetilde
\let\dpl\displaystyle
\let\ldots\dots
\let\emptyset\varnothing
\let\skpt\skippointrait
\let\epsilon\varepsilon

\def\sfrac#1#2{{#1}/{#2}}

\DeclareMathOperator{\cd}{cd}
%\DeclareMathOperator{\Hhom}{\mathcal{H}\!\mathit{om}}
\DeclareMathOperator{\Hhom}{\mathbf{Hom}}
%\DeclareMathOperator{\Eext}{\mathcal{E}\!\mathit{xt}}
\DeclareMathOperator{\Eext}{\mathbf{Ext}}
\DeclareMathOperator{\Spec}{Spec}
\DeclareMathOperator{\Tr}{Tr}

\let\emptyset=\varnothing

%%%%%%%%%%%%%%%
%%% expose XVII
%%%%%%%%%%%%%%%

\makeatletter
\let\old@optexpose\optexpose
\newcommand{\empty@optexpose}[1]{}
\newcommand{\paren@optexpose}[1]{\old@optexpose{#1}(}
\newcommand{\refnoexpose}[1]{\let\optexpose\empty@optexpose\ref{#1}\let\optexpose\old@optexpose}%\refnoexpose{X.cor5.2} engendre "5.2" au lieu de "X 5.2"
\newcommand{\neweqref}[1]{\let\optexpose\paren@optexpose\textup{\ref{#1})}\let\optexpose\old@optexpose}%\neweqref{XII.eq4.1} engendre "XII (4.1)" au lieu de "(XII 4.1)"
\makeatother
\newcommand{\refvb}[2][VI${}_{\text{B}}$ ]{\sisi{#1\refnoexpose{#2}}{\ref{#2}}}

\newcommand{\leftexp}[2]{{\vphantom{#2}}^{#1}{#2}}
\newcommand{\AAA}{\mathbf{A}}
\newcommand{\Sf}{\mathfrak{S}}
\DeclareMathOperator*{\quotprojlim}{\text{``$\varprojlim$''}}
\DeclareMathOperator*{\quotinjlim}{\text{``$\varinjlim$''}}
\DeclareMathOperator{\Fl}{F\sisi{\ell}{l}}
\DeclareMathOperator{\Pl}{P\sisi{\ell}{l}}
%\DeclareMathOperator{\cl}{c\sisi{\ell}{l}}
\DeclareMathOperator{\tdf}{tdf}
\DeclareMathOperator{\rg}{rg}
\newcommand{\iInd}{\sisi{I}{i}nd}
\newcommand{\iind}{ind}
\newcommand{\ppro}{pro}
\newcommand{\tTor}{\sisi{T}{t}or}
\newcommand{\ttor}{tor}
\newcommand{\chap}{Chap.}
\newcommand{\Ffex}{\Ff\mathit{ex}}
\newcommand{\sheafk}{\underline{k}}
\newcommand{\sheafn}{\underline{n}}
\newcommand{\sheafGamma}{\underline{\Gamma}}

% Repeated references to an already typeset footnote must not use
% \footnotemark without a matching \footnotetext: Hyperref would then create
% a live annotation to a destination that cannot exist.  This helper keeps the
% current printed footnote mark and links it to an explicit stable note target.
\makeatletter
\newcommand{\sgarepeatedfootnotemark}[1]{%
  \begingroup
  \edef\@thefnmark{\thefootnote}%
  \hyperlink{#1}{\@makefnmark}%
  \endgroup}
\makeatother

\sisi{
\DeclareMathOperator{\simarrow}{\rule{.15mm}{1.8mm}\kern -3.4pt\rightsquigarrow}
\newcommand{\sheafell}{\underline{\ell}}
\newcommand{\sheafZZ}{\underline{\ZZ}}
}{
\newcommand{\nderef}[1]{la N.D.E.~(\ref{#1}) page~\pageref{#1}}
}

%%%%%%%%%%%%%%%%%
%%%%% expose XVII
%%%%%%%%%%%%%%%%%

% END INCLUDED SGA4_English_current_buildable_source_20260906/source/fichiers-configuration-latex/macros.tex


\externaldocument{../01/01_English.tex}
\externaldocument{../02/02_English.tex}
\externaldocument{../03/03_English.tex}
\externaldocument{../04/04_English.tex}
%\externaldocument{../05/05_English.tex}

\setboolean{original}{false}

\begin{document}

\begin{center}
V. Cohomology in Topoi\\
J.-L. Verdier
\end{center}

\selectlanguage{english}version: \showgitstatus
\tableofcontents

\else
\chapter{Cohomology in Topoi}\MakeLinkTarget*{sga4.exp.V}\label{sga4.exp.V}
%\addtocontents{toc}{J.-L. Verdier}
\begin{center}
J.-L. Verdier
\end{center}
\fi

\setcounter{pageoriginal}{1}

% BEGIN INCLUDED SGA4_English_current_buildable_source_20260906/source/05/drafts/SGA4_V_LINES_57_369_English.texfrag
\section*{Introduction}

In\pageoriginale this exposé we present the commutative and
elementary cohomological invariants of topoi. In ${\rm n^{\circ}}$
1, we study flat modules and flat morphisms of ringed topoi. The
proofs are given using the hypothesis, most often satisfied in
practice, that the topoi have enough points (IV 6). These proofs are
repeated in the general case in Appendix nᵒ 8, where Deligne, by
means of the technique of local inductive limits, also generalizes
to topoi D. Lazard's theorem on the structure of flat modules. The
theorems in this exposé\sisi{,}{} are existence theorems for spectral
sequences relating the various cohomological invariants (nᵒ 3, 5,
6). It is known that, even for topological spaces, Čech cohomology
does not in general coincide with sheaf cohomology \cite{sga4.exp.V.bib.11}. In
Appendix nᵒ 7, we introduce a modified Čech computation which makes
it possible to obtain the sheaf cohomology of an arbitrary topos by
means of coverings. In that appendix we are led to use simplicial
coverings (hypercoverings), whose homotopy invariants were studied
in \cite{sga4.exp.V.bib.1} (\cf also \cite{sga4.exp.V.bib.17}).

The cohomological invariants introduced here are elementary in the
sense that we do not use derived categories \cite{sga4.exp.V.bib.12}. A reader
familiar with this language will immediately translate the various
statements of this exposé and will then be able to generalize them
to complexes and hypercohomology. The language of derived
categories is, moreover, used later in this seminar.

Here we restrict ourselves to commutative cohomology. For
noncommutative $H^1$, which is used in this seminar, and for
noncommutative $H^2$, we refer to \cite{sga4.exp.V.bib.9}. The functors that we
derive are additive; we say nothing about multiplicative structures
(\cf \cite{sga4.exp.V.bib.7}).

\setcounter{section}{-1}
\section{Generalities on abelian categories}\etiquette[0]{V.sec0}

In\pageoriginale this section we recall several lemmas, most of
which can be found in \cite{sga4.exp.V.bib.11}.

\begin{enonce*}{Proposition 0.1}\etiquette[0.1]{V.prop0.1}
Let $\mathfrak{A}$ be an abelian category having a generator. The
following conditions are equivalent:
\begin{enumerate}
\renewcommand{\theenumi}{\roman{enumi}}
\renewcommand{\labelenumi}{\theenumi)}
\item The category $\mathfrak{A}$ satisfies axiom $AB$ $5$): Small
  direct sums are representable, and if $(X_i)$, $i \in I$, is a
  small increasing filtered family of subobjects of an object $X$
  of $\mathfrak{A}$ and $Y$ is a subobject of $X$, then
$$
(\displaystyle{\sup_i}\ X_i) \cap Y = \displaystyle{\sup_i}\ (X_i \cap Y) {\quoi}.
$$
\item Small pseudo-filtered inductive limits \emph{(\hyperref[I.def2.7]{I.2.7})} %% label \XXX
are representable and commute with finite projective limits.
\end{enumerate}

Moreover, if the preceding conditions hold, small filtered
inductive limits are universal \emph{(I \hyperref[I.prop2.6]{2.6})}.
\end{enonce*}

\begin{enonce*}[remark]{Proof}
It is clear that (ii) $\Rightarrow$ (i). To show that
\sisi{i)}{(i)} $\Rightarrow$ (ii), and to prove the additional
assertion, it suffices to use the fact that $\mathfrak{A}$ is a full
subcategory of a category of modules $\Mm$ over a suitable ring,
such that the inclusion functor $u : \mathfrak{A} \rightarrow \Mm$
admits an exact left adjoint $v$ \cite{sga4.exp.V.bib.5}. The verification is then
immediate.
\end{enonce*}

\setcounter{subsection}{1}
\subsubsection{}\etiquette[0.1.1]{V.subsubsec0.1.1}

It is known \cite{sga4.exp.V.bib.11} that an abelian category $\mathfrak{A}$
having a generator and satisfying axiom $AB$ 5) has enough
injectives, that is, every object embeds in an injective object.
Moreover, by the result already cited \cite{sga4.exp.V.bib.5}, small products are
representable in $\mathfrak{A}$ (axiom $AB$ 3) $^{\ast}$).

\begin{enonce*}{Proposition 0.2}\etiquette[0.2]{V.prop0.2}
Let $\mathfrak{A}$ and $\mathfrak{B}$ be two abelian categories and
let
$\xymatrix{\mathfrak{A} \ar@<-2pt>[r]_{u}& \mathfrak{B} \ar@<-2pt>[l]_{v}
}
$
be two adjoint functors ($u$ is left adjoint to $v$). Consider the
following two properties:
\begin{enumerate}
\renewcommand{\theenumi}{\roman{enumi}}
\renewcommand{\labelenumi}{\theenumi)}
\item The functor $u$ is exact.

\item The\pageoriginale functor $v$ sends injective objects of
$\mathfrak{B}$ to injective objects of $\mathfrak{A}$.
\begin{enumerate}
\renewcommand{\theenumii}{\arabic {enumii}}
\renewcommand{\labelenumii}{\theenumii)}
\item The implication \emph{(i) $\Rightarrow$ (ii)} always holds.

If, in addition, every nonzero object of $\mathfrak{B}$ is the
source of a nonzero morphism into an injective object, then
\emph{(ii) $\Leftrightarrow$ (i)}.
\item Suppose that:
\begin{enumerate}
\renewcommand{\theenumiii}{\alph {enumiii}}
\renewcommand{\labelenumiii}{\theenumiii)}
\item the category $\mathfrak{B}$ has enough injectives;
\item one of the two equivalent conditions {\rm(i)} and {\rm(ii)}
  above holds;
\item the functor $u$ is faithful.
\end{enumerate}
\end{enumerate}
\end{enumerate}
Then the category $\mathfrak{A}$ has enough injectives.
\end{enonce*}

\begin{enonce*}[remark]{Proof}
The proof is left to the reader.
\end{enonce*}

\begin{enonce*}[remark]{Remark 0.2.1}\etiquette[0.2.1]{V.rem0.2.1}
The category of commutative groups has enough injectives. Applying
Proposition \hyperref[V.prop0.2]{0.2}, we deduce that every category
of unital modules over a ring with identity has enough injectives.
Applying the result of \cite{sga4.exp.V.bib.5} (used in the proof of
\hyperref[V.prop0.1]{0.1}) together with
\hyperref[V.prop0.2]{0.2}, we then deduce that every abelian category
having a generator and exact small filtered inductive limits has
enough injectives; this gives a new proof of that fact.
\end{enonce*}

\begin{enonce*}{Proposition 0.3}\etiquette[0.3]{V.prop0.3}
Let $\mathfrak{A}$, $\mathfrak{B}$, and $\mathfrak{C}$ be three
abelian categories, and let $u : \mathfrak{A} \rightarrow
\mathfrak{B}$ and $v : \mathfrak{B} \rightarrow \mathfrak{C}$ be
two left exact additive functors. Suppose that $\mathfrak{A}$ and
$\mathfrak{B}$ have enough injective objects. The following two
properties are equivalent:
\begin{enumerate}
\renewcommand{\theenumi}{\roman{enumi}}
\renewcommand{\labelenumi}{\theenumi)}
\item There exists a spectral functor whose $E^{p, q}_2$ term is
$$
R^p vR^qu
$$
and which converges to $R^{p+q} vu$ (suitably filtered).
\item The functor $u$ sends injective objects of $\mathfrak{A}$ to
  objects that are acyclic for the functor $v$.
\end{enumerate}
\end{enonce*}

\begin{enonce*}[remark]{Proof}
(i) $\Rightarrow$ (ii) is immediate, since it suffices to apply the
spectral functor to an\pageoriginale injective object. The
implication (ii) $\Rightarrow$ (i) is proved in \cite{sga4.exp.V.bib.11}.
\end{enonce*}

\begin{enonce*}{Proposition 0.4}\etiquette[0.4]{V.prop0.4}
Let $\mathfrak{A}$ and $\mathfrak{B}$ be two abelian categories and
let $u : \mathfrak{A} \rightarrow \mathfrak{B}$ be a %% correct word break?
left exact additive functor. Let $M$ be a subset of the set of
objects of $\mathfrak{A}$ having the following properties:
\begin{enumerate}
\renewcommand{\labelenumi}{\theenumi)}
\item Every object of $\mathfrak{A}$ embeds in an element of $M$.

\item If $X \oplus Y$ belongs to $M$, then the object $X$ belongs to
  $M$.

\item If
$$
0 \longrightarrow X' \longrightarrow X \longrightarrow X''
\longrightarrow 0
$$
is an exact sequence and $X'$ and $X$ belong to $M$, then $X''$
belongs to $M$, and the sequence
$$
0 \longrightarrow u(X') \longrightarrow u(X) \longrightarrow u(X'')
\longrightarrow 0
$$
is exact. Zero objects belong to $M$.
\end{enumerate}

Then every injective belongs to $M$, and the objects of $M$ are
acyclic for $u$, that is, for every $q \neq 0$ and every object $X$
of $M$, one has $R^q u(X) = 0$. (In particular, resolutions by
objects of $M$ can be used to compute the derived functors of $u$.)
\end{enonce*}

For the proof\sisi{}{,} see \cite{sga4.exp.V.bib.11} 3.3.1.

\begin{enonce*}{Proposition 0.5}\etiquette[0.5]{V.prop0.5}
Let $𝒰 \subset 𝒱$ be two universes, let $\mathfrak{A}$ (\resp
$\mathfrak{B}$) be a $𝒰$-category (\resp $𝒱$-category) which is
abelian, satisfies axiom $AB$ 5) relative to $𝒰$ (\resp to $𝒱$),
and has a $𝒰$-small (\resp $𝒱$-small) generating family. Let
$\epsilon : \mathfrak{A} \to \mathfrak{B}$ be an exact fully
faithful functor. The following conditions are equivalent:
\begin{itemize}
\item[1)] There exists a generating family $(X_i)_{i\in I}$ of
  $\mathfrak{A}$ such that the family $(\epsilon(X_i))_{i\in I}$ is
  generating in $\mathfrak{B}$.

\item[1′)] Every object of $\mathfrak{B}$ is isomorphic to a
  quotient of an object of the form
  $\displaystyle{⨁_{\alpha \in A}} \epsilon (Y_\alpha)$, where
  $A \in 𝒱$.

Under these conditions, the following property holds:

\item[2)] $\epsilon$ sends $𝒰$-small products to products (and thus
  commutes with $𝒰$-small projective limits).
\end{itemize}

Moreover, under the equivalent conditions 1) or 1′), the following
conditions are equivalent\pageoriginale:
\begin{itemize}
\item[a)] For every object $Y$ of $\mathfrak{A}$, every subobject of
  $\epsilon(Y)$ is isomorphic to the image under $\epsilon$ of a
  subobject of $Y$.

\item[a′)] There exists a generating family $(X_i)_{i \in I}$ of
  $\mathfrak{A}$ such that the family $(\epsilon(X_i))_{i \in I}$
  is generating in $\mathfrak{B}$ and such that, for every $i \in
  I$, every subobject of $\epsilon(X_i)$ is isomorphic to the image
  under $\epsilon$ of a subobject of $X_i$.

\item[b)] Every object of $\mathfrak{B}$ is isomorphic to a
  subobject of an object of the form
  $\prod_{\alpha \in A} \epsilon(Y_\alpha)$, where $A \in 𝒱$.

\item[c)] $\epsilon$ commutes with $𝒰$-small direct sums (and thus
  commutes with $𝒰$-small inductive limits).

Moreover, under conditions 1) and a′), one has:

\item[d)] $\epsilon$ sends injective objects to injective objects.
\end{itemize}
\end{enonce*}

\begin{enonce*}[remark]{Remark 0.5.1}\etiquette[0.5.1]{V.rem0.5.1}
When in \hyperref[V.prop0.5]{0.5} one has $𝒰 = 𝒱$, conditions 1) and
a′) imply that $\epsilon$ is an equivalence of categories (since b),
2), and a) then hold).
\end{enonce*}

\setcounter{subsection}{5}
\subsubsection{}\etiquette[0.5.2]{V.subsubsec0.5.2}
We will confine ourselves to giving indications for the proof. The
implications 1) $\Leftrightarrow$ $1'$) and 1) $\Rightarrow$ 2) are
left to the reader. It is clear that a) $\Rightarrow$ a′). We show
that a′) implies d). Let $M$ be injective in $\mathfrak{A}$. For
every $i \in I$ and every subobject $Y \hookrightarrow X_i$ of
$X_i$, the homomorphism $\Hom(X_i, M) \rightarrow \Hom(Y, M)$ is
surjective. Consequently, by a′) and the full faithfulness of
$\epsilon$, for every subobject $U \hookrightarrow \epsilon(X_i)$,
the homomorphism $\Hom(\epsilon (X_i), \epsilon(M)) \to \Hom(U,
\epsilon(M))$ is surjective. Since the $\epsilon(X_i)$ form a
generating family of $\mathfrak{B}$, $\epsilon(M)$ is injective
\cite{sga4.exp.V.bib.11}.

We show that a′) $\Rightarrow$ b). By enlarging the family of the
$X_i$ if necessary, we may suppose that, for every $i \in I$, every
quotient of $X_i$ is isomorphic to an $X_j$ for a suitable $j$. For
every $i \in I$, let $X_i \hookrightarrow M_i$ then be a
monomorphism into an injective object. The family
$(\epsilon (X_i))_{i\in I}$ is stable under quotients, and for $i
\in I$, the morphism $\epsilon(X_i) \hookrightarrow
\epsilon(M_i)$ is, by d), a monomorphism into an injective object.
It is then immediately verified that, since the family
$\epsilon(X_i)$ is generating, the family $(\epsilon(M_i))_{i \in
I}$ is cogenerating, whence b).

We show that b) $\Rightarrow$ c). Let $(Z_\alpha)_{\alpha \in A}$
be a $𝒰$-small family of objects of $\mathfrak{A}$, and let us show
that the canonical morphism
$$
⨁_{\alpha\in A} \epsilon(Z_\alpha) \to \epsilon(⨁_{\alpha \in A} Z_\alpha)
$$
is\pageoriginale an isomorphism. Since the family of objects
$\epsilon(Y)$ is cogenerating, it suffices to show that, for every
object $Y$ of $\mathfrak{A}$, the homomorphism
$$
\Hom(\epsilon(⨁_a Z_a), \epsilon(Y)) \to
\Hom(⨁_a \epsilon(Z_a), \epsilon(Y))
$$
is an isomorphism, which follows from the full faithfulness of
$\epsilon$.

It remains to show that c) $\Rightarrow$ a). Let $Z$ be a subobject
of $\epsilon(Y)$. By $1'$), there exists a $𝒰$-small family
$Y_\alpha$ of objects of $\mathfrak{A}$ and an epimorphism from
$\bigoplus_{\alpha} \epsilon(Y_\alpha)$ onto $Z$. Since $\epsilon$
commutes with $𝒰$-small direct sums, there is therefore an
epimorphism $\epsilon(Y') \to Z$. Denote by $u : \epsilon(Y') \to Z
\to \epsilon(Y)$ the composite morphism. We have $Z =
\text{Im}(u)$. Since \sisi{}{$\epsilon$} is fully faithful, we have
$u \sisi{:}{=} \epsilon(v)$, and consequently
$Z = \text{Im}(\epsilon(\sisi{V}{v})) = \epsilon(\text{Im}(v))$.


\begin{enonce*}[remark]{Exercise 0.5.2}\etiquette[0.5.2]{V.exe0.5.2}
Let $\Sex_{𝒱} (\mathfrak{A})$ be the category of contravariant
functors from $\mathfrak{A}$ to the category of $𝒱$-commutative
groups which commute with $𝒰$-small inductive limits. Show that,
under conditions 1) and a′), the canonical functor from
$\mathfrak{B}$ to $\Sex_{𝒱}(\mathfrak{A})$ is an equivalence.
\end{enonce*}


% END INCLUDED SGA4_English_current_buildable_source_20260906/source/05/drafts/SGA4_V_LINES_57_369_English.texfrag


% BEGIN INCLUDED SGA4_English_current_buildable_source_20260906/source/05/drafts/SGA4_V_LINES_370_602_English.texfrag
\section{Flat modules}\etiquette[1]{V.sec1}

\begin{enonce*}{Definition 1.1}\etiquette[1.1]{V.def1.1}
Let $(E,A)$ be a ringed topos \emph{(IV \hyperref[IV.subsubsec11.1.1]{11.1.1})}. %% label
A right (\resp left) $A$-module $M$ is said to be flat if the functor
$M\otimes_A⋅$ (\resp $⋅\otimes_A M$) from the category of left
(\resp right) $A$-modules to the category of abelian sheaves of
$E$\sisi{,}{} is exact.
\end{enonce*}

\begin{enonce*}{Proposition 1.2}\etiquette[1.2]{V.prop1.2}
Let $M$ be a $B$-$A$-bimodule.
\begin{enumerate}
\renewcommand{\labelenumi}{\theenumi)}
\item The following properties are equivalent:
\renewcommand{\theenumii}{\roman{enumii}}
\renewcommand{\labelenumii}{\theenumii)}
\begin{enumerate}
\item The module $M$ is flat as a right $A$-module.
\item For every injective $B$-module $I$, the left $A$-module
  $\Hhom_B(M,I)$ is injective.
\end{enumerate}
\item A module $M$ which is a pseudo-filtered inductive limit
\emph{(I \hyperref[I.exem2.7.1]{2.7.1})} %% label
of flat modules is flat.

\item Finally, if
$M_∙=\cdots\sisi{}{\to}M_{i+1}\to M_i\sisi{}{\to}\cdots$ is an
acyclic complex of flat modules $(M_i=0$ for $i<i_0)$, then, for every
module $F$, the complex
$$
M_∙\otimes_A F=\cdots\sisi{}{\to}M_{i+1}\otimes_A F\to
M_i\otimes_A F\sisi{}{\to}\cdots
$$
is\pageoriginale acyclic.
\end{enumerate}
\end{enonce*}

\begin{enonce*}[remark]{Proof}
By (IV \hyperref[IV.cor12.12]{12.12}), %% label
there is an adjunction isomorphism:
$$
\Hom_B(M\otimes_A\mathord\cdot,\mathord\cdot)\xrightarrow{\sim}
\Hom_A(\mathord\cdot,\Hhom_B(M,\mathord\cdot)){\quoi}.
$$
It is then enough to apply \hyperref[V.prop0.2]{0.2} to obtain the
equivalence (i) $\Leftrightarrow$ (ii). This adjunction isomorphism
also shows that the tensor product commutes with inductive limits.
The fact that pseudo-filtered inductive limits are exact
(\hyperref[V.prop0.1]{0.1}) implies the second assertion. To show that
the complex $M_∙\otimes_A F$ is acyclic, it is enough to show that,
for every injective abelian sheaf $I$, the complex
$\Hom^{∙}_{\ZZ}(M_∙\otimes_A F,I)$ is acyclic. By the adjunction
formulas, this complex is isomorphic to
$\Hom^{∙}_A(F,\Hhom_{\ZZ}(M_∙,I))$. But, by the equivalence
(i) $\Leftrightarrow$ (ii), the complex of sheaves
$$
\Hhom_{\ZZ}(M_∙,I)
$$
is an acyclic complex whose objects are injective, whence the
conclusion.
\end{enonce*}

\begin{enonce*}{Proposition 1.3.1}\etiquette[1.3.1]{V.prop1.3.1}
Let $(E,A)$ be a ringed topos, let $H$ be an object of $E$, and let
$M$ be a flat $A_{/H}$-module. Then $j_{H!}M$ is a flat $A$-module.
In particular, $A_H$ is flat on both the right and the left.
\end{enonce*}

Suppose, to fix ideas, that $M$ is a right $A_{/H}$-module. For every
left $A$-module $P$, there is a canonical isomorphism (IV 12) %% label
$$
j_{H!}M\otimes_A P\simeq
j_{H!}(M\otimes_{A_{/H}}j_H^{\ast}P){\quoi}.
$$
The functors $j_{H!}$ and $j_H^{\ast}$ are exact (IV \hyperref[IV.prop11.3.1]{11.3.1} and
\hyperref[IV.subsubsec11.2.2]{11.12.2}), and, by hypothesis, the functor
$M\otimes_{A_{/H}}-$ is exact. Consequently, the functor
$P\longmapsto j_{H!}M\otimes_A P$ is exact, and $j_{H!}M$ is flat.

\begin{enonce*}{Proposition 1.3.2}\etiquette[1.3.2]{V.prop1.3.2}
(Projection formula for closed immersions):

Let $(E,A)$ be a ringed topos, and let $i:F\to E$ be a closed subtopos
of $E$. Set $i^{\ast}A=A_{/F}$. For every right $A_{/F}$-module $M$
and every left $A$-module $P$, there\pageoriginale is a functorial
isomorphism
$$
i_\ast(M\otimes_{A_{/F}}i^\ast P)\simeq(i_\ast M)\otimes_A P.
$$
\end{enonce*}

Let $U$ be the open complement of $F$, and let $j:U\to E$ be the
canonical immersion morphism. One has
$j^{\ast}(i_\ast M\otimes_A P)\simeq
0\otimes_{A_{/U}}j^\ast(P)$ (IV 12); consequently,
$i_\ast M\otimes_A P$ has its support in $F$. Hence the adjunction
morphism
$i_\ast M\otimes_A P\to i_\ast i^\ast(i_\ast M\otimes_A P)$ is an
isomorphism (IV 14). One has
$i^{\ast}(i_\ast M\otimes_A P)\simeq
i^\ast i_\ast M\otimes_{A_{/F}}i^\ast P$ (IV 12), and, since
$i:F\to E$ is a closed immersion,
$i^{\ast}i_\ast M\simeq M$ (IV 14); this gives the asserted
isomorphism.

\begin{enonce*}{Corollary 1.3.3}\etiquette[1.3.3]{V.cor1.3.3}
For every flat $A_{/F}$-module $M$, $i_\ast M$ is flat.
\end{enonce*}

It follows from \hyperref[V.prop1.3.2]{1.3.2} and from (IV 14) that
the functor $P\longmapsto(i_\ast M)\otimes_A P$ is exact.

\setcounter{subsection}{3}
\subsection{}\etiquette[1.4]{V.subsec1.4}
Let $(E,A)$ be a ringed topos, let $x:P\to E$ be a point of $E$
(IV \hyperref[IV.def6.1]{6.1}), and let $\Vois(x)$ be the category of neighborhoods of $x$
(IV \hyperref[IV.subsec6.8]{6.8}). To every object $V$ of $\Vois(x)$ there corresponds an
object of $E$, also denoted by $V$, and a point
$x_V:P\to E/V$ of $E/V$. Moreover, to every morphism
$u:V\to W$ in $\Vois(x)$ there corresponds an essentially
commutative diagram of morphisms of topoi (IV \hyperref[IV.subsec6.7]{6.7})
\[\vcenter{
\xymatrix@C=2.5cm@R=1.5cm{
P \ar[r]^{x_V} \ar[rd]_{x_W}& E/V \ar[d]^{j_u}\\
& E/W\quad .
}}\leqno{(1.4.1)}
\]\etiquette[1.4.1]{V.eq1.4.1}
Let $N$ be an $A_x$-module (\resp a set). From the diagram
\eqhyperref[V.eq1.4.1]{1.4.1}, one obtains a diagram of $A_x$-modules
(\resp sets):
\[\vcenter{
\xymatrix@C=2cm@R=1.5cm{
x^{\ast}_W x_{W\ast}N \ar[r]^{\ad_W} \ar[d]_{\phi(u)}& N\\
x^{\ast}_V x_{V\ast}N \ar[ru]_{\ad_V}&
}}\leqno{(1.4.2)}
\]\etiquette[1.4.2]{V.1.4.2}
where\pageoriginale $\ad_W$ and $\ad_V$ are the adjunction morphisms.
Moreover, one immediately checks that
$\phi(uv)=\phi(v)\phi(u)$. Thus there is a functor from the filtered
category $\Vois(x)^{\circ}$ to the category of $A_x$-modules
(\resp sets), and a homomorphism
\begin{equation}
\Lambda(N):\varinjlim_{V\in\Vois(x)^{\circ}}
x^{\ast}_Vx_{V\ast}N\to N{\quoi}.\tag{1.4.3}
\end{equation}\etiquette[1.4.3]{V.1.4.3}

\begin{enonce*}{Proposition 1.5}\etiquette[1.5]{V.prop1.5}
For every $A_x$-module (\resp set) $N$, $\Lambda(N)$ is an
isomorphism.
\end{enonce*}

This proposition is a special case of a proposition due to Deligne
(\hyperref[V.prop8.2.6]{8.2.6}). Let us give a direct proof. By passing
to the underlying sets, it is enough to prove the proposition when
$N$ is a set (I \hyperref[I.prop2.8]{2.8}). %% label
The cofiltered category $\mathrm{Fl}(\Vois(x))$ of arrows of
$\Vois(x)$ is fibered over the category $\Vois(x)$ by the functor
$p:\mathrm{Fl}(\Vois(x))\to\Vois(x)$ which assigns to an arrow its
target. Let $D$ be a category and
$F:\mathrm{Fl}(\Vois(x))\to D$ a pro-object (I \hyperref[I.sec8.10]{8.10}). %% label
For every object $V$ of $\Vois(x)$, denote by $F_V$ the pro-object
obtained by restricting the functor $F$ to the fiber category
$\Vois(x)/V$. To every morphism $m:U\to V$ of $\Vois(x)$, the
base-change functor by $m$ associates a morphism from the pro-object
$F_U$ to the pro-object $F_V$, and the canonical morphisms of
pro-objects $F\to F_V$ determine a morphism of pro-objects
\begin{equation*}
F\to\varprojlim_{V\in\Vois(x)^{\circ}}F_V\tag{1.5.1}
\end{equation*}\etiquette[1.5.1]{V.1.5.1}
which is immediately checked to be an isomorphism. Apply this
observation to the pro-object of pointed sets
\begin{equation*}
(U,V,m)\longmapsto x^{\ast}_V(m)=F(U,V,m){\quoi}.\tag{1.5.2}
\end{equation*}\etiquette[1.5.2]{V.1.5.2}
One obtains an isomorphism of pro-objects
\begin{equation*}
F\xrightarrow{\sim}\varprojlim_{\Vois(x)}
\varprojlim_{\Vois(x)/V}x^{\ast}_V(m){\quoi};\tag{1.5.3}
\end{equation*}\etiquette[1.5.3]{V.1.5.3}
hence, for every set $N$, a bijection
\begin{equation*}
\varinjlim_{\mathrm{Fl}(\Vois(x))^{\circ}}
\Hom(x^{\ast}_V(m),N)\simeq
\varinjlim_{\Vois(x)^{\circ}}
\varinjlim_{(\Vois(x)/V)^{\circ}}
\Hom(x^{\ast}_V(m),N).\tag{1.5.4}
\end{equation*}\etiquette[1.5.4]{V.1.5.4}
But,\pageoriginale for fixed $V$,
$\varinjlim_{(\Vois(x)/V)^{\circ}}\Hom(x^{\ast}_V(m),N)$ is
identified with $x^{\ast}_Vx_{V\ast}N$. Indeed, by IV \hyperref[IV.eq6.8.1]{6.8.1} one has
$x^{\ast}_Vx_{V\ast}N\simeq
\varinjlim_{\Vois(x_V)^\circ}\Hom_{E/V}(W,x_{V\ast}N)$, and hence,
by adjunction,
$x^{\ast}_Vx_{V\ast}N\simeq
\varinjlim_{\Vois(x_V)^\circ}\Hom(x^{\ast}_V(W),N)$; moreover, the
category $\Vois(x_V)$ is equivalent to the category $\Vois(x)/V$
(IV \hyperref[IV.eq6.7.2]{6.7.2}). Finally, the canonical transition maps
\[
\varinjlim_{(\Vois(x)/U)^{\circ}}\Hom(x^{\ast}_U(m),N)\to
\varinjlim_{(\Vois(x)/V)^{\circ}}\Hom(x^{\ast}_V(n),N)
\]
in
\eqhyperref[V.1.5.4]{1.5.4} are identified with the canonical maps
$x^{\ast}_Ux_{U\ast}N\to x^{\ast}_Vx_{V\ast}N$ of
\eqhyperref[V.1.4.3]{1.4.3}, as the reader may verify. Thus there is a
bijection
\begin{equation*}
\varinjlim_{\mathrm{Fl}(\Vois(x))^{\circ}}
\Hom(x^{\ast}_V(m),N)\simeq
\varinjlim_{\Vois(x)^{\circ}}x^{\ast}_Vx_{V\ast}N\tag{1.5.5}
\end{equation*}\etiquette[1.5.5]{V.1.5.5}
and the map
$\Lambda(N):\varinjlim_{\Vois(x)^{\circ}}x^{\ast}_Vx_{V\ast}N\to N$
of \eqhyperref[V.1.4.3]{1.4.3}, composed with the bijection
\eqhyperref[V.1.5.5]{1.5.5}, comes from the maps
$\Hom(x^{\ast}_V(m),N)\to N$ which assign to a map
$r:x^{\ast}_V(m)\to N$ the image under $r$ of the distinguished point
of $x^{\ast}_V(m)$. To prove the proposition, it is then enough to
observe that every object $(U,V,m)$ of
$\mathrm{Fl}(\Vois(x))$ is bounded below by $(U,U,\id_U)$, and that
$x^{\ast}_U(\id_U)$ consists of a single element or, in other words,
that the canonical morphism from the constant pro-object consisting
of a single element to $F$ is an isomorphism of pro-objects.

\begin{enonce*}{Proposition 1.6}\etiquette[1.6]{V.prop1.6}
Let $(E,A)$ be a ringed topos, and let $M$ be an $A$-module.
\begin{enumerate}
\renewcommand{\labelenumi}{\theenumi)}
\item If $M$ is flat, then, for every point $x:P\to E$ of $E$, the
  $A_x$-module $M_x$ is flat.
\item Let $(x_i)_{i\in I}$ be a conservative family of points of $E$
  such that, for every $i\in I$, $M_{x_i}$ is a flat $A_{x_i}$-module.
\end{enumerate}
Then $M$ is flat.
\end{enonce*}

\begin{enonce*}{Lemma 1.6.1}\etiquette[1.6.1]{V.lem1.6.1}
Let $M$ be a flat module and let $V$ be an object of $E$. The
$A/V$-module $j^{\ast}_VM$ is flat.
\end{enonce*}

We must show that the functor
$P\longmapsto P\otimes_{A/V}j^{\ast}_VM$ is exact. Since the
extension-by-zero functor $j_{V!}$ is exact and faithful
(IV \hyperref[IV.prop11.3.1]{11.3.1}), it is enough to show that the functor
$P\longmapsto j_{V!}(P\otimes_{A/V}j^{\ast}_VM)$ is exact. There is a
functorial isomorphism\pageoriginale
$j_{V!}(P\otimes_{A/V}j^{\ast}_VM)\simeq
j_{V!}(P)\otimes_A M$ (IV \hyperref[IV.prop12.11]{12.11}); consequently, the functor
$P\longmapsto j_{V!}(P)\otimes_A M$ is exact (IV \hyperref[IV.prop11.3.1]{11.3.1}).

\setcounter{subsection}{6}
\setcounter{subsubsection}{1}

% END INCLUDED SGA4_English_current_buildable_source_20260906/source/05/drafts/SGA4_V_LINES_370_602_English.texfrag


% BEGIN INCLUDED SGA4_English_current_buildable_source_20260906/source/05/drafts/SGA4_V_LINES_603_925_English.texfrag
\subsubsection{}\etiquette[1.6.2]{V.subusbsec1.6.2}
Let us prove 1). To fix ideas, suppose that $M$ is a flat left
$A$-Module, and let $x:P\to E$ be a point of $E$. We shall show
that the functor $N\mapsto N\otimes_{A_x}M_x$ from the category of
right $A_x$-modules to the category of commutative groups is exact.
For this it suffices to show that this functor is left exact. With
the notation of \hyperref[V.subsec1.4]{1.4}, let $V$ be a
neighborhood of $x$, and let $j_V:E/V\to E$ be the localization
morphism. For every $A_x$-module $N$, we have
$x_V^\ast x_{V\ast}N\otimes_{A_x}M_x\simeq
x_V^\ast x_{V\ast}N\otimes_{A_x}x_V^\ast j_V^\ast M\simeq
x_V^\ast(x_{V\ast}N\otimes_{A/V}j_V^\ast M)$ (IV \hyperref[IV.prop12.11]{12.11}).
Consequently, the functor
$N\mapsto x_V^\ast x_{V\ast}N\otimes_{A_x}M_x$ is left exact. It
follows that the functor $N\mapsto N\otimes_{A_x}M_x$, being a
filtered inductive limit of left exact functors
(\hyperref[V.prop1.5]{1.5}), is left exact.

Let us prove 2). Let
$0\to P'\to P\to P''\to0$ be an exact sequence of $A$-Modules, and
let us show that the sequence
$0\to P'\otimes_A M\to P\otimes_A M\to P''\otimes_A M\to0$ is
exact. It suffices to show that, for every $i\in I$, the sequence
obtained by passing to the fibers at $x_i$ is exact (IV 6). But the
sequence of fibers at $x_i$ is the sequence (IV \hyperref[IV.cor13.5]{13.5})
$$
0\to P'_{x_i}\otimes_{A_{x_i}}M_{x_i}\to
P_{x_i}\otimes_{A_{x_i}}M_{x_i}\to
P''_{x_i}\otimes_{A_{x_i}}M_{x_i}\to0,
$$
and since $M_{x_i}$ is flat, this sequence is exact.

\begin{enonce*}{Proposition 1.7}\etiquette[1.7]{V.prop1.7}
Let $(E,A)$ be a ringed topos, let $u:E'\to E$ be a morphism of
topoi, and put $A'=u^{\ast}A$. Let $M$ be a right (\resp left)
$A'$-Module. The following conditions are equivalent:
\begin{enumerate}
\renewcommand{\theenumi}{\roman{enumi}}
\renewcommand{\labelenumi}{\theenumi)}
\item The functor $N\mapsto M\otimes_{A'}u^{\ast}N$ (\resp
$N\mapsto u^{\ast}N\otimes_{A'}M$) from the category of left
(\resp right) $A$-Modules to the category of abelian sheaves on $E'$
is exact.
\item $M$ is a flat $A'$-Module.
\end{enumerate}
\end{enonce*}

It is clear that ii) $\Rightarrow$ i).

Let us prove that i) $\Rightarrow$ ii). We give the proof only in the
case where $E'$ possesses a conservative family of points
$(x_i)_{i\in I}$. The general case is treated in
(\hyperref[V.prop8.2.7]{8.2.7}). In this particular case, which
covers most applications, we reduce\pageoriginale to the case where
$E'$ is the punctual topos and hence where the morphism $u$, which we
then denote by $x$, is a point of $E$
(\hyperref[V.prop1.6]{1.6} and IV \hyperref[IV.prop11.3.1]{11.3.1}). We restrict ourselves to
the case where $M$ is a right $A'$-Module. The case where $M$ is a
left $A'$-Module reduces to this one by passing to opposite rings.
Let $V$ be a neighborhood of $x$, and let $V_x$ be its fiber at $x$.
With the notation of \hyperref[V.subsec1.4]{1.4}, there is an
essentially commutative diagram of morphisms of topoi:
\[
\xymatrix@C=2.25cm@R=1.5cm{
P \ar[r]^{j} \ar[rd]_{x_V}& P/V_{x}\ar[r]^{j_{V_x}} \ar[d]^{x/V}& P\ar[d]^{x}\\
  & E/V \ar[r] & E\quad \sisi{,}{}
}
\]
where $x/V$ is the morphism deduced from $x$ by localization over
$V$ (IV \hyperref[IV.subsec5.10]{5.10}), and $j$ is the morphism deduced from the marked point
of $V_x$. Let us show that the functor
$N'\mapsto M\otimes_{A'}x_V^{\ast}N'$ is exact. We have
$x_V^{\ast}=j^{\ast}(x/V)^{\ast}$ and
$\id=j^{\ast}j_{V_x}^{\ast}$; consequently, there is a canonical
isomorphism
$$
M\otimes_{A'}x_V^{\ast}N'\simeq
j^{\ast}\bigl(j_{V_x}^{\ast}M\otimes_{A'/V_x}(x/V)^{\ast}N'\bigr).
$$
Since $j^{\ast}$ is exact and the functor $j_{V_x!}$ is exact and
faithful (IV \hyperref[IV.prop11.3.1]{11.3.1}), it suffices to show that the functor
$$
N'\longmapsto j_{V_x!}\Bigl(j_{V_x}^{\ast}M\otimes_{A'/V_x}
(x/V)^{\ast}N'\bigr)
$$
is exact; consequently (IV \hyperref[IV.prop12.11]{12.11}), it suffices to show that the
functor $N'\mapsto M\otimes_{A'}x^{\ast}j_{V!}N'$ is exact. But it
follows from the definition of the inverse image of an induced topos
(\hyperrefext[IV][IV.eq5.10.2]{5.10.2}) that
$j_{V_x!}(x/V)^{\ast}N'$ is equal to $x^{\ast}j_{V!}N'$, and since
the functor $j_{V!}$ is exact
(\hyperrefext[IV][IV.prop11.3.1]{11.3.1}), the functor
$N'\mapsto M\otimes_{A'}j_{V_x!}(x/V)^{\ast}N'$ is exact by
hypothesis.

For every $A'$-module $Q$, there is a functorial isomorphism
(\hyperref[V.prop1.5]{1.5}):
$$
Q\simeq \varinjlim_{\Vois(x)^{\circ}}x_V^{\ast}x_{V\ast}Q{\quoi}.
$$
By what precedes, the functor $Q\mapsto M\otimes_{A'}Q$ is a
filtered inductive limit of left exact functors. It is therefore
left exact, and hence $M$ is flat.

\begin{enonce*}{Corollary 1.7.1}\etiquette[1.7.1]{V.cor1.7.1}
Let $u:(E',A')\to(E,A)$ be a morphism of ringed topoi, and let $M$
be a flat $A$-Module. Then $u^{\ast}M$ is a flat $A'$-Module.
\end{enonce*}

This follows from the criterion given in \hyperref[V.prop1.7]{1.7}
and formula (IV \hyperlink{sga4.exp.IV.equation.13.4.5}{13.4.5}).

\begin{enonce*}{Corollary 1.7.2}\etiquette[1.7.2]{V.cor1.7.2}
Let\pageoriginale $u:(E',A')\to(E,A)$ be a morphism of ringed
topoi. The following conditions are equivalent:
\begin{enumerate}
\renewcommand{\theenumi}{\roman{enumi}}
\renewcommand{\labelenumi}{\theenumi)}
\item The right (\resp left) $u^{-1}(A)$-Module $A'$ is flat.
\item The functor $u^{\ast}$ from the category of left (\resp right)
$A$-Modules to the category of left (\resp right) $A'$-Modules is
exact.
\end{enumerate}
\end{enonce*}

The equivalence follows from \hyperref[V.prop1.7]{1.7} and the
definition of $u^{\ast}$ (IV \hyperref[IV.subsubsec13.2.1]{13.2.1}).

\begin{enonce*}{Definition 1.8}\etiquette[1.8]{V.def1.8}
A morphism of ringed topoi having the equivalent properties of
\hyperref[V.cor1.7.2]{1.7.2} is called a morphism of ringed topoi
which is left (\resp right) flat.
\end{enonce*}

\begin{enonce*}{Proposition 1.9}\etiquette[1.9]{V.prop1.9}
Let $𝒰\subset 𝒱$ be two universes, let $C$ be a $𝒰$-site, let
$C^{\sim}_{𝒰}$ and $C^{\sim}_{𝒱}$ be the categories of $𝒰$- and
$𝒱$-sheaves of sets, respectively, let
$\epsilon:C^{\sim}_{𝒰}\to C^{\sim}_{𝒱}$ be the canonical inclusion
functor, and let $A$ be a $𝒰$-sheaf of rings on $C$.
\begin{enumerate}
\renewcommand{\labelenumi}{\theenumi)}
\item The functor $\epsilon$ commutes with $𝒰$-small inductive and
projective limits of Modules. The functor $\epsilon$ is conservative
and fully faithful on the categories of Modules.

\item For every object $H$ of $C^{\sim}_{𝒰}$, there is a canonical
isomorphism $\epsilon(A_H)\simeq \epsilon(A)_{\epsilon(H)}$.

\item Every injective left or right $A$-Module is carried by
$\epsilon$ to an injective $\epsilon(A)$-Module.

\item Every flat right or left $A$-Module is carried by $\epsilon$
to a flat $\epsilon(A)$-Module.

\item The functor $\epsilon$ commutes with the formation of tensor
products and sheaves of homomorphisms.
\end{enumerate}
\end{enonce*}

The formation of the associated sheaf is independent of the
universe (II \hyperref[II.prop3.6]{3.6}). It then follows from the construction of
inductive and projective limits in categories of sheaves (II \hyperref[II.theo4.1]{4.1}
and \hyperref[II.prop6.4]{6.4}) that the functor $\epsilon$ commutes with $𝒰$-small
projective limits of sets or Modules, which gives 1). Let us prove
2). By taking a $𝒰$-small generating subcategory of
$C^{\sim}_{𝒰}$, we may reduce to the case where $C$ is
$𝒰$-small\pageoriginale (III \hyperref[III.theo4.1]{4.1}). It then follows from (II \hyperref[II.prop6.5]{6.5})
that the free $A$-Module generated by $H$ is obtained by forming the
presheaf of free $A$-modules generated by $H$, a construction that
commutes with enlargement of universes, and then forming the
associated sheaf, an operation which likewise commutes with
enlargement of universes (II \hyperref[II.prop3.6]{3.6}). To prove 3), it suffices, by
\hyperref[V.prop0.5]{0.5}, to show that every $𝒱$-subsheaf of a
$𝒰$-sheaf is a $𝒰$-sheaf, which is clear.

Let us prove 4). Let $M$ be a flat $A$-Module. It suffices to show
that, for every injective $\epsilon(A)$-Module $J$, the abelian sheaf
$\Hhom_{\epsilon(A)}(\epsilon(M),J)$ is injective
(\hyperref[V.prop1.2]{1.2}). By \hyperref[V.prop0.5]{0.5}, $J$ is a
subobject, hence a direct factor, of an object of the form
$\prod_{\alpha}\epsilon(I_{\alpha})$, where the $I_{\alpha}$ are
injective. We may therefore suppose that $J=\epsilon(I)$, where $I$
is an injective object. If the homomorphism
$$
\epsilon(\Hhom_A(M,I))\longrightarrow
\Hhom_{\epsilon(A)}(\epsilon(M),\epsilon(I))
$$
is an isomorphism, it follows from 3) that
$\Hhom_{\epsilon(A)}(\epsilon(M),\epsilon(I))$ is injective. It
therefore remains to prove 5). The fact that formation of the tensor
product commutes with the functor $\epsilon$ follows from IV \hyperref[IV.cor12.6]{12.6} and
the fact that formation of the associated sheaf commutes with
enlargement of the universe (II \hyperref[II.prop3.6]{3.6}). For every object $X$ of
$C^{\sim}_{𝒰}$ and every pair of $A$-Modules $M,N$ of
$C^{\sim}_{𝒰}$, we therefore have
$$
\Hom_{C^{\sim}_{𝒱}}(\epsilon(X),\epsilon(\Hhom_A(M,N)))\simeq
\Hom_{C^{\sim}_{𝒰}}(X,\Hhom_A(M,N))\simeq
\Hom_A(\ZZ_X\otimes_{\ZZ}M,N)
$$
by the full faithfulness of $\epsilon$, and then
$$
\Hom_A(\ZZ_X\otimes_{\ZZ}M,N)\simeq
\Hom_{\epsilon(A)}(\epsilon(\ZZ_X\otimes_{\ZZ}M),\epsilon(N))
\simeq
\Hom_{\epsilon(A)}(\epsilon(\ZZ_X)\otimes_{\ZZ}\epsilon(M),
\epsilon(N))
$$
by the full faithfulness of $\epsilon$ and what precedes. Using 2)
and IV \hyperref[IV.prop12.14]{12.14}, we finally obtain an isomorphism
$$
\Hom_{C^{\sim}_{𝒱}}(\epsilon(X),\epsilon(\Hhom_A(M,N)))\simeq
\Hom_{C^{\sim}_{𝒱}}(\epsilon(X),
\Hhom_{\epsilon(A)}(\epsilon(M),\epsilon(N)));
$$
hence the asserted isomorphism.

\setcounter{subsection}{9}
\subsection{}\etiquette[1.10]{V.subsec1.10}
Let $E$ be a topos and let $\Uu=(U_i\to X)_{i\in I}$ be a small
family of morphisms with the same target. For every finite set
$\Delta_n=\sisi{[0,\ldots,n]}{\{0,\ldots,n\}}$, put
$$
S_n(\Uu)=\coprod_{f:\Delta_n\to I}
U_{f(0)}\times_X\cdots\times_X U_{f(n)}{\quoi},
$$
the coproduct being taken over the set of maps from $\Delta_n$ to
$I$. Let $m,n$ be two integers, and let\pageoriginale
$g:\Delta_m\to\Delta_n$ be a map. We define a morphism
\begin{equation*}
s(g):S_n(\Uu)\to S_m(\Uu){\quoi}\sisi{,}{}\tag{1.10.1}
\end{equation*}\etiquette[1.10.1]{V.1.10.1}
as follows: for every $f:\Delta_n\to I$, the restriction of $s(g)$
to
\mbox{$U_{f(0)}\times_X\cdots\times_X U_{f(n)}$}
is the composite of the canonical inclusion
$$
U_{f(g(0))}\times_X\cdots\times_X U_{f(g(m))}
\hookrightarrow S_m(\Uu){\quoi},
$$
and the unique morphism
$$
s_f(g):U_{f(0)}\times_X\cdots\times_X U_{f(n)}\to
U_{f(g(0))}\times_X\cdots\times_X U_{f(g(m))}
$$
such that, for every $i\in\Delta_m$,
$$
\pr_{f(g(i))}s_f(g)=\pr_{f(g(i))}
$$
($\pr_a$ denotes the $a$-th projection). We thus obtain a
contravariant functor $\Delta_n\mapsto S_n$ from the category of
finite sets to $E$; in other words, a semi-simplicial complex
$S_∙(\Uu)$ of objects of $E$. Note that this complex is canonically
augmented toward $X$. Every functor from $E$ to a category $C$
carries $S_∙(\Uu)$ to a simplicial object of $C$. In particular, if
$A$ is a Ring of $E$, the ``free $A$-Module generated by'' functor
carries $S_∙(\Uu)$ to a simplicial complex of $A$-biModules,
augmented toward $A_X$ and denoted by $A_∙(\Uu)$. We have
\begin{equation*}
A_n(\Uu)=\bigoplus_{f:\Delta_n\to I}
A_{U_{f(0)}\times_X\cdots\times_X U_{f(n)}}{\quoi}.\tag{1.10.2}
\end{equation*}\etiquette[1.10.2]{V.1.10.2}
This complex will often be written
\[
\xymatrix@C=1cm@M=.5cm{
  \cdots  \ar@<12pt>[r]^(.4){s_0} \ar[r]^(.4){s_1} \ar@<-12pt>[r]^(.4){s_2}&
\displaystyle{\mathop{\bigoplus}_{i,j}} \quad A_{U_i ×_X U_j}
 \ar@<6pt>[r]^(.55){s_0}\ar@<-6pt>[r]^(.55){s_1} &
\displaystyle{\mathop{\bigoplus}_{i}}\quad A_{U_i} \ar@<2pt>[r]& A_X \\
}\tag{1.10.3} %% problem (numbering in the wrong place) \XXX
\]\etiquette[1.10.3]{V.1.10.3}
where the arrows in \eqhyperref[V.1.10.3]{1.10.3} (except the last,
which is the augmentation) represent the face operators of the
simplicial complex $A_∙(\Uu)$, that is, the morphisms corresponding
to the increasing injective maps from $\Delta_n$ into
$\Delta_{n+1}$ ($s_i$ omits the integer $i$). To the complex
$A_∙(\Uu)$ we associate a differential complex augmented toward
$A_X$:
\begin{equation}
\cdots\xrightarrow{d}\bigoplus_{i,j}A_{U_i\times_XU_j}
\xrightarrow{d}\bigoplus_iA_{U_i}\to A_X{\quoi},
\sisi{\tag*{1.10.4}}{\tag{1.10.4}}
\end{equation}\etiquette[1.10.4]{V.1.10.4}
by\pageoriginale putting
\begin{equation}
d=\sum(-1)^is_i.\sisi{\tag*}{\tag}{1.10.5}
\end{equation}\etiquette[1.10.5]{V.1.10.5}

\begin{enonce*}{Proposition 1.11}\etiquette[1.11]{V.prop1.11}
When the family $\Uu=(U_i\to X)_{i\in I}$ is epimorphic, the
differential complex \eqhyperref[V.1.10.4]{1.10.4} is acyclic and
gives a resolution of $A_X$.
\end{enonce*}

Denote by $\ZZ$ the constant sheaf with value $\ZZ$. By the
definition of the ``free $A$-Module generated by'' functor, for every
object $H$ of $E$ we have
$$
A_H\simeq {\ZZ}_H\otimes_{\ZZ}A{\quoi},
$$
hence
$$
A_∙(\Uu)\simeq \ZZ_∙(\Uu)\otimes_{\ZZ}A{\quoi}.
$$
Since the components of $\ZZ_∙(\Uu)$ are flat $\ZZ$-Modules
(\hyperref[V.prop1.3.1]{1.3.1}), it suffices to prove the proposition
when $A=\ZZ$.

Suppose first that $E$ is the topos of sets. The augmented complex
$S_∙(\Uu)$ is then a direct sum of augmented complexes of the form
\[
\xymatrix@C=.6cm@M=.18cm@R=.3cm{
\cdots \ar@<9pt>[r] \ar@<3pt>[r] \ar@<-3pt>[r] \ar@<-9pt>[r]&
S\times S\times S \ar@<5pt>[r]\ar[r]\ar@<-5pt>[r]&S\times S
\ar@<2pt>[r]\ar@<-2pt>[r]&S\ar[r]&\{e\}&
\txt{($\{e\}$ a one-element set).}
}
\]
Each of these augmented complexes is homotopically trivial. Thus
$S_∙(\Uu)$ is a homotopically trivial augmented complex, and its
homology is therefore trivial, proving the proposition in this case.

Now let $p:\text{Ens}\to E$ be a point of $E$. Since formation of
the complex $\ZZ_∙(\Uu)$ commutes with the inverse-image functors of
morphisms of topoi,
$p^{\ast}(\ZZ_∙(\Uu))\simeq\ZZ_∙(p^{\ast}(\Uu))$ is a resolution of
${\ZZ}_{p^{\ast}X}\simeq p^{\ast}({\ZZ}_X)$; this proves the
proposition when $E$ has enough fiber functors (IV \hyperref[IV.subsec4.6]{4.6}). This holds,
in particular, when $E$ is a topos of presheaves $C^{\wedge}$,
since for every object $X$ of $C$, $\Gamma(X,-)$ is a fiber functor.
In the general case, $E$ is equivalent to a topos of sheaves on a
small site $C$ (IV 1), and the epimorphic family
$\Uu=(U_i\to X)_{i\in I}$ is the image, under the ``associated
sheaf'' functor, of an epimorphic family
$\Uu'=(U'_i\to X')_{i\in I}$. Consequently,\pageoriginale
${\ZZ}_∙(\Uu)\simeq\underline{a}{\ZZ}_∙(\Uu')$ is a resolution of
$\underline{a}{\ZZ}_{X'}\simeq{\ZZ}_X$.


% END INCLUDED SGA4_English_current_buildable_source_20260906/source/05/drafts/SGA4_V_LINES_603_925_English.texfrag


% BEGIN INCLUDED SGA4_English_current_buildable_source_20260906/source/05/drafts/SGA4_V_LINES_926_1361_English.texfrag
\section{Čech cohomology. Cohomological notation}\etiquette[2]{V.sec2}

\subsection{General notation}\etiquette[2.1]{V.subsec2.1}

\subsubsection{}\etiquette[2.1.1]{V.subsubsec2.1.1}
Let $(E, A)$ be a ringed topos, and let $M$, $N$ be two $A$-Modules
(left modules, to fix ideas). We denote by $\Ext^q_A (E; M, N)$ the
value at $N$ of the $q$-th right derived functor of the functor
$\Hom_A(M,⋅)$ \cite{sga4.exp.V.bib.11}. The functors $\Ext_A^q(E; M, N)$, $q \in
\NN$, form a $\delta$-functor in the variable $N$. They also form a
contravariant $\delta$-functor in the variable $M$. By definition,
$\Ext_A^0(E; M, N) = \Hom_A(M, N)$.

\subsubsection{}\etiquette[2.1.2]{V.subsubsec2.1.2}

Let $X$ be an object of $E$. When $M = A_X$, the free $A$-Module
generated by $X$ (IV 12), we put $\Ext_A^q(E; A_X, N) = H^q(X,
N)$. Notice that the ring $A$ no longer appears in this new
notation. This cannot cause confusion, since we will show
(\hyperref[V.cor3.5]{3.5}) that the formation of $H^q(X,⋅)$
\sisi{commutes}{commutes} with restriction of scalars. The functor
$H^q(X,⋅)$ is the $q$-th right derived functor of
$\Hom_A(A_X,⋅) = \Hom_E(X,⋅)$, also denoted by $\Gamma(X,⋅)$. When
$M = A$, we put $\Ext^q_A(E; A, N) = H^q(E, N)$.

\subsection{Localization}\etiquette[2.2]{V.subsec2.2}

Let $X$ be an object of $E$, and let $j : E_{/X} \to E$ be the
localization morphism (IV 8). The functor $j^\ast_X$ for Modules
\sisi{,}{} is exact (\hyperref[V.prop4.11]{4.11}) and admits an exact
left adjoint $j_{X!}$ (\hyperref[V.prop4.11]{4.11}). Consequently,
it sends injective Modules to injective Modules. Thus, for a
variable $A$-Module $N$ of $E$ and a variable $A | X$-Module $M$ of
$E_{/X}$, there are functorial isomorphisms
\begin{equation*}
\Ext^q_{A | X}(E_{/X} ; M, j^\ast_X N) \simeq \Ext^q_A(E; j_{X!}M,
N){\quoi}.\tag{2.2.1}
\end{equation*}\etiquette[2.2.1]{V.2.2.1}
In particular, there are canonical isomorphisms
$$
H^q(E_{/X}, j^\ast_X N) \simeq H^q (X, N){\quoi}.
$$

For every object $X$ of $E$ and every pair $M$, $N$ of $A$-Modules,
we put
\begin{equation*}
\Ext^q_A(X; M, N) \simeq \Ext^q_{A|X}(E_{/X}; M | X, N|X){\quoi}.\tag{2.2.2}
\end{equation*}\etiquette[2.2.2]{V.2.2.2}
By\pageoriginale what precedes, the functors $\Ext^q_A(X; M, ⋅)$
are the derived functors of the functors $N \mapsto
\Hom_{A|X}(M|X, N|X)$. \sisi{The functors}{The functors}
$(M, N) \mapsto \Ext^q_A(X; M, N)$, $q > 0$, form a
$\delta$-functor with respect to each variable.

\subsection{The case of presheaf topoi. Cohomology of a
  covering}\etiquette[2.3]{V.subsec2.3}

\subsubsection{}\etiquette[2.3.1]{V.subsubsec2.3.1}
Let $C$ be a small category endowed with a presheaf of rings $A$,
let $C^{\wedge}$ be the topos of presheaves on $C$, and let $X$ be
a representable object of $C^{\wedge}$. The functor which associates
to an $A$-Module $M$ the group $\Gamma(X, M) = M(X)$ is exact (I 3).
Consequently, $H^q(X, M)=0$ for every $q > 0$ and every $A$-Module
$M$, or, equivalently, $A_X$ is a projective $A$-Module.

\subsubsection{}\etiquette[2.3.2]{V.subsubsec2.3.2}
Let $S$ be a presheaf on $C$. For every $A$-Module $M$, there is a
canonical isomorphism (I 2):
$$
\Gamma(S, M) \simeq \varprojlim_{C/S} M|S
$$
Moreover, for every injective $A$-Module $M$, the $A | S$-Module
$M | S$ is injective \hyperref[V.subsec2.2]{(2.2)}. Consequently, the groups $H^{q}(S, M)$
are the values at $M | S$ of the right derived functors of
$\varprojlim_{C/S}$. Denoting these derived functors by
$\varprojlim_{C/S}^{q}$, there are canonical isomorphisms
$$
H^q (S, M) \simeq {\varprojlim_{C/S}}^{q} M | S{\quoi}.
$$
In particular, there are canonical isomorphisms
$$
H^q (C^{\wedge}, M) \simeq {\varprojlim_{C}}^{q} M{\quoi}.
$$

\subsubsection{}\etiquette[2.3.3]{V.subsubsec2.3.3}
Let $X$ be an object of $C$, and let $\Uu = (U_i \to X)$, $i \in
I$, be a family of morphisms of $C$ such that, for every $i \in I$,
$U_i \to X$ is pullbackable (I 10). Denote by $A_∙$ the simplicial
complex studied in \hyperref[V.subsec1.10]{1.10}:
\[
\xymatrix@M=.25cm@C=1cm{
A_∙ = \cdots \ar@<7pt>[r] \ar@<1pt>[r]  \ar@<-5pt>[r] &
\displaystyle{\mathop{\coprod}_{(i,j)\in I \times I}}\quad
A_{U_i \times_X U_j} \ar@<3pt>[r]\ar@<-3pt>[r]&
\displaystyle{\mathop{\coprod}_{i\in I}} \quad A_{U_i}.
}
\]
% Why ∐ rather than ⨁ as above?
For every $A$-Module $M$, put $C^{∙}(\Uu, M) = \Hom_A (A_∙, M)$:
\[\vcenter{
\xymatrix@C=.75cm@M=.25cm{
C^{∙}(\mathcal{U},M): \displaystyle{\mathop{\prod}_{i\in I}}
M(U_i) \ar@<6pt>[r]\ar@<-2pt>[r]& \displaystyle{\mathop{\prod}_{(i,j)\in
    I\times I}} M(U_i ×_X U_j) \ar@<10pt>[r] \ar@<3pt>[r]
\ar@<-5pt>[r]& \cdots .
}}\leqno{(2.3.3.1)}
\]\etiquette[2.3.3.1]{V.2.3.3.1}
Put $H^q(\Uu, M) = H^q(C^{∙} (\Uu, M))$.


\begin{enonce*}{Proposition 2.3.4}\etiquette[2.3.4]{V.prop2.3.4}
\mbox{ }
\begin{enumerate}
\renewcommand{\labelenumi}{\theenumi)}
\item With\pageoriginale the notation of
  \hyperref[V.subsubsec2.3.3]{2.3.3}, let $R \hookrightarrow X$ be
  the sieve generated by the family $(U_i \to X)$, $i \in I$. There
  is a canonical isomorphism
$$
H^{q}(\Uu, M) \simeq H^q(R, M){\quoi}.
$$

\item The functors $H^q(\Uu,⋅)$ commute with restrictions of
  scalars.
\end{enumerate}
\end{enonce*}

Since $R$ is a subobject of $X$ in $C^{\wedge}$, the fiber products
$U_{i_1} \times_R U_{i_2} \sisi{\times}{\times_R} \cdots
\times_{R} U_{i_n}$ and $U_{i_1} \times_X U_{i_2}
\sisi{\times}{\times_X} \cdots \times_{X} U_{i_n}$ are canonically
isomorphic. It follows from \hyperref[V.prop1.11]{1.11} that the
complex $A_{𝒰 ∙}$ is a resolution of $A_R$, and from
\hyperref[V.subsubsec2.3.1]{2.3.1} that the components of $A_{𝒰 ∙}$
are projective Modules. The \sisi{cohomology}{cohomology} groups of
$C^{∙}(\mathcal{U}, M)$ are therefore canonically isomorphic to
$\Ext^{q}_A(C^{\wedge}; A_R, M)$, whence the isomorphism. Assertion
2) follows immediately from the description of
$C^{∙}(\mathcal{U}, M)$ \eqhyperref[V.2.3.3.1]{2.3.3.1}.

\begin{enonce*}{Corollary 2.3.5}\etiquette[2.3.5]{V.cor2.3.5}
Let $\mathcal{U} = (U_i \to X, i \in I)$ and
$\mathcal{U}' = (U'_j \to X, j \in J)$ be two families of morphisms
with target $X$; let $\phi = (\phi : I \to J, m_i : U_i \to
U'_{\phi(i)})$ and $\phi' = (\phi' : I \to J, m'_i : U_i \to
U'_{\phi(i)})$ be two morphisms (over $X$) from $\mathcal{U}$ to
$\mathcal{U}'$. The morphisms $\phi$ and $\phi'$ induce the
morphisms $\phi^{q}$ and ${\phi'}^{q} :
\sisi{{\phi'}^{q} :}{} H^q(\mathcal{U}, M) \to H^q(\mathcal{U}',
M)$. The morphisms $\phi^q$ and ${\phi'}^{q}$ are equal. In
particular, if the families $\mathcal{U}$ and $\mathcal{U}'$ are
equivalent (\ie if there is a morphism from $\mathcal{U}$ to
$\mathcal{U}'$ and a morphism from $\mathcal{U}'$ to $\mathcal{U}$),
the $A$-Modules $H^{q} (\mathcal{U}, M)$ and $H^q(\mathcal{U}', M)$
are canonically isomorphic.
\end{enonce*}

\begin{enonce*}[remark]{Proof}
Indeed, let $R$ and $R'$ be the sieves of $X$ generated by the
families $\mathcal{U}$ and $\mathcal{U}'$. The morphisms $\phi$ and
$\phi'$ define the same morphism from $R$ to $R'$, and therefore
induce two homotopic morphisms between the projective resolutions
$A_{\mathcal{U}∙}$ and $A_{\mathcal{U}'∙}$.
\end{enonce*}

\begin{enonce*}[remark]{Exercise 2.3.6}\etiquette[2.3.6]{V.exer2.3.6}
(Standard resolution)
\begin{enumerate}
\renewcommand{\theenumi}{\alph{enumi}}
\renewcommand{\labelenumi}{\theenumi)}
\item Let $C$ be a small category. For every integer $n > 0$, denote
  by $\mathrm{Fl}^{n}(C)$ the set of sequences of morphisms of $C :
  (u_1, \ldots, u_n)$ such that, for every $i$, $0 < i < n$, the
  target of $u_i$ equals the source of $u_{i + 1}$, so that the
  morphisms $u_i$ and $u_{i +1}$ are\pageoriginale composable. Define
  a semisimplicial set
\[
\xymatrix@M=.25cm@C=.75cm{
ES(C) = \ob \;C  & \mathrm{Fl}^{1}(C) \ar@<3pt>[l]\ar@<-3pt>[l]
 & \mathrm{Fl}^{2}(C) \ar@<6pt>[l] \ar[l] \ar@<-6pt>[l] & \mathrm{Fl}^{3}(C)\ar@<9pt>[l]
\ar@<3pt>[l] \ar@<-3pt>[l] \ar@<-9pt>[l] \quad \cdots ,
}
\]
whose face operators $s_i : \mathrm{Fl}^n(C) \to
\mathrm{Fl}^{n-1}(C)$ are as follows:
\begin{align*}
s_{0}(u_1, \ldots, u_n) &= (u_2, \ldots, u_n),\\
s_{i}(u_1, \ldots, u_n) &= (u_1, \ldots, u_{i+1} u_i, \ldots, u_n), \hskip1em 0
< i < n ,\\
s_{n}(u_1, \ldots, u_n) &= (u_1, \ldots, u_{n-1}).
\end{align*}

\item Show that when $C$ has an initial object or a final object,
  the complex $ES(C)$ is homotopically trivial.

\item For every object $X$ of $C$, denote by ${}_{X\backslash}C$ the
  category of arrows with source $X$. Define a semisimplicial
  presheaf $ES(C)$ whose value at every object $X$ of $C$ is
  $ES({}_{X\backslash}C)$. Denote by $\ZZ_{ES(C)}$ the free abelian
  semisimplicial presheaf generated by $ES(C)$. It has a canonical
  augmentation $\ZZ_{ES(C)}\to \ZZ$ to the constant presheaf with
  value $\ZZ$. Show that, upon passing to the associated differential
  complex, the complex $\ZZ_{ES(C)}$ is a resolution of $\ZZ$ and
  that its components are projective abelian presheaves.

\item For every abelian presheaf $M$, put $ST^{∙}(M) =
  \Hom(\ZZ_{ES(C)}, M)$. Describe the components of $ST^{∙}(M)$
  explicitly. Show that, for every integer $q \geq 0$,
  $H^q(ST^{∙}(M)) = H^q(C^{\wedge}, M)$.

\item Show that, for every sheaf $S$ on $C$, the functors $H^q(S,⋅)$
  commute with restrictions of scalars.

\item Observe that when $C$ is a group, $\ZZ_{ES(C)}$ is the
  standard resolution of the trivial module $\ZZ$, \cite{sga4.exp.V.bib.3}.

\item Show that, for every abelian presheaf $M$ on $C$,
  $\ZZ_{ES(C)} \otimes M$ is a (left) resolution of $M$. Define the
  functor $\varinjlim_{C}$ on presheaves. Show that it is right exact.
  Denote its left derived functors by $H_q(C, M)$. Show that the
  components of the complex $\ZZ_{ES(C)} \otimes M$ are acyclic for
  $\varinjlim_{C}$. Deduce, denoting by $ST_{∙}(M)$ the complex
  $\varinjlim_{C}(\ZZ_{ES(C)} \otimes M)$, isomorphisms
  $H_q(ST_{∙}(M)) \simeq H_q(C, M)$.

\item Let\pageoriginale $u : C^{\wedge} \to \Ens$ be the unique
  morphism from the topos $C^{\wedge}$ to the punctual topos. Denote
  by $u_{!} : C^{\wedge}_{\ZZ} \to \Ab$ the left adjoint of the
  functor $u^\ast : \Ab \to C^{\wedge}_{\ZZ}$. Show that for every
  abelian presheaf $M$, $u_{!}M =\varinjlim_{C} M$.

\item Henceforth denote by $ST^{∙}(C, M)$ and $ST_{∙}(C, M)$ the
  complexes denoted above by $ST^{∙}(M)$ and $ST_{∙}(M)$. A
  presheaf $M$ on $C$ is called locally constant if it sends all
  morphisms of $C$ to isomorphisms. Associate to every locally
  constant presheaf $M$ a locally constant presheaf $\check{M}$ on
  the category $C^{\circ}$ (the category opposite to $C$) such that,
  for every morphism $u$ of $C$, $\check{M}(u) = M(u)^{-1}$. Find a
  canonical isomorphism between the complexes $ST^{∙}(C, M)$ and
  $ST^{∙}(C^{\circ}, \check{M})$ (\resp $ST_{∙}(C, M)$ and
  $ST_{∙}(C^{\circ}, \check{M})$).

\item Let $C$ be a small category having an initial object (\resp a
  small filtered category). Show that, for every constant presheaf
  $M$, $H^{q} (C^{\wedge}, M) = 0$ for $q > 0$ (\resp $H_q(C, M) =
  0$ for $q > 0$).
\end{enumerate}
\end{enonce*}

\subsection{The case of small sites. Čech cohomology}\etiquette[2.4]{V.subsec2.4}

\subsubsection{}\etiquette[2.4.1]{V.subsubsec2.4.1}
Let $(C, A)$ be a ringed $𝒰$-site, let $C^{\sim}$ be the topos of
sheaves on $C$, and let $\epsilon: C \to C^{\sim}$ be the canonical
functor which associates to an object of $C$ the sheaf associated to
the presheaf represented by that object. By abuse of notation, for
every object $X$ of $C$ and every sheaf of $A$-modules $M$, we put
$H^q(X, M) = H^q(\epsilon(X), M)$ (\cf \hyperref[V.subsubsec2.3.1]{2.3.1}). Recall that when the
topology of $C$ is coarser than the canonical topology, as is always
the case in practice, the functor $\epsilon$ is fully faithful and
allows us to identify $C$ with its image under $\epsilon$.

\subsubsection{}\etiquette[2.4.2]{V.subsubsec2.4.2}
Denote by $\Hh^0 : C^{\sim{\phantom{a}}}_A \to C^{\wedge}_A$ the
inclusion functor from sheaves of $A$-modules to the category of
presheaves of $A$-modules. Thus, for every sheaf of $A$-modules $M$,
one has by definition
\begin{equation*}
\Hh^{0}(M)(X) = H^{0}(X, M) = M(X){\quoi},\tag{2.4.2.1}
\end{equation*}\etiquette[2.4.2.1]{V.2.4.2.1}
for every object $X$ of $C$. The functor $\Hh^{0}$ is left exact.
Its right derived functors are denoted by $\Hh^{q}$. Since, for
every object $X$ of $C$, the ``sections over $X$'' functor is exact
in the category of presheaves, one has
\begin{equation*}
\mathcal{H}^q(M)(X) = H^q(X, M){\quoi}, \tag{2.4.2.2}
\end{equation*}\etiquette[2.4.2.2]{V.2.4.2.2}
for\pageoriginale every object $X$ of $C$ and every sheaf of
$A$-modules $M$, so that the presheaf $\mathcal{H}^q(M)$ is simply
the presheaf $X \mapsto H^q(X, M)$.

\subsubsection{}\etiquette[2.4.3]{V.subsubsec2.4.3}
Suppose that $(C, A)$ is a ringed \emph{small site}, so that
$C^{\wedge}$ is a \emph{topos} to which the results of
\hyperref[V.subsec2.3]{2.3} may be applied. Let $X$ be an object of
$C$, and let $R \hookrightarrow X$ be a covering sieve. For every
\emph{presheaf} of $A$-modules $G$, the groups $H^{q}(R, G)$ (which
are thus computed in the topos $C^{\wedge}$) are called the
\emph{Čech cohomology groups of the presheaf $G$ relative to the
  covering sieve $R$.} When $R \hookrightarrow X$ is the sieve
generated by a covering family $\mathcal{U} = (U_i \to X)_{i \in I}$
of pullbackable morphisms, these groups can be computed by means of
the complex $C^{∙}(\mathcal{U}, G)$
(\hyperref[V.subsubsec2.3.3]{2.3.3}), called the \emph{Čech complex
  of $G$ relative to the covering family $\mathcal{U}$} (or relative
to the covering $\Uu$). The groups $H^{q}(\Uu,G)=H^{q}(R,G)$
(\hyperref[V.prop2.3.4]{2.3.4}) are then called the \emph{Čech
cohomology groups of $G$ relative to the covering family $\Uu$}.

\subsubsection{}\etiquette[2.4.4]{V.subsubsec2.4.4}
Let $M$ be a sheaf of $A$-modules on $C$. The groups
$H^q(\mathcal{U}, \Hh^{0}(M))$ are most often denoted, by abuse, by
$H^q(\mathcal{U}, M)$ and called the \sisi{Čech}{Čech} cohomology groups of the
sheaf $M$ relative to the covering family $\mathcal{U}$.

\subsubsection{}\etiquette[2.4.5]{V.subsubsec2.4.5}
Denote by $\check{\Hh}^0 : C^{\wedge}_A \to C^{\wedge}_A$ the
natural extension to presheaves of $A$-modules of the functor $L$
described in II. Thus, by definition, for a presheaf $G$ and an
object $X$ of $C$,
\begin{equation}
\check{\Hh}^0 (G) (X) = \varinjlim_{R \hookrightarrow X}
G(R){\quoi},\tag{2.4.5.1}
\end{equation}\etiquette[2.4.5.1]{V.2.4.5.1}
where the inductive limit is taken over the covering sieves of $X$.
It follows from \eqhyperref[V.2.4.5.1]{2.4.5.1} that the functor
$\check{\Hh}^{0}$ is left exact. The right derived functors of
$\check{\Hh}^{0}$ are denoted by $\check{\Hh}^{q}$. Since the
``sections over $X$'' and ``filtered inductive limit'' functors are
exact, it follows from \eqhyperref[V.2.4.5.1]{2.4.5.1} that
\begin{equation}
\check{\Hh}^{q} (G) (X) = \varinjlim_{R \hookrightarrow
  X} H^q(R, G){\quoi},\tag{2.4.5.2}
\end{equation}\etiquette[2.4.5.2]{V.2.4.5.2}
where the limit is taken over the covering sieves of $X$.

The\pageoriginale presheaves $\check{\Hh}^{q}$ are called the
\emph{Čech cohomology presheaves}. For every object $X$ of $C$, put
\begin{equation*}
{\check{H}}^q (X, G) = \check{\Hh}^{q} (G) (X){\quoi}.\tag{2.4.5.3}
\end{equation*}\etiquette[2.4.5.3]{V.2.4.5.3}
The groups ${\check{H}}^q(X, G)$ are called the \emph{Čech
cohomology groups}. When the topology of $C$ is defined by a
pretopology, as is most often the case in practice, it follows from
\hyperref[V.prop2.3.4]{2.3.4} that
\begin{equation*}
{\check{H}}^q (X, G) \simeq \varinjlim_{\mathcal{U}} {H}^{q}
(\mathcal{U}, G){\quoi},\tag{2.4.5.4}
\end{equation*}\etiquette[2.4.5.4]{V.2.4.5.4}
where the inductive limit is taken over the pullbackable covering
families preordered by the natural order relation on their
corresponding sieves (\hyperref[V.cor2.3.5]{2.3.5}).

\subsubsection{}\etiquette[2.4.6]{V.subsubsec2.4.6}
Let $M$ be a sheaf of $A$-modules. By abuse of notation, put
\begin{equation}
{\check{H}}^q (X, M) = {\check{H}}^q (X, \check{\Hh}^{0}(M)),
\qquad \check{\Hh}^{q} (M) = \check{\Hh}^{q}(\check{\Hh}^{0} (M)){\quoi},\tag{2.4.6.1}
\end{equation}\etiquette[2.4.6.1]{V.2.4.6.1}
and the groups ${\check{H}}^q(X, M)$ are called the \sisi{Čech}{Čech} cohomology
groups of the sheaf $M$. Note that, although the functors
${\check{H}}^q$ are derived functors on the category of presheaves,
they do not in general form a $\delta$-functor on the category of
sheaves.

\subsection{Change of universe. Čech cohomology for $𝒰$-sites}\etiquette[2.5]{V.subsec2.5}

\subsubsection{}\etiquette[2.5.1]{V.subsubsec2.5.1}
Let $(C, A)$ be a ringed $𝒰$-site, and let $𝒱$ be a universe
containing $𝒰$. The site $(C, A)$ is then a $𝒱$-site, and there are
a $𝒰$-topos $C^{\sim}_{𝒰}$, a $𝒱$-topos $C^{\sim}_{𝒱}$, and a
canonical inclusion functor $\epsilon : C^{\sim}_{𝒰}
\hookrightarrow C^{\sim}_{𝒱}$. The functor $\epsilon$ is exact and
fully faithful on the categories of Modules, and sends injective
Modules to injective Modules (\hyperref[V.prop1.9]{1.9}). Thus, for
every pair of $𝒰$-sheaves of $A$-modules, there are canonical
isomorphisms
\begin{equation*}
\Ext^q_A(C^{\sim}_{𝒰}; M, N) \simeq
\Ext^q_{\epsilon A}(C^{\sim}_{𝒱}; \epsilon M,
\epsilon N),\quad q \geq 0{\quoi}.\tag{2.5.1.1}
\end{equation*}\etiquette[2.5.1.1]{V.2.5.1.1}
In particular, for every $𝒰$-sheaf of sets $R$ on $C$, there are
canonical isomorphisms (\hyperref[V.subsubsec2.1.1]{2.1.1})
\begin{equation*}
H^q(R, M) \simeq H^q(\epsilon R, \epsilon M),\quad q \geq
0{\quoi},\tag{2.5.1.2}
\end{equation*}\etiquette[2.5.1.2]{V.2.5.1.2}
and, even more particularly, for\pageoriginale every object $X$ of
$C$, there are canonical isomorphisms
(\hyperref[V.subsubsec2.4.1]{2.4.1})
\begin{equation*}
H^q(X, M) \simeq H^q(X, \epsilon M){\quoi}.\tag{2.5.1.3}
\end{equation*}\etiquette[2.5.1.3]{V.2.5.1.3}
Loosely speaking, we may therefore say that sheaf cohomology does
not depend on the choice of universes, and, when required for a
proof \sisi{or}{or} a construction, one may always enlarge the
universe in order to compute the cohomology of a sheaf.

\subsubsection{}\etiquette[2.5.2]{V.subsubsec2.5.2}
Let $(C, A)$ be a ringed $𝒰$-site, and let $𝒱$ be a universe
containing $𝒰$. Denote by $C^{\sim}_{A}$ the category of
$𝒰$-sheaves of modules, and let
$\widehat{\epsilon} : C^{\wedge}_{A, 𝒰} \to
C^{\wedge}_{A, 𝒱}$ be the inclusion functor from $𝒰$-presheaves to
$𝒱$-presheaves of $A$-modules. The functor $\widehat{\epsilon}$ is
exact, and therefore the derived functors of the functor
$\widehat{\epsilon} \Hh^{0} : C^{\sim}_A \to C^{\wedge}_{A,
𝒱} (\phantom{ABCD})$ are the functors $\widehat{\epsilon} \Hh^q$,
$q \geq 0$. By abuse of notation, we will again denote by
$\Hh^q :C^{\sim}_A \to C^{\wedge}_{A, 𝒱}$, $q \geq 0$, the
functors $\widehat{\epsilon} \Hh^q$. When $C$ is $𝒱$-small, this
enlargement of the universe has the following advantage: the
category $C^{\wedge}_{𝒰}$ is not in general a $𝒰$-topos, and
$𝒰$-presheaves of $A$-modules are not necessarily submodules of
injective $𝒰$-presheaves, whereas the category of $𝒱$-presheaves is
a topos (a $𝒱$-topos), and consequently every $𝒱$-presheaf of
$A$-modules is a subobject of an injective $𝒱$-presheaf
(\hyperref[V.subsubsec0.1.1]{0.1.1}). Thus, for every $𝒱$-presheaf
of sets $R$ (and in particular when $R$ is a $𝒰$-presheaf) and every
$𝒰$-sheaf of $A$-modules $M$, the groups $H^q(R, \Hh^q (M))$ are
defined by \hyperref[V.subsubsec2.1.1]{2.1.1}, and it follows from
\eqhyperref[V.2.5.1.2]{2.5.1.2} that these groups do not depend on
the universe $𝒱$ under consideration. Likewise, for every pair of
integers $\geq 0$, denoted by $p$ and $q$, the presheaves
$\check{\Hh}^p (\Hh^q (M))$ are defined by
\hyperref[V.subsubsec2.4.5]{2.4.5}, and it follows from
\eqhyperref[V.2.5.1.2]{2.5.1.2} and
\eqhyperref[V.2.4.5.4]{2.4.5.4} that these presheaves do not depend
on the universe $𝒱$ used to define them.

% END INCLUDED SGA4_English_current_buildable_source_20260906/source/05/drafts/SGA4_V_LINES_926_1361_English.texfrag


% BEGIN INCLUDED SGA4_English_current_buildable_source_20260906/source/05/drafts/SGA4_V_LINES_1362_1552_English.texfrag
\section{The Cartan--Leray spectral sequence associated with a
covering}\etiquette[3]{V.sec3}\pageoriginale

\begin{enonce*}{Proposition 3.1}\etiquette[3.1]{V.prop3.1}
Let $(C,A)$ be a ringed $𝒰$-site, and let $𝒱$ be a universe
containing $𝒰$. The functor
$\Hh^{0}:C^{\sim}_A\to C^{\wedge}_{A,𝒱}$
(\hyperref[V.subsubsec2.5.2]{2.5.2}) sends injective $A$-modules to
injective presheaves. For every integer $q>0$ and every $A$-module
$M$, the sheaf associated with the presheaf $\Hh^q(M)$ is zero.
\end{enonce*}

Denote by $\underline{a}_{𝒱}$ the ``associated sheaf'' functor for
$𝒱$-presheaves, and by
$\epsilon:C^{\sim}_A\hookrightarrow C^{\sim}_{A,𝒱}$ the inclusion
functor from $𝒰$-sheaves to $𝒱$-sheaves. One has
$\underline{a}_{𝒱}\Hh^{0}=\epsilon$ (II \hyperref[II.prop3.6]{3.6}), and, since the functors
$\underline{a}_{𝒱}$ and $\epsilon$ are exact, one has
$\underline{a}_{𝒱}\Hh^q=0$ for every integer $q>0$. For every
$𝒰$-sheaf $M$ and every $𝒱$-presheaf $N$, there is a functorial
isomorphism
$\Hom_{C^{\wedge}_{A,𝒱}}(N,\Hh^{0}(M))\simeq
\Hom_{C^{\sim}_{A,𝒱}}(\underline{a}_{𝒱}N,\epsilon M)$. When $M$ is
injective, $\epsilon M$ is injective (\hyperref[V.prop1.9]{1.9}), and,
since the functor $\underline{a}_{𝒱}$ is exact, the functor
$N\longmapsto\Hom_{C^{\wedge}_{A,𝒱}}(N,\Hh^{0}(M))$ is exact.
Consequently, $\Hh^{0}(M)$ is injective.

\begin{enonce*}{Theorem 3.2}\etiquette[3.2]{V.thm3.2}
Let $(C,A)$ be a ringed $𝒰$-site, let $R$ be a $𝒰$-presheaf of
sets on $C$, and let $M$ be a sheaf of $A$-modules. There is a
spectral sequence, functorial in $R$ and $M$:
\begin{equation*}
H^{p+q}(\underline{a}R,M)\Longleftarrow E^{p,q}_2=
H^p(R,\Hh^q(M)){\quoi}.
\tag{3.2.1}
\end{equation*}\etiquette[3.2.1]{V.3.2.1}
(When $C$ is not $𝒰$-small, the term $H^p(R,\Hh^q(M))$ must be
interpreted as the cohomology of the presheaf $\Hh^q(M)$ in the topos
$C^{\wedge}_{𝒱}$, where $𝒱$ is a universe containing $𝒰$ such
that $C$ is $𝒱$-small (\hyperref[V.subsubsec2.5.2]{2.5.2}).)
\end{enonce*}

By the definition of the ``associated sheaf'' functor (\cf II), there
is an isomorphism of functors
$H^{0}(\underline{a}R,M)\simeq H^0(R,\Hh^{0}(M))$. The functor
$\Hh^{0}$ sends injective objects to injective objects. The spectral
sequence of composite functors (\hyperref[V.prop0.3]{0.3}) is the
desired spectral sequence.

\begin{enonce*}{Corollary 3.3}\etiquette[3.3]{V.cor3.3}
Let $X$ be an object of $C$, and let
$\mathcal{U}=(U_i\to X)$, $i\in I$, be a covering family such that,
for every $i\in I$, the morphism $U_i\to X$ is pullbackable. Then
there is a spectral sequence (called the Cartan--Leray spectral
sequence):
\begin{equation}
H^{p+q}(X,M)\Longleftarrow E^{p,q}_2=
H^p(\mathcal{U},\Hh^q(M)){\quoi}.\tag{3.3.1}
\end{equation}\etiquette[3.3.1]{V.3.3.1}
\end{enonce*}

Let\pageoriginale $R\hookrightarrow X$ be the sieve generated by
$\mathcal{U}$. Since this sieve is covering, the sheaf associated with
$R$ is the sheaf associated with $X$ (II \hyperref[II.def5.2]{5.2}). Taking
\hyperref[V.subsubsec2.4.1]{2.4.1} into account, one has
$H^{p+q}(\underline{a}R,M)=H^{p+q}(X,M)$. The corollary then follows
from \hyperref[V.prop2.3.4]{2.3.4}.

\begin{enonce*}{Corollary 3.4}\etiquette[3.4]{V.cor3.4}
There is a spectral sequence, functorial in the sheaf $M$ and in the
object $X$ of $C$,
\begin{equation}
H^{p+q}(X,M)\Longleftarrow E^{p,q}_2=
\check{H}^p(X,\Hh^q(M)){\quoi}.\tag{3.4.1}
\end{equation}\etiquette[3.4.1]{V.3.4.1}
As $X$ varies in $C$, this spectral sequence gives the spectral
sequence of presheaves
\begin{equation}
\Hh^{p+q}(M)\Longleftarrow E^{p,q}_2=
\check{\Hh}^p(\Hh^q(M)){\quoi}.
\tag{3.4.2}
\end{equation}\etiquette[3.4.2]{V.3.4.2}
These spectral sequences give functorial morphisms (edge morphisms):
\hypertarget{sga4.exp.V.equation.3.4.3}{}\begin{align*}%\etiquette[3.4.3]{V.3.4.3}
\Phi^q(M)&:\check{\Hh}^q(M)\to\Hh^q(M){\quoi},\tag{3.4.3}\\
\Phi^q_X(M)&:\check{H}^q(X,M)\to H^q(X,M){\quoi}.\tag{3.4.4}
\end{align*}\hypertarget{sga4.exp.V.equation.3.4.4}{}
The morphisms $\Phi^q$ and $\Phi^q_X$ are isomorphisms when $q$ is
$0$ or $1$. The morphisms $\Phi^2$ and $\Phi^2_X$ are monomorphisms.
More generally, when the presheaves $\Hh^i(M)$ vanish for
$0<i<n$, the morphisms $\Phi^q(M)$ and $\Phi^q_X(M)$ are isomorphisms
for $0\leq q\leq n$ and monomorphisms for $q=n+1$.
\end{enonce*}

The first spectral sequence is obtained by taking the inductive limit
in the spectral sequence \eqhyperref[V.3.2.1]{3.2.1} over the covering
sieves $R\hookrightarrow X$ of $X$ \eqhyperref[V.2.4.5.2]{2.4.5.2}.
The second spectral sequence follows immediately from it
\eqhyperref[V.2.4.5.3]{2.4.5.3}. The sheaves associated with the
presheaves $\Hh^q(M)$ are zero when $q>0$
(\hyperref[V.prop3.1]{3.1}). Thus
$\check{\Hh}^{0}\check{\Hh}^{0}\check{\Hh}^{q}(M)=0$ for every
$q>0$ (\hyperref[II.theo3.4]{II.3.4}). Consequently,
$\check{\Hh}^{0}\check{\Hh}^{q}(M)=0$ for $q>0$ (II \hyperref[II.prop3.2]{3.2}). The
assertions concerning the morphisms $\Phi^q$ and $\Phi^q_X$ follow
immediately.

\begin{enonce*}{Corollary 3.5}\etiquette[3.5]{V.cor3.5}
Let $(E,A)$ be a ringed topos, and let $M$ be an $A$-module (a left
module, to fix ideas). Denote by $ℳ$ the underlying abelian sheaf of
$M$. The functor $M\longmapstoℳ$ is exact; consequently, for every
object $X$ of $E$, the canonical isomorphism
$$
H^{0}(X,M)\xrightarrow{\sim}H^{0}(X,ℳ)
$$
extends\pageoriginale to morphisms
$$
H^q(X,M)\to H^q(X,ℳ)\qquad q\geq0{\quoi}.
$$
These morphisms are isomorphisms.
\end{enonce*}

For every object $Y$ of $E$, one has \eqhyperref[V.2.4.5.4]{2.4.5.4}:
$$
H^p(Y,M)=\varinjlim_{\mathcal{U}}H^p(\mathcal{U},M)\quav
H^p(Y,ℳ)=\varinjlim_{\mathcal{U}}H^p(\mathcal{U},ℳ),
$$
where the inductive limit is taken over the epimorphic families
$\mathcal{U}=(Y_i\to Y)$, $i\in I$.

Consequently, the canonical homomorphism
$\check{H}^p(Y,M)\to\check{H}^p(Y,ℳ)$ is an isomorphism
(\hyperref[V.prop2.3.4]{2.3.4}). Hence, by
\eqhyperref[V.2.4.5.3]{2.4.5.3}, the canonical homomorphism
$\check{\Hh}^p(M)\to\check{\Hh}^p(ℳ)$ is an isomorphism. If $M$ is
an injective $A$-module, then $\check{\Hh}^p(ℳ)=0$ for $p>0$; hence
$\Hh^1(ℳ)=0$ and, by induction on $p$, $\Hh^p(ℳ)=0$ for $p>0$
(\hyperref[V.cor3.4]{3.4}). Thus $H^p(X,ℳ)=0$ for $p>0$
\eqhyperref[V.2.4.2.2]{2.4.2.2} and \hyperref[V.subsec2.5]{2.5}.
Consequently, the functor $M\longmapstoℳ$ sends injective objects to
objects acyclic for the functor $H^{0}(X,⋅)$, which gives the asserted
isomorphism.

\begin{enonce*}[remark]{Exercise 3.6}\etiquette[3.6]{V.exer3.6}
Let $G$ be a topological group and let $B_G$ be its classifying topos
(IV \hyperref[IV.subsec2.5]{2.5}). Denote by $E_G$ the object of $B_G$ consisting of the
underlying topological space of $G$, endowed with the action of $G$ by
left translations. The canonical morphism from $E_G$ to the final
object $e_G$ of $B_G$ is an epimorphism; hence there is a covering
$\mathcal{U}=(E_G\to e_G)$ and, for every abelian sheaf $F$ of $B_G$,
a spectral sequence
\begin{equation*}
E^{p,q}_2=H^p(\mathcal{U},\Hh^q(F))\Longrightarrow
H^{p+q}(B_G,F){\quoi},\tag{3.6.1}
\end{equation*}\etiquette[3.6.1]{V.3.6.1}
which we propose to study.
\end{enonce*}

\begin{enumerate}
\renewcommand{\labelenumi}{\theenumi)}
\item Show that, for every integer $n$, the topos
$$
B_{G/E_G\times E_G\times\ldots\times E_G}\quad
(n\quad\mbox{factors})
$$
is canonically equivalent to the topos (IV \hyperref[IV.subsec2.5]{2.5})
$$
\TOP(G\times G\times G\times\ldots\times G)\quad
(n-1\quad\mbox{factors}).
$$
For\pageoriginale every integer $n$, denote by $F_n$ the sheaf on the
topological space $G\times\ldots\times G$ ($n$ factors) induced by
$F$.

\item Deduce that the term $E^{p,q}_2$ of
\eqhyperref[V.3.6.1]{3.6.1} is canonically isomorphic to the $p$-th
cohomology group of a complex of the form
$$
H^{q}(e,F_0)\to H^{q}(G,F_1)\to H^{q}(G\times G,F_2)\to
\cdots{\quoi},
$$
which is to be made explicit ($e$ denotes the one-point topological
space).

\item Deduce that the cohomology of the classifying topos $B_G$ with
values in locally constant sheaves is isomorphic to the corresponding
singular cohomology of a classifying space, provided that, on the
spaces of the form $G\times G\times\ldots\times G$, cohomology with
coefficients in locally constant sheaves coincides with the
corresponding singular cohomology (as is the case, for example, when
$G$ is locally contractible). (For classifying spaces, one may consult
\cite{sga4.exp.V.bib.15}.)
\end{enumerate}


% END INCLUDED SGA4_English_current_buildable_source_20260906/source/05/drafts/SGA4_V_LINES_1362_1552_English.texfrag


% BEGIN INCLUDED SGA4_English_current_buildable_source_20260906/source/05/drafts/SGA4_V_LINES_1553_1946_English.texfrag
\section{Acyclic sheaves}\etiquette[4]{V.sec4}

\begin{enonce*}{Definition 4.1}\etiquette[4.1]{V.def4.1}
\emph{Let $(E,A)$ be a ringed topos, let $F$ be an $A$-Module, and
let $S$ be a topologically generating family of $E$. The sheaf $F$
is said to be $S$-acyclic if, for every object $X$ of $S$ and every
integer $q>0$, one has $H^q(X,F)=0$. When $S$ is equal to $\ob E$,
the $S$-acyclic sheaves are called} flasque sheaves.
\end{enonce*}

\setcounter{subsection}{1}
\subsection{}\etiquette[4.2]{V.subsec4.2}
Let $(C,A)$ be a $𝒰$-ringed site, and let $F$ be a sheaf of
$A$-modules on $C$. The sheaf $F$ is said to be $C$-\emph{acyclic}
if it is $S$-acyclic, where $S$ is the family of sheaves associated
with the objects of $C$.

\begin{enonce*}{Proposition 4.3}\etiquette[4.3]{V.prop4.3}
Let $(C,A)$ be a $𝒰$-ringed site, and let $F$ be a sheaf of
$A$-modules. Denote by
$\mathcal{H}^{0}:C^{\sim}_{A}\to C^{\wedge}_{A}$ the inclusion
functor from the $A$-Modules into the category of presheaves of
$A$-modules (here these are $𝒱$-presheaves, where $𝒱$ is a universe
such that $C$ is $𝒱$-small). The following conditions are equivalent:
\begin{enumerate}
\renewcommand{\theenumi}{\roman{enumi}}
\renewcommand{\labelenumi}{\theenumi)}
\item $F$\pageoriginale is $C$-acyclic.

\item For every object $X$ of $C$ and every covering sieve
$R\hookrightarrow X$, one has
$$
H^q(R,\Hh^{0}(F))=0\quad\mbox{for every }q>0{\quoi}.
$$

\item For every object $X$ of $C$, one has
$$
\check{H}^q(X,F)=0\quad\mbox{for every }q>0{\quoi}.
$$
\end{enumerate}
\end{enonce*}

\begin{itemize}
\item[i $\Rightarrow$ ii)]: Since $F$ is $C$-acyclic, the presheaves
$\Hh^q(F)$ vanish for $q>0$. The spectral sequence
\eqhyperref[V.3.2.1]{3.2.1} then gives an isomorphism
$H^q(R,\Hh^{0}(F))\simeq H^q(X,F)$ for every $q$. Consequently,
$H^q(R,\Hh^{0}(F))=0$ for $q>0$.
\item[ii $\Rightarrow$ iii)]: This is clear upon passing to the
inductive limit.
\item[iii $\Rightarrow$ i)]: By induction on $q$, one proves that
the presheaves $\Hh^q(F)$ vanish for $q>0$. For this, use
\eqhyperref[V.3.4.2]{3.4.2}.
\end{itemize}

\setcounter{subsection}{3}
\subsection{}\etiquette[4.4]{V.subsec4.4}
It follows from criterion \hyperref[V.prop4.3]{4.3} ii) and
\hyperref[V.cor3.5]{3.5} that the property of being $S$-acyclic
depends only on the underlying abelian sheaf. In particular, a sheaf
of $A$-modules is flasque if and only if its underlying abelian sheaf
is flasque.

\begin{enonce*}{Corollary 4.5}\etiquette[4.5]{V.cor4.5}
Let $(E,A)$ be a ringed topos, and let $F$ be an $A$-Module. The
following properties are equivalent:
\begin{enumerate}
\renewcommand{\theenumi}{\roman{enumi}}
\renewcommand{\labelenumi}{\theenumi)}
\item $F$ is flasque;

\item For every epimorphic family
$𝒳=(X_i\to X)_{i\in I}$,
$$
H^q(𝒳,F)=0\quad\mbox{for every }q>0{\quoi}.
$$
\end{enumerate}
\end{enonce*}

In \hyperref[V.subsec4.2]{4.2}, take $C$ to be the topos $E$ itself.
The corollary then follows from the equivalence i)
$\Leftrightarrow$ ii) of \hyperref[V.prop4.3]{4.3}.

\setcounter{subsection}{5}
\subsection{}\etiquette[4.6]{V.subsec4.6}
Injective sheaves are flasque. Flasque sheaves are $S$-acyclic for
every topologically generating family $S$. Flasque sheaves are not
necessarily injective (take for $E$ the topos of sets). $S$-acyclic
sheaves are not necessarily flasque (Exercise \hyperref[V.exer4.13]{4.1.3}).

\begin{enonce*}{Proposition 4.7}\etiquette[4.7]{V.prop4.7}
Let $(E,A)$ be a ringed topos, let $F$ be a flasque $A$-Module, and
let $X$ be an object of $E$. For every subobject $Y$ of $X$, the
canonical homomorphism
$H^{0}(X,F)\rightarrow H^{0}(Y,F)$ is\pageoriginale surjective.
\end{enonce*}

Suppose that $Y$ is a subobject of $X$ such that the morphism
$H^{0}(X,F)\to H^{0}(Y,F)$ is not surjective. Let $Z$ be the object
obtained by gluing two copies $X_1$ and $X_2$ of $X$ along $Y$. The
object $Z$ is covered by the two subobjects $X_1$ and $X_2$, and one
has $X_1\times_ZX_2=Y$. Denote the resulting covering by
$𝒳=(X_1,X_2)$. One sees immediately that
$H^1(𝒳,F)\neq0$, a contradiction.

\setcounter{subsection}{7}
\subsection{}\etiquette[4.8]{V.subsec4.8}
For general topoi, the property described in
\hyperref[V.prop4.7]{4.7} does not characterize flasque sheaves
(Exercise \hyperref[V.exer4.15]{4.15}). It does, however,
characterize them for topoi generated by their open subobjects, and
in particular for topoi associated with topological spaces (Exercise
\hyperref[V.exer4.16]{4.16}). In the case of topological spaces, the
terminology ``flasque'' adopted here agrees with that of \cite{sga4.exp.V.bib.7}
(Exercise \hyperref[V.exer4.16]{4.16}).

\begin{enonce*}{Proposition 4.9}\etiquette[4.9]{V.prop4.9}
Let $u:(E,A)\to(E',A')$ be a morphism of ringed topoi.
\begin{enumerate}
\renewcommand{\labelenumi}{\theenumi)}
\item The functor $u_\ast$ carries flasque $A$-Modules to flasque
$A'$-Modules.

\item Let $S$ and $S'$ be topologically generating families of $E$
and $E'$, respectively, such that $u^\ast(S')\subset S$. The functor
$u_\ast$ carries $S$-acyclic $A$-Modules to $S'$-acyclic
$A'$-Modules.

\item When $u$ is a flat morphism (\hyperref[V.def1.8]{1.8}), the
functor $u_\ast$ carries injective $A$-Modules to injective
$A'$-Modules.
\end{enumerate}
\end{enonce*}

Let $X$ be an object of $E'$, let
$𝒳=(X_i\to X)_{i\in I}$ be an epimorphic family, let $F$ be a
flasque $A$-Module, and let $C^{∙}(𝒳,u_\ast F)$ be the Čech complex
of the covering $𝒳$. Using the adjunction between $u_\ast$ and
$u^\ast$, and the fact that $u^\ast$ commutes with fiber products,
one obtains a canonical isomorphism
$$
C^{∙}(𝒳,u_\ast F)\simeq C^{∙}(u^\ast(𝒳),F).
$$
Moreover, $u^\ast$ commutes with inductive limits, and hence
$u^\ast(𝒳)=(u^\ast X_i\to u^\ast X)_{i\in I}$ is an epimorphic
family. Since $F$ is flasque,
$H^q(u^\ast(𝒳),F)=0$ for $q>0$, and consequently
$$
H^q(𝒳,u_\ast F)=H^q(C^{∙}(𝒳,u_\ast F))
=H^q(C^{∙}(u^\ast(𝒳),F))=H^q(u^\ast(𝒳),F)=0
$$
for $q>0$. Thus $u_\ast F$ is flasque.

Let us prove 2). When $S'$ is stable under fiber products, an
argument analogous to the preceding one proves 2). In the general
case, we use the spectral sequence of the morphism $u$
\eqhyperref[V.5.3.2]{5.3.2}. Assertion 2) will not be used
before\pageoriginale \eqhyperref[V.5.3.2]{5.3.2}. Let $F$ be an
$S$-acyclic sheaf. The sheaves $R^qu_\ast F$ are the sheaves
associated with the presheaves $X\mapsto H^q(u^\ast X,F)$. Since
$S'$ is a topologically generating family and $F$ is $S$-acyclic,
one has $R^qu_\ast F=0$ for every $q>0$. The spectral sequence
\eqhyperref[V.5.3.2]{5.3.2} then gives, for every object $X$ of $E'$,
an isomorphism
$H^q(X,u_\ast F)\simeq H^q(u^\ast X,F)$; consequently, for every
$X$ in $S'$, $H^q(X,u_\ast F)=0$.

Let us prove 3). When $u$ is flat, the inverse-image functor
$u^\ast$ for Modules is exact (\hyperref[V.def1.8]{1.8}). Therefore
its right adjoint $u_\ast$ carries injective objects to injective
objects (\hyperref[V.prop0.2]{0.2}).

\begin{enonce*}{Proposition 4.10}\etiquette[4.10]{V.prop4.10}
Let $F$ be an $A$-Module of the ringed topos $(E,A)$, and let $G$ be
an injective $A$-Module.
\begin{enumerate}
\renewcommand{\labelenumi}{\theenumi)}
\item The functor $F\mapsto\Hhom_A(F,G)$ is exact.

\item The abelian sheaf $\Hhom_A(F,G)$ is flasque.
\end{enumerate}
\end{enonce*}

\begin{enonce*}[remark]{Proof}
Let us prove the assertions in turn.
\begin{enumerate}
\renewcommand{\labelenumi}{\theenumi)}
\item Let
$$
0\to F'\to F\to F''\to0
$$
be an exact sequence. We must show that the sequence
$$
0\to\Hhom_A(F'',G)\to\Hhom_A(F,G)\to
\Hhom_A(F',G)\to0
$$
is exact; for this it suffices to show that, for every object $H$ of
$E$, the sequence
\begin{align*}
0\to \Hom_E(H,\Hhom_A(F'',G))\to \Hom_E(H,
\Hhom_A(F,G))\hspace{3cm} \\
\to \Hom_E(H,\Hhom_A(F',G))\to0
\end{align*}
is exact. By (IV \hyperref[IV.cor12.6]{6.12}), the latter sequence is isomorphic to
$$
0\to\Hom_A(A_H\otimes_AF'',G)\to
\Hom_A(A_H\otimes_AF,G)\to
\Hom_A(A_H\otimes_AF',G)\to0,
$$
and the $A$-Module $A_H$ is flat
(\hyperref[V.prop1.3.1]{1.3.1}).

\item Let us show that $\Hhom_A(F,G)$ is flasque. Let
$$
𝒳=(X_i\to X)_{i\in I}{\quoi}
$$
be an epimorphic family of $E$, and let
\[
\xymatrix@M=.25cm@C=1cm{
C_∙(𝒳):\cdots \ar@<5pt>[r] \ar[r] \ar@<-5pt>[r] &
\displaystyle{\mathop{\coprod}_{i,j}}\ZZ_{X_i\times_X X_j}
\ar@<3pt>[r]\ar@<-3pt>[r]&
\displaystyle{\mathop{\coprod}_{i}}\ZZ_{X_i}
}
\]
be the\pageoriginale complex defined in \hyperref[V.prop1.5]{1.5}.
This complex is a flat resolution of the object $\ZZ_X$, which is
itself flat (\hyperref[V.subsec1.4]{1.4} and \hyperref[V.prop1.3.1]{1.3}).

\noindent
We then compute the groups:
$$
H^q(𝒳,\Hhom_A(F,G))\xrightarrow{\sim}
H^q\bigl(\Hom_{\ZZ}^{∙}(C_∙(𝒳),\Hhom_A(F,G))\bigr)
\xrightarrow{\sim}
H^q\bigl(\Hom_A(C_∙(𝒳)\otimes_{\ZZ}F,G)\bigr){\quoi}.
$$
The complex
$$
C_∙(𝒳)\otimes_{\ZZ}F
$$
is acyclic in degrees different from $0$, because $C_∙(𝒳)$ is a
flat resolution of a flat module. Consequently,
$\Hom_A(C_∙(𝒳)\otimes_{\ZZ}F,G)$ is acyclic in degrees different
from $0$.
\end{enumerate}
\end{enonce*}

\begin{enonce*}{Proposition 4.11}\etiquette[4.11]{V.prop4.11}
Let $(E,A)$ be a ringed topos, and let $F$ be a flasque (\resp
injective) sheaf of $A$-modules.
\begin{enumerate}
\renewcommand{\labelenumi}{\theenumi)}
\item For every object $X$ of $E$, the $A|X$-Module $j_X^\ast F$ is
flasque (\resp injective).

\item For every closed subtopos $Z$ of $E$, the sheaf of sections of
$F$ with support in $Z$ \emph{(IV 14)} is flasque (\resp injective).
\end{enumerate}
\end{enonce*}

Assertion 1) follows from \hyperref[V.cor4.5]{4.5} when $F$ is
flasque. The functor $j_X^\ast$ admits an exact left adjoint
$j_{X!}$ (IV 11). Consequently, when $F$ is injective, $j_X^\ast F$
is injective (\hyperref[V.prop0.2]{0.2}). Let us prove 2). Denote the
inclusion morphism by $i:Z\to E$. The sheaf of sections of $F$ with
support in $Z$ is the sheaf $i_\ast i^{!}F$ (IV 14). The functor
$i_\ast i^{!}$ is right adjoint to the exact functor
$i_\ast i^\ast$ (IV 14). It therefore carries injective sheaves to
injective sheaves. Let $U$ be the open complement of $Z$, let
$j:U\to E$ be the inclusion morphism, and let $F$ be a flasque
sheaf. There is an exact sequence (IV 14):
\begin{equation*}
0\to i_\ast i^{!}F\to F\to j_\ast j^\ast F{\quoi}.\tag{4.11.1}
\end{equation*}\etiquette[4.11.1]{V.4.11.1}
For every object $X$ of $E$, one has
$j_\ast j^\ast F(X)=F(X\times U)$, and the morphism
$F(X)\to j_\ast j^\ast F(X)$ induced by the last arrow of
\eqhyperref[V.4.11.1]{4.11.1} comes from the canonical injection
$X\times U\hookrightarrow X$. Since $F$ is flasque, this morphism is
surjective (\hyperref[V.prop4.7]{4.7}), and hence the last arrow of
\eqhyperref[V.4.11.1]{4.11.1} is an epimorphism of presheaves. For
every object $X$ of $E$, the exact cohomology sequence deduced from
\eqhyperref[V.4.11.1]{4.11.1} gives
$H^q(X,i_\ast i^{!}F)=0$ for $q>0$, and hence $i_\ast i^{!}F$ is
flasque.

\setcounter{subsection}{11}
\subsection{}\etiquette[4.12]{V.subsec4.12}
The\pageoriginale property of a sheaf being flasque (\resp
injective) localizes \eqhyperref[V.4.11.1]{4.11.1}. It is not, in
general, a property of local nature (Exercise
\hyperref[V.exer4.15]{4.15}). It is, however, a property of local
nature for topoi generated by their open subobjects, and in
particular for topoi associated with topological spaces (Exercise
\hyperref[V.exer4.16]{4.16}).

\begin{enonce*}[remark]{Exercise 4.13}\etiquette[4.13]{V.exer4.13}
Let $X$ be a locally compact space, and let $F$ be a $c$-soft sheaf
\cite{sga4.exp.V.bib.7}. Using the local nature of softness (\loccit), show that
the restriction of $F$ to every paracompact open subset of $X$ is a
soft sheaf. Show that the paracompact open subsets form a basis for
the topology of $X$. Deduce that, if $S$ denotes the family of
paracompact open subsets of $X$, then $F$ is $S$-acyclic. Show that
the sheaf of continuous functions on $X$ is $S$-acyclic but is not
flasque when $X$ is not discrete.
\end{enonce*}

\begin{enonce*}[remark]{Problem 4.14}\etiquette[4.14]{V.prob4.14}
Study totally acyclic topoi, \ie topoi such that
$H^q(X,F)=0$ for every $q>0$, every object $X$, and every abelian
sheaf $F$.
\end{enonce*}

\begin{enonce*}[remark]{Exercise 4.15}\etiquette[4.15]{V.exer4.15}
Let $G$ be a discrete group, and let $B_G$ be its classifying topos
(IV \hyperref[IV.subsec2.4]{2.4}). Show that for every abelian sheaf $F$ and every
monomorphism $X\hookrightarrow Y$, the homomorphism $F(Y)\to F(X)$
is surjective. Let $E(G)$ be the group $G$ regarded as a homogeneous
space under itself. It is an object of $B_G$. The topos
$B_{G/E(G)}$ is equivalent to the punctual topos (IV 8). The
morphism $E(G)\to e$ ($e$ the final object of $B_G$) is an
epimorphism. For every abelian sheaf $F$ on $B_G$, the sheaf
$F|E(G)$ is flasque. Deduce that the property of being flasque or
injective is not of local nature.
\end{enonce*}

\begin{enonce*}[remark]{Exercise 4.16}\etiquette[4.16]{V.exer4.16}
A topos $E$ is said to be generated by its open subobjects if the
open subobjects of $E$ (\ie the subobjects of the final object $e$ of
$E$) form a generating family (I 7). Such a topos has the following
property:

\noindent
(P) Every epimorphic family $X_i\to X$, $i\in I$, admits a
refinement by an epimorphic family $U_j\to X$, $j\in J$, in which
the $U_j\to X$ are monomorphisms.

\begin{enumerate}
\renewcommand{\theenumi}{\alph{enumi}}
\renewcommand{\labelenumi}{\theenumi)}
\item Are there\pageoriginale topoi which have property (P) but are
not generated by their open subobjects? Topoi associated with
topological spaces are generated by their open subobjects. If $E$
is generated by its open subobjects (\resp has (P)), then for every
object $X$ of $E$, $E_{/X}$ is generated by its open subobjects
(\resp has (P)). Every subtopos of a topos generated by its open
subobjects is generated by its open subobjects. Property (P) is not
a property of local nature.

\item Let $E$ be a topos, and let $\Ouv(E)$ be the category of open
subobjects of $E$, endowed with the induced topology. The inclusion
functor $\Ouv(E)\to E$ is a morphism of sites, and hence gives a
morphism of topoi $\Pi:E\to\Ouv(E)^{\sim}$. The functor $\Pi^{\ast}$
is fully faithful. With respect to the 2-category of topoi generated
by their open subobjects, the morphism $\Pi$ has a universal property
which the reader is invited to make explicit.

\item Let $E$ be a topos generated by its open subobjects, and let
$\Point(E)$ be the set of points of $E$ up to isomorphism. Show that
$\Point(E)$ is small. Put a topology on $\Point(E)$ and define a
morphism $\Top(\Point(E))\to\Ouv(E)^{\sim}$ making
$\Top(\Point(E))$ a subtopos of $\Ouv(E)^{\sim}$. For every
topological space $X$, show that every morphism from $\Top(X)$ to
$E$ gives a morphism from $\Top(X)$ to $\Top(\Point(E))$. Show that
a topos generated by its open subobjects is equivalent to a topos
$\Top(X)$ ($X$ a topological space) if and only if it has enough
points. Show that there exist topoi generated by their open
subobjects which are not equivalent to any topos $\Top(X)$, where
$X$ is a topological space.

\item Let $E$ be a topos having property (P). An abelian sheaf $F$
on $E$ is flasque if and only if, for every monomorphism
$X\hookrightarrow Y$, the homomorphism $F(Y)\to F(X)$ is
surjective.

\item Let $E$ be a topos having property (P). Show that the property
of a sheaf $F$ being flasque is a property of local nature.

\item Let $E$ be a topos generated by its open subobjects, and let
$e$ be the final object of $E$. An abelian sheaf $F$ is flasque if
and only if, for every open subobject $U$ of $e$, the canonical
morphism $F(e)\to F(U)$ is surjective.
\end{enumerate}
\end{enonce*}

\begin{enonce*}[remark]{Exercise 4.17}\etiquette[4.17]{V.exer4.17}
(Flasque\pageoriginale sheaves and change of universe)

Let $C$ be a $𝒰$-site (for example, a $𝒰$-topos), and let $𝒱$ be a
universe containing $𝒰$. Denote the corresponding categories of
sheaves by $C^{\sim}_{𝒰}$ and $C^{\sim}_{𝒱}$, and denote the
canonical inclusion by
$\epsilon:C^{\sim}_{𝒰}\to C^{\sim}_{𝒱}$. Let $F$ be a flasque
abelian $𝒰$-sheaf on $C$. We shall show that $\epsilon F$ is
flasque. First note that, for every object $X$ of
$C^{\sim}_{𝒰}$, one has $H^q(\epsilon X,\epsilon F)=0$ for $q>0$.
Every object $Y$ of $C^{\sim}_{𝒱}$ admits an epimorphic family
$\epsilon X_i\to Y$, $i\in I$, in which the $\epsilon X_i\to Y$
are monomorphisms. It follows that
$$
H^q(Y,\epsilon F)=displaystyle
\varprojlim_{\epsilon X\hookrightarrow Y}^{q}F(X).
$$
To show that these $\varprojlim^q$ vanish for $q\neq0$, one may
reduce to the topos $\Ouv(C_{𝒱/Y})^{\sim}$ and use the fact that, for
this topos, a sheaf is flasque if it is so locally.
\end{enonce*}


% END INCLUDED SGA4_English_current_buildable_source_20260906/source/05/drafts/SGA4_V_LINES_1553_1946_English.texfrag


% BEGIN INCLUDED SGA4_English_current_buildable_source_20260906/source/05/drafts/SGA4_V_LINES_1947_2047_English.texfrag
\section{The $R^qu_*$ and the Cartan--Leray spectral sequence for a morphism of topoi}\etiquette[5]{V.sec5}

\setcounter{subsection}{-1}
\subsection{}\etiquette[5.0]{V.subsec5.0}
Let $u:(E,A)\to(E',A')$ be a morphism of ringed topoi, and let
$u_\ast:E\to E'$ be the direct-image functor. The notation $u_\ast$
will also denote the extension of the direct-image functor to
modules. The functor $u_\ast$ on modules (left modules, to fix ideas)
is left exact. Its right derived functors are denoted by
$R^qu_\ast$, $q\geq 0$.

\begin{enonce*}{Proposition 5.1}\etiquette[5.1]{V.prop5.1}
\begin{enumerate}
\renewcommand{\labelenumi}{\theenumi)}
\item For every $A$-module $M$, the sheaf $R^qu_\ast M$ is the
sheaf associated with the presheaf $X'\mapsto H^q(u^\ast X',M)$
$(X'\in\ob E')$.

\item Formation of the functors $R^qu_\ast$ commutes with
restriction of scalars.

\item Formation of the functors $R^qu_\ast$ commutes with
localization. More precisely, for every object $X'$ of $E'$, if
$u_{/X'}:E_{/u^\ast X'}\to E'_{/X'}$ denotes the morphism deduced
from $u$ by localization \emph{(IV 8)}, then, for every $A$-module
$M$, there is a canonical isomorphism
\begin{equation*}
R^q(u_{/X'})_\ast(M|u^\ast X')\simeq R^qu_\ast(M)|X'\qquad q\geq
0{\quoi}.\tag{5.1.1}
\end{equation*}\etiquette[5.1.1]{V.5.1.1}
\end{enumerate}
\end{enonce*}

Let $u^{\wedge}_\ast:E^{\wedge}\to E^{\prime\wedge}$ denote the
direct-image functor for $𝒰$-presheaves
$(u^{\wedge}_\ast M=M\circ u^\ast)$. Since $u^\ast$ and $u_\ast$
are adjoint, there is an isomorphism
$$
u_\ast\simeq\underline{a}u^{\wedge}_\ast{}^{0}{\quoi},
$$
where\pageoriginale $\underline{a}$ is the associated-sheaf functor
for $E'$. Since the functors $\underline{a}$ and $u^{\wedge}_\ast$
are exact, we have
$$
R^qu_\ast\simeq\underline{a}u^{\wedge}_\ast\Hh^q{\quoi},
$$
which is another way of stating 1). Assertion 2) then follows from
1) and \hyperref[V.cor3.5]{3.5}. By definition of the morphism
$u_{/X'}$, there is a canonical isomorphism
\eqhyperref[V.5.1.1]{5.1.1} for $q=0$. The general case follows by
noting that the localization functors (IV 8) are exact and take
injective objects to injective objects (\hyperref[V.prop4.11]{4.11}).

\begin{enonce*}{Proposition 5.2}\etiquette[5.2]{V.prop5.2}
Let $u:E\to E'$ be a morphism of topoi and let $S'$ be a generating
family of $E'$. The sheaves $M$ which are acyclic for the functors
$H^0(u^\ast X',\mathord\cdot)$, $X'\in S'$, are acyclic for the
functor $u_\ast$. In particular, flabby sheaves are acyclic for
$u_\ast$.
\end{enonce*}

This follows from \hyperref[V.prop5.1]{5.1}, 1).

\begin{enonce*}{Proposition 5.3}\etiquette[5.3]{V.prop5.3}
Let $u:E\to E'$ be a morphism of topoi and let $M$ be an abelian
group of $E$. There is a spectral sequence
\begin{equation*}
E^{p,q}_2=H^p(E',R^qu_\ast M)\Longrightarrow H^{p+q}(E,M){\quoi}.
\tag{5.3.1}
\end{equation*}\etiquette[5.3.1]{V.5.3.1}
More generally, for every object $X'$ of $E'$, there is a spectral
sequence
\begin{equation}
E^{p,q}_2=H^p(X',R^qu_\ast M)\Rightarrow H^{p+q}(u^\ast X',M){\quoi}.
\tag{5.3.2}
\end{equation}\etiquette[5.3.2]{V.5.3.2}
\end{enonce*}

By definition of the direct- and inverse-image functors, there is an
isomorphism $H^0(X',u_\ast M)\simeq H^0(u^\ast X',M)$. The functor
$u_\ast$ takes injective objects to flabby sheaves
(\hyperref[V.subsec4.6]{4.6} and \hyperref[V.prop4.9]{4.9}). The
asserted spectral sequences are therefore spectral sequences of
composite functors (\hyperref[V.prop0.3]{0.3}).

\begin{enonce*}{Proposition 5.4}\etiquette[5.4]{V.prop5.4}
Let $u:E\to E'$ and $v:E'\to E''$ be two morphisms of topoi. There
is a spectral sequence
$$
R^pv_\ast R^qu_\ast\Longrightarrow R^{p+q}(v\circ u)_\ast{\quoi}.
$$
\end{enonce*}

We have $v_\ast u_\ast\simeq(vu)_\ast$, and the functor $u_\ast$
takes injective objects to\pageoriginale flabby sheaves
(\hyperref[V.prop4.9]{4.9}), hence to objects acyclic for $v_\ast$
(\hyperref[V.prop5.2]{5.2}). Thus there is a spectral sequence of
composite functors (\hyperref[V.prop0.3]{0.3}).


% END INCLUDED SGA4_English_current_buildable_source_20260906/source/05/drafts/SGA4_V_LINES_1947_2047_English.texfrag


% BEGIN INCLUDED SGA4_English_current_buildable_source_20260906/source/05/drafts/SGA4_V_LINES_2048_2300_English.texfrag
\section{Local Ext and cohomology with supports}\etiquette[6]{V.sec6}

\setcounter{subsection}{-1}
\subsection{}\etiquette[6.0]{V.subsec6.0}

Let $(E,A)$ be a ringed topos, and let $M$ be an $A$-module, a left
module to fix ideas. The functor
$N\longmapsto\Hhom_A(M,N)$, from the category of left $A$-modules to
the category of abelian groups, is left exact (IV 12). Its right
derived functors are denoted by
\begin{equation*}
\Eext_A^q(M,N){\quoi}.\tag{6.0.1}
\end{equation*}\etiquette[6.0.1]{V.6.0.1}
In particular, one has
\begin{equation*}
\Eext_A^{0}(M,N)=\Hhom_A(M,N){\quoi}.\tag{6.0.2}
\end{equation*}\etiquette[6.0.2]{V.6.0.2}
By definition, there are canonical isomorphisms
(\hyperref[V.subsec2.1]{2.1} and IV 12)
\begin{equation*}
H^{0}(E,\Eext_A^{0}(M,N))=\Ext^{0}_A(E;M,N)
=\Hom_A(M,N){\quoi}.\tag{6.0.3}
\end{equation*}\etiquette[6.0.3]{V.6.0.3}
More generally, for every object $X$ of $E$, one has
(\hyperref[V.subsec2.2]{2.2})
\begin{equation*}
H^{0}(X,\Eext_A^{0}(M,N))=\Ext^{0}_A(X;M,N)
=\Hom_{A|X}(M|X,N|X){\quoi}.\tag{6.0.4}
\end{equation*}\etiquette[6.0.4]{V.6.0.4}

\begin{enonce*}{Proposition 6.1}\etiquette[6.1]{V.prop6.1}
\begin{enumerate}
\renewcommand{\labelenumi}{\theenumi)}
\item Formation of the functors $\Eext^q_A$ commutes with
localization. More precisely, for every object $X$ of $E$, there are
functorial isomorphisms
\begin{equation*}
\Eext^q_A(M,N)|X\simeq
\Eext^q_{A|X}(M|X,N|X){\quoi}.\tag{6.1.1}
\end{equation*}\etiquette[6.1.1]{V.6.1.1}

\item The sheaf $\Eext^q_A(M,N)$ is isomorphic to the sheaf associated
with the presheaf $X\longmapsto\Ext^q_A(X;M,N)$.

\item There is a spectral sequence
\begin{equation*}
\Ext^{p+q}_A(E;M,N)\Longleftarrow E^{p,q}_2
=H^p(E,\Eext^q_A(M,N)){\quoi}.\tag{6.1.2}
\end{equation*}\etiquette[6.1.2]{V.6.1.2}
\end{enumerate}
More\pageoriginale generally, for every object $X$ of $E$, there is a
spectral sequence functorial in $X$ and in the arguments $M$ and $N$:
\begin{equation*}
\Ext^{p+q}_A(X;M,N)\Longleftarrow E^{p,q}_2
=H^p(X,\Eext^q_A(M,N)){\quoi}.\tag{6.1.3}
\end{equation*}\etiquette[6.1.3]{V.6.1.3}
\end{enonce*}

By definition, \eqhyperref[V.6.1.1]{6.1.1} is an isomorphism when
$q=0$ (IV 12). The general case follows, taking into account the fact
that the localization functors are exact and send injective modules to
injective modules (\hyperref[V.subsec2.2]{2.2}). The functor
$N\longmapsto\Hhom_A(M,N)=\Eext^0_A(M,N)$ sends injective modules to
flasque sheaves, and hence to sheaves acyclic for $H^{0}(X,⋅)$
(\hyperref[V.prop4.10]{4.10}). The spectral sequences
\eqhyperref[V.6.1.2]{6.1.2} and \eqhyperref[V.6.1.3]{6.1.3} are the
spectral sequences of composite functors, in view of
\eqhyperref[V.6.0.3]{6.0.3} and \eqhyperref[V.6.0.4]{6.0.4}. As $X$
varies in $E$, the spectral sequence \eqhyperref[V.6.1.3]{6.1.3}
gives a spectral sequence of presheaves and hence, upon passing to the
associated sheaves, a spectral sequence of sheaves. Since the sheaves
associated with the presheaves $X\longmapsto H^p(X,⋅)$ are zero when
$p\neq0$ (\hyperref[V.prop3.1]{3.1}), this spectral sequence
degenerates and gives the isomorphism asserted in 2).

\begin{enonce*}{Proposition 6.2}\etiquette[6.2]{V.prop6.2}
The functors $(M,N)\longmapsto\Eext^q_A(M,N)$, $q\geq0$, form a
$\delta$-functor in the variable $M$ and in the variable $N$.
Explicitly, for every exact sequence
$0\to N'\to N\to N''\to0$
(\resp $0\to M'\to M\to M''\to0$), there are long exact sequences
\begin{equation*}
\cdots\to\Eext^q_A(M,N')\to\Eext^q_A(M,N)\to\Eext^q_A(M,N'')
\xrightarrow{\delta}\Eext^{q+1}_A(M,N')\to\cdots\tag{6.2.1}
\end{equation*}\etiquette[6.2.1]{V.6.2.1}
(\resp
\begin{equation*}
\cdots\to\Eext^q_A(M'',N)\to\Eext^q_A(M,N)\to\Eext^q_A(M',N)
\xrightarrow{\delta}\Eext^{q+1}_A(M'',N)\to\cdots){\quoi}.
\tag{6.2.2}
\end{equation*}\etiquette[6.2.2]{V.6.2.2}
\end{enonce*}

It follows from the general properties of the functors $\Ext^q$ that,
for every object $X$ of $E$, the functors
$(M,N)\longmapsto\Ext^q_A(X;M,N)$, $q\geq0$, form a
$\delta$-functor in each argument. The assertion follows by allowing
$X$ to vary and taking the associated sheaf
(\hyperref[V.prop6.1]{6.1}).

\setcounter{subsection}{2}
\subsection{}\etiquette[6.3]{V.subsec6.3}
Let $(E,A)$ be a ringed topos, let $M$ be an $A$-module, let $Z$ be a
closed subtopos of $E$ (IV 9), and let $U$ be its open complement.
Denote by $H^{0}_Z(E,M)$ the group of sections of $M$ whose
support\pageoriginale is contained in $Z$ (IV 14), and by
$ℋ^{0}_Z(M)$ the subsheaf of $M$ defined by the sections of $M$
``with supports in $Z$'' (IV 14). The functors $H^{0}_Z(E,⋅)$ and
$ℋ^{0}_Z(⋅)$ are left exact (IV 14). Their derived functors are denoted
by $H^q_Z(E,⋅)$ and $ℋ^q_Z(⋅)$ respectively, and are called the
\emph{cohomology groups (\emph{\resp} sheaves) of $M$ with supports in
$Z$}\footnote{Compare with \SGA 2 I for the case of ordinary
topological spaces, as well as with the exposé of \textsc{Hartshorne}
cited on p.~80 below.}.

There are canonical isomorphisms (IV 14)
\hypertarget{sga4.exp.V.equation.6.3.1}{}\begin{align*}
H^{0}_Z(E,M)&\simeq\Hom_A(A_Z,M)\simeq
\Ext^{0}_A(E;A_Z,M){\quoi},\tag{6.3.1}\\
ℋ^{0}_Z(M)&\simeq\Hhom_A(A_Z,M)\simeq
\Eext^{0}_A(A_Z,M){\quoi},\tag{6.3.2}
\end{align*}\hypertarget{sga4.exp.V.equation.6.3.2}{}
and hence, for every $q\geq0$, there are isomorphisms
\hypertarget{sga4.exp.V.equation.6.3.3}{}\begin{align*}%\etiquette[6.3.3]{V.6.3.3}
H^q_Z(E,M)&\simeq\Ext^q_A(E;A_Z,M){\quoi},\tag{6.3.3}\\
ℋ^q_Z(M)&\simeq\Eext^q_A(A_Z,M){\quoi}.\tag{6.3.4}
\end{align*}\hypertarget{sga4.exp.V.equation.6.3.4}{}
Note that, since $A_Z$ is a bimodule, the sheaves
$\Eext^q_A(A_Z,M)\simeqℋ^q_Z(M)$ are canonically endowed with
$A$-module structures.

For every object $X$ of $E$, denote by $Z_{/X}$ the closed subtopos of
$E_{/X}$ obtained from $Z$ by localization (it is the complement of
$U\times X$). By definition, there are canonical isomorphisms
\begin{equation*}
H^{0}(X,ℋ^{0}_Z(M))\simeq
H^{0}_{Z_{/X}}(E_{/X},M|X){\quoi}.\tag{6.3.5}
\end{equation*}\etiquette[6.3.5]{V.6.3.5}
Set
\begin{equation*}
H^q_Z(X,M)\simeq H^q_{Z_{/X}}(E_{/X},M|X){\quoi}.\tag{6.3.6}
\end{equation*}\etiquette[6.3.6]{V.6.3.6}
Taking \hyperlink{sga4.exp.V.equation.6.3.2}{(6.3.2)} into account, there are canonical isomorphisms
\begin{equation*}
H^q_Z(X,M)\simeq\Ext^q_A(X;A_Z,M){\quoi}.\tag{6.3.7}
\end{equation*}\etiquette[6.3.7]{V.6.3.7}

\begin{enonce*}{Proposition 6.4}\etiquette[6.4]{V.prop6.4}
\begin{enumerate}
\renewcommand{\labelenumi}{\theenumi)}
\item Formation\pageoriginale of the functors $ℋ^q_Z$ commutes with
localization. More precisely, for every object $X$ of $E$ and every
$A$-module $M$, there are canonical isomorphisms
\begin{equation*}
ℋ^{q}_Z(M)|X\simeqℋ^q_{Z_{/X}}(M|X){\quoi}.\tag{6.4.1}
\end{equation*}\etiquette[6.4.1]{V.6.4.1}

\item The sheaf $ℋ^q_Z(M)$ is the sheaf associated with the presheaf
$X\longmapsto H^q_Z(X,M)$.

\item There is a spectral sequence
\begin{equation*}
H^{p+q}_Z(E,M)\Longleftarrow H^p(E,ℋ^q_Z(M)){\quoi}.\tag{6.4.2}
\end{equation*}\etiquette[6.4.2]{V.6.4.2}
More generally, for every object $X$ of $E$, there is a spectral
sequence
\begin{equation*}
H^{p+q}_Z(X,M)\Longleftarrow H^p(X,ℋ^q_Z(M)){\quoi}.\tag{6.4.3}
\end{equation*}\etiquette[6.4.3]{V.6.4.3}
\end{enumerate}
\end{enonce*}

This is simply Proposition \hyperref[V.prop6.1]{6.1} translated by
means of the dictionary in \hyperref[V.subsec6.3]{6.3}.

\begin{enonce*}{Proposition 6.5}\etiquette[6.5]{V.prop6.5}
With the notation of \hyperref[V.subsec6.3]{6.3}, denote by
$j:U\to E$ the canonical morphism. For every $A$-module $M$, there is
an exact sequence of sheaves
\begin{equation*}
0\toℋ^{0}_Z(M)\to M\to j_\ast(M|U)\to
ℋ^1_Z(M)\to0{\quoi},\tag{6.5.1}
\end{equation*}\etiquette[6.5.1]{V.6.5.1}
and, for $q\geq2$, there are isomorphisms
\begin{equation*}
ℋ^q_Z(M)\simeq\Eext^{q-1}_A(A_U,M)\simeq
R^{q-1}j_\ast(M|U){\quoi}.\tag{6.5.2}
\end{equation*}\etiquette[6.5.2]{V.6.5.2}
There is, in addition, a long exact sequence
\begin{equation*}
\cdots\to H^q_Z(E,M)\to H^q(E,M)\to H^q(U,M)\to
H^{q+1}_Z(E,M)\to\cdots
\footnote{This exact sequence makes precise the role of the global
cohomological invariants $H^q_Z(E,M)$ as ``cohomology groups of $E$
\emph{modulo} the open subtopos $U$, with coefficients in $M$.''}
\tag{6.5.3}
\end{equation*}\etiquette[6.5.3]{V.6.5.3}
and, more generally, for every object $X$ of $E$, there is a long
exact sequence
\begin{equation*}
\cdots\to H^q_Z(X,M)\to H^q(X,M)\to H^q(X\times U,M)\to
H^{q+1}_Z(X,M)\to\cdots.\tag{6.5.4}
\end{equation*}\etiquette[6.5.4]{V.6.5.4}
\end{enonce*}

By the definition of $A_Z$, there is an exact sequence (IV 14)
\begin{equation*}
0\to A_U\to A\to A_Z\to0{\quoi}.\tag{6.5.5}
\end{equation*}\etiquette[6.5.5]{V.6.5.5}
The\pageoriginale functors $\Eext^q_A(A,⋅)$ vanish for $q>0$. One has
$\Hhom_A(A_U,M)\simeq j_\ast(M|U)$ (IV 14), and hence
(\hyperref[V.subsec2.2]{2.2})
$\Eext^q_A(A_U,M)\simeq R^qj_\ast(M|U)$. Finally,
$\Eext^q_A(A_Z,M)\simeqℋ^q_Z(M)$ \hyperlink{sga4.exp.V.equation.6.3.4}{(6.3.4)}. The long exact sequence
\eqhyperref[V.6.2.2]{6.2.2} gives, in this case,
\eqhyperref[V.6.5.1]{6.5.1} and \eqhyperref[V.6.5.2]{6.5.2}. The exact
sequences \eqhyperref[V.6.5.3]{6.5.3} and
\eqhyperref[V.6.5.4]{6.5.4} follow from the dictionary in
\hyperref[V.subsec6.3]{6.3} and from the long exact sequence of the
$\delta$-functor $\Ext^q_A(X;\mathord\cdot,M)$, $q\geq0$, associated
with \eqhyperref[V.6.5.5]{6.5.5}.

\begin{enonce*}{Proposition 6.6}\etiquette[6.6]{V.prop6.6}
\begin{enumerate}
\renewcommand{\labelenumi}{\theenumi)}
\item Flasque sheaves are acyclic for the functors $ℋ^{0}_Z$ and
$H^{0}_Z(X,⋅)$.

\item For every $q\geq0$, the functors $ℋ^{q}_Z$ and $H^{q}_Z(X,⋅)$
commute with restriction of scalars.
\end{enumerate}
\end{enonce*}

Let $M$ be a flasque sheaf. The exact sequence
\eqhyperref[V.6.5.4]{6.5.4} gives the equalities
$H^q_Z(X,M)=0$ for every $X$ and every $q\geq2$, together with an
exact sequence
$$
0\to H^{0}_Z(X,M)\to H^{0}(X,M)\to H^{0}(X\times U,M)
\to H^1_Z(X,M)\to0{\quoi}.
$$
But, since the sheaf $M$ is flasque, the morphism
$H^{0}(X,M)\to H^{0}(X\times U,M)$ is surjective
(\hyperref[V.sec4]{4}). Consequently, $H^1_Z(X,M)=0$, and $M$ is
acyclic for $H^{0}_Z(X,⋅)$. Passing to the associated sheaves, one
deduces that $M$ is acyclic for the functor $ℋ^{0}_Z$
(\hyperref[V.prop6.4]{6.4}). It is clear that the functors $ℋ^{0}_Z$
and $H^{0}_Z(X,⋅)$ commute with restriction of scalars; since flasque
sheaves are acyclic for both functors, assertion 2) follows.

\begin{enonce*}{Proposition 6.7}\etiquette[6.7]{V.prop6.7}
Let $(E,A)$ be a ringed topos, let $Z$ be a closed subtopos of $E$,
let $U$ be its open complement, and let $M$ and $N$ be two
$A$-modules. There are isomorphisms, functorial in $M$ and $N$ and
compatible with changes of closed subtopos:
$$
ℋ^{0}_Z(\Hhom_A(M,N))\simeq\Hhom_A(M,ℋ^0_Z(N))
\simeq\Hhom_A(M\otimes_AA_Z,N){\quoi}.
$$
\end{enonce*}

This follows from (IV 14), taking \hyperlink{sga4.exp.V.equation.6.3.2}{(6.3.2)} into account.

\setcounter{subsection}{7}
\subsection{}\etiquette[6.8]{V.subsec6.8}

% END INCLUDED SGA4_English_current_buildable_source_20260906/source/05/drafts/SGA4_V_LINES_2048_2300_English.texfrag


% BEGIN INCLUDED SGA4_English_current_buildable_source_20260906/source/05/drafts/SGA4_V_LINES_2301_2629_English.texfrag
Set
\begin{equation*}
\Hhom_{A,Z}(M,N)=ℋ^{0}_Z(\Hhom_A(M,N)){\quoi}.\tag{6.8.1}
\end{equation*}\etiquette[6.8.1]{V.6.8.1}
The functor $N\longmapsto\Hhom_{A,Z}(M,N)$ is left exact. Its derived
functors are denoted by $\Eext^q_{A,Z}(M,N)$: these are the
\emph{$\Eext$ sheaves with supports in $Z$}. Thus
\begin{equation*}
\Eext^{0}_{A,Z}(M,N)=\Hhom_{A,Z}(M,N){\quoi}.\tag{6.8.2}
\end{equation*}\etiquette[6.8.2]{V.6.8.2}
Set\pageoriginale $M_Z=M\otimes_AA_Z$. It follows from IV 14 that
$M_Z$ is the direct image on $E$ of the inverse image of $M$ on $Z$.
Thus, taking \hyperref[V.prop6.7]{6.7} into account, there are
isomorphisms
\begin{equation*}
\Eext^{q}_{A,Z}(M,N)\simeq\Eext^{q}_{A}(M_Z,N){\quoi}.\tag{6.8.3}
\end{equation*}\etiquette[6.8.3]{V.6.8.3}
We now pass to the global invariants. Set
\begin{equation*}
\Hom_{A,Z}(M,N)=\Ext^{0}_{A,Z}(E;M,N)
=H^{0}_Z(E,\Hhom_A(M,N)){\quoi}.\tag{6.8.4}
\end{equation*}\etiquette[6.8.4]{V.6.8.4}
The group $\Hom_{A,Z}(M,N)$ is the subgroup of the group of
morphisms from $M$ to $N$ whose support is in $Z$ (IV 14), \ie those
which vanish on $U$. More generally, for every object $X$ of $E$, set
\begin{equation*}
\Ext^{0}_{A,Z}(X;M,N)=H^{0}_Z(X,\Hhom_A(M,N)){\quoi}.\tag{6.8.5}
\end{equation*}\etiquette[6.8.5]{V.6.8.5}
The functors $N\longmapsto\Ext^{0}_{A,Z}(X;M,N)$ are left exact.
Their derived functors are denoted by $\Ext^q_{A,Z}(X;M,N)$. These
are the \emph{$\Ext$ groups with supports in $Z$}. Definitions
\eqhyperref[V.6.3.6]{6.3.6}, \eqhyperref[V.6.8.1]{6.8.1}, and
\eqhyperref[V.6.8.5]{6.8.5}, together with the isomorphisms of
\hyperref[V.prop6.7]{6.7}, give isomorphisms
\begin{equation*}
\Ext^{0}_{A,Z}(X;M,N)\simeq
\begin{cases}
H^{0}(X,\Hhom_{A,Z}(M,N)){\quoi},\\[3pt]
\Ext^{0}_A(X;M,ℋ^{0}_Z(N)){\quoi},\\[3pt]
\Ext^{0}_A(X;M_Z,N){\quoi}.
\end{cases}\tag{6.8.6}
\end{equation*}\etiquette[6.8.6]{V.6.8.6}
The last isomorphism in \eqhyperref[V.6.8.6]{6.8.6} gives
isomorphisms
\begin{equation*}
\Ext^{q}_{A,Z}(X;M,N)\simeq
\Ext^{q}_{A}(X;M_Z,N){\quoi}.\tag{6.8.7}
\end{equation*}\etiquette[6.8.7]{V.6.8.7}

\begin{enonce*}{Proposition 6.9}\etiquette[6.9]{V.prop6.9}
\begin{enumerate}
\renewcommand{\labelenumi}{\theenumi)}
\item There are two spectral sequences, functorial in $M$ and $N$ and
compatible with changes of closed subtopos:
\begin{equation*}
\Eext^{p+q}_{A,Z}(M,N)\Longleftarrow\tag{6.9.1}
\begin{cases}
E^{p,q}_2=ℋ^p_Z(\Eext^q_A(M,N)){\quoi},\\[3pt]
'E^{p,q}_2=\Eext^p_A(M,ℋ^q_Z(N)){\quoi}.
\end{cases}
\end{equation*}\etiquette[6.9.1]{V.6.9.1}

\item There\pageoriginale are three spectral sequences, functorial in
$X$, $M$, and $N$ and compatible with changes of closed subtopos:
\begin{equation*}
\Ext^{p+q}_{A,Z}(X;M,N)\Longleftarrow\tag{6.9.2}
\begin{cases}
E^{p,q}_2=H^p_Z(X,\Eext^q_A(M,N)){\quoi},\\[3pt]
'E^{p,q}_2=H^p(X,\Eext^q_{A,Z}(M,N)){\quoi},\\[3pt]
``E^{p,q}_2=\Ext^p_A(X;M,ℋ^q_Z(N)).
\end{cases}
\end{equation*}\etiquette[6.9.2]{V.6.9.2}

\item The sheaves $\Eext^q_{A,Z}(M,N)$ are canonically isomorphic to
the sheaves associated with the presheaves
$X\longmapsto\Ext^q_{A,Z}(X;M,N)$.
\end{enumerate}
\end{enonce*}

When $N$ is injective, the sheaf $\Hhom_A(M,N)$ is flasque
(\hyperref[V.prop4.10]{4.10}), and hence acyclic for $ℋ^{0}_Z$
(\hyperref[V.prop6.6]{6.6}). The first spectral sequence in
\eqhyperref[V.6.9.1]{6.9.1} is a spectral sequence of composite
functors \eqhyperref[V.6.8.1]{6.8.1}. Likewise, when $N$ is injective,
$ℋ^{0}_Z(N)$ is injective (\hyperref[V.prop4.11]{4.11}), and the
second spectral sequence in \eqhyperref[V.6.9.1]{6.9.1} is a spectral
sequence of composite functors deduced from
\hyperref[V.prop6.7]{6.7}. The spectral sequences in
\eqhyperref[V.6.9.2]{6.9.2} are spectral sequences of composite
functors deduced from definition \eqhyperref[V.6.8.5]{6.8.5} and the
first two isomorphisms in \eqhyperref[V.6.8.6]{6.8.6}. Finally, by
allowing $X$ to vary in the second spectral sequence in
\eqhyperref[V.6.9.2]{6.9.2} and taking the associated sheaves, one
obtains a spectral sequence of sheaves which degenerates by
\hyperref[V.prop3.1]{3.1} and gives the isomorphisms in 3).

\begin{enonce*}{Proposition 6.10}\etiquette[6.10]{V.prop6.10}
The functors $(M,N)\longmapsto\Eext^q_{A,Z}(M,N)$, $q\geq0$, and
$(M,N)\longmapsto\Ext^q_{A,Z}(X;M,N)$, $q\geq0$, are
$\delta$-functors in each of the variables $M$ and $N$. Denoting by
$M_U$ the sheaf $M\otimes_AA_U$ (\cf IV 11), there is a long exact
sequence
\begin{equation}
\ldots\xrightarrow{\delta}\Eext^q_{A,Z}(M,N)\to\Eext^q_A(M,N)
\to\Eext^q_A(M_U,N)\xrightarrow{\delta}
\Eext^{q+1}_{A,Z}(M,N)\to\ldots\tag{6.10.1}
\end{equation}\etiquette[6.10.1]{V.6.10.1}
and an analogous exact sequence for the $\Ext$ groups.
\end{enonce*}

The functors $\Eext^q_A(⋅,⋅)$ and $\Ext^q_A(X;⋅,⋅)$ form
$\delta$-functors with respect to each variable
(\hyperref[V.prop6.2]{6.2}), and the functor
$M\longmapsto M_Z$ is exact because $A_Z$ is flat
(\hyperref[V.cor1.3.3]{1.3.3}). The first assertion therefore follows
from the isomorphisms \eqhyperref[V.6.8.3]{6.8.3} and
\eqhyperref[V.6.8.7]{6.8.7}. The exact sequence
\eqhyperref[V.6.10.1]{6.10.1}, and the analogous exact sequence for
the $\Ext$ groups, is the long exact sequence of the $\delta$-functor
$$
\Eext^q_A(⋅,N),\ q\geq0
(\resp\ \Ext^q_A(X;⋅,N),\ q\geq0)
$$
associated\pageoriginale with the exact sequence
$0\to M_U\to M\to M_Z\to0$, taking
\eqhyperref[V.6.8.3]{6.8.3} and \eqhyperref[V.6.8.7]{6.8.7} into
account.

\setcounter{subsection}{10}
\subsection{}\etiquette[6.11]{V.subsec6.11}
Let us indicate briefly how these results may be extended to the case
of families of supports. Let $E$ be a topos. For every object $X$ of
$E$, denote by $\Fer(X)$ the set of closed subtopoi of the topos
$E_{/X}$. By passage to complements, this set is in bijective
correspondence with the set of open subtopoi of $E_{/X}$, \ie with the
set of subobjects of $X$ (IV 8). The presheaf
$\Fer:X\longmapsto\Fer(X)$ is in fact a $𝒰$-sheaf, as is immediately
verified, and is therefore representable. In other words, there is an
object of $E$, again denoted by $\Fer$, and a closed subtopos
$Z_{\Fer}$ of $E_{/\Fer}$ such that, for every $X$, every element of
$\Fer(X)$ is obtained from $Z_{\Fer}$ by base change along a morphism
$u_Z:X\to\Fer$ uniquely determined by $Z$.

\begin{enonce*}{Definition 6.12}\etiquette[6.12]{V.def6.12}
A family of supports of $E$ is a subset $\Phi$ of the set of closed
subtopoi of $E$ having the following properties:
\begin{itemize}
\item[{\rm (S1)}] The union of a finite family of elements of $\Phi$
belongs to $\Phi$.
\item[{\rm (S2)}] Every closed subtopos of $E$ contained in an element
of $\Phi$ is an element of $\Phi$.
\end{itemize}
\end{enonce*}

\setcounter{subsection}{12}
\setcounter{subsubsection}{0}
\subsubsection{}\etiquette[6.12.1]{V.subsubsec6.12.1}
Let $E$ be a topos, let $\Phi$ be a family of supports of $E$, and let
$X$ be an object of $E$. Denote by $\Phi(X)$ the smallest family of
supports of $E_{/X}$ containing the closed subtopoi of $E_{/X}$
obtained from closed subtopoi in $\Phi$ by the base change $X\to e$
($e$ being the final object of $E$). The functor
$X\longmapsto\Phi(X)$ is a subpresheaf of the sheaf $\Fer$. It is
therefore separated.

\subsubsection{}\etiquette[6.12.2]{V.subsubsec6.12.2}
A family $\Phi$ of supports of $E$ is said to be \emph{of local
character} if it has the following property:

\noindent
(LC) For every epimorphic family $(X_i\to e)$, $i\in I$ (where $e$ is
the final object of $E$), the sequence of sets
$$
\Phi\to\prod_i\Phi(X_i)\rightrightarrows
\prod_{i,j}\Phi(X_i\times X_j)
$$
is exact.

Let $\vieuxsheaf{a\Phi}$ be the sheaf associated with the presheaf
$X\longmapsto\Phi(X)$ (\hyperref[V.subsubsec6.12.1]{6.12.1}).
Condition (LC) is equivalent to the condition that the canonical
morphism $\Phi\to\vieuxsheaf{a\Phi}(e)$ be\pageoriginale a bijection
(\cf the construction of the associated sheaf in II in the case of a
separated presheaf). To verify (LC), one may therefore restrict to a
final family of epimorphic families $(X_i\to e)$, $i\in I$.

\begin{enonce*}[remark]{Example 6.12.3}\etiquette[6.12.3]{V.exem6.12.3}
\begin{enumerate}
\renewcommand{\labelenumi}{\theenumi)}
\hypertarget{sga4.exp.V.example.6.12.3.item.1}{}\item Let $Z$ be a closed subtopos of $E$. The set of closed subtopoi
of $E$ contained in $Z$ is a family of supports of $E$. It is of local
character.

\item Let $T$ be a topological space and let $\Phi$ be a
paracompactifying family of supports of $T$ \cite{sga4.exp.V.bib.7}. In general, the
family $\Phi$ is not of local character.

\item Let $T$ be a topological space and let $p$ be an integer. The
family $\Phi_p$ of closed subsets of $T$ of Krull codimension
$\geq p$ is a family of supports of $T$. It is of local character.

\item Let $E$ be a topos with the following property: every epimorphic
family $(X_i\to e)$, $i\in I$, is dominated by a finite epimorphic
family. Then every family of supports of $E$ is of local character.
Indeed, using \hyperlink{sga4.exp.V.example.6.12.3.item.1}{example 1}), it is enough to show that a filtered
inductive limit $\Phi_\lambda$ of families of local character is a
family of local character. This follows immediately by taking the
inductive limit in the sequence of sets
$$
\Phi_\lambda\to\prod_i\Phi_\lambda(X_i)
\rightrightarrows\prod_{i,j}\Phi_\lambda(X_i\times X_j){\quoi},
$$
where $(X_i\to e)$, $i\in I$, is a finite epimorphic family. Indeed,
by \hyperref[V.subsubsec6.12.2]{6.12.2} and the hypothesis on $E$, it
is enough to consider finite epimorphic families in verifying
condition (LC); and, since filtered inductive limits commute with
finite projective limits (I 2), exactness of these sequences of sets
is preserved upon passage to the inductive limit.

\item Let $E$ be a coherent topos (\hyperrefext[VI][VI.def2.3]{2.3}).
Then, by the preceding example, every family of supports of $E$ is of
local character.
\end{enumerate}
\end{enonce*}

\setcounter{subsection}{12}
\subsection{}\etiquette[6.13]{V.subsec6.13}
Let $(E,A)$ be a ringed topos, let $\Phi$ be a family of supports of
$E$, let $N$ be an $A$-module (a left module, to fix ideas), and let
$X$ be an object of $E$. Set
\hypertarget{sga4.exp.V.equation.6.13.1}{}\begin{align*}
ℋ^{0}_{\Phi}(N)&=\varinjlim_{Z\in\Phi}ℋ^{0}_Z(N){\quoi},\tag{6.13.1}\\
H^{0}_{\Phi}(E,N)&=\varinjlim_{Z\in\Phi}H^{0}_Z(E,N){\quoi}.
\tag{6.13.2}
\end{align*}\hypertarget{sga4.exp.V.equation.6.13.2}{}
Let\pageoriginale $M$ be a left $A$-module, and set
\hypertarget{sga4.exp.V.equation.6.13.3}{}\begin{gather*}
\Eext^{0}_{A,\Phi}(M,N)=\Hhom_{A,\Phi}(M,N)
=\varinjlim_{Z\in\Phi}\Eext^{0}_{A,Z}(M,N){\quoi},\tag{6.13.3}\\
\Ext^{0}_{A,\Phi}(X;M,N)=
\Hom_{A|X,\Phi(X)}(M|X,N|X)
=\varinjlim_{Z\in\Phi}\Ext^{0}_{A,Z}(X;M,N){\quoi}.
\tag{6.13.4}
\end{gather*}\hypertarget{sga4.exp.V.equation.6.13.4}{}
There are canonical isomorphisms
\begin{equation*}%
\Eext^{0}_{A,\Phi}(M,N)\simeq
ℋ^{0}_{\Phi}(\Eext^{0}_A(M,N)){\quoi},\tag{6.13.5}
\end{equation*}\etiquette[6.13.5]{V.6.13.5}

\begin{equation*}%
\Ext^{0}_{A,\Phi}(X;M,N)\simeq
H^{0}_{\Phi}(X,\Eext^{0}_A(M,N)){\quoi}.\tag{6.13.6}
\end{equation*}\etiquette[6.13.6]{V.6.13.6}

When $\Phi$ is the family of closed subtopoi contained in a fixed
closed subtopos $Z$, there are isomorphisms
\begin{equation*}
\begin{cases}
ℋ^{0}_{\Phi}\simeqℋ^{0}_Z,\\
H^{0}_{\Phi}\simeq H^{0}_Z,\\
\Eext^{0}_{A,\Phi}(⋅,⋅)\simeq\Eext^{0}_{A,Z}(⋅,⋅),\\
\Ext^{0}_{A,\Phi}(X;⋅,⋅)\simeq\Ext^{0}_{A,Z}(X;⋅,⋅){\quoi}.
\end{cases}\tag{6.13.7}
\end{equation*}\etiquette[6.13.7]{V.6.13.7}
The inductive limits defining the preceding functors are filtered.
Consequently, these functors are left exact. Their right derived
functors are denoted by
\begin{equation}
ℋ^q_{\Phi},\ H^q_{\Phi},\ \Eext^q_{A,\Phi}(⋅,⋅),\
\Ext^q_{A,\Phi}(X;⋅,⋅){\quoi}.\tag{6.13.8}
\end{equation}\etiquette[6.13.8]{V.6.13.8}
They too are computed by inductive limits of the functors $ℋ^q_Z$,
$H^q_Z$, \etcn as $Z$ ranges through $\Phi$.

From the isomorphisms \eqhyperref[V.6.13.5]{6.13.5} and \hyperref[V.6.13.6]{(6.13.6)}, one
obtains two spectral sequences by taking inductive limits in the first
spectral sequence in \eqhyperref[V.6.9.1]{6.9.1} and the first spectral
sequence in \eqhyperref[V.6.9.2]{6.9.2}:
\hypertarget{sga4.exp.V.equation.6.13.9}{}\begin{gather*}
\Eext^{p+q}_{A,\Phi}(M,N)\Leftarrow E^{p,q}_2=
ℋ^p_{\Phi}(\Eext^q_A(M,N)){\quoi},\tag{6.13.9}\\
\Ext^{p+q}_{A,\Phi}(X;M,N)\Leftarrow E^{p,q}_2=
H^p_{\Phi}(X,\Eext^q_A(M,N)){\quoi}.\tag{6.13.10}
\end{gather*}\hypertarget{sga4.exp.V.equation.6.13.10}{}

\subsection{}\etiquette[6.14]{V.subsec6.14}\pageoriginale
With the notation of \hyperref[V.subsec6.13]{6.13}, and without any
hypothesis on $\Phi$, there is a functorial morphism
\begin{equation*}
\theta:H^{0}_{\Phi}(N)\to H^{0}(E,ℋ^{0}_{\Phi}(N)){\quoi}.
\tag{6.14.1}
\end{equation*}\etiquette[6.14.1]{V.6.14.1}
This morphism is always injective, but is not in general an
isomorphism. However, if $\Phi$ is \emph{of local character}, the
morphism $\theta$ is an isomorphism, as is immediately verified.
Likewise, if $\Phi$ is \emph{of local character}, there is an
isomorphism
\begin{equation}
\theta':\Ext^{0}_{A,\Phi}(E;M,N)\simeq
H^{0}(E,\Eext^{0}_{A,\Phi}(M,N)){\quoi}.\tag{6.14.2}
\end{equation}\etiquette[6.14.2]{V.6.14.2}

\begin{enonce*}{Proposition 6.15}\etiquette[6.15]{V.prop6.15}
Let $(E,A)$ be a ringed topos such that $E$ is coherent \emph{(VI)},
let $\Phi$ be a family of supports of $E$, and let $M$ and $N$ be two
$A$-modules. There are two spectral sequences:
\hypertarget{sga4.exp.V.equation.6.15.1}{}\begin{align*}%\etiquette[6.15.1]{V.6.15.1}
H^{p+q}_{\Phi}(E,N)\Leftarrow E^{p,q}_2
&=H^p(E,ℋ^q_{\Phi}(N)){\quoi},\tag{6.15.1}\\
\Ext^{p+q}_{A,\Phi}(E;M,N)\Leftarrow E^{p,q}_2
&=H^p(E,\Eext^q_{A,\Phi}(M,N)){\quoi}.\tag{6.15.2}
\end{align*}\hypertarget{sga4.exp.V.equation.6.15.2}{}
\end{enonce*}

These spectral sequences are deduced from the spectral sequence
\eqhyperref[V.6.4.2]{6.4.2} and the second spectral sequence in
\eqhyperref[V.6.9.2]{6.9.2} by taking the inductive limit over the
closed subtopoi $Z$ in $\Phi$, taking into account that the cohomology
of a coherent topos commutes with inductive limits of sheaves (IV 5).

\setcounter{subsection}{15}
\subsection{}\etiquette[6.16]{V.subsec6.16}
We mention, without proof, that the second spectral sequence in
\eqhyperref[V.6.9.1]{6.9.1} and the third spectral sequence in
\eqhyperref[V.6.9.2]{6.9.2} (with $X$ equal to the final object of
$E$) can be extended to the case of families of supports, provided
that the topos $E$ is coherent and the module $M$ is perfect
\cite{sga4.exp.V.bib.13}.

Finally, one can also generalize the notion of families of supports
by introducing presheaves of families of supports, sheaves of families
of supports, and the corresponding cohomology groups\pageoriginale
and sheaves\footnote{\cf \cite{sga4.exp.V.bib.12}, Chap.~IV, \S 1 (Lecture Notes 20,
Springer), for the development of these notions in the form of a fugue
with variations.}.

% END INCLUDED SGA4_English_current_buildable_source_20260906/source/05/drafts/SGA4_V_LINES_2301_2629_English.texfrag


% BEGIN INCLUDED SGA4_English_current_buildable_source_20260906/source/05/drafts/SGA4_V_LINES_2630_2810_English.texfrag
\section{Appendix: Čech cohomology}\etiquette[7]{V.sec7}

Developing an idea due to P. Cartier, we show in this section how,
in an arbitrary topos, one can compute the cohomology of a sheaf by
means of coverings. For the classical theory of paracompact spaces,
we refer to \cite{sga4.exp.V.bib.7}; for another method which applies to certain
topoi, and in particular to the étale topos, see \cite{sga4.exp.V.bib.14}.

\subsection{Skeleton and coskeleton}\etiquette[7.1]{V.subsec7.1}

\setcounter{subsubsection}{-1}
\subsubsection{}\etiquette[7.1.0]{V.subsubsec7.1.0}

Let $𝚫$ be the category of standard simplices (the objects of
$𝚫$ are the ordered sets $\Delta_n=[0,\ldots,n]$, and the
morphisms are the increasing maps). Let $E$ be a $𝒰$-topos, and
let $𝚫(E)$ be the category of presheaves on $𝚫$ with values in
$E$; in other words, the category of \emph{semi-simplicial objects}
of $E$. Denote by $𝚫[n]$ the full subcategory of $𝚫$ defined by
the objects $\Delta_p$, $p\leq n$, by
$i_n:𝚫[n]\to𝚫$ the inclusion functor, and by $𝚫 E[n]$ the
category of presheaves on $𝚫[n]$ with values in $E$; in other
words, the category of \emph{semi-simplicial objects of $E$
truncated at order} $n$. By I \hyperref[I.prop5.1]{5.1}, there is a sequence of three
adjoint functors (I \hyperref[I.sec5.3]{5.3})
\[
\xymatrix@C=3cm@M=.5cm{
𝚫 E[n] \ar@<17pt>[r]^{i_{n!}} \ar@<-19pt>[r]^{i_{n\ast}}&
𝚫 E \ar@<1pt>[l]_{i^{\ast}_n}\\
}
\]
where $i_n^\ast$ is the restriction functor to the category
$𝚫[n]$, and $i_{n\ast}$ and $i_{n!}$ are its right and left
adjoints, respectively. Note that $𝚫 E$ and $𝚫 E[n]$ are
$𝒰$-topoi (IV \hyperref[IV.theo1.2]{1.2}), and that $(i_n^\ast,i_{n\ast})$ is a morphism
of topoi (IV \hyperref[IV.def3.1]{3.1}) which is an embedding of $𝚫 E[n]$ into
$𝚫 E$ (I \hyperref[I.prop5.6]{5.6} and IV \hyperref[IV.def9.1.1]{9.1.1}).

\begin{enonce*}{Definition 7.1.1}\etiquette[7.1.1]{V.def7.1.1}
The\pageoriginale functor $i_{n!}i_n^\ast$ (\resp
$i_{n\ast}i_n^\ast$) from
$𝚫_E$ to $𝚫_E$ is denoted by $\sk_n$ (\resp $\cosk_n$) and is
called the \emph{skeleton functor of order $n$} (\resp the
\emph{coskeleton functor of order $n$}).
\end{enonce*}

The skeleton and coskeleton functors are used constantly in the
theory of semi-simplicial sets \cite{sga4.exp.V.bib.3}.

\begin{enonce*}{Proposition 7.1.2}\etiquette[7.1.2]{V.prop7.1.2}
\begin{enumerate}
\renewcommand{\labelenumi}{\theenumi)}
\item The adjunction morphism $\sk_n\to\id$ is a monomorphism. The
canonical morphisms $\sk_n\sk_m\to\sk_n$ (\resp
$\cosk_n\to\cosk_n\cosk_m$) are isomorphisms when $n\leq m$
(\resp $n\geq m$).

\item Let $u:E\to E'$ be a morphism of topoi. The functor $u^\ast$,
extended to semi-simplicial and truncated semi-simplicial objects,
commutes with the functors $i_{n!}$, $i_n^\ast$, $\sk_n$, and
$\cosk_n$.
\end{enumerate}
\end{enonce*}

The first assertion follows immediately from the definitions. To
prove 2), it suffices to observe that, by I \hyperref[I.prop5.1]{5.1}, the functors
$i_{n!},\ldots,\cosk_n$ are computed by means of inductive limits
and finite projective limits.

\subsection{An acyclicity lemma}\etiquette[7.2]{V.subsec7.2}

\setcounter{subsubsection}{-1}
\subsubsection{}\etiquette[7.2.0]{V.subsubsec7.2.0}

Let $M$ be an abelian group of $E$. The \emph{homology} of a
semi-simplicial object $K$ of $E$ with coefficients in $M$ is
defined in the usual way. First consider $A(K)$, the free abelian
semi-simplicial group generated by $K$, and then form the tensor
product $M\otimes_{\ZZ}A(K)$; this gives an abelian semi-simplicial
group of $E$. One then considers the associated complex of abelian
groups (formed by taking the alternating sum of the boundary
operators) and takes its homology. The homology objects are denoted
by $H_i(K,M)$. They are abelian groups of $E$. When $M$ is the group
$\mathbb{Z}_E$ (the free abelian group generated by the final object
of $E$), we write simply $H_i(K)$. The groups $H_i(K)$ are called
the \emph{homology groups} of $K$. A semi-simplicial object $K$ is
said to be \emph{acyclic} if, for every abelian group $M$ of $E$,
the groups $H_i(K,M)$ vanish for $i>0$.

Formation of homology commutes with inverse images under morphisms
of topoi\pageoriginale.

The purpose of this section is to prove the following lemma:

\begin{enonce*}{Lemma 7.2.1}\etiquette[7.2.1]{V.lem7.2.1}
Let $E$ be a topos, let $K_{∙}$ and $L_{∙}$ be two semi-simplicial
objects of $E$, let $v_{∙}:K_{∙}\to L_{∙}$ be a morphism of
semi-simplicial objects, and let $n$ be an integer $\geq0$. Suppose
that $v$ has the following properties:
\begin{enumerate}
\renewcommand{\labelenumi}{\theenumi)}
\item For every integer $p<n$ $(p\geq0)$, the morphism
$v_p:K_p\to L_p$ is an isomorphism.

\item The morphism $K_n\xrightarrow{v_n}L_n$ is an epimorphism.

\item The canonical morphisms
$K_{∙}\to\cosk_n(K_{∙})$ and $L_{∙}\to\cosk_n(L_{∙})$ are
isomorphisms.
\end{enumerate}
Then, for every integer $p$, the morphism $v_p$ is an epimorphism,
and the morphism $v_{∙}$ induces an isomorphism on homology objects.
\end{enonce*}

\begin{enonce*}[remark]{Remark}
The proof in fact shows that, after adopting a suitable definition
of homotopy in topoi, the morphism $v_{∙}$ is a homotopy equivalence.
\end{enonce*}

\begin{enonce*}[remark]{Proof}
We may suppose that $E$ is the topos of sheaves on a small site $C$
(IV 1), and hence that $E$ is a subtopos of a topos $E'$ having
enough fiber functors (for example, the topos $C^{\wedge}$). Denote
by $a^\ast:E'\to E$ the inverse-image functor of the embedding
morphism $E\hookrightarrow E'$. Let
$$
v_{∙}[n]:K_{∙}[n]\to L_{∙}[n]
$$
be the morphism of truncated semi-simplicial objects obtained by
truncating $v_{∙}$ at order $n$. Denote by $L_{∙}[n]'$ the image (in
$E'$) of $K_{∙}[n]$ under $v_{∙}[n]$, and by
$v_{∙}[n]':K_{∙}[n]\to L_{∙}[n]'$ the morphism induced by
$v_{∙}[n]$. Put
$L!=i_{n\ast}L_{∙}[n]'$, $K!=i_{n\ast}K_{∙}[n]$, and
$v!=i_{n\ast}v_{∙}[n]'$ (here $i_{n\ast}$ is relative to $E'$).
By 1), the morphism $v!:K!\to L!$ has properties 1), 2), and 3) of
the lemma. Moreover, by 2), 3), and \hyperref[V.prop7.1.2]{7.1.2},
the morphism $a^\ast v!$ is none other than $v_{∙}$. It therefore
suffices to prove the lemma for $v!$, and hence we may suppose that
$E$ has enough fiber functors. Using these fiber functors, we may
then reduce to the case where $E$ is the category of sets.
\end{enonce*}

First observe that the hypotheses of the lemma on the morphism
$v_{∙}$ are\pageoriginale stable under every base change
$M_{∙}\to L_{∙}$ for which $M_{∙}$ is a coskeleton of order $n$.
Consequently, the fiber $P_{∙}$ of $v_{∙}$ at any base point of
$L_{∙}$ is of the form $i_{n\ast}P_{∙}[n]$, where $P_{∙}[n]$ is a
nonempty complex truncated at order $n$ of the form
\begin{equation*}
P_n\to e\to e\to e\to\cdots\to e{\quoi}.\tag{7.2.2}
\end{equation*}\etiquette[7.2.2]{V.7.2.2}
For every integer $i$, denote by $\delta\sheaf{\Delta}_i$ the
semi-simplicial set forming the boundary of the simplex
$\sheaf{\Delta}_i$. Every morphism from
$\delta\sheaf{\Delta}_i$ to $P_{∙}$ extends to a morphism from
$\sheaf{\Delta}_i$ to $P_{∙}$. Indeed, this property is immediate
for $i\geq n$, because $P_{∙}$ is a coskeleton of order $n$, and it
is clear for $i<n$ from the description \eqhyperref[V.7.2.2]{7.2.2}.
Since $P_{∙}$ is a Kan complex, it is homotopically trivial
\cite{sga4.exp.V.bib.6}. Note that the morphism $v_{∙}$ is a fibration in the sense
of Kan \cite{sga4.exp.V.bib.6}. Since the fibers of this fibration are
homotopically trivial, $v_{∙}$ induces an isomorphism on the homotopy
groups of $K_{∙}$ and $L_{∙}$ at all base points of $K_{∙}$. By the
Hurewicz theorem, this implies that $v_{∙}$ induces an isomorphism on
homology \cite{sga4.exp.V.bib.6}. This proves the second assertion of the lemma. To
prove the first, observe that every morphism from a complex truncated
at order $n$ to $L_{∙}[n]$ lifts to a morphism to $K_{∙}[n]$.
Consequently, by 3), every morphism from a complex to $L_{∙}$ lifts
to a morphism to $K_{∙}$. Thus $v_{∙}$ is surjective.

% END INCLUDED SGA4_English_current_buildable_source_20260906/source/05/drafts/SGA4_V_LINES_2630_2810_English.texfrag


% BEGIN INCLUDED SGA4_English_current_buildable_source_20260906/source/05/drafts/SGA4_V_LINES_2811_3070_English.texfrag
\subsection{Hypercoverings}\etiquette[7.3]{V.subsec7.3}

\setcounter{subsubsection}{-1}
\subsubsection{}\etiquette[7.3.0]{V.subsubsec7.3.0}
Let $C$ be a site belonging to the universe $𝒰$, in which fiber
products and products of two objects are representable. Denote by
$C^{\wedge}$ the topos of $𝒰$-presheaves on $C$, and by
$C^{\sim}$ the topos of $𝒰$-sheaves on $C$. An object $K$ of
$C^{\wedge}$ is called \emph{semi-representable} if it is isomorphic
to a coproduct of representable presheaves.

Let $SR(C)$ be the full subcategory of $C^{\wedge}$ defined by the
semi-representable objects. In the category $SR(C)$, projective
limits indexed by finite nonempty categories are representable, and
the inclusion functor $SR(C) \hookrightarrow C^{\wedge}$ commutes
with these finite projective limits. By a remark already made, it
follows that if $K_∙$ is a semi-simplicial object truncated in degree
$n$ in $C^{\wedge}$, all of whose objects\pageoriginale are
semi-representable, then the extension $i^+_n(K_∙)$ has
the same property. Thus, if $K_∙$ is a semi-simplicial object of
$C^{\wedge}$ all of whose objects are semi-representable, then, for
every $n$, the object $\cosk_n (K_∙)$ has the same property.

\subsubsection{Definitions and notation}\etiquette[7.3.1]{V.subsubsec7.3.1}

\paragraph{}\etiquette[7.3.1.1]{V.par7.3.1.1}

Let $p > 0$ be an integer. A semi-simplicial object $K_∙$ of
$C^{\wedge}$ is called a \emph{hypercovering of type} $p$ of $C$,
or, when no confusion can result, a hypercovering of type $p$, if it
has the following properties:
\begin{itemize}
\item[HR1)] For every integer $n \geq 0$, $K_n$ is semi-representable.

\item[HR2)] The canonical morphism $K_∙ \to \cosk_p (K_∙)$ is an
isomorphism.

\item[HR3)] For every integer $n \geq 0$, the canonical morphism of
presheaves $K_{n + 1} \to (\cosk_n(K_∙))_{n + 1}$ is a covering
morphism of presheaves (II \hyperref[II.rem6.2]{6.2}). The canonical morphism of
presheaves $K_{0} \to e$, where $e$ is the final object of
$C^{\wedge}$, is a covering morphism of presheaves.
\end{itemize}

\paragraph{}\etiquette[7.3.1.2]{V.par7.3.1.2}
A hypercovering of type $p$ is a hypercovering of type $q$ for every
$q \geq p$. A semi-simplicial object $K_∙$ will be called a
\emph{hypercovering} (of $C$) if it has properties HR 1) and HR 3).

\paragraph{}\etiquette[7.3.1.3]{V.par7.3.1.3}
Denote by $HR_p$ (\resp $HR$), and call the category of
hypercoverings of type $p$ (\resp the category of hypercoverings),
the following category:
\begin{enumerate}
\renewcommand{\theenumi}{\alph {enumi}}
\renewcommand{\labelenumi}{\theenumi)}
\item The objects of $HR_p$ (\resp $HR$) are the hypercoverings of
  type $p$ (\resp the hypercoverings).

\item Let $K_∙$ and $L_∙$ be two objects of $HR_p$ (\resp $HR$). A
  morphism of $HR_p$ (\resp $HR$) with source $K_∙$ and target $L_∙$
  is a morphism $v_∙ : K_∙ \to L_∙$ of semi-simplicial objects of
  $C^{\wedge}$.
\end{enumerate}

\paragraph{}\etiquette[7.3.1.4]{V.par7.3.1.4}
Let $C$ be a site in which fiber products are representable, and let
$X$ be an object of $C$. A hypercovering of $X$ (\resp a
hypercovering of type $p$ of $X$) means a hypercovering of the site
$C/X$ (\resp a hypercovering of type $p$ of $C/X$). The category of
hypercoverings of $X$ (\resp of hypercoverings of type $p$ of $X$)
is defined similarly\pageoriginale. There is a sequence of inclusion
functors
$$
\cdots HR_p \hookrightarrow HR_{p + 1} \hookrightarrow \cdots
\hookrightarrow HR{\quoi}.
$$
These functors are fully faithful. We denote by $HR_{\infty}$ the
inductive limit of the categories $HR_p$.

\setcounter{paragraph}{5}
\paragraph{}\etiquette[7.3.1.6]{V.par7.3.1.6}
For every integer $n \geq 0$, denote by $\sheaf{\Delta}_n^c$ the
standard semi-simplicial object of dimension $n$ with values in the
category of sets, \ie the functor on $𝚫$ represented by the
ordered set $[0, n]$. Also denote by $\sheaf{\Delta}_n^c$ the
presheaf on $C$, with values in the category of semi-simplicial sets,
that is constant with value $\sheaf{\Delta}_n$. Thus the presheaf
$\sheaf{\Delta}_n^c$ is a presheaf of semi-simplicial sets. The two
canonical injections from $\Delta_0$ into $\Delta_1$ define two
morphisms of semi-simplicial presheaves:
\[
\xymatrix@C=2cm@M=.25cm{
\underline{\Delta}_0^{c} \ar@<3pt>[r]^{e_{0}} \ar@<-3pt>[r]_{e_{1}} &
\underline{\Delta}_1^c ,
}
\]
and consequently define, for every semi-simplicial presheaf $K_∙$,
two canonical injections:
\[
\xymatrix@C=2cm@M=.25cm{
K_{∙} \ar@<3pt>[r]^(.45){e_{0}} \ar@<-3pt>[r]_(.45){e_{1}}&  K_{∙}
\times \underline{\Delta}_1^c .
}
\]
Two morphisms
$
\xymatrix@C=1cm@M=.25cm{
K_{∙} \ar@<3pt>[r]^(.45){u_{0}} \ar@<-3pt>[r]_(.45){u_{1}}&  L_{∙}
}
$
of semi-simplicial presheaves are called \emph{homotopic morphisms}
if there exists a morphism
$$
v : K_∙ \times \sheaf{\Delta}_1^c \to L_∙
$$
such that the diagrams
\[
\xymatrix@R=.5cm@C=1.25cm{
K_∙ \ar[rr]^{e_i} \ar[rdd]_{u_i}&   &K_{∙}\times \sheaf{\Delta}_1^c
\ar[ldd]^v\\
                                    &   & &\txt{i = 0,1}\\
                       & L_{∙} &
}
\]
are\pageoriginale commutative. The morphism $v$ is called a
\emph{homotopy joining $u_0$ to $u_1$}. The relation ``$u_0$ and
$u_1$ are two homotopic morphisms from $K_∙$ to $L_∙$'' is not, in
general, an equivalence relation. Nevertheless, the equivalence
relation generated by this relation is compatible with composition
of morphisms. This therefore allows us to define the category of
semi-simplicial presheaves up to homotopy by passing to the quotient
by the equivalence relation generated by the homotopy relation.

\paragraph{}\etiquette[7.3.1.7]{V.par7.3.1.7}
In this way we define the categories $\sheaf{HR}_p$,
$\sheaf{HR}_{\infty}$, and $\sheaf{HR}$, the \emph{categories of
hypercoverings up to homotopy}.

\begin{enonce*}{Theorem 7.3.2}\etiquette[7.3.2]{V.thm7.3.2}
Let $C$ be a site satisfying the conditions of
\hyperref[V.subsubsec7.3.0]{7.3.0}.
\begin{enumerate}
\renewcommand{\labelenumi}{\theenumi)}
\item The category $\sheaf{HR}^{0}_p$ (\resp $\sheaf{HR}^{0}$) is
  filtered \emph{(I \hyperref[I.def2.7]{2.7})}.

\item Let $K_∙$ be an object of $HR_p$ (\resp of $HR$), let $n$ be
  an integer such that $0 \leq n \leq p$ (\resp $0 \leq n$), and
let
$$
u : X \to K_n
$$
be a covering morphism of presheaves. There exists an object $L_∙$
of $HR_p$ (\resp of $HR$) and a morphism $f : L_∙ \to K_∙$ such that
the morphism
$$
f_n : L_n \to K_n
$$
factors through $u$.

\item The associated semi-simplicial sheaves \emph{(II \hyperref[II.def3.5]{3.5})} of
hypercoverings are acyclic (\hyperref[V.subsubsec7.2.0]{7.2.0}) in
strictly positive degrees. The $0$-th homology sheaf is isomorphic
to the sheaf associated to the constant presheaf $𝒵$.
\end{enumerate}
\end{enonce*}

\begin{enonce*}[remark]{Proof}
We prove 3). Let ${K}^{\sim}_{∙}$ be the semi-simplicial sheaf
associated to $K_∙$. The ``associated sheaf'' functor commutes with
the functor $\cosk_n$ (\hyperref[V.prop7.1.2]{7.1.2}. 2)). It follows
that the semi-simplicial sheaf ${K}^{\sim}_{∙}$ has the following
properties:
\begin{enumerate}
\renewcommand{\theenumi}{\alph {enumi}}
\renewcommand{\labelenumi}{\theenumi)}
\item The canonical morphism $K_0 \to e$ is an epimorphism of
  sheaves.

\item For every integer $n \geq 0$, the canonical morphism
  $K_{n + 1} \to (\cosk_n(K_∙))_{n + 1}$ is an epimorphism of
  sheaves.
\end{enumerate}
\end{enonce*}

For $n \geq 0$, put ${L_∙}_n = \cosk_n {K}^{\sim}_{∙}$. Then, for
every $n$, there is a sequence of morphisms of semi-simplicial
sheaves
$$
{K_∙}^{\sim} \xrightarrow{u_n}{L_∙}_n\xrightarrow{v_n} {L_∙}_{n - 1}
\cdots \to {L_∙}_{0} \xrightarrow{v_0} e_∙{\quoi},
$$
and\pageoriginale, for every $j$, $0 \leq j \leq n$, $v_j$ has the
properties of \hyperref[V.lem7.2.1]{7.2.1}. Moreover, the morphism
$u_n$ induces an isomorphism on the components of degree $\leq n$.
It then follows from \hyperref[V.lem7.2.1]{7.2.1}, by induction on
$n$, that $K_∙$ is acyclic.

To prove 1) and 2), we introduce some terminology.

\begin{enonce*}{Definition 7.3.3}\etiquette[7.3.3]{V.def7.3.3}
Let $f : X_{∙} \to Y_∙$ be a morphism of semi-simplicial presheaves.
We say that $f$ is \emph{special of type} $p$ (\resp
\emph{special}) if:
\begin{enumerate}
\renewcommand{\labelenumi}{\theenumi)}
\item The objects $X_{∙}$ and $Y_∙$ are canonically isomorphic to
  their coskeleta of order $p$, and $\cosk_p (f) = f$ (\resp no
  conditions are imposed on $f$, \sisi{$X^{∙}$}{$X_{∙}$}, or
  \sisi{$Y^{∙}$}{$Y_{∙}$}).

\item For every integer $n$ such that $0 \leq n \leq p$ (\resp
  $0 \leq n$), the morphism $\Phi_{n + 1}$ appearing in the diagram
  below is covering:
\[
\xymatrix@C=1cm@R=1.25cm{
X_{n+1} \ar[rrr]^{f_{n+1}} \ar[rd]^{\Phi_{n+1}} \ar[dd]       & & & Y_{n+1}\ar[dd]\\
                                    &P_{n+1} \ar[ld] \ar[rru]& &\\
(\cosk_n X_{∙})_{n+1}\ar[rrr]_{(\cosk_n f)_{n+1}}& & &
(\cosk_n Y_{∙})_{n+1} .
}
\]
(The vertical arrows are defined by the canonical morphisms
$X_{∙} \to \cosk_n X_{∙}$ and $Y_∙ \to \cosk_n Y_∙$. The object
$P_{n+1}$ is the fiber product, and $\Phi_{n+1}$ is the unique arrow
making the diagram commutative.)

\item The morphism $f_0 : X_0 \to Y_0$ is covering.
\end{enumerate}

A semi-simplicial presheaf $K_∙$ is called special of type $p$
(\resp special) if the canonical morphism
$$
K_∙ \to e_{∙}\quad (e_{∙}\quad \mbox{\rm is the final semi-simplicial object})
$$
is special of type $p$ (\resp special), \ie if $K_∙$ satisfies
conditions \emph{HR 2)} and \emph{HR 3)} (\resp \emph{HR 3)}) of
\hyperref[V.par7.3.1.1]{7.3.1.1}.
\end{enonce*}

\begin{enonce*}{Lemma 7.3.4}\etiquette[7.3.4]{V.lem7.3.4}
\begin{enumerate}
\renewcommand{\labelenumi}{\theenumi)}
\item The\pageoriginale composite of two special morphisms (\resp
  special morphisms of type $p$) is a special morphism (\resp a
  special morphism of type $p$).

\item Let $K_∙$ be a special semi-simplicial presheaf (\resp special
  of type $p$), let $X_{∙}\xrightarrow{f}Y_∙$ be a special morphism
  (\resp special of type $p$), and let $u : K_∙ \to Y_∙$ be a
  morphism of complexes. Then the fiber product $P_{∙} =
  K_∙ ×_{Y_∙} X_∙$ is special (\resp special of type $p$).
\end{enumerate}
\end{enonce*}

The proof of this lemma is left to the reader.

\setcounter{subsubsection}{4}

% END INCLUDED SGA4_English_current_buildable_source_20260906/source/05/drafts/SGA4_V_LINES_2811_3070_English.texfrag


% BEGIN INCLUDED SGA4_English_current_buildable_source_20260906/source/05/drafts/SGA4_V_LINES_3071_3430_English.texfrag
\subsubsection{}\etiquette[7.3.5]{V.subsubsec7.3.5}
\emph{Proof of assertion} 2) \emph{of} \hyperref[V.thm7.3.2]{7.3.2}.
We may first suppose that $X$ is semi-representable. For every
semi-simplicial object $L_∙$, denote by $\Hom_{(X)}(L_∙,K_∙)$ the
set of morphisms of semi-simplicial objects equipped with a
factorization $L_n\to X\xrightarrow{u}K_n$. The functor
$\Hom_{(X)}(⋅,K_∙)$ is representable. Indeed, this functor carries
inductive limits to projective limits, and the category
$𝚫(C^{\wedge})$ is a topos. Denote by $P_{∙}$ an object
representing this functor. For an object $Z$ of $C^{\wedge}$,
likewise denote by $j_n^{+}(Z)$ the object representing the functor
$L_∙\mapsto\Hom(L_n,Z)$ on $𝚫 C^{\wedge}$, whose construction by
$\varprojlim$ is made explicit in III \hyperref[III.def1.1]{1.1}. Observe here that this
construction involves only finite projective limits. It follows
immediately that $j_n^{+}$ carries semi-representable presheaves to
semi-simplicial presheaves having semi-representable components, and
covering morphisms $Z'\to Z$ to special morphisms
(\hyperref[V.def7.3.3]{7.3.3}).

By the definition of $P_{∙}$, there is a cartesian square:
\[
\xymatrix@C=2cm@R=1.5cm{
P_{∙} \ar[r]\ar[d]& j_n^{+}(X) \ar[d]^{j_n^{+}(u)}\\
K_{∙} \ar[r]      &  j_n^{+}(K_n)\quad.
}
\]
By what has just been observed, the semi-simplicial objects
$j_n^{+}(X)$ and $j_n^{+}(K_n)$ are semi-representable, and the
morphism $j_n^{+}(u)$ is a special morphism of type $n$. The result
now follows by applying \hyperref[V.cor3.4]{3.4} and the sorites
\hyperref[V.subsubsec7.3.0]{7.3.0}.

\setcounter{subsubsection}{6}
\setcounter{paragraph}{-1}
\paragraph{}\etiquette[7.3.6.0]{V.par7.3.6.0}
Let\pageoriginale $M_∙$ and $N_∙$ be two semi-simplicial
presheaves. The functor
$$
L_∙\longmapsto\Hom(L_∙\times M_∙,N_∙)
$$
is representable. The semi-simplicial presheaf representing it will
be denoted by
$$
\Hhom_∙(M_∙,N_∙){\quoi}.
$$
The degree-$n$ component presheaf of this object will be denoted by
$\Hhom_n(M_∙,N_∙)$. There is a canonical isomorphism
$$
\Hhom_n(M_∙,N_∙)\xrightarrow{\sim}
\Hhom_0(M_∙\times\sheaf{\Delta}^c_n,N_∙){\quoi}.
$$

\begin{enonce*}{Lemma 7.3.6}\etiquette[7.3.6]{V.lem7.3.6}
Let
\[
\xymatrix@C=2cm@R=1.5cm{
M_{∙} \ar@{^{(}->}[r]& N_{∙}\\
\sk_n\sheaf{\Delta}_{n+1}^c \ar[u] \ar@{^{(}->}[r]^{u}&
\sheaf{\Delta}_{n+1}^c \ar[u] &
\kern -1.5cm \txt{\rm ($u$ the canonical injection)}
}
\]
be a cocartesian diagram of semi-simplicial presheaves.

For every hypercovering $L_∙$, the morphism
$$
\Hhom_0(N_∙,L_∙)\to\Hhom_0(M_∙,L_∙)
$$
is covering.
\end{enonce*}

\begin{enonce*}[remark]{Proof}
There is a cartesian diagram:
\[
\xymatrix{
\Hhom_0(N_{∙},L_{∙})\ar[d] \ar[r]&
\Hhom_0(M_{∙},L_{∙}) \ar[d]\\
\Hhom_0(\sheaf{\Delta}^c_{n+1},L_{∙})\ar[r]^{(\ast)} &
\Hhom_0(\sk_n\sheaf{\Delta}^c_{n+1},L_{∙}).
}
\]
It therefore suffices to show that the morphism $(\ast)$ is
covering. But $(\ast)$ is isomorphic to the morphism
$$
L_{n+1}\to(\cosk_nL_∙)_{n+1},
$$
which is covering by hypothesis. \cqfd
\end{enonce*}

\subsubsection{}\etiquette[7.3.7]{V.subsubsec7.3.7}
Proof of assertion 1) of \hyperref[V.thm7.3.2]{7.3.2}. The product
of two objects of $HR_p$ (\resp of $HR$) is again an object of
$HR_p$ (\resp of $HR$) (\hyperref[V.lem7.3.4]{7.3.4}). To prove
that the category $\sheaf{HR}^{0}_p$ (\resp $\sheaf{HR}^{0}$) is
filtered, it suffices to show that, given two morphisms in $HR_p$
(\resp $HR$)
\[
\xymatrix@C=2cm@M=.2cm{
K_{∙} \ar@<3pt>[r]^{u_0} \ar@<-3pt>[r]_{u_1}&L_∙
}
\]
there exists a morphism $v:M_∙\to K_∙$ in $HR_p$ (\resp in $HR$)
such that the morphisms
\[
\xymatrix@C=2cm@M=.2cm{
M_{∙} \ar@<3pt>[r]^{u_0v} \ar@<-3pt>[r]_{u_1v}&L_∙
}
\]
are homotopic,\pageoriginale \ie such that there exists
$w:M_∙\times\sheaf{\Delta}_1^c\to L_∙$ making the diagrams
\[
\vcenter{
\xymatrix@C=2cm@R=1.5cm{
M_{∙} \ar[d]_{v} \ar[r]^(.4){e_i}&
M_{∙}\times\underline{\Delta}^c_1 \ar[d]^{w}\\
K_{∙} \ar[r]^{u_i}&L_{∙}&\kern -1.5cm\txt{$i=0,1$.}
}}\leqno({\rm x})
\]
commute. For every semi-simplicial object $M_∙$ of $C^{\wedge}$
(not necessarily a hypercovering), let $F(M_∙)$ be the set of pairs
$(v,w)$, where $v:M_∙\to K_∙$ and
$w:M_∙\times\sheaf{\Delta}_1^c\to L_∙$, such that the diagrams (x)
commute. The functor $M_∙\mapsto F(M_∙)$ is contravariant in $M_∙$,
carries inductive limits to projective limits, and is consequently
represented by a semi-simplicial object $F_{∙}$. The object $F_{∙}$
is clearly the upper-left vertex of a cartesian diagram
\[
\xymatrix@=2cm{
F_{∙}\ar[r]\ar[d]&
\Hhom_{∙}(\sheaf{\Delta}^{c}_{1},L_{∙})\ar[d]^{\pi}\\
K_{∙}\ar[r]_{u_0,u_1}&L_{∙}\times L_{∙},
}
\]
where $\pi$ is defined by the two inclusions of
$\sheaf{\Delta}^c_0$ into $\sheaf{\Delta}^c_1$. To prove the
assertion, it therefore suffices to show that $F_{∙}$ is an object
of $HR_p$ (\resp of $HR$). For this, by
\hyperref[V.lem7.3.4]{7.3.4} 2) and the sorites
\hyperref[V.subsubsec7.3.0]{7.3.0}, it suffices to show that

\begin{enumerate}
\renewcommand{\theenumi}{\alph{enumi}}
\renewcommand{\labelenumi}{\theenumi)}
\item $\Hhom_∙(\sheaf{\Delta}^c_1,L_∙)$\pageoriginale is
semi-representable,

\item the morphism $\pi$ is special of type $p$ (\resp special).
\end{enumerate}

First, it follows immediately that
$\Hhom_∙(\sheaf{\Delta}^c_1,L_∙)$ is a semi-representable
semi-simplicial object. Indeed, its components are computed by
finite projective limits from the $L_n$, and are therefore
semi-representable. Let us now verify b). First one shows that, when
$L_∙$ is special of type $p$, the morphism
$$
\Hhom_∙(\sheaf{\Delta}_1^c,L_∙)\to
\cosk_p\Hhom_∙(\sheaf{\Delta}^c_1,L_∙)
$$
is an isomorphism; this verifies property 1) of
\hyperref[V.def7.3.3]{7.3.3}. Next, the morphism
$$
\pi_0:\Hhom_0(\sheaf{\Delta}^c_1,L_∙)\to L_0\times L_0
$$
is isomorphic to the canonical morphism
$$
L_1\xrightarrow{(d_0,d_1)}L_0\times L_0,
$$
which is covering by hypothesis ($L_∙$ is a hypercovering).

It remains to verify property 2) of \hyperref[V.def7.3.3]{7.3.3},
\ie to verify that, for every $n$, in the diagram
\[
\xymatrix@C=.75cm@R=1.25cm{
\Hhom_{n+1}(\sheaf{\Delta}_1^c,L_∙)
\ar[rr]^{\pi_{n+1}} \ar[dd] \ar[rd]_{\Phi_{n+1}}&&
L_{n+1}\times L_{n+1}\ar[dd]\\
&P_{n+1}\ar[ru]\ar[ld]&\\
(\cosk_n\Hhom_∙(\sheaf{\Delta}_1^c,L_∙))_{n+1}
\ar[rr]^{(\cosk_n\pi)_{n+1}}&&
(\cosk_nL_∙)_{n+1}\times(\cosk_nL_∙)_{n+1}
}
\]
($P_{n+1}$ is the fiber product), the morphism $\Phi_{n+1}$ is
covering. A simple ``adjoint-functor chase'' shows that
\begin{align*}
P_{n+1}&\xrightarrow{\sim}
\Hhom_0\bigl(\sk_{n+1}(\sheaf{\Delta}_{n+1}^c\times
\sheaf{\Delta}_1^c),L_∙\bigr),\\
\Hhom_{n+1}(\sheaf{\Delta}_1^c,L_∙)&\xrightarrow{\sim}
\Hhom_0(\sheaf{\Delta}_{n+1}^c\times
\sheaf{\Delta}_1^c,L_∙),
\end{align*}
and\pageoriginale that the morphism $\Phi_{n+1}$ is isomorphic to
the morphism
$$
\Hhom_0(\sheaf{\Delta}_{n+1}^c\times\sheaf{\Delta}_1^c,L_∙)
\to
\Hhom_0\bigl(\sk_{n+1}(\sheaf{\Delta}_{n+1}^c\times
\sheaf{\Delta}_1^c),L_∙\bigr)
$$
arising from the injection
$$
\sk_{n+1}(\sheaf{\Delta}_{n+1}^c\times\sheaf{\Delta}_1^c)
\hookrightarrow
\sheaf{\Delta}_{n+1}^c\times\sheaf{\Delta}_1^c{\quoi}.
$$
There is a sequence of subobjects of
$\sheaf{\Delta}_{n+1}^c\times\sheaf{\Delta}_1^c$:
$$
\sk_{n+1}(\sheaf{\Delta}_{n+1}^c\times\sheaf{\Delta}_1^c)
=M_0\hookrightarrow M_1\hookrightarrow\cdots\hookrightarrow M_k
=\sheaf{\Delta}_{n+1}^c\times\sheaf{\Delta}_1^c{\quoi},
$$
such that, for every $0\leq i<k$, the morphism
$$
M_i\hookrightarrow M_{i+1}
$$
fits into a cocartesian diagram:
\[
\xymatrix@C=1.5cm@R=1.5cm{
M_i \ar@{^{(}->}[r]&M_{i+1}&\\
\sk_{n+1}\underline{\Delta}{}^c_{n+2}
\ar@{^{(}->}[r]\ar[u]&
\underline{\Delta}{}^c_{n+2}\ar[u]\;.
}
\]
(The nondegenerate simplices of dimension $n+2$ in
$\sheaf{\Delta}_{n+1}\times\sheaf{\Delta}_1$ are added one after
another.) By \hyperref[V.lem7.3.6]{7.3.6}, the morphism
$\Phi_{n+1}$ is a composite of covering morphisms and is therefore
itself covering. This completes the proof of
\hyperref[V.thm7.3.2]{7.3.2}.

\subsection{The isomorphism theorem}\etiquette[7.4]{V.subsec7.4}

\setcounter{subsubsection}{-1}
\subsubsection{}\etiquette[7.4.0]{V.subsubsec7.4.0}

Let $F$ be a presheaf of abelian groups on a site $C$ satisfying
condition \hyperref[V.subsubsec7.3.0]{7.3.0}. For every
hypercovering $K_∙$, put
$$
H^q(K_∙,F)=H^q(\Hom_{C^{\wedge}}(K_∙,F)){\quoi}.
$$
The $H^q(K_∙,F)$ form a $\delta$-functor on the category of abelian
presheaves on $C$, but they are not, in general, the derived
functors of the functor $H^0(K_∙,F)$.

Now\pageoriginale put
\begin{equation*}
\Hvr{q}(C,F)=\varinjlim_{HR_r}H^q(⋅,F)
\text{ }(r\mbox{ may be infinite}){\quoi},\tag{7.4.0.1}
\end{equation*}\etiquette[7.4.0.1]{V.7.4.0.1}
and
\begin{equation*}
\Hvv{q}(C,F)=\varinjlim_{HR}H^q(⋅,F){\quoi}.\tag{7.4.0.2}
\end{equation*}\etiquette[7.4.0.2]{V.7.4.0.2}
The $\Hvr{q}(C,F)$ (\resp $\Hvv{q}(C,F)$) again form a
$\delta$-functor, because the category $\sheaf{HR}^0_r$ (\resp
$\sheaf{HR}^0$) is filtered (\hyperref[V.thm7.3.2]{7.3.2}). Since
the categories $\sheaf{HR}_r$ are not necessarily $𝒰$-categories,
these $\delta$-functors take values in the category of abelian
$𝒱$-groups, for a suitable universe $𝒱$.

Suppose now that the presheaf $F$ is a sheaf. Since the
semi-simplicial sheaf associated with a hypercovering is acyclic
(\hyperref[V.thm7.3.2]{7.3.2}), there is a spectral sequence,
functorial in $F$ and in $K_∙$:
\begin{equation*}
H^p(K_∙,\Hh^q(F))\Rightarrow
H^{p+q}(C^{\sim},F){\quoi}.\tag{7.4.0.3}
\end{equation*}\etiquette[7.4.0.3]{V.7.4.0.3}
Passing to the limit gives the spectral sequences
\begin{equation*}
\begin{split}
\Hvr{p}(C,\Hh^q(F))&\Rightarrow H^{p+q}(C^{\sim},F){\quoi},\\
\Hvv{p}(C,\Hh^q(F))&\Rightarrow H^{p+q}(C^{\sim},F){\quoi}.
\end{split}\tag{7.4.0.4}
\end{equation*}\etiquette[7.4.0.4]{V.7.4.0.4}

\begin{enonce*}{Theorem 7.4.1}\etiquette[7.4.1]{V.thm7.4.1}
Let $C$ be a site satisfying condition
\hyperref[V.subsubsec7.3.0]{7.3.0}, and let $F$ be an abelian sheaf
on $C$.
\begin{enumerate}
\renewcommand{\labelenumi}{\theenumi)}
\item The spectral sequences \eqhyperref[V.7.4.0.4]{7.4.0.4} define
isomorphisms
$$
\Hvr{q}(C,F)\xrightarrow{\sim}H^q(C^{\sim},F),\quad q\leq r+1,
$$
and a monomorphism
$$
\Hvr{r+2}(C,F)\longrightarrow H^{r+2}(C^{\sim},F){\quoi}.
$$

\item There is\pageoriginale an isomorphism of $\delta$-functors:
$$
\Hvi{q}(C,F)\xrightarrow{\sim}\Hvv{q}(C,F)
\xrightarrow{\sim}H^q(C^{\sim},F)\quad(\text{for every }q){\quoi}.
$$

\item Let $G$ be a presheaf of abelian groups, and let $F$ be its
associated sheaf. There are functorial isomorphisms
$$
\Hvr{q}(C,G)\xrightarrow{\sim}H^q(C^{\sim},F),\quad q\leq r-1,
$$
and a monomorphism
$$
\Hvr{r}(C,G)\longrightarrow H^r(C^{\sim},F){\quoi}.
$$
When $G$ is a separated presheaf, the latter morphism is an
isomorphism, and there is a monomorphism
$$
\Hvr{r+1}(C,G)\to H^{r+1}(C^{\sim},F){\quoi}.
$$

\item There are isomorphisms
$$
\Hvi{q}(C,G)\xrightarrow{\sim}\Hvv{q}(C,G)
\xrightarrow{\sim}H^q(C^{\sim},F)\quad(\text{for every }q){\quoi}.
$$
\end{enumerate}
\end{enonce*}

\begin{enonce*}[remark]{Proof}
It suffices to show that, if $N$ is an abelian presheaf whose
associated sheaf is zero, then
$\Hvr{q}(C,N)=0$ for $q\leq r$ (\resp
$\Hvv{q}(C,N)=0$ for every $q$); this follows immediately by using
\hyperref[V.thm3.2]{3.2}.
\end{enonce*}

\begin{enonce*}[remark]{Remark 7.4.2}\etiquette[7.4.2]{V.rem7.4.2}
\begin{enumerate}
\renewcommand{\labelenumi}{\theenumi)}
\item Note that, for sites satisfying condition
\hyperref[V.subsubsec7.3.0]{7.3.0}, hypercoverings of type $0$ are
ordinary coverings, and the functors $\Hvr{q}$ for $r=0$ are none
other than the functors $\Hv{q}$ introduced in
\eqhyperref[V.2.4.5.1]{2.4.5.1}.

\item Let $C$ be a site in which finite projective limits are
representable, and suppose that for every object $X$ of $C$ and
every covering family $(X_i\to X)$, $i\in I$, there exists an
$i_0\in I$ such that $X_{i_0}\to X$ is covering. One can then show
that hypercoverings of type $p$ whose components are representable
are cofinal in $\sheaf{HR}_p$. Likewise, hypercoverings $K_∙$ for
which there exists an integer $p$ such that $K_∙$ is of type $p$,
and whose components are representable, are cofinal
in\pageoriginale $\sheaf{HR}$. (This may be proved by following the
argument of \hyperref[V.cor3.5]{3.5}.) Thus, in the case under
consideration, these hypercoverings with representable components
suffice to compute the cohomology of the sheaves associated with
presheaves.
\end{enumerate}
\end{enonce*}


% END INCLUDED SGA4_English_current_buildable_source_20260906/source/05/drafts/SGA4_V_LINES_3071_3430_English.texfrag


% BEGIN INCLUDED SGA4_English_current_buildable_source_20260906/source/05/drafts/SGA4_V_LINES_3431_3814_English.texfrag
\section{Appendix. Local direct limits (by P. Deligne)}\etiquette[8]{V.sec8}

The editor recommends that the reader should, as a rule, avoid
reading this appendix. It presents a technique which sometimes
makes it possible to extend to topoi without enough points
assertions which the existence of points makes trivial. This
technique yields a sheaf-theoretic variant of D. Lazard's theorem
asserting that flat modules over a ring are direct limits of finite
free modules.

\subsection{Locally filtering categories}\etiquette[8.1]{V.subsec8.1}

\setcounter{subsubsection}{-1}
\subsubsection{}\etiquette[8.1.0]{V.subsubsec8.1.0}
If $\Bb$ is a category and $p:\Aa\to\Bb$ is a category over
$\Bb$, we shall use the following notation:
\begin{itemize}
\item For $U\in\Ob\Bb$, $\Aa_U$ is the fiber category $p^{-1}(U)$.

\item For $f:U\to V$ in $\Bb$ and $\lambda,\mu\in\Ob\Aa$ such
that $p(\lambda)=U$ and $p(\mu)=V$, put
$$
\Hom_f(\lambda,\mu)=p^{-1}(f),\quad\text{where }p:\Hom(\lambda,\mu)
\to\Hom(U,V){\quoi}.
$$
\end{itemize}

\begin{enonce*}{Definition 8.1.1}\etiquette[8.1.1]{V.def8.1.1}
Let $\Ss$ be a site. A \emph{locally cofiltering category (or
locally left filtering category)} over $\Ss$ is a category
$p:\Ll\to\Ss$ over $\Ss$ such that:
\begin{itemize}
\item[($\Ll0$)] For every $f:U\to V$ in $\Ss$ and every
$\mu\in\Ob\Ll_V$, there exists $\lambda\in\Ob\Ll_U$ such that
$\Hom_f(\lambda,\mu)\neq\varnothing$;

\item[($\Ll1$)] For every $U\in\Ob\Ss$ and every finite family
$(\lambda_i)$ of objects of $\Ll_U$, there exist a covering
$f_j:U_j\to U$ of $U$ and objects $\mu_j\in\Ob\Ll_{U_j}$ such that
$\Hom_{f_j}(\mu_j,\lambda_i)\neq\varnothing$ for all $i$ and $j$;

\item[($\Ll2$)] For every $f:U\to V$ in $\Ss$ and every pair of
parallel arrows $(\phi_0,\phi_1):\lambda\rightrightarrows\mu$ over
$f$, there exist a covering $f_j:U_j\to U$ of $U$ and arrows
$\psi_j$ with target $\lambda$ such that $p(\psi_j)=f_j$ and
$\phi_0\psi_j=\phi_1\psi_j$.
\end{itemize}
\end{enonce*}

\setcounter{subsubsection}{1}
\paragraph{}\etiquette[8.1.1.1]{V.par8.1.1.1}
This definition simplifies when $\Ll$ is fibered over $\Ss^{\ast}$:
the axiom
\begin{itemize}
\item[$(\Ll_\ast0)$] is then satisfied, and $(\Ll1)$, $(\Ll2)$ may
be stated as follows:

\item[$(\Ll'1)$] For every $U\in\Ob\Ss$ and every finite family
$\lambda_i$ of objects of $\Ll_U$, locally on $U$ there exists a
$\mu$ mapping to all the $\lambda_i$.

\item[$(\Ll'2)$] For every $U\in\Ob\Ss$, every pair of parallel
arrows $(\phi_0,\phi_1):\lambda\rightrightarrows\mu$ in $\Ll_U$ can,
locally on $U$, be equalized by an arrow with target $\lambda$.
\end{itemize}

\subsubsection{}\etiquette[8.1.2]{V.subsubsec8.1.2}
Let $\Ss$ be a site, let $\Ll$ be a category over $\Ss$, and let
$\mathsf{e}$ be a stack over $\Ss$. Denote by
$\sheaf{\Gamma}\mathsf{e}$ the category of global sections
$\Hhom^{\cart}_{\Ss}(\Ss,\mathsf{e})$ of $\mathsf{e}$. If $\Ss$
has a final object $S$, then $\sheaf{\Gamma}\mathsf{e}$ ``is none
other than'' $\mathsf{e}_S$.

\begin{enonce*}{Definition 8.1.2.1}\etiquette[8.1.2.1]{V.def8.1.2.1}
The category of \emph{local inverse systems of objects of
$\mathsf{e}$ indexed by $\Ll$} is the category
$\Hhom_{\Ss}(\Ll,\mathsf{e})$.
\end{enonce*}

Denote by $c$ (``associated constant local inverse system'') the
composite functor
$$
\sheaf{\Gamma}\mathsf{e}=\Hhom^{\cart}_{\Ss}(\Ss,\mathsf{e})
\hookrightarrow\Hhom_{\Ss}(\Ss,\mathsf{e})
\xrightarrow{u\mapsto u\circ p}\Hhom_{\Ss}(\Ll,\mathsf{e}).
$$
When the dependence on $\Ll$ must be made explicit, we shall write
$c_{\Ll}$ instead.

\begin{enonce*}{Definition 8.1.3}\etiquette[8.1.3]{V.def8.1.3}
The local inverse-limit functor, denoted by
$\varprojlim_{\lambda\in\Ll}$, is the partially defined functor
right adjoint to the functor $c$.
\end{enonce*}

Thus the functor
$$
\varprojlim:\Hhom_{\Ss}(\Ll,\mathsf{e})\to
\sheaf{\Gamma}\mathsf{e}
$$
satisfies
$$
\Hom(X,\varprojlim_{\lambda\in\Ll}X_\lambda)\simeq
\Hom(c(X),(X_\lambda)_{\lambda\in\Ll}){\quoi}.
$$

\begin{enonce*}{Definition 8.1.4}\etiquette[8.1.4]{V.def8.1.4}
An $\Ss$-functor $F:\Ll\to\Mm$ between locally cofiltering
categories over $\Ss$ is said to be cofinal if it satisfies the
following two conditions:
\begin{itemize}
\item[$(C\mathfrak{F}1)$] For every\pageoriginale $U\in\Ob\Ss$ and
$\mu\in\Ob\Mm_U$, there exist a covering $f_j:U_j\to U$ of $U$ and
objects $\lambda_j\in\Ob\Ll_{U_j}$ such that
$\Hom_{f_j}(F(\lambda_j),\mu)\neq\varnothing$.

\item[$(C\mathfrak{F}2)$] For every $f:U\to V$ in $\Ss$ and every
pair of parallel arrows $(\phi_0,\phi_1):F(\lambda)
\rightrightarrows\mu$ over $f$, there exist a covering
$f_j:U_j\to U$ of $U$ and arrows $\psi_j$ with target $\lambda$ in
$\Ll$, such that $p(\psi_j)=f_j$ and
$\phi_0F(\psi_j)=\phi_1F(\psi_j)$.
\end{itemize}
\end{enonce*}

The composite of two cofinal functors is cofinal; equivalences of
locally filtering categories are cofinal (the proof is left to the
reader).

\begin{enonce*}{Proposition 8.1.5}\etiquette[8.1.5]{V.prop8.1.5}
Let $F:\Ll\to\Mm$ be a cofinal functor between locally cofiltering
categories over $\Ss$, let $\mathsf{e}$ be a stack over $\Ss$, and
let $(X_\mu)_{\mu\in\Mm}$ be a local inverse system indexed by
$\Mm$. Then the morphism, functorial in $X\in\sheaf{\Gamma}\mathsf{e}$,
\[
F^{\ast}:\Hom(c_{\Mm}(X),(X_\mu)_{\mu\in\Mm})\longrightarrow
\Hom(c_{\Ll}(X),(X_{F(\lambda)})_{\lambda\in\Ll})
\]
is an isomorphism, so that
$$
F^{\ast}:\varprojlim_{\Mm}X_\mu\longrightarrow
\varprojlim_{\Ll}X_{F(\lambda)}
$$
is an isomorphism, the two sides being simultaneously defined or
undefined.
\end{enonce*}

The proof uses the following lemma, which is left to the reader:

\begin{enonce*}{Lemma 8.1.6}\etiquette[8.1.6]{V.lem8.1.6}
If $F:\Ll\to\Mm$ is cofinal, then for every $f:U\to V$ in $\Ss$
and arrows over $f$, $\phi_i:F(\lambda_i)\to\mu$ $(i=1,2)$, there
exist a covering $f_j:U_j\to U$ of $U$, objects
$\lambda'_j\in\Ob\Ll_{U_j}$, and commutative diagrams
\[
\xymatrix@=1.25cm{
              & F(\lambda_1)   \ar[rd]^{\phi_1}&\\
F(\lambda'_j) \ar@{--}[ru] \ar@{--}[rd]&                & \mu\\
              & F(\lambda_2) \ar[ru]_{\phi_2}  &
}
\]
$$
U_j\xrightarrow{f_j}U\to V
$$
whose projection in $\Ss$ is $(f_j,f)$.
\end{enonce*}

For every $U\in\Ob\Ss$ and $\mu\in\Mm_U$, condition
$(C\mathfrak{F}1)$ gives a covering $f_j:U_j\to U$ of $U$ and
arrows $\phi_j:F(\lambda_j)\to\mu$ over the $f_j$. If
$$
\psi=(\psi_\lambda)\in\Hom(c_{\Ll}X,
(X_{F(\lambda)})_{\lambda\in\Ll}):\quad
\psi_\lambda:X|p(\lambda)\to X_{F(\lambda)}
$$
is the image of
$$
\phi=(\phi_\mu)\in\Hom(c_{\Mm}X,(X_\mu)_{\mu\in\Mm}):\quad
\phi_\mu:X|p(\mu)\to X_\mu,
$$
then the diagrams
\[
\xymatrix@=1.5cm{
                               & X_{F(\lambda_j)}
                               \ar[ld]_{\psi_{\lambda_j}} \ar[rd]^{\phi_j}&        \\
X|U_j \ar[rr]_{\phi_\mu|U_j} &                  & X_\mu|U_j
}
\]
commute. Thus the $\phi_\mu$ are locally determined by $\psi$; since
$\mathsf{e}$ is a stack, it follows that $F^{\ast}$ is injective.
To prove that $F^{\ast}$ is surjective, start with $\psi$ and
construct a $\phi$ such that $\psi=F^{\ast}\phi$. The preceding
construction gives the composite arrows
$$
\phi_{\mu,j}=\phi_j\circ\psi_{\lambda_j}:X|U_j\to X_\mu|U_j{\quoi}.
$$
Let us verify that these arrows glue, so that there exists
$\phi_\mu:X|U\to X_\mu$ such that
$\phi_\mu|U_j=\phi_{\mu,j}$. Consider a commutative diagram in
\[
\xymatrix@=1cm{
                 & U_{ij}\ar[ld]_{g_i} \ar[rd]^{g_j}     &            \\
U_i  \ar[rd]_{f_i}            &            & U_j \ar[ld]^{f_j}\\
                 & U     &
}
\]
To\pageoriginale verify that
$\phi_{\mu,j}|U_{ij}=\phi_{\mu,i}|U_{ij}$, it is enough to do so
locally on $U_{ij}$; by \hyperref[V.lem8.1.6]{8.1.6} and $(\Ll0)$ this reduces us to the
case where there is a commutative diagram
\[
\xymatrix@=1.5cm{
             & F(\lambda_{ij}) \ar[ld]_{F(a_{ij})} \ar[rd]^{F(a_{ji})} &\\
F(\lambda_i) \ar[rd]_{\phi_i}&                 & F(\lambda_j) \ar[ld]^{\phi_j}\\
             & \mu             &
}
\]
over the preceding one. The diagram
\[
\xymatrix@R=1.5cm@C=2cm{
  & X_{F(\lambda_{ij})} \ar[ld] \ar[rd] \ar@/_1.68pc/[dd]&\\
X_{F(\lambda_i)}\ar[rd] & X|U_{ij}\ar[u]\ar[l]\ar[r] & X_{F(\lambda_j)} \ar[ld]\\
                 & X_\mu &
}
\]
is then commutative, and there exists a family $(\phi_\mu)$ such
that, for every arrow $\phi:F(\lambda)\to\mu$ over
$f:V\to U$, the diagram
\[
\xymatrix@=1.5cm{
                              & X_{F(\lambda)} \ar[rd]^{\phi}      &    \\
X|V \ar[rr]_{\phi_\mu|V} \ar[ru]^{\psi_\lambda} &
                              & X_\mu|V
}
\]
commutes. If $\sigma:\mu\to\mu'$ is an arrow of $\Mm$, then
$(\sigma|V)(\phi_\mu|V)=(\sigma|V)\phi\cdot\psi_\lambda
=(\sigma\phi|V)\cdot\psi_\lambda=\phi_{\mu'}|V$; hence
$\phi=(\phi_\mu)$ is a morphism of functors such that
$\psi=F^{\ast}\phi$, as required.

If\pageoriginale $E$ is an ordered set, we shall also denote by $E$
the category whose set of objects is $E$ and in which $\Hom(i,j)$
is a singleton or empty according as $i\leq j$ or not.

\begin{enonce*}{Proposition 8.1.7}\etiquette[8.1.7]{V.prop8.1.7}
Let $\Ll$ be a locally cofiltering category over a site $\Ss$.
There exist an ordered set $E$ and a functor $F$ from $E$ to $\Ll$
such that:
\begin{enumerate}
\renewcommand{\theenumi}{\roman{enumi}}
\renewcommand{\labelenumi}{(\theenumi)}
\item $E$, regarded as a category over $\Ss$ by means of
$pF:E\to\Ss$, is locally cofiltering.
\item The functor $F$ is cofinal.
\end{enumerate}
\end{enonce*}

Let $\Ll'$ be the locally cofiltering category which is the product
of $\Ll$ with the category defined by the ordered set $\ZZ$. The
projection from $\Ll'$ to $\Ll$ is cofinal, so it suffices to prove
\hyperref[V.prop8.1.7]{8.1.7} for $\Ll'$.

If $X$ is a finite set of arrows in $\Ll'$, denote by $X^0$ the set
of endpoints of the arrows in $X$. Let $E$ be the set of nonempty
finite subsets $X$ of $\mathrm{Fl}(\Ll')$ for which there exists
$\lambda_X\in\Ob\Ll'$ satisfying:
\begin{enumerate}
\renewcommand{\theenumi}{\alph{enumi}}
\renewcommand{\labelenumi}{\theenumi)}
\item no arrow of $X$, except $1_{\lambda_X}$, has target
$\lambda_X$;

\item if $\mu\in X^0$, then $1_\mu\in X$ and
$\Hom(\lambda_X,\mu)\cap X\neq\varnothing$;

\item every diagram of the form
\[
\xymatrix@=1.75cm{
\lambda_X \ar[r] \ar[rd]& \mu \ar[d]\\
            & \mu'
}
\]
consisting of arrows of $X$ is commutative.
\end{enumerate}

Order $E$ by the relation opposite to inclusion. If $X\in E$ and
$\mu\in X^0$, there is a unique arrow from $\lambda_X$ to $\mu$
which belongs to $X$. If $X,Y\in E$ and $X\supset Y$, let
$\lambda_{X,Y}$ be the unique arrow in $X$ from $\lambda_X$ to
$\lambda_Y$. The assignments $X\mapsto\lambda_X$ and
$(X\supset Y)\mapsto\lambda_{X,Y}$ form a functor from $E$ to
$\Ll'$.

\begin{enonce*}{Lemma 8.1.8}\etiquette[8.1.8]{V.lem8.1.8}
The category\pageoriginale $E$ is locally cofiltering over $\Ss$,
and the functor $\lambda:E\to\Ll'$ is cofinal.
\end{enonce*}

We must verify the following axioms successively:
\begin{itemize}
\item[(i) \emph{axiom} ($\Ll0$)] Let $f:V\to U$ and $X\in E_U$.
There exist $\lambda\in\Ob\Ll'_V$ and
$\phi\in\Hom_f(\lambda,\lambda_X)$. Moreover, $\lambda$ may be
chosen so that $\lambda\notin X^0$ (this is where
$\Ll'=\Ll\times\ZZ$ is used). Let
$X'=X\cup\{1_\lambda\}\cup Y$, where $Y$ is the set of composite
arrows
$\lambda\xrightarrow{\phi}\lambda_X\xrightarrow{\psi}\mu$
$(\psi\in X)$. Then $X'\in E$, $\lambda_{X'}=\lambda$, and
$\Hom(X',X)\neq\varnothing$.

\item[(ii) \emph{axiom} ($\Ll1$)] Let $(X_i)$ be a finite family of
objects of $E_U$. There exist a covering $f_j:U_j\to U$, objects
$\lambda_j\in\Ob\Ll'_{U_j}$, and arrows
$\phi_{ji}\in\Hom_{f_j}(\lambda_j,\lambda_{X_i})$. If
$x\in X_a^0\cap X_b^0$, then for each $j$ there exist a covering
$g_{jk}:U_{jk}\to U_j$ and arrows
$\psi_{jk}:\lambda_{jk}\to\lambda_j$ such that
$p(\psi_{jk})=g_{jk}$ and which equalize the pair of arrows
$\lambda_j\rightrightarrows x$, one belonging to $X_a$ and the
other to $X_b$. We may arrange that the $\lambda_{jk}$ belong to
none of the $X_i^0$. Replace the $(f_j,\lambda_j,\phi_{ji})$ by the
$(f_jg_{jk},\lambda_{jk},\phi_{ji}\psi_{jk})$, and repeat this
construction for every $x$ in an intersection of two of the $X_i$.
We obtain a new system $(f_j,\lambda_j,\phi_{ji})$; put
$$
X_j=\bigcup_iX_i\cup\{1_{\lambda_j}\}\cup Y_j,
$$
where $Y_j$ is the set of composite arrows
$$
\lambda_j\xrightarrow{\phi_{ji}}\lambda_{X_i}
\xrightarrow{\psi}x\quad(\psi\in X_i){\quoi}.
$$
Then $X_j\in E_{U_j}$, and these $X_j$ satisfy $(\Ll1)$.

\item[(iii) \emph{axiom} ($\Ll2$)] Trivial, since there are no
nontrivial parallel arrows!

\item[(iv) \emph{axiom} $(C\mathfrak{F}1)$] Trivial, since the
functor $\lambda$ is surjective.

\item[(v) \emph{axiom} $(C\mathfrak{F}2)$] Let $f:U\to V$ and let
$(\phi_0,\phi_1):\lambda_X\rightrightarrows\lambda$ be a pair of
parallel arrows over $f$. There exist a covering $f_j:U_j\to U$ of
$U$ and arrows $\psi_j:\lambda_j\to\lambda_X$ such that
$p(\psi_j)=f_j$, $\phi_0\psi_j=\phi_1\psi_j$, and
$\lambda_j\notin X^0$. Let
$X_j=X\cup\{1_{\lambda_j}\}\cup Y_j$, where $Y_j$ is the set of
composite arrows
$$
\lambda_j\xrightarrow[]{\psi_j}\lambda_X\xrightarrow{\psi}\mu
\quad(\psi\in X){\quoi}.
$$
Then\pageoriginale $X_j\in E_{U_j}$, and
$F(\lambda_{X_j,X})=\psi_j$ equalizes $(\phi_0,\phi_1)$.
\end{itemize}

\setcounter{subsubsection}{9}
\setcounter{paragraph}{-1}
\paragraph{}\etiquette[8.1.9.0]{V.par8.1.9.0}

Suppose that $\Ss$ is (the site of nonempty open subsets of) the
one-point topological space. A locally cofiltering category over
$\Ss$ is then simply a left filtering category (\ie one for which
$\Ss^0$ is filtering in the sense of I \hyperref[I.def2.7]{2.7}), and cofinal functors
correspond to cofinal functors between left filtering categories
(I 8). Proposition \hyperref[V.prop8.1.7]{8.1.7} may then be restated
as follows.

\begin{enonce*}{Proposition 8.1.9}\etiquette[8.1.9]{V.prop8.1.9}
For every cofiltering category $\Ll$, there exist a cofiltering
ordered set $E$ and a cofinal functor from $E$ to $\Ll$.
\end{enonce*}

In this special case, the proof of \hyperref[V.prop8.1.7]{8.1.7}
reduces essentially to the proof of \hyperref[V.prop8.1.9]{8.1.9}
given in I 8.

% END INCLUDED SGA4_English_current_buildable_source_20260906/source/05/drafts/SGA4_V_LINES_3431_3814_English.texfrag


% BEGIN INCLUDED SGA4_English_current_buildable_source_20260906/source/05/drafts/SGA4_V_LINES_3815_4065_English.texfrag
\subsubsection{}\etiquette[8.1.10]{V.subsubsec8.1.10}
Let $q:\Tt\to\Ss$ be a morphism of sites, and let
$p:\Ll\to\Ss$ be a locally cofiltered category over $\Ss$. Define
the category $\Ll^{\ast}$ as follows:
\begin{enumerate}
\renewcommand{\theenumi}{\roman{enumi}}
\renewcommand{\labelenumi}{(\theenumi)}
\item An object of $\Ll^{\ast}$ is a quadruple $(V,U,\phi,\lambda)$
such that $V\in\Ob\Tt$, $U\in\Ob\Ss$,
$\phi\in\Hom(V,q^{\ast}U)$, and $\lambda\in\Ob\Ll_U$.

\item An arrow
$f:(V,U,\phi,\lambda)\to(V',U',\phi',\lambda')$ is a triple
$(f_1,f_2,f_3)$, where $f_1\in\Hom(V,V')$,
$f_2\in\Hom(U,U')$, and $f_3\in\Hom(\lambda,\lambda')$, such that
$q^{\ast}f_2\circ\phi=\phi'\circ f_1$ and $p(f_3)=f_2$.

The functor $f\mapsto f_1$ makes $\Ll^{\ast}$ a category over
$\Tt$.
\end{enumerate}

\begin{enonce*}{Lemma 8.1.11}\etiquette[8.1.11]{V.lem8.1.11}
The category $\Ll^{\ast}$ is locally cofiltered over $\Tt$.
\end{enonce*}

Axiom $(\Ll0)$ is trivial. Let us verify $(\Ll1)$.

Let $L_i=(V,U_i,\phi_i,\lambda_i)$ be a finite family. The $\phi_i$
define a morphism from $V$ to the sheaf
$q^{\ast}a\prod U_i=a\prod q^{\ast}U_i$. Consequently, there are
commutative diagrams
\[
\xymatrix@C=2.5cm@R=1.5cm{
V_j \ar[r]^{f_j}\ar[d]^{\psi_j}&V\ar[d]^{\phi_i}\\
q^{\ast}U_j\ar[r]_{q^{\ast}(p_{ij})}&q^{\ast}U_i
}
\]
such\pageoriginale that the $f_j$ cover $V$.

Choose $\lambda_{ji}$ such that
$\Hom_{p_{ij}}(\lambda_{ji},\lambda_i)\neq\varnothing$. By
$(\Ll1)$, after refining $U_j$ (and $V_j$), we may take
$\lambda_{ji}=\lambda_j$ independently of $i$, and the
$L_j=(V_j,U_j,\psi_j,\lambda_j)$ satisfy $(\Ll1)$.

Let us verify $(\Ll2)$. Let
$(f_1,f_2):L^1\rightrightarrows L^2$, with
$L^i=(V^i,U^i,\phi^i,\lambda^i)$ and
$f_i=(f,g_i,h_i)$. The $\phi^i$ define a morphism from $V^1$ to the
sheaf
$q^{\ast}a\Ker(U^1\rightrightarrows U^2)
=a\Ker(q^{\ast}U^1\rightrightarrows q^{\ast}U^2)$, so there are
commutative diagrams
\[
\xymatrix@C=2cm@R=1.5cm{
V_j \ar[r]^{f_j}\ar@{-}[d]^{\psi_j}&
V^1\ar[r]\ar@{-}[d]^{\phi^1}&V^2\ar@{-}[d]^{\phi^2}\\
q^{\ast}U_j\ar[r]_{q^{\ast}p_j}&q^{\ast}U^1
\ar@<2pt>[r]^{q^{\ast}g_1}\ar@<-2pt>[r]_{q^{\ast}g_2}&
q^{\ast}U^2&\kern-1cm(g_1p_j=g_2p_j),
}
\]
such that the $f_j$ cover $V^1$. Let
$\chi_j:\lambda_j\to\lambda^1$ satisfy $p(\chi_j)=p_j$. Applying
$(\Ll2)$ to $(h_1\chi_j,h_2\chi_j)$, and refining $U_j$ (and
$V_j$) if necessary, we may arrange that
$h_1\chi_j=h_2\chi_j$; the
$L_j=(V_j,U_j,\psi_j,\lambda_j)$ then satisfy $(\Ll2)$.

\begin{enonce*}{Proposition 8.1.12}\etiquette[8.1.12]{V.prop8.1.12}
Let $t^{\ast}:\Ss^{\ast}\to\Ss$ be a locally cofiltered category
over $\Ss$, and endow $\Ss^{\ast}$ with the topology induced from
that of $\Ss$ by $t^{\ast}$. Then $t^{\ast}$ is an equivalence of
sites $t:\Ss\to\Ss^{\ast}$.
\end{enonce*}

By construction, the functor
$t^\ast:\Ff\mapsto\Ff\circ t^{\ast}$ carries sheaves on $\Ss$ to
sheaves on $\Ss^{\ast}$. By $(\Ll0)$ and $(\Ll1)$ applied to the
empty family, every ``sufficiently small'' object of $\Ss$ belongs
to the image of $t^{\ast}$ (that is, the image of $t^{\ast}$ is a
\emph{covering sieve} in $\Ss$). Hence, by the ``comparison lemma''
(III 4), the inclusion into $\Ss$ of the full subcategory defined by
$t^{\ast}(\Ob\Ss^{\ast})$ is an equivalence of sites. We may
therefore suppose that $t^{\ast}$ is surjective on objects.

Let $\Ff$ be a sheaf on $\Ss^{\ast}$.

\begin{enumerate}
\renewcommand{\theenumi}{\roman{enumi}}
\renewcommand{\labelenumi}{(\theenumi)}
\item Let $(f_j:U_j\to U)$ be a covering in $\Ss$, let
$\lambda^1,\lambda^2\in\Ob\Ss^{\ast}_U$, and let
$f_j^i\in\Hom_{f_j}(\mu_j,\lambda^i)$ for $i=1,2$. Then
$\Ff(\lambda^1)$ and $\Ff(\lambda^2)$ have the same image in
$\displaystyle\mathop{\prod}_j\Ff(\mu_j)$.

Indeed, let\pageoriginale $x=(x_j)$ belong to the image of
$\Ff(\lambda^1)$. For $x$ to belong to the image of
$\Ff(\lambda^2)$, it suffices that, for every commutative diagram
\[
\xymatrix@C=2.5cm@R=1.5cm{
\mu\ar[r]\ar[d]&\mu_j\ar[d]\\
\mu_k\ar[r]&\lambda^2,
}
\]
$x_j$ and $x_k$ have the same image in $\Ff(\mu)$. Since $\Ff$ is
a sheaf, it suffices to verify this after replacing $\mu$ by the
objects of one of its coverings. By $(\Ll2)$, we may then suppose
that the diagrams analogous to the preceding one, with $\lambda^2$
replaced by $\lambda^1$, also commute, in which case the assertion is
clear.

\item Let $f:U\to V$ be an arrow of $\Ss$, let
$\lambda\in\Ob\Ss^{\ast}_U$, and let
$\mu\in\Ob\Ss^{\ast}_V$. There then exist a covering
$f_j:U_j\to U$, an object $\mu'\in\Ob\Ss^{\ast}_U$, and arrows in
the diagram
\[
\xymatrix@C=2.5cm@R=1cm{
&\lambda&\\
\lambda_j\ar[ru]^{\phi_j}\ar[rd]_{\psi_j}&&\mu\\
&\mu'\ar[ru]_{\psi}&
}
\]
lying over
$$
U_j\xrightarrow{f_j}U\xrightarrow{f}V{\quoi}.
$$
By (i), there is a unique arrow $\overline f$ making the following
diagram commute:
\[
\xymatrix@C=2.5cm@R=1.5cm{
\Ff(\mu)\ar[r]^{\overline f}\ar[d]_{\psi}&
\Ff(\lambda)\ar[d]^{\phi_j}\\
\Ff(\mu')\ar[r]_{\psi_j}&\Ff(\lambda_j).
}
\]
The arrow $\overline f$ does not change if $(f_j)$ is replaced by a
finer covering $f_k:U_k\to U_{j(k)}$, if $\lambda_j$ is replaced by
$\lambda_k$ equipped with
$\phi_k\in\Hom_{f_k}(\lambda_k,\lambda_{j(k)})$, and if
$\phi_j$ (respectively $\psi_j$) is replaced by
$\phi_{j(k)}\circ\phi_k$ (respectively
$\psi_{j(k)}\circ\phi_k$).\pageoriginale

The reader may deduce that $\overline f$ depends only on $f$.

\item Let us prove that $\overline f$ is functorial in $f$. Thus let
$$
W\xrightarrow{h}V\xrightarrow{g}U
$$
be arrows in $\Ss$, and let $\lambda^0$, $\mu^0$, and $\nu^0$
belong respectively to $\Ob\Ss^{\ast}_U$, $\Ob\Ss^{\ast}_V$, and
$\Ob\Ss^{\ast}_W$.
There exists a commutative diagram
\[
\xymatrix@C=.5cm@R=.5cm{
\nu^{0}& &&&&\mu^{0}&&\lambda^{0}\\
&&\nu^{1}\ar[rrru]&&&&\mu^{1}\ar[ru]&\\
&&&{\nu^{2}}\ar[rrru]&&&&\\
&&&&&&&\\
&\nu_{jk}\ar[luuuu]^(.7){e^{0}_{jk}}
\ar[uuur]^(.7){e^{1}_{jk}}\ar[rruu]_{e^{2}_{jk}}\ar[rrrr]&&&&
\mu_j\ar[uuur]_{f^{1}_j}\ar[uuuu]^(.3){f^{0}_j}&&
}
\]
such that

\begin{enumerate}
\renewcommand{\theenumii}{\alph{enumii}}
\renewcommand{\labelenumii}{\theenumii)}
\item $t^{\ast}(\nu^i)$, $t^{\ast}(\mu^i)$,
$t^{\ast}(e^i_{jk})$, and $t^{\ast}(f^i_j)$ do not depend on $i$;

\item for fixed $i$, the $e^i_{jk}$ cover $W$, and the $f^i_j$
cover $V$.
\end{enumerate}

First construct $\mu^1$, $\nu^1$, and $\nu^2$ by $(\Ll0)$, and
the $\mu_j$ as in (ii). Let $(W'_{jk})$ be a covering of $W$ mapping
to the covering $V_j=t^{\ast}(\mu_j)$ as required by the diagram.
After refining $(W'_{jk})$ and applying $(\Ll0)$ and $(\Ll1)$, one
obtains a diagram of the required type, satisfying a) and b), but
not necessarily commutative. Refining once more and applying
$(\Ll2)$ makes it commutative.

It remains only to contemplate the following commutative diagram:
\[
\xymatrix@C=1.5cm@R=.25cm{
\Ff(\lambda^{0})\ar[r]^{\bar g}&\Ff(\mu^0)\ar[rd]\ar[rr]^{\bar h}&&
\Ff(\nu^{0})\ar[dddd]\\
&&\Ff(\nu^1)\ar[rddd]&\\
\Ff(\mu^{1})\ar[r]&\prod\Ff(\mu_j)\ar[rrdd]&&&\\
&&&\\
\Ff(\nu^2)\ar[rrr]&&&\prod\Ff(\nu_{jk})
}
\]

\item It\pageoriginale is clear that $\overline f$ is functorial
in $\Ff$. Thus every sheaf $\Ff$ on $\Ss^{\ast}$ is the direct
image under $t$ of a uniquely determined presheaf $\Ff_0$ on
$\Ss$, and every morphism of sheaves $f:\Ff\to\mathfrak g$ is the
image of a unique morphism of presheaves
$f_0:\Ff_0\to\mathfrak g_0$.
\end{enumerate}

It already follows that $t_{\ast}$ is fully faithful; it remains to
prove that $\Ff_0$ is a sheaf. If $(f_j:U_j\to U)$ is a covering,
then this covering is the image of a covering in $\Ss^{\ast}$, so
$\Ff_0(U)$ injects into $\prod_j\Ff_0(U_j)$. If elements
$x_j\in\Ff_0(U_j)$ glue, then a fortiori they glue in $\Ss^{\ast}$,
and hence come from an element of $\Ff_0(U)$. This completes the
proof of \hyperref[V.prop8.1.12]{8.1.12}.

\setcounter{subsubsection}{13}
\setcounter{paragraph}{-1}
\paragraph{}\etiquette[8.1.13.0]{V.par8.1.13.0}

Let $q:\Tt\to\Ss$ and $p:\Ll\to\Ss$ be as in
(\hyperref[V.subsubsec8.1.10]{8.1.10}), let $\Ll^{\ast}$ be the
category defined there, and let $\mathsf e$ be a stack on $\Tt$.
The stack $q_{\ast}\mathsf e$ on $\Ss$ is the stack defined by
$(q_{\ast}\mathsf e)_U=\mathsf e_{q^{\ast}U}$.

If $(X_{\lambda})_{\lambda\in\Ll}$ is a local projective system
indexed by $\Ll$ and taking values in $q_{\ast}\mathsf e$, it
defines a local projective system
$(X_L)_{L\in\Ll^{\ast}}$ taking values in $\mathsf e$ by the formula
$$
X_L=\phi^{\ast}X_{\lambda}quad\text{for }
L=(V,U,\phi,\lambda){\quoi}.
$$

The reader may verify the following:

\begin{enonce*}{Proposition 8.1.13.1}\etiquette[8.1.13.1]{V.prop8.1.13.1}
With the preceding notation, one has
$$
\varprojlim_{\lambda\in\Ll}X_{\lambda}
=\varprojlim_{L\in\Ll^{\ast}}X_L{\quoi},
$$
the two sides being simultaneously defined or undefined.
\end{enonce*}

Let $\Ss$ be a site, and let $p^{\ast}:\Ss^{\ast}\to\Ss$ be a
locally filtered category over $\Ss$. If $\mathsf e$ is a stack on
$\Ss$, then $p_{\ast}\mathsf e$ is a stack on $\Ss^{\ast}$
(\hyperref[V.par8.1.13.0]{8.1.13.0}). Moreover, the identity functor
$\Ss^{\ast}\to\Ss^{\ast}$ is a locally filtered category over
$\Ss^{\ast}$, endowed with the induced topology, and every local
inductive system $(X_{\lambda})$ indexed by $\Ss^{\ast}$ defines a
local inductive system over $\Ss^{\ast}$ indexed by $\Ss^{\ast}$.

One verifies easily:

\begin{enonce*}{Proposition 8.1.14}\etiquette[8.1.14]{V.prop8.1.14}
With\pageoriginale the preceding notation, one has
$$
\varprojlim_{\lambda\in\Ss^{\ast}/\Ss}X_{\lambda}
=\varprojlim_{\lambda\in\Ss^{\ast}/\Ss^{\ast}}X_{\lambda}{\quoi},
$$
the two sides being simultaneously defined or undefined.
\end{enonce*}

% END INCLUDED SGA4_English_current_buildable_source_20260906/source/05/drafts/SGA4_V_LINES_3815_4065_English.texfrag


% BEGIN INCLUDED SGA4_English_current_buildable_source_20260906/source/05/drafts/SGA4_V_LINES_4066_4380_English.texfrag
\subsection{Local inductive limits in categories of
sheaves}\etiquette[8.2]{V.subsec8.2}

\setcounter{subsubsection}{-1}
\subsubsection{}\etiquette[8.2.0]{V.subsubsec8.2.0}

Let $p:\Tt\to\Ss$ be a morphism of sites. Let $\mathsf{e}$ be the
following stack on $\Ss$:
\begin{enumerate}
\renewcommand{\theenumi}{\roman{enumi}}
\renewcommand{\labelenumi}{(\theenumi)}
\item An object of $\mathsf{e}$ is a pair $(U,\Ff)$ consisting of an
object $U$ of $\Ss$ and a sheaf $\Ff$ on $p^{\ast}U$.

\item An arrow $f:(U,\Ff)\to(V,\Gg)$ is a pair consisting of an arrow
$f_0:U\to V$ and a morphism of sheaves on $p^{\ast}U$:
$$
f_1:p^{\ast}(f_0)^{\ast}\Gg\to\Ff.
$$

\item The functor from $\mathsf{e}$ to $\Ss$ is given by
$f\longmapsto f_0$.
\end{enumerate}

Care must be taken that $\mathsf{e}_U$ is the \emph{opposite category}
of the category of sheaves on $p^{\ast}U$.

Let $\Ll$ be a locally cofiltered category (on the left) over $\Ss$.

\begin{enonce*}{Definition 8.2.1}\etiquette[8.2.1]{V.def8.2.1}
With the preceding notation, the category of local inductive systems,
indexed by $\Ll$, of sheaves on $\Tt$ is the category
$\Hom_{\Ss}(\Ll,\mathsf{e})^{\circ}$.
\end{enonce*}

The ``projective limit'' functor of \hyperref[V.def8.1.3]{8.1.3} is
called here the \emph{``local inductive limit'' functor}. It is a
functor from $\Hom_{\Ss}(\Ll,\mathsf{e})^{\circ}$ to
$(\Gamma\mathsf{e})^{\circ}=(\Tt)^{\sim}$.

\begin{enonce*}{Theorem 8.2.2}\etiquette[8.2.2]{V.thm8.2.2}
The ``local inductive limit'' functor from the category of inductive
systems of sheaves on $\Tt$, indexed by $\Ll$, to the category of
sheaves on $\Tt$ is defined everywhere, commutes with finite
projective limits, and commutes with arbitrary inductive limits.
\end{enonce*}

The\pageoriginale \hyperref[V.prop8.1.13.1]{propositions 8.1.13},
\hyperref[V.prop8.1.14]{8.1.14}, and \hyperref[V.prop8.1.12]{8.1.12}
reduce the assertion to the case $\Ss=\Tt=\Ll$. A local inductive
system then consists of a sheaf $\Ff_U$ on $\Ss/U$ for every
$U\in\Ob\Ss$, together, for every arrow $f:V\to U$ in $\Ss$, with an
arrow $f^{\ast}\Ff_U\to\Ff_V$, functorial in $f$.

For such inductive systems, the $\varinjlim$ are computed as follows:

\begin{enonce*}{Lemma 8.2.3}\etiquette[8.2.3]{V.lem8.2.3}
Let $U\longmapsto\Ff_U$ be a section of the category of presheaves on
the objects of $\Ss$. Then
$$
a(U\longmapsto\Gamma(U,\Ff_U))\simeq
\varinjlim_U a\Ff_U{\quoi}.
$$
\end{enonce*}

By definition, the right-hand side represents the functor which
assigns to every sheaf $\Ff$ on $\Ss$ the set of coherent systems of
arrows $\phi_U:\Ff_U\to\Ff|U$. In turn, an arrow $\phi_U$ is a
coherent system of arrows $\phi_g$, one for every $g:V\to U$, where
$\phi_g:\Ff_U(V)\to\Ff(V)$. Apply \hyperref[V.lem8.1.11]{8.1.11} to
the identity functor of $\Ss$, thereby obtaining $\Ss^{\ast}$ locally
filtered over $\Ss$. We have seen that
$$
\varinjlim_{g\in\Ss^{\ast}}\Ff_U(V_g)^{\sim}
\simeq\varinjlim_U a\Ff_U{\quoi},
$$
where $V_g$ is the source of $g$ and ${}^{\sim}$ denotes the
``generated constant sheaf'' functor. The functor
$U\longmapsto1_U$ from $\Ss$ to $\Ss^{\ast}$ is cofinal, and hence,
again by \hyperref[V.prop8.1.5]{8.1.5},
$$
\varinjlim_U\Ff_U(U)^{\sim}\simeq
\varinjlim_U a\Ff_U{\quoi}.
$$
Returning to the definitions, one finally finds
$$
\varinjlim_U\Ff_U(U)^{\sim}
=a(U\longmapsto\Gamma(U,\Ff_U)){\quoi}.
$$

This shows that, in the case to which we have reduced, the functor
$\varinjlim$ is defined everywhere and is computed as the composite
of the ``sheaf generated by a presheaf'' functor and the functor
$(\Ff_U)\longmapsto(U\longmapsto\Gamma(U,\Ff_U))$. Both of these
functors commute with finite projective limits, and the functor
$\varinjlim$ also commutes with arbitrary inductive limits by its
definition as an adjoint functor.

Lemma\pageoriginale \hyperref[V.lem8.2.3]{8.2.3} provides a systematic
method for proving properties of the local-inductive-limit functor
from the analogous properties of the ``sheaf associated with a
presheaf'' functor. One obtains in this way:

\begin{enonce*}{Proposition 8.2.4}\etiquette[8.2.4]{V.prop8.2.4}
Let $\Tt$, $\Ss$, and $\Ll$ be as above, let $\Aa_\lambda$ be a local
inductive system of sheaves of rings on $\Tt$, and let $\Ll_\lambda$
(\resp $\Mm_\lambda$) be a local inductive system of sheaves of right
(\resp left) modules over $\Aa_\lambda$, all three indexed by $\Ll$.

Then
$$
\varinjlim(\Ll_\lambda\otimes_{\Aa_\lambda}\Mm_\lambda)
\simeq\varinjlim\Ll_\lambda
\mathop{\otimes}_{\varinjlim\Aa_\lambda}
\varinjlim\Mm_\lambda{\quoi}.
$$
\end{enonce*}

(Reduce to the case $\Tt=\Ss=\Ll$, and note that both sides coincide
with the sheaf generated by the presheaf
$$
\Ll_U(U)\mathop{\otimes}_{\Aa_U(U)}\Mm_U(U){\quoi}.
$$
)

\begin{enonce*}{Proposition 8.2.5}\etiquette[8.2.5]{V.prop8.2.5}
Let $\Tt_1\xrightarrow{q}\Tt_2\to\Ss$ be two morphisms of sites, and
let $(F_\lambda)$ be a local inductive system of sheaves on $\Tt_2$
indexed by a locally cofiltered category $\Ll$ over $\Ss$. Then
$$
q^{\ast}\varinjlim_{\lambda\in\Ll}F_\lambda\simeq
\varinjlim_{\lambda\in\Ll}q^{\ast}F_\lambda{\quoi}.
$$
\end{enonce*}

This follows immediately from the fact that $q^{\ast}$ has a right
adjoint $q_\ast$.

Let $\Ss$ and $\Tt$ be two sites in which finite $\varprojlim$ are
representable, let $p:\Tt\to\Ss$ be a morphism of sites such that
$p^{\ast}:\Ss\to\Tt$ commutes with finite $\varprojlim$, and let
$\Ll$ be the locally cofiltered category over $\Tt$ whose objects are
triples $(V,U,\phi)$, where $V\in\Ob\Tt$, $U\in\Ob\Ss$, and
$\phi\in\Hom(V,p^{\ast}U)$. It is obtained by applying
\hyperref[V.lem8.1.11]{8.1.11} to $p:\Tt\to\Ss$.

\begin{enonce*}{Proposition 8.2.6}\etiquette[8.2.6]{V.prop8.2.6}
For every sheaf $\Ff$ on $\Tt$, one has
$$
\Ff=\varinjlim_{\lambda\in\Ll}\phi^{\ast}\phi_\ast\Ff,
$$
where, by abuse of notation, $\phi$ denotes the morphism induced by
$p$ from $\Tt/V$ to $\Ss/U$.
\end{enonce*}

Let\pageoriginale $\phi'$ be the inverse-image functor in the sense of
presheaves. From \hyperref[V.lem8.2.3]{8.2.3} one obtains
$$
\varinjlim_{\lambda\in\Ll}{\phi'}^{\ast}\phi_\ast\Ff
=\varinjlim_{\lambda\in\Ll}a{\phi'}^{\ast}\phi_\ast\Ff
=\varinjlim_{\lambda\in\Ll}\Ff(V)^{\sim}
=\varinjlim_V\Ff(V)^{\sim}=\Ff{\quoi}.
$$

\begin{enonce*}{Proposition 8.2.7}\etiquette[8.2.7]{V.prop8.2.7}
Let $p:T\to S$ be a morphism of topoi, let $\Oo_S$ be a sheaf of
rings on $S$, and let $\Oo_T$ be its inverse image on $T$. A right
$\Oo_T$-module $\Mm$ is flat over $\Oo_T$ if and only if the functor
\begin{equation*}
\Nn\longmapsto\Mm\mathop{\otimes}_{\Oo_T}p^{\ast}\Nn
\tag{8.2.7.1}
\end{equation*}\etiquette[8.2.7.1]{V.8.2.7.1}
from the category of left $\Oo_S$-modules to the category of abelian
sheaves on $T$ is exact.
\end{enonce*}

The ``only if'' assertion is trivial; suppose, therefore, that the
functor \eqhyperref[V.8.2.7.1]{8.2.7.1} is exact, and let us prove that
$\Mm$ is flat.

We may suppose that $p$ is defined by a morphism of sites
$p:\Tt\to\Ss$ satisfying the hypotheses made in \hyperref[V.prop8.2.6]{8.2.6}. Let
$V\in\Ob\Tt$, $U\in\Ob\Ss$, and
$\phi:V\to p^{\ast}U$. We again denote by $\phi$ the induced morphism
of sites $\phi:\Tt/V\to\Ss/U$.

\begin{enonce*}{Lemma 8.2.8}\etiquette[8.2.8]{V.lem8.2.8}
The functor
$\Nn\longmapsto\Mm|V\mathop{\otimes}_{\Oo_V}\phi^{\ast}\Nn$,
from the category of $\Oo_S|U$-modules on $\Ss/U$ to that of
$\Oo_T|V$-modules on $\Tt/V$, is exact.
\end{enonce*}

We reduce to taking $V=p^{\ast}U$. Let $j$ and $j'$ be the
``inclusion morphisms'' $j:\Ss/U\to\Ss$ and
$j':\Tt/p^{\ast}U\to\Tt$. One has
$$
j'_!(\Mm|V\mathop{\otimes}_{\Oo_V}\phi^{\ast}\Nn)
=\Mm\mathop{\otimes}_{\Oo_T}p^{\ast}j_!\Nn{\quoi},
$$
which gives \hyperref[V.lem8.2.8]{8.2.8}, since $j_!$ and $j'_!$ are exact and faithful.

Let now $Q$ be a left $\Oo_T$-module. By \hyperref[V.prop8.2.6]{(8.2.6)}, one has
$$
Q=\varinjlim_{\phi:V\to p^{\ast}U}
\phi^{\ast}\phi_\ast(Q|V),
$$
so that (\hyperref[V.prop8.2.4]{8.2.4})
$$
\Mm\mathop{\otimes}_{\Oo_T}Q
=\varinjlim_{\phi:V\to p^{\ast}U}
\Mm|V\mathop{\otimes}_{\Oo_V}
\phi^{\ast}\phi_\ast(Q|V){\quoi}.
$$
By\pageoriginale \hyperref[V.lem8.2.8]{8.2.8}, the functors
$Q\longmapsto\Mm|V\mathop{\otimes}_{\Oo_V}
\phi^{\ast}\phi_\ast(Q|V)$ are left exact; therefore
$\Mm\mathop{\otimes}_{\Oo_T}Q$ is left exact in $Q$, and hence exact.

\begin{enonce*}{Corollary 8.2.9}\etiquette[8.2.9]{V.cor8.2.9}
Let $f:(T,\Oo_T)\to(S,\Oo_S)$ be a morphism of ringed topoi. Then the
inverse image by $f$ of a sheaf of flat right modules is a sheaf of
flat modules.
\end{enonce*}

Factoring $f$ through $(T,f^{\ast}\Oo_S)$ reduces us to the case
$\Oo_T=f^{\ast}\Oo_S$. Let $\Mm$ then be a sheaf of flat modules on
$S$. By \hyperref[V.prop8.2.7]{8.2.7}, to verify that $f^{\ast}\Mm$
is flat, it is enough to verify the exactness of the functor
$$
\Nn\longmapsto f^{\ast}\Mm\mathop{\otimes}_{\Oo_T}f^{\ast}\Nn
=f^{\ast}(\Mm\mathop{\otimes}_{\Oo_S}\Nn),
$$
which proves the assertion, since $f^{\ast}$ is exact.

\begin{enonce*}{Lemma 8.2.10}\etiquette[8.2.10]{V.lem8.2.10}
Let $(S,\Oo_S)$ be a ringed topos, let $\Ff$ be a sheaf of left
modules locally of finite presentation, let $\Gg$ be a sheaf of
bimodules, and let $\Hh$ be a flat sheaf of left modules. Then the
canonical arrow
$$
\Hhom(\Ff,\Gg)\otimes\Hh\to
\Hhom(\Ff,\Gg\otimes\Hh)
$$
is an isomorphism.
\end{enonce*}

The question is local on $S$, which allows us to suppose that $\Ff$
has a finite presentation
$$
\Ff_1\to\Ff_0\to\Ff\to0{\quoi}.
$$
The first row of the following diagram is exact, since $\Hh$ is flat:
\[
\xymatrix{
0 \to \Hhom(\Ff,\Gg) \otimes \Hh \ar[r] \ar[d]&
\Hhom(\Ff_0,\Gg) \otimes \Hh \ar[r]\ar[d]^{S}&
\Hhom(\Ff_1,\Gg) \otimes \Hh \ar[d]^{S}\\
0 \to \Hhom(\Ff,\Gg \otimes \Hh) \ar[r]&
\Hhom(\Ff_0,\Gg \otimes \Hh) \ar[r]&
\Hhom(\Ff_1,\Gg \otimes \Hh)
}
\]
which proves the assertion. Taking $\Gg=\Oo_S$, one deduces from
\hyperref[V.lem8.2.10]{8.2.10}:

\begin{enonce*}{Lemma 8.2.11}\etiquette[8.2.11]{V.lem8.2.11}
Every\pageoriginale morphism from a sheaf of modules locally of finite
presentation to a flat sheaf of modules factors locally on $S$ through
a finite free sheaf $\Oo_S^n$ $(n\in\NN)$.
\end{enonce*}

\begin{enonce*}{Theorem 8.2.12}\etiquette[8.2.12]{V.thm8.2.12}
\emph{(D. Lazard).} A sheaf of modules $\Mm$ on a ringed site
$(S,\Oo_S)$ is flat if and only if it is a local inductive limit of
finite free modules.
\end{enonce*}

The ``if'' assertion follows from \hyperref[V.thm8.2.2]{8.2.2} and
\hyperref[V.prop8.2.4]{8.2.4}.

For every $U\in\Ob\Ss$, let $\Ll_0(U)$ (\resp $\Ll_1(U)$) be the
category of finite free sheaves of modules (\resp sheaves of modules
of finite presentation) on $U$, equipped with a linear map to
$\Mm|U$. For every arrow $f:V\to U$ in $\Ss$, let $f^{\ast}$ be the
restriction functor from $\Ll_i(U)$ to $\Ll_i(V)$. These functors
define a category $\Ll_i$, fibered over $\Ss$, whose fibers are the
opposite categories of the categories $\Ll_i(U)$.

The reader will verify that the category $\Ll_1$ is locally filtered
over $\Ss$ and that
$$
\Mm=\varinjlim_{(M,f)\in\Ll_1}M.
$$

If $\Mm$ is flat, the inclusion of $\Ll_0$ in $\Ll_1$ is fully
faithful and, by \hyperref[V.lem8.2.11]{8.2.11}, satisfies
$(C\mathfrak{F}1)$ of \hyperref[V.def8.1.4]{8.1.4}. It then
automatically satisfies $(C\mathfrak{F}2)$, and $\Ll_0$ is locally
filtered. By \hyperref[V.prop8.1.5]{8.1.5}, one has
$$
\Mm=\varinjlim_{(M,f)\in\Ll_0}M{\quoi},
$$
which proves \hyperref[V.thm8.2.12]{8.2.12}.

% END INCLUDED SGA4_English_current_buildable_source_20260906/source/05/drafts/SGA4_V_LINES_4066_4380_English.texfrag


% BEGIN INCLUDED SGA4_English_current_buildable_source_20260906/source/05/drafts/SGA4_V_LINES_4381_4410_English.texfrag
\begin{thebibliography}{99}\pageoriginale
\bibitem{sga4.exp.V.bib.1} M. Artin and B. Mazur: Etale Homotopy Theory, Lecture Notes no.~100, Springer-Verlag.
\bibitem{sga4.exp.V.bib.2} Blum and Herrera: Article forthcoming in Inventiones.
\bibitem{sga4.exp.V.bib.3} H. Cartan and S. Eilenberg: Homological Algebra.
\bibitem{sga4.exp.V.bib.4} P. Deligne: Théorie de Hodge (Publications de l'I.H.E.S.).
\bibitem{sga4.exp.V.bib.5} P. Gabriel and N. Popescu: CRAS.
\bibitem{sga4.exp.V.bib.6} P. Gabriel and M. Zisman: Homotopy Theory and Calculus of
  Fractions, Ergebnisse der Mathematik, Bd 35.
\bibitem{sga4.exp.V.bib.7} R. Godement: Théorie des faisceaux. Hermann.
\bibitem{sga4.exp.V.bib.8} J. Giraud: Méthode de la descente. Mémoire de la S.M.F.
\bibitem{sga4.exp.V.bib.9} J. Giraud: Algèbre homologique non commutative (forthcoming).
\bibitem{sga4.exp.V.bib.10} A. Grothendieck: On the De Rham Cohomology of Algebraic
  Varieties, I.H.E.S. no.~29.
\bibitem{sga4.exp.V.bib.11} A. Grothendieck: Sur quelques points d'Algèbre Homologique,
  Tohoku Math. Journal.
\bibitem{sga4.exp.V.bib.12} R. Hartshorne: Residues and Duality. Lecture Notes
  no.~20, Springer-Verlag.
\bibitem{sga4.exp.V.bib.13} L. Illusie: \SGA 6 I.
\bibitem{sga4.exp.V.bib.14} S. Lubkin: On a Conjecture of A. Weil. Am. J. of
  Math., p.~456, 1967.
\bibitem{sga4.exp.V.bib.15} G. Segal: Classifying Spaces and Spectral Sequences,
  I.H.E.S. no.~34.
\bibitem{sga4.exp.V.bib.16} Séminaire Cartan 1957--1958.
\bibitem{sga4.exp.V.bib.17} D. Sullivan: Geometric Topology, Part I, mimeographed
  notes, M.I.T., 1970.
\end{thebibliography}

\ifx\danslelivre\undefined
\end{document}
\fi

% END INCLUDED SGA4_English_current_buildable_source_20260906/source/05/drafts/SGA4_V_LINES_4381_4410_English.texfrag


% END INCLUDED SGA4_English_current_buildable_source_20260906/source/05/05_English.tex


\renewcommand*{\thechapter}{Vbis}
\makeatletter
\def\input@path{{../05bis/}}
\makeatother
% BEGIN INCLUDED SGA4_English_current_buildable_source_20260906/source/05bis/05bis_English.tex
\ifx\danslelivre\undefined
% Release engine: XeLaTeX (the LuaTeX compatibility layer breaks Xy-pic output).
\documentclass[10pt]{../fichiers-configuration-latex/smfart}
% BEGIN INCLUDED SGA4_English_current_buildable_source_20260906/source/fichiers-configuration-latex/paquets.tex
\usepackage{polyglossia}
%\usepackage[francais]{babel}
\usepackage{fontspec}
\usepackage[math-style=ISO,bold-style=upright,sans-style=upright]{unicode-math}
%\usepackage[mathletters]{ucs}
%\usepackage[utf8]{inputenc}
%\usepackage[T1]{fontenc}
\usepackage[unicode,hypertexnames=false,colorlinks=true,linkcolor=blue,citecolor=Sepia,urlcolor=magenta]{hyperref}
\setmainlanguage{french}
\setmainfont[Ligatures=TeX]{Linux Libertine O}
% variante possible : TeXGyrePagella ?
\setmathfont{XITS Math}

% maths
\usepackage{amsmath}
%\usepackage{amssymb}
%\usepackage{mathrsfs}
%\usepackage{stmaryrd}
%\usepackage{euscript}

% couleurs
\usepackage[usenames,dvipsnames]{xcolor}

% graphiques
\usepackage{graphicx}
\usepackage[all]{xy}
%\usepackage{tikz}
%\usetikzlibrary{matrix,arrows,calc}
%\tikzset{column sep=3em,row sep=2em,text height=1.5ex, text depth=0.25ex,>=angle 90}
%\tikzset{description/.style={fill=white,inner sep=1pt}}

% hyperliens, 'etiquettes
\usepackage{xr-hyper}
\usepackage{../fichiers-configuration-latex/smfthm}
%\ifx\cacheretiquettes\undefined
%\usepackage[inline]{showlabels} \else \fi % pour rendre visibles les 'etiquettes
%\usepackage{../fichiers-configuration-latex/makeidx} % pour faire un index
%\usepackage[pagebackref,unicode]{hyperref}  % pour voir o`u sont cit'ees les r'ef'erences
% hyperref is loaded once above with the cumulative-reader destination policy.
%\hypersetup{breaklinks=true} %permet le retour `a la ligne dans les liens trop longs. Incompatible avec pdflatex
%\renewcommand*{\backref}[1]{{\small [cit\'e page(s)~#1.]}}

\usepackage{xspace}
\usepackage{pifont} % pour expose XVII

% END INCLUDED SGA4_English_current_buildable_source_20260906/source/fichiers-configuration-latex/paquets.tex

% BEGIN INCLUDED SGA4_English_current_buildable_source_20260906/source/fichiers-configuration-latex/numerotation.tex
\renewcommand{\thefootnote}{\roman{footnote}} % pour mettre les notes en bas de page sous la forme (ii) et non (2), qui ressemble trop (?) a un exposant.

\newcounter{pageoriginal}
\newcommand{\pageoriginale}{\refstepcounter{pageoriginal}\marginpar{\footnotesize\textbf{\thepageoriginal}}\xspace}
\let\pageoriginaled\pageoriginale

\usepackage{ifthen}
\newboolean{original}
\def\sisi{\ifthenelse{\boolean{original}}}
\setboolean{original}{false}

\newcounter{cnote}
\def\nde#1{\ifthenelse{\boolean{original}}{}{\stepcounter{cnote}\def\thefootnote{\arabic{cnote}}\protect\footnote[\value{cnote}]{N.D.E.: #1}}}
\def\ndemark{\ifthenelse{\boolean{original}}{}{\stepcounter{cnote}\def\thefootnote{\arabic{cnote}}\protect\footnotemark[\value{cnote}]}}
\def\ndetext#1{\ifthenelse{\boolean{original}}{}{\def\thefootnote{\arabic{cnote}}\protect\footnotetext[\value{cnote}]{N.D.E.: #1}}}

% END INCLUDED SGA4_English_current_buildable_source_20260906/source/fichiers-configuration-latex/numerotation.tex

% BEGIN INCLUDED SGA4_English_current_buildable_source_20260906/source/fichiers-configuration-latex/macros.tex
\setmainlanguage{french}
\setotherlanguage{english}


%% Fixed, engine-independent release identifier.  The inherited dynamic Git
%% probe used Lua io.popen, exposed build-state-dependent output, and printed
%% a visible question mark outside a Git checkout.
\def\showgitstatus{English release r2, 2026-07-23}


% macros pour LaTeX
\renewcommand{\mathsf}{}

% Police différente pour U+2113 SCRIPT SMALL L
\setmathfont[range="2113]{latinmodern-math.otf}

% Give every inherited semantic label its own stable PDF destination instead
% of inheriting whichever automatic equation/page anchor happened to be live.
% The legacy label remains as a compatibility alias used by existing edges.
\makeatletter
\def\etiquette[#1]#2{%
  \Hy@raisedlink{\hyper@anchorstart{sga4.semantic.#2}\hyper@anchorend}%
  \def\@currentHref{sga4.semantic.#2}%
  \def\@currentlabel{#1}%
  \label{#2}%
}
\makeatother
\def\eqhyperref[#1]#2{(\hyperref[#1]{#2})}
\def\pagehyperref[#1]#2{(page \hyperref[#1]{#2})}
\def\hyperrefext[#1][#2]#3{\hyperref[#2]{#1 #3}}
\def\eqhyperrefext[#1][#2]#3{\hyperref[#2]{#1 #3}}     %cf. exposes 17
\def\neweqhyperrefext[#1][#2]#3{\hyperref[#2]{#1 #3}}  %cf. exposes 17

\def\shat#1{\protect\widehat{#1}}
\def\stilde#1{\protect\widetilde{#1}}
%%%%%%%%%%%%%%%%%%


\def\etcn{etc.\xspace}
\def\etco{etc...\xspace}
\def\etc{\sisi{\etco}{\etcn}}
\def\ldotsn{\sisi{...}{\etcn}}

\def\cf{{cf.}\kern.3em}
\def\Cf{{Cf.}\kern.3em}
\def\ie{{i.e.}\kern.3em}
\def\iev{{i.e.},\xspace}
\def\loccito{loc.\kern2pt cit.\xspace}
\def\loccitn{\textit{loc.\kern2pt cit.}\xspace}
\def\loccit{\sisi{\loccito}{\loccitn}}
\def\EGA{\textup{EGA}\kern.2em}
\def\SGA{\textup{SGA}\kern.3em}

\def\quoi{}

\def\quav{\quad,\quad}
\def\quaet{\quad \text{et}\quad}
\def\fibreproduct#1#2#3{#1 ×_{#3} #2}
\let\fprod\fibreproduct
\def\ccheck#1#2{\mbox{${\check{\check{\rm #1}}}^{\lower2pt\hbox{$\scriptstyle #2$}}$}}
\def\sheaf#1{\protect\mathcal{#1}}
\def\vieuxsheaf#1{{#1}} % cf. problème avec sheaf{a\Phi_i} par exemple...
\DeclareMathOperator{\Hom}{Hom}
\DeclareMathOperator{\can}{can}
\DeclareMathOperator{\Ind}{Ind}
\DeclareMathOperator{\ind}{ind}
%\DeclareMathOperator{\Hhomtop}{\mathcal{H}\!\mathit{omtop}}
\DeclareMathOperator{\Hhomtop}{\mathbf{Homtop}}
\DeclareMathOperator{\Point}{Point}
\DeclareMathOperator{\Pro}{Pro}
\DeclareMathOperator{\Top}{Top}
\DeclareMathOperator{\Mod}{Mod}
\DeclareMathOperator{\Cart}{Cart}
%\DeclareMathOperator{\Mmorsite}{\mathcal{M}\!\mathit{orsite}}
\DeclareMathOperator{\Mmorsite}{\mathbf{Morsite}}
\DeclareMathOperator{\loc}{loc}
\DeclareMathOperator{\Loc}{Loc}
\DeclareMathOperator{\Pic}{Pic}
\DeclareMathOperator{\Et}{Et}
\DeclareMathOperator{\card}{card}
\DeclareMathOperator{\Ens}{Ens}
\DeclareMathOperator{\h}{\underline{H}}
\DeclareMathOperator{\Fonct}{Fonct}
\DeclareMathOperator{\Ker}{Ker}
\DeclareMathOperator{\but}{but}
\DeclareMathOperator{\Ob}{Ob}
\DeclareMathOperator{\ob}{ob}
\DeclareMathOperator{\Lib}{Lib}
\DeclareMathOperator{\esj}{esj}
\DeclareMathOperator{\id}{id}
\DeclareMathOperator{\Id}{Id}
\DeclareMathOperator{\ad}{ad}
\DeclareMathOperator{\Iim}{Im}
\DeclareMathOperator{\pr}{pr}
\DeclareMathOperator{\Cat}{Cat}
\DeclareMathOperator{\car}{car}
\DeclareMathOperator{\SysInd}{SysInd}
\DeclareMathOperator{\Sup}{Sup}
\DeclareMathOperator{\Coker}{Coker}
\DeclareMathOperator{\Filt}{Filt}
\DeclareMathOperator{\acc}{acc}
\DeclareMathOperator{\sur}{sur}
\DeclareMathOperator{\cat}{cat}
\DeclareMathOperator{\Ord}{Ord}
\DeclareMathOperator{\Gr}{Gr}
\DeclareMathOperator{\Cov}{Cov}
\DeclareMathOperator{\Caf}{Caf}
\DeclareMathOperator{\coker}{coker}
\DeclareMathOperator{\coim}{coim}
\DeclareMathOperator{\im}{im}
\DeclareMathOperator{\Ouv}{Ouv}
\DeclareMathOperator{\diag}{diag}
\DeclareMathOperator{\Esp}{Esp}
\DeclareMathOperator{\TOP}{TOP}
\DeclareMathOperator{\Unif}{Unif}
\DeclareMathOperator{\Iind}{\mathbf{Ind}}
\DeclareMathOperator{\Hhomcart}{\mathbf{Homcart}}
\DeclareMathOperator{\Hhomtopess}{\mathbf{Homtopess}}
\DeclareMathOperator{\Ppoints}{\mathbf{Points}}
\DeclareMathOperator{\Ppoint}{\mathbf{Point}}
\DeclareMathOperator{\Ttors}{\mathbf{Tors}}
\DeclareMathOperator{\Ssex}{\mathbf{Sex}}
\DeclareMathOperator{\Ffib}{\mathbf{Fib}}
\DeclareMathOperator{\Vvois}{\mathbf{Vois}}
\DeclareMathOperator{\Ppointess}{\mathbf{Pointess}}
\DeclareMathOperator{\scrHom}{\mathbf{Hom}}
\DeclareMathOperator{\sheafPic}{\mathbf{Pic}}
\DeclareMathOperator{\sheafAut}{\mathbf{Aut}}
\DeclareMathOperator{\Hhomtoplocan}{\mathbf{Homtoplocan}}
\DeclareMathOperator{\Hhomtopan}{\mathbf{Homtopan}}
\DeclareMathOperator{\Pptgeom}{\mathbf{Ptgéom}}
\DeclareMathOperator{\Morsite}{Morsite}
\DeclareMathOperator{\Fib}{Fib}
\DeclareMathOperator{\Fais}{Fais}
\DeclareMathOperator{\kar}{kar}
\DeclareMathOperator{\Groupes}{Groupes}
\DeclareMathOperator{\Restr}{Restr}
\DeclareMathOperator{\Isom}{Isom}
\DeclareMathOperator{\Prol}{Prol}
\DeclareMathOperator{\Res}{Res}
\DeclareMathOperator{\Tors}{Tors}
\DeclareMathOperator{\esp}{esp}
\DeclareMathOperator{\es}{es}
\DeclareMathOperator{\od}{od}
\DeclareMathOperator{\gr}{gr}
\DeclareMathOperator{\sob}{sob}
\DeclareMathOperator{\dfn}{dfn}
\DeclareMathOperator{\Sch}{Sch}
\DeclareMathOperator{\ouv}{ouv}
\DeclareMathOperator{\cons}{cons}
\DeclareMathOperator{\Nerf}{Nerf}
\DeclareMathOperator{\Points}{Points}
\DeclareMathOperator{\Espsob}{Espsob}
\DeclareMathOperator{\Homtop}{Homtop}
\DeclareMathOperator{\Pt}{Pt}
\DeclareMathOperator{\Quot}{Quot}
\DeclareMathOperator{\glob}{glob}
\DeclareMathOperator{\Homfib}{Homfib}
\DeclareMathOperator{\Qouv}{Qouv}
\DeclareMathOperator{\Quotouv}{Quotouv}
\DeclareMathOperator{\ens}{ens}
\DeclareMathOperator{\ab}{ab}
\DeclareMathOperator{\an}{an}
\DeclareMathOperator{\fs}{fs}
\DeclareMathOperator{\red}{red}
\DeclareMathOperator{\Ab}{Ab}
\DeclareMathOperator{\cosq}{cosq}
\DeclareMathOperator{\Ann}{Ann}
\DeclareMathOperator{\Zar}{Zar}


%%%%%%%%%%%%
%%%%% SGA4-2
\DeclareMathOperator{\Sex}{Sex}
\DeclareMathOperator{\vois}{vois}
\DeclareMathOperator{\Vois}{Vois}
\DeclareMathOperator{\Ext}{Ext}
\DeclareMathOperator{\ext}{ext}
\DeclareMathOperator{\Fer}{Fer}
\DeclareMathOperator{\Add}{Add}
\DeclareMathOperator{\Cocart}{Cocart}
\DeclareMathOperator{\cosk}{cosk}
\DeclareMathOperator{\plat}{plat}

\DeclareMathOperator{\faisc}{faisc}
\DeclareMathOperator{\sk}{sk}
\DeclareMathOperator{\ch}{ch}

\def\HH{\mathbf{H}}
\def\qc{\mathrm{qc}}
\def\et{\mathrm{et}}
\def\coh{\mathrm{coh}}

\def\Hvr#1{%
\begin{array}[b]{l}
{\kern2pt\scriptstyle r}\\[-8pt]
{\kern1pt\scriptstyle \vee}\\[-5pt]
H^{#1}
\end{array}}

\def\Hvv#1{%
\begin{array}[b]{l}
\kern2pt\rotatebox{90}{${\scriptscriptstyle\ll}$}\\[-6pt]
H^{#1}
\end{array}}

\def\Hvi#1{%
\begin{array}[b]{l}
{\kern0.2pt\scriptscriptstyle \infty}\\[-8pt]
{\kern1pt\scriptstyle \vee}\\[-5pt]
H^{#1}
\end{array}}

\def\Hv#1{%
\begin{array}[b]{l}
{\kern1pt\scriptstyle \vee}\\[-5pt]
H^{#1}
\end{array}}

\DeclareMathOperator{\cart}{cart}
\DeclareMathOperator{\Ffaisc}{\mathbf{Faisc}}


\def\Ev#1{%
\begin{array}[b]{l}
{\kern1pt\scriptstyle \vee}\\[-5pt]
E'_{#1}
\end{array}}

\DeclareMathOperator{\Cont}{Cont}
\DeclareMathOperator{\Ccont}{\mathscr{C}\!\mathit{ont}}
\DeclareMathOperator{\Site}{Site}
\newcommand\MM{\mathbf{M}}
\def\limtop#1{{\displaystyle{\varprojlim_{#1}\!\mbox{\rm top }}}}
\def\Limtop{\mathop{\varprojlim\!\mbox{\rm top}}}
\DeclareMathOperator{\spec}{spec}
\def\ss{\scriptstyle{\mathscr{S}}}
\DeclareMathOperator{\Gal}{Gal}
\def\smallm{\scriptstyle{\mathscr{M}}}
\DeclareMathOperator{\End}{End}
\DeclareMathOperator{\ass}{ass}
\DeclareMathOperator{\zar}{zar}
\DeclareMathOperator{\Ogg}{Ogg}
\DeclareMathOperator{\res}{res}
\DeclareMathOperator{\tr}{tr}
\DeclareMathOperator{\Supp}{Supp}
\DeclareMathOperator{\codim}{codim}
\DeclareMathOperator{\rk}{rk}
\DeclareMathOperator{\cur}{cur}
\DeclareMathOperator{\cdqc}{cdqc}
\newcommand\DD{\mathbf{D}}
\DeclareMathOperator{\cl}{cl}
\def\JJ{\mathbf{J}}
\def\FF{\mathbf{F}}
\def\BB{\mathbf{B}}
\DeclareMathOperator{\Alget}{Alget}
\DeclareMathOperator{\Ass}{Ass}
\DeclareMathOperator{\supp}{supp}
\DeclareMathOperator{\rad}{rad}
\DeclareMathOperator{\prof}{prof}
\DeclareMathOperator{\cint}{int}
\DeclareMathOperator{\sic}{sic}
\DeclareMathOperator{\Tor}{Tor}
\DeclareMathOperator{\tors}{tors}
\DeclareMathOperator{\Sym}{Sym}
\DeclareMathOperator{\parf}{parf}
\DeclareMathOperator{\tor}{tor}
\DeclareMathOperator{\defn}{defn}
\DeclareMathOperator{\Div}{Div}
\DeclareMathOperator{\Pol}{Pol}
\DeclareMathOperator{\Fex}{Fex}
\DeclareMathOperator{\Lie}{Lie}
\DeclareMathOperator{\tot}{tot}
\DeclareMathOperator{\Aut}{Aut}
\DeclareMathOperator{\HOM}{HOM}
\DeclareMathOperator{\pch}{pch}
\DeclareMathOperator{\Ch}{Ch}
\DeclareMathOperator{\EXT}{EXT}
\DeclareMathOperator{\PIC}{PIC}
\DeclareMathOperator{\Rf}{Rf}
\DeclareMathOperator{\inj}{inj}
\newcommand\VV{\mathbf{V}}

\def\Xv{%
\begin{array}[b]{l}
{\kern1pt\scriptstyle \vee}\\[-5pt]
X
\end{array}}

%%%%%%%
\newcommand\TT{\mathbf{T}}
\newcommand\XX{\mathbf{X}}
\newcommand\EE{\mathbf{E}}
%%%%%%

\def\qedsymbol{\textup{C.Q.F.D.}}
\let\CQFD\qed
\let\cqfd\qed

%%%%%\bigdot
\font\bigsymb=cmsy10 at 4pt
\def\bigdot{{\kern1.2pt\raise 1.5pt\hbox{\bigsymb\char15}}}
%%%%%

\def\Lright{\mathop{\vcenter{\hbox{\hbox{$\longrightarrow$}\lower8pt \hbox{\kern-16pt\hbox{$\longrightarrow$}}}}}}

\def\Aa{{\mathcal A}}
\def\Bb{{\mathcal B}}
\def\Cc{{\mathcal C}}
\def\Dd{{\mathcal D}}
\def\Ee{{\mathcal E}}
\def\Ff{{\mathcal F}}
\def\Gg{{\mathcal G}}
\def\Hh{{\mathcal H}}
\def\Ii{{\mathcal I}}
\def\Jj{{\mathcal J}}
\def\Kk{{\mathcal K}}
\def\Ll{{\mathcal L}}
\def\Mm{{\mathcal M}}
\def\Nn{{\mathcal N}}
\def\Oo{{\mathcal O}}
\def\Pp{{\mathcal P}}
\def\Qq{{\mathcal Q}}
\def\Rr{{\mathcal R}}
\def\Ss{{\mathcal S}}
\def\Tt{{\mathcal T}}
\def\Uu{{\mathcal U}}
\def\Vv{{\mathcal V}}
\def\Ww{{\mathcal W}}
\def\Xx{{\mathcal X}}
\def\Yy{{\mathcal Y}}
\def\Zz{{\mathcal Z}}

\newcommand{\CC}{\mathbf{C}}
\newcommand{\GG}{\mathbf{G}}
\newcommand{\KK}{\mathbf{K}}
\newcommand{\LL}{\mathbf{L}}
\newcommand{\NN}{\mathbf{N}}
\newcommand{\PP}{\mathbf{P}}
\newcommand\QQ{\mathbf{Q}}
\newcommand\RR{\mathbf{R}}
\newcommand\ZZ{\mathbf{Z}}

\def\to{\mathchoice{\longrightarrow}{\rightarrow}{\rightarrow}{\rightarrow}}
\def\from{\mathchoice{\longleftarrow}{\leftarrow}{\leftarrow}{\leftarrow}}
\def\mto{\mathchoice{\longmapsto}{\mapsto}{\mapsto}{\mapsto}}
\def\mfrom{\mathrel{\from\joinrel\mapstochar\,}}
\def\hto{\mathrel{\lhook\joinrel\to}}
\def\hfrom{\mathrel{\from\joinrel\rhook}}

\def\implies{\mathchoice{\Longrightarrow}{\Rightarrow}{\Rightarrow}{\Rightarrow}}
\def\implied{\mathchoice{\Longleftarrow}{\Leftarrow}{\Leftarrow}{\Leftarrow}}
\def\iff{\mathchoice{\Longleftrightarrow}{\Leftrightarrow}{\Leftrightarrow}{\Leftrightarrow}}

\def\To#1{\mathchoice{\xrightarrow{\labelstyle~~#1~}}{\stackrel{#1}{\longrightarrow}}{}{}}
\def\Mto#1{\mathrel{\mapstochar\To{#1}}}
\def\From#1{\mathchoice{\xleftarrow{\labelstyle~~#1~}}{\stackrel{#1}{\longleftarrow}}{}{}}

\newcommand{\isomto}{\overset{\sim}{\longrightarrow}}
\newcommand{\isomfrom}{\overset{\sim}{\longleftarrow}}

\newcommand{\IFF}{\unskip~$\iff$~\ignorespaces}
\newcommand{\THEN}{\unskip~$\implies$~\ignorespaces}

\let\leq\leqslant
\let\geq\geqslant
\let\le\leqslant
\let\ge\geqslant

\def\cf{{cf.}\kern.3em}
\def\Cf{{Cf.}\kern.3em}
\def\ie{{i.e.}\kern.3em}
\def\iev{{i.e.},\xspace}
%\def\loccito{loc.\kern2pt cit.\xspace}
%\def\loccitn{\textit{loc.\kern2pt cit.}\xspace}
%\def\loccit{\sisi{\loccito}{\loccitn}}
\def\EGA{\textup{EGA}\kern.2em}
\def\SGA{\textup{SGA}\kern.3em}

\let\oldS\S
\def\S{\oldS\kern .15em}
\def\ptbl{.\kern .2em}
\def\resp{resp.\kern.3em}
\def\Exp{Exp.\kern.3em}

\let\bar\overline
\let\hat\widehat
\let\tilde\widetilde
\let\dpl\displaystyle
\let\ldots\dots
\let\emptyset\varnothing
\let\skpt\skippointrait
\let\epsilon\varepsilon

\def\sfrac#1#2{{#1}/{#2}}

\DeclareMathOperator{\cd}{cd}
%\DeclareMathOperator{\Hhom}{\mathcal{H}\!\mathit{om}}
\DeclareMathOperator{\Hhom}{\mathbf{Hom}}
%\DeclareMathOperator{\Eext}{\mathcal{E}\!\mathit{xt}}
\DeclareMathOperator{\Eext}{\mathbf{Ext}}
\DeclareMathOperator{\Spec}{Spec}
\DeclareMathOperator{\Tr}{Tr}

\let\emptyset=\varnothing

%%%%%%%%%%%%%%%
%%% expose XVII
%%%%%%%%%%%%%%%

\makeatletter
\let\old@optexpose\optexpose
\newcommand{\empty@optexpose}[1]{}
\newcommand{\paren@optexpose}[1]{\old@optexpose{#1}(}
\newcommand{\refnoexpose}[1]{\let\optexpose\empty@optexpose\ref{#1}\let\optexpose\old@optexpose}%\refnoexpose{X.cor5.2} engendre "5.2" au lieu de "X 5.2"
\newcommand{\neweqref}[1]{\let\optexpose\paren@optexpose\textup{\ref{#1})}\let\optexpose\old@optexpose}%\neweqref{XII.eq4.1} engendre "XII (4.1)" au lieu de "(XII 4.1)"
\makeatother
\newcommand{\refvb}[2][VI${}_{\text{B}}$ ]{\sisi{#1\refnoexpose{#2}}{\ref{#2}}}

\newcommand{\leftexp}[2]{{\vphantom{#2}}^{#1}{#2}}
\newcommand{\AAA}{\mathbf{A}}
\newcommand{\Sf}{\mathfrak{S}}
\DeclareMathOperator*{\quotprojlim}{\text{``$\varprojlim$''}}
\DeclareMathOperator*{\quotinjlim}{\text{``$\varinjlim$''}}
\DeclareMathOperator{\Fl}{F\sisi{\ell}{l}}
\DeclareMathOperator{\Pl}{P\sisi{\ell}{l}}
%\DeclareMathOperator{\cl}{c\sisi{\ell}{l}}
\DeclareMathOperator{\tdf}{tdf}
\DeclareMathOperator{\rg}{rg}
\newcommand{\iInd}{\sisi{I}{i}nd}
\newcommand{\iind}{ind}
\newcommand{\ppro}{pro}
\newcommand{\tTor}{\sisi{T}{t}or}
\newcommand{\ttor}{tor}
\newcommand{\chap}{Chap.}
\newcommand{\Ffex}{\Ff\mathit{ex}}
\newcommand{\sheafk}{\underline{k}}
\newcommand{\sheafn}{\underline{n}}
\newcommand{\sheafGamma}{\underline{\Gamma}}

% Repeated references to an already typeset footnote must not use
% \footnotemark without a matching \footnotetext: Hyperref would then create
% a live annotation to a destination that cannot exist.  This helper keeps the
% current printed footnote mark and links it to an explicit stable note target.
\makeatletter
\newcommand{\sgarepeatedfootnotemark}[1]{%
  \begingroup
  \edef\@thefnmark{\thefootnote}%
  \hyperlink{#1}{\@makefnmark}%
  \endgroup}
\makeatother

\sisi{
\DeclareMathOperator{\simarrow}{\rule{.15mm}{1.8mm}\kern -3.4pt\rightsquigarrow}
\newcommand{\sheafell}{\underline{\ell}}
\newcommand{\sheafZZ}{\underline{\ZZ}}
}{
\newcommand{\nderef}[1]{la N.D.E.~(\ref{#1}) page~\pageref{#1}}
}

%%%%%%%%%%%%%%%%%
%%%%% expose XVII
%%%%%%%%%%%%%%%%%

% END INCLUDED SGA4_English_current_buildable_source_20260906/source/fichiers-configuration-latex/macros.tex


\externaldocument{../01/01_English.tex}
\externaldocument{../02/02_English.tex}
\externaldocument{../03/03_English.tex}
\externaldocument{../04/04_English.tex}
\externaldocument{../05/05_English.tex}
%\externaldocument{../05bis/05bis_English.tex}

\setboolean{original}{false}

\begin{document}

\begin{center}
Vbis. Cohomological Descent Techniques\\
B. Saint-Donat
\end{center}

\selectlanguage{english}version: \showgitstatus
\tableofcontents

\else
\chapter{Cohomological Descent Techniques}
%\addtocontents{toc}{B. Saint-Donat}
\begin{center}
B. Saint-Donat
\end{center}
\fi
\phantomsection\MakeLinkTarget*{sga4.exp.Vbis}\label{sga4.exp.Vbis}

\setcounter{pageoriginal}{83}

%\footnotetext[1]{From concise notes by P. Deligne}

\let\cosk=\cosq
\let\Ysheaf=\sheaf

% BEGIN INCLUDED SGA4_English_current_buildable_source_20260906/source/05bis/drafts/SGA4_VBIS_LINES_59_205_English.texfrag
\section*{Introduction}

\begin{enumerate}
\renewcommand{\labelenumi}{\theenumi.}
\item Let\pageoriginale $X$ be a topological space and let
$\Uu=(U_i)_{i\in I}$ be a covering of $X$, assumed either open or
closed and locally finite. If $\mathfrak{F}$ is an abelian sheaf on
$X$, the Leray spectral sequence
\begin{equation*}
\Hv{p}(\Uu,\mathcal{H}^q(\mathfrak{F}))\Rightarrow H^{\ast}
(X,\mathfrak{F})\tag{1.1}
\end{equation*}\etiquette[1.1]{Vbis.eq1.1}
defined by $\Uu$ [\cite{sga4.exp.Vbis.bib.5} II (5.2.4) and (5.4.1)] may be described
as follows.

The covering $\Uu$ defines a ``\v{C}ech'' resolution
$\mathcal{C}^{\ast}(\Uu,\mathfrak{F})$, functorial in
$\mathfrak{F}$ (ibid. (\hyperref[Vbis.subsubsec5.2.1]{5.2.1})). On
the other hand, for every $\mathfrak{F}$ there is a ``canonical
flasque'' resolution $C^{\ast}(\mathfrak{F})$, functorial in
$\mathfrak{F}$ (ibid. (\hyperref[Vbis.subsec4.3]{4.3})). With this
notation, when $\Uu$ is open the spectral sequence
(\hyperref[Vbis.subsec1.1]{1.1}) is obtained from the double complex
\begin{equation*}
\Gamma(X,\mathcal{C}^{\ast}(\Uu,C^{\ast}(\mathfrak{F}))){\quoi}.
\tag{1.2}
\end{equation*}\etiquette[1.2]{Vbis.eq1.2}
When $\Uu$ is closed and locally finite, one considers the double
complex
\begin{equation*}
\Gamma(X,C^{\ast}(\mathcal{C}^{\ast}(\Uu,\mathfrak{F}))){\quoi}.
\tag{1.3}
\end{equation*}\etiquette[1.3]{Vbis.eq1.3}

\item Let us seek a unified description of these double complexes.
Let $X_0$ be the topological disjoint union of the $U_i$, and for
$n\geq0$ let $X_n$ be the $(n+1)$-fold fiber product of $X_0$ with
itself over $X$:
\begin{equation*}
X_n=\mathop{\coprod}_{i_0,\ldots,i_n\in I}U_{i_0}\cap\cdots\cap U_{i_n}
=\mathop{\coprod}_{\sigma\in\Hom([n],I)}
\mathop{\bigcap}_{i=0}^{n}U_{\sigma(i)}{\quoi}.\tag{2.1}
\end{equation*}\etiquette[2.1]{Vbis.eq2.1}
The $X_n$ form a simplicial system of topological spaces, and if
$j_n$ denotes the projection from $X_n$ to $X$, then
\begin{equation*}
\mathcal{C}^n(\Uu,\mathfrak{F})=j_{n\ast}j_n^{\ast}(\mathfrak{F}){\quoi}.
\tag{2.2}
\end{equation*}\etiquette[2.2]{Vbis.eq2.2}

Note that formation of the canonical flasque resolutions commutes
with\pageoriginale restriction to an open subset and with direct
image by a closed immersion. Consequently:
\begin{enumerate}
\renewcommand{\theenumi}{\alph{enumi}}
\renewcommand{\labelenumi}{\theenumi)}
\item if $\Uu$ is open,
\begin{equation*}
\mathcal{C}^q(\Uu,C^p(\mathfrak{F}))=j_{q\ast}j_q^{\ast}
C^p(\mathfrak{F})=j_{q\ast}C^p(j_q^{\ast}\mathfrak{F});\tag{2.3}
\end{equation*}\etiquette[2.3]{Vbis.eq2.3}
\item if $\Uu$ is closed and locally finite,
\begin{equation*}
C^p(\mathcal{C}^q(\Uu,\mathfrak{F}))=C^p(j_{q\ast}j_q^{\ast}
\mathfrak{F})=j_{q\ast}C^p(j_q^{\ast}\mathfrak{F}){\quoi}.\tag{2.4}
\end{equation*}\etiquette[2.4]{Vbis.eq2.4}
\end{enumerate}
Thus, to obtain a unified description of
\eqhyperref[Vbis.eq1.2]{1.2} and \eqhyperref[Vbis.eq1.3]{1.3}, it is
enough to take the canonical flasque resolution of
$j_q^{\ast}(\mathfrak{F})$ on $X_q$ for every $q$, and then apply
the functor $j_{q\ast}$ to this resolution.

\item The preceding description remains meaningful for any
simplicial system of topological spaces over $X$:
\hypertarget{sga4.exp.Vbis.equation.3.1}{}\begin{gather*}
\Delta^0\longrightarrow\Top/X\tag{3.1}\\
[n]\longmapsto X_n.\notag
\end{gather*}
Such a system need not be of the form \eqhyperref[Vbis.eq2.1]{2.1}.
The double complex, coaugmented by $\mathfrak{F}$,
\begin{equation*}
\mathfrak{F}\longrightarrow
(j_{q\ast}C^p(j_q^{\ast}(\mathfrak{F})))_{p,q}\tag{3.2}
\end{equation*}\etiquette[3.2]{Vbis.eq3.2}
will not, however, in general define a resolution of
$\mathfrak{F}$.

This work is devoted to finding sufficient conditions for
\eqhyperref[Vbis.eq3.2]{3.2} to define a resolution of
$\mathfrak{F}$. In that case, the spectral sequence
(\hyperref[Vbis.subsec1.1]{1.1}) generalizes to a spectral sequence
\begin{equation*}
\Hv{p}(H^q(X_p,j_p^{\ast}(\mathfrak{F})))\Rightarrow
H^{p+q}(X,\mathfrak{F}),\tag{3.3}
\end{equation*}\etiquette[3.3]{Vbis.eq3.3}
called the ``descent spectral sequence''.

In\pageoriginale the case of ``constant coefficients'', analogous
spectral sequences were obtained by Segal (\cf \cite{sga4.exp.Vbis.bib.8}), by other
methods and for other ``cohomology theories'', such as $K$-theory.
Segal works in the category of CW-complexes: he uses a ``geometric
realization'' functor which associates a new topological space with
a semi-simplicial complex of topological spaces; this new space is
to be compared with the topos associated with a simplicial topos
[\cf \hyperref[Vbis.prop1.2.12]{(1.2.12)}].

\item In \hyperref[Vbis.sec5]{Section 5} we illustrate the criteria obtained by using
resolution of singularities to construct, for every analytic space
$X$ over $\CC$, a simplicial system of nonsingular analytic spaces
over $X$,
$$
[n]\longmapsto X_n,
$$
such that \eqhyperref[Vbis.eq3.2]{3.2} defines a resolution of
$\mathfrak{F}$. Taking the constant sheaf $\CC$ for
$\mathfrak{F}$ gives in particular a spectral sequence
\begin{equation*}
H^q(X_p,\CC)\Longrightarrow H^{p+q}(X,\CC)\tag{4.1}
\end{equation*}\etiquette[4.1]{Vbis.eq4.1}
which expresses the complex cohomology of $X$ in terms of the
complex cohomology of nonsingular analytic spaces. Moreover, if $X$
is projective, all the $X_p$ may be assumed projective; this is the
essential ingredient which yields a kind of ``Hodge theory'' for
$X$ (\cf \cite{sga4.exp.Vbis.bib.2}).

\item The preceding constructions extend verbatim when the sheaf
$\mathfrak{F}$ is replaced by a bounded-below complex of sheaves.
They lead to ``localization'' techniques in derived categories.

It is known that, for $X_0$ given by \eqhyperref[Vbis.eq2.1]{2.1},
the arrow
$$
j_0^{\ast}:D^b(X)\longrightarrow D^b(X_0)
$$
is not faithful in general. We shall show that data more precise
than $j_0^{\ast}(K^{\bigdot})$\pageoriginale (for
$K^{\bigdot}\in D^+(X)$), involving the $X_n$, sometimes make it
possible to reconstruct the complex $K^{\bigdot}$.

The resulting statements will be used in the appendix to Exposé
XVII to extend the definition of the functor $Rf_!$ (where $f$ is a
separated morphism of finite type between schemes) to the case in
which $f$ is not assumed compactifiable.

For this application it is not possible to consider only
topological spaces, which are replaced here by the étale sites of
schemes. Moreover, to carry out the proofs it will be necessary to
consider both simplicial systems of spaces and multisimplicial
systems. This explains, justifies, or excuses the degree of
hypergenerality from which we shall begin.
\end{enumerate}


% END INCLUDED SGA4_English_current_buildable_source_20260906/source/05bis/drafts/SGA4_VBIS_LINES_59_205_English.texfrag


% BEGIN INCLUDED SGA4_English_current_buildable_source_20260906/source/05bis/drafts/SGA4_VBIS_LINES_206_516_English.texfrag
\section{Preliminaries}\etiquette[1]{Vbis.sec1}

\subsection{Notation}\etiquette[1.1]{Vbis.subsec1.1}

\subsubsection{}\etiquette[1.1.1]{Vbis.subsubsec1.1.1}
Throughout what follows, $\Uu$ is a universe such that $\ZZ \in
\Uu$: \emph{all the topoi considered will be
$\Uu\text{-}topoi$.}

Let $T$ and $T'$ be two topoi: a morphism $\varphi : T \to T'$
consists of a pair of functors $\varphi_{\ast} : T \to T'$ and
$\varphi^{\ast} : T'\to T$, endowed with an adjunction
$$\Hom_{T'}(.,\varphi_{\ast}.) \isomto{\Hom}_{T}(\varphi^{\ast} . ,
.),$$
such that $\varphi^{\ast}$ is left exact (\ie preserves finite
projective limits).

Let $(T, \mathcal{O}_{T})$ and $(T', \mathcal{O}_{T'})$ be two
ringed topoi: a morphism from $(T, \mathcal{O}_{T})$ to
$(T', \mathcal{O}_{T'})$ is a pair $(\varphi, \theta)$, where
$\varphi : T \to T'$ is a morphism of topoi and
$\theta : \mathcal{O}_{T'} \to \varphi_{\ast} (\mathcal{O}_T)$ is a
morphism of rings.

\subsubsection{}\etiquette[1.1.2]{Vbis.subsubsec1.1.2}
We will make constant use of the language of fibered categories as
presented in [\SGA 1 VI]; the reader may also consult \cite{sga4.exp.Vbis.bib.4}. Let
us merely fix some notation: if $E \to B$ is a fibered (\resp
cofibered) functor, then, for a morphism $m : i \to j$ in $B$, we
denote by $m^{\ast} : E_{j} \to E_i$ (\resp
$m_{\ast} : E_i \to E_j$) the associated ``inverse image'' (\resp
\pageoriginale ``direct image'') functor; each of these functors is
defined up to a unique functorial isomorphism. If
$\varphi : E \to E'$ is a $B$-functor, then, for every object $i$ of
$B$, we denote by $\varphi_i : E_i \to E'_i$ the restriction of
$\varphi$ to $E_i$.

\subsubsection{}\etiquette[1.1.3]{Vbis.subsubsec1.1.3}

Finally, we denote by $\Delta$ the following category: the objects
of $\Delta$ are the ordered sets $[n] = \{ 0, 1, \ldots, n \}$, and
the morphisms of $\Delta$ are all weakly increasing maps.
$\Delta^{+}$ (\resp $\Delta^{-}$) denotes the category whose objects
are those of $\Delta$ and whose arrows are the monomorphisms (\resp
the epimorphisms) of $\Delta$. We will freely use, as needed, the
classical notation introduced in connection with $\Delta$ [\cf
\cite{sga4.exp.Vbis.bib.3} II 2].

\subsection{D-topoi}\etiquette[1.2]{Vbis.subsec1.2}

\begin{enonce*}{Definition 1.2.1}\etiquette[1.2.1]{Vbis.def1.2.1}
Let $D$ be a category, and let $E \to D$ be a fibered and cofibered
functor. We say that $E$ is bifibered in topoi over $D$, or that $E$
is a $D$-topos, if the following conditions hold:
\begin{enumerate}
\renewcommand{\theenumi}{\alph{enumi}}
\renewcommand{\labelenumi}{(\theenumi)}
\item For every object $i$ of $D$, the fiber category $E_i$ is a
  topos.

\item For every morphism $m : i \to j$ in $D$, there exists a
  morphism of topoi $f : E_j \to E_i$ such that
  $f_{\ast} = m^{\ast}$ and $f^{\ast} = m_{\ast}$.
\end{enumerate}
\end{enonce*}

\begin{enonce*}[remark]{Remark}
In view of (a), condition (b) can equivalently be expressed by saying
that the functor $m_{\ast}$ is left exact\nde{Care must be taken here:
in a fibered category, $m_\star$ is left adjoint to $m^\star$ and is
therefore tautologically right exact; cf. [\SGA 1 VI.10]. This
explains the formulas $f_{\ast} = m^{\ast}$ and
$f^{\ast}=m_{\ast}$.}.
\end{enonce*}

When $D = \Delta$ (\resp $\Delta \times \Delta$), we speak of a
\emph{simplicial topos} (\resp a \emph{double simplicial topos}) to
mean a $\Delta$-topos (\resp a $\Delta \times \Delta$-topos).

In practice, we will encounter $D$-topoi by means of the following
considerations.

\begin{enonce*}{Definition 1.2.2}\etiquette[1.2.2]{Vbis.def1.2.2}
Let\pageoriginale $\mathcal{E}$ be a fibered and cofibered category
over $D$. We say that $\mathcal{E}$ is \emph{bifibered in duals of
topoi over $D$} if $\mathcal{E}^{\circ}$ is a $D^{\circ}$-topos.
\end{enonce*}

\hypertarget{sga4.exp.Vbis.remark.1.2.3}{}\begin{enonce*}[remark]{Remark 1.2.3}
Explicitly, $\mathcal{E}$ is bifibered in duals of topoi over $D$ if
and only if the following conditions hold:
\begin{enumerate}
\renewcommand{\theenumi}{\alph{enumi}}
\renewcommand{\labelenumi}{(\theenumi)}
\item For every object $i$ of $D$, the opposite category
  \sisi{}{${\mathcal{E}_i}^\circ$} of the fiber category
  $\mathcal{E}_i$ is a topos.

\item For every morphism $m : i \to j$ in $D$, the functor
  $m^{\ast}$ is left exact\nde{Cf. the preceding note. In other
  words, the family of functors
  $((m_{\ast})^\circ,(m^{\ast})^\circ)$ defines a compatible system
  of morphisms of topoi
  $\mathcal{E}_i^\circ\rightarrow \mathcal{E}_j^\circ$.}.
\end{enumerate}
\end{enonce*}

\setcounter{subsection}{2}
\setcounter{subsubsection}{3}
\subsubsection{}\etiquette[1.2.4]{Vbis.subsubsec1.2.4}
The reader will find examples of categories bifibered in duals of
topoi in \hyperref[Vbis.sec4]{paragraph 4}.

\subsubsection{}\etiquette[1.2.5]{Vbis.subsubsec1.2.5}
Let $\mathcal{E} \to B$ be a category bifibered in duals of topoi
over $B$, and let $X : D^{\circ} \to B$ be a functor. We leave it to
the reader to verify that
$(\fprod{D^{\circ}}{\mathcal{E}}{B})^{\circ}$ is a $D$-topos, which
we denote by $\bar{X}$.

\subsubsection{}\etiquette[1.2.6]{Vbis.subsubsec1.2.6}
We will often say that a functor $X : D^{\circ} \to B$ is a
$D$-\emph{object of} $B$, and we denote by $X_i$ the image under
this functor of an object $i$ of $D$; the $D$-objects of $B$ form a
category denoted by $D^{\circ}B$. If $S$ is an object of $B$, a
$D$-object of $B/S$ will be called a \emph{$D$-object of $B$
augmented toward} $S$. The datum of a $D$-object of $B$ augmented
toward $S$ is trivially equivalent to that of a functorial morphism
$X \to C^{D}_S$, where $X$ is a $D$-object of $B$ and $C^D_S$ is
the constant $D$-object of $B$ with value $S$.

When $D = \Delta$, we speak of a \emph{simplicial object} (\resp an
\emph{augmented simplicial object}).

We will also use \emph{double simplicial objects} (taking
$D = \Delta \times \Delta$).

\subsubsection{}\etiquette[1.2.7]{Vbis.subsubsec1.2.7}
Now suppose\pageoriginale that the category $B$ has finite fiber
products.

Let $f : R \to S$ be an arrow in $B$: the bifunctor
$$
([n], X) \rightsquigarrow \Hom_{\Ens} ([n], \Hom_S (X, R))
$$
from $\Delta^{\circ} \times (B / S)^{\circ}$ to $\Ens$ defines a
functor
$$
[n] \rightsquigarrow X_{[n]} = X_n
$$
by taking as $X_n$ a representative of the functor
$$
Z \rightsquigarrow \Hom_{\Ens} ([n], \Hom_S (Z, R))
$$
$(X_n \simeq \displaystyle\mathop{\overbrace{R \times_S R
    \times \ldots \times_S R}}^{n+1 \text{ times }})$.

We denote by $[R |_{f} S]$ or $[R | S]$ the semi-simplicial object
augmented toward $S$ thus constructed.

Finally, if $X$ and $X'$ are two semi-simplicial objects of $B$ (or
of $B/S$), and if $$u : X \to X'$$ is a functorial morphism, then,
for technical reasons, we introduce the object $$[X |_{u} X']$$
computed in the category $\Delta^{\circ}B$. This object can be
interpreted as a double simplicial object\nde{Explicitly, with the
usual ordering of the indices, one has
$$[[X |_u X']]_{(n,m)}=\displaystyle\mathop{\overbrace{X_m \times_{X'_m} X_m
    \times \ldots \times_{X'_m} X_m}}^{n+1 \text{ times }}$$
with the functorialities deduced from those on the factors.}
of $B$, which we then denote by $[[X |_u X']]$: indeed, there is a
canonical isomorphism of categories
$$
\Delta^{\circ} (\Delta^{\circ} B) \xrightarrow{\sim} (\Delta
\times \Delta)^{\circ} B{\quoi}.
$$

We now return to the general notion of a $D$-topos.

\subsubsection{}\etiquette[1.2.8]{Vbis.subsubsec1.2.8}
Let $E$ be a $D$-topos: we denote by $\Ysheaf{\Gamma}(E)$ the
category $\Hom_D (D, E)$. Let $f : D'\to D$ be a functor: the
category $\fprod{D'}{E}{D}$ is a $D'$-topos and, by composition with
$f$, we obtain a functor
$$
f^{\ast} : \Ysheaf{\Gamma} (E) \longrightarrow \Ysheaf{\Gamma} (\fprod{D'}{E}{D}){\quoi}.
$$

In\pageoriginale the case where $D'$ is reduced to an object $i$ of
$D$ (with the identity of $i$ as its only morphism), and where $f$
is the canonical inclusion denoted by $e_i$, we may take
$\fprod{D'}{E}{D}$ to be the fiber category $E_i$, and $e^{\ast}_i$ is then
identified with the ``\emph{evaluation at $i$}'' functor.

\begin{enonce*}{Proposition 1.2.9}\etiquette[1.2.9]{Vbis.prop1.2.9}
If $D'$ is a $\Uu$-small category, the functor $f^{\ast}$ has a
right adjoint and a left adjoint (denoted by $f_{\ast}$ and $f_{!}$,
respectively).
\end{enonce*}

This follows from a slight generalization of Kan's lemma [III 1.1],
which we will use repeatedly.

\begin{enonce*}{Lemma 1.2.10}\etiquette[1.2.10]{Vbis.lem1.2.10}
Let $I$, $J$, and $A$ be three categories over the same category
$B$. Suppose that $I$ is $\Uu$-small and that $A$ is fibered and
cofibered over $B$. Let a $B$-functor $f : I \to J$ be given, and
denote by $f^{\ast}$ the functor
$$
\Hom_B (J, A) \longrightarrow \Hom_B (I, A)
$$
defined by composition with $f$. If, in every fiber of $A$ over
$B$, $\Uu$-inductive limits (\resp projective limits) exist, then
$f^{\ast}$ has a left adjoint (\resp a right adjoint).
\end{enonce*}

\begin{enonce*}[remark]{Proof}
We give only the proof of the existence of the left adjoint, the
\resp part of the lemma following by duality. We use the following
fact, whose verification is left to the reader.
\end{enonce*}

\begin{enonce*}{Lemma 1.2.10.1}\etiquette[1.2.10.1]{Vbis.lem1.2.10.1}
Let $A$ be a bifibered category over a category $B$, let $b$ be an
object of $B$, and let $(F_{\lambda})_{\lambda \in \Lambda}$ be a
functor from a category $\Lambda$ to the fiber $A_b$. For every $G$
in $A_{b'}$, and every $u : b'\to b$ (\resp $u : b \to b'$) in $B$,
there is a bijection
$$
\Hom_u(G, \varprojlim F_{\lambda}) \xrightarrow{\sim}
\varprojlim \Hom_u (G, F_{\lambda})
$$
(\resp $\Hom_u (\varinjlim F_{\lambda}, G)
\xrightarrow{\sim}\varprojlim \Hom_u (F_{\lambda}, G)$)
whenever the first member is defined.
\end{enonce*}

\sisi{$+$}{} With this understood\pageoriginale, let
$I \amalg_f J$ be the category over $B$ defined by
$$
\ob (I \textstyle{\amalg}_f J) = \ob (I) \amalg \ob (J)
$$
\begin{equation*}
\Hom(x, y)=
\begin{cases}
\Hom_I(x, y) & \text{ if } x, y \in \ob (I)\\
\Hom_J(x, y) & \text{ if } x, y \in \ob (J)\\
\Hom_J(f(x), y) & \text{ if } x \in \ob(I) \text{ and } y \in \ob(J)\\
∅ & \text{ if } x \in \ob(J) \text{ and } y \in
\ob(I){\quoi}.\\
\end{cases}
\end{equation*}

\sisi{$+$}{} The category $\Hom_B(I \amalg_f J, A)$ is equivalent
to the category of triples consisting of a B-functor $F$ from $I$
to $A$, a $B$-functor $G$ from $J$ to $A$, and a morphism
$\varphi$ of $B$-functors from $F$ to
$G \circ f = f(G)$.

\sisi{$+$}{} For every object $j$ of $J$, denote by $I /j$ the
category of objects of $I$ ``placed by $f$ over $j$'': the objects
of $I /j$ are the pairs $(i, \alpha)$, where $i$ is an object of $I$
and $\alpha : i \to j$ is an arrow in $I \amalg_f J/j$; the
morphisms of $I/j$ are those of $I\amalg_{f}J/j$. If $p_I$ and $p_J$
are the projections of $I$ and $J$ onto $B$, and if $(i, \alpha)$ is
an object of $I/j$, denote by
$\alpha_{\ast} : A_{p_{I}(i)} \to A_{p_{J}(j)}$ the functor
$p_{J}(\alpha)_{\ast}$.

\sisi{$+$}{} To give a $B$-functor from $I \amalg_f J$ to $A$ is
equivalently to give $F \in \Hom_B (I, A)$,
$G \in \Hom_B (J, A)$, and a morphism, functorial in $j$,
$$
\psi_j : \varinjlim_{(i, \alpha)\in I/j} \alpha_{\ast}(F(i))
\longrightarrow G(j){\quoi}.
$$
(The functoriality in $j$ of the left-hand member follows from
\hyperref[Vbis.lem1.2.10.1]{1.2.10.1}.) The left adjoint of
$f^{\ast}$ is therefore given by the formula
$$
f_{!}(F) (j) = \varinjlim_{(i, \alpha)\in I/j}
\alpha_{\ast}(F(i)){\quoi}.
$$

\begin{enonce*}{Corollary 1.2.11}\etiquette[1.2.11]{Vbis.cor1.2.11}
Let\pageoriginale $E$ be a $D$-topos and let $i$ be an object of
$D$; the functor $e^{\ast}_i$ has a left adjoint (\resp a right
adjoint), given by
$$
e_{i!}(a)(j) = \coprod_{\alpha \in \Hom(i, j)} \alpha_{\ast}(a)
\left( \emph{\resp} e_{i\ast} (a) (j) = \prod_{\alpha \in \Hom(j,
  i)} \alpha^{\ast}(a)\right)
$$
where a is an object of $E$ over $i$.
\end{enonce*}

\begin{enonce*}{Proposition 1.2.12}\etiquette[1.2.12]{Vbis.prop1.2.12}
Let $D$ be a $\Uu$-small category and let $E$ be a $D$-topos. Then
the category $\Ysheaf{\Gamma}(E)$ is a $\Uu$-topos.
\end{enonce*}

\sisi{$+$}{} One verifies ``fiber by fiber,'' using
\hyperref[Vbis.lem1.2.10.1]{1.2.10.1}, that the category
$\Ysheaf{\Gamma}(E)$ has the following properties:
\begin{enumerate}
\renewcommand{\theenumi}{\alph{enumi}}
\renewcommand{\labelenumi}{\theenumi)}
\item Finite projective limits are representable.
\item Coproducts indexed by an element of $\Uu$ are representable.
  They are disjoint and universal.
\item Equivalence relations are universally effective.
\end{enumerate}

\sisi{$+$}{} It remains to show that $\Ysheaf{\Gamma}(E)$ has a
system of generators indexed by an element of $\Uu$. But if, for
every object $i$ of $D$, $(G_{i\lambda})_{\lambda \in \Lambda_{i}}$
is a system of generators of $E_i$ (where $\Lambda_i$ is a
$\Uu$-small set), then the family
$(e_{i!}(G_{i\lambda}))_{i, \lambda}$ is a system of generators of
$\Ysheaf{\Gamma}(E)$.

\setcounter{subsection}{2}
\setcounter{subsubsection}{12}

% END INCLUDED SGA4_English_current_buildable_source_20260906/source/05bis/drafts/SGA4_VBIS_LINES_206_516_English.texfrag


% BEGIN INCLUDED SGA4_English_current_buildable_source_20260906/source/05bis/drafts/SGA4_VBIS_LINES_517_829_English.texfrag
\subsubsection{}\etiquette[1.2.13]{Vbis.subsubsec1.2.13}
We now introduce the notion of a morphism between $D$-topoi. First,
let us specify that if $F$ and $F'$ are two categories over the same
category $B$ and $T:F\to F'$ is a $B$-functor, a
\emph{$B$-left adjoint} to $T$ means a functor $S:F'\to F$ left
adjoint to $T$ such that the canonical morphisms
$1\to T\circ S$ and $S\circ T\to1$ are $B$-morphisms of functors.
Under these conditions, one verifies immediately that $T$ is
cartesian and $S$ is cocartesian.

\begin{enonce*}{Definition 1.2.14}\etiquette[1.2.14]{Vbis.def1.2.14}
Let\pageoriginale $E$ and $E'$ be two $D$-topoi. A morphism from
$E$ to $E'$ is a pair of $D$-functors
$\Phi_{\ast}:E\to E'$ and $\Phi^{\ast}:E'\to E$, equipped with a
$D$-adjunction
$$
\Hom(\Phi^{\ast}(-),-)\xrightarrow[\xi]{\sim}
\Hom(-,\Phi_{\ast}(-)),
$$
such that, for every object $i$ of $D$, the pair
$(\Phi_{\ast i},\Phi_i^{\ast})$, equipped with the adjunction induced
by $\xi$, is a morphism of topoi from $E_i$ to $E'_i$.
\end{enonce*}

\begin{enonce*}{Proposition 1.2.15}\etiquette[1.2.15]{Vbis.prop1.2.15}
Let $E$ and $E'$ be two $D$-topoi, and let
$(\Phi_{\ast},\Phi^{\ast}):E\to E'$ be a morphism. If $D$ is a
$\Uu$-small category, then the pair
$(\Gamma(\Phi_{\ast}),\Gamma(\Phi^{\ast})):
\Ysheaf{\Gamma}(E)\to\Ysheaf{\Gamma}(E')$ is a morphism of topoi.
\end{enonce*}

This follows immediately from the preceding definition.

The following lemma, whose proof is left to the reader, provides a
way of constructing morphisms of $D$-topoi.

\begin{enonce*}{Lemma 1.2.16}\etiquette[1.2.16]{Vbis.lem1.2.16}
Let $E$ and $E'$ be two bifibered categories over the same category
$D$, and let $\Phi$ be a \emph{cartesian} $D$-functor from $E$ to
$E'$ such that, for every object $i$ of $D$,
$\Phi_i:E_i\to E'_i$ has a left adjoint. Then the choice, for every
$i$, of a left adjoint to $\Phi_i$ canonically determines a
$D$-functor $\Psi:E'\to E$ which is $D$-left adjoint to $\Phi$.
\end{enonce*}

\begin{enonce*}[remark]{Scholium 1.2.17}\etiquette[1.2.17]{Vbis.sch1.2.17}
Under the conditions of (\hyperref[Vbis.lem1.2.16]{1.2.16}), suppose
that $E$ and $E'$ are two $D$-topoi and that, for every object $i$
of $D$, every left adjoint to $\Phi_i$ is left exact. Then, if
$\Psi:E'\to E$ is a $D$-left adjoint to $\Phi$, the pair
$(\Phi,\Psi):E\to E'$ is a morphism of $D$-topoi.
\end{enonce*}

Now let $\mathcal E\to B$ be a category bifibered in duals of topoi
over $B$, let $X$ and $X'$ be two $D$-objects of $B$, and let
$\alpha:X\to X'$ be a functorial morphism. A choice of normalized
cleavages for $\mathcal E$ and $\mathcal E^{\circ}$ then canonically
determines a morphism
$(\alpha_{\ast},\alpha^{\ast}):\bar X\to\bar X'$ of $D$-topoi such
that, for every object\pageoriginale $i$ of $D$,
$(\alpha_{\ast i},\alpha_i^{\ast}):\bar X_i\to\bar X'_i$ is equal
to $(\alpha_{i\ast}^{\circ},\alpha_i^{\ast\circ})$, where
$\alpha_i:X_i\to X'_i$ is the arrow of $B$ given by $\alpha$. For
two different choices of cleavages for $\mathcal E$ and
$\mathcal E^{\circ}$, there is a unique $D$-isomorphism between the
resulting morphisms. (For the verification of these facts, the
reader may consult \cite{sga4.exp.Vbis.bib.4}, (1.17).)

\subsection{Ringed $D$-topoi}\etiquette[1.3]{Vbis.subsec1.3}

\begin{enonce*}{Definition 1.3.1}\etiquette[1.3.1]{Vbis.def1.3.1}
A ringed $D$-topos is a pair $(E,A)$, where $E$ is a $D$-topos and
$A$ is a ring of $\Ysheaf{\Gamma}(E)$.
\end{enonce*}

One then verifies that, for every object $i$ of $D$, $A_i$ is a ring
of the topos $E_i$, and that, for every morphism $m:i\to j$, the
canonical arrow $A_i\to m^{\ast}(A_j)$ is a ring homomorphism.

\begin{enonce*}{Definition 1.3.2}\etiquette[1.3.2]{Vbis.def1.3.2}
Let $(E,A)$ and $(E',A')$ be two ringed $D$-topoi. A morphism from
$(E,A)$ to $(E',A')$ is a pair $(\Phi,\theta)$, where
$\Phi:E\to E'$ is a morphism of $D$-topoi and
$\theta:A'\to\Gamma(\Phi_{\ast})(A)$ is a ring homomorphism.
\end{enonce*}

\begin{enonce*}[remark]{Remark 1.3.3}\etiquette[1.3.3]{Vbis.rem1.3.3}
A morphism $\Phi:(E,A)\to(E',A')$ of ringed $D$-topoi induces a
morphism
$(\Gamma(\Phi),\theta):(\Ysheaf{\Gamma}(E),A)\to
(\Ysheaf{\Gamma}(E'),A')$ of ringed $\Uu$-topoi when $D$ is a
$\Uu$-small category (\cf \hyperref[Vbis.prop1.2.15]{1.2.15}).
\end{enonce*}

\setcounter{subsubsection}{3}
\subsubsection{}\etiquette[1.3.4]{Vbis.subsubsec1.3.4}
Let $\mathcal E\to B$ be a category bifibered in duals of topoi over
$B$, and let $\mathcal O$ be a ring of
$\Ysheaf{\Gamma}(\mathcal E^{\circ})=
\Hom_{B^{\circ}}(B^{\circ},\mathcal E^{\circ})$. If
$X:D^{\circ}\to B$ is a $D$-object of $B$, the $D$-topos $\bar X$
(\cf \hyperref[Vbis.subsubsec1.2.5]{(1.2.5)}) is naturally ringed by
$(\mathcal O\mathbin{.}X^{\circ}):D\to\mathcal E$; the resulting
ringed $D$-topos is denoted by $(\bar X,\mathcal O)$. If
$\alpha:X\to X'$ is a functorial morphism, the morphism
$(\alpha_{\ast},\alpha^{\ast}):\bar X\to\bar X'$ (\cf
\hyperref[Vbis.sch1.2.17]{1.2.17}) canonically induces a morphism
$(\bar X,\mathcal O)\to(\bar X',\mathcal O)$ of ringed $D$-topoi,
again denoted by $\alpha$.

\subsubsection{}\etiquette[1.3.5]{Vbis.subsubsec1.3.5}
A ringed $D$-topos $(E,A)$ canonically defines a category
$\Ysheaf{\Mod}(E,A)$ bifibered in abelian categories over $D$, whose
fiber at an object $i$ of\pageoriginale $D$ is the category
$\Mod(E_i,A_i)$ of modules on the ringed topos $(E_i,A_i)$. With
this notation, the category of modules of $\Ysheaf{\Gamma}(E)$ over
$A$, denoted by $\Mod(\Ysheaf{\Gamma}(E),A)$, identifies with the
category $\Hom_D(D,\Ysheaf{\Mod}(E,A))$.

\subsubsection{}\etiquette[1.3.6]{Vbis.subsubsec1.3.6}
Let $\varphi=(\Phi,\theta):(E,A)\to(E',A')$ be a morphism of ringed
$D$-topoi. It defines two functors
$\varphi_{\ast}:\underline{\Mod}(E,A)\to
\underline{\Mod}(E',A')$ and
$\varphi^{\ast}:\underline{\Mod}(E',A')\to\Mod(E,A)$ between the
corresponding categories of modules:

\begin{itemize}
\item Let $M$ be an object of $\Ysheaf{\Mod}(E,A)$ over an object
$i$ of $D$. Then $\Phi_{\ast}(M)$ is a module over
$\Phi_{\ast}(A_i)$, and the morphism
$\theta_i:A'_i\to\Phi_{\ast}(A_i)$ makes it an $A'_i$-module,
denoted by $\varphi_{\ast}(M)$. The functor $\varphi_{\ast}$ is
called the \emph{direct-image functor under the morphism} $\varphi$.

\item Let $M'$ be an object of $\Ysheaf{\Mod}(E',A')$ over an object
$i$ of $D$. Then $\Phi^{\ast}(M')$ is a module over
$\Phi^{\ast}(A'_i)$, and
$\varphi^{\ast}(M')=\Phi^{\ast}(M')
\otimes_{\Phi^{\ast}(A'_i)}A_i$ is canonically endowed with an
$A_i$-module structure. By means of
(\hyperref[Vbis.lem1.2.16]{1.2.16}), this defines a functor
$\varphi^{\ast}$ left adjoint to $\varphi_{\ast}$, called the
\emph{inverse-image functor under the morphism} $\varphi$.
\end{itemize}

\paragraph{}\etiquette[1.3.6.1]{Vbis.par1.3.6.1}
We say that $\varphi$ is \emph{flat} if the functor
$\Gamma(\varphi^{\ast})$ is exact.

\begin{enonce*}{Proposition 1.3.7}\etiquette[1.3.7]{Vbis.prop1.3.7}
Let $f:D'\to D$ be a functor, and let $(E,A)$ be a ringed $D$-topos.
Then the canonical functor
$$
f^{\ast}:\Mod(\Ysheaf{\Gamma}(E),A)\longrightarrow
\Mod(\Ysheaf{\Gamma}(\fprod{E}{D'}{D}),A\circ f)
$$
has both a right and a left adjoint if $D'$ is a $\Uu$-small
category. In particular, it is exact.
\end{enonce*}

This follows immediately from (\hyperref[Vbis.lem1.2.10]{1.2.10})
and the identification
$\Hom_D(D,\Ysheaf{\Mod}(E,A))\simeq
\Mod(\Ysheaf{\Gamma}(E),A)$.

In accordance\pageoriginale with the general notation, we denote the
left (\resp right) adjoint to $f^{\ast}$ by $f_!$ (\resp
$f_{\ast}$).

\setcounter{subsubsection}{7}
\subsubsection{}\etiquette[1.3.8]{Vbis.subsubsec1.3.8}
The considerations below provide a convenient computational method
for derived functors of the form
$$
\sisi{\RR^{+}}{R}\Gamma(\varphi_{\ast}):
D^{+}(\Ysheaf{\Gamma}(E),A)\to
D^{+}(\Ysheaf{\Gamma}(E'),A'),
$$
where $\varphi:(E,A)\to(E',A')$ is a morphism of ringed $D$-topoi
(\cf (\hyperref[Vbis.subsubsec1.3.6]{1.3.6})). [If
$(S,\mathcal O_S)$ is a ringed topos, we denote by
$D^{+}(S,\mathcal O_S)$ the bounded-below derived
category\nde{If $f:(S,\mathcal O_S)\rightarrow(T,\mathcal O_T)$ is
a flat morphism of ringed topoi, we simply write $f^*$ (and not
$L^+(f^*)$) for the derived functor
$$f^*:D^+(T,\mathcal O_T)\rightarrow D^+(S,\mathcal O_S),$$
as in the original version.} of the category of
$\mathcal O_S$-modules on $S$.]

\subsubsection{}\etiquette[1.3.9]{Vbis.subsubsec1.3.9}
This formal calculation may moreover be applied in other contexts,
such as ``descent in $\ell$-adic cohomology'' (\cf \SGA 5).

\begin{enonce*}{Proposition 1.3.10}\etiquette[1.3.10]{Vbis.prop1.3.10}
Let $D$ be a $\Uu$-small category. If $(E,A)$ is a ringed
$D$-topos, denote by $I_{(E,A)}$ the class of objects of
$\Mod(\Ysheaf{\Gamma}(E),A)$ isomorphic to an object of the form
$\prod_{i\in\ob(D)}e_{i\ast}(Q_i)$, where, for every object $i$ of
$D$, $Q_i$ is totally acyclic \emph{(\cf V \hyperref[V.def4.1]{4.1})}\footnote{``Flasque''
in the terminology of V.}; $I_{(E,A)}$ has the following properties:
\begin{enumerate}
\renewcommand{\theenumi}{\roman{enumi}}
\renewcommand{\labelenumi}{(\theenumi)}
\item For every object $F$ of $I_{(E,A)}$ and every object $i$ of
$D$, $e_i^{\ast}(F)$ is totally acyclic.

\item Every object of $\Mod(\Ysheaf{\Gamma}(E),A)$ embeds into an
object of $I_{(E,A)}$.

\item $I_{(E,A)}$ is stable under finite direct sums.

\item For every morphism
$\varphi:(E,A)\to(E',A')$ of ringed $D$-topoi, the functor
$\Gamma(\varphi_{\ast})$ carries every acyclic complex in
$C^{+}(\Ysheaf{\Gamma}(E),A)$ formed of objects of $I_{(E,A)}$ to
an acyclic complex formed of objects of $I_{(E',A')}$.
\end{enumerate}
\end{enonce*}

\noindent{\underline{\emph{Proof}}}

\begin{enumerate}
\renewcommand{\theenumi}{\roman{enumi}}
\renewcommand{\labelenumi}{(\theenumi)}
\item In view of the explicit expression for $e_i^{\ast}$ (\cf
\hyperref[Vbis.cor1.2.11]{(1.2.1.1)}), it suffices to prove the following lemma:

\begin{enonce*}{Lemma 1.3.10.1}\etiquette[1.3.10.1]{Vbis.lem1.3.10.1}
Let\pageoriginale $(S,\mathcal O_S)$ be a ringed topos, and let
$(F_t)_{t\in T}$ be a family of totally acyclic
$\mathcal O_S$-modules indexed by a $\Uu$-small set $T$. Then
$\prod_{t\in T}F_t$ is totally acyclic.
\end{enonce*}

Let $X$ be an object of $S$. There is a spectral sequence
\begin{equation*}
H^p\left(X,{\mathop{\prod}_{t\in T}}^{(q)}F_t\right)Rightarrow
\mathop{\prod}_{t\in T}H^{p+q}(X,F_t)\tag{1.3.10.2}
\end{equation*}\etiquette[1.3.10.2]{Vbis.eq1.3.10.2}
where $\prod_{t\in T}^{(q)}$ denotes the $q$-th derived functor of
the ``product indexed by $T$'' functor. Since
$\prod_{t\in T}^{(q)}F_t$ is the sheaf associated with the presheaf
$U\mapsto\prod_{t\in T}H^q(U,F_t)$,
\eqhyperref[Vbis.eq1.3.10.2]{1.3.10.2} degenerates and gives
isomorphisms
\begin{equation*}
H^n\left(X,\prod_{t\in T}F_t\right)\xrightarrow{\sim}
\prod_{t\in T}H^n(X,F_t)\tag{1.3.10.3}
\end{equation*}\etiquette[1.3.10.3]{Vbis.eq1.3.10.3}
and hence, finally,
$H^n(X,\prod_{t\in T}F_t)=0$ for every $n>0$.

\item Let $F$ be an object of $\Mod(\Ysheaf{\Gamma}(E),A)$, and put
$F_i=e_i^{\ast}(F)$. For every $i$, choose a monomorphism
$F_i\to Q_i$, where $Q_i$ is
totally acyclic in $\Mod(E_i,A_i)$. Then
$e_{i\ast}(F_i)\to e_{i\ast}(Q_i)$ is a monomorphism (because
$e_{i\ast}$ has a left adjoint), and the canonical arrow
\begin{equation*}
F\to\prod_{i\in\Ob(D)}e_{i\ast}(F_i)\to
\prod_{i\in\Ob(D)}e_{i\ast}(Q_i)\tag{1.3.10.4}
\end{equation*}\etiquette[1.3.10.4]{Vbis.eq1.3.10.4}
is again a monomorphism, as is immediately verified.

\item The proof is left to the reader.

\item The reader may also verify that $\Gamma(\varphi_{\ast})$
carries an object of $I_{(E,A)}$ to an object of $I_{(E',A')}$. Once
this point is established, in view of (i), it suffices to verify the
following lemma:
\end{enumerate}

\begin{enonce*}{Lemma 1.3.10.5}\etiquette[1.3.10.5]{Vbis.lem1.3.10.5}
If\pageoriginale $0\to F\to G\to H\to0$ is a short exact sequence
in $\Mod(\Ysheaf{\Gamma}(E),A)$, with $F,G\in I_{(E,A)}$, then the
sequence
$$
0\to\Gamma(\varphi_{\ast})(F)\to
\Gamma(\varphi_{\ast})(G)\to
\Gamma(\varphi_{\ast})(H)\to0
$$
is exact in $\Mod(\Ysheaf{\Gamma}(E'),A')$.
\end{enonce*}

It suffices to observe that $e_i^{\ast}(H)$ is acyclic for every
object $i$ of $D$, and that $\Gamma(\varphi_{\ast})$ is computed
fiber by fiber. This completes the proof of
(\hyperref[Vbis.prop1.3.10]{1.3.10}).

\begin{enonce*}{Corollary 1.3.11}\etiquette[1.3.11]{Vbis.cor1.3.11}
Let $\varphi:(E,A)\to(E',A')$ be a morphism of ringed $D$-topoi.
The functor $\sisi{\RR^+}{R}\Gamma(\varphi_{\ast})$ can be computed
by means of resolutions formed of objects of $I_{(E,A)}$.
\end{enonce*}

Let $L\subset K^{+}(\Ysheaf{\Gamma}(E),A)$ be the full subcategory
of closed complexes of objects of $I_{(E,A)}$. By
((\hyperref[Vbis.prop1.3.10]{1.3.10}), (ii)), $L$ is a triangulated
subcategory, and the theorem (\hyperref[Vbis.subsec5.1]{5.1}) of
[\cite{sga4.exp.Vbis.bib.7}, Chapter 1] applies to it.

\begin{enonce*}{Corollary 1.3.12}\etiquette[1.3.12]{Vbis.cor1.3.12}
Let $\varphi:(E,A)\to(E',A')$ be a morphism of ringed $D$-topoi.
For every object $i$ of $D$, denote by
$\varphi_i:(E_i,A_i)\to(E'_i,A'_i)$ the induced morphism of ringed
topoi over $i$. Then the diagram
\[
\xymatrix@C=2.75cm@R=1.75cm{
D^{+}(\Gamma(E),A) \ar[r]^{\sisi{\RR^{+}(e_i^{\ast})}{e_i^{\ast}}}
\ar[d]_{\sisi{\RR^{+}}{R}\Gamma(\varphi_{\ast})}&
D^{+}(E_i,A_i) \ar[d]^{\sisi{\RR^{+}}{R}\varphi_{i\ast}}\\
D^{+}(\Gamma(E'),A')
\ar[r]_{\sisi{\RR^{+}(e_i^{\ast})}{e_i^{\ast}}}&D^{+}(E'_i,A'_i)
}
\]
is essentially commutative.
\end{enonce*}

There is a morphism
$$
\sisi{\RR^{+}(e_i^{\ast})}{e_i^{\ast}}\circ
\sisi{\RR^{+}}{R}\Gamma(\varphi_{\ast})\to
\sisi{\RR^{+}}{R}\varphi_{i\ast}\circ
\sisi{\RR^{+}(e_i^{\ast})}{e_i^{\ast}}
$$
which is verified to be an isomorphism by
(\hyperref[Vbis.prop1.3.10]{1.3.10}).

\begin{enonce*}[remark]{Remark 1.3.13}\etiquette[1.3.13]{Vbis.rem1.3.13}
Denote\pageoriginale by $\mathcal A_{(E,A)}$ the class of objects
$F$ of $\Mod(\Ysheaf{\Gamma}(E),A)$ such that $e_i^{\ast}(F)$ is
totally acyclic for every object $i$ of $D$. It follows from the
proof of (\hyperref[Vbis.prop1.3.10]{1.3.10}) that
$\sisi{\RR^{+}}{R}\Gamma(\varphi_{\ast})$ can be computed by means
of resolutions formed of objects of $\mathcal A_{(E,A)}$; this is
what is done, in particular, in the introduction by using the
``canonical flasque resolution''\footnote{See Exposé XVII for the
generalization of this notion.}. By
[(\hyperref[Vbis.prop1.3.10]{1.3.10})(i)], one has
$I_{(E,A)}\subset\mathcal A_{(E,A)}$, but we will use $I_{(E,A)}$
explicitly in \hyperref[Vbis.sec2]{Section 2}.
\end{enonce*}

% END INCLUDED SGA4_English_current_buildable_source_20260906/source/05bis/drafts/SGA4_VBIS_LINES_517_829_English.texfrag


% BEGIN INCLUDED SGA4_English_current_buildable_source_20260906/source/05bis/drafts/SGA4_VBIS_LINES_830_1073_English.texfrag
\section{The method of cohomological descent}\etiquette[2]{Vbis.sec2}

In this section, $D$ is a $\Uu$-small category.

\subsection{Generalities. Notation}\etiquette[2.1]{Vbis.subsec2.1}

\subsubsection{}\etiquette[2.1.1]{Vbis.subsubsec2.1.1}
Let $(S,\mathcal{O}_S)$ be a ringed topos. The category $S\times D$,
equipped with the canonical projection $S\times D\to D$, is a
$D$-topos; moreover, the constant section with value $\mathcal{O}_S$
defines a ringed $D$-topos $(S\times D,\mathcal{O}_S)$, called the
\emph{constant ringed $D$-topos with value} $(S,\mathcal{O}_S)$. With
this notation, the category
$\Mod(\Ysheaf{\Gamma}(S\times D),\mathcal{O}_S)$ is identified with
the category of covariant functors from $D$ to
$\Mod(S,\mathcal{O}_S)$.

Define an exact functor
$$
\epsilon^{\ast}:\Mod(S,\mathcal{O}_S)\to
\Mod(\Ysheaf{\Gamma}(S\times D),\mathcal{O}_S)
$$
by associating with every module $F$ on $S$ the constant functor with
value $F$. The functor $\epsilon^{\ast}$ has a right adjoint
$\epsilon_\ast$, which associates with every functor
$H:D\to\Mod(\mathcal{O}_S)$ its projective limit; the adjunction
morphism $F\to\epsilon_\ast\epsilon^{\ast}(F)$ is the morphism from
$F$ to the projective limit of the constant functor with value $F$.
Thus $\epsilon^{\ast}$ is fully faithful if and only if $D$ is
connected.

\begin{enonce*}{Definition 2.1.2}\etiquette[2.1.2]{Vbis.def2.1.2}
Let\pageoriginale $(E,A)$ be a ringed $D$-topos. An augmentation of
$(E,A)$ is a morphism (of ringed $D$-topoi) from $(E,A)$ to a constant
ringed $D$-topos. A ringed $D$-topos equipped with an augmentation will
be called an augmented ringed $D$-topos.
\end{enonce*}

\setcounter{subsubsection}{2}
\subsubsection{}\etiquette[2.1.3]{Vbis.subsubsec2.1.3}
Let $\mathcal{E}\to B$ be a category bifibered in duals of topoi, and
let $\mathcal{O}$ be a ring in
$\Ysheaf{\Gamma}(\mathcal{E}^{\circ})=
\Hom_{B^{\circ}}(B^{\circ},\mathcal{E}^{\circ})$. Let $S$ be an
object of $B$. A functor $D^{\circ}\to B/S$, that is, a functorial
morphism $X\xrightarrow{\theta}C^D_S$, where $X$ is a functor
$D^{\circ}\to B$ (\cf (\hyperref[Vbis.subsubsec1.2.6]{1.2.6})),
induces an augmentation
$$
\displaylines{\hfill
\theta:(\bar{X},\mathcal{O})\longrightarrow
(\mathcal{E}^{\circ}_S\times D,\mathcal{O}_S)
\hfill\llap{(\cf \hyperref[Vbis.subsubsec1.3.4]{(1.3.4)}){\quoi}.}}
$$
which we shall denote by the same letter as the functorial morphism
from which it arises.


\subsection{Cohomological descent}\etiquette[2.2]{Vbis.subsec2.2}

\subsubsection{}\etiquette[2.2.1]{Vbis.subsubsec2.2.1}
Let $\theta:(E,A)\to(S\times D,\mathcal{O}_S)$ be an augmented ringed
$D$-topos. With the notation of (\hyperref[Vbis.subsubsec1.3.6]{1.3.6})
and (\hyperref[Vbis.subsubsec2.1.1]{2.1.1}), set
$$
\bar{\theta}^{\ast}=\Gamma(\theta^{\ast})\circ\epsilon^{\ast}:
\Mod(S,\mathcal{O}_S)\to\Mod(\Ysheaf{\Gamma}(E),A)
$$
and
$\bar{\theta}_\ast=\epsilon_\ast\circ\Gamma(\theta_\ast):
\Mod(\Ysheaf{\Gamma}(E),A)\to\Mod(S,\mathcal{O}_S)$.

The image of $\bar{\theta}^{\ast}$ lies in the full subcategory of
$\Mod(\Ysheaf{\Gamma}(E),A)$ consisting of the cocartesian sections
of $\Ysheaf{\Gamma}(D,\Ysheaf{\Mod}(E,A))$. We shall denote this
category by
$\Ysheaf{\Gamma}^{\text{cocart}}(D,\Ysheaf{\Mod}(E,A))$
\nde{In other words, an augmented ringed $D$-topos is the datum of a
compatible system of morphisms of $S$-topoi
$$(m^\ast,m_\ast):E_j\longrightarrow E_i,$$
where $m\in\Hom_D(i,j)$. The category of cocartesian sections is the
category of descent data; its objects are families
$F_i\in\text{Ob}(E_i)$ equipped with cocycles of isomorphisms
$$F_j\isomto m^\ast F_i.$$
The functor which associates with every object $F$ on $S$ the constant
functor with value $F$ is the inverse image of a morphism of topoi
$$\epsilon:\Gamma(S\times D)\longrightarrow S,$$
inducing the morphism of ringed topoi of
\hyperref[Vbis.subsubsec2.1.3]{2.1.3}. The pair
$(\bar\theta_\ast,\bar\theta^\ast)$ is therefore simply induced by the
composite of morphisms of topoi
$\bar\theta=\epsilon\circ\Gamma(\theta):\Gamma(E)\to S$
(cf. \hyperref[Vbis.prop1.2.15]{1.2.15}).}

\begin{enonce*}{Definition 2.2.2}\etiquette[2.2.2]{Vbis.def2.2.2}
The morphism $\theta$ is said to be an \emph{effective descent
augmentation} if
$$
\sisi{}{\bar\theta^{\ast}=}
\Gamma(\theta^{\ast})\circ\epsilon^{\ast}:
\Mod(S,\mathcal{O}_S)\to
\Ysheaf{\Gamma}^{\text{\rm cocart}}(D,\Ysheaf{\Mod}(E,A))
$$
is an equivalence of categories.
\end{enonce*}

\setcounter{subsubsection}{2}
\subsubsection{}\etiquette[2.2.3]{Vbis.subsubsec2.2.3}
With\pageoriginale the notation of
(\hyperref[Vbis.subsubsec2.2.1]{2.2.1}), suppose that $\theta$ is
flat, so that $\bar\theta^{\ast}$ is exact and passes trivially to the
derived categories; \sisi{let $\LL^{+}(\bar\theta^{\ast})$}{we again
denote by $\bar\theta^{\ast}$} the functor thereby obtained. With this
notation:

\begin{enonce*}{Lemma 2.2.3.1}\etiquette[2.2.3.1]{Vbis.lem2.2.3.1}
There are two functorial morphisms
$$
\alpha:\id_{D^{+}(S,\mathcal{O}_S)}\to
\sisi{\RR^{+}}{R}\bar\theta_\ast\circ
\sisi{\LL^{+}(\bar\theta^{\ast})}{\bar\theta^{\ast}}
\quad\mbox{and}\quad
\beta:\sisi{\LL^{+}(\bar\theta^{\ast})}{\bar\theta^{\ast}}\circ
\sisi{\RR^{+}}{R}\bar\theta_\ast\to
\id_{D^{+}(\Ysheaf{\Gamma}(E),A)}
$$
which make these two functors adjoint.
\end{enonce*}

(\cf \cite{sga4.exp.Vbis.bib.9} (3.3)).

\begin{enonce*}{Definition 2.2.4}\etiquette[2.2.4]{Vbis.def2.2.4}
The morphism $\theta$ is said to be a \emph{1-cohomological-descent
augmentation} if
\begin{enumerate}
\renewcommand{\theenumi}{\arabic{enumi}}
\renewcommand{\labelenumi}{\theenumi$^{\circ}$)}
\item $\theta$ is flat;
\item
\sisi{$\LL^{+}(\bar\theta^{\ast})$}
{$\bar\theta^{\ast}:D^{+}(S,\mathcal{O}_S)\to
D^{+}(\Ysheaf{\Gamma}(E),A)$} is fully faithful.
\end{enumerate}
\end{enonce*}

\begin{enonce*}[remark]{Remark 2.2.5}\etiquette[2.2.5]{Vbis.rem2.2.5}
Condition $2^{\circ}$) of the preceding definition can also be
expressed by saying that the morphism
$\alpha\sisi{}{:\id_{D^{+}(S,\mathcal{O}_S)}\to
R\bar\theta_\ast\circ\bar\theta^{\ast}}$ in
(\hyperref[Vbis.lem2.2.3.1]{2.2.3.1}) is an isomorphism.
\sisi{}{This is what one proves in practice and what is useful.}
\end{enonce*}

\begin{enonce*}{Definition 2.2.6}\etiquette[2.2.6]{Vbis.def2.2.6}
The morphism $\theta$ is said to be a \emph{2-cohomological-descent
augmentation} (or an \emph{effective cohomological-descent
augmentation}) if $\theta$ is both an effective descent augmentation
and a 1-cohomological-descent augmentation.
\end{enonce*}

The preceding terminology is justified by the following result:

\begin{enonce*}{Proposition 2.2.7}\etiquette[2.2.7]{Vbis.prop2.2.7}
Let $\theta$ be an effective cohomological-descent augmentation. Then
the essential image of
$\sisi{\LL^{+}(\bar\theta^{\ast})}{\bar\theta^{\ast}}$ is the full
subcategory of $D^{+}(\Ysheaf{\Gamma}(E),A)$ consisting of the
complexes $F^{\bigdot}$ such that, for every $i$,
$H^{i}(F^{\bigdot})$ is a cocartesian section of
$\Ysheaf{\Mod}(E,A)$.
\end{enonce*}

For brevity, we shall call a complex $F^{\bigdot}$ satisfying the
preceding conditions a \emph{cohomological descent datum}.

\sisi{}{\begin{proof}}
It is clear that, if $K^{\bigdot}\in D^{+}(S,\mathcal{O}_S)$,
\sisi{}{then}
$\sisi{\LL^+(\bar\theta^{\ast})}
{\bar\theta^{\ast}(K^{\bigdot})}$ is a cohomological descent datum. It
is therefore enough to show that, for every cohomological descent
datum $F^{\bigdot}$, the canonical morphism
$$
\beta(F^{\bigdot}):
\sisi{\LL^{+}(\bar\theta^{\ast})}{\bar\theta^{\ast}}\circ
\sisi{\RR^{+}(\bar\theta_\ast)}{R\bar\theta_\ast}(F^{\bigdot})
\to F^{\bigdot}
$$
is\pageoriginale an isomorphism in
$D^{+}(\Ysheaf{\Gamma}(E),A)$.

\begin{enumerate}
\renewcommand{\theenumi}{\alph{enumi}}
\renewcommand{\labelenumi}{\theenumi)}
\item \emph{The case where $F^{\bigdot}$ is bounded}: we argue by
induction on the length $\ell$ of the interval of integers $i$ for
which $H^{i}(F^{\bigdot})\neq0$:
\begin{itemize}
\item for $\ell\leq1$, the assertion is true because $\theta$ is of
effective descent.

\item suppose $\ell>1$, and let $n$ be the greatest integer such that
$H^n(F^{\bigdot})\neq0$. There is a distinguished
triangle\footnote{For every complex $K^{\bigdot}$,
$K^{\bigdot}[-n]$ denotes the complex $T^{-n}(K^{\bigdot})$, where
$T$ is the translation functor.}
\[
\xymatrix@R=1.5cm{
          & H^{n}(F^{\bigdot})[-n] \ar[ld]&\\
F'^{\bigdot} \ar[rr] & & F^{\bigdot} \ar[lu]
}
\]
such that the induction hypothesis applies to $F'^{\bigdot}$. Since
$\beta(H^{n}(F^{\bigdot})[-n])$ and $\beta(F'^{\bigdot})$ are
isomorphisms, the same is true of $\beta(F^{\bigdot})$.
\end{itemize}

\item \emph{The general case}: denote by
$\sisi{\sigma}{\tau}_{\leq n}(F^{\bigdot})$ the complex
$$
\ldots\to F^{n-2}\to F^{n-1}\to\Ker d^n\to0\to0\ldots{\quoi},
$$
(\cf \cite{sga4.exp.Vbis.bib.7} (7.1)), so that, for every $n$, there is a commutative
diagram
\[
\xymatrix@R=1.5cm@C=2cm{
\bar{\theta}^{\ast}\circ R\bar{\theta}_{\ast}
(\tau_{\leq n}(F^{\bigdot}))
\ar[r]^{\sim}\ar[d]&
\tau_{\leq n}(F^{\bigdot})\ar[d]\\
\bar{\theta}^{\ast}\circ R\bar{\theta}_{\ast}
(F^{\bigdot})\ar[r]&F^{\bigdot}
}
\]
and it is enough to see that the morphism
$$
\sisi{\RR^{+}(\bar\theta_\ast)}{R\bar\theta_\ast}
(\sisi{\sigma}{\tau}_{\leq n}(F^{\bigdot}))
\to
\sisi{\RR^{+}(\bar\theta_\ast)}{R\bar\theta_\ast}(F^{\bigdot})
$$
induces\pageoriginale an isomorphism on $H^{i}$ for $i\leq n$. But
there is a triangle
\[
\xymatrix{
&F''^{\bigdot}\ar[ld]&\\
\tau_{\leq n}(F^{\bigdot})\ar[rr]&&
F^{\bigdot}\ar[lu]
}
\]
with $H^i({F''}^{\bigdot})=0$ for $i\leq n$. It follows that
$H^{i}(\sisi{\RR^{+}(\bar\theta_\ast)}{R\bar\theta_\ast}
(F''^{\bigdot}))=0$ for $i\leq n$, and then that
$H^{i}(\sisi{\RR^{+}(\bar\theta_\ast)}{R\bar\theta_\ast}
(\sisi{\sigma}{\tau}_{\leq n}(F^{\bigdot})))\to
H^{i}(\sisi{\RR^{+}(\bar\theta_\ast)}{R\bar\theta_\ast}
(F^{\bigdot}))$ is an isomorphism for $i\leq n$, which completes the
proof.
\end{enumerate}

\sisi{}{\end{proof}}

We shall now present an explicit method for computing
$$
\sisi{\RR^{+}}{R}\bar\theta_\ast
\sisi{}{=R\epsilon_\ast\circ R\Gamma(\theta)_\ast},
$$
which is very useful in proving certain criteria for cohomological
descent and in exploiting that notion.

\sisi{\subsection{A method for computing
$R^+\epsilon_\ast$}\etiquette[2.3]{Vbis.subsec2.3}}
{\subsection{A method for computing
$R\epsilon_\ast$}\etiquette[2.3]{Vbis.subsec2.3}}

We shall begin with two sorites.

% END INCLUDED SGA4_English_current_buildable_source_20260906/source/05bis/drafts/SGA4_VBIS_LINES_830_1073_English.texfrag


% BEGIN INCLUDED SGA4_English_current_buildable_source_20260906/source/05bis/drafts/SGA4_VBIS_LINES_1074_1378_English.texfrag
\subsubsection{}\etiquette[2.3.1]{Vbis.subsubsec2.3.1}
Let $D$ be a $\Uu$-small category, and let $\Ab$ be the category of
abelian groups belonging to $\Uu$. Denote by $r_D$ the functor
$D\to(\Hom(D,\Ab))^{\circ}$ defined by
$$
r_D(i)(j)=\ZZ^{\Hom(i,j)}
$$
for every pair $(i,j)$ of objects of $D$. For every pair $(i,X)$
consisting of an object of $D$ and an object of $\Hom(D,\Ab)$, one
then constructs a canonical and functorial isomorphism
$$
\Hom(r_D(i),X)\xrightarrow{\sim}X(i).
$$

\paragraph{}\etiquette[2.3.1.1]{Vbis.par2.3.1.1}
Let $A$ be an additive category (\cf \sisi{Tohokû}{Tôhoku}), and let
$F$ be a functor $D\to A$. Define a functor
$\bar F:(\Hom(D,\Ab))^{\circ}\to\Hom(A^{\circ},\Ens)$ by
$$
\bar F(X)(a)=\Hom_{\Hom(D,\Ab)}
(X,\Hom_A(a,F(.))),
$$
so\pageoriginale that the following diagram, in which $h_A$ denotes
the canonical functor, is essentially commutative:
\[
\xymatrix@C=3cm@R=1.5cm{
D \ar[r]^{r_D}\ar[d]_{F}&(\Hom(D,\Ab))^{\circ}\ar[d]^{\bar F}\\
A \ar[r]_{h_A}&\Hom(A^{\circ},\Ens).
}
\]

We leave it to the reader to verify that $\bar F$ sends finite direct
sums in $\Hom(D,\Ab)$ to products and that, if $\ZZ$ denotes the
constant functor $D\to\Ab$ with value $\ZZ$, then
$\bar F(\ZZ)=\varprojlim F$.

Finally, the correspondence $F\rightsquigarrow\bar F$ defines a
covariant functor
$$
\Hom(D,A)\to\Hom((\Hom(D,\Ab))^{\circ},
\Hom(A^{\circ},\Ens)){\quoi}.
$$

\paragraph{}\etiquette[2.3.1.2]{Vbis.par2.3.1.2}
Denote by $\Add(D)$ the full subcategory of
$(\Hom(D,\Ab))^{\circ}$ defined by the objects of the form $r_D(i)$,
where $i$ is an object of $D$, and their finite direct sums. The
category $\Add(D)$ is additive, and the functor
$r_D:D\to\Add(D)$ has the following universal property:

\paragraph{}\etiquette[2.3.1.3]{Vbis.par2.3.1.3}
\emph{For every additive category $A$, the functor
$G\longmapsto G\circ r_D$ induces an equivalence from the full
subcategory of $\Hom(\Add(D),A)$ consisting of the additive functors}
onto \emph{the category $\Hom(D,A)$.}

The category $\Add(D)$, characterized up to equivalence by
(\hyperref[Vbis.par2.3.1.3]{2.3.1.3}), will be called the
\emph{additive category generated by $D$}.

%% compilation problem!

%\nde{In more concrete terms, used implicitly below, the objects of
%$\Add(D)$ are the finite sequences
%$$[i_1]\oplus\cdots\oplus[i_n]:=(i_1,\ldots,i_n)$$
%of objects of $D$, with morphisms in $\Add(D)$ defined by the rules
%\begin{itemize}
%  \item the empty sequence is the zero object;
%  \item for all $i,j\in\text{ob}(D)$,
%      $\Hom([i],[j])=\ZZ^{\Hom_D(i,j)}$;
%  \item for all finite families $(i),(j)$ of objects of $D$,
%      $\Hom(\oplus[i],\oplus[j])=\prod\Hom([i],[j])$,
%\end{itemize}
%while the composition rules are deduced from those of $D$ by
%bilinearity. With this presentation, $r_D$ is the tautological
%functor $i\rightsquigarrow[i]$; and, if $F$ is a functor from $D$ to
%an additive category $A$, it extends uniquely to an additive functor
%$$\widetilde F:\Add(D)\longrightarrow A$$
%defined by $\widetilde F(\oplus[i])=\oplus F_i$, with the evident
%functoriality rules.}

\subsubsection{}\etiquette[2.3.2]{Vbis.subsubsec2.3.2}
Let $A$ be an abelian category, and let $K^{\bigdot\bigdot}$ be a
double complex, regarded as an object of $C^{+}(C^+(A))$, with the
first index corresponding to the first\pageoriginale symbol $C^+$.
Denote by $(K^{\bigdot\bigdot})_s$ the \emph{associated simple
complex} (\cf \sisi{Tohokû}{Tôhoku}
(\hyperref[Vbis.subsec2.4]{2.4}), whose notation we shall retain).
This defines a functor
\begin{equation*}
(\quad)_s:C^+(C^+(A))\to C^+(A)\tag{2.3.2.1}
\end{equation*}\etiquette[2.3.2.1]{Vbis.eq2.3.2.1}
and we leave it to the reader to verify the following lemma:

\begin{enonce*}{Lemma 2.3.2.2}\etiquette[2.3.2.2]{Vbis.lem2.3.2.2}
The functor $(\quad)_s$ defines a triangulated functor
$$
K^+(C^+(A))\to K^+(A)
$$
which preserves quasi-isomorphisms; it therefore defines a
triangulated functor
$$
D^+(C^+(A))\to D^+(A){\quoi},
$$
again denoted by $(\quad)_s$.
\end{enonce*}

\begin{enonce*}[remark]{Remark}
The same result holds for the functor
$(\quad)_s:C(C^b(A))\to C(A)$, since the spectral sequence under
consideration is biregular by (\EGA III (11.3.3)).
\end{enonce*}

\subsubsection{}\etiquette[2.3.3]{Vbis.subsubsec2.3.3}
Now let $(S,\mathcal{O}_S)$ be a ringed topos, let $D$ be a
$\Uu$-small category, and let
$Z^{\bigdot}\in C^{+}(\Add(D))$ be a cochain complex such that
$Z^n=0$ for $n<0$.

\sisi{Every object $F$ of
$\Mod(\Ysheaf{\Gamma}(S\times D),\mathcal{O}_S)$ (which is identified
with a functor $D\to\Mod(S,\mathcal{O}_S)$ by
(\hyperref[Vbis.subsubsec2.1.1]{2.1.1})) defines, by the universal
property of $\Add(D)$, an object
$\varepsilon_{Z^{\bigdot}\ast}(F)^{\bigdot}$ of
$C^{+}(S,\mathcal{O}_S)$ which varies functorially with $F$ by
(\hyperref[Vbis.par2.3.1.3]{2.3.1.3})}
{Every object $F$ of
$\Mod(\Ysheaf{\Gamma}(S\times D),\mathcal{O}_S)$ (which is identified
with a functor $D\to\Mod(S,\mathcal{O}_S)$ by
(\hyperref[Vbis.subsubsec2.1.1]{2.1.1})) defines, by the universal
property of $\Add(D)$ (\hyperref[Vbis.par2.3.1.3]{2.3.1.3}), an
additive functor
$$\widetilde F:\Add(D)\longrightarrow\Mod(S,\mathcal{O}_S).$$
The image complex
$$\varepsilon_{Z^{\bigdot}\ast}(F):=\widetilde F(Z^{\bigdot})$$
is an object of $C^{+}(S,\mathcal{O}_S)$ which varies functorially
with $F$. We have thus defined an additive functor
$$
\varepsilon_{Z^{\bigdot}\ast}:
\Mod(\Ysheaf{\Gamma}(S\times D),\mathcal{O}_S)
\longrightarrow C^{+}(S,\mathcal{O}_S),
$$
which is visibly exact}.

This defines a triangulated functor
\nde{$K^+$ denotes the category of bounded-below complexes up to
homotopy.}:
\begin{equation*}
K^{+}(\varepsilon_{Z^{\bigdot}\ast}):
K^+(\Ysheaf{\Gamma}(S\times D),\mathcal{O}_S)
\to K^+(C^+(S,\mathcal{O}_S)){\quoi}.\tag{2.3.3.1}
\end{equation*}\etiquette[2.3.3.1]{Vbis.eq2.3.3.1}
which preserves quasi-isomorphisms. Moreover, it follows from
(\hyperref[Vbis.lem2.3.2.2]{2.3.2.2}) that the composite functor
$$
(\quad)_s\circ K^{+}(\varepsilon_{Z^{\bigdot}\ast}):
K^+(\Ysheaf{\Gamma}(S\times D),\mathcal{O}_S)
\to K^{+}(S,\mathcal{O}_S)
$$
sends acyclic objects to acyclic objects. It therefore defines a
triangulated\pageoriginale functor
\begin{equation*}
\sisi{\RR^+}{R}\varepsilon_{Z^{\bigdot}\ast}:
D^+(\Ysheaf{\Gamma}(S\times D),\mathcal{O}_S)
\to D^{+}(S,\mathcal{O}_S)\tag{2.3.3.2}
\end{equation*}\etiquette[2.3.3.2]{Vbis.eq2.3.3.2}
satisfying the relation
\nde{In other words, $\varepsilon_{Z^{\bigdot}\ast}$ derives
trivially as
$F^{\bigdot}\rightsquigarrow
\varepsilon_{Z^{\bigdot}\ast}(F^{\bigdot})$.}
$Q\circ(\quad)_s\circ K^+(\varepsilon_{Z^{\bigdot}\ast})
=\sisi{\RR^{+}}{R}\varepsilon_{Z^{\bigdot}\ast}\circ Q$
\footnote{For every abelian category $A$, the letter $Q$ denotes the
canonical functor $K^{+}(A)\to D^{+}(A)$.}

\subsubsection{}\etiquette[2.3.4]{Vbis.subsubsec2.3.4}
Retain the preceding data\sisi{.}{ and set
$A=C^+(S,\mathcal{O}_S)$.
%%%%%%%
%%%%%%
First, $\widetilde F$ and $\bar F$ are compatible in the sense of
commutativity of the Yoneda-type diagram
$$
\xymatrix{
&A\ar[r]^h&\Hom(A^\circ,\Ens)\\
D\ar[rd]_{{r_D}}\ar[r]\ar[ru]^{F}&
\Add(D)\ar[d]_{\widetilde{r_D}}\ar[u]_{\widetilde F}\\
&\Hom(D,\Ab)\ar[ruu]_{\bar F}
}
$$
where $h$ and $r_D$ are respectively the Yoneda-type arrows defined
by
$$
h(a)=(\alpha\longmapsto\Hom_A(\alpha,a))
\quad\text{and}\quad
r_D(d)\stackrel{\hyperref[Vbis.subsubsec2.3.1]{2.3.1}}{=}
(\delta\longmapsto\Hom_D(d,\delta)),
$$
that is, $h\circ\widetilde F=\bar F\circ\widetilde{r_D}$. It follows
that
$$
h(\varepsilon_{Z^{\bigdot}\ast}(F))
=\bar F(\widetilde{r_D}(Z^{\bigdot})).
$$

By the definition of a projective limit,
$h(\epsilon_\ast F)=\varprojlim F$, and, as already observed, one has
tautologically $\bar F(\ZZ)=\varprojlim F$; hence
$$\bar F(\ZZ)=h(\epsilon_\ast(F)).$$
%%%%%%
%%%%%%%
}


Regard $Z^{\bigdot}$ as a chain complex in $\Hom(D,\Ab)$
\sisi{.}{, that is, as a complex of the abelian category
$\Hom(D,\Ab)^{\circ}$ by means of the functor $\widetilde r_D$. We
may therefore speak of exact sequences and resolutions. Explicitly, a
sequence in $\Add(D)$
$$Z\longrightarrow Z'\longrightarrow Z''$$
is exact if, for every object $j$ of $D$, the sequence
$$
\Hom_{\Add(D)}(Z'',[j])\longrightarrow
\Hom_{\Add(D)}(Z',[j])\longrightarrow
\Hom_{\Add(D)}(Z,[j])
$$
is exact in $\Ab$.}
Let $t:Z^0\to\ZZ$ be a morphism from $Z^0$ to the constant functor
with value $\ZZ$ such that the composite morphism
\begin{equation*}
Z^1\to Z^0\xrightarrow{t}\ZZ\tag{2.3.4.1}
\end{equation*}\etiquette[2.3.4.1]{Vbis.eq2.3.4.1}
is zero.

By functoriality (\cf (\hyperref[Vbis.par2.3.1.1]{2.3.1.1})), one
deduces a morphism
$\varepsilon_\ast(F)\to\varepsilon_{Z^{\bigdot}\ast}(F)^0$ such that
the composite morphism
\begin{equation*}
\varepsilon_\ast(F)\to\varepsilon_{Z^{\bigdot}\ast}(F)^0\to
\varepsilon_{Z^{\bigdot}\ast}(F)^1\tag{2.3.4.2}
\end{equation*}\etiquette[2.3.4.2]{Vbis.eq2.3.4.2}
is zero.

There is then a canonical morphism of triangulated functors
\begin{equation*}
Q\circ K^+(\varepsilon_\ast)\xrightarrow{j}
Q\circ(\quad)_s\circ K^+(\varepsilon_{Z^{\bigdot}\ast})
=\sisi{\RR^{+}}{R}\varepsilon_{Z^{\bigdot}\ast}\circ Q
\tag{2.3.4.3}
\end{equation*}\etiquette[2.3.4.3]{Vbis.eq2.3.4.3}

\begin{enonce*}{Proposition 2.3.5}\etiquette[2.3.5]{Vbis.prop2.3.5}
Retain the conditions and notation of
(\hyperref[Vbis.subsubsec2.3.3]{2.3.3}) and
{(\hyperref[Vbis.subsubsec2.3.4]{2.3.4})}; suppose, in addition, that
the augmented complex
$$
\ldots\to Z^n\to Z^{n-1}\to\ldots\to Z^1\to Z^0
\xrightarrow{t}\ZZ
$$
defines a resolution of $\ZZ$.

Then, for every complex $F^{\bigdot}$ consisting of objects of
$I_{(S\times D,\mathcal{O}_S)}$
(\cf \hyperref[Vbis.prop1.3.10]{1.3.10}),
$$
j(F^{\bigdot}):
Q\circ K^+(\varepsilon_\ast)(F^{\bigdot})\to
(\sisi{\RR^+}{R}\varepsilon_{Z^{\bigdot}\ast}\circ Q)(F^{\bigdot})
$$
is\pageoriginale an isomorphism.
\end{enonce*}

\sisi{$+$}{} Let $i$ be an object of $D$, and let $Q_i$ be a module on
$(S,\mathcal{O}_S)$. We show that $j(e_{i\ast}(Q_i))$ is an
isomorphism. It is enough to see that the sequence
\begin{equation*}
0\to Q_i\to\varepsilon_{Z^{\bigdot}\ast}(e_{i\ast}(Q_i))^0\to
\varepsilon_{Z^{\bigdot}\ast}(e_{i\ast}(Q_i))^1\to\ldots
\tag{2.3.5.1}
\end{equation*}\etiquette[2.3.5.1]{Vbis.eq2.3.5.1}
is exact.

For this, it is enough to show that, for every object $X$ of
$\Mod(S,\mathcal{O}_S)$, the sequence
\begin{equation*}
0\to\Hom(X,Q_i)\to
\Hom(X,\varepsilon_{Z^{\bigdot}\ast}(e_{i\ast}(Q_i))^0)\to
\Hom(X,\varepsilon_{Z^{\bigdot}\ast}(e_{i\ast}(Q_i))^1)\to\ldots
\tag{2.3.5.2}
\end{equation*}\etiquette[2.3.5.2]{Vbis.eq2.3.5.2}
is exact. An immediate calculation shows that
\eqhyperref[Vbis.eq2.3.5.2]{2.3.5.2} reduces to
\hypertarget{sga4.exp.Vbis.equation.2.3.5.3}{}\begin{gather*}
0\to\Hom(\ZZ,\Hom(X,Q_i))\to
\Hom(\Hom_{\Add(D)}(Z^0,[i]),\Hom(X,Q_i))\to\tag{2.3.5.3}\\
\Hom(\Hom_{\Add(D)}(Z^1,[i]),\Hom(X,Q_i))\to\ldots{\quoi}.
\end{gather*}
\sisi{}{By hypothesis, the sequence
$$
0\leftarrow\ZZ\leftarrow\Hom_{\Add(D)}(Z^0,[i])
\leftarrow\Hom_{\Add(D)}(Z^1,[i])\leftarrow\ldots
$$
is an exact sequence of free modules and is therefore null-homotopic.
It follows that \hyperlink{sga4.exp.Vbis.equation.2.3.5.3}{(2.3.5.3)} is exact.}

\sisi{$+$}{} If, for every object $i$ of $D$, $Q_i$ is a totally
acyclic module, then, by applying the functor
$\prod_{i\in\Ob(D)}$ to the complexes
\eqhyperref[Vbis.eq2.3.5.1]{2.3.5.1}, one sees that the resulting
complex is still acyclic by (\hyperref[Vbis.lem1.3.10.1]{1.3.10.1}).
Thus $j(\prod_{i\in\Ob(D)}e_{i\ast}(Q_i))$ is an isomorphism.

\sisi{$+$}{} We leave it to the reader to deduce that
$j(F^{\bigdot})$ is an isomorphism when
$F^n\in I_{(S\times D,\mathcal{O}_S)}$ for every $n$ (for example,
by spectral sequences).

\begin{enonce*}{Corollary 2.3.6}\etiquette[2.3.6]{Vbis.cor2.3.6}
Under\pageoriginale the conditions of
(\hyperref[Vbis.prop2.3.5]{2.3.5}), the morphism $j$ identifies
$\sisi{\RR^+}{R}\varepsilon_{Z^{\bigdot}\ast}$ with the derived
functor of $\varepsilon_\ast$.
\end{enonce*}

The subcategory of
$K^{+}(\Ysheaf{\Gamma}(S\times D),\mathcal{O}_S)$ defined by the
complexes consisting of objects of
$I_{(S\times D,\mathcal{O}_S)}$ satisfies the conditions of theorem
(\hyperref[Vbis.subsec5.1]{5.1}) of [\cite{sga4.exp.Vbis.bib.7}, I] for the functor
$K^{+}(\varepsilon_\ast)$, by
(\hyperref[Vbis.prop1.3.10]{1.3.10}) and
(\hyperref[Vbis.prop2.3.5]{2.3.5}). The corollary follows immediately
from this observation.

\begin{enonce*}{Corollary 2.3.7}\etiquette[2.3.7]{Vbis.cor2.3.7}
Under the conditions of (\hyperref[Vbis.prop2.3.5]{2.3.5}), let $(E,A)$
be a ringed $D$-topos and let
$\theta:(E,A)\to(S\times D,\mathcal{O}_S)$ be an augmentation. There
is a canonical isomorphism
$$
\sisi{\RR^+}{R}\bar\theta_\ast\xrightarrow{\sim}
\sisi{\RR^+}{R}\varepsilon_{Z^{\bigdot}\ast}\circ
\sisi{\RR^+}{R}\Gamma(\theta_\ast){\quoi}.
$$
\end{enonce*}

This follows from [\cite{sga4.exp.Vbis.bib.7}, I, 5.4].

\begin{enonce*}[remark]{Remark 2.3.8}\etiquette[2.3.8]{Vbis.rem2.3.8}
We knew a priori that $\sisi{\RR^+}{R}\varepsilon_\ast$ exists and
that there is an isomorphism
$\sisi{\RR^+}{R}\bar\theta_\ast\xrightarrow{\sim}
\sisi{\RR^+}{R}\varepsilon_\ast\circ
\sisi{\RR^+}{R}\Gamma(\theta_\ast)$, since
$\Gamma(\theta_\ast)$ and $\varepsilon_\ast$ are induced by morphisms
of ringed topoi (\cf (\hyperref[Vbis.subsubsec1.3.6]{1.3.6}) and
(\hyperref[Vbis.subsubsec2.1.1]{2.1.1})). Nevertheless, the preceding
calculation will prove very useful and applicable in other contexts,
since in practice the conditions of
(\hyperref[Vbis.prop2.3.5]{2.3.5}) will always be satisfied.
\end{enonce*}

We shall now give the fundamental examples to which
(\hyperref[Vbis.prop2.3.5]{2.3.5}) can be applied.

\setcounter{subsubsection}{8}

% END INCLUDED SGA4_English_current_buildable_source_20260906/source/05bis/drafts/SGA4_VBIS_LINES_1074_1378_English.texfrag


% BEGIN INCLUDED SGA4_English_current_buildable_source_20260906/source/05bis/drafts/SGA4_VBIS_LINES_1379_1421_English.texfrag
\subsubsection{}\etiquette[2.3.9]{Vbis.subsubsec2.3.9}
In $\Add(\Delta)$, put $Z^n=[n]$ and define
$d^n:Z^n\to Z^{n+1}$ by the formula
\begin{equation*}
d^n=\sum_{i=0}^{n+1}(-1)^i\partial^i.\tag{2.3.9.1}
\end{equation*}\etiquette[2.3.9.1]{Vbis.eq2.3.9.1}
Here $\partial^i:[n]\to[n+1]$ is the $i$th face
[\cf \cite{sga4.exp.Vbis.bib.3} II 2]. For $t:Z^0\to\ZZ$, take the evident natural
augmentation\nde{Explicitly, if $F$ is a simplicial complex with
differential $\mathrm d_F$, then $\varepsilon_{Z^{\bigdot}*}(F)$ is
the simple complex associated with the double complex $F^j(i)$ whose
horizontal differential is the simplicial differential $\mathrm d^i$
and whose vertical differential is $\mathrm d_F^j$.}.

The condition of \hyperref[Vbis.subsubsec2.3.4]{2.3.4} is trivially
satisfied. The condition of \hyperref[Vbis.prop2.3.5]{2.3.5} follows
from \cite{sga4.exp.Vbis.bib.5} I (3.7.4).

\subsubsection{}\etiquette[2.3.10]{Vbis.subsubsec2.3.10}
In\pageoriginale $\Add(\Delta\times\Delta)$, consider the simple
complex associated with the double complex
\hypertarget{sga4.exp.Vbis.equation.2.3.10.1}{}\[\vcenter{
\xymatrix@C=2cm@R=1.75cm{
\big([n],[m+1]\big)  \ar[r]^{s_1^{n,m+1}} & \big([n+1],[m+1]\big)\\
\big([n],[m]\big) \ar[r]^{s_1^{n,m}} \ar[u]^{s_2^{n,m}} &
\big([n+1],[m]\big)\ar[u]_{s_2^{n+1,m}}
}}\leqno{(2.3.10.1)}
\]
with
\begin{equation*}
\tag{2.3.10.2}
\begin{cases}
S^{n,m}_1=\displaystyle\mathop{\sum}_{i=0}^{n+1}(-1)^i
(\partial^i,\id_{[m]})\,,\\[15pt]
S^{n,m}_2=\displaystyle\mathop{\sum}_{j=0}^{m+1}(-1)^j
(\id_{[n]},\partial^j){\quoi}.
\end{cases}
\end{equation*}\etiquette[2.3.10.2]{Vbis.eq2.3.10.2}

The reader may verify (using \hyperref[Vbis.subsubsec2.3.9]{2.3.9}) that the complex so defined,
equipped with the evident natural augmentation, satisfies the
conditions of \hyperref[Vbis.prop2.3.5]{2.3.5}.

\subsubsection{}\etiquette[2.3.11]{Vbis.subsubsec2.3.11}
The multisimplicial case is treated similarly (with
$\underbrace{\Delta\times\Delta\times\cdots\times\Delta}_{p\text{ times}}$).


% END INCLUDED SGA4_English_current_buildable_source_20260906/source/05bis/drafts/SGA4_VBIS_LINES_1379_1421_English.texfrag


% BEGIN INCLUDED SGA4_English_current_buildable_source_20260906/source/05bis/drafts/SGA4_VBIS_LINES_1422_1690_English.texfrag
\subsection{Relative cohomological
  descent}\etiquette[2.4]{Vbis.subsec2.4}\pageoriginale

\setcounter{subsubsection}{-1}
\subsubsection{}\etiquette[2.4.0]{Vbis.subsubsec2.4.0}
Let $\Ee\to B$ be a category bifibered in duals of topoi, and let
$\mathcal O$ be a ring of $\Ysheaf{\Gamma}(\Ee^{\circ})$.

These data canonically define a category fibered and cofibered in
abelian categories over $B^{\circ}$, denoted by
$\Ysheaf{\Mod}(\Ee^{\circ},\mathcal O)$ (\cf
(\hyperref[Vbis.rem1.3.3]{1.3.3}) and
(\hyperref[Vbis.subsubsec1.3.4]{1.3.4})). Throughout what follows,
the ``direct-image'' and ``inverse-image'' functors will always be
taken with respect to the fibered and cofibered functor
$\Ysheaf{\Mod}(\mathcal E^{\circ},\mathcal O)\to B^{\circ}$.

\begin{enonce*}{Definition 2.4.0.1}\etiquette[2.4.0.1]{Vbis.def2.4.0.1}
A morphism $f:T\to S$ in $B$ is said to be \emph{flat} (relative to
$(\mathcal E,\mathcal O)$) if
$f_{\ast}:\Mod(\mathcal E^{\circ}_S,\mathcal O_S)\to
\Mod(\mathcal E^{\circ}_T,\mathcal O_T)$ is an exact
functor\nde{This notation may cause confusion: one should probably
write $(f^\circ)_\star$, which formally coincides with
$(f^\star)^\circ$, where $f^\star$ is the fibered functor associated
with $\Ee\rightarrow B$. For example, for the category of sheaves on
topological spaces fibered over the category of topological spaces
(cf. \hyperref[Vbis.subsec4.1]{4.1}), one has
$(f^\circ)_\star(S,F)=(T,f^{-1}F)=(f^\star)^\circ(S,F)$, where
$f:T\rightarrow S$ is a continuous map of topological spaces, $F$
is a sheaf on $S$, and $f^{-1}$ is the usual inverse image. One then
recognizes the usual flatness conditions.}.
\end{enonce*}

\subsubsection{}\etiquette[2.4.1]{Vbis.subsubsec2.4.1}
Let $G$ be a fibered and cofibered subcategory of
$\Ysheaf{\Mod}(\mathcal E^{\circ},\mathcal O)$ such that, for every
object $S$ of $B$, $G_S$ is a \emph{thick} subcategory (\cf
\sisi{Tohokû}{Tôhoku} (1.11)) of
$(\Ysheaf{\Mod}(\mathcal E^{\circ},\mathcal O))_S$. The reader will
find examples of such situations in \hyperref[Vbis.sec4]{Section 4}.

If $S$ is an object of $B$, denote by $D^+(S)$ the derived category
of $\Mod(\mathcal E^{\circ}_S,\mathcal O_S)$. Following
[\cite{sga4.exp.Vbis.bib.7} (I \S\ 4)], introduce the category $D^+_{G_S}(S)$, which
identifies with the full subcategory of $D^+(S)$ formed by the
objects $X^{\bigdot}$ such that $H^n(X^{\bigdot})\in G_S$ for every
$n$.

\emph{We assume throughout that the following condition holds.}

\paragraph{}\etiquette[2.4.1.1]{Vbis.par2.4.1.1}
\emph{For every morphism $h:T\to S$ in $B$, the functor
$\sisi{\RR^+}{Rh^{\ast}}:D^+(T)\to D^+(S)$ has the property that
$\sisi{\RR^+}{Rh^{\ast}}(F)\in D^+_{G_S}(S)$ whenever
$F\in G_T$}\nde{As above, it would be more correct to write
$R(h^\circ)^\ast=R(h_\ast)^\circ$ rather than $Rh^\ast$, which would
make the notation more consistent with the usual geometric
intuition. One could equally well have taken the fibered and
cofibered functors with respect to
$\Mod(\mathcal E,\mathcal O)\rightarrow B$ rather than the opposite
categories. The fibers, however, would then have been only duals of
topoi.}.

By [\cite{sga4.exp.Vbis.bib.7} (I. (7.3). (ii))], this condition implies that
$\sisi{\RR^+}{Rh^{\ast}}(D^+_{G_T}(T))\subset D^+_{G_S}(S)$.

\subsubsection{}\etiquette[2.4.2]{Vbis.susubsec2.4.2}
Let\pageoriginale $D$ be a $\Uu$-small category. \emph{Throughout
what follows, we suppose given a complex $Z^{\bigdot}$ of
$\Add(D)$ satisfying the conditions of}
(\hyperref[Vbis.prop2.3.5]{2.3.5}) (in applications, one will in fact
have $D=\Delta$ or $D=\Delta\times\Delta$).

If $X:D^{\circ}\to B$ is a $D$-object of $B$, denote by
$\Mod_G(\Ysheaf{\Gamma}(\bar X),\mathcal O)$ the full subcategory of
$\Mod(\Ysheaf{\Gamma}(\bar X),\mathcal O)$ formed by the objects $F$
such that $e_i^{\ast}(F)\in G_i$ $(=G_{X_i})$ for every object $i$
of $D$. The category $\Mod_G(\Ysheaf{\Gamma}(\bar X),\mathcal O)$
is a thick subcategory of
$\Mod(\Ysheaf{\Gamma}(\bar X),\mathcal O)$. Denote by
$K_G^+(\Ysheaf{\Gamma}(\bar X),\Oo)$ (\resp
$D_G^+(\Ysheaf{\Gamma}(\bar X),\Oo)$) the full subcategory of
$K^+(\Ysheaf{\Gamma}(\bar X),\Oo)$ (\resp
$D^+(\Ysheaf{\Gamma}(\bar X),\Oo)$) formed by the objects
$X^{\bigdot}$ such that
$H^n(X^{\bigdot})\in\Mod_G(\Ysheaf{\Gamma}(\bar X),\Oo)$ for every
$n$.

\subsubsection{}\etiquette[2.4.3]{Vbis.subsubsec2.4.3}
Let $X$ and $X'$ be two $D$-objects of $B$, and let
$\alpha:X\to X'$ be a functorial morphism. Recall (\cf
(\hyperref[Vbis.subsubsec1.3.4]{1.3.4})) that the corresponding
morphism of ringed $D$-topoi is again denoted by
$\alpha:(\bar X,\mathcal O)\to(\bar X',\mathcal O)$.

\begin{enonce*}{Definition 2.4.3.1}\etiquette[2.4.3.1]{Vbis.def2.4.3.1}
The morphism $\alpha:X\to X'$ is said to be \emph{flat} if the
corresponding morphism of ringed $D$-topoi is flat in the sense of
(\hyperref[Vbis.par1.3.6.1]{1.3.6.1}).
\end{enonce*}

We then have the evident lemma:

\begin{enonce*}{Lemma 2.4.3.2}\etiquette[2.4.3.2]{Vbis.lem2.4.3.2}
The morphism $\alpha:X\to X'$ is flat if and only if, for every
object $i$ of $D$, $\alpha_i:X_i\to X'_i$ is a flat morphism
\emph{(\cf (\hyperref[Vbis.def2.4.0.1]{2.4.0.1}))}.
\end{enonce*}

\subsubsection{}\etiquette[2.4.4]{Vbis.subsubsec2.4.4}
A morphism $\alpha:X\to X'$ defines a functor
\begin{equation*}
K^+(\Gamma(\alpha_*)):K^+(\Ysheaf{\Gamma}(\bar X),\Oo)
\longrightarrow K^+(\Ysheaf{\Gamma}(\bar X'),\Oo){\quoi}.
\tag{2.4.4.1}
\end{equation*}\etiquette[2.4.4.1]{Vbis.eq2.4.4.1}

\begin{enonce*}{Lemma 2.4.4.2}\etiquette[2.4.4.2]{Vbis.lem2.4.4.2}
The functor
$$
K^+(\Gamma(\alpha_*))|_{K_G^+(\Ysheaf{\Gamma}(\bar X),\Oo)}:
K_G^+(\Ysheaf{\Gamma}(\bar X),\Oo)\longrightarrow
K^+(\Ysheaf{\Gamma}(\bar X'),\Oo)
$$
has a right derived functor, denoted by
$\sisi{\RR^+}{R}_G\Gamma(\alpha_{\ast})$, and the canonical
morphism
$$
\sisi{\RR^+}{R}_G\Gamma(\alpha_{\ast})\longrightarrow
\sisi{\RR^+}{R}\Gamma(\alpha_{\ast})|_{D_G^+
(\Ysheaf{\Gamma}(\bar X),\Oo)}
$$
is\pageoriginale an isomorphism.

Moreover,
$$
\sisi{\RR^+}{R}_G\Gamma(\alpha_{\ast})
(D_G^+(\Ysheaf{\Gamma}(\bar X),\Oo))
\subset D_G^+(\Ysheaf{\Gamma}(\bar X'),\Oo).
$$
\end{enonce*}

The first part follows from (\hyperref[Vbis.prop1.3.10]{1.3.10}) and
[\cite{sga4.exp.Vbis.bib.7}, (I. (5.2))]. To verify the last assertion, by
[\cite{sga4.exp.Vbis.bib.7}, (I. (7.3). (ii))], it suffices to show that
$\sisi{\RR^+}{R}\Gamma(\alpha_{\ast})(F)\in
D_G^+(\Ysheaf{\Gamma}(\bar X'),\Oo)$ whenever
$F\in\Mod_G(\Ysheaf{\Gamma}(\bar X),\Oo)$. This follows from the
explicit calculation of
$\sisi{\RR^+}{R}\Gamma(\alpha_\ast)$ given by \hyperref[Vbis.cor1.3.11]{(1.3.1.1)} and from
hypothesis (\hyperref[Vbis.par2.4.1.1]{2.4.1.1}).

\subsubsection{}\etiquette[2.4.5]{Vbis.subsubsec2.4.5}
Let $S$ be an object of $B$. The ringed $D$-topos
$(\bar{C^D_S},\mathcal O)$ (\cf
(\hyperref[Vbis.subsubsec1.2.6]{1.2.6})) identifies with the constant
ringed $D$-topos
$(\mathcal E^{\circ}_S\times D,\mathcal O_S)$, and we have the
following lemma, which follows from [\cite{sga4.exp.Vbis.bib.7} (I. (7.3). (ii))]:

\begin{enonce*}{Lemma 2.4.5.1}\etiquette[2.4.5.1]{Vbis.lem2.4.5.1}
The functor
$\sisi{\RR^+}{R}\varepsilon_{Z^{\bigdot}\ast}|_{
D_G^+(\Ysheaf{\Gamma}(\bar{C^D_S}),\mathcal O)}$, denoted by
$\sisi{\RR^+}{R}_G\varepsilon_{Z^{\bigdot}\ast}$, takes values in
$D^+_{G_S}(S)$ and identifies with a right derived functor of
$$
K^+(\varepsilon_{\ast})|_{K_G^+
(\Ysheaf{\Gamma}(\bar{C^D_S}),\mathcal O)}:
K_G^+(\Ysheaf{\Gamma}(\bar{C^D_S}),\mathcal O)
\longrightarrow K^+(S){\quoi}.
$$
\end{enonce*}

\begin{enonce*}{Proposition 2.4.6}\etiquette[2.4.6]{Vbis.prop2.4.6}
Let $\theta:X\to C^D_S$ be an augmented $D$-object of $B$. The
functor
$K^+(\bar\theta_{\ast})|_{K_G^+
(\Ysheaf{\Gamma}(\bar X),\mathcal O)}$ has a right derived functor,
denoted by $\sisi{\RR^+}{R}_G\bar\theta_{\ast}$, taking values in
$D^+_{G_S}(S)$. Moreover, there is a canonical isomorphism
$$
\sisi{\RR^+}{R}_G\bar\theta_{\ast}\xrightarrow{\sim}
\sisi{\RR^+}{R}_G\varepsilon_{Z^{\bigdot}\ast}\circ
\sisi{\RR^+}{R}_G\Gamma(\theta_{\ast}){\quoi}.
$$
\end{enonce*}

The proof, from what precedes, is left to the reader.

We are now in a position to give the definitions concerning
cohomological descent.

\setcounter{subsubsection}{6}
\subsubsection{}\etiquette[2.4.7]{Vbis.subsubsec2.4.7}
Let\pageoriginale $\theta:X\to C^D_S$ be an augmented $D$-object of
$B$. Suppose that $\theta$ is flat, so that the restriction to
$D^+_{G_S}(S)$ of the functor $\bar\theta^{\ast}$ takes values in
$D_G^+(\Ysheaf{\Gamma}(\bar X),\mathcal O)$; denote this restriction
by $\bar\theta_G^{\ast}$. With this notation:

\begin{enonce*}{Lemma 2.4.7.1}\etiquette[2.4.7.1]{Vbis.lem2.4.7.1}
There are two functorial morphisms
$\alpha:\id_{D^+_{G_S}(S)}\to
\sisi{\RR^+}{R}_G\bar\theta_{\ast}\circ\bar\theta_G^{\ast}$ and
$\beta:\bar\theta_G^{\ast}\circ
\sisi{\RR^+}{R}_G\bar\theta_{\ast}\to
\id_{D_G^+(\Ysheaf{\Gamma}(\bar X),\mathcal O)}$ which make these
two functors adjoint.
\end{enonce*}

This follows immediately from
(\hyperref[Vbis.lem2.2.3.1]{2.2.3.1}) and
(\hyperref[Vbis.prop2.4.6]{2.4.6}).

\begin{enonce*}{Definition 2.4.8}\etiquette[2.4.8]{Vbis.def2.4.8}
The morphism $\theta$ is called an augmentation of cohomological
$1$-descent relative to $G$ if
\begin{enumerate}
\renewcommand{\theenumi}{\arabic{enumi}}
\renewcommand{\labelenumi}{\theenumi$^{\circ}$)}
\item $\theta$ is flat.

\item The functor $\bar\theta_G^{\ast}$ is fully faithful.
\end{enumerate}
\end{enonce*}

\begin{enonce*}{Remark 2.4.9}\etiquette[2.4.9]{Vbis.rem2.4.9}
Condition $2^{\circ})$ of the preceding definition may equivalently
be expressed by saying that the morphism $\alpha$ in
(\hyperref[Vbis.lem2.4.7.1]{2.4.7.1}) is an isomorphism.
\end{enonce*}

\begin{enonce*}{Definition 2.4.10}\etiquette[2.4.10]{Vbis.def2.4.10}
Let $\theta:X\to C^D_S$ be an augmented $D$-object of $B$. The
morphism $\theta$ is called an \emph{augmentation of effective
descent relative to} $G$ if the functor
$$
\Gamma(\theta^{\ast})\circ\varepsilon^{\ast}:G_S\longrightarrow
\Ysheaf{\Gamma}^{\Cocart}(D,G)
$$
is an equivalence of categories.
\end{enonce*}

\begin{enonce*}{Definition 2.4.11}\etiquette[2.4.11]{Vbis.def2.4.11}
The morphism $\theta$ is called an \emph{augmentation of
cohomological $2$-descent} (or of \emph{effective cohomological
descent}) \emph{relative to} $G$ if $\theta$ is both an augmentation
of cohomological $1$-descent and an augmentation of effective
descent relative to $G$.
\end{enonce*}

The proof of (\hyperref[Vbis.prop2.2.7]{2.2.7}) then gives the
following result:

\begin{enonce*}{Theorem 2.4.12}\etiquette[2.4.12]{Vbis.thm2.4.12}
Let\pageoriginale $\theta:X\to C^D_S$ be an augmentation of
cohomological $2$-descent relative to $G$. Then the essential image
of $\bar\theta_G^{\ast}$ is the full subcategory of
$D_G^+(\Ysheaf{\Gamma}(\bar X),\mathcal O)$ formed by the complexes
$F^{\bigdot}$ such that, for every $n$, $H^n(F^{\bigdot})$ is a
cocartesian section of $G$.
\end{enonce*}

The following proposition, whose verification is left to the reader,
translates Definition (\hyperref[Vbis.def2.4.10]{2.4.10}) into the
language of (\cite{sga4.exp.Vbis.bib.4}, 6):

\begin{enonce*}{Proposition 2.4.13}\etiquette[2.4.13]{Vbis.prop2.4.13}
The morphism $\theta:X\to C^D_S$ is an augmentation of effective
descent relative to $G$ if and only if the functor
$D^{\circ}\to B/S$ over $B$ defined by $\theta$ induces an
equivalence between the category of cartesian sections of
$G^{\circ}$ over $B/S$ and the category of cartesian sections of
$G^{\circ}$ over $D^{\circ}$.
\end{enonce*}

\begin{enonce*}{Corollary 2.4.14}\etiquette[2.4.14]{Vbis.cor2.4.14}
Suppose that finite fiber products are representable in $B$, and let
$f:R\to S$ be a morphism. The augmented semi-simplicial object
$[R|_fS]$ \emph{(\cf
(\hyperref[Vbis.subsubsec1.2.7]{1.2.7}))} is of effective descent
relative to $G$ if and only if $f$ is an effective descent morphism
for the fibered category $G^{\circ}$ over $B$.
\end{enonce*}

This follows from \hyperref[Vbis.prop2.4.13]{(2.4.13)} and (\cite{sga4.exp.Vbis.bib.4}, 9).

\begin{enonce*}{Definition 2.4.15}\etiquette[2.4.15]{Vbis.def2.4.15}
Suppose that finite fiber products are representable in $B$, and let
$f:R\to S$ be a morphism. We say that $f$ is a morphism of
cohomological $i$-descent relative to $G$ $(i=1,2)$ if the augmented
semi-simplicial object $[R|_fS]$ is of cohomological $i$-descent
relative to $G$.
\end{enonce*}

\begin{enonce*}{Definition 2.4.16}\etiquette[2.4.16]{Vbis.def2.4.16}
Suppose that finite fiber products are representable in $B$, and let
$\theta:X\to C^D_S$ be an augmentation. We say that $\theta$ is an
\emph{augmentation of universal cohomological $i$-descent
$(i=1,2)$ relative to $G$} if, for every morphism $T\to S$, the
augmentation $X_T\xrightarrow{\theta_T}C^D_T$ obtained by base
change\pageoriginale is an augmentation of cohomological
$i$-descent relative to $G$.
\end{enonce*}

The notion of a \emph{morphism of universal cohomological
$i$-descent relative to} $G$ follows immediately from the two
preceding definitions.

We will often use the terminology ``augmentation of
$G$-cohomological $i$-descent'' in place of ``augmentation of
cohomological $i$-descent relative to $G$''.

\subsection{The descent spectral sequence}\etiquette[2.5]{Vbis.subsec2.5}


% END INCLUDED SGA4_English_current_buildable_source_20260906/source/05bis/drafts/SGA4_VBIS_LINES_1422_1690_English.texfrag


% BEGIN INCLUDED SGA4_English_current_buildable_source_20260906/source/05bis/drafts/SGA4_VBIS_LINES_1691_1929_English.texfrag

\subsubsection{}\etiquette[2.5.1]{Vbis.subsubsec2.5.1}
Let $A$ be an abelian category. Denote by $F(A)$ the category whose
objects are the objects of $A$ equipped with a discrete and codiscrete
filtration, and whose morphisms are the filtered morphisms of $A$:
the category $F(A)$ is an additive category (but is not abelian).

The category $K(F(A))$ of filtered complexes of $A$, whose filtration
is discrete and codiscrete degree by degree, is a triangulated
category, and the canonical functors
\begin{equation*}
\Gr_n : K(F(A)) \longrightarrow K(A)\tag{2.5.1.1}
\end{equation*}\etiquette[2.5.1.1]{Vbis.eq2.5.1.1}
are triangulated. The set $\Sigma$ of morphisms $f$ of $K(F(A))$
such that $\Gr_n(f)$ is a quasi-isomorphism for every $n$ is therefore
a saturated multiplicative system [\cf \cite{sga4.exp.Vbis.bib.9}. \S~2. ${\rm n^\circ}$ 1].

By inverting the arrows of this system, one obtains a new
triangulated category, denoted by $D F(A)$, and the functors $\Gr_n$
extend to functors
\begin{equation*}
\Gr_n : D F(A) \longrightarrow D(A){\quoi}.\tag{2.5.1.2}
\end{equation*}\etiquette[2.5.1.2]{Vbis.eq2.5.1.2}

In what follows, we will consider only complexes bounded below: this
naturally gives rise to the notation $K^+(F(A))$ and $D^+ F(A)$.

\subsubsection{}\etiquette[2.5.2]{Vbis.subsubsec2.5.2}
Let\pageoriginale $B$ be a thick subcategory of $A$, and denote by
$K^+_B(F(A))$ the full subcategory of $K^+(F(A))$ whose objects are
the complexes $X^{\bigdot}$ such that $H^{i}(\Gr_j(X^{\bigdot})) \in B$
for every pair of integers $(i, j): K_B^+(F(A))$ is a
\emph{localizing triangulated subcategory} of $K(F(A))$ [\cf
  \cite{sga4.exp.Vbis.bib.7}. (I. \S\ 5)]. We denote by $D_B^+ F(A)$ the full subcategory
of $D F(A)$ whose objects are the complexes bounded below
$X^{\bigdot}$ such that $H^i(
\Gr_j(X^{\bigdot})) \in B$ for every pair $(i, j)$: by (\loccit),
$D^+_B F(A)$ identifies with the category of fractions
$K^+_B(F(A))_{\Sigma \cap K^{+}_{B} (F(A))}$.

The ``forgetful functor for filtrations''
$$
\iota : K^+(F(A)) \longrightarrow K^+(A)
$$
is triangulated. Moreover, it follows from [\cite{sga4.exp.Vbis.bib.5}. I. (4.7)] that
$\iota(f)$ is a quasi-isomorphism if $f \in \Sigma$ and that
$\iota(K^+_B(F(A))) \subset K^+_B(A)$. The functor $\iota$ thus
extends to a triangulated functor
\begin{equation*}
\iota : D^+ F(A) \longrightarrow D^+ (A) \tag{2.5.2.1}
\end{equation*}\etiquette[2.5.2.1]{Vbis.eq2.5.2.1}
such that $\iota (D^+_B F(A)) \subset D^+_B (A)$.

Finally, note that there is a canonical \emph{spectral functor}
$D^+_B F(A) \to B$, converging to $H^n \circ \iota$, whose $E_1$ term
is given by $H^{p+q} \circ \Gr_{p}$ [\cf \cite{sga4.exp.Vbis.bib.5} I. (4.2)].

\emph{We now return to the notation of} (2.3), \emph{assuming that}
$D = \Delta$: $Z^{\bigdot}$ will denote the complex of $\Add(\Delta)$
defined by \eqhyperref[Vbis.eq2.3.9.1]{2.3.9.1}.

\begin{enonce*}{Proposition 2.5.3}\etiquette[2.5.3]{Ii.prop2.5.3}
Let $(S, \mathcal{O}_S)$ be a ringed topos. The functor
$$\sisi{\RR^+}{R}
\varepsilon_{Z^{\bigdot} \ast} : D^{+} (\Ysheaf{\Gamma} (S \times D),
\mathcal{O}_S) \to D^{+} (S, \mathcal{O}_S)$$ admits a canonical
factorization
$$
D^{+} (\Gamma (S \times D), \mathcal{O}_S) \xrightarrow{F^{\sisi{\RR^+}{R}
\varepsilon} Z^{\bigdot}\ast} D^{+} F(S, \mathcal{O}_S)
\xrightarrow{\iota} D^{+} (S, \mathcal{O}_S){\quoi}.
$$
Moreover, for every integer $i$, the following diagram is commutative:
\[
\xymatrix@=2.5cm{
D^{+}(\Ysheaf{\Gamma}(S\times D),\Oo_S)
\ar[r]^(.55){\sisi{\RR^{+}}{R}e^{\ast}_i} \ar[d]_{F^{\sisi{\RR^{+}}{R}\epsilon_{Z^{\bigdot}\ast}}}& D^{+}(S,\Oo_S)\\
D^{+}F(S,\Oo_S) \ar[ru]_{\Gr_i}& .
}
\]
\end{enonce*}

Consider\pageoriginale the functor $C^+ (C^+(S, \mathcal{O}_S))
\xrightarrow{(\quad)_s} C^+(S, \mathcal{O}_S)$: if $X^{\cdot\cdot}$ is
an object of $C^+ (C^+(S, \mathcal{O}_S))$, we equip
$(X^{\cdot\cdot})_s$ with its second canonical filtration [\cf
  \cite{sga4.exp.Vbis.bib.5}. I.4.8]; this gives a factorization
\begin{equation*}
(\quad)_s : C^+ (C^+(S, \mathcal{O}_S)) \longrightarrow C^+ (F(S,
  \mathcal{O}_S)) \xrightarrow{\iota} C^+(S, \mathcal{O}_S)\tag{2.5.3.1}
\end{equation*}\etiquette[2.5.3.1]{Vbis.eq2.5.3.1}
which passes to the categories $K^+$:
\begin{equation*}
(\quad)_s : K^+ (C^+(S, \mathcal{O}_S)) \longrightarrow K^+ (F(S,
  \mathcal{O}_S)) \xrightarrow{\iota} K^+(S, \mathcal{O}_S)\tag{2.5.3.2}
\end{equation*}\etiquette[2.5.3.2]{Vbis.eq2.5.3.2}
because a homotopy in $C^+ (C^+(S, \mathcal{O}_S))$ induces a
filtered homotopy.

We then have a diagram
\[\vcenter{
\xymatrix{
D^{+}(\Ysheaf{\Gamma}(S\times\Delta),\Oo_S) \ar[rr]^{R\epsilon_{Z^{\bigdot}}\ast} \ar@{-->}[rd]&                   &D^{+}(S,\Oo_S)\\
                                              & D^{+}F(S,\Oo_S)
\ar[ru]_{\imath}&\\
                                  & K^{+}(F(S,\Oo_S))
\ar[u] \ar[rd]^{\imath}&\\
K^{+}(\Ysheaf{\Gamma}(S\times \Delta),\Oo_S) \ar[uuu]
\ar[rr]_{(\;)_s\circ K^{+}(\epsilon_{Z^{\bigdot}}\ast)} \ar[ru]& & K^{+}(S,\Oo_S)\ar[uuu]
}}\leqno{(2.5.3.3)}
\]\etiquette[2.5.3.3]{Vbis.eq2.5.3.3}
and one verifies that there is a unique functor
$$
{}_F \sisi{\RR^{+}}{R} \epsilon_{Z^{\bigdot} \ast} : D^{+}(\Ysheaf{\Gamma}(S \times
\Delta), \mathcal{O}_S) \longrightarrow D^{+}F(S, \mathcal{O}_S)
$$
making\pageoriginale \eqhyperref[Vbis.eq2.5.3.3]{2.5.3.3} commutative.

The rest of the proposition is evident.

\begin{enonce*}{Proposition 2.5.4}\etiquette[2.5.4]{Vbis.prop2.5.4}
Let $(E, A)$ be a semi-simplicial ringed topos and
$$\theta : (E, A) \to (S \times \Delta,
\mathcal{O}_S)$$ an augmentation. For every integer $i$, denote by
$$\theta_i : (E_i, A_i) \to (S, \mathcal{O}_S)$$ the morphism of
ringed topoi induced by $\theta$ over $i$.

The functor $\sisi{\RR^+}{R} \bar{\theta}_\ast : D^+ (\Ysheaf{\Gamma}(E), A)
\to D^+ (S, \mathcal{O}_S)$ admits a canonical factorization
$$
D^+ (\Ysheaf{\Gamma}(E), A) \xrightarrow{{}_F\sisi{\RR^{+}}{R} \bar{\theta}_{\ast}}
D^{+} F(S, \mathcal{O}_S) \xrightarrow{\iota} D^+ (S, \mathcal{O}_S)
$$
such that, for every integer $i$, the canonical diagram
\[
\xymatrix@C=2.5cm@R=1.75cm{
D^{+}(\Ysheaf{\Gamma}(E),A)  \ar[r]^{\sisi{\RR^{+}}{R}e_i^{\ast}}
\ar[d]_{{}_F \sisi{\RR^{+}}{R}
\bar{\theta}_\ast}&
D^{+}(E_i,A_i) \ar[d]^{\sisi{\RR^{+}}{R}\theta_{i\ast}}\\
D^{+}F(S,\Oo_S)  \ar[r]_{\Gr_i} & D^{+}(S,\Oo_S)\\
}
\]
is essentially commutative.
\end{enonce*}

This follows from what precedes and from \hyperref[Vbis.cor1.3.12]{(1.3.1.2)}.

\begin{enonce*}{Proposition 2.5.5}\etiquette[2.5.5]{Vbis.prop2.5.5}
Let $(E, A)$ be a semi-simplicial ringed topos and let
$$\theta : (E, A) \to (S \times \Delta, \mathcal{O}_S)$$ be an
augmentation of cohomological $1$-descent. Let
$$H : (S, \mathcal{O}_S) \to (R, \mathcal{O}_R)$$ be a morphism of
ringed topoi. Then there is a spectral functor from
$D^+(S, \mathcal{O}_S)$ to $\Mod(R, \mathcal{O}_R)$:
$$
E^{pq}_1 = R^q(H \circ \theta_p)_{\ast} \circ
\theta^{\ast}_p \longrightarrow R^{p+q}H_{\ast}
$$
called\pageoriginale the \emph{descent spectral functor}.
\end{enonce*}

Let $H \circ \theta : (E, A) \to (R \times \Delta, \mathcal{O}_R)$
be the augmentation canonically deduced from $\theta$ by composition
with $H$; we have
\begin{equation*}
\bar{(H \circ \theta)}_{\ast} = H_{\ast} \circ \bar{\theta}_{\ast}\tag{2.5.5.1}
\end{equation*}\etiquette[2.5.5.1]{Vbis.eq2.5.5.1}
so that the diagram
\[\vcenter{
\xymatrix@C=2.25cm@R=2.2cm{
D^{+}(S,\Oo_S) \ar[r]^{\sisi{\LL^{+}(\bar{\theta}^{\ast})}{\bar{\theta}^{\ast}}} \ar[rd]_{\sisi{\RR^+}{R}H_{\ast}} &
D^{+}(\Ysheaf{\Gamma}(E),A)
\ar[d]^{\sisi{\RR^{+}}{R}(\overline{H\circ
\theta})_\ast}\\
                  & D^{+}(R,\Oo_R)
}}\leqno{(2.5.5.2)}
\]\etiquette[2.5.5.2]{Vbis.eq2.5.5.2}
is essentially commutative, since
$$
\sisi{\RR^+}{R} \bar{(H \circ \theta)}_{\ast} \xrightarrow{\sim} \sisi{\RR^+}{R}
H_\ast \circ \sisi{\RR^+}{R} \bar{\theta}_{\ast} \mbox{ and } \id_{D^{+}(S,
  \mathcal{O}_S)} \xrightarrow{\sim} \sisi{\RR^+}{R} \bar{\theta}_\ast \circ
\bar{\theta}^{\ast}{\quoi}.
$$
By (\hyperref[Vbis.prop2.5.4]{2.5.4}), for every $i$ we have an
essentially commutative diagram
\[\vcenter{
\xymatrix@C=2.25cm@R=1.75cm{
D^{+}(S,\Oo_S) \ar[r]^{\bar{\theta}^{\ast}} \ar[rdd]_{\sisi{\RR^{+}}{R}H_{\ast}}&
D^{+}(\Ysheaf{\Gamma}(E),A)
\ar[r]^{\sisi{\RR^{+}}{R}e^{\ast}_i} \ar[d]_{{}_F\sisi{\RR^{+}}{R}(\bar{H\circ\theta})_\ast}&
D^{+}(E_i,A_i) \ar[d]^{\sisi{\RR^{+}}{R}(H\circ\theta_{i})_{\ast}}\\
                  & D^{+}F(R,\Oo_R)   \ar[d]_{\imath}\ar[r]_{\Gr_i}      & D^{+}(R,\Oo_R)\\
                  & D^{+}(R,\Oo_R)          &   .
}}\leqno{(2.5.5.3)}
\]\etiquette[2.5.5.3]{Vbis.eq2.5.5.3}
The\pageoriginale asserted spectral functor is obtained by taking the
canonical spectral functor on $D^+ F(R, \mathcal{O}_R)$ and composing
it with ${}_F R(\bar{H \circ \theta})_{\ast}\circ\bar{\theta}^{\ast}$.

We leave it to the reader to verify the following proposition:

\begin{enonce*}{Proposition 2.5.6}\etiquette[2.5.6]{Vbis.prop2.5.6}
With the preceding notation, the $E^{pq}_2$ term of the preceding
spectral functor is
$$
\check{H}^p \circ R^q \Gamma ((H \circ \theta)_{\ast}) \circ
\bar{\theta}^\ast
$$
where $\check{H}^p : \Mod(\Ysheaf{\Gamma}(R \times \Delta), \mathcal{O}_R) \to
\Mod(R, \mathcal{O}_R)$ is the functor that assigns to every functor
$\Delta \xrightarrow{F} \Mod(R, \mathcal{O}_R)$ the homology object
in degree $p$ of the associated complex (denoted by
$\varepsilon_{Z^{\bigdot}\ast} (F)^{\bigdot}$ in
(\hyperref[Vbis.subsubsec2.3.3]{2.3.3})).
\end{enonce*}

We now return to the terminology introduced in
(\hyperref[Vbis.subsec2.4]{2.4}): by the sorites
(\hyperref[Vbis.subsubsec2.5.2]{2.5.2}), all the preceding
considerations carry over word for word by placing the letter $G$ as
a subscript wherever this makes sense. The details are left to the
reader; in particular, one obtains:

\begin{enonce*}{Proposition 2.5.7}\etiquette[2.5.7]{Vbis.prop2.5.7}
Let\pageoriginale $\theta : X \to C^D_S$ be an augmentation of
cohomological $1$-descent relative to $G$, and let $h : S \to R$ be a
morphism of $B$; there is a spectral functor from $D^+_{G_S}(S)$ to
$G_R$:
$$
E^{pq}_1 = R^q_G(h \circ \theta)_\ast \circ \theta_{p,G}^\ast
\to R^{p + q}_G h_{\ast}
$$
with $E^{pq}_2 = \check{H}^p \circ R^q_G \Gamma((h \circ \theta)_\ast)
\circ \bar\theta_G^{\star}$.
\end{enonce*}

% END INCLUDED SGA4_English_current_buildable_source_20260906/source/05bis/drafts/SGA4_VBIS_LINES_1691_1929_English.texfrag


% BEGIN INCLUDED SGA4_English_current_buildable_source_20260906/source/05bis/drafts/SGA4_VBIS_LINES_1930_2298_English.texfrag
\section{Descent criteria}\etiquette[3]{Vbis.sec3}

\setcounter{subsection}{-1}
\subsection{Notation}\etiquette[3.0]{Vbis.subsec3.0}

Throughout what follows, we retain the notation of
(\hyperref[Vbis.subsec2.4]{2.4}). We shall always assume \emph{that
finite fiber products are representable in} $B$.

We shall also assume that $B$ satisfies the following condition:

\setcounter{subsubsection}{-1}
\subsubsection{}\etiquette[3.0.0]{Vbis.subsubsec3.0.0}
\emph{Coproducts exist in $B$, are disjoint and universal, and are
families of effective $\mathcal{E}$-descent and effective
$G^{\circ}$-descent} [\cf \cite{sga4.exp.Vbis.bib.4} (9.23), (9.25), and (9.27)].

\subsubsection{}\etiquette[3.0.1]{Vbis.subsubsec3.0.1}
We now recall some classical notation concerning simplicial objects
of a category [\cf \cite{sga4.exp.Vbis.bib.3}, Chap.~II, and (V, Appendix)].

Let $E$ be a category \emph{having finite projective limits}.
$\Delta^{\circ}E$ denotes the category of simplicial objects of $E$
(\cf (\hyperref[Vbis.subsubsec1.2.6]{1.2.6})).

For an integer $n$, denote by $\mathop{\Delta}_n$ (\resp
${\mathop{\Delta}_n}^{+}$, ${\mathop{\Delta}_n}^{-}$) the full
subcategory of $\Delta$ (\resp $\Delta^+$, $\Delta^-$) consisting of
the objects $[p]$ such that $p\leq n$.

The restriction functor
\begin{equation*}
i_n^{\ast}:\Delta^{\circ}E\longrightarrow
{\mathop{\Delta}_n}^{\circ}E\tag{3.0.1.1}
\end{equation*}\etiquette[3.0.1.1]{Vbis.eq3.0.1.1}
has a right adjoint $i_{n\ast}$, since finite projective limits exist
in\pageoriginale $E$ (\cf (\hyperref[Vbis.lem1.2.10]{1.2.10})). The
functor $i_{n\ast}\circ i_n^{\ast}$ is denoted by $\cosk_n$ and is
called the \emph{coskeleton functor of order $n$}; by abuse of
notation, we shall also use $\cosk_n$ for the functor $i_{n\ast}$.

For every integer $p$, denote by $\mathop{\Delta}_{n[p]}$ (\resp
${\mathop{\Delta}_{n[p]}}^+$) the full subcategory of
$\Delta_{[p]}$ (\resp $\Delta_{[p]}^+$), the category of objects of
$\Delta$ (\resp $\Delta^+$) over $[p]$, defined by the objects $[q]$
over $[p]$ such that $q\leq n$. We leave it to the reader to verify
that
\begin{equation*}
(\cosk_n(X))_p=\varprojlim_{\mathop{\Delta}_{n[p]}}X_q
\xrightarrow{\sim}
\varprojlim_{\mathop{\Delta}_{n[p]}^{+}}X_q{\quoi}.
\tag{3.0.1.2}
\end{equation*}\etiquette[3.0.1.2]{Vbis.eq3.0.1.2}

\subsubsection{}\etiquette[3.0.2]{Vbis.subsubsec3.0.2}
Let $X$ and $X'$ be two objects of $\Delta^{\circ}E$, and let
$f,g:X\rightrightarrows X'$ be two morphisms. A \emph{homotopy from
$f$ to $g$} consists of a map
$h_n:\Hom_{\Delta}([n],[1])\to\Hom_E(X_n,X'_n)$ for every $n$,
satisfying the following two conditions:
\begin{equation*}
h_n(\partial^1)=f_n\quad\mbox{and}\quad
h_n(\partial^\circ)=g_n\quad\mbox{for every }n{\quoi}.
\tag{3.0.2.1}
\end{equation*}\etiquette[3.0.2.1]{Vbis.eq3.0.2.1}

\setcounter{paragraph}{1}
\paragraph{}\etiquette[3.0.2.2]{Vbis.par3.0.2.2}
For every arrow $[n]\to[p]$ in $\Delta$, the following diagram is
commutative:
\[
\xymatrix@C=2.25cm{
\Hom_{\Delta}\big([p],[1]\big) \ar[r] \ar[dd]&
\Hom_E(X_p,X'_p) \ar[rd]&\\
& & \Hom_E(X_p,X'_p)\\
\Hom_{\Delta}\big([n],[1]\big) \ar[r]&
\Hom_E(X_n,X'_n) \ar[ru]&.
}
\]
N.B.: The relation thus introduced on the set of morphisms from $X$
to $X'$ is not an equivalence relation. This inconvenience can be
remedied by introducing the notion of \emph{composite homotopy}
[\cf \cite{sga4.exp.Vbis.bib.3}, Chap.~III], but we shall not use that notion.

\begin{enonce*}{Lemma 3.0.2.3}\etiquette[3.0.2.3]{Vbis.lem3.0.2.3}
Let\pageoriginale $X$ and $X'$ be two objects of $\Delta^{\circ}E$,
let $f,g:X\rightrightarrows X'$ be two morphisms, and let
$F:E\to C$ be a functor, where $C$ is an abelian category. A
simplicial homotopy from $f$ to $g$ induces a homotopy between the
morphisms of cochain complexes canonically associated with $F(f)$ and
$F(g)$.
\end{enonce*}

This reduces to the case where $C$ is the category of modules over a
ring, and one applies [\cite{sga4.exp.Vbis.bib.5}, I, (3.7.1)].

\begin{enonce*}{Lemma 3.0.2.4}\etiquette[3.0.2.4]{Vbis.lem3.0.2.4}
Let $n$ be an integer, and let $X$ and $X'$ be two objects of
$\Delta^{\circ}E$. Let $f$ and $g$ be two morphisms from $X$ to $X'$
such that $f_p=g_p$ for $p<n$. Then there exists a simplicial homotopy
from $\cosk_n(f)$ to $\cosk_n(g)$.
\end{enonce*}

For $p<n$, define
$$h_p:\Hom_{\Delta}([p],[1])\to\Hom(X_p,X'_p)$$
to be the constant map with value $f_p=g_p$. Next define
$h_n:\Hom_{\Delta}([n],[1])\to\Hom(X_n,X'_n)$ by sending every
element of $\Hom_{\Delta}([n],[1])$, except $\partial^{\circ}$, to
$f_n$, and $\partial^{\circ}$ to $g_n$. Finally, observe that
$$
\Hom_{\Delta}([k],[1])=
\varprojlim_{\substack{[q]\to[k]\\q\leq n}}\Hom([q],[1])
$$
for $k>n$, which makes it possible to define $h_k$ canonically.

\subsection{Comparison of two augmentations from the viewpoint of
1-cohomological descent}\etiquette[3.1]{Vbis.subsec3.1}

Let $S$ be an object of $B$. Denote by
$\Hom_{\plat}(D^{\circ},B/S)$ the full subcategory of
$\Hom(D^{\circ},B/S)$ whose objects are the flat augmentations
[\cf (\hyperref[Vbis.subsubsec1.2.6]{1.2.6}) and
(\hyperref[Vbis.def2.4.3.1]{2.4.3.1})].

\begin{enonce*}{Proposition 3.1.1}\etiquette[3.1.1]{Vbis.prop3.1.1}
Let $X\xrightarrow{u}C^D_S$ and $X'\xrightarrow{u'}C^D_S$ be two
objects of $\Hom_{\plat}(D^{\circ},B/S)$. Let $f:X\to X'$ be a
morphism over $S$ (\ie a morphism in $\Hom(D^{\circ},B/S)$). There is
naturally associated with it a functorial morphism
$$
\eta^f_G:\sisi{\RR^+}{R}_G\bar u'_\ast\circ{\bar u'}_G^\ast
\longrightarrow
\sisi{\RR^+}{R}_G\bar u_\ast\circ\bar u_G^\ast
$$
such that the following diagram is commutative:
\[
\xymatrix@=1.5cm{
&R_G\bar u'_\ast\circ{\bar u'}_G^{\ast}
\ar[dd]^{\eta^f_G}\\
\id_{D_{G_S}^{+}(S)}\ar[ru]\ar[rd]&\\
&R_G\bar u_\ast\circ\bar u_G^{\ast}.
}
\]

Moreover,\pageoriginale if $D=\Delta$, then $\eta^f_G$ depends only on the homotopy
class \emph{(\cf (\hyperref[Vbis.subsubsec3.0.2]{3.0.2}))} of $f$ in
$\Hom(\Delta^{\circ},B/S)$.
\end{enonce*}

It is clear that, for the construction of $\eta_G^f$, we may suppose
that $G=\Ysheaf{\Mod}(\mathcal{E}^{\circ},\mathcal{O})$.

Let $I$ be the category defined by the graph
\begin{equation*}
{0}\longrightarrow{1}{\quoi}.\tag{3.1.1.1}
\end{equation*}\etiquette[3.1.1.1]{Vbis.eq3.1.1.1}
Denote by $r_0:D\to I\times D$ (\resp $r_1$) the fully faithful
functor defined by $r_0(i)=(0,i)$ (\resp $r_1(i)=(1,i)$).

By virtue of the canonical isomorphism
\begin{equation*}
\Hom((I\times D)^{\circ},B/S)\xrightarrow{\sim}
\Hom(I^{\circ},\Hom(D^{\circ},B/S))\tag{3.1.1.2}
\end{equation*}\etiquette[3.1.1.2]{Vbis.eq3.1.1.2}
the data of (\hyperref[Vbis.prop3.1.1]{3.1.1}) define a flat
augmentation
\nde{Explicitly, one has $(XfX')_{(0,i)}=X_i$ and
$(XfX')_{(1,j)}=X'_j$; to
$\alpha\in\Hom_D(i,j)=\Hom((0,i),(0,j))$ is associated
$\alpha\circ f_i:(XfX')_{(0,i)}\to(XfX')_{(1,j)}$, and the
endomorphisms of $(\varepsilon,i)$, $\varepsilon=0,1$, act in the
evident way.}
\begin{equation*}
XfX'\xrightarrow{ufu'}C^{I\times D}_S{\quoi}.\tag{3.1.1.3}
\end{equation*}\etiquette[3.1.1.3]{Vbis.eq3.1.1.3}

Let $F$ be an object of $\Mod(\Ee^\circ_S,\Oo_S)$. The canonical
morphism
\begin{equation*}
\Ee_{I\times D}^{\ast}(F)\longrightarrow
\Gamma((ufu')_\ast)\circ\Gamma((ufu')^{\ast})
(\Ee_{I\times D}^{\ast}(F)){\quoi}\tag{3.1.1.4}
\end{equation*}\etiquette[3.1.1.4]{Vbis.eq3.1.1.4}
may be interpreted as a commutative triangle in
$\Mod(\Ysheaf{\Gamma}(\bar{C^D_S}),\Oo)$:
\[\vcenter{
\xymatrix@=1.5cm{
&\Gamma(u'_\ast)\circ\Gamma({u'}^{\ast})(\Ee_D^{\ast}(F))\ar[dd]\\
\Ee_D^{\ast}(F)\ar[ru]\ar[rd]&\\
&\Gamma(u_\ast)\circ\Gamma(u^{\ast})(\Ee_D^{\ast}(F))
}}\leqno{(3.1.1.5)}
\]\etiquette[3.1.1.5]{Vbis.eq3.1.1.5}
and hence\pageoriginale gives a functorial morphism
$\bar u'_\ast\circ{\bar u'}^{\ast}
\xrightarrow{\alpha_f}\bar u_\ast\circ\bar u^{\ast}$ such that the
diagram
\[\vcenter{
\xymatrix@=1.5cm{
&\bar u'_\ast\circ{\bar u'}^{\ast}\ar[dd]^{\alpha_f}\\
\id_{\Mod(\Ee^{\circ}_S,\Oo_S)}\ar[ru]\ar[rd]&\\
&\bar u_\ast\circ\bar u^{\ast}
}}\leqno{(3.1.1.6)}
\]\etiquette[3.1.1.6]{Vbis.eq3.1.1.6}
is commutative.

One consequently obtains a functorial morphism such that the diagram
\[\vcenter{
\xymatrix@=1.5cm{
&K^{+}(\bar u'_\ast)\circ K^{+}({\bar u'}^{\ast})
\ar[dd]^{K^{+}(\alpha_f)}\\
\id_{K^{+}(\Ee^{\circ}_S,\Oo_S)}\ar[rd]\ar[ru]&\\
&K^{+}(\bar u_\ast)\circ K^{+}(\bar u^{\ast})
}}\leqno{(3.1.1.7)}
\]\etiquette[3.1.1.7]{Vbis.eq3.1.1.7}
is commutative.

\noindent
(Observe that $K^+(\alpha_f)$ can also be obtained formally by
adjunction, using the isomorphism
$K^+(\Gamma(f^{\ast}))\circ K^+({\bar u'}^{\ast})
\xrightarrow{\sim}K^+(\bar u^{\ast})$.)

Let $F^{\bigdot}$ be a complex in $K^+(S)$, and let
$\xi:K^+(\overline{ufu'}^{\ast})(F^{\bigdot})\to J^{\bigdot}$ be a
quasi-isomorphism such that, for every integer $n$, $J^n$ is totally
acyclic objectwise (\cf (\hyperref[Vbis.rem1.3.13]{1.3.13})). By
(\loccit), the canonical arrows
$$
K^+({\bar u'}^\ast)(F^{\bigdot})\to r_1^{\ast}(J^{\bigdot})
\quad\mbox{and}\quad
K^+(\bar u^\ast)(F^{\bigdot})\to r_0^{\ast}(J^{\bigdot})
$$
are\pageoriginale quasi-isomorphisms (\cf
(\hyperref[Vbis.subsubsec1.2.8]{1.2.8}) for the notation). One
therefore obtains a morphism
$$
\sisi{\RR^+}{R}\bar u'_\ast\circ{\bar u'}^\ast(F^{\bigdot})
\xrightarrow{\eta^f(F^{\bigdot})}
\sisi{\RR^+}{R}\bar u_\ast\circ\bar u^\ast(F^{\bigdot})
$$
such that the following diagram is commutative in $D^+(S)$:
{\def\objectstyle{\scriptstyle}
\def\labelstyle{\scriptstyle}

\[\vcenter{
\xymatrix@C=0cm@R=1cm{
&Q\circ K^{+}(\bar u'_\ast)(r^{\ast}_{0}(J^{\bigdot}))
\ar[rr]^{\sim}\ar@{->}'[d][ddd]&&
R\bar u'_\ast\circ
Q(r^{\ast}_0(J^{\bigdot}))\ar[ddd]\\
Q\circ K^{+}(\bar u'_\ast)\circ K^{+}({\bar u'}^{\ast})(F^{\bigdot})
\ar[rr]\ar[ru]\ar[ddd]_{Q(K^{+}(\alpha_f)(F^{\bigdot}))}&&
R\bar u'_\ast\circ
{\bar u'}^{\ast}(F^{\bigdot})\ar[ru]^{\sim}
\ar@{-->}[ddd]^(.35){\eta^{f}(F^{\bigdot})}&\\
&&&\\
&Q\circ K^{+}(\bar u_\ast)(r^{\ast}_1(J^{\bigdot}))
\ar@{->}'[r][rr]^(.25){\sim}&&
R\bar u_\ast\circ Q(r_1^{\ast}(J^{\bigdot}))\\
Q\circ K^{+}(\bar u_\ast)\circ K^{+}(\bar u^{\ast})(F^{\bigdot})
\ar[ru]\ar[rr]&&
R\bar u_\ast\circ\bar u^{\ast}(F^{\bigdot})
\ar[ru]&
}}\leqno{(3.1.1.8)}
\]\etiquette[3.1.1.8]{Vbis.eq3.1.1.8}
}
\noindent
It is clear that $\eta^f(F^{\bigdot})$ is independent of the chosen
resolution, and one verifies that it is functorial in $F^{\bigdot}$
as the latter varies in $D^{+}(S)$.

To complete the proof of \hyperref[Vbis.prop3.1.1]{(3.1.1)}, it is enough to prove the following
lemma:

\begin{enonce*}{Lemma 3.1.1.9}\etiquette[3.1.1.9]{Vbis.lem3.1.1.9}
Let\pageoriginale $X\xrightarrow{u}C^{\Delta}_S$ and
$X'\xrightarrow{u'}C^{\Delta}_S$ be two objects of
$\Hom_{\plat}(\Delta^{\circ},B/S)$. Let $f,g:X\to X'$ be two
morphisms in $\Hom(\Delta^{\circ},B/S)$. If there exists a simplicial
homotopy from $f$ to $g$, then $\eta^f=\eta^g$.
\end{enonce*}

Let
$h_n:\Hom_{\Delta}([n],[1])\to\Hom_{B/S}(X_n,{X'}_n)$ be a
simplicial homotopy from $f$ to $g$. Let $\overline{I\times D}$ be the
category obtained from $I\times D$ by adjoining the arrows
$(1,n)\to(0,n)$ in bijective correspondence with
$h_n(\Hom([n],[1]))$, together with the relations imposed by the
compatibility of the $h_n$ as $n$ varies. The data of the lemma define
a flat augmentation
$\overline{XfX'}\xrightarrow{t}C^{\overline{I\times D}}_S$. Let
$F^{\bigdot}$ be a complex in $C^+(S)$, and let $J^{\bigdot}$ be a
resolution of $C^+(t^{\ast})(F^{\bigdot})$ such that, for every
integer $n$, $J^n$ is totally acyclic objectwise. There exists a
simplicial homotopy between
$$
C^{+}(\Gamma(u'_\ast))(r^{\ast}_0(J^{\bigdot}))
\rightrightarrows
C^{+}(\Gamma(u_\ast))(r^\ast_1(J^{\bigdot})),
$$
and hence a homotopy after passage to the total complexes
(\cf (\hyperref[Vbis.lem3.0.2.3]{3.0.2.3})).

\begin{enonce*}{Definition 3.1.2}\etiquette[3.1.2]{Vbis.def3.1.2}
With the notation of \emph{\hyperref[Vbis.prop3.1.1]{(3.1.1)}}, we say that $f$ is an
\emph{equivalence for $G$-1-cohomological descent} if $\eta^f_G$ is an
isomorphism.
\end{enonce*}

\setcounter{subsubsection}{2}
\subsubsection{}\etiquette[3.1.3]{Vbis.subsubsec3.1.3}
In what follows, denote by $L_i:\Delta\to\Delta\times\Delta$
(\resp $C_i:\Delta\to\Delta\times\Delta$) the canonical functor
defined by $[n]\longmapsto[n]\times[i]$ (\resp
$[n]\longmapsto[i]\times[n]$).

\begin{enonce*}{Proposition 3.1.4}\etiquette[3.1.4]{Vbis.prop3.1.4}
Let $X\xrightarrow{u}C^{\Delta\times\Delta}_S$ and
$X'\xrightarrow{u'}C^{\Delta\times\Delta}_S$ be two objects of
$\Hom_{\plat}((\Delta\times\Delta)^{\circ},B/S)$. Let $f:X\to X'$ be
a morphism over $S$. Suppose that
$f\ast\id_{L_i}:X\circ L_i\to X'\circ L_i$ (\resp
$f\ast\id_{C_i}:X\circ C_i\to X'\circ C_i$) is an equivalence for
$G$-1-cohomological descent for every integer $i$. Then $f$ is an
equivalence for $G$-1-cohomological descent.
\end{enonce*}

By the description of $\eta^f$ (\cf the proof of
(\hyperref[Vbis.prop3.1.1]{3.1.1})), it is enough to show that a
certain morphism of triple complexes induces an isomorphism on the
cohomology of the total complexes. A standard spectral-sequence
argument then gives the conclusion.

\begin{enonce*}{Proposition 3.1.5}\etiquette[3.1.5]{Vbis.prop3.1.5}
Let\pageoriginale $X$ and $X'$ be two objects of
$\Hom_{\plat}(\Delta^{\circ},B/S)$. Let $f:X\to X'$ be a functorial
morphism such that $f_n:X_n\to X'_n$ is a morphism of
$G$-1-cohomological descent. Then the canonical morphism
$$
[[X|_fX']]\to[[X'|_{\id}X']]
$$
\emph{(\cf (\hyperref[Vbis.subsubsec1.2.7]{1.2.7}))} is an equivalence
for $G$-1-cohomological descent.
\end{enonce*}

This follows from (\hyperref[Vbis.prop3.1.4]{3.1.4}) and the following
lemma, whose proof is left to the reader:

\begin{enonce*}{Lemma 3.1.5.1}\etiquette[3.1.5.1]{Vbis.lem3.1.5.1}
Let $f:R\to S$ be a flat morphism, and let
$Y\xrightarrow{u}C^{\Delta}_S$ and
$Y\xrightarrow{v}C^{\Delta}_R$ be two \emph{flat} augmentations such
that the following diagram is commutative in
$\Hom(\Delta^{\circ},B)$:
\[
\xymatrix@=1.5cm{
Y\ar[rr]^{v}\ar[rd]_{u}&&C^{\Delta}_R\ar[ld]^f\\
&C^{\Delta}_S&
}
\]
If $v$ is an augmentation of $G$-1-cohomological descent, then it is
an equivalence for $G$-1-cohomological descent.
\end{enonce*}

In the applications, we shall combine
(\hyperref[Vbis.prop3.1.4]{3.1.4}) and
(\hyperref[Vbis.prop3.1.5]{3.1.5}) with the following result:

\begin{enonce*}{Proposition 3.1.6}\etiquette[3.1.6]{Vbis.prop3.1.6}
Let $X\xrightarrow{u}C^{\Delta\times\Delta}_S$ be a flat
augmentation. For $u$ to be an augmentation of
$G$-1-cohomological descent, it is enough that
$$
X\circ L_i\xrightarrow{u\ast\id_{L_i}}C^{\Delta}_S
\quad(\resp\ X\circ C_i\xrightarrow{u\ast\id_{C_i}}C^{\Delta}_S)
$$
be so for every integer $i$.
\end{enonce*}

This follows from the explicit construction of
$\sisi{\RR^+}{R}_G\bar u_\ast$
(\cf (\hyperref[Vbis.subsubsec2.3.10]{2.3.10})).

\begin{enonce*}{Corollary 3.1.7}\etiquette[3.1.7]{Vbis.cor3.1.7}
Let $X\xrightarrow{u}C^{\Delta}_S$ be an augmentation. Then $u$ is of
$G$-1-cohomological descent if and only if
$[[X|_{\id}X]]\to C^{\Delta\times\Delta}_S$ is so.
\end{enonce*}

% END INCLUDED SGA4_English_current_buildable_source_20260906/source/05bis/drafts/SGA4_VBIS_LINES_1930_2298_English.texfrag


% BEGIN INCLUDED SGA4_English_current_buildable_source_20260906/source/05bis/drafts/SGA4_VBIS_LINES_2299_2585_English.texfrag
\subsection{Localization criteria}\etiquette[3.2]{Vbis.subsec3.2}\pageoriginale

\begin{enonce*}{Proposition 3.2.1}\etiquette[3.2.1]{Vbis.prop3.2.1}
Let $Y\xrightarrow{v}C^{\Delta}_S$ be an augmentation of
$G$-cohomological $1$-descent. For a flat augmentation
$X\xrightarrow{u}C^{\Delta}_S$ to be of $G$-cohomological
$1$-descent, it suffices that it become so after all base changes
$Y_n\to S$.
\end{enonce*}

By means of the base changes $Y_n\to S$, one constructs a double
semi-simplicial object $\fprod{X}{Y}{S}$ augmented toward $S$. By
(\hyperref[Vbis.prop3.1.4]{3.1.4}) and
(\hyperref[Vbis.lem3.1.5.1]{3.1.5.1}), the canonical morphism
$$
\fprod{X}{Y}{S}\to[[X|_{\id}X]]
$$
is an equivalence for $G$-cohomological $1$-descent. By
(\hyperref[Vbis.cor3.1.7]{3.1.7}), it therefore suffices to show that
the augmentation of $\fprod{X}{Y}{S}$ toward $S$ is of
$G$-cohomological $1$-descent. This follows from the fact that the
canonical morphism
$$
\fprod{X}{Y}{S}\to[[Y|_{\id}Y]]
$$
is an equivalence for $G$-cohomological $1$-descent; again one uses
(\hyperref[Vbis.prop3.1.4]{3.1.4}),
(\hyperref[Vbis.lem3.1.5.1]{3.1.5.1}), and
(\hyperref[Vbis.cor3.1.7]{3.1.7}).

\setcounter{subsubsection}{1}
\subsubsection{}\etiquette[3.2.2]{Vbis.subsubsec3.2.2}
Let
\[
\vcenter{
\xymatrix@C=3cm@R=1.5cm{
Y'\ar[r]^{\mathsf g'}\ar[d]_{f'}&Y\ar[d]^{f}\\
X'\ar[r]_{\mathsf g}&X
}}\leqno{(3.2.2.1)}
\]\etiquette[3.2.2.1]{Vbis.eq3.2.2.1}
be a commutative diagram in $B$, and let $F$ be an object of
$\Mod(\mathcal E^{\circ}_Y,\mathcal O_Y)$. There is a unique
morphism
$$
\varphi(F):g_{\ast}(f^{\ast}(F))\to
f'^{\ast}(g'_{\ast}(F))
$$
such that the following diagram (in
$\Ysheaf{\Mod}(\mathcal E^{\circ},\mathcal O)$) commutes\nde{As
above, one must regard the commutative diagram as lying in
$B^\circ$ and use the fibered and cofibered structure of
$\Mod(\mathcal E^\circ,\mathcal O)\rightarrow B^\circ$, which
explains the unusual appearance of the base-change morphism. Under
the (adjunction) arrows in the following diagram we have written
their images in $B^\circ$}:
\[
\vcenter{
\xymatrix@R=1.85cm@C=1.75cm{
&f^{\ast}(F)\ar[r]^{f^\circ}\ar[ld]_{g^\circ}&F\ar[d]^{g'^\circ}\\
g_{\ast}(f^{\ast}(F))\ar[rd]_{\Id}&&g'_{\ast}(F)\\
&f'^{\ast}(g'_{\ast}(F))\ar[ru]_{f'^\circ}&
}}\leqno{(3.2.2.2)}
\]\etiquette[3.2.2.2]{Vbis.eq3.2.2.2}

If\pageoriginale all the arrows of
\eqhyperref[Vbis.eq3.2.2.1]{3.2.2.1} are now assumed flat, one
defines in the same way a morphism, called the \emph{base-change}
morphism\nde{Here again, it would be more correct to write
$f^\circ$, $g^\circ$, $f'^\circ$, $g'^\circ$ instead of $f$, $g$,
$f'$, $g'$. If one considers the nonopposite fibered and cofibered
structure, $\xi_G$ becomes
$$g_G^\ast\circ R_Gf_{\ast}\longrightarrow
R_Gf'_{\ast}\circ g_G'^\ast.$$}
\begin{equation*}
\xi_G:g_{G\ast}\circ\sisi{\RR^+}{R}_Gf^{\ast}
\longrightarrow\sisi{\RR^+}{R}_Gf'^{\ast}\circ g'_{G\ast}{\quoi}.
\tag{3.2.2.3}
\end{equation*}\etiquette[3.2.2.3]{Vbis.eq3.2.2.3}

\begin{enonce*}{Definition 3.2.2.4}\etiquette[3.2.2.4]{Vbis.def3.2.2.4}
The diagram \emph{\hyperref[Vbis.eq3.2.2.1]{(3.2.2.1)} is said to satisfy the base-change
theorem relative to $G$ if $f$, $g$, $f'$, and $g'$} are flat
morphisms and $\xi_G$ is an isomorphism.
\end{enonce*}

\setcounter{subsubsection}{2}
\subsubsection{}\etiquette[3.2.3]{Vbis.subsubsec3.2.3}
Let $h:S\to S'$ be a \emph{flat} morphism in $B$, let
$X\xrightarrow{u}C^{\Delta}_S$ and
$X'\xrightarrow{u'}C^{\Delta}_{S'}$ be two \emph{flat}
augmentations, and let $f:X\to X'$ be a functorial morphism such that
the diagram (in $\Hom(\Delta^0,B)$)
\[
\vcenter{
\xymatrix@C=3cm@R=1.5cm{
X\ar[r]^{f}\ar[d]_{u}&X'\ar[d]^{u'}\\
C^{\Delta}_S\ar[r]^{h_c}&C^{\Delta}_{S'}
}}\leqno{(3.2.3.1)}
\]\etiquette[3.2.3.1]{Vbis.eq3.2.3.1}
commutes. (Here $h_c$ denotes the constant functorial morphism
defined by $h$.)

Let $K^{\bigdot}$ be a complex in $D_G^+(S')$. The morphism
$$
\eta_G^f(K^{\bigdot}):
\sisi{\RR^+}{R}\bar u'_{\ast}\circ\bar u'^{\ast}(K^{\bigdot})
\longrightarrow
\sisi{\RR^+}{R}(\overline{h_c\circ u})_{\ast}\circ
(\overline{h_c\circ u})^{\ast}(K^{\bigdot})
$$
induces, in view of the identities
$$
\sisi{\RR^+}{R}(\overline{h_c\circ u})_{\ast}\simeq
\sisi{\RR^+}{R}h_c^{\ast}\circ
\sisi{\RR^+}{R}\bar u_{\ast}
\quad\mbox{and}\quad
(\overline{h_c\circ u})^{\ast}\simeq
\bar u^{\ast}\circ h_{c\ast},
$$
a morphism $\ch_G^{f,h}(K^{\bigdot})$, functorial in $K^{\bigdot}$,
such that the following diagram commutes\nde{For readability, we
have suppressed the reference to the thick category $G$ in the
derived functors. The same unusual derivations $h_{c\ast}$ reappear,
owing to the choice of the cofibered structure.}:
\[
\vcenter{
\xymatrix@C=1cm@R=2.25cm{
{}_{}h_{c\ast}\circ
R_{}\bar u'_{\ast}{}_{}\circ
\bar u'^{\ast}(K^{\bigdot})
\ar[rr]^{\ch_G^{f,h}{(K^{\bigdot})}}&&
R_{}\bar u_{\ast}\circ
{}_{}\bar u^{\ast}\circ{}_{}h_{c\ast}(K^{\bigdot})\\
&h_{c\ast}(K^{\bigdot})
\ar[lu]^{{}_{}h_{c\ast}(\alpha(K^{\bigdot}))}\ar[ru]&.
}}\leqno{(3.2.3.2)}
\]\etiquette[3.2.3.2]{Vbis.eq3.2.3.2}

\begin{enonce*}{Lemma 3.2.3.3}\etiquette[3.2.3.3]{Vbis.lem3.2.3.3}
Suppose\pageoriginale that, for every integer $i$, the diagram
\[
\xymatrix@C=3cm@R=1.5cm{
X_i\ar[d]_{u_i}\ar[r]^{f_i}&X'_i\ar[d]^{u'_i}\\
S\ar[r]^{h}&S'
}
\]
satisfies the base-change theorem relative to $G$. Then
$\ch_G^{f,h}$ is an isomorphism.
\end{enonce*}

The morphism $\ch_G^{f,h}(K^{\bigdot})$ may be computed as follows.
Consider the augmentation
$$
X f X'\xrightarrow{(h_c\circ u)f u'}C^{I\times\Delta}_S
$$
(\cf \eqhyperref[Vbis.eq3.1.1.2]{3.1.1.2}), and choose a resolution
$J^{\bigdot}$ of
$K^+(\overline{(h_c\circ u)f u'}^{\ast})(K^{\bigdot})$ such that,
for every integer $n$, $J^n$ is totally acyclic object by object.

One obtains a quasi-isomorphism
$$
\theta:K^+(\Gamma(f^{\ast}))(r_0^{\ast}J^{\bigdot})
\longrightarrow r_1^{\ast}(J^{\bigdot})
$$
and a morphism
$$
t:K^+(\Gamma(h_c^{\ast}))\circ K^+(\Gamma(u'_{\ast}))
(r_0^{\ast}(J^{\bigdot}))\longrightarrow
K^+(\Gamma(u_{\ast}))(r_1^{\ast}(J^{\bigdot}))
$$
which has the factorization
\begin{align*}
&K^+(\Gamma(h_c^{\ast}))\circ K^+(\Gamma(u'_{\ast}))
(r_0^{\ast}(J^{\bigdot}))\xrightarrow{\varphi}
K^+(\Gamma(u_{\ast}))\circ K^+(\Gamma(f^{\ast}))
(r_0^{\ast}(J^{\bigdot}))
\xrightarrow{K^+(\Gamma(u_{\ast}))(\theta)}\\
&\longrightarrow K^+(\Gamma(u_{\ast}))(r_1^{\ast}(J^{\bigdot})),
\end{align*}
and\pageoriginale $\ch^{f,h}(K^{\bigdot})$ is obtained by taking the
image of $t$ under the functor
$\sisi{\RR^+}{R}\epsilon_{Z^{\bigdot\ast}}$.

The preceding factorization shows that $t$ identifies ``object by
object'' with the base-change morphisms associated with the diagrams
\[
\xymatrix@C=3cm@R=1.5cm{
X_i\ar[d]_{u_i}\ar[r]^{f_i}&X'_i\ar[d]^{u'_i}\\
S\ar[r]^{h}&S'
}
\]
It follows, by spectral sequences, that
$\ch_G^{f,h}(K^{\bigdot})$ is an isomorphism, completing the proof.

This leads to a \emph{second localization criterion}:

\begin{enonce*}{Proposition 3.2.4}\etiquette[3.2.4]{Vbis.prop3.2.4}
Let $X\xrightarrow{u}C^{\Delta}_S$ be an object of
$\Hom_{\plat}(\Delta^{\circ},B/S)$. For $u$ to be an augmentation of
$G$-cohomological $1$-descent, it suffices that it become so after
``enough for $G$''\footnote{The expression ``enough for $G$'' means
that there exists a family
$(S_{\alpha}\xrightarrow{h_{\alpha}}S)_{\alpha\in A}$ of flat
morphisms satisfying the preceding conditions and such that the
family of functors
$(h_{\alpha\ast}:G_S\to G_{S_\alpha})_{\alpha\in A}$ is
\emph{conservative}.} flat base changes $h:S'\to S$ for which the
cartesian diagrams
\[
\xymatrix@C=3cm@R=1.5cm{
X'_i\ar[d]\ar[r]&X_i\ar[d]\\
S'\ar[r]^{h}&S
}
\]
satisfy the base-change theorem relative to $G$.
\end{enonce*}

Proposition\pageoriginale (\hyperref[Vbis.prop3.2.4]{3.2.4}) follows
immediately from (\hyperref[Vbis.lem3.2.3.3]{3.2.3.3}), in view of
diagram \eqhyperref[Vbis.eq3.2.3.2]{3.2.3.2}.

\subsection{Properties of morphisms of cohomological descent}
\etiquette[3.3]{Vbis.subsec3.3}

In this section, we assume that the \emph{class of flat morphisms in
$B$ is stable under base change}.

\begin{enonce*}{Proposition 3.3.1}\etiquette[3.3.1]{Vbis.prop3.3.1}
Let $i=1,2$.
\begin{enumerate}
\renewcommand{\theenumi}{\alph{enumi}}
\renewcommand{\labelenumi}{\theenumi)}
\item Every flat morphism which has a section is a morphism of
universal $G$-cohomological $2$-descent.

\item Let $f:X\to S$ be a flat morphism, let $g:S'\to S$ be a
morphism of $G$-cohomological $i$-descent, put
$X'=\fprod{X}{S'}{S}$, and let $f'=f_{(S')}:X'\to S'$. Then $f$ is
of universal $G$-cohomological $i$-descent if and only if $f'$ is.

\item If the composite of two morphisms $f:X\to Y$ and $g:Y\to Z$
is of universal $G$-cohomological $i$-descent, then $g$ is of
universal $G$-cohomological $i$-descent.

\item The composite of two morphisms of universal
$G$-cohomological $i$-descent is a morphism of universal
$G$-cohomological $i$-descent.

\item If $f:X\to X'$ and $g:Y\to Y'$ are two $S$-morphisms of
universal $G$-cohomological $i$-descent, then $\fprod{f}{g}{S}$ is
of universal $G$-cohomological $i$-descent.

\item Let $(u_{\alpha}:X_{\alpha}\to Y_{\alpha})_{\alpha\in A}$
be a family of morphisms of $G$-cohomological $i$-descent. Then
$\amalg u_{\alpha}:\amalg_{\alpha\in A}X_{\alpha}\to
\amalg_{\alpha\in A}Y_{\alpha}$ is a morphism of
$G$-cohomological $i$-descent.
\end{enumerate}
\end{enonce*}

\noindent\emph{Proof.}
For\pageoriginale effective descent relative to $G$, we refer to
\cite{sga4.exp.Vbis.bib.4}.
\begin{enumerate}
\renewcommand{\theenumi}{\alph{enumi}}
\renewcommand{\labelenumi}{\theenumi)}
\item Using a section $s$ of $f$, compare the semi-simplicial objects
augmented toward $S$, $[X|_fS]$ and $[S|_{\id}S]$. By
(\hyperref[Vbis.lem3.0.2.4]{3.0.2.4}), one obtains an equivalence for
$G$-cohomological $1$-descent\nde{For this, note that $[X|_fS]$ and
$[S|_{\id}S]$ are simply the coskeleta of order $0$ of $f$ and
$\id$, respectively, so that $s$ and $f$ induce inverse
isomorphisms in cohomology by
\hyperref[Vbis.prop3.1.1]{3.1.1}.}.

\item This follows from (\hyperref[Vbis.prop3.2.1]{3.2.1}).

\item This follows from a) and b).

\item Let $f:X\to Y$ and $g:Y\to S$ be two morphisms of universal
$G$-cohomological $1$-descent. Put $h=g\circ f$ and consider the
fiber product
\[
\xymatrix@C=3cm@R=1.5cm{
Y\ar[d]_{g}&R\ar[d]^{g'}\ar[l]_{h'}\\
S&X\ar[l]^{h}
}
\]
The morphism $g'$ has a section $s:X\to R$ such that
$h'\circ s=f$. By c), $h'$ is of universal $G$-cohomological
$1$-descent, and the same holds for $h$ by b).

\item This follows formally from b) and d).

\item Let $(Z_{\lambda})_{\lambda\in\Lambda}$ be a family of
objects of $B$. It follows from hypotheses
(\hyperref[Vbis.subsubsec3.0.0]{3.0.0}) that there is a canonical
equivalence
$$
\Mod\left(\mathcal E^{\circ}_{\coprod_{\lambda}Z_{\lambda}},
\mathcal O_{\coprod_{\lambda}Z_{\lambda}}\right)
\xrightarrow{\sim}\displaystyle\mathop{\prod}_{\lambda}
\Mod(\mathcal E^{\circ}_{Z_\lambda},\mathcal O_{Z_\lambda})
$$
\end{enumerate}
which induces an equivalence
$$
G_{\coprod_{\lambda}Z_{\lambda}}
\xrightarrow{\sim}\prod_{\lambda}G_{Z_{\lambda}}{\quoi}.
$$
Moreover, a family
$(Q_{\lambda})_{\lambda}\in
\prod_{\lambda}\Mod(\mathcal E^{\circ}_{Z_\lambda},
\mathcal O_{Z_\lambda})$ corresponds to a totally acyclic object if
and only if $Q_{\lambda}$ is totally acyclic for every $\lambda$.

Now let
$(u_{\lambda}:X_{\lambda}\to Y_{\lambda})_{\lambda\in\Lambda}$ be a
family of morphisms of $G$-cohomological $1$-descent. Since direct
sums in $B$ are disjoint and\pageoriginale universal, for every $n$
there is an identification
$$
\left[\displaystyle\mathop{\coprod}_{\lambda}X_{\lambda}
\mathrel{|}_{\coprod u_{\lambda}}\coprod Y_{\lambda}\right]_n
\simeq\displaystyle\mathop{\coprod}_{\lambda}
[X_{\lambda}|_{u_{\lambda}}Y_{\lambda}]_n
$$
which is functorial in $n$. It then follows from the preceding
remarks and the explicit calculation of
$\sisi{\RR^+}{R}_G\theta_{\ast}$ (\cf \hyperref[Vbis.subsec2.3]{(2.3)}) that
$\coprod_{\lambda}u_{\lambda}$ is a morphism of
$G$-cohomological $1$-descent; the details are left to the reader.
The reader may also verify that
$\coprod_{\lambda}u_{\lambda}$ \emph{remains of $G$-cohomological
$1$-descent after every base change if the same is true of every
$u_{\lambda}$}.

\setcounter{subsubsection}{1}

% END INCLUDED SGA4_English_current_buildable_source_20260906/source/05bis/drafts/SGA4_VBIS_LINES_2299_2585_English.texfrag


% BEGIN INCLUDED SGA4_English_current_buildable_source_20260906/source/05bis/drafts/SGA4_VBIS_LINES_2586_2816_English.texfrag
\subsubsection{}\etiquette[3.3.2]{Vbis.subsubsec3.3.2}
Associate with each object $X$ of $B$ the set of families of
morphisms $(X_\alpha\to X)_{\alpha\in A}$ such that
$\coprod_\alpha X_\alpha\to X$ is a morphism of universal
$G$-$i$-cohomological descent. By \hyperref[Vbis.prop3.3.1]{3.3.1},
this defines a pretopology on $B$.

The topology generated by this pretopology is called the
\emph{topology of $G$-$i$-cohomological descent}. It follows from
\hyperref[Vbis.prop3.3.1]{3.3.1} c) that the covering morphisms for
this topology are precisely the morphisms of universal
$G$-$i$-cohomological descent. Finally, note that the topology of
$G$-$2$-cohomological descent is coarser than the topology of
$G^{\circ}$-descent [\cf \cite{sga4.exp.Vbis.bib.4} (6.23)].

\begin{enonce*}{Theorem 3.3.3}\etiquette[3.3.3]{Vbis.thm3.3.3}
Suppose that \emph{all arrows of $B$ are flat}. Let $S$ be an object
of $B$. Every hypercovering of $S$ for the topology of
$G$-$1$-cohomological descent whose objects are all representable
\emph{(\cf V, Appendix)} defines an augmentation of universal
$G$-$1$-cohomological descent.
\end{enonce*}

The proof proceeds in two steps; we specify that all coskeleta are
to be computed in the category $B/S$.

\begin{enonce*}{Lemma 3.3.3.1}\etiquette[3.3.3.1]{Vbis.lem3.3.3.1}
Let\pageoriginale $X$ be an object of
$\Hom(\Delta^{\circ},B/S)$. For $X$ to be of
$1$-cohomological descent (resp. universal $1$-cohomological
descent), it is enough that $\cosk_n(X)$ have this property for all
sufficiently large $n$.
\end{enonce*}

Indeed, let $\alpha_n:X\to\cosk_n(X)$ be the canonical morphism. One
checks that if $K^{\bigdot}$ is a complex of $D^+_{G_S}(S)$ such
that $H^j(K^{\bigdot})=0$ for $j<N$, then
$H^i(\eta_G^{\alpha_n})$ is an isomorphism for $i<N+n$, whence the
assertion.

\begin{enonce*}{Lemma 3.3.3.2}\etiquette[3.3.3.2]{Vbis.lem3.3.3.2}
Let $n\geq0$, let $X$ and $X'$ be two simplicial objects of $B/S$,
and let $f:X\to X'$ be a morphism. Suppose that:
\begin{enumerate}
\renewcommand{\theenumi}{\roman{enumi}}
\renewcommand{\labelenumi}{(\theenumi)}
\item $X\to\cosk_{n+1}(X)$ is an isomorphism;

\item $X'\to\cosk_{n+1}(X')$ is an isomorphism;

\item $f_i:X_i\to X'_i$ is an isomorphism for $i\leq n$;

\item $f_{n+1}$ is a morphism of universal
$G$-$1$-cohomological descent.
\end{enumerate}
Then, if $X$ is of universal $G$-$1$-cohomological descent, the same
is true of $X'$.
\end{enonce*}

Lemmas \hyperref[Vbis.lem3.3.3.1]{3.3.3.1} and
\hyperref[Vbis.lem3.3.3.2]{3.3.3.2} clearly prove Theorem
\hyperref[Vbis.thm3.3.3]{3.3.3} by induction.

\begin{enonce*}{Lemma 3.3.3.3}\etiquette[3.3.3.3]{Vbis.lem3.3.3.3}
Under the hypotheses of \hyperref[Vbis.lem3.3.3.2]{3.3.3.2}, all
the morphisms $f_p$ are morphisms of universal
$G$-$1$-cohomological descent.
\end{enonce*}

This is trivial for $p\leq n+1$. For $p>n+1$, $X_p$ (resp. $X'_p$)
may be written as
$\varprojlim_{\substack{\Delta^+\\ n+1[p]}}X_q$
(resp. $\varprojlim_{\substack{\Delta^+\\ n+1[p]}}X'_q$).

We may suppose that $p>n+1$. Write
$$
\cosq_{n+1}(X)_p=
\varprojlim_{\substack{[q]\rightarrow[p]\\q\leq n+1}}X_q
$$
as the kernel of the pair of arrows
$$
\Pi X=\prod_{\substack{[q]\rightarrow[p]\\q\leq n+1}}X_q
\rightrightarrows
\prod_{\substack{[i]\stackrel{\alpha}{\hookrightarrow}[j]\\
\searrow\swarrow\\{}[p]\\j\leq n+1}}X_i.
$$
The component $\alpha_X$ indexed by
$\alpha\in\Hom_{[p]}([i],[j])$ of this pair consists, on the one
hand, of the projection $\Pi X\to X_i$ indexed by
$[i]\to[p]$, and, on the other hand, of the composite
$$
\Pi X\longrightarrow X_j\longrightarrow X_i
$$
of the projection indexed by $[j]\to[p]$ with
$\alpha\in\Hom(X_j,X_i)$.

It follows that the diagram
$$
\xymatrix{\cosq_{n+1}(X)_p\ar[r]\ar[d]&\Pi X\ar[d]\\
\cosq_{n+1}(X')_p\ar[r]&\Pi X'}
$$
is Cartesian. Indeed, by the preceding description, it is enough to
prove that for every \textit{increasing injection}
$\alpha\in\Hom_{[p]}([i],[j])$, the diagram
$$
\xymatrix{\Ker(\alpha_X:\Pi X\rightrightarrows X_i)\ar[r]\ar[d]&
\Pi X\ar[d]\\
\Ker(\alpha_{X'}:\Pi X'\rightrightarrows X'_i)\ar[r]&\Pi X'}
$$
is Cartesian. If $i=n+1$, then
$\alpha:[i]\hookrightarrow[j]$ is the identity of $[n+1]$, so the
pair of arrows imposes no condition. The respective kernels are
$\Pi X$ and $\Pi X'$, which settles this case. Otherwise $i<n+1$
and $X_i\to X'_i$ is an isomorphism, so the diagram is Cartesian in
this case as well (testing an equality in $X_i$ or $X'_i$ amounts to
the same thing)\nde{The proof of Lemma
\hyperref[Vbis.lem3.3.3.3]{3.3.3.3} has been slightly rewritten.}.

Using \hyperref[Vbis.prop3.3.1]{3.3.1}, the relevant morphisms are
therefore morphisms of universal $G$-$1$-cohomological descent. Let
$[X'/X]^p$ be the $(p+1)$-fold fiber product of $X'$ over $X$. By
\hyperref[Vbis.lem3.3.3.3]{3.3.3.3}, it is enough to verify
\hyperref[Vbis.lem3.3.3.2]{3.3.3.2} after a base change
$[X'/X]^p\to X$. After such a base change, the hypotheses of
\hyperref[Vbis.lem3.3.3.2]{3.3.3.2} still hold and, moreover,
$f_{n+1}$ admits a section. Lemma
\hyperref[Vbis.lem3.0.2.4]{3.0.2.4} then completes the proof.

For effective descent, we have the following result.

\begin{enonce*}{Proposition 3.3.4}\etiquette[3.3.4]{Vbis.prop3.3.4}
Let $X:\Delta^{\circ}\to B/S$ be a functor. Suppose that
$X_0\to S$ is a morphism of universal $G^{\circ}$-$2$-descent, as
are the morphisms $X_{n+1}\to(\cosk_n(X))_{n+1}$ for $n=0,1$. Then
the functor $X$ is $G^{\circ}$-$2$-faithful and remains so after
every base change (in other words, $X$ defines an augmentation of
universal effective descent in the sense of
\hyperref[Vbis.def2.4.11]{2.4.11}).
\end{enonce*}

By \cite{sga4.exp.Vbis.bib.4} (7.12), it is enough to show that $i_2^{\ast}(X)$ is
$G^{\circ}$-$2$-faithful; for this we use the following lemmas.

\begin{enonce*}{Lemma 3.3.4.1}\etiquette[3.3.4.1]{Vbis.lem3.3.4.1}
Let\pageoriginale $X$ and $X'$ be two functors
$\mathop{\Delta}_2^{\circ}\to B/S$, and consider a commutative
diagram
\[
\xymatrix@R=1.5cm@C=2cm{
X'_2 \ar[r]\ar@<8pt>[r]\ar@<-8pt>[r]\ar[d]_{k}&X'_1\ar@<4pt>[l]
\ar@<-4pt>[l]\ar@<4pt>[r]\ar@<-4pt>[r]\ar@{=}[d]&X'_0\ar[l]
\ar[r]\ar@{=}[d]&S\ar@{=}[d]\\
X_2 \ar[r]\ar@<8pt>[r]\ar@<-8pt>[r]&X_1\ar@<4pt>[l]
\ar@<-4pt>[l]\ar@<4pt>[r]\ar@<-4pt>[r]&X_0\ar[l]\ar[r]&S
}
\]
such that $k$ is a morphism of $G^{\circ}$-$0$-descent. If $X$ is
$G^{\circ}$-$2$-faithful, then so is $X'$.
\end{enonce*}

This is evident.

\begin{enonce*}{Lemma 3.3.4.2}\etiquette[3.3.4.2]{Vbis.lem3.3.4.2}
Let $X$ and $X'$ be two functors
$\mathop{\Delta}_2^{\circ}\to B/S$, and consider a commutative
diagram
\def\labelstyle{\scriptstyle}
\[
\xymatrix@R=2cm@C=2cm{
X'_2 \ar[r]^{\partial^{1}} \ar@<17pt>[r]^{\partial^{0}}
\ar@<-14pt>[r]_{\partial^{2}} \ar[d]_{f_2}&X'_1\ar@<8pt>[l]
\ar@<-11pt>[l]\ar@<11pt>[r]^{\partial^{0}}\ar@<-8pt>[r]_{\partial^{1}}
\ar[d]^{f_1}&X'_0\ar[l]_{\sigma^{0}}\ar[r]\ar@{=}[d]&S\ar@{=}[d]\\
X_2 \ar[r]\ar@<12pt>[r]\ar@<-12pt>[r]&X_1\ar@<6pt>[l]
\ar@<-6pt>[l]\ar@<6pt>[r]\ar@<-6pt>[r]&X_0\ar[l]\ar[r]&S
}
\]
such that $f_1$ is of $G^{\circ}$-$1$-descent and $f_2$ is of
$G^{\circ}$-$0$-descent. Suppose in addition that
$X'_2\xrightarrow{\sim}(\cosk_1(X'))_2$. If $X$ is
$G^{\circ}$-$2$-faithful, then so is $X'$.
\end{enonce*}

Let $X''_1=\fprod{X'_1}{X'_1}{X_1}$, and let
$\Ysheaf{\partial}_0:X''_1\to X'_1$ and
$\Ysheaf{\partial}_1:X''_1\to X'_1$ be the two projections.

The definition of a coskeleton gives an arrow
$\varphi:X''_1\to X'_2$ such that
\begin{equation*}
\begin{cases}
\partial_0\circ\varphi=\Ysheaf{\partial}_0,\\
\partial_1\circ\varphi=\Ysheaf{\partial}_1,\\
\partial_2\circ\varphi=\sigma_0\partial_1\Ysheaf{\partial}_0
=\sigma_0\partial_1\Ysheaf{\partial}_1{\quoi}.
\end{cases}
\end{equation*}

Now take a descent datum on $X'$: it gives an object on $X_0$, two
objects on $X_1$, and an isomorphism between them on $X'_1$. The
arrow $\varphi$ shows that the two inverse images of these
isomorphisms on $X''_1$ are equal. Since $f_1$ is of
$1$-descent, we obtain an isomorphism between the two objects on
$X_1$; this isomorphism is a descent datum (because
$X'_2\to X_2$ is of $0$-descent), and the result follows.

\begin{enonce*}{Corollary 3.3.5}\etiquette[3.3.5]{Vbis.cor3.3.5}
Suppose that \emph{all arrows of $B$ are flat}. Every hypercovering
of $S$ for the topology of $G$-$2$-cohomological descent whose
objects are representable defines an augmentation of universal
$G$-$2$-cohomological descent.
\end{enonce*}

% END INCLUDED SGA4_English_current_buildable_source_20260906/source/05bis/drafts/SGA4_VBIS_LINES_2586_2816_English.texfrag


% BEGIN INCLUDED SGA4_English_current_buildable_source_20260906/source/05bis/drafts/SGA4_VBIS_LINES_2817_2992_English.texfrag
\section{Examples}\etiquette[4]{Vbis.sec4}

\subsection{Sheaves of abelian groups on topological
  spaces}\etiquette[4.1]{Vbis.subsec4.1}

\setcounter{subsubsection}{-1}
\subsubsection{}\etiquette[4.1.0]{Vbis.subsubsec4.1.0}
In this section, Top denotes the category of topological spaces
[elements of the fixed universe $\Uu$]. We define a category
$\mathcal{E}$ bifibered in duals of topoi over Top by taking as
objects the pairs $(X, F)$, where $X$ is a topological space and $F$
is a sheaf of sets on $X$, and as morphisms the pairs
$(f, \varphi) : (X, F) \to (Y, H)$, where $f : X \to Y$ is a
continuous map and $\varphi : H \to f_{\ast} (F)$ is a morphism of
sheaves. We take for $\mathcal{O}$ the section of $\mathcal{E}^0$
over $\Top^0$ that associates with each topological space the
constant sheaf $\ZZ_X$: we set
$G = \Ysheaf{\Mod}(\mathcal{E}^0, \mathcal{O})$, so that, for every
topological space $X$, $G_X$ identifies with the category of sheaves
of abelian groups on $X$.

Recall that a morphism $f : X \to Y$ is \emph{separated} if the
diagonal in $\fprod{X}{X}{Y}$ is closed. A \emph{proper} morphism is
a separated and universally closed morphism (note that this
definition is more restrictive than Bourbaki's).

The\pageoriginale proof of the ``base-change theorem'' below is
inspired by (\cite{sga4.exp.Vbis.bib.5} II(4.11.1)).

\begin{enonce*}{Theorem 4.1.1}\etiquette[4.1.1]{Vbis.thm4.1.1}
Let $f : X \to Y$ be a proper morphism, let $F$ be an abelian sheaf
on $X$, and consider a Cartesian diagram
\[
\xymatrix@C=2.5cm@R=1.5cm{
X' \ar[r]^{g'} \ar[d]^{f'}& X\ar[d]^{f}\\
Y' \ar[r]^{g}& Y.
}
\]
Then the canonical map \emph{(XII \hyperref[XII.eq4.2]{4.2})} or
\emph{\eqhyperref[Vbis.eq3.2.2.3]{3.2.2.3}}
$$
g^{\ast} R^i f_{\ast} F \xrightarrow{\sim} R^i f'_{\ast} {g'}^{\ast}F
$$
is an isomorphism\nde{Unlike above, these are the derived functors
of the usual functors, and not those arising from the fibered-category
structure $G^\circ\rightarrow\Top^\circ$.}.
\end{enonce*}

This theorem is equivalent to the following corollary (the corollary
is obtained by taking $Y'= (\Point)$; the theorem is obtained by
applying the corollary to $f$ and $f'$).

\begin{enonce*}{Corollary 4.1.2}\etiquette[4.1.2]{Vbis.cor4.1.2}
Let $f : X \to Y$ be a proper morphism, let $F$ be an abelian sheaf
on $X$, and let $y \in Y$. The canonical map
$$
(R^i f_{\ast} F)_y \longrightarrow H^i(f^{-1} (y), F)
$$
is bijective.
\end{enonce*}

By definition, as $U$ ranges over the neighborhoods of $y$, we have
$$
(R^if_{\ast} F)_y = \varinjlim H^i(f^{-1} (U), F){\quoi}.
$$
Since $f$ is closed, the $f^{-1}(U)$ form a fundamental system of
neighborhoods of $f^{-1}(y)$, and the result follows from the
following lemma, (\hyperref[Vbis.lem4.1.3]{4.1.3}).

\begin{enonce*}{Lemma 4.1.3}\etiquette[4.1.3]{Vbis.lem4.1.3}
Let\pageoriginale $X$ be a topological space, let $K \subset X$, and
let $F$ be an abelian sheaf on $X$. Suppose that $K$ is compact and
that any two distinct points of $K$ have disjoint neighborhoods in
$X$. Then, as $U$ ranges over the neighborhoods of $K$, we have
\begin{equation*}
\varinjlim H^i(U, F) \xrightarrow{\sim} H^i(K, F){\quoi}.\tag{4.1.3.1}
\end{equation*}\etiquette[4.1.3.1]{Vbis.eq4.1.3.1}
\end{enonce*}

We first treat the case $i = 0$. In this case, it is clear that
\eqhyperref[Vbis.eq4.1.3.1]{4.1.3.1} is injective. For surjectivity,
we use the following result.

\begin{enonce*}{Lemma 4.1.4}\etiquette[4.1.4]{Vbis.lem4.1.4}
Under the hypotheses of \emph{(\hyperref[Vbis.lem4.1.3]{4.1.3})}, if
$A$ and $B$ are two closed subsets of $K$ and $W$ is a neighborhood
of $A \cap B$ in $X$, then there are neighborhoods $U$ and $V$ of
$A$ and $B$ in $X$ such that $U \cap V \subset W$.
\end{enonce*}

We first treat the case where $A$ consists of a single point $a$.
The assertion is trivial if $a \in B$ (take $U = W$). If
$a \not \in B$, then for every $b \in B$ there are disjoint open
neighborhoods $U_b$ and $V_b$ of $a$ and $b$ in $X$. We take for
$V$ a finite union of the $V_b$ that covers $B$, and for $U$ the
intersection of the corresponding $U_b$.

In the general case, for every $a \in A$, there are open
neighborhoods $U_a$ and $V_a$ of $a$ and $B$ in $X$ such that
$U_a \cap V_a \subset W$. We take for $U$ a finite union of the
$U_a$ that covers $A$, and for $V$ the corresponding intersection
of the $V_a$.

Return to (\hyperref[Vbis.lem4.1.3]{4.1.3}). If
$s \in H^{\circ}(K, F)$, there is an open covering $U_i$ of $K$ in
$X$ and there are sections $s_i \in H^{\circ}(U_i, F)$ such that
$s_i = s$ on $U_i \cap K$. We may suppose that there are only
finitely many $U_i$. Let $(K_i)$ be a closed covering of $K$, with
$K_i \subset K \cap U_i$. Let $W_y$ be the open subset of
$U_i \cap U_j$ on which $s_i = s_j$; we have
$K_i \cap K_j \subset W_{ij}$. Apply
(\hyperref[Vbis.lem4.1.4]{4.1.4}) to all pairs $(K_i, K_j)$ and to
the $W_{ij}$: we find open subsets $V_{ij}$ with
$K_i \subset V_{ij} \subset U_i$ and
$V_{ij} \cap V_{ji} \subset W_{ij}$. Let
$U'_i = \displaystyle\mathop{\cap}_j V_{ij}$: then
$K_i \subset U'_i \subset U_i$, and $s_i = s_j$ on
$U'_i \cap U'_j$. The $s_i$ therefore glue\pageoriginale on the
neighborhood of $K$ given by the union of the $U'_i$, and
\eqhyperref[Vbis.eq4.1.3.1]{4.1.3.1} is surjective (hence bijective)
for $i = 0$\nde{The two meanings of the symbol $i$ should not be
confused.}.

\begin{enonce*}{Lemma 4.1.5}\etiquette[4.1.5]{Vbis.lem4.1.5}
If $F$ is flabby, the sheaf $F | K$ on $K$ is soft.
\end{enonce*}

Let $A \subset K$ be closed in $K$. Every section of $F$ on $A$
extends to a neighborhood, by what precedes. Since $F$ is flabby, it
extends to $X$, and hence in particular to $K$.

We now prove (\hyperref[Vbis.lem4.1.3]{4.1.3}). Let $F^{\ast}$ be a
flabby resolution of $F$. We have
\begin{eqnarray*}
\varinjlim H^i (U, F) &\stackrel{\ndemark}{=}& \varinjlim H^i (\Gamma(U, F^{\ast})) = H^i\varinjlim \Gamma(U,F^{\ast})
\stackrel{\ndemark}{=} H^i \Gamma (K, F^{\ast})\\
& \stackrel{\text{\hyperref[Vbis.lem4.1.5]{(4.1.5)}}}{=} & H^i (K, F){\quoi}.
\end{eqnarray*}
\addtocounter{cnote}{-1}\ndetext{$H^i (\Gamma(U, F^{\ast}))$ denotes here the $i$-th cohomology group of the complex $\Gamma(U,F^\ast)$.}\addtocounter{cnote}{1}
\ndetext{By \hyperref[Vbis.lem4.1.3]{(4.1.3)}, applied with $\Gamma=H^0$}
\begin{enonce*}{Corollary 4.1.6}\etiquette[4.1.6]{Vbis.cor4.1.6}
Let $f : X \to Y$ be a proper surjective morphism of topological
spaces. Then $f$ is of $G$-$1$-cohomological descent.
\end{enonce*}

This follows from (\hyperref[Vbis.cor4.1.2]{4.1.2}) and
(\hyperref[Vbis.prop3.2.4]{3.2.4}).

\begin{enonce*}{Corollary 4.1.7}\etiquette[4.1.7]{Vbis.cor4.1.7}
Let $\Uu = (U_i)_{i \in I}$ be a locally finite closed covering of a
topological space $X$. Then the canonical morphism
$f : \coprod_{i\in I} U_i \to X$ is a morphism of
$G$-$1$-cohomological descent.
\end{enonce*}

Indeed, $f$ is proper and surjective. Applying
(\hyperref[Vbis.prop2.5.6]{2.5.6}), we recover the spectral sequence
of [\cite{sga4.exp.Vbis.bib.5}. II (5.2.4)] (\cf the Introduction).

\begin{enonce*}{Proposition 4.1.8}\etiquette[4.1.8]{Vbis.prop4.1.8}
Let $f : X \to Y$ be a surjective local homeomorphism of topological
spaces. Then $f$ is a morphism of universal
$G$-$1$-cohomological descent.
\end{enonce*}

By localization on $Y$ (\cf (\hyperref[Vbis.prop3.2.4]{3.2.4})), we
are reduced to the case where $f$ has a section.

\begin{enonce*}{Corollary 4.1.9}\etiquette[4.1.9]{Vbis.cor4.1.9}
Let\pageoriginale $\Uu = (U_i)_{i \in I}$ be an open covering of a
topological space $X$. Then the canonical morphism
$f : \coprod_{i \in I} U_i \to X$ is a morphism of universal
$G$-$1$-cohomological descent.
\end{enonce*}

The descent spectral sequence then gives [\cite{sga4.exp.Vbis.bib.5} II (5.4.1)]
(\cf the Introduction).

\begin{enonce*}[remark]{Remark 4.1.10}\etiquette[4.1.10]{Vbis.rem4.1.10}
%%%%%%%%%%%%%%%%%%%
%%%%%%%%%%%%Check that this compiles
%%%%%%%%%%%%%%%%%%%%
The proofs of (\hyperrefext[VIII][VIII.sec9]{9}) show that
``1-cohomological descent'' may be replaced by
``2-cohomological descent'' in all the preceding statements.
\end{enonce*}

\subsection{Quasi-coherent modules on schemes}\etiquette[4.2]{Vbis.subsec4.2}

\setcounter{subsubsection}{-1}

% END INCLUDED SGA4_English_current_buildable_source_20260906/source/05bis/drafts/SGA4_VBIS_LINES_2817_2992_English.texfrag


% BEGIN INCLUDED SGA4_English_current_buildable_source_20260906/source/05bis/drafts/SGA4_VBIS_LINES_2993_3150_English.texfrag
\subsubsection{}\etiquette[4.2.0]{Vbis.subsubsec4.2.0}
Let $(\Sch)_s$ be the category whose objects are the schemes that are
elements of $\Uu$ and whose arrows are the quasi-compact and quasi-separated
morphisms. We define a category $\mathcal{E}$ bifibered in duals of topoi over
$(\Sch)_s$ by taking as objects the pairs $(X,F)$, where $X$ is a scheme and
$F$ a sheaf on $X$ (for the Zariski topology), and as morphisms the pairs
$(f,\varphi):(X,F)\to(Y,G)$, where $f:X\to Y$ is a continuous map and
$\varphi:G\to f_{\ast}(F)$ is a morphism of sheaves. For $\mathcal{O}$ we
take the section of $\mathcal{E}^0$ over $(\Sch)^0_s$ that associates with
every scheme its structure sheaf; for $G$ we take the subcategory of
$\Ysheaf{\Mod}(\mathcal{E}^0,\mathcal{O})$ such that, for every scheme $X$,
$G_X$ is the category of quasi-coherent modules on $X$.

\begin{enonce*}{Proposition 4.2.1}\etiquette[4.2.1]{Vbis.prop4.2.1}
Every faithfully flat, quasi-compact, and quasi-separated morphism
$f:X\to Y$ is a morphism of universal $G$-$2$-cohomological descent.
\end{enonce*}

By [\SGA I (VIII.5.2)], $f$ is a morphism of effective
$G^0$-descent. By [\EGA IV (1.7.21)], we may apply
(\hyperref[Vbis.prop3.2.4]{3.2.4}) to the Cartesian square
\[
\xymatrix@C=3cm@R=1.5cm{
  \fprod{X}{X}{Y} \ar[r] \ar[d]& X \ar[d]^{f}\\
X   \ar[r]_{f}            & Y
}
\]

\subsection{Étale sheaves on schemes}\etiquette[4.3]{Vbis.subsec4.3}
\pageoriginale

\setcounter{subsubsection}{-1}
\subsubsection{}\etiquette[4.3.0]{Vbis.subsubsec4.3.0}
Let $(\Sch)$ be the category of all schemes belonging to $\Uu$. We thus
define a category $\mathcal{E}$ bifibered in duals of topoi over $(\Sch)$ by
taking as objects the pairs $(X,\mathcal{F})$, where $X$ is a scheme and
$\mathcal{F}$ an étale sheaf on $X$, and as morphisms the pairs
$(f,\varphi):(X,\mathcal{F})\to(Y,\mathcal{G})$, where $f:X\to Y$ is a
morphism of schemes and $\varphi:\mathcal{G}\to f_{\ast}(\mathcal{F})$ is a
morphism of sheaves.

Let $n$ be an integer $\geq0$. We denote by $\mathcal{O}_n$ the section of
$\mathcal{E}^0$ over $(\Sch)^0$ that associates with every scheme $X$ the
constant sheaf $(\ZZ/n\ZZ)_X$. Put
$F_n=\Ysheaf{\Mod}(\mathcal{E}^0,\mathcal{O}_n)$, so that, for every scheme
$X$, $F_{n_X}$ identifies with the category of sheaves of $\ZZ/n\ZZ$-modules
on $X$.

\subsubsection{}\etiquette[4.3.1]{Vbis.subsubsec4.3.1}
Now let $(\Sch)_s$ be the subcategory of $(\Sch)$ whose objects are the
schemes that are elements of $\Uu$ and whose arrows are the quasi-compact and
quasi-separated morphisms. Denote by $\mathcal{E}_s$ the category bifibered
in duals of topoi over $(\Sch)_s$ deduced from $\mathcal{E}$ by the base
change $(\Sch)_s\to(\Sch)$; we again denote by $\mathcal{O}_0$ the
restriction of $\mathcal{O}_0$ to $(\Sch)_s$. Let
$G\subset\Ysheaf{\Mod}(\mathcal{E}^0_s,\mathcal{O}_0)$ be the fibered and
cofibered subcategory of
$\Ysheaf{\Mod}(\mathcal{E}^0_s,\mathcal{O}_0)$ such that, for every scheme
$X$, $G_X$ is the category of torsion abelian sheaves on $X$. The reader
will note that condition (\hyperref[Vbis.par2.4.1.1]{2.4.1.1}) is satisfied.

\begin{enonce*}{Proposition 4.3.2}\etiquette[4.3.2]{Vbis.prop4.3.2}
Let $f:X\to Y$ be a \emph{proper surjective} morphism. Then $f$ is a
morphism of universal $F_n$-$2$-cohomological descent for $n\geq1$.
Considered as a morphism of $(\Sch)_s$, $f$ is a morphism of universal
$G$-$2$-cohomological descent.
\end{enonce*}

By\pageoriginale (\hyperrefext[VIII][VIII.thm9.4]{9.4}), $f$ is a morphism
of effective $F^0$-descent, and by (XII.\hyperref[XII.thm5.1]{(5.1)}) we may apply
(\hyperref[Vbis.prop3.2.4]{3.2.4}) to the Cartesian diagram
\[
\xymatrix@C=3cm@R=1.5cm{
  \fprod{X}{X}{Y} \ar[r] \ar[d]& X \ar[d]^{f}\\
X   \ar[r]_{f}            & Y.
}
\]

\begin{enonce*}{Proposition 4.3.3}\etiquette[4.3.3]{Vbis.prop4.3.3}
Let $f:X\to Y$ be a \emph{surjective} morphism of schemes. Suppose that one
of the following hypotheses is satisfied:
\begin{enumerate}
\renewcommand{\theenumi}{\alph {enumi}}
\renewcommand{\labelenumi}{\theenumi)}
\item $f$ is integral.

\item $Y$ is locally Noetherian, and $f$ is universally open and locally of
  finite presentation.

\item $f$ is flat and locally of finite presentation.
\end{enumerate}
Then $f$ is a morphism of $F_n$-$2$-cohomological descent for every
$n\geq0$.
\end{enonce*}

\begin{enumerate}
\renewcommand{\theenumi}{\alph {enumi}}
\renewcommand{\labelenumi}{\theenumi)}
\item Repeat the proof of (\hyperref[Vbis.prop4.3.2]{4.3.2}), replacing
  (XII.\hyperref[XII.thm5.1]{(5.1)}) by ((\hyperrefext[VIII][VIII.cor5.6]{5.6})).

\item Since the question is local on $Y$, we may suppose that $Y$ is
  Noetherian. Apply (\EGA IV (14.5.10)). By
  (\hyperref[Vbis.prop3.3.1]{3.3.1}) and a), we reduce to the case in which
  every point $y$ of $Y$ has a neighborhood $U$ such that
  $f^{-1}(U)\to U$ has a section, and the result then follows from
  (\hyperref[Vbis.prop3.2.4]{3.2.4}).

\item We must show that $f$ is a morphism of
  \emph{$F$-$1$-cohomological descent} (we know from
  ((\hyperrefext[VIII][VIII.thm9.4]{9.4})) that $f$ is a morphism of
  effective $F^0$-descent).
\end{enumerate}

By localization on $Y$ (\cf (\hyperref[Vbis.prop3.2.4]{3.2.4})), we may
suppose that $Y$ is affine, with ring $A$. We may then write
$\Spec A=\varprojlim_{i\in I}\Spec A_i$, where $I$ is a filtered ordered set
and $A_i$ is a Noetherian ring for every $i$, with
$X=\varprojlim X_i$, where, for every $i$, $X_i$ is a scheme \emph{flat} and
\emph{locally of finite presentation} over $\Spec A_i$ such that
$f_i:X_i\to\Spec A_i$ is \emph{surjective}; moreover, the diagrams
\[
\xymatrix@C=3cm@R=1.5cm{
X \ar[r]\ar[d]_{f}& X_i \ar[d]^{f_i}\\
Y \ar[r]^{u_i}& \Spec A_i
}
\]
are\pageoriginale Cartesian.

Let $\mathfrak{F}$ be an abelian sheaf on $Y$. We have
$\mathfrak{F}={\varinjlim}_i u^{\ast}_i u_{i\ast}(\mathfrak{F})$, so that,
by a standard passage to the limit, we reduce to the case in which
$\mathfrak{F}$ comes from a sheaf on $\Spec A_i$ for some $i$, and b) gives
the conclusion.

\begin{enonce*}[remark]{Remark 4.3.4}\etiquette[4.3.4]{Vbis.rem4.3.4}
If $f:X\to Y$ is integral, then $R^qf_{\ast}=0$ for $q>0$. It follows
readily that it is unnecessary to restrict to complexes bounded below when
performing descent by means of $f$.
\end{enonce*}

\begin{enonce*}{Corollary 4.3.5}\etiquette[4.3.5]{Vbis.cor4.3.5}
Let $(X_{\alpha}\xrightarrow{u_\alpha}X)_{\alpha\in A}$ be a covering family
for the étale topology. Then
$\coprod_{\alpha\in A}X_{\alpha}\to X$ is a morphism of universal
$F_n$-$2$-cohomological descent for every $n\geq0$.
\end{enonce*}

Indeed, there exists a factorization
\[
\xymatrix@=1.5cm{
  &\displaystyle{\mathop{\coprod}_{a\in A} X_{\alpha}} \ar[rd]&\\
X \ar[ru]         \ar[rr]^{h} &                              & X
}
\]
where $h$ is a surjective étale morphism.


% END INCLUDED SGA4_English_current_buildable_source_20260906/source/05bis/drafts/SGA4_VBIS_LINES_2993_3150_English.texfrag


% BEGIN INCLUDED SGA4_English_current_buildable_source_20260906/source/05bis/drafts/SGA4_VBIS_LINES_3151_3452_English.texfrag
\section{Applications}\etiquette[5]{Vbis.sec5}\pageoriginale

\subsection{Construction of simplicial objects}\etiquette[5.1]{Vbis.subsec5.1}

\setcounter{subsubsection}{-1}
\subsubsection{}\etiquette[5.1.0]{Vbis.subsubsec5.1.0}

Up to and including \hyperref[Vbis.prop5.1.3]{5.1.3}, $H$ is a
category in which nonempty finite inverse limits exist and in which
finite coproducts exist and are disjoint and universal. If
$X:\Delta^0\to H$ is a simplicial object of $H$, we denote by
$d^k_i:X_k\to X_{k-1}$ (\resp by $s^k_i:X_k\to X_{k+1}$), for
$i\in[k]$ and $0\leq k$, the arrows corresponding to the face
operators $\partial^i_k:[k-1]\to[k]$ (\resp to the degeneracy
operators $\sigma^i_k:[k]\to[k+1]$) (\cf \cite{sga4.exp.Vbis.bib.3} II 2).

\begin{enonce*}{Definition 5.1.1}\etiquette[5.1.1]{Vbis.def5.1.1}
A simplicial object $X:\Delta^0\to H$ of $H$ is
$\sigma$-split if there exists a family of subobjects $NX_j$ of
$X_j$ $(j\geq0)$ such that the morphisms
$$
h_i:\coprod_{\substack{\Hom_{\Delta^{-}}([i], [j])\\
j\leq i}} NX_j \longrightarrow X_i
$$
are isomorphisms.
\end{enonce*}

Likewise, a $k$-truncated simplicial object
$X:\Delta^0_k\to H$ is $\sigma$-split if there exist $NX_j$
$(0\leq j\leq k)$ satisfying the preceding condition for
$i\leq k$.

The $NX_j$ are uniquely determined by $X$. Indeed, if $X$ is
$\sigma$-split, the degeneracy operators
$s^{k-1}_\ell:X_{k-1}\to X_k$ are isomorphisms from $X_{k-1}$
onto direct summands of $X_k$, and
$$
NX_k = \mathop{\cap}_\ell\quad (\text{complement of } s^{k-1}_\ell(X_{k-1})) {\quoi}.
$$
For $k=0$, $NX_0=X_0$.

\setcounter{subsubsection}{1}
\subsubsection{}\etiquette[5.1.2]{Vbis.subsubsec5.1.2}
For a $\sigma$-split $(n+1)$-truncated simplicial object $X$ of
$H$, we denote by $\alpha_n(X)$ the triple consisting of
\begin{enumerate}
\renewcommand{\theenumi}{\alph {enumi}}
\renewcommand{\labelenumi}{\theenumi)}
\item the restriction $i^{\ast}_n(X)$ of $X$ to $\Delta^0_n$;

\item $NX_{n+1}$;\pageoriginale

\item the evident map from $NX_{n+1}$ to $(\cosq_n(X))_{n+1}$.
\end{enumerate}

This triple $(Y,N,\beta)$ satisfies the following condition:

\noindent
$(\ast)$ $Y$ is a $\sigma$-split $n$-truncated simplicial object of
$H$, and $\beta$ is a map from $N$ to $(\cosq Y)_{n+1}$.

\begin{enonce*}{Proposition 5.1.3}\etiquette[5.1.3]{Vbis.prop5.1.3}
Let $(Y,N,\beta)$ be a triple satisfying $(\ast)$ above.
\begin{enumerate}
\renewcommand{\theenumi}{\roman {enumi}}
\renewcommand{\labelenumi}{(\theenumi)}
\item Up to unique isomorphism, there exists a unique $\sigma$-split
  $X:\Delta^0_{n+1}\to H$ such that
  $\alpha_n(X)\simeq(Y,N,\beta)$.

\item To give a morphism $f$ from $X$ to a truncated simplicial
  object $Z:\Delta^0_{n+1}\to H$ is equivalent to giving
\begin{enumerate}
\renewcommand{\labelenumii}{\theenumii)}
\item a morphism $f':Y\to i^{\ast}_n(Z)$;

\item a morphism $f'':N\to Z_{n+1}$ such that the diagram
\[
\xymatrix@C=4cm@R=2cm{
N  \ar[r] \ar[d]^{f''}& (\cosq_n Y)_{n+1} \ar[d]^{f'}\\
Z_{n+1} \ar[r]& (\cosq_n Z)_{n+1}
}
\]
commutes.
\end{enumerate}
\end{enumerate}
\end{enonce*}

Let us construct $X$.

Set $Y'_{n+1}=(\cosk_nY)_{n+1}$ and denote by
$$
{d'}^{n+1}_i:Y'_{n+1}\to Y_n
\quad(\textup{resp. }{s'}^n_i:Y_n\to Y'_{n+1})
$$
the face operators (\resp the degeneracy operators) obtained by
functoriality.

\noindent \sisi{$+$}{} Set
$Y_{n+1}=N\coprod(\coprod_{\Hom_{\Delta^{-}_{\ell\leq n}}([n+1],
  [\ell])}NY_{\ell})$, and let $\alpha:Y_{n+1}\to Y'_{n+1}$ be
the arrow whose components are $\beta$ and the arrows
$NY_{\ell}\to Y_{\ell}\to Y'_{n+1}$ for every epimorphism
$[n+1]\to[\ell]$ with $\ell\leq n$.

\noindent \sisi{$+$}{} Let\pageoriginale
$s^n_i:Y_n\simeq
\coprod_{\Hom_{\Delta^{-}_{\ell'\leq n}}([n],[{\ell}'])}NY_{{\ell}'}
\to Y_{n+1}$ be the arrow induced by $\sigma^i_n$, so that
\begin{enumerate}
\renewcommand{\theenumi}{\Roman{enumi}}
\renewcommand{\labelenumi}{(\theenumi)}
\item $\alpha\circ s^n_i={{s}'}^n_i$.

Set

\item $d^{n+1}_i={d'}^{n+1}_i\circ\alpha$.
\end{enumerate}

\noindent \sisi{$+$}{} To show that this defines an object of
$\Hom(\Delta^0_{n+1},H)$, it is enough, by [\cite{sga4.exp.Vbis.bib.3}. II (2.4)], to
verify the formulas
\begin{enumerate}
\renewcommand{\theenumi}{\alph{enumi}}
\renewcommand{\labelenumi}{\theenumi)}
\item $d^n_i\cdot d^{n+1}_j=d^n_{j-1}\cdot d^{n+1}_i$ \qquad $i<j$;

\item $s^n_i\cdot s^{n-1}_j=s^n_{j+1}\cdot s^{n-1}_i$ \qquad $i\leq j$;

\item $d^{n+1}_i\cdot s^n_j=
\begin{cases}
s^{n-1}_{j-1}\cdot d^n_i & i<j,\\
\id & i=j \mbox{ or } i=j+1,\\
s^{n-1}_{j}\cdot d^n_{i-1} & i>j+1.
\end{cases}$
\end{enumerate}

Formulas a) and c) are evident by (I) and (II), since they are known
to hold if $d^{n+1}_i$ is replaced by ${d'}^{n+1}_i$ and $s^n_i$
by ${s'}^n_i$.

Formula b) holds by construction.

We leave it to the reader to verify that the object $X$ thus obtained
satisfies \hyperref[Vbis.prop5.1.3]{5.1.3}.

\noindent \sisi{$+$}{} Note that the functor constructed in the proof
of \hyperref[Vbis.prop5.1.3]{5.1.3}, from the category of triples
satisfying (\hyperref[Vbis.subsubsec5.1.2]{5.1.2}) $(\ast)$ to the
category of $(n+1)$-truncated simplicial objects, is faithful but not
full.

Proposition \hyperref[Vbis.prop5.1.3]{5.1.3} provides a very convenient
method for constructing simplicial objects of $H$ by induction. In
the remainder of this section, we formalize this observation in the
form that will be used in \hyperref[Vbis.subsec5.2]{5.2} and
\hyperref[Vbis.subsec5.3]{5.3}.

\setcounter{subsubsection}{3}
\subsubsection{}\etiquette[5.1.4]{Vbis.subsubsec5.1.4}

Let\pageoriginale $E\xrightarrow{\pi}B$ be a functor such that, for
every object $b$ of $B$, the fiber category $E_b$ is nonempty and
satisfies the conditions of \hyperref[Vbis.subsubsec5.1.0]{5.1.0}.

Let $Q$ be a property of an object of $E$, stable under isomorphism.
Suppose that the following conditions hold:
\begin{enumerate}
\renewcommand{\theenumi}{\alph{enumi}}
\renewcommand{\labelenumi}{\theenumi)}
\item \emph{For every object $b$ of $B$, there exists an object $x$ of
  $E_b$ satisfying $Q$.}

\item \emph{For every object $b$ of $B$ and every pair $(x,y)$ of
  objects of $E_b$ satisfying $Q$, there exists an object $z$ of $E_b$
  satisfying $Q$ and morphisms $z\to x$, $z\to y$ in $E_b$.}

\item \emph{For every morphism $b'\xrightarrow{t}b$ of $B$ and every
  object $x$ of $E_b$ satisfying $Q$, there exists an object $x'$ of
  $E_{b'}$ satisfying $Q$ and a morphism $x'\to x$ over $t$.}

Let, on the other hand, $P$ be a property of an arrow of $E$ over an
identity, stable under isomorphism. Suppose that the following
conditions hold:

\item \emph{For every object $b$ of $B$ and every object $x$ of $E_b$,
  there exists an arrow $y\to x$ in $E_b$ satisfying $P$.}

\item \emph{For every object $b$ of $B$, every diagram in $E_b$ of
  the form}
\[
\xymatrix{
u \ar[rrrr]^{f} & & &                 &x\\
                & & &z\ar[ru]\ar[rd] &\\
v \ar[rrrr]_{g} & & &               & y
}
\]
\emph{in which $f$ and $g$ satisfy $P$, there is a completion to a
commutative diagram in $E_b$}
\[
\xymatrix@R=1.25cm{
  & u\ar[rrrr]^{f} & & &  & x\\
w\ar[rrrr]^{h}\ar[ru]\ar[rd] &   & & &z\ar[ru]\ar[rd] & \\
  & v\ar[rrrr]_{g} & & &  & y
}
\]
\emph{in which}\pageoriginale $h$ \emph{satisfies $P$.}

\item \emph{For every morphism $b'\xrightarrow{t}b$ of $B$ and every
  morphism $x\xrightarrow{f}y$ of $E_b$ satisfying $P$, there exists
  a morphism $x'\xrightarrow{f'}y'$ of $E_{b'}$ satisfying $P$ and a
  commutative diagram}
\[
\xymatrix@C=3cm@R=1.5cm{
x' \ar[r]^{\alpha} \ar[d]_{f'}& x\ar[d]^{f}\\
y' \ar[r]_{\beta}& y
}
\]
\emph{in which $\alpha$ and $\beta$ are morphisms over $t$.}
\end{enumerate}

\begin{enonce*}{Definition 5.1.5}\etiquette[5.1.5]{Vbis.def5.1.5}
With the notation of \hyperref[Vbis.subsubsec5.1.4]{5.1.4}, we say
that an object $X$ of $\Hom(\Delta^0,E)$ satisfies condition
\emph{(APQ)} if the following conditions hold:
\begin{enumerate}
\item[(i)] The image of $\pi\circ X$ is the constant simplicial object
  defined by an object $b$ of $B$.

\item[(ii)] As an object of $\Hom(\Delta^0,E_b)$, $X$ is
  $\sigma$-split.

\item[(iii)] $X_0$ satisfies condition $Q$.

\item[(iv)] For every integer $n$, the arrow
  $NX_{n+1}\to(\cosq_nX)_{n+1}$ satisfies condition $P$.
\end{enumerate}
\end{enonce*}

\setcounter{subsubsection}{5}
\subsubsection{}\etiquette[5.1.6]{Vbis.subsubsec5.1.6}
Let $E_{(APQ)}$ denote the subcategory of $\Hom(\Delta^0,E)$ whose
objects are the simplicial objects satisfying $(APQ)$ and whose
morphisms are the morphisms $X\xrightarrow{\alpha}Y$ of
$\Hom(\Delta^0,E)$ such that
$\pi(\alpha_n):\pi(X_n)\to\pi(Y_n)$ is the same morphism of $B$ for
every $n$. There is then a canonical functor
$\bar{\pi}:E_{(APQ)}\to B$.

\begin{enonce*}{Proposition 5.1.7}\etiquette[5.1.7]{Vbis.prop5.1.7}
The\pageoriginale functor $\bar{\pi}$ is surjective. Moreover, for
every object $b$ of $B$ and every pair $(X,Y)$ of objects of
$E_{(APQ),b}$, there exists a diagram
\[
\xymatrix@=1.5cm{
  & Z \ar[ld]\ar[rd]&\\
X &   & Y
}
\]
in $E_{(APQ),b}$.
\end{enonce*}

The proof follows immediately from \hyperref[Vbis.prop5.1.3]{5.1.3},
by virtue of the hypotheses made in
\hyperref[Vbis.subsubsec5.1.4]{5.1.4}.

\begin{enonce*}{Proposition 5.1.8}\etiquette[5.1.8]{Vbis.prop5.1.8}
Let $C$ be a category and $F$ a functor $E_{(APQ)}\to C$. We introduce
the following conditions on $F$:
\begin{enumerate}
\renewcommand{\theenumi}{\roman{enumi}}
\renewcommand{\labelenumi}{(\theenumi)}
\item For every object $b$ of $B$, the restriction of $F$ to
  $E_{(APQ),b}$ factors through the \emph{coarse} category
  \footnote{A category is called \emph{coarse} if, for every pair
  $(x,y)$ of objects, $\Hom(x,y)$ consists of exactly one element.}
  defined by the objects of $E_{(APQ),b}$.

\item For every morphism $t:b'\to b$ of $B$ and every pair
  $(X'\to X,Y'\to Y)$ of morphisms over $t$, the following diagram
  commutes:
\[
\xymatrix@C=3cm@R=2cm{
F(X')         \ar[r]  \ar[d]_{\alpha'}& F(X)\ar[d]^{\alpha}\\
F(Y')          \ar[r] & F(Y)
}
\]
where $\alpha$ and $\alpha'$ are the arrows determined by condition
\emph{(i)}.
\end{enumerate}

Then the functor $G\mapsto G\circ\bar{\pi}$ induces an equivalence
between the category $\Hom(B,C)$ and the full subcategory of
$\Hom(E_{(APQ)},C)$ defined by the functors satisfying conditions
\emph{(i) and (ii)}.
\end{enonce*}

The\pageoriginale details of the proof are left to the reader. We
indicate only how to prove that the preceding functor is essentially
surjective.

Let $F:E_{(APQ)}\to C$ be a functor satisfying conditions (i) and
(ii). For every object $b$ of $B$, choose an object $X$ of
$E_{(APQ),b}$ and set $\bar{F}(b)=F(X)$; one then verifies that
$\bar{F}(b)$ varies functorially with $b$.


% END INCLUDED SGA4_English_current_buildable_source_20260906/source/05bis/drafts/SGA4_VBIS_LINES_3151_3452_English.texfrag


% BEGIN INCLUDED SGA4_English_current_buildable_source_20260906/source/05bis/drafts/SGA4_VBIS_LINES_3453_3612_English.texfrag
\subsection{Singular cohomology of a scheme of finite type over a
  field of characteristic zero}\etiquette[5.2]{Vbis.subsec5.2}

In this section, $k$ will denote a field of characteristic zero.

\setcounter{subsubsection}{-1}
\subsubsection{}\etiquette[5.2.0]{Vbis.subsubsec5.2.0}
We shall denote by $\Sch_f /k$ (\resp Ann) the category of schemes
\emph{locally of finite type} over $k$ (\resp complex analytic
spaces). All simplicial objects of $\Sch_f /k$ (\resp Ann) will be
regarded as simplicial objects of Top by means of the canonical
functor $\Sch_f /k \to  \Top$ (\resp $\Ann \to \Top$). We retain the
data of (\hyperref[Vbis.subsubsec4.1.0]{4.1.0}) concerning cohomological descent.

Following the usual notation, we write $S \to S^{\an}$ for the
canonical functor $\Sch_f / \CC \to$ Ann defined in G.A.G.A.

\subsubsection{}\etiquette[5.2.1]{Vbis.subsubsec5.2.1}
Let $X$ be a simplicial object of $\Sch / k$ (\resp $\Ann$): the
usual compatibilities in the construction of the de Rham complex
of a morphism of schemes (\resp of analytic spaces) make it possible
to construct a complex $\Omega_{X/k}^{\Zar.}$ (\resp $\Omega_{X /
  \CC}^{\bigdot}$) of $D^+(\Ysheaf{\Gamma} (\bar{X},
\ZZ))$ such that, for every integer $i$, $\sisi{\RR^+}{R} e^{\ast}_i
(\Omega_{X/k}^{\Zar.})$ (\resp $\sisi{\RR^+}{R} e^{\ast}_i$
$(\Omega^{\bigdot}_{X /\CC})$) is the de Rham complex
$\Omega_{X/k}^{\Zar.}$ (\resp $\Omega^{\bigdot}_{X_i/ \CC}$).

Let $X \xrightarrow{u} C^{\Delta}_S$ be an augmented simplicial
object of $\Sch_f / k$ (\resp Ann). We define $H^{\Zar.}_{DR}(u)$
(\resp $H^{\bigdot}_{DR}(u)$ and $H^{\bigdot}(u, \CC)$) to be the
hypercohomology groups of the complex $\sisi{\RR^+}{R} \bar{u}_+
(\Omega_{X/k}^{\Zar.})$ (\resp $\sisi{\RR^+}{R} \bar{u}_{\ast}
(\Omega_{X/\CC}^{\bigdot})$ and $\sisi{\RR^+}{R} \bar{u} (\CC)$).

\subsubsection{}\etiquette[5.2.2]{Vbis.subsubsec5.2.2}
If\pageoriginale $X \xrightarrow{u} C^{\Delta}_S$ is a semisimplicial
object of $\Sch_f / \CC$, there are three canonical morphisms:
\begin{align*}
H^{\Zar.}_{DR} (u) &\xrightarrow{\alpha} H^{\bigdot}_{DR}(u^{\an})\\
H^{\bigdot} (u^{\an}, \CC) &\xrightarrow{\beta}
H^{\bigdot}_{DR}(u^{\an})\\
H^{\bigdot} (S^{\an}, \CC) &\xrightarrow{\gamma}
H^{\bigdot}(u^{\an}, \CC){\quoi}.
\end{align*}

\begin{enonce*}{Proposition 5.2.3}\etiquette[5.2.3]{Vbis.prop5.2.3}
With the notation of (\hyperref[Vbis.subsubsec5.2.2]{5.2.2}), suppose
that $X_n$ is smooth for every $n$: then $\alpha$ and $\beta$ are
isomorphisms. If, moreover, $X^{\an} \xrightarrow{u^{\an}}
C^{\Delta}_{S^{\an}}$ is a hypercovering of $S^{\an}$ for the
topology of universal $1$-cohomological descent, then $\gamma$ is
an isomorphism.
\end{enonce*}

The assertion that $\alpha$ is an isomorphism follows from
[\cite{sga4.exp.Vbis.bib.6}, Th. $1'$]. The Poincar\'e lemma shows that $\beta$ is an
isomorphism. The final part of the proposition follows from
(\hyperref[Vbis.thm3.3.3]{3.3.3}).

\setcounter{subsubsection}{3}
\subsubsection{}\etiquette[5.2.4]{Vbis.subsubsec5.2.4}
We shall now apply (\hyperref[Vbis.subsubsec5.1.4]{5.1.4}) in the case
where $B = \Sch_f /k$ and $E$ is the category over $B$ such that, for
every scheme $S$, the fiber category $E_S$ is the category of schemes
of $\Sch_f /k$ over $S$, the arrows over an arrow of $B$ being the
obvious ones. We shall say that a morphism $h : X \to Y$ over $S$
satisfies property $P$ if $X$ is smooth over $k$ and if $h$ is a
composite of finitely many morphisms satisfying one of the following
conditions:
\begin{enumerate}
\renewcommand{\theenumi}{\roman{enumi}}
\renewcommand{\labelenumi}{(\theenumi)}
\item it is the coproduct of a family of \'etale surjective morphisms;
\item it is the coproduct of a family of proper surjective morphisms.
\end{enumerate}
We shall say that a scheme $X$ of $E_S$ satisfies property $Q$ if the
structural morphism $X \to S$ satisfies property $P$.

We leave it to the reader to verify, using resolution of singularities,
that all the conditions of (\hyperref[Vbis.subsubsec5.1.4]{5.1.4}) are satisfied.

The\pageoriginale groups $H^{\Zar.}_{DR}(u)$ define a contravariant
functor from the category $E_{(APQ)}$ to the category of graded
abelian groups, denoted by $H^{\Zar.}_{DR}$.

\begin{enonce*}{Proposition 5.2.5}\etiquette[5.2.5]{Vbis.prop5.2.5}
The functor $H^{\Zar.}_{DR}$ satisfies the conditions of
(\hyperref[Vbis.prop5.1.8]{5.1.8}).
\end{enonce*}

By the Lefschetz principle, we are reduced to the case $k = \CC$.
In that case, for every morphism $S'\xrightarrow{t} S$ and every pair
$(u'\to u$, $v'\to v)$ of morphisms over $t$, there is a commutative
diagram:
\[
\xymatrix@C=3.5cm@R=2cm{
H^{\Zar.}_{DR}(u') \ar[r] \ar[d]_{\gamma^{-1}.\beta^{-1}.\alpha}&  H_{DR}^{\Zar.}(u)\ar[d]^{\gamma^{-1}.\beta^{-1}.\alpha}\\
H^{\bigdot}(S'^{\an},\CC)
\ar[r]^{H^{\bigdot}(t^{\an},\CC)}
& H^{\bigdot}(S^{\an},
\CC) \\
H^{\Zar.}_{DR}(v') \ar[r] \ar[u]^{\gamma^{-1}.\beta^{-1}.\alpha}& H_{DR}^{\Zar.}(v) \ar[u]_{\gamma^{-1}.\beta^{-1}.\alpha}
}
\]
which completes the proof.

\begin{enonce*}{Definition 5.2.6}\etiquette[5.2.6]{Vbis.def5.2.6}
The contravariant functor defined by
(\hyperref[Vbis.prop5.1.8]{5.1.8}) from $H^{\Zar.}_{DR}$ will be
denoted by $H^{\bigdot} (\ast, k)$. For every scheme $S$ locally of
finite type over $k$, the graded abelian group $H^{\bigdot}(S, k)$
will be called the \emph{singular cohomology of} $S$.
\end{enonce*}

\setcounter{subsubsection}{6}
\subsubsection{}\etiquette[5.2.7]{Vbis.subsubsec5.2.7}
We conclude this section with an illustration of the general
principles introduced in it. Recall first that if $S$ is a complex
analytic space, the \emph{de Rham cohomology} of $S$, denoted by
$\HH^{\bigdot}_{DR}(S)$, is by definition the hypercohomology of the
de Rham complex $\Omega^{\bigdot}_{S / \CC}$.

\begin{enonce*}{Proposition 5.2.8}\etiquette[5.2.8]{Vbis.prop5.2.8}
\emph{(Blum--Herrera)}.\pageoriginale Let $S$ be a scheme locally of
finite type over $\CC$. The canonical arrow
$$
H^{\bigdot} (S^{\an}, \CC) \longrightarrow
\HH^{\bigdot}_{DR}(S^{\an})
$$
has a retraction (in the category of graded abelian groups) that is
functorial in $S$.
\end{enonce*}

By (\hyperref[Vbis.prop5.1.8]{5.1.8}), it is enough to consider the
morphism $H^{\bigdot}(+, \CC) \circ \bar{\pi} \to
\HH^{\bigdot}_{DR} \circ \bar{\pi}$ obtained from the preceding one
by composition with $\bar{\pi}$.

Using, for every augmentation $X \xrightarrow{u} C^{\Delta}_S$, the
canonical morphism $\LL \bar{u}^{+}(\Omega^{\bigdot}_{S/C}) \to
\Omega^{\bigdot}_{X/C}$, one defines a functorial morphism
$\HH^{\bigdot}_{DR} \circ \bar{\pi} \to \HH^{\bigdot}_{DR}$ so that
the following diagram is commutative:
\[
\xymatrix@=2cm{
H^{\bigdot}_{DR}  & H^{\bigdot}_{DR} \circ \bar{\pi} \ar[l]\\
H^{\bigdot}(\ast,\CC)_{\circ}\bar{\pi}
\ar[u]_{\text{\rotatebox{90}{$\backsim$}}}^{\beta \circ \gamma}
\ar[ru]&
}
\]
which gives the desired retraction.

\begin{enonce*}[remark]{Remark 5.2.9}\etiquette[5.2.9]{Vbis.rem5.2.9}
Statement (\hyperref[Vbis.prop5.2.8]{5.2.8}) remains valid for every
analytic space once resolution of singularities is available in the
analytic setting.
\end{enonce*}


% END INCLUDED SGA4_English_current_buildable_source_20260906/source/05bis/drafts/SGA4_VBIS_LINES_3453_3612_English.texfrag


% BEGIN INCLUDED SGA4_English_current_buildable_source_20260906/source/05bis/drafts/SGA4_VBIS_LINES_3613_3730_English.texfrag
\subsection{Mixed Hodge theories}\etiquette[5.3]{Vbis.subsec5.3}

We give here a technical result that plays a key role in \cite{sga4.exp.Vbis.bib.2}.

$k$ is a field of characteristic zero.

\subsubsection{}\etiquette[5.3.1]{Vbis.subsubsec5.3.1}
Let $\Sch_{\fs} /k$ be the category of separated schemes of finite
type over $k$; denote by $E \xrightarrow{\pi} \Sch_{\fs}/k$ the
category over $\Sch_{\fs}/k$ defined as follows:

\noindent \sisi{$+$}{} If\pageoriginale $S$ is an object of
$\Sch_{\fs} /k$, an object of $E_S$ is a triple $(\bar{X}, X, i)$,
where $\bar{X}$ is a reduced projective scheme over $k$, $X$ is a
proper scheme over $S$, and $i : X \to \bar{X}$ is an open immersion
over $k$ such that $i(X)$ is dense in $\bar{X}$.

\noindent \sisi{$+$}{} If $t : S' \to S$ is a morphism of schemes,
$(\bar{X}, X, i)$ is an object over $S$, and $(\bar{X}', X', i')$ is
an object over $S'$, a morphism from $(\bar{X}, X, i)$ to
$(\bar{X}', X', i')$ over $t$ is a pair $(\bar{h}, h)$, where
$\bar{h} : \bar{X}'\to \bar{X}$ is a morphism over $k$ and
$h : X'\to X$ is a morphism over $S$, such that the following diagram
is commutative.
\[
\xymatrix@C=3.5cm@R=2cm{
X' \ar[r]^{i'} \ar[d]_{h}& \bar{X}'\ar[d]^{\bar{h}}\\
X  \ar[r]^{i}& \bar{X}.
}
\]
N.B.: the properness hypotheses imply that
$\bar{h}^{-1}(i(X)) = i(X')$.

\begin{enonce*}{Lemma 5.3.2}\etiquette[5.3.2]{Vbis.lem5.3.2}
For every object $S$ of $\Sch_{\fs} / k$, the category $E_S$ is
nonempty, has nonempty projective limits, and has finite coproducts
that are disjoint and universal.
\end{enonce*}

The fact that the category $E_S$ is nonempty follows from Chow's
lemma.

We shall describe only the construction of the direct product of two
objects $(\bar{X}, X, i)$ and $(\bar{X}', X', i')$ of $E_S$.

The canonical arrow $j : \fprod{X}{X'}{S} \to
\fprod{\bar{X}}{\bar{X}'}{k}$ is an immersion. Let $T$ be the reduced
closed subscheme of $\fprod{\bar{X}}{\bar{X}'}{k}$ whose underlying
space is the closure of the image of $j$. Then
$(T, {\fprod{X}{X'_{\red}}{S}}, j_{\red})$ has the required universal
property.

\setcounter{subsubsection}{2}
\subsubsection{}\etiquette[5.3.3]{Vbis.subsubsec5.3.3}
Let\pageoriginale $S$ be an object of $\Sch_{\fs} /k$. We shall say
that an object $(\bar{X}, X, i)$ of $E_S$ satisfies property $Q$ if
$\bar{X}$ is smooth over $k$ and $i(X)$ is the complement of a normal
crossings divisor. We shall say that a morphism
$(\bar{h}, h) : (\bar{X}', X', i') \to (\bar{X}, X, i)$ satisfies
property $P$ if $\bar{h}$ is surjective and $(\bar{X}', X', i')$
satisfies property $Q$.

\begin{enonce*}{Proposition 5.3.4}\etiquette[5.3.4]{Vbis.prop5.3.4}
With the notation of (\hyperref[Vbis.subsubsec5.3.3]{5.3.3}), all the
conditions of (\hyperref[Vbis.subsubsec5.1.4]{5.1.4}) are satisfied.
\end{enonce*}

Conditions a), b), d), and e) are readily verified using resolution
of singularities. Moreover, c) follows from f), so it remains to
verify f).

We must show that every commutative diagram
\[
\xymatrix@C=3cm@R=1cm{
              &Z         \ar[r]^{i_Z} \ar[d]_{h}& \bar{Z}\ar[d]^{\bar{h}}\\
              &X \ar[r]^{i_X} \ar[d]& \bar{X}\ar[d]\\
 S'\ar[r]^{t} &S  \ar[r]            & k
}
\]
where $(\bar{X}, X, i_X)$ and $(\bar{Z}, Z, i_Z)$ are objects of
$E_S$ and $\bar{h}$ is surjective can be completed to a commutative
diagram
\[
\xymatrix@C=1.7cm@R=.45cm{
Z' \ar[rr] \ar[rd]_{i_{Z'}} \ar[dd]_{h'}&    & Z \ar[rd]^{i_Z} \ar@{->} '[d][dd]_{h} &\\
         & \bar{Z}' \ar[rr] \ar@{->} '[d][dd]&    & \bar{Z}\ar[dd]^{\bar{h}}\\
X'  \ar[rr]^(.65){h'} \ar[dd]\ar[rd]_{i_{X'}}   &    & X \ar[rd]_{i_X} \ar@{->} '[d][dd]&\\
         & \bar{X}'  \ar@{->} '[d][dd] \ar[rr]&    & \bar{X} \ar[dd]\\
S'\ar[rr]^(.65){t}\ar[rd]        &    & S \ar[rd] &\\
         & k  \ar@{=}[rr]&    & k
}
\]
where\pageoriginale $(\bar{X}', X', i_{X'})$ and
$(\bar{Z}', Z', i_{Z'})$ are objects of $E_{S'}$ and $\bar{h'}$ is
surjective.

Consider a diagram
\[
\xymatrix@C=2cm@R=1.5cm@M=.2cm{
T \ar[r]^{\varphi} \ar@{_{(}->}[d]_{i_T}& S'\ar[r]^{t} & S\ar[d]\\
\bar{T} \ar[rr]_{\psi}&  & k\\
}
\]
where $\varphi$ is proper and $\psi$ is projective. By base change we
obtain a diagram
\[
\xymatrix@C=1cm@R=.6cm{
\fprod{Z}{T_{\red}}{S} \ar[rrr] \ar[rd]_{p} \ar[dd]&&    & Z \ar[rd] \ar@{->} '[d][dd] &\\
         & \fprod{\bar{Z}}{\bar{T}_{\red}}{k} \ar[rrr] \ar@{->} '[d][dd]& &   & \bar{Z}\ar[dd]\\
\fprod{X}{T_{\red}}{S}  \ar[rrr] \ar[dd]\ar[rd]_{q} &  &    & X \ar[rd] \ar@{->} '[d][dd]&\\
         &\fprod{\bar{X}}{\bar{T}_{\red}}{k}  \ar@{->} '[d][dd] \ar[rrr]& &   & \bar{X} \ar[dd]\\
T \ar[rd]\ar[r]& S'\ar[rr]        &    & S \ar[rd] &\\
         & \bar{T}  \ar[rrr]&  &  & k
}
\]
where\pageoriginale $p$ and $q$ are immersions. It then suffices to
replace $\fprod{\bar{Z}}{\bar{T}_{\red}}{k}$ and
$\fprod{\bar{X}}{\bar{T}_{\red}}{k}$ by the closed images of $p$ and
$q$ to obtain the required diagram. This completes the proof of
(\hyperref[Vbis.prop5.3.4]{5.3.4}).


% END INCLUDED SGA4_English_current_buildable_source_20260906/source/05bis/drafts/SGA4_VBIS_LINES_3613_3730_English.texfrag


% BEGIN INCLUDED SGA4_English_current_buildable_source_20260906/source/05bis/drafts/SGA4_VBIS_LINES_3731_3755_English.texfrag
\begin{thebibliography}{99}\pageoriginale
\bibitem{sga4.exp.Vbis.bib.1} T. Blum -- M. Herrera : De Rham Cohomology of an
  analytic Space. Inv. Math. vol 7. Fas 4-1969. p.275-296.
\bibitem{sga4.exp.Vbis.bib.2} P. Deligne : Théorie de Hodge III (I.H.E.S).
\bibitem{sga4.exp.Vbis.bib.3} P. Gabriel -- M. Zisman : Homotopy theory and calculus of
  fractions Ergebnisse der Mathematik -- Band 35 -- Springer 1967.
\bibitem{sga4.exp.Vbis.bib.4} J. Giraud : Méthode de la descente : Mémoires de la
  S.M.F. 2. 1964.
\bibitem{sga4.exp.Vbis.bib.5} R. Godement : Théorie des faisceaux.
                                Publications de l'Institut de
                                Mathématique de l'Université de
                                Strasbourg. Hermann 1964.
\bibitem{sga4.exp.Vbis.bib.6} A. Grothendieck : On the De Rham Cohomology of Algebraic
  Varieties, Publications de 1' I.H.E.S ${\rm n^{\circ}}$ 29.
\bibitem{sga4.exp.Vbis.bib.7} R. Hartshorne : Residues and Duality. Springer
  Lecture Notes ${\rm n^{\circ}}$ 20.
\bibitem{sga4.exp.Vbis.bib.8} G. Segal : Classifying spaces and Spectral sequences
  Publications de 1'I.H.E.S. ${\rm n^{\circ}}$ 34.
\bibitem{sga4.exp.Vbis.bib.9} J.L.Verdier : Catégories dérivées - Etat 0 -
  (I.H.E.S)
\end{thebibliography}

\ifx\danslelivre\undefined
\end{document}
\fi

% END INCLUDED SGA4_English_current_buildable_source_20260906/source/05bis/drafts/SGA4_VBIS_LINES_3731_3755_English.texfrag


% END INCLUDED SGA4_English_current_buildable_source_20260906/source/05bis/05bis_English.tex

\setcounter{chapter}{5}
\def\thechapter{\Roman{chapter}}

\makeatletter
\def\input@path{{../06/}}
\makeatother
% BEGIN INCLUDED SGA4_English_current_buildable_source_20260906/source/06/06_English.tex
\ifx\danslelivre\undefined
% Release engine: XeLaTeX (the LuaTeX compatibility layer breaks Xy-pic output).
\documentclass[10pt]{../fichiers-configuration-latex/smfart}
% BEGIN INCLUDED SGA4_English_current_buildable_source_20260906/source/fichiers-configuration-latex/paquets.tex
\usepackage{polyglossia}
%\usepackage[francais]{babel}
\usepackage{fontspec}
\usepackage[math-style=ISO,bold-style=upright,sans-style=upright]{unicode-math}
%\usepackage[mathletters]{ucs}
%\usepackage[utf8]{inputenc}
%\usepackage[T1]{fontenc}
\usepackage[unicode,hypertexnames=false,colorlinks=true,linkcolor=blue,citecolor=Sepia,urlcolor=magenta]{hyperref}
\setmainlanguage{french}
\setmainfont[Ligatures=TeX]{Linux Libertine O}
% variante possible : TeXGyrePagella ?
\setmathfont{XITS Math}

% maths
\usepackage{amsmath}
%\usepackage{amssymb}
%\usepackage{mathrsfs}
%\usepackage{stmaryrd}
%\usepackage{euscript}

% couleurs
\usepackage[usenames,dvipsnames]{xcolor}

% graphiques
\usepackage{graphicx}
\usepackage[all]{xy}
%\usepackage{tikz}
%\usetikzlibrary{matrix,arrows,calc}
%\tikzset{column sep=3em,row sep=2em,text height=1.5ex, text depth=0.25ex,>=angle 90}
%\tikzset{description/.style={fill=white,inner sep=1pt}}

% hyperliens, 'etiquettes
\usepackage{xr-hyper}
\usepackage{../fichiers-configuration-latex/smfthm}
%\ifx\cacheretiquettes\undefined
%\usepackage[inline]{showlabels} \else \fi % pour rendre visibles les 'etiquettes
%\usepackage{../fichiers-configuration-latex/makeidx} % pour faire un index
%\usepackage[pagebackref,unicode]{hyperref}  % pour voir o`u sont cit'ees les r'ef'erences
% hyperref is loaded once above with the cumulative-reader destination policy.
%\hypersetup{breaklinks=true} %permet le retour `a la ligne dans les liens trop longs. Incompatible avec pdflatex
%\renewcommand*{\backref}[1]{{\small [cit\'e page(s)~#1.]}}

\usepackage{xspace}
\usepackage{pifont} % pour expose XVII

% END INCLUDED SGA4_English_current_buildable_source_20260906/source/fichiers-configuration-latex/paquets.tex

% BEGIN INCLUDED SGA4_English_current_buildable_source_20260906/source/fichiers-configuration-latex/numerotation.tex
\renewcommand{\thefootnote}{\roman{footnote}} % pour mettre les notes en bas de page sous la forme (ii) et non (2), qui ressemble trop (?) a un exposant.

\newcounter{pageoriginal}
\newcommand{\pageoriginale}{\refstepcounter{pageoriginal}\marginpar{\footnotesize\textbf{\thepageoriginal}}\xspace}
\let\pageoriginaled\pageoriginale

\usepackage{ifthen}
\newboolean{original}
\def\sisi{\ifthenelse{\boolean{original}}}
\setboolean{original}{false}

\newcounter{cnote}
\def\nde#1{\ifthenelse{\boolean{original}}{}{\stepcounter{cnote}\def\thefootnote{\arabic{cnote}}\protect\footnote[\value{cnote}]{N.D.E.: #1}}}
\def\ndemark{\ifthenelse{\boolean{original}}{}{\stepcounter{cnote}\def\thefootnote{\arabic{cnote}}\protect\footnotemark[\value{cnote}]}}
\def\ndetext#1{\ifthenelse{\boolean{original}}{}{\def\thefootnote{\arabic{cnote}}\protect\footnotetext[\value{cnote}]{N.D.E.: #1}}}

% END INCLUDED SGA4_English_current_buildable_source_20260906/source/fichiers-configuration-latex/numerotation.tex

% BEGIN INCLUDED SGA4_English_current_buildable_source_20260906/source/fichiers-configuration-latex/macros.tex
\setmainlanguage{french}
\setotherlanguage{english}


%% Fixed, engine-independent release identifier.  The inherited dynamic Git
%% probe used Lua io.popen, exposed build-state-dependent output, and printed
%% a visible question mark outside a Git checkout.
\def\showgitstatus{English release r2, 2026-07-23}


% macros pour LaTeX
\renewcommand{\mathsf}{}

% Police différente pour U+2113 SCRIPT SMALL L
\setmathfont[range="2113]{latinmodern-math.otf}

% Give every inherited semantic label its own stable PDF destination instead
% of inheriting whichever automatic equation/page anchor happened to be live.
% The legacy label remains as a compatibility alias used by existing edges.
\makeatletter
\def\etiquette[#1]#2{%
  \Hy@raisedlink{\hyper@anchorstart{sga4.semantic.#2}\hyper@anchorend}%
  \def\@currentHref{sga4.semantic.#2}%
  \def\@currentlabel{#1}%
  \label{#2}%
}
\makeatother
\def\eqhyperref[#1]#2{(\hyperref[#1]{#2})}
\def\pagehyperref[#1]#2{(page \hyperref[#1]{#2})}
\def\hyperrefext[#1][#2]#3{\hyperref[#2]{#1 #3}}
\def\eqhyperrefext[#1][#2]#3{\hyperref[#2]{#1 #3}}     %cf. exposes 17
\def\neweqhyperrefext[#1][#2]#3{\hyperref[#2]{#1 #3}}  %cf. exposes 17

\def\shat#1{\protect\widehat{#1}}
\def\stilde#1{\protect\widetilde{#1}}
%%%%%%%%%%%%%%%%%%


\def\etcn{etc.\xspace}
\def\etco{etc...\xspace}
\def\etc{\sisi{\etco}{\etcn}}
\def\ldotsn{\sisi{...}{\etcn}}

\def\cf{{cf.}\kern.3em}
\def\Cf{{Cf.}\kern.3em}
\def\ie{{i.e.}\kern.3em}
\def\iev{{i.e.},\xspace}
\def\loccito{loc.\kern2pt cit.\xspace}
\def\loccitn{\textit{loc.\kern2pt cit.}\xspace}
\def\loccit{\sisi{\loccito}{\loccitn}}
\def\EGA{\textup{EGA}\kern.2em}
\def\SGA{\textup{SGA}\kern.3em}

\def\quoi{}

\def\quav{\quad,\quad}
\def\quaet{\quad \text{et}\quad}
\def\fibreproduct#1#2#3{#1 ×_{#3} #2}
\let\fprod\fibreproduct
\def\ccheck#1#2{\mbox{${\check{\check{\rm #1}}}^{\lower2pt\hbox{$\scriptstyle #2$}}$}}
\def\sheaf#1{\protect\mathcal{#1}}
\def\vieuxsheaf#1{{#1}} % cf. problème avec sheaf{a\Phi_i} par exemple...
\DeclareMathOperator{\Hom}{Hom}
\DeclareMathOperator{\can}{can}
\DeclareMathOperator{\Ind}{Ind}
\DeclareMathOperator{\ind}{ind}
%\DeclareMathOperator{\Hhomtop}{\mathcal{H}\!\mathit{omtop}}
\DeclareMathOperator{\Hhomtop}{\mathbf{Homtop}}
\DeclareMathOperator{\Point}{Point}
\DeclareMathOperator{\Pro}{Pro}
\DeclareMathOperator{\Top}{Top}
\DeclareMathOperator{\Mod}{Mod}
\DeclareMathOperator{\Cart}{Cart}
%\DeclareMathOperator{\Mmorsite}{\mathcal{M}\!\mathit{orsite}}
\DeclareMathOperator{\Mmorsite}{\mathbf{Morsite}}
\DeclareMathOperator{\loc}{loc}
\DeclareMathOperator{\Loc}{Loc}
\DeclareMathOperator{\Pic}{Pic}
\DeclareMathOperator{\Et}{Et}
\DeclareMathOperator{\card}{card}
\DeclareMathOperator{\Ens}{Ens}
\DeclareMathOperator{\h}{\underline{H}}
\DeclareMathOperator{\Fonct}{Fonct}
\DeclareMathOperator{\Ker}{Ker}
\DeclareMathOperator{\but}{but}
\DeclareMathOperator{\Ob}{Ob}
\DeclareMathOperator{\ob}{ob}
\DeclareMathOperator{\Lib}{Lib}
\DeclareMathOperator{\esj}{esj}
\DeclareMathOperator{\id}{id}
\DeclareMathOperator{\Id}{Id}
\DeclareMathOperator{\ad}{ad}
\DeclareMathOperator{\Iim}{Im}
\DeclareMathOperator{\pr}{pr}
\DeclareMathOperator{\Cat}{Cat}
\DeclareMathOperator{\car}{car}
\DeclareMathOperator{\SysInd}{SysInd}
\DeclareMathOperator{\Sup}{Sup}
\DeclareMathOperator{\Coker}{Coker}
\DeclareMathOperator{\Filt}{Filt}
\DeclareMathOperator{\acc}{acc}
\DeclareMathOperator{\sur}{sur}
\DeclareMathOperator{\cat}{cat}
\DeclareMathOperator{\Ord}{Ord}
\DeclareMathOperator{\Gr}{Gr}
\DeclareMathOperator{\Cov}{Cov}
\DeclareMathOperator{\Caf}{Caf}
\DeclareMathOperator{\coker}{coker}
\DeclareMathOperator{\coim}{coim}
\DeclareMathOperator{\im}{im}
\DeclareMathOperator{\Ouv}{Ouv}
\DeclareMathOperator{\diag}{diag}
\DeclareMathOperator{\Esp}{Esp}
\DeclareMathOperator{\TOP}{TOP}
\DeclareMathOperator{\Unif}{Unif}
\DeclareMathOperator{\Iind}{\mathbf{Ind}}
\DeclareMathOperator{\Hhomcart}{\mathbf{Homcart}}
\DeclareMathOperator{\Hhomtopess}{\mathbf{Homtopess}}
\DeclareMathOperator{\Ppoints}{\mathbf{Points}}
\DeclareMathOperator{\Ppoint}{\mathbf{Point}}
\DeclareMathOperator{\Ttors}{\mathbf{Tors}}
\DeclareMathOperator{\Ssex}{\mathbf{Sex}}
\DeclareMathOperator{\Ffib}{\mathbf{Fib}}
\DeclareMathOperator{\Vvois}{\mathbf{Vois}}
\DeclareMathOperator{\Ppointess}{\mathbf{Pointess}}
\DeclareMathOperator{\scrHom}{\mathbf{Hom}}
\DeclareMathOperator{\sheafPic}{\mathbf{Pic}}
\DeclareMathOperator{\sheafAut}{\mathbf{Aut}}
\DeclareMathOperator{\Hhomtoplocan}{\mathbf{Homtoplocan}}
\DeclareMathOperator{\Hhomtopan}{\mathbf{Homtopan}}
\DeclareMathOperator{\Pptgeom}{\mathbf{Ptgéom}}
\DeclareMathOperator{\Morsite}{Morsite}
\DeclareMathOperator{\Fib}{Fib}
\DeclareMathOperator{\Fais}{Fais}
\DeclareMathOperator{\kar}{kar}
\DeclareMathOperator{\Groupes}{Groupes}
\DeclareMathOperator{\Restr}{Restr}
\DeclareMathOperator{\Isom}{Isom}
\DeclareMathOperator{\Prol}{Prol}
\DeclareMathOperator{\Res}{Res}
\DeclareMathOperator{\Tors}{Tors}
\DeclareMathOperator{\esp}{esp}
\DeclareMathOperator{\es}{es}
\DeclareMathOperator{\od}{od}
\DeclareMathOperator{\gr}{gr}
\DeclareMathOperator{\sob}{sob}
\DeclareMathOperator{\dfn}{dfn}
\DeclareMathOperator{\Sch}{Sch}
\DeclareMathOperator{\ouv}{ouv}
\DeclareMathOperator{\cons}{cons}
\DeclareMathOperator{\Nerf}{Nerf}
\DeclareMathOperator{\Points}{Points}
\DeclareMathOperator{\Espsob}{Espsob}
\DeclareMathOperator{\Homtop}{Homtop}
\DeclareMathOperator{\Pt}{Pt}
\DeclareMathOperator{\Quot}{Quot}
\DeclareMathOperator{\glob}{glob}
\DeclareMathOperator{\Homfib}{Homfib}
\DeclareMathOperator{\Qouv}{Qouv}
\DeclareMathOperator{\Quotouv}{Quotouv}
\DeclareMathOperator{\ens}{ens}
\DeclareMathOperator{\ab}{ab}
\DeclareMathOperator{\an}{an}
\DeclareMathOperator{\fs}{fs}
\DeclareMathOperator{\red}{red}
\DeclareMathOperator{\Ab}{Ab}
\DeclareMathOperator{\cosq}{cosq}
\DeclareMathOperator{\Ann}{Ann}
\DeclareMathOperator{\Zar}{Zar}


%%%%%%%%%%%%
%%%%% SGA4-2
\DeclareMathOperator{\Sex}{Sex}
\DeclareMathOperator{\vois}{vois}
\DeclareMathOperator{\Vois}{Vois}
\DeclareMathOperator{\Ext}{Ext}
\DeclareMathOperator{\ext}{ext}
\DeclareMathOperator{\Fer}{Fer}
\DeclareMathOperator{\Add}{Add}
\DeclareMathOperator{\Cocart}{Cocart}
\DeclareMathOperator{\cosk}{cosk}
\DeclareMathOperator{\plat}{plat}

\DeclareMathOperator{\faisc}{faisc}
\DeclareMathOperator{\sk}{sk}
\DeclareMathOperator{\ch}{ch}

\def\HH{\mathbf{H}}
\def\qc{\mathrm{qc}}
\def\et{\mathrm{et}}
\def\coh{\mathrm{coh}}

\def\Hvr#1{%
\begin{array}[b]{l}
{\kern2pt\scriptstyle r}\\[-8pt]
{\kern1pt\scriptstyle \vee}\\[-5pt]
H^{#1}
\end{array}}

\def\Hvv#1{%
\begin{array}[b]{l}
\kern2pt\rotatebox{90}{${\scriptscriptstyle\ll}$}\\[-6pt]
H^{#1}
\end{array}}

\def\Hvi#1{%
\begin{array}[b]{l}
{\kern0.2pt\scriptscriptstyle \infty}\\[-8pt]
{\kern1pt\scriptstyle \vee}\\[-5pt]
H^{#1}
\end{array}}

\def\Hv#1{%
\begin{array}[b]{l}
{\kern1pt\scriptstyle \vee}\\[-5pt]
H^{#1}
\end{array}}

\DeclareMathOperator{\cart}{cart}
\DeclareMathOperator{\Ffaisc}{\mathbf{Faisc}}


\def\Ev#1{%
\begin{array}[b]{l}
{\kern1pt\scriptstyle \vee}\\[-5pt]
E'_{#1}
\end{array}}

\DeclareMathOperator{\Cont}{Cont}
\DeclareMathOperator{\Ccont}{\mathscr{C}\!\mathit{ont}}
\DeclareMathOperator{\Site}{Site}
\newcommand\MM{\mathbf{M}}
\def\limtop#1{{\displaystyle{\varprojlim_{#1}\!\mbox{\rm top }}}}
\def\Limtop{\mathop{\varprojlim\!\mbox{\rm top}}}
\DeclareMathOperator{\spec}{spec}
\def\ss{\scriptstyle{\mathscr{S}}}
\DeclareMathOperator{\Gal}{Gal}
\def\smallm{\scriptstyle{\mathscr{M}}}
\DeclareMathOperator{\End}{End}
\DeclareMathOperator{\ass}{ass}
\DeclareMathOperator{\zar}{zar}
\DeclareMathOperator{\Ogg}{Ogg}
\DeclareMathOperator{\res}{res}
\DeclareMathOperator{\tr}{tr}
\DeclareMathOperator{\Supp}{Supp}
\DeclareMathOperator{\codim}{codim}
\DeclareMathOperator{\rk}{rk}
\DeclareMathOperator{\cur}{cur}
\DeclareMathOperator{\cdqc}{cdqc}
\newcommand\DD{\mathbf{D}}
\DeclareMathOperator{\cl}{cl}
\def\JJ{\mathbf{J}}
\def\FF{\mathbf{F}}
\def\BB{\mathbf{B}}
\DeclareMathOperator{\Alget}{Alget}
\DeclareMathOperator{\Ass}{Ass}
\DeclareMathOperator{\supp}{supp}
\DeclareMathOperator{\rad}{rad}
\DeclareMathOperator{\prof}{prof}
\DeclareMathOperator{\cint}{int}
\DeclareMathOperator{\sic}{sic}
\DeclareMathOperator{\Tor}{Tor}
\DeclareMathOperator{\tors}{tors}
\DeclareMathOperator{\Sym}{Sym}
\DeclareMathOperator{\parf}{parf}
\DeclareMathOperator{\tor}{tor}
\DeclareMathOperator{\defn}{defn}
\DeclareMathOperator{\Div}{Div}
\DeclareMathOperator{\Pol}{Pol}
\DeclareMathOperator{\Fex}{Fex}
\DeclareMathOperator{\Lie}{Lie}
\DeclareMathOperator{\tot}{tot}
\DeclareMathOperator{\Aut}{Aut}
\DeclareMathOperator{\HOM}{HOM}
\DeclareMathOperator{\pch}{pch}
\DeclareMathOperator{\Ch}{Ch}
\DeclareMathOperator{\EXT}{EXT}
\DeclareMathOperator{\PIC}{PIC}
\DeclareMathOperator{\Rf}{Rf}
\DeclareMathOperator{\inj}{inj}
\newcommand\VV{\mathbf{V}}

\def\Xv{%
\begin{array}[b]{l}
{\kern1pt\scriptstyle \vee}\\[-5pt]
X
\end{array}}

%%%%%%%
\newcommand\TT{\mathbf{T}}
\newcommand\XX{\mathbf{X}}
\newcommand\EE{\mathbf{E}}
%%%%%%

\def\qedsymbol{\textup{C.Q.F.D.}}
\let\CQFD\qed
\let\cqfd\qed

%%%%%\bigdot
\font\bigsymb=cmsy10 at 4pt
\def\bigdot{{\kern1.2pt\raise 1.5pt\hbox{\bigsymb\char15}}}
%%%%%

\def\Lright{\mathop{\vcenter{\hbox{\hbox{$\longrightarrow$}\lower8pt \hbox{\kern-16pt\hbox{$\longrightarrow$}}}}}}

\def\Aa{{\mathcal A}}
\def\Bb{{\mathcal B}}
\def\Cc{{\mathcal C}}
\def\Dd{{\mathcal D}}
\def\Ee{{\mathcal E}}
\def\Ff{{\mathcal F}}
\def\Gg{{\mathcal G}}
\def\Hh{{\mathcal H}}
\def\Ii{{\mathcal I}}
\def\Jj{{\mathcal J}}
\def\Kk{{\mathcal K}}
\def\Ll{{\mathcal L}}
\def\Mm{{\mathcal M}}
\def\Nn{{\mathcal N}}
\def\Oo{{\mathcal O}}
\def\Pp{{\mathcal P}}
\def\Qq{{\mathcal Q}}
\def\Rr{{\mathcal R}}
\def\Ss{{\mathcal S}}
\def\Tt{{\mathcal T}}
\def\Uu{{\mathcal U}}
\def\Vv{{\mathcal V}}
\def\Ww{{\mathcal W}}
\def\Xx{{\mathcal X}}
\def\Yy{{\mathcal Y}}
\def\Zz{{\mathcal Z}}

\newcommand{\CC}{\mathbf{C}}
\newcommand{\GG}{\mathbf{G}}
\newcommand{\KK}{\mathbf{K}}
\newcommand{\LL}{\mathbf{L}}
\newcommand{\NN}{\mathbf{N}}
\newcommand{\PP}{\mathbf{P}}
\newcommand\QQ{\mathbf{Q}}
\newcommand\RR{\mathbf{R}}
\newcommand\ZZ{\mathbf{Z}}

\def\to{\mathchoice{\longrightarrow}{\rightarrow}{\rightarrow}{\rightarrow}}
\def\from{\mathchoice{\longleftarrow}{\leftarrow}{\leftarrow}{\leftarrow}}
\def\mto{\mathchoice{\longmapsto}{\mapsto}{\mapsto}{\mapsto}}
\def\mfrom{\mathrel{\from\joinrel\mapstochar\,}}
\def\hto{\mathrel{\lhook\joinrel\to}}
\def\hfrom{\mathrel{\from\joinrel\rhook}}

\def\implies{\mathchoice{\Longrightarrow}{\Rightarrow}{\Rightarrow}{\Rightarrow}}
\def\implied{\mathchoice{\Longleftarrow}{\Leftarrow}{\Leftarrow}{\Leftarrow}}
\def\iff{\mathchoice{\Longleftrightarrow}{\Leftrightarrow}{\Leftrightarrow}{\Leftrightarrow}}

\def\To#1{\mathchoice{\xrightarrow{\labelstyle~~#1~}}{\stackrel{#1}{\longrightarrow}}{}{}}
\def\Mto#1{\mathrel{\mapstochar\To{#1}}}
\def\From#1{\mathchoice{\xleftarrow{\labelstyle~~#1~}}{\stackrel{#1}{\longleftarrow}}{}{}}

\newcommand{\isomto}{\overset{\sim}{\longrightarrow}}
\newcommand{\isomfrom}{\overset{\sim}{\longleftarrow}}

\newcommand{\IFF}{\unskip~$\iff$~\ignorespaces}
\newcommand{\THEN}{\unskip~$\implies$~\ignorespaces}

\let\leq\leqslant
\let\geq\geqslant
\let\le\leqslant
\let\ge\geqslant

\def\cf{{cf.}\kern.3em}
\def\Cf{{Cf.}\kern.3em}
\def\ie{{i.e.}\kern.3em}
\def\iev{{i.e.},\xspace}
%\def\loccito{loc.\kern2pt cit.\xspace}
%\def\loccitn{\textit{loc.\kern2pt cit.}\xspace}
%\def\loccit{\sisi{\loccito}{\loccitn}}
\def\EGA{\textup{EGA}\kern.2em}
\def\SGA{\textup{SGA}\kern.3em}

\let\oldS\S
\def\S{\oldS\kern .15em}
\def\ptbl{.\kern .2em}
\def\resp{resp.\kern.3em}
\def\Exp{Exp.\kern.3em}

\let\bar\overline
\let\hat\widehat
\let\tilde\widetilde
\let\dpl\displaystyle
\let\ldots\dots
\let\emptyset\varnothing
\let\skpt\skippointrait
\let\epsilon\varepsilon

\def\sfrac#1#2{{#1}/{#2}}

\DeclareMathOperator{\cd}{cd}
%\DeclareMathOperator{\Hhom}{\mathcal{H}\!\mathit{om}}
\DeclareMathOperator{\Hhom}{\mathbf{Hom}}
%\DeclareMathOperator{\Eext}{\mathcal{E}\!\mathit{xt}}
\DeclareMathOperator{\Eext}{\mathbf{Ext}}
\DeclareMathOperator{\Spec}{Spec}
\DeclareMathOperator{\Tr}{Tr}

\let\emptyset=\varnothing

%%%%%%%%%%%%%%%
%%% expose XVII
%%%%%%%%%%%%%%%

\makeatletter
\let\old@optexpose\optexpose
\newcommand{\empty@optexpose}[1]{}
\newcommand{\paren@optexpose}[1]{\old@optexpose{#1}(}
\newcommand{\refnoexpose}[1]{\let\optexpose\empty@optexpose\ref{#1}\let\optexpose\old@optexpose}%\refnoexpose{X.cor5.2} engendre "5.2" au lieu de "X 5.2"
\newcommand{\neweqref}[1]{\let\optexpose\paren@optexpose\textup{\ref{#1})}\let\optexpose\old@optexpose}%\neweqref{XII.eq4.1} engendre "XII (4.1)" au lieu de "(XII 4.1)"
\makeatother
\newcommand{\refvb}[2][VI${}_{\text{B}}$ ]{\sisi{#1\refnoexpose{#2}}{\ref{#2}}}

\newcommand{\leftexp}[2]{{\vphantom{#2}}^{#1}{#2}}
\newcommand{\AAA}{\mathbf{A}}
\newcommand{\Sf}{\mathfrak{S}}
\DeclareMathOperator*{\quotprojlim}{\text{``$\varprojlim$''}}
\DeclareMathOperator*{\quotinjlim}{\text{``$\varinjlim$''}}
\DeclareMathOperator{\Fl}{F\sisi{\ell}{l}}
\DeclareMathOperator{\Pl}{P\sisi{\ell}{l}}
%\DeclareMathOperator{\cl}{c\sisi{\ell}{l}}
\DeclareMathOperator{\tdf}{tdf}
\DeclareMathOperator{\rg}{rg}
\newcommand{\iInd}{\sisi{I}{i}nd}
\newcommand{\iind}{ind}
\newcommand{\ppro}{pro}
\newcommand{\tTor}{\sisi{T}{t}or}
\newcommand{\ttor}{tor}
\newcommand{\chap}{Chap.}
\newcommand{\Ffex}{\Ff\mathit{ex}}
\newcommand{\sheafk}{\underline{k}}
\newcommand{\sheafn}{\underline{n}}
\newcommand{\sheafGamma}{\underline{\Gamma}}

% Repeated references to an already typeset footnote must not use
% \footnotemark without a matching \footnotetext: Hyperref would then create
% a live annotation to a destination that cannot exist.  This helper keeps the
% current printed footnote mark and links it to an explicit stable note target.
\makeatletter
\newcommand{\sgarepeatedfootnotemark}[1]{%
  \begingroup
  \edef\@thefnmark{\thefootnote}%
  \hyperlink{#1}{\@makefnmark}%
  \endgroup}
\makeatother

\sisi{
\DeclareMathOperator{\simarrow}{\rule{.15mm}{1.8mm}\kern -3.4pt\rightsquigarrow}
\newcommand{\sheafell}{\underline{\ell}}
\newcommand{\sheafZZ}{\underline{\ZZ}}
}{
\newcommand{\nderef}[1]{la N.D.E.~(\ref{#1}) page~\pageref{#1}}
}

%%%%%%%%%%%%%%%%%
%%%%% expose XVII
%%%%%%%%%%%%%%%%%

% END INCLUDED SGA4_English_current_buildable_source_20260906/source/fichiers-configuration-latex/macros.tex


\externaldocument{../01/01_English.tex}
\externaldocument{../02/02_English.tex}
\externaldocument{../03/03_English.tex}
\externaldocument{../04/04_English.tex}
\externaldocument{../05/05_English.tex}
\externaldocument{../05bis/05bis_English.tex}

\setboolean{original}{false}

\begin{document}

\begin{center}
VI. Finiteness Conditions. Fibered Topoi and Sites. Applications
to Questions of Passage to the Limit.\\
A. Grothendieck and J.-L. Verdier
\end{center}

\selectlanguage{english}version: \showgitstatus
\tableofcontents

\else
\chapter[Finiteness conditions and passage to the limit]{Finiteness Conditions. Fibered Topoi and Sites. Applications
to Questions of Passage to the Limit.}
%\addtocontents{toc}{A. Grothendieck and J.-L. Verdier}
\begin{center}
A. Grothendieck and J.-L. Verdier
\end{center}
\fi
\phantomsection\MakeLinkTarget*{sga4.exp.VI}\label{sga4.exp.VI}

% BEGIN INCLUDED SGA4_English_current_buildable_source_20260906/source/06/drafts/SGA4_VI_LINES_44_390_English.texfrag
\setcounter{pageoriginal}{164}

\setcounter{section}{-1}
\section{Introduction}\etiquette[0]{VI.sec0}

In\pageoriginale the present exposé, we study two closely related
types of questions.

First, we undertake a systematic study of finiteness conditions for
topoi, successively examining the notions of \emph{quasi-compactness},
\emph{quasi-separatedness}, and \emph{coherence}
(= quasi-compactness $+$ quasi-separatedness). These are inspired by
the familiar analogous notions from scheme theory. We study them
first for an object $X$ of a topos and for an arrow $u : X \to Y$ of
a topos (\hyperref[VI.sec1]{\S 1}), then for the topos $E$ itself (\hyperref[VI.sec2]{\S 2}), and finally for
a morphism of topoi $f : E \to E'$ (\hyperref[VI.sec3]{\S 3}). \hyperref[VI.sec4]{Section 4} studies these
notions in the special case of a topos obtained by gluing two topoi
(IV 9). For the subsequent sections, the essential facts to retain
from these somewhat technical developments are the following: a
\emph{coherent topos} is a topos equivalent to one of the form
$C^{\sim}$, where $C$ is a small site in which finite projective
limits are representable and every covering family admits a finite
covering subfamily; and a \emph{coherent morphism}
$f : E \to E'$ of coherent topoi is a morphism that can be described
(up to equivalence) as the morphism associated with a morphism of
sites $f : C \to C'$, where $C$ and $C'$ are as above and the
corresponding functor $f^{\ast} : C'\to C$ is left exact.

It is no doubt useful to note that, with the finiteness notions used
here, and especially that of a coherent topos, we resolutely depart
(and for the first time) from the topological spaces familiar to
analysts and ``topologists''\footnote{The quotation marks indicate
that this is a ``\ldots'' unfamiliar with the notion of a topos
(\cf IV \hyperref[IV.subsec0.4]{0.4} for the meaning of the word ``topology'').}. Thus the
topos associated (IV \hyperref[IV.subsec2.1]{2.1}) with a separated topological space $X$ is
coherent only if $X$ is a finite set! In other words, the present
exposé is directed entirely toward the types of topoi arising from
algebraic geometry and algebra. In fact, all the topoi used thus far
in these disciplines (except, of course, in algebraic geometry by
transcendental methods over the field of complex numbers!) turn out
to be locally coherent.

Second, we develop two \emph{theorems on passage to the inductive
limit} for the cohomology of topoi. One
(\hyperref[VI.cor5.2]{5.2}) tells us that the functors
$H^{i}(E, -)$\pageoriginale on a coherent topos commute with
inductive limits of abelian sheaves. The other
(\hyperref[VI.subsec8.7]{8.7}) says essentially that if a topos $X$
is represented as a filtered projective limit of coherent topoi
$X_i$, with coherent transition morphisms
$f_{ij} : X_j \to X_i$, then the cohomology of $X$ (with coefficients
in an arbitrary abelian sheaf) is computed as the inductive limit of
the cohomologies of the $X_i$. To give this statement a precise
meaning, the main point is to specify the notion of a ``projective
limit of topoi,'' which is done in \S\hyperref[VI.sec7]{\S 7}, 8. This is a very useful
geometric notion, undoubtedly also in contexts other than that of
projective limits of schemes (which had served as the model for
M.\textsc{Artin} in his initial definition of this notion). We show,
for example, in (\hyperref[VI.subsec8.4]{8.4}) how the various notions
of being ``localized at a point'' (for a Noetherian space, for a
scheme with its usual topology, or even with its étale topology) can
be unified by defining the localization as the projective limit of
the ``neighborhoods'' of that point---in perfect agreement with the
intuitive idea of localization at a point.

The two limit-passage theorems will be stated again explicitly in the
next exposé in the special setting of étale cohomology of schemes
(\hyperrefext[VII][VII.prop3.3]{3.3} and
\hyperref[VI.prop5.7]{5.7}); in that case, they will moreover be used
constantly throughout the remainder of this Seminar, and practically
everywhere étale cohomology has been used up to the present. The
reader willing to accept these two theorems in the special case in
which they are stated in Exposé VII, and interested exclusively in
applications of étale cohomology, may therefore omit the present
exposé. It is true that the notion of a \emph{constructible object} of
a topos (\ie one whose structural morphism $X \to e$ to the final
object $e$ is coherent) will likewise play an important role
throughout the rest of the Seminar (whereas it plays only a very
secondary role in the present exposé, if only because, for a coherent
topos, it coincides with the notion of a coherent object, which takes
center stage here). But this notion is developed independently in
Exposé IX, in the special case of étale sheaves on schemes, without
using the general study made in the present exposé; in that case it
also exhibits very useful special phenomena that served as the basis
for its treatment in IX (beginning with \hyperref[IX.def2.3]{Definition IX 2.3}). Since the
text of IX predates the present\pageoriginale text of Exposé VI by
five years, yet remains quite usable for the reader, we have confined
ourselves at the end of Exposé IX to proving the equivalence between
the notion of constructibility used in IX and the one introduced
here.

\section{Finiteness conditions for the objects and arrows of a
  topos}\etiquette[1]{VI.sec1}

\begin{enonce*}{Definition 1.1}\etiquette[1.1]{VI.def1.1}
Let $C$ be a site. An object $X$ of $C$ is called quasi-compact if,
for every covering family $(X_i \to X)_{i\in I}$, there is a finite
subset $J$ of the index set $I$ such that the subfamily
$(X_i \to X)_{i \in J}$ is still covering.
\end{enonce*}

Since a $𝒰$-topos $E$ is always regarded as a site (IV \hyperref[IV.subsubsec1.1.1]{1.1.1}), the
preceding definition applies in particular to an object of $E$. This
case is moreover the general one, by the following proposition.

\begin{enonce*}{Proposition 1.2}\etiquette[1.2]{VI.prop1.2}
Let $C$ be a $𝒰$-site, let $\epsilon : C \to C^{\sim}$ be the
canonical functor, and let $X$ be an object of $C$. Then $X$ is a
quasi-compact object of the site $C$ if and only if $\epsilon(X)$ is
a quasi-compact object of the topos $C^{\sim}$.
\end{enonce*}

Suppose that $X$ is quasi-compact; we prove that $\epsilon(X)$ is as
well. Let $(F_i \to \epsilon(X))_{i \in I}$ be an epimorphic family.
For every $i \in I$, choose an epimorphic family
$\epsilon(X_{ij}) \to F_i$, $j \in J_i$. The composite
$\epsilon(X_{ij}) \to \epsilon(X)$ identifies with an element
$\alpha_{ij}$ of $\epsilon(X)(X_{ij})$. In general one cannot assert
that it comes from a morphism $X_{ij} \to X$, but, by the construction
of the associated sheaf $\epsilon(X)$, one knows that, for fixed $i$
and $j$, a covering family $Y_k \to X_{ij}$ can be found such that
the images of $\alpha_{ij}$ in the $\epsilon(X)(Y_k)$ come from
elements of $X(Y_k)$, \ie from morphisms $Y_k \to X$. Thus, after
refining the epimorphic family $\epsilon(X_{ij}) \to F_i$, we may
suppose that the composites
$\epsilon(X_{ij}) \to \epsilon(X)$ come from morphisms
$X_{ij} \to X$.\pageoriginale As $i$ and $j$ vary, the family of
$\epsilon(X_{ij}) \to \epsilon(X)$ is epimorphic, so (II
\hyperref[VI.cor4.4]{4.4}) the family of $X_{ij} \to X$ is covering.
By the hypothesis on $X$, it admits a finite covering subfamily, and
the corresponding finite family of morphisms
$\epsilon(X_{ij}) \to \epsilon(X)$ is again epimorphic. Since these
factor through finitely many of the morphisms
$F_i \to \epsilon(X)$, that finite family is already epimorphic.
This proves that $\epsilon(X)$ is quasi-compact.

Conversely, the fact that quasi-compactness of $\epsilon(X)$ implies
quasi-compactness of $X$ is trivial, by \loccit

The preceding result therefore shows that the notion of
quasi-compactness on sites reduces to the corresponding notion in
topoi.

\begin{enonce*}{Proposition 1.3}\etiquette[1.3]{VI.prop1.3}
Let $C$ be a site and let $(X_i \to X)_{i \in I}$ be a finite
covering family, with the $X_i$ quasi-compact. Then $X$ is
quasi-compact.
\end{enonce*}

To see this, we may assume, by \hyperref[VI.prop1.2]{1.2}, that $C$
is a topos. Let $(X'_j \to X)$ be a covering, \ie epimorphic, family
in $C$. Then, for every $i \in I$, the family obtained by the base
change $X_{i} \to X$ is epimorphic and hence (since $X_i$ is
quasi-compact) admits a finite epimorphic subfamily, corresponding to
a finite subset $J_i \subset J$ of the index set $J$. Set
$J'= \bigcup_{i\in I} J_i$. Then $J'$ is a finite subset of $J$
because $I$ is finite, and the family
$(X'_j \to X)_{j \in {J'}}$ is already epimorphic, since it is so
locally, \cqfd

\begin{enonce*}{Corollary 1.4}\etiquette[1.4]{VI.cor1.4}
Let\pageoriginale $(X_i)_{i \in I}$ be a family of objects of a topos
$E$, and let $X$ be their sum. Then $X$ is quasi-compact if and only
if all the $X_i$ are quasi-compact and all but finitely many of the
$X_i$ are isomorphic to the ``empty object'' $∅_{E}$ of $E$.
\end{enonce*}

Sufficiency follows immediately from \hyperref[VI.prop1.3]{1.3}
(since the sum does not change when the empty summands are omitted).
For necessity, note that a finite subfamily $(X_i)_{i \in J}$ must
cover $X$, which implies that, for $i \in I-J$, one has
$X_i = \fprod{X_{i}}{\amalg_{j=J}}{X}X_{j} =
\fprod{\amalg_{j=J}X_{i}}{X_j}{X} = ∅$ (II
\hyperref[VI.lem4.5.1]{4.5.1}). On the other hand, to show that every
$X_i$ is quasi-compact, note that $X \simeq X_i \amalg {X'}_i$, and
that every covering family $Y_j \to X_i$ of $X_i$ defines a covering
family $Y_j \amalg {X'}_i \to X_i \amalg {X'}_i = X$ of $X$. The
latter admits a finite covering subfamily, which plainly implies that
the corresponding finite family of $Y_j \to X_i$ is also covering;
thus $X_i$ is quasi-compact.

\begin{enonce*}[remark]{Remark 1.5.1}\etiquette[1.5.1]{VI.rem1.5.1}
The property for an object $X$ of a topos $E$ of being quasi-compact
depends only on the induced topos $E_{/X}$, and means that the final
object of the latter is quasi-compact. Likewise, an object $X$ of a
site $C$ is quasi-compact if and only if the final object of the
induced site $C_{/X}$(III \hyperref[III.subsec5.1]{5.1}) is quasi-compact.
\end{enonce*}

\setcounter{subsection}{5}
\setcounter{subsubsection}{1}
\subsubsection{}\etiquette[1.5.2]{VI.subsubsec1.5.2}
Let $X$ be an object of a topos $E$. Then $X$ is quasi-compact if and
only if every filtered increasing family $(X_i)_{i \in I}$ of
subobjects of $X$ that covers $X$ contains some $X_i$ equal to $X$.
This condition therefore depends\pageoriginale only on the ordered
set of subobjects of $X$ (\ie the open subsets of the induced topos
$E_{/X}$).

\subsubsection{}\etiquette[1.5.3]{VI.subsubsec1.5.3}
Consider a generating family $(X_i)_{i \in I}$ of $E$. It follows
from the definitions and from \hyperref[VI.prop1.3]{1.3} that, for an
object $X$ of $E$ to be quasi-compact, it is necessary that it admit
a covering family consisting of finitely many of the $X_i$, \ie that
$X$ be isomorphic to a quotient object of a finite sum of objects of
$E$; this condition is sufficient if the $X_i$ are quasi-compact.
Note also that, under the hypothesis made on $E$, after replacing the
generating family under consideration by a subfamily if necessary,
we may clearly suppose that $I \in 𝒰$. Using the fact that the set of
quotients of an object of $E$ is small (I \hyperref[I.cor7.5]{7.5}), we then find that the
full subcategory $E_{\qc}$ of $E$ formed by the quasi-compact objects
of $E$ is equivalent to a category $\in 𝒰$, \ie the cardinal of the set of
isomorphism classes of objects of $E_{\qc}$ is $\in 𝒰$.

\begin{enonce*}[remark]{Example 1.6.1}\etiquette[1.6.1]{VI.exem1.6.1}
Let $X$ be a topological space. An open subset $U$ of $X$ is a
quasi-compact object of the site $\Ouv(X)$ of open subsets of $X$ if
and only if $U$ is a quasi-compact space for the topology induced by
$X$. More generally, a sheaf $F$ on $X$ is a quasi-compact object of
the topos $\Top(X)$ if and only if the étalé space $X'$ over $X$
associated with $F$ is quasi-compact. (This follows from the preceding
assertion, from \hyperref[VI.rem1.5.1]{1.5}, and from the fact that $\Top (X)_{/F}$ is
equivalent to $\Top(X')$ (IV \hyperref[IV.subsec5.7]{5.7})).
\end{enonce*}

\setcounter{subsection}{6}
\setcounter{subsubsection}{1}
\subsubsection{}\etiquette[1.6.2]{VI.subsubsec1.6.2}
Consider\pageoriginale on the category ($\Sch$) of schemes (elements
of $𝒰$) one of the topologies $T_i$ of \cite{sga4.exp.VI.bib.6}
\hyperref[VI.def6.3]{6.3} (for example, the Zariski topology, the
étale topology, the fppf topology, or the fpqc topology). A scheme is
quasi-compact for this topology if and only if it is a quasi-compact
scheme in the usual sense. (With the notation of \loccit, only the
fact that every family $\in P'$ is finite is used.)

\begin{enonce*}{Definition 1.7}\etiquette[1.7]{VI.def1.7}
Let $E$ be a topos and let $f : X \to Y$ be an arrow of $E$. We say
that $f$ is a quasi-compact morphism if, for every arrow $Y'\to Y$
with $Y'$ quasi-compact, the object
$X'= \fprod{X}{Y'}{Y}$ is also quasi-compact. We say that $f$ is
quasi-separated if the diagonal morphism
$X \to \fprod{X}{X}{Y}$ is quasi-compact. We say that $f$ is coherent
if $f$ is quasi-compact and quasi-separated.
\end{enonce*}

\setcounter{subsection}{7}
\setcounter{subsubsection}{0}
\subsubsection{}\etiquette[1.7.1]{VI.subsubsec1.7.1}
Unwinding the definition of a quasi-separated morphism
$f : X \to Y$, we see that it also means that, for every pair of
arrows $g_1, g_2 : X'\rightrightarrows X$ ``over $Y$,'' \ie such that
$fg_1 = fg_2$, quasi-compactness of $X'$ implies quasi-compactness of
$\Ker(g_1, g_2)$.

\begin{enonce*}{Proposition 1.8}\etiquette[1.8]{VI.prop1.8}
Let $E$ be a topos.
\begin{enumerate}
\renewcommand{\theenumi}{\roman{enumi}}
\renewcommand{\labelenumi}{(\theenumi)}
\item Every isomorphism in $E$ is a coherent morphism, \ie is
  quasi-compact and quasi-separated. The composite of two
  quasi-compact (\resp quasi-separated, \resp coherent) morphisms is
  quasi-compact (\resp quasi-separated, \resp coherent).

\item Let\pageoriginale $f : X \to Y$ and $g : Y' \to Y$ be
  morphisms in $E$, and let
  $f': X'= \fprod{X}{Y'}{Y} \to Y'$ be deduced from $f$ by the base
  change $g$. If $f$ is quasi-compact (\resp quasi-separated, \resp
  coherent), then so is $f'$.

\item Let $f : X \to Y$ and $g : Y \to Z$ be
  morphisms in $E$. If
  $gf$ is quasi-separated, then $f$ is quasi-separated.

\item Let $f : X \to Y$ and $g : Y \to Z$ be morphisms in $E$, with
  $g$ quasi-separated. If $gf$ is quasi-compact (\resp coherent),
  then $f$ is quasi-compact (\resp coherent).
\end{enumerate}
\end{enonce*}

The proofs are well known (\Cf \EGA IV 1 or EGA I, 2nd edition).
Statements (i) and (ii) in the quasi-compact case follow trivially
from the definitions. In the quasi-separated case, stability under
base change in (ii) follows immediately from the definition and the
preceding case, as does the fact that every isomorphism is
quasi-separated. Stability under composition $X \to Y \to Z$ follows
from the diagram
\[\vcenter{
\xymatrix@R=1.5cm@C=1.5cm{
X \ar[d] \ar@/_3.5pc/[dd]^(.2){\Delta_{X/Y}}_{\Delta_{X/Z}}&\\
\fprod{X}{X}{Y} \ar[d]^{u} \ar[rr]&  &Y \ar[d]^{\Delta_{Y/Z}}\\
\fprod{X}{X}{Z}  \ar[rr]&        & \fprod{Y}{Y}{Z}
}}\leqno{(*)}
\]
whose square\pageoriginale is Cartesian. By hypothesis, $\Delta_{Y/Z}$ is
quasi-compact, and hence so is $u$, which is obtained from it by base
change; moreover, $\Delta_{X/Y}$ is quasi-compact, and therefore so is
the composite $u\Delta_{X/Y} = \Delta_{X/Z}$. This proves (i) and
(ii), the coherent case following by combining the two preceding
cases. To prove the first case considered in (iv), consider the
following diagram with Cartesian squares:
\[\vcenter{
\xymatrix@C=2.75cm@R=1.5cm{
  &     X \ar[r]^{gf}          & Z\\
  X \ar[r]^{\Gamma_f=(\id_X,f)}  \ar[d]
  \ar@/_2pc/[rr]_(.55){f\qquad\qquad \fprod{f}{\id_{Y}}{Z}}  &\fprod{X}{Y}{Z} \ar[r]^{\pr_2}\ar[d]\ar[u]^{\pr_1}& Y\ar[u]_{g}\\
Y \ar[r]^{\Delta_{Y/Z}}&  \fprod{Y}{Y}{Z}& \quad .
}}\leqno{(**)}
\]
By hypothesis, $g$ is quasi-separated, so $\Delta_{Y/Z}$ is
quasi-compact and hence $\Gamma_f$ is quasi-compact by base change.
Likewise, by hypothesis $gf$ is quasi-compact, so $\pr_2$ is also
quasi-compact by base change. Consequently,
$f = \pr_2 \Gamma_f$ is quasi-compact as a composite of two
quasi-compact morphisms. The corresponding case of (iv) follows
immediately from the preceding case and (iii), which remains to be
proved. For this, return to diagram $(\ast)$. The hypothesis that
$gf$ is quasi-separated means that
$\Delta_{X/Z} = u \Delta_{X/Y}$ is quasi-compact, while the conclusion
that $f$ is quasi-separated means that $\Delta_{X/Y}$ is
quasi-compact. This follows from (iv) once we have proved that $u$ is
quasi-separated. But $u$ is\pageoriginale plainly a monomorphism, and indeed we
have the following trivial result:

\begin{enonce*}{Corollary 1.8.1}\etiquette[1.8.1]{VI.cor1.8.1}
Every monomorphism is quasi-separated.
\end{enonce*}

Indeed, if $u : X \to Y$ is a monomorphism, then $\Delta_{X/Y}$ is an
isomorphism and is therefore quasi-compact, \cqfd

\begin{enonce*}{Corollary 1.8.2}\etiquette[1.8.2]{VI.cor1.8.2}
Let $f:X \to Y$ and $f':X' \to Y'$ be two $S$-morphisms in $E$. If
$f$ and $f'$ are quasi-compact (\resp quasi-separated, \resp
coherent), then so is
$\fprod{f}{f'}{S} : \fprod{X}{X'}{S} \to \fprod{Y}{Y'}{S}$.
\end{enonce*}

This follows in the usual way by combining (i) and (ii).

\begin{enonce*}{Remark 1.9.1}\etiquette[1.9.1]{VI.rem1.9.1}

% END INCLUDED SGA4_English_current_buildable_source_20260906/source/06/drafts/SGA4_VI_LINES_44_390_English.texfrag


% BEGIN INCLUDED SGA4_English_current_buildable_source_20260906/source/06/drafts/SGA4_VI_LINES_391_677_English.texfrag
Let $E$ be any category, and let $Q$ be a subset of $\ob E$ stable
under isomorphism (thus, giving $Q$ amounts to giving a strictly full
subcategory of $E$), playing the role of the ``quasi-compact''
objects. Suppose that fiber products are representable in $E$. One
can then paraphrase Definition \hyperref[VI.def1.7]{1.7} to define the
notions of $Q$-quasi-compact, $Q$-quasi-separated, and $Q$-coherent
morphisms. Proposition \hyperref[VI.prop1.8]{1.8} and Corollary
\hyperref[VI.cor1.8.1]{1.8.1} then hold, as does
\hyperref[VI.cor1.8.2]{1.8.2} when products of two objects are
representable in $E$.
\end{enonce*}

\setcounter{subsection}{9}
\setcounter{subsubsection}{1}
\subsubsection{}\etiquette[1.9.2]{VI.subsubsec1.9.2}
Let $S$ be an object of $E$, and let $f:X\to Y$ be a morphism of
$E_{/S}$. It follows immediately from \hyperref[VI.rem1.5.1]{Remark 1.5} that $f$ is
quasi-compact (\resp quasi-separated, \resp coherent) if and only if
the corresponding morphism\pageoriginale in $E$ (obtained from $f$
by the ``forget $S$'' functor) has the same property. In particular,
if $f:X\to Y$ is a morphism in $E$, then, by taking $S=Y$ above, we
see that $f$ is quasi-compact (\resp quasi-separated, \resp coherent)
if and only if the structural morphism of the object $(X,f)$ of the
induced topos $E_{/Y}$, with target the final object of that topos, is
quasi-compact (\resp quasi-separated, \resp coherent). One also says
that (by means of the structural morphism $f$) $X$ \emph{is
quasi-compact (\emph{\resp} quasi-separated, \emph{\resp} coherent)
over} $Y$, an expression that may equally well be interpreted as
referring to the topos $E$ or to the induced topos $E_{/Y}$ (in which
$Y$ is regarded as a final object).

\subsubsection{}\etiquette[1.9.3]{VI.subsubsec1.9.3}
An object $X$ of a topos $E$ is called \emph{constructible} if it is
coherent over the final object. Every morphism from one constructible
object to another is coherent (\hyperref[VI.prop1.8]{1.8} (iv)). The
full subcategory $E_{\cons}$ of $E$ formed by the constructible
objects of $E$ is stable under finite inverse limits, as follows
immediately from \hyperref[VI.cor1.8.2]{1.8.2} and
\hyperref[VI.prop1.8]{1.8} (ii).

\subsubsection{}\etiquette[1.9.4]{VI.subsubsec1.9.4}
The notions of \hyperref[VI.def1.7]{1.7} are of real interest only in
topoi $E$ that admit a generating family formed by quasi-compact
objects. In that case, these notions are ``local on the base'' (IV
\hyperref[IV.subsubsec8.5.4]{8.5.4}), as we now show.

\begin{enonce*}{Proposition 1.10}\etiquette[1.10]{VI.prop1.10}
Suppose\pageoriginale that the topos $E$ admits a generating
subcategory $C$ formed by quasi-compact objects, and let $f:X\to Y$
be a morphism in $E$.
\begin{enumerate}
\renewcommand{\theenumi}{\roman{enumi}}
\renewcommand{\labelenumi}{(\theenumi)}
\item The morphism $f$ is quasi-compact if and only if, for every
  morphism $S\to Y$ with $S\in\Ob C$, the object
  $\fprod{X}{S}{Y}$ is quasi-compact.

\item Let $Y_i\to Y$, $i\in I$, be a covering family. The morphism
  $f$ is quasi-compact (\emph{\resp} quasi-separated,
  \emph{\resp} coherent) if and only if, for every $i\in I$, the
  morphism $f_i:X_i=\fprod{X}{Y_i}{Y}\to Y_i$ has the same property.
\end{enumerate}
\end{enonce*}

Let us prove (i). Necessity is immediate from the definition. For
sufficiency, we must prove that, for every quasi-compact object $Y'$
and every morphism $Y'\to Y$, the object
$X'=\fprod{X}{Y'}{Y}$ is quasi-compact. There is a covering family
$S_i\to Y'$ with the $S_i$ in $\Ob C$, and since $Y'$ is
quasi-compact, this covering family may be taken to be \emph{finite}.
By hypothesis the objects
${X'}_i=\fprod{X}{S_i}{Y}\simeq\fprod{X'}{S_i}{Y'}$ are
quasi-compact, and since they cover $X'$ (the $S_i$ covering $Y'$),
$X'$ is quasi-compact by \hyperref[VI.prop1.3]{1.3}.

Let us prove (ii). It is enough to treat the quasi-compact case,
since the quasi-separated case follows from Definition
\hyperref[VI.def1.7]{1.7}, and the coherent case likewise follows by
combining the two preceding cases. The necessity has already been
proved in \hyperref[VI.prop1.8]{1.8} (ii). It remains to show that if
the $f_i$ are quasi-compact, then $f$ is quasi-compact, \ie that, for
every quasi-compact $Y'$ and every morphism $Y'\to Y$, the object
$X'=\fprod{X}{Y'}{Y}$ is quasi-compact. Since $(Y_i\to Y)$ is a
covering family,\pageoriginale there exist objects $S_{\alpha}$ over
$Y$ that cover $Y'$ and such that each composite
$S_{\alpha}\to Y'\to Y$ factors through some $Y_i$. Since $C$ is
generating, the $S_{\alpha}$ may be taken in $\Ob C$, and since $Y'$
is quasi-compact, the covering family may be taken finite. The
$\fprod{X'}{S_{\alpha}}{Y'}$ then form a finite covering family of
$X'$, and by \hyperref[VI.prop1.3]{1.3} we are reduced to proving
that each
$\fprod{X'}{S_{\alpha}}{Y'}=\fprod{X}{S_{\alpha}}{Y}$ is
quasi-compact. Choosing a factorization of $S_{\alpha}\to Y$ through
some $Y_i$, the object in question is also isomorphic to
$\fprod{X_i}{S_{\alpha}}{Y_i}$, which is indeed quasi-compact since
$X_i$ is quasi-compact over $Y_i$ and $S_{\alpha}$ is quasi-compact.

\begin{enonce*}{Corollary 1.11}\etiquette[1.11]{VI.cor1.11}
Let $g:Y\to Z$ be a morphism in the topos $E$, and let
$(f_i:X_i\to Y)_{i\in I}$ be a covering family of morphisms with
target $Y$. Then:
\begin{enumerate}
\renewcommand{\theenumi}{\roman{enumi}}
\renewcommand{\labelenumi}{(\theenumi)}
\item Suppose that $I$ is finite. If the $gf_i$ are quasi-compact,
  then so is $g$.

\item Suppose that the $f_i$ are quasi-compact. If the $gf_i$ are
  quasi-separated, then so is $g$. If $I$ is finite and the $gf_i$
  are coherent, then $g$ is coherent.
\end{enumerate}
\end{enonce*}

Case (i) follows immediately from the definitions and
\hyperref[VI.prop1.3]{1.3}, without any hypothesis on $E$. To prove
(ii), it is enough, by (i), to treat the nonparenthetical case.
Consider diagram \hyperref[VI.prop1.8]{1.8} $(\ast)$ with $X$
replaced by some $X_i$. Since the family of the $f_i$ is covering,
so is the family of the $\fprod{f_i}{\id_Y}{Z}$ obtained from it by
base change. Hence,\pageoriginale by \hyperref[VI.prop1.10]{1.10}
(ii), to prove that $\Delta_{Y/Z}$ is quasi-compact, it is enough to
prove that the $\Gamma_{f_i}$ are quasi-compact. Since $gf_i$ is
quasi-separated by hypothesis, so is
$\pr_2:\fprod{X_i}{Y}{Z}\to Y$, which is obtained from it by base
change. On the other hand, $\pr_2\Gamma_{f_i}=f_i$ is quasi-compact
by hypothesis. Thus $\Gamma_{f_i}$ is quasi-compact by
\hyperref[VI.prop1.8]{1.8} (iv), applied to the composite of
$\Gamma_{f_i}$ and $\pr_2$, \cqfd

\begin{enonce*}{Corollary 1.12}\etiquette[1.12]{VI.cor1.12}
Let $(X_i)_{i\in I}$ be a finite family of objects of $E$, let $X$ be
the coproduct of the $X_i$, let $f_i:X_i\to Y$ be morphisms in $E$,
and let $f:X\to Y$ be the corresponding morphism. The morphism $f$
is quasi-compact (\resp quasi-separated, \resp coherent) if and only
if the $f_i$ have the same property; in the quasi-separated case, the
hypothesis that $I$ be finite may be omitted.
\end{enonce*}

Sufficiency is a special case of \hyperref[VI.cor1.11]{1.11}, taking
into account that the canonical morphisms $X_i\to X$ form a finite
covering family and are quasi-compact morphisms, as follows
immediately from the definition and \hyperref[VI.cor1.4]{1.4}.
Necessity follows from transitivity, \hyperref[VI.prop1.8]{1.8} (i),
taking into account that the morphisms $X_i\to X$ are even coherent
(for they are quasi-separated by \hyperref[VI.cor1.8.1]{1.8.1}).

\begin{enonce*}{Definition 1.13}\etiquette[1.13]{VI.def1.13}
Let $X$ be an object of a topos $E$. We say that $X$ is a
\emph{quasi-separated object of the topos} $E$ if, for every
quasi-compact object $S$ of $E$, every morphism from $S$ to $X$ is
quasi-compact, \ie (\hyperref[VI.def1.7]{1.7}) if, for any two
quasi-compact objects $S,T$ of $E$ and morphisms $S\to X$, $T\to X$,
the fiber product $\fprod{S}{T}{X}$ is always quasi-compact. We say
that $X$ is a coherent object of the topos $E$ if $X$ is
quasi-compact and quasi-separated.
\end{enonce*}

\begin{enonce*}{Proposition 1.14}\etiquette[1.14]{VI.prop1.14}
Let\pageoriginale $f:X\to Y$ be a morphism in a topos $E$.
\begin{enumerate}
\renewcommand{\theenumi}{\roman{enumi}}
\renewcommand{\labelenumi}{(\theenumi)}
\item If $Y$ is quasi-compact, then $f$ quasi-compact implies that
  $X$ is quasi-compact. If $Y$ is quasi-separated, then $X$
  quasi-compact implies that $f$ is quasi-compact. If $Y$ is
  coherent, $f$ is quasi-compact if and only if $X$ is quasi-compact.

\item If $Y$ is quasi-compact (\emph{\resp} quasi-separated,
  \emph{\resp} coherent) and $f$ is quasi-compact
  (\emph{\resp} quasi-separated, \emph{\resp} coherent), then $X$
  is quasi-compact (\emph{\resp} quasi-separated,
  \emph{\resp} coherent).
\end{enumerate}
\end{enonce*}

The first two assertions of (i) are contained immediately in the
definitions of a quasi-compact morphism
(\hyperref[VI.def1.7]{1.7}) and, respectively, a quasi-separated
object (\hyperref[VI.def1.13]{1.13}); the third follows by combining
the first two. The first case of (ii) is merely a restatement, and the
third again follows by combining the first two. It remains to treat
the quasi-separated case. Let $S$ be a quasi-compact object of $E$
and $g:S\to X$ a morphism; we prove that $g$ is quasi-compact. Since
$Y$ is quasi-separated, $fg$ is quasi-compact, and since $f$ is
quasi-separated, it follows (\hyperref[VI.prop1.8]{1.8} (iii)) that
$g$ is quasi-compact, \cqfd

\begin{enonce*}{Corollary 1.15}\etiquette[1.15]{VI.cor1.15}
Every subobject of a quasi-separated object of $E$ is quasi-separated.
Let $(X_i)_{i\in I}$ be a family of objects of $E$, with $I\in 𝒰$;
then $X=\amalg_iX_i$ is quasi-separated (\emph{\resp} coherent) if
and only if the $X_i$ are.
\end{enonce*}

The\pageoriginale first assertion follows from
\hyperref[VI.prop1.14]{1.14} (ii) and
\hyperref[VI.cor1.8.1]{1.8.1}, but is also immediate from the
definition. For the second assertion in the nonparenthetical case,
since the $X_i\to X$ are monomorphisms, it remains only to prove
sufficiency, which follows easily from the definitions and
\hyperref[VI.prop1.3]{1.3}; the parenthetical case follows from the
case just treated and \hyperref[VI.cor1.4]{1.4}.

\begin{enonce*}{Corollary 1.16}\etiquette[1.16]{VI.cor1.16}
Under the conditions of \hyperref[VI.prop1.10]{1.10} (ii), if the
$Y_i$ are coherent, then $f$ is quasi-compact if and only if the
$X_i=\fprod{X}{Y_i}{Y}$ are quasi-compact objects of $E$.
\end{enonce*}

Indeed, by the criterion \hyperref[VI.prop1.10]{1.10} (ii), one must
express that the $f_i:X_i\to Y_i$ are quasi-compact; by
\hyperref[VI.prop1.14]{1.14} (i), this is equivalent to the $X_i$
being quasi-compact.

\begin{enonce*}{Corollary 1.17}\etiquette[1.17]{VI.cor1.17}
Let $E$ be a topos admitting a generating family formed by
quasi-compact objects, and let $(g_i:Y_i\to Y)_{i\in I}$ be a
covering (\resp finite covering) family in $E$, with the $Y_i$
coherent. The object $Y$ is quasi-separated (\resp coherent) if and
only if, for every $i\in I$, $g_i$ is quasi-compact, or equivalently
if, for every pair of indices $i,j\in I$,
$\fprod{Y_i}{Y_j}{Y}$ is a quasi-compact object of $E$.
\end{enonce*}

The parenthetical case follows immediately from the nonparenthetical
case and \hyperref[VI.prop1.3]{1.3}. Necessity of the first condition
is immediate from the definition; let us prove sufficiency. Let
$f:X\to Y$ be a morphism, with\pageoriginale $X$ quasi-compact; we
prove that $f$ is quasi-compact. Since the $g_i$ are quasi-compact,
the $X_i=\fprod{X}{Y_i}{Y}$ are quasi-compact, which implies by
\hyperref[VI.cor1.16]{1.16} that $f$ is quasi-compact. Finally,
expressing the quasi-compactness of the $g_i$ by the same criterion
\hyperref[VI.cor1.16]{1.16} gives the second criterion of
\hyperref[VI.cor1.17]{1.17}. In particular:

\begin{enonce*}{Corollary 1.17.1}\etiquette[1.17.1]{VI.cor1.17.1}
Let $E$ be as in \hyperref[VI.cor1.17]{1.17}, let $X$ be a coherent
object of $E$, let $\Rr\rightrightarrows X$ be an equivalence relation
in $X$, and let $Y=X/\Rr$. Then the following conditions are
equivalent:
\begin{enumerate}
\renewcommand{\theenumi}{\roman{enumi}}
\renewcommand{\labelenumi}{(\theenumi)}
\item $Y$ is coherent;
\item $f:X\to Y$ is quasi-compact;
\item $\Rr$ is quasi-compact.
\end{enumerate}
\end{enonce*}

\begin{enonce*}{Corollary 1.18}\etiquette[1.18]{VI.cor1.18}
Let $E$ be a topos admitting a generating subcategory $C$ formed by
quasi-compact objects, and let $Y$ be an object of $E$. The object
$Y$ is quasi-separated if and only if every morphism $X\to Y$ with
$X\in\ob C$ is quasi-compact.
\end{enonce*}

Necessity is immediate from the definition, since the objects
$X\in\ob C$ are quasi-compact by hypothesis. To prove sufficiency,
let $Z$ be a quasi-compact object of $E$; we prove that every
morphism $Z\to Y$ is quasi-compact. By \hyperref[VI.prop1.10]{1.10}
(i), it is enough to verify that, for every $X\to Y$ with
$X\in\ob C$, the object $\fprod{Z}{X}{Y}$ is quasi-compact. This
follows from the fact that the $X\to Y$ are quasi-compact by
hypothesis and that $Z$ is quasi-compact.

\begin{enonce*}{Corollary 1.19}\etiquette[1.19]{VI.cor1.19}
Let\pageoriginale $E$ be a topos admitting a generating subcategory
formed by quasi-compact objects, and let
$(g_i:Y_i\to Y)_{i\in I}$ be a covering family in $E$, with the $Y_i$
quasi-separated and the $g_i$ quasi-compact. Then $Y$ is
quasi-separated.
\end{enonce*}

Indeed, let $X$ be a quasi-compact object of $E$; we show that every
morphism $f:X\to Y$ is quasi-compact. By
\hyperref[VI.prop1.10]{1.10} (ii), it is enough to prove that the
induced morphisms $f_i:X_i=\fprod{X}{Y_i}{Y}\to Y_i$ are
quasi-compact. Since $g_i$ is quasi-compact, $X_i$ is quasi-compact
(because $X$ is), and since $Y_i$ is quasi-separated, it follows that
$f_i$ is quasi-compact.

\begin{enonce*}[remark]{Remark 1.20.1}\etiquette[1.20.1]{VI.rem1.20.1}
The reader accustomed to the use of the terms ``quasi-compact'' and
``quasi-separated'' in algebraic geometry will expect certain other
relations between the ``absolute'' notions (concerning objects of the
topos $E$) and the ``relative'' notions (concerning arrows of $E$)
considered here; for example, if $f:X\to Y$ is a morphism in $E$ such
that $X$ is quasi-compact and quasi-separated, then $f$ is
quasi-separated. Such properties hold only under restrictive
finiteness conditions on the topos $E$, stronger than those
considered in \hyperref[VI.prop1.10]{1.10}; we shall study these
conditions in the next section.
\end{enonce*}

\setcounter{subsection}{20}
\setcounter{subsubsection}{1}
\subsubsection{}\etiquette[1.20.2]{VI.subsubsec1.20.2}
It follows again from \hyperref[VI.rem1.5.1]{1.5} that whether an object $X$ of the topos $E$
is quasi-separated depends only on the induced topos $E_{/X}$, and
means that the final object of that topos is quasi-separated, \ie
that in $E_{/X}$ the product of two quasi-compact objects is
quasi-compact.


% END INCLUDED SGA4_English_current_buildable_source_20260906/source/06/drafts/SGA4_VI_LINES_391_677_English.texfrag


% BEGIN INCLUDED SGA4_English_current_buildable_source_20260906/source/06/drafts/SGA4_VI_LINES_678_1095_English.texfrag
\subsection{Case of sites}\etiquette[1.21]{VI.subsec1.21}\pageoriginale

Let $C$ be a $𝒰$-site, $E$ the topos of $𝒰$-sheaves on $C$, $f$ an
arrow of $C$, and $X$ an object of $C$. We say that $f$ is a
\emph{quasi-compact morphism} (\resp quasi-separated morphism) \emph{of the
site} $C$ if $\epsilon(f)$ is a quasi-compact morphism (\resp
quasi-separated morphism) of the topos $E$; likewise, we say that $X$ is a
\emph{quasi-separated object of the site} $C$ if $\epsilon(X)$ is a
quasi-separated object of the topos $E$ (compare \hyperref[VI.prop1.2]{1.2}). When $C$ is a
$𝒰$-topos endowed with its canonical topology, the canonical functor
$\epsilon:C\to E$ is an equivalence of categories, and consequently the
preceding definition is compatible with the earlier \hyperref[VI.def1.7]{definitions 1.7} and
\hyperref[VI.def1.13]{1.13}. Note also that the definition under
consideration does not change when the universe $𝒰$ is replaced by another
universe $𝒱$ such that $C$ is also a $𝒱$-site, at least when $C$ admits a
topologically generating family consisting of quasi-compact objects. To see
this, we are immediately reduced to the case $𝒰\subset 𝒱$, so that $E$ is
a $𝒱$-site, and to proving that, for every arrow $f$ of the $𝒰$-topos $E$,
$f$ is quasi-compact (\resp quasi-separated) if and only if it defines a
quasi-compact arrow in the topos $E'$ of $𝒱$-sheaves on $E$, and likewise
that, for every object $X$ of $E$, $X$ is quasi-separated if and only if the
corresponding object of $E'$ is quasi-separated. We already know (without
any hypothesis on $E$) the analogous statement for a quasi-compact object
(\hyperref[VI.prop1.2]{1.2}). The case ``$f$ quasi-compact'' follows from
this, by the criterion \hyperref[VI.prop1.10]{1.10} (i), applied successively
in $E$ and\pageoriginale in $E'$ to the same generating family
$(X_i)_{i\in I}$ of $E$ consisting of quasi-compact objects. This also gives
the case ``$f$ quasi-separated'', which means that the corresponding
diagonal morphism is quasi-compact. Finally, in the case ``$X$
quasi-separated'', we are again reduced to the case ``$X$ quasi-compact''
by the criterion \hyperref[VI.cor1.18]{1.18}.

\subsection{}\etiquette[1.22]{VI.subsec1.22}
{\textit{\textbf{Examples.}}} Return to the example
\hyperref[VI.subsubsec1.6.2]{1.6.2} of the category $(\Sch)$ with a topology
$T_i$. This is a $𝒱$-site for a suitable $𝒱$, and it admits the subcategory
of affine schemes as a topologically generating category consisting of
quasi-compact objects. We leave it to the reader to verify that an arrow
$f:X\to Y$ of this site is quasi-compact (\resp quasi-separated) if and only
if it is a quasi-compact (\resp quasi-separated) morphism of schemes in the
usual sense (\cite{sga4.exp.VI.bib.3} or \cite{sga4.exp.VI.bib.4}); likewise, an object of the site is
quasi-compact (\resp quasi-separated) if and only if it is a quasi-compact
(\resp quasi-separated) scheme in the usual sense of \loccit. It is in fact
appropriate to say, as we do here, \emph{coherent morphism of schemes,
coherent scheme}, instead of ``quasi-compact and quasi-separated morphism of
schemes'', ``quasi-compact and quasi-separated scheme'', as is also done in
\EGA I, $2^{\text{nd}}$ edition.

\setcounter{subsection}{22}
\subsubsection{}\etiquette[1.22.1]{VI.subsubsec1.22.1}

For a given scheme $X\in 𝒰$, one may also consider the topos $\Top(X)$
associated with the underlying topological space, or the topos
$\Top(X_{\et})$ (\resp $\Top(X_{\text{fppf}})$) associated with the site of
étale $X$-schemes (\resp $X$-schemes locally\pageoriginale of finite
presentation) that are elements of $𝒰$, endowed with the étale topology
(\resp fppf topology). It is a $𝒰$-topos admitting a generating family
consisting of coherent objects, corresponding to the affine objects of the
site of definition.

With this understood, the final object of this topos is quasi-compact
(\resp quasi-separated, \resp coherent) if and only if the scheme $X$ is
quasi-compact (\resp quasi-separated, \resp coherent). Likewise, a morphism
of the site of étalé spaces over $X$ (\resp of the site $X_{\text{et}}$ of
étale schemes over $X$, \resp of the site $X_{\text{fppf}}$ of schemes
locally of finite presentation over $X$) is quasi-compact (\resp
quasi-separated, \resp coherent) if and only if it is a morphism of schemes
that is, in the usual sense of \loccit, quasi-compact (\resp quasi-separated,
\resp coherent).

\begin{enonce*}{Theorem 1.23}\etiquette[1.23]{VI.thm1.23}
Let $E$ be a $𝒰$-topos and $X$ an object of $E$.
\begin{enumerate}
\renewcommand{\theenumi}{\roman{enumi}}
\renewcommand{\labelenumi}{(\theenumi)}
\item For $X$ to be quasi-compact, it is necessary and sufficient that, for
  every filtered inductive system $(Y_i)_{i\in I}$ of $E$, with $I\in 𝒰$,
  the natural map
\begin{equation*}
\varinjlim  \Hom(X, Y_i) \to \Hom(X, \varinjlim Y_i)\tag{1.23.1}
\end{equation*}\etiquette[1.23.1]{VI.eq1.23.1}
be injective.

\item Suppose that $E$ admits a generating subcategory $C$ consisting of
  quasi-compact objects. A necessary condition for $X$ to be coherent is
  that, for every datum as in \emph{(i)}, the map \emph{\hyperref[VI.eq1.23.1]{(1.23.1)}} be
  bijective.
\end{enumerate}
\end{enonce*}

\begin{enumerate}
\renewcommand{\theenumi}{\roman{enumi}}
\renewcommand{\labelenumi}{(\theenumi)}
\item \underline{Necessity}.
Suppose\pageoriginale that $X$ is quasi-compact. We must prove that if
  $i\in I$ and $f_i,g_i:X\rightrightarrows Y_i$ are such that the composites
  $f,g$ of $f_i,g_i$ with $Y_i\to Y=\varinjlim Y_j$ are equal, then there
  exists $j\geq i$ such that the composites $f_j,g_j$ with $Y_i\to Y_j$ are
  already equal. Since filtered inductive limits in $E$ commute with finite
  projective limits, and in particular with kernels, we have
$$
\Ker (f, g) = \mathop{\lim}_{j\geq l} \Ker (f_j, g_j){\quoi}.
$$
By hypothesis $\Ker(f,g)=X$, so $X$ is the limit of the filtered increasing
family of its subobjects $\Ker(f_j,g_j)$. Since $X$ is quasi-compact, it is
equal to one of these subobjects; hence there exists $j\geq i$ such that
$f_j=g_j$.

\underline{Sufficiency}. To prove that $X$ is quasi-compact, it is equivalent
(\hyperref[VI.subsubsec1.5.2]{1.5.2}) to prove that, for every filtered
increasing family $(X_i)_{i\in I}$ of subobjects of $X$ whose $\varinjlim$ is
$X$, there exists $i\in I$ such that $X_i=X$. Since every monomorphism in
$E$ is effective (II \hyperref[II.sec4]{4.0}), it follows that
$$
X_i = \Ker (f_i, g_i : X \rightrightarrows X \amalg_{X_{i}} X){\quoi}.
$$
Let $Y_i=X\amalg_{X_i}X$, and consider the inductive system formed by the
$Y_i$. If $Y$ is its inductive limit, then
$Y\simeq X\amalg_{\varinjlim X_i}X=X\amalg_X X=X$.
Consequently, for a fixed $i\in I$, $f_i,g_i$ give the same element of
$\Hom(X,Y)$, so by hypothesis there exists $j\geq i$ such that $f_j=g_j$,
\ie $X_j=X$, \cqfd

\item It remains\pageoriginale to prove that, under the stated conditions,
  every morphism $X\to Y$ comes from a morphism $X\to Y_i$ for a suitable
  $i\in I$. Since $C$ generates $E$ and the $Y_i$ cover $Y$, we may cover
  $X$ by objects $X_{\alpha}$ of $C$ in such a way that each composite
  $X_{\alpha}\to X\to Y$ factors through one of the $Y_i$. Since $X$ is
  quasi-compact, we may suppose that the family of the $X_{\alpha}$ is
  finite, and hence that the morphisms $X_{\alpha}\to Y$ factor through a
  fixed $Y_i$ as morphisms $g_{\alpha i}:X_{\alpha}\to Y_i$. Since $X$ is
  quasi-separated, the $\fprod{X_{\alpha}}{X_{\beta}}{X}$ are quasi-compact.
  Moreover, for every $(\alpha,\beta)$, the composites of $\pr_1$ and
  $\pr_2$ on $\fprod{X_{\alpha}}{X_{\beta}}{X}$ with
  $X_{\alpha},X_{\beta}\to Y_i$ are
\[
\xymatrix{
\fprod{X_{\alpha}}{X_{\beta}}{X} \ar[rd] \ar[dd]& &\\
                            &X_{\alpha} \ar[ld] \ar[rrrr]&   & & & Y_i
\ar[lld]\\
X    \ar[rrr]                       &          & & Y  &
}
\]
such that their composites with $Y_i\to Y$ are equal. By (i), there
  therefore exists $j\geq i$ such that the composites with $Y_i\to Y_j$ are
  already equal. Since the set of pairs $(\alpha,\beta)$ is finite, $j$ may
  be chosen independently of the pair. Consequently, the morphisms
  $g_{\alpha j}$ arise by composition from a morphism $g_j:X\to Y_j$, whose
  composite with $Y_j\to Y$ is moreover the given morphism $X\to Y$, since
  this is so after composition with the covering maps $X_{\alpha}\to X$.
  This completes the proof.
\end{enumerate}

\begin{enonce*}[remark]{Remarks 1.23.2}\etiquette[1.23.2]{VI.rem1.23.2}
The hypothesis on $E$ in (ii) may be replaced by an additional hypothesis
\emph{on} $X$, namely that $X$ admits a fundamental system of covering
families by quasi-compact objects.
\end{enonce*}

\begin{enonce*}{Corollary 1.24}\etiquette[1.24]{VI.cor1.24}
Let\pageoriginale $E$ be a topos, and let $C$ be a strictly full subcategory
of $E$ consisting of coherent objects. Using the fact that the
$𝒰$-inductive limits in $E$ are representable, one obtains \emph{(I \hyperref[I.sec8.6.1]{8.6.1})}
a canonical functor
\begin{equation*}
\Ind (C) \to E{\quoi}.\tag{1.24.1}
\end{equation*}\etiquette[1.24.1]{VI.eq1.24.1}
This functor is fully faithful. If $E$ admits a generating subcategory
consisting of coherent objects, and if $C$ is stable under direct factors,
then the following conditions are equivalent:
\begin{enumerate}
\renewcommand{\theenumi}{\roman{enumi}}
\renewcommand{\labelenumi}{(\theenumi)}
\item The functor \emph{\hyperref[VI.eq1.24.1]{(1.24.1)}} is essentially surjective, \ie it is an
  equivalence of categories.

\item $C$ is a generating subcategory of $E$, stable under finite inductive
  limits.

\item The category $C$ is equal to the full subcategory $E_{\coh}$ of $E$
  consisting of all coherent objects of $E$, and the necessary condition
  \emph{\hyperref[VI.thm1.23]{1.23} (ii)} for an object of $E$ to be coherent
  is also sufficient.
\end{enumerate}
\end{enonce*}

Let $E_{PF}$ be the full subcategory of $E$ defined in I \hyperref[I.cor8.7.6]{8.7.6}. By I \hyperref[I.prop8.7.5]{8.7.5}
a), statement \hyperref[VI.thm1.23]{1.23} (ii) is equivalent to the second
inclusion in
$$
C \subset E_{\coh} \subset E_{PF}{\quoi},
$$
and the composite inclusion $C\supset E_{PF}$ means that \hyperref[VI.eq1.24.1]{(1.24.1)} is fully
faithful. Since\pageoriginale, by hypothesis, the subcategory $E_{\coh}$ of
$E$ is generating, the same is a fortiori true of $E_{PF}$; moreover, finite
projective limits are representable in $E$, so that we may apply I \hyperref[VI.cor8.7.7]{8.7.7}.
Condition (iii) of \hyperref[VI.cor1.24]{1.24} may also be written
$C=E_{\coh}=E_{PF}$ and is equivalent to the condition $C=E_{PF}$, which is
precisely condition (ii bis) of \loccit (taking into account that, by
hypothesis, $E_{\coh}$ and hence $E_{PF}$ is a generating subcategory of
$E$). Likewise, condition (ii) of \hyperref[VI.cor1.24]{1.24} is precisely
condition (iii bis) of \loccit. Thus the equivalence of conditions (i),
(ii bis), and (iii bis) in \loccit gives the equivalence of conditions (i),
(ii), and (iii) of \hyperref[VI.cor1.24]{1.24}, \cqfd

The proof (or, more precisely, the invocation of I \hyperref[I.cor8.7.7]{8.7.7} (ii bis)) also gives
part a) of the

\begin{enonce*}{Corollary 1.24.2}\etiquette[1.24.2]{VI.cor1.24.2}
Let $E$ be a topos admitting a generating subfamily consisting of coherent
objects, and let $E_{PF}$ be the strictly full subcategory of objects $X$ of
$E$ such that the covariant functor $\Hom(X,-)$ represented by $X$ commutes
with filtered inductive limits. Then:
\begin{enumerate}
\renewcommand{\theenumi}{\alph{enumi}}
\renewcommand{\labelenumi}{\theenumi)}
\item The natural functor
$$
\Ind (E_{PF}) \to E
$$
is an equivalence of categories.

\item For an object $X$ of $E$ to belong to $E_{PF}$, it is necessary and
  sufficient that it be isomorphic to a direct factor ($=$ image of a
  projector (\hyperref[I.sec10.6]{I, 10.6})) of the cokernel\pageoriginale of a double arrow
  $X_1\rightrightarrows X_0$, with $X_1$ and $X_0$ coherent.
\item The subcategory $E_{PF}$ of $E$ is stable under finite inductive
  limits, and is equivalent to a small category.
\end{enumerate}
\end{enonce*}

It remains to prove b) and c). The first assertion of c) is immediate from
I \hyperref[I.prop2.8]{2.8}, and implies the ``if'' part of b) by \hyperref[VI.thm1.23]{1.23} (ii).
For the ``only if'' part, use \hyperref[VI.thm1.23]{1.23} (i), which shows
that $X$ is quasi-compact, whence the existence of an epimorphism
$X_0\to X$, with $X_0$ coherent, in view of the fact that $E_{\coh}$
generates $E$ and is stable under finite sums. Let
$R=\fprod{X_0}{X_0}{X}$, so that $X$ identifies with the cokernel of
$R\rightrightarrows X_0$. By a), we have $R=\varinjlim R_i$, the filtered
inductive limit of the coker $(R_i\rightrightarrows X_0)=X_i$. By the
hypothesis $X\in\Ob E_{PF}$, for sufficiently large $i$ the morphism
$u_i:X'_i\to X$ admits a section $v:X\to X'_i$. Let
$w=vu_i:X'_i\to X'_i$, so that $w^2=w$ and
$X\simeq\Coker(w,\id_{X'_i}:X'_i\rightrightarrows X'_i)$. Since $R_i$ is
again quasi-compact, cover it by $X_1\to R_i$, with $X_i$ coherent, whence
$X_i\simeq\Coker(X_1\rightrightarrows X_0)$, which proves b).

Finally, since $E_{PF}\subset E_{\text{qu.cpct}}$, the last assertion of c)
follows from \hyperref[VI.subsubsec1.5.3]{1.5.3}.

\begin{enonce*}{Corollary 1.25}\etiquette[1.25]{VI.cor1.25}
Let $E$ be a topos such that the full subcategory $C$ of $E$ consisting of
coherent objects is generating. Consider the canonical functor
$$
\varphi : \Ind(c) \to E{\quoi}.
$$
The following\pageoriginale conditions are equivalent:
\begin{itemize}
\item[(i)] The functor $\varphi$ is essentially surjective, \ie every object
  of $E$ is a filtered inductive limit of coherent objects.

\item[(i bis)] The functor $\varphi$ is an equivalence of categories.

\item[(ii)] Every finite inductive limit in $E$ of coherent objects is a
  coherent object.

\item[(ii bis)] The cokernel in $E$ of a double arrow of coherent objects is
  a coherent object.

\item[(iii)] The converse of \emph{\hyperref[VI.thm1.23]{1.23} (ii)} holds.
\end{itemize}
\end{enonce*}

This is a special case of \hyperref[VI.cor1.24]{1.24}, taking into account
that $C$ is stable under finite sums (\hyperref[VI.cor1.4]{1.4}), which
implies the equivalence of (ii) and (ii bis), and taking into account the

\begin{enonce*}{Lemma 1.25.1}\etiquette[1.25.1]{VI.lem1.25.1}
Let $E$ be a topos. Every direct factor $X'$ of a coherent object $X$ of $E$
is coherent.
\end{enonce*}

Indeed, $X'$ is quasi-compact as a quotient of $X$
(\hyperref[VI.prop1.3]{1.3}), and quasi-separated as a subobject of $X$
(\hyperref[VI.cor1.15]{1.15}).

\begin{enonce*}{Corollary 1.26}\etiquette[1.26]{VI.cor1.26}
Let $E$ be a topos satisfying the equivalent conditions
\emph{\hyperref[VI.cor1.25]{1.25}.} Then every induced topos does likewise.
\end{enonce*}

This follows immediately from criterion (ii), taking into account
\hyperref[VI.subsubsec1.20.2]{1.20.2} and the fact that the existence of a
generating subcategory of $E$ consisting of coherent objects plainly
implies the same property for an induced topos.

\begin{enonce*}[remark]{Remark 1.27}\etiquette[1.27]{VI.rem1.27}
Among\pageoriginale the important examples of topoi satisfying the
conditions of \hyperref[VI.cor1.25]{1.25}, let us mention Noetherian topoi
(\hyperref[VI.subsec2.14]{2.14}), and the Zariski or étale topos
(\hyperrefext[VII][VII.sec1]{1}) of a scheme $X$ (\cf IX, note p.42). But
even when the topos $E$ is \emph{coherent} (\hyperref[VI.def2.3]{2.3}) (\ie
the full subcategory $C$ of coherent objects is generating and stable under
finite projective limits), it is not always true that $C$ is \emph{perfect}
(\hyperref[VI.def2.9.1]{2.9.1}), \ie satisfies the equivalent conditions of
\hyperref[VI.cor1.25]{1.25}; \cf \hyperref[VI.exer1.28]{1.28} below.
\end{enonce*}

\begin{enonce*}[remark]{Exercise 1.28}\etiquette[1.28]{VI.exer1.28}
Let $E$ be a topos and $p_1,p_2:X_1\rightrightarrows X_0$ a double arrow in
$E$. Define an increasing sequence of subobjects $R_n$ $(n\geq0)$ of
$X_0\times X_0$ by the condition
$R_0=\Sup(\text{Im}(p_1,p_2),\text{Im}(p_2,p_1))$, and
$R_n=\pr_{13}(\pr_{12}^{-1}(R_{n-1}),\pr_{23}^{-1}(R_{n-1}))$ for every
$n\geq1$, where $\pr_{ij}$ $(1\leq i<j\leq3)$ denote the evident
projections from the triple product of $X_0$ to the double product.
\begin{enumerate}
\renewcommand{\theenumi}{\alph{enumi}}
\renewcommand{\labelenumi}{\theenumi)}
\item Show that one thus obtains an increasing sequence of subobjects of
  $X_0\times X_0$, that $R=\varinjlim R_n$ is an equivalence graph (I \hyperref[I.sec10.4]{10.4})
  in $X_0$, and that the canonical morphism
$$
X = \Coker (p_1, p_2) \to X_0 /R
$$
is an isomorphism.

\item Deduce that if $X_0$ is quasi-compact and $X=\Coker(p_1,p_2)$ is
  quasi-separated (hence coherent), then the sequence $(R_n)$ is stationary.
  Conversely,\pageoriginale if this sequence is stationary, and if $X_0$ is
  assumed coherent and $X_1$ quasi-compact, then $X$ is coherent. (For this
  last point, write
  $R_n\simeq\fprod{(R_{n-1},\pr_2)}{(R_{n-1},\pr_1)}{X_0}$ and deduce that
  $R_n$ is quasi-compact over $X_0$ for every $n$, then use \hyperref[VI.cor1.19]{1.19}.)

\item Construct an example of two maps of sets
  $p_1,p_2:X_1\rightrightarrows X_0$ such that the sequence $(R_n)$ is not
  stationary.

\item Let $C$ be a small category in which finite inductive limits and
  finite projective limits are representable, and in which every morphism
  factors as an effective epimorphism followed by a monomorphism. Repeat the
  construction of the $R_n$ for a double arrow $(p_1,p_2)$ of $C$. Endowing
  $C$ with the canonical topology, show that the canonical functor
  $\epsilon:C\to C^{\sim}$ is exact, and consequently commutes with the
  formation of the $R_n$.

\item Take for $C$ a small full subcategory of $(\Ens)$, stable up to
  isomorphism under finite projective limits and finite inductive limits, and
  containing the sets $X_0,X_1$ of c). Deduce that the topos $C^{\sim}$
  (which is a \emph{coherent} topos (\hyperref[VI.def2.3]{2.3}), \ie is
  generated by the subcategory of its coherent objects, which is stable
  under finite projective limits) does not satisfy the equivalent conditions
  of \hyperref[VI.cor1.25]{1.25}.

\item Let $k$ be a field, and let $C$ be the fppf site of schemes locally of
  finite presentation over $k$ (\SGA 3 IV 6.3). Show that there exists a
  reduced affine scheme $X$ over $k$ admitting an automorphism $f$ that is
  not of finite order (take, for example, the affine plane $E^2$ endowed with
  the automorphism $f(x)=x$, $f(y)=x+y$).\pageoriginale Show that if
  $p_1,p_2:X_1\rightrightarrows X_0$ is a double arrow in $C$ such that the
  corresponding double arrow in $(\Ens) p_1(\bar{k}):p_2(\bar{k}):
  X_1(\bar{k})\rightrightarrows X_0(\bar{k})$ gives a nonstationary sequence
  $(R_n)$, then the same is true of the double arrow of $C^{\sim}$ defined by
  $p_1,p_2$. Thus, if $X_0$ is of finite type over $k$, the cokernel in
  $C^{\sim}$ of $(p_1,p_2)$ is not coherent. Deduce that the cokernel in
  $C^{\sim}$ of the pair $(f,\id_X)$ is not coherent, and hence that the
  topos $C^{\sim}$ does not satisfy the equivalent conditions of
  \hyperref[VI.cor1.25]{1.25}. Deduce more generally that if $S$ is a
  nonempty scheme and $C$ is the fppf site of schemes locally of finite
  presentation over $S$, then the topos $C^{\sim}$ does not satisfy the
  conditions of \hyperref[VI.cor1.25]{1.25}. (Care must be taken that,
  contrary to appearances, $C^{\sim}$ is therefore Noetherian
  (\hyperref[VI.def2.11]{2.11}) only if $S$ is empty, as follows from the
  preceding discussion and \hyperref[VI.subsec2.14]{2.14}.)
\end{enumerate}
\end{enonce*}

\begin{enonce*}{Definition 1.30}\etiquette[1.30]{VI.def1.30}
Let $E$ be a topos. An object $X$ of $E$ is called a \emph{pre-Noetherian
object}\footnote{The stronger notion of a \emph{Noetherian} object is given
  in \hyperref[VI.def2.11]{2.11} below.} of the topos $E$ if it satisfies the
following two equivalent conditions:
\begin{enumerate}
\renewcommand{\theenumi}{\roman{enumi}}
\renewcommand{\labelenumi}{(\theenumi)}

\item Every subobject of $X$ is quasi-compact.

\item Every increasing sequence of subobjects of $X$ is stationary.
\end{enumerate}
\end{enonce*}

\setcounter{subsection}{30}
\setcounter{subsubsection}{0}
\subsubsection{}\etiquette[1.30.1]{VI.subsubsec1.30.1}
The equivalence of conditions (i) and (ii) is clear. Note that these
conditions again depend only on the induced topos $E_{/X}$, and indeed only
on the ordered set of opens of $E_{/X}$, which is isomorphic to the ordered
set of subobjects of $X$. They are stable under passage to a subobject of
$X$.



% END INCLUDED SGA4_English_current_buildable_source_20260906/source/06/drafts/SGA4_VI_LINES_678_1095_English.texfrag


% BEGIN INCLUDED SGA4_English_current_buildable_source_20260906/source/06/drafts/SGA4_VI_LINES_1096_1382_English.texfrag
\subsection{}\etiquette[1.31]{VI.subsec1.31}
One proves,\pageoriginale as in \hyperref[VI.prop1.3]{1.3} and
\hyperref[VI.cor1.4]{1.4}, the following facts: if $(X_i\to X)$ is a
finite covering family, with the $X_i$ pre-Noetherian, then $X$ is
pre-Noetherian; in particular, a quotient of a pre-Noetherian object
of $E$ is pre-Noetherian. If $(X_i)_{i\in I}$ is a family of objects
of $E$, then their coproduct $X$ is a pre-Noetherian object of $E$ if
and only if all but finitely many of the $X_i$ are isomorphic to
$∅_E$, and all the $X_i$ are pre-Noetherian.

\subsection{}\etiquette[1.32]{VI.subsec1.32}
Clearly, a pre-Noetherian object of a topos $E$ is quasi-compact. When
$E$ admits a generating family formed by pre-Noetherian objects, the
converse holds (as follows immediately from
\hyperref[VI.subsec1.31]{1.31}): the pre-Noetherian objects of $E$ are
then precisely its quasi-compact objects.

Suppose that $E$ admits a generating family $(X_i)_{i\in I}$ formed
by quasi-compact objects. Then the following three conditions are
equivalent:
\begin{enumerate}
\renewcommand{\theenumi}{\roman{enumi}}
\renewcommand{\labelenumi}{(\theenumi)}
\item The $X_i$ are pre-Noetherian.

\item Every quasi-compact object of $E$ is pre-Noetherian.

\item Every monomorphism in $E$ is quasi-compact.
\end{enumerate}
\noindent
(We have just proved the equivalence of (i) and (ii); (ii)
$\Rightarrow$ (iii) is immediate from the definitions, and (iii)
$\Rightarrow$ (ii) follows from Definition
\hyperref[VI.def1.30]{1.30} (i).)

Note\pageoriginale that (iii) implies that \emph{every morphism of
$E$ is quasi-separated}; hence (\hyperref[VI.prop1.14]{1.14} (ii))
every object of $E$ lying over a quasi-separated object of $E$ is
itself quasi-separated.

\begin{enonce*}[remark]{Example 1.33}\etiquette[1.33]{VI.exem1.33}
(\emph{Classifying topos of a group}). Let $G$ be a group
$\in 𝒰$, and let $B_G$ be its classifying topos (IV \hyperref[IV.subsec2.4]{2.4}), \ie the
category of sets $\in 𝒰$ on which $G$ acts on the left. The nonempty
connected objects of $B_G$ (IV \hyperref[IV.subsubsec4.3.5]{4.3.5}) are the nonempty $G$-sets on
which $G$ acts transitively, and every object $E$ of $B_G$ is the
coproduct of its connected components, which are the minimal nonempty
subobjects of $E$ (namely, the orbits of $G$ in $E$). Consequently,
$E$ is a quasi-compact object of $B_G$ if and only if the set
$\pi_0(E)$ of orbits of $G$ in $E$ is finite, and in that case $E$ is
even a pre-Noetherian object of $B_G$. These objects (and even the
single object $G_s$, which is nonempty and connected) form a
generating subcategory of $B_G$. The final object $e$ of $B_G$ is
quasi-separated, \ie the product of two quasi-compact objects of
$B_G$ is quasi-compact, if and only if the product $G_s\times G_s$ is
quasi-compact, which is plainly equivalent to $G$ being finite. More
generally, one concludes that an object $E$ of $B_G$ is
quasi-separated if and only if the stabilizers of its points are
\emph{finite} subgroups of $G$, \ie if it is isomorphic to a coproduct
of homogeneous spaces $G/H$, with $H$ \emph{finite}. Indeed, by
\hyperref[VI.cor1.15]{1.15}, one may assume in this criterion that
$E$ is connected; but if $x\in E$ has stabilizer group $G_x$, we know
(IV \hyperref[IV.subsec5.8]{5.8})\pageoriginale that $B_G/E$ is equivalent to the topos
$B_{G_x}$, and hence its final object is quasi-separated if and only
if $G_x$ is finite. This criterion shows that every object of $B_G$
lying over a quasi-separated object is quasi-separated; a fortiori,
the product of two quasi-separated objects of $B_G$ is
quasi-separated. It also shows that the coherent objects of $B_G$ are
the objects isomorphic to finite coproducts of homogeneous spaces
$G/H$, with $H$ \emph{finite}. Since $G_s$ is coherent, the coherent
objects of $B_G$ already form a generating subcategory of $B_G$.
Moreover, this subcategory is stable under fiber products (as follows
from the fact that every object lying over a quasi-separated object
is quasi-separated). (In the terminology of
\hyperref[VI.def2.11]{2.11} below, $B_G$ is a \emph{locally
Noetherian} topos, but it is quasi-separated only when $G$ is finite,
in which case $B_G$ is even Noetherian.)

It is immediate that a morphism $f:E'\to E$ is quasi-compact if and
only if the induced map $\pi_0(f):\pi_0(E')\to\pi_0(E)$ on the sets
of orbits is proper, \ie has finite fibers. We recover the fact that
a monomorphism is quasi-compact (\hyperref[VI.subsec1.32]{1.32}), and
hence that every morphism is quasi-separated. Thus the coherent
morphisms in $B_G$ are simply the quasi-compact morphisms, which have
just been characterized.

Let $E$ be an object of the topos $B_G$. One verifies easily that $E$
is a filtered inductive limit of coherent objects of $B_G$ if and
only if the stabilizers of the points of $E$ are ind-finite groups
(\ie filtered\pageoriginale inductive limits of their finite
subgroups). In particular, if $G$ is not ind-finite (for example, if
$G=\ZZ$), the final object of the topos $B_G$ is not an inductive
limit of coherent objects. More precisely, the topos $B_G$ satisfies
the equivalent conditions of \hyperref[VI.cor1.25]{1.25} if and only
if its final object is a filtered inductive limit of coherent objects,
or equivalently if the group $G$ is ind-finite.
\end{enonce*}


\section{Finiteness conditions for a topos}\etiquette[2]{VI.sec2}

\begin{enonce*}{Proposition 2.1}\etiquette[2.1]{VI.prop2.1}
Let $E$ be a $𝒰$-topos. The following conditions are equivalent:
\begin{enumerate}
\renewcommand{\theenumi}{\roman{enumi}}
\renewcommand{\labelenumi}{(\theenumi)}
\item There exists a full generating subcategory $C$ of $E$, formed
  by quasi-compact objects and stable under fiber products.

\item There exists a $𝒰$-site $C$ such that every object of $C$ is
  quasi-compact, fiber products are representable in $C$, and $E$ is
  equivalent to the category of sheaves $C^{\sim}$.
\end{enumerate}

Suppose that these conditions hold. Then in \emph{(i) (\resp (ii))}
the category $C$ may even be chosen to be $𝒰$-small. Moreover, for
every $X\in\ob C$, the object $X$ \emph{(\resp $\epsilon(X)$)} of the
topos $E$ is \emph{coherent}.
\end{enonce*}

The implication (i) $\Rightarrow$ (ii) follows from IV \hyperref[IV.cor1.2.1]{1.2.1}, and the
fact that in (i) (\resp (ii)) one may take $C$ to be $𝒰$-small
follows immediately from the argument of IV \hyperref[IV.subsubsec1.2.3]{1.2.3}, applied to a small
full generating subcategory in the topological sense (II \hyperref[II.def3.0.1]{3.0.1}) of
$C$. It remains to prove (ii) $\Rightarrow$ (i) and the last assertion
of \hyperref[VI.prop2.1]{2.1}; these follow immediately from the
following result.

\begin{enonce*}{Corollary 2.1.1}\etiquette[2.1.1]{VI.cor2.1.1}
Let\pageoriginale $C$ be a $𝒰$-site in which fiber products are
representable and every object is quasi-compact, let $E=C^{\sim}$,
and let $\varepsilon:C\to E$ be the canonical functor. Then, for every
$X\in\ob C$, $\varepsilon(X)$ is a coherent object of $E$. The full
subcategory of $E$ formed by the coherent objects of $E$ lying over
some $\varepsilon(X)$ is stable under fiber products (and is plainly
generating in $E$).
\end{enonce*}

We already know that $\epsilon(X)$ is quasi-compact
(\hyperref[VI.prop1.2]{1.2}); let us prove that it is quasi-separated.
Let $F,G$ be two quasi-compact objects of $E$, and let
$F\to\varepsilon(X)$, $G\to\epsilon(X)$ be two morphisms. We must
prove that the fiber product $\fprod{F}{G}{\epsilon(X)}$ is
quasi-compact. Covering $F$ and $G$ by finitely many objects of the
form $\epsilon(Y_i)$ and, respectively, $\epsilon(Z_j)$, and
proceeding as in \hyperref[VI.prop1.2]{1.2} to reduce to the case in
which the composites $\epsilon(Y_i)\to\epsilon(X)$ and
$\epsilon(Z_j)\to\epsilon(X)$ come from morphisms $Y_i\to X$ and,
respectively, $Z_j\to X$, we are reduced by
\hyperref[VI.prop1.3]{1.3} to the case in which the morphisms
$F\to\epsilon(X)$ and $G\to\epsilon(X)$ are obtained by applying
$\epsilon$ to morphisms $Y\to X$ and $Z\to X$. But since fiber
products are representable in $C$ and $\epsilon$ commutes with them,
we are reduced to proving that $\epsilon(\fprod{Y}{Z}{X})$ is
quasi-compact, which was already noted above for every object of the
form $\epsilon(U)$.

Now let $G\to F$, $H\to F$ be morphisms between coherent objects of
$E$, and suppose that there exists a morphism $F\to\epsilon(X)$ for
some $X\in\ob C$. We prove that $\fprod{G}{H}{F}$ is coherent. Since
$G$ and $H$ are quasi-compact and $F$ is quasi-separated, it already
follows from the definitions that the fiber product is quasi-compact.
It remains\pageoriginale to see that it is quasi-separated, or,
equivalently (\hyperref[VI.cor1.15]{1.15}), that
$\fprod{G}{H}{\epsilon(X)}$ is quasi-separated. This follows from the
more general fact:

\begin{enonce*}{Corollary 2.1.2}\etiquette[2.1.2]{VI.cor2.1.2}
Under the conditions of \hyperref[VI.cor2.1.1]{2.1.1}, for every
quasi-separated object $G$ of $E$, every morphism
$G\to\epsilon(X)$ is quasi-separated. If $H\to\epsilon(X)$ is a
second morphism, with $H$ quasi-separated, then
$\fprod{G}{H}{\epsilon(X)}$ is quasi-separated.
\end{enonce*}

The first assertion implies the second: by base change, it implies
that $\fprod{G}{H}{\epsilon(X)}\to H$ is quasi-separated, and since
$H$ is quasi-separated, $\fprod{G}{H}{\epsilon(X)}$ is likewise
quasi-separated (\hyperref[VI.prop1.14]{1.14} (ii)). To prove that
$G\to\epsilon(X)$ is quasi-separated, we shall use the following
lemma.

\begin{enonce*}{Lemma 2.1.3}\etiquette[2.1.3]{VI.lem2.1.3}
Let $E$ be a topos and $f:X\to Y$ a morphism in $E$. Suppose that $E$
admits a generating family of quasi-compact objects, and that there
exists a covering family $(g_i:X_i\to X)_{i\in I}$ of $X$ such that
the $X_i$ are quasi-compact and, for every pair of indices
$i,j\in I$, $\fprod{X_i}{X_j}{Y}$ is quasi-separated. Then, if $X$
is a quasi-separated object of $E$, $f$ is a quasi-separated
morphism.
\end{enonce*}

Indeed, the family
$(\fprod{g_i}{g_j}{Y}:\fprod{X_i}{X_j}{Y}\to
\fprod{X}{X}{Y})_{(i,j)\in I\times I}$ is covering. Therefore, to see
that the morphism $\Delta_{X/Y}:X\to\fprod{X}{X}{Y}$ is
quasi-compact, it is enough to see that, for every
$(i,j)\in I\times I$, the morphism obtained from it\pageoriginale by
base change along $\fprod{g_i}{g_j}{Y}$ is quasi-compact
(\hyperref[VI.prop1.10]{1.10} (ii)). The latter is precisely the
canonical morphism
$\fprod{X_i}{X_j}{X}\to\fprod{X_i}{X_j}{Y}$. By hypothesis $X$ is
quasi-separated and the $X_i$ are quasi-compact, so
$\fprod{X_i}{X_j}{X}$ is quasi-compact; and since
$\fprod{X_i}{X_j}{Y}$ is quasi-separated by hypothesis, the morphism
$\fprod{X_i}{X_j}{X}\to\fprod{X_i}{X_j}{Y}$ is indeed quasi-compact,
\cqfd

We can now complete the proof of \hyperref[VI.cor2.1.2]{2.1.2} by
proving that every morphism $f:G\to\epsilon(X)$, with $G$
quasi-separated, is quasi-separated. To do this, cover $G$ by objects
$\epsilon(X_i)$ such that the composite morphisms
$\epsilon(X_i)\to\epsilon(X)$ come from morphisms $X_i\to X$. By
\hyperref[VI.lem2.1.3]{2.1.3}, since the $\epsilon(X_i)$ are
quasi-compact, it is enough to verify that the fiber products
$\fprod{\epsilon(X_i)}{\epsilon(X_j)}{\epsilon(X)}=
\epsilon(\fprod{X_i}{X_j}{X})$ are quasi-separated. We have already
seen that every object of $E$ of the form $\epsilon(U)$ is coherent.
This proves \hyperref[VI.cor2.1.2]{2.1.2} and completes the proof of
\hyperref[VI.prop2.1]{2.1}.

\begin{enonce*}{Proposition 2.2}\etiquette[2.2]{VI.prop2.2}
Let $E$ be a topos containing a full generating subcategory $C$
formed by coherent objects. The following conditions on $E$ are
equivalent:
\begin{itemize}
\item[(i)] Every quasi-separated object of $E$ is quasi-separated over
  the final object.

\item[(i bis)] For every morphism $f:X\to Y$ in $E$, if $X$ is
  quasi-separated, then $f$ is quasi-separated.

\item[(i ter)] Every object $X$ of $C$ is quasi-separated over the
  final object of $E$, \ie the diagonal morphism
  $X\to X\times X$ is quasi-compact.

\item[(ii)] The product\pageoriginale of two quasi-separated objects
  of $E$ is quasi-separated.

\item[(ii bis)] The product in $E$ of two objects of $C$ is
  quasi-separated.
\end{itemize}

These conditions imply that the full subcategory $E_{\coh}$ of $E$
formed by the coherent objects of $E$ is stable under fiber products
(and, a fortiori, that $E$ satisfies the equivalent conditions of
\hyperref[VI.prop2.1]{2.1}).
\end{enonce*}

Clearly (i bis) $\Rightarrow$ (i), and the converse follows from
\hyperref[VI.prop1.8]{1.8} (iii). We have (i) $\Rightarrow$ (ii) by
an argument already used: if $X$ and $Y$ are quasi-separated, then by
hypothesis (i), $X$ is quasi-separated over $e$, so $X\times Y$ is
quasi-separated over $Y$ by base change, and hence $X\times Y$ is
quasi-separated by \hyperref[VI.prop1.14]{1.14} (ii). The same
argument shows that (i ter) $\Rightarrow$ (ii bis), while (i)
$\Rightarrow$ (i ter) and (ii) $\Rightarrow$ (ii bis) are immediate.
It therefore remains to prove (ii bis) $\Rightarrow$ (i). If $X$ is a
quasi-separated object of $E$, cover $X$ by objects $X_i$ of $C$
(which is possible since $C$ is generating). Since the $X_i$ are
quasi-compact by the hypothesis on $C$, we may apply
\hyperref[VI.lem2.1.3]{2.1.3} to the morphism $X\to e$; thus, to
prove that this morphism is quasi-separated, we are reduced precisely
to verifying that the $X_i\times X_j$ are quasi-separated, which is
hypothesis (ii bis).

\begin{enonce*}[remark]{Remark 2.2.1}\etiquette[2.2.1]{VI.rem2.2.1}
Condition (i ter) is also equivalent to the following condition (i
quater): for every double arrow $g_1,g_2:Y\rightrightarrows X$ in
$C$, its kernel in $E$ is quasi-compact (or equivalently
(\hyperref[VI.cor1.15]{1.15}) coherent).
\end{enonce*}

To see\pageoriginale this, it is enough to apply criterion
\hyperref[VI.prop1.10]{1.10} (i) to the diagonal morphism
$X\to X\times X$ considered in (i ter).

\begin{enonce*}{Definition 2.3}\etiquette[2.3]{VI.def2.3}
A topos $E$ satisfying the equivalent conditions of
\emph{\hyperref[VI.prop2.2]{2.2} (\resp
\hyperref[VI.prop2.1]{2.1})} is called an \emph{algebraic topos
(\resp} a \emph{locally algebraic topos}, or \emph{locally
coherent)}. An object $X$ of a topos $E$ is called an
\emph{algebraic object of the topos} $E$ if the induced topos $E_{/X}$
is algebraic. A topos $E$ is called \emph{quasi-separated
(\resp coherent)} if it is algebraic and its final object is
quasi-separated (\resp coherent).
\end{enonce*}

\setcounter{subsection}{3}
\subsubsection{}\etiquette[2.3.1]{VI.subsubsec2.3.1}
It is in algebraic topoi that the finiteness notions developed in
${\rm n^{\circ}}$~1 have fully satisfactory properties, analogous to
those of the notions bearing the same names in the category of
schemes. All locally algebraic (= locally coherent) topoi encountered
in practice are in fact algebraic, so the practical interest of this
notion seems rather limited for the present. Nevertheless, one can
construct locally algebraic topoi that are not algebraic
(\hyperref[VI.exer2.17]{2.17} f).


% END INCLUDED SGA4_English_current_buildable_source_20260906/source/06/drafts/SGA4_VI_LINES_1096_1382_English.texfrag


% BEGIN INCLUDED SGA4_English_current_buildable_source_20260906/source/06/drafts/SGA4_VI_LINES_1383_1678_English.texfrag
\subsection{}\etiquette[2.4]{VI.subsec2.4}
Here are some general remarks concerning the notions introduced in
\hyperref[VI.def2.3]{2.3}, which will convince the reader (or so we hope)
that the terminology introduced there is coherent.

\subsubsection{}\etiquette[2.4.1]{VI.subsubsec2.4.1}
Under the conditions of \hyperref[VI.prop2.1]{2.1} (i) \resp (ii), for every
$X\in C$, $X$ (\resp $\epsilon(X)$) is an algebraic (coherent) object of $E$
(\hyperref[VI.cor2.1.2]{2.1.2} and 1.19.2); in other\pageoriginale words, the
induced topos $E_{/X}$ (\resp
$E_{\epsilon(X)}\approxeq(C_{/X})^{\sim}$) is algebraic, hence
\emph{coherent}. Equivalently, \emph{when $C$ admits a final object, $E$
itself is algebraic, hence coherent.}

\subsubsection{}\etiquette[2.4.2]{VI.subsubsec2.4.2}
If $E$ is an algebraic topos, then every induced topos $E_{/X}$ is also
algebraic (\hyperref[VI.prop2.2]{2.2} (i bis)). For a topos $E$ to be locally
algebraic $=$ locally coherent, it is necessary and sufficient that there
exist a family $(X_i)_{i\in I}$ of objects of $E$ covering the final object,
such that the $X_i$ are algebraic (\resp coherent), \ie such that the
induced topoi $E_{/X_i}$ are algebraic (\resp coherent). Sufficiency follows
immediately from the definitions, and necessity from the observations in
\hyperref[VI.subsubsec2.4.1]{2.4.1}. Of course, an algebraic topos is locally
algebraic, and if $E$ is a locally algebraic topos, every induced topos
$E_{/X}$ is locally algebraic.

\subsubsection{}\etiquette[2.4.3]{VI.subsubsec2.4.3}
If $X$ is an algebraic object of a topos $E$, then every object $X'$ of $E$
lying over $X$ is again algebraic (\cf
\hyperref[VI.subsubsec2.4.2]{2.4.2}). In particular, every object of an
algebraic topos is algebraic, and conversely, of course.

\subsubsection{}\etiquette[2.4.4]{VI.subsubsec2.4.4}
The full subcategory of $E$ consisting of the \emph{quasi-separated} objects
that are \emph{algebraic} is stable under fiber products and subobjects
(\hyperref[VI.cor1.15]{1.15}), and hence also under kernels. If $E$ is
algebraic, this category (which is then simply\pageoriginale the category of
quasi-separated objects of $E$) is also stable under products of two
objects. If $E$ is quasi-separated (and only then), this category is stable
under arbitrary finite projective limits, or equivalently contains the
final object of $E$.

The full subcategory $C$ of $E$ consisting of the \emph{coherent} objects
that are \emph{algebraic} is stable under fiber products. If $E$ is locally
algebraic, \ie locally coherent, saying that $E$ is algebraic amounts to
saying that $C$ is also stable under kernels of double arrows
(\hyperref[VI.rem2.2.1]{2.2.1}); and saying that $E$ is quasi-separated
amounts to saying that, in addition, $C$ is stable under products of two
objects (\cf \hyperref[VI.cor1.17]{1.17}). Finally, saying that $E$ is
coherent amounts to saying that $C$ is stable under arbitrary finite
projective limits, or equivalently that it contains the final object of
$E$.

\subsubsection{}\etiquette[2.4.5]{VI.subsubsec2.4.5}
Let $E$ be a topos. For $E$ to be locally coherent (\resp algebraic, \resp
quasi-separated, \resp coherent), it is necessary and sufficient that it
admit a full generating subcategory $C$ consisting of quasi-compact objects
(which will automatically be coherent (\hyperref[VI.prop2.1]{2.1})) and
\emph{stable} under fiber products (\resp under fiber products and kernels,
\resp under fiber products and products of two objects, \resp under
arbitrary finite projective limits); or equivalently that $E$ be equivalent
to a topos $C^{\sim}$, where $C$ is a $𝒰$-site every object of which is
quasi-compact, and in which fiber products and products of two objects
(\resp all finite projective limits) are representable. (For necessity,
take\pageoriginale for $C$ the category of objects of $E$ that are coherent
and \emph{algebraic}, and apply \hyperref[VI.subsec2.4]{2.4}; for
sufficiency, apply \hyperref[VI.prop2.2]{2.2} (i ter).) In this statement,
one may of course also suppose that $C$ is $𝒰$-small.

\subsubsection{}\etiquette[2.4.6]{VI.subsubsec2.4.6}
Let $E$ be a topos. If $E$ is algebraic, an object $X$ of $E$ is
quasi-separated (\resp coherent) if and only if the induced topos $E_{/X}$
is so. When $E$ is no longer assumed algebraic, one can only say that
$E_{/X}$ is quasi-separated (\resp coherent) if and only if $X$ is
quasi-separated (\resp coherent) and \emph{algebraic}. This is of little
consequence, since we will probably not have to use these notions outside
the case in which $E$ is algebraic.

\subsubsection{}\etiquette[2.4.7]{VI.subsubsec2.4.7}
For the notion of a coherent (\resp quasi-separated) object of a topos $E$
to have satisfactory properties, one must restrict to objects that are also
algebraic, a condition automatically satisfied if $E$ is algebraic. Since
we will probably never have to work with locally algebraic topoi $E$ that
are not algebraic, the question of revising the terminology introduced in
\hyperref[VI.def1.13]{1.13}, possibly reserving it for algebraic objects
alone, therefore does not arise with any great urgency!

Moreover, we have refrained from defining the notion of a quasi-compact
topos. In the spirit of the present paragraph, this name would naturally be
given to \emph{algebraic} topoi whose final object is quasi-compact.

We\pageoriginale hesitated to introduce such a usage, since if $X$ is a
topological space, it would not be true that $X$ is quasi-compact if and
only if the associated topos $\Top(X)$ is so. In this connection, note that
$\Top(X)$ is locally coherent if and only if $X$ admits a basis of opens
consisting of quasi-compact opens $U$ such that, for
$U_j,U_k\subset U_i$, $U_j\cap U_k$ is again quasi-compact (and then
$\Top(X)$ is even algebraic). This condition is satisfied for the
topological spaces underlying schemes, but rarely for spaces encountered by
analysts, even when the spaces themselves are compact. For ease of
reference, let us state explicitly:

\begin{enonce*}{Proposition 2.5}\etiquette[2.5]{VI.prop2.5}
Let $E$ be an algebraic topos and $f:X\to Y$ a morphism in $E$. If $Y$ is a
quasi-separated (\resp coherent) object of $E$, then, for $X$ to be a
quasi-separated (\resp coherent) object of $E$, it is necessary and
sufficient that the morphism $f$ be quasi-separated (\resp coherent).
\end{enonce*}

Sufficiency was already seen in \hyperref[VI.prop1.14]{1.14} (ii), and holds
without any hypothesis on $E$. Necessity in the non-parenthetical case holds
even without supposing that $Y$ is quasi-separated, and is simply Definition
\hyperref[VI.def2.3]{2.3}, via \hyperref[VI.prop2.2]{2.2} i bis). It remains
to prove that if $Y$ is coherent and $X$ is coherent, then $f$ is coherent.
Since we already know that $f$ is quasi-separated, it is enough to see that
$f$ is quasi-compact, which follows from the fact that $X$ is quasi-compact
and $Y$ quasi-separated (\hyperref[VI.def1.13]{1.13}).

\begin{enonce*}{Corollary 2.6}\etiquette[2.6]{VI.cor2.6}
Let\pageoriginale $E$ be an algebraic topos, $f:X\to Y$ a morphism in $E$,
and $(Y_i\to Y)_{i\in I}$ a covering family, with the $Y_i$ quasi-separated
(\resp coherent). For $f$ to be quasi-separated (\resp quasi-compact,
\resp coherent), it is necessary and sufficient that, for every $i\in I$,
$X_i=\fprod{X}{Y_i}{Y}$ be a quasi-separated (\resp quasi-compact, \resp
coherent) object of $E$.
\end{enonce*}

It is enough to combine \hyperref[VI.prop1.10]{1.10} (ii) and
\hyperref[VI.prop2.5]{2.5} in the ``quasi-separated'' or ``coherent'' case;
the ``quasi-compact'' case was already seen in
\hyperref[VI.cor1.16]{1.16}. Note that the necessity of the condition stated
in \hyperref[VI.cor2.6]{2.6} also holds without any hypothesis on the topos
$E$.

As a special case of \hyperref[VI.cor2.6]{2.6}, one obtains:

\begin{enonce*}{Corollary 2.7}\etiquette[2.7]{VI.cor2.7}
Let $E$ be a coherent topos and $X$ an object of $E$. For $X$ to be a
quasi-compact (\resp quasi-separated, \resp coherent) object of $E$, it is
necessary and sufficient that it be quasi-compact (\resp quasi-separated,
\resp coherent) over the final object
(\hyperref[VI.subsubsec1.9.2]{1.9.2}).
\end{enonce*}

In particular, $X$ is constructible (\hyperref[VI.subsubsec1.9.2]{1.9.2})
if and only if it is coherent.

\begin{enonce*}{Corollary 2.8}\etiquette[2.8]{VI.cor2.8}
Let $E$ be a locally coherent topos and $f:X\to Y$ a morphism in $E$. For
$f$ to be quasi-separated, it is necessary and sufficient that, for every
quasi-separated object $Y'$ of $E_{/Y}$,
$f^{\ast}(Y')=\fprod{X}{Y'}{Y}$ be a quasi-separated object of $E_{/X}$ (or
of $E$, which amounts to the same thing
(\hyperref[VI.subsubsec1.20.2]{1.20.2})).
\end{enonce*}

The\pageoriginale ``only if'' part holds without any condition on $E$,
since if $Y'$ is quasi-separated, then
$f':X'=\fprod{X}{Y'}{Y}\to Y'$ is quasi-separated by base change, and hence
$X'$ is quasi-separated (\hyperref[VI.prop1.14]{1.14} (ii)). For the
converse, note that there exists a covering family $Y_i\to Y$ of $Y$ by
quasi-separated algebraic objects. By \hyperref[VI.prop1.10]{1.10} (ii), to
verify that $f$ is quasi-separated it suffices to verify that the
$f_i:X_i=\fprod{X}{X_i}{Y}\to Y_i$ are so. But the hypothesis on $f^{\ast}$
implies that the $X_i$ are quasi-separated, and therefore so are the $f_i$,
since $E_{/Y_i}$ is an algebraic topos.

\begin{enonce*}[remark]{Remark 2.8.1}\etiquette[2.8.1]{VI.rem2.8.1}
Statements \hyperref[VI.prop2.5]{2.5} and \hyperref[VI.cor2.6]{2.6} remain
valid if the hypothesis ``$E$ algebraic'' is replaced by the more general
hypothesis ``$Y$ is algebraic'', as is seen immediately by applying the
original statement to the induced topos $E_{/Y}$, which is algebraic.
Likewise, in \hyperref[VI.cor2.8]{2.8} it is enough to suppose that $E_{/Y}$
(rather than $E$) is locally coherent.
\end{enonce*}

\begin{enonce*}{Proposition 2.9}\etiquette[2.9]{VI.prop2.9}
Let $E$ be a topos and $C$ a subcategory of $E$. The following conditions
are equivalent:
\begin{enumerate}
\renewcommand{\theenumi}{\roman{enumi}}
\renewcommand{\labelenumi}{(\theenumi)}
\item $E$ is locally coherent (\resp coherent), satisfies the equivalent
  conditions of \hyperref[VI.cor1.25]{1.25}, and $C$ is the full subcategory
  $E_{\coh}$ of $E$ consisting of the coherent objects.

\item $C$ is a strictly full generating subcategory of $E$, stable under
  finite inductive limits and fiber products (\resp and under finite
  projective limits), and consists of quasi-compact objects.
\end{enumerate}
\end{enonce*}

The implication\pageoriginale (i) $\Rightarrow$ (ii) follows immediately
from the definitions and from formulation \hyperref[VI.cor1.25]{1.25} (ii)
of the conditions under consideration in \hyperref[VI.cor1.25]{1.25}. Let
us prove the converse. The fact that $C$ is a full generating subcategory
consisting of quasi-compact objects and stable under fiber products (\resp
under finite projective limits) implies, by definition, that $E$ is locally
coherent (\resp coherent). The additional fact that $C$ is strictly full
and stable under finite inductive limits then implies that $C=E_{\coh}$ (by
\hyperref[VI.cor1.24]{1.24} (ii) $\Rightarrow$ (iii)), and that condition
\hyperref[VI.cor1.25]{1.25} (ii) is satisfied, \cqfd

\begin{enonce*}{Definition 2.9.1}\etiquette[2.9.1]{VI.def2.9.1}
A $𝒰$-topos $E$ is called perfect if it is coherent
(\hyperref[VI.def2.3]{2.3}) and satisfies the equivalent conditions of
\hyperref[VI.cor1.25]{1.25}, \ie if the subcategory $E_{\coh}$ of $E$
consisting of the coherent objects of $E$ is stable under finite inductive
limits. We say that $E$ is locally perfect if there exists a family
$(X_i)_{i\in I}$ of objects of $E$ covering the final object such that, for
every $i\in I$, the induced topos $E_{/X_i}$ is perfect.
\end{enonce*}

\setcounter{subsection}{9}
\setcounter{subsubsection}{1}
\subsubsection{}\etiquette[2.9.2]{VI.subsubsec2.9.2}

It follows immediately from this definition that a perfect (\resp locally
perfect) topos is coherent (\resp locally coherent). As examples of perfect
(\resp locally perfect) topoi, let us mention the Noetherian (\resp locally
Noetherian) topoi introduced in \hyperref[VI.def2.11]{2.11} below (\cf
\hyperref[VI.subsec2.14]{2.14}), and the Zariski and étale topoi of a
coherent scheme (\resp an arbitrary scheme) (IX, note page 42). As a
counterexample,\pageoriginale let us mention the fppf topos of a scheme $S$,
which is locally coherent (and even coherent if $S$ is a coherent scheme,
\ie quasi-compact and quasi-separated)
(\hyperrefext[VII][VII.lem5.6]{5.6}), but is not locally perfect if
$S\neq∅$ (\hyperref[VI.exer1.28]{1.28} f)).

In a locally perfect topos, the notion of constructible object is
particularly stable (\hyperref[VI.prop2.9.3]{2.9.3} below), which is one
reason why the notion seems useful. Another reason is that, like the notion
of a coherent or locally coherent topos, the notion of a perfect or locally
perfect topos is stable under formation of the closed subtopos complementary
to an open (\hyperref[VI.prop4.6]{4.6}), and behaves particularly simply
under the operation of gluing topoi (\hyperref[VI.cor4.11]{4.11},
\hyperref[VI.cor4.14]{4.14}).

\begin{enonce*}{Proposition 2.9.3}\etiquette[2.9.3]{VI.prop2.9.3}
Let $E$ be a locally perfect topos (\hyperref[VI.def2.9.1]{2.9.1}). Then the
full subcategory $E_{\cons}$ of $E$ consisting of the constructible objects
of $E$ (\hyperref[VI.subsubsec1.9.3]{1.9.3}) is stable under finite
inductive limits (and also, of course, under finite projective limits, by
(\hyperref[VI.subsubsec1.9.3]{1.9.3})).
\end{enonce*}

Since the property of an object of $E$ being constructible is local on $E$
(\hyperref[VI.prop1.10]{1.10} (ii)), we are reduced to the case in which $E$
is a perfect topos. But then $E_{\cons}=E_{\coh}$
(\hyperref[VI.cor2.7]{2.7}), and the conclusion follows from Definition
\hyperref[VI.def2.9.1]{2.9.1}.

\begin{enonce*}{Corollary 2.9.4}\etiquette[2.9.4]{VI.cor2.9.4}
Let\pageoriginale $E$ be a locally perfect topos. Then $E$ is perfect if and
only if $E$ is coherent. For every object $X$ of $E$, the induced topos
$E_{/X}$ is locally perfect.
\end{enonce*}

\begin{enonce*}[remark]{Problem 2.9.5}\etiquette[2.9.5]{VI.prob2.9.5}
It is conceivable that perfect topoi are sufficiently special to admit a
\emph{structure theory} that is equally explicit (following a model
proposed by M.\textsc{Artin} at the time of the oral seminar). Let us call a
topos \emph{finite} if it is equivalent to a topos of the form $\widehat{C}$,
where $C$ is a finite category (\cf Exercise
\hyperref[VI.exer3.11]{3.11}), and call a topos \emph{profinite} if it is a
filtered projective limit of finite topoi (in the sense of 6 below, which
applies by \hyperref[VI.exer3.12]{3.12} c)). Since a finite topos is
Noetherian (\hyperref[VI.exer2.17]{2.17} g)), hence perfect, and a filtered
projective limit of perfect topoi is perfect (6), it follows that a profinite
topos is perfect. The question is whether, conversely, every perfect topos
is profinite. This is equivalent to asking whether every finite subcategory
of $E_{\coh}$ is contained in a full subcategory $F_0$ of $E_{\coh}$ that is
equivalent to a category $F_{\coh}$, where $F$ is a finite topos (\cf
\hyperref[VI.exer3.11]{3.11}). (If the answer were negative, one would need
to find manageable additional intrinsic conditions on a perfect topos that
ensure that it is profinite.) Finally, it would remain to study the
structure of profinite topoi in terms of a suitable notion of a
``Karoubian profinite category'', inspired by IV \hyperref[IV.exer7.6]{7.6} h). This\pageoriginale
study has not yet been carried out even in the special case in which $E$ is
the Zariski or étale topos of a good coherent scheme $X$, and should then
give invariants finer than the compact space $X_{\cons}$ (\EGA IV 1.)
associated with $X$.

In the same vein, let us mention the following question: Let $E$ be a
perfect topos, $E'$ the topos induced on an open of $E$, and $E''$ the
corresponding closed subtopos (IV \hyperref[IV.subsubsec9.3.5]{9.9}). If $E'$ and $E''$ are finite topoi,
is the same true of $E$?
\end{enonce*}

% END INCLUDED SGA4_English_current_buildable_source_20260906/source/06/drafts/SGA4_VI_LINES_1383_1678_English.texfrag


% BEGIN INCLUDED SGA4_English_current_buildable_source_20260906/source/06/drafts/SGA4_VI_LINES_1679_2073_English.texfrag
\begin{enonce*}{Proposition 2.10}\etiquette[2.10]{VI.prop2.10}
Let $E$ be a topos. The following conditions are equivalent:
\begin{itemize}
\item[(i)] $E$ admits a full generating subcategory, stable under
  fiber products, consisting of pre-Noetherian objects of $E$.

\item[(ii)] $E$ admits a full generating subcategory consisting of
  quasi-separated pre-Noetherian objects (\ie coherent
  pre-Noetherian objects).

\item[(iii)] $E$ is locally coherent (\hyperref[VI.def2.3]{2.3}), and
  every quasi-compact object of $E$ is pre-Noetherian.

\item[(iii bis)] $E$ is locally coherent and admits a generating
  family consisting of pre-Noetherian objects of $E$.
\end{itemize}
\end{enonce*}

The equivalence of (iii) and (iii bis) follows from
\hyperref[VI.subsec1.32]{1.32}, and it is trivial that
(i) $\Rightarrow$ (iii bis) and (iii) $\Rightarrow$ (i); hence (i) is
equivalent to (iii) and (iii bis). Finally, (i) $\Rightarrow$ (ii)
follows from the last assertion of \hyperref[VI.prop2.1]{2.1},
\pageoriginale and it remains to prove (ii) $\Rightarrow$ (i). This
implication follows from the fact that the full subcategory of $E$
formed by the quasi-separated Noetherian objects, if it is generating,
is stable under fiber products: such a fiber product is quasi-compact
by the definitions, hence Noetherian by the first assertion of
\hyperref[VI.subsec1.32]{1.32}; and it is quasi-separated by the last
assertion of \hyperref[VI.subsec1.32]{1.32}.

\begin{enonce*}{Definition 2.11}\etiquette[2.11]{VI.def2.11}
A topos $E$ is called locally Noetherian if it satisfies the
equivalent conditions of \hyperref[VI.prop2.10]{2.10}, and Noetherian
if, in addition, its final object is coherent (or, equivalently
(\hyperref[VI.subsec1.32]{1.32}), pre-Noetherian and quasi-separated).
An object $X$ of a topos $E$ is called a Noetherian object of $E$ if
the induced topos $E_{/X}$ is a Noetherian topos.
\end{enonce*}

\setcounter{subsection}{11}
\subsection{}\etiquette[2.12]{VI.subsec2.12}
If $E$ is a locally Noetherian topos, then so is every induced topos
$E_{/X}$. Conversely, if one can find a family $(X_i)$ covering the
final object such that the $E_{/X_i}$ are locally Noetherian, then so
is $E$. If $C$ is a full generating subcategory of $E$ consisting of
pre-Noetherian objects and stable under fiber products
(\hyperref[VI.def2.11]{2.11} (i)), then for every $X \in C$,
$E_{/X}$ is a Noetherian topos, \ie the objects of $C$ are even
\emph{Noetherian} (for, if $E$ is locally Noetherian and
$X \in \Ob E$, then $E_{/X}$ is Noetherian if and only if $X$ is
coherent). It follows that a topos $E$ is locally Noetherian if and
only if the final object of $E$ can be covered by Noetherian objects
$X_i$.

\subsection{}\etiquette[2.13]{VI.subsec2.13}
In a\pageoriginale locally Noetherian topos $E$, every monomorphism is
quasi-compact and every morphism is quasi-separated
(\hyperref[VI.subsec1.32]{1.32}). A fortiori, $E$ is an
\emph{algebraic} topos (\hyperref[VI.def2.3]{2.3}), and not merely
locally coherent, \ie locally algebraic. Thus a topos $E$ is
Noetherian if and only if it is locally Noetherian and \emph{coherent}
(\hyperref[VI.def2.3]{2.3}). In other words, if $E$ is a locally
Noetherian topos, then the Noetherian objects of $E$ are precisely the
\emph{coherent} objects of $E$.

\subsection{}\etiquette[2.14]{VI.subsec2.14}
Let $E$ be a locally Noetherian topos. If $E$ is quasi-separated, \ie
if its final object is quasi-separated, then \emph{every} object of
$E$ is quasi-separated (\hyperref[VI.subsec1.32]{1.32}); hence the
pre-Noetherian objects of $E$ are the same as the Noetherian (\ie
coherent) objects of $E$. Since a quotient object of a pre-Noetherian object
is pre-Noetherian, it follows that $E$ satisfies the equivalent
conditions of \hyperref[VI.cor1.25]{1.25} (because it satisfies
condition \hyperref[VI.cor1.25]{1.25} (ii bis)). In particular, if
$C$ is the full subcategory of $E$ consisting of the Noetherian
(coherent) objects of $E$, then the natural functor
$$
\Ind(C) \to E
$$
is an equivalence of categories. We conclude that \emph{a Noetherian
topos is perfect} (\hyperref[VI.def2.9.1]{2.9.1}), and hence that a
locally Noetherian topos is locally perfect.

If one assumes only that $E$ is locally Noetherian, it remains true
that every object of $E$ is the inductive limit of its pre-Noetherian
subobjects (for, in a topos admitting a generating subcategory
consisting\pageoriginale of quasi-compact objects, every object is the
filtered inductive limit of its quasi-compact subobjects). But since a
pre-Noetherian object of $E$ is no longer necessarily coherent, this
statement then has only very limited usefulness, and we know from
(\hyperref[VI.exem1.33]{1.33}) that the conditions of
\hyperref[VI.cor1.25]{1.25} are no longer necessarily satisfied.

\begin{enonce*}[remark]{Example 2.15}\etiquette[2.15]{VI.exem2.15}
(\emph{Topological spaces}.) Let $X$ be a topological space. Then the
following conditions are equivalent: (i) the topological space $X$ is
Noetherian (\ie every open subset of $X$ is quasi-compact, \ie every
increasing sequence of open subsets of $X$ is stationary); (ii) the
topos $\Top(X)$ is Noetherian; (iii) the final object $X$ of
$\Top(X)$ is Noetherian; (iv) the final object $X$ of $\Top(X)$ is
pre-Noetherian.

Indeed, trivially (ii) $\Leftrightarrow$ (iii) $\Rightarrow$ (iv)
$\Leftrightarrow$ (i). On the other hand, (i) implies that the
generating subcategory Ouv(X) of $\Top(X)$, which is stable under
finite projective limits, consists of pre-Noetherian objects, whence
(ii).

Likewise, one sees that the topological space $X$ is locally
Noetherian (\ie is a union of Noetherian open subsets) if and only if
the topos $\Top(X)$ is locally Noetherian.
\end{enonce*}

\setcounter{subsection}{15}
\setcounter{subsubsection}{0}
\subsubsection{}\etiquette[2.15.1]{VI.subsubsec2.15.1}
Thus our Definitions \hyperref[VI.def1.30]{1.30} and
\hyperref[VI.def2.11]{2.11} agree with the accepted terminology for
topological spaces. On the other hand, in the general case we have
chosen to give the term ``Noetherian object'' a\pageoriginale stronger
meaning than the more naive notion of \hyperref[VI.def1.30]{1.30}, so
that all the properties associated with this notion (rather than
merely the grammatical structure of the definition) agree fully with
the intuition attached to \emph{Noetherian} topological spaces and
schemes.

The reader will find further examples in the following exercises.

\begin{enonce*}[remark]{Exercise 2.16}\etiquette[2.16]{VI.exer2.16}
(\emph{Spaces with operators and classifying topoi}.)

Let $X$ be a topological space on which a discrete group $G$ acts,
giving (IV \hyperref[IV.subsec2.3]{2.3}) a topos $E = \Top(X, G)$. The objects $X'$ of this
topos are interpreted as étalé spaces over $X$ with operator group
$G$, and we note that the induced topos $E_{/X'}$ is canonically
equivalent to $\Top(X', G)$.
\begin{enumerate}
\renewcommand{\theenumi}{\alph{enumi}}
\renewcommand{\labelenumi}{\theenumi)}
\item The final object of the topos $E$ is quasi-compact if and only
  if the quotient topological space $X/G$ is quasi-compact (use IV
  \hyperref[IV.subsubsec8.4.1]{8.4.1}).

\item Let $P$ be the object of $E$ defined by the trivial étalé space
  $X \times G$ with the diagonal action of $G$. Show that the induced
  topos $E_{/P}$ is equivalent to the topos $\Top(X)$ (compare IV
  \hyperref[IV.subsubsec5.8.3]{5.8.3}). Conclude that $E$ is locally coherent (\resp locally
  Noetherian) if and only if $\Top(X)$ is locally coherent (\cf
  \hyperref[VI.subsubsec2.4.7]{2.4.7}) (\resp locally Noetherian (\cf
  \hyperref[VI.exem2.15]{2.15})). Show that if $E$ is locally coherent,
  \ie locally algebraic, it is in fact algebraic, using the generating
  family formed\pageoriginale by the objects $T_U$, where $U$ is a
  coherent open subset of $X$, $T_U = G \times U$ with the action
  $g \cdot (u, g') = (u, gg')$ of $G$, and the structural morphism
  $p_U : T_U \to X$ defined by $p_U(u, g) = g\cdot u$.

\item Suppose that $E$ is algebraic, \ie that $\Top(X)$ is algebraic
  (\hyperref[VI.subsubsec2.4.7]{2.4.7}). Show that the structural
  morphism $P \to e$ is quasi-separated. Show that $E$ is
  quasi-separated if and only if $\Top(X)$ is quasi-separated (or,
  equivalently, the topological space $X$ is quasi-separated, \ie the
  intersection of two quasi-compact open subsets of $X$ is a
  quasi-compact open subset), \emph{and} $G$ is finite or $X$ is empty.
  (Compare \hyperref[VI.exem1.33]{1.33}.)

\item Show that the structural morphism $P \to e$ is quasi-compact if
  and only if $G$ is finite. Show that $E$ is coherent (\resp
  Noetherian) if and only if $\Top(X)$ is so, \emph{and} $G$ is finite
  or $X$ is empty.

\item Let $E$ be a topos and $G$ a group of $E$, giving a classifying
  topos $B_G$ (IV \hyperref[IV.subsec2.4]{2.4}). Study the finiteness conditions in $B_G$, using
  what precedes as a guide. Ask the same question starting from a
  strict pro-group $𝒢 = (G_i)_{i \in I}$ of $E$. In particular, show
  that if $E$ is coherent (\resp Noetherian) and the $G_i$ are
  coherent, then $B_{𝒢}$ is coherent (\resp Noetherian).
\end{enumerate}
\end{enonce*}

\begin{enonce*}[remark]{Exercise 2.17}\etiquette[2.17]{VI.exer2.17}
(\emph{Topoi of the form $\widehat{C}$}.)
Let $C$ be a category equivalent to a category $\in 𝒰$, and let
$E = \widehat{C}$ be the category of presheaves on $C$. By means of
the canonical functor $\varepsilon : C \to \widehat{C}$, we identify
$C$ with a full subcategory of $E$.
\begin{enumerate}
\renewcommand{\theenumi}{\alph {enumi}}
\renewcommand{\labelenumi}{(\theenumi)}
\item $C$ is\pageoriginale a generating subcategory of $E$ consisting
  of quasi-compact objects. These objects are pre-Noetherian (\cf
  \hyperref[VI.subsec1.32]{1.32}) if and only if, for every object $X$
  of $C$, every sieve on $X$ is generated by a finite family of
  morphisms $X_i \to X$ of $C$. The final object of $E$ is
  quasi-compact (\resp pre-Noetherian) if and only if there is a final
  subcategory of $C$ whose set of objects is finite (\resp if every
  sieve on $C$ is generated by a subcategory of $C$ whose set of
  objects is finite).

\item The topos $E$ admits a generating subcategory consisting of
  coherent (or, equivalently, quasi-separated) objects if and only if
  the objects of $C$ are coherent in $E$, or equivalently if, for two
  arrows $Y \xrightarrow{f} X$, $Z \xrightarrow{g} X$ in $C$, the
  object $\fprod{Y}{Z}{X}$ of $E$ is quasi-compact, \ie if one can find
  a finite family of commutative squares in $C$
\[
\xymatrix@=1cm{
 & T_i \ar[ld] \ar[rd]& \\
Y\ar[rd]_{f}&    & Z\ar[ld]^{g}\\
 & X  &,
}
\]
such that every other commutative square in $C$
\[
\xymatrix@=1cm{
 & T \ar[ld] \ar[rd]& \\
Y\ar[rd]_{f}&    & Z\ar[ld]^{g}\\
 & X  &,
}
\]
comes\pageoriginale from a morphism $T \to T_i$. (Note that if
$(F_i)_{i \in I}$ is a generating family of $E$, then every
$X \in \ob C$ is isomorphic to a direct factor of one of the $F_i$.)

\item The topos $E$ is locally coherent, \ie is generated by a family
  of coherent and \emph{algebraic} objects, if and only if the objects
  $X$ of $C$ are coherent and algebraic objects of $E$, \ie if
  condition b) and the following condition hold: for every pair of
  arrows $f$, $g : Y \rightrightarrows Z$ in $C$ \emph{over an object
  $X$ of} $C$, the kernel of this pair of arrows in $E$ is a
  quasi-compact object of $E$, \ie there is a finite family of
  commutative diagrams in $C$
\[
\xymatrix@C=1.5cm{
T_i \ar[r]& Y \ar@<2pt>[r]^{f,g}\ar@<-2pt>[r]&Z
}
\]
such that every other commutative diagram in $C$
\[
\xymatrix@C=1.5cm{
T \ar[r]& Y \ar@<2pt>[r]^{f,g}\ar@<-2pt>[r]&Z
}
\]
comes from a morphism $T \to T_i$. The topos $E$ is algebraic if and
only if it satisfies condition b) and the preceding condition, but
with an arbitrary pair of arrows $(f, g)$ of $C$ (not necessarily
over an object $X$ of $C$). The topos $E$ is quasi-separated if and
only if it satisfies the conditions of b), and if, in addition, for
two objects $X$, $Y\in\ob C$, the product $X\times Y$ in $E$ is
quasi-compact. The topos $E$ is coherent if and only if it satisfies
the two preceding conditions and there is a final subcategory of $C$
whose underlying set is finite.

\item Give\pageoriginale an example in which the objects of the
  generating subcategory $C$ of $E$ are pre-Noetherian, but $E$ admits
  no generating family consisting of coherent objects (\ie b) fails,
  or the objects of $C$ are not all coherent). For this, take for $C$
  the category having as objects $X$, $T_i$ $(i = 0, 1, \ldots)$, and
  $e$ (the final object), the only morphisms between these objects,
  besides the structural morphisms to $e$ and the identity morphisms,
  being morphisms $u_i : T_i \to X$ and $v_i : X \to T_i$ subject to
  the conditions $u_i v_i = \id_X$, and finally set
  $p_i = v_i u_i$ (which necessarily satisfies
  $p_i^2 = p_i$). One verifies that the only sieves on $X$ or on any
  $T_i$ are the two trivial sieves, and that $e$ admits exactly one
  nontrivial sieve. Thus, by a), the objects of $C$ are
  pre-Noetherian objects of $E$. Nevertheless, the product
  $X \times X$ in $E$ is not quasi-compact, since it does not satisfy
  the criterion of b).

\item Give an example in which the generating subcategory $C$ of $E$
  consists of coherent objects of $E$, but $E$ is not locally
  coherent. For this, take for $C$ the category whose set of objects
  consists of the distinct objects $e$ (the final object), $Y$, $Z$,
  and $T_i$ $(i=0, 1, \ldots)$, with the only morphisms between these
  objects, besides the morphisms to the final object $e$ and the
  identity morphisms, being morphisms $f$, $g : Y \rightrightarrows Y$
  subject to the conditions $fu_i = gu_i$. One verifies that $C$
  satisfies condition b) (with a little patience; one finds that all
  fiber products of objects of $C$ are isomorphic to $∅_E$ or belong
  to $C$,\pageoriginale with the exception of
  $\fprod{Z}{Z}{e}$ and $\fprod{Y}{Z}{e}$, which are covered by two
  elements of $C$), but clearly $\Ker(f, g)$ does not satisfy
  condition c).

\item Give an example in which $C$ is stable under fiber products (a
  fortiori, $E$ is a locally coherent, \ie locally algebraic, topos),
  but $E$ is not an algebraic topos. (Take the example given in d),
  and the full subcategory $C'$ of $E$ consisting of the objects $Y$,
  $Z$, $T_{i}$ $(i > 0)$, and $∅_{E}$.)

\item Suppose that the category $C$ is finite. Prove that the topos
  $E$ is Noetherian; that its Noetherian (\ie quasi-compact) objects
  are the contravariant functors $F : C^{\circ} \to (\Ens)$ such that,
  for every object $X$ of $C$, $F(X)$ is a \emph{finite} set; and that
  the fiber functors on $E$ take Noetherian objects to finite sets
  (\cf IV \hyperref[IV.exer7.6]{7.6}. h)).

\item Suppose that every morphism $f: X \to X$ of $C$ factors as a
  composite $ip$, where $i$ is a monomorphism and $p$ has a right
  inverse. Let $X$ be an object of $C$, let $I(X)$ be the set of its
  subobjects in the sense of $C$, regarded as an ordered set, and let
  $S(I(X))$ be the set of subsets $U$ of $I(X)$ such that, for two
  elements $X'$, $X'' \in I(X)$, the conditions $X'\in U$ and
  $X''\leq X'$ imply $X''\in U$. Show that the ordered set of
  subobjects in $E$ of $X$ is isomorphic to the set $S(I(X))$ (ordered
  by inclusion). Conclude that if $I(X)$ is finite, then $X$ is a
  pre-Noetherian object of $E$. In particular, if (in addition to the
  factorization condition above) $C$ is stable under fiber products
  (\resp under finite projective limits), and\pageoriginale if, for
  every $X \in \ob C$, the set $I(X)$ of subobjects of $X$ in $C$ is
  finite, then the topos $E$ is locally Noetherian (\resp Noetherian).

\item Let $C$ be the category of finite sets (or, alternatively, of
  nonempty finite sets) $\in 𝒰$, so that $E = \widehat{C}$ is the
  category of augmented simplicial sets (\resp of simplicial sets
  simpliciter). Show that $E$ is a Noetherian topos. (Use h).) Taking
  a nonessentially constant pro-object $(X_i)_{i \in I}$ of $C$, show
  that $E$ admits fiber functors that do not take Noetherian objects
  to Noetherian objects (\ie to finite sets).

\item Consider the following diagram of implications among properties
  of a topos $E$:
\[
\xymatrix@R=.75cm{
**[r]\txt{E is Noetherian} \ar@{=>}[d] \ar@{=>}[r]& **[r]\txt{E is
  locally Noetherian} \ar@{=>}[d] & \ar@{=>}[rdd]\\
**[r]\txt{E is coherent} \ar@{=>}[r] & **[r]\txt{E
is algebraic} \ar@{=>}[d]\\
 & **[r]\txt<9.5pc>{E is loc. coherent, i.e.\\E \raisebox{-1ex}{cohalg}
  generates E} \ar@{=>}[d]& & \txt{E \raisebox{-1ex}{pre-Noeth}
  generates E} \ar@{=>}[ldd] \\
 & **[r]\txt{E \raisebox{-1ex}{coh} generates E} \ar@{=>}[d]&\\
 & **[r]\txt{E \raisebox{-1ex}{qucpct} generates E}  &&.\\
}
\]
Show that every implication in this diagram is strict, and that among
the notions considered there are no implications other than the
composites of the implications in the preceding diagram. Here
$E_{\coh}, E_{\text{qucpct}}$, $E_{\text{prénoeth}}$,
$E_{\text{cohalg}}$\pageoriginale denote respectively the full
subcategories of $E$ consisting of the coherent, \resp quasi-compact,
\resp pre-Noetherian, \resp coherent and algebraic objects of $E$.
(One may restrict attention to topoi of the form $\widehat{C}$, using
the results stated in d), e), and f).)

\item Answer the following question (to which the author of these
  lines does not know the answer): can $E = \widehat{C}$ be locally
  Noetherian without the $\Hom(X, Y)$ (for $X$, $Y \in \ob C$) being
  finite?
\end{enumerate}
\end{enonce*}

\begin{enonce*}[remark]{Exercise 2.18}\etiquette[2.18]{VI.exer2.18}
Let $E$ be a topos.
\begin{enumerate}
\renewcommand{\theenumi}{\alph{enumi}}
\renewcommand{\labelenumi}{\theenumi)}
\item Let $X$ be a pre-Noetherian object of $E$. Show that $X$ is
  isomorphic to a finite sum of connected objects (IV \hyperref[IV.subsubsec4.3.5]{4.3.5}) of $E$.
  (Prove by contradiction that an increasing sequence of partitions
  of $X$ (IV \hyperref[IV.exer8.7]{8.7}) is stationary.)

\item Let $X$ be an object of $E$ and let
  $(f_i : X_i \to X)_{i \in I}$ be a covering family of $X$. Show that
  if each $X_i$ is isomorphic to a sum of connected objects of $E$,
  then so is $X$. (Reduce to the case in which the $X_i$ are connected,
  then to the case in which they are subobjects of $X$, and consider
  on the index set $I$ the equivalence relation $R$ generated by the
  relation $X_i \cap X_j \neq \phi_E$. Show that if $J = I/R$, then
  $X$ is the sum of the
  $X(i) = \displaystyle{\mathop{\Sup}_{i \in I}} X_i .)$

\item Conclude from b) that if $(X_i)_{i \in I}$ denotes a generating
  family of $E$, then $E$ is locally connected (IV \hyperref[IV.exer8.7]{8.7}. $\ell$)) if
  and\pageoriginale only if every $X_i$ $(i \in I)$ is isomorphic to a
  sum of connected objects of $E$.

\item Conclude from a) and c) that if $E$ admits a generating
  subcategory consisting of pre-Noetherian objects, in particular if
  $E$ is locally Noetherian, then $E$ is locally connected, and a
  fortiori is isomorphic to the sum topos (IV \hyperref[IV.exer8.7]{8.7} b)) of a family of
  connected topoi (IV \hyperref[IV.exer8.7]{8.7} e)), \ie topoi whose final object is
  connected. (In particular, one can associate with $E$, for every
  morphism $f : P \to E$, where $P$ is a ``nonempty connected and
  simply connected'' topos (IV \hyperref[IV.exer2.7.5]{2.7.5}), a fundamental pro-group
  $\pi_1(E, f)$. Note that if $E$ is locally Noetherian and ``nonempty,''
  such an $f$ can always be found, with $P$ the punctual topos, by a
  result of \textsc{Deligne} ([9]).
\end{enumerate}
\end{enonce*}

% END INCLUDED SGA4_English_current_buildable_source_20260906/source/06/drafts/SGA4_VI_LINES_1679_2073_English.texfrag


% BEGIN INCLUDED SGA4_English_current_buildable_source_20260906/source/06/drafts/SGA4_VI_LINES_2074_2418_English.texfrag
\section{Finiteness conditions for a morphism of topoi}\etiquette[3]{VI.sec3}

\begin{enonce*}{Definition 3.1}\etiquette[3.1]{VI.def3.1}
Let $f:E'\to E$ be a morphism of topoi. We say that $f$ is
\emph{quasi-compact (\resp quasi-separated)} if, for every
quasi-compact (\resp quasi-separated) object $X$ of $E$,
$f^{\ast}(X)$ is quasi-compact (\resp quasi-separated). We say that
$f$ is coherent if $f$ is quasi-compact and quasi-separated.
\end{enonce*}

\setcounter{subsection}{1}
\setcounter{subsubsection}{0}
\subsubsection{}\etiquette[3.1.1]{VI.subsubsec3.1.1}
Note that if $f$ has any one of the preceding properties, then every
morphism of topoi isomorphic to $f$ has the same property. It is also
immediate that the composite of two quasi-compact (\resp
quasi-separated, \resp coherent) morphisms of topoi is again
quasi-compact (\resp quasi-separated, \resp coherent).

\begin{enonce*}{Proposition 3.2}\etiquette[3.2]{VI.prop3.2}
With the\pageoriginale notation of \hyperref[VI.def3.1]{3.1}, let
$(X_i)_{i\in I}$ be a generating family of quasi-compact (\resp
coherent) objects of $E$. In the parenthetical case, suppose that
$E'$ admits a generating family formed by quasi-compact objects. Then
$f$ is quasi-compact (\resp coherent) if and only if, for every
$i\in I$, $f^{\ast}(X_i)$ is quasi-compact (\resp coherent).
\end{enonce*}

Necessity is immediate. For sufficiency in the first case, note that,
for every quasi-compact object $X$ of $E$, there is a finite
epimorphic family of morphisms $X_i\to X$. Hence there is a finite
epimorphic family $f^{\ast}(X_i)\to f^{\ast}(X)$, with the
$f^{\ast}(X_i)$ quasi-compact by hypothesis, and therefore
$f^{\ast}(X)$ is quasi-compact (\hyperref[VI.prop1.3]{1.3}). In the
second case, it remains to see that if $X$ is a quasi-separated
object of $E$, then $f^{\ast}(X)$ is quasi-separated. The hypothesis
on $X$ may be expressed (\hyperref[VI.cor1.17]{1.17}) by the
existence of an epimorphic family of morphisms $X_i\to X$ such that
the $\fprod{X_i}{X_j}{X}$ are quasi-compact. We then have a covering
family $f^{\ast}(X_i)\to f^{\ast}(X)$ such that the fiber products
$\fprod{f^{\ast}(X_i)}{f^{\ast}(X_j)}{f^{\ast}(X)}$ are
quasi-compact (since, by the first case, $f^{\ast}$ carries
quasi-compact objects to quasi-compact objects). As the
$f^{\ast}(X_i)$ are coherent, the conclusion follows again from
\hyperref[VI.cor1.17]{1.17}.

\begin{enonce*}{Corollary 3.3}\etiquette[3.3]{VI.cor3.3}
Let $C$ and $C'$ be two $𝒰$-sites in which fiber products are
representable and every covering family admits a finite covering
subfamily, and let $g:C\to C'$ be a morphism of sites \emph{(IV
\hyperref[IV.subsubsec4.9.1]{4.9.1})}. Then the morphism of topoi
$f:C^{\sim}\to{C'}^{\sim}$ defined by $g$ is coherent.
\end{enonce*}

Indeed,\pageoriginale $\epsilon_C(C)$ (\resp
$\epsilon_{C'}(C')$) is a generating family of $C^{\sim}$ (\resp
${C'}^{\sim}$) satisfying the corresponding conditions of
\hyperref[VI.prop3.2]{3.2} (\hyperref[VI.cor2.1.1]{2.1.1}).

\begin{enonce*}{Proposition 3.4}\etiquette[3.4]{VI.prop3.4}
Let $E$ be a locally coherent topos (\hyperref[VI.def2.3]{2.3}), and
let $f:X\to Y$ be an arrow of $E$, whence a morphism of induced topoi
\emph{(IV \hyperref[IV.eq5.5.2]{(5.5.2)})} $E_{/X}\to E_{/Y}$. This morphism of topoi is
quasi-compact (\resp quasi-separated, \resp coherent) if and only if
$f$ is a quasi-compact (\resp quasi-separated, \resp coherent)
morphism.
\end{enonce*}

The quasi-compact case follows immediately from the definitions
(without any condition on $E$); the quasi-separated case is precisely
\hyperref[VI.cor2.8]{2.8}; and the coherent case follows by combining
the two preceding cases.

\begin{enonce*}{Proposition 3.5}\etiquette[3.5]{VI.prop3.5}
Let $f:E'\to E$ be a morphism of locally coherent topoi. Let
$(X_i)_{i\in I}$ be a generating family in $E$ formed by coherent
algebraic objects. The following conditions are equivalent:
\begin{itemize}
\item[(i)] For every quasi-compact arrow $u$ of $E$, $f^{\ast}(u)$ is
  quasi-compact.

\item[(i bis)] For every arrow $u:X_i\to X_j$, $f^{\ast}(u)$ is
  quasi-compact.

\item[(ii)] For every coherent object $Y'$ of $E'$, every coherent
  algebraic object $Y$ of $E$, and every arrow
  $v:Y'\to f^{\ast}(Y)$, the corresponding morphism of topoi
$$
f_v:{E'}_{/Y'}\to E_{/Y}
$$
obtained by composing\pageoriginale the localization morphism
${E'}_{/Y'}\to{E'}_{/f^{\ast}(Y)}$ deduced from $v$
\emph{(IV \hyperref[IV.eq5.5.2]{(5.5.2)})} with the morphism
${E'}_{/f^{\ast}(Y)}\to E_{/Y}$ induced by $f$
\emph{(IV \hyperref[IV.eq5.10.1]{(5.10.1)})} is coherent.

\item[(ii bis)] The same condition as in \emph{(ii)}, but restricted
  to a collection of data
  $v_{\alpha}:{Y'}_{\alpha}\to f^{\ast}(Y_{\alpha})$ such that the
  ${Y'}_{\alpha}$ are algebraic and cover the final object of $E'$.

\item[(ii ter)] (When a generating family
  $(X'_{i'})_{i'\in I'}$ in $E'$ formed by coherent objects is given,
  together with, for every $i'\in I'$, an $i\in I$ and an arrow
  $v_{i'}:X'_{i'}\to f^{\ast}(X_i)$.) For every $i'\in I'$, the
  morphism of topoi $E'_{/{X'}_{i'}}\to E_{/X_i}$ defined by
  $v_{i'}$ is coherent.
\end{itemize}
\end{enonce*}

Since every arrow $X_i\to X_j$ is quasi-compact, it is clear that (i)
$\Rightarrow$ (i bis). Conversely, suppose (i bis) holds and let us
prove (i). Let $u:X\to Y$ be a quasi-compact arrow of $E$. In terms
of an epimorphic family $X_i\to Y$, the hypothesis on $u$ says that
the $\fprod{X}{X_i}{Y}$ are quasi-compact
(\hyperref[VI.cor1.16]{1.16}), \ie that for every $i$ there exists a
finite epimorphic family of morphisms
$X_j\to\fprod{X}{X_i}{Y}$ (\hyperref[VI.prop1.3]{1.3}). We then have
the corresponding epimorphic family
$f^{\ast}(X_i)\to f^{\ast}(Y)$, such that for every $i$ there is a
finite epimorphic family
$f^{\ast}(X_j)\to f^{\ast}(\fprod{X}{X_i}{Y})\simeq
\fprod{f^{\ast}(X)}{f^{\ast}(X_i)}{f^{\ast}(Y)}$, whose composite
with the projection $\pr_2$ to $f^{\ast}(X_i)$ is a quasi-compact
morphism $f^{\ast}(X_j)\to f^{\ast}(X_i)$ by hypothesis (i bis). It
follows (\hyperref[VI.cor1.11]{1.11} (i)) that, for every $i$,
$\pr_2:\fprod{f^{\ast}(X)}{f^{\ast}(X_i)}{f^{\ast}(Y)}$ is
quasi-compact, and hence (\hyperref[VI.prop1.10]{1.10} (ii)) that
$f^{\ast}(u):f^{\ast}(X)\to f^{\ast}(Y)$ is quasi-compact.

It remains\pageoriginale to prove (i) $\Rightarrow$ (ii) and
(ii bis) $\Rightarrow$ (i); the implication (ii) $\Rightarrow$
(ii bis) is immediate, and (ii ter) is a special case of (ii bis).
The implication (i) $\Rightarrow$ (ii) is immediate: indeed, an
object $X$ of $E_{/Y}$ is quasi-compact (\resp quasi-separated) if
and only if its structural morphism $X\to Y$ (\resp its diagonal
morphism $X\to\fprod{X}{X}{Y}$) is quasi-compact
(\hyperref[VI.cor2.7]{2.7}), and the same criterion of
quasi-compactness and quasi-separatedness applies to
$f^{\ast}_Y(X')$. Finally, let us prove (ii bis) $\Rightarrow$ (i).
Let $u:X\to Y$ be a quasi-compact arrow in $E$; we must show that
$f^{\ast}(u)$ is quasi-compact. Since the $Y'_{\alpha}$ cover the
final object of $E'$, it is enough to prove that the
$f^{\ast}(u)|Y'_{\alpha}$ are quasi-compact
(\hyperref[VI.prop1.10]{1.10} (ii)). These morphisms are precisely
the $f^{\ast}_{v_\alpha}(u|Y_{\alpha})$, so we are reduced to proving
the following result.

\begin{enonce*}{Corollary 3.6}\etiquette[3.6]{VI.cor3.6}
Let $f:E'\to E$ be a coherent morphism of locally coherent topoi.
Then, for every quasi-compact arrow $u$ of $E$, $f^{\ast}(u)$ is
quasi-compact.
\end{enonce*}

By the implication (i bis) $\Rightarrow$ (i) of
\hyperref[VI.prop3.5]{3.5} already proved, it is enough to show that,
if $u:X\to Y$ is an arrow of $E$ with $X,Y$ coherent, then
$f^{\ast}(u)$ is a quasi-compact morphism. This follows immediately
from the hypothesis on $f$, which implies that $f^{\ast}(X)$ and
$f^{\ast}(Y)$ are coherent as well.

\begin{enonce*}{Definition 3.7}\etiquette[3.7]{VI.def3.7}
Let $f:E'\to E$\pageoriginale be a morphism of locally coherent
topoi. We say that $f$ is \emph{locally coherent} if it satisfies the
equivalent conditions of \hyperref[VI.prop3.5]{3.5}.
\end{enonce*}

Observe that ``locally'' means locally on the source.

\setcounter{subsection}{7}
\setcounter{subsubsection}{0}
\subsubsection{}\etiquette[3.7.1]{VI.subsubsec3.7.1}
This notion again depends only on the isomorphism class of the
morphism of topoi $f$, and it is plainly stable under composition of
morphisms of topoi. Moreover, by \hyperref[VI.prop3.5]{3.5}, if $f$
is coherent then it is locally coherent; the converse holds if $E'$
and $E$ are coherent (by \hyperref[VI.prop3.5]{3.5} (ii bis)).
Criterion \hyperref[VI.prop3.5]{3.5} (ii bis) immediately gives the
following criterion:

\begin{enonce*}{Corollary 3.7}\etiquette[3.7]{VI.cor3.7}
Let $f:E'\to E$ be a morphism of locally coherent topoi, let
$(Y_i)_{i\in I}$ and $({Y'}_i)_{i\in I}$ be two families of objects
of $E$ and $E'$, and, for every $i\in I$, let
$v_i:Y'_i\to f^{\ast}(Y_i)$ be a morphism, whence a morphism of topoi
$$
f_{v_i}:E'_{/Y'_i}\to E_{/Y_i}{\quoi}.
$$
Suppose that the ${Y'}_i$ cover the final object of $E'$. Then $f$ is
locally coherent if and only if, for every $i\in I$, $f_{v_i}$ is
locally coherent.
\end{enonce*}

In particular, applying this to the identity morphism of an induced
topos gives:

\begin{enonce*}{Corollary 3.8}\etiquette[3.8]{VI.cor3.8}
Let $E$ be a locally coherent topos, and let $v:X\to Y$ be an arrow of
$E$. Then the morphism of induced topoi $E_{/X}\to E_{/Y}$ defined by
$v$ is locally coherent.
\end{enonce*}

\begin{enonce*}{Remark 3.9}\etiquette[3.9]{VI.rem3.9}
It seems\pageoriginale that all morphisms of algebraic topoi
encountered in practice are locally coherent (which is why the term
``locally coherent'' is probably not destined to be used very often).
Equivalently, the morphisms between coherent topoi encountered in
practice are coherent. Nevertheless, one can construct noncoherent
morphisms between coherent topoi and, more particularly, a
noncoherent morphism from the punctual topos \emph{(IV \hyperref[IV.subsec2.2]{2.2})} to a
Noetherian topos $E$, \emph{\cf \hyperref[VI.exer2.17]{2.17} i)}.
\end{enonce*}

\begin{enonce*}[remark]{Examples 3.10}\etiquette[3.10]{VI.exem3.10}
(\emph{Topoi associated with schemes}.)
Return to the example \hyperref[VI.subsec1.22]{1.22} of the category
$(\Sch)$ with the topologies $T_i$ of \SGA IV 6.3. For every scheme
$X$, let $X_{T_i}$ be the $𝒱$-topos defined by the induced site
$(\Sch)_{/X}$. A morphism of schemes $f:X\to Y$ then defines a
morphism of topoi $f_{T_i}:X_{T_i}\to Y_{T_i}$, and it follows from
\hyperref[VI.prop3.4]{3.4} and \hyperref[VI.subsec1.22]{1.22} that
this morphism of topoi is quasi-compact (\resp quasi-separated,
\resp coherent) if and only if the morphism of schemes $f$ is
quasi-compact (\resp quasi-separated, \resp coherent) in the usual
sense of \EGA IV 1. In any event, $f_{T_i}$ is locally coherent by
\hyperref[VI.cor3.8]{3.8}.

One may also associate with a scheme $X$ the $𝒰$-topoi $\Top(X)$ and
$\Top(X_{\text{ét}})$ as in \hyperref[VI.subsubsec1.22.1]{1.22.1},
and to every morphism of schemes $f:X\to Y$ there are then associated
morphisms $\Top(f)$ and $\Top(f_{\text{ét}})$ of the corresponding
topoi. Let $T(f)$ denote either of these two morphisms of
topoi.\pageoriginale One verifies immediately by
\hyperref[VI.prop3.2]{3.2} that this morphism of topoi is always
locally coherent, and that it is quasi-compact (\resp quasi-separated,
\resp coherent) if and only if the morphism of schemes $f$ is
quasi-compact (\resp quasi-separated, \resp coherent) in the usual
sense.

These observations show once again that the terminology introduced
in the present section is compatible with the accepted terminology
of scheme theory, and therefore deserves to be accepted by even the
most recalcitrant reader.
\end{enonce*}

\begin{enonce*}[remark]{Exercise 3.11}\etiquette[3.11]{VI.exer3.11}
(\emph{Coherent topoi and pretopoi}.)
\begin{enumerate}
\renewcommand{\theenumi}{\alph {enumi}}
\renewcommand{\labelenumi}{\theenumi)}
\item A $𝒰$-pretopos (or simply a pretopos) is a category $C$
  satisfying the following conditions:
\begin{enumerate}
\renewcommand{\theenumii}{\arabic{enumii}}
\renewcommand{\labelenumii}{\theenumii)}
\item Finite projective limits are representable in $C$.

\item Finite coproducts are representable in $C$, and they are
  disjoint and universal.

\item Equivalence relations in $C$ are effective, and every
  epimorphism in $C$ is universally effective.

\item $C$ is equivalent to a category $\in 𝒰$.
\end{enumerate}
Show that if $E$ is a coherent $𝒰$-topos, then the full subcategory
$E_{\coh}$ of $E$ formed by its coherent objects is a $𝒰$-pretopos,
and that the inclusion functor $E_{\coh}\to E$ is left exact,
commutes with\pageoriginale finite coproducts, and commutes with
passage to the quotient by an equivalence relation. (Use
\hyperref[VI.subsubsec1.5.3]{1.5.3}, \hyperref[VI.cor1.15]{1.15},
and \hyperref[VI.cor1.17.1]{1.17.1}.) Show that the topology induced
by $E$ on $C$ is the ``precanonical'' topology, \ie the topology whose
covering families $X_i\to X$ are those that admit a \emph{finite}
subfamily covering for the canonical topology of $C$. Consequently,
$E$ can be recovered up to equivalence from the pretopos
$C=E_{\coh}$ as the topos $C^{\sim}$ (with $C$ endowed with its
precanonical topology).

\item Let $C$ be a $𝒰$-pretopos. Endow $C$ with the precanonical
  topology. Show that every morphism $f$ of $C$ factors as $f''f'$,
  with $f'$ an epimorphism and $f''$ a monomorphism. Show that the
  topos $E=C^{\sim}$ is coherent and that the canonical functor
  $\epsilon:C\to E$ induces an equivalence from $C$ to the
  subcategory $E_{\coh}$ of $E$. (First show that $\varepsilon$ is a
  fully faithful left exact functor commuting with finite coproducts
  and with passage to the quotient by an equivalence relation, and
  then that a coherent subobject in $E$ of an object of $C$ lies in
  the essential image of $C$.) Consequently, the $𝒰$-pretopos $C$
  can be recovered up to equivalence of categories from its associated
  $𝒰$-topos $E=C^{\sim}$.

\item Let $C$ and $C'$ be two $𝒰$-pretopoi. Show that a functor
  $\varphi:C'\to C$ is a morphism of sites from $C$ to $C'$ (for the
  precanonical topologies) if and only if $\varphi$ is left exact and
  commutes with finite coproducts and with passage to a quotient by an
  equivalence relation.

\item With the\pageoriginale notation of c), suppose that $C$ and
  $C'$ are associated with two coherent $𝒰$-topoi $E,E'$ as in a).
  Show that the canonical functor (IV \hyperref[IV.subsubsec4.9.3]{4.9.3})
$$
\Mmorsite(C,C')\to\Hhomtop(E,E')
$$
induces an equivalence from its source (described explicitly in c))
to the full subcategory of its target formed by the \emph{coherent}
morphisms of topoi $E\to E'$. A quasi-inverse is obtained by
associating with every coherent morphism of topoi $f:E\to E'$ the
functor $E'_{\coh}=C'\to E_{\coh}=C$ induced by $f^{\ast}$.

\item Let $C$ be a $𝒰$-pretopos and $E=C^{\sim}$. Express directly in
  terms of exactness properties of $C$ the equivalent conditions of
  \hyperref[VI.cor1.25]{1.25} (consult
  \hyperref[VI.exer1.28]{1.28} b). Show that an object $X$ of $C$ is
  Noetherian in $E$ if and only if every increasing sequence of
  subobjects of $X$ in $C$ is stationary, and hence that $E$ is
  Noetherian if and only if every object of $C$ satisfies the
  preceding condition.
\end{enumerate}
\end{enonce*}

\begin{enonce*}[remark]{Exercise 3.12}\etiquette[3.12]{VI.exer3.12}
A topos $E$ is called a \emph{finite} topos if there exists a finite
category $C$ such that $E$ is equivalent to $\widehat{C}$.
\begin{enumerate}
\renewcommand{\theenumi}{\alph{enumi}}
\renewcommand{\labelenumi}{\theenumi)}
\item The topos $E$ is finite if and only if $E$ has enough essential
  points (IV \hyperref[IV.exer7.6]{7.6} b)) and the category $\Ppointess(E)$ of essential
  points of $E$ is equivalent to a finite category. (Use IV \hyperref[IV.exer7.6]{7.6} d)
  and h).)

\item Suppose\pageoriginale that $E$ is finite. Prove that every point
  of $E$ is essential. Prove that the 2-functors
  $E\mapsto\Ppoint(E)$ and $C\mapsto\widehat{C}$ establish
  2-equivalences between the 2-category of finite topoi $E$ and the
  2-category of Karoubian categories (IV \hyperref[IV.exer7.5]{7.5} a)) $C$ that are
  equivalent to finite categories, and that every morphism of finite
  topoi is essential (IV \hyperref[IV.exer7.6]{7.6} a)). (Use a) and IV \hyperref[IV.exer7.6]{7.6} h).)

\item Suppose that $E=\widehat{C}$, with $C$ equivalent to a finite
  category, and let $F$ be a coherent topos. Recall (IV \hyperref[IV.eq4.6.3.1]{(4.6.3.1)})
  that morphisms of topoi $f:E\to F$ correspond to functors
  $u:C\to\Ppoint(F)$. Show that $f$ is coherent if and only if $f$
  carries every coherent point of $E$ (\hyperref[VI.def3.1]{3.1}) to
  a coherent point of $F$, or equivalently if, for every $X\in\ob C$,
  $u(X)$ is a \emph{coherent} point of $F$. (Use
  \hyperref[VI.prop3.2]{3.2} and \hyperref[VI.exer2.17]{2.17} g).) Conclude that every
  morphism $f:E\to F$ from a finite topos to a finite topos is
  coherent.

\item Let $X$ be a sober topological space (IV \hyperref[IV.subsubsec4.2.1]{4.2.1}). The topos
  $\Top(X)$ (IV \hyperref[IV.subsec2.1]{2.1}) is finite if and only if $X$ is a finite set.
  (Use IV \hyperref[IV.subsubsec7.1.6]{7.1.6}.)
\end{enumerate}
\end{enonce*}


% END INCLUDED SGA4_English_current_buildable_source_20260906/source/06/drafts/SGA4_VI_LINES_2074_2418_English.texfrag


% BEGIN INCLUDED SGA4_English_current_buildable_source_20260906/source/06/drafts/SGA4_VI_LINES_2419_2936_English.texfrag
\section{Finiteness conditions in a topos obtained by
  gluing}\etiquette[4]{VI.sec4}

The present section will not be used later in the Seminar.

\subsection{}\etiquette[4.1]{VI.subsec4.1}
Let $E$ be a topos and $U$ an open of $E$ (\ie a subobject of the final
object of $E$), and consider the corresponding open and closed subtopoi
of\pageoriginale $E$ (IV 9)
$$
E'= E_{/U}, E'' = E_{/U}{\quoi}.
$$
We denote by
$$
j : E'\to E, i : E''\to E
$$
the canonical morphisms of topoi. Recall (\hyperref[VI.prop3.4]{3.4}) that
for $j$ to be a quasi-compact (\resp coherent) morphism of topoi, it is
necessary and sufficient that the inclusion
$$
j : U \to e
$$
of $U$ into the final object of $E$ (an inclusion that we shall also denote
by $j$) be a quasi-compact (\resp coherent) morphism in $E$; moreover,
$j:E'\to E$ is always quasi-separated (because $j:U\to e$ is so
(\hyperref[VI.cor1.8.1]{1.8.1})). We propose to give criteria for the topos
$E''$ to be coherent and for the morphism of topoi $i:E''\to E$ to be
coherent. First note:

\begin{enonce*}{Proposition 4.2}\etiquette[4.2]{VI.prop4.2}
With the notation of \hyperref[VI.subsec4.1]{4.1}, the inclusion morphism
$i:E''\to E$ is quasi-compact, \ie for every quasi-compact object $X$ of
$E$, the object $i^{\ast}(X)$ of $E''$ is quasi-compact; moreover, if $X$ is
pre-Noetherian (\hyperref[VI.def1.30]{1.30}), then $i^{\ast}(X)$ is
pre-Noetherian.
\end{enonce*}

This follows from \hyperref[VI.def1.1]{Definition 1.1} and the

\begin{enonce*}{Lemma 4.2.1}\etiquette[4.2.1]{VI.lem4.2.1}
The map\pageoriginale $Y\mapsto i^{\ast}(Y)$ induces an isomorphism of
ordered sets between the set of subobjects of $X$ containing the subobject
$X_U=X\times U$ of $X$, and the set of subobjects of $i^{\ast}(X)$.
\end{enonce*}

Note that since $i_{\ast}$ is conservative (because it is fully faithful)
and left exact (in particular, it commutes with fiber products), it follows
immediately that a morphism $u:Y\to Z$ in $E''$ is a monomorphism if and only
if $i_{\ast}(u):i_{\ast}(Y)\to i_{\ast}(Z)$ is one. Consequently, on
identifying (by $i_{\ast}$) $E''$ with the full subcategory of $E$ consisting
of objects $Z$ satisfying the equivalent conditions of IV \hyperref[IV.subsec9.7]{9.7} 1), the
subobjects in $E''$ of the object $Z$ of $E''$ identify with the subobjects
$Y$ of $Z$ in $E$ that belong to $E''$, or, equivalently, that contain
$Z_U\xrightarrow{\sim}U$. For $Z$ of the form
$i_{\ast}i^{\ast}(X)=X_{\complement U}$ given by the amalgamated sum
IV \hyperref[IV.subsubsec9.5.1]{(9.5.1)},
$$
i_{\ast} i^{\ast} (X) = X_{\complement U} = X \amalg_{X_U} U {\quoi},
$$
the datum of a subobject $Y$ of the latter in $E$ is equivalent to the datum
of a pair consisting of a subobject $Y_1$ of $X$ and a subobject $Y_2$ of
$U$, subject to the condition that the inverse images of these subobjects
in $X_U$ coincide. The condition that $Y$ contain $U$ then means that
$Y_2=U$, and the remaining condition on $Y_1$ is that it contain $X_U$.
This therefore establishes a bijection between the set of subobjects of
$i^{\ast}(X)$ in $E''$ and the set of subobjects of $X$ that contain $X_U$.
It is clear\pageoriginale that this is an isomorphism of ordered sets and
that its inverse map is indeed the one asserted in \hyperref[VI.lem4.2.1]{4.2.1}.

\begin{enonce*}{Corollary 4.3}\etiquette[4.3]{VI.cor4.3}
\begin{enumerate}
\renewcommand{\theenumi}{\alph{enumi}}
\renewcommand{\labelenumi}{\theenumi)}
\item Let $X$ be an object of $E$. For $X$ to be quasi-compact, it is
  sufficient that $i^{\ast}(X)$ and $j^{\ast}(X)$ be so, and this condition
  is also necessary if $j:U\to e$ is quasi-compact.

\item Let $Y$ be an object of $E''$. For $Y$ to be quasi-compact, it is
  sufficient that $i_{\ast}(Y)$ be so, and this condition is also necessary
  if $U$ is quasi-compact.
\end{enumerate}
\end{enonce*}

\noindent
\underline{Proof}
\begin{enumerate}
\renewcommand{\theenumi}{\alph{enumi}}
\renewcommand{\labelenumi}{\theenumi)}
\item Sufficiency follows from \hyperref[VI.def1.1]{Definition 1.1} and the fact that the pair
  $(i^{\ast},j^{\ast})$ is conservative (IV \hyperlink{sga4.exp.IV.proposition.9.4.1.item.c}{9.11} 3)). Necessity follows from
  the fact that $i$ and $j$ are quasi-compact (by \hyperref[VI.prop4.2]{3.14} and the hypothesis
  that $j$ is quasi-compact).

\item Since $Y\simeq i^{\ast}i_{\ast}(Y)$, sufficiency follows from
  \hyperref[VI.prop4.2]{4.2}. Necessity follows from the sufficiency in a),
  taking into account that $i^{\ast}i_{\ast}(Y)\simeq Y$ and
  $j^{\ast}i_{\ast}(Y)\simeq U$.
\end{enumerate}

\begin{enonce*}{Corollary 4.4}\etiquette[4.4]{VI.cor4.4}
Let $Y$ be an object of $E''$. If $U$ is quasi-compact or if $E$ admits a
generating family consisting of quasi-compact objects, then for $Y$ to be
quasi-separated (\resp coherent), it is \emph{sufficient} that
$i_{\ast}(Y)$ be so. If $U$ is quasi-separated (\resp coherent) and if
$j:U\to e$ is a quasi-compact morphism, then for $Y$ to be quasi-separated
(\resp coherent), it is necessary that $i_{\ast}(Y)$ be so.
\end{enonce*}

The parenthetical\pageoriginale case follows from the non-parenthetical
case, taking \hyperref[VI.cor4.3]{4.3} b) into account. Suppose that
$i_{\ast}(Y)$ is quasi-separated, and let us prove that $Y$ is so under one
of the two stated hypotheses. Thus we must prove that, for two quasi-compact
objects $Y'$ and $Y''$ over $Y$, the fiber product
$\fprod{Y'}{Y''}{Y}$ is quasi-compact. When $U$ is assumed quasi-compact,
note that it is enough to prove that
$i_{\ast}(\fprod{Y'}{Y''}{Y})$ is quasi-compact
(\hyperref[VI.cor4.3]{4.3} b)); this follows from the fact that
$i_{\ast}(Y')$ and $i_{\ast}(Y'')$ are so
(\hyperref[VI.cor4.3]{4.3} b)), using here the quasi-compactness of $U$, the
fact that $i_{\ast}$ commutes with fiber products, and the hypothesis that
$i_{\ast}(Y)$ is quasi-separated. When $E$ is assumed to admit a generating
family consisting of quasi-compact objects, $i_{\ast}(Y')$ is a filtered
inductive limit of its quasi-compact subobjects ${X'}_{\alpha}$; hence
$Y'\simeq i^{\ast}i_{\ast}(Y')$ is a filtered inductive limit of the
subobjects $i^{\ast}({X'}_{\alpha})$, and since $Y'$ is quasi-compact, it is
equal to one of the $i^{\ast}({X'}_{\alpha})$. Likewise, $Y''$ is of the
form $i^{\ast}({X''}_{\beta})$, where ${X''}_{\beta}$ is a quasi-compact
subobject of $i_{\ast}(Y'')$. But then
$\fprod{Y'}{Y''}{Y}=
i^{\ast}(\fprod{X'_{\alpha}}{X''_{\beta}}{i_{\ast}Y})$, and since
$\fprod{X'_{\alpha}}{X''_{\beta}}{i_{\ast}(Y)}$ is quasi-compact by the
hypothesis that $i_{\ast}(Y)$ is quasi-separated, it follows by
\hyperref[VI.prop4.2]{4.2} that $\fprod{Y'}{Y''}{Y}$ is also quasi-compact.

Conversely, suppose that $U$ is quasi-separated and $j:U\to e$ is
quasi-compact, and assume that $Y$ is quasi-separated. Let us show that
$i_{\ast}(Y)$ is also quasi-separated, \ie that for two quasi-compact
objects $X',X''$ over $i_{\ast}(Y)$, the fiber product
$\fprod{X'}{X''}{i_{\ast}(Y)}$ is quasi-compact. By
\hyperref[VI.cor4.3]{4.3} a), it is enough to prove that its image under
$i^{\ast}$ is quasi-compact\pageoriginale (which follows from
\hyperref[VI.prop4.2]{4.2} and the hypothesis that $Y$ is quasi-separated)
\emph{and} that its image under $j^{\ast}$ is quasi-compact. The latter is
$j^{\ast}(X')\times j^{\ast}(X'')$, and is indeed quasi-compact because
$j^{\ast}(X')$ and $j^{\ast}(X'')$ are so ($j$ being quasi-compact) and $U$
is quasi-separated.

\begin{enonce*}{Corollary 4.5}\etiquette[4.5]{VI.cor4.5}
Suppose that $E$ admits a generating subcategory consisting of quasi-compact
(\emph{\resp} pre-Noetherian) objects. Then the same is true of $E''$.
\end{enonce*}

This follows from \hyperref[VI.prop4.2]{4.2} and the

\begin{enonce*}{Lemma 4.5.1}\etiquette[4.5.1]{VI.lem4.5.1}
If $(X_{\alpha})_{\alpha}$ is a generating family in $E$, then
$(i^{\ast}(X_{\alpha}))_{\alpha}$ is a generating family in $E''$.
\end{enonce*}

Indeed, this means that the family of functors
$$
Y \to \Hom (i^{\ast}(X_{\alpha}), Y)
$$
on $E''$ is conservative. But
$\Hom(i^{\ast}(X_{\alpha}),Y)\simeq\Hom(X_{\alpha},i_{\ast}(Y))$, and it is
enough to use the fact that the functor $i_{\ast}$ is conservative.

\begin{enonce*}{Proposition 4.6}\etiquette[4.6]{VI.prop4.6}
The notation is that of \hyperref[VI.subsec4.1]{4.1}.
\begin{enumerate}
\renewcommand{\theenumi}{\alph{enumi}}
\renewcommand{\labelenumi}{\theenumi)}
\item If $E$ admits a generating family consisting of coherent objects,
  then the same is true of the closed subtopos $E''$, and the inclusion
  morphism $i:E''\to E$ is coherent.

\item Suppose that the morphism $j:U\to e$ in $E$ is quasi-compact.
If\pageoriginale $E$ is locally coherent (\emph{\resp} coherent,
\emph{\resp} quasi-separated, \emph{\resp} locally Noetherian,
\emph{\resp} Noetherian, \emph{\resp} locally perfect, \emph{\resp}
perfect), then the same is true of $E''$, and the inclusion morphism
$i:E''\to E$ is coherent.
\end{enumerate}
\end{enonce*}

\begin{enumerate}
\renewcommand{\theenumi}{\alph{enumi}}
\renewcommand{\labelenumi}{\theenumi)}
\item The first assertion will follow from the second and
  \hyperref[VI.lem4.5.1]{4.5.1}. In any event,
  \hyperref[VI.cor4.5]{4.5} implies that $E''$ admits a generating family
  consisting of quasi-compact objects, so \hyperref[VI.prop3.2]{3.2}
  applies and shows that it is enough to verify that, for every coherent
  object $X$ of $E$, the object $i^{\ast}(X)$ of $E''$ is coherent. Using
  the compatibility of the formation of the topos complementary to an open
  with localization (IV \hyperref[IV.cor9.4.3]{9.13} b)), we reduce to the case in which $X$ is the
  final object of $E$, and hence to proving that the final object of $E''$
  is coherent. Now \hyperref[VI.cor4.4]{4.4} applies, so it is enough to
  prove that $i_{\ast}(e_{E''})=e_E$ is coherent, which is indeed the case.

\item Under each of the hypotheses made in b), $E$ admits a generating
  family consisting of coherent objects, so it already follows from a) that
  $i:E''\to E$ is coherent. Moreover, \hyperref[VI.lem4.5.1]{4.5.1} then
  implies that $E''$ admits as a generating subcategory the full subcategory
  ${E''}_{\coh}$ consisting of its coherent objects. Suppose that $E$ is
  coherent, and let us show that $E''$ is also coherent, \ie
  (\hyperref[VI.subsubsec2.4.5]{2.4.5}) that the subcategory
  ${E''}_{\coh}$ of $E''$ is stable under finite projective limits, knowing
  that this is so for the subcategory $E_{\coh}$ of $E$. This follows
  immediately from the parenthetical criterion
  \hyperref[VI.cor4.4]{4.4} (which applies because $U$ is now coherent, $e$
  being\pageoriginale coherent and $j:U\to e$ being quasi-compact), taking
  into account that $i_{\ast}$ is left exact. By localization, again using
  IV \hyperref[IV.cor9.4.3]{9.13} b), we conclude that if $E$ is locally coherent, then the same is
  true of $E''$. If $E$ is quasi-separated, then $E''$ is also
  quasi-separated: indeed, its final object is quasi-separated by
  \hyperref[VI.cor4.4]{4.4}, since that of $E$ is, and it remains to prove
  that the product of two quasi-separated objects of $E''$ is
  quasi-separated (knowing that this is so in $E$); this again follows from
  the criterion \hyperref[VI.cor4.4]{4.4} and the fact that $i_{\ast}$
  commutes with products (taking into account that the subobject $U$ of $e$
  is quasi-separated, since $e$ is, so that \hyperref[VI.cor4.4]{4.4}
  applies).
\end{enumerate}

Since a topos is locally Noetherian (\resp Noetherian) if and only if it is
locally coherent (\resp coherent) and admits a generating family consisting
of pre-Noetherian objects (\hyperref[VI.prop2.10]{2.10} (iii bis)), it
follows from the preceding discussion and \hyperref[VI.cor4.5]{4.5} that if
$E$ is locally Noetherian (\resp Noetherian), then the same is true of
$E''$. Now suppose that $E$ is perfect, and let us prove that $E''$ is
perfect. Since we already know that it is coherent, it remains to verify
that it satisfies criterion \hyperref[VI.cor1.25]{1.25} (i), \ie that every
object $Y$ of $E''$ is a filtered inductive limit of coherent objects,
knowing that the analogous statement is true in $E$. Since $i_{\ast}(Y)$ is
a filtered inductive limit of coherent objects $X_{\alpha}$,
$Y\simeq i^{\ast}i_{\ast}(Y)$ is a filtered inductive limit of the coherent
objects $i^{\ast}(X_{\alpha})$. By localization (IV \hyperref[IV.cor9.4.3]{9.13} b)), we conclude
that if $E$ is locally perfect, then the same is true of $E''$. This
completes the proof of \hyperref[VI.prop4.6]{4.6}.

\begin{enonce*}[remark]{Remark 4.6.1}\etiquette[4.6.1]{VI.rem4.6.1}
In fact,\pageoriginale in \hyperref[VI.prop4.6]{4.6} b), the hypothesis that
$j:U\to e$ is quasi-compact is unnecessary, except perhaps in the case
``$E$ quasi-separated''. Indeed, by the preceding proof it is enough to see
that if $E$ is coherent, then $E''$ is coherent. Now $U$ is the filtered
inductive limit of its quasi-compact subobjects $U_{\alpha}$, and in
\hyperref[VI.sec7]{7} we shall verify that $E''$ then identifies with the
projective-limit topos (in the sense of \hyperref[VI.sec7]{7}) of the closed
subtopoi $E''_i$ complementary to the $E_{/U_i}$, which by
\hyperref[VI.prop4.6]{4.6} b) are coherent topoi with coherent transition
morphisms. It then follows that $E''$ is a coherent topos
(\hyperref[VI.sec7]{7}).
\end{enonce*}

\begin{enonce*}{Corollary 4.7}\etiquette[4.7]{VI.cor4.7}
Suppose that the subcategory $E_{\coh}$ of $E$ consisting of coherent
objects is generating, and let $X$ be an object of $E$.
\begin{enumerate}
\renewcommand{\theenumi}{\alph{enumi}}
\renewcommand{\labelenumi}{\theenumi)}
\item If $X$ is quasi-separated, then $i^{\ast}(X)$ and $j^{\ast}(X)$ are
  quasi-separated; this condition is also sufficient if $j:U\to e$ is
  quasi-compact.

\item Suppose that $j:U\to e$ is quasi-compact. Then $X$ is coherent if and
  only if $i^{\ast}(X)$ and $j^{\ast}(X)$ are coherent.

\item Suppose that $E$ is coherent and $j:U\to e$ quasi-compact. Then $E$
  is perfect (\hyperref[VI.def2.9.1]{2.9.1}) if and only if $E'$ and $E''$
  are perfect.
\end{enumerate}
\end{enonce*}

\begin{enumerate}
\renewcommand{\theenumi}{\alph{enumi}}
\renewcommand{\labelenumi}{\theenumi)}
\item We noted in \hyperref[VI.subsec4.1]{4.1} that $j:E'\to E$ is
  quasi-separated, and we know from \hyperref[VI.prop4.6]{4.6} a) that the
  same is true of $i$, which gives necessity. For sufficiency, let $X'$ and
  $X''$ be quasi-compact objects over $X$. We must prove that
  $\fprod{X'}{X''}{X}$ is quasi-compact, and for this it\pageoriginale is
  enough to prove that its images under $i^{\ast}$ and $j^{\ast}$ are so
  (\hyperref[VI.cor4.3]{4.3} a)). Since $i$ and $j$ are quasi-compact, by
  \hyperref[VI.prop4.2]{4.2} and the hypothesis that $j:U\to e$ is
  quasi-compact, the conclusion follows from the fact that $i^{\ast}$ and
  $j^{\ast}$ commute with fiber products and from the hypothesis that
  $i^{\ast}(X)$ and $j^{\ast}(X)$ are quasi-separated.

\item This follows by combining a) with \hyperref[VI.cor4.3]{4.3} a).

\item Necessity was already seen (\hyperref[VI.prop4.6]{4.6} b)).
  Sufficiency follows from the fact that, since $E$ is coherent, $E'$ and
  $E''$ are coherent (\hyperref[VI.prop4.6]{4.6} b)), so that for any of the
  topoi under consideration, being perfect means that the category of its
  coherent objects is stable under finite $\varinjlim$. We therefore
  conclude by b).
\end{enumerate}

\begin{enonce*}[remark]{Exercise 4.8}\etiquette[4.8]{VI.exer4.8}
\begin{enumerate}
\renewcommand{\theenumi}{\alph{enumi}}
\renewcommand{\labelenumi}{\theenumi)}
\item Solve the following question (to which the writer confesses, to his
  embarrassment, that he does not know the answer): with the notation of
  \hyperref[VI.subsec4.1]{4.1}, if $E$ is quasi-separated
  (\resp algebraic), is the same true of $E''$?

\item Suppose that $E$ is equivalent to a topos of the form $\widehat{C}$,
  where $C$ is a category equivalent to a category $\in 𝒰$. Prove directly,
  using \hyperref[VI.exer2.17]{2.17} c) and IV 9.24, that if $E$ is locally
  coherent (\resp coherent, \resp algebraic, \resp quasi-separated, \resp
  locally Noetherian, \resp Noetherian), then the same is true of $E''$.
\end{enumerate}
\end{enonce*}

\setcounter{subsection}{8}
\subsection{}\etiquette[4.9]{VI.subsec4.9}
Keep the notation of \hyperref[VI.subsec4.1]{4.1}, and recall (IV \hyperref[IV.theo9.5.4]{9.16}) that
the datum of a situation $(E,U)$ consisting of a topos $E$ and an open $U$
of $E$ is essentially equivalent to that of a triple $(E',E'',f)$, where
$E'$ and $E''$ are\pageoriginale topoi and
\begin{equation*}
f : E'\to E''\tag{4.9.1}
\end{equation*}\etiquette[4.9.1]{VI.eq4.9.1}
is a ``gluing functor'', \ie a functor that is \emph{left exact and
accessible}; starting from $(E,U)$ as in \hyperref[VI.subsec4.1]{4.1}, the
associated gluing functor is given by
\begin{equation*}
  f = i^{\ast} j_{\ast}{\quoi}.\tag{4.9.2}
\end{equation*}\etiquette[4.9.2]{VI.eq4.9.2}

We propose to express, in terms of the data $E'$, $E''$, and $f$, the fact
that $E$ is a coherent (\resp perfect) topos and that $U$ is a quasi-compact
object, or equivalently, that $E$ and $E'$ are coherent. We already know
that this implies that $E''$ is also coherent
(\hyperref[VI.prop4.6]{4.6} b)), so the question becomes the following:
given two \emph{coherent} topoi $E'$ and $E''$ and a gluing functor $f$
\eqhyperref[VI.eq4.9.1]{4.9.1}, under what conditions on $f$ is the glued
topos $E$ coherent (\resp perfect)? We also know
(\hyperref[VI.cor4.7]{4.7} c)) that $E$ is perfect if and only if $E$ is
coherent and $E'$ and $E''$ are perfect; thus the parenthetical case of the
problem reduces to the non-parenthetical case. Let us nevertheless note that
the solution (\hyperref[VI.prop4.10]{4.10}) of the latter is more elegant if
$E'$ is already assumed perfect.

\setcounter{subsection}{9}
\setcounter{subsubsection}{2}
\subsubsection{}\etiquette[4.9.3]{VI.subsubsec4.9.3}

First note that if $E$ is coherent, then the full subcategory $E_{\coh}$ of
$E$ consisting of the objects $X=(X',X'',u:X''\to f(X'))$ of $E$ that are
coherent is reconstructed as the subcategory $E_0$ of $E$ consisting of the
$X$ for\pageoriginale which $X'=j^{\ast}(X)$ and $X''=i^{\ast}(X)$ are
coherent (\hyperref[VI.cor4.7]{4.7} b)). A necessary condition for $E$ to be
coherent is therefore that the preceding subcategory $E_0$ be generating.
This condition is also sufficient, because by \hyperref[VI.cor4.3]{4.3} a),
$E_0$ consists of quasi-compact objects; moreover, it is clear that $E_0$ is
stable under finite projective limits, and we conclude by \hyperref[VI.subsubsec2.4.5]{3.4.5}.

Recall now (IV \hyperref[IV.subsec9.6]{9.18}) that the datum of a gluing functor
\eqhyperref[VI.eq4.9.1]{4.9.1} is equivalent to that of a \emph{sheaf} on
$E''$ with values in $\Pro(E')^0$:
\begin{equation*}
G \in \Ffaisc (E'', \Pro(E')^{0}), G : {E''}^{0} \to \Pro({E'}^{0})
\tag{4.9.4}{\quoi},
\end{equation*}\etiquette[4.9.4]{VI.eq4.9.4}
or equivalently to that of a functor $g=G^0$,
\begin{equation*}
g:E''\to \Pro (E')\tag{4.9.5}
\end{equation*}\etiquette[4.9.5]{VI.eq4.9.5}
that \emph{commutes with inductive limits}. If $C$ is a generating
subcategory of $E''$, endowed with the induced topology, we know
(II \hyperref[VI.subsec6.10]{6.10}) that the datum of $G$ is also equivalent
(up to unique isomorphism) to that of a sheaf on $C$ with values in
$\Pro(E')^0$
\begin{equation*}
G_C \in \Ffaisc (C, \Pro(E')^0), G_C : C^{\circ} \to \Pro (E')^{0}{\quoi},\tag{4.9.6}
\end{equation*}\etiquette[4.9.6]{VI.eq4.9.6}
or, equivalently, to that of a functor $g_C=G_C^0$,
\begin{equation*}
g_C : C \to \Pro (E'){\quoi},\tag{4.9.7}
\end{equation*}\etiquette[4.9.7]{VI.eq4.9.7}
satisfying\pageoriginale the familiar left-exactness conditions
(II \hyperlink{sga4.exp.II.remark.6.2.item.1}{6.2.1})). Of course, $g_C$ is simply the \emph{restriction of $g$
$\hyperref[VI.eq4.9.5]{(4.9.5)}$ to $C$}. The case of greatest interest to us is that in which we
take $C=E''_{\coh}$, giving the objects
\hypertarget{sga4.exp.VI.equation.4.9.8}{}\begin{align*}
& G_0 \in \Ffaisc (E''_{\coh}, \Pro (E')^{0}), G_0 : {E''}_{\coh}^0 \to
 \Pro (E')^0 {\quoi},\tag{4.9.8}\\
& g_0 = G^0_0 : E''_{\coh} \to \Pro(E'){\quoi},\tag{4.9.9}
\end{align*}\hypertarget{sga4.exp.VI.equation.4.9.9}{}
each of which is again equivalent to the datum of $f$.

\setcounter{subsubsection}{9}
\subsubsection{}\etiquette[4.9.10]{VI.subsubsec4.9.10}

Suppose that the full generating subcategory $C$ of $E''$ consists of
quasi-compact objects and is stable in $E''$ under finite sums, passage to
the quotient by equivalence relations, and fiber products; this is the case,
for example, for $C=E''_{\coh}$
(\hyperref[VI.cor1.15]{1.15} and \hyperref[VI.cor1.17.1]{1.17.1}). It is then
immediate, if $P$ is a category in which finite projective limits are
representable (for example $P=\Pro(E')^0$), that a functor
$G_C:C^{\circ}\to P$ is a sheaf with values in $P$ if and only if the functor
$g_C=G_C^0:C\to P^0$ commutes with finite sums and with passage to the
quotient by an equivalence relation. This specifies, in particular, which
objects are the sheaves \eqhyperref[VI.eq4.9.6]{4.9.6} (and thus expresses
the gluing functors $f:E'\to E''$).

With these reminders in place, we can give the solution to the problem posed
in \hyperref[VI.subsec4.9]{4.9}:

\begin{enonce*}{Proposition 4.10}\etiquette[4.10]{VI.prop4.10}
Let $E'$ and $E''$\pageoriginale be two coherent topoi, and let
$f:E'\to E''$ be a gluing functor \emph{(IV \hyperref[IV.eq9.3.6.3]{9.10})}, \ie a left exact and
accessible functor. Let $E$ be the topos deduced from it by gluing
\emph{(IV \hyperref[IV.theo9.5.4]{9.16})}, and denote by $E'_{\coh}$ (\resp $E'_{PF}$) the strictly
full subcategory of $E'$ consisting of the coherent objects (\resp the
objects $X$ such that the covariant functor $\Hom(X,-)$ represented by $X$
commutes with small filtered inductive limits). Consider the following
conditions:
\begin{itemize}
\item[(i)] $E$ is coherent.

\item[(ii)] For every object $X'$ of $E'$, there exists a family of
  morphisms
$$
v_{\alpha} : Y'_{\alpha} \to X'
$$
with target $X'$ and with coherent objects of $E'$ as sources, such that the
family of the
$$
f(v_{\alpha}) : f(Y'_{\alpha}) \to f(X')
$$
in $E''$ is covering.

\item[(ii bis)] $f$ commutes with (small) filtered inductive limits, and
  \emph{(ii)} holds for every $X'\in\Ob E'_{PF}$.

\item[(ii ter)] The canonical functor \emph{(IV \hyperref[IV.eq9.7.2.3]{9.20.1})}
\begin{equation*}
\pi : \Ppoint(E'') \to \Pro(E')\tag{4.10.1}
\end{equation*}\etiquette[4.10.1]{VI.eq4.10.1}
defined by the gluing functor $f$ factors (up to isomorphism) through
$\Pro(E'_{\coh})$ (via the fully faithful functor
$\Pro(E'_{\coh})\to\Pro(E')$ arising from the inclusion
$E'_{\coh}\to E'$):
\begin{equation*}
\pi_0 : \Ppoint (E'') \to \Pro(E'_{\coh}){\quoi}.\tag{4.10.2}
\end{equation*}\etiquette[4.10.2]{VI.eq4.10.2}

\item[(iii)] The\pageoriginale functor $f$ commutes with (small) filtered
  inductive limits.

\item[(iii bis)] The functor $g_0$ \hyperlink{sga4.exp.VI.equation.4.9.9}{(4.9.9)} factors (up to isomorphism)
  through a functor
\begin{equation*}
E''_{\coh} \to \Pro(E'_{PF}){\quoi},\tag{4.10.3}
\end{equation*}\etiquette[4.10.3]{VI.eq4.10.3}
via the fully faithful functor $\Pro(E'_{PF})\to\Pro(E')$ deduced from the
inclusion $E'_{PF}\to E'$.

\item[(iii ter)] The canonical functor $\pi$
  \eqhyperref[VI.eq4.10.1]{4.10.1} factors (up to isomorphism) as
\begin{equation*}
\pi_1 : \Ppoint(E'') \to \Pro(E'_{PF}){\quoi}.\tag{4.10.4}
\end{equation*}\etiquette[4.10.4]{VI.eq4.10.4}

\item[(iv)] The functor $g_0$ \hyperlink{sga4.exp.VI.equation.4.9.9}{(4.9.9)} factors (up to isomorphism) through a
  functor
\begin{equation*}
E''_{\coh} \to \Pro(E'_{\coh}){\quoi}.\tag{4.10.5}
\end{equation*}\etiquette[4.10.5]{VI.eq4.10.5}
We then have the following diagram of implications:
\begin{equation*}
\text{\rm (iv) $\Rightarrow$ (i) $\Leftrightarrow$ (ii) $\Leftrightarrow$ (ii bis)
  $\Leftrightarrow$ (ii ter) $\Rightarrow$ (iii) $\Leftrightarrow$ (iii bis)
  $\Leftrightarrow$ (iii ter)}{\quoi}.\tag{4.10.6}
\end{equation*}\etiquette[4.10.6]{VI.eq4.10.6}
\end{itemize}
\end{enonce*}

Let us immediately record the

\begin{enonce*}{Corollary 4.11}\etiquette[4.11]{VI.cor4.11}
\begin{enumerate}
\renewcommand{\theenumi}{\alph{enumi}}
\renewcommand{\labelenumi}{\theenumi)}
\item Suppose that $E'$ is perfect. Then all the conditions considered in
  \hyperref[VI.prop4.10]{4.10} are equivalent; in particular, $E$ is
  coherent if and only if $f$ commutes with filtered inductive limits, or
  equivalently if and only if the functor $g_0$ \hyperlink{sga4.exp.VI.equation.4.9.9}{(4.9.9)} factors (up to
  isomorphism)\pageoriginale through $\Pro(E'_{\coh})$.

\item For $E$ to be perfect, it is necessary and sufficient that $E'$ and
  $E''$ be perfect and that $f$ (\resp $g_0$) satisfy the condition stated
  in {\rm a)}.
\end{enumerate}
\end{enonce*}

Indeed, a) follows from the fact that $E'$ perfect means
$E'_{\coh}=E'_{PF}$, so that (iii bis) implies (iv). Assertion b) follows,
as is clear from the preliminary remarks of \hyperref[VI.subsec4.9]{4.9}.


\setcounter{subsection}{11}

% END INCLUDED SGA4_English_current_buildable_source_20260906/source/06/drafts/SGA4_VI_LINES_2419_2936_English.texfrag


% BEGIN INCLUDED SGA4_English_current_buildable_source_20260906/source/06/drafts/SGA4_VI_LINES_2937_3239_English.texfrag
\subsection{Proof of 4.10}\etiquette[4.12]{VI.subsec4.12}

\begin{enumerate}
\renewcommand{\theenumi}{\alph{enumi}}
\renewcommand{\labelenumi}{\theenumi)}
\item Let us spell out the condition equivalent to (i) obtained in
  \hyperref[VI.subsubsec4.9.3]{4.9.3}, namely that the full subcategory
  $E_0$ of $E$ is generating, \ie that for every object
  $X = (X', X'', u : X''\to f(X'))$ of $E$, the family of all
  morphisms
$$
v = (v', v'') : Y = (Y', Y'', v : Y''\to f(Y')) \to X {\quoi},
$$
with $Y \in \Ob E_0$, \ie with $Y'$ and $Y''$ coherent, is
epimorphic. Since the pair of functors $(i^{\ast}, j^{\ast})$ is
conservative, this also means that the family of corresponding
morphisms $Y' \to X'$ is epimorphic in $E'$, and that the family of
corresponding morphisms $Y''\to X''$ is epimorphic in $E''$. This is
clear for the first family, as is seen by taking $Y \in \Ob E_0$
\emph{over} $U$, \ie such that $Y'' = \varnothing_{E''}$: one obtains
the family of all arrows in $E'$ with target $X'$ and coherent source,
which is indeed epimorphic because $E'$ is coherent and hence the
subcategory ${E'}_{\coh}$ is generating. It therefore remains to
express the condition that the family $Y''\to X''$ be epimorphic for
every given object $X$ of $E$.

\item Now suppose\pageoriginale that condition (ii) is satisfied for
  the object $X'$ under consideration. It follows that one can find
  an epimorphic family of morphisms $Y''_{\beta} \to X''$ whose
  composites with $u : X''\to f(X')$ each lift to a morphism
  $Y''_{\beta} \to f({Y'_{\alpha}})$ for a suitable
  $\alpha = \phi(\beta)$. Since $E''_{\coh}$ is a generating
  subcategory of $E''$, we may moreover suppose the $Y''_{\beta}$
  coherent. The commutative diagrams
\[
\xymatrix@C=2.25cm@R=1.25cm{
Y''_{\beta}  \ar[r] \ar[d]& f(Y'_{\varphi(\beta)}) \ar[d]\\
X'' \ar[r] & f(X')
}
\]
then provide the required family of morphisms in $E$, giving an
epimorphic family $Y''_{\beta} \to X''$. Thus (ii) $\Rightarrow$ (i).

Conversely, suppose (i), and let us prove (ii). Apply the condition
made explicit in a) to the object
$$
X = j_{\ast}(X') = (X', f(X'), \id = f(X') \xrightarrow{\sim} f(X')){\quoi}.
$$
Since the $Y''\to X''\simeq f(X')$ form an epimorphic family, a
fortiori so does the family of $f(v') : f(Y') \to f(X')$. The
conclusion follows, since the $Y'$ are coherent by hypothesis. Thus
$$
{\rm(i)} \Leftrightarrow {\rm(ii)}{\quoi}.
$$

\item Note\pageoriginale now that condition (iv) also means that for
  every object $X''$ of $E''$, the pro-representable functor
$$
g(X'') : X'\longmapsto \Hom (X'', f(X'))
$$
on $E'$ is pro-represented by a pro-object whose components belong to
${E'}_{\coh}$. Equivalently, in the category $E'_{/g(X'')}$ formed by
pairs consisting of an object $X'$ of $E'$ and an element
$u \in g(X'')(X') =\Hom(X'', f(X'))$, $u : X''\to f(X')$, there is a
cofinal set of objects $(X'_{\alpha},u_{\alpha})$, with
$u_{\alpha} : X'' \to f(X'_{\alpha})$ and
$X'_{\alpha} \in \Ob E'_{\coh}$. Since the category
$E'_{/g(X'')}$ is filtered, the cofinality condition simply means
that for every object $(X',u)$, $u : X''\to f(X')$, of
$E'_{/g(X'')}$, there is a morphism from one of the
$(X'_{\alpha},u_{\alpha})$ to this object, \ie a morphism
$v'_{\alpha} : X'_{\alpha} \to X'$ making the triangle
\[\vcenter{
\xymatrix@C=2cm@R=1.5cm{
                      &f(X'_{\alpha}) \ar[d]^{f(v'_\alpha)}\\
X''\ar@{-}[ru]^{u_{\alpha}} \ar[r]& f(X').
}}\leqno{(*)}
\]
commutative. Thus condition (iv) is equivalent to the following:

\begin{itemize}
\item[(iv bis)] For every arrow $u : X''\to f(X')$, with
  $X'\in \Ob E'$ and $X''\in \Ob E''_{\coh}$, there is an arrow
  $v' : Y'\to X'$ in $E'$, with $Y'$ coherent, such that $u$ factors
  through $f(v') : f(Y') \to f(X')$.

It is clear that condition (iv bis) implies (ii), since for fixed
$X'$ the family of $u : X''\to f(X')$ whose source is coherent is
epimorphic. A fortiori, the same is therefore true of the family of
all $f(v'_{\alpha})$ in the corresponding diagrams $(\ast)$. Thus
$$
{\rm(iv)} \Rightarrow {\rm(i)}{\quoi}.
$$
\end{itemize}

\item Let us spell out condition (ii ter). Let $p$ be a point of
  $E''$. By definition, $\pi(p)$ is the pro-object of $E'$ which
  pro-represents the functor $h : X'\longmapsto f(X')_{p}$, or,
  equivalently, the pro-object defined by the canonical functor
  $E'_{/h} \to E'$, where $E'_{/h}$ denotes the category of pairs
  $(X',u)$ with $X'\in \Ob E'$ and $u\in h(X')$, \ie
  $u\in f(X')_p$. To say that this pro-object is isomorphic to a
  pro-object coming from $\Pro(E'_{\coh})$ means that the subcategory
  of $E'_{/h}$ formed by the $(X',u)$ with
  $X'\in\Ob E'_{\coh}$ is cofinal. In other words, for every object
  $(X',u)$, $u\in f(X')_p$, of $E'_{/h}$, there is an object
  $(X'_0,u_0)$ mapping to it, with $X'_0$ coherent; that is, there is
  a morphism $X'_0\to X'$ ($X'_0$ coherent) such that
  $u\in\text{Im}(f(X'_0)\to f(X'))$. To say that this holds for every
  point $p$ means that, for every object $X'$ of $E'$, the family of
  maps $f(X'_0)\to f(X')$, with $X'_0$ coherent, is surjective on
  every stalk. As we shall see in the appendix (10) that a coherent
  topos $E''$ has enough points, the resulting condition also says
  that, for every $X'\in\Ob E'$, the family
  $f(X'_0)\to f(X')$ is epimorphic. This is precisely condition (ii),
  and hence
$$
{\rm(ii)} \Leftrightarrow \text{(ii ter)}{\quoi}.
$$

\item The\pageoriginale condition that $f$ commute with filtered
  inductive limits (condition (iii)) is equivalent, since the family
  of fiber functors of $E''$ is conservative, to the condition that
  the composite functors $h : X'\longmapsto f(X')_p$ do so. Let us
  prove that this means that the pro-object $\pi(p)$ which
  pro-represents this functor lies in the essential image of
  $\Pro(E'_{PF})$. This will prove the implications
$$
{\rm(iii)} \Leftrightarrow \text{(iii ter)} \Leftarrow \text{(ii ter)}
\Rightarrow \text{(iii bis)}{\quoi}.
$$
We are therefore led to prove the following lemma.

\begin{enonce*}{Lemma 4.12.1}\etiquette[4.12.1]{VI.lem4.12.1}
Let $E'$ be a topos in which the subcategory $E'_{\coh}$ is
generating, and let $X' \in \Ob \Pro(E')$. Then $X'$ belongs to the
essential image of $\Pro(E'_{PF})$ if and only if the functor
$h : E'\to (\Ens)$ which it pro-represents commutes with filtered
inductive limits.
\end{enonce*}

Since a filtered inductive limit of functors commuting with filtered
inductive limits has the same property, sufficiency follows from the
definition of $E'_{PF}$. For necessity, we must prove that if the
functor $h$ commutes with filtered inductive limits, then in the
filtered category $E'_{/h}$ of pairs $(X',u)$ with $u\in h(X')$, the
subcategory formed by the pairs $(X',u)$ with
$X'\in\Ob E'_{PF}$ is cofinal. That is, for every $(X',u)$ in
$E'_{/h}$, there is a morphism $X'_0\to X'$, with
$X'_0\in\Ob E'_{PF}$, such that
$u\in\text{Im}(h(X'_0)\to h(X'))$.\pageoriginale But $X'$ is known to
be a filtered inductive limit of objects $X'_i$ of $E'_{PF}$
(\hyperref[VI.cor1.24.2]{1.24.2} a)), so $h(X')$ is the filtered
inductive limit of the $h(X'_i)$, which proves the assertion.

\item It remains only to prove the implications
$$
\text{(ii bis)} \Rightarrow \text{(ii)\quad and\quad (iii)}
\Leftrightarrow \text{(iii bis)}{\quoi}.
$$
The first implication follows trivially from the fact that every
object $X'$ of $E'$ is a filtered inductive limit of objects $X'_i$
of $E'_{PF}$. For the second, it is enough to note that the family of
functors
$$
Y'' \longmapsto \Hom(X'', Y'') {\quoi},
$$
for $X''\in\Ob {E''}_{\coh}$, is conservative and consists of
functors commuting with filtered inductive limits
(\hyperref[VI.thm1.23]{1.23} (ii)). Consequently, the functor
$f : E'\to E''$ commutes with filtered inductive limits if and only
if the composite functors
$X'\longmapsto \Hom(X'',f(X'))=
\Hom_{\Pro(E')}(g_0(X''),X')$ do so. By
\hyperref[VI.lem4.12.1]{4.12.1}, this is equivalent to the fact that
the $g_0(X'')$ lie in the essential image of $\Pro(E'_{PF})$. \cqfd
\end{enumerate}

\begin{enonce*}[remark]{Remark 4.13}\etiquette[4.13]{VI.rem4.13}
It is useful to state in sheaf-theoretic language the meaning of
condition \hyperref[VI.prop4.10]{4.10} (iv). Consider the commutative
diagram of functors
\[\vcenter{
\xymatrix@C=2.5cm@R=1.5cm{
\Pro(E'_{\coh})^{\circ}  \ar@{^{(}->}[r]^{\alpha}  \ar[d]^{\beta_0}&
\Pro(E')^{\circ}_{\beta} \ar[d]^{\beta}\\
\text{$\Ev{\coh}$} & \text{$\Ev{}$}\ar[l]^{\gamma} ,
}}\leqno{(4.13.1)}
\]\etiquette[4.13.1]{VI.eq4.13.1}
where\pageoriginale the superscript circles denote opposite
categories, and the second horizontal arrow $\gamma$ is restriction
of functors, all the other arrows denoting the familiar fully
faithful functors. Note that the inclusion arrow $\alpha$ \emph{does
not in general commute with projective limits}, even finite ones.
Thus, in general, a sheaf $F$ on a site $C$ (such as $E''$) defines a
presheaf $\alpha\circ F$ on the same site, with values in
$\Pro(E'')^{0}$, which need not be a sheaf. Nevertheless, if a
\emph{presheaf} $F$ on $C$ with values in
$\Pro(E'_{\coh})^{0}$ is such that the corresponding presheaf
$\alpha F$ with values in $\Pro(E')^{0}$ is a sheaf, then $F$ itself
is a sheaf. Indeed, the sheaf condition is a left-exactness property
(II \hyperlink{sga4.exp.II.remark.6.2.item.1}{6.2.1})), and the functors $\beta$, $\beta_0$, and $\gamma$ commute
with small projective limits.
\end{enonce*}

On the other hand, the characterization in
\hyperref[VI.subsubsec4.9.10]{4.9.10} of sheaves on $E''_{\coh}$ with
values in an arbitrary category $P$ shows that, for every functor
$\alpha:P\to Q$ such that $\alpha^{0}:P^0\to Q^0$ commutes with
\emph{finite coproducts} and with \pageoriginale
\emph{passage to quotients by equivalence relations}, the functor
$g\mapsto\alpha\circ g$ from $\Hom({E''}^0_{\coh},P)$ to
$\Hom({E''}^0_{\coh},Q)$ takes sheaves to sheaves. We also know that
in $\Pro(E'_{\coh})$ and in $\Pro(E')$, finite coproducts and
quotients by equivalence relations are representable, and that the
inclusion $E'_{\coh}\to E'$ commutes with them
(\hyperref[VI.cor1.15]{1.15} and
\hyperref[VI.cor1.17.1]{1.17.1}). The same is therefore true for the
categories $\Pro(E'_{\coh})$ and $\Pro(E')$ and for the inclusion
between them (I \hyperref[I.prop8.9.5]{8.9.5} b) and \hyperref[I.prop8.9.7]{8.9.7}), which is precisely $\alpha^0$,
where $\alpha$ is the functor occurring in
\eqhyperref[VI.eq4.13.1]{4.13.1}. It follows that if
\begin{equation*}
\varphi:{E''}^0_{\coh}\to\Pro(E'_{\coh})^0\tag{4.13.2}
\end{equation*}\etiquette[4.13.2]{VI.eq4.13.2}
is a functor, then $\varphi$ \emph{is a sheaf on $E''_{\coh}$ with
values in $\Pro(E'_{\coh})^0$ if and only if $\alpha\circ\varphi$ is
a sheaf on $E''_{\coh}$ with values in $\Pro(E')^0$}. Thus the
category of gluing functors $f:E'\to E''$ satisfying condition
\hyperref[VI.prop4.10]{4.10} (iv) is canonically equivalent to the
category of sheaves $\varphi$ on $E''_{\coh}$ with values in
$\Pro(E'_{\coh})^0$. Taking \hyperref[VI.prop4.10]{4.10} into account,
we conclude that such a sheaf $\varphi$ canonically defines a gluing
functor (up to unique isomorphism), and that the topos $E$ obtained
from the latter is \emph{coherent}. Moreover
(\hyperref[VI.cor4.11]{4.11}), \emph{if $E'$ is perfect}, one thereby
obtains \emph{all} coherent topoi that can be obtained from $E'$ and
$E''$ by gluing.

Note\pageoriginale also that $E'_{\coh}$ is equivalent to a small
category and is stable under finite projective limits. Consequently
(I \hyperref[I.sec8.10.14]{8.10.14}), by the fully faithful functor $\beta_0$ of
\eqhyperref[VI.eq4.13.1]{4.13.1}, $\Pro(E'_{\coh})^0$ is equivalent
to the full subcategory of
$\widehat{E'_{\coh}}=\Hom(E'_{\coh},(\Ens))$ formed by the left exact
functors. Thus the data of a sheaf
\eqhyperref[VI.eq4.13.2]{4.13.2} are equivalent to those of a functor
\begin{equation*}
F:{E''}_{\coh}^0\times E'_{\coh}\to(\Ens)\tag{4.13.3}
\end{equation*}\etiquette[4.13.3]{VI.eq4.13.3}
which is a sheaf in the first argument and left exact in the second;
or, equivalently, to those of a functor
\begin{equation*}
f_0:E'_{\coh}\to E''\tag{4.13.4}
\end{equation*}\etiquette[4.13.4]{VI.eq4.13.4}
which is left exact. If $f$ is the gluing functor satisfying condition
\hyperref[VI.prop4.10]{4.10} (iv) associated with $\varphi$, one
immediately verifies that $f_0$ is canonically isomorphic to the
restriction of $f:E'\to E''$. Thus \emph{under the conditions being
considered, $f_0$ determines $f$} (up to canonical isomorphism). To
summarize the content of \hyperref[VI.cor4.11]{4.11}, we finally state:

\begin{enonce*}{Corollary 4.14}\etiquette[4.14]{VI.cor4.14}
When $E'$ is a perfect topos, the category of gluing functors
$f:E'\to E''$ which give a coherent topos $E$ by gluing is equivalent,
under $f\mapsto f_0=f|E'_{\coh}$, to the category of left exact
functors $f_0:E'_{\coh}\to E''$, or equivalently to the\pageoriginale
category of sheaves $\varphi$ on the site $E''_{\coh}$ (or, what
amounts to the same thing, on the topos $E''$) with values in the
category $\Pro(E'_{\coh})^0$ \emph{(\cf
\hyperref[VI.subsubsec4.9.10]{4.9.10})}.
\end{enonce*}

The first assertion can also be deduced from
\hyperref[VI.prop4.10]{4.10} by observing that
$E'\approxeq\Ind(E'_{\coh})$ (\hyperref[VI.cor1.25]{1.25}). Indeed,
restriction gives an equivalence between the category of functors
$f:E'\to P$ commuting with filtered inductive limits (where $P$ is a
category in which small filtered inductive limits are representable)
and the category of arbitrary functors $f:E'_{\coh}\to P$. It is
immediate that, under these conditions, if finite projective limits
are representable in $P$, then $f$ is left exact if and only if $f_0$
is.

\begin{enonce*}[remark]{Remarks 4.15}\etiquette[4.15]{VI.rem4.15}
\begin{enumerate}
\renewcommand{\theenumi}{\alph{enumi}}
\renewcommand{\labelenumi}{\theenumi)}
\item The proof given for \hyperref[VI.prop4.10]{4.10} shows that
  conditions equivalent to (ii ter) \resp (iii ter) are obtained by
  replacing the category $\Ppoint(E'')$ by a subcategory which
  defines a \emph{conservative} family of fiber functors.

\item Taking for $E''$ the point topos (IV \hyperref[IV.subsec2.2]{2.2}), one sees immediately
  that conditions (ii ter) and (iii ter) are equivalent \emph{only if}
  $E'_{\coh}=E'_{PF}$, \ie if $E'$ is perfect. Thus the condition
  ``$E$ is coherent'' is equivalent to the condition ``$f$ commutes
  with filtered inductive limits'' only if $E'$ is perfect. On the
  other hand, taking more generally for $E''$ a coherent topos of the
  form $\widehat{C}$, one sees that \emph{in this case condition} (i)
  \emph{is equivalent to} (iv) \emph{for every $E'$}. In fact, in the
  general case the editor has been unable to construct an example in
  which $E$ is coherent without condition (iv) being satisfied. This
  is shameful!
\end{enumerate}
\end{enonce*}


% END INCLUDED SGA4_English_current_buildable_source_20260906/source/06/drafts/SGA4_VI_LINES_2937_3239_English.texfrag


% BEGIN INCLUDED SGA4_English_current_buildable_source_20260906/source/06/drafts/SGA4_VI_LINES_3240_3413_English.texfrag
\section{Commutation of the functors $H^i(X,-)$ with filtered
  inductive limits}\etiquette[5]{VI.sec5}\pageoriginale

\begin{enonce*}{Theorem 5.1}\etiquette[5.1]{VI.thm5.1}
Let $E'$ be an algebraic topos (\hyperref[VI.def2.3]{2.3}), let $E$ be
a locally coherent topos (\hyperref[VI.def2.3]{2.3}), and let
$f:E'\to E$ be a coherent morphism of topoi
(\hyperref[VI.def3.1]{3.1}). Then, for every integer $q$, the functors
$R^qf_{\ast}$ \emph{(V \hyperref[V.subsec5.0]{5.0})} commute with filtered inductive limits
of abelian sheaves.
\end{enonce*}

\begin{enonce*}{Corollary 5.2}\etiquette[5.2]{VI.cor5.2}
Let $E'$ be a coherent topos (\hyperref[VI.def2.3]{2.3}). For every
integer $q$, the functor $H^q(E',-)$ commutes with filtered inductive
limits of abelian sheaves.
\end{enonce*}

\begin{enonce*}{Corollary 5.3}\etiquette[5.3]{VI.cor5.3}
Let $E$ be a topos (\resp an algebraic topos)
(\hyperref[VI.def2.3]{2.3}), and let $X$ be an algebraic and coherent
(\resp coherent) object of $E$ (\hyperref[VI.def2.3]{2.3}). For every
integer $q$, the functor $H^q(X,-)$ commutes with filtered inductive
limits of abelian sheaves.
\end{enonce*}

Corollary \hyperref[VI.cor5.2]{5.2} follows from
\hyperref[VI.thm5.1]{5.1} by taking for $E$ the point topos and for
$f:E'\to E$ the unique morphism (IV \hyperref[IV.subsec2.2]{2.2}). Let us show that
\hyperref[VI.cor5.2]{5.2} implies \hyperref[VI.cor5.3]{5.3}. Let
$j_X:E_{/X}\to E$ be the localization morphism. There is a canonical
isomorphism, functorial in the abelian sheaf $F$ (V \hyperref[V.2.2.1]{2.2.1}),
$H^q(X,F)\simeq H^q(E_{/X},j_X^{\ast}F)$. The functor $j_X^{\ast}$
commutes with inductive limits, and the topos $E_{/X}$ is coherent
(\hyperref[VI.subsubsec2.4.6]{2.4.6}), whence
\hyperref[VI.cor5.3]{5.3}. Let us show that
\hyperref[VI.cor5.3]{5.3} implies \hyperref[VI.thm5.1]{5.1}. Let
$E_{\text{cohalg}}$ be the full subcategory of $E$ defined by the
algebraic and coherent objects of $E$. It is a generating category
(\hyperref[VI.subsubsec2.4.5]{2.4.5}). For every abelian sheaf $F$,
$R^qf_{\ast}F$ is the sheaf associated with the presheaf
$X\mapsto H^q(f^{\ast}X,F)$ $(X\in\ob C)$
(\hyperrefext[V][V.prop5.1]{5.1}). Since $f$ is coherent,
$f^{\ast}X$ is coherent for every object $X$ of $C$
(\hyperref[VI.def3.1]{3.1}). By \hyperref[VI.cor5.3]{5.3}, the
presheaf $X\mapsto H^q(f^{\ast}X,F)$ commutes with filtered inductive
limits in the argument $F$, and \hyperref[VI.thm5.1]{5.1} follows,
since the associated-sheaf functor commutes with inductive limits.
(Only the following property of the morphism $f$ is used: there is a
generating family $S$ of $E$ such that, for every $X$ of $S$,
$f^{\ast}(X)$ is algebraic and coherent.) It therefore remains to
prove \hyperref[VI.cor5.2]{5.2}. We already know that the functor
$H^{\circ}(E',-)$ commutes with filtered inductive limits of abelian
sheaves \hyperref[VI.thm1.23]{(1.2.3)}. By a classical result of \cite{sga4.exp.VI.bib.7}, to prove
\hyperref[VI.cor5.2]{5.2} it is enough to show that a filtered
inductive limit of sheaves acyclic for the functor $H^{\circ}(E',-)$
is acyclic for $H^{\circ}(E',-)$. This follows from the next lemma,
applied with $C=E'_{\coh}$.

\begin{enonce*}{Lemma 5.4}\etiquette[5.4]{VI.lem5.4}
Let\pageoriginale $E$ be a locally coherent topos
(\hyperref[VI.def2.3]{2.3}), and let $C$ be a full generating
subcategory of $E$, stable under fiber products, all of whose objects
are quasi-compact (\hyperref[VI.def1.1]{1.1}). A filtered inductive
limit of $C$-acyclic sheaves \emph{(\hyperrefext[V][V.subsec4.2]{4.2})}
is a $C$-acyclic sheaf.
\end{enonce*}

Let $(F_i)_{i\in I}$ be a filtered inductive system of $C$-acyclic
sheaves, and denote its inductive limit by $F$. Every object $Y$ of
$C$ is coherent (\hyperref[VI.prop2.1]{2.1}), and consequently
$F(Y)=\displaystyle\mathop{\varinjlim}_iF_i(Y)$ \hyperref[VI.thm1.23]{(1.2.3)}. To show that
$F$ is $C$-acyclic, it is enough to show that
$\check H^q(Y,F)=0$ for every $q>0$ and every $Y$ of $C$
(\hyperrefext[V][V.prop4.3]{4.3}). Let $X$ be an object of $C$, and
let $\Xx=(X_{\alpha}\to X)_{\alpha\in A}$ be a finite covering
family by objects of $C$. We have
$H^q(\Xx,F)=H^q(C^{\bigdot}(\Xx,F))$
(\hyperrefext[V][V.subsubsec2.4.3]{2.4.3}). By the preceding and the
fact that $C^{\bigdot}(\Xx,F)$ involves only finite products of groups
(\hyperrefext[V][V.subsubsec2.3.3]{2.3.3}), we have
$C^{\bigdot}(\Xx,F)=\varinjlim_iC^{\bigdot}(\Xx,F_i)$. Hence
$$
H^q(\Xx,F)
=H^q\bigl(\varinjlim_iC^{\bigdot}(\Xx,F_i)\bigr)
=\varinjlim_iH^q(C^{\bigdot}(\Xx,F_i))=0
$$
for $q>0$ (\hyperrefext[V][V.prop4.3]{4.3}). Consequently
$\check H^q(X,F)=\varinjlim_{\mathfrak X}H^q(\mathfrak X,F)=0$.

\begin{enonce*}{Corollary 5.5}\etiquette[5.5]{VI.cor5.5}
Let $E$ be a coherent topos (\hyperref[VI.def2.3]{2.3}), and let
${(X_i)_{i\in I}}$ be a decreasing filtered family of coherent
subobjects of the final object. Denote by $\Phi$ the family of closed
subtopoi of $E$ \emph{(IV 9)} contained in the closed complement of
one of the $X_i$. Then $\Phi$ is a family of supports of $E$
\emph{(\hyperrefext[V][V.def6.12]{6.12})}, and for every integer $q$
the functors $H^q_{\Phi}(E,-)$ and $H^q_{\Phi}$
\emph{(\hyperrefext[V][V.subsec6.13]{6.13})} commute with filtered
inductive limits.
\end{enonce*}

Let $Z_i$ be the closed complement of $X_i$. For every abelian sheaf
$F$, there is an exact sequence (\hyperrefext[V][V.prop6.5]{6.5}):
$$
0 \to H^{0}_{Z_i} (E, F) \to H^{0} (E, F) \to H^{0} (X_i, F) \to
H^{1}_{Z_i} (E, F) \to \ldots\quad {\quoi}.
$$
The functors $H^q(E,F)$ and $H^q(X_i,F)$ commute with filtered
inductive limits in $F$ (\hyperref[VI.cor5.2]{5.2},
\hyperref[VI.cor5.3]{5.3}). Consequently $H^q_{Z_i}(E,F)$ does so as
well, whence the analogous property for $H^q_{\Phi}(E,F)$ after
taking the inductive limit over the $Z_i$
(\hyperrefext[V][V.subsec6.13]{6.13}). The assertion concerning the
functors $\mathcal H^q_{\Phi}$ follows by localization and passage to
the associated sheaf.

\begin{enonce*}[remark]{Definition 5.6}\etiquette[5.6]{VI.def5.6}
Let\pageoriginale $E$ be a ringed topos with ring $A$, and let $q$ be
an integer. An $A$-module $G$ is said to be of finite
$q$-presentation if there is a family $X_i$ of objects of $E$ covering
the final object of $E$ such that, for every $i$, there is an exact
sequence of modules on $E_{/X_i}$:
\begin{equation*}
L_q \to L_{q-1} \to \cdots L_1 \to L_0 \to G_{/X_i} \to 0\tag{5.5.1}
\end{equation*}\etiquette[5.5.1]{VI.eq5.5.1}
where, for every $p$, $L_p=A^{n_p}_{/X_i}$, with $n_p$ an integer.
\end{enonce*}

\begin{enonce*}{Proposition 5.7}\etiquette[5.7]{VI.prop5.7}
Let $G$ be a sheaf of finite $q$-presentation. For every $i\leq q-1$,
the functor $\Eext^i_A(G,-)$ commutes with filtered inductive limits
of $A$-modules.
\end{enonce*}

The problem is local on $E$ (\hyperrefext[V][V.prop6.1]{6.1}). Thus,
after localizing to the $E_{/X_i}$, we may suppose that there is an
exact sequence of the form \eqhyperref[VI.eq5.5.1]{5.5.1}. Put
$L^{\bigdot}=L_q\to\cdots\to L_1\to L_0$. We have
$\Eext^i_A(G,F)\simeq\Hh^i(\Hhom_A(L^{\bigdot},F))$, and the functor
$F\mapsto\Hhom_A(L^{\bigdot},F)$ commutes with inductive limits,
which proves the proposition.

\begin{enonce*}{Corollary 5.8}\etiquette[5.8]{VI.cor5.8}
Let $(E,A)$ be a ringed topos, let $G$ be a sheaf of $A$-modules of
finite $q$-presentation, and let $X$ be an algebraic and coherent
object of $E$. The functors $F\mapsto\Ext^i_A(X;G,F)$, for
$\frac{i(i-1)}{2}\leq q-1$, commute with filtered inductive limits
in $F$.
\end{enonce*}

This follows from \hyperref[VI.prop5.7]{5.7} and
\hyperref[VI.cor5.3]{5.3} by the spectral sequence of V \hyperref[V.6.1.3]{6.1.3}.

\setcounter{subsection}{8}
\subsection{}\etiquette[5.9]{VI.subsec5.9}

Let $(E,A)$ be a ringed topos, let $\Phi$ be a family of supports of
$E$ (\hyperrefext[V][V.def6.12]{6.12}), let $G$ be an $A$-module of
finite $r$-presentation (\hyperref[VI.def5.6]{5.6}), and let $X$ be
an algebraic and coherent object of $E$ (\hyperref[VI.def2.3]{2.3}).
Taking the inductive limit over the closed subtopoi in $\Phi$ in the
last spectral sequences of V \hyperref[V.6.9.1]{6.9.1} and V \hyperref[V.6.9.2]{6.9.2}, respectively, gives
two spectral sequences (\hyperrefext[V][V.subsec6.13]{6.13}):
\begin{equation*}
\begin{cases}
'E^{pq}_2 = \varinjlim_{Z \in \Phi} \Ext^p_A (X; G, H^q_Z F)
\Rightarrow \Ext^{p+q}_{A, \Phi} (X; G, F){\quoi},\\
''E^{pq}_2 = \varinjlim_{Z \in \Phi} \Eext^p_A (G, H^q_Z F)
\Rightarrow \Eext^{p+q}_{A, \Phi} (G, F){\quoi}.
\end{cases}\tag{5.9.1}
\end{equation*}\etiquette[5.9.1]{VI.eq5.9.1}

It follows from \hyperref[VI.prop5.7]{5.7} and
\hyperref[VI.cor5.8]{5.8} that there are canonical isomorphisms
\begin{equation*}
\begin{cases}
'E^{pq}_2 \simeq \Ext^p_A (X; G, H^q_\Phi F)&\quav
\frac{p(p-1)}{2} \leq r - 1{\quoi},\\
''E^{pq}_2 \simeq \Eext^p_A (G, H^q_\Phi F) &\quav p \leq r-1{\quoi}.
\end{cases}\tag{5.9.2}
\end{equation*}\etiquette[5.9.2]{VI.eq5.9.2}

\begin{enonce*}{Corollary 5.10}\etiquette[5.10]{VI.cor5.10}
Use\pageoriginale the hypotheses and notation of
\hyperref[VI.cor5.5]{5.5}. Let $A$ be a ring of $E$, and let $G$ be an
$A$-module of finite $q$-presentation. For every integer $i$ such
that $\dfrac{i(i-1)}{2}\leq q-1$, the functors
$F\mapsto\Ext^i_{A,\Phi}(E;G,F)$ and
$F\mapsto\Eext^i_{A,\Phi}(G,F)$
\emph{(\hyperrefext[V][V.subsec6.13]{6.13})} commute with filtered
inductive limits.
\end{enonce*}

This follows from \hyperref[VI.cor5.5]{5.5} and
\hyperref[VI.prop5.7]{5.7} by the first spectral sequences of V \hyperref[V.6.9.1]{6.9.1}
and V \hyperref[V.6.9.2]{6.9.2}, respectively.


% END INCLUDED SGA4_English_current_buildable_source_20260906/source/06/drafts/SGA4_VI_LINES_3240_3413_English.texfrag


% BEGIN INCLUDED SGA4_English_current_buildable_source_20260906/source/06/drafts/SGA4_VI_LINES_3414_3872_English.texfrag
\section{Inductive and projective limits of a fibered category}\etiquette[6]{VI.sec6}

\setcounter{subsection}{-1}
\subsection{}\etiquette[6.0]{VI.subsec6.0}

We assume that the reader is familiar with the theory of fibered
categories \cite{sga4.exp.VI.bib.5}. The essential purpose of this section is to fix
the terminology and notation concerning fibered categories. Unless
otherwise stated, all categories considered in this section belong to
a fixed universe $\Uu$.

Let $\Ff$, $\Gg$, and $\Ee$ be three categories, and let
$\pi:\Ff\to\Ee$ and ${\pi}':\Gg\to\Ee$ be two functors. We denote by
$\Hhom_{\Ee}(\Ff,\Gg)$ the category whose objects are the functors
$u:\Ff\to\Gg$ such that the diagram
\[
\xymatrix@=1.75cm{
\Ff \ar@{-}[rr]^{u} \ar[rd]_{\pi}&             &\Gg \ar[ld]^{\pi'}\\
           & \Ee &
}\]
commutes, and whose morphisms are the $\Ee$-morphisms of functors,
\ie the morphisms of functors carried by $\pi'$ to identity
morphisms.

Let $\xi$ be an object of $\Ee$. The objects $X$ of $\Ff$ such that
$\pi(X)=\xi$ are called \emph{the objects of $\Ff$ over} $\xi$. Let
$f:\xi\to\eta$ be a morphism of $\Ee$. The morphisms $m$ of $\Ff$
such that $\pi(m)=f$ are called the morphisms of $\Ff$
\emph{over}\pageoriginale $f$. Let $f:\xi\to\eta$ be a morphism of
$\Ee$, let $X$ be an object of $\Ff$ over $\xi$, and let $Y$ be an
object of $\Ff$ over $\eta$; we denote by $\Hom_f(X,Y)$ the set of
morphisms from $X$ to $Y$ over $f$.

\begin{enonce*}{Definition 6.1}\etiquette[6.1]{VI.def6.1}
\begin{enumerate}
\renewcommand{\labelenumi}{\theenumi)}
\item A morphism $m:X\to Y$ of $\Ff$ is called \emph{Cartesian} if,
  for every morphism $p:Z\to Y$ over $\pi(m)$, there exists a unique
  morphism $q:Z\to X$ over the identity of $\pi(X)$ such that
  $mq=p$.

\item The category $\Ff$ is called \emph{prefibered} by $\pi$ over
  $\Ee$ if, for every morphism $f:\xi\to\eta$ of $\Ee$, every object
  of $\Ff$ over $\eta$ is the target of a Cartesian morphism over
  $f$.

\item The category $\Ff$ is called \emph{fibered} by $\pi$ over
  $\Ee$ if it is prefibered and the composite of any two composable
  Cartesian morphisms of $\Ff$ is Cartesian.
\end{enumerate}
\end{enonce*}

\setcounter{subsection}{1}
\subsubsection{}\etiquette[6.1.1]{VI.subsubsec6.1.1}

Let $\xi$ be an object of $\Ee$. The \emph{fiber category} of $\Ff$
at $\xi$, denoted by $\Ff_{\xi}$, is the subcategory of $\Ff$ whose
objects are the objects of $\Ff$ over $\xi$ and whose morphisms are
the morphisms of $\Ff$ over $\id_{\xi}$. For two objects $X$ and $Y$
of $\Ff_{\xi}$, we denote by $\Hom_{\xi}(X,Y)$ the set of morphisms
from $X$ to $Y$ in $\Ff_{\xi}$.

\subsubsection{}\etiquette[6.1.2]{VI.subsubsec6.1.2}

Let $f:\xi\to\eta$ be a morphism of $\Ee$. An object $Y$ over
$\eta$ is the target of a Cartesian morphism over $f$ if and only if
the set-valued functor on the fiber $\Ff_{\xi}$
$$
Z \mapsto \Hom_f(Z, Y)\quad Z \in \ob (\Ff_{\xi}){\quoi},
$$
is representable. When the category $\Ff$ is prefibered, the functor
$$
Y \mapsto \Hom_f(., Y)\quad Y \in \ob (\Ff_{\eta}){\quoi},
$$
with values in the category of presheaves of sets on $\Ff_{\xi}$,
in fact takes its values in the category of representable functors on
$\Ff_{\xi}$, and therefore defines, up to unique isomorphism, a
functor $f^{\ast}:\Ff_{\eta}\to\Ff_{\xi}$, called the \emph{inverse
image\pageoriginale functor for} $f$. Suppose that $\Ff$ is
prefibered and choose, for every morphism $f$ of $\Ee$, a base-change
functor $f^{\ast}$. The category $\Ff$ is then fibered over $\Ee$ if
and only if, for every pair $f,g$ of composable morphisms of $\Ee$,
the composite functor $f^{\ast}g^{\ast}$ is an inverse image functor.
Let
\begin{equation*}
C_{f, g} : f^{\ast} g^{\ast} \to (gf)^{\ast}\tag{6.1.2.1}
\end{equation*}\etiquette[6.1.2.1]{VI.eq6.1.2.1}
be the canonical isomorphism. The isomorphisms $C_{f,g}$ satisfy a
cocycle condition arising from the associativity of composition of
morphisms in $\Ee$:
\begin{equation*}
C_{f, hg}(f^{\ast} \circ C_{g, h}) = C_{gh, f}(C_{f, g} \circ h^{\ast})\tag{6.1.2.2}
\end{equation*}\etiquette[6.1.2.2]{VI.eq6.1.2.2}

\subsubsection{}\etiquette[6.1.3]{VI.subsubsec6.1.3}

Conversely, suppose that for every object $\xi$ of $\Ee$ one is given
a category $\Ff_{\xi}$, for every morphism $f:\xi\to\eta$ a functor
$f^{\ast}:\Ff_\eta\to\Ff_{\xi}$, and for every pair $(f,g)$ of
morphisms of $\Ee$ an isomorphism of functors
$C_{f,g}:f^{\ast}g^{\ast}\to(fg)^{\ast}$, such that the $C_{f,g}$
satisfy condition \eqhyperref[VI.eq6.1.2.2]{6.1.2.2}. One can then
construct, in an essentially unique way, a category $\Ff$ fibered over
$\Ee$ whose fibers ``are'' the categories $\Ff_{\xi}$ and whose
base-change functors may be chosen to be ``equal'' to the functors
given in advance \cite{sga4.exp.VI.bib.5}.

\subsubsection{}\etiquette[6.1.4]{VI.subsubsec6.1.4}
Let $\Ff$ and $\Gg$ be two categories over $\Ee$. We denote by
$$
\Hhom_{\cart/\Ee}(\Ff, \Gg)
$$
the full subcategory of $\Hhom_{\Ee}(\Ff,\Gg)$ defined by the
functors that carry Cartesian morphisms of $\Ff$ to Cartesian
morphisms of $\Gg$. The objects of \\
$\Hhom_{\cart/\Ee}(\Ff,\Gg)$ are called \emph{Cartesian functors}.

Let $C$ and $C'$ be two categories, and let $S$ be a set of morphisms
of the category $C$. We denote by $\Hom_{S^{-1}}(C,C')$ the set of
functors from $C$ to $C'$ that carry the morphisms of $S$ to
isomorphisms.

Let $(\Cat)$ denote the category whose objects are the categories
belonging to the universe and whose morphisms are the functors between
these categories (the category $(\Cat)$ does not belong to the
universe).

Let $\Ff$ be a category over $\Ee$, and let $S$ be the set of
Cartesian morphisms of $\Ff$. The category $\Ff$ defines two functors
on $(\Cat)$ with values in the category\pageoriginale of sets
belonging to the universe:
$$
C \mapsto \Hom_{S^{-1}}(\Ff,C)\quad C \in \ob (\Cat)
$$

$C \mapsto \Hom_{\Cart /\Ee}(\Ff,C\times\Ee) =$ \{Cartesian
functors from $\Ff$ to $C\times\Ee$\}

(the category $C\times\Ee$ is regarded as a category over $\Ee$ by
the second projection functor).

\begin{enonce*}{Proposition 6.2}\etiquette[6.2]{VI.prop6.2}
The functors $(\Cat)\to(\Ens):$
\begin{align*}
C &\mapsto \Hom_{S^{-1}} (\Ff, C)\\
C &\mapsto \Hom_{\Cart/\Ee} (\Ff, C \times \Ee)
\end{align*}
are canonically isomorphic. They are representable.
\end{enonce*}

\begin{enonce*}[remark]{Proof}
To prove the first assertion, it is enough to observe that the
Cartesian morphisms of $C\times\Ee$ are the morphisms of the form
$m\times f$, where $m$ is an isomorphism of $C$. To prove the second
assertion, it is enough to prove that the functor
$\Hom_{S^{-1}}(\Ff,.)$ is representable. This follows from \cite{sga4.exp.VI.bib.1}.
We indicate only the idea of the proof. One formally adjoins to the
morphisms of $\Ff$ inverses of the morphisms of $S$. One considers
the free category generated by the objects of $\Ff$, the morphisms of
$\Ff$, and the formal inverses of the morphisms of $S$ (the path
category). One then takes the quotient by the relations arising from
the relations among morphisms of $\Ff$ and relations of the form
$$
s^{-1}s=\id\quav \text{s}s^{-1}=\id\quad s\in S{\quoi}.
$$
The category thus obtained is denoted by $\Ff(S^{-1})$; equipped with
the canonical functor $Q:\Ff\to\Ff(S^{-1})$, it represents the
functor $\Hom_{S^{-1}}(\Ff,.)$.
\end{enonce*}

\begin{enonce*}{Definition 6.3}\etiquette[6.3]{VI.def6.3}
When $\Ff$ is fibered over $\Ee$, the functor
$\Hom_{S^{-1}}(\Ff,.)$ is denoted by
$\varinjlim_{\mathcal{E}^{\circ}}\Ff$ and is called the inductive-limit
functor of $\Ff$ over $\mathcal{E}^{\circ}$. The category representing
it is likewise denoted by
$\varinjlim_{\mathcal{E}^{\circ}}\Ff$ and is called the inductive-limit
category of $\Ff$ over $\Ee^{\circ}$.
\end{enonce*}

\setcounter{subsection}{4}
\setcounter{subsubsection}{-1}
\subsubsection{}\etiquette[6.4.0]{VI.subsubsec6.4.0}

Suppose that $\Ff$ is fibered over $\Ee$. For every morphism $f$ of
$\Ee$, choose a base-change functor over $f$, and let
$$
Q :\Ff \to \varinjlim_{\mathcal{E}^{\circ}} \Ff
$$
be the\pageoriginale canonical functor. Composing the inclusion
functors of the fibers of $\Ff$ into $\Ff$ with $Q$ gives, for every
object $\xi$ of $\Ee$, a functor
$$
u_{\xi} : \Ff_{\xi} \to \varinjlim_{\mathcal{E}} \Ff
$$
and, for every morphism $f:\xi\to\eta$ of $\Ee$, a diagram that
commutes up to canonical isomorphism:
\[
\xymatrix@R=1cm@C=3cm{
\Ff_{\eta} \ar[rd]^{u_{\eta}} \ar[dd]_{f^\ast}&\\
                   &\displaystyle\mathop{\varinjlim}_{\Ee^{0}}\Ff\\
\Ff_{\xi} \ar[ru]_{u_{\xi}} &
}
\]
The category $\varinjlim_{\mathcal{E}^{\circ}}\Ff$ then appears as
the inductive limit \emph{in the sense of pseudofunctors} \cite{sga4.exp.VI.bib.5} of
the pseudofunctor $\Ee^{\circ}\to\Cat$ that associates
$\Ff_{\xi}$ with every $\xi\in\ob$, and the functor
$f^{\ast}:\Ff_{\eta}\to\Ff_{\xi}$ with every
$f:\xi\to\eta$. Note, however, that even when
$\xi\mapsto\Ff_{\xi}$ is an actual functor, the category
$\varinjlim_{\mathcal{E}^{\circ}}\Ff$ is not in general the
inductive limit in the sense of I 2 of the functor
$\xi\mapsto\Ff_{\xi}$ (\cf \hyperref[VI.exe6.8]{6.8}).

\begin{enonce*}{Proposition 6.4}\etiquette[6.4]{VI.prop6.4}
Let $\Ff$ be a category fibered over $\Ee$. Suppose that the category
$\Ee$ has the following properties:
\begin{itemize}
\item[L1)] Every diagram
\[
\xymatrix@C=1.5cm@R=.65cm{
\bigdot \ar[rd]& \\
        & \bigdot\\
\bigdot \ar[ru]&
}
\]
can be embedded in a commutative diagram
\[
\xymatrix@C=1.25cm@R=1cm{
        & \bigdot \ar[rd]&\\
\bigdot \ar[ru]\ar[rd]&         & \bigdot\\
        & \bigdot \ar[ru]&
}
\]

\item[L2)] For\pageoriginale every pair of morphisms
$u,v:\cdot\rightrightarrows\cdot$ such that there exists a morphism
$t$ satisfying $tu=tv$, there exists a morphism $w$ satisfying
$uw=vw$.

(Note that if the category $\Ee^{\circ}$ is pseudofiltered
\emph{(I \hyperref[I.def2.7]{2.7})}, it has properties \emph{L1)} and \emph{L2)}.) The set
$S$ of Cartesian morphisms of $\Ff$ then has the following
properties:

\item[Fr1)] The set $S$ is stable under composition. Isomorphisms
  belong to~$S$.

\item[Fr2)] Every diagram
\[
\xymatrix@C=1.5cm@R=.65cm{
\bigdot \ar[rd]& \\
        & \bigdot\\
\bigdot \ar[ru]^{s}&,
}
\]
where $s$ belongs to $S$, can be completed to a commutative diagram
\[
\xymatrix@C=1.25cm@R=1cm{
        & \bigdot \ar[rd]&\\
\bigdot \ar[ru]^{t}\ar[rd]&         & \bigdot\\
        & \bigdot \ar[ru]^{s}&
}
\]
where $t$ belongs to $S$.

\item[Fr3)] For every pair of morphisms
$u,v:\cdot\rightrightarrows\cdot$ such that there exists a morphism
$s\in S$ satisfying $su=sv$, there exists a morphism $t\in S$ such
that $ut=vt$.
\end{itemize}
\end{enonce*}

\begin{enonce*}[remark]{Proof}
Left to the reader as an exercise.
\end{enonce*}

\begin{enonce*}{Proposition 6.5}\etiquette[6.5]{VI.prop6.5}
Let $\Ff$\pageoriginale be a category, and let $S$ be a set of
morphisms of $\Ff$ having properties \emph{Fr1), Fr2), and Fr3)
(\hyperref[VI.prop6.4]{6.4})}. For every object $X$ of $\Ff$, let
$S(X)$ denote the category of morphisms of $S$ with target $X$. The
category $S(X)$ is cofiltered \emph{(I \hyperref[I.def2.7]{2.7})}. The category
$\Ff(S^{-1})$ (which represents the functor
$\Hom_{S^{-1}}(\Ff,\cdot)$) may then be described as follows:
\begin{enumerate}
\renewcommand{\theenumi}{\alph{enumi}}
\renewcommand{\labelenumi}{\theenumi)}
\item The objects of $\Ff(S^{-1})$ are the objects of $\Ff$.

\item Let $X$ and $Y$ be two objects of $\Ff(S^{-1})$. We have
$$
\Hom_{\Ff(S^{-1})} (X, Y)= \varinjlim_{S(X)}
\Hom_{\Ff} (., Y){\quoi}.
$$

\item Let $f:X\to Y$ and $g:Y\to Z$ be two morphisms of
$\Ff(S^{-1})$. Let
\[
\xymatrix{
\bigdot  \ar[d]_{t} \ar[rrd]^{f'}&\\
X        & &Y & & t\in S
}
\]
(resp.
\[
\xymatrix{
\bigdot \ar[d]_{t'} \ar[rrd]^{g'}&\\
Y                          & & Z & & t'\in S)
}
\]
be a diagram in $\Ff$ whose image is $f$ (\resp $g$). Let
\[
\xymatrix@C=2.5cm@R=1cm{
\bigdot \ar[rd]^{f''} \ar[d]_{t''}& \\
\bigdot \ar[rd]^{f'}        & \bigdot \ar[d]^{t'}\\
                            & Y      & t''\in S
}
\]
be a commutative diagram whose existence is ensured by Fr2). The image
in $\Hom_{\Ff(S^{-1})}(X,Z)$ of the diagram
\[
\xymatrix@C=2cm{
\bigdot \ar[d]_{tt''} \ar[rd]^{g'f''}& \\
X       & Z
}
\]
is the composite morphism $gf$.

Let\pageoriginale
$$
Q : \Ff \to \Ff (S^{-1})
$$
be the canonical functor.

\item The functor $Q^{\ast}:\Ff(S^{-1})^{\widehat{\phantom{,}}}\to
  \Ff^{\widehat{\phantom{,}}}$ ($F\mapsto F\circ Q$ \emph{I \hyperref[I.sec5.0]{5.0}})
  is fully faithful and injective on objects.

\item The functor $Q_{!}:\Ff^{\widehat{\phantom{,}}}\to
  \Ff(S^{-1})^{\widehat{\phantom{,}}}$ (left adjoint to the functor
  $Q^{\ast}$ \emph{(I \hyperref[I.prop5.1]{5.1})}) is left exact. (Consequently, the same is
  true of $Q$ when finite projective limits are representable in
  $\Ff$.)
\end{enumerate}
\end{enonce*}

\begin{enonce*}[remark]{Proof}
For the proof, we refer to [1].
\end{enonce*}

\begin{enonce*}[remark]{Exercise 6.6}\etiquette[6.6]{VI.exe6.6}
Let $\Ff$ be a category, and let $S$ be a set of morphisms of $\Ff$
having properties Fr1), Fr2), and Fr3). For every object $X$ of
$\Ff$, let $S(X)$ denote the category of morphisms with target $X$.
\begin{enumerate}
\renewcommand{\theenumi}{\alph{enumi}}
\renewcommand{\labelenumi}{\theenumi)}
\item The category $S(X)$ is cofiltered (I 2).

\item For every object $X$ of $\Ff$, let $J_S(X)$ denote the set of
  sieves of $X$ that contain a sieve generated by a morphism of $S$.
  The $J_S(X)$ define a topology $T$ on $\Ff$.

\item For the topology $T$, the morphisms of $S$ are bicovering
  (II \hyperref[VI.cor5.2]{5.2}).

\item For every object $X$ of $\Ff$, let $𝒥^2(X)$ denote the category
  of bicovering morphisms with target $X$. The category $S(X)$ is
  cofinal in $𝒥^2(X)$ (I 8).

\item Let $G$ be a presheaf of sets on $\Ff$, and let $\sheaf{a}$ be
  the associated-sheaf functor for the topology $T$. For every object
  $X$ of $\Ff$, there is a canonical isomorphism
$$
\sheaf{a} G(X) \xrightarrow{\sim} \frac{\lim}{S(X)} G( \cdot ){\quoi}.
$$
A presheaf is a sheaf if and only if it carries every morphism of $S$
to an isomorphism.

\item Choose the associated-sheaf functor so that its restriction to
$\Ff$ is injective on objects. Let $\underrightarrow{\Ff}$ denote the
full subcategory of the category of sheaves of sets on $\Ff$ defined
by the sheaves associated with representable presheaves.
\pageoriginale The category $\underrightarrow{\Ff}$ is isomorphic to
the category $\Ff(S^{-1})$ described in \hyperref[VI.prop6.5]{6.5}.
The objects of $\underrightarrow{\Ff}$ form a generating family of
the topos $\Ff^{\sim}$ of sheaves of sets on $\Ff$. The topology on
$\underrightarrow{\Ff}$ induced by the canonical topology of
$\Ff^{\sim}$ is the coarse topology (II \hyperref[II.subsubsec1.1.4]{1.1.4}).

\item Let $Q:\Ff\to\underrightarrow{\Ff}$ be the canonical functor.
The functor $Q$ is a morphism from the site $\underrightarrow{\Ff}$
(with the coarse topology) to the site $\Ff$ (with the topology $T$
of b)) (IV \hyperref[IV.exer5.9]{5.9}), and induces an isomorphism on the corresponding
topoi.

\item Let $u:\Ff\to C$ be a functor. The functor $u$ is continuous if
and only if it carries the morphisms of $S$ to isomorphisms (III \hyperref[III.def1.1]{1.1}).

\item Let $u:\Ff\to C$ be a functor carrying the morphisms of $S$ to
isomorphisms. The functor $u$ factors uniquely as
\[
\xymatrix@=1.5cm{
  \Ff\ar[rr]^{u} \ar[rd]_{Q}&  & C\\
            & \underrightarrow{\Ff}{} \ar[ru]_{v}&
}
\]
(Use III \hyperref[III.rem1.4]{1.4}.)

\item Let $\Pro_S\Ff$ be the full subcategory of $\Pro\Ff$ (I 8)
defined by the pro-objects whose transition morphisms belong to $S$.
Show that the functor $Q$ factors as
\[
\xymatrix@M=.25cm@C=2cm{
\Ff\ar@{^{(}->}[r]^{i}& \Pro_S\Ff\ar[r]^{j}&
\Ff_{S^{-1}}
}
\]
where $j$ is the restriction of $\Pro Q$ to $\Pro_S\Ff$. Show that
the functor $j$ admits a fully faithful left adjoint
$$
A : \Ff_{S^{-1}} \to \Pro_S \Ff{\quoi}.
$$
Show that, for every object $X$ of $\Ff$, $AQ(X)$ is the pro-object
$$
Y \to X \in S(X) \mapsto Y{\quoi}.
$$
\end{enumerate}
\end{enonce*}


% END INCLUDED SGA4_English_current_buildable_source_20260906/source/06/drafts/SGA4_VI_LINES_3414_3872_English.texfrag


% BEGIN INCLUDED SGA4_English_current_buildable_source_20260906/source/06/drafts/SGA4_VI_LINES_3873_4230_English.texfrag
\begin{enonce*}{Proposition 6.7}\etiquette[6.7]{VI.prop6.7}
Let\pageoriginale $\Ee$ be a cofiltered category (I \hyperref[I.sec8.7]{8.7}) and
$\Ff\to\Ee$ a fibered category \emph{(\hyperref[VI.subsubsec6.1.3]{6.1.3})}.
\begin{enumerate}
\renewcommand{\labelenumi}{\theenumi)}
\item Let $\xi$ be an object of $\Ee$, and let $\Ff/\xi$ be the category
  fibered over $\Ee/\xi$ obtained by the base change
  $\Ee/\xi\to\Ee$. The canonical functor
  $\varinjlim_{\Ee/\xi}\Ff/\xi\to\varinjlim_{\Ee}\Ff$ is an equivalence of
  categories.

\item More generally, let $\Ee'\to\Ee$ be a cofinal functor \emph{(I 8)}.
  Let $\Ff'\to\Ee'$ be the fibered category deduced from
  $\Ff\to\Ee$ by the base change $\Ee'\hookrightarrow\Ee$. The canonical
  functor
$$
\varinjlim_{\Ee'} \Ff' \to \varinjlim_{\Ee} \Ff
$$
is an equivalence of categories.
\end{enumerate}
\end{enonce*}

\begin{enonce*}[remark]{Proof}
Exercise.
\end{enonce*}

\begin{enonce*}[remark]{Exercise 6.8}\etiquette[6.8]{VI.exe6.8}
\begin{enumerate}
\renewcommand{\labelenumi}{\theenumi)}
\item Let $g:{\Ee}^{\circ}\to\Cat$ be a functor and $\Gg\to\Ee$ the
  corresponding fibered category \cite{sga4.exp.VI.bib.5}. Show that there exists a
  canonical functor
  $\varinjlim_{\Ee^{\circ}}\Gg\to\varinjlim_{\Ee^{\circ}}g$. Show that
  this functor is not necessarily an equivalence of categories. (One may
  take for $g$ the constant functor whose value is the final object of Cat,
  and for $\Ee$ the category associated with a group.)

\item Suppose that the category $\Ee^{\circ}$ is pseudo-filtered. Show that
  the canonical functor
  $\varinjlim_{\Ee^{\circ}}\Gg\to\varinjlim_{\Ee^{\circ}}g$ is then an
  equivalence of categories.
\end{enumerate}
\end{enonce*}

\begin{enonce*}{Proposition 6.9}\etiquette[6.9]{VI.prop6.9}
Let $\Ff\to\Ee$ be a fibered category (\hyperref[VI.def6.1]{6.1}). The
functor $(\Cat)\to(\Ens)$
$$
C \mapsto \Hom_{\Cart/\Ee} (C \times \Ee, \Ff)
$$
is represented by the category $\Hhom_{\Cart/\Ee}(\Ee,\Ff)$.
\end{enonce*}

Let $C$ be a category, let $F:C\times\Ee\to\Ff$ be a Cartesian
$\Ee$-functor (\hyperref[VI.subsubsec6.1.4]{6.1.4}), and let $X$ be an
object of $C$. The functor $\xi\mapsto F((X,\xi))$ is a Cartesian functor
from $\Ee$ to $\Ff$, denoted $F'(X)$, that depends functorially on $X$.
Hence there is a map $F\mapsto F'$ from
$\Hom_{\Cart/\Ee}(C\Ee,\Ff)$ to
$\Hom(C,\Hhom_{\Cart/\Ee}(\Ee,\Ff))$, functorial in $C$, which is
immediately verified to be a bijection.

\setcounter{subsection}{9}
\subsection{}\etiquette[6.10]{VI.subsec6.10}

The category\pageoriginale $\Hhom_{\Cart/}(\Ee,\Ff)$ is called the category
of \emph{Cartesian sections} of the fibered category $\Ff$. It is also
sometimes called the \emph{projective-limit} category of $\Ff$ over
$\Ee^{\circ}$. It is then denoted
$\varprojlim_{\Ee^{\circ}}\Ff$ \cite{sga4.exp.VI.bib.2}.

\subsection{}\etiquette[6.11]{VI.subsec6.11}

Choose a splitting of the fibered category $\Ff\to\Ee$, \ie for every arrow
$f:\xi\to\eta$ of $\Ee$, choose a base-change functor
$f^{\ast}:\Ff_{\eta}\to\Ff_{\xi}$. For every object $\xi$ of $\Ee$, denote
by
$$
v_{\xi} : \varprojlim_{\Ee^{\circ}} \Ff \to \Ff_{\xi}
$$
the evaluation functor at $\xi$. For every morphism $f:\xi\to\eta$, there
is a diagram of functors commutative up to canonical isomorphism:
\[
\xymatrix@C=2.25cm@R=1.25cm{
\displaystyle\mathop{\varprojlim}_{\Ee^{\circ}} \Ff
\ar@{-}[r]^{v_{\xi}} \ar[rd]^{v_\eta}& \Ff_{\xi} \ar[d]^{f^\ast}\\
                                 & \Ff_\eta.
}
\]
The category $\varprojlim_{\Ee^{\circ}}\Ff$ thus appears as the projective
limit, \emph{in the sense of pseudofunctors}, of
$\xi\mapsto\Ff_{\xi}$. Even when $\xi\mapsto\Ff_{\xi}$ is an actual
functor (\ie even when the chosen splitting is a cleavage \cite{sga4.exp.VI.bib.5}), the
categories $\varprojlim_{\Ee^{\circ}}\Ff$
(\hyperref[VI.subsec6.10]{6.10}) and
$\varprojlim_{\Ee^{\circ}}\Ff_{\xi}$ (I 2) are not equivalent in general.

\section{Fibered topoi and sites}\etiquette[7]{VI.sec7}

\subsection{Fibered topoi}\etiquette[7.1]{VI.subsec7.1}

\begin{enonce*}{Definition 7.1.1}\etiquette[7.1.1]{VI.def7.1.1}
Let $𝒰$ be a universe. A $𝒰$-topos fibered over a category $I$ is a category
fibered over $I$ (\hyperref[VI.def6.1]{6.1})
$$
p : F \to I
$$
whose fibers are $𝒰$-topoi and whose inverse-image
functors\footnote{For the fibered structure
  (\hyperref[VI.subsubsec6.1.2]{6.1.2}).} relative to the morphisms
$f:i\to j$ of $I$ are functors $f^{\ast}:F_j\to F_i$ that are the
inverse-image functors of morphisms of topoi $f\bigdot:F_i\to F_j$
\emph{(IV \hyperref[IV.subsubsec3.1.2]{3.1.2})}.
\end{enonce*}

\setcounter{subsubsection}{1}
\subsubsection{}\etiquette[7.1.2]{VI.subsubsec7.1.2}

Equivalently, a $𝒰$-topos fibered over a category $I$ is a fibered category
$p:F\to I$ whose fibers are $𝒰$-topoi and whose inverse-image
functors\pageoriginale are exact and commute with inductive limits
(IV \hyperref[IV.cor1.6]{1.6}).


\subsubsection{}\etiquette[7.1.3]{VI.subsubsec7.1.3}

For every morphism $f:i\to j$ of $I$, choose an inverse-image functor
$f^{\ast}:F_j\to F_i$. For every pair
$i\xrightarrow{f}j\xrightarrow{g}k$ of composable morphisms of $I$, there is
then a canonical isomorphism
$c_{f,g}:f^{\ast}g^{\ast}\to(gf)^{\ast}$
\eqhyperref[VI.eq6.1.2.1]{6.1.2.1}, and the $c_{f,g}$ have the cocycle
property \eqhyperref[VI.eq6.1.2.2]{6.1.2.2}. In addition, for every
$f:i\to j$, choose a right adjoint $f_{\ast}:F_i\to F_j$ to the functor
$f^{\ast}:F_j\to F_i$. The functor $f_{\ast}:F_i\to F_j$ is defined up to
canonical isomorphism by the property of being right adjoint to $f^{\ast}$.
The general results on adjoint functors then give canonical isomorphisms
\begin{equation*}
{c'}_{g, f} : g_{\ast} f_{\ast} \to (gf)_{\ast}\tag{7.1.3.1}{\quoi}.
\end{equation*}\etiquette[7.1.3.1]{VI.eq7.1.3.1}
These have a cocycle property analogous to property \hyperref[VI.eq6.1.2.2]{6.2.2.2}, which we leave
to the reader to make explicit. Applying the results of
\hyperref[VI.subsubsec6.1.3]{6.1.3}, we then obtain a category fibered over
$I^{\circ}$:
\begin{equation*}
p': F'\to I^{\circ}{\quoi}.\tag{7.1.3.2}
\end{equation*}\etiquette[7.1.3.2]{VI.eq7.1.3.2}
It is readily verified that the fibered category $p':F'\to I^{\circ}$ is
independent, up to unique $I^{\circ}$-isomorphism, of the various choices
used to construct it\footnote{\etiquette[112]{page112}The writer apologizes
  for the nonintrinsic character of this construction.}. The fiber at every
object $i$ of $I^{\circ}$ of the fibered category $p':F'\to I^{\circ}$ is
canonically isomorphic to the $𝒰$-topos $F_i$ and is identified with it. The
inverse-image functors of the fibered category $p':F'\to I^{\circ}$ are
\emph{direct-image} functors of morphisms of topoi.

\subsubsection{}\etiquette[7.1.4]{VI.subsubsec7.1.4}

Let $p:F\to I$ be a fibered $𝒰$-topos. For every $f:i\to j$, choose a
morphism of topoi (IV \hyperref[IV.subsubsec3.1.2]{3.1.2}) $f\bigdot:F_i\to F_j$ such that the
inverse-image functor $f^{\ast}$ (under $f\bigdot$) is the inverse-image
functor for the fibered structure (\hyperref[VI.def6.1]{6.1}). We then
obtain cocycles $c_{f,g}$ relating these morphisms of topoi
\eqhyperref[VI.eq7.1.3.1]{7.1.3.1} and
\eqhyperref[VI.eq7.1.3.2]{7.1.3.2}. The collection of the morphisms
$f.$, $f\in\text{Fl}(I)$, and the cocycles $c_{f,g}$ is called a
\emph{bisplitting} of the fibered topos $F\to I$. Choosing a bisplitting
makes it possible to associate with the fibered topos $F\to I$ a
pseudofunctor $i\mapsto F_i$ from the category $I$ to the 2-category
of\pageoriginale $𝒰$-topoi. Conversely, one can associate with it in an
essentially unique manner (\cf \cite{sga4.exp.VI.bib.5}) a fibered $𝒰$-topos $F\to I$
endowed with a bisplitting such that the corresponding pseudofunctor is
canonically isomorphic to $i\mapsto F_i$. In practice, a fibered $𝒰$-topos
$F\to I$ will most often be denoted by $(F_i)_{i\in I}$, with the implicit
assumption that a bisplitting has been chosen. Note, however, that the
constructions carried out on fibered topoi (such as the associated total
topos (\hyperref[VI.subsec7.4]{7.4}) and the projective limit
(\hyperref[VI.subsec8.1]{8.1})) depend only on the fibered topoi and not on
any bisplittings that may have been chosen.

\begin{enonce*}{Definition 7.1.5}\etiquette[7.1.5]{VI.def7.1.5}
Let $p:F\to I$ and $q:G\to I$ be two $𝒰$-topoi fibered over $I$. A morphism
$m$ of fibered topoi from $(F,p)$ to $(G,q)$ consists of the datum of an
$I$-functor $m^{\ast}:G\to F$ \emph{(\hyperref[VI.subsubsec6.1.4]{6.1.4})}
and, for every object $i$ of $I$, the datum of a right adjoint
$m_{i\ast}:F_i\to G_i$ to the functor $m^{\ast}_i:G_i\to F_i$ obtained by
restricting $m^{\ast}$ to the fibers at $i$. These data must moreover be
subject to the following condition: for every object $i$ of $I$, the pair
of adjoint functors $(m^{\ast}_i,m_{i\ast})$ is a morphism from the topos
$F_i$ to the topos $G_i$.
\end{enonce*}

\setcounter{subsubsection}{5}
\subsubsection{}\etiquette[7.1.6]{VI.subsubsec7.1.6}

Let $m:(F,p)\to(G,q)$ be a morphism of $𝒰$-topoi fibered over $I$, and let
$p':F'\to I^{\circ}$ and $q':G'\to I^{\circ}$ be the fibered categories
associated by the construction of \hyperref[VI.subsubsec7.1.3]{7.1.3} with
the fibered topoi $(F,p)$ and $(G,q)$, respectively. Choose bisplittings of
$F$ and $G$ (\hyperref[VI.subsubsec7.1.4]{7.1.4}), denoted in both cases by
$(f:i\to j)\mapsto f.:(f^{\ast},f_{\ast})$, from $F_i$ to $F_j$ and from
$G_i$ to $G_j$, respectively. Since $f^{\ast}:G\to F$ is an $I$-functor,
for every morphism $f:i\to j$ there is a morphism of functors (transition
morphism)
\begin{equation*}
b_f : f^{\ast}m^{\ast}_j \to m^{\ast}_i f^{\ast}\tag{7.1.6.1}{\quoi},
\end{equation*}\etiquette[7.1.6.1]{VI.eq7.1.6.1}
and for every pair $i\xrightarrow{f}j\xrightarrow{g}k$ of composable
morphisms, we have
\begin{equation*}
m^{\ast}_i (c_{f, g}) b^{\ast}_f f^{\ast} (b_g) = b_{gf} c_{f, g}\tag{7.1.6.2}{\quoi}.
\end{equation*}\etiquette[7.1.6.2]{VI.eq7.1.6.2}
Passing to adjoint functors gives morphisms
\begin{equation*}
b'_f : f_{\ast}m_{i \ast} \to m_{j\ast} f_{\ast}\tag{7.1.6.3}
\end{equation*}\etiquette[7.1.6.3]{VI.eq7.1.6.3}
that satisfy relations analogous to \eqhyperref[VI.eq7.1.6.2]{7.1.6.2}.
Since the functors $f_{\ast}$ are\pageoriginale the inverse-image functors
of the fibered categories $(F',p')$ and $(G',q')$
(\hyperref[VI.subsubsec7.1.3]{7.1.3}), one can construct an
$I^{\circ}$-functor
\begin{equation*}
m_{\ast} : F'\to G'\tag{7.1.6.4}{\quoi},
\end{equation*}\etiquette[7.1.6.4]{VI.eq7.1.6.4}
whose restrictions to the fibers are the functors $m_{i\ast}$ and whose
transition isomorphisms are the $b'_f$. It is verified, with a little
patience, that the functor $m_{\ast}$ so constructed does not depend on the
choice of the morphisms of topoi $f$.\footnote{See the note at the bottom of
  page \pagehyperref[page112]{112}.}

\subsubsection{}\etiquette[7.1.7]{VI.subsubsec7.1.7}

Denote by $\Homtop_I(F,G)$ the set of morphisms of fibered topoi between the
fibered topos $(F,p)$ and the fibered topos $(G,q)$. The set
$\Homtop_I(F,G)$ is the set of objects of a category denoted
$\Hhomtop_I(F,G)$: if $m$ and $n$ are two morphisms of fibered topoi, a
morphism from $m$ to $n$ is an $I$-morphism of functors from $n^{\ast}$ to
$m^{\ast}$. By adjunction, one also obtains an $I^{\circ}$-morphism from the
functor $m_{\ast}$ to the functor $n_{\ast}$
(\hyperref[VI.subsubsec7.1.6]{7.1.6}). Thus, in the end, there are two
functors
\hypertarget{sga4.exp.VI.equation.7.1.7.1}{}\begin{align*}
\Hhomtop_I (F, G) & \to \Hhom_I (G, F)^{\circ}\tag{7.1.7.1}\\
\Hhomtop_I (F, G) &\to \Hhom_{I^{\circ}} (F', G')\tag{7.1.7.2}
\end{align*}\hypertarget{sga4.exp.VI.equation.7.1.7.2}{}
that associate with every morphism
$(m^{\ast},(m_{i\ast})_{i\in\Ob I})$, on the one hand the functor
$m^{\ast}\in\Hom_I(G,F)$ and, on the other hand, the functor
$m_{\ast}\in\Hom_{I^\circ}(F',G')$ obtained by gluing the $m_{i\ast}$ as in
(\hyperref[VI.subsubsec7.1.6]{7.1.6}). The functor $m^{\ast}:G\to F$
\hyperlink{sga4.exp.VI.equation.7.1.7.1}{(7.1.7.1)} is called the \emph{inverse-image functor} under the morphism of
fibered topoi $m:F\to G$; the functor $m_{\ast}:F'\to G'$ \hyperlink{sga4.exp.VI.equation.7.1.7.2}{(7.1.7.2)} is
called the \emph{direct-image functor} under the morphism of fibered topoi
$m:F\to G$. It follows immediately from the definitions that the functors
$m\mapsto m^{\ast}$ and $m\mapsto m_{\ast}$ are \emph{fully faithful}, and
that their essential images are, on the one hand, the set of $I$-functors
from $G$ to $F$ that, fiber by fiber, are exact and commute with inductive
limits and, on the other hand, the set of $I^0$-functors from $F'$ to $G'$
that, fiber by fiber, are exact and commute with inductive limits and, on the
other hand, the set of $I^0$-functors from $F'$ to $G'$ that, fiber by fiber,
possess an exact left adjoint. This justifies the abuses of language in
which a morphism of fibered topoi is defined by its direct image or its
inverse image. In fact, such a morphism of fibered\pageoriginale topoi is
then defined only up to unique isomorphism.

Denote by $\Hhomtop_{\Cart/I}(F,G)$ the full subcategory of
$\Hhomtop_I(F,G)$ generated by the morphisms $m$ such that $m_{\ast}$ (or,
equivalently, $m^{\ast}$) is a Cartesian functor
(\hyperref[VI.subsubsec6.1.4]{6.1.4}). Such morphisms are called
\emph{Cartesian morphisms}. Thus there are fully faithful functors induced
by \hyperlink{sga4.exp.VI.equation.7.1.7.1}{(7.1.7.1)} and \hyperlink{sga4.exp.VI.equation.7.1.7.2}{(7.1.7.2)}:
\hypertarget{sga4.exp.VI.equation.7.1.7.3}{}\begin{align*}
\Hhomtop_{\Cart/I} (F, G) & \to \Hhom_{\cart /I} (G,
  F)^{\circ}{\quoi},\tag{7.1.7.3}\\
\Hhomtop_{\Cart/I} (F, G) &\to \Hhom_{\cart /I^{\circ}} (F',
  G'){\quoi}.\tag{7.1.7.4}
\end{align*}\hypertarget{sga4.exp.VI.equation.7.1.7.4}{}

\subsubsection{}\etiquette[7.1.8]{VI.subsubsec7.1.8}

Let $p:F\to I$, $q:G\to I$, and $r:H\to I$ be three $𝒰$-topoi fibered over
$I$, and let $m:F\to G$ and $n:G\to H$ be two morphisms of fibered
$𝒰$-topoi. As in IV \hyperlink{sga4.exp.IV.subsection.3.3}{3.3}, define the composite of the two morphisms $m$ and
$n$, denoted $nm:F\to H$, and if $𝒱$ is a universe such that $𝒰\in 𝒱$,
define a category ($𝒱$-$𝒰$-$\Top_{/I}$) whose objects are the $𝒰$-topoi
fibered over $I$ with fiber categories that are elements of $𝒱$, and whose
morphisms are the morphisms of $𝒰$-topoi fibered over $I$. Moreover, as in
IV \hyperref[IV.subsubsec3.3.2]{3.3.2}, the composition map
$$
\Homtop_{/I} (F, G) \times \Homtop_{/I} (G, F) \to \Homtop_{/I} (F, H)
$$
is the map induced on objects by a ``functor of composition of morphisms'':
$$
\Hhomtop_{/I} (F, G) \Hhomtop_{/I} (G, H) \to \Hhomtop_{/I} (F, H)
$$
and the considerations of IV \hyperlink{sga4.exp.IV.subsection.3.3}{3.3} and IV \hyperref[IV.subsec3.4]{3.4} extend without change to the
case of fibered $U$-topoi. The reader will moreover have noticed that when
$I$ is a one-point category (one object and one identity morphism), a
$𝒰$-topos fibered over $I$ is ``nothing but'' a $𝒰$-topos, and a morphism
of fibered $𝒰$-topoi is ``nothing but'' a morphism of $𝒰$-topoi; the
preceding constructions then reduce to those of IV 3.

\subsubsection{}\etiquette[7.1.9]{VI.subsubsec7.1.9}

Let $p:F\to I$ be a fibered $𝒰$-topos and $\phi:J\to I$ a functor. The
fibered category $p_J:F_J\to J$ deduced from $(F,p)$ by the base change
$\phi:J\to I$ is a $𝒰$-topos fibered over $J$. Let $p':F'\to I^{\circ}$ and
$(p_J)':(F_J)'\to J^{\circ}$ be the fibered categories deduced respectively
from $(F,p)$ and $(F_J,p_J)$ by the construction of
\hyperref[VI.subsubsec7.1.3]{7.1.3}. It is immediately verified that the
category $((F_J)',(p_J)')$ is canonically isomorphic to the category
$(p')_J:(F')_J\to J^{\circ}$ deduced\pageoriginale from $p':F'\to I^{\circ}$
by the base change $\phi:J^{\circ}\to I^{\circ}$. These two categories
fibered over $J^{\circ}$ will henceforth be identified and denoted
$p'_j:F'_j\to J^{\circ}$. The base-change operation is functorial, and even
2-functorial, in its argument: for every pair $p:F\to I$ and $G\to I$ of
$𝒰$-topoi fibered over $I$, there are two commutative diagrams of categories
and functors:
\[\vcenter{
\xymatrix{
\Hhomtop_I(F,G) \ar[r] \ar[d]& \Hhom_I(G,F)^{\circ} \ar[d]\\
\Hhomtop_J(F_J,G_J) \ar[r]& \Hhom_J(G_J,F_J)^{\circ},
}}\leqno(7.1.9.1)
\]\etiquette[7.1.9.1]{VI.eq7.1.9.1}
\[\vcenter{
\xymatrix{
\Hhomtop_I(F,G) \ar[r] \ar[d]& \Hhom_{I^{\circ}}(F',G') \ar[d]\\
\Hhomtop_J(F_J,G_J) \ar[r]&
\Hhom_{J^{\circ}}(F'_{J^{\circ}},G_{J^{\circ}}),
}}\leqno{(7.1.9.2)}
\]\etiquette[7.1.9.2]{VI.eq7.1.9.2}
where the vertical functors are base-change functors, and the horizontal
functors in \eqhyperref[VI.eq7.1.9.1]{7.1.9.1} are respectively the functors
``morphism $\mapsto$ inverse image'' \hyperlink{sga4.exp.VI.equation.7.1.7.1}{(7.1.7.1)}, while those in
\eqhyperref[VI.eq7.1.9.2]{7.1.9.2} are the functors ``morphism $\to$ direct
image'' \hyperlink{sga4.exp.VI.equation.7.1.7.2}{(7.1.7.2)}.


% END INCLUDED SGA4_English_current_buildable_source_20260906/source/06/drafts/SGA4_VI_LINES_3873_4230_English.texfrag


% BEGIN INCLUDED SGA4_English_current_buildable_source_20260906/source/06/drafts/SGA4_VI_LINES_4231_4582_English.texfrag
\subsection{Fibered sites}\etiquette[7.2]{VI.subsec7.2}

\subsubsection{}\etiquette[7.2.1]{VI.subsubsec7.2.1}

A \emph{fibered $𝒰$-site} $C$ over a category $I$ is a fibered
category $p : C \to I$ whose fibers are endowed with topologies making
them $𝒰$-sites (IV \hyperref[II.def3.0.2]{3.0.2}), such that for every morphism
$f : i \to j$ of $I$, the inverse-image functor
$f^{\ast} : C_j \to C_i$ is a morphism from the site $C_i$ to the site
$C_j$ (IV \hyperref[IV.subsubsec4.9.1]{4.9.1}).

\subsubsection{}\etiquette[7.2.2]{VI.subsubsec7.2.2}

Let $p : C \to I$ and $q : D \to I$ be two fibered $𝒰$-sites over
$I$. A \emph{morphism of fibered $𝒰$-sites from $C$ to} $D$ is an
$I$-functor (\hyperref[VI.subsec6.0]{6.0}) $m : D \to C$ such that,
for every object $i$ of $I$, the functor $m_i : D_i \to C_i$ is a
morphism from the site $C_i$ to the site $D_i$. Denote by
$\Morsite_I(C, D)$ the set of morphisms of fibered sites from $C$ to
$D$. A \emph{continuous functor from the fibered site $D$ to the
fibered site} $C$ is an $I$-functor from $D$ to $C$ that induces
continuous functors on the fibers (III 1). Denote by
$\Cont_{I}(D, C)$ the set of continuous functors from the fibered site
$D$\pageoriginale to the fibered site $C$. Denote by
$\Mmorsite_I (C, D)$ the full subcategory of
$\Hhom_I (D, C)^{\circ}$ defined by the set of objects
$\Morsite_I (C, D) \subset \Hom_I (D, C)$. Likewise, denote by
$\Ccont_I(D, C)$ the full subcategory of $\Hhom_I (D, C)$ defined by
the subset of objects
$\Ccont_I (D, C) \subset \Hom_I (D, C)$. Thus we have two fully
faithful functors, injective on objects,
\begin{equation*}
\Mmorsite_I (C, D) \hookrightarrow \Ccont_I (D, C)^{\circ} \hookrightarrow
\Hhom_I (D, C)^{\circ}{\quoi}.\tag{7.2.2.1}
\end{equation*}\etiquette[7.2.2.1]{VI.eq7.2.2.1}
As in \hyperref[VI.subsubsec7.1.7]{7.1.7}, one defines
\emph{cartesian} morphisms from the fibered site $C$ to the fibered
site $D$. We have a diagram of full subcategories
\[\vcenter{
\xymatrix@M=.2cm{
\Mmorsite_{\Cart/I}(C,D)\ar@{^{(}->}[r] \ar@{_{(}->}[d]&
\Hhom_{\Cart/I}(D,C)^{\circ} \ar@{_{(}->}[d]\\
\Mmorsite_I(C,D)\ar@{^{(}->}[r]& \Hhom_I(D,C)^{\circ}.
}}\leqno{(7.2.2.2)}
\]\etiquette[7.2.2.2]{VI.eq7.2.2.2}

\subsubsection{}\etiquette[7.2.3]{VI.subsubsec7.2.3}

Morphisms of fibered $𝒰$-sites can be composed. This makes it
possible to define the category ($𝒱$-$𝒰$-Site/$I$), whose objects are
the $𝒰$-sites whose fibers belong to a universe $𝒱$ (of which $𝒰$ is
an element). The category $𝒱$-$𝒰$-site is moreover endowed with a
2-category structure; in particular, composition of morphisms of
fibered $𝒰$-sites is functorial with respect to morphisms between
morphisms. For every functor $\phi : J \to I$ and every fibered site
$p : C \to I$, the category $p_J : C_J \to J$ deduced from $(C, p)$
by the base change $\phi : J \to I$ is canonically endowed with a
fibered $𝒰$-site structure over $J$. The base-change operation is
2-functorial.

\subsubsection{}\etiquette[7.2.4]{VI.subsubsec7.2.4}

Let $p : C \to I$ be a fibered category whose fibers are endowed with
topologies making them $𝒰$-sites. Suppose that projective limits are
representable in the fiber categories. For $p : C \to I$ to be a
fibered $𝒰$-site, it is necessary that the inverse-image functors be
left exact and take covering families to covering families; this
condition is also sufficient when the topologies of the fibers are
coarser than the canonical topology (IV \hyperref[IV.subsubsec4.9.2]{4.9.2}). All the fibered sites
used in this Seminar will be of the type described above. Likewise,
let $p : C \to I$ and $q : D \to I$ be two fibered sites such that
finite projective limits are representable in the fiber categories.
\pageoriginale A cartesian functor $m : D \to C$ that induces left
exact functors on the fiber categories and takes covering families of
the fiber categories to covering families is a morphism of fibered
sites from $C$ to $D$; the converse holds if the fiber topologies of
$C$ are coarser than the canonical topology (IV \hyperref[IV.subsubsec4.9.2]{4.9.2}). All the
morphisms of fibered sites used in this Seminar will be of the type
described above.

\subsubsection{}\etiquette[7.2.5]{VI.subsubsec7.2.5}

A fibered $𝒰$-topos $p : F \to I$ is a fibered $𝒰$-site when its
fibers are endowed with the canonical topology, as we shall always do
from now on. Let $m: F \to G$ be a morphism of fibered $𝒰$-topoi. The
inverse-image functor $m^{\ast} : G \to F$
(\hyperref[VI.subsubsec7.1.7]{7.1.7}) is a morphism from the fibered
$𝒰$-site $F$ to the fibered $𝒰$-site $G$. Thus we have a functor
$(𝒱\text{-}𝒰\text{-}\Top_{/I}) \to
(𝒱\text{-}𝒰\text{-}\Site_{/I})$, which is not faithful in general.
But this functor extends naturally to a 2-functor which, for two
fibered $𝒰$-topoi $F$ and $G$ over $I$, induces a functor
\begin{equation*}
\Hhomtop_{I}(F, G) \to \Mmorsite_I (F, G)\tag{7.2.5.1}
\end{equation*}\etiquette[7.2.5.1]{VI.eq7.2.5.1}
that is an \emph{equivalence of categories}, as follows immediately
from the definitions. Moreover, the functor
\eqhyperref[VI.eq7.2.5.1]{7.2.5.1}, composed with the canonical
inclusion $\Mmorsite_I (F, G) \hookrightarrow
\Hhom_{I}(G, F)^{\circ}$, is precisely the functor \hyperlink{sga4.exp.VI.equation.7.1.7.1}{(7.1.7.1)} that
associates its inverse image with a morphism of fibered topoi.

\subsubsection{}\etiquette[7.2.6]{VI.subsubsec7.2.6}

Let $p : C \to I$ be a fibered $𝒰$-site. For every morphism
$f : i \to j$ of $I$, choose a base-change functor
$f^{\ast} : C_j \to C_{i}$. Passing to the categories of
$\sheaf{\Uu}$-sheaves, the functor $f^{\ast}$ extends to a functor
$\Top(f)^{\ast} : C^{\sim}_{j} \to C^{\sim}_{i}$, where
$\Top(f) : C^{\sim}_{i} \to C^{\sim}_{j}$ is a morphism of topoi (IV
\hyperref[IV.eq4.9.1.1]{4.9.1.1}), making the following diagram commutative up to isomorphism:
\[\vcenter{
\xymatrix@C=2.5cm@R=1cm{
C_j \ar[r]^{f^{\ast}} \ar[d]_{\varepsilon_j}& C_i\ar[d]^{\varepsilon_i}\\
C^{\widetilde{\phantom{a}}}_j \ar[r]^{\Top(f)^{\ast}}&
C^{\widetilde{\phantom{a}}}_i
}}\leqno{(7.2.6.1)}
\]\etiquette[7.2.6.1]{VI.eq7.2.6.1}
Let
\begin{equation}
c_{f, g} : f^{\ast} g^{\ast} \to (gf)^{\ast}\tag{7.2.6.2}
\end{equation}\etiquette[7.2.6.2]{VI.eq7.2.6.2}
be the\pageoriginale canonical isomorphisms
\eqhyperref[VI.eq6.1.2.1]{6.1.2.1}, and let
\begin{equation*}
b_f : \Top(f)^{\ast} \epsilon_j \to \epsilon_i f^{\ast}\tag{7.2.6.3}
\end{equation*}\etiquette[7.2.6.3]{VI.eq7.2.6.3}
be the isomorphisms that ``make'' the diagrams
\eqhyperref[VI.eq7.2.6.1]{7.2.6.1} commutative. Using the fact that the
functors $\Top(f)^{\ast}$ commute with inductive limits, one verifies
immediately that, for every pair of composable morphisms
$i \xrightarrow{f} j \xrightarrow{g} k$, there is a unique
isomorphism
\begin{equation*}
\Top(c_{f, g}) : \Top(f)^{\ast} \Top(g)^{\ast} \to \Top(gf)^{\ast}\tag{7.2.6.4}
\end{equation*}\etiquette[7.2.6.4]{VI.eq7.2.6.4}
such that
\begin{equation*}
\epsilon_i (c_{f, g})b_f \Top(f)^{\ast}(b_g) = b_{gf} \Top(c_{f,g})\tag{7.2.6.5}{\quoi}.
\end{equation*}\etiquette[7.2.6.5]{VI.eq7.2.6.5}
Moreover, the isomorphisms \eqhyperref[VI.eq7.2.6.4]{7.2.6.4}
automatically satisfy the cocycle condition of \hyperref[VI.eq6.1.2.2]{6.2.2.2}. We can
therefore construct a fibered category
\begin{equation*}
p^{\sim} : C^{\sim / I} \to I\tag{7.2.6.6}
\end{equation*}\etiquette[7.2.6.6]{VI.eq7.2.6.6}
whose fibers are canonically isomorphic to the $𝒰$-topoi
$C^{\sim}_i$ (\hyperref[VI.subsubsec6.1.3]{6.1.3}), and a cartesian
functor
\begin{equation*}
\epsilon_{C/I} : C \to C^{\sim /I}\tag{7.2.6.7}
\end{equation*}\etiquette[7.2.6.7]{VI.eq7.2.6.7}
that induces on the fibers the functors
$\epsilon_i : C_{i} \to C^{\sim}_i$. It follows immediately from the
definitions that $p^{\sim} : C^{\sim/I} \to I$ is a fibered
$𝒰$-topos over $I$, that $\epsilon_{C/I}$ is a cartesian morphism of
fibered $𝒰$-sites from $C^{\sim/I}$ to $C$, and that the fibered
$𝒰$-topos $C^{\sim /I} \to I$ and the cartesian functor
$\epsilon_{C/I}$ are independent, up to canonical isomorphism, of the
various choices used to construct them (the choices of the base-change
functors $f^{\ast}$ and of the extensions $\Top(f)^{\ast}$). The
fibered $𝒰$-topos $C^{\sim/I}$ is called the \emph{associated
fibered $𝒰$-topos} of the fibered $𝒰$-site.

\setcounter{paragraph}{7}
\paragraph{}\etiquette[7.2.6.8]{VI.par7.2.6.8}

One verifies that formation of the fibered $𝒰$-topos associated with
a fibered site ``commutes'' with base-category change operations.

\begin{enonce*}{Proposition 7.2.7}\etiquette[7.2.7]{VI.prop7.2.7}
Let\pageoriginale $E$ be a fibered $𝒰$-topos over $I$, and let $C$
be a fibered $𝒰$-site over $I$. The functor
$m \mapsto m^{\ast} \epsilon_{C/I}$, which associates with every
morphism of fibered topoi $m : E \to C^{\sim /I}$ the composite with
$\epsilon_{C/I}$ of the associated inverse-image functor
$m^{\ast} : C^{\sim / I} \to E$, induces an equivalence of categories
preserving cartesian objects
$$
\Hhomtop_I (E, C^{\sim / I}) \xrightarrow{\approx} \Mmorsite_I(E,
C){\quoi}.
$$
When finite projective limits are representable in the fiber
categories of $C$, the corresponding fully faithful functor
$$
\Hhomtop_{I} (E, C^{\sim/ I}) \to \Hhom_{I} (C, E)^{\circ}
$$
has as its essential image the set of $I$-functors $g : C \to E$ that
are left exact fiber by fiber and take covering families of the fiber
categories to epimorphic families of the corresponding fiber
categories.
\end{enonce*}

This proposition merely generalizes \hyperref[IV.prop4.9.4]{Proposition IV 4.9.4} to the
fibered context, and the proof given in IV \hyperref[IV.prop4.9.4]{4.9.4} generalizes
immediately.

\subsection{Example}\etiquette[7.3]{VI.subsec7.3}

\begin{enonce*}[remark]{Example 7.3.1}\etiquette[7.3.1]{VI.exem7.3.1}
(\emph{The fibered topos of arrows of a topos})

Let $E$ be a $𝒰$-topos, let $\text{Fl}(E)$ be the category of arrows
of $E$, and let $p : \text{Fl} (E) \to E$ be the functor that assigns
to an arrow its target. The category $(\text{Fl} (E), p)$ is a
fibered $𝒰$-topos over $E$ (\hyperref[VI.def7.1.1]{7.1.1} and IV
\hyperref[VI.thm5.1]{5.1}). The fiber over every object $X$ of $E$ is
the topos $E_{/X}$. The inverse-image functor
$f^{\ast} : E_{/Y} \to E_{/X}$ for a morphism $f : X \to Y$ is the
fiber-product functor. Thus, with every functor $\phi : I \to E$, one
can associate the fibered $𝒰$-topos over $I$
\begin{equation*}
p_I : \text{Fl}(E)_I \to I{\quoi},\tag{7.3.1.1}
\end{equation*}\etiquette[7.3.1.1]{VI.eq7.3.1.1}
obtained by the base change $\phi : I \to E$, whose fiber over every
object $i$ of $I$ is $E_{/\phi (i)}$
(\hyperref[VI.subsubsec7.1.9]{7.1.9}). If $𝒱$ is a universe of which
$E$ is an element, then, with the notation of
\hyperref[VI.subsubsec7.1.8]{7.1.8}, we have a map
\begin{equation*}
\Hom(I, E) \to \ob (𝒱\text{-}𝒰\text{-}\text{\rm top}_{/I}){\quoi},\tag{7.3.1.2}
\end{equation*}\etiquette[7.3.1.2]{VI.eq7.3.1.2}
which associates with every functor $\phi \in \Hom (I, E)$ the
fibered $𝒰$-topos over $I$ deduced from $(\text{Fl}(E), p)$ by the
\pageoriginale base change $\phi : I \to E$. Let $\phi_1$ and
$\phi_2$ be two functors from $I$ to $E$, and let
$m : \phi_1 \to \phi_2$ be a morphism of functors. Denote by
$p_1 : \text{Fl}(E)_1 \to I$ and $p_2 : \text{Fl} (E)_2 \to I$ the
fibered $𝒰$-topoi over $I$ deduced from
$p : \text{Fl}(E) \to E$ by the base changes $\phi_1 : I \to E$ and
$\phi_2 : I \to E$, respectively. For every object $i$ of $I$,
choose a morphism of topoi (IV \hyperref[IV.eq5.5.2]{5.5.2})
\begin{equation*}
\loc(m(i)) : E_{/{\phi}_{1}(i)} \to E_{/{\phi}_{2}(i)}{\quoi}.\tag{7.3.1.3}
\end{equation*}\etiquette[7.3.1.3]{VI.eq7.3.1.3}
One then verifies easily that there is a unique\footnote{Uniqueness
comes from requiring, with the notation of (IV \hyperref[IV.eq5.5.3]{5.5.3}), the formula
$(gf)_! = g_! f_!$.} cartesian morphism of fibered topoi over $I$,
\begin{equation*}
\loc(m) : \text{Fl}(E)_1 \to \text{Fl} (E)_2{\quoi},\tag{7.3.1.4}
\end{equation*}\etiquette[7.3.1.4]{VI.eq7.3.1.4}
that induces the morphisms $\loc(m(i))$ on the fiber topoi. The
morphism $\loc(m) : \text{Fl}(E)_1 \to \text{Fl} (E)_2$ is determined
up to canonical isomorphism by the morphism of functors
$m : \phi_1 \to \phi_2$. If $n : \phi_2 \to \phi_3$ is a morphism of
functors, we have an isomorphism
\begin{equation*}
c(m, n) : \loc (n) \loc(m) \simeq \loc(nm){\quoi},\tag{7.3.1.5}
\end{equation*}\etiquette[7.3.1.5]{VI.eq7.3.1.5}
and the morphisms $c(m, n)$ satisfy a cocycle property analogous to
\eqhyperref[VI.eq6.1.2.2]{6.1.2.2}.
\end{enonce*}

\begin{enonce*}[remark]{Example 7.3.2}\etiquette[7.3.2]{VI.exem7.3.2}
(The fibered site of arrows of a site).

Let $C$ be a $𝒰$-site in which \emph{fiber products are
representable}, let $\text{Fl}(C)$ be the category of arrows of $C$,
and let $p : \text{Fl}(C) \to C$ be the functor that assigns to an
arrow its target. The fibered category $(\text{Fl}(C), p)$ is a
fibered $𝒰$-site over $C$ when the fiber categories are endowed with
the induced topologies (II \hyperref[VI.cor5.2]{5.2}). Let
$\epsilon_C : C \to C^{\sim}$ be the canonical functor from $C$ to
the category of $𝒰$-sheaves on $C$. Denoting by
$\text{Fl}(\epsilon_C)$ the natural extension of the latter functor to
the arrow categories, we have a commutative diagram of categories and
functors:
\[\vcenter{
\xymatrix@C=2cm{
\text{Fl}(C) \ar[r]^{\text{Fl}(\varepsilon_C)}\ar[d]_{p}& \text{Fl}(C^{\widetilde{\phantom{a}}}) \ar[d]^{p}\\
C       \ar[r]^{\varepsilon_C} & C^{\widetilde{\phantom{a}}}
}}\leqno{(7.3.2.1)}
\]\etiquette[7.3.2.1]{VI.eq7.3.2.1}
and hence a\pageoriginale cartesian functor between categories
fibered over $C$:
\[\vcenter{
\xymatrix{
\text{Fl}(C) \ar[rr] \ar[rd]_{p}&     &
\text{Fl}(C^{\widetilde{\phantom{a}}})_C \ar[ld]^{p_{/C}}\\
                   & C     &\\
}}\leqno{(7.3.2.2)}
\]\etiquette[7.3.2.2]{VI.eq7.3.2.2}
where $p_{/C} : \text{Fl}(C^{\sim})_{C} \to C$ is deduced from
$\text{Fl}(C^{\sim}) \to C^{\sim}$ by the base change
$\epsilon_{C} : C \to C^{\sim}$. The functor
$\text{Fl}(C) \to \text{Fl}(C^{\sim})_{C}$ of
\eqhyperref[VI.eq7.3.2.2]{7.3.2.2} is a morphism from the fibered site
$\text{Fl}(C^{\sim})_{C}$ to the fibered site $\text{Fl}(C)$
(\hyperref[VI.subsubsec7.2.2]{7.2.2}); hence, by
\hyperref[VI.prop7.2.7]{7.2.7}, there is a morphism of fibered topoi
over $C$,
\begin{equation*}
\can : \text{Fl}(C^{\sim})_{C} \to \text{Fl} (C)^{\sim/C}{\quoi},\tag{7.3.2.3}
\end{equation*}\etiquette[7.3.2.3]{VI.eq7.3.2.3}
where $\text{Fl}(C)^{\sim/C}$ is the fibered topos over $C$
associated with the fibered site $\text{Fl}(C) \to C$
(\hyperref[VI.subsubsec7.2.6]{7.2.6}). It follows from the definition
of the associated fibered topos
(\hyperref[VI.subsubsec7.2.6]{7.2.6}) and from II
\hyperref[VI.cor5.5]{5.5} that the morphism can of
\eqhyperref[VI.eq7.3.2.3]{7.3.2.3} is a $C$-equivalence of fibered
$𝒰$-topoi over $C$, \ie there is a morphism
$\Theta : \text{Fl}(C)^{\sim/C} \to \text{Fl} (C^{\sim})_C$ of
fibered $𝒰$-topoi over $C$ such that the composites
$\Theta \circ \can$ and $\can \circ \Theta$ are $C$-isomorphic to the
identity.

Let $\phi : I \to C$ be a functor. By base change, we deduce a
fibered site over $I$,
\begin{equation*}
p_I : \text{Fl} (C)_I \to I\tag{7.3.2.4}
\end{equation*}\etiquette[7.3.2.4]{VI.eq7.3.2.4}
and, since formation of the associated fibered topos commutes with
base change (\hyperref[VI.subsubsec7.2.6]{7.2.6}), we have an
$I$-equivalence of fibered $𝒰$-topoi
\begin{equation*}
\can_I : \text{Fl} (C^{\sim})_{I} \to \text{Fl}(C)^{\sim/I}_{I}\tag{7.3.2.5}
\end{equation*}\etiquette[7.3.2.5]{VI.eq7.3.2.5}
where $\text{Fl}(C^{\sim})_{I} \to I$ is the fibered $𝒰$-topos
obtained by the base change
$\epsilon_C \circ \phi : I \to C^{\sim}$.

Finally, let $\phi_1 : I \to C$ and $\phi_2 : I \to C$ be two
functors, and let $m : \phi_1 \to \phi_2$ be a morphism of functors.
Denote by $\text{Fl}(C)_1$ and $\text{Fl}(C)_2$ the fibered sites over
$I$ deduced from $p : \text{Fl}(C) \to C$ by the base changes
$\phi_1$ and $\phi_2$, respectively. As in
\hyperref[VI.exem7.3.1]{7.3.1}, one constructs a morphism of fibered
sites over $I$,
\begin{equation*}
\loc(m) : \text{Fl} (C)_1 \to \text{Fl}(C)_2{\quoi}.\tag{7.3.2.6}
\end{equation*}\etiquette[7.3.2.6]{VI.eq7.3.2.6}
Denote\pageoriginale by $\text{Fl}(C^{\sim})_1 \to I$ and
$\text{Fl}(C^{\sim})_2 \to I$ the fibered $𝒰$-topoi over $I$
obtained by the base changes $\epsilon_C \circ \phi_1$ and
$\epsilon_C \circ \phi_2$, respectively. We have a diagram of
morphisms over $I$ that is commutative up to isomorphism:
\[\vcenter{
\xymatrix@C=2.5cm{
\text{Fl}(C^{\sim})_1 \ar[d]\ar[r]^{\loc(\varepsilon_C\ast m)}&
\text{Fl}(C^{\sim})_2 \ar[d]_{\bigdot}\\
\text{Fl}(C)_1       \ar[r]^{\loc(m)} & \text{Fl}(C)_2
}}\leqno{(7.3.2.7)}
\]\etiquette[7.3.2.7]{VI.eq7.3.2.7}
where the vertical morphisms are deduced by base change from the
morphisms of \eqhyperref[VI.eq7.3.2.2]{7.3.2.2}.
\end{enonce*}

\begin{enonce*}[remark]{Exercise 7.3.3}\etiquette[7.3.3]{VI.exe7.3.3}
(\emph{Small and big fibered sites}).

This exercise continues Exercise IV \hyperref[IV.exer4.10.6]{4.10.6}, whose hypotheses and
notation we use. Denote by $\MM$ the full subcategory of
$\text{Fl}(S)$ defined by the morphisms of $ℳ$, and let
$p : \MM \to S$ be the functor that assigns to a morphism its target.

$6^{\circ}$) Show that $(\MM, p)$ is a fibered site over $S$ and that,
when fiber products are representable in $S$, the inclusion functor
$\MM \to \text{Fl}(S)$ is a cartesian morphism from the fibered site
$\text{Fl} (S) \to S$ to the fibered site $\MM \to S$.

$7^{\circ}$) Let $\MM^{\sim/S} \to S$ be the fibered topos associated
with $\MM \to S$. Show that, with the notation of
\hyperref[VI.exem7.3.2]{7.3.2}, there is a canonical morphism
$g : \text{Fl} (S^{\sim})_{/S} \to \MM^{\sim/S}$ of fibered topoi
over $S$ which induces on each fiber over $X \in \ob (S)$ the
morphism $(\Prol_{X}, \Res_X)$ (IV \hyperref[IV.exer4.10.6]{4.10.6} $4^{\circ})$). Show that,
although for every $X \in \ob (S)$ the morphism
$g_X : S^{\sim}_{/X} \to S(X)^{\sim}$ admits a section (\loccit), the
morphism $g$ does not in general admit a section.
\end{enonce*}

% END INCLUDED SGA4_English_current_buildable_source_20260906/source/06/drafts/SGA4_VI_LINES_4231_4582_English.texfrag


% BEGIN INCLUDED SGA4_English_current_buildable_source_20260906/source/06/drafts/SGA4_VI_LINES_4583_4929_English.texfrag
\subsection{The total topology of a fibered topos or site}\etiquette[7.4]{VI.subsec7.4}

\begin{enonce*}{Definition 7.4.1}\etiquette[7.4.1]{VI.def7.4.1}
Let $p:C\to I$ be a fibered $𝒰$-site and, for every
$i\in\ob(I)$, denote by $\alpha_{i!}:C_i\to C$ the inclusion functor
of the fiber category $C_i$ into $C$. The \emph{total topology} on
$C$ is the coarsest topology on $C$ for which the functors
$\alpha_{i!}$, $i\in\ob(I)$, are continuous \emph{(III \hyperref[III.prop3.6]{3.6})}. The
category $C$ equipped with the total topology is called the
\emph{total site}.
\end{enonce*}

\begin{enonce*}{Proposition 7.4.2}\etiquette[7.4.2]{VI.prop7.4.2}
Let\pageoriginale $p:C\to I$ be a fibered $𝒰$-site, let $T$ be the
total topology on $C$, and, for every $i\in\ob(I)$, let $T_i$ be the
topology of the fiber $C_i$.
\begin{enumerate}
\renewcommand{\labelenumi}{\theenumi)}
\item Let $X$ be an object of $C$ over an object $i$ of $I$. A family
  $(X_\beta\to X)_{\beta\in B}$ is covering for $T$ if and only if it
  is refined by a covering family for $T_i$,
  $(Y_{\gamma}\to X)_{\gamma\in\Gamma}$, in $C_i$.

\item The topology $T$ induces the topology $T_i$ on each category
  $C_i$.

\item For every $i\in\ob(I)$, the functor $\alpha_{i!}:C_i\to C$ is
  cocontinuous.
\end{enumerate}
\end{enonce*}

\setcounter{subsubsection}{2}
\setcounter{paragraph}{0}
\paragraph{}\etiquette[7.4.2.1]{VI.par7.4.2.1}

Property 3) follows from 1) and III \hyperref[III.def2.1]{2.1}. Let us prove 1). For every
object $X$ of $C$, let $J(X)$ be the set of sieves of $C/X$ described
in 1). One verifies that the $J(X)$ have properties $T$ 1), $T$ 2),
and $T$ 3) of II \hyperref[II.def1.1]{1.1}, and hence that they define a topology $T'$ on
$C$. The topology $T'$ is the coarsest topology on $C$ for which the
covering families for the topologies $T_i$ are covering. Thus $T'$ is
coarser than $T$ (III \hyperref[III.prop1.6]{1.6}). To prove that $T'=T$, it is therefore
enough to show that, for every object $i$ of $I$, the functor
$\alpha_{i!}:C_i\to C$ is continuous when $C$ is equipped with the
topology $T'$; equivalently, it is enough to show that the topology
induced by $T'$ on the category $C_i$ is $T_i$, which will prove
property 2) at the same time. Let $T'_i$ be the topology on $C_i$
induced by $T'$, and let $u:R\hookrightarrow X$ be a covering sieve
for $T'_i$. By III \hyperref[III.prop3.2]{3.2}, the morphism
$\alpha_{i!}(u):\alpha_{i!}(R)\to\alpha_{i!}(X)$ is bicovering and in
particular covering. By the definition of $T'$, this means that the
sieve $R\hookrightarrow X$ is covering for $T_i$. Thus $T'_i$ is
coarser than $T_i$. To prove that $T'_i$ is finer than $T_i$, by the
definition of the induced topology (III \hyperref[III.subsec3.1]{3.1}), it is enough to show
that the functor $\alpha_{i!}:C_i\to C$ is continuous when $C_i$ and
$C$ are equipped with $T_i$ and $T'$, respectively. In summary, by
III \hyperref[III.prop1.2]{1.2} and \hyperref[III.prop1.5]{1.5}, to prove properties 1) and 2), it is enough to prove
the following lemma:

\begin{enonce*}{Lemma 7.4.2.2}\etiquette[7.4.2.2]{VI.lem7.4.2.2}
Under the hypotheses and with the notation of
\hyperref[VI.prop7.4.2]{7.4.2} and
\hyperref[VI.par7.4.2.1]{7.4.2.1}, let $X$ be an object of $C$ over
$i$, and let $R\xrightarrow{u}X$ be a morphism of
$C^{\widehat{\phantom{,}}}_i$. The following properties are
equivalent:
\begin{enumerate}
\renewcommand{\theenumi}{\roman{enumi}}
\renewcommand{\labelenumi}{\theenumi)}
\item The morphism\pageoriginale $R\to X$ is bicovering for $T_i$.

\item The morphism
  $\alpha_{i!}(u):\alpha_{i!}(R)\to\alpha_{i!}(X)$ is bicovering for
  $T'$.
\end{enumerate}
\end{enonce*}

i) $\Rightarrow$ ii): It is clear that if $R\to X$ is bicovering for
$T_i$, then the morphism $\alpha_{i!}(u)$ is covering for $T'$. It
remains to show that it is bicovering (II
\hyperref[VI.cor5.2]{5.2}). Let $Y$ be an object of $C$ over an object
$j$ of $I$, let $m,n:Y\rightrightarrows\alpha_{i!}(R)$ be two
morphisms of $C^{\widehat{\phantom{,}}}$ such that
$\alpha_{i!}(u)m=\alpha_{i!}(u)n$, and set
$f=p(\alpha_{i!}(u)m)$. Since the functor $\alpha_{i!}$ commutes with
inductive limits (I \hyperlink{sga4.exp.I.proposition.5.4.item.3}{5.4.3}), we have (I \hyperref[I.prop3.4]{3.4}):
$$
\Hom_{C^{\widehat{\phantom{,}}}} (Y, \alpha_{i!} (R)) \simeq
\varinjlim_{C_{i}/R} \Hom_C (Y, \alpha_{i!} (.)){\quoi}.
$$
For every object $Z$ of $C_i$, with the notation of
\hyperref[VI.subsubsec6.1.1]{6.1.1}, we have
$$
\Hom_{C} (Y, Z) \simeq \coprod_{w \in \Hom(j, i)} \Hom_w (Y,
Z){\quoi},
$$
and therefore
$$
\varinjlim_{C_{i}/R} \Hom (Y, \alpha_{i!} (.)) \simeq
\coprod_{w \in \Hom(j, i)} \varinjlim_{C_{i}/R}
\Hom_{w} (Y, \alpha_{i!} (.)){\quoi}.
$$
For every $w\in\Hom(j,i)$, however, there is an inverse image functor
$w^{\ast}:C_i\to C_j$, and
$\Hom_w(Y,\alpha_{i!}(.))\simeq\Hom_{C_j}(Y,w^{\ast}(.))$. Thus,
still denoting by $w^{\ast}:C^{\widehat{\phantom{,}}}_i\to
C^{\widehat{\phantom{,}}}_j$ the extension of $w^{\ast}$ to
presheaves (which, according to I \hyperref[VI.lem5.4]{5.4}, should be denoted
$(w^{\ast})_{!}$), and observing that the functor
$w^{\ast}:C^{\widehat{\phantom{,}}}_i\to
C^{\widehat{\phantom{,}}}_j$ commutes with inductive limits (I
\hyperlink{sga4.exp.I.proposition.5.4.item.3}{5.4.3}), we obtain an isomorphism
$$
\Hom_{C^{\widehat{\phantom{,}}}} (Y, \alpha_{i!} (R)) =
\coprod_{w \in \Hom(j, i)}
\Hom_{C^{\widehat{\phantom{,}}}_{j}}(Y,w^{\ast}(R)){\quoi}.
$$
Under this isomorphism, the morphisms
$m,n:Y\rightrightarrows\alpha_{i!}(R)$ correspond to two morphisms
$m',n':Y\rightrightarrows f^{\ast}R$ such that
$f^{\ast}(u)m'=f^{\ast}(u)n'$. Since the inverse image functor is a
morphism of topoi, it is in particular continuous, and consequently
$f^{\ast}(u)$ is bicovering (III \hyperref[III.prop3.2]{3.2}). There is therefore a covering
family $(Y_{\delta}\xrightarrow{v_{\delta}}Y)_{\delta}$ in $C_j$
such that $m'v_{\delta}=n'v_{\delta}$ for every $\delta$. It follows
that $mv_{\delta}=nv_{\delta}$ for every $\delta$, and hence that
$\alpha_{i!}(u)$ is bicovering for $T'$.

ii) $\Rightarrow$ i):\pageoriginale Suppose that
$\alpha_{i!}(R)\xrightarrow{\alpha_{i!}(u)}\alpha_{i!}(X)$ is
bicovering. The functor $\alpha_{i!}:C_i\to C$ is cocontinuous (III
\hyperref[III.def2.1]{2.1}). Hence $\alpha^{\ast}_{i}\alpha_{i!}(u)$ is bicovering (III
\hyperref[III.eq2.3.2]{2.3.2} and II \hyperref[VI.cor5.3]{5.3} ii)). The functor
$\alpha_{i!}:C_i\to C$ has property (PPF) of I \hyperref[I.prop5.14]{5.14}. We therefore
have a Cartesian diagram in $C^{\widehat{\phantom{,}}}_i$ (I \hyperref[I.prop5.14]{5.14}
3)):
\[
\xymatrix@C=2.25cm@R=1.25cm{
R  \ar[r] \ar[d]_{u}& \alpha_i^{\ast}\alpha_{i!}(R) \ar[d]^{\alpha_i^{\ast}\alpha_{i!}(u)}\\
X     \ar[r]          & \alpha_i^{\ast}\alpha_{i!}(X).
}
\]
It follows that $u$ is bicovering (II \hyperref[VI.cor5.2]{5.2}).

\begin{enonce*}[remark]{Remark 7.4.3}\etiquette[7.4.3]{VI.rem7.4.3}
\begin{enumerate}
\renewcommand{\labelenumi}{\theenumi)}
\item Proposition \hyperref[VI.prop7.4.2]{7.4.2} can be stated and
  proved in a slightly more general setting than that of fibered
  sites, but apparently without useful application: instead of a
  fibered site $p:C\to I$, one considers a fibered category
  $p:C\to I$ whose fibers are equipped with topologies and whose
  inverse image functors are continuous for these topologies (whereas
  in the case of fibered sites one requires in addition that these
  inverse image functors be morphisms of sites (IV \hyperref[IV.subsec4.9]{4.9})).

\item One can prove that the total topology of a fibered site is the
  finest topology making the functors $\alpha_{i!}:C_i\to C$
  cocontinuous (III 2).

\hypertarget{sga4.exp.VI.remark.7.4.3.item.3}{}\item Let $p:C\to I$ be a fibered $𝒰$-site over a category
  equivalent to a $𝒰$-small category. Then the total site $C$ is a
  $𝒰$-site (II \hyperref[II.def3.0.2]{3.0.2}). Indeed, if, for every object $i$ of $I$,
  $(X_{\beta})_{\beta\in B_i}$ is a small topologically generating
  family of $C_i$, then the family
  $(X_{\beta})_{\beta\in\amalg_iB_i}$ is topologically generating for
  $C$ and is $𝒰$-small. The topos of sheaves on the total site is
  called the \emph{total topos} and is denoted by $\Top(C)$.

\hypertarget{sga4.exp.VI.remark.7.4.3.item.4}{}\item Let $p:C\to I$ be a fibered $𝒰$-site over a small category.
  Since the functor $\alpha_{i!}:C_i\to C$ is cocontinuous, it gives
  rise to a morphism of topoi
  $\alpha_i=(\alpha^{\ast}_i,\alpha^{i\ast})$ from $C^{\sim}_i$ to
  $\Top(C)$ (IV \hyperref[IV.subsec4.7]{4.7}). Since, moreover, the functor
  $\alpha_{i!}:C_i\to C$ is continuous, the functor
  $\alpha^{\ast}_i:\Top(C)\to C^{\sim}_i$ admits a left adjoint,
  again denoted by\pageoriginale
  $\alpha_{i!}:C^{\sim}_i\to\Top(C)$, which extends the functor
  $\alpha_{i!}:C_i\to C$ (III \hyperref[III.prop1.2]{1.2} iv)).
\end{enumerate}
\end{enonce*}

\begin{enonce*}{Proposition 7.4.4}\etiquette[7.4.4]{VI.prop7.4.4}
Let $p:C\to I$ be a fibered $𝒰$-site over a small category $I$.
\begin{enumerate}
\renewcommand{\labelenumi}{\theenumi)}
\item A presheaf $F$ on $C$ is a sheaf for the total topology if and
  only if, for every object $i$ of $I$, the presheaf
  $F\circ\alpha_{i!}=F|C_i$ is a sheaf on $C_i$.

\item A functor $m$ from $C$ to a site $D$ is a continuous functor
  from the total site $C$ to the site $D$ if and only if, for every
  object $i$ of $I$, the functor
  $m|C_i=m\circ\alpha_{i!}:C_i\to D$ is continuous.
\end{enumerate}
\end{enonce*}

If $F$ is a sheaf for the total topology, the presheaf
$F\circ\alpha_{i!}$ is a sheaf on $C_i$ because
$\alpha_{i!}:C_i\to C$ is continuous. Conversely, suppose that
$F\circ\alpha_{i!}$ is a sheaf on $C_i$ for every $i$. Let $X$ be an
object of $C$ over an object $i$ of $I$. For every covering sieve $R$
of $C_i/X$, we have $F(X)\simeq F\circ\alpha_{i!}(R)$ and hence
$F(X)\simeq F(\alpha_{i!}(R))$. Let $R'\hookrightarrow X$ be the
sieve of $C/X$ that is the image of the canonical morphism of
presheaves $\alpha_{i!}(R)\to X$. The morphism
$\alpha_{i!}(R)\to R'$ is an epimorphism of presheaves, and therefore
$F(R')\to F(\alpha_{i!}(R))$ is injective. Since
$F(R')\to F(\alpha_{i!}(R))$ is also surjective, it is bijective, and
hence $F(X)\to F(R')$ is bijective. Since the covering sieves of $X$
are described by \eqhyperref[VI.par7.4.2.1]{7.4.2.1}, the presheaf
$F$ is a sheaf: returning to the construction of the associated
sheaf (II 3), one sees that $F$ is isomorphic to its associated sheaf
(II \hyperref[II.rem3.3]{3.3}). The second assertion follows immediately from the first and
the definition of continuity (III \hyperref[III.def1.1]{1.1}).

\setcounter{subsubsection}{4}
\subsubsection{}\etiquette[7.4.5]{VI.subsubsec7.4.5}

Let $p:C\to I$ be a fibered $𝒰$-site over a category $I$ equivalent
to a small category, and let
$$
f : i \to j
$$
be an arrow of the index category $I$. Again denoting by
$f:C^{\sim}_i\to C^{\sim}_j$ the morphism of topoi determined by the
fibered-site structure, we have (with the notation of \hyperlink{sga4.exp.VI.remark.7.4.3.item.4}{7.4.3.4}) a
diagram of morphisms of $𝒰$-topoi
\[\vcenter{
\xymatrix@C=2.25cm@R=1.25cm{
C^{\sim}_i \ar[d]_{f} \ar[r]^{\alpha_i} & \Top(C)\\
C^{\sim}_j  \ar[ru]_{\alpha_j}&
}}\leqno{(7.4.5.1)}
\]\etiquette[7.4.5.1]{VI.eq7.4.5.1}
where $\Top(C)$\pageoriginale is the total topos \hyperlink{sga4.exp.VI.remark.7.4.3.item.3}{(7.4.3.3)}. The
diagram \eqhyperref[VI.eq7.4.5.1]{7.4.5.1} \emph{is not commutative in
general}. We shall define a canonical morphism of morphisms of topoi:
\begin{equation*}
\rho_f : \alpha_i \to \alpha_j \circ f{\quoi}.\tag{7.4.5.2}
\end{equation*}\etiquette[7.4.5.2]{VI.eq7.4.5.2}
It is enough to define such a morphism at the level of inverse images,
\ie to define a morphism of functors
\begin{equation*}
\beta_f^{\ast} : f^{\ast}\circ \alpha_j^{\ast} \to
\alpha_i^{\ast}{\quoi}, \tag{7.4.5.3}
\end{equation*}\etiquette[7.4.5.3]{VI.eq7.4.5.3}
or, equivalently, using the adjunction between $f^{\ast}$ and
$f_{\ast}$, it is enough to define a morphism of functors
\begin{equation*}
\gamma_f : \alpha^{\ast}_j \to f_{\ast} \alpha^{\ast}_i{\quoi}.\tag{7.4.5.4}
\end{equation*}\etiquette[7.4.5.4]{VI.eq7.4.5.4}
Let $Y$ be an object of $C_j$ and $F$ a sheaf on the total site $C$.
By definition,
$\alpha^{\ast}_jF(Y)=F(\alpha_{j!}(Y))$ and
$f_{\ast}\alpha^{\ast}_iF(Y)=F(\alpha_{i!}f^{\ast}(Y))$, where in the
second member $f^{\ast}:C_j\to C_i$ denotes the inverse image functor
for the fibered structure. By the definition of this inverse image
functor, however, there is a canonical Cartesian morphism
\begin{equation*}
m_Y : \alpha_{i!} f^{\ast} \to \alpha_{j!}{\quoi};\tag{7.4.5.5}
\end{equation*}\etiquette[7.4.5.5]{VI.eq7.4.5.5}
hence, on applying the functor $F$, a map functorial in $Y$:
\begin{equation*}
\gamma_f (F) (Y) = F(m_Y) : \alpha_{j}^{\ast} F(Y) \to f_{\ast}
\alpha^{\ast}_i F(Y){\quoi},\tag{7.4.5.6}
\end{equation*}\etiquette[7.4.5.6]{VI.eq7.4.5.6}
which defines a morphism of sheaves on $C_j$:
\begin{equation*}
\gamma_f (F) : \alpha_{j}^{\ast} F \to f_{\ast} \alpha^{\ast}_i
F{\quoi};\tag{7.4.5.7}
\end{equation*}\etiquette[7.4.5.7]{VI.eq7.4.5.7}
whence the morphism of functors $\gamma_f$ of
\eqhyperref[VI.eq7.4.5.4]{7.4.5.4}.

Now let\pageoriginale $f:i\to j$ and $g:j\to k$ be two composable
morphisms of $I$. There is a canonical isomorphism
$$
c'_{g, f} : g_{\ast} f_{\ast} \to (gf)_{\ast}
$$
and one immediately verifies the compatibility formula
\begin{equation*}
c'_{g, f} \circ g_{\ast} (\gamma_f) \circ \gamma_g =
\gamma_{gf}{\quoi}.\tag{7.4.5.8}
\end{equation*}\etiquette[7.4.5.8]{VI.eq7.4.5.8}

\subsubsection{}\etiquette[7.4.6]{VI.subsubsec7.4.6}

Introduce the fibered $𝒰$-topos
$p^{\sim/I}:C^{\sim/I}\to I$ associated with $p:C\to I$
(\hyperref[VI.subsubsec7.2.6]{7.2.6}) and the corresponding category
fibered over $I^{\circ}$,
$(p^{\sim/I})':(C^{\sim/I})'\to I^{\circ}$
\eqhyperref[VI.eq7.1.3.2]{7.1.3.2}. The preceding considerations
associate with every object $F$ of $\Top(C)$ a family
$F_i=(\alpha^{\ast}_iF)_{i\in\ob(I)}$ of objects of the
$C^{\sim}_i$, together with a family of morphisms
$$
\gamma_f F : F_j \to f_{\ast} F_i,\quad f\in F1(I),\quad f : i \to j{\quoi},
$$
subject to the compatibility conditions of
\eqhyperref[VI.eq7.4.5.8]{7.4.5.8}. Thus with $F$ we have associated
an $I^{\circ}$-functor from the category $I^{\circ}$ to the category
$(C^{\sim/I})'$, \ie an object $\Theta_C(F)$ of the category
$\Hhom_{I^{\circ}}(I^{\circ},(C^{\sim/I})')$. Since this construction
is functorial in $F$, we obtain a functor
\begin{equation*}
\Theta_C : \Top(C) \to \Hhom_{I^{\circ}} (I^{\circ},
(C^{\sim/I})'){\quoi}.\tag{7.4.6.1}
\end{equation*}\etiquette[7.4.6.1]{VI.eq7.4.6.1}

\begin{enonce*}{Proposition 7.4.7}\etiquette[7.4.7]{VI.prop7.4.7}
Let $C\to I$ be a fibered $\mathcal{U}$-site, with $I$ equivalent to
a small category. The functor $\Theta_C$
\eqhyperref[VI.eq7.4.6.1]{7.4.6.1} is an equivalence of categories.
\end{enonce*}

We shall content ourselves with describing a quasi-inverse to
$\Theta_C$. Let \\
$\sigma\in\Hom_{I^{\circ}}(I^{\circ},(C^{\sim/I})')$, \ie an
$I^{\circ}$-functor $i\mapsto\sigma(i)$ from $I^{\circ}$ to
$(C^{\sim/I})'$. For every object $X$ of $C$ over
$p(X)=i\in\ob(I)$, there is a sheaf
$\sigma(p(X))\in\ob C^{\sim}_i$, and hence a set
$\sigma(p(X))(X)$. Let $m:X\to Y$ be a morphism of $C$ over the
morphism $f:i\to j$ of $I$. The morphism $m$ factors uniquely as a
morphism $m':X\to f^{\ast}Y$ of $C_i$ followed by the canonical
Cartesian morphism $f^{\ast}Y\to Y$. Likewise, the morphism
$\sigma(f)$ factors uniquely as a morphism
$\gamma_f(\sigma):\sigma(p(Y))\to f_{\ast}\sigma(p(X))$ followed by
the canonical $I^{\circ}$-Cartesian morphism
$f_{\ast}\sigma(p(X))\to\sigma(p(X))$. We therefore obtain a map
from $\sigma(p(Y))(Y)$ to $\sigma(p(X))(X)$,\pageoriginale by
composing the maps
$$
\sigma(p(Y))(Y) \xrightarrow{\gamma_f(\sigma)(Y)}
f_{\ast}\sigma(p(X))(Y) =
\sigma(p(X))(f^{\ast}Y) \xrightarrow{\sigma(p(X))(m')}
\sigma(p(X))(X){\quoi}.
$$
Thus with every object $X$ of $C$ we have associated a set
$\sigma(p(X))(X)$, and with every morphism $m:X\to Y$ a map
$\sigma(p(Y))(Y)\to\sigma(p(X))(X)$. One verifies that this indeed
defines a presheaf on $C$, and that this presheaf does not depend on
the choice of inverse image functors used in its construction. It
follows from \hyperref[VI.prop7.4.4]{7.4.4} that the presheaf
$X\mapsto\sigma(p(X))(X)$ is a sheaf on the total site $C$, and it is
clear that this sheaf depends functorially on the object $\sigma$ of
$\Hhom_{I^{\circ}}(I^{\circ},(C^{\sim/I})')$. Returning to the
definition of $\Theta_C$, it is also clear that the functor from
$\Hhom_{I^{\circ}}(I^{\circ},(C^{\sim/I})')$ to $\Top(C)$ just
constructed is a quasi-inverse to $\Theta_C$.

\begin{enonce*}[remark]{Remark 7.4.8}\etiquette[7.4.8]{VI.rem7.4.8}
It follows in particular from \hyperref[VI.prop7.4.7]{7.4.7} that the
category $\Hhom_{I^{\circ}}(I^{\circ},(C^{\sim/I})')$ is a
$𝒰$-topos. This latter fact can be seen directly. Using I \hyperref[I.lem9.21.10]{9.21.10},
one sees immediately that conditions a), b), and c) of IV \hyperref[IV.subsubsec1.1.2]{1.1.2} are
satisfied, and the existence of a small generating category follows
from I \hyperref[I.cor9.25]{9.25}.
\end{enonce*}

\setcounter{subsubsection}{8}

% END INCLUDED SGA4_English_current_buildable_source_20260906/source/06/drafts/SGA4_VI_LINES_4583_4929_English.texfrag


% BEGIN INCLUDED SGA4_English_current_buildable_source_20260906/source/06/drafts/SGA4_VI_LINES_4930_5251_English.texfrag
\subsubsection{}\etiquette[7.4.9]{VI.subsubsec7.4.9}

Let $p:C\to I$ and $q:D\to I$ be two $𝒰$-sites fibered over a small
category, and let $m$ be a morphism from the fibered site $C$ to the fibered
site $D$ (\hyperref[VI.subsubsec7.2.2]{7.2.2}). Denote by
$m^{\sim/I}:C^{\sim/I}\to D^{\sim/I}$ the corresponding morphism between
the associated fibered $𝒰$-topoi. The morphism of fibered topoi
$m^{\sim/I}$ is defined up to canonical isomorphism
(\hyperref[VI.prop7.2.7]{7.2.7}). It has a corresponding direct-image
functor $m^{\sim/I}_{\ast}:(C^{\sim/I})'\to(D^{\sim/I})'$
(\hyperref[VI.subsubsec7.1.7]{7.1.7}), which, on passing to sections over
$I^{\circ}$, gives a functor
\begin{equation*}
\Hom_{I^{\circ}} (I^{\circ}, m^{\sim/I}_{\ast}) : \Hhom_{I^{\circ}}
(I^{\circ}, (C^{\sim/I})') \to \Hhom_{I^{\circ}} (I^{\circ},
(D^{\sim/I})'){\quoi}.\tag{7.4.9.1}
\end{equation*}\etiquette[7.4.9.1]{VI.eq7.4.9.1}

Moreover, the functor $m$ is a continuous functor between the total sites
(\hyperref[VI.prop7.4.4]{7.4.4}), and hence, by composition, gives a functor
$m_{\ast}:\Top(C)\to\Top(F)$.

In summary, there is a diagram of functors
\[\vcenter{
\xymatrix@C=2.75cm@R=1.5cm{
\Top(C) \ar[r]^{\theta_C} \ar[d]_{m_{\ast}}& \Hhom_I^{\circ}(I^{\circ},(C^{\sim/I})') \ar[d]_{\Hom_{I^\circ}(I^{\circ},m_{\ast}^{\sim/I})}\\
\Top(D)  \ar[r]^{\theta_D} & \Hhom_I^{\circ}(I^{\circ},(D^{\sim/I})')
}}\leqno{(7.4.9.2)}
\]\etiquette[7.4.9.2]{VI.eq7.4.9.2}

\begin{enonce*}{Proposition 7.4.10}\etiquette[7.4.10]{VI.prop7.4.10}
The diagram\pageoriginale \eqhyperref[VI.eq7.4.9.2]{7.4.9.2} is commutative
up to isomorphism. When $m$ is Cartesian, $m$ is a morphism from the total
site $C$ to the total site $D$.
\end{enonce*}

We leave it to the reader to verify that diagram
\eqhyperref[VI.eq7.4.9.2]{7.4.9.2} is commutative up to isomorphism. To show
that $m$ is a morphism between the total sites, it is enough, in view of
\hyperref[VI.prop7.4.7]{7.4.7}, to show that the functor left adjoint to
$\Hom_{I^{\circ}}(I^{\circ},m^{\sim/I}_{\ast})$ is left exact. We first
describe the inverse-image functor of the morphism $m^{\sim/I}$
(\hyperref[VI.prop7.2.7]{7.2.7}); for every object $i$ of $I$, denote by
$m^{\ast}_i:D^{\sim}_i\to C^{\sim}_i$ the functor induced by
$(m^{\sim/i})^{\ast}$ on the fibers at $i$. For every morphism $f:i\to j$
of $I$, choose morphisms of topoi, denoted in both cases by
$(f^{\ast},f_{\ast}):C^{\sim}_i\to C^{\sim}_j$ and
$(f^{\ast},f_{\ast}):D^{\sim}_i\to D^{\sim}_j$. Since
$(m^{\sim/I})^{\ast}$ is a Cartesian $I^{\circ}$-functor, for every
$f:i\to j$ there is a canonical isomorphism
\begin{equation*}
m^{\ast}_i f_{\ast} \to f^{\ast} m^{\ast}_j{\quoi},\tag{7.4.10.1}
\end{equation*}\etiquette[7.4.10.1]{VI.eq7.4.10.1}
and since the functors $f^{\ast}$ and $f_{\ast}$ are adjoint, there are
canonical morphisms (adjunction morphisms)
\begin{align*}
& \phi : \id \to f_{\ast} f^{\ast}\\
& \psi : f^{\ast} f_{\ast} \to \id{\quoi}. \tag{7.4.10.2}
\end{align*}\etiquette[7.4.10.2]{VI.eq7.4.10.2}
Now consider the following sequence of functorial morphisms:
\begin{equation*}
m^{\ast}_j f^{\ast} \xrightarrow{(1)} f_{\ast} f^{\ast} m^{\ast}_j
f_{\ast} \xrightarrow{(2)} f_{\ast} m^{\ast}_i f^{\ast} f_{\ast}
\xrightarrow{(3)} f_{\ast} m^{\ast}_i{\quoi},\tag{7.4.10.3}
\end{equation*}\etiquette[7.4.10.3]{VI.eq7.4.10.3}
where morphism (1) is deduced from $\phi$, morphism (2) is deduced from the
isomorphism \eqhyperref[VI.eq7.4.10.1]{7.4.10.1}), and morphism (3 is
deduced from $\psi$. By composing the various morphisms\pageoriginale in
the sequence \eqhyperref[VI.eq7.4.10.3]{7.4.10.3}, we obtain a functorial
morphism
\begin{equation*}
\can_f : m^{\ast}_j f_{\ast} \to f_{\ast} m^{\ast}_i{\quoi}.\tag{7.4.10.4}
\end{equation*}\etiquette[7.4.10.4]{VI.eq7.4.10.4}
Now let $\tau$ be an object of
$\Hhom_{I^{\circ}}(I^{\circ}(D^{\sim/I})')$. Thus, for every object $i$ of
$I$, there is an object $\tau(i)$ of $D^{\sim}_i$, and for every morphism
$f:i\to j$, a morphism
\begin{equation*}
\gamma_{f} (\tau) : \tau(j) \to f_{\ast} \tau(i){\quoi},\tag{7.4.10.5}
\end{equation*}\etiquette[7.4.10.5]{VI.eq7.4.10.5}
where the family of the $\gamma_f(\tau)$ has a compatibility property
analogous to \eqhyperref[VI.eq7.4.5.8]{7.4.5.8}). Put
$\sigma(i=m^{\ast}_i(\tau(i))$ and denote by
\begin{equation*}
\gamma_{f} (\sigma) : \sigma(j) \to f_{\ast} \sigma(i)\tag{7.4.10.6}
\end{equation*}\etiquette[7.4.10.6]{VI.eq7.4.10.6}
the composite morphism
$m^{\ast}_j\tau(j)\xrightarrow{m_j^{\ast}\gamma_f(\tau)}
m_j^{\ast}f_{\ast}\tau(i)\xrightarrow{\can_f(\tau(i))}
f_{\ast}m^{\ast}_i\tau(i)$. One verifies that the $\gamma_f(\sigma)$ have
the compatibility property of \eqhyperref[VI.eq7.4.5.8]{7.4.5.8}, and hence
that an object of $\Hhom_{I^{\circ}}(I^{\circ},(C^{\sim/I})')$ has thus been
defined. An ``adjoint functor chasing'' then shows that the functor from
$\Hhom_{I^{\circ}}(I^{\circ},(D^{\sim/I})')$ to
$\Hhom_{I^{\circ}}(I^{\circ},(C^{\sim/I})')$ just constructed is left
adjoint to the functor
$\Hom_{I^{\circ}}(I^{\circ},m^{\sim/I}_{\ast})$. It remains to show that
this left adjoint commutes with finite projective limits, which follows
immediately from the fact that projective limits in the categories of
sections are computed fiber by fiber, and that all the functors used to
construct this left adjoint commute with finite projective limits.

\begin{enonce*}{Corollary 7.4.11}\etiquette[7.4.11]{VI.cor7.4.11}
Let $p:C\to I$ be a $𝒰$-site fibered over a small category, let
$p^{\sim/I}:C^{\sim/I}\to I$ be the associated fibered $𝒰$-topos
(\hyperref[VI.subsubsec7.2.6]{7.2.6}), and let
$\epsilon_I:C\to C^{\sim/I}$ be the canonical morphism from the fibered
$𝒰$-site $C^{\sim/I}$ to the fibered $𝒰$-site $C$. The morphism
$\epsilon_I$ induces an equivalence
$\Top(\epsilon_I):\Top(C)\to\Top(C^{\sim/I})$ between the corresponding
total topoi.
\end{enonce*}

This follows from the commutativity of diagram
\eqhyperref[VI.eq7.4.9.2]{7.4.9.2} and from the fact that the morphism
$\epsilon_I$ gives an equivalence between the associated fibered
$𝒰$-topoi.

\begin{enonce*}[remark]{Remark 7.4.11.1}\etiquette[7.4.11.1]{VI.rem7.4.11.1}
It follows from \hyperref[VI.cor7.4.11]{7.4.11} that in questions concerning
the total topos one may replace fibered sites by the corresponding fibered
topoi.
\end{enonce*}

\begin{enonce*}{Lemma 7.4.12}\etiquette[7.4.12]{VI.lem7.4.12}
Let\pageoriginale $p:C\to I$ be a site fibered over a small category $I$,
and let $i$ be a final object of $I$ \emph{(I 10)}. The functor
$\alpha_{i!}:C^{\sim}_i\to\Top(C)$ \hyperlink{sga4.exp.VI.remark.7.4.3.item.4}{(7.4.3.4)} is exact and fully faithful.
The pairs of functors
$\beta_i=(\alpha_{i!},\alpha^{\ast}_i):\Top(C)\to C^{\sim}_i$ and
$\alpha_i=(\alpha_i^{\ast},\alpha_{i\ast}):C^{\sim}_i\to\Top(C)$ are
morphisms of topoi. The composite morphism
$\beta_i\alpha_i:C^{\sim}_i\to C^{\sim}_i$ is isomorphic to the identity
morphism.
\end{enonce*}

For every object $j$ of $I$, denote by $f_j:j\to i$ the unique morphism from
$j$ to $i$. Using the equivalence
$\Theta_C:\Top(C)\xrightarrow{\sim}\Hhom_{I^{\circ}}(I^{\circ},
(C^{\sim/I})')$ (\hyperref[VI.prop7.4.7]{7.4.7}), one sees that for every
object $X$ of $C_i$, $\alpha_{i!(X)}$ is the section
$j\to f^{\ast}_i(X)$, and that for every section
$j\mapsto\sigma(j)\in\Hom_{I^{\circ}}(I^{\circ},(C^{\sim/I})')$,
$\alpha^{\ast}_i(j\mapsto\sigma(j))$ is the object $\sigma(i)$. Hence
$\alpha_i^{\ast}\alpha_{i!}$ is isomorphic to the identity. Therefore
$\alpha_{i!}$ is fully faithful. Since the functors $f_j^{\ast}$ are left
exact, the functor $\alpha_{i!}$ is left exact. Thus $\beta_i$ is a morphism
of topoi, and since $\alpha_i^{\ast}\alpha_{i!}$ is isomorphic to the
identity, the morphism $\beta_i\alpha_i$ is isomorphic to the identity
morphism.


\begin{enonce*}{Theorem 7.4.13.1}\etiquette[7.4.13.1]{VI.thm7.4.13.1}
Let $p:C\to I$ and $q:D\to I$ be two $𝒰$-sites fibered over a small
category $I$, and let $m:D\to C$ be an $I$-functor. For $m$ to be a morphism
from the fibered site $C$ to the fibered site $D$
\emph{(\hyperref[VI.subsubsec7.2.2]{7.2.2})}, it is sufficient that $m$ be a
morphism from the total site $C$ to the total site $D$.
\end{enonce*}

\setcounter{subsubsection}{13}
\setcounter{paragraph}{1}
\paragraph{}\etiquette[7.4.13.2]{VI.par7.4.13.2}

We must show that if an $I$-functor $m:D\to C$ is a morphism from the total
site $C$ to the total site $D$, then $m$ is a morphism from the fibered site
$C$ to the fibered site $D$, \ie for every object $i$ of $I$, the functor
$m_i:D_i\to C_i$ induced by $m$ on the fibers at $i$ is a morphism from the
site $C_i$ to the site $D_i$. The proof of this last assertion occupies
paragraphs \hyperref[VI.par7.4.13.3]{7.4.13.3} through
\hyperref[VI.par7.4.13.7]{7.4.13.7}.

\paragraph{}\etiquette[7.4.13.3]{VI.par7.4.13.3}
Let us show that, for every object $i$ of $I$, the functor $m_i:D_i\to C_i$
is continuous. Indeed, there is a commutative diagram:\etiquette[133]{page133}
\[
\xymatrix@C=2cm@R=1.5cm{
D_i \ar[r]^{m_i} \ar[d]_{\alpha_{i!}}& C_i \ar[d]^{\alpha_{i!}}\\
D   \ar[r]^{m}& C
}
\]
and it\pageoriginale is enough to show that, for every covering sieve
$R\hookrightarrow X$ of $C_i/X$, the morphism
$m_{i!}(R)\to m_i(X)$ is bicovering (III \hyperref[III.prop1.2]{1.2} and \hyperref[III.prop1.5]{1.5}). But by the
commutativity of the diagram on page \pagehyperref[page133]{133} and the
fact that the functors $\alpha_{i!}$ and $m$ are continuous, the morphism
$\alpha_{i!}m_{i!}(R)\to\alpha_{i!}m_i(X)$ is bicovering. The assertion
therefore follows from \hyperref[VI.lem7.4.2.2]{7.4.2.2}.

\paragraph{}\etiquette[7.4.13.4]{VI.par7.4.13.4}
\emph{Reduction to the case in which $C$ and $D$ are fibered $𝒰$-topoi.}
Since the functors $m_i$ are continuous, they admit natural extensions to
the categories of sheaves, which we denote by
$m_i^{\ast}:D_i^{\sim}\to C_i^{\sim}$ (III \hyperref[III.prop1.2]{1.2} iv)). Let $D^{\sim/I}$ and
$C^{\sim/I}$ be the fibered $𝒰$-topoi associated with the fibered sites $D$
and $C$, respectively, and for every morphism $f$ of $I$, denote by
$f^{\ast}$ the inverse-image functors for the four fibered categories $C$,
$D$, $C^{\sim/I}$, and $D^{\sim/I}$. Since the functor $m$ is an
$I$-functor, for every $f:i\to j$ there is a canonical morphism
$$
f^{\ast} mj \leftarrow m_i f^{\ast}{\quoi}.
$$
Since the functors $f^{\ast}$, $m_i$, and $m_j$ are continuous, by passing
to the categories of sheaves we deduce canonical isomorphisms
$$
f^{\ast} m_j^{\ast} \longleftarrow m_i^{\ast} f^{\ast}{\quoi}.
$$
These isomorphisms have compatibility properties that make it possible to
construct a Cartesian functor
$m^{\ast}:D^{\sim/I}\to C^{\sim/I}$ inducing the functors $m_i^{\ast}$ on
the fiber topoi. Moreover, the diagram
\[
\xymatrix@C=2.5cm@R=1.5cm{
D \ar[r]^{m} \ar[d]_{\varepsilon_I}& C \ar[d]^{\varepsilon_I}\\
D^{\sim/I}   \ar[r]^{m^{\ast}}& C^{\sim/I}
}
\]
is commutative up to isomorphism. Since the functors $\epsilon_I$ induce
equivalences on the total topoi (\hyperref[VI.cor7.4.11]{7.4.11}), and since
$m$ is a morphism between the total sites, the functor
$m^{\ast}:D^{\sim/I}\to C^{\sim/I}$ is a morphism between the total sites.
We are therefore reduced to proving the assertion of
\hyperref[VI.par7.4.13.2]{7.4.13.2} when $D$ and $C$ are fibered $𝒰$-topoi,
which we assume henceforth.

\paragraph{}\etiquette[7.4.13.5]{VI.par7.4.13.5}
\emph{For every\pageoriginale object $i$ of $I$, the functor $m_i:D_i\to C_i$
takes the final object of $D_i$ to the final object of $C_i$}. Denote by
$m:D^{\sim}\to C^{\sim}$ the natural extension of $m$ to the total topoi
(III \hyperref[III.prop1.2]{1.2}). Using the equivalences
$\Theta_D:D^{\sim}\to\Hhom_{I^{\circ}}(I^{\circ},D')$
(\hyperref[VI.prop7.4.7]{7.4.7}), one immediately verifies that $m$
associates with every section
$\sigma\in\Hom_{I^{\circ}}({I^{\circ}},D')$ the section
$(i\mapsto m_i\sigma(i))\in\Hom_{I^{\circ}}({I^{\circ}},C')$. Since $m$ is
a morphism of sites, the functor $m$ is left exact and in particular takes
the final object of $\Hhom_{I^{\circ}}({I^{\circ}},D)$ to the final object
of $\Hhom_{I^{\circ}}({I^{\circ}},C')$. For every object $i$ of $I$, denote
by $e_i$ a final object of $D_i$. A final object of
$\Hhom_{I^{\circ}}({I^{\circ}},D')$ is the section $i\mapsto e_i$. Thus the
section $i\mapsto m_i(e_i)$ is a final object of
$\Hhom_{I^{\circ}}({I^{\circ}},C')$, and consequently, for every object
$i$ of $I$, $m_i(e_i)$ is a final object of $C_i$.

\paragraph{}\etiquette[7.4.13.6]{VI.par7.4.13.6}
\emph{Reduction to the case in which $i$ is a final object of} $I$. Let $i$
be an object of $I$. The functor $\alpha_{i!}:D_i\to D$ factors as
$$
D_i \xrightarrow{\bar{\alpha}_{i!}} D_{/e_i} \xrightarrow{\loc(e_i)} D{\quoi},
$$
where $\bar{\alpha}_i$ is the evident functor and $\loc(e_i)$ the forgetful
functor. Denote by $m/i:D_{/e_i}\to D_{m(e_i)}$ the functor that associates
with every object $u:X\to e_i$ of $D_{/e_i}$ the object
$m(u):m(X)\to m(e_i)$ of $C_{/m(e_i)}$. There is a commutative diagram
\[
\xymatrix@C=2.25cm@R=1.5cm{
D_i \ar[r]^{m_i} \ar[d]_{\bar{\alpha}_{i!}}& C_i\ar[d]^{\bar{\alpha}_{i!}}\\
D_{/e_i}  \ar[r]^{m/i} & C_{/m(e_i)}
}
\]
Since $e_i$ and $m(e_i)$ are final objects of $D_i$ and $C_i$, respectively
(\hyperref[VI.par7.4.13.5]{7.4.13.5}), the categories $D_{/e_i}$ and
$C_{/m(e_i)}$ are $𝒰$-topoi fibered over $I_{/i}$, and the functor $m_{/i}$
is an $I_{/i}$-functor. Since the property of a functor being a morphism of
sites is local (IV \hyperref[IV.subsec4.9]{4.9}), the functor $m_{/i}$ is a morphism from the total
site $C_{/m(e_i)}$ to the total site $D_{/(e_i)}$. We are therefore reduced
to proving that $m_i$ is a morphism of topoi when $i$ is a final object of
$I$, which we assume henceforth.\pageoriginale

\paragraph{}\etiquette[7.4.13.7]{VI.par7.4.13.7}
\emph{End of the proof}. Denote by $\alpha_{i!}:D_i\to D^{\sim}$,
$m^{\sim}:D^{\sim}\to C^{\sim}$, and $\alpha_{i!}:C_i\to C^{\sim}$ the
natural extensions to the categories of sheaves of the functors
$\alpha_{i!}$ and $m$. There is a diagram commutative up to isomorphism:
\[
\xymatrix@C=2.5cm@R=1.5cm{
D_i \ar[r]^{m_i} \ar[d]_{\alpha_{i!}}& C_i\ar[d]^{\alpha_{i!}}\\
D^{\sim}  \ar[r]^{m^{\sim}} & C .
}
\]
Since $i$ is a final object, the functor
$\alpha_i^{\ast}\alpha_{i!}:C_i\to C_i$ is isomorphic to the identity
(\hyperref[VI.lem7.4.12]{7.4.12}), and consequently $m_i$ is isomorphic to
$\alpha_i^{\ast}m^{\sim}\alpha_{i!}$. Moreover, the functors
$\alpha_i^{\ast}$, $m^{\sim}$, and $\alpha_i^{\ast}$ are left exact
(\hyperref[VI.lem7.4.12]{7.4.12}). Consequently, $m_i$ is left exact and is
therefore a morphism from the site $C_i$ to the site $D_i$.

\begin{enonce*}[remark]{Exercise 7.4.14}\etiquette[7.4.14]{VI.exe7.4.14}
Let $X=(X_i)_{i\in I}$ be a topos fibered over $I$ and let $T$ be a topos.
Show that the categories $\Hhomtop(X_i,T)$ are the fibers of a fibered
category denoted $\Hhomtop_I(X,T\times I)$. Show that the category of
sections over $I$ of $\Hhomtop_I(X,T\times I)$ is canonically equivalent to
$\Hhomtop_I(\Top(X),T)$, where $\Top(X)$ is the total topos of $X$.
\end{enonce*}

\begin{enonce*}[remark]{Exercise 7.4.15}\etiquette[7.4.15]{VI.exe7.4.15}
Let $X=(X_i)_{i\in I}$ be a topos fibered over $I$. Define a Cartesian
morphism $m$ from $X$ to the constant fibered topos with fiber $(\Ens)$.
Show that the total topos of the constant fibered topos with fiber $(\Ens)$
is canonically equivalent to $I^{\widehat{\phantom{,}}}$. Deduce a morphism
of topoi
$$
\Top(m) :\Top(X) \to I^{\widehat{\phantom{,}}}{\quoi}.
$$
Let $F=(F_i)_{i\in I}$ (where $F_i=F|X_i$) be an object of $\Top(X)$
(\hyperref[VI.prop7.4.7]{7.4.7}). Show that
$$
\Top(m)(F) = (i \mapsto \Gamma (X_i, F_i)){\quoi}.
$$
Deduce from \hyperref[VI.lem7.4.12]{7.4.12} that if $F$ is injective
abelian, then $F_i$ is injective abelian for every $i\in\ob I$. Deduce that
$$
R^q \Top(m)(F) = (i \mapsto H^{q} (X_i, F_i)){\quoi}.
$$
Show\pageoriginale that the Cartan--Leray spectral sequence of the morphism
$m$ (\hyperrefext[V][V.sec5]{5}) is
$$
H^{p+q} (\Top(X), F) \Leftarrow {\varprojlim_{I}}^{(p)} H^{q}
(X_i, F|X_i){\quoi}.
$$
\end{enonce*}


% END INCLUDED SGA4_English_current_buildable_source_20260906/source/06/drafts/SGA4_VI_LINES_4930_5251_English.texfrag


% BEGIN INCLUDED SGA4_English_current_buildable_source_20260906/source/06/drafts/SGA4_VI_LINES_5252_5495_English.texfrag
\section{Projective limits of fibered topoi}\etiquette[8]{VI.sec8}

\subsection{Generalities}\etiquette[8.1]{VI.subsec8.1}

\begin{enonce*}{Definition 8.1.1}\etiquette[8.1.1]{VI.def8.1.1}
Let $𝒰$ be a universe and let $p:F \to I$ be a fibered $𝒰$-topos. A
pair $(C, m)$ consisting of a $𝒰$-topos $C$ and a cartesian morphism
$m : C \times I \to F$ of fibered $𝒰$-topoi over $I$ is called a
projective limit of the fibered topos $F$ if, for every $𝒰$-topos
$D$, the functor
\begin{equation*}
\Hhomtop(D, C) \to \Hhomtop_{\cart/I}(D\times I, F)\tag{8.1.1.1}
\end{equation*}\etiquette[8.1.1.1]{VI.eq8.1.1.1}
obtained by composing the base-change functor
$$
\Hhomtop(D, C) \to \Hhomtop_{I}(D\times I, C \times I)
$$
with the functor of composition with $m$
(\hyperref[VI.subsubsec7.1.8]{7.1.8}) is an equivalence of categories.
\emph{(\cf \hyperref[VI.par8.1.3.3]{8.1.3.3} for a more
``geometric'' formulation)}.
\end{enonce*}

\setcounter{subsubsection}{1}
\subsubsection{}\etiquette[8.1.2]{VI.subsubsec8.1.2}

We use the notation of \hyperref[VI.def8.1.1]{8.1.1}. Let $D$ be a
$𝒰$-topos and let $m_1 : D \times I \to F$ be a morphism of fibered
$𝒰$-topoi. Translating Definition \hyperref[VI.def8.1.1]{8.1.1}, one
immediately sees that there is a morphism of topoi $n : D \to C$ and
an isomorphism of morphisms of fibered topoi over $I$,
$\gamma : m_{1} \xrightarrow{\sim} m \circ (n \times \id_I)$, making
the diagram commutative:
\[\vcenter{
\xymatrix@C=.75cm@R=.5cm{
D \times I  \ar[rrrr]^{m_1} \ar[dddd]_{n\times \id_I} &&{\phantom{.}}\ar@/_.5pc/@{-->}[dr]_{\gamma}  &&F\\
                         & &{\phantom{a}}& &\\
                         & && & \\
                         & && & \\
C \times I \ar[rrrruuuu]_{m}&& &&{\phantom{a}}\\
}}\leqno{(8.1.2.1)}
\]\etiquette[8.1.2.1]{VI.eq8.1.2.1}
Moreover, if $(n', \gamma')$ is another pair with the property above,
there is a \emph{unique} isomorphism
$\eta : n \xrightarrow{\sim} n'$ such that
\begin{equation*}
\gamma \circ m (\eta \times \id) = \gamma'{\quoi}.\tag{8.1.2.2}
\end{equation*}\etiquette[8.1.2.2]{VI.eq8.1.2.2}
It\pageoriginale follows that if $(D, m_1)$ is a projective limit of
$F$, then $n$ is an equivalence of $𝒰$-topoi (IV \hyperref[IV.subsec3.4]{3.4}), and the
additional datum of the isomorphism $\gamma$ making diagram
\eqhyperref[VI.eq8.1.2.1]{8.1.2.1} commutative determines the pair
$(n, \gamma)$ up to canonical isomorphism. Existence of the
projective limit will be proved in \hyperref[VI.thm8.2.3]{8.2.3} when
$I$ is cofiltered.

\subsubsection{}\etiquette[8.1.3]{VI.subsubsec8.1.3}
We denote by
$$
\limtop{I} F \mbox{ or } \limtop{I}
F_i \mbox{ or also } \underleftarrow{F}
$$
a projective limit of the fibered topos $F$, by
$\mu : (\limtop{I} F) \times I \to F$ the canonical morphism, and,
for every object $i$ of $I$, by
$$
\mu_i : \underleftarrow{F} \simeq \underleftarrow{F} \times {i} \to F_{i}
$$
the morphism of topoi induced by $\mu$ on the fiber over $i$.

Choose a bisplitting of $F$ (\hyperref[VI.subsubsec7.1.4]{7.1.4}),
\ie for every $f : i \to j$ a morphism of topoi
$f\bigdot = (f^{\ast}, f_{\ast}) : F_i \to F_j$ such that $f^{\ast}$
is an inverse-image functor for the fibered structure. Then, for every
pair $i \xrightarrow{f} j \xrightarrow{g} k$ of composable morphisms
of $I$, we have an isomorphism of morphisms of topoi
(\hyperref[VI.subsubsec7.1.3]{7.1.3})
$$
c_{f, g} : g.f. \simeq (gf).\quad {\quoi}.
$$
Since $\mu$ is a cartesian morphism of fibered topoi, for every
morphism $f : i \to j$ of $I$ we have an isomorphism of morphisms of
topoi (\hyperref[VI.subsubsec7.1.6]{7.1.6})
\begin{equation*}
b_f  : f. \mu_i \xrightarrow{\sim} \mu_j{\quoi},\tag{8.1.3.1}
\end{equation*}\etiquette[8.1.3.1.]{VI.eq8.1.3.1}
and the family of the $b_f$ satisfies relations analogous to
\eqhyperref[VI.eq7.1.6.2]{7.1.6.2}.

\setcounter{paragraph}{1}
\paragraph{}\etiquette[8.1.3.2]{VI.par8.1.3.2}

Let $D$ be a $𝒰$-topos, let $m : D \times I \to F$ be a morphism of
fibered topoi, let $m_i : D = D \times \{i\} \to F_i$ be the
morphisms induced by $m$ on the fibers, and, for every $f : i \to j$,
let
$$
b'_f : f . m_i \xrightarrow{\sim} m_j{\quoi},
$$
be the transition isomorphism. It follows from
\hyperref[VI.subsubsec8.1.2]{8.1.2} that there is a morphism of topoi
$n : D \to \underleftarrow{F}$\pageoriginale and a family of
isomorphisms
$$
\gamma_i : m_i \xrightarrow{\sim} \mu_i n \quad i \in \ob I {\quoi},
$$
such that, for every $f : i \to j$,
$$
b_f \circ f. (\gamma_i) = \gamma_j \circ b'_f{\quoi}.
$$
Moreover, the morphism $n$ and the family of the $\gamma_i$,
$i \in \ob (I)$, are determined up to unique isomorphism
(\hyperref[VI.subsubsec8.1.2]{8.1.2}).

Conversely, suppose we are given a topos $\underleftarrow{F}$,
morphisms $\mu_i : \underleftarrow{F} \to F_i$, and isomorphisms $b_f$
as in \eqhyperref[VI.eq8.1.3.1]{8.1.3.1.}, satisfying the usual
compatibility conditions, and suppose all these data have the
universal property described above. Then there is a unique morphism
$\mu : \underleftarrow{F} \times I \to F$ of fibered topoi giving
rise to the $\mu_{i}$ and the $b_f$, and $(\underleftarrow{F}, \mu)$
is a projective limit of the fibered topos $F$.

\paragraph{}\etiquette[8.1.3.3]{VI.par8.1.3.3}

Let $𝒱$ be a universe such that $\sheaf{\Uu} \in 𝒱$. The preceding
considerations therefore allow projective limits of fibered topoi to
be interpreted as ``projective limits,'' in the sense of
2-categories, of the pseudofunctors with values in the 2-category of
fibered $𝒰$-topoi that are elements of $𝒱$, deduced from the fibered
structure by choosing the morphisms $f.$. The reader should take care,
however, not to confuse this notion of ``projective limit'' in the
sense of 2-categories (with commutation up to isomorphism) with the
strict (or naive) notion of projective limit (strict commutation of
diagrams). Even when the strict notion makes sense (when the
pseudofunctor is a genuine functor), generalized ``projective
limits'' and strict projective limits do not coincide in general.

\subsubsection{}\etiquette[8.1.4]{VI.subsubsec8.1.4}

Let $p : F \to I$ and $q : G \to I$ be two fibered $𝒰$-topoi, and
let $m : F \to G$ be a morphism of fibered topoi. Suppose that the
fibered topoi $F$ and $G$ have projective limits
(\hyperref[VI.def8.1.1]{8.1.1}). It follows immediately from
Definition \hyperref[VI.def8.1.1]{8.1.1} that there is a morphism of
topoi
\begin{equation*}
\underleftarrow{m} : \underleftarrow{F} \to
\underleftarrow{G}{\quoi},\tag{8.1.4.1}
\end{equation*}\etiquette[8.1.4.1]{VI.eq8.1.4.1}
and an isomorphism of morphisms of fibered topoi
$\gamma : m \circ \mu \xrightarrow{\sim}
\mu \circ \underleftarrow{m} \times \id_{I}$, making the diagram
\pageoriginale commutative:
\[\vcenter{
  \xymatrix@C=.6cm@R=.45cm{
\underleftarrow{F}\times I \ar[rrr]^{\mu}
\ar[ddd]_{\underleftarrow{m}\times \id_I}   && & F\ar[ddd]^{m}\\
                          &            & &\\
                          &            & &\ar@/_.55pc/@{-->}[ld]_{\gamma} \\
\underleftarrow{G}\times I \ar[rrr]^{\mu} &       & &  G&.
}}\leqno{(8.1.4.2)}
\]\etiquette[8.1.4.2]{VI.eq8.1.4.2}
Moreover, if $(\underleftarrow{m}', \gamma')$ is another pair with the
property above, there is a \emph{unique} isomorphism
$\eta: \underleftarrow{m} \xrightarrow{\sim} \underleftarrow{m}'$ of
morphisms of topoi such that
$$
\gamma \circ \mu (\eta \times \id) = \gamma'{\quoi}.
$$

\subsubsection{}\etiquette[8.1.5]{VI.subsubsec8.1.5}

Let $p : F \to I$ be a fibered $𝒰$-topos that is an element of $𝒱$,
let $\phi : I'\to I$ be a functor that is an element of $𝒱$, and let
$p_{\phi} : F_{\phi} \to I'$ be the fibered $𝒰$-topos deduced from
$(F, p)$ by the base change $\phi$
(\hyperref[VI.subsubsec7.1.9]{7.1.9}). Let
$(\underleftarrow{F}, \mu)$ be a projective limit of $F$. By base
change, we obtain a morphism of fibered topoi over $I'$:
$$
\mu' : \underleftarrow{F} \times I' \to F_{\phi}{\quoi}.
$$
Let $(\underleftarrow{F}_{\phi}, \mu_{\phi})$ be a projective limit of
$F_{\phi}$. By the universal property of the projective limit
(\hyperref[VI.subsubsec8.1.2]{8.1.2}), we then have a morphism of
topoi
$$
m_{\phi} : \underleftarrow{F} \to \underleftarrow{F}_{\phi}{\quoi},
$$
and an isomorphism of morphisms of fibered topoi over $I'$:
$$
\gamma : \mu'\xrightarrow{\sim} \mu_{\phi} \circ (m_{\phi} \times \id_{I'}){\quoi}.
$$
The pair $(m_{\phi}, \gamma)$ is determined up to unique isomorphism.

\subsubsection{}\etiquette[8.1.6]{VI.subsubsec8.1.6}

Let $p : F \to I$ be a fibered $𝒰$-site over a small category $I$,
let $F^{\sim/I} \to I$ be the associated fibered $𝒰$-topos
(\hyperref[VI.subsubsec7.2.6]{7.2.6}), and let $D$ be a $𝒰$-topos.
First, we have a functor
\begin{equation*}
\Hhomtop_{\cart / I} (D \times I, F^{\sim/I}) \to
\Mmorsite_{\cart/I}(D \times I, F){\quoi},\tag{8.1.6.1}
\end{equation*}\etiquette[8.1.6.1]{VI.eq8.1.6.1}
which associates with every $m$ the composite of $m^{\ast}$ with
$\epsilon_{F /I}$ (\hyperref[VI.prop7.2.7]{7.2.7}). Moreover, by the
definition of the category $\Mmorsite_{\cart /I} (D \times I, F)$
\eqhyperref[VI.eq7.2.2.2]{7.2.2.2}, we have a functor
\begin{equation*}
\Mmorsite_{\cart/I}(D \times I, F) \to \Hhom_{\cart / I} (F, D \times
I)^{\circ}{\quoi}.\tag{8.1.6.2}
\end{equation*}\etiquette[8.1.6.2]{VI.eq8.1.6.2}
Finally,\pageoriginale by the definition of the inductive limit
(\hyperref[VI.def6.3]{6.3}), we have a functor
\begin{equation*}
\Hhom_{\cart / I} (F, D \times I)^{\circ} \to
\Hhom(\varinjlim_{I^{\circ}}F, D)^{\circ}{\quoi}.\tag{8.1.6.3}
\end{equation*}\etiquette[8.1.6.3]{VI.eq8.1.6.3}
Denote by
\begin{equation*}
\Theta : \Hhomtop_{\cart / I}(D \times I, F^{\sim / I}) \to
\Hom(\varinjlim_{I^{\circ}}F, D)^{\circ}\tag{8.1.6.4}
\end{equation*}\etiquette[8.1.6.4]{VI.eq8.1.6.4}
the composite of these last three functors, and by
\begin{equation*}
\Pi : F \to \varinjlim_{I^{\circ}}F{\quoi},\tag{8.1.6.5}
\end{equation*}\etiquette[8.1.6.5]{VI.eq8.1.6.5}
the canonical functor.

\begin{enonce*}{Proposition 8.1.7}\etiquette[8.1.7]{VI.prop8.1.7}
The functor $\Theta$ is fully faithful. Its essential image is the set
of functors $m : \varinjlim_{I^{\circ}}F \to D$ such that the functor
$$
(m \circ \Pi, p) : F \to D \times I
$$
is a morphism from the total site $D \times I$ to the total site $F$.
\end{enonce*}

It follows from \hyperref[VI.prop7.2.7]{7.2.7} that the functor
\eqhyperref[VI.eq8.1.6.1]{8.1.6.1} is an equivalence of categories,
from \hyperref[VI.subsubsec7.2.2]{7.2.2} that the functor \hyperref[VI.eq8.1.6.2]{(8.1.6.2)}
is fully faithful, and from \hyperref[VI.def6.3]{6.3} that the functor
\eqhyperref[VI.eq8.1.6.3]{8.1.6.3} is an equivalence of categories.
Consequently, $\Theta$ is fully faithful. A functor
$m : \varinjlim_{I^{\circ}}F \to D$ is isomorphic to the image, under
$\Hhom_{\cart / I} (F, D \times I)^{\circ} \to
\Hhom(\varinjlim_{I^{\circ}}F, D)^{\circ}$, of the functor
$(m \circ \Pi, p)$, and such a functor belongs
to the essential image of \eqhyperref[VI.eq8.1.6.2]{8.1.6.2} if and
only if it is a morphism from the total site $D \times I$ to the total
site $F$ (\hyperref[VI.thm7.4.13.1]{7.4.13.1}); hence the proposition.

% END INCLUDED SGA4_English_current_buildable_source_20260906/source/06/drafts/SGA4_VI_LINES_5252_5495_English.texfrag


% BEGIN INCLUDED SGA4_English_current_buildable_source_20260906/source/06/drafts/SGA4_VI_LINES_5496_5856_English.texfrag
\subsection{Construction of the projective limit when the index
  category is cofiltered.}\etiquette[8.2]{VI.subsec8.2}

\begin{enonce*}{Lemma 8.2.1}\etiquette[8.2.1]{VI.lem8.2.1}
Let $I$ and $I'$ be two cofiltered categories \emph{(I \hyperref[I.sec8.1]{8.1})}, and let
$\phi:I'\to I$ be a functor such that
$\phi^{\circ}:{I'}^{\circ}\to I^{\circ}$ is cofinal \emph{(I \hyperref[I.sec8.1]{8.1})}.
Let, moreover, $p:F\to I$ be a fibered $𝒰$-topos, let $D$ be a
$𝒰$-topos, let $m:D\times I\to F$ be a Cartesian morphism of fibered
topoi, and let $m':D\times I'\to F'$ be the morphism of fibered topoi
deduced from $m$ by the base change $\phi$
(\hyperref[VI.subsubsec7.1.9]{7.1.9}). The pair $(D,m)$ is a
projective limit of $F$ if and only if $(D,m')$ is a projective limit
of $F'$.
\end{enonce*}

The proof\pageoriginale is merely a rather lengthy routine
verification. It is left to the patient reader, who may use the
viewpoint of projective limits of pseudofunctors
(\hyperref[VI.par8.1.3.3]{8.1.3.3}).

\begin{enonce*}{Lemma 8.2.2}\etiquette[8.2.2]{VI.lem8.2.2}
Let $p:F\to I$ be a fibered $𝒰$-site over a small cofiltered
category, let $\underrightarrow{F}$ be the inductive limit of the
fibered category $F$ \hyperref[VI.def6.3]{(6.3.9)}, and let
$\pi:F\to\underrightarrow{F}$ be the canonical functor. Equip
$\underrightarrow{F}$ with the coarsest topology making $\pi$
continuous \emph{(III 1)}. Let $D$ be a $𝒰$-site and
$n:\underrightarrow{F}\to D$ a functor. The following properties are
equivalent:
\begin{enumerate}
\renewcommand{\theenumi}{\roman{enumi}}
\renewcommand{\labelenumi}{\theenumi)}
\item $n$ is a morphism of sites $D\to\underrightarrow{F}$.

\item $n\circ\pi$ is a morphism of sites (where $F$ is equipped with
  the total topology).

\item $(n\circ\pi,p):F\to D\times I$ is a morphism from the total
  site $D\times I$ to the total site $F$.

\item For every object $i$ of $I$, the composite functor
  $F_i\xrightarrow{\alpha_{i!}}F\xrightarrow{n\circ\pi}D$ is a
  morphism of sites.
\end{enumerate}
\end{enonce*}

The equivalence of i) and ii) follows from
\hyperref[VI.exe6.6]{6.6}, and that of iii) and iv) follows from
\hyperref[VI.thm7.4.13.1]{7.4.13}. Let us prove the equivalence of iii) and ii). For this it is
enough to show that the first projection functor
$\pr_1:D\times I\to D$ is a morphism of sites. Denote by $D^{\sim}$
the topos of sheaves on $D$. It follows from
\hyperref[VI.prop7.4.7]{7.4.7} that $(D\times I)^{\sim}$ is equivalent
to the topos
$\Hhom_{I^{\circ}}(I^{\circ},D^{\sim}\times I^{\circ})=
\Hhom(I^{\circ},D^{\sim})$. Moreover, one sees immediately that the
natural extension of $\pr_1:D\times I\to D$ to sheaves (III \hyperref[III.prop1.2]{1.2}) is
the functor from $\Hhom(I^{\circ},D^{\sim})$ to $D^{\sim}$ that
associates with the functor $i\mapsto\sigma(i)$ the object
$\varinjlim_{I^{\circ}}\sigma(i)$. Since $I^{\circ}$ is filtered,
inductive limits over $I^{\circ}$ are exact (II 4), and hence
$\pr_1$ is a morphism of sites (IV \hyperref[IV.subsec4.9]{4.9}). It remains to prove the
equivalence of ii) and iv). First treat the case in which $D$ is a
$𝒰$-topos and $F$ a fibered $𝒰$-topos over $I$.

Let $i$ be an object of $I$. Arguing as in
\hyperref[VI.par7.4.13.6]{7.4.13.6}, one sees that the functor
$\alpha_{i!}:F_i\to F$ factors as a morphism
$\bar{\alpha}_{i!}:F_i\to F/\alpha_{i!}(e_i)$ followed by the
localization functor $F/\alpha_{i!}(e_i)\to F$, where $e_i$ is a
final object of $F_i$. It is therefore enough to show that the
composite functor
$F/\alpha_{i!}(e_i)\xrightarrow{n\circ\pi}D$ is a morphism of sites.
But there is a diagram commutative up to isomorphism:
\[
\xymatrix@C=2.5cm@R=1.25cm{
F/\alpha_{i!}(e_i) \ar[d]\ar[r]^{n\circ
  \pi/\alpha_{i!}(e_i)}& D/n\circ \pi \alpha_{i!}(e_i) \ar[d]\\
F \ar[r]^{n\circ \pi}& D
}
\]
Since\pageoriginale the property of being a morphism of sites is
local (IV \hyperref[IV.subsec5.10]{5.10}), it is enough to show that
$n\circ\pi\circ\alpha_{i!}(e_i)$ is a final object of $D$. For every
morphism $f:j\to k$ of $I$, there exists a unique morphism
$e_f:\alpha_{j!}(e_j)\to\alpha_{k!}(e_k)$ over $f$, and this
morphism is Cartesian. Thus $j\mapsto\alpha_{j!}(e_j)$ is a Cartesian
section of $F$ over $I$. Consequently, the canonical morphism
$n\circ\pi\circ\alpha_{i!}(e_i)\to
\varinjlim_I n\circ\pi\circ\alpha_{j!}(e_j)$ is an isomorphism,
because $\pi$ carries Cartesian morphisms to isomorphisms. Moreover,
$\varinjlim_I\alpha_{j!}(e_j)$, where the inductive limit is taken in
$F^{\widehat{\phantom{,}}}$, is a final object of
$F^{\widehat{\phantom{,}}}$. Since $n\circ\pi$ is a morphism of
sites, it carries the final object of $F^{\sim}$ to a final object of
$D$. Hence $\varinjlim_I n\circ\pi\circ\alpha_{j!}(e_j)$ is a final
object of $D$, completing the proof in this case.

In the general case, the functor $(n\pi,p):F\to D\times I$ is
Cartesian and continuous (\hyperref[VI.prop7.4.4]{7.4.4}). Moreover,
$n\circ\pi$ is the composite functor
$$
F \xrightarrow{(n \pi, p)} D \times I \xrightarrow{\pr_1} D{\quoi}.
$$
Passing to the associated fibered $𝒰$-topoi gives a diagram
commutative up to isomorphism:
\[
\xymatrix@C=2cm@R=1.75cm{
F       \ar@{-}[r]^{(n\pi,p)} \ar[d]_{\varepsilon_{F/I}}& D\times I
\ar[r]^{\pr_1}\ar[d]^{\varepsilon_D\times \id} & D \ar[d]^{\varepsilon_D}\\
F^{\sim/I}  \ar[r]^{(n\pi,p)^{\sim/I}} & D^{\sim}\times I \ar[r]^{\pr_1} & D^{\sim}&.
}
\]
Since $(n\pi,p)^{\sim/I}$ is Cartesian, it is of the form
$(n'\pi',p')$, where $p':F^{\sim/I}\to I$ and
$\pi':F^{\sim/I}\to\varinjlim_{I^{\circ}}F^{\sim}_i$ are the
canonical functors. Thus
$\pr_{1^{\circ}}(n\pi,p)^{\sim/I}=n'$. Since the functors
$\epsilon_{F/I}$ and $\epsilon_D$ induce equivalences on the topoi of
sheaves (\hyperref[VI.cor7.4.11]{7.4.11}), the functors $n$ and $n'$
induce isomorphic functors after passage to associated sheaves (III
\hyperref[III.prop1.2]{1.2}), and hence $n'$ is a morphism of sites. By what precedes,\pageoriginale
the functor $(n\pi,p)^{\sim/I}$ induces morphisms of topoi on the
fibers. It is therefore a morphism from the total site $F^{\sim/I}$
to the total site $D^{\sim}\times I$ \hyperref[VI.thm7.4.13.1]{(7.4.13)}. Since the functors
$\epsilon_{F/I}$ and $\epsilon_D\times\id$ induce equivalences on the
total topoi (\hyperref[VI.cor7.4.11]{7.4.11}), the functor $(n\pi,p)$
is a morphism of fibered sites \hyperref[VI.thm7.4.13.1]{(7.4.13)}, which proves the assertion.

\begin{enonce*}{Theorem 8.2.3}\etiquette[8.2.3]{VI.thm8.2.3}
Let $p:F\to I$ be a fibered $𝒰$-site over an essentially small
cofiltered category \emph{(I \hyperref[I.sec8.1]{8.1})}.
\begin{enumerate}
\renewcommand{\labelenumi}{\theenumi)}
\item Let $\underrightarrow{F}$ denote the site whose underlying
  category is the inductive limit of the fibered category $F\to I$
  \emph{(\hyperref[VI.def6.3]{6.3})} and whose topology is the
  coarsest topology making the canonical functor
  $\pi:F\to\underrightarrow{F}$ continuous. Then
  $\underrightarrow{F}$ is a $𝒰$-site and
$$
(\pi, p) : F \to \underrightarrow{F} \times I
$$
is a morphism from the fibered site $\underrightarrow{F}\times I$ to
the fibered site $F$.

\item Let $\mu:\underrightarrow{F}^{\sim}\times I\to F^{\sim/I}$ be
  the morphism of fibered $𝒰$-topoi deduced from $(\pi,p)$ by passing
  to the associated fibered $𝒰$-topoi. The pair
  $(\underrightarrow{F}^{\sim},\mu)$ is a projective limit of
  $F^{\sim/I}$.
\end{enumerate}
\end{enonce*}

\begin{enonce*}[remark]{Remark 8.2.4}\etiquette[8.2.4]{VI.rem8.2.4}
\begin{enumerate}
\renewcommand{\labelenumi}{\theenumi)}
\item In view of the notation introduced in
  \hyperref[VI.subsubsec8.1.3]{8.1.3}, one may set
$$
\limtop{I} F^{\sim} = \underrightarrow{F}^{\sim}.
$$

\item It follows from \hyperref[VI.prop7.4.4]{7.4.4} that the topology
  on $\underrightarrow{F}$ is also the coarsest topology making the
  family of functors
  $\pi_j:F_j\to F\xrightarrow{\pi}\underrightarrow{F}$,
  $j\in\ob I$, continuous.
\end{enumerate}
\end{enonce*}

\begin{enonce*}{Definition 8.2.5}\etiquette[8.2.5]{VI.def8.2.5}
The site $\underrightarrow{F}$ introduced in
\hyperref[VI.thm8.2.3]{8.2.3} is called the projective-limit site of
the fibered site $F$.
\end{enonce*}

As an exercise, the reader may, as in \hyperref[VI.lem8.2.1]{8.2.1},
define the notion of a projective limit of a fibered site and verify
that the site $\underrightarrow{F}$ is indeed a projective limit of
the site $F$. The second assertion of the theorem may therefore be
paraphrased as follows: the topos of sheaves on the projective-limit
site is equivalent to the projective limit of the associated fibered
topos.

\setcounter{subsubsection}{5}
\subsubsection{Proof of the theorem: Reduction to the case of a small
  index category}\etiquette[8.2.6]{VI.subsubsec8.2.6}

We\pageoriginale shall give only indications. Let $I'$ be a full
subcategory of $I$ such that ${I'}^{\circ}$ is cofinal in
$I^{\circ}$. Denote by $F'\to I'$ the fibered category deduced from
$F$ by the base change $I'\to I$, and by
$\pi':F'\to\underrightarrow{F}'$ the canonical functor. One first
observes that the categories $\underrightarrow{F}$ and
$\underrightarrow{F}'$ are equivalent and that the topology on
$\underrightarrow{F}'$ deduced from the topology of
$\underrightarrow{F}$ by this equivalence is the coarsest topology
making the functor $\pi':F'\to\underrightarrow{F}'$ continuous.
Assertion 1) of the theorem follows when it has been proved in the
case where $I$ is small. Assertion 2) in the general case then
follows from assertion 2) in the case where $I$ is small by Lemma
\hyperref[VI.lem8.2.1]{8.2.1}.

\subsubsection{End of the proof of the theorem}\etiquette[8.2.7]{VI.subsubsec8.2.7}
The first assertion follows from Lemma \hyperref[VI.lem8.2.2]{8.2.2},
applied with $D=\underrightarrow{F}$ and $n=\id$. Let $D$ be a
$𝒰$-topos. There is a diagram commutative up to isomorphism:
{\def\objectstyle{\scriptstyle}
\def\labelstyle{\scriptstyle}
\[
\xymatrix@C=.8cm@R=1.15cm{
\Hhomtop(D,{(\underrightarrow{F})}^{\sim}) \ar[r]^{(1)}\ar[d]_{(3)}&
\Mmorsite(D,\underrightarrow{F}) \ar@{^{(}->}[r]^{(2)} \ar[d]_{(3)}&
\Hhom(\underrightarrow{F},D)^{\circ}\ar[dd]_{(5)}\\
\Hhomtop_{\Cart/I}(D\times I,(\underrightarrow{F})^{\sim}\times
I)  \ar[r]^{(1)} \ar[d]_{(4)}& \Mmorsite_{\Cart/I}(D\times I,
(\underrightarrow{F})\times I) \ar[d]&\\
\Hhomtop_{\Cart/I}(D\times I,F^{\sim/I})
\ar[r]^{(1)} & \Mmorsite_{\Cart/I}(D\times I,F)  \ar@{^{(}->}[r]^{(2)} &
\Hhom_{\Cart/I}(F,D\times I)^{\circ},
}
\]}

\noindent where the functors of type (1) are equivalences (IV \hyperref[IV.prop4.9.4]{4.9.4}
and \hyperref[VI.prop7.2.7]{7.2.7}), the functors of type (2) are fully faithful (IV \hyperref[IV.subsubsec4.9.1]{4.9.1} and
\hyperref[IV.eq7.2.2]{7.2.2}), the functors of type (3) are base-change functors, the
functors of type (4) are composition functors (with $\mu$, \resp
$(\pi,p)$), and functor (5) is the equivalence deduced from the
universal property of the inductive limit
(\hyperref[VI.def6.3]{6.3}). To prove that the functor
$\Hhomtop(D,(\underrightarrow{F})^{\sim})\xrightarrow{(4)\circ(3)}
\Hhomtop_{/I}(D\times I,F^{\sim/I})$ is an equivalence, it is enough
to show that the functor
$$
\Mmorsite(D, \underrightarrow{F})
\xrightarrow{(4) \circ (3)} \Mmorsite_{/I}(D \times I,F)
$$
is an\pageoriginale equivalence of categories. Equivalently, it is
enough to show that the two fully faithful functors
\begin{align*}
& \Mmorsite(D, \underrightarrow{F}) \xrightarrow{(5) \circ (2)} \Hhomcart_I(F, D \times I)^{\circ}\\
& \Mmorsite_{/I}(D \times I, F) \xrightarrow{(2)} \Hhomcart_I(F, D \times
I)^{\circ}
\end{align*}
have the same essential images, which follows immediately from Lemma
\hyperref[VI.lem8.2.2]{8.2.2}.

\subsubsection{}\etiquette[8.2.8]{VI.subsubsec8.2.8}

Let $p:F\to I$ be a fibered $𝒰$-topos over a category $I$, and let
$(\limtop{I}F_i,\mu)$ be a projective limit of $F$. Thus there is a
Cartesian morphism of fibered $𝒰$-topoi over $I$:
$$
\mu:\limtop{I} F_i \times I \to F{\quoi},
$$
whence, on passing to the categories fibered over $I^{\circ}$ by the
direct image functors (\hyperref[VI.subsubsec7.1.3]{7.1.3}), a
Cartesian functor
$$
\mu_{\ast}: \limtop{I} F_i \times I^{\circ} \to F'
$$
which is the direct image functor for $\mu$
(\hyperref[VI.subsubsec7.1.7]{7.1.7}). By the definition of the
projective limit of a fibered category (\hyperref[VI.subsec6.10]{6.10}),
there is therefore a canonical functor
\begin{equation*}
\Theta : \limtop{I} F_i  \to
\varprojlim_{I, f_{\ast}} F_i = \Hhom_{\cart/I^{\circ}} (I^{\circ},
F'){\quoi}.\tag{8.2.8.1}
\end{equation*}\etiquette[8.2.8.1]{VI.eq8.2.8.1}

\begin{enonce*}{Theorem 8.2.9}\etiquette[8.2.9]{VI.thm8.2.9}
Use the notation of \hyperref[VI.subsubsec8.2.8]{8.2.8}, and suppose
that $I$ is an essentially small cofiltered category. Then the functor
$\Theta$ is an equivalence of categories. Suppose that $I$ is small.
Set $\limtop{I}F_i=(\underrightarrow{F})^{\sim}$ (which is permitted
by \hyperref[VI.thm8.2.3]{8.2.3} and
\hyperref[VI.rem8.2.4]{8.2.4}), and interpret the category
$\Hhom_{I^{\circ}}(I^{\circ},F')$ as the category of sheaves on the
total site $F$ (\hyperref[VI.prop7.4.7]{7.4.7}). Then the composite
functor
$$
(\underrightarrow{F})^{\sim} = \limtop{I} F_i
\xrightarrow{\Theta} \Hhom_{\cart/I^{\circ}} (I^{\circ}, F') \hookrightarrow \Hhom_{I^{\circ}}
(I^{\circ}, F') \approxeq F^{\sim}{\quoi},
$$
is precisely the direct image functor for the morphism of sites
defined by the functor
$$
\pi : F \to \underrightarrow{F}{\quoi}.
$$
\end{enonce*}

\setcounter{subsubsection}{9}
\subsubsection{}\etiquette[8.2.10]{VI.subsubsec8.2.10}

Some\pageoriginale comments before the proof of Theorem
\hyperref[VI.thm8.2.9]{8.2.9}. Besides solving the problem of the
existence of projective limits when the index category is cofiltered,
a solution which is curiously nontrivial in this theory, Theorems
\hyperref[VI.thm8.2.3]{8.2.3} and \hyperref[VI.thm8.2.9]{8.2.9}
provide two further pieces of information essential for working with
projective-limit topoi in practice. First, Theorem
\hyperref[VI.thm8.2.3]{8.2.3} interprets the projective-limit topos as
the topos of sheaves on a site whose underlying category is the
inductive-limit category of the categories $F_i$ under the system of
inverse image functors. In applications, one will be able to
interpret this category and its topology concretely. Second, Theorem
\hyperref[VI.thm8.2.9]{8.2.9} asserts that a sheaf on the projective
limit is determined by the system of its direct images in the topoi
$F_i$ (which gives a Cartesian section over $I^{\circ}$ of the
category $(F^{\sim/I})'$). The theorem further asserts conversely
that, when one is given for every object $i$ of $I$ a sheaf $X_i$ on
$F_i$ and, for every morphism $f:i\to j$ of $I$, an isomorphism
$X_j\simeq f_{\ast}X_i$, these isomorphisms being subject to
compatibility conditions, there exists essentially a unique sheaf on
the projective limit whose direct images give the system of the
$X_i$.

\subsubsection{}\etiquette[8.2.11]{VI.subsubsec8.2.11}

Assertion 2) of \hyperref[VI.thm8.2.3]{8.2.3} and the first assertion
of \hyperref[VI.thm8.2.9]{8.2.9} may be summarized by the striking
formula
\begin{equation*}
\varprojlim_{I, f_{\ast}}F_i^{\sim} \approxeq
\limtop{I, f.} F_i^{\sim} \approxeq
(\varinjlim_{I^{\circ}, f_{\ast}} F_i)^{\sim}{\quoi}. \tag{8.2.11.1}
\end{equation*}\etiquette[8.2.11.1]{VI.eq8.2.11.1}

\subsubsection{}\etiquette[8.2.12]{VI.subsubsec8.2.12}
\emph{Proof of Theorem} \hyperref[VI.thm8.2.9]{8.2.9}.
To prove the first assertion, reduce to the case where the category
$I$ is small by using \hyperref[VI.lem8.2.1]{8.2.1} and an analogous
lemma on projective limits of fibered categories. We leave the details
to the reader. Suppose henceforth that $I$ is a small category. We
know that the functor
$$
\pi : F \to \varinjlim_{I^{\circ}} F_i
$$
is a morphism of sites (\hyperref[VI.lem8.2.2]{8.2.2}). The direct
image functor
$$
\pi_{\ast} : (\varinjlim_{I^{\circ}} F_i)^{\sim} \to F^{\sim}
$$
that it\pageoriginale defines on the topoi is then fully faithful,
and its essential image is the set of sheaves on $F$ that carry the
Cartesian morphisms of $F$ to isomorphisms
(\hyperref[VI.exe6.6]{6.6}). These sheaves correspond, under the
equivalence
$F^{\sim}\simeq\Hhom_{I^{\circ}}(I^{\circ},F')$, to the Cartesian
sections of $F'$ over $I^{\circ}$
(\hyperref[VI.prop7.4.7]{7.4.7}). Moreover, for every
$j\in\ob(I)$, denote by $\alpha_{j!}:F_j\to F$ the inclusion functor,
which gives rise to the three functors
$\xymatrix{(\alpha_{j!},\alpha^{\ast}_j,\alpha_{j\ast}):F_j
\ar@<4pt>[r]\ar@<-4pt>[r]&F\ar[l]}$
(\hyperref[VI.rem7.4.3]{7.4.3}). It follows from the definition of the
equivalence $F^{\sim}\simeq\Hhom_{I^{\circ}}(I^{\circ},F')$
(\hyperref[VI.prop7.4.7]{7.4.7}) that a sheaf $X$ on
$\varinjlim_{I^{\circ}}F_i$ is carried by the functor
$(\varinjlim_{I^{\circ}}F_i)^{\sim}\xrightarrow{\pi_{\ast}}F^{\sim}
\simeq\Hhom_{I^{\circ}}(I^{\circ},F')$ to the Cartesian section
$j\mapsto\alpha^{\ast}_j\pi_{\ast}(X)$. But the description in
\hyperref[VI.thm8.2.3]{8.2.3} of the functor
$\mu:(\varinjlim_{I^{\circ}}F_i)^{\sim}\times I\to F$ shows that,
for every object $j$ of $I$, the direct image functor
$\mu_{j\ast}:(\varinjlim_{I^{\circ}}F_i)^{\sim}\to F_j$ for the
morphism $\mu_j:(\varinjlim_{I^{\circ}}F_i)^{\sim}\to F_j$ obtained
from $\mu$ by passing to the fiber at $j$ is precisely the functor
$\alpha^{\ast}_j\pi_{\ast}$. Consequently, the functor
$(\varinjlim_{I^{\circ}}F_i)^{\sim}\xrightarrow{\pi_{\ast}}F^{\sim}
\simeq\Hhom_{I^{\circ}}(I^{\circ},F')$ is precisely the functor
$(\varinjlim_{I^{\circ}}F_i)^{\sim}\xrightarrow{\Theta}
\Hhom_{\cart/I^{\circ}}(I^{\circ},F')\hookrightarrow
\Hhom_{I^{\circ}}(I^{\circ},F')$, which proves the theorem.


% END INCLUDED SGA4_English_current_buildable_source_20260906/source/06/drafts/SGA4_VI_LINES_5496_5856_English.texfrag


% BEGIN INCLUDED SGA4_English_current_buildable_source_20260906/source/06/drafts/SGA4_VI_LINES_5857_6042_English.texfrag
\subsection{Topology of the projective-limit site: the case of
  coherent topoi}\etiquette[8.3]{VI.subsec8.3}

\subsubsection{}\etiquette[8.3.1]{VI.subsubsec8.3.1}

In this section, $p:F\to I$ is a fibered $\Uu$-site over an
essentially small cofiltered category, whose inverse-image functors
$f^{\ast}:F_i\to F_j$ commute with fiber products. In addition, for
every object $X$ of a fiber category $F_i$, we are given a set
$\Cov(X)$ of covering families for the topology of $F_i$ which have
properties (PTO) and (PT1) of II \hyperref[II.def1.3]{1.3} and generate the topology of
$F_i$. Finally, suppose that for every $f:i\to j$, every
$Y\in\ob F_j$, and every family
$(Y_{\alpha}\to Y)_{\alpha\in A}$ in $\Cov(Y)$, the family
$(f^{\ast}(Y_{\alpha})\to f^{\ast}(Y))_{\alpha\in A}$ belongs to
$\Cov(f^{\ast}(Y))$. These conditions on $\Cov(X)$ are satisfied, for
example, if one takes $\Cov(X)=$ \{all covering families of $X$ in
$F_i$\}. In\pageoriginale this section we shall study the
projective-limit site of $F$ under suitable finiteness hypotheses
(\hyperref[VI.thm8.3.13]{8.3.13}).

\subsubsection{}\etiquette[8.3.2]{VI.subsubsec8.3.2}

Let $\pi:F\to\underrightarrow{F}$ be the canonical functor, and let
$Y$ be an object of $\underrightarrow{F}$. Denote by $\Cov(Y)$ the set
of families $(Y_{\alpha}\xrightarrow{n_{\alpha}}Y)_{\alpha\in A}$
of the following form:

There exist $i\in I$, an object $X$ of $F_i$, a family
$(X_{\alpha}\xrightarrow{{n'}_\alpha}X)_{\alpha\in A}\in\Cov(X)$,
an isomorphism $\phi:\pi(X)\simeq Y$, and a family of isomorphisms
$\phi_{\alpha}:\pi(X_{\alpha})\simeq Y_{\alpha}$, $\alpha\in A$,
such that for every $\alpha\in A$ the diagram
\[
\xymatrix@C=2.5cm@R=2cm{
\pi(X_{\alpha})
\ar@{-}[r]^{\displaystyle{\mathop{\sim}^{\phi_{\alpha}}}} \ar[d]^{\pi(n'_{\alpha})}& Y_{\alpha}\ar[d]^{n_{\alpha}} \\
\pi(X)  \ar@{-}[r]^{\displaystyle{\mathop{\sim}^{\phi}}} & Y
}
\]
is commutative.

\begin{enonce*}{Proposition 8.3.3}\etiquette[8.3.3]{VI.prop8.3.3}
\begin{enumerate}
\renewcommand{\labelenumi}{\theenumi)}
\item The families belonging to $\Cov(Y)$, for
  $Y\in\ob\underrightarrow{F}$, generate the topology of the
  projective-limit site (\hyperref[VI.def8.2.5]{8.2.5}).

\item The family $Y\mapsto\Cov(Y)$, $Y\in\ob\underrightarrow{F}$,
  has properties \emph{(PTO)} and \emph{(PT1)} of \emph{II \hyperref[II.def1.3]{1.3}}.
\end{enumerate}
\end{enonce*}

We shall need the following lemma.

\begin{enonce*}{Lemma 8.3.4}\etiquette[8.3.4]{VI.lem8.3.4}
Let $i$ be an object of $I$, and let $n:X'\to X$ be a base-changeable
morphism of $F_i$ such that $f^{\ast}(n)$ is base-changeable for every
$f:j\to i$. Then $\pi(n)$ is base-changeable. The functors
$$
\pi_j:F_j\hookrightarrow F\xrightarrow{\pi}\underrightarrow{F},
\quad j\in\ob I{\quoi},
$$
commute with fiber products.
\end{enonce*}

Recall that every object $Y$ of $\underrightarrow{F}$ is equal to an
object $\pi(Z)$, $Z\in\ob F$, and that every morphism
$m:Y\to\pi(X)$ is of the form
$$
\pi(Z)\xrightarrow{\pi(s)^{-1}}\pi(Z')
\xrightarrow{\pi(m')}\pi(X){\quoi},
$$
where $s$ is a cartesian morphism of $F$. Consequently, to show that
$\pi(n)$ is base-changeable,\pageoriginale it is enough to show that
the fiber product of every diagram
\[
\xymatrix@C=2.5cm@R=1.25cm{
                           &\pi(X')\ar[d]^{\pi(n)}\\
\pi(Z')\ar[r]^{\pi(m')}&\pi(X)
}
\]
is representable. The morphism $m':Z'\to X$ factors as
$Z'\xrightarrow{m''}X''\xrightarrow{s'}X$, where $p(m'')$ is the
identity of $p(Z')=j$ and $s'$ is cartesian over $p(s')=f$. Put
$X'''=f^{\ast}(X')$ and $n'=f^{\ast}(n)$. We then have a commutative
diagram
\[
\xymatrix@C=2.5cm@R=1.5cm{
X'''\ar[r]^{s''}\ar[d]^{n'}&X'\ar[d]^{n}\\
X''\ar[r]^{s'}&X,
}
\]
and hence a commutative diagram in $\underrightarrow{F}$:
\[
\xymatrix@C=2.5cm@R=1.5cm{
        &\pi(X''')\ar[r]^{\sim}\ar[d]_{\pi(n')}&\pi(X')\ar[d]^{\pi(n)}\\
\pi(Z')\ar@{-}[r]^{\pi(m'')}&\pi(X'')\ar[r]^{\sim}&\pi(X)
}
\]
where $n'$ is a base-changeable morphism of $F_j$. Thus, to show that
$\pi(n)$ is base-changeable, it is enough to show that the product
$\pi(Z')\times_{\pi(X'')}\pi(X''')$ is representable, and therefore
it is enough to show that
$\pi_j:F_j\hookrightarrow F\to\underrightarrow{F}$ commutes with
fiber products.

Let
\[
\xymatrix@C=2.5cm@R=1.25cm{
Y'\ar[r]\ar[d]&X'\ar[d]\\
Y\ar[r]&X
}
\]
be a cartesian diagram in $F_j$, and let $W$ be an object of $F$ over
$k\in\ob I$. We have (\hyperref[VI.prop6.5]{6.5}):
\def\labelstyle{\scriptstyle}
$$\displaystyle{\Hhom(\pi(W),\pi(Y'))\simeq\mathop{\varinjlim_{
\xymatrix@R=.3cm@C=.3cm@M=.05cm{
            &j\\
        l\ar[ru]^{f}\ar[rd]_{g}&\\
            &k
}}}\Hhom_{F_{\ell}}(g^{\ast}(W),f^{\ast}(Y'))}$$

Since\pageoriginale $f^{\ast}:F_j\to F_{\ell}$ commutes with fiber
products, we have
$$
\Hom(g^{\ast}(W),f^{\ast}(Y'))\simeq
\Hom(g^{\ast}(W),f^{\ast}(Y))
\times_{\Hom(g^{\ast}(W),f^{\ast}(X))}
\Hom(g^{\ast}(W),f^{\ast}(X')).
$$
Using the commutation of filtered limits with fiber products (I 2)
and formula (\hyperref[VI.prop6.5]{6.5}), we obtain an isomorphism
$$
\Hom(\pi(W),\pi(Y'))\simeq\Hom(\pi(W),\pi(Y))
\times_{\Hom(\pi(W),\pi(X))}\Hom(\pi(W),\pi(X')){\quoi},
$$
which shows that the diagram
\[
\xymatrix@C=2.5cm@R=1.5cm{
\pi(Y')\ar[r]\ar[d]&\pi(X')\ar[d]\\
\pi(Y)\ar[r]&\pi(X)
}
\]
is cartesian in $\underrightarrow{F}$.

\setcounter{subsubsection}{4}
\subsubsection{}\etiquette[8.3.5]{VI.subsubsec8.3.5}
\emph{Proof of} \hyperref[VI.prop8.3.3]{8.3.3}. We first prove the
second assertion. It follows from \hyperref[VI.lem8.3.4]{8.3.4} that
the families in $\Cov(Y)$, $Y\in\underrightarrow{F}$, consist of
base-changeable morphisms, whence property (PTO). To show that these
families have property (PT1) (stability under base change), a series
of reductions like those in the proof of
\hyperref[VI.prop8.3.3]{8.3.3} reduces us to proving the following
assertion:

For every $i\in\ob I$, every $Y\in\ob F_i$, every
$(Y_{\alpha}\to Y)_{\alpha\in A}\in\Cov(Y)$, and every morphism
$u:Z\to Y$ of $F_i$, the family
$$
(\pi(Z)\times_{\pi(Y)}\pi(Y_{\alpha})\to\pi(Z))_{\alpha\in A}
$$
belongs to $\Cov(\pi(Z))$.

This assertion follows from the fact that the
$\pi_i:F_i\hookrightarrow F\xrightarrow{\pi}\underrightarrow{F}$
commute with fiber products (\hyperref[VI.lem8.3.4]{8.3.4}) and that
the families $X\mapsto\Cov(X)$ of $F_i$ are stable under base change.
Let us prove the first assertion. It is clear that the families
belonging to the $\Cov(Y)$, $Y\in\underrightarrow{F}$, are covering
for the coarsest topology $T$ making the functor $\pi$ continuous
(III 1). Thus the topology $T'$ generated by these families is
coarser than $T$. To show that it is finer than\pageoriginale $T$
(and hence equal to $T$), it is enough to show that every sheaf $M$
for $T'$ is such that $M\circ\pi$ is a sheaf on $F$, or equivalently
(\hyperref[VI.rem8.2.4]{8.2.4}) that $M\circ\pi_i$ is a sheaf on
$F_i$ for every $i\in\ob I$. Since the functors
$\pi_i:F_i\to\underrightarrow{F}$ commute with fiber products
(\hyperref[VI.lem8.3.4]{8.3.4}), $M$ is a sheaf for $T'$ if and only
if (II \hyperref[II.cor2.3]{2.3} and I \hyperref[I.prop2.12]{2.12}), for every $i\in\ob I$, every $Y\in\ob F_i$,
and every $(Y_{\alpha}\to Y)_{\alpha\in A}\in\Cov(Y)$, the sequence
of sets
$$
M(\pi_i(Y))\to\prod_{\alpha\in A}M(\pi_i(Y_{\alpha}))
\rightrightarrows\prod_{(\alpha,\beta)\in A\times A}
M\bigl(\pi_i(\fprod{Y_{\alpha}}{Y_{\beta}}{Y})\bigr)
$$
is exact, \ie (\loccit) if and only if $M\circ\pi_i$ is a sheaf on
$F_i$ for every $i\in\ob I$.


% END INCLUDED SGA4_English_current_buildable_source_20260906/source/06/drafts/SGA4_VI_LINES_5857_6042_English.texfrag


% BEGIN INCLUDED SGA4_English_current_buildable_source_20260906/source/06/drafts/SGA4_VI_LINES_6043_6336_English.texfrag
\begin{enonce*}{Proposition 8.3.6}\etiquette[8.3.6]{VI.prop8.3.6}
We use the notation and hypotheses of \hyperref[VI.subsubsec8.3.1]{8.3.1}
and \hyperref[VI.subsubsec8.3.2]{8.3.2}. If, for every $i\in\ob I$,
$X\mapsto\Cov(X)$ is a pretopology on $F_i$ \emph{(II \hyperref[II.def1.3]{1.3})}, and if for
every $X\in\ob F_i$ the covering families in $\Cov(X)$ are finite, then for
every $Y\in\underrightarrow{F}$ the covering families in $\Cov(Y)$ are
finite and $Y\mapsto\Cov(Y)$ is a pretopology on $\underrightarrow{F}$.
\end{enonce*}

\setcounter{subsubsection}{6}
\subsubsection{}\etiquette[8.3.7]{VI.subsubsec8.3.7}
The only thing to prove is that the families $Y\mapsto\Cov(Y)$ have property
(PT2) of (II \hyperref[II.def1.3]{1.3}) (stability under composition). Let $i\in\ob I$, let $Y$
be an object of $F_i$, and let
$(Y_\alpha\to Y)_{\alpha\in A}\in\Cov(Y)$. In addition, let
$i_\alpha$, $\alpha\in A$, be a family of objects of $I$; for every
$\alpha$, let $Z_\alpha$ be an object of $F_i$; and let
$(Z_{\alpha\beta}\xrightarrow{n_{\alpha\beta}}Z_\alpha)_{\alpha\beta\in
B_\alpha}$ be a family in $\Cov(Z_\alpha)$. Finally, for every
$\alpha\in A$, let an isomorphism
$\phi:\pi(Z_\alpha)\xrightarrow{\sim}\pi(Y_\alpha)$ be given. Returning to
the definition of the families in $\Cov(X)$ for
$X\in\ob\underrightarrow{F}$
(\hyperref[VI.subsubsec8.3.2]{8.3.2}), we see that we must prove that the
family
$$
(\pi (Z_{\alpha \beta}) \xrightarrow{\pi(n_{\alpha}) \circ \phi_{\alpha} \circ
  \pi(n_{\alpha \beta})} \pi(Y))_{\alpha \beta \in \mathop{\coprod}_A
  B_{\alpha}}
$$
belongs to $\Cov(\pi(Y))$. We shall do this after five reductions.

\setcounter{subsubsection}{7}
\subsubsection{}\etiquette[8.3.8]{VI.subsubsec8.3.8}
\emph{Reduction to the case in which, for every $\alpha\in A$, we have}
$\phi_\alpha=\pi(m_\alpha)$, $m_\alpha:Z_\alpha\to Y_\alpha$.
\emph{For every $\alpha\in A$, we have
$\phi_\alpha=\pi(m_\alpha)\pi(s_\alpha)^{-1}$
$(\hyperref[VI.prop6.5]{6.5})$, where $s_\alpha$ is a Cartesian morphism
over $f_\alpha:i'_\alpha\to i_\alpha$. Thus there are commutative diagrams}
\[
\xymatrix@C=2.5cm@R=1.25cm{
Z_{\alpha\beta} \ar[d]_{n_{\alpha\beta}} & f^{\ast}(Z_{\alpha\beta})
\ar[l]_{s_{\alpha\beta}} \ar[d]^{f^{\ast}_{\alpha}(n_{\alpha\beta})}\\
Z_{\alpha} &   f^{\ast}_{\alpha}(Z_{\alpha}) \ar[l]_{s_{\alpha}},  &
**[l]\alpha\beta\in \displaystyle{\coprod_A} B_{\alpha},
}
\]
which\pageoriginale give commutative diagrams
\[
\xymatrix@C=2.5cm@R=1.25cm@H=.2cm{
\pi(Z_{\alpha\beta})  \ar[r]^{\sim} \ar[d]^{\pi(n_{\alpha\beta})}&
\pi(f^{\ast}_{\alpha}(Z_{\alpha\beta})) \ar[d]^{\pi(f^{\ast}_{\alpha}(m_{\alpha\beta}))}\\
\pi(Z_{\alpha})  \ar[r]^{\displaystyle\mathop{\sim}^{\pi(S_\alpha)^{-1}}} & \pi(f^{\ast}_{\alpha}(Z_{\alpha}))
\ar[r]^{\pi(m_{\alpha})} & \pi(Y_\alpha),\quad
\alpha\beta\in \displaystyle{\coprod_A} B_{\alpha}.}
\]
Since the families
$(f^{\ast}_{\alpha}(Z_{\alpha\beta})\xrightarrow{f^{\ast}(n_{\alpha})}
f^{\ast}(Z_{\alpha}))$ belong to
$\Cov(f^{\ast}_{\alpha}(Z_{\alpha}))$, we may reduce to the case in which
$\phi_\alpha=\pi(m_\alpha)$ for every $\alpha$.

\subsubsection{}\etiquette[8.3.9]{VI.subsubsec8.3.9}

\emph{Reduction to the case in which, for every $\alpha\in A$,
$i_\alpha=j$ and} $p(m_\alpha)=k:j\to i$. Put
$g_\alpha=p(m_\alpha)$. Since $I$ is cofiltered and $A$ is finite, there
exist an object $j$ of $I$ and morphisms $h_\alpha:j\to i_\alpha$ such that,
for every pair $(\alpha,\alpha')$, we have
$g_\alpha h_\alpha=g_{\alpha'}h_{\alpha'}=k$. Using the inverse-image
functors $h_\alpha^{\ast}$ as in \hyperref[VI.subsubsec8.3.8]{8.3.8}, we
reduce to the stated case.

\subsubsection{}\etiquette[8.3.10]{VI.subsubsec8.3.10}

\emph{Reduction to the case in which} $i=j$ and $k:j\to i$ \emph{is the
identity}. For every $\alpha$, there is a commutative diagram
\[
\xymatrix@C=2.5cm@R=1.5cm{
Z_{\alpha}  \ar[r]^{m'_\alpha} &  k^{\ast}(Y_{\alpha}) \ar[d]_{k^{\ast}(n_{\alpha})}
\ar[r]^{t_{\alpha}} & Y_{\alpha} \ar[d]^{n_{\alpha}}\\
                    & k^{\ast}(Y)  \ar[r]^{t} & Y,
}
\]
where the morphisms $t$ and $t_\alpha$, $\alpha\in A$, are Cartesian,
$m'_\alpha$ lies over the identity of $j$, and
$t_\alpha\circ m'_\alpha=m_\alpha$. Replacing the family
$(Y_\alpha\to Y)\alpha\in A$ by the\pageoriginale family
$(k^{\ast}(Y_\alpha)\to k^{\ast}(Y))_{\alpha\in A}\in
\Cov(k^{\ast}(Y_\alpha))$, and the morphisms $m_\alpha$ by the morphisms
$m'_\alpha$, we reduce to the stated case.

\begin{enonce*}{Lemma 8.3.11}\etiquette[8.3.11]{VI.lem8.3.11}
Let $j$ be an object of $I$ and let $m:X'\to X$ be a morphism of $F_j$ such
that $\pi(m)$ is an isomorphism. There exists a morphism $\ell:j'\to j$ such
that $\ell^{\ast}(m)$ is an isomorphism.
\end{enonce*}

Let us show that there exist a morphism $f:i\to j$ and a morphism
$q:f^{\ast}(X)\to f^{\ast}(X')$ such that
$f^{\ast}(m)q=\id_{f^{\ast}(X)}$ and such that $\pi(q)$ is an isomorphism.
Indeed, there exists an isomorphism $n:\pi(X)\to\pi(X')$ such that
$\pi(m)\circ n=\id_{\pi(X)}$. There consequently exist
(\hyperref[VI.prop6.5]{6.5}) two morphisms
$X\xleftarrow{s_1}Y_1\xrightarrow{q_1}X'$, where $s_1$ is Cartesian, such
that $n=\pi(q_1)\pi(s_1)^{-1}$. Thus
$\pi(m)\circ\pi(q_1)=\pi(s_1)$. The morphisms $mq_1$ and $s_1$ from $Y_1$
to $X$ have the same image in $\underrightarrow{F}$. Consequently (\cf the
description of the morphisms in the inductive limit
(\hyperref[VI.prop6.5]{6.5})), there exists a Cartesian morphism
$s_2:Y_2\to Y_1$ such that $mq_1s_2=s_1s_2$. Put $q_1s_2=q_2$ and denote by
$s_3$ the Cartesian morphism $s_1s_2$. Thus $mq_2=s_3$. The morphism
$q_2:Y_2\to X'$ factors in an essentially unique manner as
$Y_2\xrightarrow{q}Y_3\xrightarrow{s_4}X'$, where $s_4$ is Cartesian and
$q$ lies over the identity of $p(Y_2)$. Thus there is a commutative diagram
\[\vcenter{
\xymatrix@C=2.25cm@R=1.5cm{
Y_3 \ar[d]_{s_4} & Y_2 \ar[l]_{q}\ar[d]^{s_3}\\
X'  \ar[r]^{m} & X.          .
}}\leqno{(8.3.11.1)}
\]\etiquette[8.3.11.1]{VI.eq8.3.11.1}
Put $p(s_4)=p(s_3)=f:i\to j$. We have
$f^{\ast}(X')\simeq Y_3$, $f^{\ast}(X')\simeq Y_2$, and
$f^{\ast}(m):Y_2\to Y_3$ makes the following diagram commutative:
\[
\xymatrix@C=2.25cm@R=1.5cm{
Y_3  \ar[r]^{f^{\ast}(m)} \ar[d]_{s_4} & Y_2 \ar[d]^{s_3}\\
X'  \ar[r]^{m} & X.
}
\]
Thus\pageoriginale $s_3=s_3(f^{\ast}(m)\circ q)$; since $s_3$ is Cartesian
and $f^{\ast}(m)\circ q$ lies over the identity of $p(Y_2)$, we have
$f^{\ast}(m)\circ q=\id_{Y_2}$. Moreover, the commutativity of diagram
\eqhyperref[VI.eq8.3.11.1]{8.3.11.1} shows that $\pi(q)$ is an isomorphism.

Now apply the preceding result to the morphism
$q:f^{\ast}(X)\to f^{\ast}(X')$: there exist a morphism $f':j'\to i$ and a
morphism $q':f'^{\ast}(X')\to f'^{\ast}(X)$ such that
$f'^{\ast}(q)q'=\id_{f'^{\ast}(f'^{\ast}(X'))}$. But on putting
$\ell^{\ast}=ff'$, there is an isomorphism
$f'^{\ast}f^{\ast}\simeq\ell^{\ast}$. Thus there is a morphism
$f'^{\ast}(q):\ell^{\ast}(X)\to\ell^{\ast}(X')$ and an isomorphism
$f'^{\ast}f^{\ast}\simeq\ell^{\ast}$. Thus there is a morphism
$f'^{\ast}(q):\ell^{\ast}(X)\to\ell^{\ast}(X')$ and a morphism
$q':\ell^{\ast}(X')\to\ell^{\ast}(X)$ such that
$\ell^{\ast}(m){f'}^{\ast}(q)=\id_{\ell^{\ast}(X)}$ and
${f'}^{\ast}(q)q'=\id_{\ell^{\ast}(X')}$. Consequently
$\ell^{\ast}(m)$ is an isomorphism.

\setcounter{subsubsection}{11}
\subsubsection{}\etiquette[8.3.12]{VI.subsubsec8.3.12}
\emph{End of the proof}. Since $I$ is cofiltered and $A$ finite, it follows
from \hyperref[VI.lem8.3.11]{8.3.11} that there exists an inverse-image
functor $\ell^{\ast}$ that takes all the morphisms $m_\alpha$,
$\alpha\in A$, to isomorphisms. We are thus reduced to the case in which the
$m_\alpha$ are isomorphisms. But then the covering families
$(Z_{\alpha\beta}\to Z)_{\alpha\beta\in B_\alpha}$ are deduced by the base
change $m_\alpha:Z_\alpha\to Y_\alpha$ from covering families
$(Y_{\alpha\beta}\to Y_\alpha)_{\alpha\beta\in B_\alpha}\in
\Cov(Y_\alpha)$. We are therefore reduced to the case in which
$Z_\alpha=Y_\alpha$ and $\Phi_\alpha=\id_{\pi(Y_\alpha)}$. But in this case
the family
$(Y_{\alpha\beta}\xrightarrow{n_\alpha\circ n_{\alpha\beta}}Y)_{
\alpha\beta\in\mathop{\coprod}_A B_\alpha}$ belongs to $\Cov(Y)$ (property
(PT2)). Hence its image under $\pi$ belongs to $\Cov(\pi(Y))$
(\hyperref[VI.subsubsec8.3.2]{8.3.2}).

\begin{enonce*}{Theorem 8.3.13}\etiquette[8.3.13]{VI.thm8.3.13}
Let $p:F\to I$ be a $𝒰$-topos fibered over an essentially small cofiltered
category. If, for every object $i$ of $I$, the fiber topos $F_i$ is coherent
$(\hyperref[VI.def2.3]{2.3})$, and if, for every morphism $f:i\to j$, the
morphism of topoi
$f\bigdot=(f^{\ast},f_{\ast}):F_i\to F_j$ is coherent
(\hyperref[VI.def3.1]{3.1}), then the topos $\limtop{I}F_i$ is coherent and,
for every object $i$ of $I$, the canonical morphism of topoi
$$
\mu_i : \limtop{I} F_i \to F_i
$$
is coherent.

Moreover, if $F_{\coh}$ denotes the full subcategory of $F$ defined by the
objects of $F$ that are coherent in their fiber
(\hyperref[VI.def1.13]{1.13}), then the category $F_{\coh}$ is fibered over
$I$, and $\underrightarrow{F}_{\coh}$ is canonically equivalent to the
category of coherent objects of $\limtop{I}F_i$.
\end{enonce*}

For every\pageoriginale $i\in\ob I$, denote by $(F_i)_{\coh}$ the $𝒰$-site
(endowed with the induced topology) of coherent objects of $F_i$
(\hyperref[VI.def1.13]{1.13}). Since $F_i$ is coherent, $F_i$ is the topos
of sheaves on $(F_i)_{\coh}$. For every morphism $f:i\to j$, the morphisms
$f\bigdot=(f^{\ast},f_{\ast})$ are coherent, and consequently
$f^{\ast}(F_j)_{\coh}\subset(F_i)_{\coh}$. Thus there is a fibered $𝒰$-site
$p_{\coh}:F_{\coh}\to I$ whose fiber sites are the sites $(F_i)_{\coh}$ and
whose associated fibered topos is equivalent to $p:F\to I$. Consequently
(\hyperref[VI.thm8.2.3]{8.2.3}),
\begin{equation*}
\limtop{I, f.} F_i \approxeq
(\underrightarrow{F}_{\coh})^{\sim}\tag{8.3.13.1}
\end{equation*}\etiquette[8.3.13.1]{VI.eq8.3.13.1}
Since the topoi $F_i$ are coherent, finite products and fiber products are
representable in $(F_i)_{\coh}$ (\hyperref[VI.prop2.2]{2.2}). Moreover, the
functors $f^{\ast}:(F_i)_{\coh}\to(F_j)_{\coh}$ are left exact. Finally,
since the objects of $(F_i)_{\coh}$ are quasi-compact, the finite covering
families form a pretopology. It then follows from
(\hyperref[VI.lem8.3.4]{8.3.4}) that finite products and fiber products are
representable in $\varprojlim_{I^{\circ},F}(F_i)_{\coh}$ (addendum to Lemma
\hyperref[VI.lem8.3.4]{8.3.4}: show that the canonical functors
$(F_i)_{\coh}\to\underrightarrow{F}_{\coh}$ take the final object to the
final object), and from (\hyperref[VI.prop8.3.6]{8.3.6}) that the objects of
$\underrightarrow{F}_{\coh}$ are quasi-compact. Consequently
(\hyperref[VI.subsubsec2.4.5]{2.4.5}), $\limtop{I}(F_i)$ is coherent.
Finally, the canonical morphisms
$\mu_i:\limtop{I}(F_i)\to F_i$ are deduced, by passage to the topoi of
sheaves, from the morphisms of sites (\hyperref[VI.thm8.2.3]{8.2.3})
$$
(F_i)_{\coh} \to \underrightarrow{F}_{\coh}
$$
and hence (\hyperref[VI.cor3.3]{3.3}) the morphisms $\mu_i$ are coherent.

Let us prove the last assertion. We have \eqhyperref[VI.eq8.3.13.1]{8.3.13.1}
$$
\limtop{ℐ} (F_i) \simeq
(\underrightarrow{F}_{\coh})^{\sim}{\quoi}.
$$
Let $X$ be a coherent object of $(F_{\coh})^{\sim}$. Since $X$ is
quasi-compact, there exists a finite covering family
$(Y_\alpha\to X)$, where
$Y_\alpha\in\ob\underrightarrow{F}_{\coh}$
(\hyperref[VI.def1.1]{1.1}). On taking an index $i\in\ob I$ sufficiently
small, we may suppose that the $Y_\alpha$ are the images under
$\mu_i:F_{i\coh}\to\underrightarrow{F}_{\coh}$ of a family $Y'_\alpha$;
hence, by taking the direct sum of the $Y'_\alpha$
(\hyperref[VI.cor1.15]{1.15}), we see that there exist an index $i\in\ob I$,
an object $Y'\in F_{i\coh}$, and a surjective morphism
$\mu_i(Y')\to X$. The equivalence relation
$R=\fprod{\mu_i(Y')}{\mu_i(Y')}{X}$ is then a coherent
object\pageoriginale because $(\underrightarrow{F}_{\coh})^{\sim}$ is a
coherent topos (\hyperref[VI.prop2.2]{2.2}). Let us first show that $R$ is an
object of $\underrightarrow{F}_{\coh}$. By definition, $R$ is a subobject of
$\mu_i(Y'\times Y')$, and since $R$ is coherent, we may, after changing the
index $i$, find an arrow $m:Z\to(Y'\times Y')$ of $F_{i\coh}$ such that
$\text{Im}(\mu_i(m))=R$. It follows that, since $F_i$ is coherent,
$\text{Im}(m)$ is coherent in $F_i$ (\hyperref[VI.cor1.17.1]{1.17.1}), and
hence belongs to $F_{i\coh}$. There consequently exist an index $i$ and a
diagram
$\def\labelstyle{\scriptstyle} \xymatrix@C=1cm{
R \ar@<3pt>[r]^{p_1}\ar@<-3pt>[r]_{p_2} & Y'
}
$
of $F_{i\coh}$ such that
$\def\labelstyle{\scriptstyle}
\xymatrix@C=1cm{
  \mu_i(R)\ar@<3pt>[r]^{\mu_i(p_1)} \ar@<-3pt>[r]_{\mu_i(p_2)}& \mu_i(Y')
}
$
is an equivalence relation and such that
$X=\mu_i(Y')/\mu_i(R)$. Put $X_i=\coker(p_1,p_2)$ and
$R'=\fprod{Y'}{Y'}{X'_i}$. There is a canonical morphism in $F_i$,
$u:R\to R'$. Moreover, $\mu_i^{\ast}(u)$ is an isomorphism. Thus, after
changing the index $i$, we may suppose that $R=R'$, \ie that $R$ is an
equivalence relation. But then $X_i$ belongs to $F_{i\coh}$
(\hyperref[VI.cor1.17.1]{1.17.1}), and consequently
$X=\mu_i(X_i)$ belongs to $\underrightarrow{F}_{\coh}$.

\begin{enonce*}{Corollary 8.3.14}\etiquette[8.3.14]{VI.cor8.3.14}
Let $I$ be an essentially small cofiltered category, and let $p:F\to I$ and
$q:G\to I$ be two $𝒰$-topoi fibered over $I$ such that, for every object
$i$ of $I$, the topoi $F_i$ and $G_i$ are coherent, and such that, for every
morphism $f:i\to j$, the morphisms $f.:F_i\to F_j$ and $G_i\to G_j$ are
coherent. In addition, let $m:F\to G$ be a Cartesian morphism of fibered
$𝒰$-topoi such that, for every $i\in\ob(I)$,
$m_i:F_i\to G_i$ is coherent. Then the morphism $\underleftarrow{m}$ deduced
from $m$ by passage to the projective limit is coherent.
\end{enonce*}

For every object $i$ of $I$, the morphism $m_i$ induces a functor
$m^{\ast}_{i,\coh}:G_{i,\coh}\to F_{i,\coh}$, whence a Cartesian functor
$m^{\ast}_{\coh}=G_{\coh}\to F_{\coh}$ that is a morphism from the fibered
site $F_{\coh}$ to the fibered site $G_{\coh}$ \hyperref[VI.thm7.4.13.1]{(7.4.13)}. The functor
$\varinjlim_{I^{\circ}}m^{\ast}_{\coh}:
\varinjlim_{I^{\circ}}G_{\coh}\to\varinjlim_{I^{\circ}}F_{\coh}$ is a
morphism of sites that, on passage to the corresponding topoi, gives the
morphism $\underleftarrow{m}$. The assertion then follows from
\hyperref[VI.thm8.3.13]{8.3.13} and \hyperref[VI.cor3.3]{3.3}.

\begin{enonce*}[remark]{Exercise 8.3.15}\etiquette[8.3.15]{VI.exe8.3.15}
With the notation of \hyperref[VI.thm8.3.13]{8.3.13}, show that if the $F_i$,
$i\in I$, are algebraic (\hyperref[VI.def2.3]{2.3}) and if the morphisms
$f.:F_i\to F_j$, $f\in\text{\rm Fl}(I)$, are coherent,\pageoriginale then
$\limtop{I}F_i$ is algebraic and
$(\limtop{I}F_i)_{\coh}\xrightarrow{\sim}\underrightarrow{F}_{\coh}$.
\end{enonce*}


% END INCLUDED SGA4_English_current_buildable_source_20260906/source/06/drafts/SGA4_VI_LINES_6043_6336_English.texfrag


% BEGIN INCLUDED SGA4_English_current_buildable_source_20260906/source/06/drafts/SGA4_VI_LINES_6337_6444_English.texfrag
\subsection{Example: local topoi}\etiquette[8.4]{VI.subsec8.4}

\subsubsection{}\etiquette[8.4.1]{VI.subsubsec8.4.1}

Let $X=(X_i)_{i\in I}$ be a topos fibered over a category $I$. One
deduces from it a topos fibered over the category $\Pro(I)$ (I \hyperref[I.sec8.10]{8.10})
by the formula
\begin{equation*}
X_{\sheaf{\alpha}}=\limtop{I_{\alpha}}X_{i\alpha}\tag{8.4.1.1}
\end{equation*}\etiquette[8.4.1.1]{VI.eq8.4.1.1}
for every pro-object $\sheaf{\alpha}=(i_{\alpha})$ of $I$ (to see
this, one may generalize Exercise I
\hyperref[VI.subsubsec8.2.8]{8.2.8} to pseudofunctors).

\subsubsection{}\etiquette[8.4.2]{VI.subsubsec8.4.2}

Let $E$ be a topos, and let $p:\text{Fl}(E)\to E$ be the fibered topos
considered in \hyperref[VI.exem7.3.1]{7.3.1}. Extending it as in
\hyperref[VI.subsubsec8.4.1]{8.4.1}, one obtains a fibered topos
$\bar{\text{Fl}(E)}\to\Pro(E)$. The category $\Point(E)$ maps by a
fully faithful functor into $\Pro(E)$ (IV \hyperref[IV.subsubsec6.8.5]{6.8.5}), whence, by base
change (\hyperref[VI.subsubsec7.1.9]{7.1.9}), a fibered topos
$\Loc(E)\to\Point(E)$. The fiber $\Loc_p(E)$ at a point $p$ of $E$
is called the \emph{localization of $E$ at the point} $p$. By the
preceding, this topos depends covariantly on the point $p$. By
definition,
\begin{equation*}
\Loc_p(E)=\Limtop_{X\in\Vois(p)}E_{/X}{\quoi},\tag{8.4.2.1}
\end{equation*}\etiquette[8.4.2.1]{VI.eq8.4.2.1}
where $\Vois(p)$ is the category of neighborhoods of $p$ (IV \hyperref[IV.subsec6.8]{6.8}).

\subsubsection{}\etiquette[8.4.3]{VI.subsubsec8.4.3}

Let $m:p'\to p$ be a morphism of $\Point(E)$ (IV 6). For every
$X\in\Vois(p)$ one thereby obtains a point $p'_X$ of $E_{/X}$ and
hence, by the universal property of $\limtop{}$
\hyperref[VI.subsec8.1]{8.1}, a point $\theta_p(m)$ of $\Loc_p(E)$.
This defines a functor
\begin{equation*}
\theta_p:\Point(E)_{/p}\to\Point(\Loc_p(E)){\quoi},\tag{8.4.3.1}
\end{equation*}\etiquette[8.4.3.1]{VI.eq8.4.3.1}
which is immediately seen to be an equivalence of categories. Thus
$\Point(\Loc_p(E))$ is canonically equivalent to the category of
generalizations of $p$ (IV \hyperref[IV.subsubsec7.1.8]{7.1.8}). There is moreover a canonical
equivalence of topoi
$$
\can\bigdot:\Loc_{\theta_p(m)}(\Loc_p(E))\to\Loc_{p'}(E){\quoi},
$$
which fits into a diagram commutative up to isomorphism
\[
\xymatrix@R=1.5cm{
\Loc_{\theta_p(m)}(\Loc_p(E))\ar[rr]^{\can\bigdot}\ar[rd]_{\phi}
&&\Loc_{p'}(E)\ar[ld]^{m}\\
&\Loc_p(E)&
}
\]
where $\phi$\pageoriginale is the canonical morphism of
\eqhyperref[VI.eq8.4.2.1]{8.4.2.1}.

\subsubsection{}\etiquette[8.4.4]{VI.subsubsec8.4.4}

Localized topoi seem to be of interest only for topoi arising from
algebraic geometry, or at least for topoi having suitable finiteness
properties. Thus, when $E$ is the topos of sheaves on a separated
topological space, one sees immediately (using, for example,
\hyperref[VI.thm8.2.9]{8.2.9}) that the localized topoi are point
topoi. Moreover, in this case the category $\Point(E)$ is discrete
and the situation described in \hyperref[VI.subsubsec8.4.2]{8.4.2}
is trivial. On the other hand, let $E=\Top(X)$ be the topos of sheaves
for the Zariski topology on a scheme $X$, let $x$ be a point of $X$,
let $O_x$ be the local ring of $X$ at $x$, and let
$Y=\spec(O_x)$. One finds that
$$
\Loc_x(E)=\Top(Y){\quoi}.
$$
For other examples arising from the \'etale topology, we refer to
Expos\'e VIII of this seminar.

\subsubsection{}\etiquette[8.4.5]{VI.subsubsec8.4.5}

When $E$ is locally coherent (\hyperref[VI.def2.3]{2.3}), the
localized topoi $\Loc_p(E)$ are coherent (in
\eqhyperref[VI.eq8.4.2.1]{8.4.2.1} one may restrict to the
$X\in\Vois(p)$ which are algebraic and coherent, and then apply
\hyperref[VI.thm8.3.13]{8.3.13}). For every morphism of points
$m:p'\to p$, the morphism of topoi
$m\bigdot:\Loc_{p'}(E)\to\Loc_p(E)$ is coherent
(\hyperref[VI.cor8.3.14]{8.3.14}).

\subsubsection{}\etiquette[8.4.6]{VI.subsubsec8.4.6}

A \emph{local topos} is a topos $X$ such that the functor
$\Gamma(X,-)=$ ``sections over $X$'' (IV \hyperref[IV.subsec4.3]{4.3}) is a fiber functor
(IV 6). The point corresponding to this fiber functor is called the
\emph{center} of the local topos. When $E$ is locally coherent, its
localized topoi are local topoi (\hyperref[VI.thm8.2.3]{1.2.3},
\hyperref[VI.prop8.5.2]{8.5.2}, and \hyperref[VI.cor8.5.7]{8.5.7}).
The center of $\Loc_p(E)$ is canonically isomorphic to
$\theta_p(\id_p)$ (\eqhyperref[VI.eq8.4.3.1]{8.4.3.1}). The image in
$E$ of the center of $\Loc_p(E)$ is isomorphic to $p$. The localized
topos $\Loc_p(E)$, equipped with the canonical morphism\pageoriginale
$\Loc_p(E)\to E$, is the solution of the universal (2-universal!)
problem of mapping local topoi into $E$ in such a way that the center
maps ``to'' $p$.

\subsubsection{}\etiquette[8.4.7]{VI.subsubsec8.4.7}

A number of questions concerning local topoi have not been addressed
by the editors. Thus, if $X$ is a local topos, the final object $e$ of
$X$ has a maximal open subtopos $U\ne e$. The complement $Y$ of $U$
is a local topos having no nontrivial open subtopos. Is such a topos a
point topos? Let $U$ be a topos and let
$X=(U,\Phi:U\to(\Ens))$ be a topos obtained by gluing. What conditions
must the gluing functor $\Phi$ satisfy for $X$ to be a local topos?


% END INCLUDED SGA4_English_current_buildable_source_20260906/source/06/drafts/SGA4_VI_LINES_6337_6444_English.texfrag


% BEGIN INCLUDED SGA4_English_current_buildable_source_20260906/source/06/drafts/SGA4_VI_LINES_6445_6610_English.texfrag
\subsection{Structure of the sheaves of a filtered projective limit
  of topoi}\etiquette[8.5]{VI.subsec8.5}

\subsubsection{}\etiquette[8.5.1]{VI.subsubsec8.5.1}

In this section, $F=(F_i)_{i\in I}$ is a $\Uu$-topos fibered over a
small cofiltered category $I$. Choose a bicleavage
(\hyperref[VI.subsubsec7.1.4]{7.1.4}) of $F$, \ie for every
$f\in\text{Fl}(I):i\to j$, choose morphisms of topoi
$f\bigdot:F_i\to F_j$ such that $f^{\ast}:F_j\to F_i$ is the
inverse-image functor for the fibered structure. There are then
canonical isomorphisms $c_{f,g}$ satisfying a cocycle condition
(\hyperref[VI.subsubsec7.1.3]{7.1.3}). Denote by $\underleftarrow{F}$
the projective limit of the fibered topos $F$
(\hyperref[VI.thm8.2.3]{8.2.3}), and for every $i\in\ob I$, denote by
\begin{equation*}
\mu_i:\underleftarrow{F}\to F_i{\quoi},\tag{8.5.1.1}
\end{equation*}\etiquette[8.5.1.1]{VI.eq8.5.1.1}
the canonical morphism (\hyperref[VI.subsubsec8.1.3]{8.1.3}). Denote
by $\Top(F)$ the total topos of $F$ \hyperlink{sga4.exp.VI.remark.7.4.3.item.3}{(7.4.3.3)}. By
\hyperref[VI.thm8.2.9]{8.2.9}, the canonical functor
$F\to\underrightarrow{F}$ defines a morphism of topoi
\begin{equation*}
Q:\underleftarrow{F}\to\Top(F){\quoi}.\tag{8.5.1.2}
\end{equation*}\etiquette[8.5.1.2]{VI.eq8.5.1.2}
The topos $\underleftarrow{F}$ is identified with
$\Hhom_{\Cart/I^{\circ}}(I^{\circ},F')$
(\hyperref[VI.thm8.2.9]{8.2.9}), and the topos $\Top(F)$ is identified
with $\Hhom_{I^{\circ}}(I^{\circ},F')$
(\hyperref[VI.prop7.4.7]{7.4.7}). Under these identifications, $Q$ is
the embedding morphism of
$\Hhom_{\cart/I^{\circ}}(I^{\circ},F')$ into
$\Hhom_{I^{\circ}}(I^{\circ},F')$
(\hyperref[VI.thm8.2.9]{8.2.9}), and for every object $M$ of
$\underleftarrow{F}$ we have
\begin{equation*}
Q_{\ast}(M)=(i\longmapsto\mu_{i\ast}(M)){\quoi}.\tag{8.5.1.3}
\end{equation*}\etiquette[8.5.1.3]{VI.eq8.5.1.3}
By \hyperlink{sga4.exp.VI.remark.7.4.3.item.4}{7.4.3.4}, for every $i\in\ob I$ there is a morphism of topoi
\begin{equation*}
\alpha_i:F_i\to\Top(F){\quoi}.\tag{8.5.1.4}
\end{equation*}\etiquette[8.5.1.4]{VI.eq8.5.1.4}
The diagram\pageoriginale
\[\vcenter{
\xymatrix@R=1.25cm{
{\underleftarrow{F}}\ar[rr]^{\mu_i}\ar[rd]_{Q}&&F_i\ar[ld]^{\alpha_i}\\
&\Top(F)&
}}\leqno{(8.5.1.5)}
\]\etiquette[8.5.1.5]{VI.eq8.5.1.5}
is not commutative in general (even up to isomorphism). But the
definition of the morphisms involved gives a canonical isomorphism
\begin{equation*}
\alpha_i^{\ast}Q_{\ast}\simeq\mu_{i\ast}.\tag{8.5.1.6}{\quoi}
\end{equation*}\etiquette[8.5.1.6]{VI.eq8.5.1.6}
When $i$ is a \emph{final object} of $I$, $\alpha_i$ is an embedding
having a retraction $\beta_i:\Top(F)\to F_i$
(\hyperref[VI.lem7.4.12]{7.4.12}), and the diagram
\[\vcenter{
\xymatrix@R=1.25cm{
\underleftarrow{F}\ar[rr]^{\mu_i}\ar[rd]_{Q}&&F_i\ar[ld]^{\beta_i}\\
&\Top(F)&,
}}\leqno{(8.5.1.7)}
\]\etiquette[8.5.1.7]{VI.eq8.5.1.7}
is commutative up to canonical isomorphism.

\begin{enonce*}{Proposition 8.5.2}\etiquette[8.5.2]{VI.prop8.5.2}
Let $j\mapsto M_j$ be an object of $\Top(F)$. There is a functorial
isomorphism
\begin{equation*}
Q^{\ast}(j\mapsto M_j)\simeq\mathop{\varinjlim}_{I^{\circ}}
\mu_j^{\ast}M_j{\quoi}.\tag{8.5.2.1}
\end{equation*}\etiquette[8.5.2.1]{VI.eq8.5.2.1}
\end{enonce*}
\noindent
An object of $\Top(F)$ (\hyperref[VI.subsubsec7.1.3]{7.1.3}) consists
of the data of an assignment $j\mapsto M_j$, $M_j\in\ob F_j$, and,
for every morphism $f:i\to j$, a morphism
\begin{equation*}
\beta_f:M_j\to f_{\ast}M_i\tag{8.5.2.2}
\end{equation*}\etiquette[8.5.2.2]{VI.eq8.5.2.2}
or, equivalently by adjunction, a morphism
\begin{equation*}
\beta'_f:f^{\ast}M_j\to M_i{\quoi},\tag{8.5.2.3}
\end{equation*}\etiquette[8.5.2.3]{VI.eq8.5.2.3}
where the morphisms $\beta'_f$ are subject to the condition that, for
every pair of composable morphisms
$i\xrightarrow{f}j\xrightarrow{g}k$, the following diagram be
commutative:
\[\vcenter{
\xymatrix@C=2.25cm@R=1.5cm{
f^{\ast}g^{\ast}M_k\ar[r]^{f^{\ast}\beta'_g}\ar@{-}[d]_{c^{\ast}_{g,f}}&
f^{\ast}M_j\ar@{->}[d]^{\beta'_f}\\
(gf)^{\ast}M_k\ar[r]^{\beta'_{gf}}&M_i
}}\leqno{(8.5.2.4)}
\]\etiquette[8.5.2.4]{VI.eq8.5.2.4}
Recall\pageoriginale moreover from
\eqhyperref[VI.eq8.1.3.1]{8.1.3.1} that, for every morphism
$f:i\to j$, there is an isomorphism
$b_f:f\bigdot\mu_i\xrightarrow{\sim}\mu_j$, and that for every pair
of composable morphisms $i\xrightarrow{f}j\xrightarrow{g}k$ there is
a commutative diagram
\[
\xymatrix@C=2cm@R=1.25cm{
g\bigdot f\bigdot\mu_i\ar[r]^{c_{g,f}}\ar[d]^{g\bigdot(b_f)}&
(gf)\bigdot\mu_i\ar[d]^{b_{gf}}\\
g\bigdot\mu_j\ar[r]^{b_g}&\mu_k
}\leqno{(8.5.2.5)}
\]\etiquette[8.5.2.5]{VI.eq8.5.2.5}
Let $f:i\to j$ be a morphism of $I$. Denote by
\begin{equation*}
t_f:\mu_j^{\ast}(M_j)\to\mu_i^{\ast}(M_i)\tag{8.5.2.6}
\end{equation*}\etiquette[8.5.2.6]{VI.eq8.5.2.6}
the composite
$\mu_j^{\ast}(M_j)\xrightarrow{b_f^{\ast}}
\mu_i^{\ast}(f^{\ast}(M_j))
\xrightarrow{\mu_i^{\ast}(\beta'_f)}\mu_i^{\ast}(M_i)$.
We leave it to the reader to verify that the commutativity of diagrams
\eqhyperref[VI.eq8.5.2.4]{8.5.2.4} and
\eqhyperref[VI.eq8.5.2.5]{8.5.2.5} implies that for every pair of
composable morphisms $i\to j\to k$, one has
\begin{equation*}
t_ft_g=t_{gf}{\quoi}.\tag{8.5.2.7}
\end{equation*}\etiquette[8.5.2.7]{VI.eq8.5.2.7}
We have therefore defined a functor
$I^{\circ}\to\underleftarrow{F}$, whose inductive limit
$\varinjlim_{I^{\circ}}\mu_i^{\ast}(M_i)$ can be considered. Let $N$
be an object of $\underleftarrow{F}$. We have
\begin{equation*}
\Hom(\varinjlim_{I^{\circ}}\mu_i^{\ast}(M_i),N)\simeq
\varprojlim_{I^{\circ}}\Hom(\mu_i^{\ast}(M_i),N){\quoi},\tag{8.5.2.8}
\end{equation*}\etiquette[8.5.2.8]{VI.eq8.5.2.8}
and hence, by adjunction, an isomorphism
\begin{equation*}
\Hom(\varinjlim_{I^{\circ}}\mu_i^{\ast}(M_i),N)\simeq
\varprojlim_{I^{\circ}}\Hom_{F_i}(M_i,\mu_{i\ast}(N)){\quoi}.
\tag{8.5.2.9}
\end{equation*}\etiquette[8.5.2.9]{VI.eq8.5.2.9}
We\pageoriginale now interpret the right-hand side of
\eqhyperref[VI.eq8.5.2.9]{8.5.2.9}. For every morphism $f:i\to j$,
denote by
$$
d_f:\Hom_{F_i}(M_i,\mu_{i\ast}(N))\to
\Hom_{F_j}(M_j,\mu_{j\ast}(N)){\quoi}
$$
the transition map of the projective system occurring in
\eqhyperref[VI.eq8.5.2.9]{8.5.2.9}. It follows from the definition of
the morphisms $t_f$ (\eqhyperref[VI.eq8.5.2.6]{8.5.2.6}) that $d_f$
associates with a morphism
$$
u_i:M_i\to\mu_{i\ast}(N)
$$
the morphism
$$
d_f(u_i):M_j\to\mu_{j\ast}(N){\quoi},
$$
obtained by composing the morphisms
$$
M_j\xrightarrow{\beta_f}f_{\ast}(M_i)
\xrightarrow{f_{\ast}(u_i)}f_{\ast}\mu_{i\ast}(N)
\xrightarrow{b_f}\mu_{j\ast}(N){\quoi}.
$$
Thus an element of the right-hand side of
\eqhyperref[VI.eq8.5.2.9]{8.5.2.9} is a family of morphisms
$u_i:M_i\to\mu_{i\ast}(N)$, $i\in\ob I$, such that for every
$f:i\to j$ the diagram
\[
\xymatrix@C=2cm@R=1.25cm{
M_j\ar[r]\ar[d]_{u_j}&f_{\ast}(M_i)\ar[d]^{f_{\ast}(u_i)}\\
\mu_{j\ast}(N)\ar[r]^{\sim}&f_{\ast}\mu_{i\ast}(N)
}
\]
is commutative, that is, a morphism from the section
$(i\mapsto M_i)\in\Hhom_{I^{\circ}}(I^{\circ},F')$ to the section
$Q_{\ast}(N)$. We therefore have a functorial isomorphism
$$
\Hom(\varinjlim_{I^{\circ}}\mu_i^{\ast}(M_i),N)
\simeq\Hom((i\mapsto M_i),Q_{\ast}(N)){\quoi};
$$
the proposition follows by adjunction.


% END INCLUDED SGA4_English_current_buildable_source_20260906/source/06/drafts/SGA4_VI_LINES_6445_6610_English.texfrag


% BEGIN INCLUDED SGA4_English_current_buildable_source_20260906/source/06/drafts/SGA4_VI_LINES_6611_6839_English.texfrag
\begin{enonce*}{Proposition 8.5.3}\etiquette[8.5.3]{VI.prop8.5.3}
If, for every $f:i\to j\in\text{\rm Fl}(I)$, the functors
$f_{\ast}:F_i\to F_j$ commute with small filtered inductive limits, then for
every section $(i\mapsto M_i)\in\Hom_I^{\circ}(I^{\circ},F')$ we have
\begin{equation*}
Q_{\ast} Q^{\ast} (i \mapsto M_i) \simeq i \mapsto
\varinjlim_{f : j \to i} f_{\ast} M_j{\quoi}.\tag{8.5.3.1}
\end{equation*}\etiquette[8.5.3.1]{VI.eq8.5.3.1}
\end{enonce*}

\noindent We\pageoriginale leave it to the reader to describe explicitly,
in terms of the $\beta_f$ \eqhyperref[VI.eq8.5.2.2]{8.5.2.2} and the
$c_{f,g}$ (\hyperref[VI.subsubsec8.5.1]{8.5.1}), the transition morphisms of
the inductive systems $(f:j\to i)\mapsto f_{\ast}(M_j)$ and the transition
morphisms of the section
$i\mapsto\varinjlim_{f:j\to i}f_{\ast}M_j$. Since the functors $f_{\ast}$
commute with filtered inductive limits, the section
$i\mapsto\varinjlim_{f:j\to i}f_{\ast}M_j$ is Cartesian. Moreover, there is
a natural morphism $u$ from $(i\mapsto M_i)$ to the section
$(i\mapsto\varinjlim_{f:j\to i}f_{\ast}M_j)$, and it is clear that every
morphism from $(i\mapsto M_i)$ to a Cartesian section factors uniquely
through $u$, giving the isomorphism \eqhyperref[VI.eq8.5.3.1]{8.5.3.1}.

\begin{enonce*}{Corollary 8.5.4}\etiquette[8.5.4]{VI.cor8.5.4}
Under the hypotheses of \hyperref[VI.prop8.5.3]{8.5.3}, the functor
$Q_{\ast}$ commutes with small filtered inductive limits.
\end{enonce*}

Let $\alpha\mapsto N^{\alpha}$ be a small filtered inductive system of
$\underleftarrow{F}$. Put $M^{\alpha}=Q_{\ast}(N^{\alpha})$, so that
$M^{\alpha}$ is a section $i\mapsto M_i^{\alpha}$. We have
$N^{\alpha}\simeq Q^{\ast}Q_{\ast}(N^{\alpha})$, and since $Q^{\ast}$
commutes with inductive limits, we have
$\varinjlim_{\alpha}N^{\alpha}\simeq
Q^{\ast}(\varinjlim_{\alpha}M^{\alpha})$. Inductive limits in
$\Top(F)=\Hom_{I^{\circ}}(I^{\circ},F')$ are computed fiber by fiber. Thus
$\varinjlim_{\alpha}M^{\alpha}\simeq
(i\mapsto\varinjlim_{\alpha}M_i^{\alpha})$. By
\eqhyperref[VI.eq8.5.3.1]{8.5.3.1}, we have
$$
Q_{\ast} (\varinjlim_{\alpha} {N}^{\alpha}) \simeq Q_{\ast}
Q^{\ast} (\varinjlim_{\alpha} {M}^{\alpha}) \simeq i \mapsto
\varinjlim_{f : j \to i} f_{\ast}
(\varinjlim_{\alpha} {M}^{\alpha}_j){\quoi}.
$$
Since the functors $f_{\ast}$ commute with filtered inductive limits, it
follows that
$$
Q_{\ast} (\varinjlim_{\alpha} {N}^{\alpha}) \simeq i \mapsto
\varinjlim_{f: j\to i} \varinjlim_{\alpha} f_{\ast}
(M^{\alpha}_j) \simeq i \mapsto \varinjlim_{\alpha} \varinjlim_{f:
  j\to i} f_{\ast}
(M^{\alpha}_j) \simeq \varinjlim_{\alpha} Q_{\ast} (N^{\alpha}){\quoi}.
$$

\begin{enonce*}{Corollary 8.5.5}\etiquette[8.5.5]{VI.cor8.5.5}
If the functors $f_{\ast}:F_i\to F_j$ commute with small filtered inductive
limits, then for every object $i$ of $I$ and every $M_i\in\ob F_i$ we have
\begin{equation*}
\mu_{i\ast} \mu_{i}^{\ast} (M_i) \simeq \varinjlim_{f : j \to
      i} f_{\ast} f^{\ast}(M_i){\quoi}.\tag{8.5.5.1}
\end{equation*}\etiquette[8.5.5.1]{VI.eq8.5.5.1}
\end{enonce*}
\noindent
After making the base change $I/i\to I$, we may suppose that $i$ is a final
object of $I$ (\hyperref[VI.lem8.2.1]{8.2.1}). The morphism
$\mu_i:\underleftarrow{F}\to F_i$ is then the composite of
the\pageoriginale morphisms (\hyperref[VI.thm8.2.9]{8.2.9}) and
(\hyperref[VI.lem7.4.12]{7.4.12}):
\begin{align*}
Q &: \underleftarrow{F} \to F^{\sim}{\quoi},\\
\beta_i &: (\alpha_{i!}, \alpha^{\ast}_i) : F^{\sim} \to F_i{\quoi}.
\end{align*}
Thus $\mu_i^{\ast}(M_i)\simeq Q^{\ast}\beta_i^{\ast}(M_i)$, where
$\beta_i^{\ast}(M_i)$ is the section
$(f:i\to j)\mapsto f^{\ast}(M_i)$ whose transition morphisms are deduced
from the $c_{f,g}$. The assertion then follows from
\hyperref[VI.prop8.5.3]{8.5.3}.

\begin{enonce*}{Corollary 8.5.6}\etiquette[8.5.6]{VI.cor8.5.6}
Under the hypotheses of \hyperref[VI.cor8.5.5]{8.5.5}, the functors
$$
\mu_{i\ast} : \underleftarrow{F} \to F_i
$$
commute with filtered inductive limits.
\end{enonce*}

The functor $\mu_{i\ast}$ is the composite of the functor $Q_{\ast}$, which
commutes with filtered inductive limits, and the functor ``restriction to
the fiber at $i$'', which commutes with inductive limits.

\begin{enonce*}{Corollary 8.5.7}\etiquette[8.5.7]{VI.cor8.5.7}
Under the hypotheses of \hyperref[VI.prop8.5.3]{8.5.3}, let $i$ be an object
of $I$ and $X$ an object of $F_i$ such that the functor $\Hom(X,-)$ on $F_i$
commutes with filtered inductive limits. Then the functor
$\Hom(\mu_i^{\ast}(X),-)$ commutes with filtered inductive limits, and for
every object $M_i$ of $F_i$ we have
\begin{equation*}
\Hom(\mu^{\ast}_i(X), \mu^{\ast}_i(M_i)) \simeq \varinjlim_{f
: j \to i} \Hom(f^{\ast}(X), f^{\ast}(M_i))\tag{8.5.7.1}{\quoi}.
\end{equation*}\etiquette[8.5.7.1]{VI.eq8.5.7.1}
\end{enonce*}
\noindent
For every object $(j\mapsto M_j)$ of $\Top(F)$, we have
\begin{equation*}
\Hom(\mu^{\ast}_i(X), Q^{\ast}(j \mapsto M_j)) \simeq \varinjlim_{f
: j \to i} \Hom(f^{\ast}(X), M_{j}){\quoi}.\tag{8.5.7.2}
\end{equation*}\etiquette[8.5.7.2]{VI.eq8.5.7.2}

The functor $\Hom(\mu_i^{\ast}(X),-)$ is isomorphic to the functor
$\Hom(X,\mu_{i\ast}(-))$, and the latter commutes with filtered inductive
limits (\hyperref[VI.cor8.5.4]{8.5.4}). By
\eqhyperref[VI.eq8.5.5.1]{8.5.5.1}, we have
$$
\Hom(\mu^{\ast}_i(X), \mu^{\ast}_i (M_i)) \simeq \Hom (X, \mu_{i\ast}\mu^{\ast}_i (M_i)) \simeq  \Hom(X, \varinjlim_{f : j \to
  i} f_{\ast} f^{\ast} (M_i)){\quoi},
$$
which gives formula \eqhyperref[VI.eq8.5.7.1]{8.5.7.1}. Likewise, by
\eqhyperref[VI.eq8.5.3.1]{8.5.3.1}, we have
$$
\Hom(\mu^{\ast}_i (X), Q^{\ast}(j \mapsto M_j)) \simeq \Hom(X,
\mu_{i\ast} Q^{\ast} (j \mapsto M_j)) \simeq \Hom(X,
\varinjlim_{f : j \to i} f_{\ast} (M_j)){\quoi},
$$
which\pageoriginale gives formula \eqhyperref[VI.eq8.5.7.2]{8.5.7.2}.

\begin{enonce*}{Proposition 8.5.8}\etiquette[8.5.8]{VI.prop8.5.8}
Let $G\to I$ be a fibered $𝒰$-topos and $m:F\to G$ a Cartesian morphism of
fibered topoi \hyperref[VI.def7.1.5]{(7.1.5)}. For $G$ we use the notation introduced in
\hyperref[VI.def8.1.1]{8.1.1}. Let
$\underleftarrow{m}:\underleftarrow{F}\to\underleftarrow{G}$ be the
morphism deduced from $m$ by passage to the projective limit
(\hyperref[VI.subsubsec8.1.4]{8.1.4}). The diagram of topoi and morphisms of
topoi
\[\vcenter{
\xymatrix@C=2.25cm@R=1.25cm{
\underset{\leftarrow}{F}  \ar[r]^{Q} \ar[d]^{\underset{\leftarrow}{m}}& \Top(F)\ar[d]^{m^{\sim}}\\
\underset{\leftarrow}{G}  \ar[r]^{Q} & \Top(G)
}}\leqno{(8.5.8.1)}
\]\etiquette[8.5.8.1]{VI.eq8.5.8.1}
is commutative up to isomorphism. For every object $N$ of
$\underleftarrow{F}$, we have
\begin{equation*}
\underleftarrow{m}_{\ast} (N) \simeq \varinjlim_{I^{\circ}} \mu_j^{\ast}
m_{j\ast} \mu_{j\ast} (N){\quoi}.\tag{8.5.8.2}
\end{equation*}\etiquette[8.5.8.2]{VI.eq8.5.8.2}
\end{enonce*}

The commutativity of diagram \eqhyperref[VI.eq8.5.8.1]{8.5.8.1} follows
immediately from the definitions (\hyperref[VI.subsubsec8.2.8]{8.2.8} and
\hyperref[VI.subsubsec8.5.1]{8.5.1}). We have
$\underleftarrow{m}_{\ast}(N)\simeq
Q^{\ast}Q_{\ast}\underleftarrow{m}_{\ast}(N)$ and, by the commutativity of
\eqhyperref[VI.eq8.5.8.1]{8.5.8.1},
$\underleftarrow{m}_{\ast}(N)\simeq
Q^{\ast}m^{\sim}_{\ast}Q_{\ast}(N)$. Consequently, there is an isomorphism
$\underleftarrow{m}_{\ast}(N)\simeq
Q^{\ast}m^{\sim}_{\ast}(j\mapsto\mu_{j\ast}(N))$. The functor
$m_{\ast}^{\sim}$ is simply the functor
$\Hom_{I^{\circ}}(I^{\circ},m_{\ast})$
(\hyperref[VI.prop7.4.10]{7.4.10}), and consequently there is a canonical
isomorphism
$\underleftarrow{m}_{\ast}(N)\simeq
Q^{\ast}(j\mapsto m_{j\ast}\mu_{j\ast}(N))$. Formula
\eqhyperref[VI.eq8.5.8.2]{8.5.8.2} then follows from
\hyperref[VI.prop8.5.2]{8.5.2}.

\begin{enonce*}{Proposition 8.5.9}\etiquette[8.5.9]{VI.prop8.5.9}
Under the hypotheses and with the notation of \hyperref[VI.prop8.5.8]{8.5.8},
suppose in addition that, for every object $i$ of $I$, the functor
$m_{i\ast}:F_i\to G_i$ commutes with filtered inductive limits, and that for
every morphism $f:i\to j$ of $I$, the functor $f_{\ast}:F_i\to F_j$ commutes
with filtered inductive limits. Then the functor
$\underleftarrow{m}_{\ast}:\underleftarrow{F}\to\underleftarrow{G}$ commutes
with filtered inductive limits, and for every object $(i\mapsto M_i)$ of
$\Top(F)$ there is a canonical isomorphism
\begin{equation*}
\underleftarrow{m}_{\ast} Q^{\ast} (i \mapsto M_i) \simeq
\varinjlim_{I^{\circ}} \mu^{\ast}_j m_{j\ast} (M_j){\quoi}.\tag{8.5.9.1}
\end{equation*}\etiquette[8.5.9.1]{VI.eq8.5.9.1}
\end{enonce*}

The first\pageoriginale assertion follows from
\eqhyperref[VI.eq8.5.8.2]{8.5.8.2} and
\hyperref[VI.cor8.5.6]{8.5.6}. By \eqhyperref[VI.eq8.5.8.2]{8.5.8.2} and
\eqhyperref[VI.eq8.5.3.1]{8.5.3.1}, there is a functorial isomorphism
$$
\underleftarrow{m}_{\ast} Q^{\ast} (i \mapsto M_i) \simeq
\varinjlim_i \mu^{\ast}_{i} m_{i\ast} (\varinjlim_{f
  : j \to i} f_{\ast} (M_{j})){\quoi}.
$$
Using the commutation of the $m_{i\ast}$ with filtered inductive limits and
the isomorphisms $m_{i\ast}f_{\ast}\simeq f_{\ast}m_{j\ast}$
(\hyperref[VI.subsubsec7.1.6]{7.1.6}), we obtain
$$
\underleftarrow{m}_{\ast} Q^{\ast} (i \mapsto M_i) \simeq
\varinjlim_i \mu^{\ast}_{i} (\varinjlim_{f : j \to
  i} f_{\ast} m_{j\ast} (M_j)){\quoi}.
$$
Since the functors $\mu_i^{\ast}$ commute with inductive limits, it follows
that
$$
\underleftarrow{m}_{\ast}Q^{\ast}(i\mapsto M_{i})\simeq \varinjlim_{i}
\varinjlim_{f:j\to i} \mu^{\ast}_{i}f_{\ast}m_{j\ast}(M_{j}){\quoi}.
$$
This last object can be interpreted as an inductive limit over the category
$\text{Fl}(I)^{\circ}$, where $\text{Fl}(I)$ is the category of morphisms of
$I$. Let $\phi:I\to\text{Fl}(I)$ be the functor that associates with every
object $i$ of $I$ the identity morphism of $i$. The functor
$\phi^{\circ}:I^{\circ}\to\text{Fl}(I^{\circ})$ is cofinal, and consequently
(I \hyperref[I.sec8.1]{8.1}) there is a canonical isomorphism
$$
\underleftarrow{m}_{\ast} Q^{\ast} (i \mapsto M_i) \simeq
\varinjlim_{I^{\circ}} \mu^{\ast}_{j} m_{j\ast}(M_j){\quoi}.
$$

\begin{enonce*}{Corollary 8.5.10}\etiquette[8.5.10]{VI.cor8.5.10}
Under the conditions of \hyperref[VI.prop8.5.9]{8.5.9}, for every object $i$
of $I$ and every object $M_i$ of $F_i$, there is a canonical isomorphism
\begin{equation*}
\underleftarrow{m}_{\ast} \mu^{\ast}_i (M_i) \simeq
\varinjlim_{f:j\to i} \mu^{\ast}_{j} m_{j\ast}f^{\ast}
(M_{i}){\quoi}.\tag{8.5.10.1}
\end{equation*}\etiquette[8.5.10.1]{VI.eq8.5.10.1}
\end{enonce*}

The proof is analogous to that of Corollary \hyperref[VI.cor8.5.5]{8.5.5}.

\begin{enonce*}[remark]{Remark 8.5.11}\etiquette[8.5.11]{VI.rem8.5.11}
Recall that the direct-image functor under a coherent morphism between
coherent topoi commutes with filtered inductive limits
(\hyperref[VI.thm5.1]{5.1}). Consequently, Propositions
\hyperref[VI.prop8.5.3]{8.5.3} through \hyperref[VI.cor8.5.7]{8.5.7}, and
\hyperref[VI.prop8.5.9]{8.5.9}, \hyperref[VI.cor8.5.10]{8.5.10}, apply when
the fibered topoi under consideration are coherent and the morphisms of
fibered topoi under consideration are coherent.
\end{enonce*}


% END INCLUDED SGA4_English_current_buildable_source_20260906/source/06/drafts/SGA4_VI_LINES_6611_6839_English.texfrag


% BEGIN INCLUDED SGA4_English_current_buildable_source_20260906/source/06/drafts/SGA4_VI_LINES_6840_6915_English.texfrag
\subsection{Ringed fibered $𝒰$-topoi}\etiquette[8.6]{VI.subsec8.6}

\subsubsection{}\etiquette[8.6.1]{VI.subsubsec8.6.1}

A fibered $𝒰$-topos is called \emph{ringed} if it is equipped with a
sheaf of rings on the total site (\hyperref[VI.def7.4.1]{7.4.1}). Let
$p : F \to I$ be a fibered $𝒰$-topos. Choose a bisplitting
\pageoriginale (\hyperref[VI.subsubsec7.1.4]{7.1.4}). It follows from
\hyperref[VI.prop7.4.7]{7.4.7} that giving a sheaf of rings on $F$ is
equivalent to giving, for every object $i$ of $I$, a ring $A_i$ of
$F_i$, and, for every morphism $f : i \to j$, a morphism
$\phi_f : A_j \to f_{\ast} A_i$, the family of the $\phi_f$ being
subject to the compatibility conditions made explicit in \loccit.
The morphisms of topoi $f. : F_i \to F_j$ are therefore morphisms of
topoi ringed respectively by $A_i$ and $A_j$ (IV 11), and the ringed
structure on $F$, together with the choice of the morphisms $f$,
gives a pseudofunctor from $I$ to the 2-category of ringed topoi.
Conversely, given such a pseudofunctor, one can reconstruct the ringed
fibered $𝒰$-topos from which it arises.

\subsubsection{}\etiquette[8.6.2]{VI.subsubsec8.6.2}

As in \hyperref[VI.subsec8.1]{8.1}, one can define the notion of the
projective limit of a ringed $𝒰$-topos $(F_i, A_i, i \in I)$. We
shall suppose that this immediate generalization has been made, and
we shall freely apply to this situation the language introduced in
\hyperref[VI.subsec8.1]{8.1}. When it exists, this projective limit is
a ringed $𝒰$-topos $(H, B)$, and, for every object $i$ of $I$, we
have morphisms of ringed topoi
$\mu_i : (H, B) \to (F_i, A_i)$. If the underlying (unringed) fibered
topos $F$ admits a projective limit $\limtop{I} F_i$, and if the
inductive system of rings $i \mapsto \mu_i^{\ast}(A_i)$ admits an
inductive limit $\varinjlim_{I^{\circ}}\mu_i^{\ast} A_i$ in the
category of rings of $\limtop{I} F_i$ (as is always the case if $I$
is small), then the ringed topos
($\limtop{I} F_i$, $\varinjlim_{I^{\circ}} \mu^{\ast}_i A_i$) is
evidently a projective limit of the ringed fibered topos under
consideration.

\subsubsection{}\etiquette[8.6.3]{VI.subsubsec8.6.3}

The formulas of \hyperref[VI.subsec8.5]{8.5}, established for sheaves
of sets, are valid for abelian sheaves. They are also valid for
sheaves of modules, provided the inverse-image functors occurring in
them are interpreted as inverse-image functors in the sense of
sheaves of modules (IV \hyperref[IV.def1.1]{1.1}). We leave the verification to the reader.

\subsubsection{}\etiquette[8.6.4]{VI.subsubsec8.6.4}

Let $(p : F \to I, A)$ be a ringed fibered $𝒰$-topos over a
cofiltered category. We call $(p : F \to I, A)$ a ringed fibered
$𝒰$-topos flat on the right (\resp on the left) if, for every morphism
$f : i \to j$, the
morphism $f. : (F_i, A_i) \to (F_j, A_j)$ is a morphism of ringed
topoi flat on the right (\resp on the left)
(\hyperrefext[V][V.def1.8]{1.8}). Suppose that $I$ is essentially
small. Then the morphism of ringed topoi
$Q : (\limtop{I} F_i,
\varinjlim_{I^{\circ}} \mu^{\ast}_i (A_i)) \to
(\Hom_{I^{\circ}}(I^{\circ}, F'), A)$\pageoriginale
(\hyperref[VI.subsubsec8.5.1]{8.5.1}) is flat on the right (\resp on
the left), as follows from the definition and from
\hyperref[VI.prop8.5.2]{8.5.2}; this holds without any flatness
hypothesis on the transition morphisms $f.$. If, in addition,
$(p : F \to I, A)$ is a ringed fibered $𝒰$-topos flat on the right
(\resp on the left), the morphisms
$\mu_i : (\limtop{I} F_i,
\varinjlim_{I^{\circ}} \mu^{\ast}_i A_i) \to (F_i, A_i)$ are flat on
the right (\resp on the left). This follows from the fact that the
inverse image of a flat module is flat
(\hyperrefext[V][V.cor1.7.1]{1.7.1}), and that a filtered inductive
limit of flat modules is flat. The same arguments prove the following
fact: if $m : (F_i, A_i, i \in I) \to (G_i, A_i, i \in I)$ is a
cartesian morphism of ringed fibered $𝒰$-topoi over an essentially
small cofiltered category, and if, for every object $i \in \ob (I)$,
the morphism $m_i$ is flat on the right (\resp on the left), then the
morphism $\underleftarrow{m}$ deduced from $m$ by passage to the
projective limit is flat on the right (\resp on the left).

% END INCLUDED SGA4_English_current_buildable_source_20260906/source/06/drafts/SGA4_VI_LINES_6840_6915_English.texfrag


% BEGIN INCLUDED SGA4_English_current_buildable_source_20260906/source/06/drafts/SGA4_VI_LINES_6916_7181_English.texfrag
\subsection{Cohomology of sheaves on a projective limit of
  topoi}\etiquette[8.7]{VI.subsec8.7}

\subsubsection{}\etiquette[8.7.1]{VI.subsubsec8.7.1}

In this section, $I$ denotes a small cofiltered category,
$(p:F\to I,A)$ and $(q:G\to I,B)$ two ringed fibered $𝒰$-topoi, and
$m:(p:F\to I,A)\to(q:G\to I,B)$ a morphism of ringed fibered
$𝒰$-topoi. We suppose that, for every morphism $f:i\to j$ of $I$, the
derived functors $R^nf_{\ast}$, $n\in\NN$, of the functor
$f_{\ast}:\Mod(F_i,A_i)\to\Mod(F_j,A_j)$ on modules commute with
filtered inductive limits, and that, for every object $i\in I$, the
derived functors $R^nm_{i\ast}$ of the functor
$m_{i\ast}:(F_i,A_i)\to(G_i,B_i)$ on modules\footnote{To fix ideas,
we shall take left modules.} commute with filtered inductive limits.
Recall that these hypotheses hold when the categories $F_i,G_i$,
$(i\in\ob(I))$, are \emph{coherent} and when the morphisms
$f:F_i\to F_j$, $(f:i\to j\in \text{Fl}(I))$, and
$m_i:F_i\to G_i$, $(i\in I)$, are \emph{coherent}
(\hyperref[VI.cor5.2]{5.2}). We use the notation of
\hyperref[VI.subsubsec8.5.1]{8.5.1}. In addition, we denote by
$\underrightarrow{A}$ (\resp $\underrightarrow{B}$) the sheaf
$\varinjlim\mu^{\ast}_i(A_i)$ (\resp
$\varinjlim\mu^{\ast}_i(B_i)$). Finally, for modules, we use the
notation $\mu^{-1}_i$ for inverse image in the sense of abelian
sheaves, reserving $\mu_i^{\ast}$ for inverse image in the sense of
modules (IV 11).

\begin{enonce*}{Lemma 8.7.2}\etiquette[8.7.2]{VI.lem8.7.2}
Let\pageoriginale $j\mapsto M_j$ be an injective $A$-Module of
$\Top(F)$. Then $Q^{\ast}(j\mapsto M_j)$ is an $A$-Module acyclic for
$\underleftarrow{m}_{\ast}$, and, for every $i\in\ob(I)$, $M_i$ is
flasque.
\end{enonce*}

Let $i\in\ob(I)$, and let $e_i$ be a final object of $F_i$. The
fibered $𝒰$-topos $F_{/e_i}\to I_{/i}$ is deduced from $F\to I$ by
the base change $I_{/i}\to I$. Since the sheaf $j\mapsto M_j$ is
injective, it is flasque (\hyperrefext[V][V.subsec4.6]{4.6}). Its
restriction to the localized topos $F_{/e_i}$ is flasque
(\hyperrefext[V][V.subsec4.12]{4.12}). The restriction functor from
$F_{/e_i}$ to $F_i$ is a direct image functor for a morphism of topoi
(\hyperref[VI.lem7.4.12]{7.4.12}). It therefore carries flasque
Modules to flasque Modules (\hyperrefext[V][V.prop5.2]{5.2}). Thus
$M_i$ is flasque.

Let us prove the first assertion of the lemma. Set
$N'=Q^{\ast}(j\mapsto M_j)$. By
\eqhyperref[VI.eq8.5.3.1]{8.5.3.1}, for every $i\in\ob(I)$ we have
$$
\mu_{i\ast}(N') = \mu_{i\ast} Q^{\ast} (j \mapsto M_j) \simeq
\varinjlim_{f : j \to i} f_{\ast} (M_j){\quoi}.
$$
Using the hypothesis of \hyperref[VI.subsubsec8.7.1]{8.7.1}, we see
that the sheaf $N'$ has the following property:

(P)~For every object $i\in\ob(I)$, $\mu_{i\ast}(N')$ is
$m_{i\ast}$-acyclic, and, for every morphism $f:i\to j$ of $I$,
$\mu_{i\ast}(N')$ is $f_{\ast}$-acyclic.

It is enough to show that every object $N'$ having property (P) is
$\underleftarrow{m}_{\ast}$-acyclic. Since injectives have property
(P), it is enough to verify the following two properties
(\hyperrefext[V][V.prop0.4]{0.4}): for every exact sequence
$0\to N'\to N\to N''\to0$, where $N'$ and $N$ have property (P),
(a) $N''$ has property (P), and (b)
$m_{\ast}(N)\to m_{\ast}(N'')$ is an epimorphism. Verification of (a)
is immediate. Let us verify (b). Set
$K=\coker(Q_{\ast}(N')\to Q_{\ast}(N))$, so that $K$ is the section
$(i\mapsto K_i=\coker(\mu_{i\ast}(N')\to\mu_{i\ast}(N)))$. For every
$f:i\to j$, there is a commutative diagram
\[
\xymatrix{
0  \ar[r] & f_{\ast}\mu_{i\ast}(N') \ar[r] & f_{\ast}\mu_{i\ast}(N)\ar[r]
 & f_{\ast}(K_i) \ar[r]& 0\\
0  \ar[r] & \mu_{j\ast}(N')\ar[u] \ar[r] & \mu_{j\ast}(N) \ar[u] \ar[r] & K_j
\ar[r] \ar[u]&0.
}
\]
Since\pageoriginale $\mu_{i\ast}(N')$ is $f_{\ast}$-acyclic, the
top row is exact. Since $i\mapsto\mu_{i\ast}(N')$ and
$i\mapsto\mu_{i\ast}(N)$ are Cartesian sections, the first two
vertical morphisms are isomorphisms. Consequently, the canonical
morphism $K_j\to f_{\ast}(K_i)$ is an isomorphism, and
$i\mapsto K_i$ is a Cartesian section. Thus $K$ is of the form
$Q_{\ast}(K')$, and the canonical morphisms $Q_{\ast}(N)\to K$ and
$K\to Q_{\ast}(N'')$ are of the form $Q_{\ast}(u)$ and
$Q_{\ast}(v)$, respectively, because $Q_{\ast}$ is fully faithful.
Since $Q^{\ast}$ is exact and $Q^{\ast}Q_{\ast}$ is isomorphic to the
identity, $v$ is an isomorphism, and hence $Q_{\ast}(v)$ is an
isomorphism. Thus the sequence
$0\to Q_{\ast}(N')\to Q_{\ast}(N)\to Q_{\ast}(N'')\to0$ is exact,
and consequently, for every object $i$ of $I$, the sequence
$0\to\mu_{i\ast}(N')\to\mu_{i\ast}(N)\to
\mu_{i\ast}(N'')\to0$ is exact. Since $\mu_{i\ast}(N')$ is
$m_{i\ast}$-acyclic, the sequence
$0\to m_{i\ast}\mu_{i\ast}(N')\to
m_{i\ast}\mu_{i\ast}(N)\to
m_{i\ast}\mu_{i\ast}(N'')\to0$ is exact. Moreover, there is an
isomorphism
$\underleftarrow{m}_{\ast}\simeq
\varinjlim_{I^{\circ}}\mu^{-1}_jm_{j\ast}\mu_{j\ast}$
\eqhyperref[VI.eq8.5.8.2]{8.5.8.2}. The functor $\mu^{-1}_j$ is exact,
and filtered inductive limits are exact. Therefore the sequence
$0\to\underleftarrow{m}_{\ast}(N')\to
\underleftarrow{m}_{\ast}(N)\to
\underleftarrow{m}_{\ast}(N'')\to0$ is exact.

\begin{enonce*}{Theorem 8.7.3}\etiquette[8.7.3]{VI.thm8.7.3}
The notation and hypotheses are those of
\hyperref[VI.subsubsec8.7.1]{8.7.1}. For every integer $n\in\NN$,
the derived functor $R^n\underleftarrow{m}_{\ast}$ commutes with
filtered inductive limits, and for every $A$-Module
$j\mapsto M_j$ there is a functorial isomorphism
\begin{equation*}
R^{n} \underleftarrow{m}_{\ast} Q^{\ast} (j \mapsto M_j) \simeq
\varinjlim_{I^{\circ}}
\mu^{\ast}_j R^n m_{j\ast} (M_j)\tag{8.7.3.1}{\quoi}.
\end{equation*}\etiquette[8.7.3.1]{VI.eq8.7.3.1}
\end{enonce*}

We have $Q^{-1}(A)=\underrightarrow{A}$, and hence inverse image for
Modules is isomorphic to inverse image for abelian sheaves. It follows
from \hyperref[VI.eq8.5.8.2]{8.5.8.2} that, for every section
$j\mapsto M_j$, there is a canonical isomorphism
$\varinjlim_{I^{\circ}}\mu^{-1}_j(M_j)\simeq
\varinjlim_{I^{\circ}}\mu^{\ast}_j(M_j)$. Moreover, under the
hypotheses of \hyperref[VI.subsubsec8.7.1]{8.7.1}, there is a
canonical isomorphism \eqhyperref[VI.eq8.5.9.1]{8.5.9.1}
$\underleftarrow{m}_{\ast}Q^{\ast}(j\mapsto M_j)\simeq
\varinjlim_{I^{\circ}}\mu^{-1}_jm_{j\ast}(M_j)$. Since, for every
injective $j\mapsto M_j$, the $M_j$ are flasque and
$Q^{\ast}(j\mapsto M_j)$ is acyclic for
$\underleftarrow{m}_{\ast}$ (\hyperref[VI.lem8.7.2]{8.7.2}), there is
an isomorphism
$R^n\underleftarrow{m}_{\ast}(Q^{\ast}(j\mapsto M_j))\simeq
\varinjlim_{I^{\circ}}\mu^{-1}_jR^nm_{j\ast}(M_j)$, which gives
formula \eqhyperref[VI.eq8.7.3.1]{8.7.3.1} by what precedes. The first
assertion follows immediately from formula
\eqhyperref[VI.eq8.7.3.1]{8.7.3.1}.

\begin{enonce*}{Corollary 8.7.4}\etiquette[8.7.4]{VI.cor8.7.4}
For\pageoriginale $n\in\ZZ$, there are functorial isomorphisms
\begin{equation*}
R^n \underleftarrow{m}_{\ast} (N) \simeq \varinjlim_{I^{\circ}}
\mu^{\ast}_j R^n m_{j\ast} \mu_{j\ast} (N){\quoi}.\tag{8.7.4.1}
\end{equation*}\etiquette[8.7.4.1]{VI.eq8.7.4.1}
\end{enonce*}

Indeed, $N$ is isomorphic to
$Q^{\ast}(j\mapsto\mu_{j\ast}(N))$.

\begin{enonce*}{Corollary 8.7.5}\etiquette[8.7.5]{VI.cor8.7.5}
For every object $i$ of $I$, every integer $n$, and every
$A_i$-Module $M_i$, there is a functorial isomorphism
\begin{equation*}
R^n \underleftarrow{m}_{\ast} \mu^{\ast}_i (M_i) \simeq
\varinjlim_{f : j \to i} \mu^{\ast}_j R^n m_{j\ast} f^{\ast}
(M_i){\quoi}.\tag{8.7.5.1}
\end{equation*}\etiquette[8.7.5.1]{VI.eq8.7.5.1}
\end{enonce*}

After making the base change $I_{/i}\to I$, we may assume that $i$
is a final object of $I$ (\hyperref[VI.lem8.2.1]{8.2.1}). The object
$\mu_i^{\ast}(M_i)$ is then isomorphic to
$Q^{\ast}((f:j\to i)\mapsto f^{\ast}(M_i))$
(\cf \hyperref[VI.cor8.5.5]{8.5.5}).

\begin{enonce*}{Corollary 8.7.6}\etiquette[8.7.6]{VI.cor8.7.6}
Let $i$ be an object of $I$ and $n$ an integer. There is an
isomorphism, functorial in the section $j\mapsto M_j$:
\begin{equation*}
R^n \mu_{i\ast} Q^{\ast} (j \mapsto M_j) \simeq \varinjlim_{f :
  j \to i} R^n f_{\ast} (M_j){\quoi}.\tag{8.7.6.1}
\end{equation*}\etiquette[8.7.6.1]{VI.eq8.7.6.1}
There is an isomorphism functorial in the $\underrightarrow{A}$-Module
$N$:
\begin{equation*}
R^n \mu_{i\ast} (N) \simeq \varinjlim_{f:j\to i}
R^n f_{\ast} \mu_{j\ast} (N){\quoi}.\tag{8.7.6.2}
\end{equation*}\etiquette[8.7.6.2]{VI.eq8.7.6.2}
There is an isomorphism functorial in the $A_i$-Module $M_i$:
\begin{equation*}
R^n \mu_{i\ast} \mu^{\ast}_i (M_i) \simeq \varinjlim_{f:j\to i}
R^n f_{\ast} f^{\ast} (M_i){\quoi}.\tag{8.7.6.3}
\end{equation*}\etiquette[8.7.6.3]{VI.eq8.7.6.3}
The functors $R^n\mu_{i\ast}$ commute with filtered inductive limits.
\end{enonce*}

Making the base change $I_{/i}\to I$ reduces us to the case where
$i$ is a final object of $I$. Formulas
\eqhyperref[VI.eq8.7.6.1]{8.7.6.1} and
\eqhyperref[VI.eq8.7.6.2]{8.7.6.2} are then special cases of formulas
\eqhyperref[VI.eq8.7.3.1]{8.7.3.1},
\eqhyperref[VI.eq8.7.4.1]{8.7.4.1}, and
\eqhyperref[VI.eq8.7.5.1]{8.7.5.1}, respectively, obtained by taking
for $G$ the constant fibered topos with fiber $F_i$ and for $m$ the
morphism $(f:F_j\to F_i, f\in\ob(I_{/i}))$. The last assertion
follows from \eqhyperref[VI.eq8.7.6.2]{8.7.6.2}, taking
\hyperref[VI.cor8.5.4]{8.5.4} into account.

\begin{enonce*}{Corollary 8.7.7}\etiquette[8.7.7]{VI.cor8.7.7}
Let\pageoriginale $i$ be an object of $I$, and let $X$ be an object
of $F_i$ such that the functors $H^n(X,-)$, $n\in\NN$, commute with
filtered inductive limits. For every $n$, there is a canonical
isomorphism, functorial in the $A_i$-Module $M_i$:
\begin{equation*}
H^n (\mu_{i}^{\ast}(X), \mu^{\ast}_i (M_i)) \simeq \varinjlim_{f:j\to i}
H^n (f^{\ast} (X), f^{\ast} (M_i))\tag{8.7.7.1}
\end{equation*}\etiquette[8.7.7.1]{VI.eq8.7.7.1}
and the functors $H^n(\mu_i^{\ast}(X),-)$ commute with filtered
inductive limits. In particular, if the functors $H^n(F_i,-)$ commute
with filtered inductive limits, then the functors
$H^n(\underleftarrow{F},-)$ commute with filtered inductive limits,
and there are canonical isomorphisms, functorial in the $A_i$-Modules
$M_i$:
\begin{equation*}
H^n (\underleftarrow{F}, \mu_{i}^{\ast} (M_i)) \simeq \varinjlim_{f:j\to i}
H^n (F_j, f^{\ast} (M_i)){\quoi}.\tag{8.7.7.2}
\end{equation*}\etiquette[8.7.7.2]{VI.eq8.7.7.2}
\end{enonce*}

\noindent
It follows from \hyperrefext[V][V.prop5.3]{5.3} that there are two
spectral sequences
\begin{align*}
'E_{2}^{p, q} &= \varinjlim_{f:j\to i} H^p(X, R^q f_{\ast}
f^{\ast} (M_i)) \Rightarrow \varinjlim_{f:j\to i}
H^{p+q}(f^{\ast} (X), f^{\ast} (M_i))\\
''E_{2}^{p, q} &=  H^p(X, R^q \mu_{i\ast} \mu^{\ast}_i (M_i))
\Rightarrow H^{p+q} (\mu^{\ast}_i (X), \mu^{\ast}_i (M_i)){\quoi},
\end{align*}
and from \eqhyperref[VI.eq8.5.7.1]{8.5.7.1} that there is a morphism
between these two spectral sequences. Since the $H^{p+q}(X,-)$
commute with filtered inductive limits, it follows from
\eqhyperref[VI.eq8.7.6.2]{8.7.6.2} that this morphism of spectral
sequences is an isomorphism at the $E_2^{p,q}$ level. It therefore
induces an isomorphism on the abutments, giving
\eqhyperref[VI.eq8.7.7.1]{8.7.7.1}. For every
$\underrightarrow{A}$-Module $N$, there is a spectral sequence
(\hyperrefext[V][V.prop5.3]{5.3})
$$
E^{p, q}_2 = H^{p} (X, R^q \mu_{i\ast} (N)) \Rightarrow H^{p+q}
(\mu^{\ast}_i (X), N){\quoi}.
$$
The assertion that the functors $H^n(\mu_i^{\ast}(X),-)$ commute with
filtered inductive limits then follows from the analogous properties
of the functors $H^n(X,-)$ and $R^n\mu_{i\ast}$
(\hyperref[VI.cor8.7.6]{8.7.6}).

\begin{enonce*}[remark]{Remark 8.7.8}\etiquette[8.7.8]{VI.rem8.7.8}
When the topoi $F_j$, $j\in\ob(I)$, are coherent and the transition
morphisms are coherent, the hypothesis made on $X$ (\resp $F_i$) in
\hyperref[VI.cor8.7.7]{8.7.7} holds when $X$ (\resp $F_i$) is
coherent (\hyperref[VI.cor5.2]{5.2}). Moreover, we know that in this
case $\mu_i^{\ast}(X)$ (\resp $\underleftarrow{F}$) is coherent
(\hyperref[VI.thm8.3.13]{8.3.13}).
\end{enonce*}

\begin{enonce*}{Corollary 8.7.9}\etiquette[8.7.9]{VI.cor8.7.9}
Let\pageoriginale $i$ be an object of $I$, let $M_i$ be a left
$A_i$-Module, and let $L_i$ be a right $A_i$-Module admitting a
resolution of the form
$$
P_{i, k} \to P_{i, k-1} \to \cdots \to P_{i, 0} \to L_i \to 0
$$
where, for every integer $k$, $P_{i,k}$ is isomorphic to a finite
direct sum of objects of the form $A_{i|X}$ \emph{(IV \hyperref[IV.prop11.3.3]{11.3.3})}, with
$X$ satisfying the hypotheses of \hyperref[VI.cor8.7.7]{8.7.7}. If,
in addition to the hypotheses of \hyperref[VI.subsubsec8.7.1]{8.7.1},
the fibered topos $F$ is right flat (\hyperref[VI.subsec8.6]{8.6}),
then for $n\leq k-1$ there are canonical isomorphisms
\begin{equation*}
\Ext^n_{\underrightarrow{A}}(\underleftarrow{F}, \mu^{\ast}_i(L_i),
\mu^{\ast}_i (M_i)) \simeq \varinjlim_{f : j \to i}
\Ext^{n}_{A_j} (F_j, f^{\ast}(L_i), f^{\ast} (M_i)){\quoi}.\tag{8.7.9.1}
\end{equation*}\etiquette[8.7.9.1]{VI.eq8.7.9.1}
\end{enonce*}

\noindent
Denote by $P_{i,.}$ the resolution of $L_i$. Since $F$ is right flat,
for every $f:j\to i$, $f^{\ast}(P_{i,.})$ is a resolution of
$f^{\ast}(L_i)$, and $\mu_i^{\ast}(P_{i,.})$ is a resolution of
$\mu_i^{\ast}(L_i)$ (\hyperref[VI.subsec8.6]{8.6}). These resolutions
give two spectral sequences converging respectively to the two sides
of \eqhyperref[VI.eq8.7.9.1]{8.7.9.1}, together with a morphism between
the two spectral sequences. At the $E_1^{p,q}$ level this morphism is
an isomorphism by \eqhyperref[VI.eq8.7.7.1]{8.7.7.1}. It therefore
induces an isomorphism on the abutments.


% END INCLUDED SGA4_English_current_buildable_source_20260906/source/06/drafts/SGA4_VI_LINES_6916_7181_English.texfrag


% BEGIN INCLUDED SGA4_English_current_buildable_source_20260906/source/06/drafts/SGA4_VI_LINES_7182_7367_English.texfrag
\section{Appendix. Criterion for the existence of points}\etiquette[9]{VI.sec9}

\centerline{\textit{by} P. \textsc{Deligne}}

\medskip
\begin{enonce*}{Proposition 9.0}\etiquette[9.0]{VI.prop9.0}
Every locally coherent topos $S$ has enough points.
\end{enonce*}

Since the question is local on $S$, we may suppose that the topos $S$
is defined by a site $\Ss$ in which finite projective limits are
representable and such that every covering $f_i:U_i\to U$ admits a
finite subcovering. It is enough to prove that if
$f:\Ff\to\Gg$ is not a monomorphism, then there is a point $x$ of
$\Ss$ such that $f_x$ is not injective. By hypothesis, there are
$U\in\ob\Ss$ and $\ss,\ss'\in\Ff(U)$ such that $\ss\ne\ss'$ and
$f(\ss)=f(\ss')$. Replacing $\Ss$ by $\Ss/U$, we reduce to the
following lemma.

\begin{enonce*}{Lemma 9.1}\etiquette[9.1]{VI.lem9.1}
If\pageoriginale $\ss$ and $\ss'$ are two distinct global sections of
a sheaf $\mathfrak F$ on a site $\Ss$ satisfying the preceding
hypotheses, then there is a point $x$ of $\Ss$ such that
$\ss_x\ne\ss'_x$.
\end{enonce*}

Let $P=(U_i)_{i\in I}$ be a projective system in $\Ss$, indexed by a
filtered ordered set $I$. For every sheaf $\Hh$ on $\Ss$, put
$P(\Hh)=\varinjlim\Hh(U_i)$, and for every $V\in\Ob\Ss$, put
$P(V)=\varinjlim\Hom(U_i,V)$. The functors $P(\Hh)$ and $P(V)$
commute with fiber products. If the functor $P(V)$ takes coverings to
surjective families of maps, it is a morphism of sites from the point
topos to $\Ss$, and $P(\Hh)$ is the corresponding fiber functor.

If $P=(U_i)_{i\in I}$ and $Q=(V_j)_{j\in J}$ are two projective
systems in $\Ss$, we say that $Q$ refines $P$ if $I$ is a subset of
$J$ with the induced order and $P$ is the restriction of $Q$ to $I$.
There are then morphisms of functors from $P(\Hh)$ to $Q(\Hh)$ and
from $P(V)$ to $Q(V)$.

For every $P$ as above, denote by $\ss_P$ and $\ss'_P$ the images of
$\ss$ and $\ss'$ in $P(\Ff)$. Lemma \hyperref[VI.lem9.1]{9.1}
follows from the next lemma, in which $P$ is taken to be the
projective system indexed by $\{0\}$ and consisting of the final
object of $\Ss$.

\begin{enonce*}{Lemma 9.2}\etiquette[9.2]{VI.lem9.2}
For every $P$ as above satisfying $\ss_P\ne\ss'_P$, there is a
refinement $Q$ of $P$ satisfying $\ss_Q\ne\ss'_Q$ and such that, for
every covering $f_i:V_i\to V$ and every $f\in Q(V)$, $f$ belongs to
the image of one of the $Q(V_i)$.
\end{enonce*}

We first prove:

\begin{enonce*}{Lemma 9.3}\etiquette[9.3]{VI.lem9.3}
Let $P=(U_i)_{i\in I}$ satisfy the hypotheses of Lemma
\hyperref[VI.lem9.2]{9.2}, let $f_i:V_i\to V$ be a finite covering in
$\Ss$, and let $f\in P(V)$. There is a refinement $Q$ of $P$
satisfying $\ss_Q\ne\ss'_Q$ and such that the image of $f$ in $Q(V)$
belongs to the image of one of the $Q(V_i)$.
\end{enonce*}

There are $i_{\circ}\in I$ and
$f'\in\Hom(U_{i_{\circ}},V)$ which define $f$. For
$i\geq i_{\circ}$, put
$U_{i,k}=\fprod{U_i}{V_k}{V}$, this fiber product being defined by
$f'$. Since $(V_k)$ is a covering of $V$, $(U_{i,k})$ is a covering
of $U_i$; the following arrow is injective (for $i>i_{\circ}$):
$$
\Ff(U_i)\to\prod_k\Ff(U_{i,k}).
$$
Passing\pageoriginale to the limit, and using the commutation of
\emph{finite} products with filtered inductive limits, we see that
$$
P(\Ff)\to\prod_k\varinjlim\Ff(U_{i,k})
$$
is injective. Hence there is a $k$ such that $\ss_P$ and $\ss'_P$
have distinct images in
$\varinjlim_{i>i_{\circ}}\Ff(U_{i,k})$. Put
$I_{\circ}=\{i\mid i\geq i_{\circ}\text{ and }i\in I\}$ and
$J=I\coprod I_{\circ}$. Order $J$ by declaring $j'\leq j''$ if the
images of $j'$ and $j''$ in $I$ satisfy $j'\leq j''$ and it is not
the case that $j'\in I_{\circ}$ and $j''\in I$. The $U_i$ and
$U_{i,k}$, indexed by $J$, form a refinement $Q$ of $P$ such that
$\ss_Q\ne\ss'_Q$ and the image of $f$ in $Q(V)$ belongs to the image
of $Q(V_k)$.

\begin{enonce*}{Lemma 9.4}\etiquette[9.4]{VI.lem9.4}
Let $P=(U_i)_{i\in I}$ satisfy the hypotheses of Lemma
\hyperref[VI.lem9.2]{9.2}. There is a refinement $Q$ of $P$ such that
$\ss_Q\ne\ss'_Q$ and such that, for every finite covering $(V_k)$ of
an object $V$ of $\Ss$ and every $f\in P(V)$, the image of $f$ in
$Q(V)$ belongs to the image of one of the $Q(V_k)$.
\end{enonce*}

Let $E$ be the set of triples consisting of an object $V$ of $\Ss$, a
finite covering $(V_k)$ of $V$, and an element $f\in P(V)$. Let
$\leq$ be a well-ordering of $E$, and let $\bar E$ be the set obtained
from $E$ by adjoining a greatest element, denoted by $\infty$. By
transfinite induction on $e\in\bar E$, we shall define refinements
$Q_e$ of $P$ satisfying $\ss_{Q_e}\ne\ss'_{Q_e}$ and such that
\begin{enumerate}
\renewcommand{\theenumi}{\roman{enumi}}
\renewcommand{\labelenumi}{(\theenumi)}
\item if $e'\leq e$, then $Q_e$ refines $Q_{e'}$;

\item if $e=(V,(V_k),f)$, then the image of $f$ in $Q_e(V)$ belongs
  to the image of one of the $Q_e(V_k)$.
\end{enumerate}

Suppose the $Q_{e'}$ already defined for $e'<e$. If $e$ is the first
element of $E$, put $Q'_e=P$. If $e$ has a predecessor $e-1$, put
$Q'_e=Q_{e-1}$. Otherwise, let $Q'_e$ be the projective system with
index set $\bigcup_{e'<e}I_{e'}$ which refines the $Q_{e'}$ for
$e'<e$. In this case,
$$
Q'_e(\Ff)=\varinjlim_{e'<e}Q_{e'}(\Ff),
$$
so that, in every case,\pageoriginale $Q'_e$ refines the $Q_{e'}$ for
$e'<e$ and satisfies $\ss_{Q'_e}\ne\ss'_{Q'_e}$.

The required projective system $Q_e$ is obtained by applying Lemma
\hyperref[VI.lem9.3]{9.3} to $Q'_e$ and
$e=(V,(V_k),f)$ (\resp by taking $Q_e=Q'_e$ if $e=\infty$).

The projective system $Q_{\infty}$ satisfies Lemma
\hyperref[VI.lem9.4]{9.4}.

Lemma \hyperref[VI.lem9.4]{9.4} makes it possible to define,
inductively on $n$, a sequence $Q_n=(U_i)_{i\in I_n}$ of projective
systems in $\Ss$ such that
\begin{enumerate}
\renewcommand{\theenumi}{\roman{enumi}}
\renewcommand{\labelenumi}{(\theenumi)}
\item $Q_0=P$;

\item $Q_{n+1}$ refines $Q_n$;

\item $\ss_{Q_n}\ne\ss'_{Q_n}$;

\item for every covering $f_k:V_k\to V$ and every $f\in Q_n(V)$,
  the image of $f$ in $Q_{n+1}(V)$ belongs to the image of one of the
  $Q_{n+1}(V_k)$.
\end{enumerate}

The projective system $Q$, whose index set is the union of the $I_n$
and which extends the various $Q_n$, then satisfies Lemma
\hyperref[VI.lem9.2]{9.2}. The proof also shows:

\begin{enonce*}{Corollary 9.5}\etiquette[9.5]{VI.cor9.5}
Let $\Ss$ be a site in which fiber products are representable and
such that every covering in $\Ss$ admits a finite subcovering. Put
$c=\sup(\aleph_0,\card(\text{Fl }\Ss))$. Then there is a very dense
set $X$ of points of $\Ss$ such that
\begin{enumerate}
\renewcommand{\theenumi}{\roman{enumi}}
\renewcommand{\labelenumi}{(\theenumi)}
\item $\card(X)\leq2^c$;

\item if $x\in X$ and $U\in\Ob\Ss$, then $\card(U_x)\leq c$.
\end{enumerate}
\end{enonce*}

\begin{thebibliography}{99}\pageoriginale
\bibitem{sga4.exp.VI.bib.1} P.\textsc{Gabriel} and M.\textsc{Zisman}: Homotopy Theory
  and Calculus of Fractions, Ergebnisse der Mathematik, Bd 35.

\bibitem{sga4.exp.VI.bib.2} J. \textsc{Giraud}: M\'ethode de la descente. M\'emoire de la S.M.F.

\bibitem{sga4.exp.VI.bib.3} A. \textsc{Grothendieck} and J. \textsc{Dieudonn\'e}: \'El\'ements de
  G\'eom\'etrie Alg\'ebrique I, 2nd edition.

\bibitem{sga4.exp.VI.bib.4} A. \textsc{Grothendieck} and J. \textsc{Dieudonn\'e}: \'El\'ements de
  G\'eom\'etrie Alg\'ebrique IV I, I.H.E.S.~${\rm n^{\circ}}$

\bibitem{sga4.exp.VI.bib.5} \SGA. 1 VI, by A. \textsc{Grothendieck}, I.H.E.S.

\bibitem{sga4.exp.VI.bib.6} \SGA. 3 IV 6.3 by M. \textsc{Demazure}, Lecture Notes
  ${\rm n^{\circ}}$, Springer Verlag.

\bibitem{sga4.exp.VI.bib.7} A. \textsc{Grothendieck}: Sur quelques points d'alg\`ebre
  homologique, Tohoku Math. J.
\end{thebibliography}

\ifx\danslelivre\undefined
\end{document}
\fi

% END INCLUDED SGA4_English_current_buildable_source_20260906/source/06/drafts/SGA4_VI_LINES_7182_7367_English.texfrag


% END INCLUDED SGA4_English_current_buildable_source_20260906/source/06/06_English.tex


\makeatletter
\def\input@path{{../07/}}
\makeatother
% BEGIN INCLUDED SGA4_English_current_buildable_source_20260906/source/07/07_English.tex
\ifx\danslelivre\undefined
% Release engine: XeLaTeX (the LuaTeX compatibility layer breaks Xy-pic output).
\documentclass[10pt]{../fichiers-configuration-latex/smfart}
% BEGIN INCLUDED SGA4_English_current_buildable_source_20260906/source/fichiers-configuration-latex/paquets.tex
\usepackage{polyglossia}
%\usepackage[francais]{babel}
\usepackage{fontspec}
\usepackage[math-style=ISO,bold-style=upright,sans-style=upright]{unicode-math}
%\usepackage[mathletters]{ucs}
%\usepackage[utf8]{inputenc}
%\usepackage[T1]{fontenc}
\usepackage[unicode,hypertexnames=false,colorlinks=true,linkcolor=blue,citecolor=Sepia,urlcolor=magenta]{hyperref}
\setmainlanguage{french}
\setmainfont[Ligatures=TeX]{Linux Libertine O}
% variante possible : TeXGyrePagella ?
\setmathfont{XITS Math}

% maths
\usepackage{amsmath}
%\usepackage{amssymb}
%\usepackage{mathrsfs}
%\usepackage{stmaryrd}
%\usepackage{euscript}

% couleurs
\usepackage[usenames,dvipsnames]{xcolor}

% graphiques
\usepackage{graphicx}
\usepackage[all]{xy}
%\usepackage{tikz}
%\usetikzlibrary{matrix,arrows,calc}
%\tikzset{column sep=3em,row sep=2em,text height=1.5ex, text depth=0.25ex,>=angle 90}
%\tikzset{description/.style={fill=white,inner sep=1pt}}

% hyperliens, 'etiquettes
\usepackage{xr-hyper}
\usepackage{../fichiers-configuration-latex/smfthm}
%\ifx\cacheretiquettes\undefined
%\usepackage[inline]{showlabels} \else \fi % pour rendre visibles les 'etiquettes
%\usepackage{../fichiers-configuration-latex/makeidx} % pour faire un index
%\usepackage[pagebackref,unicode]{hyperref}  % pour voir o`u sont cit'ees les r'ef'erences
% hyperref is loaded once above with the cumulative-reader destination policy.
%\hypersetup{breaklinks=true} %permet le retour `a la ligne dans les liens trop longs. Incompatible avec pdflatex
%\renewcommand*{\backref}[1]{{\small [cit\'e page(s)~#1.]}}

\usepackage{xspace}
\usepackage{pifont} % pour expose XVII

% END INCLUDED SGA4_English_current_buildable_source_20260906/source/fichiers-configuration-latex/paquets.tex

% BEGIN INCLUDED SGA4_English_current_buildable_source_20260906/source/fichiers-configuration-latex/numerotation.tex
\renewcommand{\thefootnote}{\roman{footnote}} % pour mettre les notes en bas de page sous la forme (ii) et non (2), qui ressemble trop (?) a un exposant.

\newcounter{pageoriginal}
\newcommand{\pageoriginale}{\refstepcounter{pageoriginal}\marginpar{\footnotesize\textbf{\thepageoriginal}}\xspace}
\let\pageoriginaled\pageoriginale

\usepackage{ifthen}
\newboolean{original}
\def\sisi{\ifthenelse{\boolean{original}}}
\setboolean{original}{false}

\newcounter{cnote}
\def\nde#1{\ifthenelse{\boolean{original}}{}{\stepcounter{cnote}\def\thefootnote{\arabic{cnote}}\protect\footnote[\value{cnote}]{N.D.E.: #1}}}
\def\ndemark{\ifthenelse{\boolean{original}}{}{\stepcounter{cnote}\def\thefootnote{\arabic{cnote}}\protect\footnotemark[\value{cnote}]}}
\def\ndetext#1{\ifthenelse{\boolean{original}}{}{\def\thefootnote{\arabic{cnote}}\protect\footnotetext[\value{cnote}]{N.D.E.: #1}}}

% END INCLUDED SGA4_English_current_buildable_source_20260906/source/fichiers-configuration-latex/numerotation.tex

% BEGIN INCLUDED SGA4_English_current_buildable_source_20260906/source/fichiers-configuration-latex/macros.tex
\setmainlanguage{french}
\setotherlanguage{english}


%% Fixed, engine-independent release identifier.  The inherited dynamic Git
%% probe used Lua io.popen, exposed build-state-dependent output, and printed
%% a visible question mark outside a Git checkout.
\def\showgitstatus{English release r2, 2026-07-23}


% macros pour LaTeX
\renewcommand{\mathsf}{}

% Police différente pour U+2113 SCRIPT SMALL L
\setmathfont[range="2113]{latinmodern-math.otf}

% Give every inherited semantic label its own stable PDF destination instead
% of inheriting whichever automatic equation/page anchor happened to be live.
% The legacy label remains as a compatibility alias used by existing edges.
\makeatletter
\def\etiquette[#1]#2{%
  \Hy@raisedlink{\hyper@anchorstart{sga4.semantic.#2}\hyper@anchorend}%
  \def\@currentHref{sga4.semantic.#2}%
  \def\@currentlabel{#1}%
  \label{#2}%
}
\makeatother
\def\eqhyperref[#1]#2{(\hyperref[#1]{#2})}
\def\pagehyperref[#1]#2{(page \hyperref[#1]{#2})}
\def\hyperrefext[#1][#2]#3{\hyperref[#2]{#1 #3}}
\def\eqhyperrefext[#1][#2]#3{\hyperref[#2]{#1 #3}}     %cf. exposes 17
\def\neweqhyperrefext[#1][#2]#3{\hyperref[#2]{#1 #3}}  %cf. exposes 17

\def\shat#1{\protect\widehat{#1}}
\def\stilde#1{\protect\widetilde{#1}}
%%%%%%%%%%%%%%%%%%


\def\etcn{etc.\xspace}
\def\etco{etc...\xspace}
\def\etc{\sisi{\etco}{\etcn}}
\def\ldotsn{\sisi{...}{\etcn}}

\def\cf{{cf.}\kern.3em}
\def\Cf{{Cf.}\kern.3em}
\def\ie{{i.e.}\kern.3em}
\def\iev{{i.e.},\xspace}
\def\loccito{loc.\kern2pt cit.\xspace}
\def\loccitn{\textit{loc.\kern2pt cit.}\xspace}
\def\loccit{\sisi{\loccito}{\loccitn}}
\def\EGA{\textup{EGA}\kern.2em}
\def\SGA{\textup{SGA}\kern.3em}

\def\quoi{}

\def\quav{\quad,\quad}
\def\quaet{\quad \text{et}\quad}
\def\fibreproduct#1#2#3{#1 ×_{#3} #2}
\let\fprod\fibreproduct
\def\ccheck#1#2{\mbox{${\check{\check{\rm #1}}}^{\lower2pt\hbox{$\scriptstyle #2$}}$}}
\def\sheaf#1{\protect\mathcal{#1}}
\def\vieuxsheaf#1{{#1}} % cf. problème avec sheaf{a\Phi_i} par exemple...
\DeclareMathOperator{\Hom}{Hom}
\DeclareMathOperator{\can}{can}
\DeclareMathOperator{\Ind}{Ind}
\DeclareMathOperator{\ind}{ind}
%\DeclareMathOperator{\Hhomtop}{\mathcal{H}\!\mathit{omtop}}
\DeclareMathOperator{\Hhomtop}{\mathbf{Homtop}}
\DeclareMathOperator{\Point}{Point}
\DeclareMathOperator{\Pro}{Pro}
\DeclareMathOperator{\Top}{Top}
\DeclareMathOperator{\Mod}{Mod}
\DeclareMathOperator{\Cart}{Cart}
%\DeclareMathOperator{\Mmorsite}{\mathcal{M}\!\mathit{orsite}}
\DeclareMathOperator{\Mmorsite}{\mathbf{Morsite}}
\DeclareMathOperator{\loc}{loc}
\DeclareMathOperator{\Loc}{Loc}
\DeclareMathOperator{\Pic}{Pic}
\DeclareMathOperator{\Et}{Et}
\DeclareMathOperator{\card}{card}
\DeclareMathOperator{\Ens}{Ens}
\DeclareMathOperator{\h}{\underline{H}}
\DeclareMathOperator{\Fonct}{Fonct}
\DeclareMathOperator{\Ker}{Ker}
\DeclareMathOperator{\but}{but}
\DeclareMathOperator{\Ob}{Ob}
\DeclareMathOperator{\ob}{ob}
\DeclareMathOperator{\Lib}{Lib}
\DeclareMathOperator{\esj}{esj}
\DeclareMathOperator{\id}{id}
\DeclareMathOperator{\Id}{Id}
\DeclareMathOperator{\ad}{ad}
\DeclareMathOperator{\Iim}{Im}
\DeclareMathOperator{\pr}{pr}
\DeclareMathOperator{\Cat}{Cat}
\DeclareMathOperator{\car}{car}
\DeclareMathOperator{\SysInd}{SysInd}
\DeclareMathOperator{\Sup}{Sup}
\DeclareMathOperator{\Coker}{Coker}
\DeclareMathOperator{\Filt}{Filt}
\DeclareMathOperator{\acc}{acc}
\DeclareMathOperator{\sur}{sur}
\DeclareMathOperator{\cat}{cat}
\DeclareMathOperator{\Ord}{Ord}
\DeclareMathOperator{\Gr}{Gr}
\DeclareMathOperator{\Cov}{Cov}
\DeclareMathOperator{\Caf}{Caf}
\DeclareMathOperator{\coker}{coker}
\DeclareMathOperator{\coim}{coim}
\DeclareMathOperator{\im}{im}
\DeclareMathOperator{\Ouv}{Ouv}
\DeclareMathOperator{\diag}{diag}
\DeclareMathOperator{\Esp}{Esp}
\DeclareMathOperator{\TOP}{TOP}
\DeclareMathOperator{\Unif}{Unif}
\DeclareMathOperator{\Iind}{\mathbf{Ind}}
\DeclareMathOperator{\Hhomcart}{\mathbf{Homcart}}
\DeclareMathOperator{\Hhomtopess}{\mathbf{Homtopess}}
\DeclareMathOperator{\Ppoints}{\mathbf{Points}}
\DeclareMathOperator{\Ppoint}{\mathbf{Point}}
\DeclareMathOperator{\Ttors}{\mathbf{Tors}}
\DeclareMathOperator{\Ssex}{\mathbf{Sex}}
\DeclareMathOperator{\Ffib}{\mathbf{Fib}}
\DeclareMathOperator{\Vvois}{\mathbf{Vois}}
\DeclareMathOperator{\Ppointess}{\mathbf{Pointess}}
\DeclareMathOperator{\scrHom}{\mathbf{Hom}}
\DeclareMathOperator{\sheafPic}{\mathbf{Pic}}
\DeclareMathOperator{\sheafAut}{\mathbf{Aut}}
\DeclareMathOperator{\Hhomtoplocan}{\mathbf{Homtoplocan}}
\DeclareMathOperator{\Hhomtopan}{\mathbf{Homtopan}}
\DeclareMathOperator{\Pptgeom}{\mathbf{Ptgéom}}
\DeclareMathOperator{\Morsite}{Morsite}
\DeclareMathOperator{\Fib}{Fib}
\DeclareMathOperator{\Fais}{Fais}
\DeclareMathOperator{\kar}{kar}
\DeclareMathOperator{\Groupes}{Groupes}
\DeclareMathOperator{\Restr}{Restr}
\DeclareMathOperator{\Isom}{Isom}
\DeclareMathOperator{\Prol}{Prol}
\DeclareMathOperator{\Res}{Res}
\DeclareMathOperator{\Tors}{Tors}
\DeclareMathOperator{\esp}{esp}
\DeclareMathOperator{\es}{es}
\DeclareMathOperator{\od}{od}
\DeclareMathOperator{\gr}{gr}
\DeclareMathOperator{\sob}{sob}
\DeclareMathOperator{\dfn}{dfn}
\DeclareMathOperator{\Sch}{Sch}
\DeclareMathOperator{\ouv}{ouv}
\DeclareMathOperator{\cons}{cons}
\DeclareMathOperator{\Nerf}{Nerf}
\DeclareMathOperator{\Points}{Points}
\DeclareMathOperator{\Espsob}{Espsob}
\DeclareMathOperator{\Homtop}{Homtop}
\DeclareMathOperator{\Pt}{Pt}
\DeclareMathOperator{\Quot}{Quot}
\DeclareMathOperator{\glob}{glob}
\DeclareMathOperator{\Homfib}{Homfib}
\DeclareMathOperator{\Qouv}{Qouv}
\DeclareMathOperator{\Quotouv}{Quotouv}
\DeclareMathOperator{\ens}{ens}
\DeclareMathOperator{\ab}{ab}
\DeclareMathOperator{\an}{an}
\DeclareMathOperator{\fs}{fs}
\DeclareMathOperator{\red}{red}
\DeclareMathOperator{\Ab}{Ab}
\DeclareMathOperator{\cosq}{cosq}
\DeclareMathOperator{\Ann}{Ann}
\DeclareMathOperator{\Zar}{Zar}


%%%%%%%%%%%%
%%%%% SGA4-2
\DeclareMathOperator{\Sex}{Sex}
\DeclareMathOperator{\vois}{vois}
\DeclareMathOperator{\Vois}{Vois}
\DeclareMathOperator{\Ext}{Ext}
\DeclareMathOperator{\ext}{ext}
\DeclareMathOperator{\Fer}{Fer}
\DeclareMathOperator{\Add}{Add}
\DeclareMathOperator{\Cocart}{Cocart}
\DeclareMathOperator{\cosk}{cosk}
\DeclareMathOperator{\plat}{plat}

\DeclareMathOperator{\faisc}{faisc}
\DeclareMathOperator{\sk}{sk}
\DeclareMathOperator{\ch}{ch}

\def\HH{\mathbf{H}}
\def\qc{\mathrm{qc}}
\def\et{\mathrm{et}}
\def\coh{\mathrm{coh}}

\def\Hvr#1{%
\begin{array}[b]{l}
{\kern2pt\scriptstyle r}\\[-8pt]
{\kern1pt\scriptstyle \vee}\\[-5pt]
H^{#1}
\end{array}}

\def\Hvv#1{%
\begin{array}[b]{l}
\kern2pt\rotatebox{90}{${\scriptscriptstyle\ll}$}\\[-6pt]
H^{#1}
\end{array}}

\def\Hvi#1{%
\begin{array}[b]{l}
{\kern0.2pt\scriptscriptstyle \infty}\\[-8pt]
{\kern1pt\scriptstyle \vee}\\[-5pt]
H^{#1}
\end{array}}

\def\Hv#1{%
\begin{array}[b]{l}
{\kern1pt\scriptstyle \vee}\\[-5pt]
H^{#1}
\end{array}}

\DeclareMathOperator{\cart}{cart}
\DeclareMathOperator{\Ffaisc}{\mathbf{Faisc}}


\def\Ev#1{%
\begin{array}[b]{l}
{\kern1pt\scriptstyle \vee}\\[-5pt]
E'_{#1}
\end{array}}

\DeclareMathOperator{\Cont}{Cont}
\DeclareMathOperator{\Ccont}{\mathscr{C}\!\mathit{ont}}
\DeclareMathOperator{\Site}{Site}
\newcommand\MM{\mathbf{M}}
\def\limtop#1{{\displaystyle{\varprojlim_{#1}\!\mbox{\rm top }}}}
\def\Limtop{\mathop{\varprojlim\!\mbox{\rm top}}}
\DeclareMathOperator{\spec}{spec}
\def\ss{\scriptstyle{\mathscr{S}}}
\DeclareMathOperator{\Gal}{Gal}
\def\smallm{\scriptstyle{\mathscr{M}}}
\DeclareMathOperator{\End}{End}
\DeclareMathOperator{\ass}{ass}
\DeclareMathOperator{\zar}{zar}
\DeclareMathOperator{\Ogg}{Ogg}
\DeclareMathOperator{\res}{res}
\DeclareMathOperator{\tr}{tr}
\DeclareMathOperator{\Supp}{Supp}
\DeclareMathOperator{\codim}{codim}
\DeclareMathOperator{\rk}{rk}
\DeclareMathOperator{\cur}{cur}
\DeclareMathOperator{\cdqc}{cdqc}
\newcommand\DD{\mathbf{D}}
\DeclareMathOperator{\cl}{cl}
\def\JJ{\mathbf{J}}
\def\FF{\mathbf{F}}
\def\BB{\mathbf{B}}
\DeclareMathOperator{\Alget}{Alget}
\DeclareMathOperator{\Ass}{Ass}
\DeclareMathOperator{\supp}{supp}
\DeclareMathOperator{\rad}{rad}
\DeclareMathOperator{\prof}{prof}
\DeclareMathOperator{\cint}{int}
\DeclareMathOperator{\sic}{sic}
\DeclareMathOperator{\Tor}{Tor}
\DeclareMathOperator{\tors}{tors}
\DeclareMathOperator{\Sym}{Sym}
\DeclareMathOperator{\parf}{parf}
\DeclareMathOperator{\tor}{tor}
\DeclareMathOperator{\defn}{defn}
\DeclareMathOperator{\Div}{Div}
\DeclareMathOperator{\Pol}{Pol}
\DeclareMathOperator{\Fex}{Fex}
\DeclareMathOperator{\Lie}{Lie}
\DeclareMathOperator{\tot}{tot}
\DeclareMathOperator{\Aut}{Aut}
\DeclareMathOperator{\HOM}{HOM}
\DeclareMathOperator{\pch}{pch}
\DeclareMathOperator{\Ch}{Ch}
\DeclareMathOperator{\EXT}{EXT}
\DeclareMathOperator{\PIC}{PIC}
\DeclareMathOperator{\Rf}{Rf}
\DeclareMathOperator{\inj}{inj}
\newcommand\VV{\mathbf{V}}

\def\Xv{%
\begin{array}[b]{l}
{\kern1pt\scriptstyle \vee}\\[-5pt]
X
\end{array}}

%%%%%%%
\newcommand\TT{\mathbf{T}}
\newcommand\XX{\mathbf{X}}
\newcommand\EE{\mathbf{E}}
%%%%%%

\def\qedsymbol{\textup{C.Q.F.D.}}
\let\CQFD\qed
\let\cqfd\qed

%%%%%\bigdot
\font\bigsymb=cmsy10 at 4pt
\def\bigdot{{\kern1.2pt\raise 1.5pt\hbox{\bigsymb\char15}}}
%%%%%

\def\Lright{\mathop{\vcenter{\hbox{\hbox{$\longrightarrow$}\lower8pt \hbox{\kern-16pt\hbox{$\longrightarrow$}}}}}}

\def\Aa{{\mathcal A}}
\def\Bb{{\mathcal B}}
\def\Cc{{\mathcal C}}
\def\Dd{{\mathcal D}}
\def\Ee{{\mathcal E}}
\def\Ff{{\mathcal F}}
\def\Gg{{\mathcal G}}
\def\Hh{{\mathcal H}}
\def\Ii{{\mathcal I}}
\def\Jj{{\mathcal J}}
\def\Kk{{\mathcal K}}
\def\Ll{{\mathcal L}}
\def\Mm{{\mathcal M}}
\def\Nn{{\mathcal N}}
\def\Oo{{\mathcal O}}
\def\Pp{{\mathcal P}}
\def\Qq{{\mathcal Q}}
\def\Rr{{\mathcal R}}
\def\Ss{{\mathcal S}}
\def\Tt{{\mathcal T}}
\def\Uu{{\mathcal U}}
\def\Vv{{\mathcal V}}
\def\Ww{{\mathcal W}}
\def\Xx{{\mathcal X}}
\def\Yy{{\mathcal Y}}
\def\Zz{{\mathcal Z}}

\newcommand{\CC}{\mathbf{C}}
\newcommand{\GG}{\mathbf{G}}
\newcommand{\KK}{\mathbf{K}}
\newcommand{\LL}{\mathbf{L}}
\newcommand{\NN}{\mathbf{N}}
\newcommand{\PP}{\mathbf{P}}
\newcommand\QQ{\mathbf{Q}}
\newcommand\RR{\mathbf{R}}
\newcommand\ZZ{\mathbf{Z}}

\def\to{\mathchoice{\longrightarrow}{\rightarrow}{\rightarrow}{\rightarrow}}
\def\from{\mathchoice{\longleftarrow}{\leftarrow}{\leftarrow}{\leftarrow}}
\def\mto{\mathchoice{\longmapsto}{\mapsto}{\mapsto}{\mapsto}}
\def\mfrom{\mathrel{\from\joinrel\mapstochar\,}}
\def\hto{\mathrel{\lhook\joinrel\to}}
\def\hfrom{\mathrel{\from\joinrel\rhook}}

\def\implies{\mathchoice{\Longrightarrow}{\Rightarrow}{\Rightarrow}{\Rightarrow}}
\def\implied{\mathchoice{\Longleftarrow}{\Leftarrow}{\Leftarrow}{\Leftarrow}}
\def\iff{\mathchoice{\Longleftrightarrow}{\Leftrightarrow}{\Leftrightarrow}{\Leftrightarrow}}

\def\To#1{\mathchoice{\xrightarrow{\labelstyle~~#1~}}{\stackrel{#1}{\longrightarrow}}{}{}}
\def\Mto#1{\mathrel{\mapstochar\To{#1}}}
\def\From#1{\mathchoice{\xleftarrow{\labelstyle~~#1~}}{\stackrel{#1}{\longleftarrow}}{}{}}

\newcommand{\isomto}{\overset{\sim}{\longrightarrow}}
\newcommand{\isomfrom}{\overset{\sim}{\longleftarrow}}

\newcommand{\IFF}{\unskip~$\iff$~\ignorespaces}
\newcommand{\THEN}{\unskip~$\implies$~\ignorespaces}

\let\leq\leqslant
\let\geq\geqslant
\let\le\leqslant
\let\ge\geqslant

\def\cf{{cf.}\kern.3em}
\def\Cf{{Cf.}\kern.3em}
\def\ie{{i.e.}\kern.3em}
\def\iev{{i.e.},\xspace}
%\def\loccito{loc.\kern2pt cit.\xspace}
%\def\loccitn{\textit{loc.\kern2pt cit.}\xspace}
%\def\loccit{\sisi{\loccito}{\loccitn}}
\def\EGA{\textup{EGA}\kern.2em}
\def\SGA{\textup{SGA}\kern.3em}

\let\oldS\S
\def\S{\oldS\kern .15em}
\def\ptbl{.\kern .2em}
\def\resp{resp.\kern.3em}
\def\Exp{Exp.\kern.3em}

\let\bar\overline
\let\hat\widehat
\let\tilde\widetilde
\let\dpl\displaystyle
\let\ldots\dots
\let\emptyset\varnothing
\let\skpt\skippointrait
\let\epsilon\varepsilon

\def\sfrac#1#2{{#1}/{#2}}

\DeclareMathOperator{\cd}{cd}
%\DeclareMathOperator{\Hhom}{\mathcal{H}\!\mathit{om}}
\DeclareMathOperator{\Hhom}{\mathbf{Hom}}
%\DeclareMathOperator{\Eext}{\mathcal{E}\!\mathit{xt}}
\DeclareMathOperator{\Eext}{\mathbf{Ext}}
\DeclareMathOperator{\Spec}{Spec}
\DeclareMathOperator{\Tr}{Tr}

\let\emptyset=\varnothing

%%%%%%%%%%%%%%%
%%% expose XVII
%%%%%%%%%%%%%%%

\makeatletter
\let\old@optexpose\optexpose
\newcommand{\empty@optexpose}[1]{}
\newcommand{\paren@optexpose}[1]{\old@optexpose{#1}(}
\newcommand{\refnoexpose}[1]{\let\optexpose\empty@optexpose\ref{#1}\let\optexpose\old@optexpose}%\refnoexpose{X.cor5.2} engendre "5.2" au lieu de "X 5.2"
\newcommand{\neweqref}[1]{\let\optexpose\paren@optexpose\textup{\ref{#1})}\let\optexpose\old@optexpose}%\neweqref{XII.eq4.1} engendre "XII (4.1)" au lieu de "(XII 4.1)"
\makeatother
\newcommand{\refvb}[2][VI${}_{\text{B}}$ ]{\sisi{#1\refnoexpose{#2}}{\ref{#2}}}

\newcommand{\leftexp}[2]{{\vphantom{#2}}^{#1}{#2}}
\newcommand{\AAA}{\mathbf{A}}
\newcommand{\Sf}{\mathfrak{S}}
\DeclareMathOperator*{\quotprojlim}{\text{``$\varprojlim$''}}
\DeclareMathOperator*{\quotinjlim}{\text{``$\varinjlim$''}}
\DeclareMathOperator{\Fl}{F\sisi{\ell}{l}}
\DeclareMathOperator{\Pl}{P\sisi{\ell}{l}}
%\DeclareMathOperator{\cl}{c\sisi{\ell}{l}}
\DeclareMathOperator{\tdf}{tdf}
\DeclareMathOperator{\rg}{rg}
\newcommand{\iInd}{\sisi{I}{i}nd}
\newcommand{\iind}{ind}
\newcommand{\ppro}{pro}
\newcommand{\tTor}{\sisi{T}{t}or}
\newcommand{\ttor}{tor}
\newcommand{\chap}{Chap.}
\newcommand{\Ffex}{\Ff\mathit{ex}}
\newcommand{\sheafk}{\underline{k}}
\newcommand{\sheafn}{\underline{n}}
\newcommand{\sheafGamma}{\underline{\Gamma}}

% Repeated references to an already typeset footnote must not use
% \footnotemark without a matching \footnotetext: Hyperref would then create
% a live annotation to a destination that cannot exist.  This helper keeps the
% current printed footnote mark and links it to an explicit stable note target.
\makeatletter
\newcommand{\sgarepeatedfootnotemark}[1]{%
  \begingroup
  \edef\@thefnmark{\thefootnote}%
  \hyperlink{#1}{\@makefnmark}%
  \endgroup}
\makeatother

\sisi{
\DeclareMathOperator{\simarrow}{\rule{.15mm}{1.8mm}\kern -3.4pt\rightsquigarrow}
\newcommand{\sheafell}{\underline{\ell}}
\newcommand{\sheafZZ}{\underline{\ZZ}}
}{
\newcommand{\nderef}[1]{la N.D.E.~(\ref{#1}) page~\pageref{#1}}
}

%%%%%%%%%%%%%%%%%
%%%%% expose XVII
%%%%%%%%%%%%%%%%%

% END INCLUDED SGA4_English_current_buildable_source_20260906/source/fichiers-configuration-latex/macros.tex


\externaldocument{../01/01_English.tex}
\externaldocument{../02/02_English.tex}
\externaldocument{../03/03_English.tex}
\externaldocument{../04/04_English.tex}
\externaldocument{../05/05_English.tex}
\externaldocument{../05bis/05bis_English.tex}
\externaldocument{../06/06_English.tex}

\setboolean{original}{false}

\begin{document}

\begin{center}
VII. The étale site and topos of a scheme\\
A. Grothendieck
\end{center}

\selectlanguage{english}version: \showgitstatus
\tableofcontents

\else
\chapter{The étale site and topos of a scheme}
%\addtocontents{toc}{A. Grothendieck}
\begin{center}
A. Grothendieck
\end{center}
\fi
\phantomsection\MakeLinkTarget*{sga4.exp.VII}\label{sga4.exp.VII}

% BEGIN INCLUDED SGA4_English_current_buildable_source_20260906/source/07/drafts/SGA4_VII_LINES_43_396_English.texfrag
\setcounter{pageoriginal}{341}

In\pageoriginale the present exposé and the next, we develop the most
elementary properties concerning étale topology and cohomology. The
developments of the present exposé concern certain properties that are, for
the most part, valid for other, very different topologies, such as the
``fppf topology''. In the next exposé, we shall develop properties rather
special to the étale topology, arising from the very particular nature of
étale morphisms.

Here we partly follow three oral exposés by J.E. \textsc{Roos} (which he had
not been able to write up), notably in the proof of
\hyperrefext[VIII][VIII.thm6.3]{6.3}.

Unless expressly stated otherwise, it will be understood that the schemes
considered in the present exposé and the following ones are elements of the
fixed universe $𝒰$.

\section{The étale topology}\etiquette[1]{VII.sec1}

\subsection{}\etiquette[1.1]{VII.subsec1.1}

We denote by $(\Sch)$ the category of schemes (elements of the fixed
universe $𝒰$). Recall that a morphism $f:X\to Y$ in $(\Sch)$ is
called\pageoriginale \emph{étale} if it is locally of finite presentation
and flat, and if for every $y\in Y$, the fiber $X_y$ is discrete and its
local rings are finite \emph{separable} extensions of $k(y)$. Equivalently,
$f$ is locally of finite presentation and, for every affine $Y$-scheme $Y'$
(in the absolute sense) and every subscheme $Y'_0$ defined by a nilpotent
ideal, the map
$$
\Hom_Y (Y', X) \to \Hom_Y (Y'_{0}, X)
$$
is \emph{bijective}. For the most important properties of this notion, I
refer to \EGA IV \S\S 17 and 18, and, in the meantime, to \SGA 1 I and IV.


\subsection{}\etiquette[1.2]{VII.subsec1.2}

The \emph{étale topology} on $(\Sch)$ is the topology generated by the
pretopology for which, for every $X\in(\Sch)$, the set $\Cov X$ consists of
the families $(X_i\xrightarrow{u_i}X)_{i\in I}$ (indexed by an $I\in U$)
such that the $u_i$ are étale and
$X=\bigcup_i u_i(X_i)$ (in the set-theoretic sense).

It is convenient to associate with every $X$ the subcategory $\Et_{/X}$ of
$(\Sch)_{/X}$ consisting of the arrows $X'\to X$ that are étale. We endow it
with the ``induced topology'' (III) from the étale topology of $(\Sch)$
(also called the ``étale topology on $X$''), and denote the resulting site
by $X_{\text{ét}}$ (the ``étale site'' of $X$)
\footnote{Note, however, that it would be preferable to denote by
  $X_{\text{ét}}$ the topos $\widetilde{X_{\text{ét}}}$ defined by the
  étale site of $X$. For practical reasons, in this Seminar we shall retain
  the notation introduced here (which is that of the original seminar).}.
Note that every morphism of $\Et_{/X}$ is étale (\SGA 1 I 4.8), whence a
family $(X'_i\xrightarrow{u_i}X')_{i\in I}$ in $X_{\text{ét}}$ is covering
if and only if it is \emph{surjective}, \ie
$\bigcup_i u_i(X'_i)=X'$. One sees immediately (\cf
\hyperref[VII.prop3.1]{3.1}) that $X_{\text{ét}}$ is a $\Uu$-site, so the
results of Exposés I through VI apply. The topos
$\widetilde{X_{\text{ét}}}$ of $𝒰$-sheaves on $X_{\text{ét}}$ is the
\emph{étale topos} of $X$.

\subsection{}\etiquette[1.3]{VII.subsec1.3}

In what\pageoriginale follows, unless expressly stated otherwise, all the
topological notions considered in $\Et_{/X}$ or $(\Sch)$ are understood
with respect to the étale topology. Moreover, in language and notation we
shall commonly write $X$ in place of $X_{\text{ét}}$ or $Et_{/X}$. Thus a
``sheaf on $X$'' (understood: for the étale topology) means a sheaf on the
site $X_{\text{ét}}$. We denote by $\widetilde{X_{\text{ét}}}$ the category
of these sheaves, which is a $\Uu$-topos. If $F$ is an abelian sheaf on $X$,
we denote its cohomology groups by $H^i(X,F)$; these would be denoted
$H^i(X_{\text{ét}},F)$ in Exposé V.

\subsection{}\etiquette[1.4]{VII.subsec1.4}

Let $f:X\to Y$ be a morphism of schemes. The inverse-image functor
$$
f^{\ast} : Y' \mapsto \fprod{Y'}{X}{Y}
$$
induces a functor
\begin{equation*}
f^{\ast}_{\text{ét}} : \Et_{/Y} \to \Et_{/X}\tag{$\ast$}
\end{equation*}
that commutes with finite $\varprojlim$ and takes covering families to
covering families; it is therefore a morphism of sites
$$
f^{\ast}_{\text{ét}} : X_{\text{ét}} \to Y_{\text{ét}}{\quoi},
$$
and thus induces a functor on the categories of sheaves:
$$
f^{\text{ét}}_{\ast} : \widetilde{X_{\text{ét}}} \to \widetilde{Y_{\text{ét}}}
$$
given by
$$
f^{\text{ét}}_{\ast}  (F) = F \circ f^{\ast}_{\text{ét}}{\quoi}.
$$
Moreover, $f^{\text{ét}}_{\ast}$ admits a left adjoint
$$
f^{\ast}_{\text{ét}} : \widetilde{Y_{\text{ét}}} \to
\widetilde{X_{\text{ét}}}
$$
extending\pageoriginale the functor considered in $(\ast)$ and commuting
with arbitrary $\varinjlim$ and finite $\varprojlim$, \ie
$f^{\ast}_{\text{ét}}$ defines a \emph{morphism of topoi}
$$
f_{\text{ét}} : \widetilde{X_{\text{ét}}} \to \widetilde{Y_{\text{ét}}}{\quoi}.
$$

Clearly, for a composite of two morphisms of schemes
$$
X \xrightarrow{f} Y \xrightarrow{g} Z
$$
there are canonical isomorphisms
\begin{equation*}
\begin{cases}
(gf)^{\text{ét}}_{\ast} \simeq g_{\ast}^{\text{ét}}
  f_{\ast}^{\text{ét}} &\quav \mbox{ whence}\\
(gf)^{\ast}_{\text{ét}} \simeq f^{\ast}_{\text{ét}}
  g^{\ast}_{\text{ét}} &\quav \mbox{ \ie there is an
    isomorphism}\\
(gf)_{\text{ét}} \simeq g_{\text{ét}} f_{\text{ét}}&\quav
\end{cases}
\end{equation*}
(so that one obtains a ``pseudofunctor'' (\SGA 1 VI 8) from the category
$(\Sch)$ to the category $(\Top)$ of topoi $\in𝒰'$, where $𝒰'$ is the
smallest universe such that $𝒰\in𝒰'$; compare IV).

\subsection{Notation}\etiquette[1.5]{VII.subsec1.5}

In what follows, we shall often write $f^{\ast}$ and $f_{\ast}$ in place of
$f^{\ast}_{\text{ét}}$ and $f^{\text{ét}}_{\ast}$. If $f:X\to Y$ is a
morphism of schemes, the functors $R^if_{\ast}^{\text{ét}}$ will therefore
simply be denoted $R^if_{\ast}$. In this notation, recall the Leray spectral
sequence for $f$ and an abelian sheaf $F$ on $X$
(\hyperrefext[V][V.sec5]{5}):
$$
H^{\ast} (X, F) \Longleftarrow E^{p. q}_2 = H^{p}(Y, R^q f_\ast (F)){\quoi}.
$$
Likewise, if there are two morphisms $f:X\to Y$ and $g:Y\to Z$, there is a
Leray spectral sequence (the spectral sequence of composite functors)
$$
R^{\ast} (gf)_{\ast} (F) \Longleftarrow E^{p. q}_2 = R^{p} g_{\ast} (R^q
f_\ast (F)){\quoi}.
$$

\subsection{}\etiquette[1.6]{VII.subsec1.6}

When\pageoriginale $f:X\to Y$ is itself an étale morphism, $X$ may be
regarded as an object of $Y_{\text{ét}}$, and there is a canonical
\emph{isomorphism of sites}
$$
X_{\text{ét}} \xrightarrow{\sim} (Y_{\text{ét}})_{/X}{\quoi}.
$$
Under this isomorphism, the functor
$f^{\ast}_{\text{ét}}:Y'\mapsto\fprod{Y'}{X}{Y}$ of $(\ast)$ identifies with
the internal base-change functor in the category $\Et_{/Y}$. It follows that
the functor
$$
f^{\ast} : \widetilde{Y_{\text{ét}}} \to \widetilde{X_{\text{ét}}}
$$
is then isomorphic to the functor ``\emph{restriction} to $Y$''
$$
f^{\ast} (F) \simeq F \circ i_{X/Y}{\quoi},
$$
where $F$ is a sheaf on $Y_{\text{ét}}$ and
$i_{X/Y}=X_{\text{ét}}\to Y_{\text{ét}}$ is the evident functor (which
regards an étale scheme over $X$, $u:X'\to X$, as an étale scheme over $Y$
via $fu:X'\to Y$).

\subsection{Questions of universes}\etiquette[1.7]{VII.subsec1.7}

Note that if $X\neq∅$, then $X_{\text{ét}}$ \emph{is not} an element of the
chosen universe $𝒰$. Nevertheless, as we have noted, it is easy to see
(\hyperref[VII.prop3.1]{3.1}) that one can find a full subcategory $C$ of
$X_{\text{ét}}$, \emph{an element of} $𝒰$, satisfying the conditions of
the ``comparison lemma'' (III), so that the restriction functor induces an
equivalence $\widetilde{C}\to\widetilde{X_{\text{ét}}}$. Thus the
\emph{étale topos} is equivalent to a topos defined in terms of a site
$C\in𝒰$. In other words, $X_{\text{ét}}$ is a $U$-site, and hence by
\loccit the results of Exposés I through VI apply to this site.

\section{Examples of sheaves}\etiquette[2]{VII.sec2}\pageoriginale

\begin{enumerate}
\renewcommand{\theenumi}{\alph{enumi}}
\renewcommand{\labelenumi}{\theenumi)}
\item Let $F\in\Ob(\Sch)_{/X}$ be a scheme over $X$, and for every $X'$ over
  $X$ put
$$
F(X') = \Hom_X (X', F){\quoi}.
$$
The functor $(\Sch)^{\circ}_{/X}\to(\Ens)$ thus defined is a sheaf for the
étale topology (and even for the finer fpqc topology studied in \SGA 3 IV).
In other words, the étale topology on $(\Sch)$ is coarser than the canonical
topology. This is, essentially, the content of \SGA 1 VIII 5.1 (\cf also
\SGA 3 IV 6.3.1). A fortiori, the restriction of this sheaf to
$X_{\text{ét}}$ is a sheaf. It will again be denoted by $F$ when there is no
risk of confusion\footnote{\etiquette[6]{page6}Care must be taken, however,
  that if $X\neq∅$, the functor $\varphi$ that associates with
  $F\in\Ob(\Sch)_{/X}$ the corresponding sheaf $\varphi(F)$ is \emph{not}
  fully faithful, nor even faithful, and that $F$ cannot be reconstructed
  (up to isomorphism) from $\varphi(F)$. For example, if $S=\Spec k$, with
  $k$ an algebraically closed field, knowledge of $\varphi(F)$ is equivalent
  merely to knowledge of the underlying set of $F$ (by
  \hyperrefext[VIII][VIII.cor2.4]{2.4}). Compare IV}. Note also that the
functor thus obtained
\begin{equation*}
(\Sch)_{/X} \to \widetilde{X_{\text{ét}}}\tag{$\ast$}
\end{equation*}
commutes with finite $\varprojlim$ (this is immediate). This implies, for
example, that when $F$ is a group scheme (\resp\ldots) over $X$, the sheaf
it defines is a sheaf of groups (\resp\ldots). Note that the functor
$\Et_{/X}\to\widetilde{X_{\text{ét}}}$ induced by (*) is simply the
canonical functor, associating with every
$X'\in\Ob X_{\text{ét}}$ the functor on $X_{\text{ét}}$ that it
represents. It is therefore an isomorphism of the category $\Et_{/X}$ onto a
full subcategory of the étale topos $\widetilde{X_{\text{ét}}}$, by which we
generally identify an $X\in\Ob X_{\text{ét}}$ with the corresponding sheaf,
denoted $\widetilde{X}$ or simply $X$.

Clearly, $H^{\circ}(X,F)=\Hom_{(\Sch)_{/X}}(X,F)$.

Let us also give an interpretation of $H^1(X,F)$ when $F$ is a group
prescheme over $X$ (commutative, if desired, so that the definition of
$H^1(X,F)$ falls within Exposé V;\pageoriginale for the general case one may
consult J. Giraud's thesis\footnote{\etiquette[7]{page7}}). Familiar
arguments, which I shall not repeat here, then show that $H^1(X,F)$ is
canonically isomorphic to the group of classes (up to isomorphism) of
\emph{sheaves of sets on $X_{\text{ét}}$ that are principal homogeneous
under $F$\footnote{or $F$-torsors in the terminology of \loccit.}}. When
$F$ is \emph{affine} over $X$, \SGA 1 VIII 2.1 implies (again by standard
arguments, \cf J. \textsc{Giraud}'s thesis) that $H^1(X,F)$ is also the
group of classes of \emph{schemes} $P$ over $X$ on which $F$ acts on the
right and which are ``principal homogeneous bundles over $X$ for the étale
topology'', \ie locally trivial for the étale topology\footnote{or
  representable $F$-\emph{torsors} in the now accepted terminology}.
\end{enumerate}

\begin{enonce*}[remark]{Remarks 2.1}\etiquette[2.1]{VII.rem2.1}
Care must be taken that if, for every scheme $Z$ over $X$, the associated
sheaf on $X_{\text{ét}}$ is denoted by $\varphi_X(Z)$, and if $f:Y\to X$ is
a morphism of schemes, there is an evident homomorphism (functorial in $Z$)
\begin{equation*}
f^{\ast} (\varphi_X (Z)) \to \varphi_Y (Z_Y),\quad Z_Y =
\fprod{Z}{Y}{X}{\quoi},\tag{$\ast$}
\end{equation*}
\ie an evident homomorphism
$$
\varphi_X (Z) \to f_{\ast} (\varphi_Y (Z_Y)){\quoi},
$$
namely the homomorphism, functorial in $X'\in\Ob X_{\text{ét}}$,
$$
\Hom_X (X', Z) \to \Hom_Y (X'_{Y}, Z_{Y}){\quoi}.
$$
But note that \emph{in general $(\ast)$ is not an isomorphism, \ie the
formation of the étale sheaf associated with a relative scheme does not
commute with inverse-image functors}. Nevertheless, $(\ast)$ is an
isomorphism in the special case in which $Z$ is étale over $X$. Likewise,
the functor\pageoriginale
$$
\varphi_{X} : \Sch_{/X} \to \widetilde{X_{\text{ét}}}
$$
\emph{is not faithful} if $X\neq∅$\footnote{\Cf the note at the bottom of
  page \pagehyperref[page6]{6}}; nevertheless, its restriction to
$X_{\text{ét}}$, which is simply the canonical functor
$X_{\text{ét}}\to\widetilde{X_{\text{ét}}}$, is \emph{fully faithful},
because by a) the topology of $X_{\text{ét}}$ is coarser than its canonical
topology.

\begin{enumerate}
\renewcommand{\theenumi}{\alph{enumi}}
\renewcommand{\labelenumi}{\theenumi)}
\setcounter{enumi}{1}
\item Note that $(\ast)$ also commutes with direct sums indexed by the
  $I\in𝒰$, as is readily verified. In particular, if for every set $I$ we
  denote by $I_X$ the constant $X$-scheme defined by $I$ (\SGA 3 I 1.8), the
  direct sum of $I$ copies of $X$, then the associated sheaf is simply the
  constant sheaf $I_{X_{\text{ét}}}$ (the direct sum of $I$ copies of the
  final sheaf on $X_{\text{ét}}$). Since the functor $I\mapsto I_X$ also
  commutes with finite $\varprojlim$, we see that it takes a group to a
  group, a commutative group to a commutative group, \etc If $G$ is an
  ordinary commutative group, we shall simply write $H^i(X,G)$ in place of
  $H^i(X,G_X)$. Suppose, for example, that $G$ is a finite group. Then $G_X$
  is finite, hence affine over $X$, and thus, using the final remark of a),
  we obtain an interpretation of $H^1(X,G)$ as the set \emph{of classes of
  principal Galois coverings with group $G$} (\SGA 1 V 2.7). When $X$ is
  connected and endowed with a geometric point $a$, the ``fundamental
  group'' $\pi_1(X,a)$ (\SGA 1 V 7) gives the canonical isomorphism
\begin{equation*}
H^1(X, G) \simeq \Hom(\pi_1(X, a), G){\quoi},\tag{${\ast}{\ast}$}
\end{equation*}
(where $G$ has been assumed commutative).

One may say that in passing from Zariski cohomology to the étale topology,
``we have done what was needed'' to obtain the ``right'' $H^1$ (which occurs
in the second\pageoriginale member of $({\ast}{\ast})$) for a finite constant
coefficient group $G$. \emph{It is a remarkable fact, to be proved in the
remainder of this seminar, that this also suffices to obtain the ``right''
$H^i(X,G)$ for every torsion coefficient group (at least if $G$ is prime to
the residual characteristics of $X$)}.

\item Let $F$ be an $\Oo_X$-module on $X$, for the Zariski topology. Then
  (with the notation of \SGA 3 I 4.6) define a functor
$$
W(F) : (\Sch)_{/X}^\circ \to (\Ens)
$$
by
$$
W(F)(X') = \Gamma (X', F \otimes_{\Oo_{X}} \Oo_{X'}){\quoi}.
$$
This functor is again a sheaf for the étale topology (and even for the fpqc
topology), as follows again from \SGA 1 VIII 1.6 and \SGA 3 IV 6.3.1. A
fortiori, the restriction of $W(F)$ to $X_{\text{ét}}$ is again a sheaf,
which we also denote by $W(F)$. By definition, therefore,
$H^{\circ}(X,W(F))=\Gamma(F)=H^{\circ}(X,F)$. But there is a stronger result
if $F$ is quasi-coherent; \cf \hyperref[VII.prop4.3]{4.3}.
\end{enumerate}
\end{enonce*}

% END INCLUDED SGA4_English_current_buildable_source_20260906/source/07/drafts/SGA4_VII_LINES_43_396_English.texfrag


% BEGIN INCLUDED SGA4_English_current_buildable_source_20260906/source/07/drafts/SGA4_VII_LINES_397_665_English.texfrag
\section{Generators of the étale topos. Cohomology of an
  $\protect\varinjlim$ of sheaves}\etiquette[3]{VII.sec3}

\begin{enonce*}{Proposition 3.1}\etiquette[3.1]{VII.prop3.1}
Let $X$ be a scheme, and let $C$ be a full subcategory of
$X_{\text{\rm ét}}$ such that, for every affine scheme $X'$ étale
over $X$, $C$ contains an object isomorphic to $X'$. Then $C$ is a
``family of topological generators'' of the site $X_{\text{\rm ét}}$,
and hence a family of generators of the topos
$\widetilde{X_{\text{\rm ét}}}$. The restriction functor
$\widetilde{X_{\text{\rm ét}}}\to\widetilde{C}$ is an equivalence of
categories (where $C$ is equipped with the topology induced by that
of $X_{\text{\rm ét}}$).
\end{enonce*}

Immediate\pageoriginale from the ``comparison lemma.''

\begin{enonce*}{Corollary 3.2}\etiquette[3.2]{VII.cor3.2}
Suppose that $X$ is quasi-separated. Call the full subcategory $C$ of
$X_{\text{\rm ét}}$ formed by the schemes étale over $X$ that are of
finite presentation over $X$, equipped with the topology induced by
$X_{\text{\rm ét}}$, the restricted étale site of $X$. Then:
\begin{enumerate}
\renewcommand{\theenumi}{\roman{enumi}}
\renewcommand{\labelenumi}{(\theenumi)}
\item $C$ is stable under fiber products and is a site of finite type
  if $X$ is quasi-compact.

\item The restriction functor
  $\widetilde{X_{\text{\rm ét}}}\to\widetilde{C}$ is an equivalence of
  categories.
\end{enumerate}
\end{enonce*}

Assertion (i) is immediate, and (ii) follows from
\hyperref[VII.prop3.1]{3.1}, since $C$ satisfies the condition of
\hyperref[VII.prop3.1]{3.1} by the quasi-separatedness of $X$. Indeed,
this implies that, if $X'$ is quasi-compact, for example affine, then
it is quasi-compact \emph{over} $X$ (\EGA IV 1.2.4, with
$Z=\Spec\ZZ$), and hence is of finite presentation over $X$ if $X'$
is locally of finite type over $X$.

\begin{enonce*}{Proposition 3.3}\etiquette[3.3]{VII.prop3.3}
Suppose that $X$ is quasi-compact and quasi-separated. Then the
functors $H^q(X_{\text{\rm ét}},F)$ on the category of abelian sheaves
on $X_{\text{\rm ét}}$ commute with $\varinjlim$.
\end{enonce*}

Indeed, by \hyperref[VII.cor3.2]{3.2} (ii), one may replace
$X_{\text{ét}}$ by $C$. Since $X$ is quasi-compact, $X\in\Ob C$, and
the conclusion follows from \hyperref[VII.cor3.2]{3.2} (i) and
\hyperrefext[VI][VI.subsubsec6.1.2]{6.1.2} (3).

\section{Comparison with other topologies}\etiquette[4]{VII.sec4}

\setcounter{subsection}{-1}
\subsection{}\etiquette[4.0]{VII.subsec4.0}

First note that the examples of sheaves on $X_{\text{ét}}$ considered
in ${\rm n^{\circ}}$~2 are in fact naturally \emph{restrictions} of
sheaves\pageoriginale defined on $(\Sch)_{/X}$ equipped with its étale
topology (or even with the fpqc topology). More generally, let
$$
u^{\ast} : X_{\text{ét}} \to (\Sch)_{/X}
$$
be the inclusion functor. It is continuous (III) and commutes with
finite $\varprojlim$, and hence defines a \emph{morphism of sites}
$$
u : (\Sch)_{/X} \to X_{\text{ét}}{\quoi},
$$
whence a functor
$$
u_{\ast} : \widetilde{(\Sch)_{/X}} \to \widetilde{X_{\text{ét}}}
$$
on the associated categories of sheaves:
$$
u_{\ast}(F)=F\circ u{\quoi},
$$
and a left adjoint\footnote{Strictly speaking, since $(\Sch)_{/X}$ is
a $\Uu$-category but \emph{not} a $\Uu$-site, one cannot invoke III
for the existence of $u^{\ast}$. Nevertheless, one constructs
$u^{\ast}$ easily by
$$
u^{\ast} (F) (X') = (p_{X'})^{\ast}_{\text{ét}} (F) (X')\quad (X'\in \Ob \Sch_{/X}){\quoi},
$$
where $p_{X'}$ denotes the structural morphism $X'\to X$. Moreover,
to give meaning to \hyperref[VII.prop4.1]{4.1} (ii) and justify the
indicated proof of \hyperref[VII.prop4.1]{4.1}, one must introduce a
universe $\Vv$ such that $\Uu\in\Vv$, and regard the sheaves occurring
in the formulas of \hyperref[VII.prop4.1]{4.1} (i) as sheaves with
values in $\Vv\text{-}(\Ens)$. Alternatively, if one wishes to work
only with $\Uu$-sites, one may replace $(\Sch)_{/X}$ by a full
subcategory of $(\Sch)_{/X}$ stable under finite $\varprojlim$,
containing the $C$ of \hyperref[VII.prop3.1]{3.1}, and
$\Uu$-small.} to the latter:
$$
u^{\ast} : \widetilde{X_{\text{ét}}} \to \widetilde{\Sch_{/X}}.
$$
With this understood:

\begin{enonce*}{Proposition 4.1}\etiquette[4.1]{VII.prop4.1}
\begin{enumerate}
\renewcommand{\theenumi}{\roman{enumi}}
\renewcommand{\labelenumi}{(\theenumi)}
\item The functor $u^{\ast}$ is fully faithful; hence, for every sheaf
  $G$ on $X_{\text{\rm ét}}$, the canonical homomorphism
  $G\to u_{\ast}u^{\ast}(G)$ is an isomorphism.

\item Let $G$ (\resp $F$) be an abelian sheaf on
  $X_{\text{\rm ét}}$ (\resp $(\Sch)_{/X}$). Then the following
  canonical homomorphisms (defined, for example, as edge
  homomorphisms of the corresponding Leray spectral sequences) are
  isomorphisms:\pageoriginale
\begin{gather*}
H^{\ast} (X_{\text{\rm ét}}, F \circ u) \xrightarrow{\sim} H^{\ast}((\Sch)_{/X}, F)\\
H^{\ast} (X_{\text{\rm ét}}, G) \xrightarrow{\sim} H^{\ast}((\Sch)_{/X}, u^{\ast}G){\quoi}.
\end{gather*}
\end{enumerate}

\end{enonce*}

Indeed, the functor $u$ is a comorphism, and one can apply (III).

Thus (i) shows that, for the study of the cohomology of sheaves, it is
essentially equivalent to work with the ``small'' étale site
$X_{\text{ét}}$ or with the ``big'' étale site $(\Sch)_{/X}$.

\setcounter{subsection}{1}
\subsection{}\etiquette[4.2]{VII.subsec4.2}

One should also introduce on $(\Sch)$ (and hence on the
$(\Sch)_{/X}$) various topologies other than the étale topology, the
most useful of which are defined in \SGA 3 IV 6.3. The coarsest among
them is the \emph{Zariski topology} $(\Zar)$, defined by the
pretopology whose covering families are the surjective families of
open immersions; it is coarser than the étale topology. The finest of
these topologies is the \emph{``faithfully flat quasi-compact''
topology}, abbreviated (fpqc), which is the coarsest topology for
which the covering families in the Zariski sense, as well as
faithfully flat quasi-compact morphisms, are covering; the fpqc
topology is finer than the étale topology. As already observed, the
various examples of sheaves on $(\Sch)_{/X}$ considered in
${\rm n^{\circ}}$~2 are in fact already sheaves for the fpqc topology.

\subsubsection{}\etiquette[4.2.1]{VII.subsubsec4.2.1}

One must bear in mind, however, that for an abelian sheaf $F$ on
$(\Sch)_{/X}$ for the fpqc topology (or for a topology such as fppf,
considered in \SGA 3), the cohomology groups of $F$ for the fpqc
topology (\resp\ldots) are not always isomorphic to the cohomology
groups for\pageoriginale the étale topology. This is so even if $X$ is the spectrum
of a field $k$ and $F$ is representable by an algebraic group over
$k$, even for $H^1$. In general, one can show that the étale topology
gives the ``correct'' cohomology groups for coefficient groups that
are \emph{étale}, or more generally \emph{smooth}, group schemes over
$X$\footnote{\Cf A. Grothendieck, Le groupe de Brauer III, Th.~11.7,
in ``Dix exposés sur la topologie des schémas,'' North-Holland
Publishing Co.}, but this is no longer true for group schemes such as
\emph{radicial} groups over $S$, for which one must replace the étale
topology, still too coarse, by the fpqc or fppf topology.

\subsubsection{}\etiquette[4.2.2]{VII.subsubsec4.2.2}

As an example of the relations between cohomologies for different
topologies, consider the Zariski and étale topologies. Denote by
$X_{\Zar}$ the site of Zariski open subsets of $X$, so that there is a
canonical inclusion functor
$$
u^{\ast} : X_{\Zar} \to X_{\text{ét}}
$$
which defines a morphism of sites
$$
f : X_{\text{ét}} \to X_{\Zar}{\quoi},
$$
whence corresponding direct image functors
$$
f_{\ast} : \widetilde{X_{\text{ét}}} \to \widetilde{X_{\Zar}}\quad
(f_{\ast} (F) = F \circ u){\quoi},
$$
and inverse image functors
$$
f^{\ast} : \widetilde{X_{\Zar}} \to \widetilde{X_{\text{ét}}}{\quoi},
$$
adjoint to one another. Geometrically, the pair $(f_{\ast},f^{\ast})$
should be interpreted as a morphism of topoi
$$
f : \widetilde{X_{\text{ét}}} \to \widetilde{X_{\Zar}}{\quoi}.
$$
One obtains a homomorphism of cohomological functors
$$
H^{\ast} (X_{\Zar}, f_{\ast} (F)) \to H^{\ast} (X_{\text{ét}}, F)
$$
and\pageoriginale a Leray spectral sequence
$$
H^{\ast} (X_{\text{ét}}, F) \Longleftarrow H^{p} (X_{\Zar}, R^q
f_{\ast} (F)){\quoi},
$$
where $F$ is an abelian sheaf on $X_{\text{ét}}$. This spectral
sequence summarizes the general relations between étale and Zariski
cohomology. Of course, for coefficient sheaves $F$ such as constant
coefficient sheaves, this spectral sequence is in general far from
trivial, \ie in general $R^qf_{\ast}(F)\neq0$ (notably, take the case
where $X$ is the spectrum of a field). Nevertheless:

\begin{enonce*}{Proposition 4.3}\etiquette[4.3]{VII.prop4.3}
Let $F$ be a sheaf of varying modules on the scheme $X$ (a sheaf for
the Zariski topology), whence a sheaf $F_{\text{\rm ét}}$ on
$X_{\text{\rm ét}}$ (\cf 2, c) and a homomorphism of cohomological
functors
$$
H^{\ast} (X, F) \to H^{\ast} (X_{\text{\rm ét}}, F_{\text{\rm ét}}){\quoi}.
$$
When $F$ is quasi-coherent, the preceding homomorphism is an
isomorphism.
\end{enonce*}

Indeed, the first member is precisely
$H^{\ast}(X_{\Zar},f_{\ast}(F_{\text{ét}}))$, and with the preceding
notation it is enough to prove that
$$
R^q f_{\ast} (F_{\text{ét}}) = 0\quad \mbox{for}\quad q > 0{\quoi}.
$$
By the method for computing the $R^qf_{\ast}$ and the fact that the
affine open subsets form a basis for the topology of $X$, this follows
immediately from:

\begin{enonce*}{Corollary 4.4}\etiquette[4.4]{VII.cor4.4}
If $X$ is affine, then
$$
H^{q} (X_{\text{\rm ét}}, F_{\text{\rm ét}}) = 0\quad \mbox{for}\quad q > 0{\quoi}.
$$
\end{enonce*}

To see this, let $C=X_{\text{ét}}$, and let $C'$ be the full
subcategory formed by the \emph{affine} $X'$. Then, by
\hyperref[VII.prop3.1]{3.1},
$$
H^{q} (C, F_{\text{ét}}) \simeq H^{q} (C', F | C'){\quoi}.
$$

It is therefore\pageoriginale enough to prove that $F|C'$ is
$C'$-acyclic (\hyperrefext[V][V.def4.1]{4.1}), or equivalently that it
satisfies the condition of \hyperrefext[V][V.prop4.3]{4.3}, \ie that,
for every $X'\in\Ob C'$ and every covering family
$ℛ=(X'_i\to X')_i$ in $C'$, one has $H^q(ℛ,F)=0$ for $q>0$. One may
suppose the family $(X'_i)_i$ finite and then, after replacing the
$X'_i$ by their coproduct, that the covering family consists of a
single covering morphism $X'_i\to X'$, \ie an (étale and) surjective
morphism. We are thus reduced to proving that, if $f:X\to Y$ is a
surjective étale morphism of affine schemes and $F$ a quasi-coherent
module on $Y$, then
\begin{equation*}
H^q (X /Y, F) = 0\quad \mbox{for}\quad q > 0{\quoi}.\tag{$\ast$}
\end{equation*}
This is proved in TDTE I, B, \hyperref[VII.subsec1.1]{1.1} (in FGA,
\cf reference [4] of IX).

\begin{enonce*}[remark]{Remarks 4.5}\etiquette[4.5]{VII.rem4.5}
In fact, by \loccit, for $(\ast)$ it is enough that $X\to Y$ be
surjective and flat (rather than étale). The preceding proof then
shows that one still has isomorphisms analogous to that of
\hyperref[VII.prop4.3]{4.3}, with the étale site $X_{\text{ét}}$
replaced by $(\Sch)_{/X}$ equipped with any of the topologies finer
than the Zariski topology considered in \SGA 3 IV 6.3, for example
the fpqc topology.
\end{enonce*}

% END INCLUDED SGA4_English_current_buildable_source_20260906/source/07/drafts/SGA4_VII_LINES_397_665_English.texfrag


% BEGIN INCLUDED SGA4_English_current_buildable_source_20260906/source/07/drafts/SGA4_VII_LINES_666_1032_English.texfrag
\section{Cohomology of a projective limit of schemes}\etiquette[5]{VII.sec5}

\subsection{}\etiquette[5.1]{VII.subsec5.1}

Let $I$ be an increasing filtered preordered set,
$$
\mathfrak{X} = (X_i)_{i \in I}
$$
a projective system of schemes, the transition morphisms $u_{ji} :
X_i \to X_j$ $(i \geq j)$ being \emph{affine}. Recall (\EGA IV
8\ldots) that under these conditions,\pageoriginale the projective limit
$$
X  = \varprojlim X_i
$$
exists in the category of schemes (and even in the category of locally
ringed spaces) and may be constructed as follows: one chooses
$i_{0} \in I$, so that for $i \geq i_0$ one has an
$X_{i_0}$-isomorphism (\EGA II 1.2 and 1.3)
$$
X_i = \Spec 𝒜_{i}{\quoi},
$$
where $𝒜_i$ is a quasi-coherent Algebra on $X_{i_0}$, the
$𝒜_i$ $(i \geq i_0)$ thus forming an inductive system of
such Algebras. Setting
$$
𝒜 = \varinjlim_{i} 𝒜_i{\quoi},
$$
one may take
$$
X = \Spec 𝒜{\quoi}.
$$
Moreover, denoting by esp $S$ the underlying topological space of a
scheme $S$, one proves (\EGA IV 8\ldots) that the canonical map
$$
\esp X \to \varprojlim \esp (X_i)
$$
is a \emph{homeomorphism} and that the canonical homomorphism of
sheaves of rings
$$
\varinjlim_{i} \gamma^{-1}_{i} (\Oo_{X_i}) \to \Oo_X
$$
is bijective, where $\gamma_i : X \to X_i$ is the canonical continuous
map.

\subsection{}\etiquette[5.2]{VII.subsec5.2}

Informally, the content of \EGA IV 8, 9 may be summarized by saying
that, \emph{when $X_{i_0}$ is quasi-compact and quasi-separated},
then ``every datum of a schematic nature on $X$, of finite presentation
over $X$'', is equivalent\pageoriginale to a datum of the same nature
on one of the $X_i$, ``for sufficiently large $i$''. Thus one proves
(\EGA IV 8 $\ldots$):

\begin{enumerate}
\renewcommand{\theenumi}{\alph{enumi}}
\renewcommand{\labelenumi}{(\theenumi)}
\item if $i_1 \in I$ and $X'_{i_1}$, $X''_{i_1}$ are two
schemes of finite presentation over $X_{i_1}$, then, setting for every
$i \geq i_1$
$$
X'_{i} = \fprod{X'_{i_1}}{X_i}{X_{i_1}},\quad X' = \fprod{X'_{i_1}}{X}{X_{i_1}}{\quoi},
$$
and defining $X''_i$, $X''$ in the same way, the canonical map
$$
\varinjlim_{i \geq i_1} \Hom_{X_i} (X'_i, X''_i) \to
\Hom_X (X', X'')
$$
is bijective.

\item For every scheme $X'$ of finite presentation over $X$, there
exist an index $i_1 \in I$, a scheme $X'_{i_1}$ of finite presentation
over $X_{i_1}$, and an $X$-isomorphism
$$
X'\simeq \fprod{X'_{i_1}}{X}{X_{i_1}}{\quoi}.
$$
\end{enumerate}

\subsection{}\etiquette[5.3]{VII.subsec5.3}

The preceding result, and analogous results contained in EGA IV 8,
may be expressed in the language of the $\varinjlim$ of fibered
categories introduced in VI. For this purpose, consider the category
$\Ff$ of morphisms $f : T \to S$ of finite presentation of schemes,
and the ``target functor''
$$
\pi : \Ff \to (\Sch){\quoi},
$$
which is plainly a \emph{fibered functor}, its fiber categories being,
moreover, equivalent to $\Uu$-small categories. Consider the inverse
image of the fibered category $\Ff/(\Sch)$ by the functor
$$
i \mapsto X_i : \cat (I) \to (\Sch)
$$
(where\pageoriginale $\cat(I)$ is the category associated with $I$, with
$\Hom(i, j) \neq ∅$ if and only if $i \geq j$), giving a fibered
category
$$
\pi_{\mathfrak{X}} : \Ff_{\mathfrak{X}} \to \cat(I){\quoi},
$$
whose fiber at each $i \in I$ is canonically isomorphic to the full
subcategory $\Ff_{X_i}$ of $(\Sch)_{/X_i}$ consisting of schemes of
finite presentation over $X_i$, the base-change morphism
$\Ff_i \to \Ff_{j}$ relative to $j \geq i$ being none other than
$X'_i \mapsto X'_j = \fprod{X'_i}{X_j}{X_i}$. This being so, by
associating to every $X'_i$ of finite presentation over $X_i$ its
inverse image on $X$
$$
\varphi(X'_i) = \fprod{X'_i}{X}{X_i}{\quoi},
$$
one obtains a natural functor
$$
\varphi : \Ff_{\mathfrak{X}} \to \Ff_{X}{\quoi},
$$
which plainly takes every cartesian morphism of
$\Ff_{\mathfrak{X}}$ to an isomorphism, and hence, by definition,
factors uniquely through a canonical functor
\begin{equation*}
\psi : \varinjlim_{\Ff_X /\Cat I}
\Ff_{\mathfrak{X}} \to \Ff_{X}{\quoi}.\tag{$\ast$}
\end{equation*}
In view of the description of the first member, the result of \EGA IV
8 recalled above may then be stated by saying that $(\ast)$ is an
\emph{equivalence of categories}.

\subsection{}\etiquette[5.4]{VII.subsec5.4}

Statement a) of \hyperref[VII.subsec5.2]{5.2}. above may be refined in various ways,
by introducing some class $(\Mm)$ of morphisms of schemes, stable under
base change, and stating that for a given $X_{i_1}$-morphism\pageoriginale
$u_{i_1} : X'_{i_1} \to X''_{i_1}$, the corresponding morphism
$u': X'\to X''$ is $\Ee (\Mm)$ if and only if there exists
$i\geq i_{1}$ such that the $X_{i}$-morphism
$u_{i}:X'_{i}\to X''_{i}$ is $\Ee (\Mm)$. The statement obtained in
this way is true, for example, when $(\Mm)$ is one of the following
classes of arrows: proper morphisms (respectively projective, \resp
quasi-projective, \resp affine, \resp quasi-affine, \resp finite, \resp
quasi-finite, \resp radicial, \resp surjective, \resp flat, \resp
smooth, \resp étale, \resp unramified). The reader will find the
corresponding statements in \EGA IV, §8 (for the first nine), §11 (for
the flat case), and §17 (for the last three cases).

\subsection{}\etiquette[5.5]{VII.subsec5.5}

Now replace $\Ff \to (\Sch)$ by the fibered subcategory $\Gg$
consisting of \emph{étale morphisms} of finite presentation, whose
fiber category $\Gg_{S}$, for every scheme $S$, is none other than the
category of schemes étale and of finite presentation over $S$, and
similarly form the fibered category $\Gg_{\mathfrak{X}}$ over
$\cat(I)$, whose fiber category at every $i \in I$ is the category
$\Gg_{X_i}$. By \hyperref[VII.cor3.2]{3.2} (ii), the étale topoi
$\widetilde{X_{i_{\text{ét}}}}$, $\widetilde{X_{\text{ét}}}$ are also
canonically equivalent to $\widetilde{\Gg_{X_{i}}}$,
$\widetilde{\Gg_{X}}$, (where each $\Gg_S$ is endowed with the
topology induced by $S_{\text{ét}}$). Moreover, by
\hyperref[VII.cor3.2]{3.2} (i), each $\Gg_{X_i}$ is a site of finite
type, and is moreover equivalent to a $\Uu$-small site. In view of the
topologies on the $\Gg_{X_i}$, $\Gg$ becomes a \emph{fibered site}
over $\cat(I)$. This being so, we may state the following

\begin{enonce*}{Lemma 5.6}\etiquette[5.6]{VII.lem5.6}
The canonical functor
$$
\varinjlim_{\Gg_{\mathfrak{X}}/\cat(I)}
\Gg_{\mathfrak{X}} \to \Gg_X
$$
defines\pageoriginale an equivalence, respecting the topologies, between the
$\varinjlim$ site of the restricted étale sites
(\hyperref[VII.cor3.2]{3.2}) of the $X_i$ and the restricted étale site
of $X = \varprojlim_i X_i$.
\end{enonce*}

The fact that one obtains an equivalence of categories is one of the
statements recalled in \hyperref[VII.subsec5.4]{5.4} (the one in which
$(\Mm)$ is the class of étale morphisms in $(\Sch)$). The assertion
concerning the topologies is obtained in the same way by taking
$(\Mm)=$ the class of surjective morphisms in $(\Sch)$.

The preceding result allows us to apply the results of VI to the
present situation. In particular, we obtain

\begin{enonce*}{Theorem 5.7}\etiquette[5.7]{VII.thm5.7}
Let $\mathfrak{X} = (X_i)_{i \in I}$ be a projective system of
schemes, with $I$ an increasing filtered preordered set. Suppose that
the transition morphisms $u_{ij} : X_j \to X_i$ are affine (so that the
scheme $X = \varprojlim X_i$ is defined
(\hyperref[VII.subsec5.1]{5.1})), and that the $X_i$ are quasi-compact
and quasi-separated. Consider the fibered site $\Gg$ over $\cat(I)$
described in \hyperref[VII.subsec5.5]{5.5}, whose fiber $\Gg_i$ at
$i \in I$ is canonically isomorphic to the ``restricted étale site''
of $X_i$ (consisting of schemes étale and of finite presentation over
$X_i$). Let $F$ be an abelian sheaf on $\Gg$ (\ie $(3)$ a functor
$\Gg^{0} \to (\Ab)$ whose restriction $F_i$ to each $\Gg_i$ is a
sheaf, \ie a sheaf on $X_i$). Let $F_{\infty} =
\varinjlim_{i} u^{\ast}_i (F_i)$ be the induced sheaf on $X$, where
$u_i : X \to X_i$ is the canonical morphism. Then the canonical
homomorphisms
$$
\varinjlim_i H^n (X_i, F_i) \to H^n (X, F_{\infty})
$$
are isomorphisms.
\end{enonce*}

In\pageoriginale view of \hyperref[VII.lem5.6]{5.6}, this follows by
combining \hyperrefext[VI][VI.thm8.7.3]{8.7.3} and
\hyperrefext[VI][VI.prop8.5.2]{8.5.2}. We shall use
\hyperref[VII.thm5.7]{5.7} chiefly in the following particular case,
which is also made explicit in

\begin{enonce*}{Corollary 5.8}\etiquette[5.8]{VII.cor5.8}
With the preceding hypotheses and notation for $(X_i)_{i \in I}$,
$X$, $u_{ij}$, $u_i$, let $i_0 \in I$, and let $F_{i_0}$ be an
abelian sheaf on $X_{i_0}$. For every $i \geq i_0$, let
$F_i = u_{i_0 i} {}^{\ast}(F_{i_0})$, and let moreover
$F_{\infty} = u^{\ast}_{i_0} (F_{i_0})$. Under these conditions, the
canonical homomorphisms
$$
\varinjlim_i H^n (X_i, F_i) \to H^n (X, F_{\infty})
$$
are isomorphisms.
\end{enonce*}

Here is, however, a sometimes useful case that follows from
\hyperref[VII.thm5.7]{5.7} and not from \hyperref[VII.cor5.8]{5.8}:

\begin{enonce*}{Corollary 5.9}\etiquette[5.9]{VII.cor5.9}
With the hypotheses and notation of \hyperref[VII.thm5.7]{5.7} for
$(X_i)_{i \in I}$, $X$, $u_{ij}$, $u_i$, let $G_{i_0}$ be a
commutative group scheme locally of finite presentation over
$X_{i_0}$ (where $i_0 \in I$ is given). For every $i \geq i_0$, let
$G_i = \fprod{G_{i_0}}{X_i}{X_{i_0}}$, and similarly let
$G_{\infty} = \fprod{G_{i_0}}{X}{X_{i_0}}$. Then the canonical
homomorphisms
$$
\varinjlim_i H^n (X_i, G_i) \to H^n (X, G_{\infty})
$$
are isomorphisms.
\end{enonce*}

Indeed, we know that $G_{i_0}$ defines a sheaf on
$(\Sch)_{/X_{i_0}}$ (\cf 2 (a)), and, supposing that $i_0$ is an
initial object of $I$ (which we may do), hence obtaining a canonical
functor $\Gg \to (\Sch)_{/X_{i_0}}$, composition gives a contravariant
functor $F$ on $\Gg$, whose restriction to $\Gg_i$ is\pageoriginale
canonically isomorphic to the sheaf on $X_i$ defined by $G_i$. In
particular, $F$ is a sheaf on $\Gg$, and we may apply
\hyperref[VII.thm5.7]{5.7} to it. The hypothesis that $G_{i_0}$ is
locally of finite presentation over $X_{i_0}$ serves to ensure that the
sheaf $F_{\infty}$ of \hyperref[VII.thm5.7]{5.7} is isomorphic, by the natural homomorphism
$F_{\infty} \to \faisc(G_{\infty})$, to the sheaf
$\faisc(G_{\infty})$ defined by $G_{\infty}$ (as follows immediately
from \EGA IV 8.8.2 (i), by considering the values of the sheaves in
question on the objects of the restricted étale site of $X$
and applying \hyperref[VII.lem5.6]{5.6}). This completes the proof of
\hyperref[VII.cor5.9]{5.9}.

\begin{enonce*}{Corollary 5.10}\etiquette[5.10]{VII.cor5.10}
With the hypotheses and notation of \hyperref[VII.thm5.7]{5.7} for
$(X_i)_{i \in I}$, $X$, $u_{ij}$, $u_i$, let $F$ be a sheaf of
commutative groups on $X$; then the canonical homomorphisms
$$
\varinjlim_{i} H^n (X_i, u_{i\ast} (F)) \to H^n (X, F)
$$
are isomorphisms.
\end{enonce*}

The statements \hyperref[VII.thm5.7]{5.7} through
\hyperref[VII.cor5.10]{5.10} generalize to corresponding statements
for the $\Ext^i$ of sheaves of Modules, by
\hyperrefext[VI][VI.cor8.7.9]{8.7.9}; we shall not repeat them in the
present particular case. On the other hand, for later reference we
shall make explicit the variants of \hyperref[VII.cor5.8]{5.8} and \hyperref[VII.cor5.9]{5.9}
in terms of the functors $R^n f_{\ast}$.

\begin{enonce*}{Corollary 5.11}\etiquette[5.11]{VII.cor5.11}
Let $I$ be an increasing filtered preordered set, let $(X_i)_{i \in I}$
and $(Y_i)_{i \in I}$ be two projective systems of schemes with affine
transition morphisms, $X = \varprojlim_{i} X_i$,
$Y = \varprojlim_{i} Y_i$, $(f_i)_{i \in I}$ a projective system of
quasi-compact and quasi-separated morphisms $f_i : X_i \to Y_i$,
$i_0 \in I$, and $F_{i_0}$ a sheaf on $X_{i_0}$. For every
$i \geq i_0$, let $F_i$ be the sheaf\pageoriginale on $X_i$ that is
the inverse image of $F_{i_0}$. Similarly, let $f : X \to Y$ be
deduced from $(f_i)$ and let $F$ be the sheaf on $X$ that is the inverse
image of $F_{i_0}$. Under these conditions, the canonical homomorphism
$$
\varinjlim_{i} v^{\ast}_i (R^n f_{i\ast} (F_i)) \to R^n
f_{\ast} (F)
$$
is an isomorphism (where $v_i : Y \to Y_i$ is the canonical morphism).
\end{enonce*}

We may suppose that $i_0$ is an initial object of $I$, and, since the
question is local on $Y_{i_0}$, that $Y_{i_0}$ is affine, so that the
$Y_i$ and the $X_i$ are quasi-compact and quasi-separated. We are then
under the hypotheses for applying
\hyperrefext[VI][VI.thm8.7.3]{8.7.3} (where one reduces to
\hyperref[VII.cor5.8]{5.8}, in view of the computation of the
$R^{n} f_{\ast}$ as sheaves associated with presheaves). In the same
way one proves:

\begin{enonce*}{Corollary 5.12}\etiquette[5.12]{VII.cor5.12}
The same statement as \hyperref[VII.cor5.11]{5.11}, except that now
$F_{i_0}$ denotes a commutative group scheme locally of finite
presentation over $X_{i_0}$, and $F_i$, $F$ its inverse images (in the
sense of base change for schemes).
\end{enonce*}

Note that \hyperref[VII.cor5.9]{5.9} (\resp
\hyperref[VII.cor5.12]{5.12}) is not a particular case of
\hyperref[VII.cor5.8]{5.8} (\resp \hyperref[VII.cor5.11]{5.11}),
\cf \hyperref[VII.sec2]{2} a) remark \hyperref[VII.rem2.1]{2.1}.

\begin{enonce*}{Corollary 5.13}\etiquette[5.13]{VII.cor5.13}
With the notation of \hyperref[VII.cor5.11]{5.11}, let $F$ be a sheaf
of abelian groups on $X$. Then the canonical homomorphisms
$$
\varinjlim_i v_i^{\ast} (R^n f_{i\ast} (u_{i\ast} (F))) \to
R^n f_{\ast} (F)
$$
are isomorphisms (where $u_i : X \to X_i$ and $v_i : Y \to Y_i$ are
the canonical homomorphisms).
\end{enonce*}

The proof is essentially the same as before.

\begin{enonce*}[remark]{Remarks 5.14}\etiquette[5.14]{VII.rem5.14}
\begin{enumerate}
\renewcommand{\theenumi}{\alph{enumi}}
\renewcommand{\labelenumi}{\theenumi)}
\item The\pageoriginale preceding limit-passage results remain valid
for $H^{0}$ \resp $f_{\ast}$, for sheaves \emph{of sets} (instead of
abelian sheaves), and for $H^{1}$ \resp $R^1 f_{\ast}$ for sheaves of
groups (not necessarily commutative), where the noncommutative $H^1$
and $R^1 f_{\ast}$ are defined as usual in terms of torsors (or
principal homogeneous bundles) (\cf Giraud's thesis
\footnote{cited p.\pagehyperref[page7]{7} footnote $(\ast)$}). This can
be checked directly in the general context of VI. It is certainly
possible (and doubtless useful) also to give a variant for the
noncommutative $H^2$ of ``links'' studied by Giraud.

\item The limit-passage results for fibered sites developed in VI
assumed only that the base category was a filtered category, without it
being necessary to suppose that it was associated with a filtered
preordered set. The case of an arbitrary filtered index category is
also the natural framework for the statements developed in the present
${\rm n^{\circ}}$. If we have worked in an overly restrictive
framework, it was in order to be able to give correct references to
\EGA IV, where it is unfortunately assumed, in questions of passage to
the limit (from §8 onward), that the index category is defined in terms
of a filtered preordered set. We shall nevertheless assume henceforth
that all the results used from \EGA IV are valid for arbitrary filtered
index categories (the proofs given in \loccit being valid, essentially
without change, in this more general case). Thus we shall also use,
without further comment, the results of the present ${\rm n^{\circ}}$
in the case where $I$ is replaced by an arbitrary filtered index
category.
\end{enumerate}
\end{enonce*}

\ifx\danslelivre\undefined
\end{document}
\fi

% END INCLUDED SGA4_English_current_buildable_source_20260906/source/07/drafts/SGA4_VII_LINES_666_1032_English.texfrag


% END INCLUDED SGA4_English_current_buildable_source_20260906/source/07/07_English.tex


\makeatletter
\def\input@path{{../08/}}
\makeatother
% BEGIN INCLUDED SGA4_English_current_buildable_source_20260906/source/08/08_English.tex
\ifx\danslelivre\undefined
% Release engine: XeLaTeX (the LuaTeX compatibility layer breaks Xy-pic output).
\documentclass[10pt]{../fichiers-configuration-latex/smfart}
% BEGIN INCLUDED SGA4_English_current_buildable_source_20260906/source/fichiers-configuration-latex/paquets.tex
\usepackage{polyglossia}
%\usepackage[francais]{babel}
\usepackage{fontspec}
\usepackage[math-style=ISO,bold-style=upright,sans-style=upright]{unicode-math}
%\usepackage[mathletters]{ucs}
%\usepackage[utf8]{inputenc}
%\usepackage[T1]{fontenc}
\usepackage[unicode,hypertexnames=false,colorlinks=true,linkcolor=blue,citecolor=Sepia,urlcolor=magenta]{hyperref}
\setmainlanguage{french}
\setmainfont[Ligatures=TeX]{Linux Libertine O}
% variante possible : TeXGyrePagella ?
\setmathfont{XITS Math}

% maths
\usepackage{amsmath}
%\usepackage{amssymb}
%\usepackage{mathrsfs}
%\usepackage{stmaryrd}
%\usepackage{euscript}

% couleurs
\usepackage[usenames,dvipsnames]{xcolor}

% graphiques
\usepackage{graphicx}
\usepackage[all]{xy}
%\usepackage{tikz}
%\usetikzlibrary{matrix,arrows,calc}
%\tikzset{column sep=3em,row sep=2em,text height=1.5ex, text depth=0.25ex,>=angle 90}
%\tikzset{description/.style={fill=white,inner sep=1pt}}

% hyperliens, 'etiquettes
\usepackage{xr-hyper}
\usepackage{../fichiers-configuration-latex/smfthm}
%\ifx\cacheretiquettes\undefined
%\usepackage[inline]{showlabels} \else \fi % pour rendre visibles les 'etiquettes
%\usepackage{../fichiers-configuration-latex/makeidx} % pour faire un index
%\usepackage[pagebackref,unicode]{hyperref}  % pour voir o`u sont cit'ees les r'ef'erences
% hyperref is loaded once above with the cumulative-reader destination policy.
%\hypersetup{breaklinks=true} %permet le retour `a la ligne dans les liens trop longs. Incompatible avec pdflatex
%\renewcommand*{\backref}[1]{{\small [cit\'e page(s)~#1.]}}

\usepackage{xspace}
\usepackage{pifont} % pour expose XVII

% END INCLUDED SGA4_English_current_buildable_source_20260906/source/fichiers-configuration-latex/paquets.tex

% BEGIN INCLUDED SGA4_English_current_buildable_source_20260906/source/fichiers-configuration-latex/numerotation.tex
\renewcommand{\thefootnote}{\roman{footnote}} % pour mettre les notes en bas de page sous la forme (ii) et non (2), qui ressemble trop (?) a un exposant.

\newcounter{pageoriginal}
\newcommand{\pageoriginale}{\refstepcounter{pageoriginal}\marginpar{\footnotesize\textbf{\thepageoriginal}}\xspace}
\let\pageoriginaled\pageoriginale

\usepackage{ifthen}
\newboolean{original}
\def\sisi{\ifthenelse{\boolean{original}}}
\setboolean{original}{false}

\newcounter{cnote}
\def\nde#1{\ifthenelse{\boolean{original}}{}{\stepcounter{cnote}\def\thefootnote{\arabic{cnote}}\protect\footnote[\value{cnote}]{N.D.E.: #1}}}
\def\ndemark{\ifthenelse{\boolean{original}}{}{\stepcounter{cnote}\def\thefootnote{\arabic{cnote}}\protect\footnotemark[\value{cnote}]}}
\def\ndetext#1{\ifthenelse{\boolean{original}}{}{\def\thefootnote{\arabic{cnote}}\protect\footnotetext[\value{cnote}]{N.D.E.: #1}}}

% END INCLUDED SGA4_English_current_buildable_source_20260906/source/fichiers-configuration-latex/numerotation.tex

% BEGIN INCLUDED SGA4_English_current_buildable_source_20260906/source/fichiers-configuration-latex/macros.tex
\setmainlanguage{french}
\setotherlanguage{english}


%% Fixed, engine-independent release identifier.  The inherited dynamic Git
%% probe used Lua io.popen, exposed build-state-dependent output, and printed
%% a visible question mark outside a Git checkout.
\def\showgitstatus{English release r2, 2026-07-23}


% macros pour LaTeX
\renewcommand{\mathsf}{}

% Police différente pour U+2113 SCRIPT SMALL L
\setmathfont[range="2113]{latinmodern-math.otf}

% Give every inherited semantic label its own stable PDF destination instead
% of inheriting whichever automatic equation/page anchor happened to be live.
% The legacy label remains as a compatibility alias used by existing edges.
\makeatletter
\def\etiquette[#1]#2{%
  \Hy@raisedlink{\hyper@anchorstart{sga4.semantic.#2}\hyper@anchorend}%
  \def\@currentHref{sga4.semantic.#2}%
  \def\@currentlabel{#1}%
  \label{#2}%
}
\makeatother
\def\eqhyperref[#1]#2{(\hyperref[#1]{#2})}
\def\pagehyperref[#1]#2{(page \hyperref[#1]{#2})}
\def\hyperrefext[#1][#2]#3{\hyperref[#2]{#1 #3}}
\def\eqhyperrefext[#1][#2]#3{\hyperref[#2]{#1 #3}}     %cf. exposes 17
\def\neweqhyperrefext[#1][#2]#3{\hyperref[#2]{#1 #3}}  %cf. exposes 17

\def\shat#1{\protect\widehat{#1}}
\def\stilde#1{\protect\widetilde{#1}}
%%%%%%%%%%%%%%%%%%


\def\etcn{etc.\xspace}
\def\etco{etc...\xspace}
\def\etc{\sisi{\etco}{\etcn}}
\def\ldotsn{\sisi{...}{\etcn}}

\def\cf{{cf.}\kern.3em}
\def\Cf{{Cf.}\kern.3em}
\def\ie{{i.e.}\kern.3em}
\def\iev{{i.e.},\xspace}
\def\loccito{loc.\kern2pt cit.\xspace}
\def\loccitn{\textit{loc.\kern2pt cit.}\xspace}
\def\loccit{\sisi{\loccito}{\loccitn}}
\def\EGA{\textup{EGA}\kern.2em}
\def\SGA{\textup{SGA}\kern.3em}

\def\quoi{}

\def\quav{\quad,\quad}
\def\quaet{\quad \text{et}\quad}
\def\fibreproduct#1#2#3{#1 ×_{#3} #2}
\let\fprod\fibreproduct
\def\ccheck#1#2{\mbox{${\check{\check{\rm #1}}}^{\lower2pt\hbox{$\scriptstyle #2$}}$}}
\def\sheaf#1{\protect\mathcal{#1}}
\def\vieuxsheaf#1{{#1}} % cf. problème avec sheaf{a\Phi_i} par exemple...
\DeclareMathOperator{\Hom}{Hom}
\DeclareMathOperator{\can}{can}
\DeclareMathOperator{\Ind}{Ind}
\DeclareMathOperator{\ind}{ind}
%\DeclareMathOperator{\Hhomtop}{\mathcal{H}\!\mathit{omtop}}
\DeclareMathOperator{\Hhomtop}{\mathbf{Homtop}}
\DeclareMathOperator{\Point}{Point}
\DeclareMathOperator{\Pro}{Pro}
\DeclareMathOperator{\Top}{Top}
\DeclareMathOperator{\Mod}{Mod}
\DeclareMathOperator{\Cart}{Cart}
%\DeclareMathOperator{\Mmorsite}{\mathcal{M}\!\mathit{orsite}}
\DeclareMathOperator{\Mmorsite}{\mathbf{Morsite}}
\DeclareMathOperator{\loc}{loc}
\DeclareMathOperator{\Loc}{Loc}
\DeclareMathOperator{\Pic}{Pic}
\DeclareMathOperator{\Et}{Et}
\DeclareMathOperator{\card}{card}
\DeclareMathOperator{\Ens}{Ens}
\DeclareMathOperator{\h}{\underline{H}}
\DeclareMathOperator{\Fonct}{Fonct}
\DeclareMathOperator{\Ker}{Ker}
\DeclareMathOperator{\but}{but}
\DeclareMathOperator{\Ob}{Ob}
\DeclareMathOperator{\ob}{ob}
\DeclareMathOperator{\Lib}{Lib}
\DeclareMathOperator{\esj}{esj}
\DeclareMathOperator{\id}{id}
\DeclareMathOperator{\Id}{Id}
\DeclareMathOperator{\ad}{ad}
\DeclareMathOperator{\Iim}{Im}
\DeclareMathOperator{\pr}{pr}
\DeclareMathOperator{\Cat}{Cat}
\DeclareMathOperator{\car}{car}
\DeclareMathOperator{\SysInd}{SysInd}
\DeclareMathOperator{\Sup}{Sup}
\DeclareMathOperator{\Coker}{Coker}
\DeclareMathOperator{\Filt}{Filt}
\DeclareMathOperator{\acc}{acc}
\DeclareMathOperator{\sur}{sur}
\DeclareMathOperator{\cat}{cat}
\DeclareMathOperator{\Ord}{Ord}
\DeclareMathOperator{\Gr}{Gr}
\DeclareMathOperator{\Cov}{Cov}
\DeclareMathOperator{\Caf}{Caf}
\DeclareMathOperator{\coker}{coker}
\DeclareMathOperator{\coim}{coim}
\DeclareMathOperator{\im}{im}
\DeclareMathOperator{\Ouv}{Ouv}
\DeclareMathOperator{\diag}{diag}
\DeclareMathOperator{\Esp}{Esp}
\DeclareMathOperator{\TOP}{TOP}
\DeclareMathOperator{\Unif}{Unif}
\DeclareMathOperator{\Iind}{\mathbf{Ind}}
\DeclareMathOperator{\Hhomcart}{\mathbf{Homcart}}
\DeclareMathOperator{\Hhomtopess}{\mathbf{Homtopess}}
\DeclareMathOperator{\Ppoints}{\mathbf{Points}}
\DeclareMathOperator{\Ppoint}{\mathbf{Point}}
\DeclareMathOperator{\Ttors}{\mathbf{Tors}}
\DeclareMathOperator{\Ssex}{\mathbf{Sex}}
\DeclareMathOperator{\Ffib}{\mathbf{Fib}}
\DeclareMathOperator{\Vvois}{\mathbf{Vois}}
\DeclareMathOperator{\Ppointess}{\mathbf{Pointess}}
\DeclareMathOperator{\scrHom}{\mathbf{Hom}}
\DeclareMathOperator{\sheafPic}{\mathbf{Pic}}
\DeclareMathOperator{\sheafAut}{\mathbf{Aut}}
\DeclareMathOperator{\Hhomtoplocan}{\mathbf{Homtoplocan}}
\DeclareMathOperator{\Hhomtopan}{\mathbf{Homtopan}}
\DeclareMathOperator{\Pptgeom}{\mathbf{Ptgéom}}
\DeclareMathOperator{\Morsite}{Morsite}
\DeclareMathOperator{\Fib}{Fib}
\DeclareMathOperator{\Fais}{Fais}
\DeclareMathOperator{\kar}{kar}
\DeclareMathOperator{\Groupes}{Groupes}
\DeclareMathOperator{\Restr}{Restr}
\DeclareMathOperator{\Isom}{Isom}
\DeclareMathOperator{\Prol}{Prol}
\DeclareMathOperator{\Res}{Res}
\DeclareMathOperator{\Tors}{Tors}
\DeclareMathOperator{\esp}{esp}
\DeclareMathOperator{\es}{es}
\DeclareMathOperator{\od}{od}
\DeclareMathOperator{\gr}{gr}
\DeclareMathOperator{\sob}{sob}
\DeclareMathOperator{\dfn}{dfn}
\DeclareMathOperator{\Sch}{Sch}
\DeclareMathOperator{\ouv}{ouv}
\DeclareMathOperator{\cons}{cons}
\DeclareMathOperator{\Nerf}{Nerf}
\DeclareMathOperator{\Points}{Points}
\DeclareMathOperator{\Espsob}{Espsob}
\DeclareMathOperator{\Homtop}{Homtop}
\DeclareMathOperator{\Pt}{Pt}
\DeclareMathOperator{\Quot}{Quot}
\DeclareMathOperator{\glob}{glob}
\DeclareMathOperator{\Homfib}{Homfib}
\DeclareMathOperator{\Qouv}{Qouv}
\DeclareMathOperator{\Quotouv}{Quotouv}
\DeclareMathOperator{\ens}{ens}
\DeclareMathOperator{\ab}{ab}
\DeclareMathOperator{\an}{an}
\DeclareMathOperator{\fs}{fs}
\DeclareMathOperator{\red}{red}
\DeclareMathOperator{\Ab}{Ab}
\DeclareMathOperator{\cosq}{cosq}
\DeclareMathOperator{\Ann}{Ann}
\DeclareMathOperator{\Zar}{Zar}


%%%%%%%%%%%%
%%%%% SGA4-2
\DeclareMathOperator{\Sex}{Sex}
\DeclareMathOperator{\vois}{vois}
\DeclareMathOperator{\Vois}{Vois}
\DeclareMathOperator{\Ext}{Ext}
\DeclareMathOperator{\ext}{ext}
\DeclareMathOperator{\Fer}{Fer}
\DeclareMathOperator{\Add}{Add}
\DeclareMathOperator{\Cocart}{Cocart}
\DeclareMathOperator{\cosk}{cosk}
\DeclareMathOperator{\plat}{plat}

\DeclareMathOperator{\faisc}{faisc}
\DeclareMathOperator{\sk}{sk}
\DeclareMathOperator{\ch}{ch}

\def\HH{\mathbf{H}}
\def\qc{\mathrm{qc}}
\def\et{\mathrm{et}}
\def\coh{\mathrm{coh}}

\def\Hvr#1{%
\begin{array}[b]{l}
{\kern2pt\scriptstyle r}\\[-8pt]
{\kern1pt\scriptstyle \vee}\\[-5pt]
H^{#1}
\end{array}}

\def\Hvv#1{%
\begin{array}[b]{l}
\kern2pt\rotatebox{90}{${\scriptscriptstyle\ll}$}\\[-6pt]
H^{#1}
\end{array}}

\def\Hvi#1{%
\begin{array}[b]{l}
{\kern0.2pt\scriptscriptstyle \infty}\\[-8pt]
{\kern1pt\scriptstyle \vee}\\[-5pt]
H^{#1}
\end{array}}

\def\Hv#1{%
\begin{array}[b]{l}
{\kern1pt\scriptstyle \vee}\\[-5pt]
H^{#1}
\end{array}}

\DeclareMathOperator{\cart}{cart}
\DeclareMathOperator{\Ffaisc}{\mathbf{Faisc}}


\def\Ev#1{%
\begin{array}[b]{l}
{\kern1pt\scriptstyle \vee}\\[-5pt]
E'_{#1}
\end{array}}

\DeclareMathOperator{\Cont}{Cont}
\DeclareMathOperator{\Ccont}{\mathscr{C}\!\mathit{ont}}
\DeclareMathOperator{\Site}{Site}
\newcommand\MM{\mathbf{M}}
\def\limtop#1{{\displaystyle{\varprojlim_{#1}\!\mbox{\rm top }}}}
\def\Limtop{\mathop{\varprojlim\!\mbox{\rm top}}}
\DeclareMathOperator{\spec}{spec}
\def\ss{\scriptstyle{\mathscr{S}}}
\DeclareMathOperator{\Gal}{Gal}
\def\smallm{\scriptstyle{\mathscr{M}}}
\DeclareMathOperator{\End}{End}
\DeclareMathOperator{\ass}{ass}
\DeclareMathOperator{\zar}{zar}
\DeclareMathOperator{\Ogg}{Ogg}
\DeclareMathOperator{\res}{res}
\DeclareMathOperator{\tr}{tr}
\DeclareMathOperator{\Supp}{Supp}
\DeclareMathOperator{\codim}{codim}
\DeclareMathOperator{\rk}{rk}
\DeclareMathOperator{\cur}{cur}
\DeclareMathOperator{\cdqc}{cdqc}
\newcommand\DD{\mathbf{D}}
\DeclareMathOperator{\cl}{cl}
\def\JJ{\mathbf{J}}
\def\FF{\mathbf{F}}
\def\BB{\mathbf{B}}
\DeclareMathOperator{\Alget}{Alget}
\DeclareMathOperator{\Ass}{Ass}
\DeclareMathOperator{\supp}{supp}
\DeclareMathOperator{\rad}{rad}
\DeclareMathOperator{\prof}{prof}
\DeclareMathOperator{\cint}{int}
\DeclareMathOperator{\sic}{sic}
\DeclareMathOperator{\Tor}{Tor}
\DeclareMathOperator{\tors}{tors}
\DeclareMathOperator{\Sym}{Sym}
\DeclareMathOperator{\parf}{parf}
\DeclareMathOperator{\tor}{tor}
\DeclareMathOperator{\defn}{defn}
\DeclareMathOperator{\Div}{Div}
\DeclareMathOperator{\Pol}{Pol}
\DeclareMathOperator{\Fex}{Fex}
\DeclareMathOperator{\Lie}{Lie}
\DeclareMathOperator{\tot}{tot}
\DeclareMathOperator{\Aut}{Aut}
\DeclareMathOperator{\HOM}{HOM}
\DeclareMathOperator{\pch}{pch}
\DeclareMathOperator{\Ch}{Ch}
\DeclareMathOperator{\EXT}{EXT}
\DeclareMathOperator{\PIC}{PIC}
\DeclareMathOperator{\Rf}{Rf}
\DeclareMathOperator{\inj}{inj}
\newcommand\VV{\mathbf{V}}

\def\Xv{%
\begin{array}[b]{l}
{\kern1pt\scriptstyle \vee}\\[-5pt]
X
\end{array}}

%%%%%%%
\newcommand\TT{\mathbf{T}}
\newcommand\XX{\mathbf{X}}
\newcommand\EE{\mathbf{E}}
%%%%%%

\def\qedsymbol{\textup{C.Q.F.D.}}
\let\CQFD\qed
\let\cqfd\qed

%%%%%\bigdot
\font\bigsymb=cmsy10 at 4pt
\def\bigdot{{\kern1.2pt\raise 1.5pt\hbox{\bigsymb\char15}}}
%%%%%

\def\Lright{\mathop{\vcenter{\hbox{\hbox{$\longrightarrow$}\lower8pt \hbox{\kern-16pt\hbox{$\longrightarrow$}}}}}}

\def\Aa{{\mathcal A}}
\def\Bb{{\mathcal B}}
\def\Cc{{\mathcal C}}
\def\Dd{{\mathcal D}}
\def\Ee{{\mathcal E}}
\def\Ff{{\mathcal F}}
\def\Gg{{\mathcal G}}
\def\Hh{{\mathcal H}}
\def\Ii{{\mathcal I}}
\def\Jj{{\mathcal J}}
\def\Kk{{\mathcal K}}
\def\Ll{{\mathcal L}}
\def\Mm{{\mathcal M}}
\def\Nn{{\mathcal N}}
\def\Oo{{\mathcal O}}
\def\Pp{{\mathcal P}}
\def\Qq{{\mathcal Q}}
\def\Rr{{\mathcal R}}
\def\Ss{{\mathcal S}}
\def\Tt{{\mathcal T}}
\def\Uu{{\mathcal U}}
\def\Vv{{\mathcal V}}
\def\Ww{{\mathcal W}}
\def\Xx{{\mathcal X}}
\def\Yy{{\mathcal Y}}
\def\Zz{{\mathcal Z}}

\newcommand{\CC}{\mathbf{C}}
\newcommand{\GG}{\mathbf{G}}
\newcommand{\KK}{\mathbf{K}}
\newcommand{\LL}{\mathbf{L}}
\newcommand{\NN}{\mathbf{N}}
\newcommand{\PP}{\mathbf{P}}
\newcommand\QQ{\mathbf{Q}}
\newcommand\RR{\mathbf{R}}
\newcommand\ZZ{\mathbf{Z}}

\def\to{\mathchoice{\longrightarrow}{\rightarrow}{\rightarrow}{\rightarrow}}
\def\from{\mathchoice{\longleftarrow}{\leftarrow}{\leftarrow}{\leftarrow}}
\def\mto{\mathchoice{\longmapsto}{\mapsto}{\mapsto}{\mapsto}}
\def\mfrom{\mathrel{\from\joinrel\mapstochar\,}}
\def\hto{\mathrel{\lhook\joinrel\to}}
\def\hfrom{\mathrel{\from\joinrel\rhook}}

\def\implies{\mathchoice{\Longrightarrow}{\Rightarrow}{\Rightarrow}{\Rightarrow}}
\def\implied{\mathchoice{\Longleftarrow}{\Leftarrow}{\Leftarrow}{\Leftarrow}}
\def\iff{\mathchoice{\Longleftrightarrow}{\Leftrightarrow}{\Leftrightarrow}{\Leftrightarrow}}

\def\To#1{\mathchoice{\xrightarrow{\labelstyle~~#1~}}{\stackrel{#1}{\longrightarrow}}{}{}}
\def\Mto#1{\mathrel{\mapstochar\To{#1}}}
\def\From#1{\mathchoice{\xleftarrow{\labelstyle~~#1~}}{\stackrel{#1}{\longleftarrow}}{}{}}

\newcommand{\isomto}{\overset{\sim}{\longrightarrow}}
\newcommand{\isomfrom}{\overset{\sim}{\longleftarrow}}

\newcommand{\IFF}{\unskip~$\iff$~\ignorespaces}
\newcommand{\THEN}{\unskip~$\implies$~\ignorespaces}

\let\leq\leqslant
\let\geq\geqslant
\let\le\leqslant
\let\ge\geqslant

\def\cf{{cf.}\kern.3em}
\def\Cf{{Cf.}\kern.3em}
\def\ie{{i.e.}\kern.3em}
\def\iev{{i.e.},\xspace}
%\def\loccito{loc.\kern2pt cit.\xspace}
%\def\loccitn{\textit{loc.\kern2pt cit.}\xspace}
%\def\loccit{\sisi{\loccito}{\loccitn}}
\def\EGA{\textup{EGA}\kern.2em}
\def\SGA{\textup{SGA}\kern.3em}

\let\oldS\S
\def\S{\oldS\kern .15em}
\def\ptbl{.\kern .2em}
\def\resp{resp.\kern.3em}
\def\Exp{Exp.\kern.3em}

\let\bar\overline
\let\hat\widehat
\let\tilde\widetilde
\let\dpl\displaystyle
\let\ldots\dots
\let\emptyset\varnothing
\let\skpt\skippointrait
\let\epsilon\varepsilon

\def\sfrac#1#2{{#1}/{#2}}

\DeclareMathOperator{\cd}{cd}
%\DeclareMathOperator{\Hhom}{\mathcal{H}\!\mathit{om}}
\DeclareMathOperator{\Hhom}{\mathbf{Hom}}
%\DeclareMathOperator{\Eext}{\mathcal{E}\!\mathit{xt}}
\DeclareMathOperator{\Eext}{\mathbf{Ext}}
\DeclareMathOperator{\Spec}{Spec}
\DeclareMathOperator{\Tr}{Tr}

\let\emptyset=\varnothing

%%%%%%%%%%%%%%%
%%% expose XVII
%%%%%%%%%%%%%%%

\makeatletter
\let\old@optexpose\optexpose
\newcommand{\empty@optexpose}[1]{}
\newcommand{\paren@optexpose}[1]{\old@optexpose{#1}(}
\newcommand{\refnoexpose}[1]{\let\optexpose\empty@optexpose\ref{#1}\let\optexpose\old@optexpose}%\refnoexpose{X.cor5.2} engendre "5.2" au lieu de "X 5.2"
\newcommand{\neweqref}[1]{\let\optexpose\paren@optexpose\textup{\ref{#1})}\let\optexpose\old@optexpose}%\neweqref{XII.eq4.1} engendre "XII (4.1)" au lieu de "(XII 4.1)"
\makeatother
\newcommand{\refvb}[2][VI${}_{\text{B}}$ ]{\sisi{#1\refnoexpose{#2}}{\ref{#2}}}

\newcommand{\leftexp}[2]{{\vphantom{#2}}^{#1}{#2}}
\newcommand{\AAA}{\mathbf{A}}
\newcommand{\Sf}{\mathfrak{S}}
\DeclareMathOperator*{\quotprojlim}{\text{``$\varprojlim$''}}
\DeclareMathOperator*{\quotinjlim}{\text{``$\varinjlim$''}}
\DeclareMathOperator{\Fl}{F\sisi{\ell}{l}}
\DeclareMathOperator{\Pl}{P\sisi{\ell}{l}}
%\DeclareMathOperator{\cl}{c\sisi{\ell}{l}}
\DeclareMathOperator{\tdf}{tdf}
\DeclareMathOperator{\rg}{rg}
\newcommand{\iInd}{\sisi{I}{i}nd}
\newcommand{\iind}{ind}
\newcommand{\ppro}{pro}
\newcommand{\tTor}{\sisi{T}{t}or}
\newcommand{\ttor}{tor}
\newcommand{\chap}{Chap.}
\newcommand{\Ffex}{\Ff\mathit{ex}}
\newcommand{\sheafk}{\underline{k}}
\newcommand{\sheafn}{\underline{n}}
\newcommand{\sheafGamma}{\underline{\Gamma}}

% Repeated references to an already typeset footnote must not use
% \footnotemark without a matching \footnotetext: Hyperref would then create
% a live annotation to a destination that cannot exist.  This helper keeps the
% current printed footnote mark and links it to an explicit stable note target.
\makeatletter
\newcommand{\sgarepeatedfootnotemark}[1]{%
  \begingroup
  \edef\@thefnmark{\thefootnote}%
  \hyperlink{#1}{\@makefnmark}%
  \endgroup}
\makeatother

\sisi{
\DeclareMathOperator{\simarrow}{\rule{.15mm}{1.8mm}\kern -3.4pt\rightsquigarrow}
\newcommand{\sheafell}{\underline{\ell}}
\newcommand{\sheafZZ}{\underline{\ZZ}}
}{
\newcommand{\nderef}[1]{la N.D.E.~(\ref{#1}) page~\pageref{#1}}
}

%%%%%%%%%%%%%%%%%
%%%%% expose XVII
%%%%%%%%%%%%%%%%%

% END INCLUDED SGA4_English_current_buildable_source_20260906/source/fichiers-configuration-latex/macros.tex


\externaldocument{../01/01_English.tex}
\externaldocument{../02/02_English.tex}
\externaldocument{../03/03_English.tex}
\externaldocument{../04/04_English.tex}
\externaldocument{../05/05_English.tex}
\externaldocument{../05bis/05bis_English.tex}
\externaldocument{../06/06_English.tex}
\externaldocument{../07/07_English.tex}

\setboolean{original}{false}

\begin{document}

\begin{center}
VIII. Fiber functors, supports, cohomological study of finite morphisms\\
A. Grothendieck
\end{center}

\selectlanguage{english}version: \showgitstatus
\tableofcontents

\else
\chapter{Fiber functors, supports, cohomological study of finite morphisms}
%\addtocontents{toc}{A. Grothendieck}
\begin{center}
A. Grothendieck
\end{center}
\fi
\phantomsection\MakeLinkTarget*{sga4.exp.VIII}\label{sga4.exp.VIII}

% BEGIN INCLUDED SGA4_English_current_buildable_source_20260906/source/08/drafts/SGA4_VIII_LINES_43_576_English.texfrag
\setcounter{pageoriginal}{366}

% remarks:
% -- title too long at the top of the page

\section{Topological invariance of the étale
  topos}\etiquette[1]{VIII.sec1}

\begin{enonce*}{Theorem 1.1}\etiquette[1.1]{VIII.thm1.1}
Let \pageoriginale  $f : X'\to X$ be an \emph{integral surjective
radicial} morphism. Then the base-change functor
$$
f^{\ast} : X_{\text{\rm ét}} \to X'_{\text{\rm ét}}
$$
is an equivalence of sites (\ie an equivalence of categories that is
continuous, as is every quasi-inverse functor). Consequently, the
functors
\begin{align*}
f_{\ast} : \widetilde{X'_{\text{\rm ét}}} &\to \widetilde{X_{\text{\rm ét}}}\\
f^{\ast} : \widetilde{X_{\text{\rm ét}}} &\to \widetilde{X'_{\text{\rm
      ét}}}
\end{align*}
are equivalences of categories, quasi-inverse to one another.
\end{enonce*}

The first assertion is well known (\SGA\sisi{}{ 1} IX 4.10 and 4.11) when
$f$ is of finite presentation or when $f$ is flat; the general case
reduces to the case in which $f$ is of finite presentation. Indeed, it
is already known that the functor $f^{\ast}$ is fully faithful (\SGA 1 IX
3.4). It is enough to prove that $f^{\ast}$\pageoriginale  is
essentially surjective, \ie that every $Z'$ étale over $X'$ ``comes
from'' a $Z$ étale over $X$. Using the fact that $f$ is a universal
homeomorphism and the full faithfulness of $f^{\ast}$, we are reduced
to the case where $X$, $X'$, and $Z'$ are affine, the spectra of rings
$A$, $A'$, and $B'$. Writing $A'$ as the inductive limit of its
finitely generated $A$-subalgebras $A'_{i}$, $B'$ comes from an étale
algebra $B'_{i}$ over some $A'_{i}$ (\cf \EGA IV 8), which reduces us
to the case in which $A'$ is integral and of \emph{finite type} over
$A$, hence finite over $A$. We then have an $A$-isomorphism $A'\simeq
A''/J$, where $A'' \simeq A [t_1, \ldots t_n]$ and $J$ is an ideal of
$A''$. Writing $J$ as the inductive limit of its finitely generated
subideals $J_i$, and setting $A'_i = A'' /J_i$, we again have $A' =
\varinjlim_i A'_{i}$, so $B'$ comes from an étale algebra $B'_i$ over
some $A'_i$ (\emph{loco citato}). On the other hand, since $\Spec
A'\to \Spec A$ is surjective, the same holds for the $\Spec A'_i \to
\Spec A$; moreover, one easily verifies that since $\Spec A'\to
\Spec A$ is radicial, the same holds for $\Spec A'_i \to \Spec A$
for sufficiently large $i$ (apply \EGA IV 1.9.9 to $\Spec
A$\footnote{Let $S_i \subset S = \Spec A$ be the set of $s \in S$
such that the fiber of $X_i = \Spec A'_i$ \sisi{as sets}{at $s$} is
of separable rank 1. Then the $S_i$ are \emph{constructible} subsets
of $S$ (\EGA IV 9.7.9) forming an increasing family; their union is
$S$ by the hypothesis on $S'=\Spec A'\to S$ and the fact that the
fibers of $S'_i$ are \emph{noetherian} schemes. Thus $S=S_i$ for
sufficiently large $i$, by \EGA IV 1.9.9, \cqfd.}). This reduces us
to the case in which $A'$ is one of the $A'_i$, hence is of finite
presentation over $A$.

This proves the first assertion of
\sisi{\hyperref[VIII.prop2.1]{2.1}}{\hyperref[VIII.thm1.1]{1.1}}.
The assertion that $f_{\ast}$ and $f^{\ast}$ are mutually
quasi-inverse equivalences follows at once.

\begin{enonce*}{Corollary 1.2}\etiquette[1.2]{VIII.cor1.2}
Let $F$ be an abelian sheaf on $X$, and let $F'$ be its inverse image
on $X'$. Then the canonical homomorphisms
\begin{equation*}
H^{i} (X, F) \to H^{i} (X', F')\tag{$\ast$}
\end{equation*}
are isomorphisms. Likewise, if $F'$ is an abelian sheaf on $X'$ and
$F$ its direct image on $X$, the canonical homomorphisms $(\ast)$ are
isomorphisms.
\end{enonce*}

\begin{enonce*}[remark]{Examples 1.3}\etiquette[1.3]{VIII.exem1.3}
Here\pageoriginale  are some frequently used applications of
\hyperref[VIII.thm1.1]{1.1}:
\begin{enumerate}
\renewcommand{\theenumi}{\alph{enumi}}
\renewcommand{\labelenumi}{\theenumi)}
\item $X'$ is a subscheme of $X$ having the same underlying set as
  $X$, \ie defined by a nil-ideal $I$.

\item $X$ is a scheme over a separably closed field $k$, $k'$ is an
  algebraic closure of $k$, and $X'=X\otimes_k k'$.

\item Let $X$ be a geometrically unibranch scheme (for example, an
  algebraic curve over a field $k$ whose only singularities are
  ``cuspidal'' singularities, thus excluding ordinary double points in
  particular). Then, if $X'$ is the normalization of $X_{\red}$, by
  definition $X'\to X$ is radicial, so
  \sisi{\hyperref[VIII.prop2.1]{2.1}}{\hyperref[VIII.thm1.1]{1.1}}
  applies.
\end{enumerate}
\end{enonce*}

\begin{enonce*}[remark]{Remarks 1.4}\etiquette[1.4]{VIII.rem1.4}
A morphism $f$ as in \hyperref[VIII.thm1.1]{1.1} is a \emph{universal
  homeomorphism}, \ie it is a homeomorphism and remains so after every
  base extension. Conversely, if $f$ is a universal homeomorphism, one
  proves that $f$ satisfies the hypotheses of
  \sisi{\hyperref[VIII.prop2.1]{2.1}}{\hyperref[VIII.thm1.1]{1.1}}
  \footnote{\Cf \EGA IV 8.11.6 if $f$ is of finite presentation, and
  \EGA IV 18.12.11 in the general case.}, which explains the title of
  this section.

Let us point out explicitly, in connection with Example
\sisi{\hyperref[VIII.cor2.3]{2.3} c)}{\hyperref[VIII.exem1.3]{1.3} c)},
that if $X$ is an algebraic curve (over an algebraically closed field
$k$, to fix ideas) having at least one singular point that is not
``unibranch'' (for example, an ordinary double point), then (with $X'$
denoting the normalization of $X_{\red}$) the conclusion of
\sisi{\hyperref[VIII.cor2.2]{2.2}}{\hyperref[VIII.cor1.2]{1.2}}
already fails for $H^1$ with constant coefficients (compare \SGA 1 I
11 a) and \SGA 1 IX 5).
\end{enonce*}

\section{Sheaves on the spectrum of a field}\etiquette[2]{VIII.sec2}\pageoriginale

\begin{enonce*}{Proposition 2.1}\etiquette[2.1]{VIII.prop2.1}
Let $k$ be a field, $\bar{k}$ a separable closure of $k$, $\pi =
\Gal(\bar{k}/k)$ its topological Galois group, $X=\Spec k$, and
$\bar{X}=\Spec\bar{k}$ (on which $\pi$ acts on the left). Consider the
canonical functor
$$
i : X_{\text{\rm ét}} \to \widetilde{X_{\text{\rm ét}}}
$$
(defined by $i(X')(X'')=\Hom_X(X'',X')$), and the canonical functor
$$
j : \widetilde{X_{\text{\rm ét}}} \to \Bb_{\pi}
$$
from the category of sheaves on $X_{\text{\rm ét}}$ to the category
$\Bb_{\pi}$ \emph{(\cf IV \hyperref[IV.subsec2.7]{2.7})} of (discrete) sets on which $\pi$ acts
continuously on the left, defined by
$$
j(F) = \varinjlim_{\alpha} F(\Spec(k_{\alpha})) {\quoi},
$$
where $k_{\alpha}$ ranges over the finite subextensions of $\bar{k}$.
Then $i$ and $j$ are equivalences of categories. (Consequently, the
topos $\widetilde{X_{\text{\rm ét}}}$ is equivalent \emph{(IV \hyperref[IV.subsec3.4]{3.4})}
to the ``classifying topos'' $\Bb_{\pi^{\bigdot}}$.)
\end{enonce*}

The composite functor
$$
ji : X_{\text{\rm ét}} \to \widetilde{X_{\text{ét}}} \to (\pi\text{-}\Ens)
$$
is an equivalence of categories by the fundamental theorem of Galois
theory (in the form of \SGA I V). Since $\Bb_{\pi}$ is plainly a topos
(II 4.14), the same is therefore true of $X_{\text{ét}}$. Moreover, a
family $(X_i\to X)$ of morphisms in $X_{\text{ét}}$ is covering if and
only if it is surjective, \ie its image in $\Bb_{\pi}$ is surjective,
or, equivalently, covering for the canonical topology of $\Bb_{\pi}$.
This shows that the topology of $X_{\text{ét}}$ is indeed its canonical
topology. Consequently $i$ is an equivalence, and hence so is $j$.

\begin{enonce*}{Corollary 2.2}\etiquette[2.2]{VIII.cor2.2}
The functor\pageoriginale  $j$ of \hyperref[VIII.prop2.1]{2.1}
induces an equivalence \sisi{of}{between} the category of abelian
sheaves on $X=\Spec k$ and the category of Galois $\pi$-modules.
\end{enonce*}

Indeed, the latter are precisely the ``abelian groups'' of the topos
$\Bb_{\pi}$. The same argument, by ``transport of structure,'' gives:

\begin{enonce*}{Corollary 2.3}\etiquette[2.3]{VIII.cor2.3}
Let $F$ be an abelian sheaf on $X=\Spec k$, and let $M=j(F)$ be the
associated Galois $\pi$-module. Then there is a canonical isomorphism
of $\partial$-functors in $F$:
$$
H^{\ast} (X, F) \simeq H^{\ast} (\pi, M){\quoi},
$$
(where the right-hand side denotes Galois cohomology, studied, for
example, in Serre's course \emph{C.G.}).
\end{enonce*}

\begin{enonce*}{Corollary 2.4}\etiquette[2.4]{VIII.cor2.4}
Suppose that $k$ is separably closed. Then the functor
$$
\Gamma : \widetilde{X_{\text{\rm ét}}} \to (\Ens)
$$
is an equivalence of categories. If $F$ is an abelian sheaf on
$X=\Spec k$, then $H^{i}(X,F)=0$ for $i\neq0$. (In other words, the
topos $\widetilde{X_{\text{\rm ét}}}$ is a ``punctual topos''
\emph{(IV \hyperref[IV.subsec2.2]{2.2})}.)
\end{enonce*}

\begin{enonce*}{Corollary 2.5}\etiquette[2.5]{VIII.cor2.5}
Let $k'$ be a separably closed extension of a separably closed field
$k$, let $X=\Spec k$, $X'=\Spec k'$, and let $u:X'\to X$ be the
canonical morphism. Then the functor
$$
u^{\ast} : \widetilde{X_{\text{\rm ét}}} \to \widetilde{X'_{\text{\rm ét}}}
$$
is an equivalence of categories.
\end{enonce*}

\section{Fiber functors relative to geometric points of a
  scheme}\etiquette[3]{VIII.sec3}\pageoriginale

\subsection{}\etiquette[3.1]{VIII.subsec3.1}

Let $X$ be a scheme. We shall call a \emph{geometric point} of $X$
any $X$-scheme $\xi$ that is the spectrum of a \emph{separably closed}
field $\Omega$. The datum of such a point is therefore equivalent to
that of an ordinary point $x\in X$ together with a \emph{separably
closed} extension $\Omega$ of the residue field $k(x)$ of $x$. The
case most frequent for us will be that in which $\Omega$ is a
\emph{separable closure} of $k(x)$, which we shall then generally
denote by $\bar{k(x)}$; the corresponding geometric point of $X$ is
then denoted by $\bar{x}$. For a given $x\in X$, the $\bar{k(x)}$
(resp. $\bar{x}$) are determined up to \emph{nonunique}
$k$-isomorphism (resp. $X$-isomorphism).

\begin{enonce*}[remark]{Remarks 3.2}\etiquette[3.2]{VIII.rem3.2}
In most geometric questions, it is more natural to restrict to the
case in which $\Omega$ is \emph{algebraically closed}. The different
convention used here is specific to the study of the étale topology.
It is the convention we adopted in the definition of the fundamental
group in \SGA 1 V 7. Note, however, that the developments in
\emph{loco citato} are equally valid with the convention adopted here
(because the property of $\Omega$ used there is that every étale
covering of the spectrum of $\Omega$ is trivial, \ie precisely that
$\Omega$ is separably closed).
\end{enonce*}

\begin{enonce*}{Definition 3.3}\etiquette[3.3]{VIII.def3.3}
Let $X$ be a scheme, let $\xi$ be a geometric point of $X$, and let
$u:\xi\to X$ be its structural morphism. The (geometric) \emph{fiber
functor} relative to $\xi$, denoted $F\mapsto F_{\xi}$, is the
composite functor
$$
\widetilde{X_{\text{\rm ét}}} \to (\Ens)
$$
of
$$
\widetilde{X_{\text{\rm ét}}} \xrightarrow{u^{\ast}}
\widetilde{\xi_{\text{\rm ét}}} \xrightarrow{\Gamma_{\xi}} (\Ens){\quoi}.
$$
\end{enonce*}

Equivalently, by \hyperref[VIII.cor2.4]{2.4}, one may say that a
geometric point $\xi$ of $X$ defines a point of the topos
$X^{\sim}_{\text{ét}}$ (IV \hyperref[IV.def6.1]{6.1}), whose associated fiber functor
(\loccit) is simply being taken.

\setcounter{subsection}{3}
\subsection{}\etiquette[3.4]{VIII.subsec3.4}

Since the\pageoriginale  functor $\Gamma_{\xi}$ (the global-sections
functor on $\xi$) is an \emph{equivalence} of categories \hyperref[VIII.cor2.4]{(2.4)},
knowing the fiber functors $F\mapsto F_{\xi}$ is essentially
equivalent to knowing the inverse-image functors $u^{\ast}$. Moreover,
it follows at once from \hyperref[VIII.cor2.5]{2.5} that if $\xi$ is a
geometric point of $X$ for which there exists an $X$-morphism
$\xi'\to\xi$ (\ie $\xi'$ corresponds to a separably closed extension
$\Omega'$ of $\Omega$), then the canonical homomorphism
$F_{\xi}\to F_{\xi'}$ is a \emph{functorial isomorphism}
$$
F_{\xi} \simeq F_{\xi'}{\quoi}.
$$
This shows that, for the study of fiber functors, one may (after
replacing $\Omega$ by the separable algebraic closure of $k$ in
$\Omega$) restrict oneself to the case where $\Omega = k(\xi)$ is a
separable closure $\bar{k(x)}$ of $k(x)$, and that the fiber functors
corresponding to geometric points of $X$ lying over the same $x
\in X$ are isomorphic to one another (noncanonically).

In accordance with the conventions of \hyperref[VIII.subsec3.1]{3.1},
we shall denote by $F_{\bar{x}}$
a fiber functor corresponding to the choice of a separable closure
$\bar{k(x)}$ of $k(x)$.

Let us point out an evident transitivity property (which shows the
technical usefulness of not restricting oneself exclusively to geometric
points defined by separable closures of residue fields): Let $f :
X'\to X$ be a morphism of schemes, and let $u' : \xi'\to X'$ be a
geometric point of $X'$. Then $u = f u': \xi'\to X$ is a geometric
point of $X$, which we shall denote by $\xi$; relative to these geometric
points, there is an isomorphism, functorial in $F \in \Ob
\widetilde{X_{\text{ét}}}$:
$$
f^{\ast} (F)_{\xi'} \simeq F_{\xi}{\quoi}.
$$
Indeed, this\pageoriginale  follows from the transitivity formula for
inverse-image
functors
$$
(fu')^{\ast} \xrightarrow{\sim} u'^{\ast} f^{\ast}{\quoi}.
$$

\begin{enonce*}{Theorem 3.5}\etiquette[3.5]{VIII.thm3.5}
Let $X$ be a scheme.
\begin{enumerate}
\renewcommand{\theenumi}{\alph{enumi}}
\renewcommand{\labelenumi}{\theenumi)}
\item For every geometric point $\xi$ of $X$, the fiber functor $F
  \mapsto F_{\xi}$, $\widetilde{X_{\text{\rm ét}}} \to (\Ens)$ commutes with
  arbitrary $\varinjlim$ (indexed by categories $\in
  \Uu$) and with finite $\varprojlim$\footnote{Thus it is indeed
  a ``fiber functor'' in the sense of Exposé IV.}. In particular, it
  takes sheaves of groups (\resp abelian sheaves, \etc)
  to groups (\resp abelian groups, \etc).

\item As $x$ ranges over the points of $X$, the fiber functors
  $F_{\bar{x}}$ form a ``conservative'' family of functors, \ie if $u : F \to G$ is a homomorphism of
  sheaves on $X$, then $u$ is an isomorphism if and only if, for
  every $x \in X$, the corresponding homomorphism $u_{\bar{x}} :
  F_{\bar{x}} \to G_{\bar{x}}$ is an isomorphism\footnote{Thus there are
  ``enough fiber functors'' of the form $F \mapsto
  F_{\bar{x}}$, in the terminology of IV \hyperref[IV.prop6.5]{6.5}.}.
\end{enumerate}
\end{enonce*}

Taking into account the exactness properties in a), let us immediately
note the following formal consequences of b):

\begin{enonce*}{Corollary 3.6}\etiquette[3.6]{VIII.cor3.6}
Let $u : F \to G$ be a homomorphism of sheaves. Then $u$ is a
monomorphism (\resp an epimorphism) if and only if, for every $x
\in X$, $u_{\bar{x}}:F_{\bar{x}}\to G_{\bar{x}}$ is injective
(\resp surjective).
\end{enonce*}

\begin{enonce*}{Corollary 3.7}\etiquette[3.7]{VIII.cor3.7}
Let $u$, $v : F \to G$ be two morphisms of sheaves on $X$. Then
$u = v$ if and only if, for every $x \in X$, one has $u_{\bar{x}} =
v_{\bar{x}} : F_{\bar{x}} \to G_{\bar{x}}$. In particular, if $u$,
$v$ are two sections of $F$, one has\pageoriginale  $u = v$ if and
only if
$u_{\bar{x}} = v_{\bar{x}}$ in $F_{\bar{x}}$ for every $x \in X$.
\end{enonce*}

\begin{enonce*}{Corollary 3.8}\etiquette[3.8]{VIII.cor3.8}
Let $F \to G \to H$ be a sequence of two homomorphisms of sheaves on
$X$. Then the sequence is exact if and only if, for every $x \in X$,
the corresponding sequence $F_{\bar{x}} \to G_{\bar{x}} \to H_{\bar{x}}$
is exact.
\end{enonce*}

We leave to the reader the task of deducing the corollaries
\footnote{\Cf IV \hyperref[IV.exer6.5.2]{6.5.2}.}, %% issue: check that (**) is indeed present upon compilation
which will moreover be proved in part
in the proof that follows.

Assertion a) follows from the exactness properties already noted
in (\hyperrefext[VII][VII.subsec1.4]{1.4}) for every inverse-image functor, such as
$u^{\ast} : \widetilde{X_{\text{ét}}} \to
\widetilde{\xi_{\text{ét}}}$, and from the fact that
$\widetilde{\xi_{\text{ét}}} \to (\Ens)$ is an equivalence. To
prove b), let us note the following formula, which gives a procedure for
\emph{computing} fiber functors:

\begin{enonce*}{Proposition 3.9\footnotemark}\etiquette[3.9]{VIII.prop3.9}
\footnotetext{\Cf IV \hyperref[IV.eq6.8.3]{6.8.3}} %% issue: check the number of * upon compilation
Let $\xi \xrightarrow{u} X$ be a geometric
point of $X$, and let $C_{\xi}$ be the category of ``pointed étale
$X$-schemes at $\xi$'', \ie of pairs $(X', u')$, where
$X'\in \Ob X_{\text{\rm ét}}$ is a scheme étale over $X$, and $u'$ is an
$X$-morphism $\xi \to X'$. Then the opposite category $C^{0}_{\xi}$
is filtered and, for a varying sheaf $F$ (\resp presheaf $P$)
on $X$, there is a canonical functorial isomorphism
$$
F_{\xi} \simeq \varinjlim_{C^{0}_{\xi}} F(X'), (\emph{\resp}
(aP)_{\xi} \simeq \varinjlim_{C^{0}_{\xi}} P(X')){\quoi},
$$
(where $aP$ = the sheaf associated with $P$).
\end{enonce*}

Indeed, let us first note that for every presheaf $P$ on
$\xi_{\text{ét}}$,
if $aP$ is the associated sheaf, the natural homomorphism
$$
P(\xi) \to aP(\xi)
$$
is an isomorphism, as follows, for example, immediately from the
explicit computation of\pageoriginale  $aP$ as $L(LP)$ (II \hyperref[II.theo3.4]{3.19}). Let
$\varphi :
\widetilde{X_{\text{ét}}} \to \widetilde{\xi_{\text{ét}}}$ be the
``inverse-image'' functor. Then $u^{\ast} = \varphi^{\ast}$ is nothing
other (III \hyperref[III.prop2.3]{2.3}. (d)) than the composite $a\varphi^{\bigdot} i$, where $i :
\widetilde{X_{\text{ét}}} \to \widehat{X_{\text{ét}}}$
is the inclusion functor, and $a :
\widehat{\xi_{\text{ét}}} \to
\widetilde{\xi_{\text{ét}}}$ is the ``associated
sheaf'' functor. By the preceding remark we find
$$
u^{\ast} (F) (\xi) = (\varphi^{\bigdot} i(F)) (\xi){\quoi},
$$
so it is enough to apply the explicit expression for
$\varphi^{\bigdot}$ given in the relevant proof.\SGAArchiveOnly{\nde{The frozen source
leaves the exact internal reference unspecified.}}
The corresponding case, where one starts with a presheaf $P$ on $X$, is proved in the same way, by
using the functorial isomorphism $\varphi^{\ast} a(P) \simeq a
\varphi^{\ast} (P)$.\SGAArchiveOnly{\nde{The frozen source leaves the exact internal
reference unspecified.}}
The fact that the category $C_{\xi}$ is filtered (which remains true for \emph{every} $X$-scheme $\xi$, not
necessarily reduced to a geometric point)
follows immediately from the
fact that finite projective limits exist in $C_{\xi}$, which
follows from the fact that $X_{\text{ét}}$ is a subcategory of
$(\Sch) /X$ stable under finite projective limits.

Now let $u : F \to G$ be a homomorphism of sheaves on $X$
such that, for every $x \in X$, $u_{\bar{x}} : F_{\bar{x}} \to
G_{\bar{x}}$ is a monomorphism; let us prove that $u$ is a
monomorphism. For this, we must prove that if $X'\in \Ob
X_{\text{ét}}$, $\varphi$, $\psi \in F(X')$ are such that $u
(\varphi) = u (\psi)$, then $\varphi = \psi$. Replacing $X$ by
$X'$ and using the transitivity of fiber functors, we may
suppose $X = X'$. For every $x \in X$, we then have $u
(\varphi)_{\bar{x}} = u (\psi)_{\bar{x}}$, \ie $u_{\bar{x}}
(\varphi_{\bar{x}}) = u_{\bar{x}} (\psi_{\bar{x}})$, hence
$\varphi_{\bar{x}} = \psi_{\bar{x}}$ since $u_{\bar{x}}$ is a
monomorphism. Using \hyperref[VIII.prop3.9]{3.9}, %% XXX Reference to correct}
it follows that there exists an $X'_x \in \Ob X_{\text{ét}}$, whose image contains $x$, such
that $\varphi$ and $\psi$ have the same inverse image on $X'_x$. Since,
as $x$ varies in $X$, the $X'_x$ form a covering family of
$X$, it follows that $\varphi = \psi$.

This\pageoriginale proves that if the $u_{\bar{x}}$ are
monomorphisms, then so
is $u$. Suppose in addition that the $u_{\bar{x}}$ are
epimorphisms; let us prove that the same holds for $u$, and hence that $u$ is
an isomorphism. It remains to prove that for every $X'\in \Ob X_{\text{ét}}$
and every $\psi \in G (X')$, there exists $\varphi \in F(X')$ such that
$u(X') (\varphi) = \psi$. We may again suppose $X'= X$, and
using \hyperref[VIII.prop3.9]{3.9}, %% XXX Reference to correct
we see that, for every $x \in X$, \sisi{}{there} exists $X'_x \in \Ob
X_{\text{ét}}$ such that the image of $X'_x$ in $X$ contains $x$, and a
$\varphi_x \in F(X'_x)$ whose image in $G (X'_x)$ is
the inverse image of $\psi$. Using the fact that $u$ is a
monomorphism, we see that the $\varphi_{x}$ agree on the
$\fprod{X'_x}{X'_y}{X}$ $(x, y \in X)$ and therefore come from a
uniquely determined $\varphi \in F (X)$; one then has $u(\varphi) =
\psi$, since the inverse images on the $X'_{x}$
agree. \cqfd.

\begin{enonce*}[remark]{Scholium 3.10}\etiquette[3.10]{VIII.sch3.10}
Theorem \hyperref[VIII.thm3.5]{3.5} implies that the collection of
``fiber functors''
for the étale topology enjoys the same essential properties as
in the usual theory of sheaves on a topological
space. We shall use \hyperref[VIII.thm3.5]{3.5}
and its corollaries very frequently,
and for this reason we \sisi{}{shall} generally dispense with referring to them explicitly.
\end{enonce*}

\setcounter{subsection}{10}
\subsection{}\etiquette[3.11]{VIII.subsec3.11}

For every $x \in X$ we shall also, by abuse of language, denote by
$x$ the $X$-scheme $\Spec k(x)$, for the usual structural morphism
$$
i_x : x = \Spec k(x) \to X{\quoi}.
$$
It gives rise to a canonical functor
$$
i_x^{\ast} : \widetilde{X_{\text{ét}}} \to \widetilde{x_{\text{ét}}}{\quoi}.
$$
If $F$ is a sheaf on $X$, we shall put $F_x = i^{\ast}_x
(F)$; thus it is a sheaf on $x = \Spec k(x)$ (and as such
may also be identified with an\pageoriginale étale scheme over $x$, by
virtue of \hyperref[VIII.prop2.1]{2.1}). It depends functorially on $F$, and
the functor $i^{\ast}_x : F \mapsto F_x$ still commutes with
arbitrary $\varinjlim$ and finite $\varprojlim$, by virtue of
\hyperrefext[VII][VII.subsec1.4]{1.4}.

If $\bar{x} = \Spec \bar{k(x)}$, the fiber functor $F \mapsto
F_{\bar{x}}$ is canonically isomorphic to the composite of the functor $F
\mapsto F_x$ and the functor $j$ of \sisi{\hyperref[VIII.thm1.1]{1.1}}{\hyperref[VIII.prop2.1]{2.1}} (taking into account that the latter is
isomorphic to the fiber functor relative to the geometric point $\Spec
\bar{k}$ of $\Spec k$). It follows immediately from this and from
\hyperref[VIII.thm3.5]{3.5} that the
system of functors $(F \mapsto F_x)$ $(x \in X)$ is likewise
conservative. Moreover, if $k(x)$ is separably closed, \ie $x =
\bar{x}$, the notations $F_x$ and $F_{\bar{x}}$ are, strictly speaking,
in conflict (the first denoting an object of
$\widetilde{x_{\text{ét}}}$, the second a set), but by virtue of
\sisi{\hyperref[VIII.thm1.1]{1.1}}{\hyperref[VIII.prop2.1]{2.1}} this conflict is inessential.

\subsection{}\etiquette[3.12]{VIII.subsec3.12}

The preceding remarks show that if $\bar{x} = \Spec
\bar{k(x)}$, then, for every sheaf $F$ on $X$, the group $\pi_x = \Gal
(\bar{k(x)}/ k(x))$ acts naturally (on the left) on
the fiber set $F_{\bar{x}}$, so that $F \mapsto F_{\bar{x}}$
may in fact be regarded as a functor
$$
\widetilde{X_{\text{ét}}} \to \Bb_{\pi_x}
$$
whose knowledge is essentially equivalent (thanks to
\sisi{\hyperref[VIII.thm1.1]{1.1}}{\hyperref[VIII.prop2.1]{2.1}}) to that of the
functor $F \mapsto F_x$.

This shows, in an important special case, that contrary to what
happens for sheaves on topological spaces, ``fiber
functors'' generally admit nontrivial automorphisms.
In general, we shall determine in \hyperref[VIII.thm7.9]{7.9} all the
homomorphisms from one fiber functor to another.

\begin{enonce*}[remark]{Remark 3.13}\etiquette[3.13]{VIII.rem3.13}
Here is\pageoriginale a sometimes useful generalization
of \hyperref[VIII.thm3.5]{3.5} b). Consider a subset $E$ of $X$ such that,
for every étale morphism $f : X' \to X$, $E' = f^{-1} (E)$ is ``very dense'' in $X'$ (EGA IV 10.1.3), \ie for every $f$
as above, every open subset $U$ of $X'$ which contains $E'$ is equal
to $X'$. Then the fiber functors $F_{\bar{x}}$ on
$\widetilde{X_{\text{\rm ét}}}$, for $x \in E$, already form a
``conservative family'' (and consequently Corollaries
\hyperref[VIII.cor3.6]{3.6}, \hyperref[VIII.cor3.7]{3.7}, and \hyperref[VIII.cor3.8]{3.8} remain
valid when one restricts oneself to choosing $x \in E$ in them). This follows
immediately by repeating the proof given for
\hyperref[VIII.thm3.5]{3.5} b). The preceding result applies in particular
in the following situations:
\begin{enumerate}
\renewcommand{\theenumi}{\alph{enumi}}
\renewcommand{\labelenumi}{\theenumi)}
\item $X$ is a Jacobson scheme (EGA IV 10.4.1.), for example
  locally of finite type over a field, or over $\Spec(𝒵)$, and $E$
  is the set of closed points of $X$.

\item $X$ is locally of finite type over a scheme $S$, and $E$ is
  the set of points of $X$ which are closed in their fiber.

\item $X$ is locally noetherian, and $E$ is the set of $x \in
  X$ such that $\bar{x}$ is a finite set (\ie a
  semilocal scheme of dimension $\leq 1$) \cf \EGA IV 10.5.3 and 10.5.5.
\end{enumerate}
\end{enonce*}

% END INCLUDED SGA4_English_current_buildable_source_20260906/source/08/drafts/SGA4_VIII_LINES_43_576_English.texfrag


% BEGIN INCLUDED SGA4_English_current_buildable_source_20260906/source/08/drafts/SGA4_VIII_LINES_577_804_English.texfrag
\section{Strictly local rings and schemes}\etiquette[4]{VIII.sec4}

\subsection{}\etiquette[4.1]{VIII.subsec4.1}

Recall that, following Azumaya, a local ring $A$ is said to be
\emph{henselian} if it satisfies the following equivalent conditions:
\begin{enumerate}
\renewcommand{\theenumi}{\roman{enumi}}
\renewcommand{\labelenumi}{(\theenumi)}
\item Every finite $A$-algebra $B$ is a direct product
$\prod_i B_i$ of local rings (the $B_i$ are then necessarily
isomorphic to the $B_{\mathfrak m_i}$, where\pageoriginale
$\mathfrak m_i$ runs through the maximal ideals of $B$).

\item Every $A$-algebra $B$ which is the localization at a prime ideal
$\mathfrak p$ lying over the maximal ideal $\mathfrak m$ of $A$ of
an $A$-algebra $C$ of finite type, and such that
$B/\mathfrak mB$ is finite over $k=A/\mathfrak m$, is finite over
$A$.

\item The same as (ii), but with $C$ assumed étale over $A$ at
$\mathfrak p$ and $k\to B/\mathfrak mB$ an isomorphism.

\item \lq\lq Hensel's lemma'' in its classical form holds for $A$, that is,
for every monic polynomial $F\in A[T]$, writing
$\overset{\bigdot}{F}\in k[T]$ for the corresponding reduced
polynomial, and for every factorization
$$
\overset{\bigdot}{F}
=\overset{\bigdot}{F}_1\overset{\bigdot}{F}_2
$$
of $\overset{\bigdot}{F}$ as a product of two relatively prime monic
polynomials, $\overset{\bigdot}{F}_1$ and
$\overset{\bigdot}{F}_2$ lift (necessarily uniquely) to
$F_1,F_2\in A[T]$ such that $F=F_1F_2$.

For the proof of these equivalences and the general properties of
henselian rings, see \EGA IV 18.5.11. Recall that if $A$ is any local
ring, one proves (Nagata) that there exists a local homomorphism
$$
A\to A^h
$$
with $A^h$ \emph{henselian}, which is \emph{universal} for local
homomorphisms from $A$ to local
\sisi{Noetherian}{henselian} rings. The ring $A^h$ is called the
henselization of $A$. The construction given in \EGA IV 18.6
(different from Nagata's) consists in setting
$$
A^h=\varinjlim_i A_i{\quoi},
$$
where $A_i$ runs through the $A$-algebras essentially of finite type
and essentially étale\pageoriginale over $A$ (that is, localizations
of an étale $A$-algebra) such that $A\to A_i$ is a local
homomorphism and the residue extension $k(A_i)/k(A)$ is trivial. This
construction shows in particular that $A^h$ is flat over $A$, and
that if $A$ is Noetherian, then so is its henselization.
\end{enumerate}

\begin{enonce*}{Definition 4.2}\etiquette[4.2]{VIII.def4.2}
A local ring $A$ is said to be strictly local \emph{if it satisfies
one of the following equivalent conditions}:
\begin{enumerate}
\renewcommand{\theenumi}{\roman{enumi}}
\renewcommand{\labelenumi}{(\theenumi)}
\item $A$ is henselian (\hyperref[VIII.subsec4.1]{4.1}) and its
residue field is separably closed.

\item Every local homomorphism $A\to B$, where $B$ is a localization
of an étale algebra over $A$, is an isomorphism.
\end{enumerate}

A scheme is called a \emph{strictly local scheme} if it is isomorphic
to the spectrum of a strictly local ring.
\end{enonce*}

The equivalence of (i) and (ii) follows at once from
\hyperref[VIII.subsec4.1]{4.1}. Note that, in geometric form,
\hyperref[VIII.def4.2]{4.2} (ii) (respectively
\hyperref[VIII.subsec4.1]{4.1} (iii)) says that for every étale
$Y$-scheme $X$ and every point $x$ of the fiber $X_y$ (respectively
every point of $X_y$ \emph{rational over} $k(y)$), where $y$ is the
closed point of $Y$, there exists a section of $X$ over $Y$ passing
through $x$ (and this section is necessarily unique).

\begin{enonce*}{Definition 4.3}\etiquette[4.3]{VIII.def4.3}
Let $X$ be a scheme, let $\Oo_{X_{\text{\rm ét}}}$ be the sheaf on
$X_{\text{\rm ét}}$ defined by
$$
\Oo_{X_{\text{\rm ét}}}(X')=\Gamma(X',\Oo_{X'})
$$
(compare \emph{\hyperrefext[VII][VII.sec2]{2}. c)}), and let
$u:\xi=\Spec(\Omega)\to X$ be a geometric point of $X$. The
\emph{strictly local ring of $X$ at $\xi$}, denoted
$\Oo_{{X,}_{\xi}}$ or simply\pageoriginale $\Oo_\xi$, is the
fiber of the sheaf $\Oo_{X_{\text{\rm ét}}}$ at the point $\xi$.
Its spectrum is called the \emph{strict localization of $X$ at
$\xi$}.
\end{enonce*}

By \hyperref[VIII.prop3.9]{3.9}, the fiber $\Oo_{{X,}_{\xi}}$ has
the description
\begin{equation*}
\Oo_{{X,}_{\xi}}=\varinjlim\Gamma(X',\Oo_{X'}){\quoi},
\tag{$\ast$}
\end{equation*}
where $X'$ runs through the opposite filtered category of étale
$X$-schemes pointed by $\xi$, that is, equipped with an $X$-morphism
$\xi\to X'$. By transitivity of inductive limits, one may replace
$\Gamma(X',\Oo_{X'})$ in the right-hand side by $\Oo_{X',x'}$,
where $x'$ is the image of $\xi$ in $X'$, and one obtains
\begin{equation*}
\Oo_{{X,}_{\xi}}=\varinjlim_i A_i
\tag{\ast\ast}
\end{equation*}
where $A_i$ runs through the algebras essentially of finite type and
étale over $A=\Oo_{X,x}$, equipped with a $k$-homomorphism
$k(A_i)\to\Omega$. (Here $x$ denotes the image of $\xi$ in $X$.)
The transition homomorphisms between the $A_i$ are, of course, the
local homomorphisms of $A$-algebras $A_i\to A_j$ that make the
corresponding triangle commute:
\[
\xymatrix@R=1.25cm@C=1.75cm{
k(A_i)  \ar[r] \ar[rd]& k(A_j)\ar[d]    \\
                      &  \Omega.
}
\]
This implies that if there exists an admissible homomorphism from
$A_i$ to $A_j$, it is unique. Thus the inductive limit
$({\ast}{\ast})$ is a limit over an upward-filtered ordered index
set.

\setcounter{subsection}{3}
\subsection{}\etiquette[4.4]{VIII.subsec4.4}

The description\pageoriginale $({\ast}{\ast})$ shows that the
strictly local ring of $X$ at $\xi$ depends essentially only on the
usual local ring $A=\Oo_{X,x}$ ($x=u(\xi)$) and on the extension
$\Omega$ of $k=k(A)$. It is also called the \emph{strict
henselization} of $A$ (relative to the given separably closed
extension $\Omega$ of $k$), and is denoted $A^{hs}$. We refer to
\EGA IV 18.8 for the general properties of $A^{hs}$. We merely point
out that the construction again shows that $A^{hs}$ is flat over
$A$, and that it is Noetherian if $A$ is. Moreover, one proves easily
(\loccit) that the local homomorphism
$$
A\to A^{hs}
$$
together with the $k$-homomorphism
$$
k(A^{hs})\to\Omega
$$
solves the universal problem, relative to the data of $A$ and the
extension $\Omega$ of $k$, of finding local homomorphisms
$$
A\to A'{\quoi},
$$
where $A'$ is \emph{strictly local}, equipped with a
$k$-homomorphism
$$
k(A')\to\Omega{\quoi}.
$$
In particular, $A^{hs}$ \emph{is indeed a strictly local ring}
(which justifies the terminology of \hyperref[VIII.def4.3]{4.3}).

\subsection{}\etiquette[4.5]{VIII.subsec4.5}

Keep the notation of \hyperref[VIII.def4.3]{4.3}, and choose an
affine open neighborhood $U$ of $x$. Note that in
\hyperref[VIII.prop3.9]{3.9} one may plainly restrict to those $X'$
which are affine over $U$. These then form a pseudo-filtered inverse
system of affine schemes, so that we are\pageoriginale in the general
situation of \hyperrefext[VII][VII.sec5]{5} (see VI 5.12 b).
By the expression $(\ast)$ for $\Oo_{X,\xi}$, there is an
$X$-isomorphism
\begin{equation*}
\varprojlim_{\substack{X'\ \mbox{\rm étale}\\ \mbox{affine over }U,\\
X'_{\xi}\mbox{\rm -pointed}}}X'\simeq\Spec\Oo_{X,\xi}
\tag{\ast\ast\ast}
\end{equation*}
which makes precise the intuitive geometric meaning of
$\Spec(\Oo_{X,\xi})$ as the \lq\lq\emph{limit of the étale
neighborhoods of $\xi$}''.

\begin{enonce*}{Proposition 4.6}\etiquette[4.6]{VIII.prop4.6}
Let $X$ be a strictly local scheme, and let $x$ be its closed point,
which is therefore a geometric point of $X$. Then the functor
$F\mapsto\Gamma(X,F):\widetilde{X_{\text{\rm ét}}}\to(\Ens)$
is canonically isomorphic to the fiber functor $F\mapsto F_x$, and
hence commutes with arbitrary $\varinjlim$\footnote{not necessarily
filtered!}.
\end{enonce*}

Indeed, by Definition \hyperref[VIII.def4.2]{4.2} (ii), the final
object $X$ of the category of étale $X$-schemes pointed by $x$
considered in \hyperref[VIII.prop3.9]{3.9} dominates every other
object; hence $\varinjlim F(X')$ is canonically isomorphic to $F(X)$,
as required.

In particular, the functor induced by $\Gamma$ on the category of
abelian sheaves is exact, and therefore:

\begin{enonce*}{Corollary 4.7}\etiquette[4.7]{VIII.cor4.7}
Under the hypotheses of \hyperref[VIII.prop4.6]{4.6}, for every
abelian sheaf $F$ on $X$ one has
$$
H^q(X,F)=0\quad\mbox{\rm for}\quad q\neq0{\quoi}.
$$
\end{enonce*}

\begin{enonce*}{Corollary 4.8}\etiquette[4.8]{VIII.cor4.8}
Let $X$ be a scheme, let $\xi$ be a geometric point of $X$, let
$\bar X=\Spec\Oo_{X,\xi}$ be the corresponding strict
localization, let $F$ be a sheaf on $X$, and let $\bar F$ be its
inverse image on $\bar X$. Then there is a functorial isomorphism
$$
F_\xi\simeq\Gamma(\bar X,\bar F){\quoi}.
$$
\end{enonce*}

Indeed,\pageoriginale $\xi$ may also be regarded as a geometric point
of $\bar X$, and it suffices to apply
\hyperref[VIII.prop4.6]{4.6} and the transitivity of fibers
(\hyperref[VIII.subsec3.4]{3.4}).


% END INCLUDED SGA4_English_current_buildable_source_20260906/source/08/drafts/SGA4_VIII_LINES_577_804_English.texfrag


% BEGIN INCLUDED SGA4_English_current_buildable_source_20260906/source/08/drafts/SGA4_VIII_LINES_805_1011_English.texfrag
\section{Application to the computation of the fibers of the
$R^qf_*$}\etiquette[5]{VIII.sec5}

\subsection{}\etiquette[5.1]{VIII.subsec5.1}

Let
$$
f:X\to Y
$$
be a morphism of schemes and let $F$ be an abelian sheaf on $X$. We
wish to determine the $R^qf_*(F)$. By
\hyperref[VIII.thm3.5]{3.5}, this amounts, up to very little, to
knowing the geometric fibers $R^qf_*(F)_{\bar y}$ for $y\in Y$. Now
(\hyperrefext[V][V.prop5.1]{5.1}) $R^qf_*(F)$ is the sheaf associated
with the presheaf
$$
\Hh^q:Y'\longmapsto H^q(\fprod{X}{Y'}{Y},F)
$$
on $Y$, and hence (\hyperref[VIII.prop3.9]{3.9})
$$
R^qf_*(F)_{\bar y}=\varinjlim_{Y'}\Hh^q(Y'){\quoi},
$$
where the inductive limit is taken over the opposite filtered
category of the étale $Y$-schemes pointed by $\bar y$. After choosing
an affine open neighborhood $U$ of $y$, one may restrict the limit on
the right to the cofinal category of those $Y'$ which are affine and
lie over $U$. Thus one obtains
\begin{equation*}
R^qf_*(F)_{\bar y}\simeq
\varinjlim H^q(\fprod{X}{Y'}{Y},F){\quoi},
\tag{5.1.1}
\end{equation*}\etiquette[5.1.1]{VIII.eq5.1.1}
where $Y'$ varies in the preceding category.

Put
$$
\bar Y=\Spec(\Oo_{Y,\bar y})=\varprojlim Y'
$$
(see \hyperref[VIII.subsec4.5]{4.5}), and
$$
\bar X=\fprod{X}{\bar Y}{Y}{\quoi}.
$$
Observe that\pageoriginale in the inverse system
$\fprod{X}{Y'}{Y}$ (deduced from that of the $Y'$ by the base change
$X\to Y$), the transition morphisms are affine. We are therefore in
the general situation of \hyperrefext[VII][VII.subsec5.1]{5.1}, and
the construction of $\varprojlim$ in \loccit immediately gives a
canonical isomorphism
$$
\bar X=\varprojlim\fprod{X}{Y'}{Y}{\quoi}.
$$
Let $\bar F$ denote the sheaf on $\bar X$ that is the inverse image of
$F$. We then obtain a canonical homomorphism
$$
\varinjlim_{Y'}H^q(\fprod{X}{Y'}{Y},F)
\longrightarrow H^q(\bar X,\bar F),
$$
and hence, by comparison with \eqhyperref[VIII.eq5.1.1]{5.1.1}, a
canonical homomorphism
\begin{equation*}
R^qf_*(F)_{\bar y}\longrightarrow H^q(\bar X,\bar F)
\tag{5.1.2}{\quoi},
\end{equation*}\etiquette[5.1.2]{VIII.eq5.1.2}
which is plainly functorial in $F$.

Suppose now that $f:X\to Y$ is quasi-compact and quasi-separated.
Then the same is true of $\fprod{X}{Y'}{Y}\to Y'$, and since $Y'$ is
affine, the schemes $\fprod{X}{Y'}{Y}$ are quasi-compact and
quasi-separated. Applying \eqhyperref[VIII.eq5.1.1]{5.1.1} and the
limit-passage theorem \hyperrefext[VII][VII.cor5.8]{5.8}, one obtains:

\begin{enonce*}{Theorem 5.2}\etiquette[5.2]{VIII.thm5.2}
Let $f:X\to Y$ be a quasi-compact and quasi-separated morphism of
schemes, let $F$ be an abelian sheaf on $X$, let $y$ be a point of
$Y$, and let $\bar y$ be the geometric point over $y$ associated with
a separable closure $k(\bar y)$ of $k(y)$. Let
$\bar Y=\Spec(\Oo_{Y,\bar y})$ be the corresponding strict
localization, let $\bar X=\fprod{X}{\bar Y}{Y}$, and let $\bar F$ be
the inverse image of $F$ on $\bar X$. Then the canonical
homomorphism \eqhyperref[VIII.eq5.1.2]{5.1.2} is an isomorphism:
$$
R^qf_*(F)_{\bar y}\simeq H^q(\bar X,\bar F){\quoi}.
$$
\end{enonce*}

This statement practically reduces the determination of the fibers
of a sheaf\pageoriginale $R^qf_*(F)$ to the determination of the
cohomology of a scheme over a strictly local scheme, and it will be
used constantly below. Technically, \hyperref[VIII.thm5.2]{5.2}
explains the important role of henselian and strictly local rings in
the study of étale cohomology.

\begin{enonce*}[remark]{Remark 5.3}\etiquette[5.3]{VIII.rem5.3}
The statement remains valid for $R^0f_*(F)=f_*(F)$, for an arbitrary
sheaf of sets. Moreover, \hyperref[VIII.thm5.2]{5.2} also remains
valid for $R^1f_*(F)$ when $F$ is a sheaf of groups (not necessarily
commutative), with the usual definition of $R^1f_*(F)$ (see Giraud's
thesis).
\end{enonce*}

\setcounter{subsection}{3}
\subsection{}\etiquette[5.4]{VIII.subsec5.4}

Suppose now that $f$ is a \emph{finite} morphism. Then
$\bar f:\bar X\to\bar Y$ is finite, and since $\bar Y$ is strictly
local, it follows that $\bar X$ is a finite disjoint union of strictly
local schemes $\bar X_i$. Hence, using
\hyperref[VIII.cor4.7]{4.7}, one finds
\begin{equation*}
H^q(\bar X,\bar F)=0\quad\mbox{for}\quad q\neq0{\quoi}.
\tag{5.4.1}
\end{equation*}\etiquette[5.4.1]{VIII.eq5.4.1}
Furthermore, the components $\bar X_i$ of $\bar X$ correspond to the
points $\bar x_i$ of $\bar X$ above the closed point $\bar y$ of
$\bar Y$, that is, to the points of
$\bar X_{\bar y}=X_y\otimes_{k(y)}k(\bar y)$. These points may also
be regarded as geometric points of $X$, and $\bar X_i$ is then
precisely the strict localization of $X$ at $\bar x_i$. Thus
(\hyperref[VIII.cor4.8]{4.8})
$$
H^0(\bar X_i,\bar F)\simeq F_{\bar x_i}{\quoi},
$$
whence
\begin{equation*}
H^0(\bar X,\bar F)=\mathop{\prod}_iF_{\bar x_i}{\quoi}.
\tag{5.4.2}
\end{equation*}\etiquette[5.4.2]{VIII.eq5.4.2}
This is\pageoriginale an isomorphism functorial in the sheaf of sets
$F$ (there is no need here to restrict to abelian sheaves). Taking
\hyperref[VIII.thm5.2]{5.2} and \hyperref[VIII.rem5.3]{5.3} into
account, formulas \eqhyperref[VIII.eq5.4.1]{5.4.1} and
\eqhyperref[VIII.eq5.4.2]{5.4.2} give:

\begin{enonce*}{Proposition 5.5}\etiquette[5.5]{VIII.prop5.5}
Let $f:X\to Y$ be a finite morphism of schemes, and let $y$ be a
point of $Y$. Then for every sheaf $F$ on $X$, there is an
isomorphism, functorial in $F$,
$$
f_*(F)_{\bar y}\simeq
\prod_{\bar x\in X_y\otimes_{k(y)}k(\bar y)}F_{\bar x}{\quoi}.
$$
Consequently, the formation of $f_*(F)$ commutes with every base
change $Y'\to Y$; and if $F$ is an abelian sheaf, then
$$
R^qf_*(F)=0\quad\mbox{if}\quad q>0{\quoi}.
$$
\end{enonce*}

Note that the first formula of \hyperref[VIII.prop5.5]{5.5} is in
fact independent of \hyperref[VIII.thm5.2]{5.2} and of the general
limit-passage theorem \hyperrefext[VII][VII.thm5.7]{5.7}; together
with \hyperref[VIII.thm3.5]{3.5}, it implies that
$F\mapsto f_*(F)$ is an \emph{exact functor} on the category of
abelian sheaves, and hence again that $R^qf_*(F)=0$ for $q\neq0$.
Here is a slight variant of \hyperref[VIII.prop5.5]{5.5}:

\begin{enonce*}{Corollary 5.6}\etiquette[5.6]{VIII.cor5.6}
Let $f:X\to Y$ be an integral morphism. Then for every abelian sheaf
$F$ on $X$ one has
$$
R^qf_*(F)=0\quad\mbox{for}\quad q\neq0{\quoi}.
$$
Moreover, the functor $f_*$ on sheaves of sets commutes with every
base change $Y'\to Y$.
\end{enonce*}

Indeed, by \hyperref[VIII.thm5.2]{5.2} it is enough to prove that when
$Y$ is strictly local, $H^q(X,F)=0$ for every abelian sheaf $F$ on
$X$. We then have $Y=\Spec A$, $X=\Spec B$, where $B$ is an integral
$A$-algebra; writing\pageoriginale $B=\varinjlim B_i$, where $B_i$
runs through the finite-type subalgebras of $B$ (which are in fact
finite over $A$), one has $X=\varprojlim_iX_i$, where
$X_i=\Spec B_i$. By \hyperrefext[VII][VII.cor5.13]{5.13}, it is
enough to prove that $R^qf_*(F_i)=0$ for $q\neq0$, where $F_i$ is an
abelian sheaf on $X_i$; this follows from
\hyperref[VIII.prop5.5]{5.5}. The last assertion of
\hyperref[VIII.cor5.6]{5.6} is proved by the same reduction to
\hyperref[VIII.prop5.5]{5.5}.

\begin{enonce*}{Corollary 5.7}\etiquette[5.7]{VIII.cor5.7}
Let $f:X\to Y$ be an immersion. Then for every sheaf $F$ on $X$, the
canonical morphism
$$
f^*f_*(F)\to F
$$
is an isomorphism.
\end{enonce*}

Since $f$ factors as the composite of an open immersion and a closed
immersion, and \hyperref[VIII.cor5.7]{5.7} is trivial for an open
immersion (IV no.~3), it remains to consider the case where $f$ is a
closed immersion. We must prove that for every $x\in X$, the
corresponding homomorphism
$$
(f^*f_*(F))_{\bar x}\to F_{\bar x}
$$
is bijective. By transitivity of fibers
(\hyperref[VIII.subsec3.4]{3.4}), the left-hand side is precisely the
fiber $f_*(F)_{\bar x}$; it is therefore enough to verify that the
canonical homomorphism
$$
f_*(F)_{\bar x}\to F_{\bar x}
$$
is bijective, which follows at once from
\hyperref[VIII.prop5.5]{5.5}.

\begin{enonce*}[remark]{Remark 5.8}\etiquette[5.8]{VIII.rem5.8}
Using \hyperref[VIII.prop5.5]{5.5} and arguing again as in
\hyperref[VIII.cor5.6]{5.6}, one easily proves that if
$f:X\to Y$ is an integral morphism and $F$ is a sheaf of groups on
$X$ (not necessarily commutative), then $R^1f_*(F)$ is the terminal
sheaf on $Y$ (compare Remark \hyperref[VIII.rem5.3]{5.3}).
\end{enonce*}


% END INCLUDED SGA4_English_current_buildable_source_20260906/source/08/drafts/SGA4_VIII_LINES_805_1011_English.texfrag


% BEGIN INCLUDED SGA4_English_current_buildable_source_20260906/source/08/drafts/SGA4_VIII_LINES_1012_1556_English.texfrag
\section{Supports}\pageoriginale\etiquette[6]{VIII.sec6}

Let $U$ be a Zariski open subset of the scheme $X$. Then $U \in \Ob
X_{\text{\rm ét}}$,
and in fact $U$ is a subobject of the final object $X$ of $X_{\text{\rm
    ét}}$,
hence defines a subobject $\widetilde{U}$ of the final object of
$\widetilde{X_{\text{\rm ét}}}$, \ie an ``open subset'' of the étale
topos $X_{\text{\rm ét}}$
of $X$ (IV \hyperref[IV.def8.3]{8.3}).

\begin{enonce*}{Proposition 6.1}\etiquette[6.1]{VIII.prop6.1}
The preceding map $U \mapsto \overset{\sim}{U}$ is an isomorphism
from the ordered set of open subsets (in the Zariski sense) of $X$
onto the set of open subsets of the étale topos
$\widetilde{X_{\text{\rm ét}}}$.
\end{enonce*}

Since $X_{\text{\rm ét}} \to \widetilde{X_{\text{\rm ét}}}$ is fully
faithful, it follows immediately that the map $U \mapsto \widetilde{U}$
preserves the order structures, \ie $(U \subset V) \Leftrightarrow
(\widetilde{U} \subset \widetilde{V})$; in particular, the preceding
map is injective. To prove that it is surjective, consider a subsheaf
$F$ of the final sheaf $\widetilde{X}$, and consider the objects of
$X_{\text{ét}/ F}$, \ie the schemes $X'$ étale over $X$ such that
$F(X') \neq ∅$. Since $X' \to X$ is étale, it is an open map
(\EGA IV 2.4.6); in particular, its image (in the set-theoretic sense)
$\text{Im}(X')$ is open. Let $U$ be the open subset which is the union
of the $\text{Im}(X')$ $(X'\in \Ob X_{\text{ét} /  F})$. Since the
family of the $X'\to U$ $(X'\in \Ob X_{\text{ét}/F})$ is
surjective, hence covering, we conclude that $U \in \Ob
X_{\text{ét}/F}$, hence
$ X_{\text{ét}/F} =  X_{\text{ét}/U}$, hence $F = U$, \cqfd.

In view of \hyperref[VIII.prop6.1]{6.1}, we may therefore speak
unambiguously of an ``open subset'' of $X$, without specifying whether
we mean this notion in the usual Zariski sense or in the sense of the
étale topology.

\begin{enonce*}{Corollary 6.2}\etiquette[6.2]{VIII.cor6.2}
Let $U$ be an open subset of $X$, and let $j : U \to X$ be the
canonical immersion. Then the functor
$$
j^{\ast} :  \widetilde{X_{\text{\rm ét}}} \to \widetilde{U_{\text{\rm ét}}}
$$
induces\pageoriginale an equivalence of categories
$$
\widetilde{X_{\text{\rm ét}}} / \widetilde{U} \to
\widetilde{U_{\text{\rm ét}}}{\quoi}.
$$
\end{enonce*}

Indeed, one checks immediately that for every site $C$ in which finite
$\varprojlim$ exist, and for every subobject $U$ of the final object
$e$ of $C$, if one considers the functor $j : C \to C_{/U}$ defined by
$j(S) = S \times U$, then $j$ is a morphism of sites and $j^{\ast} :
\widetilde{C} \to \widetilde{C/U}$ induces an equivalence
$\widetilde{C}/\widetilde{U} \to \widetilde{C / U}$. It then suffices
to combine this general fact with the fact that $U_{\text{ét}}$ is
canonically isomorphic to $X_{\text{ét}/U}$.

In pictorial terms, one can express \hyperref[VIII.cor6.2]{6.2} by
saying that the operations of ``restriction to an open subset,'' in
the usual sense for schemes on the one hand and in the sense of topoi
on the other, are compatible. Here is an analogous compatibility for
the operations of ``restriction to a closed subset'':

\begin{enonce*}{Theorem 6.3}\etiquette[6.3]{VIII.thm6.3}
Let $X$ be a scheme, let $Y$ be a closed subscheme of $X$, let $U
= X - Y$ be endowed with the induced structure, and let $i : Y \to X$
and $j : U \to X$ be the canonical immersions. Then the functor
$$
i_{\ast} : \widetilde{Y_{\text{\rm ét}}} \to \widetilde{X_{\text{\rm ét}}}
$$
is fully faithful, and if $F \in \Ob \widetilde{X_{\text{\rm ét}}}$,
then $F$ is isomorphic to a sheaf of the form $i_{\ast} (G)$ if and
only if $j^{\ast}(F)$ is isomorphic to the final sheaf
$\widetilde{U}$ on $U$.
\end{enonce*}

Proof. Since $i_{\ast}$ and $i^{\ast}$ are adjoint to one another, the
fact that $i_{\ast}$ is fully faithful is also equivalent to the fact
that the functorial homomorphism
$$
i^{\ast} (i_{\ast}(G)) \to G
$$
is an\pageoriginale isomorphism, which is precisely
\hyperref[VIII.cor5.7]{5.7}. On the other hand, if $G \in \Ob
\widetilde{Y_{\text{ét}}}$, then one checks trivially, using
\hyperref[VIII.cor6.2]{6.2}, that $j^{\ast} (i_{\ast}(G))$ is the final
sheaf on $U$. Conversely, if $F \in \Ob
\widetilde{X_{\text{ét}}}$ is such that $j^{\ast}(F)$ is the final
sheaf, let us prove that $F$ is isomorphic to a sheaf of the form
$i_{\ast}(G)$, or, what amounts to the same thing, that the canonical
homomorphism
$$
F \to i_{\ast} i^{\ast} F
$$
is an isomorphism. It is enough once again to check that for every $x
\in X$, the homomorphism induced on the geometric fibers at $\bar{x}$
is an isomorphism. When $x \in U$, this is precisely the hypothesis
made on $F$. When $x \in Y$, by transitivity of fibers we are reduced
to checking that the homomorphism on the fibers at $\bar{x}$ induced
by
$$
i^{\ast} (F) \to i^{\ast} (i_{\ast} i^{\ast} (F))
$$
is an isomorphism; but the preceding homomorphism is an isomorphism by
\hyperref[VIII.cor5.7]{5.7}, applied to $F = i^{\ast} F$ and to $i$.
This completes the proof of \hyperref[VIII.thm6.3]{6.3}.

\begin{enonce*}{Corollary 6.4}\etiquette[6.4]{VIII.cor6.4}
The functor $i_{\ast}$ induces an equivalence from the category of
abelian sheaves on $Y$ to the category of abelian sheaves on $X$ whose
restriction to $U = X - Y$ is zero.
\end{enonce*}

In accordance with common usage, we shall therefore often identify an
abelian sheaf on $Y$ with an abelian sheaf on $X$ which is zero on $U$.

\setcounter{subsection}{4}
\subsection{}\etiquette[6.5]{VIII.subsec6.5}

By virtue of \hyperref[VIII.prop6.1]{6.1}, we thus know (if $X$ is a
scheme) how to interpret the ``open subsets'' of the topos $ℰ =
\widetilde{X_{\text{ét}}}$ as open subsets $U$ of $X$ in the usual
sense, and by virtue of \hyperref[VIII.thm6.3]{6.3}, under this
identification the ``residual topos''\pageoriginale $ℰ_\complement U$
of IV 3 is canonically equivalent to the topos
$\widetilde{Y_{\text{ét}}}$, where $Y = X - U$ is endowed with any
induced structure making $Y$ a scheme. We may consequently apply to
this situation the results of IV 3, in particular the description IV
\hyperref[IV.theo9.5.4]{3.3} of sheaves $F$ on $X$ in terms of triples
$(F', F'', u)$, where $F'$ is a sheaf on $Y$, $F''$ is a sheaf on $U$,
and $u$ is a homomorphism from $F'$ to $i^{\ast} j_{\ast} (F'')$. In
what follows we shall freely use the notation $i^{!}, j_{!}$ of IV.3,
which denotes functors
\begin{align*}
j_{!} : \widetilde{(U_{\text{ét}})}_{\ab} &\to
\widetilde{(X_{\text{ét}})}_{\ab}{\quoi},\\
i^{!} : \widetilde{(X_{\text{ét}})}_{\ab} &\to
\widetilde{(Y_{\text{ét}})}_{\ab}{\quoi},
\end{align*}
where the subscript ``$\ab$'' denotes the category of abelian sheaves.
These functors give rise to the two exact sequences IV \hyperref[IV.prop14.6]{3.7}.

\subsection{}\etiquette[6.6]{VIII.subsec6.6}

In view of the preceding developments and of the general terminology
introduced in IV \hyperref[IV.subsec8.5]{8.5}, it is appropriate, for an abelian sheaf $F$ on
$X$, or a section $\varphi$ of such a sheaf, to introduce the notion of
the \emph{support of} $F$, \resp of $\varphi$, as the closed subset (in
the usual, \ie Zariski, sense) complementary to the largest open subset
on which the object in question vanishes.

In the present situation, one should also replace the general notation
\hyperrefext[V][V.prop4.3]{4.3} $H^q (ℰ_{\complement U} F)$,
$ℋ^q_{\complement U} (F)$ by the notation $H^{q}_Y (X, F)$,
$ℋ^{q}_Y (F)$, the first denoting an abelian group and the second an
abelian sheaf on $Y$ (or equivalently an abelian sheaf on $X$ which is
zero on $U$). Thus, one has $\mathcal H^q_Y(F)=R^q i^!(F)$, \etc

We refer to \hyperrefext[V][V.sec6]{6} for the general properties of
the preceding functors.

\section{Specialization morphisms of fiber
  functors}\pageoriginale
\etiquette[7]{VIII.sec7}

\subsection{}\etiquette[7.1]{VIII.subsec7.1}

We saw in ${\rm N^{\circ}}$ 4 how one associates, with every geometric
point $\xi$ of the scheme $X$, a strictly local $X$-scheme
$$
\bar{X}(\xi) = \Spec \Oo_{X, \xi}{\quoi},
$$
which in fact depends only on the geometric point above the image $x$
of $\xi$ defined by the separable closure $\bar{k(x)}$ of $k(x)$ in
$\Omega = k(\xi)$. In what follows we shall often restrict ourselves
to geometric points $\xi = \Spec \Omega$ which are \emph{algebraic
  separable} over $X$, \ie such that $\Omega$ is a separable closure
of $k(x)$, \ie $\xi = \bar{x}$. Let us call an $X$-scheme $Z$ a
\emph{strict localization of} $X$ if it is $X$-isomorphic to a scheme
of the form $\Spec \Oo_{X, \bar{x}}$; thus one sees that
$$
\xi \mapsto \Spec \sheaf{\Oo}_{X, \xi} = \bar{X}(\xi)
$$
is an equivalence from the category of geometric points
\emph{algebraic separable} over $X$ to the category of $X$-schemes
which are strict localizations of $X$, when only isomorphisms are taken
as morphisms in both categories.

The last assertion is no longer correct when all $X$-morphisms are
taken as morphisms, since in the second category there may be
$X$-morphisms which are not isomorphisms. We then make the following

\begin{enonce*}{Definition 7.2}\etiquette[7.2]{VIII.def7.2}
Let $\xi$, $\xi'$ be two geometric points of the scheme $X$. A
specialization arrow from $\xi'$ to $\xi$ is any $X$-morphism between
the corresponding strict localizations
$$
\bar{X} (\xi') \to \bar{X} (\xi).
$$
We\pageoriginale say that $\xi$ is a specialization of $\xi'$, or
that $\xi'$ is a generization of $\xi$, if there exists a
specialization arrow from $\xi'$ to $\xi$.
\end{enonce*}

Note that specialization arrows compose in the evident way, so that,
with specialization arrows as morphisms, \emph{the geometric points of
  $X$ form a category}, equivalent (as is the full subcategory formed
by the geometric points algebraic and separable over $X$) to the full
subcategory of $(\Sch)_{/X}$ formed by the strict localizations of $X$.

\begin{enonce*}{Lemma 7.3}\etiquette[7.3]{VIII.lem7.3}
Let $X$ be a scheme, let $Z$ be an $X$-scheme which is isomorphic to a
pseudo-filtered projective limit of $X$-schemes $X_{i}$ which are étale,
with affine transition morphisms
\emph{(\hyperrefext[VII][VII.subsec5.1]{5.1})}, let $\xi'$ be a
geometric point of $X$, and let $Z' = \bar{X} (\xi')$ be the
corresponding strict localization. Then:
\begin{enumerate}
\renewcommand{\theenumi}{\alph{enumi}}
\renewcommand{\labelenumi}{\theenumi)}
\item The restriction map
$$
\Hom_X (Z', Z) \to \Hom_X(\xi', Z)
$$
is bijective.

\item In order that the two sides be nonempty, it is necessary and
  sufficient that the image $x'$ of $\xi'$ in $X$ lie in the image of
  $Z$.

\item Let $T$ be a $Z$-scheme. In order that $T$ be a strict
  localization of $Z$, it is necessary and sufficient that it be a
  strict localization of $X$.
\end{enumerate}
\end{enonce*}

Proof.
\begin{enumerate}
\renewcommand{\theenumi}{\alph{enumi}}
\renewcommand{\labelenumi}{\theenumi)}
\item We are immediately reduced to the case where $Z$ is one of the
  $X_i$, \ie where $Z$ is étale over $X$, and by the base change
  $Z'\to X$ we may suppose that $Z'\xrightarrow{\sim} X$, whence the
  conclusion follows from \hyperref[VIII.def4.2]{4.2} (ii).

\item Note that if $x'$ is the image of a point $z$ of $Z$, then
  $k(z)$ is necessarily a separable algebraic extension of $k(x')$;
  hence there exists a $k(x')$-homomorphism from the latter into
  $k(\xi')$, whence the conclusion.

\item The proof\pageoriginale is an easy exercise left to the reader.
\end{enumerate}

From \hyperref[VIII.def7.2]{7.2} and \hyperref[VIII.lem7.3]{7.3} we
conclude:

\begin{enonce*}{Proposition 7.4}\etiquette[7.4]{VIII.prop7.4}
Let $\xi$, $\xi'$ be two geometric points of the scheme $X$. Then the
restriction map defines a bijection from the set
$\Hom_{X} (\bar{X}(\xi'), \bar{X}(\xi))$ of specialization arrows from
$\xi'$ to $\xi$ to the set of $X$-morphisms from $\xi'$ to
$\bar{X}(\xi)$.
\end{enonce*}

\begin{enonce*}{Corollary 7.5}\etiquette[7.5]{VIII.cor7.5}
In order that $\xi$ be a specialization of $\xi'$, it is necessary and
sufficient that the same hold for the images $x$, $x'$ of $\xi$,
$\xi'$ in $X$, \ie that $x$ belong to the closure of $\{x'\}$.
\end{enonce*}

Since $\Spec \Oo_{X, \xi} \to \Spec \Oo_{X, x}$ is faithfully flat
(\hyperref[VIII.subsec4.4]{4.4}), it is surjective, and it is therefore
enough to apply the second assertion of \hyperref[VIII.lem7.3]{7.3}.

\begin{enonce*}{Corollary 7.6}\etiquette[7.6]{VIII.cor7.6}
For every scheme $X$, let $\sheaf{\Pt}(X)$ be the category of geometric
points over $X$ (or equivalently of geometric points algebraic
separable over $X$), the morphisms being specialization arrows. Let
$\xi$ be a geometric point of a scheme $X$, and let $\bar{X}(\xi)$ be
the corresponding strict localization. Then there is an equivalence of
categories
$$
\sheaf{\Pt}(\bar{X}(\xi)) \xrightarrow{\sim} \sheaf{\Pt}(X)_{/\xi}{\quoi},
$$
obtained by associating with every geometric point $\xi'$ of
$\bar{X}(\xi)$ the corresponding geometric point over $X$, together
with the specialization morphism to $\xi$ deduced from the structural
morphism $\xi' \to \bar{X}(\xi)$ by \hyperref[VIII.prop7.4]{7.4}.
\end{enonce*}

In other words, \emph{the datum of a generization $\xi'$ of the
geometric point $\xi$ is essentially equivalent to the datum of a
geometric point of $\bar{X}(\xi)$}. Note moreover that, by
\hyperref[VIII.lem7.3]{7.3} c), the corresponding
homomorphism\pageoriginale $\bar{X}(\xi') \to \bar{X}(\xi)$ makes
$\bar{X}(\xi')$ the strict localization of $\bar{X}(\xi)$ relative to
$\xi'$.

\setcounter{subsection}{6}
\subsection{}\etiquette[7.7]{VIII.subsec7.7}

Now interpret, for every geometric point $\xi$ of $X$ and every sheaf
$F$ on $X$, the fiber $F_{\xi}$ as $\Gamma (\bar{X} (\xi),
\bar{F}(\xi))$, where $\bar{F}(\xi)$ is the inverse image of $F$ on
$\bar{X}(\xi)$ (\hyperref[VIII.cor4.8]{4.8}). Then one sees that every
specialization arrow
$$
u : \xi'\to \xi
$$
induces a homomorphism, functorial in $F$:
$$
u^{\ast} : F_{\xi} \to F_{\xi'}{\quoi},
$$
called the \emph{specialization homomorphism} associated with the
specialization arrow $u$. It is evident from the transitivity of
inverse images of sheaves that, for a composite specialization arrow,
one has:
$$
(wu)^{\ast} = u^{\ast} w^{\ast}{\quoi}.
$$

\subsection{}\etiquette[7.8]{VIII.subsec7.8}

% SGA-SRC-000768: the scan supplies no exact citation; reconstructed from the mathematics.
If $ℰ$ is a topos, recall from Exposé IV, 6.1 and 6.2,
that a ``fiber functor,'' or (by abuse of language) a ``point'' of the
topos $ℰ$, is any morphism from the ``final topos'' $(\Ens)$
(isomorphic to the category of sheaves on a space reduced to one
point!) to $ℰ$, \ie any functor
$$
\varphi : ℰ \to (\Ens)
$$
which commutes with arbitrary inductive \sisi{lim.}{limits}, and with
finite projective \sisi{lim.}{limits}. The set of fiber functors of
$ℰ$ should be regarded as the set of objects of a full subcategory of
$\Hhom(ℰ, (\Ens))$, called the \emph{category of fiber functors} of the
topos $ℰ$, whose opposite is called the \emph{category of points of
  $ℰ$}, and is denoted $\sheaf{\Pt}(ℰ)$, \cf (IV \hyperref[IV.def6.1]{6.1}).

When $ℰ$ is of the form $\widetilde{X_{\text{ét}}}$, where $X$ is
a scheme, we have\pageoriginale defined in
\hyperref[VIII.def3.3]{3.3} and \hyperref[VIII.subsec7.7]{7.7} a
functor
\begin{equation*}
\sheaf{\Pt} (X) \to \sheaf{\Pt} (\widetilde{X_{\text{ét}}}){\quoi},\tag{$\ast$}
\end{equation*}
where the first member is defined in \hyperref[VIII.cor7.6]{7.6}. With
this understood, we have the following

\begin{enonce*}{Theorem 7.9}\etiquette[7.9]{VIII.thm7.9}
Let $X$ be a scheme. The preceding functor $(\ast)$ is an equivalence
from the category of geometric points over $X$ (with specialization
arrows as morphisms) to the category of points of the étale topos
$\widetilde{X_{\text{\rm ét}}}$ (the \emph{opposite} of that of fiber
functors on $\widetilde{X_{\text{\rm ét}}}$).
\end{enonce*}

Since this theorem will not be used later in the seminar, we confine
ourselves here to a sketch of the proof, in which we shall allow
ourselves certain liberties with questions of universes.
\begin{enumerate}
\renewcommand{\theenumi}{\alph{enumi}}
\renewcommand{\labelenumi}{\theenumi)}
\item \emph{The functor under consideration is fully faithful}. For
  every $X \in \Ob X_{\text{ét}}$, let $\Xv :
  \widetilde{X_{\text{ét}}} \to (\Ens)$ be defined by
$$
\Xv (F) = \Hom(X, F) = F(X){\quoi}.
$$
Note that the fiber functor
$$
\mathcal{E}_{\xi} : F \mapsto F_{\xi}
$$
can be written
$$
\mathcal{E}_{\xi} \simeq \varinjlim_i X^{\vee}_i{\quoi},
$$
where the $X_i$ are affine schemes étale over $X$, indexed by a
certain filtered category (\cf \hyperref[VIII.def4.3]{4.3} and
\hyperref[VIII.subsec4.5]{4.5}). Hence, for every functor $\varphi :
\widetilde{X_{\text{ét}}} \to (\Ens)$, one has:
$$
\Hom(\mathcal{E}_{\xi}, \varphi) \simeq \varprojlim_{i}
\Hom(X^{\vee}_i, \varphi){\quoi},
$$
and on the other hand it is well known that there is a bifunctorial
isomorphism
$$
\Hom(X^{\vee}_i, \varphi) \simeq \varphi (X_i){\quoi}.
$$

When\pageoriginale $\varphi$ is of the form $\Ee_{\xi'}$, hence
$$
\varphi \simeq \varinjlim_j {X'}_j^{\vee}{\quoi},
$$
there is therefore a natural bijection
\begin{equation*}
\Hom (\mathcal{E}_{\xi}, \mathcal{E}_{\xi'}) \simeq
\varprojlim_{i} \varinjlim_{j} \Hom_X (X'_j, X_i){\quoi}.\tag{$\ast$}
\end{equation*}
On the other hand, one has (in the category of schemes)
$$
\bar{X}(\xi) = \varprojlim_{i} X_i,\quad \bar{X} (\xi') =
\varprojlim_{j} X'_j{\quoi},
$$
whence
$$
\Hom_{X} (\bar{X} (\xi'), \bar{X}(\xi)) \simeq \varprojlim_{i}
\Hom_X (\bar{X} (\xi'), X_i){\quoi},
$$
and on the other hand, since $X_i$ is locally of finite presentation
over $X$, one has
$$
\Hom_X (\bar{X} (\xi'), X_i) \simeq \varinjlim_{j}
\Hom_X (X'_j, X_i){\quoi},
$$
whence
\begin{equation*}
\Hom_X (\bar{X} (\xi'), \bar{X} (\xi) ) \simeq
\varprojlim_{i} \varinjlim_{j} \Hom_X (X'_j,
X_i){\quoi}.\tag{${\ast}{\ast}$}
\end{equation*}
Comparison of $(\ast)$ and $({\ast}{\ast})$ gives the desired
conclusion (subject to checking compatibilities, which is left to the
reader).

\item \emph{The functor under consideration is essentially
  surjective}.
\end{enumerate}

\setcounter{subsection}{9}
\subsubsection{}\etiquette[7.9.1]{VIII.subsubsec7.9.1}

First note that if $C$ is a site in which finite $\varprojlim$ are
representable, and $\widetilde{C} = ℰ$ is the corresponding topos, then
every fiber functor $\varphi$ on $ℰ$ can be represented as a ``filtered
inductive limit'' of functors of the form
$$
\check{Y} (F) = F(Y) = \Hom(\widetilde{Y}, F), Y \in \Ob C{\quoi},
$$
(where $\widetilde{Y}$ is the sheaf associated with $Y$), as follows
\footnote{This is a special case of IV \hyperref[IV.eq6.8.3]{6.8.3}.}. Consider the
category\pageoriginale $C_{/\varphi}$ of pairs $(Y, \xi)$, with
$Y \in \Ob C$, $\xi \in \varphi (\widetilde{Y})$ (the morphisms being
defined in the evident way), and note that there is an evident
homomorphism
\begin{equation*}
\varinjlim_{(Y, \xi)\in C^0_{/\varphi}} \check{Y} \to
\varphi\tag{$\ast$}
\end{equation*}
by observing that
$$
\Hom (\varinjlim_{C^0_{/\varphi}} \check{Y},\varphi) \simeq
\varprojlim_{C_{/ \varphi}} \Hom(\check{Y}, \varphi) \simeq
\varprojlim_{C_{/ \varphi}} \varphi (\widetilde{Y}){\quoi}.
$$
There is a canonical element in $\varprojlim_{C_{/
    \varphi}} \varphi(\widetilde{Y}),$ obtained by associating with
every $(Y, \xi)$ the element $\xi$ of $\varphi (\widetilde{Y})$. Using the fact
that finite $\varprojlim$ exist in $C$, and that the functor $Y
\mapsto \varphi(\widetilde{Y})$ commutes with them, one sees that in
$C_{/\varphi}$ finite $\varprojlim$ exist; a fortiori
$C^0_{/\varphi}$ is filtered. Using the fact that every sheaf $F$ is
an inductive limit of sheaves of the form $\widetilde{Y}$, and using
the fact that $\varphi$ commutes with inductive limits, one easily
concludes that the homomorphism of functors $(\ast)$ is bijective and
thus gives the asserted representation.

\subsubsection{}\etiquette[7.9.2]{VIII.subsubsec7.9.2}

Note also that if $ℰ$ is a topos and $Y$ is an object of $ℰ$, hence
defines a topos $ℰ_{/Y}$, then the datum of a fiber functor $\varphi$
for $ℰ_{/Y}$ is equivalent to the datum of a pair $(\varphi, \xi)$,
where $\varphi$ is a fiber functor for $ℰ$, and $\xi \in \varphi(Y)$. With
the fiber functor $\psi : ℰ_{/Y} \to (\Ens)$ one associates the pair
$(\varphi, \xi)$, where $\varphi$ is the composite $Z \mapsto
\varphi(Z) = \psi (Z \times Y)$, and where $\xi \in \varphi(Y) =
\psi(Y \times Y)$ is the image of the unique element of $\psi(Y)$
under the diagonal morphism of $Y \times Y$.

\subsubsection{}\etiquette[7.9.3]{VIII.subsubsec7.9.3}

We now return to the case of the site $C =
X_{\text{ét}}$, and let $\varphi :
\widetilde{X_{\text{ét}}} \to (\Ens)$ be a fiber functor on the
étale topos of $X$. Considering\pageoriginale the restriction of
$\varphi$ to the full subcategory of
$\widetilde{X_{\text{ét}}}$ formed by the open subsets of this
topos, or, what amounts to the same thing
(\hyperref[VIII.prop6.1]{6.1}), by the open subsets $U$ of $X$, one
obtains a map
$$
\varphi_0 : \Uu (X) \to \{0, 1\}
$$
which commutes with arbitrary suprema and finite infima. (Note that for
$U \in \Uu (X)$, $\varphi(U)$ is either empty or consists of a single
point, and one takes $\varphi_{0} (U) = 0$ or 1 according as one is in
the former or the latter case.)

One easily sees, using the fact that every irreducible closed subset
of $X$ has one and only one generic point \footnote{\ie $X$ is
\emph{sober} in the terminology of IV \hyperref[IV.subsubsec4.2.1]{4.2.1}. The assertion made is a
special case of IV \hyperref[IV.subsubsec4.2.3]{4.2.3}.}, that $\varphi_0$ is defined by means of a
unique $x \in X$, by the condition
$$
\varphi_0 (U) = 1\quad \Longleftrightarrow\quad x \in U {\quoi},
$$
\ie
$$
\varphi (U) \ne\varnothing \Leftrightarrow  x \in U {\quoi}.
$$
Let $F\in \Ob \widetilde{X_{\text{ét}}}$. Then the image of $F$ in
the final sheaf $X$ is an open subset; and since
$$
\varphi (F) \to \varphi(U)
$$
is surjective, one sees that $\varphi(F) \neq ∅$ if and only if
$\varphi(U) \neq ∅$, \ie if and only if $x \in U$. Let
$C_{/\varphi}$ be the category of pairs $(X', \xi)$ where $X'\in \Ob
X_{\text{ét}}$, and $\xi \in \varphi (X')$. It was pointed out in
\hyperref[VIII.subsubsec7.9.1]{7.9.1} that $C^{0}_{/ \varphi}$ is
filtered and connected. Moreover, by
\hyperref[VIII.subsubsec7.9.2]{7.9.2} and by the preceding discussion
applied to $X'$ instead of $X$, if $(X', \xi) \in \Ob
C_{/\varphi}$, there is associated with it a unique point $x' \in
X'$, such that for an étale morphism $X'' \to X'$, $\xi$ lies in the
image of $\varphi (X'') \to \varphi(X')$ if and only if $x'$ lies in the
image of $X''$. It follows that the scheme $\varprojlim$ of the $X'$
over the category $C$ of the $(X', \xi)$ exists and is a strict
localization $Z$ of $X$, thus corresponding to a geometric point $\xi$
and a fiber functor $\mathcal{E}_{\xi}$. For\pageoriginale every sheaf
$F$, one then has
$$
\mathcal{E}_{\xi} (F) \simeq \varinjlim_{C^0_{/\varphi}} F(X'){\quoi}.
$$
In view of \hyperref[VIII.subsubsec7.9.1]{7.9.1}, we conclude that
$\varphi$ is isomorphic to $\mathcal{E}_{\xi}$, which completes the
proof of \hyperref[VIII.thm7.9]{7.9}.

% END INCLUDED SGA4_English_current_buildable_source_20260906/source/08/drafts/SGA4_VIII_LINES_1012_1556_English.texfrag


% BEGIN INCLUDED SGA4_English_current_buildable_source_20260906/source/08/drafts/SGA4_VIII_LINES_1557_1976_English.texfrag
\section{Two spectral sequences for integral morphisms}\etiquette[8]{VIII.sec8}

\begin{enonce*}{Proposition 8.1}\etiquette[8.1]{VIII.prop8.1}
Let $f : X \to Y$ be a surjective integral morphism, and let $F$ be an
abelian sheaf on $Y$. For every integer $i$, let $\Hh^{i} (F)$ be the
presheaf on $(\Sch)_{/Y}$ defined by
$$
\Hh^{i} (F) (Z) = H^i (Z, F_Z){\quoi},
$$
where $F_Z$ is the inverse image of $F$ on $Z$. Then there exists a
spectral sequence (functorial in $F$)
$$
H^{\ast} (X, F) \Longleftarrow E^{pq}_2 = H^p (X /Y, \Hh^q (F))
$$
(``descent spectral sequence'').
\end{enonce*}

Of course, the symbol $H^p (X / Y, G)$, for a presheaf $G$ on
$(\Sch)_{/Y}$, denotes the $p$-th relative \v{C}ech cohomology group,
defined by means of the complex of the $C^n(X / Y, G)
= G(X^{n+1})$, where $X^m$ denotes the $m$-th power in
$(\Sch)_{/Y}$. To prove \hyperref[VIII.prop8.1]{8.1}, let $p_n :
X^{n+1} \to Y$ be the projection, and set
$$
A^n = p_{n\ast} (𝒵_{X^{n + 1}}){\quoi},
$$
where $𝒵_{X^{n + 1}}$ denotes the constant sheaf $𝒵$ on
$X^{n+1}$. Then the $A^n$ are the components of a simplicial abelian
sheaf on $Y$, hence of a complex of abelian sheaves $A^{\ast}$ on $Y$.
Note that there is an evident homomorphism
\begin{equation*}
𝒵_Y \to A^{0}{\quoi}.\tag{$\ast$}
\end{equation*}
With this\pageoriginale understood, we have the following

\begin{enonce*}{Lemma 8.2}\etiquette[8.2]{VIII.lem8.2}
The complex $(A^{\ast})$, endowed with the homomorphism $(\ast)$, is a
resolution of $𝒵_Y$. More generally, for every abelian sheaf $F$ on
$Y$, $A^{\ast} \otimes_{𝒵}F$ is a resolution of $F$.
\end{enonce*}

\noindent
Proof: We may suppose that $Y$ is affine, hence $Y = \Spec A$, $X
= \Spec B$, where $B$ is an integral $A$-algebra. We have $B =
\varinjlim_i B_i$, where $B_i$ ranges over the finitely generated, hence
finite, subalgebras of $B$, whence $X=\varprojlim X_{i}$, where
$X_{i}=\Spec B_{i}$. Using \hyperrefext[VII][VII.cor5.11]{5.11}, one
sees that the augmented complex $A^{\ast}(X/Y)$ is then the inductive
limit of the augmented complexes $A^{\ast}(X_{i}/Y)$. This reduces us
to the case where $f$ is \emph{finite}. It suffices to prove that for
every geometric point $\bar{y}$ of $Y$, the complex $A^{\ast}_{\bar y}$ is
a resolution of $𝒵_{\bar{y}}$, and likewise after tensoring with an
$F$. But if $X_{\bar{y}}$ is the fiber of $X$ at $\bar{y}$, it follows
from \hyperref[VIII.prop5.5]{5.5} that the complex
$A^{\ast}_{\bar{y}}$ is none other than the analogous complex
$A^{\ast}(X_{\bar{y}}/\bar{y})$. This reduces us to the case where $Y$
is the spectrum of a separably closed field $k$. Using
\hyperref[VIII.exem1.3]{1.3} a) and b), we may even suppose that $k$ is
algebraically closed and $X$ is reduced, hence that $X$ is of the form
$I_{Y}$, where $I$ is a finite set. But then the complex
$A^{\ast}\otimes F$ is identified with the trivial cochain complex of
the index set $I$ with coefficients in $F_{\bar{y}}$, and is therefore
indeed a resolution of $F_{\bar{y}}$.

Thus, as a consequence of \hyperref[VIII.lem8.2]{8.2}, we obtain a
spectral sequence
$$
H^{\ast}(Y,F)\Longleftarrow E^{pq}_{2}=H^{p}(H^{q}(Y,A^{\ast}\otimes
F)){\quoi},
$$
and it remains to make the initial term explicit. There is a canonical
homomorphism\pageoriginale
$$
A^{n}\otimes F\longrightarrow p_{n_{\ast}}p^{\ast}_{n}(F){\quoi},
$$
and the latter is an isomorphism, as one sees once again by reduction
to the case where $f$ is finite and passage to fibers. Using
\hyperref[VIII.prop5.5]{5.5}, one concludes that
$$
H^{q}(Y,A^{n}\otimes F)\simeq
H^{q}(X^{n+1},p^{\ast}_{n}(F))=
\Hh^{q}(F)(X^{n+1}){\quoi},
$$
which indeed gives the initial term stated in
\hyperref[VIII.prop8.1]{8.1}.

\begin{enonce*}[remark]{Remark 8.3}\etiquette[8.3]{VIII.rem8.3}
\begin{enumerate}
\renewcommand{\theenumi}{\alph{enumi}}
\renewcommand{\labelenumi}{\theenumi)}
\item When one is given a \emph{locally finite} covering of $Y$ by
  closed subsets $Y_{i}$, and one sets $X=\coprod_{i}Y_{i}$, the
  canonical morphism $f:X\rightarrow Y$ is finite, and the initial
  term of the spectral sequence \hyperref[VIII.prop8.1]{8.1} can be
  made explicit as the cohomology of the complex defined by the
  $H^{q}(Y_{i_{0}\ldots i_{p}},F|Y_{i_{0}\ldots i_{p}})$, where
  $Y_{i_{0}\ldots i_{p}}=Y_{i_{0}}\cap\ldots\cap Y_{i_{p}}$. This is
  therefore the analogue of the Leray spectral sequence for a locally
  finite closed covering of an ordinary topological space (TF Chapter
  II 5.2.4). The latter can moreover also be generalized to a spectral
  sequence relative to a ``finite'' morphism, \ie one which is proper
  with finite fibers, by proceeding as in
  \hyperref[VIII.prop8.1]{8.1}.

\item Note the analogy between the spectral sequence of
  \hyperref[VIII.prop8.1]{8.1} and the Leray spectral sequence of a
  covering $(X_{i})$ of $Y$ (in this case by $X_{i}$ étale over $Y$);
  the latter would formally be obtained by writing the spectral
  sequence \hyperref[VIII.prop8.1]{8.1} for
  $X=\coprod_{i}X_{i}$. It is in fact likely that these two spectral
  sequences admit a common generalization, valid whenever one has a
  family of morphisms $(X_{i}\rightarrow Y)$ which is a ``universally
  effective descent family'' for the fibered category of étale sheaves
  over a variable base scheme\pageoriginale (\cf ${\rm n^{\circ}}$ 9
  below). The analogous question also arises in ordinary topology,
  concerning a common generalization of the two kinds of Leray
  spectral sequences of a covering, assumed either open or locally
  finite closed \footnote[1]{Since these lines were written, P.~Deligne
  has developed a general theory of ``descent spectral sequences,''
  \cf \Exp V bis.}.
\end{enumerate}
\end{enonce*}

\begin{enonce*}{Proposition 8.4}\etiquette[8.4]{VIII.prop8.4}
Let $Y$ be a scheme, let $\pi$ be a profinite group, let
$(\pi_{i})_{i}$ be the projective system of finite discrete quotient
groups of $\pi$, let $(X_{i})_{i\in I}$ be a projective system of
principal coverings of $Y$, with groups $\pi_{i}$, the transition
homomorphisms $X_{j}\rightarrow X_{i}$ being compatible with the
homomorphisms $\pi_{j}\rightarrow \pi_{i}$ on the groups of operators,
and let $X=\varprojlim X_{i}$
\emph{(\cf \hyperrefext[VII][VII.subsec5.1]{5.1})}, so that the group
$\pi$ acts on the $Y$-scheme $X$. Let $F$ be an abelian sheaf on $Y$.
Then there is a ``\emph{Hochschild--Serre}'' spectral sequence
(functorial in $F$)
$$
H^{\ast}(Y,F)\Longleftarrow
E^{pq}_{2}=H^{p}(\pi,\varinjlim_{i}H^{q}(X_{i},F_{X_{i}})){\quoi},
$$
where $F_{X_{i}}$ is the inverse image of $F$ on $X_{i}$, and
$H^{p}(\pi,-)$ denotes Galois cohomology.
\end{enonce*}

The term $E^{pq}_{2}$ written here is also (by the definition of
$H^{p}(\pi,-)$ and the transitivity of inductive limits) isomorphic to
$$
E^{pq}_{2}=\varinjlim_{i}
H^{p}(\pi_{i},H^{q}(X_{i},F_{X_{i}})){\quoi},
$$
which shows that it is enough to find an inductive system of spectral
sequences (depending on the index $i$)
$$
H^{\ast}(Y,F)\Longleftarrow
{}^{i}E^{pq}_{2}=H^{p}(\pi_{i},H^{q}(X_{i},F_{X_{i}})){\quoi}.
$$
This reduces us to defining the spectral sequence in the case where
$\pi$ is finite. Then $f:X\rightarrow Y$ is a covering morphism in
$Y_{\text{\rm ét}}$, so one may write\pageoriginale the Leray
spectral sequence of this morphism. In view of the canonical
isomorphisms
$$
X^{n}\simeq X\mathop×^{\pi}\pi^{n}{\quoi},
$$
a well-known calculation then shows that for every presheaf $\Hh$ on
$X_{\text{\rm ét}}$ which transforms sums into products,
$C^{\ast}(X/Y,\Hh)$ is none other than the complex of homogeneous
cochains of $\pi$ with coefficients in $\Hh(X)$, whence the stated form
of the initial term.

A variant of this proof, avoiding any calculation, consists in
% SGA-SRC-000769: corrupt source notation reconstructed as a contracted product.
regarding $E\mapsto X\mathop{\times}^{\pi}E$ as a morphism from the site of
finite $\pi$-sets to the étale site of $Y$, and writing the Leray
spectral sequence of this morphism. Finally, when $\pi$ is finite, one
may also regard the Hochschild--Serre spectral sequence as a special
case of \hyperref[VIII.prop8.1]{8.1} relative to the morphism
$f:X\rightarrow Y$.

\begin{enonce*}{Corollary 8.5}\etiquette[8.5]{VIII.cor8.5}
Suppose that $Y$ is quasi-compact and quasi-separated. Then the spectral
sequence of \hyperref[VIII.prop8.4]{8.4} takes the form
$$
H^{\ast}(Y,F)\Longleftarrow E^{pq}_{2}=H^{p}(\pi,
H^{q}(X,F_{X})){\quoi}.
$$
\end{enonce*}

Indeed, by \hyperrefext[VII][VII.cor5.8]{5.8}, there are then canonical
isomorphisms
$$
\varinjlim_{i} H^{q}(X_{i},F_{X_{i}})\simeq H^{q}(X,F_{X}){\quoi}.
$$

\begin{enonce*}{Corollary 8.6}\etiquette[8.6]{VIII.cor8.6}
Let $Y$ be a local henselian scheme with closed point $y$, let $F$ be
an abelian sheaf on $Y$, and let $F_{0}$ be the induced sheaf on
$Y_{0}=\Spec (k(y))$. Then the canonical homomorphisms
$$
H^{n}(Y,F)\longrightarrow H^{n}(Y_{0},F_{0})
$$
are\pageoriginale isomorphisms.
\end{enonce*}

Indeed, let $\bar{y}$ be a geometric point over $y$, corresponding to
a separable closure $\bar{k(y)}=k(\bar{y})$ of $k(y)$; since $Y$ is
henselian, the strict localization $X$ of $Y$ at $\bar{y}$ is the
projective limit of the connected étale Galois coverings pointed at
$\bar{y}$, denoted $X_{i}$, so that the hypotheses of
\hyperref[VIII.cor8.5]{8.5} apply. Since $H^{q}(X,F_{X})=0$ for $q>0$
by \hyperref[VIII.cor4.7]{4.7}, one obtains isomorphisms
$$
H^{n}(Y,F)\xleftarrow{\sim} H^{n}(\pi, F(X)){\quoi},
$$
where $\pi$ is the ``Galois group'' of $X$ over $Y$, isomorphic to that
of $\bar{k(y)}$ over $k=k(y)$. Similarly (or by
\hyperref[VIII.cor2.3]{2.3}) one finds
$$
H^{n}(Y_{0},F_{0})\xleftarrow{\sim} H^{n}(\pi,F_{0}(X_{0})){\quoi},
$$
where $X_{0}=\fprod{X}{Y}{Y_{0}}\simeq \Spec (k(\bar{y}))$. But the
restriction homomorphism $F(X)\rightarrow F_{0}(X_{0})$ is an
isomorphism by \hyperref[VIII.cor4.8]{4.8}, from which the conclusion
\hyperref[VIII.cor8.6]{8.6} follows immediately.

\section{Descent of étale sheaves}\etiquette[9]{VIII.sec9}

This section will not be used later in this Seminar, and may be omitted
on a first reading.

\begin{enonce*}{Proposition 9.1}\etiquette[9.1]{VIII.prop9.1}
Let\pageoriginale $f:X'\rightarrow X$ be a surjective (\resp
universally submersive \emph{(\SGA I IX 2.1)}) morphism of schemes.
Then the functor $f^{\ast}:X_{\text{\rm ét}}\rightarrow
X'_{\text{\rm ét}}$ is faithful and ``conservative''
\emph{(\cf \hyperref[VIII.cor3.6]{3.6})} (\resp induces a
\emph{fully faithful} functor from the category of sheaves on
$X_{\text{\rm ét}}$ to the category of sheaves on
$X'_{\text{\rm ét}}$ endowed with descent data relative to
$f:X'\rightarrow X$).
\end{enonce*}

The first assertion follows immediately from the transitivity of fiber
functors (\hyperref[VIII.subsec3.4]{3.4}) and from
(\hyperref[VIII.cor3.7]{3.7}); the second (which can also be stated by
saying that $f$ is a \emph{descent morphism} relative to the fibered
category of étale sheaves on variable schemes) also means that for two
sheaves $F$, $G$ on $X$, the natural diagram
\begin{equation*}
\Hom(F,G)\rightarrow \Hom (F',G')\rightrightarrows \Hom (F'',G'')\tag{$\ast$}
\end{equation*}
is exact, where $F'$ and $G'$ (\resp $F''$ and $G''$) are the inverse
images of $F$ and $G$ on $X'$ (\resp on
$X''=\fprod{X'}{X'}{X}$). Taking the final sheaf for $F$, the assertion
gives the following

\begin{enonce*}{Corollary 9.2}\etiquette[9.2]{VIII.cor9.2}
Let $f:X'\rightarrow X$ be a universally submersive morphism. Then for
every sheaf $F$ on $X$, the diagram
$$
\Gamma(X,F)\rightarrow \Gamma(X',F')\rightrightarrows \Gamma(X'',F'')
$$
is exact, where $X''=\fprod{X'}{X'}{X}$, and where $F'$, $F''$ are the
inverse images of $F$ on $X'$, $X''$.
\end{enonce*}

We shall prove \hyperref[VIII.prop9.1]{9.1} using the following

\begin{enonce*}{Lemma 9.3}\etiquette[9.3]{VIII.lem9.3}
Suppose that $f$ is universally submersive. Let $F$ be a sheaf on $X$,
denote by $S(F)$ the set of subsheaves of \sisi{$X$}{$F$},\pageoriginale
and define $S(F')$, $S(F'')$ analogously. Then the diagram of natural
maps
$$
S(F)\rightarrow S(F')\rightrightarrows S(F'')
$$
is exact.
\end{enonce*}

\noindent
Proof. The fact that $S(F)\rightarrow S(F')$ is injective follows
immediately from the fact that the functor $f^{\ast}$ is conservative;
for if $F_{i}$ $(i=1,2)$ are two subsheaves of $F$ such that
$f^{\ast}(F_{1})=f^{\ast}(F_{2})$, then the inclusions
$F_{3}=F_{1}\cap F_{2}\rightarrow F_{i}$ $(i=1,2)$ are isomorphisms,
since they become so after applying the functor $f^{\ast}$, hence
$F_{1}=F_{2}$. It remains to prove that if $G'$ is a subsheaf of $F'$
such that $\pr^{\ast}_{1}(G')=\pr^{\ast}_{2}(G')$, then $G'$ is the
inverse image of a subsheaf of $F$. Note that one can find an
epimorphism $F_{1}\rightarrow F$ in $X_{\text{\rm ét}}$, where
$F_{1}$ is \emph{representable} by a scheme étale over $X$.
Introduce $F_{2}=\fprod{F_{1}}{F_{1}}{F}$, and likewise $F'_{1}$,
$F''_{1}$, $F'_{2}$, $F''_{2}$, sheaves giving rise to a diagram of
sets
\[
\xymatrix@R=.75cm{
  S(F)  \ar[r] \ar[d]& S(F') \ar[r]\ar[d] & S(F'') \ar[d]\\
S(F_1) \ar[r] \ar@<3pt>[d] \ar@<-3pt>[d]& S(F'_1) \ar@<3pt>[r]
  \ar@<-3pt>[r] \ar@<3pt>[d] \ar@<-3pt>[d]& S(F''_{1})\ar@<3pt>[d] \ar@<-3pt>[d] \\
S(F_2)  \ar[r] & S(F'_2)\ar@<3pt>[r]
  \ar@<-3pt>[r] & S(F''_2),
}
\]
in which the columns are exact, by the exactness properties in the
category of sheaves on $X$, $X'$, $X''$, respectively. A standard
diagram chase then shows that, to prove that the first row is exact,
it is enough to prove this for rows 2 and 3. For row 2, this follows
from the fact that $F_{1}$ is representable, that $X'\rightarrow X$ is
universally\pageoriginale submersive, and from
\SGA\sisi{}{ 1} IX 2.3. On the other hand, $F_{2}$ is also
representable, since it is a subsheaf of
$\fprod{F_{1}}{F_{1}}{X}$, which is representable, and one applies
\hyperref[VIII.prop6.1]{6.1}. Thus row 3 is also exact, \cqfd

We now prove \hyperref[VIII.prop9.1]{9.1}, \ie that every morphism
$u':F'\rightarrow G'$ compatible with the descent data on $F'$, $G'$,
the morphism $u'$ comes from a morphism $u:F\rightarrow G$. Let
$H=F\times G$, hence
$H'=F'\times G'$, $H''=F''\times G''$; then the \emph{graph} of $u'$
is a subsheaf $\Gamma'$ of $H'$, whose two inverse images are equal to
the graph of the same morphism $u'':F''\rightarrow G''$. Hence, by
\hyperref[VIII.lem9.3]{9.3}, $\Gamma'$ comes from a subsheaf $\Gamma$
of $H$. I claim that $\Gamma$ is the graph of a morphism
$u:F\rightarrow G$, \ie that the morphism $p:\Gamma\rightarrow F$
induced by $\pr_{1}$ is an isomorphism: indeed, it becomes an
isomorphism after the base change $X'\rightarrow X$, and we apply the
already proved part of \hyperref[VIII.prop9.1]{9.1}. The morphism
$u:F\rightarrow G$ then has the required property, \cqfd

\begin{enonce*}{Theorem 9.4}\etiquette[9.4]{VIII.thm9.4}
Let $f:X'\rightarrow X$ be a surjective morphism of schemes. Suppose
that one of the following hypotheses holds:
\begin{enumerate}
\renewcommand{\theenumi}{\alph{enumi}}
\renewcommand{\labelenumi}{\theenumi)}
\item $f$ is integral.

\item $f$ is proper.

\item $f$ is flat and locally of finite presentation
\footnote[1]{It is likely that this second hypothesis is in fact
superfluous.}.

\item $X$ is discrete (e.g. the spectrum of a field).
\end{enumerate}

Then $f$ is an \emph{effective descent morphism} for the fibered
category $\Ff$ over $(\Sch)$ of étale sheaves on variable schemes,
\ie the functor $f^{\ast}$ induces an \emph{equivalence} between the
category of\pageoriginale sheaves on $X$ and the category of sheaves on
$X'$ endowed with descent data relative to $f:X'\rightarrow X$.
\end{enonce*}

\setcounter{subsection}{4}
\setcounter{subsubsection}{0}
\subsubsection{}\etiquette[9.4.1]{VIII.subsubsec9.4.1}

Let $F'$ be a sheaf on $X'$ endowed with descent data relative to
$f:X'\rightarrow X$. Define a functor
$$
G:X^{0}_{\text{\rm ét}}\longrightarrow (\Ens)
$$
by the formula
$$
G(Y)=\Ker (F'(Y')\rightrightarrows F''(Y'')){\quoi}.
$$
One sees immediately that $G$ is a \emph{sheaf} for the étale
topology. On the other hand, there is a canonical injective
homomorphism $G\rightarrow f_{\ast}(F')$, hence a homomorphism
$f^{\ast}(G)\rightarrow F'$, and the latter is evidently compatible
with the descent data. It remains to determine whether this
homomorphism is an \emph{isomorphism}. Note that $G$ can also be defined
by the exactness of
$$
G\rightarrow f_{\ast}(F')\rightrightarrows g_{\ast}(F''){\quoi},
$$
so for every base change $Y\xrightarrow{h}X$ \emph{which commutes with
  the formation of} $f_{\ast}(F')$ and $g_{\ast}(F'')$,
$h^{\ast}(G)=G_{Y}$ is identified with the sheaf ``descended'' from
$F'_{Y'}$ by $Y'\rightarrow Y$, and of course the homomorphism
$f_{Y}{}^{\ast}(G_{Y})\rightarrow (F_{Y})'$ is the one deduced from
$f^{\ast}(G)\rightarrow F'$ by base change. Now \emph{suppose that the
  base change} $f:Y=X'\rightarrow X$ \emph{commutes with the formation
  of} $f_{\ast}(F')$ and $g_{\ast}(F'')$, and note that since
$Y'=\fprod{X'}{X'}{X}$ has a section over $Y=X'$, hence $Y'\rightarrow
Y$ is an \emph{effective} descent morphism for every fibered category,
in particular for $\mathcal{F}$, it follows that
$f^{\ast}_{Y}(G_{Y})\rightarrow (F_{Y})'$ is an isomorphism. Thus
$f^{\ast}(G)\rightarrow F'$ becomes an isomorphism after the base
change $Y'=\fprod{X'}{X'}{X}\rightarrow X'$. Since the latter has a
section, $f^{\ast}(G)\rightarrow F'$ is an isomorphism, and hence the
descent datum\pageoriginale under consideration on $F'$ is effective.

\subsubsection{}\etiquette[9.4.2]{VIII.subsubsec9.4.2}

This proves Theorem \hyperref[VIII.thm9.4]{9.4} in case a), by
\hyperref[VIII.cor5.6]{5.6}. Case b) also follows from the fact that if
$f$ is proper, then $f_{\ast}$ commutes with \emph{every} base change
$Y'\rightarrow Y$ (which will be proved in a later exposé (XII \hyperref[XII.thm5.1]{5.1}
(i))).

\subsubsection{}\etiquette[9.4.3]{VIII.subsubsec9.4.3}

In case c), the question being local on $X$, we may suppose that $X$
is affine and that there exists an affine scheme $X'_{1}$, of finite
presentation, quasi-finite and faithfully flat over $X$, together with
an $X$-morphism $X'_{1}\rightarrow X'$, which reduces us to the case
where $f$ is quasi-finite. After localizing further on $X$, this time
in the sense of the étale topology, one sees that there will exist an
open subset $X'_{1}$ of $X'$ such that $X'_{1}\rightarrow X$ is finite
and surjective, which reduces us to the case where $f$ is finite and
surjective, already treated in a).

\subsubsection{}\etiquette[9.4.4]{VIII.subsubsec9.4.4}
In case d), we may suppose (by localizing on $X$) that $X$ consists of
a single point, and even, in view of \hyperref[VIII.thm1.1]{1.1}, that
it is the spectrum of a field $k$. Moreover, $f$ is universally open
(\EGA IV 2.4.9), so the hypotheses of \hyperref[VIII.prop9.1]{9.1}
hold, and the argument of \hyperref[VIII.subsubsec9.4.1]{9.4.1} applies
again, making a base change with $Y=\Spec k$, where $k=k(x)$, for an
$x\in X$. But then it is still true that $f_{\ast}$ commutes with the
base change $Y\rightarrow X$, as we shall see later (XVI \hyperref[XVI.compl1.4]{1.4} and \hyperref[XVI.cor1.5]{1.5}).

This completes the proof of \hyperref[VIII.thm9.4]{9.4}.


\ifx\danslelivre\undefined
\end{document}
\fi

% END INCLUDED SGA4_English_current_buildable_source_20260906/source/08/drafts/SGA4_VIII_LINES_1557_1976_English.texfrag


% END INCLUDED SGA4_English_current_buildable_source_20260906/source/08/08_English.tex


\makeatletter
\def\input@path{{../09/}}
\makeatother
% BEGIN INCLUDED SGA4_English_current_buildable_source_20260906/source/09/09_English.tex
\ifx\danslelivre\undefined
% Release engine: XeLaTeX (the LuaTeX compatibility layer breaks Xy-pic output).
\documentclass[10pt]{../fichiers-configuration-latex/smfart}
% BEGIN INCLUDED SGA4_English_current_buildable_source_20260906/source/fichiers-configuration-latex/paquets.tex
\usepackage{polyglossia}
%\usepackage[francais]{babel}
\usepackage{fontspec}
\usepackage[math-style=ISO,bold-style=upright,sans-style=upright]{unicode-math}
%\usepackage[mathletters]{ucs}
%\usepackage[utf8]{inputenc}
%\usepackage[T1]{fontenc}
\usepackage[unicode,hypertexnames=false,colorlinks=true,linkcolor=blue,citecolor=Sepia,urlcolor=magenta]{hyperref}
\setmainlanguage{french}
\setmainfont[Ligatures=TeX]{Linux Libertine O}
% variante possible : TeXGyrePagella ?
\setmathfont{XITS Math}

% maths
\usepackage{amsmath}
%\usepackage{amssymb}
%\usepackage{mathrsfs}
%\usepackage{stmaryrd}
%\usepackage{euscript}

% couleurs
\usepackage[usenames,dvipsnames]{xcolor}

% graphiques
\usepackage{graphicx}
\usepackage[all]{xy}
%\usepackage{tikz}
%\usetikzlibrary{matrix,arrows,calc}
%\tikzset{column sep=3em,row sep=2em,text height=1.5ex, text depth=0.25ex,>=angle 90}
%\tikzset{description/.style={fill=white,inner sep=1pt}}

% hyperliens, 'etiquettes
\usepackage{xr-hyper}
\usepackage{../fichiers-configuration-latex/smfthm}
%\ifx\cacheretiquettes\undefined
%\usepackage[inline]{showlabels} \else \fi % pour rendre visibles les 'etiquettes
%\usepackage{../fichiers-configuration-latex/makeidx} % pour faire un index
%\usepackage[pagebackref,unicode]{hyperref}  % pour voir o`u sont cit'ees les r'ef'erences
% hyperref is loaded once above with the cumulative-reader destination policy.
%\hypersetup{breaklinks=true} %permet le retour `a la ligne dans les liens trop longs. Incompatible avec pdflatex
%\renewcommand*{\backref}[1]{{\small [cit\'e page(s)~#1.]}}

\usepackage{xspace}
\usepackage{pifont} % pour expose XVII

% END INCLUDED SGA4_English_current_buildable_source_20260906/source/fichiers-configuration-latex/paquets.tex

% BEGIN INCLUDED SGA4_English_current_buildable_source_20260906/source/fichiers-configuration-latex/numerotation.tex
\renewcommand{\thefootnote}{\roman{footnote}} % pour mettre les notes en bas de page sous la forme (ii) et non (2), qui ressemble trop (?) a un exposant.

\newcounter{pageoriginal}
\newcommand{\pageoriginale}{\refstepcounter{pageoriginal}\marginpar{\footnotesize\textbf{\thepageoriginal}}\xspace}
\let\pageoriginaled\pageoriginale

\usepackage{ifthen}
\newboolean{original}
\def\sisi{\ifthenelse{\boolean{original}}}
\setboolean{original}{false}

\newcounter{cnote}
\def\nde#1{\ifthenelse{\boolean{original}}{}{\stepcounter{cnote}\def\thefootnote{\arabic{cnote}}\protect\footnote[\value{cnote}]{N.D.E.: #1}}}
\def\ndemark{\ifthenelse{\boolean{original}}{}{\stepcounter{cnote}\def\thefootnote{\arabic{cnote}}\protect\footnotemark[\value{cnote}]}}
\def\ndetext#1{\ifthenelse{\boolean{original}}{}{\def\thefootnote{\arabic{cnote}}\protect\footnotetext[\value{cnote}]{N.D.E.: #1}}}

% END INCLUDED SGA4_English_current_buildable_source_20260906/source/fichiers-configuration-latex/numerotation.tex

% BEGIN INCLUDED SGA4_English_current_buildable_source_20260906/source/fichiers-configuration-latex/macros.tex
\setmainlanguage{french}
\setotherlanguage{english}


%% Fixed, engine-independent release identifier.  The inherited dynamic Git
%% probe used Lua io.popen, exposed build-state-dependent output, and printed
%% a visible question mark outside a Git checkout.
\def\showgitstatus{English release r2, 2026-07-23}


% macros pour LaTeX
\renewcommand{\mathsf}{}

% Police différente pour U+2113 SCRIPT SMALL L
\setmathfont[range="2113]{latinmodern-math.otf}

% Give every inherited semantic label its own stable PDF destination instead
% of inheriting whichever automatic equation/page anchor happened to be live.
% The legacy label remains as a compatibility alias used by existing edges.
\makeatletter
\def\etiquette[#1]#2{%
  \Hy@raisedlink{\hyper@anchorstart{sga4.semantic.#2}\hyper@anchorend}%
  \def\@currentHref{sga4.semantic.#2}%
  \def\@currentlabel{#1}%
  \label{#2}%
}
\makeatother
\def\eqhyperref[#1]#2{(\hyperref[#1]{#2})}
\def\pagehyperref[#1]#2{(page \hyperref[#1]{#2})}
\def\hyperrefext[#1][#2]#3{\hyperref[#2]{#1 #3}}
\def\eqhyperrefext[#1][#2]#3{\hyperref[#2]{#1 #3}}     %cf. exposes 17
\def\neweqhyperrefext[#1][#2]#3{\hyperref[#2]{#1 #3}}  %cf. exposes 17

\def\shat#1{\protect\widehat{#1}}
\def\stilde#1{\protect\widetilde{#1}}
%%%%%%%%%%%%%%%%%%


\def\etcn{etc.\xspace}
\def\etco{etc...\xspace}
\def\etc{\sisi{\etco}{\etcn}}
\def\ldotsn{\sisi{...}{\etcn}}

\def\cf{{cf.}\kern.3em}
\def\Cf{{Cf.}\kern.3em}
\def\ie{{i.e.}\kern.3em}
\def\iev{{i.e.},\xspace}
\def\loccito{loc.\kern2pt cit.\xspace}
\def\loccitn{\textit{loc.\kern2pt cit.}\xspace}
\def\loccit{\sisi{\loccito}{\loccitn}}
\def\EGA{\textup{EGA}\kern.2em}
\def\SGA{\textup{SGA}\kern.3em}

\def\quoi{}

\def\quav{\quad,\quad}
\def\quaet{\quad \text{et}\quad}
\def\fibreproduct#1#2#3{#1 ×_{#3} #2}
\let\fprod\fibreproduct
\def\ccheck#1#2{\mbox{${\check{\check{\rm #1}}}^{\lower2pt\hbox{$\scriptstyle #2$}}$}}
\def\sheaf#1{\protect\mathcal{#1}}
\def\vieuxsheaf#1{{#1}} % cf. problème avec sheaf{a\Phi_i} par exemple...
\DeclareMathOperator{\Hom}{Hom}
\DeclareMathOperator{\can}{can}
\DeclareMathOperator{\Ind}{Ind}
\DeclareMathOperator{\ind}{ind}
%\DeclareMathOperator{\Hhomtop}{\mathcal{H}\!\mathit{omtop}}
\DeclareMathOperator{\Hhomtop}{\mathbf{Homtop}}
\DeclareMathOperator{\Point}{Point}
\DeclareMathOperator{\Pro}{Pro}
\DeclareMathOperator{\Top}{Top}
\DeclareMathOperator{\Mod}{Mod}
\DeclareMathOperator{\Cart}{Cart}
%\DeclareMathOperator{\Mmorsite}{\mathcal{M}\!\mathit{orsite}}
\DeclareMathOperator{\Mmorsite}{\mathbf{Morsite}}
\DeclareMathOperator{\loc}{loc}
\DeclareMathOperator{\Loc}{Loc}
\DeclareMathOperator{\Pic}{Pic}
\DeclareMathOperator{\Et}{Et}
\DeclareMathOperator{\card}{card}
\DeclareMathOperator{\Ens}{Ens}
\DeclareMathOperator{\h}{\underline{H}}
\DeclareMathOperator{\Fonct}{Fonct}
\DeclareMathOperator{\Ker}{Ker}
\DeclareMathOperator{\but}{but}
\DeclareMathOperator{\Ob}{Ob}
\DeclareMathOperator{\ob}{ob}
\DeclareMathOperator{\Lib}{Lib}
\DeclareMathOperator{\esj}{esj}
\DeclareMathOperator{\id}{id}
\DeclareMathOperator{\Id}{Id}
\DeclareMathOperator{\ad}{ad}
\DeclareMathOperator{\Iim}{Im}
\DeclareMathOperator{\pr}{pr}
\DeclareMathOperator{\Cat}{Cat}
\DeclareMathOperator{\car}{car}
\DeclareMathOperator{\SysInd}{SysInd}
\DeclareMathOperator{\Sup}{Sup}
\DeclareMathOperator{\Coker}{Coker}
\DeclareMathOperator{\Filt}{Filt}
\DeclareMathOperator{\acc}{acc}
\DeclareMathOperator{\sur}{sur}
\DeclareMathOperator{\cat}{cat}
\DeclareMathOperator{\Ord}{Ord}
\DeclareMathOperator{\Gr}{Gr}
\DeclareMathOperator{\Cov}{Cov}
\DeclareMathOperator{\Caf}{Caf}
\DeclareMathOperator{\coker}{coker}
\DeclareMathOperator{\coim}{coim}
\DeclareMathOperator{\im}{im}
\DeclareMathOperator{\Ouv}{Ouv}
\DeclareMathOperator{\diag}{diag}
\DeclareMathOperator{\Esp}{Esp}
\DeclareMathOperator{\TOP}{TOP}
\DeclareMathOperator{\Unif}{Unif}
\DeclareMathOperator{\Iind}{\mathbf{Ind}}
\DeclareMathOperator{\Hhomcart}{\mathbf{Homcart}}
\DeclareMathOperator{\Hhomtopess}{\mathbf{Homtopess}}
\DeclareMathOperator{\Ppoints}{\mathbf{Points}}
\DeclareMathOperator{\Ppoint}{\mathbf{Point}}
\DeclareMathOperator{\Ttors}{\mathbf{Tors}}
\DeclareMathOperator{\Ssex}{\mathbf{Sex}}
\DeclareMathOperator{\Ffib}{\mathbf{Fib}}
\DeclareMathOperator{\Vvois}{\mathbf{Vois}}
\DeclareMathOperator{\Ppointess}{\mathbf{Pointess}}
\DeclareMathOperator{\scrHom}{\mathbf{Hom}}
\DeclareMathOperator{\sheafPic}{\mathbf{Pic}}
\DeclareMathOperator{\sheafAut}{\mathbf{Aut}}
\DeclareMathOperator{\Hhomtoplocan}{\mathbf{Homtoplocan}}
\DeclareMathOperator{\Hhomtopan}{\mathbf{Homtopan}}
\DeclareMathOperator{\Pptgeom}{\mathbf{Ptgéom}}
\DeclareMathOperator{\Morsite}{Morsite}
\DeclareMathOperator{\Fib}{Fib}
\DeclareMathOperator{\Fais}{Fais}
\DeclareMathOperator{\kar}{kar}
\DeclareMathOperator{\Groupes}{Groupes}
\DeclareMathOperator{\Restr}{Restr}
\DeclareMathOperator{\Isom}{Isom}
\DeclareMathOperator{\Prol}{Prol}
\DeclareMathOperator{\Res}{Res}
\DeclareMathOperator{\Tors}{Tors}
\DeclareMathOperator{\esp}{esp}
\DeclareMathOperator{\es}{es}
\DeclareMathOperator{\od}{od}
\DeclareMathOperator{\gr}{gr}
\DeclareMathOperator{\sob}{sob}
\DeclareMathOperator{\dfn}{dfn}
\DeclareMathOperator{\Sch}{Sch}
\DeclareMathOperator{\ouv}{ouv}
\DeclareMathOperator{\cons}{cons}
\DeclareMathOperator{\Nerf}{Nerf}
\DeclareMathOperator{\Points}{Points}
\DeclareMathOperator{\Espsob}{Espsob}
\DeclareMathOperator{\Homtop}{Homtop}
\DeclareMathOperator{\Pt}{Pt}
\DeclareMathOperator{\Quot}{Quot}
\DeclareMathOperator{\glob}{glob}
\DeclareMathOperator{\Homfib}{Homfib}
\DeclareMathOperator{\Qouv}{Qouv}
\DeclareMathOperator{\Quotouv}{Quotouv}
\DeclareMathOperator{\ens}{ens}
\DeclareMathOperator{\ab}{ab}
\DeclareMathOperator{\an}{an}
\DeclareMathOperator{\fs}{fs}
\DeclareMathOperator{\red}{red}
\DeclareMathOperator{\Ab}{Ab}
\DeclareMathOperator{\cosq}{cosq}
\DeclareMathOperator{\Ann}{Ann}
\DeclareMathOperator{\Zar}{Zar}


%%%%%%%%%%%%
%%%%% SGA4-2
\DeclareMathOperator{\Sex}{Sex}
\DeclareMathOperator{\vois}{vois}
\DeclareMathOperator{\Vois}{Vois}
\DeclareMathOperator{\Ext}{Ext}
\DeclareMathOperator{\ext}{ext}
\DeclareMathOperator{\Fer}{Fer}
\DeclareMathOperator{\Add}{Add}
\DeclareMathOperator{\Cocart}{Cocart}
\DeclareMathOperator{\cosk}{cosk}
\DeclareMathOperator{\plat}{plat}

\DeclareMathOperator{\faisc}{faisc}
\DeclareMathOperator{\sk}{sk}
\DeclareMathOperator{\ch}{ch}

\def\HH{\mathbf{H}}
\def\qc{\mathrm{qc}}
\def\et{\mathrm{et}}
\def\coh{\mathrm{coh}}

\def\Hvr#1{%
\begin{array}[b]{l}
{\kern2pt\scriptstyle r}\\[-8pt]
{\kern1pt\scriptstyle \vee}\\[-5pt]
H^{#1}
\end{array}}

\def\Hvv#1{%
\begin{array}[b]{l}
\kern2pt\rotatebox{90}{${\scriptscriptstyle\ll}$}\\[-6pt]
H^{#1}
\end{array}}

\def\Hvi#1{%
\begin{array}[b]{l}
{\kern0.2pt\scriptscriptstyle \infty}\\[-8pt]
{\kern1pt\scriptstyle \vee}\\[-5pt]
H^{#1}
\end{array}}

\def\Hv#1{%
\begin{array}[b]{l}
{\kern1pt\scriptstyle \vee}\\[-5pt]
H^{#1}
\end{array}}

\DeclareMathOperator{\cart}{cart}
\DeclareMathOperator{\Ffaisc}{\mathbf{Faisc}}


\def\Ev#1{%
\begin{array}[b]{l}
{\kern1pt\scriptstyle \vee}\\[-5pt]
E'_{#1}
\end{array}}

\DeclareMathOperator{\Cont}{Cont}
\DeclareMathOperator{\Ccont}{\mathscr{C}\!\mathit{ont}}
\DeclareMathOperator{\Site}{Site}
\newcommand\MM{\mathbf{M}}
\def\limtop#1{{\displaystyle{\varprojlim_{#1}\!\mbox{\rm top }}}}
\def\Limtop{\mathop{\varprojlim\!\mbox{\rm top}}}
\DeclareMathOperator{\spec}{spec}
\def\ss{\scriptstyle{\mathscr{S}}}
\DeclareMathOperator{\Gal}{Gal}
\def\smallm{\scriptstyle{\mathscr{M}}}
\DeclareMathOperator{\End}{End}
\DeclareMathOperator{\ass}{ass}
\DeclareMathOperator{\zar}{zar}
\DeclareMathOperator{\Ogg}{Ogg}
\DeclareMathOperator{\res}{res}
\DeclareMathOperator{\tr}{tr}
\DeclareMathOperator{\Supp}{Supp}
\DeclareMathOperator{\codim}{codim}
\DeclareMathOperator{\rk}{rk}
\DeclareMathOperator{\cur}{cur}
\DeclareMathOperator{\cdqc}{cdqc}
\newcommand\DD{\mathbf{D}}
\DeclareMathOperator{\cl}{cl}
\def\JJ{\mathbf{J}}
\def\FF{\mathbf{F}}
\def\BB{\mathbf{B}}
\DeclareMathOperator{\Alget}{Alget}
\DeclareMathOperator{\Ass}{Ass}
\DeclareMathOperator{\supp}{supp}
\DeclareMathOperator{\rad}{rad}
\DeclareMathOperator{\prof}{prof}
\DeclareMathOperator{\cint}{int}
\DeclareMathOperator{\sic}{sic}
\DeclareMathOperator{\Tor}{Tor}
\DeclareMathOperator{\tors}{tors}
\DeclareMathOperator{\Sym}{Sym}
\DeclareMathOperator{\parf}{parf}
\DeclareMathOperator{\tor}{tor}
\DeclareMathOperator{\defn}{defn}
\DeclareMathOperator{\Div}{Div}
\DeclareMathOperator{\Pol}{Pol}
\DeclareMathOperator{\Fex}{Fex}
\DeclareMathOperator{\Lie}{Lie}
\DeclareMathOperator{\tot}{tot}
\DeclareMathOperator{\Aut}{Aut}
\DeclareMathOperator{\HOM}{HOM}
\DeclareMathOperator{\pch}{pch}
\DeclareMathOperator{\Ch}{Ch}
\DeclareMathOperator{\EXT}{EXT}
\DeclareMathOperator{\PIC}{PIC}
\DeclareMathOperator{\Rf}{Rf}
\DeclareMathOperator{\inj}{inj}
\newcommand\VV{\mathbf{V}}

\def\Xv{%
\begin{array}[b]{l}
{\kern1pt\scriptstyle \vee}\\[-5pt]
X
\end{array}}

%%%%%%%
\newcommand\TT{\mathbf{T}}
\newcommand\XX{\mathbf{X}}
\newcommand\EE{\mathbf{E}}
%%%%%%

\def\qedsymbol{\textup{C.Q.F.D.}}
\let\CQFD\qed
\let\cqfd\qed

%%%%%\bigdot
\font\bigsymb=cmsy10 at 4pt
\def\bigdot{{\kern1.2pt\raise 1.5pt\hbox{\bigsymb\char15}}}
%%%%%

\def\Lright{\mathop{\vcenter{\hbox{\hbox{$\longrightarrow$}\lower8pt \hbox{\kern-16pt\hbox{$\longrightarrow$}}}}}}

\def\Aa{{\mathcal A}}
\def\Bb{{\mathcal B}}
\def\Cc{{\mathcal C}}
\def\Dd{{\mathcal D}}
\def\Ee{{\mathcal E}}
\def\Ff{{\mathcal F}}
\def\Gg{{\mathcal G}}
\def\Hh{{\mathcal H}}
\def\Ii{{\mathcal I}}
\def\Jj{{\mathcal J}}
\def\Kk{{\mathcal K}}
\def\Ll{{\mathcal L}}
\def\Mm{{\mathcal M}}
\def\Nn{{\mathcal N}}
\def\Oo{{\mathcal O}}
\def\Pp{{\mathcal P}}
\def\Qq{{\mathcal Q}}
\def\Rr{{\mathcal R}}
\def\Ss{{\mathcal S}}
\def\Tt{{\mathcal T}}
\def\Uu{{\mathcal U}}
\def\Vv{{\mathcal V}}
\def\Ww{{\mathcal W}}
\def\Xx{{\mathcal X}}
\def\Yy{{\mathcal Y}}
\def\Zz{{\mathcal Z}}

\newcommand{\CC}{\mathbf{C}}
\newcommand{\GG}{\mathbf{G}}
\newcommand{\KK}{\mathbf{K}}
\newcommand{\LL}{\mathbf{L}}
\newcommand{\NN}{\mathbf{N}}
\newcommand{\PP}{\mathbf{P}}
\newcommand\QQ{\mathbf{Q}}
\newcommand\RR{\mathbf{R}}
\newcommand\ZZ{\mathbf{Z}}

\def\to{\mathchoice{\longrightarrow}{\rightarrow}{\rightarrow}{\rightarrow}}
\def\from{\mathchoice{\longleftarrow}{\leftarrow}{\leftarrow}{\leftarrow}}
\def\mto{\mathchoice{\longmapsto}{\mapsto}{\mapsto}{\mapsto}}
\def\mfrom{\mathrel{\from\joinrel\mapstochar\,}}
\def\hto{\mathrel{\lhook\joinrel\to}}
\def\hfrom{\mathrel{\from\joinrel\rhook}}

\def\implies{\mathchoice{\Longrightarrow}{\Rightarrow}{\Rightarrow}{\Rightarrow}}
\def\implied{\mathchoice{\Longleftarrow}{\Leftarrow}{\Leftarrow}{\Leftarrow}}
\def\iff{\mathchoice{\Longleftrightarrow}{\Leftrightarrow}{\Leftrightarrow}{\Leftrightarrow}}

\def\To#1{\mathchoice{\xrightarrow{\labelstyle~~#1~}}{\stackrel{#1}{\longrightarrow}}{}{}}
\def\Mto#1{\mathrel{\mapstochar\To{#1}}}
\def\From#1{\mathchoice{\xleftarrow{\labelstyle~~#1~}}{\stackrel{#1}{\longleftarrow}}{}{}}

\newcommand{\isomto}{\overset{\sim}{\longrightarrow}}
\newcommand{\isomfrom}{\overset{\sim}{\longleftarrow}}

\newcommand{\IFF}{\unskip~$\iff$~\ignorespaces}
\newcommand{\THEN}{\unskip~$\implies$~\ignorespaces}

\let\leq\leqslant
\let\geq\geqslant
\let\le\leqslant
\let\ge\geqslant

\def\cf{{cf.}\kern.3em}
\def\Cf{{Cf.}\kern.3em}
\def\ie{{i.e.}\kern.3em}
\def\iev{{i.e.},\xspace}
%\def\loccito{loc.\kern2pt cit.\xspace}
%\def\loccitn{\textit{loc.\kern2pt cit.}\xspace}
%\def\loccit{\sisi{\loccito}{\loccitn}}
\def\EGA{\textup{EGA}\kern.2em}
\def\SGA{\textup{SGA}\kern.3em}

\let\oldS\S
\def\S{\oldS\kern .15em}
\def\ptbl{.\kern .2em}
\def\resp{resp.\kern.3em}
\def\Exp{Exp.\kern.3em}

\let\bar\overline
\let\hat\widehat
\let\tilde\widetilde
\let\dpl\displaystyle
\let\ldots\dots
\let\emptyset\varnothing
\let\skpt\skippointrait
\let\epsilon\varepsilon

\def\sfrac#1#2{{#1}/{#2}}

\DeclareMathOperator{\cd}{cd}
%\DeclareMathOperator{\Hhom}{\mathcal{H}\!\mathit{om}}
\DeclareMathOperator{\Hhom}{\mathbf{Hom}}
%\DeclareMathOperator{\Eext}{\mathcal{E}\!\mathit{xt}}
\DeclareMathOperator{\Eext}{\mathbf{Ext}}
\DeclareMathOperator{\Spec}{Spec}
\DeclareMathOperator{\Tr}{Tr}

\let\emptyset=\varnothing

%%%%%%%%%%%%%%%
%%% expose XVII
%%%%%%%%%%%%%%%

\makeatletter
\let\old@optexpose\optexpose
\newcommand{\empty@optexpose}[1]{}
\newcommand{\paren@optexpose}[1]{\old@optexpose{#1}(}
\newcommand{\refnoexpose}[1]{\let\optexpose\empty@optexpose\ref{#1}\let\optexpose\old@optexpose}%\refnoexpose{X.cor5.2} engendre "5.2" au lieu de "X 5.2"
\newcommand{\neweqref}[1]{\let\optexpose\paren@optexpose\textup{\ref{#1})}\let\optexpose\old@optexpose}%\neweqref{XII.eq4.1} engendre "XII (4.1)" au lieu de "(XII 4.1)"
\makeatother
\newcommand{\refvb}[2][VI${}_{\text{B}}$ ]{\sisi{#1\refnoexpose{#2}}{\ref{#2}}}

\newcommand{\leftexp}[2]{{\vphantom{#2}}^{#1}{#2}}
\newcommand{\AAA}{\mathbf{A}}
\newcommand{\Sf}{\mathfrak{S}}
\DeclareMathOperator*{\quotprojlim}{\text{``$\varprojlim$''}}
\DeclareMathOperator*{\quotinjlim}{\text{``$\varinjlim$''}}
\DeclareMathOperator{\Fl}{F\sisi{\ell}{l}}
\DeclareMathOperator{\Pl}{P\sisi{\ell}{l}}
%\DeclareMathOperator{\cl}{c\sisi{\ell}{l}}
\DeclareMathOperator{\tdf}{tdf}
\DeclareMathOperator{\rg}{rg}
\newcommand{\iInd}{\sisi{I}{i}nd}
\newcommand{\iind}{ind}
\newcommand{\ppro}{pro}
\newcommand{\tTor}{\sisi{T}{t}or}
\newcommand{\ttor}{tor}
\newcommand{\chap}{Chap.}
\newcommand{\Ffex}{\Ff\mathit{ex}}
\newcommand{\sheafk}{\underline{k}}
\newcommand{\sheafn}{\underline{n}}
\newcommand{\sheafGamma}{\underline{\Gamma}}

% Repeated references to an already typeset footnote must not use
% \footnotemark without a matching \footnotetext: Hyperref would then create
% a live annotation to a destination that cannot exist.  This helper keeps the
% current printed footnote mark and links it to an explicit stable note target.
\makeatletter
\newcommand{\sgarepeatedfootnotemark}[1]{%
  \begingroup
  \edef\@thefnmark{\thefootnote}%
  \hyperlink{#1}{\@makefnmark}%
  \endgroup}
\makeatother

\sisi{
\DeclareMathOperator{\simarrow}{\rule{.15mm}{1.8mm}\kern -3.4pt\rightsquigarrow}
\newcommand{\sheafell}{\underline{\ell}}
\newcommand{\sheafZZ}{\underline{\ZZ}}
}{
\newcommand{\nderef}[1]{la N.D.E.~(\ref{#1}) page~\pageref{#1}}
}

%%%%%%%%%%%%%%%%%
%%%%% expose XVII
%%%%%%%%%%%%%%%%%

% END INCLUDED SGA4_English_current_buildable_source_20260906/source/fichiers-configuration-latex/macros.tex


\externaldocument{../01/01_English.tex}
\externaldocument{../02/02_English.tex}
\externaldocument{../03/03_English.tex}
\externaldocument{../04/04_English.tex}
\externaldocument{../05/05_English.tex}
\externaldocument{../05bis/05bis_English.tex}
\externaldocument{../06/06_English.tex}
\externaldocument{../07/07_English.tex}
\externaldocument{../08/08_English.tex}

\setboolean{original}{false}

\begin{document}

\begin{center}
IX. Constructible sheaves; cohomology of an algebraic curve\\
M. Artin
\end{center}

\selectlanguage{english}version: \showgitstatus
\tableofcontents

\else
\chapter{Constructible sheaves; cohomology of an algebraic curve}
%\addtocontents{toc}{M. Artin}
\begin{center}
M. Artin
\end{center}
\fi
\phantomsection\MakeLinkTarget*{sga4.exp.IX}\label{sga4.exp.IX}

% BEGIN INCLUDED SGA4_English_current_buildable_source_20260906/source/09/drafts/SGA4_IX_LINES_43_278_English.texfrag
\setcounter{pageoriginal}{0}

\setcounter{section}{-1}
\section{Introduction}\etiquette[0]{IX.sec0}\pageoriginale

Note that, beginning with the present exposé, in contrast to the
preceding exposés and to the classical cohomological theory of
topological spaces, we are obliged, for the validity of the essential
statements, to restrict ourselves to \emph{torsion} sheaves (and often,
more precisely, to torsion sheaves prime to the residue
characteristics).

Let us point out that in fact there is no ``reasonable'' theory of
cohomology with integral (or even real) coefficients for varieties over
an algebraically closed field of characteristic $p\neq 0$, analogous
to the classical theory for the case $k=\CC$ (as Serre observed). More
precisely, there is no contravariant functor $H^1$, defined, say, on
the category of smooth projective schemes over an algebraically closed
field $K$ of characteristic $p>0$, with values in the category of
finite-dimensional vector spaces over $\RR$ (or a subfield of $\RR$),
``commuting with products,'' \ie such that $H^1(X\times
Y)\xleftarrow{\sim}H^1(X)\times H^1(Y)$, and such that $\dim H^1(X)=2$
if $X$ is a connected curve of genus 1. Indeed, it would follow that
if $X$ is an abelian variety over $K$, then the map $u\mapsto u^1$
from $\End(X)$ to $\End(H^1(X))$ is \emph{additive}, hence a
\emph{representation} of the opposite ring of $\End(X)$. But in
characteristic $p$ there exist elliptic\pageoriginale curves $X$ whose
endomorphism ring is a maximal order $A$ in a definite quaternion
algebra over $\ZZ$ ([3], p.198), and such an algebra plainly has no
representation in a vector space of dimension 2 over $\RR$, since
$A\otimes_{\ZZ}\RR$ is a division algebra (the quaternion algebra).

\section{The formalism of torsion sheaves}\etiquette[1]{IX.sec1}

Let $T$ be a topos and $F$ an abelian sheaf on $T$. One may multiply
by $n$ in $F$ ($n\in \ZZ$ an integer), this multiplication being
induced by the analogous operation in (Ab). Denote by ${}_nF$ the
kernel of this multiplication; thus ${}_nF(X)$, for $X\in \Ob T$, is
the group of sections of $F(X)$ whose order divides $n$.

We denote by $\PP$ the set of all prime numbers.

\begin{enonce*}{Definition 1.1}\etiquette[1.1]{IX.def1.1}
Let $p$ be a set of prime numbers. We say that $F$ is a $p$-torsion
sheaf, \sisi{where}{or} a $p$-sheaf, if the canonical morphism
$$
\varinjlim_n {}_nF\longrightarrow F\text{ is bijective}{\quoi},
$$
where $n$ ranges over the set of integers such that $\ass n\subset p$,
that is, such that the prime numbers dividing $n$ belong to $p$. If
$p=\PP$ is the set of all prime numbers, we simply say that $F$ is a
torsion sheaf.
\end{enonce*}

\begin{enonce*}{Proposition 1.2}\etiquette[1.2]{IX.prop1.2}
\begin{enumerate}
\renewcommand{\theenumi}{\roman{enumi}}
\renewcommand{\labelenumi}{(\theenumi)}
\item $F$ is a $p$-torsion sheaf if and only if $F$ is the sheaf
  associated with a presheaf taking values in $p$-torsion abelian
  groups; one may take
$$
P=\frac{\lim(\text{\rm presheaf})}{\ass n\subset p} {}_nF{\quoi}.
$$

\item If\pageoriginale $F$ is a $p$-torsion sheaf and if $X\in \Ob T$
  is a quasi-compact object\footnote[1]{An object $X$ of a topos is
  called quasi-compact if every covering family $\{X_i\rightarrow X\}$
  can be refined by a finite covering family. A morphism $u$ of
  topoi is quasi-compact if the inverse image of every quasi-compact
  object is again quasi-compact.}, then $F(X)$ is a $p$-torsion group.

\item If $T$ is locally of finite type {\rm (VI \hyperref[VI.def1.1]{1.1})}, then $F$ is a
  $p$-torsion sheaf if and only if $F(X)$ is a $p$-torsion group for
  every quasi-compact $X\in \Ob T$. In this case, one has
$$
\varinjlim_{\ass n\subset p} H^q(T/X; {}_nF)\xrightarrow{\sim}
H^q(T/X;F)
$$
for every quasi-compact $X$ and every $q$.

\item If $u:T\rightarrow T'$ is a morphism of topoi and if $F'$ is a
  $p$-torsion sheaf on $T'$, then the inverse image $u^{\ast}F'$ is a
  $p$-torsion sheaf on $T$.

\item If $u:T\rightarrow T'$ is a morphism of topoi, with $T$ and
  $T'$ locally of finite type and $u$ quasi-compact${}^\ast$, and if
  $F$ is a $p$-torsion sheaf on $T$, then the $R^qu_{\ast}F$ are
  $p$-torsion sheaves on $T'$ for every $q$.
\end{enumerate}
\end{enonce*}

\noindent
\emph{Proof}:
\begin{enumerate}
\renewcommand{\theenumi}{\roman{enumi}}
\renewcommand{\labelenumi}{(\theenumi)}
\item Clearly, if a sheaf $F$ is $p$-torsion, then $F$ is the sheaf
  associated with the presheaf $P$ whose group of sections $P(X)$ is
  the subgroup of $p$-torsion elements of $F(X)$. Conversely, let
  $F=\underline{a}P$ (II) be the sheaf associated with $P$,
  where $P$ takes values in the category of abelian
  $p$-torsion groups. From the exact sequence
$$
0\rightarrow {}_nP \rightarrow P\xrightarrow{n}P
$$
one deduces an exact sequence
$$
0\rightarrow \underline{a}({}_nP)\rightarrow F\xrightarrow{n}F{\quoi},
$$
whence\pageoriginale $\underline{a}({}_nP)={}_nF$. Since the functor $\underline{a}$
commutes with inductive limits, it follows from
$$
\frac{\lim(\text{\rm presheaf})}{\ass n\subset p}{}_nP\xrightarrow{\sim}P
$$
that one also has
$$
\varinjlim_{\ass n\subset p} {}_nF\xrightarrow{\sim}F{\quoi}.
$$

\item Put $P=\underrightarrow{\lim(\text{\rm presheaf})}{}_n F$. Then one has
  $P\subset F$, so $P$ is a separated presheaf, and consequently
$$
\underline{a}P(X)=\varinjlim (\ker(\prod_i P(X_i)\rightrightarrows)){\quoi},
$$
where, as usual, the limit is taken over the covering families
$\{X_i\rightarrow X\}$. But since $X$ is quasi-compact,
it is enough to take finite covering families. For these families,
$\prod_iP(X_i)$ is $p$-torsion; therefore so is the $\ker$. It follows
that $\underline{a}P(X)=F(X)$ is $p$-torsion.

\item To verify that $\varinjlim {}_nF\rightarrow F$ is bijective,
  it is enough to do so on the values at $X\in \Ob T$ for every
  $X$ in a family of generators. We are therefore reduced to (ii) for the
  first assertion. The second assertion is a consequence of

\item A consequence of (i), since $u^{\ast}(\underline{a}P')=\underline{a}(u\cdot P')$

\item An easy consequence of (iii).
\end{enumerate}

We shall also need a nonabelian variant. But there are technical
problems in the definition (I do not know whether they are serious),
and we shall therefore content ourselves with giving a definition for the
étale topology of a scheme.

\begin{enonce*}{Definition 1.3}\etiquette[1.3]{IX.def1.3}
Let\pageoriginale $p$ be a set of prime numbers. A group $G$ is called
an \emph{ind-p-group} if every finite subset of $G$
generates a finite subgroup of order $n$, with $\ass n\subset p$. This
is equivalent to saying that the finite subgroups $G_{\alpha}$
of order $n_{\alpha}$, with $\ass n_{\alpha}\subset p$, form a
filtered set, and that $G=\varinjlim G_{\alpha}$. If
$p=\PP$, one says \emph{ind-finite group} instead of ind-p-group.
\end{enonce*}

One immediately verifies the following
\begin{enonce*}{Proposition 1.4}\etiquette[1.4]{IX.prop1.4}
\begin{enumerate}
\renewcommand{\theenumi}{\roman{enumi}}
\renewcommand{\labelenumi}{(\theenumi)}
\item A subgroup as well as a quotient of an ind $p$-group is
  again an ind $p$-group.

\item A filtered inductive limit of ind-$p$-groups is an
  ind-$p$-group.

\item A finite projective limit of ind-p-groups is an ind-p-group.
\end{enumerate}
\end{enonce*}

\begin{enonce*}{Definition 1.5}\etiquette[1.5]{IX.def1.5}
Let $X$ be a scheme. A sheaf $F$ of groups on $X$ is called
a \emph{sheaf of ind $p$-groups} (\resp of \emph{ind-finite
  groups})
if, for every étale $U\rightarrow X$, with $U$ quasi-compact, $F(U)$ is an
ind $p$-group (\resp an ind-finite group).
\end{enonce*}

\begin{enonce*}{Proposition 1.6}\etiquette[1.6]{IX.prop1.6}
Let $F$ be a sheaf of groups on the scheme $X$.
\begin{enumerate}
\renewcommand{\theenumi}{\roman{enumi}}
\renewcommand{\labelenumi}{(\theenumi)}
\item $F$ is a sheaf of ind-$p$-groups if and only if, for
  every geometric point $\xi$ of $X$, the fiber $F_{\xi}$ is an
  ind-p-group.

\item If $f:X\rightarrow X'$ is a morphism and $F'$ is a
  sheaf of ind-p-groups on $X'$, then $f^{\ast}F'$ is a sheaf of
  ind-p-groups.

\item If $f:X\rightarrow X'$ is a quasi-compact morphism and $F$ is
  a sheaf of ind-p-groups on $X$, then $f_{\ast}F$ is a
  sheaf of ind-p-groups.
\end{enumerate}
\end{enonce*}

\noindent
\emph{Proof.}
(i) Suppose that $F$ is a sheaf of ind-p-groups, and let
  $\xi$ be a geometric point of $X$. Then $F_{\xi}=\varinjlim F(X')$, where
  $X'$ ranges over a\pageoriginale pseudo-filtered system of schemes
étale over $X$ (VIII \hyperref[VIII.prop3.9]{3.9}). It is evident that the affine $X'$ (hence
quasi-compact) form a cofinal set, and consequently that $F_{\xi}$ is an
inductive limit of ind $p$-groups, and hence is an ind-p-group by
\hyperref[IX.prop1.4]{1.4} (ii). Conversely, suppose that for every $\xi$ the fiber
$F_{\xi}$ is an ind-p-group, and let $U\rightarrow X$ be étale, with $U$
quasi-compact. Let $S\subset F(U)$ be a finite subset. For every
geometric point $\xi$ of $U$, the image of $S$ in $F_{\xi}$ generates
a finite group. Since a finite group is finitely presented, one
easily concludes that there exists a $U_{\xi}$ étale over $U$,
pointed at $\xi$, such that $S$ generates a finite group in $F(U_{\xi})$
(\cf VIII 4). A finite number of such $U_{\xi}$ form a covering
of $U$, say $\{U_1,\ldots,U_n\}$. Then the image of $S$ in $\prod
F(U_i)$ generates a finite subgroup, and since $F(U)\subset \prod
F(U_i)$, the result follows.

Assertion (ii) is trivial from (i), and (iii) is immediate.

% END INCLUDED SGA4_English_current_buildable_source_20260906/source/09/drafts/SGA4_IX_LINES_43_278_English.texfrag


% BEGIN INCLUDED SGA4_English_current_buildable_source_20260906/source/09/drafts/SGA4_IX_LINES_279_650_English.texfrag
\section{Constructible sheaves}\etiquette[2]{IX.sec2}

\setcounter{subsection}{-1}
\subsection{}\etiquette[2.0]{IX.subsec2.0}
Let $T$ be a topos. Recall that with every element $S\in \Ob (\Ens)$
one associates the sheaf $S_T$ associated with the presheaf $P$ such
that $P(X)=S$ for every $X\in \Ob (T)$. The object of $T$ representing
$S_T$ is
$$
\coprod_{s\in S}e= ``S\times e \mbox{''}{\quoi},
$$
where $e$ is the final object of $T$. The sheaf $S_T$ is called the
\emph{constant sheaf} with value $S$ (IV). In the evident way one
defines the notion of a \emph{constant sheaf of groups, constant sheaf
of $A$-modules} ($A$ a ring), and \emph{constant abelian sheaf}, with
value $M$. A morphism $f_T:S_T\rightarrow S'_T$ of constant sheaves is
said to be constant if it comes from a morphism of sets
$f:S\rightarrow S'$.

A sheaf $F$ is said to be \emph{locally constant} if there exists a
covering $\{e_i\rightarrow e\}$ of the final object such that $F$
becomes constant on each $e_i$. One\pageoriginale defines analogously
the notion of a locally constant \emph{sheaf of groups} or sheaf of
$A$-\emph{modules}, and of a \emph{locally constant morphism}
$f:F\rightarrow G$, when $F$ and $G$ are locally constant sheaves.
Finally, a locally constant sheaf of groups or of $A$-modules is said
to be \emph{of finite type} (\resp \emph{of finite presentation}) if
its local values are groups or $A$-modules of finite type (\resp of
finite presentation).

\begin{enonce*}{Lemma 2.1}\etiquette[2.1]{IX.lem2.1}
\begin{enumerate}
\renewcommand{\theenumi}{\roman{enumi}}
\renewcommand{\labelenumi}{(\theenumi)}
\item Let $f:F\rightarrow G$ be a morphism of locally constant sheaves
of sets on $T$, and suppose that the local values of $F$ are finite
sets. Then $f$ is locally constant.

\item Let $f:F\rightarrow G$ be a morphism of locally constant sheaves
of $A$-modules, with $F$ of finite type. Then $f$ is locally constant,
and the kernel and cokernel of $f$ are locally constant.

\item Let $0\rightarrow F'\rightarrow F\rightarrow F''\rightarrow 0$
be an exact sequence of sheaves of $A$-modules (\resp of groups), with
$F'$ and $F''$ locally constant, and $F''$ of finite presentation.
Then $F$ is locally constant.
\end{enumerate}
\end{enonce*}

Proof of (i) and (ii). We reduce without difficulty to the case in
which $F$ and $G$ are constant. Consider (ii): let $F=M_T$, $G=N_T$
($M$, $N$ being $A$-modules), and let $S\subset M$ be a finite set of
generators of $M$. Then $f$ is determined by its restriction to $S_T$,
and $f$ will be locally constant if $f|S_T$ is. Thus we are reduced to
(i) for the first assertion of (ii), and the second assertion is a
trivial consequence of the first.

It therefore remains to prove (i). Suppose that $F=S_T$, $G={S'}_T$
are constant, and also denote by $f$ the morphism of objects
$f:S\times e\rightarrow S'\times e$ (notation as above). Let
$f_s:e\rightarrow S'\times e$ $(s\in S)$ be the components of $f$:
since $S$ is finite, it plainly suffices to treat the case $f=f_s$,
that is, the case in which $S$ is a one-element set. Let
$X_{s'}\rightarrow e$ $(s'\in S')$ be the family of morphisms that make
the diagrams
$$
\xymatrix@C=2cm@R=1.5cm{
X_{s'} \ar[r] \ar[d]& e\ar[d]^{i_{s'}}\\
e   \ar[r]^{f}   & S'\times e
}
$$
cartesian, where\pageoriginale $i_{s'}:e\rightarrow S'\times e$ is
the ``inclusion'' into the $s'$-th component. The family $\{i_{s'}\}$
is trivially covering, and hence so is $\{X_{s'}\rightarrow e\}$. One
checks immediately (since coproducts in a topos are disjoint)
%% missing reference
that $f$ becomes constant on $X_{s'}$, and consequently $f$ is locally
constant.

Proof of (iii). Consider the case of sheaves of $A$-modules. We may
suppose $F'$ and $F''$ constant, with $F''$ having value $M''$, an
$A$-module of finite presentation. If $M''$ is free, and hence has a
finite basis $(e_i)_{1\leq i\leq n}$, then the $e_i$ define sections of
$F''$ and hence lift locally to sections of $F$; this shows that the
extension $F$ of $F''$ by $F'$ splits locally, which proves that $F$
is locally isomorphic to $F'\times F''$, and is therefore locally
constant. In the general case, choose a surjective homomorphism
$\varphi:L''\rightarrow M''$, with $L''$ free of finite type; the
kernel $R$ of $\varphi$ is of finite type, and setting
$L=\fprod{L''_T}{F}{M''_T}=\fprod{L''_T}{F}{F''}$, one sees on the one
hand that $L$ is an extension of $L''_T$ by $F'$, and is therefore
locally constant by what precedes, and on the other hand that there is
an exact sequence $0\rightarrow R_T\rightarrow L\rightarrow
F\rightarrow 0$, which, by (ii), implies that $F$ is locally constant.

From now on we restrict ourselves to the étale topoi of schemes.

\begin{enonce*}{Lemma 2.2}\etiquette[2.2]{IX.lem2.2}
Let $X$ be a scheme and let $F$ be a locally constant sheaf on $X$.
Then $F$ is represented by an étale $Y/X$. If, moreover, the fibers of
$F$ are finite (\resp and nonempty), then $Y$ is an étale covering
(\resp and is surjective).
\end{enonce*}

Proof. Descent (\SGA 1 IX 4.1 if $F$ has finite fibers, \SGA 3 X 5.4
in the general case).

\begin{enonce*}{Definition 2.3}\etiquette[2.3]{IX.def2.3}
A sheaf of sets (\resp of groups, \resp of $A$-modules) is said to be
constructible if, for every affine open $U\subset X$, there exists a
decomposition of $U$ as a union of finitely many constructible
{\rm (\EGA $0_{{\rm III}}$ 9.1.2)} reduced locally closed subschemes
$U_i$ such that the sheaf induced by $F$\pageoriginale on each $U_i$
is locally constant, with finite value (\resp finite, \resp of finite
presentation\footnote[1]{When $A$ is not assumed noetherian, for the
purposes of Homological Algebra it is preferable to require here,
instead of finite presentation alone for the $A$-module $M$, that it
have a left resolution by free $A$-modules of finite type. The notion
introduced in (\hyperref[IX.def2.3]{2.3}) should then be called:
$F$ is 1-constructible (more generally, one would define
$n$-constructibility in the evident way, for every integer $n\geq 0$).
For the present Exposé we shall provisionally retain the terminology in
the form (\hyperref[IX.def2.3]{2.3}), which is chiefly reasonable when
$A$ is noetherian, in which case it coincides with that just indicated.
Compare \SGA 6 I for finiteness notions in nonnoetherian cases.}).
\end{enonce*}

\setcounter{subsection}{3}
\subsubsection{}\etiquette[2.3.1]{IX.subsubsec2.3.1}
The finiteness hypothesis on the local values amounts to saying that
for every geometric point $\xi$, the fiber $F_{\xi}$ is finite (\resp
finite, \resp of finite presentation). It follows immediately from
\hyperref[IX.lem2.1]{2.1} (i) that a sheaf of groups $F$ is
constructible if and only if the underlying sheaf of sets of $F$ is
constructible.

Care must be taken that if $F$ is an abelian sheaf, it is constructible
as a sheaf of groups (or equivalently, as a sheaf of sets) if and only
if it is constructible as a sheaf of $\ZZ$-modules, and if
\emph{in addition} its fibers are finite. Thus, $\ZZ_X$ is
constructible as a $\ZZ$-module, but not as a sheaf of groups.

\begin{enonce*}{Proposition 2.4}\etiquette[2.4]{IX.prop2.4}
\begin{enumerate}
\renewcommand{\theenumi}{\roman{enumi}}
\renewcommand{\labelenumi}{(\theenumi)}
\item Let $X$ be quasi-compact and quasi-separated. Then a sheaf \\
(\resp sheaf of groups, \resp sheaf of $A$-modules) is constructible if
and only if there exists a decomposition of $X$ as a union of
constructible locally closed subsets such that $F$ becomes locally
constant and finite (\resp finite, \resp of finite presentation) on
each $X_i$.

\item To verify that $F$ is constructible, it is enough to verify that
$F| U_i$ is constructible for the $U_i$ of a given open covering of
$X$.

\item Let $f:Y\rightarrow X$ be a morphism. Then $f^{\ast}F$ is
constructible if $F$ is.

\item The\pageoriginale set of points at which the fiber of a
constructible sheaf is nonempty (\resp nonzero) is a locally
constructible subset {\rm (\EGA $0_{\rm III}$ 9.1.2)}.

\item Let $X$ be locally noetherian. Then a sheaf of sets
(\resp\ldots) $F$ on $X$ is constructible if and only if for every
$x\in X$ there exists a nonempty open subset $U$ of the closure $\bar{x}$
of $x$ such that $F$ induces a locally constant sheaf with finite value
(\resp\ldots) on $U$.
\end{enumerate}
\end{enonce*}

Proof. Plainly, if $f:Y\rightarrow X$ is a morphism and if there exists
a decomposition of $X$ as a union of constructible locally closed
subsets $X_i$ such that $F$ becomes locally constant on each $X_i$,
the same is true of $Y$ and $f^{\ast}F$ (\cf \EGA IV 1.8.2). This
proves the implication $\Leftarrow$ of (i). Now suppose that $X$ is
quasi-compact and quasi-separated, and let $\{U_1,\ldots,U_n\}$ be an
affine open covering of $X$ such that $F|U_i$ is constructible for
every $i$. To find a decomposition of $X$ as in (i), it is enough to
do so for $U_n$ and for $Y=X-U_n$ (with the induced reduced structure),
since $U_n$ and $Y$ are constructible in $X$ by virtue of the
hypothesis that ``$X$ is quasi-separated''. Now a decomposition exists
for
$U_n$ by the definition of constructibility, and $Y$ is the union of
$n-1$ affine opens $V_i=Y\cap U_i$ $(i=1,\ldots,n-1)$ with
$F| V_i$ constructible, whence result (i) follows by induction on $n$.
The same reasoning proves (ii); indeed, the hypothesis implies that
$F| U$ is constructible for every ``sufficiently small'' affine open
$U$, and an arbitrary affine open $V$ can be covered by finitely many
such opens. Assertion (ii) shows that the notion of constructibility is
local on $X$, and (iii), (iv) follow immediately. The implication
$\Rightarrow$ of (v) is immediate and does not depend on the
noetherian hypothesis. If $X$ is noetherian, one proves the implication
$\Leftarrow$ by the usual ``noetherian induction''.

\begin{enonce*}{Proposition 2.5}\etiquette[2.5]{IX.prop2.5}
Let $X$ be quasi-compact and quasi-separated, and let $F$ be a
constructible sheaf of groups or of $A$-modules on $X$. Then there
exists a finite filtration of $F$ whose successive quotients are of the
form $i_{!}G$, where $i:U\rightarrow X$ is the inclusion of a
constructible locally closed subset and where $G$ is locally constant
and constructible on $U$. If $X$ is noetherian, there exists such a
filtration with the $U$ irreducible.
\end{enonce*}

Proof.\pageoriginale By \hyperref[IX.prop2.4]{2.4} (i), there exists a
decomposition of $X$ as a finite union of constructible locally closed
subsets $X_i$ such that each $F|X_i$ is locally constant and finite
(\resp of finite presentation). Write $X_i=U_i\cap \complement V_i$,
where $U_i$ and $V_i$ are constructible opens, and let $N$ be the
number of opens of $X$ in the subtopology $T$ generated by
$\{U_i,V_i\}$. We argue by induction on $N$: if $W$ is a minimal
nonempty element of $T$, it is plain that $F|W$ is locally constant and
of finite presentation. Let $i:W\rightarrow X$ be the inclusion
morphism and let $Y=X-W$. We have the exact sequence
$$
0\longrightarrow i_{!}i^{\ast}F\longrightarrow F\longrightarrow F|
Y\longrightarrow 0{\quoi},
$$
and are thereby reduced to the same assertion for $Y$ and the sheaf
$F| Y$, where the induction hypothesis applies.

\begin{enonce*}{Proposition 2.6}\etiquette[2.6]{IX.prop2.6}
Suppose that $A$ is a noetherian ring.
\begin{itemize}
\item[(i)] A finite projective (\resp inductive) limit of constructible
sheaves of $A$-modules or sets is constructible. In particular, the
kernel, cokernel, and image of a morphism of constructible sheaves of
$A$-modules are constructible.

\item[(i bis)] A finite projective limit of constructible sheaves of
sets (\resp of groups) is constructible, and if $f:F\rightarrow G$ is
a morphism of constructible sheaves of groups, the kernel, image, and
cokernel (if $\Iim(f)$ is a normal subgroup) are constructible. A finite
inductive limit of constructible sheaves of sets is constructible.

\item[(ii)] Let $0\rightarrow M'\rightarrow M\rightarrow
M''\rightarrow 0$ be an exact sequence of sheaves of $A$-modules
(\resp of groups), with $M'$, $M''$ constructible. Then $M$ is
constructible.

\item[(iii)] Let $M_1\rightarrow M_2\rightarrow M_3\rightarrow
M_4\rightarrow M_5$ be an exact sequence of sheaves of $A$-modules.
Then if $M_i$ is constructible for $i=1,2,4,5$, the same holds for
$i=3$.
\end{itemize}
\end{enonce*}

\emph{Proof.} (i) We\pageoriginale leave the assertions for sets and
groups to the reader. For sheaves of $A$-modules, it is enough to prove
that the kernel and cokernel of a morphism $f:F\rightarrow G$ of
constructible sheaves of $A$-modules are constructible. The assertion
is local on $X$, and we may therefore suppose that $X$ is affine, hence
quasi-compact and quasi-separated. By \hyperref[IX.prop2.4]{2.4} (i),
we then reduce to the case in which $F$ and $G$ are locally constant,
and conclude by using \hyperref[IX.lem2.1]{2.1} (ii) and the noetherian
hypothesis on $A$. Assertion (ii) is proved analogously using
\hyperref[IX.lem2.1]{2.1} (iii). Assertion (iii) follows immediately
from (i) and (ii).

\begin{enonce*}{Proposition 2.7}\etiquette[2.7]{IX.prop2.7}
Let $X$ be a quasi-compact and quasi-separated scheme, and let $F$ be a
sheaf of sets (\resp of $A$-modules). In order that $F$ be
constructible, it is necessary and sufficient that it be isomorphic to
the cokernel of a pair of morphisms $H\rightrightarrows G$, where $H$
and $G$ are sheaves of sets representable by schemes étale and of
finite presentation over $X$ (\resp that $F$ be isomorphic to the
cokernel of a homomorphism $𝒜_{V,X}\rightarrow 𝒜_{U,X}$, where $U$ and
$V$ are two schemes étale and of finite presentation over $X$).
\end{enonce*}

%% A_V [and not 𝒜] in 2.9 \XXX

Sufficiency follows easily from \hyperref[IX.prop2.6]{2.6} (i bis).

Suppose that $F$ is a constructible sheaf of sets. Using the existence
of infinite coproducts in the étale site over $X$, one sees that one can
find an epimorphism $G\rightarrow F$, with $G$ representable by a
scheme étale over $X$, which may moreover be supposed to be a
coproduct of affine schemes $G_i$ $(i\in I)$, and hence separated over
$X$. For every finite subset $J$ of the index set $I$, let $G_J$ be
the coproduct of the $G_i$ for $i\in J$, and let $F_J$ be its image in
$F$. Since $G_J$ and $F$ are constructible, so is $F_J$
(\hyperref[IX.prop2.6]{2.6} (i bis)); hence the set $X_J$ of $x\in X$
such that the fibers of $F$ and $F_J$ at a geometric point of $X$ over
$x$ are equal is locally constructible (\hyperref[IX.prop2.4]{2.4}
(iv). Since the family of $X_J$ is increasing and has union $X$, and
since $X$ is quasi-compact, one of the $X_J$ equals $X$
(\EGA IV 1.9.9). This shows, after replacing $G$ by $G_J$, that we may
suppose $G$ affine, and hence separated over $X$, and quasi-compact over
$X$ since $X$ is quasi-separated (\EGA IV 1.2.4). Let then
$H=\fprod{G}{G}{F}$. It is a subsheaf of $\fprod{G}{G}{X}$, which is
representable, and is therefore itself representable (VIII \hyperref[VIII.prop6.1]{6.1}).
Since $H$ is moreover constructible (\hyperref[IX.prop2.6]{2.6}
(i bis)),\pageoriginale it follows again from the preceding argument
that it is quasi-compact, and therefore (being separated over $X$) of
finite presentation over $X$. This proves the first assertion of
\hyperref[IX.prop2.7]{2.7}, and the second is proved by the same method
(NB it is not necessary that $A$ be noetherian). Let us note in passing
that the proof also proves the following corollary:

\begin{enonce*}{Corollary 2.7.1}\etiquette[2.7.1]{IX.cor2.7.1}
Let $X$ be a scheme and let $F$ be a scheme étale over $X$. In order
that the corresponding étale sheaf on $X$ be constructible, it is
necessary and sufficient that $F$ be of finite presentation over $X$.
\end{enonce*}

Proposition \hyperref[IX.prop2.7]{2.7} will be useful here chiefly for
deducing certain limit-passage properties for constructible sheaves:

\begin{enonce*}{Corollary 2.7.2}\etiquette[2.7.2]{IX.cor2.7.2}
Let $X$ be a quasi-compact and quasi-separated scheme. Then every
sheaf of sets (\resp of ind-p-groups, \resp of $A$-modules) on $X$ is
the inductive limit of a filtered inductive system of constructible
sheaves of sets (\resp\ldots).
\end{enonce*}

If $F$ is a sheaf of sets, retain the notation of the proof of
\hyperref[IX.prop2.7]{2.7}, and let $L$ be the set of pairs $(J,H')$,
where $J$ is a finite subset of $I$ and $H'$ is a quasi-compact open
subset of $H_J=\fprod{G_J}{G_J}{F}$. Order $L$ in the evident way, and
observe that, since $F$ is the limit of the $F_J$, and each $F_J$ is
the limit of the $\Coker(H'\rightarrow G_J)$ for the quasi-compact
opens $H'$ of $H_J$, $F$ is the inductive limit of the inductive system
of the $F_{\ell}=\Coker(H'\rightarrow G_J)$ for
$\ell=(H',J)\in L$, proving the case of sets. The case of sheaves of
$A$-modules is treated in essentially the same way, beginning by
representing $F$ as the cokernel of a homomorphism
$𝒜_{V,X}\rightarrow 𝒜_{U,X}$, with $U$, $V$ schemes étale and
separated over $X$. The case of ind-$\sheaf{p}$-groups requires a little
more attention. Retaining the preceding notation, denote by $L(G_J)$
the free group sheaf generated by the sheaf of sets $G_J$, by $N_J$
the kernel of the homomorphism $L(G_J)\rightarrow F$ deduced from
$G_J\rightarrow F$, and by $L$ the set of pairs $(J,N')$, where $N'$
is a subsheaf of \emph{sets} of $N_J$ which is ``quasi-compact'', \ie
such that there exists an epimorphism $P\rightarrow N'$, with $P$
representable by a quasi-compact $X$-scheme. Order $L$ by declaring
that $(J,N')$ is smaller than $(J_1, N'_1)$ if $J\subset J_1$ and if
the canonical\pageoriginale homomorphism
$L(G_J)\rightarrow L(G_{J_1})$ maps $N'$ into $N'_1$. It is immediate
that $L$ is filtered, and that one obtains an inductive system of
sheaves of groups $F_{\ell}$, indexed by $L$, by taking, for
$\ell=(J,N')$, $F_{\ell}=$ the quotient sheaf of groups of $L(G_J)$ by
the invariant subsheaf of groups generated by $N'$. It is moreover
immediate that $F$ is isomorphic to the inductive limit of the
$F_{\ell}$, and it remains to prove that there exists a subset $L'$ of
$L$, cofinal in $L$, such that for $\ell\in L'$, $F_{\ell}$ is a
constructible sheaf of $\sheaf{p}$-groups. For this, it is enough to
prove that for every $J$ one can find a subsheaf of sets $N'$ of $N_J$
which is quasi-compact, contains a given quasi-compact subsheaf of sets
$N'_0$, and is such that the corresponding $F_{\ell}$ is a
constructible sheaf of $\sheaf{p}$-groups. Note first that it is enough
for this that the fibers of $F_{\ell}$ be finite $\sheaf{p}$-groups:
then $F_{\ell}$ will automatically be constructible, as the reader
will readily verify. On the other hand, the reader may also verify that
the set $X_{N'}$ of points $x$ of $X$ such that the fiber of
$F_{\ell}$ at a geometric point over $x$ is a finite
$\sheaf{p}$-group is locally constructible. Moreover, the subsets
$X_{N'}$ of $X$ are increasing functions of $N'$, and their union is
all of $X$: this last point follows from the fact that the inductive
limit of the $F_{\ell}$, for $\ell=(J,N')$ with $J$ fixed, is the
subsheaf of groups $\bar{F_J}$ of $F$ generated by the subsheaf of sets
$F_J$, which is quasi-compact, and hence $\bar{F_J}$ has fibers that
are finite $\sheaf{p}$-groups. On the other hand, in a free group on a
finite number of generators, every subgroup of finite index is a free group on a
finite number of generators; hence, if it is expressed as a filtered union of its
finite subsets, one of these generates the whole subgroup. It
follows readily that the union of the

% END INCLUDED SGA4_English_current_buildable_source_20260906/source/09/drafts/SGA4_IX_LINES_279_650_English.texfrag


% BEGIN INCLUDED SGA4_English_current_buildable_source_20260906/source/09/drafts/SGA4_IX_LINES_651_1136_English.texfrag
$X_{N'}$ is indeed $X$. Using (\EGA IV 1.9.9) once again, one finds
that one of the $X_{N'}$ is equal to $X$, which completes the proof of
\hyperref[IX.cor2.7.2]{2.7.2}.

\setcounter{subsection}{7}
\setcounter{subsubsection}{2}
\paragraph{}\etiquette[2.7.2.1]{IX.par2.7.2.1}
Note that the somewhat laborious proof in the case of group schemes
brings to light the lack of a convenient sorites for constructibility
conditions on sheaves of groups with fibers not necessarily finite; we
apologize to the reader, inviting them to fill this gap by their own
means if necessary, and point out by way of consolation that in
\hyperref[IX.prop2.9]{2.9} (iii) we obtain a simpler proof when $X$ is
Noetherian.

\begin{enonce*}{Corollary 2.7.3}\etiquette[2.7.3]{IX.cor2.7.3}
Let\pageoriginale $X$ be a quasi-compact and quasi-separated scheme,
and let $F$ be a constructible sheaf of sets (\resp of finite
ind-groups, \resp of $A$-modules). Then the functor $\Hom(F,G)$, where
$G$ is a variable sheaf of sets (\resp\ldots), commutes with inductive
limits in $G$. When $F$ is a sheaf of $A$-modules, with $A$
Noetherian, the functors $\Ext^{i}(X;F,G)$ in the $A$-module $G$ also
commute with inductive limits.
\end{enonce*}

This follows from \hyperref[IX.prop2.7]{2.7} and (VII \hyperref[VII.prop3.3]{3.3}) by the usual
arguments, which are left to the reader. Similarly, the conjunction of
\hyperref[IX.prop2.7]{2.7} and the results of (VII 5) easily gives:

\begin{enonce*}{Corollary 2.7.4}\etiquette[2.7.4]{IX.cor2.7.4}
Let $I$ be a filtered category, and let $i\mapsto
X_i:I^0\rightarrow (\Sch)$ be a functor which transforms arrows into
affine morphisms and objects into quasi-compact and quasi-separated
schemes. For every scheme $Y$, let $F_{c}(Y)$ be the category of
constructible sheaves of sets (\resp of ind-$\sheaf{p}$-groups, \resp
of $A$-modules) on $Y$. Then there is an equivalence of categories
$$
F_{c}(X)\xleftarrow{\approx}\varinjlim_i F_{c}(X_i){\quoi},
$$
where $X=\varprojlim_i X_i$ {\rm (VII 5)} and where the second member
denotes the $\varinjlim$ category of fibered categories, defined
in \hyperref[VI.def6.3]{VI 6.3}. In particular, every constructible sheaf of sets
(\resp\ldots) on $X$ is isomorphic to the inverse image of a sheaf of
the same kind on one of the $X_i$.
\end{enonce*}

\begin{enonce*}{Proposition 2.8}\etiquette[2.8]{IX.prop2.8}
Let $f:X\rightarrow Y$ be a surjective morphism locally of finite
presentation, and let $F$ be a sheaf of sets (\resp\ldots) on $Y$.
Then $F$ is constructible if and only if $f^{\ast}(F)$ is
constructible.
\end{enonce*}

We shall use the following

\begin{enonce*}{Lemma 2.8.1}\etiquette[2.8.1]{IX.lem2.8.1}
Let $f:X\rightarrow Y$ be a surjective morphism locally of finite
presentation, with $Y$ quasi-compact and quasi-separated. Then there
exists a finite partition of $Y$ into subschemes $Y_i$ of finite
presentation over $Y$ (in particular,\pageoriginale each $Y_i$ is
constructible in $Y$), and, for every $i$, finite surjective morphisms
$Y''_i\xrightarrow{g_i}Y'_i\xrightarrow{h_i}Y_i$, with $h_i$ étale
and $g_i$ locally free (\ie flat and of finite presentation, in
addition to the condition that $g_i$ be finite) and radicial, and
finally a $Y$-morphism $Y''_i\rightarrow X$
\footnote[1]{This assertion is also given in \EGA IV 17.16.4.}.
\end{enonce*}

Using the fact that $Y$ is quasi-compact and quasi-separated, we are
easily reduced to the case where $Y$ is quasi-affine, hence separated
(by using a finite affine open covering $(U_i)_{1\leq i\leq n}$ of $Y$
and considering the $Y_i=U_i-U_i\cap \bigcup_{j< i} U_j$), and then to
the case where $Y$ is affine (by the same argument). The scheme $X$ is
a union of affine open subsets $X_i$, and the image $T_i$ of $X_i$ in
$Y$ is constructible (\EGA IV 1.8.4); the union of the $T_i$ is $Y$
since $f$ is surjective, whence it follows that $Y$ is the union of a
finite number of the $T_i$ (\EGA IV 1.9.9), so that, replacing $X$ by
the sum scheme of the corresponding $X_i$, we may suppose that $X$ is
affine. The standard limit procedure (\EGA IV 8.1.2 c)) then allows
us to suppose in addition that $Y$ is Noetherian. The usual Noetherian
induction, together with the limit procedure (\EGA IV 8.1.2 a)),
reduces us to the case where $Y$ is the spectrum of a field $k$. Since
then $X\neq ∅$, there exists a closed point $x$ in $X$, corresponding
to a finite extension $k'$ of $k$, itself a radicial extension of a
separable extension $k'_s$ of $k$. We then take
$Y'=\Spec(k'_s)$, $Y''=\Spec(k')$, which completes the proof of
\hyperref[IX.lem2.8.1]{2.8.1}.

Necessity in \hyperref[IX.prop2.8]{2.8} being trivial
(\hyperref[IX.prop2.4]{2.4} (iii)), let us prove sufficiency. By the
preceding lemma we are reduced to the case where $f$ is of the form
$gh$, where $g$ and $h$ satisfy the conditions stated for $g_i$, $h_i$
in the lemma. We may therefore suppose successively either that $f$ is
finite radicial, a trivial case by (VIII \hyperref[VIII.thm1.1]{1.1}), or that $f$ is finite
étale. Note moreover that, using the hypothesis that $f^{\ast}(F)$ is
constructible and the definition of constructibility, we could suppose
(after replacing $X$ by an $X$-prescheme
$X'=\coprod_i X_i$ of finite presentation and with surjective
structural morphism) that $f^{\ast}(F)$ is already locally constant.
But when $f$ is étale and surjective, this implies that $F$ is locally
constant. Since in addition the hypotheses made on the fibers of
$f^{\ast}(F)$ remain valid for those of $F$, with $f$ surjective, it
indeed follows that $F$ is constructible, which completes the proof of
\hyperref[IX.prop2.8]{2.8}.

\begin{enonce*}{Proposition 2.9}\etiquette[2.9]{IX.prop2.9}
Let\pageoriginale $X$ be a Noetherian scheme, let $A$ be a Noetherian
ring, and let $p$ be a set of prime numbers.
\begin{enumerate}
\renewcommand{\theenumi}{\roman{enumi}}
\renewcommand{\labelenumi}{(\theenumi)}
\item The category of sheaves of sets (\resp of ind $p$-groups, \resp
of $A$-modules) on $X$ is locally Noetherian, that is, it has a set of
generators formed of Noetherian objects.

\item A sheaf $F$ of sets (\resp of ind $p$-groups, \resp of
$A$-modules) on $X$ is constructible if and only if it is Noetherian.
If $F$ is a sheaf of $A$-modules, then $F$ is constructible if and only
if it is a quotient of a finite sum of sheaves of the form $A_{U/X}$,
where $U\rightarrow X$ is étale and of finite type (notation of
{\rm \hyperref[IV.subsec2.4]{IV.2.4}}: $A_{U/X}=A_U$).

%% check the reference IV.2.4 [seems incorrect to me] \XXX

\item The constructible subsheaves of a sheaf $F$ of sets
(\resp\ldots) form an inductive system, and $F$ is the inductive limit
of these subsheaves.
\end{enumerate}
\end{enonce*}

Proof. We leave the assertion for sets to the reader. First consider
the case of a sheaf of $A$-modules. For (i), one must find a set of
generators whose elements are Noetherian objects. But the sheaves of
the form $A_{U/X}$, with $U\rightarrow X$ étale and of finite type,
form a set of generators which are easily seen to be constructible,
and it therefore suffices to prove the following lemma:

\begin{enonce*}{Lemma 2.10}\etiquette[2.10]{IX.lem2.10}
Let $X$ be Noetherian. Then every constructible $A$-module is
Noetherian.
\end{enonce*}

Proof of the lemma.

By Noetherian induction, we suppose the lemma true for every subscheme
$X'$ of $X$ distinct from $X$. Choose a nonempty open subset $X_0$ of
$X$ such that $F$ induces a locally constant sheaf $F_0$ on $X_0$.

Now let $\{F_i\}$, $i\in \NN$, be an increasing sequence of subsheaves
of $F_0$. Since, for a geometric point $P$ which is generic for an
irreducible component of $X_0$, the fiber of $F_0$ at $P$ is of finite
type, the sequence of fibers $(F_i)_P$\pageoriginale is stationary,
and we may suppose that it is constant. Now let $Q$ be a geometric
point of $X_0$ which is a \emph{specialization} of $P$ (VIII \hyperref[VIII.def7.2]{7.2}).
Since $F_0$ is locally constant on $X_0$, one sees immediately that the
specialization morphism $F_{0_Q}\rightarrow F_{0_P}$ (VIII.7) is
bijective, hence $(F_i)_Q\rightarrow (F_i)_P$ is \emph{injective} for
every $i$.

Let $V/X$ be étale and pointed at $P$, and let $s_1,\ldots,s_n$ be
sections $\in F_1(V)$ which generate $(F_1)_P$. Then the $s_i$ generate
$(F_i)_Q$ for every $i$ and every geometric point $Q$ which is a
specialization of $P$ above $V$. The sequence of sheaves $F_i$ is thus
constant in a neighborhood of the image of $P$, say in $X'\neq ∅$.
Let $Y=X-X'$; it therefore suffices to prove that the sequence of the
$F_i| Y$ is stationary, which follows from the Noetherian induction
hypothesis.

The assertions (ii) for $A$-modules are now immediate: since the
$A_{U/X}$ are Noetherian generators, one plainly has $F$ Noetherian
$\Leftrightarrow F$ is a quotient of a finite sum of sheaves
$A_{U/X}$. But $A_{U/X}$ is constructible, and consequently, in view
of \hyperref[IX.prop2.6]{2.6}, $F$ Noetherian $\Rightarrow F$
constructible; the converse implication having already been proved
(\hyperref[IX.lem2.10]{2.10}), this establishes (ii). Finally,
assertion (iii) for $A$-modules is a property of locally Noetherian
categories (\cf [1], Ch. IV).

Now consider the case of ind-$p$-groups. First note that a
constructible sheaf of groups $F$ is certainly Noetherian. Indeed,
since the fibers are finite, we may use the same argument as for
$A$-modules. We shall now prove assertion (iii) directly, which will
imply (i)---indeed, it will show that constructible sheaves form a set
of generators. Moreover, any Noetherian sheaf will then necessarily be
constructible, whence (ii).

To prove (iii), it plainly suffices to prove that a subsheaf (of
groups) $S$ of a sheaf $F$ of ind $p$-groups generated by a finite
number of sections $s_i\in F(U_i)$, with $U_i\rightarrow X$ étale and
of finite type, is constructible \footnote[1]{Recall the definition of
this sheaf: let $L_{U/X}$ be the sheaf representing the functor
$F\mapsto F(U)$ in the category of groups. The section
$s_i\in F(U_i)$ induces a morphism $L_{U_i/X}\rightarrow F$, and $S$ is
the image of the sum morphism $\amalg_i L_{U_i/X}\rightarrow F$. It is
the smallest subsheaf which ``contains'' the sections $s_i$.}. (Since
a\pageoriginale constructible sheaf is Noetherian, it is generated by
a finite number of sections.) By Noetherian induction, we may suppose
that this is true for every closed subscheme $Y$ of $X$ distinct from
$X$. Now the notion of a subsheaf generated by sections commutes with
the inverse image of sheaves; indeed, let $L_{U/X}$ be the sheaf which
represents the functor $F\mapsto F(U)$ in the category of groups. Then
if $f:Y\rightarrow X$, one has
$f^{\ast}(L_{U/X})=L_{\fprod{U}{Y}{X}/Y}$. Since $S$ is the image of
the sum of the $L_{U_i/X}$, and since $f^{\ast}$ commutes with sums, it
is clear that $f^{\ast}S$ is the subsheaf of $f^{\ast}F$ generated by
the sections $f^{\ast}(s_i)$.

We may therefore suppose that $S$ induces a constructible sheaf on
every closed subset $Y$ distinct from $X$, and it thus suffices
(\hyperref[IX.prop2.4]{2.4} (i)) to prove that $S$ induces a
constructible sheaf on a suitable nonempty open subset of $X$.
Replacing $X$ by a sufficiently small open subset, we may suppose that
every $U_i$ is an étale covering of $X$. It suffices to prove that $S$
is then locally constant, with finite fibers, on a suitable nonempty
open subset. This assertion is local on $X$ for the étale topology, so
we may suppose that every $U_i$ is completely split, say
$U_i=\amalg_{n_i}X$ (the sum of $n_i$ copies of $X$). Let $s_{ij}$
$(j=1,\ldots,n_i)$ be the sections of $F(X)$ such that
$s_i$ $(s_{i1},\ldots,s_{in_i})$. Plainly, $S$ is none other than the
subsheaf of $F$ generated by the $s_{ij}\in F(X)$. Since $F(X)$ is an
ind $p$-group (\hyperref[IX.def1.5]{1.5}), the $s_{ij}$ generate a
finite subgroup \hyperref[IX.def1.3]{1.3}. Thus we may suppose that the
set of the $s_i$ is a finite subgroup of $F(X)$; but then $S$ is
identical to the sheaf \emph{of sets} generated by the $s_i$, which is
constructible by \hyperref[IX.prop2.6]{2.6} (i) as a sheaf of sets,
hence as a sheaf of groups, \cqfd

In the following two propositions, we give criteria for a sheaf to be
locally constant or constructible, using specialization homomorphisms
(VIII \hyperref[VIII.subsec7.7]{7.7}).

\begin{enonce*}{Proposition 2.11}\etiquette[2.11]{IX.prop2.11}
Let $X$ be a scheme, and let $F$ be a constructible sheaf of sets
(\resp of ind-finite groups, \resp of $A$-modules); in the case of
sets,\pageoriginale suppose in addition that $F$ has finite fibers. Let
$x\in X$, and let $\bar{x}$ be a geometric point over $x$. In order that $F$
be locally constant in a neighborhood of $x$, it is necessary and
sufficient that, for every geometric point $\bar{x}'$ which is a
generization of $\bar{x}$, and every specialization arrow
$\bar{x'}\rightarrow \bar{x}$ {\rm (VIII \hyperref[VIII.def7.2]{7.2})}, the specialization
homomorphism {\rm (VIII \hyperref[VIII.subsec7.7]{7.7})} $F_{\bar{x}}\rightarrow F_{\bar{x'}}$ be
an isomorphism.
\end{enonce*}

The case of a sheaf of groups being a special case of that of a sheaf
of sets, we confine ourselves to the latter; the case of a sheaf of
modules is treated in essentially the same way. Necessity of the
condition being trivial, we need only establish sufficiency. Since the
fiber $I=F_{\bar{x}}$ is finite, after replacing $X$ by a suitable
scheme $X'$ étale over $X$, we may suppose that one can find a
homomorphism $u:I_X\rightarrow F$ from the constant sheaf $I_X$ to
$F$, inducing an isomorphism on the fibers at $\bar{x}$. Using the
constructibility of both sheaves and \hyperref[IX.prop2.6]{2.6}, one
easily sees that the set $Z$ of points $z$ of $X$ such that the
homomorphism induced on the fibers at a geometric point $\bar{z}$ over
$z$ is bijective is a locally constructible subset of $X$. The
hypothesis immediately implies that it contains the generizations of
$x$, hence (\EGA IV 1.10.1) it is a neighborhood of $x$, which proves
that $F$ is locally constant (indeed, even constant in this case) in a
neighborhood of $x$,\cqfd

\begin{enonce*}{Definition 2.12}\etiquette[2.12]{IX.def2.12}
\footnote[1]{The terminology is due to S. Lubkin.}: A scheme $X$ is
said to be path-connected if for every pair $P$, $Q$ of geometric
points of $X$, there exist geometric points $P=P_0,\ldots,P_n$;
$Q_1,\ldots,Q_n=Q$ and specializations {\rm (VIII \hyperref[VIII.def7.2]{7.2})}
$P_i\rightarrow Q_i$ and $P_{i-1}\rightarrow Q_i$
$(i=1,\ldots,n)$. One says that $X$ is locally path-connected (for the
étale topology) if for every étale $U\rightarrow X$ there exists an
étale covering $\{U_i\rightarrow U\}$ of $U$ such that $U_i$ is
path-connected for every $i$.
\end{enonce*}

One checks immediately that a locally Noetherian prescheme is locally
path-connected.

\begin{enonce*}{Proposition 2.13}\etiquette[2.13]{IX.prop2.13}
\begin{enumerate}
\renewcommand{\theenumi}{\roman{enumi}}
\renewcommand{\labelenumi}{(\theenumi)}
\item Let\pageoriginale $X$ be locally path-connected and let $F$ be a
sheaf of sets (\resp of groups, \resp of $A$-modules) on $X$. Suppose
that the fibers of $F$ are finite (\resp finite, \resp of finite
presentation). Then $F$ is locally constant if and only if, for every
specialization $P\rightarrow Q$ of geometric points, the specialization
morphism {\rm (VIII \hyperref[VIII.subsec7.7]{7.7})} $F_Q\rightarrow F_P$ is bijective.

\item Suppose that $X$ is locally Noetherian and that the fibers of $F$
are finite (\resp finite, \resp of finite presentation). Then, if for
every specialization $P\rightarrow Q$ the specialization morphism
$F_Q\rightarrow F_P$ is injective, $F$ is constructible.

\item Suppose that $X$ is locally Noetherian and that the fibers of the
sheaf of sets $F$ are finite. Let $c:X\rightarrow \NN$ be the function
which associates with $x\in X$ the number of elements of the fiber
$F_{\bar{x}}$ of $F$ at a geometric point $\bar{x}$ above $x$. Then
$F$ is constructible if and only if $c$ is a constructible function,
\ie if and only if $c^{-1}(n)$ is constructible for every
$n\in \NN$.
\end{enumerate}
\end{enonce*}

\emph{Proof.}
\begin{enumerate}
\renewcommand{\theenumi}{\roman{enumi}}
\renewcommand{\labelenumi}{(\theenumi)}
\item Let $P$ be a geometric point of $X$. We shall find an étale
neighborhood of $P$, that is, an étale $U\rightarrow X$ pointed at
$P$, such that $F$ becomes constant on $U$. Let $V'$ be an étale
neighborhood of $P$ such that $F(V')\rightarrow F_P$ is surjective.
Such a $V'$ exists by the finiteness hypothesis. Let $U\rightarrow V'$
be an étale neighborhood of $P$ in $V'$ such that $U$ is
path-connected. Then $F(U)\rightarrow F_P$ is still surjective. Choose
a subset $S\subset F(U)$ which maps bijectively onto $F_P$. If $F$ is
a sheaf of groups (\resp of $A$-modules), one easily sees that, after
changing $U$ if necessary, one may find such a set which is moreover a
subgroup (\resp submodule) of $F(U)$. This gives a morphism
$S_U\rightarrow F| U$, where $S_U$ is the constant sheaf with value
$S$. One has
$$
(S_U)_P\xrightarrow{\sim}F_P{\quoi}.
$$
Since\pageoriginale every specialization morphism on the fibers of $F$
(\resp $S_U$) is bijective, and since $U$ is path-connected, it follows
that $S_U\xrightarrow{\sim}F| U$, whence the result.

\item The assertion is local, and we may therefore suppose that $X$ is
Noetherian. Let $P$ be a geometric point above a maximal point of $X$.
There exists an irreducible étale neighborhood, hence a path-connected
one, $V\rightarrow X$ of $P$ such that $F(V)$ maps surjectively onto
$F_P$. It then maps surjectively onto $F_Q$ for every geometric point
$Q$ of $V$, or, what amounts to the same thing, for every geometric
point of the image $U$ of $V$ in $X$, since such a point is a
specialization of $P$ and $F_Q\rightarrow F_P$ is injective. Thus every
specialization morphism in $U$ is bijective, hence $F| U$ is locally
constant. By Noetherian induction we may suppose that $F| X-U$ is
constructible, whence the result.

\item It plainly suffices to treat the case where the function $c$ is
constant. Let $V\rightarrow X$ be an étale neighborhood of a geometric
point $P$ above a maximal point of $X$, such that $F(V)$ maps
surjectively onto $F_P$. Let $S$ be a subset of $F(V)$ such that
$S\xrightarrow{\sim}F_P$. Then, for the constant sheaf $S_V$ and every
specialization $P\rightarrow Q$ of geometric points, one has a
commutative diagram
$$
\xymatrix@C=2cm@R=1.25cm{
(S_V)_P   \ar[r]^{\sim}& F_P\\
(S_V)_Q   \ar[r] \ar[u]^{\text{\rotatebox{90}{$\backsim$}}}  & F_Q. \ar[u]
}
$$
Since $\lvert F_Q\rvert=\lvert F_P\rvert$, the arrow $F_Q\rightarrow F_P$ is bijective, so
$F$ is locally constant in a neighborhood of $P$ by (i). By
Noetherian induction, we are done.
\end{enumerate}

The following result (ii) will be of technical use later; it will allow
us to reduce certain verifications to the case of a constant sheaf.

\begin{enonce*}{Proposition 2.14}\etiquette[2.14]{IX.prop2.14}
Let\pageoriginale $X$ be a scheme.
\begin{enumerate}
\renewcommand{\theenumi}{\roman{enumi}}
\renewcommand{\labelenumi}{(\theenumi)}
\item For every finite morphism of finite presentation
$f:X'\rightarrow X$, and every constructible sheaf of sets (\resp of
ind-finite groups, \resp of $A$-modules) on $X'$, $f_{\ast}(F)$ is
constructible.

\item Suppose that $X$ is quasi-compact and quasi-separated, and let
$F$ be a constructible sheaf of sets (\resp\ldots) on $X$. Then one can
find a finite family $(p_i:X'_i\rightarrow X_i)$ of finite morphisms,
for every $i$ a constant constructible sheaf of sets (\resp\ldots)
$C_i$ on $X'_i$, and a monomorphism
$$
F\hookrightarrow \prod_i p_{i\ast}(C_i){\quoi}.
$$
\end{enumerate}
\end{enonce*}

First prove (i). By definition, we may suppose that $X$ is affine, and
then, using \hyperref[IX.cor2.7.4]{2.7.4}, the standard limit procedure
reduces us to the case where $X$ is Noetherian. A further passage to
the limit, again using \hyperref[IX.cor2.7.4]{2.7.4} and also
\hyperref[IX.prop2.4]{2.4}. (v), reduces us to the case where $X$ is the
spectrum of a field. The conclusion then reduces to the corresponding
assertion concerning the fiber of $f_{\ast}(F)$ at a geometric point
over $X$, which follows immediately from (VIII \hyperref[VIII.prop5.5]{5.5}, \hyperref[VIII.rem5.8]{5.8}).

Let us prove (ii). By \hyperref[IX.prop2.4]{2.4} (i), there exists a
partition of $X$ into subschemes $Y_i$, with inclusion morphisms
$u_i:Y_i\rightarrow X$ of finite presentation, such that the
restriction $F_i$ of $F$ to every $Y_i$ is a locally constant sheaf. It
is then immediate that the canonical homomorphisms $F\rightarrow
u_{i\ast}(F_i)$ define a homomorphism $F\rightarrow \prod_i
u_{i\ast}(F_i)$ which is a monomorphism.

In the preceding construction, one could even arrange that every $F_i$
be not only locally constant, but also \emph{isotrivial,} \ie that there
exist a surjective \emph{finite} étale morphism
$q_i:Y'_i\rightarrow Y_i$ such that $q^{\ast}_i(F_i)$ is a constant
sheaf on $Y'_i$, say $C_{iY'_i}$. This follows easily, for example, from
the definition of \hyperref[IX.lem2.8.1]{2.8.1} and from (VIII \hyperref[VIII.thm1.1]{1.1}). On
the other hand, it\pageoriginale follows from the ``Main Theorem'' in
the form (\EGA IV 8.12.6) that there exists a finite morphism
$p_i:Z_i\rightarrow X$ and an open immersion
$v_i:Y'_i\rightarrow Z_i$ of $X$-schemes.

Since $F_i$ embeds into
$q_{i\ast}(q^{\ast}_i(F_i))=q_{i\ast}(C_{iY'_i})$, $F$ embeds into the
product of the
$u_{i\ast}(q_{i\ast}(C_{iY'_i}))=p_{i\ast}(v_{i\ast}(C_{iY'_i}))$. If
every $Z_i$ is normal, then Lemma (\hyperref[IX.lem2.14.1]{2.14.1})
below implies that $v_{i\ast}(C_{iY'_i})\simeq C_{iZ_i}$, hence $F$
embeds into the product of the $p_{i\ast}(C_{iZ_i})$, and we are done.
In the general case, introduce the normalization $\bar{Z}_i$ of $Z_i$,
the inverse image $\bar{Y}'_i$ of $Y'_i$ in $\bar{Z}_i$, and the
immersion $\bar{v}_i:\bar{Y}'_i\rightarrow \bar{Z}_i$. One sees
immediately that $\bar{Y}'_i$ contains the maximal points of
$\bar{Z}_i$, and is therefore dense in $\bar{Z}_i$. If
$r_i:\bar{Z}_i\rightarrow Z_i$ and $s_i:\bar{Y}'_i\rightarrow Y'_i$
are the projections, since the latter is surjective, it follows that
$C_{iY'_i}$ embeds into $s_{i\ast}(C_{i\bar{Y}'_i})$, hence
$v_{i\ast}(C_{iY'_i})$ embeds into
$v_{i\ast}(s_{i\ast}(C_{i\bar{Y}'_i}))=r_{i\ast}(\bar{v}_{i\ast}(C_{i\bar{Y}'_i}))$,
which is isomorphic to $r_{i\ast}(C_{i\bar{Z}_i})$ by
(\hyperref[IX.lem2.14.1]{2.14.1}) below. Thus $F$ embeds into the
product of the
$p_{i\ast}(r_{i\ast}(C_{\bar{Z}_i}))=t_{i\ast}(C_{i\bar{Z}_i})$, where
$t_i=p_ir_i$ is the canonical projection $\bar{Z}_i\rightarrow X$.
When this morphism is finite for every $i$, we are done. In the general
case, one can only say that $\bar{Z}_i$ is \emph{integral} over $X$,
hence of the form $\sheaf{\Spec}(𝒜_i)$, where $𝒜_i$ is a
quasi-coherent sheaf of $𝒪_X$-algebras which is integral over $𝒪_X$.
Using (\EGA IV 1.7.7), one sees that $𝒜_i$ is the inductive limit of
its subalgebras $𝒜_{i,j}$ finite over $X$, hence $\bar{Z}_i$ is the
projective limit (VII 5) of the preschemes
$Z_{ij}=\Spec(A_{ij})$, finite over $X$. Applying (VII \hyperref[VII.cor5.11]{5.11}), one finds
that $t_{i\ast}(C_{i\bar{Z}_i})$ is the inductive limit of the
$t_{ij^{\ast}}(C_{iZ_{ij}})$, where $t_{ij}:Z_{ij}\rightarrow X$ is the
canonical projection. By (\hyperref[IX.cor2.7.3]{2.7.3}), the
homomorphism $F\rightarrow t_{i\ast}(C_{i\bar{Z}_i})$ therefore factors
through a homomorphism $F\rightarrow t_{ij\ast}(C_{iZ_{ij}})$ for a
suitable $j=j(i)$. We now set $X'_i=Z_{i,j(i)}$, and obtain an embedding
as stated in \hyperref[IX.prop2.14]{2.14}. It remains to prove the
following

\begin{enonce*}{Lemma 2.14.1}\etiquette[2.14.1]{IX.lem2.14.1}
Let $X$ be a geometrically unibranch scheme (for example, a normal
scheme), and let $i:U\rightarrow X$ be a quasi-compact open immersion.
Then the closure $Z$\pageoriginale of $U$, endowed with the induced
reduced structure, is also geometrically unibranch; more precisely, for
every $z\in Z$, one has
$𝒪_{Z,z}\xleftarrow{\sim}𝒪_{X_{\red},z}$. Suppose on the other hand
that $i$ is dominant, and let $C$ be a set, with $C_U$ the constant
sheaf on $U$ defined by $C$. Then the natural homomorphism
$C_X\rightarrow i_{\ast}(C_U)$ is an isomorphism. The same conclusions
hold if the open immersion $i$ is replaced by an inclusion
$i:x\rightarrow X$, where $x$ is a maximal point of $X$.
\end{enonce*}

Using (\EGA IV 2.3.11), we are reduced for the first assertion to the
case where $X$ is local, with closed point $z$; but then
$X_{\text{ réd}}$ is integral and $U$ is nonempty, hence
$Z=X_{\text{ réd}}$, and the assertion is evident (thus, instead of
assuming $X$ geometrically unibranch, it would have sufficed here to
suppose that the local rings of $X_{\text{réd}}$ are integral
domains). For the second assertion, looking fiber by fiber, we are
reduced by (VIII \hyperref[VIII.rem5.3]{5.3}) \loccit to the case where $X$ is strictly local,
and then to proving that $H^0(X,C_X)\rightarrow H^0(U,C_U)$ is
bijective. Here again, $X$ is irreducible, and $U$ is a nonempty open
subset of it and hence irreducible, and the two sides of the preceding
map are identified with $C$, whence the desired conclusion. (Here the
notion ``geometrically unibranch'' has been used in the form of
(\EGA IV 18.8.14), as meaning that the strict localizations of the
rings of $X_{\text{ réd}}$ are integral domains.) The case of the
inclusion of a maximal point (stated here for ease of future reference)
is treated in exactly the same way.

\begin{enonce*}[remark]{Remarks 2.14.2}\etiquette[2.14.2]{IX.rems2.14.2}
When $X$ is Noetherian, the proof given (or an immediate deduction)
shows that in \hyperref[IX.prop2.14]{2.14} (ii), the $X'_i$ may be
assumed integral; if in addition $X$ is assumed universally Japanese,
\ie the rings of its affine open subsets are universally Japanese
(\EGA ${\rm 0_{IV}}$ 23.1.1, IV 7.7.2), the $X'_i$ may moreover be
assumed normal.
\end{enonce*}

% END INCLUDED SGA4_English_current_buildable_source_20260906/source/09/drafts/SGA4_IX_LINES_651_1136_English.texfrag


% BEGIN INCLUDED SGA4_English_current_buildable_source_20260906/source/09/drafts/SGA4_IX_LINES_1137_1339_English.texfrag
\section{Kummer and Artin--Schreier theories}\etiquette[3]{IX.sec3}\pageoriginale

\setcounter{subsection}{-1}
\subsection{}\etiquette[3.0]{IX.subsec3.0}
We denote by $(\GG_m)_X$ the sheaf represented by the multiplicative
group $\Spec\Oo_X[t,t^{-1}]$ over $X$. Thus, for $U\to X$ étale, the
sections of $(\GG_m)_X$ over $U$ are the invertible elements
$\Gamma(U,\Oo_U^*)$ of $\Gamma(U,\Oo_U)$. Let $n\in\NN$. The kernel
of the $n$-th power map on $(\GG_m)_X$ is the \lq\lq sheaf of $n$-th
roots of unity'', denoted $(𝛍_n)_X$. Thus there is an exact sequence
\begin{equation*}
0\to(𝛍_n)_X\to(\GG_m)_X\xrightarrow{n}(\GG_m)_X{\quoi}.
\tag{3.1}
\end{equation*}\etiquette[3.1]{IX.eq3.1}
If $n$ is invertible on $X$ (that is, if $n$ is prime to every
residue characteristic of $X$), the $n$-th power map is surjective
\emph{as a morphism of sheaves}:

\begin{enonce*}{Kummer theory 3.2}\etiquette[3.2]{IX.Thr3.2}
If $n\in\NN$ is invertible on $X$, there is an exact sequence
$$
0\to(𝛍_n)_X\to(\GG_m)_X\xrightarrow{n}(\GG_m)_X\to0.
$$
\end{enonce*}

Indeed, let $u\in(\GG_m)_X(U)$, with $U\to X$ étale. We must find an
étale covering $\{U_i\to U\}$ such that $u$ becomes an $n$-th power
on every $U_i$. Since $n$ is invertible on $U$, the equation
$$
T^n-u=0
$$
is \lq\lq separable'' over $U$, that is,
$$
U'=\Spec\Oo_U[T]/(T^n-u)
$$
is étale over $U$. Since $U'\to U$ is surjective, it is a covering,
and $u$ becomes an $n$-th power on $U'$.

We point out\pageoriginale that when $n$ is invertible on $X$,
$(𝛍_n)_X$ is a locally constant sheaf, in fact locally isomorphic to
the sheaf $(\ZZ/n\ZZ)_X$, and is constant in important cases. Kummer
theory therefore gives information about the cohomology of $X$ with
values in certain constant sheaves. By definition,
$$
H^0(X,(\GG_m)_X)=\Gamma(X,\Oo_X^*){\quoi},
$$
and in dimension one we have:

\begin{enonce*}{Theorem 3.3}\etiquette[3.3]{IX.thm3.3}
(\lq\lq Hilbert's Theorem 90'') There is an isomorphism
$$
H^1(X,(\GG_m)_X)\xrightarrow{\sim}\Pic X{\quoi},
$$
where $\Pic X$ is the group of isomorphism classes of invertible
sheaves on $X$.
\end{enonce*}

\emph{Proof.} This says that the canonical morphism
$$
H^1(X_{\zar},(\GG_m)_X)\longrightarrow
H^1(X_{\text{\rm ét}},(\GG_m)_X)
$$
is bijective. Equivalently,
$$
R^1\epsilon_*(\GG_m)_X=0,
$$
where $\epsilon_*$ is direct image for the evident morphism of sites
$\epsilon:X_{\et}\to X_{\zar}$ (VIII 4). We are therefore reduced to
the case where $X$ is affine. Since $H^1$ may be computed by
\v{C}ech cohomology (V \hyperref[V.subsec2.5]{2.5}), and since for quasi-compact $X$ it is
enough to use quasi-compact coverings, this follows immediately from
descent theory for sheaves (see FGA 190 or \SGA VIII 1).

\setcounter{subsection}{3}
\setcounter{subsubsection}{0}
\subsubsection{}\etiquette[3.3.1]{IX.subsubsec3.3.1}
Suppose that $X$ is Noetherian and has no embedded component, and let
$x_j$ be the maximal points of $X$. Let $R_j$ be the local Artinian
ring of $X$ at $x_j$, and let $i_j:\Spec R_j\to X$ be the inclusion.
There is\pageoriginale a canonical injection
$$
(\GG_m)_X\longrightarrow
\prod_j i_{j^{\ast}}(\GG_m)_{\Spec R_j}
$$
of sheaves on $X_{\text{\rm ét}}$ (respectively $X_{\zar}$). Its
cokernel $D$ is called the \emph{sheaf of Cartier divisors} on
$X_{\text{\rm ét}}$ (respectively $X_{\zar}$)\footnote[1]{In
\EGA IV 21, the word \lq\lq divisors'' is used instead of
\lq\lq Cartier divisors''.}.

\begin{enonce*}{Corollary 3.4}\etiquette[3.4]{IX.cor3.4}
Let $X$ be Noetherian and without embedded components, let
$\epsilon:X_{\text{\rm ét}}\to X_{\zar}$ be the canonical morphism of
sites, and let $D$ be the sheaf of Cartier divisors on
$X_{\text{\rm ét}}$. Then $\epsilon_*D$ is the sheaf $D_{\zar}$ of
Cartier divisors on $X_{\zar}$, and in particular
$$
H^0(X_{\text{\rm ét}},D)\simeq H^0(X_{\zar},D_{\zar}){\quoi}.
$$
\end{enonce*}

\emph{Proof.} Applying $\epsilon_*$ to the exact sequence
$$
0\to(\GG_m)_X\longrightarrow
\prod_j i_{j^{\ast}}(\GG_m)_{\Spec R_j}
\longrightarrow D\to0
$$
gives an exact sequence, since
$R^1\epsilon_*(\GG_m)_X=0$ (\hyperref[IX.thm3.3]{3.3}); the first
assertion follows.

When $X$ has characteristic $p\neq0$, the following result sometimes
replaces \hyperref[IX.Thr3.2]{3.2} in the study of $p$-torsion
coefficients:

\begin{enonce*}{Artin--Schreier theory 3.5}\etiquette[3.5]{IX.Thr3.5}
If $X$ has characteristic $p>0$, there is an exact sequence
$$
0\to(\ZZ/p\ZZ)_X\to\Oo_X\xrightarrow{\mathfrak p}\Oo_X\to0{\quoi},
$$
where $\mathfrak p$ is the morphism of additive groups defined by
$\mathfrak p(f)=f^p-f$.
\end{enonce*}

The kernel of $\mathfrak p$ is the sheaf of sections of $\Oo_X$ which
locally belong to the prime field, and is therefore a constant sheaf.
The morphism $\mathfrak p$ is\pageoriginale surjective for the étale
topology because the equation $T^p-T-a$ ($a\in\Gamma(U,\Oo_U)$) is
always \lq\lq separable'' in characteristic $p$, that is, defines an
étale covering of $U$. In applying this exact sequence, we shall of
course use VII \hyperref[VII.subsec4.2]{4.2}.

\begin{enonce*}{Proposition 3.6}\etiquette[3.6]{IX.prop3.6}
\begin{enumerate}
\renewcommand{\theenumi}{\roman{enumi}}
\renewcommand{\labelenumi}{(\theenumi)}
\item Let $X$ be a scheme and let $i:x\to X$ be the inclusion of a
point. Then
$$
H^1(X,i_*(\ZZ)_x)=0,
$$
where $(\ZZ)_x$ is the \lq\lq constant sheaf'' of value $\ZZ$
(see \hyperref[IX.subsec2.0]{2.0}).

\item If $X$ is irreducible and if, for every geometric point
$\bar x$ of $X$, the strict localization $X_{\bar x}$ of $X$ at
$\bar x$ is irreducible (one should say that \lq\lq $X$ is
geometrically unibranch'', see \EGA IV 18.8.14), then
$$
H^1(X,(\ZZ)_X)=0{\quoi}.
$$
\end{enumerate}
\end{enonce*}

\emph{Proof.}
\begin{enumerate}
\renewcommand{\theenumi}{\roman{enumi}}
\renewcommand{\labelenumi}{(\theenumi)}
\item The Leray spectral sequence
$$
E^{pq}_2=H^p(X,R^qi_*(\ZZ)_x)\Longrightarrow
H^{p+q}(x,(\ZZ)_x)
$$
gives an injection
$H^1(X,i_*(\ZZ)_x)\to H^1(x,(\ZZ)_x)$. Since
$x=\Spec k(x)$ is the spectrum of a field, it is well known that
$H^1(x,(\ZZ)_x)=0$. (This follows from the fact that
$H^q(x,F)$ is torsion for $q>0$ in any case (see VII \hyperref[VII.rem2.1]{2.2} and CG),
and from the exact sequences
$0\to\ZZ\xrightarrow{n}\ZZ\to\ZZ/n\to0$.) This proves (i).

\item Let $i:x\to X$ be the inclusion of the generic point of $X$.
Under the hypotheses of (ii), the canonical morphism
$(\ZZ)_X\to i_*(\ZZ)_x$ is immediately seen to be bijective (and the
converse is also true); thus (ii) is reduced to (i).
\end{enumerate}

\begin{enonce*}[remark]{Remarks 3.7}\etiquette[3.7]{IX.rems3.7}
In \hyperref[IX.prop3.6]{3.6} (i) and (ii), $\ZZ$ may be replaced by
any torsion-free abelian group. On the other hand, (ii) becomes false
if the hypothesis that $X$ is geometrically unibranch is omitted, as
one sees, for example, by taking $X$ to be an algebraic curve over a
field $K$ having an ordinary double point (compute $H^1(X,\ZZ)$ in
this case!).
\end{enonce*}


% END INCLUDED SGA4_English_current_buildable_source_20260906/source/09/drafts/SGA4_IX_LINES_1137_1339_English.texfrag


% BEGIN INCLUDED SGA4_English_current_buildable_source_20260906/source/09/drafts/SGA4_IX_LINES_1340_1459_English.texfrag
\section{The case of an algebraic curve}\etiquette[4]{IX.sec4}

\setcounter{subsection}{-1}
\subsection{}\etiquette[4.0]{IX.subsec4.0}
Let $X$ be a Noetherian scheme of dimension $1$. Let $R(X)$ be the
ring of rational functions on $X$, and let
$i:\Spec R(X)\to X$ be the inclusion. The ring $R(X)$ is Artinian,
the product of the local rings of $X$ at its maximal points, and the
results of VIII 2 may be applied.

Let $F$ be an abelian sheaf on $\Spec R(X)$, and consider the sheaves
$R^qi_*F$, $q>0$, on $X$. It follows at once from VIII \hyperref[VIII.thm5.2]{5.2} that the
fiber of $R^qi_*F$ at a geometric point lying over a maximal point of
$X$ vanishes. Since $X$ has dimension $1$, such a sheaf is of a very
special kind:

\begin{enonce*}{Lemma 4.1}\etiquette[4.1]{IX.lem4.1}
Let $X$ be Noetherian and let $F$ be a sheaf on $X$. The following
conditions are equivalent (a sheaf satisfying them will be called a
\lq\lq skyscraper sheaf''):
\begin{enumerate}
\renewcommand{\theenumi}{\roman{enumi}}
\renewcommand{\labelenumi}{(\theenumi)}
\item The fiber of $F$ at every geometric point $\bar x$ of $X$
lying over a nonclosed point $x\in X$ vanishes.

\item Every section $s\in F(X')$, where $X'\to X$ is étale of finite
type, vanishes away from finitely many closed points of $X$.

\item There is an isomorphism
$F\simeq\bigoplus_x i_{x*}i_x^*F$, where $x$ runs through the closed
points of $X$ and $i_x:x\to X$ is the inclusion.
\end{enumerate}
\end{enonce*}

\emph{Proof.}\pageoriginale Suppose (i) holds and let
$z\in F(X')$, where $X'\to X$ is étale of finite type. The support of
$z$ (VIII \hyperref[VIII.subsec6.6]{6.6}) is then a closed subset of $X'$ consisting of
\emph{closed} points, and is therefore finite because $X'$ is
Noetherian. Thus (i) implies (ii). The implication (ii) to (i)
follows from the computation of fibers (VIII \hyperref[VIII.prop3.9]{3.9}), and (iii) implies
(i), since for every family $(G_x)$ of sheaves on the closed points
$x$ of $X$, the sheaf $F=\bigoplus i_{x*}(G_x)$ on $X$ satisfies
(i): fiber functors commute with direct sums, and the support of
$i_{x*}(G_x)$ is plainly contained in $\{x\}$.

It remains to prove that (ii) implies (iii). Consider the canonical
morphism
$$
F\longrightarrow\prod_x i_{x*}i_x^*F{\quoi},
$$
where $x$ runs through the closed points of $X$. By (ii), it factors
through
$$
F\longrightarrow\bigoplus_x i_{x*}i_x^*F{\quoi},
$$
which gives the desired morphism. We leave to the reader the
verification that it is bijective.

\setcounter{subsection}{1}
\subsubsection{}\etiquette[4.1.1]{IX.subsubsec4.1.1}
Apply the lemma to the sheaf $R^qi_*F$ introduced above. One has
$$
R^qi_*F\xrightarrow{\sim}
\bigoplus_x i_{x*}i_x^*(R^qi_*F),
$$
where $x$ runs through the closed points of $X$. By VIII 5, the fiber
of $R^qi_*F$ at a geometric point $\bar x$ of $X$ is
$$
(R^qi_*F)_{\bar x}\simeq H^q(\Spec R(X_{\bar x}),F),
$$
where
$R(X_{\bar x})\simeq R(X)\otimes_{\Oo_X}\Oo_{X_{\bar x}}$,
$X_{\bar x}$ is the strict localization of $X$ at $\bar x$ (VIII 4),
and $F$ also denotes the sheaf induced on
$\Spec R(X_{\bar x})$. Suppose now that the residue field of every
closed point $x$ of $X$ is separably closed. Then (see also VI 1)
\begin{align*}
H^p(X,R^qi_*F)
&\simeq H^p\left(X,\bigoplus_x i_{x*}i_x^*(R^qi_*F)\right)\\
&\simeq \bigoplus_x H^p\left(X,i_{x*}i_x^*(R^qi_*F)\right)\\
&\simeq \bigoplus_x H^p\left(x,i_x^*(R^qi_*F)\right)\\
&\simeq
\begin{cases}
\displaystyle\bigoplus_x H^q(\Spec R(X_x),F)
  &\mbox{if }p=0,\\
0 &\mbox{if }p>0{\quoi}.
\end{cases}
\end{align*}
Thus:

\begin{enonce*}{Corollary 4.2}\etiquette[4.2]{IX.cor4.2}
Let $X$ be a Noetherian scheme of dimension $1$ such that, for every
closed point $x$ of $X$, the residue field $k(x)$ is separably
closed. Let $F$ be an abelian sheaf on $\Spec R(X)$, where $R(X)$ is
the ring of rational functions on $X$. Consider the Leray spectral
sequence
$$
E^{pq}_2=H^p(X,R^qi_*F)\Longrightarrow
H^{p+q}(\Spec R(X),F){\quoi}.
$$
One has
\begin{align*}
E^{pq}_2&=0\quad\mbox{if }p>0\mbox{ and }q>0{\quoi},\\
E^{0q}_2&=\bigoplus_x H^q(\Spec R(X_x),F)
  \quad\mbox{if }q>0{\quoi},
\end{align*}
where $x$ runs through the closed points of $X$, and $X_x$ is the
strict localization of $X$ at $x$.

The\pageoriginale spectral sequence therefore gives an exact sequence
\begin{align*}
0\longrightarrow H^1(X,i_*F)&\longrightarrow H^1(\Spec R(X),F)
\longrightarrow\bigoplus_xH^1(\Spec R(X_x),F)\\
&\longrightarrow H^2(X,i_*F)\longrightarrow H^2(\Spec R(X),F)
\longrightarrow\bigoplus_xH^2(\Spec R(X_x),F)\\
&\longrightarrow\cdots .
\end{align*}
\end{enonce*}


% END INCLUDED SGA4_English_current_buildable_source_20260906/source/09/drafts/SGA4_IX_LINES_1340_1459_English.texfrag


% BEGIN INCLUDED SGA4_English_current_buildable_source_20260906/source/09/drafts/SGA4_IX_LINES_1460_1606_English.texfrag
\setcounter{subsection}{2}
\subsection{}\etiquette[4.3]{IX.subsec4.3}
Suppose moreover that the $\ell$-cohomological dimension of
$\Spec R(X)$, for a given prime number $\ell$, is at most $1$ (that
is, $H^q(\Spec R(X),F)=0$ for $q>1$ and every $\ell$-torsion sheaf
$F$ on $\Spec R(X)$). Let $K^1,\ldots,K^m$ be the residue fields of
the Artinian ring $R(X)$, that is, the residue fields of the maximal
points of $X$. By the dictionary of VIII 2, the
$\ell$-cohomological dimension of $\Spec R(X)$ is
$$
\cd_\ell(\Spec R(X))
=\mathop{\sup}_i\{\cd_\ell(G(\overline K^i/K^i))\}{\quoi},
$$
where $\cd_\ell(G)$ denotes the $\ell$-cohomological dimension of the
profinite group $G$ (see CG).

Every residue field $K_x$ of $R(\bar X(x))$, for a strict
localization $\bar X(x)$ of $X$, is plainly a separable (possibly
infinite) extension of one of the $K^i$, and hence
$$
\cd_\ell(G(\bar K_x/K_x))
\leq\cd_\ell(G(\bar K^i/K^i))
$$
by CG II 4.1, Proposition 10. Thus
$\cd_\ell(\Spec R(\bar X(x)))\leq1$ as well.

Applying the exact sequence of \hyperref[IX.cor4.2]{4.2} to an
$\ell$-torsion sheaf $F$, one obtains
\begin{equation*}
\begin{cases}
0\to H^1(X,i_*F)\to H^1(\Spec R(X),F)\to\\
\qquad\displaystyle\bigoplus_xH^1(\Spec R(\bar X(x)),F)
\to H^2(X,i_*F)\to0{\quoi},\\
H^q(X,i_*F)=0\quad\mbox{for }q>2{\quoi}.
\end{cases}
\tag{4.3.1}
\end{equation*}\etiquette[4.3.1]{IX.eq4.3.1}

Note\pageoriginale that the hypothesis
$\cd_\ell(\Spec R(X))\leq1$ holds if $X$ is an algebraic curve over a
separably closed field (Tsen's theorem [6, Theorem
\hyperref[IX.subsec4.3]{4.3}]). Analogous exact sequences in this case
were studied by $\Ogg[2]$ and \v{S}afarevi\v{c}.

\subsection{}\etiquette[4.4]{IX.subsec4.4}
Suppose that $\ell$ is invertible on $X$ and take
$F=(\GG_m)_{R(X)}$. Applying \hyperref[IX.thm3.3]{3.3} and
[CG, Chapter I, Proposition 12], still under the hypothesis
$\cd_\ell(\Spec R(X))\leq1$, one finds
\begin{equation*}
\begin{cases}
H^1(\Spec R(X),\GG_m)=0,\\
H^q(\Spec R(X),\GG_m)\;\displaystyle{\mathop{=}_{\ell}}\;0
\quad(q\geq2){\quoi},
\end{cases}
\tag{4.4.1}
\end{equation*}\etiquette[4.4.1]{IX.eq4.4.1}
where the $\ell$ under the equality sign means that the group on the
left, which is torsion in any case, has no $\ell$-torsion.

The same result holds with $R(X)$ replaced by $R(\bar X(x))$. Since
the sheaf $\GG_m$ on $\Spec R(X)$ plainly induces the sheaf $\GG_m$
on $\Spec R(\bar X(x))$, \hyperref[IX.cor4.2]{4.2} gives
\begin{equation*}
\begin{cases}
H^1(X,i_*(\GG_m)_{R(X)})=0,\\
H^q(X,i_*(\GG_m)_{R(X)})
\;\displaystyle{\mathop{=}_{\ell}}\;0
\quad\mbox{if }q>1{\quoi}.
\end{cases}
\tag{4.4.2}
\end{equation*}\etiquette[4.4.2]{IX.eq4.4.2}
For every sheaf $F$ on $X$, \hyperref[IX.lem4.1]{4.1} shows that the
kernel $P$ and cokernel $Q$ of the canonical homomorphism
$F\to i_*i^*F$ are skyscraper sheaves. Since the residue fields at
the closed points of $X$ are assumed separably closed, one therefore
has $H^j(X,P)=H^j(X,Q)=0$ for $j>0$. The long exact cohomology
sequence immediately implies that
$$
H^j(X,F)\longrightarrow H^j(X,i_*i^*F)
$$
is\pageoriginale an isomorphism for $j\geq2$. Taking
\eqhyperref[IX.eq4.4.2]{4.4.2} into account, one obtains
\begin{equation*}
H^j(X,\GG_m)\displaystyle{\mathop{\;=\;}_{\ell}}0
\quad\mbox{for }j\geq2{\quoi},
\tag{4.5}
\end{equation*}\etiquette[4.5]{IX.eq4.5}
valid when $X$ is a Noetherian scheme of dimension $1$ satisfying
$\cd_\ell(R(X))\leq1$ and whose residue fields at closed points are
separably closed. Kummer theory (\hyperref[IX.Thr3.2]{3.2}) and
\hyperref[IX.thm3.3]{3.3} now give:

\begin{enonce*}{Theorem 4.6}\etiquette[4.6]{IX.thm4.6}
Let $X$ be Noetherian of dimension $1$, and let $n$ be an integer
invertible on $X$. Suppose that $\cd_\ell(R(X))\leq1$ for every prime
number $\ell$ dividing $n$, and that the residue field $k(x)$ is
separably closed for every closed point $x$ of $X$. Then
$$
H^q(X,(𝛍_n)_X)=0\quad\mbox{if }q>2{\quoi},
$$
and the cohomology for $q\leq2$ is given by the exact sequence
\begin{align*}
0\to H^0(X,(𝛍_n)_X)&\to\Gamma(X,\Oo_X^*)
\xrightarrow{n}\Gamma(X,\Oo_X^*)\to\\
H^1(X,(𝛍_n)_X)&\to\Pic X\xrightarrow{n}\Pic X
\to H^2(X,(𝛍_n)_X)\to0{\quoi}.
\end{align*}
\end{enonce*}

\begin{enonce*}{Corollary 4.7}\etiquette[4.7]{IX.cor4.7}
Let $X$ be a complete connected curve over a separably closed field
$k$ whose characteristic is prime to $n$. Then
\begin{align*}
H^0(X,(𝛍_n)_X)&\simeq𝛍_n(k){\quoi},\\
H^1(X,(𝛍_n)_X)&\simeq{}_n(\Pic X)
  =\{\text{elements of $\Pic X$ whose order divides $n$}\},\\
H^2(X,(𝛍_n)_X)&\simeq(\Pic X)/n
  \simeq(\ZZ/n\ZZ)^c{\quoi},
\end{align*}
where\pageoriginale $c$ is the number of irreducible components of
$X$.
\end{enonce*}

For the last assertion, recall that the Picard scheme [TDTE V] of
such an $X/\Spec k$ has a decomposition
\begin{equation*}
0\longrightarrow A\longrightarrow\sheaf{\Pic}_{X/k}
\longrightarrow\ZZ^c\longrightarrow0
\tag{4.8}
\end{equation*}\etiquette[4.8]{IX.eq4.8}
where $A$ is a connected algebraic group over $\Spec k$. The
\lq\lq multiplication by $n$'' map on such an $A$ is étale, hence
surjective on $k$-valued points. It follows that
$$
{}_n(\Pic X)\simeq{}_n(A(k))
$$
and
$$
(\Pic X)/n\simeq(\ZZ/n\ZZ)^c{\quoi}.
$$

If $X$ is smooth over $\Spec k$, so that its Jacobian $A$ is an
abelian variety, the group \({}_n(\Pic X)\) is isomorphic to
$(\ZZ/n)^{2g}$, where $g$ is the genus of $X$. Of course,
$𝛍_n(k)\simeq\ZZ/n\ZZ$ as well. Thus the ranks of the cohomology
groups as $\ZZ/n$-modules are $1$, $2g$, and $1$, as expected.


% END INCLUDED SGA4_English_current_buildable_source_20260906/source/09/drafts/SGA4_IX_LINES_1460_1606_English.texfrag


% BEGIN INCLUDED SGA4_English_current_buildable_source_20260906/source/09/drafts/SGA4_IX_LINES_1607_1883_English.texfrag
\section{The trace method}\etiquette[5]{IX.sec5}

\subsection{}\etiquette[5.1]{IX.subsec5.1}
Let $X$ be a scheme, let $f:X'\rightarrow X$ be an étale covering, and
let $F$ be a sheaf on $X$. Then the adjunction formulas give a
\emph{restriction} morphism
$$
F\xrightarrow{\res}f_{\ast}f^{\ast}F{\quoi}.
$$
Since $f$ is finite, one has $H^q(X,f_{\ast}f^{\ast}F)\simeq H^q(X',
f^{\ast}F)$, whence a morphism, also called restriction:
\begin{equation*}
H^q(X,F)\xrightarrow{\res}H^q(X',f^{\ast} F){\quoi},\tag{5.1.1}
\end{equation*}\etiquette[5.1.1]{IX.eq5.1.1}
which\pageoriginale is moreover the evident morphism. But when $f$ is
étale, there is also another morphism, called the \emph{trace
morphism},
\begin{equation*}
f_{\ast}f^{\ast} F\xrightarrow{\tr} F{\quoi},\tag{5.1.2}
\end{equation*}\etiquette[5.1.2]{IX.eq5.1.2}
which induces morphisms on cohomology
\begin{equation*}
H^q(X', f^{\ast} F)\xrightarrow{\tr}H^q
(X,F){\quoi}.\tag{5.1.3}
\end{equation*}\etiquette[5.1.3]{IX.eq5.1.3}
The trace morphism is characterized by the following two properties:
\begin{enumerate}
\renewcommand{\theenumi}{\roman{enumi}}
\renewcommand{\labelenumi}{(\theenumi)}
\item It is local for the étale topology.

\item If $X'$ is the disjoint sum of $d$ copies of $X$, so that
$f_{\ast}f^{\ast} F\simeq F^d$, the trace is the summation map
$F^d\rightarrow F$.
\end{enumerate}

Uniqueness follows trivially from (i), (ii), since every étale
covering is locally constant. Once uniqueness is known, existence
follows by the usual descent argument.

From (ii) one deduces the following formula:
\begin{equation*}
\tr\circ \res:F\longrightarrow F\mbox{ is multiplication by the
local degree of $f$.}\tag{5.1.4}
\end{equation*}\etiquette[5.1.4]{IX.eq5.1.4}
If the degree is a constant $d$, it follows that the morphism
$\tr\circ \res: H^q(X,F)\rightarrow H^q(X,F)$ is multiplication by
$d$.

\begin{enonce*}{Corollary 5.2}\etiquette[5.2]{IX.cor5.2}
Let $f:X'\rightarrow X$ be an étale covering of constant degree $d$.
If $F$ is an $r$-torsion sheaf on $X$ with $(r,d)=1$, then
$H^q(X',f^{\ast} F)=0$ implies that $H^q(X,F)=0$.
\end{enonce*}

Indeed, multiplication by $d$ induces an isomorphism of $F$, hence
an\pageoriginale isomorphism on cohomology. If the latter isomorphism
factors through zero, then $H^q(X,F)=0$.

\setcounter{subsection}{2}
\subsubsection{}\etiquette[5.2.1]{IX.subsubsec5.2.1}
Now let $i:U\rightarrow X$ be an immersion, let $f:U'\rightarrow U$ be
an étale covering of degree $d$, and suppose that $f$ is induced by a
finite morphism $g:X'\rightarrow X$ by base change, \ie that one has
the cartesian diagram
$$
\xymatrix@C=2cm@R=1.25cm{
U' \ar[r]^{i'}\ar[d]_{f}& X'\ar[d]^{g}\\
U  \ar[r]^{i}& X
}
$$
Let $F$ be a sheaf on $U$. One has
$$
\xymatrix@C=2.25cm@M=.25cm{
F \ar[r]^{\txt{res}} \ar@/_1.5pc/[rr]_{d}& f_{\ast}f^{\ast}F
\ar[r]^{\txt{tr}}& F,
}
$$
hence, since $i_{\ast}f_{\ast}=g_{\ast}i'_{\ast}$,
\hypertarget{sga4.exp.IX.equation.5.3}{}$$
\xymatrix@C=2.25cm{
i_{\ast}F \ar@/_2pc/[rr]_{d} \ar[r]^{\txt{res}}& g_{\ast}i' {}_{\ast}f^{\ast}F \ar[r]^{\txt{tr}}& i_{\ast}F,
}\leqno{(5.3)}
$$
and since $i_{!}f_{\ast}\simeq g_{\ast}i'_{!}$ (which follows, for
example, from (VIII \hyperref[VIII.prop5.5]{5.5})),
$$
\xymatrix@C=2cm{
i_{!}F \ar@/_2pc/[rr]_{d} \ar[r]^{\txt{res}}& g_{\ast}i' {}_{!}f^{\ast}F \ar[r]^{\txt{tr}}& i_{!}F,
}\leqno{(5.4)}
$$\etiquette[5.4]{IX.eq5.4}
whence
$$
\xymatrix@C=2cm{
H^{q}(X,i_{\ast} F) \ar@/_2pc/[rr]_{d'} \ar[r]^{\txt{res}}& H^{q}(X', i' _{\ast}f^{\ast}F) \ar[r]^{\txt{tr}}& H^q(X,i_{\ast}F)
}\leqno{(5.3')}
$$\MakeLinkTarget*{sga4.exp.IX.equation.5.3-prime}\label{5.3'}
and\pageoriginale
$$
\xymatrix@C=2cm{
H^{q}(X,i_{!} F) \ar[r]^{\txt{res}} \ar@/_2.5pc/[rr]_{d}& H^{q}(X', i' _{!}f^{\ast}F)
\ar[r]^{\txt{tr}}& H^q(X,i_{!}F) ;
}\leqno{(5.4')}
$$\MakeLinkTarget*{sga4.exp.IX.equation.5.4-prime}\label{5.4'}
From \eqhyperref[IX.eq5.4]{5.4} one deduces the following corollaries:

\begin{enonce*}{Proposition 5.5}\etiquette[5.5]{IX.prop5.5}
Let $X$ be a quasi-compact and quasi-separated scheme, let $\ell \in
\PP$, and let $H$ be a half-exact functor with values
in {\rm (Ab)}, defined on the category of $\ell$-sheaves on $X$, and
compatible with filtered inductive limits. The following assertions
are equivalent:
\begin{enumerate}
\renewcommand{\theenumi}{\roman{enumi}}
\renewcommand{\labelenumi}{(\theenumi)}
\item $H=0$.

\item $H(f_{\ast}i_{!}(\ZZ/\ell))=0$ for every finite
$f:Y\rightarrow X$, and every open subscheme $i:U\rightarrow Y$ of
$Y$, such that $U$ is of finite presentation over $X$. If $X$ is
noetherian, one may even restrict to integral $Y$.
\end{enumerate}
\end{enonce*}

Proof. By \hyperref[IX.cor2.7.2]{2.7.2} and \hyperref[IX.prop2.5]{2.5}
it is enough to verify $H(F)=0$ for $F$ of the form $F=i$, $G$, where
$i:U\rightarrow X$ is the inclusion of a subscheme of $X$ of finite
presentation over $X$, and where $G$ is a locally constant
``isotrivial'' $\ell$-sheaf on $U$. Let $U''\rightarrow U$ be a
principal covering with group $g$, such that $g$ becomes constant on
$U''$. Let $A$ be the $\ell$-group of values of $G$ on $U''$. The
group $g$ acts on $A$, and this action determines the structure of
$G$. Let $\sheaf{h}$ be an $\ell$-Sylow subgroup of $g$, and let
$U'=U''/\sheaf{h}$ be the covering of $U$ corresponding to
$\sheaf{h}$. The degree $d$ of $U'$ over $U$ is prime to $\ell$.
Choose an extension of $f:U'\rightarrow U$ to a finite morphism
$g:X'\rightarrow X$, which always exists (\EGA IV 8.12.6), and let
$i':U'\rightarrow X'$ be the inclusion. Applying
\eqhyperref[IX.eq5.4]{5.4} and the fact that $d$ is prime to $\ell$,
we are reduced to proving
that $H(g_{\ast}i'_{!}f^{\ast} G)=0$.

Now\pageoriginale it is well known that the only abelian
$\ell$-torsion group that is simple under an action of an $\ell$-group
$\sheaf{h}$ is $\ZZ/\ell$ with trivial action, and it follows that
there exists a filtration of $A$ as an $\sheaf{h}$-module whose
successive quotients are isomorphic to $\ZZ/\ell$ with trivial action.
This filtration gives a corresponding filtration of $f^{\ast} G$, and
therefore, in order to prove that $H(g_{\ast}i'_{!} f^{\ast} G)=0$, it
is enough to prove that $H(g_{\ast}i'_{!}\ZZ/r)=0$, whence the result.

\begin{enonce*}{Proposition 5.6}\etiquette[5.6]{IX.prop5.6}
Let $X$ be a noetherian scheme, let $\ell \in \PP$, and
let $\varphi$ be a function with values in $\NN$, defined on the set
of integral schemes finite over $X$. Suppose that
$\varphi(Y_1)\leq \varphi(Y_2)$ whenever there exists an $X$-morphism
$Y_1\rightarrow Y_2$, and that the inequality is strict if $Y_1$ is a
closed subscheme of $Y_2$, distinct from $Y_2$. For every $Y$ finite
over $X$, put $(Y)=\Sup (Y_i)$, where the $Y_i$ are the irreducible
components of $Y$, endowed with the induced reduced structure. Then
the following assertions are equivalent:
\begin{enumerate}
\renewcommand{\theenumi}{\roman{enumi}}
\renewcommand{\labelenumi}{(\theenumi)}
\item For every $Y$ finite over $X$ and every $\ell$-sheaf $F$ on
$Y$, one has $H^q(Y,F)=0$ if $q>\varphi(\Supp F)$.

\item For every integral $Y$ finite over $X$, one has
$H^q(Y,\ZZ/\ell)=0$ if $q>\varphi(Y)$.
\end{enumerate}
\end{enonce*}

Proof. Suppose that (ii) is true. We may apply the preceding
proposition to each $H^q$, taking into account that cohomology commutes
with inductive limits and with direct images under finite morphisms;
we are thus reduced, for the verification of (i), to the case where
$Y$ is integral and finite over $X$, and $F$ is of the form
$i_{!}(\ZZ/\ell)_U$, where $i:U\rightarrow X$ is a nonempty open
subscheme.

Now there is an exact sequence
$$
0\longrightarrow i_{!}(\ZZ/\ell)_U\longrightarrow
(\ZZ/\ell)_Y\longrightarrow
(\ZZ/\ell)_Z\longrightarrow 0
$$
where $Z=Y-U$, hence $Z\neq Y$. Arguing by induction on the value of
$\varphi$, we may\pageoriginale suppose that
$H^q(\ZZ, \ZZ/\ell)=0$ for $q>\varphi (Y)-1$, whence
$H^q(Y,i_{!}(\ZZ/\ell))=0$ for $q>\varphi(Y)$ by the exact
cohomology sequence.

\begin{enonce*}{Corollary 5.7}\etiquette[5.7]{IX.cor5.7}
Let $X$ be an algebraic curve over a separably closed field $k$ of
characteristic different from $\ell$, with $\ell\in \PP$.
Then the cohomology of $X$ with coefficients in a sheaf $F$ of
$\ell$-torsion vanishes in degrees $>2$. If $X$ is affine, the cohomology
vanishes in degrees $>1$.
\end{enonce*}

For integral $Y$ finite over $X$, put $\varphi(Y)=2\dim Y$ (\resp
$\varphi(Y)=\dim Y$ if $X$, hence $Y$, is affine). By \hyperref[IX.prop5.6]{7.6} it is enough
to prove that $H^q(Y,\ZZ/\ell)=0$ if $q>\varphi(Y)$, which is trivial
if $\dim Y=0$. It therefore remains to prove that
$H^q(Y,\ZZ/\ell)=0$ if $q>2$ (\resp if $q>1$ in the affine case) for
an integral curve $Y$. Since $(\ZZ/\ell)_Y\simeq (𝛍_{\ell})_Y$, this
follows from \hyperref[IX.thm4.6]{4.6} for $q>2$.

Suppose that $X$ is affine. Since radicial extensions do not affect
cohomology (VIII \hyperref[VIII.thm1.1]{1.1}), we are reduced to the case where $k$ is
algebraically closed. Let $f:X'\rightarrow X$ be the normalization of
$X$, which is a finite morphism and an isomorphism outside a finite set
of points, say $S$. There is an exact sequence of sheaves on $X$
$$
0\longrightarrow (\ZZ/n)_X\longrightarrow
f_{\ast}(\ZZ/n)_{X'}\longrightarrow \epsilon
\longrightarrow 0{\quoi},
$$
where $\epsilon$ is concentrated on $S$, so that
$H^q(X,\epsilon)=0$ for $q>0$. By (VIII \hyperref[VIII.cor5.6]{5.6}), we are reduced to proving
that
$$
H^q(X,f_{\ast}(\ZZ/n)_{X'})\simeq
H^q(X',\ZZ/n)=0
$$
if $q>1$, that is, to the case where $X$ is normal. We may also plainly
suppose that $X$ is connected. Let $i:X\rightarrow \bar{X}$ be an open
immersion with $\bar{X}$ complete and nonsingular. Since $X$ is
affine, at least one point is missing from $\bar{X}$, and it follows
immediately that the morphism
$$
A(k)\longrightarrow \Pic X
$$
is\pageoriginale surjective, where $A$ is the Jacobian of $\bar{X}$
(\cf \eqhyperref[IX.eq4.8]{4.8}). Since $A(k)$ is divisible by
$\ell\neq \car k$, the same is true of $\Pic X$, whence
$H^2(X,𝛍_{\ell})=0$ by \hyperref[IX.thm4.6]{4.6}, \cqfd

\begin{enonce*}[remark]{Remark 5.8}\etiquette[5.8]{IX.rem5.8}
Let us record, for later reference, that the argument of
\hyperref[IX.prop5.5]{5.5} in fact establishes the following result:
Let $X$ be a quasi-compact and quasi-separated scheme, let $\ell$ be a
prime number, and let $C$ be a strictly full subcategory of the
category of constructible abelian $\ell$-torsion sheaves, stable under
direct factors and extensions. Suppose that for every finite morphism
$f:Y\rightarrow X$ and every open subscheme $i:U\rightarrow Y$ of $Y$
that is of finite presentation over $X$, one has
$f_{\ast}(i_{!}(𝐙/\ell𝐙)_U)\in C$. Then $C$ contains all constructible
%% 𝒵 here is ℤ [check the macros]
$\ell$-torsion sheaves.
\end{enonce*}

\begin{thebibliography}{99}
\bibitem{sga4.exp.IX.bib.1} Gabriel, P. Des Catégories Abéliennes,
Thèse, Paris (1962).

\bibitem{sga4.exp.IX.bib.2} Ogg, A. Cohomology of abelian varieties over
functions fields, Annals of Math. Vol. 76, No. 2 (1962).

\bibitem{sga4.exp.IX.bib.3} Deuring, M. Die Typen der Multiplikatorenringe elliptischer
Funktionenkörper, Hamb. Abh. Bol. 13 (1940) p. 197.

\bibitem{sga4.exp.IX.bib.4} Grothendieck, A. Fondements de la Géométrie
Algébrique, Extraits du Séminaire Bourbaki 1957-1962), Secrétariat
mathématique, 11, rue Pierre Curie, Paris 5ème. (\emph{Cited as} FGA).

\bibitem{sga4.exp.IX.bib.5} Serre, J.P. Cohomologie Galoisienne, Lecture Notes
nᵒ 5, (Springer) (\emph{Cited as} CG).

\bibitem{sga4.exp.IX.bib.6} Douady, A. Cohomologie des groupes compacts
totalement discontinus, Séminaire Bourbaki, exp. 189.
\end{thebibliography}

\ifx\danslelivre\undefined
\end{document}
\fi

% END INCLUDED SGA4_English_current_buildable_source_20260906/source/09/drafts/SGA4_IX_LINES_1607_1883_English.texfrag


% END INCLUDED SGA4_English_current_buildable_source_20260906/source/09/09_English.tex


\makeatletter
\def\input@path{{../10/}}
\makeatother
% BEGIN INCLUDED SGA4_English_current_buildable_source_20260906/source/10/10_English.tex
\ifx\danslelivre\undefined
% Release engine: XeLaTeX (the LuaTeX compatibility layer breaks Xy-pic output).
\documentclass[10pt]{../fichiers-configuration-latex/smfart}
% BEGIN INCLUDED SGA4_English_current_buildable_source_20260906/source/fichiers-configuration-latex/paquets.tex
\usepackage{polyglossia}
%\usepackage[francais]{babel}
\usepackage{fontspec}
\usepackage[math-style=ISO,bold-style=upright,sans-style=upright]{unicode-math}
%\usepackage[mathletters]{ucs}
%\usepackage[utf8]{inputenc}
%\usepackage[T1]{fontenc}
\usepackage[unicode,hypertexnames=false,colorlinks=true,linkcolor=blue,citecolor=Sepia,urlcolor=magenta]{hyperref}
\setmainlanguage{french}
\setmainfont[Ligatures=TeX]{Linux Libertine O}
% variante possible : TeXGyrePagella ?
\setmathfont{XITS Math}

% maths
\usepackage{amsmath}
%\usepackage{amssymb}
%\usepackage{mathrsfs}
%\usepackage{stmaryrd}
%\usepackage{euscript}

% couleurs
\usepackage[usenames,dvipsnames]{xcolor}

% graphiques
\usepackage{graphicx}
\usepackage[all]{xy}
%\usepackage{tikz}
%\usetikzlibrary{matrix,arrows,calc}
%\tikzset{column sep=3em,row sep=2em,text height=1.5ex, text depth=0.25ex,>=angle 90}
%\tikzset{description/.style={fill=white,inner sep=1pt}}

% hyperliens, 'etiquettes
\usepackage{xr-hyper}
\usepackage{../fichiers-configuration-latex/smfthm}
%\ifx\cacheretiquettes\undefined
%\usepackage[inline]{showlabels} \else \fi % pour rendre visibles les 'etiquettes
%\usepackage{../fichiers-configuration-latex/makeidx} % pour faire un index
%\usepackage[pagebackref,unicode]{hyperref}  % pour voir o`u sont cit'ees les r'ef'erences
% hyperref is loaded once above with the cumulative-reader destination policy.
%\hypersetup{breaklinks=true} %permet le retour `a la ligne dans les liens trop longs. Incompatible avec pdflatex
%\renewcommand*{\backref}[1]{{\small [cit\'e page(s)~#1.]}}

\usepackage{xspace}
\usepackage{pifont} % pour expose XVII

% END INCLUDED SGA4_English_current_buildable_source_20260906/source/fichiers-configuration-latex/paquets.tex

% BEGIN INCLUDED SGA4_English_current_buildable_source_20260906/source/fichiers-configuration-latex/numerotation.tex
\renewcommand{\thefootnote}{\roman{footnote}} % pour mettre les notes en bas de page sous la forme (ii) et non (2), qui ressemble trop (?) a un exposant.

\newcounter{pageoriginal}
\newcommand{\pageoriginale}{\refstepcounter{pageoriginal}\marginpar{\footnotesize\textbf{\thepageoriginal}}\xspace}
\let\pageoriginaled\pageoriginale

\usepackage{ifthen}
\newboolean{original}
\def\sisi{\ifthenelse{\boolean{original}}}
\setboolean{original}{false}

\newcounter{cnote}
\def\nde#1{\ifthenelse{\boolean{original}}{}{\stepcounter{cnote}\def\thefootnote{\arabic{cnote}}\protect\footnote[\value{cnote}]{N.D.E.: #1}}}
\def\ndemark{\ifthenelse{\boolean{original}}{}{\stepcounter{cnote}\def\thefootnote{\arabic{cnote}}\protect\footnotemark[\value{cnote}]}}
\def\ndetext#1{\ifthenelse{\boolean{original}}{}{\def\thefootnote{\arabic{cnote}}\protect\footnotetext[\value{cnote}]{N.D.E.: #1}}}

% END INCLUDED SGA4_English_current_buildable_source_20260906/source/fichiers-configuration-latex/numerotation.tex

% BEGIN INCLUDED SGA4_English_current_buildable_source_20260906/source/fichiers-configuration-latex/macros.tex
\setmainlanguage{french}
\setotherlanguage{english}


%% Fixed, engine-independent release identifier.  The inherited dynamic Git
%% probe used Lua io.popen, exposed build-state-dependent output, and printed
%% a visible question mark outside a Git checkout.
\def\showgitstatus{English release r2, 2026-07-23}


% macros pour LaTeX
\renewcommand{\mathsf}{}

% Police différente pour U+2113 SCRIPT SMALL L
\setmathfont[range="2113]{latinmodern-math.otf}

% Give every inherited semantic label its own stable PDF destination instead
% of inheriting whichever automatic equation/page anchor happened to be live.
% The legacy label remains as a compatibility alias used by existing edges.
\makeatletter
\def\etiquette[#1]#2{%
  \Hy@raisedlink{\hyper@anchorstart{sga4.semantic.#2}\hyper@anchorend}%
  \def\@currentHref{sga4.semantic.#2}%
  \def\@currentlabel{#1}%
  \label{#2}%
}
\makeatother
\def\eqhyperref[#1]#2{(\hyperref[#1]{#2})}
\def\pagehyperref[#1]#2{(page \hyperref[#1]{#2})}
\def\hyperrefext[#1][#2]#3{\hyperref[#2]{#1 #3}}
\def\eqhyperrefext[#1][#2]#3{\hyperref[#2]{#1 #3}}     %cf. exposes 17
\def\neweqhyperrefext[#1][#2]#3{\hyperref[#2]{#1 #3}}  %cf. exposes 17

\def\shat#1{\protect\widehat{#1}}
\def\stilde#1{\protect\widetilde{#1}}
%%%%%%%%%%%%%%%%%%


\def\etcn{etc.\xspace}
\def\etco{etc...\xspace}
\def\etc{\sisi{\etco}{\etcn}}
\def\ldotsn{\sisi{...}{\etcn}}

\def\cf{{cf.}\kern.3em}
\def\Cf{{Cf.}\kern.3em}
\def\ie{{i.e.}\kern.3em}
\def\iev{{i.e.},\xspace}
\def\loccito{loc.\kern2pt cit.\xspace}
\def\loccitn{\textit{loc.\kern2pt cit.}\xspace}
\def\loccit{\sisi{\loccito}{\loccitn}}
\def\EGA{\textup{EGA}\kern.2em}
\def\SGA{\textup{SGA}\kern.3em}

\def\quoi{}

\def\quav{\quad,\quad}
\def\quaet{\quad \text{et}\quad}
\def\fibreproduct#1#2#3{#1 ×_{#3} #2}
\let\fprod\fibreproduct
\def\ccheck#1#2{\mbox{${\check{\check{\rm #1}}}^{\lower2pt\hbox{$\scriptstyle #2$}}$}}
\def\sheaf#1{\protect\mathcal{#1}}
\def\vieuxsheaf#1{{#1}} % cf. problème avec sheaf{a\Phi_i} par exemple...
\DeclareMathOperator{\Hom}{Hom}
\DeclareMathOperator{\can}{can}
\DeclareMathOperator{\Ind}{Ind}
\DeclareMathOperator{\ind}{ind}
%\DeclareMathOperator{\Hhomtop}{\mathcal{H}\!\mathit{omtop}}
\DeclareMathOperator{\Hhomtop}{\mathbf{Homtop}}
\DeclareMathOperator{\Point}{Point}
\DeclareMathOperator{\Pro}{Pro}
\DeclareMathOperator{\Top}{Top}
\DeclareMathOperator{\Mod}{Mod}
\DeclareMathOperator{\Cart}{Cart}
%\DeclareMathOperator{\Mmorsite}{\mathcal{M}\!\mathit{orsite}}
\DeclareMathOperator{\Mmorsite}{\mathbf{Morsite}}
\DeclareMathOperator{\loc}{loc}
\DeclareMathOperator{\Loc}{Loc}
\DeclareMathOperator{\Pic}{Pic}
\DeclareMathOperator{\Et}{Et}
\DeclareMathOperator{\card}{card}
\DeclareMathOperator{\Ens}{Ens}
\DeclareMathOperator{\h}{\underline{H}}
\DeclareMathOperator{\Fonct}{Fonct}
\DeclareMathOperator{\Ker}{Ker}
\DeclareMathOperator{\but}{but}
\DeclareMathOperator{\Ob}{Ob}
\DeclareMathOperator{\ob}{ob}
\DeclareMathOperator{\Lib}{Lib}
\DeclareMathOperator{\esj}{esj}
\DeclareMathOperator{\id}{id}
\DeclareMathOperator{\Id}{Id}
\DeclareMathOperator{\ad}{ad}
\DeclareMathOperator{\Iim}{Im}
\DeclareMathOperator{\pr}{pr}
\DeclareMathOperator{\Cat}{Cat}
\DeclareMathOperator{\car}{car}
\DeclareMathOperator{\SysInd}{SysInd}
\DeclareMathOperator{\Sup}{Sup}
\DeclareMathOperator{\Coker}{Coker}
\DeclareMathOperator{\Filt}{Filt}
\DeclareMathOperator{\acc}{acc}
\DeclareMathOperator{\sur}{sur}
\DeclareMathOperator{\cat}{cat}
\DeclareMathOperator{\Ord}{Ord}
\DeclareMathOperator{\Gr}{Gr}
\DeclareMathOperator{\Cov}{Cov}
\DeclareMathOperator{\Caf}{Caf}
\DeclareMathOperator{\coker}{coker}
\DeclareMathOperator{\coim}{coim}
\DeclareMathOperator{\im}{im}
\DeclareMathOperator{\Ouv}{Ouv}
\DeclareMathOperator{\diag}{diag}
\DeclareMathOperator{\Esp}{Esp}
\DeclareMathOperator{\TOP}{TOP}
\DeclareMathOperator{\Unif}{Unif}
\DeclareMathOperator{\Iind}{\mathbf{Ind}}
\DeclareMathOperator{\Hhomcart}{\mathbf{Homcart}}
\DeclareMathOperator{\Hhomtopess}{\mathbf{Homtopess}}
\DeclareMathOperator{\Ppoints}{\mathbf{Points}}
\DeclareMathOperator{\Ppoint}{\mathbf{Point}}
\DeclareMathOperator{\Ttors}{\mathbf{Tors}}
\DeclareMathOperator{\Ssex}{\mathbf{Sex}}
\DeclareMathOperator{\Ffib}{\mathbf{Fib}}
\DeclareMathOperator{\Vvois}{\mathbf{Vois}}
\DeclareMathOperator{\Ppointess}{\mathbf{Pointess}}
\DeclareMathOperator{\scrHom}{\mathbf{Hom}}
\DeclareMathOperator{\sheafPic}{\mathbf{Pic}}
\DeclareMathOperator{\sheafAut}{\mathbf{Aut}}
\DeclareMathOperator{\Hhomtoplocan}{\mathbf{Homtoplocan}}
\DeclareMathOperator{\Hhomtopan}{\mathbf{Homtopan}}
\DeclareMathOperator{\Pptgeom}{\mathbf{Ptgéom}}
\DeclareMathOperator{\Morsite}{Morsite}
\DeclareMathOperator{\Fib}{Fib}
\DeclareMathOperator{\Fais}{Fais}
\DeclareMathOperator{\kar}{kar}
\DeclareMathOperator{\Groupes}{Groupes}
\DeclareMathOperator{\Restr}{Restr}
\DeclareMathOperator{\Isom}{Isom}
\DeclareMathOperator{\Prol}{Prol}
\DeclareMathOperator{\Res}{Res}
\DeclareMathOperator{\Tors}{Tors}
\DeclareMathOperator{\esp}{esp}
\DeclareMathOperator{\es}{es}
\DeclareMathOperator{\od}{od}
\DeclareMathOperator{\gr}{gr}
\DeclareMathOperator{\sob}{sob}
\DeclareMathOperator{\dfn}{dfn}
\DeclareMathOperator{\Sch}{Sch}
\DeclareMathOperator{\ouv}{ouv}
\DeclareMathOperator{\cons}{cons}
\DeclareMathOperator{\Nerf}{Nerf}
\DeclareMathOperator{\Points}{Points}
\DeclareMathOperator{\Espsob}{Espsob}
\DeclareMathOperator{\Homtop}{Homtop}
\DeclareMathOperator{\Pt}{Pt}
\DeclareMathOperator{\Quot}{Quot}
\DeclareMathOperator{\glob}{glob}
\DeclareMathOperator{\Homfib}{Homfib}
\DeclareMathOperator{\Qouv}{Qouv}
\DeclareMathOperator{\Quotouv}{Quotouv}
\DeclareMathOperator{\ens}{ens}
\DeclareMathOperator{\ab}{ab}
\DeclareMathOperator{\an}{an}
\DeclareMathOperator{\fs}{fs}
\DeclareMathOperator{\red}{red}
\DeclareMathOperator{\Ab}{Ab}
\DeclareMathOperator{\cosq}{cosq}
\DeclareMathOperator{\Ann}{Ann}
\DeclareMathOperator{\Zar}{Zar}


%%%%%%%%%%%%
%%%%% SGA4-2
\DeclareMathOperator{\Sex}{Sex}
\DeclareMathOperator{\vois}{vois}
\DeclareMathOperator{\Vois}{Vois}
\DeclareMathOperator{\Ext}{Ext}
\DeclareMathOperator{\ext}{ext}
\DeclareMathOperator{\Fer}{Fer}
\DeclareMathOperator{\Add}{Add}
\DeclareMathOperator{\Cocart}{Cocart}
\DeclareMathOperator{\cosk}{cosk}
\DeclareMathOperator{\plat}{plat}

\DeclareMathOperator{\faisc}{faisc}
\DeclareMathOperator{\sk}{sk}
\DeclareMathOperator{\ch}{ch}

\def\HH{\mathbf{H}}
\def\qc{\mathrm{qc}}
\def\et{\mathrm{et}}
\def\coh{\mathrm{coh}}

\def\Hvr#1{%
\begin{array}[b]{l}
{\kern2pt\scriptstyle r}\\[-8pt]
{\kern1pt\scriptstyle \vee}\\[-5pt]
H^{#1}
\end{array}}

\def\Hvv#1{%
\begin{array}[b]{l}
\kern2pt\rotatebox{90}{${\scriptscriptstyle\ll}$}\\[-6pt]
H^{#1}
\end{array}}

\def\Hvi#1{%
\begin{array}[b]{l}
{\kern0.2pt\scriptscriptstyle \infty}\\[-8pt]
{\kern1pt\scriptstyle \vee}\\[-5pt]
H^{#1}
\end{array}}

\def\Hv#1{%
\begin{array}[b]{l}
{\kern1pt\scriptstyle \vee}\\[-5pt]
H^{#1}
\end{array}}

\DeclareMathOperator{\cart}{cart}
\DeclareMathOperator{\Ffaisc}{\mathbf{Faisc}}


\def\Ev#1{%
\begin{array}[b]{l}
{\kern1pt\scriptstyle \vee}\\[-5pt]
E'_{#1}
\end{array}}

\DeclareMathOperator{\Cont}{Cont}
\DeclareMathOperator{\Ccont}{\mathscr{C}\!\mathit{ont}}
\DeclareMathOperator{\Site}{Site}
\newcommand\MM{\mathbf{M}}
\def\limtop#1{{\displaystyle{\varprojlim_{#1}\!\mbox{\rm top }}}}
\def\Limtop{\mathop{\varprojlim\!\mbox{\rm top}}}
\DeclareMathOperator{\spec}{spec}
\def\ss{\scriptstyle{\mathscr{S}}}
\DeclareMathOperator{\Gal}{Gal}
\def\smallm{\scriptstyle{\mathscr{M}}}
\DeclareMathOperator{\End}{End}
\DeclareMathOperator{\ass}{ass}
\DeclareMathOperator{\zar}{zar}
\DeclareMathOperator{\Ogg}{Ogg}
\DeclareMathOperator{\res}{res}
\DeclareMathOperator{\tr}{tr}
\DeclareMathOperator{\Supp}{Supp}
\DeclareMathOperator{\codim}{codim}
\DeclareMathOperator{\rk}{rk}
\DeclareMathOperator{\cur}{cur}
\DeclareMathOperator{\cdqc}{cdqc}
\newcommand\DD{\mathbf{D}}
\DeclareMathOperator{\cl}{cl}
\def\JJ{\mathbf{J}}
\def\FF{\mathbf{F}}
\def\BB{\mathbf{B}}
\DeclareMathOperator{\Alget}{Alget}
\DeclareMathOperator{\Ass}{Ass}
\DeclareMathOperator{\supp}{supp}
\DeclareMathOperator{\rad}{rad}
\DeclareMathOperator{\prof}{prof}
\DeclareMathOperator{\cint}{int}
\DeclareMathOperator{\sic}{sic}
\DeclareMathOperator{\Tor}{Tor}
\DeclareMathOperator{\tors}{tors}
\DeclareMathOperator{\Sym}{Sym}
\DeclareMathOperator{\parf}{parf}
\DeclareMathOperator{\tor}{tor}
\DeclareMathOperator{\defn}{defn}
\DeclareMathOperator{\Div}{Div}
\DeclareMathOperator{\Pol}{Pol}
\DeclareMathOperator{\Fex}{Fex}
\DeclareMathOperator{\Lie}{Lie}
\DeclareMathOperator{\tot}{tot}
\DeclareMathOperator{\Aut}{Aut}
\DeclareMathOperator{\HOM}{HOM}
\DeclareMathOperator{\pch}{pch}
\DeclareMathOperator{\Ch}{Ch}
\DeclareMathOperator{\EXT}{EXT}
\DeclareMathOperator{\PIC}{PIC}
\DeclareMathOperator{\Rf}{Rf}
\DeclareMathOperator{\inj}{inj}
\newcommand\VV{\mathbf{V}}

\def\Xv{%
\begin{array}[b]{l}
{\kern1pt\scriptstyle \vee}\\[-5pt]
X
\end{array}}

%%%%%%%
\newcommand\TT{\mathbf{T}}
\newcommand\XX{\mathbf{X}}
\newcommand\EE{\mathbf{E}}
%%%%%%

\def\qedsymbol{\textup{C.Q.F.D.}}
\let\CQFD\qed
\let\cqfd\qed

%%%%%\bigdot
\font\bigsymb=cmsy10 at 4pt
\def\bigdot{{\kern1.2pt\raise 1.5pt\hbox{\bigsymb\char15}}}
%%%%%

\def\Lright{\mathop{\vcenter{\hbox{\hbox{$\longrightarrow$}\lower8pt \hbox{\kern-16pt\hbox{$\longrightarrow$}}}}}}

\def\Aa{{\mathcal A}}
\def\Bb{{\mathcal B}}
\def\Cc{{\mathcal C}}
\def\Dd{{\mathcal D}}
\def\Ee{{\mathcal E}}
\def\Ff{{\mathcal F}}
\def\Gg{{\mathcal G}}
\def\Hh{{\mathcal H}}
\def\Ii{{\mathcal I}}
\def\Jj{{\mathcal J}}
\def\Kk{{\mathcal K}}
\def\Ll{{\mathcal L}}
\def\Mm{{\mathcal M}}
\def\Nn{{\mathcal N}}
\def\Oo{{\mathcal O}}
\def\Pp{{\mathcal P}}
\def\Qq{{\mathcal Q}}
\def\Rr{{\mathcal R}}
\def\Ss{{\mathcal S}}
\def\Tt{{\mathcal T}}
\def\Uu{{\mathcal U}}
\def\Vv{{\mathcal V}}
\def\Ww{{\mathcal W}}
\def\Xx{{\mathcal X}}
\def\Yy{{\mathcal Y}}
\def\Zz{{\mathcal Z}}

\newcommand{\CC}{\mathbf{C}}
\newcommand{\GG}{\mathbf{G}}
\newcommand{\KK}{\mathbf{K}}
\newcommand{\LL}{\mathbf{L}}
\newcommand{\NN}{\mathbf{N}}
\newcommand{\PP}{\mathbf{P}}
\newcommand\QQ{\mathbf{Q}}
\newcommand\RR{\mathbf{R}}
\newcommand\ZZ{\mathbf{Z}}

\def\to{\mathchoice{\longrightarrow}{\rightarrow}{\rightarrow}{\rightarrow}}
\def\from{\mathchoice{\longleftarrow}{\leftarrow}{\leftarrow}{\leftarrow}}
\def\mto{\mathchoice{\longmapsto}{\mapsto}{\mapsto}{\mapsto}}
\def\mfrom{\mathrel{\from\joinrel\mapstochar\,}}
\def\hto{\mathrel{\lhook\joinrel\to}}
\def\hfrom{\mathrel{\from\joinrel\rhook}}

\def\implies{\mathchoice{\Longrightarrow}{\Rightarrow}{\Rightarrow}{\Rightarrow}}
\def\implied{\mathchoice{\Longleftarrow}{\Leftarrow}{\Leftarrow}{\Leftarrow}}
\def\iff{\mathchoice{\Longleftrightarrow}{\Leftrightarrow}{\Leftrightarrow}{\Leftrightarrow}}

\def\To#1{\mathchoice{\xrightarrow{\labelstyle~~#1~}}{\stackrel{#1}{\longrightarrow}}{}{}}
\def\Mto#1{\mathrel{\mapstochar\To{#1}}}
\def\From#1{\mathchoice{\xleftarrow{\labelstyle~~#1~}}{\stackrel{#1}{\longleftarrow}}{}{}}

\newcommand{\isomto}{\overset{\sim}{\longrightarrow}}
\newcommand{\isomfrom}{\overset{\sim}{\longleftarrow}}

\newcommand{\IFF}{\unskip~$\iff$~\ignorespaces}
\newcommand{\THEN}{\unskip~$\implies$~\ignorespaces}

\let\leq\leqslant
\let\geq\geqslant
\let\le\leqslant
\let\ge\geqslant

\def\cf{{cf.}\kern.3em}
\def\Cf{{Cf.}\kern.3em}
\def\ie{{i.e.}\kern.3em}
\def\iev{{i.e.},\xspace}
%\def\loccito{loc.\kern2pt cit.\xspace}
%\def\loccitn{\textit{loc.\kern2pt cit.}\xspace}
%\def\loccit{\sisi{\loccito}{\loccitn}}
\def\EGA{\textup{EGA}\kern.2em}
\def\SGA{\textup{SGA}\kern.3em}

\let\oldS\S
\def\S{\oldS\kern .15em}
\def\ptbl{.\kern .2em}
\def\resp{resp.\kern.3em}
\def\Exp{Exp.\kern.3em}

\let\bar\overline
\let\hat\widehat
\let\tilde\widetilde
\let\dpl\displaystyle
\let\ldots\dots
\let\emptyset\varnothing
\let\skpt\skippointrait
\let\epsilon\varepsilon

\def\sfrac#1#2{{#1}/{#2}}

\DeclareMathOperator{\cd}{cd}
%\DeclareMathOperator{\Hhom}{\mathcal{H}\!\mathit{om}}
\DeclareMathOperator{\Hhom}{\mathbf{Hom}}
%\DeclareMathOperator{\Eext}{\mathcal{E}\!\mathit{xt}}
\DeclareMathOperator{\Eext}{\mathbf{Ext}}
\DeclareMathOperator{\Spec}{Spec}
\DeclareMathOperator{\Tr}{Tr}

\let\emptyset=\varnothing

%%%%%%%%%%%%%%%
%%% expose XVII
%%%%%%%%%%%%%%%

\makeatletter
\let\old@optexpose\optexpose
\newcommand{\empty@optexpose}[1]{}
\newcommand{\paren@optexpose}[1]{\old@optexpose{#1}(}
\newcommand{\refnoexpose}[1]{\let\optexpose\empty@optexpose\ref{#1}\let\optexpose\old@optexpose}%\refnoexpose{X.cor5.2} engendre "5.2" au lieu de "X 5.2"
\newcommand{\neweqref}[1]{\let\optexpose\paren@optexpose\textup{\ref{#1})}\let\optexpose\old@optexpose}%\neweqref{XII.eq4.1} engendre "XII (4.1)" au lieu de "(XII 4.1)"
\makeatother
\newcommand{\refvb}[2][VI${}_{\text{B}}$ ]{\sisi{#1\refnoexpose{#2}}{\ref{#2}}}

\newcommand{\leftexp}[2]{{\vphantom{#2}}^{#1}{#2}}
\newcommand{\AAA}{\mathbf{A}}
\newcommand{\Sf}{\mathfrak{S}}
\DeclareMathOperator*{\quotprojlim}{\text{``$\varprojlim$''}}
\DeclareMathOperator*{\quotinjlim}{\text{``$\varinjlim$''}}
\DeclareMathOperator{\Fl}{F\sisi{\ell}{l}}
\DeclareMathOperator{\Pl}{P\sisi{\ell}{l}}
%\DeclareMathOperator{\cl}{c\sisi{\ell}{l}}
\DeclareMathOperator{\tdf}{tdf}
\DeclareMathOperator{\rg}{rg}
\newcommand{\iInd}{\sisi{I}{i}nd}
\newcommand{\iind}{ind}
\newcommand{\ppro}{pro}
\newcommand{\tTor}{\sisi{T}{t}or}
\newcommand{\ttor}{tor}
\newcommand{\chap}{Chap.}
\newcommand{\Ffex}{\Ff\mathit{ex}}
\newcommand{\sheafk}{\underline{k}}
\newcommand{\sheafn}{\underline{n}}
\newcommand{\sheafGamma}{\underline{\Gamma}}

% Repeated references to an already typeset footnote must not use
% \footnotemark without a matching \footnotetext: Hyperref would then create
% a live annotation to a destination that cannot exist.  This helper keeps the
% current printed footnote mark and links it to an explicit stable note target.
\makeatletter
\newcommand{\sgarepeatedfootnotemark}[1]{%
  \begingroup
  \edef\@thefnmark{\thefootnote}%
  \hyperlink{#1}{\@makefnmark}%
  \endgroup}
\makeatother

\sisi{
\DeclareMathOperator{\simarrow}{\rule{.15mm}{1.8mm}\kern -3.4pt\rightsquigarrow}
\newcommand{\sheafell}{\underline{\ell}}
\newcommand{\sheafZZ}{\underline{\ZZ}}
}{
\newcommand{\nderef}[1]{la N.D.E.~(\ref{#1}) page~\pageref{#1}}
}

%%%%%%%%%%%%%%%%%
%%%%% expose XVII
%%%%%%%%%%%%%%%%%

% END INCLUDED SGA4_English_current_buildable_source_20260906/source/fichiers-configuration-latex/macros.tex


\externaldocument{../01/01_English.tex}
\externaldocument{../02/02_English.tex}
\externaldocument{../03/03_English.tex}
\externaldocument{../04/04_English.tex}
\externaldocument{../05/05_English.tex}
\externaldocument{../05bis/05bis_English.tex}
\externaldocument{../06/06_English.tex}
\externaldocument{../07/07_English.tex}
\externaldocument{../08/08_English.tex}
\externaldocument{../09/09_English.tex}

\setboolean{original}{false}

\begin{document}

\begin{center}
X. Cohomological dimension: first results\\
M. Artin
\end{center}

\selectlanguage{english}version: \showgitstatus
\tableofcontents

\else
\chapter{Cohomological dimension: first results}
%\addtocontents{toc}{M. Artin}
\begin{center}
M. Artin
\end{center}
\fi
\phantomsection\MakeLinkTarget*{sga4.exp.X}\label{sga4.exp.X}

% BEGIN INCLUDED SGA4_English_current_buildable_source_20260906/source/10/drafts/SGA4_X_LINES_43_382_English.texfrag
\setcounter{pageoriginal}{42}

\section{Introduction}\etiquette[1]{X.sec1}\pageoriginale

In this Exposé, unless expressly stated otherwise, ``sheaf'' means
``abelian sheaf''.

Let $X$ be a scheme and let $\ell\in \PP$. ``Recall'' that the
\emph{$\ell$-cohomological dimension} of $X$ is the largest integer
$\cd_{\ell}X=n$ (or $=\infty$ if no such number exists) such that
$H^n(X,F)\neq 0$ for at least one $\ell$-torsion sheaf $F$ on $X$.
If $X=\Spec A$, we write $\cd_{\ell} A=\cd_{\ell}\Spec A$. Taking
into account the dictionary (VIII 2), %% label
we thus recover the ``classical'' notion of [CG] when $A=k$ is a field.
Since we plainly have the formula
$$
\cd_{\ell}(\coprod_i X_i)=\displaystyle{\mathop{\sup}_i} \cd_{\ell}X_i
$$
the classical theory may, in the evident way, be extended to an
arbitrary artinian ring.\pageoriginale This theory is set out in
$(CG)=$ ``\sisi{Galois cohomology}{Galois cohomology}'' %% Galois?
by Serre. We shall add a few refinements in ${\rm n^{\circ}}$
\hyperref[X.sec2]{2}.

Consider, for example, the case of an irreducible algebraic variety
$X$, of dimension $n$, over a separably closed field $k$, and fix
$\ell\neq \car k$. A theorem to be found in ${\rm n^{\circ}}$
\hyperref[X.sec4]{4} gives the upper bound $2n$ for its cohomological
dimension; this is the ``true'' dimension of $X$ (suggested by the
topological analogy when $k=\CC$). This result is fairly elementary
once the theory for fields is known, and in fact the upper bound $2n$
can be improved except when $X$ is complete: there is another upper
bound, $2n-1$, which holds whenever $X$ is ``missing at least one
point''. The latter is the correct ``stable upper bound'', in the sense
that every $X$ of dimension $n$ contains an open subset of
cohomological dimension $2n-1$. A third upper bound is obtained for an
affine variety, where one obtains $n$ as the cohomological dimension.
This is a generalization of one of the Lefschetz theorems for
hyperplane sections. These last two assertions are deeper and can only
be proved later (XIV).

Note that cohomological dimension is not a local notion (in contrast
to what happens for paracompact spaces). For example, a henselian ring
with separably closed residue field has cohomological dimension zero,
but if the closed point is removed, one may obtain an arbitrarily large
number, namely $2n-1$ if $n=\dim A$ in the most important cases; this
also agrees with the topological analogy supplied by the cohomology of
$U-\{x\}$, where $U$ is a small open neighborhood of a point $x$ of a
complex analytic space.

\section{Auxiliary results over a field}\etiquette[2]{X.sec2}\pageoriginale

We have the following

\begin{enonce*}{Theorem 2.1}\etiquette[2.1]{X.thm2.1}
Let $k'/k$ be an extension of finite type of a field $k$, and let
$\ell\in \PP$. Then
$$
\cd_{\ell}k'\leqq \cd_{\ell}k+\deg.\tr.(k'/k){\quoi},
$$
If the inequality is strict and if $\ell\neq \car k$, then $\ell=2$,
$k$ is orderable, and $k'$ is not orderable (\ie $-1$ is a sum of
squares in $k'$).
\end{enonce*}

Recall that when $\ell=\car k$ one has $\cd_{\ell}k'\leqq 1$ (CG II
\hyperref[X.thm2.2]{2.2}), so the question of equality in
\hyperref[X.thm2.1]{2.1} scarcely arises. Moreover,
$\cd_2k=\infty$ if $k$ is orderable. Indeed, the algebraic closure
$\bar{k}$ of $k$ contains a subextension $\widetilde{k}$ (a maximal
ordered subfield of $\bar{k}$) with
$[\bar{k}:\widetilde{k}]=2$. Plainly,
$\cd_2\widetilde{k}=\infty$, and
$\cd_2\widetilde{k}\leqq \cd_2 k$ by (CG I Prop.~14).

Proof. Almost all of this is well known. The inequality\sisi{}{,} as
well as equality if $\cd_{\ell}k<\infty$, are results of Tate (CG II
4.1.2). It remains to prove that $\cd_{\ell}k'<\infty$ and
$\cd_{\ell}k=\infty$ \sisi{implies}{imply} that $k$ is orderable and
$\ell=2$, and it is enough to treat the following two cases:
\begin{itemize}
\item[a)] $[k':k]<\infty$,

\item[b)] $k'=k(x)$ with $x$ transcendental over $k$.
\end{itemize}
In case a) we may suppose $k'/k$ separable. Then the Galois group
$H=G(\bar{k}/k')$ is an open subgroup of $G=G(\bar{k}/k)$, and by a
result of Serre [5] one has $\cd_{\ell}H=\cd_{\ell}G$ unless $G$
contains an element of order $\ell$. Since an algebraically closed
field $\bar{k}$ cannot have a subfield $\widetilde{k}$ of finite index
unless $\ell=2$ and $\widetilde{k}$ is orderable [7], the result
follows.

Consider\pageoriginale case b). Let $R$ be the henselization of
$\Spec k[x]$ at the point $x=0$, which is a discrete valuation ring,
and let $K$ be the fraction field of $R$. Then
$\cd_{\ell}k(x)<\infty$ implies $\cd_{\ell}K<\infty$. It is therefore
enough to prove that $\cd_{\ell}K<\infty$ implies
$\cd_{\ell}k<\infty$, which is a consequence of the following theorem:

\begin{enonce*}{Theorem 2.2}\etiquette[2.2]{X.thm2.2}
\footnote{This theorem was proved recently by $Ax$
    [1].}. Let $R$ be a henselian discrete valuation ring, with
    fraction field $K$ and residue field $k$. Let $\ell\in
    \PP$. We have
$$
\cd_{\ell}K=\cd_{\ell}k+1
$$
in the following two situations:
\begin{enumerate}
\renewcommand{\theenumi}{\roman{enumi}}
\renewcommand{\labelenumi}{(\theenumi)}
\item $\ell\neq \car k$.

\item $\ell \neq \car K$, $k$ perfect, and $\cd_{\ell}k<\infty$.
\end{enumerate}
\end{enonce*}

\begin{enonce*}[remark]{Remark}
Naturally, if $\ell=\car K$ one has $\cd_{\ell}K\leq 1$ in every case
(CG II \hyperref[X.thm2.2]{2.2}). It therefore remains to study the
relation of $\cd_{\ell}K$ with, perhaps, the module of absolute
differentials $\Omega^1_k$ (\EGA IV 20.4.3), in the unequal
characteristic case when $\car k=\ell$, $k$ imperfect.
\SGAArchiveOnly{\nde{The frozen source calls for a citation to work of Kato but supplies
no bibliographic target.}}
\end{enonce*}

Proof of the theorem. As for Theorem (\hyperref[X.thm2.1]{2.1}), most
of it is due to Tate. Case (ii) is treated explicitly in (CG II
Prop.~12) if $R$ is complete, and we reduce to that case by means of the
following lemma:

\begin{enonce*}{Lemma 2.2.1}\etiquette[2.2.1]{X.lem2.2.1}
Let $R$ be a henselian discrete valuation ring with fraction field $K$,
and let $\widehat{K}$ be the fraction field of the completion
$\widehat{R}$ of $R$. Then
$$
G(\bar{\widehat{K}}/\widehat{K})\simeq G(\bar{K}/K){\quoi}.
$$
\end{enonce*}

Proof. This means that the functor $L\mapsto
\widehat{L}=L\otimes_K \widehat{K}$ from\pageoriginale the category of
étale $K$-algebras to the category of étale $\widehat{K}$-algebras is
an equivalence. Now every étale algebra $\widehat{K}'$ over
$\widehat{K}$ is induced by a suitable étale algebra $K'$ over $K$.
Indeed, if $\widehat{K}'=\widehat{K}[\alpha]$ where $\alpha$ satisfies
the separable irreducible equation $f(T)=0$ with coefficients in
$\widehat{R}$, it is enough to take as $K'$ the algebra
$K[\alpha_0]$, where $\alpha_0$ is a root of an equation $f_0(T)=0$
with coefficients in $R$, with $f_0\equiv f(\widehat{m}^N)$, $N$ large
and $\widehat{m}$ the maximal ideal of $\widehat{R}$. This is an
application of the usual ``Newton method'' (Bourbaki,
Alg.~Comm.~III 4.5). It remains to prove that the functor is fully
faithful (it is plainly faithful), and this amounts to proving that if
$K'$ is a \emph{field}, then so is $\widehat{K}'$. Let $R'$ be the
normalization of $R$ in $K'$, which is finite over $R$ since $K'/K$ is
separable. The ring $R'$ is a discrete valuation ring, and the
completion of $R'$ is $\widehat{R}'=R'\otimes_R \widehat{R}$. Hence
$\widehat{R}'$ is a discrete valuation ring, and since $K'$ is its
total ring of fractions, it is indeed a field.

Consider assertion (i) of
\sisi{\eqhyperref[X.eq3.1]{3.1}}{\hyperref[X.thm2.2]{2.2}}.
If $\cd_{\ell}k<\infty$, we are again reduced by
\hyperref[X.lem2.2.1]{2.2.1}
to (CG II Prop.~12) (the hypothesis in the statement of that
proposition that $k$ be perfect is \sisi{used}{used} in the proof only
in the case of residue characteristic $\ell$). It therefore remains to
prove that $\cd_{\ell}K<\infty$ implies $\cd_{\ell}k<\infty$.

Let $\bar{G}=G(\bar{k}/k)=G(\widetilde{K}/K)$, where $\widetilde{K}$ is
the maximal unramified extension of $K$ (\ie the fraction field of the
discrete valuation ring $\widetilde{R}$ with separably closed residue
field, the strict henselization of $R$), let
$G=G(\bar{K}/K)$, $H=G(\bar{K}/\widetilde{K})$, so that
$\bar{G}=G/H$. Let $\bar{G}'\subset \bar{G}$ be an $\ell$-Sylow
subgroup, and let $k'/k$, $K'/K$ be the extensions corresponding to
$\bar{G}'$. We have $\cd_{\ell}k<\infty$ if and only if
$\cd_{\ell}k'<\infty$ (CG I Prop.~14), so by replacing $R$ with its
normalization in $K'$, we reduce to the case in which $\bar{G}$ is an
$\ell$-group. It is then enough to prove that if
$\cd_{\ell}G<\infty$, then $H^N(\bar{G},\ZZ/\ell)=0$ for sufficiently
large $N$ (CG I Prop.~20).

We\pageoriginale have $\cd_{\ell}H=1$ by the part of i) already proved,
applied to $\widetilde{R}$. Moreover, we have
\begin{equation*}
H^1(H, 𝛍_{\ell})=H^1(\Spec \widetilde{K},
𝛍_{\ell})=\widetilde{K}^{\ast}/\widetilde{K}^{\ast\ell}=\ZZ/\ell{\quoi}.\tag{2.2.2}
\end{equation*}\etiquette[2.2.2]{X.eq2.2.2}
Indeed, $H^1(\Spec \widetilde{K},
𝛍_{\ell})=\widetilde{K}^{\ast}/\widetilde{K}^{\ast\ell}$
by Kummer theory \hyperrefext[IX][IX.Thr3.2]{3.2}. We have an exact
sequence
$$
0\longrightarrow \widetilde{R}^{\ast}\longrightarrow
\widetilde{K}^{\ast}\longrightarrow
\ZZ\longrightarrow 0{\quoi},
$$
and $\widetilde{R}^{\ast}$ is divisible by $\ell$ since
$\widetilde{R}$ is strictly local and $\ell$ is invertible in
$\widetilde{R}$, whence
\sisi{\hyperref[X.lem2.2.1]{2.2.1}}{\eqhyperref[X.eq2.2.2]{2.2.2}}.

Consider the Hochschild--Serre spectral sequence
$$
E^{pq}_2=H^p(\bar{G}, H^q(H,
𝛍_{\ell}))\Longrightarrow
H^{p+q}(G,𝛍_\ell){\quoi}.
$$
We have $E^{pq}_2=0$ if $q>1$. If, moreover,
$\cd_{\ell}G<\infty$, the ``boundary'' morphism
$$
H^{n-1}(\bar{G}, H^1(H,𝛍_{\ell}))\longrightarrow
H^{n+1}(\sisi{\bar{G}}{G}, 𝛍_{\ell})
$$
must be bijective for sufficiently large $n$, and it is necessary to
prove that then
$$
H^{n-1}
(\bar{G},H^1(H,𝛍_{\ell}))=H^{n-1}(\bar{G},\ZZ/\ell)=0{\quoi}.
$$

Set $R'=R[x']$ with ${x'}^{\ell}=x$, where $x$ is the local parameter
of $R$, and let $K'$ be the fraction field of $R'$,
$G'=G(\bar{K}/\sisi{k}{K}')$, \etcn The residue field of $R'$ is $k$,
since we have supposed that $\bar{G}$ is an $\ell$-group, and hence that
the $\ell$\sisi{-th}{-th} roots of unity are in $K$. We have a diagram
$$
\xymatrix@M=.25cm@C=1.5cm{
H' \ar[r] \ar@{^{(}->}[d]& G' \ar[r] \ar@{^{(}->}[d]&\bar{G}\\
H  \ar[r]& G \ar[r] & \bar{G},
}
$$
whence\pageoriginale a morphism of Hochschild--Serre spectral
sequences, which gives a commutative diagram
$$
\xymatrix{
H^{n-1}(\bar{G},H^{1}(H',𝛍_l)) \ar[rd]& &\\
                                              &
H^{n+1}(\sisi{\bar{G}}{G},𝛍_l) &\\
H^{n-1}(\bar{G},H^{1}(H,𝛍_l))\ar[uu]^{\varepsilon}
\ar[ru] & &.
}
$$
Since $\cd_{\ell}G<\infty$ implies $\cd_{\ell}G'<\infty$, all the
morphisms in this diagram are isomorphisms for large $n$. But the
morphism
$$
\ZZ/r=H^1(H,𝛍_{\ell})\longrightarrow
H^1(H',𝛍_{\ell})=\ZZ/\ell
$$
which induces the arrow $\epsilon$ is plainly zero, and hence
$H^{n-1}(\bar{G},\ZZ/\ell)=0$,\cqfd

\begin{enonce*}{Theorem 2.3}\etiquette[2.3]{X.thm2.3}
Let $A$ be a noetherian henselian ring of dimension $1$, and let
$\ell\in \PP$ be invertible in $A$. Let $k$ be the residue field of
$A$ and let $R(A)$ be the ring of rational functions
%% say that this is the total ring of fractions? [with an EGA reference]?
on $\Spec A$. We have
$$
\cd_{\ell}R(A)\leq \cd_{\ell}k+1{\quoi},
$$
and equality holds if $\ell\neq 2$ or if $k$ is not orderable.
\end{enonce*}

Proof. Replacing $A$ by $A/P$ for a minimal prime ideal $P$ of $A$ and
applying VIII 1, %% missing reference \XXX
we see that we may suppose $A$ integral. Let $K$ be its fraction field,
and let $R$ be its normalization, which is local and henselian, and
hence a discrete valuation ring (it is noetherian by Krull--Akizuki
(Bourbaki, Alg. Comm. Chap.~7, \S 2, prop.~5)), and let $k'$ be the
residue field of $R$. It is known that $[k':k]<\infty$ (\loccit). Now
$K$ is the fraction field of $R$, and therefore
$$
\cd_{\ell}K=\cd_{\ell}k'+1\leq \cd_{\ell}k+1
$$
by\pageoriginale \hyperref[X.thm2.2]{2.2}, with equality under the last
hypothesis of \hyperref[X.thm2.3]{2.3}, by
\hyperref[X.thm2.1]{2.1}.

\begin{enonce*}[remark]{Example}
Equality does not hold for $\ell=2$, when $A$ is the henselization at
the origin of $\RR[x,y]/(x^2+y^2)$.
\end{enonce*}

\begin{enonce*}{Corollary 2.4}\etiquette[2.4]{X.cor2.4}
Let $A$ be a noetherian local ring with residue field $k$, and let
$\ell\in \PP$ be invertible in $A$. Suppose that $\ell\neq 2$
\sisi{or}{or} that no residue field of $\Spec A$ is orderable. Then
$$
\cd_{\ell}R(A)\geq \cd_{\ell}k+\dim A{\quoi},
$$
where $R(A)$ is the ring of rational functions on $\Spec A$.
\end{enonce*}

Proof. We may suppose $A$ integral, \sisi{with}{with} fraction field
$R(A)=K$. Let $x\in 𝔪_A$\sisi{ be a non-zero-divisor}{}.
%% harmonize maximal/prime ideals \XXX
Then $\dim A/(x)=\dim A-1$. By induction on $\dim A$, we have
$$
\cd_{\ell}R(A/(x))\geq \cd_{\ell}k+\dim A-1{\quoi}.
$$
Let $k'$ be a residue field of $R(A/(x))$ such that
$\cd_{\ell}k'=\cd_{\ell}R(A/(x))$ (\cf VIII 1), and let $P$ %% 𝔭
be the prime ideal that is the kernel of $A\rightarrow k'$. The ring
$A_P$ has dimension $1$. It is therefore enough to prove the
corollary for the local ring $A_P$; that is, we are reduced to the case
in which $\dim A=1$. Let $\widetilde{A}$ be the henselization of $A$.
Every residue field of $R(\widetilde{A})$ identifies with a separable
extension of $R(\widetilde{A})=K$, and therefore
$$
\cd_{\ell}R(A)=\cd_{\ell}K\geq \cd_{\ell}R(\widetilde{A}){\quoi}.
$$
We are thus reduced to the case in which $A$ is henselian of dimension
1, which follows from \hyperref[X.thm2.3]{2.3}.

\begin{enonce*}{Corollary 2.5}\etiquette[2.5]{X.cor2.5}
Let\pageoriginale $X$ be a noetherian scheme and let $\ell\in
\PP$ be invertible on $X$. Then
$$
\cd_{\ell}R(X)\geq \dim X{\quoi}.
$$
\end{enonce*}

Indeed, let $x\in X$ be a point such that $\dim X=\dim \Oo_{X,x}$ (or
such that $\dim \Oo_{X,x}$ is very large if $\dim X=\infty$). Let $A$
be the strict henselization of $\Oo_{X,x}$. Then
$\cd_{\ell}R(X)\geq \cd_{\ell}R(\Oo_{X,x})\geq \cd_{\ell}R(A)$ by
the usual argument. Moreover, $\dim A=\dim X$ (\resp is very large).
If $\ell=2$, and if $\ell$ is invertible on $X$, then the equation
$T^2+1=0$ is separable over $A$, so $-1$ is a square in $A$ because
$A$ is strictly local; hence the residue fields of $A$ are not
orderable. Thus the hypotheses of \hyperref[X.cor2.4]{2.4} are
satisfied for $A$ in every case, and
$\cd_{\ell}R(A)\geq \dim A$, whence the result.

% END INCLUDED SGA4_English_current_buildable_source_20260906/source/10/drafts/SGA4_X_LINES_43_382_English.texfrag


% BEGIN INCLUDED SGA4_English_current_buildable_source_20260906/source/10/drafts/SGA4_X_LINES_383_678_English.texfrag
\section{Field of fractions of a strictly local ring}\etiquette[3]{X.sec3}

\setcounter{subsection}{-1}
\subsection{}\etiquette[3.0]{X.subsec3.0}
In the following sections we shall have essential need of knowing
$\cd_{\ell}R(A)$, where $A$ is a strictly local ring, that is, a
henselian ring with separably closed residue field. Unfortunately,
$\cd_{\ell}R(A)$ is not known in general, for example if
$A=\widehat{\ZZ}_p[[x]](p\neq \ell)$.

It is nevertheless conjectured\SGAArchiveOnly{\nde{The frozen source states that this
was proved by O.~Gabber but supplies no bibliographic target.}} that if
$A$ is noetherian and $\ell$ is invertible in $A$, then one
\sisi{will have}{has}
\begin{equation*}
\cd_{\ell}R(A)=\dim A\quad (\ell\mbox{ invertible in } A){\quoi},\tag{3.1}
\end{equation*}\etiquette[3.1]{X.eq3.1}
at least in the most important cases, for example if $A$ is an
\emph{excellent} ring (\EGA IV 7.8.2).

Note that if $\dim A=1$, \eqhyperref[X.eq3.1]{3.1} is true by
\hyperref[X.thm2.3]{2.3}. Moreover, the inequality $\geq$ holds in any
case by \hyperref[X.cor2.5]{2.5}. We shall prove later (XIX), using
resolution of singularities \cite{sga4.exp.X.bib.3}, that \eqhyperref[X.eq3.1]{3.1}
is also true if $A$ is excellent of residual\pageoriginale
characteristic zero. Here we shall prove \eqhyperref[X.eq3.1]{3.1}
only in the easy case of \hyperref[X.prop3.2]{3.2} below.

Naturally one has $\cd_{\ell}R(A)\leq 1$ if $A$ has characteristic
$\ell$, so this case is trivial. It remains, moreover, to analyze the
case of unequal characteristics with residual characteristic $\ell$,
and one does not expect \eqhyperref[X.eq3.1]{3.1} to be true in this
\sisi{case there}{case}. It will certainly be necessary to take into
account the rank of the module $\Omega^1_k$ ($k$ the residue field),
\ie the degree $[k:k^{\ell}]$.\SGAArchiveOnly{\nde{The frozen source points to work of
Gabber and Orgogozo but supplies no bibliographic target.}}

\begin{enonce*}{Proposition 3.2}\etiquette[3.2]{X.prop3.2}
Let $X$ be a prescheme of finite type over the spectrum of a field $k$
or over the spectrum of a discrete valuation ring $R$, and let
$\ell\in \PP$ be invertible in $k$ (\resp~$R$). Then $\hyperref[X.eq3.1]{3.1}$ is true for
every ring $A$ that is a strict localization of $X$.
\end{enonce*}

\emph{Proof}. Take $X$ integral (one reduces to this case as usual),
and let $n=\dim X$.

Suppose that $X$ is defined over a \emph{separably closed} field $k$.
Then for every ring $A$ that is a strict localization of $X$ at a
\emph{closed} point, one has $\dim A=n$. Since every residue field $K$
of $R(A)$ is identified with an algebraic extension of $R(X)$, and
since
$$
\deg.\tr.(R(X)/k)=\dim X=n{\quoi},
$$
%% \tr. etc. ugly: there is one space too many (tr .) \XXX
one has $\cd_{\ell}K\leq n$ by \hyperref[X.thm2.1]{2.1}, whence
$\cd_{\ell}R(A)\leq n$. The opposite inequality follows from
\hyperref[X.cor2.5]{2.5}.

Now suppose that $X$ is of finite type over $\Spec R$, where $R$ is a
discrete valuation ring, and let $A$ be the strict localization of $X$
at a point $x$. If $x$ does not lie above the closed point of $\Spec
R$, one reduces immediately to the case where $X$ is of finite type
over a field (by localizing $R$). For the same reason, one may also
suppose that the morphism $X\rightarrow \Spec R$ is dominant.

Suppose\pageoriginale for the moment, in addition, that $R$ is a
\emph{strictly local} discrete valuation ring and that the point $x$
is \emph{closed} in $X$. Then one has $\dim A=\dim X=n$ (\EGA IV
5.6.5). Since every residue field $K$ of $R(A)$ is identified with a
separable extension of $R(X)$, one has
$$
\cd_{\ell}R(A)\leq \cd_{\ell}R(X)\leq
\cd_{\ell}L+\deg.\tr.(R(X)/L)=1+(n-1)=n
$$
by \hyperref[X.thm2.1]{2.1} and \hyperref[X.thm2.2]{2.2}, where $L$ is
the field of fractions of $R$; the result follows in this case, taking
\hyperref[X.cor2.5]{2.5} into account.

The general case is deduced from these special cases by means of the
following lemma:

\begin{enonce*}{Lemma 3.3}\etiquette[3.3]{X.lem3.3}
\begin{enumerate}
\renewcommand{\theenumi}{\roman{enumi}}
\renewcommand{\labelenumi}{(\theenumi)}
\item Let $X$ be of finite type over the spectrum of a field $k$, and
  let $A$ be a strict localization ring of $X$ at a point $x$. Then
  there exist a separably closed field $k'$, a scheme $X'$ of finite
  type over $\Spec k'$, and a point $x'$ closed in $X'$ such that $A$
  is isomorphic to a strict localization ring of $X'$ at $x'$.

\item Let $X$ be of finite type over $\Spec R$, where $R$ is a
  discrete valuation ring, let $x$ be a point of $X$ above the closed
  point of $\Spec R$, and let $A$ be a strict localization of $X$ at
  $x$. There exist a strictly local discrete valuation ring $R'$, a
  scheme $X'$ of finite type over $\Spec R'$, and a closed point $x'$
  of $X'$ above the closed point of $\Spec R'$, such that $A$ is
  isomorphic to a strict localization of $X'$ at $x'$.
\end{enumerate}
\end{enonce*}

Proof.
\begin{enumerate}
\renewcommand{\theenumi}{\roman{enumi}}
\renewcommand{\labelenumi}{(\theenumi)}
\item We may take $X$ affine.

Let $Y$ be the closure of $x$ in $X$, with the induced reduced
structure, and let $d=\dim Y$. By the normalization lemma (\cite{sga4.exp.X.bib.2}
Chap. V \S 3 ${\rm n^{\circ}}$ 1), there exists a morphism
$$
X\longrightarrow \Spec k[y_1,\ldots, y_d]=T
$$
inducing\pageoriginale a \emph{finite} morphism $Y\rightarrow T$.
Consequently, the strict localization ring of $\Oo_{X,x}$ is
isomorphic to the strict localization of $\Oo_{X',x'}$, where
$X'=\fprod{X}{\Spec\bar{R(T)}}{T}$, and where $x'$ is a closed point
of $X'$.

\item This case is treated similarly. We may suppose that $X$ is
affine, with ring $A$. Let $Y$ be the closure of $x$ in $X$ with the
induced reduced structure; it is a closed subscheme of the closed
fiber $X_0$ of $X$ over $\Spec R$, and let $d=\dim Y$. Then, lifting
the elements $y_1,\ldots,y_d$ of $A\otimes_R k$ considered in (i) to
elements of $A$, one obtains a morphism
$$
X\longrightarrow \Spec R[y_1,\ldots,y_d]=T
$$
inducing a \emph{finite} morphism $Y\rightarrow T_0$. Let $R'$ be the
strict localization of $R[y_1,\ldots,y_d]=B$ at the prime ideal
$𝔪B$. %% specify that 𝔪 is the maximal ideal of R? \XXX
Then $R'$ is a strictly local discrete valuation ring, and one may
take $X'=\fprod{X}{\Spec R'}{T}$, and $x'$ a closed point of $X'$ above
$x$.
\end{enumerate}

\section{Cohomological dimension: the case where $ℓ$ is invertible in $𝒪_X$}\etiquette[4]{X.sec4}

\begin{enonce*}{Theorem 4.1}\etiquette[4.1]{X.thm4.1}
Let $X$ be a noetherian scheme, let $\ell\in \PP$, and let $\varphi$
be a function on $X$ with values in $\NN$ and satisfying the following
conditions:
\begin{enumerate}
\renewcommand{\theenumi}{\roman{enumi}}
\renewcommand{\labelenumi}{(\theenumi)}
\item For every $x\in X$, $\cd_{\ell}k(x)\leq \varphi (x)$.

\item Let $x\in X$ and let $y\in Y=$ the closure of $x$, with $y\neq
x$. Let $K$ be the ring of rational functions of a strict localization
of $Y$ above $y$. One has $\cd_{\ell}K<\varphi(x)-\varphi(y)$.
\end{enumerate}
Then for every $\ell$-torsion sheaf $F$ on $X$, one has
\begin{gather*}
H^q(X,F)=0\\
q>\varphi(F)\mathop{=}_{\text{\rm def}}\Sup\{\varphi(x)\mid x\in X,
F_{\bar{x}}\neq 0\}{\quoi}.\tag*{if}
\end{gather*}
\end{enonce*}

Proof.\pageoriginale
The hypotheses will also be satisfied by the restriction of
$\varphi$\sisi{,}{} to any subprescheme of $X$, so by noetherian
induction we may suppose the theorem true for every closed subprescheme
of $X$ distinct from $X$. Let $F$ be given, and write
$F=\varinjlim F_i$, where the $F_i$ are the constructible subsheaves
of $F$ (\hyperrefext[IX][IX.prop2.9]{2.9} (iii)).

For every $F_i$, the set of $x\in X$ such that $F_x\neq 0$ is
constructible in $X$, and we may apply the induction hypothesis to
every $F_i$ whose support is not dense. We thus see that the induction
hypothesis in fact implies that the theorem is true for every $F$ that
vanishes at at least one of the maximal points of $X$. We shall
therefore suppose that $F$ does not satisfy this condition.

Now let $i:\Spec R(X)\rightarrow X$ be the inclusion, where $R(X)$ is
the ring of rational functions of $X$, and consider the following
exact sequence, where $K$ and $C$ are defined as kernel and cokernel,
respectively:
$$
0\longrightarrow K\longrightarrow F\longrightarrow i_{\ast}
i^{\ast}F\longrightarrow C\longrightarrow 0{\quoi}.
$$
Here $K$ and $C$ vanish at every maximal point of $X$, so by (ii),
$\varphi(K)$ (\resp $\varphi(C)$) is strictly less than
$n=\varphi(F)$, whence $H^q(X,K)=H^q(X,C)=0$ for $q\geq n$. Thus,
\sisi{}{,} in order to prove $H^q(X,F)=0$ for $q>n$, it is enough, by
the cohomology exact sequences, to prove
$H^q(X,i_{\ast}i^{\ast}F)=0$; hence we may suppose that $F$ is of the
form $i_{\ast}G$, for a suitable $\ell$-sheaf $G$ on $\Spec R(X)$.

Consider the Leray spectral sequence\sisi{.}{}
$$
E^{p,q}_2=H^p(X,R^qi_{\ast}G)\Longrightarrow H^{p+q}(\Spec
R(X),G){\quoi}.
$$
Now $R(X)_{\text{\rm réd}}$ is the product of the fields $k(x)$ for
$x$ maximal in $X$, so $\cd_{\ell}R(X)\leq n$ by (i); therefore the
abutment of the spectral sequence vanishes in
\sisi{$\dim > n$}{dimension strictly greater than $n$}.
We\pageoriginale wish to prove that $E^{p0}_2=0$ for $p>n$. Examining
the spectral sequence, one sees that it is enough to prove that
$E^{pq}_2=0$ if $q>0$ and if $p>n-q-1$. Thus, by induction, it is
enough to prove that $\varphi(R^qi_{\ast}G)<n-q$ for $q>0$. But by
VIII \hyperref[VIII.thm5.2]{5.2}, the fiber of $R^qi_{\ast}G$ at a geometric point $\bar{y}$
above a point $y\in X$ is precisely $H^q(\Spec K, G_K)$, where $K$ is
the ring of rational functions of the strict localization of $X$ at
$\bar{y}$, and $G_K$ is the induced sheaf on $\Spec K$. By (ii), one
has $\cd_{\ell}K<n-\varphi(y)$, whence
$(R^qi_{\ast}G)_{\bar{y}}=0$ if $q\geq n-\varphi(y)$, which gives the
desired result.

\begin{enonce*}{Corollary 4.2}\etiquette[4.2]{X.cor4.2}
Let $X$ be a noetherian scheme, let $\ell\in \PP$, let
$n\in \NN$, and suppose that the following conditions are satisfied:
\begin{enumerate}
\renewcommand{\theenumi}{\roman{enumi}}
\renewcommand{\labelenumi}{(\theenumi)}
\item For every point $y$ of $X$ of codimension $c$, one has
  $\cd_{\ell}k(y)\leq n-2c$.

\item For every ring $A$ that is a strict localization of an
  irreducible closed subscheme $Y$ of $X$, one has
  $\cd_{\ell}R(A)\leq \dim A$.

Then
$$
\cd_{\ell}X\leq n{\quoi}.
$$
More precisely, if $F$ is an $\ell$-sheaf that vanishes at every point
$y$ of codimension $<s$, one has
$$
H^q(X,F)=0\mbox{ if } q>n-2s{\quoi}.
$$
\end{enumerate}
\end{enonce*}

Indeed, one may take $\varphi(y)=n-2\codim y$ and apply
\hyperref[X.thm4.1]{4.1}, whose condition (i) (\resp (ii)) follows
from the corresponding condition (i) (\resp (ii)) of
(\hyperref[X.cor4.2]{4.2}), taking into account for (ii) the inequality
$\dim \Oo_{Y,y}+\dim \Oo_{X,x}\leq \dim \Oo_{X,y}$.

\begin{enonce*}{Corollary 4.3}\etiquette[4.3]{X.cor4.3}
Let $X$ be of finite type and dimension $n$ over a field $k$, and let
$\ell\neq \car k$. Then $\cd_{\ell}X\leq 2n+\cd_{\ell}k$.
\end{enonce*}

One\pageoriginale applies the preceding corollary, replacing $n$ there
by $2n+cd_{\ell}k$, and using \hyperref[X.thm2.1]{2.1} and
\hyperref[X.prop3.2]{3.2} to verify conditions (i) and (ii) of
\hyperref[X.cor4.2]{4.2}, taking into account for (i) the inequality
$$
\deg.\tr.K(y)/k+\dim \Oo_{X,y}\leq \dim X{\quoi}.
$$

\begin{enonce*}{Corollary 4.4}\etiquette[4.4]{X.cor4.4}
Let $X$ be noetherian, and let $\ell\in \PP$ be invertible on $X$.
Suppose that condition {\rm (ii)} of $\hyperref[X.cor4.2]{4.2}$ is satisfied, and moreover
that $\ell\neq 2$ \sisi{where}{or} that no residue field of $X$ is
orderable. Then one has
$$
\cd_{\ell}X\leq \cd_{\ell}R(X)+\dim X\leq 2\cd_{\ell}R(X){\quoi}.
$$
\end{enonce*}

Put $n=\cd_{\ell}R(X)+\dim X$ in \hyperref[X.cor4.2]{4.2}, and apply
\hyperref[X.cor2.4]{2.4}, which implies condition (i) of
\hyperref[X.cor4.2]{4.2}, since
\[\cd_{\ell}k(y)\leq \cd_{\ell}R(\Oo_{X,y})-\dim \Oo_{X,y}\leq
\cd_{\ell}R(X)-\dim\Oo_{X,y}\leq n-2\dim \Oo_{X,y}.
\]
The second inequality in \hyperref[X.cor4.4]{4.4} follows from
\hyperref[X.cor2.5]{2.5}.

\begin{enonce*}[remark]{Example 4.5}\etiquette[4.5]{X.exem4.5}
The inelegant hypothesis on residue fields made in
\hyperref[X.cor4.4]{4.4} is necessary, as is seen by taking
$X=\Spec \RR[x,y,z]/(x^2+y^2+z^2)$, where all the other conditions are
satisfied, with $\dim X=\cd_{\ell}R(X)=\ell=2$, but the sheaf $\ZZ/2$
concentrated at the point $(0,0,0)$ has nonzero cohomology in every
degree, so $\cd_{\ell}X=+\infty$.
\end{enonce*}

\setcounter{subsection}{5}
\subsection{}\etiquette[4.6]{X.subsec4.6}
One is also tempted to try to impose hypotheses in
\hyperref[X.cor4.2]{4.2} only on the closed points of $X$, but this is
not possible. For example, there exist discrete valuation rings with
algebraically closed residue field and with function field of
arbitrarily large cohomological dimension\footnote[1]{To see this, it
  is enough to note that if $X$ is a smooth algebraic scheme over a
  field $k$, admitting a point $x\in X(k)$, then one can find a
  ``formal curve arc'' passing through $x$ whose ``algebraic envelope''
  is $X$, and whose function field therefore contains that $K$ of $X$
  and induces on it a discrete valuation whose residue field is $k$.}.
The spectrum of such a ring will likewise have very large
cohomological dimension.
%% to do (for the note): i -> (i) \XXX

% END INCLUDED SGA4_English_current_buildable_source_20260906/source/10/drafts/SGA4_X_LINES_383_678_English.texfrag


% BEGIN INCLUDED SGA4_English_current_buildable_source_20260906/source/10/drafts/SGA4_X_LINES_679_937_English.texfrag
\section{Cohomological dimension: the case $ℓ=p$}\etiquette[5]{X.sec5}\pageoriginale 

\setcounter{subsection}{-1}
\subsection{}\etiquette[5.0]{X.subsec5.0}
Using Artin\sisi{ - }{-}Schreier theory
(\hyperrefext[IX][IX.Thr3.5]{3.5}), one obtains a better upper bound
than \hyperref[X.cor4.3]{4.3} for the $p$-cohomological dimension of a
scheme $X$ of characteristic~$p$. Let $\cdqc X$ be the largest number
$n$ such that $H^n(X,F)\neq 0$ for at least one \emph{quasi-coherent}
sheaf of modules $F$ on $X$ (in the Zariski sense). We have the
following

\begin{enonce*}{Theorem 5.1}\etiquette[5.1]{X.thm5.1}
Let $X$ be a Noetherian scheme of characteristic $p>0$. Then
$$
\cd_pX\leq \cdqc X+1{\quoi}.
$$
In particular, the $p$-cohomological dimension of a Noetherian affine
scheme $X$ of characteristic $p>0$ is at most $1$.
\end{enonce*}

Proof. Apply \hyperrefext[IX][IX.prop5.5]{5.5}:

Taking account of the fact that the $R^q f_{\ast}F$, $q>0$, vanish for
a finite morphism $f$ and an abelian sheaf $F$ (VIII \hyperref[VIII.prop5.5]{5.5})
%% reference \XXX
(\resp a quasi-coherent sheaf $F$), we are reduced to proving that
$H^q(X,\sisi{i_1}{i_!}(\ZZ/p)_U)=0$ for $i:U\rightarrow X$ an open
immersion and for $q>\cdqc X$.

Let $Y$ be a closed subscheme of $X$ with support $X-U$, and let $J$
be the coherent sheaf of ideals defining $Y$. The morphism
$p:f\mapsto f^p-f$ (\hyperrefext[IX][IX.Thr3.5]{3.5}) induces a
morphism $J\rightarrow J$. Examining the exact sequence
\hyperrefext[IX][IX.Thr3.5]{3.5} and the exact sequence
$$
0\longrightarrow i_{!}(\ZZ/p)_U\longrightarrow
(\ZZ/p)_X\longrightarrow
(\ZZ/p)_Y\longrightarrow 0{\quoi},
$$
one sees that there is in fact an exact sequence
\begin{equation*}
0\longrightarrow i_{!}(\ZZ/p)_U\longrightarrow
J\longrightarrow J\longrightarrow 0{\quoi}.\tag{$\ast$}
\end{equation*}
Since $J$ is coherent, one has $H^q(X,J)=0$ for $q>\cdqc X$, whence the
result\pageoriginale by the long exact cohomology sequence.

For a scheme quasi-projective over a separably closed field $k$, this
upper bound can in certain cases be reduced \sisi{by}{by}~$1$. In fact,
we have

\begin{enonce*}{Corollary 5.2}\etiquette[5.2]{X.cor5.2}
Let $X$ be a scheme of finite type over a field $k$ of characteristic
$p>0$. One has
$$
\cd_pX\leq \dim X+1{\quoi},
$$
and if $k$ is separably closed and $X$ is quasi-projective
\footnote[1]{Using Chow's lemma, it is easy to replace the condition
  ``$X$ quasi-projective'' by ``$X$ separated.''},
one has
$$
\cd_pX\leq \dim X{\quoi}.
$$
\end{enonce*}

Indeed, the first assertion is clear, since $\cdqc X\leq \dim X$
(\cite{sga4.exp.X.bib.6} 4.15.2). Let us treat the second. For such an $X$, the
cohomology groups $H^n(X,F)$ $(n=\dim X)$, for a coherent sheaf $F$,
are finite-dimensional vector spaces over $k$. Let $V=H^n(X,J)$,
where $J$ is the sheaf of ideals occurring in the exact sequence
$(\ast)$. The morphism $\varphi:V\rightarrow V$ deduced from $(\ast)$
is evidently of the form $\varphi=\epsilon-\id$ \sisi{or}{where}
$\epsilon$ satisfies $\epsilon(rv)=r^p\epsilon(V)$ for $r\in k$ and
$v\in V$. The Jacobian of the map $\Phi$ of the affine scheme deduced
from $\Phi$ is therefore $-I$, \sisi{$I$}{where~$I$ is} the identity
matrix. Consequently $\Phi$ is étale, hence surjective because it is
additive, and it follows that $\Phi$ is surjective since $k$ is
separably closed, whence the conclusion immediately.

\section{Cohomological dimension for a prescheme of finite type over $\mathrm{Spec} ℤ$}\etiquette[6]{X.sec6}\pageoriginale

\begin{enonce*}{Proposition 6.1}\etiquette[6.1]{X.prop6.1}
Let $X$ be quasi-finite over $\Spec \ZZ$. Suppose that
$\ell\neq 2$ \sisi{where}{or} that every residue field of $X$ of
characteristic zero is totally imaginary. Then $\cd_{\ell}X\leq 3$.
\end{enonce*}

This is a consequence \sisi{of}{of} \sisi{Theorem}{Theorem}
\hyperref[X.thm2.2]{2.2} (ii), of \sisi{Corollary}{Corollary}
\hyperref[X.cor4.2]{4.2}, and of the fact that a number field (totally
imaginary if $\ell =2$) has $\ell$-cohomological dimension $\leq 2$
(CG II 4.4).

\sisi{It follows in fact from class field theory}{A consequence of
class field theory is} that equality holds if the morphism
$X\rightarrow \Spec \ZZ$ is proper and surjective.

Unfortunately, the results of ${\rm n^{\circ}}$ 4 do not apply as they
stand if $X$ is only of finite type over $\Spec \ZZ$, because
hypothesis (ii) of (\hyperref[X.cor4.2]{4.2}) is not in general
satisfied for local rings of \emph{residual characteristic} $\ell$.
One must therefore use result (\hyperref[X.thm5.1]{5.1}) to obtain the
correct upper bound; we shall now reap the benefit of Theorem
(\hyperref[X.thm5.1]{5.1}).

\begin{enonce*}{Theorem 6.2}\etiquette[6.2]{X.thm6.2}
Let $X$ be of finite type over $\Spec \ZZ$, and let $\ell\in \PP$.
Suppose, if $\ell=2$, that no residue field of $X$ is orderable. Then
$$
\cd_{\ell}X\leq 2\dim X+1{\quoi}.
$$
More precisely, let $F$ be a sheaf of $\ell$-torsion such that for
every point $x$ of $X$ such that $F_{\bar{x}}\neq 0$, one has
$$
\dim\{\bar{x}\}\leq
\begin{cases}
(N-1)/2 &\mbox{if\quad } \ell\neq \car k(x)\\
N-1     &\mbox{if\quad } \ell=\car k(x)
\end{cases}
$$
($\{\bar{x}\}$ is the closure of $\{x\}$); then
$$
H^p(X,F)=0\mbox{ if } p>N{\quoi}.
$$
\end{enonce*}

Proof.\pageoriginale By induction on $X$, we may suppose the theorem
true for every closed subscheme distinct from $X$. If $F$ is as in the
statement and we write $F=\varinjlim F_i$, where the $F_i$ are the
constructible subsheaves of $F$
(\hyperrefext[IX][IX.prop2.9]{2.9} (iii)), then every $F_i$ also
satisfies the same hypotheses. It therefore suffices to prove the
theorem when $F$ is constructible.

In this case, the set $E$ of $x\in X$ such that $F_{\bar{x}}\neq 0$ is
a constructible subset, and the induction hypothesis implies that the
theorem is true for every $F$ such that $E$ is not dense. Moreover, the
theorem is a consequence of \hyperref[X.thm5.1]{5.1} for a sheaf
concentrated on the fiber of $X$ over the point $(\ell)$ of
$\Spec \ZZ$, and of \hyperref[X.cor4.3]{4.3} for a sheaf concentrated
on the fiber over $(p)$, where $p\neq \ell$, taking account of the fact
that $\cd_{\ell}\ZZ/p \ZZ=1$ (CG II 2.2).

Take $X$ reduced, and let $X=X_{0}\cup X_1$, where $X_1$ is the closed
subset which is the union of the irreducible components of $X$ of
characteristic $\ell$, and where $X_{0}$ is the union of the other
irreducible components, $X_{0}$ and $X_1$ being endowed with their
induced reduced structures. Let $j_{\nu}:X_{\nu}\rightarrow X$ be the
inclusions, let $d_{\nu}=\dim X_{\nu}$, and let
$N=\sup \{2d_{0}+1, d_1+1\}$. For every constructible $F$ on $X$, there
is an exact sequence\sisi{.}{}
\begin{equation*}
0\longrightarrow F\longrightarrow
{j_{0}}_{\ast} j_{0}{}^{\ast}F\oplus {j_1}_{\ast} j_1
{}^{\ast}F\longrightarrow C\longrightarrow 0{\quoi},\tag{$\ast$}
\end{equation*}
where $C$ is concentrated on $X_{0}\cap X_1$, which has characteristic
$\ell$ and dimension $\leq d_1-1$. It follows that
$$
H^q(X,C)=0\mbox{ if } q>N-1{\quoi},
$$
because by \hyperref[X.thm5.1]{5.1} one has
$\cd_{\ell}(X_{0}\cap X_1)\leq \dim
\sisi{}{(}X_{0}\cap X_1\sisi{}{)}+1\leq d_1\leq N-1$. Examining the
relative long exact cohomology sequence of $(\ast)$, we are reduced to
the case $X=X_{0}$ \sisi{where}{or} $X=X_1$, and the latter is a
consequence of \hyperref[X.thm5.1]{5.1}, as already noted.

Suppose\pageoriginale therefore that every maximal point of $X$ has
characteristic different from $\ell$, and let
$i_{\nu}:x_{\nu}\rightarrow X$ be the inclusions of the maximal points.
Examining the exact sequence
$$
0\longrightarrow K\longrightarrow F\longrightarrow \bigoplus_{\nu}
i_{\nu\ast}i_{\nu}{}^{\ast}F\longrightarrow C\longrightarrow 0
$$
we are reduced by induction to the case
$F=i_{\nu\ast}i_{\nu}{}^{\ast}F$, that is, where $F$ is of the form
$i_{\ast}G$ for a suitable sheaf $G$ on $\Spec k(x)$, with $x$ a point
of $X$. We may therefore suppose that $X$ is irreducible and reduced,
with generic point $x$, and that the characteristic of $k(x)$ is zero
by \hyperref[X.cor4.3]{4.3}.

Consider the spectral sequence
\begin{equation*}
H^p(X,R^qi_{\ast}G)\Longrightarrow H^{p+q}(\Spec k(x),
G){\quoi}.\tag{${\ast}{\ast}$} 
\end{equation*}
Let $s=\dim X$. We must prove that
\begin{equation*}
E^{p0}_2=0\mbox{ for } p> 2s+1{\quoi}.\tag{$+$}
\end{equation*}
Now $k(x)[i]$ has transcendence degree $s-1$ over $\QQ[i]$
%% say that i = √(-1)? \XXX
which is a field of cohomological dimension~$2$. Since $k(x)$ is not
orderable if $\ell=2$, it follows from \hyperref[X.thm2.1]{2.1} that
$\cd_{\ell}k(x)=s+1$. This implies that the target of the spectral
sequence vanishes for $n>s+1$. It therefore suffices, for $(+)$, to
prove that in $({\ast}{\ast})$ one has
\begin{equation*}
E^{pq}_2=0\mbox{ if } p>2s-q\mbox{ and } q>0{\quoi}.\tag{${+}{+}$}
\end{equation*}

Let $y$ be a point of $X$, and let $K$ be the ring of fractions of a
strict localization at a geometric point $\widetilde{y}$ of $X$ above
$y$. Then $K$ is a direct product of fields, each of transcendence
degree $s-1$ over a strict localization of $\ZZ$, and hence, by
\hyperref[X.thm2.1]{2.1} and \hyperref[X.thm2.2]{2.2}, one has
$\cd_{\ell}K\leq s$. Moreover, one may apply
\hyperref[X.prop3.2]{3.2} if the residual characteristic of $y$ is not
$\ell$. Thus one has
\begin{equation*}
(R^qi_{\ast}G)_{\widetilde{y}}=0
\begin{cases}
\mbox{ if $q>\codim y$, \ie $\dim \bar{y}>s-q$, and $\ell\neq \car
k(y)$,}\\
\mbox{or if $q>s$ in all cases.}
\end{cases}
\tag{$''$}
\end{equation*}
In\pageoriginale particular, one has $E^{pq}_2=0$ if $q>s$. To verify
$({+}{+})$ in the case $0<q\leq s$, apply the induction hypothesis to
the sheaves $R^qi_{\ast}(G)$.

According to the statement and $('')$, one must set $N'=2s-q$ and
verify the two inequalities
\begin{align*}
& s-q\leq (N'-1)_{/2}\\
& s-1=\dim X_{\ell}\leq N'-1{\quoi},
\end{align*}
where $X_{\ell}$ is the \sisi{fibered space}{fiber} of $X$ over the
\sisi{point}{point} $(\ell)$ of $\Spec \ZZ$. They do indeed hold if
$0<q\leq s$\sisi{,}{.} \cqfd

\begin{thebibliography}{99}
\bibitem{sga4.exp.X.bib.1} 
Ax, J.  Proof of some conjectures of Serre on
cohomological dimension, Proc. Amer. Math. Society 16 (1965)
1214-1221.

\bibitem{sga4.exp.X.bib.2} Bourbaki, N. Algèbre \sisi{Commutative}{commutative}, Paris (1960)

\bibitem{sga4.exp.X.bib.3} 
Hironaka, H.  Resolution of singularities of an algebraic surface
over a field of characteristic zero, Annals of Math. Vol. 79, No. 1, 2
(1964)

 
\bibitem{sga4.exp.X.bib.4} Nagata, M. Local Rings, New York (1962)

\bibitem{sga4.exp.X.bib.5} 
Serre, J-P.  Sur la dimension cohomologique \sisi{de}{des} groupes profinis,
Topology Vol. 3, p. 413 (1965)


\bibitem{sga4.exp.X.bib.6} Godement, R. \sisi{Théorie des Faisceaux}{Topologie algébrique et théorie des faisceaux}, Paris (1958).

\bibitem{sga4.exp.X.bib.7} 
E. Artin et A.  Schreier - Eine Kennzeichnung der \sisi{reel}{reell}
abgeschlossenen Körper, Hamb. Abh. 5 (1927) pp. 225-231.


\bibitem{sga4.exp.X.bib.8} Serre, J.P. Cohomologie \sisi{Galoisienne}{galoisienne}, Lecture Notes
${\rm n^{\circ}}$ 5 (Springer).
\end{thebibliography}

\ifx\danslelivre\undefined
\end{document}
\fi

% END INCLUDED SGA4_English_current_buildable_source_20260906/source/10/drafts/SGA4_X_LINES_679_937_English.texfrag


% END INCLUDED SGA4_English_current_buildable_source_20260906/source/10/10_English.tex


\makeatletter
\def\input@path{{../11/}}
\makeatother
% BEGIN INCLUDED SGA4_English_current_buildable_source_20260906/source/11/11_English.tex
\ifx\danslelivre\undefined
% Release engine: XeLaTeX (the LuaTeX compatibility layer breaks Xy-pic output).
\documentclass[10pt]{../fichiers-configuration-latex/smfart}
% BEGIN INCLUDED SGA4_English_current_buildable_source_20260906/source/fichiers-configuration-latex/paquets.tex
\usepackage{polyglossia}
%\usepackage[francais]{babel}
\usepackage{fontspec}
\usepackage[math-style=ISO,bold-style=upright,sans-style=upright]{unicode-math}
%\usepackage[mathletters]{ucs}
%\usepackage[utf8]{inputenc}
%\usepackage[T1]{fontenc}
\usepackage[unicode,hypertexnames=false,colorlinks=true,linkcolor=blue,citecolor=Sepia,urlcolor=magenta]{hyperref}
\setmainlanguage{french}
\setmainfont[Ligatures=TeX]{Linux Libertine O}
% variante possible : TeXGyrePagella ?
\setmathfont{XITS Math}

% maths
\usepackage{amsmath}
%\usepackage{amssymb}
%\usepackage{mathrsfs}
%\usepackage{stmaryrd}
%\usepackage{euscript}

% couleurs
\usepackage[usenames,dvipsnames]{xcolor}

% graphiques
\usepackage{graphicx}
\usepackage[all]{xy}
%\usepackage{tikz}
%\usetikzlibrary{matrix,arrows,calc}
%\tikzset{column sep=3em,row sep=2em,text height=1.5ex, text depth=0.25ex,>=angle 90}
%\tikzset{description/.style={fill=white,inner sep=1pt}}

% hyperliens, 'etiquettes
\usepackage{xr-hyper}
\usepackage{../fichiers-configuration-latex/smfthm}
%\ifx\cacheretiquettes\undefined
%\usepackage[inline]{showlabels} \else \fi % pour rendre visibles les 'etiquettes
%\usepackage{../fichiers-configuration-latex/makeidx} % pour faire un index
%\usepackage[pagebackref,unicode]{hyperref}  % pour voir o`u sont cit'ees les r'ef'erences
% hyperref is loaded once above with the cumulative-reader destination policy.
%\hypersetup{breaklinks=true} %permet le retour `a la ligne dans les liens trop longs. Incompatible avec pdflatex
%\renewcommand*{\backref}[1]{{\small [cit\'e page(s)~#1.]}}

\usepackage{xspace}
\usepackage{pifont} % pour expose XVII

% END INCLUDED SGA4_English_current_buildable_source_20260906/source/fichiers-configuration-latex/paquets.tex

% BEGIN INCLUDED SGA4_English_current_buildable_source_20260906/source/fichiers-configuration-latex/numerotation.tex
\renewcommand{\thefootnote}{\roman{footnote}} % pour mettre les notes en bas de page sous la forme (ii) et non (2), qui ressemble trop (?) a un exposant.

\newcounter{pageoriginal}
\newcommand{\pageoriginale}{\refstepcounter{pageoriginal}\marginpar{\footnotesize\textbf{\thepageoriginal}}\xspace}
\let\pageoriginaled\pageoriginale

\usepackage{ifthen}
\newboolean{original}
\def\sisi{\ifthenelse{\boolean{original}}}
\setboolean{original}{false}

\newcounter{cnote}
\def\nde#1{\ifthenelse{\boolean{original}}{}{\stepcounter{cnote}\def\thefootnote{\arabic{cnote}}\protect\footnote[\value{cnote}]{N.D.E.: #1}}}
\def\ndemark{\ifthenelse{\boolean{original}}{}{\stepcounter{cnote}\def\thefootnote{\arabic{cnote}}\protect\footnotemark[\value{cnote}]}}
\def\ndetext#1{\ifthenelse{\boolean{original}}{}{\def\thefootnote{\arabic{cnote}}\protect\footnotetext[\value{cnote}]{N.D.E.: #1}}}

% END INCLUDED SGA4_English_current_buildable_source_20260906/source/fichiers-configuration-latex/numerotation.tex

% BEGIN INCLUDED SGA4_English_current_buildable_source_20260906/source/fichiers-configuration-latex/macros.tex
\setmainlanguage{french}
\setotherlanguage{english}


%% Fixed, engine-independent release identifier.  The inherited dynamic Git
%% probe used Lua io.popen, exposed build-state-dependent output, and printed
%% a visible question mark outside a Git checkout.
\def\showgitstatus{English release r2, 2026-07-23}


% macros pour LaTeX
\renewcommand{\mathsf}{}

% Police différente pour U+2113 SCRIPT SMALL L
\setmathfont[range="2113]{latinmodern-math.otf}

% Give every inherited semantic label its own stable PDF destination instead
% of inheriting whichever automatic equation/page anchor happened to be live.
% The legacy label remains as a compatibility alias used by existing edges.
\makeatletter
\def\etiquette[#1]#2{%
  \Hy@raisedlink{\hyper@anchorstart{sga4.semantic.#2}\hyper@anchorend}%
  \def\@currentHref{sga4.semantic.#2}%
  \def\@currentlabel{#1}%
  \label{#2}%
}
\makeatother
\def\eqhyperref[#1]#2{(\hyperref[#1]{#2})}
\def\pagehyperref[#1]#2{(page \hyperref[#1]{#2})}
\def\hyperrefext[#1][#2]#3{\hyperref[#2]{#1 #3}}
\def\eqhyperrefext[#1][#2]#3{\hyperref[#2]{#1 #3}}     %cf. exposes 17
\def\neweqhyperrefext[#1][#2]#3{\hyperref[#2]{#1 #3}}  %cf. exposes 17

\def\shat#1{\protect\widehat{#1}}
\def\stilde#1{\protect\widetilde{#1}}
%%%%%%%%%%%%%%%%%%


\def\etcn{etc.\xspace}
\def\etco{etc...\xspace}
\def\etc{\sisi{\etco}{\etcn}}
\def\ldotsn{\sisi{...}{\etcn}}

\def\cf{{cf.}\kern.3em}
\def\Cf{{Cf.}\kern.3em}
\def\ie{{i.e.}\kern.3em}
\def\iev{{i.e.},\xspace}
\def\loccito{loc.\kern2pt cit.\xspace}
\def\loccitn{\textit{loc.\kern2pt cit.}\xspace}
\def\loccit{\sisi{\loccito}{\loccitn}}
\def\EGA{\textup{EGA}\kern.2em}
\def\SGA{\textup{SGA}\kern.3em}

\def\quoi{}

\def\quav{\quad,\quad}
\def\quaet{\quad \text{et}\quad}
\def\fibreproduct#1#2#3{#1 ×_{#3} #2}
\let\fprod\fibreproduct
\def\ccheck#1#2{\mbox{${\check{\check{\rm #1}}}^{\lower2pt\hbox{$\scriptstyle #2$}}$}}
\def\sheaf#1{\protect\mathcal{#1}}
\def\vieuxsheaf#1{{#1}} % cf. problème avec sheaf{a\Phi_i} par exemple...
\DeclareMathOperator{\Hom}{Hom}
\DeclareMathOperator{\can}{can}
\DeclareMathOperator{\Ind}{Ind}
\DeclareMathOperator{\ind}{ind}
%\DeclareMathOperator{\Hhomtop}{\mathcal{H}\!\mathit{omtop}}
\DeclareMathOperator{\Hhomtop}{\mathbf{Homtop}}
\DeclareMathOperator{\Point}{Point}
\DeclareMathOperator{\Pro}{Pro}
\DeclareMathOperator{\Top}{Top}
\DeclareMathOperator{\Mod}{Mod}
\DeclareMathOperator{\Cart}{Cart}
%\DeclareMathOperator{\Mmorsite}{\mathcal{M}\!\mathit{orsite}}
\DeclareMathOperator{\Mmorsite}{\mathbf{Morsite}}
\DeclareMathOperator{\loc}{loc}
\DeclareMathOperator{\Loc}{Loc}
\DeclareMathOperator{\Pic}{Pic}
\DeclareMathOperator{\Et}{Et}
\DeclareMathOperator{\card}{card}
\DeclareMathOperator{\Ens}{Ens}
\DeclareMathOperator{\h}{\underline{H}}
\DeclareMathOperator{\Fonct}{Fonct}
\DeclareMathOperator{\Ker}{Ker}
\DeclareMathOperator{\but}{but}
\DeclareMathOperator{\Ob}{Ob}
\DeclareMathOperator{\ob}{ob}
\DeclareMathOperator{\Lib}{Lib}
\DeclareMathOperator{\esj}{esj}
\DeclareMathOperator{\id}{id}
\DeclareMathOperator{\Id}{Id}
\DeclareMathOperator{\ad}{ad}
\DeclareMathOperator{\Iim}{Im}
\DeclareMathOperator{\pr}{pr}
\DeclareMathOperator{\Cat}{Cat}
\DeclareMathOperator{\car}{car}
\DeclareMathOperator{\SysInd}{SysInd}
\DeclareMathOperator{\Sup}{Sup}
\DeclareMathOperator{\Coker}{Coker}
\DeclareMathOperator{\Filt}{Filt}
\DeclareMathOperator{\acc}{acc}
\DeclareMathOperator{\sur}{sur}
\DeclareMathOperator{\cat}{cat}
\DeclareMathOperator{\Ord}{Ord}
\DeclareMathOperator{\Gr}{Gr}
\DeclareMathOperator{\Cov}{Cov}
\DeclareMathOperator{\Caf}{Caf}
\DeclareMathOperator{\coker}{coker}
\DeclareMathOperator{\coim}{coim}
\DeclareMathOperator{\im}{im}
\DeclareMathOperator{\Ouv}{Ouv}
\DeclareMathOperator{\diag}{diag}
\DeclareMathOperator{\Esp}{Esp}
\DeclareMathOperator{\TOP}{TOP}
\DeclareMathOperator{\Unif}{Unif}
\DeclareMathOperator{\Iind}{\mathbf{Ind}}
\DeclareMathOperator{\Hhomcart}{\mathbf{Homcart}}
\DeclareMathOperator{\Hhomtopess}{\mathbf{Homtopess}}
\DeclareMathOperator{\Ppoints}{\mathbf{Points}}
\DeclareMathOperator{\Ppoint}{\mathbf{Point}}
\DeclareMathOperator{\Ttors}{\mathbf{Tors}}
\DeclareMathOperator{\Ssex}{\mathbf{Sex}}
\DeclareMathOperator{\Ffib}{\mathbf{Fib}}
\DeclareMathOperator{\Vvois}{\mathbf{Vois}}
\DeclareMathOperator{\Ppointess}{\mathbf{Pointess}}
\DeclareMathOperator{\scrHom}{\mathbf{Hom}}
\DeclareMathOperator{\sheafPic}{\mathbf{Pic}}
\DeclareMathOperator{\sheafAut}{\mathbf{Aut}}
\DeclareMathOperator{\Hhomtoplocan}{\mathbf{Homtoplocan}}
\DeclareMathOperator{\Hhomtopan}{\mathbf{Homtopan}}
\DeclareMathOperator{\Pptgeom}{\mathbf{Ptgéom}}
\DeclareMathOperator{\Morsite}{Morsite}
\DeclareMathOperator{\Fib}{Fib}
\DeclareMathOperator{\Fais}{Fais}
\DeclareMathOperator{\kar}{kar}
\DeclareMathOperator{\Groupes}{Groupes}
\DeclareMathOperator{\Restr}{Restr}
\DeclareMathOperator{\Isom}{Isom}
\DeclareMathOperator{\Prol}{Prol}
\DeclareMathOperator{\Res}{Res}
\DeclareMathOperator{\Tors}{Tors}
\DeclareMathOperator{\esp}{esp}
\DeclareMathOperator{\es}{es}
\DeclareMathOperator{\od}{od}
\DeclareMathOperator{\gr}{gr}
\DeclareMathOperator{\sob}{sob}
\DeclareMathOperator{\dfn}{dfn}
\DeclareMathOperator{\Sch}{Sch}
\DeclareMathOperator{\ouv}{ouv}
\DeclareMathOperator{\cons}{cons}
\DeclareMathOperator{\Nerf}{Nerf}
\DeclareMathOperator{\Points}{Points}
\DeclareMathOperator{\Espsob}{Espsob}
\DeclareMathOperator{\Homtop}{Homtop}
\DeclareMathOperator{\Pt}{Pt}
\DeclareMathOperator{\Quot}{Quot}
\DeclareMathOperator{\glob}{glob}
\DeclareMathOperator{\Homfib}{Homfib}
\DeclareMathOperator{\Qouv}{Qouv}
\DeclareMathOperator{\Quotouv}{Quotouv}
\DeclareMathOperator{\ens}{ens}
\DeclareMathOperator{\ab}{ab}
\DeclareMathOperator{\an}{an}
\DeclareMathOperator{\fs}{fs}
\DeclareMathOperator{\red}{red}
\DeclareMathOperator{\Ab}{Ab}
\DeclareMathOperator{\cosq}{cosq}
\DeclareMathOperator{\Ann}{Ann}
\DeclareMathOperator{\Zar}{Zar}


%%%%%%%%%%%%
%%%%% SGA4-2
\DeclareMathOperator{\Sex}{Sex}
\DeclareMathOperator{\vois}{vois}
\DeclareMathOperator{\Vois}{Vois}
\DeclareMathOperator{\Ext}{Ext}
\DeclareMathOperator{\ext}{ext}
\DeclareMathOperator{\Fer}{Fer}
\DeclareMathOperator{\Add}{Add}
\DeclareMathOperator{\Cocart}{Cocart}
\DeclareMathOperator{\cosk}{cosk}
\DeclareMathOperator{\plat}{plat}

\DeclareMathOperator{\faisc}{faisc}
\DeclareMathOperator{\sk}{sk}
\DeclareMathOperator{\ch}{ch}

\def\HH{\mathbf{H}}
\def\qc{\mathrm{qc}}
\def\et{\mathrm{et}}
\def\coh{\mathrm{coh}}

\def\Hvr#1{%
\begin{array}[b]{l}
{\kern2pt\scriptstyle r}\\[-8pt]
{\kern1pt\scriptstyle \vee}\\[-5pt]
H^{#1}
\end{array}}

\def\Hvv#1{%
\begin{array}[b]{l}
\kern2pt\rotatebox{90}{${\scriptscriptstyle\ll}$}\\[-6pt]
H^{#1}
\end{array}}

\def\Hvi#1{%
\begin{array}[b]{l}
{\kern0.2pt\scriptscriptstyle \infty}\\[-8pt]
{\kern1pt\scriptstyle \vee}\\[-5pt]
H^{#1}
\end{array}}

\def\Hv#1{%
\begin{array}[b]{l}
{\kern1pt\scriptstyle \vee}\\[-5pt]
H^{#1}
\end{array}}

\DeclareMathOperator{\cart}{cart}
\DeclareMathOperator{\Ffaisc}{\mathbf{Faisc}}


\def\Ev#1{%
\begin{array}[b]{l}
{\kern1pt\scriptstyle \vee}\\[-5pt]
E'_{#1}
\end{array}}

\DeclareMathOperator{\Cont}{Cont}
\DeclareMathOperator{\Ccont}{\mathscr{C}\!\mathit{ont}}
\DeclareMathOperator{\Site}{Site}
\newcommand\MM{\mathbf{M}}
\def\limtop#1{{\displaystyle{\varprojlim_{#1}\!\mbox{\rm top }}}}
\def\Limtop{\mathop{\varprojlim\!\mbox{\rm top}}}
\DeclareMathOperator{\spec}{spec}
\def\ss{\scriptstyle{\mathscr{S}}}
\DeclareMathOperator{\Gal}{Gal}
\def\smallm{\scriptstyle{\mathscr{M}}}
\DeclareMathOperator{\End}{End}
\DeclareMathOperator{\ass}{ass}
\DeclareMathOperator{\zar}{zar}
\DeclareMathOperator{\Ogg}{Ogg}
\DeclareMathOperator{\res}{res}
\DeclareMathOperator{\tr}{tr}
\DeclareMathOperator{\Supp}{Supp}
\DeclareMathOperator{\codim}{codim}
\DeclareMathOperator{\rk}{rk}
\DeclareMathOperator{\cur}{cur}
\DeclareMathOperator{\cdqc}{cdqc}
\newcommand\DD{\mathbf{D}}
\DeclareMathOperator{\cl}{cl}
\def\JJ{\mathbf{J}}
\def\FF{\mathbf{F}}
\def\BB{\mathbf{B}}
\DeclareMathOperator{\Alget}{Alget}
\DeclareMathOperator{\Ass}{Ass}
\DeclareMathOperator{\supp}{supp}
\DeclareMathOperator{\rad}{rad}
\DeclareMathOperator{\prof}{prof}
\DeclareMathOperator{\cint}{int}
\DeclareMathOperator{\sic}{sic}
\DeclareMathOperator{\Tor}{Tor}
\DeclareMathOperator{\tors}{tors}
\DeclareMathOperator{\Sym}{Sym}
\DeclareMathOperator{\parf}{parf}
\DeclareMathOperator{\tor}{tor}
\DeclareMathOperator{\defn}{defn}
\DeclareMathOperator{\Div}{Div}
\DeclareMathOperator{\Pol}{Pol}
\DeclareMathOperator{\Fex}{Fex}
\DeclareMathOperator{\Lie}{Lie}
\DeclareMathOperator{\tot}{tot}
\DeclareMathOperator{\Aut}{Aut}
\DeclareMathOperator{\HOM}{HOM}
\DeclareMathOperator{\pch}{pch}
\DeclareMathOperator{\Ch}{Ch}
\DeclareMathOperator{\EXT}{EXT}
\DeclareMathOperator{\PIC}{PIC}
\DeclareMathOperator{\Rf}{Rf}
\DeclareMathOperator{\inj}{inj}
\newcommand\VV{\mathbf{V}}

\def\Xv{%
\begin{array}[b]{l}
{\kern1pt\scriptstyle \vee}\\[-5pt]
X
\end{array}}

%%%%%%%
\newcommand\TT{\mathbf{T}}
\newcommand\XX{\mathbf{X}}
\newcommand\EE{\mathbf{E}}
%%%%%%

\def\qedsymbol{\textup{C.Q.F.D.}}
\let\CQFD\qed
\let\cqfd\qed

%%%%%\bigdot
\font\bigsymb=cmsy10 at 4pt
\def\bigdot{{\kern1.2pt\raise 1.5pt\hbox{\bigsymb\char15}}}
%%%%%

\def\Lright{\mathop{\vcenter{\hbox{\hbox{$\longrightarrow$}\lower8pt \hbox{\kern-16pt\hbox{$\longrightarrow$}}}}}}

\def\Aa{{\mathcal A}}
\def\Bb{{\mathcal B}}
\def\Cc{{\mathcal C}}
\def\Dd{{\mathcal D}}
\def\Ee{{\mathcal E}}
\def\Ff{{\mathcal F}}
\def\Gg{{\mathcal G}}
\def\Hh{{\mathcal H}}
\def\Ii{{\mathcal I}}
\def\Jj{{\mathcal J}}
\def\Kk{{\mathcal K}}
\def\Ll{{\mathcal L}}
\def\Mm{{\mathcal M}}
\def\Nn{{\mathcal N}}
\def\Oo{{\mathcal O}}
\def\Pp{{\mathcal P}}
\def\Qq{{\mathcal Q}}
\def\Rr{{\mathcal R}}
\def\Ss{{\mathcal S}}
\def\Tt{{\mathcal T}}
\def\Uu{{\mathcal U}}
\def\Vv{{\mathcal V}}
\def\Ww{{\mathcal W}}
\def\Xx{{\mathcal X}}
\def\Yy{{\mathcal Y}}
\def\Zz{{\mathcal Z}}

\newcommand{\CC}{\mathbf{C}}
\newcommand{\GG}{\mathbf{G}}
\newcommand{\KK}{\mathbf{K}}
\newcommand{\LL}{\mathbf{L}}
\newcommand{\NN}{\mathbf{N}}
\newcommand{\PP}{\mathbf{P}}
\newcommand\QQ{\mathbf{Q}}
\newcommand\RR{\mathbf{R}}
\newcommand\ZZ{\mathbf{Z}}

\def\to{\mathchoice{\longrightarrow}{\rightarrow}{\rightarrow}{\rightarrow}}
\def\from{\mathchoice{\longleftarrow}{\leftarrow}{\leftarrow}{\leftarrow}}
\def\mto{\mathchoice{\longmapsto}{\mapsto}{\mapsto}{\mapsto}}
\def\mfrom{\mathrel{\from\joinrel\mapstochar\,}}
\def\hto{\mathrel{\lhook\joinrel\to}}
\def\hfrom{\mathrel{\from\joinrel\rhook}}

\def\implies{\mathchoice{\Longrightarrow}{\Rightarrow}{\Rightarrow}{\Rightarrow}}
\def\implied{\mathchoice{\Longleftarrow}{\Leftarrow}{\Leftarrow}{\Leftarrow}}
\def\iff{\mathchoice{\Longleftrightarrow}{\Leftrightarrow}{\Leftrightarrow}{\Leftrightarrow}}

\def\To#1{\mathchoice{\xrightarrow{\labelstyle~~#1~}}{\stackrel{#1}{\longrightarrow}}{}{}}
\def\Mto#1{\mathrel{\mapstochar\To{#1}}}
\def\From#1{\mathchoice{\xleftarrow{\labelstyle~~#1~}}{\stackrel{#1}{\longleftarrow}}{}{}}

\newcommand{\isomto}{\overset{\sim}{\longrightarrow}}
\newcommand{\isomfrom}{\overset{\sim}{\longleftarrow}}

\newcommand{\IFF}{\unskip~$\iff$~\ignorespaces}
\newcommand{\THEN}{\unskip~$\implies$~\ignorespaces}

\let\leq\leqslant
\let\geq\geqslant
\let\le\leqslant
\let\ge\geqslant

\def\cf{{cf.}\kern.3em}
\def\Cf{{Cf.}\kern.3em}
\def\ie{{i.e.}\kern.3em}
\def\iev{{i.e.},\xspace}
%\def\loccito{loc.\kern2pt cit.\xspace}
%\def\loccitn{\textit{loc.\kern2pt cit.}\xspace}
%\def\loccit{\sisi{\loccito}{\loccitn}}
\def\EGA{\textup{EGA}\kern.2em}
\def\SGA{\textup{SGA}\kern.3em}

\let\oldS\S
\def\S{\oldS\kern .15em}
\def\ptbl{.\kern .2em}
\def\resp{resp.\kern.3em}
\def\Exp{Exp.\kern.3em}

\let\bar\overline
\let\hat\widehat
\let\tilde\widetilde
\let\dpl\displaystyle
\let\ldots\dots
\let\emptyset\varnothing
\let\skpt\skippointrait
\let\epsilon\varepsilon

\def\sfrac#1#2{{#1}/{#2}}

\DeclareMathOperator{\cd}{cd}
%\DeclareMathOperator{\Hhom}{\mathcal{H}\!\mathit{om}}
\DeclareMathOperator{\Hhom}{\mathbf{Hom}}
%\DeclareMathOperator{\Eext}{\mathcal{E}\!\mathit{xt}}
\DeclareMathOperator{\Eext}{\mathbf{Ext}}
\DeclareMathOperator{\Spec}{Spec}
\DeclareMathOperator{\Tr}{Tr}

\let\emptyset=\varnothing

%%%%%%%%%%%%%%%
%%% expose XVII
%%%%%%%%%%%%%%%

\makeatletter
\let\old@optexpose\optexpose
\newcommand{\empty@optexpose}[1]{}
\newcommand{\paren@optexpose}[1]{\old@optexpose{#1}(}
\newcommand{\refnoexpose}[1]{\let\optexpose\empty@optexpose\ref{#1}\let\optexpose\old@optexpose}%\refnoexpose{X.cor5.2} engendre "5.2" au lieu de "X 5.2"
\newcommand{\neweqref}[1]{\let\optexpose\paren@optexpose\textup{\ref{#1})}\let\optexpose\old@optexpose}%\neweqref{XII.eq4.1} engendre "XII (4.1)" au lieu de "(XII 4.1)"
\makeatother
\newcommand{\refvb}[2][VI${}_{\text{B}}$ ]{\sisi{#1\refnoexpose{#2}}{\ref{#2}}}

\newcommand{\leftexp}[2]{{\vphantom{#2}}^{#1}{#2}}
\newcommand{\AAA}{\mathbf{A}}
\newcommand{\Sf}{\mathfrak{S}}
\DeclareMathOperator*{\quotprojlim}{\text{``$\varprojlim$''}}
\DeclareMathOperator*{\quotinjlim}{\text{``$\varinjlim$''}}
\DeclareMathOperator{\Fl}{F\sisi{\ell}{l}}
\DeclareMathOperator{\Pl}{P\sisi{\ell}{l}}
%\DeclareMathOperator{\cl}{c\sisi{\ell}{l}}
\DeclareMathOperator{\tdf}{tdf}
\DeclareMathOperator{\rg}{rg}
\newcommand{\iInd}{\sisi{I}{i}nd}
\newcommand{\iind}{ind}
\newcommand{\ppro}{pro}
\newcommand{\tTor}{\sisi{T}{t}or}
\newcommand{\ttor}{tor}
\newcommand{\chap}{Chap.}
\newcommand{\Ffex}{\Ff\mathit{ex}}
\newcommand{\sheafk}{\underline{k}}
\newcommand{\sheafn}{\underline{n}}
\newcommand{\sheafGamma}{\underline{\Gamma}}

% Repeated references to an already typeset footnote must not use
% \footnotemark without a matching \footnotetext: Hyperref would then create
% a live annotation to a destination that cannot exist.  This helper keeps the
% current printed footnote mark and links it to an explicit stable note target.
\makeatletter
\newcommand{\sgarepeatedfootnotemark}[1]{%
  \begingroup
  \edef\@thefnmark{\thefootnote}%
  \hyperlink{#1}{\@makefnmark}%
  \endgroup}
\makeatother

\sisi{
\DeclareMathOperator{\simarrow}{\rule{.15mm}{1.8mm}\kern -3.4pt\rightsquigarrow}
\newcommand{\sheafell}{\underline{\ell}}
\newcommand{\sheafZZ}{\underline{\ZZ}}
}{
\newcommand{\nderef}[1]{la N.D.E.~(\ref{#1}) page~\pageref{#1}}
}

%%%%%%%%%%%%%%%%%
%%%%% expose XVII
%%%%%%%%%%%%%%%%%

% END INCLUDED SGA4_English_current_buildable_source_20260906/source/fichiers-configuration-latex/macros.tex


\externaldocument{../01/01_English.tex}
\externaldocument{../02/02_English.tex}
\externaldocument{../03/03_English.tex}
\externaldocument{../04/04_English.tex}
\externaldocument{../05/05_English.tex}
\externaldocument{../05bis/05bis_English.tex}
\externaldocument{../06/06_English.tex}
\externaldocument{../07/07_English.tex}
\externaldocument{../08/08_English.tex}
\externaldocument{../09/09_English.tex}
\externaldocument{../10/10_English.tex}

\setboolean{original}{false}

\begin{document}

\begin{center}
XI. Comparison with classical cohomology: the case of a smooth scheme\\
M. Artin
\end{center}

\selectlanguage{english}version: \showgitstatus
\tableofcontents

\else
\chapter{Comparison with classical cohomology: the case of a smooth scheme}
%\addtocontents{toc}{M. Artin}
\begin{center}
M. Artin
\end{center}
\fi
\phantomsection\MakeLinkTarget*{sga4.exp.XI}\label{sga4.exp.XI}

% BEGIN INCLUDED SGA4_English_current_buildable_source_20260906/source/11/drafts/SGA4_XI_LINES_43_199_English.texfrag
\setcounter{pageoriginal}{63}

\section{Introduction}\etiquette[1]{XI.sec1}\pageoriginale

In Section~\hyperref[XI.sec4]{4} we shall prove that étale cohomology
with values in a locally constant torsion sheaf agrees with classical
cohomology for a scheme $X$ smooth over $\Spec\CC$. The proof, which is
essentially the one given in \cite{sga4.exp.XI.bib.1}, is elementary once one has the
theorem of Grauert recalled below (\hyperref[XI.thm4.3]{4.3} (iii)),
which says that every finite covering of $X(\CC)$ comes from an étale
covering of $X$; this is why we give it here, although we shall prove a
more complete result later, using Hironaka's resolution of
singularities theorem \cite{sga4.exp.XI.bib.3}\footnote[1]{Note, moreover, that
  Hironaka's resolution of singularities also gives a very simple
  proof of Grauert's result, and in this capacity may in any case be
  regarded as underlying the comparison theorems. See in this
  connection the exposé by Mme M. Raynaud in \SGA 1 XII (North Holland
  Pub. Cie).}. Furthermore, the good neighborhoods constructed in
\hyperref[XI.sec3]{Section~3} could be useful in other\pageoriginale situations.

Note that étale cohomology and classical cohomology do not agree for
sheaves that are not torsion. For example, one has
$H^1(X_{\et},\ZZ)=0$ for every normal scheme $X$
(\hyperrefext[IX][IX.prop3.6]{3.6} (ii)), but classical cohomology does
not necessarily vanish.

\section{Existence of sufficiently general hyperplane sections}\etiquette[2]{XI.sec2}

The following result is classical, but a reference is lacking (pending
\EGA V):

\begin{enonce*}{Theorem 2.1}\etiquette[2.1]{XI.thm2.1} %Why is the statement
                                %not formulated for a linear subspace
                                %varying around a fixed point?
                                %That is, (i) with $P\in L$ or
                                %$(ii)$ with $d=1$
Let $k$ be a field and let $V\subset \PP^n_k$ be a projective
algebraic scheme purely of dimension $r$. Let $V'$ be the open
subscheme of points of $V$ smooth over $\Spec k$, and let $P$ be a
rational point of $\PP^n$. Then:
\begin{enumerate}
\renewcommand{\theenumi}{\roman{enumi}}
\renewcommand{\labelenumi}{(\theenumi)}
\item A ``sufficiently general'' linear subspace $L$ of a given
  dimension $d$ in $\PP^n$ meets $V$ transversely at every point of
  $L\cap V'$.

\item Let $\{H_1,\ldots,H_s\}$ $(s\leq r)$ be a system of %a
                                %'hyperplane' is not linear here
  hypersurfaces of degree $d\geq 2$ passing through $P$, and let
  $Y=H_1\cap\ldots \cap H_s$. If $\{H_1,\ldots,H_s\}$ is sufficiently
  general, then $Y$ meets $V$ transversely at every point of $Y\cap
  V'$.
\end{enumerate}
\end{enonce*}

The phrase ``for a sufficiently general linear subspace'' means that
there exists a dense open subset $U$ of the Grassmannian parametrizing
the linear subspaces of dimension % 'which parametrizes'
$d$, such that the assertion under consideration about $L$ is true
whenever the parameter of $L$ lies in $U$. Thus this does not imply
that such an $L$ (defined over $k$) exists if $k$ is not infinite.
Likewise,
% mention the Veronese embedding?
the hypersurfaces of degree $d$ passing through $P$ are parametrized
\pageoriginale by a projective space $\PP^N$, and the phrase ``for
$\{H_1,\ldots,H_s\}$ sufficiently general'' means that there exists a
dense open subset $U$ of $(\PP^N)^s$ such that the assertion under
consideration is true whenever the parameter of
$(H_1,\ldots,H_s)$ lies in $U$.

Proof of (i): By induction one is reduced immediately to the case
where $L$ is a hyperplane. Let $(V'_{\nu})$ be a finite open covering
of $V'$. If assertion (i) is true when the symbol $V'$ is replaced by
$V'_{\nu}$ for every $\nu$, then it is also true for $V'$, because the
intersection of finitely many dense open subsets is again open and
dense. We may therefore replace the situation by an affine algebraic
scheme, say $V\subset \EE^n_k=\Spec k[x_1,\ldots,x_n]$. The hyperplane
$L$ will be the zero set of a linear equation
$$
\ell(x)=a_0+\sum^n_{i=1}a_ix_i{\quoi}.
$$

Recall that a smooth scheme is locally a complete intersection. We may
therefore replace $V'$ by an open subset of points of $V$ such that
there exist polynomials $f_1(x),\ldots,f_{n-r}(x)$ vanishing on $V$ and
indices $1\leq j_{\nu}\leq n$ $(\nu=1,\ldots,n-r)$ such that the
determinant
$$
\left|\frac{\partial f_i}{\partial x_{j_{\nu}}}\right|\quad 
(i,\nu=1,\ldots,n-r)
$$
is invertible. Say $j_{\nu}=\nu$ $(\nu=1,\ldots,n-r)$. We may then
solve the equations
$$
y_j=\sum^{n-r}_{i=1}\quad c_i\quad \frac{\partial f_i}{\partial x_j}
$$
for\pageoriginale $c_i$ as rational functions of the variables $x_i$,
$y_j$ $(i=1,\ldots, n-r; j=1,\ldots,n)$, and these functions $c_i(x,y)$
are defined at every point of $V'$.

Now to say that $L$ meets $V$ transversely at a point $Q$ of $V'$ is
equivalent to saying that the value of the matrix
$$
\begin{bmatrix}
\frac{\partial f_1}{\partial x_1} & \cdots & \frac{\partial
  f_1}{\partial x_n}\\
\vdots & &\vdots\\
\frac{\partial f_{n-r}}{\partial x_1} & & \frac{\partial
  f_{n-r}}{\partial x_n}\\
a_1 & \cdots & a_n
\end{bmatrix}
$$
at $Q$ has rank $n-r+1$. Thus, for $L$ to meet $V$ transversely at every
point of $V'$, it is necessary and sufficient that the system of
$r+1$ equations
\begin{align*}
a_j &= \sum^{n-r}_{i=1} c_i (x,a)\frac{\partial f_i}{\partial x_j}\qquad
(j=n-r+1,\ldots, n)\\
a_0 &+ \sum^n_{i=1} a_i x_i =0
\end{align*}
have no solution on $V'$. Such a condition is plainly constructible in
$(a)$: Let $X=V'\times A$, where $A$ is the affine space with
coordinates $a_0,\ldots,a_n$, and let $Y\subset X$ be the closed subset
of solutions of these equations. The condition is that the fiber of
$Y/A$ be empty, which is indeed a constructible\pageoriginale
condition (\EGA IV 9.5.1). It is therefore enough % \ref
to verify that the fiber above the generic point of $A$ is empty,
which follows immediately from the form of the equations. This
completes the proof of (i). (ii): Consider the rational map
$\PP^n\rightarrow \PP^N$ given by the linear system of hypersurfaces
of degree $d$ ($d\geq 2$) passing through $P$. It is well known (and
easily verified) that this is a ``blowing up'' of $P$, and therefore
gives a locally closed immersion
$$
f:\PP^n-\{P\}\longrightarrow \PP^N{\quoi}.
$$
Let $V''=V'-(\{P\}\cap V')$, let
$\bar{V}{}''=\sisi{f}{F}(V'')$, and let $L_i$ be the hyperplane of
$\PP^N$ corresponding to the hypersurface $H_i$
$(i=1,\ldots,s)$. Then $V''$ is isomorphic to $\bar{V}{}''$, and by
the definition of $f$, to say that $Y$ meets $V''$ transversely is
equivalent to saying that $L_1\cap\ldots \cap L_s$ meets
$\bar{V}{}''$ transversely. Putting $\bar{V}=$ the closure of
$\bar{V}{}''$, we are reduced to assertion (i) for $\bar{V}$, which
implies (ii) with $V''$ in place of $V'$. It remains to consider the
point $P$, if it is a smooth point of $V$. But it is evident that if
$\{H_1,\ldots,H_s\}$ is sufficiently general and $P$ is smooth on
$V$, then $Y$ meets $V$ transversely at $P$, whence the result.

% END INCLUDED SGA4_English_current_buildable_source_20260906/source/11/drafts/SGA4_XI_LINES_43_199_English.texfrag


% BEGIN INCLUDED SGA4_English_current_buildable_source_20260906/source/11/drafts/SGA4_XI_LINES_200_353_English.texfrag
\section{Construction of good neighborhoods}\etiquette[3]{XI.sec3}

\begin{enonce*}{Definition 3.1}\etiquette[3.1]{XI.def3.1}
An elementary fibration is a morphism of schemes
${f:X\rightarrow S}$ which can be embedded in a commutative diagram
$$
\xymatrix@=2cm{
X  \ar[r]^{j} \ar[rd]_{f}& \bar{X} \ar[d]^{\bar{f}}& Y \ar[l]_{i}  \ar[ld]^{g}\\
               &    S   &      
}
$$
satisfying\pageoriginale the following conditions:
\begin{enumerate}
\renewcommand{\theenumi}{\roman{enumi}}
\renewcommand{\labelenumi}{(\theenumi)} 
% \item $i$ is a closed immersion.
\item $j$ is an open immersion, dense in every fiber, and
  $X=\bar{X}-Y$. 

\item $\bar{f}$ is smooth and projective, with geometrically
  irreducible fibers of dimension $1$.

\item $g$ is an étale covering, and every fiber of $g$ is nonempty.
\end{enumerate}
\end{enonce*}


One easily checks that such an embedding of $X$, if it exists, is
unique up to canonical isomorphism.

\begin{enonce*}{Definition 3.2}\etiquette[3.2]{XI.def3.2}
A good neighborhood relative to $S$ is an $S$-scheme $X$ for which
there exist $S$-schemes
$$
X=X_n,\ldots,X_0=S
$$
and elementary fibrations $f_i:X_i\rightarrow X_{i-1}$,
% $f_i$ is an $S$-morphism
$i=1,\ldots, n$.
\end{enonce*}

\begin{enonce*}{Proposition 3.3}\etiquette[3.3]{XI.prop3.3}
\footnote{This result extends without difficulty to a semilocal
    assertion.}
Let $k$ be an algebraically closed field, let $X/\Spec k$ be a smooth
scheme, and let $x\in X$ be a rational point. There exists an open
subset of $X$ containing $x$ which is a good neighborhood (relative to
$\Spec k$).
\end{enonce*}

Proof. Take $X$ irreducible. By induction on $\dim X=n$, it suffices to
find a neighborhood $U$ of $x$ in $X$ and an elementary fibration
$f:U\rightarrow V$, with $V$ smooth and of dimension $n-1$. Indeed,
there will exist a neighborhood $V'$ of $v=f(x)$ which is a good
neighborhood, and we may take
$U'=U\cap \sisi{f^{-1}(V')}{V'}$ as a good neighborhood of $x$.

We may suppose that $X\subset \EE^r$ is affine. Let $X_0$ be the
closure of $X$ in $\PP^r$. Let $\bar{X}$ be the normalization of $X_0$
and let $Y=\bar{X}-X$, endowed with the induced reduced structure. In
addition, let\pageoriginale $S\subset \bar{X}$ be the closed subset of
singular points. One has $S\subset Y$ and
\begin{align*}
\dim \bar{X} &= \dim X=n{\quoi},\\
\dim Y &= n-1{\quoi},\\
\dim S &\leq n-2{\quoi}.
\end{align*}

Embed $\bar{X}$ in a projective space $\PP^N$ by means of an invertible
sheaf $L$ of the form $M^{\otimes r}$, with $M$ very ample and
$r\geq 2$. There then exist hyperplanes % linear
$H_1,\ldots, H_{n-1}$ of $\PP^N$, where $H_i$ is the zero set of
$$
\sum^N_{\nu =0}a_{i\nu} x_{\nu} =0{\quoi},
$$
which contain $x$ and such that the intersection
$L=H_1\cap\ldots \cap H_{n-1}$ has dimension $N-n+1$ and meets
$\bar{X}$ and $Y$ transversely (\hyperref[XI.thm2.1]{2.1}).
% see the comment on Theorem 2.1
 %For the intersection with $\bar{X}$: $S\cap L=∅$ for sufficiently
 %general $L$, hence $\bar{X}\cap L=(\bar{X} - Y)\cap L$ and
 %$\bar{X}- Y$ is smooth.
 %For the intersection with $Y$: $Y$ is reduced and $k$ algebraically closed,
 %hence every irreducible component is generically smooth, and sufficiently
 %general $L$ does not meet the singular locus of $Y$.
The intersection $\bar{X}\cap L$ is a smooth connected curve (this
follows from Bertini's theorem), and $Y\cap L$ has dimension $0$.
Choose in addition another hyperplane
$$
H_0:\sum^N_{\nu=0}a_{0 \nu}x_{\nu}=0
$$
such that $H_0$ meets $\sisi{\bar{X}}{X}\cap L$ transversely and
$H_0\cap Y\cap L=∅$. % One must also choose $x\notin H_0$ so that
                    % the map $f'$ below is well defined at $x$

Consider the projection $\PP^N\rightarrow \PP^{n-1}$ obtained by
``introducing the projective coordinates''
$$
y_i=\sum^N_{\nu =0}a_{i_{\nu}}x_{\nu}{\quoi}.
$$
It is a rational map defined outside the center of projection
$C=H_0\cap\ldots \cap H_{n-1}$. Let $\epsilon :P'\rightarrow \PP^N$ be
the blowup of $C$, so that there is a diagram
$$
\xymatrix@=1.5cm{
\PP^{N}  \ar[rd]& &P'\ar[ld]^{\pi}\ar[ll]_{\epsilon}\\
               &\PP^{n-1} &
}
$$
where\pageoriginale $\pi$ is a morphism. Let $\bar{X}{}'\subset P'$ be
the ``proper'' inverse image of $\bar{X}$, that is, the closure of
$\epsilon^{-1}(\bar{X}-(\bar{X}\cap C))$. Since, by hypothesis, $C$
meets $\bar{X}$ transversely, the morphism
$\bar{X}{}'\rightarrow \bar{X}$ is a blowup of the finite set
$\bar{X}\cap C$. Let $X'=X-\bar{X}\cap C$, which is also identified
with an open subscheme of $\bar{X}{}'$, and let
$Y'=\bar{X}{}'-X'$ be the closed subscheme with the induced reduced
structure. There is a diagram $\DD$ of morphisms
$$
\xymatrix@=2cm{
  X'  \ar[r]^{j} \ar[rd]_{f'}& \bar{X}{}'  \ar[d]^{\bar{f}{}'} & Y'\ar[l]_{i} \ar[ld]^{g'}     &\\
                           & \PP^{n-1} &\qquad ,   
}
$$
and I claim that there exists a neighborhood $V$ of $v=f'(x)$ such
that the restriction of $\DD$ to $V$ satisfies conditions (i), (ii),
(iii) of (\hyperref[XI.def3.1]{3.1}), hence that $f'_{\mid V}$ is an
elementary fibration. This will complete the proof.

Condition (i) is trivial. For (ii), note that $\bar{X}\cap L$ is a
smooth curve by hypothesis, and one sees immediately that there is a
one-to-one morphism ${\bar{f}{}'}^{-1}(v)\rightarrow \bar{X}\cap L$
% one-to-one morphism?
induced by $\epsilon$. Hence
$\left({\bar{f}{}'}^{-1}(v)\right)_{\red}\rightarrow
\bar{X}\cap L$ is bijective. To check that $\bar{f}{}'$ is smooth
% remove the \left and \right?
over a neighborhood of $v$, it suffices, by Hironaka's lemma
(\SGA 1 II 2.6 %\ref
\footnote{and \EGA IV 5.12.10.}), to check that it is smooth at the
generic point of ${\bar{f}{}'}^{-1}(v)$. At this point, $\bar{X}{}'$ is
isomorphic to $\bar{X}$, and the morphism is smooth because $L$ meets
$\bar{X}$ transversely.

It remains to prove that $g'$ is étale in a neighborhood of $v$ (since
$Y$ has dimension $n-1$, it is evident that every fiber of $g'$ is
nonempty). One has $Y'=\epsilon^{-1}(Y)\amalg D_1
\amalg\ldots \amalg D_r$, where $\bar{X}\cap C=\{P_1,\ldots,P_r\}$ and
$D_i$ is the blowup\pageoriginale of $P_i$ in $\bar{X}$. Every $D_i$
is identified with $\epsilon^{-1}(P_i)$, and maps isomorphically onto
$\PP^{n-1}$\sisi{. Moreover, the projection}{} is defined at every
point of $Y$%since $Y\cap C=∅$
, and the induced morphism on $Y$ is none other than
$g'_{\mid \epsilon^{-1}(Y)}$. It is étale over $v$, because $L$ meets
$Y$ transversely, hence $g'$ is étale over~$v$.

% END INCLUDED SGA4_English_current_buildable_source_20260906/source/11/drafts/SGA4_XI_LINES_200_353_English.texfrag


% BEGIN INCLUDED SGA4_English_current_buildable_source_20260906/source/11/drafts/SGA4_XI_LINES_354_643_English.texfrag
\section{The comparison theorem}\etiquette[4]{XI.sec4}

\setcounter{subsection}{-1}
\subsection{}\etiquette[4.0]{XI.subsec4.0}
Let $X$ be a scheme locally of finite type over $\Spec \CC$, and
consider the usual topology on the space $X(\CC)$ of rational points
of $X$. We shall denote by $X_{\cl}$ the site of \emph{local
isomorphisms} $f:U\rightarrow X(\CC)$, that is, of morphisms $f$ of
topological spaces having the following property:

\sisi{for}{For} every $x\in U$ there exists an open neighborhood $U_x$
such that $f_{\mid U_x}$ is a homeomorphism of $U_x$ onto an open
neighborhood of $f(x)$.

A family $\{U_{\nu}\rightarrow U\}$ of morphisms of $X_{\cl}$ is said
to be \emph{covering} if $U$ is the union of the images of the
$U_{\nu}$.

Since an open immersion is a local isomorphism, we have a morphism
(inclusion) of sites (IV\quad ) % a reference is missing
\footnote{More precisely, the inclusion of categories
  $\Ouv(X(\CC))\rightarrow X_{\cl}$ defines a morphism of
  sites \hyperref[XI.eq4.1]{(4.1)} in the opposite direction.}
\begin{equation*}
\delta:X_{\cl}\longrightarrow X(\CC)\qquad
(\ast){\quoi}.\tag{4.1}
\end{equation*}\etiquette[4.1]{XI.eq4.1}
Now for every local isomorphism $U\rightarrow X(\CC)$ there exists, by
definition, a ``covering'' of $U$ by open subsets of $X(\CC)$, and one
deduces from (\quad) that the topoi % a reference is missing
associated with the two sites are equivalent (by $\delta_{\ast}$). In
particular, one may replace $X(\CC)$ by $X_{\cl}$ when computing
ordinary cohomology.

Now let $f:X'\rightarrow X$ be an étale morphism of schemes. Then
$f(\CC):X'(\CC)\rightarrow X(\CC)$ is a local isomorphism. This follows
immediately from the\pageoriginale Jacobian criterion
(\SGA 1 II 4.10) and the implicit function theorem. The functor
$X'\mapsto X'(\CC)$ therefore gives a morphism of sites
\begin{equation*}
\epsilon:X_{\cl}\longrightarrow X_{\et}{\quoi},\tag{4.2}
\end{equation*}\etiquette[4.2]{XI.eq4.2}
where $X_{\et}$ is the étale site.

Recall the following results\footnote{Which appear in the Exposé cited
  in the footnote on p.~2.}: % identify the Exposés
                                % above to identify the Exposé
                                % mentioned in the note?

\begin{enonce*}{Theorem 4.3}\etiquette[4.3]{XI.thm4.3}
\begin{enumerate}
\renewcommand{\theenumi}{\roman{enumi}}
\renewcommand{\labelenumi}{(\theenumi)}
\item $X$ is connected and nonempty if and only if $X(\CC)$ is
  connected and nonempty.

\item The functor $\epsilon$ is fully faithful.

\item (``Grauert--Remmert theorem''). The functor $\epsilon$ induces an
  equivalence between the category of finite coverings of $X(\CC)$,
  endowed with its usual topology as a locally compact space, and the
  category of étale coverings of $X$.
\end{enumerate}
\end{enonce*}

We shall accept (i) as known. It follows, for example, from (GAGA)
%\ref
\cite{sga4.exp.XI.bib.4}. Assertion (ii) follows easily from (i). Indeed, to verify the
surjectivity assertion contained in (ii), let
$\varphi:X'(\CC)\rightarrow X''(\CC)$ be an $X(\CC)$-morphism. We may
suppose that all the schemes are connected and affine, and hence
separated. Then the graph $\Gamma$ of $\varphi$ is a connected
component of $\fprod{X'}{X''}{X}(\CC)$, whence
$\Gamma=Y(\CC)$ for some connected component $Y$ of
$\fprod{X'}{X''}{X}$, by (i). Then $Y$ is the graph of a morphism
$f:X'\rightarrow X''$ inducing $\varphi$, \ie the projection
$Y\rightarrow X'$ is an isomorphism: indeed,
$Y(\CC)\rightarrow X(\CC)$ is an isomorphism, from which it follows
easily that the same is true of $Y\rightarrow X'$ (which is étale,
radicial, and surjective!). For (iii), in view of (ii), everything
amounts to proving that every finite covering of $X(\CC)$ is of the
form $X'(\CC)$, where $X'\rightarrow X$\pageoriginale is étale.

For this, descent (\SGA 1 IX 4.7) immediately reduces us to
%\ref
the case in which $X$ is connected and normal. The problem is local on
$X$ (for the Zariski topology), and we may therefore suppose $X$
affine. Let $f:X\rightarrow \EE^n$ be a finite morphism, and let
$T'\rightarrow X(\CC)$ be a connected finite covering. One checks
immediately that the composite morphism
$\eta':T'\rightarrow \EE^n(\CC)$ is an ``analytic covering'' in the
sense of (\cite{sga4.exp.XI.bib.2} \S \hyperref[XI.sec2]{2}, Def.~\hyperref[XI.sec3]{3}).
By (\cite{sga4.exp.XI.bib.2} \S \hyperref[XI.sec2]{2}, Satz 8),
\sisi{$\eta'$}{} extends to an ``analytic covering''
$\eta:T\rightarrow \PP^n(\CC)$, and it is enough to prove that $\eta$
is \sisi{induced}{induced} by a suitable morphism of schemes
$f:Y\rightarrow \PP^n$, with $Y$ normal. Now the fundamental theorem
(\cite{sga4.exp.XI.bib.2} \S 13 Satz 42) asserts that there is a canonical normal
analytic structure on $T$ such that $\eta$ induces a morphism of
analytic spaces. Since $\eta$ is then a finite morphism, the direct
image of the structure sheaf on $T$ is a sheaf $\mathfrak{a}$ of
coherent analytic algebras on $\PP^n$. It follows from
\sisi{(GAGA)}{GAGA} \cite{sga4.exp.XI.bib.4} % \ref
that $\mathfrak{a}$ is the analytic sheaf associated with an algebraic
sheaf $A$ on $\PP^n$, and we take $Y=\sheaf{\Spec}A$.

We can now prove the following result:

\begin{enonce*}{Theorem 4.4}\etiquette[4.4]{XI.thm4.4}
Let $X$ be a smooth scheme over $\Spec \CC$.
\begin{enumerate}
\renewcommand{\theenumi}{\roman{enumi}}
\renewcommand{\labelenumi}{(\theenumi)}
\item There is an equivalence between the category of constructible
  locally constant torsion sheaves on $X_{\et}$ and the category of
  locally constant torsion sheaves with finite fibers on $X_{\cl}$,
  the equivalence being given by the quasi-inverse functors
  ${}_{\ast}$ and ${}^{\ast}$.

\item Let $F$ be a locally constant torsion sheaf with finite fibers on
  $X_{\cl}$, and also denote by $F$ the induced sheaf
  $\epsilon_{\ast}F$ on $X_{\et}$. Then
$$
R^q\epsilon_{\ast}F=0\quad \mbox{ if } q>0{\quoi}.
$$

\item With\pageoriginale the notation of {\rm(i)}, the canonical
  morphism
$$
H^q(X_{\et},F)\longrightarrow H^q(X_{\cl},F)
$$
is bijective for every $q$. In particular,
$$
H^q(X_{\et}, \ZZ/\sisi{n\ZZ}{n})\xrightarrow{\sim} H^q(X_{\cl},\ZZ/\sisi{n\ZZ}{n})
$$
for every $q$.
\end{enumerate}
\end{enonce*}

Proof. For (i), note that the sheaves in question are those that are
representable by étale coverings of $X$ (\resp by finite coverings of
$X(\CC)$). Since there is an equivalence between the categories of
these coverings (\hyperref[XI.thm4.3]{4.3}), assertion (i) follows
immediately. The isomorphism of (iii) follows from (ii) and from the
Leray spectral sequence V \hyperref[V.prop5.2]{5.2} % \ref + 5.3 instead of 5.2?
for $\epsilon$. It remains to prove (ii).

Now $R^q\epsilon_{\ast}F$ is the sheaf associated with the presheaf
$\Rr^qF,$, where $(\Rr^qF)(X')=H^q(X'_{\cl},F)$ for
$X'\rightarrow X$ étale. To prove that $R^q\epsilon_{\ast}F=0$ if
$q>0$, we must therefore prove the following \sisi{lemma}{}:

\begin{enonce*}{Lemma 4.5}\etiquette[4.5]{XI.lem4.5}
Let $\xi\in H^q(X_{\cl},F)$ and $x\in X_{\cl}=X(\CC)$. There exists an
étale morphism $X'\rightarrow X$ whose image contains $x$ and such that
the image of $\xi$ in $H^q(X'_{\cl},F)$ is zero.
\end{enonce*}

Proof of the lemma. Induction on $n=\dim X$; the problem is local on
$X$ for the étale topology, \sisi{}{and} we may therefore suppose $F$
constant, and by dévissage we may even suppose that
$F=\ZZ/\sisi{n\ZZ}{n}$. Moreover, replacing $X$ by a Zariski open
neighborhood of $x$, we may suppose (\hyperref[XI.prop3.3]{3.3}) that
$X$ admits an elementary fibration (\hyperref[XI.def3.1]{3.1})
$f:X\rightarrow S$. Retain the notation of
(\hyperref[XI.def3.1]{3.1}): by\pageoriginale direct calculation, one
sees that $R^q{j}_{\ast}F=0$ if $q>1$, that
$R^1j_{\cl\!\ast}F$ is the extension by $0$ of a constant sheaf on
$Y_{\cl}$, and finally that $j_{\cl\!\ast}F$ is the constant sheaf
$\ZZ/\sisi{n\ZZ}{n}$ on $\bar{X}$. Moreover, one may compute
$R^q \bar{f}_{\cl\!\ast}(j_{\cl\!\ast}F)$ fiber by fiber. From the
Leray spectral sequence
$$
E^{pq}_2=R^p\bar{f}_{\cl\!\ast}(R^qj_{\cl\!\ast}F)\Longrightarrow
R^{p+q}f_{\cl\!\ast}F
$$
one deduces that $R^{p+q}f_{\cl\!\ast}F$ may likewise be computed fiber
by fiber, and hence that

$R^0f_{\cl\!\ast}\ZZ/\sisi{n\ZZ}{n}=\ZZ/\sisi{n\ZZ}{n}$,

$R^1f_{\cl\!\ast}\ZZ/\sisi{n\ZZ}{n}$ is a locally constant torsion
  sheaf on $S_{\cl}$,

$R^qf_{\cl\!\ast}\ZZ/\sisi{n\ZZ}{n}=0$ if $q>1$.

The Leray spectral sequence

$E^{pq}_2=H^p(S_{\cl}, R^qf_{\cl\!\ast}F)\Longrightarrow
H^{p+q}(X_{\cl},F)$

\noindent
thus reduces to the exact sequence
\begin{equation*}
\ldots\longrightarrow H^q(S_{\cl},f_{\cl\!\ast}F)\longrightarrow
H^q(X_{\cl},F)\longrightarrow H^{q-1}(S_{\cl},
R^1f_{\cl\!\ast}F)\longrightarrow\tag{$\ast$}
\end{equation*}
Let $s$ be the image of $x$ in $S$. By the induction hypothesis and by
(i), for every class $\eta\in H^q(S_{\cl},G)$ ($G$ locally constant
torsion, with finite fibers), there exists an ``étale neighborhood''
$S'\rightarrow S$ of $s$ such that the image of $\eta$ in
$H^q(S'_{\cl},G)$ is zero. Now the functors $R^qf_{\cl\!\ast}$ plainly
commute with étale base changes, and if $S'\rightarrow S$ is an étale
neighborhood of $s$, then
$X'=\fprod{X}{S'}{S}\rightarrow X$ is an étale neighborhood of $x$.
When $q\geq 1$, we therefore conclude by means of the exact sequence
$(\ast)$, first taking $G=R^1f_{\cl\!\ast}F$ and
then\pageoriginale $G=f_{\cl\!\ast}F$. % The proof ends here; the
                                % following is a remark for the case $q=1$
When $q=1$, note that $H^1(X_{\cl},\ZZ/n \ZZ)$ (\resp
$H^1(X,\ZZ/n\ZZ)$) classifies the principal étale coverings
% should one instead speak of finite coverings for $X_\cl$?
of $X_{\cl}$ (\resp $X$), with group $\ZZ/n\ZZ$; hence, by
\hyperref[XI.thm4.3]{4.3} (iii), the map
$H^1(X,\ZZ/n\ZZ)\rightarrow H^1(X_{\cl},\ZZ/n\ZZ)$ is bijective,
from which the effacement result \hyperref[XI.lem4.5]{4.5} for $q=1$
follows at once.

\begin{enonce*}[remark]{Variant 4.6}\etiquette[4.6]{XI.var4.6}
One may also prove (\hyperref[XI.lem4.5]{4.5}) in the following elegant
way, due to Serre: one notes that a connected good neighborhood
(\hyperref[XI.def3.2]{3.2}) $X$ satisfies the following two conditions:
\begin{enumerate}
\renewcommand{\theenumi}{\roman{enumi}}
\renewcommand{\labelenumi}{(\theenumi)}
\item $\pi_n(X(\CC))=0$ if $n>1$.

\item $\pi_1(X(\CC))$ is an iterated extension of free groups of finite
type.
\end{enumerate}
\end{enonce*}

Indeed, if $X\rightarrow S$ is an elementary fibration where $S$ is a
good neighborhood, we may suppose (i) and (ii) true for $S$ by
induction. Now $X(\CC)\rightarrow S(\CC)$ is a locally trivial fiber
space, as is readily seen, with fibers isomorphic to $X_0(\CC)$, where
$X_0$ is an arbitrary fiber of $X/S$. Now $X_0$ is a connected
nonsingular \emph{noncomplete} curve, and it is well known that this
implies that $\pi_n(X_0(\CC))=0$ if $n>1$, and that
$\pi_1(X_0(\CC))$ is a free group. % \ref
Then (i) and (ii) follow from the homotopy exact sequence
$$
\longrightarrow \pi_0(X_0(\CC))\longrightarrow
\pi_n (X(\CC))\longrightarrow \pi_n
(S(\CC))\longrightarrow\ldots
$$

Thus $X(\CC)$ is a $K(\pi,1)$ space, and this implies that every
cohomology class $\xi\in H^n(X(\CC),F)$ becomes zero on the universal
covering of $X(\CC)$. It is enough to prove that $\xi$ already becomes
zero on a finite covering $X'$ of $X$, which amounts to proving that
$\pi_1=\pi_1(X(\CC))$ is\pageoriginale a \emph{good group}
(\cf CG 1 16), %\ref
that is, that the morphism from $\pi_1$ to its completion
$\sisi{\widehat}{}{\pi}_1=\varprojlim_N \pi_1/N$ (where $N$ runs
through the set of invariant subgroups % 'invariant' = 'normal'?
of $\pi_1$ of finite index) induces a bijection
$$
H^q(\widehat{\pi}_1,M)\xrightarrow{\sim} H^q(\pi_1,M)
$$
for every continuous finite $\pi_1$-module $M$ (where the symbol on the
left denotes the cohomology of $\widehat{\pi}_1$ as a profinite group).
Now an iterated extension of free groups of finite type is a good group
(\cf CG I p.~15-16, exc.~1,2), whence the result. % \ref

\begin{thebibliography}{99}
\bibitem{sga4.exp.XI.bib.1} GIRAUD J., ``Analysis Situs'', Séminaire
Bourbaki, Exposé nᵒ~256.

\bibitem{sga4.exp.XI.bib.2} GRAUERT H\sisi{.}{} and REMMERT R., ``Komplexe Räume'',
Math, Ann. Bd. 136, p. 245 (1958).

\bibitem{sga4.exp.XI.bib.3} HIRONAKA H., ``Resolution of singularities of an
algebraic variety over a field of characteristic zero'' Annals of
Math. Vol. 39, nᵒ~1, 2 (1964).

\bibitem{sga4.exp.XI.bib.4} SERRE J.P., ``Géométrie Algébrique et Géométrie
Analytique'', Annales de l'Institut Fourier, t. VI, p. 1-12 (1956).
\end{thebibliography}


\ifx\danslelivre\undefined
\end{document}
\fi

% END INCLUDED SGA4_English_current_buildable_source_20260906/source/11/drafts/SGA4_XI_LINES_354_643_English.texfrag


% END INCLUDED SGA4_English_current_buildable_source_20260906/source/11/11_English.tex


\makeatletter
\def\input@path{{../12/}}
\makeatother
% BEGIN INCLUDED SGA4_English_current_buildable_source_20260906/source/12/12_English.tex
\ifx\danslelivre\undefined
% Release engine: XeLaTeX (the LuaTeX compatibility layer breaks Xy-pic output).
\documentclass[10pt]{../fichiers-configuration-latex/smfart}
% BEGIN INCLUDED SGA4_English_current_buildable_source_20260906/source/fichiers-configuration-latex/paquets.tex
\usepackage{polyglossia}
%\usepackage[francais]{babel}
\usepackage{fontspec}
\usepackage[math-style=ISO,bold-style=upright,sans-style=upright]{unicode-math}
%\usepackage[mathletters]{ucs}
%\usepackage[utf8]{inputenc}
%\usepackage[T1]{fontenc}
\usepackage[unicode,hypertexnames=false,colorlinks=true,linkcolor=blue,citecolor=Sepia,urlcolor=magenta]{hyperref}
\setmainlanguage{french}
\setmainfont[Ligatures=TeX]{Linux Libertine O}
% variante possible : TeXGyrePagella ?
\setmathfont{XITS Math}

% maths
\usepackage{amsmath}
%\usepackage{amssymb}
%\usepackage{mathrsfs}
%\usepackage{stmaryrd}
%\usepackage{euscript}

% couleurs
\usepackage[usenames,dvipsnames]{xcolor}

% graphiques
\usepackage{graphicx}
\usepackage[all]{xy}
%\usepackage{tikz}
%\usetikzlibrary{matrix,arrows,calc}
%\tikzset{column sep=3em,row sep=2em,text height=1.5ex, text depth=0.25ex,>=angle 90}
%\tikzset{description/.style={fill=white,inner sep=1pt}}

% hyperliens, 'etiquettes
\usepackage{xr-hyper}
\usepackage{../fichiers-configuration-latex/smfthm}
%\ifx\cacheretiquettes\undefined
%\usepackage[inline]{showlabels} \else \fi % pour rendre visibles les 'etiquettes
%\usepackage{../fichiers-configuration-latex/makeidx} % pour faire un index
%\usepackage[pagebackref,unicode]{hyperref}  % pour voir o`u sont cit'ees les r'ef'erences
% hyperref is loaded once above with the cumulative-reader destination policy.
%\hypersetup{breaklinks=true} %permet le retour `a la ligne dans les liens trop longs. Incompatible avec pdflatex
%\renewcommand*{\backref}[1]{{\small [cit\'e page(s)~#1.]}}

\usepackage{xspace}
\usepackage{pifont} % pour expose XVII

% END INCLUDED SGA4_English_current_buildable_source_20260906/source/fichiers-configuration-latex/paquets.tex

% BEGIN INCLUDED SGA4_English_current_buildable_source_20260906/source/fichiers-configuration-latex/numerotation.tex
\renewcommand{\thefootnote}{\roman{footnote}} % pour mettre les notes en bas de page sous la forme (ii) et non (2), qui ressemble trop (?) a un exposant.

\newcounter{pageoriginal}
\newcommand{\pageoriginale}{\refstepcounter{pageoriginal}\marginpar{\footnotesize\textbf{\thepageoriginal}}\xspace}
\let\pageoriginaled\pageoriginale

\usepackage{ifthen}
\newboolean{original}
\def\sisi{\ifthenelse{\boolean{original}}}
\setboolean{original}{false}

\newcounter{cnote}
\def\nde#1{\ifthenelse{\boolean{original}}{}{\stepcounter{cnote}\def\thefootnote{\arabic{cnote}}\protect\footnote[\value{cnote}]{N.D.E.: #1}}}
\def\ndemark{\ifthenelse{\boolean{original}}{}{\stepcounter{cnote}\def\thefootnote{\arabic{cnote}}\protect\footnotemark[\value{cnote}]}}
\def\ndetext#1{\ifthenelse{\boolean{original}}{}{\def\thefootnote{\arabic{cnote}}\protect\footnotetext[\value{cnote}]{N.D.E.: #1}}}

% END INCLUDED SGA4_English_current_buildable_source_20260906/source/fichiers-configuration-latex/numerotation.tex

% BEGIN INCLUDED SGA4_English_current_buildable_source_20260906/source/fichiers-configuration-latex/macros.tex
\setmainlanguage{french}
\setotherlanguage{english}


%% Fixed, engine-independent release identifier.  The inherited dynamic Git
%% probe used Lua io.popen, exposed build-state-dependent output, and printed
%% a visible question mark outside a Git checkout.
\def\showgitstatus{English release r2, 2026-07-23}


% macros pour LaTeX
\renewcommand{\mathsf}{}

% Police différente pour U+2113 SCRIPT SMALL L
\setmathfont[range="2113]{latinmodern-math.otf}

% Give every inherited semantic label its own stable PDF destination instead
% of inheriting whichever automatic equation/page anchor happened to be live.
% The legacy label remains as a compatibility alias used by existing edges.
\makeatletter
\def\etiquette[#1]#2{%
  \Hy@raisedlink{\hyper@anchorstart{sga4.semantic.#2}\hyper@anchorend}%
  \def\@currentHref{sga4.semantic.#2}%
  \def\@currentlabel{#1}%
  \label{#2}%
}
\makeatother
\def\eqhyperref[#1]#2{(\hyperref[#1]{#2})}
\def\pagehyperref[#1]#2{(page \hyperref[#1]{#2})}
\def\hyperrefext[#1][#2]#3{\hyperref[#2]{#1 #3}}
\def\eqhyperrefext[#1][#2]#3{\hyperref[#2]{#1 #3}}     %cf. exposes 17
\def\neweqhyperrefext[#1][#2]#3{\hyperref[#2]{#1 #3}}  %cf. exposes 17

\def\shat#1{\protect\widehat{#1}}
\def\stilde#1{\protect\widetilde{#1}}
%%%%%%%%%%%%%%%%%%


\def\etcn{etc.\xspace}
\def\etco{etc...\xspace}
\def\etc{\sisi{\etco}{\etcn}}
\def\ldotsn{\sisi{...}{\etcn}}

\def\cf{{cf.}\kern.3em}
\def\Cf{{Cf.}\kern.3em}
\def\ie{{i.e.}\kern.3em}
\def\iev{{i.e.},\xspace}
\def\loccito{loc.\kern2pt cit.\xspace}
\def\loccitn{\textit{loc.\kern2pt cit.}\xspace}
\def\loccit{\sisi{\loccito}{\loccitn}}
\def\EGA{\textup{EGA}\kern.2em}
\def\SGA{\textup{SGA}\kern.3em}

\def\quoi{}

\def\quav{\quad,\quad}
\def\quaet{\quad \text{et}\quad}
\def\fibreproduct#1#2#3{#1 ×_{#3} #2}
\let\fprod\fibreproduct
\def\ccheck#1#2{\mbox{${\check{\check{\rm #1}}}^{\lower2pt\hbox{$\scriptstyle #2$}}$}}
\def\sheaf#1{\protect\mathcal{#1}}
\def\vieuxsheaf#1{{#1}} % cf. problème avec sheaf{a\Phi_i} par exemple...
\DeclareMathOperator{\Hom}{Hom}
\DeclareMathOperator{\can}{can}
\DeclareMathOperator{\Ind}{Ind}
\DeclareMathOperator{\ind}{ind}
%\DeclareMathOperator{\Hhomtop}{\mathcal{H}\!\mathit{omtop}}
\DeclareMathOperator{\Hhomtop}{\mathbf{Homtop}}
\DeclareMathOperator{\Point}{Point}
\DeclareMathOperator{\Pro}{Pro}
\DeclareMathOperator{\Top}{Top}
\DeclareMathOperator{\Mod}{Mod}
\DeclareMathOperator{\Cart}{Cart}
%\DeclareMathOperator{\Mmorsite}{\mathcal{M}\!\mathit{orsite}}
\DeclareMathOperator{\Mmorsite}{\mathbf{Morsite}}
\DeclareMathOperator{\loc}{loc}
\DeclareMathOperator{\Loc}{Loc}
\DeclareMathOperator{\Pic}{Pic}
\DeclareMathOperator{\Et}{Et}
\DeclareMathOperator{\card}{card}
\DeclareMathOperator{\Ens}{Ens}
\DeclareMathOperator{\h}{\underline{H}}
\DeclareMathOperator{\Fonct}{Fonct}
\DeclareMathOperator{\Ker}{Ker}
\DeclareMathOperator{\but}{but}
\DeclareMathOperator{\Ob}{Ob}
\DeclareMathOperator{\ob}{ob}
\DeclareMathOperator{\Lib}{Lib}
\DeclareMathOperator{\esj}{esj}
\DeclareMathOperator{\id}{id}
\DeclareMathOperator{\Id}{Id}
\DeclareMathOperator{\ad}{ad}
\DeclareMathOperator{\Iim}{Im}
\DeclareMathOperator{\pr}{pr}
\DeclareMathOperator{\Cat}{Cat}
\DeclareMathOperator{\car}{car}
\DeclareMathOperator{\SysInd}{SysInd}
\DeclareMathOperator{\Sup}{Sup}
\DeclareMathOperator{\Coker}{Coker}
\DeclareMathOperator{\Filt}{Filt}
\DeclareMathOperator{\acc}{acc}
\DeclareMathOperator{\sur}{sur}
\DeclareMathOperator{\cat}{cat}
\DeclareMathOperator{\Ord}{Ord}
\DeclareMathOperator{\Gr}{Gr}
\DeclareMathOperator{\Cov}{Cov}
\DeclareMathOperator{\Caf}{Caf}
\DeclareMathOperator{\coker}{coker}
\DeclareMathOperator{\coim}{coim}
\DeclareMathOperator{\im}{im}
\DeclareMathOperator{\Ouv}{Ouv}
\DeclareMathOperator{\diag}{diag}
\DeclareMathOperator{\Esp}{Esp}
\DeclareMathOperator{\TOP}{TOP}
\DeclareMathOperator{\Unif}{Unif}
\DeclareMathOperator{\Iind}{\mathbf{Ind}}
\DeclareMathOperator{\Hhomcart}{\mathbf{Homcart}}
\DeclareMathOperator{\Hhomtopess}{\mathbf{Homtopess}}
\DeclareMathOperator{\Ppoints}{\mathbf{Points}}
\DeclareMathOperator{\Ppoint}{\mathbf{Point}}
\DeclareMathOperator{\Ttors}{\mathbf{Tors}}
\DeclareMathOperator{\Ssex}{\mathbf{Sex}}
\DeclareMathOperator{\Ffib}{\mathbf{Fib}}
\DeclareMathOperator{\Vvois}{\mathbf{Vois}}
\DeclareMathOperator{\Ppointess}{\mathbf{Pointess}}
\DeclareMathOperator{\scrHom}{\mathbf{Hom}}
\DeclareMathOperator{\sheafPic}{\mathbf{Pic}}
\DeclareMathOperator{\sheafAut}{\mathbf{Aut}}
\DeclareMathOperator{\Hhomtoplocan}{\mathbf{Homtoplocan}}
\DeclareMathOperator{\Hhomtopan}{\mathbf{Homtopan}}
\DeclareMathOperator{\Pptgeom}{\mathbf{Ptgéom}}
\DeclareMathOperator{\Morsite}{Morsite}
\DeclareMathOperator{\Fib}{Fib}
\DeclareMathOperator{\Fais}{Fais}
\DeclareMathOperator{\kar}{kar}
\DeclareMathOperator{\Groupes}{Groupes}
\DeclareMathOperator{\Restr}{Restr}
\DeclareMathOperator{\Isom}{Isom}
\DeclareMathOperator{\Prol}{Prol}
\DeclareMathOperator{\Res}{Res}
\DeclareMathOperator{\Tors}{Tors}
\DeclareMathOperator{\esp}{esp}
\DeclareMathOperator{\es}{es}
\DeclareMathOperator{\od}{od}
\DeclareMathOperator{\gr}{gr}
\DeclareMathOperator{\sob}{sob}
\DeclareMathOperator{\dfn}{dfn}
\DeclareMathOperator{\Sch}{Sch}
\DeclareMathOperator{\ouv}{ouv}
\DeclareMathOperator{\cons}{cons}
\DeclareMathOperator{\Nerf}{Nerf}
\DeclareMathOperator{\Points}{Points}
\DeclareMathOperator{\Espsob}{Espsob}
\DeclareMathOperator{\Homtop}{Homtop}
\DeclareMathOperator{\Pt}{Pt}
\DeclareMathOperator{\Quot}{Quot}
\DeclareMathOperator{\glob}{glob}
\DeclareMathOperator{\Homfib}{Homfib}
\DeclareMathOperator{\Qouv}{Qouv}
\DeclareMathOperator{\Quotouv}{Quotouv}
\DeclareMathOperator{\ens}{ens}
\DeclareMathOperator{\ab}{ab}
\DeclareMathOperator{\an}{an}
\DeclareMathOperator{\fs}{fs}
\DeclareMathOperator{\red}{red}
\DeclareMathOperator{\Ab}{Ab}
\DeclareMathOperator{\cosq}{cosq}
\DeclareMathOperator{\Ann}{Ann}
\DeclareMathOperator{\Zar}{Zar}


%%%%%%%%%%%%
%%%%% SGA4-2
\DeclareMathOperator{\Sex}{Sex}
\DeclareMathOperator{\vois}{vois}
\DeclareMathOperator{\Vois}{Vois}
\DeclareMathOperator{\Ext}{Ext}
\DeclareMathOperator{\ext}{ext}
\DeclareMathOperator{\Fer}{Fer}
\DeclareMathOperator{\Add}{Add}
\DeclareMathOperator{\Cocart}{Cocart}
\DeclareMathOperator{\cosk}{cosk}
\DeclareMathOperator{\plat}{plat}

\DeclareMathOperator{\faisc}{faisc}
\DeclareMathOperator{\sk}{sk}
\DeclareMathOperator{\ch}{ch}

\def\HH{\mathbf{H}}
\def\qc{\mathrm{qc}}
\def\et{\mathrm{et}}
\def\coh{\mathrm{coh}}

\def\Hvr#1{%
\begin{array}[b]{l}
{\kern2pt\scriptstyle r}\\[-8pt]
{\kern1pt\scriptstyle \vee}\\[-5pt]
H^{#1}
\end{array}}

\def\Hvv#1{%
\begin{array}[b]{l}
\kern2pt\rotatebox{90}{${\scriptscriptstyle\ll}$}\\[-6pt]
H^{#1}
\end{array}}

\def\Hvi#1{%
\begin{array}[b]{l}
{\kern0.2pt\scriptscriptstyle \infty}\\[-8pt]
{\kern1pt\scriptstyle \vee}\\[-5pt]
H^{#1}
\end{array}}

\def\Hv#1{%
\begin{array}[b]{l}
{\kern1pt\scriptstyle \vee}\\[-5pt]
H^{#1}
\end{array}}

\DeclareMathOperator{\cart}{cart}
\DeclareMathOperator{\Ffaisc}{\mathbf{Faisc}}


\def\Ev#1{%
\begin{array}[b]{l}
{\kern1pt\scriptstyle \vee}\\[-5pt]
E'_{#1}
\end{array}}

\DeclareMathOperator{\Cont}{Cont}
\DeclareMathOperator{\Ccont}{\mathscr{C}\!\mathit{ont}}
\DeclareMathOperator{\Site}{Site}
\newcommand\MM{\mathbf{M}}
\def\limtop#1{{\displaystyle{\varprojlim_{#1}\!\mbox{\rm top }}}}
\def\Limtop{\mathop{\varprojlim\!\mbox{\rm top}}}
\DeclareMathOperator{\spec}{spec}
\def\ss{\scriptstyle{\mathscr{S}}}
\DeclareMathOperator{\Gal}{Gal}
\def\smallm{\scriptstyle{\mathscr{M}}}
\DeclareMathOperator{\End}{End}
\DeclareMathOperator{\ass}{ass}
\DeclareMathOperator{\zar}{zar}
\DeclareMathOperator{\Ogg}{Ogg}
\DeclareMathOperator{\res}{res}
\DeclareMathOperator{\tr}{tr}
\DeclareMathOperator{\Supp}{Supp}
\DeclareMathOperator{\codim}{codim}
\DeclareMathOperator{\rk}{rk}
\DeclareMathOperator{\cur}{cur}
\DeclareMathOperator{\cdqc}{cdqc}
\newcommand\DD{\mathbf{D}}
\DeclareMathOperator{\cl}{cl}
\def\JJ{\mathbf{J}}
\def\FF{\mathbf{F}}
\def\BB{\mathbf{B}}
\DeclareMathOperator{\Alget}{Alget}
\DeclareMathOperator{\Ass}{Ass}
\DeclareMathOperator{\supp}{supp}
\DeclareMathOperator{\rad}{rad}
\DeclareMathOperator{\prof}{prof}
\DeclareMathOperator{\cint}{int}
\DeclareMathOperator{\sic}{sic}
\DeclareMathOperator{\Tor}{Tor}
\DeclareMathOperator{\tors}{tors}
\DeclareMathOperator{\Sym}{Sym}
\DeclareMathOperator{\parf}{parf}
\DeclareMathOperator{\tor}{tor}
\DeclareMathOperator{\defn}{defn}
\DeclareMathOperator{\Div}{Div}
\DeclareMathOperator{\Pol}{Pol}
\DeclareMathOperator{\Fex}{Fex}
\DeclareMathOperator{\Lie}{Lie}
\DeclareMathOperator{\tot}{tot}
\DeclareMathOperator{\Aut}{Aut}
\DeclareMathOperator{\HOM}{HOM}
\DeclareMathOperator{\pch}{pch}
\DeclareMathOperator{\Ch}{Ch}
\DeclareMathOperator{\EXT}{EXT}
\DeclareMathOperator{\PIC}{PIC}
\DeclareMathOperator{\Rf}{Rf}
\DeclareMathOperator{\inj}{inj}
\newcommand\VV{\mathbf{V}}

\def\Xv{%
\begin{array}[b]{l}
{\kern1pt\scriptstyle \vee}\\[-5pt]
X
\end{array}}

%%%%%%%
\newcommand\TT{\mathbf{T}}
\newcommand\XX{\mathbf{X}}
\newcommand\EE{\mathbf{E}}
%%%%%%

\def\qedsymbol{\textup{C.Q.F.D.}}
\let\CQFD\qed
\let\cqfd\qed

%%%%%\bigdot
\font\bigsymb=cmsy10 at 4pt
\def\bigdot{{\kern1.2pt\raise 1.5pt\hbox{\bigsymb\char15}}}
%%%%%

\def\Lright{\mathop{\vcenter{\hbox{\hbox{$\longrightarrow$}\lower8pt \hbox{\kern-16pt\hbox{$\longrightarrow$}}}}}}

\def\Aa{{\mathcal A}}
\def\Bb{{\mathcal B}}
\def\Cc{{\mathcal C}}
\def\Dd{{\mathcal D}}
\def\Ee{{\mathcal E}}
\def\Ff{{\mathcal F}}
\def\Gg{{\mathcal G}}
\def\Hh{{\mathcal H}}
\def\Ii{{\mathcal I}}
\def\Jj{{\mathcal J}}
\def\Kk{{\mathcal K}}
\def\Ll{{\mathcal L}}
\def\Mm{{\mathcal M}}
\def\Nn{{\mathcal N}}
\def\Oo{{\mathcal O}}
\def\Pp{{\mathcal P}}
\def\Qq{{\mathcal Q}}
\def\Rr{{\mathcal R}}
\def\Ss{{\mathcal S}}
\def\Tt{{\mathcal T}}
\def\Uu{{\mathcal U}}
\def\Vv{{\mathcal V}}
\def\Ww{{\mathcal W}}
\def\Xx{{\mathcal X}}
\def\Yy{{\mathcal Y}}
\def\Zz{{\mathcal Z}}

\newcommand{\CC}{\mathbf{C}}
\newcommand{\GG}{\mathbf{G}}
\newcommand{\KK}{\mathbf{K}}
\newcommand{\LL}{\mathbf{L}}
\newcommand{\NN}{\mathbf{N}}
\newcommand{\PP}{\mathbf{P}}
\newcommand\QQ{\mathbf{Q}}
\newcommand\RR{\mathbf{R}}
\newcommand\ZZ{\mathbf{Z}}

\def\to{\mathchoice{\longrightarrow}{\rightarrow}{\rightarrow}{\rightarrow}}
\def\from{\mathchoice{\longleftarrow}{\leftarrow}{\leftarrow}{\leftarrow}}
\def\mto{\mathchoice{\longmapsto}{\mapsto}{\mapsto}{\mapsto}}
\def\mfrom{\mathrel{\from\joinrel\mapstochar\,}}
\def\hto{\mathrel{\lhook\joinrel\to}}
\def\hfrom{\mathrel{\from\joinrel\rhook}}

\def\implies{\mathchoice{\Longrightarrow}{\Rightarrow}{\Rightarrow}{\Rightarrow}}
\def\implied{\mathchoice{\Longleftarrow}{\Leftarrow}{\Leftarrow}{\Leftarrow}}
\def\iff{\mathchoice{\Longleftrightarrow}{\Leftrightarrow}{\Leftrightarrow}{\Leftrightarrow}}

\def\To#1{\mathchoice{\xrightarrow{\labelstyle~~#1~}}{\stackrel{#1}{\longrightarrow}}{}{}}
\def\Mto#1{\mathrel{\mapstochar\To{#1}}}
\def\From#1{\mathchoice{\xleftarrow{\labelstyle~~#1~}}{\stackrel{#1}{\longleftarrow}}{}{}}

\newcommand{\isomto}{\overset{\sim}{\longrightarrow}}
\newcommand{\isomfrom}{\overset{\sim}{\longleftarrow}}

\newcommand{\IFF}{\unskip~$\iff$~\ignorespaces}
\newcommand{\THEN}{\unskip~$\implies$~\ignorespaces}

\let\leq\leqslant
\let\geq\geqslant
\let\le\leqslant
\let\ge\geqslant

\def\cf{{cf.}\kern.3em}
\def\Cf{{Cf.}\kern.3em}
\def\ie{{i.e.}\kern.3em}
\def\iev{{i.e.},\xspace}
%\def\loccito{loc.\kern2pt cit.\xspace}
%\def\loccitn{\textit{loc.\kern2pt cit.}\xspace}
%\def\loccit{\sisi{\loccito}{\loccitn}}
\def\EGA{\textup{EGA}\kern.2em}
\def\SGA{\textup{SGA}\kern.3em}

\let\oldS\S
\def\S{\oldS\kern .15em}
\def\ptbl{.\kern .2em}
\def\resp{resp.\kern.3em}
\def\Exp{Exp.\kern.3em}

\let\bar\overline
\let\hat\widehat
\let\tilde\widetilde
\let\dpl\displaystyle
\let\ldots\dots
\let\emptyset\varnothing
\let\skpt\skippointrait
\let\epsilon\varepsilon

\def\sfrac#1#2{{#1}/{#2}}

\DeclareMathOperator{\cd}{cd}
%\DeclareMathOperator{\Hhom}{\mathcal{H}\!\mathit{om}}
\DeclareMathOperator{\Hhom}{\mathbf{Hom}}
%\DeclareMathOperator{\Eext}{\mathcal{E}\!\mathit{xt}}
\DeclareMathOperator{\Eext}{\mathbf{Ext}}
\DeclareMathOperator{\Spec}{Spec}
\DeclareMathOperator{\Tr}{Tr}

\let\emptyset=\varnothing

%%%%%%%%%%%%%%%
%%% expose XVII
%%%%%%%%%%%%%%%

\makeatletter
\let\old@optexpose\optexpose
\newcommand{\empty@optexpose}[1]{}
\newcommand{\paren@optexpose}[1]{\old@optexpose{#1}(}
\newcommand{\refnoexpose}[1]{\let\optexpose\empty@optexpose\ref{#1}\let\optexpose\old@optexpose}%\refnoexpose{X.cor5.2} engendre "5.2" au lieu de "X 5.2"
\newcommand{\neweqref}[1]{\let\optexpose\paren@optexpose\textup{\ref{#1})}\let\optexpose\old@optexpose}%\neweqref{XII.eq4.1} engendre "XII (4.1)" au lieu de "(XII 4.1)"
\makeatother
\newcommand{\refvb}[2][VI${}_{\text{B}}$ ]{\sisi{#1\refnoexpose{#2}}{\ref{#2}}}

\newcommand{\leftexp}[2]{{\vphantom{#2}}^{#1}{#2}}
\newcommand{\AAA}{\mathbf{A}}
\newcommand{\Sf}{\mathfrak{S}}
\DeclareMathOperator*{\quotprojlim}{\text{``$\varprojlim$''}}
\DeclareMathOperator*{\quotinjlim}{\text{``$\varinjlim$''}}
\DeclareMathOperator{\Fl}{F\sisi{\ell}{l}}
\DeclareMathOperator{\Pl}{P\sisi{\ell}{l}}
%\DeclareMathOperator{\cl}{c\sisi{\ell}{l}}
\DeclareMathOperator{\tdf}{tdf}
\DeclareMathOperator{\rg}{rg}
\newcommand{\iInd}{\sisi{I}{i}nd}
\newcommand{\iind}{ind}
\newcommand{\ppro}{pro}
\newcommand{\tTor}{\sisi{T}{t}or}
\newcommand{\ttor}{tor}
\newcommand{\chap}{Chap.}
\newcommand{\Ffex}{\Ff\mathit{ex}}
\newcommand{\sheafk}{\underline{k}}
\newcommand{\sheafn}{\underline{n}}
\newcommand{\sheafGamma}{\underline{\Gamma}}

% Repeated references to an already typeset footnote must not use
% \footnotemark without a matching \footnotetext: Hyperref would then create
% a live annotation to a destination that cannot exist.  This helper keeps the
% current printed footnote mark and links it to an explicit stable note target.
\makeatletter
\newcommand{\sgarepeatedfootnotemark}[1]{%
  \begingroup
  \edef\@thefnmark{\thefootnote}%
  \hyperlink{#1}{\@makefnmark}%
  \endgroup}
\makeatother

\sisi{
\DeclareMathOperator{\simarrow}{\rule{.15mm}{1.8mm}\kern -3.4pt\rightsquigarrow}
\newcommand{\sheafell}{\underline{\ell}}
\newcommand{\sheafZZ}{\underline{\ZZ}}
}{
\newcommand{\nderef}[1]{la N.D.E.~(\ref{#1}) page~\pageref{#1}}
}

%%%%%%%%%%%%%%%%%
%%%%% expose XVII
%%%%%%%%%%%%%%%%%

% END INCLUDED SGA4_English_current_buildable_source_20260906/source/fichiers-configuration-latex/macros.tex


\externaldocument{../01/01_English.tex}
\externaldocument{../02/02_English.tex}
\externaldocument{../03/03_English.tex}
\externaldocument{../04/04_English.tex}
\externaldocument{../05/05_English.tex}
\externaldocument{../05bis/05bis_English.tex}
\externaldocument{../06/06_English.tex}
\externaldocument{../07/07_English.tex}
\externaldocument{../08/08_English.tex}
\externaldocument{../09/09_English.tex}
\externaldocument{../10/10_English.tex}
\externaldocument{../11/11_English.tex}

\setboolean{original}{false}

\begin{document}

\begin{center}
XII. The base-change theorem for a proper morphism\\
M. Artin
\end{center}

\selectlanguage{english}version: \showgitstatus
\tableofcontents

\else
\chapter{The base-change theorem for a proper morphism}
%\addtocontents{toc}{M. Artin}
\begin{center}
M. Artin
\end{center}
\fi
\phantomsection\MakeLinkTarget*{sga4.exp.XII}\label{sga4.exp.XII}

% BEGIN INCLUDED SGA4_English_current_buildable_source_20260906/source/12/drafts/SGA4_XII_LINES_43_233_English.texfrag
\setcounter{pageoriginal}{78}

\section{Introduction}\etiquette[1]{XII.sec1}\pageoriginale

In this exposé we prove the fundamental theorem which says, in a
slightly more precise form, that for a \emph{proper} morphism
$f:X\to Y$ of schemes and a torsion abelian sheaf $F$ on $X$, the
fiber of $R^qf_*F$ at a geometric point $\xi$ of $Y$ is isomorphic to
the cohomology $H^q(X_\xi,F|X_\xi)$ of the fiber $X_\xi$.

Recall that the analogous result for a proper morphism of paracompact
spaces is well known (Godement II 4.11) and easy. It is considerably
more delicate for schemes. For example, in contrast with the case of
paracompact spaces, the hypothesis that $F$ be torsion is essential.

The proof proceeds in several stages. By more or less formal
dévissage arguments, one reduces to the case where $Y$ is
Noetherian, the relative dimension of $f$ is at most $1$, and the
sheaf $F$ is constant. To handle this special case, one uses the
specialization theorem for the fundamental group
(\hyperref[XII.thm5.9]{5.9}), which is a nonabelian variant of the
theorem stated above and can be proved directly in the case under
consideration (\hyperrefext[XIII][XIII.sec2]{2}).

\section{An example}\etiquette[2]{XII.sec2}\pageoriginale

We show that the hypothesis that $F$ be torsion is necessary. Let
$X$ be a nonsingular surface over an algebraically closed field $k$,
let $Y$ be a nonsingular curve over $k$, and let $f:X\to Y$ be a
proper morphism whose generic fiber is smooth and whose geometric
fiber at a geometric point $\bar y$ of $Y$ is irreducible and has an
ordinary double point. Let $F$ be the constant sheaf $\ZZ_X$.

Recall the following fact. For any point $z$ of a scheme $Z$, let
$i:z\to Z$ be the inclusion. Then
$H^1(Z,i_*\ZZ_z)=0$ (\hyperrefext[IX][IX.prop3.6]{3.6} (i)). Since
$X$ is normal, hence unibranch, it follows that
$\ZZ_X\simeq i_*\ZZ_x$, where $i:x\to X$ is the inclusion of the
generic point. Thus $H^1(X,\ZZ_X)=0$ and likewise
$R^1f_*\ZZ_X=0$. But for the geometric fiber $X_\circ$ of $X/Y$
having the ordinary double point, I claim that
$H^1(X_\circ,\ZZ_{X_\circ})\simeq\ZZ$\footnote[1]{Compare, from the
point of view of $\Pi_1$, with \SGA 3 X (6, Examples) and (1.6).};
this gives the counterexample, since
$\ZZ_{X_\circ}\simeq\ZZ_X|X_\circ$.

Let $i_\circ:x_\circ\to X_\circ$ be the inclusion of the generic
point of $X_\circ$, and let $Q$ be the singular point of $X_\circ$.
The curve $X_\circ$ has two \lq\lq branches'' at $Q$, and it follows
that the fiber of $i_{\circ*}(\ZZ_{x_\circ})$ at $Q$ is isomorphic to
$\ZZ\times\ZZ$. Away from $Q$, the curve $X_\circ$ is normal, hence
$i_{\circ*}(\ZZ_{x_\circ})\simeq\ZZ_{X_\circ}$ on
$X_\circ-Q$. Therefore there is an exact sequence
$$
0\longrightarrow\ZZ_{X_\circ}\longrightarrow
i_{\circ*}(\ZZ_{x_\circ})\longrightarrow(\ZZ)_Q
\longrightarrow0{\quoi},
$$
where the last term is extension by zero of the sheaf $\ZZ$ at the
point $Q$. All three terms have $H^0$ isomorphic to $\ZZ$, and
$H^1(X_\circ,i_{\circ*}(\ZZ_{x_\circ}))=0$. The long exact
cohomology sequence therefore gives
$H^1(X_\circ,\ZZ_{X_\circ})=\ZZ$\pageoriginale.

\section{Review of nonabelian $H^1$}\etiquette[3]{XII.sec3}

We briefly recall some results on nonabelian cohomology that will be
needed below. We omit most proofs, which are well known in various
special cases and may be found in Giraud's thesis [1].

Let $X$ be a scheme and $F$ a sheaf of groups on $X$. For every
étale $X'/X$, one defines a pointed set $H^1(X',F)$, covariantly
functorial in $F$ and contravariantly functorial in $X'$. The functor
$H^1$ commutes with finite products of sheaves. If $f:X\to Y$ is a
morphism, one defines a sheaf of pointed sets (that is, with a
distinguished section) $R^1f_*F$. It is the sheaf associated with the
presheaf of pointed sets $\Rr^1F$ which assigns to every étale
$Y'/Y$ the set $H^1(\fprod{X}{Y'}{Y},F)$.

There are the usual exact cohomology sequences; we record one of
them:

\begin{enonce*}{Proposition 3.1}\etiquette[3.1]{XII.prop3.1}
Let $0\to F\xrightarrow{u}G$ be an injection of sheaves of groups,
and let $C=G/F$ be the sheaf of homogeneous spaces under $G$. There
is an exact sequence of sheaves of pointed sets
$$
0\to f_*F\to f_*G\to f_*C\xrightarrow{\partial}
R^1f_*F\xrightarrow{u^1}R^1f_*G{\quoi}.
$$
\end{enonce*}

This exact sequence is obtained, of course, by passing to the
associated sheaves in the exact sequence of presheaves whose value
on an étale $Y'/Y$ is
$$
0\to F(X')\to G(X')\to C(X')\xrightarrow{\partial}
H^1(X',F)\to H^1(X',G){\quoi},
$$
where\pageoriginale $X'=\fprod{X}{Y'}{Y}$. The equivalence relation
induced by $\partial$ is the one given by the action of $f_*G$ on
$f_*C$. As usual, the equivalence relation induced by $u^1$ may be
made explicit using the notion of \lq\lq twisting'' [1].

\begin{enonce*}{Proposition 3.2}\etiquette[3.2]{XII.prop3.2}
Let $X\xrightarrow{f}Y\xrightarrow{g}Z$ be morphisms of schemes. For
every étale $Z'/Z$, putting
$Y'=\fprod{Y}{Z'}{Z}$ and $X'=\fprod{X}{Z'}{Z}$, there is an exact
sequence
$$
0\longrightarrow H^1(Y',f_*F)\longrightarrow H^1(X',F)
\longrightarrow H^0(Y',R^1f_*F){\quoi}.
$$
Passing to the associated sheaves gives an exact sequence of pointed
sheaves on $Z$:
$$
0\longrightarrow (R^1g_*)f_*F\longrightarrow R^1(gf)_*F
\longrightarrow g_*(R^1f_*)F{\quoi}.
$$
\end{enonce*}

\begin{enonce*}{Proposition 3.3}\etiquette[3.3]{XII.prop3.3}
The family of functors of the form $H^1(X',F)$, for $X'$ étale over
$X$, and of the form $R^1f_*F$, for an arbitrary morphism
$f:X\to Y$, is effaceable: there exists an injection
$F\xrightarrow{u}G$ such that all the induced morphisms
$H^1(X',F)\to H^1(X',G)$ and $R^1f_*F\to R^1f_*G$ are zero. More
precisely, there exists such an injection for which every
$H^1(X',G)$, with $X'$ étale over $X$, and every pointed sheaf
$R^1f_*G$ vanish. If $X$ is Noetherian and $F$ is a constructible
sheaf of ind-$\LL$-groups, then such an injection exists with $G$
also a sheaf of ind-$\LL$-groups.
\end{enonce*}

\emph{Proof.} The first assertion is presumably true in an arbitrary
topos, but for the étale topology it is convenient to perform the
effacement using the \lq\lq Godement resolution''
\begin{equation*}
F\longrightarrow G=\prod_{x\in X}i_{x*}i_x^*F{\quoi}.
\tag{3.4}
\end{equation*}\etiquette[3.4]{XII.eq3.4}
Here\pageoriginale $x$ runs through the points of $X$ and
$i_x:\bar x\to X$ is a geometric point over $x$. For $X'/X$ étale,
one has
$$
G(X')=\prod_{x\in X}(i_x^*F)(\fprod{X'}{\bar x}{X})
=\prod_{x\in X}\left(\prod_{\bar x'/\bar x}F_{\bar x'}\right){\quoi},
$$
where the $\bar x'$ are geometric points of $X'$. For each point of
$X'$ there is plainly at least one $\bar x'$ above it; hence
$F(X')\subset G(X')$ (VIII 3), and therefore $F\subset G$.
Moreover, every étale $\bar u/\bar x$ is a disjoint union of spectra
of separably closed fields, so $H^1(\bar u,\bar G)=0$ for every sheaf
$\bar G$ on $\bar x$. It follows from
\hyperref[XII.prop3.2]{3.2} that
$H^1(X',i_{x*}i_x^*F)=0$ for every étale $X'/X$. Since $H^1$
commutes with products, $H^1(X',G)=0$, and consequently
$R^1f_*G=0$ for every $f:X\to Y$, by the description of $R^1f_*G$
as the sheaf associated with the presheaf
$Y'\mapsto H^1(\fprod{X}{Y'}{Y},G)$.

Suppose in addition that $F$ is constructible and $X$ is Noetherian.
It is immediate that there is a finite set of points
$\{x_1,\ldots,x_n\}$ of $X$ such that
$$
F\longrightarrow\prod_{1\leq i\leq n}i_{x_i*}i_{x_i}^*F
$$
is already injective. The sheaf $i_{x_i*}i_{x_i}^*F$ is a sheaf of
ind-finite groups (IX \hyperref[IX.def1.5]{1.5}), and hence so is the finite product
(\hyperrefext[IX][IX.prop1.4]{1.4} and
\hyperrefext[IX][IX.prop1.6]{1.6} (iii)). Replacing $G$ by this sheaf
gives the last assertion.

\setcounter{subsection}{4}
\subsection{}\etiquette[3.5]{XII.subsec3.5}
As in the case of abelian cohomology (treated in V \hyperref[V.subsec2.4]{2.4}), $H^1$ may be
computed by the \v{C}ech procedure [1].

% END INCLUDED SGA4_English_current_buildable_source_20260906/source/12/drafts/SGA4_XII_LINES_43_233_English.texfrag


% BEGIN INCLUDED SGA4_English_current_buildable_source_20260906/source/12/drafts/SGA4_XII_LINES_234_382_English.texfrag
\section{The base-change morphism}\etiquette[4]{XII.sec4}\pageoriginale

Let
$$
\xymatrix@C=2.75cm@R=1.5cm{
X \ar[d]_{f}&X'\ar[l]_{g'}\ar[d]^{f'}\\
S&S'\ar[l]_{g}
}
$$
be a Cartesian square of morphisms of schemes. From it one obtains a
morphism of functors $\varphi=\varphi_g$
\begin{equation*}
g^{\ast}f_{\ast}\xrightarrow{\varphi}f'_{\ast}{g'}^{\ast}\tag{4.1}
\end{equation*}\etiquette[4.1]{XII.eq4.1}
as follows. Giving such a morphism is equivalent, by the adjunction
between $g^{\ast}$ and $g_{\ast}$, to giving a morphism of functors
\begin{equation*}
f_{\ast}\longrightarrow g_{\ast}{f'}_{\ast}{g'}^{\ast}{\quoi}.
\tag{4.1 adj}
\end{equation*}\MakeLinkTarget*{sga4.exp.XII.equation.4.1-adj}\label{4.1 adj}
We take this to be the morphism
$$
f_{\ast}\xrightarrow{(\id\rightarrow {g'}_{\ast}{g'}^{\ast})}
f_{\ast}{g'}_{\ast}{g'}^{\ast}
=g_{\ast}{f'}_{\ast}{g'}^{\ast}{\quoi},
$$
using
$$
g_{\ast}{f'}_{\ast}=(gf')_{\ast}=(fg')_{\ast}
=f_{\ast}{g'}_{\ast}{\quoi}.
$$

There is another candidate for \eqhyperref[XII.eq4.1]{4.1}. Indeed,
giving such a morphism is also equivalent to giving
$$
{f'}^{\ast}g^{\ast}f_{\ast}\longrightarrow {g'}^{\ast},
$$
and one may take
$$
{f'}^{\ast}g^{\ast}f_{\ast}\simeq
{g'}^{\ast}f^{\ast}f_{\ast}
\xrightarrow{(f^{\ast}f_{\ast}\rightarrow\id)}{g'}^{\ast}{\quoi}.
$$

I do not know\pageoriginale whether these two choices always coincide
\footnote[1]{Mr. \textsc{Deligne} resolved this uncertainty in the
affirmative, in a considerably more general context; \cf Exposé XVII.};
in any event, we shall use only the first. We call it the
\emph{base-change morphism.}

Restricting \eqhyperref[XII.eq4.1]{4.1} to abelian sheaves
(respectively, to sheaves of groups) gives a morphism
$\varphi^q=\varphi^q_g$ for every $q\geq0$ (respectively, for
$q=0,1$):
\begin{equation*}
g^{\ast}(R^qf_{\ast})\xrightarrow{\varphi^q}
(R^q{f'}_{\ast}){g'}^{\ast}{\quoi}.
\tag{4.2}
\end{equation*}\etiquette[4.2]{XII.eq4.2}
Indeed, this is equivalent to giving
\begin{equation*}
(R^qf_{\ast})\longrightarrow
g_{\ast}(R^q{f'}_{\ast}){g'}^{\ast}{\quoi},
\tag{4.2 adj}
\end{equation*}\MakeLinkTarget*{sga4.exp.XII.equation.4.2-adj}\label{4.2 adj}
and there are the following \lq\lq evident'' morphisms:
\begin{gather*}
(R^qf_{\ast})\xrightarrow{a}
(R^qf_{\ast}){g'}_{\ast}{g'}^{\ast}
\xrightarrow{b}R^q(fg')_{\ast}{g'}^{\ast}\\
=R^q(gf')_{\ast}{g'}^{\ast}
\xrightarrow{c}g_{\ast}(R^q{f'}_{\ast}){g'}^{\ast}{\quoi}.
\end{gather*}

\begin{enonce*}[remark]{Remark 4.3}\etiquette[4.3]{XII.rem4.3}
Suppose that $g$ (and hence $g'$) is a closed immersion. Then the
morphisms $b$ and $c$ above are isomorphisms (VIII \hyperref[VIII.cor5.7]{5.7}). It follows
that \eqhyperref[XII.eq4.2]{4.2} is essentially just the morphism $a$
above, obtained by applying $R^qf_{\ast}$ to the canonical morphism
$F\rightarrow {g'}_{\ast}{g'}^{\ast}F$.
\end{enonce*}

\begin{enonce*}{Proposition 4.4}\etiquette[4.4]{XII.prop4.4}
(Composition of base-change morphisms).
\begin{enumerate}
\renewcommand{\theenumi}{\roman{enumi}}
\renewcommand{\labelenumi}{(\theenumi)}
\item Let
$$
\xymatrix@C=2.5cm@R=1.5cm{
X\ar[d]_{f}&X'\ar[l]_{g'}\ar[d]_{f'}&
X''\ar[l]_{h'}\ar[d]_{f''}\\
S&S'\ar[l]_{g}&S''\ar[l]_{h}
}
$$
be a\pageoriginale Cartesian diagram of morphisms of schemes. Then
the following diagram is commutative:
$$
\xymatrix@C=2cm@R=1.5cm{
(gh)^{\ast}f_{\ast}\ar[rr]^{\varphi_{gh}}
\ar[d]^{\text{\rotatebox{90}{$\backsim$}}}&&
f''_{\ast}(g'h')^{\ast}
\ar[d]^{\text{\rotatebox{90}{$\backsim$}}}\\
h^{\ast}g^{\ast}f_{\ast}
\ar[r]^{h^{\ast}\varphi_g}&
h^{\ast}f'_{\ast}g'^{\ast}
\ar[r]^{\varphi_h g'^{\ast}}&
f''_{\ast}h'^{\ast}g'^{\ast}.}
$$
The same is true of the diagram obtained by replacing
$f_{\ast}$, $f'_{\ast}$, and $f''_{\ast}$ by
$R^qf_{\ast}$, $R^qf'_{\ast}$, and $R^qf''_{\ast}$, in the natural
domains of definition of the functors in question
\textup{(\cf $\textup{n}^{\circ}$ 3)}.

\item Let
$$
\xymatrix@C=3.25cm@R=1.25cm{
Y\ar[d]_{e}&Y'\ar[l]_{g''}\ar[d]^{e'}\\
X\ar[d]_{f}&X'\ar[l]_{g'}\ar[d]^{f'}\\
S&S'\ar[l]_{g}
}
$$
be a Cartesian diagram of morphisms of schemes. Then the following
diagram is commutative:
$$
\xymatrix@C=2.75cm@R=1.5cm{
g^{\ast}(fe)_{\ast}\ar[rr]^{\varphi_g^{fe}}\ar@{=}[d]&&
(f'e')_{\ast}g''^{\ast}\ar@{=}[d]\\
g^{\ast}f_{\ast}e_{\ast}
\ar[r]^{\varphi_g^f e_{\ast}}&
f'_{\ast}g'{}^{\ast}e_{\ast}
\ar[r]^{f'_{\ast}\varphi_{g'}^e}&
f'_{\ast}e'_{\ast}g''^{\ast},
}
$$
where the three horizontal arrows arise respectively from the
base-change homomorphisms for $fe$, $f$, and $e$.

In the abelian case, the base-change morphisms for $fe$ are induced by
a morphism of spectral sequences
$$
\xymatrix{
E^{p,q}_2=g^{\ast}(R^pf_{\ast})(R^qe_{\ast})F
\ar@{=>}[r]\ar[d]_{\varphi}&
g^{\ast}R^{p+q}(fe)_{\ast}F\ar[d]_{\varphi}\\
(R^pf'_{\ast})(R^qe'_{\ast})g''^{\ast}F
\ar@{=>}[r]&R^{p+q}(f'e')_{\ast}g''^{\ast}F,
}
$$
which, on the initial terms, is induced by the base-change morphisms
for $e$ and $f$. For a sheaf of groups, there is a morphism of exact
sequences
$$
\xymatrix{
0\ar[r]&g^{\ast}(R^1f_{\ast})e_{\ast}F
\ar[r]\ar[d]^{\varphi}&
g^{\ast}R^1(fe)_{\ast}F\ar[r]\ar[d]^{\varphi}&
g^{\ast}f_{\ast}(R^1e_{\ast})F\ar[d]^{\varphi}\\
0\ar[r]&(R^1f'_{\ast})e'_{\ast}g''^{\ast}F\ar[r]&
R^1(f'e')_{\ast}g''^{\ast}F\ar[r]&
f'_{\ast}(R^1e'_{\ast})g''^{\ast}F.
}
$$
\end{enumerate}
\end{enonce*}

We leave the proof to the reader.

% END INCLUDED SGA4_English_current_buildable_source_20260906/source/12/drafts/SGA4_XII_LINES_234_382_English.texfrag


% BEGIN INCLUDED SGA4_English_current_buildable_source_20260906/source/12/drafts/SGA4_XII_LINES_383_784_English.texfrag
\section{Statement of the main theorem and some variants}\etiquette[5]{XII.sec5}

The theorem is as follows:

\begin{enonce*}{Theorem 5.1}\etiquette[5.1]{XII.thm5.1}
(Proper base-change theorem). Let
$$
\xymatrix@C=3.25cm@R=1.5cm{
X \ar[d]_{f} & X' \ar[l]_{g'} \ar[d]^{f'}\\
S & S' \ar[l]_{g}
}
$$
be a Cartesian diagram of schemes, with $f$ proper.
\begin{enumerate}
\renewcommand{\theenumi}{\roman{enumi}}
\renewcommand{\labelenumi}{(\theenumi)}
\item Let $F$ be a sheaf of sets on $X$. Then the morphism
  \eqhyperref[XII.eq4.1]{4.1}
$$
\varphi:g^{\ast}f_{\ast}F\longrightarrow {f'}_{\ast}{g'}^{\ast}F
$$
is bijective.

\item Let\pageoriginale $F$ be a sheaf of groups (resp. of
  ind-finite groups) on $X$. Then the morphism
  \eqhyperref[XII.eq4.2]{4.2}
$$
\varphi^1:g^{\ast}(R^1f_{\ast})F\longrightarrow
(R^1{f'}_{\ast}){g'}^{\ast}F
$$
is injective (resp. bijective).

\item Let $F$ be a torsion abelian sheaf on $X$. Then the morphism
  \eqhyperref[XII.eq4.2]{4.2}
$$
\varphi^q:g^{\ast}(R^qf_{\ast})F\longrightarrow
(R^q{f'}_{\ast}){g'}^{\ast}F
$$
is bijective for every $q\geq0$.
\end{enumerate}
\end{enonce*}

We first state some immediate consequences of the theorem. If
$S'=\xi$ is a geometric point of $S$, so that a sheaf on $\xi$ is
determined by its global sections, then
$H^0(\xi,R^q{f'}_{\ast}G)=H^q(X',G)$ for every sheaf $G$ on $X'$.
This gives the following corollary (which is, moreover, essentially
equivalent to \hyperref[XII.thm5.1]{5.1}):

\begin{enonce*}{Corollary 5.2}\etiquette[5.2]{XII.cor5.2}
Let $f:X\longrightarrow S$ be proper, let $\xi\longrightarrow S$ be
a geometric point, and let $X_{\xi}$ be the fiber of $f$ at $\xi$.
\begin{enumerate}
\renewcommand{\theenumi}{\roman{enumi}}
\renewcommand{\labelenumi}{(\theenumi)}
\item If $F$ is a sheaf of sets on $X$, the canonical base-change
  morphism
$$
(f_{\ast}F)_{\xi}\longrightarrow H^0(X_{\xi},F|X_{\xi})
$$
is bijective.

\item If $F$ is a sheaf of groups (resp. of ind-finite groups), the
  base-change morphism
$$
(R^1f_{\ast}F)_{\xi}\longrightarrow H^1(X_{\xi},F|X_{\xi})
$$
is injective (resp. bijective).

\item If\pageoriginale $F$ is a torsion abelian sheaf, then for every
  $q\geq0$ the base-change morphism
$$
(R^qf_{\ast}F)_{\xi}\longrightarrow H^q(X_{\xi},F|X_{\xi})
$$
is bijective.
\end{enumerate}
\end{enonce*}

Taking \hyperrefext[X][X.cor4.3]{4.3} and
\hyperrefext[X][X.cor5.2]{5.2} into account, we obtain:

\begin{enonce*}{Corollary 5.3}\etiquette[5.3]{XII.cor5.3}
Under the hypotheses of \hyperrefext[X][X.cor5.2]{5.2}, let $n$ be
the dimension of the fiber $X_{\xi}$. Then, for every torsion
abelian sheaf $F$ on $X$, the fiber $(R^qf_{\ast}F)_{\xi}$ vanishes
for $q>2n$. If $F$ is a sheaf of $p$-torsion groups, where $p$ is
the characteristic of $\xi$, then the fiber of $R^qf_{\ast}F$ at
$\xi$ vanishes for $q>n$.
\end{enonce*}

Recall that the \emph{relative dimension of $X/S$ is} $\leq n$ if
every fiber has dimension $\leq n$. We therefore have:

\begin{enonce*}{Corollary 5.3 bis}\etiquette[5.3 bis]{XII.corbis5.3}
Let $f:X\longrightarrow S$ be proper, with $X/S$ of relative
dimension $\leq n$. Then $R^qf_{\ast}F=0$ for $q>2n$ and every
torsion abelian sheaf $F$ on $X$. If $S$ has characteristic $p$ and
$F$ is a sheaf of $p$-torsion abelian groups, then
$R^qf_{\ast}F=0$ for $q>n$.
\end{enonce*}

Applying \hyperref[XII.thm5.1]{5.1} to a morphism between spectra of
separably closed fields gives:

\begin{enonce*}{Corollary 5.4}\etiquette[5.4]{XII.cor5.4}
(Invariance of cohomology under extension of the base field in the
  proper case). Let $k\subset K$ be separably closed fields and let
  $X$ be a proper scheme over $\Spec k$. Put $X'=X_K$. Let $F$ be a
  sheaf on $X$ and let $F'$ be its inverse-image sheaf on $X'$. Then:
\begin{enumerate}
\renewcommand{\theenumi}{\roman{enumi}}
\renewcommand{\labelenumi}{(\theenumi)}
\item $H^0(X,F)\xrightarrow{\sim}H^0(X',F')$.

\item If\pageoriginale $F$ is a sheaf of groups (resp. of
  ind-finite groups), the morphism
$$
H^1(X,F)\longrightarrow H^1(X',F')
$$
is injective (resp. bijective).

\item If $F$ is a torsion abelian sheaf, then
$$
H^q(X,F)\xrightarrow{\sim}H^q(X',F')
$$
for every $q$.
\end{enumerate}
\end{enonce*}

A slightly different variant of \hyperref[XII.cor5.2]{5.2} is the
following:

\begin{enonce*}{Corollary 5.5}\etiquette[5.5]{XII.cor5.5}
Let $S$ be the spectrum of a henselian ring, and let
$S'=s_{\circ}$ be its closed point. Let $f:X\longrightarrow S$ be
proper. Put $X_{\circ}=X'$ and $F_{\circ}={g'}^{\ast}F$. Then:
\begin{enumerate}
\renewcommand{\theenumi}{\roman{enumi}}
\renewcommand{\labelenumi}{(\theenumi)}
\item $H^0(X,F)\xrightarrow{\sim}H^0(X_{\circ},F_{\circ})$

for every sheaf of sets $F$.

\item The morphism
$H^1(X,F)\longrightarrow H^1(X_{\circ},F_{\circ})$

is injective (resp. bijective) for every sheaf of groups (resp. of
ind-finite groups) $F$.

\item $H^q(X,F)\xrightarrow{\sim}H^q(X_{\circ},F_{\circ})$

for every torsion abelian sheaf $F$ and every $q$.
\end{enumerate}
\end{enonce*}

Indeed, this statement is just a special case of
\hyperref[XII.cor5.2]{5.2} when $S$ is strictly local. Suppose that
\hyperref[XII.cor5.2]{5.2} is known. Let $\bar S$ be the strict
localization of $S$, let $\bar s$ be the closed point of $\bar S$,
and let $G$ be the Galois group of $\bar S/S$ (which is also the
Galois group of $\bar s/s_{\circ}$). Then $G$ acts on the fibers of
the $R^q$, and, with the evident notation, one has
\begin{align*}
H^0(X,F)&\simeq H^0(\bar X,\bar F)^G,\\
H^0(X_{\circ},F_{\circ})&\simeq
H^0(\bar X_{\circ},\bar F_{\circ})^G
\end{align*}
for\pageoriginale a sheaf of sets $F$. Assertion (i) follows, since
$H^0(\bar X,\bar F)\xrightarrow{\sim}
H^0(\bar X_{\circ},\bar F_{\circ})$ and this isomorphism plainly
commutes with the action of $G$.

If $F$ is a sheaf of groups, one readily defines a morphism of exact
sequences of pointed sets
$$\vcenter{\etiquette[5.5.1]{XII.eq5.5.1}
\xymatrix{
0 \ar[r] & H^1(G,H^0(\bar X,\bar F)) \ar[r]\ar[d] &
H^1(X,F) \ar[r]\ar[d] & H^1(\bar X,\bar F)^G \ar[d]\\
0 \ar[r] & H^1(G,H^0(\bar X_{\circ},\bar F_{\circ})) \ar[r] &
H^1(X_{\circ},F_{\circ}) \ar[r] &
H^1(\bar X_{\circ},\bar F_{\circ})^G ,
}}\leqno{(5.5.1)}
$$
in which the left-hand vertical arrow is an isomorphism and the
right-hand one is injective by \hyperref[XII.cor5.2]{5.2}. Let
$\alpha,\beta\in H^1(X,F)$ be
classes having the same image in $H^1(X_{\circ},F_{\circ})$, and let
$F^{\alpha}$ be the sheaf obtained by twisting $F$ by $\alpha$
\cite{sga4.exp.XII.bib.1}. Let $\beta'$ be the class in $H^1(X,F^{\alpha})$
corresponding to $\beta$. Since the class $\alpha'$ corresponding to
$\alpha$ is neutral, the image of $\beta'$ in
$H^1(X_{\circ},(F^{\alpha})_{\circ})$ is neutral.

Consider the diagram deduced from \eqhyperref[XII.eq5.5.1]{5.5.1}
by replacing $F$ with $F^{\alpha}$. It shows that the image of
$\beta'$ in $H^1(\bar X,\bar F^{\alpha})^G$ is neutral, and hence
that $\beta'$ comes from an element of
$H^1(G,H^0(\bar X,\bar F^{\alpha}))$; this element must itself be
neutral. Thus $\beta'$ is neutral, that is,
$\beta'=\alpha'$, whence $\beta=\alpha$. This proves the injectivity
of the morphism in (ii).

Now suppose that $F$ is a sheaf of ind-finite groups, and let us
prove that the arrow in (ii) is surjective. It is enough to show that
the functor $F\mapsto H^1(X_{\circ},F_{\circ})$ is effaceable.
Indeed, let $0\longrightarrow F\xrightarrow{u}G$ be an injection
which effaces this functor, and put $C=G/F$. There is a morphism of
exact sequences (\hyperref[XII.prop4.4]{4.4})
$$
\xymatrix{
H^0(X,C) \ar[r]\ar[d] & H^1(X,F) \ar[d]_{\epsilon}\ar[r] &
H^1(X,G) \ar[d]\\
H^0(X_{\circ},C_{\circ}) \ar[r] &
H^1(X_{\circ},F_{\circ}) \ar[r]^{u_{\circ}^1} &
H^1(X_{\circ},G_{\circ})
}
$$
in which $u_{\circ}^1$ is zero; this indeed implies the surjectivity
of $\epsilon$.

Let us prove\pageoriginale that the functor
$F\mapsto H^1(X_{\circ},F_{\circ})$ is effaceable for sheaves $F$ of
ind-finite groups. Let $F\xrightarrow{u}C$ be the \lq\lq Godement
resolution'' \eqhyperref[XII.eq3.4]{3.4}. We have
$$
H^1(G,H^0(\bar X_{\circ},\bar C_{\circ}))\simeq
H^1(G,H^0(\bar X,\bar C))\subset H^1(X,C)=0{\quoi}.
$$

Moreover, there is a commutative square
$$
\xymatrix{
H^1(\bar X,\bar F) \ar@{-}[r]^{\sim}\ar[d] &
H^1(\bar X_{\circ},\bar F_{\circ})\ar[d]\\
0=H^1(\bar X,\bar C) \ar[r] &
H^1(\bar X_{\circ},\bar C_{\circ})
}
$$
(indeed,
$H^1(\bar X,\bar C)=\varinjlim H^1(X',C)=0$, where $X'/X$ runs
through the étale morphisms), so its right-hand vertical arrow is
zero. We therefore have a morphism of exact sequences
$$
\xymatrix{
0 \ar[r] & H^1(G,H^0(\bar X_{\circ},\bar F_{\circ}))
\ar[r]\ar[d] & H^1(X_{\circ},F_{\circ}) \ar[r]\ar[d]^{\beta} &
H^1(\bar X_{\circ},\bar F_{\circ})^G \ar[d]^{\alpha^G}\\
&0 \ar[r] & H^1(X_{\circ},C_{\circ}) \ar[r] &
H^1(\bar X_{\circ},\bar C_{\circ})^G
}
$$
in which $\alpha^G$ is zero. It follows that $\beta$ is also zero,
and hence that the functor is effaceable.

Finally, let us prove assertion (iii) of
\hyperref[XII.cor5.5]{5.5}. There is a morphism of spectral
sequences
$$
\xymatrix{
**[l]E_2^{p,q}=H^p(G,H^q(\bar X,\bar F))
\ar@{=>}[r]\ar[d] & H^n(X,F)\ar[d]^{\varphi}\\
H^p(G,H^q(\bar X_{\circ},\bar F_{\circ}))
\ar@{=>}[r] & H^n(X_{\circ},F_{\circ}).
}
$$
By \hyperref[XII.cor5.2]{5.2}, the morphisms
$H^q(\bar X,\bar F)\longrightarrow
H^q(\bar X_{\circ},\bar F_{\circ})$ are isomorphisms. Consequently
the morphism of $E_2^{p,q}$ terms is an isomorphism, and therefore so
is the morphism of abutments. This proves the result, \cqfd

\setcounter{subsection}{5}
\subsection{}\etiquette[5.6]{XII.subsec5.6}
Let us examine\pageoriginale \hyperref[XII.cor5.5]{5.5} (i) and (ii)
when $F$ is constant. Let $F=T_X$ be the constant sheaf of sets on
$X$ with value $T$. Then $F(X)$ is simply the set of locally constant
maps from $X$ to $T$. If $X_{\circ}$ is a subscheme of $X$ and $T$
has at least two elements, it follows that
$F(X)\longrightarrow F(X_{\circ})$ is bijective if and only if the
map
$$
U\longmapsto U_{\circ}=U\cap X_{\circ}
$$
from the set $\Theta(X)$ of subsets of $X$ that are both open and
closed to the analogous set $\Theta(X_{\circ})$ is bijective. If the
set $\Pi_{\circ}(X_{\circ})$ of connected components of $X_{\circ}$
is finite, this is equivalent to saying that the map
$\Pi_{\circ}(X_{\circ})\longrightarrow\Pi_{\circ}(X)$ is bijective.
Thus \hyperref[XII.cor5.5]{5.5} (i) may also be restated as follows
(bearing in mind that in this case $X_{\circ}$ is of finite type over
a field, hence Noetherian, so that $\Pi_{\circ}(X_{\circ})$ is
finite):

\begin{enonce*}{Corollary 5.7}\etiquette[5.7]{XII.cor5.7}
Let $S$ be the spectrum of a henselian ring, let $s_{\circ}$ be its
closed point, let $f:X\longrightarrow S$ be a proper morphism, and
let $X_{\circ}$ be the closed fiber of $X/S$. Then the map
$$
\Pi_{\circ}(X_{\circ})\longrightarrow\Pi_{\circ}(X)
$$
induced by $X_{\circ}\longrightarrow X$ is bijective.
\end{enonce*}

(In particular, this implies that $\Pi_{\circ}(X)$ is finite.) We
shall prove \hyperref[XII.cor5.7]{5.7} directly in the Noetherian
case:

\begin{enonce*}{Lemma 5.8}\etiquette[5.8]{XII.lem5.8}
Statement \hyperref[XII.cor5.7]{5.7} is true if $S$ is Noetherian
\footnote[1]{\Cf \EGA IV 18.5.19.}.
\end{enonce*}

\emph{Proof.} In the Noetherian case, $\Pi_{\circ}(X)$ is certainly
finite. Replacing $X$ by a connected component, we may suppose that
$\Pi_{\circ}(X)$ has one element, that is, that $X$ is connected and
nonempty. We must show that $X_{\circ}$ is also connected and
nonempty. The image of $X$ in $S$ is closed because $f$ is proper,
and is nonempty because $X$ is nonempty; it therefore contains
$s_{\circ}$, which shows that $X_{\circ}$ is nonempty. Let
$$
\xymatrix@R=1cm@C=1.5cm{
X \ar[rd]^{f'} \ar[dd]_{f}& \\
  & S'\ar[ld]^{\pi}\\
S &
}
$$
be the\pageoriginale Stein factorization of $f$ (\EGA III 4.3.3), so
that $S'/S$ is finite and $S'$ is connected. Every fiber of $f'$ is
known to be connected and nonempty, and the underlying space of
$X_{\circ}$ is the union of the underlying spaces of the fibers over
the closed points of $S'$. It remains only to show that $S'$ is the
spectrum of a \emph{local} ring. Since $S$ is henselian, the ring
$A'$ of $S'$ is a product of local rings (VII \hyperref[VII.prop4.1]{4.1} (i)); because $S'$
is connected, $A'$ is local. This proves \hyperref[XII.lem5.8]{5.8}.

\setcounter{subsection}{9}
\setcounter{subsubsection}{-1}
\subsubsection{}\etiquette[5.9.0]{XII.subsubsec5.9.0}
Suppose now that $F=G_X$ is a constant sheaf of groups with value a
finite group $G$. Then $H^1(X,G_X)$ classifies the principal Galois
coverings of $X$ with group $G$ (VII \hyperref[VII.rem2.1]{2.1}). Corollary
\hyperref[XII.cor5.5]{5.5} (ii) says that the morphism
$X_{\circ}\longrightarrow X$ induces a bijection between the sets of
isomorphism classes of these coverings. By the definition of the
$\Pi_1$, it follows that \hyperref[XII.cor5.5]{5.5} (ii) for constant
sheaves of finite groups (in the presence of (i)) is equivalent to:

\begin{enonce*}{Theorem 5.9}\etiquette[5.9]{XII.thm5.9}
(Specialization theorem for the fundamental group). With the notation
  of \hyperref[XII.cor5.5]{5.5}, suppose that $X$ is connected and
  nonempty. Then $X_{\circ}$ is connected and nonempty; and, if
  $\bar x_{\circ}$ is a geometric point of $X_{\circ}$, the morphism
  of fundamental groups
$$
\Pi_1(X_{\circ},\bar x_{\circ})\longrightarrow
\Pi_1(X,\bar x_{\circ})
$$
is an isomorphism.
\end{enonce*}

On the other\pageoriginale hand, by Galois theory (\SGA V), statement
\hyperref[XII.thm5.9]{5.9} (in the case where $X_{\circ}$ is
connected, as we may suppose by \hyperref[XII.cor5.7]{5.7}) is
equivalent to:

\begin{enonce*}{Theorem 5.9 bis}\etiquette[5.9 bis]{XII.thmbis5.9}
Let $S$ be the spectrum of a henselian ring, let $s_{\circ}$ be the
closed point of $S$, let $f:X\longrightarrow S$ be a proper
morphism, and let $X_{\circ}$ be the closed fiber of $X/S$. For a
scheme $Z$, denote by $\Et(Z)$ the category of étale coverings of
$Z$. Then the restriction functor
$$
\Et(X)\longrightarrow\Et(X_{\circ})
$$
is an equivalence of categories.
\end{enonce*}

\begin{enonce*}[remark]{Remark 5.10}\etiquette[5.10]{XII.rem5.10}
Recall that \hyperref[XII.thm5.9]{5.9} was already proved
(\SGA 1 X 2.1) when $S$ is the spectrum of a complete Noetherian
local ring. Unfortunately, that result does not seem to yield a
proof of \hyperref[XII.thm5.9]{5.9} in general
\footnote[1]{The situation has changed since these lines were
  written; \cf Th.~\hyperref[XII.prop3.1]{3.1} in M. Artin,
  Algebraic approximation of structures over complete local rings
  (to appear).}. The proof of \hyperref[XII.thm5.9]{5.9} given in this exposé and the next is
very different from the one given in \SGA 1 X.
\end{enonce*}

\setcounter{subsection}{10}
\subsection{}\etiquette[5.11]{XII.subsec5.11}
One can give another \lq\lq noncommutative'' variant of
\hyperref[XII.cor5.5]{5.5} (ii), using J. \textsc{Giraud}'s theory of
2-cohomology (\cf J. \textsc{Giraud}'s thesis, in preparation).
Suppose that $S$ is Noetherian. Let $L$ be a \emph{band} on $X$, and
suppose that $L$ is represented locally (for the étale topology) by
a sheaf of ind-finite groups. Consider the map
$$
\varphi^{(2)}:H^2(X,L)\longrightarrow H^2(X_{\circ},L_{\circ}){\quoi},
$$
where $L_{\circ}$ is the restriction of $L$ to $X_{\circ}$. Then
\begin{equation*}
(\varphi^{(2)})^{-1}(H^2(X_{\circ},L_{\circ})')
=H^2(X,L)'{\quoi},\tag{5.11.1}
\end{equation*}\etiquette[5.11.1]{XII.eq5.11.1}
where\pageoriginale the \lq\lq primes'' denote the subsets formed by
the neutral elements.

In view of \hyperref[XII.cor5.5]{5.5} (iii), this statement is also
equivalent to the following:

\emph{For every gerbe $\mathcal G$ on $X$, with band $L$ as above,
  let $\mathcal G_{\circ}$ denote its restriction to $X_{\circ}$.
  Then the restriction functor on the categories of sections}
$$
\mathcal G(X)\longrightarrow\mathcal G_{\circ}(X_{\circ})
$$
\emph{is an equivalence of categories.}

Indeed, the fact that this functor is fully faithful follows from
\hyperref[XII.cor5.5]{5.5} (i), and the fact that it is essentially
surjective when $\mathcal G(X)\neq\varnothing$, that is, when the
class $\xi\in H^2(X,\sisi{\mathcal G}{L})$ of $\mathcal G$ is
neutral, follows from \hyperref[XII.cor5.5]{5.5} (ii). It remains to
express the fact that
$\mathcal G_{\circ}(X_{\circ})\neq\varnothing$ implies
$\mathcal G(X)\neq\varnothing$, and this is precisely
\eqhyperref[XII.eq5.11.1]{5.11.1}.

We shall not give the proof of \eqhyperref[XII.eq5.11.1]{5.11.1}
here. We merely point out that, for every Noetherian scheme $X$ and
every closed subscheme of $X$, one can in fact show that the validity of
statements \hyperref[XII.cor5.5]{5.5} (i), (ii) implies that of
\eqhyperref[XII.eq5.11.1]{5.11.1}. The proof is fairly elementary
and uses the descent result VIII \hyperref[VIII.thm9.4]{9.4} a).

% END INCLUDED SGA4_English_current_buildable_source_20260906/source/12/drafts/SGA4_XII_LINES_383_784_English.texfrag


% BEGIN INCLUDED SGA4_English_current_buildable_source_20260906/source/12/drafts/SGA4_XII_LINES_785_1200_English.texfrag
\section{First reductions}\etiquette[6]{XII.sec6}

We use the abbreviation $(f,F,g)$ for the data of
\hyperref[XII.thm5.1]{5.1} in any one of cases (i), (ii), or (iii) of
that theorem.

\begin{enonce*}{Lemma 6.1}\etiquette[6.1]{XII.lem6.1}
For \hyperref[XII.thm5.1]{5.1} to hold for the data $(f,F,g)$, where
$f$ and $F$ are fixed and $g$ is arbitrary, it is enough that the
theorem hold for $(\bar f,\bar F,\bar g)$ for every $S$-scheme
$\bar S$ that is the strict localization of a scheme $S'$ locally\pageoriginale of
finite type over $S$ at a point closed in its fiber. Here $\bar f$
and $\bar F$ are obtained from $f$ and $F$ by the base change
$\bar S\to S$, and $\bar g:\bar s\to\bar S$ is the inclusion of the
closed point of $\bar S$.
\end{enonce*}

\noindent
\emph{Proof.} We may plainly assume that $S$ and $S'$ are affine,
say the spectra of rings $A$ and $A'$. Then $A'$ is the filtered
direct limit of its finite-type $A$-subalgebras $A'_i$. Using the
theory of passage to the limit (VII \hyperref[VII.cor5.11]{5.11} and \hyperref[VII.rem5.14]{5.14}), we are reduced to
proving \hyperref[XII.thm5.1]{5.1} for the triples $(f,F,g_i)$, where
$g_i:\Spec(A'_i)=S'_i\to S$; this reduces us to the case in which
$g$ is of finite type. To prove that the base-change homomorphisms
are isomorphisms, it is enough to prove that they induce isomorphisms
on the fibers of the sheaves in question, and it is enough to
consider geometric fibers at points $s'$ of $S'$ that are closed in
their fibers over $S$ (VIII \hyperref[VIII.rem3.13]{3.13}). Let $\bar s'$ be a geometric
point of $S'$ corresponding to a separable closure of $k(s')$.
Since $s'$ is algebraic over the corresponding point $s$ of $S$,
$k(\bar s')$ is also a separable closure of $k(s)$. Denote by
$\bar s$ the prescheme $\bar s'$, regarded as lying over $S$, that
is, as a geometric point of $S$, and consider the strict
localizations $\bar S$ and $\bar S'$ of $S$ and $S'$ at $\bar s$ and
$\bar s'$, respectively. We then have a commutative square
$$
\vcenter{
\xymatrix@C=3cm@R=1.25cm{
\bar{s}'  \ar[r] \ar[d]& \bar{s}\ar[d]\\
\bar{S}'  \ar[r] &  \bar{S},
}}
$$
in which the first horizontal arrow is an isomorphism. In view of
VIII \hyperref[VIII.thm5.2]{5.2} and \hyperref[VIII.rem5.3]{5.3}, the homomorphism induced on the fibers at
$\bar s'$ by the base-change homomorphism associated with $(f,F,g)$
identifies with the homomorphism
$$
H^{i}(\bar{X},\bar{F})\longrightarrow H^{i}(\bar{X}',\bar{F}')
$$
induced\pageoriginale by $\bar X'\to\bar X$, where $\bar X$ and
$\bar X'$ denote the schemes obtained from $X$ by the base changes
$\bar S\to S$ and $\bar S'\to S$, respectively. The commutative
square above plainly gives a commutative square
$$
\xymatrix@C=3cm@R=1.5cm{
H^{i}(\bar{X}_{\bar{s}}, \bar{F}) \ar[r]&
H^{i}(\bar{X}'_{\bar{s}'},\bar{F}')\\
H^{i}(\bar{X},\bar{F})  \ar[r] \ar[u]&
H^{i}(\bar{X}',\bar{F}') \ar[u],
}
$$
in which the first horizontal arrow is again an isomorphism.
Consequently the second horizontal arrow is an isomorphism as soon
as the two vertical arrows are. This is precisely what the hypothesis
of \hyperref[XII.lem6.1]{6.1} guarantees, \cqfd

\setcounter{subsection}{1}
\subsubsection{}\etiquette[6.1.1]{XII.subsubsec6.1.1}
Thus, to prove \hyperref[XII.thm5.1]{5.1} in a situation $(F,f,-)$,
we are reduced to studying the special case
\hyperref[XII.cor5.5]{5.5}, with $S$ strictly local. This leads us to
study in general a morphism $h:Y\to X$ of schemes and to seek general
conditions under which the corresponding homomorphisms
$$
H^{i}(X,F)\longrightarrow H^{i}(Y,h^{\ast}(F))
$$
are bijective. This subsection gives some easy auxiliary results on
this situation. For convenience in later references, we state them
in greater generality and precision than is needed for the proof of
\hyperref[XII.thm5.1]{5.1}.

\begin{enonce*}{Lemma 6.2}\etiquette[6.2]{XII.lem6.2}
Let $C$ and $C'$ be two abelian categories, let
$T^{\bigdot}=(T^{i})$ and $T'^{\bigdot}=({T'}^{i})$ be two
cohomological functors from $C$ to $C'$, and let
$\varphi=(\varphi^{i})$ be a morphism of\pageoriginale cohomological functors from
$T^{\bigdot}$ to $T'^{\bigdot}$. Assume that $T^{\bigdot}$ and
$T'^{\bigdot}$ vanish in degrees $<0$, and that $T^{i}$ is
effaceable for $i>0$. Let $n$ be an integer. The following conditions
are equivalent:
\begin{enumerate}
\renewcommand{\theenumi}{\alph{enumi}}
\renewcommand{\labelenumi}{\theenumi)}
\item $\varphi^{i}$ is an isomorphism for $i\leq n$ and a
  monomorphism for $i=n+1$.

\item (If $n\geq -1$.) The homomorphism $\varphi^{i}$ is a
  monomorphism for $i=0$ and an epimorphism for $i\leq n$.

\item (If $n\geq 0$.) The homomorphism $\varphi^{\circ}$ is an
  isomorphism, and the ${T'}^{i}$ are effaceable for $0<i\leq n$.
\end{enumerate}

If these conditions hold and $A\in\Ob(C)$, then
$\varphi^{n+1}(A):T^{n+1}(A)\to {T'}^{n+1}(A)$ is an epimorphism
(hence an isomorphism) if and only if ${T'}^{n+1}(A)$ is effaceable,
that is, if there exists a monomorphism $A\to B$ such that the
corresponding homomorphism
${T'}^{n+1}(A)\to {T'}^{n+1}(B)$ is zero. If $C'=(\Ab)$, the category
of abelian groups, and $\xi\in {T'}^{n+1}(A)$, then $\xi$ lies in the
image of $T^{n+1}(A)$ if and only if $\xi$ is effaceable, that is, if
there exists a monomorphism $A\to B$ such that the image of $\xi$ in
${T'}^{n+1}(B)$ is zero.
\end{enonce*}

\noindent
\emph{Proof.} It is clear that (a) implies (b) and (c). Since the
hypotheses on $T^{\bigdot}$ imply that the $T^{i}$ for $i>0$ are the
right satellites of $T^{\circ}$, the implication (c) $\Rightarrow$
(a) follows from the axiomatic characterization of satellite
functors. It also follows, by induction on $n$, from the last
assertion of \hyperref[XII.lem6.2]{6.2}, proved by the familiar
argument that we do not repeat here (compare the proof of
\hyperref[XII.cor6.6]{6.6} below). The implication (b)
$\Rightarrow$ (a) is immediate if $n=-1$. If $n\geq0$, we proceed by
induction on $n$: the induction hypothesis allows us to suppose that
$\varphi^{i}$ is an isomorphism for $i\leq n-1$ and\pageoriginale a monomorphism
for $i=n$. The hypothesis then makes it an isomorphism also for
$i=n$, and it remains to show, using the fact that $\varphi^n$ is an
isomorphism, that it is a monomorphism for $i=n+1$. Let $A\to B$ be
a monomorphism that effaces $T^{n+1}$, giving an exact sequence
$0\to A\to B\to C\to0$ in $C$ and a morphism of exact sequences in
$C'$:
$$
\xymatrix@R=1.25cm{
T^{n}(B)  \ar[r] \ar[d]& T^{n}(C)  \ar[r]\ar[d] & T^{n+1}(A) \ar[d] \ar[r] & T^{n+1}(B)\ar[d]\\
T'^{n}(B)  \ar[r] & T'^{n}(C) \ar[r] & T'^{n+1}(A)  \ar[r] &
T'^{n+1}(B),
}
$$
where the first two vertical arrows are isomorphisms by what we have
just proved, and $T^{n+1}(A)\to T^{n+1}(B)$ is zero by construction.
The five lemma then implies that
$T^{n+1}(A)\to {T'}^{n+1}(A)$ is a monomorphism, \cqfd

\begin{enonce*}{Corollary 6.3}\etiquette[6.3]{XII.cor6.3}
Suppose that $T^{\bigdot}$ and $T'^{\bigdot}$ arise from
cohomological functors (denoted by the same symbols) on the right
derived category $D^{+}(C)$ \cite{sga4.exp.XII.bib.3}, with values in $C'$; the
$T^{i}$ for $i>0$ are not necessarily assumed effaceable. Suppose
that, for $K^{\bigdot}\in D^{+}(C)$, the condition
$H^{i}(K^{\bigdot})=0$ for $i<0$ implies
$T^{i}(K^{\bigdot})={T'}^{i}(K^{\bigdot})=0$ for $i<0$. Let
$K^{\bigdot}\in\Ob D^{+}(C)$ be a bounded-below complex in $C$ such
that $H^{i}(K^{\bigdot})=0$ for $i<0$.
\begin{enumerate}
\renewcommand{\theenumi}{\roman{enumi}}
\renewcommand{\labelenumi}{(\theenumi)}
\item If condition {\rm (a)} of \hyperref[XII.lem6.2]{6.2} is
  satisfied, then the homomorphism
  $\varphi^{i}(K^{\bigdot})=T^{i}(K^{\bigdot})\to
  {T'}^{i}(K^{\bigdot})$ is an isomorphism for $i\leq n$ and a
  monomorphism for $i=n+1$.

\item Let, in addition, $L^{\bigdot}\in\Ob D^{+}(C)$ satisfy
  $H^{i}(L^{\bigdot})=0$ for $i<0$, and let
  $u:L^{\bigdot}\to K^{\bigdot}$ be a morphism in $D^{+}(C)$ that
  induces an injective homomorphism
  $H^{\circ}(u):H^{\circ}(L^{\bigdot})\to
  H^{\circ}(K^{\bigdot})$. Then the homomorphism
  $T^{n+1}(L^{\bigdot})\to {T'}^{n+1}(L^{\bigdot})$ is surjective\pageoriginale if
  and only if the image of ${T'}^{n+1}(L^{\bigdot})$ in
  ${T'}^{n+1}(K^{\bigdot})$ is contained in that of
  $T^{n+1}(K^{\bigdot})$. If $C'=(\Ab)$ and
  $\xi\in {T'}^{n+1}(L^{\bigdot})$, then $\xi$ lies in the image of
  $T^{n+1}(L^{\bigdot})$ if and only if its image in
  ${T'}^{n+1}(K^{\bigdot})$ lies in that of
  $T^{n+1}(K^{\bigdot})$.
\end{enumerate}
\end{enonce*}

Recall that the hypothesis that $T^{\bigdot}$ and $T'^{\bigdot}$
arise from cohomological functors on $D^{+}(C)$ may be spelled out as
follows: they arise from cohomological functors on the abelian
category of bounded-below complexes of $C$ that commute with degree
shifts, send morphisms of complexes homotopic to zero to zero
morphisms in $C'$, and send every morphism
$u:K^{\bigdot}\to L^{\bigdot}$ of complexes that is a
quasi-isomorphism (that is, induces isomorphisms
$H^{i}(K^{\bigdot})\to H^{i}(L^{\bigdot})$) to an isomorphism in
$C'$. In practice, all the cohomological functors from $C$ to $C'$
that one encounters arise in this way.

\medskip
\noindent
\emph{Proof of} \hyperref[XII.cor6.3]{6.3}. We first prove (i). The
assertion is clear if $H^{i}(K^{\bigdot})=0$ for $i\leq n+1$, since
the degree properties imposed on $T^{\bigdot}$ and $T'^{\bigdot}$
then give
$T^{i}(K^{\bigdot})={T'}^{i}(K^{\bigdot})=0$ for $i\leq n+1$.
On the other hand, let $m\geq0$ be an integer such that
$H^{i}(K^{\bigdot})=0$ for $i<m$; such an integer exists, and for
example $m=0$ will do. We proceed by descending induction on $m$, the
assertion having been proved when $m=n+2$. Thus assume the assertion
proved for $m'>m$, and prove it for $m$. By the hypothesis on the
degrees of $K^{\bigdot}$, the complex $K^{\bigdot}$ is
quasi-isomorphic to a complex concentrated in degrees $\geq m$, so
we may suppose that it is itself concentrated in those degrees.
There follows an exact sequence of complexes
$$
0\longrightarrow H^m[-m]\longrightarrow K^{\bigdot}\longrightarrow
K'^{\bigdot}\longrightarrow 0{\quoi},
$$
where $H^m=H^m(K^{\bigdot})$, the notation $[-m]$ denotes a shift of
$-m$ in the degrees of a complex,\pageoriginale and $K'^{\bigdot}$ is defined as
the cokernel. The cohomology objects of $K'^{\bigdot}$ are therefore
those of $K^{\bigdot}$ in degrees $\geq m+1$ and are zero in degrees
$\leq m$. In particular, the induction hypothesis applies to
$K'^{\bigdot}$. The preceding exact sequence gives rise to a
morphism of exact sequences
$$
\xymatrix{
T^{i-1} (K'^{\;\bigdot}) \ar[d] \ar[r] & T^{i-m}(H^m)  \ar[r]\ar[d] &
T^{i}(K^{\bigdot}) \ar[r] \ar[d]& T^{i} (K'^{\;\bigdot})  \ar[r] \ar[d]&
T^{i+1-m}(H^m) \ar[d]\\
T'^{i-1}(K'^{\;\bigdot})  \ar[r] & T'^{i-m}(H^m)  \ar[r] &
T'^{i}(K^{\bigdot})  \ar[r] & T'^{i}(K'^{\;\bigdot})  \ar[r] & T'^{i+1-m}(H^m).
}
$$
Taking $i\leq n$ in this diagram and using the induction hypothesis
and $m\geq0$, we find that the vertical arrows surrounding the
middle vertical arrow are isomorphisms, hence in particular
epimorphisms, and that the rightmost vertical arrow is a
monomorphism. The five lemma therefore shows that the middle
vertical arrow is an epimorphism. Taking $i\leq n+1$, we similarly
find that the vertical arrows surrounding the middle one are
monomorphisms, while the leftmost vertical arrow is an isomorphism,
hence in particular an epimorphism. The five lemma then shows that
the middle vertical arrow is a monomorphism. This proves
\hyperref[XII.cor6.3]{6.3} (i).

To prove the second assertion, \hyperref[XII.cor6.3]{6.3} (ii), we
similarly use the exact sequence
$$
0\longrightarrow L^{\bigdot}\longrightarrow K^{\bigdot}\longrightarrow
K'^{\bigdot}\longrightarrow 0{\quoi},
$$
where $K'^{\bigdot}$ is defined as the mapping cylinder of
$L^{\bigdot}\to K^{\bigdot}$ and again has cohomology in
nonnegative degrees. The usual diagram chase in
$$
\xymatrix{
T^{n}(K^{\bigdot})  \ar[r] \ar[d]& T^{n}(K'^{\;\bigdot})  \ar[r] \ar[d]& T^{n+1}
(L^{\bigdot})  \ar[r]\ar[d] & {T^{n+1}(K^{\bigdot})}  \ar[r]
\ar[d]&{T^{n+1}(K'^{\;\bigdot})} \ar[d]\\
T'^{n}(K^{\bigdot}) \ar[r] & T'^{n}(K'^{\;\bigdot})  \ar[r] &
T'^{n+1}(L^{\bigdot}) \ar[r] & T'^{n+1}(K^{\bigdot}) \ar[r] & T'^{n+1}(K'^{\bigdot})
}
$$
gives\pageoriginale the last assertion of
\hyperref[XII.cor6.3]{6.3}, since, by the first part already proved,
the first two vertical arrows are isomorphisms and the last two are
monomorphisms. This completes the proof of
\hyperref[XII.cor6.3]{6.3}.

\begin{enonce*}{Lemma 6.4}\etiquette[6.4]{XII.lem6.4}
Theorem \hyperref[XII.thm5.1]{5.1} holds if $f$ is finite.
\end{enonce*}

Indeed, this follows at once from VIII \hyperref[VIII.prop5.5]{5.5} and \hyperref[VIII.rem5.8]{5.8}.

\begin{enonce*}{Proposition 6.5}\etiquette[6.5]{XII.prop6.5}
Let $h:Y\to X$ be a quasi-compact and quasi-separated morphism of
schemes. In {\rm (i)} below, $I$ denotes a fixed set having at least
two elements; in {\rm (ii)} and {\rm (iii)}, $\LL$ denotes a
nonempty subset of the set $\PP$ of prime numbers; and in {\rm (iii)},
$n$ denotes an integer.
\begin{enumerate}
\renewcommand{\theenumi}{\roman{enumi}}
\renewcommand{\labelenumi}{(\theenumi)}
\item The following conditions are equivalent:
\begin{itemize}
\item[a)] For every sheaf of sets $F$ on $X$, the canonical map
$$
H^{\circ}(X,F)\longrightarrow H^{\circ}(Y,h^{\ast}(F))
$$
is injective (respectively, bijective).

\item[b)] For every finite morphism $X'\to X$, let
$h':Y'\to X'$ be the morphism obtained from $h$ by the base change
$X'\to X$. The canonical map
$$
H^{\circ}(X', I_{X'})\longrightarrow H^{\circ}(Y',I_{Y'})
$$
induced by $h'$ is injective (respectively, bijective).

\item[${\rm b')}$] With the notation of {\rm (b)}, the map
$U\mapsto {h'}^{-1}(U)$ from the set of subsets of $X'$ that are both
open and closed to the set of subsets of $Y'$ that are both open and
closed is injective (respectively, bijective).

\item[${\rm b'')}$] (If\pageoriginale $X$ is locally Noetherian.)
With the notation of {\rm (b)}, the induced map on sets of connected
components
$$
\Pi_{\circ}(h'):\Pi_{\circ}(Y')\longrightarrow \Pi_{\circ}(X')
$$
is surjective (respectively, bijective). Equivalently, if $X'$ is
nonempty (respectively, connected and nonempty), then the same is true
of $Y'$.
\end{itemize}

Moreover, if these conditions hold, then for every sheaf of groups
$F$ on $X$, the functor $P\mapsto h^{\ast}(P)$ from the category of
torsors (homogeneous principal bundles) under $F$ to the category of
torsors under $h^{\ast}(F)$ is faithful (respectively, fully
faithful). In particular, in the latter case, the canonical map
$$
H^1(X,F)\longrightarrow H^1(Y,h^{\ast}(F))
$$
is injective. Finally, under the same hypotheses, the functor
$X'\mapsto Y'=\fprod{X'}{Y}{X}$ from the category of étale coverings
of $X$ to the category of étale coverings of $Y$ is faithful
(respectively, fully faithful).

\item The following conditions are equivalent:
\begin{itemize}
\item[a)] For every sheaf $F$ of ind-$\LL$-finite groups on $X$, the
canonical map
$$
H^i(X,F)\longrightarrow H^i(Y,h^{\ast}(F))
$$
is bijective for $i=0,1$.

\item[b)] For every finite morphism $X'\to X$, let
$h':Y'\to X'$ be the morphism obtained from $h$ by base change, and
let $G$ be any ordinary finite $\LL$-group. The canonical map
$$
H^i(X',G_{X'})\longrightarrow H^i(Y',G_{Y'})
$$
is bijective for $i=0,1$.

\item[${\rm b')}$] (If\pageoriginale $\LL=\PP$.) With the notation
of {\rm (b)}, the inverse-image functor under $h'$ induces an
equivalence from the category of étale coverings of $X'$ to the
category of étale coverings of $Y'$.

\item[${\rm b'')}$] (If $X$ is Noetherian and $\LL=\PP$.) With the
notation of {\rm (b)}, if $X'$ is nonempty then so is $Y'$. If $y'$ is
a geometric point of $Y'$ and $x'$ is its image in $X'$, the canonical
map
$$
\Pi_i(Y',y')\longrightarrow \Pi_i(X',x')\quav\mbox{ for }i=0,1{\quoi},
$$
is bijective.
\end{itemize}

\item The following conditions are equivalent:
\begin{itemize}
\item[a)] For every sheaf $F$ of $\LL$-torsion groups on $X$, the
homomorphism
$$
H^{i}(X,F)\longrightarrow H^{i}(Y,h^{\ast}(F))
$$
is an isomorphism for $i\leq n$ and a monomorphism for $i=n+1$.

\item[${\rm a')}$] (If $n\geq -1$.) For every $F$ as above, the
homomorphism in question is injective for $i=0$ and surjective for
$i\leq n$.

\item[b)] (If $n\geq -1$.) For every finite morphism $X'\to X$, let
$h':Y'\to X'$ be the morphism obtained from $h$ by the base change
$X'\to X$. Then for every $\ell\in\LL$ and $\nu>0$, the canonical
homomorphism
$$
H^{i}(X',(\ZZ/\ell^{\nu}\ZZ)_{X'})\longrightarrow
H^{i}(Y',(\ZZ/\ell^{\nu}\ZZ)_{Y'})
$$
is injective for $i=0$ and surjective for $i\leq n$.
\end{itemize}
\end{enumerate}
\end{enonce*}

In addition, we have the following supplement to the preceding
statement:

\begin{enonce*}{Corollary 6.6}\etiquette[6.6]{XII.cor6.6}
\begin{enumerate}
\renewcommand{\theenumi}{\roman{enumi}}
\renewcommand{\labelenumi}{(\theenumi)}
\item Suppose that the non-parenthesized condition in
\hyperref[XII.prop6.5]{6.5} {\rm (i) (a)} holds. Let $F$ be a sheaf of
sets on $X$, and let $\xi\in H^{\circ}(Y,h^{\ast}(F))$. Suppose\pageoriginale that
there exists a monomorphism of sheaves of sets $F\to G$ such that the
image of $\xi$ in $H^{\circ}(Y,h^{\ast}(G))$ lies in the image of
$H^{\circ}(X,G)$. Then $\xi$ lies in the image of
$H^{\circ}(X,F)$.

\item Suppose that the parenthesized condition in
\hyperref[XII.prop6.5]{6.5} {\rm (i) (a)} holds. Let $F$ be a sheaf of
groups on $X$, and let $\xi\in H^1(Y,h^{\ast}(F))$. Suppose that there
exists a monomorphism of sheaves of groups $F\to G$ such that the
image of $\xi$ in $H^1(Y,h^{\ast}(G))$ lies in the image of
$H^1(X,G)$. Then $\xi$ lies in the image of $H^1(X,F)$.

\item Suppose that condition
\hyperref[XII.prop6.5]{6.5} {\rm (iii) (a)} holds for a fixed value of
$n$. Let $F$ be a sheaf of $\LL$-torsion groups on $X$, and let
$\xi\in H^{n+1}(Y,h^{\ast}(F))$. Suppose that there exists a
monomorphism $F\to G$ of sheaves of $\LL$-torsion groups such that
the image of $\xi$ in $H^{n+1}(Y,h^{\ast}(G))$ lies in the image of
$H^{n+1}(X,G)$. Then $\xi$ lies in the image of $H^{n+1}(X,F)$.
\end{enumerate}
\end{enonce*}

\noindent
\emph{Proof of \hyperref[XII.prop6.5]{6.5} and
\hyperref[XII.cor6.6]{6.6}.}
We first prove \hyperref[XII.cor6.6]{6.6}. The case

% END INCLUDED SGA4_English_current_buildable_source_20260906/source/12/drafts/SGA4_XII_LINES_785_1200_English.texfrag


% BEGIN INCLUDED SGA4_English_current_buildable_source_20260906/source/12/drafts/SGA4_XII_LINES_1201_1600_English.texfrag
of \hyperref[XII.cor6.6]{6.6} (iii) is a special case of the last assertion of
\hyperref[XII.cor6.3]{6.3}, applied when $C$ is the category of abelian
$\LL$-torsion sheaves on $X$, $C'$ is the category $(\Ab)$, and
$T^{i}(F)=H^{i}(X,F)$, ${T'}^{i}(F)=H^{i}(Y,h^{\ast}(F))$ (note that
$T'^{\bigdot}$ is indeed a cohomological functor in $F$, since the functor
$h^{\ast}$ is exact). Of course, the use of complexes as in
\hyperref[XII.cor6.3]{6.3} is unnecessary here (it will be convenient later,
in \hyperref[XII.prop6.8]{6.8}); we use only the last assertion of
\hyperref[XII.cor6.3]{6.3}.

To prove (i), consider the amalgamated sum $H=G\amalg_F G$, the colimit in
the category of sheaves of sets on $X$ of the diagram
$$
\xymatrix@R=.5cm{
  & G\\
F \ar[ru] \ar[rd]&\\
  & G
}
$$
The\pageoriginale usual exactness properties of the topos of sheaves of sets
show, since $F\rightarrow G$ is a monomorphism, that there is an
\emph{exact} diagram of sheaves of sets on $X$
$$
F\to G\rightrightarrows H
$$
(this may be expressed by saying that every monomorphism in a topos is
effective, and may be proved by reducing, as usual, to the category of
sets). Since the functors $H^{0}$ and $h^{\ast}$ are left exact, we obtain a
morphism of exact diagrams of sets
$$
\xymatrix@R=1.25cm{
H^{0}(X,F) \ar[r]\ar[d]& H^{0}(X,G)  \ar[d]\ar@<2pt>[r]\ar@<-2pt>[r]& H^{0}(X,H)\ar[d]\\
H^{0}(Y,h^{\ast}(F)) \ar[r]& H^{0}(Y,h^{\ast}(G)) \ar@<2pt>[r]\ar@<-2pt>[r]& H^{0}(Y,h^{\ast}(H)),
}
$$
in which the vertical arrows are injective. An immediate diagram chase then
proves assertion \hyperref[XII.cor6.6]{6.6} (i).

To prove \hyperref[XII.cor6.6]{6.6} (ii), we first prove the second assertion
of \hyperref[XII.prop6.5]{6.5} (i), namely that, under condition a) of
\hyperref[XII.prop6.5]{6.5} (i), for every sheaf of groups $F$ on $X$, the
functor $P\mapsto h^{\ast}(P)$ from the category of torsors under $F$ to the
category of torsors under $h^{\ast}(F)$ is faithful (\resp fully faithful),
and hence, in the respective case, induces an injection on $H^1$. For this,
if $P$ and $P'$ are two torsors under $F$, denote by
$\sheaf{\Isom}_F(P,P')$ the sheaf of $F$-isomorphisms from $P$ to $P'$.
It is immediate that the formation of this sheaf commutes with every base
extension, and in particular with base extension by $h$. (This assertion is
valid for every morphism of sites.) Applying the hypothesis on the effect of
$h^{\ast}$ on $H^{0}$ to the sheaf $\sheaf{\Isom}_F(P,P')$, we find that the
map
$$
\Isom_F(P,P')\longrightarrow
\Isom_F(h^{\ast}(P),h^{\ast}(P'))
$$
is injective (\resp bijective), which says that the functor
$P\mapsto h^{\ast}(P)$ is\pageoriginale faithful (\resp fully faithful).
(Note that every morphism in the category of torsors under a sheaf of groups
is an isomorphism.)

On the other hand, note that if $F$ is a subsheaf of groups of a sheaf of
groups $G$, then giving a torsor $P$ under $F$ amounts to giving a torsor
$Q$ under $G$ (namely the one obtained from $P$ by extending the structure
group along $F\rightarrow G$), \emph{together with} a section of $Q/F$
(which is indeed interpreted as a reduction of the structure group from
$G$ to $F$). With the notation of \hyperref[XII.cor6.6]{6.6} (ii), the image
of $\xi$ in $H^1(Y,h^{\ast}(G))$ comes from an element of $H^1(X,G)$; the
latter is represented by a torsor $Q$ under $G$, so that the image of $\xi$
is the class of the torsor $h^{\ast}(Q)$ under $h^{\ast}(G)$. The fact that
this class comes from $\xi\in H^1(Y,h^{\ast}(F))$ means explicitly that
$\xi$ is the class of the $h^{\ast}(F)$-torsor defined by $h^{\ast}(Q)$ and
a suitable section of $h^{\ast}(Q)/h^{\ast}(F)$. Since the functor
$h^{\ast}$ is right exact, this last sheaf is none other than
$h^{\ast}(Q/F)$; and since $h^{\ast}$ induces a bijection on the $H^{0}$ of
sheaves of sets, the section in question of $h^{\ast}(Q)/h^{\ast}(F)$ comes
from a section of $Q/F$. This last section defines an $F$-torsor by reduction
of the structure group, and the class of this torsor in $H^1(X,F)$ plainly
has $\xi$ as its image in $H^1(Y,h^{\ast}(F))$. This proves
\hyperref[XII.cor6.6]{6.6} (ii) and completes the proof of
\hyperref[XII.cor6.6]{6.6}.

We now prove \hyperref[XII.prop6.5]{6.5}, beginning with
\hyperref[XII.prop6.5]{6.5} (iii). Clearly a) implies ${\rm a'}$), and the
converse implication is a special case of \hyperref[XII.lem6.2]{6.2}
(again applied when $C$ is the category of $\LL$-torsion sheaves on $X$ and
$C'$ the category of abelian groups), provided that one proves that the
functors $H^{i}(X,F)$, for $i>0$, are effaceable in $C$. This presents no
difficulty; alternatively, we can avoid proving that result by\pageoriginale
noting that $F$ is the filtered colimit of its subsheaves $F_i$ for which
there exists an integer $n>0$ annihilating $F_i$
(\hyperrefext[IX][IX.def1.1]{1.1}). Then $h^{\ast}(F)$ is the colimit of the
$h^{\ast}(F_i)$, and, in view of the commutation of cohomology with filtered
colimits (VII \hyperref[VII.prop3.3]{3.3}), which follows from the quasi-compactness and
quasi-separatedness hypotheses on $X$ and $Y$ (hypotheses not used for
\hyperref[XII.cor6.6]{6.6}), we are reduced to proving condition a) for the
$F_i$. Thus we are reduced to the case in which there exists an integer
$m>0$ such that $mF=0$. We may plainly suppose that every prime divisor of
$m$ belongs to $\LL$, and can then take for $C$ the category of sheaves of
$\ZZ/m\ZZ$-modules on $X$. In this category it is known that the functors
$H^{i}(X,-)$ are effaceable for $i\geq1$. This proves the equivalence of
conditions a) and ${\rm a'}$) in \hyperref[XII.prop6.5]{6.5} (iii).

On the other hand, apply a) to the direct image of a $\LL$-torsion sheaf
$F'$ on a scheme $X'$ finite over $X$; this direct image is a
$\LL$-torsion sheaf by \hyperrefext[IX][IX.prop1.2]{1.2} (v) and
\hyperrefext[IX][IX.prop1.6]{1.6} (iii). Using
\hyperref[XII.lem6.4]{6.4} and VIII \hyperref[VIII.prop5.5]{5.5}, we find that the homomorphism
$$
H^{i}(X',F')\longrightarrow H^{i}(Y',{h'}^{\ast}(F'))
$$
is an isomorphism for $i\leq n$ and a monomorphism for $i=n+1$, whence, a
fortiori, condition b). It remains to prove that this condition implies
${\rm a'}$). For this, use the fact (\hyperrefext[IX][IX.cor2.7.2]{2.7.2})
that every $\LL$-torsion sheaf $F$ on $X$ is the filtered colimit of
constructible $\LL$-torsion sheaves. The preceding colimit argument then
shows that, in ${\rm a'}$), it suffices to consider constructible $F$. By
\hyperrefext[IX][IX.prop2.14]{2.14}, there then exist finitely many finite
morphisms $p_i:X'_i\rightarrow X$, constant $\LL$-torsion sheaves
$G_{X'_i}$ on the $X'_i$, and an injective homomorphism
$$
F\longrightarrow \prod_{i}p_{i\ast}(G_{X'_i}).
$$
Now\pageoriginale if $n=-1$, that is, if the desired conclusion reduces to
the assertion that
$$
H^{0}(X,F)\longrightarrow H^{0}(Y,h^{\ast}(F))
$$
is injective, then in this question we may at once replace $F$ by a sheaf
$G$ into which $F$ embeds, and hence may replace $F$ by the product of the
$p_{i\ast}(G_{X'_i})$. In view of \hyperref[XII.lem6.4]{6.4}, we are thus
reduced to proving the injectivity assertion after replacing
$h:Y\rightarrow X$ by $h':Y'\rightarrow X'$ and $F$ by a constant sheaf on
$X'$ whose value is an ordinary finite $\LL$-torsion group $G$. But $G$ is
isomorphic to a direct sum of groups of the form
$\ZZ/\ell^{\nu}\ZZ$, with $\ell\in\LL$. We are therefore indeed reduced
to verifying b).

When $n\geq0$, we proceed similarly. By induction on $n$, we may suppose
that the desired conclusion (and hence a)) has been proved for $n-1$ in
place of $n$. Applying \hyperref[XII.cor6.6]{6.6} (iii), we see that we may
again replace $F$ by a sheaf $G$ into which $F$ embeds, and we again take
$G=p_{\ast}(G_{X'})$. In view of \hyperref[XII.lem6.4]{6.4} and the
vanishing of $R^{i}p_{\ast}$ for $i\geq1$ when $p$ is finite (VIII \hyperref[VIII.prop5.5]{5.5}),
we are again reduced to proving surjectivity for $H^n$ after replacing
$h:Y\rightarrow X$ by $h':Y'\rightarrow X'$ and $F$ by a constant sheaf
on $X'$ whose value is an ordinary finite $\LL$-torsion group. Injectivity
is already known by the induction hypothesis and the implication
${\rm a')}\Rightarrow{\rm a)}$; by devissage on $G$, we are thus again
reduced to the case in which $G$ is of the form
$\ZZ/\ell^{\nu}\ZZ$, namely the case considered in condition b) of
\hyperref[XII.prop6.5]{6.5} (iii). This proves case (iii) of
\hyperref[XII.prop6.5]{6.5}.

Cases (i) and (ii) are proved in exactly the same way as regards the
equivalence of conditions a) and b) in \hyperref[XII.prop6.5]{6.5} (i)
\resp \hyperref[XII.prop6.5]{6.5} (ii): one uses\pageoriginale
\hyperrefext[IX][IX.cor2.7.2]{2.7.2} and the commutation of $H^{0}$ and
$H^1$ with filtered colimits of sheaves (VII \hyperref[VII.rem5.14]{5.14}) to reduce to the case in
which $F$ is constructible, and then uses \hyperrefext[IX][IX.prop2.14]{2.14},
\hyperref[XII.cor6.6]{6.6} (i), (ii), and VIII \hyperref[VIII.rem5.8]{5.8} to reduce to the special
case considered in b). Moreover, in (i), the equivalence of conditions b),
${\rm b'}$), and ${\rm b''}$) is trivial and is recorded only for
reference, as is the fact that they imply the faithfulness (\resp full
faithfulness) of the functor $X'\mapsto Y'$ from the category of etale
coverings of $X$ to the category of etale coverings of $Y$. The same is true
of the equivalence of conditions b), ${\rm b'}$), and ${\rm b''}$) in
\hyperref[XII.prop6.5]{6.5} (ii), which were added for good measure. The
proof of \hyperref[XII.prop6.5]{6.5} is now complete.

\begin{enonce*}[remark]{Remark 6.7}\etiquette[6.7]{XII.rem6.7}
\footnote{The editor recommends omitting these remarks, as well as
\hyperref[XII.rems6.13]{6.13}; they were surreptitiously introduced
(together with various other more or less felicitous changes to the
original text) by an irreverent collaborator.}
\begin{enumerate}
\renewcommand{\labelenumi}{\theenumi)}
\item The arguments proving \hyperref[XII.prop6.5]{6.5} and
  \hyperref[XII.cor6.6]{6.6} are plainly very general and essentially
trivial, and it would be useful to extract them as abstract lemmas of the
same kind as \hyperref[XII.lem6.2]{6.2} (with the sheaves
$p_{\ast}(G_{X'})$ playing the role of cogenerators of the category $C$ in
which one works). We leave this pleasant exercise to the reader, merely
pointing out that essentially the same statement could be given for
functors such as $R^{i}f_{\ast}$ in place of the functors $H^{i}$. The same
remarks apply to \hyperref[XII.prop6.8]{6.8} and
\hyperref[XII.prop6.11]{6.11} below.

\item When $X$ is noetherian, the arguments given show that in conditions
b) and their variants in \hyperref[XII.prop6.5]{6.5} (i), (ii), and (iii),
one may restrict to \emph{integral} $X'$; when $X$ is universally Japanese,
one may even take them to be integral and normal. This is also possible\pageoriginale
without the Japanese hypothesis on $X$, provided that \emph{integral}
morphisms $p:X'\rightarrow X$ are used in place of finite morphisms. None of
these variants will be used later in the seminar.

\item The proof of \hyperref[XII.cor6.6]{6.6} does not use the
quasi-compactness and quasi-separatedness hypothesis on $X$. The same is
true of the fact that in (i), (ii), and (iii), condition a) implies the
other conditions, and of the mutual equivalence of conditions b),
${\rm b'}$), and ${\rm b''}$). Note also that the form b) of the conditions
in \hyperref[XII.prop6.5]{6.5} (i) shows that, for $\LL\neq\phi$, they are
special cases (when $X$ is quasi-compact and quasi-separated) of
\hyperref[XII.prop6.5]{6.5} (iii), for $n=-1$ \resp $n=0$. This remains
true, moreover (if one restricts to the conditions in form a)), without the
stated restriction on $X$. Indeed, for every sheaf of sets $F$ and every
prime number $\ell$, one can find two abelian $\ell$-torsion sheaves $G$
and $H$, and a homomorphism $u:G\rightarrow H$ of the underlying sheaves of
sets, such that $F$ is isomorphic to the inverse image of the zero section
(one takes for $G$ the free $\ZZ/\ell\ZZ$-module generated by $F$, and for
$H$ the constant sheaf $\ZZ/\ell\ZZ$ on $X$). The same argument shows that
(without any quasi-compactness or quasi-separatedness restriction)
\hyperref[XII.prop6.5]{6.5} (ii) a) (where it even suffices to take $i=0$)
implies the respective condition \hyperref[XII.prop6.5]{6.5} (i) a). Since
the latter already implies the injectivity of
$H^1(X,F)\rightarrow H^1(Y,h^{\ast}(F))$ for every sheaf of groups $F$ on
$X$, in the statement of \hyperref[XII.prop6.5]{6.5} (ii) a) it therefore
suffices, for $i=1$, to require the \emph{surjectivity} of
$H^1(X,F)\rightarrow H^1(Y,h^{\ast}(F))$. Unfortunately,
\hyperref[XII.prop6.5]{6.5} (ii) cannot be regarded as a special case of
\hyperref[XII.prop6.5]{6.5} (iii), which often forces us to repeat in the
noncommutative case an argument already carried out in essence in the
commutative case. Note also that, apart from \hyperref[XII.cor6.10]{6.10},
the results that follow are stated so as to remain valid
without\pageoriginale any quasi-compactness or quasi-separatedness
hypothesis.
\end{enumerate}
\end{enonce*}

\begin{enonce*}{Proposition 6.8}\etiquette[6.8]{XII.prop6.8}
(Descent lemma). With the general notation of \hyperref[XII.prop6.5]{6.5},
suppose in addition that a morphism $p:\bar{X}\rightarrow X$ is given. Let
$\bar{h}:\bar{Y}\rightarrow\bar{X}$ denote the morphism obtained from $h$
by the base change $p$, and let $p_Y:\bar{Y}\rightarrow Y$ denote the
canonical morphism. We suppose either that $p$ is surjective, or that $h$
is a closed immersion and $X=p(\bar{X})\cup h(Y)$.
\begin{enumerate}
\renewcommand{\theenumi}{\roman{enumi}}
\renewcommand{\labelenumi}{(\theenumi)}
\item Suppose that $\bar{h}:\bar{Y}\rightarrow\bar{X}$ satisfies the
injectivity (\resp bijectivity) condition of
\hyperref[XII.prop6.5]{6.5} {\rm(i) a)}. In the respective case, suppose in
addition that the base-change morphism
\begin{equation*}
h^{\ast}(p_{\ast}(\bar{F}))\longrightarrow
p_{Y\ast}(\bar{h}^{\ast}(\bar{F}))\tag{$\ast$}
\end{equation*}
is injective for every sheaf of sets $\bar{F}$ on $\bar{X}$. Then the
morphism $h:Y\rightarrow X$ also satisfies the injectivity (\resp
bijectivity) condition of \hyperref[XII.prop6.5]{6.5} {\rm(i) a)}.

\item Suppose that $\bar{h}:\bar{Y}\rightarrow\bar{X}$ satisfies condition
\hyperref[XII.prop6.5]{6.5} {\rm(ii) a)}, and that for every sheaf of
ind-$\LL$-groups $\bar{F}$ on $\bar{X}$ the base-change morphism $(\ast)$
is bijective. Then $h:Y\rightarrow X$ also satisfies condition
\hyperref[XII.prop6.5]{6.5} {\rm(ii) a)}.

\item Suppose that $\bar{h}:\bar{Y}\rightarrow\bar{X}$ satisfies condition
\hyperref[XII.prop6.5]{6.5} {\rm(iii) a)}, and that for every abelian
$\LL$-torsion sheaf $\bar{F}$ on $\bar{X}$ the base-change homomorphism
$$
h^{\ast}(R^{i}p_{\ast}(\bar{F}))\longrightarrow
R^{i}p_{Y\ast}(\bar{h}^{\ast}(\bar{F}))
$$
is bijective for $i\leq n-1$ and injective for $i=n$, where $n\geq1$ is a
given integer. Then $h:Y\rightarrow X$ also satisfies condition
\hyperref[XII.prop6.5]{6.5} {\rm(iii) a)}.
\end{enumerate}
\end{enonce*}

\noindent
\emph{Proof of} \hyperref[XII.prop6.8]{6.8}.\pageoriginale In the case where
$p$ is not assumed surjective, but $h$ is a closed immersion and
$X=p(\bar{X})\cup h(Y)$, consider the disjoint-union scheme $\bar{X}'$ of
$\bar{X}$ and $Y$, and the morphism $p':\bar{X}'\rightarrow X$ induced by
$p$ and $h$. Then $p'$ is surjective; moreover, in each of the three cases
(i), (ii), and (iii), it is immediate that the hypotheses on the pair
$(h,p)$ remain satisfied for the pair $(h,p')$. We are therefore reduced to
the case in which $p$ is surjective, which we assume from now on. In this
case, for every sheaf $F$ on $X$, the canonical homomorphism
$$
F\longrightarrow p_{\ast}(p^{\ast}(F))
$$
is injective. In case (i), to verify the injectivity of
$H^{0}(X,F)\rightarrow H^{0}(Y,h^{\ast}(F))$, we are immediately reduced
to the case in which $F$ is replaced by $p_{\ast}(\bar{F})$, where
$\bar{F}=p^{\ast}(F)$. Consider the commutative diagram
$$
\xymatrix@C=2.75cm@R=.75cm{
**[r]H^{0}(X,p_{\ast}(\bar{F})) \ar[dd]^{\text{\rotatebox{90}{$\backsim$}}} \ar[r]& **[r]H^{0}(Y,h^{\ast}(p_{\ast}(\bar{F}))) \ar[d]\\
                           &**[r]H^{0}(Y,p_{Y\ast}(\bar{h}^{\ast}(\bar{F})))\ar[d]_{\text{\rotatebox{90}{$\backsim$}}}\\
**[r]H^{0}(\bar{X},\bar{F})  \ar[r] &
                           **[r]H^{0}(\bar{Y},\bar{h}^{\ast}(\bar{F})),
}
$$
where the left vertical arrow and the second right vertical arrow are the
canonical isomorphisms, and the first right vertical arrow is induced by
the base-change homomorphism for $h$. By hypothesis the second horizontal
arrow is injective, and it follows immediately that the first is as well;
this proves the non-respective assertion of (i). For the respective
assertion, suppose that the second horizontal arrow is surjective. Since the
base-change homomorphism is a monomorphism by hypothesis, the first right
vertical arrow is a monomorphism, and hence so is the composite of the
two\pageoriginale right vertical arrows. It follows immediately that the
first horizontal arrow is also surjective, completing the proof of (i).

To prove (ii), what precedes and \hyperref[XII.prop6.5]{6.5} (i) reduce us
to proving the surjectivity of
$H^1(X,F)\rightarrow H^1(Y,h^{\ast}(F))$. Applying
\hyperref[XII.cor6.6]{6.6} (ii), we are further reduced to proving it for
$F$ of the form $p_{\ast}(\bar{F})$, where $\bar{F}$ is a sheaf of
ind-$\LL$-groups on $\bar{X}$ (\hyperrefext[IX][IX.prop1.6]{1.6} (ii)).
Writing the exact sequences for $H^1$ (variants of
\hyperref[XII.prop3.2]{3.2}) for the morphisms $p$ and $p_Y$, we obtain a
morphism of noncommutative exact sequences:
$$
\xymatrix@R=1.5cm{
H^{1}(X,p_{\ast}(\bar{F})) \ar[r] \ar[d] & H^{1}(\bar{X},\bar{F}) \ar[r]\ar[d] &
H^{0}(X,R^{1}p_{\ast}(\bar{F})) \ar[d]&\\
H^{1}(Y,p_{Y\ast}(\bar{h}^{\ast}(\bar{F}))) \ar[r] &
H^{1}(\bar{Y},\bar{h}^{\ast}(\bar{F})) \ar[r] &
H^{0}(Y,R^{1}p_{Y\ast}(\bar{h}^{\ast}(\bar{F}))),
}
$$
where in each row the first arrow is injective and identifies the first term
with the inverse image, under the second arrow, of the distinguished point
of the third term. By the hypothesis on $\bar{h}$, the middle vertical arrow
is bijective. The two outer vertical arrows factor, respectively, through
$$
H^1(Y,h^{\ast}(p_{\ast}(\bar{F})))
\quad\hbox{and}\quad
H^{0}(Y,h^{\ast}(R^1p_{\ast}(\bar{F}))),
$$
by composing the homomorphisms $H^{i}(X,-)\rightarrow H^{i}(Y,h^{\ast}(-))$
(for $i=1$ and $i=0$, respectively) with the homomorphisms obtained by
applying $H^{i}(Y,-)$ (for those same values of $i$) to the base-change
homomorphisms
$$
h^{\ast}(p_{\ast}(\bar{F}))\longrightarrow
p_{Y\ast}(\bar{h}^{\ast}(\bar{F}))
$$
and, respectively,
$$
h^{\ast}(R^1p_{\ast}(\bar{F}))\longrightarrow
R^1p_{Y\ast}(\bar{h}^{\ast}(\bar{F})).
$$
Since the first of these arrows is an isomorphism by hypothesis, the first
vertical arrow in the diagram above identifies with the homomorphism
$H^1(X,-)\rightarrow H^1(Y,h^{\ast}(-))$ that we wish to study. To prove
that it is bijective, it therefore remains to prove that the last vertical
arrow in the diagram is injective. By what we have just said, it is the
composite of two maps,\pageoriginale the first of which is a canonical
homomorphism $H^{0}(X,-)\rightarrow H^{0}(Y,h^{\ast}(-))$, and hence is
bijective, while the second is injective because it is obtained by applying
the left exact functor $H^{0}(Y,-)$ to the monomorphism
$$
h^{\ast}(R^1p_{\ast}(\bar{F}))\longrightarrow
R^1p_{Y\ast}(\bar{h}^{\ast}(\bar{F})).
$$
The latter fact follows readily from the bijectivity hypothesis for the
base-change homomorphism in the preceding dimension $0$, for example by
using the second assertion of \hyperref[XII.prop6.5]{6.5} (i) and the usual
calculation of the fibers of the higher direct images (VIII \hyperref[VIII.rem5.3]{5.3}). This
proves \hyperref[XII.prop6.8]{6.8} (ii).

It remains to prove \hyperref[XII.prop6.8]{6.8} (iii). For this, rather than
using the injection $F\rightarrow p_{\ast}(p^{\ast}(F))$ directly as above,
it is more convenient to use the homomorphism
$$
F\longrightarrow Rp_{\ast}(\bar{F})
$$
in the right derived category $D^{+}(C)$, where $C$ denotes the category of
abelian sheaves on $X$. Recall that, by definition,
$Rp_{\ast}(\bar{F})$ is the complex $p_{\ast}(C(\bar{F}))$, where
$C(\bar{F})$ is an injective resolution of $\bar{F}$ in the category of
abelian sheaves on $\bar{X}$. If $K^{\bigdot}$ is a complex on $X$, we
denote by $H^{i}(X,K^{\bigdot})$ the hypercohomology groups of $X$ with
coefficients in $K^{\bigdot}$, and use analogous notation on $Y$. By
\hyperref[XII.cor6.3]{6.3}, we are reduced to proving the surjectivity of
\begin{equation*}
H^n(X,Rp_{\ast}(\bar{F}))\longrightarrow
H^n(Y,h^{\ast}(Rp_{\ast}(\bar{F}))).\tag{1}
\end{equation*}\etiquette[1]{XII.eq1}
By hypothesis, the base-change homomorphism
\begin{equation*}
h^{\ast}(Rp_{\ast}(\bar{F}))\longrightarrow
Rp_{Y\ast}(\bar{h}^{\ast}(\bar{F}))\tag{2}
\end{equation*}\etiquette[2]{XII.eq2}
induces an isomorphism on the cohomology sheaves in degrees $\leq n-1$ and
a monomorphism in degree $n$. Equivalently, the ``mapping cylinder''
\pageoriginale $K^{\bigdot}$ of the preceding homomorphism has nonzero
cohomology sheaves only in degrees $\geq n$. Hence
$H^{i}(Y,K^{\bigdot})=0$ for $i<n$, and the cohomology exact sequence shows
that the homomorphisms
\begin{equation*}
H^{i}(Y,h^{\ast}(Rp_{\ast}(\bar{F})))\longrightarrow
H^{i}(Y,Rp_{Y\ast}(\bar{h}^{\ast}(\bar{F})))\tag{3}
\end{equation*}\etiquette[3]{XII.eq3}
induced by the homomorphism above are isomorphisms for $i\leq n-1$ and a
monomorphism for $i=n$. Consider the composite of (1) and (3):
\begin{equation*}
H^n(X,Rp_{\ast}(\bar{F}))\longrightarrow
H^n(Y,h^{\ast}(Rp_{\ast}(\bar{F})))\longrightarrow
H^n(Y,Rp_{Y\ast}(\bar{h}^{\ast}(\bar{F})))\tag{4}
\end{equation*}\etiquette[4]{XII.eq4}

% END INCLUDED SGA4_English_current_buildable_source_20260906/source/12/drafts/SGA4_XII_LINES_1201_1600_English.texfrag


% BEGIN INCLUDED SGA4_English_current_buildable_source_20260906/source/12/drafts/SGA4_XII_LINES_1601_1924_English.texfrag
The two outer terms are respectively isomorphic, by definition, to
$H^n(\bar{X},\bar{F})$ and
$H^n(\bar{Y},\bar{h}^{\ast}(\bar{F}))$, and the composite in (4) is
precisely the homomorphism induced by $\bar{h}$, which is bijective by the
hypothesis on $\bar{h}$. Since the second arrow in (4) is injective by what
we have just seen, the first is also bijective. This completes the proof of
\hyperref[XII.prop6.8]{6.8}.

Moreover, the preceding proof immediately gives the following slightly
more precise and more general result:

\begin{enonce*}{Corollary 6.9}\etiquette[6.9]{XII.cor6.9}
The notation is that of \hyperref[XII.prop6.8]{6.8}. Suppose in addition
that $h$ satisfies the unparenthesized condition in
\hyperref[XII.prop6.5]{6.5} \textup{(i)} (respectively, the parenthesized
condition in \hyperref[XII.prop6.5]{6.5} \textup{(i)}, respectively, the
condition in \hyperref[XII.prop6.5]{6.5} \textup{(iii)}, with $n\geq -1$), and
that, for every sheaf of sets $\bar{F}$ on $\bar{X}$, the base-change
homomorphism
$$
h^{\ast}(p_{\ast}(\bar{F}))\longrightarrow
p_{Y\ast}(\bar{h}^{\ast}(\bar{F}))
$$
is injective (respectively, bijective; respectively, that for every sheaf
of $\LL$-torsion groups $\bar{F}$ on $\bar{X}$, the base-change
homomorphism
$$
h^{\ast}R^{i}p_{\ast}(\bar{F})\longrightarrow
R^{i}p_{Y\ast}(\bar{h}^{\ast}(\bar{F}))
$$
is bijective for $i\leq n$ and injective for $i=n+1$). Let $F$ be a sheaf
of sets (respectively, of ind-$\LL$-groups, respectively, a sheaf of
$\LL$-torsion groups) on $X$, and let
$\xi\in H^{n+1}(Y,h^{\ast}(F))$ (where $n=-1$ in case\pageoriginale \textup{(i)} and $n=0$
in case \textup{(ii)}). Put $\bar{F}=p^{\ast}(F)$. Then $\xi$ lies in the image
of $H^{n+1}(X,F)$ if and only if its inverse image in
$$
H^{n+1}(\bar{Y},p_Y^{\ast}(h^{\ast}(F)))
=H^{n+1}(\bar{Y},\bar{h}^{\ast}(\bar{F}))
$$
lies in the image of $H^{n+1}(\bar{X},\bar{F})$.
\end{enonce*}

We shall use this sharpened form of \hyperref[XII.prop6.8]{6.8} to prove
the following result.

\begin{enonce*}{Corollary 6.10}\etiquette[6.10]{XII.cor6.10}
The notation is that of \hyperref[XII.prop6.5]{6.5}. For simplicity,
suppose that the unparenthesized condition in
\hyperref[XII.prop6.5]{6.5} \textup{(i)} holds.
\begin{enumerate}
\renewcommand{\theenumi}{\roman{enumi}}
\renewcommand{\labelenumi}{(\theenumi)}
\item The parenthesized condition in \hyperref[XII.prop6.5]{6.5}
\textup{(i)} holds (respectively, condition
\hyperref[XII.prop6.5]{6.5} \textup{(ii) a)} holds, respectively, condition
\hyperref[XII.prop6.5]{6.5} \textup{(iii) a)} holds) if and only if the
following condition is satisfied. For every sheaf $F$ of sets
(respectively, of ind-$\LL$-groups, respectively, of abelian
$\LL$-torsion groups) on $X$, every
$\xi\in H^{\circ}(Y,h^{\ast}(F))$ (respectively, every
$\xi\in H^{i}(Y,h^{\ast}(F))$ with $i=0,1$, respectively, every
$\xi\in H^{i}(Y,h^{\ast}(F))$ with $i\leq n$), and every nonempty closed
subset $X'$ of $X$, if $Y'$ denotes its inverse image in $Y$, there exists
a morphism $p:\bar{X}\rightarrow X$ satisfying the following conditions:
\begin{itemize}
\item[$1^{\circ}$)] The image $p(\bar{X})$ is contained in $X'$ and
contains a nonempty open subset of $X'$.

\item[$2^{\circ}$)] For every sheaf of sets $\bar{F}$ on $\bar{X}$, the
base-change homomorphism
$$
h^{\ast}(p_{\ast}(\bar{F}))\longrightarrow
p_{Y\ast}(\bar{h}^{\ast}(\bar{F}))
$$
is an isomorphism (respectively, for every sheaf of ind-$\LL$-groups
$\bar{F}$ on $\bar{X}$, the base-change homomorphism
$$
h^{\ast}(R^{i}p_{\ast}(\bar{F}))\longrightarrow
R^{i}p_{Y\ast}(\bar{h}^{\ast}(\bar{F}))
$$
is bijective for $i=0,1$; respectively, for every sheaf of abelian
$\LL$-torsion groups $\bar{F}$ on $\bar{X}$, the base-change homomorphism
is bijective for $i\leq n$).

\item[$3^{\circ}$)] The inverse image of $\xi$ in
$$
H^{i}(\bar{Y},p_Y^{\ast}(h^{\ast}(F)))
=H^{i}(\bar{Y},\bar{h}^{\ast}(p^{\ast}(F)))
$$
lies in the image of $H^{i}(\bar{X},p^{\ast}(F))$.
\end{itemize}

\item Let\pageoriginale $n\geq -1$, and suppose that, for every sheaf of
sets $F$ (respectively, every sheaf of ind-$\LL$-groups, respectively,
every sheaf of abelian $\LL$-torsion groups) on $X$, the homomorphism
$$
H^{i}(X,F)\longrightarrow H^{i}(Y,h^{\ast}(F))
$$
is bijective for $i\leq n$ and injective for $i=n+1$ (in the first case
we take $n=-1$, and in the second case $n=-1$ or $0$). Let
$\xi\in H^{n+1}(Y,h^{\ast}(F))$. Then $\xi$ lies in the image of
$H^{n+1}(X,F)$ if and only if it satisfies the condition stated in
\textup{(i)} above, with $i=n+1$.
\end{enumerate}
\end{enonce*}

\noindent
\emph{Proof of} \hyperref[XII.cor6.10]{6.10}. Using
\hyperref[XII.prop6.5]{6.5} and an immediate induction on $n$, one sees
that \hyperref[XII.cor6.10]{6.10} (i) follows from
\hyperref[XII.cor6.10]{6.10} (ii), which we now prove.

Suppose that $\xi$ does not lie in the image of $H^{n+1}(X,F)$. Let
$\Phi$ be the set of closed subsets $X'$ of $X$ such that the image of
$\xi$ in $H^{n+1}(Y',h^{\ast}(F)|Y')$ does not lie in the image of
$H^{n+1}(X',F|X')$. By hypothesis, $X\in\Phi$, so $\Phi$ is nonempty.
Order $\Phi$ by the relation $\supset$, and let us show that $\Phi$ is
inductive. It is enough to prove that, if $(X'_{\lambda})$ is a totally
ordered family of elements of $\Phi$ and $X'$ is their intersection, then
$X'\in\Phi$. Give the $X'_{\lambda}$ and $X'$ their induced reduced
structures. The morphisms $X'_{\lambda}\rightarrow X$ are affine (being
closed immersions), and hence the inverse system of the $X'_{\lambda}$
satisfies the hypotheses considered in VII 5. Moreover, one immediately
sees that $X'$ is the inverse limit of the $X'_{\lambda}$. Likewise, $Y'$
is the inverse limit of the inverse system of the $Y'_{\lambda}$, which
are affine over $Y$. By VII \hyperref[VII.cor5.8]{5.8}, therefore,
$$
\varinjlim H^{n+1}(X'_{\lambda},F|X'_{\lambda})\xrightarrow{\sim}
H^{n+1}(X',F|X')\quav
\varinjlim H^{n+1}(Y'_{\lambda},h^{\ast}(F)|Y'_{\lambda})
\xrightarrow{\sim}H^{n+1}(Y',h^{\ast}(F)|Y'){\quoi}.
$$
It follows at once that, if
$(\xi|Y')\in\Iim H^{n+1}(X',F|X')$, then there\pageoriginale would be an index $\lambda$
such that
$(\xi|Y'_{\lambda})\in\Iim H^{n+1}(X'_{\lambda},F|X'_{\lambda})$, which
is absurd. Thus $X'\in\Phi$, as asserted. Consequently $\Phi$ is
inductive and therefore contains an inclusion-minimal element, say $X'$.
Since $\varnothing\notin\Phi$, the set $X'$ is nonempty. Apply the
hypothesis of \hyperref[XII.cor6.10]{6.10} to $\xi$ and $X'$. This gives a
morphism $p:\bar{X}\rightarrow X$ satisfying conditions $1^{\circ}$)--
$3^{\circ}$) in \hyperref[XII.cor6.10]{6.10}. Observe that we may even
suppose that $p$ is \emph{surjective onto $X'$}. Indeed, let $U$ be a
nonempty open subset of $X'$ contained in $p(\bar{X})$, and let $X'_1$ be
its complement in $X'$. By construction, $X'_1\notin\Phi$. Put
$\bar{X}_1=\bar{X}\amalg X'_1$, and let
$p_1:\bar{X}_1\rightarrow X$ be the morphism defined by $p$ and the
inclusion $X'_1\hookrightarrow X$. Using the relation
$X'_1\notin\Phi$, one sees immediately that $p_1$ still satisfies the same
hypotheses as $p$; moreover, it is surjective onto $X'$. We may therefore
suppose that $p$ is surjective onto $X'$. Now apply
\hyperref[XII.cor6.9]{6.9}, replacing $h:Y\rightarrow X$ there by
$h':Y'\rightarrow X'$ (with $X'$ endowed, say, with its induced reduced
structure), and replacing $\xi$ by $\xi|Y'$. It follows from
\hyperref[XII.cor6.9]{6.9} that
$(\xi|Y')\in\Iim H^{n+1}(X',F|X')$, contrary to $X'\in\Phi$. This
completes the proof of \hyperref[XII.cor6.10]{6.10}.

We shall use \hyperref[XII.cor6.10]{6.10} to reduce
\hyperref[XII.cor5.5]{5.5} to the case in which $f$ is projective. In the
non-Noetherian case, however, this requires a non-Noetherian variant of
Chow's lemma, which will be given in the next section.

\begin{enonce*}{Proposition 6.11}\etiquette[6.11]{XII.prop6.11}
(Transitivity lemma). The notation is that of
\hyperref[XII.prop6.8]{6.8}, but no surjectivity hypothesis is imposed on
$p$. Suppose that $h$ satisfies condition \textup{(i) a)} (respectively,
\textup{(ii) a)}, respectively, \textup{(iii) a)}) of
\hyperref[XII.prop6.5]{6.5}. In case \textup{(i)}, suppose in addition that,
for every sheaf of sets $\bar{F}$ on $\bar{X}$, the base-change
homomorphism of \hyperref[XII.prop6.8]{6.8} is a monomorphism
(respectively, an isomorphism); in case \textup{(ii)}, suppose that $\bar{X}$
is Noetherian and that, for every sheaf of ind-$\LL$-groups $\bar{F}$, the
base-change homomorphism
$$
h^{\ast}(R^ip_{\ast}(\bar{F}))\longrightarrow
R^ip_{Y\ast}(\bar{h}^{\ast}(\bar{F}))
$$
is\pageoriginale bijective for $i=0,1$; and finally, in case \textup{(iii)}, suppose that the
base-change homomorphism is bijective for $i\leq n$ and injective for
$i=n+1$. Then $\bar{h}:\bar{Y}\rightarrow\bar{X}$ also satisfies
condition \textup{(i)} (respectively, \textup{(ii)}, respectively, \textup{(iii)}) of
\hyperref[XII.prop6.5]{6.5}.
\end{enonce*}

\noindent
\emph{Proof of} \hyperref[XII.prop6.11]{6.11}. In case (i), consider the
commutative diagram in the proof of \hyperref[XII.prop6.8]{6.8} (i). Here
the hypotheses imply that the first horizontal arrow and the upper right
vertical arrow are injective (respectively, bijective), and hence so is the
second horizontal arrow, as required. In case (ii), let $\bar{F}$ be a
sheaf of ind-$\LL$-groups on $\bar{X}$. We must show that every element of
$H^1(\bar{Y},\bar{h}^{\ast}(\bar{F}))$ comes from an element of
$H^1(\bar{X},\bar{F})$. By \hyperref[XII.cor6.6]{6.6} (ii), applied to
$\bar{h}$, it is enough to find a monomorphism
$\bar{F}\rightarrow\bar{G}$ of sheaves of groups on $\bar{X}$ that
\emph{effaces} $H^1(\bar{Y},\bar{h}^{\ast}(\bar{F}))$. Consider the
second row of the diagram used in the proof of
\hyperref[XII.prop6.8]{6.8} (ii). It is enough, successively, to find a
monomorphism $\bar{F}\rightarrow\bar{G}$, with $\bar{G}$ a sheaf of
ind-$\LL$-finite groups, that effaces the last term
$H^{\circ}(Y,R^1p_{Y\ast}(\bar{h}^{\ast}(\bar{F})))$, and then a
monomorphism of sheaves of groups $\bar{G}\rightarrow\bar{H}$ that
effaces $H^1(Y,p_{Y\ast}(\bar{h}^{\ast}(\bar{G})))$. As observed in the
proof of \hyperref[XII.prop6.8]{6.8} (i), the two terms to be effaced are
respectively isomorphic, by the outer vertical arrows of the diagram in
question, to $H^{\circ}(X,R^1p_{\ast}(\bar{F}))$ and
$H^1(X,p_{\ast}(\bar{F}))$, in view of the hypotheses made here on the
base-change homomorphisms for $p$ and $h$. These isomorphisms are plainly
functorial in $\bar{F}$, so it is enough to efface the latter two terms.
By \hyperref[XII.prop3.3]{3.3}, one can efface
$R^1p_{\ast}(\bar{F})$ (and, a fortiori,
$H^{\circ}(X,R^1p_{\ast}(\bar{F}))$) by embedding $\bar{F}$ into an
ind-$\LL$-group $\GG$. It remains to efface
$H^1(X,p_{\ast}(\bar{F}))$. This is possible because, by\pageoriginale the diagram in
question, the latter term maps into $H^1(\bar{X},\bar{F})$ by a
monomorphism functorial in $\bar{F}$, while
$H^1(\bar{X},\bar{F})$ is again effaceable by
\hyperref[XII.prop3.3]{3.3}.

It remains to treat case (iii), and here again we read the proof of
\hyperref[XII.prop6.8]{6.8} backwards\footnote[1]{Instead of
\hyperref[XII.cor6.3]{6.3}, one could also use the Leray spectral sequence
for $p$ and $p_Y$.}. We must prove that the homomorphism
$$
H^{i}(\bar{X},\bar{F})\longrightarrow
H^{i}(\bar{Y},\bar{h}^{\ast}(\bar{F}))
$$
is an isomorphism for $i\leq n$ and a monomorphism for $i=n+1$. These
homomorphisms are induced by a homomorphism of complexes
$$
\RR\Gamma_{\bar{X}}(\bar{F})\longrightarrow
\RR\Gamma_{\bar{Y}}(\bar{h}^{\ast}(\bar{F})){\quoi},
$$
obtained from an injective resolution $C(\bar{F})$ and a homomorphism from
$\bar{h}^{\ast}(C(\bar{F}))$ (which is a resolution of
$\bar{h}^{\ast}(\bar{F})$) to an injective resolution of
$\bar{h}^{\ast}(\bar{F})$. The homomorphism above also identifies with
the composite
$$
\RR\Gamma_X(Rp_{\ast}(\bar{F}))\longrightarrow
\RR\Gamma_Y(h^{\ast}(Rp_{\ast}(\bar{F})))\longrightarrow
\RR\Gamma_Y(Rp_{Y\ast}(\bar{h}^{\ast}(\bar{F}))){\quoi},
$$
where the first homomorphism is the homomorphism induced by $h^{\ast}$ for
$Rp_{\ast}(\bar{F})$, and the second is obtained by applying
$\RR\Gamma_Y$ to the base-change homomorphism
$$
h^{\ast}(Rp_{\ast}(\bar{F}))\longrightarrow
Rp_{Y\ast}(\bar{h}^{\ast}(\bar{F})).
$$
By the hypothesis on $h$ and \hyperref[XII.cor6.3]{6.3} (i), the first
homomorphism therefore induces isomorphisms on the cohomology groups of
the complexes in question in the required range. The same is true of the
second by our hypothesis on the base-change homomorphisms, using the
argument given in the proof of \hyperref[XII.prop6.8]{6.8} (iii). The same
conclusion thus holds for their composite, completing the proof of
\hyperref[XII.prop6.11]{6.11}.

\begin{enonce*}[remark]{Remark 6.12}\etiquette[6.12]{XII.rem6.12}
The argument\pageoriginale given for \hyperref[XII.prop6.11]{6.11} in cases (i) and
(iii) gives the following more precise result. In order to conclude, for a
\emph{given} sheaf $\bar{F}$ on $\bar{X}$, that
$$
H^{i}(\bar{X},\bar{F})\longrightarrow
H^{i}(\bar{Y},\bar{h}^{\ast}(\bar{F}))
$$
is bijective for $i\leq n$ and injective for $i=n+1$, it is enough to
assume that
\begin{itemize}
\item[$1^{\circ}$)] the base-change homomorphism
$$
h^{\ast}(R^{i}p_{\ast}(\bar{F}))\longrightarrow
R^{i}p_{Y\ast}(\bar{h}^{\ast}(\bar{F}))
$$
is an isomorphism for $i\leq n$ and a monomorphism for $i=n+1$, and

\item[$2^{\circ}$)] for every pair of integers $i,j\geq0$, the
homomorphism
$$
H^{i}(X,R^{j}p_{\ast}(\bar{F}))\longrightarrow
H^{i}(Y,h^{\ast}(R^{j}p_{\ast}(\bar{F})))
$$
induced by $h$ is an isomorphism if $i+j\leq n$ and a monomorphism if
$i+j=n+1$.
\end{itemize}

There is plausibly an analogous result in case (ii), but its statement
would necessarily involve noncommutative $2$-cohomology. To avoid having
to invoke that theory, in \hyperref[XII.prop6.11]{6.11} (ii) we used a
different argument, based on the effaceability properties of
noncommutative $H^1$; unlike the argument above, however, it yields no
result for a fixed $\bar{F}$.
\end{enonce*}

\begin{enonce*}[remark]{Remarks 6.13}\etiquette[6.13]{XII.rems6.13}
The most important case in which the conditions of
\hyperref[XII.prop6.5]{6.5} (i), (ii), and (iii) hold is that considered
in \hyperref[XII.cor5.5]{5.5}; it will be proved in the next exposé using
the reductions made in the present one. Another interesting case (in
which $h$ is no longer a closed immersion as in
\hyperref[XII.cor5.5]{5.5}) will be studied in XV. Let us also mention the
closely related case in which $X$ is the spectrum of a Noetherian ring
$A$, separated and complete for the topology defined by an ideal $J$, and
$Y=\Spec(A/J)$ is the closed subscheme of $X$ defined by $J$. One then
verifies condition \hyperref[XII.prop6.5]{6.5} (ii) directly, in the form\pageoriginale
$\textup{ b'')}$ (with $\LL=\PP$), by \EGA IV 18.3.2. It is very plausible
that the conditions of \hyperref[XII.prop6.5]{6.5} (iii) also hold for
every $n$; that is, for every torsion sheaf $F$ on $X$, the homomorphisms
$$
H^n(X,F)\longrightarrow H^n(Y,F|Y)
$$
are bijective. One can formulate considerably more general conjectures,
related to a generalization of the construction of henselization to
nonlocal schemes. First, no example is known in which $h$ is a closed
immersion and the conditions of \hyperref[XII.prop6.5]{6.5} (ii) and (iii)
(with $\LL=\PP$) fail when the parenthesized condition in
\hyperref[XII.prop6.5]{6.5} (i) holds, that is, when $(X,Y)$ is a
\emph{henselian pair} in the terminology of \EGA IV 18.5.5. One may
therefore conjecture that every such pair satisfies the conditions of
\hyperref[XII.prop6.5]{6.5} (ii) and (iii) (with $\LL=\PP$, for every
$n$).

On the other hand, one may ask under what conditions the following
property, for a pair $(X,Y)$ consisting of a quasi-compact and
quasi-separated scheme $X$ and a closed subscheme $Y$, implies that the
pair is henselian:

$(\ast)$ For every scheme $X'$ étale over $X$, put
$Y'=\fprod{X'}{Y}{X}$. The canonical map
$$
\Gamma(X'/X)\longrightarrow\Gamma(Y'/Y)
$$
is bijective (\cf \EGA IV 18.5.4 b)). This condition is plainly weaker
\emph{a priori} than the henselian condition
\hyperref[XII.prop6.5]{6.5} (i), and is in fact strictly weaker: it holds,
for example, when $X$ is a normal irreducible projective scheme of
dimension $\geq2$ over a field $k$ and $Y$ is a hyperplane section of $X$
(as follows readily from \SGA 2 XII 2.1, taking $S=\Spec(k)$), whereas
the pair $(X,Y)$ is plainly not henselian. It is possible, however, that
when $X$ is affine, condition $(\ast)$ implies that $(X,Y)$ is a henselian
pair. If so, for every affine scheme $X$ and every closed subset $Y$ of
$X$, a\pageoriginale process of étale localization along $Y$, analogous to the
henselization of local rings, ought to associate to it a \emph{henselian}
pair $(X',Y)$, where $X'$ is ``pro-étale'' over $X$. Another hypothesis on
$(X,Y)$ (suggested by results of Hironaka and Rossi on the contractibility
of certain subvarieties of algebraic varieties) under which the henselian
property might perhaps be deduced from the weaker property $(\ast)$ is the
following: the immersion $Y\rightarrow X$ is regular of codimension $1$,
and the conormal sheaf $\NN_{Y/X}=\JJ/\JJ^2$ (where $\JJ$ is the ideal
defining $Y$ in $X$) is an \emph{ample} invertible sheaf on $Y$. (In the
counterexample considered above, by contrast, this conormal sheaf was
anti-ample, that is, its inverse was ample.)
\end{enonce*}

\footnotetext{\hypertarget{sga4.exp.XII.note.section7}{}A reader interested in the statements of \hyperref[XII.sec5]{Section 5} only in
the Noetherian case may omit the present section.}

% END INCLUDED SGA4_English_current_buildable_source_20260906/source/12/drafts/SGA4_XII_LINES_1601_1924_English.texfrag


% BEGIN INCLUDED SGA4_English_current_buildable_source_20260906/source/12/drafts/SGA4_XII_LINES_1925_2169_English.texfrag
\section[A variant of Chow's lemma ${}^{(*)}$.]{A variant of Chow's lemma
\texorpdfstring{\protect\hyperlink{sga4.exp.XII.note.section7}{${}^{(*)}$}}{(*)}.}\etiquette[7]{XII.sec7}

We shall need the following variant of (\EGA II 5.6.1), stripped of
all Noetherian and irreducibility hypotheses.

\begin{enonce*}{Lemma 7.1}\etiquette[7.1]{XII.lem7.1}
Let $S$ be a quasi-compact and quasi-separated scheme, and let
$f:X\rightarrow S$ be a separated morphism of finite type, with $X$
nonempty. There exists a commutative diagram
$$
\xymatrix@=1.25cm{
X \ar[rd]_{f}&    &\bar{X}\ar[ll]_{\pi}\ar[ld]^{\bar{f}}\\
  &  S &
}
$$
with\pageoriginale $\pi$ projective and $\bar{f}$ quasi-projective,
and a nonempty open subset $U$ of $X$, such that $\pi$ induces an
isomorphism from
$\bar{U}=\fprod{U}{\bar{X}}{X}=\pi^{-1}(U)$ onto $U$.
\end{enonce*}

\noindent
\emph{Proof.} We first make the following observation. Let
$i:X'\rightarrow X$ be a closed $S$-subscheme of $X$, and let
$V\subset X'$ be a subscheme that is open in $X$ and dense in $X'$.
If the lemma holds for $X'$, then it holds for $X$. Indeed, let
$$
\xymatrix@=1.25cm{
X' \ar[rd]_{f'}&    &\bar{X}'\ar[ll]_{\pi'}\ar[ld]^{\bar{f}'}\\
  &  S &
}
$$
be a diagram as in the statement, and let $U'\neq\varnothing$ be an
open subset of $X'$ such that $\pi'$ induces an isomorphism
$\bar{U}'={\pi'}^{-1}(U')\rightarrow U'$. Since $V$ is dense in
$X'$, we have $U'\cap V\neq\varnothing$, and we may therefore replace
$U'$ by this open subset. Then $U'$ is open in $X$, and we set
$\pi=i\pi'$, $\bar{f}=\bar{f}'$, and $U=U'$.

Now $X$ admits a finite covering by nonempty affine open subsets
$U_k$, $k=1,\ldots,n$. The hypotheses on $S$ imply that every
$U_k\rightarrow S$ is quasi-affine, and hence, a fortiori,
quasi-compact. We prove the lemma for every $X$ by induction on $n$.
If the intersection $U$ of the $U_k$ is dense in $X$, we may copy the
proof of (\EGA II 5.6, B, C, D), with the necessary changes in
notation. If the intersection is nonempty, we reduce to the case in
which it is dense by replacing $X$ by the closure $X'$ of $U$ (as in
the observation above), and $U_k$ by $U'_k=U_k\cap X'$. Finally,
suppose that $U$ is empty, and choose $r$, $1\leq r<n$, such that
$$
V=\bigcap_{k=1}^{r}U_k\neq\varnothing
\qquad\hbox{but}\qquad
V\cap U_{r+1}=\varnothing;
$$
this is plainly possible since $U_1\neq\varnothing$. We may now
replace $X$ by the closure $X'$ of $V$ and each $U_k$ by
$U_k\cap X'$, discarding the empty members. Since
$U_{r+1}\cap X'=\varnothing$ (it is open in $X'$ and disjoint from
the dense open $V$), the result follows from the induction
hypothesis.

\section{Final reductions}\etiquette[8]{XII.sec8}\pageoriginale

\begin{enonce*}{Lemma 8.1}\etiquette[8.1]{XII.lem8.1}
To prove any one of the assertions
{\rm \hyperref[XII.thm5.1]{5.1} (i), (ii), (iii)}, it is enough to do
so under the assumptions that $f$ is \emph{projective} and $S$ is
Noetherian.
\end{enonce*}

Suppose first that \hyperref[XII.thm5.1]{5.1} has been proved when $f$
is projective; we show how the same assertion follows for arbitrary
$f$. Using \hyperref[XII.lem6.1]{6.1}, we are reduced to the situation
of \hyperref[XII.cor5.5]{5.5}, with $S$ strictly local. We now apply
\hyperref[XII.cor6.10]{6.10}, taking $Y=X_{\circ}$. Given
$(\xi,X')$ as in \hyperref[XII.cor6.10]{6.10}, we apply Chow's lemma
\hyperref[XII.lem7.1]{7.1} to find an $S$-morphism
$p:\bar{X}\rightarrow X'$, with $\bar{X}$ projective over $S$, such
that $p(\bar{X})$ contains a nonempty open subset of $X'$ (condition
$1^{\circ}$ of \hyperref[XII.cor6.10]{6.10}). Since $p$ is
projective, the base-change homomorphisms for $(p,h)$ are isomorphisms
by hypothesis (condition $2^{\circ}$ of
\hyperref[XII.cor6.10]{6.10}). Finally, since $\bar{X}$ is projective
over $S$, the inclusion
$\bar{Y}=\bar{X}_{\circ}\rightarrow\bar{X}$ induces, by hypothesis,
an isomorphism on cohomology groups
$$
H^{i}(\bar{X},\bar{F})\longrightarrow
H^{i}(\bar{X}_{\circ},\bar{h}^{\ast}(\bar{F})),
$$
which is condition $3^{\circ}$ of
\hyperref[XII.cor6.10]{6.10}. Thus all the conditions of
\hyperref[XII.cor6.10]{6.10} are satisfied, and the conclusion
follows.

We next show that, in order to verify any one of the three assertions
\hyperref[XII.thm5.1]{5.1} for every situation $(f,F,g)$ with $f$
projective, we may further assume that $S$ is Noetherian. By
\hyperref[XII.lem6.1]{6.1}, we may suppose that $S$ is strictly local
and that $S'\rightarrow S$ is the inclusion of the closed point; we
are thus reduced to verifying \hyperref[XII.cor5.5]{5.5} with $S$
strictly local. Moreover, since $X$ embeds into the projective scheme
$\PP^r_S$, compatibility \hyperref[XII.prop4.4]{4.4} (ii) allows us
to replace $X$ by $\PP^r_S$ and $F$ by its direct image on
$\PP^r_S$. We are therefore reduced to the case $X=\PP^r_S$.

Write $S=\Spec(A)$. The ring $A$ is the filtered direct limit of
Noetherian Henselian local subrings $A_i$ for which the inclusions
$A_i\rightarrow A$ are local: one may take the strict henselizations,
at the prime ideals induced by the maximal ideal of $A$, of the
finitely generated $\ZZ$-subalgebras of $A$. Under these conditions,
the residue field $k$ of $A$ is likewise the filtered direct limit of
the residue fields $k_i$ of the $A_i$ (\EGA IV 5.13.1). Set
$S_i=\Spec(A_i)$ and $X_i=\PP^r_{S_i}$. Suppose that a direct system
$(F_i)$ of sheaves of the type occurring in the assertion under
consideration is given on the $X_i$, as in VII \hyperref[VII.thm5.7]{5.7}, and let $F$ be
the sheaf it defines on $X$. Let $F_{i\circ}$ (respectively,
$F_{\circ}$) be the sheaf induced by $F_i$ on the fiber
$X_{i\circ}$ over the closed point $s_i$ of $S_i$ (respectively, by
$F$ on the fiber $X_{\circ}$ over the closed point $s$ of $S$). We
then have a commutative diagram
$$
\xymatrix{
\displaystyle\mathop{\varinjlim}_{i} H^{n}(X_i,F_i) \ar[r] \ar[d]&
H^{n}(X,F)\ar[d]\\
\displaystyle\mathop{\varinjlim}_{i} H^{n}(X_{i\circ},F_{i\circ}) \ar[r] &
H^{n}(X_{\circ},F_{\circ}),
}
$$
whose two horizontal arrows are isomorphisms by VII \hyperref[VII.thm5.7]{5.7}. Consequently,
to prove that the second vertical arrow is an isomorphism, it is
enough to prove this for the first one; hence it is enough to prove
that the homomorphisms
$$
H^n(X_i,F_i)\longrightarrow H^n(X_{i\circ},F_{i\circ})
$$
are isomorphisms. On the other hand, every sheaf $F$ of the type under
consideration on $X$ is isomorphic to the limit of a direct system of
the preceding kind; for example, one may take $F_i=u_{i*}(F)$, where
$u_i:X\rightarrow X_i$ is the canonical morphism, by
\hyperref[XII.cor5.2]{5.2} and
\hyperrefext[IX][IX.prop1.6]{1.6} (iii). This completes the proof of
\hyperref[XII.lem8.1]{8.1}.

\begin{enonce*}{Corollary 8.2}\etiquette[8.2]{XII.cor8.2}
Theorem {\rm \hyperref[XII.thm5.1]{5.1} (i)} is true.
\end{enonce*}

Indeed, it is enough to combine \hyperref[XII.lem8.1]{8.1}, criterion
\hyperref[XII.prop6.5]{6.5} (i) ${\rm b'')}$, and
\hyperref[XII.lem5.8]{5.8}.

\begin{enonce*}{Lemma 8.3}\etiquette[8.3]{XII.lem8.3}
To\pageoriginale prove any one of the assertions
{\rm \hyperref[XII.thm5.1]{5.1} (i), (ii), (iii)}, it is enough to do
so when $f$ is projective of relative dimension $\leq 1$ and $S$ is
Noetherian.
\end{enonce*}

Indeed, \hyperref[XII.lem8.1]{8.1} already shows that we may restrict
ourselves to the case in which $f$ is projective and $S$ is
Noetherian, and an argument used above further allows us to suppose
that $X=\PP^r_S$. We proceed by induction on $r$, noting that the
theorem holds by hypothesis for $r=1$. We may therefore suppose that
$r\geq2$ and that the theorem has been proved for
$\PP^{r-1}_S\rightarrow S$ (over every base $S$). Embed
$Z=\PP^{r-2}_S$ in $X=\PP^r_S$ in the usual way, and let
$p:\bar{X}\rightarrow X$ be the blowup of $X$ along $Z$. A direct
calculation (\cf \EGA V) gives a natural morphism
$p_1:\bar{X}\rightarrow X_1=\PP^1_S$ that makes $\bar{X}$ an
$X_1$-scheme locally isomorphic to $\PP^{r-1}_{X_1}$ (it is the
projective bundle over $X_1$ associated with a suitable locally free
sheaf of rank $r$ on $X_1$). On the other hand, the fibers of $p$ have
dimension $\leq1$, as do those of $f_1:X_1\rightarrow S$. We may
therefore apply the relative-dimension-one hypothesis to $p$ and
$f_1$, and the induction hypothesis to $p_1$, in the commutative
diagram
$$
\xymatrix{
   & \bar{X} \ar[ld]_{p} \ar[rd]^{p_1} &          \\
 X \ar[rd]_{f}&           & X_1\ar[ld]^{f_1}\\
   &    S      &
}$$
Since $p$ is surjective, \hyperref[XII.prop6.8]{6.8} reduces us to
proving \hyperref[XII.thm5.1]{5.1} for the composite
$\bar{f}=fp=f_1p_1$. But we may already apply
\hyperref[XII.thm5.1]{5.1} to both $f_1$ and $p_1$. Taking
\hyperref[XII.prop6.11]{6.11} into account, we conclude that
\hyperref[XII.thm5.1]{5.1} also holds for their composite $f_1p_1$.
This completes the proof of \hyperref[XII.lem8.3]{8.3}.

\begin{enonce*}{Lemma 8.4}\etiquette[8.4]{XII.lem8.4}
To\pageoriginale prove {\rm \hyperref[XII.thm5.1]{5.1} (ii)}, it is
enough to prove the ``specialization theorem for the fundamental
group'' {\rm \hyperref[XII.thmbis5.9]{5.9 bis}} in the case where $f$
is projective of relative dimension $\leq1$ and $S$ is Noetherian and
strictly local.
\end{enonce*}

This follows at once from \hyperref[XII.lem8.3]{8.3},
\hyperref[XII.lem6.1]{6.1}, and \hyperref[XII.prop6.5]{6.5} (ii)
(criterion b)).

The specialization theorem in the case of \hyperref[XII.lem8.4]{8.4}
will be proved directly in the next exposé
(\hyperrefext[XIII][XIII.sec2]{2}).

\begin{enonce*}{Lemma 8.5}\etiquette[8.5]{XII.lem8.5}
To prove {\rm \hyperref[XII.thm5.1]{5.1} (iii)}, it is enough to prove
{\rm \hyperref[XII.thm5.1]{5.1} (ii)} and, in addition, the following
proposition (which will be proved in the next exposé
{\rm (\hyperrefext[XIII][XIII.sec3]{3})}, or in
{\rm (\EGA IV 21.9.12)}):
\end{enonce*}

\begin{enonce*}{Proposition 8.6}\etiquette[8.6]{XII.prop8.6}
Let $S$ be Henselian and Noetherian, let $f:X\rightarrow S$ be
projective of relative dimension $\leq1$, and let $X_{\circ}$ be the
closed fiber of $X/S$. For every scheme $Z$, let
$\Pic Z=H^1(Z,(G_m)_Z)$ denote the Picard group of $Z$. The restriction
homomorphism
$$
\Pic X\longrightarrow \Pic X_{\circ}
$$
is surjective.
\end{enonce*}

\noindent
\emph{Proof of \hyperref[XII.lem8.5]{8.5}.} By
\hyperref[XII.prop6.5]{6.5} (iii) and
\hyperref[XII.cor8.2]{8.2}, it is enough to prove that
$$
\varphi^q:H^q(X,\ZZ/\ell^\nu \ZZ)\longrightarrow
H^q(X_{\circ},\ZZ/\ell^{\nu}\ZZ)
$$
is surjective for every $q\geq1$ and every prime $\ell$. For $q=1$,
surjectivity is already known from \hyperref[XII.thm5.1]{5.1} (ii).
Moreover, since $X_{\circ}$ is a projective scheme of dimension
$\leq1$ over a separably closed field $k$, we have
$$
H^q(X_{\circ},\ZZ/\ell^{\nu}\ZZ)=0
$$
for $q>2$ when $\ell\neq\operatorname{char} k$
(\hyperrefext[X][X.cor4.3]{4.3}), and for $q>1$ when
$\ell=\operatorname{char} k$ (\hyperrefext[X][X.cor5.2]{5.2}). Thus
surjectivity is automatic for these values of $q$.

It\pageoriginale remains to consider $q=2$ for a prime
$\ell\neq\operatorname{char} k$. Put $n=\ell^{\nu}$. The sheaf
$(\mu_n)_S$ of $n$-th roots of unity is locally isomorphic to
$(\ZZ/n\ZZ)_S$, and, since $S$ is strictly local, there is a
noncanonical isomorphism
$$
(\mu_n)_S\simeq(\ZZ/n\ZZ)_S,
$$
and hence $(\mu_n)_X\simeq(\ZZ/n\ZZ)_X$.

Kummer theory \hyperrefext[IX][IX.Thr3.2]{3.2}, applied to $X$ and
$X_{\circ}$, gives a morphism of exact sequences
$$
\xymatrix@R=1.5cm{
\Pic X \ar[r] \ar[d]^{\epsilon}&
H^{2}(X,(\mu_n)_X) \ar[d]^{\varphi^2}\ar[r]&
H^{2}(X,(G_m)_X)\ar[d]\\
\Pic X_{\circ} \ar[r]&
H^2(X_{\circ},(\mu_n)_{X_{\circ}}) \ar[r]&
H^2(X_{\circ},(G_m)_{X_{\circ}}).
}
$$
Now $H^2(X_{\circ},(G_m)_{X_{\circ}})=0$ by
\hyperrefext[IX][IX.thm4.6]{4.6}. Consequently, the surjectivity of
$\epsilon$ implies that of $\varphi^2$, which proves the lemma.

\begin{thebibliography}{99}
\bibitem{sga4.exp.XII.bib.1} \textsc{Giraud} J., \emph{Cohomologie non-abélienne},
thesis (forthcoming), and two notes in C.R., vol. 260, pp. 2392 and
2656 (1965).

\bibitem{sga4.exp.XII.bib.2} \textsc{Godement} R., \emph{Théorie des faisceaux}, Paris
(1958).

\bibitem{sga4.exp.XII.bib.3} \textsc{Verdier} J.L., \emph{Catégories dérivées},
Institut des Hautes Études Scientifiques (1964).
\end{thebibliography}

\ifx\danslelivre\undefined
\end{document}
\fi

% END INCLUDED SGA4_English_current_buildable_source_20260906/source/12/drafts/SGA4_XII_LINES_1925_2169_English.texfrag


% END INCLUDED SGA4_English_current_buildable_source_20260906/source/12/12_English.tex


\makeatletter
\def\input@path{{../13/}}
\makeatother
% BEGIN INCLUDED SGA4_English_current_buildable_source_20260906/source/13/13_English.tex
\ifx\danslelivre\undefined
% Release engine: XeLaTeX (the LuaTeX compatibility layer breaks Xy-pic output).
\documentclass[10pt]{../fichiers-configuration-latex/smfart}
% BEGIN INCLUDED SGA4_English_current_buildable_source_20260906/source/fichiers-configuration-latex/paquets.tex
\usepackage{polyglossia}
%\usepackage[francais]{babel}
\usepackage{fontspec}
\usepackage[math-style=ISO,bold-style=upright,sans-style=upright]{unicode-math}
%\usepackage[mathletters]{ucs}
%\usepackage[utf8]{inputenc}
%\usepackage[T1]{fontenc}
\usepackage[unicode,hypertexnames=false,colorlinks=true,linkcolor=blue,citecolor=Sepia,urlcolor=magenta]{hyperref}
\setmainlanguage{french}
\setmainfont[Ligatures=TeX]{Linux Libertine O}
% variante possible : TeXGyrePagella ?
\setmathfont{XITS Math}

% maths
\usepackage{amsmath}
%\usepackage{amssymb}
%\usepackage{mathrsfs}
%\usepackage{stmaryrd}
%\usepackage{euscript}

% couleurs
\usepackage[usenames,dvipsnames]{xcolor}

% graphiques
\usepackage{graphicx}
\usepackage[all]{xy}
%\usepackage{tikz}
%\usetikzlibrary{matrix,arrows,calc}
%\tikzset{column sep=3em,row sep=2em,text height=1.5ex, text depth=0.25ex,>=angle 90}
%\tikzset{description/.style={fill=white,inner sep=1pt}}

% hyperliens, 'etiquettes
\usepackage{xr-hyper}
\usepackage{../fichiers-configuration-latex/smfthm}
%\ifx\cacheretiquettes\undefined
%\usepackage[inline]{showlabels} \else \fi % pour rendre visibles les 'etiquettes
%\usepackage{../fichiers-configuration-latex/makeidx} % pour faire un index
%\usepackage[pagebackref,unicode]{hyperref}  % pour voir o`u sont cit'ees les r'ef'erences
% hyperref is loaded once above with the cumulative-reader destination policy.
%\hypersetup{breaklinks=true} %permet le retour `a la ligne dans les liens trop longs. Incompatible avec pdflatex
%\renewcommand*{\backref}[1]{{\small [cit\'e page(s)~#1.]}}

\usepackage{xspace}
\usepackage{pifont} % pour expose XVII

% END INCLUDED SGA4_English_current_buildable_source_20260906/source/fichiers-configuration-latex/paquets.tex

% BEGIN INCLUDED SGA4_English_current_buildable_source_20260906/source/fichiers-configuration-latex/numerotation.tex
\renewcommand{\thefootnote}{\roman{footnote}} % pour mettre les notes en bas de page sous la forme (ii) et non (2), qui ressemble trop (?) a un exposant.

\newcounter{pageoriginal}
\newcommand{\pageoriginale}{\refstepcounter{pageoriginal}\marginpar{\footnotesize\textbf{\thepageoriginal}}\xspace}
\let\pageoriginaled\pageoriginale

\usepackage{ifthen}
\newboolean{original}
\def\sisi{\ifthenelse{\boolean{original}}}
\setboolean{original}{false}

\newcounter{cnote}
\def\nde#1{\ifthenelse{\boolean{original}}{}{\stepcounter{cnote}\def\thefootnote{\arabic{cnote}}\protect\footnote[\value{cnote}]{N.D.E.: #1}}}
\def\ndemark{\ifthenelse{\boolean{original}}{}{\stepcounter{cnote}\def\thefootnote{\arabic{cnote}}\protect\footnotemark[\value{cnote}]}}
\def\ndetext#1{\ifthenelse{\boolean{original}}{}{\def\thefootnote{\arabic{cnote}}\protect\footnotetext[\value{cnote}]{N.D.E.: #1}}}

% END INCLUDED SGA4_English_current_buildable_source_20260906/source/fichiers-configuration-latex/numerotation.tex

% BEGIN INCLUDED SGA4_English_current_buildable_source_20260906/source/fichiers-configuration-latex/macros.tex
\setmainlanguage{french}
\setotherlanguage{english}


%% Fixed, engine-independent release identifier.  The inherited dynamic Git
%% probe used Lua io.popen, exposed build-state-dependent output, and printed
%% a visible question mark outside a Git checkout.
\def\showgitstatus{English release r2, 2026-07-23}


% macros pour LaTeX
\renewcommand{\mathsf}{}

% Police différente pour U+2113 SCRIPT SMALL L
\setmathfont[range="2113]{latinmodern-math.otf}

% Give every inherited semantic label its own stable PDF destination instead
% of inheriting whichever automatic equation/page anchor happened to be live.
% The legacy label remains as a compatibility alias used by existing edges.
\makeatletter
\def\etiquette[#1]#2{%
  \Hy@raisedlink{\hyper@anchorstart{sga4.semantic.#2}\hyper@anchorend}%
  \def\@currentHref{sga4.semantic.#2}%
  \def\@currentlabel{#1}%
  \label{#2}%
}
\makeatother
\def\eqhyperref[#1]#2{(\hyperref[#1]{#2})}
\def\pagehyperref[#1]#2{(page \hyperref[#1]{#2})}
\def\hyperrefext[#1][#2]#3{\hyperref[#2]{#1 #3}}
\def\eqhyperrefext[#1][#2]#3{\hyperref[#2]{#1 #3}}     %cf. exposes 17
\def\neweqhyperrefext[#1][#2]#3{\hyperref[#2]{#1 #3}}  %cf. exposes 17

\def\shat#1{\protect\widehat{#1}}
\def\stilde#1{\protect\widetilde{#1}}
%%%%%%%%%%%%%%%%%%


\def\etcn{etc.\xspace}
\def\etco{etc...\xspace}
\def\etc{\sisi{\etco}{\etcn}}
\def\ldotsn{\sisi{...}{\etcn}}

\def\cf{{cf.}\kern.3em}
\def\Cf{{Cf.}\kern.3em}
\def\ie{{i.e.}\kern.3em}
\def\iev{{i.e.},\xspace}
\def\loccito{loc.\kern2pt cit.\xspace}
\def\loccitn{\textit{loc.\kern2pt cit.}\xspace}
\def\loccit{\sisi{\loccito}{\loccitn}}
\def\EGA{\textup{EGA}\kern.2em}
\def\SGA{\textup{SGA}\kern.3em}

\def\quoi{}

\def\quav{\quad,\quad}
\def\quaet{\quad \text{et}\quad}
\def\fibreproduct#1#2#3{#1 ×_{#3} #2}
\let\fprod\fibreproduct
\def\ccheck#1#2{\mbox{${\check{\check{\rm #1}}}^{\lower2pt\hbox{$\scriptstyle #2$}}$}}
\def\sheaf#1{\protect\mathcal{#1}}
\def\vieuxsheaf#1{{#1}} % cf. problème avec sheaf{a\Phi_i} par exemple...
\DeclareMathOperator{\Hom}{Hom}
\DeclareMathOperator{\can}{can}
\DeclareMathOperator{\Ind}{Ind}
\DeclareMathOperator{\ind}{ind}
%\DeclareMathOperator{\Hhomtop}{\mathcal{H}\!\mathit{omtop}}
\DeclareMathOperator{\Hhomtop}{\mathbf{Homtop}}
\DeclareMathOperator{\Point}{Point}
\DeclareMathOperator{\Pro}{Pro}
\DeclareMathOperator{\Top}{Top}
\DeclareMathOperator{\Mod}{Mod}
\DeclareMathOperator{\Cart}{Cart}
%\DeclareMathOperator{\Mmorsite}{\mathcal{M}\!\mathit{orsite}}
\DeclareMathOperator{\Mmorsite}{\mathbf{Morsite}}
\DeclareMathOperator{\loc}{loc}
\DeclareMathOperator{\Loc}{Loc}
\DeclareMathOperator{\Pic}{Pic}
\DeclareMathOperator{\Et}{Et}
\DeclareMathOperator{\card}{card}
\DeclareMathOperator{\Ens}{Ens}
\DeclareMathOperator{\h}{\underline{H}}
\DeclareMathOperator{\Fonct}{Fonct}
\DeclareMathOperator{\Ker}{Ker}
\DeclareMathOperator{\but}{but}
\DeclareMathOperator{\Ob}{Ob}
\DeclareMathOperator{\ob}{ob}
\DeclareMathOperator{\Lib}{Lib}
\DeclareMathOperator{\esj}{esj}
\DeclareMathOperator{\id}{id}
\DeclareMathOperator{\Id}{Id}
\DeclareMathOperator{\ad}{ad}
\DeclareMathOperator{\Iim}{Im}
\DeclareMathOperator{\pr}{pr}
\DeclareMathOperator{\Cat}{Cat}
\DeclareMathOperator{\car}{car}
\DeclareMathOperator{\SysInd}{SysInd}
\DeclareMathOperator{\Sup}{Sup}
\DeclareMathOperator{\Coker}{Coker}
\DeclareMathOperator{\Filt}{Filt}
\DeclareMathOperator{\acc}{acc}
\DeclareMathOperator{\sur}{sur}
\DeclareMathOperator{\cat}{cat}
\DeclareMathOperator{\Ord}{Ord}
\DeclareMathOperator{\Gr}{Gr}
\DeclareMathOperator{\Cov}{Cov}
\DeclareMathOperator{\Caf}{Caf}
\DeclareMathOperator{\coker}{coker}
\DeclareMathOperator{\coim}{coim}
\DeclareMathOperator{\im}{im}
\DeclareMathOperator{\Ouv}{Ouv}
\DeclareMathOperator{\diag}{diag}
\DeclareMathOperator{\Esp}{Esp}
\DeclareMathOperator{\TOP}{TOP}
\DeclareMathOperator{\Unif}{Unif}
\DeclareMathOperator{\Iind}{\mathbf{Ind}}
\DeclareMathOperator{\Hhomcart}{\mathbf{Homcart}}
\DeclareMathOperator{\Hhomtopess}{\mathbf{Homtopess}}
\DeclareMathOperator{\Ppoints}{\mathbf{Points}}
\DeclareMathOperator{\Ppoint}{\mathbf{Point}}
\DeclareMathOperator{\Ttors}{\mathbf{Tors}}
\DeclareMathOperator{\Ssex}{\mathbf{Sex}}
\DeclareMathOperator{\Ffib}{\mathbf{Fib}}
\DeclareMathOperator{\Vvois}{\mathbf{Vois}}
\DeclareMathOperator{\Ppointess}{\mathbf{Pointess}}
\DeclareMathOperator{\scrHom}{\mathbf{Hom}}
\DeclareMathOperator{\sheafPic}{\mathbf{Pic}}
\DeclareMathOperator{\sheafAut}{\mathbf{Aut}}
\DeclareMathOperator{\Hhomtoplocan}{\mathbf{Homtoplocan}}
\DeclareMathOperator{\Hhomtopan}{\mathbf{Homtopan}}
\DeclareMathOperator{\Pptgeom}{\mathbf{Ptgéom}}
\DeclareMathOperator{\Morsite}{Morsite}
\DeclareMathOperator{\Fib}{Fib}
\DeclareMathOperator{\Fais}{Fais}
\DeclareMathOperator{\kar}{kar}
\DeclareMathOperator{\Groupes}{Groupes}
\DeclareMathOperator{\Restr}{Restr}
\DeclareMathOperator{\Isom}{Isom}
\DeclareMathOperator{\Prol}{Prol}
\DeclareMathOperator{\Res}{Res}
\DeclareMathOperator{\Tors}{Tors}
\DeclareMathOperator{\esp}{esp}
\DeclareMathOperator{\es}{es}
\DeclareMathOperator{\od}{od}
\DeclareMathOperator{\gr}{gr}
\DeclareMathOperator{\sob}{sob}
\DeclareMathOperator{\dfn}{dfn}
\DeclareMathOperator{\Sch}{Sch}
\DeclareMathOperator{\ouv}{ouv}
\DeclareMathOperator{\cons}{cons}
\DeclareMathOperator{\Nerf}{Nerf}
\DeclareMathOperator{\Points}{Points}
\DeclareMathOperator{\Espsob}{Espsob}
\DeclareMathOperator{\Homtop}{Homtop}
\DeclareMathOperator{\Pt}{Pt}
\DeclareMathOperator{\Quot}{Quot}
\DeclareMathOperator{\glob}{glob}
\DeclareMathOperator{\Homfib}{Homfib}
\DeclareMathOperator{\Qouv}{Qouv}
\DeclareMathOperator{\Quotouv}{Quotouv}
\DeclareMathOperator{\ens}{ens}
\DeclareMathOperator{\ab}{ab}
\DeclareMathOperator{\an}{an}
\DeclareMathOperator{\fs}{fs}
\DeclareMathOperator{\red}{red}
\DeclareMathOperator{\Ab}{Ab}
\DeclareMathOperator{\cosq}{cosq}
\DeclareMathOperator{\Ann}{Ann}
\DeclareMathOperator{\Zar}{Zar}


%%%%%%%%%%%%
%%%%% SGA4-2
\DeclareMathOperator{\Sex}{Sex}
\DeclareMathOperator{\vois}{vois}
\DeclareMathOperator{\Vois}{Vois}
\DeclareMathOperator{\Ext}{Ext}
\DeclareMathOperator{\ext}{ext}
\DeclareMathOperator{\Fer}{Fer}
\DeclareMathOperator{\Add}{Add}
\DeclareMathOperator{\Cocart}{Cocart}
\DeclareMathOperator{\cosk}{cosk}
\DeclareMathOperator{\plat}{plat}

\DeclareMathOperator{\faisc}{faisc}
\DeclareMathOperator{\sk}{sk}
\DeclareMathOperator{\ch}{ch}

\def\HH{\mathbf{H}}
\def\qc{\mathrm{qc}}
\def\et{\mathrm{et}}
\def\coh{\mathrm{coh}}

\def\Hvr#1{%
\begin{array}[b]{l}
{\kern2pt\scriptstyle r}\\[-8pt]
{\kern1pt\scriptstyle \vee}\\[-5pt]
H^{#1}
\end{array}}

\def\Hvv#1{%
\begin{array}[b]{l}
\kern2pt\rotatebox{90}{${\scriptscriptstyle\ll}$}\\[-6pt]
H^{#1}
\end{array}}

\def\Hvi#1{%
\begin{array}[b]{l}
{\kern0.2pt\scriptscriptstyle \infty}\\[-8pt]
{\kern1pt\scriptstyle \vee}\\[-5pt]
H^{#1}
\end{array}}

\def\Hv#1{%
\begin{array}[b]{l}
{\kern1pt\scriptstyle \vee}\\[-5pt]
H^{#1}
\end{array}}

\DeclareMathOperator{\cart}{cart}
\DeclareMathOperator{\Ffaisc}{\mathbf{Faisc}}


\def\Ev#1{%
\begin{array}[b]{l}
{\kern1pt\scriptstyle \vee}\\[-5pt]
E'_{#1}
\end{array}}

\DeclareMathOperator{\Cont}{Cont}
\DeclareMathOperator{\Ccont}{\mathscr{C}\!\mathit{ont}}
\DeclareMathOperator{\Site}{Site}
\newcommand\MM{\mathbf{M}}
\def\limtop#1{{\displaystyle{\varprojlim_{#1}\!\mbox{\rm top }}}}
\def\Limtop{\mathop{\varprojlim\!\mbox{\rm top}}}
\DeclareMathOperator{\spec}{spec}
\def\ss{\scriptstyle{\mathscr{S}}}
\DeclareMathOperator{\Gal}{Gal}
\def\smallm{\scriptstyle{\mathscr{M}}}
\DeclareMathOperator{\End}{End}
\DeclareMathOperator{\ass}{ass}
\DeclareMathOperator{\zar}{zar}
\DeclareMathOperator{\Ogg}{Ogg}
\DeclareMathOperator{\res}{res}
\DeclareMathOperator{\tr}{tr}
\DeclareMathOperator{\Supp}{Supp}
\DeclareMathOperator{\codim}{codim}
\DeclareMathOperator{\rk}{rk}
\DeclareMathOperator{\cur}{cur}
\DeclareMathOperator{\cdqc}{cdqc}
\newcommand\DD{\mathbf{D}}
\DeclareMathOperator{\cl}{cl}
\def\JJ{\mathbf{J}}
\def\FF{\mathbf{F}}
\def\BB{\mathbf{B}}
\DeclareMathOperator{\Alget}{Alget}
\DeclareMathOperator{\Ass}{Ass}
\DeclareMathOperator{\supp}{supp}
\DeclareMathOperator{\rad}{rad}
\DeclareMathOperator{\prof}{prof}
\DeclareMathOperator{\cint}{int}
\DeclareMathOperator{\sic}{sic}
\DeclareMathOperator{\Tor}{Tor}
\DeclareMathOperator{\tors}{tors}
\DeclareMathOperator{\Sym}{Sym}
\DeclareMathOperator{\parf}{parf}
\DeclareMathOperator{\tor}{tor}
\DeclareMathOperator{\defn}{defn}
\DeclareMathOperator{\Div}{Div}
\DeclareMathOperator{\Pol}{Pol}
\DeclareMathOperator{\Fex}{Fex}
\DeclareMathOperator{\Lie}{Lie}
\DeclareMathOperator{\tot}{tot}
\DeclareMathOperator{\Aut}{Aut}
\DeclareMathOperator{\HOM}{HOM}
\DeclareMathOperator{\pch}{pch}
\DeclareMathOperator{\Ch}{Ch}
\DeclareMathOperator{\EXT}{EXT}
\DeclareMathOperator{\PIC}{PIC}
\DeclareMathOperator{\Rf}{Rf}
\DeclareMathOperator{\inj}{inj}
\newcommand\VV{\mathbf{V}}

\def\Xv{%
\begin{array}[b]{l}
{\kern1pt\scriptstyle \vee}\\[-5pt]
X
\end{array}}

%%%%%%%
\newcommand\TT{\mathbf{T}}
\newcommand\XX{\mathbf{X}}
\newcommand\EE{\mathbf{E}}
%%%%%%

\def\qedsymbol{\textup{C.Q.F.D.}}
\let\CQFD\qed
\let\cqfd\qed

%%%%%\bigdot
\font\bigsymb=cmsy10 at 4pt
\def\bigdot{{\kern1.2pt\raise 1.5pt\hbox{\bigsymb\char15}}}
%%%%%

\def\Lright{\mathop{\vcenter{\hbox{\hbox{$\longrightarrow$}\lower8pt \hbox{\kern-16pt\hbox{$\longrightarrow$}}}}}}

\def\Aa{{\mathcal A}}
\def\Bb{{\mathcal B}}
\def\Cc{{\mathcal C}}
\def\Dd{{\mathcal D}}
\def\Ee{{\mathcal E}}
\def\Ff{{\mathcal F}}
\def\Gg{{\mathcal G}}
\def\Hh{{\mathcal H}}
\def\Ii{{\mathcal I}}
\def\Jj{{\mathcal J}}
\def\Kk{{\mathcal K}}
\def\Ll{{\mathcal L}}
\def\Mm{{\mathcal M}}
\def\Nn{{\mathcal N}}
\def\Oo{{\mathcal O}}
\def\Pp{{\mathcal P}}
\def\Qq{{\mathcal Q}}
\def\Rr{{\mathcal R}}
\def\Ss{{\mathcal S}}
\def\Tt{{\mathcal T}}
\def\Uu{{\mathcal U}}
\def\Vv{{\mathcal V}}
\def\Ww{{\mathcal W}}
\def\Xx{{\mathcal X}}
\def\Yy{{\mathcal Y}}
\def\Zz{{\mathcal Z}}

\newcommand{\CC}{\mathbf{C}}
\newcommand{\GG}{\mathbf{G}}
\newcommand{\KK}{\mathbf{K}}
\newcommand{\LL}{\mathbf{L}}
\newcommand{\NN}{\mathbf{N}}
\newcommand{\PP}{\mathbf{P}}
\newcommand\QQ{\mathbf{Q}}
\newcommand\RR{\mathbf{R}}
\newcommand\ZZ{\mathbf{Z}}

\def\to{\mathchoice{\longrightarrow}{\rightarrow}{\rightarrow}{\rightarrow}}
\def\from{\mathchoice{\longleftarrow}{\leftarrow}{\leftarrow}{\leftarrow}}
\def\mto{\mathchoice{\longmapsto}{\mapsto}{\mapsto}{\mapsto}}
\def\mfrom{\mathrel{\from\joinrel\mapstochar\,}}
\def\hto{\mathrel{\lhook\joinrel\to}}
\def\hfrom{\mathrel{\from\joinrel\rhook}}

\def\implies{\mathchoice{\Longrightarrow}{\Rightarrow}{\Rightarrow}{\Rightarrow}}
\def\implied{\mathchoice{\Longleftarrow}{\Leftarrow}{\Leftarrow}{\Leftarrow}}
\def\iff{\mathchoice{\Longleftrightarrow}{\Leftrightarrow}{\Leftrightarrow}{\Leftrightarrow}}

\def\To#1{\mathchoice{\xrightarrow{\labelstyle~~#1~}}{\stackrel{#1}{\longrightarrow}}{}{}}
\def\Mto#1{\mathrel{\mapstochar\To{#1}}}
\def\From#1{\mathchoice{\xleftarrow{\labelstyle~~#1~}}{\stackrel{#1}{\longleftarrow}}{}{}}

\newcommand{\isomto}{\overset{\sim}{\longrightarrow}}
\newcommand{\isomfrom}{\overset{\sim}{\longleftarrow}}

\newcommand{\IFF}{\unskip~$\iff$~\ignorespaces}
\newcommand{\THEN}{\unskip~$\implies$~\ignorespaces}

\let\leq\leqslant
\let\geq\geqslant
\let\le\leqslant
\let\ge\geqslant

\def\cf{{cf.}\kern.3em}
\def\Cf{{Cf.}\kern.3em}
\def\ie{{i.e.}\kern.3em}
\def\iev{{i.e.},\xspace}
%\def\loccito{loc.\kern2pt cit.\xspace}
%\def\loccitn{\textit{loc.\kern2pt cit.}\xspace}
%\def\loccit{\sisi{\loccito}{\loccitn}}
\def\EGA{\textup{EGA}\kern.2em}
\def\SGA{\textup{SGA}\kern.3em}

\let\oldS\S
\def\S{\oldS\kern .15em}
\def\ptbl{.\kern .2em}
\def\resp{resp.\kern.3em}
\def\Exp{Exp.\kern.3em}

\let\bar\overline
\let\hat\widehat
\let\tilde\widetilde
\let\dpl\displaystyle
\let\ldots\dots
\let\emptyset\varnothing
\let\skpt\skippointrait
\let\epsilon\varepsilon

\def\sfrac#1#2{{#1}/{#2}}

\DeclareMathOperator{\cd}{cd}
%\DeclareMathOperator{\Hhom}{\mathcal{H}\!\mathit{om}}
\DeclareMathOperator{\Hhom}{\mathbf{Hom}}
%\DeclareMathOperator{\Eext}{\mathcal{E}\!\mathit{xt}}
\DeclareMathOperator{\Eext}{\mathbf{Ext}}
\DeclareMathOperator{\Spec}{Spec}
\DeclareMathOperator{\Tr}{Tr}

\let\emptyset=\varnothing

%%%%%%%%%%%%%%%
%%% expose XVII
%%%%%%%%%%%%%%%

\makeatletter
\let\old@optexpose\optexpose
\newcommand{\empty@optexpose}[1]{}
\newcommand{\paren@optexpose}[1]{\old@optexpose{#1}(}
\newcommand{\refnoexpose}[1]{\let\optexpose\empty@optexpose\ref{#1}\let\optexpose\old@optexpose}%\refnoexpose{X.cor5.2} engendre "5.2" au lieu de "X 5.2"
\newcommand{\neweqref}[1]{\let\optexpose\paren@optexpose\textup{\ref{#1})}\let\optexpose\old@optexpose}%\neweqref{XII.eq4.1} engendre "XII (4.1)" au lieu de "(XII 4.1)"
\makeatother
\newcommand{\refvb}[2][VI${}_{\text{B}}$ ]{\sisi{#1\refnoexpose{#2}}{\ref{#2}}}

\newcommand{\leftexp}[2]{{\vphantom{#2}}^{#1}{#2}}
\newcommand{\AAA}{\mathbf{A}}
\newcommand{\Sf}{\mathfrak{S}}
\DeclareMathOperator*{\quotprojlim}{\text{``$\varprojlim$''}}
\DeclareMathOperator*{\quotinjlim}{\text{``$\varinjlim$''}}
\DeclareMathOperator{\Fl}{F\sisi{\ell}{l}}
\DeclareMathOperator{\Pl}{P\sisi{\ell}{l}}
%\DeclareMathOperator{\cl}{c\sisi{\ell}{l}}
\DeclareMathOperator{\tdf}{tdf}
\DeclareMathOperator{\rg}{rg}
\newcommand{\iInd}{\sisi{I}{i}nd}
\newcommand{\iind}{ind}
\newcommand{\ppro}{pro}
\newcommand{\tTor}{\sisi{T}{t}or}
\newcommand{\ttor}{tor}
\newcommand{\chap}{Chap.}
\newcommand{\Ffex}{\Ff\mathit{ex}}
\newcommand{\sheafk}{\underline{k}}
\newcommand{\sheafn}{\underline{n}}
\newcommand{\sheafGamma}{\underline{\Gamma}}

% Repeated references to an already typeset footnote must not use
% \footnotemark without a matching \footnotetext: Hyperref would then create
% a live annotation to a destination that cannot exist.  This helper keeps the
% current printed footnote mark and links it to an explicit stable note target.
\makeatletter
\newcommand{\sgarepeatedfootnotemark}[1]{%
  \begingroup
  \edef\@thefnmark{\thefootnote}%
  \hyperlink{#1}{\@makefnmark}%
  \endgroup}
\makeatother

\sisi{
\DeclareMathOperator{\simarrow}{\rule{.15mm}{1.8mm}\kern -3.4pt\rightsquigarrow}
\newcommand{\sheafell}{\underline{\ell}}
\newcommand{\sheafZZ}{\underline{\ZZ}}
}{
\newcommand{\nderef}[1]{la N.D.E.~(\ref{#1}) page~\pageref{#1}}
}

%%%%%%%%%%%%%%%%%
%%%%% expose XVII
%%%%%%%%%%%%%%%%%

% END INCLUDED SGA4_English_current_buildable_source_20260906/source/fichiers-configuration-latex/macros.tex


\externaldocument{../01/01_English.tex}
\externaldocument{../02/02_English.tex}
\externaldocument{../03/03_English.tex}
\externaldocument{../04/04_English.tex}
\externaldocument{../05/05_English.tex}
\externaldocument{../05bis/05bis_English.tex}
\externaldocument{../06/06_English.tex}
\externaldocument{../07/07_English.tex}
\externaldocument{../08/08_English.tex}
\externaldocument{../09/09_English.tex}
\externaldocument{../10/10_English.tex}
\externaldocument{../11/11_English.tex}
\externaldocument{../12/12_English.tex}

\setboolean{original}{false}

\begin{document}

\begin{center}
XIII. Base-change theorem for a proper morphism: conclusion of the proof\\
M. Artin
\end{center}

\selectlanguage{english}version: \showgitstatus
\tableofcontents

\else
\chapter{Base-change theorem for a proper morphism: conclusion of the proof}
%\addtocontents{toc}{M. Artin}
\begin{center}
M. Artin
\end{center}
\fi
\phantomsection\MakeLinkTarget*{sga4.exp.XIII}\label{sga4.exp.XIII}

% BEGIN INCLUDED SGA4_English_current_buildable_source_20260906/source/13/drafts/SGA4_XIII_LINES_43_316_English.texfrag
\setcounter{pageoriginal}{131}

\section{The projective and flat case}\etiquette[1]{XIII.sec1}\pageoriginale

Recall the statement of \hyperrefext[XII][XII.thm5.9]{5.9 bis} in the case
under consideration here:

\begin{enonce*}{Proposition 1.1}\etiquette[1.1]{XIII.prop1.1}
Let $S$ be the spectrum of a Noetherian henselian ring, and let
$f:X\rightarrow S$ be a projective and flat morphism. Let $X_{\circ}$ be
the closed fiber of $X/S$. The restriction functor
$$
\Et(X)\longrightarrow \Et(X_{\circ})
$$
is an equivalence of categories.
\end{enonce*}

As we have already observed, it follows from
\hyperrefext[XII][XII.prop6.5]{6.5} (i) that the functor is fully faithful.
It remains to prove that every étale covering $Y_{\circ}$ of $X_{\circ}$ is
induced by an étale covering $Y$ of $X$.

\begin{enonce*}{Lemma 1.2}\etiquette[1.2]{XIII.lem1.2}
Let $S$ be locally Noetherian, let $f:X\rightarrow S$ be projective and
flat, and let $Y\subset X$ be a subscheme. The subfunctor
$\epsilon=\prod_{X/S}Y/X$ of the terminal functor, \lq\lq which expresses
the condition $Y=X$\rq\rq, that is, such that for $S'\rightarrow S$ one has
$$
\epsilon(S')=
\begin{cases}
\phi & \mbox{if }\quad Y'\neq X',\\
\{\phi\} &\mbox{if }\quad Y'=X',
\end{cases}
$$
(where\pageoriginale a prime denotes the effect of the base change
$S'\rightarrow S$), is representable by a subscheme $Z$ of $S$. If
$Y$ is open (\resp closed), then so is $Z$.
\end{enonce*}

\noindent
\emph{Proof}. Let $Y$ be closed in an open subset $U$ of $X$, and let
$C$ be the closed complement of $U$. Since $f$ is proper, the image of
$C$ in $S$ is a closed subset $D$, and it is clear that the condition
$U=X$ is represented by the open subset $S-D$ of $S$. We may therefore
suppose that $U=X$ already, that is, that $Y$ is a closed subscheme.
Consider the surjective morphism
$\Oo_X\xrightarrow{\alpha}\Oo_Y$ of $\Oo_X$-modules; saying that $Y=X$
amounts to saying that there exists a morphism
$\Oo_Y\xrightarrow{\beta}\Oo_X$ such that the composite $\beta\alpha$ is
the identity, that is, that there exists a section of
$\Hhom(\Oo_Y,\Oo_X)$ lifting the identity section of
$\Hhom(\Oo_X,\Oo_X)$. Now it follows from \EGA III 7.7.8 that the functor
which assigns to $S'/S$ the set of sections of
$\Hhom(\Oo_Y,\Oo_X)$ (\resp $\Hhom(\Oo_X,\Oo_X)$) is represented by a
vector bundle $V_1$ (\resp $V_2$) over $S$. Let
$\varphi:V_1\rightarrow V_2$ be the morphism $u\mapsto u\alpha$ induced
by $\Oo_X\rightarrow\Oo_Y$, and let $s:S\rightarrow V_2$ be the section
corresponding to the identity $\Oo_X\rightarrow\Oo_X$. The inverse image
of the section $s$ under $\varphi$ is identified with a closed subscheme
$Z$ of $S$, and it is clear that $Z$ represents the functor $\epsilon$,
which proves the lemma.

\begin{enonce*}{Lemma 1.3}\etiquette[1.3]{XIII.lem1.3}
Let $S$ be a locally Noetherian scheme, let $f:X\rightarrow S$ be a
projective and flat morphism, and let $E$ be a locally free sheaf on $X$.
Let
$$
QE:(\Sch)/S\longrightarrow(\Ens)
$$
be the functor which assigns to $S'/S$ the set $QE(S')$ of quotient
sheaves $A'$ of $E'=E\otimes_{\Oo_S}\Oo_{S'}$, endowed with the structure
of an étale $\Oo_{X'}$-algebra. This functor is representable by a scheme
locally of finite type over $S$.
\end{enonce*}

\noindent
\emph{Proof}.\pageoriginale The functor $\sheaf{\Quot}(E)$, with
$\sheaf{\Quot}(E)(S')$ equal to the set of quotients of $E'$ that are flat
over $S'$, is representable by a scheme locally of finite type (TDTE IV,
exp.~221), say by $Z/S$. Let $M_u$ be the \lq\lq universal\rq\rq quotient
of $E\otimes_{\Oo_S}\Oo_Z$ over $X_Z$. The subfunctor of
$\sheaf{\Quot}(E)$ consisting of the quotients which are flat
\emph{over} $X'$ is represented by an open subset $U$ of $Z$. Indeed,
let $M'$ be a quotient of $E'$, given by an $S$-morphism $S'\rightarrow
Z$. Since $M'$ is flat over $S'$, saying that it is flat over $X'$ amounts
to saying that, for every point $s'\in S'$, the module $M'_{s'}$ is flat
over the fiber $X'_{s'}$ (\EGA IV 5.9), which is equivalent to
$M_u|_{X_z}$ being flat over the fiber $X_z$ at the point $z$ which is
the image of $s'$. But the set $U$ of $z\in Z$ such that
$M_u|_{X_z}$ is flat over $X_z$, that is, such that $M_u$ is flat over
$X_Z$ at every point of $X_z$, is open in $Z$ by \EGA IV 11.1.5,
proving the assertion.

It now suffices to represent the relative functor $QE/U[f]$. We are thus
reduced to the following lemma, after replacing $S$ by $U$ and $X$ by
$\fprod{X}{U}{S}$:

\begin{enonce*}{Lemma 1.4}\etiquette[1.4]{XIII.lem1.4}
Let $S$ be a locally Noetherian scheme, let $f:X\rightarrow S$ be a
projective and flat morphism, and let $M$ be a locally free
$\Oo_X$-module. The functor
$\Alget:(\Sch)/S\rightarrow(\Ens)$, which assigns to $S'/S$ the set
$\Alget(S')$ of algebra structures on $M'$ that are étale over
$\Oo_{X'}$, is representable by a scheme of finite type over $S$.
\end{enonce*}

\noindent
\emph{Proof}. For brevity, we shall say \lq\lq representable\rq\rq,
understanding \lq\lq by a relative scheme of \emph{finite type}\rq\rq.
All the vector bundles encountered below will be defined by
\emph{coherent} sheaves, and hence will be of finite type. A bilinear
composition law on $M'$ is given by a section of the locally free sheaf
$\Hhom(M'\otimes M',M')$. Since
$\Hhom(M'\otimes M',M')$ is flat over $S$, it follows from\pageoriginale
\EGA III 7.7.6 that the functor \lq\lq bilinear composition law\rq\rq is
represented by a vector bundle $Z$ over $S$. We must check that the
additional conditions that the law be associative, commutative, and
possess a unit section are represented by a subscheme of $Z$. By
considering the relative functor, we may replace $S$ by $Z$, and may
therefore suppose that a composition law on $M$ has been given and prove
that the subfunctor of $S$ expressing the preceding conditions is
representable by a subscheme of $S$.

For associativity, the condition is that two sections of
$\Hhom(M\otimes M\otimes M,M)$ become equal after the base change
$S'\rightarrow S$. Let $Y$ be the fiber product of these sections over
the vector bundle corresponding to
$\Hhom(M\otimes M\otimes M,M)$. The scheme $Y$ is identified with a
closed subscheme of $X$, and associativity after $S'\rightarrow S$ is
expressed by the condition $Y'=X'$. It therefore defines a subfunctor of
$S$ represented by a closed subscheme of $S$, by
\hyperref[XIII.lem1.2]{1.2}. For commutativity, the condition is that two
sections of $\Hhom(M\otimes M,M)$ become equal, which again defines a
subfunctor of $S$ represented by a closed subscheme, by
\hyperref[XIII.lem1.2]{1.2}. Finally, suppose, as we now may, that the law
is already associative and commutative, and consider the condition that a
unit section exist; such a section will then be unique. The functor whose
value on $S'/S$ is the set of sections of $M'$ is represented by a vector
bundle, by \EGA III 7.7.6, and it remains only to represent the subfunctor
of this bundle corresponding to the \emph{unit sections} (taking account
of the uniqueness just recalled). Passing to the relative situation as
usual, we may suppose that a section $\epsilon$ of $M$ has already been
given, and prove that the subfunctor of $S$ expressing the condition that
it become a\pageoriginale unit section is representable. But the condition
that $\epsilon$ be a unit is that the composite of the evident morphisms
$M'\rightarrow M'\otimes M'\rightarrow M'$ be the identity, that is, that
two specified sections of $\Hom(M,M)$ become equal. We conclude once
again by \hyperref[XIII.lem1.2]{1.2}.

Finally, suppose that a commutative unital algebra structure has been
given on the locally free module $M$. The locus where it is étale is open
in $X$ (\SGA I 4.5), and is therefore represented over $S$ by an (open)
subscheme. This completes the proof of \hyperref[XIII.lem1.4]{1.4}, and
hence of \hyperref[XIII.lem1.3]{1.3}.

\setcounter{subsection}{4}
\subsection{}\etiquette[1.5]{XIII.subsec1.5}
Let $Z/S$ now be the object representing the functor considered in
\hyperref[XIII.lem1.3]{1.3}, and consider the condition that $Z$ be smooth
over $S$ at a point $z$. Translating \SGA I III 3.1 into the language of
the functor $QE$, one obtains the following condition:

Let $B$ be a local Artinian ring, let $\Mm=\rad B$, let $J$ be an ideal of
$B$ such that $J\Mm=0$, put $S'=\Spec B$ and $S''=\Spec B/J$, and let
$S'\rightarrow S$ be a morphism. We use the evident notation for the
effects of the base changes $S'\rightarrow S$ and $S''\rightarrow S$.

\begin{itemize}
\item[$(\ast)$] For every quotient $A''$ of $E''$, endowed with the
structure of an étale $\Oo_{X''}$-algebra and corresponding to an
$S$-morphism $S''\rightarrow Z$ whose image is $z$, there exists a
quotient $A'$ of $E'$, endowed with the structure of an étale
$\Oo_{X'}$-algebra, which induces it.
\end{itemize}

Now, since $X'$ and $X''$ have the same underlying space, an algebra $A'$
inducing $A''$ is already determined up to unique isomorphism
(\SGA I 8.3). We\pageoriginale may therefore express condition
$({\ast}{\ast})$ as follows:

\begin{itemize}
\item[$({\ast}{\ast})$] Let $A'$ be an étale algebra over $\Oo_{X'}$,
and let $A''$ be the algebra induced by $A'$ on $X''$. Every surjective
homomorphism
$$
E''\xrightarrow{\varphi''}A''
$$
corresponding to a morphism $S''\rightarrow Z$ whose image is $z$ lifts
to a homomorphism (necessarily surjective by Nakayama's lemma)
$$
E'\xrightarrow{\varphi'}A'{\quoi}.
$$
\end{itemize}

From the exact sequence
$$
0\longrightarrow J\longrightarrow B\longrightarrow B/J\longrightarrow
0
$$
we obtain an exact sequence of sheaves on $X'$:
$$
0\longrightarrow J\otimes\Hhom(E',A')\longrightarrow
\Hhom(E',A')\longrightarrow\Hhom(E'',A'')\longrightarrow
0{\quoi}.
$$
The obstruction to lifting $\varphi''$ lies in
$H^1(X',J\otimes\Hhom(E',A'))$. Let
$S_{\circ}=\Spec B/\Mm$. Since $J\Mm=0$, one has
$$
J\otimes\Hhom(E',A')\simeq
J\otimes\Hhom(E_{\circ},A_{\circ}){\quoi},
$$
and it is clear that
$$
H^1(X',J\otimes\Hhom(E_{\circ},A_{\circ}))\simeq
J\otimes H^1(X_{\circ},\Hhom(E_{\circ},A_{\circ})){\quoi}.
$$
Taking account of the fact that cohomology commutes with extension of the
base fields, we obtain:

\begin{enonce*}{Corollary 1.6}\etiquette[1.6]{XIII.cor1.6}
Let $Z/S$ be the scheme representing the functor of
\hyperref[XIII.lem1.3]{1.3}, let
$X_Z=\fprod{X}{Z}{S}$, and let $A_Z$ be the (\lq\lq general\rq\rq)
quotient étale algebra of $E_Z$. Then $Z/S$ is smooth at every point
$z_{\circ}$ such that
$$
H^1(X_{\circ},\Hhom(E_{\circ},A_{\circ}))=0{\quoi},
$$
where\pageoriginale the subscript zero denotes restriction to the fiber
$X_{\circ}$ of $X_Z/Z$ at the point $z_{\circ}$.
\end{enonce*}

We can now complete the proof of \hyperref[XIII.prop1.1]{1.1}. With the
notation of the proposition, let $Y_{\circ}/X_{\circ}$ be an étale
covering, given by a sheaf of algebras $A_{\circ}$ that is étale over
$\Oo_{X_{\circ}}$. For $n$ sufficiently large, $A_{\circ}(n)$ is
generated by its sections (\EGA III 2.2.1). We thus have a surjection
$$
\Oo^N_{X_{\circ}}\longrightarrow A_{\circ}(n)\longrightarrow
0{\quoi},
$$
and hence a surjection
$$
\Oo^N_{X_{\circ}}(-n)=E_{\circ}\longrightarrow
A_{\circ}\longrightarrow 0{\quoi}.
$$
Choose $n$ sufficiently large that we also have
\begin{equation*}
H^1(X_{\circ},\Hhom(E_{\circ},A_{\circ}))=0{\quoi},\tag{$+$}
\end{equation*}
which is possible by \EGA III 2.2.1, since
$$
\Hhom(E_{\circ},A_{\circ})
=[\Hhom(\Oo^N_{X_{\circ}},A_{\circ})](n)
=[A_{\circ}(n)]^N{\quoi}.
$$
Put $E=\Oo^N_X(-n)$, and let $Z/S$ be the scheme representing the functor
considered in \hyperref[XIII.lem1.3]{1.3}. The quotient $A_{\circ}$ of
$E_{\circ}$ corresponds to a point $z_{\circ}\in Z$ in the closed fiber
of $Z/S$, rational over $k(s)$; by \hyperref[XIII.subsec1.5]{1.5} and
$(+)$, $Z/S$ is smooth at $z_{\circ}$. Since $S$ is the spectrum of a
\emph{henselian} ring, there is a section of $Z/S$ through $z_{\circ}$;
this follows, for example, immediately from the original definition
(\SGA II 1.1) of smoothness. It follows that $A_{\circ}$ is induced by an
algebra $A$ étale over $X$ (a quotient of $E$), which completes the proof
of \hyperref[XIII.prop1.1]{1.1}.

% END INCLUDED SGA4_English_current_buildable_source_20260906/source/13/drafts/SGA4_XIII_LINES_43_316_English.texfrag


% BEGIN INCLUDED SGA4_English_current_buildable_source_20260906/source/13/drafts/SGA4_XIII_LINES_317_549_English.texfrag
\section{The case of relative dimension at most one}\etiquette[2]{XIII.sec2}\pageoriginale

The assertion is the following.

\begin{enonce*}{Proposition 2.1}\etiquette[2.1]{XIII.prop2.1}
With the notation of \textup{\hyperrefext[XII][XII.thm5.9]{5.9 bis}}, the
functor $\Et(X)\rightarrow\Et(X_{\circ})$ is an equivalence of
categories if $S$ is Noetherian and strictly local and if $f$ is
projective of relative dimension $\leq1$.

By \textup{\hyperrefext[XII][XII.lem8.5]{8.5}}, this will complete the
proof of \textup{\hyperrefext[XII][XII.thm5.1]{5.1} (ii)}, which
generalizes both \hyperref[XIII.prop1.1]{1.1} and
\hyperref[XIII.prop2.1]{2.1}.
\end{enonce*}

\begin{enonce*}{Lemma 2.2}\etiquette[2.2]{XIII.lem2.2}
Let $S$ be the spectrum of a Noetherian local ring $A$, and let $X/S$
be projective of relative dimension $\leq n$. There exist a projective
and flat $S$-scheme $Z/S$ of relative dimension $\leq n$ and a closed
immersion $X\hookrightarrow Z$ over $S$.
\end{enonce*}

\noindent
\emph{Proof.} Let $X\rightarrow\PP^N_S$ be a projective embedding over
$S$, and let $J\subset A[x_{\circ},\ldots,x_N]$ be the homogeneous
ideal defining $X$. Put $k=A/\mathfrak m_A$, $S_{\circ}=\Spec k$, and
$X_{\circ}=X\otimes_Sk$. Then $X_{\circ}$ is defined by the
homogeneous ideal
$J_{\circ}=\Iim(J\rightarrow k[x_{\circ},\ldots,x_N])$. Since
$\dim X_{\circ}\leq n$, it is well known that there exist elements
$f^{\circ}_1,\ldots,f^{\circ}_{N-n}$ of $J_{\circ}$ defining a
subscheme $Z_{\circ}$ of $\PP^N_k$ of dimension $n$. Let
$f_i$ $(1\leq i\leq N-n)$ be elements of $J$ lifting the
$f^{\circ}_i$, and let $Z$ be the subscheme of $\PP^N_S$ defined by
$f_1,\ldots,f_{N-n}$. By \EGA IV 11.3, $Z$ is flat over $S$; and by
\EGA IV 13.1.5 it has relative dimension
$\dim Z_{\circ}=n$. This proves \hyperref[XIII.lem2.2]{2.2}.

\noindent
\emph{Proof of} \hyperref[XIII.prop2.1]{2.1}. Let $X/S$ be projective
of relative dimension $\leq1$, with $S$ Noetherian and strictly local.
We may assume that $X$ (and hence $X_{\circ}$) is connected. Let
$X\rightarrow Z$ be a closed immersion over $S$, where $Z$ is
projective, flat, and of relative dimension $\leq1$; see
\hyperref[XIII.lem2.2]{2.2}. Let $s_{\circ}$ be the closed point of
$S$. We have a corresponding\pageoriginale diagram of categories
$$
\xymatrix@C=2.5cm@R=1.25cm{
\Et(Z)\ar[r]\ar[d]^{a}&\Et(X)\ar[d]^{b}\\
\Et(Z_{\circ})\ar[r]^{c}&\Et(X_{\circ}).
}
$$
We wish to prove that $b$ is an equivalence. By
\hyperrefext[XII][XII.lem5.8]{5.8}, $b$ is fully faithful, so it remains
to prove that it is essentially surjective. The functor $a$ is an
equivalence by \hyperref[XIII.prop1.1]{1.1}; it therefore suffices to
prove that $c$ is essentially surjective. Dropping the subscript
${}_{\circ}$, we are reduced to the following proposition.

\begin{enonce*}{Proposition 2.3}\etiquette[2.3]{XIII.prop2.3}
Let $k$ be a separably closed field, let $Z$ be a $k$-scheme locally
of finite type and of dimension $\leq1$, and let $X\hookrightarrow Z$
be a connected closed subscheme. Then every étale covering of $X$ is
induced by an étale covering of $Z$.
\end{enonce*}

\noindent
\emph{Proof.} We may suppose that $Z$ and $X$ have dimension $1$, the
case in which $X$ is zero-dimensional being immediate from the
hypothesis on $k$. Then $Z=X\cup Y$, where $Y$ is the closed
subscheme obtained as the closure of $Z-X$. Put $V=X\cap Y$; this is
a discrete, that is, locally Artinian, scheme. The morphism
$X\amalg Y\rightarrow Z$ is finite and surjective, hence is a strict
descent morphism for the category of étale coverings
(\SGA \hyperrefext[IX][IX.cor4.7]{4.7}). Clearly
\begin{equation*}
\fprod{(X\amalg Y)}{(X\amalg Y)}{Z}
=X\amalg Y\amalg V^{(1)}\amalg V^{(2)}
\quad(V^{(i)}=V){\quoi},\tag{$\ast$}
\end{equation*}
where $\pr_i$ sends $V^{(i)}$ to the component $X$ of $X\amalg Y$
and sends $V^{(j)}$ to the component $Y$ when $j\neq i$.

Let $X'/X$ be an étale covering of degree $n$, and let $Y''/Y$ be a
covering of degree $n$ (for example, the trivial covering of degree
$n$), so that $X'\amalg Y''$ is\pageoriginale a covering of
$X\amalg Y$ of degree $n$. It is enough to find descent data for this
covering relative to $X\amalg Y\rightarrow Z$. Inspection of
$(\ast)$ shows that such descent data amount simply to an isomorphism
$$
\fprod{X'}{V}{X}=V'\xrightarrow{\sim}
V''=\fprod{Y''}{V}{Y}{\quoi}.
$$
Since $k$ is separably closed, so are the residue fields of the points
of $V$; consequently every étale covering of $V$ is completely split.
Thus $V'\simeq V''$, which proves the result.

\begin{enonce*}[remark]{Remark 2.4}\etiquette[2.4]{XIII.rem2.4}
One may avoid the delicate result from \SGA IX and descent theory by a
direct argument showing that the category of étale coverings of $Z$
is equivalent to the category of triples $(X',Y',\theta)$, where
$X'$ (respectively $Y'$) is an étale covering of $X$ (respectively
$Y$), and $\theta$ is a $V$-isomorphism
$\fprod{X'}{V}{X}\simeq\fprod{Y'}{V}{Y}$. This result, valid for any
scheme $Z$ that is the union of two closed subschemes $X$ and $Y$, is
proved directly from the analogous description of quasi-coherent
modules on $Z$ (when $X=V(\mathfrak J)$, $Y=V(\mathfrak K)$, and
$\mathfrak J\cap\mathfrak K=0$).
\end{enonce*}

\footnotetext[1]{This also appears in \EGA IV 21.9.11 and 21.9.12 in
a slightly more general form; in particular, in
\hyperref[XIII.cor3.2]{3.2} it is enough for $f$ to be \emph{separated}
rather than proper.}
\section{An auxiliary result on the Picard group ${}^{(*)}$}\etiquette[3]{XIII.sec3}

\begin{enonce*}{Proposition 3.1}\etiquette[3.1]{XIII.prop3.1}
Let $X$ be a Noetherian scheme and let $X_{\circ}$ be a closed
subscheme of $X$ such that \textup{(a)} $\dim X_{\circ}\leq1$, and
\textup{(b)} for every nonempty connected closed subset $Y$ of $X$,
$Y_{\circ}=Y\cap X_{\circ}$ is nonempty and connected.
\begin{enumerate}
\renewcommand{\theenumi}{\roman{enumi}}
\renewcommand{\labelenumi}{(\theenumi)}
\item Let $Z_{\circ}$ be the union of the sets
$\overline{\{x\}}\cap X_{\circ}$, for $x\in\Ass\Oo_X$, that consist
of a single point. For every Cartier divisor $D_{\circ}$ on
$X_{\circ}$ such that $\supp(D_{\circ})$ does not meet $Z_{\circ}$,
there exists a Cartier divisor $D$ on $X$ inducing $D_{\circ}$ (and if
$D_{\circ}\geq0$, one may take $D\geq0$).

\item Suppose there is an affine open subset $U_{\circ}$ of
$X_{\circ}$ containing the finite set
$\Ass\Oo_{X_{\circ}}\cup Z_{\circ}$ (this condition is automatically
satisfied if $X_{\circ}$ admits an ample invertible sheaf). Then
$\Pic(X)\rightarrow\Pic(X_{\circ})$ is surjective.
\end{enumerate}
\end{enonce*}

\noindent
\emph{Proof.}
\begin{enumerate}
\renewcommand{\theenumi}{\roman{enumi}}
\renewcommand{\labelenumi}{(\theenumi)}
\item Since $\dim X_{\circ}\leq1$, every Cartier divisor
$D_{\circ}$ on $X_{\circ}$ plainly has discrete support, and is
therefore the difference of two effective Cartier divisors whose
supports are contained in that of $D_{\circ}$. We are thus reduced in
(i) to the case where $D_{\circ}$ is effective. For every
$x\in\Supp D_{\circ}$ it is then defined at $x$ by a non-zero-divisor
$f_x\in\rad(\Oo_{X_{\circ},x})$. The element $f_x$ is induced by a
section $g_x$ of $\Oo_X$ on an open neighborhood $U_x$ of $x$. I claim
that, after shrinking $U_x$, the section $g_x$ is a non-zero-divisor
in $\Oo_{U_x}$, that is,
$V(g_x)\cap\Ass\Oo_{U_x}=\varnothing$.

Indeed, conditions (a) and (b) imply that, for every
$z\in\Ass\Oo_X$, the set $\overline{\{z\}}\cap X_{\circ}$ is a
connected closed subset of $X_{\circ}$ and hence is either a single
closed point or a connected union of one-dimensional irreducible
components of $X_{\circ}$. Since $x\notin Z_{\circ}$, in the first
case $x\notin\overline{\{z\}}$. In the second case, one cannot have
both $x\in\overline{\{z\}}$ and $z\in V(g_x)$, for then $g_x$, and
therefore $f_x$, would vanish on an irreducible component of
$\overline{\{z\}}\cap X_{\circ}$, which is also an irreducible
component of $X_{\circ}$; this is incompatible with $f_x$ being a
non-zero-divisor in $\Oo_{X_{\circ},x}$. After shrinking $U_x$, we may
therefore suppose\pageoriginale that
$V(g_x)\cap\Ass\Oo_{U_x}=\varnothing$, so $g_x$ is a
non-zero-divisor.

Let $Y_x$ be the closure of $V(g_x)$ in $X$. The point $x$ is isolated
in $Y_x\cap X_{\circ}\cap U_x$, hence also in
$Y_x\cap X_{\circ}$. By condition (b), $Y_x$ has a connected
component $Y'_x$ such that $Y'_x\cap X_{\circ}=\{x\}$. Moreover,
$Y'_x\subset U_x$: indeed,
$T_x=Y'_x-(Y'_x\cap U_x)$ is closed in $X$ and satisfies
$T_x\cap X_{\circ}=\varnothing$, so condition (b) gives
$T_x=\varnothing$. Let $D'_x$ be the Cartier divisor on $X$ induced
by $Y'_x$ on $U_x$ and by zero on $X-Y'_x$. Then
$D=\sum_xD'_x$ is an effective Cartier divisor on $X$ inducing
$D_{\circ}$ on $X_{\circ}$.

\item The conclusion follows from (i) and the fact that every
invertible module $\LL_{\circ}$ on $X_{\circ}$ is isomorphic to the
module defined by a Cartier divisor $D_{\circ}$ whose support does not
meet $Z_{\circ}$. This last assertion means that $\LL_{\circ}$ admits
a section on an open set $V_{\circ}$ containing
$\Ass\Oo_{X_{\circ}}\cup Z_{\circ}$ which vanishes at no point of
that set. It is enough to prove this after replacing $X_{\circ}$ by
the affine open $U_{\circ}$; there it follows from the familiar fact
that an invertible module on a semilocal scheme is isomorphic to the
structure sheaf. Finally, if $X_{\circ}$ admits an ample invertible
sheaf, every finite subset of $X_{\circ}$ is contained in an affine
open subset. This proves (ii).
\end{enumerate}

We have the following corollary, which generalizes
\hyperrefext[XII][XII.prop8.6]{8.6} and completes the proof of
\hyperrefext[XII][XII.thm5.1]{5.1} (iii).

\begin{enonce*}{Corollary 3.2}\etiquette[3.2]{XIII.cor3.2}
Let $S$ be the spectrum of a Noetherian Henselian ring, and let
$f:X\rightarrow S$ be a proper morphism of relative dimension
$\leq1$. For every closed subscheme $X'_{\circ}$ of $X$ having the
same underlying space as the closed fiber\pageoriginale $X_{\circ}$
of $X/S$, the canonical map
$\Pic(X)\rightarrow\Pic(X'_{\circ})$ is surjective.
\end{enonce*}

Indeed, apply the proposition to $(X,X'_{\circ})$. Condition (a) holds
by hypothesis, and condition (b) follows from
\hyperrefext[XII][XII.cor5.7]{5.7}. It remains to verify the existence
of $U_{\circ}$ required in (ii); this follows because $X'_{\circ}$ is
necessarily projective, as is every proper curve over a field.

\begin{enonce*}[remark]{Remark}
A reader who does not wish to assume or reconstruct the fact that
proper curves over a field are projective may use the corollary only
under the additional hypothesis that $f$ is projective; this suffices
for the application intended here. Conversely, admitting the
projectivity of proper algebraic curves over a field, one concludes
that under the hypotheses of \hyperref[XIII.cor3.2]{3.2}, $X$ is
necessarily projective over $S$. Indeed, choose an ample
$\xi_{\circ}\in\Pic(X_{\circ})$, lift it to some
$\xi\in\Pic(X)$, and apply \EGA III 4.7.1.
\end{enonce*}

\begin{thebibliography}{99}
\bibitem{sga4.exp.XIII.bib.TDTE} Grothendieck A., Technique de descente et
théorèmes d'existence en Géométrie Algébrique, in Fondements de la
Géométrie Algébrique (Extraits du Sém. Bourbaki 1957--62), Secrétariat
mathématique, 11, rue Pierre Curie.
\end{thebibliography}

\ifx\danslelivre\undefined
\end{document}
\fi

% END INCLUDED SGA4_English_current_buildable_source_20260906/source/13/drafts/SGA4_XIII_LINES_317_549_English.texfrag


% END INCLUDED SGA4_English_current_buildable_source_20260906/source/13/13_English.tex


\makeatletter
\def\input@path{{../14/}}
\makeatother
% BEGIN INCLUDED SGA4_English_current_buildable_source_20260906/source/14/14_English.tex
\ifx\danslelivre\undefined
% Release engine: XeLaTeX (the LuaTeX compatibility layer breaks Xy-pic output).
\documentclass[10pt]{../fichiers-configuration-latex/smfart}
% BEGIN INCLUDED SGA4_English_current_buildable_source_20260906/source/fichiers-configuration-latex/paquets.tex
\usepackage{polyglossia}
%\usepackage[francais]{babel}
\usepackage{fontspec}
\usepackage[math-style=ISO,bold-style=upright,sans-style=upright]{unicode-math}
%\usepackage[mathletters]{ucs}
%\usepackage[utf8]{inputenc}
%\usepackage[T1]{fontenc}
\usepackage[unicode,hypertexnames=false,colorlinks=true,linkcolor=blue,citecolor=Sepia,urlcolor=magenta]{hyperref}
\setmainlanguage{french}
\setmainfont[Ligatures=TeX]{Linux Libertine O}
% variante possible : TeXGyrePagella ?
\setmathfont{XITS Math}

% maths
\usepackage{amsmath}
%\usepackage{amssymb}
%\usepackage{mathrsfs}
%\usepackage{stmaryrd}
%\usepackage{euscript}

% couleurs
\usepackage[usenames,dvipsnames]{xcolor}

% graphiques
\usepackage{graphicx}
\usepackage[all]{xy}
%\usepackage{tikz}
%\usetikzlibrary{matrix,arrows,calc}
%\tikzset{column sep=3em,row sep=2em,text height=1.5ex, text depth=0.25ex,>=angle 90}
%\tikzset{description/.style={fill=white,inner sep=1pt}}

% hyperliens, 'etiquettes
\usepackage{xr-hyper}
\usepackage{../fichiers-configuration-latex/smfthm}
%\ifx\cacheretiquettes\undefined
%\usepackage[inline]{showlabels} \else \fi % pour rendre visibles les 'etiquettes
%\usepackage{../fichiers-configuration-latex/makeidx} % pour faire un index
%\usepackage[pagebackref,unicode]{hyperref}  % pour voir o`u sont cit'ees les r'ef'erences
% hyperref is loaded once above with the cumulative-reader destination policy.
%\hypersetup{breaklinks=true} %permet le retour `a la ligne dans les liens trop longs. Incompatible avec pdflatex
%\renewcommand*{\backref}[1]{{\small [cit\'e page(s)~#1.]}}

\usepackage{xspace}
\usepackage{pifont} % pour expose XVII

% END INCLUDED SGA4_English_current_buildable_source_20260906/source/fichiers-configuration-latex/paquets.tex

% BEGIN INCLUDED SGA4_English_current_buildable_source_20260906/source/fichiers-configuration-latex/numerotation.tex
\renewcommand{\thefootnote}{\roman{footnote}} % pour mettre les notes en bas de page sous la forme (ii) et non (2), qui ressemble trop (?) a un exposant.

\newcounter{pageoriginal}
\newcommand{\pageoriginale}{\refstepcounter{pageoriginal}\marginpar{\footnotesize\textbf{\thepageoriginal}}\xspace}
\let\pageoriginaled\pageoriginale

\usepackage{ifthen}
\newboolean{original}
\def\sisi{\ifthenelse{\boolean{original}}}
\setboolean{original}{false}

\newcounter{cnote}
\def\nde#1{\ifthenelse{\boolean{original}}{}{\stepcounter{cnote}\def\thefootnote{\arabic{cnote}}\protect\footnote[\value{cnote}]{N.D.E.: #1}}}
\def\ndemark{\ifthenelse{\boolean{original}}{}{\stepcounter{cnote}\def\thefootnote{\arabic{cnote}}\protect\footnotemark[\value{cnote}]}}
\def\ndetext#1{\ifthenelse{\boolean{original}}{}{\def\thefootnote{\arabic{cnote}}\protect\footnotetext[\value{cnote}]{N.D.E.: #1}}}

% END INCLUDED SGA4_English_current_buildable_source_20260906/source/fichiers-configuration-latex/numerotation.tex

% BEGIN INCLUDED SGA4_English_current_buildable_source_20260906/source/fichiers-configuration-latex/macros.tex
\setmainlanguage{french}
\setotherlanguage{english}


%% Fixed, engine-independent release identifier.  The inherited dynamic Git
%% probe used Lua io.popen, exposed build-state-dependent output, and printed
%% a visible question mark outside a Git checkout.
\def\showgitstatus{English release r2, 2026-07-23}


% macros pour LaTeX
\renewcommand{\mathsf}{}

% Police différente pour U+2113 SCRIPT SMALL L
\setmathfont[range="2113]{latinmodern-math.otf}

% Give every inherited semantic label its own stable PDF destination instead
% of inheriting whichever automatic equation/page anchor happened to be live.
% The legacy label remains as a compatibility alias used by existing edges.
\makeatletter
\def\etiquette[#1]#2{%
  \Hy@raisedlink{\hyper@anchorstart{sga4.semantic.#2}\hyper@anchorend}%
  \def\@currentHref{sga4.semantic.#2}%
  \def\@currentlabel{#1}%
  \label{#2}%
}
\makeatother
\def\eqhyperref[#1]#2{(\hyperref[#1]{#2})}
\def\pagehyperref[#1]#2{(page \hyperref[#1]{#2})}
\def\hyperrefext[#1][#2]#3{\hyperref[#2]{#1 #3}}
\def\eqhyperrefext[#1][#2]#3{\hyperref[#2]{#1 #3}}     %cf. exposes 17
\def\neweqhyperrefext[#1][#2]#3{\hyperref[#2]{#1 #3}}  %cf. exposes 17

\def\shat#1{\protect\widehat{#1}}
\def\stilde#1{\protect\widetilde{#1}}
%%%%%%%%%%%%%%%%%%


\def\etcn{etc.\xspace}
\def\etco{etc...\xspace}
\def\etc{\sisi{\etco}{\etcn}}
\def\ldotsn{\sisi{...}{\etcn}}

\def\cf{{cf.}\kern.3em}
\def\Cf{{Cf.}\kern.3em}
\def\ie{{i.e.}\kern.3em}
\def\iev{{i.e.},\xspace}
\def\loccito{loc.\kern2pt cit.\xspace}
\def\loccitn{\textit{loc.\kern2pt cit.}\xspace}
\def\loccit{\sisi{\loccito}{\loccitn}}
\def\EGA{\textup{EGA}\kern.2em}
\def\SGA{\textup{SGA}\kern.3em}

\def\quoi{}

\def\quav{\quad,\quad}
\def\quaet{\quad \text{et}\quad}
\def\fibreproduct#1#2#3{#1 ×_{#3} #2}
\let\fprod\fibreproduct
\def\ccheck#1#2{\mbox{${\check{\check{\rm #1}}}^{\lower2pt\hbox{$\scriptstyle #2$}}$}}
\def\sheaf#1{\protect\mathcal{#1}}
\def\vieuxsheaf#1{{#1}} % cf. problème avec sheaf{a\Phi_i} par exemple...
\DeclareMathOperator{\Hom}{Hom}
\DeclareMathOperator{\can}{can}
\DeclareMathOperator{\Ind}{Ind}
\DeclareMathOperator{\ind}{ind}
%\DeclareMathOperator{\Hhomtop}{\mathcal{H}\!\mathit{omtop}}
\DeclareMathOperator{\Hhomtop}{\mathbf{Homtop}}
\DeclareMathOperator{\Point}{Point}
\DeclareMathOperator{\Pro}{Pro}
\DeclareMathOperator{\Top}{Top}
\DeclareMathOperator{\Mod}{Mod}
\DeclareMathOperator{\Cart}{Cart}
%\DeclareMathOperator{\Mmorsite}{\mathcal{M}\!\mathit{orsite}}
\DeclareMathOperator{\Mmorsite}{\mathbf{Morsite}}
\DeclareMathOperator{\loc}{loc}
\DeclareMathOperator{\Loc}{Loc}
\DeclareMathOperator{\Pic}{Pic}
\DeclareMathOperator{\Et}{Et}
\DeclareMathOperator{\card}{card}
\DeclareMathOperator{\Ens}{Ens}
\DeclareMathOperator{\h}{\underline{H}}
\DeclareMathOperator{\Fonct}{Fonct}
\DeclareMathOperator{\Ker}{Ker}
\DeclareMathOperator{\but}{but}
\DeclareMathOperator{\Ob}{Ob}
\DeclareMathOperator{\ob}{ob}
\DeclareMathOperator{\Lib}{Lib}
\DeclareMathOperator{\esj}{esj}
\DeclareMathOperator{\id}{id}
\DeclareMathOperator{\Id}{Id}
\DeclareMathOperator{\ad}{ad}
\DeclareMathOperator{\Iim}{Im}
\DeclareMathOperator{\pr}{pr}
\DeclareMathOperator{\Cat}{Cat}
\DeclareMathOperator{\car}{car}
\DeclareMathOperator{\SysInd}{SysInd}
\DeclareMathOperator{\Sup}{Sup}
\DeclareMathOperator{\Coker}{Coker}
\DeclareMathOperator{\Filt}{Filt}
\DeclareMathOperator{\acc}{acc}
\DeclareMathOperator{\sur}{sur}
\DeclareMathOperator{\cat}{cat}
\DeclareMathOperator{\Ord}{Ord}
\DeclareMathOperator{\Gr}{Gr}
\DeclareMathOperator{\Cov}{Cov}
\DeclareMathOperator{\Caf}{Caf}
\DeclareMathOperator{\coker}{coker}
\DeclareMathOperator{\coim}{coim}
\DeclareMathOperator{\im}{im}
\DeclareMathOperator{\Ouv}{Ouv}
\DeclareMathOperator{\diag}{diag}
\DeclareMathOperator{\Esp}{Esp}
\DeclareMathOperator{\TOP}{TOP}
\DeclareMathOperator{\Unif}{Unif}
\DeclareMathOperator{\Iind}{\mathbf{Ind}}
\DeclareMathOperator{\Hhomcart}{\mathbf{Homcart}}
\DeclareMathOperator{\Hhomtopess}{\mathbf{Homtopess}}
\DeclareMathOperator{\Ppoints}{\mathbf{Points}}
\DeclareMathOperator{\Ppoint}{\mathbf{Point}}
\DeclareMathOperator{\Ttors}{\mathbf{Tors}}
\DeclareMathOperator{\Ssex}{\mathbf{Sex}}
\DeclareMathOperator{\Ffib}{\mathbf{Fib}}
\DeclareMathOperator{\Vvois}{\mathbf{Vois}}
\DeclareMathOperator{\Ppointess}{\mathbf{Pointess}}
\DeclareMathOperator{\scrHom}{\mathbf{Hom}}
\DeclareMathOperator{\sheafPic}{\mathbf{Pic}}
\DeclareMathOperator{\sheafAut}{\mathbf{Aut}}
\DeclareMathOperator{\Hhomtoplocan}{\mathbf{Homtoplocan}}
\DeclareMathOperator{\Hhomtopan}{\mathbf{Homtopan}}
\DeclareMathOperator{\Pptgeom}{\mathbf{Ptgéom}}
\DeclareMathOperator{\Morsite}{Morsite}
\DeclareMathOperator{\Fib}{Fib}
\DeclareMathOperator{\Fais}{Fais}
\DeclareMathOperator{\kar}{kar}
\DeclareMathOperator{\Groupes}{Groupes}
\DeclareMathOperator{\Restr}{Restr}
\DeclareMathOperator{\Isom}{Isom}
\DeclareMathOperator{\Prol}{Prol}
\DeclareMathOperator{\Res}{Res}
\DeclareMathOperator{\Tors}{Tors}
\DeclareMathOperator{\esp}{esp}
\DeclareMathOperator{\es}{es}
\DeclareMathOperator{\od}{od}
\DeclareMathOperator{\gr}{gr}
\DeclareMathOperator{\sob}{sob}
\DeclareMathOperator{\dfn}{dfn}
\DeclareMathOperator{\Sch}{Sch}
\DeclareMathOperator{\ouv}{ouv}
\DeclareMathOperator{\cons}{cons}
\DeclareMathOperator{\Nerf}{Nerf}
\DeclareMathOperator{\Points}{Points}
\DeclareMathOperator{\Espsob}{Espsob}
\DeclareMathOperator{\Homtop}{Homtop}
\DeclareMathOperator{\Pt}{Pt}
\DeclareMathOperator{\Quot}{Quot}
\DeclareMathOperator{\glob}{glob}
\DeclareMathOperator{\Homfib}{Homfib}
\DeclareMathOperator{\Qouv}{Qouv}
\DeclareMathOperator{\Quotouv}{Quotouv}
\DeclareMathOperator{\ens}{ens}
\DeclareMathOperator{\ab}{ab}
\DeclareMathOperator{\an}{an}
\DeclareMathOperator{\fs}{fs}
\DeclareMathOperator{\red}{red}
\DeclareMathOperator{\Ab}{Ab}
\DeclareMathOperator{\cosq}{cosq}
\DeclareMathOperator{\Ann}{Ann}
\DeclareMathOperator{\Zar}{Zar}


%%%%%%%%%%%%
%%%%% SGA4-2
\DeclareMathOperator{\Sex}{Sex}
\DeclareMathOperator{\vois}{vois}
\DeclareMathOperator{\Vois}{Vois}
\DeclareMathOperator{\Ext}{Ext}
\DeclareMathOperator{\ext}{ext}
\DeclareMathOperator{\Fer}{Fer}
\DeclareMathOperator{\Add}{Add}
\DeclareMathOperator{\Cocart}{Cocart}
\DeclareMathOperator{\cosk}{cosk}
\DeclareMathOperator{\plat}{plat}

\DeclareMathOperator{\faisc}{faisc}
\DeclareMathOperator{\sk}{sk}
\DeclareMathOperator{\ch}{ch}

\def\HH{\mathbf{H}}
\def\qc{\mathrm{qc}}
\def\et{\mathrm{et}}
\def\coh{\mathrm{coh}}

\def\Hvr#1{%
\begin{array}[b]{l}
{\kern2pt\scriptstyle r}\\[-8pt]
{\kern1pt\scriptstyle \vee}\\[-5pt]
H^{#1}
\end{array}}

\def\Hvv#1{%
\begin{array}[b]{l}
\kern2pt\rotatebox{90}{${\scriptscriptstyle\ll}$}\\[-6pt]
H^{#1}
\end{array}}

\def\Hvi#1{%
\begin{array}[b]{l}
{\kern0.2pt\scriptscriptstyle \infty}\\[-8pt]
{\kern1pt\scriptstyle \vee}\\[-5pt]
H^{#1}
\end{array}}

\def\Hv#1{%
\begin{array}[b]{l}
{\kern1pt\scriptstyle \vee}\\[-5pt]
H^{#1}
\end{array}}

\DeclareMathOperator{\cart}{cart}
\DeclareMathOperator{\Ffaisc}{\mathbf{Faisc}}


\def\Ev#1{%
\begin{array}[b]{l}
{\kern1pt\scriptstyle \vee}\\[-5pt]
E'_{#1}
\end{array}}

\DeclareMathOperator{\Cont}{Cont}
\DeclareMathOperator{\Ccont}{\mathscr{C}\!\mathit{ont}}
\DeclareMathOperator{\Site}{Site}
\newcommand\MM{\mathbf{M}}
\def\limtop#1{{\displaystyle{\varprojlim_{#1}\!\mbox{\rm top }}}}
\def\Limtop{\mathop{\varprojlim\!\mbox{\rm top}}}
\DeclareMathOperator{\spec}{spec}
\def\ss{\scriptstyle{\mathscr{S}}}
\DeclareMathOperator{\Gal}{Gal}
\def\smallm{\scriptstyle{\mathscr{M}}}
\DeclareMathOperator{\End}{End}
\DeclareMathOperator{\ass}{ass}
\DeclareMathOperator{\zar}{zar}
\DeclareMathOperator{\Ogg}{Ogg}
\DeclareMathOperator{\res}{res}
\DeclareMathOperator{\tr}{tr}
\DeclareMathOperator{\Supp}{Supp}
\DeclareMathOperator{\codim}{codim}
\DeclareMathOperator{\rk}{rk}
\DeclareMathOperator{\cur}{cur}
\DeclareMathOperator{\cdqc}{cdqc}
\newcommand\DD{\mathbf{D}}
\DeclareMathOperator{\cl}{cl}
\def\JJ{\mathbf{J}}
\def\FF{\mathbf{F}}
\def\BB{\mathbf{B}}
\DeclareMathOperator{\Alget}{Alget}
\DeclareMathOperator{\Ass}{Ass}
\DeclareMathOperator{\supp}{supp}
\DeclareMathOperator{\rad}{rad}
\DeclareMathOperator{\prof}{prof}
\DeclareMathOperator{\cint}{int}
\DeclareMathOperator{\sic}{sic}
\DeclareMathOperator{\Tor}{Tor}
\DeclareMathOperator{\tors}{tors}
\DeclareMathOperator{\Sym}{Sym}
\DeclareMathOperator{\parf}{parf}
\DeclareMathOperator{\tor}{tor}
\DeclareMathOperator{\defn}{defn}
\DeclareMathOperator{\Div}{Div}
\DeclareMathOperator{\Pol}{Pol}
\DeclareMathOperator{\Fex}{Fex}
\DeclareMathOperator{\Lie}{Lie}
\DeclareMathOperator{\tot}{tot}
\DeclareMathOperator{\Aut}{Aut}
\DeclareMathOperator{\HOM}{HOM}
\DeclareMathOperator{\pch}{pch}
\DeclareMathOperator{\Ch}{Ch}
\DeclareMathOperator{\EXT}{EXT}
\DeclareMathOperator{\PIC}{PIC}
\DeclareMathOperator{\Rf}{Rf}
\DeclareMathOperator{\inj}{inj}
\newcommand\VV{\mathbf{V}}

\def\Xv{%
\begin{array}[b]{l}
{\kern1pt\scriptstyle \vee}\\[-5pt]
X
\end{array}}

%%%%%%%
\newcommand\TT{\mathbf{T}}
\newcommand\XX{\mathbf{X}}
\newcommand\EE{\mathbf{E}}
%%%%%%

\def\qedsymbol{\textup{C.Q.F.D.}}
\let\CQFD\qed
\let\cqfd\qed

%%%%%\bigdot
\font\bigsymb=cmsy10 at 4pt
\def\bigdot{{\kern1.2pt\raise 1.5pt\hbox{\bigsymb\char15}}}
%%%%%

\def\Lright{\mathop{\vcenter{\hbox{\hbox{$\longrightarrow$}\lower8pt \hbox{\kern-16pt\hbox{$\longrightarrow$}}}}}}

\def\Aa{{\mathcal A}}
\def\Bb{{\mathcal B}}
\def\Cc{{\mathcal C}}
\def\Dd{{\mathcal D}}
\def\Ee{{\mathcal E}}
\def\Ff{{\mathcal F}}
\def\Gg{{\mathcal G}}
\def\Hh{{\mathcal H}}
\def\Ii{{\mathcal I}}
\def\Jj{{\mathcal J}}
\def\Kk{{\mathcal K}}
\def\Ll{{\mathcal L}}
\def\Mm{{\mathcal M}}
\def\Nn{{\mathcal N}}
\def\Oo{{\mathcal O}}
\def\Pp{{\mathcal P}}
\def\Qq{{\mathcal Q}}
\def\Rr{{\mathcal R}}
\def\Ss{{\mathcal S}}
\def\Tt{{\mathcal T}}
\def\Uu{{\mathcal U}}
\def\Vv{{\mathcal V}}
\def\Ww{{\mathcal W}}
\def\Xx{{\mathcal X}}
\def\Yy{{\mathcal Y}}
\def\Zz{{\mathcal Z}}

\newcommand{\CC}{\mathbf{C}}
\newcommand{\GG}{\mathbf{G}}
\newcommand{\KK}{\mathbf{K}}
\newcommand{\LL}{\mathbf{L}}
\newcommand{\NN}{\mathbf{N}}
\newcommand{\PP}{\mathbf{P}}
\newcommand\QQ{\mathbf{Q}}
\newcommand\RR{\mathbf{R}}
\newcommand\ZZ{\mathbf{Z}}

\def\to{\mathchoice{\longrightarrow}{\rightarrow}{\rightarrow}{\rightarrow}}
\def\from{\mathchoice{\longleftarrow}{\leftarrow}{\leftarrow}{\leftarrow}}
\def\mto{\mathchoice{\longmapsto}{\mapsto}{\mapsto}{\mapsto}}
\def\mfrom{\mathrel{\from\joinrel\mapstochar\,}}
\def\hto{\mathrel{\lhook\joinrel\to}}
\def\hfrom{\mathrel{\from\joinrel\rhook}}

\def\implies{\mathchoice{\Longrightarrow}{\Rightarrow}{\Rightarrow}{\Rightarrow}}
\def\implied{\mathchoice{\Longleftarrow}{\Leftarrow}{\Leftarrow}{\Leftarrow}}
\def\iff{\mathchoice{\Longleftrightarrow}{\Leftrightarrow}{\Leftrightarrow}{\Leftrightarrow}}

\def\To#1{\mathchoice{\xrightarrow{\labelstyle~~#1~}}{\stackrel{#1}{\longrightarrow}}{}{}}
\def\Mto#1{\mathrel{\mapstochar\To{#1}}}
\def\From#1{\mathchoice{\xleftarrow{\labelstyle~~#1~}}{\stackrel{#1}{\longleftarrow}}{}{}}

\newcommand{\isomto}{\overset{\sim}{\longrightarrow}}
\newcommand{\isomfrom}{\overset{\sim}{\longleftarrow}}

\newcommand{\IFF}{\unskip~$\iff$~\ignorespaces}
\newcommand{\THEN}{\unskip~$\implies$~\ignorespaces}

\let\leq\leqslant
\let\geq\geqslant
\let\le\leqslant
\let\ge\geqslant

\def\cf{{cf.}\kern.3em}
\def\Cf{{Cf.}\kern.3em}
\def\ie{{i.e.}\kern.3em}
\def\iev{{i.e.},\xspace}
%\def\loccito{loc.\kern2pt cit.\xspace}
%\def\loccitn{\textit{loc.\kern2pt cit.}\xspace}
%\def\loccit{\sisi{\loccito}{\loccitn}}
\def\EGA{\textup{EGA}\kern.2em}
\def\SGA{\textup{SGA}\kern.3em}

\let\oldS\S
\def\S{\oldS\kern .15em}
\def\ptbl{.\kern .2em}
\def\resp{resp.\kern.3em}
\def\Exp{Exp.\kern.3em}

\let\bar\overline
\let\hat\widehat
\let\tilde\widetilde
\let\dpl\displaystyle
\let\ldots\dots
\let\emptyset\varnothing
\let\skpt\skippointrait
\let\epsilon\varepsilon

\def\sfrac#1#2{{#1}/{#2}}

\DeclareMathOperator{\cd}{cd}
%\DeclareMathOperator{\Hhom}{\mathcal{H}\!\mathit{om}}
\DeclareMathOperator{\Hhom}{\mathbf{Hom}}
%\DeclareMathOperator{\Eext}{\mathcal{E}\!\mathit{xt}}
\DeclareMathOperator{\Eext}{\mathbf{Ext}}
\DeclareMathOperator{\Spec}{Spec}
\DeclareMathOperator{\Tr}{Tr}

\let\emptyset=\varnothing

%%%%%%%%%%%%%%%
%%% expose XVII
%%%%%%%%%%%%%%%

\makeatletter
\let\old@optexpose\optexpose
\newcommand{\empty@optexpose}[1]{}
\newcommand{\paren@optexpose}[1]{\old@optexpose{#1}(}
\newcommand{\refnoexpose}[1]{\let\optexpose\empty@optexpose\ref{#1}\let\optexpose\old@optexpose}%\refnoexpose{X.cor5.2} engendre "5.2" au lieu de "X 5.2"
\newcommand{\neweqref}[1]{\let\optexpose\paren@optexpose\textup{\ref{#1})}\let\optexpose\old@optexpose}%\neweqref{XII.eq4.1} engendre "XII (4.1)" au lieu de "(XII 4.1)"
\makeatother
\newcommand{\refvb}[2][VI${}_{\text{B}}$ ]{\sisi{#1\refnoexpose{#2}}{\ref{#2}}}

\newcommand{\leftexp}[2]{{\vphantom{#2}}^{#1}{#2}}
\newcommand{\AAA}{\mathbf{A}}
\newcommand{\Sf}{\mathfrak{S}}
\DeclareMathOperator*{\quotprojlim}{\text{``$\varprojlim$''}}
\DeclareMathOperator*{\quotinjlim}{\text{``$\varinjlim$''}}
\DeclareMathOperator{\Fl}{F\sisi{\ell}{l}}
\DeclareMathOperator{\Pl}{P\sisi{\ell}{l}}
%\DeclareMathOperator{\cl}{c\sisi{\ell}{l}}
\DeclareMathOperator{\tdf}{tdf}
\DeclareMathOperator{\rg}{rg}
\newcommand{\iInd}{\sisi{I}{i}nd}
\newcommand{\iind}{ind}
\newcommand{\ppro}{pro}
\newcommand{\tTor}{\sisi{T}{t}or}
\newcommand{\ttor}{tor}
\newcommand{\chap}{Chap.}
\newcommand{\Ffex}{\Ff\mathit{ex}}
\newcommand{\sheafk}{\underline{k}}
\newcommand{\sheafn}{\underline{n}}
\newcommand{\sheafGamma}{\underline{\Gamma}}

% Repeated references to an already typeset footnote must not use
% \footnotemark without a matching \footnotetext: Hyperref would then create
% a live annotation to a destination that cannot exist.  This helper keeps the
% current printed footnote mark and links it to an explicit stable note target.
\makeatletter
\newcommand{\sgarepeatedfootnotemark}[1]{%
  \begingroup
  \edef\@thefnmark{\thefootnote}%
  \hyperlink{#1}{\@makefnmark}%
  \endgroup}
\makeatother

\sisi{
\DeclareMathOperator{\simarrow}{\rule{.15mm}{1.8mm}\kern -3.4pt\rightsquigarrow}
\newcommand{\sheafell}{\underline{\ell}}
\newcommand{\sheafZZ}{\underline{\ZZ}}
}{
\newcommand{\nderef}[1]{la N.D.E.~(\ref{#1}) page~\pageref{#1}}
}

%%%%%%%%%%%%%%%%%
%%%%% expose XVII
%%%%%%%%%%%%%%%%%

% END INCLUDED SGA4_English_current_buildable_source_20260906/source/fichiers-configuration-latex/macros.tex


\externaldocument{../01/01_English.tex}
\externaldocument{../02/02_English.tex}
\externaldocument{../03/03_English.tex}
\externaldocument{../04/04_English.tex}
\externaldocument{../05/05_English.tex}
\externaldocument{../05bis/05bis_English.tex}
\externaldocument{../06/06_English.tex}
\externaldocument{../07/07_English.tex}
\externaldocument{../08/08_English.tex}
\externaldocument{../09/09_English.tex}
\externaldocument{../10/10_English.tex}
\externaldocument{../11/11_English.tex}
\externaldocument{../12/12_English.tex}
\externaldocument{../13/13_English.tex}

\setboolean{original}{false}

\begin{document}

\begin{center}
XIV. Finiteness theorem for a proper morphism; cohomological dimension of affine algebraic schemes\\
M. Artin
\end{center}

\selectlanguage{english}version: \showgitstatus
\tableofcontents

\else
\chapter[Proper finiteness and affine cohomological dimension]{Finiteness theorem for a proper morphism; cohomological dimension of affine algebraic schemes}
%\addtocontents{toc}{M. Artin}
\begin{center}
M. Artin
\end{center}
\fi
\phantomsection\MakeLinkTarget*{sga4.exp.XIV}\label{sga4.exp.XIV}

% BEGIN INCLUDED SGA4_English_current_buildable_source_20260906/source/14/drafts/SGA4_XIV_LINES_43_300_English.texfrag
\setcounter{pageoriginal}{144}

\newcommand{\degtr}{{\rm deg.tr.\ }}

This\pageoriginale exposé contains two important theorems whose proofs
use the base-change theorem for a proper morphism.

\section{Finiteness theorem for a proper morphism}\etiquette[1]{XIV.sec1}

\begin{enonce*}{Theorem 1.1}\etiquette[1.1]{XIV.thm1.1}
Let $f:X\rightarrow S$ be a proper morphism of finite presentation. Let
$A$ be a left Noetherian ring which is a $\ZZ/n$-algebra for some
suitable $n\in\NN$. Let $F$ be a \emph{constructible} sheaf of sets
(\resp groups, \resp $A$-modules) on $X$. Then the sheaf $f_*F$
(\resp $R^1f_*F$, \resp $R^qf_*F$ for every $q\geq0$) is also
constructible.
\end{enonce*}

In particular, we have:

\begin{enonce*}{Corollary 1.2}\etiquette[1.2]{XIV.cor1.2}
Let $X$ be a proper scheme over the spectrum of a separably closed
field $k$, and let $F$ be a constructible torsion abelian sheaf on
$X$. Then the groups $H^q(X,F)$ are finite for every $q\geq0$.
\end{enonce*}

\begin{enonce*}[remark]{Remark 1.3}\etiquette[1.3]{XIV.rem1.3}
The theorem is false for abelian sheaves if the condition that $f$ be
proper is omitted, as is seen by taking $F=\ZZ/p$, $p=\car k$,
$X=\EE^1_k$, the affine line, and $f$ either the inclusion of $X$ in
$\PP^1_k$ or the structural morphism $X\rightarrow\Spec k$.
Nevertheless, it is conjectured\nde{Ofer Gabber presented the following
generalization of \hyperref[XIV.thm1.1]{1.1} at the conferences held in
honor of Luc Illusie and Pierre Deligne in 2005. Let $f:X\to Y$ be a
morphism of finite type between quasi-excellent Noetherian schemes
(\EGA IV 7.8), and let $F$ be a constructible sheaf of
$\mathbf Z/n$-modules on $X_{\text{ét}}$, with $n$ invertible on $Y$.
Then the sheaves $R^if_*F$ on $Y_{\text{ét}}$ are constructible and
vanish for all but finitely many $i$. There are analogous assertions
for constructible sheaves of sets \resp groups of order invertible on
$Y$. \newline
Before Gabber's result, \hyperref[XIV.thm1.1]{Theorem 1.1} was known in
the following case. Let $S$ be a regular scheme of dimension $\leq1$,
let $f:X\to Y$ be a morphism of finite type between $S$-schemes, and
let $F$ be a sheaf of left $A$-modules on $X_{\text{ét}}$. Then the
sheaves $R^if_*F$ are constructible (\SGA 4 1/2, Finitude 1.1).}
that it is true if $F$ has order prime to the\pageoriginale residue
characteristics of $S$ and $S$ is \lq\lq excellent\rq\rq\ (\EGA IV
7.8). When resolution of singularities is available (for example, in
characteristic zero), this can be proved for an excellent
equicharacteristic scheme $S$
(\hyperrefext[XIX][XIX.thm5.1]{5.1}).
\end{enonce*}

\noindent
\emph{Proof of} \hyperref[XIV.thm1.1]{1.1}. The assertion is local on
$S$, so we may take $S$ affine. Then $S$ is a limit of spectra of rings
of finite type over $\ZZ$, and the data of the theorem can be obtained,
by base change $S\rightarrow S_0$, from a proper morphism
$f_0:X_0\rightarrow S_0$, with $S_0$ of finite type over $\ZZ$, and a
constructible sheaf $F_0$ on $X_0$; see \EGA IV 8 and
\hyperrefext[IX][IX.cor2.7.4]{2.7.4}. Since the inverse image of a
constructible sheaf is again constructible
(\hyperrefext[IX][IX.prop2.4]{2.4} (iii)), application of
\hyperrefext[XII][XII.thm5.1]{5.1} to the base-change morphism
$S_0\leftarrow S$ reduces us to proving the constructibility of
$f_{0*}F_0$ (\resp $R^qf_{0*}F_0$), whence the following lemma.

\begin{enonce*}{Lemma 1.4}\etiquette[1.4]{XIV.lem1.4}
It is enough to verify \hyperref[XIV.thm1.1]{1.1} when $S$ is of finite
type over $\Spec\ZZ$.
\end{enonce*}

\begin{enonce*}{Lemma 1.5}\etiquette[1.5]{XIV.lem1.5}
The set-theoretic assertion of \hyperref[XIV.thm1.1]{1.1} is true: that
is, $f_*F$ is a constructible sheaf of sets if $F$ is one.
\end{enonce*}

\noindent
\emph{Proof}. We may suppose that $S$ is Noetherian
(\hyperref[XIV.lem1.4]{1.4}). Let
$F\rightarrow G=\prod_i\pi_{i*}(C_i)$ be an injection, with
$\pi_i:X_i\rightarrow X$ finite and $C_i$ constant and constructible on
$X_i$ (\hyperrefext[IX][IX.prop2.14]{2.14}). We have
$f_*F\subset f_*G$; hence, by \hyperrefext[IX][IX.prop2.9]{2.9} (ii),
it is enough to treat the case $F=G$. Now
$f_*G=\prod_i(f\pi_i)_*C_i$, and a finite product of constructible
sheaves on $S$ is constructible
(\hyperrefext[IX][IX.prop2.6]{2.6} (ibis)). Replacing $X$ by $X_i$, we
are reduced to the case where $F=D_X$ is a constructible constant
sheaf with value $D$. By \hyperrefext[XII][XII.thm5.1]{5.1} (i), the
fiber $(f_*F)_{\bar s}$ is isomorphic to $D^{C(s)}$, where
$C(s)=\pi_0(X_{\bar s})$;\pageoriginale it is therefore finite. By
\hyperrefext[IX][IX.prop2.13]{2.13} (iii), it is enough to prove that
the function \lq\lq number of elements in the fiber of
$f_*F$\rq\rq\ is constructible on $S$. This function is
$\card(D^C)$, where $C$ is the function \lq\lq number of connected
components in the geometric fiber\rq\rq, and $C$ is constructible
(\EGA IV 9.7.9), which proves the result.

\begin{enonce*}{Lemma 1.6}\etiquette[1.6]{XIV.lem1.6}
Let $S$ be a Noetherian scheme, let $f:X\rightarrow S$ be proper, let
$F$ be a constructible sheaf of groups (\resp $A$-modules) on $X$, and
let $q$ be an integer equal to $0$ or $1$ (\resp an integer $\geq0$).
Let $S_i$ $(1\leq i\leq n)$ be subschemes of $S$ such that
$S=\bigcup_iS_i$, and let $F_i$ and $f_i:X_i\rightarrow S_i$ be the
\lq\lq restrictions\rq\rq\ of the data to $S_i$. If
$R^qf_{i*}F_i$ is constructible for every $i$, then $R^qf_*F$ is
constructible as well.
\end{enonce*}

\noindent
\emph{Proof}. The restriction of $R^qf_*F$ to $S_i$ is isomorphic to
$R^qf_{i*}F_i$ by \hyperrefext[XII][XII.thm5.1]{5.1}, and the lemma
follows from \hyperrefext[IX][IX.prop2.8]{2.8}.

We now proceed by arguments analogous to those of
\hyperrefext[XII][XII.sec8]{8}:

\begin{enonce*}{Lemma 1.7}\etiquette[1.7]{XIV.lem1.7}
Let $S$ be Noetherian, and let
$$
\xymatrix@=1.25cm{
X        \ar[rd]_{f} &           &\bar{X} \ar[ll]_{\pi}
 \ar[ld]^{\bar{f}}\\
                 &    S      &
}
$$
be a commutative diagram of proper morphisms, and let
$j:U\rightarrow X$ be an open subscheme such that the open subscheme
$\bar U=\fprod{U}{\bar X}{X}$ maps isomorphically onto $U$. Let
$Y=X-U$ be the reduced closed subscheme, and let $f_0:Y\rightarrow S$
and $i:Y\rightarrow X$ be the canonical morphisms. Suppose that
\hyperref[XIV.thm1.1]{Theorem 1.1}
holds\nde{The hypothesis concerning $R^q\pi_*$ occurs in the proof only
when $q=0$. In that case constructibility is already known from
\hyperref[XIV.lem1.5]{1.5}.}
for $\bar f$, $f_0$, and $\pi$. Then it also holds for $f$.
\end{enonce*}

\noindent
\emph{Proof}.\pageoriginale First consider a sheaf of $A$-modules $F$
on $X$. We have the exact sequence
$$
0\longrightarrow j_!j^*F\longrightarrow F\longrightarrow
i_*i^*F\longrightarrow0{\quoi}.
$$
In view of the corresponding cohomology exact sequence for the
$R^qf_*$ and \hyperrefext[IX][IX.prop2.6]{2.6}, it is enough to verify
constructibility for the two outer terms. Notice that $i_*i^*F$ is
constructible by \hyperrefext[IX][IX.prop2.4]{2.4} (iii) and
\hyperrefext[IX][IX.prop2.14]{2.14} (i), and hence so is $j_!j^*F$
(\hyperrefext[IX][IX.prop2.6]{2.6}). Since $i$ is a closed immersion,
$$
(R^qf_*)i_*(i^*F)\simeq(R^qf_{0*})(i^*F){\quoi},
$$
by \hyperrefext[VIII][VIII.cor5.6]{5.6}; constructibility in this case
therefore follows from the hypothesis on $f_0$. Moreover, if
$\bar j:\bar U\rightarrow\bar X$ denotes the inclusion morphism and
$\bar F=\pi^*F$, then
\begin{equation*}
j_!j^*F\simeq\pi_*\bar j_!\bar j^*\bar F{\quoi},\tag{$*$}
\end{equation*}
as is seen by first defining a morphism from left to right and then
proving it to be invertible fiber by fiber using
\hyperrefext[XII][XII.thm5.1]{5.1}. At every geometric point $\bar y$
of $Y$, the restriction of $\bar j_!\bar j^*\bar F$ to the fiber
$\bar X_{\bar y}$ of $\bar X$ over $X$ is zero; hence, by
\hyperrefext[XII][XII.thm5.1]{5.1} (iii), the fiber of
$(R^q\pi_*)(\bar j_!\bar j^*\bar F)$ at $\bar y$ is zero for every
$q\geq0$. On the other hand, $\pi$ induces an isomorphism over $X-Y$,
so
$$
(R^q\pi_*)(\bar j_!\bar j^*\bar F)=0
$$
for $q>0$. The Leray spectral sequence
$$
E^{p,q}_2=(R^pf_*)(R^q\pi_*)(\bar j_!\bar j^*\bar F)
\Longrightarrow R^n\bar f_*(\bar j_!\bar j^*\bar F)
$$
together with $(*)$ therefore gives isomorphisms
$$
(R^q\bar f_*)(\bar j_!\bar j^*\bar F)\simeq
(R^qf_*)(j_!j^*F)
$$
for $q\geq0$. Since $\bar j_!\bar j^*\bar F$ is constructible (by the
same argument as for $j_!j^*F$), the left-hand side is constructible by
the hypothesis on $\bar f$. Thus $(R^qf_*)(j_!j^*F)$ is constructible,
and consequently so is $R^qf_*F$.\pageoriginale

Now suppose that $F$ is a constructible sheaf of groups on $X$. Replace
$\pi:\bar X\rightarrow X$ by
$\pi\amalg i:(\bar X\amalg Y)\rightarrow X$. It is then enough to
prove that if $\pi:\bar X\rightarrow X$ is \emph{surjective} and
\hyperref[XIV.thm1.1]{1.1} holds for $\pi$ and $\bar f$, then it also
holds for $f$. First let $\bar F$ be a constructible sheaf of groups on
$\bar X$. Then $\pi_*\bar F$ is constructible on $X$
(\hyperref[XIV.lem1.5]{1.5}), and there is an injection
$$
(R^1f_*)\pi_*\bar F\hookrightarrow(R^1\bar f_*)\bar F.
$$
The right-hand side is constructible by hypothesis, and hence so is the
left-hand side (\hyperrefext[IX][IX.prop2.9]{2.9} (ii)). Since $\pi$ is
surjective, there is an injection
$F\rightarrow\pi_*\pi^*F=G$. The sheaf $\bar F=\pi^*F$ is
constructible, so $G$ and $R^1f_*G$ are constructible as well, by what
we have just proved. We are thus reduced to the following lemma.

\begin{enonce*}{Lemma 1.8}\etiquette[1.8]{XIV.lem1.8}
Let $S$ be Noetherian, let $f:X\rightarrow S$ be a proper morphism, and
let
$\displaystyle F\mathop{\hookrightarrow}^{u}G$
be an injection of constructible sheaves of groups on $X$. If
$R^1f_*G$ is constructible, then $R^1f_*F$ is constructible as well.
\end{enonce*}

\noindent
\emph{Proof}. Let $C=G/F$, which is a constructible sheaf of
homogeneous spaces under $G$ (apply
\hyperrefext[IX][IX.prop2.6]{2.6}). Consider the exact sequence
\begin{equation*}
f_*C\xrightarrow{\delta}R^1f_*F\xrightarrow{u^1}R^1f_*G{\quoi}.
\tag{$*$}
\end{equation*}
To prove that $R^1f_*F$ is constructible, it is enough, by Noetherian
induction and \hyperref[XIV.lem1.6]{1.6}, to prove that if
$S\neq\varnothing$, there exists a nonempty open subset of $S$ on which
$R^1f_*F$ is constructible. We may therefore suppose that $R^1f_*G$ is
locally constant. Since constructibility is local for the étale
topology, we may further suppose that $S$ is irreducible and that
$R^1f_*G$ is constant (and constructible), with value
$E=\{e_1,\ldots,e_n\}$. For $i=1,\ldots,n$, let $D_i$ be the
subsheaf\pageoriginale of $R^1f_*F$ which is the inverse image of the section $e_i$
under $u^1$. Since \lq\lq coproducts are universal\rq\rq\ in a topos,
we have $R^1f_*F=\amalg_iD_i$; it is therefore enough to prove that each
$D_i$ is constructible over a suitable nonempty scheme $S'$ étale over
$S$.

Recall that $R^1f_*F$ is the sheaf associated with the presheaf
$\Rr^1F$, where
$\Rr^1F(S')=H^1(X\times_SS',F)$. Consequently, we may suppose (after
replacing $S$ by a nonempty scheme $S'$ étale over

% END INCLUDED SGA4_English_current_buildable_source_20260906/source/14/drafts/SGA4_XIV_LINES_43_300_English.texfrag


% BEGIN INCLUDED SGA4_English_current_buildable_source_20260906/source/14/drafts/SGA4_XIV_LINES_301_601_English.texfrag
$S$) that $D_i$ is the “empty sheaf” (hence constructible),
\emph{or} that there exists an element $\alpha\in H^{1}(X,F)$ inducing
a section $\sisi{\alpha}{\bar{\alpha}}$ of $D_i$. In the latter case, let
$$
f_{\ast}C^{\alpha}\xrightarrow{\delta}R^{1}f_{\ast}F^{\alpha}\longrightarrow
R^{1}f_{\ast}G^{\alpha}
$$
be the exact sequence deduced from $(\ast)$ by twisting with
$\alpha$. We have
$$
R^{1}f_{\ast}F^{\alpha}\simeq R^{1}f_{\ast}F\quav
R^{1}f_{\ast}G^{\alpha}\simeq R^{1}f_{\ast}G{\quoi},
$$
and the section of $R^{1}f_{\ast}F^{\alpha}$ corresponding to
$\bar{\alpha}$ is the unit section. Thus the subsheaf
$D_i^{\alpha}$ of $R^{1}f_{\ast}F^{\alpha}$ corresponding to $D_i$
is the inverse image of the unit section of $R^{1}f_{\ast}G^{\alpha}$;
that is, $D_i^{\alpha}$ is the image of $f_{\ast}C^{\alpha}$ under
$\delta$. Since $f_{\ast}C^{\alpha}$ is constructible
(\hyperref[XIV.lem1.5]{1.5}), so is $D_i^{\alpha}$, and hence so is
$D_i$ (apply \hyperrefext[IX][IX.prop2.6]{2.6}). This proves the lemma.

\begin{enonce*}{Lemma 1.9}\etiquette[1.9]{XIV.lem1.9}
Let $S$ be noetherian, and let
$$
\xymatrix@=1.25cm{
X       \ar[rd]_{f}      &      &\bar{X} \ar[ll]_{\pi} \ar[ld]^{\bar{f}}\\
              & S    &
}
$$
be a commutative diagram of proper morphisms. If Theorem
\hyperref[XIV.thm1.1]{1.1} holds for\pageoriginale $f$ and $\pi$, then it also holds
for $\bar{f}$.
\end{enonce*}

\noindent
\emph{Proof}. In the case of a constructible sheaf $\bar{F}$ of
$A$-modules, we use the Leray spectral sequence
$$
E^{p,q}_2=(R^{p}f_{\ast})(R^{q}\pi_{\ast})\bar{F}\Longrightarrow
(R^{n}\bar{f}_{\ast})\bar{F}{\quoi}.
$$
By hypothesis, $E^{p,q}_2$ is constructible for every $p,q$; hence
the abutment is constructible as well, as follows from
\hyperrefext[IX][IX.prop2.6]{2.6}.

Let $\bar{F}$ be a constructible sheaf of groups. We have the exact
sequence (\hyperrefext[XII][XII.prop3.2]{3.2})
\begin{equation*}
0\longrightarrow (R^{1}f_{\ast})\pi_{\ast}\bar{F}\longrightarrow
(R^{1}\bar{f}_{\ast})\bar{F}\longrightarrow
f_{\ast}(R^{1}\pi_{\ast})\bar{F}{\quoi},\tag{$\ast$}
\end{equation*}
whose outer terms are constructible by the hypothesis and
\hyperref[XIV.lem1.5]{1.5}. We proceed as in the proof of
\hyperref[XIV.lem1.8]{1.8}: by noetherian induction and
\hyperref[XIV.lem1.6]{1.6}, it suffices to assume $S\neq ∅$ and prove
constructibility over a suitable nonempty open subset. We may thus
assume that $f_{\ast}(R^{1}\pi_{\ast})\bar{F}$ is locally constant and,
after an étale localization, even constant and constructible
(\sisi{\hyperrefext[IX][IX.prop2.8]{2.8}}{\hyperrefext[IX][IX.prop2.4]{2.4}~(i)}),
with value $E=\{e_1,\ldots,e_n\}$. Let $D_i$ be the subsheaf of
$(R^{1}\bar{f}_{\ast})\bar{F}$ that is the inverse image of $e_i$, so
that $(R^{1}\bar{f}_{\ast})\bar{F}=\coprod_iD_i$. It is enough to show
that each sheaf $D_i$ induces a constructible sheaf over some suitable
nonempty connected étale $S'/S$. If $D_i$ is not the “empty sheaf”, we
reduce to the case in which there exists an element
$\alpha\in H^{1}(\bar{X},\bar{F})$ inducing a section $\bar{\alpha}$ of
$D_i$. Let
$$
\sisi{
0\longrightarrow (R^{1}f_{\ast})\pi_{\ast}\bar{F}\longrightarrow
(R^{1}\bar{f}_{\ast})\bar{F}\longrightarrow
f_{\ast}(R^{1}\pi_{\ast})\bar{F}
}
{0\longrightarrow (R^{1}f_{\ast})\pi_{\ast}\bar{F}^{\alpha}\longrightarrow
(R^{1}\bar{f}_{\ast})\bar{F}^{\alpha}\longrightarrow
f_{\ast}(R^{1}\pi_{\ast})\bar{F}^{\alpha}}
$$
be the exact sequence deduced from $(\ast)$ by twisting $\bar{F}$ with
$\alpha$. The section of $(R^{1}\bar{f}_{\ast})\bar{F}^{\alpha}$
corresponding to $\bar{\alpha}$ is the unit section. Consequently the
subsheaf $D_i^{\alpha}$ of \hbox{$(R^{1}\bar{f}_{\ast})\bar{F}^{\alpha}$}
corresponding to $D_i$ is the inverse image of the unit section of
$f_{\ast}(R^{1}\pi_{\ast})\bar{F}^{\alpha}$,
\pageoriginale whence
$D_i^{\alpha}\simeq (R^{1}f_{\ast})\pi_{\ast}\bar{F}^{\alpha}$,
which is indeed a constructible sheaf, \cqfd

\begin{enonce*}{Lemma 1.10}\etiquette[1.10]{XIV.lem1.10}
It suffices to verify \hyperref[XIV.thm1.1]{1.1} when $S$ is of finite
type over $\Spec\ZZ$, and $f$ is projective and of relative dimension
$\leq 1$.
\end{enonce*}

\noindent
\emph{Proof}. Suppose first that $f$ is projective, and argue by
induction on the relative dimension of $f$, which may be assumed finite
after replacing $S$ by an affine open subset. If $f$ has relative
dimension $\leq n$, with $n\geq 2$, then locally on $S$ one can find a
finite morphism $\pi:X\rightarrow\PP^n_S$. By
\hyperref[XIV.lem1.9]{1.9} and
\sisi{IX 2.15}{\hyperrefext[IX][IX.prop2.14]{2.14}} (i), we reduce to
the case in which $f$ is the structural morphism of $\PP^n_S$.
Consider the commutative diagram of morphisms
$$
\xymatrix@C=2cm{
                   & \bar{X}\ar[ld]_{\sisi{}{\pi}} \ar[rd]^{a} \ar@{-}[dd]^{\bar{f}} &       \\
\PP^{n}_S \ar[rd]_{f}  &                  &\PP^{1}_S \ar[ld]^{b}\\
                   &  S                          &
}
$$
already considered in the proof of
\hyperrefext[XII][XII.lem8.3]{8.3}, where $\bar{X}$ is obtained from
$\PP^n_S$ by blowing up the closed subscheme
$Y\simeq\PP^{n-2}_S$. The relative dimension of $b$ is $\leq1$, and
that of $a$ is $\leq n-1$; therefore the theorem holds for $a$ and $b$
by the induction hypothesis, and hence for $\bar{f}$ by
\hyperref[XIV.lem1.9]{1.9}. Moreover, the relative dimension of $\pi$
is $\leq1$, while that of $f_0:Y\rightarrow S$ is $\leq n-1$. It
therefore follows from \hyperref[XIV.lem1.7]{1.7} that
\hyperref[XIV.thm1.1]{1.1} holds for $f$, and thus for every projective
morphism. Now let $f$ be an arbitrary proper morphism, and argue by
noetherian induction on $X$: we may assume the theorem holds for every
proper closed subscheme of $X$ and that $X\neq∅$. Let
$$
\xymatrix@=1.5cm{
X  \ar[rd]_{f}          &           &\bar{X} \ar[ll]_{\pi} \ar[ld]^{\bar{f}}\\
                        &    S       &
}
$$
be a\pageoriginale diagram of the type supplied by Chow's lemma
\hyperrefext[XII][XII.lem7.1]{7.1}. Applying
\hyperref[XIV.lem1.7]{1.7} to this diagram yields
\hyperref[XIV.thm1.1]{1.1} for $f$, and proves the lemma.

\setcounter{subsection}{10}
\subsection{}\etiquette[1.11]{XIV.subsec1.11}
We can now complete the proof of \hyperref[XIV.thm1.1]{1.1}. If $F$ is
a constructible sheaf of groups, one can find an injection
$F\hookrightarrow G=\prod_i\pi_{i\sisi{}{,}\ast}C_i$, where each
$\pi_i$ is finite and each $C_i$ is constant and constructible
(\sisi{IX 2.15}{\hyperrefext[IX][IX.prop2.14]{2.14}}); by
\hyperref[XIV.lem1.8]{1.8}, we may therefore assume that $F$ is
constant. Likewise, if $F$ is a constructible sheaf of $A$-modules,
then by applying \sisi{IX 2.15}{\hyperrefext[IX][IX.prop2.14]{2.14}}
one can find a resolution $F\rightarrow G^{\bigdot}$,
$G^{\bigdot}=\{0\rightarrow G^0\rightarrow G^1\rightarrow
G^2\rightarrow\ldots\}$, of $F$ such that every $G^{\nu}$ is of the
form
$G^{\nu}=\prod_i\sisi{\eta}{\pi}_{i\sisi{}{,}\ast}C_i$.
The spectral sequence of the resolution $G^{\bigdot}$,
$$
E^{p,q}_2=H^{\sisi{q}{p}}((R^{q}f_{\ast})G^{\bigdot})\Longrightarrow
(R^{n}f_{\ast})F{\quoi},
$$
then reduces us to the case in which $F$ is constant and constructible,
with value $M$.

Recall that $A$ is a $\ZZ/n$-algebra. Let $\ell$ be a prime number
dividing $n$. We have the exact sequence
$$
0\longrightarrow \ell M\longrightarrow M\longrightarrow M/\ell
M\longrightarrow 0{\quoi},
$$
and it is again enough to prove \hyperref[XIV.thm1.1]{1.1} for
$\ell M$ and $M/\ell M$.

An easy induction then reduces us to the case in which $A$ is a
$\ZZ/\ell$-algebra, $\ell\in\PP$, and $M$ is a finite-type
$A$-module. Then $M$ is a vector space over $\ZZ/\ell$, and the
canonical morphism
$$
H^q(Y,\ZZ/\ell)\otimes_{\ZZ/\ell}M\longrightarrow H^q(Y,M)
$$
is bijective for every quasi-compact and quasi-separated $Y$ (indeed,
this is trivial if $M$ has finite rank, and both sides commute with
filtered direct limits\pageoriginale
(\sisi{VI 1.6}{\hyperrefext[VI][VI.thm5.1]{5.1} and
\hyperrefext[VII][VII.prop3.3]{3.3}})). Hence
$$
[R^qf_{\ast}(\ZZ/\ell)]\otimes_{\ZZ/\ell}M\xrightarrow{\sim}
R^qf_{\ast}(M){\quoi}.
$$
It plainly suffices, therefore, to show that $R^qf_{\ast}(\ZZ/\ell)$
is constructible as a $\ZZ/\ell$-module; in other words, we are reduced
to the case $A=\ZZ/\ell$, which we assume from now on.

Observe that the case in which $f$ has relative dimension $\leq0$
(that is, $f$ is finite), with $F$ arbitrary, now follows from
\hyperrefext[VIII][VIII.cor5.6]{5.6} and
\hyperref[XIV.lem1.5]{1.5}.

We proceed by noetherian induction on $S$. By
\hyperref[XIV.lem1.6]{1.6}, we may assume $S\neq∅$ and may replace $S$
by any nonempty open subset. Moreover, we may assume that $S$ is
integral, with generic point $s$. The fiber $X_s$ is an algebraic
scheme of dimension $\leq1$; consequently there is a purely inseparable extension
$K'$ of $k(s)=K$ such that the normalization $Y$ of
$(X_s\otimes_KK')_{\text{\rm red}}$ is smooth over $\Spec K'$
\footnote[1]{\Cf \EGA IV 17.15.14.}. After replacing $S$ by a nonempty
open subset, we may assume that there exists a finite, radicial,
surjective morphism $S'\rightarrow S$ such that $S'_s$ is
$K$-isomorphic to $\Spec K'$, a smooth projective morphism
$\bar{X}'\rightarrow S'$, and a finite surjective $S'$-morphism
$\pi:\bar{X}'\rightarrow X'=\fprod{X}{S'}{S}$
inducing $Y\rightarrow X_s\otimes_KK'$ on the fibers over the generic
point $s$ of $S$ (\EGA IV 9.6.1). Now\sisi{}{,} by
\hyperrefext[VIII][VIII.thm1.1]{1.1}\sisi{}{,} making a purely inseparable
extension is harmless. We may therefore replace $S$ by $S'$ and $X$ by
$X'$; that is, we may assume that $f$ fits into a diagram
$$
\xymatrix@=1.5cm{
X   \ar[rd]_{f}& &\bar{X} \ar[ll]_{\pi} \ar[ld]^{\bar{f}}\\
         &         S        &           ,
}
$$
in which\pageoriginale $\pi$ is finite and surjective and $\bar f$ is
smooth. Moreover, by construction, we may assume that there is a closed
subscheme $Y$ of $X$ whose generic fiber is finite and such that $\pi$
induces an isomorphism over $X-Y$. Replacing $S$ once more by a
nonempty open subset, we reduce to the case in which
$f_0:Y\rightarrow S$ is also finite.

Since $F$ is constant, $\pi^{\ast}F$ is constant as well, and
$F\rightarrow\pi_{\ast}\pi^{\ast}F$ is injective because $\pi$ is
surjective. If $F$ is a sheaf of groups, it suffices to prove that
$(R^1f_{\ast})(\pi_{\ast}\pi^{\ast}F)\simeq
(R^1\bar f_{\ast})(\pi^{\ast}F)$ is constructible
(\hyperref[XIV.lem1.8]{1.8}). If $F$ is a sheaf of $\ZZ/\ell$-modules,
we have an exact sequence
$$
0\longrightarrow F\longrightarrow\pi_{\ast}\pi^{\ast}F\longrightarrow
C\longrightarrow0{\quoi},
$$
where $C$ is supported on $Y$, because $\pi$ induces an isomorphism over
$X-Y$. It follows that $R^qf_{\ast}C=R^qf_{0\sisi{}{,}\ast}C$ is
constructible (and vanishes if $q>0$), so it is once again enough to
show that
$(R^qf_{\ast})(\pi_{\ast}\pi^{\ast}F)=
(R^q\bar f_{\ast})(\pi^{\ast}F)$ is constructible. Replacing $f$ by
$\bar f$ and $F$ by $\pi^{\ast}F$, we are reduced to the case in which
$F$ is constant and constructible and $f$ is smooth.

Let us first treat a constant sheaf of groups, $\sisi{f}{F}=G_X$. All the fibers
of $(R^1f_{\ast})G_X$ are finite. Indeed, by applying
\hyperrefext[XII][XII.cor5.2]{5.2} (ii), we reduce, in order to prove
this, to the case in which $S$ is the spectrum of a separably closed
field and $X$ is smooth of dimension $\leq1$ over $S$. The result is
known in this case (\SGA 1 X 2.6). It therefore remains only to prove
that the specialization morphisms are injective
(\hyperrefext[IX][IX.prop2.13]{2.13} (ii)). For this we reduce to the
case in which $S$ is strictly local; by the base-change theorem
\hyperrefext[XII][XII.thm5.1]{5.1} (ii), we may even replace $S$ by the
spectrum of a strictly local discrete valuation ring. Moreover, we may
assume $X$ connected. Since\pageoriginale $H^1(X,G)$ classifies the
principal Galois coverings of $X$ with group $G$, it is enough to check
that if $X'\rightarrow X$ is a connected étale covering, then its
geometric fiber is connected as well. This follows readily from the
fact that $f$ is smooth and proper with geometrically connected fibers.

We now treat the case $F=(\ZZ/\ell)_X$. For $q=0$ and $q=1$, the
result is already known (\hyperref[XIV.lem1.5]{1.5} and the preceding
argument). \sisi{By \hyperrefext[X][X.cor4.3]{4.3} and}{By}
\hyperrefext[XII][XII.corbis5.3]{5.3 bis}, only the case $q=2$ remains;
furthermore, if $S$ has characteristic $\ell$, then
$R^2f_{\ast}(\ZZ/\ell)=0$. Replacing $S$ by a suitable nonempty open
subset, we therefore reduce to the case in which $\ell$ is invertible
on $S$. Then, by \hyperrefext[XII][XII.cor5.2]{5.2}
\sisi{(ii)}{(iii)} and
\sisi{X 4.7}{\hyperrefext[X][X.cor5.2]{5.2}}, the fiber of
$R^2f_{\ast}(\ZZ/\ell)$ at a geometric point $\bar s$ above $s\in S$
is isomorphic to $(\ZZ/\ell)^{c(s)}$, where $c(s)$ is the number of
one-dimensional irreducible components of the geometric fiber
$X_{\bar s}$. The function $c$ is constructible on $S$
(\EGA IV 9.7.9), and the result follows
(\hyperrefext[IX][IX.prop2.13]{2.13} (iii)).

\begin{enonce*}[remark]{Remark 1.11}\etiquette[1.11]{XIV.rem1.11}
The proof of \hyperref[XIV.thm1.1]{1.1} becomes much simpler if one
uses the formalism of cohomology with proper supports, which will be
developed in \hyperref[sga4.exp.XVII]{XVII} as a direct consequence of the base-change
theorem for a proper morphism \hyperrefext[XII][XII.thm5.1]{5.1}.
Moreover, the statement of \hyperref[XIV.thm1.1]{1.1} generalizes to
the case in which one is given an arbitrary sheaf $𝒜$ of torsion rings
on $\sisi{Y}{S}$ and a complex $K^{\bigdot}$ of $f^{\ast}(𝒜)$-modules
on $X$ satisfying a “constructibility” condition which, in the special
case of \hyperref[XIV.thm1.1]{1.1}, would simply say that the cohomology
sheaves $\Hh^i(K^{\bigdot})$ of $K^{\bigdot}$ are constructible
$A$-modules and vanish for sufficiently large $i$.
\end{enonce*}

% END INCLUDED SGA4_English_current_buildable_source_20260906/source/14/drafts/SGA4_XIV_LINES_301_601_English.texfrag


% BEGIN INCLUDED SGA4_English_current_buildable_source_20260906/source/14/drafts/SGA4_XIV_LINES_602_722_English.texfrag
\section{A variant of dimension}\etiquette[2]{XIV.sec2}\pageoriginale

\subsection{}\etiquette[2.1]{XIV.subsec2.1}
Let $X$ be a scheme. We shall write
\begin{equation*}
d(x)=\dim\bar{\{x\}}\tag{2.1.1}
\end{equation*}\etiquette[2.1.1]{XIV.eq2.1.1}
for $x\in X$, where $\bar{\{x\}}$ is the closure of $\{x\}$. If $F$
is an abelian sheaf on $X$, we put
\begin{equation*}
d(F)=\sup \{d(x)\mid x\in X \quaet F_{\bar{x}}\neq 0\}{\quoi}.\tag{2.1.2}
\end{equation*}\etiquette[2.1.2]{XIV.eq2.1.2}
We shall need a corrected notion in the local case:

\begin{enonce*}{Definition 2.2}\etiquette[2.2]{XIV.def2.2}
Let $Y$ be the spectrum of a universally catenary Noetherian local
ring (\EGA IV 5.6.2), and let $f:X\rightarrow Y$ be a morphism of
finite type. Let $x\in X$ and $y=f(x)$. We put
$$
\delta(x)=\dim\bar{\{y\}}+\degtr k(x)/k(y){\quoi}.
$$
If $F$ is an abelian sheaf on $X$, we put
$$
\delta(F)=\sup \{\delta(x)\mid x\in X\quaet F_{\bar{x}}\neq 0\}{\quoi}.
$$
\end{enonce*}

\begin{enonce*}{Proposition 2.3}\etiquette[2.3]{XIV.prop2.3}
\begin{enumerate}
%%%
%%% The typesetting of this enumerate environment is terrible.
%%% The spacing between the dash and (i) is too big.
%%% Also, the items (ii), ... should better be without indent.
%%% If you use the package enum.sty, then you can build the necessary list environment.
%%%
\renewcommand{\theenumi}{\roman{enumi}}
\renewcommand{\labelenumi}{(\theenumi)}
\item Let $i:X'\rightarrow X$ be an immersion and let $x'\in X'$.
Then $\delta(i(x'))=\delta(x')$.
%%% new line break within formula in text: spend extra line to induce new line breaks?

\item Let $X_{0}$ be the closed fiber of $X/Y$. We have
$$
\delta(x)=
\begin{cases}
d(x) & \mbox{if } \bar{\{x\}}\cap X_{0}\neq ∅\\
d(x)+1 & \mbox{if } \bar{\{x\}}\cap X_{0}=∅{\quoi};
\end{cases}
$$

\item $\delta(x)=d(x)$ if $f$ is proper.
\end{enumerate}
\end{enonce*}

\noindent
\emph{Proof}.\pageoriginale Assertion (i) follows immediately from the
definition. Since $\bar{\{x\}}\cap X_{0}$ is nonempty when $f$ is
proper, (iii) follows from (ii). For (ii), we may replace $X$ by
$\bar{\{x\}}$ and $Y$ by $\bar{\{y\}}$ (with their induced reduced
structures), as is permitted by (i). Let $x'\in X$ and $y'=f(x')$.
We have the formula (\EGA IV 5.6.1.1)
$$
\dim \Oo_{X,x'}=\dim
\Oo_{Y,y'}+\degtr k(x)/k(y) - \degtr k(x')/k(y'){\quoi}.
$$

Suppose that $X_{0}\neq ∅$. It is then clear that
$\displaystyle{\mathop{\sup}_{x'}}$
$\{\dim \Oo_{X,x'}\}=\dim X$ is attained by taking $x'$ to be a closed
point of $X_{0}$, whence
$$
\dim X=\dim \Oo_{X,x'}=\dim Y + \degtr k(x)/k(y)=\delta(x){\quoi}.
$$
If $X_{0}=∅$, the supremum is plainly attained by taking $x'$ to
be a closed point in its fiber $f^{-1}(y')$, with
$\dim\bar{\{y'\}}=1$. Such points exist because the $y'\in Y$ with
$\dim\bar{\{y'\}}=1$ form a very dense set (\EGA IV 10.1.3,
\sisi{10.5.9, 1.8.4}{10.3.1, 10.4.1, 10.5.9}). Thus
\begin{align*}
\dim X=\dim \Oo_{X,x'} &= \dim \Oo_{Y,y'}+\degtr k(x)/k(y)\\
&= \dim Y-1 +\degtr k(x)/k(y)=\delta(x) \sisi{=}{-} 1{\quoi},
\end{align*}
which proves the result.

\begin{enonce*}{Proposition 2.4}\etiquette[2.4]{XIV.prop2.4}
Let $f:X\rightarrow Y$ be a morphism of schemes of finite type over
the spectrum of a field $k$. Let $x\in X$, let $y=f(x)$, and let
$y_{0}$ be a specialization of $y$. Let $f':X'\rightarrow Y'$ be the
strict localization of $f$ at a geometric point $\bar{y}_{0}$ above
$y_{0}$, and let $x'\in X'$ be a point above $x$. Then
$$
\delta(x')=d(x)-d(y_{0}){\quoi}.
$$
\end{enonce*}

\noindent
\emph{Proof}.\pageoriginale We have
\begin{align*}
d(x) &= d(y)+\degtr k(x)/k(y)\\
\delta(x') &= d(y')+\degtr k(x')/k(y'){\quoi},
\end{align*}
where $y'=f'(x')$. Moreover,
$$
\degtr k(x)/k(y)=\degtr k(x')/k(y'){\quoi},
$$
because the extensions $k(x')/k(x)$ and $k(y')/k(y)$ are algebraic.
It therefore remains to prove that
$$
d(y')=d(y)-d(y_{0}){\quoi}.
$$

We may assume that $y$ is maximal in $Y$, and hence that $y'$ is
maximal in $Y'$. Then
\begin{align*}
d(y') &= \dim \Oo_{Y',\bar{y}_{0}}\\
      &= \dim \Oo_{Y,y_{0}}\quad \text{\rm (\EGA IV 6.1.3)}\\
      &= \sisi{\dim Y}{\dim_{y_0} Y} -d(y_{0})\quad \text{\rm (\EGA IV 5.2.3.1)}\\
      &= d(y)-d(y_{0}){\quoi},
\end{align*}
which proves the result.

% END INCLUDED SGA4_English_current_buildable_source_20260906/source/14/drafts/SGA4_XIV_LINES_602_722_English.texfrag


% BEGIN INCLUDED SGA4_English_current_buildable_source_20260906/source/14/drafts/SGA4_XIV_LINES_723_863_English.texfrag
\section{Cohomological dimension of affine algebraic schemes}\etiquette[3]{XIV.sec3}

\begin{enonce*}{Theorem 3.1~(d)}\etiquette[3.1]{XIV.thm3.1}\nde{Theorem \ref{XIV.thm3.1} is now known as the \textit{affine Lefschetz theorem}.}
Let $f:X\rightarrow Y$ be an affine morphism of schemes of finite type
over a field $k$. Let $F$ be a torsion sheaf on $X$ such that
$d(F)\leq d$ \emph{(\cf \hyperref[XIV.subsec2.1]{2.1})}.
Then $d(R^{q}f_{\ast}F)\leq d-q$.
\end{enonce*}

In particular, taking $Y=\Spec(k)$, we obtain:

\begin{enonce*}{Corollary 3.2}\etiquette[3.2]{XIV.cor3.2}
Let\pageoriginale $X$ be an affine scheme of finite type over a separably
closed field $k$.
Then\nde{Consider the analogous assertion in complex analytic geometry:
let $X$ be an affine variety over $\mathbf{C}$ and let $F$ be a constructible sheaf on $X$. Then $H^m(X^{\rm an},F)=0$ for every $m>\dim X$, where $X^{\rm an}$ is the analytic space associated with $X$. Although this assertion follows from \hyperref[XIV.cor3.2]{3.2} by the comparison theorem \hyperrefext[XVI][XVI.thm4.1]{4.1}, there are also two entirely analytic proofs, due to Pierre Deligne and Hélène Esnault; see Esnault, H., Variation on Artin's vanishing theorem, \textit{Advances in Math.} \textbf{198} (2005), 435--438.}
$$
\cd X\leq \dim X{\quoi}.
$$
\end{enonce*}

\begin{enonce*}{Corollary 3.3}\etiquette[3.3]{XIV.cor3.3}
Let $X\subset \PP^N_k$ be a projective scheme of dimension $n$ over a
separably closed field $k$, and let $Y=H\cap X$ be a ``hyperplane
section'' of $X$. Then, for every torsion sheaf $F$ on $X$, the canonical
homomorphism
$$
H^{q}_Y(X,F)\longrightarrow H^{q}(X,F)
$$
is bijective for $q>n+1$ and surjective for $q=n+1$.
\end{enonce*}

The corollary follows immediately from
\sisi{\hyperref[XIV.def2.2]{2.2}}{\hyperref[XIV.cor3.2]{3.2}} and the exact
sequence for a closed subset
(\sisi{IV 3.7.1}{\hyperrefext[V][V.6.5.4]{6.5.4}}),
%%% This last reference should read \ref{V.eq6.5.4}, but the TeX-file of chapter V contains the
%%% label without "eq".
since $X-Y=\sisi{\mathcal{U}}{U}$ is affine of dimension $\leq n$.
Notice that \hyperref[XIV.cor3.3]{3.3} is a generalization of one of
Lefschetz's theorems on hyperplane sections. Indeed, if $X$ is smooth
over $k$ and $F$ is locally constant of order prime to the characteristic,
the cohomology with support $H^{q}_Y(X,F)$ can be computed explicitly
using the relative cohomological purity theorem
(\hyperrefext[XVI][XVI.sec3]{3}) and the spectral sequence
(\sisi{\quad}{\hyperrefext[V][V.6.4.3]{6.4.3}}),
%%% This last reference should read \ref{V.eq6.4.3}, but the TeX-file of chapter V contains the
%%% label without "eq".
and one finds that
$$
H^{q}_Y(X,\ZZ/n)\simeq H^{q-2}(Y,\mu_n^{-1}) \mbox{\rm\quad for
all $q$;}
$$
the homomorphism $H^{q-2}(Y,\mu_n^{-1})\rightarrow
H^{q}(X,\ZZ/n)$ furnished by \hyperref[XIV.cor3.3]{3.3} is precisely the
``Gysin homomorphism'', which will be studied later
\footnote[1]{\Cf \SGA 2 XIV for a systematic study of Lefschetz-type
theorems.}.

Theorem \hyperref[XIV.thm3.1]{3.1}~(d) is equivalent\nde{More precisely,
the argument in the text gives the following result. Assertion
\hyperref[XIV.cor3.4]{3.4}~(d') for $d'\leq d$ implies
\hyperref[XIV.thm3.1]{3.1}~(d), and assertion
\hyperref[XIV.thm3.1]{3.1}~(d) implies
\hyperref[XIV.cor3.4]{3.4}~(d). This suffices for the conclusion below.}
to the following statement:

\begin{enonce*}{Corollary 3.4~(d)}\etiquette[3.4]{XIV.cor3.4}
Let $Y'$ be a strictly local scheme that is a strict localization of an
algebraic scheme over a field. Let $f':X'\rightarrow Y'$ be an affine
morphism\pageoriginale of finite type, and let $F'$ be a torsion sheaf on
$X'$ such that $\delta(F')\leq d$ \emph{(\cf
(\hyperref[XIV.def2.2]{2.2}))}. Then $H^{q}(X',F')=0$ for $q>d$.
\end{enonce*}

Indeed, suppose that \hyperref[XIV.cor3.4]{3.4}~(d) is true. Let
$f:X\rightarrow Y$ be an affine morphism of algebraic schemes over $k$,
and let $F$ be a torsion sheaf on $X$ such that $d(F)\leq d$. Let $y$ be
a point of $Y$, and let $\xi$ be a geometric point over $y$. To prove
\sisi{\ref{XIV.subsec3.1}}{\hyperref[XIV.thm3.1]{3.1}}~(d), it is enough
to prove that the stalk $(R^{q}f_{\ast}F)_{\xi}$ vanishes when
$d(y)>d-q$.

Let $f':X'\rightarrow Y'$ be the strict localization of $f$ at the point
$\xi$; that is, let $Y'$ be the strict localization of $Y$ at $\xi$, and
let $X'=\fprod{X}{Y'}{Y}$.

Let $F'$ be the sheaf induced by $F$ on $X'$, so that
$(R^{q}f_{\ast}F)_{\xi}=H^{q}(X',F')$
(\hyperrefext[VIII][VIII.thm5.2]{5.2}).
\sisi{We have
$$
\delta(F')\leq d(F)-d(y)\leq d-d(y)\quad (\hyperref[XIV.prop2.4]{2.4}){\quoi},
$$}
{By (\hyperref[XIV.prop2.4]{2.4}), we have
$$
\delta(F')\leq d(F)-d(y)\leq d-d(y){\quoi},
$$}
%%% I find it irritating that the reference to the proposition from which the above follows,
%%% is referred to in a style that might be confused with an equation number
%%% (although the correct equation number is 3.4.1)
and applying \hyperref[XIV.cor3.4]{3.4}~(d), we obtain
$H^{q}(X',F')=0$ if $q>d-d(y)$, which proves the result.

Conversely, suppose that
\sisi{\hyperref[XIV.subsec2.1]{2.1}}{\hyperref[XIV.thm3.1]{3.1}}~(d) is
true, and let $F'$ and $f':X'\rightarrow Y'$ be as in the statement of
\hyperref[XIV.cor3.4]{3.4}~(d). Since cohomology commutes with filtered
inductive limits (\hyperrefext[VII][VII.prop3.3]{3.3}) and $F'$ is the
inductive limit of its constructible subsheaves
(\hyperrefext[IX][IX.prop2.9]{2.9}), we may suppose that $F'$ is
constructible. Now $Y'$ is the strict localization of an algebraic
scheme $Y$ at a geometric point $\xi$ over a \emph{closed} point $y_0$
(\hyperrefext[X][X.lem3.3]{3.3}). Write
$Y'=\varprojlim_{\alpha}Y_{\alpha}$, where $(Y_{\alpha})$ is a filtered
inverse system of affine schemes étale over $Y$
(\hyperrefext[VIII][VIII.subsec4.5]{4.5}). By
\hyperrefext[IX][IX.cor2.7.4]{2.7.4}, $F'$ is induced by a constructible
torsion sheaf \sisi{}{$F_\alpha$} on
$\sisi{X=\fprod{X}{Y_K}{Y}}{X_\alpha=\fprod{X}{Y_\alpha}{Y}}$ for
sufficiently large $\alpha$. Using \hyperref[XIV.prop2.4]{2.4} and\pageoriginale \EGA IV 9.5.5,
we see immediately that we may further suppose that
$d(F_{\alpha})=\delta(F')\leq d$. Replacing $Y$ by $Y_{\alpha}$, we are
reduced to the case in which $F'$ is induced by a sheaf $F$ on $X$, with
$d(F)\leq d$. Then $R^{q}f_{\ast}F=0$ for $q>\sisi{D}{d}$ by
\sisi{\hyperref[XIV.subsec2.1]{2.1}}{\hyperref[XIV.thm3.1]{3.1}}~(d),
and hence in particular
$$
(R^{q}f_{\ast}F)_{\xi}=H^{q}(X',F')=0\quad{\mbox{\rm if}}\quad
q>d{\quoi},
$$
which proves \hyperref[XIV.cor3.4]{3.4}~(d).

As a special case of \hyperref[XIV.cor3.4]{3.4}~(d), we have:

\begin{enonce*}{Corollary 3.5}\etiquette[3.5]{XIV.cor3.5}
Let $Y'$ be a strictly local scheme of dimension $\leq d$, which is a
strict localization of an algebraic scheme over a field $k$. Let
$U\subset Y'$ be an affine open subset. Then $\cd U\leq d$.
\end{enonce*}

\begin{enonce*}[remark]{Remark 3.6}\etiquette[3.6]{XIV.rem3.6}
It seems very plausible\nde{Ofer Gabber has succeeded in proving the following analogue of \ref{XIV.thm3.1}: let $f:X\to Y$ be an affine morphism of finite type, with $Y$ quasi-excellent (\EGA IV 7.8) and endowed with a dimension function $\delta_Y$. Let $\delta_X$ be the associated rectified dimension function on $X$. If $n$ is invertible on $Y$, then for every sheaf $F$ of $\mathbf{Z}/n$-modules on $X$ one has $\delta_Y(R^q f_\ast F)\leq\delta_X(F)-q$ for every $q\geq0$. Here, as above, for a sheaf $G$ we define $\delta(G)=\sup\{\delta(x)\mid x\in X\quaet G_{\bar{x}}\neq0\}$.
}
that \hyperref[XIV.cor3.4]{3.4}, and hence
\hyperref[XIV.cor3.5]{3.5}, remains valid for \emph{every} local
Noetherian scheme $Y'$ (not necessarily a strict localization of an
algebraic scheme), at least when $Y'$ is \emph{excellent}
(\EGA IV 7.8). This will be proved in
(\sisi{XX 6}{\hyperrefext[XIX][XIX.sec6]{6}}) under the further
assumption that $Y$ has characteristic zero.
\end{enonce*}

% END INCLUDED SGA4_English_current_buildable_source_20260906/source/14/drafts/SGA4_XIV_LINES_723_863_English.texfrag


% BEGIN INCLUDED SGA4_English_current_buildable_source_20260906/source/14/drafts/SGA4_XIV_LINES_864_1076_English.texfrag
\section[Proof of Theorem 3.1]{Proof of Theorem \texorpdfstring{\hyperref[XIV.thm3.1]{3.1}}{3.1}}\etiquette[4]{XIV.sec4}

The proof proceeds by induction on $d$.

\begin{enonce*}{Lemma 4.1}\etiquette[4.1]{XIV.lem4.1}
\emph{\hyperref[XIV.thm3.1]{3.1}~(0)} and
\emph{\hyperref[XIV.thm3.1]{3.1}~(1)} hold.
\end{enonce*}

\noindent
\emph{Proof}. This is trivial for $d=0$, in which case $\supp F$ is
finite over $k$, and we apply
(\hyperrefext[VIII][VIII.prop5.5]{5.5}). Let $F$ and
$f:X\rightarrow Y$ be as in the statement of
\hyperref[XIV.thm3.1]{3.1}, and suppose that $d(F)\leq1$. Replacing
$F$ by a constructible subsheaf
(\hyperrefext[IX][IX.prop2.9]{2.9} and
\hyperrefext[VII][VII.prop3.3]{3.3}), we reduce to the case in which
$F$ is constructible. The support\pageoriginale of $F$ is then
contained in a closed subscheme of $X$ of dimension $\leq1$
(\hyperrefext[IX][IX.def2.3]{2.3}\sisi{ (ii)}{}), and we may replace
$X$ by this subscheme; that is, we may assume that $\dim X\leq1$.
By \hyperrefext[IX][IX.cor5.7]{5.7}, the cohomological dimension of an
affine algebraic scheme of dimension $1$ over a separably closed field
is $\leq1$. It is enough, moreover, to treat the case in which $k$ is
separably closed. Indeed, let $\bar f:\bar X\rightarrow\bar Y$ be the
morphism obtained from $f$ by the base change
$\Spec\bar k\rightarrow\Spec k$, where $\bar k$ is the separable
closure of $k$. For a geometric point $\xi$ of $Y$, the strict
localization $f':X'\rightarrow Y'$ of $f$ at $\xi$ identifies with a
strict localization of $\bar f$, and consequently the fiber
$(R^{q}f_{\ast}F)_\xi$ identifies with a fiber of
$R^{q}\bar f_{\ast}\bar F$. Since $f$ is affine, it follows that
$R^{q}f_{\ast}F=0$ for $q>1$. It is also clear that
$d(f_{\ast}F)\leq1$. For
$q=1$, note that the hypothesis $\dim X\leq1$ implies that there is a
finite closed subset $Z$ of $Y$ such that $f$ is finite over $Y-Z$.
It follows that $R^{1}f_{\ast}F$ vanishes on $Y-Z$
(\hyperrefext[VIII][VIII.cor5.6]{5.6}), and hence
$d(R^{1}f_{\ast}F)\leq1$. This proves the lemma.

Let $d\geq2$, and now suppose that
\hyperref[XIV.thm3.1]{3.1}~(d') (or, equivalently,
\hyperref[XIV.cor3.4]{3.4}~(d')) holds for $d'<d$.

\begin{enonce*}{Lemma 4.2}\etiquette[4.2]{XIV.lem4.2}
To prove \emph{\hyperref[XIV.cor3.4]{3.4}~(d)} in general, it is
enough to treat the case in which $X'=\EE^1_{Y'}$ is affine
one-space over $Y'$ and $f$ is the structural morphism.
\end{enonce*}

\noindent
\emph{Proof}. We may plainly assume that $X'=\EE^N_{Y'}$ for a
suitable $N$, since $X'$ can be embedded in some $\EE^N_{Y'}$.
Suppose that $N>1$ and that the result is known, for every $Y'$, for
$\EE^r_{Y'}$ with $r<N$. Let $F'$ be a torsion sheaf on
$\EE^N_{Y'}$ with\pageoriginale $d(F')\leq d$. Consider the diagram
$$
\xymatrix@=1.5cm{
\EE^{N}_{Y'} \ar[rd]_{f'}       \ar[rr]^{g'}&      &
\EE^{N-1}_{Y'} \ar[ld]^{h'}\\
                       & Y' & &,
}
$$
where $g'$ and $h'$ are the canonical projections, so that
$\EE^N_{Y'}$ is isomorphic to affine one-space over
$\EE^{N-1}_{Y'}$. Since every strict localization of
$\EE^{N-1}_{Y'}$ is plainly also the strict localization of an
algebraic scheme, we may apply the induction hypothesis to the fibers
of $R^{q}g'_{\ast}F'$, and obtain
\begin{equation*}
\delta(R^{q}g'_{\ast}F')\leq d-q{\quoi}.\tag{$\ast$}
\end{equation*}
Moreover, \hyperref[XIV.cor3.4]{3.4}~(d') holds for $h'$ and
$d'\leq d$, by the induction hypothesis on $N$.

Consider the Leray spectral sequence
$$
E^{p,q}_2=H^{p}(\EE^{N-1}_{Y'},R^{q}g'_{\ast}(F'))
\Longrightarrow H^{n}(\EE^N_{Y'},F'){\quoi}.
$$
We have $E^{p,q}_2=0$ if $p>d-q$, as follows by applying $(\ast)$ and
\hyperref[XIV.cor3.4]{3.4}~(d') to the morphism $h'$. Consequently the
abutment vanishes for $n>d$, proving the lemma.

Consider the following additional statement, which is a special case
of \hyperref[XIV.cor3.5]{3.5}:

\begin{enonce*}{Statement 4.3~(d)}\etiquette[4.3]{XIV.enon4.3}
Let $X'$ be strictly local of dimension $\leq d$, and the strict
localization of an algebraic scheme over a field. Let
$f\in\Gamma(X',\Oo_{X'})$, and let $U\subset X'$ be the open subset
$U=X'-V(f)$. Then $\cd U\leq d$.
\end{enonce*}

\begin{enonce*}{Lemma 4.4}\etiquette[4.4]{XIV.lem4.4}
\emph{\hyperref[XIV.enon4.3]{4.3}~(d)} implies
\emph{\hyperref[XIV.cor3.4]{3.4}~(d)}.
\end{enonce*}

\noindent
\emph{Proof}.\pageoriginale Consider the diagram
$$
\xymatrix@=1.5cm{
\EE^1_{Y'} \ar[rr]^i \ar[rd]_{f'}&            &
\PP^1_{Y'} \ar[ld]^{g'}\\
                    &       Y'   &
}
$$
where $Y'$ is as in the statement of
\sisi{2.6}{\hyperref[XIV.cor3.4]{3.4}}, and where $\PP^1_{Y'}$ is the
projective line obtained from $\EE^1$ ``by adjoining the section
$Y'_\infty$ at infinity.'' Let $F'$ be a torsion sheaf on $\EE^1$
with $d(\supp F')\leq d$. We have the spectral sequence
$$
H^{p}(\PP^1,R^{q}i_{\ast}F')\Longrightarrow
H^{n}(\EE^1,F'){\quoi}.
$$

Now the $R^{q}i_{\ast}F'$, for $q>0$, are supported on $Y'_\infty$, which is
$Y'$-isomorphic to $Y'$. Since $Y'$ is strictly local, it follows that
$H^{p}(\PP^1,R^{q}i_{\ast}F')=0$ if $p,q>0$. Moreover, by the \emph{proper
base-change theorem} (\hyperrefext[XII][XII.thm5.1]{5.1}), the
cohomology of a torsion sheaf on $\PP^1$ can be computed on the closed
fiber, which is a scheme of cohomological dimension $\leq2$
(\hyperrefext[X][X.cor4.3]{4.3}). Thus
$H^{p}(\PP^1,i_{\ast}F')=0$ if $p>2$. Since we wish to prove that
$H^{n}(\EE^1,F')=0$ for $n>d$, and since $d\geq2$, we are reduced to
proving that $H^{0}(\PP^1,R^{q}i_{\ast}F')=0$ for $q>d$. This group, which is
also isomorphic to $H^{0}(Y'_\infty,R^{\sisi{1}{q}}i_{\ast}F')$, is, by
(\hyperrefext[VIII][VIII.prop4.6]{4.6}), isomorphic to the fiber of
$R^{q}i_{\ast}F'$ at the point $Q$, where $Q$ is the point at infinity of
$\PP^1$ in the closed fiber.

It is enough (\hyperrefext[VII][VII.prop3.3]{3.3} and
\hyperrefext[IX][IX.prop2.9]{2.9}) to treat the case in which $F'$ is
constructible. Since $\delta(\supp F')\leq d$, there is then a closed
subscheme $X'$ of $\PP^1$, with $\dim X'\leq d$, which contains the
support of $F'$. Let $\widetilde X'$ be the localization of $X'$ at
the geometric point $Q$, put
$$
\widetilde U=\widetilde X'-
\fprod{\widetilde X'}{Y'_\infty}{\PP^1},
$$
and let $\widetilde F$ be the sheaf\pageoriginale on $\widetilde U$ induced by $F'$
via the morphism $\widetilde U\rightarrow\PP^1$. The fiber of
$R^{q}i_{\ast}F'$ at $Q$ is then precisely
$H^{\sisi{p}{q}}(\widetilde U,\widetilde F)$
(\hyperrefext[VIII][VIII.thm5.2]{5.2}), and we apply
\hyperref[XIV.enon4.3]{4.3}~(d) to the strictly local scheme
$\widetilde X'$ and the open subset $\widetilde U$.

The following lemma, together with \hyperref[XIV.lem4.4]{4.4}, will
complete the proof.

\begin{enonce*}{Lemma 4.5}\etiquette[4.5]{XIV.lem4.5}
\emph{\hyperref[XIV.cor3.4]{3.4}~(d')} for $d'<d$ implies
\emph{\hyperref[XIV.enon4.3]{4.3}~(d')} for $d'\leq d$.
\end{enonce*}

\noindent
\emph{Proof}. Suppose that \hyperref[XIV.cor3.4]{3.4}~(d') holds for
$d'<d$. By induction, we may suppose that
\hyperref[XIV.enon4.3]{4.3}~(d') holds for $d'<d$. Let $X'$, $f$, and
$U$ be as in the statement \hyperref[XIV.enon4.3]{4.3}~(d). Then
$X'$ is the strict localization, at a \emph{closed} geometric point,
of an algebraic scheme $X_0$ over a field $k$. After replacing $X_0$
by a neighborhood of $x_0$, we may take $X_0$ to be affine and of
dimension $\leq d$. Then $X'$ is the inverse limit of affine schemes
$X_i$ which are étale over $X_0$
(\hyperrefext[VIII][VIII.subsec4.5]{4.5}).

Let $F$ be a torsion sheaf on $U$. We must prove that $H^q(U,F)=0$ for
$q>d$, and once again it is enough to treat the case in which $F$ is
constructible. The data $f$, $F$, and $U$ then arise from analogous
data on one of the $X_i$
(\sisi{\hyperrefext[IX][IX.prop2.7]{2.7}}{\hyperrefext[IX][IX.cor2.7.4]{2.7.4}}).
Thus there is a sheaf $F_i$ on $U_i=X_i-V(f_i)$ which induces $F$ on
$U$. Since $U=\varprojlim U_i$, we have
(\hyperrefext[VII][VII.cor5.8]{5.8})
$$
H^q(U,F)=\varinjlim H^q(U_i,F_i){\quoi}.
$$
It is therefore enough to prove the following fact: for every $i$,
there is a commutative diagram
$$
\xymatrix@=1.5cm{
 U \ar[rr] \ar[rd]&        & U''\ar[ld]\\
   &  U_i   &
}
$$
in which $U\rightarrow U_i$ is the given arrow and
$\cd(U'')\sisi{<}{\leq}d$. Indeed, this implies that\pageoriginale every element of
$H^q(U_i,F_i)$, for $q>d$, has zero image in $H^q(U,F)$.

Suppose that $i=0$, and consider the morphism
$g:X_0\rightarrow\EE^1_k=Y_0$ defined by
$f_0\in\Gamma(X_0,\Oo_{X_0})$. Let $Y''$ be the strict localization
of $\EE^1_k$ at a geometric point above the origin. Thus $Y''$ is the
spectrum of a discrete valuation ring. Put
$X''=\fprod{X_0}{Y''}{Y_0}$.

The morphism $X'\rightarrow X_0$ factors through $X''$, since $X'$ is
strictly local. Moreover, the open subset
$U''=\fprod{X''}{U_0}{X_0}$ is precisely the generic fiber of
$X''/Y''$. It is therefore an affine algebraic scheme over the generic
point of $Y''$, and plainly $\dim U''\leq d-1$.

Let us verify that $\cd U''\leq d$, which will complete the proof. Let
$K$ be the residue field of $Y''$ at its generic point. This field has
cohomological dimension $1$; that is, $\cd(G)=1$, where $G$ is the
Galois group of the separable closure $\bar K$ of $K$
(\hyperrefext[X][X.thm2.2]{2.2}). Put
$\bar U''=\fprod{U''}{(\Spec\bar K)}{\Spec K}$. By the induction
hypothesis on $d$, we have $\cd\bar U''\leq d-1$ (this is
\sisi{\hyperref[XIV.def2.2]{2.2} for a scheme of dimension $\leq d$}
{\hyperref[XIV.cor3.2]{3.2} for a scheme of dimension $<d$}). It then
follows from the Hochschild--Serre spectral sequence
(\hyperrefext[VIII][VIII.prop8.4]{8.4}) that $\cd U''\leq d$:
$$
H^p(G,H^q(\bar U\sisi{'}{''},\bar F))\Longrightarrow H^n(U'',F){\quoi},
$$
which proves Lemma \hyperref[XIV.lem4.5]{4.5}.

\ifx\danslelivre\undefined
\end{document}
\fi

% END INCLUDED SGA4_English_current_buildable_source_20260906/source/14/drafts/SGA4_XIV_LINES_864_1076_English.texfrag


% END INCLUDED SGA4_English_current_buildable_source_20260906/source/14/14_English.tex


\makeatletter
\def\input@path{{../15/}}
\makeatother
% BEGIN INCLUDED SGA4_English_current_buildable_source_20260906/source/15/15_English.tex
\ifx\danslelivre\undefined
% Release engine: XeLaTeX (the LuaTeX compatibility layer breaks Xy-pic output).
\documentclass[10pt]{../fichiers-configuration-latex/smfart}
% BEGIN INCLUDED SGA4_English_current_buildable_source_20260906/source/fichiers-configuration-latex/paquets.tex
\usepackage{polyglossia}
%\usepackage[francais]{babel}
\usepackage{fontspec}
\usepackage[math-style=ISO,bold-style=upright,sans-style=upright]{unicode-math}
%\usepackage[mathletters]{ucs}
%\usepackage[utf8]{inputenc}
%\usepackage[T1]{fontenc}
\usepackage[unicode,hypertexnames=false,colorlinks=true,linkcolor=blue,citecolor=Sepia,urlcolor=magenta]{hyperref}
\setmainlanguage{french}
\setmainfont[Ligatures=TeX]{Linux Libertine O}
% variante possible : TeXGyrePagella ?
\setmathfont{XITS Math}

% maths
\usepackage{amsmath}
%\usepackage{amssymb}
%\usepackage{mathrsfs}
%\usepackage{stmaryrd}
%\usepackage{euscript}

% couleurs
\usepackage[usenames,dvipsnames]{xcolor}

% graphiques
\usepackage{graphicx}
\usepackage[all]{xy}
%\usepackage{tikz}
%\usetikzlibrary{matrix,arrows,calc}
%\tikzset{column sep=3em,row sep=2em,text height=1.5ex, text depth=0.25ex,>=angle 90}
%\tikzset{description/.style={fill=white,inner sep=1pt}}

% hyperliens, 'etiquettes
\usepackage{xr-hyper}
\usepackage{../fichiers-configuration-latex/smfthm}
%\ifx\cacheretiquettes\undefined
%\usepackage[inline]{showlabels} \else \fi % pour rendre visibles les 'etiquettes
%\usepackage{../fichiers-configuration-latex/makeidx} % pour faire un index
%\usepackage[pagebackref,unicode]{hyperref}  % pour voir o`u sont cit'ees les r'ef'erences
% hyperref is loaded once above with the cumulative-reader destination policy.
%\hypersetup{breaklinks=true} %permet le retour `a la ligne dans les liens trop longs. Incompatible avec pdflatex
%\renewcommand*{\backref}[1]{{\small [cit\'e page(s)~#1.]}}

\usepackage{xspace}
\usepackage{pifont} % pour expose XVII

% END INCLUDED SGA4_English_current_buildable_source_20260906/source/fichiers-configuration-latex/paquets.tex

% BEGIN INCLUDED SGA4_English_current_buildable_source_20260906/source/fichiers-configuration-latex/numerotation.tex
\renewcommand{\thefootnote}{\roman{footnote}} % pour mettre les notes en bas de page sous la forme (ii) et non (2), qui ressemble trop (?) a un exposant.

\newcounter{pageoriginal}
\newcommand{\pageoriginale}{\refstepcounter{pageoriginal}\marginpar{\footnotesize\textbf{\thepageoriginal}}\xspace}
\let\pageoriginaled\pageoriginale

\usepackage{ifthen}
\newboolean{original}
\def\sisi{\ifthenelse{\boolean{original}}}
\setboolean{original}{false}

\newcounter{cnote}
\def\nde#1{\ifthenelse{\boolean{original}}{}{\stepcounter{cnote}\def\thefootnote{\arabic{cnote}}\protect\footnote[\value{cnote}]{N.D.E.: #1}}}
\def\ndemark{\ifthenelse{\boolean{original}}{}{\stepcounter{cnote}\def\thefootnote{\arabic{cnote}}\protect\footnotemark[\value{cnote}]}}
\def\ndetext#1{\ifthenelse{\boolean{original}}{}{\def\thefootnote{\arabic{cnote}}\protect\footnotetext[\value{cnote}]{N.D.E.: #1}}}

% END INCLUDED SGA4_English_current_buildable_source_20260906/source/fichiers-configuration-latex/numerotation.tex

% BEGIN INCLUDED SGA4_English_current_buildable_source_20260906/source/fichiers-configuration-latex/macros.tex
\setmainlanguage{french}
\setotherlanguage{english}


%% Fixed, engine-independent release identifier.  The inherited dynamic Git
%% probe used Lua io.popen, exposed build-state-dependent output, and printed
%% a visible question mark outside a Git checkout.
\def\showgitstatus{English release r2, 2026-07-23}


% macros pour LaTeX
\renewcommand{\mathsf}{}

% Police différente pour U+2113 SCRIPT SMALL L
\setmathfont[range="2113]{latinmodern-math.otf}

% Give every inherited semantic label its own stable PDF destination instead
% of inheriting whichever automatic equation/page anchor happened to be live.
% The legacy label remains as a compatibility alias used by existing edges.
\makeatletter
\def\etiquette[#1]#2{%
  \Hy@raisedlink{\hyper@anchorstart{sga4.semantic.#2}\hyper@anchorend}%
  \def\@currentHref{sga4.semantic.#2}%
  \def\@currentlabel{#1}%
  \label{#2}%
}
\makeatother
\def\eqhyperref[#1]#2{(\hyperref[#1]{#2})}
\def\pagehyperref[#1]#2{(page \hyperref[#1]{#2})}
\def\hyperrefext[#1][#2]#3{\hyperref[#2]{#1 #3}}
\def\eqhyperrefext[#1][#2]#3{\hyperref[#2]{#1 #3}}     %cf. exposes 17
\def\neweqhyperrefext[#1][#2]#3{\hyperref[#2]{#1 #3}}  %cf. exposes 17

\def\shat#1{\protect\widehat{#1}}
\def\stilde#1{\protect\widetilde{#1}}
%%%%%%%%%%%%%%%%%%


\def\etcn{etc.\xspace}
\def\etco{etc...\xspace}
\def\etc{\sisi{\etco}{\etcn}}
\def\ldotsn{\sisi{...}{\etcn}}

\def\cf{{cf.}\kern.3em}
\def\Cf{{Cf.}\kern.3em}
\def\ie{{i.e.}\kern.3em}
\def\iev{{i.e.},\xspace}
\def\loccito{loc.\kern2pt cit.\xspace}
\def\loccitn{\textit{loc.\kern2pt cit.}\xspace}
\def\loccit{\sisi{\loccito}{\loccitn}}
\def\EGA{\textup{EGA}\kern.2em}
\def\SGA{\textup{SGA}\kern.3em}

\def\quoi{}

\def\quav{\quad,\quad}
\def\quaet{\quad \text{et}\quad}
\def\fibreproduct#1#2#3{#1 ×_{#3} #2}
\let\fprod\fibreproduct
\def\ccheck#1#2{\mbox{${\check{\check{\rm #1}}}^{\lower2pt\hbox{$\scriptstyle #2$}}$}}
\def\sheaf#1{\protect\mathcal{#1}}
\def\vieuxsheaf#1{{#1}} % cf. problème avec sheaf{a\Phi_i} par exemple...
\DeclareMathOperator{\Hom}{Hom}
\DeclareMathOperator{\can}{can}
\DeclareMathOperator{\Ind}{Ind}
\DeclareMathOperator{\ind}{ind}
%\DeclareMathOperator{\Hhomtop}{\mathcal{H}\!\mathit{omtop}}
\DeclareMathOperator{\Hhomtop}{\mathbf{Homtop}}
\DeclareMathOperator{\Point}{Point}
\DeclareMathOperator{\Pro}{Pro}
\DeclareMathOperator{\Top}{Top}
\DeclareMathOperator{\Mod}{Mod}
\DeclareMathOperator{\Cart}{Cart}
%\DeclareMathOperator{\Mmorsite}{\mathcal{M}\!\mathit{orsite}}
\DeclareMathOperator{\Mmorsite}{\mathbf{Morsite}}
\DeclareMathOperator{\loc}{loc}
\DeclareMathOperator{\Loc}{Loc}
\DeclareMathOperator{\Pic}{Pic}
\DeclareMathOperator{\Et}{Et}
\DeclareMathOperator{\card}{card}
\DeclareMathOperator{\Ens}{Ens}
\DeclareMathOperator{\h}{\underline{H}}
\DeclareMathOperator{\Fonct}{Fonct}
\DeclareMathOperator{\Ker}{Ker}
\DeclareMathOperator{\but}{but}
\DeclareMathOperator{\Ob}{Ob}
\DeclareMathOperator{\ob}{ob}
\DeclareMathOperator{\Lib}{Lib}
\DeclareMathOperator{\esj}{esj}
\DeclareMathOperator{\id}{id}
\DeclareMathOperator{\Id}{Id}
\DeclareMathOperator{\ad}{ad}
\DeclareMathOperator{\Iim}{Im}
\DeclareMathOperator{\pr}{pr}
\DeclareMathOperator{\Cat}{Cat}
\DeclareMathOperator{\car}{car}
\DeclareMathOperator{\SysInd}{SysInd}
\DeclareMathOperator{\Sup}{Sup}
\DeclareMathOperator{\Coker}{Coker}
\DeclareMathOperator{\Filt}{Filt}
\DeclareMathOperator{\acc}{acc}
\DeclareMathOperator{\sur}{sur}
\DeclareMathOperator{\cat}{cat}
\DeclareMathOperator{\Ord}{Ord}
\DeclareMathOperator{\Gr}{Gr}
\DeclareMathOperator{\Cov}{Cov}
\DeclareMathOperator{\Caf}{Caf}
\DeclareMathOperator{\coker}{coker}
\DeclareMathOperator{\coim}{coim}
\DeclareMathOperator{\im}{im}
\DeclareMathOperator{\Ouv}{Ouv}
\DeclareMathOperator{\diag}{diag}
\DeclareMathOperator{\Esp}{Esp}
\DeclareMathOperator{\TOP}{TOP}
\DeclareMathOperator{\Unif}{Unif}
\DeclareMathOperator{\Iind}{\mathbf{Ind}}
\DeclareMathOperator{\Hhomcart}{\mathbf{Homcart}}
\DeclareMathOperator{\Hhomtopess}{\mathbf{Homtopess}}
\DeclareMathOperator{\Ppoints}{\mathbf{Points}}
\DeclareMathOperator{\Ppoint}{\mathbf{Point}}
\DeclareMathOperator{\Ttors}{\mathbf{Tors}}
\DeclareMathOperator{\Ssex}{\mathbf{Sex}}
\DeclareMathOperator{\Ffib}{\mathbf{Fib}}
\DeclareMathOperator{\Vvois}{\mathbf{Vois}}
\DeclareMathOperator{\Ppointess}{\mathbf{Pointess}}
\DeclareMathOperator{\scrHom}{\mathbf{Hom}}
\DeclareMathOperator{\sheafPic}{\mathbf{Pic}}
\DeclareMathOperator{\sheafAut}{\mathbf{Aut}}
\DeclareMathOperator{\Hhomtoplocan}{\mathbf{Homtoplocan}}
\DeclareMathOperator{\Hhomtopan}{\mathbf{Homtopan}}
\DeclareMathOperator{\Pptgeom}{\mathbf{Ptgéom}}
\DeclareMathOperator{\Morsite}{Morsite}
\DeclareMathOperator{\Fib}{Fib}
\DeclareMathOperator{\Fais}{Fais}
\DeclareMathOperator{\kar}{kar}
\DeclareMathOperator{\Groupes}{Groupes}
\DeclareMathOperator{\Restr}{Restr}
\DeclareMathOperator{\Isom}{Isom}
\DeclareMathOperator{\Prol}{Prol}
\DeclareMathOperator{\Res}{Res}
\DeclareMathOperator{\Tors}{Tors}
\DeclareMathOperator{\esp}{esp}
\DeclareMathOperator{\es}{es}
\DeclareMathOperator{\od}{od}
\DeclareMathOperator{\gr}{gr}
\DeclareMathOperator{\sob}{sob}
\DeclareMathOperator{\dfn}{dfn}
\DeclareMathOperator{\Sch}{Sch}
\DeclareMathOperator{\ouv}{ouv}
\DeclareMathOperator{\cons}{cons}
\DeclareMathOperator{\Nerf}{Nerf}
\DeclareMathOperator{\Points}{Points}
\DeclareMathOperator{\Espsob}{Espsob}
\DeclareMathOperator{\Homtop}{Homtop}
\DeclareMathOperator{\Pt}{Pt}
\DeclareMathOperator{\Quot}{Quot}
\DeclareMathOperator{\glob}{glob}
\DeclareMathOperator{\Homfib}{Homfib}
\DeclareMathOperator{\Qouv}{Qouv}
\DeclareMathOperator{\Quotouv}{Quotouv}
\DeclareMathOperator{\ens}{ens}
\DeclareMathOperator{\ab}{ab}
\DeclareMathOperator{\an}{an}
\DeclareMathOperator{\fs}{fs}
\DeclareMathOperator{\red}{red}
\DeclareMathOperator{\Ab}{Ab}
\DeclareMathOperator{\cosq}{cosq}
\DeclareMathOperator{\Ann}{Ann}
\DeclareMathOperator{\Zar}{Zar}


%%%%%%%%%%%%
%%%%% SGA4-2
\DeclareMathOperator{\Sex}{Sex}
\DeclareMathOperator{\vois}{vois}
\DeclareMathOperator{\Vois}{Vois}
\DeclareMathOperator{\Ext}{Ext}
\DeclareMathOperator{\ext}{ext}
\DeclareMathOperator{\Fer}{Fer}
\DeclareMathOperator{\Add}{Add}
\DeclareMathOperator{\Cocart}{Cocart}
\DeclareMathOperator{\cosk}{cosk}
\DeclareMathOperator{\plat}{plat}

\DeclareMathOperator{\faisc}{faisc}
\DeclareMathOperator{\sk}{sk}
\DeclareMathOperator{\ch}{ch}

\def\HH{\mathbf{H}}
\def\qc{\mathrm{qc}}
\def\et{\mathrm{et}}
\def\coh{\mathrm{coh}}

\def\Hvr#1{%
\begin{array}[b]{l}
{\kern2pt\scriptstyle r}\\[-8pt]
{\kern1pt\scriptstyle \vee}\\[-5pt]
H^{#1}
\end{array}}

\def\Hvv#1{%
\begin{array}[b]{l}
\kern2pt\rotatebox{90}{${\scriptscriptstyle\ll}$}\\[-6pt]
H^{#1}
\end{array}}

\def\Hvi#1{%
\begin{array}[b]{l}
{\kern0.2pt\scriptscriptstyle \infty}\\[-8pt]
{\kern1pt\scriptstyle \vee}\\[-5pt]
H^{#1}
\end{array}}

\def\Hv#1{%
\begin{array}[b]{l}
{\kern1pt\scriptstyle \vee}\\[-5pt]
H^{#1}
\end{array}}

\DeclareMathOperator{\cart}{cart}
\DeclareMathOperator{\Ffaisc}{\mathbf{Faisc}}


\def\Ev#1{%
\begin{array}[b]{l}
{\kern1pt\scriptstyle \vee}\\[-5pt]
E'_{#1}
\end{array}}

\DeclareMathOperator{\Cont}{Cont}
\DeclareMathOperator{\Ccont}{\mathscr{C}\!\mathit{ont}}
\DeclareMathOperator{\Site}{Site}
\newcommand\MM{\mathbf{M}}
\def\limtop#1{{\displaystyle{\varprojlim_{#1}\!\mbox{\rm top }}}}
\def\Limtop{\mathop{\varprojlim\!\mbox{\rm top}}}
\DeclareMathOperator{\spec}{spec}
\def\ss{\scriptstyle{\mathscr{S}}}
\DeclareMathOperator{\Gal}{Gal}
\def\smallm{\scriptstyle{\mathscr{M}}}
\DeclareMathOperator{\End}{End}
\DeclareMathOperator{\ass}{ass}
\DeclareMathOperator{\zar}{zar}
\DeclareMathOperator{\Ogg}{Ogg}
\DeclareMathOperator{\res}{res}
\DeclareMathOperator{\tr}{tr}
\DeclareMathOperator{\Supp}{Supp}
\DeclareMathOperator{\codim}{codim}
\DeclareMathOperator{\rk}{rk}
\DeclareMathOperator{\cur}{cur}
\DeclareMathOperator{\cdqc}{cdqc}
\newcommand\DD{\mathbf{D}}
\DeclareMathOperator{\cl}{cl}
\def\JJ{\mathbf{J}}
\def\FF{\mathbf{F}}
\def\BB{\mathbf{B}}
\DeclareMathOperator{\Alget}{Alget}
\DeclareMathOperator{\Ass}{Ass}
\DeclareMathOperator{\supp}{supp}
\DeclareMathOperator{\rad}{rad}
\DeclareMathOperator{\prof}{prof}
\DeclareMathOperator{\cint}{int}
\DeclareMathOperator{\sic}{sic}
\DeclareMathOperator{\Tor}{Tor}
\DeclareMathOperator{\tors}{tors}
\DeclareMathOperator{\Sym}{Sym}
\DeclareMathOperator{\parf}{parf}
\DeclareMathOperator{\tor}{tor}
\DeclareMathOperator{\defn}{defn}
\DeclareMathOperator{\Div}{Div}
\DeclareMathOperator{\Pol}{Pol}
\DeclareMathOperator{\Fex}{Fex}
\DeclareMathOperator{\Lie}{Lie}
\DeclareMathOperator{\tot}{tot}
\DeclareMathOperator{\Aut}{Aut}
\DeclareMathOperator{\HOM}{HOM}
\DeclareMathOperator{\pch}{pch}
\DeclareMathOperator{\Ch}{Ch}
\DeclareMathOperator{\EXT}{EXT}
\DeclareMathOperator{\PIC}{PIC}
\DeclareMathOperator{\Rf}{Rf}
\DeclareMathOperator{\inj}{inj}
\newcommand\VV{\mathbf{V}}

\def\Xv{%
\begin{array}[b]{l}
{\kern1pt\scriptstyle \vee}\\[-5pt]
X
\end{array}}

%%%%%%%
\newcommand\TT{\mathbf{T}}
\newcommand\XX{\mathbf{X}}
\newcommand\EE{\mathbf{E}}
%%%%%%

\def\qedsymbol{\textup{C.Q.F.D.}}
\let\CQFD\qed
\let\cqfd\qed

%%%%%\bigdot
\font\bigsymb=cmsy10 at 4pt
\def\bigdot{{\kern1.2pt\raise 1.5pt\hbox{\bigsymb\char15}}}
%%%%%

\def\Lright{\mathop{\vcenter{\hbox{\hbox{$\longrightarrow$}\lower8pt \hbox{\kern-16pt\hbox{$\longrightarrow$}}}}}}

\def\Aa{{\mathcal A}}
\def\Bb{{\mathcal B}}
\def\Cc{{\mathcal C}}
\def\Dd{{\mathcal D}}
\def\Ee{{\mathcal E}}
\def\Ff{{\mathcal F}}
\def\Gg{{\mathcal G}}
\def\Hh{{\mathcal H}}
\def\Ii{{\mathcal I}}
\def\Jj{{\mathcal J}}
\def\Kk{{\mathcal K}}
\def\Ll{{\mathcal L}}
\def\Mm{{\mathcal M}}
\def\Nn{{\mathcal N}}
\def\Oo{{\mathcal O}}
\def\Pp{{\mathcal P}}
\def\Qq{{\mathcal Q}}
\def\Rr{{\mathcal R}}
\def\Ss{{\mathcal S}}
\def\Tt{{\mathcal T}}
\def\Uu{{\mathcal U}}
\def\Vv{{\mathcal V}}
\def\Ww{{\mathcal W}}
\def\Xx{{\mathcal X}}
\def\Yy{{\mathcal Y}}
\def\Zz{{\mathcal Z}}

\newcommand{\CC}{\mathbf{C}}
\newcommand{\GG}{\mathbf{G}}
\newcommand{\KK}{\mathbf{K}}
\newcommand{\LL}{\mathbf{L}}
\newcommand{\NN}{\mathbf{N}}
\newcommand{\PP}{\mathbf{P}}
\newcommand\QQ{\mathbf{Q}}
\newcommand\RR{\mathbf{R}}
\newcommand\ZZ{\mathbf{Z}}

\def\to{\mathchoice{\longrightarrow}{\rightarrow}{\rightarrow}{\rightarrow}}
\def\from{\mathchoice{\longleftarrow}{\leftarrow}{\leftarrow}{\leftarrow}}
\def\mto{\mathchoice{\longmapsto}{\mapsto}{\mapsto}{\mapsto}}
\def\mfrom{\mathrel{\from\joinrel\mapstochar\,}}
\def\hto{\mathrel{\lhook\joinrel\to}}
\def\hfrom{\mathrel{\from\joinrel\rhook}}

\def\implies{\mathchoice{\Longrightarrow}{\Rightarrow}{\Rightarrow}{\Rightarrow}}
\def\implied{\mathchoice{\Longleftarrow}{\Leftarrow}{\Leftarrow}{\Leftarrow}}
\def\iff{\mathchoice{\Longleftrightarrow}{\Leftrightarrow}{\Leftrightarrow}{\Leftrightarrow}}

\def\To#1{\mathchoice{\xrightarrow{\labelstyle~~#1~}}{\stackrel{#1}{\longrightarrow}}{}{}}
\def\Mto#1{\mathrel{\mapstochar\To{#1}}}
\def\From#1{\mathchoice{\xleftarrow{\labelstyle~~#1~}}{\stackrel{#1}{\longleftarrow}}{}{}}

\newcommand{\isomto}{\overset{\sim}{\longrightarrow}}
\newcommand{\isomfrom}{\overset{\sim}{\longleftarrow}}

\newcommand{\IFF}{\unskip~$\iff$~\ignorespaces}
\newcommand{\THEN}{\unskip~$\implies$~\ignorespaces}

\let\leq\leqslant
\let\geq\geqslant
\let\le\leqslant
\let\ge\geqslant

\def\cf{{cf.}\kern.3em}
\def\Cf{{Cf.}\kern.3em}
\def\ie{{i.e.}\kern.3em}
\def\iev{{i.e.},\xspace}
%\def\loccito{loc.\kern2pt cit.\xspace}
%\def\loccitn{\textit{loc.\kern2pt cit.}\xspace}
%\def\loccit{\sisi{\loccito}{\loccitn}}
\def\EGA{\textup{EGA}\kern.2em}
\def\SGA{\textup{SGA}\kern.3em}

\let\oldS\S
\def\S{\oldS\kern .15em}
\def\ptbl{.\kern .2em}
\def\resp{resp.\kern.3em}
\def\Exp{Exp.\kern.3em}

\let\bar\overline
\let\hat\widehat
\let\tilde\widetilde
\let\dpl\displaystyle
\let\ldots\dots
\let\emptyset\varnothing
\let\skpt\skippointrait
\let\epsilon\varepsilon

\def\sfrac#1#2{{#1}/{#2}}

\DeclareMathOperator{\cd}{cd}
%\DeclareMathOperator{\Hhom}{\mathcal{H}\!\mathit{om}}
\DeclareMathOperator{\Hhom}{\mathbf{Hom}}
%\DeclareMathOperator{\Eext}{\mathcal{E}\!\mathit{xt}}
\DeclareMathOperator{\Eext}{\mathbf{Ext}}
\DeclareMathOperator{\Spec}{Spec}
\DeclareMathOperator{\Tr}{Tr}

\let\emptyset=\varnothing

%%%%%%%%%%%%%%%
%%% expose XVII
%%%%%%%%%%%%%%%

\makeatletter
\let\old@optexpose\optexpose
\newcommand{\empty@optexpose}[1]{}
\newcommand{\paren@optexpose}[1]{\old@optexpose{#1}(}
\newcommand{\refnoexpose}[1]{\let\optexpose\empty@optexpose\ref{#1}\let\optexpose\old@optexpose}%\refnoexpose{X.cor5.2} engendre "5.2" au lieu de "X 5.2"
\newcommand{\neweqref}[1]{\let\optexpose\paren@optexpose\textup{\ref{#1})}\let\optexpose\old@optexpose}%\neweqref{XII.eq4.1} engendre "XII (4.1)" au lieu de "(XII 4.1)"
\makeatother
\newcommand{\refvb}[2][VI${}_{\text{B}}$ ]{\sisi{#1\refnoexpose{#2}}{\ref{#2}}}

\newcommand{\leftexp}[2]{{\vphantom{#2}}^{#1}{#2}}
\newcommand{\AAA}{\mathbf{A}}
\newcommand{\Sf}{\mathfrak{S}}
\DeclareMathOperator*{\quotprojlim}{\text{``$\varprojlim$''}}
\DeclareMathOperator*{\quotinjlim}{\text{``$\varinjlim$''}}
\DeclareMathOperator{\Fl}{F\sisi{\ell}{l}}
\DeclareMathOperator{\Pl}{P\sisi{\ell}{l}}
%\DeclareMathOperator{\cl}{c\sisi{\ell}{l}}
\DeclareMathOperator{\tdf}{tdf}
\DeclareMathOperator{\rg}{rg}
\newcommand{\iInd}{\sisi{I}{i}nd}
\newcommand{\iind}{ind}
\newcommand{\ppro}{pro}
\newcommand{\tTor}{\sisi{T}{t}or}
\newcommand{\ttor}{tor}
\newcommand{\chap}{Chap.}
\newcommand{\Ffex}{\Ff\mathit{ex}}
\newcommand{\sheafk}{\underline{k}}
\newcommand{\sheafn}{\underline{n}}
\newcommand{\sheafGamma}{\underline{\Gamma}}

% Repeated references to an already typeset footnote must not use
% \footnotemark without a matching \footnotetext: Hyperref would then create
% a live annotation to a destination that cannot exist.  This helper keeps the
% current printed footnote mark and links it to an explicit stable note target.
\makeatletter
\newcommand{\sgarepeatedfootnotemark}[1]{%
  \begingroup
  \edef\@thefnmark{\thefootnote}%
  \hyperlink{#1}{\@makefnmark}%
  \endgroup}
\makeatother

\sisi{
\DeclareMathOperator{\simarrow}{\rule{.15mm}{1.8mm}\kern -3.4pt\rightsquigarrow}
\newcommand{\sheafell}{\underline{\ell}}
\newcommand{\sheafZZ}{\underline{\ZZ}}
}{
\newcommand{\nderef}[1]{la N.D.E.~(\ref{#1}) page~\pageref{#1}}
}

%%%%%%%%%%%%%%%%%
%%%%% expose XVII
%%%%%%%%%%%%%%%%%

% END INCLUDED SGA4_English_current_buildable_source_20260906/source/fichiers-configuration-latex/macros.tex


\externaldocument{../01/01_English.tex}
\externaldocument{../02/02_English.tex}
\externaldocument{../03/03_English.tex}
\externaldocument{../04/04_English.tex}
\externaldocument{../05/05_English.tex}
\externaldocument{../05bis/05bis_English.tex}
\externaldocument{../06/06_English.tex}
\externaldocument{../07/07_English.tex}
\externaldocument{../08/08_English.tex}
\externaldocument{../09/09_English.tex}
\externaldocument{../10/10_English.tex}
\externaldocument{../11/11_English.tex}
\externaldocument{../12/12_English.tex}
\externaldocument{../13/13_English.tex}
\externaldocument{../14/14_English.tex}

\setboolean{original}{false}

\begin{document}

\begin{center}
XV. Acyclic morphisms\\
M. Artin
\end{center}

\selectlanguage{english}version: \showgitstatus
\tableofcontents

\else
\chapter{Acyclic morphisms}
%\addtocontents{toc}{M. Artin}
\begin{center}
M. Artin
\end{center}
\fi
\phantomsection\MakeLinkTarget*{sga4.exp.XV}\label{sga4.exp.XV}

% BEGIN INCLUDED SGA4_English_current_buildable_source_20260906/source/15/drafts/SGA4_XV_LINES_43_400_English.texfrag
\setcounter{pageoriginal}{167}

% To do:
% standardize notation for restrictions, e.g. F|X -> F_{|X}.
% add ~ to avoid short mathematical terms standing alone at the start of a line.
% replace \ref{...} by (\ref{}); cf. for example the statement of 1.15.
% do not set remarks in italics (cf. Remark 1.18).

\section*{Introduction}\pageoriginale

Let $g:X\rightarrow Y$ be a morphism of schemes. In the first section
we study conditions under which $g$ is acyclic, that is, such that for
every torsion sheaf $F$ on $Y$ one has
$H^{q}(Y,F)\simeq H^{q}(X,g^{\ast}F)$, and such that this remains true
when $Y$ is replaced by a $Y'$ étale over $Y$. A natural necessary
condition for acyclicity is that the geometric fibers be acyclic, that
is, have trivial cohomology for constant torsion sheaves. We shall see
that this condition, together with a local condition called local
acyclicity of $g$ (\hyperref[XV.def1.11]{1.11}), implies that $g$ is
acyclic.

The fundamental theorem is that a smooth morphism is locally acyclic
for torsion sheaves of order prime to the residue characteristics
\hyperref[XV.thm2.1]{2.1}. This will imply the base-change theorem for
a smooth morphism \hyperrefext[XVI][XVI.thm1.1]{1.1}, which will be
developed, together with its first consequences, in the next exposé.

The present exposé is independent of Exposés XI to XIV, and in
particular of the ``base-change theorem for a proper morphism.'' The
latter will again enter in an essential way, but in conjunction with
the fundamental theorem of the present exposé, beginning with the next
exposé.

\section{Generalities on globally and locally acyclic
  morphisms}\etiquette[1]{XV.sec1}\pageoriginale

% g' is not defined as g x_Y Y' here, whereas it is in 1.3.

\begin{enonce*}{Proposition 1.1}\etiquette[1.1]{XV.prop1.1}
Let $g:X\rightarrow Y$ be a morphism of schemes. The following
assertions {\rm(i)} to {\rm(iii)} are equivalent:
\begin{itemize}
\item[(i)] For every étale morphism $Y'\rightarrow Y$ (where $Y'$ may
  be taken affine over an affine open subset of $Y$) and every sheaf of
  sets $F$ on $Y'$, the canonical morphism $F\rightarrow
  g'_{\ast}{g'}^{\ast}F$ is injective (\resp bijective).

\item[(ii)] For every étale morphism $Y'\rightarrow Y$ (where $Y'$ may
  be taken affine over an affine open subset of $Y$) and every sheaf of
  sets $F$ on $Y'$, the canonical morphism $H^{0}(Y',F)\rightarrow
  H^{0}(X',{g'}^{\ast}F)$ is injective (\resp bijective).

\item[(iii)] For every locally quasi-finite morphism $Y'\rightarrow Y$
  and every sheaf of sets $F$ on $Y'$, the canonical morphism
  $H^{0}(Y',F)\rightarrow H^{0}(X',{g'}^{\ast}F)$ is injective
  (\resp bijective).

\item[(ii bis)] (If $g$ is quasi-compact and quasi-separated.) As in
  {\rm(ii)}, but with $F$ a constant sheaf of the form $I_{Y'}$, where
  $I$ is a prescribed set with $\card I\geq 2$.
\end{itemize}
\end{enonce*}

\noindent
\emph{Proof}. The equivalence of (i) and (ii) follows immediately from
the definitions, while (iii) $\Rightarrow$ (ii) $\Rightarrow$ (ii bis)
are trivial. We have (ii bis) $\Rightarrow$ (ii), since we may suppose
$Y$ affine, so that $X$ and $Y$ are quasi-compact and quasi-separated,
and may then apply \hyperrefext[XII][XII.prop6.5]{6.5} (i). For the
implication (i) $\Rightarrow$ (iii), we first record the following
lemma:

\begin{enonce*}{Lemma 1.2}\etiquette[1.2]{XV.lem1.2}
Let\pageoriginale $f:Y'\rightarrow Y$ be a locally quasi-finite
morphism. Then $f$ is ``locally finite on $Y'$ for the étale
topology,'' that is, there are commutative diagrams
$$\vcenter{
\xymatrix@=1.75cm{
Y'_i        \ar[r]^{a_i}   \ar[d]^{b_i}& Y'\ar[d]^{f}\\
Y_i          \ar[r]^{c_i}              & Y
}}\leqno{(*)}
$$
in which the $a_i$ and $c_i$ are étale, the $b_i$ are finite, and the
$a_i$ form a covering of $Y'$ \footnote[1]{\Cf \EGA IV 18.12.1.}.
\end{enonce*}

Indeed, let $y'$ be a point of $Y'$ and let $y=f(y')$. It is enough to
find a diagram $({\ast})$ such that $y'$ lies in the image of $Y'_i$.
Let $\widetilde{Y}$ be the spectrum of the henselization of the ring
$\Oo_{Y,y}$ \hyperrefext[VIII][VIII.subsec4.1]{4.1}, and let $Z$ be
the (ordinary) localization of
$\widetilde{Y}'=\fprod{Y'}{\widetilde{Y}}{Y}$ at the unique point
lying over $y'$ and the closed point $\widetilde{y}$ of
$\widetilde{Y}$. By \hyperrefext[VIII][VIII.subsec4.1]{4.1} (ii),
% label for (ii)?
$Z\rightarrow \widetilde{Y}$ is finite, and $Z$ is an induced open
subscheme of $\widetilde{Y}'$. Since $\widetilde{Y}$ is a limit of
schemes $Y_i$ étale over $Y$, one sees immediately from \EGA IV 9 that,
for a suitable $Y_i$, there is an open subscheme $Z_i$ of
$\fprod{Y'}{Y_i}{Y}$ such that
$Z=\fprod{Z_i}{\widetilde{Y}}{Y_i}$; in $({\ast})$ it is then enough to
take $Y'_i=Z_i$. \cqfd


Now suppose that (i) of \hyperref[XV.prop1.1]{1.1} holds. It is clear
that assertion (iii) of \hyperref[XV.prop1.1]{1.1} is equivalent to the
assertion obtained from (i) by replacing the étale morphism
$Y'\rightarrow Y$ by a locally quasi-finite morphism. In this form,
the assertion (for given $Y$ and $Y'$) is local on $Y$ and $Y'$ for
the étale topology; hence the lemma reduces us to the case of a
\emph{finite} morphism $f:Y'\rightarrow Y$.\pageoriginale Let
$f':X'\rightarrow X$ be the morphism obtained from $f$ by base change
along $g$. Then the base-change morphism
$g^{\ast}f_{\ast}F\rightarrow {f'}_{\ast}{g'}^{\ast}F$ is bijective
\hyperrefext[VIII][VIII.cor5.6]{5.6}. By (i),
\sisi{the adjunction morphism
$f_{\ast}F\rightarrow g_{\ast}g^{\ast}f_{\ast}F$ is injective
(resp.~bijective)}{the adjunction morphism
$f_{\ast}F\rightarrow g_{\ast}g^{\ast}f_{\ast}F$ is injective
(resp.~bijective)}.
Consequently, the morphism $H^{0}(Y',F)\xrightarrow{\epsilon}
H^{0}(Y',g'_{\ast}{g'}^{\ast}F)$ in (iii) is the composite
$H^{0}(Y',F)=H^{0}(Y,f_{\ast}F)\xrightarrow{\epsilon'}
H^{0}(Y,g_{\ast}g^{\ast}f_{\ast}F)=H^{0}(X,g^{\ast}f_{\ast}F)\simeq
H^{0}(X,{f'}_{\ast}{g'}^{\ast}F)=H^{0}(X',{g'}^{\ast}F)$.
Since $\epsilon'$ is injective (\resp bijective), the same is true of
$\epsilon$. This completes the proof of \hyperref[XV.prop1.1]{1.1}.

\begin{enonce*}{Definition 1.3}\etiquette[1.3]{XV.def1.3}
A morphism $g:X\rightarrow Y$ having one of the equivalent properties
{\rm(i)} to {\rm(iii)} (\resp the corresponding parenthesized
properties) of \hyperref[XV.prop1.1]{1.1} is called a
\emph{$(-1)$-acyclic morphism} (\resp a \emph{$0$-acyclic morphism}).
The morphism is said to be \emph{universally} $(-1)$-acyclic
(\resp $0$-acyclic) if, for every base-change morphism
$Y'\rightarrow Y$, the morphism
$g'=\fprod{g}{Y'}{Y}:\fprod{X}{Y'}{Y}\rightarrow Y'$ is
$(-1)$-acyclic (\resp $0$-acyclic). Finally, a scheme $X$ is said to
be $(-1)$-acyclic (\resp $0$-acyclic) if it is nonempty (\resp
connected and nonempty).
\end{enonce*}

\begin{enonce*}{Corollary 1.4}\etiquette[1.4]{XV.cor1.4}
Let $g:X\rightarrow Y$ be a $0$-acyclic morphism. Then for every
locally quasi-finite morphism $Y'\rightarrow Y$ and every sheaf of
groups $F$ on $Y'$, the canonical morphism
$$
H^{1}(Y',F)\longrightarrow H^{1}(X',{g'}^{\ast}F)
$$
is injective.
\end{enonce*}

This follows immediately from \hyperrefext[XII][XII.prop6.5]{6.5} (i).

\begin{enonce*}{Corollary 1.5}\etiquette[1.5]{XV.cor1.5}
Let\pageoriginale $g:X\rightarrow Y$ be a morphism. If $g$ is
surjective (\ie its fibers are nonempty, \ie its fibers are
$(-1)$-acyclic), then $g$ is $(-1)$-acyclic. The converse holds if $g$
is quasi-compact and quasi-separated.
\end{enonce*}

The first assertion is clear in the form \hyperref[XV.prop1.1]{1.1}
(ii), by considering the fibers of the sheaves in question. The second
follows from the definition in the form \hyperref[XV.prop1.1]{1.1}
(iii) and the limit argument of \hyperrefext[VII][VII.sec5]{5}, taking
into account that, for every $y\in Y$, $y=\Spec k(y)$ is the inverse
limit of affine schemes locally quasi-finite over $Y$, namely the
affine neighborhoods of $y$ in $\bar{\{y\}}$ (endowed with the induced
reduced structure). \sisi{}{--} Note that \hyperref[XV.cor1.5]{1.5}
implies that, when $g$ is quasi-compact and quasi-separated, the
$(-1)$-acyclicity of $g$ already implies that $g$ is
\emph{universally} $(-1)$-acyclic.

\begin{enonce*}{Proposition 1.6}\etiquette[1.6]{XV.prop1.6}
Let $g:X\rightarrow Y$ be a morphism of schemes, let
$\LL\subset\PP$ be a set of prime numbers, and let $n\in\NN$ be a
natural number (\resp let $n=1$). The following assertions {\rm(i)} to
{\rm(iii)} (\resp the corresponding parenthesized assertions) are
equivalent. If, moreover, $g$ is quasi-compact and quasi-separated,
then {\rm(i)} to {\rm(iii)} (\resp the corresponding parenthesized
assertions) are also equivalent to {\rm(iv)} (\resp the corresponding
parenthesized assertion).
\begin{enumerate}
\renewcommand{\theenumi}{\roman{enumi}}
\renewcommand{\labelenumi}{(\theenumi)}
\item For every étale morphism $Y'\rightarrow Y$ and every abelian
  $\LL$-torsion sheaf (\resp sheaf of
  $\text{\rm ind}$-$\LL$-groups) \emph{\hyperrefext[IX][IX.def1.5]{1.5}}
  $F$ on $Y'$, one has
$$
\begin{cases}
F\xrightarrow{\sim} {g'}_{\ast}{g'}^{\ast}F, &\\
(R^{q}{g'}_{\ast}){g'}^{\ast}F=0 & \mbox{if}\quad 1\leq q\leq n\quad
(\resp \mbox{if\; } q=1){\quoi}.
\end{cases}
$$

\item For\pageoriginale every étale morphism $Y'\rightarrow Y$ and
every abelian $\LL$-torsion sheaf (\resp sheaf of
ind-$\LL$-groups) $F$ on $Y'$, the canonical morphism
$$
H^{q}(Y',F)\longrightarrow H^{q}(X',{g'}^{\ast}F)
$$
is bijective if $q\leq n$ and injective if $q=n+1$ (\resp is
bijective if $q\leq1$).

\item For every locally quasi-finite morphism $Y'\rightarrow Y$ and
every abelian $\LL$-torsion sheaf (\resp sheaf of
ind-$\LL$-groups) $F$ on $Y'$, the morphism
$$
H^{q}(Y',F)\longrightarrow H^{q}(X', {g'}^{\ast}F)
$$
is bijective if $q\leq n$ and injective if $q=n+1$ (\resp is
bijective if $q\leq1$).

\item The morphism $g$ is surjective, and for every locally
quasi-finite morphism $Y'\rightarrow Y$, every $\nu>0$, and every
$\ell\in\LL$, the canonical morphism
$$
H^{q}(Y',\ZZ/\ell^{\nu})\longrightarrow H^{q}(X',\ZZ/\ell^{\nu})
$$
is \emph{surjective} for $0\leq q\leq n$ (\resp for every locally
quasi-finite morphism $Y'\rightarrow Y$ and every ordinary finite
$\LL$-group $G$, the canonical morphism
$$
H^{i}(Y',G)\longrightarrow H^{i}(X',G)
$$
is \emph{surjective} for $i=0,1$).
\end{enumerate}
\end{enonce*}

\noindent
\emph{Proof of}\pageoriginale \hyperref[XV.prop1.6]{1.6}. The
implications (ii) $\Leftarrow$ (iii) $\Rightarrow$ (iv) are trivial.
% trivial surjectivity? (g is not necessarily coherent)
Now $(R^{q}{g'}_{\ast}){g'}^{\ast}F$ is the sheaf associated with the
presheaf $R^{q}(Y'')=H^{q}(X'',{g''}^{\ast}F)$
($Y''\rightarrow Y'$ étale). If (ii) holds, we have
$H^{q}(Y'',F)\xrightarrow{\sim}H^{q}(X'',{g''}^{\ast}F)$, so the
associated sheaf is zero\sisi{ if $1\leq q\leq n$}{}, \ie
(ii) $\Rightarrow$ (i). The implication (i) $\Rightarrow$ (ii) follows
from \hyperref[XV.prop1.1]{1.1} (i) (which corresponds to the case
$n=0$), together with the Leray spectral sequence
$$
H^{p}(Y',(R^{q}{g'}_{\ast}){g'}^{\ast}F)\Longrightarrow
H^{p+q}(X',{g'}^{\ast}F)
$$
in the abelian case, and the exact sequence
$$
0\longrightarrow H^{1}(Y',{g'}_{\ast}{g'}^{\ast}F)\longrightarrow
H^{1}(X',{g'}^{\ast}F)\longrightarrow
H^{0}(Y',(R^{1}{g'}_{\ast}){g'}^{\ast}F){\quoi}
$$
in the case of sheaves of groups.

(i) $\Rightarrow$ (iii): Since (i) and (ii) are equivalent, it is
clear that (iii) is equivalent to the statement obtained from (i) by
replacing the word ``étale'' by the words ``locally quasi-finite.'' In
this form, for given $Y$ and $Y'$, it is an assertion local on $Y'$ and
on $Y$ for the étale topology, so \hyperref[XV.lem1.2]{1.2} reduces us
to the case in which $f:Y'\rightarrow Y$ is a \emph{finite} morphism.
Let $f':X'\rightarrow X$ be the morphism
\sisi{obtained from~$f$ by base change along~$g$}{obtained from~$f$ by base change along~$g$}.
By \hyperrefext[VIII][VIII.prop5.5]{5.5}
(\resp \hyperrefext[VIII][VIII.rem5.8]{5.8}) (applied three times), we
have
$$
f_{\ast}(R^{q}{g'}_{\ast}){g'}^{\ast}F\simeq
R^{q}(fg')_{\ast}{g'}^{\ast}F\simeq
R^{q}(gf')_{\ast}{g'}^{\ast}F\simeq
(R^{q}g_{\ast}){f'}_{\ast}{g'}^{\ast}F\simeq
(R^{q}g_{\ast})g^{\ast}(f_{\ast}F)
$$
and the last of these is zero, for the values of $q$ in question, if
(i) holds. Thus (i) $\Rightarrow$ (iii).

It remains to prove (iv) $\Rightarrow$ (iii), and this follows
immediately\pageoriginale from \hyperrefext[XII][XII.prop6.5]{6.5}.

\begin{enonce*}{Definition 1.7}\etiquette[1.7]{XV.def1.7}
Let $n\in\NN$ and $\LL\subset\PP$. A morphism
$g:X\rightarrow Y$ satisfying one of the equivalent conditions
{\rm(i)} to {\rm(iii)} (\resp the corresponding parenthesized
conditions) of \hyperref[XV.prop1.6]{1.6} is called an
\emph{$n$-acyclic morphism} (\resp a \emph{$1$-aspherical morphism})
for $\LL$. We say that $g$ is \emph{acyclic} for $\LL$ if it is
$n$-acyclic for $\LL$ for every $n$. We say that $g$ is
\emph{universally $n$-acyclic} (\resp \emph{universally
$1$-aspherical}) for $\LL$ if, for every $Y'\rightarrow Y$, the
morphism $g'=\fprod{g}{Y'}{Y}$ is $n$-acyclic (\resp $1$-aspherical)
for $\LL$. Finally, a scheme $X$ is said to be $n$-\emph{acyclic}
(\resp $1$-\emph{aspherical}) for $\LL$ if it is $0$-acyclic in the
sense of \hyperref[XV.def1.3]{1.3} and if, for every $\ell\in\LL$ and
every $1\leq q\leq n$, one has $H^{q}(X,\ZZ/\ell)=0$ (\resp and if,
for every ordinary finite $\LL$-group $G$, one has $H^{1}(X,G)=0$).
\end{enonce*}

\begin{enonce*}[remark]{Remarks 1.8}\etiquette[1.8]{XV.rem1.8}
\begin{enumerate}
\renewcommand{\theenumi}{\alph{enumi}}
\renewcommand{\labelenumi}{\theenumi)}
\item If $g$ is quasi-compact and quasi-separated, then, in order for
  $g$ to be universally $n$-acyclic, it is enough that $g'$ be
  $n$-acyclic whenever $Y'\rightarrow Y$ is locally of finite
  presentation. We do not need this fact and leave its proof, which is
  an exercise in passage to the limit, to the reader.

\item We do not know whether a morphism that is $n$-acyclic for
  $\LL$ is necessarily universally $n$-acyclic for $\LL$. The same
  question arises for the local variants below
  \hyperref[XV.def1.11]{1.11}.

\item When $n=0$, we recover the notion of
  \hyperref[XV.def1.3]{1.3}. Of course,
  \hyperref[XV.prop1.6]{1.6} and \hyperref[XV.def1.7]{1.7} could have
  been formulated so as to include the case $n=-1$. The following
  proposition is also valid in that case.\pageoriginale
\end{enumerate}
\end{enonce*}

\begin{enonce*}{Proposition 1.9}\etiquette[1.9]{XV.prop1.9}
\begin{enumerate}
\renewcommand{\theenumi}{\roman{enumi}}
\renewcommand{\labelenumi}{(\theenumi)}
\item Let $X\xrightarrow{g}Y\xrightarrow{h}Z$ be morphisms of schemes,
  let $\LL\subset\PP$, and let $n\in\NN$. Suppose that $g$ is
  $0$-acyclic (\resp $1$-aspherical for $\LL$, \resp $n$-acyclic for
  $\LL$). Put $f=hg$. Then, for every sheaf of sets (\resp of
  ind\sisi{-}{ }$\LL$-groups, \resp abelian $\LL$-torsion sheaf) $F$
  on $Z$, there is an isomorphism
  $f_{\ast}f^{\ast}F\xrightarrow{\sim}h_{\ast}h^{\ast}F$ (\resp
  isomorphisms $(R^{p}f_{\ast})f^{\ast}F\xrightarrow{\sim}
  (R^{p}h_{\ast})h^{\ast}F$ for $p\leq1$, \resp for $p\leq n$). In
  particular, $h$ is $0$-acyclic (\resp $1$-aspherical for $\LL$,
  \resp $n$-acyclic for $\LL$) if and only if $f$ is.

\item Let $\{g_{\alpha}:X_{\alpha}\rightarrow Y_{\alpha}\}$ be an
  inverse system of morphisms of schemes such that the morphisms
  $X_{\alpha}\rightarrow X_{\beta}$ and
  $Y_{\alpha}\rightarrow Y_{\beta}$ are affine and the $g_{\alpha}$
  are quasi-compact and quasi-separated. Let $g:X\rightarrow Y$ be the
  inverse limit of the $g_{\alpha}$
  \hyperrefext[VII][VII.subsec5.1]{5.1}. Then, if all the $g_{\alpha}$
  are $0$-acyclic (\resp $1$-aspherical for $\LL$, \resp $n$-acyclic
  for $\LL$), the same is true of $g$.
\end{enumerate}
\end{enonce*}

\noindent
\emph{Proof}. For (i), first observe that
$$
h_{\ast}h^{\ast}F\xrightarrow{\sim}
h_{\ast}(g_{\ast}g^{\ast})h^{\ast}F\simeq f_{\ast}f^{\ast}F{\quoi}.
$$
The second assertion of (i) is an immediate consequence of the Leray
spectral sequence
$$
(R^{p}h_{\ast})(R^{q}g_{\ast})f^{\ast}F\Longrightarrow
(R^{p+q}f_{\ast})f^{\ast}F{\quoi}.
$$
Finally, the nonabelian assertion follows from the exact sequence
\hyperrefext[XII][XII.prop3.2]{3.2}.

% END INCLUDED SGA4_English_current_buildable_source_20260906/source/15/drafts/SGA4_XV_LINES_43_400_English.texfrag


% BEGIN INCLUDED SGA4_English_current_buildable_source_20260906/source/15/drafts/SGA4_XV_LINES_401_700_English.texfrag
Assertion\pageoriginale (ii) follows readily from the theory of passage
to the limit \hyperrefext[VII][VII.sec5]{5} and VI. % label for chapter?

\begin{enonce*}{Corollary 1.9.1}\etiquette[1.9.1]{XV.cor1.9.1}
Let $g:X\rightarrow Y$ be a quasi-compact and quasi-separated morphism,
and let $n\in \NN$, $\LL\subset \PP$. If $g$ is
$n$-acyclic for $\LL$ (\resp $1$-aspherical for
$\LL$), then, for every geometric point $\bar{y}$ of $Y$
algebraic over a point $y$ of $Y$, the geometric fiber $X_{\bar{y}}$
is $n$-acyclic for $\LL$ (\resp $1$-aspherical for $\LL$).
\end{enonce*}

Indeed, $\bar{y}$ is the inverse limit of affine schemes $Y_i$ that
are locally quasi-finite over $Y$, and since the
$X_i=\fprod{X}{Y_i}{Y}\rightarrow Y_i$ are $n$-acyclic for
$\LL$ (\resp\ldots), the same holds for
$X_{\bar{y}}\rightarrow \bar{y}$ by \hyperref[XV.prop1.9]{1.9} (ii), which gives the
asserted conclusion.

\begin{enonce*}{Proposition 1.10}\etiquette[1.10]{XV.prop1.10}
Let $g:X\rightarrow Y$ be a morphism of schemes, let $\LL\subset
\PP$, and let $n\in \NN$. Then conditions {\rm(i)} through
{\rm(iv)} below are equivalent.
\begin{enumerate}
\renewcommand{\theenumi}{\roman{enumi}}
\renewcommand{\labelenumi}{(\theenumi)}
\item Let $\bar{y}$ be a geometric point of $Y$ and $\bar{x}$ a geometric
point of $X$ lying over $\bar{y}$. Let $\widetilde{X}$ and
$\widetilde{Y}$ be the strict localizations of $X$ and $Y$ at
$\bar{x}$ and $\bar{y}$, respectively, and let
$\widetilde{g}:\widetilde{X}\rightarrow
\widetilde{Y}$ be the morphism induced by $g$. Then $\widetilde{g}$ is
$0$-acyclic (\resp $1$-aspherical for $\LL$, \resp
$n$-acyclic for $\LL$).

\item For every diagram with Cartesian squares
$$
\vcenter{
\xymatrix@C=2cm@R=1.25cm{
X   \ar[d]^{g}& X'\ar[l]_{j}  \ar[d]^{g'} & X'' \ar[l]_{j'} \ar[d]^{g''}\\
Y    & Y'\ar[l]_{i} &   Y'' \ar[l]_{i'} 
}}\leqno{(\ast)}
$$
in which\pageoriginale $i$ is etale and $i'$ is etale and of finite
presentation, and for every sheaf of sets (\resp of ind-$\LL$-groups,
\resp an abelian $\LL$-torsion sheaf) $F$ on $Y''$, the base-change
morphism \hyperrefext[XII][XII.eq4.1]{4.1}, \hyperrefext[XII][XII.eq4.2]{4.2}
$$
{g'}^{\ast}i'_{\ast}F\longrightarrow {j'}_{\ast}{g''}^{\ast}F
$$
is bijective (\resp the morphisms
$$
{g'}^{\ast}(R^{q}i'_{\ast})F\longrightarrow
(R^{q}j'_{\ast}){g''}^{\ast}F
$$
are bijective for $q=0$, $1$, \resp are bijective for $0\leq q\leq
n$ and injective for $q=n+1$).

\item The same statement as {\rm(ii)}, except that the morphism $i'$ is
assumed only quasi-finite and quasi-separated.

\item The same statement as {\rm(ii)}, but with $i$ assumed only
locally quasi-finite and, on the other hand, $i'$ a quasi-compact open
immersion.
\end{enumerate}
\end{enonce*}

\noindent
\emph{Proof}. The implication (iii) $\Rightarrow$ (ii) is
trivial. Let us verify that (i) $\Rightarrow$ (iii). It is enough to check
the morphism in (iii) fiber by fiber. Let $\bar{y}$ be a geometric point
of $Y'$ and $\bar{x}$ a geometric point of $X'$ lying over
$\bar{y}$. Let
$$
\xymatrix@=1.75cm{
\widetilde{X}'   \ar[d]^{\widetilde{g}'}& \widetilde{X}''
\ar[l]_{\widetilde{j}'} \ar[d]^{\widetilde{g}''}\\
 \widetilde{Y}'           & \widetilde{Y}'' \ar[l]_{\widetilde{i}'}
}
$$
be the Cartesian diagram deduced from $({\ast})$, where
$\widetilde{X}'$ and $\widetilde{Y}'$ are the strict localizations of
$X'$ and $Y'$ at $\bar{x}$ and $\bar{y}$, respectively, and where
$\widetilde{Y}''=\fprod{Y''}{\widetilde{Y}'}{Y'}$,
$\widetilde{X}''=\fprod{X''}{\widetilde{X}'}{X'}$.\pageoriginale By
\hyperrefext[VIII][VIII.thm5.2]{5.2} and \hyperrefext[VIII][VIII.rem5.3]{5.3}, we have
$$
H^{q}(\widetilde{Y}'',\widetilde{F})\simeq
(R^{q}i'{}_{\ast}F)_{\bar{y}}\simeq
({g'}^{\ast}(R^{q}i'_{\ast})F)_{\bar{x}}{\quoi},
$$
where $\widetilde{F}$ is the inverse image of $F$ on
$\widetilde{Y}''$, and
$$
H^{q}(\widetilde{X}'',{\widetilde{g}''}{}^{\ast}\widetilde{F})\simeq
((R^{q}j'_{\ast}){g''}{}^{\ast}F)_{\bar{x}}{\quoi}.
$$
Since $\widetilde{i}'$ is quasi-finite, we may apply
\hyperref[XV.prop1.6]{1.6} (iii) (\resp \hyperref[XV.cor1.5]{1.5} (iii)) to the morphism
$$
H^{q}(\widetilde{Y}'',\widetilde{F})\longrightarrow
H^{q}(\widetilde{X}'',{\widetilde{g}''}{}^{\ast}\widetilde{F})
$$
for the values of $q$ under consideration, which proves (i)
$\Rightarrow$ (iii).

(ii) $\Rightarrow$ (i). Assume (ii) and let us verify the acyclicity
criterion \hyperref[XV.prop1.6]{1.6} (ii) for a morphism
$\widetilde{g}:\widetilde{X}\rightarrow \widetilde{Y}$. Let
$\widetilde{Y}'\rightarrow \widetilde{Y}$ be etale, with
$\widetilde{Y}'$ affine, and let $\widetilde{F}$ be a sheaf of sets
(\resp an abelian $\LL$-torsion sheaf, \resp a sheaf of
ind-$\LL$-groups) on $\widetilde{Y}'$. We must prove that
$$
H^{q}(\widetilde{Y}',\widetilde{F})\xrightarrow{\sim}
H^{q}(\widetilde{X}', \widetilde{g}'{}^{\ast}\widetilde{F}){\quoi},
$$
where $\widetilde{g}':\widetilde{X}'\rightarrow \widetilde{Y}'$ is the
morphism obtained from $\widetilde{g}$ by base change. Now write
$\widetilde{Y}=\varprojlim Y_{\alpha}$, where the $Y_{\alpha}$ are
affine and etale over $Y$ (assuming $Y$ affine, as we may). We can
``descend'' the etale morphism
$\widetilde{Y}'\rightarrow \widetilde{Y}$ to one of the $Y_{\alpha}$,
say to $Y'_{0}\rightarrow Y_{0}$, and we have
$\widetilde{Y}'=\varprojlim
Y'_{\alpha}=\fprod{Y}{Y'_{0}}{Y_{0}}$. Every sheaf
$\widetilde{F}$ on $\widetilde{Y}'$ is the direct limit of the sheaves
${f'_{\alpha}}^{\ast}{f'_{\alpha}}_{\ast}\widetilde{F}$, where
$f'_{\alpha}:\widetilde{Y}'\rightarrow Y'_{\alpha}$ is the canonical
morphism. By \sisi{XII 1.2 (v) and 1.6 (iii)}{\hyperrefext[IX][IX.prop1.2]{1.2} and \hyperrefext[IX][IX.prop1.6]{1.6}}, the sheaves
$f'_{\alpha}{}_{\ast}\widetilde{F}$ are sheaves of sets (\resp abelian
$\LL$-torsion sheaves, \resp sheaves of
ind-$\LL$-groups). In view\pageoriginale of \hyperrefext[VII][VII.prop3.3]{3.3},
it is therefore enough to treat the case where $\widetilde{F}$ descends
to one of the $Y'_{\alpha}$, say to a sheaf $F_{0}$ on $Y'_{0}$. Applying
(ii) to the diagram
$$
\xymatrix@C=2.25cm@R=1.5cm{
X \ar[d]_{g}& X_0 \ar[l] \ar[d]_{g_0}& X'_0 \ar[l]_{j'_0} \ar[d]_{g'_0}\\
Y & Y_0 \ar[l]& Y'_0 \ar[l]_{i'_0},
}
$$
we deduce that
${g_{0}}^{\ast}(R^{q}{i'_0}_{\ast})F_{0}\xrightarrow{\sim}
(R^{q}{j'_0}_{\ast}){g'_0}^{\ast}F_{0}$ for the applicable
values of $q$. The geometric points $\bar{y}$ and $\bar{x}$ lift
canonically to geometric points of $Y_{0}$ and $X_{0}$ (which we denote
by the same symbols), and by \hyperrefext[VIII][VIII.thm5.2]{5.2} and
\hyperrefext[VIII][VIII.rem5.3]{5.3} we have
$$
H^{q}(\widetilde{Y}',\widetilde{F})=((R^{q}{i'_0}_{\ast})F_{0})_{\bar{y}}=({g_0}^{\ast}
(R^{q}{i'_0}_{\ast})F_0)_{\bar{x}}
$$
and
$$
H^{q}(\widetilde{X}',{\widetilde{g}'}{}^{\ast}\widetilde{F})=((R^{q}j'_{0\ast}){g'_0}^{\ast}F_{0})_{\bar{x}}
$$
for the values of $q$ under consideration, which proves the result.

It remains to show that conditions (i) through (iii) are equivalent to
(iv). Using form (i), we see that these conditions are stable under a
locally quasi-finite base change $Y'\rightarrow Y$; in form (ii), this
immediately shows that they imply (iv). Conversely, (iv) implies (iii),
since, in order to verify (iii), it follows at once---using, for example,
the Leray spectral sequence or the descent result
\hyperrefext[XII][XII.prop6.8]{6.8}---that we may assume $Y'$ and $Y''$ affine. But then, by
the ``Main Theorem'' \EGA IV 8.12.8, $Y''\rightarrow Y'$ factors as
$Y''\xrightarrow{i'_1}Y'_1\xrightarrow{i'_2}Y'$, with $i'_1$ an open
immersion\pageoriginale and $i'_2$ a \emph{finite} morphism. Using
\hyperrefext[VIII][VIII.prop5.5]{5.5}, \hyperrefext[VIII][VIII.rem5.8]{5.8} for $i'_2$, we are reduced to
verifying (ii) for $i'_1$, which follows from (iv). This completes the
proof of \hyperref[XV.prop1.10]{1.10}.

\begin{enonce*}{Definition 1.11}\etiquette[1.11]{XV.def1.11}
A morphism $g:X\rightarrow Y$ of preschemes satisfying any one of the
equivalent conditions {\rm(i)} through {\rm(iv)} of $\hyperref[XV.prop1.10]{1.10}$ is called a
\emph{locally $0$-acyclic} morphism (\resp \emph{locally
$1$-aspherical} for $\LL$, \resp \emph{locally
$n$-acyclic} for $\LL$). In the obvious way (compare
\hyperref[XV.def1.7]{1.7}), one defines the notions of a \emph{locally acyclic} morphism
and of a \emph{universally locally $0$-acyclic} morphism
(\resp\ldots) for $\LL$.
\end{enonce*}

The statements (i) through (iii) of \hyperref[XV.prop1.10]{1.10} can of course be varied.
Thus, in (iii) and (iv), it is enough to restrict to a constant sheaf $F$
of the form $G_{I}$,
% what is this notation?
with $G$ a finite set (\resp a finite $\LL$-group, \resp $G$ of the
form $\ZZ/\ell^{\nu}\ZZ$, with $\ell\in\LL$, $\nu\in \NN$); this
follows readily from \hyperrefext[VIII][VIII.thm5.2]{5.2},
\hyperrefext[VIII][VIII.rem5.3]{5.3}, and \hyperrefext[XII][XII.prop6.5]{6.5}. We also record the following
variant:

\begin{enonce*}{Corollary 1.12}\etiquette[1.12]{XV.cor1.12}
Suppose that $g:X\rightarrow Y$ is locally of finite type. To verify
local $0$-acyclicity (\resp\ldots), it is enough to verify $\hyperref[XV.prop1.10]{1.10}$
{\rm(i)} in the case where the geometric point $\bar{x}$ is
\emph{closed} in the geometric fiber $X_{\bar{y}}$.
\end{enonce*}

Indeed, in the proof of (i) $\Rightarrow$ (iii) of \hyperref[XV.prop1.10]{1.10}, we
verified (iii) fiber by fiber, and it was enough to do so for the fibers
at geometric points $\bar{x}$ that are closed in their geometric
fiber,\pageoriginale since $g$ is locally of finite type
\hyperrefext[VIII][VIII.rem3.13]{3.13} b).

\begin{enonce*}{Proposition 1.13}\etiquette[1.13]{XV.prop1.13}
\begin{enumerate}
\renewcommand{\theenumi}{\roman{enumi}}
\renewcommand{\labelenumi}{(\theenumi)}
\item Let $X\xrightarrow{g}Y\xrightarrow{h}Z$ be morphisms of schemes,
let $\LL\subset \PP$, and let $n\in \NN$. Suppose that $g$ is
surjective and locally $0$-acyclic (\resp locally $1$-aspherical for
$\LL$, \resp locally $n$-acyclic for $\LL$). Then $h$ is
locally $0$-acyclic (\resp\ldots) if and only if $f=hg$ is.

\item Let $\{g_{\alpha}:X_{\alpha}\rightarrow Y_{\alpha}\}$ be an
inverse system of morphisms of preschemes such that the transition
morphisms $X_{\alpha}\rightarrow X_{\beta}$ and
$Y_{\alpha}\rightarrow Y_{\beta}$ are affine. Let $g$ be the inverse
limit of the $\{g_{\alpha}\}$ \hyperrefext[VII][VII.subsec5.1]{5.1}.
If every $g_{\alpha}$ is locally $0$-acyclic (\resp\ldots), then so is
$g$.
\end{enumerate}
\end{enonce*}

This follows by a chain of reasoning analogous to \hyperref[XV.prop1.9]{1.9}.

\begin{enonce*}[remark]{Remarks 1.14}\etiquette[1.14]{XV.rems1.14}
To simplify his task, the author omitted from \hyperref[XV.prop1.10]{1.10} through
\hyperref[XV.prop1.13]{1.13} a treatment of the case $n=-1$, and leaves it to the
reader to verify that the statements, definitions, and proofs apply to
this case essentially without change. Thus, with the notation of
\hyperref[XV.prop1.10]{1.10}, condition (i) says that the induced local morphisms
$\widetilde{X}\rightarrow \widetilde{Y}$ are $(-1)$-acyclic, i.e., by
\hyperref[XV.cor1.5]{1.5}, \emph{surjective}, and conditions (ii) through (iv) say
that the base-change morphism
${g'}^{\ast}({i'}_{\ast}(F))\rightarrow {j'}_{\ast}({g''}^{\ast}(F))$
is a \emph{monomorphism}. Under these conditions, we shall therefore
say that $g$ is \emph{locally $(-1)$-acyclic}. This condition is thus
stable\pageoriginale under locally quasi-finite base change. As an
example, a flat morphism is locally $(-1)$-acyclic, since for such a
morphism the surjectivity of the maps
$\widetilde{X}\rightarrow \widetilde{Y}$ is a well-known consequence
of their flatness. We also note that, when $f$ is locally of finite
presentation, $f$ is locally $(-1)$-acyclic if and only if it is
\emph{universally open}, at least when $Y$ is locally Noetherian; this
follows readily from \EGA IV 14.4.9, 1.10.3. For such a morphism, local
$(-1)$-acyclicity consequently implies \emph{universal} local
$(-1)$-acyclicity. (It is moreover very likely\footnote[1]{This was in
fact verified by M. Raynaud. \Cf the new edition of \EGA IV 14!} that,
in these last statements, the hypothesis that $Y$ be locally
Noetherian is unnecessary; this is at any rate so if the set of
irreducible components of $Y$ is locally finite \EGA IV 14.4.10.)
\end{enonce*}

\begin{enonce*}{Theorem 1.15}\etiquette[1.15]{XV.thm1.15}
Let $g:X\rightarrow Y$ be a quasi-compact and quasi-separated morphism,
let $\LL\subset \PP$ be a set of prime numbers, and let $n$ be a
natural number (\resp $n=1$). Suppose that $g$ is $(n-1)$-acyclic for
$\LL$ and locally $(n-1)$-acyclic for $\LL$
$\hyperref[XV.def1.3]{1.3}$, $\hyperref[XV.def1.7]{1.7}$, $\hyperref[XV.def1.11]{1.11}$, $\hyperref[XV.rems1.14]{1.14}$. Let $F$ be a sheaf on $Y$, which, for
$n\geq 1$, is assumed to have the structure of an abelian
$\LL$-torsion sheaf (\resp the structure of a sheaf of groups), and
let $\xi$ be an element of $H^{n}(X,g^{\ast}(F))$. In order that
$\xi$ come from an element of $H^{n}(Y,F)$ (which is then uniquely
determined, thanks to the $(n-1)$-acyclicity hypothesis on $g$), it is
necessary and sufficient that, for every geometric point $\bar{y}$ of
$Y$ algebraic over a point $y$ of $Y$, the image $\xi_{\bar{y}}$ of
$\xi$ in $H^{n}(X_{\bar{y}},F|X_{\bar{y}})$ lie in the image of
$H^{n}(\bar{y},F_{\bar{y}})$ (which, for $n\geq 1$,
means\pageoriginale that $\xi_{\bar{y}}=0$).
\end{enonce*}

We note at once that \hyperref[XV.thm1.15]{1.15} shows that, under the preliminary
conditions imposed on $g$, the morphism $g$ is $n$-acyclic for
$\LL$ (\resp $1$-aspherical for $\LL$) if and only if, for
every $\bar{y}$ as above, $X_{\bar{y}}$ is $n$-acyclic for
$\LL$ (\resp is $1$-aspherical for $\LL$): the ``if'' direction
follows directly from \hyperref[XV.thm1.15]{1.15} and the definitions, while the ``only
if'' direction is in any case immediate from \hyperref[XV.cor1.9.1]{1.9.1}. With this
said, an immediate induction on $n$ gives the following result.

\begin{enonce*}{Corollary 1.16}\etiquette[1.16]{XV.cor1.16}
Let $g$ be a quasi-compact and quasi-separated morphism that is locally
$(n-1)$-acyclic for $\LL$, where $n$ is a given natural number
(\resp where $n=1$). Then $g$ is $n$-acyclic for $\LL$
(\resp is $1$-aspherical for $\LL$) if and only if, for every
geometric point $\bar{y}$ of $Y$ algebraic over a point $y$ of $Y$,
the fiber $X_{\bar{y}}$ is $n$-acyclic for $\LL$ (\resp
$1$-aspherical for $\LL$).
\end{enonce*}

\begin{enonce*}{Corollary 1.17}\etiquette[1.17]{XV.cor1.17}
Let $g:X\rightarrow Y$ be a morphism of preschemes, let
$\LL\subset \PP$ be a set of prime numbers, and let $n$ be a
natural number (\resp $n=1$). In order that $g$ be locally

% END INCLUDED SGA4_English_current_buildable_source_20260906/source/15/drafts/SGA4_XV_LINES_401_700_English.texfrag


% BEGIN INCLUDED SGA4_English_current_buildable_source_20260906/source/15/drafts/SGA4_XV_LINES_701_969_English.texfrag
$n$-acyclic for $\LL$ (\resp locally $1$-aspherical for
$\LL$), it is necessary and sufficient that, for every geometric
point $\bar{x}$ of $X$, if $\widetilde{g}:\widetilde{X}\rightarrow
\widetilde{Y}$ is the morphism between the strict localizations of $X$
and $Y$ (at $\bar{x}$ and the corresponding geometric point $\bar{y}$
of $Y$) induced by $g$, every geometric fiber $X_{\bar{z}}$ of $g$,
with respect to a geometric point $\bar{z}$ of $Y$ algebraic over a
point $z$ of $Y$, is $n$-acyclic for $\LL$ (\resp
$1$-aspherical for $\LL$).
\end{enonce*}

Necessity follows from criterion \hyperref[XV.prop1.10]{1.10} (i) and
\hyperref[XV.cor1.9.1]{1.9.1}. For\pageoriginale sufficiency, we proceed by induction
on $n$, which allows us, when $n\geq 1$, to suppose that $g$ is already
locally $(n-1)$-acyclic for $\LL$. The same then holds for the
morphisms $\widetilde{g}$, and \hyperref[XV.cor1.16]{1.16} implies that
$\widetilde{g}$ is $n$-acyclic (\resp $1$-aspherical for
$\LL$); by criterion \hyperref[XV.prop1.10]{1.10} (i), this implies that
$g$ is locally $n$-acyclic (\resp locally $1$-aspherical) for
$\LL$.

\begin{enonce*}{Corollary 1.18}\etiquette[1.18]{XV.cor1.18}
With the notation of \hyperref[XV.cor1.17]{1.17}, if $g$ is locally of
finite type, then in the stated criterion it is enough to consider
geometric points $\bar{x}$ whose underlying point $x$ is closed in its
fiber.
\end{enonce*}

This follows from the proof just given and from \hyperref[XV.cor1.12]{1.12}.

\noindent
\emph{Proof of} \hyperref[XV.thm1.15]{1.15}\footnote[1]{The reader
interested only in the Noetherian case is invited to consult
\sisi{Remark }{}\hyperref[XV.rem1.18]{1.18}.}.
% There are two statements numbered 1.18 in this exposé.
We first show that we may suppose $Y$ strictly local. In the abelian
case, with $n\geq 1$, the Leray spectral sequence for $g$ and
$g^{\ast}(F)$, together with
$R^{i}g_{\ast}(g^{\ast}(F))=0$ for $0<i\leq n-1$ and
$F\xrightarrow{\sim}g_{\ast}g^{\ast}(F)$, which express the
$(n-1)$-acyclicity hypothesis on $g$, yields a canonical exact
sequence
$$
0\longrightarrow H^{n}(Y,F)\longrightarrow
H^{n}(X,g^{\ast}(F))\longrightarrow
H^{0}(Y,R^{n}g_{\ast}(g^{\ast}(F))){\quoi},
$$
which shows that it is enough to prove that the image of $\xi$ in the
third term is zero. This can be checked fiber by fiber, and in view of
\hyperrefext[VIII][VIII.thm5.2]{5.2} we are reduced to the case where
$Y$ is strictly local. The case $n=0$, with $F$ a sheaf of sets, is
handled in the same way: the hypothesis on $g$ implies that
$F\rightarrow g_{\ast}g^{\ast}(F)$ is injective, and $\xi$ is
identified with a section of the second sheaf that must be shown to
come from a section of the first. This too can be checked
fiber by\pageoriginale fiber, and one applies
\hyperrefext[VIII][VIII.rem5.3]{5.3}. Essentially the same argument
works in the noncommutative case, with $n=1$, using the exact sequence
\hyperrefext[XII][XII.prop3.2]{3.2}
$$
H^{1}(Y,F)\longrightarrow H^{1}(X,g^{\ast}(F))\longrightarrow
H^{0}(Y,R^{1}g_{\ast}(g^{\ast}(F))){\quoi},
$$
and using \hyperrefext[VIII][VIII.rem5.3]{5.3}.

We therefore suppose that $Y$ is strictly local, and hence that $X$ is
quasi-compact and quasi-separated. The morphism $g:X\rightarrow Y$
satisfies the conditions of \hyperrefext[XII][XII.cor6.10]{6.10},
which reduces us to the case in which, for every closed subset $Y_1$
of $Y$ distinct from $Y$, the restriction of $\xi$ to
$X_1=\fprod{X}{Y_1}{Y}$ lies in the image of
$H^{n}(Y_1,F|Y_1)\rightarrow H^{n}(X_1,g^{\ast}(F)|X_1)$, and to
checking in this case that one can find a finite morphism
$f:Y'\rightarrow Y$ whose image in $Y$ contains a nonempty open subset
of $Y$, and such that the inverse image $\xi'$ of $\xi$ on
$X'=\fprod{X}{Y'}{Y}$ lies in the image of
$H^{n}(Y',F')\rightarrow H^{n}(X',{g'}^{\ast}(F'))$. We note that, in
view of \hyperrefext[VIII][VIII.prop5.5]{5.5} and
\hyperrefext[VIII][VIII.rem5.8]{5.8}, condition $2^{\circ}$) of
\loccit is automatically satisfied for this morphism
$f:Y'\rightarrow Y$.

To find such an $f$, we may suppose that $Y$ is reduced, and we
consider a maximal point $y$ of $Y$. Let $\bar{y}$ be the geometric
point corresponding to a separable closure of $k(y)$, so that
$\bar{y}$ is the filtered inverse limit of open subsets $U_i$ of
finite integral $Y$-schemes $Y_i$, each of which has an image in $Y$
containing a neighborhood of $y$. Using the hypothesis on $\xi$, we
see that there is an index $i$ such that the inverse image of $\xi$ on
$\fprod{X}{U_i}{Y}$ is zero if $n\geq 1$, \resp lies in the image of
$H^{0}(U_i,F_{U_i})$ if $n=0$.

We\pageoriginale take $Y'$ to be $Y_i$, and note that we may suppose
that $U'=U_i$ is the inverse image of an affine open neighborhood $U$
of $y$ in $Y$. We then see that the inverse image of $\xi$ on
$\fprod{X}{Z'}{Y}$ (where $Z'=Y'-U'$ is the inverse image in $Y'$ of
$Z=Y-U$) lies in the image of $H^{n}(Z',F_{Z'})$. After replacing $Y$
by $Y'$, we are therefore reduced to the following case: there is an
affine open subset $U$ of $Y$, with complement $Z$, such that the
restrictions of $\xi$ to $X_U$ and $X_Z$ lie respectively in the image
of $H^{n}(U,F_U)$ and in the image of $H^{n}(Z,F_Z)$; moreover, if
$n\geq 1$, these restrictions may be assumed to be \emph{zero}.

When $n\geq 1$, we shall show that under these conditions $\xi=0$. To
simplify notation, we write $Y'$ in place of $X$, denote by $U'$ and
$Z'$ the inverse images of $U$ and $Z$ in $Y'$, and denote the inverse
image $g^{\ast}(F)$ by $F'$. Let
$i:U\rightarrow Y$, $j:Z\rightarrow Y$,
$i':U'\rightarrow Y'$, and $j':Z'\rightarrow Y'$ be the inclusions,
$$
\xymatrix@C=2cm@R=1.5cm{
U'  \ar[r]^{i'} \ar[d]^{g_U}& Y' = X  \ar[d]^{g} & Z'  \ar[l]_{j'} \ar[d]^{g_Z}\\
U   \ar[r]^{i} & Y & Z \ar[l]_{j},
}
$$
and consider the exact sequence of sheaves
$$
0\longrightarrow F'_{U',Y'}\longrightarrow F'\longrightarrow
F'_{Z',Y'}\longrightarrow 0{\quoi}.
$$
% Reference for this notation?
Since the image of $\xi$ in
$H^{n}(Y',F'_{Z',Y'})=H^{n}(Z',F'_{Z'})$ is zero, $\xi$ comes from an
element of $H^{n}(Y',F'_{U',Y'})$, and the restriction of the latter to
$U'$ is equal to that of $\xi$, and is therefore zero. Since
$F'_{U',Y'}$ is isomorphic to the inverse image of $F_{U,Y}$, we are
thus reduced\pageoriginale to the case of a sheaf of the form
$F_{U,Y}$; hence we may suppose that $F_Z=0$ (which takes account of
the hypothesis $\xi|Z'=0$). Note that if $F\rightarrow G$ is a
monomorphism of abelian $\LL$-torsion sheaves (\resp of sheaves of
groups), it is enough to prove that the image $\eta$ of $\xi$ in
$H^{n}(Y',G')$ lies in the image of $H^{n}(Y,G)$, \ie is zero
\hyperrefext[XII][XII.cor6.6]{6.6}. This allows us to replace $F$ by
$i_{\ast}(F_U)$, since the hypothesis $F_Z=0$ implies that the
canonical homomorphism $F\rightarrow i_{\ast}(F_U)$ is injective. Now,
since $g$ is locally $(n-1)$-acyclic for $\LL$, and a fortiori
$0$-acyclic for $\LL$ (because $n\geq 1$), the inverse image of
$i_{\ast}(F_U)$ on $Y'$ is identified with
$i'_{\ast}(F'_{U'})$. We are therefore reduced to showing that the
image $\eta$ of $\xi$ in $H^{n}(Y',i'_{\ast}(F'_{U'}))$ is zero. The
restriction of $\eta$ to $U'$ is zero, so it is enough to check that
the homomorphism
\begin{equation*}
H^{n}(Y',i'_{\ast}(F'_{U'}))\longrightarrow H^{n}(U',F'_{U'})\tag{$\ast$}
\end{equation*}
is a monomorphism. In the abelian case, we use the Leray spectral
sequence for $i'$ and $F'_{U'}$:
$$
E^{p,q}_2=H^p(Y',R^qi'_{\ast}(F'_{U'}))\Longrightarrow
H^{\ast}(U',F'_{U'}){\quoi},
$$
and the relations
\begin{equation*}
E^{p,q}_2=0\; \mbox{\rm for } 0<p\leq n-1\quav q\leq n-1{\quoi}.\tag{${\ast}{\ast}$}
\end{equation*}
To verify the latter, note that the local $(n-1)$-acyclicity
hypothesis on $g$ for $\LL$ implies that there are isomorphisms
$$
R^{q}i'_{\ast}(F'_{U'})\simeq g^{\ast}(R^{q}i_{\ast}(F_U))\;\mbox{\rm for } q\leq n-1{\quoi},
$$
and,\pageoriginale in view of the fact that $g$ is $(n-1)$-acyclic
for $\LL$, so that
$H^{p}(Y,G)\xrightarrow{\sim}H^{p}(Y',G')$ for $p\leq n-1$ and every
abelian $\LL$-torsion sheaf $G$ on $Y$, we obtain
$({\ast}{\ast})$, since $H^{p}(Y,G)=0$ for $p>0$. In the same way we
obtain
\begin{equation*}
E^{0,q}_2\simeq H^{0}(Y,R{}^qi_{\ast}(F_U))\;\mbox{\rm
for}\quad q\leq n-1{\quoi}.\tag{${\ast}{\ast}{\ast}$}
\end{equation*}
The Leray spectral sequence under consideration then gives rise to an
exact sequence that forms the second row of the following commutative
diagram:
$$
\xymatrix{
**[r]H^{0}(Y,R^{n-1} i_{\ast}(F_U))  \ar[r] \ar[d]^{\text{\rotatebox{90}{$\sim$}}}&
**[r]H^n(Y,i_{\ast} (F_U)) = 0 \ar[d]\\
**[r]H^{0}(Y', R^{n-1} i'_{\ast}(F'_{U'}))  \ar[r] &
**[r]H^{n}(Y',i'_{\ast}(F'_{U'}))  \ar[r] & **[r]H^{n}(U',F'_{U'}),
}
$$
where the first vertical arrow is bijective by
$({\ast}{\ast}{\ast})$. The diagram shows that $(\ast)$ is injective,
which completes the proof in this case. In the nonabelian case,
$n=1$, we use the same diagram with $n=1$:
$$
\xymatrix{
**[r]H^{0}(Y, i_{\ast}(F_U))  \ar[r] \ar[d]^{\text{\rotatebox{90}{$\sim$}}}&
**[r]H^1(Y,i_{\ast} (F_U)) = 0 \ar[d]\\
**[r]H^{0}(Y', i'_{\ast}(F'_{U'}))  \ar[r] &
**[r]H^{1}(Y',i'_{\ast}(F'_{U'}))  \ar[r] & **[r]H^{1}(U',F'_{U'}),&
}
$$
\hyperrefext[XII][XII.prop3.2]{3.2}, where the first vertical arrow is
again bijective by the hypothesis of local and global $0$-acyclicity
on $g$; once again we conclude that $(\ast)$ is injective in this
case.

Finally, consider the case $n=0$. By virtue of
\sisi{\qquad}{\qquad}, specifying a section $\xi$ of $F'$ amounts to
% Missing reference
specifying a section $\xi_{U'}$ of $F'_{U'}$ and a section
$\xi_{Z'}$ of $F'_{Z'}$ such that the image of $\xi_{Z'}$ under the
natural homomorphism
$$
F'_{Z'}\longrightarrow {j'}^{\ast}(i'_{\ast}(F'_{U'}))
$$
is\pageoriginale the section of the right-hand side induced by
$\xi_{U'}$. By hypothesis, there are sections $\eta_Z$ and $\eta_U$ of
$F_Z$ and $F_U$, respectively, that induce $\xi_{Z'}$ and
$\xi_{U'}$. It remains to see that they satisfy the analogous
compatibility condition, which says that they are the restrictions of
a single section $\eta$ of $F$. We must verify the equality of two
sections of $j^{\ast}(i_{\ast}(F_U))$. Since $Z'\rightarrow Z$ is
surjective ($g$ being surjective because it is $(-1)$-acyclic by
hypothesis), it is enough to check that the two corresponding sections
of
$$
g_Z{}^{\ast}(j^{\ast}(i_{\ast}(F_U)))
=j'{}^{\ast}(g^{\ast}(i_{\ast}(F_U)))
$$
are equal. Through the base-change homomorphism
\begin{equation*}
g^{\ast}(i_{\ast}(F_U))\longrightarrow
i'_{\ast}(g_U{}^{\ast}(F_U))=i'_{\ast}(F'_{U'}){\quoi},\tag{$\ast$}
\end{equation*}
the sheaf under consideration maps to
${j'}^{\ast}(i'_{\ast}(F'_{U'}))$, and this latter homomorphism is
injective because $(\ast)$ is injective, by the hypothesis that $g$ is
locally $(-1)$-acyclic. It therefore suffices to verify that the two
sections of ${j'}^{\ast}(i'_{\ast}(F'_{U'}))$ that are the images of
the preceding two sections are equal. But these are immediately seen
to be precisely the sections considered above, which are equal by
hypothesis. This completes the proof.

\begin{enonce*}{Remark 1.18}\etiquette[1.18]{XV.rem1.18}
There is a refinement of \hyperref[XV.thm1.15]{1.15} that gives a much
simpler proof of \hyperref[XV.thm1.15]{1.15} when $Y$ is assumed to be
locally Noetherian. In addition to the preliminary hypotheses on $g$
and $F$, suppose that there is a finite set of geometric points
$f_i:\bar{y}_i\rightarrow Y$ of $Y$ such that the morphism
$$
F\longrightarrow \prod_if_{i\ast}(f_i{}^{\ast}(F))
$$
is\pageoriginale injective (a hypothesis automatically satisfied if
$Y$ is Noetherian and $F$ is constructible
\hyperrefext[IX][IX.prop2.14]{2.14}). I claim that under these
conditions criterion \hyperref[XV.thm1.15]{1.15} remains valid when
one restricts to the geometric points $\bar{y}_i$ alone. Indeed, by
\hyperrefext[XII][XII.cor6.6]{6.6}, we may suppose that $F$ itself is
of the form $f_{\ast}(M_{\bar{y}})$, where $M$ is a set \resp a group
\resp an abelian $\LL$-torsion group, and
$f:\bar{y}\rightarrow Y$ is a geometric point. Let
$f':X'=X_{\bar{y}}\rightarrow X$ be the canonical morphism. Assuming,
for definiteness, that $n\geq 1$, the local hypothesis on $g$ implies
that
$$
g^{\ast}(f_{\ast}(M_{\bar{y}}))\simeq f'_{\ast}(M_{X'})\sisi{,\;}{\quav}
R^{i}f'_{\ast}(M_{X'})=0\quad\mbox{\rm for}\quad 1\leq i\leq
n-1{\quoi},
$$
the second relation following from the fact that its left-hand side is
isomorphic to $g^{\ast}(R^{i}f_{\ast}(M_{\bar{y}}))$, which is zero
since manifestly $R^{i}f_{\ast}(M_{\bar{y}})=0$ for $i\geq 1$. From
these relations we conclude that
$$
H^{n}(X,g^{\ast}(F))\simeq H^{n}(X,f'_{\ast}(M_{X'}))\longrightarrow
H^{n}(X',M_{X'})
$$
is injective, using the Leray spectral sequence for $f'$. This proves
that an element $\xi$ of $H^{n}(X,g^{\ast}(F))$ whose induced element
in $H^{n}(X_{\bar{y}},M_{X_{\bar{y}}})$ is zero is itself zero. To
deduce from this a simplified proof of \hyperref[XV.thm1.15]{1.15} in
the case where $Y$ is locally Noetherian, note that we may restrict to
the case where $Y$ is Noetherian; an easy passage to the limit, using
\hyperrefext[IX][IX.cor2.7.2]{2.7.2} and
\hyperrefext[VII][VII.sec5]{5}, then reduces us to the case where $F$
is also constructible, to which the statement just proved applies.
\end{enonce*}

% END INCLUDED SGA4_English_current_buildable_source_20260906/source/15/drafts/SGA4_XV_LINES_701_969_English.texfrag


% BEGIN INCLUDED SGA4_English_current_buildable_source_20260906/source/15/drafts/SGA4_XV_LINES_970_1136_English.texfrag
\section{Local acyclicity of a smooth morphism}\etiquette[2]{XV.sec2}\pageoriginale

The theorem is the following:

\begin{enonce*}{Theorem 2.1}\etiquette[2.1]{XV.thm2.1}
Let $g:X\rightarrow Y$ be a smooth morphism, and let $\LL\subset
\PP$ be the complement of the set of residual characteristics of
$X$. Then $g$ is locally acyclic and locally $1$-aspherical for
$\LL$.
\end{enonce*}

Let us first record the corollary:

\begin{enonce*}{Corollary 2.2}\etiquette[2.2]{XV.cor2.2}
Let $g:\EE^m_Y\rightarrow Y$ be the structural morphism of affine
$m$-space. Then $g$ is acyclic and $1$-aspherical for the set $\LL$
complementary to the set of residual characteristics of $Y$.
\end{enonce*}

Indeed, induction on $m$ and \hyperref[XV.prop1.9]{1.9} (i) reduce us
to the case $m=1$. By \hyperref[XV.thm2.1]{2.1}, $g$ is locally
$n$-acyclic and locally $1$-aspherical for $\LL$, so the \emph{local}
condition of \hyperref[XV.cor1.16]{1.16} is satisfied. It therefore
remains only to prove that the geometric fibers of $\EE^1_Y$ are
acyclic and $1$-aspherical for $\LL$; that is, that $\Spec k[t]$ is
acyclic and $1$-aspherical for $\LL$ when $k$ is a separably closed
field of characteristic $p\notin\LL$. The acyclicity assertion follows
from \hyperrefext[IX][IX.thm4.6]{4.6}, and the $1$-asphericity assertion
is well known. More generally, let $X'$ be a normal irreducible Galois
covering of $X=\PP^1_k$, unramified away from the point $\infty$ and
tamely ramified at that point. We prove that the degree $n$ of the
covering is $1$. The Hurwitz formula gives
$$
2(g'-1)=2n(g-1)+\deg \mathfrak{d}_{X'/X},
$$
where $g$ and $g'$ are the genera of $X$ and $X'$, and
$\mathfrak{d}_{X'/X}$ is the different divisor. We have $g=0$, while
tame ramification implies $\deg \mathfrak{d}_{X'/X}\leq n-1$; hence
$2(g'-1)\leq-(n+1)$, so $n\leq2(1-g')-1$. Since $g'\geq0$, it follows
that $n\leq1$, and therefore $n=1$, \cqfd

\medskip
\noindent
\emph{Proof of} \hyperref[XV.thm2.1]{2.1}. The assertion is local on
$X$ and $Y$ for the étale topology, so we may assume
\sisi{(}{}\SGA 1 II 1.1\sisi{)}{} that $g$ is the structural morphism
of affine $m$-space $\EE^m_Y$ and that $Y$ is affine. By the
transitivity of local acyclicity, \hyperref[XV.prop1.13]{1.13} (i),
and the identity $\EE^m_Y=\EE^{m-1}_{\EE^1_Y}$, we are immediately
reduced to $m=1$, that is, to $X=\Spec\Oo_Y[t]$.

Apply \hyperref[XV.cor1.18]{1.18}. It remains to prove that the
geometric fibers of the morphism of strictly local schemes
$\widetilde g:\widetilde X\rightarrow\widetilde Y$ associated with a
pair of geometric points $\bar x,\bar y$, where $\bar x$ is closed in
the fiber $X_{\bar y}$, are acyclic for $\LL$ and $1$-aspherical for
$\LL$. We may plainly assume that $Y=\widetilde Y$. Let
$Y'\rightarrow Y$ be a finite, local, surjective morphism, so that
$Y'$ is strictly local, and let $\bar y'$ be the point of $Y'$ above
$\bar y$ and $\bar x'$ the point of $X'$ above both $\bar x$ and
$\bar y'$. Since $X'\rightarrow X$ is finite, the strict localization
of $X'$ at $\bar x'$ is precisely
$\widetilde X'=\fprod{\widetilde X}{Y'}{Y}$. Consequently every
geometric fiber of $\widetilde X'/Y'$ is obtained from a geometric
fiber of $\widetilde X/Y$ by an algebraic (necessarily radicial)
extension of the base field; moreover, at least one geometric fiber of
$\widetilde X'/Y'$ lies above each geometric fiber of
$\widetilde X/Y$. Thus, by \hyperrefext[VIII][VIII.thm1.1]{1.1}, we
may replace $Y$ by $Y'$. Choosing $Y'$ suitably reduces us to the case
$k(\bar x)=k(\bar y)$. Then $\bar x$ is a rational point of the fiber
$X_{\bar y}=\Spec k(\bar y)[t]$, say the point $t=a^0$. Lift $a^0$ to
a section $a$ of $\Oo_{\sisi{Y}{\widetilde Y}}$ and make the
``translation'' $t_1=t-a$; we are reduced to the case in which
$\bar x$ is the point $\{t=0\}$ of $\Spec k(\bar y)[t]$. Put
$A=\Gamma(\sisi{Y}{\widetilde Y},
\Oo_{\sisi{Y}{\widetilde Y}})$, so that $A$ is a strictly local
ring.

\begin{enonce*}[remark]{Notation 2.3}\etiquette[2.3]{XV.not2.3}
Let $A$ be a Henselian local ring. We denote by $A\{t\}$ the
Henselization of $A[t]$ at the prime ideal generated by $\rad A$ and
$t$.
\end{enonce*}

\begin{enonce*}{Corollary 2.4}\etiquette[2.4]{XV.cor2.4}
Let $A$ be Henselian and let $A'$ be a finite $A$-algebra. Then
$$
A'\{t\}\simeq A'\otimes_A A\{t\}{\quoi}.
$$
\end{enonce*}

This follows immediately from \hyperrefext[VIII][VIII.subsec4.1]{4.1}.
Note that if $A$ is strictly local, then so is $A\{t\}$.

The proof of the theorem is reduced to the following lemma:

\begin{enonce*}{Lemma 2.5}\etiquette[2.5]{XV.lem2.5}
Let $A$ be a Henselian ring of residual characteristic $p$. Then every
geometric fiber of $\Spec A\{t\}$ over $\Spec A$ is acyclic and
$1$-aspherical for $\LL=\PP-\{p\}$.
\end{enonce*}

Since $A\{t\}$ is the Henselization of $A[t]$, a fiber
$\Spec(A\{t\}\otimes_A K)$, where $K$ is a residue field of $A$, is
an inverse limit\pageoriginale of affine schemes $X_i$ of finite type
over the ``line'' $\Spec K[t]$, that is, a limit of affine algebraic
curves. Hence every geometric fiber of $\Spec A\{t\}$ is a limit of
affine curves $(X_i)_{\bar K}$ over the separably closed field
$\bar K$. It follows from \hyperrefext[IX][IX.cor5.7]{5.7} and
\hyperrefext[VII][VII.thm5.7]{5.7} that the cohomological dimension of
the geometric fibers is at most $1$. It is therefore enough to verify
that the geometric fibers are $1$-\emph{aspherical} for $\LL$; that
is, that such a geometric fiber $\bar Z$ is connected and nonempty and
that $H^1(\bar Z,G)=0$ for every finite $\LL$-group $G$.

Write $A=\varinjlim B_\alpha$, where the $B_\alpha$ are rings of
finite type over $\ZZ$. The morphism $B_\alpha\rightarrow A$ factors
through the Henselization $A_\alpha$ of $B_\alpha$ at the prime ideal
induced by $\rad A$, whence $A=\varinjlim A_\alpha$. One readily
checks that $A\{t\}=\varinjlim A_\alpha\{t\}$ as well. Let $\bar y$ be
a geometric point of $\Spec A$, let $X_{\bar y}$ be the geometric
fiber of $S=\Spec A\{t\}$ over $\Spec A$ at $\bar y$, let
$\bar y_\alpha$ be the point of $\Spec A_\alpha$ induced by $\bar y$,
and let $X_{\alpha,\bar y_\alpha}$ be the geometric fiber of
$\Spec A_\alpha\{t\}$ over $\Spec A_\alpha$ at $\bar y_\alpha$. We
have $X_{\bar y}=\varprojlim X_{\alpha,\bar y_\alpha}$. By
\hyperref[XV.prop1.9]{1.9} (ii), applied to the morphisms
$X_{\alpha,\bar y_\alpha}\rightarrow\bar y_\alpha$, the proof of
\hyperref[XV.lem2.5]{2.5} is therefore reduced to the case
$A=A_\alpha$, that is, to the case in which $A$ is the Henselization
of an algebra of finite type over $\ZZ$ at a given prime ideal. We
assume this from now on. Recall that such an algebra is
\emph{excellent} \sisi{(}{}\EGA IV 7.8.3 (iii), 7.8.6 (i)\sisi{)}{},
and hence that $A_\alpha$ is excellent as well (as follows readily
from the definition; \cf \EGA IV 18.7.6).

Let us verify $0$-acyclicity. (For a more general result, see the
Appendix.) Since $A\{t\}/A$ has the section ``$t=0$'', a geometric
fiber\pageoriginale $\bar Z$ is nonempty. To see that it is connected,
we reduce, by replacing $A$ by a quotient
\sisi{(}{}\hyperref[XV.cor2.4]{2.4}\sisi{)}{}, to the case in which
$A$ is integral and $\bar Z$ is the geometric fiber over the generic
point of $\Spec A$. It is enough to prove that for every finite
separable extension $K'$ of the fraction field $K$ of $A$, the scheme
$Z_{K'}$, where $Z=\Spec(A\{t\}\otimes_A K)$, is connected. Let $A'$
be the normalization of $A$ in $K'$ (a finite $A$-algebra,
\sisi{(}{}\EGA IV 7.8.2\sisi{)}{}) and let $Z'$ be the generic fiber
of $\Spec A'\{t\}$ over $A'$. By \hyperref[XV.cor2.4]{2.4},
$A'\{t\}\simeq A'\otimes_A A\{t\}$, and hence
$Z_{K'}\simeq Z'$. Replace $A$ by $A'$. It is then enough to prove
that $Z$ is integral, and for this it is enough to prove that
$A\{t\}$ is integral. Since $A$ is normal, $A[t]$ is normal, and so
is $A\{t\}$ \sisi{(}{}\SGA 1 I 9.5 (i)\sisi{)}{}. This proves the
assertion.

The proof of \hyperref[XV.lem2.5]{2.5}, and hence of
\hyperref[XV.thm2.1]{2.1}, is thus reduced to the following lemma,
which will be proved in the next section:

\begin{enonce*}{Lemma 2.6}\etiquette[2.6]{XV.lem2.6}
Let $A$ be a Henselian excellent ring, and let $\bar Z$ be a geometric
fiber of $\Spec A\{t\}/\Spec A$. Then every principal Galois covering
$\bar Z'$ of $\bar Z$ with group $G$ of order prime to the residual
characteristic $p$ of $A$ is trivial.
\end{enonce*}

% END INCLUDED SGA4_English_current_buildable_source_20260906/source/15/drafts/SGA4_XV_LINES_970_1136_English.texfrag


% BEGIN INCLUDED SGA4_English_current_buildable_source_20260906/source/15/drafts/SGA4_XV_LINES_1137_1315_English.texfrag
\section{Proof of the main lemma}\etiquette[3]{XV.sec3}\pageoriginale

Replacing $A$ by a finite $A$-algebra $A'$ (\cf
\hyperref[XV.cor2.4]{2.4}), assertion \hyperref[XV.lem2.6]{2.6} is reduced
to the case where $A$ is integral and $\bar{Z}$ is the geometric generic
fiber of $\Spec A\{t\}$ over $\Spec A$. Let $K$ be the fraction field of
$A$, and let $Z=\Spec A\{t\}\otimes_A K$. If $\bar{Z}'/\bar{Z}$ is a
principal Galois covering, it can be descended to some $Z_{K'}$, where
$K'$ is a suitable finite separable extension of $K$. Replacing $A$
(\hyperref[XV.cor2.4]{2.4}) by its normalization in $K'$, which is a
finite $A$-algebra \sisi{(}{}\EGA IV 7.8.3 (vi)\sisi{)}{}, we are reduced
to proving the lemma in the following form:

\begin{enonce*}{Lemma 3.1}\etiquette[3.1]{XV.lem3.1}
With the notation of \hyperref[XV.lem2.6]{2.6}, suppose that $A$ is
normal, and let $Z'/Z$ be a principal étale Galois covering, with group
$G$ of order prime to the residual characteristic $p$ of $A$. Then $Z'$
is induced by a Galois extension $K'$ of $K$; that is, \emph{(\SGA 1 IX
6.2)}, the covering $\bar{Z}'$ of $\bar{Z}$ deduced from $Z'$ is trivial.
\end{enonce*}

\noindent
\emph{Proof}. Let $B'$ be the normalization of $A\{t\}$ in the ring of
rational functions of $Z'$, and consider the height-$1$ prime ideals $P$
of $A\{t\}$ at which the $A\{t\}$-algebra $B'$ is ramified. Since $Z'/Z$
is étale, $P$ does not lie over the generic point of $\Spec A$; that is,
$\mathfrak{p}=P\cap A\neq 0$. Thus
$A\{t\}/\mathfrak{p}A\{t\}=(A/\mathfrak{p})\{t\}$ is integral and has
dimension at most $\dim A$. Since $A\{t\}$ is excellent
\sisi{(}{}\EGA IV 7.8.6 (i) and 18.7.6\sisi{)}{}, hence catenary, we have
$\dim \Spec A\{t\}/P=(\dim A+1)-1=\dim A\geq
\dim A\{t\}/\mathfrak{p}A\{t\}$, whence\pageoriginale
$P=\mathfrak{p}A\{t\}$, since $P\supset \mathfrak{p}A\{t\}$.

Let $\mathfrak{p}_1,\ldots,\mathfrak{p}_m$ therefore be the height-$1$
prime ideals of $A$ such that $B'/A\{t\}$ is ramified at the ideals
$\mathfrak{p}_iA\{t\}$, and let $s$ be the least common multiple of
the orders of the inertia groups \sisi{(}{}\SGA IV 2\sisi{)}{}
$G_i\subset G$ for $B'$ at the prime ideals lying above the
$\mathfrak{p}_iA\{t\}$. This is an integer dividing the order of $G$,
and is therefore prime to $p$. Let $f\in A$ be an element that vanishes
to order~$1$ at every $\mathfrak{p}_i$ (such an $f$ exists, as follows,
for example, from Bourbaki, Alg. Comm., Chap. II, \S 3, cor. to prop.
17). Replacing $A$ \sisi{(}{}\hyperref[XV.cor2.4]{2.4}\sisi{)}{} by its
normalization $A_1$ in the finite extension
$K_1=K[x]/(x^s-f)$ of $K$, and replacing $B'$ by the normalization of
$A_1\otimes_A B'$, Abhyankar's lemma \sisi{(}{}\SGA 1 X
3.6\sisi{)}{} reduces us to the case where $B'$ is ramified at no
height-$1$ prime ideal of $A\{t\}$.

The algebra $B'/A\{t\}$ is then also étale over the fiber of
$\Spec A\{t\}$ at every codimension-$1$ point of $\Spec A$. Indeed, let
$\mathfrak{p}$ be a height-$1$ prime ideal of $A$. Since $A$ is normal,
$A_{\mathfrak{p}}$ is a discrete valuation ring, hence regular, and so
$A\{t\}\otimes_A A_{\mathfrak{p}}$ is regular as well---indeed, it is a
limit of schemes étale over $A_{\mathfrak{p}}[t]$. By hypothesis, the
algebra $B'/A\{t\}$ is étale at every point of codimension $\leq 1$ of
$\Spec A\{t\}\otimes_A A_{\mathfrak{p}}$, hence it is étale everywhere
by the Zariski--Nagata purity theorem \sisi{(}{}\SGA 1 X
3.1\sisi{)}{}; for a proof, see for example \SGA 2, X 3.4.

Replacing $A$ \sisi{(}{}\hyperref[XV.cor2.4]{2.4}\sisi{)}{} by a
suitable finite $A$-algebra, we are reduced to the case where, in
addition, the algebra $B'/tB'$ over
$A\{t\}/tA\{t\}\simeq A$ is completely split over the generic point of
$\Spec A$. It therefore suffices to prove that, under these conditions,
the algebra $B'$ over $A\{t\}$\pageoriginale is ``trivial,'' and, since
$A\{t\}$ is Henselian, it suffices to prove that $B'$ is étale over
$A\{t\}$. Indeed, the fact that $B'/tB'$ splits over the generic point
of $\Spec A$ implies that it splits over $\Spec(A)$
\sisi{(}{}\SGA 1 I 10.1\sisi{)}{}; thus the residual extensions of $B'$
over the maximal ideal of $A\{t\}$ are trivial, and we apply \EGA IV
18.5.14. We are therefore reduced to the following lemma:

\begin{enonce*}{Lemma 3.2}\etiquette[3.2]{XV.lem3.2}
Let $B$ be an excellent normal local ring, and let $0\neq t\in \rad B$
be such that $B/tB$ is normal. Let $B'$ be a finite integral normal
$B$-algebra containing $B$, such that:
\begin{enumerate}
\renewcommand{\theenumi}{\roman{enumi}}
\renewcommand{\labelenumi}{(\theenumi)}
\item The algebra $B'/tB'$ over $B/tB$ is completely split at the
maximal points of $\Spec(B/tB)$.

\item The algebra $B'$ over $B$ is étale at the codimension-$1$ points
of $\Spec(B/tB)$.
\end{enumerate}
Then $B'$ is étale over $B$.
\end{enonce*}

\noindent
\emph{Proof}. By (ii), the assertion is trivial for $\dim B\leq 2$,
that is, for $\dim B/tB\leq 1$. Suppose that $B'/B$ is not étale; let
$x$ be a maximal point of the set of points of $\Spec B/tB$ over which
$B'/B$ is not étale, and let $C$ be the local ring of $\Spec B$ at $x$.
Then $C$ is again an excellent normal ring (\EGA IV 7.8.2), we have
$0\neq t\in\rad C$, and (i) and (ii) hold for the $C$-algebra
$C'=C\otimes {}_BB'$. We may therefore suppose that $B=C$; that is,
$B'$ is étale over $B$ above every point of $\Spec B/tB$, except
possibly at the origin. Since $B/tB$ is normal\pageoriginale (in fact,
unibranchness would suffice), condition (i) implies that the algebra $B'/tB'$ is
completely split away from the origin. Moreover, we may suppose that
$\dim B\geq 3$. We are thus reduced to the following lemma:

\begin{enonce*}{Lemma 3.3}\etiquette[3.3]{XV.lem3.3}
Let $B$ be an excellent normal local ring, let $0\neq t\in \rad B$,
and suppose that $\dim B\geq 3$. Let $B'$ be a finite integral normal
$B$-algebra containing $B$, such that the $B/tB$-algebra $B'/tB'$ is
completely split away from the origin. Then $B'$ is étale over $B$.
\end{enonce*}

\noindent
\emph{Proof}. Let $X=\Spec B$, let $X_1=X-\{x_0\}$, where $x_0$ is the
closed point, and let $i:U\rightarrow X$ be the open immersion whose
image consists of the points $x\in X$ at which the algebra $B'/B$ is
étale. Then
$(X-U)\cap V(t)=\{x_0\}$ has dimension $0$. Since $V(t)$ is defined by
\emph{one} equation, it follows that $\dim(X-U)\leq 1$ (\EGA
$0_{\text{\rm IV}}$ 16.2.5); consequently (since $X$ is catenary),
$x\in X-U\Rightarrow \dim \mathcal{O}_{X,x}\geq \dim B-1\geq 2$; that
is, $\codim(X-U,X)\geq 2$.

Let $X'=\Spec(B')$, and let $U'$ be the inverse image of $U$ in $X'$.
Since $X$ is universally catenary, the relation
$\codim(X-U,X)\geq 2$ implies
$\codim(X'-U',X')\geq 2$ (\EGA IV 5.6.10). Since $X'$ is normal, we
therefore have
$$
B'=\Gamma(X',\Oo_{X'})\xrightarrow{\sim}
\Gamma(U',\Oo_{U'}){\quoi}.
$$
Let $\mathcal{B}'=\widetilde{B}'$ be the sheaf of $\Oo_X$-algebras that
defines $X'$ over $X$, and let $\mathcal{C}$ be its restriction to $U$.
The preceding formula is then equivalent to
$$
\mathcal{B}'\xrightarrow{\sim}i_{\ast}(\mathcal{C}){\quoi},
$$
and\pageoriginale \hyperref[XV.lem3.3]{3.3} is consequently reduced to
the following statement (in which $B'$ no longer occurs): Let $B$, $t$ be
as in \hyperref[XV.lem3.3]{3.3}, let $X=\Spec(B)$, let $Y=\Spec(B/tB)$, let $U$ be an open
neighborhood of $Y_1=Y-\{x_0\}$ in $X_1=X-\{x_0\}$, and let $U'$ be an
étale covering of $U$ whose restriction to $Y_1$ is completely split.
If $\mathcal{C}$ denotes the sheaf of $\Oo_U$-algebras defining $U'$,
then (with $i:U\rightarrow X$ denoting the canonical immersion)
$i_{\ast}(\mathcal{C})$ is a coherent \emph{étale} algebra on $X$ (or,
equivalently, in view of the relation $\prof \Oo_{X,x}\geq 2$ for
$x\in X-U$ noted above and \EGA IV 5.10.5, $U'$ extends to an
\emph{étale} covering $X'$ of $X$).

Let $\widehat{X}=\Spec\widehat{B}$, where $\widehat{B}$ is the
completion of $B$. Then $\widehat{B}$ is normal (\EGA IV 7.8.3 (v)) and
$\widehat{B}/B$ is faithfully flat. Let
$$
\xymatrix@C=1.5cm@R=1cm{
\hat{U}      \ar[r]^{\sisi{\hat{i}}{\hat{i}}} \ar[d]^{f_1}& \hat{X} \ar[d]^{f}\\
U            \ar[r]^{i} &  X
}
$$
be the Cartesian diagram deduced from $\widehat{X}\rightarrow X$.
Since $f$ is faithfully flat, the functor $i_{\ast}$ commutes with base
change by $f$ \sisi{(}{}\EGA IV 2.3.1\sisi{)}{}, and we are reduced
\sisi{(}{}\EGA IV 17.7.1\sisi{)}{} to proving that
$\hat{i}_{\ast}(\widehat{\mathcal{C}})$ is an étale algebra on
$\widehat{X}$. We may therefore replace $B$ by $\widehat{B}$; that is,
we may suppose that $B$ is \emph{complete}. In this case we shall prove
that $U'$ is a completely split covering of $U$ (which will plainly
imply that it extends to an étale covering of $X$, and will complete the
proof). This assertion is now an immediate consequence of ``local
Lefschetz theory'' \sisi{(}{}\SGA\pageoriginale 2 X 2.1 (i) and
2.5\sisi{)}{}, as is seen at once, taking account once again of the
relation $\prof \Oo_{X,x}\geq 2$ for $\dim \{x\}=1$.

\begin{enonce*}[remark]{Remark 3.4}\etiquette[3.4]{XV.rem3.4}
One can prove \hyperref[XV.lem3.2]{3.2} and
\hyperref[XV.lem3.3]{3.3} under more general conditions, by following
essentially the same proof: instead of assuming that $B$ is excellent
and normal, it is enough to suppose that the completion $\widehat{B}$
is normal (a condition stable under localization, by the results of
\EGA IV 7), or merely that the completion of $B$ for the $(tB)$-adic
topology is normal and is a quotient of a regular ring.
\end{enonce*}

% END INCLUDED SGA4_English_current_buildable_source_20260906/source/15/drafts/SGA4_XV_LINES_1137_1315_English.texfrag


% BEGIN INCLUDED SGA4_English_current_buildable_source_20260906/source/15/drafts/SGA4_XV_LINES_1316_1417_English.texfrag
\section{Appendix: A criterion for local $0$-acyclicity}\etiquette[4]{XV.sec4}\pageoriginale 

The results of this section will not be used later in the seminar.

\begin{enonce*}{Theorem 4.1}\etiquette[4.1]{XV.thm4.1}
Let $g:X\rightarrow Y$ be a morphism of schemes. Assume that $g$ is flat
with separable fibers, and that $X$ and $Y$ are locally noetherian
\sisi{(}{}respectively, that $g$ is locally of finite presentation\sisi{)}{}.
Then $g$ is universally locally $0$-acyclic.
\end{enonce*}

We may clearly assume that $X$ and $Y$ are affine; in the second case,
the standard method of \EGA IV 8, using \EGA IV 9.7.7, then reduces us to
the case in which $Y$, and hence $X$, is noetherian. We may therefore
confine ourselves to proving the first assertion. Since the hypothesis
is stable under base change of finite type, the passage to the limit in
\hyperref[XV.prop1.13]{1.13}~(ii) reduces us to proving that $g$ is locally
$0$-acyclic. For this, \hyperref[XV.cor1.17]{1.17} reduces us to proving that,
if $g$ is moreover assumed to be strictly local, then the geometric
fibers of $g$ are connected. This is precisely \EGA IV 18.9.8.

\begin{enonce*}{Corollary 4.2}\etiquette[4.2]{XV.cor4.2}
Let $g:X\rightarrow Y$ be a morphism of schemes, with $Y$ discrete.
Then $g$ is universally locally $0$-acyclic.
\end{enonce*}

By \hyperrefext[VIII][VIII.thm1.1]{1.1}, we may assume that $Y$ is the
spectrum of a perfect field $k$. We may also assume that $X$ is affine,
and hence is a projective limit of affine schemes of finite type over
$k$; in view of \hyperref[XV.prop1.13]{1.13}~(ii), this reduces us to the
case in which $X$ is of finite type over~$k$. Finally,\pageoriginale we may
clearly assume that $X$ is reduced, and it is then separable over $k$
because $k$ is perfect. The conclusion now follows from
\hyperref[XV.thm4.1]{4.1}.

\begin{enonce*}{Corollary 4.3}\etiquette[4.3]{XV.cor4.3}
Assume that $g$ is surjective and satisfies the hypotheses of
\hyperref[XV.thm4.1]{4.1} or \hyperref[XV.cor4.2]{4.2}; in the
non-parenthetical case of \hyperref[XV.thm4.1]{4.1}, assume in addition
that $g$ is either quasi-compact and quasi-separated or universally
open. Then $g$ is a morphism of universal effective descent for the
fibered category of étale sheaves on varying preschemes.
\end{enonce*}

In the cases under consideration, $g$ is universally submersive, and it
follows from \hyperrefext[VIII][VIII.prop9.1]{9.1} that $g$ is a morphism
of universal descent for the fibered category in question. Thus, for
the question of effectivity, we may confine ourselves to the case in
which $Y$ is affine. Moreover, in each of the cases under consideration,
$X$ can be covered by a \emph{finite} number of affine open subsets
$X_i$ whose images cover $Y$. Replacing $X$ by the disjoint union of the
$X_i$, we may therefore assume that $X$ is affine, and hence a fortiori
quasi-compact and quasi-separated over $Y$. We may then conclude by the
argument of \hyperrefext[VIII][VIII.subsubsec9.4.1]{9.4.1}, using the fact
\hyperrefext[XVI][XVI.thm1.1]{1.1} that, for a quasi-compact and
quasi-separated morphism $g$, the formation of $g_{\ast}(F)$ (where $F$
is an étale sheaf on $X$) commutes with every base change
$Y'\rightarrow Y$ that is locally $0$-acyclic.

\begin{enonce*}[remark]{Remarks 4.4}\etiquette[4.4]{XV.rems4.4}
\begin{enumerate}
\renewcommand{\theenumi}{\alph{enumi}}
\renewcommand{\labelenumi}{\theenumi)}
\item This gives the proof, which had been left pending, of
\hyperrefext[VIII][VIII.thm9.4]{9.4}~d), as a special case of
\hyperref[XV.cor4.3]{4.3}. % label for d) ?
Another, easier proof, not using the rather delicate result
\pageoriginale of \EGA IV 18.9.8, is obtained by observing that, by
\hyperrefext[VIII][VIII.prop9.1]{9.1}, we may confine ourselves to the
case of the morphism $\Spec(K)\rightarrow \Spec(k)$ induced by a field
extension $K/k$. But $\Spec(K)$ is universally locally $0$-acyclic over
$\Spec(k)$, since it is even universally acyclic for
$\LL=\PP-\{p\}$, $p=\car k$, as we shall see in
\hyperrefext[XVI][XVI.cor1.5]{1.5}.

\item It is not known whether every quasi-compact and quasi-separated,
surjective, locally $(-1)$-acyclic (or universally locally
$(-1)$-acyclic) morphism is a morphism of effective descent for the
fibered category of étale sheaves on varying preschemes. This statement
seems quite plausible and would strengthen \hyperref[XV.cor4.3]{4.3} and
\hyperrefext[VIII][VIII.thm9.4]{9.4}~c), d). % labels for c) and d) ?
Compare the effectivity statement in
\textsc{Murre}, Sém. Bourbaki ${\rm n^{\circ}}$~293, p.~17.

\item It is plausible that, under the hypotheses of
\hyperref[XV.cor4.2]{4.2}, $g$ is even locally acyclic for $\LL$ and
locally $1$-aspherical for $\LL$, where $\LL$ is the set of prime numbers
distinct from the residue characteristics of $Y$. For this question we
may clearly assume that $Y$ is the spectrum of an algebraically closed
field $k$, and it can be shown (\SGA 1964/65) that the answer is
affirmative whenever resolution of singularities is available for
schemes of finite type over $k$. Thus the answer is affirmative when
$Y$ has characteristic zero, as follows by using the results of
\textsc{Hironaka}~\footnote[1]{\Cf \SGA 1 XIII.}.
\end{enumerate}
\end{enonce*}


%XII 4.1, 4.2 equation ref.
%XII 1.2, 1.6 no ref. at 501

\ifx\danslelivre\undefined
\end{document}
\fi

% END INCLUDED SGA4_English_current_buildable_source_20260906/source/15/drafts/SGA4_XV_LINES_1316_1417_English.texfrag


% END INCLUDED SGA4_English_current_buildable_source_20260906/source/15/15_English.tex


\makeatletter
\def\input@path{{../16/}}
\makeatother
% BEGIN INCLUDED SGA4_English_current_buildable_source_20260906/source/16/16_English.tex
\ifx\danslelivre\undefined
% Release engine: XeLaTeX (the LuaTeX compatibility layer breaks Xy-pic output).
\documentclass[10pt]{../fichiers-configuration-latex/smfart}
% BEGIN INCLUDED SGA4_English_current_buildable_source_20260906/source/fichiers-configuration-latex/paquets.tex
\usepackage{polyglossia}
%\usepackage[francais]{babel}
\usepackage{fontspec}
\usepackage[math-style=ISO,bold-style=upright,sans-style=upright]{unicode-math}
%\usepackage[mathletters]{ucs}
%\usepackage[utf8]{inputenc}
%\usepackage[T1]{fontenc}
\usepackage[unicode,hypertexnames=false,colorlinks=true,linkcolor=blue,citecolor=Sepia,urlcolor=magenta]{hyperref}
\setmainlanguage{french}
\setmainfont[Ligatures=TeX]{Linux Libertine O}
% variante possible : TeXGyrePagella ?
\setmathfont{XITS Math}

% maths
\usepackage{amsmath}
%\usepackage{amssymb}
%\usepackage{mathrsfs}
%\usepackage{stmaryrd}
%\usepackage{euscript}

% couleurs
\usepackage[usenames,dvipsnames]{xcolor}

% graphiques
\usepackage{graphicx}
\usepackage[all]{xy}
%\usepackage{tikz}
%\usetikzlibrary{matrix,arrows,calc}
%\tikzset{column sep=3em,row sep=2em,text height=1.5ex, text depth=0.25ex,>=angle 90}
%\tikzset{description/.style={fill=white,inner sep=1pt}}

% hyperliens, 'etiquettes
\usepackage{xr-hyper}
\usepackage{../fichiers-configuration-latex/smfthm}
%\ifx\cacheretiquettes\undefined
%\usepackage[inline]{showlabels} \else \fi % pour rendre visibles les 'etiquettes
%\usepackage{../fichiers-configuration-latex/makeidx} % pour faire un index
%\usepackage[pagebackref,unicode]{hyperref}  % pour voir o`u sont cit'ees les r'ef'erences
% hyperref is loaded once above with the cumulative-reader destination policy.
%\hypersetup{breaklinks=true} %permet le retour `a la ligne dans les liens trop longs. Incompatible avec pdflatex
%\renewcommand*{\backref}[1]{{\small [cit\'e page(s)~#1.]}}

\usepackage{xspace}
\usepackage{pifont} % pour expose XVII

% END INCLUDED SGA4_English_current_buildable_source_20260906/source/fichiers-configuration-latex/paquets.tex

% BEGIN INCLUDED SGA4_English_current_buildable_source_20260906/source/fichiers-configuration-latex/numerotation.tex
\renewcommand{\thefootnote}{\roman{footnote}} % pour mettre les notes en bas de page sous la forme (ii) et non (2), qui ressemble trop (?) a un exposant.

\newcounter{pageoriginal}
\newcommand{\pageoriginale}{\refstepcounter{pageoriginal}\marginpar{\footnotesize\textbf{\thepageoriginal}}\xspace}
\let\pageoriginaled\pageoriginale

\usepackage{ifthen}
\newboolean{original}
\def\sisi{\ifthenelse{\boolean{original}}}
\setboolean{original}{false}

\newcounter{cnote}
\def\nde#1{\ifthenelse{\boolean{original}}{}{\stepcounter{cnote}\def\thefootnote{\arabic{cnote}}\protect\footnote[\value{cnote}]{N.D.E.: #1}}}
\def\ndemark{\ifthenelse{\boolean{original}}{}{\stepcounter{cnote}\def\thefootnote{\arabic{cnote}}\protect\footnotemark[\value{cnote}]}}
\def\ndetext#1{\ifthenelse{\boolean{original}}{}{\def\thefootnote{\arabic{cnote}}\protect\footnotetext[\value{cnote}]{N.D.E.: #1}}}

% END INCLUDED SGA4_English_current_buildable_source_20260906/source/fichiers-configuration-latex/numerotation.tex

% BEGIN INCLUDED SGA4_English_current_buildable_source_20260906/source/fichiers-configuration-latex/macros.tex
\setmainlanguage{french}
\setotherlanguage{english}


%% Fixed, engine-independent release identifier.  The inherited dynamic Git
%% probe used Lua io.popen, exposed build-state-dependent output, and printed
%% a visible question mark outside a Git checkout.
\def\showgitstatus{English release r2, 2026-07-23}


% macros pour LaTeX
\renewcommand{\mathsf}{}

% Police différente pour U+2113 SCRIPT SMALL L
\setmathfont[range="2113]{latinmodern-math.otf}

% Give every inherited semantic label its own stable PDF destination instead
% of inheriting whichever automatic equation/page anchor happened to be live.
% The legacy label remains as a compatibility alias used by existing edges.
\makeatletter
\def\etiquette[#1]#2{%
  \Hy@raisedlink{\hyper@anchorstart{sga4.semantic.#2}\hyper@anchorend}%
  \def\@currentHref{sga4.semantic.#2}%
  \def\@currentlabel{#1}%
  \label{#2}%
}
\makeatother
\def\eqhyperref[#1]#2{(\hyperref[#1]{#2})}
\def\pagehyperref[#1]#2{(page \hyperref[#1]{#2})}
\def\hyperrefext[#1][#2]#3{\hyperref[#2]{#1 #3}}
\def\eqhyperrefext[#1][#2]#3{\hyperref[#2]{#1 #3}}     %cf. exposes 17
\def\neweqhyperrefext[#1][#2]#3{\hyperref[#2]{#1 #3}}  %cf. exposes 17

\def\shat#1{\protect\widehat{#1}}
\def\stilde#1{\protect\widetilde{#1}}
%%%%%%%%%%%%%%%%%%


\def\etcn{etc.\xspace}
\def\etco{etc...\xspace}
\def\etc{\sisi{\etco}{\etcn}}
\def\ldotsn{\sisi{...}{\etcn}}

\def\cf{{cf.}\kern.3em}
\def\Cf{{Cf.}\kern.3em}
\def\ie{{i.e.}\kern.3em}
\def\iev{{i.e.},\xspace}
\def\loccito{loc.\kern2pt cit.\xspace}
\def\loccitn{\textit{loc.\kern2pt cit.}\xspace}
\def\loccit{\sisi{\loccito}{\loccitn}}
\def\EGA{\textup{EGA}\kern.2em}
\def\SGA{\textup{SGA}\kern.3em}

\def\quoi{}

\def\quav{\quad,\quad}
\def\quaet{\quad \text{et}\quad}
\def\fibreproduct#1#2#3{#1 ×_{#3} #2}
\let\fprod\fibreproduct
\def\ccheck#1#2{\mbox{${\check{\check{\rm #1}}}^{\lower2pt\hbox{$\scriptstyle #2$}}$}}
\def\sheaf#1{\protect\mathcal{#1}}
\def\vieuxsheaf#1{{#1}} % cf. problème avec sheaf{a\Phi_i} par exemple...
\DeclareMathOperator{\Hom}{Hom}
\DeclareMathOperator{\can}{can}
\DeclareMathOperator{\Ind}{Ind}
\DeclareMathOperator{\ind}{ind}
%\DeclareMathOperator{\Hhomtop}{\mathcal{H}\!\mathit{omtop}}
\DeclareMathOperator{\Hhomtop}{\mathbf{Homtop}}
\DeclareMathOperator{\Point}{Point}
\DeclareMathOperator{\Pro}{Pro}
\DeclareMathOperator{\Top}{Top}
\DeclareMathOperator{\Mod}{Mod}
\DeclareMathOperator{\Cart}{Cart}
%\DeclareMathOperator{\Mmorsite}{\mathcal{M}\!\mathit{orsite}}
\DeclareMathOperator{\Mmorsite}{\mathbf{Morsite}}
\DeclareMathOperator{\loc}{loc}
\DeclareMathOperator{\Loc}{Loc}
\DeclareMathOperator{\Pic}{Pic}
\DeclareMathOperator{\Et}{Et}
\DeclareMathOperator{\card}{card}
\DeclareMathOperator{\Ens}{Ens}
\DeclareMathOperator{\h}{\underline{H}}
\DeclareMathOperator{\Fonct}{Fonct}
\DeclareMathOperator{\Ker}{Ker}
\DeclareMathOperator{\but}{but}
\DeclareMathOperator{\Ob}{Ob}
\DeclareMathOperator{\ob}{ob}
\DeclareMathOperator{\Lib}{Lib}
\DeclareMathOperator{\esj}{esj}
\DeclareMathOperator{\id}{id}
\DeclareMathOperator{\Id}{Id}
\DeclareMathOperator{\ad}{ad}
\DeclareMathOperator{\Iim}{Im}
\DeclareMathOperator{\pr}{pr}
\DeclareMathOperator{\Cat}{Cat}
\DeclareMathOperator{\car}{car}
\DeclareMathOperator{\SysInd}{SysInd}
\DeclareMathOperator{\Sup}{Sup}
\DeclareMathOperator{\Coker}{Coker}
\DeclareMathOperator{\Filt}{Filt}
\DeclareMathOperator{\acc}{acc}
\DeclareMathOperator{\sur}{sur}
\DeclareMathOperator{\cat}{cat}
\DeclareMathOperator{\Ord}{Ord}
\DeclareMathOperator{\Gr}{Gr}
\DeclareMathOperator{\Cov}{Cov}
\DeclareMathOperator{\Caf}{Caf}
\DeclareMathOperator{\coker}{coker}
\DeclareMathOperator{\coim}{coim}
\DeclareMathOperator{\im}{im}
\DeclareMathOperator{\Ouv}{Ouv}
\DeclareMathOperator{\diag}{diag}
\DeclareMathOperator{\Esp}{Esp}
\DeclareMathOperator{\TOP}{TOP}
\DeclareMathOperator{\Unif}{Unif}
\DeclareMathOperator{\Iind}{\mathbf{Ind}}
\DeclareMathOperator{\Hhomcart}{\mathbf{Homcart}}
\DeclareMathOperator{\Hhomtopess}{\mathbf{Homtopess}}
\DeclareMathOperator{\Ppoints}{\mathbf{Points}}
\DeclareMathOperator{\Ppoint}{\mathbf{Point}}
\DeclareMathOperator{\Ttors}{\mathbf{Tors}}
\DeclareMathOperator{\Ssex}{\mathbf{Sex}}
\DeclareMathOperator{\Ffib}{\mathbf{Fib}}
\DeclareMathOperator{\Vvois}{\mathbf{Vois}}
\DeclareMathOperator{\Ppointess}{\mathbf{Pointess}}
\DeclareMathOperator{\scrHom}{\mathbf{Hom}}
\DeclareMathOperator{\sheafPic}{\mathbf{Pic}}
\DeclareMathOperator{\sheafAut}{\mathbf{Aut}}
\DeclareMathOperator{\Hhomtoplocan}{\mathbf{Homtoplocan}}
\DeclareMathOperator{\Hhomtopan}{\mathbf{Homtopan}}
\DeclareMathOperator{\Pptgeom}{\mathbf{Ptgéom}}
\DeclareMathOperator{\Morsite}{Morsite}
\DeclareMathOperator{\Fib}{Fib}
\DeclareMathOperator{\Fais}{Fais}
\DeclareMathOperator{\kar}{kar}
\DeclareMathOperator{\Groupes}{Groupes}
\DeclareMathOperator{\Restr}{Restr}
\DeclareMathOperator{\Isom}{Isom}
\DeclareMathOperator{\Prol}{Prol}
\DeclareMathOperator{\Res}{Res}
\DeclareMathOperator{\Tors}{Tors}
\DeclareMathOperator{\esp}{esp}
\DeclareMathOperator{\es}{es}
\DeclareMathOperator{\od}{od}
\DeclareMathOperator{\gr}{gr}
\DeclareMathOperator{\sob}{sob}
\DeclareMathOperator{\dfn}{dfn}
\DeclareMathOperator{\Sch}{Sch}
\DeclareMathOperator{\ouv}{ouv}
\DeclareMathOperator{\cons}{cons}
\DeclareMathOperator{\Nerf}{Nerf}
\DeclareMathOperator{\Points}{Points}
\DeclareMathOperator{\Espsob}{Espsob}
\DeclareMathOperator{\Homtop}{Homtop}
\DeclareMathOperator{\Pt}{Pt}
\DeclareMathOperator{\Quot}{Quot}
\DeclareMathOperator{\glob}{glob}
\DeclareMathOperator{\Homfib}{Homfib}
\DeclareMathOperator{\Qouv}{Qouv}
\DeclareMathOperator{\Quotouv}{Quotouv}
\DeclareMathOperator{\ens}{ens}
\DeclareMathOperator{\ab}{ab}
\DeclareMathOperator{\an}{an}
\DeclareMathOperator{\fs}{fs}
\DeclareMathOperator{\red}{red}
\DeclareMathOperator{\Ab}{Ab}
\DeclareMathOperator{\cosq}{cosq}
\DeclareMathOperator{\Ann}{Ann}
\DeclareMathOperator{\Zar}{Zar}


%%%%%%%%%%%%
%%%%% SGA4-2
\DeclareMathOperator{\Sex}{Sex}
\DeclareMathOperator{\vois}{vois}
\DeclareMathOperator{\Vois}{Vois}
\DeclareMathOperator{\Ext}{Ext}
\DeclareMathOperator{\ext}{ext}
\DeclareMathOperator{\Fer}{Fer}
\DeclareMathOperator{\Add}{Add}
\DeclareMathOperator{\Cocart}{Cocart}
\DeclareMathOperator{\cosk}{cosk}
\DeclareMathOperator{\plat}{plat}

\DeclareMathOperator{\faisc}{faisc}
\DeclareMathOperator{\sk}{sk}
\DeclareMathOperator{\ch}{ch}

\def\HH{\mathbf{H}}
\def\qc{\mathrm{qc}}
\def\et{\mathrm{et}}
\def\coh{\mathrm{coh}}

\def\Hvr#1{%
\begin{array}[b]{l}
{\kern2pt\scriptstyle r}\\[-8pt]
{\kern1pt\scriptstyle \vee}\\[-5pt]
H^{#1}
\end{array}}

\def\Hvv#1{%
\begin{array}[b]{l}
\kern2pt\rotatebox{90}{${\scriptscriptstyle\ll}$}\\[-6pt]
H^{#1}
\end{array}}

\def\Hvi#1{%
\begin{array}[b]{l}
{\kern0.2pt\scriptscriptstyle \infty}\\[-8pt]
{\kern1pt\scriptstyle \vee}\\[-5pt]
H^{#1}
\end{array}}

\def\Hv#1{%
\begin{array}[b]{l}
{\kern1pt\scriptstyle \vee}\\[-5pt]
H^{#1}
\end{array}}

\DeclareMathOperator{\cart}{cart}
\DeclareMathOperator{\Ffaisc}{\mathbf{Faisc}}


\def\Ev#1{%
\begin{array}[b]{l}
{\kern1pt\scriptstyle \vee}\\[-5pt]
E'_{#1}
\end{array}}

\DeclareMathOperator{\Cont}{Cont}
\DeclareMathOperator{\Ccont}{\mathscr{C}\!\mathit{ont}}
\DeclareMathOperator{\Site}{Site}
\newcommand\MM{\mathbf{M}}
\def\limtop#1{{\displaystyle{\varprojlim_{#1}\!\mbox{\rm top }}}}
\def\Limtop{\mathop{\varprojlim\!\mbox{\rm top}}}
\DeclareMathOperator{\spec}{spec}
\def\ss{\scriptstyle{\mathscr{S}}}
\DeclareMathOperator{\Gal}{Gal}
\def\smallm{\scriptstyle{\mathscr{M}}}
\DeclareMathOperator{\End}{End}
\DeclareMathOperator{\ass}{ass}
\DeclareMathOperator{\zar}{zar}
\DeclareMathOperator{\Ogg}{Ogg}
\DeclareMathOperator{\res}{res}
\DeclareMathOperator{\tr}{tr}
\DeclareMathOperator{\Supp}{Supp}
\DeclareMathOperator{\codim}{codim}
\DeclareMathOperator{\rk}{rk}
\DeclareMathOperator{\cur}{cur}
\DeclareMathOperator{\cdqc}{cdqc}
\newcommand\DD{\mathbf{D}}
\DeclareMathOperator{\cl}{cl}
\def\JJ{\mathbf{J}}
\def\FF{\mathbf{F}}
\def\BB{\mathbf{B}}
\DeclareMathOperator{\Alget}{Alget}
\DeclareMathOperator{\Ass}{Ass}
\DeclareMathOperator{\supp}{supp}
\DeclareMathOperator{\rad}{rad}
\DeclareMathOperator{\prof}{prof}
\DeclareMathOperator{\cint}{int}
\DeclareMathOperator{\sic}{sic}
\DeclareMathOperator{\Tor}{Tor}
\DeclareMathOperator{\tors}{tors}
\DeclareMathOperator{\Sym}{Sym}
\DeclareMathOperator{\parf}{parf}
\DeclareMathOperator{\tor}{tor}
\DeclareMathOperator{\defn}{defn}
\DeclareMathOperator{\Div}{Div}
\DeclareMathOperator{\Pol}{Pol}
\DeclareMathOperator{\Fex}{Fex}
\DeclareMathOperator{\Lie}{Lie}
\DeclareMathOperator{\tot}{tot}
\DeclareMathOperator{\Aut}{Aut}
\DeclareMathOperator{\HOM}{HOM}
\DeclareMathOperator{\pch}{pch}
\DeclareMathOperator{\Ch}{Ch}
\DeclareMathOperator{\EXT}{EXT}
\DeclareMathOperator{\PIC}{PIC}
\DeclareMathOperator{\Rf}{Rf}
\DeclareMathOperator{\inj}{inj}
\newcommand\VV{\mathbf{V}}

\def\Xv{%
\begin{array}[b]{l}
{\kern1pt\scriptstyle \vee}\\[-5pt]
X
\end{array}}

%%%%%%%
\newcommand\TT{\mathbf{T}}
\newcommand\XX{\mathbf{X}}
\newcommand\EE{\mathbf{E}}
%%%%%%

\def\qedsymbol{\textup{C.Q.F.D.}}
\let\CQFD\qed
\let\cqfd\qed

%%%%%\bigdot
\font\bigsymb=cmsy10 at 4pt
\def\bigdot{{\kern1.2pt\raise 1.5pt\hbox{\bigsymb\char15}}}
%%%%%

\def\Lright{\mathop{\vcenter{\hbox{\hbox{$\longrightarrow$}\lower8pt \hbox{\kern-16pt\hbox{$\longrightarrow$}}}}}}

\def\Aa{{\mathcal A}}
\def\Bb{{\mathcal B}}
\def\Cc{{\mathcal C}}
\def\Dd{{\mathcal D}}
\def\Ee{{\mathcal E}}
\def\Ff{{\mathcal F}}
\def\Gg{{\mathcal G}}
\def\Hh{{\mathcal H}}
\def\Ii{{\mathcal I}}
\def\Jj{{\mathcal J}}
\def\Kk{{\mathcal K}}
\def\Ll{{\mathcal L}}
\def\Mm{{\mathcal M}}
\def\Nn{{\mathcal N}}
\def\Oo{{\mathcal O}}
\def\Pp{{\mathcal P}}
\def\Qq{{\mathcal Q}}
\def\Rr{{\mathcal R}}
\def\Ss{{\mathcal S}}
\def\Tt{{\mathcal T}}
\def\Uu{{\mathcal U}}
\def\Vv{{\mathcal V}}
\def\Ww{{\mathcal W}}
\def\Xx{{\mathcal X}}
\def\Yy{{\mathcal Y}}
\def\Zz{{\mathcal Z}}

\newcommand{\CC}{\mathbf{C}}
\newcommand{\GG}{\mathbf{G}}
\newcommand{\KK}{\mathbf{K}}
\newcommand{\LL}{\mathbf{L}}
\newcommand{\NN}{\mathbf{N}}
\newcommand{\PP}{\mathbf{P}}
\newcommand\QQ{\mathbf{Q}}
\newcommand\RR{\mathbf{R}}
\newcommand\ZZ{\mathbf{Z}}

\def\to{\mathchoice{\longrightarrow}{\rightarrow}{\rightarrow}{\rightarrow}}
\def\from{\mathchoice{\longleftarrow}{\leftarrow}{\leftarrow}{\leftarrow}}
\def\mto{\mathchoice{\longmapsto}{\mapsto}{\mapsto}{\mapsto}}
\def\mfrom{\mathrel{\from\joinrel\mapstochar\,}}
\def\hto{\mathrel{\lhook\joinrel\to}}
\def\hfrom{\mathrel{\from\joinrel\rhook}}

\def\implies{\mathchoice{\Longrightarrow}{\Rightarrow}{\Rightarrow}{\Rightarrow}}
\def\implied{\mathchoice{\Longleftarrow}{\Leftarrow}{\Leftarrow}{\Leftarrow}}
\def\iff{\mathchoice{\Longleftrightarrow}{\Leftrightarrow}{\Leftrightarrow}{\Leftrightarrow}}

\def\To#1{\mathchoice{\xrightarrow{\labelstyle~~#1~}}{\stackrel{#1}{\longrightarrow}}{}{}}
\def\Mto#1{\mathrel{\mapstochar\To{#1}}}
\def\From#1{\mathchoice{\xleftarrow{\labelstyle~~#1~}}{\stackrel{#1}{\longleftarrow}}{}{}}

\newcommand{\isomto}{\overset{\sim}{\longrightarrow}}
\newcommand{\isomfrom}{\overset{\sim}{\longleftarrow}}

\newcommand{\IFF}{\unskip~$\iff$~\ignorespaces}
\newcommand{\THEN}{\unskip~$\implies$~\ignorespaces}

\let\leq\leqslant
\let\geq\geqslant
\let\le\leqslant
\let\ge\geqslant

\def\cf{{cf.}\kern.3em}
\def\Cf{{Cf.}\kern.3em}
\def\ie{{i.e.}\kern.3em}
\def\iev{{i.e.},\xspace}
%\def\loccito{loc.\kern2pt cit.\xspace}
%\def\loccitn{\textit{loc.\kern2pt cit.}\xspace}
%\def\loccit{\sisi{\loccito}{\loccitn}}
\def\EGA{\textup{EGA}\kern.2em}
\def\SGA{\textup{SGA}\kern.3em}

\let\oldS\S
\def\S{\oldS\kern .15em}
\def\ptbl{.\kern .2em}
\def\resp{resp.\kern.3em}
\def\Exp{Exp.\kern.3em}

\let\bar\overline
\let\hat\widehat
\let\tilde\widetilde
\let\dpl\displaystyle
\let\ldots\dots
\let\emptyset\varnothing
\let\skpt\skippointrait
\let\epsilon\varepsilon

\def\sfrac#1#2{{#1}/{#2}}

\DeclareMathOperator{\cd}{cd}
%\DeclareMathOperator{\Hhom}{\mathcal{H}\!\mathit{om}}
\DeclareMathOperator{\Hhom}{\mathbf{Hom}}
%\DeclareMathOperator{\Eext}{\mathcal{E}\!\mathit{xt}}
\DeclareMathOperator{\Eext}{\mathbf{Ext}}
\DeclareMathOperator{\Spec}{Spec}
\DeclareMathOperator{\Tr}{Tr}

\let\emptyset=\varnothing

%%%%%%%%%%%%%%%
%%% expose XVII
%%%%%%%%%%%%%%%

\makeatletter
\let\old@optexpose\optexpose
\newcommand{\empty@optexpose}[1]{}
\newcommand{\paren@optexpose}[1]{\old@optexpose{#1}(}
\newcommand{\refnoexpose}[1]{\let\optexpose\empty@optexpose\ref{#1}\let\optexpose\old@optexpose}%\refnoexpose{X.cor5.2} engendre "5.2" au lieu de "X 5.2"
\newcommand{\neweqref}[1]{\let\optexpose\paren@optexpose\textup{\ref{#1})}\let\optexpose\old@optexpose}%\neweqref{XII.eq4.1} engendre "XII (4.1)" au lieu de "(XII 4.1)"
\makeatother
\newcommand{\refvb}[2][VI${}_{\text{B}}$ ]{\sisi{#1\refnoexpose{#2}}{\ref{#2}}}

\newcommand{\leftexp}[2]{{\vphantom{#2}}^{#1}{#2}}
\newcommand{\AAA}{\mathbf{A}}
\newcommand{\Sf}{\mathfrak{S}}
\DeclareMathOperator*{\quotprojlim}{\text{``$\varprojlim$''}}
\DeclareMathOperator*{\quotinjlim}{\text{``$\varinjlim$''}}
\DeclareMathOperator{\Fl}{F\sisi{\ell}{l}}
\DeclareMathOperator{\Pl}{P\sisi{\ell}{l}}
%\DeclareMathOperator{\cl}{c\sisi{\ell}{l}}
\DeclareMathOperator{\tdf}{tdf}
\DeclareMathOperator{\rg}{rg}
\newcommand{\iInd}{\sisi{I}{i}nd}
\newcommand{\iind}{ind}
\newcommand{\ppro}{pro}
\newcommand{\tTor}{\sisi{T}{t}or}
\newcommand{\ttor}{tor}
\newcommand{\chap}{Chap.}
\newcommand{\Ffex}{\Ff\mathit{ex}}
\newcommand{\sheafk}{\underline{k}}
\newcommand{\sheafn}{\underline{n}}
\newcommand{\sheafGamma}{\underline{\Gamma}}

% Repeated references to an already typeset footnote must not use
% \footnotemark without a matching \footnotetext: Hyperref would then create
% a live annotation to a destination that cannot exist.  This helper keeps the
% current printed footnote mark and links it to an explicit stable note target.
\makeatletter
\newcommand{\sgarepeatedfootnotemark}[1]{%
  \begingroup
  \edef\@thefnmark{\thefootnote}%
  \hyperlink{#1}{\@makefnmark}%
  \endgroup}
\makeatother

\sisi{
\DeclareMathOperator{\simarrow}{\rule{.15mm}{1.8mm}\kern -3.4pt\rightsquigarrow}
\newcommand{\sheafell}{\underline{\ell}}
\newcommand{\sheafZZ}{\underline{\ZZ}}
}{
\newcommand{\nderef}[1]{la N.D.E.~(\ref{#1}) page~\pageref{#1}}
}

%%%%%%%%%%%%%%%%%
%%%%% expose XVII
%%%%%%%%%%%%%%%%%

% END INCLUDED SGA4_English_current_buildable_source_20260906/source/fichiers-configuration-latex/macros.tex


\externaldocument{../01/01_English.tex}
\externaldocument{../02/02_English.tex}
\externaldocument{../03/03_English.tex}
\externaldocument{../04/04_English.tex}
\externaldocument{../05/05_English.tex}
\externaldocument{../05bis/05bis_English.tex}
\externaldocument{../06/06_English.tex}
\externaldocument{../07/07_English.tex}
\externaldocument{../08/08_English.tex}
\externaldocument{../09/09_English.tex}
\externaldocument{../10/10_English.tex}
\externaldocument{../11/11_English.tex}
\externaldocument{../12/12_English.tex}
\externaldocument{../13/13_English.tex}
\externaldocument{../14/14_English.tex}
\externaldocument{../15/15_English.tex}

\setboolean{original}{false}

\begin{document}

\begin{center}
XVI. Smooth base change theorem and applications\\
M. Artin
\end{center}

\selectlanguage{english}version: \showgitstatus
\tableofcontents

\else
\chapter{Smooth base change theorem and applications}
%\addtocontents{toc}{M. Artin}
\begin{center}
M. Artin
\end{center}
\fi
\phantomsection\MakeLinkTarget*{sga4.exp.XVI}\label{sga4.exp.XVI}

% BEGIN INCLUDED SGA4_English_current_buildable_source_20260906/source/16/drafts/SGA4_XVI_LINES_43_274_English.texfrag
\setcounter{pageoriginal}{206}

\section{The smooth base change theorem}\etiquette[1]{XVI.sec1}\pageoriginale

\begin{enonce*}{Theorem 1.1}\etiquette[1.1]{XVI.thm1.1}
Let $\LL\subset \PP$, let $n\in \NN$, and let
$$\vcenter{
\xymatrix@C=1.5cm@R=1cm{
X   \ar[d]^{f}& X' \ar[l]_{g'} \ar[d]^{f'}\\
Y     & \ar[l]_{g} Y'
}}\leqno{(*)}
$$
be a cartesian diagram, with $f$ quasi-compact and quasi-separated\nde{That is,
coherent in the terminology of the source.}, and $g$
universally locally $0$-acyclic (\resp
\sisi{univ. loc.}{universally locally} $1$-aspherical for $\LL$, \resp
\sisi{univ. loc.}{universally locally} $n$-acyclic for $\LL$) \emph{(\hyperrefext[XV][XV.def1.11]{1.11}).} Then,
for every sheaf $F$ of sets (\resp of
\emph{ind}-$\LL$-groups, \resp an abelian sheaf of
$\LL$-torsion) on $X$, the base-change morphism
\emph{(XII \hyperref[XII.eq4.2]{4.1.2})}
$$
\varphi^{q}:g^{\ast}(R^{q}f_{\ast})F\longrightarrow
(R^{q}f'_{\ast}){g'}^{\ast}F
$$
is bijective for $q=0$ (\resp $q\leq 1$, \resp $q\leq n$).
\end{enonce*}

In particular, applying \hyperrefext[XV][XV.thm2.1]{2.1}, we obtain the following.

\begin{enonce*}{Corollary 1.2}\etiquette[1.2]{XVI.cor1.2}
(Smooth base change theorem). Suppose that
  the morphism $g$ in $(\ast)$ is smooth, that $f$ is quasi-compact
  and quasi-separated, and that $\LL\subset \PP$ is
  the complement of the set of residue characteristics
  of $Y$. Then, for every sheaf $F$ of sets
  (\resp of \emph{ind}-$\LL$-groups, \resp an abelian sheaf of
  $\LL$-torsion) on $X$, the base-
  change\pageoriginale morphism $\varphi^{q}$ above is
  bijective for $q=0$ (\resp for $q=0$, $1$, \resp for every $q$).
\end{enonce*}

\noindent
\emph{Proof of} \hyperref[XVI.thm1.1]{1.1}. We first treat the case in which
$f:X\rightarrow Y$ is quasi-projective: this is simply a combination of the
proper base change theorem \hyperrefext[XII][XII.thm5.1]{5.1} and
Definition \hyperrefext[XV][XV.prop1.9]{1.9}. Indeed, we have a commutative diagram
$$
\xymatrix@=1.25cm{
 X \ar[rr]^{i} \ar[rd]_{f}&        & \bar{X}\ar[ld]^{\bar{f}}\\
                         &   Y    &
}
$$
where $\bar{f}$ is projective, hence proper, and $i$ is an open
immersion. We know from \hyperrefext[XII][XII.thm5.1]{5.1} that the base change theorem holds
for $\bar{f}$, whatever the morphism $g:Y'\rightarrow Y$. Since
$\fprod{g}{\bar{X}}{Y}:\fprod{Y'}{\bar{X}}{Y}\rightarrow \bar{X}$ is
again universally $0$-acyclic (\resp\ldots), it follows immediately
from (\hyperrefext[XII][XII.prop4.4]{4.4} (ii)) that it suffices to prove the theorem for the
morphism $i$; that is, we are reduced to the case in which $f$ is an
open immersion. The theorem then follows from Definition
\hyperrefext[XV][XV.prop1.10]{1.10} (ii).

We now treat the general case. The assertion is local on $Y$, and
we may therefore suppose that $Y$ is affine.

\begin{enonce*}{Lemma 1.3}\etiquette[1.3]{XVI.lem1.3}
We may moreover suppose that $X$ is affine.
\end{enonce*}

\noindent
\emph{Proof}. Arguing as in \hyperrefext[XII][XII.lem6.1]{6.1}, in order to prove
\hyperref[XVI.thm1.1]{1.1} for a given $f$, we are reduced to proving it in the situation
deduced\pageoriginale from the given situation by the base changes
$\widetilde{Y}\rightarrow Y$, $\widetilde{Y}'\rightarrow Y'$
for a local morphism $\widetilde{g}:\widetilde{Y}'\rightarrow
\widetilde{Y}$ between strict localizations of $Y'$ and $Y$ induced by
$g$, and to proving that the homomorphisms
\begin{equation*}
H^{q}(\widetilde{X},F)\longrightarrow H^{q}(\widetilde{X}',F')\quad
(\mbox{\rm where } F'={g'}^{\ast}(F))\tag{$\ast$}
\end{equation*}
are bijective for $q\leq n$, where $\widetilde{X}=\fprod{X}{\widetilde{Y}}{Y}$,
$\widetilde{X}'=\fprod{X'}{\widetilde{Y}'}{Y'}$. We are thus reduced to the case in which
$Y$, $Y'$, and $g$ are strictly local, and to proving in this case the
bijectivity of the preceding maps, assuming that
\hyperref[XVI.thm1.1]{1.1} has been proved
for affine $f$. Let $(X_i)$ be a finite covering of $X$ by affine open
subschemes $X_i$, and let $Z$ be the disjoint union of the $X_i$, which is
therefore affine and is equipped with a surjective morphism $h:Z\rightarrow X$, whence
a morphism $h':Z'\rightarrow X'$. Note that $h$ is separated, and
that it is affine when $f$ is separated. Applying the descent lemma
\hyperrefext[XII][XII.prop6.8]{6.8}, we see that, in order to prove the bijectivity of the
maps $(\ast)$ for $q\leq n$ and every $F$, it suffices to
prove the bijectivity of the corresponding maps
$H^{q}(Z,h^{\ast}(F))\rightarrow H^{q}(Z',{h'}^{\ast}(F'))$,
provided that the base-change homomorphisms for
$Z\rightarrow X$ and the base change $X'\rightarrow X$ are
isomorphisms in dimensions $\leq n$. If $f$ is separated, then
$h$ is affine, and this is so by hypothesis; this proves \hyperref[XVI.thm1.1]{1.1}
when $f$ is assumed separated. In the general case, we may then
apply the preceding result to $h$, which is always separated, and again
conclude that \hyperref[XVI.thm1.1]{1.1} holds for $f$, proving \hyperref[XVI.lem1.3]{1.3}.

\begin{enonce*}[remark]{Remark}
One could also invoke the Leray spectral sequence for the
open covering $(X_i)$ of $X$, and its noncommutative
variants,\pageoriginale but this method, which does not differ
essentially from using \hyperrefext[XII][XII.prop6.8]{6.8}, has the disadvantage of
forcing us once again to distinguish the usual three cases, work
that has already been done in \loccit
\end{enonce*}

To complete the proof of the theorem in the general case, it
suffices to reduce the proof to the case in which $f:X\rightarrow Y$ is
affine and of finite type, hence quasi-projective. By \hyperref[XVI.lem1.3]{1.3}, we may
suppose that $X$ is affine. Then let $X=\varprojlim X_{\alpha}$, where the
morphisms $f_{\alpha}:X_{\alpha}\rightarrow Y$ are affine and of finite
type. The usual limit argument
(\hyperrefext[VII][VII.sec5]{5}) completes the proof.

\begin{enonce*}[remark]{Complements 1.4}\etiquette[1.4]{XVI.compl1.4}
\begin{enumerate}
\renewcommand{\theenumi}{\alph{enumi}}
\renewcommand{\labelenumi}{\theenumi)}
\item We leave it to the reader to check that the same
  proof still yields a conclusion when $g$ is assumed
  universally locally $(-1)$-acyclic (\hyperrefext[XV][XV.rems1.14]{1.14}): in this
  case, the base-change homomorphism
$$
\phi^{0}:g^{\ast}(f_{\ast}(F))\longrightarrow f'_{\ast}({g'}^{\ast}(F))
$$
is \emph{injective} for every sheaf of sets $F$ on
$X$. On the other hand, suppose again that $n\geq 0$. One may supplement
\hyperref[XVI.thm1.1]{1.1} with the following injectivity statement, which is also contained in the
preceding proof: if $g$ is universally locally
$0$-acyclic, then for every sheaf of groups $F$ on $X$,
the base-change homomorphism $\phi^{1}$ of \hyperref[XVI.thm1.1]{1.1} is a
monomorphism; if $g$ is universally locally $n$-acyclic
for $\LL$, then the homomorphism $\phi^{n+1}$ of \hyperref[XVI.thm1.1]{1.1} is a
monomorphism for every abelian sheaf of $\LL$-torsion $F$
on $X$.

\item It\pageoriginale is probably also possible to obtain an
  analogous injectivity conclusion from the hypothesis that $g$ is
  locally $1$-aspherical for $\LL$, by introducing the
  nonabelian $2$-cohomology invariants from \textsc{Giraud}'s thesis;
  compare \hyperrefext[XII][XII.subsec5.11]{5.11} \footnote[1]{This is indeed what is
  established in J. Giraud's book.}; the same question also arises
  for (\hyperref[XVI.cor1.6]{1.6}) and (\hyperref[XVI.cor2.3]{2.3}) below.
\end{enumerate}
\end{enonce*}

%raghu
\begin{enonce*}{Corollary 1.5}\etiquette[1.5]{XVI.cor1.5}
Let $K/k$ be a field extension, and let
$\LL=\PP-\{p\}$, where $p= \car k$. Then $\Spec
(K)\rightarrow \Spec (k)$ is universally locally acyclic
for $\LL$ and universally locally $1$-aspherical for
$\LL$; if $k$ and $K$ are separably closed, then the
preceding morphism is also universally acyclic for
$\LL$ and universally $1$-aspherical for $\LL$.
\end{enonce*}

We immediately note the following consequence of (\hyperref[XVI.cor1.5]{1.5}).

\begin{enonce*}{Corollary 1.6}\etiquette[1.6]{XVI.cor1.6}
Let $K/k$ be an extension of separably closed fields, let $X$ be a scheme,
and let $F$ be a sheaf of \sisi{set}{sets} (\resp of ind-$\LL$-groups, \resp
of abelian groups of $\LL$-torsion) on $X$, where
$\LL=\PP-\{p\}$, $p=\car k$. Then the map
$$
H^{q}(X,F)\longrightarrow H^{q}(X_{K},F_{K})
$$
is bijective for $q=0$ (\resp for $q\leq 1$, \resp for every $q$).
\end{enonce*}

Indeed, this follows at once from the fact that $\Spec(K)\rightarrow
\Spec(k)$ is universally acyclic for $\LL$ and
universally $1$-aspherical for $\LL$, and from the
definitions.\sisi{---~}{~}Note that, by \hyperrefext[VIII][VIII.cor2.4]{2.4}, the preceding homomorphism
on the $H^{q}$ may also be interpreted as being identified\sisi{}{,}
with the base-change homomorphism
$$
g^{\ast}(R^{q}f_{\ast}(F))\longrightarrow
R^{q}{f'}_{\ast}({g'}^{\ast}(F)){\quoi},
$$
where\pageoriginale $f:X\rightarrow \Spec(k)$, $f':X_K\rightarrow
\Spec(K)$, $g:\Spec(K)\rightarrow \Spec(k)$, and $g':X_K\rightarrow X$
are the evident morphisms. Thus, when $X$ is quasi-compact and
quasi-separated, the bijectivity of these homomorphisms may also be
regarded, by \hyperref[XVI.thm1.1]{1.1}, as a consequence of the fact that
$\Spec(K)\rightarrow \Spec(k)$ is universally \emph{locally}
acyclic for $\LL$ and universally locally
$1$-aspherical for $\LL$. This latter fact therefore already implies
\hyperref[XVI.cor1.6]{1.6} for $X$ quasi-compact and quasi-separated, which plainly suffices
to imply that $\Spec (K)\rightarrow \Spec(k)$ is
universally acyclic for $\LL$ and universally
$1$-aspherical for $\LL$ (\hyperrefext[XV][XV.def1.7]{1.7}, \hyperrefext[XV][XV.prop1.6]{1.6} (i)).

This shows that, to prove \hyperref[XVI.cor1.5]{1.5}, it suffices to prove the
first assertion of \hyperref[XVI.cor1.5]{1.5}. After passing to the perfect closure of
$k$, which is permissible by \hyperrefext[VIII][VIII.thm1.1]{1.1},
we may then suppose that $k$ is \emph{perfect}. Every extension $K$ of $k$ is the filtered direct limit of its
finitely generated $k$-subalgebras $A$, and since $k$ is perfect, after
localizing such an $A$, we may suppose it smooth, so that $K$
appears as a filtered direct limit of smooth subalgebras. Since the
latter are universally locally acyclic for
$\LL$ and universally locally $1$-aspherical for
$\LL$ by \hyperrefext[XV][XV.thm2.1]{2.1}, the conclusion follows from \hyperrefext[XV][XV.prop1.13]{1.13} (ii),
\cqfd

\begin{enonce*}[remark]{Remark 1.7}\etiquette[1.7]{XVI.rem1.7}
The invariance theorem \hyperref[XVI.cor1.5]{1.5} should be compared with
\hyperrefext[XII][XII.cor5.4]{5.4}, where it was not necessary to suppose that the torsion sheaf
under consideration is prime to the residue characteristics, but where, on the other hand,
$X$ must be assumed \emph{proper} over $k$. Already in the case of
$H^{1}$ $(X,\ZZ/p\ZZ)$, where $X$\pageoriginale is an affine algebraic
curve (for example, the affine line) over an algebraically closed
field $k$ of characteristic $p>0$, the analogue of the preceding statements becomes false,
because of the phenomena of ``wild ramification\nde{That is, ``wild.''}'' at infinity,
which imply that, in such a case, the classification of principal étale covers
with group $\ZZ/p\ZZ$ is essentially ``continuous.''
\end{enonce*}

% END INCLUDED SGA4_English_current_buildable_source_20260906/source/16/drafts/SGA4_XVI_LINES_43_274_English.texfrag


% BEGIN INCLUDED SGA4_English_current_buildable_source_20260906/source/16/drafts/SGA4_XVI_LINES_275_440_English.texfrag
\section{Specialization theorem for cohomology groups}\etiquette[2]{XVI.sec2}

\begin{enonce*}{Theorem 2.1}\etiquette[2.1]{XVI.thm2.1}
Let $\LL\subset \PP$, let $n\in \NN$, and let
$f:X\rightarrow Y$ be a proper morphism of finite presentation
\footnote[1]{For a less restrictive hypothesis than properness that yields
  the same conclusions, \cf \SGA 5 II 3 and \SGA 1 XIII.} that is locally
$0$-acyclic (\resp locally $1$-aspherical for $\LL$, \resp locally
$n$-acyclic for $\LL$). Let $F$ be a sheaf of sets (\resp of
\emph{ind}-$\LL$-groups, \resp an abelian sheaf of $\LL$-torsion) that is
constructible and locally constant on $X$. Then the $R^{q}f_{\ast}F$ are
constructible and \emph{locally constant} for $q=0$ (\resp for $q=0$, $1$,
\resp for $q\leq n$), and for every geometric point $\bar{y}$ of $Y$ one has
$$
(R^{q}f_{\ast}F)_{\bar{y}}\simeq
H^{q}(X_{\bar{y}},F_{\bar{y}}){\quoi},
$$
for the same values of $q$.
\end{enonce*}

\begin{enonce*}[remark]{Remark}
Note that the last assertion is simply the proper base-change theorem
\hyperrefext[XII][XII.cor5.2]{5.2}.
\end{enonce*}

In view of \hyperrefext[XV][XV.thm2.1]{2.1}, we immediately obtain the
following corollary:\pageoriginale

\begin{enonce*}{Corollary 2.2}\etiquette[2.2]{XVI.cor2.2}
(Specialization theorem for cohomology groups). Let
$f:X\rightarrow Y$ be a proper smooth morphism, and let
$\LL\subset \PP$ be the complement of the set of residual characteristics
of $Y$. Let $F$ be a sheaf of sets (\resp of ind-$\LL$-groups, \resp an
abelian sheaf of $\LL$-torsion) that is constructible on $X$. Then the
$R^{q}f_{\ast}(F)$ are constructible and \emph{locally constant} for $q=0$
(\resp for $q=0$, $1$, \resp for every $q$), and for every geometric point
$\bar{y}$ of $Y$ one has
$$
R^{q}f_{\ast}(F)_{\bar{y}}\simeq H^{q}(X_{\bar{y}},F_{\bar{y}})
$$
for the same values of $q$.
\end{enonce*}

\begin{enonce*}[remark]{Remark}
Taking $n=1$ and translating into the language of fundamental groups, we
recover (\SGA 1 X 3.8), which compares the fundamental groups of the fibers
of a proper smooth morphism.
\end{enonce*}

\noindent
\emph{Proof of} \hyperref[XVI.thm2.1]{2.1}. Since we already know that the
$R^{q}f_{\ast}(F)$ are constructible \hyperrefext[XIV][XIV.thm1.1]{1.1}, it
suffices, by \hyperrefext[IX][IX.prop2.11]{2.11}, to prove that for every
specialization morphism $\bar{y}_1\rightarrow \bar{y}_{0}$ of geometric
points of $Y$, the corresponding specialization morphism
(\hyperrefext[VIII][VIII.subsec7.7]{7.7})
\begin{equation*}
R^{q}f_{\ast}(F)_{\bar{y}_{0}}\longrightarrow
R^{q}f_{\ast}(F)_{\bar{y}_1}\tag{$\ast$}
\end{equation*}
is bijective for the values of $q$ under consideration. For this we have the
following slightly more general result (in which it is unnecessary to
assume\pageoriginale that $f$ is of finite presentation or that $F$ is
constructible):

\begin{enonce*}{Corollary 2.3}\etiquette[2.3]{XVI.cor2.3}
Let $\LL\subset \PP$, let $n\in \NN$, let
$f:X\rightarrow Y$ be a proper morphism that is locally $0$-acyclic (\resp
locally $1$-aspherical for $\LL$, \resp locally $n$-acyclic for $\LL$), and
let $F$ be a sheaf of sets (\resp of ind-$\LL$-groups, \resp an abelian
sheaf of $\LL$-torsion) on $X$ that is locally constant. Then for every
specialization morphism $\bar{y}_1\rightarrow \bar{y}_0$ of geometric points
of $Y$, the corresponding specialization homomorphism $(\ast)$ above is
bijective for $q=0$ (\resp for $q=0$, $1$, \resp for $q\leq n$). Moreover,
when $n=0$, \ie when $f$ is locally $0$-acyclic, for every locally constant
sheaf of groups $F$ on $X$, the homomorphism
$R^{1}f_{\ast}(F)_{\bar{y}_0}\rightarrow R^{1}f_{\ast}(F)_{\bar{y}_1}$ is
injective; for arbitrary $n$, and for every locally constant abelian sheaf
of $\LL$-torsion $F$ on $X$, the specialization homomorphism $(\ast)$ is
injective for $q=n+1$.
\end{enonce*}

There exists an integral strictly local scheme $Y'$ and a morphism
$g:Y'\rightarrow Y$ that maps the closed point $y'_0$ to $y_{0}$ and the
generic point $y'_1$ to $y_1$: for example, it is enough first to take the
strict localization of $Y$ at $\bar{y}_0$, and then its integral closed
subscheme defined by a point lying above $y_1$. Moreover, after replacing
$Y'$ by its normalization in an algebraic closure of its function field, we
may assume that $Y'$ is normal and that $k(y'_1)$ is algebraically closed,
so that $y'_0$ and $y'_1$ are geometric points of $Y'$. Using the
base-change theorem (\hyperrefext[XII][XII.thm5.1]{5.1}) for $(f,F,g)$, the
proof of \hyperref[XVI.cor2.3]{2.3} is reduced, after replacing\pageoriginale
$Y$, $\bar{y}_0$, $\bar{y}_1$ by $Y'$, $y'_0$, $y'_1$, to proving \hyperref[XVI.cor2.3]{2.3} when
$Y$ is normal and strictly local and when $\bar{y}_0$ and $\bar{y}_1$ are,
respectively, its closed point $y_0$ and its generic point $y_1$. Moreover,
using \hyperrefext[XII][XII.sec4]{4}, we see that the specialization
homomorphism $(\ast)$ is then identified with the homomorphism
$$
H^{q}(X,F)\longrightarrow H^{q}(X_1,g^{\ast}_1(F)){\quoi},
$$
where $X_1=X_{y_1}$ and where $g:y_1\rightarrow Y$ and
$g_1:X_1\rightarrow X$ are the canonical morphisms. Our assertion now
follows from the following lemma (in which $f$ is no longer assumed proper):

\begin{enonce*}{Lemma 2.4}\etiquette[2.4]{XVI.lem2.4}
With the preliminary notation of $\hyperref[XVI.cor2.3]{2.3}$, drop the properness assumption on
$f$, but assume instead that $Y$ is normal and integral, with generic point
$y_1$ such that $k(y_1)$ is separably closed, and let
$X_1=X_{y_1}$ and $g_1:X_1\rightarrow X$ be the canonical morphism. Then the
homomorphism
\begin{equation*}
H^{q}(X,F)\longrightarrow H^{q}(X_1,g^{\ast}_1(F))\tag{${\ast}{\ast}$}
\end{equation*}
is bijective for $q=0$ (\resp bijective for $q=0$, $1$, \resp bijective for
$q\leq n$ and injective for $q=n+1$). Moreover, when $n=0$, \ie when $f$ is
locally $0$-acyclic, for every sheaf of groups $F$ on $X$, the homomorphism
$({\ast}{\ast})$ is injective for $q=1$.
\end{enonce*}

An immediate argument (\hyperrefext[XV][XV.prop1.6]{1.6} (i) $\Rightarrow$
(ii)) shows that it is enough to prove the relation
$$
F\xrightarrow{\sim} g_{1\ast}g_1{}^{\ast}(F){\quoi},
$$
and,\pageoriginale for $n\geq 1$, the additional relations
$$
R^{q}g_{1\ast}g_1{}^{\ast}(F)=0\quad\mbox{\rm for}\quad 1\leq
q\leq n{\quoi}.
$$
These relations are local on $X$ for the etale topology, and $F$ is locally
constant for the etale topology, so we may assume that $F$ is constant and
hence of the form $G_X$, where $G$ is a set (\resp an ind-$\LL$-group,
\resp an abelian group of $\LL$-torsion). Note also that the assumption that
$Y$ is normal implies
$g_{\ast}(G_{y_1})\xleftarrow{\sim}G_Y$
(\hyperrefext[IX][IX.lem2.14.1]{2.14.1}), and the assumption that $k(y_1)$
is separably closed implies that the $R^{q}g_{\ast}(G_{y_1})$ vanish for
$q\geq 1$. On the other hand, the assumption that $f$ is locally
$0$-acyclic (\resp\ldots) implies that the formation of the
$R^{q}g_{\ast}(G_{y_1})$ commutes with the base change $f:X\rightarrow Y$
for $q\leq n$. Here we use the fact that we may assume $Y$ affine and that
$y_1$ then appears as the inverse limit of its affine open neighborhoods,
which are quasi-compact in $Y$, so that Definition
\hyperrefext[XV][XV.def1.11]{1.11} and the theory of passage to the limit
\hyperrefext[VII][VII.sec5]{5} apply. We conclude that indeed
$$
G_X\xrightarrow{\sim} g_{1\ast}(G_{X_1})\quav
R^{q}g_{1\ast}(G_{X_{1}})=0\quad\mbox{\rm for}\quad 1\leq q\leq
n{\quoi},
$$
which completes the proof of \hyperref[XVI.lem2.4]{2.4}, and hence that of
\hyperref[XVI.thm2.1]{2.1}.

\begin{enonce*}[remark]{Remark}
The preceding proof, via \hyperref[XVI.lem2.4]{2.4}, which is simpler than
our original proof, is due to M. Lubkin.
\end{enonce*}

\begin{enonce*}{Corollary 2.5}\etiquette[2.5]{XVI.cor2.5}
Under the hypotheses of $\hyperref[XVI.thm2.1]{2.1}$, assume that $Y$ is connected, and let
$q\leq n$. Then the groups $H^{q}(X_{\bar{y}},F)$, for geometric
points\pageoriginale $\bar{y}$ of $Y$, are all isomorphic to one another.
\end{enonce*}

% END INCLUDED SGA4_English_current_buildable_source_20260906/source/16/drafts/SGA4_XVI_LINES_275_440_English.texfrag


% BEGIN INCLUDED SGA4_English_current_buildable_source_20260906/source/16/drafts/SGA4_XVI_LINES_441_700_English.texfrag
\section{The relative cohomological purity theorem}\etiquette[3]{XVI.sec3}

\begin{enonce*}{Definition 3.1}\etiquette[3.1]{XVI.def3.1}
Let $S$ be a scheme. A \emph{smooth $S$-pair} $(Y,X)$ is a closed
$S$-immersion $i:Y\rightarrow X$ of smooth $S$-preschemes, that is, a
commutative diagram of preschemes
$$
\xymatrix@=1.25cm{
Y \ar[rr]^{i} \ar[rd]_{h}&     &X\ar[ld]^{f}\\
  & S   &
}
$$
such that $f$ and $h$ are smooth and $i$ is a closed immersion. We
write $U=X-Y$, denote by $j:U\rightarrow X$ the open immersion, and
by $g:U\rightarrow S$ the structural morphism. The \emph{codimension}
of $(Y,X)$ at a point $y\in Y$ is the codimension at $y$ of the fiber
$Y_s$ in $X_s$, where $s=h(y)$.
\end{enonce*}

Thus the codimension is a locally constant function on $Y$. We shall
express the fact that this function is in fact constant,
\sisi{with}{of} value~$c$, by saying that ``$(Y,X)$ has
codimension~$c$.''

Let $(Y,X)$ be a smooth $S$-pair and let $y\in Y$. Recall (\SGA II
4.10) that there exist integers $m$ and $n$, a neighborhood $X'$ of
$y$ in $X$, and an étale morphism
$$
\phi:X'\longrightarrow \EE^n_S=\Spec \Oo_S[t_1,\ldots,t_n]
$$
such\pageoriginale that $Y'=Y\cap X'$ is the inverse image of the
closed subscheme $\EE^m_S$ defined by the equations
$$
t_{m+1}=\ldots=t_n=0{\quoi}.
$$
This result can be expressed by saying that \emph{locally for the
  étale topology, every smooth $S$-pair $(Y,X)$ is isomorphic to the
  standard pair $(\EE^m_S, \EE^n_S)$ for suitable $m$ and $n$,} where
$$
\EE^r_S=\Spec \Oo_S[t_1,\ldots,t_r]{\quoi}.
$$

In the present section, we propose to calculate the local cohomology
sheaves $\HH^q_Y(X,F)=(R^{q}i^{!})F$
(\sisi{\hyperrefext[V][V.sec4]{4}}{\hyperrefext[V][V.sec6]{6}} and
\hyperrefext[VIII][VIII.subsec6.6]{6.6}) for a smooth $S$-pair $(Y,X)$,
with coefficients in a locally constant abelian torsion sheaf $F$
prime to the residue characteristics. We begin with some preliminary
considerations:

\begin{enonce*}{Proposition 3.2}\etiquette[3.2]{XVI.prop3.2}
Let
$$
\xymatrix@=1.25cm{
U \ar[rr]^{j} \ar[rd]_{g}&     &X\ar[ld]^{f}\\
  & S   &
}
$$
be a commutative diagram in which $f$ is smooth and $j$ is an open
immersion such that the fiber $U_s$ is dense in $X_s$ for every
$s\in S$. Let $F$ be a sheaf of sets on $S$. Then the canonical
morphism
$$
f^{\ast}F\longrightarrow j_{\ast}g^{\ast}F=j_{\ast}j^{\ast}f^{\ast}F
$$
is bijective.
\end{enonce*}

\noindent
\emph{Proof}.\pageoriginale We readily reduce to the case where $U$ is
quasi-compact in $X$. Let $\bar{s}$ be a geometric point of $S$, let
$\bar{x}$ be a geometric point of $X$ lying over $\bar{s}$, let
$\widetilde{f}:\widetilde{X}\rightarrow \widetilde{S}$ be the morphism
of the strict localizations of $X$ and $S$ at the corresponding
geometric points, and let
$\widetilde{U}=\fprod{U}{\widetilde{X}}{X}$. With the evident notation,
we have
$$
(f^{\ast}F)_{\bar{x}}=H^{0}(\widetilde{X},\widetilde{f}^{\ast}\widetilde{F})\simeq
H^{0}(\widetilde{S},\widetilde{F}){\quoi}.
$$
Consequently, it suffices to prove that
$\widetilde{g}:\widetilde{U}\rightarrow \widetilde{S}$ is
$0$-acyclic, which will imply
$$
(j_{\ast}g^{\ast}F)_{\bar{x}}\simeq
H^{0}(\widetilde{U},\widetilde{g}^{\ast}\widetilde{F})\simeq
H^{0}(\widetilde{S},\widetilde{F}){\quoi}.
$$
Since $\widetilde{g}$ is a limit of smooth morphisms, it is locally
$0$-acyclic (\hyperrefext[XV][XV.thm2.1]{2.1},
\hyperrefext[XV][XV.def1.11]{1.11}), and by
\hyperrefext[XV][XV.cor1.12]{1.12} it suffices to prove that the
geometric fibers of $\widetilde{U}/\widetilde{S}$ are $0$-acyclic,
that is, connected and nonempty. Let $\bar{s}'$ be a geometric point
of $\widetilde{S}$. Then the fiber $\widetilde{X}_{\bar{s}'}$ is
regular, connected, and nonempty, and the hypothesis implies that
$\widetilde{U}_{\bar{s}'}$ is a dense open subset of
$\widetilde{X}_{\bar{s}'}$, hence connected and nonempty, which proves
the result.

\begin{enonce*}{Corollary 3.2.1}\etiquette[3.2.1]{XVI.cor3.2.1}
Under the hypotheses of \hyperref[XVI.prop3.2]{3.2}, the functor
$X'\mapsto X'|U$ from the category of étale coverings of $X$ to the
category of étale coverings of $U$ is fully faithful. In particular,
for every finite group $G$, the restriction map
$$
H^{1}(X,G)\longrightarrow H^{1}(U,G)
$$
is injective.
\end{enonce*}

This\pageoriginale follows from \hyperref[XVI.prop3.2]{3.2}, by a
standard argument which, for ease of reference, we shall spell out in
a

\begin{enonce*}{Lemma 3.2.2}\etiquette[3.2.2]{XVI.lem3.2.2}
Let $f:X\rightarrow Y$ be a morphism of schemes, and let
$f^{\ast}:\Et(Y)\rightarrow \Et(X)$ be the inverse-image functor
$Y'\mapsto \fprod{Y'}{X}{Y}$ from the category $\Et(Y)$ of étale
coverings of $Y$ to $\Et(X)$. The following conditions are equivalent:
\begin{itemize}
\item[(i)] $f^{\ast}$ is faithful (\resp fully faithful).

\item[(ii)] For every étale covering $Y'$ of $Y$, denoting by
  $f':X'\rightarrow Y'$ the morphism deduced from $f$ by the base
  change $Y'\rightarrow Y$, the map $U\mapsto {f'}^{-1}(U)$ from the
  set of clopen subsets of $Y'$ to the set of clopen subsets of $X'$
  is injective (\resp bijective).

\item[(ii bis)] With the notation of {\rm(ii)}, for every constant
  sheaf of sets $C$ on $Y'$, the canonical map
$$
H^{0}(Y',C)\longrightarrow H^{0}(X',{f'}^{\ast}C)
$$
is injective (\resp bijective).

\item[(iii)] With the notation of {\rm(ii)}, the canonical map
$$
\Gamma(Y'/Y)\longrightarrow \Gamma(X'/X)
$$
is injective (\resp bijective).
\end{itemize}
\end{enonce*}

The proof is left to the reader (\cf \SGA 2, IX 3.1 and 3.2).

\begin{enonce*}[remark]{Remarks 3.2.3}\etiquette[3.2.3]{XVI.rems3.2.3}
One\pageoriginale can generalize Proposition
\hyperref[XVI.prop3.2]{3.2} considerably, both for sheaves of sets, by
weakening the hypotheses on
$f$, and for sheaves of groups, with conclusions concerning the
$R^{q}g_{\ast}(f^{\ast}F)$ (\cf \SGA 2 XIV 1 and \SGA 1 XIII). The
analogous remark applies to \hyperref[XVI.cor3.4]{3.4} through
\hyperref[XVI.cor3.6]{3.6}.
\end{enonce*}

Let us also ``recall'' the Zariski--Nagata purity theorem in its
relative form:

\begin{enonce*}{Theorem 3.3}\etiquette[3.3]{XVI.thm3.3}
(Relative purity theorem). Let
$$
\xymatrix@=1.25cm{
U \ar[rr]^{j} \ar[rd]_{g}&     &X\ar[ld]^{f}\\
  & S   &
}
$$
be a commutative diagram in which $f$ is smooth and $j$ is an open
immersion such that, for every $s\in S$, $X_s-U_s$ has codimension
$\geq 2$ everywhere in $X_s$. Then the functor $X'\mapsto X'|U$ from
the category of étale coverings of $X$ to the category of étale
coverings of $U$ is an \emph{equivalence of categories.}
\end{enonce*}

\noindent
\emph{Proof}. We saw in \hyperref[XVI.cor3.2.1]{3.2.1} that the
functor is fully faithful. The fact that it is
\sisi{ess.}{essentially} surjective will follow from the classical
theorem. Given an étale covering $U'$ of $U$, we must
show that it extends to an étale covering $X'$ of $X$. There is an
immediate and easy reduction first to the case where $U$ is
quasi-compact in $X$, and then to the case where $S$ is affine and of
finite type over $\Spec \ZZ$, $X$ is connected, and $f$ is of finite
type.

We\pageoriginale proceed by induction on $n=\dim S$, the assertion
being a consequence of the classical theorem (\SGA 2 X 3.4) when
$n=0$. Since uniqueness is available, we can apply the descent
technique of \sisi{\SGA IX 4.7}{\SGA 1 IX 4.7} to the normalization
$S\leftarrow \bar{S}$ and thereby reduce to the case where $S$ (and
hence $X$) is normal ($\bar{S}\rightarrow S$ being finite (\EGA IV
7.8.3 (iii) (vi)). Moreover, after enlarging $U$, we may suppose that
$U$ is maximal among the open subsets over which $U'$ can be extended,
and we must then prove that $U=X$. Otherwise, let $x$ be a maximal
point of $X-U$, and put $s=f(x)$. It is enough to prove that $U'$ can
be extended over a neighborhood of $x$. We may suppose that
$\dim \Oo_{S,s}=n\geq 1$ and that the theorem has already been proved
in dimension $<n$. Put $B=\Oo_{X,x}$, and let $B'$ be the normalization
of $B$ in the ring of rational functions on the covering $U'$ of $U$;
thus $B'$ is a finite $B$-algebra (\loccit). It plainly suffices to
prove that $B'$ is an étale $B$-algebra.

Let $A=\Oo_{S,s}$, $0\neq t\in \max A$, $A_{0}=A/tA$, and
$B_{0}=B/tB$. By the induction hypothesis, the restriction of the
covering to $U_{0}=U\otimes_A A_{0}$ extends to an étale covering of
$\Spec B_{0}$, which we denote by $\Spec B^{\ast}_{0}$. After replacing
$X$ by a suitable étale covering of a neighborhood of $x$ (which is
permissible), we may suppose that $\Spec B^{\ast}_{0}$ is completely
split over $\Spec B_{0}$. We are thereby reduced to the case where the
$B_{0}$-algebra $B'_{0}$ is completely split away from the closed
point of $\Spec B_{0}$. The hypotheses of
\hyperrefext[XV][XV.lem3.3]{3.3} are then satisfied, which proves the
result.

\begin{enonce*}{Corollary 3.4}\etiquette[3.4]{XVI.cor3.4}
Under the same hypotheses as in \hyperref[XVI.thm3.3]{3.3}, we have
$(R^{1}j_{\ast})g^{\ast}F=0$ for every sheaf $F$ of ind-finite groups
on $S$.
\end{enonce*}

We reduce\pageoriginale as usual to the case where $S$ is noetherian,
$f$ is of finite type, and $F$ is constructible. Let
$0\rightarrow F\rightarrow G$ be an injection of sheaves of
ind-finite groups on $S$, and put $C=G/F$. It follows from
\hyperref[XVI.prop3.2]{3.2} that the functor $j_{\ast}g^{\ast}$ is
exact, so that $j_{\ast}g^{\ast}G\rightarrow j_{\ast}g^{\ast}C$ is
surjective, and hence (\hyperrefext[XII][XII.prop3.1]{3.1}) that
$(R^{1}j_{\ast})g^{\ast}F\rightarrow
(R^{1}j_{\ast})g^{\ast}G$ is injective. We can therefore replace $F$
by $G$, and thus reduce, by \hyperrefext[IX][IX.prop2.14]{2.14} and
\hyperrefext[VIII][VIII.prop5.5]{5.5}, to the case where $F=G$ is
constant, with value an ordinary finite group $G$.
%% This reduction is not perfectly clear in the nonabelian case:
%% give a reference for R^1 j_* p_*G = p_* R^1 j_*G
Let $\bar{x}$ be a geometric point of $X$, let $\widetilde{X}$ be the
strict localization of $X$ at $\bar{x}$, and let
$\widetilde{U}=\fprod{\widetilde{X}}{U}{X}$. It then follows from
\hyperref[XVI.thm3.3]{3.3} that
$$
(R^{q}j_{\ast}g^{\ast}F)_{\bar{x}}=H^{q}(\widetilde{U},G_{\widetilde{U}})\simeq
H^{q}(\widetilde{X},G_{\widetilde{X}}){\quoi},
$$
and the latter is zero, which proves the result.

\begin{enonce*}{Lemma 3.5}\etiquette[3.5]{XVI.lem3.5}
(Relative Abhyankar's lemma). Let $(Y,X)$ be a smooth $S$-pair of
  codimension $1$, such that $Y$ is defined in $X$ by an equation
  $t=0$. Let $U=X-Y$, and let $V/U$ be a principal Galois covering of
  order $n$ prime to the residue characteristics of $Y$. Put
  $X'=\Spec \Oo_X[z]/(z^n-t)$, $U'=\fprod{U}{X'}{X}$, and
  $V'=\fprod{V}{X'}{X}$. Then the covering $V'$ extends uniquely to an
  étale covering of $X'$.
\end{enonce*}

\noindent
\emph{Proof}. We immediately reduce to the case where $S$, $X$, and
$Y$ are affine and of finite type over $\Spec \ZZ$. Uniqueness is a
special case of \hyperref[XVI.cor3.2.1]{3.2.1}. We shall deduce
existence from \hyperrefext[XV][XV.cor1.12]{1.12}: let $\bar{s}$ be a
geometric point of $S$ and let $\bar{x}$ be a geometric point of $Y$
lying over $\bar{s}$. Let
$\widetilde{f}:\widetilde{X}\rightarrow \widetilde{S}$ be the morphism
of the corresponding strict localizations\pageoriginale induced by
$f$, let $\widetilde{U}=\fprod{U}{\widetilde{X}}{X}$, etc\ldots By
descent, in view of uniqueness, it plainly suffices to prove that
$\widetilde{V}'$ extends to an étale covering of $\widetilde{X}'$, that
is (since $\widetilde{X}'$ admits no nontrivial covering), that
$\widetilde{V}'$ is a trivial covering of $\widetilde{U}'$.

Let

% END INCLUDED SGA4_English_current_buildable_source_20260906/source/16/drafts/SGA4_XVI_LINES_441_700_English.texfrag


% BEGIN INCLUDED SGA4_English_current_buildable_source_20260906/source/16/drafts/SGA4_XVI_LINES_701_956_English.texfrag
$\bar{U}=\varprojlim_m\widetilde{U}_{m}=\Spec \Oo_{\widetilde{U}}[t^{1/m}]$,
where $m$ runs through the set of integers prime to the residue
characteristic of $\widetilde{X}$. Then
$\bar{U}/\widetilde{U}$ is an “infinite Galois covering” with abelian
group $G=\varprojlim_m \ZZ/m\ZZ$, and one sees immediately that it is
enough to prove that $\bar{U}$ is $1$-aspherical for the set
$\LL\subset \PP$ complementary to the residue characteristic of
$\widetilde{X}$.

Consider the morphism $\bar{f}:\bar{U}\rightarrow
\widetilde{S}$. It is a limit of smooth morphisms, and is therefore
locally $1$-aspherical for $\LL$ (\hyperrefext[XV][XV.thm2.1]{2.1}, \hyperrefext[XV][XV.def1.11]{1.11}
(ii)). By \hyperrefext[XV][XV.cor1.12]{1.12}, it is enough to prove that the geometric
fibres of $\bar{U}/\widetilde{S}$ are $1$-aspherical for $\LL$. Let
$\bar{s}$ be a geometric point of $\widetilde{S}$. The geometric fibre
$\bar{U}_{\bar{s}}$ is the limit of the fibres
$(\widetilde{U}_m)_{\bar{s}}$. Consequently, it is enough to prove that
every principal Galois covering of some
$(\widetilde{U}_m)_{\bar{s}}$, of order prime to the
\sisi{res.\ char.}{residue characteristics}, induces a trivial covering
of $\bar{U}_{\bar{s}}$. Replacing $\widetilde{X}$ by
$\widetilde{X}_m=\Spec \Oo_{\widetilde{X}}[t^{1/m}]$ (which is still
smooth over $\widetilde{S}$), we reduce to the case $m=1$, that is,
$\widetilde{U}_m=\widetilde{U}$. Let $V$ be such a covering of
$\widetilde{U}_{\bar{s}}$. Now $\widetilde{X}_{\bar{s}}=Z$ is a
regular strictly local scheme, and $W=\widetilde{U}_{\bar{s}}$ is the
open complement of a regular subset of codimension $1$. We may
therefore apply the usual form of Abhyankar's lemma
(\SGA 1 X 3.6) to conclude that $V$ induces a
covering\pageoriginale $V_m$ of $W_m$ which extends to an étale
covering of $Z_m$, for a suitable $m$. But since $Z$, and hence $Z_m$,
is strictly local, this covering is trivial; therefore $V_m$ is a
trivial covering of $W_m$, \cqfd

\begin{enonce*}[remark]{Remark 3.5.1}\etiquette[3.5.1]{XVI.rem3.5.1}
When the order $n$ of $G$ is not assumed prime to the residue
characteristics of $X$, one can still prove the following: let $m$ be
the largest integer prime to the residue characteristics of $Y$ which
divides the order of $G$, and suppose that for every point $s\in S$,
the covering $V_s$ of $U_s=X_s-Y_s$ has, at the maximal points of
$Y_s$, inertia groups whose orders are prime to the residue
characteristics of $Y$. Then the conclusion of
\hyperref[XVI.lem3.5]{3.5} remains valid on taking
$X'=\Spec \Oo_X[z]/(z^m-t)$. More generally, one may also take for $m$
any common multiple of the orders of the inertia groups just
considered, provided these orders are prime to the corresponding
residue characteristics.
\end{enonce*}

\begin{enonce*}{Corollary 3.6}\etiquette[3.6]{XVI.cor3.6}
(Cohomological purity in dimension $1$). Let $(Y,X)$ be a smooth
  $S$-pair of codimension $1$, and let $F$ be a constructible locally
  constant sheaf of groups on $X$, whose orders are prime to the
  residue characteristics of $X$. Then, with the notation of $\hyperref[XVI.def3.1]{3.1}$,
  $(R^{1}j_{\ast})j^{\ast}F$ is locally isomorphic, as a sheaf of
  pointed sets, to $i_{\ast}(F_Y/\cint(F_Y))$, where $F_Y=F|Y$ and
  where $F_Y/\cint(F_Y)$ denotes the quotient of $F_Y$ by its action
  on itself through inner automorphisms.
\end{enonce*}

\begin{enonce*}[remark]{Remark}
The isomorphism of Corollary \hyperref[XVI.cor3.6]{3.6} is not
canonical.
\end{enonce*}

\noindent
\emph{Proof}.\pageoriginale We immediately reduce to the case where
$X$ is noetherian and connected. Since the assertion is local on $X$
for the étale topology, we may suppose that $F$ is a constant sheaf,
with value an ordinary finite group $G$ of order $m$ prime to the
residue characteristics of $X$, and moreover that
$𝛍_{mX}\simeq (\ZZ/m\ZZ)_X$. Let $U=X-Y$ and let
$U'=\Spec \Oo_U[t^{1/m}]$, which is an “essential” principal Galois
covering of $U$ with group $\ZZ/m\ZZ$, since
$𝛍_{mX}\simeq (\ZZ/m\ZZ)_X$. Principal coverings of $U$ with group
$G$ which are trivialized over $U'$ are classified locally on $Y$ by
$H=\Hom(\ZZ/m\ZZ,G)$ modulo $\cint(G)$, which is a pointed set
isomorphic to the underlying set of $G/\cint(G)$. This gives a
morphism from the constant sheaf $H_Y$ on $Y$ to
$(R^{1}j_{\ast})G_U$, and it is enough to prove that this morphism is
bijective, which follows immediately from \hyperref[XVI.lem3.5]{3.5}.

\begin{enonce*}{Theorem 3.7}\etiquette[3.7]{XVI.thm3.7}
\emph{(Cohomological purity)}. Let $(Y,X)$ be a smooth $S$-pair of
codimension $c>0$. Let $F$ be an abelian sheaf on $X$ locally
isomorphic (for the étale topology) to a sheaf of the form
$f^{\ast}(G)$, where $G$ is a torsion sheaf on $S$, prime to the
residue characteristics (for example, $F$ may be a locally constant
sheaf of finite groups whose orders are prime to the residue
characteristics of $X$). Then, with the notation of $\hyperref[XVI.def3.1]{3.1}$, one has
\begin{align*}
&\HH^{q}_Y(X,F)=(R^{q}i^{!})F=0\qquad\mbox{if}\quad q\neq 2c,\\
&\begin{cases}
 F\xrightarrow{\sim} j_{\ast}j^{\ast}F\quad\text{and} &\tag{\ie}\\
(R^{q}j_{\ast})j^{\ast}F=0 &\hspace{1.3cm} \mbox{if}\quad q\neq 0,\quad 2c-1,
\end{cases}
\end{align*}
and\pageoriginale
$$
\HH^{2c}_Y(X,F)=i^{\ast}[(R^{2c-1}j_{\ast})j^{\ast}F]
$$
is a sheaf locally isomorphic to $i^{\ast}(F)$.
\end{enonce*}

\noindent
\emph{Proof}. The assertion is local on $X$ for the étale topology;
we may therefore \footnote[1]{At least if $F$ were assumed locally
  constant. The reader will convince himself that the proof which
  follows also applies, essentially, to the general case.}
suppose that $F$ is constant, and hence, if desired, that
$F\simeq f^{\ast}(\ZZ/n\ZZ)_S$ for the constant sheaf
$(\ZZ/n\ZZ)_S$ on $S$.

First treat the case $c=1$. (The assertion concerning
$(R^{1}j_{\ast})j^{\ast}F$, at least in the case where $F$ is locally
constant, would follow immediately from \hyperref[XVI.cor3.6]{3.6},
but we shall derive it again.) Recall that $(Y,X)$ is locally
isomorphic to the pair $(\EE^{n-1}_S,\EE^n_S)$, and that, upon
replacing $S$ by $\EE^{n-1}_S$, we reduce to the case where
$X=\EE^{1}_S=\Spec \Oo_S[t]$ and $Y$ is the section $t=0$ of
$X/S$. This situation is locally isomorphic to the inclusion of the
section at infinity in the projective line $\PP^{1}_S$. We may
therefore take $X=\PP^{1}_S$, take $Y$ to be the section at infinity,
and put $U=\EE^{1}_S=X-Y$.

Consider the Leray spectral sequence
$$
(R^{p}f_{\ast})(R^{q}j_{\ast})(\ZZ/n)\Longrightarrow
(R^{p+q}g_{\ast})(\ZZ/n){\quoi}.
$$
Its abutment is zero for $p+q>0$, by \hyperrefext[XV][XV.cor2.2]{2.2}.
Moreover, the $R^{q}j_{\ast}$ are supported on $Y$ when $q>0$, and
$Y$ maps isomorphically onto $S$ under $f$; on the other hand,
$j_{\ast}(F|U)\xleftarrow{\sim}F$ by \hyperref[XVI.prop3.2]{3.2}. It
follows that $(R^{p}f_{\ast})(R^{q}j_{\ast})=0$ if $p,q>0$, that
$f_{\ast}(R^{q}j_{\ast})$ determines the sheaf $(R^{q}j_{\ast})$ when
$q>0$, and finally that
$\sisi{\RR}{R}^{p}f_{\ast}j_{\ast}(F|U)=
\sisi{\RR}{R}^{p}f_{\ast}(F)$. The spectral sequence therefore
reduces to isomorphisms\pageoriginale
$$
f_{\ast}(R^{q}j_{\ast})(F|U)\xrightarrow{\sim}
(R^{q+1}f_{\ast})(F)\quad\mbox{\rm if}\quad q>0{\quoi}.
$$
But $(R^{q}f_{\ast})(F)$ can be computed fibre by fibre
(\hyperrefext[XII][XII.cor5.2]{5.2}); it is therefore zero unless
$q=0,2$, and is a sheaf locally isomorphic to $G$ if $q=2$, as follows
easily from \hyperrefext[XII][XII.cor5.2]{5.2} and
\hyperrefext[IX][IX.cor4.7]{4.7}. Thus $(R^{q}j_{\ast})(F|U)=0$ if
$q\neq 1$, and it is a sheaf locally isomorphic to $i^{\ast}(F)$ if
$q=1$.

The case $c>1$ is now handled easily by induction: let $(Y,X)$ be a
smooth $S$-pair of codimension $c>1$. It is clear that, locally on
$X$, one can find a subscheme $Z$ of $X$ of codimension $1$ which
contains $Y$ and is smooth over $S$, so that there is a commutative
diagram
$$
\xymatrix@=1.25cm{
Y       \ar[rr]^{u} \ar[rd]_{i}&    & Z\ar[ld]^v            \\
                  & X &,
}
$$
in which each pair is a smooth $S$-pair (one ought to call
$(Y,Z,X)$ a smooth $S$-triple), the codimension of $(Z,X)$ is $1$,
and that of $(Y,Z)$ is $c-1$.

Now there is a spectral sequence of composite functors
$$
(R^{p}u^{!})(R^{q}v^{!})F\Longrightarrow (R^{p+q}i^{!})F{\quoi}.
$$
By the case already proved,
$(R^{q}v^{!})F=0$ if $q\neq 2$ and is locally isomorphic to
$v^{\ast}(F)$ if $q=2$. The induction hypothesis, applied to
$(R^{q}v^{!})F$, then shows that
$(R^{p}u^{!})(R^{q}v^{!})F$ is zero if
$(p,q)\neq(2c-2,2)$\pageoriginale and is locally isomorphic to
$u^{\ast}v^{\ast}(F)$ if $(p,q)=(2c-2,2)$. The result for
$(R^{p+q}i^{!})F$ follows immediately, \cqfd

\begin{enonce*}{Corollary 3.8}\etiquette[3.8]{XVI.cor3.8}
Let $(Y,X)$ be a smooth $S$-pair of codimension $c$, let $n$ be a
natural number prime to the residue characteristics of $X$, and put
$$
T_{Y/X}=\HH^{2c}_Y(\ZZ/n\ZZ){\quoi},
$$
which, by $\hyperref[XVI.thm3.7]{(3.7)}$, is a sheaf on $Y$ locally isomorphic (for the
étale topology) to the sheaf $(\ZZ/n\ZZ)_Y$. Let $F$ be an abelian
sheaf on $X$, annihilated by $n$, satisfying the condition stated in
$\hyperref[XVI.thm3.7]{3.7}$. Then (for the record) $\HH^{i}_Y(F)=0$ for $i\neq 2c$, and
moreover there is a canonical isomorphism
$$
\HH^{2c}_Y(F)\simeq i^{\ast}(F)\otimes T_{Y/X}{\quoi}.
$$
\end{enonce*}

\noindent
\emph{Proof}. One easily defines a canonical homomorphism
$$
\varphi:i^{\ast}(F)\otimes T_{Y/X}\longrightarrow \HH^{2c}_Y(F)
$$
that is, a homomorphism
$$
i^{\ast}(F)\longrightarrow
\sheaf{\Hom}(\HH^{2c}_Y(\ZZ/n\ZZ),\HH^{2c}_Y(F)){\quoi},
$$
using the canonical homomorphism
$F\rightarrow\sheaf{\Hom}(\ZZ/n\ZZ,F)$. It remains to prove that
$\varphi$ is an isomorphism, which is easily checked by following the
proof given for \hyperref[XVI.thm3.7]{3.7}.

\begin{enonce*}{Corollary 3.8.1}\etiquette[3.8.1]{XVI.cor3.8.1}
Under the conditions of \hyperref[XVI.cor3.8]{3.8}, the formation of
the $\HH^{i}_Y(F)$\pageoriginale commutes with every base change
$S'\rightarrow S$.
\end{enonce*}

This is clear.

\begin{enonce*}{Corollary 3.9}\etiquette[3.9]{XVI.cor3.9}
(Absolute cohomological purity theorem). Let $X$ be a regular scheme,
and let $Y$ be a regular closed subprescheme of codimension $c$ at
every point. Suppose moreover that $X$ is locally of finite type over
a perfect field $k$. Then, if $n$ is prime to the residue
characteristics of $X$, one has
$$
\HH^{i}_Y(\ZZ/n\ZZ)=0\mbox{ if } i\neq 2c{\quoi},
$$
and $\HH^{2c}_Y(\ZZ/n\ZZ)$ is a sheaf on $Y$ locally isomorphic (for
the étale topology) to the constant sheaf $\ZZ/n\ZZ$.
\end{enonce*}

This is a special case of \hyperref[XVI.thm3.7]{3.7}, since the
hypothesis implies that $X$ and $Y$ are smooth over $k$.

\begin{enonce*}[remark]{Remarks 3.10}\etiquette[3.10]{XVI.rems3.10}
\begin{enumerate}
\renewcommand{\theenumi}{\alph{enumi}}
\renewcommand{\labelenumi}{\theenumi)}
\item In \hyperref[XVI.cor3.8]{3.8} one can make the structure of the
sheaf $T_{Y/X}$ explicit: one finds that it is canonically isomorphic
to the sheaf $(𝛍_n)_Y^{\otimes -c}$, where $𝛍_n$ denotes the sheaf of
$n$th roots of unity (locally isomorphic to the constant sheaf
$\ZZ/n\ZZ$). The isomorphism in question is studied, under the name
of the local \emph{fundamental class} of $Y$ in $X$, in the slightly
more general setting of relative complete intersections, in Exposé
XVIII, which is devoted to duality.

\item It\pageoriginale is very plausible that the conclusions of
\hyperref[XVI.cor3.9]{3.9} and \hyperref[XVI.rems3.10]{3.10} a) remain
valid without assuming the existence of a base field $k$, at least
when $X$ is an “excellent” prescheme\SGAArchiveOnly{\nde{The frozen source attributes
this later development to work of Gabber but supplies no bibliographic
target.}}. This
will in any case be established in Exposé XIX when $X$ is of
characteristic zero, using Hironaka's resolution of singularities. As
we shall see in \SGA 5, the “absolute purity theorem” (as well as
resolution of singularities) is needed in particular to establish the
“local duality formula”, which itself is one of the principal
ingredients of the \textsc{Lefschetz--Verdier} formula.
\end{enumerate}
\end{enonce*}

% END INCLUDED SGA4_English_current_buildable_source_20260906/source/16/drafts/SGA4_XVI_LINES_701_956_English.texfrag


% BEGIN INCLUDED SGA4_English_current_buildable_source_20260906/source/16/drafts/SGA4_XVI_LINES_957_1200_English.texfrag
\section{Comparison theorem for the cohomology of algebraic schemes
over $\CC$.}\etiquette[4]{XVI.sec4}

Let $X$ be an algebraic scheme over $\Spec \CC$. We shall compare
the étale and classical sites of $X$, and we retain the notation of
\hyperrefext[XI][XI.sec4]{4}. The comparison theorem, whose general
form is \hyperrefext[XI][XI.thm4.4]{4.4}, is as follows:

\begin{enonce*}{Theorem 4.1}\etiquette[4.1]{XVI.thm4.1}
Let $f:X\rightarrow S$ be a morphism of finite type between schemes
locally of finite type over $\Spec \CC$, so that there is a
commutative diagram of morphisms of sites
$$
\xymatrix@C=2cm@R=1.5cm{
X_{\et}  \ar[d]_{f_{\et}}                   &   X_{\cl}
\ar[l]_{\varepsilon} \ar[d]^{f_{\cl}} \\
S_{\et}                    &   S_{\cl}\ar[l]_{\varepsilon}
}
$$
Let\pageoriginale $F$ be a sheaf of sets (\resp of ind-finite
groups, \resp a torsion abelian sheaf), and suppose that one of the
following cases holds:
\begin{enumerate}
\renewcommand{\theenumi}{\roman{enumi}}
\renewcommand{\labelenumi}{(\theenumi)}
\item $f$ is proper,

\item $F$ is constructible.
\end{enumerate}
Then the base-change morphisms \sisi{\emph{XII 4.1.2}}{\hyperrefext[XII][XII.eq4.2]{4.2}}
\footnote[1]{In fact, we use the base-change homomorphism under more
general conditions than those considered in Exposé XII, where we
restricted ourselves to a \emph{Cartesian} square of morphisms of
\emph{preschemes}. The reader will readily verify that the definition
extends verbatim to an \emph{essentially commutative} square of
morphisms of \emph{sites}. See the following exposé for further
developments on this point.}
$$
\varphi:\epsilon^{\ast}(R^{q}f_{\et\ast})F\longrightarrow
(R^{q}f_{\cl\ast})\epsilon {}^{\ast}F
$$
are bijective for $q=0$ (\resp for $q=0$, $1$, \resp for $q\geq
0$).
\end{enonce*}

Observe that the Grauert--Remmert theorem
(\hyperrefext[XI][XI.thm4.3]{4.3} (iii)), which we used in the proof
of \hyperrefext[XI][XI.thm4.4]{4.4}, is a special case of
\hyperref[XVI.thm4.1]{4.1}.

For\pageoriginale the proof, we first treat the case in which $f$ is
proper. In this case, the result will follow from GAGA and the proper
base-change theorem, by an elementary argument that could have
appeared in Exposé XIV. To treat the general case, we shall also use
resolution of singularities [1] and \hyperref[XVI.thm3.7]{3.7}.

We may clearly assume that $S$, and hence also $X$, is of finite type
over $\CC$.

\emph{The case in which $f$ is proper.} We recall the following
elementary facts:
\begin{enumerate}
\renewcommand{\theenumi}{\alph{enumi}}
\renewcommand{\labelenumi}{(\theenumi)}
\item The topological space $X_{\cl}$ is locally compact.

\item If $f:X\rightarrow S$ is proper, then
$f_{\cl}:X_{\cl}\rightarrow S_{\cl}$ is a \emph{proper} map of
topological spaces.
\end{enumerate}

Indeed, (a) is trivial, since locally $X_{\cl}$ is a closed subset of
some $\CC^n$. For (b), Chow's lemma reduces us to the case in which
$f$ is projective, and this case is trivial.

Let $F$ be a sheaf of sets (\resp\ldots) on $X_{\et}$.
For both morphisms $f_{\et}$ and $f_{\cl}$, formation of the $R^{q}$
commutes with formation of the fibers (\hyperrefext[XII][XII.cor5.2]{5.2}
and \cite{sga4.exp.XVI.bib.2}, 4.11.1). It is therefore enough to verify the theorem
fiber by fiber, at the closed points of $S_{\et}$.

In the following lemma, we do not assume that $g$ is proper:

\begin{enonce*}{Lemma 4.2}\etiquette[4.2]{XVI.lem4.2}
Let\pageoriginale
$$
\xymatrix@=1.5cm{
X  \ar[rd]_{g}& &\bar{X} \ar[ll]_{h}\ar[ld]^{\bar{g}}\\
& S &
}
$$
be a commutative diagram with $h$ proper, and let $U$ be an open
subset of $X$ such that $h$ induces an isomorphism from the open
subset $\bar{U}=\fprod{U}{\bar{X}}{X}$ of $\bar{X}$ onto $U$. Let
$j:U\rightarrow X$ and $\bar{j}:\bar{U}\rightarrow \bar{X}$ be the
inclusions, let $F$ be a sheaf of sets (\resp\ldots) on $U$, and let
$\bar{F}$ be the induced sheaf on $\bar{U}$. If Theorem $\hyperref[XVI.thm4.1]{4.1}$ is true
for $(\bar{g},\bar{j}_{!}\bar{F})$, then it is also true for
$(g,j_{!}F)$.
\end{enonce*}

\noindent
\emph{Proof}. The theorem is true for
$(h,\bar{j}_{!}\bar{F})$, because we are reduced to checking it fiber
by fiber, and there are then two cases: either the restriction of the
sheaf to the fiber is zero, or the morphism $h$ is an isomorphism. In
both cases the result is trivial. We find that
$$
\epsilon^{\ast}j_{!}F\simeq
\epsilon^{\ast}h_{\text{\rm et\,}\ast}(\bar{j}_{!}\bar{F})\simeq
h_{\cl\ast}\epsilon^{\ast}(\bar{j}_{!}\bar{F})
$$
and
$$
(R^{q}h_{\text{\rm et\,}\ast})\bar{j}_{!}\bar{F}
=(R^{q}h_{\cl\ast})\epsilon^{\ast}(\bar{j}_{!}\bar{F})=0
\text{ if }q>0
$$
for the values of $q$ under consideration, and hence that
$$
(R^{q}g_{\text{\rm et\,}\ast})j_{!}F\simeq
(R^{q}\bar{g}_{\text{\rm et\,}\ast})\bar{j}_{!}\bar{F}
$$
and
$$
(R^{q}g_{\cl \ast})j_{!}F\simeq (R^{q}\bar{g}_{\cl
\ast})\bar{j}_{!}\bar{F}{\quoi},
$$
which proves \hyperref[XVI.lem4.2]{Lemma 4.2}.\pageoriginale

\begin{enonce*}{Lemma 4.3}\etiquette[4.3]{XVI.lem4.3}
To verify \hyperref[XVI.thm4.1]{4.1} for all data $(f,F)$ with $f$
proper, it is enough to do so when, in addition, $f$ has relative
dimension $\leq 1$ and $S=\Spec \CC$ is a point.
\end{enonce*}

\noindent
\emph{Proof}. We proceed by induction on the relative dimension.
Suppose the result known in relative dimension $<n$, where $n>1$.
As we have observed, we may check it fiber by fiber; that is, we may
suppose that $S=\Spec \CC$ is a point, and hence that $X$ has
dimension $\leq n$. Moreover, we may assume that $X$ is reduced. Let
$\varphi$ be a rational function on $X$ that is nonconstant on
\sisi{every irreducible component}{every irreducible component} of
$X$, let $U$ be a dense open subset of $X$
on which $\varphi$ is defined, and let
$\bar{X}\subset X\times \PP^{1}$ be the closure of the graph of the
morphism $\varphi:U\rightarrow \EE^{1}$. We have a commutative
diagram
$$
\xymatrix@R=1.5cm@C=1.5cm{
&\bar{X}\ar[ld]_{h} \ar[rd]^{a} \ar[dd]^{\bar{g}}   &\\
X  \ar[rd]_{g}  &                                  &\PP^1 \ar[ld]_{b}\\
            & **[r]Y              =\Spec \CC&
}
$$
in which the relative dimensions of $a$ and $b$ are $<n$. By the
induction hypothesis, the theorem is true for $(a,.)$ and
$(b,.)$, and it follows that it is also true for
$(\bar{g},.)$. Indeed, this is trivial for $q=0$; if $F$
is a torsion abelian sheaf on $\bar{X}$, it follows\pageoriginale from
the morphism of Leray spectral sequences
$$
\xymatrix{
\varepsilon^{\ast}(R^{p}b_{\text{et}\ast})(R^{q}a_{\text{et}\ast})F \ar@{=>}[r] \ar[d]&
\varepsilon^{\ast}(R^{p+q}\bar{g}_{\text{et}\ast})F  \ar[d]\\
(R^{p}b_{\text{cl}\ast}(R^{q}{a_{\text{cl}\ast}}))
\varepsilon^{\ast}F  \ar@{=>}[r] & (R^{p+q}\bar{g}_{\text{cl}\ast})\varepsilon^{\ast}F;
}
$$
finally, if $F$ is a sheaf of ind-finite groups on $\bar{X}$, note
that since $\epsilon^{\ast}(R^{1}\bar{g}_{\text{\rm et\,}\ast})$ is
an effaceable functor, it is enough to prove (in view of the result for
$q=0$) that $(R^{1}\bar{g}_{\cl\ast})\epsilon^{\ast}$ is also
effaceable (\cf \hyperrefext[XII][XII.cor8.2]{8.2}). This follows
from the exact sequence
$$
0\longrightarrow (R^{1}b_{\cl\ast})a_{\cl
\ast}\epsilon^{\ast}F\longrightarrow (R^{1}\bar{g}_{\cl
\ast})\epsilon^{\ast}F\longrightarrow b_{\cl \ast}(R^{1}a_{\cl
\ast})\epsilon^{\ast}F{\quoi},
$$
and from the fact that $(R^{1}b_{\cl\ast})\epsilon^{\ast}$ and
$(R^{1}a_{\cl\ast})\epsilon^{\ast}$ are effaceable by the induction
hypothesis.

Now let $F$ be a sheaf of sets on $X$. We have
$F\hookrightarrow h_{\text{\rm et\,}\ast}h_{\et}{}^{\ast}F=G$.
The exact sequence
\begin{equation*}
F\rightarrow G\rightrightarrows G\amalg_F G{\quoi},\tag{4.3.1}
\end{equation*}\etiquette[4.3.1]{XVI.eq4.3.1}
reduces us (for $q=0$) to proving \hyperref[XVI.thm4.1]{4.1} for $G$,
and hence for a sheaf of the form
$h_{\text{\rm et\,}\ast}\bar{F}$. For such a sheaf, however, the
result follows from the theorem for $(\bar{g},.)$ and for
$(h,.)$; and $h$ has relative dimension $<n$, as is
immediate.

Let $F$ be a torsion abelian sheaf on $X$. We have the exact sequence
\begin{equation*}
0\longrightarrow j_{!}j^{\ast}F\longrightarrow F\longrightarrow
i_{\ast}i^{\ast}F\longrightarrow 0{\quoi}.\tag{$\ast$}
\end{equation*}
Since $U$ is dense in $X$, we have $\dim Y<n$, so
\hyperref[XVI.thm4.1]{4.1} is true\pageoriginale for
$(g|Y,i^{\ast}F)$, and hence for $(g,i_{\ast}i^{\ast}F)$. It is
therefore enough, by the five lemma, to verify
\hyperref[XVI.thm4.1]{4.1} for $(g,j_{!}j^{\ast}F)$, and this follows
from the theorem for $\bar{g}$ (\hyperref[XVI.lem4.2]{4.2}).

Let $F$ be a sheaf of ind-finite groups on $X$. It is enough to prove
that the functor $(R^{1}g_{\cl\ast})\epsilon^{\ast}$ is effaceable for
the sheaf $F$. This follows from ($\ast$) and from effaceability for
$j_{!}j^{\ast}F$ and for $i_{\ast}i^{\ast}F$, that is, from the
theorem for $(g,j_{!}j^{\ast}F)$ (\hyperref[XVI.lem4.2]{4.2}) and for
$(g,i_{\ast}i^{\ast}F)$ (the induction hypothesis), \cqfd

\begin{enonce*}{Lemma 4.4}\etiquette[4.4]{XVI.lem4.4}
To prove $\hyperref[XVI.thm4.1]{(4.1)}$ when $f$ is proper, it is enough
to prove the following assertion: let $X$ be a complete regular curve
over $\Spec \CC$, and let $F$ be a constant constructible sheaf of
sets (\resp\ldots) on $X$. Then the morphisms
$$
H^{q}(X_{\et},F)\longrightarrow H^{q}(X_{\cl},F)
$$
are bijective for the values of $q$ under consideration.
\end{enonce*}

\noindent
\emph{Proof}. Since it is enough to verify the theorem fiber by fiber,
\hyperref[XVI.lem4.3]{4.3} reduces us to the case in which $X$ is
proper and has dimension $\leq 1$. The case of dimension $0$ is trivial,
and it follows that the theorem is true for a finite morphism. Since
$f$ is proper, the cohomology of $X_{\et}$ and of $X_{\cl}$ commutes
with filtered direct limits (\hyperrefext[VII][VII.prop3.3]{3.3} and
[2], 4.12.1), and we are therefore reduced to the case in which $F$
is constructible (\hyperrefext[IX][IX.prop2.9]{2.9} (iii)).

We apply \hyperrefext[IX][IX.prop2.14]{2.14}. For a sheaf of sets, the
exact sequence \eqhyperref[XVI.eq4.3.1]{4.3.1} and
\hyperrefext[IX][IX.prop2.14]{2.14} show that we may take
$F=\pi{}_{\ast}C$, where\pageoriginale $\pi:X'\rightarrow X$ is
finite, $X'$ is normal and hence regular, and $C$ is a constant sheaf
on $X'$. For a torsion abelian sheaf, we take a resolution of $F$ by
finite products of sheaves of the form $\pi_{\ast}C$ and apply the
spectral sequence of a resolution, again reducing to the case
$F=\pi_{\ast}C$. For a sheaf of ind-finite groups, recall that it is
enough to prove effaceability of
$(R^{1}g_{\cl\ast})\epsilon^{\ast}$ for the sheaf $F$, hence for a
sheaf $G$ into which $F$ can be embedded; thus we are again reduced to
the case $F=\pi_{\ast}C$. It is now clear from
(\hyperrefext[VIII][VIII.prop5.5]{5.5} and \hyperrefext[VIII][VIII.rem5.8]{5.8}) that

% END INCLUDED SGA4_English_current_buildable_source_20260906/source/16/drafts/SGA4_XVI_LINES_957_1200_English.texfrag


% BEGIN INCLUDED SGA4_English_current_buildable_source_20260906/source/16/drafts/SGA4_XVI_LINES_1201_1426_English.texfrag
we may now replace $X$ by $X'$, which proves the lemma.

To treat this last case, it is plainly enough (in view of the
cohomological-dimension results
(\sisi{IX 7.7}{\hyperrefext[IX][IX.cor5.7]{5.7}}) and
\cite{sga4.exp.XVI.bib.2}, 4.14.1) to prove that for a complete regular curve one has
\begin{itemize}
\item[$(q=0)$:] $X$ is connected and nonempty if and only if
$X_{\cl}$ is connected and nonempty.

\item[$(q=1)$:] The base-change morphism induces an equivalence between
the category of étale coverings of $X$ and the category of finite
analytic étale coverings of $X_{\cl}$.

\item[$(q=2)$:] There is an isomorphism
$\epsilon:H^2(X_{\et},\mu_n)\xrightarrow{\sim}
H^2(X_{\cl},\mu_n)$.
\end{itemize}

The first two assertions follow immediately from GAGA. For the last,
note that $H^2(X_{\et},\GG_m)=0$
(\hyperrefext[IX][IX.eq4.5]{4.5}). Likewise, if $\Oo^*$ denotes the
sheaf of invertible holomorphic functions on $X_{\cl}$, then
$H^2(X_{\cl},\Oo^*)=0$. This follows from the exact sequence
$$
0\longrightarrow\ZZ\longrightarrow\Oo\xrightarrow{\exp}
\Oo^*\longrightarrow0
$$
and\pageoriginale from the fact that
$H^2(X_{\cl},\Oo)=H^2(X_{\et},\Oo_X)=0$ by GAGA. Consequently the
Kummer exact sequence \hyperrefext[IX][IX.Thr3.2]{3.2} (respectively,
the exact sequence on $X_{\cl}$
$$
0\rightarrow\mu_n\rightarrow
\Oo^*\xrightarrow{n}\Oo^*\rightarrow0)
$$
gives
\begin{align*}
&H^1(X_{\et},\GG_m)/n\xrightarrow{\sim}
H^2(X_{\et},\mu_n)\\
(\resp\quad &H^1(X_{\cl},\Oo^*)/n\xrightarrow{\sim}
H^2(X_{\cl},\mu_n)){\quoi}.
\end{align*}
By GAGA, the canonical morphism
$$
H^1(X_{\et},\GG_m)=\Pic X\longrightarrow
\Pic X_{\cl}=H^1(X_{\cl},\Oo^*)
$$
is bijective. This proves the result, and hence proves
\hyperref[XVI.thm4.1]{4.1} when $f$ is proper.

\medskip
\noindent
\emph{The case where $F$ is constructible}. The proof proceeds by
induction on $\dim X$. Suppose the result known in dimensions $<n$,
and prove it when $\dim X\leq n$.

\begin{enonce*}{Lemma 4.5}\etiquette[4.5]{XVI.lem4.5}
We may suppose that $f:X\rightarrow S$ is a dense open immersion and
that $F$ is constant.
\end{enonce*}

\noindent
\emph{Proof}. As in the proof of \hyperref[XVI.lem1.3]{1.3}, we may
suppose that $X$ and $S$ are affine, and hence that $f$ is
quasi-projective. We then apply the argument of
\hyperref[XVI.lem4.4]{4.4} to reduce to $F=\pi_*C$, where
$\pi:X'\rightarrow X$ is finite and $C$ is constant. Replacing $X$ by
$X'$ therefore reduces us\pageoriginale to the case where $F$ is
constant. We may replace $X$ by $X_{\red}$. Let
$$
\xymatrix@=1.25cm{
  X  \ar[rr]^{h}\ar[rd]_{f} &      &\bar{X} \ar[ld]^{\bar{f}}\\
             &     S  &
}
$$
be a commutative diagram in which $\bar X$ is reduced, $\bar f$ is
projective, and $h$ is a dense open immersion. It is enough to verify
the theorem for the two morphisms $h$ and $\bar f$ (compare the proof
of \hyperref[XVI.lem4.3]{4.3}); hence it is enough to verify it for
$h$, since $\bar f$ is proper. We are thus reduced to a dense open
immersion.

\begin{enonce*}{Lemma 4.6}\etiquette[4.6]{XVI.lem4.6}
We may further suppose that $S$ is regular.
\end{enonce*}

\noindent
\emph{Proof}. Let $S'\xrightarrow{h}S$ be a ``resolution of the
singularities of $S$'', that is, a surjective, proper, birational
morphism with $S'$ regular, and let
$$
\xymatrix@C=2cm@R=1.25cm{
X'     \ar[r]^{f'} \ar[d]^{h'}&     S'\ar[d]^{h}\\
X   \ar[r]^{f}  & S
}
$$
be the Cartesian diagram obtained from $h$.

Let $F$ be a constant sheaf of sets on $X$. We have
$F\hookrightarrow G=h'_*{h'}^*F$, and it follows that, for $q=0$, it
is enough to verify \hyperref[XVI.thm4.1]{4.1} for $G$ (compare
\eqhyperref[XVI.eq4.3.1]{4.3.1}), that is, for a sheaf of the form
${h'}_*F'$, where $F'$ is constant on $X'$. We have a
commutative\pageoriginale diagram (in which the symbol ``ét'' is
suppressed)
$$
\xymatrix@C=2cm@R=1.5cm{
\varepsilon^*f_* (h'_*F') = \varepsilon^*h_*f'_*F'  \ar[d]^{d}\ar[r]^(.6){a} & h_{\cl *}
	    \varepsilon^*f'_*F' \ar[d]^{b}\\
	    f_{\cl *}\varepsilon^*(h'_*F')  \ar[r]^(.4){c} &
	    f_{\cl *} h'_{\cl *}\varepsilon^*F' = h_{\cl *}f'_{\cl *}
	    \varepsilon^*F',
}
$$
and $a$ and $c$ are bijective by the result for a proper morphism. To
prove that $d$ is bijective, it is enough to prove that $b$ is, that
is, that
$$
\epsilon^*f'_*F'\xrightarrow{\sim}
f'_{\cl *}\epsilon^*F'.
$$
We are therefore reduced to $S=S'$, with $S$ regular.

\noindent
(Remark: \sisi{We}{we} have used, without mentioning it explicitly,
the finiteness theorem for a proper morphism
\hyperrefext[XIV][XIV.thm1.1]{1.1}.)

Let $F$ be a constant torsion abelian sheaf on $X$. It is enough to
treat $F=(\ZZ/n)_X$. The morphism $h'$ is an isomorphism above a dense
open subset $j:U\rightarrow X$; let $j':U'\rightarrow X'$ be the open
subset ${h'}^{-1}(U)\xrightarrow{\sim}U$, and put $Y=X-U$. Since $U$
is dense, $\dim Y<n$, so the theorem holds for
$(f,(\ZZ/n)_Y)$. The exact sequence
$$
0\longrightarrow j_!(\ZZ/n)_U\longrightarrow
(\ZZ/n)_X\longrightarrow(\ZZ/n)_Y\longrightarrow0
$$
and the five lemma reduce us to proving the theorem for
$j_!(\ZZ/n)_U$ and $f$, and hence, by
\hyperref[XVI.lem4.2]{4.2}, to proving it for
$(hf',{j'_!}(\ZZ/n)_{U'})$. We have an exact sequence
$$
0\longrightarrow {j'_!}(\ZZ/n)_{U'}\longrightarrow
(\ZZ/n)_{X'}\longrightarrow(\ZZ/n)_{Y'}\longrightarrow0{\quoi},
$$
and\pageoriginale the theorem holds for $(hf',(\ZZ/n)_{Y'})$ by the
induction hypothesis. It is therefore enough to prove it for
$(hf',(\ZZ/n)_{X'})$. But the proper case gives the theorem for
$(h,\cdot)$, and the Leray spectral sequence for the pair of morphisms
$h,f'$ shows that it is enough to prove it for $(f',\ZZ/n)$. This
proves the lemma in the abelian case.

Finally, suppose that $F$ is a constant sheaf of finite groups. Recall
that it is enough to prove that the functor
$(R^1f_{\cl *})\epsilon^*(\mathord\cdot)$ is effaceable for $F$, hence
for a sheaf $G$ into which $F$ embeds, and therefore for
$G={h'}_*{h'}^*F$, that is, for a sheaf of the form ${h'}_*F'$, where
$F'$ is constant on $X'$. We have
$$
\epsilon^*({h'}_*F')\xrightarrow{\sim}
h'_{\cl *}(\epsilon^*F')
$$
and
$$
0\longrightarrow
(R^1f_{\cl *})h'_{\cl *}\epsilon^*(\mathord\cdot)\longrightarrow
(R^1(f_{\cl}h'_{\cl})_*)\epsilon^*(\mathord\cdot)
=(R^1(h_{\cl}f'_{\cl})_*)\epsilon^*(\mathord\cdot){\quoi}.
$$
Thus it is enough to prove that the last functor is effaceable for a
constant sheaf. But the exact sequence
$$
0\longrightarrow (R^1h_{\cl *})f'_{\cl *}
\epsilon^*(\mathord\cdot)\longrightarrow
(R^1(h_{\cl}f'_{\cl})_*)\epsilon^*(\mathord\cdot)\longrightarrow
h_{\cl *}(R^1f'_{\cl *})\epsilon^*(\mathord\cdot){\quoi}
$$
shows that the first term is effaceable, for example by the Godement
resolution (\hyperrefext[XII][XII.prop3.3]{3.3}), because $h$ is
proper. It is therefore enough to prove that the right-hand term is
effaced by the Godement resolution, hence that
$(R^1f'_{\cl *})\epsilon^*(\mathord\cdot)$ is so effaced, and thus
that the theorem holds for $(f',F')$, where $F'$ is constant. This
proves the lemma.

\noindent
\emph{End of the proof}.\pageoriginale We now suppose that $F$ is a
constant sheaf and that $f:X\rightarrow S$ is a dense open immersion,
with $S$ regular. The assertion for sheaves of sets now follows
immediately from \hyperref[XVI.prop3.2]{3.2} and the analogous trivial
lemma for the classical topology\sisi{---indeed}{. Indeed}, these show
that $\epsilon^*f_*F$ and $f_{\cl *}\epsilon^*F$ are constant with the
same value and hence are isomorphic under $\varphi$.

For $q>0$, we first reduce further to the case $X=S-Y$ with $Y$
regular, that is, to the case in which $(Y,S)$ is a smooth
$\Spec\CC$-pair. Indeed, write $X=S-C_1$, where $C_1$ is given its
induced reduced structure, and define recursively

$Y_\nu=$ the dense open subset of regular points of $C_\nu$,

$C_{\nu+1}=C_\nu-Y_\nu$, with its induced reduced structure.

\noindent
Since $Y_\nu$ is dense in $C_\nu$, we have
$\dim C_{\nu+1}<\dim C_\nu$, and hence a sequence
$$
X=X_1\subset X_2=S-C_2\subset X_3=S-C_3\subset\cdots\subset X_r=S{\quoi},
$$
where the inclusions $i_\nu:X_\nu\rightarrow X_{\nu+1}$ are open
immersions and
$$
X_\nu=X_{\nu+1}-Y_\nu,\qquad
Y_\nu\ \mbox{closed in $X_{\nu+1}$ and regular}{\quoi}.
$$
The direct image $i_{\nu *}F$ of a constant sheaf on $X$ is constant
(\hyperref[XVI.prop3.2]{3.2}), and the $R^qi_{\nu *}F$ are
constructible and concentrated on the subvarieties of $Y$ of
dimension at most $n-1$ (\hyperref[XVI.cor3.4]{3.4},
\hyperref[XVI.cor3.6]{3.6}, \hyperref[XVI.thm3.7]{3.7}) for the values
of $q$ under consideration. The induction hypothesis and the Leray
spectral sequence (respectively, the exact sequence
\hyperrefext[XII][XII.prop3.2]{3.2}) therefore reduce us\pageoriginale
to the case of $i_\nu$, which proves the asserted reduction.

For a smooth pair $(Y,S)$, however, the values of $(R^qf_*)F$ have
been computed explicitly for the étale topology
(\hyperref[XVI.cor3.4]{3.4}, \hyperref[XVI.cor3.6]{3.6},
\hyperref[XVI.thm3.7]{3.7}); an analogous and easy computation, which
we leave to the reader, gives the corresponding results for the
classical topology. Direct comparison proves the theorem for $f$.


% END INCLUDED SGA4_English_current_buildable_source_20260906/source/16/drafts/SGA4_XVI_LINES_1201_1426_English.texfrag


% BEGIN INCLUDED SGA4_English_current_buildable_source_20260906/source/16/drafts/SGA4_XVI_LINES_1427_1560_English.texfrag
\section{The finiteness theorem for algebraic schemes in
characteristic zero}\etiquette[5]{XVI.sec5}

\begin{enonce*}{Theorem 5.1}\etiquette[5.1]{XVI.thm5.1}
Let $k$ be a field of characteristic zero, and let $f:X\rightarrow S$
be a morphism of finite type between schemes locally of finite type
over $k$. Let $F$ be a constructible sheaf of sets (\resp of finite
groups, \resp a torsion abelian sheaf) on $X$. Then the
$R^qf_*F$ are constructible for $q=0$ (\resp for $q=0,1$, \resp for
every $q$).
\end{enonce*}

The hypothesis that $k$ have characteristic zero is imposed because
we shall use the resolution-of-singularities theorem \cite{sga4.exp.XVI.bib.1}. The
proof is also valid in characteristic $p>0$ if resolution of
singularities is assumed, \emph{provided that the orders of $F$ are
prime to $p$}. Thus in nonzero characteristic one obtains a result for
$\dim X\leq2$ by applying Abhyankar's results \cite{sga4.exp.XVI.bib.3}. In addition to
the ever-present\pageoriginale proper base-change theorem, another
ingredient of the proof is the ``absolute cohomological purity
theorem'' \hyperref[XVI.cor3.9]{3.9}\SGAArchiveOnly{\nde{The frozen source attributes
this later development to work of Deligne and Gabber but supplies no
bibliographic target.}}.

Recall that the theorem has already been proved for a proper morphism
and arbitrary $S$ (\hyperrefext[XIV][XIV.thm1.1]{1.1}). Moreover, in
Exposé XIX we shall obtain (again in characteristic zero) the finiteness
theorem under the much weaker hypotheses that $S$ be an
\emph{excellent} scheme and $f$ be of finite type, by proving for such
schemes the absolute-purity theorem in the form indicated in
\hyperref[XVI.rems3.10]{3.10}.

In nonzero characteristic and dimension at least $3$, for the time
being we have only the following easy consequence of
\hyperrefext[XI][XI.prop3.3]{3.3}:

\begin{enonce*}{Theorem 5.2}\etiquette[5.2]{XVI.thm5.2}
Let $X$ be a smooth algebraic scheme over $\Spec k$, where $k$ is a
separably closed field, and let $F$ be a locally constant sheaf of
finite abelian groups (\resp finite groups), of orders prime to the
characteristic of $k$. Then $H^q(X,F)$ is a finite group (\resp a
finite pointed set) for every $q$ (\resp for $q=0,1$).
\end{enonce*}

\noindent
\emph{Proof of} \hyperref[XVI.thm5.2]{5.2}. We may suppose that $k$ is
algebraically closed (\hyperrefext[VIII][VIII.thm1.1]{1.1}). We argue
by induction on $n=\dim X$. By \hyperrefext[XI][XI.prop3.3]{3.3}, there
is a hypercovering of $X$ by open subsets that are ``good
neighborhoods''. The spectral sequence of a hypercovering immediately
reduces us to the case in which $X$ is a good neighborhood and hence
admits an elementary fibration (\hyperrefext[XI][XI.def3.1]{3.1})
$$
\xymatrix@=1.5cm{
X \ar[rr]^{j} \ar[rd]_{f}&         &
\bar{X}\ar[ld]^{\bar{f}} \\
             &      S  & .
}
$$
By\pageoriginale \hyperref[XVI.prop3.2]{3.2},
\hyperref[XVI.cor3.6]{3.6}, and \hyperref[XVI.thm3.7]{3.7}, the sheaf
$j_*F$ is locally constant on $\bar X$, $(R^1j_*)F$ is locally
constant on $Y=\bar X-X$, and $(R^qj_*)F=0$ for $q>1$. It follows from
\hyperref[XVI.thm2.1]{2.1} that
$(R^p\bar f_*)(R^qj_*)F$ is a locally constant sheaf on $S$ for the
values of $q$ under consideration. The result follows from the
induction hypothesis and the Leray spectral sequence (\resp the exact
sequence \hyperrefext[XII][XII.prop3.2]{3.2}).

The proof of \hyperref[XVI.thm5.1]{5.1} is very close to the proof of
\hyperref[XVI.thm4.1]{4.1} in the constructible case. We merely
indicate the main steps. Argue by induction on
$\dim X=n$. First reduce to $X$ and $S$ affine, and hence $f$
quasi-projective, by applying the method used in the proof of \hyperref[XVI.lem1.3]{1.3}.
Next apply \sisi{IX 2.15}{\hyperrefext[IX][IX.prop2.14]{2.14}} and
\hyperrefext[VIII][VIII.prop5.5]{5.5},
\hyperrefext[VIII][VIII.rem5.8]{5.8} to reduce to the case in which
$F$ is constant. There is a commutative diagram
$$
\xymatrix@=1.5cm{
X \ar[rr]^{h} \ar[rd]_{f}&         & \bar{X}\ar[ld]^{\bar{f}}\\
             &      S  & ,
}
$$
in which $h$ is a dense open immersion and $\bar f$ is proper, so the
theorem holds for $\bar f$
(\hyperrefext[XIV][XIV.thm1.1]{1.1}). It is therefore enough to prove
the theorem for $(h,F)$, that is, for a dense open
immersion\pageoriginale and a constant sheaf. Let
$h:S'\rightarrow S$ be a resolution of the singularities of $S$, and
let
$$
\xymatrix@C=1.75cm@R=1cm{
X  \ar[d]^{f} & X' \ar[l]_{h'} \ar[d]^{f'}\\
S          & S' \ar[l]_{h}
}
$$
be the Cartesian diagram obtained from $h$. We reduce to proving the
theorem for $f'$, and hence to the case in which $S=S'$ is regular.
The inductive method used at the end of the proof of
\hyperref[XVI.thm4.1]{4.1} reduces us to the case in which $Y=S-X$ is
nonsingular, that is, $(Y,S)$ is a smooth pair. The proof is completed
by applying \hyperref[XVI.prop3.2]{3.2},
\hyperref[XVI.cor3.4]{3.4}, \hyperref[XVI.cor3.6]{3.6}, and
\hyperref[XVI.thm3.7]{3.7}.

\begin{enonce*}[remark]{Remark 5.3}\etiquette[5.3]{XVI.rem5.3}
The proof of \hyperref[XVI.thm5.1]{5.1} just outlined is also valid
when one is given a left Noetherian ring $A$ annihilated by an integer
$n>0$ (which must be assumed prime to the characteristic when the
latter is not assumed to be zero; compare the remarks following
\hyperref[XVI.thm5.1]{5.1}) and considers constructible sheaves of left
$A$-modules.
\end{enonce*}

\begin{thebibliography}{99}
\bibitem{sga4.exp.XVI.bib.1} Hironaka, H. Resolution of singularities of an
algebraic variety over a field of characteristic zero, Annals of
Math., vol. 79 (1964), p. 109.

\bibitem{sga4.exp.XVI.bib.2} Godement,\pageoriginale P. Théorie des
\sisi{Faisceaux}{faisceaux}, Paris, 1958\sisi{}{.}

\bibitem{sga4.exp.XVI.bib.3} Abhyankar, S. Local uniformization on algebraic
surfaces over ground fields of characteristic $p\neq0$, Ann. of
Math., vol. 63 (1956), p. 491.

\bibitem{sga4.exp.XVI.bib.4} J.-P. Serre, \sisi{Géométrie Algébrique et Géométrie
Analytique}{Géométrie algébrique et géométrie analytique},
Annales Inst. Fourier 1956, pp. 1--42 (\sisi{Cited}{cited} GAGA).
\end{thebibliography}

% XII 4.1 no ref. 78
% XII 4.1.2 no ref. 992
% 4.11.1 no ref. 1040
% XI 7.7 no ref. 1210
% IX 4.5 is an equation ref.
% IX 2.15 no ref. 1504

\ifx\danslelivre\undefined
\end{document}
\fi

% END INCLUDED SGA4_English_current_buildable_source_20260906/source/16/drafts/SGA4_XVI_LINES_1427_1560_English.texfrag


% END INCLUDED SGA4_English_current_buildable_source_20260906/source/16/16_English.tex


\makeatletter
\def\input@path{{../17/}}
\makeatother
% BEGIN INCLUDED SGA4_English_current_buildable_source_20260906/source/17/17_English.tex
\ifx\danslelivre\undefined
% Release engine: XeLaTeX (the LuaTeX compatibility layer breaks Xy-pic output).
\documentclass[10pt]{../fichiers-configuration-latex/smfart}
% BEGIN INCLUDED SGA4_English_current_buildable_source_20260906/source/fichiers-configuration-latex/paquets.tex
\usepackage{polyglossia}
%\usepackage[francais]{babel}
\usepackage{fontspec}
\usepackage[math-style=ISO,bold-style=upright,sans-style=upright]{unicode-math}
%\usepackage[mathletters]{ucs}
%\usepackage[utf8]{inputenc}
%\usepackage[T1]{fontenc}
\usepackage[unicode,hypertexnames=false,colorlinks=true,linkcolor=blue,citecolor=Sepia,urlcolor=magenta]{hyperref}
\setmainlanguage{french}
\setmainfont[Ligatures=TeX]{Linux Libertine O}
% variante possible : TeXGyrePagella ?
\setmathfont{XITS Math}

% maths
\usepackage{amsmath}
%\usepackage{amssymb}
%\usepackage{mathrsfs}
%\usepackage{stmaryrd}
%\usepackage{euscript}

% couleurs
\usepackage[usenames,dvipsnames]{xcolor}

% graphiques
\usepackage{graphicx}
\usepackage[all]{xy}
%\usepackage{tikz}
%\usetikzlibrary{matrix,arrows,calc}
%\tikzset{column sep=3em,row sep=2em,text height=1.5ex, text depth=0.25ex,>=angle 90}
%\tikzset{description/.style={fill=white,inner sep=1pt}}

% hyperliens, 'etiquettes
\usepackage{xr-hyper}
\usepackage{../fichiers-configuration-latex/smfthm}
%\ifx\cacheretiquettes\undefined
%\usepackage[inline]{showlabels} \else \fi % pour rendre visibles les 'etiquettes
%\usepackage{../fichiers-configuration-latex/makeidx} % pour faire un index
%\usepackage[pagebackref,unicode]{hyperref}  % pour voir o`u sont cit'ees les r'ef'erences
% hyperref is loaded once above with the cumulative-reader destination policy.
%\hypersetup{breaklinks=true} %permet le retour `a la ligne dans les liens trop longs. Incompatible avec pdflatex
%\renewcommand*{\backref}[1]{{\small [cit\'e page(s)~#1.]}}

\usepackage{xspace}
\usepackage{pifont} % pour expose XVII

% END INCLUDED SGA4_English_current_buildable_source_20260906/source/fichiers-configuration-latex/paquets.tex

% BEGIN INCLUDED SGA4_English_current_buildable_source_20260906/source/fichiers-configuration-latex/numerotation.tex
\renewcommand{\thefootnote}{\roman{footnote}} % pour mettre les notes en bas de page sous la forme (ii) et non (2), qui ressemble trop (?) a un exposant.

\newcounter{pageoriginal}
\newcommand{\pageoriginale}{\refstepcounter{pageoriginal}\marginpar{\footnotesize\textbf{\thepageoriginal}}\xspace}
\let\pageoriginaled\pageoriginale

\usepackage{ifthen}
\newboolean{original}
\def\sisi{\ifthenelse{\boolean{original}}}
\setboolean{original}{false}

\newcounter{cnote}
\def\nde#1{\ifthenelse{\boolean{original}}{}{\stepcounter{cnote}\def\thefootnote{\arabic{cnote}}\protect\footnote[\value{cnote}]{N.D.E.: #1}}}
\def\ndemark{\ifthenelse{\boolean{original}}{}{\stepcounter{cnote}\def\thefootnote{\arabic{cnote}}\protect\footnotemark[\value{cnote}]}}
\def\ndetext#1{\ifthenelse{\boolean{original}}{}{\def\thefootnote{\arabic{cnote}}\protect\footnotetext[\value{cnote}]{N.D.E.: #1}}}

% END INCLUDED SGA4_English_current_buildable_source_20260906/source/fichiers-configuration-latex/numerotation.tex

% BEGIN INCLUDED SGA4_English_current_buildable_source_20260906/source/fichiers-configuration-latex/macros.tex
\setmainlanguage{french}
\setotherlanguage{english}


%% Fixed, engine-independent release identifier.  The inherited dynamic Git
%% probe used Lua io.popen, exposed build-state-dependent output, and printed
%% a visible question mark outside a Git checkout.
\def\showgitstatus{English release r2, 2026-07-23}


% macros pour LaTeX
\renewcommand{\mathsf}{}

% Police différente pour U+2113 SCRIPT SMALL L
\setmathfont[range="2113]{latinmodern-math.otf}

% Give every inherited semantic label its own stable PDF destination instead
% of inheriting whichever automatic equation/page anchor happened to be live.
% The legacy label remains as a compatibility alias used by existing edges.
\makeatletter
\def\etiquette[#1]#2{%
  \Hy@raisedlink{\hyper@anchorstart{sga4.semantic.#2}\hyper@anchorend}%
  \def\@currentHref{sga4.semantic.#2}%
  \def\@currentlabel{#1}%
  \label{#2}%
}
\makeatother
\def\eqhyperref[#1]#2{(\hyperref[#1]{#2})}
\def\pagehyperref[#1]#2{(page \hyperref[#1]{#2})}
\def\hyperrefext[#1][#2]#3{\hyperref[#2]{#1 #3}}
\def\eqhyperrefext[#1][#2]#3{\hyperref[#2]{#1 #3}}     %cf. exposes 17
\def\neweqhyperrefext[#1][#2]#3{\hyperref[#2]{#1 #3}}  %cf. exposes 17

\def\shat#1{\protect\widehat{#1}}
\def\stilde#1{\protect\widetilde{#1}}
%%%%%%%%%%%%%%%%%%


\def\etcn{etc.\xspace}
\def\etco{etc...\xspace}
\def\etc{\sisi{\etco}{\etcn}}
\def\ldotsn{\sisi{...}{\etcn}}

\def\cf{{cf.}\kern.3em}
\def\Cf{{Cf.}\kern.3em}
\def\ie{{i.e.}\kern.3em}
\def\iev{{i.e.},\xspace}
\def\loccito{loc.\kern2pt cit.\xspace}
\def\loccitn{\textit{loc.\kern2pt cit.}\xspace}
\def\loccit{\sisi{\loccito}{\loccitn}}
\def\EGA{\textup{EGA}\kern.2em}
\def\SGA{\textup{SGA}\kern.3em}

\def\quoi{}

\def\quav{\quad,\quad}
\def\quaet{\quad \text{et}\quad}
\def\fibreproduct#1#2#3{#1 ×_{#3} #2}
\let\fprod\fibreproduct
\def\ccheck#1#2{\mbox{${\check{\check{\rm #1}}}^{\lower2pt\hbox{$\scriptstyle #2$}}$}}
\def\sheaf#1{\protect\mathcal{#1}}
\def\vieuxsheaf#1{{#1}} % cf. problème avec sheaf{a\Phi_i} par exemple...
\DeclareMathOperator{\Hom}{Hom}
\DeclareMathOperator{\can}{can}
\DeclareMathOperator{\Ind}{Ind}
\DeclareMathOperator{\ind}{ind}
%\DeclareMathOperator{\Hhomtop}{\mathcal{H}\!\mathit{omtop}}
\DeclareMathOperator{\Hhomtop}{\mathbf{Homtop}}
\DeclareMathOperator{\Point}{Point}
\DeclareMathOperator{\Pro}{Pro}
\DeclareMathOperator{\Top}{Top}
\DeclareMathOperator{\Mod}{Mod}
\DeclareMathOperator{\Cart}{Cart}
%\DeclareMathOperator{\Mmorsite}{\mathcal{M}\!\mathit{orsite}}
\DeclareMathOperator{\Mmorsite}{\mathbf{Morsite}}
\DeclareMathOperator{\loc}{loc}
\DeclareMathOperator{\Loc}{Loc}
\DeclareMathOperator{\Pic}{Pic}
\DeclareMathOperator{\Et}{Et}
\DeclareMathOperator{\card}{card}
\DeclareMathOperator{\Ens}{Ens}
\DeclareMathOperator{\h}{\underline{H}}
\DeclareMathOperator{\Fonct}{Fonct}
\DeclareMathOperator{\Ker}{Ker}
\DeclareMathOperator{\but}{but}
\DeclareMathOperator{\Ob}{Ob}
\DeclareMathOperator{\ob}{ob}
\DeclareMathOperator{\Lib}{Lib}
\DeclareMathOperator{\esj}{esj}
\DeclareMathOperator{\id}{id}
\DeclareMathOperator{\Id}{Id}
\DeclareMathOperator{\ad}{ad}
\DeclareMathOperator{\Iim}{Im}
\DeclareMathOperator{\pr}{pr}
\DeclareMathOperator{\Cat}{Cat}
\DeclareMathOperator{\car}{car}
\DeclareMathOperator{\SysInd}{SysInd}
\DeclareMathOperator{\Sup}{Sup}
\DeclareMathOperator{\Coker}{Coker}
\DeclareMathOperator{\Filt}{Filt}
\DeclareMathOperator{\acc}{acc}
\DeclareMathOperator{\sur}{sur}
\DeclareMathOperator{\cat}{cat}
\DeclareMathOperator{\Ord}{Ord}
\DeclareMathOperator{\Gr}{Gr}
\DeclareMathOperator{\Cov}{Cov}
\DeclareMathOperator{\Caf}{Caf}
\DeclareMathOperator{\coker}{coker}
\DeclareMathOperator{\coim}{coim}
\DeclareMathOperator{\im}{im}
\DeclareMathOperator{\Ouv}{Ouv}
\DeclareMathOperator{\diag}{diag}
\DeclareMathOperator{\Esp}{Esp}
\DeclareMathOperator{\TOP}{TOP}
\DeclareMathOperator{\Unif}{Unif}
\DeclareMathOperator{\Iind}{\mathbf{Ind}}
\DeclareMathOperator{\Hhomcart}{\mathbf{Homcart}}
\DeclareMathOperator{\Hhomtopess}{\mathbf{Homtopess}}
\DeclareMathOperator{\Ppoints}{\mathbf{Points}}
\DeclareMathOperator{\Ppoint}{\mathbf{Point}}
\DeclareMathOperator{\Ttors}{\mathbf{Tors}}
\DeclareMathOperator{\Ssex}{\mathbf{Sex}}
\DeclareMathOperator{\Ffib}{\mathbf{Fib}}
\DeclareMathOperator{\Vvois}{\mathbf{Vois}}
\DeclareMathOperator{\Ppointess}{\mathbf{Pointess}}
\DeclareMathOperator{\scrHom}{\mathbf{Hom}}
\DeclareMathOperator{\sheafPic}{\mathbf{Pic}}
\DeclareMathOperator{\sheafAut}{\mathbf{Aut}}
\DeclareMathOperator{\Hhomtoplocan}{\mathbf{Homtoplocan}}
\DeclareMathOperator{\Hhomtopan}{\mathbf{Homtopan}}
\DeclareMathOperator{\Pptgeom}{\mathbf{Ptgéom}}
\DeclareMathOperator{\Morsite}{Morsite}
\DeclareMathOperator{\Fib}{Fib}
\DeclareMathOperator{\Fais}{Fais}
\DeclareMathOperator{\kar}{kar}
\DeclareMathOperator{\Groupes}{Groupes}
\DeclareMathOperator{\Restr}{Restr}
\DeclareMathOperator{\Isom}{Isom}
\DeclareMathOperator{\Prol}{Prol}
\DeclareMathOperator{\Res}{Res}
\DeclareMathOperator{\Tors}{Tors}
\DeclareMathOperator{\esp}{esp}
\DeclareMathOperator{\es}{es}
\DeclareMathOperator{\od}{od}
\DeclareMathOperator{\gr}{gr}
\DeclareMathOperator{\sob}{sob}
\DeclareMathOperator{\dfn}{dfn}
\DeclareMathOperator{\Sch}{Sch}
\DeclareMathOperator{\ouv}{ouv}
\DeclareMathOperator{\cons}{cons}
\DeclareMathOperator{\Nerf}{Nerf}
\DeclareMathOperator{\Points}{Points}
\DeclareMathOperator{\Espsob}{Espsob}
\DeclareMathOperator{\Homtop}{Homtop}
\DeclareMathOperator{\Pt}{Pt}
\DeclareMathOperator{\Quot}{Quot}
\DeclareMathOperator{\glob}{glob}
\DeclareMathOperator{\Homfib}{Homfib}
\DeclareMathOperator{\Qouv}{Qouv}
\DeclareMathOperator{\Quotouv}{Quotouv}
\DeclareMathOperator{\ens}{ens}
\DeclareMathOperator{\ab}{ab}
\DeclareMathOperator{\an}{an}
\DeclareMathOperator{\fs}{fs}
\DeclareMathOperator{\red}{red}
\DeclareMathOperator{\Ab}{Ab}
\DeclareMathOperator{\cosq}{cosq}
\DeclareMathOperator{\Ann}{Ann}
\DeclareMathOperator{\Zar}{Zar}


%%%%%%%%%%%%
%%%%% SGA4-2
\DeclareMathOperator{\Sex}{Sex}
\DeclareMathOperator{\vois}{vois}
\DeclareMathOperator{\Vois}{Vois}
\DeclareMathOperator{\Ext}{Ext}
\DeclareMathOperator{\ext}{ext}
\DeclareMathOperator{\Fer}{Fer}
\DeclareMathOperator{\Add}{Add}
\DeclareMathOperator{\Cocart}{Cocart}
\DeclareMathOperator{\cosk}{cosk}
\DeclareMathOperator{\plat}{plat}

\DeclareMathOperator{\faisc}{faisc}
\DeclareMathOperator{\sk}{sk}
\DeclareMathOperator{\ch}{ch}

\def\HH{\mathbf{H}}
\def\qc{\mathrm{qc}}
\def\et{\mathrm{et}}
\def\coh{\mathrm{coh}}

\def\Hvr#1{%
\begin{array}[b]{l}
{\kern2pt\scriptstyle r}\\[-8pt]
{\kern1pt\scriptstyle \vee}\\[-5pt]
H^{#1}
\end{array}}

\def\Hvv#1{%
\begin{array}[b]{l}
\kern2pt\rotatebox{90}{${\scriptscriptstyle\ll}$}\\[-6pt]
H^{#1}
\end{array}}

\def\Hvi#1{%
\begin{array}[b]{l}
{\kern0.2pt\scriptscriptstyle \infty}\\[-8pt]
{\kern1pt\scriptstyle \vee}\\[-5pt]
H^{#1}
\end{array}}

\def\Hv#1{%
\begin{array}[b]{l}
{\kern1pt\scriptstyle \vee}\\[-5pt]
H^{#1}
\end{array}}

\DeclareMathOperator{\cart}{cart}
\DeclareMathOperator{\Ffaisc}{\mathbf{Faisc}}


\def\Ev#1{%
\begin{array}[b]{l}
{\kern1pt\scriptstyle \vee}\\[-5pt]
E'_{#1}
\end{array}}

\DeclareMathOperator{\Cont}{Cont}
\DeclareMathOperator{\Ccont}{\mathscr{C}\!\mathit{ont}}
\DeclareMathOperator{\Site}{Site}
\newcommand\MM{\mathbf{M}}
\def\limtop#1{{\displaystyle{\varprojlim_{#1}\!\mbox{\rm top }}}}
\def\Limtop{\mathop{\varprojlim\!\mbox{\rm top}}}
\DeclareMathOperator{\spec}{spec}
\def\ss{\scriptstyle{\mathscr{S}}}
\DeclareMathOperator{\Gal}{Gal}
\def\smallm{\scriptstyle{\mathscr{M}}}
\DeclareMathOperator{\End}{End}
\DeclareMathOperator{\ass}{ass}
\DeclareMathOperator{\zar}{zar}
\DeclareMathOperator{\Ogg}{Ogg}
\DeclareMathOperator{\res}{res}
\DeclareMathOperator{\tr}{tr}
\DeclareMathOperator{\Supp}{Supp}
\DeclareMathOperator{\codim}{codim}
\DeclareMathOperator{\rk}{rk}
\DeclareMathOperator{\cur}{cur}
\DeclareMathOperator{\cdqc}{cdqc}
\newcommand\DD{\mathbf{D}}
\DeclareMathOperator{\cl}{cl}
\def\JJ{\mathbf{J}}
\def\FF{\mathbf{F}}
\def\BB{\mathbf{B}}
\DeclareMathOperator{\Alget}{Alget}
\DeclareMathOperator{\Ass}{Ass}
\DeclareMathOperator{\supp}{supp}
\DeclareMathOperator{\rad}{rad}
\DeclareMathOperator{\prof}{prof}
\DeclareMathOperator{\cint}{int}
\DeclareMathOperator{\sic}{sic}
\DeclareMathOperator{\Tor}{Tor}
\DeclareMathOperator{\tors}{tors}
\DeclareMathOperator{\Sym}{Sym}
\DeclareMathOperator{\parf}{parf}
\DeclareMathOperator{\tor}{tor}
\DeclareMathOperator{\defn}{defn}
\DeclareMathOperator{\Div}{Div}
\DeclareMathOperator{\Pol}{Pol}
\DeclareMathOperator{\Fex}{Fex}
\DeclareMathOperator{\Lie}{Lie}
\DeclareMathOperator{\tot}{tot}
\DeclareMathOperator{\Aut}{Aut}
\DeclareMathOperator{\HOM}{HOM}
\DeclareMathOperator{\pch}{pch}
\DeclareMathOperator{\Ch}{Ch}
\DeclareMathOperator{\EXT}{EXT}
\DeclareMathOperator{\PIC}{PIC}
\DeclareMathOperator{\Rf}{Rf}
\DeclareMathOperator{\inj}{inj}
\newcommand\VV{\mathbf{V}}

\def\Xv{%
\begin{array}[b]{l}
{\kern1pt\scriptstyle \vee}\\[-5pt]
X
\end{array}}

%%%%%%%
\newcommand\TT{\mathbf{T}}
\newcommand\XX{\mathbf{X}}
\newcommand\EE{\mathbf{E}}
%%%%%%

\def\qedsymbol{\textup{C.Q.F.D.}}
\let\CQFD\qed
\let\cqfd\qed

%%%%%\bigdot
\font\bigsymb=cmsy10 at 4pt
\def\bigdot{{\kern1.2pt\raise 1.5pt\hbox{\bigsymb\char15}}}
%%%%%

\def\Lright{\mathop{\vcenter{\hbox{\hbox{$\longrightarrow$}\lower8pt \hbox{\kern-16pt\hbox{$\longrightarrow$}}}}}}

\def\Aa{{\mathcal A}}
\def\Bb{{\mathcal B}}
\def\Cc{{\mathcal C}}
\def\Dd{{\mathcal D}}
\def\Ee{{\mathcal E}}
\def\Ff{{\mathcal F}}
\def\Gg{{\mathcal G}}
\def\Hh{{\mathcal H}}
\def\Ii{{\mathcal I}}
\def\Jj{{\mathcal J}}
\def\Kk{{\mathcal K}}
\def\Ll{{\mathcal L}}
\def\Mm{{\mathcal M}}
\def\Nn{{\mathcal N}}
\def\Oo{{\mathcal O}}
\def\Pp{{\mathcal P}}
\def\Qq{{\mathcal Q}}
\def\Rr{{\mathcal R}}
\def\Ss{{\mathcal S}}
\def\Tt{{\mathcal T}}
\def\Uu{{\mathcal U}}
\def\Vv{{\mathcal V}}
\def\Ww{{\mathcal W}}
\def\Xx{{\mathcal X}}
\def\Yy{{\mathcal Y}}
\def\Zz{{\mathcal Z}}

\newcommand{\CC}{\mathbf{C}}
\newcommand{\GG}{\mathbf{G}}
\newcommand{\KK}{\mathbf{K}}
\newcommand{\LL}{\mathbf{L}}
\newcommand{\NN}{\mathbf{N}}
\newcommand{\PP}{\mathbf{P}}
\newcommand\QQ{\mathbf{Q}}
\newcommand\RR{\mathbf{R}}
\newcommand\ZZ{\mathbf{Z}}

\def\to{\mathchoice{\longrightarrow}{\rightarrow}{\rightarrow}{\rightarrow}}
\def\from{\mathchoice{\longleftarrow}{\leftarrow}{\leftarrow}{\leftarrow}}
\def\mto{\mathchoice{\longmapsto}{\mapsto}{\mapsto}{\mapsto}}
\def\mfrom{\mathrel{\from\joinrel\mapstochar\,}}
\def\hto{\mathrel{\lhook\joinrel\to}}
\def\hfrom{\mathrel{\from\joinrel\rhook}}

\def\implies{\mathchoice{\Longrightarrow}{\Rightarrow}{\Rightarrow}{\Rightarrow}}
\def\implied{\mathchoice{\Longleftarrow}{\Leftarrow}{\Leftarrow}{\Leftarrow}}
\def\iff{\mathchoice{\Longleftrightarrow}{\Leftrightarrow}{\Leftrightarrow}{\Leftrightarrow}}

\def\To#1{\mathchoice{\xrightarrow{\labelstyle~~#1~}}{\stackrel{#1}{\longrightarrow}}{}{}}
\def\Mto#1{\mathrel{\mapstochar\To{#1}}}
\def\From#1{\mathchoice{\xleftarrow{\labelstyle~~#1~}}{\stackrel{#1}{\longleftarrow}}{}{}}

\newcommand{\isomto}{\overset{\sim}{\longrightarrow}}
\newcommand{\isomfrom}{\overset{\sim}{\longleftarrow}}

\newcommand{\IFF}{\unskip~$\iff$~\ignorespaces}
\newcommand{\THEN}{\unskip~$\implies$~\ignorespaces}

\let\leq\leqslant
\let\geq\geqslant
\let\le\leqslant
\let\ge\geqslant

\def\cf{{cf.}\kern.3em}
\def\Cf{{Cf.}\kern.3em}
\def\ie{{i.e.}\kern.3em}
\def\iev{{i.e.},\xspace}
%\def\loccito{loc.\kern2pt cit.\xspace}
%\def\loccitn{\textit{loc.\kern2pt cit.}\xspace}
%\def\loccit{\sisi{\loccito}{\loccitn}}
\def\EGA{\textup{EGA}\kern.2em}
\def\SGA{\textup{SGA}\kern.3em}

\let\oldS\S
\def\S{\oldS\kern .15em}
\def\ptbl{.\kern .2em}
\def\resp{resp.\kern.3em}
\def\Exp{Exp.\kern.3em}

\let\bar\overline
\let\hat\widehat
\let\tilde\widetilde
\let\dpl\displaystyle
\let\ldots\dots
\let\emptyset\varnothing
\let\skpt\skippointrait
\let\epsilon\varepsilon

\def\sfrac#1#2{{#1}/{#2}}

\DeclareMathOperator{\cd}{cd}
%\DeclareMathOperator{\Hhom}{\mathcal{H}\!\mathit{om}}
\DeclareMathOperator{\Hhom}{\mathbf{Hom}}
%\DeclareMathOperator{\Eext}{\mathcal{E}\!\mathit{xt}}
\DeclareMathOperator{\Eext}{\mathbf{Ext}}
\DeclareMathOperator{\Spec}{Spec}
\DeclareMathOperator{\Tr}{Tr}

\let\emptyset=\varnothing

%%%%%%%%%%%%%%%
%%% expose XVII
%%%%%%%%%%%%%%%

\makeatletter
\let\old@optexpose\optexpose
\newcommand{\empty@optexpose}[1]{}
\newcommand{\paren@optexpose}[1]{\old@optexpose{#1}(}
\newcommand{\refnoexpose}[1]{\let\optexpose\empty@optexpose\ref{#1}\let\optexpose\old@optexpose}%\refnoexpose{X.cor5.2} engendre "5.2" au lieu de "X 5.2"
\newcommand{\neweqref}[1]{\let\optexpose\paren@optexpose\textup{\ref{#1})}\let\optexpose\old@optexpose}%\neweqref{XII.eq4.1} engendre "XII (4.1)" au lieu de "(XII 4.1)"
\makeatother
\newcommand{\refvb}[2][VI${}_{\text{B}}$ ]{\sisi{#1\refnoexpose{#2}}{\ref{#2}}}

\newcommand{\leftexp}[2]{{\vphantom{#2}}^{#1}{#2}}
\newcommand{\AAA}{\mathbf{A}}
\newcommand{\Sf}{\mathfrak{S}}
\DeclareMathOperator*{\quotprojlim}{\text{``$\varprojlim$''}}
\DeclareMathOperator*{\quotinjlim}{\text{``$\varinjlim$''}}
\DeclareMathOperator{\Fl}{F\sisi{\ell}{l}}
\DeclareMathOperator{\Pl}{P\sisi{\ell}{l}}
%\DeclareMathOperator{\cl}{c\sisi{\ell}{l}}
\DeclareMathOperator{\tdf}{tdf}
\DeclareMathOperator{\rg}{rg}
\newcommand{\iInd}{\sisi{I}{i}nd}
\newcommand{\iind}{ind}
\newcommand{\ppro}{pro}
\newcommand{\tTor}{\sisi{T}{t}or}
\newcommand{\ttor}{tor}
\newcommand{\chap}{Chap.}
\newcommand{\Ffex}{\Ff\mathit{ex}}
\newcommand{\sheafk}{\underline{k}}
\newcommand{\sheafn}{\underline{n}}
\newcommand{\sheafGamma}{\underline{\Gamma}}

% Repeated references to an already typeset footnote must not use
% \footnotemark without a matching \footnotetext: Hyperref would then create
% a live annotation to a destination that cannot exist.  This helper keeps the
% current printed footnote mark and links it to an explicit stable note target.
\makeatletter
\newcommand{\sgarepeatedfootnotemark}[1]{%
  \begingroup
  \edef\@thefnmark{\thefootnote}%
  \hyperlink{#1}{\@makefnmark}%
  \endgroup}
\makeatother

\sisi{
\DeclareMathOperator{\simarrow}{\rule{.15mm}{1.8mm}\kern -3.4pt\rightsquigarrow}
\newcommand{\sheafell}{\underline{\ell}}
\newcommand{\sheafZZ}{\underline{\ZZ}}
}{
\newcommand{\nderef}[1]{la N.D.E.~(\ref{#1}) page~\pageref{#1}}
}

%%%%%%%%%%%%%%%%%
%%%%% expose XVII
%%%%%%%%%%%%%%%%%

% END INCLUDED SGA4_English_current_buildable_source_20260906/source/fichiers-configuration-latex/macros.tex


\externaldocument{../01/01_English.tex}
\externaldocument{../02/02_English.tex}
\externaldocument{../03/03_English.tex}
\externaldocument{../04/04_English.tex}
\externaldocument{../05/05_English.tex}
\externaldocument{../05bis/05bis_English.tex}
\externaldocument{../06/06_English.tex}
\externaldocument{../07/07_English.tex}
\externaldocument{../08/08_English.tex}
\externaldocument{../09/09_English.tex}
%\externaldocument{../10/10_English.tex}
\externaldocument{../11/11_English.tex}
\externaldocument{../12/12_English.tex}
\externaldocument{../13/13_English.tex}
\externaldocument{../14/14_English.tex}
\externaldocument{../15/15_English.tex}
\externaldocument{../16/16_English.tex}

\setboolean{original}{false}

\begin{document}

\begin{center}
XVII. Cohomology with proper supports\\
P. Deligne
\end{center}

\selectlanguage{english}version: \showgitstatus
\tableofcontents
\mainmatter

\else
\chapter{Cohomology with proper supports}\MakeLinkTarget*{sga4.exp.XVII}\label{sga4.exp.XVII}
%\addtocontents{toc}{P. Deligne}
\begin{center}
P. Deligne
\end{center}
\fi

% The edition is fixed to the corrected (second) reading.  Keeping the
% generic ifthen-based selector inside Xy-pic matrices breaks current
% Xy-pic; this equivalent expandable selector avoids that toolchain bug.
\renewcommand{\sisi}[2]{#2}

% BEGIN INCLUDED SGA4_English_current_buildable_source_20260906/source/17/drafts/SGA4_XVII_LINES_44_260_English.texfrag
\setcounter{pageoriginal}{251}

\footnotetext[1]{The present exposé and the following one, written in
  1968 and 1969, reproduce and supplement the oral lectures of
  A.~\textsc{Grothendieck} (spring 1964). The author of this written
  version, who did not attend the oral seminar, drew in part on notes
  by A.~\textsc{Grothendieck}.}


\section*{Introduction}\pageoriginale

In this exposé we develop the formalism of cohomology with proper
support. The questions of variance have been treated with considerable
care, whereas they were only made plausible in the oral seminar. This
explains the length of the exposé, whose Sections
\hyperref[XVII.sec1]{1} to \hyperref[XVII.sec4]{4} are devoted to
categories and general topology. The reader is very strongly advised
to read only Sections \hyperref[XVII.sec5]{5} and
\hyperref[XVII.sec6]{6}, where the geometric substance of the exposé
is concentrated (\S \hyperref[XVII.sec5]{5}: construction and
variance of cohomology with proper supports, base-change, finiteness,
finite \tTor-dimension, and comparison theorems; \S
\hyperref[XVII.sec6]{6}: theory of the \sisi{trace}{trace} morphism for
a quasi-finite flat morphism).

In Sections \hyperref[XVII.sec1]{1} and \hyperref[XVII.sec2]{2}, we
``recall'' some results on derived categories. The \S
\hyperref[XVII.sec1]{1} ${\rm n^{\circ}}$ 1 is devoted to the sign
formalism. In the \S \hyperref[XVII.sec3]{3}, we treat the problem of
gluing two variance formalisms. In the \S \hyperref[XVII.sec4]{4}, we
introduce flat resolutions and study the special properties of
\textsc{Godement}'s flasque resolutions.

In\pageoriginale this exposé, the higher direct-image functors with
proper supports $R^{q}f_{!}$, and the functor $Rf_{!}$ from which they
arise, are defined only when $f$ is a compactifiable morphism
(\hyperref[XVII.def3.2.1]{3.2.1}). In the appendix, written by
\textsc{B. Saint-Donat}, it is shown how to extend the definition to
separated morphisms of finite type with quasi-compact quasi-separated
target.\nde{\MakeLinkTarget*{sga4.exp.XVII.editor-note.introduction.nagata-compactification}\label{XVII.ndeConrad}The Nagata compactification theorem
ensures that every separated morphism of finite type with
quasi-compact quasi-separated target is compactifiable. See B.~Conrad,
``Deligne's notes on Nagata compactifications'', \emph{J. Ramanujan
Math.\ Soc.} \textbf{22} (2007), 205--257, cited as [C] below, \hyperref[XVII.subsec4.1]{4.1}, as
well as ``Erratum for `Deligne's notes on Nagata compactifications''',
\emph{J. Ramanujan Math.\ Soc.} \textbf{24} (2009), 427--428.}

The \S \hyperref[XVII.sec5]{5} ${\rm n^{\circ}}$
\hyperref[XVII.sec5]{5} (cohomology of a symmetric product) will not be
used again in this seminar. It will be used in \SGA 5 to refine the
rationality theorem for $L$-functions.\nde{See \SGA
4$\frac{1}{2}$, $L$-functions modulo $\ell^n$ and modulo $p$.}

The \S \hyperref[XVII.sec6]{6} ${\rm n^{\circ}}$
\hyperref[XVII.sec3]{3} (the theory of the trace for \emph{continuous}
coefficients) will be used in Exposé XVIII only in the relatively easy
case of smooth groups; a reader interested in the Poincaré duality
theorem (global duality), and willing to accept a transcendental
argument, may even dispense entirely with reading this \S
\hyperref[XVII.sec6]{6} ${\rm n^{\circ}}$
\hyperref[XVII.sec3]{3}, as well as most of the \S
\refnoexpose{XVIII.sec1} of XVIII.

\setcounter{section}{-1}
\section{Terminological preliminaries}\etiquette[0]{XVII.sec0}

\subsection{}\etiquette[0.1]{XVII.subsec0.1}
The sign $=$ placed between two groups of symbols denoting objects of a
category will sometimes mean (by abuse of notation) that these objects
are canonically isomorphic. The category and the canonical isomorphism
should in principle have been defined beforehand. In a diagram, the
sign $=$ will then denote the isomorphism itself.

\hypertarget{sga4.exp.XVII.subsection.0.2}{}\setcounter{subsection}{2}
\setcounter{subsubsection}{-1}
\subsubsection{}\etiquette[0.2.0]{XVII.subsubsec0.2.0}
Let $f:S\rightarrow S'$ be a morphism of sites
(\hyperrefext[IV][IV.subsubsec4.9.3]{4.9.3}). If $U$ and $U'$ are
objects of $S$ and $S'$, an $f$-morphism from $U$ to $U'$ will by
definition be a morphism from $U$ to $f^{\ast}U'$. For the étale sites
of schemes, this recovers the usual notion.

\subsubsection{}\etiquette[0.2.1]{XVII.subsubsec0.2.1}
If\pageoriginale $\mathcal{F}$ and $\mathcal{F}'$ are sheaves on $S$
and $S'$, an $f$-morphism from $\mathcal{F}$ to $\mathcal{F}'$ will
mean, equivalently,
\begin{enumerate}
\renewcommand{\theenumi}{\roman{enumi}}
\renewcommand{\labelenumi}{(\theenumi)}
\item a morphism from $f^{\ast}\mathcal{F}'$ to $\mathcal{F}$;

\item a morphism from $\mathcal{F}'$ to $f_{\ast}\mathcal{F}$;

\item a function which, to every $f$-morphism $\varphi$ from an object
  $U$ of $S$ to an object $U'$ of $S'$, assigns a function from
  $\mathcal{F}'(U')$ to $\mathcal{F}(U)$, compatibly with the
  composition of $\varphi$ with an arrow of $S$ or of $S'$
  \footnote{These terminologies are not compatible with identifying
  the objects of $S$ with the associated sheaves.}.
\end{enumerate}

It follows from (i), (ii), and (iii) that sheaves of sets form a
fibered and cofibered category over the category of sites.

Likewise for sheaves of left modules on ringed sites. Likewise for
étale sheaves on schemes; this is not immediately a special case of
the preceding statement, since sites in fact form a 2-category and
``the étale site of $X$'' is only a pseudofunctor in $X$
(\hyperrefext[VII][VII.subsec1.4]{1.4}).

The preceding definition makes sheaves of sets (\resp\ldots) on
varying sites into a fibered category over the category of sites,
whose fiber categories are the \emph{opposite categories} of the
usual categories of sheaves of sets (\resp\ldots).

\subsection{}\etiquette[0.3]{XVII.subsec0.3}
Let $I$ be a finite set, let $\epsilon$ be a function on $I$ with
values in $\{+1,-1\}$, let $(\Aa_i)_{i\in I}$ be a family of additive
categories graded by translation functors $T_i$ (\cite{sga4.exp.XVII.bib.CD} I
\sisi{1.10}{1.1.0}), and let $F$ be a multiadditive multifunctor from
the categories $\Aa_i$ to an additive category $\Aa$ graded by the
translation functor $T$. We suppose that $F$ is covariant (\resp
contravariant) in the variables $i$ for which $\epsilon_i=1$ (\resp
$\epsilon_i=-1$). Modifying and supplementing\pageoriginale
(\cite{sga4.exp.XVII.bib.CD} \sisi{I 1.1}{II 2.1.1}), we shall say that $F$ is a
\emph{graded functor} if we are given a family
$(\varphi_i)_{i\in I}$ of isomorphisms of functors between $T\circ F$
and $F\circ T_i^{\varepsilon_{i}}$, such that the following diagrams are
anticommutative for $i\neq j$ in $I$:
\[
\xymatrix{
F \circ (T_i^{\varepsilon_i}, T_j^{\varepsilon_j})  \ar[rr]^{\varphi_j
  \ast T_i^{\varepsilon_i}} \ar[dd]^{\varphi_i\ast T_j^{\varepsilon_j}}& &T\circ F \circ T_i^{\varepsilon_i} \ar[dd]^{T\ast
  \varphi_i}\\
 & - & \\
 T\circ F \circ T_j^{\varepsilon_j}  \ar[rr]^{T\ast \varphi_j} & &T\circ T
  \circ F.
}
\]
We leave it to the reader to define the composite of two graded
functors and verify that it is again a graded functor.

\subsection{}\etiquette[0.4]{XVII.subsec0.4}
Let $\Aa$ be an additive category and let $I$ be a finite set. For
$i\in I$, denote by $1_i$ the $i$\textsuperscript{th} basis vector of
$\ZZ^{I}$. An \emph{$I$-fold complex} (\resp a
\sisi{naive}{\emph{naive}} \emph{$I$-fold complex}) of $\Aa$ consists of
\begin{enumerate}
\renewcommand{\theenumi}{\alph{enumi}}
\renewcommand{\labelenumi}{(\theenumi)}
\item a family $(K^{k})_{k\in \ZZ^{I}}$ of objects of $\Aa$;

\item for every $i\in I$ and every $k\in \ZZ^{I}$, an arrow
  $d^{k}_i:K^{k}\rightarrow K^{k+1_i}$, these arrows satisfying
  $d^{k+1_i}_i d^{k}_i=0$ and, for $i\neq j$,
  $d^{k+1_j}_i d^{k}_j+d^{k+1_i}_j d^{k}_i=0$ (\resp
  $d^{k+1_j}_i d^{k}_j=d^{k+1_i}_j d^{k}_i$).
\end{enumerate}

\subsection{}\etiquette[0.5]{XVII.subsec0.5}
A\pageoriginale \emph{\ppro-object} of a category $\sisi{C}{\Cc}$ is a functor from
$\sisi{C}{\Cc}$ to $(\Ens)$ which is a filtered direct limit (over a
small filtered category) of representable functors (\cf
\hyperrefext[I][I.sec8.10]{8.10}). Every \ppro-object is a direct limit,
over a small filtered ordered set, of representable functors. If
$X_{i}$ is an inverse system of objects of $\sisi{C}{\Cc}$, indexed by a
small filtered category, we denote by $\quotprojlim X_{i}$ the
\ppro-object $\varinjlim h^{X_i}$.

\subsection{}\etiquette[0.6]{XVII.subsec0.6}
The reader may dualize \hyperref[XVII.subsec0.5]{0.5} to the case of
\iind-objects (\Cf \hyperrefext[I][I.sec8.2]{8.2}).

\subsection{}\etiquette[0.7]{XVII.subsec0.7}
In accordance with the new terminology, what was formerly called a
prescheme is now called a \emph{scheme}, and what was formerly called
a scheme is now called a \emph{separated scheme}.

\subsection{}\etiquette[0.8]{XVII.subsec0.8}
The category of \emph{ringed schemes} is the category whose objects are
schemes whose étale site is endowed with a sheaf of rings, an arrow
from $(S,\Aa)$ to $(T,\Bb)$ being a pair consisting of a morphism of
schemes $f$ from $S$ to $T$ and a homomorphism $\varphi$ of sheaves of
rings from $f^{\ast}\Bb$ to $\Aa$. The morphism of schemes $f$ is
said to be \emph{induced by} $(f,\varphi)$.

\subsection{}\etiquette[0.9]{XVII.subsec0.9}
If $X$ is a scheme over $Y$, we denote by $(X/Y)^n$ the $n$-fold fiber
product of $X$ over $Y$.

\subsection{}\etiquette[0.10]{XVII.subsec0.10}
Let $S$ be a scheme. The \emph{big étale site} (\resp the \emph{big
fppf site}, \resp the \emph{big fpqc site}) of $S$ is the site
$\Sch/{S}$, endowed with the étale topology (\resp the fppf topology,
\resp the fpqc topology) (\SGA 3 IV 6.3). Note carefully that this is
not a $\sisi{U}{\Uu}$-site ($\sisi{U}{\Uu}$ being the
\sisi{finite}{fixed} universe, and $\Sch/{S}$ consisting of the schemes
$\in \sisi{U}{\Uu}$).

The\pageoriginale \emph{small fpqc site} (\resp the \emph{small fppf
site}) of $S$ is the site of schemes flat over $S$ (\resp flat and of
finite presentation), endowed with the fpqc topology (\resp the fppf
topology). When it is necessary to avoid ambiguity, the \emph{small
étale site} of $S$ will mean the site $S_{\et}$
(\hyperrefext[VII][VII.subsec1.2]{1.2}).

\subsection{}\etiquette[0.11]{XVII.subsec0.11}
In many cases, the author has allowed himself to speak of commutative
diagrams of morphisms of sites where one should have spoken of
essentially commutative diagrams, \ie diagrams commutative up to
isomorphism (\cf IV \hyperref[IV.subsubsec3.2.2]{3.2.2}). The reader may verify that the arguments we
give also apply to the general situation.

\subsection{}\etiquette[0.12]{XVII.subsec0.12}
The terminology ``\emph{coherent scheme}'' for ``quasi-compact
quasi-separated scheme'' has been surreptitiously introduced here and
there by A.~Grothendieck.

\subsection{}\etiquette[0.13]{XVII.subsec0.13}
A torsion sheaf $F$ on a scheme $S$ will be said to be \emph{prime to
the residual characteristics of} $S$ if it is the direct limit of its
subsheaves annihilated by integers $n$ invertible on $S$.

% END INCLUDED SGA4_English_current_buildable_source_20260906/source/17/drafts/SGA4_XVII_LINES_44_260_English.texfrag

% BEGIN INCLUDED SGA4_English_current_buildable_source_20260906/source/17/drafts/SGA4_XVII_LINES_261_600_English.texfrag
\section{Derived categories}\etiquette[1]{XVII.sec1}

\subsection{Exact functors (the sign rules)}\etiquette[1.1]{XVII.subsec1.1}

\subsubsection{}\etiquette[1.1.1]{XVII.subsubsec1.1.1}
For the fundamental theorems concerning derived categories, I refer
to \textsc{Verdier} \cite{sga4.exp.XVII.bib.CD} and \cite{sga4.exp.XVII.bib.V2}. In this subsection,
particularly careful attention has been paid to questions of signs.

Recall that if $f:X\rightarrow Y$ is a morphism of complexes
(in an additive category, which will henceforth be understood), its
\emph{cone} $C(f)$\pageoriginale is defined by
$$\etiquette[1.1.1.1]{XVII.eq1.1.1.1}
C(f)^{n}=X^{n+1}\oplus Y^{n}\quad d^{n}=-d^{n+1}_X+f^{n+1}+d^{n}_Y{\quoi}.\leqno(1.1.1.1)
$$
When $X=0$ (\resp $Y=0$), we have $C(f)=Y$ (\resp $C(f)=X[1]$), so
that the commutative diagram
\[
\xymatrix{
0        \ar@{=}[r] \ar[d]&   0   \sisi{\ar@{=}[r]}{\ar[r]} \ar[d]&  X
\ar@{=}[r]\ar[d]^{f}  & X \ar[d]\\
X        \ar[r]^{f}   &  Y    \ar@{=}[r] & Y     \ar[r]        & 0
}
\]
defines a ``triangle''
\[
X\xrightarrow{f}Y\rightarrow C(f)\xrightarrow{j}X[1]{\quoi}.
\]

In the category of complexes up to homotopy, a triangle is called
\emph{distinguished} if it is isomorphic to a triangle of this type,
and \emph{antidistinguished} if it becomes distinguished when the
signs of its arrows are changed. One checks that the triangle defined
by an exact sequence of complexes which is split in each degree
(\cite{sga4.exp.XVII.bib.CD} p.~10, \cf the proof of
\hyperref[XVII.lem1.1.5.4]{1.1.5.4}) is antidistinguished and that, up
to isomorphism, every antidistinguished triangle is obtained in this
way.

\subsubsection*{${}^{\ast}$1.1.2.}\etiquette[*1.1.2]{XVII.subsubsec*1.1.2}
Sign calculations are most easily carried out using the following
``formal rules''. Write an element of $C(f)^{n}$ in the form
$1\otimes x+y$ $(x\in X^{n+1}, y\in Y^{n})$, where $1$ is a
``cell'' of dimension $1$, and set $d(1)\otimes x=f(x)$. Thus
\[
d(1\otimes x+y)=d1\otimes x-1\otimes dx+dy=-1\otimes
dx+f(x)+dy{\quoi}._{\ast}
\]

\begin{enonce*}{Definition 1.1.3}\etiquette[1.1.3]{XVII.def1.1.3}
\begin{enumerate}
\renewcommand{\theenumi}{\roman{enumi}}
\renewcommand{\labelenumi}{(\theenumi)}
\item A\pageoriginale graded functor (\hyperref[XVII.subsec0.3]{0.3})
  from a triangulated category $\Aa$ to a triangulated category $\Bb$
  is said to be exact if it sends distinguished triangles to
  distinguished triangles.

\item A contravariant graded functor from a triangulated category
  $\Aa$ to a triangulated category $\Bb$ is said to be
  \emph{exact} if, for every distinguished triangle
  $(X,Y,Z,u,v,w)$ of $\Aa$, the triangle
  $(F(Z),F(Y), F(X), F(v),F(u),TF(w))$ is distinguished (in this
  definition, $FX$ is identified with $TFTX$ by means of the grading
  of $F$).

\item Let $(\Aa_i)_{i\in I}$ be a finite family of triangulated
  categories, let $\epsilon$ be a function from $I$ to
  $\{+1,-1\}$, and let $\Aa$ be a triangulated category. A graded
  functor $F$ from the product of the $\Aa_i$ to $\Aa$, covariant
  in those $i$ for which $\epsilon(i)=1$ and contravariant in those
  for which $\epsilon(i)=-1$, is said to be \emph{exact} if, for every
  $i\in I$, the functors from $\Aa_i$ to $\Aa$ obtained from $F$ by
  fixing all variables except the $i$-th are exact functors.
\end{enumerate}
\end{enonce*}

A composite of exact functors is again an exact functor.

\setcounter{subsubsection}{3}
\subsubsection{}\etiquette[1.1.4]{XVII.subsubsec1.1.4}
If $(K^{n_1\ldots n_p}, d_1,\ldots,d_p)$ is a multicomplex
(\hyperref[XVII.subsec0.4]{0.4}), its associated simple complex is
defined by
$$\etiquette[1.1.4.0]{XVII.eq1.1.4.0}
K^{n}=\prod_{\sum n_i=n}K^{(n_i)}\quav d=\sum d_i{\quoi}.\leqno(1.1.4.0)
$$

When, exceptionally, it is defined using a direct sum rather than a
product, this fact will be stated explicitly.

Let $I$ be a finite set. Following \textsc{Cartan--Eilenberg} and
\textsc{J.~P. Serre}, we shall explain how every naive $I$-fold
complex $K$ has canonically\pageoriginale associated with it an
$I$-fold complex (\hyperref[XVII.subsec0.4]{0.4}).

For every total order $<$ on $I$, let $K(<)$ be the following
$I$-fold complex:
$$\etiquette[1.1.4.1]{XVII.eq1.1.4.1}
\begin{cases}
K(<)^{\sheafk }=K^{\sheafk }\\
d(<)_i^{\sheafk }=(-1)^{\sum\limits_{j<i} k_j} d_i{\quoi}.
\end{cases}\leqno{(1.1.4.1)}
$$

If $<_1$ and $<_2$ are two \sisi{total}{total} orders on $I$, the
\emph{canonical isomorphism $\tau$ between $K(<_1)$ and $K(<_2)$} is,
by definition, the one given by
$$\etiquette[1.1.4.2]{XVII.eq1.1.4.2}
\tau^{\sheafk }=(-1)^{\epsilon(<_1,<_2)},\quad\mbox{\rm where}\quad
\epsilon(<_1,<_2)=\sum_{\substack{i<_1j\\ \sisi{i<_2 j}{j<_2 i}}}k_ik_j{\quoi}.\leqno{(1.1.4.2)}
$$

\subsubsection*{1.1.4.2.1}\etiquette[1.1.4.2.1]{XVII.par1.1.4.2.1}
One checks that these canonical isomorphisms form a transitive system
of isomorphisms among the $K(<)$, as $<$ ranges over the total orders
on $I$; these $K(<)$ therefore determine an $I$-fold complex, unique
up to a unique isomorphism, which in this subsection will be denoted by
$K(\ast)$ and is called the \emph{$I$-fold complex associated with the
naive $I$-fold complex $K$}. Note that $K(\ast)^{\sheafk}$ is
\emph{not canonically isomorphic to $K^{\sheafk}$: only the choice of
a total order on $I$ makes it possible to identify these two objects.}

If $K$ is a naive $I$-fold complex and $i\in I$, let $K[1_i]$ denote
the following naive complex, whose differentials are denoted by
$d_j[1_i]$:
$$\etiquette[1.1.4.3]{XVII.eq1.1.4.3}
\begin{cases}
K[1_i]^{\sheafk } = K^{\sheafk +1_i}\\
d_j[1_i]^{\sheafk } = d_j^{\sheafk +1_{i}}\quad{\mbox{\rm if }} i\neq
j\\
d_i[1_i]^{\sheafk }= -d_i^{\sheafk +1_i}{\quoi}.
\end{cases}\leqno{(1.1.4.3)}
$$

If\pageoriginale $K$ is an $I$-fold complex, set
$$\etiquette[1.1.4.4]{XVII.eq1.1.4.4}
\begin{cases}
K[1_i]^{\sheafk } = K^{\sheafk +1_i}\\
d_j[1_i]^{\sheafk } = -d_j[1_i]^{\sheafk +1_i}\sisi{}{.}
\end{cases}\leqno{(1.1.4.4)}
$$

If $K$ is a naive $I$-fold complex and $<$ is a total order on $I$,
so that $K(\ast)$ is given by \eqhyperref[XVII.eq1.1.4.1]{1.1.4.1},
define an isomorphism $\sigma$ between $K[1_i](\ast)$ and
$K(\ast)[1_i]$ by the formula
$$\etiquette[1.1.4.5]{XVII.eq1.1.4.5}
\sigma^{\sheafk }=(-1)^{\sum\limits_{j<i}k_j} \id_{K^{\sheafk +1_i}}{\quoi}.\leqno{(1.1.4.5)}
$$

\setcounter{subsubsection}{4}
\setcounter{paragraph}{5}
\paragraph{}\etiquette[1.1.4.6]{XVII.par1.1.4.6}
One checks that this isomorphism is independent of the choice of the
total order $<$; it therefore defines an \emph{isomorphism, called the
canonical isomorphism, between} $K[1_i](\ast)$ and $K(\ast)[1_i]$.

\paragraph{}\etiquette[1.1.4.7]{XVII.par1.1.4.7}
For any $i,j\in I$, if $K$ is an $I$-fold complex or a naive $I$-fold
complex, we identify $K[1_i][1_j]$ and $K[1_j][1_i]$ in the evident
way.

For $i\neq j$ in $I$, the following diagram of canonical
isomorphisms of type \hyperref[XVII.par1.1.4.6]{1.1.4.6} or
\hyperref[XVII.par1.1.4.7]{1.1.4.7} is then \emph{anticommutative}:
$$\etiquette[1.1.4.7.1]{XVII.eq1.1.4.7.1}
\vcenter{
\xymatrix{
K[1_i][1_j](\ast) \ar@{=}[r] \ar@{<->}[d]& K[1_j][1_i](\ast)
\ar@{<->}[r] & K[1_j](\ast)[1_i]\ar@{<->}[d] \\
K[1_i](\ast)[1_j] \ar@{<->}[r] & K(\ast)[1_i][1_j]  \ar@{=}[r]& K(\ast)[1_j][1_i]\rlap{.}
}}\leqno{(1.1.4.7.1)}
$$


\subsubsection{}\etiquette[1.1.5]{XVII.subsubsec1.1.5}
Let\pageoriginale $F:\Aa\rightarrow \Bb$ be a contravariant additive
functor between additive categories. Its \emph{extension to complexes}
$F:C(\Aa)\rightarrow C(\Bb)$ is defined as follows. If $K$ is a
complex of $\Aa$, set
$$\etiquette[1.1.5.1]{XVII.eq1.1.5.1}
\begin{cases}
F(K)^{k}=F(K^{-k})\\
d^{k}_{F(K)}=(-1)^{k+1}F(d^{-k-1})\sisi{}{.}
\end{cases}\leqno(1.1.5.1)
$$

The \emph{canonical isomorphism $\sigma$ between $F(K)[1]$ and
$F(K[-1])$} is defined by
$$\etiquette[1.1.5.2]{XVII.eq1.1.5.2}
\sigma^{\sisi{\sheafk }{k}}=(-1)^{k+1}\id_{F(K^{-k-1})}{\quoi}.\leqno(1.1.5.2)
$$

If $F:\Aa\rightarrow \Bb$ and $G:\Bb\rightarrow
\sisi{C}{\Cc}$ are two contravariant additive functors, define the
canonical isomorphism $\rho$ between $G{\circ}F(K)$ \emph{and}
$G(F(K))$ by
$$\etiquette[1.1.5.3]{XVII.eq1.1.5.3}
\rho^{k}=(-1)^{k}\sisi{}{\id_{G\circ F(K^k)}}{\quoi}.\leqno(1.1.5.3)
$$

\begin{enonce*}{Lemma 1.1.5.4}\etiquette[1.1.5.4]{XVII.lem1.1.5.4}
\begin{enumerate}
\renewcommand{\theenumi}{\roman{enumi}}
\renewcommand{\labelenumi}{(\theenumi)}
\item For every contravariant additive functor
  $F:\Aa\rightarrow \Bb$, formulas
  \eqhyperref[XVII.eq1.1.5.1]{1.1.5.1} and
  \eqhyperref[XVII.eq1.1.5.2]{1.1.5.2} define an exact contravariant
  functor from $K(\Aa)$ to $K(\Bb)$.

\item For every pair $(G,F)$ of composable contravariant additive
  functors, formula \eqhyperref[XVII.eq1.1.5.3]{1.1.5.3} defines an
  isomorphism of graded functors between the extension of $GF$ to
  complexes and the composite of the extensions of $G$ and $F$ to
  complexes.
\end{enumerate}
\end{enonce*}

\begin{enonce*}[remark]{Proof}
A\pageoriginale morphism of degree $i$, $f:K\rightarrow L$, may be
regarded as any one of the following:
\begin{enumerate}
\renewcommand{\theenumi}{\alph{enumi}}
\renewcommand{\labelenumi}{(\theenumi)}
\item a morphism from $K$ to $L[i]$;

\item a morphism from $K[-i]$ to $L$;

\item a family $f^{n}:K^{n}\rightarrow L^{n+i}$ which commutes or
  anticommutes with the differentials according to the parity of $i$.
\end{enumerate}
\end{enonce*}

If $f:K\rightarrow L[i]$ is a morphism of degree $i$, then
$F(f):F(L)[-i]\rightarrow F(K)$ is again a morphism of degree $i$;
for $i=1$, we have
$$\etiquette[1.1.5.5]{XVII.eq1.1.5.5}
F(f)^{n}=(-1)^{n+1}F(f^{-n-1}){\quoi}.\leqno(1.1.5.5)
$$

If $0\rightarrow K\xrightarrow{u} L\xrightarrow{v} M\rightarrow 0$
is an exact sequence of complexes split degree by degree, the choice
of a splitting allows the differential of $L$ to be written in the
form
\[
d_L=d_K+d_M+f{\quoi},
\]
where $f$ is a morphism of degree one from $M$ to $K$. Likewise, the
dual exact sequence
\[
0\longrightarrow F(M)\longrightarrow F(L)\longrightarrow
F(K)\longrightarrow 0
\]
defines a morphism of degree one $f^{\ast}$ from $F(K)$ to
$F(\sisi{L}{M})$. We have $f^{\ast}=F(f)$. It follows that the image
under $F$ of a triangle $(K,L,M,u,v,f)$ defined by a short exact
sequence of complexes is again a triangle of the same type and, by
\hyperref[XVII.subsubsec1.1.1]{1.1.1}, this proves (i).

Under the hypotheses of (ii), for $f$ of degree one, we have
$G(F(f))^{n}=(-1)^{n+1}(-1)^{-n}G\circ F(f)^{n}=-G\circ F(f)^{n}$; in
particular, the grading isomorphism $G(F(K[1]))\simeq
G(F(K))[1]$ is the isomorphism whose components are $-1$.

The\pageoriginale diagram
\[
\xymatrix{
G\circ F(K)[1]  \ar@{<->}[r]^{\rho[1]} \ar@{=}[d] & G(F(K))[1] \ar@{<->}[d]\\
G\circ F(K[1])  \ar@{<->}[r]^{\rho} & G(F(K[1]))
}
\]
is therefore commutative, which proves (ii).

\subsubsection{}\etiquette[1.1.6]{XVII.subsubsec1.1.6}
Let $s=+$ or $s=-$; let $I$ be a finite set, let
$\epsilon:I\rightarrow \{+,-\}$ be a function, let
$(\Aa_i)_{i\in I}$ and $\Aa$ be additive categories, and let $F$ be
a multiadditive multifunctor from the categories $\Aa_i$ to $\Aa$,
covariant in the arguments indexed by those $i$ for which
$\epsilon(i)=+$ and contravariant in the others. We shall define the
extension $K^{s}(F)$ (or simply $F$) of $F$ to complexes as an exact
multifunctor from the categories $K^{s\epsilon(i)}(\Aa_i)$
($s\epsilon(i)=\pm$, according to the sign rule) to $K^{s}(\Aa)$,
\footnote[1]{The construction also provides an extension of $F$ to a
  multifunctor from the categories $C^{s\epsilon(i)}(\Aa_i)$ to
  $C^{s}(\Aa)$.} having the same variance as $F$. We shall then explain
in what sense this construction is compatible with the composition of
functors. Other conventions for \emph{degrees} are possible and will
be used; that issue is independent of the questions of \emph{sign}
considered here.

\paragraph{}\etiquette[1.1.6.1]{XVII.par1.1.6.1}
If $F$ is contravariant in one variable, $K^{s}(F)$ is defined by the
formulas of \hyperref[XVII.subsubsec1.1.5]{1.1.5}.

\paragraph{}\etiquette[1.1.6.2]{XVII.par1.1.6.2}
If $F$ is covariant, the functor $K^{s}(F)$ is defined, by passage to
the quotient, from the composite of the evident functor from the
product of the categories $C^{s}(\Aa_i)$ to the category of naive
$I$-fold complexes, the functor
(\hyperref[XVII.subsubsec1.1.4]{1.1.4}) from the category of naive
$I$-fold complexes to that of\pageoriginale $I$-fold complexes
(\hyperref[XVII.par1.1.4.2.1]{1.1.4.2.1}), and the associated simple
complex functor \eqhyperref[XVII.eq1.1.4.0]{1.1.4.0}:
\[
K^{s}(F)=C^{s}(F)(\ast)_{s.a}{\quoi}.
\]

Let $K_j\in K^{s}(\Aa_j)$ be complexes, and let $i\in I$; set
$K_j[1_i]=K_j$ for $j\neq i$, and $K_i[1_i]=K_{i}[1]$. We then have
the evident isomorphisms
\[
C^{s}(F)(K_{j}[1_i])\simeq C^{s}(F)(K_j)[1_i]{\quoi},
\]
and
\[
C^{s}(F)(K_j)(\ast)[1_i]_{s.a}\simeq
C^{s}(F)(K_j)(\ast)_{s.a}[1]{\quoi},
\]
whence, via \eqhyperref[XVII.eq1.1.4.5]{1.1.4.5}, a
\emph{canonical isomorphism} $\sigma$ between
$K^{s}(F)(K_{j}[1_i])$ and $K^{s}(F)(K_{j})[1]$. It follows from
\hyperref[XVII.par1.1.4.7]{1.1.4.7} that these isomorphisms define a
grading of $K^{s}(F)$.

\paragraph{}\etiquette[1.1.6.3]{XVII.par1.1.6.3}
In the general case, $F$ is the composite of the canonical
contravariant functors $\Aa_i\rightarrow \Aa^{\circ}_i$ (for
$\epsilon(i)=-$) with a covariant functor $F^{+}$; define the functor
$K^{s}(F)$ by composition using
\hyperref[XVII.par1.1.6.1]{1.1.6.1} and
\hyperref[XVII.par1.1.6.2]{1.1.6.2}.

\begin{enonce*}{Proposition 1.1.7}\etiquette[1.1.7]{XVII.prop1.1.7}
Under the preceding hypotheses, the functor $K^{s}(F)$ is an exact
multiadditive functor.
\end{enonce*}

It suffices to verify \hyperref[XVII.prop1.1.7]{1.1.7} under the
hypotheses of \hyperref[XVII.par1.1.6.1]{1.1.6.1} (where the
proposition is \hyperref[XVII.lem1.1.5.4]{1.1.5.4}~(i)) or under those
of \hyperref[XVII.par1.1.6.2]{1.1.6.2}. To verify that the functors
deduced from $K^{s}(F)$ by fixing all variables except the $i$-th are
exact, one may describe $K^{s}(F)$ explicitly in terms of an order on
$I$ for which $i$ is the least element; a distinguished triangle of
the standard\pageoriginale form
\eqhyperref[XVII.eq1.1.1.1]{1.1.1.1} is then sent to a distinguished
triangle of the same form.

\setcounter{subsubsection}{7}
\subsubsection{}\etiquette[1.1.8]{XVII.subsubsec1.1.8}
Let $I$ and $J$ be two finite sets, let $\Psi:I\rightarrow J$ be a
surjective map, and let $\epsilon_{I}:I\rightarrow \{+,-\}$ and

% END INCLUDED SGA4_English_current_buildable_source_20260906/source/17/drafts/SGA4_XVII_LINES_261_600_English.texfrag

% BEGIN INCLUDED SGA4_English_current_buildable_source_20260906/source/17/drafts/SGA4_XVII_LINES_601_1000_English.texfrag
$\epsilon_{J}:J\rightarrow \{+,-\}$ be functions, and let us be given
multiadditive multifunctors, covariant in the arguments for which
$\epsilon=+$ and contravariant in the others:
\begin{align*}
G_j&: \prod_{\Psi(i)=j}\Aa_i\longrightarrow
\Bb_j\qquad (j\in J)\\
F&:  \prod_{j\in J} \Bb_j\longrightarrow \sisi{C}{\Cc.}
\end{align*}

Choose total orders on $I$ and $J$ such that $\Psi$ is an increasing
map. We then define the \emph{canonical isomorphism $\rho$ between the
  extension to complexes of $F\circ (G_j)$ and the composite of the
  extensions to complexes of $F$ and of the $G_j$} by the formula
$$\etiquette[1.1.8.1]{XVII.eq1.1.8.1}
\rho^{\sheafk }=\sisi{\prod_i \eta_i^{k_i}}{(-1)^{A(\sheafk )}}:\; \mbox{\rm automorphism of }
F\circ (G_j)(K_{i}^{k_i\epsilon\sisi{}{_I}(i)\epsilon\sisi{}{_J}(\Psi(i))}){\quoi}.\leqno(1.1.8.1)
$$%
where \sisi{$\eta_i=-1$ if $\epsilon_I(i)=\epsilon_{J}(\Psi(i))=-$,
$\eta_i=1$ otherwise.}{\[
A(\sheafk )=\sum_{\epsilon_I(i)=\epsilon_J(\Psi(i))=-} k_i + \sum_{\epsilon_J(j)=-} \sum_{\substack{\Psi(a)=\Psi(b)=j\\a<b}} k_a k_b.
\]}%See the Erratum in SGA 4 1/2, p. 312

This formula generalizes \eqhyperref[XVII.eq1.1.5.3]{1.1.5.3}; one verifies that $\rho$ does not depend
on the orders chosen on $I$ and $J$, and is an isomorphism of graded
functors (\cf \hyperref[XVII.lem1.1.5.4]{1.1.5.4} (ii)).

These isomorphisms $\rho$ satisfy the evident cocycle condition for a
triple composite.

\subsubsection{}\etiquette[1.1.9]{XVII.subsubsec1.1.9}
Let\pageoriginale $I$ be a finite set, let $i$ and $j$ be two distinct
elements of $I$, let $I'=I\backslash\{i,j\}$, let $(\Aa_i)_{i\in I}$ and
$\Aa$ be additive categories, and let $F:(\Aa_i)_{i\in
  I}\rightarrow \Aa$ and $G:(\Aa_i)_{i\in
  I'}\rightarrow \Aa$ be multiadditive multifunctors. Suppose that
$\Aa_i=\Aa_j$ and that $F$ is contravariant in the $i$-th variable and
covariant in the $j$-th. For $X$, $Y\in \Ob \Aa_i$, denote by
$F_{X,Y}$ the multifunctor from $(\Aa_i)_{i\in I'}$ to $\Aa$
obtained by fixing the $i$-th and $j$-th arguments:
\[
F_{X,Y}(\ldots)=F(\ldots\sisi{}{,} X,\ldots,Y,\ldots){\quoi}.
\]
A \emph{contraction morphism} $c:F\rightarrow G$ consists of, for
every $X\in \Aa_i$, a morphism of functors
\[
c_{X}:F_{X,X}\longrightarrow G{\quoi},
\]
such that, for every arrow $f:X\rightarrow Y$ of $\Aa_i$, the diagram
\[
\xymatrix{
F_{Y,X} \ar[d] \ar[r] & F_{X,X} \ar[d]^{c_X}\\
F_{Y,Y}  \ar[r]_{c_Y}     & G
}
\]
is commutative.

The conventions of \hyperref[XVII.subsubsec1.1.6]{1.1.6} are motivated by the

\begin{enonce*}{Lemma 1.1.9.1}\etiquette[1.1.9.1]{XVII.lem1.1.9.1}
Let $c:F\rightarrow G$ be a contraction morphism, and let
$(K_{\ell})_{\ell \in I}$ be complexes in the $\Aa_{\ell}$, with
$K_{i}=K_{j}$. Let $<$ be a total order on $I$ for which $i$ is the
predecessor of $j$. With the notation of
\hyperref[XVII.subsubsec1.1.6]{1.1.6}, for $k\in \ZZ$, let
\[
c^{k}:K^{s}(F)(K_{\ell})^{k}\longrightarrow K^{s}(G)(K_{\ell})^{k}
\]
be\pageoriginale the morphism whose restriction to $F(\ldots
K_{i}^{k_i}K_{j}^{k_j}\ldots)$ is $0$ if $k_i\neq k_j$, and is
$c_{K_i^{k_i}}$ if $k_{i}=k_{j}$. Then $(c^{k})_{k\in \ZZ}$ is a
morphism of complexes.
\end{enonce*}

We leave to the reader the verification of this lemma, as well as the
fact that the \emph{contraction morphism} $(c^{k})$ does not depend on
the order chosen on $I$.

The standard example of a contraction morphism is the composition law
\[
\Hom(Y,Z)\otimes \Hom(X,Y)\longrightarrow \Hom(X,Z)\sisi{}{.}
\]

\begin{enonce*}[remark]{Examples 1.1.10}\etiquette[1.1.10]{XVII.exem1.1.10}
\begin{enumerate}
\renewcommand{\theenumi}{\roman{enumi}}
\renewcommand{\labelenumi}{(\theenumi)}
\item Applying the definitions of \hyperref[XVII.subsubsec1.1.6]{1.1.6} to the
  \emph{tensor-product} functor gives the usual notion of the tensor
  product of two complexes, with differential
\[
d(x\otimes y)=dx\otimes y+(-1)^{\deg x}x\otimes dy{\quoi}.
\]

\item Apply the definitions of \hyperref[XVII.subsubsec1.1.6]{1.1.6} to the
  functor $\Hom(X,Y)$, \emph{regarding $Y$ as the first variable and
  $X$ as the second variable}, a convention that we shall always
  follow. This convention identifies, for two complexes $X$, $Y$, the
  $k$\sisi{-th}{-th} component of the complex $\Hom(X,Y)$ defined in
  \hyperref[XVII.subsubsec1.1.6]{1.1.6} with $\prod_n\Hom(X^{n},Y^{n+k})$, and gives the
  usual differential: for $f\in \Hom(X^{n},Y^{n+k})$, one has
\[
df=d\circ f-(-1)^{k}f\circ d{\quoi}.
\]

\item The preceding conventions have the virtue that the following
  systems of arrows are morphisms of complexes, without any signs
  having\pageoriginale to be inserted (\hyperref[XVII.subsubsec1.1.9]{1.1.9}):
\begin{align*}
& \Hom^{\bigdot}(L,M)\otimes \Hom^{\bigdot}(K,L)\longrightarrow
  \Hom^{\bigdot}(K,M):f\otimes g\mapsto f\circ g\\
& \Hom^{\bigdot}(L,M)\otimes L\longrightarrow M: f\otimes x\mapsto f(x){\quoi}.
\end{align*}
\end{enumerate}
\end{enonce*}

\subsubsection*{${}^{\ast}$1.1.10.1}\etiquette[*1.1.10.1]{XVII.par1.1.10.1}
The convention \eqhyperref[XVII.eq1.1.4.5]{1.1.4.5}--(\hyperref[XVII.par1.1.6.2]{1.1.6.2}) is justified as follows for
$F=\otimes$: an element of $((K\otimes L)[1])^{n}$ is written in the
form (\hyperref[XVII.subsubsec*1.1.2]{*1.1.2}) (recalling that $(K\otimes L)[1]=C(K\otimes
L\rightarrow 0)$ \eqhyperref[XVII.eq1.1.1.1]{1.1.1.1}) $\sum 1\otimes(x^{p}\otimes y^{q})$, an
element of $(K[1]\otimes L)^{n}$ in the form $\sum (1\otimes
x^{p})\otimes y^{q}$, and an element of $(K\otimes L[1])^{n}$ in the
form $\sum x^{p}\otimes (1\otimes y^{q})$; the sign appears when
$x^{p}$ and $1$ are interchanged. For $F=$ dual, the convention
\eqhyperref[XVII.eq1.1.5.2]{1.1.5.2} is justified as follows: an element of $K^{\ast}[1]$ is written
$1\otimes \omega$; its value on an element $1'\otimes x$ of $K[-1]$
($1'$ a cell of dimension $-1$) is $<1\otimes \omega, 1'\otimes
x>=(-1)^{\deg\omega}<1, 1'><\omega,x>=(-1)^{\deg\omega}<\omega,x>$;
this is the convention adopted.${}_{\ast}$

\begin{enonce*}[remark]{Remark 1.1.11}\etiquette[1.1.11]{XVII.rem1.1.11}
``Recall'' that a \emph{preadditive category} is a category whose sets
of arrows $\Hom(X,Y)$ are endowed with abelian-group structures, the
composition of morphisms being biadditive. An \emph{additive functor}
$F$ between preadditive categories is a functor satisfying the formula
\[
F(f+g)=F(f)+F(g)
\]
for $f$, $g\in \Hom(X,Y)$. The additive functors from a preadditive
category $\Aa$ to another $\Bb$ form a category denoted by
$\Hhom_{\text{\rm add}}(\Aa,\Bb)$.
\end{enonce*}

Let\pageoriginale $(D)$ be the unique preadditive category whose set of objects is
indexed by $\ZZ$, $(X\sisi{_n}{^n})_{n\in \ZZ}$, and for which
$$\etiquette[1.1.11.1]{XVII.eq1.1.11.1}
\begin{cases}
\Hom(X^{n},X^{n})=\ZZ,\;\mbox{\rm generated by the identity}\\
\Hom(X^{n},X^{n+1})=\ZZ,\;\mbox{\rm generated by an arrow
  denoted $d_n$}\\
\Hom(X^{i},X^{j})=0\;\mbox{if } j-i\neq 0\;\mbox{\rm or } 1{\quoi}.
\end{cases}\leqno(1.1.11.1)
$$

Let $D$ be the complex of $(D)$ defined by $D^{n}=X^{n}$ and
$d^{n}=d_n$.

For every preadditive category $\Aa$, the functor
$G\mapsto G(D)$ from $\Hhom_{\text{\rm add}}((D),\Aa)$ to
$C(\Aa)$ is an isomorphism.

The construction of \hyperref[XVII.subsubsec1.1.6]{1.1.6} extends to the case of a multiadditive functor
from preadditive categories $\Aa_i$ to an additive category $\Aa$.

Let
\[
F^{\bigdot}=(F^{n},d^{n})_{n\in \ZZ}
\]
be a complex of multiadditive functors from additive categories
$\Aa_i$ to an additive category $\Aa$. One may regard
$F^{\bigdot}$ as a functor from $(D)$ to a category of multifunctors,
or equivalently as a multiadditive functor
\[
F_1:\prod\Aa_i\times (D)\longrightarrow \Aa{\quoi}.
\]
The \emph{extension of $F\sisi{}{^{\bigdot}}$ to complexes} is defined by the formula
$$\etiquette[1.1.11.2]{XVII.eq1.1.11.2}
F^{\bigdot}(K_{1}\ldots K_{n})=F_{1}(K_{1}\ldots K_{n},D){\quoi}.\leqno(1.1.11.2)
$$

\setcounter{subsubsection}{11}
\subsubsection{}\etiquette[1.1.12]{XVII.subsubsec1.1.12}
Let $\Aa$ and $\Bb$ be two additive categories and let
\[
F^{\bigdot}=(F^{n},d^{n})_{n\in \ZZ}
\]
be a complex of additive functors from $\Aa$ to $\Bb$.
Suppose that the functors\pageoriginale $F^{i}$ admit right adjoints $F_i$; let
$\leftexp{t}{d}^{i}:F_{i+1}\rightarrow F_i$ be the transpose of
$d^{i}$, and let $F_{\bigdot}$ be the complex of additive functors from
$\Bb$ to $\Aa$ whose components are the $F_{i}$ and whose differential
operator is
\[
(d_\bigdot)^{i}=(-1)^{i+1}\,\leftexp{t}{d}^{-i-1}:F_{-i}\longrightarrow
F_{-i-1}{\quoi}.
\]
Let ${}^{t}$ be the \emph{contravariant} functor from $(D)$ to $(D)$
given by
\[
\leftexp{t}{X}^{n}=X^{-n}\quav \leftexp{t}{d}^{n}=d^{-n-1}{\quoi}.
\]
The following diagram is commutative up to canonical isomorphism:
$$\etiquette[1.1.12.1]{XVII.eq1.1.12.1}
\vcenter{
\xymatrix@C=6em{
\mathcal{A} \times (D) \times \mathcal{B}  \ar[d]^{(\Id,{}^t,\Id)}\ar[r]^(.55){(F^{n}(A), \Id)} &
\mathcal{B}\times \mathcal{B}  \ar[r]^{\Hom_{\mathcal{B}}}&
(\Ab)\ar@{=}[d] \\
\mathcal{A} \times (D) \times \mathcal{B}
\ar[r]^(.6){(\Id,(F_{\bigdot})^n(B))} & \mathcal{A} \times \mathcal{A}
\ar[r]^{\Hom_{\mathcal{A}}} & (\Ab) \rlap{.}
}}\leqno{(1.1.12.1)}
$$
This expresses the adjunction formula
\begin{equation*}
\Hom(F^{n}(A),B)\overset{\sim}{\longleftrightarrow}\Hom(A,F_{n}(B)){\quoi}.
\end{equation*}
The general conventions of \hyperref[XVII.def1.1.3]{1.1.3} therefore give an isomorphism
\begin{equation*}
\Hom^{\bigdot}(F^{\bigdot}(K),L)\overset{\sim}{\longleftrightarrow}\Hom\sisi{}{^\bigdot}(K,F_{\bigdot}(L)){\quoi}.
\end{equation*}

\begin{enonce*}{Definition 1.1.13}\etiquette[1.1.13]{XVII.def1.1.13}
If $K$ is a complex in an abelian category $\Aa$, the \emph{right
truncation (\resp left truncation)} of $K$ in dimension $n$, denoted
by $\tau_{\leq n}(K)$ \emph{(\resp $\tau_{\geq n}(K)$)}, means the
following complexes:
\begin{align*}
\tau_{\leq n}(K) &: \ldots K^{n-2}\longrightarrow
K^{n-1}\longrightarrow \Ker(d^n)\longrightarrow 0\\
\tau_{\geq n}(K) &: 0\longrightarrow \coker (d^{n-1})\longrightarrow
K^{n+1}\longrightarrow K^{n+2}\ldots.
\end{align*}
\end{enonce*}

\begin{enonce*}{Proposition 1.1.14}\etiquette[1.1.14]{XVII.prop1.1.14}
\begin{enumerate}
\renewcommand{\theenumi}{\roman{enumi}}
\renewcommand{\labelenumi}{(\theenumi)}
\item $H^{i}(\tau_{\leq n}(K))=0$\pageoriginale for $i>n$, and $H^{i}(\tau_{\geq
n}(K))=0$ for $i<n$.

\item The canonical morphism $\tau_{\leq n}(K)\rightarrow K$ (\resp
$K\rightarrow \tau_{\geq n}K$) induces isomorphisms on the $H^{i}$
for $i\leq n$ (\resp $i\geq n$).
\end{enumerate}
\end{enonce*}

The functor $\tau_{\leq n}$ (\resp $\tau_{\geq n}$) sends homotopic
morphisms to homotopic morphisms. It therefore follows from
\sisi{\hyperref[XVII.subsubsec1.1.4]{1.1.4}}{\hyperref[XVII.prop1.1.14]{1.1.14}} that the functor of \hyperref[XVII.def1.1.13]{1.1.13} remains
meaningful in the derived category.

\setcounter{subsubsection}{14}
\subsubsection{}\etiquette[1.1.15]{XVII.subsubsec1.1.15}
If $K$ is a double complex, with first (\resp second) differential
$d'$ (\resp $d''$), denote by $\tau''_{\leq n}(K)$ the double
subcomplex of $K$ such that
\begin{equation*}
\tau''_{\leq n}(K)^{pq}=
\begin{cases}
K^{pq}&\mbox{\rm if } q<n\\
\Ker(d'') &\mbox{\rm if } q=n\\
0 &\mbox{\rm if } q>n
\end{cases}
\end{equation*}
and define $\tau''_{\geq n}(K)$, $\tau'_{\leq n}(K)$, and
$\tau'_{\geq n}(K)$ similarly.

\subsubsection{}\etiquette[1.1.16]{XVII.subsubsec1.1.16}
We shall also have occasion to consider the ``stupid truncations'' of
a complex $K$, defined by
\begin{xalignat*}{2}
\sigma_{\leq n}(K)^{i}=
\begin{cases}
K^{i} &\mbox{\rm if } i\leq n\\
0 &\mbox{\rm if } i>n
\end{cases}
&&
\sigma_{\geq n}(K)^{i}=
\begin{cases}
K^{i} &\mbox{\rm if } i\geq n\\
0 &\mbox{\rm if } i<n{\quoi}.
\end{cases}
\end{xalignat*}

\subsection{Derived functors}\etiquette[1.2]{XVII.subsec1.2}\pageoriginale

Let $\Aa$ be a triangulated category and $S$ a saturated
multiplicative system (\cite{sga4.exp.XVII.bib.CD} \sisi{}{I }2.1.2). If $A\in \Ob \Aa$, denote
by $A^{+}$ and $A^{-}$ the \iind- and \ppro-objects
((\hyperref[XVII.subsec0.5]{0.5}) and (\hyperref[XVII.subsec0.6]{0.6}))
\begin{equation*}
A^{+}=\quotinjlim_{A\xrightarrow{s}A'} A' \qquad A^{-}=\quotprojlim_{A'\xrightarrow{s}A}A'{\quoi},
\end{equation*}
where $s$ ranges over the filtered category of arrows in $S$ with
source or target $A$, respectively.

The functors $A\mapsto A^{+}$ and $A\mapsto A^{-}$ make the elements
of $S$ invertible and define fully faithful functors from
$\Aa(S^{-1})$ to the categories of Ind- and \ppro-objects of $\Aa$,
respectively.

\begin{enonce*}{Definition 1.2.1}\etiquette[1.2.1]{XVII.def1.2.1}
Let $F$ be an exact covariant multifunctor from triangulated categories
$(\Aa_i)_{0<i\leq n}$ to the triangulated category $\Aa$, and
let $S_i$ and $S$ be saturated multiplicative systems in the
categories $\Aa_i$ and $\Aa$.
\begin{enumerate}
\renewcommand{\theenumi}{\roman{enumi}}
\renewcommand{\labelenumi}{(\theenumi)}
\item The right derived functor $RF$ of $F$ (relative to the $S_i$ and
$S$) is the functor from the categories $\Aa_i(S^{-1}_i)$ to the
category of \iInd-objects of $\Aa(S^{-1})$ that makes the following
diagram commutative.
\[
\xymatrix@C=5em{%csy
  (\Aa_i)  \ar[d]\ar[r]^{(A_i)\sisi{\simarrow}{\mapsto}  F(A^{+}_i)} & \Ind
  \Aa\ar[d] \\
\Aa_i(S^{-1}_i)\ar[r]^{RF} & \Ind\Aa(S^{-1}).
}
\]

\item We\pageoriginale say that $RF$ is defined at $(A_i)$ if the
\iInd-object $RF(A_i)$ of $\Aa(S^{-1})$ is ``essentially constant,''
\ie comes from an object of $\Aa(S^{-1})$; this object is called
\emph{the value of} $RF$ at $(A_i)$.

\item We say that $F$ is \emph{right derivable} if $RF$ is defined
everywhere; in that case we shall again denote by $RF$ the resulting
functor from the $\Aa_i(S^{-1}_i)$ to $\Aa(S^{-1})$.

\item A family $(A_i)$ of objects of the $\Aa_i$ is said to be
\emph{adapted to} $RF$ if the canonical morphism from $F(A_i)$ to
$RF(A_i)$ is an isomorphism.
\end{enumerate}
\end{enonce*}

We leave to the reader the task of defining by duality the left
derived functor $LF$ of $F$ (with values in $\Pro \Aa(S^{-1})$) and
of extending the preceding definitions to cases in which $F$ is
covariant in some arguments and contravariant in others.

When $RF$ is defined everywhere, the definition given here agrees
with Verdier's (\cite{sga4.exp.XVII.bib.CD} p.~39), as is readily verified.

\begin{enonce*}{Proposition 1.2.2}\etiquette[1.2.2]{XVII.prop1.2.2}
Let $(A'_1, A_1,A''_1)$ be a distinguished triangle in $\Aa_1$, and
let\break $A_i\in \Ob \Aa_i$ $(i\neq 1)$.
\begin{enumerate}
\renewcommand{\theenumi}{\roman{enumi}}
\renewcommand{\labelenumi}{(\theenumi)}
\item If $(A'_{1},(A_{i}))$ and $(A''_{1},(A_{i}))$ are adapted to
  $RF$, then the family $(A_{1},(A_{i}))$ is also adapted to $RF$.

\item If $RF$ is defined at $(A'_1, (A_i))$ and at $(A''_1, (A_i))$,
it is also defined at $(A_1, (A_i))$, and the triangle
$(RF(A'_1,(A_i)), RF(A_1,(A_i)),RF(A''_1,(A_i)))$ is distinguished in
$\Aa(S^{-1})$.
\end{enumerate}
\end{enonce*}

Let us prove\pageoriginale (ii) $\Rightarrow$ (i). Since $F$ is exact,
the arrow $$(F(A'_1,(A_i)), F(A_1,(A_i)), F(A''_1,(A_i)))\rightarrow
(RF(A'_1,(A_i)), RF(A_1,(A_i)), RF(A''_1,(A_i)))$$ is a morphism of
distinguished triangles. By hypothesis this arrow is an isomorphism
at two vertices; it is therefore an isomorphism at the third, as
desired.

\begin{enonce*}{Lemma 1.2.2.1}\etiquette[1.2.2.1]{XVII.lem1.2.2.1}
If $\Delta=(X,Y,Z)$ is a distinguished triangle in a triangulated
category $\Aa$ and $S$ is a saturated multiplicative system in $\Aa$,
then the triangle $\Delta^{+}=(X^{+},Y^{+},Z^{+})$ in $\Ind\Aa$ is a
limit of distinguished triangles in $\Aa$.
\end{enonce*}

Let $\mathcal{L}$ denote the category of morphisms of distinguished
triangles with source $\Delta$ whose arrows belong to $S$.
\begin{enumerate}
\renewcommand{\theenumi}{\alph{enumi}}
\renewcommand{\labelenumi}{(\theenumi)}
\item If $\Delta=(X,Y,Z)$ and $\Delta'=(X',Y',Z')$ are two
distinguished triangles in $\Aa$, and if $f:\Delta\rightarrow \Delta'$
is an $\Aa(S^{-1})$-morphism of triangles, there exist in $\Aa$
morphisms of triangles $s:\Delta'\rightarrow\Delta''$ and
$f_1:\Delta\rightarrow \Delta''$ such that the components of $s$
belong to $S$ and the diagram
\[
\vcenter{
\xymatrix{
\Delta                \ar[rd]^{f_1}  \ar[d]_{f}  &\\
\Delta'          \ar[r]^{s}    & \Delta''
}}\tag{1}
\]
is commutative in $\Aa(S^{-1})$.

% END INCLUDED SGA4_English_current_buildable_source_20260906/source/17/drafts/SGA4_XVII_LINES_601_1000_English.texfrag

% BEGIN INCLUDED SGA4_English_current_buildable_source_20260906/source/17/drafts/SGA4_XVII_LINES_1001_1400_English.texfrag

\sisi{\noindent}{}%
There is a commutative diagram
\[
\vcenter{
\xymatrix{
X \ar[rd]^{f_1} \ar[d]_{f}&  &\\
X' \ar[r]^{s\sisi{}{_X}} & X''\sisi{}{_0}&
}}(s\sisi{}{_X} \in S, f_1 \in \Fl \Aa)\tag{2}
\]
and\pageoriginale a morphism $\Delta'\xrightarrow{s\sisi{_1}{}}\Delta''\sisi{_1}{_0}$
in $S$, having $s\sisi{}{_X}$ as a component. Argue similarly at the
other vertices: one obtains a diagram of type (1), in which, however,
$f_1$ is not yet a morphism of triangles \emph{in}
$\Aa$, since the diagrams
\[
\vcenter{
\xymatrix{
Z  \ar[r]^{f\sisi{}{_1}} \ar[d]_{w}& Z''_{0} \ar[d]_{w\sisi{'}{''}}\\
X         \ar[r]^{f\sisi{}{_1}} & X''_{0}
}}\tag{3}
\]
may, for example, fail to be commutative \emph{in} $\Aa$,
being commutative only in $\Aa(S^{-1})$. There exists
$t\sisi{_0}{_X}:X''_0\rightarrow X''_1$ in $S$ such that $t\sisi{}{_X}w''f\sisi{}{_1}=t\sisi{}{_X}f\sisi{}{_1}w$
and $t:\Delta''_0\rightarrow \Delta''_1$ having $t\sisi{_0}{_X}$
as a component. Replace $\Delta''_0$ by $\Delta''_1$,
$s$ by $ts$, and $f_1$ by $tf_1$; in the resulting diagram of type (1),
(3) is now commutative; proceeding similarly at the other
vertices, one obtains (1). We leave it to the reader to verify
similarly:

\item If $f$, $g:\Delta\rightrightarrows \Delta'$ are two morphisms
of distinguished triangles in $\Aa$, equal in
$\Aa(S^{-1})$, then there exists a morphism of distinguished
triangles whose arrows are in $S$, $s:\Delta'\rightarrow \Delta''$, such
that $\sisi{S}{s}f=\sisi{S}{s}g$.

\item The category $\Ll$ is filtering on the right.

\item The proof of the lemma is completed by noting that
\[
\Delta^{+}=\quotinjlim_{s\in \Ll}{\mbox{\rm (target
of $s$).}}
\]
\end{enumerate}

Let us prove\pageoriginale \hyperref[XVII.prop1.2.2]{1.2.2} (ii). Let $X'$ and $X''$ be objects of
$\Aa(S^{-1})$ representing $RF(A'_1,(A_i))$ and $RF(A''_1,
(A_i))$, and complete the morphism of degree 1 from $X''$ to $X'$ to
a distinguished triangle $(X',X,X'')$. Replacing the $A_i$, $A'_1$,
$A_1$, and $A''_1$ by the target of an arrow in \sisi{$S$}{$S_i$ or $S_1$} with source $A_i$,
$A'_1$, $A_1$, or $A''_1$, we reduce to the case where there is a
commutative diagram
\[
\xymatrix{
F(A''_1, (A_i)) \ar[d] & X''  \ar[l]_-{a''} \ar[d]\\
F(A'_1, (A_i))    & X'  \ar[l]_-{a'}\rlap{,}
}
\]
where $a''$ and $a'$ are sections of arrows from $F(A''_1,(A_i))$ to
$X''$ and from $F(A'_1,(A_i))$ to $X'$. By axiom $TR$ 3 of
\cite{sga4.exp.XVII.bib.CD}, this diagram can be completed to a morphism of triangles
\begin{align*}
(X'XX'')& \longrightarrow
(F(A'_1,(A_i)),F(A_1,(A_i)),F(A''_1,(A_i)))\\
& \longrightarrow
(RF(A'_1,(A_i)),RF(A_1,(A_i)),RF(A''_1,(A_i))){\quoi}.
\end{align*}

The composite arrow is a morphism of triangles whose source is a
distinguished triangle and whose target is a limit (in
$\Ind\Aa(S^{-1})$) of distinguished triangles (\hyperref[XVII.lem1.2.2.1]{1.2.2.1}). For any
$Y$ in $\Aa(S^{-1})$, the following diagram is
commutative and its rows are exact:
\[
\text{\tiny
$\xymatrix@C=1em{
\Hom\sisi{'}{(}Y, X''[\sisi{}{-}1]) \ar[r]\ar[d]& \Hom(Y,X') \ar[r] \ar[d]& \Hom(Y,X)
\ar[r]\ar[d]& \Hom(Y,X'') \ar[r] \ar[d]&
\Hom(Y,X'[1]) \ar[d]\\
\Hom(Y,RF(A''\sisi{}{_1},A_i)[\sisi{}{-}1])\ar[r] & \Hom(Y,RF(A'_1,A_i)\sisi{}{)}\ar[r] &  \Hom(Y,RF(A_1,A_i))\ar[r] &
\Hom(Y,RF(A''_1,A_i)\sisi{}{)}\ar[r] & \Hom(Y,RF(A'_1,A_i)[1])
}$}
\]
By\pageoriginale assumption, all arrows other than the middle arrow
are isomorphisms. The \hyperref[XVII.sec5]{5}-lemma shows that the middle arrow
is also an isomorphism, for every $Y$, and hence that the distinguished
triangle $(X',X,X'')$ is isomorphic to the triangle ($RF(A'_1,(A_i))$,
$RF(A_1,(A_i))$, $RF(A''_1,(A_i))$). This completes the proof.

If $\Aa$ and $\Bb$ are two abelian categories and
$F$ is an additive covariant functor from $\Aa$ to $\Bb$,
$F$ defines exact functors from $K(\Aa)$ to
$K(\Bb)$, from $K^{+}(\Aa)$ to $K^{+}(\Bb)$, and
from $K^{-}(\Aa)$ to $K^{-}(\Bb)$ (\hyperref[XVII.prop1.1.7]{1.1.7}).

\begin{enonce*}{Definition 1.2.3}\etiquette[1.2.3]{XVII.def1.2.3}
\begin{enumerate}
\renewcommand{\theenumi}{\roman{enumi}}
\renewcommand{\labelenumi}{(\theenumi)}
\item The \emph{right derived functors of $F$} are the functors
from $D^{+}(\Aa)$ to $\Ind D^{+}\sisi{\Bb}{(\Bb)}$ and from
$D(\Aa)$ to $\Ind D(\Bb)$ right derived from the
extensions of $F$ to complexes $K^{+}(\Aa)\rightarrow
K^{+}(\Bb)$ and $K(\Aa)\rightarrow K(\Bb)$.

\item An object $A$ of $\Aa$ is said to be \emph{right acyclic}
for $F$, or, by abuse of language, \emph{acyclic for} $RF$, if the
complex concentrated at $A$ in degree $0$ is adapted to $RF$.
\end{enumerate}
\end{enonce*}

We shall be interested mainly in the case where $RF$ is defined everywhere; it
may then be regarded as a functor from $D^{+}(\Aa)$ to
$D^{+}(\Bb)$, or from $D(\Aa)$ to $D(\Bb)$, as the case may be.

\begin{enonce*}{Proposition 1.2.4}\etiquette[1.2.4]{XVII.prop1.2.4}
The following diagram is commutative.
\[
\xymatrix{
D^{+}(\mathcal{A})  \ar[r] \ar[d]^{RF}& D(\mathcal{A})\ar[d]^{RF}        \\
 \Ind \;D^{+}(\mathcal{B}) \ar[r] & \Ind \; D(\mathcal{B}).
}
\]
\end{enonce*}

If\pageoriginale $K$ is a complex that is zero in degrees $<n$ and
$s:K\rightarrow L$ is a quasi-isomorphism, the composite morphism
$K\xrightarrow{s}L\rightarrow \sigma_{\geq n}(L)$
(\hyperref[XVII.subsubsec1.1.6]{1.1.6}) is again a quasi-isomorphism. The
quasi-isomorphisms from $K$ to a complex bounded below therefore
form a cofinal (full) subcategory of the category of
all quasi-isomorphisms with source $K$. The proposition follows
formally.

In particular, if the functor $RF$\sisi{,}{} is defined everywhere on
$D(\Aa)$, it is defined everywhere on $D^{+}(\sisi{\Bb}{\Aa})$, and
the derived functor $RF:D^{+}(\Aa)\rightarrow
D^{+}(\Bb)$ is the restriction of the derived functor
$RF:D(\Aa)\rightarrow D(\Bb)$.

\setcounter{subsubsection}{4}
\subsubsection{}\etiquette[1.2.5]{XVII.subsubsec1.2.5}
We leave it to the reader to:
\begin{enumerate}
\renewcommand{\theenumi}{\roman{enumi}}
\renewcommand{\labelenumi}{(\theenumi)}
\item Generalize Definition \hyperref[XVII.def1.2.3]{1.2.3} to the case of a complex of
multifunctors, covariant in some variables and contravariant in
others;

\item Dualize Definition \hyperref[XVII.subsubsec1.2.5]{1.2.5} (i) to define the left derived
functors;

\item Extend Proposition \hyperref[XVII.prop1.2.4]{1.2.4} to the case of a bounded complex of
multifunctors;

\item Dualize Proposition \hyperref[XVII.subsubsec1.2.5]{1.2.5} (iii) by
considering left derived functors and the inclusion of
$D^{-}(\Bb)$ into $D(\Bb)$.
\end{enumerate}

\begin{enonce*}{Proposition 1.2.6}\etiquette[1.2.6]{XVII.prop1.2.6}
Let $\Aa$ and $\Bb$ be two abelian categories,
let $𝒫$ be a subset of $\Ob(\Aa)$, and let $F$ be an additive functor from
$\Aa$ to $\Bb$. Assume that
\begin{enumerate}
\renewcommand{\theenumi}{\alph{enumi}}
\renewcommand{\labelenumi}{(\theenumi)}
\item every object of $\Aa$ is a quotient of an element of $𝒫$;

\item for every short exact sequence
\[
0\longrightarrow P\longrightarrow Q\longrightarrow R\longrightarrow
0{\quoi},
\]
where $Q$ and $R$ belong to $𝒫$, $P$ also belongs to $𝒫$, and the
$0$-sequence $0\rightarrow FP\rightarrow FQ\rightarrow FR\rightarrow
0$\pageoriginale is exact. Then every object of $𝒫$ is
acyclic for $LF$.
\end{enumerate}
\end{enonce*}

If $A\in \Ob \Aa$, the left resolutions of $A$ by
objects of $𝒫$ form a cofinal system in the filtering
category of left resolutions of $A\sisi{_0}{}$. It therefore suffices to
prove that if a complex
\[
\ldots\longrightarrow A_2\longrightarrow A_1\longrightarrow A_0
\longrightarrow A\longrightarrow 0
\]
is acyclic and has its objects in $𝒫$, its image is acyclic. Splitting it
into short exact sequences
$$\etiquette[1.2.6.i]{XVII.eq1.2.6.i}
0\longrightarrow \Iim(d_{i-1})\longrightarrow A_i\longrightarrow
\Iim(d_i)\longrightarrow 0\qquad\sisi{(1.2.6. i)\quad}{}(i\geq 0)\sisi{}{\leqno(1.2.6.i)}%\tag{1.2.6.$i$}
$$
one sees by induction on $i$ that $\Iim(d_i)\in \PP$ and that the image
of \sisi{(\hyperref[XVII.prop1.2.6]{1.2.6} $i$)}{\eqhyperref[XVII.eq1.2.6.i]{1.2.6.i}} under $F$ is a short exact sequence. The image of the complex
is therefore acyclic.

\begin{enonce*}{Proposition 1.2.7}\etiquette[1.2.7]{XVII.prop1.2.7}
Let $F$ be an additive functor from an abelian category to an
abelian category $\Bb$. If every object of $\Aa$ is a
quotient (\resp subobject) of an object acyclic for $LF$ (\resp
$RF$) (\hyperref[XVII.def1.2.3]{1.2.3}), the derived functor $LF:D^{-}(\Aa)\rightarrow
D^{-}(\Bb)$ (\resp $RF:D^{+}(\Aa)\rightarrow
D^{+}(\Bb)$) is defined everywhere, and every complex bounded
above (\resp below) of objects acyclic for $LF$
(\resp $RF$) is adapted.
\end{enonce*}

It suffices to consider the left derived functor $LF$. If $K$ is a complex
bounded above, the quasi-isomorphisms $s:L\rightarrow K$ with
$L$ bounded above and consisting of acyclic objects form a
cofinal system in the category of quasi-isomorphisms with target $K$. It
suffices to show that, when $K$ consists of acyclic objects,
$F(s):F(L)\rightarrow F(K)$ is a quasi-isomorphism;\pageoriginale
then, in the general case, one will have $F(L)=LF(K)$. Let $M$ be the cone of
$s$. The complex consists of acyclic object\sisi{}{s} and $H^{\ast}(M)=0$; one must
prove that $H^{\ast}(F(M))=0$.

Let $i$ be an integer. The adapted complexes form a triangulated
subcategory, so that every bounded complex of acyclic objects,
in particular $\tau_{\geq i}(M)$, is adapted. Since $H^{j}(\tau_{\geq
i}(\sisi{L}{M}))=0$ for $j>i$, one has, for $j>i$,
$H^{j}(F(M))=H^{j}(F(\tau_{\geq i}(M)))=H^{j}LF(\tau_{\geq i}(M))=0$.

\begin{enonce*}{Proposition 1.2.8}\etiquette[1.2.8]{XVII.prop1.2.8}
Let $F^{\bigdot}$ be a complex bounded below of functors from
$\Aa$ to $\Bb$. If for every $K\in
K^{+}(\Aa)$ there exists a quasi-isomorphism $L\xrightarrow{s}K$,
where $L\in K^{+}(\Aa)$ is adapted to all the $RF^{i}$, then
the derived functor $RF^{\bigdot}:D^{+}(\Aa)\rightarrow
D^{+}(\Bb)$ exists, and every complex \emph{bounded
below} that is adapted to all the $RF^{i}$ is adapted to $RF^{\bigdot}$.
\end{enonce*}

\sisi{\noindent}{}%
The proof is left to the reader, as is that of the dual statement,
or of the variants obtained by taking $F$ contravariant, or by taking
a bounded complex of functors and working in $D(\Aa)$
and $D(\Bb)$.

\begin{enonce*}{Definition 1.2.9}\etiquette[1.2.9]{XVII.def1.2.9}
Let $F$ be a functor from $D^{b}(\Aa)$ to $D(\Bb)$, and let
$d$ be an interval in $\ZZ$. We say that $F$ has
\emph{cohomological amplitude $\subset d$} if, for every $A\in \Ob
\Aa$, $H^{i}F(A)=0$ for $i\not\in d$.
\end{enonce*}

Very often, $d$ will be an interval of the form $[0,x]$ or $[-x,0]$,
and we then say that $F$ has cohomological \resp
homological dimension $\leq x$. We say that $F$ has finite dimension if
there exists a finite interval $d$ in $\ZZ$ such that $F$ has
amplitude $\subset d$.

\begin{enonce*}{Proposition 1.2.10}\etiquette[1.2.10]{XVII.prop1.2.10}
Assume\pageoriginale that the hypotheses of \hyperref[XVII.prop1.2.7]{1.2.7} hold and that the
functor $RF:D^{+}(\Aa)\rightarrow D^{+}(\Bb)$ (\resp
$LF:D^{-}(\Aa)\rightarrow D^{-}(\Bb)$) has
finite dimension. Then every complex of objects acyclic for $RF$
(\resp $LF$) is adapted, and the functor $RF:D(\Aa)\rightarrow
D(\Bb)$ (\resp $LF:D(\Aa)\rightarrow D(\Bb)$) is defined everywhere.
\end{enonce*}

For the proof, we refer to Verdier \cite{sga4.exp.XVII.bib.CD}. There is an analogous statement
for bounded complexes of functors. The superscript $+$ (\resp $-$)
may also be replaced by $b$.

\section{Categories fibered in derived categories}\etiquette[2]{XVII.sec2}

\subsection{Introduction}\etiquette[2.1]{XVII.sec2.1}

The editor insists that the reader refrain from reading this \S. Here
we give the formalism of the trivial duality theorem and
resolve the perplexity raised by Artin in \hyperrefext[XII][XII.sec4]{4}.

For the notions of a category over another category, a fibered category, a
cofibered category, and a cleavage, we refer to \SGA 1 VI. Among other things,
that exposé establishes an equivalence between the notions of
normalized split fibered categories over a category $\sisi{C}{\Cc}$, and
pseudofunctors from $\sisi{C}{\Cc}^{\circ}$ to $(\Cat)$.

\begin{enonce*}{Definition 2.1.1}\etiquette[2.1.1]{XVII.def2.1.1}
If $X$ is a ringed site, $D(X)$ denotes the derived category
of the abelian category of sheaves of left modules on $X$.
\end{enonce*}

When confusion must be avoided, we instead write
$D(X,\Aa)$, where $\Aa$ denotes the sheaf of rings.

If\pageoriginale $f:X\rightarrow Y$ is a morphism of ringed sites,
$f$ defines derived functors
\begin{itemize}
\item[(1)] $Rf_{\ast}:D^{+}(X)\rightarrow D^{+}(Y)$ and

\item[(2)] $Lf^{\ast}:D^{-}(Y)\rightarrow D^{-}(X)$,
\end{itemize}
and, on passing to the opposite categor\sisi{y}{ies}, the functors
\begin{itemize}
\item[$\rm(1^{\circ})$] $Rf_{\ast}:D^{+}(X)^{\circ}\rightarrow
D^{+}(Y)^{\circ}$

\item[$\rm(2^{\circ})$] $Lf^{\ast}:D^{-}(Y)^{\circ}\rightarrow
D^{-}(X)^{\circ}$,
\end{itemize}
``functorial'' in $f$. More precisely, $(1^{\circ})$ and
$(2^{\circ})$ define two pseudofunctors from the category of
sites to that of triangulated categories, one covariant and the other
contravariant. One thereby obtains two categories over the category of sites,
one cofibered, with fibers $D^{+}(X)^{\circ}$ and direct-image functors
$Rf_{\ast}$, and the other fibered, with fibers
$D^{-}(X)^{\circ}$ and inverse-image functors
$\sisi{\LL}{L} f^{\ast}$. If the functor $f^{\ast}$ has finite
homological dimension (\ie has finite ``\tTor-dimension''), then in
(2) the superscript may be replaced by a $+$. Thus one obtains two
categories over the category $\Ss_{\tdf}$ o\sisi{f}{f} ringed sites and
morphisms of ringed sites of finite \tTor-dimension; the fibers of these
categories are the categories $D^{+}(X)^{\circ}$. The ``trivial
duality theorem'' of \textsc{Verdier} (\cite{sga4.exp.XVII.bib.CD} p.~48) states that these two
categories are canonically isomorphic; they are therefore fibered and
cofibered; the ``direct-image'' functors are the functors
$Rf_{\ast}$\sisi{}{ }(sic), and the ``inverse-image'' functors are the
functors $Lf^{\ast}$\sisi{}{ }(sic).

\setcounter{subsection}{1}
\setcounter{subsubsection}{1}
\subsubsection{}\etiquette[2.1.2]{XVII.subsubsec2.1.2}

The language of fibered categories readily resolves the
perplexity raised by \textsc{Artin} in \hyperrefext[XII][XII.sec4]{4}: it suffices to apply the
following proposition to the category fibered and cofibered over the category of
sites\pageoriginale whose fibers are the opposites of the categories of sheaves
on the corresponding sites (\cf \hyperlink{sga4.exp.XVII.subsection.0.2}{0.2}).

Let $\Ee$ be a category over $\Bb$, let $f:y\rightarrow x$
be an arrow of $\Bb$, and let $F\in \Ob(\Ee_x)$. Recall
that a pair $(G,\varphi)$ $(G\in \Ob \Ee_y, \varphi\in
\Hom_f(G,F))$ is a \emph{strict inverse image} of $F$
under $f$ if, for all $g:\sisi{Z}{z}\rightarrow y$ and $H\in \Ob
\Ee\sisi{_Z}{_z}$, one has
\[
\Hom_g(H,G)\xrightarrow[\sim]{\varphi}\Hom_{fg}(H,F){\quoi}.
\]
Similarly for direct images. If $\Ee$ is fibered or cofibered
over $\Bb$, this notion reduces to that of inverse image.

\begin{enonce*}{Proposition 2.1.3}\etiquette[2.1.3]{XVII.prop2.1.3}
Let $\epsilon$ be a category over $\Bb$, let $y\in \Ob
\Bb$, let $F\in \Ob \Ee_y$, and consider a commutative diagram in
$\Bb$:
$$\etiquette[2.1.3.1]{XVII.eq2.1.3.1}
\vcenter{
\xymatrix{
y'  \ar[r]^{g'} \ar[d]_{f'}& y\ar[d]^{f}         \\
x'          \ar[r]    & x.
}}\leqno{(2.1.3.1)}
$$
Assume that the direct and inverse images $f_{\ast}F$,
$g^{\ast}f_{\ast}F$, ${g'}^{\ast}F$, and $f'_{\ast}{g'}^{\ast}F$
exist in the strict sense. Then there exists a unique arrow
$\varphi \in \Fl \Ee_{x'}$, called the \emph{base-change
arrow}, making the diagram
$$\etiquette[2.1.3.2]{XVII.eq2.1.3.2}
\vcenter{
\xymatrix{
                             &  g'^{\ast}F \ar[ld]  \ar[r]    & F \ar[d]\\
f'_{\ast}g'^{\ast}F \ar[rd]_{\varphi}  &                                & f_{\ast}F\\
  & g^{\ast}f_{\ast}F \ar[ru]&.
}}\leqno{(2.1.3.2)}
$$
commutative. If\pageoriginale, in addition, the direct images $g'_{\ast}{g'}^{\ast
}F$ and $g'_{\ast}f'_{\ast}{g'}^{\ast}F$ exist in the strict sense, then the arrow $\varphi$ coincides with
the one defined in \emph{\hyperrefext[XII][XII.sec4]{4}}. If ${f'}^{\ast}f'_{\ast}{g'}^{\ast}F$
and ${f'}^{\ast}f'_{\ast}F$ exist in the strict sense, it also coincides
with the $2$\sisi{nd}{\textsuperscript{nd}} arrow defined in \emph{\hyperrefext[XII][XII.sec4]{4}}.

In particular, the two base-change arrows of
\emph{\hyperrefext[XII][XII.sec4]{4}} coincide.
\end{enonce*}

By hypothesis, the arrows in diagram \eqhyperref[XVII.eq2.1.3.2]{2.1.3.2} induce
bijections
\sisi{\begin{align*}
& \Hom_y(F,F)=\Hom_{g'}({g'}^{\ast}F,F)=\Hom_{fg'}({g'}^{\ast}F,f_{\ast}F)
\quad\mbox{\rm and}\\
&
\Hom_{x'}(f'_{\ast}{g'}^{\ast}F,g^{\ast}f_{\ast}F)=\Hom_{f'}({g'}^{\ast}F,
g^{\ast}f_{\ast}F)=\Hom_{gf'}({g'}^{\ast}F, f_{\ast}F){\quoi}.
\end{align*}}{\[\Hom_{x'}(f'_{\ast}{g'}^{\ast}F,g^{\ast}f_{\ast}F)=\Hom_{f'}({g'}^{\ast}F,
g^{\ast}f_{\ast}F)=\Hom_{gf'}({g'}^{\ast}F, f_{\ast}F).\]}
Since $fg'=gf'$, $\Hom_{fg'}({g'}^{\ast}F,f_{\ast}F)$ is identical
to $\Hom_{gf'}({g'}^{\ast}F,f_{\ast}F)$, and the arrows of diagram
\sisi{(2.1.2.2)}{\eqhyperref[XVII.eq2.1.3.2]{2.1.3.2}} induce \sisi{bijections}{maps}
\[
\Hom_y(F,F)\sisi{=}{\to}\Hom_{fg'}({g'}^{\ast}F,f_{\ast}F)=\Hom_{x'}(f'_{\ast}{g'}^{\ast}F,g^{\ast}f_{\ast}F){\quoi},
\]
and the image of $1_F$ is the unique arrow of $\Ee_{x'}$ making

% END INCLUDED SGA4_English_current_buildable_source_20260906/source/17/drafts/SGA4_XVII_LINES_1001_1400_English.texfrag

% BEGIN INCLUDED SGA4_English_current_buildable_source_20260906/source/17/drafts/SGA4_XVII_LINES_1401_1800_English.texfrag
\eqhyperref[XVII.eq2.1.3.2]{2.1.3.2} commutative.%See the Erratum in SGA 4 1/2, p. 312.

By definition, in \hyperrefext[XII][XII.sec4]{4}, the base change arrow $\psi$
was the one making the following diagram commutative
$$\etiquette[2.1.3.3]{XVII.eq2.1.3.3}
\vcenter{
\xymatrix{
f'_{\ast}g'{}^{\ast}F \ar[d]^{\psi}  & g'{}^{\ast}F \ar[d]^{u_4} \ar[l] \ar@{=}[r] &
g'{}^{\ast}F \ar[r] \ar[d]^{u_3}   & F \ar@{=}[r] \ar[d]^{u_2}&  F \ar[d]^{u_1}\\
g'_{\ast}f_{\ast}F   & {f'}^{\ast}g^{\ast}f_{\ast}F  \ar[l] \ar@{=}[r] &
g'f^{\ast}f_{\ast}F   \ar[r] & f^{\ast}f_{\ast}F \ar[r] &  f_{\ast}F.
}}\leqno{(2.1.3.3)}
$$
The horizontal arrows and $u_1$ are the canonical arrows; $u_2$
is therefore the adjunction arrow, $u_3={g\sisi{^{1}}{'}}^{\ast}(u_2)$, $u_4$ is
deduced from the isomorphism between ${f'}^{\ast}g^{\ast}$ and
${g'}^{\ast}f^{\ast}$, and the arrow $\psi$ is deduced from $u_4$ by
adjunction.

The composite\pageoriginale map with source ${g'}^{\ast}F$ and
target $f_{\ast}F$ is the one that appears in \eqhyperref[XVII.eq2.1.3.2]{2.1.3.2}, and the
same therefore holds in the following commutative diagram, whose left-hand
square is taken from \sisi{(2.1.2.3)}{\eqhyperref[XVII.eq2.1.3.3]{2.1.3.3}}
\[
\xymatrix{
g'{}^{\ast}F \ar[r]^{u_4} \ar[d]& f'{}^{\ast}g{}^{\ast}f_{\ast}F
\ar[r] \ar[d]&
f_{\ast}F \ar@{=}[d]\\
f'_{\ast}g'{}^{\ast}F  \ar[r]^{\psi} & g^{\ast}f_{\ast}F \ar[r] &
f_{\ast}F\rlap{.}
}
\]
This means that $\psi$ satisfies the defining property of
$\varphi$.

The definition of $\varphi$ is self-dual, so there is no need to prove
the second part of \hyperref[XVII.prop2.1.3]{2.1.3}.

\subsection{Categories fibered in triangulated categories}\etiquette[2.2]{XVII.subsec2.2}

\begin{enonce*}{Definition 2.2.1}\etiquette[2.2.1]{XVII.def2.2.1}
An “additive category over a category $\Bb$”
is a category $\Aa$ over $\Bb$ whose sets of
arrows $\Hom_f(X,Y)$ (for $f\in \Fl (\Bb)$, with $X$ (\resp
$Y$) over the source (\resp the target) of $f$) have been endowed with
abelian group structures, in such a way that:
\begin{enumerate}
\renewcommand{\theenumi}{\roman{enumi}}
\renewcommand{\labelenumi}{(\theenumi)}
\item For any $x\xrightarrow{f}y\xrightarrow{g}z$ in
$\Bb$ and $X$, $Y$, $Z$ in $\Aa$ lying over $\sisi{s}{x}$,
$y$, and $z$, respectively, the composition law
\[
\Hom_g(Y,Z)\times \Hom_f
(X,Y)\xrightarrow{\circ}\Hom_{gf}(X,Z)
\]
is biadditive.

\item The fiber categories are additive.
\end{enumerate}
\end{enonce*}

We shall not need here the fact that these abelian group structures, when
they exist, are uniquely determined by
$\Aa$, $\Bb$, and $\rho$.\nde{here $\rho$ denotes the structural functor $\Aa\to \Bb$.} A $\Bb$-functor
$F:\Aa_1\rightarrow \Aa_2$ between additive categories
over $\Bb$ is said to be \emph{additive} if it induces
homomorphisms
\[
F:\Hom_{\Aa_1,f}(X,Y)\longrightarrow
\Hom_{\Aa_{2},f}(X,Y){\quoi}.
\]

\begin{enonce*}{Definition 2.2.2}\etiquette[2.2.2]{XVII.def2.2.2}
An\pageoriginale “abelian category over a category
$\Bb$” is an additive category $\Aa$ over
$\Bb$ whose fibers are abelian, and such that for
every arrow $f:x\rightarrow y$ of $\Bb$ the bifunctor
\[
\Hom_f(X,Y):\Aa^{\circ}_x\times \Aa_y\longrightarrow
(\Ab)
\]
is left exact in $X$ and $Y$.
\end{enonce*}

We shall not need here the fact that if $\Aa$ is a category
fibered and cofibered over $\Bb$ (\SGA 1 VI 6), it is abelian
over $\Bb$ as soon as its fibers are abelian categories.

\begin{enonce*}{Definition 2.2.3}\etiquette[2.2.3]{XVII.def2.2.3}
A “triangulated category over a category $\Bb$” is an additive category $\Aa$ over $\Bb$, endowed
with a functor $T:\Aa\rightarrow \Aa$, and whose
fibers $\Aa_x$, for $x\in \Ob \Bb$, are endowed
with sets $\Delta_x$ of triangles, in such a way that:
\begin{enumerate}
\renewcommand{\theenumi}{\roman{enumi}}
\renewcommand{\labelenumi}{(\theenumi)}
\item $T$ is an additive $\Bb$-automorphism of $\Aa$.

\item Each fiber category $\Aa_x$, endowed with the functor $T_x$
induced by $T$ and with the set $\Delta_x$, is a triangulated
category \emph{(\cite{sga4.exp.XVII.bib.CD} \sisi{}{I 1.}1.1)}.

\item If the triangles $(X,Y,Z,u,v,w)$ and $(X',Y',Z',u',v',w')$
belong to $\Delta_x$ and $\Delta_{x'}$, respectively, every
diagram
\[
\xymatrix{
            & Z   \ar[ldd]_{w}&            &  &           & Z'   \ar[ldd]|\hole_(.3){w'} & \\
            &   & Y\ar[lu]_{v} \ar[rrrr]_(.4){\sisi{j}{g}}&  &     &       &Y'\ar[lu]_{v'} \\
X   \ar[rru]^{u} \ar[rrrr]_{f}&             &          &  &   X'   \ar[rru]^{u'} &              &
}
\]
in which $u'f=gu$ can be completed to a morphism of triangles
$(f,g,h):(X,Y,Z)\rightarrow (X',Y',Z')$.
\end{enumerate}
\end{enonce*}

$T$\pageoriginale will be called the translation functor; we shall
often write $A[1]$ rather than $T(A)$, with $A[n]$ then denoting $T^{n}(A)$
$(n\in\ZZ)$.

The triangles belonging to the sets $\Delta_x$ will be called
\emph{distinguished} or \emph{exact}.

\begin{enonce*}{Proposition 2.2.4}\etiquette[2.2.4]{XVII.prop2.2.4}
Let $\Aa$ be a triangulated category over $\Bb$, let $x$
be an object of $\Bb$, and let $(X,Y,Z,u,v,w)$ be a distinguished triangle of
$\Aa_x$. For every $f:y\rightarrow x$ and every object $A\in
\Ob \Aa_y$, the following infinite sequence is exact:
\begin{align*}
\ldots  \longrightarrow \Hom_f(A,X) \longrightarrow \Hom_f(A,Y)&\longrightarrow
\Hom_f(A,Z)\\
 & \longrightarrow \Hom_f(A,X[1])\longrightarrow
\Hom_f(A,Y[1])\longrightarrow\ldots
\end{align*}
\end{enonce*}

There is an analogous statement for $f:x\rightarrow y$ and $A\in \Ob \Aa_y$.

We leave it to the reader to prove this proposition by
paraphrasing \textsc{Verdier} \cite{sga4.exp.XVII.bib.CD} p.~4. We also leave it to the reader to give
a meaning to the following proposition and to verify it:

\begin{enonce*}{Proposition 2.2.5}\etiquette[2.2.5]{XVII.prop2.2.5}
Let $\Aa$ be an additive (\resp abelian, \resp
triangulated) category over $\Bb$. For every functor
$\Bb'\rightarrow \Bb$ with target $\Bb$, the
category $\Aa'=\fprod{\Aa}{\Bb'}{\Bb}$ is
additive (\resp abelian, \resp triangulated) over $\Bb$.
\end{enonce*}

\begin{enonce*}[remark]{Example 2.2.6}\etiquette[2.2.6]{XVII.exem2.2.6}
The category of sheaves of modules on varying ringed sites
is abelian over the category of ringed sites.
\end{enonce*}

\begin{enonce*}[remark]{Example 2.2.7}\etiquette[2.2.7]{XVII.exem2.2.7}
Let $\Aa$ be an additive category over $\Bb$. The
category $C(\Aa/\Bb)$ of complexes of $\Aa$
over $\Bb$ has as objects the complexes of the fiber
categories. An arrow $f$ from $(X^n,d)$ to $(Y^n,d)$ is a system
of\pageoriginale arrows $f^{n}:X^{n}\rightarrow Y^{n}$ such that
$\sisi{fd}{df}^{n}=f^{n+1} d$. We set $p(f)=p(f^{n})$ (this arrow does not depend
on $n$, since the $p(d)$ are
identities). $C(\Aa/\Bb)$ is an additive category
over $\Bb$. A “translation” functor is defined by setting
\[
(K[1])^{n}=K^{n+1}\quav (d[1])\sisi{_n}{^n}=-d\sisi{_{n+1}}{^{n+1}}{\quoi}.
\]
\end{enonce*}

\begin{enonce*}[remark]{Example 2.2.8}\etiquette[2.2.8]{XVII.exem2.2.8}
Let $\Aa$ be an additive category over $\Bb$. \emph{The
category $K(\Aa/\Bb)$ of complexes of $\Aa$
over $\Bb$ modulo homotopy} has the same objects as
$C(\Aa/\Bb)$. It is obtained by declaring the arrows homotopic
to zero to be zero in $C(\Aa/\Bb)$; in
other words, for $f:x\rightarrow y$ in $\Bb$,
$X^{\bigdot}\in \Ob C(\Aa/\Bb)_x$ and $Y^{\bigdot}\in
\Ob C(\Aa/\Bb)_y$, we set
\[
\Hom_f(X^{\bigdot},
Y^{\bigdot})=H^{0}\Hom^{\bigdot}_f(X^{\bigdot},Y^{\bigdot})
\]
\end{enonce*}

The translation functor descends to the quotient.

If $F:\Bb'\rightarrow \Bb$ is a functor with target
$\Bb$, and if
$\Aa'=\fprod{\Aa}{\Bb'}{\Bb}$, we
leave it to the reader to identify $C(\Aa'/\Bb')$
with $\fprod{C(\Aa/\Bb)}{\Bb'}{\Bb}$ and
$K(\Aa'/\Bb')$ with
$\fprod{K(\Aa/\Bb)}{\Bb'}{\Bb}$. In
particular, the fibers of the category $K(\Aa/\Bb)$
are the categories $K(\Aa_x)$ and, as such, are
triangulated.

\begin{enonce*}{Proposition 2.2.9}\etiquette[2.2.9]{XVII.prop2.2.9}
\begin{enumerate}
\renewcommand{\theenumi}{\roman{enumi}}
\renewcommand{\labelenumi}{(\theenumi)}
\item With the additional structures defined in \hyperref[XVII.exem2.2.8]{2.2.8}, the
category $K(\Aa/\Bb)$ is triangulated over $\Bb$.

\item The formation of the additive (\resp triangulated)
category $C(\Aa/\Bb)$ (\resp $K(\Aa/\Bb)$) is
compatible with every change of base category
$\Bb'\rightarrow \Bb$.
\end{enumerate}
\end{enonce*}

Conditions\pageoriginale (i) and (ii) of Definition \hyperref[XVII.def2.2.3]{2.2.3} are
verified fiber by fiber. Condition (iii) is verified by paraphrasing \textsc{Verdier}
\cite{sga4.exp.XVII.bib.CD} p.~4.

\begin{enonce*}[remark]{Remark 2.2.10}\etiquette[2.2.10]{XVII.rem2.2.10}
The preceding discussion also applies to the categor\sisi{y}{ies}
$K^{+}(\Aa/\Bb)$, $K^{-}(\Aa/\Bb)$, and
$K^{b}(\Aa/\Bb)$ obtained by restricting
respectively to complexes bounded below, bounded above, or bounded.
\end{enonce*}

\begin{enonce*}[remark]{Remark 2.2.11}\etiquette[2.2.11]{XVII.rem2.2.11}
Taking $\Bb$ to be the final category, one recovers the
usual categories of complexes.
\end{enonce*}

\setcounter{subsection}{2}
\setcounter{subsubsection}{11}
\subsubsection{}\etiquette[2.2.12]{XVII.subsubsec2.2.12}
Let $F$ be a pseudofunctor from a category $\Bb^{\circ}$ to
$(\Cat)$. It is known (\SGA 1 VI) that $F$ defines a category fibered over
$\Bb$, whose fibers are the categories $F(x)$. For any
$f:x\rightarrow y$ in $\Bb$ and $X\in \Ob F(\sisi{n}{x})$, $Y\in \Ob
F(y)$, we have by definition
$$\etiquette[2.2.12.1]{XVII.eq2.2.12.1}
\Hom_f(X,Y)=\Hom_{F(x)}(X,F(f)(Y)){\quoi}.\leqno(2.2.12.1)
$$
It is known (ibid.) that the preceding construction defines an
equivalence between the 2-category of pseudofunctors from
$\Bb^{\circ}$ to $(\Cat)$ and the 2-category of categories
fibered over $\Bb$.

Let us start with a pseudofunctor from $\Bb^{\circ}$ to the
2-category whose objects are additive (\resp
abelian, \resp triangulated) categories and whose 1-arrows are additive
(\resp left exact, \resp triangulated) functors. One immediately verifies that the
category defined by \sisi{(2.1.12.1)}{\eqhyperref[XVII.eq2.2.12.1]{2.2.12.1}} is additive (\resp abelian, \resp
triangulated). Conversely:

\begin{enonce*}{Proposition 2.2.13}\etiquette[2.2.13]{XVII.prop2.2.13}
The\pageoriginale preceding construction establishes an equivalence between
\begin{enumerate}
\renewcommand{\theenumi}{\alph{enumi}}
\renewcommand{\labelenumi}{(\theenumi)}
\item The $2$-category of pseudofunctors from $\Bb^{\circ}$
to the $2$-category whose objects are additive
(\resp abelian, \resp triangulated) categories and whose $1$-arrows are
additive (\resp left exact, \resp triangulated) functors.

\item The $2$-category of additive (\resp abelian,
\resp triangulated) categories over $\Bb$ that are fibered over $\Bb$.
\end{enumerate}
\end{enonce*}

The theory developed in \SGA 1 VI shows that it suffices to verify the
following point:

\setcounter{subsubsection}{13}
\setcounter{paragraph}{0}
\paragraph{}\etiquette[2.2.13.1]{XVII.par2.2.13.1}
Let $\Aa$ be an additive (\resp abelian, \resp
triangulated) category over $\Bb$, let $f:x\rightarrow y$ be an arrow of
$\Bb$, and suppose that an inverse image functor
$f^{\ast}:\Aa_y\rightarrow \Aa\sisi{_n}{_x}$ exists; then $f^{\ast}$
is additive (\resp left exact, \resp triangulated).

The functor $f^{\ast}$ gives rise to isomorphisms
\[
\Hom_f(X,Y)\simeq \Hom_x (X,f^{\ast}(Y)){\quoi}.
\]
In the triangulated case\sisi{}{}, it remains to prove that $f^{\ast}$ is triangulated
for the isomorphism $Tf^{\ast}=f^{\ast}T$ (where $T$ is the
translation functor) that makes the following diagrams commutative:
$$\etiquette[2.2.13.2]{XVII.eq2.2.13.2}
\vcenter{
\xymatrix{
\Hom_f(X,Y)  \ar@{=}[r] \ar@{=}[d]& \Hom_x(X,f^{\ast}Y)\ar@{=}[rd]\\
\Hom_f(TX,TY)\ar@{=}[rd]           &      & \Hom_x(TX,Tf^{\ast}Y) \\
             &      \Hom_x(TX,f^{\ast}TY)\ar@{=}[ru] &
}}\leqno{(2.2.13.2)}
$$

\begin{itemize}
\item[\emph{Additive case}:] Consider\pageoriginale a diagram
$$\etiquette[2.2.13.3]{XVII.eq2.2.13.3}
\vcenter{
\xymatrix@C=2.5cm@R=1.75cm{
f^{\ast} X \ar@<-3pt>[d]_{f^{*}(u)}\ar@<3pt>^{f^{\ast}(v)}[d]        \ar[r]^{p} & X
\ar@<-3pt>_{u}[d] \ar@<3pt>^{v}[d]\\
f^{\ast} X' \ar[r]^{p'}& X'\rlap{.}
}}\leqno{(2.2.13.3)}
$$
We have
$p'f^{\ast}(u+v)=(u+v)p=up+vp=p'f^{\ast}(u)+p'f^{\ast}(v)=p'(f^{\ast}(u)+f^{\ast}(v))$
and hence $f^{\ast}(u+v)=f^{\ast}(u)+f^{\ast}(v)$.

\item[\emph{Abelian case}:] by definition (\hyperref[XVII.subsec2.2]{2.2}), the functor
$\Hom(X,f^{\ast}Y)=\Hom_f(X,Y)$ is left exact in $Y$, for every
$X$, whence $f^{\ast}$ is left exact.

\item[\emph{Triangulated case}:] let $(Y',Y,Y'')$ be a distinguished triangle in
$\Aa_y$ and let \sisi{$(f^{\ast}Y',f^{\ast}Y)$}{$(f^{\ast}Y',f^{\ast}Y, X'')$ be a distinguished triangle in $\Aa_x$}. By \hyperref[XVII.def2.2.3]{2.2.3},
the canonical morphism between the bases of these triangles can be
completed to a morphism of triangles, giving rise to the following
diagram:
\[
\xymatrix@C=2em{
&  X''       \ar[rrr]^{\sisi{}{s}}\ar[ldd]   &     & &**[r]f^{\ast}Y''\ar[rrr]\ar[ldd]   &  & & Y''\ar[ldd] \\                      &     & f^{\ast}Y   \sisi{\ar[rrr]|\hole}{\ar@{=}[rrr]|\hole}\ar[lu]& &      &f^{\ast}Y \ar[rrr]|\hole\sisi{}{\ar[ul]}&&& Y\ar[lu]\\
  f^{\ast}Y' \sisi{\ar[rrr]}{\ar@{=}[rrr]}\ar[rru]         &&& f^{\ast}Y'\sisi{}{\ar[rru]}   \ar[rrr] && & Y'\ar[rru]&&.
}
\]

Let us prove that $s$ is an isomorphism, which will complete the
proof. For every $X$ in $\Ob \Aa_x$, we have the
commutative diagram:
\[
\text{\scriptsize
$\xymatrix@C=1em{
\Hom_x(X,f^{\ast}Y') \ar[r] \ar@{=}[d]& \Hom_x(X,f^{\ast}Y) \ar@{=}[d]\ar[r]& \Hom_x(X,X'')\ar[d]
\ar[r]& \Hom_x(X,f^{\ast}Y'[1])\ar[r]\ar@{=}[d]
& \Hom_x(X,f^{\ast}Y[1]) \ar@{=}[d]\\
\Hom_x(X,f^{\ast}Y') \ar[r]\ar@{=}[d]&  \Hom_x(X,f^{\ast}Y)\ar[r] \ar@{=}[d]&
\Hom_x(X,f^{\ast}Y'') \ar[r]\ar@{=}[d]&  \Hom_x(X,f^{\ast}Y'[1])\ar@{=}[d] \ar[r]&
\Hom_x(X,f^{\ast}Y[1]) \ar@{=}[d]\\
\Hom_f(X,Y)  \ar[r]& \Hom_f(X,Y)\ar[r]  &
\Hom_f(X,Y'') \ar[r]& \Hom_f(X,Y'[1])\ar[r] &
\Hom_x(X,Y[1]).
}$}
\]

The\pageoriginale first and last rows are exact by
\hyperref[XVII.prop2.2.4]{2.2.4}. By the five lemma, the map induced by $s$ from
$\Hom(X,X'')$ to $\Hom(X,f^{\ast}Y'')$ is always bijective, and $s$
is therefore an isomorphism.
\end{itemize}

There are, of course, dual results for cofibered categories.

\subsection{Trivial duality formula}\etiquette[2.3]{XVII.subsec2.3}

Let $\Aa$ be a category over a base category~$\Bb$,
and let $S$ be a set of arrows of $\Aa$, each of which
projects onto an identity of $\Bb$. The category
$\Aa(S^{-1})$ obtained from $\Aa$ by the calculus of fractions
is a category over $\Bb$ (by the universal property).

Let $\Bb'\rightarrow \Bb$ be a functor; denote by
$\Aa'$ the fiber product
$\fprod{\Aa}{\Bb'}{\Bb}$ and by $S'$ the inverse
image of $S$ in $\Aa'$. The composite functor
$\Aa'\rightarrow \Aa\rightarrow \Aa(S^{-1})$
makes the arrows of $S'$ invertible and hence factors through
$\Aa'({S'}^{-1})$. We thus obtain a canonical $\Bb'$-functor
as follows:
$$\etiquette[2.3.1]{XVII.eq2.3.1}
\Aa'(S^{-1})\longrightarrow
\fprod{\Bb'}{\Aa(\sisi{S}{{S'}}^{-1})}{\Bb}{\quoi}.\leqno(2.3.1)
$$

\begin{enonce*}{Lemma 2.3.2}\etiquette[2.3.2]{XVII.lem2.3.2}
If $\Bb'$ is a subcategory of $\Bb$, and if
whenever a composite arrow $h=f\circ g$ of $\Bb$
belongs to $\Bb'$, both $f$ and $g$ belong to
$\Bb'$, then the functor \eqhyperref[XVII.eq2.3.1]{2.3.1} is an isomorphism.
\end{enonce*}

Note\pageoriginale that \eqhyperref[XVII.eq2.3.1]{2.3.1} always induces a bijection on
objects, and that, in the present case,
$\fprod{\Bb'}{\Aa(S^{-1})}{\Bb}$ is a
subcategory of $\Aa(S^{-1})$.

Recall the “explicit” construction of $\Aa(S^{-1})$
\cite[\sisi{Chap.}{\chap}~I]{sga4.exp.XVII.bib.GZ}. Given two objects $X$ and $Y$ of $\Aa$, one
considers the “formal composites” $f_1\ldots f_n$, where each
$f_i$ is either an arrow of $\Aa$ or the inverse of an
arrow of $S$. The set $\Hom_{\Aa(S^{-1})}(X,Y)$ is the
quotient of the set of these “formal composites” by the
equivalence relation generated by the relations
\begin{alignat*}{2}
& f_1\ldots uv\ldots f_n\equiv f_1\ldots (u\circ v)\ldots f_n &\quad &
\mbox{\rm for } u,v\;\mbox{\rm in } \Fl(\Aa)\\
& f_1\ldots ss^{-1}\ldots f_n\equiv f_1\ldots f_n && \mbox{\rm for }
s\in S\\
& f_1\ldots s^{-1}s\ldots f_n\equiv f_1\ldots f_n && \mbox{\rm for }
s\in S{\quoi}.
\end{alignat*}

% END INCLUDED SGA4_English_current_buildable_source_20260906/source/17/drafts/SGA4_XVII_LINES_1401_1800_English.texfrag

% BEGIN INCLUDED SGA4_English_current_buildable_source_20260906/source/17/drafts/SGA4_XVII_LINES_1801_2200_English.texfrag

If these constructions are applied to $\Aa$ and to the
subcategory $\Aa'$, one checks that $\Aa'(S^{-1})$
is identical to the inverse-image subcategory of $\Aa(S^{-1})$
of $\Bb'$, since if a ``formal composite'' $f_1\ldots f_n$ has as
its composite in $\Bb$ an arrow of $\Bb'$, the image of each $f_i$
will also be an arrow of $\Bb'$.

Suppose now that $\Aa$ is a triangulated category over $\Bb$, and
that the sets $S_x=S\cap \Fl(\Aa_x)$, for $x\in \Ob(\Bb)$, are
saturated multiplicative systems (\cite{sga4.exp.XVII.bib.CD} \sisi{Ch.}{\chap}~I
\S~2\sisi{. 1}{ nᵒ 1}).

\begin{enonce*}{Definition 2.3.3}\etiquette[2.3.3]{XVII.def2.3.3}
Under the preceding hypotheses, the category $\Aa/\Bb$ is said to be
almost derivable (\resp derivable, \resp coderivable) relative to $S$
if, for every functor $F:\Bb'\rightarrow\Bb$, the functor
\eqhyperref[XVII.eq2.3.1]{2.3.1} is an isomorphism (\resp and if
$\Aa(S^{-1})$ is fibered, \resp cofibered over $\Bb$).
\end{enonce*}

We\pageoriginale denote by $I$ the category of the following diagram:
$$\etiquette[2.3.4]{XVII.eq2.3.4}
I:\underset{0}{\times} \xrightarrow{F} \underset{1}{\times}{\quoi}.\leqno{(2.3.4)}
$$

By \hyperref[XVII.prop2.2.13]{2.2.13}, to give a triangulated category
$\Aa$ over $I$, fibered (\resp cofibered) over $\sisi{\Aa}{I}$, is
equivalently to give its fibers $\Aa_0$ and $\Aa_1$ and the
triangulated inverse-image functor
$F^{\ast}:\Aa_1\rightarrow\Aa_0$ (\resp and the triangulated
direct-image functor $F_{\ast}:\Aa_0\rightarrow\Aa_1$). If
$A_0\in\Ob\Aa_0$ and $A_1\in\Ob\Aa_1$, then by
\eqhyperref[XVII.eq2.2.12.1]{2.2.12.1},
\[
\Hom_F(A_0,A_1)=\Hom(A_0,F^{\ast}A_1)
\]
(\resp $\Hom_F(A_0,A_1)=\Hom(F_{\ast}A_0,A_1)$).

\begin{enonce*}{Proposition 2.3.5}\etiquette[2.3.5]{XVII.prop2.3.5}
Let $\Aa$ be a triangulated category over $I$ and let $S$ be as in
(\hyperref[XVII.def2.3.3]{2.3.3}).
\begin{enumerate}
\renewcommand{\theenumi}{\roman{enumi}}
\renewcommand{\labelenumi}{(\theenumi)}
\item $\Aa$ is almost derivable relative to $S$, and
$\Aa(S^{-1})$ is triangulated over $I$.

\item If $\Aa$ is fibered (\resp cofibered), the functor $F^{\ast}$
(\resp $F_{\ast}$) is right derivable (\resp left derivable) relative
to $S_0$ and $S_1$ (\hyperref[XVII.def1.2.1]{1.2.1}) if and only if
$\Aa$ is derivable (\resp coderivable) over $I$.
\end{enumerate}
\end{enonce*}

By Lemma \hyperref[XVII.lem2.3.2]{2.3.2}, the fiber categories of
$\Aa(S^{-1})$ are $\Aa_0(S^{-1}_0)$ and $\Aa_1(S^{-1}_1)$, where
$S_i=S\cap\Fl(\Aa_i)$ $(i=0,1)$. With the notation of \S
\hyperref[XVII.sec1]{1} ${\rm n^{\circ}}$ \hyperref[XVII.sec2]{2}, the
functors $A_0\rightarrow A^{-}_0$ and $A_1\rightarrow A^{+}_1$
identify $\Aa_0(S^{-1}_0)$ and $\Aa_1(S^{-1}_1)$ with categories of
\ppro- and \iind-objects of $\Aa_0$ and $\Aa_1$. This makes it possible
to define a category over $I$, with fibers $\Aa_0(S^{-1}_0)$ and
$\Aa_1(S^{-1}_1)$, by setting, for $A_0\in\Aa_0$ and $A_1\in\Aa_1$,
\[
\Hom_F(A_0,A_1)=\Hom_F(A^{-}_0,A^{+}_1)=\varinjlim_{A'_0{\xrightarrow[s]{}}A_{\sisi{p}{0}}}
\varinjlim_{A_1{\xrightarrow[t]{}}{A'_1}}\Hom(A'_0,A'_1){\quoi},
\]
where\pageoriginale $s$ and $t$ range over the elements of $S_0$ and
$S_1$ with target (respectively source) $A_0$ (respectively $A_1$).
The category $\Aa$ maps to this category, and one checks immediately
that it satisfies the universal property of $\Aa(S^{-1})$, over $I$
and after every base change.

To verify (i), it suffices to prove that this category, namely
$\Aa(S^{-1})$, is triangulated, and it suffices to verify condition
\hyperref[XVII.def2.2.3]{2.2.3}~(iii). Let $(X_0,Y_0,Z_0)$ and
$(X_1,Y_1,Z_1)$ therefore be distinguished triangles of $\Aa_0$ and
$\Aa_1$, giving rise to the distinguished triangles
$(X^{-}_0,Y^{-}_0,Z^{-}_0)$ and $(X^{+}_1,Y^{+}_1,Z^{+}_1)$ of
$\Aa_0(S^{-1}_0)$ and $\Aa_1(S^{-1}_1)$.
$$\etiquette[2.3.5.1]{XVII.eq2.3.5.1}
\vcenter{
\xymatrix{
X^{-}_{0}  \ar[r]^{u} \ar[d]& X^{+}_1  \ar[d]\\
Y^{-}_{0}  \ar[r]^{v} & Y^{+}_1\rlap{.}\\
}}\leqno{(2.3.5.1)}
$$

By Lemma \hyperref[XVII.lem1.2.2.1]{1.2.2.1}, the triangles
$(X^{-}_0,Y^{-}_0,Z^{-}_0)$ and $(X^{+}_1,Y^{+}_1,Z^{+}_1)$ are limits
of distinguished triangles; there therefore exist morphisms of
distinguished triangles
\begin{equation*}
\begin{split}
(X^{-}_0,Y^{-}_0,Z^{-}_0) &\longrightarrow (X'_0, Y'_0,Z'_0)\\
(X'_1, Y'_1, Z'_1) &\longrightarrow (X^{+}_1, Y^{+}_1, Z^{+}_1){\quoi},
\end{split}\tag{2.3.5.2}
\end{equation*}\etiquette[2.3.5.2]{XVII.eq2.3.5.2}%Removing this comment would introduce unwanted white space.
and\pageoriginale a diagram
\[
\xymatrix{
X'_{0}  \ar[r]^{u'} \ar[d]& X'_1 \ar[d] \\
Y'_{0}  \ar[r]^{v'} & Y'_1\rlap{,}
}
\]
compatible with \eqhyperref[XVII.eq2.3.5.1]{2.3.5.1}. Since $\Aa$ is
triangulated over $I$, this diagram can be extended to a morphism of
triangles
\[
(X'_0,Y'_0,Z'_0)\longrightarrow (X'_1, Y'_1, Z'_1){\quoi}.
\]
Composing the latter with the morphisms
\eqhyperref[XVII.eq2.3.5.2]{2.3.5.2}, we see that diagram
\eqhyperref[XVII.eq2.3.5.1]{2.3.5.1} can be extended to a morphism of
triangles, which proves \hyperref[XVII.def2.2.3]{2.2.3}~(iii).

Let us verify (ii), assuming for example that $\Aa$ is fibered over
$I$. If $A_1\in\Ob(\Aa_1(S^{-1}_1))$, an inverse image of $A_1$ in
$\Aa(S^{-1})$ is an object $B$ of $\Aa_0(S^{-1}_0)$ such that, in
$\Aa_0(S^{-1}_0)$,
\begin{align*}
\Hom_{\Aa_0(S^{-1}_0)}(A_0,B) & =\Hom_F(A^{-}_0,
A^{+}_1)=\Hom(A^{-}_0,
F^{\ast}(A^{+}_1))\\
&=\Hom_{\Aa_0(S^{-1}_0)}
(A_0,F^{\ast}(A^{+}_1)){\quoi}.
\end{align*}
Such an object $B$ exists if and only if the \iInd-object
$F^{\ast}(A^{+}_1)$ of $\Aa_0(S^{-1}_0)$ is essentially constant
(\ie representable), and $B$ then represents this \iind-object, which
is precisely $RF(A_1)$ (\hyperref[XVII.def1.2.1]{1.2.1}); this proves
the assertion.

\begin{enonce*}{Corollary 2.3.6}\etiquette[2.3.6]{XVII.cor2.3.6}
With the notation of \hyperref[XVII.def2.3.3]{2.3.3}, if $\Aa$ is
almost derivable over $\Bb$, then $\Aa(S^{-1})$ is a triangulated
category over $\Bb$.
\end{enonce*}

First, the fiber category $\Aa(S^{-1})_x$ can be identified with
$\Aa_x(S^{-1}_x)$: apply the hypothesis to the base change
$\{x\}\rightarrow\Bb$. This\pageoriginale category is triangulated,
equipped with a set $\Delta_x$ of triangles, which gives a meaning to
the preceding statement. To verify condition \hyperref[XVII.def2.2.3]{2.2.3}~(iii), it suffices
to do so after every base change $I\rightarrow\Bb$, and the assertion
follows from \hyperref[XVII.prop2.3.5]{2.3.5}~(i).

It follows immediately from the proof of
\hyperref[XVII.prop2.3.5]{2.3.5}~(ii) that:

\begin{enonce*}{Theorem 2.3.7}\etiquette[2.3.7]{XVII.thm2.3.7}
\emph{(trivial duality formula).}

Let $\Aa$ be a triangulated category fibered and cofibered over $I$,
and denote by $F_{\ast}$ and $F^{\ast}$ the direct- and inverse-image
functors. Let $S=S_0\cup S_1$ be a multiplicative system as in
\hyperref[XVII.def2.3.3]{2.3.3}. Suppose that $LF_{\ast}$ is defined
at $A_0\in\Ob\Aa(S^{-1}_0)$ and that $RF^{\ast}$ is defined at
$A_1\in\Ob\Aa_1(S^{-1}_1)$. Then
\[
\Hom_{\Aa_0(S^{-1}_0)}(A_0,RF^{\ast}A_1)=\Hom_{\Aa\sisi{_{(S^{-1})}}{(S^{-1})}}(A_0,A_1)=\Hom_{\Aa_1(S^{-1}_1)}(LF_{\ast}A_0,A_1){\quoi}.
\]
\end{enonce*}

The positions of the symbols $R$, $L$, ${}^{\ast}$, and ${}_{\ast}$
in \hyperref[XVII.thm2.3.7]{2.3.7} are only apparently anomalous: in
fact, in the category \hyperlink{sga4.exp.XVII.subsection.0.2}{0.2} of sheaves on variable sites, fibered and
cofibered, the fiber categories are the opposites of the usual
categories of sheaves.

\subsection{Categories fibered in derived categories}\etiquette[2.4]{XVII.subsec2.4}

\setcounter{subsubsection}{-1}
\subsubsection{}\etiquette[2.4.0]{XVII.subsubsec2.4.0}
In this ${\rm n^{\circ}}$, $\Aa$ denotes a triangulated category
cofibered over $\Bb$, and $S$ a multiplicative system of morphisms of
$\Aa$, all lying over an identity, such that
$S_x=S\cap\Fl(\Aa_x)$ is a saturated multiplicative system of $\Aa_x$
\cite[\sisi{}{I }1.2.1]{sga4.exp.XVII.bib.CD} for every $x\in\Ob(\Bb)$.

We\pageoriginale seek a criterion for $\Aa$ to be coderivable relative
to $S$ (\hyperref[XVII.def2.3.3]{2.3.3}); this implies that
$\Aa(S^{-1})$ is triangulated over $\Bb$ when it is equipped with the
set of triangles isomorphic to the image of a distinguished triangle
of $\Aa$ (\hyperref[XVII.prop2.3.5]{2.3.5}).

\begin{enonce*}{Proposition 2.4.1}\etiquette[2.4.1]{XVII.prop2.4.1}
Let $\Aa$, $\Bb$, and $S$ be as above. For $\Aa$ to be coderivable
relative to $S$, it suffices that the following condition hold:
\begin{itemize}
\item[$(\ast)$] $\forall x\in\Ob\Bb$, $\forall A\in\Ob(\Aa_x)$,
$\exists s:A'\rightarrow A$ in $S$\sisi{}{,} $\forall g$,
$f:x\xrightarrow{f}y\xrightarrow{g}\sisi{Z}{z}$, $f_{\ast}A'$ is
left adapted to the functor $g_{\ast}$.
\end{itemize}
\end{enonce*}

Let $\Aa^{1}$ denote the full subcategory of $\Aa$ consisting of the
objects all of whose direct images are adapted to all direct-image
functors. Then \hyperref[XVII.prop1.2.2]{1.2.2} shows that $\Aa^{1}$ is
a triangulated subcategory of $\Aa$, evidently cofibered. Set
$S^{1}=S\cap\text{Fl}(\Aa^{1})$.

\begin{enonce*}{Lemma 2.4.1.1}\etiquette[2.4.1.1]{XVII.lem2.4.1.1}
Suppose that $S$ is stable under direct images. Then $\Aa/\Bb$ is
coderivable and $\Aa(S^{-1})$ can be computed by a calculus of right
fractions.
\end{enonce*}

Let us verify condition c) of \cite[I 2.2]{sga4.exp.XVII.bib.GZ}, possibly the least
trivial of the conditions. We are given $f$ and $s$ $(s\in S)$ and
must find
$$\etiquette[2.4.1.2]{XVII.eq2.4.1.2}
\vcenter{
\xymatrix{
A'        \ar@{-->}[r]      & B'          \\
A  \ar[u]^{s}   \sisi{\ar[r]}{\ar[r]_f}          & B\rlap{,}  \ar@{-->}[u]^{t}
}}\leqno{(2.4.1.2)}
$$
a\pageoriginale commutative diagram \eqhyperref[XVII.eq2.4.1.2]{2.4.1.2}
with $t\in S$. Let $F$ be the image of $f$ in $\Bb$. By hypothesis,
$F_{\ast}(s)$ still belongs to $S$, so by \textsc{Verdier} there is a
commutative diagram
\[
\xymatrix{
F_{\ast}(A')  \ar[r] & B'\\
F_{\ast}(A)    \ar[r]_{f}\sisi{\ar[u]_{s}}{\ar[u]^{F_\ast(s)}}   & B\ar[u]_{t}
}
\]
and hence also a diagram \eqhyperref[XVII.eq2.4.1.2]{2.4.1.2}.

Since $\Aa(S^{-1})$ is computed by a calculus of right fractions, it
is immediate that its formation commutes with base change
$c:\Bb'\rightarrow\Bb$.

This lemma applies in particular to $\Aa^{1}$ and $S^{1}$; the
following lemma therefore implies \hyperref[XVII.prop2.4.1]{2.4.1}:

\begin{enonce*}{Lemma 2.4.1.3}\etiquette[2.4.1.3]{XVII.lem2.4.1.3}
Over $\Bb$, and after every base change $c:\Bb'\rightarrow\Bb$, the
canonical functor
$\Phi:\Aa\sisi{_1}{^1}(\sisi{{S^1}}{(S^1)}^{-1})\rightarrow\Aa(S^{-1})$
is an equivalence of categories.
\end{enonce*}

Let us construct a quasi-inverse to $\Phi$. For each $A\in\Ob\Aa$,
choose $s_A\in S$ with target $A$ and source $A^1\in\Ob\Aa^{1}$,
with $s_A=\Id_A$ when $A\in\Ob\Aa^{1}$. Let $f$ be an arrow of $\Aa$
with target $A$, source $B$, and image $F$ in $\Bb$. Since
$s_B\in S$ and $F_{\ast}A^{1}\in\Ob\Aa^{1}$, one can find
$B\sisi{^{1}}{'}$, \sisi{$v$ and $s\in S^{1}$}{$s\in S^1$ and $u$}
making the following diagram commutative:
$$\etiquette[2.4.1.4]{XVII.eq2.4.1.4}
\vcenter{
\xymatrix{
A      \ar[r]^{f} &  B      \ar[r]_{s_B} & B\sisi{_1}{^1}       \\
A\sisi{_1}{^1} \ar[u]_{s_A}\ar[r]^{a} & F_{\ast}A\sisi{_1}{^1}  \ar[u]& B'\rlap{.} \ar[l]_{s} \ar[u]_u
}}\leqno{(2.4.1.4)}
$$

The\pageoriginale arrow $f^{1}=us^{-1}a$ of
$\Aa^{1}(\sisi{{S^1}}{(S^1)}^{-1})$ then depends only on $f$.

Let $f$ and $g$ be two composable arrows, with images $F$ and $G$ in
$\Bb$. To prove that $(g\circ f)^{1}=g^{1}\circ f^{1}$, consider the
following diagram, which contains diagrams
\eqhyperref[XVII.eq2.4.1.4]{2.4.1.4} for $f$ and $g$:
\[
\xymatrix@C=2em{
                      & A \ar[rr]^{f}&      & B\ar[rrr]^{g}     & &
                      & C  &       \\
 A^{1} \ar[ru] \ar[rr]^{a}&          &F_{\ast}A^{1}\ar[ru] \ar[rrr]^(.6){a'}& & & (GF)_{\ast}A^1\ar[ru] &
              & &C^{1}\ar[llu]\\
& & & B^1 \ar[rrr]^(.4){b}|(.58)\hole\ar[uu]|\hole & & & G_{\ast}B^1\ar[uu] &     \\
&  & B'\ar[uu]^{s}\ar[rrr]^{b'}\ar[ru]_{u}       &   & &G_{\ast}B'\ar[ru]\ar[uu]^(.3){G_{\ast}s}    &  &C''\ar[ll]^{r}\ar[r]_{w} & C'\ar[llu]_t\ar[uu]_{v}\rlap{.}
}
\]

Since $t\in S$, there exist $C''$, $r\in S^{1}$, and $w$ making the
diagram commutative. By definition,\break
$g^{1}\circ f^{1}=vt^{-1}b\,u\,s^{-1}a$ and
$(g\circ f)^{1}=vw(G_{\ast}(s)r)^{-1}a'a$, and the conclusion follows
from the commutativity of the diagram.

We have thus defined a functor from $\Aa$ to
$\Aa^{1}(\sisi{{S^1}}{(S^1)}^{-1})$; by the universal property of
$\Aa(S^{-1})$, it factors through
$\Psi:\Aa(S^{-1})\rightarrow\Aa^{1}(\sisi{{S^1}}{(S^1)}^{-1})$.
The composite $\Phi\Psi$ is isomorphic to the identity because its
composite with the functor $J:\Aa\rightarrow\Aa(S^{-1})$ is
isomorphic to $J$. One checks similarly that $\Psi\Phi$ is isomorphic
to the identity.

\begin{enonce*}{Corollary 2.4.2}\etiquette[2.4.2]{XVII.cor2.4.2}
With the notation of \hyperref[XVII.prop2.4.1]{2.4.1}, $\Aa$ is
coderivable relative to $S$ provided that, for every finite set $B$
of arrows of $\Bb$, condition $(\ast)$ holds when only direct images
under the arrows of $B$ are considered.
\end{enonce*}

For\pageoriginale every category $\Bb'$ with a finite set of arrows, Proposition
\hyperref[XVII.prop2.4.1]{2.4.1} is applicable after every base change
$\Bb'\rightarrow\Bb$. We shall take for $\Bb'$ the categories of the
following diagrams:
\[
  \begin{gathered}
    *\\
    (0)
  \end{gathered}\qquad
  \begin{gathered}
    *\to*\\
    (1)
  \end{gathered}\qquad
  \begin{gathered}
    *\to*\to*\\
    (2)
  \end{gathered}\qquad
  \begin{gathered}
    *\to*\to*\to*\\
    (3)
  \end{gathered}
\]
Consider the categories $\Aa_x(S^{-1}_x)$ $(x\in\Ob\Bb)$. Given
$f\in\Fl(\Bb)$ and $A_1,A_2$ in these categories over the source and
target of $f$, define $\Hom_f(A_1,A_2)$ to be the analogous set defined
after the base change of type (1) determined by $f$. Composition of
morphisms is defined using base changes of type (3). In this way one
constructs a category $\Aa(S^{-1})^{\ast}$ over $\Bb$, whose formation
is compatible with every base change; using base changes of types
(0), (1), and (2), one checks that it solves the universal problem
defining $\Aa(S^{-1})$.

The following result is dual to \hyperref[XVII.cor2.4.2]{2.4.2}.

\begin{enonce*}{Proposition 2.4.3}\etiquette[2.4.3]{XVII.prop2.4.3}
Let $\Aa$ be a triangulated category fibered over $\Bb$, and let $S$
be as in \hyperref[XVII.prop2.4.1]{2.4.1}. For $\Aa$ to be derivable
relative to $S$, it suffices that, for every finite set $B$ of arrows
of $\Bb$,
\begin{itemize}
\item[$(\ast')$] $\forall x\in\Ob(\Bb)$, $\forall A\in\Ob(\Aa_x)$,
$\exists s:A\rightarrow A'$ in $S$,
$\forall g,f:Z\xrightarrow{g}Y\xrightarrow{f}X$ in $B$, $f^{\ast}A'$
is right adapted to the functor $g^{\ast}$.
\end{itemize}
\end{enonce*}

\section{Gluing fibered or cofibered categories}\etiquette[3]{XVII.sec3}\pageoriginale

\subsection{Introduction}\etiquette[3.1]{XVII.subsec3.1}

Let $f:Y\rightarrow X$ be a separated morphism of finite type from a
scheme $Y$ to a scheme $X$, or a continuous map from a locally compact
space to a locally compact space. One can then define a ``direct image
with proper supports'' functor, denoted $f_{!}$, from the category of
abelian sheaves on $Y$ to the category of abelian sheaves on $X$ (\cf
\sisi{\hyperref[XVII.subsec5.3]{5.3}}{\hyperref[XVII.subsec6.1]{6.1}}
in the case of schemes, and \textsc{Verdier} \cite{sga4.exp.XVII.bib.V3} for separated
locally compact spaces). When $Y$ and $X$ are schemes of finite type
over $\CC$, these two definitions are compatible in the evident sense.
Unfortunately, for schemes, unlike for locally compact spaces, the
derived functors of $f_{!}$ are pathological; one can nevertheless
define a reasonable ``direct image with proper supports'' functor
directly from $D^{+}(Y)$ to $D^{+}(X)$. To define it, factor $f$ as an
open immersion followed by a proper morphism: $f=\bar f j$. The desired
functor $Rf_{!}$ is then defined as the composite of ``extension by
zero'' $j_{!}$ and the functor $R\bar f_{\ast}$, derived from direct
image under $\bar f$.

These ``direct image with proper supports'' functors $Rf_{!}$ are
usable only after one verifies for them a ``variance formalism'' that
includes, for a composite $f=gh$, a formula
\[
Rf_{!}=Rg_{!}Rh_{!}
\]
which is best expressed in terms of cofibered categories. Such a
formalism is available separately for ``extensions by zero'' and for
\pageoriginale direct images under proper morphisms, and the problem is to glue these
formalisms. To avoid later being bogged down by pointless computations,
we propose in this \S to list the elementary verifications required to
carry out such a gluing. We do so in a slightly more abstract setting
than is strictly necessary, in the hope that the result may be useful
again.

The construction of $Rf_{!}$ sketched above requires that the morphism
$f$ be factorable as an open immersion followed by a proper morphism.
Following \textsc{Hartshorne} \cite[III 8 p.~189]{sga4.exp.XVII.bib.RD}, in nᵒ
\hyperref[XVII.sec2]{2} we introduce the appropriate notion of a
compactifiable morphism.

\subsection{Compactifiable morphisms}\etiquette[3.2]{XVII.subsec3.2}

% END INCLUDED SGA4_English_current_buildable_source_20260906/source/17/drafts/SGA4_XVII_LINES_1801_2200_English.texfrag

% BEGIN INCLUDED SGA4_English_current_buildable_source_20260906/source/17/drafts/SGA4_XVII_LINES_2201_2600_English.texfrag
\begin{enonce*}{Definition 3.2.1}\etiquette[3.2.1]{XVII.def3.2.1}
A morphism $f$ from an $S$-scheme $X$ to a quasi-compact,
quasi-separated $S$-scheme $Y$ is said to be \emph{$S$-compactifiable} if there exists an
$S$-scheme $P$ proper over $S$ and a factorization of $f$ as a
separated quasi-finite morphism $j$ from $X$ to $\fprod{P}{Y}{S}$ followed by
the projection from $\fprod{P}{Y}{S}$ to $Y$
\[
\xymatrix{
                                & &  &\fprod{P}{Y}{S} \ar[rd]\ar[ld]& \\
X \ar[rr]^(.7){f} \ar[rrru]^{j}\ar[rrrd]& &Y \ar[rd]&                & P\ar[ld]\\
                                & &  &       S        &.
}
\]
\end{enonce*}

When $S=Y$, we shall simply speak of a compactifiable morphism.

\begin{enonce*}{Proposition 3.2.2}\etiquette[3.2.2]{XVII.prop3.2.2}
\begin{enumerate}
\renewcommand{\theenumi}{\roman{enumi}}
\renewcommand{\labelenumi}{(\theenumi)}
\item An\pageoriginale $S$-compactifiable morphism is separated and of finite
type.

\item If $Y$ is quasi-compact and quasi-separated, a separated quasi-finite
morphism with target $Y$ is \sisi{}{$S$-}compactifiable.

\item The composite of two $S$-compactifiable morphisms is again
$S$-compactifiable.

\item Let there be a diagram of quasi-compact, quasi-separated $S$-schemes
\[
\xymatrix{
\fprod{X}{Y'}{Y} \ar[d]_{f'} \ar[r] & X \ar[d]^{f}      \\
Y' \ar[r]& Y\rlap{.}
}
\]
If $f$ is $S$-compactifiable, then $f'$ is $S$-compactifiable.

\item Let $g$ be an $S$-compactifiable morphism;\nde{it suffices to assume that $g$ is a quasi-compact, quasi-separated morphism of $S$-schemes.} for a composite morphism
$g\circ f$ to be $S$-compactifiable, it is necessary that $\sisi{h}{f}$ be so.
\end{enumerate}
\end{enonce*}

Assertions (i), (ii), and (iv) are immediate. To prove (iii) and
(v), consider a composite $f=gh$, with $g$ \sisi{}{$S$-}compactifiable.

Suppose that $gh$ is \sisi{}{$S$-}compactifiable, whence a diagram
\[
\xymatrix{
                     & \fprod{P}{Y}{S} \ar[d] \ar[r]^{v} & \fprod{P}{Z}{S}\ar[d] \\  X  \ar[ru]^{u} \ar[r]^{h}   & Y                \ar[r]^{g} &    Z\rlap{.}
}
\]
The composite $vu$ is quasi-finite; hence $u$ is quasi-finite, and (v) is
proved.

If $h$ is compactifiable, we have a diagram
\[
\xymatrix{
  &&                & &\sisi{Q\times P \times Z}{\fprod{\fprod{Q}{P}{S}}{Z}{S}} \ar[d]&\\
  &&\fprod{Q}{Y}{S} \ar[d]\ar[rru]\ar[rdd]& & \fprod{P}{Z}{S} \ar[d]\ar[rdd] & \\
X \ar[rrrdd]\ar[rru] \ar[rr]^{h}&& Y \sisi{\ar[rru]}{\ar[rru]|(.2)\hole} \ar[rdd]\sisi{\ar[rr]^{g}}{\ar[rr]^g|(.25)\hole}            & & Z\ar[ldd]                & \\
  &&                &Q\ar[d]&                  &P\ar[lld] \\
  &&                &S&                   &,
}
\]
and\pageoriginale $f=fh$ factors through a separated quasi-finite morphism
to $\fprod{(\fprod{Q}{P}{S})}{Z}{S}$.

Nagata asserts in \sisi{[5]}{\cite{sga4.exp.XVII.bib.Na}} that every morphism of finite type between separated
integral noetherian schemes is compactifiable.\nde{more generally, every separated morphism of finite type whose target is quasi-compact and quasi-separated is compactifiable. See \nderef{XVII.ndeConrad}.} The editor admits to
not having understood the proof.

Now fix a scheme $S$ and denote by $(S)$ the
subcategory of the category of $S$-schemes whose objects are
the quasi-compact, quasi-separated $S$-schemes and whose arrows are
the $S$-compactifiable $S$-morphisms.\nde{Thanks to the compactification theorem, one may also work in the subcategory of the category of (absolute) schemes whose objects are the quasi-compact, quasi-separated schemes and whose arrows are the separated morphisms of finite type (\ie compactifiable morphisms). Note that the stability of compactifiable morphisms under composition is equivalent to the general form of the compactification theorem.}

\begin{enonce*}{Proposition 3.2.3}\etiquette[3.2.3]{XVII.prop3.2.3}
In the category $(S)$:
\begin{enumerate}
\renewcommand{\theenumi}{\roman{enumi}}
\renewcommand{\labelenumi}{(\theenumi)}
\item The open immersions (\resp the proper morphisms)
define a subcategory $(S,i)$ (\resp $(S,p)$) of $(S)$ having
the same objects as $(S)$.

\item Fiber products exist in $(S,p)$ and are fiber
products in $(S)$.

\item Every morphism $f$ is the composite $f=pj$ of a proper morphism
$p$ and an open immersion.
\end{enumerate}
\end{enonce*}

By\pageoriginale \hyperref[XVII.prop3.2.2]{3.2.2} (iii), (iv), and (v), fiber products
exist in $(S)$ and are fiber products in the category of
all schemes. Since the fiber product $X=\fprod{X_1}{X_2}{Y}$ of
two schemes proper over $Y$ is again proper over $Y$, and a
scheme is proper over $X$ if and only if it is proper over $X_1$
and $X_2$, condition (ii) is satisfied.

If $f$ is a morphism of $(S)$, by hypothesis it admits, in $(S)$,
a factorization $f=pj'$, where $p$ is proper and $j\sisi{}{'}$ is separated and quasi-finite.
By \EGA IV 18.12.13, $j'$ admits a factorization
$j'=qj$, where $q$ is finite and $j$ is an open immersion; moreover, by
\hyperref[XVII.prop3.2.2]{3.2.2} (ii), $q$ and $j$ belong to $(S)$. We then have $f=(pq)j$,
which proves (iii). Assertion (i) follows from \hyperref[XVII.prop3.2.2]{3.2.2} (iii).

\setcounter{subsubsection}{3}
\subsubsection{}\etiquette[3.2.4]{XVII.subsubsec3.2.4}
In the rest of this \S, suppose given a category $(S)$, equipped
with two subcategories $(S,i)$ and $(S,p)$, such that
\begin{enumerate}
\renewcommand{\theenumi}{\roman{enumi}}
\renewcommand{\labelenumi}{(\theenumi)}
\item $\Ob(S)=\Ob(S,i)=\Ob(S,p)$;

\item fiber products exist in $(S,p)$ and are fiber
products in $(S)$;

\item $\forall f\in \Fl(S)$, $\exists p\in \Fl(S,p)$, $\exists
j\in \Fl(S,i)$, $f=pj$.
\end{enumerate}

The arrows of $(S,i)$ (\resp of $(S,p)$) will be called open
immersions (\resp proper morphisms).

\begin{enonce*}{Definition 3.2.5}\etiquette[3.2.5]{XVII.def3.2.5}
\begin{enumerate}
\renewcommand{\theenumi}{\roman{enumi}}
\renewcommand{\labelenumi}{(\theenumi)}
\item A \emph{compactification} of a morphism $f$ of $(S)$ (\resp
of a pair $(f,g)$, \resp of a triple $(f,g,h)$ of composable
morphisms) is a commutative diagram \eqhyperref[XVII.eq3.2.5.1]{3.2.5.1} (\resp
\eqhyperref[XVII.eq3.2.5.2]{3.2.5.2}, \resp \eqhyperref[XVII.eq3.2.5.3]{3.2.5.3}) whose horizontal arrows are
open immersions and whose oblique arrows are proper.

\item \emph{A\pageoriginale morphism of compactifications} is a
  morphism of
diagrams \eqhyperref[XVII.eq3.2.5.1]{3.2.5.1} (\resp\ldots) that is the identity on the 1\sisi{st}{\textsuperscript{st}}
vertical column and is proper at every vertex
$$
\etiquette[3.2.5.1]{XVII.eq3.2.5.1}
\etiquette[3.2.5.2]{XVII.eq3.2.5.2}
\etiquette[3.2.5.3]{XVII.eq3.2.5.3}
\begin{gathered}[b]
\xymatrix{
Y        \ar[d]_{f} \ar@{^{(}->}[r]^{i}                & \bar{Y} \ar[ld]^{\bar{f}}\\
X      & &}\\
\hyperref[XVII.eq3.2.5.1]{(3.2.5.1)}
\end{gathered}\qquad
\begin{gathered}[b]
\xymatrix{
Z        \ar@{^{(}->}[r]^{j} \sisi{\ar[d]}{\ar[d]_g}& \bar{Z}\ar@{^{(}->}[r]^{j'}\sisi{\ar@{-}[ld]^{\bar{g}}}{\ar[ld]^{\bar{g}}}
&\bar{\bar{Z}} \ar[ld]^{\bar{\bar{g}}}\\
Y \ar[d]_{f}\ar@{^{(}->}[r]^{i} & {\bar{Y}} \ar[ld]^{\bar{f}}&\\
X      & &}\\
\hyperref[XVII.eq3.2.5.2]{(3.2.5.2)}
\end{gathered}\qquad
\begin{gathered}[b]
\xymatrix@=3em{
\bigdot \ar@{^{(}->}[r]_(.4){\sisi{h}{k}} \ar[d]_{h}& \bigdot \ar@{^{(}->}[r]_(.4){\sisi{h}{k}'}\ar[ld]^(.4){\bar{h}} & \bigdot
\ar@{^{(}->}[r]_(.4){\sisi{h}{k}''} \ar[ld]^(.4){\bar{\bar{h}}}&\bigdot \ar[ld]^(.4){\bar{\bar{\bar{h}}}}\\
\bigdot \ar@{^{(}->}[r]_(.4){j}\ar[d]_{g} & \bigdot \ar@{^{(}->}[r]_(.4){j'}
\ar[ld]^(.4){\bar{g}}& \bigdot \ar[ld]^(.4){\bar{\bar{g}}}\\
\bigdot
\ar@{^{(}->}[r]_(.4){i} \ar[d]_{f}&\bigdot \ar[ld]^(.4){\bar{f}}\\
\bigdot}\\
\hyperref[XVII.eq3.2.5.3]{(3.2.5.3)}
\end{gathered}
$$
\end{enumerate}
\end{enonce*}

\begin{enonce*}{Proposition 3.2.6}\etiquette[3.2.6]{XVII.prop3.2.6}
\begin{enumerate}
\renewcommand{\theenumi}{\roman{enumi}}
\renewcommand{\labelenumi}{(\theenumi)}
\item Every morphism (\resp every pair $(f,g)$, \resp every triple
$(f,g,h)$ of composable morphisms) admits a compactification.

\item The category of compactifications of $f$ is left
filtering.\nde{``left filtering'' means ``cofiltered'' (\hyperrefext[I][I.def2.7]{2.7}). For axioms $L_1$ and $L_2$ in the proof of \hyperref[XVII.prop3.2.6]{3.2.6}, \cf \hyperrefext[V][V.par8.1.1.1]{8.1.1.1}.}

\item If diagrams $\sisi{(3.2.4.2)}{\eqhyperref[XVII.eq3.2.5.2]{3.2.5.2}}_i$ $(i=1,2)$ are
compactifications of $(f,g)$, there exists a compactification
$\sisi{(3.2.4.2)}{\eqhyperref[XVII.eq3.2.5.2]{3.2.5.2}}_3$ of $(f,g)$ and morphisms of compactifications from
$\sisi{(3.2.4.2)}{\eqhyperref[XVII.eq3.2.5.2]{3.2.5.2}}_3$ to the compactifications $\sisi{(3.2.4.2)}{\eqhyperref[XVII.eq3.2.5.2]{3.2.5.2}}_i$.
\end{enumerate}
\end{enonce*}

\begin{enumerate}
\renewcommand{\theenumi}{\roman{enumi}}
\renewcommand{\labelenumi}{(\theenumi)}
\item The case of a morphism is precisely \hyperref[XVII.subsubsec3.2.4]{3.2.4}
(iii). Let us treat the case of a triple; with the notation of \eqhyperref[XVII.eq3.2.5.3]{3.2.5.3},
it suffices to construct successively compactifications of $f$,
$g$, and $h$, of $i\bar{g}$ and $j\bar{h}$, and finally of
$j'\bar{\bar{h}}$.

\item Let\pageoriginale $\eqhyperref[XVII.eq3.2.5.1]{3.2.5.1}_1$ and $\eqhyperref[XVII.eq3.2.5.1]{3.2.5.1}_2$ be two
compactifications of $f$. The projection $\sisi{\rho}{p}$ from
$\fprod{\bar{Y}_1}{\bar{Y}_2}{S}$ to $X$ is proper, and $Y$ maps
to it by $(i_1,i_2)$, which (\hyperref[XVII.subsubsec3.2.4]{3.2.4} (iii)) can be compactified as
$(i_1,i_2)=qj$. The compactification $f=(p\sisi{\cdot}{\circ} q)j$ dominates
$\sisi{(3.2.4.1)}{\eqhyperref[XVII.eq3.2.5.1]{3.2.5.1}}_1$ and $\sisi{(3.2.4.1)}{\eqhyperref[XVII.eq3.2.5.1]{3.2.5.1}}_2$.

If $r$ and $s$ are two morphisms from $\sisi{(3.2.4.1)}{\eqhyperref[XVII.eq3.2.5.1]{3.2.5.1}}_1$ to
$\sisi{(3.2.4.1)}{\eqhyperref[XVII.eq3.2.5.1]{3.2.5.1}}_2$, the kernel $K$ of the pair of arrows
$\bar{Y}_{1} \overset{r}{\underset{s}\rightrightarrows} \bar{Y}_{2}$ maps by
a proper morphism to $\bar{Y}_1$, and $i_1$ factors through $K$:
write $i_1=ki$. If $i=pi_3$ is a compactification of $i$,
the existence of the compactification $f=(\bar{f}_1kp)i_3$ shows that
axiom $L_2$ is satisfied.

\item Apply axiom $L_1$ for filtering categories to the
compactifications of $f$ and $g$ deduced from the $\sisi{(3.2.4.2)}{\eqhyperref[XVII.eq3.2.5.2]{3.2.5.2}}_i$
$(i=1,2)$. We thereby obtain $\bar{Y}_3$ and $\bar{Z}_3$ giving rise to
commutative diagrams $\eqhyperref[XVII.eq3.2.5.1]{3.2.5.1}_i$ $(i=1,2)$
\[
\xymatrix{
&         &\bar{Z}_i \ar[d] \ar@{^{(}->}[r]& \bar{\bar{Z}}_i\ar[d]\\
\bar{Z}_3 \ar[rru] \ar[d]^{\bar{g}_3} &&Y   \ar@{^{(}->}[r]      & \bar{Y}_i\\
Y \ar@{=}[rru]\ar[r]_{i_3} & \bar{Y}_3 \ar[rru]&&\sisi{(3.2.5.1)_i}{},
}
\]
and it remains to find a compactification of $i_3\bar{g}_3$ by
$\bar{\bar{Z}}_3$, such that $\bar{\bar{Z}}_3$ can map by a
proper morphism to $\bar{\bar{Z}}_i$, so as to give rise to
a commutative cube $(i=1,2)$. If the target of a morphism between
compactifications of $i_3\bar{g}_3$ has this property, its source has it
as well, and by (ii), it will suffice to consider separately the cases
$i=1$ and $i=2$. Let\pageoriginale $T$ be the fiber product of $\bar{Y}_3$
and $\bar{\bar{Z}}_i$ over $\bar{Y}_i$: then $T$ occurs in a
commutative cube analogous to the desired one, except that the arrow $k$ from
$\bar{Z}_3$ to $T$ might not be an open
immersion. Writing $k=pj$ ($p$ proper, $j$ an open immersion) and repla\c{}cing $T$
by the source $\bar{\bar{Z}}_3$ of $p$, we obtain the desired
cube.
\end{enumerate}

\begin{enonce*}{Remark 3.2.7}\etiquette[3.2.7]{XVII.rem3.2.7}
Proposition \hyperref[XVII.prop3.2.3]{3.2.3} would still have been valid if, in Definition
\hyperref[XVII.def3.2.1]{3.2.1} of compactifiable morphisms, we had omitted the requirement that $Y$
be quasi-compact and quasi-separated and required instead that $j$
be a quasi-compact open immersion and not merely a
quasi-finite morphism.
\end{enonce*}

\subsection{Gluing}\etiquette[3.3]{XVII.subsec3.3}

Recall that the hypotheses of \hyperref[XVII.subsubsec3.2.4]{3.2.4} are satisfied.

\subsubsection{}\etiquette[3.3.1]{XVII.subsubsec3.3.1}
Suppose given:
\begin{itemize}
\item[$\alpha$)] for each $X\in \Ob (S)$, a category $F(X)$,

\item[$\beta$)] an $(S,i)$-cofibered category $F_i$ and an $(S,p)$-
cofibered category $F_{p}$,

\item[$\gamma$)] for each $X\in \Ob(S)$, isomorphisms between
$F(X)$, $(F_{i})_X$, and $(F_{p})_X$,

\item[$\delta$)] a normalized cleavage of $F_{i}$ and a normalized cleavage
of $F_{p}$.
\end{itemize}

By (\SGA 1 VI), it ``amounts to the same thing'' to give instead $\alpha$) and:
\begin{itemize}
\item[(i)] For every proper morphism $p:X\rightarrow Y$, a functor
$p_{\ast}:F(X)\rightarrow F(Y)$,

\item[(ii)] for every composable pair $(p,q)$ of proper morphisms,
an\pageoriginale isomorphism of functors
$c_{p,q}:p_{\ast}q_{\ast}\leftrightarrow (pq)_{\ast}$,

\item[(i\sisi{}{'}) and (ii\sisi{)'}{')},] analogous data for $(S,i)$,
\end{itemize}
these data satisfying the following conditions:
\begin{itemize}
\item[(a)] if $p$ is an identity, then $p_{\ast}$, $c_{p,q}$, and
$c_{q,p}$ are identities,

\item[(b)] for every composable triple of proper morphisms, we have
\[
c_{p,qr}\circ (p_{\ast}\ast c_{q,r})=c_{pq,r}\circ (c_{p,q}\ast
r_{\ast}){\quoi},
\]

\item[(a\sisi{)'}{')} and (b\sisi{)'}{')},] analogous conditions for $(S,p)$.
\end{itemize}

Suppose \emph{given} in addition,
\begin{itemize}
\item[(iii)] for every commutative diagram $D$
\[
\vcenter{
\xymatrix{
X \ar[r]^{j'} \ar[d]^{p'}& Y\ar[d]^{p} \\
Z \ar[r]_{j}             &           T
}}\quad j,j' \in \Fl(S,i)\text{; }
  p,p' \in \Fl(S,p),
\]
an isomorphism of functors $d(D):p_{\ast}j'_{\ast}\leftrightarrow
j_{\ast}p'_{\ast}$.
\end{itemize}

\begin{enonce*}{Proposition 3.3.2}\etiquette[3.3.2]{XVII.prop3.3.2}
If data {\rm (i) (ii) ${\rm (i')}$ ${\rm (ii')}$ (iii)}
satisfy conditions {\rm (a) (b) ${\rm (a')}$ ${\rm (b')}$} as well
as {\rm(c)} and ${\rm (c')}$ stated below, then there exists one
and ``essentially'' only one category $F$ cofibered over $S$,
equipped with isomorphisms $\fprod{F}{(S,i)}{S}=F_i$ and
$\fprod{F}{(S,p)}{S}=F_p$, compatible with datum \hyperref[XVII.subsubsec3.3.1]{3.3.1} $\gamma$)
and such that the isomorphisms $d(D)$ of {\rm(iii)} are the
composite isomorphisms
$j_{\ast}p'_{\ast}=(jp')_{\ast}=(pj')_{\ast}=p_{\ast}j'_{\ast}$.
\end{enonce*}

\begin{itemize}
\item[(c)] for every commutative diagram of the following type
\[
\xymatrix{
X_1 \ar@{^{(}->}[r]^{i} \ar[d]^{p_1}& X_2 \ar@{^{(}->}[r]^{i'}
\ar[d]^{p_2}& X_3 \ar[d]^{p_3}\\
Y_1 \ar@{^{(}->}[r]^{j}& Y_2 \ar@{^{(}->}[r]^{j'}& Y_3,\\
}
\]
whose\pageoriginale horizontal arrows are open
immersions and whose vertical arrows are proper, the following
triangle of isomorphisms $d$ from (iii) is commutative:
\[
\xymatrix@R=1cm{
p_{3{\ast}} i'_{\ast} i_{\ast} \ar[rr] \ar@{-}[rd]&  & j'_{\ast}j_{\ast}p_{1{\ast}}\ar@{-}[ld]\\
                               &j'_{\ast}p_{2{\ast}}
                                i_{\ast}&.
}
\]

\item[(c')] For every commutative diagram of the following type
\[
\xymatrix{
X_1  \ar[r]^{p} \ar[d]^{i_1}& X_2  \ar[r]^{p'} \ar[d]^{i_2}& X_3 \ar[d]^{i_3}          \\
Y_1         \ar[r]^{q} & Y_2  \ar[r]^{q'} & Y_3\rlap{,}
}
\]
whose horizontal arrows are proper and whose vertical
arrows are open immersions, the following triangle
of isomorphisms $d$ from (iii) is commutative:
\[
\xymatrix{
i_{3{\ast}} p'_{\ast} p_{\ast} \ar[rr] \ar@{-}[rd]&    &q'_{\ast} q_{\ast}
i_{1{\ast}} \\
&q'_{\ast} i_{2{\ast}} p_{\ast}  \ar@{-}[ru]&.
}
\]
\end{itemize}

\setcounter{subsubsection}{2}
\paragraph{Definition of $f_{\ast}$}\etiquette[3.3.2.1]{XVII.par3.3.2.1}\pageoriginale
Let $f$ be a morphism from $Y$ to $X$. Each compactification of $f$,
with diagram \sisi{(\hyperref[XVII.def3.2.5]{3.2.5} (i))}{\eqhyperref[XVII.eq3.2.5.1]{3.2.5.1}}, defines a composite functor $\bar{f}_{\ast}i_{\ast}$ from $F(Y)$ to $F(X)$. A morphism of compactifications
defines a commutative diagram
$$\etiquette[3.3.2.2]{XVII.eq3.3.2.2}
\vcenter{
\xymatrix{
Y  \ar@{^{(}->}[r]^{i_1} \ar@{=}[d]& \bar{Y}_1  \sisi{\ar[d]}{\ar[d]^p}\\
Y  \ar@{^{(}->}[r]^{i_2} \ar[d]        & \bar{Y}_2 \ar[ld]^{\bar{f}_2}\\
X                         &.
}}\leqno{(3.3.2.2)}
$$
Datum (iii) therefore provides us with an isomorphism
\[
\bar{f}_{2{\ast}} i_{2{\ast}}=\bar{f}_{2\ast} \circ
(i_{2\ast}\circ 1_{Y{\ast}})=\bar{f}_{2{\ast}}\circ
(p_{\ast}\circ i_{1{\ast}})=\bar{f}_{1{\ast}}\circ
i_{1{\ast}}{\quoi}.
\]
Axiom ${\rm(c')}$ also ensures that the isomorphism associated with a
composite of morphisms of compactifications is the composite of the
isomorphisms associated with each of them. Since the category of morphisms of
compactifications is cofiltered (\hyperref[XVII.prop3.2.6]{3.2.6} (ii)), we thereby obtain a
transitive system of isomorphisms between the functors associated with the
various compactifications of $f$. Choosing one of them, we
define the functor $f_{\ast}$.

\setcounter{paragraph}{2}
\paragraph{Transitivity homomorphisms}\etiquette[3.3.2.3]{XVII.par3.3.2.3}

Every compactification of a pair $(f,g)$ of composable morphisms
(diagram \eqhyperref[XVII.eq3.2.5.2]{3.2.5.2}) defines a compactification of $f$, one of $g$,
one of $f\circ g=(\bar{f}\bar{\bar{g}})(j'j)$, and an isomorphism
$c_{f,g}:(fg)_{\ast}\isomto f_{\ast}g_{\ast}$, the composite of the
isomorphisms
\[
(fg)_{\ast}=(\bar{g}\,\bar{\bar{g}})_{\ast}(j'j)_{\ast}=\bar{f}_{\ast}(\bar{\bar{g}}_{\ast}j'_{\ast})j_{\ast}=\bar{f}_{\ast}(i_{\ast}\bar{g}_{\ast})j_{\ast}=(\bar{f}_{\ast}i_{\ast})(\bar{g}_{\ast}j_{\ast})=f_{\ast}g_{\ast}
\]
(where\pageoriginale the middle isomorphism comes from (iii)).

% END INCLUDED SGA4_English_current_buildable_source_20260906/source/17/drafts/SGA4_XVII_LINES_2201_2600_English.texfrag

% BEGIN INCLUDED SGA4_English_current_buildable_source_20260906/source/17/drafts/SGA4_XVII_LINES_2601_3000_English.texfrag

Let us prove that this isomorphism does not depend on the compactification
chosen for $(f,g)$. By \hyperref[XVII.prop3.2.6]{3.2.6} (iii), it will suffice to compare
the isomorphisms $c_{f,g}$ obtained from two compactifications
$\sisi{(3.2.4.2)}{\eqhyperref[XVII.eq3.2.5.2]{3.2.5.2}}_i$ $(i=1,2)$ such that there exists a morphism of
compactifications from the first to the second. The arrows in these
morphisms will be denoted by $p$.

Consider the following diagram of isomorphisms of functors. For
readability, the ${\ast}$ have been omitted.
\[
\xymatrix{
\bar{f}_2 \bar{\bar{g}}_2j'_2j_2\sisi{}{\ar@{}[rrd]|{(1)}} \ar@{-}[rr] \ar@{-}[ddd] \ar@{-}[rd]& &
\bar{f}_2\bar{\bar{g}}_2pj'_1j_1
\ar@{-}[rr] \ar@{-}[ld] \ar@{-}[rd]& &\bar{f}_1\bar{\bar{g}}_1j'_1j_1 \ar@{-}[ddd]\sisi{}{\ar@{}[lld]|{(2)}}\\
 \ar@{}[ddr]|{(3)} & \bar{f}_2\bar{\bar{g}}_2j'_2pj_1\ar@{-}[ru] \ar@{-}[dd]& &
\bar{f}_2p\bar{\bar{g}}_1j'_1j_1 \ar@{-}[dd]\ar@{-}[ru]& \ar@{}[ldd]|{(5)}\\
 &  &  (4)  & &\\
\bar{f}_2i_2\bar{g}_2j_2 \ar@{-}[r] & \bar{f}_2i_2\bar{g}_2pj_1 & &
\bar{f}_2pi_1\bar{g}_1j_1 \ar@{-}[r]&  \bar{f}_1i_1\bar{g}_1j_1\\
& &  \bar{f}_2i_21_Y\bar{g}_1j_1 \ar@{-}[ru]\ar@{-}[lu]& & .\\
}
\]

\begin{enumerate}
\renewcommand{\theenumi}{\alph{enumi}}
\renewcommand{\labelenumi}{\theenumi)}
\item The composite isomorphism\pageoriginale along the 1\sisi{st}{\textsuperscript{st}} row is the isomorphism
\hyperref[XVII.par3.3.2.1]{3.3.2.1} between two definitions of $(fg)_{\ast}$.

\item The composite isomorphism along the last row is the isomorphism
\hyperref[XVII.par3.3.2.1]{3.3.2.1} between two definitions of $f_{\ast}g_{\ast}$.

\item The two outer vertical isomorphisms are the isomorphisms
$c_{f,g}$ between $(fg)_{\ast}$ and $f_{\ast}g_{\ast}$ deduced from one
or the other compactification.
\end{enumerate}
The problem is therefore to prove that the
outer boundary of the diagram is commutative. A commutative diagram of
functors remains commutative when each of them is composed with the same
functor. Bearing this in mind, the commutativity of (\sisi{i}{1}) follows from (c)
applied to the diagram
\[
\xymatrix{
Z  \ar@{^{(}->}[r]^{j_1} \ar@{=}[d]& \bar{Z}_1  \ar@{^{(}->}[r]^{j'_1} \ar[d]& \bar{\bar{Z}}_1 \ar[d]\\
Z  \ar@{^{(}->}[r]^{j_2} & \bar{Z}_2  \ar@{^{(}->}[r]^{j'_2} & \bar{\bar{Z}}_2\rlap{.}
}
\]
The commutativity of (3) and (5) is trivial; that of (2), which
essentially concerns proper morphisms, is also trivial. It remains to
consider the hexagon (4).

Consider the following commutative cube
\[
\xymatrix{
                              &**[r]\bar{Z}_2  \ar@{^{(}->}[rrr]^{ j'_2}
                              \ar[dd]^(.35){\bar{g}_2}|\hole &    &   &
                              **[r]\bar{\bar{Z}}_2 \ar[dd]^{\bar{\bar{g}}_2}\\
  \bar{Z}_1\sisi{}{\ar@{-->}[rd]}  \ar@{^{(}->}[rrr]^(.65){ j'_1} \ar[ru] \ar[dd]^{\bar{g}_1}   &           &               &
                              \bar{\bar{Z}}_1 \ar[ru]\ar@{-->}[rd] \ar[dd]^(.35){\bar{\bar{g}}_1}|\hole&\\
& **[r] Y\ar[rrr]_{i_2}  \ar@{^{(}->}[rrr]  & && **[r]\bar{Y}_2 \\
Y  \ar@{^{(}->}[rrr]_{i_1}   \ar@{=}[ru]  &           &               &
                              \bar{Y}_1\ar[ru]  &.\\
}
\]

There\pageoriginale are six ways to follow the edges of this cube from
$\bar{Z}_1$ to $\bar{Y}_2$, each ``adjacent'' to two others;
there is an isomorphism between the composite functors associated with
two adjacent paths (\ie paths separated by a single face), and the
commutativity of (4) will follow from the commutativity of the resulting
hexagon.

The moral of this kind of diagram is that an edge represents a
functor, a face an isomorphism between composite functors, and a
volume a compatibility condition (between the faces forming its
boundary). We decompose the cube into two triangular prisms (along
the ``plane'' $\bar{Z}_1 \bar{\bar{Z}}_1 Y \bar{Y}_2$) in order
to reduce to the hypotheses ${\rm(c')}$.

More precisely, we complete the hexagon to form the following diagram
of isomorphisms (for readability, the ${\ast}$ have
been omitted):
\[
\xymatrix{
                   & & \bar{\bar{g}}_2 j'_2p      & &   \\
i_2\bar{g}_2 p  \ar@{-}[rru] \ar@{-}[rd]\ar@{-}[dd] &   &   &&   \bar{\bar{g}}_2 p j'_1 \ar@{-}[ld]\ar[dd]\ar@{-}[llu]\\
 & i_2(g_2p) \ar@{-}[rr]&      & (\bar{\bar{g}}_2p)j'_1 & \\
i_21_{Y}\bar{g}_1 \ar@{-}[ru] \ar@{-}[rrd]&      &&   &
  p\bar{\bar{g}}_1 j'_1 \ar@{-}[lu] \ar@{-}[lld]\\
&  &  pi_1 \bar{g}_1  & &
}
\]
The\pageoriginale triangles are trivially commutative, and the
pentagons are commutative by ${\rm(c')}$.

\paragraph{Cocycle condition}\etiquette[3.3.2.4]{XVII.par3.3.2.4}

Let us verify that the isomorphisms $c_{f,g}$ constructed in \hyperref[XVII.par3.3.2.3]{3.3.2.3}
satisfy a cocycle condition for a composable triple $(f,g,h)$ of
morphisms. Let us introduce a compactification of
$(f,g,h)$, with diagram \sisi{(3.2.4.3)}{\eqhyperref[XVII.eq3.2.5.3]{3.2.5.3}}, and consider the following diagram
of isomorphisms of functors:
\[
\sisi{\xymatrix{
\bar{f} i \bar{g} j \bar{h} k \ar@{-}[rr] \ar@{-}[dd]&
&\bar{f}\bar{\bar{g}} j' j\bar{h} k \ar@{-}[dd]\\
 & \bar{f}\bar{\bar{g}}j' \bar{\bar{h}}
k' k \ar@{-}[rd]&\\
\bar{f} i \bar{g}\bar{\bar{h}}k' k \ar@{-}[ru] \ar@{-}[rr]&    &
\bar{f}\bar{\bar{g}}\bar{\bar{h}}k'' k' k.
}}{\xymatrix{
\bar{f} i \bar{g} j \bar{h} k \ar@{-}[rr]^{c_{f,g}*h_*} \ar@{-}[dd]_{f_**c_{g,h}}&&
\bar{f}\bar{\bar{g}} j' j\bar{h} k \ar@{-}[dd]\ar@{-}[rddd]^{c_{fg,h}}\\
&\\
\bar{f} i \bar{g}\bar{\bar{h}}k' k \ar@{-}[rrrd]_{c_{f,gh}} \ar@{-}[rr] && \bar{f}\bar{\bar{g}}j' \bar{\bar{h}}
k' k \ar@{-}[rd]\\
& &   & \bar{f}\bar{\bar{g}}\bar{\bar{\bar{h}}}k'' k' k.
}}
\]
The commutativity of the square is trivial, and that of the triangles follows from
(c) and ${\rm(c')}$. Since the commutativity of the outer boundary is what
had to be proved, this completes the proof of \hyperref[XVII.prop3.3.2]{3.3.2}.

\subsubsection{}\etiquette[3.3.3]{XVII.subsubsec3.3.3}

Suppose that over $(S,i)$ and $(S,p)$ there are given, respectively, fibered
categories equipped with normalized cleavages $F_i$ and $F_p$, and suppose
that $F_i$ and $F_p$ have the same fiber $F(X)$ for every $X\in \Ob S$.

According to \SGA 1 VI, it ``amounts to the same thing'' to give instead:
\begin{itemize}
\item[(i bis)] for every proper morphism $p:X\rightarrow Y$, a
functor $p^{\ast}:F(Y)\rightarrow F(X)$,

\item[(ii bis)] for every composable pair $(p,q)$ of proper
morphisms, an isomorphism of functors
$c_{p,q}:q^{\ast}p^{\ast}\leftrightarrow (pq)^{\ast}$,

\item[(${\rm i'}$ bis)] and\pageoriginale (${\rm ii'}$ bis), analogous data for
$(S,i)$, these data satisfying the conditions (a bis) (b bis) (${\rm a'}$ bis)
(${\rm b'}$ bis) dual to \hyperref[XVII.subsubsec3.3.1]{3.3.1} (a) (b) ${\rm (a')}$ ${\rm (b')}$.

Suppose \emph{given} in addition,

\item[(iii bis)] for every commutative diagram $D$
\[
\vcenter{
\xymatrix{
X \ar[r]^{j'} \ar[d]^{p'}&  Y\ar[d]^{p} \\
Z \ar[r]^{j}& T
}}\quad j,j' \in \Fl(S,i)\text{; }p\sisi{}{,} p' \in \Fl(S,p)
\]
an isomorphism of functors $d(D):{j'}^{\ast}p^{\ast}\leftrightarrow
{p'}^{\ast}j^{\ast}$.
\end{itemize}

We leave to the reader the formulation of the conditions (c bis) and (${\rm c'}$
bis) dual to \hyperref[XVII.prop3.3.2]{3.3.2} (c) and ${\rm (c')}$.

Arguments parallel to the preceding ones then prove the following
proposition.

\begin{enonce*}{Proposition 3.3.4}\etiquette[3.3.4]{XVII.prop3.3.4}
If data {\rm (i bis) (ii bis) (${\rm i'}$ bis) (${\rm ii'}$
bis) (iii bis)} satisfy the conditions {\rm (a bis) (b bis) (${\rm
a'}$ bis) (${\rm b'}$ bis) (c bis) (${\rm c'}$ bis)}, there exists one and
essentially only one category $F$ cofibered over $S$ and equipped
with isomorphisms $\fprod{F}{(S,i)}{S}=F_i$, $\fprod{F}{(S,p)}{S}=F_p$,
these isomorphisms being compatible with the identifications
$F_{iX}=F_{pX}$ $(X\in \Ob S)$ and such that the isomorphisms $d(D)$ of
{\rm (iii bis)} become the composite isomorphisms
${p'}^{\ast}j^{\ast}=(jp')^{\ast}=(pj')^{\ast}={j'}^{\ast}p^{\ast}$.
\end{enonce*}

It should be noted, however, that \hyperref[XVII.prop3.3.4]{3.3.4} is not dual to \hyperref[XVII.prop3.3.2]{3.3.2},
since the hypotheses imposed on $(S,i)$ and $(S,p)$ are not self-dual.

\section{Resolutions. Application to the base change
morphism}\etiquette[4]{XVII.sec4}\pageoriginale

\subsection{Flat resolutions}\etiquette[4.1]{XVII.subsec4.1}

Recall that if $f:(X,\Aa)\rightarrow (Y,\Bb)$ is a
morphism of ringed sites (or ringed topoi), the inverse image under $f$
of a sheaf of left $\Bb$-modules $\Mm$ is the
sheaf of left $\Aa$-modules
$\Aa\otimes_{f^{\ast}\Bb} f^{\ast}\Mm$.

\begin{enonce*}{Proposition 4.1.1}\etiquette[4.1.1]{XVII.prop4.1.1}
Let $f:(X,\Aa)\rightarrow (Y,\Bb)$ be a morphism of
ringed sites, let $K$ be a sheaf of right $\Aa$-modules, and let
$0\rightarrow L\rightarrow M\rightarrow N\rightarrow 0$ be an exact sequence
of sheaves of left $\Bb$-modules. If $M$ and $N$
are flat, then
\begin{enumerate}
\renewcommand{\theenumi}{\roman{enumi}}
\renewcommand{\labelenumi}{(\theenumi)}
\item $L$ is flat,

\item the following sequence is exact:
\[
0\longrightarrow K\otimes_{\Aa}f^{\ast}L\longrightarrow
K\otimes_{\Aa}f^{\ast}M\longrightarrow
K\otimes_{\Aa}f^{\ast}N\longrightarrow 0{\quoi}.
\]
\end{enumerate}
\end{enonce*}

Factor $f$ into the arrows
$(X,\Aa)\xrightarrow{u}(X,f^{\sisi{\ast}{-1}}\Bb)\xrightarrow{v}(Y,\Bb)$.
The functor $v^{\ast}$ is exact and takes flat modules to
flat modules (\sisi{IV 15}{\hyperrefext[V][V.sec1]{1}}), so it suffices to prove \hyperref[XVII.prop4.1.1]{4.1.1} for
$u$. The isomorphism
$$\etiquette[4.1.1.1]{XVII.eq4.1.1.1}
K\otimes_{\Aa}(\Aa\sisi{_{f^{\ast}\Bb}\otimes}{\otimes_{f^{-1}\Bb}} P)=K\otimes_{f^{\sisi{\ast}{-1}}\Bb}P\leqno(4.1.1.1)
$$
then reduces us to the case where $f$ is the identity.

Take an exact sequence $0\rightarrow R\rightarrow P\rightarrow
K\rightarrow 0$ with $P$ flat\sisi{ (\hyperref[XVII.subsubsec4.1.2]{4.1.2})}{}. The rows and columns of the
following diagram are exact:
\[
\xymatrix{
 &      &                          &  0\ar[d]                        &  &\\
 &      & R\otimes L  \ar[r]\ar[d] & P\otimes L  \ar[r]\ar[d] & K\otimes L
 \ar[r]\ar[d]& 0\\
 & 0\ar[r]     & R\otimes M  \ar[r]\ar[d] & P\otimes M  \ar[r]\ar[d] & K\otimes M
 \ar[r]\ar[d]& 0\\
 & 0\ar[r]     & R\otimes N  \ar[r] \ar[d]& P\otimes N  \ar[r]\ar[d] & K\otimes M
 \ar[r] \ar[d]& 0\\
 &      & 0         &   0   &   0   &.
}
\]

Apply\pageoriginale the snake lemma to the first two
columns: we obtain
\[
0\xrightarrow{j}K\otimes L\rightarrow K\otimes M{\quoi},
\]
which proves (ii). Now drop the assumption that $P$ is flat. In the
preceding diagram the three columns and two rows are short exact
sequences; the first row is therefore exact, \ie $L$ is flat.

\setcounter{subsubsection}{1}
\subsubsection{}\etiquette[4.1.2]{XVII.subsubsec4.1.2}
Let $\Ss$ be a subcategory of the category of ringed
sites.\nde{\MakeLinkTarget*{sga4.exp.XVII.editor-note.4.1.2.ringed-site-category}\label{XVII.ndesous}the theory is valid more generally for categories $\Ss$ consisting of ringed sites.} The sheaves of modules on the ringed sites in
$\Ss$ form a category $\Mm$, abelian over
$\Ss$ (\hyperref[XVII.def2.2.2]{2.2.2}), fibered and cofibered, whose fibers are the
opposites of the usual categories of sheaves of modules on a
site. We denote by $K(\Ss)$ the category $K(\Mm)$
(\hyperref[XVII.exem2.2.8]{2.2.8}) and by $K^{+}(\Ss)$ (\resp $K^{-}(\Ss)$) the
full subcategory of $K(\Ss)$ whose fiber categories
are the opposites of the categories $K^{+}(S)$ (\resp $K^{-}(S)$) for
$S\in \Ob \Ss$. The categories $K(\Ss)$,
$K^{+}(\Ss)$, and $K^{-}(\Ss)$ are triangulated over
$\Ss$, fibered and cofibered.

Let\pageoriginale $f:S\rightarrow S'$ be an arrow of $\Ss$. Every
injective object is acyclic (\hyperref[XVII.def1.2.3]{1.2.3}) for the functor
$Rf_{\ast}$, and every flat sheaf of modules is acyclic for the
functor $Lf^{\ast}$ (\hyperref[XVII.prop4.1.1]{4.1.1} and \hyperref[XVII.prop1.2.6]{1.2.6}),
whence the derived functors (\hyperref[XVII.prop1.2.7]{1.2.7})
\begin{align*}
Rf_{\ast}:D^{+}(S) &\longrightarrow D^{+}(S')\tag{4.1.2.1}\\
Lf^{\ast}:D^{-}(S\sisi{}{'}) &\longrightarrow D^{-}(S\sisi{'}{}){\quoi}.\tag{4.1.2.2}
\end{align*}\etiquette[4.1.2.1]{XVII.eq4.1.2.1}
\etiquette[4.1.2.2]{XVII.eq4.1.2.2}%The vertical space following this paragraph is too large.

If $Rf_{\ast}$ (\resp $Lf^{\ast}$) has finite dimension, then
in \eqhyperref[XVII.eq4.1.2.1]{4.1.2.1} (\resp \eqhyperref[XVII.eq4.1.2.2]{4.1.2.2}) one may replace the exponent $+$ (\resp $-$) by
$b$, by $-$ (\resp $+$), or omit it (\hyperref[XVII.prop1.2.10]{1.2.10}).\nde{\MakeLinkTarget*{sga4.exp.XVII.editor-note.4.1.2.spaltenstein-resolutions}\label{XVII.ndeKS}using Spaltenstein resolutions, the exponents may be omitted unconditionally. In the commutative case, see M.~Kashiwara and P.~Schapira, \emph{Categories and Sheaves}, Springer, 2010, cited as [KS] below, 18.6.}

We denote by $D(\Ss)$ (\resp $D^{+}(\Ss)$, \resp
$D^{-}(\Ss)$) the category obtained from $K(\Ss)$ (\resp
$K^{+}(\Ss)$, \resp $K^{-}(\Ss)$) by inverting the
quasi-isomorphisms.

\begin{enonce*}{Scholium 4.1.3}\etiquette[4.1.3]{XVII.sch4.1.3}
(``Trivial duality theorem'').
\begin{enumerate}
\renewcommand{\theenumi}{\alph{enumi}}
\renewcommand{\labelenumi}{(\theenumi)}
\item Assume that all the functors $Rf_{\ast}$ (\resp
$Lf^{\ast}$) have finite dimension. Then:
\begin{itemize}
\item[(i)] The formation of $D(\Ss)$ commutes with every change
of base category $\Ss'\rightarrow \Ss$. The category
$D(\Ss)$ is triangulated over $\Ss$ \emph{(\hyperref[XVII.def2.2.3]{2.2.3})}, and its
fibers identify with the opposites of the categories $D(S)$ for $S\in \Ob
\Ss$.

\item[(ii)] The category $D(\Ss)$ is cofibered (\resp fibered)
over $\Ss$. Let $f:S\rightarrow S'$; if $Lf^{\ast}$ (\resp
$Rf_{\ast}$) has finite dimension, then every object of $D(S')$
(\resp $D(S)$) admits an inverse image (\resp direct image) under $f$ in the sense
of the category $D(\Ss)$ over $\Ss$. The
inverse image and direct image functors identify with $Lf^{\ast}$ and
$Rf_{\ast}$.

\item[(iii)] Every\pageoriginale object of $D^{-}(S')$ (\resp $D^{+}(S)$) admits an
inverse image (\resp direct image) in the sense of $D(\Ss)$, and the
inverse image (\resp direct image) functor identifies with $Lf^{\ast}$
(\resp $Rf_{\ast}$).\nde{the finite-dimension hypotheses in (a) are superfluous. \Cf the preceding N.D.E.}
\end{itemize}

\item \begin{itemize}
\item[(i)] The formation of $D^{-}(\Ss)$ (\resp
$D^{+}(\Ss)$) commutes with every change of base category
$\Ss'\rightarrow \Ss$. The category
$D^{-}(\Ss)$ (\resp $D^{+}(\Ss)$) is triangulated over
$\Ss$, and its fibers identify with the opposites of the categories
$D^{-}(S)$ (\resp $D^{+}(S)$).

\item[(ii)] The category $D^{-}(\Ss)$ (\resp
$D^{+}(\Ss)$) is fibered (\resp cofibered) over $\Ss$;
the ``inverse image'' (\resp ``direct image'')
functors identify with the derived functors $Lf^{\ast}$ (\resp $Rf_{\ast}$).

\item[(iii)] If $f:S\rightarrow S'$ is such that $Rf_{\ast}$ (\resp
$Lf^{\ast}$) has finite dimension, every object of $D^{-}(S)$ (\resp
$D^{+}(S')$) admits a direct image in the sense of $D^{-}(\Ss)$
(\resp an inverse image in the sense of $D^{+}(\Ss)$), and the
``direct image'' (\resp ``inverse image'')
functor identifies with $Rf_{\ast}$ (\resp $Lf^{\ast}$).
\end{itemize}
\end{enumerate}
\end{enonce*}

It suffices to verify that in each case the hypotheses of
\hyperref[XVII.prop2.4.1]{2.4.1} or
\hyperref[XVII.prop2.4.3]{2.4.3} are satisfied, the other assertions following, after a
base change $\{S\}\rightarrow \Ss$, or $I\rightarrow
\Ss$, from \hyperref[XVII.prop2.3.5]{2.3.5} and \hyperref[XVII.prop1.2.10]{1.2.10}.
In case (b), complexes of
flat (\resp flasque) sheaves satisfy the hypotheses of
\sisi{\hyperref[XVII.prop2.4.1]{2.4.1}}{\hyperref[XVII.prop2.4.3]{2.4.3} (\resp \hyperref[XVII.prop2.4.1]{2.4.1})}. In case (a), we use \sisi{\hyperref[XVII.prop2.4.3]{2.4.3}}{\hyperref[XVII.prop2.4.1]{2.4.1} (\resp \hyperref[XVII.prop2.4.3]{2.4.3})}
so that we need only
consider those $f$ for which the dimension of $Rf_{\ast}$ (\resp
$Lf^{\ast}$) is bounded by a fixed $N$. Let us prove that there are ``enough'' complexes $K$ whose components have their
direct (\resp inverse) images all acyclic for the
direct image (respectively, inverse image) functors under consideration. Let $K$ be a
complex\pageoriginale and let $K^{\ast}$ (\resp $K_{\ast}$) be a Cartan-Eilenberg resolution
of $K$ by complexes with flasque components
(\resp flat components).\nde{it is required only that the augmentation $K^p\to (K^*)^{p\bigdot}$ (\resp $(K_*)^{p\bigdot}\to K^p$) be a flasque (\resp flat) resolution for every $p$.} If $n>N$, we see that
$f_{\ast}\tau''_{<n}(K^{\ast})=\tau''_{<\sisi{}{n}}f_{\ast}(K^{\ast})$ (\resp
$f^{\ast}\tau''_{>-n}(K_{\ast})=\tau''_{>-n}f^{\ast}(K_{\ast})$), where
$\tau\sisi{^{\ast}}{''}$ denotes a truncation with respect to the second
differential. For $n>2N$, it follows that $\tau''_{<n}(K^{\ast})$
(\resp $\tau''_{>-n}(K_{\ast})$) satisfies the hypothesis of \sisi{\hyperref[XVII.prop2.4.3]{2.4.3}}{\hyperref[XVII.prop2.4.1]{2.4.1} (\resp \hyperref[XVII.prop2.4.3]{2.4.3})}, and every
complex $K$ is the source (\resp target) of a quasi-isomorphism
whose target (\resp source) is such a complex.

\setcounter{subsubsection}{3}
\subsubsection{}\etiquette[4.1.4]{XVII.subsubsec4.1.4}
Consider a commutative square of ringed sites \eqhyperref[XVII.eq4.1.4.1]{4.1.4.1}
(\cf~\hyperref[XVII.subsec0.11]{0.11})
$$\etiquette[4.1.4.1]{XVII.eq4.1.4.1}
\vcenter{
\xymatrix{
  X' \ar[r]^{g'}\ar[d]_{f'}& X\ar[d]^{f}\\
S' \ar[r]^{g} & S\rlap{.}
}}\leqno{(4.1.4.1)}
$$
We may consider the subcategory of the category of sites consisting
of this diagram;\nde{one must instead consider the category of ringed sites defined by this diagram, which is not strictly speaking a subcategory of the category of ringed sites in general. \Cf \nderef{XVII.ndesous}.} we then use \sisi{\hyperref[XVII.subsubsec2.1.2]{2.1.2}}{\hyperref[XVII.prop2.1.3]{2.1.3}} to
define a ``base change morphism'' from $Lg^{\ast}Rf_{\ast}$ to $Rf'_{\ast}L{g'}^{\ast}$ whenever \hyperref[XVII.sch4.1.3]{4.1.3} is
applicable:

\begin{enonce*}{Proposition 4.1.5}\etiquette[4.1.5]{XVII.prop4.1.5}
The base change morphism of \sisi{\hyperref[XVII.subsubsec2.1.2]{2.1.2}}{\hyperref[XVII.prop2.1.3]{2.1.3}}, from $Lg^{\ast}Rf_{\ast}$
to $Rf'_{\ast}L{g'}^{\ast}$, is defined in each of the
following cases:
\begin{enumerate}
\renewcommand{\theenumi}{\roman{enumi}}
\renewcommand{\labelenumi}{(\theenumi)}
\item We work in the categories $D^{+}$, and $Lg^{\ast}$ and
$L{g'}^{\ast}$ have finite dimension.

\item We work in the categories $D^{-}$, and $Rf_{\ast}$ and
$Rf'_{\ast}$ have finite dimension.

\item We work in the unbounded derived categories, and
$Rf_{\ast}$, $Rf'_{\ast}$, $Lg^{\ast}$, and $L{g'}^{\ast}$ have
finite dimension, as do $Lf^{\ast}$ and $L{f'}^{\ast}$ (or: as do
$Rg_{\ast}$ and $Rg'_{\ast}$).\nde{the finite-dimension hypotheses in (iii) are superfluous. \Cf \nderef{XVII.ndeKS}.}
\end{enumerate}
\end{enonce*}

It\pageoriginale should be noted that the last hypothesis in (iii) looks
rather like a practical joke. The editor does not know how to dispense with it except when
the sites under consideration have enough points (\hyperref[XVII.subsubsec4.2.12]{4.2.12}). The functorial
properties of this base change morphism follow
at once from the description in \sisi{\hyperref[XVII.subsubsec2.1.2]{2.1.2}}{\hyperref[XVII.prop2.1.3]{2.1.3}}.

\setcounter{subsubsection}{5}
\subsubsection{}\etiquette[4.1.6]{XVII.subsubsec4.1.6}
Let $(S,\Aa)$ be a ringed site, and consider the ``tensor product'' functor
\[

% END INCLUDED SGA4_English_current_buildable_source_20260906/source/17/drafts/SGA4_XVII_LINES_2601_3000_English.texfrag

% BEGIN INCLUDED SGA4_English_current_buildable_source_20260906/source/17/drafts/SGA4_XVII_LINES_3001_3400_English.texfrag
\otimes:\Mod\sisi{_d}{}(S,\Aa^{\circ})\times
\Mod\sisi{_g}{}(S,\Aa)\longrightarrow \Mod(S,C)
\]%For the first component of the source, take either right Aa-modules or left Aa-circ-modules.
where $C$ is the center of $\Aa$.

The recipe of \hyperref[XVII.subsubsec1.1.6]{1.1.6}, in applying which
the simple complex associated with a double complex is defined using a
sum, extends this functor to a functor still denoted $\otimes$:
\[
\otimes:K(S,\Aa^{\circ})\times K(S,\Aa)\longrightarrow
K(S,C){\quoi}.
\]

\begin{enonce*}{Proposition 4.1.7}\etiquette[4.1.7]{XVII.prop4.1.7}
With the preceding notation:
\begin{enumerate}
\renewcommand{\theenumi}{\roman{enumi}}
\renewcommand{\labelenumi}{(\theenumi)}
\item The tensor-product functor is left derivable
(\hyperref[XVII.def1.2.1]{1.2.1}) to a ``total tensor product'' functor
\[
{}\overset{L}{\otimes}{}:D^{-}(S,\Aa^{\circ})\times
D^{-}(S,\Aa)\longrightarrow D^{-}(S,C){\quoi}\text{;}
\]

\item a pair $(K,L)$ of bounded-above complexes is adapted to
$\overset{L}{\otimes}$ whenever either $K$ or $L$ has flat components.
\end{enumerate}
\end{enonce*}

The proof is standard.

\begin{enonce*}{Notation 4.1.8}\etiquette[4.1.8]{XVII.not4.1.8}
If\pageoriginale $K\in D^{-}(S,\Aa^{\circ})$ and
$L\in D^{-}(S,\Aa)$, denote by $K\overset{L}{\otimes}L$ the total
tensor product of $K$ and $L$. The sheaf
$H_{k}(K\overset{L}{\otimes}L)=H^{-k}(K\overset{L}{\otimes}L)$ is
denoted by $\Tor_k(K,L)$ and is called the $k$\textsuperscript{th}
local hypertor of $K$ and $L$.
\end{enonce*}

Let $K\in D^{-}(S,\Aa^{\circ})$ and $L\in D^{-}(S,\Aa)$. If
$H^{i}(K)=0$ for $i>k$ and $H^{i}(L)=0$ for $i>\ell$, one checks that
$H^{i}(K\overset{L}{\otimes}L)=0$ for $i>k+\ell$ and that
$H^{k+\ell}(K\overset{L}{\otimes}L)=H^{k}(K)\otimes H^{\ell}(L)$.
Fix $K$ and regard $K\overset{L}{\otimes}L$ as a functor of $L$. If
Definition \hyperref[XVII.def1.2.9]{1.2.9} of cohomological amplitude
as an interval in $\ZZ$ is applied, the upper endpoint of this interval
is trivial in nature: it is the largest integer $i$ such that
$H^{i}(K)\neq0$. The lower endpoint is more interesting:

\begin{enonce*}{Definition 4.1.9}\etiquette[4.1.9]{XVII.def4.1.9}
A complex $K\in\Ob D^{b}(S,\Aa^{\circ})$ is said to have
\emph{\tTor-dimension $\leq n$} if the following equivalent conditions
hold:
\begin{enumerate}
\renewcommand{\theenumi}{\roman{enumi}}
\renewcommand{\labelenumi}{(\theenumi)}
\item For every $\Aa$-module $L$, $\Tor_k(K,L)=0$ for $k>n$.

\item For every complex of $\Aa$-modules $L$ such that $H^i(L)=0$ for
$i<\ell$, one has $H^i(K\overset{L}{\otimes}L)=0$ for $i<\ell-n$.

\item There is a quasi-isomorphism with target $K$ and source a
\sisi{}{ bounded} complex $K'$ whose components are flat and vanish for
$i<-n$.
\end{enumerate}
\end{enonce*}

\setcounter{subsubsection}{9}
\subsubsection{}\etiquette[4.1.10]{XVII.subsubsec4.1.10}
Here are some possible variants:
\begin{enumerate}
\renewcommand{\theenumi}{\roman{enumi}}
\renewcommand{\labelenumi}{(\theenumi)}
\item If $K\in\Ob D^{b}(S,\Aa^{\circ})$ has finite \ttor-dimension
and $L\in\Ob D(S,\Aa)$, then the functor $\overset{L}{\otimes}$ is
defined (\hyperref[XVII.def1.2.1]{1.2.1}) at $(K,L)$.

\item If $K\in\Ob K^{b}(S,\Aa^{\circ})$ has components of finite
\tTor-dimension, the functor $L\mapsto K\otimes L$
$(L\in K(S,\Aa))$ is left derivable.

\item If every module has finite \tTor-dimension,\nde{this condition is
superfluous. In the case where $\Aa$ is commutative, see [KS], 18.6,
cited in \nderef{XVII.ndeKS}.} the functor $\otimes$ derives to a
functor
\[
{}\overset{L}{\otimes}{}:D(S,\Aa^{\circ})\times
D(S,\Aa)\longrightarrow D(S,\Cc){\quoi}.
\]
\end{enumerate}

\subsection{Godement flasque resolutions}\etiquette[4.2]{XVII.subsec4.2}\pageoriginale

\subsubsection{}\etiquette[4.2.1]{XVII.subsubsec4.2.1}
Let $S$ be a topos and $X$ a set. To give a family of point\sisi{}{s}
of $S$ indexed by $X$ is equivalently to give a morphism
$k:\Top(X)\rightarrow S$ from the topos defined by the discrete
topological space $X$ to $S$. For this family to be conservative
(IV \hyperref[IV.def6.4.1]{6.4.1}), it is necessary and sufficient that, for every sheaf $F$
on $S$, the adjunction arrow from $F$ to $k_{\ast}k^{\ast}F$ be a
monomorphism (\cf IV \hyperref[IV.subsubsec6.4.0]{6.4.0}).

\begin{enonce*}{Definition 4.2.2}\etiquette[4.2.2]{XVII.def4.2.2}
Let $k:\Top(X)\rightarrow S$ be a conservative family of points of a
topos $S$. The \emph{Godement resolution}, or \emph{canonical flasque
resolution} (relative to $X$), of an abelian sheaf $F$ on $S$, denoted
$\Cc_X^{\ast}(F)$ (or simply $\Cc^{\ast}(F)$), is the following right
resolution of $F$:
\begin{enumerate}
\renewcommand{\labelenumi}{\theenumi)}
\item $\Cc^{0}(F)=k_{\ast}k^{\ast}F$ and
$\epsilon:F\rightarrow\Cc^{0}(F)$ is the adjunction arrow.

\item Set $d^{-1}=\epsilon$; recursively, for $n\geq0$, define

$\Cc^{n+1}(F)=\Cc^{0}(\coker d^{n-1})$, and define $d^n$ as the
composite arrow
\[
d^{n}:\Cc^n(F)\rightarrow
\coker(d^{n-1})\rightarrow\Cc^{0}(\coker(d^{n-1})){\quoi}.
\]
\end{enumerate}
\end{enonce*}

\begin{enonce*}{Proposition 4.2.3}\etiquette[4.2.3]{XVII.prop4.2.3}
Under the hypotheses of \emph{\hyperref[XVII.def4.2.2]{4.2.2}},
\begin{enumerate}
\renewcommand{\theenumi}{\roman{enumi}}
\renewcommand{\labelenumi}{(\theenumi)}
\item $\Cc^n(F)$ is a flasque sheaf
\emph{\sisi{(V 3.7)}{\hyperrefext[V][V.def4.1]{4.1}}},

\item $\Cc^n(F)$ is an exact functor of $F$,

\item the stalk at a point $x\in X$ of the complex $\Cc^{\ast}(F)$ is
a canonically split resolution of $F_x$.
\end{enumerate}
\end{enonce*}

The\pageoriginale sheaves $\Cc^n(F)$ are flasque because they are
direct images of sheaves (automatically flasque) on the discrete
topological space $X$.

The functors $k^{\ast}$ and $k_{\ast}$ are exact. Hence the functor
$\Cc^0$ is exact. We prove by induction on $n\geq0$ that the functor
$\mathcal{Z}^{n}(F)=\coker(d^{n-1})$ is exact, and therefore so is
$\Cc^{n+1}=\Cc^0\mathcal{Z}^{n}$. Set $\Zz^{-1}(F)=F$. An exact
sequence of sheaves
\[
0\longrightarrow F'\longrightarrow F\longrightarrow F''\longrightarrow0
\]
defines, for $n\geq0$, a diagram
\[
\xymatrix{
        &   0   \ar[d]             &  0   \ar[d]            &  0 \ar[d]\\
0  \ar[r] & \Zz^{n-1}(F\sisi{^1}{'})  \ar[r] \ar[d]& \Zz^{n-1}(F)  \ar[r] \ar[d]& \Zz^{n-1}(F'')
\ar[r] \ar[d]& 0 \\
0  \ar[r] & \Cc^n(F')  \ar[r] \ar[d]& \Cc^n(F)  \ar[r] \ar[d]&
\Cc^n (F'')  \ar[r] \ar[d]& 0 \\
 0  \ar[r] & \Zz^{n}(F')  \ar[r]\ar[d] & \Zz^{n}(F)  \ar[r] \ar[d]& \Zz^{n}(F'')  \ar[r] \ar[d]&
0\\
        &  0               &  0             &  0      &,
}
\]
whose columns are exact, as are, by the induction hypothesis, the
first two rows. The 3\sisi{rd}{\textsuperscript{rd}} row is therefore
exact, proving (ii).

The short exact sequences
\[
0\longrightarrow\mathcal{Z}^{n+1}(F)\longrightarrow
\Cc^{n}(F)\longrightarrow\mathcal{Z}^{n}(F)\longrightarrow0
\]
are special cases of the short exact sequence
$$\etiquette[4.2.3.1]{XVII.eq4.2.3.1}
0\rightarrow F\xrightarrow{\epsilon}\Cc^{0}(F)\rightarrow
\mathcal{Z}^{0}(F)\rightarrow0\leqno(4.2.3.1)
$$
(make\pageoriginale the substitution $F\mapsto\mathcal{Z}^{n-1}(F)$).
To prove (iii), it remains to show that the inverse-image sequence of
\eqhyperref[XVII.eq4.2.3.1]{4.2.3.1} under $k$ is canonically split.
Indeed, the composite arrow
\[
k^{\ast}F\xrightarrow{k^{\ast}(\epsilon)}k^{\ast}\Cc^{0}F=k^{\ast}(k_{\ast}k^{\ast})F=(k^{\ast}k_{\ast})k^{\ast}F\longrightarrow
k^{\ast}F
\]
is the identity.

\begin{enonce*}[remark]{Remark 4.2.4}\etiquette[4.2.4]{XVII.rem4.2.4}
If $(x_i)_{i\in X}$ is a conservative family of points of a site, this
family corresponds to a conservative family of points of the
associated topos, which again makes it possible to define canonical
flasque resolutions.
\end{enonce*}

\begin{enonce*}[remark]{Remark 4.2.5}\etiquette[4.2.5]{XVII.rem4.2.5}
If $S$ is a ringed topos, the functors $\Cc^n$ carry modules to modules.
\end{enonce*}

\begin{enonce*}[remark]{Remark 4.2.6}\etiquette[4.2.6]{XVII.rem4.2.6}
Let $\Ss$ be a category consisting of sites. Suppose that, for each
$S\in\Ob\Ss$, a set $S^d$ and a conservative family
$k_S:\Top(S^d)\rightarrow S$ of points of $S$ indexed by $S^d$ are
given. Suppose that $S^d$ is functorial in $S$ and that the diagrams
\[
\xymatrix{
\Top(S^d)   \ar[r]^{f^d} \ar[d]^{k_S}& \Top(T^d)\ar[d]^{k_T}\\
S        \ar[r]^{f}  & T
}
\]
are commutative. The corresponding canonical flasque resolutions are
then ``functorial in $S$.''
\end{enonce*}

If\pageoriginale $\Ss$ is the category of topological spaces, one may
take $S^d$ to be the underlying set.

If $\Ss$ is the category of the étale sites of schemes $S$ locally of
finite type over a scheme $S_0$, and if, for each $s\in S_0$,
$\bar{k(s)}$ is an algebraic closure of $k(s)$, one may take for $S^d$
the set of points of $S$ with values in one of the $\bar{k(s)}$.

The condition ``of finite type'' may be replaced by a cardinality
condition if the algebraic closures $\bar{k(s)}$ are replaced by
sufficiently large algebraically closed extensions.

\setcounter{subsubsection}{6}
\subsubsection{}\etiquette[4.2.7]{XVII.subsubsec4.2.7}
One advantage of canonical flasque resolutions is that they make it
possible to construct complexes adapted simultaneously to functors of
``tensor-product'' type and to functors of ``direct-image'' type.
\sisi{It}{They} thereby give a perhaps more comprehensible
construction of the base-change morphisms
(\hyperref[XVII.prop4.1.5]{4.1.5}).

\begin{enonce*}{Lemma 4.2.8}\etiquette[4.2.8]{XVII.lem4.2.8}
Let $f:(S,\Aa)\rightarrow(S',\Bb)$ be a morphism of ringed sites, let
$K\in\Ob K(S,\Aa^{\circ})$ and $L\in\Ob K(S',\Bb)$ be complexes, and
let $k:\Top(X)\rightarrow S$ be a conservative set of points of $S$.
Suppose that one of the following conditions holds:
\begin{enumerate}
\renewcommand{\theenumi}{\alph{enumi}}
\renewcommand{\labelenumi}{(\theenumi)}
\item The cohomologies of $K$ and $L$ are bounded above;

\item $\Aa$ has finite \tTor-dimension over
$f^{\sisi{\ast}{-1}}\Bb$, and $K\in\Ob D^b(S,\Aa^{\circ})$ has finite
\tTor-dimension;

\item $L\in\Ob D^b(S',\Bb)$ has finite \tTor-dimension.
\end{enumerate}

Then\pageoriginale the functor $(K,L)\mapsto K\otimes f^{\ast}L$ is
derivable\sisi{}{ on the left} \emph{(\hyperref[XVII.def1.2.1]{1.2.1})} at
$(K,L)$; denote its derived functor by
${}\overset{\LL}{\otimes}\LL f^{\ast}$.
\begin{enumerate}
\renewcommand{\theenumi}{\roman{enumi}}
\renewcommand{\labelenumi}{(\theenumi)}
\item If $x$ is a point of $S$, then
$(K\overset{\LL}{\otimes}\LL f^{\ast}L)_x\simeq
K_x\overset{\LL}{\otimes}_{\Bb_{f(x)}}L_{f(x)}$.

\item The pair $(K,L)$ is left adapted to
${}\overset{\LL}{\otimes}f^{\ast}$
\eqhyperref[XVII.def1.2.1]{1.2.1} if and only if, for every $x\in X$,
the pair $(K_x,L_{f(x)})$ is left adapted to the functor
$\otimes_{\Bb_{f(x)}}$.

\item Under hypotheses $(a)$ or $(b)$, the functor
$K\overset{\LL}{\otimes}f^{\ast}L$ in $L$ has cohomological amplitude
$\subset d$ \emph{(\hyperref[XVII.def1.2.9]{1.2.9})} if and only if,
for every $x\in X$, the complex $K_x$ of $\Bb_{f(x)}$-modules has
\tTor-amplitude $\subset d$.\nde{hypotheses (a), (b), and (c) are all
superfluous. \Cf \nderef{XVII.ndeKS}.}
\end{enumerate}
\end{enonce*}

This lemma, whose proof is standard, expresses the ``pointwise nature''
of the functor under consideration.

\setcounter{subsubsection}{8}
\subsubsection{}\etiquette[4.2.9]{XVII.subsubsec4.2.9}

Let $L$ be a complex. The \emph{Godement resolution} (relative to $X$)
of $L$ is the simple complex (\emph{defined using sums}) associated
with the double complex $\sisi{C}{\Cc}^{\ast}L$
(\hyperref[XVII.subsubsec1.1.4]{1.1.4}). The \emph{Godement resolution
truncated at order} $n$, $\tau''_{\leq n}\sisi{C}{\Cc}^{\ast}L$, is
the simple complex associated with the double complex obtained from
$\sisi{C}{\Cc}^{\ast}L$ by truncation
(\hyperref[XVII.subsubsec1.1.15]{1.1.15}) in the direction of the
Godement differential. By \hyperref[XVII.prop4.2.3]{4.2.3}~(iii), $L$
is, point by point on $X$, homotopy equivalent to the complexes
$\tau''_{\leq n}\sisi{C}{\Cc}^{\ast}L$ and
$\sisi{C}{\Cc}^{\ast}L$; these complexes are therefore ``as good as
$L$'' with respect to functors of ``tensor-product'' type
(\hyperref[XVII.lem4.2.8]{4.2.8}).

\begin{enonce*}{Proposition 4.2.10}\etiquette[4.2.10]{XVII.prop4.2.10}
Let $(S,\Aa)$ be a ringed site, let $X$ be a conservative set of
points of $S$, and let $N,M\in\NN\cup\{\infty\}$, with $N$ or $M$
finite. Let $A$ be the subcategory of $K(S)$ consisting of complexes
$L$ such that:
\begin{enumerate}
\renewcommand{\theenumi}{\alph{enumi}}
\renewcommand{\labelenumi}{(\theenumi)}
\item If $N=\infty$, then $H^i(L)=0$ for $i$ sufficiently large; if
$M=\infty$, then $H^i(L)=0$ for $i$ sufficiently small.

\item For every\pageoriginale $x\in X$ and every $\Aa_x$-module $Q$
of \tTor-dimension $\leq N$, the complex $L_x$ is left adapted to the
functor $\sisi{M_x}{Q}\otimes_{\Aa_x}\ast$.

\item For every morphism of topoi
$f:\sisi{X^{\sim}}{S}\rightarrow E$ of cohomological dimension
$\leq M$, the complex $L$ is right adapted to the functor $f_{\ast}$.
\end{enumerate}

Then the inclusion of $A$ into $K(S)$ induces an equivalence between
the category obtained from $A$ by inverting quasi-isomorphisms and the
category $D^{\ast}(S)$, where $\ast=+$ if $N<M=\infty$, $\ast=-$ if
$M<N=\infty$, and $\ast$ is blank if $N,M<\infty$.
\end{enonce*}

Let $K$ be a complex satisfying (a). If $N=\infty$, then the canonical
arrow $a:K_1=\sisi{\sigma}{\tau}_{\leq i}K\rightarrow K$ is a
quasi-isomorphism for $i$ sufficiently large. If $N<\infty$, set
$K_1=K$. Let $K_1^{\ast}$ be a flat Cartan--Eilenberg resolution and
set $K_2=\tau''_{\geq\sisi{}{-}N}K_1$ (truncation in the direction of
the new differential). The composite arrow from $K_2$ to $K$ is a
quasi-isomorphism, and $K_2$ satisfies (b).

Let $K$ be a complex satisfying (a) and (b). If $M=\infty$, then
$N<\infty$; for $i$ sufficiently large,
$b:K\rightarrow\sisi{\sigma}{\tau}_{\geq i}K_1=K_3$ is a
quasi-isomorphism, \sisi{from}{and} $K_3$ still satisfies (a) and (b).
If $M<\infty$, set $K=K_3$. Let
$K_4=\sisi{\sigma}{\tau}''_{\leq M}\sisi{C}{\Cc}^{\ast}K_3$. Then,
by \hyperref[XVII.prop4.2.3]{4.2.3}~(iii), $K_4$ satisfies
\sisi{(a), (c), and (b)}{(a), (b), and (c)}.

The proposition therefore follows from two successive applications of
\cite[I 2.4.2]{sga4.exp.XVII.bib.CD}.

\begin{enonce*}[remark]{Variants 4.2.11}\etiquette[4.2.11]{XVII.var4.2.11}
If $N<\infty$ (\resp if $M<\infty$, \resp if $N,M<\infty$), one could
equally have worked with bounded-below (\resp bounded-above, \resp
bounded) complexes.
\end{enonce*}

\setcounter{subsubsection}{11}
\subsubsection{}\etiquette[4.2.12]{XVII.subsubsec4.2.12}
Here\pageoriginale is how these constructions may be used to define
the \emph{base-change morphism} of
\hyperref[XVII.subsubsec4.1.4]{4.1.4} under hypotheses more general
than those of \loccit, when the sites under consideration have enough
points.

Let there be a commutative diagram (\cf
(\hyperref[XVII.subsec0.11]{0.11})) of sites
$$\etiquette[4.2.12.1]{XVII.eq4.2.12.1}
\vcenter{
\xymatrix{
X'          \ar[r]^{g'} \ar[d]^{f'}& X\ar[d]^{f}\\
S'          \ar[r]^{g}  & S\rlap{,}
}}\leqno{(4.2.12.1)}
$$
let $\Aa$ be a sheaf of rings on $S$, $\Aa'$ a sheaf of rings on
$S'$, and $L$ a complex of $(\Aa',g^{\ast}\Aa)$-bimodules on $S'$.
Suppose that:
\begin{enumerate}
\renewcommand{\theenumi}{\roman{enumi}}
\renewcommand{\labelenumi}{(\theenumi)}
\item the sites $X$ and $X'$ admit conservative families of points,
$k:Q\rightarrow X$ and $k':Q'\rightarrow X'$, fitting into a
commutative diagram
$$\etiquette[4.2.12.2]{XVII.eq4.2.12.2}
\vcenter{
\sisi{\xymatrix{
X'         \ar[r]^{g'}  \ar[d]^{k'} & X     \ar[d]^{k}  & \\
Q'         \ar[r]^{g'_0} & Q &.
}}{\xymatrix{
Q'         \ar[r]^{g'_0}  \ar[d]^{k'} & Q     \ar[d]^{k}  & \\
X'         \ar[r]^{g'} & X \rlap{.}
}}}\leqno{(4.2.12.2)}
$$

\item One of the following hypotheses holds:
\begin{enumerate}
\renewcommand{\theenumi}{\alph{enumi}}
\renewcommand{\labelenumi}{(\theenumi)}
\item The functors $Rf_{\ast}$ and $Rf'_{\ast}$ have finite
cohomological dimension, and $L$ is bounded above. Set $s=-$.

\item $L$ is bounded below, with components of uniformly bounded
\tTor-dimension (as $g^{\ast}\Aa$-modules). Set
$s=+$.

% END INCLUDED SGA4_English_current_buildable_source_20260906/source/17/drafts/SGA4_XVII_LINES_3001_3400_English.texfrag

% BEGIN INCLUDED SGA4_English_current_buildable_source_20260906/source/17/drafts/SGA4_XVII_LINES_3401_3800_English.texfrag

\item $L$ is bounded, and the hypotheses of (a) and (b) hold. Set
$s=$ blank.
\end{enumerate}
\end{enumerate}

We\pageoriginale denote by
$L\overset{\LL}{\otimes}\LL g^{\ast}:D^{s}(S,\Aa)\rightarrow
D^{s}(S', \Aa')$ the left derived functor (\hyperref[XVII.def1.2.1]{1.2.1}) of the functor
$K\mapsto L\otimes_{g^{\ast}\Aa}g^{\ast}K$ from
$K^{s}(S,\Aa)$ to $K^{s}(S',\Aa')$.

We shall define a base-change morphism
$$\etiquette[4.2.12.3]{XVII.eq4.2.12.3}
\varphi:L\overset{\LL}{\otimes}\LL g^{\ast}Rf_{\ast}K\longrightarrow
Rf'_{\ast}({f'}^{\ast}L\overset{\LL}{\otimes}\LL{g'}^{\ast}K)\leqno(4.2.12.3)
$$
for $K$ in $D^{s}(X,\Aa)$.\nde{\MakeLinkTarget*{sga4.exp.XVII.editor-note.4.2.12.3.unbounded-adjunction}\label{XVII.ndecb}Adjunctions in unbounded derived categories make it possible to define such a morphism without hypotheses (i) and (ii). \Cf \nderef{XVII.ndeKS}.}

By \hyperref[XVII.prop4.2.10]{4.2.10}, it suffices to define it, functorially in $K$,
for the complexes belonging to the category $A$ of \hyperref[XVII.prop4.2.10]{4.2.10}, with
\begin{align*}
N & \geq\sup_i \text{ (tor\sisi{}{-}dimension of $L^{i}$ over
$g^{\ast}$)}\\
M & \geq\text{cohomological dimension of $Rf_{\ast}$}.
\end{align*}
Such a complex is adapted both to the functor $Rf_{\ast}$ and to
the functor ${f'}^{\ast}L\overset{\LL}{\otimes}_{\Aa'}\sisi{}{\LL}{g'}^{\ast}$
(\hyperref[XVII.lem4.2.8]{4.2.8}). The canonical morphisms
\begin{align*}
\alpha &: f_{\ast}K\longrightarrow Rf_{\ast}K\\
\beta &:
{f'}^{\ast}L\overset{\LL}{\otimes}\LL \sisi{g}{{g'}}^{\ast}K\longrightarrow
\sisi{}{{f'}^*}L\otimes \sisi{g}{{g'}}^{\ast}K
\end{align*}
are therefore quasi-isomorphisms. To compute the left-hand side
$I$ of \eqhyperref[XVII.eq4.2.12.3]{4.2.12.3}, it therefore suffices to take a left resolution
(\ie a quasi-isomorphism)
\[
\gamma:K_{1}\longrightarrow f_{\ast}K{\quoi},
\]
with $K_{1}$ adapted to
$L\overset{\LL}{\otimes}\LL g^{\ast}$; we
have
\[
I\simeq L\otimes \sisi{\LL}{} g^{\ast}K_{1}{\quoi}.
\]

Similarly,\pageoriginale to compute the right-hand side $II$ of
\eqhyperref[XVII.eq4.2.12.3]{4.2.12.3}, it suffices to take a right resolution
\[
\delta:{f'}^{\ast}L\otimes {g'}^{\ast}K\longrightarrow K_2{\quoi},
\]
with $K_{2}$ adapted to $Rf'_{\ast}$; then
\[
II\simeq f'_{\ast}K_{2}{\quoi}.
\]
We then define the base-change morphism as the composite
\[
I\simeq L\otimes g^{\ast}K_{1}\xrightarrow{\gamma}L\otimes
g^{\ast}f_{\ast}K\xrightarrow{c}f'_{\ast}({f'}^{\ast}L\otimes
{g'}^{\ast}K)\xrightarrow{\sisi{\alpha}{\delta}}f'_{\ast}K_{2}\simeq II{\quoi},
\]
where $c$ is the usual base-change morphism
at the level of actual complexes.

This base-change morphism has the following characteristic
property, which in principle makes it easy to verify its compatibility
properties.

\begin{enonce*}{Proposition 4.2.13}\etiquette[4.2.13]{XVII.prop4.2.13}
Under the preceding hypotheses, let $K\in \Ob
C^{s}(X\sisi{'}{},\Aa)$, $K_{1}\in \Ob C^{s}(\sisi{X}{S},\Aa)$ and
$K_{2}\in \Ob C^{s}(\sisi{S}{X'},\Aa)$. Moreover, let $u$ be an
$f$-morphism from $K$ to $K_{1}$, \ie $u\in
\Hom(K_{1},f_{\ast}K)\simeq \Hom(f^{\ast}K_{1},K)$, and let $v$ be a
morphism from ${f'}^{\ast}L\otimes {g'}^{\ast}K$ to $K_{2}$. From
$u$, $v$, and the usual base-change morphism, one obtains a
morphism of complexes $c_{u,v}:L\otimes g^{\ast}K_{1}\rightarrow
f'_{\ast}\sisi{k}{K}_2$. Let $u'$, $v'$, and $c'_{uv}$ denote the composite morphisms
\begin{align*}
u' &: K_{1}\xrightarrow{u} f_{\ast}K\rightarrow Rf_{\ast}K\\
v' &: L\overset{\LL}{\otimes}\LL{g'}^{\ast}K\rightarrow
L\otimes {g'}^{\ast}K\xrightarrow{v} K_{2}\\
c'_{u,v} &:
L\overset{\LL}{\otimes}\LL{g'}^{\ast}K_{1}\rightarrow
L\otimes {g'}^{\ast}K_{1}
\xrightarrow{c_{uv}}f'_{\ast}K_{2}\rightarrow Rf'_{\ast}K_{2}{\quoi}.
\end{align*}
The\pageoriginale following diagram is commutative:
$$\etiquette[4.2.13.1]{XVII.eq4.2.13.1}
\vcenter{
\xymatrix{
L \overset{\LL}{\otimes} \sisi{L}{\LL}g^{\ast}Rf_{\ast}K
\ar[r]^{\varphi} & Rf'_{\ast}(f'^{\ast}L
\overset{\LL}{\otimes}\sisi{L}{\LL} g'^{\ast}K) \ar[d]^{Rf'_{\ast}(v')}&\\
L \overset{\LL}{\otimes} \sisi{L}{\LL}g^{\ast}K_1 \ar[u]_{L \overset{\LL}{\otimes} \sisi{L}{\LL}g^{\ast}(u')}
\ar[r]^{c'_{uv}} & Rf'_{\ast}K_2.
}}\leqno{(4.2.13.1)}
$$
\end{enonce*}

Suppose first that $K_{1}=f_{\ast}K$ and that
$K_{2}={f'}^{\ast}L\otimes {g'}^{\ast}K$, and consider
the following commutative diagram of complexes:
\[
\xymatrix{
K^{a}  \ar[r] \ar[d]& K\ar[d]\\
K^{c}  \ar[r] & K^{b}
}
\]
such that $K^{a}$ is adapted to
$\sisi{L}{}{f'}^{\ast}L\overset{\LL}{\otimes}\sisi{L}{\LL}{g'}^{\ast}$,
$K^{b}$ is adapted to $Rf_{\ast}$, and $K^{c}$ is adapted
to both functors. To construct this diagram, we first
construct $K^{a}$, then apply a canonical flasque resolution
(truncated if necessary) to $K^{a}$ and $K$.

Apply the functor $f_{\ast}$ (\resp~$L\otimes {g'}^{\ast}$) to this
diagram, and resolve the complexes in the diagram
$(K_{1},K^{a}_{1},K^{b}_{1}, K^{c}_{1})$ (\resp $(K_{2}, K^{a}_{2},
K^{b}_{2}, K^{c}_{2})$) so as to obtain a commutative cube, in which
$\approx$ denotes a quasi-isomorphism:
\[
\xymatrix@C=2em{
          &         &K'_1 \ar[rdd] \ar[rrrrrr]^{\approx}&        && &         &
          & K_1\ar[rdd]\\  K_1^{a'}\ar[rru]\ar[rdd]\ar[rrrrrr]^(.75){\approx} &         &      &        && & K_1^{a}\ar[rru] \ar[rdd]&       &\\
  &         &      &K^{b'}_1 \ar[rrrrrr] ^(.25){\approx}&& &         &         &   & K^{b}_1\\
          &K_1^{c'} \ar[rru]^{\approx}\ar[rrrrrr] ^(.65){\approx}&      &        && &         &K_1^{c}\ar[rru]^{\approx} \\
}
\]
where\pageoriginale $K'_{1}$, $K^{a'}_{1}$, $K^{b'}_{1}$, and $K^{c'}_{1}$
are adapted to the functor
$L\overset{\LL}{\otimes}\sisi{L}{\LL}g^{\ast}$ (\resp
\[
\xymatrix@C=2em{
        &         &K_2 \ar[rdd] \ar[rrrrrr]^{\approx}&        && &         &
          & K'_2\ar[rdd]\\
K_2^{a}\ar[rru]\ar[rdd]^{\approx}\ar[rrrrrr]^(.75){\approx} &         &      &        && & K_2^{a'}\ar[rru] \ar[rdd]^(.65){\approx}&       &\\
  &         &      &K^{b}_2 \ar[rrrrrr] ^(.25){\approx}&& &         &         &   & K^{b'}_2\\
          &K_2^{c} \ar[rru]\ar[rrrrrr] ^(.65){\approx}&      &
        && &         &K_2^{c'}\ar[rru] & &, \\
}
\]
where $K'_{2}$, $K^{a'}_{2}$, $K^{b'}_{2}$, and $K^{c'}_{2}$ are adapted to
the functor $Rf'_{\ast}$). Using the usual
base-change morphisms, one obtains a commutative prism of complexes over
$S$.
%%%%%%
\newcommand{\circlenum}[1]{\bigcirc \kern -.255cm \lower
  .05em\hbox{${#1}$}}

\[
\def\objectstyle{\scriptstyle}
\xymatrix@C=-.27cm{
        &         & \scriptstyle{L\otimes g^{\ast}K'_1}\ar[rrrrrr]_*+[o][F-]{\scriptscriptstyle{4}} \ar[rdd]^(.35)*+[o][F-]{\scriptscriptstyle{2}}&        && &         &
          & \scriptstyle{L\otimes g^{\ast}K_1} \ar[rdd] \ar[rrrrr]_*+[o][F-]{\scriptscriptstyle{4}}& & & & & \scriptstyle{f'_{\ast}K_2} \ar[rrrrr]_*+[o][F-]{\scriptscriptstyle{4}}\ar[rdd] & & & && \scriptscriptstyle{f'_{\ast}K'_2} \ar[rdd]\\
\scriptstyle{L\otimes g^{\ast}K_1^{a'}} \ar[rru]\ar[rdd]\ar[rrrrrr] &         &
        &        && & \scriptstyle{L\otimes g^{\ast}K_1\sisi{}{^a}} \ar[rru] \ar[rdd]\ar[rrrrr]&       &
        &&& \scriptstyle{f'_{\ast}K^{a}_2}\ar[rdd]\ar[rru] \ar[rrrrr]&
        & & & &
        \scriptstyle{f'_{\ast}K^{a'}_2}\ar[rru]_*+[o][F-]{\scriptscriptstyle{3}}\ar[rdd]^(.7)*+[o][F-]{\scriptscriptstyle{3}}^{\quad \raisebox{1cm}{}{\scriptscriptstyle{-1}}}_(.7){\text{\rotatebox{110}{$\approx$}}}\\
  &         &      & \scriptstyle{L\otimes g^{\ast}K^{\sisi{c}{b}'}_1}  \ar[rrrrrr]
        && &         &         &   & \scriptstyle{L\otimes g^{\ast}K^{b}_2}
        \ar[rrrrr]& & & & & \scriptstyle{f'_{\ast}K^{b}_2} \ar[rrrrr]& & & & & \scriptstyle{f'_{\ast}K^{b'}_2}\\
          &\scriptstyle{L\otimes g^{\ast}K^{c'}_1}
 \ar[rrrrrr]_*+[o][F-]{\scriptscriptstyle{1}}\ar[rru]_{\text{\rotatebox{50}{$\approx$}}}^*+[o][F-]{\scriptscriptstyle{2}}^{\qquad\qquad \raisebox{-.3cm}{}{\scriptscriptstyle{-1}}}
&      &        && &         &\scriptstyle{L\otimes g^{\ast}K^{c}_2} \ar[rru] \ar[rrrrr]_*+[o][F-]{\scriptscriptstyle{1}}& & & & &\scriptstyle{f'_{\ast}K^{c}_2}\ar[rru]\ar[rrrrr]_*+[o][F-]{\scriptscriptstyle{1}}&
        & & & & \scriptstyle{f'_{\ast}K^{c'}_{2'}}\ar[rru]\\
}
\]
and\pageoriginale one concludes by noting that the morphisms in diagram
\eqhyperref[XVII.eq4.2.13.1]{4.2.13.1} are identified with the following composites:

for $\varphi$: the composite of the morphisms \ding{192};

for $L\overset{\LL}{\otimes}\sisi{L}{\LL}g^{\ast}(u')$: the
composite of the morphisms \ding{193}, one of which is the inverse of a
quasi-isomorphism;

for $Rf'_{\ast}(v')$: the composite of the morphisms \ding{194}, one of which is
the inverse of a quasi-isomorphism;

for $\sisi{C}{c}'_{uv}$: the composite of the morphisms \ding{195}.

In the general case, set $K'_{1}=f_{\ast}K$ and
$K'_{2}={f'}^{\ast}L\otimes g^{\ast}K$. We then have a
diagram
\[
\xymatrix{
&L \overset{\LL}{\otimes} \sisi{L}{\LL}g^{\ast}Rf_{\ast}K
\ar[r]^\varphi&
Rf'_{\ast}(f'^{\ast}L \overset{\LL}{\otimes}\sisi{L}{\LL}
g'^{\ast}K) \ar[d] \ar[rd]& \\
L \overset{\LL}{\otimes} \sisi{L}{\LL}g^{\ast}K_1 \ar[r]\ar[ru]& L
\overset{\LL}{\otimes} \sisi{L}{\LL}g^{\ast}K'_1 \ar[r] \ar[u]&
Rf'_{\ast}K'_2 \ar[r]&  Rf'_{\ast}K_2,
}
\]
whose triangles are commutative, as is the inner rectangle
by the preceding argument. The boundary is therefore commutative, and this proves
\hyperref[XVII.prop4.2.13]{4.2.13}.

\subsection{The base-change theorem}\etiquette[4.3]{XVII.subsec4.3}\pageoriginale

We shall reformulate in terms of derived categories the
base-change theorem for proper morphisms \hyperrefext[XII][XII.thm5.1]{5.1}
(iii).

\begin{enonce*}{Theorem 4.3.1}\etiquette[4.3.1]{XVII.thm4.3.1}
Let
\[
\xymatrix{
X' \ar[rr]^{g'} \ar[d]^{f'}&& X\ar[d]^{f}\\
S' \ar[rr]^{g}& & S
}
\]
be a Cartesian diagram of schemes, with $f$ proper. Let $\Aa$ and $\Aa'$ be sheaves
of torsion rings on $S$ and $S'$, and let $L$ be a bounded-above complex
of $(\Aa',g^{\ast}\Aa)$-bimodules. Assume that
the dimensions of the fibers of $f$ are bounded. Hypothesis
\hyperref[XVII.subsubsec4.2.12]{4.2.12} {\rm(a)} is then satisfied, and the base-change morphism
\eqhyperref[XVII.eq4.2.12.3]{4.2.12.3}, between functors from $D^{-}(X,f^{\ast}\Aa)$
to $D^{-}(S',\Aa')$,
\[
\varphi_K:L\overset{\LL}{\otimes}\sisi{L}{\LL}g^{\ast}\sisi{}{Rf_*}K\longrightarrow
Rf'_{\ast}({f'}^{\ast}L\overset{\LL}{\otimes}\sisi{L}{\LL}{g'}^{\ast}K){\quoi},
\]
is an isomorphism.
\end{enonce*}

The functors under consideration are triangulated and ``way-out'' \cite[I 7]{sga4.exp.XVII.bib.RD},
so, to verify that $\varphi_{K}$ is an isomorphism, it
suffices to verify that the $\varphi_{ℋ^{i}(K)}$ are
isomorphisms. We may therefore assume that $K$ is a sheaf
of torsion $f^{\ast}\Aa$-modules concentrated in degree $0$.

To\pageoriginale verify that $\varphi_{K}$ is an isomorphism,
it suffices to do so after applying the forgetful functor from the category of
$\Aa'$-modules to that of
$\ZZ/n\ZZ$-modules, where $n\Aa=n\Aa'=0$.

This reduces us to the case
$\Aa'=\ZZ/n\ZZ$, with $n\Aa=0$. We may
then replace $L$ by an isomorphic complex in the derived
category, bounded above and with components flat over $g^{\ast}\Aa$.
Since the functors under consideration are triangulated and ``way-
out'' in $L$, it suffices to verify that $\varphi_{K}$ is an
isomorphism after replacing $L$ by one of its components. We
may therefore assume that $L$ is a sheaf of flat
$g^{\ast}\Aa$-modules placed in degree $0$. We must
prove that the morphisms
\[
\varphi^{q}:L\otimes g^{\ast}R^{q}f_{\ast}K\longrightarrow
R^{q}f'_{\ast}({f'}^{\ast}L\otimes g^{\ast}K)
\]
are isomorphisms. Let $s$ be a geometric point of $S'$, let $X_s$
be the fiber of $X'$ at $s$, and let $K_s$ be the inverse image of $K$ on
$X_s$. By \hyperrefext[XII][XII.thm5.1]{5.1} (iii), the fiber at $s$ of $\varphi^{q}$ is
the morphism
\[
\varphi^{q}_s:L_{s}\otimes H^{q}(X_s,K)\longrightarrow H^{q}(X_s,
L_s\otimes K){\quoi}.
\]
By $D$. Lazard's theorem \cite{sga4.exp.XVII.bib.La}, the flat $\Aa_s$-module
$L_s$ is the filtered direct limit of free $\Aa_s$-modules
$L_{s,i}$; the functor $H^{q}(X_s,\ast)$ commutes with filtered
direct limits\sisi{}{ (\hyperrefext[VII][VII.prop3.3]{3.3})}. The morphism $\varphi^{q}$ is therefore the direct
limit of the morphisms
\[
\varphi^{q}_i:L_{s,i}\otimes H^{q}(X_s,K)\longrightarrow H^{q}(X_s,
L_{s,i}\otimes K){\quoi}.
\]
The latter are plainly isomorphisms; hence so is
$\varphi^{q}_s$, and $\varphi\sisi{}{^q}$ is an isomorphism.

\begin{enonce*}[remark]{Variants 4.3.2}\etiquette[4.3.2]{XVII.var4.3.2}
\begin{enumerate}
\renewcommand{\theenumi}{\roman{enumi}}
\renewcommand{\labelenumi}{(\theenumi)}
\item When\pageoriginale $L$ is a bounded complex whose components
have finite \tTor-dimension over $g^{\ast}\Aa$, hypotheses (c) of
\hyperref[XVII.subsubsec4.2.12]{4.2.12} are satisfied, and one may work in
the unbounded derived category.\nde{One may work in
the unbounded derived category without these hypotheses. \Cf \nderef{XVII.ndeKS}.}

\item If $L$ is bounded below, with components whose \tTor-dimension is \sisi{finite}{uniformly bounded}
over $g^{\ast}\Aa$, hypotheses (b) of \hyperref[XVII.subsubsec4.2.12]{4.2.12} are satisfied,
even if $\Aa$ and $\Aa'$ are not torsion; the
base-change morphism $\varphi_{K}$ $(K\in
D^{+}(X,f^{\ast}\Aa)$) will be an isomorphism if the sheaves
$ℋ^{i}(K)$ are torsion.
\end{enumerate}
\end{enonce*}

\section{The direct-image functors with proper supports}\etiquette[5]{XVII.sec5}

\subsection{The fundamental construction}\etiquette[5.1]{XVII.subsec5.1}

We shall follow the program set out in \hyperref[XVII.subsec3.1]{3.1}
and define the
functor $Rf_{!}$ for a compactifiable morphism $f$
(\hyperref[XVII.subsec3.2]{3.2}). In the appendix we shall
show how to treat the case of an arbitrary separated morphism of
finite type.\nde{\Cf \nderef{XVII.ndeConrad}.}

\subsubsection{}\etiquette[5.1.1]{XVII.subsubsec5.1.1}

Let $X$ be a topos, $U$ an open subtopos of $X$ (\hyperrefext[IV][IV.def8.3]{8.3}), and $j$ the
inclusion morphism from $U$ to $X$. In \hyperrefext[IV][IV.subsec5.2]{5.2} we defined an ``\emph{extension by the empty set}'' functor from the category of sheaves
of sets on $U$ to that of sheaves of sets on $X$. This
functor $j_{!}$ is left adjoint to the restriction functor
$j^{\ast}$.

We also denote by $j_{!}$ the functor from the category of
sheaves of pointed sets on $U$ (\resp sheaves of groups,
\resp sheaves of $j^{\ast}\Aa$-modules, for
$\Aa$ a sheaf of rings on $X$) to\pageoriginale the
analogous category on $X$, left adjoint to the restriction
functor $j^{\ast}$ (\sisi{IV\quad}{\hyperrefext[IV][IV.prop11.3.1]{11.3.1}}). These three functors $j_{!}$ are
compatible with the forgetful functors from the category of modules to
that of groups, and from the category of groups to that of
pointed sets \footnote[1]{But not to that of sets!
}. If $i:F\rightarrow X$ is the inclusion of the closed complement of
$U$, then $j^{\ast}j_{!}$ is (canonically isomorphic to) the identity
and $i^{\ast}j_{!}$ is the constant functor with value the initial object
(or the final object, which here amounts to the same thing). This also characterizes
$j_{!}$. It follows that the three functors $j_{!}$ under consideration are
exact and embed into the corresponding functors $j_{\ast}$. They are
called the \emph{extension by zero} functors. The formula
$(j_1j_2)^{\ast}=j^{\ast}_2 j^{\ast}_1$ transposes into a
transitivity isomorphism
\[
(j_1 j_2)_{!}=j_{1!}j_{2!}{\quoi}.
\]
The functors $j_{!}$ are compatible with base change:

\begin{enonce*}{Lemma 5.1.2}\etiquette[5.1.2]{XVII.lem5.1.2}
Let $f:X\rightarrow Y$ be a morphism of topoi, let $V$ be an open subtopos of $Y$,
and let $U$ be its inverse image in $X$:
\[
\xymatrix{
U \ar[rr]^{f'}  \ar[d]^{j'}& & V\ar[d]^{j}\\
X \ar[rr]^{f}   && Y
}
\]
There is then a unique base-change isomorphism
$f^{\ast}j_{!}\isomto j'_{!}{f'}^{\ast}$ making the
diagram
$$\etiquette[5.1.2.1]{XVII.eq5.1.2.1}
\vcenter{
\xymatrix{
f^{\ast}j_{!}    \ar[r]^{\sim} \ar@{^{(}->}[d] & j'_{!}f'^{\ast} \ar@{^{(}->}[d]     \\
f^{\ast}j_{\ast}  \ar[r]^{\alpha}&j'_{\ast}f'^{\ast},
}}\leqno{(5.1.2.1)}
$$
commutative, where $\alpha$ is the base-change morphism.
\end{enonce*}

\begin{enonce*}[remark]{Proof}
The\pageoriginale base-change morphism
$\varphi:j^{\ast}f_{\ast}\rightarrow f'_{\ast}{j'}^{\ast}$ is
trivially an isomorphism. It transposes into a base-change isomorphism
$\varphi':j'_{!}{f'}^{\ast}\isomto
f^{\ast}j_{!}$. To verify that diagram \eqhyperref[XVII.eq5.1.2.1]{5.1.2.1} is
commutative, it suffices to see this after restriction to $U$, where it
is trivial.
\end{enonce*}

\setcounter{subsubsection}{2}
\subsubsection{}\etiquette[5.1.3]{XVII.subsubsec5.1.3}
Under the hypotheses of \hyperref[XVII.subsubsec5.1.1]{5.1.1}, with $X$ ringed by a sheaf
of rings $\Aa$, the functor $j_{!}$ on modules is
exact, hence passes trivially to the derived categories and defines a
functor, still denoted $j_{!}$, from $D(U,j^{\ast}\Aa)$ to
$D(X,\Aa)$. We still have transitivity isomorphisms
$$\etiquette[5.1.3.1]{XVII.eq5.1.3.1}
(j_1 j_2)\sisi{}{!}=j_{1!}j_{2!}{\quoi}.\leqno(5.1.3.1)
$$

At the level of derived categories, we still have an
adjunction formula
$$\etiquette[5.1.3.2]{XVII.eq5.1.3.2}
\Hom_{D(X)}(j_{!}K,L)\isomto \Hom_{D(U)}
(K,j^{\ast}L){\quoi}.\leqno(5.1.3.2)
$$

\subsubsection{}\etiquette[5.1.4]{XVII.subsubsec5.1.4}
Let $f:X\rightarrow Y$ be a proper morphism, let $\Aa$ be a sheaf of
\emph{torsion} rings on $Y$, and endow $X$ with
the inverse-image sheaf of rings of $\Aa$. Assume that the
dimensions of the fibers of $f$ are bounded above by an integer $n$ (as is
the case if $Y$ is quasi-compact). We then know that, for every
sheaf of $f^{\ast}\Aa$-modules $F$ on $X$, the sheaves
$R^{q}f_{\ast}F$ vanish for $q>2n$ (\hyperrefext[XII][XII.corbis5.3]{5.3 bis}). The functor
$f_{\ast}$ therefore admits a derived functor (\hyperref[XVII.prop1.2.10]{1.2.10})
\[
Rf_{\ast}:D(X,f^{\ast}\Aa)\longrightarrow
D(Y,\Aa){\quoi}.
\]

% END INCLUDED SGA4_English_current_buildable_source_20260906/source/17/drafts/SGA4_XVII_LINES_3401_3800_English.texfrag

% BEGIN INCLUDED SGA4_English_current_buildable_source_20260906/source/17/drafts/SGA4_XVII_LINES_3801_4200_English.texfrag

If\pageoriginale $f$ is the composite $gh$ of two proper morphisms,
there is a transitivity formula
$Rf_{\ast}=Rg_{\ast}Rh_{\ast}$, which, as usual, can be expressed
in terms of cofibered categories (Scholium \sisi{\hyperref[XVII.subsubsec5.1.3]{5.1.3}}{\hyperref[XVII.sch4.1.3]{4.1.3}}).

\subsubsection{}\etiquette[5.1.5]{XVII.subsubsec5.1.5}
Consider a commutative diagram of schemes
$$\etiquette[5.1.5.1]{XVII.eq5.1.5.1}
\vcenter{
\xymatrix{
U        \ar@{^{(}->}[r]_{j_1} \ar[d]^{g}      & X  \ar[d]^{f}\\
V        \ar@{^{(}->}[r]_{j_2}       & Y,
}}\leqno{(5.1.5.1)}
$$
in which $j_1$ and $j_2$ are open immersions and $f$, $g$ are
proper morphisms whose fiber dimensions are bounded above by an
integer $n$. Let $\Aa$ be a sheaf of torsion rings on
$Y$, and equip $U$, $V$, and $X$ with the appropriate inverse images of
$\Aa$. By \eqhyperref[XVII.eq5.1.3.2]{5.1.3.2}, if $K\in D(U)$, defining a morphism
$$\etiquette[5.1.5.2]{XVII.eq5.1.5.2}
d:j_{2!}Rg_{\ast}K\longrightarrow Rf_{\ast}j_{1!}K\leqno{(5.1.5.2)}
$$
amounts to defining a morphism
\[
d':Rg_{\ast}K\longrightarrow j^{\ast}_2 Rf_{\ast}j_{1!}K{\quoi}.
\]

The fiber product $X'=V\times {}_YX$ fits into a commutative
diagram
$$\etiquette[5.1.5.3]{XVII.eq5.1.5.3}
\vcenter{
\xymatrix{
U \ar@{^{(}->}[r]^{k} \ar[d]^{g} & X'  \ar@{^{(}->}[r]_{j'_2} \ar[d]^{f'}& X \ar[d]^{f}          &\\
V         \ar@{=}[r] & V         \ar@{^{(}->}[r]_{j_2}  & Y   \rlap{,}
}}\leqno{(5.1.5.3)}
$$
and,\pageoriginale trivially, we have $j^{\ast}_2Rf_{\ast}j_{1!}K\simeq
Rf'_{\ast}(j'_2{}^\ast j_{1!}K)\simeq Rf'_{\ast}k_{!}K$.

The morphism $k$ is a proper open immersion; hence
$k_{!}=k_{\ast}$, giving an isomorphism $Rf'_{\ast}k_{!}K\simeq
Rf'_{\ast}k_{\ast}K\simeq Rg_{\ast}K$, which provides the desired isomorphism
$d'$ and the morphism $d$. Moreover, the restriction
of $d$ to $V$ is an isomorphism.

\begin{enonce*}{Lemma 5.1.6}\etiquette[5.1.6]{XVII.lem5.1.6}
The morphism \eqhyperref[XVII.eq5.1.5.2]{5.1.5.2} is an isomorphism.
\end{enonce*}

It remains to verify that the restriction of $d$ to $Y-V=F$ is an
isomorphism. Let $X''=\fprod{X}{F}{Y}$ and apply the proper
base-change theorem (\hyperref[XVII.thm4.3.1]{4.3.1}) to the Cartesian square
\[
\xymatrix{
X'' \ar[r]_{i'} \ar[d]^{f''}& X\ar[d]^{f\sisi{'}{}}\\
F   \ar[r]_{i}& Y\rlap{,}
}
\]
with $\Aa'=i^{\ast}\Aa$ and $L=\Aa'[0]$. We
obtain:
$i^{\ast}Rf_{\ast}j_{1!}K\simeq Rf''_{\ast}{i'}^{\ast}j_{1!}K\simeq
Rf''_{\ast}0\simeq 0$, so that both sides of \eqhyperref[XVII.eq5.1.5.2]{5.1.5.2} are
zero outside $V$ and $d$ is an isomorphism.

\setcounter{subsubsection}{6}
\subsubsection{}\etiquette[5.1.7]{XVII.subsubsec5.1.7}
Let $S$ be a scheme ringed by a sheaf of torsion rings
$\Aa$. For every $S$-scheme $\varphi:X\rightarrow S$, set
$\Aa_X=\varphi^{\ast}\Aa$. Work in the category $(S)$ of
quasi-compact, quasi-separated $S$-schemes and
$S$-compactifiable morphisms (\hyperref[XVII.subsec3.2]{3.2}). By
\hyperref[XVII.prop3.2.3]{3.2.3}, this category, the subcategory $(S,i)$ of open
immersions, and the subcategory $(S,p)$ of proper morphisms
satisfy the hypotheses of \hyperref[XVII.subsubsec3.2.4]{3.2.4}.

We\pageoriginale propose to apply the theory of \hyperref[XVII.prop3.3.2]{3.3.2} to the following data
(notation of \hyperref[XVII.subsubsec3.3.1]{3.3.1})

\begin{itemize}
\item[$\alpha$)] $F(X)=D(X,\Aa_X)$ for $\varphi:X\rightarrow
S$ in $S$:
\begin{itemize}
\item[(i)] For $p:X\rightarrow Y$ proper, take the functor
$p_{\ast}$ to be the functor $Rp_{\ast}$.

\item[(ii)] Use the usual transitivity isomorphisms
$R(pq)_{\ast}=Rp_{\ast}Rq_{\ast}$.

\item[${\rm(i')}$] For an open immersion $i:X\rightarrow Y$, take
the functor $i_{\ast}$ to be the functor $i_{!}$.

\item[${\rm(ii')}$] Use the transitivity isomorphisms
\eqhyperref[XVII.eq5.1.3.1]{5.1.3.1}.

\item[${\rm(iii)}$] For every commutative diagram \eqhyperref[XVII.eq5.1.5.1]{5.1.5.1}, take
the isomorphism \eqhyperref[XVII.eq5.1.5.2]{5.1.5.2} for $d$.
\end{itemize}
\end{itemize}

The conditions (a) (b) (c) ${\rm(a')}$ ${\rm(b')}$
${\rm(c')}$ of \hyperref[XVII.prop3.3.2]{3.3.2} are verified, proving part (a) of:

\begin{enonce*}{Theorem 5.1.8}\etiquette[5.1.8]{XVII.thm5.1.8}
\begin{enumerate}
\renewcommand{\theenumi}{\alph{enumi}}
\renewcommand{\labelenumi}{(\theenumi)}
\item In the category $(S)$ defined in \hyperref[XVII.subsubsec5.1.7]{5.1.7}, there exists one, and
essentially only one, way to define, for every arrow
$f:X\rightarrow Y$, a functor
\[
Rf_{!}:D(X,\Aa_X)\longrightarrow D(Y,\Aa_Y){\quoi},
\]
and, for every composite morphism
$f=gh:X\xrightarrow{h}Y\xrightarrow{g}Z$, transitivity
isomorphisms $c_{\sisi{f,g}{g,h}}$ between $Rf_{!}$ and $Rg_{!}$ $Rh_{!}$ such
that
\begin{enumerate}
\renewcommand{\theenumii}{\roman{enumii}}
\renewcommand{\labelenumii}{(\theenumii)}
\item The $c_{f,g}$ satisfy the ``cocycle condition''
\[
(c_{f,g}\ast Rh_{!})\circ c_{fg,h}=(Rf_{!}\ast c_{g,h})\circ
c_{f,gh}{\quoi}.
\]

\item When restricted to proper morphisms (\resp to
open immersions), one recovers the variance theory of
\hyperref[XVII.subsubsec5.1.4]{5.1.4} (\resp~\hyperref[XVII.subsubsec5.1.3]{5.1.3}).

\item For\pageoriginale every diagram \eqhyperref[XVII.eq5.1.5.1]{5.1.5.1}, the isomorphism
$d$ of \eqhyperref[XVII.eq5.1.5.2]{5.1.5.2} is the composite isomorphism $c_{f,j\sisi{'}{_1}} \circ c^{-1}_{j\sisi{,f'}{_2,g}}$.
\end{enumerate}

\item The functors $Rf_{!}$ are ``way out'' \emph{\cite[I 7]{sga4.exp.XVII.bib.RD}} and
triangulated.
\end{enumerate}
\end{enonce*}

Every functor $Rf_{!}$ is the composite of a functor $Rg_{\ast}$, for $g$
proper with quasi-compact target, hence with fibers of bounded dimension, and
of a functor $j_{!}$ for $j$ an open immersion. These functors are
way out and triangulated, and hence so is $Rf_{!}$.

\setcounter{subsubsection}{8}
\paragraph{}\etiquette[5.1.8.1]{XVII.par5.1.8.1}

The transitivity isomorphisms of \hyperref[XVII.thm5.1.8]{5.1.8} (a) are often used via
the ``spectral sequence of composite functors'' deduced from them,
valid for $K\in \Ob D(X,\Aa_X)$:
$$\etiquette[5.1.8.2]{XVII.eq5.1.8.2}
E^{pq}_2=R^{p}g_! R^{q}h_{!}(K)\Longrightarrow
R^{p+q}f_{!}(K){\quoi}.\leqno{(5.1.8.2)}
$$

To obtain it, note that by \hyperref[XVII.thm5.1.8]{5.1.8} b) and \textsc{Verdier}
\cite{sga4.exp.XVII.bib.V2}, for every
$L\in \Ob D(Y,\Aa_Y)$, there is a spectral sequence
$$\etiquette[5.1.8.3]{XVII.eq5.1.8.3}
E^{pq}_2=R^{p} g_!(ℋ^{q}(L))\Longrightarrow
R^{p+q}g_!(L){\quoi},\leqno{(5.1.8.3)}
$$
and \eqhyperref[XVII.eq5.1.8.2]{5.1.8.2} is the special case of \eqhyperref[XVII.eq5.1.8.3]{5.1.8.3} with $L=Rh_{!}K$.

\begin{enonce*}{Definition 5.1.9}\etiquette[5.1.9]{XVII.def5.1.9}
\begin{enumerate}
\renewcommand{\theenumi}{\roman{enumi}}
\renewcommand{\labelenumi}{(\theenumi)}
\item We call the ``\emph{
\sisi{total}{total} direct-image functors with proper supports}'' the functors constructed in \hyperref[XVII.thm5.1.8]{5.1.8}. We shall omit
the adjective ``total'' when there is no risk of
confusion.

\item The $q$\textsuperscript{th} \emph{direct-image functor with proper
support,} denoted $R^{q}f_{!}$, is the composite functor
$\Hh^{q}\circ Rf_{!}$.

\item If $X$ is a scheme over a separably closed field $k$ that is
embeddable in a proper scheme over $k$, identify sheaves of
modules on\pageoriginale $\Spec (k)$ with ordinary modules, and denote
the functor $R^{q}f_{!}$ by one of the notations $H^{q}_c(X,\quad)$,
$H^{q}_{!}(X,\quad)$, or simply $H^{q}_c(\quad)$ if there is no risk
of confusion. It is then called the ``$q$\textsuperscript{th} \emph{hypercohomology
group with proper supports},'' or the ``$q$\textsuperscript{th} \emph{cohomology group with proper supports}'' if it
is applied to a complex concentrated in degree $0$.
\end{enumerate}
\end{enonce*}

\setcounter{subsubsection}{9}
\setcounter{paragraph}{0}
\paragraph{}\etiquette[5.1.9.1]{XVII.par5.1.9.1}
The notation $Rf_{!}$ is abusive in that \emph{the functor $Rf_{!}$
is not the right derived functor of the functor $f_{!}$} that will be defined in
\hyperref[XVII.def6.1.2]{6.1.2}. Some prefer the notation $R_{!}f$.

\subsubsection{}\etiquette[5.1.10]{XVII.subsubsec5.1.10}
The ingredients in the proof of \hyperref[XVII.thm5.1.8]{5.1.8} are the
proper
base-change theorem and the identity between $j_{\ast}$ and $j_{!}$
when $j$ is a proper open immersion. This makes it possible to state
various variants of \hyperref[XVII.thm5.1.8]{5.1.8}; it is left to the reader
to give them a precise meaning and to verify them.

\begin{enonce*}{Variant 5.1.11}\etiquette[5.1.11]{XVII.var5.1.11}
In the category $(S)$, there exists one, and essentially only one,
way to define, for every arrow $f:X\rightarrow Y$, a functor
$f_{!}$ from the category of sheaves of pointed sets on $X$ to
the category of sheaves of pointed sets on $Y$, and, for every
composite morphism $f=gh$, transitivity isomorphisms satisfying
conditions {\rm (i)\sisi{,}{} (ii) (iii)} analogous to those of \sisi{\hyperref[XVII.subsubsec5.1.4]{5.1.4}}{\hyperref[XVII.thm5.1.8]{5.1.8}}.
\end{enonce*}

See \hyperref[XVII.def6.1.2]{6.1.2} for a direct description of this functor.

\begin{enonce*}{Variant 5.1.12}\etiquette[5.1.12]{XVII.var5.1.12}
For\pageoriginale every compactifiable morphism
$f:X\rightarrow Y$, one can define a functor $R^{1}f_{!}$ from the category
of sheaves of \iInd-finite groups on $X$ to the category of
sheaves of pointed sets on $Y$, together with, for every
compactification $f=\bar{f}j$ of $f$, an isomorphism
$R^{1}f_{!}\simeq R^{1}\bar{f}_{\ast}\circ j_{!}$.
\end{enonce*}

No transitivity formula is postulated here; the only problem is
to show that the functor $R^{1}f_{!}$ is independent of the chosen
compactification of $f$. Let $p:(\bar{f}_1, i_{1})\rightarrow
(\bar{f}_{2}, i_{2})$ be a morphism of compactifications, \ie a
diagram
$$\etiquette[5.1.12.1]{XVII.eq5.1.12.1}
\vcenter{
\xymatrix{
X   \ar[rr]^{i_1} \ar[rrd]^{i_2} \ar[dd]_{f}& & \bar{X}_1 \ar[d]^{p}\\
    & &  \bar{X}_2 \ar[lld]^{\bar{f}_2}\\
Y   & & .
}}\leqno{(5.1.12.1)}
$$
where $\bar{f}_{1}=\bar{f}_2\circ p$, with $p$ proper. If $G$ is a
sheaf \sisi{of}{of} groups on $X$, there is an exact sequence of pointed
sets (with $\alpha$ injective)
\[
0\rightarrow
R^{1}\sisi{f}{\bar{f}}_{2\ast}(p_{\ast}i_{1!}G)\xrightarrow{\alpha}R^{1}\sisi{f}{\bar{f}}_{1\ast}(i_{1!}G)\rightarrow
\bar{f}_{2\ast}R^{1}p_{\ast}(i_{1!}G){\quoi},
\]
and a morphism from $i_{2!}G$ to $p_{\ast}i_{1!}G$, whence, by
composition, a morphism $\beta$ from $R^{1}\sisi{f}{\bar{f}}_{2\ast}i_{2!}G$ to
$R^{1}\sisi{f}{\bar{f}}_{1\ast}i_{1!}G$. Reasoning as in \hyperref[XVII.par3.3.2.1]{3.3.2.1}, one sees that the
problem is to prove that this morphism is always an isomorphism
when $G$ is \iInd-finite.

Applying\pageoriginale the proper base-change theorem
to $p$ and the sheaf of sets $i_{1!}G$, one checks that
$p_{\ast}i_{1!}G=i_{2!}G$, yielding the exact sequence
\[
0\rightarrow R^{1}
\sisi{f}{\bar{f}}_{2\ast}(i_{2!}G)\xrightarrow{\beta}R^{1}\sisi{f}{\bar{f}}_{1\ast}(i_{1!}G)\rightarrow
\bar{f}_{2\ast}R^{1}p_{\ast}i_{1!}G{\quoi}.
\]
The morphism $\beta$ is injective (since $\alpha$ is). Applying
the proper base-change theorem to $p$ and the sheaf of
groups $G$, assumed \iInd-finite, one finds that
$R^{1}p_{\ast}i_{1!}G=0$. Thus $\beta$ is an epimorphism, which
completes the proof of \hyperref[XVII.var5.1.12]{5.1.12}.

\begin{enonce*}{Definition 5.1.13}\etiquette[5.1.13]{XVII.def5.1.13}
For every site $X$, denote by $D_{\tors}(X)$ the full
subcategory of the derived category of the category of abelian sheaves
on $X$ consisting of the complexes whose cohomology sheaves are
torsion.
\end{enonce*}

\begin{enonce*}{Variant 5.1.14}\etiquette[5.1.14]{XVII.var5.1.14}
In the category $(S)$, there exists one, and essentially only one,
way to define, for every arrow $f:X\rightarrow Y$, a
triangulated and ``way out'' functor \emph{\cite[I 7]{sga4.exp.XVII.bib.RD}}
\[
Rf_{!}:D^{+}_{\tors}(X)\longrightarrow D^{+}_{\tors}(Y)
\]
and, for every composite morphism $f=gh$, transitivity isomorphisms
satisfying conditions \mbox{\rm (i) (ii) (iii)}
analogous to those of \hyperref[XVII.subsubsec5.1.4]{5.1.4}.
\end{enonce*}

\begin{enonce*}{Variant 5.1.\sisi{}{1}5}\etiquette[5.1.15]{XVII.var5.1.15}
Work in the category of $S$-schemes ringed by
sheaves of torsion rings, the arrows being morphisms of ringed schemes
inducing (\hyperref[XVII.subsec0.8]{0.8}) an $S$-compactifiable morphism of
schemes. For every arrow
\[
f:(X,\Aa)\longrightarrow (Y,\Bb),
\]
let\pageoriginale $Rf_{!}$ be the composite of the forgetful functor from the category of
$\Aa$-modules to that of $f^{\sisi{\ast}{-1}}\Bb$-modules, followed
by the usual $Rf_{!}$. These functors have the variance properties of the
usual functors $Rf_{!}$.
\end{enonce*}

Note, however, that the notation $Rf_{!}$ is reasonable here only
if $\Aa$ is finite over $f^{\sisi{\ast}{-1}}\Bb$, as one
sees by taking $X=Y=\Spec(k)$, with $k$ algebraically closed.

\setcounter{subsubsection}{15}
\subsubsection{}\etiquette[5.1.16]{XVII.subsubsec5.1.16}
Let $f:X\rightarrow S$ be a compactifiable morphism, let $j:U\rightarrow
X$ be an open immersion into $X$, let $i:F\rightarrow X$ be the closed complement,
let $\Aa_S$ be a sheaf of torsion rings on $S$,
let $\Aa_X$ be its inverse image on $X$, and let $K\in \Ob
D(X,\Aa_X)$.

The exact sequence of complexes
\[
0\longrightarrow j_{!}j^{\ast}K\longrightarrow K\longrightarrow
i_{!}i^{\ast}K\longrightarrow 0
\]
gives rise to a distinguished triangle in $D(X,\Aa_X)$ by
\textsc{Verdier} \cite[\sisi{I 2.4}{II 1.1.5}]{sga4.exp.XVII.bib.CD}. Its image under $Rf_{!}$ is a distinguished triangle
(\hyperref[XVII.thm5.1.8]{5.1.8} (b))
$$\etiquette[5.1.16.1]{XVII.eq5.1.16.1}
R(fj)_{!}j^{\ast}K\longrightarrow Rf_{!}K\longrightarrow
R(fi)_{!}i^{\ast}K\longrightarrow R(fj)_{!}j^{\ast}K[1]\ldots\leqno{(5.1.16.1)}
$$
which defines a long exact cohomology sequence
$$\etiquette[5.1.16.2]{XVII.eq5.1.16.2}
\ldots \sisi{}{\to}R^{q}(fj)_{!}j^{\ast}K\rightarrow
R^{q}f_{!}K\rightarrow R^{q}(fi)_{!}i^{\ast}K \xrightarrow{j} R^{q+1}(fj)_{!}j^{\ast}K\rightarrow\ldots.\leqno{(5.1.16.2)}
$$

In the special case where $S$ is the spectrum of a separably
closed field and $K$ is a sheaf $𝒢$ concentrated in degree $0$, the
sequence \eqhyperref[XVII.eq5.1.16.2]{5.1.16.2} becomes
$$\etiquette[5.1.16.3]{XVII.eq5.1.16.3}
\ldots H^{q}_c(U,𝒢)\rightarrow H^{q}_c (X,𝒢)\rightarrow
H^{q}_c(F,𝒢)\xrightarrow{j} H^{q+1}_c(U,𝒢)\ldots\leqno{(5.1.16.3)}
$$

\subsubsection{}\etiquette[5.1.17]{XVII.subsubsec5.1.17}

Consider\pageoriginale a diagram of $S$-compactifiable schemes
\[
\xymatrix{
X \ar[rr]^{u} \ar[rd]_{f}& & Y\ar[ld]^{g}\\
  & S &,
}
\]
with $u$ \emph{proper}. Equip $X$ and $Y$ with the inverse images
of a sheaf of torsion rings $\Aa_S$ on $S$. For $K\in
D(Y,\Aa_Y)$, the evident (adjunction) morphism $K\rightarrow
Ru_{\ast}u^{\ast}K$ defines, upon applying the functor $Rg_{!}$, a
morphism
$$\etiquette[5.1.17.1]{XVII.eq5.1.17.1}
u^{\ast}:Rg_{!}K\longrightarrow
Rg_{!}Ru_{\ast}u^{\ast}K=Rg_{!}Ru_{!}u^{\ast}K\sisi{\eqsim}{\simeq} Rf_{!}u^{\ast}K\leqno(5.1.17.1)
$$
(\emph{contravariance of cohomology with proper supports with respect to
proper morphisms}).

When $S$ is the spectrum of an algebraically closed field and $K$ is
the sheaf $𝒢$ placed in degree $0$, this morphism becomes
$$\etiquette[5.1.17.2]{XVII.eq5.1.17.2}
u^{\ast}:H^{q}_c(Y,𝒢)\longrightarrow H^{q}_c(X,𝒢)\leqno(5.1.17.2)
$$

\subsection{The base-change theorem}\etiquette[5.2]{XVII.subsec5.2}

\subsubsection{}\etiquette[5.2.1]{XVII.subsubsec5.2.1}

Let $f:X\rightarrow Y$ be a morphism of topoi, let $V$ be an open of $Y$,
let $U=f^{-1}(V)$ be its inverse image in $X$, let $\Aa$ be a sheaf
of rings on $Y$, let $\Aa'$ be a sheaf of rings on $X$, and let $L$
be a bounded-above complex of $(\Aa',
f^{\ast}\Aa)$-bimodules:
$$\etiquette[5.2.1.1]{XVII.eq5.2.1.1}
\vcenter{
\xymatrix{
U \ar[rr]^{f'} \ar@{^{(}->}[d]^{j'}& & V \ar@{^{(}->}[d]^{j}\\
X \ar[rr]_{f} & & Y.
}}\leqno{(5.2.1.1)}
$$

There\pageoriginale is then a base-change isomorphism
\eqhyperref[XVII.eq5.1.2.1]{5.1.2.1} between functors from $D(V,\Aa)$ to
$D(X,f^{\ast}\Aa):f^{\ast}j_{!}\isomto
j'_{!}{f'}^{\ast}$. Moreover, there is an evident isomorphism
between functors from $D^{-}(U,{f'}^{\ast}\Aa)$ to
$D^{-}(X,\Aa'):L\overset{\LL}{\otimes}_{f^{\ast}\Aa}j'_{!}K\simeq
j'_{!}({j'}^{\ast}L\overset{\LL}{\otimes}_{f^{\ast}\Aa}K)$.
To construct it, it suffices to note that if a complex $K$ is
adapted to the functor
${j'}^{\ast}L\overset{\LL}{\otimes}_{f^{\ast}\Aa}$,
then the complex $j'_{!}K$ is adapted to the functor
$L\overset{\LL}{\otimes}_{f^{\ast}\Aa}$,
and that there is an isomorphism of the preceding type at the level of
complexes. Composing these two isomorphisms gives a
base-change isomorphism between functors from
$D^{-}(V,\Aa)$ to $D^{-}(X,\Aa')$:
$$\etiquette[5.2.1.2]{XVII.eq5.2.1.2}
L\overset{\LL}{\otimes}_{f^{\ast}\Aa}f^{\ast}j_{!}K\isomto j'_{!}({j'}^{\ast}L\overset{\LL}{\otimes}_{{f\sisi{'}{}}^{\ast}\Aa}{f'}^{\ast}K){\quoi}.\leqno(5.2.1.2)
$$

If $L$ is bounded with components of finite \ttor-dimension over

% END INCLUDED SGA4_English_current_buildable_source_20260906/source/17/drafts/SGA4_XVII_LINES_3801_4200_English.texfrag

% BEGIN INCLUDED SGA4_English_current_buildable_source_20260906/source/17/drafts/SGA4_XVII_LINES_4201_4600_English.texfrag
$f^{\ast}\Aa$,\nde{this condition is superfluous. \Cf
\nderef{XVII.ndeKS}.} this isomorphism is still defined for
$K\in\Ob D(V,\Aa)$.

\subsubsection{}\etiquette[5.2.2]{XVII.subsubsec5.2.2}

Let $f:X\rightarrow S$ and $g:S'\rightarrow S$ be morphisms of
schemes, and suppose that $f$ is the composite of an open immersion
$j$ and a proper morphism $p$ whose fibers have dimension bounded by
an integer $n$ (this is automatic when $S$ is quasi-compact). These
morphisms fit into a diagram with cartesian squares
\[
\xymatrix{
X' \ar[rrrr]_{g'} \ar@{^{(}->}[rd]^{j'} \ar[dd]_{f'}& &         &   &
X\ar@{^{(}->}[rd]^{j} \ar[dd]_(.35){f} & \\
             &\bar{X}'\ar[rrrr]|(.75)\hole_{\bar{g}}\ar[ld]^{p'}  &&   &   & \bar{X}\ar[ld]^{p}\\
S' \ar[rrrr]_{g}  & & &&  S &.
}
\]
Let\pageoriginale $\Aa$ be a sheaf of torsion rings on $S$, $\Aa'$ a
sheaf of torsion rings on $S'$, and $L$ a bounded-above complex of
$(\Aa',g^{\ast}\Aa)$-bimodules. We ring $X$ and $\bar X$
(\resp $X'$ and $\bar X'$) by the inverse images $\Aa_X$ and
$\Aa_{\bar X}$ (\resp $\Aa'_{X'}$ and $\Aa'_{\bar X'}$) of $\Aa$
(\resp $\Aa'$) on $X$ and $\bar X$ (\resp $X'$ and $\bar X'$). We
then have base-change isomorphisms
\eqhyperref[XVII.eq5.2.1.2]{5.2.1.2} between the functors
${p'}^{\ast}L\overset{\LL}{\otimes}\bar g^{\ast}j_{!}K$ and
$j'_{!}({f'}^{\ast}L\sisi{\otimes}{\overset{\LL}{\otimes}}
{g'}^{\ast}K)$ from $D^{-}(X,\Aa_X)$ to
$D^{-}(\bar X',\Aa'_{\bar X'})$, and base-change isomorphisms
(\hyperref[XVII.thm4.3.1]{4.3.1}) between the functors
$L\overset{\LL}{\otimes}g^{\ast}Rp_{\ast}K$ and
$Rp'_{\ast}({p'}^{\ast}L\overset{\LL}{\otimes}\bar g^{\ast}K)$ from
$D^{-}(\bar X,\Aa_{\bar X})$ to $D^{-}(S',\Aa')$. By composition,
one obtains a base-change isomorphism between functors from
$D^{-}(X,\Aa_X)$ to $D^{-}(S',\Aa')$:
$$\etiquette[5.2.2.1]{XVII.eq5.2.2.1}
L\overset{\LL}{\otimes}g^{\ast}Rp_{\ast}j_{!}K\isomto
Rp'_{\ast}j'_{!}({f'}^{\ast}L\overset{\LL}{\otimes}{g'}^{\ast}K){\quoi}.\leqno(5.2.2.1)
$$

When $L$ is bounded with components of finite \ttor-dimension, the
preceding arguments and \hyperref[XVII.var4.3.2]{4.3.2}~(i) define the
base-change isomorphism for $K\in\Ob D(X,\Aa_X)$.\nde{by combining the
preceding N.D.E. and \nderef{XVII.ndecb}, this isomorphism can be
defined without any condition on $L$ or the hypothesis that the
dimension of the fibers of $p$ is bounded.}

In order to refer to \hyperref[XVII.thm5.1.8]{5.1.8}, suppose that $S$
is quasi-compact and quasi-separated. The functors
$Rp_{\ast}j_{!}=Rf_{!}$ and $Rp'_{\ast}j'_{!}=Rf'_{!}$ then depend only
on $f$ and $f'$, giving a meaning to the following statement.

\begin{enonce*}{Lemma 5.2.3}\etiquette[5.2.3]{XVII.lem5.2.3}
The base-change isomorphism \eqhyperref[XVII.eq5.2.2.1]{5.2.2.1} does
not depend on the chosen compactification $f=pj$ of $X$.
\end{enonce*}

Let $q:\bar X_1\rightarrow\bar X_2$ be a morphism of compactifications
$$\etiquette[5.2.3.1]{XVII.eq5.2.3.1}
\vcenter{
\xymatrix{
X \ar[rr]^{j_1} \ar[rrd]|(.65)\hole^{j_2} \ar[dd]_{f}& & \bar{X}_1 \ar[d]_{q} \ar[lldd]^{p_1}\\
  & &  \bar{X}_2 \ar[lld]^{p_2}\\
S & & .
}}\leqno{(5.2.3.1)}
$$
The\pageoriginale base-change isomorphism $c_1$ relative to $\bar X_1$
is obtained by ``composing'' those relative to $j_1$ and $p_1$, or,
equivalently, by composing those relative to $j_1$, $q$, and $p_2$.
The isomorphism $c_2$ relative to $\bar X_2$ is obtained by composing
those relative to $j_2$ and $p_2$. To verify that $c_1=c_2$, it
suffices to verify that the isomorphism $c'_1$ relative to $j_2$ is the
composite $c'_2$ of those relative to $j_1$ and $q$. Use a prime to
denote the arrows obtained from the arrows
\eqhyperref[XVII.eq5.2.3.1]{5.2.3.1} by the base change $g$; $c'_1$ and
$c'_2$ are isomorphisms
\[
c'_{1}, c'_{2}:L\overset{\LL}{\otimes}\bar{g}_{2}^*\ast
j_{2!}K\isomto  j'_{2!}({j'}_{2}^*\ast
L\overset{\LL}{\otimes} {g'}^{\ast}K){\quoi}.
\]
By \eqhyperref[XVII.eq5.1.3.2]{5.1.3.2}, to verify that $c'_1=c'_2$ it
suffices to verify it after restricting to $\sisi{X}{\bar X}'_2$,
which is immediate.

The base-change isomorphisms relative to two compactifications between
which a morphism exists are therefore equal, and the conclusion
follows from \hyperref[XVII.prop3.2.6]{3.2.6}~(ii).

\begin{enonce*}{Lemma 5.2.4}\etiquette[5.2.4]{XVII.lem5.2.4}
Let $f=uv:X\rightarrow Y\rightarrow S$ be a composite of
$S$-compactifiable morphisms. The base-change isomorphisms
\eqhyperref[XVII.eq5.2.2.1]{5.2.2.1} relative to $u$, $v$, and $f$
give rise to a commutative diagram (notation as in
\eqhyperref[XVII.eq5.2.4.1]{5.2.4.1} below):
\[
\xymatrix{
L\overset{\LL}{\otimes} g^{\ast} R\sisi{uv}{(uv)}_{!} K
\ar@<-4pt>[rr]^{\sim} \ar@{=}[d]& & R\sisi{u'v'}{(u'v')}_{!}((u'v')^{\ast}L
\overset{\LL}{\otimes}g''^{\ast}K)\ar@{=}[d]\\
L\overset{\LL}{\otimes} g^{\ast}Ru_{!}Rv_{!}K \ar@<-4pt>[r]^(.45){\sim}& Ru'_{!}(u'^{\ast}L
\overset{\LL}{\otimes}g'^{\ast}Rv_{!}\sisi{}{K}) \ar@<-4pt>[r]^(.45){\sim}& Ru'_{!}Rv'_{!}(v'^{\ast}u'^{\ast}L\overset{\LL}{\otimes}g''^{\ast}K).
}
\]
\end{enonce*}

Choose\pageoriginale a compactification of $(u,v)$
(\hyperref[XVII.def3.2.5]{3.2.5}):
$$\etiquette[5.2.4.1]{XVII.eq5.2.4.1}
\vcenter{
\xymatrix{
X' \ar[rr]_{g''} \ar[d]_{\sisi{u}{v}'}& & X \ar@{^{(}->}[r]^{j_1} \ar[d]_{\sisi{u}{v}}& \bar{X}
\ar[ld]_{p_1}\ar@{^{(}->}[r]^{j_2} \ar[d]& \bar{\bar{X}} \ar[ld]^{p_3}& \\
Y' \ar[rr]_{g'} \ar[d]_{\sisi{v}{u}'}& & Y \ar@{^{(}->}[r]^{j_3}\ar[d]_{\sisi{v}{u}}& \bar{Y}\ar[ld]^{p_2} &  &\\
S' \ar[rr]_{g} &  & S &      &  .
}}\leqno{(5.2.4.1)}
$$

Lemma \hyperref[XVII.lem5.2.4]{5.2.4} says that the base-change
isomorphism $c_1$ relative to $f$ is the composite $c_2$ of those
relative to $u$ and $v$. The isomorphism $c_1$ is obtained by composing
the base-change isomorphisms relative to $j_1$, $j_2$, $p_3$, and
$p_2$. For $c_2$, one composes those relative to $j_1$, $p_1$, $j_3$,
and $p_2$. It therefore suffices to compare the isomorphisms $c'_1$,
the composite of those relative to $p_3$ and $j_2$, and $c'_2$, the
composite of those relative to $j_3$ and $p_1$.

Consider the cartesian diagram
\[
\xymatrix{
  \bar{X}' \ar[ddd]^{p'_1} \ar[rrrr]^{\bar{g}''} \ar@{^{(}->}[rrd]^{j'_2} &    &       &    & \bar{X}\ar@{^{(}->}[rrd]^{j_2} \ar[ddd]^{\sisi{}{p_1}}|(.38)\hole\\
                          & &\bar{\bar{X}}' \ar[rrrr]_(.3){\bar{\bar{g}}} \ar[ddd]^{p'_3}&   &&
& \bar{\bar{X}} \ar[ddd]^{\sisi{}{p_3}}\\
 & & &\\
  Y'  \ar@{^{(}->}[rrd]_{j'_3}\ar[rrrr]^(.7){g'}|(.52)\hole \ar[rdd]_{\sisi{v}{u}'} &    &&
& Y\ar[rdd]_(.65){\sisi{v}{u}} \ar@{^{(}->}[rrd]^{j_3}\\
  & &\bar{Y}' \ar[rrrr]_{\bar{g}'} \ar[ld]^{p'_2}&        &  & &\bar{Y}\ar[ld]^{p_2}\\
 &S'  \ar[rrrr]_{g} &&& & S &,
}
\]
$L'={p'_2}^{\ast}L$\sisi{}{, $K'=j_{1!}K$}\pageoriginale and
\[
c'_{1}, c'_{2}: L'\overset{\LL}{\otimes} {\bar{g}'}^{\ast}
R(j_{3}p_{1})_{!} K\sisi{}{'}\isomto  R(j'_{3} p'_{1})_{!} ((j'_{3}
p'_{1})^{\ast} L'\otimes {\bar{g}''}^{\ast}K\sisi{}{'}){\quoi}.
\]
By \eqhyperref[XVII.eq5.1.3.2]{5.1.3.2}, to verify that $c'_1=c'_2$ it
suffices to verify that ${j'_3}^{\ast}(c'_1)={j'_3}^{\ast}(c'_2)$,
which is immediate.

\begin{enonce*}{Lemma 5.2.5}\etiquette[5.2.5]{XVII.lem5.2.5}
Let there be morphisms of schemes ringed by sheaves of torsion rings
\[
(S'', \Aa'')\xrightarrow{g_{2}} (S',
\Aa')\xrightarrow{g_{1}} (S,\Aa){\quoi},
\]
and let $f:X\rightarrow S$ be a compactifiable morphism, giving a
diagram:
\[
\xymatrix{
X'' \ar[r]_{g'_2} \ar[d]^{f''}& X' \ar[r]_{g'_1}\ar[d]^{f'}& X \ar[d]^{f}\\
S'' \ar[r]_{g_2}& S' \ar[r]_{g_1}& S \rlap{.}
}
\]
Ring\pageoriginale $X$, $X'$, and $X''$ by $f^{\ast}\Aa$,
${f'}^{\ast}\Aa'$, and ${f''}^{\ast}\Aa''$, respectively. Then the
diagram
\[
\xymatrix{
L(g_1g_2)^{\ast}Rf_{!}  \ar[rr]^{\sim} \ar@{=}[d]&& Rf''_{!}L(g'_1g'_2)^{\ast}\ar@{=}[d]\\
Lg_{\sisi{1}{2}}^{\ast}Lg^{\ast}_{\sisi{2}{1}} R f_{!} \ar[r]^{\sim}& Lg^{\ast}_{\sisi{1}{2}} Rf'_{!}
L\sisi{g^\ast_2}{g'^\ast_1}  \ar[r]^{\sim} & Rf''_{!} Lg'^{\ast}_{\sisi{1}{2}}  L g'^{\ast}_{\sisi{2}{1}}
}
\]
is commutative.
\end{enonce*}

This lemma follows immediately from the analogous results for the
functors $j_{!}$ and $Rp_{\ast}$ with $p$ proper. We leave it to the
reader to formulate the precise version with modules.

The preceding statements are summarized in the following theorem.

\begin{enonce*}{Theorem 5.2.6}\etiquette[5.2.6]{XVII.thm5.2.6}
The functor $Rf_{!}$ commutes with base change.
\end{enonce*}

\begin{enonce*}[remark]{Variants 5.2.7}\etiquette[5.2.7]{XVII.var5.2.7}
\begin{enumerate}
\renewcommand{\theenumi}{\roman{enumi}}
\renewcommand{\labelenumi}{(\theenumi)}
\item The functors $f_{!}$ (for sheaves of pointed sets:
\hyperref[XVII.var5.1.11]{5.1.11}) and the functors $R^1f_{!}$ (for
sheaves of groups: \hyperref[XVII.var5.1.12]{5.1.12}) commute with
base change.

\item The functors $Rf_{!}$ (\hyperref[XVII.var5.1.14]{5.1.14})
commute with base change.

\item Under the hypotheses of (\hyperref[XVII.subsubsec5.2.2]{5.2.2}),
if $L$ is bounded and has components of finite \ttor-dimension,\nde{as
observed in \hyperref[XVII.subsubsec5.2.2]{5.2.2}, these hypotheses are
superfluous.} the functors $Rf_{!}:D(X)\rightarrow D(S)$ commute with
base change.
\end{enumerate}
\end{enonce*}

The\pageoriginale conjunction of the following two Propositions
\hyperref[XVII.prop5.2.8]{5.2.8} and
\hyperref[XVII.prop5.2.9]{5.2.9} is essentially equivalent to Theorem
\hyperref[XVII.thm5.2.6]{5.2.6}.

\begin{enonce*}{Proposition 5.2.8}\etiquette[5.2.8]{XVII.prop5.2.8}
Let $f:X\rightarrow S$ be a compactifiable morphism and
$\mathcal F$ a torsion sheaf on $X$. For every geometric point $s$ of
$S$, let $X_s$ be the geometric fiber of $X$ at $s$ and let
$\mathcal F_s$ be the induced sheaf. Then canonically
\[
(R^{q}f_{!}(\mathcal{F}))_s=H^{q}_c(X_s,\mathcal{F}_s){\quoi}.
\]
\end{enonce*}

\begin{enonce*}{Corollary 5.2.8.1}\etiquette[5.2.8.1]{XVII.cor5.2.8.1}
If the fibers of $f$ have dimension $\leq d$, the functor $Rf_{!}$ has
cohomological dimension $\leq2d$
(\hyperref[XVII.def1.2.9]{1.2.9}): for every torsion sheaf
$\sisi{F}{ℱ}$, one has $R^if_{!}\sisi{F}{ℱ}=0$ for $i>2d$. The functor
$R^{2d}f_{!}$ is therefore right exact on the category of torsion
sheaves.
\end{enonce*}

By base change (\hyperref[XVII.prop5.2.8]{5.2.8}), we reduce to the
case where $S$ is the spectrum of an algebraically closed field. If
$\bar X$ is then a proper scheme over $S$ containing $X$ as a dense
Zariski open subscheme, with $j:X\rightarrow\bar X$, then
$\dim(\bar X)=\dim(X)\leq\sisi{D}{d}$. Hence, for $ℱ$ a torsion sheaf
on $X$ and $i>2d$, by \hyperrefext[X][X.cor4.3]{4.3},
\refnoexpose{X.cor5.2}
\[
H^{i}_c(X,ℱ)=H^{i}(\bar{X},j_{!}\sisi{F}{ℱ})=0{\quoi}.
\]

\begin{enonce*}{Proposition 5.2.9}\etiquette[5.2.9]{XVII.prop5.2.9}
Let $f:X\rightarrow S$ be a compactifiable morphism, let $\Aa$ be a
sheaf of torsion rings on $S$, let $K$ be a bounded-above complex of
left $f^{\ast}\Aa$-modules, and let $L$ be a bounded-above complex of
right $\Aa$-modules. Then canonically
\[
L\overset{\LL}{\underset{\Aa}{\otimes}}
Rf_{!}(K)=Rf_{!}(f^{\ast}L\overset{\LL}{\underset{f^{\ast}\Aa}{\otimes}}K){\quoi}.
\]
\end{enonce*}

This\pageoriginale proposition formally implies the following result:

\begin{enonce*}{Theorem 5.2.10}\etiquette[5.2.10]{XVII.thm5.2.10}
Let $f:X\rightarrow S$ be a compactifiable morphism, let $\Aa$ be a
sheaf of torsion rings on $S$, let $K\in D^{-}(X,f^{\ast}\Aa)$, and
let $d\in\ZZ$. If $K$ has finite \tTor-dimension $\leq d$
(\hyperref[XVII.def4.1.9]{4.1.9}), then $Rf_{!}K$ has finite
\tTor-dimension $\leq d$.
\end{enonce*}

It suffices to verify that, for every sheaf $L$ of right $\Aa$-modules
on $S$,
\[
\sisi{H}{\Hh}^{-i}(L\overset{\LL}{\otimes}Rf_{!}K)=0\quad\mbox{\rm
for}\quad i>d{\quoi}.
\]
By \hyperref[XVII.prop5.2.9]{5.2.9}, we know that
\[
\sisi{H}{\Hh}^{-i}(L\overset{\sisi{L}{\LL}}{\otimes}Rf_{!}K)=R^{-i}f_{!}(f^{\ast}L\overset{\sisi{L}{\LL}}{\otimes}K){\quoi},
\]
and, by hypothesis,
\[
\sisi{H}{\Hh}^{-i}(f^{\ast}L\overset{\sisi{L}{\LL}}{\otimes}K)=0\quad\mbox{\rm for}\quad
i>d{\quoi}.
\]
The conclusion follows by observing that if a complex $M$ of sheaves
of modules on $X$ satisfies $\sisi{H}{\Hh}^{-i}(M)=0$ for $i>d$, then
one still has $\sisi{H}{\Hh}^{-i}(Rf_{!}M)=0$ for $i>d$.

For constant noetherian sheaves of rings, one has the following more
general result:

\begin{enonce*}{Theorem 5.2.11}\etiquette[5.2.11]{XVII.thm5.2.11}
Let $f:X\rightarrow S$ be a morphism of sites, let $A$ be a ring, let
$C$ be its center, and let $d\in\ZZ$. Suppose that
\begin{enumerate}
\renewcommand{\theenumi}{\roman{enumi}}
\renewcommand{\labelenumi}{(\theenumi)}
\item the functor $Rf_{\ast}:D^{+}(X,C)\rightarrow D^{+}(S,C)$ has
finite cohomological dimension;

\item $A$ is right noetherian;

\item $S$ has enough points.
\end{enumerate}
For\pageoriginale every complex $K\in\Ob D^{-}(X,A)$ of sheaves of
left $A$-modules, if $K$ has finite \tTor-dimension $\leq d$, then the
complex $Rf_{\ast}K\in\Ob D^{-}(S,A)$ has \tTor-dimension $\leq d$.
\end{enonce*}

\begin{enonce*}[remark]{Proof}
By (i), the functor $f_{\ast}$ is derivable to functors
\[
Rf_{\ast}:D(X,A)\longrightarrow D(S,A)
\]
and
\[
Rf_{\ast}:D(X,C)\longrightarrow D(S,C){\quoi}.
\]
\end{enonce*}

By \hyperref[XVII.lem4.2.8]{4.2.8}, it suffices to verify that, for
every point $s$ of $S$, the complex $(Rf_{\ast}K)_s$ of $A$-modules
has \tTor-dimension $\leq d$; for this it suffices that, for every
finitely generated right $A$-module $L$,
\[
H^{-i}(L\overset{L}{\otimes}Rf_{\ast}K)=0\quad\mbox{\rm for}\quad
i>d{\quoi}.
\]

Reasoning as in \hyperref[XVII.thm5.2.10]{5.2.10}, we see that it
suffices for the ``base-change arrow''
$$\etiquette[5.2.11.1]{XVII.eq5.2.11.1}
L\overset{L}{\otimes}Rf_{\ast}K\longrightarrow Rf_{\ast}(L\overset{L}{\otimes}K)\leqno(5.2.11.1)
$$
to be an isomorphism.

Let $L^{\ast}$ be a resolution of $L$ by finitely generated free
$A$-modules. Since the functors under consideration are way-out to the
right, it suffices to prove that
\eqhyperref[XVII.eq5.2.11.1]{5.2.11.1} is an isomorphism after replacing
$L$ by any component of $L\sisi{_{\ast}}{^\ast}$, which reduces us to
the immediate case where $L$ is finitely generated and free.

\begin{enonce*}[remark]{Remarks 5.2.12}\etiquette[5.2.12]{XVII.rems5.2.12}
\begin{enumerate}
\renewcommand{\theenumi}{\roman{enumi}}
\renewcommand{\labelenumi}{(\theenumi)}
\item The\pageoriginale hypothesis (iii) of
\hyperref[XVII.thm5.2.11]{5.2.11} is probably unnecessary.

\item If the noetherian hypothesis on $A$ is omitted,
\hyperref[XVII.thm5.2.11]{5.2.11} remains valid when the functors
$R^if_{\ast}$ commute with filtered direct limits.

\item The hypothesis that $S$ is ringed by a \emph{constant} sheaf of
rings is essential. A counterexample is obtained by taking for $S$ the
spectrum of a strictly local discrete valuation ring, for $f$ the
inclusion of the generic point $\eta$, for the sheaf of rings $\Aa$ the
sheaf equal to $\ZZ/n\ZZ$ at the generic point and to $\ZZ/n^2\ZZ$ at
the closed point $s$ ($n>1$ prime to the residue characteristics), and
for the sheaf the module $f^{\ast}\Aa$.
\end{enumerate}
\end{enonce*}

\subsection{Comparison theorem and finiteness theorem}\etiquette[5.3]{XVII.subsec5.3}

\subsubsection{}\etiquette[5.3.1]{XVII.subsubsec5.3.1}
To a scheme $X$ of finite type over $\CC$ are associated its étale
site $X_{\text{\rm et}}$, its transcendental site $X_{\cl}$, and a
morphism of sites (\hyperrefext[XI][XI.sec4]{4})
\[
\epsilon:X_{\cl}\longrightarrow X_{\et}{\quoi}.
\]

If $j:U\rightarrow X$ is an open immersion, the commutative diagram
$$
\xymatrix{
U_{\cl}  \ar[r]\ar[d]   & U_{\et} \ar[d]\\
X_{\cl}  \ar[r] & X_{\et}
}
$$
identifies $U^{\sim}_{\cl}$ with the open subtopos of
$X^{\sim}_{\cl}$ that is the inverse image of the open subtopos
$U_{\widetilde{\text{\rm et}}}$ of
$X_{\widetilde{\text{\rm et}}}$.

If $p:Y\rightarrow X$ is a proper morphism, the continuous map
$p^{\an}:Y^{\an}\rightarrow X^{\an}$ is proper and separated. To see
this, Chow's lemma reduces us to the case where $p$ is projective.

Let\pageoriginale $f:X\rightarrow S$ be a compactifiable morphism of
schemes of finite type over $\CC$, and let $f=pj$ be a compactification
of $f$.
$$\etiquette[5.3.1.1]{XVII.eq5.3.1.1}
\vcenter{
\xymatrix{
X_{\cl} \ar[rrr]_{\epsilon} \ar[rd]^{j_{\cl}} \ar[dd]_{f_{\cl}}&  & & X \ar[rd]^{j}\ar[dd]_(.6){f}|\hole & \\
            &\bar{X}_{\cl}\ar[ld]^{p_{\cl}}  \ar[rrr]_{\epsilon} &  & &\bar{X}\ar[ld]^{p} \\
S_{\cl}  \ar[rrr]_{\epsilon} &                 &   &  S & .
}}\leqno{(5.3.1.1)}
$$
The functor $p_{\cl\ast}j_{\cl !}$ is precisely the direct-image
functor with proper supports $f_{\cl !}$ (see \textsc{Verdier}
\cite{sga4.exp.XVII.bib.V3}). If $S_{\cl}$, and hence $X_{\cl}$, is separated, the
functor $j_{\cl !}$ carries soft sheaves on $X_{\cl}$ to soft sheaves
on $\bar X_{\cl}$ \cite[II 3.5.5]{sga4.exp.XVII.bib.TF}. Since the question is local on
$S$, it in any case carries flasque sheaves on $X_{\cl}$ to sheaves on
$\sisi{X}{\bar X}_{\cl}$ that are acyclic for the functor
$p_{\cl\ast}$. Hence (\cf \textsc{Verdier} \cite{sga4.exp.XVII.bib.V3})
$$
Rf_{\cl !}=Rp_{\cl\ast}j_{\cl !}{\quoi}.
$$

\subsubsection{}\etiquette[5.3.2]{XVII.subsubsec5.3.2}

% END INCLUDED SGA4_English_current_buildable_source_20260906/source/17/drafts/SGA4_XVII_LINES_4201_4600_English.texfrag

% BEGIN INCLUDED SGA4_English_current_buildable_source_20260906/source/17/drafts/SGA4_XVII_LINES_4601_5000_English.texfrag

By \hyperref[XVII.lem5.1.2]{5.1.2}, there is a base-change
isomorphism between functors from $D(X)$ to $D(\bar X_{\cl})$:
$$
\epsilon^{\ast}j_{!}\isomto j_{\cl !}\epsilon^{\ast}{\quoi}.
$$

By\pageoriginale \hyperref[XVII.subsubsec4.2.12]{4.2.12} or
\hyperref[XVII.prop4.1.5]{4.1.5}, there is a base-change morphism
between functors from $D^{+}(\bar X)$ to $D^{+}(S_{\cl})$:
$$
\epsilon^{\ast}Rp_{\ast}\longrightarrow Rp_{\cl
\ast}\epsilon^{\ast}{\quoi}.
$$
By composition one obtains a \emph{base-change morphism between
functors from} $D^{+}_{\tors}(X)$ to $D^{+}_{\tors}(S_{\cl})$
(\hyperref[XVII.def5.1.13]{5.1.13}):
$$\etiquette[5.3.2.1]{XVII.eq5.3.2.1}
\epsilon^{\ast}Rf_{!}\longrightarrow Rf_{\cl
!}\epsilon^{\ast}{\quoi}.\leqno(5.3.2.1)
$$

As in \hyperref[XVII.lem5.2.3]{5.2.3}, one checks that it does not
depend on the chosen compactification of $f$; it satisfies a
transitivity formula of the type in \hyperref[XVII.lem5.2.4]{5.2.4}.

\begin{enonce*}{Theorem 5.3.3}\etiquette[5.3.3]{XVII.thm5.3.3}
\emph{(Comparison theorem).} The morphism
\eqhyperref[XVII.eq5.3.2.1]{5.3.2.1} is an isomorphism.
\end{enonce*}

By construction, it suffices to treat the case where $f$ is proper.
The functors under consideration are triangulated and ``way-out''
\cite[I 7]{sga4.exp.XVII.bib.RD}. It therefore suffices to treat the case where $K$ is
a torsion sheaf placed in degree $0$, and the result follows from
\hyperrefext[XVI][XVI.thm4.1]{4.1}~(i).

\begin{enonce*}[remark]{Variants 5.3.4}\etiquette[5.3.4]{XVII.var5.3.4}
\begin{enumerate}
\renewcommand{\theenumi}{\roman{enumi}}
\renewcommand{\labelenumi}{(\theenumi)}
\item If $\Aa$ is a sheaf of torsion rings on $S_{\text{\rm et}}$,
one similarly obtains an isomorphism
\eqhyperref[XVII.eq5.3.2.1]{5.3.2.1} between functors from
$D(X,f^{\ast}\Aa)$ to $D(S,\epsilon^{\ast}\Aa)$.

\item In low degrees one obtains the usual variants for sheaves of
pointed sets and sheaves of \iInd-finite groups.
\end{enumerate}
\end{enonce*}

\begin{enonce*}{Corollary 5.3.5}\etiquette[5.3.5]{XVII.cor5.3.5}
Let\pageoriginale $X$ be a separated scheme of finite type over $\CC$
which admits a compactification in a proper scheme over $\CC$, and let
$F$ be a torsion sheaf on $X$. Then
$$
H^{q}_{!}(X,F)\isomto H^{q}_c(X_{\cl},\epsilon^{\ast}F){\quoi}.
$$
\end{enonce*}

\begin{enonce*}{Theorem 5.3.6}\etiquette[5.3.6]{XVII.thm5.3.6}
Let $f:X\rightarrow S$ be a compactifiable morphism of finite
presentation and let $A$ be a left noetherian torsion ring. The
functor $Rf_{!}:D(X,A)\rightarrow D(S,A)$ carries complexes with
constructible cohomology to complexes with constructible cohomology.
\end{enonce*}

The assertion is local on $S$, which may be assumed affine. By passage
to the limit, we reduce to the case where $S$ is noetherian. Complexes
with constructible cohomology form a triangulated subcategory of
$D(S,A)$. Since the functor $Rf_{!}$ is ``way-out,'' we reduce to
verifying that if $\mathcal F$ is a constructible sheaf of $A$-modules,
then the sheaves $R^qf_{!}\mathcal F$ are constructible. Write $f=pj$,
with $p$ proper and $j$ an open immersion. Then
$R^qf_{!}\mathcal F=R^qp_{\ast}(j_{!}\mathcal F)$. The sheaf
$j_{!}\mathcal F$ is constructible, and the conclusion follows from
\hyperrefext[XIV][XIV.thm1.1]{1.1}.

\setcounter{subsubsection}{6}
\subsubsection{}\etiquette[5.3.7]{XVII.subsubsec5.3.7}
In low degrees, there are analogous statements for direct images with
proper supports of a constructible sheaf of \iInd-finite sets or
groups.

\begin{enonce*}{Corollary 5.3.8}\etiquette[5.3.8]{XVII.cor5.3.8}
Let $X$ be a separated scheme of finite type over a separably closed
field $k$, and let $\Ff$ be a constructible torsion sheaf on $X$.
Suppose that $X$ can be compactified in a proper scheme over $k$.\nde{this
hypothesis is automatically satisfied. See \nderef{XVII.ndeConrad}.}
Then the groups $H^q_c(X,\mathcal F)$ are finite.
\end{enonce*}

\subsection{Künneth formula}\etiquette[5.4]{XVII.subsec5.4}\pageoriginale

\subsubsection{}\etiquette[5.4.1]{XVII.subsubsec5.4.1}
Let there be a commutative diagram (\hyperref[XVII.subsec0.11]{0.11})
of sites
$$\etiquette[5.4.1.1]{XVII.eq5.4.1.1}
\vcenter{
\xymatrix@R=.5cm@C=1cm{
             &  Z'\ar[ld]_{p'}\ar[rd]^{q'} \ar[dd]^{h}&    \\
X' \ar[dd]_{f}&                   &  Y' \ar[dd]^{g}\\
             &       Z\ar[ld]_{p} \ar[rd]^{q}   &    \\
X \ar[rd]           &                    & Y \ar[ld]\\
             &      S             & &.
}}\leqno{(5.4.1.1)}
$$
Suppose $S$ is ringed by a sheaf of rings $\Aa$ with center $\Cc$; we
again denote by $\Aa$ and $\Cc$ the inverse images of these sheaves on
the sites of \eqhyperref[XVII.eq5.4.1.1]{5.4.1.1}.

If $K$ is a sheaf of right $\Aa$-modules on $X$ and $L$ a sheaf of
left $\Aa$-modules on $Y$, set
\[
K\mathop{\boxtimes}_{\Aa}L=p^{\ast}\sisi{}{K}
\underset{\Aa}{\otimes}q^{\ast}L
\]
(external tensor product). Similarly for $K$ (\resp $L$) on $X'$
(\resp $Y'$). Now let $K$ be a sheaf of right $\Aa$-modules on $X'$
and $L$ a sheaf of left $\Aa$-modules on $Y'$. The base-change arrows
(\sisi{}{\hyperref[XVII.prop2.1.3]{2.1.3}, }
\sisi{\hyperrefext[XII][XII.eq4.1]{4.1}}{\neweqhyperrefext[XII][XII.eq4.1]{4.1}})
define
$$\etiquette[5.4.1.2]{XVII.eq5.4.1.2}
\begin{cases}
p^{\ast}f_{\ast}K\longrightarrow h_{\ast}{p'}^{\ast}K\\
q^{\ast}g_{\ast}L\longrightarrow h_{\ast}{q'}^{\ast}L{\quoi},
\end{cases}\leqno(5.4.1.2)
$$
and hence,\pageoriginale by taking their tensor product, an arrow
$$\etiquette[5.4.1.3]{XVII.eq5.4.1.3}
f_{\ast}K\boxtimes_{\Aa} g_{\ast}L\longrightarrow
h_{\ast}{p'}^{\ast}K\otimes_{\Aa} h_{\ast}
{q'}^{\ast}L\longrightarrow h_{\ast}(K\boxtimes_{\Aa}L){\quoi}.\leqno(5.4.1.3)
$$

Suppose that the functors
\begin{align*}
Rf_{\ast} &:D^{+}(X',\Aa^{\circ})\longrightarrow D^{+}(X,\Aa^{\circ})\\
Rg_{\ast} &:D^{+}(Y',\Aa)\longrightarrow D^{+}(Y,\Aa)\\
Rh_{\ast} &:D^{+}(Z',\Cc)\longrightarrow D^{+}(Z,\Cc)
\end{align*}
have finite cohomological dimension and therefore extend to the
unbounded derived categories. For $K\in\Ob D^{-}(X',\Aa^{\circ})$ and
$L\in\Ob D^{-}(Y',\Aa)$, one is then led to define a
\emph{Künneth arrow}
$$\etiquette[5.4.1.4]{XVII.eq5.4.1.4}
Rf_{\ast}K\overset{\LL}{\boxtimes}_{\Aa}
Rg_{\ast}L\longrightarrow Rh_{\ast}(K
\overset{\LL}{\boxtimes}_{\Aa}L){\quoi}.\leqno(5.4.1.4)
$$

We define it only when the sites under consideration have enough
points, so that \hyperref[XVII.prop4.2.10]{4.2.10} and
\hyperref[XVII.lem4.2.8]{4.2.8} may be used. By
\hyperref[XVII.prop4.2.10]{4.2.10}, it suffices to define it for a pair
$(K,L)$ left adapted to the functor
$(K,L)\mapsto K\boxtimes_{\Aa}L$, where $K$ and $L$ are bounded above
and right adapted to the functors $f_{\ast}$ and $g_{\ast}$.

Let $K_1$ be a flat resolution of $f_{\ast}K$, let $L_1$ be a flat
resolution of $g_{\ast}L$, and let
$u:K\overset{\LL}{\boxtimes}\sisi{}{_\Aa}L\rightarrow M$ be a
resolution of $K\overset{\LL}{\boxtimes}\sisi{}{_\Aa}L$ adapted to
the functor $Rh_{\ast}$. We have morphisms
\begin{gather*}
Rf_{\ast}K\overset{\LL}{\boxtimes}_{\Aa}Rg_{\ast}L\simeq
K_1\boxtimes_{\Aa}L_1\longrightarrow
f_{\ast}K\boxtimes_{\Aa}g_{\ast}L\longrightarrow\\
\longrightarrow h_{\ast}(K\boxtimes_{\Aa}L)\longrightarrow
h_{\ast}M\simeq Rh_{\ast}(K\overset{\LL}{\boxtimes}_{\Aa}L)
\end{gather*}
and define \eqhyperref[XVII.eq5.4.1.4]{5.4.1.4} to be their composite.

This\pageoriginale arrow has the following characteristic property,
whose proof is entirely analogous to that of
\hyperref[XVII.prop4.2.13]{4.2.13}:

\setcounter{paragraph}{4}
\paragraph{}\etiquette[5.4.1.5]{XVII.par5.4.1.5}
Let $K\in\Ob C^{-}(X',\Aa^{\circ})$ and
$L\in\Ob C^{-}(Y',\Aa)$. Set $K_1=f_{\ast}K$,
$L_1=g_{\ast}L$, and $M_2=K\boxtimes\sisi{}{_\Aa}L$. The following
diagram is then commutative:
$$
\xymatrix@R=1.25cm{
Rf_{\ast}K
\overset{\LL}{\boxtimes}_{\Aa}Rg_{\ast}
L \ar[rrr]& & & Rh_{\ast}(K
\overset{\LL}{\boxtimes}_{\Aa}L) \ar[d]\\
K_1\overset{\LL}{\boxtimes}_{\Aa}L_{\sisi{k}{1}}
\ar[r] \ar[u]&
K_1\boxtimes\sisi{}{_\Aa}L_1 \ar[r]& h_{\ast}M_2 \ar[r]& Rh_{\ast}M_2.
}
$$

There is an analogous commutativity (\cf
\hyperref[XVII.prop4.2.13]{4.2.13}) starting instead from morphisms
$K_1\rightarrow f_{\ast}K$, $L_1\rightarrow g_{\ast}L$, and
$K\boxtimes\sisi{}{_\Aa}L\rightarrow M_2$.

The commutativity in \hyperref[XVII.par5.4.1.5]{5.4.1.5} makes it, in
principle, easy to verify the compatibility properties of the Künneth
morphism.

\subsubsection{}\etiquette[5.4.2]{XVII.subsubsec5.4.2}
Let \eqhyperref[XVII.eq5.4.1.1]{5.4.1.1} be a commutative diagram of
schemes, and let $\Aa$ be a torsion sheaf on $S$. Suppose that $f$ and
$g$ are compactifiable (\hyperref[XVII.def3.2.1]{3.2.1}), that
$Z=\fprod{X}{Y}{S}$, and that $Z'=\fprod{X'}{Y'}{S}$. Factor $f$ and
$g$ as a proper morphism and an open immersion; this gives a diagram
$$
\xymatrix@C=1.75cm@R=1cm{
  X'      \sisi{\ar[r]^{p'}}{} \ar[d]^{j_1} \ar@/_1.75pc/[dd]_{f}  & Z'   \ar[d]^{j_3}\sisi{}{\ar[l]_{p'}\ar[r]^{q'}}
& \sisi{Y'\ar[l]_{q'}}{} \ar[d]^{j_2} \ar@/^1.75pc/[dd]^{g}\\
\bar{X}' \sisi{\ar[r]^{\bar{p}}}{} \ar[d]^{\bar{f}}  & \bar{Z}' \ar[d]^{\bar{h}}\sisi{}{\ar[l]_{\bar{p}}\ar[r]^{\bar{q}}} & \bar{Y}' \sisi{\ar[l]_{\bar{q}}}{}\ar[d]^{\bar{g}}\\
X            &    Z \ar[l]_{p} \ar[r]^{q}  & Y
}
$$
in\pageoriginale which $\bar Z'=\fprod{\bar X'}{\bar Y'}{S}$. For
complexes $K$ and $L$ on $X'$ and $Y'$, there is an isomorphism of an
immediate nature
$$\etiquette[5.4.2.1]{XVII.eq5.4.2.1}
j_{3!}(K\boxtimes_{\Aa}L)\isomto
j_{1!}K\boxtimes_{\Aa}j_{2!}L{\quoi}.\leqno(5.4.2.1)
$$

By ``composing'' the inverse of this isomorphism with the Künneth arrow
\eqhyperref[XVII.eq5.4.1.4]{5.4.1.4} relative to
$(\bar f,\bar g,\bar h,j_{1!}K,j_{2!}L)$, one obtains, \emph{for
$K\in\Ob D^{-}(X',\Aa)$ and $L\in\Ob D^{-}(Y',\Aa)$, a Künneth arrow}
$$\etiquette[5.4.2.2]{XVII.eq5.4.2.2}
Rf_{!}K\overset{\LL}{\boxtimes}\sisi{}{_\Aa}Rg_{!}L\longrightarrow
Rh_{!}(K\overset{\LL}{\boxtimes}\sisi{}{_\Aa}L){\quoi},\leqno(5.4.2.2)
$$
which one checks is independent of the chosen compactifications.

\begin{enonce*}{Theorem 5.4.3}\etiquette[5.4.3]{XVII.thm5.4.3}
Under the hypotheses of \hyperref[XVII.subsubsec5.4.2]{5.4.2}, the
Künneth arrow for cohomology with proper supports
\eqhyperref[XVII.eq5.4.2.2]{5.4.2.2} is an isomorphism.
\end{enonce*}

\begin{enonce*}[remark]{Proof}
\begin{enumerate}
\renewcommand{\theenumi}{\alph{enumi}}
\renewcommand{\labelenumi}{(\theenumi)}
\item By construction, it suffices to verify the result when $f$ and
$g$ (and hence also $h$) are proper.

\item Let there be a commutative diagram
$$\etiquette[5.4.3.1]{XVII.eq5.4.3.1}
\vcenter{
\xymatrix@C=1.75cm@R=1.25cm{
X_1 \ar[r]^{p_1} \ar[d]& S_1 \sisi{\ar[r]^{q_1}}{}\ar[d]& Y_1\ar[d]\sisi{}{\ar[l]_{q_1}}\\
X   \ar[r]^{p}& S  \sisi{\ar[r]^{q}}{} &  Y\sisi{}{\ar[l]_q}.
}}\leqno{(5.4.3.1)}
$$
and let $Z_1=\fprod{X_1}{Y_1}{S_1}$,
$X'_1=\fprod{X'}{X_1}{X}$, $Y'_1=\fprod{Y'}{Y_1}{Y}$, and
$Z'_1=\fprod{X'_1}{Y'_1}{S_1}=\fprod{Z'}{Z\sisi{_{1}}{}}{Z_1}$, so
that there is a morphism of diagrams
\sisi{(\hyperref[XVII.subsubsec5.4.1]{5.4.1})}{\eqhyperref[XVII.eq5.4.1.1]{5.4.1.1}},
\[
(S_1,X_1,Y_1,Z_1,X'_1,Y'_1,Z'_1)\longrightarrow
(S,X,Y,Z,X',Y',Z')
\]
all of whose arrows will be denoted uniformly by $C$.

We\pageoriginale shall admit that:

\begin{enonce*}{Lemma 5.4.3.2}\etiquette[5.4.3.2]{XVII.lem5.4.3.2}
The following diagram is commutative:
$$
\xymatrix{
C^{\ast}(Rf_{!} K \overset{L}{\boxtimes}_{\Aa}
Rg_{!} L) \ar@{=}[r] \ar[d]& C^{\ast}Rf_{!}K \overset{L}{\boxtimes}_{\Aa}
C^{\ast} Rg_{!} L \ar[r]^{\sim}&  Rf_{1!} C^{\ast}K
\overset{L}{\boxtimes}_{\Aa} Rg_{1!} C^{\ast}
L \ar[d]\\
C^{\ast}(Rh_{!}(K \overset{L}{\boxtimes}_{\Aa}
L)) \ar[r]^{\sim}& Rh_{\sisi{}{1}!}(C^{\ast}(K
 \overset{L}{\boxtimes}_{\Aa}  L))\ar@{=}[r] &Rh_{\sisi{}{1}!}(C^{\ast}K
 \overset{L}{\boxtimes}_{\Aa}  C^*L ).
}
$$
\end{enonce*}

In this diagram, the horizontal arrows are base-change arrows and the
vertical arrows are the Künneth arrows for the two diagrams
\eqhyperref[XVII.eq5.4.1.1]{5.4.1.1} under consideration.

To prove \hyperref[XVII.thm5.4.3]{5.4.3}, it suffices to prove that the
Künneth arrow is an isomorphism at every geometric point of $Z$. By
\hyperref[XVII.lem5.4.3.2]{5.4.3.2}, it suffices to verify this after
any base change \eqhyperref[XVII.eq5.4.3.1]{5.4.3.1} such that $S_1$
is the spectrum of an algebraically closed field and $p_1,q_1$ are
isomorphisms.

After such a base change, diagram
\eqhyperref[XVII.eq5.4.1.1]{5.4.1.1} reduces to a diagram
$$\vcenter{\etiquette[5.4.3.3]{XVII.eq5.4.3.3}
\xymatrix@R=1.5cm@C=2cm{
    & \fprod{X}{Y}{S} \ar[ld]_{p} \ar[rd]^{q} \ar[dd]^{f\times g}&     \\
X \ar[rd]_{f}  &                   &    Y \ar[ld]^{g}\\
    &      S             &\sisi{.}{}
}}\leqno{(5.4.3.3)}
$$
and the Künneth arrow reduces to an arrow
\eqhyperref[XVII.eq5.4.4.1]{5.4.4.1}
$$\etiquette[5.4.3.4]{XVII.eq5.4.3.4}
k:Rf_{!}K\overset{\LL}{\otimes}\sisi{}{_\Aa} Rg_{!} L\longrightarrow
R(f\times g)_{!}(K\overset{\LL}{\boxtimes}\sisi{}{_\Aa}L){\quoi}.\leqno(5.4.3.4)
$$

\item Let\pageoriginale \eqhyperref[XVII.eq5.4.3.3]{5.4.3.3} be a
diagram in which $f$ and $g$ are proper. We shall admit that:

\begin{enonce*}{Lemma 5.4.3.5}\etiquette[5.4.3.5]{XVII.lem5.4.3.5}
The following diagram is commutative:
$$
\xymatrix@C=1.75cm{
Rf_{!}K \overset{\sisi{L}{\LL}}{\underset{\Aa}{\otimes}} Rg_{!}L
\ar[r] \ar[d]& R(f\times g)_{!}(K
\overset{\sisi{L}{\LL}}{\underset{\Aa}{\otimes}} L) \ar@{=}[d]\\
Rf_{!}(K \overset{\sisi{L}{\LL}}{\sisi{\boxtimes}{\otimes}}_{\Aa}
f^{\ast}Rg_{!}L) \ar[r]& Rf_{!} Rp_{!}(p^{\ast}K \overset{\sisi{}{\LL}}{\otimes}\sisi{}{_\Aa} q^{\ast}L).
}
$$
\end{enonce*}

In this diagram, the horizontal arrows are \sisi{the arrows}{the arrow}
of Künneth and the base-change arrow. The vertical arrows are a
base-change arrow by $f$ and a transitivity isomorphism. The theorem
\sisi{\hyperref[XVII.subsubsec5.5.9]{5.5.9}}{\hyperref[XVII.thm5.4.3]{5.4.3}}
follows from this lemma and \hyperref[XVII.thm4.3.1]{4.3.1}.

The diligent reader may verify that this lemma remains valid if $f$
and $g$ are not necessarily proper.
\end{enumerate}
\end{enonce*}

\setcounter{subsubsection}{3}
\subsubsection{}\etiquette[5.4.4]{XVII.subsubsec5.4.4}
Let $S$ be a scheme, let $f_i:X'_i\rightarrow X_i$ $(i\in I)$ be a
finite family of compactifiable morphisms of $S$-schemes, let $X$ be
the product over $S$ of the $X_i$, let $X'$ be the product over $S$ of
the $X'_i$, and let $f:X'\rightarrow X$ be the product of the $f_i$.
If $S$ is ringed by $\Aa$ as in
\hyperref[XVII.subsubsec5.4.1]{5.4.1}, and the $X_i$ and $X'_i$ are
ringed by the inverse image of $\Aa$, there is more generally an
isomorphism
$$\etiquette[5.4.4.1]{XVII.eq5.4.4.1}
\mathop{\boxtimes}^{\LL}_{i\in I} Rf_{i!}(K_i)
\isomto Rf_{!}(\mathop{\boxtimes}^{\LL}_{i\in I}K_i){\quoi},\leqno(5.4.4.1)
$$
for $K_i\in\Ob D^{-}(X'_i,\Aa)$. Neither the definition of this
isomorphism nor that of the two sides of
\eqhyperref[XVII.eq5.4.4.1]{5.4.4.1} \emph{requires} the choice of a
total order on $I$.

\subsection{Cohomology of symmetric powers}\etiquette[5.5]{XVII.subsec5.5}
\pageoriginale

The main result of this ${\rm n^{\circ}}$ is Theorem
\hyperref[XVII.thm5.5.21]{5.5.21}, which in particular gives the
cohomology with proper supports of a symmetric power of a scheme $X$
over an algebraically closed field $k$, for suitable coefficients (for
example $\ZZ/n\ZZ$), in terms of the cohomology with proper supports
of $X$. The analogue of this theorem in semisimplicial topology is
well known, at least for cohomology with coefficients in a commutative
ring $A$; it can be verified at the level of semisimplicial cochains.
Indeed, if $X_{\ast}$ is a semisimplicial set, there is a canonical
continuous bijection from the symmetric $n\sisi{th}{-th}$ power
$(X_{\ast})^n/\Sf_n$ of $X_{\ast}$ to the
$n$\textsuperscript{th} symmetric power of the geometric realization
of $X_{\ast}$, and this bijection induces an isomorphism on cohomology.

\subsubsection{}\etiquette[5.5.1]{XVII.subsubsec5.5.1}
Recall some properties (due to \textsc{D. Lazard}) of flat modules,
some properties (due to \textsc{Roby}) of the functors $TS^n$ and
$T^n$, and some properties (due to \textsc{Dold} and \textsc{Puppe})
of the derived functors of nonadditive functors.

\paragraph{}\etiquette[5.5.1.1]{XVII.par5.5.1.1}
\emph{Every flat module over a ring $A$ is a filtered direct limit of
finitely generated free modules} (\textsc{Lazard} \cite[I 1.2]{sga4.exp.XVII.bib.La}).
This result implies the following statements:

\paragraph{}\etiquette[5.5.1.2]{XVII.par5.5.1.2}
The functor $\varinjlim$ is an equivalence from the category of
\iInd-objects (\hyperrefext[I][I.sec8.1]{8.1}) of the category of
finitely generated free (\resp finitely generated projective)
$A$-modules to the category of flat $A$-modules.

\paragraph{}\etiquette[5.5.1.3]{XVII.par5.5.1.3}
If\pageoriginale $A$ and $B$ are two rings, every functor
$T:\Mod(A)\rightarrow\Mod(B)$ that carries finitely generated free
modules to flat modules and commutes with filtered direct limits
carries flat modules to flat modules (\textsc{Lazard}
\cite[I 1.5]{sga4.exp.XVII.bib.La}).

\paragraph{}\etiquette[5.5.1.4]{XVII.par5.5.1.4}
Every exact sequence of $A$-modules, with $L''$ flat,

% END INCLUDED SGA4_English_current_buildable_source_20260906/source/17/drafts/SGA4_XVII_LINES_4601_5000_English.texfrag

% BEGIN INCLUDED SGA4_English_current_buildable_source_20260906/source/17/drafts/SGA4_XVII_LINES_5001_5400_English.texfrag
$$
0\longrightarrow L'\longrightarrow L\longrightarrow L''\longrightarrow
0
$$
is a filtered inductive limit of split exact sequences (\textsc{Lazard} \cite[I 2.3]{sga4.exp.XVII.bib.La}).

\subsubsection{}\etiquette[5.5.2]{XVII.subsubsec5.5.2}
Let $A$ be a commutative ring and $M$ an $A$-module.

\paragraph{}\etiquette[5.5.2.1]{XVII.par5.5.2.1}
For every integer $n$, denote by $TS^{n}(M)$ (or $TS^{n}_A(M)$) the
module of symmetric tensors of degree $n$ of $M$. Denoting by
$\Sf_n$ the symmetric group of degree $n$, by definition one has
$$
TS^{n}(M)=(\bigotimes^{n}M)^{\Sf_n}{\quoi}.
$$

If $B$ is an $A$-algebra, then $TS^{n}(B)$ is a subalgebra of
$\bigotimes\limits^{n}B$. If, moreover, $B$ is commutative,
then $\Spec (TS^{n}(B))$ is by definition the $n$\textsuperscript{th}
symmetric power of $\Spec(B)$ over $\Spec(A)$:
$$
\Spec(TS^{n}_{A}(B))=\Sym^{n}_{\Spec(A)}(\Spec(B)\underset{\dfn}{=}(\Spec(B)/\Spec(A))^{n}/\Sf_n{\quoi}.
$$

\paragraph{}\etiquette[5.5.2.2]{XVII.par5.5.2.2}
Recall (\textsc{Roby} \cite[I \sisi{1 pg.226}{2 p.~219}]{sga4.exp.XVII.bib.Ro}) that a \emph{polynomial law} $f$ from an
$A$-module $M$ to an $A$-module $N$ is the datum, for every
commutative $A$-algebra $A'$, of a map
$f_{A'}:M\otimes_{A}A'\rightarrow N\otimes_{A}A'$, these maps
being such that, for every homomorphism of $A$-algebras
$u:A'\rightarrow A''$, the diagram
\[
\xymatrix{
M \otimes_A A'  \ar[r]^{f_{A'}} \ar[d]^{u}& N \otimes_A A' \ar[d]^{u}\\
M \otimes_A A'' \ar[r]^{f_{A''}}& N \otimes_A A''
}
\]
is\pageoriginale commutative.

A polynomial law $f$ is said to be \emph{homogeneous of degree} $n$ (\textsc{Roby} \cite[I
\sisi{1}{8} p\sisi{g}{}.~226]{sga4.exp.XVII.bib.Ro}) if, for every $a'\in A'$, one has identically
\[
f_{A'}(a'x)={a'}^{n}f_{A'}(x){\quoi}.
\]
Such a polynomial law will also be called an \emph{$n$\textsuperscript{ic} map}
(quadratic, cubic, \ldots); their set will be denoted by
\[
\Hom n^{\text{\rm ique}} (M,N){\quoi}.
\]

More generally, if $(M_{i})_{i\in I}$ is a finite family of
$A$-modules and $\sheafn =(n_{i})_{i\in I}$ a family of integers, a
multihomogeneous polynomial law of multidegree $\sheafn $ from the modules $M_{i}$
to $N$ is a polynomial law from $\sisi{\sum}{\oplus} M_{i}$ to $N$ satisfying
identically, for $a_{i}\in A'$,
\[
f_{A'}((a_{i}x_{i})_{i\in I})=\prod_{i\in I} a^{n_{i}}_i\cdot
f_{A'}((x_{i})_{i\in I}){\quoi}.
\]
We shall denote by $\Hom \sheafn ^{\text{\rm ique}}((M_{i}),N)$ the set of these
polynomial laws.

\paragraph{}\etiquette[5.5.2.3]{XVII.par5.5.2.3}
For every module $M$, the covariant functor from $A$-modules to
sets
$$
N\longmapsto \Hom n^{\text{\rm ique}} (M,N)
$$
is represented by an $A$-module $\Gamma^{n}(M)$ (or
$\Gamma^{n}_{A}(M)$), equipped with a \emph{universal\pageoriginale
$n$\textsuperscript{ic} map} $\gamma^{n}:M\rightarrow
\Gamma^{n}(M)$ (\textsc{Roby} \cite[IV 1 \sisi{pg.226}{p.~265}]{sga4.exp.XVII.bib.Ro}). Given $(M_{i})_{i\in I}$
and $\sheafn $ as in \hyperref[XVII.par5.5.2.2]{5.5.2.2}, the covariant functor
$$
N\longmapsto \Hom \sheafn ^{\text{\rm ique}} ((M_{i}),N)
$$
is represented by the module $\bigotimes\limits_{i\in
I}\Gamma^{n_i}(M_{i})$, equipped with the \emph{universal multihomogeneous
polynomial law (of multidegree $\sheafn $)} $(m_{i})\mapsto
\bigotimes\limits_{i\in I}\gamma^{n_{i}}(m_{i})$.

\paragraph{}\etiquette[5.5.2.4]{XVII.par5.5.2.4}
The functors $TS^{n}$ and $\Gamma^{n}$ commute with filtered
inductive limits (for the latter, \textsc{Roby} \cite[IV \sisi{6 pg.275}{7 p.~277}]{sga4.exp.XVII.bib.Ro}). They
take free modules to free modules, and hence flat modules to
flat modules (\hyperref[XVII.par5.5.1.2]{5.5.1.2}).

\paragraph{}\etiquette[5.5.2.5]{XVII.par5.5.2.5}
The $n^{\text{\rm ique}}$ map from $M$ to $TS^{n}(M)$ given by
$m\mapsto \bigotimes\limits^{n}m$ determines a
map $\varphi$, called canonical, from $\Gamma^{n}(M)$ to
$TS^{n}(M)$, satisfying
$$
\varphi(\gamma^{n}(m))=\bigotimes^{n}m{\quoi}.
$$
This map is an isomorphism when $M$ is finite free
(\textsc{Roby} \cite[IV \sisi{pg.272}{5 p.~272}]{sga4.exp.XVII.bib.Ro}), and hence also when $M$ is flat (\hyperref[XVII.par5.5.2.4]{5.5.2.4}
and \hyperref[XVII.par5.5.1.1]{5.5.1.1}).

Let $M$ be an $A$-module and $P$ an exact sequence
$$
P:P_{1}\xrightarrow{d} P_{0} \xrightarrow{\epsilon} M\longrightarrow
0{\quoi}.
$$
Let $P'_{1}=P_{1}\times P_{0}$, and let $j_{0}$ and $j_{1}$ be the
morphisms from $P'_{1}$ to $P_{0}$ with coordinates $(d,\Id)$ and
$(0,\Id)$. The sequence
$$
\xymatrix@C=1.75cm{
\Gamma^{n}(P'_1) \ar@<2pt>[r]^{\Gamma^{n}(j_0)} \ar@<-2pt>[r]_{\Gamma^{n}(j_1)}& \Gamma^{n}(P_0) \ar[r]& \Gamma^{n}(M) \ar[r]& 0
}
$$
is then exact (\textsc{Roby} \cite[\sisi{pg.284}{IV 10 p.~284}]{sga4.exp.XVII.bib.Ro}). \emph{When $P$ is a flat presentation
of} $M$,\pageoriginale this yields a description of the functor
$\Gamma^{n}$ in terms of the functor $TS^{n}$
$$\etiquette[5.5.2.5.1]{XVII.eq5.5.2.5.1}
TS^{n}(P'_1)\rightrightarrows TS^{n}(P_{0})\rightarrow
\Gamma^{n}(M)\rightarrow 0{\quoi}.\leqno(5.5.2.5.1)
$$

\paragraph{}\etiquette[5.5.2.6]{XVII.par5.5.2.6}
If $M=M'\oplus M''$, the $n^{\text{\rm ique}}$ map
\[
m\longmapsto (\gamma^{n'}m\otimes \gamma^{n''}m)_{n'+n''=n}
\]
defines an isomorphism, called canonical (\textsc{Roby} \cite[III \sisi{4}{9 p.~262}]{sga4.exp.XVII.bib.Ro})
$$\etiquette[5.5.2.6.1]{XVII.eq5.5.2.6.1}
\Gamma^{n}(M)\isomto \sisi{\sum}{\bigoplus}_{n'+n''=n} \Gamma^{n'}(M)\otimes \Gamma^{n''}(M){\quoi}.\leqno(5.5.2.6.1)
$$

If $\Sigma$ is an exact sequence
\[
\Sigma:0\longrightarrow M'\longrightarrow M\longrightarrow
M''\longrightarrow 0{\quoi},
\]
the module $\Gamma^{n}(M)$ is equipped with a canonical
(decreasing) filtration and morphisms
$$\etiquette[5.5.2.6.2]{XVII.eq5.5.2.6.2}
\Gamma^{i}(M'')\otimes \Gamma^{n-i}(M')\longrightarrow \Gr^{i}(\Gamma^{n}(M)){\quoi}.\leqno(5.5.2.6.2)
$$
These data are characterized by being functorial
in $\Sigma$ and satisfying the evident compatibility with \eqhyperref[XVII.eq5.5.2.6.1]{5.5.2.6.1}
when $\Sigma$ is split. The arrows \eqhyperref[XVII.eq5.5.2.6.2]{5.5.2.6.2} are
isomorphisms when $\Sigma$ is split, and hence also when $M''$ is flat
(\hyperref[XVII.par5.5.1.4]{5.5.1.4}).

\paragraph{}\etiquette[5.5.2.7]{XVII.par5.5.2.7}
The functor $\Gamma^{n}$ commutes with extension of scalars (\textsc{Roby} \cite[III \sisi{3 pg.262}{8 p.~262}]{sga4.exp.XVII.bib.Ro}) (the analogous property for $TS^{n}$ is false in
general).

\paragraph{}\etiquette[5.5.2.8]{XVII.par5.5.2.8}
If $u:M_{1}\times M_{2}\rightarrow N$ is a bilinear map,
then $\gamma^{n}\circ u$ is a bihomogeneous polynomial law of bidegree
$(n,n)$ and therefore determines (\hyperref[XVII.par5.5.2.3]{5.5.2.3}) $u^{n}:\Gamma^{n}(M_{1})\otimes
\Gamma^{n}(M_{2})\rightarrow \Gamma^{n}(N)$ satisfying
$u^{n}(\gamma^{n}x\otimes
\gamma^{n}y)=\gamma^{n}u(x,y)$. In\pageoriginale particular, if $B$
is an $A$-algebra, multiplication $B\times B\rightarrow B$ induces a
product on $\Gamma^{n}(B)$. This product makes $\Gamma^{n}(B)$ an
$A$-algebra, and the canonical map from $\Gamma^{n}(B)$ to
$TS^{n}(B)$ is an algebra homomorphism.

\subsubsection{}\etiquette[5.5.3]{XVII.subsubsec5.5.3}
Let $\Aa$ be an abelian category. We shall denote by
$\text{\rm ssimpl}^{+}(\Aa)$ the category of
co-semisimplicial objects\nde{In modern terminology they are called \emph{cosimplicial} objects.} in $\Aa$, by $C^{\geq 0}(\Aa)$ the category of
differential complexes $K$ with $K^{i}=0$ for $i<0$, and by $K^{\geq
  0}(\Aa)$ the image of $C^{\geq 0}(\Aa)$ in $K(\Aa)$. If $K\in \Ob
\text{ssimpl}^{+}(\Aa)$, the associated normalized complex
$N(K)\in \Ob C^{\geq 0}(\Aa)$ is defined by
$$\etiquette[5.5.3.1]{XVII.eq5.5.3.1}
\begin{cases}
N(K)^{n} =\text{the intersection of the kernels of the degeneracy
operators}\\
\qquad\qquad\text{from $K^{n}$ to $K^{n-1}$}\\
d=\sum(-1)^{i}\partial_{i}{\quoi}.
\end{cases}\leqno(5.5.3.1)
$$

If $K$ also denotes the complex with components $K^{n}$ and
differential $d=\sum(-1)^{i}\partial_{i}$, then $N(K)$ is a
direct-summand subcomplex of $K$, and $K/N(K)$ is homotopic to
zero.

By \textsc{Dold-Puppe} \cite{sga4.exp.XVII.bib.DP}, the functor $N$ is an equivalence of
categories
$$
N:\;\text{\rm ssimpl}^{+}(\Aa)\longrightarrow C^{\geq
0}(\Aa){\quoi}.
$$
This equivalence and its inverse take homotopic morphisms to
homotopic morphisms. Moreover, for every $K\in \Ob\text{\rm
ssimpl}^{+}(\Aa)$, $K^{n}$ is a direct sum of copies of $N(K)^{i}$
$(i\leq n)$; if $\PP$ is a property of objects of
$\Aa$, stable under finite direct sums and direct summands, it
follows that the $K^{i}$ have property $\PP$ for $i\leq n$ if
and only if the $N(K)^{i}$ have this property for $i\leq n$.

The equivalence\pageoriginale $N$ is used by \textsc{Dold} and
\textsc{Puppe} to
extend every functor (not necessarily additive)
$T:\Aa\rightarrow \Bb$ from $\Aa$ to an
arbitrary category $\Bb$ to a functor $C^{\geq 0}
(T)$ (or simply $T$) from $C^{\geq 0}(\Aa)$ to $C^{\geq
0}(\Bb)$
$$\etiquette[5.5.3.2]{XVII.eq5.5.3.2}
C^{\geq 0}(T)=NTN^{-1}{\quoi}.\leqno(5.5.3.2)
$$

The functor $C^{\geq 0}(T)$ passes to the quotient to define
$$\etiquette[5.5.3.3]{XVII.eq5.5.3.3}
T\quad\mbox{\rm or}\quad K^{\geq 0}(T):K^{\geq
0}(\Aa)\longrightarrow K^{\geq 0}(\Bb){\quoi}.\leqno(5.5.3.3)
$$

\setcounter{paragraph}{3}
\paragraph{}\etiquette[5.5.3.4]{XVII.par5.5.3.4}
If $T$ is \emph{of degree} $\leq n$, \ie if the $(n+1)$\textsuperscript{th} “cross-effect” vanishes (\textsc{Dold-Puppe} \cite[4.18 \sisi{pg.232}{p.~232}]{sga4.exp.XVII.bib.DP})\sisi{}{,}
then, for every $K\in \Ob C^{\geq 0}(\Aa)$ such that $K^{i}=0$
for $i>m$, one has $[T(K)]^{i}=0$ for $i>nm$ (\textsc{Dold-Puppe} \cite[4.2\sisi{.}{}3
\sisi{pg.233}{p.~233}]{sga4.exp.XVII.bib.DP}).

\paragraph{}\etiquette[5.5.3.5]{XVII.par5.5.3.5}
In particular, if $T$ is a finite-degree functor from the category
of $A$-modules to that of $B$-modules, which takes flat objects
to flat objects, then the image of a bounded flat complex is a
bounded flat complex.

\begin{enonce*}{Lemma 5.5.4}\etiquette[5.5.4]{XVII.lem5.5.4}
Let $\Aa$ be an abelian category.
\begin{enumerate}
\renewcommand{\theenumi}{\roman{enumi}}
\renewcommand{\labelenumi}{(\theenumi)}
\item The quasi-isomorphisms in $K^{\geq 0}(\Aa)$
admit a calculus of left and right fractions.

\item The evident functor is an equivalence between the category
obtained from $K^{\geq 0}(\Aa)$ by inverting the
quasi-isomorphisms and the category $D^{\geq 0}(\Aa)$.
\end{enumerate}
\end{enonce*}

The\pageoriginale functor $\tau_{\geq 0}:C(\Aa)\rightarrow
C^{\geq 0}(\Aa)$ (\hyperref[XVII.def1.1.13]{1.1.13}) passes to the quotient to define a
functor $\tau_{\geq 0}:K(\Aa)\rightarrow K^{\geq
0}(\Aa)$, left adjoint to the inclusion of $K^{\geq
0}(\Aa)$ in $\sisi{\Aa}{K(\Aa)}$. The functor $\tau_{\geq 0}$
takes quasi-isomorphisms to quasi-isomorphisms, and the adjunction arrow
$X\rightarrow \tau_{\geq 0}(X)$ is a quasi-isomorphism
if and only if $H^{i}(X)=0$ for $i<0$. Lemma
\hyperref[XVII.lem5.5.4]{5.5.4} therefore follows
from \textsc{Verdier} \cite{sga4.exp.XVII.bib.CD} and parts (i) and (ii) of the following lemma, whose
proof is left to the reader:

\begin{enonce*}{Lemma 5.5.4.1}\etiquette[5.5.4.1]{XVII.lem5.5.4.1}
Let $K$ be a category, $K'$ a full subcategory, $S$ a
set of morphisms of $K$ equal to the inverse image of the set of
isomorphisms of $K(S^{-1})$ (\cf \cite[\sisi{chap 1}{\chap~1}]{sga4.exp.XVII.bib.GZ}), and $S'$ its
trace in $K'$. Suppose that the inclusion $i$ of $K'$ in $K$ has
a left adjoint $\tau$, and that $\tau(S)\subset S'$.
\begin{enumerate}
\renewcommand{\theenumi}{\roman{enumi}}
\renewcommand{\labelenumi}{(\theenumi)}
\item The evident functor from $K'({S'}^{-1})$ to $K(S^{-1})$ is
fully faithful, and has $\tau$ as a left adjoint.

\item If $K(S^{-1})$ admits a calculus of left fractions
(\resp right fractions), then $K'({S'}^{-1})$ has the same property.

If the inclusion $i$ does not necessarily have a left adjoint, and if
$\tau$ denotes the \sisi{partially}{not everywhere} defined functor left adjoint to $i$,
one has:

\item If $K(S^{-1})$ admits a calculus of left fractions
(\resp right fractions) and if, for every $u\in S$ with target (\resp source)
in $K'$, $\tau(u)$ is defined and belongs to $S'$, then $K'({S'}^{-1})$
admits a calculus of left fractions (\resp right fractions), and the functor
$K'({S'}^{-1})\rightarrow K(S^{-1})$ is fully faithful.
\end{enumerate}
\end{enonce*}

\setcounter{subsubsection}{4}
\subsubsection{}\etiquette[5.5.5]{XVII.subsubsec5.5.5}
Let\pageoriginale $T:\Aa\rightarrow \Bb$ be a functor
between abelian categories. As in \hyperref[XVII.def1.2.1]{1.2.1}, one defines the derived functors of
$T$
\begin{align*}
LT &: D^{\geq 0}(\Aa)\longrightarrow \Pro(D^{\geq
0}(\Bb))\\
RT &: D^{\geq 0}(\Aa)\longrightarrow \Ind(D^{\geq 0}(\Bb))
\end{align*}
by the formulas
\begin{align*}
LT(K) &= \quotprojlim_{s:K'\rightarrow K}T(K')\\
RT(K) &= \quotinjlim_{s:K\rightarrow K'}T(K'){\quoi},
\end{align*}
where $s$ ranges over the filtered category of quasi-isomorphisms with target
(\resp source) $K$.

We say that $T$ is \emph{left (\resp right) derivable at $K$ if
$LT(K)$} $(RT(K))$ is representable. We say that $K$ is \emph{left
(\resp right) adapted} to $T$ (or, by abuse of language,
adapted to $LT$ (\resp $RT$)) if the following canonical arrow is an
isomorphism:
\[
LT(K)\isomto  TK\text{ (\resp $TK \isomto  RT(K)$)}{\quoi}.
\]

\begin{enonce*}{Proposition 5.5.6}\etiquette[5.5.6]{XVII.prop5.5.6}
Let $F$ be a functor from the category of modules over a ring $A$
to an abelian category $\Bb$. Suppose that filtered
inductive limits in $\Bb$ are representable and
exact, and that $F$ commutes with filtered
inductive limits. Then every bounded complex $M^{\bigdot}\in \Ob K^{\geq
0}(A)$ of flat modules is adapted to $LF$.
\end{enonce*}

The\pageoriginale proposition follows immediately from the definitions, the
following lemma, which will be proved in \hyperref[XVII.par5.5.6.4]{5.5.6.4}, and the fact that the
quasi-isomorphisms $u:N^{\bigdot}\rightarrow M^{\bigdot}$, with
$N^{\bigdot}\in \sisi{B}{C}^{b,\geq 0}(A)$ having flat components, form a
cofinal set in the category of quasi-isomorphisms with target
$M^{\bigdot}$.

\begin{enonce*}{Lemma 5.5.6.1}\etiquette[5.5.6.1]{XVII.lem5.5.6.1}
Let $u:N^{\bigdot}\rightarrow M^{\bigdot}$ be a quasi-isomorphism
between flat complexes that are objects of $C^{b,\geq 0}(A)$. Then $F(u)$ is
a quasi-isomorphism.
\end{enonce*}

\begin{enonce*}{Lemma 5.5.6.2}\etiquette[5.5.6.2]{XVII.lem5.5.6.2}
Let $\Aa$ be an additive category and $d$ an interval in
$\ZZ$ that is bounded above. Denote by $C^{d}$ (\resp
$C^{d,b}$, \resp $K^{d}$, \resp $K^{d,b}$) the category of complexes
that vanish in degrees $i\not\in d$ (\resp and are bounded, \resp up to
homotopy). There is an equivalence
$$
\sisi{\sheafell }{\ell}:\Ind(C^{d,b}(\Aa)\sisi{}{)}\isomto
C^{d}(\Ind\Aa){\quoi};
$$
the inverse of this equivalence induces a functor
$$
\sisi{\sheafell }{\ell}':K^{d}(\Ind\Aa)\longrightarrow
\Ind(K^{d,b}(\Aa)){\quoi}.
$$
\end{enonce*}

\begin{enumerate}
\renewcommand{\theenumi}{\Alph{enumi}}
\renewcommand{\labelenumi}{\theenumi)}
\item Let us prove that $\sisi{\sheafell }{\ell}$ is fully faithful. Thus, let $K$ be a
bounded complex of objects of $\Aa$ and $L=\quotinjlim
L_{\alpha}$ an \iInd-complex. One has
\begin{align*}
\Hom(K,L)\overset{\dfn}{=} \varinjlim \Hom(K,L_{\alpha})  = & \varinjlim
Z^{0}(\Hom^{\bigdot} (K,L_{\alpha}))\simeq Z^{0}
\varinjlim_{\alpha} \Hom^{\bigdot} (K,L_{\alpha})\\
& \isomto
Z^{0} \Hom^{\bigdot} (K,\ell L)=\Hom(K,\ell L){\quoi},
\end{align*}
the penultimate isomorphism following from the fact that $K$ is finite, so
that $\varinjlim$ commutes with passage from the double complex to the total
complex. This proves the assertion.

\item We\pageoriginale leave it to the reader to verify (rather laboriously) that
$\sisi{\sheafell }{\ell}$ is essentially surjective by constructing,
from left to right and starting in an arbitrarily chosen degree, “enough” bounded complexes mapping to a
bounded-above complex of \iInd-objects $K$; and likewise to develop the
following point:

\item If $\ell f:\ell \varinjlim K_{\alpha}\rightarrow \ell \varinjlim
L_{\beta}$ is null-homotopic, $f=dH+Hd$, then, for every $\alpha$,
$H$ lifts to $H^{i}_{\alpha}:K^{i}_{\alpha}\rightarrow
L^{i\sisi{}{-1}}_{\beta(\alpha)}$ which, for $\beta$ sufficiently large, defines a
null-homotopy of $f_{\alpha}:K_{\alpha}\rightarrow L_{\beta(\alpha)}$.
\end{enumerate}

\setcounter{subsubsection}{6}
\setcounter{paragraph}{2}
\paragraph{}\etiquette[5.5.6.3]{XVII.par5.5.6.3}
Let $A$ be a ring, $d$ an interval in $\ZZ$ bounded above,
$P(A)$ the category of finite projective $A$-modules,
and $\Pl(A)$ the category of flat $A$-modules. Denote by
$D^{\Tor\subset d}(A)$ the image of $C^{d}(\Pl(A))$ in $D(A)$ and by
$D^{\parf\subset d}(A)$ the image of $C^{d,b}(P(A))$ in
$D(A)$. By \hyperref[XVII.par5.5.1.2]{5.5.1.2}, one has $\Pl(A)\simeq \Ind
P(A)$. On the
other hand, $K^{d,b}(P(A))\simeq D^{\parf\subset d}(A)$. It then follows
from \hyperref[XVII.lem5.5.6.2]{5.5.6.2}, and from the fact that $H^{i}$ commutes with filtered
inductive
limits, that there exists a commutative diagram of functors
$$
\xymatrix@C=2.25cm{
C^{d}(\text{Pl}(A)) \ar[d]& \Ind C^{d,b}(P(A)) \ar[d]\ar[l]_{\substack{\varinjlim\\\sim}}\\
K^{d}(\text{Pl}(A)) \ar[d]\ar[r]^{\ell'}& \Ind K^{d,b}(P(A)) \ar[d]^{\sisi{}{\simeq}}\\
D^{\Tor\subset d}(A) \ar[r]^{\ell''}& \Ind D^{\parf \subset d}(A).
}
$$

% END INCLUDED SGA4_English_current_buildable_source_20260906/source/17/drafts/SGA4_XVII_LINES_5001_5400_English.texfrag

% BEGIN INCLUDED SGA4_English_current_buildable_source_20260906/source/17/drafts/SGA4_XVII_LINES_5401_5800_English.texfrag

\setcounter{subsubsection}{6}
\paragraph{}\etiquette[5.5.6.4]{XVII.par5.5.6.4}
Let us prove\pageoriginale \hyperref[XVII.lem5.5.6.1]{5.5.6.1}. With
the notation of \hyperref[XVII.par5.5.6.3]{5.5.6.3}, if $d$ is an
interval $[0,k]$, then (because $F$ commutes with $\varinjlim$) there
are commutative diagrams
$$
\xymatrix@C=2.5cm{
K^{d}(\text{Pl}(A))          \ar[r]^{\ell'} \ar[d]_{K^{\geq 0}(F)} & \Ind K^d(P(A))  \ar[d]_{\Ind(K^{\geq 0}(F))}             \\
K^{+}(\Bb) \ar[d]_{H^i}                        & \Ind K^{+}(\Bb)     \ar[d]_{\Ind(H^i)}       \\
\Bb                  & \Ind \Bb  \ar[l]_{\varinjlim}.
}
$$
If $u$ is as in \hyperref[XVII.lem5.5.6.1]{5.5.6.1}, then by
\hyperref[XVII.par5.5.6.3]{5.5.6.3}, $\ell'(u)$ is an isomorphism, and
$H^iK^{\geq0}(F)(u)$ is therefore an isomorphism, as required.

\subsubsection{}\etiquette[5.5.7]{XVII.subsubsec5.5.7}
Let $X$ be a scheme over a scheme $S$, and let $k$ be an integer.
Suppose that the symmetric power $\Sym^k_S(X)$ is defined. A reader
who wishes not to make this hypothesis may work in the category of
algebraic spaces (\textsc{Artin}\sisi{ [1]}{\nde{the original text
cites an unspecified work here, probably one of the two works of Artin
cited in the notes on pages \pageref{XVII.noteArtin1} and
\pageref{XVII.noteArtin2}.}}). Ring $S$ by a sheaf $\Aa$ of
\emph{commutative} rings. Beginning with
\hyperref[XVII.subsubsec5.5.19]{5.5.19}, $\Aa$ will be assumed to be
\emph{torsion}. By abuse of notation, we shall \emph{still denote by}
$\Aa$ the inverse image of $\Aa$ on any $S$-scheme.

Let $ℱ$ be a sheaf of modules on $X$; the sheaf
$\mathop{\boxtimes}\limits^kℱ$, the external tensor product of $k$
copies of $ℱ$, is a sheaf on $X^k/S$, equivariant under the action of
the symmetric group $\Sf_k$ on $X^k/S$. Passing\pageoriginale to
invariants, it defines a sheaf $TS^k_{\ext}(ℱ)$ on $\Sym^k_S(X)$,
called the \emph{external symmetric tensor power of $ℱ$}\sisi{}{.}

To reason ``point by point,'' it is useful to work with the points of
$\Sym^k_S(X)$ described as follows.

\paragraph{}\etiquette[5.5.7.1]{XVII.par5.5.7.1}
Let $s$ be a geometric point of $S$. A point $x$ of $\Sym^k_S(X)$
with values in $k(s)$ defines a geometric point of $\Sym^k_S(X)$. Let
$(x_1\ldots x_k)$ denote a point of $(X/S)^k$ with values in $k(s)$
whose image is $x$, let $y_1\ldots y_\ell$ be the distinct elements of
the sequence $x_1\ldots x_k$, and let $n_i$ $(1\leq i\leq\ell)$ be the
multiplicity of $y_i$ in this sequence $(\sum n_i=k)$.

With the notation of \hyperref[XVII.par5.5.7.1]{5.5.7.1}, and writing
$\Sf$ for a symmetric group, one has
$$\etiquette[5.5.7.2]{XVII.eq5.5.7.2}
TS^{k}_{\ext}(F)_x\simeq (\bigotimes^{\ell}_{1}
(\bigotimes^{n_i} F_{y_i}))^{\prod^{\ell}_{1}\Sf_{n_i}}{\quoi},\leqno(5.5.7.2)
$$
where the tensor products are taken over $\Aa_S$. In particular,
\emph{for $F$ flat}, by \hyperref[XVII.par5.5.2.5]{5.5.2.5} one has
$$\etiquette[5.5.7.3]{XVII.eq5.5.7.3}
TS^{k}_{\ext}(F)_x\isomfrom \bigotimes^{\ell}_{1}
\Gamma^{n_i}F_{y_i}{\quoi},\leqno(5.5.7.3)
$$
using induction on $\ell$, the flatness of the
$\Gamma^{n_i}F_{y_i}$ (\hyperref[XVII.par5.5.2.4]{5.5.2.4}), and the
following fact: if $G$ (\resp $H$) is a finite group acting on an
$A$-module $M$\sisi{}{ }(\resp $N$), and if $M$ and $N^H$ are flat,
then $M^G\otimes_A N^H\isomto(M\otimes_A N)^{G\times H}$.

\subsubsection{}\etiquette[5.5.8]{XVII.subsubsec5.5.8}
Let $F$ be a module on $X$, and let $P$ be a flat presentation of $F$:
$$
P:P_1\longrightarrow P_0\longrightarrow F\longrightarrow0{\quoi}.
$$
Define\pageoriginale $P'_1$, $j_0$, and $j_1$ as in
\hyperref[XVII.par5.5.2.5]{5.5.2.5}, and define
$\Gamma^k_{\ext}(F)$ (the \emph{external $\Gamma$ functor}) to be the
cokernel of the pair of arrows
$(TS^k_{\ext}(j_0),TS^k_{\ext}(j_1))$. With the notation of
\hyperref[XVII.par5.5.7.1]{5.5.7.1}, one deduces from
\eqhyperref[XVII.eq5.5.7.3]{5.5.7.3} and
\hyperref[XVII.par5.5.2.5]{5.5.2.5} isomorphisms
$$\etiquette[5.5.8.1]{XVII.eq5.5.8.1}
\Gamma^k_{\ext}(F)_x\simeq\bigotimes^{\ell}_{1}
\Gamma^{n_i}(F_{y_i}){\quoi}.\leqno(5.5.8.1)
$$

This formula shows that $\Gamma^k_{\ext}(F)$, which a priori depends
functorially on $P$, depends, up to a unique isomorphism, only on $F$.

\subsubsection{}\etiquette[5.5.9]{XVII.subsubsec5.5.9}
There is a functorial morphism
$$
\alpha:\Gamma^k_{\ext}(F)\longrightarrow TS^k_{\ext}(F){\quoi}.
$$
The pair $(\Gamma,\alpha)$ is characterized by the following two
properties:
\begin{enumerate}
\renewcommand{\theenumi}{\roman{enumi}}
\renewcommand{\labelenumi}{(\theenumi)}
\item For $F$ flat, $\alpha$ is an isomorphism.

\item With the notation of \hyperref[XVII.par5.5.2.5]{5.5.2.5}, for
every presentation
$$
P_1\xrightarrow{d}P_0\xrightarrow{\epsilon}F\longrightarrow0
$$
of $F$, the sequence
$$
\Gamma^k_{\ext}(P'_1)\rightrightarrows
\Gamma^k_{\ext}(P_0)\rightarrow\Gamma^k_{\ext}(F)\rightarrow0
$$
is exact.

Moreover:

\item The functors $TS^k$ and $\Gamma^k$ commute with filtered direct
limits (by
\sisi{\hyperref[XVII.eq5.5.7.2]{5.5.7.2}}{\eqhyperref[XVII.eq5.5.7.2]{5.5.7.2}},
\sisi{\hyperref[XVII.eq5.5.8.1]{5.5.8.1}}{\eqhyperref[XVII.eq5.5.8.1]{5.5.8.1}},
\hyperref[XVII.par5.5.2.4]{5.5.2.4}).

\item The functor $\Gamma^k$ commutes with every extension
$\Aa\rightarrow\Aa'$ of the base ring.

\item The\pageoriginale functor $\Gamma^k$ commutes with every base
change $S'\rightarrow S$: for every cartesian square
$$
\xymatrix@C=2.5cm@R=1.5cm{
X' \ar[r]^{g'}\ar[d]^{f'}& X \ar[d]^{f}\\
S' \ar[r]^{g}& S,
}
$$
one has
$$
\Sym^k_g(g')^{\ast}\Gamma^k_{\ext}\simeq
\Gamma^k_{\ext}{g'}^{\ast}{\quoi}.
$$
(\sisi{\hyperref[XVII.eq5.5.8.1]{5.5.8.1}}{\eqhyperref[XVII.eq5.5.8.1]{5.5.8.1}}
and \hyperref[XVII.par5.5.2.7]{5.5.2.7}).

\item The functor $\Gamma^k$ commutes with inverse images under every
\emph{radicial} morphism $i:X'\rightarrow X$ (for example a closed
immersion): \sisi{O}{o}ne has
$$
\Sym^k_S(i)^{\ast}\Gamma^k_{\ext}\simeq
\Gamma^k_{\ext}i^{\ast}
$$
\eqhyperref[XVII.eq5.5.8.1]{5.5.8.1}.

\item For an \emph{immersion} $j:X'\rightarrow X$, one has
$$
\Sym^k_S(j)_{!}\Gamma^k_{\ext}\simeq\Gamma^k_{\ext}j_{!}
$$
\eqhyperref[XVII.eq5.5.8.1]{5.5.8.1}.
\end{enumerate}

\subsubsection{}\etiquette[5.5.10]{XVII.subsubsec5.5.10}

Let $k_1,k_2$, and $k=k_1+k_2$ be three integers $\geq0$. There is
then a finite morphism (\quad):\nde{I have not found references for the
finiteness of $\vee_{k_1k_2}$ in general.}
$$
\vee_{k_1k_2}:\Sym^{k_1}_S(X)\times_S\Sym^{k_2}_S(X)
\longrightarrow\Sym^k_S(X){\quoi}.
$$
If $F_i$ $(i=1,2)$ is a module on $\Sym_S^{k_i}(X)$, denote by
$F_1\vee F_2$ the sheaf on $\Sym_S^k(X)$
$$
F_1\vee F_2=\vee_{k_1k_2\ast}(F_1\boxtimes F_2){\quoi}.
$$

\subsubsection{}\etiquette[5.5.11]{XVII.subsubsec5.5.11}
Let\pageoriginale there be an exact sequence of modules on $X$
\[
\Sigma:0\longrightarrow F'\longrightarrow F\longrightarrow
F''\longrightarrow0{\quoi}.
\]
With the notation of \hyperref[XVII.par5.5.7.1]{5.5.7.1}, by
\eqhyperref[XVII.eq5.5.8.1]{5.5.8.1} and
\hyperref[XVII.par5.5.2.6]{5.5.2.6} there is a filtration on
$\Gamma^k_{\ext}(F)_x$, and there are morphisms
\[
\sum_{\sum^{\ell}_{1}i_\alpha=i}
\bigotimes^{\ell}_{1}\Gamma^{i_\alpha}(F'_{y_\alpha})\otimes
\bigotimes^{\ell}_{1}\Gamma^{n_i-i_\alpha}(F''_{y_\alpha})
\longrightarrow\Gr^i(\Gamma^k_{\ext}(F)_x){\quoi}.
\]
The first member is precisely
$(\Gamma^i_{\ext}(F')\vee\Gamma^{k-i}_{\ext}(F''))_x$. One readily
checks that these data, constructed point by point, arise from a
(unique) filtration with $k+1$ steps on $\Gamma^k_{\ext}(F)$ and from
(unique) morphisms
$$\etiquette[5.5.11.1]{XVII.eq5.5.11.1}
\Gamma^i_{\ext}(F')\vee\Gamma^{k-i}_{\ext}(F'')
\longrightarrow\Gr^i(\Gamma^k_{\ext}(F)){\quoi}.\leqno(5.5.11.1)
$$

If the exact sequence $\Sigma$ is split, one has more precisely
(verification point by point from
\sisi{\hyperref[XVII.eq5.5.2.6.2]{5.5.2.6.2}}{\eqhyperref[XVII.eq5.5.2.6.1]{5.5.2.6.1}}):
$$\etiquette[5.5.11.2]{XVII.eq5.5.11.2}
\Gamma^k_{\ext}(F'\oplus F'')\simeq
\sisi{\sum}{\bigoplus}_{k_1+k_2=k}\Gamma^{k_1}_{\ext}(F')\vee
\Gamma^{k_2}_{\ext}(F''){\quoi}.\leqno(5.5.11.2)
$$

If the exact sequence $\Sigma$ is split point by point, or if $F''$ is
flat, the morphisms \eqhyperref[XVII.eq5.5.11.1]{5.5.11.1} are
isomorphisms (\cf \hyperref[XVII.par5.5.1.4]{5.5.1.4}).

\subsubsection{}\etiquette[5.5.12]{XVII.subsubsec5.5.12}
Let $u:X\rightarrow Y$ be a morphism of $S$-schemes. Suppose that the
symmetric powers $\Sym^k_S(X)$ and $\Sym^k_S(Y)$ exist. Let $F$ be a
module on $X$, let $G$ be a module on $Y$, and let
$v\in\Hom(u^{\ast}G,F)\underset{\sisi{n}{}}{\simeq}\Hom(G,u_{\ast}F)$
be a $u$-homomorphism from $F$ to $G$ \hyperlink{sga4.exp.XVII.subsection.0.2}{(0.2)}. If $\prod^n u$ is the
$n$\textsuperscript{th} cartesian power of $u$, the homomorphism $v$
induces a $\prod^n u$-homomorphism $\mathop{\boxtimes}\limits^n v$
from $\mathop{\boxtimes}\limits^nF$ to
$\mathop{\boxtimes}\limits^nG$. Passing\pageoriginale to invariants,
the latter defines
\[
TS^n(v):TS^n_{\ext}(F)\xrightarrow[\Sym_S^n(u)]{}TS^n_{\ext}(G){\quoi}.
\]

It follows from Definition \hyperref[XVII.subsubsec5.5.8]{5.5.8} that
there exists a unique $\Sym^n(u)$-homomorphism
$\Gamma^n_{\ext}(v)$, functorial in $(F,G,v)$,
$$\etiquette[5.5.12.0]{XVII.eq5.5.12.0}
\Gamma^n_{\ext}(v):\Gamma^n_{\ext}(F)
\xrightarrow[\Sym^n_S(u)]{}\Gamma^n_{\ext}(G)\leqno(5.5.12.0)
$$
making the following diagram commutative (\cf \hyperlink{sga4.exp.XVII.subsection.0.2}{(0.2)} for the unusual
direction of the arrows):
$$
\xymatrix@C=3.5cm@R=2cm{
\Gamma^n_{\ext}(F) \ar[r]^{\Gamma^n_{\ext}(v)}_{\text{over $\Sym^n_S(u)$}}& \Gamma^n_{\ext}(G)\\
TS^n_{\ext}(F) \ar[r]^{TS^n_{\ext}(v)}_{\text{over $\Sym^n_S(u)$}}\ar[u]^{\alpha}_{\text{over $\Id$}} & TS^n_{\ext}(G)\ar[u]^{\alpha}_{\text{over $\Id$}}.
}
$$

We shall use this arrow mainly as a ``symmetric Künneth arrow'': for
$G=u_{\ast}F$ and $v$ the canonical arrow, the arrow
$\Gamma^n_{\ext}(v)$ identifies with a homomorphism\sisi{s}{} of
modules
$$\etiquette[5.5.12.1]{XVII.eq5.5.12.1}
k^n:\Gamma^n_{\ext}(u_{\ast}F)\longrightarrow
\Sym^n(u)_{\ast}\Gamma^n_{\ext}(F){\quoi}.\leqno(5.5.12.1)
$$

\subsubsection{}\etiquette[5.5.13]{XVII.subsubsec5.5.13}
Consider the condition:

\paragraph{}\etiquette[5.5.13.1]{XVII.par5.5.13.1}
$K\in\Ob K^{b,\geq0}(X,\Aa)$ and, for every geometric point $x$ of
$X$, the complex $K_x$ is homotopy equivalent to a complex of flat
modules $L_x\in\Ob C^{b,\geq0}(\Aa_x)$.

\begin{enonce*}{Lemma 5.5.13.2}\etiquette[5.5.13.2]{XVII.lem5.5.13.2}
The\pageoriginale subcategory $D^{b,\tor\leq0}(X,\Aa)$ of
$D^b(X,\Aa)$ consisting of complexes of \tTor-dimension $\leq0$
\emph{\hyperref[XVII.def4.1.9]{4.1.9}} is equivalent to the category
obtained from the category of complexes
$K\in\Ob K^{b,\geq0}(X,\Aa)$ satisfying
\emph{(\hyperref[XVII.par5.5.13.1]{5.5.13.1})} by inverting
quasi-isomorphisms. This construction may be performed indifferently
by a calculus of left or right fractions.
\end{enonce*}

Let $P$ be the subcategory of $K^{-}(X,\Aa)$ consisting of complexes
that are pointwise homotopy equivalent to a complex of flat modules.
It is known (\textsc{Verdier}
\sisi{\cite{sga4.exp.XVII.bib.CD}}{\cite[I 2.4.2]{sga4.exp.XVII.bib.CD}}) that $D^{-}(X,\Aa)$ is obtained
from $P$ by inverting quasi-isomorphisms, using a calculus of right or
left fractions. Apply \hyperref[XVII.lem5.5.4.1]{5.5.4.1} to the
inclusion of $K^{b,\geq0}(X,\Aa)\cap P$ into $P$. Indeed, if $K$ has
\tTor-dimension $\leq0$ and is pointwise homotopy equivalent to
complexes of flat modules $L_x$, then $\tau_{\geq0}(K)$ is pointwise
homotopy equivalent to the $\tau_{\geq0}(L_x)$, and the latter are
still complexes of flat modules, so $\tau_{\geq0}(K)$ satisfies
\hyperref[XVII.par5.5.13.1]{5.5.13.1}.

\subsubsection{}\etiquette[5.5.14]{XVII.subsubsec5.5.14}
By \hyperref[XVII.lem5.5.6.1]{5.5.6.1},
\sisi{\hyperref[XVII.eq5.5.7.2]{5.5.7.2}}{\eqhyperref[XVII.eq5.5.7.2]{5.5.7.2}}
or
\sisi{\hyperref[XVII.eq5.5.8.1]{5.5.8.1}}{\eqhyperref[XVII.eq5.5.8.1]{5.5.8.1}}
and \hyperref[XVII.par5.5.2.4]{5.5.2.4}, every quasi-isomorphism
$v:K'\rightarrow K$ between complexes satisfying
(\hyperref[XVII.par5.5.13.1]{5.5.13.1}) induces quasi-isomorphisms
\begin{align*}
TS^k(v)&:TS^k_{\ext}(K')\xrightarrow{\approx}TS^k_{\ext}(K)\\
\Gamma^k(v)&:\Gamma^k_{\ext}(K')\xrightarrow{\approx}\Gamma^k_{\ext}(K){\quoi}.
\end{align*}
It follows that the functors $\Gamma^k$ and $TS^k$ are derivable
(\hyperref[XVII.subsubsec5.5.5]{5.5.5}) on the subcategory
$D^{b,\tor\leq0}(X,\Aa)$ of $D^b(X,\Aa)$, and that complexes
satisfying \hyperref[XVII.par5.5.13.1]{5.5.13.1} are left adapted to
these functors. By \hyperref[XVII.par5.5.2.4]{5.5.2.4},\pageoriginale
$\Gamma^k$ and $TS^k$ induce functors
\begin{align*}
LTS^k &:D^{b,\tor\leq0}(X,\Aa)\longrightarrow
D^{b,\tor\leq0}(\Sym^k_S(X),\Aa)\tag{5.5.14.1}\\
L\Gamma^k &:D^{b,\tor\leq0}(X,\Aa)\longrightarrow
D^{b,\tor\leq0}(\Sym^k_S(X),\Aa){\quoi}.\tag{5.5.14.2}
\end{align*}\etiquette[5.5.14.1]{XVII.eq5.5.14.1}
\etiquette[5.5.14.2]{XVII.eq5.5.14.2}%Removing this comment would introduce unwanted white space.
Moreover, by \hyperref[XVII.subsubsec5.5.9]{5.5.9}~(i), the canonical
morphism $\alpha$ (\hyperref[XVII.subsubsec5.5.9]{5.5.9}) induces an
isomorphism between the functors \eqhyperref[XVII.eq5.5.14.1]{5.5.14.1}
and \eqhyperref[XVII.eq5.5.14.2]{5.5.14.2}:
$$\etiquette[5.5.14.3]{XVII.eq5.5.14.3}
\alpha:L\Gamma^k\isomto LTS^k{\quoi}.\leqno(5.5.14.3)
$$

\setcounter{paragraph}{3}
\paragraph{}\etiquette[5.5.14.4]{XVII.par5.5.14.4}
The isomorphisms of \hyperref[XVII.subsubsec5.5.9]{5.5.9}~(iv), (v),
(vi), and (vii) pass to the derived category and define isomorphisms:
\begin{itemize}
\item[${\rm(iv')}$] For an $\Aa$-algebra $\Aa'$, one has
$$
L\Gamma^k_{\ext_{\Aa}}(K)\overset{\LL}{\otimes}\Aa'\simeq
L\Gamma^k_{\ext_{\Aa'}}(K\overset{\LL}{\underset{\Aa}{\otimes}}\Aa'){\quoi}.
$$

\item[${\rm(v')}$] For every base change $g:S'\rightarrow S$, one has
$$
\Sym^k_g(g')^{\ast}L\Gamma^k_{\ext}K\simeq
L\Gamma^k_{\ext}({g'}^{\ast}K){\quoi}.
$$

\item[${\rm(vi')}$] For every \emph{radicial} morphism
$i:X'\rightarrow X$, one has
$$
\Sym^k_S(i)^{\ast}L\Gamma^k_{\ext}K\simeq
L\Gamma^k_{\ext}i^{\ast}K{\quoi}.
$$

\item[${\rm(vii')}$] For an \emph{immersion} $j:X'\rightarrow X$,
one has
$$
\Sym^k_S(j)_{!}L\Gamma^k_{\ext}\simeq
L\Gamma^k_{\ext}j_{!}{\quoi}.
$$
\end{itemize}

\subsubsection{}\etiquette[5.5.15]{XVII.subsubsec5.5.15}
Let there be an exact sequence of flat complexes in
$C^{\geq0}(X,\Aa)$,
$$
0\longrightarrow K'\longrightarrow K\longrightarrow K''\longrightarrow0{\quoi},
$$
and apply the functor $\Gamma^k_{\ext}$
(\hyperref[XVII.subsubsec5.5.3]{5.5.3}). One finds (as in
\hyperref[XVII.subsubsec5.5.11]{5.5.11}) that the semisimplicial module
$\Gamma^k_{\ext}(N^{-1}K)$ on $\Sym^k_S(X)$ has a filtration
whose\pageoriginale successive quotients are the semisimplicial
modules $\Gamma^{k'}_{\ext}(N^{-1}K')\vee
\Gamma^{k''}_{\ext}(N^{-1}K'')$ $(k'+k''=k)$.

Passing to the derived category, one finds that
$L\Gamma^k_{\ext}(K)$ is represented by a filtered complex whose
successive quotients represent the objects
$L\Gamma^i_{\ext}(K'\sisi{}{)}\overset{\LL}{\vee}
L\Gamma^{k-i}_{\ext}(K'')$ of the derived category (apply
\hyperref[XVII.par5.5.2.4]{5.5.2.4} and the
\textsc{Eilenberg--Zilber--Cartier} theorem (isomorphism in
$\mathcal K$ between a bisemisimplicial object and its diagonal:
\textsc{Dold\sisi{ }{-}Puppe} \cite[2.9
\sisi{pg.213}{p.~213}]{sga4.exp.XVII.bib.DP}).

In the derived category, this construction can already be made from
an exact sequence of complexes satisfying
\hyperref[XVII.par5.5.13.1]{5.5.13.1}.

${}^{\ast}$ In more intrinsic terms, let $\Delta$ be a ``\emph{true
triangle}'' of objects of $D^{b,\tor\leq0}(X,\Aa)$, \ie an object of
the ``filtered derived category'' with $\Gr^i(\Delta)=0$ for
$i\neq0,1$ and $\Gr^i(\Delta)\in\Ob D^{b,\tor\leq0}(X,\Aa)$
(\textsc{Deligne}, in the Appendix to \textsc{Verdier} \cite{sga4.exp.XVII.bib.V2}, or
\refvb[VI B nᵒ]{Vbis.subsubsec2.5.1},
\refnoexpose{Vbis.subsubsec2.5.2}).

The preceding construction then gives an object
$L\Gamma^k_{\ext}(\Delta)$ of the filtered derived category of
$\Sym^k_S(X)$, isomorphisms
\[
\Gr^i(L\Gamma^k_{\ext}(\Delta))\simeq
L\Gamma^i_{\ext}(\Gr^1(\Delta))\overset{\LL}{\vee}
L\Gamma^{k-i}_{\ext}(\Gr^0(\Delta)\sisi{}{)}
\]
and, denoting by $\lvert\quad\rvert$ the underlying complex of a
filtered complex, an isomorphism
\[
\lvert L\Gamma^k_{\ext}(\Delta)\rvert\simeq
L\Gamma^k_{\ext}(\lvert\Delta\rvert){\quoi}.{}_{\ast}
\]

\begin{enonce*}{Lemma 5.5.16}\etiquette[5.5.16]{XVII.lem5.5.16}
Let $N$ be an integer. The category $D^{b,\tor\leq0}(X,\Aa)$ is
obtained from the category of complexes
$K\in\Ob K^{b,\geq0}(X,\Aa)$ satisfying conditions {\rm(i)} and
{\rm(ii)} below by inverting quasi-isomorphisms:
\begin{enumerate}
\renewcommand{\theenumi}{\roman{enumi}}
\renewcommand{\labelenumi}{(\theenumi)}
\item condition\pageoriginale
(\hyperref[XVII.par5.5.13.1]{5.5.13.1}).

\item The complex $K$ is adapted to every direct-image functor of
cohomological dimension $\leq N$, relative to an isomorphism of topoi
$f:\widetilde X_{\text{\rm et}}\rightarrow T$.
\end{enumerate}
\end{enonce*}

The existence of a calculus of right fractions follows from
\hyperref[XVII.subsubsec5.5.13]{5.5.13} and from the fact that, by
truncated canonical flasque resolutions
(\hyperref[XVII.subsubsec4.2.9]{4.2.9}), every complex $K$ satisfying
(i) maps to a complex satisfying (i) and (ii).

\setcounter{subsubsection}{16}


% END INCLUDED SGA4_English_current_buildable_source_20260906/source/17/drafts/SGA4_XVII_LINES_5401_5800_English.texfrag

% BEGIN INCLUDED SGA4_English_current_buildable_source_20260906/source/17/drafts/SGA4_XVII_LINES_5801_6200_English.texfrag
\subsubsection{}\etiquette[5.5.17]{XVII.subsubsec5.5.17}
Let $u:X\rightarrow Y$ be a morphism of $S$-schemes. Suppose that
$\Sym^n_S(X)$ and $\Sym^n_S(Y)$ exist, that there is an integer $N$
such that the functor $u_{\ast}$ has dimension $\leq N$ on the category
of $\Aa$-modules, and that the functor $Ru_{\ast}$ carries bounded
complexes of \tTor-dimension $\leq0$ to complexes of \tTor-dimension
$\leq0$. The latter two conditions hold automatically if $u$ is
proper, $\Aa$ is torsion, and $Y$ is coherent (\ie quasi-compact and
quasi-separated) (\hyperref[XVII.cor5.2.8.1]{5.2.8.1},
\hyperref[XVII.thm5.2.10]{5.2.10}).\nde{the coherence of $Y$ may be
replaced by the weaker condition that the dimensions of the fibers of
$u$ are bounded.}

Let $K\in\Ob K^{b,\geq0}(X,\Aa)$ satisfy conditions (i) and (ii) of
\hyperref[XVII.lem5.5.16]{5.5.16}, with $N$ as above. By
\hyperref[XVII.subsubsec5.5.14]{5.5.14}, the arrow
$$
\beta:L\Gamma^n_{\ext}(K)\isomto\Gamma^n_{\ext}(K)
$$
is a quasi-isomorphism. Similarly, the arrow
$$
\gamma:u_{\ast}K\isomto Ru_{\ast}K
$$
is a quasi-isomorphism.

Define\pageoriginale the \emph{symmetric Künneth arrow} from
$L\Gamma^n_{\ext}Ru_{\ast}(K)$ to
$R\Sym^n(u)_{\ast}L\Gamma^n_{\ext}(K)$ to be the composite
\begin{multline*}
k^{n}: L\Gamma^{n}_{\ext}Ru_{\ast}K\xleftarrow[\gamma]{\sim}
L\Gamma^{n}_{\ext}u_{\ast}K\rightarrow
\Gamma^{n}_{\ext}u_{\ast}K\xrightarrow[\protect\hyperlink{sga4.semantic.XVII.eq5.5.12.1}{5.5.12.1}]{}\Sym^{n}_{S}(u)_{\ast}\Gamma^{n}_{\ext}(K)\\
\longrightarrow R
\Sym^{n}_{S}(u)_{\ast}\Gamma^{n}_{\ext}(K)\xleftarrow[\beta]{\sim} R \Sym^{n}_{S}(u)_{\ast}L\Gamma^{n}_{\ext}(K){\quoi}.\tag{5.5.17.1}
\end{multline*}\etiquette[5.5.17.1]{XVII.eq5.5.17.1}%The vertical space following this paragraph is too large.

It follows from \hyperref[XVII.lem5.5.16]{5.5.16} that the morphisms
\eqhyperref[XVII.eq5.5.17.1]{5.5.17.1} define a morphism between
functors from $D^{b,\tor\leq0}(X,\Aa)$ to
$D^{+}(\Sym^n_S(Y),\Aa)$:
$$\etiquette[5.5.17.2]{XVII.eq5.5.17.2}
k^n:L\Gamma^n_{\ext}Ru_{\ast}\longrightarrow
R\Sym^n_S(u)_{\ast}L\Gamma^n_{\ext}{\quoi},\leqno(5.5.17.2)
$$
also called the \emph{symmetric Künneth morphism}.

This morphism is characterized by Proposition
\hyperref[XVII.prop5.5.18]{5.5.18} below.

\setcounter{paragraph}{2}
\paragraph{}\etiquette[5.5.17.3]{XVII.par5.5.17.3}
Let $S$, $X$, $Y$, and $u$ be as above, let
$K\in\Ob C^{b,\geq0}(X,\Aa)$ be a complex of \tTor-dimension
$\leq0$, let $K_1\in\Ob C^{b,\geq0}(Y,\Aa)$, and let
$K_2\in\Ob C^b(\Sym^n_S(X),\Aa)$. Let $v$ be a $u$-homomorphism
\hyperlink{sga4.exp.XVII.subsection.0.2}{(0.2)} from $K$ to $K_1$, and let $t$ be a homomorphism (in the usual
sense, as are all those that follow) from $\Gamma^n_{\ext}(K)$ to
$K_2$. There are morphisms
\begin{enumerate}
\renewcommand{\theenumi}{\alph{enumi}}
\renewcommand{\labelenumi}{\theenumi)}
\itemsep=8pt
\item $\Gamma^n_{\ext}(v):\Gamma^n_{\ext}(K_1)\rightarrow
\Sym^n_S(u)_{\ast}\Gamma^n_{\ext}(K)$\qquad
\eqhyperref[XVII.eq5.5.12.0]{5.5.12.0}

\item
$\Sym^n_S(u)_{\ast}(t):\Sym^n_S(u)_{\ast}\Gamma^n_{\ext}(K)
\rightarrow\Sym^n_S(u)_{\ast}K_2$,
\end{enumerate}
whose composite
$$
\varphi':\Gamma^n_{\ext}(K_1)\longrightarrow
\Sym^n_S(u)_{\ast}K_2{\quoi},
$$
defines, by composition,
$$
\varphi:L\Gamma^n_{\ext}(K_1)\longrightarrow
R\Sym^n(u)_{\ast}K_2{\quoi}.
$$
Moreover,\pageoriginale $v$ and $t$ induce
$$
\begin{cases}
v':K_1\longrightarrow Ru_{\ast}K\\
t':L\Gamma^n_{\ext}K\longrightarrow K_2{\quoi}.
\end{cases}
$$

Then:

\begin{enonce*}{Proposition 5.5.18}\etiquette[5.5.18]{XVII.prop5.5.18}
Under the hypotheses of \hyperref[XVII.par5.5.17.3]{5.5.17.3}, the
following diagram is commutative in $D^{+}(\Sym^n_S(Y),\Aa)$:
\[
\xymatrix@C=4cm@R=1.75cm{
  L\Gamma^n_{\ext}Ru_{\ast}K \ar[r]^{k^n}_{\text{\textsuperscript{(symmetric Künneth)}}}& R\Sym^n_S(u)_{\ast}L
\Gamma^n_{\ext}K \ar[d]_{R
  \Sym^n_S(u)_{\ast}(t')}\\
L\Gamma^n_{\ext}K_1 \ar[r]^{\varphi}\ar[u]^{L\Gamma^n_{\ext}(v')} & R\Sym^n_S(u)_{\ast}K_2.
}
\]
\end{enonce*}

The proof is entirely similar to that of
\hyperref[XVII.par5.4.1.5]{5.4.1.5} and is left to the reader.

\setcounter{subsubsection}{18}
\subsubsection{Henceforth $\Aa$ is assumed to be
  torsion}\etiquette[5.5.19]{XVII.subsubsec5.5.19}

Let $u:X\rightarrow Y$ be a compactifiable morphism of $S$-schemes:
\[
\xymatrix@C=3cm@R=1.75cm{
   & \bar{X} \ar[rd]^{\bar{u}}&\\
X \ar@{^{(}->}[ru]^{j} \ar[rr]^{u} \ar[rd]&         & Y \ar[ld]\\
   &   S     & .
}
\]
Suppose\pageoriginale that $\Sym^n_S(X)$,
$\Sym^n_S(\bar X)$, and $\Sym^n_S(Y)$ exist. For every module $F$ on
$X$, one then has
$$\etiquette[5.5.19.1]{XVII.eq5.5.19.1}
\Gamma^n_{\ext}(j_{!}F)\simeq
\Sym^n_S(j)_{!}\Gamma^n_{\ext}(F){\quoi},\leqno(5.5.19.1)
$$
and this isomorphism passes immediately to the derived categories
\sisi{}{(\hyperref[XVII.par5.5.14.4]{5.5.14.4}~(vii'))}. Define the
\emph{symmetric Künneth morphism} between functors from
$D^{b,\tor\leq0}(X,\Aa)$ to
$D^{b,\tor\leq0}(\Sym^n_S(Y),\Aa)$\nde{How can one show that
$\Sym^n_S(u)$ is of finite type in general? This condition is
necessary for $R\Sym^n_S(u)_!$ to be defined.}
$$\etiquette[5.5.19.2]{XVII.eq5.5.19.2}
k^n_u\quad\mbox{\rm or}\quad
k^n:L\Gamma^n_{\ext}Ru_{!}\longrightarrow
R\Sym^n_S(u)_{!}L\Gamma^n_{\ext}\leqno(5.5.19.2)
$$
to be the composite
\begin{align*}
L\Gamma^n_{\ext}Ru_{!}\stackrel{\dfn}{\simeq}L\Gamma^n_{\ext}R\bar u_{\ast}j_{!}
&\xrightarrow{\eqhyperref[XVII.eq5.5.17.2]{5.5.17.2}}
R\Sym^n_S(\bar u)_{\ast}L\Gamma^n_{\ext}j_!\\
&\xrightarrow[\sisi{}{\text{(\hyperref[XVII.par5.5.14.4]{5.5.14.4} (vii'))}}]{\sisi{\substack{\eqhyperref[XVII.eq5.5.19.1]{5.5.19.1}\\\sim}}{\sim}}
R\Sym^n_S(\bar u)_*\Sym^n_S(j)_{!}L\Gamma^n_{\ext}\\
&\xrightarrow{\sisi{\text{(5.5.13.5 (vii'))}}{}}
R\Sym^n_S(u)_{!}L\Gamma^n_{\ext}.
\end{align*}

\setcounter{subsubsection}{19}
\subsubsection{}\etiquette[5.5.20]{XVII.subsubsec5.5.20}
We leave to the reader the following verifications, which are easy in
principle using \hyperref[XVII.prop5.5.18]{5.5.18}.

\paragraph{}\etiquette[5.5.20.1]{XVII.par5.5.20.1}
(Intrinsic character). The morphism
\eqhyperref[XVII.eq5.5.19.2]{5.5.19.2} is independent of the choice of
the compactification $\bar u$ of $u$.

\paragraph{}\etiquette[5.5.20.2]{XVII.par5.5.20.2}
(Transitivity). Let there be a diagram of $S$-schemes
$$
X\xrightarrow{u}Y\xrightarrow{v}Z{\quoi},
$$
with\pageoriginale $u$ and $v$ quasi-projective; suppose that
$\Sym^n_S(Z)$ exists (and hence also $\Sym^n_S(Y)$ and $\Sym^n_S(X)$
(\SGA 1 V 1.7, \EGA II 4.5.4)). The following diagram of functors is
then commutative:
$$
\xymatrix@C=1.5cm@R=1.75cm{
L\Gamma^n_{\ext}R(vu)_{!} \ar@{=}[d] \ar[rr]^{k^n_{vu}}& & R\Sym^n_S(vu)_{!}L\Gamma^n_{\ext}\ar@{=}[d]\\
L\Gamma^n_{\ext}Rv_{!}Ru_{!} \ar[r]^{k^n_v}& R\Sym^n_S(v)_!
L\Gamma^n_{\ext}Ru_{!} \ar[r]^{k^n_u}& R\Sym^n_S(v)_{!}R\Sym^n_S(u)_{!}L\Gamma^n_{\ext}.
}
$$

\paragraph{}\etiquette[5.5.20.3]{XVII.par5.5.20.3}
The symmetric Künneth morphism is compatible with changes of base or
space of the types \hyperref[XVII.subsubsec5.5.9]{5.5.9}~(iv), (v), and (vi). More precisely, one has
the following compatibilities with the morphisms
\sisi{5.5.13.5}{\hyperref[XVII.par5.5.14.4]{5.5.14.4}}
${\rm(iv')}$, ${\rm(v')}$, and ${\rm(vi')}$.
\begin{itemize}
\item[${\rm(iv'')}$] For every $\Aa$-algebra $\Aa'$ and every
quasi-projective $S$-morphism $u:X\rightarrow Y$, the diagram
\[\def\objectstyle{\scriptstyle}
\xymatrix@C=.8cm@R=1.75cm{
L\Gamma^n_{\ext}Ru_{!}(K)\overset{\LL}{\otimes}_{\Aa}\Aa'
\ar@{-}[r]^{\sim} \ar[d]^{k^n_u}&
L\Gamma^n_{\ext}(Ru_{!}(K)\overset{\LL}{\otimes}_{\Aa}\Aa')
\ar@{->}[r]^{\sim}&
L\Gamma^n_{\ext}Ru_{!}(K\overset{\LL}{\otimes}_{\Aa}\Aa') \ar[d]^{k^n_u}\\
R\Sym^n_S(u)_{!}L\Gamma^n_{\ext}(K)\overset{\LL}{\otimes}_{\Aa}\Aa'
\ar@{->}[r]^{\sim}&
R\Sym^n_S(u)_{!}(L\Gamma^n_{\ext}(K)\overset{\LL}{\otimes}_{\Aa}\Aa')
\ar@{-}[r]^{\sim}&
R\Sym^n_S(u)_{!}L\Gamma^n_{\ext}(K\overset{\LL}{\otimes}_{\Aa}\Aa')
}
\]
is commutative.

\item[${\rm(v'')}$] For\pageoriginale every base change
$g:S'\rightarrow S$, giving rise to a diagram
\sisi{cartesian}{with cartesian squares}
$$
\xymatrix{
X' \ar[rrr]^{g'}\ar[rd]^{u'}\ar[dd]&   & & X \ar[rd]^{u}\ar[dd]|\hole &   \\
   & Y' \ar[ld]\ar[rrr]^{g''}&  & & Y \ar[ld]\\
S' \ar[rrr]_{g}&  &  & S &   ,
}
$$
the diagram of functors, entirely analogous to ${\rm(iv'')}$, is
commutative:
$$\def\objectstyle{\scriptstyle}
\xymatrix@R=2.25cm{
\Sym^n_g(g'')^{\ast}L\Gamma^n_{\ext}Ru_{!} \ar@{-}[r]^{\sim} \ar[d]^{k^n_u}&
L\Gamma^n_{\ext}g''^{\ast}Ru_{!} \ar[r]^{\sim}& L\Gamma^n_{\ext}Ru'_{!}
g'^{\ast}\ar[d]^{k^n_{u'}} \\
\Sym^n_g(g'')^{\ast}R\Sym^n_S(u)_{!}L\Gamma^n_{\ext} \ar[r]^{\sim}& R
\Sym^n_S(u')_{!}\Sym^n_g(g')^{\ast}L\Gamma^n_{\ext} \ar@{-}[r]^\sim& R\Sym^n_S(u')_{!}L\Gamma^n_{\ext}g'^{\ast}.
}
$$

\item[${\rm(vi'')}$] For every \emph{radicial} base change
$g:Y'\rightarrow Y$, giving rise to a cartesian diagram of $S$-schemes
$$
\xymatrix@C=2.5cm@R=1.5cm{
X' \ar[r]^{g'}\ar[d]^{u'}& X \ar[d]^{u}\\
Y' \ar[r]^{g}& Y,
}
$$
the following diagram of functors is commutative:
\[
\xymatrix{
\Sym^n_S(g)^{\ast}L\Gamma^n_{\ext}Ru_{!} \ar@{-}[r]^{\sim} \ar[d]^{k^n_u}&
L\Gamma^n_{\ext}g^{\ast}Ru_{!} \ar[r]^{\sim}& L\Gamma^n_{\ext}
Ru'_{!}g'^{\ast} \ar[d]^{k^n_{u'}}\\
\Sym^n_S(g)^{\ast}R\Sym^n_S(u\sisi{_!)}{)_!L\Gamma^n_{\ext}} \ar[r]^{\sim}& R
\Sym^n_S(u')_{!}{g'}^{\ast}L\Gamma^n_{\ext} \ar@{-}[r]^{\sim}& R\Sym^n_S(u')_{!}L\Gamma^n_{\ext}g'^{\ast}
}
\]
\end{itemize}

\begin{enonce*}{Theorem 5.5.21}\etiquette[5.5.21]{XVII.thm5.5.21}
{\rm(Symmetric\pageoriginale Künneth formula).} Let $S$ be a coherent
scheme (= quasi-compact and quasi-separated) ringed by a sheaf $\Aa$
of commutative torsion rings whose stalks are noetherian rings. Ring
each $S$-scheme by the inverse image of $\Aa$. Let $u:X\rightarrow Y$
be a morphism of quasi-projective $S$-schemes. For every\sisi{}{ }
integer $n\geq0$ and every $K\in\Ob D^{b,\tor\leq0}(X,\Aa)$, the
symmetric Künneth arrow \eqhyperref[XVII.eq5.5.17.2]{5.5.17.2}
$$
k^n:L\Gamma^n_{\ext}Ru_{!}K\longrightarrow
R\Sym^n_S(u)_{!}L\Gamma^n_{\ext}K
$$
is an isomorphism in
$D^{b,\tor\leq0}(\Sym^n_S(Y),\Aa)$.
\end{enonce*}

The proof of this theorem will occupy the rest of \hyperref[XVII.subsec5.5]{5.5}.

\begin{enonce*}[remark]{Remark 5.5.21.1}\etiquette[5.5.21.1]{XVII.rem5.5.21.1}
The quasi-projectivity hypothesis serves only to ensure that $u$
admits a compactification $\bar u$
$$
X\xrightarrow{j}\bar X\xrightarrow{\bar u}Y
$$
such that the schemes $\Sym^n_S(X)$, $\Sym^n_S(\bar X)$, and
$\Sym^n_S(Y)$ exist for every $n$. The language of algebraic spaces
should make it possible to replace this hypothesis by the sole
hypothesis that $u$ is compactifiable.
\end{enonce*}

The\pageoriginale first part of the proof consists in reducing
\hyperref[XVII.thm5.5.21]{5.5.21}, by dévissage, to an
``irreducible'' case:

\begin{enonce*}{Lemma 5.5.22}\etiquette[5.5.22]{XVII.lem5.5.22}
Theorem \hyperref[XVII.thm5.5.21]{5.5.21} follows from the following
special case:
\begin{itemize}
\item[${(\ast)}$] For every algebraically closed field $k$, every
smooth projective curve $X$ over $k$, and every prime number $\ell$,
the symmetric Künneth arrows (obtained for
$Y=S=\Spec(k)$, $\Aa=\ZZ/\ell$, $K=\ZZ/\ell$; see the explanations
below)
$$\etiquette[5.5.22.1]{XVII.eq5.5.22.1}
L\Gamma^n(R\Gamma(X,\ZZ/\ell))\longrightarrow
R\Gamma(\Sym^n_k(X),\ZZ/\ell)\leqno(5.5.22.1)
$$
are isomorphisms $(n\geq0)$.
\end{itemize}
\end{enonce*}

In \sisi{\hyperref[XVII.eq5.5.22.1]{5.5.22.1}}{\eqhyperref[XVII.eq5.5.22.1]{5.5.22.1}},
the symbol $\Gamma$ has an unfortunate double use: in $R\Gamma$ it
denotes the ``global sections'' functor, whereas in $L\Gamma^n$ it
denotes the functor of \hyperref[XVII.par5.5.2.3]{5.5.2.3}. The
isomorphism $\Gamma^n_{\ext}(\ZZ/\ell)\simeq\ZZ/\ell$ is used in this
formula.

Let $S$, $X$, $Y$, $u$, and $\Aa$ be as in
\hyperref[XVII.thm5.5.21]{5.5.21}.

\begin{enonce*}{Lemma 5.5.22.2}\etiquette[5.5.22.2]{XVII.lem5.5.22.2}
Let there be an exact sequence of complexes in
$C^{b,\geq0}(X,\Aa)$, of \ttor-dimension $\leq0$:
$$
0\longrightarrow K'\longrightarrow K\longrightarrow K''\longrightarrow0{\quoi}.
$$
If all the Künneth arrows $k^n$ $(n\geq0)$ relative to two of the
complexes $K'$, $K$, and $K''$ are isomorphisms, then the Künneth
arrows relative to the third are also isomorphisms.
\end{enonce*}

We\pageoriginale reduce to the case where $u$ is projective and the
complexes $K'$, $K$, and $K''$ satisfy conditions (i) and (ii) of
\hyperref[XVII.lem5.5.16]{5.5.16}, for $N$ at least twice the fiber
dimension of $u$. The symmetric Künneth arrows are then given by
\eqhyperref[XVII.eq5.5.17.1]{5.5.17.1}.

We prove by induction on $n$ that if, for $i\leq n$ and two of the
complexes $K'$, $K$, and $K''$, the arrows $k^i$ are isomorphisms,
then the arrows $k^i$ $(i\leq n)$ are isomorphisms for the third.

To see this, note that by
\sisi{\hyperref[XVII.subsubsec5.5.14]{5.5.14}}{\hyperref[XVII.subsubsec5.5.15]{5.5.15}},
the morphism
\begin{equation*}
k^n:L\Gamma^n_{\ext}Ru_{\ast}K\longrightarrow
R\Sym^n(u)_{\ast}L\Gamma^n_{\ext}K\tag{a}
\end{equation*}
can be represented by a morphism of filtered complexes
$$
{k'}^n:K_0\longrightarrow K_1
$$
with finite filtration; on the successive quotients this morphism
induces arrows isomorphic to the composite arrows
\begin{equation*}
L\Gamma^i_{\ext}Ru_{\ast}K'\overset{\LL}{\vee}
L\Gamma^{n-i}_{\ext}Ru_{\ast}K''\longrightarrow
R\Sym^n(u)_{\ast}(L\Gamma^i_{\ext}K'\overset{\sisi{}{\LL}}{\vee}
L\Gamma^{n-i}_{\ext}K'')\tag{$\text{b}_i$}
\end{equation*}
as follows:
\begin{gather*}
L\Gamma^i_{\ext}Ru_{\ast}K'\overset{\LL}{\vee}
L\Gamma^{n-i}_{\ext}Ru_{\ast}K''\xrightarrow{k^i\overset{\LL}{\vee}k^{n-i}}
R\Sym^i(u)_{\ast}L\Gamma^i_{\ext}K'\overset{\LL}{\vee}
R\Sym^{n-i}(u)_{\ast}L\Gamma^{n-i}_{\ext}(K'')\\
\underset{\defn}{=}
\vee_{i,n-i{\ast}}(R\Sym^i(u)_{\ast}L\Gamma^i_{\ext}K'\overset{\LL}{\boxtimes}
R\Sym^{n-i}(u)_{\ast}L\Gamma^{n-i}_{\ext}(K''))\\
\xrightarrow{\substack{\hyperref[XVII.thm5.4.3]{5.4.3}\\\sim}}
\vee_{i,n-i{\ast}}\sisi{(R\Sym^i(u)_{\ast}\times R\Sym^{n-i}(u)_{\ast})}{(R\Sym^i(u)\times R\Sym^{n-i}(u))_\ast}
(L\Gamma^i_{\ext}K'\overset{\LL}{\boxtimes}
L\Gamma^{n-i}_{\ext}K'')\\
=R\Sym^n(u)_{\ast}\vee_{i,n-i{\ast}}(L\Gamma^i_{\ext}K'\overset{\LL}{\boxtimes}L\Gamma^{n-i}_{\ext}K'')
\overset{\dfn}{=}R\Sym^n(u)_{\ast}(L\Gamma^i_{\ext}K'\overset{\LL}{\vee}L\Gamma^{n-i}_{\ext}K''){\quoi}.
\end{gather*}

By\pageoriginale the induction hypothesis (and the Künneth formula),
these arrows are quasi-isomorphisms for $i\neq0,n$. Any two of the
following assertions therefore imply the third: (a) is a
quasi-isomorphism; $(\text{b}_0)$ is a quasi-isomorphism;
$(\text{b}_n)$ is a quasi-isomorphism. These are precisely the three
arrows in question, \cqfd

\setcounter{subsubsection}{22}
\setcounter{paragraph}{2}
\paragraph{}\etiquette[5.5.22.3]{XVII.par5.5.22.3}
Assume assertion $(\ast)$ of \hyperref[XVII.lem5.5.22]{5.5.22} and
prove \hyperref[XVII.thm5.5.21]{5.5.21}. Let $S$, $X$, $Y$, $\Aa$,
$K$, and $u$ therefore be as in \hyperref[XVII.thm5.5.21]{5.5.21}.
The problem is local on $S$, which may be assumed affine. Every complex
$K$ is a successive extension of complexes of the form $i_{!}i^{\ast}K$
for $i$ the inclusion in $X$ of a locally closed affine subset. By
dévissage (\hyperref[XVII.lem5.5.22.2]{5.5.22.2}), we may therefore
suppose that $K$ has the form $i_{!}K_1$, for an immersion
$i:X_1\rightarrow X$ with affine source. Compatibility
\hyperref[XVII.par5.5.20.2]{5.5.20.2} permits us to replace $X$ by
$X_1$ and $K$ by $K_1$, and hence to suppose that $u$ is affine.
Factor $u$ as a composite $u=u_1\ldots u_k$ of morphisms with fibers
of dimension $\leq1$. Compatibility
\hyperref[XVII.par5.5.20.2]{5.5.20.2} shows that it suffices to prove
that the arrows $k^n$ relative to the $u_i$ and to any complex $K$ are
isomorphisms. We are thus reduced to the case where $S$ is
\emph{affine and $u$ has relative dimension $\leq1$}.

\paragraph{}\etiquette[5.5.22.4]{XVII.par5.5.22.4}
For $k^n$ to be an isomorphism, it suffices that its restrictions to
the geometric fibers of $\Sym^n_S(Y)$ be isomorphisms; compatibility
\hyperref[XVII.par5.5.20.3]{5.5.20.3} ${\rm(v'')}$ therefore reduces
us to the case where $S$ is the \emph{spectrum of an algebraically
closed field $k$}.

\paragraph{}\etiquette[5.5.22.5]{XVII.par5.5.22.5}
For\pageoriginale every rational point $y$ of $\Sym^n_k(Y)$, there is
a finite reduced subscheme $Y'$ of $Y$ such that $\sisi{x}{y}$ belongs
to the image of $\Sym^n_k(Y')$. Compatibility
\hyperref[XVII.par5.5.20.3]{5.5.20.3} ${\rm(vi'')}$ therefore permits
us to suppose that $Y$ \emph{is a disjoint union of finitely many
copies of} $\Spec(k)$.

\paragraph{}\etiquette[5.5.22.6]{XVII.par5.5.22.6}
If $Y=\amalg Y_i$, then every complex $K$ on $X$ is a sum of complexes
supported on one of the schemes $X_i=u^{-1}(Y_i)$. By dévissage
(\hyperref[XVII.lem5.5.22.2]{5.5.22.2}) and compatibility
\hyperref[XVII.par5.5.20.3]{5.5.20.3} ${\rm(vi'')}$, it suffices to
prove \hyperref[XVII.thm5.5.21]{5.5.21} for each morphism
$u_i:X_i\rightarrow Y_i$. By \hyperref[XVII.par5.5.22.5]{5.5.22.5},
this reduces us to the case where
$$
\text{\rm $Y=S=\Spec(k)$, $k$ algebraically closed, $\dim X\leq1$.}
$$

The sheaf $\Aa$ now reduces to a noetherian ring $A$. The following
proposition will allow us to reduce to the case where $A$ is a field.

\begin{enonce*}{Proposition 5.5.23}\etiquette[5.5.23]{XVII.prop5.5.23}
Let $A$ be a (commutative) noetherian ring, let $K\in\Ob D^{-}(A)$ be
a bounded-above complex of $A$-modules, and let $n$ be an integer. The
following conditions are equivalent:

% END INCLUDED SGA4_English_current_buildable_source_20260906/source/17/drafts/SGA4_XVII_LINES_5801_6200_English.texfrag

% BEGIN INCLUDED SGA4_English_current_buildable_source_20260906/source/17/drafts/SGA4_XVII_LINES_6201_6600_English.texfrag
\begin{enumerate}
\renewcommand{\theenumi}{\roman{enumi}}
\renewcommand{\labelenumi}{(\theenumi)}
\item $K$ has \tTor-dimension $\leq n$

\item For every $x\in \Spec (A)$ and every $i<-n$,
$H^{i}(K\overset{\LL}{\underset{A}{\otimes}} k(x))=0$.
\end{enumerate}
\end{enonce*}

One has (i) $\Rightarrow$ (ii) trivially\sisi{,}{.} Suppose (ii). To prove
(i), we prove by noetherian induction on the closed subset $F$ of
$\Spec(A)$ that if a module $N$ satisfies $N_{x}=0$ for $x\not\in F$,
then $H^{i}(K\overset{\LL}{\otimes}_{A}N)=0$
for $i<-n$.

Let\pageoriginale $I$ be the unique ideal defining $F$ such that $A/I$
is reduced, and let $N_{k}$ be the submodule of $N$ annihilated by
$I^{k}$. One has $N=\varinjlim N_{k}$. Since the Tor functors
commute with filtered direct limits, we reduce to the case where there
exists $k$ such that $N=N_{k}$, so that $N$ is a successive extension
of modules annihilated by $I$; the long exact cohomology sequence
reduces us to the case where $N$ is already an $A/I$-module.

Let $\eta$ be a maximal point of $F$, and define $A/I$-modules
$M_{1}$ and $M_{2}$ by the exact sequence
$$\etiquette[5.5.23.1]{XVII.eq5.5.23.1}
0\longrightarrow M_{1}\longrightarrow N \xrightarrow{i} N\otimes
k(\eta)\longrightarrow M_{2} \longrightarrow 0{\quoi}.\leqno(5.5.23.1)
$$

One has $(M_{j})_{\eta}=0$ $(j=1,2)$. If, for every closed subset $G\subset F$
with $\eta\not\in G$, we denote by $M_{j,G}$ the submodule of
$M_{j}$ consisting of the sections of $M_{j}$ supported in $G$, then
$$
M_{j}=\varinjlim_{\eta\not\in G} M_{j,G}{\quoi}.
$$
By the induction hypothesis and passage to the limit, one therefore has
$$
H^{i}(K\overset{\LL}{\otimes}_{A} M_j)=0\quad\mbox{\rm
for}\quad i< -n{\quoi}.
$$

The module $N\otimes k(\eta)$ is a vector space over $k(\eta)$, hence a
direct sum of copies of $k(\eta)$. By hypothesis, one has
$$
H^{i}(K\overset{\LL}{\otimes} (N\otimes
k(\eta))\sisi{}{)}=0\quad\mbox{\rm for}\quad i<-n{\quoi}.
$$

If $I$ is the image of the arrow $i$ in \eqhyperref[XVII.eq5.5.23.1]{5.5.23.1}, there are
exact sequences
\begin{align*}
& H^{i-1}(K\overset{\LL}{\otimes}M_{2})\longrightarrow
H^{i}(K\overset{\LL}{\otimes}I)\longrightarrow
H^{i}(K\overset{\LL}{\otimes}(N\otimes k(\eta)))\\
& H^{i}(K\overset{\LL}{\otimes}M_{1})\longrightarrow
H^{i}(K\overset{\LL}{\otimes}N)\longrightarrow
H^{i}(K\overset{\LL}{\otimes}I){\quoi},
\end{align*}
from which\pageoriginale it follows that
$H^{i}(K\overset{\LL}{\otimes}N)=0$ for $i<-n$,
\cqfd

\begin{enonce*}{Corollary 5.5.24}\etiquette[5.5.24]{XVII.cor5.5.24}
For a morphism $u:K\rightarrow L$ in $D^{-}(A)$ to be an
isomorphism, it suffices that for every $x\in \Spec(A)$, the morphism
$$
u_{x}:K\overset{\LL}{\otimes}_{A}k(x)\longrightarrow
L\overset{\LL}{\otimes}_{A}k(x)
$$
be an isomorphism in $D^{-}(k(x))$.
\end{enonce*}

\sisi{\noindent}{}Let $M$ be the cone (mapping cylinder) of $u$. By hypothesis, one has
$M\overset{\LL}{\otimes}_{A}k(x)=0$ for every
$x\in\Spec(A)$. By \hyperref[XVII.prop5.5.23]{5.5.23}, one therefore has $M=0$, and $u$ is an
isomorphism.

\setcounter{subsubsection}{25}
\setcounter{paragraph}{0}
\paragraph{}\etiquette[5.5.25.1]{XVII.par5.5.25.1}
Let us continue the proof of \hyperref[XVII.lem5.5.22]{5.5.22}, interrupted
at \hyperref[XVII.prop5.5.23]{5.5.23}. By the compatibility
\hyperref[XVII.par5.5.20.3]{5.5.20.3} ${\rm(iv'')}$ and
\hyperref[XVII.cor5.5.24]{5.5.24}, the morphism $k^{n}$
associated with $K\in \Ob D^{b,\tor\geq 0}(X,A)$ is an isomorphism if and
only if the morphisms $k^{n}$ associated with
$K\overset{\LL}{\otimes}_{A}k(x)\in \Ob
D^{b,\tor\geq 0}(X,k(x))$ are isomorphisms, for $x\in
\Spec(A)$. This reduces us to the case where $A$ \emph{is a field of
characteristic} $\ell > 0$.

\paragraph{}\etiquette[5.5.25.2]{XVII.par5.5.25.2}
The exact sequence (where $C$ denotes a cone)
$$
0\longrightarrow K[-1]\longrightarrow CK[-1]\longrightarrow
K\longrightarrow 0
$$
shows that the morphisms $k^{n}$ are isomorphisms for $K$ if and
only if they are so for $K[-1]$ (\hyperref[XVII.lem5.5.22.2]{5.5.22.2}).

\paragraph{}\etiquette[5.5.25.3]{XVII.par5.5.25.3}
The\pageoriginale exact sequences
$$
0\longrightarrow \sigma_{\geq n} K\longrightarrow K\longrightarrow
\sigma_{\leq n} K\longrightarrow 0
$$
show that, for the Künneth morphisms to be isomorphisms when
$K$ is a flat complex, it suffices that they be so for the components
of $K$, placed in degree $0$ (\hyperref[XVII.lem5.5.22.2]{5.5.22.2}\sisi{)}{} and \sisi{(}{}\hyperref[XVII.par5.5.25.2]{5.5.25.2}).

\paragraph{}\etiquette[5.5.25.4]{XVII.par5.5.25.4}
Thus let $\sisi{F}{K}$ be a sheaf of $A$-vector spaces. This sheaf is a filtered
direct limit of constructible sheaves (\hyperrefext[IX][IX.cor2.7.2]{2.7.2}), automatically
flat since $A$ is a field, and, by passage to the limit, we
reduce to the case where $K$ is \emph{constructible} and concentrated in degree zero.

\paragraph{}\etiquette[5.5.25.5]{XVII.par5.5.25.5}
If a complex $K$ is a direct summand of a complex $L$, the morphism
$k^{n}$ associated with $K$ is a direct summand of the morphism $k^{n}$ associated
with $L$, hence is an isomorphism if the latter is one. The trace
method (\hyperrefext[IX][IX.rem5.8]{5.8}) therefore applies, and it suffices to treat the case where $F$
is of the form $j_{!}A_{X'}$, for $j:X'\rightarrow X$ quasi-finite. The
compatibility \hyperref[XVII.par5.5.20.2]{5.5.20.2} allows us to replace $X$ by $X'$, whence the
reduction
$$
\mbox{$X$ \emph{has dimension $\leq 1$ and} $K=A_{X}${\quoi}.}
$$

\paragraph{}\etiquette[5.5.25.6]{XVII.par5.5.25.6}
We may assume that $X$ is reduced. If $j:U\rightarrow X$ is the smooth
locus of $X$, with complement $i:F\rightarrow X$, the exact sequence
$$
0\longrightarrow j_{!}j^{\ast}A_{X}\longrightarrow
A_{X}\longrightarrow i_{!}i^{\ast}A_{X}\longrightarrow 0
$$
reduces us by dévissage to the cases where $X$ \emph{is either a smooth
curve or a finite set.} Moreover, the formula
$$
A_{X}=A\overset{\LL}{\otimes}_{\ZZ/\ell}(\ZZ/\ell)_{X}
$$
reduces us by \hyperref[XVII.par5.5.20.3]{5.5.20.3} ${\rm(iv'')}$ to the case where $A_{X}=\ZZ/\ell$.

If\pageoriginale $X$ is a curve, we embed $X$ in the corresponding
smooth projective curve, and a further dévissage reduces us to
one of the following two cases:
\begin{enumerate}
\renewcommand{\theenumi}{\roman{enumi}}
\renewcommand{\labelenumi}{(\theenumi)}
\item $X$ is a smooth projective curve, $Y=S=\Spec(k)$ with $k$
algebraically closed, $\Aa=\ZZ/\ell$, $K=\ZZ/\ell$.

\item $X$ is a finite disjoint union of copies of $\Spec(k)$,
$Y=S=\Spec(k)$ ($k$ algebraically closed),
$\Aa=\ZZ/\ell$, $K=\ZZ/\ell$.
\end{enumerate}

A final dévissage of $\sisi{}{(}\ZZ/\ell\sisi{}{)}_{X}$ into the sheaves
$\sisi{}{(}\ZZ/\ell\sisi{}{)}_{X\sisi{}{'}}$, for $X'$ a connected component of $X$, reduces us in
(i) and (ii) to assuming that $X$ is connected. Case (ii) is then
trivial. This proves \hyperref[XVII.lem5.5.22]{5.5.22}.

\begin{enonce*}{Lemma 5.5.26}\etiquette[5.5.26]{XVII.lem5.5.26}
Theorem \emph{\hyperref[XVII.thm5.5.21]{5.5.21}} follows from the following two special cases
of assertion $(\ast)$ of \emph{\hyperref[XVII.lem5.5.22]{5.5.22}}:
\begin{itemize}
\item[\emph{Case} 1.] Assume that $k=\CC$.

\item[\emph{Case} 2.] Assume that $k$ has characteristic $p>0$
and that $\ell=p$.
\end{itemize}
\end{enonce*}

With the notation of \hyperref[XVII.lem5.5.22]{5.5.22}, denote by $p$ the characteristic of
$k$. If $p=0$, the Lefschetz principle (relying here on
compatibility \hyperref[XVII.par5.5.20.3]{5.5.20.3} ${\rm(v'')}$ and on XII \hyperref[XII.cor5.4]{5.4}) allows us to
reduce to the case where $k=\CC$.

It remains to treat the case where $p>0$ and $\ell\neq p$. The curve $X$
lifts to a curve $X'$ over the ring $W(k)$ of Witt vectors
over $k$ (\SGA 1 III 7.4). Put $S=\Spec(W(k))$, and consider the diagram
$$
\xymatrix@C=2.25cm@R=1.5cm{
X'_{\eta} \ar[d]^{f_\eta}          \ar[r]              & X' \ar[d]^{f'}& X\ar[l]\ar[d]^{f}\\
\eta               \ar[r]         & S   & \Spec(k), \ar[l],
}
$$
and\pageoriginale let $f_{n}$, $f'_{n}$, and $f_{\eta n}$ be the
respective projections of $\Sym^{n}_{k}(X)$, $\Sym^{n}_{S}(X')$, and
$\Sym^{n}_{\eta}(X'_{\eta})$ onto $\Spec(k)$, $S$, and $\eta$. It is known
that $f'_{n}$ is a smooth projective morphism (projectivity is
clear; for smoothness, the simplest approach is to reduce to the case where
$X'=\PP^{1}_{S}$, in which case $\Sym^{n}_{S}(X')\simeq
\PP^{n}$, as follows from the theorem on elementary symmetric
functions). By the specialization theorem \hyperrefext[XVI][XVI.cor2.2]{2.2}, the
sheaves $R^{i}f'_{n{\ast}}$ $(\ZZ/\ell)$ are locally
constant on $S$. The same holds for the sheaves
$L^{i}\Gamma^{n}(Rf'_{\ast}\ZZ/\ell)$. To verify that the
morphism of \eqhyperref[XVII.eq5.5.22.1]{5.5.22.1} associated with $X$ over $k$ is an isomorphism, it
therefore suffices, by \hyperref[XVII.par5.5.20.3]{5.5.20.3} ${\rm(v'')}$, to prove the analogous statement
for the geometric generic fiber $X_{\bar{\eta}}$ of $X$ over
$S$. The field $k(\bar{\eta})$ has characteristic $0$, and, in view of the
first reduction, this proves \hyperref[XVII.lem5.5.26]{5.5.26}.

To prove \hyperref[XVII.thm5.5.21]{5.5.21} in the two special cases of
\hyperref[XVII.lem5.5.26]{5.5.26}, we
compare the symmetric Künneth morphism on the one hand with its
transcendental analogue (\hyperref[XVII.subsubsec5.5.27]{5.5.27} and
\hyperref[XVII.subsubsec5.5.28]{5.5.28}) and on the other hand with its coherent analogue
(\hyperref[XVII.subsubsec5.5.29]{5.5.29} through \hyperref[XVII.lem5.5.37]{5.5.37}).

\setcounter{subsubsection}{26}
\subsubsection{}\etiquette[5.5.27]{XVII.subsubsec5.5.27}
If $p:X\rightarrow S$ is a continuous map of topological
spaces, with $S$ ringed by a sheaf of commutative rings
$\Aa$, the constructions of \hyperref[XVII.subsubsec5.5.7]{5.5.7} and
\hyperref[XVII.subsubsec5.5.8]{5.5.8} apply unchanged
to define a functor $\Gamma^{n}_{\ext}$ from modules on
$X$ to modules on $\Sym^{n}_{S}(X)=X^{n}/\Sf_{n}$. The
formal arguments of \hyperref[XVII.subsubsec5.5.7]{5.5.7} through \hyperref[XVII.prop5.5.18]{5.5.18}
carry over. Likewise, if $S$, $X$, and $Y$
are locally compact Hausdorff topological spaces
$$
\xymatrix@C=1.5cm@R=1.25cm{
X \ar[rr]^{u}\ar[rd]&  &  Y\ar[ld]\\
            & S &,
}
$$
and\pageoriginale if the fibers of $u$ have bounded cohomological
dimension, the formalism of \hyperref[XVII.subsubsec5.5.19]{5.5.19} and
\hyperref[XVII.subsubsec5.5.20]{5.5.20} for the symmetric Künneth morphism in
cohomology with proper supports carries over.

If $p:X\rightarrow S$ is a morphism of schemes of finite type over
$\CC$, the functors $\Gamma^{n}$ and $L\Gamma^{n}$ commute with the
functors of ``passage to the transcendental setting''. Finally, if $X$, $Y$, $S$,
$\Aa$, and $K$ are as in \hyperref[XVII.thm5.5.21]{5.5.21}, with $S$ separated and of finite type
over $\CC$, and if we denote generically by $\epsilon$ the
morphisms of sites of type $\epsilon:X_{\cl}\rightarrow
X_{\et}$, the following diagram is commutative:
$$
\vcenter{\etiquette[5.5.27.1]{XVII.eq5.5.27.1}
\xymatrix@C=1.75cm@R=1.5cm{
\epsilon^{\ast} L\Gamma^{n}_{\ext} Ru_{!}K  \ar[r]^{k^{\sisi{k}{n}}} \ar[d]&
\epsilon^{\ast} R \Sym^{n}_{S}(u)_{!} L \Gamma^{n}_{\ext} K \ar[d]\\
L\Gamma^{n}_{\ext} R u^{\an}_{!} \epsilon^{\ast}K  \ar[r]^{k^{n \an}}
& R \Sym^{n}_{S}(u^{\an})_{!} L\Gamma^{n}_{\ext} \epsilon^{\ast}K.
}}\leqno{(5.5.27.1)}
$$

\subsubsection{}\etiquette[5.5.28]{XVII.subsubsec5.5.28}
By the compatibility \eqhyperref[XVII.eq5.5.27.1]{5.5.27.1} and an elementary
case of the comparison theorem
(\hyperrefext[XI][XI.thm4.4]{4.4} or \hyperrefext[XVI][XVI.thm4.1]{4.1}),
to treat Case 1 of \hyperref[XVII.lem5.5.26]{5.5.26},
it suffices to verify that if $X$ is a smooth complex projective
curve, then the topological analogue of \eqhyperref[XVII.eq5.5.22.1]{5.5.22.1} is an
isomorphism.

The topological space $X^{\an}$ is triangulable (this fact is
trivial here, since $X$ is a branched covering of the Riemann
sphere). The topological analogue of the dévissage lemma \hyperref[XVII.lem5.5.22.2]{5.5.22.2}
then reduces us to proving
\begin{itemize}
\item[$(\ast)$] If\pageoriginale $D$ is a closed
simplex (of dimension 0, 1, or 2), the topological symmetric Künneth
morphism
$$\etiquette[5.5.28.1]{XVII.eq5.5.28.1}
k^{n}:L\Gamma^{n} R\Gamma(D,\ZZ/\ell)\longrightarrow R\Gamma(\Sym^{n}(D\sisi{}{)},\ZZ/\ell)\leqno(5.5.28.1)
$$
is an isomorphism.
\end{itemize}

One has
$$
R\Gamma(D,\ZZ/\ell)=\ZZ/\ell\;\mbox{\rm placed in degree }
0{\quoi},
$$
and hence
$$
L\Gamma^{n}R\Gamma(D,\ZZ/\ell)=\ZZ/\ell\; \mbox{\rm
placed in degree $0$}{\quoi}.
$$

On the other hand, $\Sym^{n}(D)$ is contractible, and again one has
$$
R\Gamma(\Sym^{n}(D), \ZZ/\ell)=\ZZ/\ell\;\mbox{\rm placed
in degree $0$}{\quoi}.
$$

The $H^{i}$ of the two sides of \eqhyperref[XVII.eq5.5.28.1]{5.5.28.1} therefore vanish in nonzero degrees, and the morphism
induced on $H^0$ identifies with the identity of
$\ZZ/\ell$.

\subsubsection{}\etiquette[5.5.29]{XVII.subsubsec5.5.29}
For the application we have in mind, it is useful to regard a
quasi-coherent sheaf on a scheme $X$ as a
quasi-coherent sheaf of $\Oo$-modules on the étale site of $X$.

Let $p:X\rightarrow \sisi{X}{S}$ be a quasi-projective morphism of schemes (cf.
\hyperref[XVII.rem5.5.21.1]{5.5.21.1} \footnote[1]{Everything that follows, including
  \hyperref[XVII.prop5.5.34]{5.5.34}, remains valid
under the sole assumption that $X$ and $Y$ admit invertible sheaves
relatively ample over $S$ (without assuming that $X$ and $Y$ are of
finite type over $S$).}). If $ℱ$ is a sheaf of $\Oo_{X}$-modules
on $X_{\text{\rm et}}$, then
\begin{enumerate}
\renewcommand{\theenumi}{\alph{enumi}}
\renewcommand{\labelenumi}{\theenumi)}
\item We denote by $\bigotimes\limits^{n}ℱ$ the
sheaf of modules on $(X/S)^{n}_{\text{\rm et}}$ given by the tensor product
of the inverse-image $\Oo_{(X/S)^n}$-modules of $ℱ$ on
$(X/S)^{n}_{\et}$;

\item the sheaf $\bigotimes\limits^{n}ℱ$ is a
$\mathcal{G}_n$-equivariant sheaf. Passing to invariants, it
defines a sheaf $TS^{n}_{\text{\rm ext}}(ℱ)$ on $\Sym^{n}_{S}(X)$:
if $\sigma$ is the projection of $(X/S)^{n}$ onto $\Sym^{n}_{S}(X)$, one has
$$\etiquette[5.5.29.1]{XVII.eq5.5.29.1}
TS^{n}_{\ext}(ℱ)=(\sigma_{\ast}\bigotimes^{n}ℱ)^{\Sf^{n}}{\quoi}.\leqno(5.5.29.1)
$$
\end{enumerate}

\emph{Assume\pageoriginale that $p$ is flat}; as in
\hyperref[XVII.subsubsec5.5.8]{5.5.8}, one then defines the functor
$\Gamma^{m}_{\text{\rm ext}}$ to be the $0$\sisi{th}{\textsuperscript{th}} left derived
functor of $TS^{n}_{\ext}$. With the notation of \hyperref[XVII.par5.5.2.5]{5.5.2.5} and
\hyperref[XVII.subsubsec5.5.8]{5.5.8}, if
$$
P_{1}\longrightarrow P_{0}\longrightarrow ℱ\longrightarrow 0
$$
is a flat presentation of $ℱ$, then the sequence
$$
\xymatrix@C=1.75cm@M=.25cm{
TS^{n}_{\ext} (P'_1) \ar@<3pt>[r]^{j_0} \ar@<-3pt>[r]_{j_1}& TS^{n}_{\ext}(P_0)\ar[r] & \Gamma^{n}_{\ext}(\sisi{F}{ℱ})
}
$$
is exact. If $x$ is a geometric point of $\Sym^{n}_{S}(X)$ as
in \hyperref[XVII.par5.5.7.1]{5.5.7.1}, one has (\cf \eqhyperref[XVII.eq5.5.8.1]{5.5.8.1})
\begin{align*}
& \Gamma^{n}_{\text{\rm
ext}}(ℱ)_{x}=\Oo_{x}\otimes_{C}(\bigotimes^{\ell}_{1}
\Gamma_{\Oo_{S}}^{n_{i}}(\sisi{F}{ℱ}_{y_{i}})){\quoi},\tag{5.5.29.2}\\
& C=\bigotimes^{\ell}_{1} \Sym^{n_{i}}_{\Oo_{S}}(\Oo_{y_{i}}){\quoi}.\tag*{where}
\end{align*}\etiquette[5.5.29.2]{XVII.eq5.5.29.2}%Removing this comment would introduce unwanted white space.
Indeed, it suffices to verify this formula in the trivial case where
the fibers of $\sisi{F}{ℱ}$ are free. Note that $\Oo_{x}$ is flat over
$C$, being the strict henselization of $C$ at the closed point that is the image of $x$.

\subsubsection{}\etiquette[5.5.30]{XVII.subsubsec5.5.30}
Assume that $X$ and $S$ are affine, $p$ is flat, and $ℱ$ is quasi-coherent. Let
$S=\Spec(A)$, $X=\Spec(B)$, and let $F=H^0(X,ℱ)$, so that
$ℱ$ is the sheaf $F^{\sim}$ defined by the $B$-module $F$. It follows
from \eqhyperref[XVII.eq5.5.29.1]{5.5.29.1} that
$$
TS^{n}_{\ext}(ℱ)=TS^{n}_{A}(F)^{\sim}{\quoi},
$$
and, on deriving, one therefore has
$$\etiquette[5.5.30.1]{XVII.eq5.5.30.1}
\Gamma^{n}_{\text{\rm
ext}}(ℱ)=\Gamma^{n}_{A}(F)^{\sim}{\quoi}.\leqno(5.5.30.1)
$$
The\pageoriginale functor $\Gamma^{n}_{\text{\rm ext}}$ thus appears
as a functor $\Gamma^{n}$ ``relative to $A$'', as
\eqhyperref[XVII.eq5.5.29.2]{5.5.29.2} already expressed locally.

\subsubsection{}\etiquette[5.5.31]{XVII.subsubsec5.5.31}
We denote by $D^{b,\tor_{S}\leq 0}(X,\Oo_{X})$ the subcategory
of $D^{b}(X,\Oo_{X})$ formed by the complexes of $\Oo_{X}$-modules
which, as complexes of $p^{-1}(\Oo_{S})$-modules, have
\tTor-dimension $\leq 0$.

Consider the condition

\paragraph{}\etiquette[5.5.31.1]{XVII.par5.5.31.1}
$K\in \Ob C^{b,\geq 0}(X,\sisi{O}{\Oo}_{X})$, and, at every geometric point $x$
of $X$, $K_{x}$ is homotopic to a complex of $\Oo_{X,x}$-modules, in
nonnegative degrees, flat over $\Oo_{S,p(x)}$.

Since the morphism $p$ is flat, so that there are ``enough''
$\Oo_{X}$-modules flat over $S$, one checks as in \hyperref[XVII.lem5.5.13.2]{5.5.13.2} that
$D^{b,\tor_{S}\leq 0}(X,\Oo_{X})$ is obtained from the category of
complexes $K\in \Ob K^{b,\geq 0}(X,\Oo_{X})$ satisfying (\hyperref[XVII.par5.5.31.1]{5.5.31.1}) by
inverting quasi-isomorphisms.

It follows from \eqhyperref[XVII.eq5.5.29.2]{5.5.29.2} and \hyperref[XVII.lem5.5.6.1]{5.5.6.1} that if $u:K\rightarrow L$ is a
quasi-isomorphism between complexes satisfying (\hyperref[XVII.par5.5.31.1]{5.5.31.1}), then
$\Gamma^{n}_{\text{\rm ext}}(u)$ (in the sense of \textsc{Dold-Puppe}) is again a
quasi-isomorphism. The functor $\Gamma^{n}_{\text{\rm ext}}$ can
therefore be derived to give
$$\etiquette[5.5.31.2]{XVII.eq5.5.31.2}
L\Gamma^{n}_{\text{\rm ext}}:D^{b,\tor_{S}\leq
0}(X,\Oo_{X})\longrightarrow D^{b,\tor_{S}\leq 0}(\Sym^{n}_{S}(X), \Oo_{\Sym^{n}_{S}(X)})\leqno(5.5.31.2)
$$

\subsubsection{}\etiquette[5.5.32]{XVII.subsubsec5.5.32}
We denote by $D^{b,\tor_{S}\leq 0}_{\text{\rm qcoh}}(X,\Oo_{X})$
the subcategory of $D^{b,\tor_{S}\leq 0}(X,\Oo_{X})$ formed by
complexes with quasi-coherent cohomology.

% END INCLUDED SGA4_English_current_buildable_source_20260906/source/17/drafts/SGA4_XVII_LINES_6201_6600_English.texfrag

% BEGIN INCLUDED SGA4_English_current_buildable_source_20260906/source/17/drafts/SGA4_XVII_LINES_6601_7000_English.texfrag

Standard\pageoriginale arguments make it possible, as in
\hyperref[XVII.subsubsec5.5.17]{5.5.17}, to
define a symmetric Künneth arrow of the following type:

Let $u:X\rightarrow Y$ be a morphism of quasi-projective $S$-schemes
flat over $S$; for $K\in \Ob D^{b,\tor_{S}\leq 0}_{\text{\rm
qcoh}}(X,\Oo_{X})$, the symmetric Künneth arrow $k^{n}$ in
$D^{b,\tor_{S}\leq 0}_{\text{\rm qcoh}}(\Sym^{n}_{S}(Y),\Oo)$ is
an arrow
$$\etiquette[5.5.32.1]{XVII.eq5.5.32.1}
k^{n}:L\Gamma^{n}_{\text{\rm ext}}Ru_{\ast}K\longrightarrow
R\Sym^{n}_{S}u_{\ast}L\Gamma^{n}_{\text{\rm ext}} K{\quoi}.\leqno(5.5.32.1)
$$

We leave it to the reader to verify the formal properties of
\eqhyperref[XVII.eq5.5.32.1]{5.5.32.1} that will be used, notably the following one:

\begin{enonce*}{Lemma 5.5.33}\etiquette[5.5.33]{XVII.lem5.5.33}
If $S$ is of characteristic $p>0$, and if $u:X\rightarrow S$ is
proper and flat, the diagram of symmetric Künneth arrows
\eqhyperref[XVII.eq5.5.17.2]{5.5.17.2} and \eqhyperref[XVII.eq5.5.32.1]{5.5.32.1} is commutative ($Y=S, K=\sisi{\sheafZZ}{\ZZ}/p$ or
$\Oo_{X}$):
\[
\xymatrix{
L\Gamma^{n}_{\sisi{\sheafZZ}{\ZZ}/p}Ru_{\ast}(\sisi{\sheafZZ}{\ZZ}/p)
\ar[r]^{k^n} \ar[d]&R
\Sym^{n}_S(u)_{\ast} (\sisi{\sheafZZ}{\ZZ}/p) \ar[d]&\\
L\Gamma^{n}_{\Oo} Ru_{\ast}(\Oo) \ar[r]^{k^n}& R
\Sym^{n}_{S}(u)_{\ast}\Oo &.
}
\]
\end{enonce*}


\begin{enonce*}{Proposition 5.5.34}\etiquette[5.5.34]{XVII.prop5.5.34}
Under the hypotheses of \hyperref[XVII.subsubsec5.5.32]{5.5.32}, the symmetric Künneth arrow
\eqhyperref[XVII.eq5.5.32.1]{5.5.32.1} is an isomorphism.
\end{enonce*}

\begin{enumerate}
\renewcommand{\theenumi}{\alph{enumi}}
\renewcommand{\labelenumi}{\theenumi)}
\item The question is local on $S$, which may be assumed affine.

\item For every point $x$ of $\Sym^{n}_{S}(Y)$, there exists an affine
open $Y_{1}$ of $Y$ such that $x$ lies in $\Sym^{n}_{S}(Y_{1})$. By
localization, we may therefore assume that $Y$ is affine.

\item Every\pageoriginale complex $K$ of the type under consideration can be
represented by a complex $K_{0}$ of quasi-coherent modules on
$S$, zero in degrees $<0$. If $\mathcal{U}$ is an affine covering
of $X$, then $K$ admits a finite resolution by sums of
complexes of the form $j_{\ast}j^{\ast}K$, where $j$ is the inclusion of an
affine open (alternating Čech complex). Using
the quasi-coherent analogue of the dévissage lemma \hyperref[XVII.lem5.5.22.2]{5.5.22.2}, this
reduces us to the case where there exist an affine open $j:X_{1}\hookrightarrow X$
of $X$ and $K_{1}\in D^{b,\tor_{S}\leq 0}_{\text{\rm
qcoh}}(X_{1},\Oo)$ such that $K=Rj_{\ast}K_{1}$.

\item In view of the quasi-coherent analogue of transitivity \hyperref[XVII.par5.5.20.2]{5.5.20.2}, it
suffices to prove \hyperref[XVII.prop5.5.34]{5.5.34} for $uj$ and $j$. Reduction b), applied
to $j$, finally reduces us to the case where $X$, $Y$, and $S$ are affine. We
may then represent $K$ by a complex $K_{2}\in C^{b,\geq
0}(X,\Oo_{X})$ with \sisi{}{$S$-}flat quasi-coherent components (\sisi{\quad}{\hyperref[XVII.def4.1.9]{4.1.9}}). The
functors $Ru_{\ast}$ and $L\Gamma^{n}$ can then be computed without
having to resolve $K_{2}$ further, and it remains to note that, by \eqhyperref[XVII.eq5.5.30.1]{5.5.30.1},
for $ℱ$ quasi-coherent on $X$, we have
\[
\Gamma^{n}_{\text{\rm ext}} u_{\ast}ℱ\isomto
\Sym^{n}_{S}(u)_{\ast} \Gamma^{n}_{\text{\rm ext}}ℱ{\quoi},
\]
the two sides being identified with
\[
\Gamma^{n}_{A}(F)^{\sim} \quad\mbox{\rm for}\quad
ℱ=F^{\sim}{\quoi},
\]
\cqfd
\end{enumerate}

\setcounter{subsubsection}{34}
\subsubsection{}\etiquette[5.5.35]{XVII.subsubsec5.5.35}
Let us treat case $2$ of \hyperref[XVII.lem5.5.26]{5.5.26}.
With the notation of \hyperref[XVII.lem5.5.22]{5.5.22}, the
following diagram is commutative (\hyperref[XVII.lem5.5.33]{5.5.33}):
$$\vcenter{\etiquette[5.5.35.1]{XVII.eq5.5.35.1}
\xymatrix@C=2cm{
  L\Gamma^{n} (R\Gamma(X,\ZZ/p)) \ar[r]^{k^n}_*+[o][F-]{\scriptstyle{1}}
  \ar[d]& R\Gamma(\Sym^{n}_{k}(X),\ZZ/p)\ar[d] \\
L\Gamma^{n}(R\Gamma(X,\Oo)) \ar[r]^{k^n}_*+[o][F-]{\scriptstyle{2}} & R\Gamma(\Sym^{n}_{k}(X),\Oo),
}}\leqno{(5.5.35.1)}
$$
and\pageoriginale by \hyperref[XVII.prop5.5.34]{5.5.34}, arrow \ding{193} is an
isomorphism. The Frobenius endomorphism $F:x\rightarrow x^{p}$ of
$\Oo$ acts $p$-linearly on $R\Gamma(X,\Oo)$ and on
$R\Gamma(\Sym^{n}_{k}(X),\Oo)$. It therefore also acts $p$-linearly
on $L\Gamma^{n}(R\Gamma(X,\Oo))$. One checks that $k^{n}$ and $F$
commute.

Artin--Schreier theory provides exact sequences
$$
0\longrightarrow H^{i}(X,\ZZ/p)\longrightarrow
H^{i}(X,\Oo)\xrightarrow{F-I} H^{i}(X,\Oo)\longrightarrow
\sisi{\Oo}{0}{\quoi},
$$
and similarly for $\Sym^{n}_{S}(X)$. To deduce that \ding{192} is
an isomorphism from the fact that \ding{193} is one, it remains to
verify the following

\begin{enonce*}{Lemma 5.5.36}\etiquette[5.5.36]{XVII.lem5.5.36}
Let $k$ be a perfect field of characteristic $p>0$, let $K\in C^{\geq
0}(\ZZ/p)$ be a complex of finite-dimensional $\ZZ/p$-vector spaces,
let $L\in C^{\geq 0}(k)$ be a complex of finite-dimensional $k$-vector spaces,
let $F$ be a $p$-linear endomorphism of $L$, and let
$u:K\rightarrow L$ be a morphism. Suppose that
$H^{\ast}(Fu)=H^{\ast}(u)$, and that the sequences
$$
0\longrightarrow H^{i}(K)\longrightarrow
H^{i}(L)\xrightarrow{F-I}H^{i}(L)\longrightarrow 0
$$
are exact. Then\pageoriginale the sequences
$$
0\longrightarrow H^{i}(\Gamma^{n}K)\longrightarrow
H^{i}(\Gamma^{n}L)\xrightarrow{F-I}H^{i}(\Gamma^{n}L)\longrightarrow 0
$$
are exact.
\end{enonce*}

Since $H^{\ast}(Fu)=H^{\ast}(u)$ and $\ZZ/p$ is a field,
there exists a homotopy
$$
Fu-u=dH+Hd{\quoi}.
$$
Since the field $k$ is perfect, we can write
$$
FH_{0}-H_{0}=H{\quoi}.
$$
Replacing $u$ by $u-(dH_{0}+H_{0}d)$, we may assume that
$Fu=u$.

The functor $S$ that assigns to every $k$-vector space $V$ equipped with a
$p$-linear endomorphism $F$ the $\ZZ/p$-vector space $V^{S}$ of those $v\in
V$ such that $Fv=v$ is an exact functor. The hypothesis therefore means
that the morphism
$$
K\longrightarrow L^{S}
$$
is a quasi-isomorphism. It remains to prove that the functor $S$ commutes
with the functor $\Gamma^{n}$. This follows readily from the following well-known
lemma, whose proof is left to the reader:

\begin{enonce*}{Lemma 5.5.37}\etiquette[5.5.37]{XVII.lem5.5.37}
A finite-dimensional vector space $V$ over a perfect field $k$, equipped
with a $p$-linear endomorphism $F$, admits a unique
decomposition
$$
V=V_{0}+V_{n}
$$
with $V^{S}\otimes k\isomto  V_{0}$, $F(V_{n})\subset V_{n}$,
and $F|V_{n}$ nilpotent.
\end{enonce*}

This completes the proof of \hyperref[XVII.thm5.5.21]{5.5.21}.

\section[The functor $f_{!}$]{The functor $f_{!}$}\etiquette[6]{XVII.sec6}\pageoriginale

We propose to give a direct construction of the functor $f_{!}$
of \hyperref[XVII.var5.1.11]{5.1.11}. We will then use this construction to define the ``trace morphism'' (\hyperref[XVII.thm6.2.3]{6.2.3}) for quasi-finite flat morphisms of
finite presentation.

\subsection[New construction of the functor $f_{!}$]{New construction of the functor $f_{!}$}\etiquette[6.1]{XVII.subsec6.1}

Let $f:X\rightarrow S$ be a morphism locally of finite type. A
subset $P$ of $X$ is said to be \emph{proper} over $S$ if it is the image
of a subscheme of $X$ proper over $S$.

\begin{enonce*}{Proposition 6.1.1}\etiquette[6.1.1]{XVII.prop6.1.1}
Let $f:X\rightarrow S$ be a morphism locally of finite type and
quasi-separated.
\begin{enumerate}
\renewcommand{\theenumi}{\roman{enumi}}
\renewcommand{\labelenumi}{(\theenumi)}
\item A subset of $X$ proper over $S$ is locally closed, and is
closed if $f$ is separated.

\item If $f$ is separated, the union of two subsets of $X$ proper over
$S$ is proper over $S$.

\item For every $S$-scheme $U$, let $\Phi_{X}(U)$ denote the set
of subsets of $\fprod{U}{X}{S}$ proper over $U$. The functor
$\Phi_{X}$ is a sheaf for the \emph{fpqc} topology. If $f$ is
separated, $\Phi\sisi{}{_X}$ is a sheaf of filtered ordered sets for the
inclusion relation.
\end{enumerate}
\end{enonce*}

Assertions\pageoriginale (i) and (ii) are trivial, and the second
assertion of (iii) follows from (ii). The functor $U\mapsto
\mathcal{P}(\fprod{U}{X}{S})$ is a sheaf for every topology
whose coverings are surjective. It follows from \EGA
IV 2.3.14 and \EGA IV 2.7.1 (vii) that $\Phi_{X}$ is a
subsheaf thereof.

If $\Ff $ is an étale sheaf of pointed sets on a
scheme $X$, recall (IV \hyperref[IV.subsubsec8.5.2]{8.5.2}) that the support of a section $s$ of
$\Ff $ is the complement of the largest open set on which $s=0$,
where $0$ denotes the distinguished section.

\begin{enonce*}{Definition 6.1.2}\etiquette[6.1.2]{XVII.def6.1.2}
Let $f:X\rightarrow S$ be a separated morphism locally of finite type,
and let $\Ff $ be a sheaf of pointed sets on $X$. We
denote by $f_{!}\Ff $ the subsheaf of
$f_{\ast}\Ff $ whose sections over $U$ (étale over $S$) are
the sections of $\Ff $ over $\fprod{U}{X}{S}$ with support proper
over $U$.
\end{enonce*}

It follows from \hyperref[XVII.prop6.1.1]{6.1.1} that $f_{!}\Ff $ is indeed a
sheaf.

One readily verifies that:

\begin{enonce*}{Lemma 6.1.3}\etiquette[6.1.3]{XVII.lem6.1.3}
If $f$ is the composite $gh$ of two separated morphisms locally of
finite type, there exists a unique isomorphism of functors
$f_{!}\simeq g_{!}h_{!}$ making the following diagram commutative
$$
\xymatrix@C=1cm@R=1.25cm@M=.2cm{
f_{!}  \ar@{-}[r]^{\sim} \ar@{^{(}->}[d]& g_{!}h_{!} \ar@{^{(}->}[d]\\
f_{\ast} \ar@{=}[r]      & g_{\ast} h_{\ast}.
}
$$
\end{enonce*}

If\pageoriginale $f$ is an open immersion (\resp a proper morphism),
the functor $f_{!}$ is the extension-by-zero functor
(\resp the functor $f_{\ast}$). By \hyperref[XVII.lem6.1.3]{6.1.3}, definitions
\hyperref[XVII.def6.1.2]{6.1.2} and \hyperref[XVII.var5.1.11]{5.1.11} therefore agree on their
common domain of validity.

The properties of the functor $f_{!}$ verified in \hyperref[XVII.sec5]{\S 5} remain
valid in the more general case \hyperref[XVII.def6.1.2]{6.1.2}.

\begin{enonce*}{Proposition 6.1.4}\etiquette[6.1.4]{XVII.prop6.1.4}
Let $f:X\rightarrow S$ be a separated morphism locally of finite type.
\begin{enumerate}
\renewcommand{\theenumi}{\roman{enumi}}
\renewcommand{\labelenumi}{(\theenumi)}
\item The functor $f_{!}$ is left exact (\ie commutes with finite
projective limits) and commutes with filtered inductive limits.

\item If $f$ is quasi-finite, the functor $f_{!}$ is an exact functor
from the category of abelian sheaves on $X$ to that of
abelian sheaves on $S$.

\item For every Cartesian diagram
$$
\xymatrix@C=1.5cm@R=1.25cm{
X'          \ar[r]^{g'} \ar[d]^{f'}& X \ar[d]^{f}\\
S'          \ar[r]^{g}  & S
}
$$
there exists a base-change isomorphism making the
diagram
$$
\xymatrix@C=1.5cm@R=1.25cm{
g^{\ast}f_{!}  \sisi{\ar@{=}[r]}{\ar@{-}[r]^\sim} \ar@{^{(}->}[d]& f'_{!}g'^{\ast}\ar@{^{(}->}[d]\\
g^{\ast}f_{\ast}  \ar[r]    & f'_{\ast}g'^{\ast},
}
$$
commute, in which the second horizontal arrow is the base-change
arrow \emph{\neweqhyperrefext[XII][XII.eq4.1]{4.1}}.
\end{enumerate}
\end{enonce*}

\begin{enonce*}[remark]{Proof}
For\pageoriginale every $U\in \Ob S_{\et}$, we have
\[
f_{!}\Ff (U)=\varinjlim_{P\in
\Phi(U)}\Gamma_{\sisi{p}{P}}(\Ff ){\quoi},
\]
and for $U$ quasi-compact and quasi-separated, the functor
$\Ff \mapsto f_{!}\Ff (U)$ commutes with finite
projective limits and filtered inductive limits, as a filtered
inductive limit (\hyperref[XVII.prop6.1.1]{6.1.1}\sisi{)}{} (iii)\sisi{}{)} of functors having these properties.

This proves (i). Assertion (ii) follows at once from (iii) (take
the geometric points of $S$ as base changes).

Assertion (iii) is local on $S$, which may be assumed
affine. Let $I$ be the filtered ordered set of open subsets of $X$
of finite type over $S$. If
$U\overset{j}{\hookrightarrow}X$
is an open subset of $X$,
$(fj)_{!}(\Ff |U)=f_{!}(j_{!}(\Ff |U))$ maps to
$f_{!}(\Ff\sisi{|U}{})$. Moreover
$$
f_{!}(\Ff )=\varinjlim_{U\in I}
(fj\sisi{}{)}_{!}(\Ff |U){\quoi}.
$$
This allows us to assume that $X$ is of finite type over $S$ and to apply
Chow's lemma (\EGA II 2\sisi{nd}{\textsuperscript{nd}} ed.)\nde{see [C, 2.6], cited in \nderef{XVII.ndeConrad}.} to find a projective and
surjective morphism $p^{0}:X^{0}\rightarrow X$ such that $fp\sisi{_{0}}{^0}$
is compactifiable. Let $X^{1}=\fprod{X^{0}}{X^{0}}{X}$ and let $p^{1}$
be the natural arrow from $X^{1}$ to $X$. The sequence
$$
0\rightarrow \Ff \rightarrow
p^{0}_{\ast}p^{0{\ast}}\Ff \rightrightarrows
p^{1}_{\ast}p^{1{\ast}}\Ff
$$
is exact (\sisi{VII 9.2}{\hyperrefext[VIII][VIII.cor9.2]{9.2}}).

Taking into account the functoriality of the base-change arrow and
(i), it suffices to verify (iii) for the sheaves
$p^{0}_{\ast}p^{0^{\ast}}\Ff $ and
$p^{1}_{\ast}p^{1^{\ast}}\Ff $. By the change of notation
$(p^{1}\mapsto p^{0})$, we\pageoriginale are thus reduced to assuming
that $\Ff $ is of the form $p^{0}_{\ast}\mathcal{G}$. Assertion
(iii) for the pair $(f,p^{0}_{\ast}\mathcal{G})$ is equivalent to
assertion (iii) for the pairs $(p^{0}_{\ast}\sisi{}{,}\mathcal{G})$ and
$(fp^{0}_{\ast},\mathcal{G})$. We conclude by \hyperref[XVII.var5.2.7]{5.2.7} (i), which applies
because $p^{0}$ and $fp^{0}$ are compactifiable.
\end{enonce*}

\begin{enonce*}[remark]{Counterexample 6.1.6}\etiquette[6.1.6]{XVII.contrexem6.1.6}
The functor $Rf_{!}$ does not in general agree with the derived functor
of the functor $f_{!}$. Indeed, if $X$ is affine over an
algebraically closed field $k$, we have
$$
f_{!}\Ff =\sisi{\sum}{\bigoplus}_{x\in X(k)}H^{0}_{x}(\Ff ){\quoi},
$$
and the $i$\textsuperscript{th} derived functor of $f_{!}$ therefore takes the value
at $\Ff $
$$
\sisi{\sum}{\bigoplus}_{x\in X(k)}H^{i}_{x}(\Ff ){\quoi}.
$$
On the other hand, if $X$ is smooth, irreducible, and of dimension 1,
$$
H^{2}_{c}(X,\mu_{\ell^{n}})\simeq \ZZ/\ell^{n}
$$
for $\ell$ prime to the characteristic of $k$.
\end{enonce*}

\subsection{The trace morphism in relative dimension 0}\etiquette[6.2]{XVII.subsec6.2}

The main result of this ${\rm n^{\circ}}$ is Theorem
\hyperref[XVII.thm6.2.3]{6.2.3}. The reader is advised not to read the technical preliminaries
\hyperref[XVII.subsubsec6.2.1]{6.2.1} and
\hyperref[XVII.prop6.2.2]{6.2.2}. In the proof of
\hyperref[XVII.thm6.2.3]{6.2.3}, these allow one to consider only finite morphisms $f$, a reduction that,
via \EGA IV 18.12\sisi{.}{,} is in any event fairly evident.

\subsubsection{}\etiquette[6.2.1]{XVII.subsubsec6.2.1}
Let\pageoriginale $f:X\rightarrow S$ be a \emph{quasi-finite
separated} morphism. For every $U\in \Ob S^{\et}$, let
$\Psi_{f}(U)$ denote the category of pairs consisting of a finite pointed
set $I$, with distinguished point denoted $0$, and a decomposition of
$f^{\ast}U$ into open and closed subsets $f^{\ast}U=\prod_{i\in
I}V_{i}$, with $V_{i}$ finite over $U$ for $i\neq 0$. A morphism from
$\varphi=(V_{i})_{i\in I}$ to $\varphi'=(V'_{i})_{i\in I'}$ is a
map $\sigma$ from $I'$ to $I$ such that, for $i\in I$, one has
$V_{i}=\bigcup_{j\in \sigma^{-1}(i)}V'_{j}$. We say that $\varphi'$
\emph{refines} $\varphi$ if there exists a morphism from $\varphi$ to
$\varphi'$.

Let $\Ff $ be a sheaf of abelian groups on $S$, let $U\in
\Ob S_{\text{\rm et}}$, let $\varphi=(V_{i})_{i\in I}\in \Ob\Psi_{f}(U)$,
let $I^{\ast}=I\backslash \{0\}$, and let $f_{i}$ be the projection of $V_{i}$ onto
$U$. We have
$$
f_{!}f^{\ast}\Ff |U=f_{0!}f_{0}^{\ast}(\Ff |U)\oplus
\bigoplus_{i\neq 0} f_{i\ast}f_{i}^{\ast}(\Ff |U){\quoi},
$$
hence a composite arrow
$$\etiquette[6.2.1.1]{XVII.eq6.2.1.1}
\tau_{\varphi}:\Ff ^{I^{\ast}}|U\longrightarrow
\bigoplus_{i\neq 0}
f_{i\ast}f_{i}^{\ast}(\Ff |U)\longrightarrow f_{!}f^{\ast}\Ff |U{\quoi}.\leqno(6.2.1.1)
$$

If $\sigma:I\rightarrow J$ is a morphism of pointed sets, we
denote by $\leftexp{t}{\sigma}:\Ff ^{J^{\ast}}\rightarrow
\Ff ^{I^{\ast}}$ the arrow with coordinates $\leftexp{t}{\sigma}^i_{\phantom{i}j}=\Id$
if $\sigma(i)=j$, and $\leftexp{t}{\sigma}^i_{\phantom{i}j}=0$ if $\sigma(i)\neq
j$. If $\sigma:I'\rightarrow I$ is a morphism from $\varphi$ to
$\varphi'$, the diagram
$$
\vcenter{\etiquette[6.2.1.2]{XVII.eq6.2.1.2}
\xymatrix@R=1.25cm{
\Ff^{I^{\ast}}|U \ar[rr]^{\leftexp{t}{\sigma}}   \ar[rd] \ar[rd]_{\tau_{\varphi}}&&
\Ff^{I'^{\ast}}|U     \ar[ld]^{\tau_{\varphi'}}\\
      &          f_{!}f^{\ast}\Tt|U  &
}}\leqno{(6.2.1.2)}
$$
is\pageoriginale commutative. Moreover, ${}^{t}(\sigma\rho)={}^{t}\rho
{}^{t}\sigma$. The $\Psi_{f}(U)$ form a fibered category
$\Psi_{f}$ over $S_{\text{\rm et}}$; the preceding constructions

% END INCLUDED SGA4_English_current_buildable_source_20260906/source/17/drafts/SGA4_XVII_LINES_6601_7000_English.texfrag

% BEGIN INCLUDED SGA4_English_current_buildable_source_20260906/source/17/drafts/SGA4_XVII_LINES_7001_7400_English.texfrag
commute with localization and therefore define an $S_{\text{\rm
et}}$-functor $\varphi=(V_{i})_{i\in I} \in \Psi_{f}(U)\mapsto
\Ff ^{I^{\ast}}|U$ from $\Psi$ to the $S_{\text{\rm
et}}$-category of sheaves on étale open subschemes of $S$, and a
morphism of functors
$\tau_{\varphi}:\Ff ^{I^{\ast}}|U\rightarrow
f_{!}f^{\ast}\Ff |U$.

\begin{enonce*}{Proposition 6.2.2}\etiquette[6.2.2]{XVII.prop6.2.2}
\begin{enumerate}
\renewcommand{\theenumi}{\roman{enumi}}
\renewcommand{\labelenumi}{(\theenumi)}
\item The category $\Psi_{f}$ is locally filtering on the right
\emph{(\sisi{IV 14}{\hyperrefext[V][V.def8.1.1]{8.1.1}})} over $S_{\text{\rm et}}$.

\item The arrows \eqhyperref[XVII.eq6.2.1.1]{6.2.1.1} define an isomorphism\sisi{s}{}
\[
\tau:\varinjlim_{\varphi=(V_{i})_{i\in I}\in
\Psi_{f}(U)}\Ff ^{I^{\ast}}|U\isomto  f_{!}f^{\ast}\Ff\sisi{}{|U} {\quoi}.
\]
\end{enumerate}
\end{enonce*}

If $S$ is strictly local with closed point $s$, the $S$-scheme $X$
admits a decomposition $\varphi\in\Psi_{f}(S)$ into open and closed
parts $X=\coprod_{i\in I}X_{i}$ indexed by a pointed set $I$
such that the fiber of $X_{0}$ at $s$ is empty, the $X_{i}$ $(i\neq
0)$ are finite over $S$, and the fiber of $X_{i}$ $(i\neq 0)$ at $s$
consists of a single point (\EGA IV 18.12.1). Such a decomposition is
a final object of $\Psi\sisi{}{_f}(S)$. Moreover, by \hyperref[XVII.prop6.1.4]{6.1.4} (iii) (or
because this is evident), the fiber at $s$ of the arrow $\tau_{\varphi}$
\eqhyperref[XVII.eq6.2.1.1]{6.2.1.1} is an isomorphism.

In the general case, let $x$ be a geometric point of $S$ and let
$\Psi_{x}$ be the fiber of $\Psi_{f}$ at $x$,
\[
\Psi_{x}=\varinjlim_{U(x)} \Psi_{f}(U){\quoi},
\]
where $U$ ranges over the étale neighborhoods of $x$. Let $S_{x}$ be the strict localization of
$S$ at $x$, $X_{x}=\fprod{X}{S_{x}}{S}$, and $f_{x}:X_{x}\rightarrow
S_{x}$. The standard limit arguments show\pageoriginale that
the category $\Psi_{x}$ is equivalent to the category
$\Phi_{f_{x}}(S_{x})$, hence is filtering on the right. It follows that
$\Psi\sisi{}{_f}$ is locally filtering on the right. Moreover, the fiber at $x$ of
$\tau$ identifies with the arrow
$$
\tau_{x}:\varinjlim_{\Psi_{x}}\Ff _{x}^{I_{x}}\longrightarrow
(f_{!}f^{\ast}\Ff )_{x}
$$
and, by the preceding discussion, this arrow is an isomorphism.

\begin{enonce*}{Theorem 6.2.3}\etiquette[6.2.3]{XVII.thm6.2.3}
There is a unique way to define, for every morphism
$f:X\rightarrow S$ that is separated, flat, of finite presentation, and quasi-finite,
and every abelian sheaf $\Ff $ on $S$, a trace morphism
$$
\Tr_{f}:f_{!}f^{\ast}\Ff \longrightarrow \Ff
$$
such that the following conditions hold:
\begin{itemize}
\item[\textsuperscript{(Var 1)}] $\Tr_{f}$ is functorial in $\Ff $.

\item[\textsuperscript{(Var 2)}] $\Tr_{f}$ is compatible with every base change, \ie
for every cartesian diagram
$$
\xymatrix{
X' \ar[rr]^{g'} \ar[d]^{f'}& & X\ar[d]^{f}\\
S' \ar[rr]^{g} & &  S
}
$$
the following diagram of functors is commutative:
$$
\xymatrix{
f'_{!} f'^{\ast} g^{\ast} \ar@{=}[r] \ar[d]^{\Tr_{f'}\ast g^{\ast}}& f'_{!}g'^{\ast} f^{\ast}\ar@{=}[r] &
g^{\ast}f_{!}f^{\ast} \ar[d]^{g^{\ast}\ast
  \Tr_f}\\
g^{\ast} \ar@{=}[rr]&  & g^{\ast} &.
}
$$

\item[\textsuperscript{(Var 3)}] $\Tr_{f}$\pageoriginale is compatible with
composition, \ie if $f$ is the composite $f=gh$ of two morphisms
that are separated, flat, of finite presentation, and quasi-finite, the following diagram
is commutative:
$$
\xymatrix{
f_{!}f^{\ast}\ar[rrdd]_{\Tr_f} \ar@{=}[rr]& & g_{!}h_{!} h^{\ast}g_{\ast}\ar[d]^{g_{!}\ast\mathord{\Tr_{h}}\ast g^\ast}\\
              & &  g_{!}g^{\ast} \ar[d]^{\Tr_g}\\
              & & \Id .
}
$$

\item[\textsuperscript{(Var 4)}] (Normalization). If $f$ is finite of constant rank $n$,
the composite homomorphism
$$
\Ff \rightarrow
f_{\ast}f^{\ast}\Ff =f_{!}f^{\ast}\Ff \xrightarrow{\Tr_{f}}\Ff
$$
is multiplication by $n$.
\end{itemize}
\end{enonce*}

\sisi{\emph{First note that conditions \emph{(Var 2)} through \emph{(Var 4)}
imply that}}{First note that conditions (Var 2) through (Var 4)
imply that}

\begin{enonce*}{Lemma 6.2.3.1}\etiquette[6.2.3.1]{XVII.lem6.2.3.1}
If $X=X'\amalg X''$\sisi{}{,} then, setting $f'=f|X'$ and $f''=f|X''$, the
morphism $\Tr_{f}$:
$$
f_{!}f^{\ast}\Ff =f'_{!}{f'}^{\ast}\Ff \oplus
f''_{!}{f''}^{\ast}\Ff \longrightarrow \Ff
$$
is the sum of $\Tr_{f!}$ and $\Tr_{f''}$.
$$
\xymatrix@C=2.25cm@R=.75cm{
  X' \ar@{^{(}->}[rd]^{j'}&           &\\
   & X   \ar[r]^{f}  \ar[r]^{f}    & S\\
X'' \ar@{^{(}->}[ru]_{j''}& &.
}
$$
\end{enonce*}

If\pageoriginale one applies (Var 2) to $j'$, for base change by
$j'$, and (Var 4) to the identity map, one sees that
$$
\Tr_{j'}:j'_{!}{j'}^{\ast}f^{\ast}\Ff \longrightarrow
f^{\ast}\Ff
$$
is the canonical map defined by the decomposition of
$f^{\ast}\Ff $ as the direct sum
$$
f^{\ast}\Ff =j'_{!}{j'}^{\ast}f^{\ast}\Ff \oplus
j''_{!}{j''}^{\ast}f^{\ast}\Ff {\quoi}.
$$

If one applies (Var 3) to $f'=f\circ j'$, one sees that the following diagram
is commutative
\[
\xymatrix@R=1.25cm{
f'_{!} f'^{\ast}\Ff  \ar[rr] \ar[rd]_{\Tr_{f'}}& &f_{!}f^{\ast}\Ff \ar[ld]^{\Tr_f}\\
                                   & \Ff    &    ,
}
\]
which, together with the analogous statement for $f''$, is precisely \hyperref[XVII.lem6.2.3.1]{6.2.3.1}.

Let $f:X\rightarrow S$ be a separated, flat, quasi-finite morphism of
finite presentation. For every $U\in \Ob S_{\text{\rm et}}$, let
$\Psi'_{f}(U)$ be the full subcategory of $\Psi_{f}(U)$ (\sisi{\hyperref[XVII.prop6.1.1]{6.1.1}}{\hyperref[XVII.subsubsec6.2.1]{6.2.1}})
consisting of the decompositions $(V_{i})_{i\in I}$ such that, for $i\neq
0$, $V_{i}$ is locally free of constant rank over $U$. The
categories\pageoriginale $\Psi'_{f}(U)$ form a fibered subcategory
of $\Psi_{f}$, cofinal in $\Psi_{f}$.

Let $\Ff $ and $\mathcal{G}$ be abelian sheaves on
$S$. By \hyperref[XVII.prop6.2.2]{6.2.2} (ii), giving a morphism
$g:f_{!}f^{\ast}\Ff \rightarrow \mathcal{G}$ is equivalent to giving,
for each decomposition $\varphi=(V_{i})_{i\in I}\in\Psi'_{f}(U)$,
a morphism $g_{\varphi}:\Ff ^{I^{\ast}}|U\rightarrow
\mathcal{G}|U$, compatible with the morphisms ${}^{t}\sigma$ of
\eqhyperref[XVII.eq6.2.1.2]{6.2.1.2}. If $f_{i}$ is the projection from $V_{i}$ to $U$,
$g_{\varphi}$ is the arrow whose coordinate arrows are
the composites:
\[
g^{i}_{\varphi}:\Ff |U\rightarrow
f_{i\ast}f^{\ast}_{i}\Ff |U\rightarrow
f_{!}f^{\ast}\Ff |U\xrightarrow{g} \mathcal{G}|U{\quoi}.
\]

It follows that there exists a unique arrow $\Tr_{f}$, compatible
with étale localization and satisfying (Var 4) and \hyperref[XVII.lem6.2.3.1]{6.2.3.1}. It is clear
that this arrow satisfies (Var 1) and (Var~2).

To verify (Var 3), by base change one reduces to the
easy case in which $S$ is the spectrum of an algebraically closed field.

\setcounter{subsubsection}{3}
\subsubsection{}\etiquette[6.2.4]{XVII.subsubsec6.2.4}
A \emph{weighting} of a separated quasi-finite morphism
$f:X\rightarrow S$ is the datum, for each geometric point $x$ of
$X$, of an integer $n(x)$ depending only on the point of $X$ underlying $x$,
such that:
\begin{itemize}
\item[$(\ast)$] \emph{For every $U\in \Ob S_{\text{\rm et}}$ and
every open and closed subspace $X_{1}$ of $f^{-1}(U)$ that is finite over $U$, the
function assigning to each geometri\sisi{c}{cal} point $s$ of $U$ the number
$\sum_{x\in X_{1}(s)}n(x)$ is locally constant.}
\end{itemize}

If $f=gh$ is the composite of two weighted morphisms, we weight $f$
by setting $n_{f}(x)=n_{h}(x)\sisi{.}{\cdot}n_{g}(h(x))$.

If\pageoriginale the morphism $f'$ is obtained by a base change
$g$ from a weighted morphism $f$, we weight $f'$ by setting
$n_{f'}(x)=n_{f}(g'(x))$.

The preceding arguments in fact prove:

\begin{enonce*}{Proposition 6.2.5}\etiquette[6.2.5]{XVII.prop6.2.5}
There is a unique way to define, for every weighted, separated,
quasi-finite morphism $f:X\rightarrow S$ and every abelian sheaf
$\Ff $ on $S$, a trace morphism
$\Tr_{f}:f_{!}f^{\ast}\Ff \rightarrow \Ff $ that is
functorial in $\Ff $, compatible with base change and with
composition, and satisfies:
\begin{itemize}
\item[\textsuperscript{(Var $\sisi{\leftexp{\prime}{4}}{4'}$)}] If $f$ is finite and there exists an integer $n$ such
that, for every geometric point $s$ of $S$, one has $\sum_{x\in
X(s)}n(x)=n$, then the composite arrow
$$
\Ff \longrightarrow
f_{\ast}f^{\ast}\Ff =f_{!}f^{\ast}\Ff \xrightarrow{\Tr_{f}}\Ff
$$
is multiplication by $n$.
\end{itemize}
\end{enonce*}

\begin{enonce*}[remark]{Example 6.2.6}\etiquette[6.2.6]{XVII.exem6.2.6}
Let $S$ be a \sisi{}{reduced }noetherian geometrically unibranch scheme (for
example, a normal scheme), and let $f:X\rightarrow S$ be a separated
quasi-finite morphism. For every geometric point $x$ of $X$ with image $s$ in $S$,
let $\Oo_{x}$ (\resp $\Oo_{s}$) be the strictly local ring of $X$
at $x$ (\resp of $S$ at $s$). The scheme $\Spec(\Oo_{s})$ is reduced
and has a unique maximal point $\eta$. Moreover, $\Oo_{x}$ is finite over
$\Oo_{s}$. Set $n(x)=\rg_{\eta}(\Oo_{x})$. One verifies that
$n(x)$ is a weighting of $f$, whence a trace morphism $\Tr_{f}$.
\end{enonce*}

\setcounter{subsubsection}{6}
\subsubsection{}\etiquette[6.2.7]{XVII.subsubsec6.2.7}
Let\pageoriginale $u:X\rightarrow Y$ be a separated, quasi-finite,
flat morphism of finite presentation. Since the trace morphism (\hyperref[XVII.thm6.2.3]{6.2.3}) is
functorial, it defines, for every complex of abelian sheaves $K$
on $Y$, a morphism of complexes
$$\etiquette[6.2.7.1]{XVII.eq6.2.7.1}
\Tr_{u}:u_{!} u^{\ast}K\longrightarrow K{\quoi}.\leqno(6.2.7.1)
$$
We shall use the terms \emph{trace morphism, Gysin
morphism} and \emph{morphism} $u_{\ast}$ interchangeably for any morphism obtained from
\eqhyperref[XVII.eq6.2.7.1]{6.2.7.1} or \hyperref[XVII.thm6.2.3]{6.2.3} by applying a functor; for example, the
following ones.

Consider a commutative diagram
$$
\xymatrix@C=1.5cm@R=1.75cm{
X  \ar[rr]^{u} \ar[rd]_{f}&   & Y\ar[ld]^{g}\\
                          & S &
}
$$
with $g$ separated and locally of finite type. For every abelian sheaf
$\Ff $ on $Y$, the trace morphism defines a morphism
$$\etiquette[6.2.7.2]{XVII.eq6.2.7.2}
\Tr_{u}\quad\mbox{\rm or}\quad u_{\ast}:f_{!}u^{\ast}\Ff \simeq
g_{!}u_{!}u^{\ast}\Ff \xrightarrow{g_{!}(\Tr_{u})} g_{!}\Ff {\quoi}.\leqno(6.2.7.2)
$$

If $g$ is compactifiable, and if $K$ is a complex bounded
below with torsion cohomology, or a complex of
$g^{\ast}\Aa$-modules for $\Aa$ a sheaf of torsion rings on
$S$, the trace morphism \eqhyperref[XVII.eq6.2.7.1]{6.2.7.1} likewise induces
\begin{align*}
\Tr_{u}\quad\mbox{\rm or}\quad u_{\ast} &: Rf_{!}u^{\ast}K\simeq
Rg_{!}u_{!}u^{\ast}K\xrightarrow{Rg_{!}(\Tr_{u})}Rg_{!}K{\quoi},\tag{6.2.7.3}\\
\Tr_{u}\quad\mbox{\rm or}\quad u_{\ast} &: R^{q}f_{!}u^{\ast}K\simeq
R^{q}g_{!}u_{!}u^{\ast}K\xrightarrow{R^{q}g_{!}(\Tr_{u})}
R^{q}g_{!}K{\quoi}.\tag{6.2.7.4}
\end{align*}\etiquette[6.2.7.3]{XVII.eq6.2.7.3}
\etiquette[6.2.7.4]{XVII.eq6.2.7.4}%The vertical space following this paragraph is too large.

For\pageoriginale $S$ the spectrum of an algebraically closed field,
\eqhyperref[XVII.eq6.2.7.4]{6.2.7.4} is a morphism
$$\etiquette[6.2.7.5]{XVII.eq6.2.7.5}
u_{\ast}:\HH^{q}_{c}(X,u^{\ast}K)\longrightarrow \HH^{q}_{c}(Y,K){\quoi}.\leqno(6.2.7.5)
$$

Thus one sees that \emph{cohomology with proper supports is
covariant} with respect to quasi-finite morphisms.

\subsubsection{}\etiquette[6.2.8]{XVII.subsubsec6.2.8}
Let $u:X\rightarrow Y$ be as in \hyperref[XVII.subsubsec6.2.7]{6.2.7}, let $\Delta_{n}=[0,n]$
be the standard set with $\sisi{}{(}n+1\sisi{}{)}$ elements, and let $X_{n}$ be
the $(n+1)$\textsuperscript{fold} iterated fiber product of $X$ over $Y$. The $X_{n}$ form
a simplicial scheme augmented to $Y$ (\cf \sisi{VI${}_{\text{B}}$}{VBIS}). We
denote by $u_{n}:X_{n}\rightarrow Y$ the augmentation
morphism. If $\Ff $ is an abelian sheaf on $Y$, the
sheaves $\Ff _{n}=u_{n!}u_{n}^{\ast}\Ff $ form,
by \eqhyperref[XVII.eq6.2.7.2]{6.2.7.2} (applied after the change of notation
$S\sisi{\rightarrow}{\mapsto} Y$, $X\sisi{\rightarrow}{\mapsto} X_{n}$, $Y\sisi{\rightarrow}{\mapsto} Y$ or $X_{m}$),
a \emph{simplicial} sheaf augmented to $\Ff $. By definition, the
differential complex of ``\emph{nondegenerate chains}'' of this
simplicial sheaf has as its $n$\textsuperscript{th}
component the quotient of $\Ff _{n}$, namely the cokernel of the direct
sum of the degeneracy operators with target $\Ff _{n}$ and
of the morphisms $\sigma-\epsilon(\sigma)$ for $\sigma$ a
symmetry operator.

\begin{enonce*}{Proposition 6.2.9}\etiquette[6.2.9]{XVII.prop6.2.9}
Under the preceding hypotheses, if $u$ is separated, étale, of finite type,
and surjective, then for every abelian sheaf $\Ff $ on $Y$, one has
\begin{enumerate}
\renewcommand{\theenumi}{\roman{enumi}}
\renewcommand{\labelenumi}{(\theenumi)}
\item The sheaves $\Ff _{n}$ form a (left)
resolution of $\Ff $.

\item The differential complex of nondegenerate chains of the simplicial
sheaf $(\Ff _{n})$ is a (left) resolution of
$\Ff $. If all the geometric fibers of $u$ have at most
$d$ points, this resolution has length $\leq d$.
\end{enumerate}
\end{enonce*}

By\pageoriginale base change, it suffices to verify this statement
when $Y$ is the spectrum of an algebraically closed field, in which case
\hyperref[XVII.subsubsec6.2.8]{6.2.8} reduces to computing the simplicial homology
of a simplex.

Let $(U_{i})_{i\in I}$ be a finite open covering of the scheme $Y$,
let $\Ff $ be an abelian sheaf on $Y$, and, for each subset $P$
of $I$, let $j_{P}$ be the inclusion of $U_{P}=\cap_{i\in P}U_{i}$ into
$\sisi{X}{Y}$. Suppose that a total order on $I$ has been chosen. If $X=\coprod_{i\in
I}U_{i}$ and if $u$ is the projection from $X$ to $Y$, the resolution
\hyperref[XVII.prop6.2.9]{6.2.9} (ii) of $\Ff $ is given by
$$\etiquette[6.2.9.1]{XVII.eq6.2.9.1}
(\Ff _{n})_{n\,\text{deg}}=\bigoplus_{|P|=n+1}j_{P!}j^{\ast}_{P}\Ff {\quoi}.\leqno(6.2.9.1)
$$

This resolution still makes sense when $I$ is infinite.

\setcounter{subsubsection}{9}
\subsubsection{}\etiquette[6.2.10]{XVII.subsubsec6.2.10}
Under the hypotheses of \hyperref[XVII.prop6.2.9]{6.2.9}, let $\Aa$ be a sheaf
of torsion rings on $Y$ and let $K\in \Ob^{-}(Y,\Aa)$. We
regard $X$ as ringed by $u^{\ast}\Aa$. Since resolution \hyperref[XVII.prop6.2.9]{6.2.9} (i) (\resp
\hyperref[XVII.prop6.2.9]{6.2.9} (ii)) is functorial in $\Ff $, it defines a
resolution of $K$ by a double complex $K_{1}$. This resolution
is called the extraordinary \emph{Čech resolution} (\resp the
\emph{extraordinary alternating} Čech resolution) of $K$.

Consider a commutative diagram
$$
\xymatrix@C=1.5cm@R=1.75cm{
X  \ar[rr]^{u} \ar[rd]_{f}&   & Y\ar[ld]^{g}\\
                          & S & ,
}
$$
with $g$ compactifiable. With the notation of
\hyperref[XVII.subsubsec6.2.8]{6.2.8}, if $f_{n}$ is the projection from $X_{n}$ to
$S$, one has $Rf_{n!}=Rg_{!}\sisi{U}{u}_{n!}$. The complex $K_{1}$, filtered
by\pageoriginale the semisimplicial degree, therefore defines a spectral
sequence
$$\etiquette[6.2.10.1]{XVII.eq6.2.10.1}
E^{-p,q}_{1}=R^{q}f_{p!}(u^{\ast}_{p}K)\Longrightarrow R^{-p+q} g_{!}
K{\quoi}.\leqno(6.2.10.1)
$$
The differentials $d_{1}$ arise from the simplicial arrows, so
that
$$\etiquette[6.2.10.2]{XVII.eq6.2.10.2}
E^{-p,q}_{2}=\check{\Hh}^{p}(R^{q}f_{p!}(u^{\ast}_{p}K)\sisi{}{)}\Longrightarrow
R^{-p+q}g_{!}K{\quoi}.\leqno(6.2.10.2)
$$

Under the hypotheses of \eqhyperref[XVII.eq6.2.9.1]{6.2.9.1}, one similarly obtains a spectral
sequence
$$\etiquette[6.2.10.3]{XVII.eq6.2.10.3}
E^{-p,q}_{1}=\bigoplus_{|P|=p+1}R^{q}f_{p!}(j^{\ast}_{p}K)\Longrightarrow R^{-p+q}g_{!}K{\quoi}.\leqno(6.2.10.3)
$$

The spectral sequences \eqhyperref[XVII.eq6.2.10.1]{6.2.10.1} and \eqhyperref[XVII.eq6.2.10.3]{6.2.10.3} are called the
\emph{upper localization spectral sequences.}

\begin{enonce*}{Proposition 6.2.11}\etiquette[6.2.11]{XVII.prop6.2.11}
If $f$ is a separated étale morphism of finite type, the functor
$f_{!}$ is left adjoint to the functor $f^{\ast}$, and its
adjunction arrow is the trace morphism.
\end{enonce*}

It is known that if $X$ is an object of a site $S$, the functor $f^{\ast}$
restricting abelian sheaves to $X$ has a left adjoint
$f_{!}$. If $F$ is an abelian sheaf on $X$, $f_{!}F$ is the
sheaf generated by the presheaf
$$
V\longrightarrow \bigoplus_{\xi \in \Hom(V,X)}F(\xi){\quoi}.
$$
In the étale case, provisionally denote this functor by $f\sisi{''}{^0}_{!}$. Its
adjunction property immediately implies that its formation commutes with
every base change. The trace morphism defines a morphism $t$
from $f_{!}\Ff $ to $f^{0}_{!}\Ff $,
compatible\pageoriginale with every base change. It suffices to see
that $t$ is an isomorphism fiber by fiber, \ie when $S$ is the spectrum of a
separably closed field, where this is trivial.

\setcounter{subsubsection}{11}
\paragraph{}\etiquette[6.2.11.1]{XVII.par6.2.11.1}
One should note that this interpretation of $f^{0}_{!}$ is
specific to the case of the étale site and does not extend to the fpqf site of a
prescheme, for example.

\subsection{Variant of the trace for continuous coefficients}\etiquette[6.3]{XVII.subsec6.3}

In this no. we prove the assertions of \textsc{Grothendieck}
\cite[no. 6]{sga4.exp.XVII.bib.Gr}, and apply these results to the definition
of a ``trace morphism'' for suitable fppf sheaves.

% END INCLUDED SGA4_English_current_buildable_source_20260906/source/17/drafts/SGA4_XVII_LINES_7001_7400_English.texfrag

% BEGIN INCLUDED SGA4_English_current_buildable_source_20260906/source/17/drafts/SGA4_XVII_LINES_7401_7807_English.texfrag

\subsubsection{}\etiquette[6.3.1]{XVII.subsubsec6.3.1}
Let $A$ be a commutative ring. If $L$ and $M$ are two finitely
generated projective $A$-modules, the map $u\mapsto
\bigwedge\limits^{m}u$ from $\Hom(L,M)$ to
$\Hom(\bigwedge\limits^{m}L,\bigwedge\limits^{m}M)$ is
compatible with every base change and is $m$\sisi{-homogeneous}{\textsuperscript{homogeneous}} (\sisi{\hyperref[XVII.subsubsec5.5.2]{5.5.2}}{\hyperref[XVII.par5.5.2.2]{5.5.2.2}}). By
\sisi{\hyperref[XVII.subsubsec5.5.2]{5.5.2}}{\hyperref[XVII.par5.5.2.3]{5.5.2.3}}, it defines a morphism
$$\etiquette[6.3.1.1]{XVII.eq6.3.1.1}
\bigwedge^{m}:\Gamma^{m}\Hom(L,M)\longrightarrow
\Hom(\bigwedge^{m} L, \bigwedge^{m}M)\leqno(6.3.1.1)
$$
such that, after every base change, for $u\in \Hom(L,M)$ one has
$\bigwedge\limits^{m}(\gamma^{m}(u))=\bigwedge\limits^{m}u$.
Since $\Hom(L,M)$ is flat, this morphism is identified (\sisi{\hyperref[XVII.subsubsec5.5.2]{5.5.2}}{\hyperref[XVII.par5.5.2.5]{5.5.2.5}}) with a
morphism
$$\etiquette[6.3.1.2]{XVII.eq6.3.1.2}
\bigwedge^{m}:TS^{m}\Hom(L,M)\longrightarrow
\Hom(\bigwedge^{m}L, \bigwedge^{m}M){\quoi}.\leqno(6.3.1.2)
$$
N.B. Such a homomorphism can also be defined whenever $2$ is
invertible in $A$: a symmetric element of
$\bigotimes\limits^{m}\Hom(L,M)$ defines a symmetric
homomorphism from $\bigotimes\limits^{m}L$ to
$\bigotimes\limits^{m}M$, and this descends to the
quotient to define a homomorphism from
$\bigwedge\limits^{m}L$ to
$\bigwedge\limits^{m}M$.

Given\pageoriginale three finitely generated projective $A$-modules,
the diagram
$$
\vcenter{\etiquette[6.3.1.3]{XVII.eq6.3.1.3}
\xymatrix{
\Gamma^{m} \Hom(L,M) \otimes \Gamma^{m} \Hom(M,N)  \ar[r] \ar[d]^{\bigwedge\limits^{m}\otimes \bigwedge\limits^{m}}& \Gamma^{m}
\Hom(L,N) \ar[d]^{\bigwedge\limits^{m}}\\
\Hom(\bigwedge\limits^{m} L,
\bigwedge\limits^{m} M) \otimes
\Hom(\bigwedge\limits^{m} M,
\bigwedge\limits^{m} N)         \ar[r]     &
\Hom(\bigwedge\limits^{m} L,
\bigwedge\limits^{m} N),
}}\leqno{(6.3.1.3)}
$$
in which the horizontal arrows are deduced from the composition maps
(and, for the first, from (\hyperref[XVII.par5.5.2.8]{5.5.2.8})), is commutative. This
expresses the formula
$\bigwedge\limits^{m}(vu)=\bigwedge\limits^{m}v\circ
\bigwedge\limits^{m}u$.

If $L$ is a free $A$-module of rank $n$, \eqhyperref[XVII.eq6.3.1.2]{6.3.1.2} defines a
canonical morphism
$$\etiquette[6.3.1.4]{XVII.eq6.3.1.4}
\det:\Gamma^{n}\End_{A}(L)=TS^{n}\End_{A}(L)\longrightarrow
A{\quoi},\leqno(6.3.1.4)
$$
which is an algebra homomorphism by
\eqhyperref[XVII.eq6.3.1.3]{6.3.1.3}. By definition,
one has
$$\etiquette[6.3.1.5]{XVII.eq6.3.1.5}
\det (u^{\otimes n})=\det (u\text{; }L){\quoi}.\leqno(6.3.1.5)
$$

If, in addition, $L$ is a module over an $A$-algebra $B$, \eqhyperref[XVII.eq6.3.1.4]{6.3.1.4}
defines by composition an algebra homomorphism
$$\etiquette[6.3.1.6]{XVII.eq6.3.1.6}
{\det}_{L}:TS^{n}_{A}(B)\longrightarrow A{\quoi}.\leqno(6.3.1.6)
$$

\begin{enonce*}{Lemma 6.3.2}\etiquette[6.3.2]{XVII.lem6.3.2}
Suppose that $L$ occurs in an exact sequence
\[
\sisi{F}{E}:0\longrightarrow L'\longrightarrow L\longrightarrow
L''\longrightarrow 0
\]
with $L'$ (\resp $L''$) locally free of rank $n_{1}$ (\resp
$n_{2}$). Let $n=n_{1}+n_{2}$, let $\alpha$ be the
``associated graded'' morphism from the group $\End(E)$ of
endomorphisms\pageoriginale of the extension $E$ to
$\End(L')\oplus \End(L'')$, and, for two flat $A$-modules $X$ and $Y$,
let $\beta$ be the morphism from $TS^{n}(X\oplus Y)$ to
$TS^{n_{1}}(X)\otimes TS^{n_{2}}(Y)$ deduced from the isomorphism
(\hyperref[XVII.par5.5.2.6]{5.5.2.6})
\[
TS^{n}(X\oplus Y)\simeq \sisi{\sum}{\bigoplus}_{n'+n''=n} TS^{n'\sisi{_{1}}{}}(X)\otimes
TS^{n''}(Y){\quoi}.
\]
The following diagram is commutative:{\footnotesize
\[
\xymatrix@C=1cm@R=1.5cm{
  TS^{n}(\End\sisi{)}{} (E)) \ar[d] \ar[r]^(.4){TS^{n}(\alpha)} & TS^{n}(\End(L')\oplus
\Ens(L''))  \ar[r]^(.45){\beta} & TS^{n_1} (\End(L')) \otimes TS^{n_2}
(\End(L'')) \ar[d]^{\det \otimes \det}\\
TS^{n}(\sisi{L}{\End(L)})         \ar[rr]^{\det} &                 & A.
}
\]
}
\end{enonce*}

By the universal property \hyperref[XVII.subsubsec5.5.2]{5.5.2} of $TS^{n}$, it suffices to
check that the two images of $u^{\otimes n}$ coincide for $u\in
\End(E)$. If $\alpha{u}=(u',u'')$, these images are respectively
$\det(u)$ and $\det(u').\det(u'')$, which proves the assertion.

\begin{enonce*}{Lemma 6.3.3}\etiquette[6.3.3]{XVII.lem6.3.3}
Let $A$ be a\sisi{}{ commutative} ring, $B$ a\sisi{}{ commutative} $A$-algebra, $C$ a
not necessarily commutative $B$-algebra, $E$ a $B$-module that is free
of rank $n$ over $A$, and $F$ a $C$-module that is free of rank $m$
over $B$. Suppose that $C$ is flat over $A$. By the universal property
\hyperref[XVII.par5.5.2.2]{5.5.2.2} of $\Gamma^{nm}$ and \hyperref[XVII.par5.5.2.3]{5.5.2.3}, there is a unique arrow
$\sigma:TS^{nm}_{A}(C)\rightarrow TS^{n}_{A}TS^{\sisi{n}{m}}_{B}(C)$ such
that, after every ring extension $A\rightarrow A'$, one has
$\sigma(c^{\otimes nm})=(c^{\otimes m})^{\otimes n}$. The following
diagram is then commutative:
$$
\xymatrix@C=2cm@R=1.5cm{
TS^n_A (B) \ar[d]^{\det_E}  &  & TS^{n}_A TS^{m}_B (C) \ar[ll]_{TS^n_A
  (\det_F)}\\
A    & & TS^{nm}_A (C) \ar[ll]_{\det_{E\otimes_B F}} \ar[u]_{\sigma}.
}
$$
\end{enonce*}

\begin{enonce*}[remark]{Proof}
By\pageoriginale \hyperref[XVII.par5.5.2.2]{5.5.2.2} and \hyperref[XVII.par5.5.2.3]{5.5.2.3}, it suffices to show that the
two images of an element $u^{\otimes nm}$ of $TS^{nm}_{A}(C)$
coincide. These images are $\det_{E}(\det_{F}(u))$ and
$\det_{E\otimes_{B}F}(u)$; they are expressed in terms of $A$, $B$,
$E$, $F$, and $u\in \End_{B}(F)$, and, to check that they coincide,
one may first replace $B$ by its (commutative) image $B_{1}$ in
$\End(E)$ and $F$ by $F\otimes_{B}B_{1}$. The formula to be proved is
then the block-determinant formula (\textsc{Bourbaki},
\sisi{Alg.}{\emph{Alg.}} \sisi{Ch.}{\chap}~8 \S~12, Lemma 2, \sisi{p.140}{p.~140}).
\end{enonce*}

When $C$ is no longer necessarily flat over $A$, the same proof gives,
more generally, a commutative diagram of algebras
$$\vcenter{\etiquette[6.3.3.1]{XVII.eq6.3.3.1}
\xymatrix@C=3cm@R=1.5cm{
\Gamma^{n}_A (B)  \ar[d]^{\det_E}& \Gamma^{n}_A
\Gamma^n_{B}(C) \ar[l]_{\Gamma_A^{n}(\det_F)}\\
A            & \ar[l]_{\det_{E\otimes_B F}}  \Gamma^{nm}_A (C) \ar[u]_{\sigma}.
}}\leqno{(6.3.3.1)}
$$

\setcounter{subsubsection}{3}
\subsubsection{}\etiquette[6.3.4]{XVII.subsubsec6.3.4}
If $B$ is a commutative $A$-algebra, $\Spec(TS^{n}(B))$ is precisely
the $n$\textsuperscript{th} symmetric power of $\Spec(B)$ over
$\Spec(A)$, \ie the quotient of $(\Spec(B)/\Spec(A))^{n}$ by the
symmetric group (\SGA 1 \sisi{IV 3}{V 1}).

If\pageoriginale $X$ is a scheme separated\sisi{}{ over $S$}, a
quasi-coherent $\Oo_{X}$-module $\Ff $ on $X$ will be called
\emph{finite locally free} (\resp of rank $n$) \emph{over} $S$
if its support is finite over $S$ and its direct image is locally free
of finite rank (\resp of rank $n$) over $S$. By the
\emph{support} of $\Ff $ we shall again mean the \emph{subscheme} of
$X$ defined by the annihilator of $\Ff $ (which is a quasi-coherent
ideal, since $\Ff $ is quasi-coherent of finite type).

If $\Ff $ is finite locally free over $S$, of rank $n$ and with support
$E$, then $E$ is affine over $S$; by globalizing over $S$, $\Ff $
defines, by means of \eqhyperref[XVII.eq6.3.1.6]{6.3.1.6}, a section
$\sigma_{E}(\Ff )$ of $\Sym^{n}_{S}(E)$, hence, by functoriality, a
section $\sigma(\Ff )$ of the algebraic space
\footnote[1]{\MakeLinkTarget*{sga4.exp.XVII.footnote.6.3.4.artin-1}\label{XVII.noteArtin1}M. \textsc{Artin}, Algebraization of formal moduli
I\sisi{}{, in \emph{Global Analysis: Papers in Honor of K.~Kodaira}, Princeton University Press, 1969, 21--71}. We shall use the preceding construction only when
$\Sym^{n}_{S}(X)$ exists as a scheme, which is the case when $X$ is
quasi-projective over $S$.} $\Sym^{n}_{S}(X)$,
$$\etiquette[6.3.4.1]{XVII.eq6.3.4.1}
\sigma(\Ff )\in \Gamma(\Sym^{n}_{S}(X)/S){\quoi}.\leqno(6.3.4.1)
$$

In particular, if $f:S'\rightarrow S$ is finite locally free of rank
$n$, then, for $\Ff =\Oo_{S'}$, one obtains a canonical section
$$\etiquette[6.3.4.2]{XVII.eq6.3.4.2}
t\in\Gamma(\Sym^{n}_{S}(S')/S){\quoi}.\leqno(6.3.4.2)
$$

\begin{enonce*}{Proposition 6.3.5}\etiquette[6.3.5]{XVII.prop6.3.5}
Let $X$ be a scheme separated over $S$, and let $\Ff _{1}$ and
$\Ff _{2}$ be two $\Oo_{X}$-modules finite locally free of ranks
$n_{1}$ and $n_{2}$ over $S$. Put $n=n_{1}+n_{2}$. Let $d$ be the
canonical arrow $d:\Sym^{n_{1}}_{S}(X)\times
\Sym^{n_{2}}_{S}(X)\rightarrow \Sym^{n}_{S}(X)$. If $\Ff $ is an
extension of $\Ff _{1}$ by $\Ff _{2}$, then
$$
\sigma(\Ff )=d(\sigma(\Ff _{1}),\sigma(\Ff _{2})){\quoi}.
$$
\end{enonce*}

One\pageoriginale may replace $X$ by the support (finite over $S$) of
$\Ff _{1}\oplus \Ff _{2}$ and suppose that $S$ is affine. Translated
into ring-theoretic terms, the proposition then follows from
\hyperref[XVII.lem6.3.2]{6.3.2}.

\begin{enonce*}{Lemma 6.3.6}\etiquette[6.3.6]{XVII.lem6.3.6}
If $X$ is a scheme separated over $S$, $\Ff $ an $\Oo_{X}$-module
finite locally free of rank $n$ over $S$, and $\Ll$ an invertible
sheaf on $X$, then
$$
\sigma(\Ff )=\sigma(\Ff \otimes\Ll){\quoi}.
$$
\end{enonce*}

Replace $X$ by the support (\hyperref[XVII.subsubsec6.3.4]{6.3.4}) of
$\Ff $\sisi{ l}{. L}ocally over $S$, $\Ll$ is then isomorphic to
$\Oo_{X}$, and the assertion is trivial.

The reader may verify:

\begin{enonce*}{Proposition 6.3.7}\etiquette[6.3.7]{XVII.prop6.3.7}
For every diagram of $S$-schemes separated over $S$
\[
\xymatrix{
X        \ar[rr]^{u} \ar[rd]_{f} &        &  Y  \ar[ld]^{g}\\
                                & S      &,
}
\]
and every $\Oo_{X}$-module $\Ff $ finite locally free of rank $n$ over
$S$, the $\Oo_{Y}$-module $u_{\ast}\Ff $ is finite locally free of
rank $n$ over $S$, and
\[
\sisi{u}{(\Sym^n_S u)}(\sigma(\Ff ))=\sigma(u_{\ast}\Ff ){\quoi}.
\]
\end{enonce*}

\setcounter{subsubsection}{7}
\paragraph{}\etiquette[6.3.7.1]{XVII.par6.3.7.1}
Finally, note that the construction of the morphisms $\sigma$
\eqhyperref[XVII.eq6.3.4.1]{6.3.4.1} and $t$ \eqhyperref[XVII.eq6.3.4.2]{6.3.4.2} is \emph{compatible with every
base change}, which makes sense by \hyperref[XVII.par5.5.2.7]{5.5.2.7}.

\subsection*{Application 1. The case of smooth curves}\pageoriginale

\subsubsection{}\etiquette[6.3.8]{XVII.subsubsec6.3.8}
Let $f:X\rightarrow S$ be a smooth quasi-projective curve over $S$,
\ie a smooth quasi-projective morphism purely of relative dimension
$1$. The theory of Hilbert schemes \sisi{[}{}\textsc{Grothendieck} \cite{sga4.exp.XVII.bib.Gr}) shows that the
fpqc sheaf $\sheaf{\Div}^{n}_{X/S}$ which assigns to each $S$-scheme
$S'$ the set of relative divisors on $X'=\fprod{X}{S'}{S}$ that are
\emph{finite} over $S$ and of degree $n$ is representable.

For each relative divisor $E$ on $X$, finite of degree $n$ over $S$,
there is a canonical element \sisi{(6.3.4.3)}{\eqhyperref[XVII.eq6.3.4.2]{6.3.4.2}}
$t(E)\in \Sym^{n}_{S}(E)$, with image $\sigma(E)$ in
$\Sym^{n}_{S}(X)$. Since this construction is compatible with every
base change, it defines
$$\etiquette[6.3.8.1]{XVII.eq6.3.8.1}
\sigma:\sheaf{\Div}^{n}_{X/S}\longrightarrow
\Sym^{n}_{S}(X):E\longrightarrow \sigma(\Oo_{E}){\quoi}.\leqno(6.3.8.1)
$$

Each section $s$ of $f$ defines a relative divisor finite over $S$ of
degree $1$, namely the divisor $s(S)$. The morphism
\[
(X/S)^{n}\longrightarrow \sheaf{\Div}^{n}_{X/S}:(s_{i})_{0<i\leq
n}\longmapsto \sum s_{i}(S)
\]
is symmetric, and hence induces a morphism
$$\etiquette[6.3.8.2]{XVII.eq6.3.8.2}
\alpha:\Sym^{n}_{S}(X)\longrightarrow \sheaf{\Div}^{n}_{X/S}{\quoi}.\leqno(6.3.8.2)
$$

\begin{enonce*}{Proposition 6.3.9}\etiquette[6.3.9]{XVII.prop6.3.9}
The morphisms $\sigma$ and $\alpha$ of \hyperref[XVII.subsubsec6.3.8]{6.3.8} are inverse
isomorphisms.
\end{enonce*}

\begin{enonce*}[remark]{Proof}
\begin{enumerate}
\renewcommand{\theenumi}{\alph{enumi}}
\renewcommand{\labelenumi}{\theenumi)}
\item $\sigma\alpha=\Id$. It suffices to prove that the composite arrow
$$
(X/S)^{n}\xrightarrow{\pi}\Sym^{n}_{S}(X)\xrightarrow{\alpha}\sheaf{\Div}^{n}_{X/S}\xrightarrow{\sigma}\Sym^{n}_{S}(X)
$$
is\pageoriginale $\pi$. For $n=1$, the assertion is clear; the general
case follows from the

\begin{enonce*}{Lemma 6.3.9.1}\etiquette[6.3.9.1]{XVII.lem6.3.9.1}
With the notation of \hyperref[XVII.prop6.3.5]{6.3.5} and \eqhyperref[XVII.eq6.3.8.1]{6.3.8.1}, one has
\[
\sigma(E+F)=d(\sigma(E),\sigma(F)){\quoi}.
\]
\end{enonce*}

This lemma follows from \hyperref[XVII.prop6.3.5]{6.3.5}, \hyperref[XVII.lem6.3.6]{6.3.6}, and the exact sequence
\[
0\longrightarrow \Oo(-F)\otimes \Oo_{E}\longrightarrow
\Oo_{E+F}\longrightarrow \Oo_{F}\longrightarrow 0{\quoi}.
\]

\item It is known that $\sheaf{\Div}^{n}_{X/S}$ is smooth over $S$
(there is no obstruction to infinitesimally lifting divisors on a
curve (\cf \SGA 3 XI 1)). The morphism
$\alpha\pi:(X/S)^{n}\rightarrow \sheaf{\Div}^{n}_{X/S}$ is a
quasi-finite surjective morphism between smooth $S$-scheme\sisi{}{s} whose
source has equidimensional fibers. It is therefore flat (\EGA ${\rm
0_{IV}}$ 17.3.5), so $\alpha\pi$ is an epimorphism of fppf sheaves, and
hence so is $\alpha$. In view of a), this completes the proof of
\hyperref[XVII.prop6.3.9]{6.3.9}.
\end{enumerate}
\end{enonce*}

\begin{enonce*}{Lemma 6.3.10}\etiquette[6.3.10]{XVII.lem6.3.10}
Let $k$ be a commutative ring, $M$ a free $k$-module, and $u$ an
endomorphism of $M$. Denote by $N_{\det(u-T)=0}$ the norm from
$k[T]/\det(u-T)$ to $k$. Then, for $P\in k[T]$, one has
\[
N_{\det(u-T)=0} (P(T))=\det(P(u),M){\quoi}.
\]
\end{enonce*}

This amounts to checking an algebraic identity between the entries of
the matrix of $u$ and the coefficients of $P$. To prove
\hyperref[XVII.lem6.3.10]{6.3.10}, it therefore suffices (by familiar
arguments) to treat the case in which $k$ is algebraically closed and
$u$ is semisimple, with eigenvalues $a_{i}$.

One\pageoriginale then has
\begin{align*}
& \det(u-T)=\Pi(a_{i}-T)\\
& N_{\det(u-T)=0}(P(T))=\Pi P(a_{i})\quad\mbox{\rm and}\\
& \det(P(u),M)=\Pi P(a_{i}){\quoi}.
\end{align*}

\setcounter{subsubsection}{11}
\setcounter{paragraph}{-1}
\paragraph{}\etiquette[6.3.11.0]{XVII.par6.3.11.0}
Let $f:X\rightarrow S$ be a quasi-projective morphism of finite
presentation, flat with purely one-dimensional Cohen--Macaulay fibers;
these conditions are thus satisfied if $X$ is a smooth quasi-projective
curve over $S$. Let $\Ff $ be an $\Oo_{X}$-module finite locally free
\emph{over} $S$. Since $\Ff $ vanishes on a fiberwise dense open subset
of $X$, there is a canonical arrow $\psi:\Oo_{X}\rightarrow
\det^{\ast}(\Ff )$ (\SGA 6 XI or \textsc{Mumford} \sisi{[1 Ch.~5 \S~3]}{\cite[\chap~5 \S~3]{sga4.exp.XVII.bib.Mu}}). We
denote by $\delta(\Ff )$ the divisor defined by the annihilator ideal of
$\coker(\psi)$ (ibid.); this divisor is finite over $S$.

\begin{enonce*}{Proposition 6.3.11.1}\etiquette[6.3.11.1]{XVII.prop6.3.11.1}
Under the hypotheses of \hyperref[XVII.par6.3.11.0]{6.3.11.0} and with the notation of
\hyperref[XVII.subsubsec6.3.8]{6.3.8} and \eqhyperref[XVII.eq6.3.4.1]{6.3.4.1}, if $\Ff $ has rank $n$, then
$\delta(\Ff )$ has degree $n$ and
\[
\sigma(\Ff )=\sigma(\delta(\Ff )){\quoi}.
\]
\end{enonce*}

\begin{enonce*}[remark]{Proof}
\begin{enumerate}
\renewcommand{\theenumi}{\alph{enumi}}
\renewcommand{\labelenumi}{\theenumi)}
\item \emph{The case where $f$ is smooth.} By a standard reduction, we
may assume that $S$ is an Artinian local scheme with algebraically
closed residue field. The additivity (\hyperref[XVII.prop6.3.5]{6.3.5})
of $\sigma$ allows us to suppose that the support of $\Ff $ is
concentrated at a rational point $x$ of $X$. One checks at once that if
$T$ is a uniformizer at $x$, defining $p:X\rightarrow
\sisi{A}{\AAA}^{1}_{S}$, then $p_{\ast}(\Oo_{\sisi{S}{\delta}(\Ff )})\simeq
\Oo_{\delta(p_{\ast}\Ff )}$ (because the constructions in question
commute with passage to the completion). By this and
\hyperref[XVII.prop6.3.7]{6.3.7}, we reduce to the case
$X=\sisi{A'}{\AAA^1}_{S}$.

If\pageoriginale $X=\AAA\sisi{'}{^1}_{S}$ and $S$ is affine,
$S=\Spec(k)$, then $\Ff $ is identified with a locally free $k$-module
$M$ equipped with an endomorphism $u$ (the action of $T$). As a
$k[T]$-module, $M$ has the resolution
\[
0\rightarrow
k[T]\otimes_{k}M\xrightarrow{\varphi}k[T]\otimes_{k}M\xrightarrow{\epsilon}
M\rightarrow 0{\quoi},
\]
where $\epsilon(P\otimes m)=P(u)(m)$ and
$\varphi(1\otimes m)=1\otimes um-T\otimes m$. From the formula
$\varphi=u-T$, it follows that the divisor $\delta(\Ff )$ has equation
$\det(u-T,M)=0$.

To check \sisi{6.3.11}{\hyperref[XVII.prop6.3.11.1]{6.3.11.1}}, it suffices to show that, after every base
change and for every section $\varphi$ of $\Oo_{X}$, one has
\[
{\det}_{k}(\varphi,M)={\det}_{k}(\varphi,\Oo_{\delta(\Ff )})(=N_{\delta(\Ff )/S}(\varphi)){\quoi}.
\]
For $\sisi{f}{\varphi}=P(T)$, this is precisely the assertion of
\hyperref[XVII.lem6.3.10]{6.3.10}.

\item The general case. We first state more explicitly the assertion
to be proved, in the following more concrete form:

\begin{enonce*}{Corollary 6.3.11.2}\etiquette[6.3.11.2]{XVII.cor6.3.11.2}
Under the hypotheses of \hyperref[XVII.par6.3.11.0]{6.3.11.0} on
$f:X\rightarrow S$, let $M$ be a finite locally free module on $X$, and
let $u$ be an endomorphism of $M$ such that $\det u$ is a section of
$\Oo_{X}$ that is regular on every fiber $X_{s}$ (whence
$\Coker u=M/u(M)$ and $\Coker \det u=\Oo_{X}/\det u$ are both flat and
of finite presentation over $S$, with \sisi{quasi-finite fibers}{quasi-finite supports}). Suppose that the common support\nde{the support of $M/uM$ in the sense of \hyperref[XVII.subsubsec6.3.4]{6.3.4} is a subscheme of the support of $\Oo/\det u$ defined by a nilideal.} of these
sheaves is finite over $S$, so that they define finite locally free
sheaves over $S$. Then, for every section $\varphi$ of $\Oo_{X}$, the
determinants over $\Oo_{S}$ of $\varphi$ acting on $\Coker u$ and on
$\Coker \det u$ are equal:
$$\etiquette[6.3.11.3]{XVII.eq6.3.11.3}
{\det}_{S}(\varphi,M/u(M))={\det}_{S}(\varphi,\Oo_{X}/\det
u){\quoi}.\leqno(6.3.11.3)
$$
\end{enonce*}

% END INCLUDED SGA4_English_current_buildable_source_20260906/source/17/drafts/SGA4_XVII_LINES_7401_7807_English.texfrag

% BEGIN INCLUDED SGA4_English_current_buildable_source_20260906/source/17/drafts/SGA4_XVII_LINES_7808_8200_English.texfrag

We\pageoriginale note that this statement implies the equality of the ranks
of the two locally free modules under consideration over $S$, by
homogeneity (when $\varphi$ is replaced by $T\varphi$, with $T$ an
``indeterminate'' section of $\Oo_{S}$). Moreover, %\c
replacing $\varphi$ by $1+T\varphi$, where $\sisi{j}{T}$ is such an indeterminate (more
precisely, one makes a base change $S'=S[T]\rightarrow
S$), one deduces from \eqhyperref[XVII.eq6.3.11.3]{6.3.11.3} the apparently more general formula
$$\etiquette[6.3.11.4]{XVII.eq6.3.11.4}
\Pol_{S}(\varphi,M/uM)(T)=\Pol_{S}(\varphi,\Oo_{X}/\det u)(T)\leqno(6.3.11.4)
$$
for the characteristic polynomials.

To prove \eqhyperref[XVII.eq6.3.11.3]{6.3.11.3}, we shall reduce to the already proved smooth
case. By étale localization on $S$, we reduce to the case where $S$ is
strictly local with closed point $s$; then, by additivity, we reduce
to the case where the support of $M/u(M)$ has a single point on $X_{s}$, say
$x$. After making a flat extension of the base, we may assume
the residue field extension $k(x)/k(s)$ is trivial; and after replacing $X$
by a sufficiently small neighborhood of $x$, and adding to $\varphi$ a
function that induces the zero operation on $M/u(M)$ and on
$\Oo_{X}/ \det u$ (which does not change the formula to be proved), one
readily sees that we may assume that $\varphi$ induces a section of
$\Oo_{X_{s}}$ regular at $x$, and hence (after shrinking $X$)
defines a quasi-finite flat morphism
$$
p:X\longrightarrow X'=S[T]{\quoi}.
$$
By étale localization on $X$ and on $X'$, we may therefore assume
that we have a \emph{finite} and \emph{flat} morphism $p$ as above, with
$X'$ smooth over $S$, and $\varphi=p^{\ast}(\varphi')$, where $\varphi'$ is
a section of $\Oo_{X'}$. Set
\begin{align*}
& M'=p_{\ast}(M)\quav u'=p_{\ast}(u){\quoi},\\
& P=\Coker u\quav Q=\Coker \det u{\quoi},\\
& P'=\Coker u'\quav Q'=\Coker \det u'{\quoi}.
\end{align*}
We\pageoriginale know that
\[
\det u'=N_{\sisi{X'/X}{X/X'}}\det u,\quad\mbox{\rm hence}\quad Q'=\Coker
N_{\sisi{X'/X}{X/X'}}(\det u){\quoi}.
\]
We trivially have
\[
{\det}_{S}(\varphi,P)={\det}_{S}(\varphi',P\sisi{}{'}){\quoi},
\]
and by formula \eqhyperref[XVII.eq6.3.11.3]{6.3.11.3}, proved in the smooth case,
\[
{\det}_{S}(\varphi',P')={\det}_{S}(\varphi',Q'){\quoi},
\]
so it remains to prove the formula
\[
{\det}_{S}(\varphi',Q')={\det}_{S}(\varphi,Q){\quoi},
\]
which is \emph{equivalent} to formula \eqhyperref[XVII.eq6.3.11.3]{6.3.11.3} to be established. We thus see
that the validity of the latter depends only on the section $\det u$ of
$\Oo_{X}$, or equivalently that it is equivalent to the same formula in the case
where $M$ is replaced by $\Oo_{X}$, and $\varphi$ by $\det u$. But
in the case $M=\Oo_{X}$, formula \eqhyperref[XVII.eq6.3.11.3]{6.3.11.3} is tautological!
This completes the proof.
\end{enumerate}
\end{enonce*}

\begin{enonce*}[remark]{Remark 6.3.11.5}\etiquette[6.3.11.5]{XVII.rem6.3.11.5}
In\pageoriginale the case where $S$ is Artinian, the equality over $S$ of the ranks
of $\Coker u$ and $\Coker \det u$ (notation of \hyperref[XVII.cor6.3.11.2]{6.3.11.2}) is a
special case of \EGA IV 21.10.17.3, of which \eqhyperref[XVII.eq6.3.11.3]{6.3.11.3} may be
regarded as a ``relative'' version.
\end{enonce*}

\begin{enonce*}{Proposition 6.3.12}\etiquette[6.3.12]{XVII.prop6.3.12}
Let $X$ and $Y$ be two quasi-projective curves of finite
presentation, flat with Cohen-Macaulay fibers over $S$, and let $u:X\rightarrow Y$
be a quasi-finite flat morphism.

The following diagram
\[
\xymatrix@C=4em{
\sheaf{\Div}^{n}_{X/S}  \ar[r]^{u_{\ast}} \ar[d]_{\sigma_X} &
\sheaf{\Div}^{n}_{Y/S} \ar[d]_{\sigma_Y}\\
\Sym^{n}_S(X)     \ar[r]^{\Sym^{n}(u)} & \Sym^{n}_S(Y)
}
\]
is then commutative \footnote[1]{\MakeLinkTarget*{sga4.exp.XVII.footnote.6.3.12.1}\label{XVII.note192}The morphism $u_{\ast}$ is defined
because we restrict ourselves to divisors \emph{finite} over $S$. We take, for
example, $u_{\ast}(D)=\delta(u_{\ast}\Oo_{D})$ as the definition, a
formula that agrees with \EGA IV 21.5.5 when $u$ is finite.}.
\end{enonce*}

\sisi{\noindent}{}%
Since $u_{\ast}(D)=\delta(u_{\ast}\Oo_{D})$, \hyperref[XVII.prop6.3.12]{6.3.12} follows from
\hyperref[XVII.prop6.3.7]{6.3.7} and \hyperref[XVII.prop6.3.11.1]{6.3.11.1}.

\subsection*{Application 2. The trace morphism}

\setcounter{subsubsection}{12}
\subsubsection{}\etiquette[6.3.13]{XVII.subsubsec6.3.13}
Let $G$ be an abelian sheaf on the large fppf site of a scheme $S$
(\hyperref[XVII.subsec0.10]{0.10}). Suppose that $G$ satisfies the following condition.

\paragraph{}\etiquette[6.3.13.1]{XVII.par6.3.13.1}
\emph{For every affine flat morphism of finite presentation
$f:X\rightarrow Y$ of $S$-schemes, and for every integer $n\geq 1$, the
morphism}
\[
\Hom_{S}(\Sym^{n}_{Y}(X),G)\longrightarrow \Hom_{S}((X/Y)^{n},G)
\]
\emph{is\pageoriginale injective and has as its image the set of
symmetric morphisms from $(X/Y)^{n}$ to $G$.}

This condition is stable under base change.

Let $f:T\rightarrow S$ be a finite locally free morphism of rank
$n$. Let $\pr_{i}$ be the $i$\textsuperscript{th} projection from $(T/S)^{n}$
to $T$, and let $t$ be the section \sisi{(6.3.4.3)}{\eqhyperref[XVII.eq6.3.4.2]{6.3.4.2}} of $\Sym^{n}_{S}(T)$. If $u\in
G(T)$, then $\sum \pr^{\ast}_{i}(u)$ is a symmetric element of
$G\sisi{}{(}(T/S)^{n}$) and therefore defines
\[
u_{n}\in \sisi{}{G(}\Sym^{n}_{S}(T)\sisi{}{)}{\quoi}.
\]
We set
\[
\Tr_{f}(u)=t^{\ast}(u_{n}){\quoi}.
\]

When $f$ is finite locally free, we define $\Tr_{f}(u)$,
locally on $S$, by the preceding formula. This construction is
compatible with every base change and therefore defines a
\emph{trace morphism}
$$\etiquette[6.3.13.2]{XVII.eq6.3.13.2}
\Tr_{f}:f_{\ast}f^{\ast}G\longrightarrow G{\quoi}.\leqno(6.3.13.2)
$$

\subsubsection{}\etiquette[6.3.14]{XVII.subsubsec6.3.14}
Condition (\hyperref[XVII.par6.3.13.1]{6.3.13.1}) is satisfied in the following cases:
\begin{enumerate}
\renewcommand{\theenumi}{\alph{enumi}}
\renewcommand{\labelenumi}{\theenumi)}
\item If $G$ is representable, by the definition of symmetric
powers.

\item For $G$ the inverse image on $S_{\text{\rm fppf}}$ of a
sheaf on the small étale site of $S$ (\hyperref[XVII.subsec0.10]{0.10}).

\item For $G$ defined by a quasi-coherent $\Oo_{S}$-module $G_{0}$.
\end{enumerate}

Assertion c) reduces to (\hyperref[XVII.par5.5.2.7]{5.5.2.7}).

\begin{enonce*}{Proposition 6.3.15}\etiquette[6.3.15]{XVII.prop6.3.15}
For\pageoriginale $G$ satisfying (\hyperref[XVII.par6.3.13.1]{6.3.13.1}), we have:
\begin{enumerate}
\renewcommand{\theenumi}{\roman{enumi}}
\renewcommand{\labelenumi}{(\theenumi)}
\item The morphism $\Tr_{f}$ is functorial in $G$, and hence additive.

\item Its formation is compatible with every base change
$S'\rightarrow S$.

\item If $T=\coprod_{i}T_{i}$, so that $f_{\ast}G\simeq \sum
f_{i\ast}G$, then $\Tr_{f}=\sum \Tr_{f_{i}}$.

\item If $T$ has constant rank $n$ over $S$, the composite morphism
\[
G\longrightarrow f_{\ast}f^{\ast}G\xrightarrow{\Tr_{f}} G
\]
is multiplication by $n$.

\item The trace morphism is compatible with composition \emph{(\cf
\hyperref[XVII.thm6.2.3]{6.2.3} Var 3).}
\end{enumerate}
\end{enonce*}

Assertions (i) and (ii) are clear, and (iii) is an easy consequence
of \hyperref[XVII.prop6.3.5]{6.3.5}. With the notation
of \hyperref[XVII.subsubsec6.3.13]{6.3.13}, if $u\in G(S)$, then
$u_{n}$ is $n$ times the pullback of $u$, so that
$\Tr_{f}(f^{\ast}u)=t^{\ast}(u_{n})=nu$.

To prove (v), we reduce, by étale localization on $S$ and
using (iii), to proving that for
\[
f=gh:S''\xrightarrow{h} S'\xrightarrow{g}S{\quoi},
\]
with $\sisi{f}{g}$ and $\sisi{g}{h}$ finite locally free of \emph{constant ranks} $n$
and $m$, the diagram
$$\vcenter{\etiquette[6.3.15.1]{XVII.eq6.3.15.1}
\xymatrix{
G(S'')    \ar[rr]^{\Tr_h}\ar[rd]_{\Tr_f} &        &G(S') \ar[ld]^{\Tr_g}\\
                &  G(S)         &
}}\leqno{(6.3.15.1)}
$$
is commutative. For $u\in G(S'')$, consider the diagram
$$\vcenter{\etiquette[6.3.15.2]{XVII.eq6.3.15.2}
\xymatrix@C=2.5cm@R=1.75cm{
\Sym^{n}_S(S')  \ar[rr]^{\Sym^{n}(t)} \ar[rd]^{\Tr_{\sisi{g}{h}}(u)_n} & & \Sym^{n}_S
\Sym^{m}_{S'} (S'')  \ar[dd] \ar[ld]^{(u_m)_n}\\
                                          & G &\\
          S   \ar@<3pt>[ru]^{\Tr_{\sisi{f}{g}}\Tr_{\sisi{g}{h}}(u)} \ar@<-3pt>[ru]_{\Tr_{\sisi{fg}{gh}}(u)}  \ar[rr]^{t}  \ar[uu]_{t}   &        & \Sym^{nm}_S(S'') \ar[lu]_{u_{nm}}.
}}\leqno{(6.3.15.2)}
$$
The\pageoriginale right vertical arrow is defined because $S\sisi{'}{''}$ is
flat over $S$; every triangle is commutative, and (v) will therefore be a
consequence of the commutativity of the outer boundary.

Translating into ring-theoretic terms, this follows from the
special case $B=E$ and $C=F$ of \hyperref[XVII.lem6.3.3]{6.3.3}.

\setcounter{subsubsection}{15}
\subsubsection{}\etiquette[6.3.16]{XVII.subsubsec6.3.16}
Now let $f:T\rightarrow S$ be a \sisi{separated }{}quasi-finite flat morphism of
finite presentation. If $ℱ$ is an fppf sheaf on $T$, we define
its direct image with proper supports $f_{!}\sisi{T}{ℱ}$ on $S$ as the
presheaf that assigns to $S'/S$ the set of sections of $\sisi{F}{ℱ}$ on
$T'$ whose support is proper over $S'$. This presheaf is a
sheaf by \hyperref[XVII.prop6.1.1]{6.1.1}. With the notation of
\hyperref[XVII.prop6.1.1]{6.1.1}, it may also be
expressed as a local inductive limit (\sisi{IV 14.\quad}{\hyperrefext[V][V.subsec8.2]{8.2}}):
\[
f_{!}\sisi{F}{ℱ(U)}=\varinjlim_{V\in\Phi (U)}f_{\ast}(\sisi{F}{ℱ}|V){\quoi},
\]
a formula\pageoriginale that accounts for the fairly evident fact that the
properties of $f_{!}$ reduce to those of $f'_{\ast}$ when $f'$ is
finite. We thus obtain:

\begin{enonce*}{Proposition 6.3.17}\etiquette[6.3.17]{XVII.prop6.3.17}
There is a unique way to define, for every
$f:T\rightarrow S$ \sisi{}{separated }quasi-finite flat morphism of finite presentation and every
fppf sheaf $G$ on $S$ satisfying (\hyperref[XVII.par6.3.13.1]{6.3.13.1}), a trace morphism
$$
\Tr_{f}:f_{!}f^{\ast}G\longrightarrow G
$$
satisfying:
\begin{enumerate}
\renewcommand{\theenumi}{\roman{enumi}}
\renewcommand{\labelenumi}{(\theenumi)}
\item For finite $f$, one recovers the morphism \eqhyperref[XVII.eq6.3.13.2]{6.3.13.2}.

\item Conditions \emph{(Var 1)} through \emph{(Var 4)} of \hyperref[XVII.thm6.2.3]{6.2.3} are
satisfied, mutatis mutandis.
\end{enumerate}
\end{enonce*}

The quasi-finite morphism $f$ induces a commutative diagram (\cf
(\hyperref[XVII.subsec0.11]{0.11})) of sites (in which $\alpha_{\ast}$ is the ``restriction'' functor)
$$
\xymatrix{
T_{\text{\rm fppf}} \ar[rr]^{\alpha} \ar[d]^{f}& & T_{\text{et}} \ar[d]^{f_{\et}}\\
S_{\text{\rm fppf}} \ar[rr]^{\alpha}& & S_{\et}.
}
$$
For $G$ on $T_{\text{\rm fppf}}$, we have
$$
f_{!}\alpha_{\ast}\simeq \alpha_{\ast}f_{!}{\quoi}.
$$
We\pageoriginale also have a ``base-change''
morphism
$$
c:f^{\ast}_{\text{\rm et}}\alpha_{\ast}\longrightarrow
\alpha_{\ast}f^{\ast}{\quoi}.
$$
We leave it to the reader to verify the following compatibility
between the trace morphisms \hyperref[XVII.thm6.2.3]{6.2.3} and
\hyperref[XVII.prop6.3.17]{6.3.17}. For $G$ on $S_{\text{\rm
fppf}}$ satisfying (\hyperref[XVII.par6.3.13.1]{6.3.13.1}), the following diagram is commutative:
$$\MakeLinkTarget*{sga4.exp.XVII.equation.6.3.17.1}\label{XVII.eq6.3.17.1}
\vcenter{\xymatrix{
f_{\et!}f^{\ast}_{\et}\alpha_{\ast} G \ar[r]^{f_{!}c} \ar[d]_{\Tr_f}^{\hyperref[XVII.thm6.2.3]{(6.2.3)}}&
f_{\et!}\alpha_{\ast}f^{\ast}G  \ar@{<->}[r]^{\sim}&
\alpha_{\ast}f_!f^{\ast}G \ar[d]_{\alpha_{\ast}\Tr_f}^{\hyperref[XVII.prop6.3.17]{(6.3.17)}}\\
\alpha_{\ast}G                   \ar@{=}[rr] &        & \alpha_{\ast}G.
}}\leqno(6.3.17.1)
$$

In more concrete terms, this means that if $G$ is an fppf sheaf
on $S$ whose restriction to the étale site defines an étale sheaf
$\alpha_{\ast}G$, and if $u$ is a properly supported section of
the inverse image $f^{\ast}_{\text{\rm et}}(\alpha_{\ast}G)$, with image
$u_{1}$ in $f^{\ast}G$, then the elements $\Tr_{f}(u)$ (in the sense of \hyperref[XVII.thm6.2.3]{6.2.3}) and
$\Tr_{f}(u_{1})$ (in the sense of \hyperref[XVII.prop6.3.17]{6.3.17}) agree in $G(S)$.

\begin{enonce*}[remark]{Example 6.3.18}\etiquette[6.3.18]{XVII.exem6.3.18}
Let $f:S'\rightarrow S$ be a finite locally free morphism of rank
$n$ and let $u\in \Gamma(S',\Oo_{S'})=\Hom_{S}(S',\EE^{1}_{S})$.\nde{Here $\EE^1_S$ denotes the affine line over $S$.} The morphism $u$
defines a morphism
\[
\Sym(u):\Sym^{n}_{S}(S')\longrightarrow
\Sym^{n}_{S}(\EE^{1}_{S})=\Spec(S[\sigma_{1}\ldots
\sigma_{n}]\sisi{)}{\quoi},
\]
where the $\sigma_{i}$ are the elementary symmetric functions (\cf
\textsc{Bourbaki}, \emph{Alg.} 5, App. I no.~2\nde{reference to be checked. Probably IV \S~6 no.~1 in the new version.}). We then have (by an
easy calculation left to the reader):
$$\etiquette[6.3.18.1]{XVII.eq6.3.18.1}
[\Sym(u)\circ t]^{\ast}(\sigma_{i})=\sigma_{i}(u){\quoi},\leqno(6.3.18.1)
$$
where\pageoriginale $t\in \Gamma$ $(\Sym^{n}_{S}(S')/S)$ is the canonical
element \sisi{(6.3.4.3)}{\eqhyperref[XVII.eq6.3.4.2]{6.3.4.2}}, and where $\sigma_{i}(u)$ denotes the
$i$\textsuperscript{th} coefficient of the characteristic polynomial of
multiplication by $u$ on $f_{\ast}\Oo_{S'}$. It follows that the
``trace'' maps
\[
\begin{cases}%raghu
\GG_{a}(S')\longrightarrow \GG_{a}(S)\\
\GG_{m}(S')\longrightarrow \GG_{m}(S)
\end{cases}
\]
are respectively the usual \emph{trace} and \emph{norm}\sisi{}{}; indeed,
the morphism $\Sym^{n}_{S}(\EE^{1}_{S})\rightarrow
\EE^{1}_{S}$ induced by the additive law (\resp multiplicative law)
on $\EE^{1}_{S}$ is given by the function $\sigma_{1}$ (\resp
$\sigma_{n}$) on $\Sym^{n}_{S}(\EE^{1}_{S})\simeq \Spec
(S[\sigma_{1},\ldots,\sigma_{n}])$.

\sisi{\noindent}{}%
It then follows from \hyperref[XVII.eq6.3.17.1]{(6.3.17.1)} that, when $n$ is invertible on $S$, the
diagram of étale sheaves
$$\vcenter{\etiquette[6.3.18.2]{XVII.eq6.3.18.2}
\xymatrix{
0 \ar[r]& f_{\ast}\mu_n \ar[r] \ar[d]_{\Tr_f}^{\hyperref[XVII.thm6.2.3]{(6.2.3)}}& f_{\ast}\sisi{G}{\GG}_m
\ar[r] \ar[d]^{N_f}
& f_{\ast}\sisi{G}{\GG}_m \ar[d]^{N_f} \ar[r]& 0\\
0 \ar[r]& \mu_n \ar[r]        & \relax\sisi{G}{\GG}_m    \ar[r]     & \relax\sisi{G}{\GG}_m  \ar[r]       & 0
}}\leqno{(6.3.18.2)}
$$
is commutative.
\end{enonce*}

\begin{enonce*}{Proposition 6.3.19}\etiquette[6.3.19]{XVII.prop6.3.19}
Consider a commutative diagram
\[
\xymatrix{
X \ar[rr]^{u} \ar[rd]_{p}&   & Y\ar[ld]^{q}\\
  & S &
}
\]
in\pageoriginale which $X$ and $Y$ are smooth quasi-projective
curves over $S$ and $u$ is a quasi-finite morphism.
\begin{enumerate}
\renewcommand{\theenumi}{\roman{enumi}}
\renewcommand{\labelenumi}{(\theenumi)}
\item If $D$ is a divisor on $X$ finite over $S$ (see note
\emph{p.~\pageref{XVII.note192}}) and $g$ is a section of $G_{Y}$ on $Y$, then
\[
\Tr_{D/S}(u^{\ast}g)=\Tr_{(u_{\ast}D)/S}(g){\quoi}.
\]

\item If $E$ is a divisor on $Y$ finite over $S$, if $g$ is a section of
$G_{X}$ on $X$, and if $u$ is finite, then
\[
\Tr_{E/S}(\Tr_{X/Y}(g))=\Tr_{(u\sisi{_{\ast}}{^\ast}E)/S}(\sisi{G}{g}){\quoi}.
\]
\end{enumerate}
\end{enonce*}

\begin{enonce*}[remark]{Proof}
\begin{enumerate}
\renewcommand{\theenumi}{\roman{enumi}}
\renewcommand{\labelenumi}{(\theenumi)}
\item Let $g_{n}$ be the section of $G$ on $\Sym^{n}_{S}(Y)$ obtained from
$g$ (\hyperref[XVII.subsubsec6.3.13]{6.3.13}). By \hyperref[XVII.prop6.3.12]{6.3.12}, we have
\begin{align*}
\Tr_{D/S}(u^{\ast}g)\overset{\dfn}{=}
\sigma(D)^{\ast}(u^{\ast}g)_{n} & =\sigma(D)^{\ast}(\sisi{u}{\Sym^n_S(u)}^{\ast}g_{n})=\sisi{u}{(\Sym^n_S(u)}\circ
\sigma(D))^\ast
(g_{n})\\
& =\sigma(u_{\ast}D)^\ast(g_{n})=\Tr_{u_{\ast}D}(g_{n}){\quoi}.
\end{align*}

\item Formula (ii) follows from \hyperref[XVII.prop6.3.15]{6.3.15} (ii) and (v):
\[
\Tr_{E/S}(\Tr_{X/Y}(g))=\Tr_{E/S}(\Tr_{u^{\ast}E/E}(g))=\Tr_{u^{\ast}E/S}(g){\quoi}.
\]
\end{enumerate}
\end{enonce*}

\subsection*{Application 3. Trace of a torsor}

\setcounter{subsubsection}{19}
\subsubsection{}\etiquette[6.3.20]{XVII.subsubsec6.3.20}
Let $G$ be a sheaf on the large fppf site of $S$ that satisfies
(\hyperref[XVII.par6.3.13.1]{6.3.13.1}). Let $u:X\rightarrow Y$ be a finite locally free morphism
of $S$-schemes, and let $K$ be a $G_{X}$-torsor on $X$. Suppose that,
locally for the fppf topology on $Y$, the torsor $K$ on $X$ is trivial. This is
the case if $K$ is locally trivial for the étale topology.

\begin{enonce*}{Proposition 6.3.21}\etiquette[6.3.21]{XVII.prop6.3.21}
Under\pageoriginale the hypotheses of \hyperref[XVII.subsubsec6.3.20]{6.3.20}, there exists, up to unique
isomorphism, exactly one torsor $\Tr_{u}(K)$ under $G_{Y}$ on $Y$,
equipped with a morphism $\Tr_{u}$ or $\Tr^{K}_{u}:u_{\ast}K\rightarrow
\Tr_{u}(K)$, which, for $g$ a local section (in the fppf sense) of $G$ on
$Y$ and $k$ a local section of $u_{\ast}K$, satisfies
$$

% END INCLUDED SGA4_English_current_buildable_source_20260906/source/17/drafts/SGA4_XVII_LINES_7808_8200_English.texfrag

% BEGIN INCLUDED SGA4_English_current_buildable_source_20260906/source/17/drafts/SGA4_XVII_LINES_8201_8600_English.texfrag
\Tr^{K}_{u}(gk)=\Tr_{u}(g)\cdot\Tr_{u}\sisi{}{^K}(k){\quoi}.
$$
\end{enonce*}

The problem is local on $Y$; we may therefore assume that $K$ is trivial:
$K=G_{X}$. One solution to the problem is to take $K=G_{Y}$ and as
trace morphism the morphism \eqhyperref[XVII.eq6.3.13.2]{6.3.13.2}. One immediately verifies that it is
the only one.

The trace functor \sisi{\hyperref[XVII.subsubsec6.3.20]{6.3.20}}{\hyperref[XVII.prop6.3.21]{6.3.21}} is defined in particular
when $G$ satisfies the following condition

\setcounter{subsubsection}{21}
\paragraph{For every $S$-scheme $X$, every torsor under $G_{X}$ on $X$
is locally trivial for the étale topology}\etiquette[6.3.21.1]{XVII.par6.3.21.1}

This condition is satisfied when $G$ is a smooth group, when $G$ is the inverse
image of a sheaf on the small étale site of $S$, or when $G$ is
defined by a quasi-coherent sheaf on $S$.

By construction one has
$$\etiquette[6.3.21.2]{XVII.eq6.3.21.2}
\Tr_{u}(K)=u_{\ast}(K) \times_{u_{\ast}G_{X}}G_{Y}{\quoi}.\leqno(6.3.21.2)
$$

\subsubsection{}\etiquette[6.3.22]{XVII.subsubsec6.3.22}
To define a trace functor under more general hypotheses, the
key point is, for $f:X\rightarrow T$ a quasi-projective, flat
morphism of finite presentation of $S$-schemes, to define a
reasonable functor $K\mapsto K_{n}$ from the $G_{X}$-torsors on $X$ to the
$G_{Z}$-torsors on $Z=\Sym^{n}_{T}(X)$. We\pageoriginale first
have the following lemma, where $f_{n}$ denotes the projection from $\Sym^{n}_{T}(X)$ to $T$;
it is proved in the same way as \hyperref[XVII.prop6.3.21]{6.3.21}:

\hypertarget{sga4.exp.XVII.lemma.6.3.22.1}{}\begin{enonce*}{Lemma 6.3.22.1}\etiquette[6.3.22.1]{XVII.lem6.3.22.1}
With the preceding notation, and if, locally on $T$, the torsor
$K$ is trivial, then, up to unique isomorphism, there exists a unique
torsor $K_{n}$ on $\Sym^{n}_{T}(X)$, equipped with a morphism
$\alpha:f_{\ast}K\rightarrow f_{n{\ast}}(K_{n})$, such that for $g$ a
local \emph{(fppf)} section of $f_{\ast}G_{X}$ and $k$ a local section
of $\sisi{u}{f}_{\ast}K$, one has $\alpha(gk)=g_{n}\sisi{\circ}{\cdot} \alpha(k)$ (notation of
\hyperref[XVII.subsubsec6.3.13]{6.3.13} for $g_{n}$; the definition of $g_{n}$ here uses
the flatness hypothesis \etc).
\end{enonce*}

The general case will be handled by a somewhat contrived descent argument. First
of all:

\hypertarget{sga4.exp.XVII.proposition.6.3.23}{}\begin{enonce*}{Proposition 6.3.23}\etiquette[6.3.23]{XVII.prop6.3.23}
Consider a commutative diagram
$$\vcenter{\etiquette[6.3.23.1]{XVII.eq6.3.23.1}
\xymatrix{
X \ar[rr]^{u} \ar[rd]& & Y\ar[ld]\\
  & T &
}}\leqno{(6.3.23.1)}
$$
of flat, finitely presented, quasi-projective $T$-schemes, with
$u$ étale and surjective. Then the semisimplicial system of schemes
$\Sym^{n}_{T}((X/Y)^{k+1})$ is a hypercover for the étale
topology of the scheme $\Sym^{n}_{T}(Y)$.
\end{enonce*}

\begin{enonce*}[remark]{Proof}
\begin{enumerate}
\renewcommand{\theenumi}{\alph{enumi}}
\renewcommand{\labelenumi}{\theenumi)}
\item We shall use the following auxiliary result.

\setcounter{subsubsection}{23}\setcounter{paragraph}{1}
\paragraph{}\etiquette[6.3.23.2]{XVII.par6.3.23.2}
\emph{Let $(X,Y,T)$ be as above, let $t\in T$, and let $Z$ be a
$T$-scheme that is the disjoint union of finitely many $T$-schemes $\Spec(k_{i})$, where
$k_{i}$ is a finite extension of $k_{t}$,\nde{Here $k_t$ denotes the residue field of $T$ at $t$.} and\pageoriginale consider a
commutative diagram of $T$-schemes}
$$\vcenter{\etiquette[6.3.23.3]{XVII.eq6.3.23.3}
\xymatrix{
              & Z \ar@{^{(}->}[ld]_{x} \ar@{^{(}->}[rd]^{y}     \\
X_t \ar[rr]^{u}&    & Y_t,
}}\leqno{(6.3.23.3)}
$$
\emph{such that $x$ and $y$ are closed immersions. Let $z\in
\Sym^{n}_{T}(Z)$, and let $x(\sisi{Z}{z})$ and $y(\sisi{Z}{z})$ be its images in
$\Sym^{n}_{T}(X)$ and $\Sym^{n}_{T}(Y)$. Then $\Sym^{n}_{T}(u)$ is
étale at $x(\sisi{Z}{z})$.}

The formation of the scheme $\Sym^{n}_{T}(X)$ is compatible with every
base change $\sisi{S'}{T'}\rightarrow \sisi{S}{T}$ (\hyperref[XVII.par5.5.2.7]{5.5.2.7}). By reduction to the
noetherian case, it follows that $\Sym^{n}_{T}(X)$ is of finite
presentation over $\sisi{S}{T}$; moreover, it is flat over $\sisi{S}{T}$ (\hyperref[XVII.par5.5.2.4]{5.5.2.4}). By the
fiberwise flatness criterion, to prove
(\hyperref[XVII.par6.3.23.2]{6.3.23.2}) we reduce to the case where $\sisi{S}{T}$ is the spectrum of a field. The completion of
$\Sym^{n}_{T}(X)$ at $\sisi{Z}{z}$ then has as its affine algebra the inverse
limit of the affine algebras localized at $x(z)$ of the schemes
$\Sym^{n}_{T}(Z_{k})$, where $Z_{k}$ ranges over the infinitesimal neighborhoods of $Z$
in $X$. By hypothesis, $u$ induces isomorphisms between the
infinitesimal neighborhoods of $Z$ in $X$ and those in $Y$. The completion of
$\Sym^{n}_{T}(u)$ at $x(z)$ and $y(z)$ is therefore an isomorphism, and
this proves (\hyperref[XVII.par6.3.23.2]{6.3.23.2}).

It follows from \hyperref[XVII.par6.3.23.2]{6.3.23.2} that the scheme given by the disjoint union, over $z\in
\Sym^{n}_{T}(Z)$, of the henselizations at $y(z)$ of $\Sym^{n}_{T}(Y)$
depends only on $Z$ and on the $T$-scheme $X^{h}_{Z}$ that is the disjoint union of the henselizations of
$X$ at the various points of $Z$. We shall denote this scheme by
$\Sym^{h\,n}_{T}(X^{h}_{Z})$; it is functorial in
$X^{h}_{Z}$.

\item By\pageoriginale passing to the limit, to prove \hyperref[XVII.prop6.3.23]{6.3.23}, it
suffices to show that for $y\in \Sym^{n}_{T}(Y)$, closed in its
fiber, the pullback to the henselization of $\Sym^{n}_{T}(Y)$ at that point of the
simplicial scheme $\Sym^{n}_{T}((X/Y)^{k+1})$ is a hypercover
for the étale topology. This problem is local on $T$ for the
étale topology; we thereby reduce to the case where there exists a diagram
\eqhyperref[XVII.eq6.3.23.3]{6.3.23.3} with $y=x(z)$ (after killing the residue-field extensions by a
separable extension of $k_{t}$). The semisimplicial system under consideration
is then obtained from the system
$$
\Sym^{h\,n}_{T}((X^{h}_{u^{-1}(Z)}/Y^{h}_{Z})^{k+1}),\quad\mbox{\rm
augmented to}\quad \Sym^{h\,n}_{T}(Y^{h}_{Z})
$$
by base change, so it suffices to consider the
latter. The semisimplicial system of
$(X^{h}_{u^{-1}(Z)}/Y^{h}_{Z})^{k+1}$ is semisimplicially
homotopic to the constant system $Y^{h}_{Z}$, because
$X^{h}_{u^{-1}(Z)}/Y^{h}_{Z}$ admits a section. The system
$\Sym^{h\,n}_{T}((X^{h}_{u^{-1}(Z)}/Y^{h}_{Z})^{h+1})$ is therefore likewise
homotopic to $\Sym^{h\, n}_{T} (Y^{h}_{Z})$, and in particular is
a hypercover of it. This proves \sisi{6.3.2.3}{\hyperref[XVII.prop6.3.23]{6.3.23}}.
\end{enumerate}
\end{enonce*}

\subsubsection{}\etiquette[6.3.24]{XVII.subsubsec6.3.24}
Let $f:Y\rightarrow T$ be a flat, quasi-projective morphism of
finite presentation between $S$-schemes, and let $K$ be a $G_{Y}$-torsor on
$Y$. Consider the following problem:

\setcounter{paragraph}{-1}
\paragraph{}\etiquette[6.3.24.0]{XVII.par6.3.24.0}
Define, for every morphism $u:X\rightarrow Y$, with $X$
flat, quasi-projective, and of finite presentation over $T$, a torsor
$(u^{\ast}K)_{n}$ on $\Sym^{n}_{T}(X)$, and for every $v:X'\rightarrow
X$ an isomorphism $c_{v}:(v^{\ast}u^{\ast}K)_{n}\simeq
\Sym^{n}_{T}(v)^{\ast}\sisi{K}{(u^K)}_{n}$, such that
\begin{enumerate}
\renewcommand{\theenumi}{\roman{enumi}}
\renewcommand{\labelenumi}{(\theenumi)}
\item the\pageoriginale isomorphisms $c_{v}$ are compatible with
  composition;

\item when $u^{\ast}K$ is locally trivial on $T$, one recovers the
torsor of \hyperref[XVII.lem6.3.22.1]{6.3.22.1} and its functoriality.
\end{enumerate}

By \hyperref[XVII.prop6.3.23]{6.3.23}, the existence and uniqueness of a solution to
\hyperref[XVII.par6.3.24.0]{6.3.24.0}, for every torsor obtained from $K$ by an inverse image of the
type considered in \hyperref[XVII.par6.3.24.0]{6.3.24.0}, is a local question for the
étale topology on $Y$. In particular, if $G$ satisfies (\hyperref[XVII.par6.3.21.1]{6.3.21.1}), the problem
\hyperref[XVII.par6.3.24.0]{6.3.24.0} has a unique solution.

Now suppose that $G$ satisfies:

\paragraph*{(6.3.24.1)}\MakeLinkTarget*{sga4.exp.XVII.paragraph.6.3.24.1}\label{XVII.par6.3.24.1}
\emph{For every integer $n$ and every diagram of affine
$S$-schemes
\[
\xymatrix{
X \ar[rr]^{u} \ar[rd]_{f}& & Y\ar[ld]^{g}\\
  &T&,
}
\]
with $f$ and $g$ flat and of finite presentation and $u$ finite,
locally free, and surjective, the truncated semisimplicial scheme
$(\Sym^{n}_{T}((X/Y)^{k+1}))_{0\leq k\leq 2}$, augmented to
$\Sym^{n}_{T}(Y)$, is a descent diagram \emph{(\textsc{Giraud} \cite{sga4.exp.XVII.bib.Gi})} for the
category of $G$-torsors.}

\paragraph*{6.3.24.1 bis}\etiquette[6.3.24.1 bis]{XVII.par6.3.24.1bis}
It follows in the standard way from \hyperref[XVII.prop6.3.23]{6.3.23} that
condition (\ref{XVII.par6.3.24.1})
implies the analogous condition in which one assumes only that $f$ and $g$ are
flat, quasi-projective, and of finite presentation, and that $u$ is flat and
surjective. Indeed, locally for the étale topology, $u$ admits
finite locally free quasi-sections, and one concludes by \hyperref[XVII.prop6.3.23]{6.3.23} and
\textsc{Giraud} \cite{sga4.exp.XVII.bib.Gi}, applied to a diagram forming a square of descent data.

Given\pageoriginale a $G_{Y}$-torsor $K$ on a $T$-scheme $Y$ that is
flat, quasi-projective, and of finite presentation, there exists by definition
a faithfully flat morphism of finite presentation $u:X\rightarrow Y$
such that $u^{\ast}K$ is trivial.\nde{One may assume that $T$ is affine. Then there exists $u$ as stated, which may even be assumed affine so that the preceding paragraph can subsequently be applied.} If $u_{i}$ denotes the projection from
$(X/Y)^{i+1}$ to $Y$ $(i\geq 0)$, the torsors $(u^{\ast}_{i}K)_{n}$
(\hyperref[XVII.lem6.3.22.1]{6.3.22.1}) on the schemes $\Sym^{n}_{T}((X/Y)^{i+1})$ form a
descent datum $(i=0,1,2)$ on this semisimplicial scheme augmented
to $\Sym^{n}_{T}(Y)$. Suppose that $G$ satisfies the condition:

\paragraph*{(6.3.24.2)}\MakeLinkTarget*{sga4.exp.XVII.paragraph.6.3.24.2}\label{XVII.par6.3.24.2}
\emph{The descent data constructed above are all
effective} (\cf \ref{XVII.par6.3.24.1}).

Then we denote by $K_{n}$ the torsor on $\Sym^{n}_{T}(Y)$
(determined up to unique isomorphism) defined by this descent
datum.

\subsubsection{}\etiquette[6.3.25]{XVII.subsubsec6.3.25}
Suppose that $G$ satisfies (\hyperref[XVII.par6.3.21.1]{6.3.21.1}) or conditions
(\ref{XVII.par6.3.24.1}) and (\ref{XVII.par6.3.24.2}). We leave it to the
reader to make sense of the following formalism and to verify it:

\paragraph{}\etiquette[6.3.25.1]{XVII.par6.3.25.1}
\emph{Base change}. Let
\[
\xymatrix{
X' \ar[rr]^{g'} \ar[d]_{f'}& & X \ar[d]^{f}\\
T' \ar[rr]_{g}& & T.
}
\]
be a commutative square. Denote by $g'_{n}$ the arrow from $\Sym^{n}_{T'}(X')$ to
$\Sym^{n}_{T}(X)$. For $X$ a torsor under $G$ on $X$, one has
\[
{g'_n}^{\ast} K_{n}\simeq ({g'}^{\ast}K)_{n}{\quoi}.
\]

Now let\pageoriginale $f:X\rightarrow T$ be a morphism, and let $K$ be a
$G_{X}$-torsor on $X$.

\paragraph{}\etiquette[6.3.25.2]{XVII.par6.3.25.2}
\emph{Change of group.} Let $r:G_{1}\rightarrow G_{2}$ be a
morphism between groups both satisfying (\hyperref[XVII.par6.3.21.1]{6.3.21.1})
or (\ref{XVII.par6.3.24.1}), (\ref{XVII.par6.3.24.2}). One has
\[
r(K_{n})\simeq r(K)_{n}{\quoi}.
\]

\paragraph{}\etiquette[6.3.25.3]{XVII.par6.3.25.3}
\emph{Additivity in $K$}\sisi{:}{.} One has
\[
(K'+K'')_{n}\simeq K'_{n}+K''_{n}{\quoi}.
\]

\paragraph{}\etiquette[6.3.25.4]{XVII.par6.3.25.4}
\emph{Additivity in $n$}. Let $n_{1}$ and $n_{2}$ be two integers,
$n=n_{1}+n_{2}$, and let
\[
\alpha:\fprod{\Sym_{T}{}^{n_{1}}(X)}{\Sym_{T}{}^{n_{2}}(X)}{T}\longrightarrow
\Sym^{n}_{T}(X)
\]
be the canonical arrow; then
\[
\alpha^{\ast}K_{n}\simeq \pr^{\ast}_{1}
K_{n_{1}}+\pr^{\ast}_{2}K_{n_{2}}{\quoi}.
\]

\paragraph{}\etiquette[6.3.25.5]{XVII.par6.3.25.5}
\emph{Multiplicativity in $n$.} Let $n$ and $m$ be two integers, and let
$\alpha$ be the canonical arrow from $\Sym^{n}_{T}(\Sym^{m}_{T}(X))$ to
$\Sym^{nm}_{T}(X)$ (\cf \hyperref[XVII.lem6.3.3]{6.3.3}). One has
\[
\alpha^{\ast}(K_{nm})\simeq (K_{m})_{n}{\quoi}.
\]

\subsubsection{}\etiquette[6.3.26]{XVII.subsubsec6.3.26}
Suppose that $G$ satisfies (\hyperref[XVII.par6.3.21.1]{6.3.21.1}) or
(\ref{XVII.par6.3.24.1}),
(\ref{XVII.par6.3.24.2}). If
$f:X\rightarrow T$ is a finite locally free morphism of rank $n$
between $S$-schemes, if $t$ is the section \hyperref[XVII.eq6.3.4.2]{(6.3.4.3)} of
$\Sym^{n}_{T}(X)$, and if $K$ is a $G_{X}$-torsor on $X$, we
define a $G_{T}$-torsor $\Tr_{f}(K)$ on $T$ by the formula
$$\etiquette[6.3.26.1]{XVII.eq6.3.26.1}
\Tr_{f}(K)=t^{\ast}(K_{n}){\quoi}.\leqno(6.3.26.1)
$$

If\pageoriginale $f$ is locally free, we define $\Tr_{f}(K)$
locally on $T$ by the same formula. This definition, via a
canonical and functorial isomorphism, is equivalent to the definition
\hyperref[XVII.prop6.3.21]{6.3.21} on their common domain of validity
(the definition \hyperref[XVII.prop6.3.21]{6.3.21}
providing a solution to the descent problem posed). It follows
immediately from (\hyperref[XVII.par6.3.25.1]{6.3.25.1}), (\hyperref[XVII.par6.3.25.2]{6.3.25.2}),
and (\hyperref[XVII.par6.3.25.3]{6.3.25.3}) that the functor
$\Tr_{f}$ is \emph{compatible with base changes $T'\rightarrow
T$, compatible with changes of group $r:G_{1}\rightarrow G_{2}$,
and compatible with addition of torsors.}

The commutativity of the outer boundary of diagram \eqhyperref[XVII.eq6.3.15.2]{6.3.15.2}, together with
(\hyperref[XVII.par6.3.25.1]{6.3.25.1}) and (\hyperref[XVII.par6.3.25.4]{6.3.25.4}), yields
\emph{compatibility with composition of morphisms}:
$$
\Tr_{fg}(K)\simeq \Tr_{f}(\Tr_{g}(K)\sisi{}{)}{\quoi}.
$$

\subsubsection{}\etiquette[6.3.27]{XVII.subsubsec6.3.27}
In the particular case of relative curves, one also has:

\begin{enonce*}{Construction 6.3.27.1}\etiquette[6.3.27.1]{XVII.const6.3.27.1}
\emph{If} $f:X\rightarrow S$ is a smooth quasi-projective curve
over $S$, $D_{1}$ and $D_{2}$ are two relative divisors finite over $S$, and
$K$ is a $G$-torsor on $X$, then, setting
$\Tr_{D/S}(K)=\Tr_{D/S}(K|D)$, there is a canonical isomorphism
$$
\Tr_{D_{1}+D_{2}}(K)\simeq \Tr_{D_{1}}(K)+\Tr_{D_{2}}(K){\quoi}.
$$
\end{enonce*}

\sisi{\noindent}{}%
This isomorphism follows immediately from the isomorphism (\hyperref[XVII.par6.3.25.4]{6.3.25.4}) and
\hyperref[XVII.lem6.3.9.1]{6.3.9.1}. Finally, the arguments of \hyperref[XVII.prop6.3.19]{6.3.19} also give:

\begin{enonce*}{Construction 6.3.27.2}\etiquette[6.3.27.2]{XVII.const6.3.27.2}
Consider a commutative diagram
$$
\xymatrix{
X   \ar[rr]^{u} \ar[rd]_{p}&     &  Y \ar[ld]^{q}\\
     &   S &,
}
$$
in\pageoriginale which $X$ and $Y$ are smooth
quasi-projective curves over $S$ and $u$ is a quasi-finite morphism.
\begin{enumerate}
\renewcommand{\theenumi}{\roman{enumi}}
\renewcommand{\labelenumi}{(\theenumi)}
\item If $D$ is a divisor on $X$ finite over $S$ and $K$ is a torsor
under $G$ on $Y$, one has \emph{(\cf note p.~\pageref{XVII.note192})}:
$$
\Tr_{D/S}(u^{\ast}K)\simeq \Tr_{(u_{\ast}D)/S}(K){\quoi}.
$$

\item If $E$ is a divisor on $Y$ finite over $S$, $K$ is a torsor under
$G$ on $X$, and $u$ is finite, one has
$$
\Tr_{E/S}(\Tr_{X/Y}(K))\simeq \Tr_{u^{\ast}E/S}(K){\quoi}.
$$
\end{enumerate}
\end{enonce*}

\hypertarget{sga4.exp.XVII.proposition.6.3.28}{}\begin{enonce*}{Proposition 6.3.28}\etiquette[6.3.28]{XVII.prop6.3.28}
Let $A$ and $B$ be two (commutative) $R$-algebras and let $u:A\rightarrow
B$ be a homomorphism making $B$ a faithful, finite-type, locally free
$A$-module. Then the \sisi{}{co-}semisimplicial complex of algebras
$(TS^{n}_{R}(\bigotimes\limits^{k+1}_{A}B))_{k\geq 0}$
is, as a differential \sisi{}{co-}semisimplicial $TS^{n}_{R}(A)$-module,
a split resolution of $TS^{n}_{R}(A)$.
\end{enonce*}

\begin{enonce*}[remark]{Proof}
The complex of modules associated with the \sisi{}{co-}semisimplicial module
$TS^{n}_{R}(\bigotimes\limits^{k+1}_{A}B)$ depends only
on the $A$-module $B$, equipped with a module homomorphism $u:A\rightarrow
B$. As an $A$-module, $B$ can be written $B=A\oplus L$, and this
makes it possible to write down an explicit homotopy operator.
\end{enonce*}

Recall that:

\hypertarget{sga4.exp.XVII.proposition.6.3.29}{}\begin{enonce*}{Proposition 6.3.29}\etiquette[6.3.29]{XVII.prop6.3.29}
Let $A$ be a ring, let $B^{\ast}$ be a co-semisimplicial $A$-algebra,
and let $c:A\rightarrow B^{\ast}$ be the structure morphism. For the
inverse-image functor on modules ($=$ extension of scalars)
$\epsilon^{\ast}$ to be faithful (\resp fully faithful), it suffices that
the sequence of $A$-modules
$$
0\longrightarrow A\longrightarrow B^{0}
$$
(\resp\pageoriginale $0\rightarrow A\rightarrow
B^{0}\xrightarrow{j_{0}-j_{1}}B^{1}$)
be exact, and remain so after tensoring with any $A$-module.
\end{enonce*}

The proof is immediate and is left to the reader.

It follows in the standard way from (\hyperref[XVII.prop6.3.28]{6.3.28}) and (\hyperref[XVII.prop6.3.29]{6.3.29}) that

\hypertarget{sga4.exp.XVII.corollary.6.3.30}{}\begin{enonce*}{Corollary 6.3.30}\etiquette[6.3.30]{XVII.cor6.3.30}
Condition (\ref{XVII.par6.3.24.1}) is satisfied by every affine group $G$.
\end{enonce*}

(Use the fact that a torsor under $G$ is representable by a scheme
affine over $T$.)

It should not be difficult to extend \hyperref[XVII.cor6.3.30]{6.3.30} to every group $G$
representable by an algebraic space \footnote[1]{\MakeLinkTarget*{sga4.exp.XVII.footnote.6.3.30.artin-2}\label{XVII.noteArtin2}Algebraic Spaces
by M.~\textsc{Artin}. Yale University, Whittemore Lectures, May 1969.}. I
conjecture that condition (\ref{XVII.par6.3.24.2}) is satisfied by every group
representable by a flat algebraic space. We have the tongue-in-cheek
statement:

\hypertarget{sga4.exp.XVII.proposition.6.3.31}{}\begin{enonce*}{Proposition 6.3.31}\etiquette[6.3.31]{XVII.prop6.3.31}
Let the following sequence be exact on $S$:
$$
0\rightarrow G\xrightarrow{r}G'\xrightarrow{u}H\rightarrow 0{\quoi}.
$$ Suppose that $G'$ and $H$, and hence $G$, satisfy
(\hyperref[XVII.par6.3.13.1]{6.3.13.1}), that $G'$ satisfies (\hyperref[XVII.par6.3.21.1]{6.3.21.1}), and
that $G$ satisfies (\ref{XVII.par6.3.24.1}). Then $G$ satisfies
(\ref{XVII.par6.3.24.2}).
\end{enonce*}

This applies, for example, to an affine group that is the kernel of an epimorphism
of representable smooth groups.

% END INCLUDED SGA4_English_current_buildable_source_20260906/source/17/drafts/SGA4_XVII_LINES_8201_8600_English.texfrag

% BEGIN INCLUDED SGA4_English_current_buildable_source_20260906/source/17/drafts/SGA4_XVII_LINES_8601_9000_English.texfrag
\begin{enonce*}{Proof of 6.3.31}
\emph{Let\pageoriginale $K$ be a $G$-torsor on
$X$. By \hyperref[XVII.prop6.3.23]{6.3.23}, the question is local on $X$ for the
étale topology; it therefore suffices to treat the case where $r(K)$ is trivial. The
$G$-torsors $K$ such that $r(K)$ is trivial are identified with
trivial $G'$-torsors $K'$ equipped with an isomorphism
$k:H\isomto u(K')$. Put $Z=\Sym^{n}_{T}(X)$. The morphism
of trivial torsors $k$ induces $k_{n}:H_{Z}\isomto
u(K'_{n})$, and the corresponding $G$-torsor on $\Sym^{n}\sisi{}{_T}(X)$
represents the descent datum under consideration.}
\end{enonce*}

\section{Appendix}\etiquette[7]{XVII.sec7}\pageoriginale

\setcounter{subsection}{-1}
\subsection{Preliminaries}\etiquette[7.0]{XVII.subsec7.0}

In this appendix we shall use the results of \sisi{${\rm VI_{B}}$}{VBIS},
whose notation we retain.

\subsubsection{}\etiquette[7.0.1]{XVII.subsubsec7.0.1}
Let $\Sch_{!}$ be the category whose objects are quasi-compact,
quasi-separated schemes and whose morphisms are separated morphisms
of finite type. We denote by $\Ee$ the bifibred category over
$\Sch_{!}$, with fibre categories opposite to topoi, defined by étale
sheaves on schemes \sisi{${\rm VI_{B}}$ (4.3.0)}{(\hyperrefext[Vbis][Vbis.subsubsec4.3.0]{4.3.0})}. Fix once and for all an
object $R$ of $\Sch_{!}$, and denote by $\Ee_{R}$ the
bifibred category over $\Sch_{!R}$ deduced from
$\Ee$. Finally, let $\mathcal{A}$ be a ring in the topos
$\sheafGamma(\Ee^{\circ}_{R})$ satisfying the following conditions:
\begin{enumerate}
\renewcommand{\theenumi}{\roman{enumi}}
\renewcommand{\labelenumi}{(\theenumi)}
\item $\mathcal{A}$ is a \emph{cocartesian} section of
$\Ee^{\circ}_{R}$ over $\Sch^{\circ}_{!R}$.

\item For every scheme $X$ in $\Sch_{!R}$, $\mathcal{A}_{X}$ is a
\emph{torsion} sheaf.
\end{enumerate}

With this understood, we shall work with the category
$\sheaf{\Mod}(\sisi{}{\sheafGamma(}\Ee^{\circ}_{R}\sisi{}{)},\mathcal{A})$ (\refvb[VI B ]{Vbis.subsubsec1.3.5}).

Our aim is to define, for every morphism $f:X\rightarrow Y$ in
$\Sch_{!R}$, a functor $Rf_{!}:D(X,\mathcal{A}_X)\rightarrow
D(Y,\mathcal{A}_{Y})$, agreeing with the previously defined functor
when $f$ is compactifiable, and making the categories
$D(X,\mathcal{A}_{X})$, for varying $X$, into a cofibred category
over $\Sch_{!R}$.\nde{\MakeLinkTarget*{sga4.exp.XVII.editor-note.7.0.1.compactification}\label{XVII.ndeapp}Every morphism in $\Sch_{!}$ is compactifiable by Nagata's compactification theorem (see \nderef{XVII.ndeConrad}). This appendix shows how to avoid using the compactification theorem in the definition of $Rf_!$ by relying on Chow's lemma (\hyperref[XVII.lem7.0.6]{7.0.6}).}

\subsubsection{}\etiquette[7.0.2]{XVII.subsubsec7.0.2}
Although all the results of this appendix are stated in the setting
specified in \hyperref[XVII.subsubsec7.0.1]{7.0.1}, they could be transcribed, with obvious
modifications, into the following setting: take as a section of
$\sheafGamma(\Ee^{\circ}_{R})$ the constant section whose value is
the constant sheaf $\ZZ$, and, for every scheme $X$, restrict to the
derived category $D^{\ast}_{\tors}(X)$ defined by complexes of sheaves
of abelian groups whose cohomology sheaves are torsion sheaves
(\refvb[VI B ]{Vbis.subsubsec4.3.1}).

\begin{enonce*}{Lemma 7.0.3}\etiquette[7.0.3]{XVII.lem7.0.3}
Let\pageoriginale there be given a commutative diagram of schemes
$$\vcenter{\etiquette[7.0.3.1]{XVII.eq7.0.3.1}
\xymatrix{
X' \ar[rr]^{j'} \ar[d]_{f'}& & X \ar[d]^{f}\\
Y' \ar[rr]_{j}& & Y.
}}\leqno{(7.0.3.1)}
$$
where $f$ and $f'$ are proper and $j$ and $j'$ are immersions. Then,
for every complex $K^{\bigdot}$ in $D^{+}_{\tors}(X')$, the canonical
arrow\nde{The definition of this arrow is similar to that of \eqhyperref[XVII.eq5.1.5.2]{5.1.5.2}, up to an easy dévissage.}
$$
j_{!}\circ R^{+}f'\sisi{}{_*}(K^{\bigdot})\longrightarrow R^{+}f\sisi{}{_*}\circ
j_{!}(K^{\bigdot})
$$
is an isomorphism.
\end{enonce*}

By the ``lemma on way-out functors'' (\cite{sga4.exp.XVII.bib.RD} I \sisi{(}{}7.1), it
suffices to show that, for every integer $i$ and every torsion sheaf
$F$, the arrow
$$\etiquette[7.0.3.2]{XVII.eq7.0.3.2}
j_{!}\circ R^{i}f'\sisi{}{_*}(F)\longrightarrow R^{i}f\sisi{}{_*}\circ j_{!}(F)\leqno(7.0.3.2)
$$
is an isomorphism.

By replacing $j$ and $j'$ with their closed images, we are reduced to
the case where $j$ and $j'$ are open immersions with everywhere dense
images. The properness of $f$ and $f'$ then implies that diagram
\eqhyperref[XVII.eq7.0.3.1]{7.0.3.1} is cartesian. Under these conditions, the proper base-change
theorem, in the form \hyperrefext[XII][XII.cor5.2]{5.2}, shows that, for every geometric point
$\xi$ of $Y$, the arrow
$(j_{!}\circ R^{i}f'\sisi{}{_*}(F))_{\xi}\rightarrow (R^{i}f\sisi{}{_*}\circ
j_{!}(F))_{\xi}$ is an isomorphism.

\begin{enonce*}{Construction 7.0.4}\etiquette[7.0.4]{XVII.const7.0.4}
\emph{Let $D$ be a category, and let $X$ and $Y$ be two $D$-objects in
$\Sch_{!R}$ such that, for every arrow $\alpha\rightarrow \beta$ in
$D$, the morphisms $X_{\beta}\rightarrow X_{\alpha}$ and
$Y_{\beta}\rightarrow Y_{\alpha}$ are proper. Let $j:X\rightarrow
Y$ be a morphism such that, for every object $\alpha$ of $D$,
$j_{\alpha}:X_{\alpha}\rightarrow Y_{\alpha}$ is an immersion. The
isomorphism \eqhyperref[XVII.eq7.0.3.2]{7.0.3.2} for $i=0$ satisfies the usual compatibility
conditions and allows one to construct an \emph{exact} functor:}
\[
j_{!}:\Mod(\sheafGamma(\bar{X}),\mathcal{A})\longrightarrow
\Mod(\sheafGamma(\bar{Y}),\mathcal{A})
\]
\emph{(\cf \sisi{(}{}\refvb{Vbis.subsubsec1.2.5} for the notation)\sisi{.}{}
which\pageoriginale induces a triangulated functor, again denoted $j_{!}$:}
\[
j_{!}:D(\sheafGamma(\bar{X}),\mathcal{A})\longrightarrow
D(\sheafGamma(\bar{Y},\mathcal{A}))
\]

\emph{Then, for every object $\alpha$ of $D$, there is a commutative
diagram, where $j_{\alpha !}$ denotes the usual ``extension by $0$''
functor:}
\[
\xymatrix{
D(\sheafGamma(\bar{X}), \sisi{}{\Aa}) \ar[rr]^{j_!} \ar[d]_{R e^{\sisi{+}{*}}_{\alpha}}&
&D(\sheafGamma(\bar{Y}),\Aa) \ar[d]^{R e^{\sisi{+}{*}}_{\alpha}}\\
D(X_{\alpha}, \Aa_{X_{\alpha}})  \ar[rr]_{j_{\alpha !}} & &
  D(Y_{\alpha},\Aa_{Y_\alpha}).
}
\]
\end{enonce*}

\begin{enonce*}{Construction 7.0.5}\etiquette[7.0.5]{XVII.constru7.0.5}
Let there be given a commutative diagram in $\Sch_{!R}$
$$\vcenter{\etiquette[7.0.5.1]{XVII.eq7.0.5.1}
\xymatrix{
X'   \ar[rr]^{i'} \ar[d]_{f'}&    & X \ar[d]^{f}\\
Y'   \ar[rr]^{i} &    & Y
}}\leqno{(7.0.5.1)}
$$
where $i$ and $i'$ are immersions, $f$ is a proper morphism, and $f'$
is a proper \emph{surjective} morphism. It defines a commutative diagram
of simplicial objects in $\Sch_{!R}$
$$\vcenter{\etiquette[7.0.5.2]{XVII.eq7.0.5.2}
\xymatrix{
[X'|Y']      \ar[d]_{u_{f'}}  \ar[r]^{j} & [X|Y] \ar[d]^{u_f}\\
C^{\Delta}_{Y'}  \ar[r]_{i} & C^{\Delta}_{Y}
}}\leqno{(7.0.5.2)}
$$
where $j_{n}:[X'|Y']_{n}\rightarrow [X|Y]_{n}$ is an immersion
for every $n$. With this notation, there is a canonical isomorphism
\[
i_{!}\isomto  R^{+}\bar{u}_{f\sisi{+}{*}}\circ j_{!}\circ
L^{+}\bar{u}_{f'}\sisi{}{^*}{\quoi}.
\]
\end{enonce*}

By replacing $i$ and $i'$ with their closed images, we may suppose that
$i$ and $i'$ are open immersions with everywhere dense images. Diagram
\eqhyperref[XVII.eq7.0.5.1]{7.0.5.1} is then cartesian and $f$ is \emph{surjective}. Moreover,
for every integer $n$, $j_{n}$ is an\pageoriginale open immersion and
the diagram
\[
\xymatrix{
[X'|Y']_n        \ar[r]^{j_n} \ar[d]& [X|Y]_n \ar[d]\\
Y'  \ar[r]_{i} & Y
}
\]
is cartesian.

Consequently, $L^{+}\bar{u}_{f}^{\sisi{+}{*}}\circ i_{!}$ is isomorphic to $j_{!}\circ
L^{+}\bar{u}_{f'}^{\ast}$. Using the fact that $u_{f}$ is a
cohomological-descent augmentation (\sisi{${\rm VI_{B}}$ (\refnoexpose{Vbis.prop4.3.2})}{\refvb{Vbis.prop4.3.2}}), we
obtain a sequence of isomorphisms
\[
i_{!}\isomto R^{+}\bar{u}_{f\sisi{+}{*}}\circ
L^{+}(\bar{u}^{\sisi{+}{*}}_{f})\circ i_{!}\isomto
R^{+}\bar{u}_{f\ast}\circ j_{!}\circ L^{+}\bar{u}_{f'}^{\ast}
\]
which gives the desired isomorphism.\nde{More directly, without assuming $f'$ surjective, there is still a sequence of canonical arrows
\[i_!\to i_! R^+\bar{u}_{f'*}L^+\bar{u}^*_{f'}\isomto  R^+\bar{u}_{f*} j_! L^+u^*_{f'},\]
of which the second is an isomorphism (\hyperref[XVII.lem7.0.3]{7.0.3}). When $f'$ is surjective, the first is also an isomorphism by cohomological descent, and the composite agrees with the isomorphism constructed in the text.
}

Finally, we shall use the following variant of Chow's lemma, whose
proof will appear in \EGA II, 2\sisi{nd}{\textsuperscript{nd}} edition,\nde{See [C, 2.6], cited in \nderef{XVII.ndeConrad}.} and which is an
improvement of (\sisi{\hyperref[XVII.subsec7.1]{7.1}}{\hyperrefext[XII][XII.lem7.1]{7.1}}):

\begin{enonce*}{Lemma 7.0.6}\etiquette[7.0.6]{XVII.lem7.0.6}
Let $S$ be a quasi-compact and quasi-separated scheme and let
$f:X\rightarrow S$ be a separated morphism of finite type. Then there
exists a commutative diagram
\[
\xymatrix{
X' \ar[rr]^{p} \ar[d]_{f'}& & X \ar[lld]^{f}\\
S                         & &
}
\]
satisfying the following conditions:
\begin{enumerate}
\renewcommand{\theenumi}{\roman{enumi}}
\renewcommand{\labelenumi}{(\theenumi)}
\item $p$ is proper and \emph{surjective};

\item there is a dense open subset $U$ of $X$ such that $p$ induces an
isomorphism from $p^{-1}(U)$ onto $U$;

\item $f'$ is quasi-projective.
\end{enumerate}
\end{enonce*}


\subsection{Definition of $R^{+}f_{!}$}\etiquette[7.1]{XVII.subsec7.1}\pageoriginale

\begin{enonce*}{Definition 7.1.1}\etiquette[7.1.1]{XVII.def7.1.1}
Let $f:X\rightarrow S$ be a morphism of schemes. A
\emph{pseudo-compactification} of $f$ consists of a commutative diagram
$$\vcenter{
\xymatrix{
X_1 \ar[r]^{\sisi{}{j_1}} \ar[d]_{p_1}& T_1\ar[ldd]^{f_1}\\
X   \ar[d]_{f}& \\
S   &
}}\leqno{(7.1.1.0)}
$$\etiquette[7.1.1.0]{XVII.eq7.1.1.0}
where $p_{1}$ is proper and surjective, $j_{1}$ is an immersion, and
$f_{1}$ is proper.
\end{enonce*}

\setcounter{subsubsection}{1}
\subsubsection{}\etiquette[7.1.2]{XVII.subsubsec7.1.2}
Let $f:X\rightarrow S$ be a morphism in $\Sch_{!R}$. By Chow's lemma
(\hyperref[XVII.lem7.0.6]{7.0.6}), $f$ has a pseudo-compactification
$(p_{1},j_{1},f_{1})$. Diagram \eqhyperref[XVII.eq7.1.1.0]{7.1.1.0} defines a morphism
\[
[X_{1}|X]\xrightarrow{j_{1}}[T_{1}|S]
\]
again denoted by $j_{1}$, such that $j_{1n}$ is an immersion for every
$n$. We then put $R^{+(1)}f_{!}=\sisi{L^{+}\bar{u}_{p_{1}}\cdot
j_{1!}\circ R^{+}\bar{u}_{f_{1}}}{R^+\bar{u}_{f_1*}\circ j_{1!} \circ L^+\bar{u}_{p_1}^*}$, thereby obtaining a
\emph{triangulated} functor from $D^{+}(X,\mathcal{A}_{X})$ to
$D^{+}(S,\mathcal{A}_{S})$.

\begin{enonce*}{Construction 7.1.3}\etiquette[7.1.3]{XVII.constru7.1.3}
For every pair $((p_{1},j_{1},f_{1}),(p_{2},j_{2},f_{2}))$ of
pseudo-compactifications of $f$, there is an isomorphism
$\alpha_{12}:R^{+(1)}f_{!}\rightarrow R^{+(2)}f_{!}$ such that, for
every other pseudo-compactification $(p_{3}, j_{3},f_{3})$, one has
$\alpha_{23}\circ \alpha_{12}=\alpha_{13}$.
\end{enonce*}

\begin{enumerate}
\renewcommand{\theenumi}{\alph{enumi}}
\renewcommand{\labelenumi}{\theenumi)}
\item Define a bisimplicial system $[[X_{1},X_{2}|X]]$ by setting
\[
[[X_{1},X_{2}|X]]_{[n]\times [m]}=[X_{1}|X]_{n}\times
[X_{2}|X]_{n}
\]
Similarly, set
\[
[[X_{1}|X]]_{[n]\times[m]}=[X_{1}|S]_{n}{\quoi},
\]
all the morphisms $[[X_{1}|X]]_{[n]\times [m]}\rightarrow
[[X_{1}|X]]_{[n]\times [m']}$ induced by morphisms
$[m']\rightarrow [m]$ being equal to the identity. Define the
bisimplicial systems $[[T_{1},T_{2}|S]]$ and $[[T_{1}|S]]$ in the same
way, so that there is a commutative diagram
\[
\xymatrix@C=1em{
  & & [[T_1,T_2|S]] \ar[rrrr]\ar'[dr][rrrddd]_(.2){u_{f_1f_2}}|(.38)\hole& &   & &  [[T_1|S]]\ar[lddd]^{u_{f_1}}\\
  [[X_1,X_2|X]]
  \ar[rrrr]\ar[rru]^{j_{12}}\ar[rrrddd]_{u_{p_1 p_2}} & &  & & [[X_1|X]]\ar[rru]_{j_1}\ar[lddd]_(.75){u_{p_1}}  \\
  & & & & &   &\\
  & & & & & C^{\Delta\times\Delta}_S\\
  &&& C^{\Delta\times\Delta}_X \ar[rru]
}
\]
where\pageoriginale $j_{12}$ is an immersion degree by degree. The
computational procedure described in \sisi{VI b (\refnoexpose{Vbis.subsec2.3})}{\hyperrefext[Vbis][Vbis.subsec2.3]{2.3}} makes it possible to construct a
morphism
\[
R^{+(1)} f_{!}\xrightarrow{\alpha_{1}}R^{+}\bar{u_{f_{1}f_{2}}}\sisi{}{_*}\circ
j_{12!} \circ L^{+}\bar{u_{p_{1}p_{2}}}^{\sisi{+}{\ast}}=R^{{+}(1,2)}f_{!}
\]
which is obtained from a morphism of triple complexes of sheaves on
$S$.

\item $\alpha_{1}$ \emph{is an isomorphism}:

To prove this, one may fix the first index $n$ coming from
$\bigl[[T_{1},T_{2}]\bigr]_{[n]\times[m]}$ in the preceding morphism of
triple complexes and consider the resulting morphism of double
complexes. One checks that this morphism is precisely the isomorphism
defined in Lemma \hyperref[XVII.constru7.0.5]{7.0.5} for the diagram
\[
\xymatrix{
**[r][X_1|X]_n \times X_2  \ar[r] \ar[d] & **[r][T_1|S]_n \times
T_2 \ar[d]\\
**[r][X_1|X]_n             \ar[r] & **[r][T_1|S]_n
}
\]

\item In the same way, construct an isomorphism
\[
R^{+(2)}f_{!}\xrightarrow{\alpha_{2}}R^{+}\bar{u_{f_{1}f_{2}}}_{\sisi{+}{*}}\circ
j_{12!} \circ L^{+}\bar{u_{p_{1}p_{2}}}\sisi{}{^*}
\]
and set $\alpha_{12}=\alpha^{-1}_{2}\circ \alpha_{1}$.
\end{enumerate}

To verify the relation $\alpha_{23}\circ \alpha_{12}=\alpha_{13}$,
one uses a trisimplicial object. We leave it to the reader to check
that the resulting diagram is commutative:\pageoriginale
$$
\xymatrix{
  R^{+(2,3)}{f_!} \ar[rd]& R^{+(2)}f_! \ar[r] \ar[l]& R^{+(1,2)}f_! \ar[ld]\\
                &         R^{+(1,2,3)}f_!  &                  \\
R^{+(\sisi{2}{3})}f_! \ar[uu]\sisi{}{\ar[rd]}     &                         &  R^{+(1)}f_! \ar[uu]\sisi{}{\ar[ld]}\\
                &    R^{+(1,3)}f_! \ar[uu]        &
}
$$

\begin{enonce*}{Construction 7.1.4}\etiquette[7.1.4]{XVII.constru7.1.4}
Let $Z\xrightarrow{g}Y$ and $Y\xrightarrow{f}X$ be two morphisms in
$\Sch_{!R}$. For every triple $((p_{1},j_{1},f_{1}),
(p_{2},j_{2}g_{2}), (p_{3},j_{3},(fg)_{3}))$, where
$(p_{1},j_{1},f_{1})$ is a pseudo-compactification of $f$,
$(p_{2},j_{2},g_{2})$ a pseudo-compactification of $g$, and
$(p_{3},j_{3},(fg)_{3})$ a pseudo-compactification of $fg$, there is an
isomorphism $\sisi{R^{+(3)}fg_{!}\isomfrom  R^{+(2)}g_{!}\circ
R^{+(1)}f_{!}}{R^{+(1)}f_!\circ R^{+(2)}g_!\isomto R^{+(3)}(fg)_!}$ compatible with the previously defined
isomorphisms $\alpha_{ij}$.
\end{enonce*}

We proceed in several steps.

\setcounter{subsubsection}{4}
\paragraph{}\etiquette[7.1.4.1]{XVII.par7.1.4.1}
Let $(p_{1},j_{1},f_{1})$ (\resp $(p_{2},j_{2},f_{2})$) be a
pseudo-compactification of $f:Y\rightarrow X$ (\resp of
$g:Z\rightarrow Y$):
$$\vcenter{\etiquette[7.1.4.1.1]{XVII.eq7.1.4.1.1}
\xymatrix{
  &     & Z\ar[d]^{g}  & Z_2\ar[l]_{p_2}\ar@{^{(}->}[r]^{j_2}  & T_2 \ar[lld]^{g_2}\\
T_1  \ar[rrd]_{f_1}& Y_1 \ar[r]^{p_1} \ar@{_{(}->}[l]_{j_1}& Y \ar[d]^{f} &      &    \\
    &     & X        &      &
}}\leqno{(7.1.4.1.1)}
$$

The morphism $Z'=\fprod{Z}{Y_{1}}{Y}\rightarrow Y_{1}$ obtained from
$g$ by base change is equipped with a pseudo-compactification:
$$\vcenter{\etiquette[7.1.4.1.2]{XVII.eq7.1.4.1.2}
\xymatrix{
  T'_2 = \fprod{T_2}{Y_1}{Y}   \ar[rdd]_{g'_2}& \fprod{Z_2}{Y_1}{Y} =Z'_2 \ar[d]^{p'_2}
\ar@{_{(}->}[l]_{j'_2}\\
&  \fprod{Z}{Y_1}{Y} = Z' \ar[d]^{g'}\\
&    Y_1
}}\leqno{(7.1.4.1.2)}
$$


Now\pageoriginale choose a \sisi{}{pseudo-}compactification of
\sisi{$j_{1}\cdot g'_{2}$.
$$\etiquette[7.1.4.1.3]{XVII.eq7.1.4.1.3}
\vcenter{
\xymatrix{
T_{12}        \ar[d]_{g''_2} & \fprod{T_2}{Y_1}{Y}=T'_2
\ar@{_{(}->}[l]_(.6){j'_1} \ar[d]^{g'_2}\\
T_1            &   Y_1 \ar@{_{(}->}[l]
}}\leqno(7.1.4.1.3)
$$%
}{$j_1\circ g'_2$\nde{In the original text a compactification was chosen, which is possible by the compactification theorem. Since the purpose of the appendix is to avoid that theorem (\cf \nderef{XVII.ndeapp}), the editor has taken the liberty of replacing it with a pseudo-compactification and slightly modifying \hyperref[XVII.par7.1.4.2]{7.1.4.2}. Another solution would be to restrict to pseudo-compactifications $(p_1,j_1,f_1)$ for which $f_1$ is \emph{projective}.}
$$\etiquette[7.1.4.1.3]{XVII.eq7.1.4.1.3}
\vcenter{
\xymatrix{
T_{12}        \ar[dd]_{g''_2} &T_2''\ar[d]\ar[d]^p\ar@{_{(}->}[l]_{j'_1}\\
& \fprod{T_2}{Y_1}{Y}=T'_2\ar[d]^{g'_2}\\
T_1            &   Y_1 \ar@{_{(}->}[l]_{j_1}
}}\leqno(7.1.4.1.3)
$$}%
which defines a pseudo-compactification of the composite $fg$; denote
the corresponding functor by $R^{+(12)}\sisi{\int g}{(fg)}_{!}$.

\paragraph{}\etiquette[7.1.4.2]{XVII.par7.1.4.2}
\emph{The preceding data make it possible to define an isomorphism}
$R^{+(12)}\sisi{\int g}{(fg)}_{!}\isomfrom  R^{+(1)}f_{!}\sisi{\cdot}{\circ} R^{+(2)}g_{!}$:

There is a commutative diagram, where $Z_2''=\fprod{T''_2}{Z'_2}{T'_2}$ and all the arrows are
naturally deduced from \eqhyperref[XVII.eq7.1.4.1.1]{7.1.4.1.1}, \eqhyperref[XVII.eq7.1.4.1.2]{7.1.4.1.2}, and
\eqhyperref[XVII.eq7.1.4.1.3]{7.1.4.1.3}.
$$\etiquette[7.1.4.2.1]{XVII.eq7.1.4.2.1}
\vcenter{
\xymatrix{
[T_{12}|X] \ar[rd]_{g''_2} &   &[T\sisi{'}{''}_2|Y] \ar[rd]_{g'_2\sisi{}{p}} \ar[ll]_{j'_1}&
[Z\sisi{'}{''}_{2}|Z] \ar[r]^{r}\ar[l]_{j\sisi{'}{''}_2} \ar[d]_{l}&
C^{\Delta}_Z \ar[d]^{g}\\
            &    [T_1|X] \ar[rrrd]_{u_{f_1}} &  & [Y_1|Y]\ar[r]_{u_{p_1}}\ar[ll]_{j_1} & C_Y^{\Delta} \ar[d]^{f}\\
            &       &  &   &  C_X^{\Delta}
}}\leqno(7.1.4.2.1)
$$

We shall\pageoriginale define a natural arrow:

% END INCLUDED SGA4_English_current_buildable_source_20260906/source/17/drafts/SGA4_XVII_LINES_8601_9000_English.texfrag

% BEGIN INCLUDED SGA4_English_current_buildable_source_20260906/source/17/drafts/SGA4_XVII_LINES_9001_9400_English.texfrag
\[
\beta:R^{+(1)}f_{!}\circ R^{+(2)}g_{!}\longrightarrow R^{+(12)}\sisi{fg}{(fg)}_{!}
\]
which we shall show to be an isomorphism.

For this we shall use two lemmas:

\begin{enonce*}{Lemma 7.1.4.2.2}\etiquette[7.1.4.2.2]{XVII.lem7.1.4.2.2}
Let there be given a commutative diagram:
\[
\xymatrix{
X_1    \ar[d]_{p_1}     \ar[r]^{j_1} & T_1 \ar[d]^{h'}\\
X    \ar[d]_{f}       \ar[r]_{j}           & T \ar[ld]^{h}\\
    S       &
}
\]
where $h$ and $h'$ are proper, $p_{1}$ is proper and surjective, and
$j$ and $j_{1}$ are immersions. Then the canonical arrow
\[
R^{+}h\sisi{}{_*}\circ j_{!}\longrightarrow R^{+(1)}\sisi{(hh')}{f}_{!}
\]
is an isomorphism.
\end{enonce*}

Indeed, by (\hyperref[XVII.constru7.0.5]{7.0.5}), the arrow
\[
j_{!}\longrightarrow R^{+}\bar{u}_{h'\sisi{}{*}}\circ j_{1!}\circ
L^{+}\bar{u}_{p_{1}}\sisi{}{^*}
\]
is an isomorphism. The result follows by applying $R^{+}h\sisi{}{_*}$.

\begin{enonce*}{Lemma 7.1.4.2.3}\etiquette[7.1.4.2.3]{XVII.lem7.1.4.2.3}
With the notation of diagram \emph{\eqhyperref[XVII.eq7.1.4.2.1]{7.1.4.2.1}}, the canonical arrows
\begin{gather*}
R^{+(1)}f_{!}\circ R^{+}\bar{u}_{p_{1}\sisi{}{*}}\longrightarrow
R^{+}\bar{u}_{f_{1}\sisi{}{*}}\circ j_{1!}\\
R^{+(2)}g_{!}\circ R^{+}\bar{r}\sisi{}{_*}\longrightarrow
R^{+}\bar{u}_{p_{1}\sisi{}{*}}\circ R^{+}\sisi{g'_{2}}{(g_2'p)_*}\circ j\sisi{'}{''}_{2!}
\end{gather*}
are isomorphisms.
\end{enonce*}

We prove only the first isomorphism, leaving it to the reader to
deduce\pageoriginale the second.

With the evident notation\sisi{, put $C=\fprod{[Y_{1}|Y]}{X_{1}}{X}$
and $\bar{C}=\fprod{[T_{1}|S]}{T_{1}}{S}$, so that there is a
commutative diagram:
\[
\xymatrix{
                                & \bar{C} \ar[rdd]|\hole \ar[rr]  & & T_1 \ar[ldd]^{f_1}\\
C                 \ar[rr]  \ar[ru]^{i}  \ar[d]^{q}    &    &  Y_1  \ar[ru]_{j_1} \ar[d]^{p_1}&      \\
[Y_1|Y] \ar[rr]_{u_{p_1}}&  &      Y         &
}
\]
where $i$ is an immersion and $q$ is proper and surjective. One then
constructs a diagram
\[
\xymatrix{
\bigl[\bigl[C \underset{q}{\bigm|}[Y_1|Y]\bigr]\bigr] \ar[r]^{i} \ar[d]^{u_q}&
[[\bar{C}|C^{\Delta}_Y]] \ar[ldd]^{k}\\
|X_1 | X |                             \ar[d]         &\\
Y                       &                        &
}
\]
and it follows from (\hyperref[XVII.lem7.1.4.2.2]{7.1.4.2.2}) that the functor $R\bar{u}j_{1}\circ
j_{1}$ is identified with the functor $R^{+}\bar{k}\circ i_{!}\circ
L^{+}\bar{u}_{q}$ and, by (\hyperref[XVII.lem7.0.3]{7.0.3}) and the base-change theorem, the latter functor is identified with $R^{+(1)}f_{!}\circ
R^{+}\bar{u}_{p_{1}}$. }{ and the notation of \hyperref[XVII.constru7.1.3]{7.1.3} a), there is a commutative diagram
\[\xymatrix{&[Y_1|Y]\ar[rr]^{j_1}\ar[dd]^(.7){\bar{u}_{p_1}}|\hole &&[T_1|X]\ar[dd]^{\bar{u}_{f_1}}\\
[[Y_1,Y_1|Y]]\ar[ur]^{k}\ar[dd]_{q}\ar[rr]^(.7)i && [[T_1,T_1|X]]\ar[dd]\ar[ur]\\
&Y\ar[rr]^(.3){f}|\hole && X\\
[Y_1|Y]\ar[ur]^{\bar{u}_{p_1}}\ar[rr]^{j_1} &&[T_1|X]\ar[ur]_{\bar{u}_{f_1}} }\]
where $q$ and $k$ are the two projections from $[[Y_1,Y_1|Y]]$ to $[Y_1|Y]$. By cohomological descent, $R^+\bar{u}_{f_1 *} \circ j_{1!}$ is identified with $R^+\bar{u}_{f_1 *} \circ j_{1!} \circ R^+ q_* L^+ q^*$, which, by \hyperref[XVII.lem7.0.3]{7.0.3} applied to the front and top faces of the cube above, is identified with $R^+\bar{u}_{f_1*}\circ j_{1!}\circ R^+k_* \circ L^+q^*$. By the base-change theorem, this last functor is identified with $R^+\bar{u}_{f_1 *} \circ j_{1!} \circ L^+\bar{u}_{p_1*} \circ R^+\bar{u}_{p_1 *}$, which by definition is $R^{+(1)}f_{!}\circ
R^{+}\bar{u}_{p_{1}*}$.\nde{the proof in the original text is slightly different %and not very clear
and uses \hyperref[XVII.lem7.1.4.2.2]{7.1.4.2.2}.}

}%
Lemma (\hyperref[XVII.lem7.1.4.2.3]{7.1.4.2.3}) having been proved, there is a
sequence of isomorphisms
\begin{align*}
R^{+(\sisi{3}{12})} \sisi{fg}{(fg)}_{!} &= R^{+}\bar{u}_{f_{1}\sisi{}{*}}\circ R^{+}g''_{2}\circ
j'_{1!}\circ j\sisi{'}{''}_{2!}\circ L^{+}\bar{r}^{\sisi{+}{*}}\\
&\isomfrom  R^{+}\bar{u}_{f_{1}\sisi{+}{*}}\circ j_{1!}\circ R^{+}
\sisi{g'_{2}}{(g'_2p)_*}\circ j'_{2!}\circ L^{+}\bar{r}^{\sisi{+}{*}}\\
&\isomfrom  R^{+(1)}f_{!}\circ R^{+}\bar{u}_{p_{1}\sisi{+}{*}}\circ
R^{+}\sisi{g'_{2+}}{(g'_2p)_*}\circ j'_{2!}\circ L^{+}\bar{r}\sisi{}{^*}\\
&\isomfrom  R^{+(1)}f_{!}\circ R^{+(2)}g_{!}\circ
R^{+}\bar{r}\sisi{}{_*}\circ L^{+}\bar{r}^{\sisi{+}{*}}\\
&\isomfrom  R^{+(1)}f_{!}\circ R^{+(2)}g_{!}
\end{align*}
whose\pageoriginale composite gives $\beta$.

\paragraph{}\etiquette[7.1.4.3]{XVII.par7.1.4.3}
If $(p_{3},j_{3},(fg)_{3})$ is a pseudo-compactification of $fg$, an
isomorphism
\[
R^{+(1)}f_{!}\circ
R^{+(2)}g_{!}\xrightarrow{\gamma_{123}}R^{+(3)}fg_{!}
\]
is defined by composing $\alpha_{(12)3}$ with $\beta$.

\paragraph{}\etiquette[7.1.4.4]{XVII.par7.1.4.4}
Now let $(p_{1'},j_{1'},f_{1'})$ and $(p_{2'},j_{2'},g_{2'})$ be two
other pseudo-compactifications of $f$ and $g$, together with a diagram
\sisi{(\hyperref[XVII.par7.1.4.3]{7.1.4.3})}{\eqhyperref[XVII.eq7.1.4.1.3]{7.1.4.1.3}}; one checks that the following diagram is commutative:
\[
\xymatrix{
R^{+(12)}\sisi{fg}{(fg)}_1 \ar[d]_{\alpha_{(12)(1'2')}} & R^{+(1)}f_1 \circ
R^{+(2)}g_! \ar[d]^{\alpha_{11'}\ast \alpha_{22'}} \ar[l]_(.6){\beta}\\
R^{+(1'2')}\sisi{fg}{(fg)}_!  & R^{+(1')}f_! \circ R^{+(2')}g_! \ar[l]^(.6){\beta'}
}
\]
by considering a diagram of the type \sisi{(\hyperref[XVII.par7.1.4.4]{7.1.4.4})}{\eqhyperref[XVII.eq7.1.4.2.1]{7.1.4.2.1}} in which the
simplicial objects are replaced by bisimplicial objects.

It follows that the following diagram is commutative:
\[
\xymatrix{
R^{+(3)}\sisi{fg}{(fg)}_! \ar[d]_{\alpha_{33'}}& R^{+(1)}f_! \circ R^{+(2)}g_!
\ar[l]_(.6){\gamma_{123}} \ar[d]^{\alpha_{11'}\ast
  \alpha_{22'}}\\
R^{+(3')}\sisi{fg}{(fg)}_!  & R^{+(1')}f_! \circ R^{+(2')}g_! \ar[l]_(.6){\gamma_{1'2'3'}}
}
\]
for every pseudo-compactification $(p_{3'},j_{3'},(fg)_{3'})$ of
$fg$, which completes the proof of \hyperref[XVII.constru7.1.4]{7.1.4}.

\subsubsection{}\etiquette[7.1.5]{XVII.subsubsec7.1.5}
By choosing a pseudo-compactification for every morphism $f$ in
$\Sch_{!R}$, in such a way that the trivial pseudo-compactification is
associated with an identity morphism, one defines, for every morphism
$f:X\rightarrow Y$, a triangulated functor
\[
R^{+}f_{!}:D^{+}(X,\mathcal{A}_{X})\longrightarrow
D^{+}(Y,\mathcal{A}_{Y})
\]
and,\pageoriginale for every pair $(f,g)$ of morphisms, an isomorphism
\[
R^{+}fg_{!}\xrightarrow{\substack{c_{fg}\\\sim}}
R^{+}f_{!}\circ R^{+}g_{!}
\]
Moreover, if $f$ is an identity, $R^{+}f_{!}$ is the identity functor
and $c_{fg}$ is the identity.

\begin{enonce*}{Proposition 7.1.6}\etiquette[7.1.6]{XVII.prop7.1.6}
The correspondence $f\rightarrow R^{+}f_{!}$ defines a normalized
pseudofunctor from $\Sch_{!R}$ to the $2$-category of categories. In
other words, the categories $D^{+}(X,\mathcal{A}_{X})$, for $X$ an
object of $\Sch_{!R}$, can be organized into a cofibered category over
$\Sch_{!R}$ by setting, for every morphism $f:X\rightarrow Y$ in
$\Sch_{!R}$,
\[
\Hom_{f}(K,L)=\Hom(R^{+}f_{!}K,L)
\]
with $K\in D^{+}(X,\mathcal{A}_{X})$ and $L\in D^{+}(Y,\mathcal{A}_{Y})$.
\end{enonce*}

It remains to see that the identification
$R^{+}fg_{!}\xrightarrow{\substack{c_{fg}\\\sim}}
R^{+}f_{!}\circ R^{+}g_{!}$ is compatible with triple composites. Let
$h:T\rightarrow Z$, $g:Z\rightarrow Y$, and $f:Y\rightarrow X$ be
three morphisms in $\Sch_{!R}$; by what precedes, to verify this
compatibility it suffices to verify it after an arbitrary choice of
pseudo-compactifications for $h$, $g$, and $f$.

There is a diagram:
$$\vcenter{\etiquette[7.1.6.1]{XVII.eq7.1.6.1}
  \xymatrix{
 T_5 \ar[rd]& T_4 \ar[rd] \ar@{_{(}->}[l]& T_3 \ar[rd]\ar@{_{(}->}[l]
 & T'\ar@{_{(}->}[l] \ar[r]^{p_3} \ar[d]& T \ar[d]^{h} \\
        & T_3 \ar[rd]  & T_2 \ar@{_{(}->}[l] \ar[rd]  & Z'\ar[r]^{p_2}\ar@{_{(}->}[l] \ar[d]& Z
  \ar[d]^{g}\\
        &    & T_1 \ar[rrd]& Y_1\ar@{_{(}->}[l] \ar[r]^{p_1}& Y \ar[d]^{f}\\
        &    &     &     & X
}}\leqno{(7.1.6.1)}
$$
where all the oblique arrows are proper and $p_{1}$, $p_{2}$, and
$p_{3}$ are proper and surjective.

A diagram of the type (\hyperref[XVII.par7.1.4.4]{7.1.4.4}) constructed from
\eqhyperref[XVII.eq7.1.6.1]{7.1.6.1} then gives the conclusion.

\begin{enonce*}{Proposition 7.1.7}\etiquette[7.1.7]{XVII.prop7.1.7}
If\pageoriginale $f:X\rightarrow S$ is a compactifiable morphism, the
functor $Rf_{!}$ coincides with the one defined in Exposé \emph{XVII}.
\end{enonce*}

Indeed, there is a commutative diagram:
\[
\xymatrix{
X_1         \ar[r]^{j_1} \ar[d]_{p_1}& T_1 \ar[d]^{h'}\\
X           \ar[r]_{j}  \ar[d]_{f} & T \ar[ld]^{h}\\
S
}
\]
to which one applies (\hyperref[XVII.lem7.1.4.2.2]{7.1.4.2.2}).

\begin{enonce*}{Definition 7.1.8}\etiquette[7.1.8]{XVII.def7.1.8}
If $f:X\rightarrow Y$ is a morphism in $\Sch_{!}$, the functor
$R^{+}f_{!}$ just defined is called the \emph{direct-image functor with
proper supports} (in the sense of derived categories). The functor
$H^{q}\circ R^{+}f_{!}$, denoted $R^{q}f_{!}$, is called the
$q$\textsuperscript{th} \emph{direct-image functor with proper supports.}
\end{enonce*}

With this notation, we leave it to the reader to verify the

\begin{enonce*}{Construction 7.1.9}\etiquette[7.1.9]{XVII.constru7.1.9}
Let $(\pi_{1}, j_{1}, f_{1})$ be a pseudo-compactification of
$f:X\rightarrow Y$, and denote by $\pi^{n}_{1}$ the canonical morphism
$[X_{1}|X]_{n}\rightarrow X$. Then, for every complex $K^{\bigdot}$ in
$D^{+}(X,\mathcal{A}_{X})$, there are two biregular spectral sequences:
\begin{align*}
E^{pq}_{1} &= R^{q}f_{!}(K^{p})\Longrightarrow R^{p+q}
f_{!}(K^{\bigdot})\tag{7.1.9.1}\\
E^{pq}_{1} &= R^{q}(f\circ
\pi^{p}_{1})_{!}(\pi^{p\sisi{+}{*}}_{1}(K^{\bigdot}))\Longrightarrow
R^{p+q}f_{!}(K^{\bigdot})\tag{7.1.9.2}
\end{align*}\etiquette[7.1.9.1]{XVII.eq7.1.9.1}
\etiquette[7.1.9.2]{XVII.eq7.1.9.2}%The vertical space following this paragraph is too large.
\end{enonce*}

(Take, for example, the canonical flabby resolution of $j_{1!}\circ
L^{+}\bar{u}_{p_{1}}^{\sisi{+}{*}}(K^{\bigdot})$ (\cf XVII (\hyperref[XVII.subsec4.2]{4.2})), and then its
direct image under $u_{f_{1}}$.)

\subsection[Some properties of $R^{+}f_{!}$]{Some properties of
  \mathversion{bold}$R^{+}f_{!}$}\etiquette[7.2]{XVII.subsec7.2}
\pageoriginale

\begin{enonce*}{Proposition 7.2.1}\etiquette[7.2.1]{XVII.prop7.2.1}
Let $f:X\rightarrow Y$ be a morphism in $\Sch_{!R}$ of relative
dimension $\leq d$. Then the functor $R^{+}f_{!}$ has cohomological
dimension $\leq 2d$.
\end{enonce*}

It is enough to show that, for every sheaf $F$ of
$\mathcal{A}_{X}$-modules on $X$, one has $R^{i}f_{!}(F)=0$ for
$i>2d$. (It then follows by a standard argument that, for every complex
$K^{\bigdot}$ in $D^{+}(X,\mathcal{A}_{X})$ such that
$H^{i}(K^{\bigdot})=0$ for $i>k$, one has $R^{i}f_{!}(K^{\bigdot})=0$
for $i>k+2d$.)

Let $\mathcal{U}=(U_{i})_{i\in I}$ be a finite covering of $X$ by
affine open subsets. One defines (\hyperref[XVII.prop6.2.9]{6.2.9}) a left resolution
$K^{\bigdot}_{\mathcal{U}}\rightarrow F$ of $F$ by putting
\[
K^{n}_{\mathcal{U}}=\bigoplus_{|P|=\sisi{}{-}n+1} j_{P!}j^{*}_{P}(F)
\]
where, for every subset $P$ of $I$, $j_{P}$ is the inclusion of
$\cup_{i\in P} U_{i}$ into $X$.

From \eqhyperref[XVII.eq7.1.9.1]{7.1.9.1} one obtains a spectral sequence
\[
E^{pq}_{1}=\bigoplus_{\substack{|P|=\sisi{}{-}p+1\\ P\sisi{\in}{\subset} I\sisi{}{\text{ nonempty}}}} R^{p}(fj_{P})_{!}
j^{\ast}_{P}(F)\Longrightarrow R^{p+q}f_{!}(F)
\]
and the fact that $fj_{P}$ is compactifiable, a case covered by Exposé
XVII, shows that $E^{pq}_{1}=0$ for $p>0$ and $q>2d$; it follows that
$R^{n}f_{!}(F)=0$ for $n>2d$, as required.

\begin{enonce*}{Proposition 7.2.2}\etiquette[7.2.2]{XVII.prop7.2.2}
Let $f:X\rightarrow Y$ be a morphism in $\Sch_{!R}$ and let
$K^{\bigdot}$ be an object of $D^{b}(X,\mathcal{A}_{X})$ of finite
\tTor-dimension. Then $R^{+}f_{!}(K^{\bigdot})$ has finite
\tTor-dimension.
\end{enonce*}

As above, let $\mathcal{U}=(U_{i})_{i\in I}$ be a finite covering of
$X$ by affine open subsets. The functoriality of
$K^{n}_{\mathcal{U}}$ makes it possible to construct a bounded left
resolution of $K^{\bigdot}$ by complexes of finite \tTor-dimension,
since the latter property can be checked pointwise (\hyperref[XVII.def4.1.9]{4.1.9}). We
denote the double complex so obtained by
$K^{\bigdot\bigdot}_{\mathcal{U}}$, the \emph{second} index being the
one corresponding to $K^{\bigdot}$. With the conventions of
(\sisi{VI b.(\refnoexpose{Vbis.subsubsec2.3.2})}{\hyperrefext[Vbis][Vbis.subsubsec2.3.2]{2.3.2}}), for every integer $n$ there is an exact sequence, where the operation
$\sigma_{\leq n}$ is taken in
$C^{+}(C^{+}(\Mod(X,\mathcal{A}_{X}))$:
\[
\sisi{0\longrightarrow (\sigma_{\leq
n}(K^{\bigdot\bigdot}_{\mathcal{U}}))_{s}\longrightarrow (\sigma_{\leq
n+1}(K^{\bigdot\bigdot}_{\mathcal{U}}))_{s}\longrightarrow (-1)^{n+1}
K^{n+1,\bigdot}_{\mathcal{U}}\longrightarrow 0}{0\to (-1)^{n+1}K^{n+1,\bigdot}_{\mathcal{U}}\to(\sigma_{\leq
n+1}(K^{\bigdot\bigdot}_{\mathcal{U}}))_{s}\to(\sigma_{\leq
n}(K^{\bigdot\bigdot}_{\mathcal{U}}))_{s}\to 0}
\]
hence\pageoriginale a distinguished triangle in $D^{+}(X,\mathcal{A}_{X})$:
\[
\xymatrix@C=1em{
  & K^{n+1,\bigdot}_{\Uu} \sisi{\ar[ld]}{\ar[rd]}&    \\
  (\sigma_{\leq n}(K^{\bigdot\bigdot}_{\Uu}))_s \sisi{\ar[rr]}{\ar[ur]}&  &(\sigma_{\leq
  n+1}(K^{\bigdot\bigdot}_{\Uu}))_s \sisi{\ar[lu]}{\ar[ll]}\\
}
\]

Using the fact that $R^{+}f_{!}(K^{n+1,\bigdot}_{\mathcal{U}})$ has
finite \tTor-dimension for every $n$ (\hyperref[XVII.thm5.2.10]{5.2.10}), it follows that
$R^{+}f_{!}((\sigma_{\leq
0}(K^{\bigdot\bigdot}_{\mathcal{U}}))_{s})=R^{+}f_{!}(K^{\bigdot})$
has finite \tTor-dimension.

\begin{enonce*}[remark]{Remark 7.2.3}\etiquette[7.2.3]{XVII.rem7.2.3}
The preceding argument does not show that
$R^{+}f_{!}(K^{\bigdot})$ has \tTor-dimension $\leq k$ whenever
$K^{\bigdot}$ has \tTor-dimension $\leq k$ (\cf \sisi{(3.7)}{\hyperref[XVII.subsubsec7.3.7]{7.3.7}}).
\end{enonce*}

\begin{enonce*}{Proposition 7.2.4}\etiquette[7.2.4]{XVII.prop7.2.4}
Let there be given a Cartesian diagram in $\Sch_{!R}$:
\[
\xymatrix{
X'   \ar[r]^{g'}\ar[d]_{f'}& X \ar[d]^{f}\\
S'  \ar[r]_{g} & S
}
\]
Then, for every complex $K^{\bigdot}$ in $D^{+}(X,\mathcal{A}_{X})$,
the canonical morphism
\[
g^{\ast}\circ
R^{+}f_{!}(K^{\bigdot})\xrightarrow{\mathrm{ch}_{K^{\bigdot}}}R^{+}f_{!}\sisi{}{'}\circ
{g'}^{\ast}(K^{\bigdot})
\]
is an isomorphism.
\end{enonce*}

Let $(\pi_{1}, j_{1}, f_{1})$ be a pseudo-compactification of $f$; by
base change it gives a pseudo-compactification
$(\pi'_{1}, j'_{1}, f'_{1})$ of $f'$ and a commutative diagram:
\[
\xymatrix@R=2ex{
                 && T'_1 \ar'[ldddd]^(.8){f'_1}|(.55)\hole [lldddddddd]|(.25)\hole \ar[rrrr] &&&& T_1 \ar[lldddddddd]^{f_1}\\
                 &\\
X'_1 \ar[rruu]^{j'_1}\ar[ddd]_{\pi'_1} \ar[rrrr] &&&& X_1 \ar[rruu]_{j_1} \ar[ddd]^{\pi_1}&\\
                     & & &\\
                     & *{} & &\\
X'   \ar[rrrr]_{g'} \ar[ddd]_{f'} &&&& X  \ar[ddd]_{f}  &\\
                     & & &\\
                     & & &\\
S'   \ar[rrrr]_{g} &&&&  S         &
}
\]
which\pageoriginale makes it possible to define the arrow
$\mathrm{ch}_{K}$.

By \eqhyperref[XVII.eq7.1.9.2]{7.1.9.2}, there is then a morphism of
biregular spectral sequences:
\[
\xymatrix{
g^{\ast}\circ R^{q} (f \circ \pi^{p}_1)_! (\pi^{p\sisi{+}{\ast}}_1
(K^{\bigdot})) \ar[d]^{u^{pq}_1} \ar[r]&  g^{\ast} \circ R^{p+q}f_{!}(K^{\bigdot}) \ar[d]^{\mathrm{ch}_K}\\
R^{q}(f'\circ \pi'^{p}_1)_! (\pi'^{p\ast}_1 (g'^{\ast}(K^{\bigdot})))  \ar[r] &
  R^{+}f'_! \circ g'^{\ast}(K\sisi{}{^{\bigdot}})
}
\]
and $\mathrm{ch}_{K^{\bigdot}}$ is an isomorphism because the same is
true of the $u^{pq}_{1}$ by (\hyperref[XVII.thm5.2.6]{5.2.6}).

\subsection{Definition of $Rf_{!}$}\etiquette[7.3]{XVII.subsec7.3}

We now propose to extend the functor $R^{+}f_{!}$ ``in a reasonable
way'' to a functor
$Rf_{!}:D(X,\mathcal{A}_{X})\rightarrow D(Y,\mathcal{A}_{Y})$, for
every morphism $f$ in $\Sch_{!R}$.

\subsubsection{}\etiquette[7.3.1]{XVII.subsubsec7.3.1}
Let $(p_{1}, j_{1}, f_{1})$ be a pseudo-compactification of $f$, and
let $F$ be an $\mathcal{A}_{X}$-module. Taking the canonical flabby
resolution of $j_{1!}\circ L\bar{u}_{p_{1}}^{\ast}(F)$ (\cf
\hyperref[XVII.subsec4.2]{4.2}), then its direct image under $f_{1}$,
and then the simple complex associated with the double complex so
obtained gives a complex of $\mathcal{A}_{Y}$-modules:
$$\etiquette[7.3.1.1]{XVII.eq7.3.1.1}
0\longrightarrow
\Xx^{0}(F)\xrightarrow{S^{1}(F)}\Xx^{1}(F)\ldots
\Xx^{n}(F)\xrightarrow{S^{n}(F)}
\Xx^{n+1}(F)\ldots.\leqno(7.3.1.1)
$$
functorial\pageoriginale in $F$ and isomorphic in
$D^{+}(Y,\mathcal{A}_{Y})$ to $R^{+}f_{!}(F)$.

By \hyperref[XVII.prop7.2.1]{7.2.1}, $H^{i}(\Xx^{\bigdot}(F))=0$ for $i>2d$, where $d$
is an integer such that the relative dimension of $f$ is $\leq d$;
hence the complex
$$\etiquette[7.3.1.2]{XVII.eq7.3.1.2}
0\to \Xx^{0}(F)\to
\Xx^{1}(F)\to \ldots \sisi{}{\to}
\Xx^{2d-1}(F)\to \Ker S^{2d}(F)\to
0\leqno(7.3.1.2)
$$
is canonically isomorphic in $D^{+}(Y,\mathcal{A}_{Y})$ to the
complex \eqhyperref[XVII.eq7.3.1.1]{7.3.1.1}. Moreover:

\begin{enonce*}{Lemma 7.3.2}\etiquette[7.3.2]{XVII.lem7.3.2}
The functors $\Xx^{0}$,
$\Xx^{1},\ldots,\Xx^{2d-1}$, and $\Ker S^{2d}$ are exact
in $F$.
\end{enonce*}

% END INCLUDED SGA4_English_current_buildable_source_20260906/source/17/drafts/SGA4_XVII_LINES_9001_9400_English.texfrag

% BEGIN INCLUDED SGA4_English_current_buildable_source_20260906/source/17/drafts/SGA4_XVII_LINES_9401_9548_English.texfrag

The exactness of the functors $\Xx^{n}$ follows from the exactness
of the canonical flabby resolution and from the fact that the direct
image under a morphism of a short exact sequence of flabby sheaves is
again a short exact sequence.

To prove the exactness of $\Ker S^{2d}$, for a short exact sequence
$0\rightarrow F'\rightarrow F\rightarrow F''\rightarrow 0$ of
sheaves on $X$, consider the diagram
\[
\xymatrix{
\Ker S^{2d}(F')  \ar[r] \ar[d]& \Xx^{2d+1}(F') \ar[r] \ar[d]&
\Xx^{2d+2}(F') \ar[d] \ldots\\
\ker S^{2d}(F)   \ar[r] \ar[d]& \Xx^{2d+1}{(F)} \ar[r] \ar[d]&
\Xx^{2d+2}(F)\ar[d] \ldots \\
\ker S^{2d}(F'')   \ar[r] & \Xx^{2d+1}{(F'')} \ar[r] &
\Xx^{2d+2}{(F'')} \ldots \\
}
\]
and conclude by using the homology exact sequence associated with the
exact sequence
\[
0\longrightarrow \sigma_{\geq 2d+1}
(\Xx^{\bigdot}(F'))\longrightarrow \sigma_{\geq 2d+1}
(\Xx^{\bigdot}(F))\longrightarrow \sigma_{\geq
2d+1}(\Xx^{\bigdot}(F''))\to0.
\]

We now have the following general lemma:

\begin{enonce*}{Lemma 7.3.3}\etiquette[7.3.3]{XVII.lem7.3.3}
Let $A$ and $B$ be two abelian categories, and let
$\mathrm{Je}^{\bigdot}$ be an object of $C^{b}(\Ffex(A,B))$, where
$\Ffex(A,B)$ denotes the category of exact functors from $A$ to $B$.
Associating with every complex $K^{\bigdot}$ in $C(A)$ the simple
complex associated with the double complex
$\mathrm{Je}^{\bigdot}(K^{\bigdot})$ gives a triangulated functor
$K(A)\rightarrow K(B)$ that preserves quasi-isomorphisms.
\end{enonce*}

The functor in question may be written\pageoriginale as a composite
\[
C(A)\longrightarrow C^{b}(C(B))\xrightarrow{(\quad)_{s}}C(B),
\]
where $(\quad)_{s}$ is the associated-simple-complex functor, and one
uses (\sisi{VI b (\refnoexpose{Vbis.lem2.3.2.2})}{\hyperrefext[Vbis][Vbis.lem2.3.2.2]{2.3.2.2}}. Remark); the details are left to the
reader.

\begin{enonce*}[remark]{Exercise 7.3.4}\etiquette[7.3.4]{XVII.exer7.3.4}
Show that, for every complex $K^{\bigdot}$ in
$D^{+}(X,\mathcal{A}_{X})$, the canonical arrow
\[
\tau_{\leq 2d} \Xx^{\bigdot}(K^{\bigdot})\longrightarrow
\Xx^{\bigdot}(K^{\bigdot})
\]
induces a quasi-isomorphism on the associated simple complexes.
\end{enonce*}

For this, observe first that the canonical arrow
\[
\varinjlim_{k\geq 2d} (\tau_{\leq k}
\Xx^{\bigdot}(K^{\bigdot}))_{s}\longrightarrow
(\Xx^{\bigdot}(K^{\bigdot}))_{s}
\]
is an isomorphism, and then that the cohomology sheaves of a complex
commute with filtered direct limits.

\setcounter{subsubsection}{4}
\subsubsection{}\etiquette[7.3.5]{XVII.subsubsec7.3.5}
Let $R^{(1)}f_{!}:D(X,\mathcal{A}_{X})\rightarrow
D(Y,\mathcal{A}_{Y})$ be the functor defined by
\eqhyperref[XVII.eq7.3.1.2]{7.3.1.2} using
\sisi{\hyperref[XVII.rem7.2.3]{7.2.3}}{\hyperref[XVII.lem7.3.3]{7.3.3}}.
In view of
\sisi{\hyperref[XVII.prop7.2.4]{7.2.4}}{\hyperref[XVII.exer7.3.4]{7.3.4}},
it can be computed directly from
\eqhyperref[XVII.eq7.3.1.1]{7.3.1.1}, and one can check that it ``does
not depend'' on the chosen pseudo-compactification. One would then
have to verify statements \hyperref[XVII.constru7.1.3]{7.1.3},
\hyperref[XVII.constru7.1.4]{7.1.4}, and
\hyperref[XVII.prop7.1.6]{7.1.6} after deleting the superscripts $+$,
which the editor did not have the courage to do.

\begin{enonce*}{Definition 7.3.6}\etiquette[7.3.6]{XVII.def7.3.6}
The functor $Rf_{!}:D(X,\mathcal{A}_{X})\rightarrow
D(Y,\mathcal{A}_{Y})$ constructed in
\hyperref[XVII.subsubsec7.3.5]{7.3.5} is again called the
\emph{direct-image functor with proper supports.}
\end{enonce*}

\setcounter{subsubsection}{6}
\subsubsection{}\etiquette[7.3.7]{XVII.subsubsec7.3.7}
One can then show that the functor $Rf_{!}$ commutes with base change
and that statements (\hyperref[XVII.prop5.2.9]{5.2.9}) and
(\hyperref[XVII.thm5.2.10]{5.2.10}) remain valid for every morphism in
$\Sch_{!R}$, which justifies the constructions just made.

\begin{thebibliography}{\sisi{TF}{99}}
\sisi{\bibitem[1]{sga4.exp.XVII.bib.DP}}{\bibitem{sga4.exp.XVII.bib.DP}} A. Dold\pageoriginale and D. Puppe. Homologie nicht\sisi{ }{-}additiver
Funktoren. Anwendungen. Ann. Inst. Fourier \sisi{XI 1961}{\textbf{11} (1961), 201--312}.

\sisi{\bibitem[1]{sga4.exp.XVII.bib.GZ}}{\bibitem{sga4.exp.XVII.bib.GZ}} Gabriel-Zisman. Calculus of Fractions and Homotopy
Theory, Ergebnisse, Bd 35, Springer 1967.

\sisi{\bibitem[1]{sga4.exp.XVII.bib.Gi}}{\bibitem{sga4.exp.XVII.bib.Gi}} J. Giraud. Méthode de la descente, Mémoire nᵒ 2 de la Société Mathématique de France (1964).

\sisi{\bibitem[TF]{sga4.exp.XVII.bib.TF}}{\bibitem{sga4.exp.XVII.bib.TF}} R. Godement. Théorie des faisceaux, Hermann,
1958. Act. Scient. Ind. nᵒ 1252, (Paris)\sisi{ (cited [TF])}{}.

\sisi{\bibitem[1]{sga4.exp.XVII.bib.Gr}}{\bibitem{sga4.exp.XVII.bib.Gr}} A. Grothendieck. TDTE IV: les schémas de
Hilbert. Séminaire Bourbaki 221. May 1961. Reprinted in FGA.

\sisi{\bibitem[RD]{sga4.exp.XVII.bib.RD}}{\bibitem{sga4.exp.XVII.bib.RD}} R. Hartshorne. Residues and Duality. Lecture Notes
in Mathematics nᵒ 20. Springer 1966.

\sisi{\bibitem[1]{sga4.exp.XVII.bib.La}}{\bibitem{sga4.exp.XVII.bib.La}} D. Lazard. Thesis --- Paris 1968.

\sisi{}{\bibitem{sga4.exp.XVII.bib.Mu} D. Mumford, \emph{Geometric Invariant Theory}, Ergebnisse der Mathematik und ihrer Grenzgebiete 34, Springer-Verlag, 1965.}

\sisi{}{\bibitem{sga4.exp.XVII.bib.Na} M. Nagata, ``A generalization of the imbedding problem of an abstract variety in a complete variety'', \emph{J. Math. Kyoto Univ.} \textbf{3} (1963), 89--102}

\sisi{\bibitem[1]{sga4.exp.XVII.bib.Ro}}{\bibitem{sga4.exp.XVII.bib.Ro}} N. Roby. Lois polynômes et lois formelles en théorie
des modules. Ann. E.N.S. \textbf{80} \sisi{1963, p.}{(1963), }231--\sisi{2}{3}48.

\sisi{\bibitem[CD]{sga4.exp.XVII.bib.CD}}{\bibitem{sga4.exp.XVII.bib.CD}} J.L. Verdier. Catégories dérivées (état 0),
mimeographed by the I.H.E.S.

\sisi{\bibitem[2]{sga4.exp.XVII.bib.V2}}{\bibitem{sga4.exp.XVII.bib.V2}} J.L. Verdier. Catégories dérivées (to appear in
Grundlehren der Math. Wiss, Springer).

\sisi{\bibitem[3]{sga4.exp.XVII.bib.V3}}{\bibitem{sga4.exp.XVII.bib.V3}} J.L. Verdier. Dualité dans la cohomologie des
espaces localement compacts. Séminaire Bourbaki 300. Nov. 1965.
\end{thebibliography}


%0.2 no ref. 1957, 1329, 5575, 5594, 5821
%2.1.2.2 no ref. 1386
%2.1.2.3 no ref  1412
%2.1.12.1 no ref 1641
%3.2.4.3 no ref  2680
%5.1.8.2 no ref 3925
%XII 4.1 no ref 4690, 6830
%6.3.4.3 no ref 7561, 7876, 8060,8409
%6.3.11 no ref 7736, 7833
%6.2.17 no ref 8042
%6.3.17.1 no ref 8078
%6.3.2.3 no ref 8271
%4.3.0 no ref 8569
%3.7 no ref 9229


\ifx\danslelivre\undefined
\let\SGAXVIIenddocument\enddocument
\else
\let\SGAXVIIenddocument\relax
\fi
\SGAXVIIenddocument

% END INCLUDED SGA4_English_current_buildable_source_20260906/source/17/drafts/SGA4_XVII_LINES_9401_9548_English.texfrag


% END INCLUDED SGA4_English_current_buildable_source_20260906/source/17/17_English.tex


\makeatletter
\def\input@path{{../18/}}
\makeatother
% BEGIN INCLUDED SGA4_English_current_buildable_source_20260906/source/18/18_English.tex
\ifx\danslelivre\undefined
% Release engine: XeLaTeX (the LuaTeX compatibility layer breaks Xy-pic output).
\documentclass[10pt]{../fichiers-configuration-latex/smfart}
% BEGIN INCLUDED SGA4_English_current_buildable_source_20260906/source/fichiers-configuration-latex/paquets.tex
\usepackage{polyglossia}
%\usepackage[francais]{babel}
\usepackage{fontspec}
\usepackage[math-style=ISO,bold-style=upright,sans-style=upright]{unicode-math}
%\usepackage[mathletters]{ucs}
%\usepackage[utf8]{inputenc}
%\usepackage[T1]{fontenc}
\usepackage[unicode,hypertexnames=false,colorlinks=true,linkcolor=blue,citecolor=Sepia,urlcolor=magenta]{hyperref}
\setmainlanguage{french}
\setmainfont[Ligatures=TeX]{Linux Libertine O}
% variante possible : TeXGyrePagella ?
\setmathfont{XITS Math}

% maths
\usepackage{amsmath}
%\usepackage{amssymb}
%\usepackage{mathrsfs}
%\usepackage{stmaryrd}
%\usepackage{euscript}

% couleurs
\usepackage[usenames,dvipsnames]{xcolor}

% graphiques
\usepackage{graphicx}
\usepackage[all]{xy}
%\usepackage{tikz}
%\usetikzlibrary{matrix,arrows,calc}
%\tikzset{column sep=3em,row sep=2em,text height=1.5ex, text depth=0.25ex,>=angle 90}
%\tikzset{description/.style={fill=white,inner sep=1pt}}

% hyperliens, 'etiquettes
\usepackage{xr-hyper}
\usepackage{../fichiers-configuration-latex/smfthm}
%\ifx\cacheretiquettes\undefined
%\usepackage[inline]{showlabels} \else \fi % pour rendre visibles les 'etiquettes
%\usepackage{../fichiers-configuration-latex/makeidx} % pour faire un index
%\usepackage[pagebackref,unicode]{hyperref}  % pour voir o`u sont cit'ees les r'ef'erences
% hyperref is loaded once above with the cumulative-reader destination policy.
%\hypersetup{breaklinks=true} %permet le retour `a la ligne dans les liens trop longs. Incompatible avec pdflatex
%\renewcommand*{\backref}[1]{{\small [cit\'e page(s)~#1.]}}

\usepackage{xspace}
\usepackage{pifont} % pour expose XVII

% END INCLUDED SGA4_English_current_buildable_source_20260906/source/fichiers-configuration-latex/paquets.tex

% BEGIN INCLUDED SGA4_English_current_buildable_source_20260906/source/fichiers-configuration-latex/numerotation.tex
\renewcommand{\thefootnote}{\roman{footnote}} % pour mettre les notes en bas de page sous la forme (ii) et non (2), qui ressemble trop (?) a un exposant.

\newcounter{pageoriginal}
\newcommand{\pageoriginale}{\refstepcounter{pageoriginal}\marginpar{\footnotesize\textbf{\thepageoriginal}}\xspace}
\let\pageoriginaled\pageoriginale

\usepackage{ifthen}
\newboolean{original}
\def\sisi{\ifthenelse{\boolean{original}}}
\setboolean{original}{false}

\newcounter{cnote}
\def\nde#1{\ifthenelse{\boolean{original}}{}{\stepcounter{cnote}\def\thefootnote{\arabic{cnote}}\protect\footnote[\value{cnote}]{N.D.E.: #1}}}
\def\ndemark{\ifthenelse{\boolean{original}}{}{\stepcounter{cnote}\def\thefootnote{\arabic{cnote}}\protect\footnotemark[\value{cnote}]}}
\def\ndetext#1{\ifthenelse{\boolean{original}}{}{\def\thefootnote{\arabic{cnote}}\protect\footnotetext[\value{cnote}]{N.D.E.: #1}}}

% END INCLUDED SGA4_English_current_buildable_source_20260906/source/fichiers-configuration-latex/numerotation.tex

% BEGIN INCLUDED SGA4_English_current_buildable_source_20260906/source/fichiers-configuration-latex/macros.tex
\setmainlanguage{french}
\setotherlanguage{english}


%% Fixed, engine-independent release identifier.  The inherited dynamic Git
%% probe used Lua io.popen, exposed build-state-dependent output, and printed
%% a visible question mark outside a Git checkout.
\def\showgitstatus{English release r2, 2026-07-23}


% macros pour LaTeX
\renewcommand{\mathsf}{}

% Police différente pour U+2113 SCRIPT SMALL L
\setmathfont[range="2113]{latinmodern-math.otf}

% Give every inherited semantic label its own stable PDF destination instead
% of inheriting whichever automatic equation/page anchor happened to be live.
% The legacy label remains as a compatibility alias used by existing edges.
\makeatletter
\def\etiquette[#1]#2{%
  \Hy@raisedlink{\hyper@anchorstart{sga4.semantic.#2}\hyper@anchorend}%
  \def\@currentHref{sga4.semantic.#2}%
  \def\@currentlabel{#1}%
  \label{#2}%
}
\makeatother
\def\eqhyperref[#1]#2{(\hyperref[#1]{#2})}
\def\pagehyperref[#1]#2{(page \hyperref[#1]{#2})}
\def\hyperrefext[#1][#2]#3{\hyperref[#2]{#1 #3}}
\def\eqhyperrefext[#1][#2]#3{\hyperref[#2]{#1 #3}}     %cf. exposes 17
\def\neweqhyperrefext[#1][#2]#3{\hyperref[#2]{#1 #3}}  %cf. exposes 17

\def\shat#1{\protect\widehat{#1}}
\def\stilde#1{\protect\widetilde{#1}}
%%%%%%%%%%%%%%%%%%


\def\etcn{etc.\xspace}
\def\etco{etc...\xspace}
\def\etc{\sisi{\etco}{\etcn}}
\def\ldotsn{\sisi{...}{\etcn}}

\def\cf{{cf.}\kern.3em}
\def\Cf{{Cf.}\kern.3em}
\def\ie{{i.e.}\kern.3em}
\def\iev{{i.e.},\xspace}
\def\loccito{loc.\kern2pt cit.\xspace}
\def\loccitn{\textit{loc.\kern2pt cit.}\xspace}
\def\loccit{\sisi{\loccito}{\loccitn}}
\def\EGA{\textup{EGA}\kern.2em}
\def\SGA{\textup{SGA}\kern.3em}

\def\quoi{}

\def\quav{\quad,\quad}
\def\quaet{\quad \text{et}\quad}
\def\fibreproduct#1#2#3{#1 ×_{#3} #2}
\let\fprod\fibreproduct
\def\ccheck#1#2{\mbox{${\check{\check{\rm #1}}}^{\lower2pt\hbox{$\scriptstyle #2$}}$}}
\def\sheaf#1{\protect\mathcal{#1}}
\def\vieuxsheaf#1{{#1}} % cf. problème avec sheaf{a\Phi_i} par exemple...
\DeclareMathOperator{\Hom}{Hom}
\DeclareMathOperator{\can}{can}
\DeclareMathOperator{\Ind}{Ind}
\DeclareMathOperator{\ind}{ind}
%\DeclareMathOperator{\Hhomtop}{\mathcal{H}\!\mathit{omtop}}
\DeclareMathOperator{\Hhomtop}{\mathbf{Homtop}}
\DeclareMathOperator{\Point}{Point}
\DeclareMathOperator{\Pro}{Pro}
\DeclareMathOperator{\Top}{Top}
\DeclareMathOperator{\Mod}{Mod}
\DeclareMathOperator{\Cart}{Cart}
%\DeclareMathOperator{\Mmorsite}{\mathcal{M}\!\mathit{orsite}}
\DeclareMathOperator{\Mmorsite}{\mathbf{Morsite}}
\DeclareMathOperator{\loc}{loc}
\DeclareMathOperator{\Loc}{Loc}
\DeclareMathOperator{\Pic}{Pic}
\DeclareMathOperator{\Et}{Et}
\DeclareMathOperator{\card}{card}
\DeclareMathOperator{\Ens}{Ens}
\DeclareMathOperator{\h}{\underline{H}}
\DeclareMathOperator{\Fonct}{Fonct}
\DeclareMathOperator{\Ker}{Ker}
\DeclareMathOperator{\but}{but}
\DeclareMathOperator{\Ob}{Ob}
\DeclareMathOperator{\ob}{ob}
\DeclareMathOperator{\Lib}{Lib}
\DeclareMathOperator{\esj}{esj}
\DeclareMathOperator{\id}{id}
\DeclareMathOperator{\Id}{Id}
\DeclareMathOperator{\ad}{ad}
\DeclareMathOperator{\Iim}{Im}
\DeclareMathOperator{\pr}{pr}
\DeclareMathOperator{\Cat}{Cat}
\DeclareMathOperator{\car}{car}
\DeclareMathOperator{\SysInd}{SysInd}
\DeclareMathOperator{\Sup}{Sup}
\DeclareMathOperator{\Coker}{Coker}
\DeclareMathOperator{\Filt}{Filt}
\DeclareMathOperator{\acc}{acc}
\DeclareMathOperator{\sur}{sur}
\DeclareMathOperator{\cat}{cat}
\DeclareMathOperator{\Ord}{Ord}
\DeclareMathOperator{\Gr}{Gr}
\DeclareMathOperator{\Cov}{Cov}
\DeclareMathOperator{\Caf}{Caf}
\DeclareMathOperator{\coker}{coker}
\DeclareMathOperator{\coim}{coim}
\DeclareMathOperator{\im}{im}
\DeclareMathOperator{\Ouv}{Ouv}
\DeclareMathOperator{\diag}{diag}
\DeclareMathOperator{\Esp}{Esp}
\DeclareMathOperator{\TOP}{TOP}
\DeclareMathOperator{\Unif}{Unif}
\DeclareMathOperator{\Iind}{\mathbf{Ind}}
\DeclareMathOperator{\Hhomcart}{\mathbf{Homcart}}
\DeclareMathOperator{\Hhomtopess}{\mathbf{Homtopess}}
\DeclareMathOperator{\Ppoints}{\mathbf{Points}}
\DeclareMathOperator{\Ppoint}{\mathbf{Point}}
\DeclareMathOperator{\Ttors}{\mathbf{Tors}}
\DeclareMathOperator{\Ssex}{\mathbf{Sex}}
\DeclareMathOperator{\Ffib}{\mathbf{Fib}}
\DeclareMathOperator{\Vvois}{\mathbf{Vois}}
\DeclareMathOperator{\Ppointess}{\mathbf{Pointess}}
\DeclareMathOperator{\scrHom}{\mathbf{Hom}}
\DeclareMathOperator{\sheafPic}{\mathbf{Pic}}
\DeclareMathOperator{\sheafAut}{\mathbf{Aut}}
\DeclareMathOperator{\Hhomtoplocan}{\mathbf{Homtoplocan}}
\DeclareMathOperator{\Hhomtopan}{\mathbf{Homtopan}}
\DeclareMathOperator{\Pptgeom}{\mathbf{Ptgéom}}
\DeclareMathOperator{\Morsite}{Morsite}
\DeclareMathOperator{\Fib}{Fib}
\DeclareMathOperator{\Fais}{Fais}
\DeclareMathOperator{\kar}{kar}
\DeclareMathOperator{\Groupes}{Groupes}
\DeclareMathOperator{\Restr}{Restr}
\DeclareMathOperator{\Isom}{Isom}
\DeclareMathOperator{\Prol}{Prol}
\DeclareMathOperator{\Res}{Res}
\DeclareMathOperator{\Tors}{Tors}
\DeclareMathOperator{\esp}{esp}
\DeclareMathOperator{\es}{es}
\DeclareMathOperator{\od}{od}
\DeclareMathOperator{\gr}{gr}
\DeclareMathOperator{\sob}{sob}
\DeclareMathOperator{\dfn}{dfn}
\DeclareMathOperator{\Sch}{Sch}
\DeclareMathOperator{\ouv}{ouv}
\DeclareMathOperator{\cons}{cons}
\DeclareMathOperator{\Nerf}{Nerf}
\DeclareMathOperator{\Points}{Points}
\DeclareMathOperator{\Espsob}{Espsob}
\DeclareMathOperator{\Homtop}{Homtop}
\DeclareMathOperator{\Pt}{Pt}
\DeclareMathOperator{\Quot}{Quot}
\DeclareMathOperator{\glob}{glob}
\DeclareMathOperator{\Homfib}{Homfib}
\DeclareMathOperator{\Qouv}{Qouv}
\DeclareMathOperator{\Quotouv}{Quotouv}
\DeclareMathOperator{\ens}{ens}
\DeclareMathOperator{\ab}{ab}
\DeclareMathOperator{\an}{an}
\DeclareMathOperator{\fs}{fs}
\DeclareMathOperator{\red}{red}
\DeclareMathOperator{\Ab}{Ab}
\DeclareMathOperator{\cosq}{cosq}
\DeclareMathOperator{\Ann}{Ann}
\DeclareMathOperator{\Zar}{Zar}


%%%%%%%%%%%%
%%%%% SGA4-2
\DeclareMathOperator{\Sex}{Sex}
\DeclareMathOperator{\vois}{vois}
\DeclareMathOperator{\Vois}{Vois}
\DeclareMathOperator{\Ext}{Ext}
\DeclareMathOperator{\ext}{ext}
\DeclareMathOperator{\Fer}{Fer}
\DeclareMathOperator{\Add}{Add}
\DeclareMathOperator{\Cocart}{Cocart}
\DeclareMathOperator{\cosk}{cosk}
\DeclareMathOperator{\plat}{plat}

\DeclareMathOperator{\faisc}{faisc}
\DeclareMathOperator{\sk}{sk}
\DeclareMathOperator{\ch}{ch}

\def\HH{\mathbf{H}}
\def\qc{\mathrm{qc}}
\def\et{\mathrm{et}}
\def\coh{\mathrm{coh}}

\def\Hvr#1{%
\begin{array}[b]{l}
{\kern2pt\scriptstyle r}\\[-8pt]
{\kern1pt\scriptstyle \vee}\\[-5pt]
H^{#1}
\end{array}}

\def\Hvv#1{%
\begin{array}[b]{l}
\kern2pt\rotatebox{90}{${\scriptscriptstyle\ll}$}\\[-6pt]
H^{#1}
\end{array}}

\def\Hvi#1{%
\begin{array}[b]{l}
{\kern0.2pt\scriptscriptstyle \infty}\\[-8pt]
{\kern1pt\scriptstyle \vee}\\[-5pt]
H^{#1}
\end{array}}

\def\Hv#1{%
\begin{array}[b]{l}
{\kern1pt\scriptstyle \vee}\\[-5pt]
H^{#1}
\end{array}}

\DeclareMathOperator{\cart}{cart}
\DeclareMathOperator{\Ffaisc}{\mathbf{Faisc}}


\def\Ev#1{%
\begin{array}[b]{l}
{\kern1pt\scriptstyle \vee}\\[-5pt]
E'_{#1}
\end{array}}

\DeclareMathOperator{\Cont}{Cont}
\DeclareMathOperator{\Ccont}{\mathscr{C}\!\mathit{ont}}
\DeclareMathOperator{\Site}{Site}
\newcommand\MM{\mathbf{M}}
\def\limtop#1{{\displaystyle{\varprojlim_{#1}\!\mbox{\rm top }}}}
\def\Limtop{\mathop{\varprojlim\!\mbox{\rm top}}}
\DeclareMathOperator{\spec}{spec}
\def\ss{\scriptstyle{\mathscr{S}}}
\DeclareMathOperator{\Gal}{Gal}
\def\smallm{\scriptstyle{\mathscr{M}}}
\DeclareMathOperator{\End}{End}
\DeclareMathOperator{\ass}{ass}
\DeclareMathOperator{\zar}{zar}
\DeclareMathOperator{\Ogg}{Ogg}
\DeclareMathOperator{\res}{res}
\DeclareMathOperator{\tr}{tr}
\DeclareMathOperator{\Supp}{Supp}
\DeclareMathOperator{\codim}{codim}
\DeclareMathOperator{\rk}{rk}
\DeclareMathOperator{\cur}{cur}
\DeclareMathOperator{\cdqc}{cdqc}
\newcommand\DD{\mathbf{D}}
\DeclareMathOperator{\cl}{cl}
\def\JJ{\mathbf{J}}
\def\FF{\mathbf{F}}
\def\BB{\mathbf{B}}
\DeclareMathOperator{\Alget}{Alget}
\DeclareMathOperator{\Ass}{Ass}
\DeclareMathOperator{\supp}{supp}
\DeclareMathOperator{\rad}{rad}
\DeclareMathOperator{\prof}{prof}
\DeclareMathOperator{\cint}{int}
\DeclareMathOperator{\sic}{sic}
\DeclareMathOperator{\Tor}{Tor}
\DeclareMathOperator{\tors}{tors}
\DeclareMathOperator{\Sym}{Sym}
\DeclareMathOperator{\parf}{parf}
\DeclareMathOperator{\tor}{tor}
\DeclareMathOperator{\defn}{defn}
\DeclareMathOperator{\Div}{Div}
\DeclareMathOperator{\Pol}{Pol}
\DeclareMathOperator{\Fex}{Fex}
\DeclareMathOperator{\Lie}{Lie}
\DeclareMathOperator{\tot}{tot}
\DeclareMathOperator{\Aut}{Aut}
\DeclareMathOperator{\HOM}{HOM}
\DeclareMathOperator{\pch}{pch}
\DeclareMathOperator{\Ch}{Ch}
\DeclareMathOperator{\EXT}{EXT}
\DeclareMathOperator{\PIC}{PIC}
\DeclareMathOperator{\Rf}{Rf}
\DeclareMathOperator{\inj}{inj}
\newcommand\VV{\mathbf{V}}

\def\Xv{%
\begin{array}[b]{l}
{\kern1pt\scriptstyle \vee}\\[-5pt]
X
\end{array}}

%%%%%%%
\newcommand\TT{\mathbf{T}}
\newcommand\XX{\mathbf{X}}
\newcommand\EE{\mathbf{E}}
%%%%%%

\def\qedsymbol{\textup{C.Q.F.D.}}
\let\CQFD\qed
\let\cqfd\qed

%%%%%\bigdot
\font\bigsymb=cmsy10 at 4pt
\def\bigdot{{\kern1.2pt\raise 1.5pt\hbox{\bigsymb\char15}}}
%%%%%

\def\Lright{\mathop{\vcenter{\hbox{\hbox{$\longrightarrow$}\lower8pt \hbox{\kern-16pt\hbox{$\longrightarrow$}}}}}}

\def\Aa{{\mathcal A}}
\def\Bb{{\mathcal B}}
\def\Cc{{\mathcal C}}
\def\Dd{{\mathcal D}}
\def\Ee{{\mathcal E}}
\def\Ff{{\mathcal F}}
\def\Gg{{\mathcal G}}
\def\Hh{{\mathcal H}}
\def\Ii{{\mathcal I}}
\def\Jj{{\mathcal J}}
\def\Kk{{\mathcal K}}
\def\Ll{{\mathcal L}}
\def\Mm{{\mathcal M}}
\def\Nn{{\mathcal N}}
\def\Oo{{\mathcal O}}
\def\Pp{{\mathcal P}}
\def\Qq{{\mathcal Q}}
\def\Rr{{\mathcal R}}
\def\Ss{{\mathcal S}}
\def\Tt{{\mathcal T}}
\def\Uu{{\mathcal U}}
\def\Vv{{\mathcal V}}
\def\Ww{{\mathcal W}}
\def\Xx{{\mathcal X}}
\def\Yy{{\mathcal Y}}
\def\Zz{{\mathcal Z}}

\newcommand{\CC}{\mathbf{C}}
\newcommand{\GG}{\mathbf{G}}
\newcommand{\KK}{\mathbf{K}}
\newcommand{\LL}{\mathbf{L}}
\newcommand{\NN}{\mathbf{N}}
\newcommand{\PP}{\mathbf{P}}
\newcommand\QQ{\mathbf{Q}}
\newcommand\RR{\mathbf{R}}
\newcommand\ZZ{\mathbf{Z}}

\def\to{\mathchoice{\longrightarrow}{\rightarrow}{\rightarrow}{\rightarrow}}
\def\from{\mathchoice{\longleftarrow}{\leftarrow}{\leftarrow}{\leftarrow}}
\def\mto{\mathchoice{\longmapsto}{\mapsto}{\mapsto}{\mapsto}}
\def\mfrom{\mathrel{\from\joinrel\mapstochar\,}}
\def\hto{\mathrel{\lhook\joinrel\to}}
\def\hfrom{\mathrel{\from\joinrel\rhook}}

\def\implies{\mathchoice{\Longrightarrow}{\Rightarrow}{\Rightarrow}{\Rightarrow}}
\def\implied{\mathchoice{\Longleftarrow}{\Leftarrow}{\Leftarrow}{\Leftarrow}}
\def\iff{\mathchoice{\Longleftrightarrow}{\Leftrightarrow}{\Leftrightarrow}{\Leftrightarrow}}

\def\To#1{\mathchoice{\xrightarrow{\labelstyle~~#1~}}{\stackrel{#1}{\longrightarrow}}{}{}}
\def\Mto#1{\mathrel{\mapstochar\To{#1}}}
\def\From#1{\mathchoice{\xleftarrow{\labelstyle~~#1~}}{\stackrel{#1}{\longleftarrow}}{}{}}

\newcommand{\isomto}{\overset{\sim}{\longrightarrow}}
\newcommand{\isomfrom}{\overset{\sim}{\longleftarrow}}

\newcommand{\IFF}{\unskip~$\iff$~\ignorespaces}
\newcommand{\THEN}{\unskip~$\implies$~\ignorespaces}

\let\leq\leqslant
\let\geq\geqslant
\let\le\leqslant
\let\ge\geqslant

\def\cf{{cf.}\kern.3em}
\def\Cf{{Cf.}\kern.3em}
\def\ie{{i.e.}\kern.3em}
\def\iev{{i.e.},\xspace}
%\def\loccito{loc.\kern2pt cit.\xspace}
%\def\loccitn{\textit{loc.\kern2pt cit.}\xspace}
%\def\loccit{\sisi{\loccito}{\loccitn}}
\def\EGA{\textup{EGA}\kern.2em}
\def\SGA{\textup{SGA}\kern.3em}

\let\oldS\S
\def\S{\oldS\kern .15em}
\def\ptbl{.\kern .2em}
\def\resp{resp.\kern.3em}
\def\Exp{Exp.\kern.3em}

\let\bar\overline
\let\hat\widehat
\let\tilde\widetilde
\let\dpl\displaystyle
\let\ldots\dots
\let\emptyset\varnothing
\let\skpt\skippointrait
\let\epsilon\varepsilon

\def\sfrac#1#2{{#1}/{#2}}

\DeclareMathOperator{\cd}{cd}
%\DeclareMathOperator{\Hhom}{\mathcal{H}\!\mathit{om}}
\DeclareMathOperator{\Hhom}{\mathbf{Hom}}
%\DeclareMathOperator{\Eext}{\mathcal{E}\!\mathit{xt}}
\DeclareMathOperator{\Eext}{\mathbf{Ext}}
\DeclareMathOperator{\Spec}{Spec}
\DeclareMathOperator{\Tr}{Tr}

\let\emptyset=\varnothing

%%%%%%%%%%%%%%%
%%% expose XVII
%%%%%%%%%%%%%%%

\makeatletter
\let\old@optexpose\optexpose
\newcommand{\empty@optexpose}[1]{}
\newcommand{\paren@optexpose}[1]{\old@optexpose{#1}(}
\newcommand{\refnoexpose}[1]{\let\optexpose\empty@optexpose\ref{#1}\let\optexpose\old@optexpose}%\refnoexpose{X.cor5.2} engendre "5.2" au lieu de "X 5.2"
\newcommand{\neweqref}[1]{\let\optexpose\paren@optexpose\textup{\ref{#1})}\let\optexpose\old@optexpose}%\neweqref{XII.eq4.1} engendre "XII (4.1)" au lieu de "(XII 4.1)"
\makeatother
\newcommand{\refvb}[2][VI${}_{\text{B}}$ ]{\sisi{#1\refnoexpose{#2}}{\ref{#2}}}

\newcommand{\leftexp}[2]{{\vphantom{#2}}^{#1}{#2}}
\newcommand{\AAA}{\mathbf{A}}
\newcommand{\Sf}{\mathfrak{S}}
\DeclareMathOperator*{\quotprojlim}{\text{``$\varprojlim$''}}
\DeclareMathOperator*{\quotinjlim}{\text{``$\varinjlim$''}}
\DeclareMathOperator{\Fl}{F\sisi{\ell}{l}}
\DeclareMathOperator{\Pl}{P\sisi{\ell}{l}}
%\DeclareMathOperator{\cl}{c\sisi{\ell}{l}}
\DeclareMathOperator{\tdf}{tdf}
\DeclareMathOperator{\rg}{rg}
\newcommand{\iInd}{\sisi{I}{i}nd}
\newcommand{\iind}{ind}
\newcommand{\ppro}{pro}
\newcommand{\tTor}{\sisi{T}{t}or}
\newcommand{\ttor}{tor}
\newcommand{\chap}{Chap.}
\newcommand{\Ffex}{\Ff\mathit{ex}}
\newcommand{\sheafk}{\underline{k}}
\newcommand{\sheafn}{\underline{n}}
\newcommand{\sheafGamma}{\underline{\Gamma}}

% Repeated references to an already typeset footnote must not use
% \footnotemark without a matching \footnotetext: Hyperref would then create
% a live annotation to a destination that cannot exist.  This helper keeps the
% current printed footnote mark and links it to an explicit stable note target.
\makeatletter
\newcommand{\sgarepeatedfootnotemark}[1]{%
  \begingroup
  \edef\@thefnmark{\thefootnote}%
  \hyperlink{#1}{\@makefnmark}%
  \endgroup}
\makeatother

\sisi{
\DeclareMathOperator{\simarrow}{\rule{.15mm}{1.8mm}\kern -3.4pt\rightsquigarrow}
\newcommand{\sheafell}{\underline{\ell}}
\newcommand{\sheafZZ}{\underline{\ZZ}}
}{
\newcommand{\nderef}[1]{la N.D.E.~(\ref{#1}) page~\pageref{#1}}
}

%%%%%%%%%%%%%%%%%
%%%%% expose XVII
%%%%%%%%%%%%%%%%%

% END INCLUDED SGA4_English_current_buildable_source_20260906/source/fichiers-configuration-latex/macros.tex


\externaldocument{../01/01_English.tex}
\externaldocument{../02/02_English.tex}
\externaldocument{../03/03_English.tex}
\externaldocument{../04/04_English.tex}
\externaldocument{../05/05_English.tex}
\externaldocument{../05bis/05bis_English.tex}
\externaldocument{../06/06_English.tex}
\externaldocument{../07/07_English.tex}
\externaldocument{../08/08_English.tex}
\externaldocument{../09/09_English.tex}
%\externaldocument{../10/10_English.tex}
\externaldocument{../11/11_English.tex}
\externaldocument{../12/12_English.tex}
\externaldocument{../13/13_English.tex}
\externaldocument{../14/14_English.tex}
\externaldocument{../15/15_English.tex}
\externaldocument{../16/16_English.tex}
\externaldocument{../17/17_English.tex}

\setboolean{original}{false}

\begin{document}

\begin{center}
XVIII. The global duality formula\\
P. Deligne
\end{center}

\selectlanguage{english}version: \showgitstatus
\tableofcontents
\mainmatter

\else
\chapter{The global duality formula}
%\addtocontents{toc}{P. Deligne}
\begin{center}
P. Deligne
\end{center}
\fi
\phantomsection\MakeLinkTarget*{sga4.exp.XVIII}\label{sga4.exp.XVIII}

% The edition is fixed to the corrected (second) reading. This expandable
% selector is output-equivalent and avoids current Xy-pic conditional failures.
\renewcommand{\sisi}[2]{#2}
\let\Arrowvert\Vert

% BEGIN INCLUDED SGA4_English_current_buildable_source_20260906/source/18/drafts/SGA4_XVIII_LINES_44_450_English.texfrag
\setcounter{pageoriginal}{480}
\newcommand{\acorriger}[1]{\nde{#1}}

\setcounter{section}{-1}
\section{Introduction}\etiquette[0]{XVIII.sec0}\pageoriginale

\subsection{}\etiquette[0.1]{XVIII.subsec0.1}
The present exposé is devoted to Poincaré duality. This duality is
constructed on the model of \textsc{Verdier} [1], who treats the case of
locally compact Hausdorff topological spaces. We show a priori that the
functor $Rf_{!}$ of Exposé XVII admits a right adjoint $Rf^{!}$ and
establish many of its properties, before computing it more explicitly
in the case of smooth morphisms.

The possibility of such a construction seems to rest on the following
properties of the functor $Rf_{!}$:
\begin{enumerate}
\renewcommand{\theenumi}{\alph{enumi}}
\renewcommand{\labelenumi}{\theenumi)}
\item For the existence of an adjoint $Rf^{!}$ and its formal
  properties:
\begin{itemize}
\item[(I)] the functor $Rf_{!}$ commutes with base change;

\item[(II)] the functors $R^{i}f_{!}$ commute with filtered inductive
  limits; they vanish for sufficiently large $i$;

\item[(III)] the functor $Rf_{!}$ can be computed ``at the level of
  complexes'', \ie starting from a functor between categories of
  complexes.
\end{itemize}

\item For the computation of this adjoint when $f$ is smooth:
\begin{itemize}
\item[(IV)] the computation of $Rf_{!}((\ZZ/n\ZZ)_{U})$ for
  ``small'' $U$.
\end{itemize}
\end{enumerate}

\subsection{}\etiquette[0.2]{XVIII.subsec0.2}
In \hyperref[XVIII.sec3]{\S 3} nᵒ 1, which is independent of \S \hyperref[XVIII.sec1]{\S 1}, 2, we
construct the functor $Rf^{!}$; constructions specific to the situation
are used there, rather than only (I), (II), and (III). This makes this
nᵒ rather unpleasant; the author confesses that he did not
always fully understand what was happening.

\S \hyperref[XVIII.sec2]{2} is devoted to point (IV); it is purely formal
from\pageoriginale \hyperref[XVIII.sec1]{\S 1}, where the cohomology of curves is studied.

In \S \hyperref[XVIII.sec1]{1}, we prove substantially more than is needed for
\S \hyperref[XVIII.sec2]{2}, also studying the case of continuous
coefficients. This section was profoundly influenced by \textsc{Serre}
\cite{sga4.exp.XVIII.bib.6}. The points useful in what follows are:
\begin{enumerate}
\renewcommand{\theenumi}{\alph{enumi}}
\renewcommand{\labelenumi}{\theenumi)}
\item nᵒ \hyperref[XVIII.subsec1.1]{1.1}, devoted to the trace morphism in the
  case of curves, and generalized in \hyperref[XVIII.thm2.9]{2.9};

\item the vanishing theorem \hyperref[XVIII.lemfond1.6.9]{1.6.9}. A reader willing to accept a
  transcendental argument may read nᵒ \hyperref[XVIII.subsec1.6]{1.6}, beginning
  with \hyperref[XVIII.lem1.6.6]{1.6.6} (2\textsuperscript{nd} proof), independently of nᵒ \hyperref[XVIII.subsec1.2]{1.2}
  to \hyperref[XVIII.subsec1.5]{1.5}\nde{On a first reading, one cannot too strongly recommend
  following this route: after $\S$\hyperref[XVIII.subsec1.1]{1.1}, proceed
  directly to $\S$\hyperref[XVIII.sec2]{2} and $\S$\hyperref[XVIII.sec3]{3}.}.
  The essential ingredient in \hyperref[XVIII.lemfond1.6.9]{1.6.9} is either
  form of Poincaré duality for smooth curves over an algebraically
  closed field, together with the acyclicity lemma \hyperrefext[XV][XV.lem2.6]{2.6} (a preliminary
  to the smooth base-change theorem \hyperrefext[XVI][XVI.thm1.1]{1.1}). Its conclusion is
  generalized in \hyperref[XVIII.thm2.14]{2.14}.
\end{enumerate}

Number \hyperref[XVIII.subsec3.2]{3.2} collects the consequences of \S
\hyperref[XVIII.sec2]{2}.

\subsection{}\etiquette[0.3]{XVIII.subsec0.3}
\textsc{Verdier}'s method in Poincaré duality makes it possible to define the
functor $Rf^{!}$ for compactifiable $f$. The price of this generality is
that the compatibilities to be verified are unfamiliar and sometimes,
even for the most trivial ones, such as \hyperref[XVIII.lem3.2.3]{3.2.3}, require
outlandish proofs. The author admits that he has not proved as many
of them as he should have.

An interested reader will have no difficulty in extending, using the
appendix to Exposé XVII, the definition of $Rf^{!}$ to separated
morphisms of finite type whose target is a coherent scheme (\ie a
quasi-compact quasi-separated scheme).

\subsection{}\etiquette[0.4]{XVIII.subsec0.4}
The main result \hyperref[XVIII.thm3.2.5]{3.2.5} is stated in concrete terms that
would make it possible to return formally to the viewpoint taken in the
oral seminar, where the functor $Rf^{!}$ was defined only for\pageoriginale
smoothable morphisms, by means of a factorization through a smooth
morphism and an immersion. An outline of the proof from the oral
seminar can be found in \textsc{Verdier} \cite{sga4.exp.XVIII.bib.8}.

\section{Cohomology of curves}\etiquette[1]{XVIII.sec1}\pageoriginale

\subsection{The trace morphism}\etiquette[1.1]{XVIII.subsec1.1}

\subsubsection{}\etiquette[1.1.1]{XVIII.subsubsec1.1.1}
Let $n$ be an integer $\geq 1$ and let $X$ be a scheme on which $n$ is
invertible. The sheaf $𝛍_n$ (IX \hyperref[IX.eq3.1]{3.1}) on $X$ is then a sheaf of modules
locally free of rank 1 over the constant sheaf of rings $\ZZ/n$.
For $i\in \ZZ$, we denote by $\ZZ/n(i)$ its
$i^{\text{\rm th}}$ tensor power:
\begin{equation*}
\ZZ/n(i)=\Hhom(\ZZ/n,\GG_{m})^{\otimes
    i}{\quoi}.\tag{1.1.1.1}
\end{equation*}\etiquette[1.1.1.1]{XVIII.eq1.1.1.1}

If $n=dn'$, exponentiation by $d$ gives an isomorphism
$$
𝛍_n\otimes_{\ZZ/n}\ZZ/n'\xrightarrow{\sim}𝛍_{n'},
$$
and hence isomorphisms, called canonical:
\begin{equation*}
\ZZ/n(i)\otimes_{\ZZ/n}\ZZ/n'\xrightarrow{\sim} \ZZ/n'(i){\quoi}.\tag{1.1.1.2}
\end{equation*}\etiquette[1.1.1.2]{XVIII.eq1.1.1.2}

If $F$ is an abelian sheaf such that $nF=0$, set
\begin{equation*}
F(i)=F\otimes_{\ZZ/n}\ZZ/n(i){\quoi}.\tag{1.1.1.3}
\end{equation*}\etiquette[1.1.1.3]{XVIII.eq1.1.1.3}%
If already $n'F=0$, the isomorphism \eqhyperref[XVIII.eq1.1.1.2]{1.1.1.2} allows us to identify
$F\otimes_{\ZZ/n}\ZZ/n(i)$ with
${F\otimes_{\ZZ/n'}\ZZ/n'(i)}$, so that $F(i)$ depends only
on $F$ and not on the choice of $n$. More generally, if $F$ is a
torsion sheaf prime to the residual characteristics of X
(\hyperrefext[XVII][XVII.subsec0.13]{0.13}), and if ${}_{m}F$ is the kernel of multiplication
by $m$ on $F$, we define $F(i)$ to be the inductive limit, as $m$ tends
multiplicatively to infinity through the integers invertible on $X$:
\begin{equation*}
F(i)=\varinjlim_{m}{}_{m}F(i){\quoi}.\tag{1.1.1.4}
\end{equation*}\etiquette[1.1.1.4]{XVIII.eq1.1.1.4}

The functor $F\mapsto F(i)$ is exact. If $f:Y\rightarrow X$ is a
morphism of schemes (\resp a quasi-compact quasi-separated morphism
of schemes), and if $F$ is a torsion sheaf on $X$ (\resp on $Y$),
prime\pageoriginale to the residual characteristics of $X$, then
$f^{\ast}(F(i))\simeq(f^{\ast}F)(i)$ (\resp
$(f_{\ast}F)(i)\simeq f_{\ast}(F(i))$). If $\mathcal{A}$ is a sheaf of
rings on $X$, annihilated by an integer $n$ invertible on $X$, this
extends to an exact functor $K\mapsto K(i)$ from $D(X,\mathcal{A})$ to
$D(X,\mathcal{A})$. This functor commutes with inverse-image functors
and direct-image functors.

The functors $F\mapsto F(i)$ are sometimes called ``\emph{Tate twist}'' functors.

\begin{enonce*}{Definition 1.1.2}\etiquette[1.1.2]{XVIII.def1.1.2}
A \emph{flat curve} over a scheme $S$ is a flat, separated morphism
$f:X\rightarrow S$ of finite presentation whose fibers are of pure
dimension one.
\end{enonce*}

The expression ``of pure dimension one'' does not exclude the empty
scheme. Contrary to \EGA II 7.4.2, a curve over a field is allowed here
to be empty, but is required to be separated.

When $S$ is the spectrum of a field, we shall simply speak of a curve
over $S$; a curve over a field is quasi-projective\nde{In \EGA II
7.4.10, this statement is proved only under the additional assumption
that $X$ is normal, and it is announced (\EGA II 7.4.12) that a proof in
the general case will appear in \EGA V, which has not been published.
The reader may fill this gap as an exercise. In any case, this has no
bearing on the construction of the trace morphism, since it can be
carried out first locally (on quasi-projective open subschemes) and the
trace morphisms can then be glued as in the proof of
\hyperref[XVIII.prop1.1.6]{1.1.6}.}.

\setcounter{subsubsection}{2}
\subsubsection{}\etiquette[1.1.3]{XVIII.subsubsec1.1.3}
Let $X$ be a curve over an algebraically closed field $k$ of
characteristic exponent $p$ prime to an integer $n\geq 1$. First
suppose that $X$ is reduced, and let $\bar{X}$ be a complete curve over
$k$ in which $X$ is a dense \emph{open subscheme}.

The complement $Y=\bar{X}-X$ is of dimension $0$. The exact cohomology
sequence (\hyperrefext[XVII][XVII.eq5.1.16.3]{5.1.16.3}) therefore induces an isomorphism
\begin{equation*}
H^{2}_{c}(X,\ZZ/n(1))\xrightarrow{\sim} H^{2}(\bar{X},\ZZ/n(1)){\quoi}.\tag{1.1.3.1}
\end{equation*}\etiquette[1.1.3.1]{XVIII.eq1.1.3.1}%
We also have an isomorphism \hyperrefext[IX][IX.cor4.7]{4.7} (given by Kummer
theory)
\begin{equation*}
H^{2}(\bar{X},\ZZ/n(1))=\Pic(\bar{X})/n=(\ZZ/n)^{c}\tag{1.1.3.2}
\end{equation*}\etiquette[1.1.3.2]{XVIII.eq1.1.3.2}%
where $c$ denotes the set of irreducible components of $X$, or of
$\bar{X}$, which amounts to the same thing.

If\pageoriginale $X$ is no longer necessarily reduced, and if $c$ is
the set of its irreducible components, then for $i\in c$ let $n_{i}$
denote the multiplicity of $X$ at the generic point $\eta_{i}$ of the
$i^{\text{\rm th}}$ irreducible component of $X$:
$$
n_{i}=\lg(O_{X,\eta_{i}}){\quoi}.
$$

We know that the restriction map from $H^{2}_{c}(X,\ZZ/n(1))$ to
$H^{2}_{c}(X_{\red}, \ZZ/n(1))$ is an isomorphism. Thus,
via \eqhyperref[XVIII.eq1.1.3.1]{1.1.3.1} and \eqhyperref[XVIII.eq1.1.3.2]{1.1.3.2} applied to
$X_{\text{\rm red}}$, there is a canonical isomorphism between
$H^{2}_{c}(X,\ZZ/n(1))$ and $(\ZZ/n)^{c}$. We denote by $\Tr_{X}$,
or simply by $\Tr$, the composite map
\begin{equation*}
\Tr:H^{2}_{c}(X,\ZZ/n(1))\xrightarrow{\sim}(\ZZ/n)^{c}\xrightarrow{t}\ZZ/n{\quoi},\tag{1.1.3.3}
\end{equation*}\etiquette[1.1.3.3]{XVIII.eq1.1.3.3}%
where
$$
t((a_{i})_{i\in c})=\sum n_{i}a_{i}{\quoi}.
$$

If $F$ is an abelian group annihilated by $n$, there is an
isomorphism
$$
H^{2}_{c}(X,F(1))\mathop{\longleftrightarrow}^{\sim}F\otimes
H^{2}_{c}(X,\ZZ/n(1))
$$
and hence, via \eqhyperref[XVIII.eq1.1.3.3]{1.1.3.3}, a \emph{trace morphism}
\begin{equation*}
\Tr:H^{2}_{c}(X,F(1))\longrightarrow F{\quoi}.\tag{1.1.3.4}
\end{equation*}\etiquette[1.1.3.4]{XVIII.eq1.1.3.4}

If $n=n'd$, the diagram
$$\vcenter{\etiquette[1.1.3.5]{XVIII.eq1.1.3.5}
\xymatrix@R=1.25cm{
0  \ar[r] & 𝛍_n  \ar[r]\ar[d]_{x^d} & \GG_m \ar[d]_{x^d} \ar[r]^{x^n} & {\GG}_m 
\ar[r] \ar@{=}[d]& 0\\
0  \ar[r] & 𝛍_{n'}  \ar[r] & \GG_m  \ar[r]^{x^{n'}} & {\GG}_m 
\ar[r] & 0
}}\leqno{(1.1.3.5)}
$$
is commutative, and therefore so is the diagram
$$\vcenter{\etiquette[1.1.3.6]{XVIII.eq1.1.3.6}
  \xymatrix@C=2cm@R=1.25cm{
H^{2}_c(X,\ZZ/n(1)) \ar[d] \ar[r]^-{\sim} & (\ZZ/n)^{c} \ar[d]\\
H^{2}_c(X,\ZZ/n'(1)) \ar[r]^-{\sim}  & (\ZZ/n')^{c}
}}\leqno{(1.1.3.6)}
$$
by\pageoriginale the definition in \hyperrefext[IX][IX.cor4.7]{4.7} of the isomorphism
\eqhyperref[XVIII.eq1.1.3.2]{1.1.3.2}.

The map \eqhyperref[XVIII.eq1.1.3.4]{1.1.3.4} thus depends only on $F$, and not on the choice
of the integer $n$ prime to $p$ such that $nF=0$. By passage to the
limit, we define it for every torsion abelian group $F$ prime to $p$.

The following lemma follows immediately from the definitions.

\begin{enonce*}{Lemma 1.1.4}\etiquette[1.1.4]{XVIII.lem1.1.4}
Under the preceding hypotheses, if, for $i\in c$, $U_{i}$ is a nonempty
open subscheme of $X$ contained in its $i^{\text{\rm th}}$ irreducible
component, the diagram
$$
\xymatrix@R=1.5cm@C=1.25cm{
\displaystyle{\mathop{\oplus}_{i\in c}} H^2_c(U_i, F(1))
\ar[rr]^{\sim} \ar[rd]_{\displaystyle{\mathop{\oplus}_{i\in c}}
  Tr_{U_i}}& & H^2_c(X,F(1))\ar[ld]^{Tr_X}\\
 & F &
}
$$
is commutative.
\end{enonce*}

\begin{enonce*}{Lemma 1.1.5}\etiquette[1.1.5]{XVIII.lem1.1.5}
Let $X$ and $Y$ be two curves over an algebraically closed field $k$,
let $u:X\rightarrow Y$ be a quasi-finite flat morphism, and let $n$ be
an integer prime to the characteristic exponent $p$ of $k$.

The following diagram, in which $\Tr_{u}$ is the map
\emph{(\hyperrefext[XVII][XVII.thm6.2.3]{6.2.3}),} is then commutative:
$$
\xymatrix@R=1.5cm{
  H^2_c(X,\ZZ/n(1)) \ar[rr]^{Tr_u} \ar[rd]_{Tr_X}&              & H^2_c(Y,\ZZ/n(1)) \ar[ld]^{Tr_Y}\\
                         & \ZZ/n &. 
}
$$
\end{enonce*}

Lemma \hyperref[XVIII.lem1.1.4]{1.1.4} allows us to reduce to the case where $X$ and $Y$ are
irreducible, with $X_{\text{\rm red}}$ and $Y_{\text{\rm red}}$ smooth
as well. If $e$ (\resp $f$) is the multiplicity of $X$ (\resp $Y$) at
the generic point, and if $e=df$, the triangles marked $+$ in the
following diagram are commutative:
$$
\vcenter{\etiquette[1.1.5.1]{XVIII.eq1.1.5.1}
\xymatrix@R=1.5cm{
H^2_c(X,\ZZ/n(1))
\ar[rr]^{\sim}_(.5){{\displaystyle\mathop{\phantom{a}}_{\displaystyle\mathop{{\phantom{a}}}_{{\displaystyle\mathop{+}}}}}}  \ar[dd]_{\Tr_u} \ar[rd]_{\Tr_X}& &
H^2_c(X_{\red},\ZZ/n(1)) \ar[dd]^{d.\Tr_{u_{\red}}} \ar[ld]^{e.\Tr_{X_{\red}}}\\
                               & \ZZ/n &\\
H^2_c(Y,\ZZ/n(1)) \ar[rr]^{\displaystyle\mathop{\sim}^{\displaystyle\mathop{\phantom{a}}^{{\displaystyle\mathop{\phantom{a}}^{\textstyle{+}}}}}} \ar[ru]^{\Tr_Y}&  & H^2_c(Y_{\red},\ZZ/n(1))\ar[lu]_{f.\Tr_{Y_{\red}}}, 
}}\sisi{}{\leqno{(1.1.5.1)}}
$$

Its\pageoriginale outer boundary is commutative, as follows from the
commutativity of the diagram
$$
\xymatrix@C=2cm@R=1.5cm{
u_!\ZZ/n_X(1)           \ar[r]^{\sim} \ar[d]^{\Tr_u} &
u_{{\red}_!}\ZZ/n_{X_{\red}}(1) \ar[d]^{d.\Tr_{u_{\red}}}\\
\ZZ/n_Y(1)  \ar[r]^{\sim}& \ZZ/n_{Y_{\red}}(1).}
$$
It suffices to verify this last compatibility at the generic point of
$Y$, which is trivial.

\sisi{}{To establish the commutativity of the left-hand triangle in
\eqhyperref[XVIII.eq1.1.5.1]{1.1.5.1}, it therefore remains only to establish that of the
right-hand triangle. Hence we may now assume in addition that $X$ and $Y$
are reduced\nde{This paragraph was added in this edition, and the next
two were modified.}.}

Let $\bar{X}$ and $\bar{Y}$ be the complete \sisi{non-smooth}{smooth}
curves containing \sisi{$X_{\red}$}{$X$} and \sisi{$Y_{\red}$}{$Y$} as dense
open subschemes (\EGA II 7.4.11). The morphism \sisi{$u_{\red}$}{$u$}
extends to a flat morphism $\bar{u}:\bar{X}\rightarrow \bar{Y}$
(\EGA II 7.4.9), and by \hyperref[XVIII.lem1.1.4]{1.1.4} it suffices to verify
\hyperref[XVIII.lem1.1.5]{1.1.5} for $\bar{u}$.

By (XVII \hyperref[XVII.eq6.3.18.2]{6.3.18.2}), the diagram of sheaves on $\bar{Y}$ below, in
which $N$ denotes the norm,
$$
\xymatrix@R=1.25cm{
0  \ar[r] & \bar{u}_{\ast}𝛍_n  \ar[r] \ar[d]^{\Tr_{\bar{u}}}&
\bar{u}_{\ast}\GG_m  \ar[r]^{x^n} \ar[d]^{N}&
\bar{u}_{\ast} \GG_m  \ar[r] \ar[d]^{N}& 0 \\
0  \ar[r] & 𝛍_n  \ar[r] & \GG_m  \ar[r]^{x^n} & \GG_m 
\ar[r] & 0 ,
}
$$
is commutative, and hence so is the diagram
$$
\xymatrix@C=1.5cm@R=1.25cm{
\Pic(\bar{X})/n \ar[r]^(.45){\sim}\ar[d]^{N}& H^2_c(\bar{X},\ZZ/n(1))
\ar[d]^{\Tr_{\bar{u}}}\\ 
\Pic(\bar{Y})/n \ar[r]^(.45){\sim}& H^2_c(\bar{Y},\ZZ/n(1)),\\
}
$$
and\pageoriginale it remains to observe that if $\Ll$ is an invertible sheaf
of degree 1 on $\bar{X}$, then its norm again has degree $1$ (this fact
following from the compatibility \sisi{\EGA IV 21.10.7}{\EGA IV 21.10.17},
taking into account the definition \EGA IV 21.5.5 of the norm of a
divisor).

\begin{enonce*}{Proposition 1.1.6}\etiquette[1.1.6]{XVIII.prop1.1.6}
There is a unique way to define, for every compactifiable flat curve
$f:X\rightarrow S$ and every torsion sheaf $F$ on $S$ prime to the
residual characteristics of $S$, a trace morphism
$$
\Tr_{f}:R^{2}f_{!}(f^{\ast}F(1))\longrightarrow F{\quoi},
$$
functorial in $F$, whose formation is compatible with every base change
\emph{(\cf \hyperrefext[XVII][XVII.thm6.2.3]{6.2.3}, (VAR 2))}, and which, when $S$ is the spectrum
of an algebraically closed field, coincides with the trace morphism
\eqhyperref[XVIII.eq1.1.3.4]{1.1.3.4}.
\end{enonce*}

The morphism $\Tr_{f}$ is prescribed fiber by fiber; its uniqueness is
therefore clear.

Let $p:\PP^{1}_{S}\rightarrow S$ be the projective line over $S$, and let $n$ be
an integer invertible on $S$. We have
$\sheafPic_{S}(\PP^{1}_{S})=\ZZ_{S}$, whence Kummer theory
(\hyperrefext[IX][IX.Thr3.2]{3.2}) gives a morphism $\ZZ/n\rightarrow
R^{1}p_{\ast}\ZZ/n(1)$ which one checks fiber by fiber
(\hyperrefext[XII][XII.cor5.2]{5.2}) to be an isomorphism. Its inverse $\Tr_{p}$ induces the trace
morphism \eqhyperref[XVIII.eq1.1.3.4]{1.1.3.4} fiber by fiber.

Let $f:X\rightarrow S$ be a flat curve over $S$ such that there exists a
quasi-finite flat $S$-morphism $u$ from $X$ to $\PP^1_S$. By
\hyperref[XVIII.lem1.1.5]{1.1.5}, the composite morphism $\Tr_{p}\circ \Tr_{u}$
(\hyperrefext[XVII][XVII.thm6.2.3]{6.2.3}) has, at every geometric point of $S$, the trace morphism
\eqhyperref[XVIII.eq1.1.3.4]{1.1.3.4} as its fiber. In particular, it does not depend on the
choice of $u$.

Suppose\pageoriginale that $X$ is a union of open subschemes $U_{i}$ such
that there exists a quasi-finite flat morphism from $U_{i}$ to
$\PP^{1}_S$. If $\alpha_{i}$ (\resp $\alpha_{ij}$) is the inclusion of
$U_{i}$ (\resp $U_{ij}=U_{i}\cap U_{j}$) in $X$, the sequence
$$
\sum R^{2}(f\alpha_{ij})_{!}\ZZ/n(1)\rightrightarrows
\sum R^2(f\alpha_{i})_{!}\ZZ/n(1)\rightarrow
R^{2}f_{!}\ZZ/n(1)\rightarrow 0{\quoi}.
$$
is exact, by (\hyperrefext[XVII][XVII.prop6.2.9]{6.2.9}) and the right exactness of
$R^{2}f_{!}$ (since $R^{3}f_{!}=0$). One checks fiber by fiber, using
\hyperref[XVIII.lem1.1.5]{1.1.5}, that the sum of the trace morphisms
$\underset{i}\sum R^{2}(f\alpha_{i})_{!}\ZZ/n(1)\rightarrow
\ZZ/n$ factors through a trace morphism $\Tr_{f}$ from
$R^{2}f_{!}\ZZ/n(1)$ which induces \eqhyperref[XVIII.eq1.1.3.4]{1.1.3.4} fiber by fiber.

% END INCLUDED SGA4_English_current_buildable_source_20260906/source/18/drafts/SGA4_XVIII_LINES_44_450_English.texfrag

% BEGIN INCLUDED SGA4_English_current_buildable_source_20260906/source/18/drafts/SGA4_XVIII_LINES_451_855_English.texfrag

For a general $X$, let $j:X'\rightarrow X$ be the largest open subset
of $X$ such that $f|X'$ is a Cohen--Macaulay morphism (\cf \EGA IV
12.1.1 (vi)). Every point $x$ of $X'$ has a neighborhood $U_{x}$ in
$X'$ for which there exists a quasi-finite $S$-morphism $u_{x}$ from
$U_{x}$ to $\PP^{1}_{S}$. This morphism is automatically flat (\EGA IV
11.3.10 and 15.4.2), and $X'$ therefore satisfies the preceding
hypotheses. Moreover, $X'$ is fiberwise dense in $X$; the morphism from
$R^{2}(fj)_{!}\ZZ/n(1)$ to $R^{2}f_{!}\ZZ/n(1)$ is consequently an
isomorphism, as can be checked fiberwise, and the composite arrow
$$
\Tr_{f}:R^{2}f_{!}\ZZ/n(1)\xleftarrow{\sim}
R^{2}(fj)_{!}\ZZ/n(1)\xrightarrow{\Tr_{fj}}\ZZ/n
$$
induces \eqhyperref[XVIII.eq1.1.3.4]{1.1.3.4} on every fiber.

If $F$ satisfies $nF=0$, one has (\hyperrefext[XVII][XVII.thm5.2.6]{5.2.6})
\begin{equation*}
R^{2}f_{!}\ZZ/n(1)\otimes F\xrightarrow{\sim}
R^{2}f_{!}(f^{\ast}F(1))\tag{1.1.6.1}
\end{equation*}\etiquette[1.1.6.1]{XVIII.eq1.1.6.1}%
and $\Tr_{f}$ is defined to be the composite arrow
$$
\Tr_{f}:R^{2}f_{!}(f^{\ast}F(1))\xleftarrow{\sim}
R^{2}f_{!}\ZZ/n(1) \otimes F\rightarrow F{\quoi}.
$$
In the general case, $\Tr_{f}$ is defined as the filtered colimit of
the trace morphisms relative to the sheaves ${}_nF$, for $n$
invertible on $S$. These arrows induce \eqhyperref[XVIII.eq1.1.3.4]{1.1.3.4} on every fiber
and therefore commute with every base change, \cqfd

One\pageoriginale checks fiberwise, using \hyperref[XVIII.lem1.1.4]{1.1.4},
\hyperref[XVIII.lem1.1.5]{1.1.5}, and \eqhyperref[XVIII.eq1.1.6.1]{1.1.6.1}, the following results:

\begin{enonce*}{Lemma 1.1.7}\etiquette[1.1.7]{XVIII.lem1.1.7}
Let $f:Y\rightarrow S$ be a compactifiable flat curve over $S$, let
$u:X\rightarrow Y$ be a quasi-finite flat morphism of finite
presentation, and let $F$ be a torsion sheaf on $S$, of torsion prime
to the residue characteristics of $S$. The following diagram is
commutative:
$$
\xymatrix@C=1.5cm@R=1.25cm{
R^2(fu)_! (fu)^{\ast}F(1) \ar[d]^{\Tr_{fu}}\ar@{=}[r]& R^2f_!(u_!
u^{\ast})f^{\ast}F(1) \ar[r]^{\Tr_u}& R^2f_!f^{\ast}F(1)\ar[d]^{\Tr_f} \\
F     \ar@{=}[rr] &&  F,
}
$$
\end{enonce*}

\begin{enonce*}{Lemma 1.1.8}\etiquette[1.1.8]{XVIII.lem1.1.8}
Let $u:Y\rightarrow S$ be a quasi-finite flat morphism of finite
presentation, let $f:X\rightarrow Y$ be an $S$-compactifiable flat
curve \emph{(\hyperrefext[XVII][XVII.def3.2.1]{3.2.1})}, and let $F$ be as in \hyperref[XVIII.lem1.1.7]{1.1.7}. The
following diagram is commutative\nde{This compatibility merits a few
arguments. To verify it, one may first reduce by base change to the case
where $S$ is the spectrum of an algebraically closed field $k$. One may
then suppose that $Y$ is connected and nonempty, in which case $Y\to S$
has a (unique) section $S\to Y$ that is a closed immersion defined by a
nilpotent ideal. The arguments in the proof of \hyperref[XVIII.prop1.1.6]{1.1.6}
allow us to suppose that $X=\mathbf{P}^1_Y$. It remains only to observe
that the multiplicity of $Y$ equals the multiplicity of $X$ at its
generic point. See also the proof of the commutativity of
\eqhyperref[XVIII.eq2.9.3]{2.9.3}.}:
$$
\xymatrix@C=1.5cm@R=1.25cm{
R^{2}(uf)_! (uf)^{\ast}F(1)  \ar@{=}[r] \ar[d]^{\Tr_{uf}}&
u_!R^2f_!f^{\ast} u^{\ast}F(1) \ar[r]^{\Tr_f} &
u_!u^{\ast}F \ar[d]^{\Tr_u}\\
F        \ar@{=}[rr] & &  F
}
$$
\end{enonce*}

One similarly verifies from \eqhyperref[XVIII.eq1.1.3.2]{1.1.3.2}:

\begin{enonce*}{Lemma 1.1.9}\etiquette[1.1.9]{XVIII.lem1.1.9}
If $f:X\rightarrow S$ is a compactifiable smooth curve and the
geometric fibers of $f$ are irreducible, then the trace morphism
(\hyperref[XVIII.prop1.1.6]{1.1.6}) is an isomorphism.
\end{enonce*}

\subsection{$1$-acyclicity of projective
  space}\etiquette[1.2]{XVIII.subsec1.2}

In this section, the fppf topology is in principle always used
(\hyperrefext[XVII][XVII.subsec0.10]{0.10}).

\subsubsection{}\etiquette[1.2.1]{XVIII.subsubsec1.2.1}

Let $\epsilon$ be a locally free module on a scheme $S$, everywhere of
rank\pageoriginale $\geq 2$, let $p:\PP(\epsilon)\rightarrow S$ be the
corresponding projective bundle, and let $G$ be an abelian sheaf on the
big fppf site of $S$ (\hyperrefext[XVII][XVII.subsec0.10]{0.10}). Each pair $(K_{0},\varphi)$ consisting of a
torsor $K_{0}$ under $G$ and a homomorphism
$\varphi:\GG_{mS}\rightarrow G$ defines a $p^{\ast}G$-torsor
$K(K_{0},\varphi)$ on $\PP(\epsilon)$, the sum of $K_{0}$ and the image
under $\varphi$ of the $\GG_{m}$-torsor $\Oo(1)$\nde{In this section,
we implicitly identify line bundles and $\GG_{m}$-torsors by the
correspondence that associates with a line bundle the sheaf of its
invertible sections, on which $\GG_{m}$ acts by multiplication. A reader
concerned with details may wish to know that, following Grothendieck's
convention, the projective scheme $\PP(\epsilon)$ is the bundle of
hyperplanes in $\epsilon$, and likewise the punctured vector bundle
$\VV(\epsilon)^{\ast}\to \PP(\epsilon)$ is an open subset of
$\VV(\epsilon)$, which is the spectrum of the symmetric algebra of
$\epsilon$.}:
\begin{equation*}
K(K_{0},\varphi)=K_{0}+\varphi \Oo(1){\quoi}.\tag{1.2.1.1}
\end{equation*}\etiquette[1.2.1.1]{XVIII.eq1.2.1.1}%
The pairs $(K_{0},\varphi)$ are made into a category by setting
$$
\Hom((K_{0},\varphi),(K'_{0},\varphi'))=
\begin{cases}
∅\quad\mbox{\rm if}\quad \varphi\neq \varphi'\\
\Hom(K_{0},K'_{0})\quad\mbox{\rm if}\quad \varphi=\varphi'{\quoi}.
\end{cases}
$$
The torsor $K(K_{0},\varphi)$ then depends functorially on
$(K_{0},\varphi)$. It is compatible with localization and therefore
defines a functor from the stack over $S$ of pairs $(K_{0},\varphi)$ to
the stack over $S$ of $p^{\ast}G$-torsors on $\PP(\epsilon)$ (= the
stack whose sections over $U/S$ are the $p^{\ast}G$-torsors on
$\fprod{U}{\PP}{S}(\epsilon)$). These stacks are stacks in groupoids
(= every morphism over a $U/S$ is an isomorphism).

\begin{enonce*}{Theorem 1.2.2}\etiquette[1.2.2]{XVIII.thm1.2.2}
Under the hypotheses of \hyperref[XVIII.subsubsec1.2.1]{1.2.1}, if $G$ is a flat commutative
group scheme of finite presentation over $S$, then the functor $K$ of
\hyperref[XVIII.subsubsec1.2.1]{1.2.1} is an equivalence.
\end{enonce*}

\setcounter{subsubsection}{2}
\paragraph{}\etiquette[1.2.2.1]{XVIII.par1.2.2.1}

This assertion is equivalent to the conjunction of
\begin{gather*}
G\xrightarrow{\sim} p_{\ast}p^{\ast}G\quad\mbox{\rm and}\\
\Hhom(\GG_{m},G)\xrightarrow{\sim} R^{1}p_{\ast}p^{\ast}G{\quoi}.
\end{gather*}
Indeed, a morphism of stacks in groupoids is an equivalence if and
only if it induces an isomorphism on the sheaf generated by the
presheaf of isomorphism classes of objects (the 2\textsuperscript{nd}
formula) and an isomorphism on the sheaf of automorphisms of any local
object (the 1\textsuperscript{st} formula).

We shall give two proofs (\hyperref[XVIII.lem1.2.4]{1.2.4} and
\hyperref[XVIII.subsubsec1.2.5]{1.2.5}) of full faithfulness\nde{With regard to
essential surjectivity, the reader should be warned that the proof will
use the general proof technique explained in the introduction to the
\EGA.}.

\begin{enonce*}{Lemma 1.2.3}\etiquette[1.2.3]{XVIII.lem1.2.3}
Let\pageoriginale $f:X\rightarrow S$ be a proper flat morphism. Suppose
that $f_{\ast}\Oo_{X}=\Oo_{S}$ universally and
$R^{1}f_{\ast}\Oo_{X}=0$ universally. Every $S$-morphism from $X$ to a
finite-type $S$-group scheme $G$ then factors through $f$ and a section
of $G$; more precisely,
$G\xrightarrow{\sim} f_{\ast}f^{\ast}G$.
\end{enonce*}

It suffices to prove the 1\textsuperscript{st} assertion.

By the rigidity lemma (\textsc{Mumford} \cite[6.1. p.~115]{sga4.exp.XVIII.bib.4}; note
that \textsc{Mumford} does not use the Noetherian hypothesis imposed on
$S$), it suffices to treat the case where $S$ is the spectrum of an
algebraically closed field. Suppose first that $G$ is an abelian
scheme, and let $G^{\ast}$ be the dual abelian scheme. Since
$G=G^{{\ast}{\ast}}$, giving $\varphi:X\rightarrow G$ is equivalent to
giving an invertible sheaf $\Ll$ on $X\times G^{\ast}$, trivialized
along $X\times \{e\}$ and algebraically equivalent to zero on the
fibers of the projection $\pr_{1}$ from $X\times G^{\ast}$ to $X$. If
$x:S\to X$ is a rational point of $X$, then
$\Ll\sim \pr^{\ast}_{2}(x\times G^\ast)^{\ast}\Ll\otimes \Mm$, with
$\Mm$ trivialized along $x\times G^{\ast}$ and $X\times \{e\}$. Giving
$\Mm$ is equivalent to giving a morphism of schemes from $G^{\ast}$ to
$\sheafPic_{X/S}$ that takes $e$ to $e$. Since
$H^{1}(X,\Oo_{X})=0$, one has $\sheafPic^0_{X/S}=0$\nde{This follows
from the canonical isomorphism
$\Lie(\sheafPic_{X/S})=H^{1}(X,\Oo_{X})=0$, obtained by comparing both
groups with the kernel of $\Pic(X\times_k
\Spec (k[\epsilon]))\to \Pic(X)$ (where $S=\Spec(k)$).
For further details on the properties of Picard schemes, see
\cite{sga4.exp.XVIII.bib.2bis}.},
$\Mm\sim\Oo$, and $\Ll\sim \pr^{\ast}_{2}\Ll'$, with $\Ll'$
\sisi{equivalent}{algebraically equivalent} to zero on $G^{\ast}$,
hence defining $g\in G(S)$ such that $\varphi=g\circ f$.

In the general case, $G$ is an extension of an abelian scheme $A$ by
an affine group $G_0$. The preceding result reduces us to the case
$A=0$, and then
$$
\Hom_{S}(\Spec (H^0(X,\Oo_{X}),G)\xrightarrow{\sim}
\Hom_{S}(X,G){\quoi}.
$$

\begin{enonce*}{Lemma 1.2.4}\etiquette[1.2.4]{XVIII.lem1.2.4}
Under the hypotheses of \hyperref[XVIII.subsubsec1.2.1]{1.2.1}, if $G$ is a group scheme
of finite presentation over $S$ or is defined by a quasi-coherent sheaf
on $S$, then the functor \eqhyperref[XVIII.eq1.2.1.1]{1.2.1.1} is fully faithful.
\end{enonce*}

\begin{itemize}
\item[A.] We prove that, if $K(K_{0},\varphi)$ is isomorphic to
$K(K'_0,\varphi')$, then $\varphi=\varphi'$. For quasi-coherent $G$,
one has $\varphi=\varphi'=0$, and there is nothing to prove. For $G$ of
finite presentation, by (\SGA 3 IX 5.1), it suffices\pageoriginale to
prove the assertion when $S$ is the spectrum of an algebraically
closed field $k$. Let $T$ be the largest subtorus of $G$, and put
$H=G/T$. The cohomology exact sequence gives
$$
0\rightarrow H^{0}(\PP(\epsilon),T)\rightarrow
H^{0}(\PP(\epsilon),G) \xrightarrow{\alpha}
H^{0}(\PP(\epsilon),H)
 \xrightarrow{\delta}
H^{1}(\PP(\epsilon),T)\xrightarrow{\beta}H^{1}(\PP(\epsilon),G){\quoi}.
$$
By \hyperref[XVIII.lem1.2.3]{1.2.3}, $H^0(\PP(\epsilon),G)\sim G(k)$ and
$H^0(\PP(\epsilon),H)\sim H(k)$, so $\alpha$ is surjective and
$\delta$ is zero. If $Y(T)=\Hom(\GG_{m},T)$, then
$T\sim G_{m}\otimes Y(T)$ and $H^{1}(\PP(\epsilon),T)\sim Y(T)$.
Thus the injectivity of $\beta$ expresses the injectivity of the
canonical arrow defined by $\Oo(1)$,
$\Hom(\GG_{m},G)\simeq \Hom(\GG_{m},T)\rightarrow
H^{1}(\PP(\epsilon),G)$, proving the assertion.

\item[B.] Proving that the functor \eqhyperref[XVIII.eq1.2.1.1]{1.2.1.1} is fully faithful is
local on $S$. We may therefore suppose that $\PP(\epsilon)/S$ has a
section $s$\sisi{}{ and choose an isomorphism
$s^\ast\Oo(1)\sim \Oo$\nde{This additional necessary datum, which was
missing from the original edition, has been added in this edition.}}.
If $K(K_{0},\varphi)$ is isomorphic to $K(K'_{0},\varphi')$, then
$\varphi=\varphi'$ by A, and
$K_{0}\sim s^{\ast}K(K_{0},\varphi)\sim
s^{\ast}K(K'_{0},\varphi')\sim K'_{0}$. By the formula
\begin{equation*}
\Hom (K(K_{0},\varphi), K(K'_{0},\varphi'))\sim \Hom(K(K_{0}-K'_{0},\varphi-\varphi'),G_{\PP(\epsilon)}){\quoi},\tag{1.2.4.1}
\end{equation*}\etiquette[1.2.4.1]{XVIII.eq1.2.4.1}%
it remains only to show that
$$
\Hom_{\tors.}(G,G)\xrightarrow{\sim}
\Hom_{\tors.}(p^{\ast}G,p^{\ast}G){\quoi}.
$$
This follows from \hyperref[XVIII.lem1.2.3]{1.2.3}.
\end{itemize}

\setcounter{subsubsection}{4}
\subsubsection{}\etiquette[1.2.5]{XVIII.subsubsec1.2.5}

Here is another, more \sisi{Kunneth-like}{Künneth-like}, proof of full
faithfulness. The question is local on $S$; by \eqhyperref[XVIII.eq1.2.4.1]{1.2.4.1}, it therefore
suffices to check that
$$
\Hom_{\PP(\epsilon)}(G,K(G,\varphi))
\begin{cases}
=∅ &\text{\rm if}\quad \varphi\neq 0\\
\xleftarrow{\sim} G(S)\ &\text{\rm if}\quad \varphi= 0{\quoi}.
\end{cases}
$$
This is precisely what the following lemma says\nde{This translation
merits some explanation. The line bundle $\Oo(1)$ on
$\PP(\epsilon)$ becomes canonically trivialized after pullback by the
canonical projection $\VV(\epsilon)^{\ast}\to \PP(\epsilon)$. After
pullback to $\VV(\epsilon)^{\ast}$, the $G$-torsor $K(G,\varphi)$ is
therefore canonically trivialized as well. \emph{Via} this
identification, an isomorphism $G\to K(G,\varphi)$ thus induces an
automorphism of the trivial $G$-torsor over
$\VV(\epsilon)^{\ast}$, that is, an $S$-morphism
$\VV(\epsilon)^{\ast}\to G$. To descend from $\PP(\epsilon)$, this
morphism must satisfy a compatibility like the one stated, and the
converse follows by faithfully flat descent. Strictly speaking, it
seems to the reviewer that the correct condition should be
$\Psi(\lambda v)=\varphi(\lambda)^{-1}\Psi(v)$, but this has no effect,
since one can apply the lemma to $\Psi^{-1}$ instead of $\Psi$. The
condition $\Psi(\lambda v)=\varphi(\lambda)\Psi(v)$ in the original
edition also gives a correct statement (since $\varphi$ can be replaced
by $\varphi^{-1}$), but it has been changed so that the proof need not
be modified.}:

\begin{enonce*}{Lemma 1.2.6}\etiquette[1.2.6]{XVIII.lem1.2.6}
Let $\epsilon$ be a locally free module everywhere of rank $\geq 2$ on
$S$, let $G$ be a group scheme of finite presentation over $S$, let
$\varphi:\GG_{m}\rightarrow G$ be a group homomorphism, let
$\VV(\epsilon)$ be the vector bundle defined by $\epsilon$, let
$\VV(\epsilon)^{\ast}$\pageoriginale be the corresponding punctured
vector bundle, and let $\Psi:\VV(\epsilon)^{\ast}\rightarrow G$ be an
$S$-morphism satisfying the following identity between $S$-morphisms
from $\GG_{m}\times \VV(\epsilon)^{\ast}$ to $G$:
$$
\Psi(\lambda v)=\sisi{\varphi (\lambda)\Psi(v)}{\Psi(v)\varphi(\lambda)}{\quoi}.
$$
Then $\varphi=e$, and $\Psi$ factors through the projection $p$ from
$\PP(\epsilon)$ to $S$.
\end{enonce*}

It suffices to prove the second assertion. We reduce at once to the
case where $S$ is Noetherian and then, by (\SGA 3 IX 5.1) and the
rigidity lemma (\textsc{Mumford} \cite[6.1 p.~115]{sga4.exp.XVIII.bib.6}) applied to
$\PP(\epsilon)$, to the case where $S$ is the spectrum of a field. Let
$n$ be an integer, let $W_{n}$ and $G_{n}$ be the $n^{\text{\rm th}}$
infinitesimal neighborhoods of the zero section in $\VV(\epsilon)$ and
$G$, let $\mathfrak{g}_{n}$ be the affine algebra of $G_{n}$, and let
$P_{n}(\Psi):\VV(\epsilon)^{\ast}\rightarrow
\Hhom_{S}(W_{n},G)$ be the principal part of $n^{\text{\rm th}}$ order of
$\Psi$\nde{More concretely, $P_n(\Psi)$ is the morphism corresponding
by adjunction to the morphism $\VV(\epsilon)^{\ast}\times_S W_n\to G$
(which we also denote by $P_n$) obtained by composing $\Psi$ with the
morphism $\VV(\epsilon)^{\ast}\times_{S} W_n\to
\VV(\epsilon)^\ast$ induced on these subschemes by addition on
$\VV(\epsilon)$. In terms of the functors of points of these
$S$-schemes, this can be interpreted more concretely by saying that,
if $v$ is a section of $\VV(\epsilon)^{\ast}$ and $t$ a section of
$W_n$, then $P_n(\Psi)(v,t)=\Psi(v+t)$.}.
Every section of $\Hhom_{S}(W_{n},G)$ defines, by right translation by
the inverse of the image of the zero section, a morphism from $W_{n}$
to $G_{n}$ that takes $0$ to $e$; this defines
$$
\Psi_{1}=P_{n}(\Psi).\Psi^{-1}:\VV(\epsilon)^{\ast}\longrightarrow
\Hhom_{S}(W_{n},G_{n}){\quoi}.
$$
The hypothesis implies that $\Psi_{1}$ commutes with homotheties, which
act on the right-hand side by transport of structure from their action
on $W_{n}\subset \VV(\epsilon)$\nde{With the notation above,
$\Psi_1$ can be characterized by the relation
$\Psi_1(v,t).\Psi(v)=\Psi(v+t)$. It follows at once that, if $\lambda$
is a section of $\GG_{m}$, then
$\Psi_1(\lambda v,\lambda t)=\Psi_1(v,t)$, which is the asserted
commutation with homotheties provided that $\GG_{m}$ acts correctly on
$\Hhom_S(W_n,G_n)$, namely, that a section $\lambda$ of $\GG_{m}$ acts
by right composition with multiplication by the \emph{inverse} of
$\lambda$ on $W_n$.}. The affine algebra of $W_{n}$ is
$\sum^{n}_{i=0}\Sym^{i}(\epsilon)$. By homogeneity, $\Psi_{1}$ is
therefore defined by a section, on $\PP(\epsilon)$, of the
$\Oo_{\PP(\epsilon)}$-module
$$
\sum^{n}_{i=0}p^{\ast}\Hhom(\mathfrak{g}_{n},\Sym^{i}(\epsilon))(-i){\quoi}.\nde{To
verify this assertion in detail, one may begin by defining the pullback
of this section to $\VV(\epsilon)^{\ast}$; one then obtains
tautologically a section $s=(s_0,\dots,s_n)$ of the trivial bundle
obtained by tensoring the structure sheaf with the vector space
$\sum^{n}_{i=0}\Hhom(\mathfrak{g}_{n},\Sym^{i}(\epsilon))$. It satisfies
the relation $s_i(v)=\lambda^i s_i(\lambda v)$. This is indeed descent
data giving the desired section on $\PP(\epsilon)$. One can also argue
that the sections $s_i$ extend to all of $\VV(\epsilon)$ (since the
origin is a closed subset of codimension $\geq 2$ in
$\VV(\epsilon)$). By the density of $\GG_m$ in $\AAA^1$, the extended
sections $s_i$ satisfy the relation
$s_i(v)=\lambda^i s_i(\lambda v)$ for $\lambda\in\AAA^1$. Taking
$\lambda=0$ gives $s_i=0$ for $i>0$.}
$$
The components of index $i>0$ of this module have no section other
than the zero section, and it follows that
$$
\Psi_{1}=0.
$$
Since this holds for every $n$, it implies that $\Psi$ factors through
a section of $G$.

This second proof extends to prove the following analytic variant of
\hyperref[XVIII.lem1.2.4]{1.2.4}. When $G$ is affine, the proof of
\hyperref[XVIII.thm1.2.2]{1.2.2} can be given independently of this variant.

\setcounter{subsubsection}{6}
\subsubsection{}\etiquette[1.2.7]{XVIII.subsubsec1.2.7}

Let\pageoriginale $L$ be the fraction field of a complete discrete
valuation ring $\Oo$, let $U_{m}$ be the rigid analytic ``generic
fiber'' group of the formal completion of $\GG_{m}$, and let
$\Oo(1)^{\wedge}$ be the $U_{m}$-torsor on $\PP^{r}_{L}$ obtained by
formally completing $\Oo(1)$. If $G$ is a rigid analytic group over
$L$, the functor $(K_{0},\varphi)\mapsto
p^{\ast}K_{0}+\varphi\Oo(1)^{\wedge}$ (\cf \hyperref[XVIII.subsubsec1.2.1]{1.2.1}) is a fully
faithful functor from the category of pairs $(K_{0},\varphi)$
consisting of a $G$-torsor on $\Spec(L)$ and
$\varphi\in \Hom(U_m,G)$ to the category of $p^{\ast}G$-torsors on
$\PP^{r}_{L}$\acorriger{The reader may find an introduction to the
language of rigid analytic geometry in \cite{sga4.exp.XVIII.bib.5bis}. The nonexistence
of nontrivial sections of the sheaves $\Oo(n)$ for $n<0$ in this
context follows, for example, from a rigid version of GAGA; see
\cite{sga4.exp.XVIII.bib.3ter}.}.

\begin{enonce*}{Lemma 1.2.8}\etiquette[1.2.8]{XVIII.lem1.2.8}
Let $k$ be a field of characteristic exponent $p$, and let
$f:X\rightarrow \Spec(k)$ be a proper scheme over $k$ satisfying
$H^0(X,\Oo_{X})=k$ and $\sheafPic^{\tau}_{X}=0$\nde{A definition of
$\sheafPic^{\tau}$ in a more general setting will be found in
\hyperref[XVIII.def1.6.4]{1.6.4}.}. For every finite-type commutative
affine group scheme $G$ over $k$ containing no tori, one has
$R^{1}f_{\ast}(f^{\ast}G)=0$.
\end{enonce*}

% END INCLUDED SGA4_English_current_buildable_source_20260906/source/18/drafts/SGA4_XVIII_LINES_451_855_English.texfrag

% BEGIN INCLUDED SGA4_English_current_buildable_source_20260906/source/18/drafts/SGA4_XVIII_LINES_856_1250_English.texfrag

It suffices to show that after every base change $c:S\rightarrow
\Spec(k)$, the functor $f^{\ast}$ from the category of $G$-torsors on
$S$ to that of $G$-torsors on $X_{S}$ is an equivalence. This problem
is local on $\Spec (k)$, so that, in order to solve it, one may reduce
to the case where $X$ has a rational point, where $G$ is a successive
extension of the groups $\GG_{a}$, $\alpha_{p}$, $𝛍_{p}$,
$\ZZ/p$, and $\ZZ/\ell_{i}$, with $\ell_{i}$ prime to $p$,
and where $\ZZ/\ell_{i}\sim 𝛍_{\ell_i}$.

We know that $\Lie(\sheafPic_{X})=H^{1}(X,\Oo_{X})=0$. By
\hyperref[XVIII.lem1.2.3]{1.2.3}
and the existence of a section, the functors $f^{\ast}$ under
consideration are fully faithful; it therefore suffices to prove that
$R^{1}f_{\ast}(f^{\ast}G)=0$, or, by dévissage, that
$$
R^{1}f_{\ast}\GG_{a}=R^{1}f_{\ast}\alpha_{p}=R^{1}f_{\ast}𝛍_{p}=R^{1}f_{\ast}\ZZ/p=R^{1}f_{\ast}\ZZ/\ell_{i}=0{\quoi}.
$$

For $\GG_{a}$, this follows from $H^{1}(X,\Oo_{X})=0$. The case of
$\alpha_{p}$ (\resp $\ZZ/p$) follows from it via
\hyperref[XVIII.lem1.2.3]{1.2.3}, by the long exact cohomology sequence
defined by the exact sequence
\begin{align*}
& 0\longrightarrow \alpha_{p}\longrightarrow
  \GG_{a}\xrightarrow{x^{p}} \GG_{a}\longrightarrow 0\\
& 0\longrightarrow \ZZ/p\longrightarrow
  \GG_{a}\xrightarrow{x-x^{p}}\GG_{a}\longrightarrow
  0\qquad (p\neq 1)){\quoi}.\tag*{(\resp}
\end{align*}
The case of $\ZZ/\ell_{i}$ (\resp of $𝛍_{p}$) follows from the
isomorphism $\ZZ/\ell_{i}\sim 𝛍_{\ell_{i}}$,
from\pageoriginale \hyperref[XVIII.lem1.2.3]{1.2.3}, and from the long
exact sequence deduced from the Kummer exact sequence
$$
0\longrightarrow 𝛍_{n}\longrightarrow \GG_{m}\longrightarrow
\GG_{m}\longrightarrow 0{\quoi}.
$$

\begin{enonce*}{Lemma 1.2.9}\etiquette[1.2.9]{XVIII.lem1.2.9}
Let $S$ be a local Artinian scheme, let $\bar{s}$ be a geometric point
of $S$, let $f:X\rightarrow S$ be a proper flat morphism, and let $G$
be a flat commutative group scheme of finite type over $S$. Assume
that $f$ is smooth or that $G$ is affine. If
$H^0(X_{\bar{s}},\Oo)=k(\bar{s})$, if
$\sheafPic^{\tau}_{X_{\bar{s}}}=0$, and if $G_{\bar{s}}$ contains no
subgroup $\GG_{m}$, then the functor $f^{\ast}$ from the category of
$G$-torsors on $S$ to that of $f^{\ast}G$-torsors on $X$ is an
equivalence.
\end{enonce*}

\begin{itemize}
\item[a.] The problem in question is local on $S$ when the finite flat
  morphisms $u:S'\rightarrow S$ are taken as localization morphisms.
  This makes it possible to assume that $f$ admits a section
$x$\nde{The existence of a section after such a base change may be
deduced from \hyperref[XVIII.lem2.3]{2.3} and
\hyperref[XVIII.lem2.4]{2.4}.}. The functor $f^{\ast}$ is then fully
  faithful, by \hyperref[XVIII.lem1.2.3]{1.2.3} and the existence of
  $x$, and it remains to show that every torsor trivialized along $x$
  is trivial.

\item[b.] Suppose that $S$ is the spectrum of a field. If $G$ is
  affine, the assertion follows from
  \hyperref[XVIII.lem1.2.8]{1.2.8}. If $X$ is smooth and $G$
  is an extension of an abelian variety $A$ by an affine group
  $G_{0}$, denote by $G^{n}$ the inverse image of ${}_{n}A$ in $G$.
  The group $H^{1}(X,A)$ is torsion (\textsc{Raynaud}
  \cite[XIII 2.6]{sga4.exp.XVIII.bib.5}), and is therefore the union of the images of the
  $H^{1}(X,{}_{n}A)$. The group $H^{1}(X,G)$ is consequently the union
  of the images of the $H^{1}(X,G^{n})$. Since the groups $G^{n}$ are
  affine, the conclusion again follows from
  \hyperref[XVIII.lem1.2.8]{1.2.8}.

\item[c.] Let us prove the general case (assuming the existence of a
  section $x$) by induction on the length of $S$. Let
  $i:S_{0}\hookrightarrow S$ be a closed subscheme defined by an ideal
  $I$ of square zero and length 1.
\end{itemize}

By the induction hypothesis, the $G$-torsors $P$ on $X$ for which
$x^{\ast}P$ is trivial are deformations of the trivial $G$-torsor
$P_{0}$ on $X\times S_{0}$. If $G$ is smooth, and if $L$ is the $\Lie$
algebra of $G_{s}$, these deformations are classified by
$$
H^{1}(X_{s},f^{\ast}(L\otimes I))=0{\quoi},
$$
which\pageoriginale proves \hyperref[XVIII.lem1.2.9]{1.2.9} in this
case.

For $G$ flat over $S$, this argument generalizes in terms of the
\emph{Lie complex} $L$ of $G_{s}$.
% SGA-SRC-001128 — reconstructed by mathematical inference; scan inconclusive
If $\mathrm{L}\Lambda(G_{s}/k(s))$ denotes the relative cotangent complex
of $G_{s}$, this complex $L\in D^{b}(k(s))$ is the dual of
$\mathrm{L}e^{\ast}(\mathrm{L}\Lambda(G_{s}/k(s)))$. We have
$H^{i}(L)=0$ for $i\neq 0$, $1$\sisi{}{(\textsc{Illusie}
\cite[VII~2.4.1]{sga4.exp.XVIII.bib.3})}.

The deformations over $S$ (\resp over $X$) of the trivial $G$-torsor
over $S_{0}$ (\resp over $X_{0}$) are classified by $H^{1}(L\otimes I)$
(\resp by $H^{1}(X_{s},f^{\ast}(L\otimes I))$ (\textsc{Illusie}
\cite[VII~2.4.4~(ii)]{sga4.exp.XVIII.bib.3}).
By hypothesis, we have $H^{1}(X_{s},ℋ^0(f^{\ast}(L\otimes
I)))=0$ (because $H^{1}(X_{s},\Oo)=0$) and $H^{1}(L\otimes I)\sim
H^0(X_{s}, ℋ^{1}(f^{\ast}L\otimes I))$). We conclude that
$$
H^{1}(L\otimes I)\xrightarrow{\sim} H^{1}(X_{s},f^{\ast}(L\otimes
I)){\quoi},
$$
which means that every deformation over $X$ of the trivial $G$-torsor
over $X_{0}$ is the inverse image of one and only one deformation over
$S$ of the trivial $G$-torsor over $S_{0}$. If this deformation is
trivial along a section, it is therefore trivial, which proves
\hyperref[XVIII.lem1.2.9]{1.2.9}.

\setcounter{subsubsection}{9}
\subsubsection{}\etiquette[1.2.10]{XVIII.subsubsec1.2.10}
Let us prove \hyperref[XVIII.thm1.2.2]{1.2.2} when $S$ is Artinian. Let
$T$ be the largest subtorus of $G$; since full faithfulness is already
known (\hyperref[XVIII.lem1.2.4]{1.2.4}), the problem is local on $S$
for the étale topology: we may assume that $T$ is diagonalizable,
$T\sim \GG^{n}_{m}$. By \hyperref[XVIII.lem1.2.3]{1.2.3} and
\hyperref[XVIII.lem1.2.9]{1.2.9}, the exact cohomology sequence gives
$$
R^{1}f_{\ast}T\xrightarrow{\sim} R^{1}f_{\ast}G{\quoi},
$$
and the assertion follows, since it is trivial for $G=T$ and
$\Hom(\GG_{m},T)\xrightarrow{\sim}\Hom(\GG_{m},G)$ (\cf
\hyperref[XVIII.par1.2.2.1]{1.2.2.1}).

\subsubsection{}\etiquette[1.2.11]{XVIII.subsubsec1.2.11}
Let us prove \hyperref[XVIII.thm1.2.2]{1.2.2} when $S$ is complete
Noetherian local. Let $s$ be a section of~$p$\nde{As above, one must
also choose an isomorphism $s^\ast\Oo(1)\sim \Oo$.}.

\noindent
If $K$ is a $G$-torsor on $\PP(\epsilon)$, by
\hyperref[XVIII.subsubsec1.2.10]{1.2.10} there exists a unique pair
$(\widehat{\varphi},\widehat{\Psi})$ consisting of a formal
homomorphism $\widehat{\varphi}:\GG_{m}\rightarrow G$ and a formal
isomorphism\pageoriginale
$\widehat{\Psi}:\widehat{p^{\ast}s^{\ast}K}+
\widehat{\varphi}\widehat{\Oo(1)}\xrightarrow{\sim}
\widehat{K}$\nde{One must require this isomorphism to become the
obvious isomorphism after applying $s^\ast$.}\nde{What is meant here is
that, if $A$ is the ring of $S$, $I$ its maximal ideal, and
$S_n=\Spec(A/I^n)$, a formal morphism consists of a compatible system
of morphisms between objects over $S_n$ for every $n$. The question of
algebraicity is whether this compatible system is obtained by inverse
image from similar data over $S$.}. If $G$ is affine, then
$\widehat{\varphi}$ is algebraizable (\SGA 3 IX 7.1). Let us
temporarily admit that, as will be proved in
(\hyperref[XVIII.subsubsec1.2.13]{1.2.13}), $\widehat{\varphi}$ is
always algebraizable as $\varphi:\GG_{m}\rightarrow G$.

If $K_{0}=s^{\ast}K$, and if
$K_{1}=p^{\ast}K_{0}+\varphi\Oo(1)$, we have a formal isomorphism
$\widehat{\Psi}:\widehat{K_{1}}=\widehat{K(K_{0},\varphi)}\xrightarrow{\sim}
\widehat{K}$, \ie a formal trivialization of $K-K_{1}$. By formal
GAGA, this section is algebraic, which proves
\hyperref[XVIII.subsubsec1.2.11]{1.2.11}. More precisely, if $G$ is
quasi-affine, then $K-K_{1}$ is representable and one uses \EGA III
5.1.4. In the general case, one proceeds in the same way with the
algebraic space $K-K_{1}$.

\subsubsection{Algebraicity of $\widehat{\varphi}$: preliminary
  verification over a trait}\etiquette[1.2.12]{XVIII.subsubsec1.2.12}

Suppose that $S$ is a complete trait, $S=\Spec(V)$, and let $L$ be the
fraction field of~$V$.
\begin{itemize}
\item[a)] Over $L$, there exists a unique triple
  $(K_{0L},\varphi_{L},\Psi_{L})$:
$$
\Psi_{L}:K(K_{0L},\varphi_{L})\xrightarrow{\sim} K|\Spec (L){\quoi}.
$$

\item[b)] By (\hyperref[XVIII.subsubsec1.2.10]{1.2.10})
  (\hyperref[XVIII.subsubsec1.2.11]{1.2.11}), and with the notation of
  \hyperref[XVIII.subsubsec1.2.7]{1.2.7}, we have a rigid analytic map
  $\widehat{\varphi_{L}}:U_{m}\rightarrow \widehat{G}$, a
  $\widehat{G}$-torsor $\widehat{K_{0}}$, and a rigid analytic
  isomorphism
$$
\widehat{\Psi}:p^{\ast}\widehat{K_0}+\widehat{\varphi_{L}}
\widehat{\Oo(1)}\xrightarrow{\sim}
\widehat{K}{\quoi}.
$$
By the uniqueness assertion \hyperref[XVIII.subsubsec1.2.7]{1.2.7},
$\widehat{\varphi_{L}}$ is identified with the map deduced from
$\varphi_{L}$. Thus the formal morphism $\widehat{\Psi}$ of
\hyperref[XVIII.subsubsec1.2.11]{1.2.11} induces on $𝛍_{n}$ a map that
coincides with the map deduced from the algebraic map
$\varphi_{L}:\GG_{m}\rightarrow G$.
\end{itemize}

\subsubsection{Algebraicity of $\widehat{\varphi}$}\etiquette[1.2.13]{XVIII.subsubsec1.2.13}

The formal map $\widehat{\varphi}:\GG_{m}\rightarrow G$ of
\hyperref[XVIII.subsubsec1.2.11]{1.2.11} induces, for $\ell$ invertible
on $S$, maps
$$
\varphi_{\ell}:T_{\ell}(\widehat{\GG_{m}})=T_{\ell}(\GG_{m})\longrightarrow
T_{\ell}(G){\quoi}.
$$
By \hyperref[XVIII.subsubsec1.2.12]{(1.2.12)} and (\EGA II 7.1.9), for every point $t$ of $S$, there
exists an\pageoriginale algebraic morphism
$\varphi_{t}:\GG_{m_{t}}\rightarrow G_{t}$ such that
$T_{\ell}(\varphi_{t})=T_{\ell}(\widehat{\varphi})_{t}$. By (\SGA 3 XV 4.1
  bis), $\widehat{\varphi}$ is therefore algebraizable.

\subsubsection{}\etiquette[1.2.14]{XVIII.subsubsec1.2.14}
Let us prove \hyperref[XVIII.thm1.2.2]{1.2.2}. By passage to the limit,
we reduce to the case where $S$ is Noetherian, and then where $S$ is
Noetherian local with affine coordinate ring $A$. Let $s$ be a section
of $p$. Let $K$ be a $G$-torsor on $\PP(\epsilon)$. Over
$S'=\Spec(\widehat{A})$, there exists one and only one homomorphism
$\varphi':\GG_{mS'}\rightarrow G_{S'}$ such that the inverse image
$K'$ of $K$ on $\fprod{\PP(\epsilon)}{S'}{S}$ is of the form
$K(K'_{0},\varphi')$. This homomorphism $\varphi'$ descends to
$\varphi:\GG_{m}\rightarrow G$. To prove the essential surjectivity
of the functor \hyperref[XVIII.subsubsec1.2.1]{1.2.1}, we may replace
$K$ by $K-\varphi\Oo(1)$, \ie assume that $\varphi=0$. Over $S'$,
there then exists one and only one isomorphism
$\Psi:{p'}^{\ast}s^{\ast}K'\xrightarrow{\sim}K'$, inducing the
identity on the section $s$. This morphism descends to $S$ (fpqc
descent of morphisms of fppf torsors), and this proves
\hyperref[XVIII.thm1.2.2]{1.2.2}.

\begin{enonce*}{Variants 1.2.15}\etiquette[1.2.15]{XVIII.var1.2.15}
The conclusion of \hyperref[XVIII.thm1.2.2]{1.2.2} still holds in the
following cases.
\begin{enumerate}
\renewcommand{\theenumi}{\roman{enumi}}
\renewcommand{\labelenumi}{(\theenumi)}
\item $G$ is defined by a quasi-coherent module on $S$.

\item $G$ is the inverse image of a torsion sheaf on the small
  étale site of $S$.
\end{enumerate}
\end{enonce*}

In case (i), it is standard that the arrows
(\hyperref[XVIII.par1.2.2.1]{1.2.2.1}) are isomorphisms, taking into
account the fact that the fppf cohomology of $G$ coincides with its
cohomology for the Zariski topology.

In case (ii), the arrows (\hyperref[XVIII.par1.2.2.1]{1.2.2.1}) are
isomorphisms by the base change theorem for a proper morphism and the
simple connectedness of $\PP^{r}_{k}$ for $k$ algebraically closed.

\subsection{Integration of torsors}\etiquette[1.3]{XVIII.subsec1.3}

\subsubsection{}\etiquette[1.3.1]{XVIII.subsubsec1.3.1}
Let $S$ be a scheme and $G$ an abelian sheaf on the big fppf site of
$S$, satisfying the following three conditions.

\paragraph{}\etiquette[1.3.1.1]{XVIII.par1.3.1.1}
\begin{itemize}
\item[(i)] \emph{For\pageoriginale every projective space
$p:\PP^r_{S}\rightarrow S$, with $r\geq 1$, one has, on the big}
fppf \emph{site of} $S$ (\hyperrefext[XVII][XVII.subsec0.10]{0.10})
\begin{gather*}
G\xrightarrow{\sim} p_{\ast}p^{\ast}G\quad\mbox{and}\\
\Hhom(\GG_{m},G)\xrightarrow{\sim} R^{1}p_{\ast}p^{\ast}G
\end{gather*}

\item[(ii)] \emph{For every $n$ and every smooth affine $S$-morphism
  $f:X\rightarrow T$, if $K_{1}$ and $K_{2}$ are two $G$-torsors on
  $\Sym^{n}_{T}(X)$, the set $\Hom(K_{1},K_{2})$ is identified with
  the set of symmetric morphisms from the inverse image of $K_{1}$ on
  $(X/T)^{n}$ to the inverse image of $K_{2}$ on $(X/T)^{n}$.}

This condition is equivalent to the conjunction of:

\item[${\rm (ii')}$]
$\Hom_{S}(\Sym^{n}_{T}(X),G)\xrightarrow{\sim}\Hom_{S}((X/T)^{n},
G)^{\sigma_n}$, and

\item[${\rm (ii'')}$] If the inverse image on $(X/T)^{n}$ of a
torsor $K$ on $\Sym^{n}_{T}(X)$ admits a symmetric section, then $K$
is trivial.

\item[(iii)] $G$ \emph{satisfies either condition}
(\hyperrefext[XVII][XVII.par6.3.21.1]{6.3.21.1}), \emph{or conditions}
(\ref{XVII.par6.3.24.1}) \emph{and} (\ref{XVII.par6.3.24.2}).
\end{itemize}

The conclusion of \hyperref[XVIII.thm1.2.2]{1.2.2} therefore holds
(\cf \hyperref[XVIII.par1.2.2.1]{1.2.2.1}), as does the formula of
\hyperrefext[XVII][XVII.subsubsec6.3.26]{6.3.26}.

Condition (i) is discussed in \hyperref[XVIII.thm1.2.2]{1.2.2} and
\hyperref[XVIII.var1.2.15]{1.2.15}. Condition (ii) holds whenever $G$
is representable and of finite presentation, or is the inverse image
of a sheaf on the small étale site of $S$, or is flat
quasi-coherent. Condition \hyperref[XVIII.par1.3.1.1]{1.3.1.1}
therefore holds in each of the following cases.

\paragraph{}\etiquette[1.3.1.2]{XVIII.par1.3.1.2}
\begin{enumerate}
\renewcommand{\theenumi}{\alph{enumi}}
\renewcommand{\labelenumi}{(\theenumi)}
\item $G$ is smooth and of finite presentation
  (\hyperref[XVIII.thm1.2.2]{1.2.2} and
  \hyperrefext[XVII][XVII.par6.3.21.1]{6.3.21.1}).

\item $G$ is the inverse image of a torsion sheaf on the small
  étale site of $S$ (\hyperref[XVIII.var1.2.15]{1.2.15} (ii) and
  \hyperrefext[XVII][XVII.par6.3.21.1]{6.3.21.1}).

\item $G$ is affine, and the kernel of an epimorphism of smooth groups
  (\hyperref[XVIII.thm1.2.2]{1.2.2} and XVII \hyperref[XVII.eq6.3.3.1]{6.3.3.1}).

\item $G$ is defined by a flat quasi-coherent sheaf on $S$
  (\hyperref[XVIII.var1.2.15]{1.2.15} (i) and
  \hyperrefext[XVII][XVII.par6.3.21.1]{6.3.21.1}).
\end{enumerate}

\subsubsection{}\etiquette[1.3.2]{XVIII.subsubsec1.3.2}
Let\pageoriginale $f:X\rightarrow S$ be a smooth projective curve over
a scheme $S$, with nonempty geometrically connected fibers. Let $\Ll$
be an invertible sheaf on $X$ that is ``sufficiently'' relatively
ample in the following sense.

\paragraph{}\etiquette[1.3.2.1]{XVIII.par1.3.2.1}
For every geometric point $s$ of $S$, if $X_{s}$ has genus $g$, then
\begin{gather*}
\deg_{X_{s}}(\Ll)>2g-2,\\
\deg_{X_{s}}(\Ll)\geq 1\quad\mbox{\rm if}\quad g=0,\quaet
\deg_{X_{s}}(\Ll)\geq 2\quad\mbox{\rm if}\quad g=1
\end{gather*}

We then have $R^{1}f_{\ast}\Ll=0$, and $f_{\ast}\Ll$ is locally free,
its formation commutes with every base change, and it has rank
$\geq 2$ everywhere. If a section $s$ of $f_{\ast}\Ll$ vanishes at
no point, the scheme $D_{s}$ of zeros of $s$, regarded as a section of
$\Ll$ on $X$, is a relative divisor (\EGA IV 21.15.2), finite locally
free over $S$ since $X/S$ has relative dimension 1. We have canonically
$$
s:\Oo(D_{s})\xrightarrow{\sim}\Ll{\quoi}.
$$
The divisor $D_{s}$ depends on $s$ only up to multiplication by a
section of $\Oo^{\ast}_{S}$. The preceding construction therefore
gives us a relative divisor $Z_{\Ll}$ on the scheme obtained from $X$
by the base change
$$
p:\PP(f_{\ast}(\Ll)^{\vee})\longrightarrow S{\quoi}.
$$

Let $K$ be a torsor under $G_{X}$, and let $K|Z_{\Ll}$ be its inverse
image on $Z_{\Ll}$. The trace of this torsor
(\hyperrefext[XVII][XVII.subsubsec6.3.26]{6.3.26}), from $Z_{\Ll}$ to
$\PP(f_{\ast}(\Ll)^{\vee})$, is a torsor $K_{\Ll}$ on
$\PP(f_{\ast}(\Ll)^{\vee})$:
$$
K_{\Ll}=\Tr_{Z_{\Ll}/\PP}(K|Z_{\Ll}){\quoi}.
$$

By \hyperref[XVIII.thm1.2.2]{1.2.2}, there exists a triple, unique up
to unique isomorphism, consisting of a $G$-torsor
$\langle \Ll$, $K$] on $S$, a morphism
$\varphi_{\Ll,K}:\GG_{m}\rightarrow G$, and an isomorphism
$$
K_{\Ll}\simeq
p^{\ast}\langle \Ll,K]+\varphi_{\Ll,K}\Oo(1){\quoi}.
$$

\subsubsection{}\etiquette[1.3.3]{XVIII.subsubsec1.3.3}
If\pageoriginale the composition of torsors is written additively, we
have canonically
\begin{align*}
& \langle \Ll,K_{1}+K_{2}] \simeq
\langle \Ll,K_{1}]+\langle \Ll,K_{2}]\tag{1.3.3.1}\\
& \varphi_{\Ll, K_{1}+K_{2}}=\varphi_{\Ll,K_{1}}+\varphi_{\Ll,K_{2}}{\quoi}.\tag{1.3.3.2}
\end{align*}\etiquette[1.3.3.1]{XVIII.eq1.3.3.1}
\etiquette[1.3.3.2]{XVIII.eq1.3.3.2}

Let $\Ll_{1}$ and $\Ll_{2}$ be two invertible sheaves on $X$
satisfying (\hyperref[XVIII.par1.3.2.1]{1.3.2.1}), and let
$\Ll=\Ll_{1}\otimes \Ll_{2}$. The tensor product of sections defines
a morphism
$$
\xymatrix@R=1.5cm@C=.5cm{
 & **[r] \PP
  (f_{\ast}(\Ll_1)^{\vee}) \ar[ld]^{\pr_1}& \times_S \ar[dd]^{p_{12}}& **[l]\PP
  (f_{\ast}(\Ll_2)^{\vee})\ar[rr]^{\pi}\ar[rd]^{\pr_2} &     &
 \PP(f_{\ast}(\Ll)^{\vee}) \ar@{-}[dd]^{p}\\
**[r]\PP(f_*(\Ll_1)^{\vee})\ar[rrd]_{p_1}&          &                &                 &  **[l]\PP(f_{\ast}(\Ll_2)^{\vee}) \ar[lld]^{p_2}  &\\
 &          &  S     \ar@{=}[rrr]        &                 &      & S.
}
$$
We have $\pi^{\ast}Z_{\Ll}\simeq Z_{\Ll_{1}}+Z_{\Ll_{2}}$, whence a

% END INCLUDED SGA4_English_current_buildable_source_20260906/source/18/drafts/SGA4_XVIII_LINES_856_1250_English.texfrag

% BEGIN INCLUDED SGA4_English_current_buildable_source_20260906/source/18/drafts/SGA4_XVIII_LINES_1251_1650_English.texfrag
isomorphism (\hyperrefext[XVII][XVII.const6.3.27.1]{6.3.27.1})
\begin{equation*}
\pi^{\ast}K_{\Ll}=\pr^{\ast}_{1}K_{\Ll_{1}}+\pr^{\ast}_{2}K_{\Ll_{2}}{\quoi}.\tag{1.3.3.3}
\end{equation*}\etiquette[1.3.3.3]{XVIII.eq1.3.3.3}

There is a canonical isomorphism
$$
\pi^{\ast}\Oo(1)\simeq \pr^{\ast}_{1}\Oo(1)\otimes
\pr^{\ast}_{2}\Oo(1){\quoi},
$$
so that \eqhyperref[XVIII.eq1.3.3.3]{1.3.3.3} may be rewritten as
\begin{gather*}
\pi^{\ast}(p^{\ast}\langle \Ll,K]+\varphi_{\Ll,K}\Oo(1))=
p^{\ast}_{12}\langle \Ll,K]+\pr^{\ast}_{1}\varphi_{\Ll,K}\Oo(1)+\pr^{\ast}_{2}\varphi_{\Ll,K}\Oo(1) \\
= p^{\ast}_{12}\langle \Ll_{1},K]+p^{\ast}_{12}\langle \Ll_{2},K]+\pr^{\ast}_{1}\varphi_{\Ll_{1},K}\Oo(1)+\pr^{\ast}_{2}\varphi_{\Ll_{2},K}\Oo(1){\quoi}.
\end{gather*}

By applying \hyperref[XVIII.thm1.2.2]{1.2.2} twice, one obtains a canonical
isomorphism
\begin{gather*}
\langle \Ll_{1}\otimes \Ll_{2}, K]=\langle \Ll_{1},K]+\langle \Ll_{2},K]\tag{1.3.3.4}
\end{gather*}\etiquette[1.3.3.4]{XVIII.eq1.3.3.4}
and\pageoriginale an identity
$$
\varphi_{\Ll_{1}\otimes\Ll_{2},K}=\varphi_{\Ll_{1},K}=\varphi_{\Ll_{2},K}{\quoi}.
$$

Thus one sees that $\varphi_{\Ll,K}$ depends only on $K$, and one
sets
\begin{equation*}
\deg(K)=\varphi_{\Ll,K}{\quoi}.\tag{1.3.3.5}
\end{equation*}\etiquette[1.3.3.5]{XVIII.eq1.3.3.5}%
Consequently,
\begin{align*}
& K_{\Ll}\simeq
p^{\ast}\langle \Ll,K]+\deg(K)\Oo(1){\quoi}.\tag{1.3.3.6}\\
& \deg(K_{1}+K_{2})=\deg(K_{1})+\deg(K_{2}){\quoi}.\tag{1.3.3.7}
\end{align*}\etiquette[1.3.3.6]{XVIII.eq1.3.3.6}
\etiquette[1.3.3.7]{XVIII.eq1.3.3.7}


Let $\mathcal{U}$ be an invertible sheaf that is locally a direct
summand of $f_{\ast}\Ll$, defining a section $u$ of
$\PP((f_{\ast}\Ll)^{\vee})$. Let $D_{u}$ be the relative divisor
defined by the local invertible sections of $\mathcal{U}$.
$$
u^{\ast}\Oo(1)\simeq \mathcal{U}^{\vee}
$$
It therefore follows from \eqhyperref[XVIII.eq1.3.3.6]{1.3.3.6} that
\begin{equation*}
u^{\ast}K_{\Ll}=\Tr_{D_u/S}(K)\simeq
\langle \Ll,K]-\deg(K)\mathcal{U}{\quoi}\nde{The notation $\Tr_{D_u/S}(K)$
is introduced below in \eqhyperref[XVIII.eq1.3.5.0]{1.3.5.0}.}.\tag{1.3.3.8}
\end{equation*}\etiquette[1.3.3.8]{XVIII.eq1.3.3.8}

We leave it to the reader to verify the following compatibility:

\begin{enonce*}{Proposition 1.3.4}\etiquette[1.3.4]{XVIII.prop1.3.4}
The following diagram of the isomorphisms \eqhyperref[XVIII.eq1.3.3.1]{1.3.3.1} and \eqhyperref[XVIII.eq1.3.3.4]{1.3.3.4} is
commutative:
$$
\xymatrix@C=.1cm@R=.1cm{
\langle \Ll_{1}\otimes \Ll_{2}, K_{1}+K_{2}] & \simeq & \langle
\Ll_{1}\otimes \Ll_{2}, K_{1}]+\langle \Ll_{1} \otimes \Ll_{2}, K_{2}]\\
 & & \rotatebox{90}{$\simeq$} \\
\rotatebox{90}{$\simeq$} & & \langle \Ll_{1}, K_{1}]+\langle\Ll_{2}, K_{1}]+\langle
    \Ll_{1}, K_{2}]+\langle \Ll_{2}, K_{2}]\\
 & & \rotatebox{-50}{$\Arrowvert$}\\
\langle\Ll_{1}, K_{1}+K_{2}]+\langle\Ll_{2},K_{1}+K_{2}] & \simeq  &
    \langle\Ll_{1}, K_{1}]+\langle \Ll_{1},K_{2}]+\langle\Ll_{2},K_{1}]+\langle\Ll_{2},K_{2}]. 
}
$$%raghu

%$$
%\xymatrix{
%\langle \Ll_1\otimes \Ll_2, K_1 + K_2] \ar@{-}[dd]&\simeq& \langle \Ll_1\otimes \Ll_2, K_1] + \langle \Ll_! \otimes
%  \Ll_2, K_2]\ar@{-}[d]\\
%  & & \langle \Ll_1\otimes \Ll_2, K_1] + \langle \Ll_! \otimes
%  \Ll_2, K_2] \ar@{=}[d]\\
%\langle \Ll_1\otimes \Ll_2, K_1] + \langle \Ll_! \otimes
%  \Ll_2, K_2] & \simeq& \langle \Ll_1\otimes \Ll_2, K_1] + \langle \Ll_! \otimes
%  \Ll_2, K_2].
%}
%$$


\end{enonce*}

The isomorphisms \eqhyperref[XVIII.eq1.3.3.1]{1.3.3.1} and \eqhyperref[XVIII.eq1.3.3.4]{1.3.3.4} are moreover
compatible with the associativity and commutativity isomorphisms in a
sense whose precise formulation we leave to the reader. Formula
\eqhyperref[XVIII.eq1.3.3.4]{1.3.3.4} then makes it possible to extend by ``linearity'' the definition of
$\langle \Ll,K]$ when $\Ll$ no\pageoriginale longer necessarily
satisfies condition (\hyperref[XVIII.par1.3.2.1]{1.3.2.1}).

\setcounter{subsubsection}{4}
\subsubsection{}\etiquette[1.3.5]{XVIII.subsubsec1.3.5}
For every effective relative divisor $D$ on $X$, set
\begin{equation*}
\Tr_{D/S}(K)=\Tr_{D/S}(K|D){\quoi}.\tag{1.3.5.0}
\end{equation*}\etiquette[1.3.5.0]{XVIII.eq1.3.5.0}

If $\Oo(D)$ satisfies (\hyperref[XVIII.par1.3.2.1]{1.3.2.1}), the section
\hyperref[XVIII.sec1]{1} of $\Oo(D)$ is a
section of $f_{\ast}\Oo(D)$, and as such generates an invertible sheaf
that is locally a direct summand of $f_{\ast}\Oo(D)$, of which it
defines a trivialization. The isomorphism \eqhyperref[XVIII.eq1.3.3.8]{1.3.3.8} therefore
gives us an isomorphism
\begin{equation*}
\langle \Oo(D),K]\simeq \Tr_{D/S}(K){\quoi}.\tag{1.3.5.1}
\end{equation*}\etiquette[1.3.5.1]{XVIII.eq1.3.5.1}

One immediately verifies that the diagrams
$$\vcenter{
\xymatrix{
  \langle \Oo(D) , K_1 + K_2] \ar[r]^{\sim} \ar@{-}[d]_{\text{\rotatebox{90}{$\sim$}}}& \Tr_{D/S}(K_1 + K_2)\ar@{-}[d]^{\text{\rotatebox{90}{$\sim$}}}\\
\langle \Oo(D),K_1] + \langle \Oo(D) ,K_2] \ar[r]^{\sim}& \Tr_{D/S}(K_1) + \Tr_{D/S}(K_2)
}}\leqno{(1.3.5.2)}
$$\etiquette[1.3.5.2]{XVIII.eq1.3.5.2}%
and
$$\vcenter{
\xymatrix{
\langle \Oo(D_1 + D_2),K] \ar[r]^{\sim}
\ar@{-}[d]^{\text{\rotatebox{90}{$\sim$}}}& \Tr_{(D_1 +
  D_{2})/S}(K)\ar@{-}[d]^{\text{\rotatebox{90}{$\sim$}}\quad (XVII\
  \hyperref[XVII.const6.3.27.1]{6.3.27.1})}\\ 
\langle \Oo(D_1),K] + \langle \Oo(D_2),K] \ar[r]^{\sim}& \Tr_{D_1 /S}(K) +
    \Tr_{D_{2} /S} (K)
}}\leqno{(1.3.5.3)}
$$\etiquette[1.3.5.3]{XVIII.eq1.3.5.3}%
are commutative. By linearity, this makes it possible to define the
isomorphism \eqhyperref[XVIII.eq1.3.5.1]{1.3.5.1} for every relative divisor.

\subsubsection{}\etiquette[1.3.6]{XVIII.subsubsec1.3.6}
Let $D$ and $E$ be two effective relative Cartier divisors, and let
$f$ be a rational function such that
$$
\text{\rm div\,}(f)=D-E{\quoi},
$$
\ie an isomorphism between $\Oo(D)$ and $\Oo(E)$. The isomorphism $f$
between $\Oo(D)$ and $\Oo(E)$ defines, by \eqhyperref[XVIII.eq1.3.5.1]{1.3.5.1}, an isomorphism
\begin{equation*}
\langle f,K]:\Tr_{D/S}(K)\xrightarrow{\sim} \Tr_{E/S}(K){\quoi}.\tag{1.3.6.1}
\end{equation*}\etiquette[1.3.6.1]{XVIII.eq1.3.6.1}%
By\pageoriginale \eqhyperref[XVIII.eq1.3.5.2]{1.3.5.2} for the second formula, one has
\begin{align*}
& \langle fg,K]=\langle f,K].\langle g,K]\\
& \langle f,K_{1}+K_{2}]=\langle f,K_{1}]+\langle f,K_{2}]{\quoi}.
\end{align*}

If $\lambda$ is a section of $\Oo^{\ast}_{S}$, then, for every
invertible sheaf $\Ll$ on $X$, multiplication by
$\lambda$ is an automorphism of $\Ll$. By transport of
structure, this automorphism induces on the two sides of \eqhyperref[XVIII.eq1.3.3.8]{1.3.3.8}
automorphisms which correspond to one another under the isomorphism
\eqhyperref[XVIII.eq1.3.3.8]{1.3.3.8}. It follows that
$$
\langle \lambda, K]=\deg (K)(\lambda){\quoi}\;;
$$
if $\deg(K)\neq 0$, the isomorphism $\langle f,K]$ between $\Tr_{D/S}(K)$ and
$\Tr_{E/S}(K)$ thus depends on the choice of $f$.

Let $\lambda\in \Gamma(S,\Oo^{\ast}_{S})$, $g\in \Gamma(S,G)$,
let $\Ll$ be an invertible sheaf on $X$, and let $K$ be a $G$-torsor on
$X$. If $\langle \lambda,g]$ is the automorphism of $\langle \Ll,K]$ obtained
by transport of structure from the automorphisms $\lambda$ of
$\Ll$ and $g$ of $K$, then $\langle \lambda,g]$, identified with a
section of $G$, is given by
$\langle \lambda,g]=\langle \lambda,K]+\langle \Ll,g]$; since
$\langle \Oo(D),g]=\Tr_{D/S}(g)$, the preceding formula gives
\begin{equation*}
\langle \lambda,g]=\deg(K)(\lambda)+\deg(\Ll).g{\quoi}.\tag{1.3.6.2}
\end{equation*}\etiquette[1.3.6.2]{XVIII.eq1.3.6.2}

\subsubsection{}\etiquette[1.3.7]{XVIII.subsubsec1.3.7}
Let $f:X\rightarrow S$ be a smooth projective curve over $S$, and let
$G$ be a group over $S$ satisfying \hyperref[XVIII.par1.3.1.1]{1.3.1.1}. If
$X\xrightarrow{f_{0}}T\xrightarrow{\pi}S$ is the Stein factorization of
$f$, then $\pi$ is an étale covering of $S$, and the
sheaf $f_{\ast}G_{m}$ is the torus
$\prod_{T/S}\GG_{m}$. Locally on $S$ (for the étale topology),
$X$ is the disjoint union of a family of curves $(X_{i})_{i\in I}$ over
$S$ with nonempty geometrically connected fibers. If $K$ is a torsor
under $G$ and $\Ll$ an invertible sheaf on $X$, one then sets
$$
\langle \Ll,K]=\sum_{i}\langle \Ll|X_{i},K|X_{i}]
$$
and defines the degree of
$K:f_{\ast}\GG_{m}=\GG^{I}_{m}\rightarrow G$ to have
coordinates $\deg(K|X_{i}):\GG_{m}\rightarrow G$. The local
constructions globalize and give\pageoriginale
functors and morphisms
\begin{align*}
& \langle \Ll,K]_{X/S}=\Tr_{T/S}(\langle \Ll,K]_{X/T})\tag{1.3.7.1}\\
& \deg(K):f_{\ast}\GG_{m}=\prod_{T/S}\GG_{m}\longrightarrow G:\deg_{X/S}(K)=\prod_{T/S}\deg_{X/T}(K){\quoi}.\tag{1.3.7.2}
\end{align*}\etiquette[1.3.7.1]{XVIII.eq1.3.7.1}
\etiquette[1.3.7.2]{XVIII.eq1.3.7.2}

For an automorphism $\lambda$ of $\Ll$, defined by an
element $\lambda_{0}$ of $\Gamma(T,\Oo^{\ast}_{T})$, one has
\begin{equation*}
\langle \lambda,K]=\deg(K)(\lambda){\quoi}.\tag{1.3.7.3}
\end{equation*}\etiquette[1.3.7.3]{XVIII.eq1.3.7.3}

The total degree of $K$ is defined as the composite
$$
\deg \tot(K):\GG_{m}\longrightarrow
f_{\ast}\GG_{m}\xrightarrow{\deg(K)}G{\quoi}.
$$
The preceding results give:

\begin{enonce*}[remark]{Formulary 1.3.8}\etiquette[1.3.8]{XVIII.for1.3.8}
For every smooth projective curve $f:X\rightarrow S$, there is a
functor $\langle~,~]$, whose formation is compatible with every base
change, which associates to an invertible sheaf $\Ll$ on $X$ and a
torsor $K$ under $f^{\ast}G$ a torsor $\langle \Ll,K]$ under
$G$ on $S$. This functor is endowed with the following additional
structures and satisfies the following conditions:
\end{enonce*}

\setcounter{subsubsection}{8}
\paragraph{}\etiquette[1.3.8.1]{XVIII.par1.3.8.1}
$\langle \Ll,K]$ is biadditive in $\Ll$ and $K$; the
biadditivity isomorphisms are compatible with the associativity and
commutativity isomorphisms and with base changes, satisfy the
compatibility (\hyperref[XVIII.prop1.3.4]{1.3.4}), and are compatible with the
isomorphisms (\hyperref[XVIII.par1.3.8.2]{1.3.8.2}) below (via
\hyperrefext[XVII][XVII.par6.3.25.3]{6.3.25.3} and \hyperrefext[XVII][XVII.const6.3.27.1]{6.3.27.1}).

\paragraph{}\etiquette[1.3.8.2]{XVIII.par1.3.8.2}
For every effective relative divisor $E$ on $X/S$, there is a
canonical isomorphism
$$
\langle \Oo(E),K]\simeq \Tr_{E/S}(K|E){\quoi}.
$$

\paragraph{}\etiquette[1.3.8.3]{XVIII.par1.3.8.3}
The degree \eqhyperref[XVIII.eq1.3.7.2]{1.3.7.2} satisfies, for $\lambda\in f_{\ast}\GG_{m}$
an automorphism of $\Ll$,
$$
\langle \lambda,K]=\deg(K)(\lambda){\quoi}.
$$
For $\lambda\in \GG_{m}(S)$, the total degree satisfies
$$
\langle \lambda,K]=\deg\tot(K)(\lambda){\quoi}.
$$
Likewise,\pageoriginale for $g\in G(S)$ one has
$\langle \Ll,g]=\deg(\Ll).g$.

One deduces a canonical isomorphism:
\begin{equation*}
\langle f^{\ast}\Ll,K]\simeq \deg\tot(K)(\Ll){\quoi}.\tag{1.3.8.4}
\end{equation*}\etiquette[1.3.8.4]{XVIII.eq1.3.8.4}

It follows immediately from the definitions that, for every
homomorphism $\varphi:G\rightarrow G'$, there is a canonical
isomorphism, compatible with the preceding data:
\begin{equation*}
\langle \Ll,\varphi(K)]\simeq \varphi(\langle \Ll,K]){\quoi}.\tag{1.3.8.5}
\end{equation*}\etiquette[1.3.8.5]{XVIII.eq1.3.8.5}

\setcounter{paragraph}{5}
\paragraph{}\etiquette[1.3.8.6]{XVIII.par1.3.8.6}
Let $u:X\rightarrow Y$ be a flat morphism between smooth projective
curves over $S$. There are then isomorphisms compatible with
biadditivity and base changes:
\begin{itemize}
\item[--] For $\Ll$ an invertible sheaf on $Y$ and $K$ a
$G_{Y}$-torsor on $Y$, one has
$$
\langle \Ll, u^{\ast}K]\xrightarrow{\sim}\langle N_{X/Y}(\Ll),K]{\quoi}.
$$

\item[--] For $\Ll$ an invertible sheaf on $X$ and $K$ a
$G_{X}$-torsor on $X$, one has
$$
\langle u^{\ast}\Ll,K]\xrightarrow{\sim}
\langle \Ll,\Tr_{X/Y}K]{\quoi}.
$$
\end{itemize}
Moreover, these isomorphisms satisfy the evident compatibility for a
composite $uv$ of morphisms of curves; for $\Ll=\Oo(E)$, they
identify with the isomorphisms (\hyperrefext[XVII][XVII.const6.3.27.2]{6.3.27.2})
\begin{align*}
\Tr_{E/S}(u^{\ast}K) &\xrightarrow{\sim} \Tr_{u_{\ast}E/S}(K)\\
\Tr_{u^{\ast}E/S}(K) &\xrightarrow{\sim} \Tr_{E/S}\Tr_{X/Y}K{\quoi}.
\end{align*}
It remains to prove \hyperref[XVIII.par1.3.8.6]{1.3.8.6}. One reduces to the case in which $X$ and $Y$
have nonempty geometrically connected fibers, and notes that the
preceding explicit formulas give, for every section $s$ of $\Ll$
that does not vanish as a section of the direct image of $\Ll$, an
isomorphism $\Psi_{s}$ of the desired kind, homogeneous in $s$, and
therefore independent of $s$ when $\Ll$ satisfies
(\hyperref[XVIII.par1.3.2.1]{1.3.2.1}). The general case follows.


\subsubsection{}\etiquette[1.3.9]{XVIII.subsubsec1.3.9}
Under the hypotheses of \hyperref[XVIII.subsubsec1.3.7]{1.3.7},
we saw in (\hyperref[XVIII.for1.3.8]{1.3.8}) that every $G$-torsor $K$
on $X$ defines a functor $\langle \ast,K]$, denoted below by
$F(\ast)$,
endowed\pageoriginale with the following additional data and satisfying
the following conditions:
\begin{enumerate}
\renewcommand{\theenumi}{\roman{enumi}}
\renewcommand{\labelenumi}{(\theenumi)}
\item $F$ is a morphism of fppf stacks over $S$
  (\hyperrefext[XVII][XVII.subsec0.10]{0.10} and \textsc{Giraud}
  \cite{sga4.exp.XVIII.bib.1})
\begin{itemize}
\item[--] whose source is the stack whose $T$-objects, for an
$S$-scheme $T$, are the invertible sheaves on $X_{T}$;

\item[--] whose target is the stack of $G$-torsors.
\end{itemize}

\item The functor $F$ is endowed with additivity isomorphisms
$F(\Ll_{1}\otimes \Ll_{2})\xrightarrow{\sim}F(\Ll_{1})+F(\Ll_{2})$,
which are functorial, compatible with the associativity and
commutativity isomorphisms, and compatible with base changes.
\end{enumerate}

\begin{enonce*}{Theorem 1.3.10}\etiquette[1.3.10]{XVIII.thm1.3.10}
Under the hypotheses of \hyperref[XVIII.subsubsec1.3.7]{1.3.7}, the functor
$\Phi:K\mapsto \langle \ast,K]$ is
an equivalence between the category of $G$-torsors on $X$ and the
category of morphisms of stacks considered in
(\hyperref[XVIII.subsubsec1.3.9]{1.3.9}).
\end{enonce*}

One reduces easily to the case in which $X$ has nonempty connected fibers.

Let $F$ satisfy \hyperref[XVIII.subsubsec1.3.9]{1.3.9} (i) and (ii). For every
$S$-scheme $T$ and every section $t$ of $X_{T}/T$, denote also by $t$ the relative
divisor $t(T)$. Then $F(\Oo(t))$ is a $G_{T}$-torsor on $T$. We
denote by $\Psi(F)$ the torsor thus obtained on $X$, for $t$ the
``universal point'' of $X$ with values in $X$, defined by the
identity map from $X$ to $X$. For every section $t$ as above, one has
\begin{equation*}
F(\Oo(t))\simeq t^{\ast}\Psi(F){\quoi}.\tag{1.3.10.1}
\end{equation*}\etiquette[1.3.10.1]{XVIII.eq1.3.10.1}

The isomorphism \hyperref[XVIII.par1.3.8.2]{1.3.8.2} gives an isomorphism of functors
$\Psi\circ\Phi \simeq \Id$. We construct an isomorphism
$\Phi\circ\Psi\simeq \Id$.

For every section $t$ of $X_{T}$ over $T$, the isomorphism \eqhyperref[XVIII.eq1.3.10.1]{1.3.10.1}
is an isomorphism
$$
F(\Oo(t))\xrightarrow[\alpha]{\sim}\langle \Oo(t),\Psi(F)]{\quoi}.
$$

By additivity, every family $t_{1},\ldots, t_{n}$ of sections of
$X_{T}$ over $T$ defines an isomorphism
$$
F(\Oo(\sum t_{i}))\xrightarrow[\alpha]{\sim}\langle \Oo(\sum
t_{i}),\Psi(F)]{\quoi},
$$
and by \hyperref[XVIII.subsubsec1.3.9]{1.3.9} (ii), this isomorphism is independent
of the order of the $t_{i}$.

Take\pageoriginale $T=(X/S)^{n}$ and for $t_{1}\ldots t_{n}$ the
universal $n^{\text{\rm tuple}}$ of sections. If $T'=\Sym^{n}_{S}(X)$,
there is then on $X_{T'}$ a divisor $\sum t_{i}$ whose inverse image is
the divisor $\sum t_{i}$ on $X_{T}$. On $T$,
$\alpha$ is a symmetric isomorphism
$$
\alpha:F(\Oo(\sum t_{i}))\xrightarrow{\sim} \langle \Oo(\sum
t_{i}),\Psi(F)]
$$
between torsors that are inverse images of torsors on $T'$. By
hypothesis (\hyperref[XVIII.par1.3.1.1]{1.3.1.1}) (ii), this isomorphism comes from a unique
isomorphism on $T'$:
$$
\alpha:F(\Oo(\sum t_{i}))\xrightarrow{\sim} \langle \Oo (\sum
t_{i}),\Psi (F)]\qquad (\text{\rm on } T'){\quoi}.
$$

By \hyperrefext[XVII][XVII.prop6.3.9]{6.3.9}, one thus has, for every relative divisor $D$ on
$X$, an isomorphism
\begin{equation*}
\alpha_{D}:F(\Oo(D))\xrightarrow{\sim}
\langle \Oo(D),\Psi(F)]{\quoi},\tag{1.3.10.2}
\end{equation*}\etiquette[1.3.10.2]{XVIII.eq1.3.10.2}%
additive in $D$, whose formation is compatible with every base change.

Let $\Ll$ be a sufficiently ample invertible sheaf on $X$ in the
sense of \hyperref[XVIII.par1.3.2.1]{1.3.2.1}. Every nowhere-zero section $s$ of
$f_{\ast}\Ll$ defines, with the notation of \hyperref[XVIII.subsubsec1.3.2]{1.3.2}, an
isomorphism
$$
s:\Oo(D_{s})\xrightarrow{\sim}\Ll{\quoi},
$$
and hence, by \eqhyperref[XVIII.eq1.3.10.2]{1.3.10.2}, an isomorphism
\begin{equation*}
\alpha_{s}:F(\Ll)\xrightarrow{\sim}\langle \Ll,\Psi(F)]{\quoi}.\tag{1.3.10.3}
\end{equation*}\etiquette[1.3.10.3]{XVIII.eq1.3.10.3}

% END INCLUDED SGA4_English_current_buildable_source_20260906/source/18/drafts/SGA4_XVIII_LINES_1251_1650_English.texfrag

% BEGIN INCLUDED SGA4_English_current_buildable_source_20260906/source/18/drafts/SGA4_XVIII_LINES_1651_2050_English.texfrag

Let us prove that this isomorphism is independent of $s$. Indeed:
\begin{enumerate}
\renewcommand{\theenumi}{\roman{enumi}}
\renewcommand{\labelenumi}{(\theenumi)}
\item The construction of $\alpha_{s}$ commutes with every base change.

\item There exists a homomorphism $\varphi:\GG_{m}\rightarrow G$ such that
$$
\alpha_{\lambda s}=\varphi(\lambda).\alpha_{s}
$$
(by the functoriality of both sides of \eqhyperref[XVIII.eq1.3.10.3]{1.3.10.3}).
\end{enumerate}

One necessarily has $\varphi=0$, since otherwise the torsor on
$\PP((f_{\ast}\Ll)^{\vee})$ obtained by pulling back
$F(\Ll)-\langle \Ll,\Psi(F)]$ would have degree $\neq 0$. The
$\alpha_{s}$ therefore define a morphism from
$\PP((f_{\ast}\Ll)^{\vee})$ to $G$, and this morphism is
constant by hypothesis\pageoriginale (\hyperref[XVIII.par1.3.1.1]{1.3.1.1}) (i).
We may therefore set $\alpha_{\Ll}=\alpha_{s}$.

The isomorphisms $\alpha_{\Ll}$ thus form an isomorphism
of functors between the restrictions of the functors $F$ and $\langle *,\Psi(F)]$
to sufficiently ample sheaves. This isomorphism is compatible with the
additivity isomorphisms, which allows it to be extended by
linearity to all invertible sheaves and completes the construction
of
$$
\alpha:\Phi\circ\Psi\simeq \Id{\quoi}.
$$

\begin{enonce*}[remark]{Example 1.3.11}\etiquette[1.3.11]{XVIII.exem1.3.11}
$G=\GG_{m}$. Under the hypotheses of \hyperref[XVIII.subsubsec1.3.7]{1.3.7}, and in the
special case where $G=\GG_{m}$, we shall write $\langle~,~\rangle $ rather than
$\langle~,~]$; if $\Ll$ and $\Mm$ are two invertible sheaves
on $X$, then $\langle \Ll$, $\Mm\rangle $ is an invertible sheaf
on $S$. Moreover, there are isomorphisms (\hyperref[XVIII.par1.3.8.2]{1.3.8.2})
$$
\langle \Oo(D),\Mm\rangle \simeq N_{D/S}(\Mm){\quoi}.
$$

If $D$ and $E$ are disjoint effective relative divisors, then
$\Oo(E)|D$ is canonically isomorphic to $\Oo$, whence an isomorphism
\begin{equation}
\sigma_{D,E}:\langle \Oo(D),\Oo(E)\rangle \simeq N_{D/S}(\Oo(E))\simeq
N_{D/S}(\Oo)\simeq \Oo{\quoi}.\tag{1.3.11.1}
\end{equation}\etiquette[1.3.11.1]{XVIII.eq1.3.11.1}
\end{enonce*}

\begin{enonce*}{Proposition 1.3.12}\etiquette[1.3.12]{XVIII.prop1.3.12}
Let $(D_{i})_{i=1,2}$ and $E$ be effective relative divisors, with
$D_{i}$ disjoint from $E$ $(i=1,2)$. Let $f$ be a rational function
satisfying $\text{\rm div\,}(f)=D_{1}-D_{2}$, \ie an isomorphism
$f:\Oo(D_{1})\xrightarrow{\sim}\Oo(D_{2})$. The following diagrams
of isomorphisms are commutative:
\begin{equation*}
\vcenter{
\xymatrix@C=2.5cm@R=1.25cm{
\langle \Oo(D_1), \Oo(E)\rangle   \ar[d]^{\langle f,\Oo(E)\rangle } \ar[r]^(.65){\sigma_{D_{1},E}} & \Oo\ar[d]^{N_{E/S}(f)}\\
\langle \Oo(D_2),\Oo(E)\rangle    \ar[r]^(.65){\sigma_{D_{2},E}} & \Oo
}}\tag{1.3.12.1}
\end{equation*}\etiquette[1.3.12.1]{XVIII.eq1.3.12.1}
\begin{equation*}
\vcenter{
\xymatrix@C=2.5cm@R=1.25cm{
\langle \Oo(E), \Oo(D_1)\rangle   \ar[d]^{\langle \Oo(E),f\rangle } \ar[r]^(.65){\sigma_{E,D_{1}}} & \Oo\ar[d]^{N_{E/S}(f)}\\
\langle \Oo(E),\Oo(D_2)\rangle    \ar[r]^(.65){\sigma_{E,D_{2}}} & \Oo
}}\tag{1.3.12.2}
\end{equation*}\etiquette[1.3.12.2]{XVIII.eq1.3.12.2}
\end{enonce*}

The\pageoriginale formula \eqhyperref[XVIII.eq1.3.12.2]{1.3.12.2} simply expresses the
commutativity of the diagram
$$
\xymatrix@C=2.5cm@R=1.25cm{
N_{E/S}(\Oo(D_1))  \ar[r]^(.55){\sim} \ar[d]^{N_{E/S}(f)} & \Oo \ar[d]^{N_{E/S}(f)} \\         
N_{E/S}(\Oo(D_2))  \ar[r]^(.55){\sim} & \Oo.
}
$$

For every $S$-scheme $T$ and every invertible sheaf $\Ll$ on
$X_{T}$, set $F(\Ll)=N_{E_{T}/T}(\Ll)$. The
hypotheses of \hyperref[XVIII.thm1.3.10]{1.3.10} are satisfied by this functor, so
there exists a unique invertible sheaf $\Mm$ equipped with
an isomorphism of functors compatible with biadditivity:
$$
N_{E_{T}/T}(\Ll)\xrightarrow{\sim}\langle \Ll,\Mm\rangle {\quoi}.
$$

Moreover, if $d:T\rightarrow X_{T}$ is a section of $X$ over $T$, there
is a canonical isomorphism
$$
d^{\ast}\Mm\xrightarrow{\sim} N_{E_{T}/T}(\Oo(d)){\quoi}.
$$
Taking for $d$ the universal section of $X$, $\Delta:X\rightarrow
X_{X}$, we obtain
$$
\Mm\xrightarrow{\sim} N_{E_{X}/X}(\Oo(\Delta)){\quoi}.
$$

Locally on $S$, $E$ is a sum of sections $E=\sum t_{i}$
(\hyperrefext[XVII][XVII.subsec6.3]{6.3}); for such an $E$, one has
$$
N_{E_{X}/X}(\Oo(\Delta))=\otimes
t^{\ast}_{i}\Oo(\Delta)=\sum\Oo(t_{i})=\Oo(E){\quoi},
$$
and this isomorphism, being independent of the order of the $t_{i}$, defines
a canonical isomorphism\nde{The argument seems to assume that, locally,
the $t_i$ are well defined up to permutation, which is not true,
as can be seen in the case where $S=\Spec k[\epsilon]$ is the
spectrum of the algebra of dual numbers, $X=\PP^1_S$, and $E$ is the
closed subscheme of $\AAA^1_S=\Spec \Oo_S[T]\subset X$ defined by
the equation $(T-\varepsilon)(T+\varepsilon)=T^2$. It seems instead necessary
to apply the ``cancelled descent'' condition
\hyperref[XVIII.par1.3.1.1]{1.3.1.1}~(ii) to obtain the result in the situation of the
universal divisor on $X$ after base change
by $\Sym^n_S X\to S$, before deducing it for the divisor $E$.}
\begin{equation*}
N_{E_{X}/X}(\Oo(\Delta))\sim \Oo(E){\quoi},\tag{1.3.12.3}
\end{equation*}\etiquette[1.3.12.3]{XVIII.eq1.3.12.3}%
so that
\begin{equation*}
N_{E_{T}/T}(\Ll)\xrightarrow{\sim}\langle \Ll,\Oo(E)\rangle{\quoi}.%
\tag{1.3.12.4}
\end{equation*}\etiquette[1.3.12.4]{XVIII.eq1.3.12.4}

The isomorphism \eqhyperref[XVIII.eq1.3.12.3]{1.3.12.3} is compatible with biadditivity and
base changes and is characterized, in addition, by the fact that for
every\pageoriginale section $d$ of $X_{T}$ over $T$, the isomorphism
$$
N_{E/T}(\Oo(d))\xrightarrow{\sim}\langle \Oo(d),\Oo(E)\rangle
\xrightarrow{\sim} d^{\ast}\Oo(E)
$$
is the pullback of the universal isomorphism \eqhyperref[XVIII.eq1.3.12.3]{1.3.12.3}. In
particular, for $d$ disjoint from $E$, the diagram
$$
\xymatrix{
N_{E/T}(\Oo(d))  \ar[r]^{\sim} \ar@{=}[d]& \langle \Oo(d), \Oo(E)\rangle 
\ar[r]^{\sim} & d^\ast\Oo(E) \ar@{=}[d]\\
\Oo              \ar@{=}[rr] &          & \Oo  
}
$$
is commutative. By additivity\nde{The same objection raised above applies
here as well. One can argue in the same way by considering ``the universal
divisor disjoint from $E$,'' defined after base change by
$\Sym^n_S (X-E)\to S$.},
it follows that, for $D$ disjoint from
$E$, the diagram
\begin{equation*}
\vcenter{
\xymatrix{
N_{E/T}(\Oo(D))  \ar[r]^{\sim} \ar@{=}[d]& \langle \Oo(D), \Oo(E)\rangle 
\ar[r]^{\sim} & N_{D/T}(\Oo(E)) \ar@{=}[d]\\
\Oo              \ar@{=}[rr] &          & \Oo  
}}\tag{1.3.12.5}
\end{equation*}\etiquette[1.3.12.5]{XVIII.eq1.3.12.5}%
is commutative. The arrow $\sigma_{E,D}$ is therefore the composite of
the inverse of \hyperref[XVIII.eq1.3.12.4]{1.3.12.4} and an evident arrow
$$
\langle \Oo(D),\Oo(E)\rangle \xleftarrow{\sim} N_{E/S}(\Oo(D))\longrightarrow
\Oo
$$
and \eqhyperref[XVIII.eq1.3.12.1]{1.3.12.1} follows at once, since
\hyperref[XVIII.eq1.3.12.4]{1.3.12.4} is functorial.

We deduce from \hyperref[XVIII.prop1.3.12]{1.3.12} (\cf \textsc{Serre}
\cite[Ch.~III, prop.~7, p.~46]{sga4.exp.XVIII.bib.6})

\begin{enonce*}{Corollary 1.3.13}\etiquette[1.3.13]{XVIII.cor1.3.13}
Let $(D_{i})_{i=1,2}$ and $(E_{i})_{i=1,2}$ be effective relative
divisors, with $D_{i}$ disjoint from $E_{j}$, let $f$ be a rational
function satisfying $\text{\rm div\,}(f)=D_{1}-D_{2}$, and let $g$ be a
rational function satisfying $\text{\rm div\,}(g)=E_{1}-E_{2}$. Then
$$
N_{E_{1}/S}(f).N_{E_{2}/S}(f)^{-1}=N_{D_{1}/S}(g).N_{D_{2}/S}(g)^{-1}
$$
or, briefly,
\begin{equation*}
f(\text{\rm div\,}(g))=g(\text{\rm div\,}(f)){\quoi}.\tag{1.3.13.1}
\end{equation*}\etiquette[1.3.13.1]{XVIII.eq1.3.13.1}
\end{enonce*}

This formula, in the form
$$
N_{D_{2}/S}(g)N_{E_{1}/S}(f)=N_{E_{2}/S}(f)N_{D_{1}/S}(g){\quoi},
$$
expresses, via \eqhyperref[XVIII.eq1.3.12.1]{1.3.12.1} and \eqhyperref[XVIII.eq1.3.12.2]{1.3.12.2}, the commutativity of the
diagram
$$
\xymatrix@C=2.25cm@R=1.5cm{
\langle \Oo(D_1), \Oo(E_1)\rangle   \ar[r]^{\langle \Oo(D_1), g\rangle } \ar[d]^{\langle f,\Oo(E_1)\rangle }&
\langle \Oo(D_1),\Oo(E_2)\rangle  \ar[d]^{\langle f,\Oo(E_2)\rangle }\\
\langle \Oo(D_2), \Oo(E_1)\rangle   \ar[r]^{\langle \Oo(D_2),g\rangle } & \langle \Oo(D_2), \Oo(E_2)\rangle .
}
$$

\begin{enonce*}{Corollary 1.3.14}\etiquette[1.3.14]{XVIII.cor1.3.14}
For\pageoriginale $G=\GG_{m}$, the degree (\hyperref[XVIII.subsubsec1.3.7]{1.3.7}) coincides
with the usual degree of invertible sheaves.
\end{enonce*}

It suffices to compare formulas \eqhyperref[XVIII.eq1.3.7.3]{1.3.7.3} and \eqhyperref[XVIII.eq1.3.12.1]{1.3.12.1} for $f$
the pullback of $f_0\in \Gamma(S,\Oo^{\ast}_{S})$.

\setcounter{subsubsection}{14}
\subsubsection{}\etiquette[1.3.15]{XVIII.subsubsec1.3.15}
For any sufficiently ample $\Ll$ and $\Mm$,
there exist locally on $S$ disjoint divisors $D$ and $E$ and
isomorphisms
$$
\alpha:\Ll\longrightarrow \Oo(D)\quaet
\beta:\Mm\longrightarrow \Oo(E){\quoi}.
$$
Denote by $\tau$ the isomorphism making the diagram
$$
\xymatrix@C=1.5cm@R=1.25cm{
\langle \Ll,\Mm\rangle  \ar[r]^{\langle \alpha,\beta\rangle }
\ar[d]^{\tau}_{\text{\rotatebox{90}{$\backsim$}}}& \langle \Oo(D), \Oo(E)\rangle  \ar[r]^(.7){\sigma_{D,E}}& \Oo\ar@{=}[d]\\
\langle \Mm,\Ll\rangle  \ar[r]^{\langle \beta,\alpha\rangle }& \langle \Oo(E), \Oo(D)\rangle  \ar[r]^(.7){\sigma_{E,D}}& \Oo.
}
$$
It follows from \hyperref[XVIII.prop1.3.12]{1.3.12} that this isomorphism is independent
of the particular choice of $D$, $E$, $\alpha$, and $\beta$; it therefore globalizes to
a canonical isomorphism
\begin{equation*}
\tau:\langle \Ll,\Mm\rangle \xrightarrow{\sim}
\langle \Mm,\Ll\rangle {\quoi}.\tag{1.3.15.1}
\end{equation*}\etiquette[1.3.15.1]{XVIII.eq1.3.15.1}

This isomorphism satisfies $\tau^{2}=1$; it is compatible with
biadditivity, since \eqhyperref[XVIII.eq1.3.11.1]{1.3.11.1} is, and extends by linearity to
arbitrary invertible sheaves.

It follows from \eqhyperref[XVIII.eq1.3.12.5]{1.3.12.5} that the morphism
\eqhyperref[XVIII.eq1.3.12.4]{1.3.12.4} identifies with the composite
$$
N_{E_{T}/T}(\Ll)\xleftarrow{\sim}
\langle \Oo(E),\Ll\rangle \xrightarrow{\sim\tau}\langle \Ll,\Oo(E)\rangle 
$$
(verify this for $\Ll=\Oo(D)$ with $D$ disjoint from $E$).

\begin{enonce*}{Formulary 1.3.16}\etiquette[1.3.16]{XVIII.for1.3.16}
For\pageoriginale every smooth projective curve $f:X\rightarrow
S$, there is a functor $\langle~,~\rangle $, whose construction commutes with every
base change, that associates with two invertible sheaves $\Ll$ and
$\Mm$ on $X$ an invertible sheaf
$\langle \Ll,\Mm\rangle $ on $S$. This functor is endowed with the
following additional structures and satisfies the following
conditions:
\end{enonce*}

\setcounter{subsubsection}{16}
\paragraph{}\etiquette[1.3.16.1]{XVIII.par1.3.16.1}
With the notation of \hyperref[XVIII.for1.3.8]{1.3.8}, for $G=\GG_{m}$, one has
$$
\langle \Ll,\Mm\rangle =\langle \Ll,\Mm]{\quoi},
$$
in particular, $\langle \Ll,\Mm\rangle $ is biadditive in
$\Ll$, $\Mm$ in the sense specified in \hyperref[XVIII.par1.3.8.1]{1.3.8.1}, and this
biadditivity is compatible with that of the right-hand side of a canonical
isomorphism (functorial in $\Mm$)
$$
\langle \Oo(E),\Mm\rangle \xrightarrow{\sim} N_{E/S}(\Mm){\quoi}.
$$

\paragraph{}\etiquette[1.3.16.2]{XVIII.par1.3.16.2}
There is a symmetry isomorphism
$$
\tau:\langle \Ll,\Mm\rangle \xrightarrow{\sim}\langle \Mm,\Ll\rangle {\quoi},
$$
functorial in $\Ll$ and $\Mm$, compatible with
base changes and the biadditivity isomorphisms, and satisfying
$\tau^{2}=1$. It makes the following diagram commutative, for disjoint
divisors $D$ and $E$:
$$
\xymatrix{
N_{D/S}(\Oo(E))\ar@{-}[d]^{\text{\rotatebox{90}{$\backsim$}}}& \langle \Oo(D), \Oo(E)\rangle \ar[l]\ar[r]^{\tau} &
\langle \Oo(E), \Oo(D)\rangle  \ar[r]&
N_{E/S}(\Oo(D))\ar@{-}[d]^{\text{\rotatebox{90}{$\backsim$}}}\\
\Oo \ar@{=}[rrr]&    & & \Oo.
}
$$

These formulas imply formulas \eqhyperref[XVIII.eq1.3.12.1]{1.3.12.1}, \eqhyperref[XVIII.eq1.3.12.2]{1.3.12.2},
\eqhyperref[XVIII.eq1.3.13.1]{1.3.13.1}, and

\paragraph{}\etiquette[1.3.16.3]{XVIII.par1.3.16.3}
For $a$, $b\in \Gamma(S,\Oo^{\ast}_{S})$, denoting by
$\langle a,b\rangle $ the automorphism of $\langle \Ll,\Mm\rangle $, identified with
a section of $\GG_{m}$, obtained by functoriality from the
automorphisms $a$ and $b$ of $\Ll$ and $\Mm$, one has
$$
\langle a,b\rangle =a^{\deg(\Mm)} b^{\deg(\Ll)}{\quoi}.
$$

\paragraph{}\etiquette[1.3.16.4]{XVIII.par1.3.16.4}
If $\Ll_0$ (\resp $\Mm_0$) is an
invertible module on $S$, there is a canonical isomorphism
\begin{align*}
\langle f^{\ast}\Ll_0,\Mm\rangle  &\xrightarrow{\sim}
\Ll_0^{\otimes\deg(\Mm)}\\
\langle \Ll,f^{\ast}\Mm_0\rangle  &\xrightarrow{\sim} \Mm_0^{\otimes\deg(\Ll)}\qquad){\quoi}.\tag*{(\resp} 
\end{align*}

\paragraph{}\etiquette[1.3.16.5]{XVIII.par1.3.16.5}
If\pageoriginale $u:X\rightarrow Y$ is a finite flat morphism of
smooth projective curves over $S$, then, for $\Ll$ an
invertible sheaf on $X$ and $\Mm$ an invertible sheaf on
$Y$, the isomorphisms \hyperref[XVIII.par1.3.8.6]{1.3.8.6}
\begin{align*}
\langle \Ll,u^{\ast}\Mm\rangle  &\sim
\langle N_{X/Y}\Ll,\Mm\rangle \\
\langle u^{\ast}\Mm,\Ll\rangle  &\sim \langle \Mm,N_{X/Y}\Ll\rangle 
\end{align*}
are deduced from one another by the symmetry
\eqhyperref[XVIII.par1.3.16.2]{1.3.16.2} (verify this for
$\Ll=\Oo(E)$ and $\Mm=\Oo(F)$ with $E$ and $F$
disjoint).

\paragraph{}\etiquette[1.3.16.6]{XVIII.par1.3.16.6}
For an invertible sheaf $\Ll$ on $X$, the symmetry arrow
$$
\tau:\langle \Ll,\Ll\rangle \longrightarrow
\langle \Ll,\Ll\rangle 
$$
is multiplication by $(-1)^{\deg(\Ll)}$.

\subsubsection{}\etiquette[1.3.17]{XVIII.subsubsec1.3.17}
It remains to prove \hyperref[XVIII.par1.3.16.6]{1.3.16.6}. We shall give two proofs, the
first using the theory of the determinant (\SGA 6
XI)\nde{That exposé was not written, as explained in
\cite{sga4.exp.XVIII.bib.0ter},
which also contains a warning concerning
the multiplicativity property of the
determinant in distinguished triangles: it requires care, without
which counterexamples arise.
A reference on this subject is \cite{sga4.exp.XVIII.bib.3bis}.}.

\sisi{}{Before proceeding to each of the proofs, one may observe
that if $\Ll_1$ and $\Ll_2$ are two invertible sheaves on $X$ and
$\Ll=\Ll_1\otimes\Ll_2$, then the automorphism $\tau_{\Ll}\colon \langle
\Ll,\Ll\rangle \longrightarrow \langle \Ll,\Ll\rangle$ (identified with an
invertible function on $S$) is the product of the automorphisms
$\tau_{\Ll_1}$ and $\tau_{\Ll_2}$. In the situation where $S$ is the spectrum
of a field (and $X$ is geometrically connected),
one may therefore suppose that $\Ll=\Oo(D)$ with $D$ a
nonzero effective divisor. Since $\tau^2=\Id$, after replacing $\Ll$ by a
sufficiently large odd power (or writing it as a quotient of two ample
bundles), one may also make $\Ll$ as ample as desired,
if necessary.

The arguments below will also require reducing to the case where
$S$ is the spectrum of a field. For this, it suffices to reduce to the
case where $S$ is reduced, and this can be done by showing that fppf-locally
on $S$, the datum $(X,\Ll)$ comes from a similar datum
over a reduced base. Before making this reduction, one may reduce to the
case where $X\to S$ is equipped with a section $s\colon S\to
X$ and where the geometric fibers of $X\to S$ are connected and have some
genus $g$. Using \cite[\S 5.2]{sga4.exp.XVIII.bib.4}, one can obtain a smooth
projective curve $X_0\to S_0$ of genus $g$, with $S_0$ smooth of finite type
over $\Spec(\ZZ)$, such that \emph{locally}, every relative curve as
above is obtained from it by base change. Using the section
$s_0\colon S_0\to X_0$, one can rigidify the Picard scheme
$S_1=\sheafPic_{X_0/S_0}\to S_0$ and thus obtain a ``universal'' line bundle $\Ll_1$ on $X_1=X_0\times_{S_0}S_1$. Locally, every
relative curve as above equipped with a line bundle comes from
$(X_1,\Ll_1)$ over $S_1$, which is indeed reduced\nde{The preceding two
paragraphs were added in this edition.}.}

\emph{1\textsuperscript{st} proof}: Let $u:T\rightarrow S$ be a
finite locally free morphism. If $\Ll$ is an invertible sheaf on
$T$, there is an isomorphism
\begin{equation*}
N_{T/S}(\Ll)\simeq \det(u_{\ast}\Ll)\det(u_{\ast}\Oo_{T})^{-1}\tag{1.3.17.1}
\end{equation*}\etiquette[1.3.17.1]{XVIII.eq1.3.17.1}%
uniquely determined by the following three conditions:
\begin{enumerate}
\renewcommand{\theenumi}{\alph{enumi}}
\renewcommand{\labelenumi}{\theenumi)}
\item its construction commutes with base changes;

\item for $\Ll=\Oo_{T}$, it is the identity morphism of
$\Oo_{S}$;

\item it is functorial in $\Ll$.
\end{enumerate}

Let $D$ and $E$ be two effective relative divisors on $X$, let $i$
be the inclusion of $D$ in $X$, and let $u$ be the projection of $D$ onto $S$. The
short exact sequence
$$
0\longrightarrow \Oo (-E)\longrightarrow \Oo\longrightarrow
\Oo_{E}\longrightarrow 0
$$
allows the isomorphism
$$
\langle \Oo(D),\Oo(-E)\rangle \simeq N_{D/S}(\Oo(-E))\simeq
\det(u_{\ast}i^{\ast}\Oo(-E))\det (u_{\ast}i^{\ast}\Oo)^{-1}

% END INCLUDED SGA4_English_current_buildable_source_20260906/source/18/drafts/SGA4_XVIII_LINES_1651_2050_English.texfrag

% BEGIN INCLUDED SGA4_English_current_buildable_source_20260906/source/18/drafts/SGA4_XVIII_LINES_2051_2450_English.texfrag
$$
as\pageoriginale an isomorphism
$$
\langle \Oo(D),\Oo(-E)\rangle =\det Ru_{\ast}(Li^{\ast}\Oo_{E})^{-1}\simeq
\det Rf_{\ast}(\Oo_{D}\mathop{\otimes}^{\LL}
\Oo_{E})^{-1}{\quoi}.
$$
or
\begin{equation*}
\langle \Oo(D),\Oo(E)\rangle \simeq \det Rf_{\ast}(\Oo_{D}\mathop{\otimes}^{\LL}\Oo_{E}){\quoi}.\tag{1.3.17.2}
\end{equation*}\etiquette[1.3.17.2]{XVIII.eq1.3.17.2}

For disjoint $D$ and $E$, one has
$\Oo_{D}\displaystyle{\mathop{\otimes}^{\LL}}\Oo_{E}=0$, and
the resulting isomorphism of $\langle \Oo(D)$, $\Oo(E)\rangle $ with
$\Oo_{S}$ is the one already considered. It follows that, for disjoint
$D$ and $E$, the symmetry $\tau$ and the symmetry isomorphism
$\Oo_{D}\displaystyle{\mathop{\otimes}^{\LL}}\Oo_{E}\xrightarrow{\sim}\Oo_{E}\displaystyle{\mathop{\otimes}^{\LL}}\Oo_{D}$
are compatible via \eqhyperref[XVIII.eq1.3.17.2]{1.3.17.2}. By specialization\nde{It
seems that here one must first reduce to the case where $S$ is the
spectrum of a field, as explained above.
It is then possible to make $D$ and
$E$ occur in two families of divisors parametrized by an integral
scheme in such a way that the divisors in the two families are disjoint
over the generic point. This may be obtained, for example, by considering
the universal situation over which the condition that $D$ and $E$ be
disjoint holds on a dense open subset of $\Sym^n_S
X\times \Sym^m_S X$, which is integral.},
this remains true for arbitrary $D$ and $E$.

When $D=E$ is defined by an ideal $I$, one has
\begin{alignat*}{2}
&
ℋ^{0}(\Oo_{D}\mathop{\otimes}^{\LL}\Oo_{D})=\Oo_{D}
&\qquad& \text{\rm symmetry } = +\sisi{1}{\Id}\\
& ℋ^{-1}(\Oo_{D}\mathop{\otimes}^{\LL}\Oo_{D})=I/I^{2} &&
\text{\rm symmetry } =-\sisi{1}{\Id}\\
& ℋ^{i}(\Oo_{D}\mathop{\otimes}^{\LL}\Oo_{D})=0 &&
\text{\rm for } i\neq 0, -1{\quoi}.
\end{alignat*}
Consequently, the map
$$
\tau:\langle \Oo(D), \Oo(D)\rangle \xrightarrow{\sim}\langle \Oo(D), \Oo(D)\rangle 
$$
is $(-1)^{\rk I/I^{2}}=(-1)^{\deg(\Oo(D))}$, which proves \hyperref[XVIII.par1.3.16.6]{1.3.16.6}.

\smallskip
\noindent
\emph{2\textsuperscript{nd} proof}: By reduction to a “universal”
case, we may suppose that $S$ is reduced. This case reduces to that in which $S$
is the spectrum of an algebraically closed field $k$. In this case, let
$X_{1}$, $X_{2}$, $Y_{1}$, and $Y_{2}$ be four divisors on $X$, let $f$ and
$g$ be two rational functions, and suppose that
\begin{align*}
&\text{\rm div\,}(f)=X_{2}-X_{1}\\
&\text{\rm div\,}(g)=Y_{2}-Y_{1}\\
&(X_{1}\cup Y_{2})\cap (X_{2}\cup Y_{1})=∅{\quoi}.
\end{align*}
The maps $f$ and $g$ then define\nde{Compared with the original
edition, numerous “signs” have been modified here and there through
the end of this paragraph.}
$$
\langle f,g\rangle :\langle \Oo(X_{2}), \Oo(Y_{2})\rangle \longrightarrow
\langle \Oo(X_{1}),\Oo(Y_{1})\rangle {\quoi}.
$$
The\pageoriginale two sides are canonically isomorphic to $k$
\eqhyperref[XVIII.eq1.3.11.1]{1.3.11.1}, so that $\langle f,g\rangle $ is identified with an
element of $k^{\ast}$.

\hypertarget{sga4.exp.XVIII.lemma.1.3.18}{}\begin{enonce*}{Lemma 1.3.18}\etiquette[1.3.18]{XVIII.lem1.3.18}
One has
$$
\langle f,g\rangle =\prod_{x\in X_{1}\cup Y_{2}}(-1)^{v_{x}(f)v_{x}(g)}
(f^{v_{x}(g)}/g^{v_{x}(f)})(x){\quoi}.
$$
\end{enonce*}

Let us show that (\hyperref[XVIII.lem1.3.18]{1.3.18}) implies
(\hyperref[XVIII.par1.3.16.6]{1.3.16.6}). We may suppose that
$\Ll\sim \Oo(X)\sim \Oo(Y)$ with $X$ and $Y$ effective, and
that $X\cap Y=∅$\nde{We saw above that we could suppose that
$\Ll$ was very ample.}. Let $f$ be a rational function such that
$\text{\rm div\,}(f)=Y-X$\nde{Care should be taken not to confuse $X$ with the
ambient curve, which is also denoted by $X$.}. The commutative diagram
$$
\xymatrix@C=2.5cm@R=1.25cm{
\langle \Ll, \Ll\rangle  \ar[r]^{\tau} \ar@{=}[d]&
\langle \Ll,\Ll\rangle  \ar@{=}[d]\\
\langle \Oo(X), \Oo(Y)\rangle   \ar[r]^{\langle f^{-1},f\rangle } & \langle
\Oo(Y), \Oo(X)\rangle 
}
$$
shows that
$$
\tau=\langle f^{-1},f\rangle =\prod_{x\in
Y}(-1)^{-v_{x}(f)^{2}}=(-1)^{\deg(\Ll)}{\quoi}.
$$ Let us prove \hyperref[XVIII.lem1.3.18]{1.3.18}. We reduce to supposing that $X_{i}$ and
$Y_{i}$ are very ample. There are then factorizations $f=f_{1}f_{2}$;
$g=g_{1}g_{2}$ with
\begin{alignat*}{2}
&\text{\rm div\,}(f_{1}) = X'-X_{1} &\qquad& \text{\rm div\,}(g_{1})
= Y'-Y_{1}\\
& \text{\rm div\,}(f_{2}) = X_{2}-X' &\qquad& \text{\rm div\,}(g_{2})
= Y_{2}-Y'{\quoi},
\end{alignat*}
the divisors $X'$, $Y'$, and $X_{1}\cup X_{2}\cup Y_{1}\cup Y_{2}$
being pairwise disjoint.

Denote by $\langle~,~\rangle^{\ast}$ the right-hand side of
(\hyperref[XVIII.lem1.3.18]{1.3.18}), and let us prove that
$$ \langle f,g\rangle ^{\ast}=\langle f_{1},g_{1}\rangle ^{\ast}\cdot
\langle f_{2},g_{2}\rangle ^{\ast}{\quoi}.
$$
By \textsc{Serre} \cite[Ch.~III, prop.~6, p.~44]{sga4.exp.XVIII.bib.6}, one has
\begin{equation*}
\prod_{x}(-1)^{v_{x}(f_{1})v_{x}(g_{2})}
(f_{1}^{v_{x}(g_{2})}/g_{2}^{v_{x}(f_{1})}) (x)=1{\quoi}.\tag{1.3.18.1}
\end{equation*}\etiquette[1.3.18.1]{XVIII.eq1.3.18.1}%
It follows that
% SGA-SRC-001129 — exact continuation restored from the original scan
\begin{eqnarray*}
\langle f_{1},g_{1}\rangle ^{\ast}\cdot \langle f_{2},g_{2}\rangle
^{\ast} &= &g_{1}(X_{1})f_{1}(Y')g_{2}(X')f_{2}(Y_{2})\cdot\\
& & \prod_{X_{1}\cup X'\cup Y'\cup Y_{2}}(-1)^{v_{x}(f_{1})v_{x}(g_{2})}
(f_{1}^{v_{x}(g_{2})}\cdot g_{2}^{-v_{x}(f_{1})})(x)\\
&=&\prod_{x\in X'} g_{2}(x)^{-v_{x}(f_{2})}\cdot
g_{2}(x)^{-v_{x}(f_{1})}\cdot \prod_{x\in
Y'}f_{1}(x)^{v_{x}(g_{1})}\cdot f_{1}(x)^{v_{x}(g_{2})}\cdot\\
&&\prod_{x\in X_{1}\cup Y_{2}}(-1)^{v_{x}(f)v_{x}(g)}
((f_{1}f_{2})^{v_{x}(g_{2})}\cdot
(g_{1}g_{2})^{-v_{x}(f_{1})})(x)=\langle f_{1}f_{2},g_{1}g_{2}\rangle
^{\ast}{\quoi}. 
\end{eqnarray*}
One\pageoriginale has, by functoriality of $\langle~,~\rangle $,
$\langle f,g\rangle =\langle f_{1},g_{1}\rangle \langle
f_{2},g_{2}\rangle $, and it is readily checked that $\langle
f_{i},g_{i}\rangle =\langle f_{i},g_{i}\rangle ^{\ast}$\nde{This is the
special case in which one assumes in addition that $X_1\cap Y_2=∅$; it follows
from \hyperref[XVIII.eq1.3.12.1]{1.3.12.1} and \hyperref[XVIII.eq1.3.12.2]{1.3.12.2}.}. The lemma
follows.

\subsection{Strictly commutative Picard
  categories}\etiquette[1.4]{XVIII.subsec1.4}

Some results from “general topology” will be needed
to express the results of \hyperref[XVIII.subsec1.3]{1.3} by a “universal
coefficient formula.”

The results of this section are suggested in a letter from
A.~\textsc{Grothendieck} to J.L. \textsc{Verdier} (1966).

The notions of associativity and commutativity for bifunctors
(see below) were introduced by \textsc{Mac Lane}.

\subsubsection{}\etiquette[1.4.1]{XVIII.subsubsec1.4.1}
Let $C$ be a category, let $F:C\times C\rightarrow C$ be a functor, and let
$$
\sigma:F(F(X,Y),Z)\xrightarrow{\sim} F(X,F(Y,Z))
$$
be an isomorphism of trifunctors. The pair $(F,\sigma)$ is said to be
an \emph{associative functor}, or $\sigma$ is said to be an
\emph{associativity datum} on $F$, if the following condition holds:

\begin{itemize}
\item[(Ass)] For every family $(X_{i})_{i\in I}$ of objects of
$C$, denoting by $e:I\rightarrow M(I)$ the canonical map from
$I$ into the free monoid (without identity) generated by $I$, there exist a
map $ℱ:M(I)\rightarrow \Ob C$, isomorphisms
$a_{i}:ℱ(e_{i})\xrightarrow{\sim} X_{i}$, and isomorphisms
$a_{g,h}:ℱ(gh)\xrightarrow{\sim}F(ℱ(g),ℱ(h))$ such that the
diagrams
\begin{equation*}
\vcenter{
\xymatrix@C=1.25cm@R=1.25cm{
ℱ(f(gh))  \ar@{=}[d]\ar[r]^-{\sim}_-{a_{f,gh}} &
F(ℱ(f),ℱ(gh))
\ar[r]^-{\sim}_-{a_{g,h}}&
F(ℱ(f),F(ℱ(g),ℱ(h)))\\
ℱ((fg)h)  \ar[r]^-{\sim}_-{a_{fg,h}}& F(ℱ(fg),ℱ(h))
\ar[r]^-{\sim}_-{a_{f,g}}& F(F(ℱ(f),ℱ(g)),ℱ(h))\ar[u]_{\sigma}^{\text{\rotatebox{90}{$\backsim$}}}
}}\tag{1.4.1.1}
\end{equation*}\etiquette[1.4.1.1]{XVIII.eq1.4.1.1}%
are commutative.

We shall not need the fact that axiom (Ass) is equivalent to
the “pentagon axiom,” namely that the following diagram is
commutative:
$$
\xymatrix@C=-.5cm{
                                  & F(F(X,Y),F(Z,T))&  \\
F(X,F(Y,F(Z,T))) \ar@{<->}[ru]    &                 & F(F(F(X,Y),Z),T)\ar@{<->}[lu] \\
F(X,F(F(Y,Z),T)) \ar@{<->}[rr]\ar@{<->}[u]    &    & F(F(X,F(Y,Z)),T). \ar@{<->}[u] \\
}
$$

Now suppose\pageoriginale given functorial isomorphisms
\begin{align*}
\sigma_{X,Y,Z} &: F(F(X,Y),Z)\xrightarrow{\sim}F(X,F(Y,Z))\\
\tau_{X,Y} &: F(X,Y)\xrightarrow{\sim} F(Y,X){\quoi}.
\end{align*}
We shall say that $\sigma$ and $\tau$ make $F$ an
\emph{associative} and \emph{strictly commutative} functor if the
following condition holds.

\item[(Ass.s.Com)] For every family $(X_{i})_{i\in I}$
of objects of $C$, denoting by $e:I\rightarrow N(I)$ the canonical
map from $I$ into the free abelian monoid (without identity) generated
by $I$, there exist a map $ℱ:N(I)\rightarrow \Ob C$, isomorphisms
$a_{i}:ℱ(e_{i})\xrightarrow{\sim} X_{i}$, and isomorphisms
$a_{g,h}:ℱ(gh)\xrightarrow{\sim}
F(ℱ(g),ℱ(h))$
such that diagram \eqhyperref[XVIII.eq1.4.1.1]{1.4.1.1} is commutative, as is the
diagram
\begin{equation*}
\vcenter{
\xymatrix@C=2.25cm{
ℱ(gh) \ar@{=}[d] \ar[r]^(.4){a_{g,h}}& F(ℱ(g),ℱ(h))\ar[d]^{\tau}\\
ℱ(hg)    \ar[r]^(.4){a_{h,g}}& F(ℱ(h),ℱ(g)).        
}}\tag{1.4.1.2}
\end{equation*}\etiquette[1.4.1.2]{XVIII.eq1.4.1.2}
\end{itemize}

We shall not need the fact that this axiom is equivalent to the
conjunction of:
\begin{enumerate}
\renewcommand{\labelenumi}{\theenumi)}
\item the pentagon axiom;

\item $\tau_{X,X}:X+X\rightarrow X+X$ is the identity;

\item $\tau_{Y,X}\circ\tau_{X,Y}:X+Y\rightarrow Y+X\rightarrow X+Y$
is the identity;

\item the hexagon axiom: the diagram
$$
\xymatrix{
          &     X + (Y + Z)                 &\\
 (X+Y)+Z  \ar@{-}[ru]_{\sigma} \ar@{-}[d]^{\tau}& & X+(Z+Y)\ar@{-}[lu]^{\tau}\\
Z+(X+Y) \ar@{-}[rd]_{\sigma}  &  & (X+Z) + Y \ar@{-}[u]_{\sigma}\\
          &    (Z+X)+Y    \ar@{-}[ru]_{\tau}                  &
}
$$
is\pageoriginale commutative.
\end{enumerate}

Note that condition 2) is not very often satisfied
in practice (hence the terminology \emph{strictly} commutative).

\hypertarget{sga4.exp.XVIII.definition.1.4.2}{}\begin{enonce*}{Definition 1.4.2}\etiquette[1.4.2]{XVIII.def1.4.2}
A \emph{strictly commutative Picard category $\Pp$}
is a nonempty category all of whose morphisms are
isomorphisms, equipped with a functor $+:\Pp\times
\Pp\rightarrow \Pp:(X,Y)\mapsto X+Y$ and
functorial isomorphisms
\begin{align*}
\sigma &: (X+Y)+Z\xrightarrow{\sim} X+(Y+Z)\\
\tau &: X+Y\longrightarrow Y+X
\end{align*}
making $+$ an associative and strictly commutative functor, and
such that, moreover, for every $X\in \Ob \Pp$, the functor
$Y\mapsto X+Y$ is an equivalence of categories.
\end{enonce*}

The reader will verify:

\hypertarget{sga4.exp.XVIII.lemma.1.4.3}{}\begin{enonce*}{Lemma 1.4.3}\etiquette[1.4.3]{XVIII.lem1.4.3}
If $(X_{i})_{i\in I}$ is a family of objects of a strictly
commutative Picard category $\Pp$, there exist a map
$\Sigma:\ZZ^{(I)}\rightarrow \Ob \Pp$, isomorphisms
$a_{i}:\Sigma(e_{i})\xrightarrow{\sim} X_{i}$, and
$a_{\underline{n},\underline{m}}:\Sigma(\underline{n}+\underline{m})
\xrightarrow{\sim}\Sigma(\underline{n})+\Sigma(\underline{m})$
such that the diagrams of types \eqhyperref[XVIII.eq1.4.1.1]{1.4.1.1} and \eqhyperref[XVIII.eq1.4.1.2]{1.4.1.2}
are commutative. The system
$(\Sigma,(a_{i}),(a_{\underline{n},\underline{m}}))$ is unique up to unique
isomorphism, and is functorial in $(X_{i})_{i\in I}$.
\end{enonce*}

\setcounter{subsubsection}{3}
\subsubsection{}\etiquette[1.4.4]{XVIII.subsubsec1.4.4}
A strictly commutative Picard category admits, up to unique
isomorphism, exactly one \emph{neutral object}, which may here be defined
as a pair $(e,\varphi)$ consisting of an object $e$ and an
isomorphism $\varphi:e+e\xrightarrow{\sim}e$. If $(e,\varphi)$ is a
neutral object, there is a unique isomorphism of functors
$$
\alpha_{s}:e+X\xrightarrow{\sim} X
$$
making the diagram
$$
\xymatrix@C=2cm@R=1.25cm{
e + (e + X) \ar@{<->}[r] \ar[d]^{e+\alpha_s} & (e + e) + X \ar[d]^{\varphi}\\
e + X \ar@{=}[r]                       & e + X .
}
$$
commutative.

One\pageoriginale similarly defines $\alpha_{d}:X+e\xrightarrow{\sim}X$,
and $\varphi$ is a special case of both $\alpha_{s}$ and
$\alpha_{d}$; $\tau$ interchanges $\alpha_{s}$ and $\alpha_{d}$. The group
$\Aut(e)$ is abelian, and for every $X\in \Ob\Pp$, the
functors $X+$ and $+X$ induce the same isomorphism between $\Aut(e)$ and
$\Aut(X)$.

\hypertarget{sga4.exp.XVIII.definition.1.4.5}{}\begin{enonce*}{Definition 1.4.5}\etiquette[1.4.5]{XVIII.def1.4.5}
A \emph{strictly commutative Picard stack $\Pp$ on a}
site $\Ss$ is a stack in groupoids $\Pp$ on
$\Ss$ (\textsc{Giraud} \cite{sga4.exp.XVIII.bib.1}), equipped with a functor
$+:\fprod{\Pp}{\Pp}{\Ss}\rightarrow \Pp$ and
isomorphisms of functors
\begin{align*}
\tau_{x,y} &: x+y\longrightarrow y+x\\
\sigma_{x,y,z} &: (x+y)+z\longrightarrow x+(y+z)
\end{align*}
which, for every $U\in \Ob \Ss$, make $\Pp(U)$ a
strictly commutative Picard category.
\end{enonce*}

In what follows, we shall simply speak of a \emph{Picard stack},
without specifying “strictly commutative.” By
\hyperref[XVIII.subsubsec1.4.4]{1.4.4}, every Picard
stack $\Pp$ admits a global “neutral object”
$e$, and $\sheafAut(e)$ is an abelian sheaf.

\setcounter{subsubsection}{5}
\subsubsection{}\etiquette[1.4.6]{XVIII.subsubsec1.4.6}
An \emph{additive functor} $F:\Pp_{1}\rightarrow
\Pp_{2}$ between Picard stacks on $\Ss$ is an
$\Ss$-functor (necessarily Cartesian) equipped with an
isomorphism of functors
$$
F(x+y)\xrightarrow{\sim} F(x)+F(y)
$$
making the diagrams
$$
\xymatrix@C=2.25cm@R=1.5cm{
F(x+y)  \ar[r] \ar[d]^{F(\tau)}& F(x) + F(y) \ar[d]^{\tau}\\
F(y+x)         \ar[r] & F(y) + F(x)
}
$$
and
$$
\xymatrix@R=1.5cm{
  F((x+y)+z) \ar[r]\ar[d]^{F(\sigma)} & F(x+y)+F(z) \ar[r]&(F(x) + F(y)) + F(z)       \ar[d]^{\sigma}  &
\\
F(x + (y + z))    \ar[r]        & F(x) + F(y +z)  \ar[r] & F(x) + (F(y)
+ F(z)). 
}
$$
commutative.

An\pageoriginale \emph{morphism of additive functors}
$u:F\rightarrow G$ is a
morphism of $\Ss$-functors (automatically an isomorphism)
making the diagram
$$\vcenter{
\xymatrix@C=2cm{
F(x+y)       \ar[r]^{u_{x+y}} \ar[d]       & G(x+y) \ar[d]\\
F(x) + F(y)  \ar[r]^{u_x + u_y}     & G(x) + G(y). 
}}
$$
commutative.

\subsubsection{}\etiquette[1.4.7]{XVIII.subsubsec1.4.7}
If $\Pp_{1}$ and $\Pp_{2}$ are two Picard stacks
on $\Ss$, the Picard stack
$\HOM(\Pp_{1},\Pp_{2})$ is the following Picard stack:
\begin{enumerate}
\renewcommand{\theenumi}{\alph{enumi}}
\renewcommand{\labelenumi}{\theenumi)}
\item if $U\in \Ob \Ss$, its objects over $U$ are the additive
functors from $\Pp_{1}|U$ to $\Pp_{2}|U$; its
morphisms are morphisms of additive functors;

\item the sum of two additive functors $F_{1}$ and
$F_{2}$ is defined by the formula $(F_{1}+F_{2})(x)=F_{1}(x)+F_{2}(x)$;
the structural isomorphism is the one making the diagram
$$
\xymatrix@C=1cm{
\scriptstyle{(F_1 + F_2)(x+y)} \ar[r]^{\rm{\scriptstyle{structural\qquad}}}\ar@{=}[dd]& \scriptstyle{(F_1 + F_2)(x) + (F_1 + F_2)(y)}\ar@{=}[rd]& \\
                                 & &\scriptstyle{F_1(x) + F_2(x) +  F_1(y)+F_2(y)}\\
\scriptstyle{F_1(x+y)+F_2(x+y)}\ar[rd]                    & &\\
&\scriptstyle{F_1(x) + F_1(y) + F_2(x) + F_2(y)}\ar@{<->}[ruu]_{\tau}&
}
$$
commutative;

\item the associativity and commutativity isomorphisms are
defined by means of the analogous isomorphisms in $\Pp_{2}$.
\end{enumerate}

\subsubsection{}\etiquette[1.4.8]{XVIII.subsubsec1.4.8}
If $\Pp_{1}$, $\Pp_{2}$, and $\Pp$ are three
Picard stacks, a biadditive functor from $\Pp_{1}\times
\Pp_{2}$ to $\Pp$ is an $\Ss$-functor $F$
from $\fprod{\Pp_{1}}{\Pp_{2}}{\Ss}$ to $\Pp$,
equipped with isomorphisms of functors
\begin{align*}
F(x_{1}+y_{1},x_{2}) &\xrightarrow{\sim} F(x_{1},
x_{2})+F(y_{1},x_{2})\\
F(x_{1}, x_{2}+y_{2}) &\xrightarrow{\sim} F(x_{1}, x_{2})+F(x_{1},
y_{2}){\quoi},
\end{align*}

% END INCLUDED SGA4_English_current_buildable_source_20260906/source/18/drafts/SGA4_XVIII_LINES_2051_2450_English.texfrag

% BEGIN INCLUDED SGA4_English_current_buildable_source_20260906/source/18/drafts/SGA4_XVIII_LINES_2451_2850_English.texfrag
such that

\begin{enumerate}
\renewcommand{\theenumi}{\alph{enumi}}
\renewcommand{\labelenumi}{\theenumi)}
\item For fixed $x_{1}$ (\resp $x_{2}$), $F(x_{1},\ast)$
(\resp $F(\ast,x_{2})$) satisfies the compatibility conditions of
  \hyperref[XVIII.subsubsec1.4.6]{1.4.6}. 

\item For\pageoriginale $U\in \Ob \Ss$, $x_{1}$, $y_{1}\in
\Ob \Pp_{1}(U)$ and $x_{2}$, $y_{2}\in \Ob \Pp_{2}(U)$,
the diagram
$$
\xymatrix@C=0cm@R=.5cm{
\scriptscriptstyle{F(x_1+y_1, x_2+y_2)}\ar[dd] \ar[r] &\scriptscriptstyle{F(x_1+y_1, x_2) + F(x_1+y_1,y_2)}\ar[rd]\\
 & &\scriptscriptstyle{F(x_1,x_2) + F(y_1,x_2) +F(x_1,y_2)+F(y_1,y_2)}\\
\scriptscriptstyle{F(x_1,x_2+y_2) + F(y_1,x_2+y_2)}\ar[rd] \\
& \scriptscriptstyle{F(x_1,x_2)+F(x_1,y_2) + F(y_1,x_2)+F(y_1,y_2)}\ar@{<->}[ruu]_{\tau}
}
$$
is commutative.
\end{enumerate}

For example, the canonical functor
$$
\HOM(\Pp_{1},\Pp_{2})\times
\Pp_{1}\longrightarrow \Pp_{2}:(F,x)\longmapsto F(x)
$$
is biadditive.

As in \hyperref[XVIII.subsubsec1.4.7]{1.4.7}, one proves that the biadditive functors from
$\Pp_{1}\times \Pp_{2}$ to $\Pp$ form a
Picard stack $\HOM(\Pp_{1}, \Pp_{2};\Pp)$.

We shall see in (\hyperref[XVIII.const1.4.20]{1.4.20}) that there exists a ``$2$-universal'' biadditive functor: $\Pp_{1}\times \Pp_{2}\rightarrow
\Pp_{1}\otimes \Pp_{2}$; more precisely, there exist
a Picard stack $\Pp_{1}\otimes \Pp_{2}$ and a
biadditive functor $\otimes:\Pp_{1}\times
\Pp_{2}\rightarrow \Pp_{1}\otimes \Pp_{2}$ such
that, for every Picard stack $\Pp$, the functor defined by
$\otimes$:
\begin{equation*}
\HOM(\Pp_{1}\otimes \Pp_{2},\Pp)\longrightarrow \HOM (\Pp_{1},\Pp_{2};\Pp)\tag{1.4.8.1}
\end{equation*}\etiquette[1.4.8.1]{XVIII.eq1.4.8.1}%
is an equivalence. This stack is unique up to a unique equivalence
(itself unique up to a unique isomorphism).

\subsubsection{}\etiquette[1.4.9]{XVIII.subsubsec1.4.9}
Let $u:\Ss_{1}\rightarrow \Ss_{2}$ be a morphism of
sites. The direct image $u_{\ast}\Pp$ of a Picard stack
$\Pp$ on $\Ss_{1}$ is the Picard stack on
$\Ss_{2}$ defined by
$$
u_{\ast}\Pp(V)=\Pp(u^{\ast}V){\quoi}.
$$
Its objects over $V\in \Ob \Ss_{2}$ are the objects of
$\Pp$ over $u^{\ast}V$; its morphisms, its addition law, and
its compatibility isomorphisms are those of $\Pp$. The direct
image $u_{\ast}$ is a $2$-functor.

\subsubsection{}\etiquette[1.4.10]{XVIII.subsubsec1.4.10}
A\pageoriginale \emph{Picard prestack} $\Pp$ on a site
$\Ss$ is a prestack $\Pp$ on $\Ss$ (\textsc{Giraud}
\cite{sga4.exp.XVIII.bib.1}), equipped with a functor
$+:\fprod{\Pp}{\Pp}{\Ss}\rightarrow \Pp$ and
associativity and commutativity isomorphisms $\sigma$ and
$\tau$, such that, for every $U\in \Ob \Ss$, $\Pp(U)$
is a (strictly commutative) Picard category. If
$\Pp_{1}$ and $\Pp_{2}$ are two Picard prestacks,
one defines in the evident way (\cf \hyperref[XVIII.subsubsec1.4.6]{1.4.6},
\hyperref[XVIII.subsubsec1.4.7]{1.4.7}) the \emph{additive functors}
from $\Pp_{1}$ to $\Pp_{2}$ and the Picard prestack
$\HOM(\Pp_{1}, \Pp_{2})$ that they form. If
$\Pp$ is a Picard prestack, and $j:\Pp\rightarrow
a\Pp$ is the stack that it generates, there exists, up to an
essentially unique equivalence, one and only one pair consisting of a
Picard-stack structure on $a\Pp$ and an additive-functor structure
on $j$; equipped with these data, the pair $(j,a\Pp)$
is called \emph{the Picard stack generated by} $\Pp$. For
every Picard stack $\Pp_{1}$, one has an equivalence
\begin{equation*}
\HOM(a\Pp,\Pp_{1})\xrightarrow{\sim} \HOM(\Pp,\Pp_{1}){\quoi}.\tag{1.4.10.1}
\end{equation*}\etiquette[1.4.10.1]{XVIII.eq1.4.10.1}

\subsubsection{}\etiquette[1.4.11]{XVIII.subsubsec1.4.11}
We shall denote by $C^{[-1,0]}(\Ss)$ the category of complexes
of abelian sheaves $K$ on $\Ss$ such that $K^{i}=0$ for
$i\not\in [-1,0]$. To every complex $K\in \Ob C^{[-1,0]}(\Ss)$
$$
K:d:K^{-1}\longrightarrow K^{0}
$$
is associated the following Picard prestack $\pch(K)$:

\begin{enumerate}
\renewcommand{\theenumi}{\Roman{enumi}}
\renewcommand{\labelenumi}{(\theenumi)}
\item For $U\in \Ob \Ss$, one has $\Ob \pch(K)(U)=K^{0}(U)$.

\item If $x$, $y\in K^{0}(U)$, an arrow from $x$ to y is an element
$f$ of $K^{-1}(U)$ such that $df=y-x$.

\item The composition law for arrows is the addition law in
$K^{-1}(U)$.

\item The functor $+$ is given by the addition law in $K^{0}(U)$
and $K^{-1}(U)$.

\item The associativity and commutativity isomorphisms are
provided by the zero element of $K^{-1}(U)$.
\end{enumerate}

We shall denote by $\ch(K)$ the Picard stack generated
(\hyperref[XVIII.subsubsec1.4.10]{1.4.10}) by the Picard
prestack $\pch(K)$. The sheaf $ℋ^{0}(K)$ is interpreted
as the following associated sheaf
\begin{align*}
ℋ^{0}(K) & \simeq a\; \text{\rm ($U\longmapsto$ group of isomorphism
classes}\tag{1.4.11.1}\\
&\qquad\qquad \text{\rm of objects of $\ch(K)(U)$)}{\quoi}.
\end{align*}\etiquette[1.4.11.1]{XVIII.eq1.4.11.1}

On the\pageoriginale other hand,
\begin{equation*}
ℋ^{-1}(K)\simeq \sheafAut(e)\tag{1.4.11.2}
\end{equation*}\etiquette[1.4.11.2]{XVIII.eq1.4.11.2}

\subsubsection{}\etiquette[1.4.12]{XVIII.subsubsec1.4.12}
Every morphism of complexes $f:K\rightarrow L$ induces an additive
functor $\pch(f):\pch(K)\rightarrow \pch(L)$, hence an additive functor
$$
\ch(f):\ch(K)\longrightarrow \ch(L){\quoi}.
$$

It follows from \eqhyperref[XVIII.eq1.4.11.1]{1.4.11.1} and \eqhyperref[XVIII.eq1.4.11.2]{1.4.11.2} that $\ch(f)$ is an
equivalence if and only if $f$ is a quasi-isomorphism. From
\eqhyperref[XVIII.eq1.4.10.1]{1.4.10.1} one deduces that, for two morphisms $f$, $g:K\rightarrow L$,
one has
$$
\Hom(\ch(f),\ch(g))\xrightarrow{\sim} \Hom(\pch(f),\pch(g)){\quoi}.
$$

A morphism of \emph{functors} $h:\pch(f)\rightarrow \pch(g)$ is a
morphism of sheaves $h:K^{0}\rightarrow L^{-1}$ such that
\begin{equation*}
g(x)-f(x)=dh(x)\tag{1.4.12.1}
\end{equation*}\etiquette[1.4.12.1]{XVIII.eq1.4.12.1}%
and such that, for every triple $(x,y,u)$ with $y-x=du$, one
has
\begin{equation*}
h(y)+f(u)=g(u)+h(x){\quoi}\tag{1.4.12.2}
\end{equation*}\etiquette[1.4.12.2]{XVIII.eq1.4.12.2}

That $h$ is a morphism of additive functors means that
\begin{equation*}
h(x+y)=h(x)+h(y){\quoi}.\tag{1.4.12.3}
\end{equation*}\etiquette[1.4.12.3]{XVIII.eq1.4.12.3}%
Condition \eqhyperref[XVIII.eq1.4.12.3]{1.4.12.3} makes it possible to rewrite \eqhyperref[XVIII.eq1.4.12.2]{1.4.12.2} as
\begin{equation*}
g(u)-f(u)=h\sisi{du}{(du)} ;\tag{$1.4.12.2'$}
\NextLinkTarget{sga4.exp.XVIII.equation.1.4.12.2-prime}\label{1.4.12.2'}
\end{equation*}
conditions \eqhyperref[XVIII.eq1.4.12.1]{1.4.12.1} through \eqhyperref[XVIII.eq1.4.12.3]{1.4.12.3} therefore mean that $h$ is a
homotopy from $f$ to $g$, and
\begin{equation*}
\Hom(\ch(f),\ch(g))\simeq \{H:\;\text{\rm homotopy } K\longrightarrow
L| g-f=dH+Hd\}\tag{1.4.12.4}
\end{equation*}\etiquette[1.4.12.4]{XVIII.eq1.4.12.4}

\begin{enonce*}{Lemma 1.4.13}\etiquette[1.4.13]{XVIII.lem1.4.13}
\begin{enumerate}
\renewcommand{\theenumi}{\Roman{enumi}}
\renewcommand{\labelenumi}{(\theenumi)}
\item For every Picard stack $\Pp$, there exists a complex
$K$ such that $\Pp=\ch(K)$.

\item For every additive functor $F:\ch(K)\rightarrow \ch(L)$, there
exist a quasi-isomorphism $k:K'\rightarrow K$ and a morphism
$\ell:K'\rightarrow L$ such that $F\simeq \ch(\ell)\ch(k)^{-1}$.
\end{enumerate}
\end{enonce*}

Let us prove\pageoriginale (I). Let $(k_{i},U_{i})_{i\in I}$ be a family
such that
\begin{enumerate}
\renewcommand{\theenumi}{\alph{enumi}}
\renewcommand{\labelenumi}{(\theenumi)}
\item $k_{i}\in \Ob \Pp(U_{i})$;

\item every object of $\Pp$ is locally isomorphic to an inverse
image of one of the $k_{i}$.

Set $K^{0}=\displaystyle{\mathop{\oplus}_{i\in
I}}\ZZ_{U_{i}}$. It follows from \hyperref[XVIII.lem1.4.3]{1.4.3} that there exists an additive
functor $F$ from $\ch(0\rightarrow K^{0})$ to $\Pp$ taking
the basis $e_{i}$ of $\ZZ_{U_{i}}$ to $k_{i}$. Let $K^{-1}$ be the
sheaf of pairs consisting of a local section $x$ of $K^{0}$ and
an isomorphism $t:F(0)\xrightarrow{\sim} F(x)$. Addition on
$K^{-1}$ is defined by $(x_{1},
t_{1})+(x_{2},t_{2})=(x_{1}+x_{2},t)$, where $t$ makes the
diagram
$$
\xymatrix@C=1.5cm{
F(0) + F(0)    \ar[r]^{t_1+t_2} \ar@{<->}[d]   & F(x_1) + F(x_2)\ar@{<->}[d]\\
F(0+0)          \ar[r]^{t}         & F(x_1+x_2).
}
$$
commutative. Define $d:K^{-1}\rightarrow K^{0}$ by $d(x,t)=x$; $d$ is
additive.

If $y-x=dt$, there exists a unique morphism $F(t)$ making the
diagram
$$
\xymatrix@C=1.75cm{
F(x)           \ar[r]^{F(t)}   \ar@{<->}[d]& F(y) \ar@{<->}[d]\\
F(0)+F(x)      \ar[r]^{t+F(x)} & F(y-x)+F(x)
}
$$
commutative, and one verifies that this construction defines an equivalence
$$
F:\ch(K)\longrightarrow \Pp{\quoi}.
$$

Let us prove (II). Thus let $F:\ch(K)\rightarrow \ch(L)$ be an additive
functor. There then exists a family $(k_{i},\ell_{i}, \alpha_{i},
U_{i})_{i\in I}$ such that

\item $k_{i}\in K^{0}(U_{i})$, $\ell_{i}\in L^{0}(U_{i})$, and
$\alpha_{i}$ is a morphism from $F(k_{i})$ to $\ell_{i}$;

\item the homomorphism $k^{0}:\displaystyle{\mathop{\oplus}_{i\in
I}}\ZZ_{U_{i}}\rightarrow K^{0}$ with coordinates $k_{i}$ is an
epimorphism.

Set ${K'}^{0}=\displaystyle{\mathop{\oplus}_{i\in
I}}\ZZ_{U_{i}}$, and let $\ell^{0}:{K'}^{0}\rightarrow L^{0}$ be
the homomorphism with coordinates $\ell_{i}$. Since the functor $F$ is
additive, there exists a unique family of isomorphisms
$\alpha_{x}:Fk^{0}(x)\xrightarrow{\sim}\ell^{0}(x)$ ($x$ a local
section of ${K'}^{0}$) such that

\item for $e_{i}$ the \hyperref[XVIII.sec1]{section 1} of $\ZZ_{U_{i}}$, one has
$\alpha_{e_{i}}=\alpha_{i}$;

\item The\pageoriginale diagrams
$$
\xymatrix@C=2cm{
Fk^0(x+y)  \ar[r]^{\alpha_{x+y}} \ar@{<->}[d]& \ell^0(x+y)
\ar@{=}[d]\\
Fk^0(x)+ Fk^0(y)    \ar[r]^{\alpha_{x} + \alpha_{y}}     & \ell^0(x)+\ell^0(y)
}
$$
are commutative.

Let ${K'}^{-1}=\fprod{K^{-1}}{{K'}^{0}}{K^{0}}$; by (d), the
evident arrow is a quasi-isomorphism of complexes
$k:K'\rightarrow K$. If $x$, $y\in {K'}^{0}(U)$, and if $t\in
{K'}^{-1}(U)$ satisfies $y-x=dt$, then there exists a
unique section $\ell^{-1}(t;x,y)$ of $L^{-1}(U)$ defining the
morphism that makes the diagram

\item $\begin{array}[t]{c}
\xymatrix@C=2.75cm@R=1.75cm{
&F(k^0(x))  \ar[r]^{F(k^{-1}(t))} \ar[d]^{\alpha_x}&
  F(k^0(y))\ar[d]^{\alpha_y} \\ 
&\ell^0(x)     \ar[r]^{\ell^{-1}(t;x,y)}& \ell^0(y).
}
\end{array}$


commutative. The fact that $F$ is a functor gives

\item $\ell^{-1}(t;x,y)+\ell^{-1}(t';y,z)=\ell^{-1}(t+t';x,z)$.

The functoriality of the additivity isomorphism of $F$ gives

\item $\ell^{-1}(t_{1}+t_{2};x_{1}+x_{2},
y_{1}+y_{2})=\ell^{-1}(t_{1};x_{1},y_{1})+\ell^{-1}(t_{2};x_{2},y_{2})$.

From (h), one obtains
\begin{align*}
&\ell^{-1}(0;x,x)+\ell^{-1}(0;x,x)=\ell^{-1}(0;x,x)\qquad \text{\ie}\\
& \ell^{-1}(0;x,x)=0{\quoi}.
\end{align*}
It then follows from (i) that $\ell^{-1}(t;x,y)$ depends only on $t$:
$$
\ell^{-1}(t;x,y)=\ell^{-1}(t)
$$
and that $\ell^{-1}(t)$ is additive in $t$. One concludes that
$\ell=(\ell^{-1}, \ell^{0})$ is a morphism of complexes and,
by (f) and (g), the arrows $\alpha$ form an isomorphism of
additive functors $F\circ \ch(k)=\ch(\ell)$.
\end{enumerate}

\setcounter{subsubsection}{13}
\subsubsection{}\etiquette[1.4.14]{XVIII.subsubsec1.4.14}
Let $\Ch^{\flat}(\Ss)$ be the category whose objects
are the small Picard stacks on $\Ss$, and whose
arrows are the isomorphism classes of additive\pageoriginale
functors. The construction $\ch$ defines a functor from
$C^{[-1,0]}(\Ss)$ (\hyperref[XVIII.subsubsec1.4.11]{1.4.11}) to
$\Ch^{\flat}(\Ss)$. Let $D^{[-1,0]}(\Ss)$ be the
subcategory of the derived category of $\Ss$ formed by the
complexes $K$ such that $H^{i}(K)=0$ for $i\neq 0$ or $-1$. It follows
from \hyperref[XVIII.subsubsec1.4.12]{1.4.12} and \hyperref[XVIII.lem1.4.13]{1.4.13} that

\begin{enonce*}{Proposition 1.4.15}\etiquette[1.4.15]{XVIII.prop1.4.15}
The functor $\ch$ induces an equivalence of categories
$$
\ch:D^{[-1,0]}(\Ss)\xrightarrow{\sim}
\Ch^{\flat}(\Ss){\quoi}.
$$
\end{enonce*}

We shall denote by $\flat$ the inverse equivalence of $\ch$. For every
Picard stack $\Pp$ on $\Ss$,
$\Pp^{\flat}\in \Ob D^{[-1,0]}(\Ss)$ determines
$\Pp$ up to the isomorphism class of an equivalence.

\begin{enonce*}{Lemma 1.4.16}\etiquette[1.4.16]{XVIII.lem1.4.16}
Let $L\in \Ob C^{[-1,0]}(\Ss)$ be such that $L^{-1}$ is
injective.
\begin{enumerate}
\renewcommand{\theenumi}{\Roman{enumi}}
\renewcommand{\labelenumi}{(\theenumi)}
\item $\pch(L)$ is already a stack.

\item For $K\in \Ob C^{[-1,0]}(\Ss)$, every additive functor
$F:\ch(K)\rightarrow \ch(L)$ is isomorphic to a functor $\ch(f)$ for
$f$ a morphism from $K$ to $L$.
\end{enumerate}
\end{enonce*}

\begin{enumerate}
\renewcommand{\theenumi}{\Roman{enumi}}
\renewcommand{\labelenumi}{(\theenumi)}
\item For any $U\in \Ob \Ss$ and $x\in \Ob \ch(L)(U)$,
the sheaf of pairs consisting of a local section $\ell$ of
$L^{0}|U$ and an isomorphism between $x$ and $\ell$ is a principal
homogeneous space under $L^{-1}$; it therefore has a section, and
$j:\pch(L)\rightarrow \ch(L)$ is surjective on the isomorphism
classes of objects, hence is an equivalence.

\item Let $L'$ be a complex of injectives, equipped with a quasi-isomorphism
$\varphi:L{\xrightarrow[\approx]{}}L'$, such that $\varphi^{i}$ is an
isomorphism for $i\leq -1$. By \hyperref[XVIII.lem1.4.13]{1.4.13} and
\cite{sga4.exp.XVIII.bib.9}, there exists a diagram commutative up to homotopy
$$
\xymatrix@R=1cm@C=1.5cm{
L  \ar[r]^{\varphi}_{\approx} & L'\\
K' \ar[u]^{f_0}  \ar[r]^{\Psi}_{\approx}   & K\ar[u]_{f_1}
}
$$
such that $F\simeq \ch(f_{0})\ch(\Psi)^{-1}$. Moreover, $K=\tau_{\leq
0}K$ and $L=\tau_{\leq 0}L'$, so that $f_{1}$ factors through
$f:K\rightarrow L$, and $F\simeq \ch(f)$.
\end{enumerate}

The lemma refines \hyperref[XVIII.prop1.4.15]{1.4.15} by the following consequence
of \hyperref[XVIII.subsubsec1.4.12]{1.4.12}.

\begin{enonce*}{Corollary 1.4.17}\etiquette[1.4.17]{XVIII.cor1.4.17}
The\pageoriginale construction $\ch$ defines an equivalence of $2$-categories between
\begin{enumerate}
\renewcommand{\theenumi}{\alph{enumi}}
\renewcommand{\labelenumi}{\theenumi)}
\item the $2$-category of Picard stacks on $\Ss$.

\item the $2$-category whose objects and $1$-arrows are the
objects and arrows of the full subcategory of
$C^{[-1,0]}(\Ss)$ formed by the complexes $L$ with $L^{-1}$
injective, and whose $2$-arrows are homotopies between arrows:
$\Hom(f,g)=\{H|g-f=\text{\rm dH}+\text{\rm Hd}\}$.
\end{enumerate}
\end{enonce*}

\begin{enonce*}[remark]{Construction 1.4.18}\etiquette[1.4.18]{XVIII.const1.4.18}
One has
$$
\HOM(\Pp_{1}, \Pp_{2})^{\flat}\simeq \tau_{\leq
0}R\Hhom (\Pp^{\flat}_{1}, \Pp^{\flat}_{2}){\quoi}.
$$
\end{enonce*}

Let $K$, $L\in \Ob C^{[-1,0]}(\Ss)$, with $L^{-1}$ injective. By
\hyperref[XVIII.lem1.4.16]{1.4.16} (II) and \hyperref[XVIII.subsubsec1.4.12]{1.4.12}, one has
\begin{equation*}
\ch(\tau_{\leq
0}\Hhom(K,L))\xrightarrow{\sim}\HOM(\ch(K),\ch(L))\tag{1.4.18.1}
\end{equation*}\etiquette[1.4.18.1]{XVIII.eq1.4.18.1}%
and one deduces Construction \hyperref[XVIII.const1.4.18]{1.4.18} from this functorially.

\begin{enonce*}[remark]{Construction 1.4.19}\etiquette[1.4.19]{XVIII.const1.4.19}
For $f:\Ss_{1}\to\Ss_{2}$ a morphism of sites, one has
$$
(f_{\ast}\Pp)^{\flat}\simeq \tau_{\leq
0}\,Rf_{\ast}(\Pp^{\flat}){\quoi}.
$$
\end{enonce*}

Let $L\in \Ob C^{[-1,0]}$ with $L^{-1}$ injective. By
\hyperref[XVIII.lem1.4.16]{1.4.16} (I), one has

% END INCLUDED SGA4_English_current_buildable_source_20260906/source/18/drafts/SGA4_XVIII_LINES_2451_2850_English.texfrag

% BEGIN INCLUDED SGA4_English_current_buildable_source_20260906/source/18/drafts/SGA4_XVIII_LINES_2851_3250_English.texfrag
\begin{equation*}
\ch(f_{\ast}L)\xrightarrow{\sim} f_{\ast}\ch(L)\tag{1.4.19.1}
\end{equation*}\etiquette[1.4.19.1]{XVIII.eq1.4.19.1}%
and from this one deduces a functorial construction \hyperref[XVIII.const1.4.19]{1.4.19}.

We shall not need to use the following

\begin{enonce*}[remark]{Construction 1.4.20}\etiquette[1.4.20]{XVIII.const1.4.20}
The tensor product promised in \hyperref[XVIII.subsubsec1.4.8]{1.4.8} exists, and
$$
(\Pp_{1}\otimes \Pp_{2})^{\flat}\simeq \tau_{\geq
-1}\Pp^{\flat}_{1}\mathop{\otimes}^{\LL}\Pp^{\flat}_{2}{\quoi}.
$$
\end{enonce*}

One checks, as in \hyperref[XVIII.subsubsec1.4.12]{1.4.12} and
\hyperref[XVIII.lem1.4.13]{1.4.13} (II), that for all
$K_{1}$, $K_{2}$, and $L$ in $C^{[-1,0]}(\Ss)$, one has
\begin{enumerate}
\renewcommand{\theenumi}{\alph{enumi}}
\renewcommand{\labelenumi}{(\theenumi)}
\item Every morphism $f\in \Hom(K_{1}\otimes K_{2},L)\simeq
\Hom(\tau_{\geq-1}(K_{1}\otimes K_{2}),L)$ defines a biadditive
functor $\ch(f)\in \Hom(\ch(K_{1})$, $\ch(K_{2})$; $\ch(L)$); moreover,
morphisms of functors identify with homotopies;

\item For every biadditive functor $F:\ch(K_{1})$,
$\ch(K_{2})\rightarrow \ch(L)$, there exist quasi-isomorphisms
$k_{i}:K'_{i}\rightarrow K_{i}$ and $f:K'_{1}\otimes
K'_{2}\rightarrow L$ such that $F\simeq \ch(f)\circ
(\ch(k_{1})^{-1},\ch(k_{2})^{-1})$.
\end{enumerate}

If\pageoriginale now $K^0_{1}$ or $K^0_{2}$ is flat,
as is ${K'}^0_{1}$ or ${K'}^0_{2}$, then
$k_{1}\otimes k_{2}:\tau_{\geq -1}{K'}_{1}\otimes {K'}_{2}\rightarrow
\tau_{\geq -1}K_{1}\otimes K_{2}$ is a quasi-isomorphism, and it
follows that $\ch(\tau_{\geq -1}(K_{1}\otimes K_{2}))$ satisfies the
universal property \eqhyperref[XVIII.eq1.4.8.1]{1.4.8.1}, whence the existence of $\otimes$ and
of a functorial construction \hyperref[XVIII.const1.4.20]{1.4.20}.

\setcounter{subsubsection}{20}
\subsubsection{}\etiquette[1.4.21]{XVIII.subsubsec1.4.21}
Let $G$ be an abelian sheaf, and let $G[1]$ be the complex with value
$G$ placed in degree $-1$. The prestack $\pch(G[1])$ identifies
with the prestack of trivial torsors under $G$, and hence $\ch(G[1])$
``is none other than'' the stack of torsors under $G$, with its usual
addition law.

\subsubsection{}\etiquette[1.4.22]{XVIII.subsubsec1.4.22}
Let $K\in C^{[-1,0]}(\Ss)$, and let $G$ be an abelian sheaf. The
extensions $E$ of $K$ by $G$ form a Picard stack, $\EXT(K,G)$, under
% SGA-SRC-001136
Baer addition.
$$
E:0\rightarrow
G[0]\xrightarrow{\alpha}E\xrightarrow{\beta}K\rightarrow 0{\quoi}.
$$

To every extension $E$ one associates an additive functor from $\ch(K)$
to $\ch(G[1])$ which, when applied to a local section $x$ of $K^{0}$,
gives the torsor $\beta^{-1}(x)$. The reader will verify the following.

\begin{enonce*}{Proposition 1.4.23}\etiquette[1.4.23]{XVIII.prop1.4.23}
The construction outlined above is an equivalence of Picard stacks
$$
\EXT(K,G)\xrightarrow{\sim} \HOM(\ch(K),\ch(G[1])){\quoi}.
$$
\end{enonce*}

\setcounter{subsubsection}{23}
\subsubsection{}\etiquette[1.4.24]{XVIII.subsubsec1.4.24}
Let $F$ be an object of the site $\Ss$. We shall also write $F$ for
the sheaf represented by $F$, and write $\ZZ^{(F)}$ for the abelian
sheaf it generates. Let $C$ be a Picard stack on $\Ss$. One checks
easily (\cf \hyperref[XVIII.lem1.4.3]{1.4.3}) that giving any one of the following amounts to the
same thing:
\begin{enumerate}
\renewcommand{\theenumi}{\alph{enumi}}
\renewcommand{\labelenumi}{\theenumi)}
\item an additive functor $H:\ch(\ZZ^{(F)})\rightarrow C$;

\item a morphism of \emph{stacks}:

(stack of local sections of $F$, with only identity morphisms)
$\rightarrow C$;

\item an object of $C(F)$.
\end{enumerate}

Suppose (in order to be able to apply Definition
\hyperref[XVIII.subsubsec1.4.9]{1.4.9}) that fiber products exist in $\Ss$, and let $f$ be the
canonical morphism of sites $f:\Ss/F\rightarrow \Ss$. The preceding
construction gives an equivalence
\begin{equation*}
\HOM(\ch(\ZZ^{(F)}),C)\xrightarrow{\sim}
f_{\ast}f^{\ast}C{\quoi}.\tag{1.4.24.1}
\end{equation*}\etiquette[1.4.24.1]{XVIII.eq1.4.24.1}

The isomorphism\pageoriginale obtained from \eqhyperref[XVIII.eq1.4.24.1]{1.4.24.1} by the construction
$$
\tau_{\leq 0} R\Hhom(\ZZ^{(F)},C^{\flat})\xrightarrow{\sim}
\tau_{\leq 0} Rf_{\ast} f^{\ast}C^{\flat}{\quoi},
$$
may also be deduced from the more general isomorphism
\begin{equation*}
R\Hhom(\ZZ^{(F)},K)\xrightarrow{\sim}
Rf_{\ast}f^{\ast}K\quad (K\in \Ob D^{+}(\Ss)){\quoi}.\tag{1.4.24.2}
\end{equation*}\etiquette[1.4.24.2]{XVIII.eq1.4.24.2}

\subsection{The universal coefficient formula}\etiquette[1.5]{XVIII.subsec1.5}

\subsubsection{}\etiquette[1.5.1]{XVIII.subsubsec1.5.1}
Let $f:X\rightarrow S$ be a smooth projective curve over $S$. We
denote by $\PIC(X/S)$ the Picard stack on the big fppf site of $S$
(\hyperrefext[XVII][XVII.subsec0.10]{0.10}), namely the direct image (\hyperref[XVIII.subsubsec1.4.9]{1.4.9}) of the stack of
invertible sheaves on $X$. We have
\begin{align*}
\PIC(X/S)&= f_{\ast}\ch\GG_{m}[1]\tag{1.5.1.1}\\
\PIC(X/S)^{\flat} &=\tau_{\leq 0}
Rf_{\ast}(\GG_{m}[1]){\quoi}.\tag{1.5.1.2}
\end{align*}\etiquette[1.5.1.1]{XVIII.eq1.5.1.1}
\etiquette[1.5.1.2]{XVIII.eq1.5.1.2}

Each section $t$ of $X$ defines an invertible sheaf $\Oo(t)$ on $X$.
With the notation of \hyperref[XVIII.subsubsec1.4.24]{1.4.24} (\cf
\hyperref[XVIII.subsubsec1.4.24]{1.4.24} $a\leftrightarrow b$), this construction, being compatible with
every base change, defines a morphism of Picard stacks on $S$
\begin{equation*}
\ch \ZZ^{(X)}\longrightarrow \PIC(X/S){\quoi}.\tag{1.5.1.3}
\end{equation*}\etiquette[1.5.1.3]{XVIII.eq1.5.1.3}

If $C$ is any Picard stack on $S$, the morphisms
\eqhyperref[XVIII.eq1.5.1.3]{1.5.1.3} and \eqhyperref[XVIII.eq1.4.24.1]{1.4.24.1} define an additive functor
\begin{equation*}
\HOM(\PIC(X/S),C)\longrightarrow\HOM(\ch(\ZZ^{(X)}),
C)\xrightarrow{\sim} f_{\ast}f^{\ast}C{\quoi}.\tag{1.5.1.4}
\end{equation*}\etiquette[1.5.1.4]{XVIII.eq1.5.1.4}

Upon applying the construction $\flat$, the functor \eqhyperref[XVIII.eq1.5.1.3]{1.5.1.3}
defines a morphism\nde{Here the shift functor
$[1]$ is to be understood to have higher precedence than the truncation $\tau_{\leq 0}$.}
\begin{equation*}
\ZZ^{(X)}\longrightarrow \tau_{\leq 0}
Rf_{\ast}\GG_{m}[1]\longrightarrow
Rf_{\ast}\GG_{m}[1]{\quoi},\tag{1.5.1.5}
\end{equation*}\etiquette[1.5.1.5]{XVIII.eq1.5.1.5}%
and the analogue of the functor \eqhyperref[XVIII.eq1.5.1.4]{1.5.1.4} is, by \eqhyperref[XVIII.eq1.4.24.2]{1.4.24.2}, a morphism
\begin{equation*}
R\Hhom(\tau_{\leq 0} Rf_{\ast}(\GG_{m}[1]),K)\longrightarrow
Rf_{\ast} f^{\ast}K{\quoi}.\tag{1.5.1.6}
\end{equation*}\etiquette[1.5.1.6]{XVIII.eq1.5.1.6}

If $C$ is the stack $\ch(G[1])$ of torsors under a group $G$
satisfying (\hyperref[XVIII.par1.3.1.1]{1.3.1.1}), Theorem \hyperref[XVIII.thm1.3.10]{1.3.10}
asserts that the functor
\eqhyperref[XVIII.eq1.5.1.4]{1.5.1.4} is an equivalence,\pageoriginale the inverse equivalence being the one that
associates with a torsor $K$ under $G_{X}$ the additive functor
$\langle \ast,K]$. The assertion that \eqhyperref[XVIII.eq1.5.1.4]{1.5.1.4} is an equivalence amounts to
saying that the morphism deduced from \eqhyperref[XVIII.eq1.5.1.4]{1.5.1.4} or \eqhyperref[XVIII.eq1.5.1.6]{1.5.1.6}
\begin{equation*}
\tau_{\leq 1}R\Hhom(\tau_{\leq
0}Rf_{\ast}(\GG_{m}[1]),G)\longrightarrow \tau_{\leq 1}
Rf_{\ast}f^{\ast}G\tag{1.5.1.7}
\end{equation*}\etiquette[1.5.1.7]{XVIII.eq1.5.1.7}%
is an isomorphism.

\begin{enonce*}{Theorem 1.5.2}\etiquette[1.5.2]{XVIII.thm1.5.2}
\emph{(universal coefficient formula).}

Let $f:X\rightarrow S$ be a smooth projective curve over $S$, let $C$
be a Picard stack on $S_{\text{\rm fppf}}$, and let $K=C^{\flat}$
(\hyperref[XVIII.prop1.4.15]{1.4.15}). Suppose that $K$ is locally isomorphic (in
$D^{+}(S_{\text{\rm fppf}})$) to complexes of the form
$G_{-1}\rightarrow G_{0}$, where the sheaves $G_{i}$ satisfy
(\hyperref[XVIII.par1.3.1.1]{1.3.1.1}). Then:
\begin{enumerate}
\renewcommand{\theenumi}{\Roman{enumi}}
\renewcommand{\labelenumi}{(\theenumi)}
\item The functor \eqhyperref[XVIII.eq1.5.1.4]{1.5.1.4}
\begin{equation*}
\HOM(\PIC(X/S),C)\longrightarrow f_{\ast}f^{\ast}C\tag{1.5.2.1}
\end{equation*}\etiquette[1.5.2.1]{XVIII.eq1.5.2.1}%
is an equivalence of Picard stacks.

\item The morphism deduced from \eqhyperref[XVIII.eq1.5.1.6]{1.5.1.6} or \eqhyperref[XVIII.eq1.5.1.4]{1.5.1.4}
\begin{equation*}
\tau_{\leq 0}R\Hhom(\tau_{\leq 0}
Rf_{\ast}(\GG_{m}[1]),K)\longrightarrow \tau_{\leq 0}
Rf_{\ast}f^{\ast}K\tag{1.5.2.2}
\end{equation*}\etiquette[1.5.2.2]{XVIII.eq1.5.2.2}%
is an isomorphism.
\end{enumerate}
\end{enonce*}

This theorem makes $\tau_{\leq 0}Rf_{\ast}(\GG_{m}[1])$ play the
role played by homology in the classical universal coefficient
formula.

It is clear that, for every stack $C$, one has $(I)\leftrightarrow
(II)$, and both $(I)$ and $(II)$ are local in nature on $S$. If $C$
is of the form $\ch(G[1])$, then the theorem follows from
\hyperref[XVIII.thm1.3.10]{1.3.10}, as noted above. It remains to show that, if $K$ is of the form
$$
K:d:G_{-1}\longrightarrow G_{0}{\quoi},
$$
and if \eqhyperref[XVIII.eq1.5.2.2]{1.5.2.2} is an isomorphism for $G_{-1}[1]$ and $G_{0}[1]$,
then \eqhyperref[XVIII.eq1.5.2.2]{1.5.2.2} is an isomorphism for $K$. To verify this, it suffices to
compare the cohomology long exact sequence
$$
0\longrightarrow R^{-1}f_{\ast}K\longrightarrow
f_{\ast}G_{-1}\longrightarrow f_{\ast}G_{0}\longrightarrow
R^{0}f_{\ast}K\longrightarrow R^{1} f_{\ast}G_{-1}\longrightarrow
R^{1}f_{\ast}G_{0}
$$
with\pageoriginale the analogous long exact sequence for the left-hand
side of \eqhyperref[XVIII.eq1.5.2.2]{1.5.2.2}, and to apply the five lemma.

\setcounter{subsubsection}{2}
\subsubsection{}\etiquette[1.5.3]{XVIII.subsubsec1.5.3}
Let $u$ be a flat morphism of smooth projective curves over $S$
$$
\xymatrix@=1.25cm{
X              \ar[rr]^{u} \ar[rd]_{f_1}&       &Y \ar[ld]^{f_2}\\
                                       &  S    &           .
}
$$
The diagram of additive functors
$$\vcenter{\etiquette[1.5.3.1]{XVIII.eq1.5.3.1}
\xymatrix{
\ZZ^{(X)}  \ar[r]\ar[d]^{u} & \PIC(X/S) \ar[d]^{N_{X/Y}}\\
\ZZ^{(Y)}  \ar[r] & \PIC(Y/S)
}}\leqno{(1.5.3.1)}
$$
is then essentially commutative.

For every Picard stack $C$ on $S$, the diagram
$$\vcenter{\etiquette[1.5.3.2]{XVIII.eq1.5.3.2}
\xymatrix@R=1.75cm@C=2cm{
\HOM(\PIC(Y/S),C)         \ar[r] \ar[d]^{\HOM(N_{X/Y},C)}&
f_{2\ast}f^{\ast}_2 C \ar[d]^{u^{\ast}}\\
\HOM(\PIC(X/S), C)        \ar[r] & f_{1\ast}f^{\ast}_1 C
}}\leqno{(1.5.3.2)}
$$
is therefore essentially commutative, and for $K\in \Ob D^{+}(S)$,
the diagram
$$\vcenter{\etiquette[1.5.3.3]{XVIII.eq1.5.3.3}
\xymatrix@C=2cm@R=1.5cm{
R\Hhom(\tau_{\leq 0} Rf_{2\ast}\GG_m[1],K) \ar[r]\ar[d]^{R\Hhom(N_{X/Y},K)}&
Rf_{2\ast}f^{\ast}_2 K  \ar[d]^{u^{\ast}}\\
R\Hhom(\tau_{\leq 0} Rf_{1\ast}\GG_m[1],K) \ar[r]&
Rf_{1\ast}f^{\ast}_{1} K
}}\leqno{(1.5.3.3)}
$$
is commutative.

Now let $G$ be an abelian sheaf on $S$ satisfying
(\hyperref[XVIII.par1.3.1.1]{1.3.1.1}). The trace \hyperrefext[XVII][XVII.subsec6.3]{6.3} gives us an additive functor
\begin{equation*}
\Tr_{u}:f_{1\ast}f^{\ast}_{1}\ch(G[1])\longrightarrow f_{2\ast}f^{\ast}_{2}\ch(G[1])\tag{1.5.3.4}
\end{equation*}\etiquette[1.5.3.4]{XVIII.eq1.5.3.4}%
and\pageoriginale \hyperref[XVIII.par1.3.8.6]{1.3.8.6} gives us an isomorphism making the diagram
essentially commutative
$$\vcenter{\etiquette[1.5.3.5]{XVIII.eq1.5.3.5}
\xymatrix@R=1.75cm{
\HOM(\PIC(X/S), \ch(G[1])) \ar[d]^{\HOM(u^{\ast}, \ch(G[1]))}&  f_{1\ast}f^{\ast}_1
\ch(G[1]) \ar[l] \ar[d]^{\Tr_u}\\
\HOM(\PIC(Y/S), \ch(G[1])) & f_{2\ast}f^{\ast}_2 \ch(G[1]).\ar[l]
}}\leqno{(1.5.3.5)}
$$

The diagram
$$\vcenter{\etiquette[1.5.3.6]{XVIII.eq1.5.3.6}
\xymatrix@R=1.75cm{
\tau_{\leq 1} R\Hhom(\tau_{\leq 0} Rf_{1\ast} \GG_{m}[1],
G)  \ar[d]_{R\Hhom(u^{\ast},G)}  & \tau_{\leq 1} Rf_{1\ast}
f^{\ast}_1G  \ar[l]_-{\sim} \ar[d]^{Tr_u}\\
\tau_{\leq 1} R\Hhom(\tau_{\leq 0} Rf_{2\ast} \GG_m[1], G)&
  \tau_{\leq 1} Rf_{2\ast}f^{\ast}_2 G \ar[l]_-{\sim}
}}\leqno{(1.5.3.6)}
$$
is therefore commutative.

\subsubsection{}\etiquette[1.5.4]{XVIII.subsubsec1.5.4}
With the notation of \hyperref[XVIII.thm1.5.2]{1.5.2}, let $G$ be an abelian sheaf on
$S$ satisfying \hyperref[XVIII.par1.3.1.1]{1.3.1.1}. Locally on $S$ for the étale topology, one has
(noncanonically)
$$
\tau_{\leq 0}
Rf_{\ast}\GG_{m}[1]\simeq
f_{\ast}\GG_{m}[1]+R^{1}f_{\ast}\GG_{m}{\quoi}.
$$

At the level of Picard stacks, one indeed defines a section
$\ch(R^{1}f_{\ast}\GG_{m})\rightarrow
f_{\ast}\ch(\GG_{m}[1])$\nde{Recall that this stack identifies with
$\PIC(X/S)$, \cf \eqhyperref[XVIII.eq1.5.1.1]{1.5.1.1}.} by associating with each isomorphism class
of invertible sheaves the invertible sheaf belonging to that class,
rigidified along suitable sections of $X/S$.

The isomorphism \eqhyperref[XVIII.eq1.5.2.2]{1.5.2.2} therefore gives an isomorphism
\begin{equation*}
\Hhom(R^{1}f_{\ast}\GG_{m},G)\xrightarrow{\sim}\Hhom_{S}(X,G){\quoi}\tag{1.5.4.1}
\end{equation*}\etiquette[1.5.4.1]{XVIII.eq1.5.4.1}%
and an exact sequence
\begin{equation*}
0\longrightarrow
\Eext^{1}(R^{1}f_{\ast}\GG_{m},G)\longrightarrow
R^1f_{\ast}G\longrightarrow
\Hhom(f_{\ast}\GG_{m},G)\longrightarrow 0{\quoi}.\tag{1.5.4.2}
\end{equation*}\etiquette[1.5.4.2]{XVIII.eq1.5.4.2}

The morphisms in \eqhyperref[XVIII.eq1.5.4.1]{1.5.4.1} and \eqhyperref[XVIII.eq1.5.4.2]{1.5.4.2} are interpreted as
follows:
\begin{enumerate}
\renewcommand{\theenumi}{\alph{enumi}}
\renewcommand{\labelenumi}{\theenumi)}
\item The\pageoriginale morphism \eqhyperref[XVIII.eq1.5.4.1]{1.5.4.1} and the first
  morphism in \eqhyperref[XVIII.eq1.5.4.2]{1.5.4.2} are defined by the canonical map
  from $X$ to $\sheafPic_{X/S}$.

\item The second morphism in \eqhyperref[XVIII.eq1.5.4.2]{1.5.4.2} is the degree \eqhyperref[XVIII.eq1.3.7.2]{1.3.7.2}.
\end{enumerate}

Locally on $S$ for the étale topology, one has (noncanonically)
$\sheafPic_{X/S}\sim \sheafPic^{0}_{X/S}\times \ZZ^{k}$; hence
$$
\Eext^{1}(R^{1}f_{\ast}\GG_{m},G)\xrightarrow{\sim}\Eext^{1}(\sheafPic^{0}_{X/S}, G)
$$
and \eqhyperref[XVIII.eq1.5.4.2]{1.5.4.2} may also be written
\begin{equation*}
0 \longrightarrow \Eext^{1}(\sheafPic^{0}_{X/S},G)\longrightarrow
R^{1}f_{\ast}G\longrightarrow
\Hhom(f_{\ast}\GG_{m},G)\longrightarrow 0{\quoi}.\tag{1.5.4.3}
\end{equation*}\etiquette[1.5.4.3]{XVIII.eq1.5.4.3}

\subsubsection{}\etiquette[1.5.5]{XVIII.subsubsec1.5.5}
When, in \eqhyperref[XVIII.eq1.5.2.2]{1.5.2.2}, one takes $K=\GG_{m}[1]$, Theorem
\hyperref[XVIII.thm1.5.2]{1.5.2} appears as a self-duality theorem for $\tau_{\leq
0}(Rf_{\ast}\GG_{m}[1])$. Any complex representing
$\tau_{\leq 0} Rf_{\ast} \GG_{m}[1]$ admits a canonical three-step
filtration, whose successive quotients are, up to unique quasi-isomorphism,
$$
\sheafPic_{X/S}/\sheafPic^{0}_{X/S}\quav \sheafPic^{0}_{X/S}\quav
(f_{\ast}\GG_{m})[1]{\quoi}.
$$

The preceding self-duality corresponds to a duality (with values in
$\GG_{m}$) between $\sheafPic_{X/S}/\sheafPic^{0}_{X/S}$ and
$f_{\ast}\GG_{m}$, and to a self-duality (with values in
$\GG_{m}[1]$) on $\sheafPic^{0}_{X/S}$ (self-duality of the
Jacobian).

\subsubsection{}\etiquette[1.5.6]{XVIII.subsubsec1.5.6}
The situation is distinctly worse for smooth curves over
$S:f:X\rightarrow S$ that are not assumed proper over $S$. The
additive group seems to behave uncontrollably. The preceding arguments
nevertheless extend to the case of smooth curves obtained from a proper
curve over $S$ by removing a part finite over $S$, when $G$ is an
étale torsion sheaf whose torsion is prime to the residue
characteristics of $S$. One may hope that étale $p$-torsion sheaves
also behave reasonably; see \textsc{Serre}
\cite[Ch.~VI nᵒ 11 p.~126]{sga4.exp.XVIII.bib.6}.

The role of Theorem \hyperref[XVIII.thm1.2.2]{1.2.2} will here be played by the acyclicity
theorem\pageoriginale \hyperrefext[XV][XV.cor2.2]{2.2}, which has the following as a corollary.

\begin{enonce*}{Lemma 1.5.7}\etiquette[1.5.7]{XVIII.lem1.5.7}
Let $p:X\rightarrow S$ be the projection of a standard affine space
over $S$, and let $G$ be an abelian torsion sheaf on $S_{\text{\rm et}}$,
of torsion prime to the residue characteristics of $S$. We also write
$G$ for the inverse image of $G$ on the site $S_{\text{\rm
fppf}}$. Then the functor $p^{\ast}$ is an equivalence of categories
between the category of torsors under $G$ on $S_{\text{\rm
et}}$ (or on $S_{\text{\rm fppf}}$, which amounts to the same thing,
by {\textsc{Grothendieck} \cite[Corollaire~11.9]{sga4.exp.XVIII.bib.2}}), and the category of torsors under $G_{X}$
on $X$.
\end{enonce*}

\setcounter{subsubsection}{7}
\subsubsection{}\etiquette[1.5.8]{XVIII.subsubsec1.5.8}
Let $\bar{f}:\bar{X}\rightarrow S$ be a proper morphism of finite
presentation whose fibers are purely one-dimensional, let $Y$ be a
closed subscheme of $\bar{X}$ finite over $S$, defined by an ideal $\mathfrak{m}$,
and suppose that $X=\bar{X}-Y$ is a smooth curve over $S$. Let
furthermore $G$ be a torsion sheaf on $S_{\text{\rm et}}$, of torsion
prime to the residue characteristics of $S$. We shall also write $G$
for the sheaf on the big fppf site of $S$ (\hyperrefext[XVII][XVII.subsec0.10]{0.10}) obtained as
the inverse image of $G$.

Let $\GG_{m}(\mathfrak{m})$ be the subsheaf of $\GG_{m\bar{X}}$
consisting of sections equal to 1 on $Y$. A torsor under
$\GG_{m}(\mathfrak{m})$ identifies with an invertible sheaf on $\bar{X}$
trivialized on $Y$. Every relative divisor on $X$, finite over $S$,

% END INCLUDED SGA4_English_current_buildable_source_20260906/source/18/drafts/SGA4_XVIII_LINES_2851_3250_English.texfrag

% BEGIN INCLUDED SGA4_English_current_buildable_source_20260906/source/18/drafts/SGA4_XVIII_LINES_3251_3650_English.texfrag
defines such an invertible sheaf. Locally on $S$ for the
étale topology, there exist such divisors giving rise to relatively
ample invertible sheaves.

\begin{enonce*}{Lemma 1.5.9}\etiquette[1.5.9]{XVIII.lem1.5.9}
Under the hypotheses of \hyperref[XVIII.subsubsec1.5.8]{1.5.8}, with $S$ affine\nde{It is
in fact necessary to suppose that $S$ is local with closed point $s$. If one wishes to retain
the original hypothesis, then in the conclusion one should require that
$m$ be defined only over the inverse image of an open neighborhood of
$s$ in $S$, as occurs below in \hyperref[XVIII.par1.5.10.1]{1.5.10.1}.},
let $\Oo_{\bar{X}}(1)$ be a relatively ample invertible sheaf on
$\bar{X}$, trivialized along $Y$, and let $\Mm$ be an
invertible sheaf on $\bar{X}$, trivialized along $Y$. For
every point $s\in S$, there exists an integer $n_{0}$ such that, for every
integer $n\geq n_{0}$ and every section $m_{s}$ of
$\Mm(n)\otimes k(s)$ on the fiber $\bar{X}_{s}$ which is equal to $1$
on $Y_{s}$, there exists a section $m$ of $\Mm(n)$ which is equal to $1$
on $Y$ and lifts $m_{s}$.
\end{enonce*}

We reduce to the case where $S$ is Noetherian (since a family of
surjective morphisms remains a family of surjective morphisms after base
change\nde{Eliminating this Noetherian hypothesis is somewhat more
subtle than might appear at first sight,
because the question is affine rather than linear. If
$\bar{X}$, $Y$, and $\Mm$ are obtained by a base change $S\to S_0$
(which is assumed to be a local morphism between local schemes) from corresponding data
$\bar{X}_0$, $Y_0$, and $\Mm_0$ over a Noetherian $S_0$
(\cf \EGA IV 8), and if the desired assertion is known in this
Noetherian situation, then one can obtain the desired result over $S$ at
the cost of an unpleasant diagram chase, during which one may
have to consider a variant of the desired assertion in which the sections are required
to be equal to $0$ instead of $1$ on $Y_s$ or $Y$.}).
Let $\Nn$ be the coherent sheaf of local sections
of $\Mm$ vanishing on\pageoriginale $\bar{X}_{s}$ and $Y$. For
each $n$, the local solutions to the problem form a principal homogeneous
space under $\Nn(n)$, so that the obstruction to
lifting $m_{s}$ lies in
$H^{1}(X,\Nn(n))=H^{0}(S,R^{1}\bar{f}_{\ast}\Nn(n))$. But,
for $n$ sufficiently large, $R^{1}\bar{f}_{\ast}\Nn(n)=0$ (\EGA III 2.2.1).

\setcounter{subsubsection}{9}
\subsubsection{}\etiquette[1.5.10]{XVIII.subsubsec1.5.10}
Let $\Ll$ be an invertible sheaf on $\bar{X}$ trivialized
along $Y$, and let $s\in S$. Consider the conditions

\paragraph{}\etiquette[1.5.10.1]{XVIII.par1.5.10.1}
Every section of $\Ll\otimes k(s)$ on $X_{s}$ which is equal to 1 on
$Y_{s}$ lifts to a section of $\Ll$ on the inverse
image $U$ of a neighborhood of $s$, equal to 1 on $\fprod{Y}{U}{S}$.

\paragraph{}\etiquette[1.5.10.2]{XVIII.par1.5.10.2}
There exist four sections $m_{i}$ of $\Ll\otimes k(s)$ on $X_{s}$,
equal to 1 on $Y_s$, and linearly independent at the generic point of
every irreducible component of $X_{s}$.

Consider the following conditions, concerning respectively a
section $m$ of $\Ll$, a pair $(m_{1},m_{2})$, and a triple
$(m_{1},m_{2},m_{3})$ of sections of $\Ll$.

\paragraph{}\etiquette[1.5.10.3]{XVIII.par1.5.10.3}
$m$ is equal to 1 on $Y$, and the subscheme $Z(m)$ of $X$ defined by $m=0$
is a relative divisor over $S$ (which means that it is quasi-finite
over $S$, and hence finite over $S$, since $X=\bar{X}-Y$ is smooth over $S$).

\paragraph{}\etiquette[1.5.10.4]{XVIII.par1.5.10.4}
Let $\lambda$ be the canonical coordinate on the affine line
$g:\EE^{1}_{S}\rightarrow S$. The section $\lambda
m_{1}+(1-\lambda)m_{2}$ of ${g'}^{\ast}\Ll$ satisfies
(\hyperref[XVIII.par1.5.10.3]{1.5.10.3}).

\paragraph{}\etiquette[1.5.10.5]{XVIII.par1.5.10.5}
Similarly, $\lambda m_{1}+\mu m_{2}+(1-\lambda -\mu)m_{3}$ satisfies
(\hyperref[XVIII.par1.5.10.3]{1.5.10.3}).

We leave it to the reader to verify, by a general-position
argument\nde{It should be observed here that, in view of
(\hyperref[XVIII.par1.5.10.1]{1.5.10.1}) and the fact that the second part of condition
(\hyperref[XVIII.par1.5.10.3]{1.5.10.3}) can be checked in the local situation by
considering only the fiber over $s$, the lemma now involves only
sections of the restriction of $\Ll$ to the fiber
$X_s$. For example, for (I), condition (\hyperref[XVIII.par1.5.10.2]{1.5.10.2})
provides an \emph{affine} space $V$ of dimension $3$ of
sections of $\Ll\otimes k(s)$ on $X_s$ which are equal to $1$ on $Y_s$, and no
section belonging to this space $V$ vanishes at any generic point
of $X_s$. We seek to determine $k\in V$ such that, for
every $i$, $(m_i,k)$ satisfies (\hyperref[XVIII.par1.5.10.4]{1.5.10.4}), that is, such that the
affine line joining the points $m_i$ and $k$ does not pass through the origin
(in any of its incarnations associated with the different irreducible
components of $X_s$); in other words, it suffices that $k$ not belong
to the line joining $m_i$ and the origin. One thus obtains a conjunction of
conditions satisfied on nonempty open subsets of $V$. To be
completely rigorous and to account for the fact that the irreducible
components need not be geometrically irreducible, when one must take a point
of $V$ in an extension (if $k(s)$ is finite), it is in fact necessary to
introduce a finite Galois extension $s'\to s$ such that the irreducible
components of $X_{s'}$ are geometrically irreducible. The dense open subset of
$V_{s'}$ obtained by arguing on $X_{s'}$ is Galois-invariant and hence defines
the desired dense open subset of $V$.}, the following

\begin{enonce*}{Lemma 1.5.11}\etiquette[1.5.11]{XVIII.lem1.5.11}
Let $\Ll$ be an invertible sheaf on $X$, trivialized along
$Y$, and let $\bar{s}$ be a geometric point of $S$ with image $s$ in
$S$. If $\Ll$ satisfies (\hyperref[XVIII.par1.5.10.1]{1.5.10.1}) and
(\hyperref[XVIII.par1.5.10.2]{1.5.10.2}) at $s$, then:
\begin{enumerate}
\renewcommand{\theenumi}{\Roman{enumi}}
\renewcommand{\labelenumi}{(\theenumi)}
\item If $n$ sections $m_{i}$ of $\Ll$ satisfy (\hyperref[XVIII.par1.5.10.3]{1.5.10.3}),
there exist an étale neighborhood $U$ of $\bar{s}$ and a section $k$ of
$\Ll|\bar{f}^{-1}(U)$ such that the pairs $(m_{i},k)$ satisfy
(\hyperref[XVIII.par1.5.10.4]{1.5.10.4}).

\item Let\pageoriginale $m_{1}$, $m_{2}$, $k_{1}$, $k_{2}$ be four
sections of $\Ll$ such that the pairs $(m_{i},k_{j})$
satisfy (\hyperref[XVIII.par1.5.10.4]{1.5.10.4}). There then exist an étale neighborhood $U$ of
$\bar{s}$ and a section $k$ of $\Ll|\bar{f}^{-1}(U)$ such that the
triples $(m_{i}, k_{j}, k)$ satisfy (\hyperref[XVIII.par1.5.10.5]{1.5.10.5}). 
\end{enumerate}
\end{enonce*}

When $k(s)$ is infinite, one could simply take $U$ to be a
Zariski neighborhood.

\setcounter{subsubsection}{11}
\subsubsection{}\etiquette[1.5.12]{XVIII.subsubsec1.5.12}
Now let $\Ll$ be an invertible sheaf on $\bar{X}$
rigidified along $Y$, and let $K$ be a $G$-torsor on $X$, where $G$ is
as in \hyperref[XVIII.lem1.5.7]{1.5.7}. We propose to define a $G$-torsor
$\langle \Ll,K]$ on $S$, generalizing the one studied in \hyperref[XVIII.subsec1.3]{1.3}, and
additive in $\Ll$.

Let $(m_{1},m_{2})$ be a pair of sections of $\Ll$ that
satisfies (\hyperref[XVIII.par1.5.10.4]{1.5.10.4}). With the notation of
\hyperref[XVIII.par1.5.10.3]{1.5.10.3} and \hyperref[XVIII.par1.5.10.4]{1.5.10.4}, the
trace, from $\lambda m_{1}+(1-\lambda)m_{2}=0$ to $\EE^{1}_{S}$,
of the torsor ${g'}^{\ast}K$ is a torsor $K_{1,2}$ on
$\EE^{1}_{S}$. Its pullback by the section $\lambda =
0$ (\resp $\lambda = 1$) is the torsor $K_{2}=\Tr_{Z(m_{2})/S}(K)$
(\resp $K_{1}=\Tr_{Z(m_{1})/S}(K)$). By \hyperref[XVIII.lem1.5.7]{1.5.7}, the torsor
$K_{1,2}$ is canonically the pullback of a torsor under $G$
on $S$; hence there is a canonical isomorphism between $K_{1}$ and $K_{2}$.

If $m_{1}$, $m_{2}$, and $k$ are three sections of $\Ll$ such
that the pairs $(m_{i},k)$ satisfy (\hyperref[XVIII.par1.5.10.4]{1.5.10.4}), composition again
gives an isomorphism
\begin{equation*}
\Psi_{k}:K_{1}=\Tr_{Z(m_{1})/S}(K)\xrightarrow{\sim}
K_{2}=\Tr_{Z(m_{2})/S}(K){\quoi}.\tag{1.5.12.1}
\end{equation*}\etiquette[1.5.12.1]{XVIII.eq1.5.12.1}

If $E$ is a relative divisor on $X$, finite over $S$, the sections
$m'_{i}=m_{i}\otimes 1$ of $\Ll\otimes \Oo(E)$ still satisfy
(\hyperref[XVIII.par1.5.10.3]{1.5.10.3}), and we have (for $i=1,2$)
\begin{equation*}
K'_{i}=\Tr_{Z(m'_{i})/S}(K)\simeq K_{i}+\Tr_{E/S}(K)\tag{1.5.12.2}
\end{equation*}\etiquette[1.5.12.2]{XVIII.eq1.5.12.2}%
and the isomorphism $\Psi_{k\otimes 1}:K'_{1}\rightarrow K'_{2}$ is
deduced from $\Psi_{k}$:
\begin{equation*}
\Psi_{k\otimes 1}=\Psi_{k}+\id_{\Tr_{E/S}}(K){\quoi}.\tag{1.5.12.3}
\end{equation*}\etiquette[1.5.12.3]{XVIII.eq1.5.12.3}

To verify that $\Psi_{k}$ does not depend on $k$, it suffices to do so
in an étale neighborhood of every point $s$ of $S$, and, by \hyperref[XVIII.eq1.5.12.3]{1.5.12.3}
and \hyperref[XVIII.lem1.5.9]{1.5.9}, we\pageoriginale reduce to the case in which
conditions (\hyperref[XVIII.par1.5.10.1]{1.5.10.1}) and (\hyperref[XVIII.par1.5.10.2]{1.5.10.2}) hold. With
the notation of \hyperref[XVIII.lem1.5.11]{1.5.11} (II),
$\Psi_{\lambda k_{i}+(1-\lambda)k}-\Psi_{k_{i}}$ is a morphism from
$\EE^{1}_{S}$ to $G$, zero at $\lambda =1$, and hence zero; thus
$\Psi_{k_{i}}=\Psi_{k}$, and therefore $\Psi_{k_{1}}=\Psi_{k_{2}}$.

Let $m_{1}$ and $m_{2}$ be two sections of $\Ll$ satisfying
\hyperref[XVIII.par1.5.10.3]{1.5.10.3}, and let $K_{i}=\Tr_{Z(m_{i})/S}(K)$. To define
an isomorphism
$\Psi_{m_{2},m_{1}}:K_{1}\rightarrow K_{2}$, it suffices to do so in an
étale neighborhood of every point $s$ of $S$, compatibly with
base change. If $E$ is a relative divisor on $X$, finite over
$S$, such that $\Ll\otimes \Oo(E)$ satisfies (\hyperref[XVIII.par1.5.10.1]{1.5.10.1})
and (\hyperref[XVIII.par1.5.10.2]{1.5.10.2}) (see \hyperref[XVIII.lem1.5.9]{1.5.9}), and if the
pairs $(m_{i}\otimes 1,k)$
satisfy (\hyperref[XVIII.par1.5.10.4]{1.5.10.4}), we set (\cf \eqhyperref[XVIII.eq1.5.12.3]{1.5.12.3})
$$
\Psi_{m_{2},m_{1}}+\Tr_{E/S}(K)=\Psi_{k}{\quoi}.
$$
This defines $\Psi_{m_{2},m_{1}}$, and one verifies using
(\hyperref[XVIII.lem1.5.11]{1.5.11}. (I)) that $\Psi_{m_{3},
  m_{2}}\Psi_{m_{2},m_{1}}=\Psi_{m_{3},m_{1}}$. These
constructions are compatible with base changes. If there exists
a section $m$ of $\Ll$ satisfying (\hyperref[XVIII.par1.5.10.3]{1.5.10.3}), we set
$$
\langle \Ll,K]=\Tr_{Z(m)/S}(K)
$$
and, by what precedes, this torsor is independent, up to canonical
isomorphism, of the choice of $m$.

In the general case, note that to define
$\langle \Ll,K]$ it suffices to define it in an étale neighborhood of every point $s$
of $S$, compatibly with base change. If, as above,
$\Ll\otimes \Oo(E)$ satisfies (\hyperref[XVIII.par1.5.10.1]{1.5.10.1})
and (\hyperref[XVIII.par1.5.10.2]{1.5.10.2}),
then $\Ll\otimes \Oo(E)$ has sections satisfying
(\hyperref[XVIII.par1.5.10.3]{1.5.10.3}), and we set
$$
\langle \Ll,K]=\langle \Ll\otimes \Oo(E), K]-\Tr_{E/S}(K){\quoi}.
$$

\subsubsection{}\etiquette[1.5.13]{XVIII.subsubsec1.5.13}
Under the hypotheses of \hyperref[XVIII.subsubsec1.5.8]{1.5.8}, if $x$ and $y$ are
two sections of
$\bar{f}_{\ast}\GG_{m}(\mathfrak{m})$, then $\lambda x+(1-\lambda)y$ is
again a section for $\lambda$ in an open subset with nonempty fibers of
the affine line over $S$.

This can be expressed by saying that

\paragraph{}\etiquette[1.5.13.1]{XVIII.par1.5.13.1}
``$\bar{f}_{\ast}\GG_{m}(\mathfrak{m})$ is path-connected.''

Every\pageoriginale morphism from an open subset with nonempty fibers of the
affine line to $G$ is constant. There is therefore no
nontrivial homomorphism from $\bar{f}_{\ast}\GG_{m}(\mathfrak{m})$ to
$G$, and the automorphisms of $\Ll$ act trivially on
$\langle \Ll,K]$. The torsor $\langle \Ll,K]$ thus depends only on
$K$ and on the section of $R^{1}\bar{f}_{\ast}\GG_{m}(\mathfrak{m})$ defined
by $\Ll$. By globalization, this makes it possible to define a
symbol $\langle \lambda,K]$ for $\lambda$ a section of
$R^{1}\bar{f}_{\ast}\GG_{m}(\mathfrak{m})$. The torsor $\langle \lambda,K]$ is
additive in $\lambda$, and therefore corresponds to an extension $e(K)$ of
$R^{1}\bar{f}_{\ast}\GG_{m}(\mathfrak{m})$ by $G$ (see \hyperref[XVIII.prop1.4.23]{1.4.23}
or \SGA 7 VII 1.1.6 and 1.2).

The formation of $e(K)$ is compatible with every base change and with
addition of torsors, and hence gives an additive functor
$$
e:\bar{f}_{\ast}\ch(G[1])\longrightarrow
\EXT(R^{1}\bar{f}_{\ast}\GG_{m}(\mathfrak{m}),G){\quoi}.
$$

Let $j$ be the canonical map from $X$ into
$R^{1}\bar{f}_{\ast}\GG_{m}(\mathfrak{m})$, $j:t\rightarrow$ (class of
$\Oo(t)$). Pullback by $j$ is an additive functor
$$
j^{\ast}:\EXT(R^{1}\bar{f}_{\ast}\GG_{m}(\mathfrak{m}),G)\longrightarrow
\bar{f}_{\ast}\ch(G[1])
$$
and trivially $j^{\ast}e\sim \Id$. One verifies as in
\hyperref[XVIII.thm1.3.10]{1.3.10} and \hyperref[XVIII.thm1.5.2]{1.5.2} that, for the fppf topology,

\begin{enonce*}{Proposition 1.5.14}\etiquette[1.5.14]{XVIII.prop1.5.14}
Under the hypotheses of \hyperref[XVIII.subsubsec1.5.8]{1.5.8}, the functors $e$ and
$j$ of \hyperref[XVIII.subsubsec1.5.13]{1.5.13} are
quasi-inverse equivalences. It follows from
\hyperref[XVIII.const1.4.19]{1.4.19} and \hyperref[XVIII.prop1.4.23]{1.4.23} that there is an isomorphism
$$
\Ext^{1}(R^{1}\bar{f}_{\ast}\GG_{m}(\mathfrak{m}),G)\xrightarrow{\sim}
H^{1}(X,G){\quoi}.
$$
\end{enonce*}

\setcounter{subsubsection}{14}
\subsubsection{}\etiquette[1.5.15]{XVIII.subsubsec1.5.15}
Let $f_{i}:X_{i}\rightarrow S$ $(i=1,2)$ be two compactified curves
as in \hyperref[XVIII.subsubsec1.5.8]{1.5.8}.

Let $\bar{u}:\bar{X}_{1}\rightarrow \bar{X}_{2}$ be a morphism such that
$Y_{1}$ is a subscheme of the inverse image of $Y_{2}$ and which
induces a finite flat morphism from $X_{1}$ to $X_{2}$. Then,
for $\Ll$ an invertible sheaf on $\bar{X}_{2}$ rigidified along $Y_2$
and $K$ a torsor on $X_{1}$, we have
$$
\langle \Ll,\tr_{X_{1}/X_{2}}(K)]\xrightarrow{\sim}
\langle \bar{u}^{\ast}\Ll,K]{\quoi},
$$
and\pageoriginale we deduce the commutativity of the diagram
$$\vcenter{\etiquette[1.5.15.1]{XVIII.eq1.5.15.1}
\xymatrix@R=1.75cm{
\Ext^{1}(R^1\bar{f}_{1\ast} \GG_m (\mathfrak{m}_1), G) 
\ar[r]^-{\sim}\ar[d]^{\Ext^{1}(\bar{u}^{\ast},\Id)} &   H^{1}(X_1,G) \ar[d]^{\Tr_{X_1 /X_2}}\\
\Ext^1(R^1\bar{f}_{2\ast}\GG_m (\mathfrak{m}_2), G)  \ar[r]^-{\sim}&
H^1(X_2,G).
}}\leqno{(1.5.15.1)}
$$

If $\bar{u}:\bar{X}_{1}\rightarrow \bar{X}_{2}$ is finite and flat, and
satisfies $Y_{1}=\bar{u}^{-1}Y_{2}$ (as schemes), then, for
$\Ll$ an invertible sheaf on $X_{1}$ rigidified along $Y_1$
and $K$ a torsor on $X_{2}$, we have
$$
\langle \Ll,u^{\ast}K]\xrightarrow{\sim}
\langle N_{\bar{X}_{1}/\bar{X}_{2}}\Ll,K]
$$
and we deduce the commutativity of the diagram
\begin{equation*}
\vcenter{
\xymatrix@R=1.75cm@C=1.5cm{
\Ext^{1}(R^{1}\bar{f}_{2\ast} \GG_m (\mathfrak{m}_2) , G)
\ar[d]^{\Ext^{1}(N_{\bar{X}_1 /\bar{X}_2}, \Id)}\ar[r]^-{\sim} &
H^{1}(X_2,G) \ar[d]^{u^{\ast}}\\
\Ext^{1}(R^{1}\bar{f}_{1\ast} \GG_m (\mathfrak{m}_1), G)  \ar[r]^-{\sim} &
H^{1}(X_1, G).
}}\tag{1.5.15.2}
\end{equation*}\etiquette[1.5.15.2]{XVIII.eq1.5.15.2}

One could of course state this result more precisely in terms of Picard
stacks.

\subsection{A vanishing theorem}\etiquette[1.6]{XVIII.subsec1.6}

The present section, from \hyperref[XVIII.lem1.6.6]{1.6.6} (2\textsuperscript{nd}
proof) to the end,
can be read independently of \hyperref[XVIII.sec2]{sections 2} to 5.

\begin{enonce*}{Lemma 1.6.1}\etiquette[1.6.1]{XVIII.lem1.6.1}
Let $X$ be a scheme and $Y$ a closed subscheme of $X$ with complement
$U$, defined by an ideal $\mathfrak{m}$.
$$
U\mathop{\hookrightarrow}^{j} X\mathop{\hookleftarrow}^{i} Y{\quoi}.
$$
Let $\GG_{m}(\mathfrak{m})$ be the \emph{fpqc} sheaf whose sections
on an $X$-scheme $X_{1}$ are the sections of $\Oo^{\ast}_{X_{1}}$
equal to $1$ on the inverse image of $Y$. For $n\geq 1$ invertible
on $X$, the sequence
$$
0\rightarrow j_{!}𝛍_{n}\rightarrow
\GG_{m}(\mathfrak{m})\xrightarrow{n} \GG_{m}(\mathfrak{m})\rightarrow
0
$$
is an exact sequence of sheaves on the big étale site
\emph{(\hyperrefext[XVII][XVII.subsec0.10]{0.10})} of $X$.
\end{enonce*}

Let\pageoriginale $i$ be the morphism of big étale sites
$i:Y\rightarrow X$ induced by $i$:
$$
i^{\ast}(X_{1}/X)=\fprod{X_{1}}{Y}{X}{\quoi}.
$$
The functor $i_{\ast}$ is exact. In the diagram
$$\vcenter{\etiquette[1.6.1.1]{XVIII.eq1.6.1.1}
\xymatrix{
             &    0  \ar[d]           &  0  \ar[d]                   &    0\ar[d]    \\
0  \ar[r]    &j_!𝛍_n  \ar[r]\ar[d]  & \GG_m(\mathfrak{m})\ar[r]\ar[d]  &
             \GG_m(\mathfrak{m}) \ar[r] \ar[d] &0\\
0  \ar[r]     &𝛍_n  \ar[r]  \ar[d]  & \GG_m  \ar[r] \ar[d]  &
  \GG_m  \ar[r] \ar[d]  & 0\\
0  \ar[r]    & i_{\ast}𝛍_n  \ar[d] \ar[r] &
  i_{\ast}\GG_m  \ar[r]\ar[d]  & i_{\ast}\GG_m   \ar[r]\ar[d] & 0\\
             &    0             &  0                    &    0   \\
}}\leqno{(1.6.1.1)}
$$
the columns are exact, as are the last two rows
(\hyperrefext[IX][IX.Thr3.2]{3.2}). The first row is therefore exact.

\setcounter{subsubsection}{1}
\subsubsection{}\etiquette[1.6.2]{XVIII.subsubsec1.6.2}
Let $k$ be a field and
$$
G'_{1}\longrightarrow G'_{2}\longrightarrow G\longrightarrow
G''_{2}\longrightarrow G''_{1}
$$
an exact sequence of abelian sheaves on the big fppf site of
$\Spec (k)$. Recall that if $G'_{1}$ and $G'_{2}$ are
representable and of finite type, and $G''_{2}$ and $G''_{1}$ are representable
and locally of finite type, then $G$ is representable and locally of
finite type\acorriger{In the case where $G'_1=G'_2=0$, this is clear. In the
case where $G''_1=G''_2=0$, using \SGA 3 V 10.1.1 (new edition),
one can reduce to the case where $G'_1\to G'_2$ is a
monomorphism, which is treated in (\SGA 3 VI${}_{\text{A}}$ 3.2).
We are thus reduced to the case where $G'_1=G''_1=0$. One can then regard the
sheaf $G$ as a torsor under $G'_2$ over the scheme $G''_2$. If
$G'_2$ is assumed affine (the situation to which one is reduced,
by (\SGA 3 VI${}_{\text{B}}$ 11.17), if before making the reductions
$G'_1$ and $G'_2$ were assumed affine), then the representability of $G$
follows immediately from faithfully flat descent. This special case is
sufficient for the applications made in this paragraph. If $G'_2$ is
quasi-projective, the representability of $G$ can be obtained by using
(\SGA 1 VIII 7.7) and arguments similar to
those of \cite[Lemma~3.3.1]{sga4.exp.XVIII.bib.4bis}. In fact, every group scheme $G$
of finite type over a field is quasi-projective: one may suppose that $G$
is connected; this is then \cite[Corollary~1.2]{sga4.exp.XVIII.bib.0bis} if $G$ is smooth, and (M.
Raynaud) the general case reduces to it because, by (\SGA 3
VII${}_{\text{A}}$ 8.3), there exists
a finite morphism from $G$ to a smooth group scheme.}.

If $P$ is a connected commutative algebraic group of finite type over
$k$, and $n$ is an integer invertible in $k$, the sequence
\begin{equation*}
0\rightarrow {}_{n}P\rightarrow P\xrightarrow{n} P\rightarrow 0\tag{1.6.2.1}
\end{equation*}\etiquette[1.6.2.1]{XVIII.eq1.6.2.1}%
is exact. If $G$ is a finite étale commutative group over $k$, killed
by $n$, the first group and the last arrow of the exact sequence
$$
\Hom(P,G)\rightarrow \Hom({}_{n}P,G)\xrightarrow{j}
\Ext^{1}(P,G)\xrightarrow{n} \Ext^{1}(P,G)
$$
are zero, whence an isomorphism
\begin{equation*}
j:\Hom({}_{n}P,G)\xrightarrow{\sim} \Ext^{1}(P,G){\quoi}.\tag{1.6.2.2}

% END INCLUDED SGA4_English_current_buildable_source_20260906/source/18/drafts/SGA4_XVIII_LINES_3251_3650_English.texfrag

% BEGIN INCLUDED SGA4_English_current_buildable_source_20260906/source/18/drafts/SGA4_XVIII_LINES_3651_4050_English.texfrag
\end{equation*}\etiquette[1.6.2.2]{XVIII.eq1.6.2.2}

\subsubsection{}\etiquette[1.6.3]{XVIII.subsubsec1.6.3}
Under\pageoriginale the hypotheses of \hyperref[XVIII.lem1.6.1]{1.6.1}, suppose
that $X$ is proper over the spectrum of a field $k$
% SGA-SRC-001141
$$
p:X\longrightarrow \Spec(k){\quoi}.
$$
The exact sequence
$$
0\longrightarrow \GG_{m}(\mathfrak{m})\longrightarrow
\GG_{m}\longrightarrow i_{\ast}\GG_{m}\longrightarrow 0
$$
defines an exact sequence of fppf sheaf cohomology on $\Spec
(k)$, which is easily interpreted directly as
\begin{equation*}
 0\rightarrow p_{\ast}\GG_{m}(\mathfrak{m})\rightarrow
p_{\ast}\GG_{m_{X}}\rightarrow
(pi)_{\ast}\GG_{m_{Y}}\xrightarrow{j}\tag{1.6.3.1}
 R^{1}p_{\ast}\GG_{m}(\mathfrak{m})\rightarrow
\sheafPic_{X/k}\rightarrow \sheafPic_{Y/k}{\quoi}.
\end{equation*}\etiquette[1.6.3.1]{XVIII.eq1.6.3.1}

The sheaves $p_{\ast}\GG_{m_{X}}$ and
$(pi)_{\ast}\GG_{m_{Y}}$ are representable by group schemes of finite
type\sisi{}{ and affine}. The sheaves $\sheafPic_{X/k}$ and
$\sheafPic_{Y/k}$ are representable by group schemes locally of finite
type. By \hyperref[XVIII.subsubsec1.6.2]{1.6.2}, the sheaves
$p_{\ast}\GG_{m}(\mathfrak{m})$ and
$R^{1}p_{\ast}\GG_{m}(\mathfrak{m})$ are representable. By
\eqhyperref[XVIII.par1.5.13.1]{1.5.13.1}, the former is represented by a connected group scheme
(of finite type); the latter is locally of finite type.

\begin{enonce*}{Definition 1.6.4}\etiquette[1.6.4]{XVIII.def1.6.4}
Under the preceding hypotheses, we denote by
$\sheafPic_{\mathfrak{m},X/k}$ the group scheme representing
$R^{1}p_{\ast}\GG_{m}(\mathfrak{m})$. We denote by
$\sheafPic^{0}_{\mathfrak{m},X/k}$ its identity component and by
$\sheafPic^{\tau}_{\mathfrak{m},X/k}$ the inverse image in
$\sheafPic_{\mathfrak{m},X/k}$ of the torsion subgroup of
$\sheafPic_{\mathfrak{m},X/k}/\sheafPic^{0}_{\mathfrak{m},X/k}$.
\end{enonce*}

The scheme $\sheafPic_{\mathfrak{m},X/k}$ represents the fppf sheaf generated by
the presheaf which associates with every scheme $S$ over $k$ the set
of isomorphism classes of invertible modules on $\fprod{S}{X}{k}$,
trivialized along $\fprod{S}{Y}{k}$ (\hyperref[XVIII.subsubsec1.5.8]{1.5.8}).

The Kummer exact sequence \hyperref[XVIII.lem1.6.1]{1.6.1} gives, for $k$
algebraically closed
and $n$ invertible in $k$, an isomorphism
\begin{equation*}
H^{1}_{c}(U,𝛍_{n})\xrightarrow{\sim}
{}_{n}\sheafPic^{\tau}_{\mathfrak{m},X/k}(k){\quoi}.\tag{1.6.4.1}
\end{equation*}\etiquette[1.6.4.1]{XVIII.eq1.6.4.1}

\setcounter{subsubsection}{4}
\subsubsection{}\etiquette[1.6.5]{XVIII.subsubsec1.6.5}
Let\pageoriginale $X$ be a smooth curve over an algebraically closed
field $k$, let $\bar{X}$ be the smooth projective curve containing $X$ as a
dense open subscheme, let $Y$ be a closed subscheme of $\bar{X}$ such that $X=\bar{X}-Y$,
let $\mathfrak{m}$ be the ideal defining $Y$, let $n$ be an integer invertible in $k$, and
let $G$ be a finite abelian group killed by $n$.

The group $\sheafPic_{\mathfrak{m},X/k}$ is noncanonically the product of the
subgroup $\sheafPic^{0}_{\mathfrak{m},\bar{X}/k}$ and a free abelian
group. By \eqhyperref[XVIII.eq1.6.4.1]{1.6.4.1}, \eqhyperref[XVIII.eq1.6.2.2]{1.6.2.2}, and
\hyperref[XVIII.prop1.5.14]{1.5.14}, we therefore have a canonical composite isomorphism
\begin{gather*}
\Hom(H^{1}_{c}(X,𝛍_{n}),G)\xleftarrow{\sim}
\Hom({}_{n}\sheafPic^{\tau}_{\mathfrak{m},\bar{X}/k}(k),G)=\Hom({}_{n}\sheafPic^{0}_{\mathfrak{m},\bar{X}/k},G)\\
\xrightarrow{\sim}
\Ext^{1}(\sheafPic^{0}_{\mathfrak{m},\bar{X}/k},G)\xleftarrow{\sim}
\Ext^{1}(\sheafPic_{\mathfrak{m},\bar{X}/k},G)\xrightarrow{\sim} H^{1}(X,G)
\end{gather*}

If $u:X\rightarrow Y$ is a finite morphism between smooth curves over
$k$, we deduce from \hyperrefext[XVII][XVII.subsec6.3]{6.3}. and \eqhyperref[XVIII.eq1.5.15.2]{1.5.15.2} the commutativity of the
diagram
\begin{equation*}
\vcenter{
\xymatrix@C=2cm@R=1.5cm{
\Hom(H^{1}_c(Y,𝛍_n),G)  \ar[r]^(.55){\sim} \ar[d]^{\Hom(\Tr_u,G)} &
H^{1}(Y,G)               \ar[d]^{u^{\ast}}\\
\Hom(H^1_c(X,𝛍_n),G)    \ar[r]^(.55){\sim} & H^{1}(X,G).
}}\tag{1.6.5.1}
\end{equation*}\etiquette[1.6.5.1]{XVIII.eq1.6.5.1}

\begin{enonce*}{Lemma 1.6.6}\etiquette[1.6.6]{XVIII.lem1.6.6}
Let $X$ be a smooth curve over an algebraically closed field $k$, let $n$
be an integer invertible in $k$, and let $u:X'\rightarrow X$ be a tame
étale covering. The morphism
$$
\Tr_{u}:H^{1}_{c}(X',\ZZ/n)\longrightarrow
H^{1}_{c}(X,\ZZ/n)
$$
is then (noncanonically\nde{The isomorphism is canonical if an
isomorphism $𝛍_n\simeq \mathbf{Z}/n$ has been chosen.})
isomorphic to the transpose of the morphism
$$
u^{\ast}:H^{1}(X,\ZZ/n)\longrightarrow
H^{1}(X',\ZZ/n){\quoi}.
$$
\emph{1\textsuperscript{st} proof.} Since $𝛍_{n}$ is isomorphic to
$\ZZ/n$, \hyperref[XVIII.lem1.6.6]{1.6.6} is the special case of \eqhyperref[XVIII.eq1.6.5.1]{1.6.5.1} with
$G=\ZZ/n$.
\end{enonce*}

\noindent
\emph{2\textsuperscript{nd} proof.}
If $k$ is the field $\CC$ of complex numbers, then, by the
comparison theorems, the morphisms $\Tr_{u}$ and $u^{\ast}$
are identified with the analogous morphisms defined by the continuous map
of topological spaces $u:X'(\CC)\rightarrow
X(\CC)$:
\begin{align*}
\Tr_{u} &: H^{1}_{c}(X'(\CC),\ZZ/n) \longrightarrow
H^{1}_{c}(X(\CC), \ZZ/n)\quad \text{and}\\
u^{\ast} &: H^{1}(X(\CC),\ZZ/n)\longrightarrow H^{1}(X'(\CC),\ZZ/n){\quoi}.
\end{align*}
These\pageoriginale two homomorphisms are transposes of one another
under Poincaré duality.

By the Lefschetz principle, the assertion \hyperref[XVIII.lem1.6.6]{1.6.6} remains
true for $k$ of characteristic $0$ (the cohomology groups
under consideration are indeed invariant under extension of the algebraically closed
base field).

If $k$ has characteristic $p>0$, let $W(k)$ be the ring of Witt vectors
over $k$. Let $\bar{X}$ be the smooth projective completion
of $X$. By (\SGA 1 III 7.4), the curve $\bar{X}$ lifts to a
smooth projective curve $\bar{X}_{1}$ over $W(k)$. Since $\bar{X}$
is smooth, every point $s_{i}\in S=\bar{X}-X$ lifts to a
section $s'_{i}$ of $\bar{X}_{1}$ over $W(k)$; set
$S_{1}=⋃_{i}s'_{i}(\Spec(W(k)))$ and $X_{1}=\bar{X}_{1}-S_{1}$.

Recall (\SGA 1 XIII 2.4)
that the covering $X'$ of $X$ lifts to an
étale covering $X'_{1}$ of $X_{1}$, and that $X'_{1}$ is obtained from a
smooth projective curve $\bar{X}'_{1}$ over $W(k)$ by removing
the union $S'_{1}$ of finitely many disjoint sections.
$$
\xymatrix@=.5cm{
                  &     X_1 \ar[ddddd]^{f_1}       &                 &
                  X\ar[ll]\ar[ddddd]\\
 & & &\\
X'_1 \ar[ruu]_{u_1}\ar[rddd]_{f'_1}&                & X' \ar[ll]\ar[ruu]_{u} \ar[rddd] &   \\
& & &\\
& & &\\
                  &    \Spec(W(k)) &                 &\Spec(k).\ar[ll]
}
$$

The étale sheaves $R^{1}f_{1\ast}(\ZZ/n)$,
$R^{1}f'_{1\ast}(\ZZ/n)$, $R^{1}f_{1!}(\ZZ/n)$, and
$R^{1}f'_{1!}(\ZZ/n)$ are locally constant and their formation is
compatible with base change,\pageoriginale as follows from
the exact sequences relating the cohomology of $\bar{X}_{1}$, $X_{1}$,
and $S_{1}$, or of $\bar{X}'_{1}$, $X'_{1}$, and $S'_{1}$.

The special or generic geometric fibers of the morphisms
$$
\Tr_{u_{1}}:R^{1}f'_{1!}\ZZ/n\longrightarrow
R^{1}f_{1!}\ZZ/n
$$
or
$$
u^{\ast}_{1}:R^{1}f_{1\ast}\ZZ/n\longrightarrow
R^{1}f'_{1\ast}\ZZ/n
$$
are therefore isomorphic, and this reduces us to the case already treated in which $k$
has characteristic zero.

\begin{enonce*}{Lemma 1.6.7}\etiquette[1.6.7]{XVIII.lem1.6.7}
Let $X$ be a connected smooth curve over an algebraically closed field
$k$ of characteristic exponent $p$, and let $n$ be an integer prime to $p$. There
exists a principal covering \emph{(= finite étale Galois
(surjective))} $u:X'\rightarrow X$, of degree dividing a power of
$n$, such that the trace morphism
$$
\Tr_{u}:H^{1}_{c}(X', \ZZ/n)\longrightarrow
H^{1}_{c}(X,\ZZ/n)
$$
is zero.
\end{enonce*}

Recall that the group $H^{1}(X,\ZZ/n)$ is finite (\hyperrefext[IX][IX.thm4.6]{4.6} or
\hyperrefext[XVI][XVI.thm5.2]{5.2}). Let $u:X'\rightarrow X$ be the largest connected abelian principal
covering of $X$ whose Galois group is killed by $n$ (its Galois
group is denoted by $H_{1}(X,\ZZ/n)$). By construction, the
homomorphism
$$
u^{\ast}:H^{1}(X,\ZZ/n)\longrightarrow H^{1}(X', \ZZ/n)
$$
is zero. By \hyperref[XVIII.lem1.6.6]{1.6.6}, the homomorphism
$$
\Tr_{u}:H^{1}_{c}(X',\ZZ/n)\longrightarrow
H^{1}_{c}(X,\ZZ/n)
$$
is therefore zero.

Recall the following lemma (\EGA IV 15.6.5):

\begin{enonce*}{Lemma 1.6.8}\etiquette[1.6.8]{XVIII.lem1.6.8}
Let $f:X\rightarrow S$ be a smooth morphism and let $x$ be a section of
$f$. Then there exists a Zariski open subset $U$ of $X$, containing $x$,
such that the geometric fibers of $f|U$ are connected.
\end{enonce*}

\begin{enonce*}{Fundamental Lemma 1.6.9}\etiquette[1.6.9]{XVIII.lemfond1.6.9}
Let\pageoriginale $f:X\rightarrow S$ be a compactifiable smooth curve
(\hyperref[XVIII.def1.1.2]{1.1.2}) \emph{(\hyperrefext[XVII][XVII.subsec3.2]{3.2})},
let $x$ be a geometric point
of $X$, let $s$ be its image in $S$, and let $n\geq 1$ be an integer invertible on
$S$. There exist an étale neighborhood $V$ of $s$ in $S$ and an étale
neighborhood $U$ of $x$ in $f^{-1}(V)=\fprod{V}{X}{S}$
\begin{equation*}
\vcenter{
\xymatrix@R=1.5cm@C=1cm{
x \ar[r]\ar@{-}[d]& U\ar[rd]_{f'}\ar[r]^{u} & f^{-1}(V)\ar[r] \ar[d]^{f_v}& X\ar[d]^{f}\\
s \ar[rr]&   & V\ar[r]         & S,
}}\tag{1.6.9.1}
\end{equation*}\etiquette[1.6.9.1]{XVIII.eq1.6.9.1}%
such that $R^{0}f'_{!}\sheaf{\ZZ}/n=0$, the arrow
\emph{\hyperrefext[XVII][XVII.subsec6.2]{6.2}}
$$
\Tr_{u}:R^{1}f'_{!}\sheaf{\ZZ}/n\longrightarrow
R^{1}f_{v!}\sheaf{\ZZ}/n
$$
is zero, and the trace morphism
$$
\Tr_{f'}:R^{2}f'_{!}\sheaf{\ZZ}/n(1)\longrightarrow
\sheaf{\ZZ}/n
$$
is an isomorphism.
\end{enonce*}

For $R^{0}f'_{!}\sheaf{\ZZ}/n=0$, it suffices that $f'$ be
quasi-affine; for $\Tr_{f'}$ to be an isomorphism, it suffices that
the geometric fibers of $f'$ be connected
(\hyperref[XVIII.lem1.1.9]{1.1.9}); by
\hyperref[XVIII.lem1.6.8]{1.6.8}, these properties of $f'$ are easy to obtain
locally on
$S$ for the étale topology. It remains to prove, for $X$ with geometrically
connected fibers, the existence of a diagram \eqhyperref[XVIII.eq1.6.9.1]{1.6.9.1} for
which the arrow $\Tr_{u}$ under consideration is zero: it will suffice to
shrink $U$ and $V$ so that the other two required properties of
$f'$ hold.

The problem is local on $S$ in a neighborhood of $s$; by \hyperrefext[XVII][XVII.thm5.2.6]{5.2.6}
(base change), we immediately reduce to the case where $S$ is Noetherian. By
passage to the limit, and using \hyperrefext[XVII][XVII.thm5.3.6]{5.3.6} (constructibility), we
reduce to the case where $S$ is strictly local Noetherian, with closed point the image
of $s$.

Let $t$ be a geometric point of $S$. By \hyperref[XVIII.lem1.6.7]{1.6.7}, there
exists a principal covering of degree prime to the characteristic exponent
% SGA-SRC-001143
$p$ of $k(t)$, $u:X^{1}_{t}\rightarrow X_{t}$, such that the arrow
$$
\Tr_{u}:H^{1}_{c}(X^{1}_{t},\ZZ/n)\longrightarrow
H^{1}_{c}(X_{t},\ZZ/n)
$$
is\pageoriginale zero. By the acyclicity theorem \hyperrefext[XV][XV.thm2.1]{2.1} (or
\hyperrefext[XV][XV.lem2.6]{2.6}), there exists an étale neighborhood $v:U\rightarrow X$ of $x$ in
$X$ such that $X^{1}_{t}$ splits over $U_{t}$, so that
$v_{t}:U_{t}\rightarrow X_{t}$ factors through $X^{1}_{t}$. We therefore still have
$$
\Tr_{v_{t}}=0:H^{1}_{c}(U_{t},\ZZ/n)\longrightarrow
H^{1}_{c}(X_{t},\ZZ/n){\quoi}.
$$

For every geometric point $t$ of $S$, there therefore exists an étale
neighborhood of $x$ in $X$
$$
\xymatrix{
U \ar[r]^{v} \ar[rd]_{f'}& X \ar[d]^{f}\\
  & S
}
$$
such that the fiber at $t$ of
\begin{equation*}
\Tr_{v}:R^{1}f'_{!}\ZZ/n\longrightarrow
R^{1}f_{!}\ZZ/n\tag{1.6.9.2}
\end{equation*}\etiquette[1.6.9.2]{XVIII.eq1.6.9.2}%
is zero (\hyperrefext[XVII][XVII.thm5.2.6]{5.2.6}).

As $U$ varies, the images of the arrows \eqhyperref[XVIII.eq1.6.9.2]{1.6.9.2} form a
filtered decreasing system of subsheaves of the constructible sheaf
$R^{1}f_{!}\ZZ/n$, which becomes zero at any given geometric
point of $S$. By Noetherian induction, one verifies that such a
system is always stationary, with stationary value zero. Hence there
exists $U$ such that \eqhyperref[XVIII.eq1.6.9.2]{1.6.9.2} is zero, and this completes the
proof of \hyperref[XVIII.lemfond1.6.9]{1.6.9}.

\section{The trace morphism}\etiquette[2]{XVIII.sec2}\pageoriginale

This \S is devoted to the proof of Theorems
\hyperref[XVIII.thm2.9]{2.9} and \hyperref[XVIII.prop2.12]{2.12}. The proof of
\hyperref[XVIII.thm2.9]{2.9} uses only no.~1 of \S \hyperref[XVIII.sec1]{1}.

\begin{enonce*}{Lemma 2.1}\etiquette[2.1]{XVIII.lem2.1}
Let $f:X\rightarrow S$ be a compactifiable morphism
\emph{(\hyperrefext[XVII][XVII.def3.2.1]{3.2.1})} of relative dimension $\leq d$, and
let $j:U\hookrightarrow X$ be an open subscheme of finite type of $X$, whose
complement has fiberwise dimension $< d$ (for example, $U$ is
fiberwise dense in $X$). If $F$ is a torsion sheaf on
$X$, the canonical arrow
$$
R^{2d}(fj)_{!}j^{\ast}F\xrightarrow{\sim} R^{2d}f_{!}F
$$
is an isomorphism.
\end{enonce*}

Let $i$ be the inclusion of $Y=X-U$ into $X$. The exact sequence
(XVII \hyperref[XVII.eq5.1.16.2]{5.1.16.2}) gives us an exact sequence
$$
R^{2d-1}(fi)_{!}i^{\ast}F\longrightarrow
R^{2d}(fj)_{!}j^{\ast}F\longrightarrow R^{2d}f_{!}F\longrightarrow
R^{2d}(fi)_{!}i^{\ast}F
$$
whose outer terms are zero by (\hyperrefext[XVII][XVII.cor5.2.8.1]{5.2.8.1}).

\begin{enonce*}{Lemma 2.2}\etiquette[2.2]{XVIII.lem2.2}
Let $f:X\rightarrow S$ be a compactifiable morphism of relative
dimension $\leq d$, and let $u_{i}:X_{i}\rightarrow X$ be a family of
separated étale morphisms of finite type; let
$u_{ij}:X_{ij}=\fprod{X_{i}}{X_{j}}{X}\rightarrow X$. For every
torsion sheaf $F$ on $X$, the sequence
$$
\bigoplus_{i,j} R^{2d}(fu_{ij})_{!}u^{\ast}_{ij}F\rightrightarrows
\bigoplus_{i} R^{2d}(fu_{i})_{!}u^{\ast}_{i}F\rightarrow
R^{2d}f_{!}F\rightarrow 0
$$
of arrows \emph{(XVII \hyperref[XVII.eq6.2.7.2]{6.2.7.2})} is exact.
\end{enonce*}

This follows immediately from the right exactness of the functor $R^{2d}f_{!}$ and
from \hyperrefext[XVII][XVII.prop6.2.9]{6.2.9}.

\begin{enonce*}{Lemma 2.3}\etiquette[2.3]{XVIII.lem2.3}
Let $f:X\rightarrow S$ be a flat morphism of finite
presentation. The open subset of $X$ on which $f$ is Cohen--Macaulay \emph{(\EGA
IV 6.8.1 and 12.2.1 (vii)} is relatively of finite presentation and
fiberwise dense.
\end{enonce*}

The\pageoriginale question is local on $S$, which we reduce to assuming
Noetherian, whence the 1\textsuperscript{st} assertion. The second is checked
fiberwise, hence with $S$ the spectrum of a field. The local rings
of $X$ at the maximal points of $X$ are Artinian, hence
Cohen--Macaulay, which proves the assertion since ``Cohen--Macaulay'' is
an ``open'' property.

\begin{enonce*}{Lemma 2.4}\etiquette[2.4]{XVIII.lem2.4}
Let $f:X\rightarrow S$ be a flat morphism of finite presentation and
Cohen--Macaulay, with $S$ quasi-separated. Every point of $X$ has a
(Zariski) open neighborhood $U$ such that there exists a quasi-finite
flat $S$-morphism of finite presentation from $U$ to the standard vector bundle
$\EE^{d}_{S}$.
\end{enonce*}

This is a special case of \EGA IV 15.4.3.

\begin{enonce*}{Lemma 2.5}\etiquette[2.5]{XVIII.lem2.5}
Let $\underline{u}=(u_{1},\ldots,u_{d})$ and $\underline{v}=(v_{1},\ldots,v_{d})$
be two systems of parameters of a Cohen--Macaulay local ring $A$
of dimension $d$ and maximal ideal $\mathfrak{m}$. There exists a sequence
$\underline{w}_{0},\ldots,\underline{w}_{n}$ of systems of parameters of $A$ such that
$\underline{w}_{0}=\underline{u}$, $\underline{w}_{n}=\underline{v}$, and
each $\underline{w}_{i+1}$ is obtained from $\underline{w}_{i}$ by
changing only one of the parameters.
\end{enonce*}

We argue by induction on $d$, the assertion being vacuous for $d<2$;
we therefore suppose $d\geq 2$. For an element $x\in \mathfrak{m}$ to be
$A/u_{1}$-regular (\resp $A/v_{1}$-regular), it is necessary and
sufficient that $x$ belong to no member of $\Ass(A/u_{1})$ (\resp
$\Ass(A/v_{1})$); since no finite union of prime ideals $\neq \mathfrak{m}$
is equal to $\mathfrak{m}$, there exists $x\in \mathfrak{m}$ that is both $A/u_{1}$- and $A/v_{1}$-regular, and hence there also exist systems of parameters
$\underline{a}=(u_{1},x,\ldots)$ and $\underline{b}=(v_{1},x,\ldots)$. We conclude by
applying the induction hypothesis to $\underline{u}-\{u_{1}\}$ and
$\underline{a}-\{u_{1}\}$ in $A/u_{1}$, to $\underline{a}-\{x\}$ and
$\underline{b}-\{x\}$
in $A/x$, and to $\underline{b}-\{v_{1}\}$ and $\underline{v}-\{v_{1}\}$ in $A/v_{1}$.

\begin{enonce*}{Lemma 2.6}\etiquette[2.6]{XVIII.lem2.6}
Let $X$ be a Cohen--Macaulay scheme of finite type over an
algebraically closed field $k$, and let $u$, $v:X\rightarrow \EE^{d}_{k}$
be two quasi-finite flat morphisms from $X$ to a standard vector
bundle. Every closed point $x$ of $X$ has a (Zariski) open
neighborhood $U$ such that there exists a sequence\pageoriginale of morphisms
$w_{0},\ldots,w_{n}:U\rightarrow \EE^{d}_{k}$ such that $w_{0}=u|U$,
$w_{n}=v|U$, and $w_{i+1}$ differs from $w_{i}$ in only one
coordinate.
\end{enonce*}

Suppose first that $u(x)=v(x)=0$. The germs at $x$ of
quasi-finite flat morphisms from $X$ to $\AAA^{d}_{k}$,

% END INCLUDED SGA4_English_current_buildable_source_20260906/source/18/drafts/SGA4_XVIII_LINES_3651_4050_English.texfrag

% BEGIN INCLUDED SGA4_English_current_buildable_source_20260906/source/18/drafts/SGA4_XVIII_LINES_4051_4450_English.texfrag
sending $x$ to $0$, then identify with the systems of parameters of
the local ring $\Oo_{X,x}$, and \hyperref[XVIII.lem2.6]{2.6} follows from
\hyperref[XVIII.lem2.5]{2.5}. In the general
case, one first joins $u$ to $u-u(x)$ (\resp $v$ to $v-v(x)$)
by the chain of morphisms
$$
u-(u_{1}(x),\ldots,u_{i}(x),0,\ldots,0)(\text{\resp}
v-(v_{1}(x),\ldots,v_{i}(x),0,\ldots,0)){\quoi}.
$$

\setcounter{subsection}{6}
\subsection{}\etiquette[2.7]{XVIII.subsec2.7}
Let there be a pair of composable $S$-compactifiable morphisms
$X\xrightarrow{g} Y\xrightarrow{f}S$. Suppose that $f$ has relative
dimension $\leq d$ and $g$ has relative dimension $\leq
e$. For every torsion sheaf $F$ on $X$, the composition spectral
sequence
\begin{equation*}
R^{p}f_{!}R^{q}g_{!}F\Longrightarrow R^{p+q}(fg)_{!}F\tag{2.7.1}
\end{equation*}\etiquette[2.7.1]{XVIII.eq2.7.1}%
then gives us (by the maximum-cycle argument) an isomorphism
\begin{equation*}
R^{2d}f_{!}R^{2e}g_{!}F\simeq R^{2(d+e)}(fg)_{!}F{\quoi}.\tag{2.7.2}
\end{equation*}\etiquette[2.7.2]{XVIII.eq2.7.2}

\subsection{}\etiquette[2.8]{XVIII.subsec2.8}
Let
$$
a^{d}:\EE^{d}_{S}\longrightarrow S
$$
be the standard vector bundle over a coherent (\ie quasi-compact,
quasi-separated) scheme $S$. If $F$ is a torsion sheaf on $S$, prime
to the residual characteristics of $S$, we define as follows, by
induction on $d$, a trace morphism
\begin{equation*}
\Tr_{a^d}:R^{2d}a^{d}_{!}F(d)\xrightarrow{\sim} F{\quoi}.\tag{2.8.1}
\end{equation*}\etiquette[2.8.1]{XVIII.eq2.8.1}

For $d=0$, $\Tr_{a^0}$ is the identity. For $d=1$, $\Tr_{a^1}$
is the isomorphism (\hyperref[XVIII.prop1.1.6]{1.1.6}) (\cf
\hyperref[XVIII.lem1.1.9]{1.1.9}). Finally, there is a canonical isomorphism
$$
\EE^{d+1}_{S}=\fprod{\EE^{1}_{S}}{\EE^{d}_{S}}{S}=\EE^{1}_{\EE_{S}^d}
$$
and\pageoriginale we define $\Tr_{a^{d+1}}$ as the composite of the
isomorphisms $\Tr_{a^{d}_{S}}$ and
$Ra^{d}_{S!}(\Tr_{a^1_{\EE^d_S}})$
(\cf diagram Var 3 of \hyperref[XVIII.thm2.9]{2.9}).

The formation of the source (and the target) of \eqhyperref[XVIII.eq2.8.1]{2.8.1} is compatible with every
base change. Every algebraic group over $S$ acting on
$\EE^{d}_{S}$ therefore acts on these sheaves; a connected algebraic group
necessarily acts trivially. The group of coordinate
permutations, contained in the affine group, therefore acts
trivially on the source (and the target) of \eqhyperref[XVIII.eq2.8.1]{2.8.1}.

\setcounter{subsubsection}{1}
\subsubsection{}\etiquette[2.8.2]{XVIII.subsubsec2.8.2}
The isomorphism \eqhyperref[XVIII.eq2.8.1]{2.8.1} is invariant under permutation of the
coordinates.

This point could also be verified by interpreting
\eqhyperref[XVIII.eq2.8.1]{2.8.1} in terms
of cup products (XVII \hyperref[XVII.eq5.4.2.2]{5.4.2.2} and \hyperrefext[XVII][XVII.lem5.4.3.5]{5.4.3.5}).

\begin{enonce*}{Theorem 2.9}\etiquette[2.9]{XVIII.thm2.9}
Consider triples $(f,d,F)$ consisting of a compactifiable morphism
\emph{(\hyperrefext[XVII][XVII.subsec6.3]{6.3})} $f:X\rightarrow Y$, an integer $d$, and
a torsion sheaf
$F$ on $Y$ prime to the residual characteristics of
$Y$, where the morphism $f$ satisfies the condition $(\ast)_{d}$: there exists
an open subset $U$ of $X$ such that the restriction of $f$ to $U$ is a
flat morphism of finite presentation with fibers of dimension $\leq d$,
and such that the fibers of $X-U$ have dimension $< d$.

There is one and only one way to associate with every
triple $(f,d,F)$ as above a trace morphism
$$
\Tr_{f}:R^{2d}f_{!}f^{\ast}F(d)\longrightarrow F
$$
such that the following conditions are satisfied:
\begin{itemize}
\item[(Var 1)] (Functoriality). The trace morphism is functorial in $F$.

\item[(Var 2)] (Compatibility with base change). For every
Cartesian square
$$
\xymatrix{
X' \ar[r]_{g'}\ar[d]^{f'}& X\ar[d]^{f}\\
Y' \ar[r]_{g} & Y
}
$$
with\pageoriginale $Y'$ quasi-compact and quasi-separated, the diagram
$$
\xymatrix@R=1.25cm{
R^{2d}f'_!f'^{\ast}g^{\ast}F(d) \ar@{=}[r]\ar[d]^{Tr_{f'}}& R^{2d}f'_!
g'^{\ast} f^{\ast} F(d) & g^{\ast}
R^{2d}f_!f^{\ast} F(d) \ar[l]_{\sim} \ar[d]^{g^{\ast}(Tr_f)}\\
g^{\ast}F \ar@{=}[rr]&            & g^{\ast}F
}
$$
in which the horizontal isomorphism is the base-change arrow
\emph{\eqhyperrefext[XVII][XVII.eq5.2.2.1]{5.2.2.1}}, is commutative.

\item[(Var 3)] (Compatibility with composition). For every pair of
composable $Z$-compactifiable morphisms
$$
X\xrightarrow{g}Y\xrightarrow{f}Z{\quoi},
$$
integers $d$ and $e$ such that $f$ satisfies $(\ast)_{d}$ and $g$
satisfies $(\ast)_{e}$, and torsion sheaf $F$ on $Z$, the
composite morphism $fg$ satisfies $(\ast)_{d+e}$ and the diagram
$$
\xymatrix@C=2.25cm{
R^{2d}f_! R^{2e} g_! g^{\ast} f^{\ast}F(d)(e) \ar@{<->}[d]_{\text{\rotatebox{90}{$\backsim$}}} \ar[r]^(.6){R^{2d}f_!(\Tr_g)} &
R^{2d}f_! f^{\ast}F(d) \ar[d]^{\Tr_f}\\
R^{2(d+e)} (fg)_! (fg)^{\ast} F(d+e) \ar[r]^(.65){\Tr_{fg}} & F
}
$$
in which the vertical isomorphism is the arrow \eqhyperref[XVIII.eq2.7.2]{2.7.2}, is
commutative.

\item[(Var 4)] (Normalizations).
\begin{itemize}
\item[(I)] If $f$ is finite locally free of rank $n$ and $d=0$,
the composite morphism
$$
F\longrightarrow
f_{\ast}f^{\ast}F=f_{!}f^{\ast}F\xrightarrow{\tr_{f}}F
$$
is multiplication by $n$.

\item[(II)] If\pageoriginale $f$ is the standard affine line
$f:\EE^{1}_{Y}\rightarrow Y$, and if $d=1$, the trace morphism
$$
\Tr_{f}: R^{2}f_{!}f^{\ast}F(1)\longrightarrow F
$$
is the isomorphism \eqhyperref[XVIII.eq2.8.1]{2.8.1}.
\end{itemize}
\end{itemize}
\end{enonce*}

\begin{enumerate}
\renewcommand{\theenumi}{\alph{enumi}}
\renewcommand{\labelenumi}{\theenumi)}
\item For $d=0$, condition $(\ast)_{d}$ says that $f$ is
quasi-finite and flat of finite presentation. By (\hyperrefext[XVII][XVII.thm6.2.3]{6.2.3}), for
$d=0$, the trace morphism must coincide with the morphism constructed
in \loccit

When $f$ is the projection of standard affine space
$\AAA^{d}_{S}$ onto $S$, it follows from (Var~4)~(II) and (Var~3) that
$\Tr_{f}$ must be the isomorphism \eqhyperref[XVIII.eq2.8.1]{2.8.1}.

\item Let $f:X\rightarrow Y$ be a compactifiable morphism satisfying the
condition $(\ast)_{d}$: there exists a quasi-finite flat
$Y$-morphism of finite presentation $u$ from $X$ to the standard vector bundle
$a^{d}:\EE^{d}_{Y}\rightarrow Y$.

Let $t(f,u)$ be the ``composite'' morphism of the trace morphisms
\eqhyperref[XVIII.eq2.8.1]{2.8.1} and (\SGA \hyperrefext[XVII][XVII.thm6.2.3]{6.2.3}) (\cf (Var 3))
$$
t(f,u)=\Tr_{a^d}\circ R^{2d}a^{d}_{!}(\Tr_{u}){\quoi}.
$$

This morphism $t(f,u)$ satisfies (Var 1) and (Var 2). By base
change, to verify that $t(f,u)$ depends only on $f$ and $d$, and
not on $u$, it suffices to verify this when $Y$ is the spectrum of an
algebraically closed field $k$. By \hyperref[XVIII.lem2.2]{2.2}, the question is
local on $X$,
so by \hyperref[XVIII.lem2.6]{2.6}, it suffices to show that
$t(f,u)=t(f,v)$
when $u$ and $v$ differ in only one coordinate, which, by
\hyperref[XVIII.subsubsec2.8.2]{2.8.2}, may be assumed to be the first. Setting
$E=\EE^{d-1}_{Y}$, we thus have a commutative diagram
$$
\xymatrix{
                              &    &\EE^1_{E} \ar[rdd]_(.3){a} &  &  \\
X \ar[rrrr]_{v}\ar[rru]^{u} \ar[rdd]_{f} \ar[rrrd]_{x}&    &                 &
                              &\EE^{1}_E \ar[ld]_{a}\\
                              &    &                 &E \ar[lld]^w& \\
                              & Y  &                 &  &
}
$$
By\pageoriginale construction (\hyperref[XVIII.subsec2.8]{2.8}),
$$
\Tr_{wa}=\Tr_{w}\circ R^{2(d-1)}w_!(\Tr_{a}){\quoi}\;;
$$
we deduce that $t(f,u)$ is the ``composite'' of $\Tr_{w}$ and
$t(x,u)$:
\begin{align*}
t(f,u) &= \Tr_{wa}\circ R^{2d}(wa)_{!}(\Tr_{u})\\
       &= \Tr_{w}\circ R^{2(d-1)}w_{!}(\Tr_{a})\circ
       R^{2d}(wa)_{!}(\Tr_{u})\\
       &= \Tr_{w}\circ R^{2(d-1)}w_{!}(\Tr_{a}\circ
       R^{2}a_{!}(\Tr_{u})){\quoi}.
\end{align*}
Moreover (\hyperref[XVIII.lem1.1.7]{1.1.7}), $t(x,u)$ is precisely the morphism
$\Tr_{x}$ of \hyperref[XVIII.prop1.1.6]{1.1.6}; hence $t(x,u)=t(x,v)=\Tr_{x}$
and $t(f,u)=t(f,v)$.

We denote by $\Tr_{f}$ the morphism $t(f,u)$ constructed
above. By \hyperrefext[XVII][XVII.thm6.2.3]{6.2.3}, this trace morphism satisfies (Var 1), (Var 2),
(Var 3) for $e=0$, and (Var 4). Whenever it is defined, it necessarily
coincides with the trace morphism sought, if the latter exists.

\item Let $f:X\rightarrow Y$ be a compactifiable Cohen--Macaulay
morphism, flat of finite presentation and pure of relative dimension
$d$. Let $(U_{i})_{i\in I}$ be a cover of $X$ by quasi-compact
open subsets such that there exists a quasi-finite flat
$Y$-morphism from $U_{i}$ to $\EE^{d}_{Y}$ (\hyperref[XVIII.lem2.4]{2.4}). By
\hyperref[XVIII.lem2.2]{2.2} and b), there
is then a unique morphism $\Tr_{f}$ making the following diagram
commutative, where the notation of \hyperref[XVIII.lem2.2]{2.2} is used:
$$
\xymatrix{
\bigoplus_{i,j}
R^{2d}(fu_{ij})_!(fu_{ij})^{\ast} F(d) \ar@<2pt>[r]
\ar@<-2pt>[r]&\bigoplus_{i}
R^{2d}(fu_i)_!(fu_i)^{\ast}F(d)
\ar[rdd]_{\bigoplus_{i}\Tr_{fu_i}}\ar[r]
&
R^{2d}f_!f^{\ast}F(d) \ar[dd]^{\Tr_f}\\
& & \\
& & F
}
$$

This trace morphism does not depend on the chosen open cover
(compare the two with the sum cover); it satisfies the conditions
(Var 1), (Var 2), (Var 3) for $e=0$, and (Var 4), as follows
immediately from b).

\item Let $f:X\rightarrow Y$ be a compactifiable morphism satisfying
$(\ast)_{d}$, and let $j:U\hookrightarrow X$\pageoriginale be the open subset promised
by $(\ast)_{d}$. By \hyperref[XVIII.lem2.3]{2.3}, after replacing $U$
by a smaller open subset, we may suppose that $fj$ is Cohen--Macaulay and
pure of dimension $d$.

\noindent
By \hyperref[XVIII.lem2.1]{2.1}, the canonical morphism
$$
R^{2d}(fj)_{!}(fj)^{\ast}F(d)\xrightarrow{\sim}
R^{2d}f_{!}f^{\ast}F(d)
$$
is an isomorphism, and we define $\Tr_{f}$ to be the composite
$$
\Tr_{f}:R^{2d}f_{!}f^{\ast}F(d)\xleftarrow{\sim}
R^{2d}(fj)_{!}(fj)^{\ast}F(d)\xrightarrow{\Tr_{fj}}F{\quoi}.
$$
It is immediate, by (Var 3) for $e=0$ established in c), that
$\Tr_{f}$ does not depend on the choice of $U$. The trace morphism satisfies
(Var 1), (Var 2), (Var 3) for $e=0$, and (Var 4), and is the only one
that can satisfy the imposed conditions. It remains to prove that it
satisfies (Var 3), and for this purpose the preceding arguments reduce
us to considering only pairs $(f,d)$, $(g,e)$ satisfying
$(\ast)_{d}$ and $(\ast)_{e}$.

The reader who wishes to polish the arguments that follow will first verify the

\begin{enonce*}{Lemma 2.9.1}\etiquette[2.9.1]{XVIII.lem2.9.1}
Let $(f,g,h)$ be three composable $T$-compactifiable morphisms,
satisfying $(\ast)_{k}$, $(\ast)_{\ell}$, and $(\ast)_{m}$.
$$
X\xrightarrow{h}Y\xrightarrow{g}Z\xrightarrow{f}T{\quoi}.
$$
If the pairs $(f,g)$ and $(g,h)$ satisfy {\rm (Var 3)}, then
$(f,gh)$ satisfies {\rm (Var 3)} if and only if $(fg,h)$
satisfies {\rm (Var 3)}.
\end{enonce*}

The condition in the conclusion also says that
$\Tr_{fgh}$ is the ``composite'' of $\Tr_{f}$, $\Tr_{g}$, and $\Tr_{h}$.

\item Let there be a commutative diagram:
$$
\xymatrix@R=1.25cm@C=1.5cm{
X \ar[r]^{v} \ar[rd]_{g}& \EE^{e}_Y \ar[r]^{u'}\ar[d]^{c} & \EE_Z^{d+e}\ar[d]^{b} \\
            &     Y   \ar[r]^{u} \ar[rd]_{f}                       & \EE_Z^{d}\ar[d]^{a}  \\
            &                                  &     Z
}
$$
The\pageoriginale trace morphism relative to $fg$ is the
``composite'' of the trace morphisms relative to $ab$ and $u'v$; by
construction for $ab$, and by \hyperrefext[XVII][XVII.thm6.2.3]{6.2.3} for $u'v$, this is also the ``composite'' of the trace morphisms relative to $a$, $b$, $u'$, and
$v$. On the other hand, the ``composite'' of the trace morphisms relative to
$f$ and $g$ is the ``composite'' of those relative to $a$, $u$, $c$, and
$v$. We must therefore verify that, for every Cartesian diagram
\begin{equation*}
\vcenter{
\xymatrix@C=1.25cm@R=1cm{
\EE^{e}_S  \ar[d]^{b'} \ar[r]^{u'}   & \EE^{e}_T \ar[d]^{b}\\
    S      \ar[r]^{u}     &      T
}}\tag{2.9.2}
\end{equation*}\etiquette[2.9.2]{XVIII.eq2.9.2}%
with $u$ quasi-finite and flat of finite presentation, the diagram of
trace morphisms \eqhyperref[XVIII.eq2.8.1]{2.8.1} and \hyperrefext[XVII][XVII.subsec6.2]{6.2}.)
\begin{equation*}
\vcenter{
\xymatrix@C=2.25cm{
R^{2e}b_!u'_! u'^{\ast}b {}^{\ast} F(e) \ar[r]^(.55){R^{2e}b_!
  (\Tr_{u'})}\ar@{=}[d]& R^{2e}b_! b^{\ast} F(e)\ar[r]^{\Tr_b}
&  F \ar@{=}[d]\\
u_! R^{2e}b'_! b'^{\ast} u^{\ast} F(e)\ar[r]^(.55){u_!(\Tr_{b'})} & u_!
u^{\ast} F \ar[r]^{\Tr_u}& F
}}\tag{2.9.3}
\end{equation*}\etiquette[2.9.3]{XVIII.eq2.9.3}%
is commutative.
\end{enumerate}

By base change, it suffices to verify this when $T$ is the spectrum of an
algebraically closed field $k$. If $S$ is a sum of schemes $S_{i}$, we
are reduced to verifying the commutativity of \eqhyperref[XVIII.eq2.9.3]{2.9.3} only for the
diagrams \eqhyperref[XVIII.eq2.9.2]{2.9.2} with base $u_{i}:S_{i}\rightarrow T$. This makes it
possible to suppose that $S$ is the spectrum of a local Artinian
$k$-algebra $A$, of finite degree $n=[A:k]$. If we then identify sheaves on
$T$ (\resp on $\EE^{e}_{T}$) with sheaves on $S$ (\resp
on $\EE^{e}_{S}$) by the functor $u^{\ast}$ (\resp
${u'}^{\ast}$) (VIII \hyperref[VIII.thm1.1]{1.1}), the morphism $\Tr_{b}$ identifies with the
morphism $\Tr_{b'}$ (compatibility of \eqhyperref[XVIII.eq2.8.1]{2.8.1} with base
change), while $\Tr_{u}$ and $\Tr_{u'}$ identify with
multiplication by $n$, whence the assertion\nde{It was also possible to
reduce this compatibility to the case $e=1$, which is obtained by combining
\hyperref[XVIII.lem1.1.7]{1.1.7} and \hyperref[XVIII.lem1.1.8]{1.1.8}. The arguments given here
are moreover similar to those that were passed over in silence in
\hyperref[XVIII.lem1.1.8]{1.1.8}.}. This completes the proof
of \hyperref[XVIII.thm2.9]{2.9}.

We\pageoriginale have seen in the course of the proof that:

\begin{enonce*}{Proposition 2.10}\etiquette[2.10]{XVIII.prop2.10}
The trace morphisms \emph{\hyperrefext[XVII][XVII.thm6.2.3]{6.2.3},
  XVIII~\hyperref[XVIII.prop1.1.6]{1.1.6}}, and \eqhyperref[XVIII.eq2.8.1]{2.8.1} are special
  cases of the trace morphism \hyperref[XVIII.thm2.9]{2.9}.
\end{enonce*}

\begin{enonce*}[remark]{Remark 2.10.1}\etiquette[2.10.1]{XVIII.rem2.10.1}
It can be shown without difficulty that the trace morphism of
\hyperref[XVIII.thm2.9]{2.9} is an
isomorphism if and only if the geometric fibers of $f$ have
exactly one irreducible component of dimension $d$, and its ``multiplicity'' is prime to $n$. We shall not give here the direct proof of this fact, which will follow
immediately from the ``global duality theorem'' of \S \hyperref[XVIII.sec3]{3}.
\end{enonce*}

\setcounter{subsection}{10}
\subsection{}\etiquette[2.11]{XVIII.subsec2.11}
Let $f:X\rightarrow S$ and $g:Y\rightarrow S$ be two
$S$-compactifiable morphisms satisfying $(\ast)_{d}$ and
$(\ast)_{e}$, respectively.
$$
\xymatrix{
     & \fprod{X}{Y}{S} \ar[dd]^{f\times g}\ar[ld]_{\pr_1}\ar[rd]^{\pr_2} & \\
X\ar[rd]_{f} &                                        & Y\ar[ld]^{g} \\
             &    S                                   &
}
$$

The Künneth isomorphism (\hyperrefext[XVII][XVII.thm5.4.3]{5.4.3})
$$
Rf_{!}(\ZZ/n){\mathop{\otimes}^{\LL}}_{\ZZ/n}Rg_{!}(\ZZ/n)\xrightarrow{\sim}
R(f\times g)_{!}(\ZZ/n)
$$
induces an isomorphism
\begin{equation*}
R^{2d}f_{!}\ZZ/n\otimes
R^{2e}g_{!}\ZZ/n\xrightarrow{\sim} R^{2(d+e)}(f\times
g)_{!}(\ZZ/n)\tag{2.11.1}
\end{equation*}\etiquette[2.11.1]{XVIII.eq2.11.1}

It is clear that $f\times g$ satisfies $(\ast)_{d+e}$; we have the
compatibility:

\begin{enonce*}{Proposition 2.12}\etiquette[2.12]{XVIII.prop2.12}
The following diagram
$$\vcenter{\etiquette[2.12.1]{XVIII.eq2.12.1}
\xymatrix@R=1.25cm{
R^{2d}f_!\ZZ/n(d) \otimes R^{2e}g_!\ZZ/n (e)
\ar[r]^{\sim} \ar[d]^{\Tr_f\otimes \Tr_g} & R^{2(d+e)}(f\times g)_!
\ZZ/n(d+e) \ar[d]^{\Tr_{f\times g}}\\
\ZZ/n  \ar@{=}[r] & \ZZ/n,
}}\leqno{(2.12.1)}
$$
in which the first horizontal arrow is a twisted form of
\eqhyperref[XVIII.eq2.11.1]{2.11.1},\pageoriginale is commutative.
\end{enonce*}

To verify this, it suffices to use (Var 3) and the asymmetric
interpretation of the Künneth arrow used in the proof of
\hyperrefext[XVII][XVII.thm5.4.3]{5.4.3} (\hyperrefext[XVII][XVII.lem5.4.3.5]{5.4.3.5}).

\setcounter{subsection}{12}
\subsection{}\etiquette[2.13]{XVIII.subsec2.13}
Let $f:X\rightarrow S$ be an $S$-compactifiable morphism satisfying
$(\ast)_{d}$, let $\mathcal{A}$ be a sheaf of rings on $S$, and let $n$ be an
integer invertible on $S$ such that $n\mathcal{A}=0$.

For every $K\in D(S,\mathcal{A})$, we have (\hyperrefext[XVII][XVII.thm5.2.6]{5.2.6})
\begin{equation*}

% END INCLUDED SGA4_English_current_buildable_source_20260906/source/18/drafts/SGA4_XVIII_LINES_4051_4450_English.texfrag

% BEGIN INCLUDED SGA4_English_current_buildable_source_20260906/source/18/drafts/SGA4_XVIII_LINES_4451_4850_English.texfrag
K{\mathop{\otimes}^{\LL}}_{\ZZ/n}Rf_{!}\ZZ/n(d)\xrightarrow{\sim} Rf_{!}(f^{\ast}K(d)){\quoi}.\tag{2.13.1}
\end{equation*}\etiquette[2.13.1]{XVIII.eq2.13.1}
 
The trace morphism
$$
\Tr_{f}:R^{2d}f_{!}\ZZ/n(d)\longrightarrow \ZZ/n
$$
may also be interpreted as a morphism
$$
\Tr_{f}:Rf_{!}\ZZ/n(d)[2d]\longrightarrow \ZZ/n{\quoi}.
$$
Tensoring this morphism with $K$ defines, via
\eqhyperref[XVIII.eq2.13.1]{2.13.1}, a trace
morphism
\begin{equation*}
\Tr_{f}:Rf_{!}(f^{\ast}K(d)[2d])\longrightarrow K{\quoi}.\tag{2.13.2}
\end{equation*}\etiquette[2.13.2]{XVIII.eq2.13.2}

As we shall see in \hyperref[XVIII.thm3.2.5]{3.2.5}, the following theorem is
essentially equivalent to the global duality theorem in étale
cohomology.

\begin{enonce*}{Theorem 2.14}\etiquette[2.14]{XVIII.thm2.14}
Let $f:X\rightarrow S$ be a smooth compactifiable morphism
purely of relative dimension $d$, and let $n$ be an integer $\geq 1$
invertible on $S$. For every geometric point $x$ of $X$
with image $s$ in $S$, there exist an étale neighborhood $V$ of $s$ in $S$
and an étale neighborhood $U$ of $x$ in $f^{-1}(V)$
\begin{equation*}
\vcenter{
\xymatrix@C=1.5cm@R=1.25cm{
U \ar[rd]_{f'_V} \ar[r]^{j} & f^{-1}(V)   \ar[r] \ar[d]^{f_V}& X\ar[d]^{f}\\
                                             &     V           \ar[r]  & S
}}\tag{2.14.1}
\end{equation*}\etiquette[2.14.1]{XVIII.eq2.14.1}%
such\pageoriginale that the morphism
\sisi{\emph{(XVII 6.2.7.2)}}{\eqhyperrefext[XVII][XVII.eq6.2.7.4]{6.2.7.4}}
$$
R^{i}f'_{V!}\ZZ/n\longrightarrow R^{i}f_{V!}\ZZ/n
$$
is zero for $i<2d$, and the morphism
$$
\Tr_{f'_{V}}:R^{2d}f'_{V!}\ZZ/n(d)\longrightarrow \ZZ/n
$$
is an isomorphism.
\end{enonce*}

We first prove the following.

\begin{enonce*}{Lemma 2.14.2}\etiquette[2.14.2]{XVIII.lem2.14.2}
Let $\mathcal{A}$ be an abelian category, let $k$ be an integer $\geq 0$,
let $(K_{i})_{0\leq i\leq 2k}$ be objects of $D^{b}(\mathcal{A})$ such that
$H^{p}(K_{i})=0$ for $p\not\in [0,k]$, let $f_{i}:K_{i}\rightarrow
K_{i+1}$ $(0\leq i\leq 2k-1)$ be morphisms, and let $f$ be their composite. If
$H^{p}(f_{i})=0$ for $p<k$, then there exists in $D^{b}(\mathcal{A})$
a morphism $\varphi$ from the complex $H^{k}(K_{0})[-k]$, consisting of
$H^{k}(K_{0})$ placed in degree $k$, to $K_{2k}$, making
the diagram
$$
\xymatrix{
K_{0}  \ar[d]\ar[r]^{f} & K_{2k}\\
\sigma_{\geq k}(K_{0})  \ar@{=}[r] & H^{k}(K_{0})[-k] \ar[u]_{\varphi}.
}
$$
commute.
\end{enonce*}

The lemma is trivial for $k=0$. We prove it by induction on
$k$. Applying the induction hypothesis to the complexes $\sigma_{\geq
1}(K_{i})[1]$ gives a morphism
$\varphi':H^{k}(K_{0})[-k]\rightarrow \sigma_{\geq 1}K_{2k-2}$ making
the following diagram commute, where $f'=f_{2k-3}\ldots f_{0}$:
$$
\xymatrix{
K_{0}  \ar[d]\ar[r]^{f'} & \sigma_{\geq 1}K_{2k-2}            \\
\sigma_{\geq k}(K_{0})  \ar@{=}[r] & H^{k}(K_{0})[-k]. \ar[u]_{\varphi'}
}
$$

For every complex $L$, the distinguished triangle
$$
\xymatrix{
                            &\sigma_{\geq 1}(L)\ar[ld]_{+1}  &       \\
  \sigma_{\leq 0}(L) \ar[rr] &                    & L\ar[lu]
}
$$
gives,\pageoriginale for every complex $M$, an exact sequence
$$
\ldots\rightarrow \Hom(M,\sigma_{\leq 0}(L))\rightarrow \Hom
(M,L)\rightarrow \Hom (M,\sigma_{\geq
1}L)\xrightarrow{j}\Ext^{1}(M,\sigma_{\leq 0}L){\quoi}.
$$

In particular, we have a diagram
$${\def\objectstyle{\scriptscriptstyle}
\xymatrix@C=0.4cm{
\Hom(\sigma_{\geq k}K_{0}, \sigma_{\leq 0} K_{2k-2})\ar[r] \ar[d]&
\overset{\text{\ding{176}}}{\underset{\phantom{\text{\ding{176}}}}
{\Hom}}(\sigma_{\geq k}K_{0},
    K_{2k-2}) \ar[r]\ar[d] &
\overset{\text{\ding{174}}}{\underset{\phantom{\text{\ding{174}}}}
{\Hom}}(\sigma_{\geq
      k}K_{0}, \sigma_{\geq 1} K_{2k-2})\ar[r] \ar[d]&
\overset{\text{\ding{175}}}{\underset{\phantom{\text{\ding{175}}}}
{\Ext}^1}(\sigma_{\geq
      k}K_{0}, \sigma_{\leq 0} K_{2k-2})\ar[d]\\
\overset{\text{\ding{178}}}{\underset{\phantom{\text{\ding{178}}}}
{\Hom}}(K_0,\sigma_{\leq 0} K_{2k-2})\ar[r] &
\overset{\text{\ding{172}}}{\underset{\phantom{\text{\ding{172}}}}
{\Hom}}(K_0,K_{2k-2}) \ar[r] &
\overset{\text{\ding{173}\ding{177}}}{\underset{\phantom{\text{\ding{173}\ding{177}}}}
{\Hom}}(K_0, \sigma_{\geq
  1}K_{2k-2})\ar[r]& \Ext^{1}(K_0, \sigma_{\leq 0} K_{2k-2})
}}\leqno{(2.14.3)}\etiquette[2.14.3]{XVIII.eq2.14.3}
$$
which maps to the analogous diagrams for $K_{2k-1}$ and
$K_{2k}$. By hypothesis, the morphism $f'$ at \ding{172} has an image
at \ding{173} that lifts to \ding{174}. The obstruction
\ding{175} to the element \ding{174} lifting to \ding{176}
has zero image in diagram \eqhyperref[XVIII.eq2.14.3]{2.14.3} for
$K_{2k-1}$. In this new diagram, \ding{174} therefore lifts to
\ding{176}, and the difference between \ding{172} and the image of
\ding{176} has zero image at \ding{177}, and hence is the image of an
element \ding{178}. This element \ding{178} has zero image in
diagram \hyperref[XVIII.eq2.14.3]{2.14.3} for $K_{2k}$; in this new diagram,
$f$ is therefore the image of \ding{176}, which solves the problem.

We prove that, for a given value of $d$, assertion
\hyperref[XVIII.thm2.14]{2.14} implies the following.

\begin{enonce*}{Corollary 2.14.4}\etiquette[2.14.4]{XVIII.cor2.14.4}
Under the hypotheses of \hyperref[XVIII.thm2.14]{2.14}, there exists a diagram
\eqhyperref[XVIII.eq2.14.1]{2.14.1} such
that the morphism \emph{\hyperrefext[XVII][XVII.subsubsec6.2.4]{6.2.4}}, in
$D^{b}(V,\sheaf{\ZZ}/n)$, admits a factorization
$$
\xymatrix@C=.75cm@R=1cm{
Rf'_{V!}\sheaf{\ZZ}/n(d)  \ar[rr] \ar[rd]_{t}&  & Rf_{V!}\sheaf{\ZZ}/n(d)\\
                          &\sheaf{\ZZ}/n[-2d], \ar[ru]_{\varphi}
}
$$
where $t$ is the arrow
% SGA-SRC-001151
$$
t:Rf'_{V!}\ZZ/n(d)\rightarrow \sigma_{\geq
2d}Rf'_{V!}\ZZ/n(d)=H^{2d}(Rf'_{V!}\ZZ/n(d))[-2d]\xrightarrow{\Tr}\sheaf{\ZZ}/n[-2d]{\quoi}.
$$
\end{enonce*}

Set\pageoriginale $f'_{0}=f$ and construct, inductively, $4d$
diagrams \eqhyperref[XVIII.eq2.14.1]{2.14.1}
$$
\xymatrix@C=2cm@R=1.5cm{
U^{'}_{i+1} \ar[r]\ar[rd]_{f'_{i+1}}& f'^{-1}_{i}(V'_{i+1}) \ar[d]\ar[r]& U'_i \ar[d]^{f'_{i}}\\
            & V'_{i+1}  \ar[r]            & V'_i
}
$$
satisfying \hyperref[XVIII.thm2.14]{2.14}. Set $V=V'_{4d}$ and
$U_{i}=\fprod{U'_{i}}{V}{V'_{i}}$. We obtain a diagram
$$
\xymatrix@R=1.5cm@C=1cm{
U_{4d}\ar[r] \ar[rrrd]^{f_{4d}}& \cdots \ar[r]& U_1 \ar[r]\ar[rd]^{f_1}& U_0\ar[r]\ar[d]^{f^0} & X \ar[d]^{f}\\
            &              &           & V \ar@{-}[r]  & S ,
}
$$
with $U_{0}=f^{-1}(V)$, such that the arrows
$$
R^{i}f_{k+1!}\ZZ/n(d)\longrightarrow R^{i}f_{k!}\ZZ/n(d)
$$
are zero (for $i<2d$) and the arrow
$$
\Tr:R^{2d}f_{4d!}\sheaf{\ZZ}/n(d)\longrightarrow \sheaf{\ZZ}/n
$$
is an isomorphism. By \hyperref[XVIII.lem2.14.2]{2.14.2}, the corollary holds
for $U=U_{4d}$.

We are now \sisi{ready}{ready}
to prove \hyperref[XVIII.thm2.14]{2.14} by induction on
$d$. The case $d=0$ is trivial, and the case $d=1$ is precisely
\hyperref[XVIII.lemfond1.6.9]{1.6.9}. We may therefore suppose that $d\geq 2$.

Since the problem is local on $X$ and $S$, we may suppose
that $f$ admits a factorization $f=gh$, with $g$ and $h$ smooth
and compactifiable, purely of relative dimensions $d'$ and $d''$, and that $d'$,
$d''<d=d'+d''$. Apply \hyperref[XVIII.cor2.14.4]{2.14.4} to $g$ and to the
morphisms deduced from
$h$ by base change, so as to obtain a diagram
$$
\xymatrix@C=2.2cm@R=1.5cm{
  U  \ar[r]^{u} \ar[rd]^{h'_0} & f^{-1}W \ar[d]^{h_0}       \ar[r] & X \ar[d]^{h} \\
V  \ar[r]^{v} \ar[rd]^{g'_0} & g^{-1}W \ar[r] \ar[d]^{g_0}       & Y \ar[d]^{g}\\
               &    W   \ar[r]               & S,
}
$$
such\pageoriginale that the morphisms (XVII \hyperref[XVII.eq6.2.7.3]{6.2.7.3}) admit
factorizations
\begin{align*}
 Rh'_{0!}\ZZ/n(d'') & \xrightarrow{\Tr}
\ZZ/n[-2d'']\xrightarrow{\varphi} Rh_{0!}\ZZ/n(d'')\\
Rg'_{0!}\ZZ/n(d') & \xrightarrow{\Tr}\ZZ/n[-2d']\xrightarrow{\Psi} Rg_{0!}\ZZ/n(d')
\end{align*}
and such that $\Tr_{h'_0}$ and $\Tr_{g'_0}$ are isomorphisms.

Let $h_{1}$ and $h'_{1}$ be the morphisms deduced from $h_{0}$ and $h'_0$
by base change to $V$, and let $f'=g'_{0}h'_{1}$,
$f_{W}=g_{0}h_{0}$.

The following diagram is commutative:
$$
\xymatrix{
  Rf'_{!}\ZZ/n(d) \ar[rr]\ar@{=}[d] & & Rf_{W!}\ZZ/n(d) \ar@{=}[d]\\
Rg'_{0!} Rh'_1\ZZ/n(d) \ar[r]& Rg_{0!} Rh'_{0!}\ZZ/n(d) \ar[r]& Rg_{0!}
Rh_{0!}\ZZ/n(d),
}
$$
as is the diagram
$$
\xymatrix@C=2.1cm@R=1cm{
  Rg'_{0!} Rh'_{1!} \ZZ/n(d)   \ar@{=}[d] \ar@<-5ex>@/_2pc/[ddd]& &\\
Rg'_{0!} v^{\ast} Rh'_{0!} \ZZ/n(d)
 \ar[dd] \ar[r]_{Rg'_{0!}v^{\ast}(t)(d')}&
Rg'_{0!}\ZZ/n(d')[-2d''] \ar[d]^{t[-2d'']}& \\
   & \ZZ/n[-2d] \ar[d]^{\Psi[-2d'']}\\
Rg_{0!} Rh'_{0!} \ZZ/n(d) \ar[r]_{Rg_{0!}(t)(d')} \ar@/_2.2pc/[rr]& Rg_{0!}\ZZ/n(d')[-2d''] \ar[r]_{Rg_{0!}(\varphi)(d')}& Rg_{0!}Rh_{0!}\ZZ/n(d)
}
$$
For\pageoriginale the diagram
$$
\xymatrix@C=1.5cm@R=1.25cm{
U_1 \ar[r]\ar[rd]^{f'} & f^{-1}W \ar[r]\ar[d]^{f_W}& X\ar[d]^{f}\\
                       &  W \ar@{-}[r]     & S,
}
$$
the canonical arrow from $Rf'_{!}\ZZ/n(d)$ to
$Rf_{W!}\ZZ/n(d)$ therefore factors through $\ZZ/n[-2d]$; hence the
arrows from $R^{i}f'_{!}\ZZ/n(d)$ to
$R^{i}f_{W!}\ZZ/n(d)$ are zero for $i<2d$. Finally,
$\Tr_{f'}$ is an isomorphism, being the composite of the isomorphisms
$\Tr_{g'_{0}}$ and $R^{2d'}g'_{0!}\Tr_{h'_{1}}$.

This completes the proof of \hyperref[XVIII.thm2.14]{2.14}.

\section{The global duality theorem}\etiquette[3]{XVIII.sec3}\pageoriginale

\subsection{The functor $Rf^{!}$}\etiquette[3.1]{XVIII.sec3.1}

We shall need a variant of the canonical flasque resolutions.

\begin{enonce*}{Lemma 3.1.1}\etiquette[3.1.1]{XVIII.lem3.1.1}
Let $X$ be a coherent scheme (\ie quasi-compact and quasi-separated). The
functor $\varinjlim$ induces an equivalence between the category of
Ind-objects of the category of constructible sheaves of sets (\resp of
groups, \resp of abelian groups) and the category of sheaves
of sets (\resp of Ind-finite groups, \resp of torsion abelian groups).
\end{enonce*}

This follows from \hyperrefext[IX][IX.cor2.7.2]{2.7.2} and \hyperrefext[IX][IX.cor2.7.3]{2.7.3}.

\setcounter{subsubsection}{1}
\subsubsection{}\etiquette[3.1.2]{XVIII.subsubsec3.1.2}
Let $\Ff$ be a torsion abelian sheaf on a scheme $X$,
the filtered direct limit of constructible sheaves
$(\Ff_{i})_{i\in I}$, and let $X_{0}$ be a conservative set of
points of $X$. The \emph{modified canonical flasque
resolution} of $\mathcal{F}$ is defined by the formula
$$
C^{\ast}_{\ell}(\Ff)=\varinjlim_{i}C^{\ast}(\Ff_{i}){\quoi},
$$ where the $C^{\ast}(\Ff_{i})$ are the canonical flasque resolutions
of \hyperrefext[XVII][XVII.def4.2.2]{4.2.2}. By \hyperref[XVIII.lem3.1.1]{3.1.1}, for $X$
coherent, this definition is independent of the choice of the $\Ff_{i}$. It
also commutes with localization; this fact allows us to define
$C^{\ast}_{\ell}(\Ff)$ by globalization when $\Ff$ is a torsion
sheaf on an arbitrary scheme $X$, not necessarily coherent.

\emph{The complex $C^{\ast}_{\ell}(\Ff)$ is a resolution of
$\Ff$, functorial in $\Ff$, whose formation
commutes with localization and filtered direct limits; the
functors $C^{n}_{\ell}$ are exact.}

It follows that, for a sheaf of rings $\mathcal{A}$ on $X$,
the functors $C^{n}_{\ell}$ carry $\mathcal{A}$-Modules to
$\mathcal{A}$-Modules; they likewise carry torsion sheaves
to torsion sheaves. It is well known that

\begin{enonce*}{Lemma 3.1.3}\etiquette[3.1.3]{XVIII.lem3.1.3}
Let\pageoriginale $\mathcal{A}$ be an abelian $\sheaf{\mathcal{U}}$-category
having a small generating set and such that direct limits
indexed by a small filtered ordered set are representable in
$\mathcal{A}$, and are exact. A functor $F$ from $\mathcal{A}^{\circ}$
to $(\Ens)$ is representable if and only if it carries small direct
limits (arbitrary ones) to inverse limits.
\end{enonce*}

\Cf for example \SGA 4 I \hyperref[I.prop8.12.5]{8.12.5}, where $\mathcal{A}$ is not assumed
abelian, nor are filtered $\varinjlim$ assumed exact.

We shall apply this lemma to the category $\Mod(X,\mathcal{A})$ of
sheaves of modules on a ringed site $(X,\mathcal{A})$. In this special
case, if the functor $F$ is represented by the sheaf
$ℱ$, then
\begin{equation*}
ℱ(U)=F(\mathcal{A}_{U}){\quoi}.\tag{3.1.3.1}
\end{equation*}\etiquette[3.1.3.1]{XVIII.eq3.1.3.1}

This description of $ℱ$ gives a direct verification of
\hyperref[XVIII.lem3.1.3]{3.1.3} in the case at hand.

Recall that if $f:X\rightarrow S$ is a compactifiable morphism
(\hyperrefext[XVII][XVII.subsec3.2]{3.2}.), then $S$ is quasi-compact and quasi-separated (sic), and the
dimensions of the fibers of $f$ are bounded.

\begin{enonce*}{Theorem 3.1.4}\etiquette[3.1.4]{XVIII.thm3.1.4}
Let $f:X\rightarrow S$ be a compactifiable morphism and let $\mathcal{A}$
be a sheaf of torsion rings on $S$. Then the functor
\begin{equation*}
Rf_{!}:D(X,f^{\ast}\mathcal{A})\longrightarrow D(S,\mathcal{A})\tag{3.1.4.1}
\end{equation*}\etiquette[3.1.4.1]{XVIII.eq3.1.4.1}%
admits a partial right adjoint
\begin{equation*}
Rf^{!}:D^{+}(S,\mathcal{A})\longrightarrow D^{+}(X,f^{\ast}\mathcal{A}):\tag{3.1.4.2}
\end{equation*}\etiquette[3.1.4.2]{XVIII.eq3.1.4.2}%
for $K\in \Ob D(X,f^{\ast}\mathcal{A})$ and $L\in \Ob
D^{+}(S,\mathcal{A})$, there is a functorial isomorphism
\begin{equation*}
\Hom(Rf_{!}K,L)\simeq \Hom(K,Rf^{!}L){\quoi}.\tag{3.1.4.3}
\end{equation*}\etiquette[3.1.4.3]{XVIII.eq3.1.4.3}%
Endowed with the translation morphism making the diagram
\begin{equation*}
\vcenter{
\xymatrix{
  \Hom(Rf_!(K[1]),L)
  \ar@{<->}[d]_{\text{\rotatebox{90}{$\backsim$}}}\ar@{<->}[r]^-{\sim}&
  \Hom(K[1],Rf^{!}L)\ar@{-}[r]^-\sim&
\Hom(K,(Rf^{!}L)[-1])\ar@{<->}[d]^{\text{\rotatebox{90}{$\backsim$}}}\\
\Hom((Rf_!K)[1],L)\ar@{-}[r]^-\sim & \Hom(Rf_!K,L[-1])
\ar@{<->}[r]^-{\sim}_-{\txt{\rm adj.}} & \Hom(K,Rf^{!}(L[-1])),
}}\tag{3.1.4.4}
\end{equation*}\etiquette[3.1.4.4]{XVIII.eq3.1.4.4}%
commute,\pageoriginale this functor $Rf^{!}$ is a triangulated functor.
\end{enonce*}

N.B. The notation $Rf^{!}$ is abusive in that $Rf^{!}$ is not in
general the derived functor of a functor $f^{!}$.

\noindent
\emph{Proof} Let $d$ be an integer such that every fiber of $f$
has dimension $<d$, and choose a compactification
$$
\xymatrix@R=1.25cm@C=1.5cm{
X \ar@{^{(}->}[r]^{j} \ar[d]_{f}& \bar{X}\ar[ld]^{\bar{f}}\\
S.
}
$$

If $K$ is a complex of $f^{\ast}\mathcal{A}$-Modules on $X$, then
$$
Rf_{!}K\simeq R\bar{f}_{\ast}(j_{!}K){\quoi}.
$$

For every sheaf $F$ on $X$, the components $F^{i}$ of the complex
$\tau_{\leq 2d}C^\ast_{\ell}j_{!}F$ satisfy
\begin{equation*}
R^{k}\bar{f}_{\ast}F^{i}=0\quad \text{\rm for}\quad k>0{\quoi}.\tag{3.1.4.5}
\end{equation*}\etiquette[3.1.4.5]{XVIII.eq3.1.4.5}%
Indeed, for $i\neq 2d$, $F^{i}$ is a direct limit of flasque sheaves,
and for $i=2d$ and $k>0$, we have (\hyperrefext[XVII][XVII.cor5.2.8.1]{5.2.8.1})
$$
R^{k}\bar{f}_{\ast}F^{2d}=R^{k+2d}f_{!}F=0{\quoi}.
$$

The simple complex associated with the truncated modified canonical
flasque resolution double complex of $K$ is a resolution of $K$. By
\eqhyperref[XVIII.eq3.1.4.5]{3.1.4.5}, we therefore have (notation of \hyperrefext[XVII][XVII.subsubsec1.1.15]{1.1.15})
\begin{equation*}
R\bar{f}_{\ast}(j_{!}K)\xleftarrow{\sim} \bar{f}_{\ast}\tau''_{\leq
2d} C^{\ast}_{\ell}j_{!}K{\quoi}.\tag{3.1.4.6}
\end{equation*}\etiquette[3.1.4.6]{XVIII.eq3.1.4.6}

Denote by $f^{\bigdot}_{!}$ the functor from sheaves of modules
on $X$ to complexes\pageoriginale of sheaves on $S$ defined
by
\begin{equation*}
f^{\bigdot}_{!}(F)=\bar{f}_{\ast}\tau_{\leq
2d}C^{\ast}_{\ell}j_{!}F{\quoi}.\tag{3.1.4.7}
\end{equation*}\etiquette[3.1.4.7]{XVIII.eq3.1.4.7}

\begin{enonce*}{Lemma 3.1.4.8}\etiquette[3.1.4.8]{XVIII.lem3.1.4.8}
\begin{enumerate}
\renewcommand{\theenumi}{\roman{enumi}}
\renewcommand{\labelenumi}{(\theenumi)}
\item The functors $f^{i}_{!}$ are exact and commute with filtered
direct limits; we have $f^{i}_{!}=0$ for $i\not\in[0,2d]$.

\item Since the functor $f^{\bigdot}_{!}$ is bounded with exact
components, its derived functor \emph{(\hyperrefext[XVII][XVII.prop1.2.10]{1.2.10})} is obtained trivially and is denoted
$Rf^{\bigdot}_{!}$, or simply
$f^{\bigdot}_{!}:D(X,f^{\ast}\mathcal{A})\rightarrow
D(S,\mathcal{A})$. The arrow \eqhyperref[XVIII.eq3.1.4.6]{3.1.4.6} is an isomorphism of
functors
$$
Rf^{\bigdot}_{!}\xrightarrow{\sim} Rf_{!}{\quoi}.

% END INCLUDED SGA4_English_current_buildable_source_20260906/source/18/drafts/SGA4_XVIII_LINES_4451_4850_English.texfrag

% BEGIN INCLUDED SGA4_English_current_buildable_source_20260906/source/18/drafts/SGA4_XVIII_LINES_4851_5250_English.texfrag
$$
\end{enumerate}
\end{enonce*}

The functors $C^{i}_{\ell}$ are exact and commute with filtered
direct limits (\hyperref[XVIII.subsubsec3.1.2]{3.1.2}). They send every
sheaf to an acyclic sheaf. The functors $\bar{f}_{\ast}$
and $j_{!}$ commute with filtered direct limits. It follows that
\begin{enumerate}
\renewcommand{\theenumi}{\alph{enumi}}
\renewcommand{\labelenumi}{\theenumi)}
\item the functors $\bar{f}_{\ast}C^{i}_{\ell}j_{!}$, and hence also the
functors $f^{i}_{!}$, commute with filtered direct limits;

\item for $i<2d$, the functors
$f^{i}_{!}=\bar{f}_{\ast}C^{i}_{\ell}j_{!}$ are exact. For $i=2d$,
one verifies that $f^{i}_{!}$ is exact by \eqhyperref[XVIII.eq3.1.4.5]{3.1.4.5}.
\end{enumerate}

The remaining assertions of \hyperref[XVIII.lem3.1.4.8]{3.1.4.8} are trivial.

\hypertarget{sga4.exp.XVIII.remark.3.1.4.9}{}\begin{enonce*}[remark]{Remark 3.1.4.9}\etiquette[3.1.4.9]{XVIII.rem3.1.4.9}
The definition of $f^{\bigdot}_{!}$ is independent of the sheaf
of rings $\mathcal{A}$, and still makes sense for any abelian
sheaf. Assertion \hyperref[XVIII.lem3.1.4.8]{3.1.4.8} (i) remains valid in the
categories of torsion abelian sheaves on $X$ and $S$.
\end{enonce*}

By \hyperref[XVIII.lem3.1.4.8]{3.1.4.8} and \hyperref[XVIII.lem3.1.3]{3.1.3}, the functor
$$
f^{i}_{!}:\Mod(X,f^{\ast}\mathcal{A})\longrightarrow
\Mod(S,\mathcal{A})
$$
has a right adjoint $f^{!}_{i}$; since $f^{i}_{!}$ is exact,
$f^{!}_{i}$ sends injectives to injectives.

With the sign convention of (\hyperrefext[XVII][XVII.subsubsec1.1.12]{1.1.12}), the functors $f^{!}_{i}$
form a complex of left exact functors, zero in cohomological degree
$n\not\in [-2d,0]$:
\begin{equation*}
f^{!}_{\bigdot}:\Mod(S,\mathcal{A})\longrightarrow
C^{b}(X,f^{\ast}\mathcal{A}){\quoi}.\tag{3.1.4.10}
\end{equation*}\etiquette[3.1.4.10]{XVIII.eq3.1.4.10}

The\pageoriginale functor $f^{!}_{\bigdot}$ derives
(\hyperrefext[XVII][XVII.prop1.2.10]{1.2.10}) to a triangulated functor
\begin{equation*}
Rf^{!}:D^{+}(S,\mathcal{A})\longrightarrow
D^{+}(X,f^{\ast}\mathcal{A}){\quoi},\tag{3.1.4.11}
\end{equation*}\etiquette[3.1.4.11]{XVIII.eq3.1.4.11}%
right adjoint to the functor $Rf_{!}$
$(=Rf^{\bigdot}_{!}=Lf^{\bigdot}_{!})$ (``trivial duality
formula''). More precisely, if $I\in \Ob C^{+}(S,\mathcal{A})$
is a complex of sheaves of injective $\mathcal{A}$-modules, the
$\mathcal{A}$-modules $f^{!}_{i} I^{n}$ are injective, and there is
an isomorphism of triple complexes (\hyperrefext[XVII][XVII.subsubsec1.1.12]{1.1.12})
\begin{equation*}
\Hom^{\bigdot}(f^{\bigdot}_{!}K, I)\xrightarrow{\sim}
\Hom^{\bigdot}(K,f^{!}_{\bigdot}I){\quoi}.\tag{3.1.4.12}
\end{equation*}\etiquette[3.1.4.12]{XVIII.eq3.1.4.12}%
Passing to $H^{0}$ of the associated simple complex (computed using products)
gives the adjunction statement. We leave it to the reader to
verify the compatibility \eqhyperref[XVIII.eq3.1.4.4]{3.1.4.4}.

\hypertarget{sga4.exp.XVIII.remark.3.1.5}{}\begin{enonce*}[remark]{Remark 3.1.5}\etiquette[3.1.5]{XVIII.rem3.1.5}
Suppose that the functor $Rf^{!}$ has finite cohomological
amplitude. By \hyperrefext[XVII][XVII.prop1.2.10]{1.2.10}, the functor $f^{!}_{\bigdot}$ then admits
a derived functor
\begin{equation*}
Rf^{!}:D(S,\mathcal{A})\longrightarrow D(X,f^{\ast}\mathcal{A}){\quoi}.\tag{3.1.5.1}
\end{equation*}\etiquette[3.1.5.1]{XVIII.eq3.1.5.1}%
This derived functor is again right adjoint to the functor $Rf_{!}$
\eqhyperref[XVIII.eq3.1.4.1]{3.1.4.1}. Indeed, we have
\begin{align*}
& \Hom_{D(S)}(Rf_{!}K,L)\simeq
\varinjlim_{L\xrightarrow{\sim}L'}\Hom_{K(S)}(f^{\bigdot}_{!}K,L')\simeq
\varinjlim_{\substack{K'\xrightarrow{\sim}K\\ L\xrightarrow{\sim}
L'}}\Hom_{K(S)} (f^{\bigdot}_{!}K', L')\\
& \Hom_{D(X)}(K,Rf^{!}L)\simeq
\varinjlim_{K'\xrightarrow{\sim}K}\Hom_{K(X)}(K',Rf^{!}L)\simeq
\varinjlim_{\substack{K'\xrightarrow{\sim}K\\
L\xrightarrow{\sim}L'}}\Hom_{K(X)}(K', f^{!}_{\bigdot}L')
\end{align*}
and, by adjunction,
% SGA-SRC-001156
$$
\Hom_{K(S)}(f^{\bigdot}_{!}K',L')=H^{0}\Hom(f^{\bigdot}_{!}K',L')\simeq
H^{0}\Hom(K',f^{!}_{\bigdot}L')=\Hom_{K(X)}(K',f^{!}_{\bigdot}L'){\quoi}.
$$
\end{enonce*}

\hypertarget{sga4.exp.XVIII.definition.3.1.6}{}\begin{enonce*}{Definition 3.1.6}\etiquette[3.1.6]{XVIII.def3.1.6}
The triangulated functor
$$
Rf^{!}:D^{+}(S,\mathcal{A})\longrightarrow
D^{+}(X,f^{\ast}\mathcal{A})
$$
(or,\pageoriginale under the hypotheses of \hyperref[XVIII.rem3.1.5]{3.1.5},
$$
Rf^{!}:D(S,\mathcal{A})\longrightarrow
D(X,f^{\ast}\mathcal{A})){\quoi},
$$
right adjoint to the functor $Rf_{!}$ \eqhyperref[XVIII.eq3.1.4.1]{3.1.4.1}, is called the
\emph{exceptional inverse-image functor.}
\end{enonce*}

\hypertarget{sga4.exp.XVIII.proposition.3.1.7}{}\begin{enonce*}{Proposition 3.1.7}\etiquette[3.1.7]{XVIII.prop3.1.7}
Suppose that the compactifiable morphism $f:X\rightarrow S$ has relative
dimension $\leq d$.
\begin{enumerate}
\renewcommand{\theenumi}{\roman{enumi}}
\renewcommand{\labelenumi}{(\theenumi)}
\item If $L\in \Ob D^{+}(S,\mathcal{A})$ satisfies $H^{i}(L)=0$ for
$i\leq k$, then $H^{i}Rf^{!}L=0$ for $i\leq k-2d$.

\item For $L\in \Ob D^{b}(S,\mathcal{A})$, we have
$$
\dim \inj (Rf^{!}L)\leq \dim \inj (L){\quoi}.
$$
\end{enumerate}
\end{enonce*}

The functor $Rf^{!}$ is the right derived functor of
$f^{!}_{\bigdot}:\Mod(S,\mathcal{A})\rightarrow
C^{b}\Mod(X,f^{\ast}\mathcal{A})$. Assertion (i) follows from the fact that
$(f^{!}_{\bigdot})^{i}=0$ for $i<-2d$, and assertion (ii) follows from
the fact that $(f^{!}_{\bigdot})^{i}=0$ for $i>0$, and that the functors
$f^{!}_{i}$ send injectives to injectives.

\hypertarget{sga4.exp.XVIII.proposition.3.1.8}{}\begin{enonce*}{Proposition 3.1.8}\etiquette[3.1.8]{XVIII.prop3.1.8}
Let $f:X\rightarrow S$ be a quasi-finite compactifiable morphism.
\begin{enumerate}
\renewcommand{\theenumi}{\roman{enumi}}
\renewcommand{\labelenumi}{(\theenumi)}
\item The functor $f_{!}:\Mod(X,f^{\ast}\mathcal{A})\rightarrow
\Mod(S,\mathcal{A})$ admits a right adjoint
$f^{!}:\Mod(S,\mathcal{A})\rightarrow \Mod(X,f^{\ast}\mathcal{A})$, and
$Rf^{!}$ is the derived functor of $f^{!}$.

\item If $f$ is a closed immersion, then $f^{!}$ is the functor of
``sections with support in $X$.''

\item If $f$ is étale, then $f^{!}$ identifies with the functor
$f^{\ast}$, with
$$
\Tr_{f}:f_{!}f^{\ast}F\longrightarrow F
$$
as adjunction morphism.
\end{enumerate}
\end{enonce*}

Assertion (i) follows immediately from the special case $d=0$ of the
proof of \hyperref[XVIII.thm3.1.4]{3.1.4}: we have $f^{\bigdot}_{!}=f_{!}$ and
$f^{!}_{\bigdot}=f^{!}$. Equivalently, the existence of an
adjoint functor $f^{!}$ follows from \hyperref[XVIII.lem3.1.3]{3.1.3} and from the fact that $f_{!}$
is\pageoriginale exact and commutes with filtered direct limits;
the adjunction $(f_{!},Rf^{!})$ is a ``trivial duality
formula.''

Assertion (ii) is clear, and (iii) is \hyperrefext[XVII][XVII.prop6.2.11]{6.2.11}.

\setcounter{subsubsection}{8}
\subsubsection{}\etiquette[3.1.9]{XVIII.subsubsec3.1.9}
Let $f:X\rightarrow S$ be a compactifiable morphism and $\mathcal{A}$
a sheaf of torsion rings on $S$, and equip $X$ with the inverse-image
sheaf of rings of $\mathcal{A}$. Let $K\in \Ob
D^{-}(X,f^{\ast}\mathcal{A})$, $L\in \Ob D^{+}(X,f^{\ast}\mathcal{A})$,
and represent $L$ by a bounded-below complex of injective sheaves, so that
$$
R\Hhom(K,L)\simeq \Hhom(K,L){\quoi}.
$$
The complex $\Hhom(K,L)$ has flasque components, and therefore
\begin{equation*}
Rf_{\ast}R\Hhom(K,L)\simeq f_{\ast}\Hhom(K,L){\quoi}.\tag{3.1.9.1}
\end{equation*}\etiquette[3.1.9.1]{XVIII.eq3.1.9.1}%
Choose a compactification of $f$ and an integer $d$ as in the
proof of \hyperref[XVIII.thm3.1.4]{3.1.4}, giving a functor
$f^{\bigdot}_{!}$ \eqhyperref[XVIII.eq3.1.4.7]{3.1.4.7}.
\begin{equation*}
\vcenter{
\xymatrix@C=1.5cm@R=.5cm{
X_V \ar[rr]^{k_X} \ar@{^{(}->}[rd]_{j_V} \ar[ddd]^{f_V}&      & X\ar[ddd]^{f} \ar@{^{(}->}[rd]^{j}         \\
    & \bar{X}_V \ar[ldd]^{\bar{f}_V} \ar[rr]_(.3){k_{\bar{X}}}   &   & \bar{X} \ar[ldd]_{\bar{f}}\\
    &         &        &\\
V  \ar[rr]^{k} &      & S.
}}\tag{3.1.9.2}
\end{equation*}\etiquette[3.1.9.2]{XVIII.eq3.1.9.2}


For every $V\in \Ob S_{\text{\rm et}}$, there is an isomorphism\nde{This
morphism is given by the construction of \hyperrefext[XVII][XVII.rem4.2.6]{4.2.6}, but there is no
reason for it to be an isomorphism unless coherent choices of points have
first been made on the topoi under consideration. For this to be the case, the
square occurring in \hyperrefext[XVII][XVII.rem4.2.6]{4.2.6} must in some sense be
Cartesian. More precisely, one should fix here a family of points
$(p)_{p\in P}$
of the topos $S_{\text{\rm et}}$, and for every $V\in \Ob S_{\text{\rm et}}$ one
considers the modified canonical flasque resolution using the family of
points associated by the correspondence of \hyperrefext[IV][IV.subsec6.7]{6.7}
with the pairs $(p,\xi)$ where $p\in P$ and $\xi\in V_p$.}
\begin{equation*}
k^{\ast}f^{\bigdot}_{!}\xrightarrow{\sim}
f^{\bigdot}_{V!}k^{\ast}_{X}{\quoi},\tag{3.1.9.3}
\end{equation*}\etiquette[3.1.9.3]{XVIII.eq3.1.9.3}%
so that
\begin{equation*}
\Hhom^{\bigdot}(f^{\bigdot}_{!}K,f^{\bigdot}_{!}L)(V)\simeq
\Hhom^{\bigdot}(f^{\bigdot}_{V!}(k^{\ast}_{X}K),f^{\bigdot}_{V!}(k^{\ast}_{X}L)){\quoi}.
\end{equation*}%
The functoriality of $f^{\bigdot}_{V!}$ gives us
$$
\Hhom^{\bigdot}(K,L)(X_{V})=\Hom^{\bigdot}(k^{\ast}_{X}K,k^{\ast}_{X}L)\longrightarrow
\Hom^{\bigdot}(f^{\bigdot}_{V!}(k^{\ast}_{X}K),f^{\bigdot}_{V!}(k^{\ast}_{X}L)){\quoi},
$$
hence an arrow, at the level of complexes,
\begin{equation*}
f_{\ast}\Hhom^{\bigdot}(K,L)\longrightarrow
\Hhom^{\bigdot}(f^{\bigdot}_{!}K,f^{\bigdot}_{!}L){\quoi}.\tag{3.1.9.4}
\end{equation*}\etiquette[3.1.9.4]{XVIII.eq3.1.9.4}

By\pageoriginale \eqhyperref[XVIII.eq3.1.9.1]{3.1.9.1}, this arrow derives to
\begin{equation*}
Rf_{\ast}R\Hhom(K,L)\longrightarrow
R\Hhom(Rf_{!}K,Rf_{!}L){\quoi}.\tag{3.1.9.5}
\end{equation*}\etiquette[3.1.9.5]{XVIII.eq3.1.9.5}

Now let $K\in \Ob D^{-}(X,\mathcal{A})$ and $L\in \Ob
D^{+}(S,\mathcal{A})$. The adjunction arrow from $Rf_{!}Rf^{!}L$ to
$L$, together with \eqhyperref[XVIII.eq3.1.9.4]{3.1.9.4}, gives a composite arrow
$$
Rf_{\ast}R\Hhom(K,Rf^{!}L)\longrightarrow R\Hhom(Rf_{!}K,
Rf_{!}Rf^{!}L)\longrightarrow R\Hhom(Rf_{!}K,L){\quoi},
$$
the local form on $S$ of the adjunction isomorphism:
\begin{equation*}
Rf_{\ast}R\Hhom(K,Rf^{!}L)\longrightarrow
R\Hhom(Rf_{!}K,L){\quoi}.\tag{3.1.9.6}
\end{equation*}\etiquette[3.1.9.6]{XVIII.eq3.1.9.6}

\hypertarget{sga4.exp.XVIII.proposition.3.1.10}{}\begin{enonce*}{Proposition 3.1.10}\etiquette[3.1.10]{XVIII.prop3.1.10}
The arrow \hyperref[XVIII.eq3.1.9.6]{3.1.9.6} is an isomorphism.
\end{enonce*}

Let $V\in \Ob S_{\text{\rm et}}$ and retain the notation of
\eqhyperref[XVIII.eq3.1.9.2]{3.1.9.2}. The (``transitivity'') isomorphism
$$
k_{!}Rf_{V!}\simeq Rf_{!}k_{X!}{\quoi},
$$
whose existence is the basis of the theory of the functor $Rf_{!}$,
transposes to a \emph{localization isomorphism} (\hyperref[XVIII.prop3.1.8]{3.1.8}. (iii))
\begin{equation*}
k^{\ast}_{X}Rf^{!}\simeq Rf^{!}_{V}k^{\ast}{\quoi},\tag{3.1.10.1}
\end{equation*}\etiquette[3.1.10.1]{XVIII.eq3.1.10.1}%
represented at the level of complexes by a morphism
$$
k^{\ast}_{X}f^{!}_{\bigdot}\longrightarrow f^{!}_{V\bigdot}k^{\ast}{\quoi}.
$$

To verify that \eqhyperref[XVIII.eq3.1.9.6]{3.1.9.6} is an isomorphism, it suffices to verify
that for every $V\in \Ob S_{\text{\rm et}}$, \eqhyperref[XVIII.eq3.1.9.6]{3.1.9.6} induces an
isomorphism
\begin{align*}
\Hom &(k^{\ast}_{X}K,k^{\ast}_{X}Rf^{!}L)\simeq
 H^{0}(V,Rf_{\ast}R\Hhom (K,Rf^{!}L))\tag{3.1.10.2}\\
& \longrightarrow H^{0}(V,R\Hom(Rf_{!}K,L))\simeq
 \Hhom(k^{\ast}Rf_{!}K,k^{\ast}L)
\end{align*}\etiquette[3.1.10.2]{XVIII.eq3.1.10.2}

For $V=S$, this is precisely \eqhyperref[XVIII.eq3.1.4.3]{3.1.4.3}. The general case follows
from the following compatibility, \emph{which the editor has not verified}:

\hypertarget{sga4.exp.XVIII.lemma.3.1.10.3}{}\begin{enonce*}{Lemma 3.1.10.3}\etiquette[3.1.10.3]{XVIII.lem3.1.10.3}
The\pageoriginale diagram
% SGA-SRC-001158
$$
\xymatrix@C=2.5cm@R=1.5cm{
\Hom(k^{\ast}_XK, k^{\ast}_XRf^{!}L)
\ar[r]^{\scriptstyle{\hyperref[XVIII.eq3.1.10.2]{(3.1.10.2)}}} \ar[d]^{\scriptstyle{\hyperref[XVIII.eq3.1.10.1]{(3.1.10.1)}}}_{\text{\rotatebox{90}{$\backsim$}}}&
\Hom(k^{\ast}Rf_!K,k^{\ast}L) \ar[d]^{\txt{\rm (localization)}}_{\text{\rotatebox{90}{$\backsim$}}}\\
\Hom(k^{\ast}_XK,Rf^{!}_Vk^{\ast}L)
\ar@{<->}[r]^{\scriptstyle{\rm{\hyperref[XVIII.eq3.1.4.3]{(3.1.4.3\; on\; V)}}}}_{\sim}& \Hom(Rf_{V!}k^{\ast}_XK,k^{\ast}L)
}
$$
is commutative.
\end{enonce*}

\setcounter{subsubsection}{10}
\subsubsection{}\etiquette[3.1.11]{XVIII.subsubsec3.1.11}
Consider a Cartesian square
\begin{equation*}
\vcenter{
\xymatrix@C=2cm@R=1.5cm{
X' \ar[r]^{g'}\ar[d]^{f'}& X\ar[d]^{f}\\
S' \ar[r]^{g} & S,
}}\tag{3.1.11.1}
\end{equation*}\etiquette[3.1.11.1]{XVIII.eq3.1.11.1}%
with $f$ compactifiable, $\mathcal{A}$ a sheaf of torsion rings on
$S$, $\mathcal{B}$ a sheaf of torsion rings on
$S'$, and $M$ a bounded-above complex of
% SGA-SRC-001159
$(g^{\ast}\mathcal{A},\mathcal{B})$-bimodules. At the level of
complexes, there is then a morphism, functorial in $K$ $(K\in \Ob
C(X,f^{\ast}\mathcal{A}))$
\begin{equation*}
M\otimes_{\mathcal{A}}g^{\ast}f^{\bigdot}_{!}K\longrightarrow
f'_{!}{}^{\bigdot}
({f'}^{\ast}M\otimes_{\mathcal{A}}{g'}^{\ast}K){\quoi},\tag{3.1.11.2}
\end{equation*}\etiquette[3.1.11.2]{XVIII.eq3.1.11.2}%
which transposes to a morphism, functorial in $L$ $(L\in \Ob
C(S',\mathcal{B}))$
\begin{equation*}
g'_{\ast}\Hhom_{\mathcal{B}}({f'}^{\ast}M,f'_{\bigdot}{}^{!}L)\longrightarrow
f^{!}_{\bigdot}g_\ast\Hhom_{\mathcal{B}}(M,L){\quoi}.\tag{3.1.11.3}
\end{equation*}\etiquette[3.1.11.3]{XVIII.eq3.1.11.3}

Passing to derived functors gives a morphism, functorial in $L$, from $D^{+}$
$(S',\mathcal{B})$ to $D^{+}(X,f^{\ast}\mathcal{A})$:
\begin{equation*}
Rg'_{\ast}R\Hhom_{\mathcal{B}}({f'}^{\ast}M,Rf'{}^{!}L)\longrightarrow
Rf^{!}Rg_{\ast}R\Hom_{\mathcal{B}}(M,L){\quoi},\tag{3.1.11.4}
\end{equation*}\etiquette[3.1.11.4]{XVIII.eq3.1.11.4}%
which makes the following diagram commutative, where $K\in \Ob
D^{-}(X,f^{\ast}\mathcal{A})$ and $L\in \Ob D^{+}(S',\mathcal{B})$:
$$
\xymatrix{
\scriptscriptstyle{\Hom(K,Rg'_{\ast} R\Hhom_{\Bb}(f'^{\ast}M,Rf'^{!}L))}
\ar@{<->}[d] \ar[r]^{\text{\ding{172}}} &
\scriptscriptstyle{\Hom(K,Rf^{!}Rg_{\ast}R\Hhom_{\Bb}(M,L))}
\ar@{<->}[r]_{{\scriptscriptstyle{\rm adj}}}& \scriptscriptstyle{\Hom(Rf_!K, Rg_{\ast}
R\Hhom_{\Bb}(M,L))}\ar@{<->}[d]\\
\scriptscriptstyle{\Hom(g'^{\ast}K,R\Hhom_{\Bb}(f'^{\ast} M, Rf'^{!}L))}\ar@{<->}[d]
& &  \scriptscriptstyle{\Hom(g^{\ast}Rf_! K, R\Hhom_{\Bb} (M,L))} \ar@{<->}[d]\\
\scriptscriptstyle{\Hom(f'^{\ast}M\otimes_{\Aa}g'^{\ast}K,
Rf'^{!}L)} \ar@{<->}[r]_{\scriptscriptstyle{\rm adj}}& \scriptscriptstyle{\Hom(Rf'^{\cdot}_!(f'^{\ast}M\otimes^{L}_\Aa
g'^{\ast}K),L)}
\ar[r]^{\text{\ding{173}}} &
\scriptscriptstyle{\Hom_{\Bb}(M\mathop{\otimes}^{\LL}_\Aa g^{\ast}Rf_!K, L)}.
}
$$

In\pageoriginale this diagram, \ding{173} (deduced from the arrow
\hyperrefext[XVII][XVII.eq5.2.2.1]{5.2.2.1}) is an isomorphism (\hyperrefext[XVII][XVII.thm5.2.6]{5.2.6}). The
arrow \ding{172}
is therefore an isomorphism for every $K$, and this implies that
\eqhyperref[XVIII.eq3.1.11.4]{3.1.11.4} is an isomorphism.

The proper base change theorem thus transposes to the following

\hypertarget{sga4.exp.XVIII.proposition.3.1.12}{}\begin{enonce*}{Proposition 3.1.12}\etiquette[3.1.12]{XVIII.prop3.1.12}
The arrow \eqhyperref[XVIII.eq3.1.11.4]{3.1.11.4} is an isomorphism.
\end{enonce*}

This isomorphism is called the \emph{induction isomorphism}.

Here are three special cases.

\hypertarget{sga4.exp.XVIII.corollary.3.1.12.1}{}\begin{enonce*}{Corollary 3.1.12.1}\etiquette[3.1.12.1]{XVIII.cor3.1.12.1}
Let $f:X\rightarrow S$ be a compactifiable morphism, let $\mathcal{A}$
and $\mathcal{B}$ be two sheaves of torsion rings on $S$, and let
$\varphi:\mathcal{A}\rightarrow \mathcal{B}$ be a
homomorphism. Denote by $\rho$ the restriction-of-scalars functor from
$\mathcal{B}$ to $\mathcal{A}$. Then, for $K\in \Ob
D^{+}(S,\mathcal{B})$, we have
$$
\rho Rf^{!}K\xrightarrow{\sim} Rf^{!}(\rho K){\quoi}.
$$
\end{enonce*}

It suffices in \eqhyperref[XVIII.eq3.1.11.3]{3.1.11.3} to take $S=S'$, $X=X'$,
$\mathcal{A}=\mathcal{A}$, $\mathcal{B}=\mathcal{B}$, $M=\mathcal{B}$.

Thus the sheaf of rings plays no essential role in the
construction of $Rf^{!}$.

\hypertarget{sga4.exp.XVIII.corollary.3.1.12.2}{}\begin{enonce*}{Corollary 3.1.12.2}\etiquette[3.1.12.2]{XVIII.cor3.1.12.2}
Let $f:X\rightarrow S$ be a compactifiable morphism, let $\mathcal{A}$
be a sheaf of torsion rings on $S$, and let $K\in \Ob D^{-}$
$(S,\mathcal{A})$, $L\in \Ob D^{+}(S,\mathcal{A})$. We have the \emph{induction
formula}
\begin{gather*}
R\Hhom(f^{\ast}K,Rf^{!}L)\xrightarrow{\sim} Rf^{!}R\Hhom(K,L)\\
(\text{\rm take } S=S', X=X'){\quoi}.
\end{gather*}
\end{enonce*}

\hypertarget{sga4.exp.XVIII.corollary.3.1.12.3}{}\begin{enonce*}{Corollary 3.1.12.3}\etiquette[3.1.12.3]{XVIII.cor3.1.12.3}
% SGA-SRC-001161
For every square \hyperref[XVIII.subsubsec3.1.11]{3.1.11} and every
$L\in \Ob D^{+}(S',\mathcal{B})$, we have
\begin{gather*}
Rg'_{\ast}R{f'}^{!}L\xrightarrow{\sim} Rf^{!}Rg_{\ast}L\\
(\text{\rm take } \mathcal{B}=g^{\ast}\mathcal{A}=M){\quoi}.
\end{gather*}
\end{enonce*}

\setcounter{subsubsection}{12}
\subsubsection{}\etiquette[3.1.13]{XVIII.subsubsec3.1.13}
For\pageoriginale a composite $gh$ of $S$-compactifiable morphisms
$$
X\xrightarrow{h}Y\xrightarrow{g}Z
$$

% END INCLUDED SGA4_English_current_buildable_source_20260906/source/18/drafts/SGA4_XVIII_LINES_4851_5250_English.texfrag

% BEGIN INCLUDED SGA4_English_current_buildable_source_20260906/source/18/drafts/SGA4_XVIII_LINES_5251_5650_English.texfrag
the composition isomorphism
$$
c_{g,h}:Rg_{!}Rh_{!}\xrightarrow{\sim} R(gh)_{!}
$$
transposes to a composition morphism\nde{The transposed morphism actually goes
in the other direction, but this causes no problem since it is an
isomorphism.}
\begin{equation*}
c^{!}_{g,h}:Rh^{!}Rg^{!}\xrightarrow{\sim} R(gh)^{!}\tag{3.1.13.1}
\end{equation*}\etiquette[3.1.13.1]{XVIII.eq3.1.13.1}%
satisfying the usual cocycle condition for a triple composite. In
other words, the \emph{categories $D^{+}$
$(X,f^{\ast}\mathcal{A})$ form the fibers of a fibered and cofibered
category over the category whose objects are the $S$-schemes} $f:X\rightarrow S$, and whose arrows are
the $S$-compactifiable morphisms, \emph{the direct-image} (\resp
\emph{inverse-image}) \emph{functors being the functors} $Rf_{!}$
(\resp $Rf^{!}$). For every commutative square of compactifiable
morphisms,
$$
\xymatrix@=1.75cm{
X' \ar[r]_{g} \ar[d]^{f'}& X \ar[d]^{f}\\
S' \ar[r]_{g'}& S,
}
$$
we therefore have a cobase-change morphism (\hyperrefext[XVII][XVII.prop2.1.3]{2.1.3})
\begin{equation*}
ch^{!}:Rf'_{!}Rg^{!}\longrightarrow Rg^{!} Rf_{!}{\quoi}.\tag{3.1.13.2}
\end{equation*}\etiquette[3.1.13.2]{XVIII.eq3.1.13.2}

\subsubsection{}\etiquette[3.1.14]{XVIII.subsubsec3.1.14}
Consider a Cartesian square
$$
\xymatrix@=1.75cm{
X' \ar[r]^{g'} \ar[d]^{f'}& X \ar[d]^{f}\\
S' \ar[r]^{g}& S,
}
$$
with $f$ compactifiable, and let $\mathcal{A}$ be a sheaf of torsion
rings on $S$. Endow $S'$, $X$, and $X'$ with the inverse images of
$\mathcal{A}$ as their sheaves of rings, again denoted\pageoriginale by $\mathcal{A}$ by abuse
of notation. The base-change isomorphism
\begin{equation*}
g^{\ast}Rf_{!}\xrightarrow{\sim} Rf'_{!}{g'}^{\ast}\tag{3.1.14.1}
\end{equation*}\etiquette[3.1.14.1]{XVIII.eq3.1.14.1}%
allows us to regard the categories $D(S,\mathcal{A})$,
$D(X,\mathcal{A})$, $D(S',\mathcal{A})$, and $D(X',\mathcal{A})$ as
the fibers of a cofibered category over the category defined by the
commutative diagram
$$
\xymatrix@C=2cm@R=1cm{
X' \ar[d]& X\ar[l]\ar[d]\\
S' & S\ar[l],
}
$$
the direct-image functors being the functors $Rf_{!}$, $Rf'_{!}$,
$g^{\ast}$, and ${g'}^{\ast}$ (sic). When we restrict to the
categories $D^{+}$, we even obtain a fibered and cofibered category
over this diagram, the direct-image functors being the functors
$Rf^{!}$, $R{f'}^{!}$, $Rg_{\ast}$, and $Rg'_{\ast}$ (sic). We deduce
from this
\begin{enumerate}
\renewcommand{\theenumi}{\alph{enumi}}
\renewcommand{\labelenumi}{\theenumi)}
\item a ``transitivity isomorphism'' $(\sic)$ for the ``inverse-image''
functors
% SGA-SRC-001163
$$
Rg'_{\ast} R{f'}^{!}\xrightarrow{\sim} Rf^{!}Rg_{\ast}
$$
already obtained in \hyperref[XVIII.cor3.1.12.3]{3.1.12.3} (with the same definition by
transposition of \eqhyperref[XVIII.eq3.1.14.1]{3.1.14.1}).

\item a ``base-change morphism'' (\hyperrefext[XVII][XVII.prop2.1.3]{2.1.3})
\begin{equation*}
{g'}^{\ast}Rf^{!}\longrightarrow R{f'}^{!}g^{\ast}\tag{3.1.14.2}
\end{equation*}\etiquette[3.1.14.2]{XVIII.eq3.1.14.2}%
which, by \loccit, can be described in the following two ways:
\end{enumerate}

\begin{enumerate}
\renewcommand{\theenumi}{\alph{enumi}}
\renewcommand{\labelenumi}{\theenumi)}
\item by adjunction, giving a morphism \eqhyperref[XVIII.eq3.1.14.2]{3.1.14.2} amounts to
giving a morphism
$$
Rf'_{!}{g'}^{\ast}Rf^{!}\longrightarrow g^{\ast}{\quoi}.
$$
The base-change theorem provides such a morphism as the composite
$$
Rf'_{!}{g'}^{\ast}Rf^{!}\xleftarrow{\sim}
g^{\ast}Rf_{!}Rf^{!}\xrightarrow{g^{\ast}\ast \text{adj.}} g^{\ast}{\quoi}.
$$

\item by\pageoriginale adjunction, giving a morphism \eqhyperref[XVIII.eq3.1.14.2]{3.1.14.2}
amounts to giving a morphism
$$
Rf^{!}\longrightarrow Rg'_{\ast}R{f'}^{!}g^{\ast}{\quoi}.
$$
The isomorphism (\hyperref[XVIII.cor3.1.12.3]{3.1.12.3}) provides such a morphism as the composite
$$
Rf^{!}\xrightarrow[Rf^{!}\ast \text{adj.}]{} Rf^{!}Rg_{\ast}
g^{\ast}\xleftarrow{\sim} Rg'_{\ast}R{f'}^{!}g^{\ast}{\quoi}.
$$
\end{enumerate}

From their interpretations in terms of a fibered and cofibered
category, it is clear that \emph{the morphisms} (\hyperref[XVIII.cor3.1.12.3]{3.1.12.3}),
\eqhyperref[XVIII.eq3.1.13.2]{3.1.13.2}, \emph{and} \eqhyperref[XVIII.eq3.1.14.2]{3.1.14.2} \emph{each satisfy two
compatibilities of the type} \hyperrefext[XII][XII.prop4.4]{4.4} \emph{when $f$ or $g$ is a composite of two
morphisms\nde{These constructions can also be interpreted in the
context of exchange structures and cross functors, \cf \cite{sga4.exp.XVIII.bib.0}.}.}

The following results from general topology will be used in the
following subsection.

\subsubsection{}\etiquette[3.1.15]{XVIII.subsubsec3.1.15}
Let $f:X\rightarrow S$ be a morphism of sites. Denote by
$\Gamma(f)$ the following site:
\begin{enumerate}
\renewcommand{\theenumi}{\alph{enumi}}
\renewcommand{\labelenumi}{\theenumi)}
\item The category $\Gamma(f)$ is the category of triples
$(U,V,\varphi)$ with
$$
U\in \Ob (X), V\in \Ob S,\quad \varphi\in \Hom
(U,f^{\ast}V){\quoi}.
$$

\item A family of morphisms
$(u_{i},v_{i}):(U_{i},V_{i},\varphi_{i})\rightarrow (U,V,\varphi)$ is
covering if, in $X$, the morphisms $u_{i}:U_{i}\rightarrow U$ cover
$U$.
\end{enumerate}

We then have:

\paragraph{}\etiquette[3.1.15.1]{XVIII.par3.1.15.1}
The functor from $\Gamma(f)$ to $X$ given by
$$
(U,V,\varphi)\longmapsto U
$$
is an equivalence of sites $\gamma:X\rightarrow \Gamma(f)$.

\paragraph{}\etiquette[3.1.15.2]{XVIII.par3.1.15.2}
The functor $f^{\ast}_{\Gamma}$ from $S$ to $\Gamma(f)$
$$
V\longmapsto (f^{\ast}V,V,\Id_{f^{\ast}V})
$$
is a morphism of sites, denoted $f_{\Gamma}:\Gamma(f)\rightarrow S$, or
simply $f$ by abuse of notation. It satisfies
$f_{\Gamma}\gamma\simeq f$. The functor $f^{\ast}_{\Gamma}$ has a
left adjoint $f_{\Gamma\bigdot}$:
$$
f_{\Gamma\bigdot}:(U,V,\varphi)\longmapsto V{\quoi}.
$$
The\pageoriginale functor $f_{\Gamma\bigdot}$ is therefore a
comorphism of sites from $\Gamma(f)$ to $S$\nde{Since this terminology
is not familiar, it should be specified here that these are the notions
appearing in \hyperrefext[III][III.sec2]{2} and \hyperrefext[IV][IV.subsec4.7]{4.7} (one should bear in mind that
in the original edition, at the location of this second reference,
misprints made it impossible to understand in precisely which
direction the arrows went).},
and defines the same morphism of topoi as the morphism of sites
$f_{\Gamma}$.
\begin{equation*}
\vcenter{
\xymatrix@R=2.2cm@C=1.75cm{
X &         &\Gamma(f) \ar[ll]^{\approx}
\ar@<4pt>[ld]^{f_{\Gamma}}& \\
  &    S \ar[lu]^{f^{\ast}}  \ar@<4pt>[ru]^{f^{\ast}_{\Gamma}} & .
}}\tag{3.1.15.3}
\end{equation*}\etiquette[3.1.15.3]{XVIII.eq3.1.15.3}

We deduce the following rule for computing the inverse image of a
sheaf\nde{Here \hyperrefext[III][III.prop2.3]{2.3} or
\hyperrefext[IV][IV.eq4.7.3]{4.7.3} is being used.}

\setcounter{paragraph}{3}
\paragraph{}\etiquette[3.1.15.4]{XVIII.par3.1.15.4}
Let $\Ff$ be a sheaf on $S$, let
$f^{\bigdot}_{\Gamma}\Ff$ be the presheaf on $\Gamma(f)$ given
by
$$
f^{\bigdot}_{\Gamma}\Ff(U,V,\varphi)=\Ff(V){\quoi},
$$
and let $af^{\bigdot}_{\Gamma}\Ff$ be the associated sheaf. Then
$$
f^{\ast}\Ff(U)\simeq
af^{\bigdot}_{\Gamma}\Ff(U,V,\varphi){\quoi}.
$$

\subsubsection{}\etiquette[3.1.16]{XVIII.subsubsec3.1.16}
Let $f:X\rightarrow S$ be a morphism of sites, $\mathcal{A}$ a
sheaf of rings on $S$, $\mathcal{B}$ a sheaf of rings on $X$, and $d$
an integer. Let $K$ be an object of the following type:

\paragraph{}\etiquette[3.1.16.1]{XVIII.par3.1.16.1}
$K$ associates to every object $U$ of $X$ a complex $K(U)$ of
$(\mathcal{A},f_{|U\ast}(\mathcal{B}|U))$-modules on $S$, zero in
degrees $>d$, and $K(U)$ is a \emph{covariant} functor of $U$.

For every sheaf of $\mathcal{A}$-modules $ℱ$ on $S$, denote by
$K^{!}ℱ$ the complex of sheaves on $X$ generated by the complex of
presheaves of $\mathcal{B}$-modules
$$
U\longmapsto \Hom_{\mathcal{A}}(K(U),ℱ){\quoi}.
$$

Denote by $RK^{!}$ the derived functor of $K^{!}$
$$
RK^{!}:D^{+}(S,\mathcal{A})\longrightarrow
D^{+}(X,\mathcal{B}){\quoi}.
$$
If $U:K_{1}\rightarrow K_{2}$ is a morphism of objects \hyperref[XVIII.par3.1.16.1]{3.1.16.1},
then $U$ induces\pageoriginale morphisms of functors
$$
{}^{t}u:K^{!}_{2}\longrightarrow K^!_{1}\quaet RK^{!}_{2}\longrightarrow
RK^{!}_{1}{\quoi}.
$$

Under the preceding hypotheses, we have:

\begin{enonce*}{Proposition 3.1.17}\etiquette[3.1.17]{XVIII.prop3.1.17}
\begin{enumerate}
\renewcommand{\theenumi}{\roman{enumi}}
\renewcommand{\labelenumi}{(\theenumi)}
\item For every point $x$ of $X$ and every $L\in
\Ob D^{+}(S,\mathcal{A})$, denoting by $V(x)$ the category of
neighborhoods of $x$, we have
\begin{equation*}
ℋ^{n}(RK^{!}L)_{x}\simeq \varinjlim_{U\in \Ob V(x)}
\Ext^{n}(K(U),L){\quoi}.\tag{3.1.17.1}
\end{equation*}\etiquette[3.1.17.1]{XVIII.eq3.1.17.1}

\item We therefore have a spectral sequence
\begin{equation*}
% SGA-SRC-001166
E^{pq}_{2}=\varinjlim_{U\in \Ob V(x)}
\Ext^{p}(ℋ^{-q}(K(U)),L)\Longrightarrow
ℋ^{p+q}(RK^{!}L)_{x}{\quoi}.\tag{3.1.17.2}
\end{equation*}\etiquette[3.1.17.2]{XVIII.eq3.1.17.2}

\item If $X^{0}$ is a conservative set of points of $X$, and
$u:K_{1}\rightarrow K_{2}$ is a morphism of objects \eqhyperref[XVIII.eq3.1.11.1]{3.1.11.1}, then,
in order that
$$
{}^{t}u:RK^{!}_{2}\longrightarrow RK^{!}_{1}
$$
be an isomorphism, it is necessary and sufficient that, for every $n$
and every $x\in X^{0}$, the morphism of pro-objects of
$\Mod(S,\mathcal{A})$
\begin{equation*}
u_{x}:\mathop{``\varprojlim\mbox{''}}_{U\in \Ob V(x)}
ℋ^{n}K_{1}(U)\longrightarrow \mathop{``\varprojlim\mbox{''}}_{U\in
\Ob V(x)} ℋ^{n}K_{2}(U)\tag{3.1.17.3}
\end{equation*}\etiquette[3.1.17.3]{XVIII.eq3.1.17.3}%
be an isomorphism.
\end{enumerate}
\end{enonce*}

To prove (i), take $L$ to be a complex of injectives, so that
\eqhyperref[XVIII.eq3.1.17.1]{3.1.17.1} is immediate; it is clear that (i) $\Rightarrow$ (ii).

The sufficiency of condition \eqhyperref[XVIII.eq3.1.17.3]{3.1.17.3} in (iii) follows immediately
from \eqhyperref[XVIII.eq3.1.17.2]{3.1.17.2}. We shall not need the fact that it is necessary.
If ${}^{t}u$ is an isomorphism, then, by
(\hyperref[XVIII.par3.1.16.1]{3.1.16.1}), for every injective sheaf $I$
on $(S,\mathcal{A})$, we have
$$
\Hom(\mathop{``\varprojlim\mbox{''}}_{V(x)}ℋ^{n}K_{2}(U),I)\xrightarrow{\sim}\Hom(\mathop{``\varprojlim\mbox{''}}_{V(x)}
ℋ^{n}K_{1}(U),I){\quoi}, 
$$
and we conclude by noting that in every abelian category with enough
injectives (here, $\Mod(S,\mathcal{A})$), the functors $\Hom(\ast,I)$,
for $I$ an injective $\mathcal{A}$-module, form a conservative family
of functors on $\Pro(\mathcal{A})$.

\begin{enonce*}[remark]{Example 3.1.18}\etiquette[3.1.18]{XVIII.exem3.1.18}
Let\pageoriginale $f:X\rightarrow S$ be a compactifiable morphism of
schemes, and let $\mathcal{A}$ be a sheaf of torsion rings on $S$.
Endow $X$ with $f^{\ast}\mathcal{A}$ as its sheaf of rings. Choose an integer $d$ that bounds
the dimension of the fibers of $f$, as well as a compactification of
$f$, and set (\eqhyperref[XVIII.eq3.1.4.7]{3.1.4.7})
\begin{equation*}
K(U)=f^{\bigdot}_{!}((f^{\ast}\mathcal{A})_{U}){\quoi}.\tag{3.1.18.1}
\end{equation*}\etiquette[3.1.18.1]{XVIII.eq3.1.18.1}
\end{enonce*}

By formula \eqhyperref[XVIII.eq3.1.3.1]{3.1.3.1} for the adjoint of $f^{\bigdot}_{!}$, we have
an isomorphism
\begin{equation*}
f^{!}_{\bigdot}\simeq K^{!}:\Mod(S,\mathcal{A})\longrightarrow C^{b}\Mod(X,f^{\ast}\mathcal{A}){\quoi}.\tag{3.1.18.2}
\end{equation*}\etiquette[3.1.18.2]{XVIII.eq3.1.18.2}%
and hence
\begin{equation*}
Rf^{!}=RK^{!}:D^{+}(S,\mathcal{A})\longrightarrow
D^{+}(X,f^{\ast}\mathcal{A}){\quoi}.\tag{3.1.18.3}
\end{equation*}\etiquette[3.1.18.3]{XVIII.eq3.1.18.3}

\begin{enonce*}[remark]{Example 3.1.19}\etiquette[3.1.19]{XVIII.exem3.1.19}
Let $f:X\rightarrow S$, $\mathcal{A}$, $\mathcal{B}$, and $d$ be as
in \hyperref[XVIII.subsubsec3.1.16]{3.1.16}. Let $K$ be an object \sisi{of the}{of the}
following type.
\end{enonce*}

\setcounter{subsubsection}{19}
\setcounter{paragraph}{0}
\paragraph{}\etiquette[3.1.19.1]{XVIII.par3.1.19.1}
For every $W=(U,V,\varphi)\in \Ob \Gamma(f)$
(\hyperref[XVIII.subsubsec3.1.15]{3.1.15})\nde{By definition, $\varphi$ is an
$X$-morphism $U\to V\times_S X$. To understand this example properly, one must
adopt the equivalent viewpoint that interprets $\varphi$ as an
$S$-morphism $\varphi\colon U\to V$.}, $K(W)$ is a
complex of $(\mathcal{A}_{V},\varphi_{\ast}\mathcal{B}_{|U})$-modules on $V$,
functorial in $W$: for every morphism $(u,v):
W_{1}=(U_{1},V_{1},\varphi_{1})\to W_{2}=(U_{2},V_{2},\varphi_{2})$ it
is given by
\begin{align*}
K(u,v) &: K(W_{1})\longrightarrow v^{\ast}K(W_{2}),\quad\mbox{\ie}\\
K(u,v) &: v_{!}K(W_{1})\longrightarrow K(W_{2}){\quoi}.
\end{align*}

With every object $K$ of type (\hyperref[XVIII.par3.1.19.1]{3.1.19.1}) is associated an object $K'$
of type (\hyperref[XVIII.par3.1.16.1]{3.1.16.1})
relative to $f_{\Gamma}:\Gamma(f)\rightarrow S$, given by
\begin{equation*}
K'(U,V,\varphi)=j_{!}K(U,V,\varphi)\tag{3.1.19.2}
\end{equation*}\etiquette[3.1.19.2]{XVIII.eq3.1.19.2}%
where $j$ is the ``inclusion morphism'' of $V$ into $S$. We set
$$
K^{!}\mathop{=}_{\defn} {K'}^{!}\quaet RK^{!}\mathop{=}_{\defn}
R{K'}^{!}{\quoi}.
$$

\begin{enonce*}[remark]{Example 3.1.20}\etiquette[3.1.20]{XVIII.exem3.1.20}
With the notation of \hyperref[XVIII.exem3.1.19]{3.1.19}, take
$\mathcal{B}=f^{\ast}\mathcal{A}$ and set
\begin{equation*}
K(U,V,\varphi)=\mathcal{A}_{V}\quad\text{\rm(in degree 0)}\tag{3.1.20.1}
\end{equation*}\etiquette[3.1.20.1]{XVIII.eq3.1.20.1}
\end{enonce*}

By\pageoriginale \hyperref[XVIII.par3.1.15.4]{3.1.15.4}, we then have
isomorphisms of functors
\begin{gather*}
K^{!}\simeq f^{\ast}\tag{3.1.20.2}\\
RK^{!}\simeq f^{\ast}{\quoi}.\tag{3.1.20.3}
\end{gather*}\etiquette[3.1.20.2]{XVIII.eq3.1.20.2}
\etiquette[3.1.20.3]{XVIII.eq3.1.20.3}

\subsection{Poincaré duality}\etiquette[3.2]{XVIII.subsec3.2}

\subsubsection{}\etiquette[3.2.1]{XVIII.subsubsec3.2.1}
Let $f$ be a compactifiable flat morphism, purely of relative
dimension $d$ (or, more generally, satisfying condition
$(\ast)_{d}$ of \hyperref[XVIII.thm2.9]{2.9}):
$$
\xymatrix@C=2cm@R=1.5cm@M=.2cm{
X \ar@{^{(}->}[r]^{j} \ar[d]_{f}& \bar{X}\ar[ld]_{\bar{f}} &\\
S &.
}
$$
Let $n$ be an integer $\geq 1$ invertible on $S$, and suppose that $S$
is ringed by a sheaf of rings $\mathcal{A}$ such that
$n\mathcal{A}=0$. By \hyperref[XVIII.exem3.1.18]{3.1.18}, the functor
$Rf^{!}:D^{+}(S,\mathcal{A})\rightarrow D^{+}(X,f^{\ast}\mathcal{A})$ can be
computed by the procedure of \hyperref[XVIII.subsubsec3.1.16]{3.1.16}, for $K$ given
by \eqhyperref[XVIII.eq3.1.18.1]{3.1.18.1}.

Let $K'$ be the functor of type (\hyperref[XVIII.par3.1.19.1]{3.1.19.1}) that
associates to $(U,V,\varphi)\in \Ob \Gamma(f)$ the complex of sheaves
$\ZZ/n(-d)[-2d]$ on $V$, consisting of $\ZZ/n(-d)$ placed in degree
$2d$, and let $K''$ be the corresponding object of type
(\hyperref[XVIII.par3.1.16.1]{3.1.16.1}) under \eqhyperref[XVIII.eq3.1.19.2]{3.1.19.2}, relative to
$f:\Gamma(f)\rightarrow S$.

We may regard $K$ and $K''$ as objects of
\hyperref[XVIII.par3.1.16.1]{3.1.16.1} relative to
$\Gamma(f)\rightarrow S$. Moreover, for every $(U,V,\varphi)\in \Ob
\Gamma(f)$, the trace morphism \hyperref[XVIII.thm2.9]{2.9} defines a morphism
of sheaves on $V$
$$
\Tr:R^{2d}\varphi_{!}\ZZ/n\longrightarrow
\ZZ/n(-d){\quoi},
$$
whence, denoting by $j$ the morphism from $V$ to $S$, a morphism of
functors
\begin{gather*}
K(U)=\bar{f}_{\ast}\tau_{\leq
2d}C^{\ast}_{\ell}\ZZ/n_{U}\longrightarrow
ℋ^{2d}(K(U))[-2d]\simeq\tag{3.2.1.1}\\
R^{2d}f_{!}(\ZZ/n)_{U}\simeq j_{!}R^{2d}\varphi_{!} \ZZ/n
\xrightarrow{\Tr_{\varphi}} j_{!}
\ZZ/n(-d)=K''(U,V,\varphi){\quoi}.
\end{gather*}\etiquette[3.2.1.1]{XVIII.eq3.2.1.1}
 
A\pageoriginale twisted form of \hyperref[XVIII.exem3.1.20]{3.1.20} shows that
$$
{RK''}^{!}=Rf^{\ast}(d)[2d]{\quoi}.
$$
By \hyperref[XVIII.subsubsec3.1.16]{3.1.16}, we therefore have a composite morphism
\begin{equation*}

% END INCLUDED SGA4_English_current_buildable_source_20260906/source/18/drafts/SGA4_XVIII_LINES_5251_5650_English.texfrag

% BEGIN INCLUDED SGA4_English_current_buildable_source_20260906/source/18/drafts/SGA4_XVIII_LINES_5651_5928_English.texfrag
t_{f}:Rf^{\ast}(d)[2d]\simeq {RK''}^{!}\longrightarrow RK^{!}\simeq
Rf^{!}{\quoi}.\tag{3.2.1.2}
\end{equation*}\etiquette[3.2.1.2]{XVIII.eq3.2.1.2}

\begin{enonce*}[remark]{Remark 3.2.2}\etiquette[3.2.2]{XVIII.rem3.2.2}
If the functor $Rf^{!}$ has finite cohomological dimension, then
$t_f$ is still defined as a morphism between functors from
$D(S,\mathcal{A})$ to $D(X,f^{\ast}\mathcal{A})$ (definition in
\hyperref[XVIII.rem3.1.5]{3.1.5}).
\end{enonce*}

There is another method for constructing an arrow \eqhyperref[XVIII.eq3.2.1.2]{3.2.1.2},
by adjunction from \eqhyperref[XVIII.eq2.13.2]{2.13.2}
$$
\Tr:Rf_{!}(f^{\ast}L(d)[2d])\longrightarrow L{\quoi}.
$$
These two arrows coincide:

\begin{enonce*}{Lemma 3.2.3}\etiquette[3.2.3]{XVIII.lem3.2.3}
For $L\in \Ob D^{+}(S,\mathcal{A})$ (\resp $D(S,\mathcal{A})$ \emph{if
the functor $Rf^{!}$ has finite cohomological dimension}), the
diagram
$$
\xymatrix@C=2.5cm@R=1.75cm{
Rf_!(f^{\ast}L(d)[2d]) \ar[d]^{\Tr_f} \ar[r]^{Rf_!(t_f)}  & Rf_!
Rf^{!}L \ar[d]^{\txt{\rm adjunction}}\\
L                              \ar@{=}[r] & L
}
$$
is commutative.
\end{enonce*}

We first give a variant of the definition of $t_f$, in terms
of the functor $f^!_\bigdot$ \eqhyperref[XVIII.eq3.1.4.10]{3.1.4.10}.
Let $U\in \Ob X_{\text{\rm et}}$,
$V\in \Ob S_{\text{\rm et}}$, and $\varphi\in \Hom_{f}(U,V)$. By
definition,
$$
f^{!}_{\bigdot}(L)(U)=\Hom(f^{\bigdot}_{!}\mathcal{A}_{U},L){\quoi}\;;
$$
the trace morphism defines a morphism
$$
\Tr_{f}:f^{\bigdot}_{!}\mathcal{A}_{U}\longrightarrow
\mathcal{A}_{V}(-d)[-2d]
$$
whose transpose is
$$
t_{\varphi}:f^{!}_{\bigdot}(L)(U)\longleftarrow L(d)[2d](V){\quoi}.
$$
The\pageoriginale morphisms $t_{\varphi}$ define a morphism
$$
t_{f}:f^{\ast}L(d)[2d]\longrightarrow f^{!}_{\bigdot}L
$$
and we leave it to the reader to verify that, in the derived category,
it coincides with \eqhyperref[XVIII.eq3.2.1.2]{3.2.1.2}. This description shows that the following
diagrams are commutative at the level of complexes, for $L$ a
complex of modules:
$$
\xymatrix@R=1.25cm{
\Hom^{\bigdot}(\Aa_V,L) \ar[r] \ar[d]^{\Tr_f}& \Hom^{\bigdot}(f^{\ast}\Aa_V(d) [2d], f^{\ast}
L(d) [2d])\ar[d]^{t_f}\\
\Hom(f^{\bigdot}_!f^{\ast}\Aa_V(d)[2d],L) \ar@{-}[r]^{\sim}& \Hom^{\bigdot}(f^{\ast}\Aa_V(d)[2d], f^{!}_{\bigdot}L).
}
$$
By functoriality of the arrows in this diagram in $\mathcal{A}_{V}$,
it follows that for every sheaf, or every complex of sheaves $K$, the
following diagram is commutative:
$$
\xymatrix{
\Hom^{\bigdot}(K,L) \ar[r] \ar[d]^{\Tr_f}& \Hom^{\bigdot}(f^{\ast}K(d)[2d],
f^{\ast}L(d)[2d]) \ar[d]\\
\Hom^{\bigdot}(f^{\bigdot}_!f^{\ast}K(d)[2d],L) \ar@{<->}[r]^{\sim} &\Hom^{\bigdot}(f^{\ast}K(d)[2d], f^{!}_{\bigdot}L).
}
$$
Taking $K=L$ and $\Id:L\rightarrow L$, we obtain \hyperref[XVIII.lem3.2.3]{3.2.3}
at the level of complexes.

I would be grateful to anyone who has understood this proof if they
would explain it to me.

\setcounter{subsubsection}{3}
\subsubsection{}\etiquette[3.2.4]{XVIII.subsubsec3.2.4}
We deduce at once from this lemma and from the definition
(\hyperref[XVIII.subsubsec3.1.13]{3.1.13}) of the
composition isomorphisms by adjunction that if $f$ is the composite
of two flat, finitely presented, compactifiable morphisms that are
purely of relative dimensions $d_{1}$ and $d_{2}:f=gh$, and if
$d=d_{1}+d_{2}$, then the diagram
$$
\xymatrix{
h^{\ast}(g^{\ast}L(d_2)[2d_2])(d_1)[2d_1]
\ar@{-}[r]^(.6){\sim}
\ar[d]^{t_h\ast t_g}&(gh)^{\ast}L(d)[2d] \ar[d]^{t_{gh}}\\
Rh^{!} \; Rg^{!} \; L  \ar@{-}[r]^{\sim}&  R(gh)^{!}L
}
$$
is\pageoriginale commutative.

\begin{enonce*}{Theorem 3.2.5}[Poincaré duality]%
\etiquette[3.2.5]{XVIII.thm3.2.5}
Let $f:X\rightarrow S$ be a compactifiable smooth morphism
\emph{(thus $S$ is quasi-compact and quasi-separated (sic))}, let
$d$ be the locally constant function on $X$ ``relative dimension of
$f$'', let $n\geq 1$ be an integer invertible on $S$, and let
$\mathcal{A}$ be a sheaf of rings on $S$ satisfying
$n\mathcal{A}=0$. Then the functor
$$
Rf_{!}:D(X,f^{\ast}\mathcal{A})\longrightarrow D(S,\mathcal{A})
$$
admits as right adjoint the following functor from $D(S,\mathcal{A})$
to $D(X,f^{\ast}\mathcal{A})$:
$$
K\longmapsto f^{\ast}K(d)[2d]
$$
with the morphism \eqhyperref[XVIII.eq2.13.2]{2.13.2}
$$
\Tr_{f}:Rf_{!}f^{\ast}K(d)[2d]\longrightarrow K
$$
as adjunction arrow:
\hypertarget{sga4.exp.XVIII.equation.3.2.5.1}{}$$
\Hom(K,f^{\ast}L(d)[2d])\xrightarrow{\sim} \Hom(Rf_{!}K,L){\quoi}.
$$
\end{enonce*}

\noindent
\emph{Proof}. We reduce to the case where $d$ is constant. By the
description \hyperref[XVIII.subsubsec3.2.1]{3.2.1} of the arrow $t_{f}$
\eqhyperref[XVIII.eq3.2.1.2]{3.2.1.2}, Theorem \hyperref[XVIII.thm2.14]{2.14},
together with criterion \hyperref[XVIII.prop3.1.17]{3.1.17} (iii), implies that the arrow $t_{f}$ is an
isomorphism (for $L\in D^{+}(S,\mathcal{A})$)
$$
t_{f}:f^{\ast}L(d)[2d]\xrightarrow{\sim} Rf^{!}L{\quoi}.
$$

In particular, the functor $Rf^{!}$ has finite cohomological
dimension and is therefore defined on the unbounded derived category. The morphism
$t_{f}$ remains an isomorphism for $L\in D(S,\mathcal{A})$\nde{More
precisely, the fact
that $Rf^!$ has finite cohomological dimension makes it possible to show that one
can define the functor $Rf^!$ on the unbounded derived category
in such a way
that, for a complex $I^\bigdot$ of injective sheaves, $Rf^!I^\bigdot$
is the total complex associated with the double complex $f^!_\bigdot I^\bigdot$.
Then, for every $n\in\ZZ$, the morphism induced on the
$n$-th cohomology sheaf by $t_f$ evaluated at $I^\bigdot$
is identified with the same morphism for a stupid truncation $I^{\geq k}$ for
$k$ sufficiently small, which proves the assertion.}.

By\pageoriginale definition, $Rf^{!}$ is the right adjoint of the
functor $Rf_{!}$; by \hyperref[XVIII.lem3.2.3]{3.2.3}, the adjunction arrow
from $Rf_{!}Rf^{!}K$ to $K$ is identified,
via $t_{f}$, with $\Tr_{f}$, and this
completes the proof of \hyperref[XVIII.thm3.2.5]{3.2.5}.

Henceforth we identify, by means of $t_{f}$, the functors
$f^{\ast}L(d)[2d]$ and $Rf^{!}L$.

\setcounter{subsubsection}{5}
\subsubsection{}\etiquette[3.2.6]{XVIII.subsubsec3.2.6}
Let $S$ be the spectrum of an algebraically closed field $k$, let $X$
be a compactifiable scheme, smooth over $k$ and purely of relative
dimension $d$, let $n$ be an integer invertible in $k$, and let $F$ be a
locally constant constructible sheaf of $\ZZ/n$-modules, \ie a ``coefficient system'' killed by $n$. Denoting by
${}^\vee$ a $\ZZ/n$-dual, we have
$$
R\Hhom(F,\ZZ/n)=F^{\vee}\quad\text{\rm (in degree 0)}{\quoi}.
$$
This formula follows from the fact that $\ZZ/n$ is an
injective $\ZZ/n$-module.
% SGA-SRC-001173
The adjunction isomorphism of \hyperref[XVIII.thm3.2.5]{3.2.5}, or rather
the isomorphism deduced from it,
$$
R\Gamma\;R\Hhom(F,\ZZ/n(d)[2d])\xrightarrow{\sim}
R\Hom(R\Gamma_{c}(F),\ZZ/n){\quoi},
$$
is therefore written here as
\begin{equation*}
R\Gamma (F^{\vee}(d)[2d])\xrightarrow{\sim} R\Hom
(R\Gamma_{c}(F),\ZZ/n){\quoi}.\tag{3.2.6.1}
\end{equation*}\etiquette[3.2.6.1]{XVIII.eq3.2.6.1}%
Passing to cohomology groups, and using again the fact that
$\ZZ/n$ is an injective $\ZZ/n$-module, this gives
\begin{equation*}
H^{2d-i}(X,F^{\vee}(d))\xrightarrow{\sim}H^{i}_{c}(X,F)^{\vee}{\quoi},\tag{3.2.6.2}
\end{equation*}\etiquette[3.2.6.2]{XVIII.eq3.2.6.2}%
which is the usual form of Poincaré duality.

\begin{thebibliography}{99}\pageoriginale%

\sisi{}{%
\bibitem{sga4.exp.XVIII.bib.0} J. \textsc{Ayoub}, Les six opérations de Grothendieck et le
formalisme des cycles évanescents dans le monde motivique. I,
Astérisque 314 (2007).

\bibitem{sga4.exp.XVIII.bib.0bis} B. \textsc{Conrad}, A modern proof of Chevalley's theorem on
algebraic groups, Journal of the Ramanujan Mathematical Society 17 (2002),
p.~1--18.

\bibitem{sga4.exp.XVIII.bib.0ter} D. \textsc{Ferrand}, On the non additivity of the trace
in derived categories, \url{http://arxiv.org/abs/math/0506589}

}

\bibitem{sga4.exp.XVIII.bib.1} J. \textsc{Giraud}, Cohomologie non abélienne. Columbia
  University.

\bibitem{sga4.exp.XVIII.bib.2} A. \textsc{Grothendieck}, Le groupe de Brauer
  III. In: Dix exposés sur
la cohomologie des schémas. North Holl. Publ. Co.

\sisi{}%
{\bibitem{sga4.exp.XVIII.bib.2bis} A. \textsc{Grothendieck}, Technique de descente et théorème
d'existence en géométrie algébrique. VI. Les schémas de Picard : propriétés
générales, Séminaire Bourbaki 1961-1962, exposé nᵒ 236,
pages 221-243.}

\sisi{\bibitem{sga4.exp.XVIII.bib.3} L. \textsc{Illusie}, Thèse}%
{\bibitem{sga4.exp.XVIII.bib.3} L. \textsc{Illusie}, Complexe cotangent et déformations II.
Lecture Notes in Math. 283. Springer Verlag.}

\sisi{}{%
\bibitem{sga4.exp.XVIII.bib.3bis} F. \textsc{Knudsen}, D. \textsc{Mumford}, The
projectivity of the moduli space of stable curves. I : Preliminaries on
``det'' and ``Div'', Mathematica Scandinavica 39 (1976), p.~19--55.

\bibitem{sga4.exp.XVIII.bib.3ter} U. \textsc{Köpf},
Über eigentliche Familien algebraischer Varietäten
über affinoiden Räumen, Schriftenreihe des Mathematischen Instituts
der Universität Münster. 2. Serie, Heft 7 (1974).}

\bibitem{sga4.exp.XVIII.bib.4} D. \textsc{Mumford}, Geometric invariant
  theory. Springer-Verlag 1965.

\sisi{}{%
\bibitem{sga4.exp.XVIII.bib.4bis} F. \textsc{Orgogozo}, Isomotifs de dimension inférieure ou
égale à un, Manuscripta Mathematica 115 (2004), p.~339--360.}

\bibitem{sga4.exp.XVIII.bib.5} M. \textsc{Raynaud}, Faisceaux amples sur les schémas
  en groupes et les espaces homogènes. Lecture Notes in Math. 119. Springer Verlag.

\sisi{}{%
\bibitem{sga4.exp.XVIII.bib.5bis} M. \textsc{Raynaud}, Géométrie analytique rigide d'après
Tate, Kiehl. Mémoires de la Société Mathématique de France, tome 39--40
(1974), p.~319--327.}

\bibitem{sga4.exp.XVIII.bib.6} J.P. \textsc{Serre}, Groupes algébriques et corps de
classe. Publ. Inst. Math. Nancago. Hermann 1959.

\bibitem{sga4.exp.XVIII.bib.7} J.L. \textsc{Verdier}, Dualité dans la cohomologie des espaces
localement compacts. Sém. Bourbaki 300. Nov. 1965.

\bibitem{sga4.exp.XVIII.bib.8} J.L. \textsc{Verdier}, A duality theorem in the etale
  cohomology of schemes. Woods-hole seminar on alg. geometry. July 1964.

\bibitem{sga4.exp.XVIII.bib.9} Catégories dérivées, by J.L. \textsc{Verdier}. Notes
mimeographed by the IHÉS.

\end{thebibliography}


%2.17  - 106
%IX 3.1 - 138
%XVII  6.3.18.2 - 345
%1.2.0 - 855
%1.2.7.3 - 865
%1.2.2.11 - 892
%1.2.12 - 920
%6.3.3.1 - 1026
%1.3.12.4 - 1631, 1636, 1704
%1.3.2.7 - 1669
%1.2.9.7 - 1778
%1.2.16.2 - 1784
%2.12.3 - 4257
%3.1.9.10 -  5337
%3.2.5.1 - 5461

\ifx\danslelivre\undefined
\let\SGAXVIIIenddocument\enddocument
\else
\let\SGAXVIIIenddocument\relax
\fi
\SGAXVIIIenddocument

% END INCLUDED SGA4_English_current_buildable_source_20260906/source/18/drafts/SGA4_XVIII_LINES_5651_5928_English.texfrag


% END INCLUDED SGA4_English_current_buildable_source_20260906/source/18/18_English.tex


\makeatletter
\def\input@path{{../19/}}
\makeatother
% BEGIN INCLUDED SGA4_English_current_buildable_source_20260906/source/19/19_English.tex
\ifx\danslelivre\undefined
% Release engine: XeLaTeX (the LuaTeX compatibility layer breaks Xy-pic output).
\documentclass[10pt]{../fichiers-configuration-latex/smfart}
% BEGIN INCLUDED SGA4_English_current_buildable_source_20260906/source/fichiers-configuration-latex/paquets.tex
\usepackage{polyglossia}
%\usepackage[francais]{babel}
\usepackage{fontspec}
\usepackage[math-style=ISO,bold-style=upright,sans-style=upright]{unicode-math}
%\usepackage[mathletters]{ucs}
%\usepackage[utf8]{inputenc}
%\usepackage[T1]{fontenc}
\usepackage[unicode,hypertexnames=false,colorlinks=true,linkcolor=blue,citecolor=Sepia,urlcolor=magenta]{hyperref}
\setmainlanguage{french}
\setmainfont[Ligatures=TeX]{Linux Libertine O}
% variante possible : TeXGyrePagella ?
\setmathfont{XITS Math}

% maths
\usepackage{amsmath}
%\usepackage{amssymb}
%\usepackage{mathrsfs}
%\usepackage{stmaryrd}
%\usepackage{euscript}

% couleurs
\usepackage[usenames,dvipsnames]{xcolor}

% graphiques
\usepackage{graphicx}
\usepackage[all]{xy}
%\usepackage{tikz}
%\usetikzlibrary{matrix,arrows,calc}
%\tikzset{column sep=3em,row sep=2em,text height=1.5ex, text depth=0.25ex,>=angle 90}
%\tikzset{description/.style={fill=white,inner sep=1pt}}

% hyperliens, 'etiquettes
\usepackage{xr-hyper}
\usepackage{../fichiers-configuration-latex/smfthm}
%\ifx\cacheretiquettes\undefined
%\usepackage[inline]{showlabels} \else \fi % pour rendre visibles les 'etiquettes
%\usepackage{../fichiers-configuration-latex/makeidx} % pour faire un index
%\usepackage[pagebackref,unicode]{hyperref}  % pour voir o`u sont cit'ees les r'ef'erences
% hyperref is loaded once above with the cumulative-reader destination policy.
%\hypersetup{breaklinks=true} %permet le retour `a la ligne dans les liens trop longs. Incompatible avec pdflatex
%\renewcommand*{\backref}[1]{{\small [cit\'e page(s)~#1.]}}

\usepackage{xspace}
\usepackage{pifont} % pour expose XVII

% END INCLUDED SGA4_English_current_buildable_source_20260906/source/fichiers-configuration-latex/paquets.tex

% BEGIN INCLUDED SGA4_English_current_buildable_source_20260906/source/fichiers-configuration-latex/numerotation.tex
\renewcommand{\thefootnote}{\roman{footnote}} % pour mettre les notes en bas de page sous la forme (ii) et non (2), qui ressemble trop (?) a un exposant.

\newcounter{pageoriginal}
\newcommand{\pageoriginale}{\refstepcounter{pageoriginal}\marginpar{\footnotesize\textbf{\thepageoriginal}}\xspace}
\let\pageoriginaled\pageoriginale

\usepackage{ifthen}
\newboolean{original}
\def\sisi{\ifthenelse{\boolean{original}}}
\setboolean{original}{false}

\newcounter{cnote}
\def\nde#1{\ifthenelse{\boolean{original}}{}{\stepcounter{cnote}\def\thefootnote{\arabic{cnote}}\protect\footnote[\value{cnote}]{N.D.E.: #1}}}
\def\ndemark{\ifthenelse{\boolean{original}}{}{\stepcounter{cnote}\def\thefootnote{\arabic{cnote}}\protect\footnotemark[\value{cnote}]}}
\def\ndetext#1{\ifthenelse{\boolean{original}}{}{\def\thefootnote{\arabic{cnote}}\protect\footnotetext[\value{cnote}]{N.D.E.: #1}}}

% END INCLUDED SGA4_English_current_buildable_source_20260906/source/fichiers-configuration-latex/numerotation.tex

% BEGIN INCLUDED SGA4_English_current_buildable_source_20260906/source/fichiers-configuration-latex/macros.tex
\setmainlanguage{french}
\setotherlanguage{english}


%% Fixed, engine-independent release identifier.  The inherited dynamic Git
%% probe used Lua io.popen, exposed build-state-dependent output, and printed
%% a visible question mark outside a Git checkout.
\def\showgitstatus{English release r2, 2026-07-23}


% macros pour LaTeX
\renewcommand{\mathsf}{}

% Police différente pour U+2113 SCRIPT SMALL L
\setmathfont[range="2113]{latinmodern-math.otf}

% Give every inherited semantic label its own stable PDF destination instead
% of inheriting whichever automatic equation/page anchor happened to be live.
% The legacy label remains as a compatibility alias used by existing edges.
\makeatletter
\def\etiquette[#1]#2{%
  \Hy@raisedlink{\hyper@anchorstart{sga4.semantic.#2}\hyper@anchorend}%
  \def\@currentHref{sga4.semantic.#2}%
  \def\@currentlabel{#1}%
  \label{#2}%
}
\makeatother
\def\eqhyperref[#1]#2{(\hyperref[#1]{#2})}
\def\pagehyperref[#1]#2{(page \hyperref[#1]{#2})}
\def\hyperrefext[#1][#2]#3{\hyperref[#2]{#1 #3}}
\def\eqhyperrefext[#1][#2]#3{\hyperref[#2]{#1 #3}}     %cf. exposes 17
\def\neweqhyperrefext[#1][#2]#3{\hyperref[#2]{#1 #3}}  %cf. exposes 17

\def\shat#1{\protect\widehat{#1}}
\def\stilde#1{\protect\widetilde{#1}}
%%%%%%%%%%%%%%%%%%


\def\etcn{etc.\xspace}
\def\etco{etc...\xspace}
\def\etc{\sisi{\etco}{\etcn}}
\def\ldotsn{\sisi{...}{\etcn}}

\def\cf{{cf.}\kern.3em}
\def\Cf{{Cf.}\kern.3em}
\def\ie{{i.e.}\kern.3em}
\def\iev{{i.e.},\xspace}
\def\loccito{loc.\kern2pt cit.\xspace}
\def\loccitn{\textit{loc.\kern2pt cit.}\xspace}
\def\loccit{\sisi{\loccito}{\loccitn}}
\def\EGA{\textup{EGA}\kern.2em}
\def\SGA{\textup{SGA}\kern.3em}

\def\quoi{}

\def\quav{\quad,\quad}
\def\quaet{\quad \text{et}\quad}
\def\fibreproduct#1#2#3{#1 ×_{#3} #2}
\let\fprod\fibreproduct
\def\ccheck#1#2{\mbox{${\check{\check{\rm #1}}}^{\lower2pt\hbox{$\scriptstyle #2$}}$}}
\def\sheaf#1{\protect\mathcal{#1}}
\def\vieuxsheaf#1{{#1}} % cf. problème avec sheaf{a\Phi_i} par exemple...
\DeclareMathOperator{\Hom}{Hom}
\DeclareMathOperator{\can}{can}
\DeclareMathOperator{\Ind}{Ind}
\DeclareMathOperator{\ind}{ind}
%\DeclareMathOperator{\Hhomtop}{\mathcal{H}\!\mathit{omtop}}
\DeclareMathOperator{\Hhomtop}{\mathbf{Homtop}}
\DeclareMathOperator{\Point}{Point}
\DeclareMathOperator{\Pro}{Pro}
\DeclareMathOperator{\Top}{Top}
\DeclareMathOperator{\Mod}{Mod}
\DeclareMathOperator{\Cart}{Cart}
%\DeclareMathOperator{\Mmorsite}{\mathcal{M}\!\mathit{orsite}}
\DeclareMathOperator{\Mmorsite}{\mathbf{Morsite}}
\DeclareMathOperator{\loc}{loc}
\DeclareMathOperator{\Loc}{Loc}
\DeclareMathOperator{\Pic}{Pic}
\DeclareMathOperator{\Et}{Et}
\DeclareMathOperator{\card}{card}
\DeclareMathOperator{\Ens}{Ens}
\DeclareMathOperator{\h}{\underline{H}}
\DeclareMathOperator{\Fonct}{Fonct}
\DeclareMathOperator{\Ker}{Ker}
\DeclareMathOperator{\but}{but}
\DeclareMathOperator{\Ob}{Ob}
\DeclareMathOperator{\ob}{ob}
\DeclareMathOperator{\Lib}{Lib}
\DeclareMathOperator{\esj}{esj}
\DeclareMathOperator{\id}{id}
\DeclareMathOperator{\Id}{Id}
\DeclareMathOperator{\ad}{ad}
\DeclareMathOperator{\Iim}{Im}
\DeclareMathOperator{\pr}{pr}
\DeclareMathOperator{\Cat}{Cat}
\DeclareMathOperator{\car}{car}
\DeclareMathOperator{\SysInd}{SysInd}
\DeclareMathOperator{\Sup}{Sup}
\DeclareMathOperator{\Coker}{Coker}
\DeclareMathOperator{\Filt}{Filt}
\DeclareMathOperator{\acc}{acc}
\DeclareMathOperator{\sur}{sur}
\DeclareMathOperator{\cat}{cat}
\DeclareMathOperator{\Ord}{Ord}
\DeclareMathOperator{\Gr}{Gr}
\DeclareMathOperator{\Cov}{Cov}
\DeclareMathOperator{\Caf}{Caf}
\DeclareMathOperator{\coker}{coker}
\DeclareMathOperator{\coim}{coim}
\DeclareMathOperator{\im}{im}
\DeclareMathOperator{\Ouv}{Ouv}
\DeclareMathOperator{\diag}{diag}
\DeclareMathOperator{\Esp}{Esp}
\DeclareMathOperator{\TOP}{TOP}
\DeclareMathOperator{\Unif}{Unif}
\DeclareMathOperator{\Iind}{\mathbf{Ind}}
\DeclareMathOperator{\Hhomcart}{\mathbf{Homcart}}
\DeclareMathOperator{\Hhomtopess}{\mathbf{Homtopess}}
\DeclareMathOperator{\Ppoints}{\mathbf{Points}}
\DeclareMathOperator{\Ppoint}{\mathbf{Point}}
\DeclareMathOperator{\Ttors}{\mathbf{Tors}}
\DeclareMathOperator{\Ssex}{\mathbf{Sex}}
\DeclareMathOperator{\Ffib}{\mathbf{Fib}}
\DeclareMathOperator{\Vvois}{\mathbf{Vois}}
\DeclareMathOperator{\Ppointess}{\mathbf{Pointess}}
\DeclareMathOperator{\scrHom}{\mathbf{Hom}}
\DeclareMathOperator{\sheafPic}{\mathbf{Pic}}
\DeclareMathOperator{\sheafAut}{\mathbf{Aut}}
\DeclareMathOperator{\Hhomtoplocan}{\mathbf{Homtoplocan}}
\DeclareMathOperator{\Hhomtopan}{\mathbf{Homtopan}}
\DeclareMathOperator{\Pptgeom}{\mathbf{Ptgéom}}
\DeclareMathOperator{\Morsite}{Morsite}
\DeclareMathOperator{\Fib}{Fib}
\DeclareMathOperator{\Fais}{Fais}
\DeclareMathOperator{\kar}{kar}
\DeclareMathOperator{\Groupes}{Groupes}
\DeclareMathOperator{\Restr}{Restr}
\DeclareMathOperator{\Isom}{Isom}
\DeclareMathOperator{\Prol}{Prol}
\DeclareMathOperator{\Res}{Res}
\DeclareMathOperator{\Tors}{Tors}
\DeclareMathOperator{\esp}{esp}
\DeclareMathOperator{\es}{es}
\DeclareMathOperator{\od}{od}
\DeclareMathOperator{\gr}{gr}
\DeclareMathOperator{\sob}{sob}
\DeclareMathOperator{\dfn}{dfn}
\DeclareMathOperator{\Sch}{Sch}
\DeclareMathOperator{\ouv}{ouv}
\DeclareMathOperator{\cons}{cons}
\DeclareMathOperator{\Nerf}{Nerf}
\DeclareMathOperator{\Points}{Points}
\DeclareMathOperator{\Espsob}{Espsob}
\DeclareMathOperator{\Homtop}{Homtop}
\DeclareMathOperator{\Pt}{Pt}
\DeclareMathOperator{\Quot}{Quot}
\DeclareMathOperator{\glob}{glob}
\DeclareMathOperator{\Homfib}{Homfib}
\DeclareMathOperator{\Qouv}{Qouv}
\DeclareMathOperator{\Quotouv}{Quotouv}
\DeclareMathOperator{\ens}{ens}
\DeclareMathOperator{\ab}{ab}
\DeclareMathOperator{\an}{an}
\DeclareMathOperator{\fs}{fs}
\DeclareMathOperator{\red}{red}
\DeclareMathOperator{\Ab}{Ab}
\DeclareMathOperator{\cosq}{cosq}
\DeclareMathOperator{\Ann}{Ann}
\DeclareMathOperator{\Zar}{Zar}


%%%%%%%%%%%%
%%%%% SGA4-2
\DeclareMathOperator{\Sex}{Sex}
\DeclareMathOperator{\vois}{vois}
\DeclareMathOperator{\Vois}{Vois}
\DeclareMathOperator{\Ext}{Ext}
\DeclareMathOperator{\ext}{ext}
\DeclareMathOperator{\Fer}{Fer}
\DeclareMathOperator{\Add}{Add}
\DeclareMathOperator{\Cocart}{Cocart}
\DeclareMathOperator{\cosk}{cosk}
\DeclareMathOperator{\plat}{plat}

\DeclareMathOperator{\faisc}{faisc}
\DeclareMathOperator{\sk}{sk}
\DeclareMathOperator{\ch}{ch}

\def\HH{\mathbf{H}}
\def\qc{\mathrm{qc}}
\def\et{\mathrm{et}}
\def\coh{\mathrm{coh}}

\def\Hvr#1{%
\begin{array}[b]{l}
{\kern2pt\scriptstyle r}\\[-8pt]
{\kern1pt\scriptstyle \vee}\\[-5pt]
H^{#1}
\end{array}}

\def\Hvv#1{%
\begin{array}[b]{l}
\kern2pt\rotatebox{90}{${\scriptscriptstyle\ll}$}\\[-6pt]
H^{#1}
\end{array}}

\def\Hvi#1{%
\begin{array}[b]{l}
{\kern0.2pt\scriptscriptstyle \infty}\\[-8pt]
{\kern1pt\scriptstyle \vee}\\[-5pt]
H^{#1}
\end{array}}

\def\Hv#1{%
\begin{array}[b]{l}
{\kern1pt\scriptstyle \vee}\\[-5pt]
H^{#1}
\end{array}}

\DeclareMathOperator{\cart}{cart}
\DeclareMathOperator{\Ffaisc}{\mathbf{Faisc}}


\def\Ev#1{%
\begin{array}[b]{l}
{\kern1pt\scriptstyle \vee}\\[-5pt]
E'_{#1}
\end{array}}

\DeclareMathOperator{\Cont}{Cont}
\DeclareMathOperator{\Ccont}{\mathscr{C}\!\mathit{ont}}
\DeclareMathOperator{\Site}{Site}
\newcommand\MM{\mathbf{M}}
\def\limtop#1{{\displaystyle{\varprojlim_{#1}\!\mbox{\rm top }}}}
\def\Limtop{\mathop{\varprojlim\!\mbox{\rm top}}}
\DeclareMathOperator{\spec}{spec}
\def\ss{\scriptstyle{\mathscr{S}}}
\DeclareMathOperator{\Gal}{Gal}
\def\smallm{\scriptstyle{\mathscr{M}}}
\DeclareMathOperator{\End}{End}
\DeclareMathOperator{\ass}{ass}
\DeclareMathOperator{\zar}{zar}
\DeclareMathOperator{\Ogg}{Ogg}
\DeclareMathOperator{\res}{res}
\DeclareMathOperator{\tr}{tr}
\DeclareMathOperator{\Supp}{Supp}
\DeclareMathOperator{\codim}{codim}
\DeclareMathOperator{\rk}{rk}
\DeclareMathOperator{\cur}{cur}
\DeclareMathOperator{\cdqc}{cdqc}
\newcommand\DD{\mathbf{D}}
\DeclareMathOperator{\cl}{cl}
\def\JJ{\mathbf{J}}
\def\FF{\mathbf{F}}
\def\BB{\mathbf{B}}
\DeclareMathOperator{\Alget}{Alget}
\DeclareMathOperator{\Ass}{Ass}
\DeclareMathOperator{\supp}{supp}
\DeclareMathOperator{\rad}{rad}
\DeclareMathOperator{\prof}{prof}
\DeclareMathOperator{\cint}{int}
\DeclareMathOperator{\sic}{sic}
\DeclareMathOperator{\Tor}{Tor}
\DeclareMathOperator{\tors}{tors}
\DeclareMathOperator{\Sym}{Sym}
\DeclareMathOperator{\parf}{parf}
\DeclareMathOperator{\tor}{tor}
\DeclareMathOperator{\defn}{defn}
\DeclareMathOperator{\Div}{Div}
\DeclareMathOperator{\Pol}{Pol}
\DeclareMathOperator{\Fex}{Fex}
\DeclareMathOperator{\Lie}{Lie}
\DeclareMathOperator{\tot}{tot}
\DeclareMathOperator{\Aut}{Aut}
\DeclareMathOperator{\HOM}{HOM}
\DeclareMathOperator{\pch}{pch}
\DeclareMathOperator{\Ch}{Ch}
\DeclareMathOperator{\EXT}{EXT}
\DeclareMathOperator{\PIC}{PIC}
\DeclareMathOperator{\Rf}{Rf}
\DeclareMathOperator{\inj}{inj}
\newcommand\VV{\mathbf{V}}

\def\Xv{%
\begin{array}[b]{l}
{\kern1pt\scriptstyle \vee}\\[-5pt]
X
\end{array}}

%%%%%%%
\newcommand\TT{\mathbf{T}}
\newcommand\XX{\mathbf{X}}
\newcommand\EE{\mathbf{E}}
%%%%%%

\def\qedsymbol{\textup{C.Q.F.D.}}
\let\CQFD\qed
\let\cqfd\qed

%%%%%\bigdot
\font\bigsymb=cmsy10 at 4pt
\def\bigdot{{\kern1.2pt\raise 1.5pt\hbox{\bigsymb\char15}}}
%%%%%

\def\Lright{\mathop{\vcenter{\hbox{\hbox{$\longrightarrow$}\lower8pt \hbox{\kern-16pt\hbox{$\longrightarrow$}}}}}}

\def\Aa{{\mathcal A}}
\def\Bb{{\mathcal B}}
\def\Cc{{\mathcal C}}
\def\Dd{{\mathcal D}}
\def\Ee{{\mathcal E}}
\def\Ff{{\mathcal F}}
\def\Gg{{\mathcal G}}
\def\Hh{{\mathcal H}}
\def\Ii{{\mathcal I}}
\def\Jj{{\mathcal J}}
\def\Kk{{\mathcal K}}
\def\Ll{{\mathcal L}}
\def\Mm{{\mathcal M}}
\def\Nn{{\mathcal N}}
\def\Oo{{\mathcal O}}
\def\Pp{{\mathcal P}}
\def\Qq{{\mathcal Q}}
\def\Rr{{\mathcal R}}
\def\Ss{{\mathcal S}}
\def\Tt{{\mathcal T}}
\def\Uu{{\mathcal U}}
\def\Vv{{\mathcal V}}
\def\Ww{{\mathcal W}}
\def\Xx{{\mathcal X}}
\def\Yy{{\mathcal Y}}
\def\Zz{{\mathcal Z}}

\newcommand{\CC}{\mathbf{C}}
\newcommand{\GG}{\mathbf{G}}
\newcommand{\KK}{\mathbf{K}}
\newcommand{\LL}{\mathbf{L}}
\newcommand{\NN}{\mathbf{N}}
\newcommand{\PP}{\mathbf{P}}
\newcommand\QQ{\mathbf{Q}}
\newcommand\RR{\mathbf{R}}
\newcommand\ZZ{\mathbf{Z}}

\def\to{\mathchoice{\longrightarrow}{\rightarrow}{\rightarrow}{\rightarrow}}
\def\from{\mathchoice{\longleftarrow}{\leftarrow}{\leftarrow}{\leftarrow}}
\def\mto{\mathchoice{\longmapsto}{\mapsto}{\mapsto}{\mapsto}}
\def\mfrom{\mathrel{\from\joinrel\mapstochar\,}}
\def\hto{\mathrel{\lhook\joinrel\to}}
\def\hfrom{\mathrel{\from\joinrel\rhook}}

\def\implies{\mathchoice{\Longrightarrow}{\Rightarrow}{\Rightarrow}{\Rightarrow}}
\def\implied{\mathchoice{\Longleftarrow}{\Leftarrow}{\Leftarrow}{\Leftarrow}}
\def\iff{\mathchoice{\Longleftrightarrow}{\Leftrightarrow}{\Leftrightarrow}{\Leftrightarrow}}

\def\To#1{\mathchoice{\xrightarrow{\labelstyle~~#1~}}{\stackrel{#1}{\longrightarrow}}{}{}}
\def\Mto#1{\mathrel{\mapstochar\To{#1}}}
\def\From#1{\mathchoice{\xleftarrow{\labelstyle~~#1~}}{\stackrel{#1}{\longleftarrow}}{}{}}

\newcommand{\isomto}{\overset{\sim}{\longrightarrow}}
\newcommand{\isomfrom}{\overset{\sim}{\longleftarrow}}

\newcommand{\IFF}{\unskip~$\iff$~\ignorespaces}
\newcommand{\THEN}{\unskip~$\implies$~\ignorespaces}

\let\leq\leqslant
\let\geq\geqslant
\let\le\leqslant
\let\ge\geqslant

\def\cf{{cf.}\kern.3em}
\def\Cf{{Cf.}\kern.3em}
\def\ie{{i.e.}\kern.3em}
\def\iev{{i.e.},\xspace}
%\def\loccito{loc.\kern2pt cit.\xspace}
%\def\loccitn{\textit{loc.\kern2pt cit.}\xspace}
%\def\loccit{\sisi{\loccito}{\loccitn}}
\def\EGA{\textup{EGA}\kern.2em}
\def\SGA{\textup{SGA}\kern.3em}

\let\oldS\S
\def\S{\oldS\kern .15em}
\def\ptbl{.\kern .2em}
\def\resp{resp.\kern.3em}
\def\Exp{Exp.\kern.3em}

\let\bar\overline
\let\hat\widehat
\let\tilde\widetilde
\let\dpl\displaystyle
\let\ldots\dots
\let\emptyset\varnothing
\let\skpt\skippointrait
\let\epsilon\varepsilon

\def\sfrac#1#2{{#1}/{#2}}

\DeclareMathOperator{\cd}{cd}
%\DeclareMathOperator{\Hhom}{\mathcal{H}\!\mathit{om}}
\DeclareMathOperator{\Hhom}{\mathbf{Hom}}
%\DeclareMathOperator{\Eext}{\mathcal{E}\!\mathit{xt}}
\DeclareMathOperator{\Eext}{\mathbf{Ext}}
\DeclareMathOperator{\Spec}{Spec}
\DeclareMathOperator{\Tr}{Tr}

\let\emptyset=\varnothing

%%%%%%%%%%%%%%%
%%% expose XVII
%%%%%%%%%%%%%%%

\makeatletter
\let\old@optexpose\optexpose
\newcommand{\empty@optexpose}[1]{}
\newcommand{\paren@optexpose}[1]{\old@optexpose{#1}(}
\newcommand{\refnoexpose}[1]{\let\optexpose\empty@optexpose\ref{#1}\let\optexpose\old@optexpose}%\refnoexpose{X.cor5.2} engendre "5.2" au lieu de "X 5.2"
\newcommand{\neweqref}[1]{\let\optexpose\paren@optexpose\textup{\ref{#1})}\let\optexpose\old@optexpose}%\neweqref{XII.eq4.1} engendre "XII (4.1)" au lieu de "(XII 4.1)"
\makeatother
\newcommand{\refvb}[2][VI${}_{\text{B}}$ ]{\sisi{#1\refnoexpose{#2}}{\ref{#2}}}

\newcommand{\leftexp}[2]{{\vphantom{#2}}^{#1}{#2}}
\newcommand{\AAA}{\mathbf{A}}
\newcommand{\Sf}{\mathfrak{S}}
\DeclareMathOperator*{\quotprojlim}{\text{``$\varprojlim$''}}
\DeclareMathOperator*{\quotinjlim}{\text{``$\varinjlim$''}}
\DeclareMathOperator{\Fl}{F\sisi{\ell}{l}}
\DeclareMathOperator{\Pl}{P\sisi{\ell}{l}}
%\DeclareMathOperator{\cl}{c\sisi{\ell}{l}}
\DeclareMathOperator{\tdf}{tdf}
\DeclareMathOperator{\rg}{rg}
\newcommand{\iInd}{\sisi{I}{i}nd}
\newcommand{\iind}{ind}
\newcommand{\ppro}{pro}
\newcommand{\tTor}{\sisi{T}{t}or}
\newcommand{\ttor}{tor}
\newcommand{\chap}{Chap.}
\newcommand{\Ffex}{\Ff\mathit{ex}}
\newcommand{\sheafk}{\underline{k}}
\newcommand{\sheafn}{\underline{n}}
\newcommand{\sheafGamma}{\underline{\Gamma}}

% Repeated references to an already typeset footnote must not use
% \footnotemark without a matching \footnotetext: Hyperref would then create
% a live annotation to a destination that cannot exist.  This helper keeps the
% current printed footnote mark and links it to an explicit stable note target.
\makeatletter
\newcommand{\sgarepeatedfootnotemark}[1]{%
  \begingroup
  \edef\@thefnmark{\thefootnote}%
  \hyperlink{#1}{\@makefnmark}%
  \endgroup}
\makeatother

\sisi{
\DeclareMathOperator{\simarrow}{\rule{.15mm}{1.8mm}\kern -3.4pt\rightsquigarrow}
\newcommand{\sheafell}{\underline{\ell}}
\newcommand{\sheafZZ}{\underline{\ZZ}}
}{
\newcommand{\nderef}[1]{la N.D.E.~(\ref{#1}) page~\pageref{#1}}
}

%%%%%%%%%%%%%%%%%
%%%%% expose XVII
%%%%%%%%%%%%%%%%%

% END INCLUDED SGA4_English_current_buildable_source_20260906/source/fichiers-configuration-latex/macros.tex


\externaldocument{../01/01_English.tex}
\externaldocument{../02/02_English.tex}
\externaldocument{../03/03_English.tex}
\externaldocument{../04/04_English.tex}
\externaldocument{../05/05_English.tex}
\externaldocument{../05bis/05bis_English.tex}
\externaldocument{../06/06_English.tex}
\externaldocument{../07/07_English.tex}
\externaldocument{../08/08_English.tex}
\externaldocument{../09/09_English.tex}
%\externaldocument{../10/10_English.tex}
\externaldocument{../11/11_English.tex}
\externaldocument{../12/12_English.tex}
\externaldocument{../13/13_English.tex}
\externaldocument{../14/14_English.tex}
\externaldocument{../15/15_English.tex}
\externaldocument{../16/16_English.tex}
\externaldocument{../17/17_English.tex}
\externaldocument{../18/18_English.tex}

\setboolean{original}{false}

\begin{document}

\begin{center}
XIX. Cohomology of excellent preschemes of equal characteristic\\
M. Artin
\end{center}

\selectlanguage{english}version: \showgitstatus
\tableofcontents
\mainmatter

\else
\chapter{Cohomology of excellent preschemes of equal characteristic}
%\addtocontents{toc}{M. Artin}
\begin{center}
M. Artin
\end{center}
\fi
\phantomsection\MakeLinkTarget*{sga4.exp.XIX}\label{sga4.exp.XIX}

% The edition is fixed to the corrected (second) reading. This expandable
% selector is output-equivalent and avoids current Xy-pic conditional failures.
\renewcommand{\sisi}[2]{#2}

% BEGIN INCLUDED SGA4_English_current_buildable_source_20260906/source/19/drafts/SGA4_XIX_LINES_44_450_English.texfrag
\newcommand{\UU}{\mathbf{U}}

In\pageoriginale this exposé we develop the theory of étale
cohomology for excellent preschemes (\EGA IV 7.8.1) of equal
characteristic, \emph{assuming the theorem on resolution of
  singularities\pageoriginale for such preschemes}, in the following form:

Let $X=\Spec A$ be an excellent local scheme of equal characteristic,
and let $U\subset X$ be a regular open subset. There exist a proper
birational morphism $f:X'\rightarrow X$ and an open immersion
$U\rightarrow X'$ over $X$ such that $X'$ is regular and
$Y=X'-U$ is everywhere of codimension \hyperref[XIX.sec1]{1}, 
with normal crossings.

\emph{The proofs are therefore valid, “for the time being,” only in
  characteristic zero} \cite{sga4.exp.XIX.bib.4}. One could also deduce from them
results for preschemes of characteristic $p>0$ in low dimensions
\cite{sga4.exp.XIX.bib.1}.

Under the resolution hypothesis, we shall obtain more or less
satisfactory results for preschemes of equal characteristic and for
coefficients prime to the characteristic (\cf \hyperref[XIX.thm3.2]{3.2}, \hyperref[XIX.thm4.1]{4.1},
\hyperref[XIX.thm5.1]{5.1}, \hyperref[XIX.thm6.1]{6.1}). On the other hand, almost nothing is known about the cohomology
of general preschemes of unequal characteristic, even in dimension
\hyperref[XIX.sec2]{2}, where resolution is available \cite{sga4.exp.XIX.bib.2}. For example, the
purity theorem \hyperref[XIX.thm2.1]{2.1} is not known for any complete ring of unequal
characteristic of dimension $>1$. This result, for the formal power
series rings $k[[x_{1},\ldots, x_{n}]]$ (\hyperref[XIX.thm1.2]{1.2}), is one of the main
tools in the theory. Let us also note that the resolution theorem is
not used in the proof of \hyperref[XIX.thm1.2]{1.2}. \emph{This theorem is therefore
  proved in arbitrary characteristic.}

We shall often use without explicit mention the properties of
excellent\pageoriginale schemes developed in \EGA IV 7.8, such as
stability under extensions of finite type (\EGA IV 7.8.3(ii)) and
under strict localizations. The latter fact follows from resolution
and \EGA IV 7.9.5.

\section{Purity for the ring $k[[x_{1},\ldots,x_{n}]]$}\etiquette[1]{XIX.sec1}

Recall the following corollary of \hyperrefext[XVI][XVI.thm3.7]{3.7}:

\begin{enonce*}{Corollary 1.1}\etiquette[1.1]{XIX.cor1.1}
Let $A$ be a strictly local ring, and denote by $A\{x\}$ the strict
localization of $A[x]$ at the prime ideal generated by $\rad A$ and
$x$. Let $U=\Spec A\{x\}[1/x]$. Then, for $n$ prime to the residual
characteristic of $A$, we have
$$
H^{q}(U,𝛍_{n})\simeq
\begin{cases}
𝛍_{n}(A) & \text{\rm if}\quad q=0\\
\ZZ/n & \text{\rm if}\quad q=1\\
0            & \text{\rm if}\quad q>1{\quoi}.
\end{cases}
$$
\end{enonce*}

Indeed, let $X=\Spec A[x]$, and let $Y=V(x)$ be the zero locus of $x$
in $X$. Then $(Y,X)$ is a smooth pair (\hyperrefext[XVI][XVI.def3.1]{3.1}) over
$\Spec A$. By (V 4, VIII 5), $H^{q}(U,𝛍_{n})$
is identified with the fiber of the sheaf $ℋ^{q+1}_{Y}$
$(X,𝛍_{n})$ at a geometric point above the
point $(\rad A,x)$ of $X$, and the corollary therefore follows
immediately from (\hyperrefext[XVI][XVI.thm3.7]{3.7}). Note that $𝛍_{n}$ and $\ZZ/n$
are both cyclic groups of order $n$. We have used them here to give
the canonical\pageoriginale formulas (\cf \hyperref[XIX.thm3.4]{3.4}).

\begin{enonce*}{Theorem 1.2}\etiquette[1.2]{XIX.thm1.2}
Let $k$ be a separably closed field, and let $U=\Spec
k[[x_{1},\ldots,x_{r}]]$ $[1/x_{r}]$. Then, for $n$ prime to the
characteristic of $k$, we have
$$
H^{q}(U,𝛍_{n})=
\begin{cases}
𝛍_{n}(A) &\text{\rm if}\quad q=0\\
\ZZ/n &\text{\rm if}\quad q=1\\
0 &\text{\rm if}\quad q>1{\quoi}.
\end{cases}
$$
\end{enonce*}

\begin{enonce*}[remark]{Remark}
The proof applies equally to the ring of convergent power series when
$k=\CC$. Moreover, one easily avoids the hypothesis that $k$ be
separably closed by stating the result with a little more care. But
to establish the general purity theorem \hyperref[XIX.thm3.2]{3.2} (under the
resolution-of-singularities hypothesis, of course), we need the result
only when $k$ is separably closed, and these other assertions are
corollaries of that theorem \hyperref[XIX.thm3.2]{3.2}.
\end{enonce*}

In dimensions 0 and 1, the result can be proved for every regular
strictly local ring:

\begin{enonce*}{Lemma 1.3}\etiquette[1.3]{XIX.lem1.3}
Let $A$ be a regular strictly local ring and let $x\in \rad A$ be a
local parameter. Let $U=\Spec A[1/x]$. Then
$H^{0}(U,𝛍_{n})=𝛍_{n}(A)$, and
there are canonical isomorphisms
$$
\ZZ/n\simeq H^{0}(U,\GG_{m})/(H^{0}(U,\GG_{m}))^{n}\simeq
H^{1}(U,𝛍_{n}){\quoi},
$$
where\pageoriginale the generator of $\ZZ/n$ is the “residue” of $x\in H^{0}$ $(U,\GG_{m})$.
\end{enonce*}

\medskip
\noindent
\emph{Proof}. For $q=0$, the assertion is that $U$ is connected
and nonempty. This is true. For $q=1$, apply the Kummer exact sequence
(\hyperrefext[IX][IX.rems3.7]{3.7}). We have $\Pic U=0$ because $A$ is regular and
factorial, and $U$ is an open subset of $X=\Spec A$. It follows that
we have the exact sequence
$$
H^{0}(U,\GG_{m})\xrightarrow{n}H^{0}(U,\GG_{m})\rightarrow
H^{1}(U,𝛍_{n})\rightarrow 0{\quoi},
$$
and it remains to prove that
$H^{0}(U,\GG_{m})/(H^{0}(U,\GG_{m}))^{n}$ is isomorphic to
$\ZZ/n$ and generated by the residue of $x$. But there is an obvious
exact sequence
$$
\xymatrix@R=.1cm{
0 \ar[r]& H^0(X,\GG_m)\ar[r] & H^{0}(U,\GG_m)\ar[r] &
\ZZ \ar[r]& 0\\
& \text{\rotatebox{-90}{$\simeq$}} & \text{\rotatebox{-90}{$\simeq$}} &\\
& A^{\ast}  & A[1/x]^{\ast}   
}
$$
where the image of a section $a\in H^{0}$ $(U,\GG_{m})$ in
$\ZZ$ is the order of vanishing of $a$ along $\{x=0\}$. The
group $A^{\ast}$ is divisible by $n$, because $A$ is strictly local
and $n$ is prime to the residual characteristic. Indeed, if
$u\in H^{0}$ $(X,\GG_{m})$, the equation $Y^{n}-u=0$ defines a finite
étale extension of $A$, which is therefore completely split. It
follows that $H^{1}$ $(U,𝛍_{n})\simeq \ZZ/n$, proving the lemma.

\medskip
\noindent
\emph{Proof of} \hyperref[XIX.thm1.2]{1.2}. It remains to treat the case
$q>1$. Recall the terminology of (V, Appendix). We have
$$
H^{q}(U,𝛍_{n})=\varinjlim_{\UU}
H^{q}(\UU,𝛍_{n})
$$ 
where\pageoriginale $\xymatrix{\UU = \bigg\{U_0 & U_1
\ar@<3pt>[l] \ar@<-3pt>[l] & \ldots \ar@<5.5pt>[l]
\ar[l]\ar@<-5.5pt>[l]\bigg\}}$ ranges over the category (filtered up
to homotopy) of \emph{hypercoverings} of $U$. Let $\UU$ be such a
hypercovering, and take the $U_{i}$ to be separated and of finite type.
Since we wish to prove that $H^{q}$ $(U,𝛍_{n})=0$, it suffices to find
a hypercovering $\VV$ and a morphism $\VV\rightarrow \UU$ such that
the induced morphism
$$
H^{q}(\UU,𝛍_{n})\longrightarrow
H^{q}(\VV,𝛍_{n})
$$
is the zero morphism.

Let $\{W_{1},\ldots,W_{m}\}$ be the set of connected components of all
the $U_{i}$ for $i\leq q+1$, and let $B_{j}$ be the normalization of
$A=k[[x_{1},\ldots,x_{r}]]$ in the field $R(W_{j})$ of rational
functions on $W_{j}$. Then, since $W_{j}\rightarrow U$ is étale, the
scheme $W_{j}$ is an open subset of $\Spec B_{j}$, say the open
complement of $\Spec C_{j}$, where $C_{j}$ is the appropriate reduced
quotient ring of $B_{j}$.

Choose in addition an element $b_{j}$ of $B_{j}$ such that $b_{j}$ generates the
separable extension $R(W_{j})$ of
$k((x_{1},\ldots,x_{r}))$. Let $f_{j}(Y)\in A[Y]$ be the monic
irreducible polynomial of $b_{j}$ over $A$, so that $B_{j}$ is
isomorphic to the normalization of the ring $A[Y]/(f_{j})$. Finally,
let $d_{j}\in A$ be the discriminant of $f_{j}(Y)$, which is a nonzero
element of $A$.

Now we may replace the $x_{i}$ for $i<r$ by other elements of $A$, the
only condition being that $\{x_{1},\ldots,x_{r}\}$ form a system of
parameters of $A$.

It is easy to see that, after changing these elements if necessary,
we\pageoriginale may suppose that the following hypotheses are
satisfied:

\setcounter{subsection}{3}
\subsection{}\etiquette[1.4]{XIX.subsec1.4}
\begin{itemize}
\item [(i)] For every $j$ such that $d_{j}$ is not a unit,
$x_{1},\ldots,x_{r-1}$ and $d_{j}$ generate an ideal primary for
$\rad A$.

\item[(ii)] For every $j$, the images of $x_{1},\ldots,x_{r-1}$ in
$C_{j}$ generate an ideal primary for $\rad C_{j}$.
\end{itemize}
 We then have the following lemma:

\begin{enonce*}{Lemma 1.5}\etiquette[1.5]{XIX.lem1.5}
Let $A^{0}=k[[x_{1},\ldots,x_{r-1}]]$ $\{x_{r}\}$ be the henselization
of $k[[x_{1},\ldots,x_{r-1}]]$ $[x_{r}]$ at the point
$(x_{1},\ldots,x_{r})$, and let $X^{0}=\Spec A^{0}$ and $U^{0}=\Spec
A^{0}$ $[1/x_{r}]$. Under the conditions of $\hyperref[XIX.subsec1.4]{1.4}$ above, there exists
a hypercovering $\UU^{0}$ of $U^{0}$ (of order $q+1$) and an
isomorphism $\fprod{\UU^{0}}{X}{X^{0}}\simeq\UU$.
\end{enonce*}

Assume the lemma. The ring $A^{0}$ is excellent (\EGA IV 7.8.3
(ii), 7.9.5), and consequently the connected components of the
$U^{0}_{i}$ and of the $U_{i}$ correspond bijectively (\EGA IV
7.8.3 (vii)). The isomorphism $\fprod{\UU^{0}}{X}{X^{0}}\simeq
\UU$ therefore gives an isomorphism
$$
H^{q}(\UU^{0},𝛍_{n})\simeq
H^{q}(\UU,𝛍_{n}){\quoi}.
$$
Since $k$ is separably closed, $A^{0}$ is a strictly local ring, and
we may apply \hyperref[XIX.cor1.1]{1.1} to the ring $A^{0}$. It follows that
$H^{q}$ $(U^{0},𝛍_{n})=0$. But
$$
H^{q}(U^{0},𝛍_{n})\simeq
\varinjlim_{\VV^{0}}H^{q}(\VV^{0},𝛍_{n}){\quoi},
$$
where $\VV^{0}$ ranges over the category of hypercoverings of
$U^{0}$ of order $q+1$,\pageoriginale and $H^{q}$
$(\UU^{0},𝛍_{n})$ is a finite group. Consequently, there exists a
morphism of hypercoverings $\VV^{0}\rightarrow \UU^{0}$ such that
the induced morphism
$$
H^{q}(\UU^{0},𝛍_{n})\longrightarrow
H^{q}(\VV^{0},𝛍_{n})
$$
is the zero morphism. Put
$\VV=\fprod{\VV^{0}}{X}{X^{0}}$. Then $H^{q}$
$(\VV^{0},𝛍_{n})\simeq
H^{q}(\VV,𝛍_{n})$ (by the same argument as
for $\UU^{0}$), so
$$
H^{q}(\UU,𝛍_{n})\longrightarrow
H^{q}(\VV,𝛍_{n})
$$
is also the zero morphism, \cqfd

\medskip
\noindent
\emph{Proof of} \hyperref[XIX.lem1.5]{1.5}. It suffices to find finite
$A^{0}$-algebras $B^{0}_{j}$ and $C^{0}_{j}$ and $A$-isomorphisms
\begin{align*}
B^{0}_{j}\otimes_{A^{0}}A &\xrightarrow{\sim} B_{j}\\
C^{0}_{j}\otimes_{A^{0}}A &\xrightarrow{\sim} C_{j}{\quoi}.
\end{align*}
Indeed, by (\cite{sga4.exp.XIX.bib.3}~\hyperref[XIX.subsec1.4]{1.4}) the morphism
$B_{j}\rightarrow C_{j}$ is induced by a morphism
$B^{0}_{j}\rightarrow C^{0}_{j}$, so the open subset $W_{j}\subset
\Spec B_{j}$ is induced by an open subset $W^{0}_{j}\subset \Spec
B^{0}_{j}$. Moreover, the $U$-morphisms among the $W_{j}$ are induced
by morphisms among the $A$-algebras $B_{j}$, hence
(\cite{sga4.exp.XIX.bib.3}~\hyperref[XIX.subsec1.4]{1.4}) by morphisms among the $B^{0}_{j}$, and therefore by
$U^{0}$-morphisms among the $W^{0}_{j}$. It follows that the
$W^{0}_{j}$ form the set of connected components of a simplicial
object $\UU^{0}$ over $U^{0}$ that induces $\UU$, and one sees
immediately that it is a hypercovering of $U^{0}$.

Now \hyperref[XIX.subsec1.4]{1.4} (ii) does imply that $C_{j}$ is induced by an
$A^{0}$-algebra\pageoriginale $C^{0}_{j}$, as is easily verified
\footnote[1]{Indeed, \hyperref[XIX.subsec1.4]{1.4} (ii) implies that $C_{j}$ is finite over
$k[[X_{1},\ldots,X_{r-1}]]$, and a fortiori over $A^{0}$, and it
suffices to take $C^{0}_{j}=C_{j}$ with the $A^{0}$-module structure
obtained by restriction of scalars.}. For $B_{j}$, it suffices to find
a monic polynomial $f^{0}_{j}(Y)\in A^{0}[Y]$ such that the two
$A$-algebras
$$
A[Y]/(f_{j})\quav
A[Y]/(f^{0}_{j})=A\otimes_{A^{0}}A^{0}[Y]/(f^{0}_{j})
$$
are isomorphic. Indeed, $A^{0}[Y]/(f^{0}_{j})$ is then reduced, and
we take for $B^{0}_{j}$ the normalization (\EGA IV 7.8.3 (vi)) of
$A^{0}[Y]/(f^{0}_{j})$. The ring $B^{0}_{j}\otimes_{A^{0}}A$ is still
normal (\EGA IV 7.8.3 (vii)) and birational to $B_{j}$, hence equal to
that ring.

By \hyperref[XIX.subsec1.4]{1.4} (i), the ideal $I_{j}\subset A$ generated by $d_{j}$ is
induced by an ideal of $k[[x_{1},\ldots,x_{r-1}]]$ $[x_{r}]$
(Weierstrass preparation theorem \cite{sga4.exp.XIX.bib.5}\cite{sga4.exp.XIX.bib.6}), hence by an ideal
$I^{0}_{j}$ of $A^{0}$. Moreover, $A/I_{j}$ is a finite
$k[[x_{1},\ldots,x_{r-1}]]$-algebra, and is therefore isomorphic to
$A^{0}/I^{0}_{j}$. The latter ring is thus complete, and consequently
the $I^{0}_{j}$-adic completion of $A^{0}$ is complete and isomorphic
to $A$. We can therefore find a monic polynomial
$f^{0}_{j}\in A^{0}$ $[Y]$ such that in $A$ $[Y]$ we have
$$
f_{j}\equiv f^{0}_{j}(\text{\rm modulo } d_{j}{}^{2}\rad A){\quoi}.
$$
By (\cite{sga4.exp.XIX.bib.3}, \hyperref[XIX.lem1.3]{1.3}), $A[Y]/(f_{j})\simeq A[Y]/(f^{0}_{j})$, which
proves the lemma.

\section{The case of a strictly local ring}\etiquette[2]{XIX.sec2}

The theorem is as follows:

\begin{enonce*}{Theorem 2.1}\etiquette[2.1]{XIX.thm2.1}
\footnote[1]{Depends on resolution of singularities (\cf
Introduction).}.\pageoriginale Let $A$ be an excellent, regular, strictly
local ring of equal characteristic, and let $x$ be a “local
parameter”\footnote[2]{By which we mean here: an element belonging to a
\emph{regular} system of parameters.} of $A$. Let $U=\Spec
A[1/x]$. Then, for $n$ prime to the characteristic of $A$, we have
$$
H^{q}(U,𝛍_{n})=
\begin{cases}
𝛍_{n}(A) &\text{\rm if}\quad q=0\\
\ZZ/n &\text{\rm if}\quad q=1\\
0 &\text{\rm if}\quad q>1{\quoi}.
\end{cases}
$$
\end{enonce*}

The proof proceeds by induction on $q$. We denote by $P(N)$ the
assertion of \hyperref[XIX.thm2.1]{2.1} for the values $q\leq N$. 
By \hyperref[XIX.lem1.3]{1.3}, the assertion $P(1)$ is true.

\begin{enonce*}{Lemma 2.2}\etiquette[2.2]{XIX.lem2.2}
Suppose that $P(N)$ is true. Let $A\rightarrow A'$ be a local
morphism of excellent, regular, strictly local rings of equal
characteristic. Let $\{x_{1},\ldots,x_{r}\}$ be part of a regular
system of parameters of $A$ whose set of images in $A'$ is again part
of a regular system of parameters. Let
\begin{align*}
U &= \Spec A[1/x_{1},\ldots,1/x_{r}]\\
U'&= \Spec A'[1/x'_{1},\ldots,1/x'_{r}]{\quoi}.
\end{align*}
Then the canonical morphism
$$
H^{q}(U,𝛍_{n})\longrightarrow
H^{q}(U',𝛍_{n})
$$
is\pageoriginale bijective for every $q\leq N$.
\end{enonce*}

\medskip
\noindent
\emph{Proof}. \footnote[3]{A shorter proof than the one that follows,
but using the cup-product formalism, would consist in showing that,
under the hypothesis $P(N)$, the cohomology
$H^{\ast}(U,\ZZ/n\ZZ)$ is, in dimensions $\leq N$, the exterior
algebra (over $\ZZ/n\ZZ)$ on
$H^{1}(U,\ZZ/n\ZZ)\simeq (\ZZ/n\ZZ)^{r}$,
which reduces matters to establishing \hyperref[XIX.lem2.2]{2.2} for $q=1$, a well-known
case. The calculation indicated under the hypothesis $P(N)$ is made
explicit, for example, in \SGA 7.}
We argue by induction on $r$: the assertion is trivial for $r=0$.
Let $r>0$ and suppose the theorem has been proved for $r-1$. Let
$B=A/(x_{r})$, $V=\Spec A[1/x_{1},\ldots,1/x_{r-1}]$, and
$Y=\Spec B[1/x_{1},\ldots,1/x_{r-1}]$ (where the image of $x_{i}$ in
$B$ is denoted by the same symbol). There is an open immersion
$i:U\rightarrow V$, and $Y$ is the complementary closed subset,
defined by the element $x_{r}$.

Applying $P(N)$ to the strict localizations of $V$ at the geometric
points of $Y$, we find that
\begin{equation*}
R^{q}i_{\ast}(𝛍_{n})_{U}=
\begin{cases}
(𝛍_{n})_{V} &\text{\rm if}\quad q=0\\
(\ZZ/n)_{Y} &\text{\rm if}\quad q=1\\
0 &\text{\rm if}\quad 1<q\leq N{\quoi},
\end{cases}\tag{2.3}
\end{equation*}\etiquette[2.3]{XIX.eq2.3}
because these strict localizations are excellent regular rings of
equal characteristic, and the element $x_{r}$ is a local parameter at every
point of $Y$.

We use a “prime” to denote the analogous objects deduced from $A'$.
We obtain a morphism of spectral sequences
$$
\xymatrix@C=.4cm{
E^{pq}_2 \ar[d]^{\varphi^{pq}_2} & = \text{\phantom{ab}} H^{p}(V,R^{q}i_{\ast}
𝛍_n)\ar[d]^{\varphi} \ar@{=>}[r] & H^{p+q}(U,𝛍_n) \ar[d]^{\varphi}&\\
E'^{pq}_2 &= \text{\phantom{ab}} H^{p}(V',R^{q}i'_{\ast}𝛍_n) \ar@{=>}[r]& H^{p+q}(U',𝛍_n).
}
$$
Now\pageoriginale it is clear from \hyperref[XIX.lem1.3]{1.3} and \eqhyperref[XIX.eq2.3]{2.3}
that the base change $V\xleftarrow{g}V'$ induces an isomorphism
$$
g^{\ast}R^{q}i_{\ast}𝛍_{n}\xrightarrow{\sim}
R^{q}i'_{\ast}𝛍_{n}
$$
for $q\leq N$ (only for $q=0$, $1$ is there anything at all to
check). Consequently, the induction hypothesis, applied to the
morphism
$$
H^{q}(V,𝛍_{n})\longrightarrow
H^{q}(V',𝛍_{n}){\quoi},
$$
implies that $\varphi^{pq}_{2}$ is an isomorphism for $q=0$ and
$p\leq N$.

Similarly, we can apply the induction hypothesis to the morphism of
rings $B\rightarrow B'$; these are regular, excellent, strictly local rings of
equal characteristic, and the elements $x_{1},\ldots,x_{r-1}$ form part of a
system of parameters. Therefore
$$

% END INCLUDED SGA4_English_current_buildable_source_20260906/source/19/drafts/SGA4_XIX_LINES_44_450_English.texfrag

% BEGIN INCLUDED SGA4_English_current_buildable_source_20260906/source/19/drafts/SGA4_XIX_LINES_451_850_English.texfrag
H^{q}(Y,𝛍_{n})\xrightarrow{\sim} H^{q}(Y',
𝛍_{n})\quad\text{\rm if}\quad q\leq N{\quoi}.
$$
Since $(𝛍_{n})\simeq (\ZZ/n)_{Y}$, it follows that
$\varphi^{pq}_{2}$ is also bijective for $p\leq N$ and for $q=1$.

Since
$R^{q}i_{\ast}𝛍_{n}=R^{q}i'_{\ast}𝛍_{n}=0$
if $1<q\leq N$, we find $E^{pq}_{2}={E'}^{pq}_{2}=0$ if $1<q\leq N$,
and for $1\leq m\leq N$ the morphism of spectral sequences therefore
reduces to a morphism of exact sequences
$$
\xymatrix@C=.4cm{
H^{m-2}(Y,\ZZ/n) \ar[r] \ar[d]^{a}& H^{m}(V,𝛍_n) \ar[r]\ar[d]^{b}&
H^m(U,𝛍_n) \ar[r]^(.45){e} \ar[d]^{c}&
H^{m-1}(Y,\ZZ/n)  \ar[r] \ar[d]^{d}&
H^{m+1}(V,𝛍_n) \ar[d]\\
H^{m-2}(Y',\ZZ/n) \ar[r] & H^{m}(V',𝛍_n) \ar[r]&
H^m(U',𝛍_n) \ar[r]^(.45){e} & H^{m-1}(Y',\ZZ/n)  \ar[r]&
H^{m+1}(V',𝛍_n).
}
$$
Now\pageoriginale we have already seen that $a$, $b$, and $d$ are
bijective, and it follows that $c$ is \emph{injective}. To prove that
$c$ is \emph{surjective}, it suffices to prove that the morphism
$$
H^{m}(U,𝛍_{n})\xrightarrow{e}H^{m-1}(Y,\ZZ/n)
$$
is surjective for $m\leq N$.

% SGA-SRC-001196
Let $k\subset A$ be an arbitrary separably closed field, and let
$A^{0}=k\{x_{1},\ldots, x_{r}\}$ be the henselization of
$k[x_{1},\ldots,x_{r}]$ at the origin. Then the hypotheses of
\hyperref[XIX.thm2.1]{2.1} are satisfied for the morphism
$A^{0}\rightarrow A$ and the elements $x_{i}$. With the evident
notation, we have a commutative square
$$
\xymatrix@R=1.25cm{
H^{m}(U^{0},𝛍_n)  \ar[r]^(.45){e^0} \ar[d]^{c^0}& H^{m-1}(Y^0,
\ZZ/n) \ar[d]^{d^{0}}\\
H^{m}(U,𝛍_n)  \ar[r]^(.45){e} & H^{m-1}(Y,
\ZZ/n)
}
$$
and we already know that the arrow $d^{0}$ is bijective. It therefore
suffices to prove that the morphism $e^{0}$ is surjective. We are thus
reduced (dropping the superscript ${}^{0}$) to verifying surjectivity
in the case $A=k\{x_{1},\ldots,x_{r}\}$.

For this ring, which is a limit of schemes smooth over $k$, we proved
the purity theorem in \hyperrefext[XVI][XVI.thm3.7]{3.7}. Thus\pageoriginale
$R^{q}i_{\ast}𝛍_{n}=0$ for every $q>1$, and consequently there is an
exact sequence
$$
H^{m}(U,𝛍_{n})\xrightarrow{e}H^{m-1}(Y,\ZZ/n)\rightarrow
H^{m+1}(V,𝛍_{n})\xrightarrow{f}H^{m+1}(U,𝛍_{n})
$$
for every $m\geq 1$. The surjectivity of $e$ is therefore equivalent
to the injectivity of $f$.

Let $Z=\Spec k\{x_{1},\ldots,x_{r-1}\}$
$[1/x_{1},\ldots,1/x_{r-1}]$. The morphism $\Spec
k\{x_{1},\ldots,x_{r}\}\rightarrow \Spec k\{x_{1},\ldots,x_{r-1}\}$
given by the inclusion of rings is acyclic. This is the local
acyclicity of a smooth morphism (\hyperrefext[XVI][XVI.thm1.1]{1.1}).
Since the open subset V of $\Spec k\{x_{1},\ldots, x_{r}\}$ is the
inverse image of $Z$, the morphism
$$
H^{q}(Z,𝛍_{n})\longrightarrow
H^{q}(V,𝛍_{n})
$$
is bijective for every $q$ (\hyperrefext[XV][XV.def1.3]{1.3}). It
therefore suffices to prove that the morphism
$$
H^{q}(Z,𝛍_{n})\longrightarrow
H^{q}(U,𝛍_{n})
$$
is injective for every $q$.

Now $k\{x_{1},\ldots,x_{r}\}$ is the inductive limit of rings
$A_{\alpha}$ étale over $k\{x_{1},\ldots,x_{r-1}\}$
$[x_{r}]$, each equipped with a lifting to that ring of the point
$x_{1}=x_{2}=\ldots=x_{r}=0$. Consequently, $U$ is the projective
limit of the open subsets
$U_{\alpha}=\Spec A_{\alpha}[1/x_{1},\ldots,1/x_{r}]$, and we have
$$
H^{q}(U,𝛍_{n})=\varinjlim_{\alpha}H^{q}(U_{\alpha},
𝛍_{n}){\quoi}.
$$
% SGA-SRC-001198
But since $A_{\alpha}[1/x_{r}]$ is smooth over the henselian ring
$k\{x_{1},\ldots,x_{r-1}\}$,\pageoriginale and since the closed fiber
of $A_{\alpha}[1/x_{r}]$ is nonempty, we see immediately that the
$k\{x_{1},\ldots, x_{r-1}\}$-algebra $A_{\alpha}[1/x_{r}]$ admits
sections. Hence $U_{\alpha}$ admits sections over $Z$, and it follows
that
$$
H^{q}(Z,𝛍_{n})\longrightarrow
H^{q}(U_{\alpha},𝛍_{n})
$$
is injective for every $q$, and therefore that
$$
H^{q}(Z,𝛍_{n})\longrightarrow
H^{q}(U,𝛍_{n})
$$
is also injective, which proves the lemma.

\hypertarget{sga4.exp.XIX.lemma.2.4}{}\begin{enonce*}{Lemma 2.4}\etiquette[2.4]{XIX.lem2.4}
Suppose that $P(N)$ is true. Let $A\rightarrow A'$ be a regular
morphism \emph{(\EGA IV 6.8.1)} of excellent, strictly local rings of
equal characteristic (not necessarily regular), and let
$U\subset \Spec A$ be a regular open subset. Let
$U'=\fprod{U}{\Spec A'}{\Spec A}$. Then, for $n$ prime to the residue
characteristic of $A$, the morphism
$$
H^{q}(U,𝛍_{n})\longrightarrow
H^{q}(U',𝛍_{n})
$$
is bijective for $q\leq N$.
\end{enonce*}

\medskip
\noindent
\emph{Proof.} Let $f:X\rightarrow \Spec A$ be a resolution of
singularities of $\Spec A$ such that the open subset $U$ of $\Spec A$
lifts to an open immersion $i:U\rightarrow X$, and such that the closed
subset $Y=X-U$ is pure of codimension 1 and has normal crossings. We
denote by a prime the base change induced by $A\rightarrow A'$. The
morphism $g:X'\rightarrow X$ is regular (\EGA IV 6.8.3), so
$f':X'\rightarrow \Spec A'$ is a resolution of $\Spec A'$,
$i':U'\rightarrow X'$ is an open immersion, and\pageoriginale
$Y'=X'-U'$ is pure of codimension 1 and has normal crossings. All these
schemes are excellent.

Let $p'$ be a point of $X'$ and let $p$ be its image in $X$. Since $Y$
has normal crossings in $X$, there exist elements
$x_{1},\ldots,x_{r}\in O_{X,p}$ such that $\{x_{1},\ldots,x_{r}\}$ is
part of a system of parameters of $O_{X,p}$ and such that $\Spec
O_{X,p}\cap Y=V(x_{1},\ldots,x_{r})$. Moreover, it follows immediately
from the fact that $g:X'\rightarrow X$ is regular that the images
$x'_{i}$ of the $x_{i}$ in $O_{X',p'}$ form part of a system of
parameters of this ring, and that $\Spec O_{X',p'}\cap
Y'=V(x'_{1},\ldots,x'_{r})$.

Let $B'$ be the strict localization of $O_{X',p'}$ at a geometric point
$\bar{p}$ above $p'$, and let $B$ be the strict localization of
$O_{X,p}$ at the corresponding geometric point of $X$ (these are
excellent rings (\EGA IV 7.8.3 (ii), 7.9.5)). Let
$$
U_{\bar{p}}=\Spec B[1/x_{1},\ldots,1/x_{r}]\quav U'_{\bar{p}}=\Spec
B'[1/x'_{1},\ldots,1/x'_{r}]{\quoi}.
$$
Lemma \hyperref[XIX.lem2.2]{2.2} applies, and we deduce that the
morphisms
\begin{equation*}
H^{q}(U_{\bar{p}},𝛍_{n})\longrightarrow
H^{q}(U'_{\bar{p}},𝛍_{n})\tag{$\ast$}
\end{equation*}
are bijective for $q\leq N$.

Let us compute the sheaves $R^{q}i_{\ast}𝛍_{n}$ and
$R^{q}i'_{\ast}𝛍_{n}$ for the open immersions
$i:U\rightarrow X$ and $i':U'\rightarrow X'$: the stalk
$(R^{q}i'_{\ast}𝛍_{n})_{\bar{p}}$ at the geometric point $\bar{p}$ is
precisely $H^{q}(U'_{\bar{p}},𝛍_{n})$ (VIII 5). Likewise,
$(R^{q}i_{\ast}𝛍_{n})_{\bar{p}}=H^{q}(U_{\bar{p}},𝛍_{n})$. Thus
$(\ast)$ implies that the base-change morphisms
\begin{equation*}
g^{\ast}R^{q}i_{\ast}𝛍_{n}\longrightarrow R^{q}i'_{\ast}𝛍_{n}\tag{${\ast}{\ast}$}
\end{equation*}
are\pageoriginale bijective for $q\leq N$.

By the base-change theorem for a proper morphism
(\hyperrefext[XII][XII.cor5.5]{5.5}), we have (since $X$ and $X'$ are
proper over the spectra of strictly local rings, and the closed fibers
of $X$ and $X'$ are equal)
$$
H^{q}(X,F)\xrightarrow{\sim} H^{q}(X',g^{\ast}F)
$$
for every torsion sheaf $F$ and every $q$. Consequently,
$({\ast}{\ast})$ implies that in the morphism of spectral sequences
$$
\xymatrix@C=.5cm{
E^{pq}_2 \ar[d]& =& H^{p}(X,R^{q}i_{\ast}𝛍_n) \ar@{=>}[r]\ar[d]^{\varphi^{pq}_2}&
H^{p+q}(U,𝛍_n) \ar[d]^{\varphi}\\
E'^{pq}_2   &=&   H^{p}(X',R^{q}i'_{\ast}𝛍_n)\ar@{=>}[r]
& H^{p+q}(U',𝛍_n),
}
$$
the arrows $\varphi^{pq}_{2}$ are bijective for every $q\leq N$.
It is easily seen that this implies that $\varphi^{pq}_{\infty}$ is
bijective if $q\leq N$ and $p+q\leq N+1$. Consequently, the morphism
of the abutments
$$
H^{m}(U,𝛍_{n})\longrightarrow
H^{m}(U',𝛍_{n})
$$
is bijective for $m\leq N$, which proves the lemma.

The following lemma will complete the proof of \hyperref[XIX.thm2.1]{2.1}:

\hypertarget{sga4.exp.XIX.lemma.2.5}{}\begin{enonce*}{Lemma 2.5}\etiquette[2.5]{XIX.lem2.5}
$P(N)\Rightarrow P(N+1)$.
\end{enonce*}

\medskip
\noindent
\emph{Proof}. Let $A$ be strictly local, regular, excellent, and of
equal\pageoriginale characteristic, let $x\in \rad A$ be a local
parameter, and let $U=\Spec A[1/x]$. Let $\hat{A}$ be the completion of
$A$, let $\hat{x}$ be the image of $x$, let
$\hat{U}=\Spec \hat{A}[1/x]$, and let $g:\hat{U}\rightarrow U$ be the
canonical morphism. It suffices to prove that
\begin{equation*}
\begin{split}
g_{\ast}(𝛍_{n})_{\hat{U}}\simeq
(𝛍_{n})_{U}\quad\text{\rm and}\\
R^{q}g_{\ast}𝛍_{n}=0\quad\text{\rm if}\quad 1\leq
q\leq N{\quoi}.
\end{split}\tag{2.6}
\end{equation*}\etiquette[2.6]{XIX.eq2.6}
Indeed, the Leray spectral sequence
$$
E^{pq}_{2}=H^{p}(U,R^{q}g_{\ast}𝛍_{n})\Longrightarrow
H^{p+q}(\hat{U},𝛍_{n})
$$
will give isomorphisms
$$
H^{q}(U,𝛍_{n})\xrightarrow{\sim}
H^{q}(\hat{U},𝛍_{n})
$$
for $q\leq N$, and an exact sequence
$$
0\longrightarrow H^{N+1}(U,𝛍_{n})\longrightarrow
H^{N+1}(\hat{U},𝛍_{n})\longrightarrow
H^{0}(U,R^{N+1}g_{\ast}𝛍_{n}){\quoi}.
$$
But we already know that $P(1)$ is true (\hyperref[XIX.lem1.3]{1.3}),
and it is therefore enough to prove that $H^{q}(U,𝛍_{n})=0$ if
$1<q\leq N+1$. We may thus replace $A$ by $\hat{A}$, which is
isomorphic to a ring of formal power series
$\hat{A}\simeq k[[x_{1},\ldots,x_{r}]]$, with $k$ separably closed and
$\hat{x}=x_{r}$. This case was treated in \hyperref[XIX.thm1.2]{1.2}.

Let us therefore verify \eqhyperref[XIX.eq2.6]{2.6}: the fact that
$g_{\ast}(𝛍_{n})_{\hat{U}}\simeq
(𝛍_{n})_{U}$ is simply the assertion of \eqhyperref[XIX.eq2.3]{2.3}
for $q=0$. Recall that $R^{q}g_{\ast}𝛍_{n}$ is the sheaf associated
with the presheaf $\Rr^{q}$, where
$$
\Rr^{q}(V)=H^{q}(\hat{V},𝛍_{n})
$$
for\pageoriginale $V\rightarrow U$ étale. It suffices to take $V$
separated and of finite type; then $V$ is a disjoint union of integral
schemes, and it suffices to consider such $V$. Let $B$ be the
normalization of $A$ in the field of fractions of
$\Gamma(V,O_{V})$. Then $B$ is a finite $A$-algebra (\EGA IV 7.8.3
(vi)), and we see immediately that V is a regular open subset of
$\Spec B$ and that $\hat{V}$ is the corresponding open subset of
$\Spec \hat{B}$, where $\hat{B}=B\otimes_{A}\hat{A}$. It follows that
the hypotheses of \hyperref[XIX.lem2.4]{2.4} are satisfied for the open
subset $V\subset \Spec B$ and for the morphism
$B\rightarrow \hat{B}$. Hence
$$
\Rr^{q}(V)\simeq H^{q}(V, 𝛍_{n})\simeq
H^{q}(\hat{V}, 𝛍_{n})
$$
for $q\leq N$. But the sheaf associated with the presheaf
$V\mapsto H^{q}(V,𝛍_{n})$ is plainly zero, which proves the lemma.

\section{Purity}\etiquette[3]{XIX.sec3}

We shall adopt terminology analogous to that of
\hyperrefext[XVI][XVI.sec3]{3}:

\hypertarget{sga4.exp.XIX.definition.3.1}{}\begin{enonce*}{Definition 3.1}\etiquette[3.1]{XIX.def3.1}
A \emph{regular pair $(Y,X)$ of codimension $c$} is a closed immersion
$i:Y\rightarrow X$ of locally Noetherian regular schemes such that for
every $y\in Y$ one has $\codim_{y}(Y,X)=c$ (then there exists an open
neighborhood $X'$ of $y$ in $X$ and sections
$x_{1},\ldots,x_{c}\in \Gamma(X',O_{X})$ that generate the ideal of
$Y$, and which form part of a regular system of parameters of $X$ at
every point of $Y\cap X'$). We write $U=X-Y$, and
$j:U\rightarrow X$ for the corresponding open immersion.
\end{enonce*}

A\pageoriginale \emph{morphism} $(Y',X')\rightarrow (Y,X)$ of regular
pairs of codimension $c$ is a morphism $X'\rightarrow X$ such that
$Y'=\fprod{Y}{X'}{X}$. The notion of a regular triple, of a morphism of
regular triples, \etcn, is defined in the analogous evident way.

\hypertarget{sga4.exp.XIX.theorem.3.2}{}\begin{enonce*}{Theorem 3.2}\etiquette[3.2]{XIX.thm3.2}
\footnote[1]{Depends on resolution of singularities (\cf
Introduction).}. \emph{(Cohomological purity)} Let $(Y,X)$ be a regular
pair of codimension $c>0$, where $X$ is an excellent prescheme of equal
characteristic. Let $F$ be a locally constant sheaf on $X$ of cyclic
groups of order $n$ prime to the characteristic of $X$. Then, in the
notation of \emph{(V 4, VIII 5)},
\begin{align*}
ℋ^{q}_{Y} (X,F) &= (R^{q}i^{!})F=0\quad\text{\rm if}\quad q\neq
2c{\quoi},\\
(R^{q}j_{\ast})j^{\ast}F &= 0\quad\text{\rm if}\quad q\neq 0,
2c-1,\quad \text{and}\\
ℋ^{2c}_{Y}(X,F) &= i^{\ast}(R^{2c-1}j_{\ast})j^{\ast}F
\end{align*}
\noindent
is a locally constant sheaf of cyclic groups of order $n$ on $Y$.
\end{enonce*}

\medskip
\noindent
\emph{Proof}. (\cf \hyperrefext[XVI][XVI.thm3.7]{3.7}). We see
immediately that
$$
ℋ^{q}_{Y}(X,F)=0\quad\text{\rm if}\quad q=0,1{\quoi}.
$$
By (V \hyperref[V.cor4.5]{4.5}),
$$
ℋ^{q}_{Y}(X,F)=(R^{q}i^{!})F\simeq
i^{\ast}(R^{q-1}j_{\ast})j^{\ast}F\quad\text{\rm if}\quad q>1{\quoi},
$$
so the assertions for the three sheaves are equivalent. Since the
assertions are local on $X$ for the étale topology, we may
assume\pageoriginale that $F=𝛍_{n}$.

For the case $c=1$, the calculation of the stalks (VIII 5) of
$(R^{q}j_{\ast})j^{\ast}𝛍_{n}$ reduces us to the situation considered
in \hyperref[XIX.thm2.1]{2.1}, whence
$$
((R^{q}j_{\ast})j^{\ast}𝛍_{n})_{\bar{y}}=
\begin{cases}
\ZZ/n &\text{\rm if}\quad q=1\\
0 &\text{\rm if}\quad q>1
\end{cases}
$$
for every geometric point $\bar{y}$ of $Y$. Since these sheaves vanish
outside $Y$ in any case, this proves the assertion for $q>1$. For
$q=1$ $(=2c-1)$, we apply the criterion of
(\hyperrefext[IX][IX.prop2.13]{2.13} (i)) and the explicit calculation
of \hyperref[XIX.lem1.3]{1.3}.

The case $c>1$ is treated by induction (\cf
\hyperrefext[XVI][XVI.thm3.7]{3.7}): let $(Y,X)$ be a regular pair of
codimension $c>1$. By definition, it is clear that, locally on $X$, we
can find a regular triple $(Y,Z,X)$ in which $(Z,X)$ has codimension 1
and $(Y,Z)$ has codimension $c-1$. Let
$$
\xymatrix{
  Y \ar[rr]^{u} \ar[rd]_{i}   &   & Z\ar[ld]^{v}\\
      & X & 
}
$$
be the closed immersions. We have the spectral sequence
$$
E^{pq}_{2}=(R^{p}u^{!})(R^{q}v^{!})F\Longrightarrow
(R^{p+q}i^{!})F{\quoi}.
$$
By the induction hypothesis, $(R^{q}v^{!})F=0$ if $q\neq 2$, and is
locally constant \etc\ldots if $q=2$. Thus, again by the induction
hypothesis, applied to the morphism $u$,
$(R^{p}u^{!})(R^{2}v^{!})F=0$ if $(p,q)\neq (2c-2,2)$\pageoriginale
and is locally constant \etc\ldots if $(p,q)=(2c-2,2)$, whence the
result for the abutment $(R^{p+q}i^{!})F$, \cqfd

\setcounter{subsection}{2}
\subsection{}\etiquette[3.3]{XIX.subsec3.3}
It remains to determine the locally constant sheaves of
(\hyperref[XIX.thm3.2]{3.2}) up to \emph{canonical} isomorphism. For
each prescheme $S$ we introduce the locally constant sheaves
$(𝛍_{n}^{\otimes r})_{S}=$ the $r$-fold tensor product of the sheaf
$(𝛍_{n})_{S}$ over the constant sheaf of rings $\ZZ_{S}$. We also set
$(𝛍_{n}^{\otimes 0})_{S}=(\ZZ/n)_{S}$. It is immediately seen that if
$n$ is invertible on $S$, then $(𝛍_{n}^{\otimes r})_{S}$ is a locally
constant sheaf of cyclic groups of order $n$. We have formulas of the
form $(𝛍_{n}^{\otimes r})\otimes
(𝛍_{n}^{\otimes
s})=(𝛍_{n}^{\otimes r+s})$, and if $S'\rightarrow
S$, then $(𝛍_{n}^{\otimes r})_{S'}$ is canonically isomorphic to the
inverse image on $S'$ of $(𝛍_{n}^{\otimes r})_{S}$. We shall often
omit the symbol $S$ when no confusion can arise.

Let $(Y,X)$ be a regular pair of codimension 1, and suppose that $Y$ is
defined by an equation $x=0$. We have the exact sequence of sheaves
$$
0\longrightarrow (\GG_{m})_{X}\longrightarrow
j_{\ast}(\GG_{m})_{U}\longrightarrow \ZZ_{Y}\longrightarrow
0{\quoi},
$$
whence, by applying (\hyperref[XIX.lem1.3]{1.3}), canonical
isomorphisms
\begin{equation*}
(R^{1}j_{\ast})𝛍_{n}\simeq
j_{\ast}(\GG_{m})_{U}/(j_{\ast}(\GG_{m})_{U})^{n}\simeq
(\ZZ/n)_{Y}{\quoi},\tag{3.3.1}
\end{equation*}\etiquette[3.3.1]{XIX.eq3.3.1}
which plainly do not depend on the choice of $x$. Consequently, we have
a canonical isomorphism
\begin{equation*}
\varphi:\ZZ/n\xrightarrow{\sim}
(R^{2}i^{!})𝛍_{n}\simeq i^{\ast}(R^{1}j_{\ast})𝛍_{n}{\quoi}.\tag{3.3.2}
\end{equation*}\etiquette[3.3.2]{XIX.eq3.3.2}

\begin{enonce*}{Theorem 3.4}\etiquette[3.4]{XIX.thm3.4}

% END INCLUDED SGA4_English_current_buildable_source_20260906/source/19/drafts/SGA4_XIX_LINES_451_850_English.texfrag

% BEGIN INCLUDED SGA4_English_current_buildable_source_20260906/source/19/drafts/SGA4_XIX_LINES_851_1250_English.texfrag
\footnote[1]{Depends on resolution of singularities (\cf
Introduction). This theorem may also be regarded as the conjunction of
\hyperref[XIX.thm3.2]{3.2} and the theory of variance of \Exp XVIII by means of
``local fundamental classes,'' where homomorphisms
$\varphi(Y,X)$ (not necessarily bijective) are studied for $Y$ locally a
complete intersection in $X$.}.\pageoriginale Let $C$ be the family
of regular pairs of excellent schemes of equal characteristic. There exists a
unique function $\varphi$ assigning to every regular pair $(Y,X)$ of
codimension $c$ an isomorphism (called canonical) of sheaves
$$
\varphi(Y,X)=\varphi:\ZZ/n\xrightarrow{\sim}
ℋ^{2c}_{Y}(X,𝛍_{n}^{\otimes
c})=(R^{2c}i^{!})𝛍_{n}^{\otimes c}
$$
such that $\varphi$ satisfies the following conditions:

\begin{enumerate}
\renewcommand{\theenumi}{\roman{enumi}}
\renewcommand{\labelenumi}{(\theenumi)}
\item If $c=1$, and if $Y$ is defined in $X$ by an equation $x=0$,
% SGA-SRC-001210
then $\varphi$ is the morphism \eqhyperref[XIX.eq3.3.2]{3.3.2}.

\item $\varphi$ is compatible with smooth morphisms of pairs; that is, if
$(Y',X')\rightarrow (Y,X)$ is a morphism, let $g:Y'\rightarrow Y$ be the
induced morphism. The diagram of $Y'$-sheaves
$$
\xymatrix{
\ZZ/n  \ar[rr]^{\varphi(Y',X')} & &
(R^{2c}i^{!})𝛍_n^{{\boxtimes}c}\\
g^{\ast}(\ZZ/n) \ar[u]\ar[rr]^{g^{\ast}\varphi(Y,X)}& &
g^{\ast}(R^{2c}i^{!})𝛍_n^{{\boxtimes}c} \ar[u]
}
$$
is commutative (which implies in particular that the right-hand vertical
arrow is an isomorphism).

\item \emph{Transitivity:} let $(Z,Y,X)$ be a regular triple, that is, a
commutative diagram of closed immersions of regular preschemes \etcn,
$$
\xymatrix{
Z \ar[rr]^{v}\ar[rd]_{w}&    &Y\ar[ld]^{u}\\
            & X  & .
}
$$
Let\pageoriginale $a$, $b$, $c$ be the codimensions of $u$, $v$, $w$,
respectively \emph{(thus $c=a+b$)}. We have an isomorphism
\begin{align*}
\varphi(Y,X)\otimes 𝛍_{n} {}^{\otimes
b}:𝛍_{n} {}^{\otimes b} \longrightarrow &
 (R^{2a}u^{!})(𝛍_{n}^{\otimes a}) \otimes
𝛍_{n} {}^{\otimes b}\simeq\\
& \simeq (R^{2a}u^{!})𝛍_{n}{}^{\otimes c}{\quoi}.
\end{align*}
The spectral sequence $(R^{p}v^{!})(R^{q}u^{!})F\Rightarrow
(R^{p+q}w^{!})F$ gives an isomorphism
$$
\xi : (R^{2b}v^{!})(R^{2a}u^{!})𝛍_{n}{}^{\otimes
c}\xrightarrow{\sim} (R^{2c}w^{!})𝛍_{n}{}^{\otimes
c}{\quoi}.
$$
We thus obtain a diagram of isomorphisms
% SGA-SRC-001212 — tensor typography checked against the scan; missing \varphi reconstructed from transitivity
$$\vcenter{\etiquette[3.4.1]{XIX.eq3.4.1}
\xymatrix{
\ZZ/n \ar[r]^{\varphi(Z,Y)} \ar[d]^{\varphi(Z,X)}&
(R^{2b}v^{!})𝛍_{n}^{\otimes b} \ar[d]^{\eta =
  (R^{2b}v^{!})(\varphi(Y,X)\otimes 𝛍_{n}^{\otimes b})}\\
(R^{2c}w^{!})𝛍_{n}^{\otimes c} &
(R^{2b}v^{!})(R^{2a}u^{!})𝛍_{n}^{\otimes c}\ar[l]
}}\leqno{(3.4.1)}
$$
and the transitivity assertion is that this diagram is commutative.
\end{enumerate}
\end{enonce*}

\medskip
\noindent
\emph{Proof.} The proof of existence and uniqueness proceeds by induction on
the codimension $c$. For $c=1$, we define the morphism $\varphi$ locally by
\eqhyperref[XIX.eq3.3.1]{3.3.1}. Since this morphism does not depend on the
choice of the parameter $x$, gluing gives an isomorphism $\varphi$ for every
$(Y,X)$ of codimension 1; and it is clear from the definition of the morphism
and from \hyperref[XIX.lem1.3]{1.3} that the $\varphi$ so constructed satisfies
(ii). Transitivity does not enter when $c=1$.

Now\pageoriginale suppose that existence and uniqueness have already been
proved when the codimension is $<c$, and that $c>1$. Applying the induction
hypothesis, we immediately find that in the situation of
\hyperref[XIX.thm3.4]{3.4} (iii), the three arrows of
\eqhyperref[XIX.eq3.4.1]{3.4.1} that are already defined, \ie
$\varphi (Z,Y)$, $\xi$, $\eta$, are compatible with morphisms of regular
triples. In this situation set
$$
\Psi (Z,Y,X)=\xi \eta \varphi (Z,Y){\quoi},
$$
so $\Psi$ too is compatible with morphisms of triples. Now, if $(Z,X)$ is a
regular pair of codimension $c>1$, locally on $X$ one can find a triple
$(Z,Y,X)$ with $0<\codim(Y,X)<c$. It therefore remains only to prove that the
arrow $\Psi(Z,Y,X)$ is independent of $Y$. Indeed, this will prove the
uniqueness of $\varphi$ in codimension $c$, and existence will then follow by
gluing.

If $(Z,Y,X)$ is arbitrary with $\codim(Y,X)>1$, locally on $X$ one can find a
regular triple $(Y,W,X)$ in which $(W,X)$ has codimension 1. A transitivity
diagram chase, which we leave to the reader, therefore reduces us to checking
independence of $Y$ when $(Y,X)$ has codimension 1.

Let $(Z,Y_{0},X)$ and $(Z,Y_{1},X)$ be two such triples. It suffices to check
the assertion fiber by fiber, and we are thus reduced to the case in which
$X$ is strictly local and $Y_{i}$ is defined by an equation, say $f_{i}=0$
$(i=0,1)$. Let $\bar{X}=\sheaf{\Spec}O_{X}[t]$,
$\bar{Z}=\fprod{Z}{\bar{X}}{X}$, and let $\bar{Y}$ be the closed subscheme of
$\bar{X}$ defined by the equation
$$
(f_{1}-f_{0})t+f_{0}=0{\quoi}.
$$
Then\pageoriginale $\bar{Y}$ is plainly regular at the points lying over the
sections $\epsilon_{0}:\{t=0\}$ and $\epsilon_{1}:\{t=1\}$ of
$\bar{X}/X$. Let $C$ be the closed subset (\EGA IV 7.8.3 (iv)) of $\bar{Y}$
consisting of the points at which $\bar{Y}$ is not regular, and replace
$\bar{X}$, $\bar{Y}$, $\bar{Z}$ by $\bar{X}-C$, $\bar{Y}-C$,
$\bar{Z}-C\cap \bar{Z}$, respectively. Then $(\bar{Z},\bar{Y},\bar{X})$ is a
regular triple, and the sections $\epsilon_{i}$ give morphisms of triples
$$
(Z,Y_{i},X)\longrightarrow (\bar{Z},\bar{Y},\bar{X})\quad
(i=0,1){\quoi}.
$$

Let $F=(R^{2c}w^{!})𝛍_{n}{}^{\otimes c}$ be the sheaf on $Z$.
It is a locally constant sheaf isomorphic to $\ZZ/n$ (by either of the
isomorphisms $\Psi (Z,Y_{i},X)$), and one sees immediately that the sheaf
$\bar{F}=(R^{2c}\bar{w}^{!})𝛍_{n}{}^{\otimes c}$ on $\bar{Z}$ is the inverse
image on $\bar{Z}$ of $F$. Since the fibers of $\bar{Z}\rightarrow Z$ are
connected and nonempty and $\bar{Z}\rightarrow Z$ is smooth, hence open, it
follows that the induced morphism
$$
H^{0}(Z,F)\longrightarrow H^{0}(\bar{Z},\bar{F})
$$
is \emph{bijective}; hence each of the morphisms
$$\vcenter{
\xymatrix@R=.5cm{
                 &                                & H^{0}(Z,F)\\
H^{0}(Z,F) \ar[r]& H^{0}(\bar{Z},\bar{F})\ar[ru]^{\epsilon^{\ast}_0} \ar[rd]_{\epsilon^{\ast}_1} &           \\
                 &                                & H^{0}(Z,F)
}}\leqno{(*)}
$$
is bijective, and consequently $\epsilon^{\ast}_{0}=\epsilon^{\ast}_{1}$.

To give an isomorphism of a sheaf $F$ with $\ZZ/n$ is equivalent to giving a
global section (the image of 1) that is a generator. Let $\bar{\alpha}$,
$\alpha_{i}$ be the global sections of $\bar{F}$, $F$ given by the
isomorphisms\pageoriginale $\Psi(\bar{Z},\bar{Y},\bar{X})$,
$\Psi(Z,Y_{i},X)$, respectively. Since $\Psi(.,.,.)$ is compatible with
morphisms, the diagram $(\ast)$ implies that
$$
\alpha_{0}=\epsilon^{\ast}_{0}(\bar{\alpha})=\epsilon^{\ast}_{1}(\bar{\alpha})=\alpha_{1}{\quoi},
$$
whence $\Psi(Z,Y_{0},X)=\Psi(Z,Y_{1},X)$, \cqfd

\section{Local acyclicity of a regular morphism}\etiquette[4]{XIX.sec4}

\begin{enonce*}{Theorem 4.1}\etiquette[4.1]{XIX.thm4.1}
\footnote[1]{Depends on resolution of
singularities (\cf Introduction).} \emph{(local acyclicity).} Let
$g:X'\rightarrow X$ be a regular morphism \emph{(\EGA IV 6.8.1)} of
excellent schemes of equal characteristic. Then $g$ is universally locally
acyclic for $n$ prime to the characteristic.
\end{enonce*}

\medskip
\noindent
\emph{Proof}. Since the hypothesis is stable under base change of finite
type, it follows at once from \hyperrefext[XV][XV.prop1.13]{1.13} ii) that it
suffices to prove that $g$ is locally acyclic for $n$. Apply the criterion of
\hyperrefext[XV][XV.cor1.17]{1.17}: let $\bar{x}'$ be a geometric point of
$X'$ and $\bar{x}$ the corresponding geometric point of $X$. It follows
immediately from the definition of a regular morphism that the morphism of
strict localizations at these points induced by $g$ is again regular. We are
thus reduced to the case in which $X=\Spec A$ and $X'=\Spec A'$ are strictly
local, $g$ is a local morphism, and by \hyperrefext[XV][XV.cor1.17]{1.17} it
suffices to prove that the\pageoriginale geometric fibers of $g$ are acyclic
for $n$. Replacing $X$ by an integral closed subscheme and $X'$ by the
corresponding closed subscheme reduces us to the case of the geometric generic
fiber $X'_{\bar{x}}$. Take $\bar{x}=\Spec \bar{K}$, where $\bar{K}$ is the
separable closure of the field of rational functions on $X$. Then
$$
\bar{x}=\varprojlim_{\alpha}U_{\alpha}
$$
where $U_{\alpha}$ ranges over a filtered category of connected étale
$X$-schemes, which may be assumed \emph{regular} (\EGA IV 7.8.3 (v)). Let
$B_{\alpha}$ be the normalization of $A$ in the field of rational functions
$R(U_{\alpha})$. Set $U'_{\alpha}=\fprod{X'}{U_{\alpha}}{X}$ and
$B'_{\alpha}=A'\otimes_{A}B_{\alpha}$. Then $U_{\alpha}$ is a regular open
subscheme of $\Spec B_{\alpha}$ and $B_{\alpha}\rightarrow B'_{\alpha}$ is
% SGA-SRC-001217
regular, so by \hyperref[XIX.lem2.4]{Lemma 2.4}
$$
H^{q}(U_{\alpha},𝛍_{n})\xrightarrow{\sim}
H^{q}(U'_{\alpha},𝛍_{n})
$$
for every $q$. Now the fiber $X'_{\bar{x}}$ is the limit of the
$U'_{\alpha}$, hence
\begin{align*}
H^{q}(X'_{\bar{x}},𝛍_{n}) &= \varinjlim_{\alpha}
H^{q}(U'_{\alpha},𝛍_{n})\\
&= \varinjlim_{\alpha} H^{q}(U_{\alpha}, 𝛍_{n})\\
 &= H^{q}(\bar{x},𝛍_{n})=0\quad\mbox{\rm if}\quad q>0,\cqfd.
\end{align*}

\begin{enonce*}{Corollary 4.2}\etiquette[4.2]{XIX.cor4.2}
(base change by a regular morphism). Let
$$
\xymatrix@C=2.25cm@R=1cm{
X \ar[d]^{f}& X'\ar[l]_{g'} \ar[d]^{f'}\\
Y & Y'\ar[l]_{g}
}
$$
be a Cartesian diagram, where $g:Y'\rightarrow Y$ is a regular morphism of
excellent preschemes of equal characteristic, and $f$ is quasi-compact and
quasi-separated.\pageoriginale Then for every abelian torsion sheaf $F$ whose
torsion is prime to the characteristic, the base-change morphisms
\emph{(XII \hyperref[XII.eq4.2]{4.2})}
$$
g^{\ast}(R^{q}f_{\ast})F\longrightarrow (R^{q}f'_{\ast}){g'}^{\ast}F
$$
are bijective for every $q\geq 0$.
\end{enonce*}

This follows from \hyperref[XIX.thm4.1]{4.1} and
\hyperrefext[XVI][XVI.thm1.1]{1.1}.

From \hyperref[XIX.thm4.1]{4.1} and \hyperrefext[XV][XV.cor1.17]{1.17} we
immediately deduce

\begin{enonce*}{Corollary 4.3}\etiquette[4.3]{XIX.cor4.3}
Let $A$ be a strictly local excellent ring of equal characteristic, and let
$f:X\rightarrow \Spec A$ be a morphism. Let $A\rightarrow A'$ be a regular
morphism, where $A'$ too is strictly local, excellent, and of equal
characteristic. Denote by a $'$ the base change induced by $A\rightarrow A'$.
For every torsion sheaf $F$ on $X$, of torsion prime to the characteristic of
$A$, we have
$$
H^{q}(X,F)\xrightarrow{\sim} H^{q}(X',F')
$$
for every $q\geq 0$.
\end{enonce*}

This corollary applies in particular when $A'$ is the completion $\hat{A}$ of
$A$. (In this case, the hypothesis ``$A$ strictly local'' can moreover be
replaced by ``$A$ henselian local,'' as is seen immediately.)

\section{Finiteness theorem}\etiquette[5]{XIX.sec5}\pageoriginale

\begin{enonce*}{Theorem 5.1}\etiquette[5.1]{XIX.thm5.1}
\footnote[1]{\hypertarget{sga4.exp.XIX.note.resolution}{}Depends on resolution of singularities (\cf
Introduction).} \emph{(finiteness).} Let $f:X\rightarrow S$ be a morphism of
finite type between excellent schemes of equal characteristic, and let $F$ be
a constructible abelian torsion sheaf on $X$, of torsion prime to the
characteristic. Then $(R^{q}f_{\ast})F$ is likewise constructible for every
$q$.
\end{enonce*}

This result generalizes the theorem of \hyperrefext[XVI][XVI.thm5.1]{5.1}.
The proof is nearly the same: one applies the purity result
\hyperref[XIX.thm3.2]{3.2} instead of \hyperrefext[XVI][XVI.thm3.7]{3.7}.

\begin{enonce*}{Corollary 5.2}\etiquette[5.2]{XIX.cor5.2}
Let $A$ be an excellent strictly local ring of equal characteristic and let
$f:X\rightarrow \Spec A$ be a morphism of finite type. Let $F$ be a
constructible abelian sheaf on $X$, of torsion prime to the characteristic of
$A$. Then the groups $H^{q}(X,F)$ are finite for every $q$.
\end{enonce*}

\section{Cohomological dimension of affine morphisms}\etiquette[6]{XIX.sec6}

Let $f:X\rightarrow Y$ be a morphism of finite type between excellent
preschemes, and let $F$ be a sheaf on $X$. Recall that we defined
(\hyperrefext[XIV][XIV.def2.2]{2.2}) an integer $\delta(F)$ when $Y$ is local.
We may therefore regard\pageoriginale $\delta$ as a function $\delta(F,f)$
on an arbitrary $Y$, by setting
\begin{equation*}
\delta(F,f,y)=\delta(F'){\quoi},\tag{6.0}
\end{equation*}\etiquette[6.0]{XIX.eq6.0}
where $Y'=\Spec O_{Y,y}$, $X'\rightarrow Y'$ is the morphism obtained from
$f$ by the base change $Y'\rightarrow Y$, and $F'$ is the inverse-image sheaf
of $F$ on $X'$.

We have the following result, which generalizes
\hyperrefext[XVI][XVI.def3.1]{3.1}:

\begin{enonce*}{Theorem 6.1}\etiquette[6.1]{XIX.thm6.1}
\sgarepeatedfootnotemark{sga4.exp.XIX.note.resolution}. Let $f:X\rightarrow Y$ be an affine morphism of finite type,
where $Y$ is an excellent scheme of equal characteristic, and let $F$ be an
abelian torsion sheaf on $X$. Then $R^{q}f_{\ast}F$ vanishes at every point
$y\in Y$ such that $\delta(F,f,y)<q$.
\end{enonce*}

The local variant, plainly equivalent to \hyperref[XIX.thm6.1]{6.1}, is the
following:

\begin{enonce*}{Theorem 6.1 bis}\etiquette[6.1 bis]{XIX.thmbis6.1}
\sgarepeatedfootnotemark{sga4.exp.XIX.note.resolution}. Let $A$ be a strictly local excellent ring of equal
characteristic, and let $f:X\rightarrow Y=\Spec A$ be an affine morphism of
finite type. Let $F$ be a torsion sheaf on $X$. Then
$$
H^{q}(X,F)=0\quad\text{\rm if}\quad q>\delta(F){\quoi}.
$$
\end{enonce*}

Recall (\hyperrefext[X][X.thm5.1]{5.1}) that assertions of this kind are of no
interest for $F$ of $p$-torsion, where $p$ is the characteristic.\pageoriginale
We may therefore suppose that $F$ is of torsion prime to the characteristic.

We have the following corollary:

\begin{enonce*}{Corollary 6.2}\etiquette[6.2]{XIX.cor6.2}
Let $A$ be a henselian, excellent ring of equal characteristic, with residue
field $k$, and let $U\subset \Spec A$ be an affine open subscheme. Let
$\ell \in \PP$. We have
$$
\cd_{\ell}U\leq \cd_{\ell}k+\dim A{\quoi}.
$$
\end{enonce*}

Indeed, let $\widetilde{A}$ be the strict localization of $A$, which is an
infinite Galois étale covering with group $G=G(\bar{k}/k)$, and let
$\widetilde{U}$ be the affine open subscheme of $\Spec \widetilde{A}$ induced
by $U$. Applying the Hochschild--Serre spectral sequence (VIII \hyperref[VIII.prop8.4]{8.4})
$$
H^{p}(G,H^{q}(\widetilde{U},.))\Longrightarrow H^{p+q}(U,.)
$$ 
reduces us to treating the case $A=\widetilde{A}$, and this is then a special
case of \hyperref[XIX.thm6.1]{6.1} bis.

With the notation above, let $R(A)$ be the ring of rational functions on
$A$. We have $\Spec R(A)=\varprojlim_{\alpha}U_{\alpha}$, where $U_{\alpha}$
is an affine open subscheme of $\Spec A$. The theory of passage to the limit
therefore gives

\begin{enonce*}{Corollary 6.3}\etiquette[6.3]{XIX.cor6.3}
Let $A$ be henselian, excellent, and of equal characteristic, with residue
field $k$. Let $R(A)$ be the field of rational functions on $\Spec A$. Then,
for $\ell \in \PP$, we have
$$
\cd_{\ell}R(A)\leq \cd_{\ell}k+\dim A
$$
with\pageoriginale equality if $\ell$ is invertible in $k$, and if moreover
$\ell\neq 2$ or $k$ is not orderable.
\end{enonce*}

Equality in the latter case follows from \hyperrefext[X][X.cor2.4]{2.4}.

Finally, applying \hyperrefext[X][X.cor4.2]{4.2},
\hyperrefext[X][X.cor4.4]{4.4}, we find

\begin{enonce*}{Corollary 6.4}\etiquette[6.4]{XIX.cor6.4}
Let $X$ be a noetherian excellent scheme of equal characteristic, and let
$R(X)$ be the ring of rational functions on $X$. Let $\ell\in \PP$ be
invertible on $X$, and suppose that $\ell\neq 2$ or that no residue field of
$X$ is orderable. Then
$$
\cd_{\ell}(X)\leq \cd_{\ell}R(X)+\dim X\leq 2\cd_{\ell}R(X){\quoi}.
$$
\end{enonce*}

\medskip
\noindent
\emph{Proof of} \hyperref[XIX.thm6.1]{6.1}. We begin with a reduction
analogous to that of \hyperrefext[XIV][XIV.sec3]{3}. Induction on $\delta$:
Clearly, \hyperref[XIX.thm6.1]{6.1} for $\delta\leq d$ is equivalent to
(\hyperref[XIX.thm6.1]{6.1} bis) for $\delta\leq d$. For $\delta=0$ this is
trivial. We treat the case $\delta\leq 1$ in the form
\hyperref[XIX.thm6.1]{6.1} bis. Since $F$ is the filtered direct limit of its constructible
subsheaves (\hyperrefext[IX][IX.prop2.9]{2.9} (iii)), we may suppose that $F$
is constructible; then the support of $F$ is contained in a closed subset
$Z\subset X$ with $\delta(z)\leq 1$ for every $z\in Z$
(\hyperrefext[XIV][XIV.def2.2]{2.2}). Each reduced irreducible component
$Z_{i}$ of $Z$ is then either an affine scheme of finite type and dimension
$\leq 1$ over the closed point $y$ of $Y$, or a scheme quasi-finite over a
reduced closed subscheme $C$ of $Y$ of dimension 1. In the latter case,
$C=\Spec B$, where $B$ is a strictly local integral domain of dimension 1.
Such a scheme $Z_{i}$ is the spectrum of a strictly local integral domain of
dimension $\leq 1$, or of a field finite over the fraction field of $B$.
Thus in all cases the cohomological dimension of $Z_{i}$ is at
most\pageoriginale 1, by \hyperrefext[X][X.thm2.3]{2.3} and
\hyperrefext[IX][IX.cor5.7]{5.7}. We have an exact sequence
$$
0\longrightarrow F\longrightarrow \prod_{i} F_{i}\longrightarrow

% END INCLUDED SGA4_English_current_buildable_source_20260906/source/19/drafts/SGA4_XIX_LINES_851_1250_English.texfrag

% BEGIN INCLUDED SGA4_English_current_buildable_source_20260906/source/19/drafts/SGA4_XIX_LINES_1251_1650_English.texfrag
\epsilon \longrightarrow 0
$$
where $F_{i}$ is the restriction of $F$ to $Z_{i}$, from which the
result for $Z$ follows.

Suppose now that the theorem has been proved for the values
$\delta<d$, and let us prove it for $\delta(F)\leq d$, assuming $d\geq 2$.

\hypertarget{sga4.exp.XIX.lemma.6.5}{}\begin{enonce*}{Lemma 6.5}\etiquette[6.5]{XIX.lem6.5}
It suffices to treat the case where
$X=\EE^{1}_{Y}=\sheaf{\Spec}O_{Y}[t]$ is affine space of
dimension $1$, and where $Y$ is strictly local.
\end{enonce*}

The proof is that of \hyperrefext[XIV][XIV.lem4.2]{4.2}: It is clear that one may
suppose $Y$ strictly local, and that $\delta(F)\leq d$. We must
prove that $H^{q}(X,F)=0$ if $q>d$. We may embed $X$ in
$\EE^{N}_{Y}$, so it suffices to treat the case
$X=\EE^{N}_{Y}$. Induction on $N$: Consider the diagram
$$
\xymatrix@C=1.2cm@R=1cm{
\EE^{N}_Y \ar[rr]^{g} \ar[rd]_{f}&   & \EE^{N-1}_Y \ar[ld]^{h}\\
                     & Y   &  
}
$$
so that $\EE^{N}_{Y}$ is affine space of dimension 1
over $\EE^{N-1}_{Y}$. Applying the induction hypothesis and
the Leray spectral sequence
$$
H^{p}(\EE^{N-1},R^{q}g_{\ast}F)\Longrightarrow
H^{p+q}(\EE^{N},F)
$$
shows that it suffices to prove that
$$
\delta(R^{q}g_{\ast}F)\leq d-q{\quoi}.
$$
This\pageoriginale means that if $z\in \EE^{N-1}$ is a
point such that $\delta(z)>d-q$, then $R^{q}g_{\ast}F=0$ at the
point $z$. But the induction hypothesis on $N$ applies to the morphism
$g$, and hence $R^{q}g_{\ast}F=0$ at the point $z$ if $\delta(F,g,z)<
q$. Thus it suffices to prove the inequality
$$
\delta(F,g,z)+\delta(z)\leq \delta(F){\quoi}.
$$
We leave this pleasure to the reader.

\hypertarget{sga4.exp.XIX.lemma.6.6}{}\begin{enonce*}{Lemma 6.6}\etiquette[6.6]{XIX.lem6.6}
To prove $\hyperref[XIX.thm6.1]{6.1}$ for the values $\delta\leq d$, it suffices to
prove the following: Let $B$ be a complete normal local ring of equal
characteristic and dimension $d$. Let $U\subset \Spec B$ be an
affine open subset. Suppose that there exists an element $x\in \rad B$ which is
invertible on $U$. Then
$$
H^{q}(U,\ZZ/\ell)=0\quad\text{\rm if}\quad q>d{\quoi}.
$$
\end{enonce*}

This lemma is moreover a special case of Corollary \hyperref[XIX.cor6.2]{6.2}.

\medskip
\noindent
\emph{Proof} (\cf \hyperrefext[XIV][XIV.lem4.4]{4.4}). By \hyperref[XIX.lem6.5]{6.5},
it suffices to treat the case where $Y$ is strictly local and
$X=\EE^{1}_{Y}$. Consider the diagram
$$
\xymatrix@C=1cm@R=.8cm{
\EE^{1}_Y \ar[rr]^{i} \ar[rd]_{f}&   & \PP^{1}_Y \ar[ld]^{g}\\
          & Y &,
}
$$
where $\PP^{1}_{Y}$ is the projective line and $\EE^{1}_{Y}$
is obtained from $\PP^{1}$ by removing the section $Y_{\infty}$ at
infinity. Let $F$ be a torsion sheaf on $\EE^{1}$, with
$\delta(F)\leq d$. We have the spectral sequence 
$$
H^{p}(\PP^{1},R^{q}i_{\ast}F)\Longrightarrow
H^{p+q}(\EE^{1},F){\quoi}.
$$
Now\pageoriginale the $R^{q}i_{\ast}F$ for $q>0$ are concentrated on
the section $Y_{\infty}$, which is isomorphic to $Y$. Since $Y$ is
strictly local, it follows that $H^{p}(\PP^{1},
R^{q}i_{\ast}F)=0$ if $p>0$ and $q>0$. Furthermore, by \hyperrefext[XII][XII.cor5.3]{5.3}, we have
$H^{p}(\PP^{1}, i_{\ast}F)=0$ if $p>2$. Since we wish to prove
that $H^{n}(\EE^{1},F)=0$ if $n>d$, and $d\geq 2$, it suffices
to prove that
$$
H^{0}(\PP^{1},
R^{q}i_{\ast}F)=H^{0}(Y_{\infty},R^{q}i_{\ast}F)=0\quad\text{\rm
if}\quad q>d{\quoi}.
$$
This group is isomorphic to the stalk of $R^{q}i_{\ast}F$ at the
point $y_{\infty}$ at infinity of $\PP^{1}$ in the closed fiber.

Now it suffices (VII \hyperref[VII.prop3.3]{3.3} and \hyperrefext[IX][IX.prop2.9]{2.9} (iii)) to treat the case where $F$
is constructible. Then, since $\delta(F)\leq d$, there exists a
closed subscheme $X$ of $\PP^{1}$ with $\dim X\leq d$ such that
$X$ contains the support of $F$. Let $\Spec A=\widetilde{X}$ be the
strict localization of $X$ at the point $y_{\infty}$,
$\widetilde{U}=\widetilde{X}-\fprod{\widetilde{X}}{Y_{\infty}}{\PP^{1}}$,
and let $\widetilde{F}$ be the sheaf induced by $F$ on $\widetilde{U}$. Note that $\widetilde{U}$
is obtained from $\widetilde{X}$ by localizing with respect to an element
$x\in \Gamma(X,O_{X})$. We have
$$
(R^{q}i_{\ast}F)_{y_{\infty}}\simeq
H^{q}(\widetilde{U},\widetilde{F}){\quoi}.
$$
Thus it suffices to prove that $H^{q}(\widetilde{U},\widetilde{F})=0$, if $q>d$,
for every torsion sheaf $\widetilde{F}$ on $\widetilde{U}$.

We may now apply \hyperrefext[IX][IX.prop5.6]{5.6}, with $\varphi=\delta$, and we
find that it suffices to prove that
$$
H^{q}(V,\ZZ/\ell)=0\quad\text{\rm if}\quad q>\delta(V)
$$
for\pageoriginale every $V$ finite and integral over
$\widetilde{U}$ and for invertible $\ell$. Let $x:\bar{V}\rightarrow
V$ be the normalization of $V$. We have the exact sequence
$$
0\longrightarrow (\ZZ/\ell)_{V}\longrightarrow
x_{\ast}(\ZZ/\ell)_{\bar{V}}\longrightarrow C\longrightarrow 0
$$
with $\delta(C)<\delta(V)$, and it follows from this (by induction on
$\delta$) that it suffices to take $V$ \emph{normal}.

Let $B$ be the normalization of $A$ in the field of rational functions
on $V$. Then $B$ is an excellent strictly local ring of equal
characteristic, and $V$ is an affine open subset of $\Spec B$, obtained by
localizing with respect to $x$. Let $\hat{B}$ be the completion of $B$, and
$\hat{V}$ the induced open subset (which is obtained from $\Spec \hat{B}$ by
localizing with respect to the image of $x$ in $\hat{B}$). By \hyperref[XIX.cor4.3]{4.3},
we have
$$
H^{q}(V,\ZZ/\ell)\simeq H^{q}(\hat{V},\ZZ/\ell){\quoi}.
$$
We are therefore reduced to proving that
$$
H^{q}(\hat{V},\ZZ/\ell)=0\quad\text{\rm if}\quad q>\delta
(V)=\dim B=\dim \hat{B}{\quoi}.
$$
By the induction hypothesis, the theorem is true if $\delta(V)<d$,
so Theorem \hyperref[XIX.thm6.1]{6.1} bis applies in this situation. Consequently,
it suffices to treat the case $\dim \hat{B}=d$, and since
$\hat{B}$ is normal, we are in the situation of the lemma, whence the
result.

\setcounter{subsection}{7}
\setcounter{subsubsection}{-1}
\subsubsection{}\etiquette[6.7.0]{XIX.subsubsec6.7.0}
We shall first prove \hyperref[XIX.lem6.6]{6.6} under additional
hypotheses, when 
$B$ has characteristic $p>0$. Let $A$ be an
equidimensional, complete ring of equal characteristic, and let $k\subset A$ be a
field such that $A/\rad A$ is a finite extension of $k$. We say that
the\pageoriginale $k$-algebra $A$ is \emph{formally separable}
if there exists a $k$-subalgebra of $A$ of the form
$k[[x_{1},\ldots,x_{n}]]$ such that the morphism $\Spec A\rightarrow
\Spec k[[x]]$ is finite and étale at the generic point of $\Spec k[[x]]$
$(x=(x_{1},\ldots,x_{n}))$.

\hypertarget{sga4.exp.XIX.lemma.6.7}{}\begin{enonce*}{Lemma 6.7}\etiquette[6.7]{XIX.lem6.7}
Let $k\subset B$ be a field over which $B/\rad B$ is
finite. Under the induction hypothesis on $d$ above, the conclusion
of \hyperref[XIX.lem6.6]{6.6} holds if the hypotheses of $\hyperref[XIX.lem6.6]{6.6}$ are satisfied and if,
in addition, either of the following conditions is satisfied:
\begin{enumerate}
\renewcommand{\theenumi}{\roman{enumi}}
\renewcommand{\labelenumi}{(\theenumi)}
\item $k$ has characteristic $0$.

\item There exists an element $x\in \rad B$ which is invertible on $U$ and
such that $\bar{B}=B/(x)$ is a formally separable $k$-algebra.
\end{enumerate}
\end{enonce*}

\medskip
\noindent
\emph{Proof}. Choose elements $x_{1}=x$,
$x_{2},\ldots, x_{d}$ of $\rad B$ such that $\{x_{1},\ldots,x_{d}\}$
generates a primary ideal for $\rad B$. Then $B$ is a finite $A$
-algebra, where $A=k[[x_{1},\ldots,x_{d}]]$, and $\bar{B}$ is the
$k[[x_{2},\ldots, x_{d}]]=\bar{A}$-algebra $B/x_{1}B$. Since in
case (ii) $\bar{B}$ is formally separable, we may in that case
choose the elements $x_{2},\ldots,x_{d}$ in such a way that
$\bar{B}$ becomes a generically étale $k[[x_{2},\ldots,x_{d}]]$-algebra.
It then follows immediately that $B$ is
also a generically étale $A$-algebra, which is of course also true
in case (i).

Let $b\in B$ be an element which generates the field extension $R(B)$ over
$k((x))$. Let $\epsilon \in k[[x]]$ be the discriminant of the
monic irreducible equation of $b$ over $k[[x]]$. In case (ii),
we may choose $b$ so that $x_{1}$ does not divide
$\epsilon$. Write moreover\pageoriginale $\Spec B-U=V(x_{1})\cup Y$,
where $Y$ contains no irreducible component of $V(x_{1})$. Let
$C$ be the ring $\Gamma(Y,O_{Y})$, where $Y$ is given its induced
reduced structure. Changing the $x_{2},\ldots,x_{d}$ if necessary, we
may suppose that $V(x_{d})$ is not contained in the image of $Y$
in $\Spec A$, and that, moreover, in case (ii), $V(x_{d})$ is not
contained in $V(\epsilon)$.

\hypertarget{sga4.exp.XIX.lemma.6.8}{}\begin{enonce*}{Lemma 6.8}\etiquette[6.8]{XIX.lem6.8}
Let $A_{0}=k[[x_{1},\ldots,x_{d-1}]]$ $\{x_{d}\}$ be the strict localization
of $k[[x_{1},\ldots,x_{d-1}]]$ $[x_{d}]$ at the point
$x_{1}=x_{2}=\ldots=x_{d}=0$. Under the conditions above, there exist
an $A^{0}$-algebra $B^{0}$, an affine open subset $U^{0}\subset \Spec
B^{0}$, and an isomorphism $A\otimes_{A^{0}}B^{0}\simeq B$ such that
the open subset of $\Spec B$ induced by $U^{0}$ is $U$.
\end{enonce*}

The proof of this lemma is analogous to that of \hyperref[XIX.lem1.5]{1.5}:
it suffices to find $A^{0}$-algebras $B^{0}$, $C^{0}$ and
isomorphisms $A\otimes_{A^{0}}B^{0}\simeq B$,
$A\otimes_{A^{0}}C^{0}\simeq C$, because then there will exist
(\cite{sga4.exp.XIX.bib.3}~\hyperref[XIX.subsec1.4]{1.4}) a unique (surjective) morphism
$B^{0}\rightarrow C^{0}$ inducing the morphism $B\rightarrow C$. We
therefore take $U^{0}=\Spec B^{0}-\{(\Spec C^{0})\cup V(x_{1})\}$, and
immediately verify that $U$ is induced by $U^{0}$. The existence of
$C^{0}$, and of $B^{0}$ in case (ii), is proved as in
\hyperref[XIX.lem1.5]{1.5}. The existence of $B^{0}$ in case (i) is a
consequence of (\cite{sga4.exp.XIX.bib.3} 
\hyperref[XIX.thm5.1]{5.1}). (It may also be proved directly by
applying Abhyankar's lemma to the ramification along $\{x_{1}=0\}$,
and descent.)

We can now complete the proof of
\hyperref[XIX.lem6.7]{6.7}. Applying 
\hyperref[XIX.cor4.3]{4.3} once more, we find that
$$
H^{q}(U,\ZZ/\ell)\simeq H^{q}(U^{0},\ZZ/\ell){\quoi},
$$
from which\pageoriginale it follows that it suffices to prove that the latter
group is zero for $q>d$. Now $\Spec A^{0}$ is a limit of schemes
affine and étale over $k[[x_{1},\ldots,x_{d-1}]]$ $[x_{d}]$,
so $\Spec B^{0}$ is a limit of schemes of finite type over
$\Spec k[[x_{1},\ldots,x_{d-1}]]$, each fiber of which has dimension
$\leq 1$. Since $U^{0}\subset \Spec B^{0}[1/x_{1}]$, we have
$$
U^{0}=\varprojlim U^{0}_{\alpha}
$$
where $U^{0}_{\alpha}$ is affine of finite type over the scheme
$V^{0}=\Spec k[[x_{1},\ldots, x_{d-1}]]$ $[1/x_{1}]$ and where each
fiber of $U^{0}_{\alpha}/V^{0}$ has dimension $\leq 1$.

Let $f^{0}:U^{0}\rightarrow V^{0}$ be the structural morphism. Since
the fibers of $f^{0}_{\alpha}$ all have dimension $\leq 1$, for
$v\in V^{0}$ we have
$$
\delta(\ZZ/\ell, f^{0}_{\alpha},v)\leq \dim O_{V^{0},v}+1\leq
d-1
$$
(because $\dim O_{V^{0},v}\leq d-2$). The induction hypothesis on
% SGA-SRC-001224
$d$ is therefore applicable to the morphism $f^{0}_{\alpha}$, and, by
applying \hyperref[XIX.thm6.1]{6.1}, we find
$$
R^{q}f^{0}_{\alpha\ast}(\ZZ/\ell)=0\quad\text{\rm at the
point}\quad v\quad\text{\rm if}\quad q>\dim O_{V^{0},v}+1{\quoi},
$$
\ie if $q>d-\dim\bar{v}$. Thus $R^{q}f^{0}_{\alpha
\ast}(\ZZ/\ell)\neq 0$ at the point $v$ implies that $\dim
\bar{v}\geq d-q$. Therefore, for the inclusion $i:V^{0}\rightarrow \Spec
k[[x_{1},\ldots, x_{d-1}]]$ (\cf \eqhyperref[XIX.eq6.0]{6.0}), we have
$$
\delta (R^{q}f^{0}_{\alpha\ast}(\ZZ/\ell))\leq d-q{\quoi},
$$
and the induction hypothesis, in the form of \hyperref[XIX.thm6.1]{6.1} bis, applied to the
morphism\pageoriginale $i$, implies that
$$
H^{p}(V^{0},R^{q}f^{0}_{\alpha\ast}(\ZZ/\ell))=0\quad\text{\rm
if}\quad p+q>d{\quoi}.
$$
The Leray spectral sequence
$$
E^{pq}_{2}=H^{p}(V^{0},R^{q}f^{0}_{\alpha\ast}(\ZZ/\ell))\Longrightarrow
H^{p+q}(U^{0},\ZZ/\ell)
$$
therefore implies that $H^{n}(U^{0}_{\alpha},\ZZ/\ell)=0$ if $n>d$,
and since $U^{0}$ is the limit of the $U^{0}_{\alpha}$, we likewise find
$$
H^{n}(U^{0},\ZZ/\ell)=0\quad\text{\rm if}\quad n>d{\quoi}.
$$

This completes the proof of \hyperref[XIX.lem6.7]{6.7}, and hence of Theorem \hyperref[XIX.thm6.1]{6.1} in
characteristic 0.

\section{Affine morphisms --- end of the proof}\etiquette[7]{XIX.sec7}

Recall that to complete the proof of \hyperref[XIX.thm6.1]{6.1} in characteristic
$p>0$ (still under the resolution hypothesis), we had reduced by
\hyperref[XIX.lem6.6]{6.6} to proving the following assertion for $d\geq 2$:

\hypertarget{sga4.exp.XIX.lemma.7.1}{}\begin{enonce*}{Lemma 7.1}\etiquette[7.1]{XIX.lem7.1}
Let $B$ be a complete normal local ring of equal characteristic, with
separably closed residue field, and of dimension $d$. Let $U\subset
\Spec B$ be an affine open subset such that there exists $x\in \rad B$ which is
invertible on $U$. Then 
$$
H^{q}(U,\ZZ/\ell)=0\quad\text{\rm if}\quad q>d{\quoi}.
$$

Moreover,\pageoriginale by the induction hypothesis, we may suppose 
that $\hyperref[XIX.thm6.1]{6.1}$ is true for the values $\delta<d$, and that $\hyperref[XIX.lem7.1]{7.1}$ has already been proved
in the particular case $\hyperref[XIX.lem6.7]{6.7}$ {\rm (ii)}.
\end{enonce*}

First note that the condition that $B$ be normal is not
important. Indeed, let $B$ be arbitrary, let $\pi:\bar{B}\rightarrow
B$ be the normalization of $B$, and let $\bar{U}\subset\Spec \bar{B}$ be the inverse
image of $U$. We have an exact sequence
$$
0\longrightarrow (\ZZ/\ell)_{U}\longrightarrow
\pi_{\ast}(\ZZ/\ell)_{\bar{U}}\longrightarrow C\longrightarrow
0{\quoi},
$$
where the dimension of the support of $C$ is $<d$, so $H^{q}(U,C)=0$ if
$q>d-1$. Since $H^{q}(U,\pi_{\ast}\bigdot)\simeq
H^{q}(\bar{U},\bigdot)$ (VIII \hyperref[VIII.cor5.6]{5.6}), we see that
\begin{equation*}
\begin{split}
& H^{q}(U,\ZZ/\ell)=0\quad \text{\rm if}\quad q>d\quad
\text{\rm is equivalent to}\\
& H^{q}(\bar{U},\ZZ/\ell)=0\quad\text{\rm if}\quad q>d{\quoi}.
\end{split}\tag{7.2}
\end{equation*}\etiquette[7.2]{XIX.eq7.2}

\setcounter{subsection}{3}
\setcounter{subsubsection}{-1}
\subsubsection{}\etiquette[7.3.0]{XIX.subsubsec7.3.0}
We shall briefly recall, without proof, the properties of the notion of
formal separability that we need. The reader may consult
(\EGA IV 18.11.10, 18.11.11): Let $A$ be a complete noetherian local
$k$-algebra whose residue field is a finite extension
of $k$, and let $A'$ denote the completion of the local ring
$A\otimes_{k}k'$ (where $k'=k^{1/p}$). Then the extension $A\rightarrow
A'$ is \emph{radicial} (possibly infinite), as is easily seen
--- this is the well-known example of Nagata. It is therefore
harmless, for the calculation of étale cohomology, to replace $A$ by
$A'$ (VIII 1).

Now the condition that an equidimensional algebra $A$ of
dimension\pageoriginale $d$ be formally \emph{separable},
that is, be a finite generically étale extension of
$k[[x_{1},\ldots, x_{d}]]$, can be expressed in terms of the module
$\widehat{\Omega}^{1}_{A/k}$ (\EGA IV 18.11.10), or of the Jacobian
criterion (\EGA $0_{IV}$ 22.6, IV 7.1)\footnote[1]{The co-editor of this
seminar has not succeeded in elucidating the meaning of this
reference to a situation which seems quite different to him, nor the
justification of the assertion that follows. A. G.}, and one sees that if $A$
is not formally separable, then the $k'$-algebra $A'$ is not
reduced. For the proof, one may for example reduce to the
case where $A$ is given by an equation $f=0$, $f\in k[[x_{1},\ldots,
x_{d+1}]]$, and apply the Jacobian criterion and the following lemma,
whose pleasant proof we leave to the reader:

\hypertarget{sga4.exp.XIX.lemma.7.3}{}\begin{enonce*}{Lemma 7.3}\etiquette[7.3]{XIX.lem7.3}
Let $f(x_{1},\ldots,x_{n})\in k[[x_{1},\ldots,x_{n}]]$ be a formal
power series. Suppose that for each $i=1,\ldots,n$ there exists an
invertible element $u_{i}\in k[[x_{1},\ldots,x_{n}]]$ such that $u_{i}f$
is a series in $x_{i}^{p}$ and in the
$x_{1},\ldots,x_{i-1},x_{i+1},\ldots,x_{n}$. Then there exists an
invertible element $v\in k[[x_{1},\ldots,x_{n}]]$ such that
$$
vf\in k[[x_{1}^{p},\ldots,x_{n}^{p}]]{\quoi}.
$$
\end{enonce*}

It follows easily that, by applying a finite number of times the
% SGA-SRC-001226
procedure of replacing $(A,k)$ by $(\bar{A}',k')$, where $\bar{A}'$ is the
normalization of $A'$, one is reduced to the situation where $A$ is indeed
formally separable. Similarly, one can prove that for a given $x\in
\rad A$, the ring $(A/(x))_{\text{\rm réd}}$ becomes a formally
separable algebra by the same procedure, applied to
the algebra $A$.

Applying VIII \hyperref[VIII.thm1.1]{1.1} and \hyperref[VIII.def7.2]{7.2}, we find

\hypertarget{sga4.exp.XIX.corollary.7.4}{}\begin{enonce*}{Corollary 7.4}\etiquette[7.4]{XIX.cor7.4}
It suffices to prove $\hyperref[XIX.lem7.1]{(7.1)}$ under the additional condition that
$B$ contain a field $k$ over which $B/\rad B$ is finite, and
such that the $k$-algebra $(B/(x))_{\text{\rm réd}}$ is formally
separable.\pageoriginale
\end{enonce*}
 
The proof of \hyperref[XIX.lem7.1]{7.1} in this case consists in a
reduction to the particular
case treated in \hyperref[XIX.lem6.7]{6.7} (ii):

With the notation of \hyperref[XIX.lem7.1]{7.1}, \hyperref[XIX.cor7.4]{7.4}, let
$$
D(x)=\sum_{i}r_{i}X_{i}
$$
be the divisor of $x$, so that each $X_{i}$ is a closed irreducible
subscheme of $\Spec B$ of codimension 1 which is formally separable over
$k$ (that is, such that
$B/\mathfrak{J}(X_{i})$ is a formally separable $k$-algebra,
which is moreover equivalent to the assertion that $(B/(x))_{\text{\rm
réd}}$ is formally separable). Choose $y\in\rad B$ such that
we have
$$
D(y)=\sum_{i}(r_{i}-1)X_{i}+\Gamma

% END INCLUDED SGA4_English_current_buildable_source_20260906/source/19/drafts/SGA4_XIX_LINES_1251_1650_English.texfrag

% BEGIN INCLUDED SGA4_English_current_buildable_source_20260906/source/19/drafts/SGA4_XIX_LINES_1651_1920_English.texfrag
$$
where $\Gamma$ contains none of the $X_{i}$, and consider the blow-up
$$
f:Z\longrightarrow \Spec B
$$
of the ideal $(x,y)$. The spectra of the rings $B[y/x]$ and $B[x/y]$
form an affine open covering of $Z$.

Let $B_{1}$ be the strict localization of $B[x/y]$ at the maximal ideal
generated by $(\rad B,x/y)$. For each closed point $Q$ of the closed
fiber of $\Spec B[y/x]\rightarrow \Spec B$, let $A_{Q}$ denote the
strict localization of $B[y/x]$ at $Q$, and let $U_{Q}$ be the affine
open subset of $\Spec A_{Q}$ that is the inverse image of $U$.

\hypertarget{sga4.exp.XIX.lemma.7.5}{}\begin{enonce*}{Lemma 7.5}\etiquette[7.5]{XIX.lem7.5}
To\pageoriginale prove $\hyperref[XIX.lem7.1]{7.1}$ for the ring $B$, it suffices to prove
that for $q>d(\geq 2)$ we have
\begin{equation*}
H^{q}(\Spec B_{1}[y/x],\ZZ/\ell)=0\tag{i}
\end{equation*}
and
\begin{equation*}
H^{q}(U_{Q},\ZZ/\ell)=0\quad\text{for each $Q$}{\quoi}.\tag{ii}
\end{equation*}
\end{enonce*}

\medskip
\noindent
\emph{Proof.} We have open immersions
$$
U\hookrightarrow \Spec B[1/x]\hookrightarrow \Spec B[y/x]{\quoi},
$$
and hence a diagram
$$
\xymatrix{
U \ar[rr]^{i} \ar[rd]_{k}&         & \Spec B[y/x] \ar[ld]^j\\
                        &   Z     &      . 
}
$$
If $Q$ is a closed point of the closed fiber of $\Spec
B[y/x]\rightarrow \Spec B$, we have
$(R^{q}i_{\ast}\ZZ/\ell)_{Q}=0$ if $q>d$, by hypothesis. For
each point $P$ of $\Spec B[y/x]$ which is not closed in the closed
fiber of $\Spec B[y/x]\rightarrow \Spec B$, we have $\delta(F,i,P)<d$,
where $F=(\ZZ/\ell)_{U}$ (notation of \eqhyperref[XIX.eq6.0]{6.0}). Thus, applying
the induction hypothesis and \hyperref[XIX.thm6.1]{6.1} to the morphism $i$, we find that for
the morphism $\Spec B[y/x]\rightarrow \Spec B$ we have
$$
\delta(R^{q}i_{\ast}\ZZ/\ell)\leq d-q{\quoi}.
$$
Consequently, the induction hypothesis in the form \hyperref[XIX.thm6.1]{6.1} bis,
applied to the latter morphism, implies that
$$
H^{p}(\Spec B[y/x],R^{q}i_{\ast}\ZZ/\ell)=0
$$
if\pageoriginale $q>0$ and $p>d-q$. With the Leray spectral sequence
$$
E^{pq}_{2}=H^{p}(\Spec
B[y/x],R^{q}i_{\ast}\ZZ/\ell)\Longrightarrow
H^{p+q}(U,\ZZ / \ell){\quoi}.
$$
we are reduced, in proving \hyperref[XIX.lem7.1]{7.1}, to proving that
$$
H^{p}(\Spec B[y/x],i_{\ast}\ZZ/\ell)=0\quad\text{\rm if}\quad
p>d{\quoi}.
$$

We have an exact sequence
\begin{equation*}
0\longrightarrow \ZZ/\ell\longrightarrow
i_{\ast}(\ZZ/\ell)_{U}\longrightarrow C\longrightarrow 0\tag{7.6}
\end{equation*}\etiquette[7.6]{XIX.eq7.6}
where $\delta(C)<d$. It therefore suffices (again by induction) to prove
that
\begin{equation*}
H^{p}(\Spec B[y/x],\ZZ/\ell)=0\quad\text{\rm if}\quad
p>d{\quoi}.\tag{7.7}
\end{equation*}\etiquette[7.7]{XIX.eq7.7}

Now the open immersion $\Spec B[y/x]\xrightarrow{j}Z$ is obtained by
removing from $Z$ the closed subset $C$ ``at infinity,'' and in the
other open subset $\Spec B[x/y]$ of $Z$ it is identified with the open
subset $C=V(x/y)$. The $R^{q}j_{\ast}\ZZ/\ell$ are concentrated on
$C$, and the morphism $C\rightarrow \Spec B$ is a closed immersion.
Consequently, we have $H^{p}(Z,R^{q}j_{\ast}\ZZ/\ell)=0$
if $p>0$ and $q>0$. Since the closed fiber of $Z\rightarrow \Spec B$ has
dimension 1 and this morphism is proper, \hyperrefext[XII][XII.cor5.3]{5.3} bis implies
that $H^{p}(Z,j_{\ast}\ZZ/\ell)=0$ if $p>2$. The spectral sequence
$$
E^{pq}_{2}=H^{p}(Z,R^{q}j_{\ast}\ZZ/\ell)\Longrightarrow
H^{p+q}(\Spec B[y/x],\ZZ/\ell)
$$
therefore gives isomorphisms
$$
H^{q}(\Spec B[y/x],\ZZ/\ell)\simeq
H^{0}(Z,R^{q}j_{\ast}\ZZ/\ell)\simeq
(R^{q}j_{\ast}\ZZ/\ell)_{P}
$$
if\pageoriginale $q>2$, where $P$ is the point ``at infinity'' of
$Z$. These stalks are none other (VIII 5) than
% SGA-SRC-001231
$$
(R^{q}j_{\ast}\ZZ/\ell)_{P}\simeq H^{q}(\Spec
B_{1}[y/x],\ZZ/\ell){\quoi}.
$$

By \eqhyperref[XIX.eq7.7]{7.7}, it therefore suffices to prove that these latter groups
vanish if $q>d$ $(\geq 2)$. We have an exact sequence analogous to \eqhyperref[XIX.eq7.6]{7.6}
which, together with the induction hypothesis, shows that it suffices that
$$
H^{q}(\Spec B_{1}[y/x],\ZZ/\ell)=0\quad\mbox{\rm if}\quad
q>d{\quoi},
$$
whence the lemma.

\medskip
\noindent
\emph{Proof of} \hyperref[XIX.lem7.5]{7.5} (i).

Consider the closed subset $C=V(x/y)$ of $\Spec B[x/y]$, with its
induced reduced structure. It is immediate that the closed immersion
of $C$ into $\Spec B$ induced by the morphism $Z\rightarrow
\Spec B$ identifies $C$ with the closed subscheme $X=\bigcup X_{i}$ of
$\Spec B$. Since
\begin{align*}
D(x) &= \sum r_{i}X_{i}\\
D(y) &= \sum (r_{i}-1)X_{i}+\Gamma{\quoi},
\end{align*}
the element $x/y$ generates the ideal $\mathfrak{J}(X_{i})$ locally at
the generic point of $X_{i}$, where the local ring is a discrete
valuation ring because $B$ is normal. Consequently, the
morphism $\Spec B[x/y]\rightarrow \Spec B$ is an isomorphism in a
neighborhood of such a point, and the element $x/y$ vanishes to order 1
on each of the irreducible\pageoriginale components $C_{i}$
(corresponding to $X_{i}$) of $C$.

Since the closed subset $C$ lifts to $\Spec B_{1}$, this remains
true if the ring $B[x/y]$ is replaced by $B_{1}$, and hence
$B_{1}/(x/y)$ is a \emph{complete} ring which is a formally separable
$k$-algebra.

To prove \hyperref[XIX.lem7.5]{7.5} (i), we may first apply \hyperref[XIX.cor4.3]{4.3},
which allows us to
replace $B_{1}$ by its completion $\hat{B}_{1}$, without affecting the
ring $B_{1}/(x/y)$. By \eqhyperref[XIX.eq7.2]{7.2}, we may moreover replace
$\hat{B}_{1}$ by its normalization $\bar{B}_{1}$. Since $B_{1}$ (and hence
$\hat{B}_{1}$) is already regular at the generic points of the $C_{i}$, the
morphism $\Spec \bar{B}_{1}\rightarrow \Spec \hat{B}_{1}$ is an
isomorphism at these points, and so $\bar{B}_{1}/(x/y)$ is still a
formally separable $k$-algebra. We are thus reduced to the special case
\hyperref[XIX.lem6.7]{6.7} (ii), which proves the result.

\medskip
\noindent
\emph{Proof of} \hyperref[XIX.lem7.5]{7.5} (ii).

Choose the point $Q$ of the closed fiber of $\Spec
B[y/x]\rightarrow \Spec B$. The strict localization $A_{Q}$ may be written as
a limit
$$
A_{Q}=\varinjlim A
$$
where $A$ ranges over a filtered category of rings that are étale and
of finite type over $B[y/x]$. Let $U_{A}$ be the affine open subset of
$\Spec A$ that is the inverse image of $U$. Since $x$ is invertible on
$U_{A}$, the morphism $U_{A}\rightarrow \Spec B$ is étale, and $U_{A}$
is normal. Consequently, if $C$ denotes the normalization of $B$
in the field $R(A)$ of rational functions on $A$, we have an open
immersion
$$
U_{A}\longrightarrow \Spec C{\quoi}.
$$
It\pageoriginale suffices (VII \hyperref[VII.thm5.7]{5.7}) to prove that for each $A$,
the image of $H^{q}(U_{A},\ZZ/\ell)$ in
$H^{q}(U_{Q},\ZZ/\ell)$ is zero if $q>d$. Since $x$ is
invertible on $U_{Q}$, the morphism $U_{Q}\rightarrow \Spec A$
factors through $\Spec A[1/f]$, if $f\in B$ is any element
such that $f$ is divisible by $x$ in $B[y/x]$ and
$f/x$ does not vanish at the point $Q$. It thus suffices to prove that
for a suitable such $f$, we have
$$
H^{q}(U_{A,f},\ZZ/\ell)=0\quad\text{\rm if}\quad q>d{\quoi},
$$
where $U_{A,f}=U_{A}-V(f)$ is the affine open subset of $\Spec A[1/f]$, which is
also an affine open subset of $\Spec C$.

Let $c\in k$ be an element such that $y/x+c$ does not vanish at $Q$, and
let $J$ be the ideal in $B$ of the closed subset $\Gamma\cap X$. Then if
\begin{equation*}
f\equiv y+cx(\text{\rm mod }J^{N})\quad N>0{\quoi},\tag{7.8}
\end{equation*}\etiquette[7.8]{XIX.eq7.8}
the element $f$ satisfies the condition above: Indeed, it is
immediate that in $\Spec B[y/x]$ we have $V(J)=V(x)$. Thus \eqhyperref[XIX.eq7.8]{7.8}
does imply that $f$ is divisible by $x$, whence
$$
f/x\equiv y/x+c(\text{\rm mod }J){\quoi}.
$$
Since $J$ vanishes at $Q$, and $y/x+c$ is nonzero at $Q$, this
proves our assertion.

Now $\Gamma\cap X$ is a closed subset of $\Spec B$ of codimension
2. Consequently, we can choose an $f$ satisfying \eqhyperref[XIX.eq7.8]{7.8},
call it $z$, such that $V(z)\cap V(y)$ is likewise of
codimension\pageoriginale 2. Also choose elements
$x_{3},\ldots,x_{d}$ such that $y$, $z$, $x_{3},\ldots,x_{d}$ form a
system of parameters of the ring $B$. Then $B$, and hence $C$ as well, is
a finite $k[[y,z,x_{3},\ldots,x_{d}]]$-algebra.

Let $C^{0}$ be the normalization of $k[[y,z,x_{3},\ldots,x_{d}]]$ in the
separable closure of $k$ $((y,z,x))$ in the field of rational
functions $R(C)$, so that $C$ is a radicial extension of
$C^{0}$, and $C^{0}$ is a generically étale extension of $k$
$((y,z,x))$ (where $x=x_{3},\ldots,x_{d}$).

Choose a linear combination
$$
f=az+by\quad (a,b\in k){\quoi},
$$
such that the morphism
$$
\Spec C^{0}\longrightarrow \Spec k[[z,y,x]]
$$
is étale over the generic point of $V(f)$, and such that $f/x$ does not
vanish at $Q$ in $\Spec B[y/x]$. It suffices to prove that
$$
H^{q}(U_{A,f},\ZZ/\ell)=0\quad\text{\rm if}\quad q>d{\quoi},
$$
and hence (VIII \hyperref[VIII.thm1.1]{1.1}) to prove that
$$
H^{q}(U^{0}_{A,f},\ZZ/\ell)=0\quad\text{\rm if}\quad
q>d{\quoi},
$$
where $U^{0}_{A,f}$ is the open subset of $\Spec C^{0}$ corresponding to
the open subset $U_{A,f}$ of $\Spec C$, which is affine, as follows by
descent (\EGA IV 2.7.1 (xiii)). But by construction,
the element $f$ of $C^{0}$ is invertible on $U^{0}_{A,f}$, and
$C^{0}/(f)$ is a formally separable $k$-algebra. We are therefore
in the situation of \hyperref[XIX.lem6.7]{6.7} (ii), which completes
the proof of \hyperref[XIX.lem7.5]{7.5} (ii), and hence that of \hyperref[XIX.thm6.1]{6.1}.

\begin{thebibliography}{99}
\bibitem{sga4.exp.XIX.bib.1} \textsc{Abhyankar}, S.\pageoriginale Local uniformization on
algebraic surfaces over ground fields of characteristic $p\neq
0$. Annals of Math. vol. 63 (1963), and Resolution of singularities of
embedded algebraic surfaces, Academic Press (1966).

\bibitem{sga4.exp.XIX.bib.2} \textsc{Abhyankar}, S. Resolution of singularities of
arithmetical surfaces. Arithmetical Algebraic Geometry, New York
(1966).

\bibitem{sga4.exp.XIX.bib.3} \textsc{Artin}, M. Etale coverings of schemes over hensel
rings, Amer, Journal Math. vol 88 ${\rm n^{\circ}}$ 4 p 915-934
(1966).

\bibitem{sga4.exp.XIX.bib.4} \textsc{Hironaka}, H. Resolution of singularities of an
algebraic variety over a field of characteristic zero I, II. Annals of
Math. vol. 79, ${\rm n^{\circ}}$ 1, 2 (1964).

\bibitem{sga4.exp.XIX.bib.5} \textsc{Weierstrass}, K. Einige auf die analytischen
Functionen mehrerer Veränderlichen sich beziehende
Sätze. Mathematische Werke von K. Weierstrass, Bd 2, Berlin (1895).

\bibitem{sga4.exp.XIX.bib.6} \textsc{Bourbaki}, N. Commutative Algebra, Chap. 3.
\end{thebibliography}


%XII 4.2 no ref. 1072

\ifx\danslelivre\undefined
\let\SGAXIXenddocument\enddocument
\else
\let\SGAXIXenddocument\relax
\fi
\SGAXIXenddocument

% END INCLUDED SGA4_English_current_buildable_source_20260906/source/19/drafts/SGA4_XIX_LINES_1651_1920_English.texfrag


% END INCLUDED SGA4_English_current_buildable_source_20260906/source/19/19_English.tex


\backmatter

\end{document}
